\IfFileExists{stacks-project-book.cls}{%
\documentclass{stacks-project-book}
}{%
\documentclass{amsbook}
}

\usepackage{verbatim}
\newenvironment{reference}{\comment}{\endcomment}
\newenvironment{slogan}{\comment}{\endcomment}
\newenvironment{history}{\comment}{\endcomment}

\usepackage[all]{xy}

\xyoption{2cell}
\UseAllTwocells

\usepackage{multicol}

\usepackage{lmodern}
\usepackage[T1]{fontenc}
\usepackage[french]{babel}
\AtBeginDocument{\shorthandoff{?}}


\usepackage{hyperref}


\theoremstyle{plain}
\newtheorem{theorem}[subsection]{Théorème}
\newtheorem{proposition}[subsection]{Proposition}
\newtheorem{lemma}[subsection]{Lemme}

\theoremstyle{definition}
\newtheorem{definition}[subsection]{Définition}
\newtheorem{example}[subsection]{Exemple}
\newtheorem{exercise}[subsection]{Exercice}
\newtheorem{situation}[subsection]{Situation}

\theoremstyle{remark}
\newtheorem{remark}[subsection]{Remarque}
\newtheorem{remarks}[subsection]{Remarques}

\renewcommand{\proofname}{Démonstration}

\numberwithin{equation}{subsection}

\def\lim{\mathop{\mathrm{lim}}\nolimits}
\def\colim{\mathop{\mathrm{colim}}\nolimits}
\def\Spec{\mathop{\mathrm{Spec}}}
\def\Hom{\mathop{\mathrm{Hom}}\nolimits}
\def\Ext{\mathop{\mathrm{Ext}}\nolimits}
\def\SheafHom{\mathop{\mathcal{H}\!\mathit{om}}\nolimits}
\def\SheafExt{\mathop{\mathcal{E}\!\mathit{xt}}\nolimits}
\def\Sch{\mathit{Sch}}
\def\Mor{\mathop{\mathrm{Mor}}\nolimits}
\def\Ob{\mathop{\mathrm{Ob}}\nolimits}
\def\Sh{\mathop{\mathit{Sh}}\nolimits}
\def\NL{\mathop{N\!L}\nolimits}
\def\CH{\mathop{\mathrm{CH}}\nolimits}
\def\proetale{{pro\text{-}\acute{e}tale}}
\def\etale{{\acute{e}tale}}
\def\QCoh{\mathit{QCoh}}
\def\Ker{\mathop{\mathrm{Ker}}}
\def\Im{\mathop{\mathrm{Im}}}
\def\Coker{\mathop{\mathrm{Coker}}}
\def\Coim{\mathop{\mathrm{Coim}}}

\DeclareMathSymbol{\boxtimes}{\mathbin}{AMSa}{"02}

\def\QCohstack{\mathcal{QC}\!\mathit{oh}}
\def\Cohstack{\mathcal{C}\!\mathit{oh}}
\def\Spacesstack{\mathcal{S}\!\mathit{paces}}
\def\Quotfunctor{\mathrm{Quot}}
\def\Hilbfunctor{\mathrm{Hilb}}
\def\Curvesstack{\mathcal{C}\!\mathit{urves}}
\def\Polarizedstack{\mathcal{P}\!\mathit{olarized}}
\def\Complexesstack{\mathcal{C}\!\mathit{omplexes}}
\def\Pic{\mathop{\mathrm{Pic}}\nolimits}
\def\Picardstack{\mathcal{P}\!\mathit{ic}}
\def\Picardfunctor{\mathrm{Pic}}
\def\Deformationcategory{\mathcal{D}\!\mathit{ef}}
\begin{document}
\title{Le projet Stacks --- édition française, partie 7}
\maketitle
\tableofcontents

% OK, start here
%

\chapter{Cohomologie de de Rham}



\phantomsection
\label{derham-section-phantom}


\section{Introduction}
\label{derham-section-introduction}

\noindent
Dans ce chapitre, nous commençons par étudier le complexe de de Rham
d'un morphisme de schémas et nous terminons en démontrant que la cohomologie
de de Rham définit une théorie cohomologique de Weil lorsque le corps de base est de caractéristique nulle.




\section{Le complexe de de Rham}
\label{derham-section-de-rham-complex}

\noindent
Soit $p : X \to S$ un morphisme de schémas. Il existe un complexe
$$
\Omega^\bullet_{X/S} =
\mathcal{O}_{X/S} \to \Omega^1_{X/S} \to \Omega^2_{X/S} \to \ldots
$$
de $p^{-1}\mathcal{O}_S$-modules, où
$\Omega^i_{X/S} = \wedge^i(\Omega_{X/S})$
est placé en degré $i$ et où la différentielle est déterminée par la règle
$\text{d}(g_0 \text{d}g_1 \wedge \ldots \wedge \text{d}g_p) =
\text{d}g_0 \wedge \text{d}g_1 \wedge \ldots \wedge \text{d}g_p$
sur les sections locales.
Voir Modules, section \ref{modules-section-de-rham-complex}.

\medskip\noindent
Étant donné un diagramme commutatif
$$
\xymatrix{
X' \ar[r]_f \ar[d] & X \ar[d] \\
S' \ar[r] & S
}
$$
de schémas, on dispose de morphismes canoniques de complexes
$f^{-1}\Omega_{X/S}^\bullet \to \Omega^\bullet_{X'/S'}$ et
$\Omega_{X/S}^\bullet \to f_*\Omega^\bullet_{X'/S'}$.
Voir Modules, section \ref{modules-section-de-rham-complex}.
Par linéarisation, on obtient, pour tout $p$, un morphisme linéaire
$f^*\Omega^p_{X/S} \to \Omega^p_{X'/S'}$.

\medskip\noindent
En particulier, si $f : Y \to X$ est un morphisme de schémas au-dessus
d'un schéma de base $S$, on dispose d'un morphisme de complexes
$$
\Omega^\bullet_{X/S} \longrightarrow f_*\Omega^\bullet_{Y/S}
$$
Par linéarisation, on voit que, pour tout $p \geq 0$, on obtient un morphisme canonique
$$
\Omega^p_{X/S} \otimes_{\mathcal{O}_X} f_*\mathcal{O}_Y
\longrightarrow
f_*\Omega^p_{Y/S}
$$

\begin{lemma}
\label{derham-lemma-base-change-de-rham}
Soit
$$
\xymatrix{
X' \ar[r]_f \ar[d] & X \ar[d] \\
S' \ar[r] & S
}
$$
un diagramme cartésien de schémas. Les morphismes considérés
ci-dessus induisent alors des isomorphismes
$f^*\Omega^p_{X/S} \to \Omega^p_{X'/S'}$.
\end{lemma}

\begin{proof}
Combiner Morphismes, lemme \ref{morphisms-lemma-base-change-differentials},
avec le fait que la formation des puissances extérieures commute au changement de base.
\end{proof}

\begin{lemma}
\label{derham-lemma-etale}
Considérons un diagramme commutatif de schémas
$$
\xymatrix{
X' \ar[r]_f \ar[d] & X \ar[d] \\
S' \ar[r] & S
}
$$
Si $X' \to X$ et $S' \to S$ sont étales, les morphismes considérés
ci-dessus induisent des isomorphismes
$f^*\Omega^p_{X/S} \to \Omega^p_{X'/S'}$.
\end{lemma}

\begin{proof}
On a $\Omega_{S'/S} = 0$ et $\Omega_{X'/X} = 0$ ; voir par exemple
Morphismes, lemme \ref{morphisms-lemma-etale-at-point}. Les suites exactes courtes de
Morphismes, lemmes
\ref{morphisms-lemma-triangle-differentials} et
\ref{morphisms-lemma-triangle-differentials-smooth},
montrent alors que $\Omega_{X'/S'} = \Omega_{X'/S} = f^*\Omega_{X/S}$.
On conclut en prenant les puissances extérieures.
\end{proof}






\section{Cohomologie de de Rham}
\label{derham-section-de-rham-cohomology}

\noindent
Soit $p : X \to S$ un morphisme de schémas. On appelle
{\it cohomologie de de Rham de $X$ sur $S$} les groupes de cohomologie
$$
H^i_{dR}(X/S) = H^i(R\Gamma(X, \Omega^\bullet_{X/S}))
$$
Puisque $\Omega^\bullet_{X/S}$ est un complexe de $p^{-1}\mathcal{O}_S$-modules,
ces groupes de cohomologie sont naturellement des modules sur $H^0(S, \mathcal{O}_S)$.

\medskip\noindent
Étant donné un diagramme commutatif
$$
\xymatrix{
X' \ar[r]_f \ar[d] & X \ar[d] \\
S' \ar[r] & S
}
$$
de schémas, les morphismes canoniques de la section \ref{derham-section-de-rham-complex}
fournissent des morphismes d'image inverse
$$
f^* :
R\Gamma(X, \Omega^\bullet_{X/S})
\longrightarrow
R\Gamma(X', \Omega^\bullet_{X'/S'})
$$
et
$$
f^* : H^i_{dR}(X/S) \longrightarrow H^i_{dR}(X'/S')
$$
Ces images inverses satisfont à une loi de composition évidente.
En particulier, si l'on travaille sur un schéma de base fixé $S$, la cohomologie
de de Rham est un foncteur contravariant sur la catégorie des schémas sur $S$.

\begin{lemma}
\label{derham-lemma-de-rham-affine}
Soit $X \to S$ un morphisme de schémas affines défini par le morphisme d'anneaux
$R \to A$. Alors $R\Gamma(X, \Omega^\bullet_{X/S}) = \Omega^\bullet_{A/R}$
dans $D(R)$ et $H^i_{dR}(X/S) = H^i(\Omega^\bullet_{A/R})$.
\end{lemma}

\begin{proof}
Cela résulte de Cohomologie des schémas, lemme
\ref{coherent-lemma-quasi-coherent-affine-cohomology-zero},
et du lemme d'acyclicité de Leray
(Catégories dérivées, lemme \ref{derived-lemma-leray-acyclicity}).
\end{proof}

\begin{lemma}
\label{derham-lemma-quasi-coherence-relative}
Soit $p : X \to S$ un morphisme de schémas. Si $p$ est quasi-compact
et quasi-séparé, alors $Rp_*\Omega^\bullet_{X/S}$ est un objet
de $D_\QCoh(\mathcal{O}_S)$.
\end{lemma}

\begin{proof}
Il existe une suite spectrale de première page
$E_1^{a, b} = R^bp_*\Omega^a_{X/S}$ convergeant vers
la cohomologie de $Rp_*\Omega^\bullet_{X/S}$
(voir Catégories dérivées, lemme \ref{derived-lemma-two-ss-complex-functor}).
Ainsi, d'après Homologie, lemme \ref{homology-lemma-first-quadrant-ss},
il suffit de montrer que $R^bp_*\Omega^a_{X/S}$ est quasi-cohérent.
Cela résulte de Cohomologie des schémas, lemme
\ref{coherent-lemma-quasi-coherence-higher-direct-images}.
\end{proof}

\begin{lemma}
\label{derham-lemma-coherence-relative}
Soit $p : X \to S$ un morphisme propre de schémas, où $S$ est localement
noethérien. Alors $Rp_*\Omega^\bullet_{X/S}$ est un objet
de $D_{\textit{Coh}}(\mathcal{O}_S)$.
\end{lemma}

\begin{proof}
Dans ce cas, d'après Morphismes, lemme \ref{morphisms-lemma-finite-type-differentials},
les modules $\Omega^i_{X/S}$ sont cohérents. On peut donc reprendre exactement
l'argument de la démonstration du lemme \ref{derham-lemma-quasi-coherence-relative}
en utilisant Cohomologie des schémas, proposition
\ref{coherent-proposition-proper-pushforward-coherent}.
\end{proof}

\begin{lemma}
\label{derham-lemma-finite-de-Rham}
Soit $A$ un anneau noethérien. Soit $X$ un schéma propre sur $S = \Spec(A)$.
Alors $H^i_{dR}(X/S)$ est un $A$-module de type fini pour tout $i$.
\end{lemma}

\begin{proof}
C'est un cas particulier du lemme \ref{derham-lemma-coherence-relative}.
\end{proof}

\begin{lemma}
\label{derham-lemma-proper-smooth-de-Rham}
Soit $f : X \to S$ un morphisme propre et lisse de schémas. Alors
$Rf_*\Omega^p_{X/S}$, pour $p \geq 0$, et $Rf_*\Omega^\bullet_{X/S}$ sont
des objets parfaits de $D(\mathcal{O}_S)$ dont la formation commute
à tout changement de base.
\end{lemma}

\begin{proof}
Puisque $f$ est lisse, les modules $\Omega^p_{X/S}$ sont des
$\mathcal{O}_X$-modules localement libres de rang fini ; voir Morphismes, lemme
\ref{morphisms-lemma-smooth-omega-finite-locally-free}. Leur
formation commute à tout changement de base d'après
le lemme \ref{derham-lemma-base-change-de-rham}. Par conséquent,
$Rf_*\Omega^p_{X/S}$ est un objet parfait de $D(\mathcal{O}_S)$
dont la formation commute à tout changement de base ; voir
Catégories dérivées des schémas, lemme
\ref{perfect-lemma-flat-proper-perfect-direct-image-general}.
Cela démontre la première assertion du lemme.

\medskip\noindent
Pour démontrer que $Rf_*\Omega^\bullet_{X/S}$ est parfait sur $S$, on peut travailler
localement sur $S$. On peut donc supposer $S$ quasi-compact. Cela signifie
que l'on peut supposer $\Omega^n_{X/S}$ nul pour $n$ assez grand.
Pour tout $p \geq 0$, nous affirmons que
$Rf_*\sigma_{\geq p}\Omega^\bullet_{X/S}$ est un
objet parfait de $D(\mathcal{O}_S)$ dont la formation commute
à tout changement de base. D'après ce qui précède,
c'est vrai pour $p \gg 0$. Supposons l'assertion vraie pour $p$
et considérons le triangle distingué
$$
\sigma_{\geq p}\Omega^\bullet_{X/S} \to
\sigma_{\geq p - 1}\Omega^\bullet_{X/S} \to
\Omega^{p - 1}_{X/S}[-(p - 1)] \to
(\sigma_{\geq p}\Omega^\bullet_{X/S})[1]
$$
dans $D(f^{-1}\mathcal{O}_S)$.
En appliquant le foncteur exact $Rf_*$, on obtient un triangle distingué
dans $D(\mathcal{O}_S)$.
Comme la propriété d'être parfait satisfait à la propriété des deux sur trois
(Cohomologie, lemme \ref{cohomology-lemma-two-out-of-three-perfect}),
on en déduit que $Rf_*\sigma_{\geq p - 1}\Omega^\bullet_{X/S}$ est un
objet parfait de $D(\mathcal{O}_S)$. Il en va de même de la
commutation à tout changement de base.
\end{proof}





\section{Cup-produit}
\label{derham-section-cup-product}

\noindent
Considérons les applications
$\Omega^p_{X/S} \times \Omega^q_{X/S} \to \Omega^{p + q}_{X/S}$
définies par $(\omega , \eta) \longmapsto \omega \wedge \eta$.
En utilisant la formule pour $\text{d}$ donnée à la section \ref{derham-section-de-rham-complex}
et la règle de Leibniz pour $\text{d} : \mathcal{O}_X \to \Omega_{X/S}$,
on voit que $\text{d}(\omega \wedge \eta) = \text{d}(\omega) \wedge \eta +
(-1)^{\deg(\omega)} \omega \wedge \text{d}(\eta)$. Cela signifie que
$\wedge$ définit un morphisme
\begin{equation}
\label{derham-equation-wedge}
\wedge :
\text{Tot}(
\Omega^\bullet_{X/S} \otimes_{p^{-1}\mathcal{O}_S} \Omega^\bullet_{X/S})
\longrightarrow
\Omega^\bullet_{X/S}
\end{equation}
de complexes de $p^{-1}\mathcal{O}_S$-modules.

\medskip\noindent
En combinant le cup-produit de
Cohomologie, section \ref{cohomology-section-cup-product},
avec (\ref{derham-equation-wedge}), on obtient un cup-produit
$H^0(S, \mathcal{O}_S)$-bilinéaire
$$
\cup : H^i_{dR}(X/S) \times H^j_{dR}(X/S) \longrightarrow H^{i + j}_{dR}(X/S)
$$
Par exemple, si $\omega \in \Gamma(X, \Omega^i_{X/S})$ et
$\eta \in \Gamma(X, \Omega^j_{X/S})$ sont fermées,
le cup-produit des classes de cohomologie de de Rham de
$\omega$ et de $\eta$ est la classe de cohomologie de de Rham de $\omega \wedge \eta$ ;
voir la discussion dans Cohomologie, section \ref{cohomology-section-cup-product}.

\medskip\noindent
Étant donné un diagramme commutatif
$$
\xymatrix{
X' \ar[r]_f \ar[d] & X \ar[d] \\
S' \ar[r] & S
}
$$
de schémas, les morphismes d'image inverse
$f^* : R\Gamma(X, \Omega^\bullet_{X/S}) \to R\Gamma(X', \Omega^\bullet_{X'/S'})$
et
$f^* : H^i_{dR}(X/S) \longrightarrow H^i_{dR}(X'/S')$
sont compatibles avec le cup-produit défini ci-dessus.

\begin{lemma}
\label{derham-lemma-cup-product-graded-commutative}
Soit $p : X \to S$ un morphisme de schémas.
Le cup-produit sur $H^*_{dR}(X/S)$ est associatif et
gradué commutatif.
\end{lemma}

\begin{proof}
Cela résulte de
Cohomologie, lemmes \ref{cohomology-lemma-cup-product-associative} et
\ref{cohomology-lemma-cup-product-commutative},
et du fait que $\wedge$ est associatif et gradué commutatif.
\end{proof}

\begin{remark}
\label{derham-remark-relative-cup-product}
Soit $p : X \to S$ un morphisme de schémas. On peut considérer
$\Omega^\bullet_{X/S}$ comme un faisceau de
$p^{-1}\mathcal{O}_S$-algèbres différentielles graduées ; voir
Faisceaux différentiels gradués, définition \ref{sdga-definition-dga}.
En particulier, la discussion de
Faisceaux différentiels gradués, section \ref{sdga-section-misc},
s'applique. Par exemple, pour tout diagramme commutatif
$$
\xymatrix{
X \ar[d]_p \ar[r]_f & Y \ar[d]^q \\
S \ar[r]^h & T
}
$$
de schémas, il existe un cup-produit relatif canonique
$$
\mu :
Rf_*\Omega^\bullet_{X/S}
\otimes_{q^{-1}\mathcal{O}_T}^\mathbf{L}
Rf_*\Omega^\bullet_{X/S}
\longrightarrow
Rf_*\Omega^\bullet_{X/S}
$$
dans $D(Y, q^{-1}\mathcal{O}_T)$, qui est associatif et
qui redonne en cohomologie le cup-produit considéré ci-dessus.
\end{remark}

\begin{remark}
\label{derham-remark-cup-product-as-a-map}
Soit $f : X \to S$ un morphisme de schémas. Soit $\xi \in H_{dR}^n(X/S)$.
D'après la discussion de
Faisceaux différentiels gradués, section \ref{sdga-section-misc},
il existe un morphisme canonique
$$
\xi' : \Omega^\bullet_{X/S} \to \Omega^\bullet_{X/S}[n]
$$
dans $D(f^{-1}\mathcal{O}_S)$, caractérisé de manière unique par
(1) et (2) dans la liste de propriétés suivante :
\begin{enumerate}
\item $\xi'$ se relève en un morphisme dans la catégorie dérivée des
$\Omega^\bullet_{X/S}$-modules différentiels gradués à droite, et
\item $\xi'(1) = \xi$ dans
$H^0(X, \Omega^\bullet_{X/S}[n]) = H^n_{dR}(X/S)$,
\item le morphisme $\xi'$ envoie $\eta \in H^m_{dR}(X/S)$
sur $\xi \cup \eta$ dans $H^{n + m}_{dR}(X/S)$,
\item la construction de $\xi'$ commute aux restrictions aux
ouverts : pour un ouvert $U \subset X$, la restriction $\xi'|_U$ est
le morphisme correspondant à l'image $\xi|_U \in H^n_{dR}(U/S)$,
\item pour tout diagramme comme dans la remarque \ref{derham-remark-relative-cup-product},
on obtient un diagramme commutatif
$$
\xymatrix{
Rf_*\Omega^\bullet_{X/S}
\otimes_{q^{-1}\mathcal{O}_T}^\mathbf{L}
Rf_*\Omega^\bullet_{X/S} \ar[d]_{\xi' \otimes \text{id}}
\ar[r]_-\mu &
Rf_*\Omega^\bullet_{X/S} \ar[d]^{\xi'} \\
Rf_*\Omega^\bullet_{X/S}[n]
\otimes_{q^{-1}\mathcal{O}_T}^\mathbf{L}
Rf_*\Omega^\bullet_{X/S}
\ar[r]^-\mu &
Rf_*\Omega^\bullet_{X/S}[n]
}
$$
dans $D(Y, q^{-1}\mathcal{O}_T)$.
\end{enumerate}
\end{remark}




\section{Cohomologie de Hodge}
\label{derham-section-hodge-cohomology}

\noindent
Soit $p : X \to S$ un morphisme de schémas. On appelle
{\it cohomologie de Hodge de $X$ sur $S$} les groupes de cohomologie
$$
H^n_{Hodge}(X/S) = \bigoplus\nolimits_{n = p + q} H^q(X, \Omega^p_{X/S})
$$
considérés comme un $H^0(X, \mathcal{O}_X)$-module gradué. Le produit extérieur
des formes, combiné au cup-produit de
Cohomologie, section \ref{cohomology-section-cup-product},
définit un cup-produit $H^0(X, \mathcal{O}_X)$-bilinéaire
$$
\cup :
H^i_{Hodge}(X/S) \times H^j_{Hodge}(X/S)
\longrightarrow
H^{i + j}_{Hodge}(X/S)
$$
Bien entendu, si $\xi \in H^q(X, \Omega^p_{X/S})$ et
$\xi' \in H^{q'}(X, \Omega^{p'}_{X/S})$, alors $\xi \cup \xi' \in
H^{q + q'}(X, \Omega^{p + p'}_{X/S})$.

\begin{lemma}
\label{derham-lemma-cup-product-hodge-graded-commutative}
Soit $p : X \to S$ un morphisme de schémas.
Le cup-produit sur $H^*_{Hodge}(X/S)$ est associatif et gradué commutatif.
\end{lemma}

\begin{proof}
La démonstration est identique à celle du
lemme \ref{derham-lemma-cup-product-graded-commutative}.
\end{proof}

\noindent
Étant donné un diagramme commutatif
$$
\xymatrix{
X' \ar[r]_f \ar[d] & X \ar[d] \\
S' \ar[r] & S
}
$$
de schémas, on dispose de morphismes d'image inverse
$f^* : H^i_{Hodge}(X/S) \longrightarrow H^i_{Hodge}(X'/S')$
compatibles avec les graduations et avec le cup-produit défini ci-dessus.







\section{Deux suites spectrales}
\label{derham-section-hodge-to-de-rham}

\noindent
Soit $p : X \to S$ un morphisme de schémas. Puisque la catégorie
des $p^{-1}\mathcal{O}_S$-modules sur $X$ a suffisamment d'injectifs,
il existe une résolution de Cartan--Eilenberg de $\Omega^\bullet_{X/S}$.
Voir Catégories dérivées, lemme \ref{derived-lemma-cartan-eilenberg}.
On peut donc appliquer Catégories dérivées, lemme
\ref{derived-lemma-two-ss-complex-functor}, pour obtenir deux suites spectrales
convergeant toutes deux vers la cohomologie de de Rham de $X$ sur $S$.

\medskip\noindent
La première est habituellement appelée {\it suite spectrale de Hodge vers de Rham}.
Sa première page est donnée par
$$
E_1^{p, q} = H^q(X, \Omega^p_{X/S})
$$
qui sont les groupes de cohomologie de Hodge de $X/S$ (d'où son nom). La
différentielle $d_1$ de cette page est donnée par les applications
$d_1^{p, q} : H^q(X, \Omega^p_{X/S}) \to H^q(X, \Omega^{p + 1}_{X/S})$
induites par la différentielle
$\text{d} : \Omega^p_{X/S} \to \Omega^{p + 1}_{X/S}$.
Voici une représentation
$$
\xymatrix{
H^2(X, \mathcal{O}_X) \ar[r] \ar@{-->}[rrd] \ar@{..>}[rrrdd] &
H^2(X, \Omega^1_{X/S}) \ar[r] \ar@{-->}[rrd] &
H^2(X, \Omega^2_{X/S}) \ar[r] &
H^2(X, \Omega^3_{X/S}) \\
H^1(X, \mathcal{O}_X) \ar[r] \ar@{-->}[rrd] &
H^1(X, \Omega^1_{X/S}) \ar[r] \ar@{-->}[rrd] &
H^1(X, \Omega^2_{X/S}) \ar[r] &
H^1(X, \Omega^3_{X/S}) \\
H^0(X, \mathcal{O}_X) \ar[r] &
H^0(X, \Omega^1_{X/S}) \ar[r] &
H^0(X, \Omega^2_{X/S}) \ar[r] &
H^0(X, \Omega^3_{X/S})
}
$$
où les flèches en tirets indiquent la source et le but
des différentielles de la page $E_2$, et la flèche en pointillés une différentielle
de la page $E_3$. En considérant le degré $0$, on en déduit que
$$
H^0_{dR}(X/S) =
\Ker(\text{d} : H^0(X, \mathcal{O}_X) \to H^0(X, \Omega^1_{X/S}))
$$
Bien entendu, cela résulte aussi immédiatement du fait que le
complexe de de Rham commence en degré $0$ par $\mathcal{O}_X \to \Omega^1_{X/S}$.

\medskip\noindent
La seconde suite spectrale est habituellement appelée
{\it suite spectrale conjuguée}. Sa deuxième page est donnée par
$$
E_2^{p, q} = H^p(X, H^q(\Omega^\bullet_{X/S})) = H^p(X, \mathcal{H}^q)
$$
où $\mathcal{H}^q = H^q(\Omega^\bullet_{X/S})$ est le faisceau de cohomologie de degré $q$
du complexe de de Rham de $X/S$. Les différentielles
de cette page sont données par $E_2^{p, q} \to E_2^{p + 2, q - 1}$.
Voici une représentation
$$
\xymatrix{
H^0(X, \mathcal{H}^2) \ar[rrd] \ar@{..>}[rrrdd] &
H^1(X, \mathcal{H}^2) \ar[rrd] &
H^2(X, \mathcal{H}^2) &
H^3(X, \mathcal{H}^2) \\
H^0(X, \mathcal{H}^1) \ar[rrd] &
H^1(X, \mathcal{H}^1) \ar[rrd] &
H^2(X, \mathcal{H}^1) &
H^3(X, \mathcal{H}^1) \\
H^0(X, \mathcal{H}^0) &
H^1(X, \mathcal{H}^0) &
H^2(X, \mathcal{H}^0) &
H^3(X, \mathcal{H}^0)
}
$$
En considérant le degré $0$, on en déduit que
$$
H^0_{dR}(X/S) = H^0(X, \mathcal{H}^0)
$$
ce qui est évident à la réflexion. En degré $1$, on obtient une suite exacte
$$
0 \to H^1(X, \mathcal{H}^0) \to H^1_{dR}(X/S) \to
H^0(X, \mathcal{H}^1) \to H^2(X, \mathcal{H}^0) \to H^2_{dR}(X/S)
$$
Il se trouve que, si $X \to S$ est lisse et si $S$ est en caractéristique $p$,
les faisceaux $\mathcal{H}^q$ peuvent être calculés (en termes de certains
faisceaux de différentielles), et la suite spectrale conjuguée est un outil
précieux (insérer ici une référence ultérieure).



\section{La filtration de Hodge}
\label{derham-section-hodge-filtration}

\noindent
Soit $X \to S$ un morphisme de schémas. La filtration de Hodge sur $H^n_{dR}(X/S)$
est la filtration induite par la suite spectrale de Hodge vers de Rham
(Homologie, définition
\ref{homology-definition-filtration-cohomology-filtered-complex}).
Pour éviter toute ambiguïté, nous la définissons explicitement comme suit.

\begin{definition}
\label{derham-definition-hodge-filtration}
Soit $X \to S$ un morphisme de schémas. La {\it filtration de Hodge}
sur $H^n_{dR}(X/S)$ est la filtration dont les termes sont
$$
F^pH^n_{dR}(X/S) = \Im\left(H^n(X, \sigma_{\geq p}\Omega^\bullet_{X/S})
\longrightarrow H^n_{dR}(X/S)\right)
$$
où $\sigma_{\geq p}\Omega^\bullet_{X/S}$ est défini dans
Homologie, section \ref{homology-section-truncations}.
\end{definition}

\noindent
Bien entendu, $\sigma_{\geq p}\Omega^\bullet_{X/S}$ est un sous-complexe du
complexe de de Rham relatif, et l'on obtient une filtration
$$
\Omega^\bullet_{X/S} = \sigma_{\geq 0}\Omega^\bullet_{X/S} \supset
\sigma_{\geq 1}\Omega^\bullet_{X/S} \supset
\sigma_{\geq 2}\Omega^\bullet_{X/S} \supset
\sigma_{\geq 3}\Omega^\bullet_{X/S} \supset \ldots
$$
du complexe de de Rham relatif, avec
$\text{gr}^p(\Omega^\bullet_{X/S}) = \Omega^p_{X/S}[-p]$.
La suite spectrale construite dans
Cohomologie, lemme \ref{cohomology-lemma-spectral-sequence-filtered-object},
pour $\Omega^\bullet_{X/S}$ considéré comme un complexe filtré de faisceaux
coïncide avec la suite spectrale de Hodge vers de Rham construite à la
section \ref{derham-section-hodge-to-de-rham}, d'après
Cohomologie, exemple \ref{cohomology-example-spectral-sequence-bis}. De plus, le
produit extérieur (\ref{derham-equation-wedge}) envoie
$\text{Tot}(\sigma_{\geq i}\Omega^\bullet_{X/S} \otimes_{p^{-1}\mathcal{O}_S}
\sigma_{\geq j}\Omega^\bullet_{X/S})$ dans
$\sigma_{\geq i + j}\Omega^\bullet_{X/S}$. On obtient donc des
diagrammes commutatifs
$$
\xymatrix{
H^n(X, \sigma_{\geq i}\Omega^\bullet_{X/S}))
\times 
H^m(X, \sigma_{\geq j}\Omega^\bullet_{X/S}))
\ar[r] \ar[d] &
H^{n + m}(X, \sigma_{\geq i + j}\Omega^\bullet_{X/S})) \ar[d] \\
H^n_{dR}(X/S) \times
H^m_{dR}(X/S)
\ar[r]^\cup &
H^{n + m}_{dR}(X/S)
}
$$
En particulier, on trouve que
$$
F^iH^n_{dR}(X/S) \cup F^jH^m_{dR}(X/S) \subset F^{i + j}H^{n + m}_{dR}(X/S)
$$






\section{Formule de K\"unneth}
\label{derham-section-kunneth}

\noindent
Une propriété importante de la cohomologie de de Rham est l'existence
d'une formule de K\"unneth.

\medskip\noindent
Soient $a : X \to S$ et $b : Y \to S$ des morphismes de schémas de même
but. Soient $p : X \times_S Y \to X$ et $q : X \times_S Y \to Y$ les
projections, et posons $f = a \circ p = b \circ q$. Voici le diagramme correspondant
$$
\xymatrix{
& X \times_S Y \ar[ld]^p \ar[rd]_q \ar[dd]^f \\
X \ar[rd]_a & & Y \ar[ld]^b \\
& S
}
$$
Dans cette section, pour un $\mathcal{O}_X$-module $\mathcal{F}$
et un $\mathcal{O}_Y$-module $\mathcal{G}$, posons
$$
\mathcal{F} \boxtimes \mathcal{G} =
p^*\mathcal{F} \otimes_{\mathcal{O}_{X \times_S Y}} q^*\mathcal{G}
$$
Le bifoncteur
$(\mathcal{F}, \mathcal{G}) \mapsto \mathcal{F} \boxtimes \mathcal{G}$
sur les modules quasi-cohérents s'étend en un bifoncteur sur les modules quasi-cohérents
et les opérateurs différentiels d'ordre fini sur $S$ ; voir
Morphismes, remarque \ref{morphisms-remark-base-change-differential-operators}.
Les différentielles des complexes de de Rham $\Omega^\bullet_{X/S}$ et
$\Omega^\bullet_{Y/S}$ sont des opérateurs différentiels d'ordre $1$
sur $S$, d'après Modules, lemme
\ref{modules-lemma-differentials-relative-de-rham-complex-order-1}.
Il est donc possible de considérer le complexe
$$
\text{Tot}(\Omega^\bullet_{X/S} \boxtimes \Omega^\bullet_{Y/S})
$$
Voir la discussion dans Catégories dérivées des schémas, section
\ref{perfect-section-kunneth-complexes}.

\begin{lemma}
\label{derham-lemma-de-rham-complex-product}
Dans la situation ci-dessus, il existe un isomorphisme canonique
$$
\text{Tot}(\Omega^\bullet_{X/S} \boxtimes \Omega^\bullet_{Y/S})
\longrightarrow
\Omega^\bullet_{X \times_S Y/S}
$$
de complexes de $f^{-1}\mathcal{O}_S$-modules.
\end{lemma}

\begin{proof}
On sait que
$
\Omega_{X \times_S Y/S} = p^*\Omega_{X/S} \oplus q^*\Omega_{Y/S}
$
d'après Morphismes, lemme \ref{morphisms-lemma-differential-product}.
En prenant les puissances extérieures, on obtient
$$
\Omega^n_{X \times_S Y/S} =
\bigoplus\nolimits_{i + j = n}
p^*\Omega^i_{X/S} \otimes_{\mathcal{O}_{X \times_S Y}} q^*\Omega^j_{Y/S} =
\bigoplus\nolimits_{i + j = n}
\Omega^i_{X/S} \boxtimes \Omega^j_{Y/S}
$$
par les propriétés élémentaires des puissances extérieures. Ces identifications déterminent
des isomorphismes entre les termes des complexes à gauche et à droite
de la flèche du lemme. Nous omettons de vérifier que ces morphismes
sont compatibles avec les différentielles.
\end{proof}

\noindent
Posons $A = \Gamma(S, \mathcal{O}_S)$. En combinant le résultat du
lemme \ref{derham-lemma-de-rham-complex-product} avec le morphisme de
Catégories dérivées des schémas, équation
(\ref{perfect-equation-de-rham-kunneth}),
on obtient un cup-produit
$$
R\Gamma(X, \Omega^\bullet_{X/S})
\otimes_A^\mathbf{L}
R\Gamma(Y, \Omega^\bullet_{Y/S})
\longrightarrow
R\Gamma(X \times_S Y, \Omega^\bullet_{X \times_S Y/S})
$$
Au niveau de la cohomologie, en utilisant la discussion de
Compléments d'algèbre, section \ref{more-algebra-section-products-tor},
on obtient une application canonique
$$
H^i_{dR}(X/S) \otimes_A H^j_{dR}(Y/S)
\longrightarrow
H^{i + j}_{dR}(X \times_S Y/S),\quad
(\xi, \zeta) \longmapsto p^*\xi \cup q^*\zeta
$$
Notons que la construction ci-dessus consiste bien à
prendre d'abord les images inverses, puis leur cup-produit.

\begin{lemma}
\label{derham-lemma-kunneth-de-rham}
Supposons $X$ et $Y$ lisses, quasi-compacts et à diagonale affine sur
$S = \Spec(A)$. Alors le morphisme
$$
R\Gamma(X, \Omega^\bullet_{X/S})
\otimes_A^\mathbf{L}
R\Gamma(Y, \Omega^\bullet_{Y/S})
\longrightarrow
R\Gamma(X \times_S Y, \Omega^\bullet_{X \times_S Y/S})
$$
est un isomorphisme dans $D(A)$.
\end{lemma}

\begin{proof}
D'après Morphismes, lemme \ref{morphisms-lemma-smooth-omega-finite-locally-free},
les faisceaux $\Omega^n_{X/S}$ et $\Omega^m_{Y/S}$ sont respectivement des
$\mathcal{O}_X$-modules et des $\mathcal{O}_Y$-modules localement libres de rang fini. D'autre part, $X$ et $Y$
sont plats sur $S$ (Morphismes, lemme \ref{morphisms-lemma-smooth-flat}),
d'où la platitude de $\Omega^n_{X/S}$ et de $\Omega^m_{Y/S}$ sur $S$.
Observons aussi que $\Omega^\bullet_{X/S}$ est localement borné. Le résultat
découle donc du lemme \ref{derham-lemma-de-rham-complex-product} et de
Catégories dérivées des schémas, lemme \ref{perfect-lemma-kunneth-special}.
\end{proof}

\noindent
Il existe une version relative du cup-produit, à savoir un morphisme
$$
Ra_*\Omega^\bullet_{X/S}
\otimes_{\mathcal{O}_S}^\mathbf{L}
Rb_*\Omega^\bullet_{Y/S}
\longrightarrow
Rf_*\Omega^\bullet_{X \times_S Y/S}
$$
dans $D(\mathcal{O}_S)$. Sa construction combine
le lemme \ref{derham-lemma-de-rham-complex-product} avec le morphisme de
Catégories dérivées des schémas, équation
(\ref{perfect-equation-relative-de-rham-kunneth}).
La construction montre que ce morphisme est donné par le diagramme
$$
\xymatrix{
Ra_*\Omega^\bullet_{X/S}
\otimes_{\mathcal{O}_S}^\mathbf{L}
Rb_*\Omega^\bullet_{Y/S}
\ar[d]^{\text{units of adjunction}}  \\
Rf_*(p^{-1}\Omega^\bullet_{X/S})
\otimes_{\mathcal{O}_S}^\mathbf{L}
Rf_*(q^{-1}\Omega^\bullet_{Y/S}) \ar[r] \ar[d]^{\text{relative cup product}} &
Rf_*(\Omega^\bullet_{X \times_S Y/S})
\otimes_{\mathcal{O}_S}^\mathbf{L}
Rf_*(\Omega^\bullet_{X \times_S Y/S}) \ar[d]^{\text{relative cup product}} \\
Rf_*(p^{-1}\Omega^\bullet_{X/S}
\otimes_{f^{-1}\mathcal{O}_S}^\mathbf{L}
q^{-1}\Omega^\bullet_{Y/S})
\ar[d]^{\text{from derived to usual}} \ar[r] &
Rf_*(\Omega^\bullet_{X \times_S Y/S}
\otimes_{f^{-1}\mathcal{O}_S}^\mathbf{L}
\Omega^\bullet_{X \times_S Y/S})
\ar[d]^{\text{from derived to usual}} \\
Rf_*\text{Tot}(p^{-1}\Omega^\bullet_{X/S}
\otimes_{f^{-1}\mathcal{O}_S}
q^{-1}\Omega^\bullet_{Y/S}) \ar[r] \ar[d]^{\text{canonical map}} &
Rf_*\text{Tot}(\Omega^\bullet_{X \times_S Y/S}
\otimes_{f^{-1}\mathcal{O}_S}
\Omega^\bullet_{X \times_S Y/S})
\ar[d]^{\eta \otimes \omega \mapsto \eta \wedge \omega} \\
Rf_*\text{Tot}(\Omega^\bullet_{X/S} \boxtimes \Omega^\bullet_{Y/S})
\ar@{=}[r]
&
Rf_*\Omega^\bullet_{X \times_S Y/S}
}
$$
Ici, la première flèche utilise les unités $\text{id} \to Rp_* p^{-1}$
et $\text{id} \to Rq_* q^{-1}$ des adjonctions, ainsi que les
identifications $Rf_* p^{-1} = Ra_* Rp_* p^{-1}$ et
$Rf_* q^{-1} = Rb_* Rq_* q^{-1}$.
La deuxième flèche est le cup-produit relatif de
Cohomologie, remarque \ref{cohomology-remark-cup-product}.
La troisième flèche est le morphisme du produit tensoriel dérivé
de complexes vers le complexe total de leur produit tensoriel.
L'égalité finale est le lemme \ref{derham-lemma-de-rham-complex-product}.
Sur les sections globales, cette construction redonne celle présentée plus haut.

\begin{lemma}
\label{derham-lemma-kunneth-de-rham-relative}
Supposons $X \to S$ et $Y \to S$ lisses et quasi-compacts,
et les morphismes $X \to X \times_S X$ et $Y \to Y \times_S Y$ affines.
Alors le cup-produit relatif
$$
Ra_*\Omega^\bullet_{X/S}
\otimes_{\mathcal{O}_S}^\mathbf{L}
Rb_*\Omega^\bullet_{Y/S}
\longrightarrow
Rf_*\Omega^\bullet_{X \times_S Y/S}
$$
est un isomorphisme dans $D(\mathcal{O}_S)$.
\end{lemma}

\begin{proof}
Conséquence immédiate du lemme \ref{derham-lemma-kunneth-de-rham}.
\end{proof}







\section{Première classe de Chern en cohomologie de de Rham}
\label{derham-section-first-chern-class}

\noindent
Soit $X \to S$ un morphisme de schémas. Il existe un morphisme de complexes
$$
\text{d}\log : \mathcal{O}_X^*[-1] \longrightarrow \Omega^\bullet_{X/S}
$$
qui envoie la section $g \in \mathcal{O}_X^*(U)$ sur la section
$\text{d}\log(g) = g^{-1}\text{d}g$ de $\Omega^1_{X/S}(U)$.
On peut donc considérer l'application
$$
\Pic(X) = H^1(X, \mathcal{O}_X^*) =
H^2(X, \mathcal{O}_X^*[-1]) \longrightarrow H^2_{dR}(X/S)
$$
où la première égalité est donnée par
Cohomologie, lemme \ref{cohomology-lemma-h1-invertible}.
L'image de la classe d'isomorphisme du module inversible
$\mathcal{L}$ est notée $c^{dR}_1(\mathcal{L}) \in H^2_{dR}(X/S)$.

\medskip\noindent
On peut également utiliser le morphisme $\text{d}\log : \mathcal{O}_X^* \to \Omega^1_{X/S}$
pour définir une classe de Chern en cohomologie de Hodge
$$
c_1^{Hodge} : \Pic(X) \longrightarrow H^1(X, \Omega^1_{X/S})
\subset H^2_{Hodge}(X/S)
$$
Ces constructions sont compatibles avec les images inverses.

\begin{lemma}
\label{derham-lemma-pullback-c1}
Étant donné un diagramme commutatif
$$
\xymatrix{
X' \ar[r]_f \ar[d] & X \ar[d] \\
S' \ar[r] & S
}
$$
de schémas, les diagrammes
$$
\xymatrix{
\Pic(X') \ar[d]_{c_1^{dR}} &
\Pic(X) \ar[d]^{c_1^{dR}} \ar[l]^{f^*} \\
H^2_{dR}(X'/S') &
H^2_{dR}(X/S) \ar[l]_{f^*}
}
\quad
\xymatrix{
\Pic(X') \ar[d]_{c_1^{Hodge}} &
\Pic(X) \ar[d]^{c_1^{Hodge}} \ar[l]^{f^*} \\
H^1(X', \Omega^1_{X'/S'}) &
H^1(X, \Omega^1_{X/S}) \ar[l]_{f^*}
}
$$
sont commutatifs.
\end{lemma}

\begin{proof}
Démonstration omise.
\end{proof}

\noindent
« Calculons » l'élément $c^{dR}_1(\mathcal{L})$ en cohomologie de {\v C}ech
(avec les conventions de signe pour les différentielles de {\v C}ech
de Cohomologie, section
\ref{cohomology-section-cech-cohomology-of-complexes}).
Choisissons un recouvrement ouvert
$\mathcal{U} : X = \bigcup_{i \in I} U_i$ tel que
l'on dispose d'une section trivialisante $s_i$ de $\mathcal{L}|_{U_i}$ pour tout $i$.
Sur les intersections $U_{i_0i_1} = U_{i_0} \cap U_{i_1}$,
on dispose d'une fonction inversible $f_{i_0i_1}$ telle que
$f_{i_0i_1} = s_{i_1}|_{U_{i_0i_1}} s_{i_0}|_{U_{i_0i_1}}^{-1}$\footnote{La
différentielle de {\v C}ech d'un cycle de degré $0$, $\{a_{i_0}\}$,
vaut $a_{i_1} - a_{i_0}$ sur $U_{i_0i_1}$.}.
Bien entendu, on a
$$
f_{i_1i_2}|_{U_{i_0i_1i_2}}
f_{i_0i_2}^{-1}|_{U_{i_0i_1i_2}}
f_{i_0i_1}|_{U_{i_0i_1i_2}} = 1
$$
La classe de cohomologie de $\mathcal{L}$ dans $H^1(X, \mathcal{O}_X^*)$ est
l'image de la classe de cohomologie de {\v C}ech du cocycle $\{f_{i_0i_1}\}$ dans
$\check{\mathcal{C}}^\bullet(\mathcal{U}, \mathcal{O}_X^*)$.
On voit donc que $c_1^{dR}(\mathcal{L})$ est l'image
de la classe de cohomologie associée au cocycle de {\v C}ech
$\{\alpha_{i_0 \ldots i_p}\}$ dans
$\text{Tot}(\check{\mathcal{C}}^\bullet(\mathcal{U}, \Omega_{X/S}^\bullet))$,
de degré $2$, donné par
\begin{enumerate}
\item $\alpha_{i_0} = 0$ dans $\Omega^2_{X/S}(U_{i_0})$,
\item $\alpha_{i_0i_1} = f_{i_0i_1}^{-1}\text{d}f_{i_0i_1}$ dans
$\Omega^1_{X/S}(U_{i_0i_1})$, et
\item $\alpha_{i_0i_1i_2} = 0$ dans $\mathcal{O}_{X/S}(U_{i_0i_1i_2})$.
\end{enumerate}
Supposons donnés des modules inversibles $\mathcal{L}_k$, $k = 1, \ldots, a$,
chacun trivialisé sur $U_i$ pour tout $i \in I$, donnant des cocycles
$f_{k, i_0i_1}$ et $\alpha_k = \{\alpha_{k, i_0 \ldots i_p}\}$ comme ci-dessus.
En utilisant la règle de
Cohomologie, section \ref{cohomology-section-cech-cohomology-of-complexes},
on calcule
$$
\beta = \alpha_1 \cup \alpha_2 \cup \ldots \cup \alpha_a
$$
comme le cocycle $\beta = \{\beta_{i_0 \ldots i_p}\}$
décrit comme suit :
\begin{enumerate}
\item $\beta_{i_0 \ldots i_p} = 0$ dans
$\Omega^{2a - p}_{X/S}(U_{i_0 \ldots i_p})$ sauf si $p = a$, et
\item $\beta_{i_0 \ldots i_a} = (-1)^{a(a - 1)/2}
\alpha_{1, i_0i_1} \wedge \alpha_{2, i_1 i_2} \wedge \ldots \wedge
\alpha_{a, i_{a - 1}i_a}$ dans
$\Omega^a_{X/S}(U_{i_0 \ldots i_a})$.
\end{enumerate}
C'est donc un cocycle représentant
$c_1^{dR}(\mathcal{L}_1) \cup \ldots \cup c_1^{dR}(\mathcal{L}_a)$.
Bien entendu, le même calcul montre que le cocycle
$\{\beta_{i_0 \ldots i_a}\}$ dans
$\check{\mathcal{C}}^a(\mathcal{U}, \Omega_{X/S}^a))$
représente la classe de cohomologie
$c_1^{Hodge}(\mathcal{L}_1) \cup \ldots \cup c_1^{Hodge}(\mathcal{L}_a)$

\begin{remark}
\label{derham-remark-truncations}
Voici une reformulation des calculs ci-dessus en termes plus abstraits.
Soit $p : X \to S$ un morphisme de schémas. Soit $\mathcal{L}$ un
$\mathcal{O}_X$-module inversible. Si l'on considère $\text{d}\log$ comme un morphisme
$$
\mathcal{O}_X^*[-1] \to \sigma_{\geq 1}\Omega^\bullet_{X/S}
$$,
l'identification $\Pic(X) = H^1(X, \mathcal{O}_X^*)$ utilisée plus haut fournit une
classe de cohomologie
$$
\gamma_1(\mathcal{L}) \in H^2(X, \sigma_{\geq 1}\Omega^\bullet_{X/S})
$$
L'image de $\gamma_1(\mathcal{L})$ par le morphisme
$\sigma_{\geq 1}\Omega^\bullet_{X/S} \to \Omega^\bullet_{X/S}$
redonne $c_1^{dR}(\mathcal{L})$. En particulier, on voit que
$c_1^{dR}(\mathcal{L}) \in F^1H^2_{dR}(X/S)$ ; voir la
section \ref{derham-section-hodge-filtration}. L'image de $\gamma_1(\mathcal{L})$
par le morphisme $\sigma_{\geq 1}\Omega^\bullet_{X/S} \to \Omega^1_{X/S}[-1]$
redonne $c_1^{Hodge}(\mathcal{L})$. En prenant le cup-produit
(voir section \ref{derham-section-hodge-filtration}), on obtient
$$
\xi = \gamma_1(\mathcal{L}_1) \cup \ldots \cup \gamma_1(\mathcal{L}_a) \in
H^{2a}(X, \sigma_{\geq a}\Omega^\bullet_{X/S})
$$
Les diagrammes commutatifs de la section \ref{derham-section-hodge-filtration}
montrent que $\xi$ a pour image
$c_1^{dR}(\mathcal{L}_1) \cup \ldots \cup c_1^{dR}(\mathcal{L}_a)$
dans $H^{2a}_{dR}(X/S)$ par le morphisme
$\sigma_{\geq a}\Omega^\bullet_{X/S} \to \Omega^\bullet_{X/S}$.
Il en résulte aussi que
$c_1^{dR}(\mathcal{L}_1) \cup \ldots \cup c_1^{dR}(\mathcal{L}_a)$
appartient à $F^a H^{2a}_{dR}(X/S)$. De même, le morphisme
$\sigma_{\geq a}\Omega^\bullet_{X/S} \to \Omega^a_{X/S}[-a]$
envoie $\xi$ sur
$c_1^{Hodge}(\mathcal{L}_1) \cup \ldots \cup c_1^{Hodge}(\mathcal{L}_a)$
dans $H^a(X, \Omega^a_{X/S})$.
\end{remark}

\begin{remark}
\label{derham-remark-log-forms}
Soit $p : X \to S$ un morphisme de schémas. Pour $i > 0$,
notons $\Omega^i_{X/S, log} \subset \Omega^i_{X/S}$ le sous-faisceau abélien
engendré par les sections locales de la forme
$$
\text{d}\log(u_1) \wedge \ldots \wedge \text{d}\log(u_i)
$$
où $u_1, \ldots, u_n$ sont des sections locales inversibles de $\mathcal{O}_X$.
Pour $i = 0$, le sous-faisceau $\Omega^0_{X/S, log} \subset \mathcal{O}_X$
est l'image de $\mathbf{Z} \to \mathcal{O}_X$. Pour tout $i \geq 0$,
on dispose d'un morphisme de complexes
$$
\Omega^i_{X/S, log}[-i] \longrightarrow \Omega^\bullet_{X/S}
$$
car la différentielle d'une forme logarithmique est nulle. De plus, le produit extérieur
de formes logarithmiques est encore logarithmique ; on obtient donc des applications bilinéaires
$$
\wedge :  \Omega^i_{X/S, log} \times
\Omega^j_{X/S, log} \longrightarrow \Omega^{i + j}_{X/S, log}
$$
compatibles avec (\ref{derham-equation-wedge}) et avec les morphismes ci-dessus.
Soit $\mathcal{L}$ un $\mathcal{O}_X$-module inversible.
En utilisant le morphisme de faisceaux abéliens
$\text{d}\log : \mathcal{O}_X^* \to \Omega^1_{X/S, log}$
et l'identification $\Pic(X) = H^1(X, \mathcal{O}_X^*)$,
on obtient une classe de cohomologie canonique
$$
\tilde \gamma_1(\mathcal{L}) \in H^1(X, \Omega^1_{X/S, log})
$$
Ces classes possèdent les propriétés suivantes :
\begin{enumerate}
\item l'image de $\tilde \gamma_1(\mathcal{L})$ par le morphisme canonique
$\Omega^1_{X/S, log}[-1] \to \sigma_{\geq 1}\Omega^\bullet_{X/S}$
s'obtient en envoyant $\tilde \gamma_1(\mathcal{L})$ sur la classe
$\gamma_1(\mathcal{L}) \in 
H^2(X, \sigma_{\geq 1}\Omega^\bullet_{X/S})$
de la remarque \ref{derham-remark-truncations},
\item l'image de $\tilde \gamma_1(\mathcal{L})$ par le morphisme canonique
$\Omega^1_{X/S, log}[-1] \to \Omega^\bullet_{X/S}$
s'obtient en envoyant $\tilde \gamma_1(\mathcal{L})$ sur $c_1^{dR}(\mathcal{L})$ dans
$H^2_{dR}(X/S)$,
\item l'image de $\tilde \gamma_1(\mathcal{L})$ par le morphisme canonique
$\Omega^1_{X/S, log} \to \Omega^1_{X/S}$
s'obtient en envoyant $\tilde \gamma_1(\mathcal{L})$ sur $c_1^{Hodge}(\mathcal{L})$ dans
$H^1(X, \Omega^1_{X/S})$,
\item la construction de ces classes est compatible avec les images inverses,
\item ajouter d'autres propriétés ici.
\end{enumerate}
\end{remark}





\section{Cohomologie de de Rham d'un fibré en droites}
\label{derham-section-line-bundle}

\noindent
Un fibré en droites est un cas particulier de fibré vectoriel, lequel est
un cône muni d'une structure supplémentaire. Pour étudier convenablement
leur complexe de de Rham, il convient d'examiner le complexe de de Rham
d'un anneau gradué.

\begin{remark}[Complexe de de Rham d'un anneau gradué]
\label{derham-remark-de-rham-complex-graded}
Soit $G$ un monoïde abélien noté additivement, d'élément neutre $0$.
Soit $R \to A$ un morphisme d'anneaux, et supposons $A$ muni d'une graduation
$A = \bigoplus_{g \in G} A_g$ par des $R$-modules telle que $R$ s'envoie dans $A_0$
et que $A_g \cdot A_{g'} \subset A_{g + g'}$. Le module des différentielles
est alors muni d'une graduation
$$
\Omega_{A/R} = \bigoplus\nolimits_{g \in G} \Omega_{A/R, g}
$$
où $\Omega_{A/R, g}$ est le $R$-sous-module de $\Omega_{A/R}$
engendré par les $a_0 \text{d}a_1$ avec $a_i \in A_{g_i}$ tels que
$g = g_0 + g_1$. De même, on obtient
$$
\Omega^p_{A/R} = \bigoplus\nolimits_{g \in G} \Omega^p_{A/R, g}
$$
où $\Omega^p_{A/R, g}$ est le $R$-sous-module de $\Omega^p_{A/R}$
engendré par les $a_0 \text{d}a_1 \wedge \ldots \wedge \text{d}a_p$
avec $a_i \in A_{g_i}$ tels que $g = g_0 + g_1 + \ldots + g_p$.
Bien entendu, les différentielles respectent la graduation et le produit
extérieur est compatible avec les graduations de manière évidente.
\end{remark}

\noindent
Soit $f : X \to S$ un morphisme de schémas. Soit $\pi : C \to X$ un cône ; voir
Constructions, définition \ref{constructions-definition-abstract-cone}.
Rappelons que cela signifie que $\pi$ est affine et que l'on dispose d'une graduation
$\pi_*\mathcal{O}_C = \bigoplus_{n \geq 0} \mathcal{A}_n$ avec
$\mathcal{A}_0 = \mathcal{O}_X$.
En utilisant la discussion de la remarque \ref{derham-remark-de-rham-complex-graded}
sur les ouverts affines, on trouve que\footnote{Avec nos excuses pour la notation !}
$$
\pi_*(\Omega^\bullet_{C/S}) =
\bigoplus\nolimits_{n \geq 0} \Omega^\bullet_{C/S, n}
$$
est canoniquement une somme directe de sous-complexes. De plus, on dispose d'une factorisation
$$
\Omega^\bullet_{X/S} \to \Omega^\bullet_{C/S, 0} \to
\pi_*(\Omega^\bullet_{C/S})
$$
et l'on sait que $\omega \wedge \eta \in \Omega^{p + q}_{C/S, n + m}$
si $\omega \in \Omega^p_{C/S, n}$ et $\eta \in \Omega^q_{C/S, m}$.

\medskip\noindent
Soit $f : X \to S$ un morphisme de schémas. Soit $\pi : L \to X$ le
fibré en droites associé au $\mathcal{O}_X$-module inversible $\mathcal{L}$.
Cela signifie que $\pi$ est l'unique morphisme affine tel que
$$
\pi_*\mathcal{O}_L = 
\bigoplus\nolimits_{n \geq 0} \mathcal{L}^{\otimes n}
$$
comme $\mathcal{O}_X$-algèbres. Ainsi, $L$ est un cône sur $X$.
La discussion ci-dessus fournit une
décomposition canonique en somme directe
$$
\pi_*(\Omega^\bullet_{L/S}) =
\bigoplus\nolimits_{n \geq 0} \Omega^\bullet_{L/S, n}
$$
compatible avec le produit extérieur et avec la décomposition
de $\pi_*\mathcal{O}_L$ ci-dessus, et telle que $\Omega_{X/S}$
s'envoie dans la composante $\Omega_{L/S, 0}$ de degré $0$.

\medskip\noindent
Un autre cas nous sera utile. Considérons en effet le
complémentaire\footnote{Le schéma $L^\star$ est le $\mathbf{G}_m$-torseur
sur $X$ associé à $L$. C'est pourquoi la graduation obtenue ci-dessous est
une $\mathbf{Z}$-graduation ; comparer avec Groupoïdes,
exemple \ref{groupoids-example-Gm-on-affine} et
lemmes \ref{groupoids-lemma-complete-reducibility-Gm} et
\ref{groupoids-lemma-Gm-equivariant-module}.}
$L^\star \subset L$ de la section nulle $o : X \to L$ dans notre fibré
en droites $L$. Un calcul local donne un isomorphisme canonique
$$
(L^\star \to X)_*\mathcal{O}_{L^\star} =
\bigoplus\nolimits_{n \in \mathbf{Z}} \mathcal{L}^{\otimes n}
$$
de $\mathcal{O}_X$-algèbres. Le membre de droite est une algèbre $\mathbf{Z}$-graduée
quasi-cohérente sur $\mathcal{O}_X$. En utilisant la discussion de la
remarque \ref{derham-remark-de-rham-complex-graded} sur les ouverts affines, on obtient
$$
(L^\star \to X)_*(\Omega^\bullet_{L^\star/S}) =
\bigoplus\nolimits_{n \in \mathbf{Z}} \Omega^\bullet_{L^\star/S, n}
$$
de manière compatible avec le produit extérieur et avec la décomposition
de $(L^\star \to X)_*\mathcal{O}_{L^\star}$ ci-dessus, et telle que
$\Omega_{X/S}$ s'envoie dans la composante $\Omega_{L^\star/S, 0}$ de degré $0$.
Le complexe $\Omega^\bullet_{L^\star/S, 0}$ nous intéressera
tout particulièrement.

\begin{lemma}
\label{derham-lemma-the-complex-for-L-star}
Avec les notations ci-dessus, il existe une suite exacte courte de complexes
$$
0 \to \Omega^\bullet_{X/S} \to
\Omega^\bullet_{L^\star/S, 0} \to
\Omega^\bullet_{X/S}[-1] \to 0
$$
\end{lemma}

\begin{proof}
Nous avons construit le morphisme
$\Omega^\bullet_{X/S} \to \Omega^\bullet_{L^\star/S, 0}$ ci-dessus.

\medskip\noindent
Construction de
$\text{Res} : \Omega^\bullet_{L^\star/S, 0} \to \Omega^\bullet_{X/S}[-1]$.
Soit $U \subset X$ un ouvert, et soient $s \in \mathcal{L}(U)$
et $s' \in \mathcal{L}^{\otimes -1}(U)$ des sections telles que
$s' s = 1$. Alors $s$ donne une section inversible du faisceau
d'algèbres $(L^\star \to X)_*\mathcal{O}_{L^\star}$ sur $U$,
d'inverse $s' = s^{-1}$. On peut alors considérer la forme de degré $1$
$\text{d}\log(s) = s' \text{d}(s)$, qui appartient à
$\Omega^1_{L^\star/S, 0}(U)$ d'après notre construction de la graduation sur
$\Omega^1_{L^\star/S}$. Les calculs sur les ouverts affines effectués ci-dessous
montreront que $1$ et $\text{d}\log(s)$ engendrent librement
$\Omega^\bullet_{L^\star/S, 0}|_U$ comme module à droite sur
$\Omega^\bullet_{X/S}|_U$.
On peut donc définir $\text{Res}$ sur $U$ par la règle
$$
\text{Res}(\omega' + \text{d}\log(s) \wedge \omega) = \omega
$$
pour tous $\omega', \omega \in \Omega^\bullet_{X/S}(U)$. Ce
morphisme est indépendant du choix du générateur local $s$ ; il se
recolle donc en un morphisme global. En effet, un autre choix de $s$
serait de la forme $gs$ pour un certain élément inversible $g \in \mathcal{O}_X(U)$,
et l'on aurait $\text{d}\log(gs) = g^{-1}\text{d}(g) + \text{d}\log(s)$,
d'où l'indépendance résulte aisément.
Enfin, observons que notre règle définissant $\text{Res}$
est compatible avec les différentielles,
puisque $\text{d}(\omega' + \text{d}\log(s) \wedge \omega) =
\text{d}(\omega') - \text{d}\log(s) \wedge \text{d}(\omega)$
et que la différentielle sur $\Omega^\bullet_{X/S}[-1]$
envoie $\omega'$ sur $-\text{d}(\omega')$ selon notre convention de signe dans
Homologie, définition \ref{homology-definition-shift-cochain}.

\medskip\noindent
Calcul local. On peut recouvrir $X$ par des ouverts affines $U \subset X$
tels que $\mathcal{L}|_U \cong \mathcal{O}_U$, qui s'envoient en outre
dans un ouvert affine $V \subset S$. Écrivons $U = \Spec(A)$, $V = \Spec(R)$,
et choisissons un générateur $s$ de $\mathcal{L}$. On obtient
$$
L^\star \times_X U = \Spec(A[s, s^{-1}])
$$
Le calcul des différentielles donne
$$
\Omega^1_{A[s, s^{-1}]/R} =
A[s, s^{-1}] \otimes_A \Omega^1_{A/R} \oplus A[s, s^{-1}] \text{d}\log(s)
$$
et, en prenant les puissances extérieures, on obtient donc
$$
\Omega^p_{A[s, s^{-1}]/R} =
A[s, s^{-1}] \otimes_A \Omega^p_{A/R}
\oplus
A[s, s^{-1}] \text{d}\log(s) \otimes_A \Omega^{p - 1}_{A/R}
$$
En prenant les composantes de degré $0$, on trouve
$$
\Omega^p_{A[s, s^{-1}]/R, 0} =
\Omega^p_{A/R} \oplus \text{d}\log(s) \otimes_A \Omega^{p - 1}_{A/R}
$$
ce qui achève la démonstration du lemme.
\end{proof}

\begin{lemma}
\label{derham-lemma-the-complex-for-L-star-gives-chern-class}
Le morphisme de « bord »
$\delta : \Omega^\bullet_{X/S} \to \Omega^\bullet_{X/S}[2]$
dans $D(X, f^{-1}\mathcal{O}_S)$ provenant de
la suite exacte courte du lemme \ref{derham-lemma-the-complex-for-L-star}
est le morphisme de la remarque \ref{derham-remark-cup-product-as-a-map}
pour $\xi = c_1^{dR}(\mathcal{L})$.
\end{lemma}

\begin{proof}
Plus précisément, considérons le décalage
$$
0 \to \Omega^\bullet_{X/S}[1] \to
\Omega^\bullet_{L^\star/S, 0}[1] \to
\Omega^\bullet_{X/S} \to 0
$$
de la suite exacte courte du lemme \ref{derham-lemma-the-complex-for-L-star}.
Comme composante de degré zéro d'une graduation sur
$(L^\star \to X)_*\Omega^\bullet_{L^\star/S}$,
$\Omega^\bullet_{L^\star/S, 0}$ est une
$\mathcal{O}_X$-algèbre différentielle graduée, et le morphisme
$\Omega^\bullet_{X/S} \to \Omega^\bullet_{L^\star/S, 0}$
est un homomorphisme de $\mathcal{O}_X$-algèbres différentielles graduées.
On peut donc considérer $\Omega^\bullet_{X/S}[1] \to
\Omega^\bullet_{L^\star/S, 0}[1]$ comme un morphisme de
$\Omega^\bullet_{X/S}$-modules différentiels gradués à droite sur $X$. Le morphisme
$\text{Res} : \Omega^\bullet_{L^\star/S, 0}[1] \to \Omega^\bullet_{X/S}$
est un morphisme de $\Omega^\bullet_{X/S}$-modules différentiels gradués à droite,
car il est défini localement par la règle
$\text{Res}(\omega' + \text{d}\log(s) \wedge \omega) = \omega$ ; voir
la démonstration du lemme \ref{derham-lemma-the-complex-for-L-star}.
Ainsi, d'après la discussion de
Faisceaux différentiels gradués, section \ref{sdga-section-misc},
$\delta$ provient d'un morphisme
$\delta' : \Omega^\bullet_{X/S} \to \Omega^\bullet_{X/S}[2]$
dans la catégorie dérivée $D(\Omega^\bullet_{X/S}, \text{d})$
des modules différentiels gradués à droite sur le complexe de de Rham.
L'unicité énoncée dans la remarque \ref{derham-remark-cup-product-as-a-map}
montre qu'il suffit de prouver que $\delta(1) = c_1^{dR}(\mathcal{L})$.

\medskip\noindent
Nous affirmons qu'il existe un diagramme commutatif
$$
\xymatrix{
0 \ar[r] &
\mathcal{O}_X^* \ar[r] \ar[d]_{\text{d}\log} &
E \ar[r] \ar[d] &
\underline{\mathbf{Z}} \ar[d] \ar[r] &
0 \\
0 \ar[r] &
\Omega^\bullet_{X/S}[1] \ar[r] &
\Omega^\bullet_{L^\star/S, 0}[1] \ar[r] &
\Omega^\bullet_{X/S} \ar[r] &
0
}
$$
dont la ligne supérieure est une suite exacte courte de faisceaux abéliens dont le
morphisme de bord envoie $1$ sur la classe de $\mathcal{L}$ dans
$H^1(X, \mathcal{O}_X^*)$. Il suffit de démontrer cette assertion,
par compatibilité des morphismes de bord avec les morphismes de suites
exactes courtes. On définit $E$ comme le faisceau associé à la règle
$$
U \longmapsto \{(s, n) \mid
n \in \mathbf{Z},\ s \in \mathcal{L}^{\otimes n}(U)\text{ generator}\}
$$,
la structure de groupe étant donnée par $(s, n) \cdot (t, m) = (s \otimes t, n + m)$.
Le morphisme vertical du milieu envoie $(s, n)$ sur $\text{d}\log(s)$. On obtient ainsi
un morphisme de suites exactes courtes,
car le morphisme $Res : \Omega^1_{L^\star/S, 0} \to \mathcal{O}_X$
construit dans la démonstration du lemme \ref{derham-lemma-the-complex-for-L-star} envoie
$\text{d}\log(s)$ sur $1$ si $s$ est un générateur local de $\mathcal{L}$.
Pour calculer le bord de $1$ dans la ligne supérieure, choisissons des trivialisations locales
$s_i$ de $\mathcal{L}$ sur des ouverts $U_i$ comme à la
section \ref{derham-section-first-chern-class}. Sur les intersections
$U_{i_0i_1} = U_{i_0} \cap U_{i_1}$,
on dispose d'une fonction inversible $f_{i_0i_1}$ telle que
$f_{i_0i_1} = s_{i_1}|_{U_{i_0i_1}} s_{i_0}|_{U_{i_0i_1}}^{-1}$,
et la classe de cohomologie de $\mathcal{L}$ est donnée par le cocycle de {\v C}ech
$\{f_{i_0i_1}\}$. On a alors, bien entendu,
$$
(f_{i_0i_1}, 0) = (s_{i_1}, 1)|_{U_{i_0i_1}} \cdot
(s_{i_0}, 1)|_{U_{i_0i_1}}^{-1}
$$
comme sections de $E$, ce qui achève la démonstration.
\end{proof}

\begin{lemma}
\label{derham-lemma-push-omega-a}
Avec les notations ci-dessus, on a
\begin{enumerate}
\item $\Omega^p_{L^\star/S, n} =
\Omega^p_{L^\star/S, 0} \otimes_{\mathcal{O}_X} \mathcal{L}^{\otimes n}$
pour tout $n \in \mathbf{Z}$, comme $\mathcal{O}_X$-modules quasi-cohérents,
\item $\Omega^\bullet_{X/S} = \Omega^\bullet_{L/X, 0}$
comme complexes, et
\item pour $n > 0$ et $p \geq 0$, on a
$\Omega^p_{L/X, n} = \Omega^p_{L^\star/S, n}$.
\end{enumerate}
\end{lemma}

\begin{proof}
Dans chaque cas, il existe un morphisme canonique défini globalement,
qui est un isomorphisme par des calculs locaux que nous omettons.
\end{proof}

\begin{lemma}
\label{derham-lemma-line-bundle-characteristic-zero}
Dans la situation ci-dessus, supposons qu'il existe un morphisme $S \to \Spec(\mathbf{Q})$.
Alors $\Omega^\bullet_{X/S} \to \pi_*\Omega^\bullet_{L/S}$ est un
quasi-isomorphisme et $H_{dR}^*(X/S) = H_{dR}^*(L/S)$.
\end{lemma}

\begin{proof}
Soit $R$ une $\mathbf{Q}$-algèbre. Soit $A$ une $R$-algèbre.
L'énoncé local affine affirme que le morphisme
$$
\Omega^\bullet_{A/R} \longrightarrow \Omega^\bullet_{A[t]/R}
$$
est un quasi-isomorphisme de complexes de $R$-modules. Il s'agit en fait d'une
équivalence d'homotopie dont un inverse à homotopie près est le morphisme envoyant
$g \omega + g' \text{d}t \wedge \omega'$ sur $g(0)\omega$ pour
$g, g' \in A[t]$ et $\omega, \omega' \in \Omega^\bullet_{A/R}$.
L'homotopie envoie $g \omega + g' \text{d}t \wedge \omega'$
sur $(\int g') \omega'$, où $\int g' \in A[t]$ est le polynôme
de terme constant nul dont la dérivée par rapport à $t$
est $g'$. Bien entendu, on utilise ici le fait que $R$ contient $\mathbf{Q}$,
puisque $\int t^n = (1/n)t^{n + 1}$.
\end{proof}

\begin{example}
\label{derham-example-affine-line}
Le lemme \ref{derham-lemma-line-bundle-characteristic-zero} est
faux en caractéristique positive. Le complexe de de Rham de
$\mathbf{A}^1_k = \Spec(k[x])$ sur un corps $k$ se présente comme une somme directe
$$
k \oplus
\bigoplus\nolimits_{n \geq 1}
(k \cdot t^n \xrightarrow{n}
k \cdot t^{n - 1} \text{d}t)
$$
Ainsi, si la caractéristique de $k$ est $p > 0$,
on voit que $H^0_{dR}(\mathbf{A}^1_k/k)$ et
$H^1_{dR}(\mathbf{A}^1_k/k)$
sont tous deux de dimension infinie sur $k$.
\end{example}








\section{Cohomologie de de Rham de l'espace projectif}
\label{derham-section-projective-space}

\noindent
Soit $A$ un anneau. Soit $n \geq 1$. Le morphisme structural
$\mathbf{P}^n_A \to \Spec(A)$ est propre et lisse de dimension relative
$n$. Il est lisse de dimension relative $n$ et de type fini,
car $\mathbf{P}^n_A$ possède un recouvrement ouvert affine fini par des schémas tous
isomorphes à $\mathbf{A}^n_A$ ; voir Constructions, lemme
\ref{constructions-lemma-standard-covering-projective-space}.
Il est propre, car il est aussi séparé et universellement fermé,
d'après Constructions, lemme \ref{constructions-lemma-projective-space-separated}.
Notons $\mathcal{O}$ et $\mathcal{O}(d)$ respectivement le faisceau structural
$\mathcal{O}_{\mathbf{P}^n_A}$ et les faisceaux tordus de Serre
$\mathcal{O}_{\mathbf{P}^n_A}(d)$.
Notons $\Omega = \Omega_{\mathbf{P}^n_A/A}$ le faisceau
des différentielles relatives et $\Omega^p$ ses puissances extérieures.

\begin{lemma}
\label{derham-lemma-euler-sequence}
Il existe une suite exacte courte
$$
0 \to \Omega \to \mathcal{O}(-1)^{\oplus n + 1} \to \mathcal{O} \to 0
$$
\end{lemma}

\begin{proof}
Pour l'expliciter, rappelons que
$\mathbf{P}^n_A = \text{Proj}(A[T_0, \ldots, T_n])$,
et écrivons symboliquement
$$
\mathcal{O}(-1)^{\oplus n + 1} =
\bigoplus\nolimits_{j = 0, \ldots, n} \mathcal{O}(-1) \text{d}T_j
$$
La première flèche
$$
\Omega \to
\bigoplus\nolimits_{j = 0, \ldots, n} \mathcal{O}(-1) \text{d}T_j
$$
de la suite exacte courte ci-dessus
est donnée, sur chacun des ouverts standards
$D_+(T_i) = \Spec(A[T_0/T_i, \ldots, T_n/T_i])$
mentionnés plus haut, par la règle
$$
\sum\nolimits_{j \not = i} g_j \text{d}(T_j/T_i)
\longmapsto
\sum\nolimits_{j \not = i} g_j/T_i \text{d}T_j
- (\sum\nolimits_{j \not = i} g_jT_j/T_i^2) \text{d}T_i
$$
Cette expression a un sens, puisque $1/T_i$ est une section de $\mathcal{O}(-1)$
sur $D_+(T_i)$. Le morphisme
$$
\bigoplus\nolimits_{j = 0, \ldots, n} \mathcal{O}(-1) \text{d}T_j
\to
\mathcal{O}
$$
est défini en envoyant $\text{d}T_j$ sur $T_j$ ; plus précisément, sur
$D_+(T_i)$, on envoie la section $\sum g_j \text{d}T_j$ sur
$\sum T_jg_j$. Nous omettons de vérifier que l'on obtient ainsi
une suite exacte courte.
\end{proof}

\noindent
Étant donnés un entier $k \in \mathbf{Z}$ et un
$\mathcal{O}_{\mathbf{P}^n_A}$-module quasi-cohérent $\mathcal{F}$,
notons comme d'habitude $\mathcal{F}(k)$ le tordu de Serre d'indice $k$ de $\mathcal{F}$.
Voir Constructions, définition \ref{constructions-definition-twist}.

\begin{lemma}
\label{derham-lemma-twisted-hodge-cohomology-projective-space}
Dans la situation ci-dessus, les groupes de cohomologie vérifient
\begin{enumerate}
\item $H^q(\mathbf{P}^n_A, \Omega^p) = 0$
sauf si $0 \leq p = q \leq n$,
\item pour $0 \leq p \leq n$, le $A$-module
$H^p(\mathbf{P}^n_A, \Omega^p)$ est libre de rang $1$.
\item pour $q > 0$, $k > 0$ et $p$ quelconque, on a
$H^q(\mathbf{P}^n_A, \Omega^p(k)) = 0$, et
\item ajouter d'autres propriétés ici.
\end{enumerate}
\end{lemma}

\begin{proof}
Nous utiliserons les résultats de Cohomologie des schémas, lemme
\ref{coherent-lemma-cohomology-projective-space-over-ring},
sans plus les mentionner. En particulier, les assertions sont vraies
pour $H^q(\mathbf{P}^n_A, \mathcal{O}(k))$.

\medskip\noindent
Démonstration pour $p = 1$. Considérons la suite exacte courte
$$
0 \to \Omega \to \mathcal{O}(-1)^{\oplus n + 1} \to \mathcal{O} \to 0
$$
du lemme \ref{derham-lemma-euler-sequence}. Puisque la cohomologie de $\mathcal{O}(-1)$
est nulle en tout degré, il en résulte que
$H^q(\mathbf{P}^n_A, \Omega)$ est nul sauf en degré $1$,
où il est librement engendré par le bord de $1$ dans
$H^0(\mathbf{P}^n_A, \mathcal{O})$.

\medskip\noindent
Supposons $p > 1$. Considérons la suite exacte courte
ci-dessus comme définissant une filtration à $2$ crans sur $\mathcal{O}(-1)^{\oplus n + 1}$.
La filtration induite sur $\wedge^p\mathcal{O}(-1)^{\oplus n + 1}$ se présente
comme suit :
$$
0 \to \Omega^p \to \wedge^p\left(\mathcal{O}(-1)^{\oplus n + 1}\right)
\to \Omega^{p - 1} \to 0
$$
Observons que $\wedge^p\mathcal{O}(-1)^{\oplus n + 1}$ est isomorphe
à une somme directe comptant autant de facteurs qu'il existe de façons de
choisir, parmi $n + 1$ éléments, $p$ éléments, chaque facteur étant
isomorphe à $\mathcal{O}(-p)$,
et que sa cohomologie est donc nulle en tout degré.
Par l'hypothèse de récurrence, cela montre que $H^q(\mathbf{P}^n_A, \Omega^p)$
est nul sauf si $q = p$, et que $H^p(\mathbf{P}^n_A, \Omega^p)$ est libre
de rang $1$, engendré par le bord du générateur de
$H^{p - 1}(\mathbf{P}^n_A, \Omega^{p - 1})$.

\medskip\noindent
Soit $k > 0$. Observons que $\Omega^n = \mathcal{O}(-n - 1)$, par exemple
en utilisant la suite exacte courte ci-dessus pour $p = n + 1$.
La cohomologie de $\Omega^n(k)$ est donc nulle en degrés strictement positifs.
En utilisant les suites exactes courtes
$$
0 \to \Omega^p(k) \to \wedge^p\left(\mathcal{O}(-1)^{\oplus n + 1}\right)(k)
\to \Omega^{p - 1}(k) \to 0
$$
et une récurrence {\it descendante} sur $p$, on obtient l'annulation de la
cohomologie de $\Omega^p(k)$ en degrés strictement positifs pour tout $p$.
\end{proof}

\begin{lemma}
\label{derham-lemma-hodge-cohomology-projective-space}
On a $H^q(\mathbf{P}^n_A, \Omega^p) = 0$
sauf si $0 \leq p = q \leq n$. Pour $0 \leq p \leq n$, le $A$-module
$H^p(\mathbf{P}^n_A, \Omega^p)$ est libre de rang $1$, de base
$c_1^{Hodge}(\mathcal{O}(1))^p$.
\end{lemma}

\begin{proof}
L'annulation et la liberté résultent du
lemme \ref{derham-lemma-twisted-hodge-cohomology-projective-space}.
Pour $p = 0$, il est bien vrai que
$1 \in H^0(\mathbf{P}^n_A, \mathcal{O})$ est un générateur.

\medskip\noindent
Démonstration pour $p = 1$. Considérons la suite exacte courte
$$
0 \to \Omega \to \mathcal{O}(-1)^{\oplus n + 1} \to \mathcal{O} \to 0
$$
du lemme \ref{derham-lemma-euler-sequence}. Dans la démonstration du
lemme \ref{derham-lemma-twisted-hodge-cohomology-projective-space},
nous avons vu que le générateur de $H^1(\mathbf{P}^n_A, \Omega)$
est le bord $\xi$ de $1 \in H^0(\mathbf{P}^n_A, \mathcal{O})$.
Comme dans la démonstration du lemme \ref{derham-lemma-euler-sequence}, identifions
$\mathcal{O}(-1)^{\oplus n + 1}$ à
$\bigoplus_{j = 0, \ldots, n} \mathcal{O}(-1)\text{d}T_j$.
Considérons le recouvrement ouvert
$$
\mathcal{U} : 
\mathbf{P}^n_A =
\bigcup\nolimits_{i = 0, \ldots, n} D_{+}(T_i)
$$
La restriction de la section globale $1$ de $\mathcal{O}$
à $U_i = D_+(T_i)$ se relève en la section $T_i^{-1} \text{d}T_i$ de
$\bigoplus \mathcal{O}(-1)\text{d}T_j$ sur $U_i$. Ainsi, le cocycle
représentant $\xi$ est donné par
$$
T_{i_1}^{-1} \text{d}T_{i_1} - T_{i_0}^{-1} \text{d}T_{i_0} =
\text{d}\log(T_{i_1}/T_{i_0}) \in \Omega(U_{i_0i_1})
$$
D'autre part, pour tout $i$, la section $T_i$ trivialise
$\mathcal{O}(1)$ sur $U_i$. On voit donc que
$f_{i_0i_1} = T_{i_1}/T_{i_0} \in \mathcal{O}^*(U_{i_0i_1})$
est le cocycle représentant $\mathcal{O}(1)$ dans $\Pic(\mathbf{P}^n_A)$ ;
voir la section \ref{derham-section-first-chern-class}.
Par conséquent, $c_1^{Hodge}(\mathcal{O}(1))$
est donné par le cocycle $\text{d}\log(T_{i_1}/T_{i_0})$,
qui coïncide avec celui obtenu pour $\xi$ ci-dessus.

\medskip\noindent
Démonstration par récurrence pour $p$ quelconque. Les cas initiaux $p = 0, 1$ ont été traités
ci-dessus. Supposons $p > 1$. Dans la démonstration du
lemme \ref{derham-lemma-twisted-hodge-cohomology-projective-space},
nous avons vu que le générateur de $H^p(\mathbf{P}^n_A, \Omega^p)$
est le bord de $c_1^{Hodge}(\mathcal{O}(1))^{p - 1}$
dans la suite exacte longue de cohomologie associée à
$$
0 \to \Omega^p \to \wedge^p\left(\mathcal{O}(-1)^{\oplus n + 1}\right)
\to \Omega^{p - 1} \to 0
$$
D'après le calcul de la section \ref{derham-section-first-chern-class},
la classe de cohomologie $c_1^{Hodge}(\mathcal{O}(1))^{p - 1}$
est représentée, au signe près, par le cocycle de composantes
$$
\beta_{i_0 \ldots i_{p - 1}} =
\text{d}\log(T_{i_1}/T_{i_0}) \wedge
\text{d}\log(T_{i_2}/T_{i_1}) \wedge \ldots \wedge
\text{d}\log(T_{i_{p - 1}}/T_{i_{p - 2}})
$$
dans $\Omega^{p - 1}(U_{i_0 \ldots i_{p - 1}})$. Ces
$\beta_{i_0 \ldots i_{p - 1}}$ se relèvent en les sections
$\tilde \beta_{i_0 \ldots i_{p -1}} =
T_{i_0}^{-1}\text{d}T_{i_0} \wedge \beta_{i_0 \ldots i_{p - 1}}$
de $\wedge^p(\bigoplus \mathcal{O}(-1) \text{d}T_j)$ sur
$U_{i_0 \ldots i_{p - 1}}$. On en déduit que le générateur de
$H^p(\mathbf{P}^n_A, \Omega^p)$ est donné par le cocycle dont
les composantes sont
\begin{align*}
\sum\nolimits_{a = 0}^p (-1)^a
\tilde \beta_{i_0 \ldots \hat{i_a} \ldots i_p}
& =
T_{i_1}^{-1}\text{d}T_{i_1} \wedge \beta_{i_1 \ldots i_p}
+ \sum\nolimits_{a = 1}^p (-1)^a
T_{i_0}^{-1}\text{d}T_{i_0} \wedge
\beta_{i_0 \ldots \hat{i_a} \ldots i_p} \\
& =
(T_{i_1}^{-1}\text{d}T_{i_1} - T_{i_0}^{-1}\text{d}T_{i_0}) \wedge
\beta_{i_1 \ldots i_p} +
T_{i_0}^{-1}\text{d}T_{i_0} \wedge \text{d}(\beta)_{i_0 \ldots i_p} \\
& =
\text{d}\log(T_{i_1}/T_{i_0}) \wedge \beta_{i_1 \ldots i_p}
\end{align*}
considérées comme sections de $\Omega^p$ sur $U_{i_0 \ldots i_p}$.
C'est, au signe près, le cocycle représentant
$c_1^{Hodge}(\mathcal{O}(1))^p$, ce qui achève la démonstration.
\end{proof}

\begin{lemma}
\label{derham-lemma-de-rham-cohomology-projective-space}
Pour $0 \leq i \leq n$, le groupe de cohomologie de de Rham
$H^{2i}_{dR}(\mathbf{P}^n_A/A)$ est un $A$-module libre de rang $1$,
de base $c_1^{dR}(\mathcal{O}(1))^i$.
En tout autre degré, la cohomologie de de Rham de $\mathbf{P}^n_A$
sur $A$ est nulle.
\end{lemma}

\begin{proof}
Considérons la suite spectrale de Hodge vers de Rham de la
section \ref{derham-section-hodge-to-de-rham}.
Le calcul de la cohomologie de Hodge de $\mathbf{P}^n_A$ sur $A$
effectué dans le lemme \ref{derham-lemma-hodge-cohomology-projective-space}
montre que cette suite spectrale dégénère à la page $E_1$.
On voit ainsi que $H^{2i}_{dR}(\mathbf{P}^n_A/A)$ est un
$A$-module libre de rang $1$ pour $0 \leq i \leq n$, et qu'il est nul sinon.
Observons que $c_1^{dR}(\mathcal{O}(1))^i \in H^{2i}_{dR}(\mathbf{P}^n_A/A)$
pour $i = 0, \ldots, n$ et que, pour $i = n$, cet élément est
l'image de $c_1^{Hodge}(\mathcal{L})^n$ par le morphisme de complexes
$$
\Omega^n_{\mathbf{P}^n_A/A}[-n]
\longrightarrow
\Omega^\bullet_{\mathbf{P}^n_A/A}
$$
Cela résulte, par exemple, de la discussion de la remarque \ref{derham-remark-truncations}
ou de la description explicite des cocycles représentant ces classes à la
section \ref{derham-section-first-chern-class}.
La suite spectrale montre que le morphisme induit
$$
H^n(\mathbf{P}^n_A, \Omega^n_{\mathbf{P}^n_A/A}) \longrightarrow
H^{2n}_{dR}(\mathbf{P}^n_A/A)
$$
est un isomorphisme et, puisque $c_1^{Hodge}(\mathcal{L})^n$ engendre
la source (lemme \ref{derham-lemma-hodge-cohomology-projective-space}),
on en déduit que $c_1^{dR}(\mathcal{L})^n$ engendre
le but. La $A$-bilinéarité des cup-produits
entraîne que $c_1^{dR}(\mathcal{L})^i$
engendre également $H^{2i}_{dR}(\mathbf{P}^n_A/A)$ pour
$0 \leq i \leq n$.
\end{proof}









\section{La suite spectrale d'un morphisme lisse}
\label{derham-section-relative-spectral-sequence}

\noindent
Considérons un diagramme commutatif de schémas
$$
\xymatrix{
X \ar[rr]_f \ar[rd]_p & & Y \ar[ld]^q \\
& S
}
$$
où $f$ est un morphisme lisse. On obtient alors une suite exacte courte
localement scindée
$$
0 \to f^*\Omega_{Y/S} \to \Omega_{X/S} \to \Omega_{X/Y} \to 0
$$
d'après Morphismes, lemme \ref{morphisms-lemma-triangle-differentials-smooth}.
Considérons-la comme une filtration décroissante $F$ sur $\Omega_{X/S}$,
avec $F^0\Omega_{X/S} = \Omega_{X/S}$, $F^1\Omega_{X/S} = f^*\Omega_{Y/S}$ et
$F^2\Omega_{X/S} = 0$. En appliquant le foncteur $\wedge^p$, on obtient,
pour tout $p$, une filtration induite
$$
\Omega^p_{X/S} = F^0\Omega^p_{X/S} \supset
F^1\Omega^p_{X/S} \supset
F^2\Omega^p_{X/S} \supset \ldots \supset F^{p + 1}\Omega^p_{X/S} = 0
$$
dont les quotients successifs sont
$$
\text{gr}^k\Omega^p_{X/S} =
F^k\Omega^p_{X/S}/F^{k + 1}\Omega^p_{X/S} =
f^*\Omega^k_{Y/S} \otimes_{\mathcal{O}_X} \Omega^{p - k}_{X/Y} =
f^{-1}\Omega^k_{Y/S} \otimes_{f^{-1}\mathcal{O}_Y} \Omega^{p - k}_{X/Y}
$$
pour $k = 0, \ldots, p$. En fait, le lecteur peut vérifier, à l'aide de la
règle de Leibniz, que $F^k\Omega^\bullet_{X/S}$ est un sous-complexe de
$\Omega^\bullet_{X/S}$. On munit ainsi $\Omega^\bullet_{X/S}$
d'une structure de complexe filtré. On peut aussi le voir en observant
que
$$
F^k\Omega^\bullet_{X/S} =
\Im\left(\wedge :
\text{Tot}(
f^{-1}\sigma_{\geq k}\Omega^\bullet_{Y/S} \otimes_{p^{-1}\mathcal{O}_S}
\Omega^\bullet_{X/S})
\longrightarrow
\Omega^\bullet_{X/S}\right)
$$
est l'image d'un morphisme de complexes sur $X$. Le complexe filtré
$$
\Omega^\bullet_{X/S} = F^0\Omega^\bullet_{X/S} \supset
F^1\Omega^\bullet_{X/S} \supset F^2\Omega^\bullet_{X/S} \supset \ldots
$$
a pour composantes graduées associées
$$
\text{gr}^k\Omega^\bullet_{X/S} = 
f^{-1}\Omega^k_{Y/S}[-k] \otimes_{f^{-1}\mathcal{O}_Y} \Omega^\bullet_{X/Y}
$$
d'après ce qui précède.

\begin{lemma}
\label{derham-lemma-spectral-sequence-smooth}
Soit $f : X \to Y$ un morphisme quasi-compact, quasi-séparé et lisse
de schémas sur un schéma de base $S$. Il existe une suite spectrale bornée
de première page
$$
E_1^{p, q} =
H^q(\Omega^p_{Y/S} \otimes_{\mathcal{O}_Y}^\mathbf{L} Rf_*\Omega^\bullet_{X/Y})
$$
convergeant vers $R^{p + q}f_*\Omega^\bullet_{X/S}$.
\end{lemma}

\begin{proof}
Considérons $\Omega^\bullet_{X/S}$ comme un complexe filtré muni de la
filtration introduite ci-dessus. La suite spectrale est celle de
Cohomologie, lemme
\ref{cohomology-lemma-relative-spectral-sequence-filtered-object}.
D'après Catégories dérivées des schémas, lemme
\ref{perfect-lemma-cohomology-de-rham-base-change}, on a
$$
Rf_*\text{gr}^k\Omega^\bullet_{X/S} =
\Omega^k_{Y/S}[-k] \otimes_{\mathcal{O}_Y}^\mathbf{L} Rf_*\Omega^\bullet_{X/Y}
$$,
d'où la conclusion.
\end{proof}

\begin{remark}
\label{derham-remark-gauss-manin}
Dans le lemme \ref{derham-lemma-spectral-sequence-smooth}, considérons les faisceaux de cohomologie
$$
\mathcal{H}^q_{dR}(X/Y) = H^q(Rf_*\Omega^\bullet_{X/Y})
$$
Si $f$ est propre en plus d'être lisse et si $S$ est un schéma sur
$\mathbf{Q}$, alors $\mathcal{H}^q_{dR}(X/Y)$ est localement libre de rang fini (insérer
ici une référence ultérieure). Si l'on suppose seulement que les $\mathcal{H}^q_{dR}(X/Y)$
sont des $\mathcal{O}_Y$-modules plats, on obtient (nous omettons le bref argument)
$$
E_1^{p, q} =
\Omega^p_{Y/S} \otimes_{\mathcal{O}_Y} \mathcal{H}^q_{dR}(X/Y)
$$,
et les différentielles de la suite spectrale sont des morphismes
$$
d_1^{p, q} :
\Omega^p_{Y/S} \otimes_{\mathcal{O}_Y} \mathcal{H}^q_{dR}(X/Y)
\longrightarrow
\Omega^{p + 1}_{Y/S} \otimes_{\mathcal{O}_Y} \mathcal{H}^q_{dR}(X/Y)
$$
En particulier, pour $p = 0$, on obtient un morphisme
$d_1^{0, q} : \mathcal{H}^q_{dR}(X/Y) \to
\Omega^1_{Y/S} \otimes_{\mathcal{O}_Y} \mathcal{H}^q_{dR}(X/Y)$
qui est une connexion intégrable
$\nabla$ (insérer ici une référence ultérieure),
et le complexe
$$
\mathcal{H}^q_{dR}(X/Y) \to
\Omega^1_{Y/S} \otimes_{\mathcal{O}_Y} \mathcal{H}^q_{dR}(X/Y) \to
\Omega^2_{Y/S} \otimes_{\mathcal{O}_Y} \mathcal{H}^q_{dR}(X/Y) \to \ldots
$$
dont les différentielles sont données par $d_1^{\bullet, q}$
est le complexe de de Rham de $\nabla$.
La connexion $\nabla$ est appelée {\it connexion de Gauss--Manin}.
\end{remark}






\section{Théorèmes de type Leray--Hirsch}
\label{derham-section-leray-hirsch}

\noindent
Dans cette section, nous montrons que, pour un morphisme lisse et propre,
on peut parfois exprimer la cohomologie de de Rham de la source en fonction
de celle du but.

\begin{lemma}
\label{derham-lemma-relative-global-generation-on-fibres}
Soit $f : X \to Y$ un morphisme lisse et propre de schémas.
Soient $N$ et $n_1, \ldots, n_N \geq 0$ des entiers, et soient
$\xi_i \in H^{n_i}_{dR}(X/Y)$, $1 \leq i \leq N$.
Supposons que, pour tout point $y \in Y$, les images de $\xi_1, \ldots, \xi_N$
dans $H^*_{dR}(X_y/y)$ forment une base sur $\kappa(y)$. Alors le morphisme
$$
\bigoplus\nolimits_{i = 1}^N \mathcal{O}_Y[-n_i]
\longrightarrow
Rf_*\Omega^\bullet_{X/Y}
$$
associé à $\xi_1, \ldots, \xi_N$ est un isomorphisme.
\end{lemma}

\begin{proof}
D'après le lemme \ref{derham-lemma-proper-smooth-de-Rham},
$Rf_*\Omega^\bullet_{X/Y}$ est un objet parfait de $D(\mathcal{O}_Y)$
dont la formation commute à tout changement de base.
Le morphisme du lemme est donc un morphisme $a : K \to L$
entre objets parfaits de $D(\mathcal{O}_Y)$
dont la restriction dérivée à tout point est un isomorphisme,
d'après notre hypothèse sur les fibres. Le cône $C$ de $a$ est alors un objet parfait
de $D(\mathcal{O}_Y)$ (Cohomologie, lemme
\ref{cohomology-lemma-two-out-of-three-perfect}) dont
la restriction dérivée à tout point est nulle. Il en résulte que $C$
est nul, d'après Compléments d'algèbre, lemme
\ref{more-algebra-lemma-lift-perfect-from-residue-field},
et que $a$ est un isomorphisme. (On utilise aussi Catégories dérivées des schémas,
lemmes \ref{perfect-lemma-affine-compare-bounded} et
\ref{perfect-lemma-perfect-affine}, pour la traduction en termes algébriques.)
\end{proof}

\noindent
Nous démontrons d'abord le résultat principal de cette section dans le
cas particulier suivant.

\begin{lemma}
\label{derham-lemma-global-generation-on-fibres}
Soit $f : X \to Y$ un morphisme lisse et propre de schémas sur une base $S$.
Supposons que
\begin{enumerate}
\item $Y$ et $S$ sont affines, et
\item il existe des entiers $N$ et $n_1, \ldots, n_N \geq 0$ et des classes
$\xi_i \in H^{n_i}_{dR}(X/S)$, $1 \leq i \leq N$, tels que,
pour tout point $y \in Y$, les images de $\xi_1, \ldots, \xi_N$
dans $H^*_{dR}(X_y/y)$ forment une base sur $\kappa(y)$.
\end{enumerate}
Alors l'application
$$
\bigoplus\nolimits_{i = 1}^N H^*_{dR}(Y/S) \longrightarrow
H^*_{dR}(X/S), \quad
(a_1, \ldots, a_N) \longmapsto  \sum \xi_i \cup f^*a_i
$$
est un isomorphisme.
\end{lemma}

\begin{proof}
Écrivons $Y = \Spec(A)$ et $S = \Spec(R)$.
Dans ce cas, $\Omega^\bullet_{A/R}$ calcule
$R\Gamma(Y, \Omega^\bullet_{Y/S})$ d'après le lemme \ref{derham-lemma-de-rham-affine}.
Choisissons un recouvrement ouvert affine fini $\mathcal{U} : X = \bigcup_{i \in I} U_i$.
Considérons le complexe
$$
K^\bullet =
\text{Tot}(\check{\mathcal{C}}^\bullet(\mathcal{U}, \Omega_{X/S}^\bullet))
$$
comme dans
Cohomologie, section \ref{cohomology-section-cech-cohomology-of-complexes}.
Rassemblons quelques propriétés de ce complexe, dont la plupart
figurent dans la référence que nous venons de donner :
\begin{enumerate}
\item $K^\bullet$ est un complexe de $R$-modules dont les termes sont des
$A$-modules,
\item $K^\bullet$ représente $R\Gamma(X, \Omega^\bullet_{X/S})$ dans $D(R)$
(Cohomologie des schémas, lemme
\ref{coherent-lemma-quasi-coherent-affine-cohomology-zero} et
Cohomologie, lemme \ref{cohomology-lemma-cech-complex-complex-computes}),
\item il existe un morphisme naturel $\Omega^\bullet_{A/R} \to K^\bullet$
de complexes de $R$-modules, $A$-linéaire sur les termes, qui
induit le morphisme d'image inverse $H^*_{dR}(Y/S) \to H^*_{dR}(X/S)$
en cohomologie,
\item $K^\bullet$ possède une multiplication notée $\wedge$
qui en fait une $R$-algèbre différentielle graduée,
\item la multiplication sur $K^\bullet$
induit le cup-produit sur $H^*_{dR}(X/S)$
(Cohomologie, section \ref{cohomology-section-cup-product}),
\item la filtration $F$ sur $\Omega^*_{X/S}$ induit une filtration
$$
K^\bullet =
F^0K^\bullet \supset F^1K^\bullet \supset F^2K^\bullet \supset \ldots
$$
par des sous-complexes de $K^\bullet$ telle que
\begin{enumerate}
\item $F^kK^n \subset K^n$ est un $A$-sous-module,
\item $F^kK^\bullet \wedge F^lK^\bullet \subset F^{k + l}K^\bullet$,
\item $\text{gr}^kK^\bullet$ est un complexe de $A$-modules,
\item $\text{gr}^0K^\bullet =
\text{Tot}(\check{\mathcal{C}}^\bullet(\mathcal{U}, \Omega_{X/Y}^\bullet))$
et représente $R\Gamma(X, \Omega^\bullet_{X/Y})$ dans $D(A)$,
\item la multiplication induit un isomorphisme
$\Omega^k_{A/R}[-k] \otimes_A \text{gr}^0K^\bullet \to \text{gr}^kK^\bullet$
\end{enumerate}
\end{enumerate}
Nous omettons les démonstrations détaillées de ces assertions ; voir la discussion
qui précède la construction de la suite spectrale du
lemme \ref{derham-lemma-spectral-sequence-smooth}.

\medskip\noindent
Pour tout $i = 1, \ldots, N$, choisissons un cocycle $x_i \in K^{n_i}$
représentant $\xi_i$. Considérons ensuite le morphisme de complexes
$$
\tilde x :
M^\bullet = \bigoplus\nolimits_{i = 1, \ldots, N}
\Omega^\bullet_{A/R}[-n_i]
\longrightarrow
K^\bullet
$$
qui envoie $\omega$ dans le facteur d'indice $i$ sur $x_i \wedge \omega$.
Il reste seulement à montrer que ce morphisme est un quasi-isomorphisme.
Munissons $M^\bullet$ d'une structure de complexe filtré
par la règle
$$
F^kM^\bullet =
\bigoplus\nolimits_{i = 1, \ldots, N}
(\sigma_{\geq k}\Omega^\bullet_{A/R})[-n_i]
$$
Avec ce choix, $\tilde x$ est un morphisme de complexes filtrés.
Observons que $\text{gr}^0M^\bullet = \bigoplus A[-n_i]$
et que la multiplication induit un isomorphisme
$\Omega^k_{A/R}[-k] \otimes_A \text{gr}^0M^\bullet \to \text{gr}^kM^\bullet$.
Par construction et d'après le lemme \ref{derham-lemma-relative-global-generation-on-fibres},
$$
\text{gr}^0\tilde x :
\text{gr}^0M^\bullet \longrightarrow
\text{gr}^0K^\bullet
$$
est un isomorphisme dans $D(A)$. Il en résulte que, pour tout $k \geq 0$,
on obtient des isomorphismes
$$
\text{gr}^k \tilde x :
\text{gr}^kM^\bullet = \Omega^k_{A/R}[-k] \otimes_A \text{gr}^0M^\bullet
\longrightarrow
\Omega^k_{A/R}[-k] \otimes_A \text{gr}^0K^\bullet =
\text{gr}^kK^\bullet
$$
dans $D(A)$. En effet, le complexe
$\text{gr}^0K^\bullet =
\text{Tot}(\check{\mathcal{C}}^\bullet(\mathcal{U}, \Omega_{X/Y}^\bullet))$
est K-plat comme complexe de $A$-modules, d'après Catégories dérivées des schémas,
lemme \ref{perfect-lemma-K-flat}.
Le produit tensoriel du membre de droite est donc le
produit tensoriel dérivé, ce qui est aussi vrai par inspection pour le membre de gauche.
Enfin, le produit tensoriel dérivé
$\Omega^k_{A/R}[-k] \otimes_A^\mathbf{L} -$ est un foncteur sur $D(A)$
et envoie donc les isomorphismes sur des isomorphismes.
Par récurrence sur $k$, on en déduit que
$$
\tilde x : M^\bullet/F^kM^\bullet \to K^\bullet/F^kK^\bullet
$$
est un isomorphisme dans $D(R)$, puisqu'on dispose des suites exactes courtes
$$
0 \to F^kM^\bullet/F^{k + 1}M^\bullet \to
M^\bullet/F^{k + 1}M^\bullet \to
\text{gr}^kM^\bullet \to 0
$$
et de même pour $K^\bullet$. Cela prouve que $\tilde x$ est un
quasi-isomorphisme, les filtrations étant finies en chaque degré.
\end{proof}

\begin{proposition}
\label{derham-proposition-global-generation-on-fibres}
Soit $f : X \to Y$ un morphisme lisse et propre de schémas sur une base $S$.
Soient $N$ et $n_1, \ldots, n_N \geq 0$ des entiers, et soient
$\xi_i \in H^{n_i}_{dR}(X/S)$, $1 \leq i \leq N$.
Supposons que, pour tout point $y \in Y$, les images de $\xi_1, \ldots, \xi_N$
dans $H^*_{dR}(X_y/y)$ forment une base sur $\kappa(y)$. Le morphisme
$$
\tilde \xi = \bigoplus \tilde \xi_i[-n_i] :
\bigoplus \Omega^\bullet_{Y/S}[-n_i]
\longrightarrow
Rf_*\Omega^\bullet_{X/S}
$$
(voir la démonstration) est un isomorphisme dans $D(Y, (Y \to S)^{-1}\mathcal{O}_S)$, et
l'application correspondante
$$
\bigoplus\nolimits_{i = 1}^N H^*_{dR}(Y/S) \longrightarrow
H^*_{dR}(X/S), \quad
(a_1, \ldots, a_N) \longmapsto  \sum \xi_i \cup f^*a_i
$$
est un isomorphisme.
\end{proposition}

\begin{proof}
Notons $p : X \to S$ et $q : Y \to S$ les morphismes structuraux.
Soit $\xi'_i : \Omega^\bullet_{X/S} \to \Omega^\bullet_{X/S}[n_i]$
le morphisme de la remarque \ref{derham-remark-cup-product-as-a-map} correspondant
à $\xi_i$. Notons
$$
\tilde \xi_i :
\Omega^\bullet_{Y/S} \to Rf_*\Omega^\bullet_{X/S}[n_i]
$$
la composée de $\xi'_i$ avec le morphisme canonique
$\Omega^\bullet_{Y/S} \to Rf_*\Omega^\bullet_{X/S}$.
En utilisant
$$
R\Gamma(Y, Rf_*\Omega^\bullet_{X/S}) = R\Gamma(X, \Omega^\bullet_{X/S})
$$,
le morphisme $\tilde \xi_i$ induit en cohomologie l'application $\eta \mapsto \xi_i \cup f^*\eta$
de $H^m_{dR}(Y/S)$ dans $H^{m + n}_{dR}(X/S)$.
De plus, puisque la formation de $\xi'_i$ commute aux
restrictions aux ouverts, il en va de même de celle de $\tilde \xi_i$.

\medskip\noindent
On peut donc considérer le morphisme
$$
\tilde \xi = \bigoplus \tilde \xi_i[-n_i] :
\bigoplus \Omega^\bullet_{Y/S}[-n_i]
\longrightarrow
Rf_*\Omega^\bullet_{X/S}
$$
Pour démontrer le lemme, il suffit de montrer que c'est un isomorphisme dans
$D(Y, q^{-1}\mathcal{O}_S)$. Si l'on pouvait montrer que $\tilde \xi$
provient d'un morphisme de complexes filtrés (pour des filtrations convenables),
on pourrait appliquer la suite spectrale du
lemme \ref{derham-lemma-spectral-sequence-smooth} pour conclure.
Cela demanderait plus de travail qu'il n'est nécessaire ; notre méthode
consistera plutôt à nous ramener au cas affine (dont la démonstration utilise,
en un certain sens, la suite spectrale).

\medskip\noindent
En effet, si $Y' \subset Y$ est un ouvert quelconque d'image inverse
$X' \subset X$, alors $\tilde \xi|_{X'}$ induit l'application
$$
\bigoplus\nolimits_{i = 1}^N H^*_{dR}(Y'/S) \longrightarrow
H^*_{dR}(X'/S), \quad
(a_1, \ldots, a_N) \longmapsto  \sum \xi_i|_{X'} \cup f^*a_i
$$
en cohomologie sur $Y'$ ; voir la discussion ci-dessus.
Il suffit donc de trouver une base de la topologie
de $Y$ telle que la proposition soit vraie pour les ouverts de cette base
(en particulier, on peut alors oublier le morphisme $\tilde \xi$).
On se ramène ainsi au cas où $Y$ et $S$
sont affines, traité par le lemme \ref{derham-lemma-global-generation-on-fibres},
ce qui achève la démonstration.
\end{proof}





\section{Formule du fibré projectif}
\label{derham-section-projective-space-bundle-formula}

\noindent
Le titre dit tout.

\begin{proposition}
\label{derham-proposition-projective-space-bundle-formula}
Soit $X \to S$ un morphisme de schémas. Soit $\mathcal{E}$ un
$\mathcal{O}_X$-module localement libre de rang constant $r$. Considérons le morphisme
$p : P = \mathbf{P}(\mathcal{E}) \to X$.
Alors l'application
$$
\bigoplus\nolimits_{i = 0, \ldots, r - 1} H^*_{dR}(X/S)
\longrightarrow
H^*_{dR}(P/S)
$$
définie par la règle
$$
(a_0, \ldots, a_{r - 1}) \longmapsto
\sum\nolimits_{i = 0, \ldots, r - 1} c_1^{dR}(\mathcal{O}_P(1))^i \cup p^*(a_i)
$$
est un isomorphisme.
\end{proposition}

\begin{proof}
Choisissons un ouvert affine $\Spec(A) \subset X$ sur lequel $\mathcal{E}$ se restreint
au module localement libre trivial $\mathcal{O}_{\Spec(A)}^{\oplus r}$.
Alors $P \times_X \Spec(A) = \mathbf{P}^{r - 1}_A$. On voit ainsi que
$p$ est propre et lisse ; voir la section \ref{derham-section-projective-space}.
De plus, les classes $c_1^{dR}(\mathcal{O}_P(1))^i$, $i = 0, 1, \ldots, r - 1$,
restreintes à une fibre $X_y = \mathbf{P}^{r - 1}_y$, engendrent librement la
cohomologie de de Rham $H^*_{dR}(X_y/y)$ sur $\kappa(y)$ ; voir
le lemme \ref{derham-lemma-de-rham-cohomology-projective-space}. Nous avons donc vérifié les
hypothèses de la proposition \ref{derham-proposition-global-generation-on-fibres},
d'où le résultat.
\end{proof}

\begin{remark}
\label{derham-remark-projective-space-bundle-formula}
Dans la situation de la
proposition \ref{derham-proposition-projective-space-bundle-formula},
on obtient en outre que le morphisme
$$
\tilde \xi :
\bigoplus\nolimits_{t = 0, \ldots, r - 1}
\Omega^\bullet_{X/S}[-2t]
\longrightarrow
Rp_*\Omega^\bullet_{P/S}
$$
est un isomorphisme dans $D(X, (X \to S)^{-1}\mathcal{O}_X)$, ce qui résulte
immédiatement de l'application de la
proposition \ref{derham-proposition-global-generation-on-fibres}.
Notons que la flèche pour $t = 0$ est simplement le morphisme canonique
$c_{P/X} : \Omega^\bullet_{X/S} \to Rp_*\Omega^\bullet_{P/S}$
de la section \ref{derham-section-de-rham-complex}.
On peut en fait préciser davantage ce morphisme dans ce cas particulier.
Considérons le morphisme canonique
$$
\xi' : \Omega^\bullet_{P/S} \to \Omega^\bullet_{P/S}[2]
$$
de la remarque \ref{derham-remark-cup-product-as-a-map} correspondant à
$c_1^{dR}(\mathcal{O}_P(1))$. Alors
$$
\xi'[2(t - 1)] \circ \ldots \circ \xi'[2] \circ \xi' : 
\Omega^\bullet_{P/S} \to \Omega^\bullet_{P/S}[2t]
$$
est le morphisme de la remarque \ref{derham-remark-cup-product-as-a-map} correspondant à
$c_1^{dR}(\mathcal{O}_P(1))^t$. En suivant les choix effectués dans la
démonstration de la proposition \ref{derham-proposition-global-generation-on-fibres},
on obtient la valeur
$$
\tilde \xi|_{\Omega^\bullet_{X/S}[-2t]} =
Rp_*\xi'[-2] \circ \ldots \circ Rp_*\xi'[-2(t - 1)] \circ
Rp_*\xi'[-2t] \circ c_{P/X}[-2t]
$$
pour la restriction de notre isomorphisme au facteur direct
$\Omega^\bullet_{X/S}[-2t]$. Cela a la conséquence simple
suivante, que nous utiliserons plus bas : posons
$$
M = \bigoplus\nolimits_{t = 1, \ldots, r - 1} \Omega^\bullet_{X/S}[-2t]
\quad\text{et}\quad
K = \bigoplus\nolimits_{t = 0, \ldots, r - 2} \Omega^\bullet_{X/S}[-2t]
$$,
considérés comme des sous-complexes de la source de la flèche $\tilde \xi$.
Il résulte formellement de la discussion ci-dessus que
$$
c_{P/X} \oplus
\tilde \xi|_M :
\Omega^\bullet_{X/S} \oplus M \longrightarrow
Rp_*\Omega^\bullet_{P/S}
$$
est un isomorphisme et que le diagramme
$$
\xymatrix{
K \ar[d]_{\tilde \xi|_K} \ar[r]_{\text{id}} &
M[2] \ar[d]^{(\tilde \xi|_M)[2]} \\
Rp_*\Omega^\bullet_{P/S} \ar[r]^{Rp_*\xi'} &
Rp_*\Omega^\bullet_{P/S}[2]
}
$$
est commutatif, où $\text{id} : K \to M[2]$ identifie le facteur
correspondant à $t$ dans la décomposition de $K$ au facteur
correspondant à $t + 1$ dans la décomposition de $M$.
\end{remark}









\section{Pôles logarithmiques le long d'un diviseur}
\label{derham-section-divisor}

\noindent
Soit $X \to S$ un morphisme de schémas. Soit $Y \subset X$ un
diviseur de Cartier effectif. Si $X$ est, étale-localement le long de $Y$,
de la forme $Y \times \mathbf{A}^1$, il existe une suite exacte courte canonique
de complexes
$$
0 \to \Omega^\bullet_{X/S} \to
\Omega^\bullet_{X/S}(\log Y) \to
\Omega^\bullet_{Y/S}[-1] \to 0
$$
possédant de nombreuses propriétés utiles, que nous étudierons dans cette section. Il existe une
variante de cette construction dans laquelle on part d'un diviseur à croisements
normaux
(Morphismes étales, définition \ref{etale-definition-strict-normal-crossings}),
que nous étudierons ailleurs (insérer ici une référence ultérieure).

\begin{definition}
\label{derham-definition-local-product}
Soit $X \to S$ un morphisme de schémas. Soit $Y \subset X$ un
diviseur de Cartier effectif. On dit que le
{\it complexe de de Rham à pôles logarithmiques est défini pour $Y \subset X$ sur $S$}
si, pour tout $y \in Y$ et toute équation locale $f \in \mathcal{O}_{X, y}$
de $Y$, on a
\begin{enumerate}
\item $\mathcal{O}_{X, y} \to \Omega_{X/S, y}$, $g \mapsto g \text{d}f$,
est une injection scindée, et
\item $\Omega^p_{X/S, y}$ est sans $f$-torsion pour tout $p$.
\end{enumerate}
\end{definition}

\noindent
Un calcul local facile montre qu'il suffit, pour tout $y \in Y$,
de trouver une équation locale $f$ pour laquelle les conditions (1) et (2) sont satisfaites.

\begin{lemma}
\label{derham-lemma-log-complex}
Soit $X \to S$ un morphisme de schémas. Soit $Y \subset X$ un
diviseur de Cartier effectif.
Supposons que le complexe de de Rham à pôles logarithmiques soit défini pour $Y \subset X$ sur $S$.
Il existe une suite exacte courte canonique
de complexes
$$
0 \to \Omega^\bullet_{X/S} \to
\Omega^\bullet_{X/S}(\log Y) \to
\Omega^\bullet_{Y/S}[-1] \to 0
$$
\end{lemma}

\begin{proof}
Notre hypothèse est que, pour tout $y \in Y$ et toute équation locale
$f \in \mathcal{O}_{X, y}$ de $Y$, on a
$$
\Omega_{X/S, y} = \mathcal{O}_{X, y}\text{d}f \oplus M
\quad\text{et}\quad
\Omega^p_{X/S, y} = \wedge^{p - 1}(M)\text{d}f \oplus \wedge^p(M)
$$
pour un certain module $M$ dont les puissances extérieures sont sans $f$-torsion, à savoir $\wedge^p(M)$.
Il en résulte que
$$
\Omega^p_{Y/S, y} = \wedge^p(M/fM) = \wedge^p(M)/f\wedge^p(M)
$$
Nous utiliserons tacitement ces faits dans la suite.
En particulier, les faisceaux $\Omega^p_{X/S}$ n'ont pas de sections locales non nulles
à support dans $Y$, et l'on dispose d'une inclusion canonique
$$
\Omega^p_{X/S} \subset \Omega^p_{X/S}(Y)
$$ ;
voir Compléments de platitude, section \ref{flat-section-eta}. Soit $U = \Spec(A)$
un sous-schéma ouvert affine tel que $Y \cap U = V(f)$ pour un certain
élément non diviseur de zéro $f \in A$. Considérons le $\mathcal{O}_U$-sous-module
de $\Omega^p_{X/S}(Y)|_U$ engendré par
$\Omega^p_{X/S}|_U$ et $\text{d}\log(f) \wedge \Omega^{p - 1}_{X/S}$,
où $\text{d}\log(f) = f^{-1}\text{d}(f)$.
Il est indépendant du choix de $f$, car un autre générateur de
l'idéal de $Y$ sur $U$ est égal à $uf$ pour une unité $u \in A$, et l'on obtient
$$
\text{d}\log(uf) - \text{d}\log(f) = \text{d}\log(u) = u^{-1}\text{d}u
$$,
qui est une section de $\Omega_{X/S}$ sur $U$. Ces faisceaux locaux
se recollent en un sous-module quasi-cohérent
$$
\Omega^p_{X/S} \subset \Omega^p_{X/S}(\log Y) \subset \Omega^p_{X/S}(Y)
$$
Convenons de considérer $\Omega^p_{Y/S}$ comme un
$\mathcal{O}_X$-module quasi-cohérent. Il existe un unique morphisme surjectif
$\mathcal{O}_X$-linéaire
$$
\text{Res} : \Omega^p_{X/S}(\log Y) \to \Omega^{p - 1}_{Y/S}
$$
défini par la règle
$$
\text{Res}(\eta' + \text{d}\log(f) \wedge \eta) = \eta|_{Y \cap U}
$$
pour tout ouvert $U$ comme ci-dessus et tous
$\eta' \in \Omega^p_{X/S}(U)$ et $\eta \in \Omega^{p - 1}_{X/S}(U)$.
Si une forme $\eta$ sur $U$ se restreint à zéro sur $Y \cap U$, alors
$\eta = \text{d}f \wedge \eta' + f\eta''$ pour certaines formes $\eta'$ et $\eta''$
sur $U$. On en déduit
une suite exacte courte
$$
0 \to \Omega^p_{X/S} \to \Omega^p_{X/S}(\log Y) \to \Omega^{p - 1}_{Y/S} \to 0
$$
pour tout $p$. Il reste à définir les différentielles
$\Omega^p_{X/S}(\log Y) \to \Omega^{p + 1}_{X/S}(\log Y)$.
Sur le sous-faisceau $\Omega^p_{X/S}$, on utilise la différentielle du
complexe de de Rham de $X$ sur $S$. Enfin, on pose
$\text{d}(\text{d}\log(f) \wedge \eta) = -\text{d}\log(f) \wedge \text{d}\eta$.
Le signe est imposé par la règle de Leibniz (sur $\Omega^\bullet_{X/S}$)
et est compatible avec la différentielle sur $\Omega^\bullet_{Y/S}[-1]$,
qui est bien $-\text{d}_{Y/S}$ selon notre convention de signe dans
Homologie, définition \ref{homology-definition-shift-cochain}.
On obtient ainsi une suite exacte courte
de complexes comme dans l'énoncé du lemme.
\end{proof}

\begin{definition}
\label{derham-definition-log-complex}
Soit $X \to S$ un morphisme de schémas. Soit $Y \subset X$ un
diviseur de Cartier effectif. Supposons que le complexe de de Rham à pôles logarithmiques
soit défini pour $Y \subset X$ sur $S$. Alors le complexe
$$
\Omega^\bullet_{X/S}(\log Y)
$$
construit dans le lemme \ref{derham-lemma-log-complex} est le
{\it complexe de de Rham à pôles logarithmiques pour $Y \subset X$ sur $S$}.
\end{definition}

\noindent
Ce complexe possède de nombreuses propriétés utiles.

\begin{lemma}
\label{derham-lemma-multiplication-log}
Soit $p : X \to S$ un morphisme de schémas. Soit $Y \subset X$ un
diviseur de Cartier effectif. Supposons que le complexe de de Rham à pôles logarithmiques
soit défini pour $Y \subset X$ sur $S$.
\begin{enumerate}
\item Les applications
$\wedge : \Omega^p_{X/S} \times \Omega^q_{X/S} \to \Omega^{p + q}_{X/S}$
s'étendent de manière unique en des applications $\mathcal{O}_X$-bilinéaires
$$
\wedge : \Omega^p_{X/S}(\log Y) \times \Omega^q_{X/S}(\log Y)
\to \Omega^{p + q}_{X/S}(\log Y)
$$
satisfaisant à la règle de Leibniz
$
\text{d}(\omega \wedge \eta) = \text{d}(\omega) \wedge \eta +
(-1)^{\deg(\omega)} \omega \wedge \text{d}(\eta)$,
\item pour la multiplication de (1), le morphisme
$\Omega^\bullet_{X/S} \to \Omega^\bullet_{X/S}(\log(Y)$
est un homomorphisme de $\mathcal{O}_S$-algèbres différentielles graduées,
\item via les applications de (1), on a $\Omega^p_{X/S}(\log Y) =
\wedge^p(\Omega^1_{X/S}(\log Y))$, et
\item le morphisme
$\text{Res} : \Omega^\bullet_{X/S}(\log Y) \to \Omega^\bullet_{Y/S}[-1]$
vérifie
$$
\text{Res}(\omega \wedge \eta) = \text{Res}(\omega) \wedge \eta|_Y
$$
pour toute section locale $\omega$ de $\Omega^p_{X/S}(\log Y)$ et toute section locale $\eta$
de $\Omega^q_{X/S}$.
\end{enumerate}
\end{lemma}

\begin{proof}
Cela résulte par un calcul direct de la construction locale
du complexe dans la démonstration du lemme \ref{derham-lemma-log-complex}.
Nous omettons les détails.
\end{proof}

\noindent
Considérons un diagramme commutatif
$$
\xymatrix{
X' \ar[r]_f \ar[d] & X \ar[d] \\
S' \ar[r] & S
}
$$
de schémas. Soit $Y \subset X$ un diviseur de Cartier effectif
dont l'image inverse $Y' = f^*Y$ est définie
(Diviseurs, définition
\ref{divisors-definition-pullback-effective-Cartier-divisor}).
Supposons que
le complexe de de Rham à pôles logarithmiques soit défini pour $Y \subset X$ sur $S$
et que
le complexe de de Rham à pôles logarithmiques soit défini pour $Y' \subset X'$ sur $S'$.
On obtient alors un morphisme de suites exactes courtes de complexes
$$
\xymatrix{
0 \ar[r] &
f^{-1}\Omega^\bullet_{X/S} \ar[r] \ar[d] &
f^{-1}\Omega^\bullet_{X/S}(\log Y) \ar[r] \ar[d] &
f^{-1}\Omega^\bullet_{Y/S}[-1] \ar[r] \ar[d] &
0 \\
0 \ar[r] &
\Omega^\bullet_{X'/S'} \ar[r] &
\Omega^\bullet_{X'/S'}(\log Y') \ar[r] &
\Omega^\bullet_{Y'/S'}[-1] \ar[r] &
0
}
$$
Par linéarisation, on obtient, pour tout $p$, un morphisme linéaire
$f^*\Omega^p_{X/S}(\log Y) \to \Omega^p_{X'/S'}(\log Y')$.

\begin{lemma}
\label{derham-lemma-gysin-via-log-complex}
Soit $f : X \to S$ un morphisme de schémas. Soit $Y \subset X$ un diviseur
de Cartier effectif. Supposons que le complexe de de Rham à pôles logarithmiques soit défini pour
$Y \subset X$ sur $S$. Notons
$$
\delta : \Omega^\bullet_{Y/S} \to \Omega^\bullet_{X/S}[2]
$$
dans $D(X, f^{-1}\mathcal{O}_S)$ le morphisme de « bord » provenant de la
suite exacte courte du lemme \ref{derham-lemma-log-complex}. Notons
$$
\xi' : \Omega^\bullet_{X/S} \to \Omega^\bullet_{X/S}[2]
$$
dans $D(X, f^{-1}\mathcal{O}_S)$ le morphisme de la
remarque \ref{derham-remark-cup-product-as-a-map}
correspondant à $\xi = c_1^{dR}(\mathcal{O}_X(-Y))$. Notons
$$
\zeta' : \Omega^\bullet_{Y/S} \to \Omega^\bullet_{Y/S}[2]
$$
dans $D(Y, f|_Y^{-1}\mathcal{O}_S)$ le morphisme de la
remarque \ref{derham-remark-cup-product-as-a-map} correspondant à
$\zeta = c_1^{dR}(\mathcal{O}_X(-Y)|_Y)$. Alors le diagramme
$$
\xymatrix{
\Omega^\bullet_{X/S} \ar[d]_{\xi'} \ar[r] &
\Omega^\bullet_{Y/S} \ar[d]^{\zeta'} \ar[ld]_\delta \\
\Omega^\bullet_{X/S}[2] \ar[r] &
\Omega^\bullet_{Y/S}[2]
}
$$
est commutatif dans $D(X, f^{-1}\mathcal{O}_S)$.
\end{lemma}

\begin{proof}
Plus précisément, on définit $\delta$ comme le morphisme de bord correspondant à la
suite exacte courte décalée
$$
0 \to \Omega^\bullet_{X/S}[1] \to
\Omega^\bullet_{X/S}(\log Y)[1] \to
\Omega^\bullet_{Y/S} \to 0
$$
Il suffit de démontrer la commutativité de chaque triangle. Posons
$\mathcal{L} = \mathcal{O}_X(-Y)$. Notons $\pi : L \to X$ le fibré en droites
tel que $\pi_*\mathcal{O}_L = \bigoplus_{n \geq 0} \mathcal{L}^{\otimes n}$.

\medskip\noindent
Commutativité du triangle supérieur gauche.
D'après le lemme \ref{derham-lemma-the-complex-for-L-star-gives-chern-class},
le morphisme $\xi'$ est le morphisme de bord du triangle donné dans le
lemme \ref{derham-lemma-the-complex-for-L-star}.
Par fonctorialité, il suffit de prouver qu'il existe un morphisme de
suites exactes courtes
$$
\xymatrix{
0 \ar[r] &
\Omega^\bullet_{X/S}[1] \ar[r] \ar[d] &
\Omega^\bullet_{L^\star/S, 0}[1] \ar[r] \ar[d] &
\Omega^\bullet_{X/S} \ar[r] \ar[d] &
0 \\
0 \ar[r] &
\Omega^\bullet_{X/S}[1] \ar[r] &
\Omega^\bullet_{X/S}(\log Y)[1] \ar[r] &
\Omega^\bullet_{Y/S} \ar[r] &
0
}
$$
où les flèches verticales de gauche et de droite sont les flèches évidentes.
On peut définir la flèche verticale du milieu par la règle
$$
\omega' + \text{d}\log(s) \wedge \omega \longmapsto
\omega' + \text{d}\log(f) \wedge \omega
$$
où $\omega', \omega$ sont des sections locales de $\Omega^\bullet_{X/S}$,
où $s$ est un générateur local de $\mathcal{L}$ et
où $f \in \mathcal{O}_X(-Y)$ est la section correspondante du faisceau
d'idéaux de $Y$ dans $X$. Puisque les constructions des morphismes dans les
lemmes \ref{derham-lemma-the-complex-for-L-star} et \ref{derham-lemma-log-complex}
se correspondent exactement, cela convient.

\medskip\noindent
Commutativité du triangle inférieur droit. Notons
$\overline{L}$ la restriction de $L$ à $Y$.
D'après le lemme \ref{derham-lemma-the-complex-for-L-star-gives-chern-class},
le morphisme $\zeta'$ est le morphisme de bord du triangle donné dans le
lemme \ref{derham-lemma-the-complex-for-L-star}, appliqué au fibré en droites
$\overline{L}$ sur $Y$.
Par fonctorialité, il suffit de prouver qu'il existe un morphisme de
suites exactes courtes
$$
\xymatrix{
0 \ar[r] &
\Omega^\bullet_{X/S}[1] \ar[r] \ar[d] &
\Omega^\bullet_{X/S}(\log Y)[1] \ar[r] \ar[d] &
\Omega^\bullet_{Y/S} \ar[r] \ar[d] &
0 \\
0 \ar[r] &
\Omega^\bullet_{Y/S}[1] \ar[r] &
\Omega^\bullet_{\overline{L}^\star/S, 0}[1] \ar[r] &
\Omega^\bullet_{Y/S} \ar[r] &
0 \\
}
$$
où les flèches verticales de gauche et de droite sont les flèches évidentes.
On peut définir la flèche verticale du milieu par la règle
$$
\omega' + \text{d}\log(f) \wedge \omega \longmapsto
\omega'|_Y + \text{d}\log(\overline{s}) \wedge \omega|_Y
$$
où $\omega', \omega$ sont des sections locales de $\Omega^\bullet_{X/S}$,
où $f$ est un générateur local de $\mathcal{O}_X(-Y)$ considéré comme
une fonction sur $X$, et où $\overline{s}$ désigne $f|_Y$ considéré comme une
section de $\mathcal{L}|_Y = \mathcal{O}_X(-Y)|_Y$.
Puisque les constructions des morphismes dans les
lemmes \ref{derham-lemma-the-complex-for-L-star} et \ref{derham-lemma-log-complex}
se correspondent exactement, cela convient.
\end{proof}

\begin{lemma}
\label{derham-lemma-log-complex-consequence}
Soit $X \to S$ un morphisme de schémas. Soit $Y \subset X$ un diviseur
de Cartier effectif. Supposons que le complexe de de Rham à pôles logarithmiques soit défini pour
$Y \subset X$ sur $S$. Soit $b \in H^m_{dR}(X/S)$ une classe de cohomologie de de Rham
dont la restriction à $Y$ est nulle. Alors
$c_1^{dR}(\mathcal{O}_X(Y)) \cup b = 0$ dans $H^{m + 2}_{dR}(X/S)$.
\end{lemma}

\begin{proof}
Cela résulte immédiatement du lemme \ref{derham-lemma-gysin-via-log-complex}.
En effet, on a
$$
c_1^{dR}(\mathcal{O}_X(Y)) \cup b =
-c_1^{dR}(\mathcal{O}_X(-Y)) \cup b = -\xi'(b) = -\delta(b|_Y) = 0
$$
comme voulu. Pour la deuxième égalité, voir
la remarque \ref{derham-remark-cup-product-as-a-map}.
\end{proof}

\begin{lemma}
\label{derham-lemma-check-log-smooth}
Soient $X \to T \to S$ des morphismes de schémas. Soit $Y \subset X$ un diviseur
de Cartier effectif. Si $X \to T$ et $Y \to T$ sont tous deux lisses, alors
le complexe de de Rham à pôles logarithmiques est défini pour $Y \subset X$ sur $S$.
\end{lemma}

\begin{proof}
Soit $y \in Y$ un point.
D'après Compléments sur les morphismes, lemme \ref{more-morphisms-lemma-etale-local-structure},
il existe un entier $0 \geq m$ et un diagramme commutatif
$$
\xymatrix{
Y \ar[d] &
V \ar[l] \ar[d] \ar[r] &
\mathbf{A}^m_T
\ar[d]^{(a_1, \ldots, a_m) \mapsto (a_1, \ldots, a_m, 0)} \\
X &
U \ar[l] \ar[r]^-\pi &
\mathbf{A}^{m + 1}_T
}
$$
où $U \subset X$ est ouvert, $V = Y \cap U$,
$\pi$ est étale, $V = \pi^{-1}(\mathbf{A}^m_T)$ et $y \in V$.
Notons $z \in \mathbf{A}^m_T$ l'image de $y$. On a alors
$$
\Omega^p_{X/S, y} = \Omega^p_{\mathbf{A}^{m + 1}_T/S, z}
\otimes_{\mathcal{O}_{\mathbf{A}^{m + 1}_T, z}} \mathcal{O}_{X, x}
$$
d'après le lemme \ref{derham-lemma-etale}. Notons $x_1, \ldots, x_{m + 1}$
les fonctions coordonnées sur $\mathbf{A}^{m + 1}_T$.
Puisque les conditions (1) et (2) de la définition \ref{derham-definition-local-product}
ne dépendent pas du choix de la coordonnée locale,
il suffit de les vérifier lorsque $f$ est
l'image de $x_{m + 1}$ par l'homomorphisme local plat d'anneaux
$\mathcal{O}_{\mathbf{A}^{m + 1}_T, z} \to \mathcal{O}_{X, x}$.
Il suffit donc de vérifier les conditions (1) et (2)
pour $\mathbf{A}^m_T \subset \mathbf{A}^{m + 1}_T$ et le point $z$.
Pour traiter ce cas, on peut supposer $S = \Spec(A)$ et $T = \Spec(B)$
affines. Soit $A \to B$ le morphisme d'anneaux correspondant au morphisme
$T \to S$, et posons $P = B[x_1, \ldots, x_{m + 1}]$, de sorte que
$\mathbf{A}^{m + 1}_T = \Spec(P)$. On a
$$
\Omega_{P/A} = \Omega_{B/A} \otimes_B P \oplus
\bigoplus\nolimits_{j = 1, \ldots, m} P \text{d}x_j \oplus
P \text{d}x_{m + 1}
$$
Ainsi, l'application $P \to \Omega_{P/A}$, $g \mapsto g \text{d}x_{m + 1}$,
est une injection scindée, et $x_{m + 1}$ est non diviseur de zéro sur
$\Omega^p_{P/A}$ pour tout $p \geq 0$. La localisation en l'idéal premier
correspondant à $z$ achève la démonstration.
\end{proof}

\begin{remark}
\label{derham-remark-check-log-completion-1}
Soit $S$ un schéma localement noethérien. Soit $X$ localement de type
fini sur $S$. Soit $Y \subset X$ un diviseur de Cartier effectif.
Si le morphisme
$$
\mathcal{O}_{X, y}^\wedge \longrightarrow \mathcal{O}_{Y, y}^\wedge
$$
admet une section pour tout $y \in Y$, alors
le complexe de de Rham à pôles logarithmiques est défini pour $Y \subset X$ sur $S$.
Si nous avons besoin de ce résultat, nous en formulerons un énoncé précis et
en ajouterons ici une démonstration.
\end{remark}

\begin{remark}
\label{derham-remark-check-log-completion-2}
Soit $S$ un schéma localement noethérien. Soit $X$ localement de type
fini sur $S$. Soit $Y \subset X$ un diviseur de Cartier effectif.
Si, pour tout $y \in Y$, on peut trouver un diagramme de schémas sur $S$
$$
X \xleftarrow{\varphi} U \xrightarrow{\psi} V
$$
où $\varphi$ est étale et $\psi|_{\varphi^{-1}(Y)} : \varphi^{-1}(Y) \to V$
est étale, alors le complexe de de Rham à pôles logarithmiques est défini pour
$Y \subset X$ sur $S$. C'est notamment le cas lorsque le couple $(X, Y)$
est, étale-localement, de la forme $(V \times \mathbf{A}^1, V \times \{0\})$.
Si nous avons besoin de ce résultat, nous en formulerons
un énoncé précis et en ajouterons ici une démonstration.
\end{remark}

















\section{Calculs}
\label{derham-section-calculations}

\noindent
Dans cette section, nous calculons certains groupes de cohomologie de Hodge
et de de Rham pour un éclatement standard.

\medskip\noindent
Fixons un anneau $R$ et posons $S = \Spec(R)$. Fixons des entiers $0 \leq m$ et
$1 \leq n$. Considérons l'immersion fermée
$$
Z = \mathbf{A}^m_S \longrightarrow \mathbf{A}^{m + n}_S = X,\quad
(a_1, \ldots, a_m) \mapsto (a_1, \ldots, a_m, 0, \ldots 0).
$$
Nous allons considérer l'éclatement $L$ de $X$
le long du sous-schéma fermé $Z$. Écrivons
$$
X =
\mathbf{A}^{m + n}_S =
\Spec(A)
\quad\text{avec}\quad
A = R[x_1, \ldots, x_m, y_1, \ldots, y_n]
$$
Nous considérerons $A = R[x_1, \ldots, x_m, y_1, \ldots, y_n]$ comme une
$R$-algèbre graduée en posant $\deg(x_i) = 0$ et $\deg(y_j) = 1$.
Pour cette graduation, on a
$$
P =
\text{Proj}(A) =
\mathbf{A}^m_S \times_S \mathbf{P}^{n - 1}_S =
Z \times_S \mathbf{P}^{n - 1}_S =
\mathbf{P}^{n - 1}_Z
$$
Observons que l'idéal définissant $Z$ dans $X$ est $A_+$.
Ainsi, $L$ est le Proj de l'algèbre de Rees
$$
A \oplus A_+ \oplus (A_+)^2 \oplus \ldots =
\bigoplus\nolimits_{d \geq 0} A_{\geq d}
$$
Par conséquent, $L$ est un exemple du phénomène étudié
plus généralement dans Compléments sur les morphismes, section
\ref{more-morphisms-section-proj-spec} ;
nous utiliserons les observations qui y ont été faites sans plus les mentionner.
En particulier, on dispose d'un diagramme commutatif
$$
\xymatrix{
P \ar[r]_0 \ar[d]_p &
L \ar[r]_-\pi \ar[d]^b &
P \ar[d]^p \\
Z \ar[r]^i &
X \ar[r] &
Z
}
$$
tel que $\pi : L \to P$ soit un fibré en droites sur
$P = Z \times_S \mathbf{P}^{n - 1}_S$,
de section nulle $0$, dont l'image $E = 0(P) \subset L$
est le diviseur exceptionnel de l'éclatement $b$.

\begin{lemma}
\label{derham-lemma-comparison}
Pour $a \geq 0$, on a
\begin{enumerate}
\item le morphisme
$\Omega^a_{X/S} \to b_*\Omega^a_{L/S}$ est un isomorphisme,
\item le morphisme $\Omega^a_{Z/S} \to p_*\Omega^a_{P/S}$ est un isomorphisme,
et
\item le morphisme $Rb_*\Omega^a_{L/S} \to i_*Rp_*\Omega^a_{P/S}$ induit un isomorphisme
sur les faisceaux de cohomologie en degrés $\geq 1$.
\end{enumerate}
\end{lemma}

\begin{proof}
Démontrons d'abord (2). Puisque
$P = Z \times_S \mathbf{P}^{n - 1}_S$,
on voit que
$$
\Omega^a_{P/S} = \bigoplus\nolimits_{a = r + s}
\text{pr}_1^*\Omega^r_{Z/S} \otimes
\text{pr}_2^*\Omega^s_{\mathbf{P}^{n - 1}_S/S}
$$
En rappelant que $p = \text{pr}_1$, la formule de projection
(Cohomologie, lemme \ref{cohomology-lemma-projection-formula})
donne
$$
p_*\Omega^a_{P/S} = \bigoplus\nolimits_{a = r + s}
\Omega^r_{Z/S} \otimes
\text{pr}_{1, *}\text{pr}_2^*\Omega^s_{\mathbf{P}^{n - 1}_S/S}
$$
D'après les calculs de la section \ref{derham-section-projective-space},
en particulier ceux de la
démonstration du lemme \ref{derham-lemma-hodge-cohomology-projective-space},
on a $\text{pr}_{1, *}\text{pr}_2^*\Omega^s_{\mathbf{P}^{n - 1}_S/S} = 0$,
sauf si $s = 0$, auquel cas on obtient
$\text{pr}_{1, *}\mathcal{O}_P = \mathcal{O}_Z$.
Cela démontre (2).

\medskip\noindent
D'après les résultats de la section \ref{derham-section-line-bundle}, en particulier
le lemme \ref{derham-lemma-push-omega-a}, on a
$\pi_*\Omega^a_{L/S} = \Omega^a_{P/S} \oplus
\bigoplus_{k \geq 1} \Omega^a_{L/S, k}$.
Puisque la composée $\pi \circ 0$ dans le diagramme ci-dessus
est l'identité de $P$, il suffit, pour démontrer (3), de montrer que
les groupes de cohomologie supérieure de $\Omega^a_{L/S, k}$ sont nuls pour $k > 0$.
D'après les lemmes \ref{derham-lemma-the-complex-for-L-star} et \ref{derham-lemma-push-omega-a},
on dispose de suites exactes courtes
$$
0 \to \Omega^a_{P/S} \otimes \mathcal{O}_P(k)
\to \Omega^a_{L/S, k} \to
\Omega^{a - 1}_{P/S} \otimes \mathcal{O}_P(k) \to 0
$$,
où $\Omega^{a - 1}_{P/S} = 0$ si $a = 0$. Puisque
$P = Z \times_S \mathbf{P}^{n - 1}_S$, on a
$$
\Omega^a_{P/S} = \bigoplus\nolimits_{i + j = a}
\Omega^i_{Z/S} \boxtimes \Omega^j_{\mathbf{P}^{n - 1}_S/S}
$$
d'après le lemme \ref{derham-lemma-de-rham-complex-product}.
Comme $\Omega^i_{Z/S}$ est libre de rang fini,
on voit qu'il suffit de montrer que les groupes de cohomologie supérieure de
$\mathcal{O}_Z \boxtimes \Omega^j_{\mathbf{P}^{n - 1}_S/S}(k)$
sont nuls pour $k > 0$. Cela résulte du
lemme \ref{derham-lemma-twisted-hodge-cohomology-projective-space}
appliqué à $P = Z \times_S \mathbf{P}^{n - 1}_S = \mathbf{P}^{n - 1}_Z$,
ce qui achève la démonstration de (3).

\medskip\noindent
Il reste à démontrer (1). Si $n = 1$, on éclate
un diviseur de Cartier effectif ; $b$ est donc un isomorphisme,
d'où (1). Si $n > 1$, la composée
$$
\Gamma(X, \Omega^a_{X/S})
\to
\Gamma(L, \Omega^a_{L/S})
\to
\Gamma(L \setminus E, \Omega^a_{L/S})
=
\Gamma(X \setminus Z, \Omega^a_{X/S})
$$
est un isomorphisme, puisque $\Omega^a_{X/S}$ est libre de rang fini
(nous omettons ce détail). Ainsi, (1) ne peut être faux que s'il
existe des éléments non nuls de $\Gamma(L, \Omega^a_{L/S})$ qui s'annulent
en dehors de $E = 0(P)$. Puisque $L$ est un fibré en droites sur $P$
de section nulle $0 : P \to L$, il suffit de montrer que,
sur un fibré en droites, il n'existe pas de section non nulle d'un faisceau
de différentielles qui s'annule identiquement hors de la section nulle.
Le lecteur le vérifiera soit, de préférence, par un calcul local,
soit en utilisant $\Omega_{L/S, k} \subset \Omega_{L^\star/S, k}$
(voir les références ci-dessus).
\end{proof}

\noindent
Nous suggérons au lecteur de passer maintenant à la section suivante.

\begin{lemma}
\label{derham-lemma-comparison-bis}
Pour $a \geq 0$, il existe des morphismes canoniques
$$
b^*\Omega^a_{X/S} \longrightarrow
\Omega^a_{L/S} \longrightarrow
b^*\Omega^a_{X/S} \otimes_{\mathcal{O}_L} \mathcal{O}_L((n - 1)E)
$$
dont la composée est induite par l'inclusion
$\mathcal{O}_L \subset \mathcal{O}_L((n - 1)E)$.
\end{lemma}

\begin{proof}
La première flèche de la formule affichée est
décrite à la section \ref{derham-section-de-rham-complex}.
Pour obtenir la deuxième flèche, il faut montrer que, si l'on considère
une section locale de $\Omega^a_{L/S}$ comme une « section méromorphe »
de $b^*\Omega^a_{X/S}$, elle a un pôle d'ordre au plus
$n - 1$ le long de $E$.
Pour le voir, travaillons dans des cartes affines locales
de $L$. Rappelons que $L$ est recouvert par les spectres des
algèbres d'éclatement affines $A[\frac{I}{y_i}]$, où $I = A_{+}$
est l'idéal engendré par $y_1, \ldots, y_n$. Voir
Algèbre, section \ref{algebra-section-blow-up} et
Diviseurs, lemme \ref{divisors-lemma-blowing-up-affine}.
Par symétrie, il suffit de travailler dans la
carte correspondant à $i = 1$. Alors
$$
A[\frac{I}{y_1}] = R[x_1, \ldots, x_m, y_1, t_2, \ldots, t_n]
$$
où $t_i = y_i/y_1$ ; voir
Compléments d'algèbre, lemme \ref{more-algebra-lemma-blowup-regular-sequence}.
Ainsi, le module $\Omega^1_{L/S}$ est librement engendré, sur l'ouvert
affine correspondant, par
$\text{d}x_1, \ldots, \text{d}x_m$, $\text{d}y_1$ et
$\text{d}t_1, \ldots, \text{d}t_n$.
Bien entendu, les $m + 1$ premiers de ces générateurs proviennent de
$b^*\Omega^1_{X/S}$, et pour les $n - 1$ restants, on a
$$
\text{d}t_j =
\text{d}\frac{y_j}{y_1} =
\frac{1}{y_1}\text{d}y_j - \frac{y_j}{y_1^2}\text{d}y_1 =
\frac{\text{d}y_j - t_j \text{d}y_1}{y_1}
$$,
qui a un pôle d'ordre $1$ le long de $E$, puisque $E$ est défini par $y_1$
sur cette carte. Les produits extérieurs de $a$ de ces éléments forment une base
de $\Omega^a_{L/S}$ sur cette carte, et au plus
$n - 1$ des $\text{d}t_j$ interviennent, ce qui achève la démonstration.
\end{proof}

\begin{lemma}
\label{derham-lemma-blowup-twist-same-cohomology}
Soit $E = 0(P)$ le diviseur exceptionnel de l'éclatement $b$.
Pour tout $\mathcal{O}_X$-module localement libre $\mathcal{E}$ et
tout $0 \leq i \leq n - 1$, le morphisme
$$
\mathcal{E}
\longrightarrow
Rb_*(b^*\mathcal{E} \otimes_{\mathcal{O}_L} \mathcal{O}_L(iE))
$$
est un isomorphisme dans $D(\mathcal{O}_X)$.
\end{lemma}

\begin{proof}
D'après la formule de projection, il suffit de le montrer pour
$\mathcal{E} = \mathcal{O}_X$ ; voir Cohomologie, lemme
\ref{cohomology-lemma-projection-formula}.
Puisque $X$ est affine, il suffit de montrer que les morphismes
$$
H^0(X, \mathcal{O}_X) \to
H^0(L, \mathcal{O}_L) \to
H^0(L, \mathcal{O}_L(iE))
$$
sont des isomorphismes et que $H^j(X, \mathcal{O}_L(iE)) = 0$
pour $j > 0$ et $0 \leq i \leq n - 1$ ; voir Cohomologie des schémas, lemme
\ref{coherent-lemma-quasi-coherence-higher-direct-images-application}.
Puisque $\pi$ est affine, on peut calculer les sections globales et la
cohomologie après application de $\pi_*$ ; voir Cohomologie des schémas, lemme
\ref{coherent-lemma-relative-affine-cohomology}. Si $n = 1$,
$L \to X$ est un isomorphisme et $i = 0$, d'où la première assertion.
Si $n > 1$, considérons la composée
$$
H^0(X, \mathcal{O}_X) \to H^0(L, \mathcal{O}_L) \to
H^0(L, \mathcal{O}_L(iE)) \to H^0(L \setminus E, \mathcal{O}_L) =
H^0(X \setminus Z, \mathcal{O}_X)
$$
Puisque
$H^0(X \setminus Z, \mathcal{O}_X) = H^0(X, \mathcal{O}_X)$ dans ce
cas, car $Z$ est de codimension $n \geq 2$ dans $X$ (détails omis), on en déduit
la première assertion. Pour la seconde, rappelons que
$\mathcal{O}_L(E) = \mathcal{O}_L(-1)$ ; voir Diviseurs, lemme
\ref{divisors-lemma-blowing-up-gives-effective-Cartier-divisor}.
On a donc
$$
\pi_*\mathcal{O}_L(iE) =
\pi_*\mathcal{O}_L(-i) =
\bigoplus\nolimits_{k \geq -i} \mathcal{O}_P(k)
$$,
comme expliqué dans
Compléments sur les morphismes, section \ref{more-morphisms-section-proj-spec}.
On conclut par l'annulation des groupes de cohomologie des faisceaux tordus
du faisceau structural sur $P = \mathbf{P}^{n - 1}_Z$,
établie dans Cohomologie des schémas, lemme
\ref{coherent-lemma-cohomology-projective-space-over-ring}.
\end{proof}























\section{Éclatements et cohomologie de de Rham}
\label{derham-section-blowing-up}

\noindent
Fixons un schéma de base $S$, un morphisme lisse $X \to S$ et un sous-schéma fermé
$Z \subset X$, lui aussi lisse sur $S$. Notons $b : X' \to X$
l'éclatement de $X$ le long de $Z$. Notons $E \subset X'$ le diviseur
exceptionnel. Voici le diagramme
\begin{equation}
\label{derham-equation-blowup}
\vcenter{
\xymatrix{
E \ar[r]_j \ar[d]_p & X' \ar[d]^b \\
Z \ar[r]^i & X
}
}
\end{equation}
Notre but dans cette section est de démontrer que l'application
$b^* : H_{dR}^*(X/S) \longrightarrow H_{dR}^*(X'/S)$
est injective (bien que l'on puisse en dire beaucoup plus).

\begin{lemma}
\label{derham-lemma-comparison-general}
Avec les notations de Compléments sur les morphismes, lemme \ref{more-morphisms-lemma-blowup},
pour $a \geq 0$, on a
\begin{enumerate}
\item le morphisme
$\Omega^a_{X/S} \to b_*\Omega^a_{X'/S}$ est un isomorphisme,
\item le morphisme $\Omega^a_{Z/S} \to p_*\Omega^a_{E/S}$ est un isomorphisme,
\item le morphisme $Rb_*\Omega^a_{X'/S} \to i_*Rp_*\Omega^a_{E/S}$ induit un isomorphisme
sur les faisceaux de cohomologie en degrés $\geq 1$.
\end{enumerate}
\end{lemma}

\begin{proof}
Soit $\epsilon : X_1 \to X$ un morphisme étale surjectif. Notons
$i_1 : Z_1 \to X_1$, $b_1 : X'_1 \to X_1$, $E_1 \subset X'_1$ et
$p_1 : E_1 \to Z_1$ les changements de base des objets considérés dans
Compléments sur les morphismes, lemme \ref{more-morphisms-lemma-blowup}.
Observons que $i_1$ est une immersion fermée
de schémas lisses sur $S$, et que $b_1$ est l'éclatement de centre
$Z_1$, d'après Diviseurs, lemme \ref{divisors-lemma-flat-base-change-blowing-up}.
Supposons que l'on puisse démontrer (1), (2) et (3)
pour les morphismes $b_1$, $p_1$ et $i_1$. Alors, d'après
le lemme \ref{derham-lemma-etale}, les images inverses par $\epsilon$
des morphismes de (1), (2) et (3) sont des isomorphismes. Puisque $\epsilon$
est un morphisme plat surjectif, on conclut.
Ainsi, en travaillant étale-localement et en utilisant
Compléments sur les morphismes, lemme \ref{more-morphisms-lemma-etale-local-structure},
on peut supposer être dans la situation étudiée à la
section \ref{derham-section-calculations}. Le lemme est alors
identique au lemme \ref{derham-lemma-comparison}.
\end{proof}

\begin{lemma}
\label{derham-lemma-distinguished-triangle-blowup}
Avec les notations de Compléments sur les morphismes, lemme \ref{more-morphisms-lemma-blowup},
et en notant $f : X \to S$ le morphisme structural, il existe un
triangle distingué canonique
$$
\Omega^\bullet_{X/S} \to
Rb_*(\Omega^\bullet_{X'/S}) \oplus i_*\Omega^\bullet_{Z/S} \to
i_*Rp_*(\Omega^\bullet_{E/S}) \to
\Omega^\bullet_{X/S}[1]
$$
dans $D(X, f^{-1}\mathcal{O}_S)$, où les quatre morphismes
$$
\begin{matrix}
\Omega^\bullet_{X/S} & \to & Rb_*(\Omega^\bullet_{X'/S}), \\
\Omega^\bullet_{X/S} & \to & i_*\Omega^\bullet_{Z/S}, \\
Rb_*(\Omega^\bullet_{X'/S}) & \to & i_*Rp_*(\Omega^\bullet_{E/S}), \\
i_*\Omega^\bullet_{Z/S} & \to & i_*Rp_*(\Omega^\bullet_{E/S})
\end{matrix}
$$
sont les morphismes canoniques (section \ref{derham-section-de-rham-complex}),
au changement de signe près de l'un d'entre eux.
\end{lemma}

\begin{proof}
Choisissons un triangle distingué
$$
C \to Rb_*\Omega^\bullet_{X'/S} \oplus i_*\Omega^\bullet_{Z/S}
\to i_*Rp_*\Omega^\bullet_{E/S} \to C[1]
$$
dans $D(X, f^{-1}\mathcal{O}_S)$. Il suffit de montrer que
$\Omega^\bullet_{X/S}$ est isomorphe à $C$ d'une manière compatible
avec les morphismes canoniques. Les axiomes des catégories triangulées
fournissent un morphisme de triangles distingués
$$
\xymatrix{
C' \ar[r] \ar[d] &
b_*\Omega^\bullet_{X'/S} \oplus i_*\Omega^\bullet_{Z/S} \ar[r] \ar[d] &
i_*p_*\Omega^\bullet_{E/S} \ar[r] \ar[d] &
C'[1] \ar[d] \\
C \ar[r] &
Rb_*\Omega^\bullet_{X'/S} \oplus i_*\Omega^\bullet_{Z/S} \ar[r] &
i_*Rp_*\Omega^\bullet_{E/S} \ar[r] &
C[1]
}
$$
D'après le lemme \ref{derham-lemma-comparison-general}, partie (3), et
Catégories dérivées, proposition \ref{derived-proposition-9}, on en déduit que
$C' \to C$ est un isomorphisme. D'après le lemme \ref{derham-lemma-comparison-general}, partie (2),
le morphisme $i_*\Omega^\bullet_{Z/S} \to i_*p_*\Omega^\bullet_{E/S}$
est un isomorphisme. Ainsi, $C' = b_*\Omega^\bullet_{X'/S}$
dans la catégorie dérivée. Enfin, le lemme \ref{derham-lemma-comparison-general},
partie (1), montre que cet objet est égal à $\Omega^\bullet_{X/S}$.
Nous omettons la vérification de la compatibilité avec les morphismes canoniques.
\end{proof}

\begin{proposition}
\label{derham-proposition-blowup-split}
Avec les notations de Compléments sur les morphismes, lemme \ref{more-morphisms-lemma-blowup},
le morphisme $\Omega^\bullet_{X/S} \to Rb_*\Omega^\bullet_{X'/S}$
admet une rétraction dans $D(X, (X \to S)^{-1}\mathcal{O}_S)$.
\end{proposition}

\begin{proof}
Considérons le triangle construit dans le
lemme \ref{derham-lemma-distinguished-triangle-blowup}.
Nous affirmons que le morphisme
$$
Rb_*(\Omega^\bullet_{X'/S}) \oplus i_*\Omega^\bullet_{Z/S} \to
i_*Rp_*(\Omega^\bullet_{E/S})
$$
admet une section dont l'image contient le facteur direct $i_*\Omega^\bullet_{Z/S}$.
D'après Catégories dérivées, lemme \ref{derived-lemma-split}, cela montrera que
la première flèche du triangle admet une rétraction qui s'annule sur
le facteur direct $i_*\Omega^\bullet_{Z/S}$, ce qui démontre le lemme.
Pour prouver l'assertion, nous décomposerons $Rp_*\Omega^\bullet_{E/S}$
en une somme directe dont le premier facteur correspond à
$\Omega^\bullet_{Z/S}$ et dont le second se relève
par $Rb_*\Omega^\bullet_{X'/S}$.

\medskip\noindent
Démonstration de l'assertion. On peut décomposer $X$ en sous-schémas ouverts et fermés
de dimension relative constante sur $S$ ; voir
Morphismes, lemme \ref{morphisms-lemma-smooth-omega-finite-locally-free}.
Puisque la catégorie dérivée $D(X, f^{-1}\mathcal{O})_S)$ se décompose alors
en un produit de catégories, on peut supposer $X$ de
dimension relative constante $N$ sur $S$. On peut décomposer
$Z = \coprod Z_m$ en sous-schémas ouverts et fermés de dimension
relative $m \geq 0$ sur $S$. La restriction $i_m : Z_m \to X$ de
$i$ à $Z_m$ est une immersion régulière de codimension $N - m$ ; voir Diviseurs, lemme
\ref{divisors-lemma-immersion-smooth-into-smooth-regular-immersion}.
Soit $E = \coprod E_m$ la décomposition correspondante, c'est-à-dire
que l'on pose $E_m = p^{-1}(Z_m)$. Si $p_m : E_m \to Z_m$ désigne la
restriction de $p$ à $E_m$, on dispose d'un isomorphisme canonique
$$
\tilde \xi_m :
\bigoplus\nolimits_{t = 0, \ldots, N - m - 1}
\Omega^\bullet_{Z_m/S}[-2t]
\longrightarrow
Rp_{m, *}\Omega^\bullet_{E_m/S}
$$
dans $D(Z_m, (Z_m \to S)^{-1}\mathcal{O}_S)$,
où, en degré $0$, on a le morphisme canonique
$\Omega^\bullet_{Z_m/S} \to Rp_{m, *}\Omega^\bullet_{E_m/S}$.
Voir la remarque \ref{derham-remark-projective-space-bundle-formula}.
On dispose donc d'un isomorphisme
$$
\tilde \xi :
\bigoplus\nolimits_m
\bigoplus\nolimits_{t = 0, \ldots, N - m - 1}
\Omega^\bullet_{Z_m/S}[-2t]
\longrightarrow
Rp_*(\Omega^\bullet_{E/S})
$$
dans $D(Z, (Z \to S)^{-1}\mathcal{O}_S)$
dont la restriction au facteur direct
$\Omega^\bullet_{Z/S} = \bigoplus \Omega^\bullet_{Z_m/S}$ de la source
est le morphisme canonique $\Omega^\bullet_{Z/S} \to Rp_*(\Omega^\bullet_{E/S})$.
Considérons les sous-complexes $M_m$ et $K_m$ du complexe
$\bigoplus\nolimits_{t = 0, \ldots, N - m - 1} \Omega^\bullet_{Z_m/S}[-2t]$
introduits dans la remarque \ref{derham-remark-projective-space-bundle-formula}.
Posons
$$
M = \bigoplus M_m
\quad\text{et}\quad
K = \bigoplus K_m
$$
On a $M = K[-2]$ et, par construction, le morphisme
$$
c_{E/Z} \oplus \tilde \xi|_M :
\Omega^\bullet_{Z/S} \oplus M
\longrightarrow
Rp_*(\Omega^\bullet_{E/S})
$$
est un isomorphisme (voir la remarque citée ci-dessus).

\medskip\noindent
Considérons le morphisme
$$
\delta : \Omega^\bullet_{E/S}[-2] \longrightarrow \Omega^\bullet_{X'/S}
$$
dans $D(X', (X' \to S)^{-1}\mathcal{O}_S)$ du
lemme \ref{derham-lemma-gysin-via-log-complex},
qui a la propriété que la composée
$$
\Omega^\bullet_{E/S}[-2] \longrightarrow \Omega^\bullet_{X'/S}
\longrightarrow
\Omega^\bullet_{E/S}
$$
est le morphisme $\theta'$ de la remarque \ref{derham-remark-cup-product-as-a-map} pour
$c_1^{dR}(\mathcal{O}_{X'}(-E))|_E) = c_1^{dR}(\mathcal{O}_E(1))$.
La dernière assertion de la remarque \ref{derham-remark-projective-space-bundle-formula}
montre que le diagramme
$$
\xymatrix{
K[-2] \ar[d]_{(\tilde \xi|_K)[-2]} \ar[r]_{\text{id}} &
M \ar[d]^{\tilde x|_M} \\
Rp_*\Omega^\bullet_{E/S}[-2] \ar[r]^-{Rp_*\theta'} &
Rp_*\Omega^\bullet_{E/S}
}
$$
est commutatif. La section cherchée dans notre
assertion est donc le morphisme
\begin{align*}
Rp_*(\Omega^\bullet_{E/S})
& \xrightarrow{(c_{E/Z} \oplus \tilde \xi|_M)^{-1}}
\Omega^\bullet_{Z/S} \oplus M \\
& \xrightarrow{\text{id} \oplus \text{id}^{-1}}
\Omega^\bullet_{Z/S} \oplus K[-2] \\
& \xrightarrow{\text{id} \oplus (\tilde \xi|_K)[-2]}
\Omega^\bullet_{Z/S} \oplus Rp_*\Omega^\bullet_{E/S}[-2] \\
& \xrightarrow{\text{id} \oplus Rb_*\delta}
\Omega^\bullet_{Z/S} \oplus Rb_*\Omega^\bullet_{X'/S}
\end{align*}
La relation entre $\theta'$ et $\delta$ énoncée ci-dessus,
ainsi que le diagramme commutatif faisant intervenir $\theta'$, $\tilde \xi|_K$
et $\tilde \xi|_M$ ci-dessus, donnent exactement ce qu'il faut
pour montrer qu'il s'agit d'une section du morphisme canonique
$\Omega^\bullet_{Z/S} \oplus Rb_*(\Omega^\bullet_{X'/S}) \to
Rp_*(\Omega^\bullet_{E/S})$, ce qui achève la démonstration de l'assertion.
\end{proof}

\noindent
Le lemme \ref{derham-lemma-splitting-on-omega-a}
montre qu'il est nettement plus facile de construire la rétraction en cohomologie
de Hodge que d'établir le résultat de la
proposition \ref{derham-proposition-blowup-split}.
Nous invitons vivement le lecteur à passer à la section suivante.

\begin{lemma}
\label{derham-lemma-ext-zero}
Soit $i : Z \to X$ une immersion fermée de schémas, régulière de
codimension $c$. Alors $\Ext^q_{\mathcal{O}_X}(i_*\mathcal{F}, \mathcal{E}) = 0$
pour $q < c$, lorsque $\mathcal{E}$ est localement libre sur $X$ et $\mathcal{F}$
est un $\mathcal{O}_Z$-module quelconque.
\end{lemma}

\begin{proof}
D'après la suite spectrale du local au global pour $\Ext$, il suffit
de le démontrer localement sur les ouverts affines de $X$. Voir
Cohomologie, section \ref{cohomology-section-ext}.
On peut donc supposer $X = \Spec(A)$
et l'existence d'une suite régulière $f_1, \ldots, f_c$ dans $A$
telle que $Z = \Spec(A/(f_1, \ldots, f_c))$. On peut supposer $c \geq 1$.
On voit alors que $f_1 : \mathcal{E} \to \mathcal{E}$
est injectif. Puisque $i_*\mathcal{F}$ est annulé par $f_1$,
cela montre que le lemme est vrai pour $i = 0$ et que l'on dispose
d'une surjection
$$
\Ext^{q - 1}_{\mathcal{O}_X}(i_*\mathcal{F}, \mathcal{E}/f_1\mathcal{E})
\longrightarrow
\Ext^q_{\mathcal{O}_X}(i_*\mathcal{F}, \mathcal{E})
$$
Il suffit donc de montrer que la source de cette flèche est nulle.
Répétons l'argument : si $c \geq 2$, le morphisme
$f_2 : \mathcal{E}/f_1\mathcal{E} \to \mathcal{E}/f_1\mathcal{E}$
est injectif. Puisque $i_*\mathcal{F}$ est annulé par $f_2$,
cela montre que le lemme est vrai pour $q = 1$ et que l'on dispose d'une
surjection
$$
\Ext^{q - 2}_{\mathcal{O}_X}(i_*\mathcal{F},
\mathcal{E}/f_1\mathcal{E} + f_2\mathcal{E})
\longrightarrow
\Ext^{q - 1}_{\mathcal{O}_X}(i_*\mathcal{F}, \mathcal{E}/f_1\mathcal{E})
$$
En poursuivant de cette manière, on démontre le lemme.
\end{proof}

\begin{lemma}
\label{derham-lemma-splitting-on-omega-a}
Avec les notations de Compléments sur les morphismes, lemme \ref{more-morphisms-lemma-blowup},
pour $a \geq 0$, il existe une unique flèche
$Rb_*\Omega^a_{X'/S} \to \Omega^a_{X/S}$ dans $D(\mathcal{O}_X)$
dont la composée avec $\Omega^a_{X/S} \to Rb_*\Omega^a_{X'/S}$
est l'identité de $\Omega^a_{X/S}$.
\end{lemma}

\begin{proof}
On peut décomposer $X$ en sous-schémas ouverts et fermés
de dimension relative constante sur $S$ ; voir
Morphismes, lemme \ref{morphisms-lemma-smooth-omega-finite-locally-free}.
Puisque la catégorie dérivée $D(X, f^{-1}\mathcal{O})_S)$ se décompose alors
en un produit de catégories, on peut supposer $X$ de
dimension relative constante $N$ sur $S$. On peut décomposer
$Z = \coprod Z_m$ en sous-schémas ouverts et fermés de dimension
relative $m \geq 0$ sur $S$. La restriction $i_m : Z_m \to X$ de
$i$ à $Z_m$ est une immersion régulière de codimension $N - m$ ; voir Diviseurs, lemme
\ref{divisors-lemma-immersion-smooth-into-smooth-regular-immersion}.
Soit $E = \coprod E_m$ la décomposition correspondante, c'est-à-dire
que l'on pose $E_m = p^{-1}(Z_m)$. Nous affirmons qu'il existe des morphismes naturels
$$
b^*\Omega^a_{X/S} \to \Omega^a_{X'/S} \to
b^*\Omega^a_{X/S} \otimes_{\mathcal{O}_{X'}}
\mathcal{O}_{X'}(\sum (N - m - 1)E_m)
$$
dont la composée est induite par l'inclusion
$\mathcal{O}_{X'} \to \mathcal{O}_{X'}(\sum (N - m - 1)E_m)$.
Pour le démontrer, il suffit de montrer que le
conoyau de la première flèche est, localement sur $X'$, annulé par
une équation locale du diviseur de Cartier effectif $\sum (N - m - 1)E_m$.
Pour cela, on peut travailler étale-localement sur $X$ comme dans la
démonstration du lemme \ref{derham-lemma-comparison-general}, et appliquer
le lemme \ref{derham-lemma-comparison-bis}.
En calculant étale-localement à l'aide du lemme \ref{derham-lemma-blowup-twist-same-cohomology},
on voit que la composée induite
$$
\Omega^a_{X/S} \to Rb_*\Omega^a_{X'/S} \to
Rb_*\left(b^*\Omega^a_{X/S} \otimes_{\mathcal{O}_{X'}}
\mathcal{O}_{X'}(\sum (N - m - 1)E_m)\right)
$$
est un isomorphisme dans $D(\mathcal{O}_X)$,
d'où l'existence du morphisme du lemme.

\medskip\noindent
Pour l'unicité, il suffit de montrer qu'il n'existe pas de morphisme non nul de
$\tau_{\geq 1}Rb_*\Omega_{X'/S}$ vers $\Omega^a_{X/S}$ dans $D(\mathcal{O}_X)$.
Il suffit pour cela de montrer qu'il n'existe pas de morphisme non nul
de $R^qb_*\Omega_{X'/s}[-q]$ vers $\Omega^a_{X/S}$ dans $D(\mathcal{O}_X)$
pour $q \geq 1$ (détails omis). D'après
le lemme \ref{derham-lemma-comparison-general},
$R^qb_*\Omega_{X'/s} \cong i_*R^qp_*\Omega^a_{E/S}$
est l'image directe d'un module sur $Z = \coprod Z_m$.
Observons en outre que la restriction de $R^qp_*\Omega^a_{E/S}$
à $Z_m$ n'est non nulle que pour $q < N - m$. En effet, les fibres de
$E_m \to Z_m$ sont de dimension $N - m - 1$, et l'on peut appliquer Limites, lemme
\ref{limits-lemma-higher-direct-images-zero-above-dimension-fibre}.
L'annulation voulue résulte donc du lemme \ref{derham-lemma-ext-zero}.
\end{proof}














\section{Comparaison de faisceaux de formes différentielles}
\label{derham-section-quasi-finite-syntomic}

\noindent
Le but de cette section est de construire, pour tout morphisme syntomique
localement quasi-fini de schémas $f : Y \to X$, des morphismes
$$
c^p_{Y/X} :
\Omega^p_{Y/\mathbf{Z}}
\longrightarrow
f^*\Omega^p_{X/\mathbf{Z}} \otimes_{\mathcal{O}_Y} \det(\NL_{Y/X})
$$
pour tout $p \geq 0$, satisfaisant aux propriétés suivantes :
\begin{enumerate}
\item
\label{derham-item-degree-zero}
$c_{Y/X}^0(1) = \delta(\NL_{Y/X})$ ; voir
Discriminants, section \ref{discriminant-section-tate-map},
\item
\label{derham-item-multiplicative}
$c_{Y/X}^{q + p}(\omega \wedge \eta) = \omega \wedge c_{Y/X}^p(\eta)$
pour les sections locales $\omega$ de $f^*\Omega^q_{X/\mathbf{Z}}$
et $\eta$ de $\Omega^p_{Y/\mathbf{Z}}$,
\item
\label{derham-item-base-change}
la formation de $c^p_{Y/X}$ est compatible avec la restriction aux ouverts
et le changement de base
(au sens du lemme \ref{derham-lemma-base-change-Garel-upstairs}).
\end{enumerate}
Notons que ces conditions impliquent que les morphismes $c^p_{Y/X}$ sont
$\mathcal{O}_Y$-linéaires et que leur composée avec
$f^*\Omega^p_{X/\mathbf{Z}} \to \Omega^p_{Y/\mathbf{Z}}$
est la multiplication par $\delta(\NL_{Y/X})$.
Nous montrerons en fait qu'il existe une unique famille d'opérateurs
$c^p_{Y/X}$ possédant les propriétés (1), (2) et (3). Ce résultat sera appliqué à la
section \ref{derham-section-trace} pour construire le morphisme trace sur
les complexes de de Rham lorsque $f$ est fini.

\begin{lemma}
\label{derham-lemma-funny-map}
Soit $R$ un anneau, et considérons un diagramme commutatif
$$
\xymatrix{
0 \ar[r] &
K^0 \ar[r] &
L^0 \ar[r] &
M^0 \ar[r] & 0 \\
& & L^{-1} \ar[u]_\partial \ar@{=}[r] &
M^{-1} \ar[u]
}
$$
de $R$-modules dont la ligne supérieure est exacte, et où $M^0$ et $M^{-1}$
sont libres de même rang fini. Pour $p \geq 0$, il existe des morphismes canoniques
$$
c^p : 
\wedge^p(H^0(L^\bullet))
\longrightarrow
\wedge^p(K^0) \otimes_R \det(M^\bullet)
$$
possédant les propriétés suivantes :
\begin{enumerate}
\item $c^0(1) = \delta(M^\bullet)$, et
\item $c^{q + p}(\omega \wedge \eta) = \omega \wedge c^p(\eta)$
pour $\omega \in \wedge^q(K^0)$ et $\eta \in \wedge^p(H^0(L^\bullet))$.
\end{enumerate}
En particulier, la composée de $c^p$ avec le morphisme
$\wedge^p(K^0) \to \wedge^p(H^0(L^\bullet))$ est la multiplication
par $\delta(M^\bullet)$.
\end{lemma}

\begin{proof}
Supposons $M^0$ et $M^{-1}$ libres de rang $n$. Pour tout $p \geq 0$,
il existe une unique surjection
$$
\pi_p :
\wedge^{p + n}(L^0)
\longrightarrow
\wedge^p(K^0) \otimes \wedge^n(M^0)
$$
possédant les deux propriétés suivantes : (i) on a
$
\pi_p(k_1 \wedge \ldots \wedge k_p \wedge l_1 \wedge \ldots \wedge l_n) =
k_1 \wedge \ldots \wedge k_p \otimes
m_1 \wedge \ldots \wedge m_n
$
si $k_1, \ldots, k_p \in K^0$ et si $l_1, \ldots, l_n \in L^0$ ont pour images
$m_1, \ldots, m_n \in M^0$, et (ii) le noyau de $\pi_p$ est le
sous-module engendré par les produits extérieurs
$l_1 \wedge \ldots \wedge l_{p + n}$ tels qu'un nombre $> p$ des
$l_j$ appartiennent à $K^0$. Choisissons une base $e_1, \ldots, e_n$ de $L^{-1} = M^{-1}$.
On définit alors notre morphisme par la règle
$$
\eta \mapsto
\pi_p(\tilde \eta \wedge \partial(e_1) \wedge \ldots \wedge \partial(e_n))
\otimes (e_1 \wedge \ldots \wedge e_n)^{\otimes -1}
$$
où $\eta \in \wedge^p(H^0(L^\bullet))$ et où $\tilde \eta \in \wedge^p(L^0)$
est un relèvement de $\eta$. Le membre de droite a un sens, car
$\det(M^\bullet) = \wedge^n(M^0) \otimes \wedge^n(M^{-1})^{\otimes -1}$.
Il est indépendant du choix de $\tilde \eta$ puisque
$\partial(\xi) \wedge \partial(e_1) \wedge \ldots \wedge \partial(e_n) = 0$
pour tout $\xi \in L^{-1}$. Un calcul direct montre que le morphisme
est indépendant du choix de la base de $L^{-1} = M^{-1}$.
Il résulte immédiatement de la définition de $\delta(M^\bullet)$ dans
Compléments d'algèbre, section \ref{more-algebra-section-determinants-complexes},
que $c^0(1) = \delta(M^\bullet)$. Enfin, la règle de l'assertion (2)
se vérifie directement à partir de la définition.
\end{proof}

\begin{lemma}
\label{derham-lemma-funny-map-functorial}
Soit $R_1 \to R_2$ un homomorphisme d'anneaux. Pour $i = 1, 2$, considérons
des diagrammes commutatifs
$$
\xymatrix{
0 \ar[r] &
K_i^0 \ar[r] &
L_i^0 \ar[r] &
M_i^0 \ar[r] & 0 \\
& & L_i^{-1} \ar[u]_{\partial_i} \ar@{=}[r] &
M_i^{-1} \ar[u]
}
$$
de $R_i$-modules comme dans le lemme \ref{derham-lemma-funny-map}. Supposons donnés
des morphismes $K_1^0 \to K_2^0$, $L_1^j \to L_2^j$ et $M_1^j \to M_2^j$
compatibles avec le morphisme d'anneaux $R_1 \to R_2$ et
avec les morphismes des diagrammes affichés. Si les morphismes
$M_1^j \otimes_{R_1} R_2 \to M_2^j$ sont des isomorphismes, alors
les diagrammes
$$
\xymatrix{ 
\wedge^p(H^0(L_1^\bullet)) \ar[d]
\ar[r]_-{c^p} &
\wedge^p(K_1^0) \otimes_{R_1} \det(M_1^\bullet) \ar[d] \\
\wedge^p(H^0(L_2^\bullet))
\ar[r]^-{c^p} &
\wedge^p(K_2^0) \otimes_{R_2} \det(M_2^\bullet)
}
$$
sont commutatifs.
\end{lemma}

\begin{proof}
Cela résulte de la description explicite du morphisme donnée
dans la démonstration du lemme \ref{derham-lemma-funny-map}. Notons que l'hypothèse
selon laquelle les $M_1^j \otimes_{R_1} R_2 \to M_2^j$ sont des isomorphismes
est nécessaire pour voir que la base utilisée pour $M_1^{-1}$ s'envoie sur une base
de $M_2^{-1}$ (et donc que l'on peut utiliser la « même » base
dans les deux constructions pour en vérifier la compatibilité).
\end{proof}

\begin{remark}
\label{derham-remark-local-description}
Soit $A$ un anneau. Soit $P = A[x_1, \ldots, x_n]$. Soient
$f_1, \ldots, f_n \in P$, et posons $B = P/(f_1, \ldots, f_n)$.
Supposons $A \to B$ quasi-fini. Alors
$B$ est une intersection complète globale relative sur $A$ (Algèbre, définition
\ref{algebra-definition-relative-global-complete-intersection}), et
$(f_1, \ldots, f_n)/(f_1, \ldots, f_n)^2$ est libre, engendré par
les classes $\overline{f}_i$, d'après Algèbre, lemme
\ref{algebra-lemma-relative-global-complete-intersection-conormal}.
Considérons le diagramme suivant
$$
\xymatrix{
\Omega_{A/\mathbf{Z}} \otimes_A B \ar[r] &
\Omega_{P/\mathbf{Z}} \otimes_P B \ar[r] &
\Omega_{P/A} \otimes_P B \\
&
(f_1, \ldots, f_n)/(f_1, \ldots, f_n)^2 \ar[u] \ar@{=}[r] &
(f_1, \ldots, f_n)/(f_1, \ldots, f_n)^2 \ar[u]
}
$$
La colonne de droite représente $\NL_{B/A}$ dans $D(B)$ ; sa cohomologie est donc
$\Omega_{B/A}$ en degré $0$. La ligne supérieure est la suite exacte courte scindée
$0 \to \Omega_{A/\mathbf{Z}} \otimes_A B \to
\Omega_{P/\mathbf{Z}} \otimes_P B \to \Omega_{P/A} \otimes_P B \to 0$.
La colonne du milieu a pour cohomologie $\Omega_{B/\mathbf{Z}}$ en degré $0$,
d'après Algèbre, lemme \ref{algebra-lemma-differential-seq}.
Ainsi, le lemme \ref{derham-lemma-funny-map} fournit des morphismes canoniques de $B$-modules
$$
\Omega^p_{B/\mathbf{Z}} \longrightarrow
\Omega^p_{A/\mathbf{Z}} \otimes_A \det(\NL_{B/A})
$$
pour $p \geq 0$, satisfaisant aux propriétés (1) et (2) et, en particulier,
tels que la composée avec
$\Omega^p_{A/\mathbf{Z}} \to \Omega^p_{B/\mathbf{Z}}$
soit la multiplication par $\delta(\NL_{B/A})$.
\end{remark}

\begin{lemma}
\label{derham-lemma-base-change-Garel-upstairs}
Considérons un diagramme commutatif
$$
\xymatrix{
Y' \ar[d]_{f'} \ar[r]_b & Y \ar[d]^f \\
X' \ar[r]^a & X
}
$$
de schémas qui induit un isomorphisme de $Y'$ sur un sous-schéma ouvert de
$X' \times_X Y$. Supposons $f$ localement quasi-fini et syntomique, et
supposons donnés des morphismes $c^p_{Y/X} : \Omega^p_{Y/\mathbf{Z}} \to
f^*\Omega^p_{X/\mathbf{Z}} \otimes_{\mathcal{O}_Y} \det(\NL_{Y/X})$
pour $p \geq 0$ satisfaisant à
(\ref{derham-item-degree-zero}) et (\ref{derham-item-multiplicative}).
Il existe alors au plus une famille de morphismes
$c^p_{Y'/X'} : \Omega^p_{Y'/\mathbf{Z}} \to
(f')^*\Omega^p_{X'/\mathbf{Z}} \otimes_{\mathcal{O}_{Y'}} \det(\NL_{Y'/X'})$
pour $p \geq 0$
satisfaisant à (\ref{derham-item-degree-zero}) et (\ref{derham-item-multiplicative}),
telle que les diagrammes
$$
\xymatrix{
b^*\Omega^p_{Y/\mathbf{Z}} \ar[rr]_-{b^*c^p_{Y/X}} \ar[d] & &
b^*(f^*\Omega^p_{X/\mathbf{Z}} \otimes_{\mathcal{O}_Y} \det(\NL_{Y/X}))
\ar@{=}[r] &
(f')^*a^*\Omega^p_{X/\mathbf{Z}} \otimes_{\mathcal{O}_Y} b^*\det(\NL_{Y/X}))
\ar[d] \\
\Omega^p_{Y'/\mathbf{Z}} \ar[rrr]^-{c^p_{Y'/X'}} & & &
(f')^*\Omega^p_{X'/\mathbf{Z}} \otimes_{\mathcal{O}_{Y'}} \det(\NL_{Y'/X'})
}
$$
soient commutatifs pour tout $p \geq 0$. Les flèches verticales utilisent ici les morphismes
$b^*\Omega^p_{Y/\mathbf{Z}} \to \Omega^p_{Y'/\mathbf{Z}}$ et
$a^*\Omega^p_{X/\mathbf{Z}} \to \Omega^p_{X'/\mathbf{Z}}$
de la section \ref{derham-section-de-rham-complex},
ainsi que l'identification $b^*\det(\NL_{Y/X}) = \det(\NL_{Y'/X'})$ de
Discriminants, section \ref{discriminant-section-tate-map}.
\end{lemma}

\begin{proof}
Le morphisme
$$
(f')^*\Omega_{X'/\mathbf{Z}} \oplus b^*\Omega_{Y/\mathbf{Z}}
\longrightarrow
\Omega_{Y'/\mathbf{Z}}
$$
est surjectif, car $Y'$ est un sous-schéma ouvert du produit fibré ;
l'énoncé algébrique correspondant affirme que $\Omega_{B \otimes_A C/\mathbf{Z}}$
est engendré, comme module, par les images de $\Omega_{B/\mathbf{Z}}$ et de
$\Omega_{C/\mathbf{Z}}$. Ainsi, pour tout $p$, le morphisme
$$
\bigoplus\nolimits_{i = 0, \ldots, p}
(f')^*\Omega^i_{X'/\mathbf{Z}} \otimes
b^*\Omega^{p - i}_{Y/\mathbf{Z}}
\longrightarrow
\Omega^p_{Y'/\mathbf{Z}}
$$
est surjectif. Les conditions (\ref{derham-item-degree-zero}) et
(\ref{derham-item-multiplicative}), combinées à la commutativité
du diagramme du lemme, impliquent que le morphisme $c^p_{Y'/X'}$
est déterminé (mais son existence n'est pas immédiate).
\end{proof}

\begin{lemma}
\label{derham-lemma-existence-base-change-Garel-upstairs}
Considérons un diagramme commutatif
$$
\xymatrix{
Y' \ar[d]_{f'} \ar[r]_b & Y \ar[d]^f \\
X' \ar[r]^a & X
}
$$
de schémas qui induit un isomorphisme de $Y'$ sur un sous-schéma ouvert de
$X' \times_X Y$. Supposons $f$ localement quasi-fini et syntomique. Alors, pour
tout $y' \in Y'$, on peut trouver des ouverts $V' \subset Y'$,
$V \subset Y$, $U' \subset X$ et $U \subset X$ tels que $y' \in V'$,
$f'(V') \subset U'$, $b(V') \subset V$, $a(U') \subset U$,
$f(V) \subset U$, et tels que, pour $p \geq 0$, il existe des morphismes
$c^p_{V/U}$ et $c^p_{V'/U'}$ satisfaisant à (\ref{derham-item-degree-zero})
et (\ref{derham-item-multiplicative}), compatibles avec le diagramme
$$
\xymatrix{
V' \ar[d] \ar[r]_b & V \ar[d] \\
U' \ar[r] & U
}
$$
au sens précisé dans le lemme \ref{derham-lemma-base-change-Garel-upstairs}.
\end{lemma}

\begin{proof}
Il est clair que l'on peut remplacer $Y'$ par $X' \times_X Y$ pour le démontrer.
Choisissons des ouverts affines $V \subset Y$ et $U \subset X$, où $V$ contient
l'image de $y'$, comme dans
Discriminants, lemme \ref{discriminant-lemma-syntomic-quasi-finite}, partie (5).
Choisissons un ouvert affine $U' \subset X'$ contenant l'image de $y'$
et s'envoyant dans $U$. En remplaçant $X, Y, X', Y'$ par
$U, V, U', U' \times_U V$, on peut supposer être dans la situation
décrite au paragraphe suivant.

\medskip\noindent
Supposons que $X' = \Spec(A')$, $X = \Spec(A)$, $Y = \Spec(B)$ et
$Y' = \Spec(A' \otimes_A B)$, où $B = A[x_1, \ldots, x_n]/(f_1, \ldots, f_n)$
est une intersection complète globale relative sur $A$. Alors
$B' = A'[x_1, \ldots, x_n]/(f'_1, \ldots, f'_n)$ est aussi une intersection
complète globale relative sur $A'$, où $f'_i$ est l'image de $f_i$
dans $A'[x_1, \ldots, x_n]$ ; voir Algèbre, lemme
\ref{algebra-lemma-base-change-relative-global-complete-intersection}.
La construction de la remarque \ref{derham-remark-local-description}
fournit les morphismes $c^p_{V/U}$ et $c^p_{V'/U'}$.
Ils sont compatibles d'après le lemme \ref{derham-lemma-funny-map-functorial}
et la compatibilité des présentations de $B$ et de $B'$ sur
$A$ et $A'$. Nous omettons certains détails.
\end{proof}

\begin{lemma}
\label{derham-lemma-Garel-upstairs}
Il existe une unique règle qui, à tout morphisme de schémas syntomique
localement quasi-fini $f : Y \to X$, associe des morphismes de $\mathcal{O}_Y$-modules
$$
c^p_{Y/X} :
\Omega^p_{Y/\mathbf{Z}}
\longrightarrow
f^*\Omega^p_{X/\mathbf{Z}} \otimes_{\mathcal{O}_Y} \det(\NL_{Y/X})
$$
pour $p \geq 0$, satisfaisant à (\ref{derham-item-degree-zero}), (\ref{derham-item-multiplicative})
et (\ref{derham-item-base-change}). En particulier, la composée de
$c^p_{Y/X}$ avec $f^*\Omega^p_{X/\mathbf{Z}} \to \Omega^p_{Y/\mathbf{Z}}$
est la multiplication par $\delta(\NL_{Y/X})$.
\end{lemma}

\begin{proof}
Cette démonstration est très proche de celle de
Discriminants, proposition \ref{discriminant-proposition-tate-map},
que nous suggérons au lecteur de consulter d'abord.

\medskip\noindent
Reformulons l'énoncé. Considérons la catégorie
$\mathcal{C}$ dont les objets, notés $Y/X$, sont les morphismes syntomiques localement quasi-finis
$f : Y \to X$ de schémas, et dont les morphismes
$b/a : Y'/X' \to Y/X$ sont les diagrammes commutatifs
$$
\xymatrix{
Y' \ar[d]_{f'} \ar[r]_b & Y \ar[d]^f \\
X' \ar[r]^a & X
}
$$
qui induisent un isomorphisme de $Y'$ sur un sous-schéma ouvert de
$X' \times_X Y$. Le lemme signifie que, pour tout objet
$Y/X$ de $\mathcal{C}$, on dispose de morphismes $c^p_{Y/X}$, $p \geq 0$,
possédant les propriétés (\ref{derham-item-degree-zero}) et (\ref{derham-item-multiplicative}),
et que, pour tout morphisme $b/a : Y'/X' \to Y/X$ de $\mathcal{C}$,
$c^p_{Y/X}$ et $c^p_{Y'/X'}$ sont compatibles comme dans le
lemme \ref{derham-lemma-base-change-Garel-upstairs}.

\medskip\noindent
Étant donnés $Y/X$ dans $\mathcal{C}$ et $y \in Y$, on peut trouver
des ouverts affines $V \subset Y$ et $U \subset X$ avec $f(V) \subset U$
tels qu'il existe des morphismes
$$
\Omega^p_{Y/\mathbf{Z}}|_V
\longrightarrow
\left(
f^*\Omega^p_{X/\mathbf{Z}} \otimes_{\mathcal{O}_Y} \det(\NL_{Y/X})
\right)|_V
$$
possédant les propriétés (\ref{derham-item-degree-zero}) et (\ref{derham-item-multiplicative}).
Il suffit pour cela de choisir des ouverts affines comme dans
Discriminants, lemme \ref{discriminant-lemma-syntomic-quasi-finite}, partie (5),
puis d'utiliser la construction de la remarque \ref{derham-remark-local-description}.

\medskip\noindent
Notons que le lieu étale de $f$ est exactement l'ensemble des points où
$\delta(\NL_{Y/X})$ ne s'annule pas ; voir la discussion dans
Discriminants, section \ref{discriminant-section-tate-map}.
En particulier,
$f^*\Omega^p_{X/\mathbf{Z}} \to \Omega^p_{Y/\mathbf{Z}}$
devient un isomorphisme après inversion de $\delta(\NL_{Y/X})$.
On en déduit que, si $\Omega^p_{X/\mathbf{Z}}$ est localement libre de rang fini et si
l'annulateur de $\delta(\NL_{Y/X})$ dans $\mathcal{O}_Y$ est nul,
les morphismes locaux construits au paragraphe précédent
sont uniques et se recollent automatiquement !

\medskip\noindent
Notons $\mathcal{C}_{nice} \subset \mathcal{C}$ la sous-catégorie pleine
formée des $Y/X$ tels que
\begin{enumerate}
\item $\Omega_{X/\mathbf{Z}}$ est localement libre, et
\item l'annulateur de $\delta(\NL_{Y/X})$ dans $\mathcal{O}_Y$ est nul.
\end{enumerate}
D'après les observations du paragraphe précédent, pour tout
objet $Y/X$ de $\mathcal{C}_{nice}$, il existe des morphismes uniques
$c^p_{Y/X}$, $p \geq 0$, satisfaisant aux conditions (\ref{derham-item-degree-zero})
et (\ref{derham-item-multiplicative}). Si $b/a : Y'/X' \to Y/X$
est un morphisme de $\mathcal{C}_{nice}$, les morphismes
$c^p_{Y/X}$ et $c^p_{Y'/X'}$ sont compatibles comme dans le
lemme \ref{derham-lemma-base-change-Garel-upstairs} : en effet,
des morphismes compatibles existent localement d'après le
lemme \ref{derham-lemma-existence-base-change-Garel-upstairs},
et l'unicité que nous venons de mentionner montre qu'ils coïncident
avec $c^p_{Y/X}$ et $c^p_{Y'/X'}$.
Autrement dit, nous avons résolu le problème
sur la sous-catégorie pleine $\mathcal{C}_{nice}$. Pour $Y/X$ dans $\mathcal{C}_{nice}$,
nous continuerons à noter $c^p_{Y/X}$ la solution ainsi obtenue.

\medskip\noindent
Plus généralement, supposons donné un morphisme $b/a : Y'/X' \to Y/X$
de $\mathcal{C}$ tel que $Y/X$ soit un objet de $\mathcal{C}_{nice}$.
Le même argument, utilisant cette fois l'unicité du
lemme \ref{derham-lemma-base-change-Garel-upstairs}, montre qu'il existe des
morphismes $c^p_{Y'/X'}$, $p \geq 0$, satisfaisant aux conditions (\ref{derham-item-degree-zero})
et (\ref{derham-item-multiplicative}), compatibles avec les morphismes déjà construits
$c^p_{Y/X}$. Appelons-les les changements de base
$(b/a)^*c^p_{Y/X}$.

\medskip\noindent
Considérons des morphismes
$$
Y_1/X_1 \xleftarrow{b_1/a_1} Y/X \xrightarrow{b_2/a_2} Y_2/X_2
$$
dans $\mathcal{C}$ tels que $Y_1/X_1$ et $Y_2/X_2$ soient des objets
de $\mathcal{C}_{nice}$.
{\bf Assertion.} Les deux changements de base $(b_i/a_i)^*c^p_{Y_i/X_i}$, $i = 1, 2$,
sont égaux. Nous montrerons d'abord que cette assertion entraîne le lemme,
puis nous la démontrerons.

\medskip\noindent
Soient $d, n \geq 1$, et considérons le morphisme syntomique
localement quasi-fini $Y_{n, d} \to X_{n, d}$
construit dans Discriminants, exemple
\ref{discriminant-example-universal-quasi-finite-syntomic}.
Alors $Y_{n, d}$ et $Y_{n, d}$ sont des schémas irréductibles, de type fini et
lisses sur $\mathbf{Z}$. En effet, $X_{n, d}$ est le spectre d'un
anneau de polynômes sur $\mathbf{Z}$, et $Y_{n, d}$ est un sous-schéma ouvert
d'un tel spectre. Le morphisme $Y_{n, d} \to X_{n, d}$ est syntomique localement quasi-fini
et étale au-dessus d'un ouvert dense ; voir Discriminants, lemme
\ref{discriminant-lemma-universal-quasi-finite-syntomic-etale}.
Ainsi, $\delta(\NL_{Y_{n, d}/X_{n, d}})$ est non nul : on dispose, par exemple,
de la description locale de $\delta(\NL_{Y/X})$ dans
Discriminants, remarque \ref{discriminant-remark-local-description-delta},
et de la description locale des morphismes étales dans
Morphismes, lemme \ref{morphisms-lemma-etale-at-point}, partie (8).
Or une section non nulle d'un module inversible sur un schéma
régulier irréductible a un annulateur nul. Ainsi,
$Y_{n, d}/X_{n, d}$ est un objet de $\mathcal{C}_{nice}$.

\medskip\noindent
Soit $Y/X$ un objet quelconque de $\mathcal{C}$. Soit $y \in Y$.
D'après Discriminants, lemme \ref{discriminant-lemma-locally-comes-from-universal},
on peut trouver $n, d \geq 1$ et des morphismes
$$
Y/X \leftarrow V/U \xrightarrow{b/a} Y_{n, d}/X_{n, d}
$$
de $\mathcal{C}$ tels que $V \subset Y$ et $U \subset X$ soient ouverts.
On dispose alors du changement de base $c^p_{V/U} = (b/a)^*c^p_{Y_{n, d}/X_{n, d}}$.
L'assertion garantit que ces morphismes construits localement $c^p_{V/U}$ se recollent !
On obtient donc des morphismes globaux bien définis $c^p_{Y/X}$ possédant les propriétés
(\ref{derham-item-degree-zero}) et (\ref{derham-item-multiplicative}).
Enfin, soit $b/a : Y'/X' \to Y/X$ un morphisme quelconque de $\mathcal{C}$.
On dispose des morphismes $c^p_{Y'/X'}$, $p \geq 0$ et $c^p_{Y/X}$, $p \geq 0$,
construits dans ce paragraphe. Pour vérifier leur compatibilité au sens du
lemme \ref{derham-lemma-base-change-Garel-upstairs}, on peut travailler localement
sur $Y'/X'$ et $Y/X$. On peut donc supposer qu'il existe un morphisme
$Y/X \to Y_{n, d}/X_{n, d}$. Par construction, la famille
$c^p_{Y'/X'}$ et la famille $c^p_{X/Y}$ sont toutes deux compatibles avec
le morphisme vers $Y_{n, d}/X_{n, d}$. Il en résulte aisément
(en utilisant l'unicité du lemme \ref{derham-lemma-base-change-Garel-upstairs})
que $c^p_{Y'/X'}$ est compatible avec $c^p_{Y/X}$.
Il reste donc à démontrer l'assertion.

\medskip\noindent
Le reste de la démonstration est consacré à l'assertion. On peut choisir un point
$y \in Y$ et montrer que les morphismes coïncident dans un voisinage ouvert de $y$.
On peut donc remplacer $Y_1$ et $Y_2$ par des voisinages ouverts de
l'image de $y$ dans $Y_1$ et $Y_2$. On peut ainsi supposer
$Y, X, Y_1, X_1, Y_2, X_2$ affines, et les écrire comme spectres
d'anneaux $B, A, B_1, A_1, B_2, A_2$. Voici le diagramme
$$
\xymatrix{
B_1 \ar[r] & B & B_2 \ar[l] \\
A_1 \ar[u] \ar[r] & A \ar[u] & A_2 \ar[l] \ar[u]
}
$$
Par hypothèse, le spectre de $B$ est un ouvert affine
du spectre de $A \otimes_{A_1} B_1$ comme de celui de $A \otimes_{A_2} B_2$.
En restreignant encore les ouverts, on peut supposer qu'il existe des éléments
$g_i \in A \otimes_{A_i} B_i$ tels que nos morphismes donnent des isomorphismes
$(A \otimes_{A_i} B_i)_{g_i} = B$ ; voir
Propriétés, lemme \ref{properties-lemma-standard-open-two-affines}.
Soient $x_\alpha$ et $y_\beta$ des familles de
variables suffisamment grandes pour que l'on puisse choisir des surjections
$$
A'_1 = A_1[x_\alpha] \to A
\quad\text{et}\quad
A'_2 = A_2[y_\beta] \to A
$$
de $A_1$-algèbres et de $A_2$-algèbres, respectivement. On peut alors choisir des relèvements
$h_i \in A'_i \otimes_{A_i} B_i$ de $g_i \in A \otimes_{A_i} B_i$,
et considérer le diagramme
$$
\xymatrix{
(A'_1 \otimes_{A_1} B_1)_{h_1} \ar[r] & B &
(A'_2 \otimes_{A_1} B_2)_{h_2} \ar[l] \\
A'_1 \ar[u] \ar[r] & A \ar[u] & A'_2 \ar[l] \ar[u]
}
$$
Par construction, les deux carrés sont cocartésiens. Considérons ensuite le
morphisme d'anneaux
$$
A' = A'_1 \times_A A'_2 \longrightarrow
B' =
(A'_1 \otimes_{A_1} B_1)_{h_1} \times_B
(A'_2 \otimes_{A_1} B_2)_{h_2}
$$
D'après
Compléments d'algèbre, lemme \ref{more-algebra-lemma-module-over-fibre-product},
on a $A'_1 \otimes_{A'} B' = (A'_1 \otimes_{A_1} B_1)_{h_1}$
et $A'_2 \otimes_{A'} B' = (A'_2 \otimes_{A_2} B_2)_{h_2}$. En particulier,
les fibres du morphisme $Y' = \Spec(B') \to \Spec(A') = X'$ sont
des sous-schémas ouverts de changements de base des fibres des morphismes
$Y_i \to X_i$ ; elles sont donc finies et sont des intersections complètes
locales (voir
Discriminants, lemme \ref{discriminant-lemma-syntomic-quasi-finite}).
D'après Compléments d'algèbre, lemme
\ref{more-algebra-lemma-properties-algebras-over-fibre-product},
le morphisme d'anneaux $A' \to B'$ est plat et de présentation finie.
Ainsi, Discriminants, lemme
\ref{discriminant-lemma-syntomic-quasi-finite}, partie (3),
montre que $Y'/X'$ est un objet de $\mathcal{C}$.
Considérons maintenant le diagramme commutatif
$$
\xymatrix{
& & Y/X \ar[ld] \ar@/_2pc/[lldd]_{b_1/a_1} \ar[rd] \ar@/^2pc/[rrdd]^{b_2/a_2} \\
& Y'_1/X'_1 \ar[rd] \ar[ld] & & Y'_2/X'_2 \ar[ld] \ar[rd] \\
Y_1/X_1 & & Y'/X' & & Y_2/X_2
}
$$
où $Y'_i/X'_i$ correspond à $A'_i \to (A'_i \otimes_{A_i} B_i)_{h_i}$.
Notons que $Y'_i/X'_i$ est un objet de $\mathcal{C}_{nice}$, car il
s'obtient en prenant un sous-schéma ouvert du produit d'un
espace affine de dimension infinie avec $Y_i/X_i$ ; nous omettons ce détail.
En particulier, l'image inverse de $c^p_{Y_i/X_i}$ par $(b_i/a_i)$ coïncide
avec celle de $c^p_{Y'_i/X'_i}$ par $Y/X \to Y'_i/X'_i$.
On pourrait conclure si $Y'/X'$ était un objet de $\mathcal{C}_{nice}$,
mais ce n'est presque jamais le cas. En effet, en prenant les images inverses de $c^p_{Y'/X'}$
le long des deux chemins du carré, on obtiendrait la conclusion voulue.
Pour contourner le fait que $Y'/X'$ n'appartient pas à $\mathcal{C}_{nice}$,
notons que les arguments précédents montrent que, quitte à restreindre tous
les schémas $X, Y, X'_1, Y'_1, X'_2, Y'_2, X', Y'$ à des ouverts, on peut trouver
$n, d \geq 1$ et prolonger le diagramme comme suit :
$$
\xymatrix{
& Y/X \ar[ld] \ar[rd] \\
Y'_1/X'_1 \ar[rd] & & Y'_2/X'_2 \ar[ld] \\
& Y'/X' \ar[d] \\
& Y_{n, d}/X_{n, d}
}
$$
On peut alors reprendre l'argument déjà donné en prenant
l'image inverse de $c^p_{Y_{n, d}/X_{n, d}}$. Cela achève la démonstration.
\end{proof}








\section{Morphismes trace sur les complexes de de Rham}
\label{derham-section-trace}

\noindent
Une référence pour une partie des résultats de cette section est \cite{Garel}.
Soit $S$ un schéma. Soit $f : Y \to X$ un morphisme fini localement libre
de schémas sur $S$. Il existe alors un morphisme trace
$\text{Trace}_f : f_*\mathcal{O}_Y \to \mathcal{O}_X$ ; voir
Discriminants, section \ref{discriminant-section-discriminant}.
Dans cette situation, un morphisme trace sur les complexes de de Rham est un morphisme
de complexes
$$
\Theta_{Y/X} : f_*\Omega^\bullet_{Y/S} \longrightarrow \Omega^\bullet_{X/S}
$$
tel que $\Theta_{Y/X}$ soit égal à $\text{Trace}_f$ en degré $0$
et vérifie
$$
\Theta_{Y/X}(\omega \wedge \eta) = \omega \wedge \Theta_{Y/X}(\eta)
$$
pour les sections locales $\omega$ de $\Omega^\bullet_{X/S}$ et $\eta$
de $f_*\Omega^\bullet_{Y/S}$. Nous ne savons pas si un tel morphisme trace
$\Theta_{Y/X}$ existe pour tout morphisme fini localement libre $Y \to X$ ;
veuillez écrire à
\href{mailto:stacks.project@gmail.com}{stacks.project@gmail.com}
si vous disposez d'un contre-exemple ou d'une démonstration.

\begin{example}
\label{derham-example-no-trace}
Voici un exemple où il n'existe pas de morphisme trace sur les complexes de de Rham.
Considérons la $\mathbf{C}$-algèbre $B = \mathbf{C}[x, y]$ munie de
l'action de $G = \{\pm 1\}$ définie par $x \mapsto -x$ et $y \mapsto -y$.
Les invariants $A = B^G$ forment un anneau intègre normal de type fini sur $\mathbf{C}$,
engendré par $x^2, xy, y^2$. Nous affirmons que, pour l'inclusion $A \subset B$,
il n'existe pas de morphisme trace raisonnable
$\Omega_{B/\mathbf{C}} \to \Omega_{A/\mathbf{C}}$
sur les formes de degré $1$. Considérons en effet l'élément
$\omega = x \text{d} y \in \Omega_{B/\mathbf{C}}$.
Puisque $\omega$ est invariant sous l'action de $G$, si un morphisme trace « raisonnable »
existait, $2\omega$ devrait appartenir à l'image de
$\Omega_{A/\mathbf{C}} \to \Omega_{B/\mathbf{C}}$. Ce n'est
pas le cas : il est impossible d'écrire $2\omega$ comme combinaison linéaire
de $\text{d}(x^2)$, $\text{d}(xy)$ et $\text{d}(y^2)$,
même à coefficients dans $B$.
Cet exemple contredit le théorème principal de
\cite{Zannier}.
\end{example}

\begin{lemma}
\label{derham-lemma-Garel}
Il existe une unique règle qui, à tout morphisme syntomique fini
de schémas $f : Y \to X$, associe des morphismes de $\mathcal{O}_X$-modules
$$
\Theta^p_{Y/X} :
f_*\Omega^p_{Y/\mathbf{Z}}
\longrightarrow
\Omega^p_{X/\mathbf{Z}}
$$
satisfaisant aux propriétés suivantes :
\begin{enumerate}
\item la composée avec
$\Omega^p_{X/\mathbf{Z}} \otimes_{\mathcal{O}_X} f_*\mathcal{O}_Y
\to f_*\Omega^p_{Y/\mathbf{Z}}$ est égale à
$\text{id} \otimes \text{Trace}_f$,
où $\text{Trace}_f : f_*\mathcal{O}_Y \to \mathcal{O}_X$
est le morphisme de
Discriminants, section \ref{discriminant-section-discriminant},
\item la règle est compatible avec le changement de base.
\end{enumerate}
\end{lemma}

\begin{proof}
Supposons d'abord $X$ localement noethérien. D'après
le lemme \ref{derham-lemma-Garel-upstairs}, on dispose d'un morphisme canonique
$$
c^p_{Y/X} : \Omega_{Y/\mathbf{Z}}^p
\longrightarrow
f^*\Omega_{X/\mathbf{Z}}^p \otimes_{\mathcal{O}_Y} \det(\NL_{Y/X})
$$
D'après Discriminants, proposition \ref{discriminant-proposition-tate-map},
on dispose d'un isomorphisme canonique
$$
c_{Y/X} : \det(\NL_{Y/X}) \to \omega_{Y/X}
$$
envoyant $\delta(\NL_{Y/X})$ sur $\tau_{Y/X}$. En combinant ces morphismes, on obtient
$$
c^p_{Y/X} \otimes c_{Y/X} :
\Omega_{Y/\mathbf{Z}}^p
\longrightarrow
f^*\Omega_{X/\mathbf{Z}}^p \otimes_{\mathcal{O}_Y} \omega_{Y/X}
$$
D'après Discriminants, section \ref{discriminant-section-finite-morphisms},
cela équivaut à la donnée d'un morphisme
$$
\Theta_{Y/X}^p :
f_*\Omega_{Y/\mathbf{Z}}^p
\longrightarrow
\Omega_{X/\mathbf{Z}}^p
$$
Rappelons que la relation entre $c^p_{Y/X} \otimes c_{Y/X}$
et $\Theta_{Y/X}^p$ utilise le morphisme d'évaluation
$f_*\omega_{Y/X} \to \mathcal{O}_X$
qui envoie $\tau_{Y/X}$ sur $\text{Trace}_f(1)$ ; voir
Discriminants, section \ref{discriminant-section-finite-morphisms}.
La propriété (1) est donc satisfaite. La propriété (2) l'est pour les changements de base par
$X' \to X$ avec $X'$ localement noethérien, car $c^p_{Y/X}$ et
$c_{Y/X}$ sont tous deux compatibles avec ces changements de base. Pour $f : Y \to X$
syntomique fini et $X$ localement noethérien,
nous continuerons à noter $\Theta^p_{Y/X}$ la solution ainsi obtenue.

\medskip\noindent
Unicité. Supposons donné un morphisme syntomique fini
$f: Y \to X$ tel que $X$ soit lisse sur $\Spec(\mathbf{Z})$
et que $f$ soit étale au-dessus d'un ouvert dense de $X$. Nous affirmons que,
dans ce cas, $\Theta^p_{Y/X}$ est déterminé de manière unique par la propriété (1).
Considérons en effet les morphismes
$$
\Omega^p_{X/\mathbf{Z}} \otimes_{\mathcal{O}_X} f_*\mathcal{O}_Y \to
f_*\Omega^p_{Y/\mathbf{Z}} \to
\Omega^p_{X/\mathbf{Z}}
$$
Le faisceau $\Omega^p_{X/\mathbf{Z}}$ est sans torsion (par l'hypothèse de
lissité) ; il suffit donc de vérifier que la restriction de
$\Theta^p_{Y/X}$ est déterminée de manière unique sur l'ouvert dense au-dessus
duquel $f$ est étale, autrement dit, on peut supposer $f$ étale.
Or, si $f$ est étale, alors
$f^*\Omega_{X/\mathbf{Z}} = \Omega_{Y/\mathbf{Z}}$,
de sorte que la première flèche de l'équation affichée est un isomorphisme.
Puisque la composée est déterminée, cela garantit l'unicité.

\medskip\noindent
Soit $f : Y \to X$ un morphisme syntomique fini de schémas localement noethériens.
Soit $x \in X$. D'après Discriminants, lemme
\ref{discriminant-lemma-locally-comes-from-universal-finite},
on peut trouver $d \geq 1$ et un diagramme commutatif
$$
\xymatrix{
Y \ar[d] &
V \ar[d] \ar[l] \ar[r] &
V_d \ar[d] \\
X &
U \ar[l] \ar[r] &
U_d
}
$$
tel que $x \in U \subset X$ soit ouvert, $V = f^{-1}(U)$
et $V = U \times_{U_d} V_d$. Ainsi, $\Theta^p_{Y/X}|_V$
est l'image inverse du morphisme $\Theta^p_{V_d/U_d}$.
Or, d'après la discussion précédente sur l'unicité et
Discriminants, lemmes
\ref{discriminant-lemma-universal-finite-syntomic-smooth} et
\ref{discriminant-lemma-universal-finite-syntomic-etale},
le morphisme $\Theta^p_{V_d/U_d}$ est déterminé de manière unique
par la condition (1). L'unicité est donc établie.

\medskip\noindent
Nous avons maintenant établi l'existence et l'unicité
pour tout morphisme syntomique fini $Y \to X$ avec $X$ localement noethérien.
On pourrait donner un argument analogue à celui de la démonstration du
lemme \ref{derham-lemma-Garel-upstairs} pour passer à $X$ quelconque.
On peut toutefois utiliser directement
l'approximation noethérienne absolue pour conclure.
En effet, pour construire $\Theta^p_{Y/X}$,
il suffit de le faire Zariski-localement sur $X$ (à condition de
démontrer aussi l'unicité). On peut donc supposer $X$ affine (nous omettons
ce détail). On peut alors écrire $X = \lim_{i \in I} X_i$
comme limite de schémas affines noethériens indexée par un ensemble ordonné filtrant $I$.
D'après Algèbre, lemme \ref{algebra-lemma-colimit-category-fp-algebras},
on peut trouver $0 \in I$ et un morphisme de présentation
finie entre schémas affines $f_0 : Y_0 \to X_0$ dont le changement de base à
$X$ est $Y \to X$. En augmentant $0$, on peut supposer $Y_0 \to X_0$
fini et syntomique ; voir
Algèbre, lemmes \ref{algebra-lemma-colimit-lci} et
\ref{algebra-lemma-colimit-finite}. Pour $i \geq 0$, le
changement de base $f_i : Y_i = Y_0 \times_{X_0} X_i \to X_i$ est aussi syntomique fini.
Alors
$$
\Gamma(X, f_*\Omega^p_{Y/\mathbf{Z}}) =
\Gamma(Y, \Omega^p_{Y/\mathbf{Z}}) =
\colim_{i \geq 0} \Gamma(Y_i, \Omega^p_{Y_i/\mathbf{Z}}) =
\colim_{i \geq 0} \Gamma(X_i, f_{i, *}\Omega^p_{Y_i/\mathbf{Z}})
$$
On peut donc, et l'on doit, définir $\Theta^p_{Y/X}$ comme la limite inductive
des morphismes $\Theta^p_{Y_i/X_i}$. Ce morphisme est compatible avec tout
diagramme cartésien
$$
\xymatrix{
Y' \ar[r] \ar[d] & Y \ar[d] \\
X' \ar[r] & X
}
$$
où $X'$ est affine, car nous le savons dans le cas des schémas affines noethériens
d'après les arguments précédents (nous omettons ce détail ; indication : si l'on
écrit aussi $X' = \lim_{j \in J} X'_j$, alors, pour tout $i \in I$, il existe $j \in J$
et un morphisme $X'_j \to X_i$ compatible avec le morphisme $X' \to X$).
Cela achève la démonstration.
\end{proof}

\begin{proposition}
\label{derham-proposition-Garel}
\begin{reference}
\cite{Garel}
\end{reference}
Soit $f : Y \to X$ un morphisme syntomique fini de schémas.
Les morphismes $\Theta^p_{Y/X}$ du lemme \ref{derham-lemma-Garel} définissent un morphisme de complexes
$$
\Theta_{Y/X} :
f_*\Omega^\bullet_{Y/\mathbf{Z}}
\longrightarrow
\Omega^\bullet_{X/\mathbf{Z}}
$$
possédant les propriétés suivantes :
\begin{enumerate}
\item en degré $0$, on obtient
$\text{Trace}_f : f_*\mathcal{O}_Y \to \mathcal{O}_X$ ; voir
Discriminants, section \ref{discriminant-section-discriminant},
\item on a
$\Theta_{Y/X}(\omega \wedge \eta) = \omega \wedge \Theta_{Y/X}(\eta)$
pour $\omega$ dans $\Omega^\bullet_{X/\mathbf{Z}}$ et $\eta$
dans $f_*\Omega^\bullet_{Y/\mathbf{Z}}$,
\item si $f$ est un morphisme sur un schéma de base $S$,
$\Theta_{Y/X}$ induit un morphisme de complexes
$f_*\Omega^\bullet_{Y/S} \to \Omega^\bullet_{X/S}$.
\end{enumerate}
\end{proposition}

\begin{proof}
D'après Discriminants, lemme
\ref{discriminant-lemma-locally-comes-from-universal-finite},
pour tout $x \in X$, on peut trouver $d \geq 1$ et un diagramme commutatif
$$
\xymatrix{
Y \ar[d] &
V \ar[d] \ar[l] \ar[r] &
V_d \ar[d] \ar[r] &
Y_d = \Spec(B_d) \ar[d] \\
X &
U \ar[l] \ar[r] &
U_d \ar[r] &
X_d = \Spec(A_d)
}
$$
tel que $x \in U \subset X$ soit un ouvert affine, $V = f^{-1}(U)$
et $V = U \times_{U_d} V_d$. Écrivons $U = \Spec(A)$ et $V = \Spec(B)$,
observons que $B = A \otimes_{A_d} B_d$, et rappelons que
$B_d = A_d e_1 \oplus \ldots \oplus A_d e_d$. Supposons donnés
$a_1, \ldots, a_r \in A$ et $b_1, \ldots, b_s \in B$.
On peut écrire $b_j = \sum a_{j, l} e_l$ avec $a_{j, l} \in A$.
Posons $N = r + sd$ et considérons les factorisations
$$
\xymatrix{
V \ar[r] \ar[d] &
V' = \mathbf{A}^N \times V_d \ar[r] \ar[d] &
V_d \ar[d] \\
U \ar[r]&
U' = \mathbf{A}^N \times U_d \ar[r] &
U_d
}
$$
Ici, la flèche horizontale inférieure droite est donnée par le morphisme
$U \to U_d$ (du diagramme précédent) et le morphisme
$U \to \mathbf{A}^N$ défini par $a_1, \ldots, a_r, a_{1, 1}, \ldots, a_{s, d}$.
On voit alors que les fonctions $a_1, \ldots, a_r$ appartiennent à l'image de
$\Gamma(U', \mathcal{O}_{U'}) \to \Gamma(U, \mathcal{O}_U)$,
et que les fonctions $b_1, \ldots, b_s$ appartiennent à l'image de
$\Gamma(V', \mathcal{O}_{V'}) \to \Gamma(V, \mathcal{O}_V)$.
On voit ainsi que, pour toute famille finie d'éléments\footnote{Ces
éléments sont en effet des sommes finies d'éléments de la forme
$a_0 \text{d}a_1 \wedge \ldots \wedge \text{d}a_i$ avec
$a_0, \ldots, a_i \in A$, ou des sommes finies d'éléments de la forme
$b_0 \text{d}b_1 \wedge \ldots \wedge \text{d}b_j$ avec
$b_0, \ldots, b_j \in B$.} des groupes
$$
\Gamma(V, \Omega^i_{Y/\mathbf{Z}}),\quad i = 0, 1, 2, \ldots
\quad\text{et}\quad
\Gamma(U, \Omega^j_{X/\mathbf{Z}}),\quad j = 0, 1, 2, \ldots
$$,
on peut trouver des factorisations $V \to V' \to V_d$ et
$U \to U' \to U_d$ avec $V' = \mathbf{A}^N \times V_d$ et
$U' = \mathbf{A}^N \times U_d$ comme ci-dessus,
telles que ces sections soient les images inverses de sections de
$$
\Gamma(V', \Omega^i_{V'/\mathbf{Z}}),\quad i = 0, 1, 2, \ldots
\quad\text{et}\quad
\Gamma(U', \Omega^j_{U'/\mathbf{Z}}),\quad j = 0, 1, 2, \ldots
$$
Il en résulte que, pour vérifier
$\text{d} \circ \Theta_{Y/X} = \Theta_{Y/X} \circ \text{d}$,
il suffit de vérifier cette égalité pour $\Theta_{V'/U'}$.
Il en va de même de la propriété (2) du lemme.

\medskip\noindent
D'après Discriminants, lemmes
\ref{discriminant-lemma-universal-finite-syntomic-smooth} et
\ref{discriminant-lemma-universal-finite-syntomic-etale},
le schéma $U_d$ est lisse et le morphisme $V_d \to U_d$
est étale au-dessus d'un ouvert dense de $U_d$.
Il en va donc de même du morphisme
$V' \to U'$. Puisque $\Omega_{U'/\mathbf{Z}}$ est localement libre et donc
$\Omega^p_{U'/\mathbf{Z}}$ sans
torsion, il suffit de vérifier les relations voulues
après restriction à l'ouvert au-dessus duquel $V'$ est fini étale.
On peut ensuite vérifier ces relations après un changement de base étale
surjectif. On peut donc décomposer le revêtement fini étale
et supposer que l'on considère un morphisme de la forme
$$
\coprod\nolimits_{i = 1, \ldots, d} W \longrightarrow W
$$
où $W$ est lisse sur $\mathbf{Z}$.
Dans ce cas, toutes les propriétés locales de notre construction se vérifient trivialement
(pourvu qu'elles soient vraies). Cela achève la démonstration de (1) et (2).

\medskip\noindent
Enfin, (3) résulte de (2), car $\Omega_{Y/S}$
est le quotient de $\Omega_{Y/\mathbf{Z}}$ par le sous-module
engendré par les images inverses de sections locales de $\Omega_{S/\mathbf{Z}}$.
\end{proof}

\begin{example}
\label{derham-example-Garel}
Soit $A$ un anneau. Soit $f = x^d + \sum_{0 \leq i < d} a_{d - i} x^i \in A[x]$.
Soit $B = A[x]/(f)$. D'après la proposition \ref{derham-proposition-Garel},
on dispose d'un morphisme de complexes
$$
\Theta_{B/A} : \Omega^\bullet_B \longrightarrow \Omega^\bullet_A
$$
En particulier, si $t \in B$ désigne l'image de $x \in A[x]$,
on peut considérer les éléments
$$
\Theta_{B/A}(t^i\text{d}t) \in \Omega^1_A,\quad i = 0, \ldots, d - 1
$$
Quels sont ces éléments ? Par le même principe que dans la démonstration de la
proposition \ref{derham-proposition-Garel}, il suffit de les calculer
dans le cas universel, c'est-à-dire lorsque $A = \mathbf{Z}[a_1, \ldots, a_d]$,
ou même après remplacement de $A$ par le corps des fractions
$\mathbf{Q}(a_1, \ldots, a_d)$. En écrivant symboliquement
$$
f = \prod\nolimits_{i = 1, \ldots, d} (x - \alpha_i)
$$,
on voit que, sur $\mathbf{Q}(\alpha_1, \ldots, \alpha_d)$,
l'algèbre $B$ se décompose :
$$
\mathbf{Q}(a_0, \ldots, a_{d - 1})[x]/(f) 
\longrightarrow
\prod\nolimits_{i = 1, \ldots, d} \mathbf{Q}(\alpha_1, \ldots, \alpha_d),
\quad
t \longmapsto (\alpha_1, \ldots, \alpha_d)
$$
Ainsi, par exemple,
$$
\Theta(\text{d}t) = \sum \text{d} \alpha_i = - \text{d}a_1
$$
On a ensuite
$$
\Theta(t\text{d}t) = \sum \alpha_i \text{d}\alpha_i =
a_1 \text{d} a_1 - \text{d}a_2
$$
Puis on a
$$
\Theta(t^2\text{d}t) = \sum \alpha_i^2 \text{d}\alpha_i =
- a_1^2 \text{d} a_1 + a_1 \text{d}a_2 + a_2 \text{d}a_1 - \text{d}a_3
$$
(sauf erreur de calcul), et ainsi de suite. Cela suggère que,
si $f(x) = x^d - a$, alors
$$
\Theta_{B/A}(t^i\text{d}t) =
\left\{
\begin{matrix}
0 & \text{if} & i = 0, \ldots, d - 2 \\
\text{d}a & \text{if} & i = d - 1
\end{matrix}
\right.
$$
dans $\Omega_A$. C'est vrai, car, dans ce cas particulier, on peut effectuer
le calcul pour l'extension $\mathbf{Q}(a)[x]/(x^d - a)$
afin de le vérifier directement.
\end{example}

\begin{lemma}
\label{derham-lemma-Garel-map-frobenius-smooth-char-p}
Soit $p$ un nombre premier. Soit $X \to S$ un morphisme lisse
de dimension relative $d$ de schémas de caractéristique $p$.
Le Frobenius relatif $F_{X/S} : X \to X^{(p)}$ de $X/S$
(Variétés, définition \ref{varieties-definition-relative-frobenius})
est syntomique fini, et le morphisme correspondant
$$
\Theta_{X/X^{(p)}} :
F_{X/S, *}\Omega^\bullet_{X/S} \to \Omega^\bullet_{X^{(p)}/S}
$$
est nul en tout degré sauf en degré $d$, où il définit une
surjection.
\end{lemma}

\begin{proof}
Observons que $F_{X/S}$ est un morphisme fini d'après
Variétés, lemme \ref{varieties-lemma-relative-frobenius-finite}.
Pour montrer que $F_{X/S}$ est plat, il suffit de montrer que
le morphisme $F_{X/S, s} : X_s \to X^{(p)}_s$ entre fibres
est plat pour tout $s \in S$ ; voir Compléments sur les morphismes, théorème
\ref{more-morphisms-theorem-criterion-flatness-fibre}.
La platitude de $X_s \to X^{(p)}_s$ résulte de
Algèbre, lemme \ref{algebra-lemma-CM-over-regular-flat}
(et de la finitude déjà établie).
D'après Compléments sur les morphismes, lemme
\ref{more-morphisms-lemma-lci-permanence},
le morphisme $F_{X/S}$ est localement d'intersection complète.
Ainsi, $F_{X/S}$ est syntomique fini (voir
Compléments sur les morphismes, lemme \ref{more-morphisms-lemma-flat-lci}).

\medskip\noindent
Pour tout point $x \in X$, on peut choisir un diagramme commutatif
$$
\xymatrix{
X \ar[d] & U \ar[l] \ar[d]_\pi \\
S & \mathbf{A}^d_S \ar[l]
}
$$
où $\pi$ est étale et où $x \in U$ est un ouvert de $X$ ; voir
Morphismes, lemme \ref{morphisms-lemma-smooth-etale-over-affine-space}.
Observons que
$\mathbf{A}^d_S \to \mathbf{A}^d_S$, $(x_1, \ldots, x_d) \mapsto
(x_1^p, \ldots, x_d^p)$, est le Frobenius relatif de $\mathcal{A}^d_S$
sur $S$. Le diagramme commutatif
$$
\xymatrix{
U \ar[d]_\pi \ar[r]_{F_{X/S}} & U^{(p)} \ar[d]^{\pi^{(p)}} \\
\mathbf{A}^d_S \ar[r]^{x_i \mapsto x_i^p} & \mathbf{A}^d_S
}
$$
de
Variétés, lemme \ref{varieties-lemma-relative-frobenius-endomorphism-identity},
pour $\pi : U \to \mathbf{A}^d_S$ est cartésien d'après
Morphismes étales, lemme
\ref{etale-lemma-relative-frobenius-etale}.
Puisque la construction de $\Theta$ est compatible avec le changement de base
et que $\Omega_{U/S} = \pi^*\Omega_{\mathbf{A}^d_S/S}$
(lemme \ref{derham-lemma-etale}),
il suffit de démontrer le lemme pour
$\mathbf{A}^d_S$.

\medskip\noindent
Soit $A$ un anneau de caractéristique $p$. Considérons l'unique homomorphisme de $A$-algèbres
$A[y_1, \ldots, y_d] \to A[x_1, \ldots, x_d]$
envoyant $y_i$ sur $x_i^p$. Les arguments précédents
nous ramènent au calcul du morphisme
$$
\Theta^i : \Omega^i_{A[x_1, \ldots, x_d]/A} \to
\Omega^i_{A[y_1, \ldots, y_d]/A}
$$
Nous invitons vivement le lecteur à effectuer lui-même ce calcul.
Comme dans l'exemple \ref{derham-example-Garel}, on peut se ramener au calcul
d'une formule pour $\Theta^i$ dans le cas universel
$$
\mathbf{Z}[y_1, \ldots, y_d] \to \mathbf{Z}[x_1, \ldots, x_d],\quad
y_i \mapsto x_i^p
$$
On peut à son tour trouver la formule pour $\Theta^i$ en calculant dans le cas
complexe, c'est-à-dire pour le morphisme de $\mathbf{C}$-algèbres
$$
\mathbf{C}[y_1, \ldots, y_d] \to \mathbf{C}[x_1, \ldots, x_d],\quad
y_i \mapsto x_i^p
$$
On peut même inverser $x_1, \ldots, x_d$ et $y_1, \ldots, y_d$.
Dans ce cas, on a $\text{d}x_i = p^{-1} x_i^{- p + 1}\text{d}y_i$.
On voit donc que
\begin{align*}
\Theta^i(
x_1^{e_1} \ldots x_d^{e_d} \text{d}x_1 \wedge \ldots \wedge \text{d}x_i)
& =
p^{-i} \Theta^i(
x_1^{e_1 - p + 1} \ldots x_i^{e_i - p + 1} x_{i + 1}^{e_{i + 1}} \ldots
x_d^{e_d} \text{d}y_1 \wedge \ldots \wedge \text{d}y_i ) \\
& =
p^{-i} \text{Trace}(x_1^{e_1 - p + 1} \ldots x_i^{e_i - p + 1}
x_{i + 1}^{e_{i + 1}} \ldots x_d^{e_d})
\text{d}y_1 \wedge \ldots \wedge \text{d}y_i
\end{align*}
d'après les propriétés de $\Theta^i$. Un calcul élémentaire montre
que la trace figurant dans cette expression est nulle, sauf si
$e_1, \ldots, e_i$ sont congrus à $-1$ modulo $p$
et si $e_{i + 1}, \ldots, e_d$ sont divisibles par $p$.
De plus, dans ce cas, on obtient
$$
p^{d - i} y_1^{(e_1 - p + 1)/p} \ldots y_i^{(e_i - p + 1)/p}
y_{i + 1}^{e_{i + 1}/p} \ldots y_d^{e_d/p}
\text{d}y_1 \wedge \ldots \wedge \text{d}y_i
$$
On en déduit que le résultat est nul en caractéristique $p$, sauf si $d = i$,
auquel cas on obtient toutes les formes de degré $d$.
\end{proof}








\section{Dualité de Poincaré}
\label{derham-section-poincare-duality}

\noindent
Dans cette section, nous démontrons la dualité de Poincaré pour la cohomologie de de Rham
d'un schéma propre et lisse sur un corps. Expliquons d'abord comment elle
se présente pour la cohomologie de Hodge.

\begin{lemma}
\label{derham-lemma-duality-hodge}
Soit $k$ un corps. Soit $X$ un schéma non vide, lisse et propre sur $k$,
équidimensionnel de dimension $d$. Il existe une application $k$-linéaire
$$
t : H^d(X, \Omega^d_{X/k}) \longrightarrow k
$$,
unique à précomposition près par la multiplication par une unité de
$H^0(X, \mathcal{O}_X)$, possédant la propriété suivante : pour tous $p, q$, l'accouplement
$$
H^q(X, \Omega^p_{X/k}) \times H^{d - q}(X, \Omega^{d - p}_{X/k})
\longrightarrow
k, \quad
(\xi, \xi') \longmapsto t(\xi \cup \xi')
$$
est parfait.
\end{lemma}

\begin{proof}
D'après Dualité pour les schémas, lemme \ref{duality-lemma-duality-proper-over-field},
on a $\omega_X^\bullet = \Omega^d_{X/k}[d]$.
Puisque $\Omega_{X/k}$ est localement libre de rang $d$
(Morphismes, lemme \ref{morphisms-lemma-smooth-omega-finite-locally-free}),
on a
$$
\Omega^d_{X/k} \otimes_{\mathcal{O}_X} (\Omega^p_{X/k})^\vee
\cong
\Omega^{d - p}_{X/k}
$$
On obtient donc une application $k$-linéaire $t : H^d(X, \Omega^d_{X/k}) \to k$
pour laquelle l'énoncé est vrai, d'après Dualité pour les schémas, lemme
\ref{duality-lemma-duality-proper-over-field-perfect}.
En particulier, l'accouplement
$H^0(X, \mathcal{O}_X) \times H^d(X, \Omega^d_{X/k}) \to k$
est parfait, ce qui implique que toute application $k$-linéaire
$t' : H^d(X, \Omega^d_{X/k}) \to k$ est de la forme
$\xi \mapsto t(g\xi)$ pour un certain $g \in H^0(X, \mathcal{O}_X)$.
Bien entendu, pour que $t'$ définisse encore une dualité
entre $H^0(X, \mathcal{O}_X)$ et $H^d(X, \Omega^d_{X/k})$,
il faut que $g$ soit une unité. Notons $\langle -, - \rangle_{p, q}$
l'accouplement construit à l'aide de $t$, et $\langle -, - \rangle'_{p, q}$
celui construit à l'aide de $t'$. On a clairement
$$
\langle \xi, \xi' \rangle'_{p, q} =
\langle g\xi, \xi' \rangle_{p, q}
$$
pour $\xi \in H^q(X, \Omega^p_{X/k})$ et
$\xi' \in H^{d - q}(X, \Omega^{d - p}_{X/k})$. Puisque $g$ est une unité, donc
inversible, on voit qu'en utilisant $t'$ au lieu de $t$, on obtient encore des accouplements
parfaits pour tous $p, q$.
\end{proof}

\begin{lemma}
\label{derham-lemma-bottom-part-degenerates}
Soit $k$ un corps. Soit $X$ un schéma lisse et propre sur $k$. L'application
$$
\text{d} : H^0(X, \mathcal{O}_X) \to H^0(X, \Omega^1_{X/k})
$$
est nulle.
\end{lemma}

\begin{proof}
Puisque $X$ est lisse sur $k$, il est géométriquement réduit sur $k$ ; voir
Variétés, lemme \ref{varieties-lemma-smooth-geometrically-normal}.
Ainsi, $H^0(X, \mathcal{O}_X) = \prod k_i$
est un produit fini d'extensions de corps finies
séparables $k_i/k$ ; voir Variétés, lemme
\ref{varieties-lemma-proper-geometrically-reduced-global-sections}.
Il en résulte que $\Omega_{H^0(X, \mathcal{O}_X)/k} = \prod \Omega_{k_i/k} = 0$
(voir par exemple Algèbre, lemme
\ref{algebra-lemma-characterize-separable-algebraic-field-extensions}).
Puisque l'application du lemme se factorise sous la forme
$$
H^0(X, \mathcal{O}_X) \to
\Omega_{H^0(X, \mathcal{O}_X)/k} \to
H^0(X, \Omega_{X/k})
$$
par fonctorialité du complexe de de Rham
(voir la section \ref{derham-section-de-rham-complex}), on conclut.
\end{proof}

\begin{lemma}
\label{derham-lemma-top-part-degenerates}
Soit $k$ un corps. Soit $X$ un schéma lisse et propre sur $k$,
équidimensionnel de dimension $d$. L'application
$$
\text{d} : H^d(X, \Omega^{d - 1}_{X/k}) \to H^d(X, \Omega^d_{X/k})
$$
est nulle.
\end{lemma}

\begin{proof}
Il est tentant de penser que cela résulte de la combinaison des
lemmes \ref{derham-lemma-bottom-part-degenerates} et \ref{derham-lemma-duality-hodge}.
Cela ne fonctionne toutefois pas, car les morphismes $\mathcal{O}_X \to \Omega^1_{X/k}$
et $\Omega^{d - 1}_{X/k} \to \Omega^d_{X/k}$ ne sont pas $\mathcal{O}_X$-linéaires.
On ne peut donc utiliser la fonctorialité décrite dans
Dualité pour les schémas, remarque
\ref{duality-remark-coherent-duality-proper-over-field},
pour conclure que l'application du lemme \ref{derham-lemma-bottom-part-degenerates}
est duale de celle du présent lemme.

\medskip\noindent
On peut remplacer $X$ par une composante connexe de $X$. On peut donc supposer
$X$ irréductible. D'après
Variétés, lemmes \ref{varieties-lemma-smooth-geometrically-normal} et
\ref{varieties-lemma-proper-geometrically-reduced-global-sections},
$k' = H^0(X, \mathcal{O}_X)$ est une extension finie
séparable $k'/k$. Puisque $\Omega_{k'/k} = 0$
(voir par exemple Algèbre, lemme
\ref{algebra-lemma-characterize-separable-algebraic-field-extensions}),
on voit que $\Omega_{X/k} = \Omega_{X/k'}$
(voir Morphismes, lemme \ref{morphisms-lemma-triangle-differentials}).
On peut donc remplacer $k$ par $k'$ et supposer $H^0(X, \mathcal{O}_X) = k$.

\medskip\noindent
Supposons $H^0(X, \mathcal{O}_X) = k$. On en déduit que
$\dim H^d(X, \Omega^d_{X/k}) = 1$ d'après le lemme \ref{derham-lemma-duality-hodge}.
Supposons d'abord que la caractéristique de $k$ soit un nombre premier $p$.
Notons $F_{X/k} : X \to X^{(p)}$ le Frobenius relatif de $X$ sur $k$ ;
gardons à l'esprit les propriétés de ce morphisme établies dans le
lemme \ref{derham-lemma-Garel-map-frobenius-smooth-char-p}.
Considérons le diagramme commutatif
$$
\xymatrix{
H^d(X, \Omega^{d - 1}_{X/k}) \ar[d] \ar[r] &
H^d(X^{(p)}, F_{X/k, *}\Omega^{d - 1}_{X/k}) \ar[d] \ar[r]_{\Theta^{d - 1}} &
H^d(X^{(p)}, \Omega^{d - 1}_{X^{(p)}/k}) \ar[d] \\
H^d(X, \Omega^d_{X/k}) \ar[r] &
H^d(X^{(p)}, F_{X/k, *}\Omega^d_{X/k}) \ar[r]^{\Theta^d} &
H^d(X^{(p)}, \Omega^d_{X^{(p)}/k})
}
$$
Les deux flèches horizontales de gauche sont des isomorphismes, car $F_{X/k}$ est fini ; voir
Cohomologie des schémas, lemme \ref{coherent-lemma-relative-affine-cohomology}.
Le carré de droite est commutatif, car $\Theta_{X^{(p)}/X}$ est un morphisme de
complexes et $\Theta^{d - 1}$ est nul. Il suffit donc de montrer que
$\Theta^d$ est non nul (car la dimension de la source de
$\Theta^d$ est $1$ d'après la discussion précédente). Or on sait que
$$
\Theta^d : F_{X/k, *}\Omega^d_{X/k} \to \Omega^d_{X^{(p)}/k}
$$
est surjectif, donc encore surjectif après application du foncteur exact à
droite $H^d(X^{(p)}, -)$ (l'exactitude à droite résulte de l'annulation de la cohomologie
au-delà de $d$, conséquence de
Cohomologie, proposition \ref{cohomology-proposition-vanishing-Noetherian}).
Enfin, $H^d(X^{(d)}, \Omega^d_{X^{(d)}/k})$ est non nul, par exemple parce
qu'il est dual de $H^0(X^{(d)}, \mathcal{O}_{X^{(p)}})$ d'après le
lemme \ref{derham-lemma-duality-hodge} appliqué à $X^{(p)}$ sur $k$.
Cela achève la démonstration dans ce cas.

\medskip\noindent
Supposons enfin que la caractéristique de $k$ soit $0$.
On peut écrire $k$ comme limite inductive filtrante de ses
$\mathbf{Z}$-sous-algèbres de type fini $R$. Pour l'une d'elles, on peut trouver un
diagramme cartésien de schémas
$$
\xymatrix{
X \ar[d] \ar[r] & Y \ar[d] \\
\Spec(k) \ar[r] & \Spec(R)
}
$$
tel que $Y \to \Spec(R)$ soit lisse de dimension relative $d$ et propre.
Voir Limites, lemmes \ref{limits-lemma-descend-finite-presentation},
\ref{limits-lemma-descend-smooth}, \ref{limits-lemma-descend-dimension-d} et
\ref{limits-lemma-eventually-proper}.
Les modules $M^{i, j} = H^j(Y, \Omega^i_{Y/R})$ sont des $R$-modules de type fini ; voir
Cohomologie des schémas, lemme
\ref{coherent-lemma-proper-over-affine-cohomology-finite}.
Ainsi, en remplaçant $R$ par un localisé, on peut supposer tous ces
modules libres de rang fini. On a
$M^{i, j} \otimes_R k = H^j(X, \Omega^i_{X/k})$
par changement de base plat (Cohomologie des schémas, lemme
\ref{coherent-lemma-flat-base-change-cohomology}).
Il suffit donc de montrer que $M^{d - 1, d} \to M^{d, d}$
est nul. C'est un morphisme entre modules libres de rang fini sur un anneau intègre ;
il suffit donc de trouver un ensemble dense d'idéaux premiers $\mathfrak p \subset R$
tel que le morphisme devienne nul après tensorisation par $\kappa(\mathfrak p)$.
Puisque $R$ est de type fini sur $\mathbf{Z}$, on peut prendre
l'ensemble des idéaux premiers $\mathfrak p$ dont le corps résiduel
est de caractéristique positive (détails omis). Observons que
$$
M^{d - 1, d} \otimes_R \kappa(\mathfrak p) =
H^d(Y_{\kappa(\mathfrak p)},
\Omega^{d - 1}_{Y_{\kappa(\mathfrak p)}/\kappa(\mathfrak p)})
$$,
par exemple d'après Limites, lemme
\ref{limits-lemma-higher-direct-images-zero-above-dimension-fibre}.
Il en va de même pour $M^{d, d}$. On voit donc que
$M^{d - 1, d} \otimes_R \kappa(\mathfrak p) \to
M^{d, d} \otimes_R \kappa(\mathfrak p)$
est nul d'après le cas de caractéristique positive traité ci-dessus.
\end{proof}

\begin{proposition}
\label{derham-proposition-poincare-duality}
Soit $k$ un corps. Soit $X$ un schéma non vide, lisse et propre sur $k$,
équidimensionnel de dimension $d$. Il existe une application $k$-linéaire
$$
t : H^{2d}_{dR}(X/k) \longrightarrow k
$$,
unique à précomposition près par la multiplication par une unité de
$H^0(X, \mathcal{O}_X)$, possédant la propriété suivante : pour tout $i$, l'accouplement
$$
H^i_{dR}(X/k) \times H_{dR}^{2d - i}(X/k)
\longrightarrow
k, \quad
(\xi, \xi') \longmapsto t(\xi \cup \xi')
$$
est parfait.
\end{proposition}

\begin{proof}
La suite spectrale de Hodge vers de Rham
(section \ref{derham-section-hodge-to-de-rham}), l'annulation
de $\Omega^i_{X/k}$ pour $i > d$, le théorème d'annulation de
Cohomologie, proposition \ref{cohomology-proposition-vanishing-Noetherian},
et les résultats des lemmes \ref{derham-lemma-bottom-part-degenerates} et
\ref{derham-lemma-top-part-degenerates}
montrent que $H^0_{dR}(X/k) = H^0(X, \mathcal{O}_X)$
et $H^d(X, \Omega^d_{X/k}) = H_{dR}^{2d}(X/k)$.
Plus précisément, ces identifications proviennent des morphismes
de complexes
$$
\Omega^\bullet_{X/k} \to \mathcal{O}_X[0]
\quad\text{et}\quad
\Omega^d_{X/k}[-d] \to \Omega^\bullet_{X/k}
$$
Choisissons $t : H_{dR}^{2d}(X/k) \to k$ correspondant, par cette identification,
à une application $t$ comme dans le lemme \ref{derham-lemma-duality-hodge}.
On voit alors, en tout cas, que l'accouplement affiché dans le lemme
est parfait pour $i = 0$.

\medskip\noindent
Notons $\underline{k}$ le faisceau constant de valeur $k$ sur $X$.
Abrégeons $\Omega^\bullet = \Omega^\bullet_{X/k}$.
Considérons le morphisme (\ref{derham-equation-wedge}), qui s'écrit ici
$$
\wedge :
\text{Tot}(\Omega^\bullet \otimes_{\underline{k}} \Omega^\bullet)
\longrightarrow
\Omega^\bullet
$$
Pour tout entier $p = 0, 1, \ldots, d$, ce morphisme
annule le sous-complexe
$\text{Tot}(\sigma_{> p} \Omega^\bullet \otimes_{\underline{k}}
\sigma_{\geq d - p} \Omega^\bullet)$ pour des raisons de degré.
Ainsi, la restriction de $\wedge$ au sous-complexe
$\text{Tot}(\Omega^\bullet \otimes_{\underline{k}}
\sigma_{\geq d - p}\Omega^\bullet)$ se factorise par un morphisme de complexes
$$
\gamma_p :
\text{Tot}(\sigma_{\leq p} \Omega^\bullet \otimes_{\underline{k}}
\sigma_{\geq d - p} \Omega^\bullet)
\longrightarrow
\Omega^\bullet
$$
Par le même procédé qu'à la section \ref{derham-section-cup-product}, on obtient des
cup-produits
$$
H^i(X, \sigma_{\leq p} \Omega^\bullet) \times
H^{2d - i}(X, \sigma_{\geq d - p}\Omega^\bullet)
\longrightarrow
H_{dR}^{2d}(X, \Omega^\bullet)
$$
Nous montrerons par récurrence sur $p$ que ces cup-produits, composés avec $t$,
induisent des accouplements parfaits entre $H^i(X, \sigma_{\leq p} \Omega^\bullet)$
et $H^{2d - i}(X, \sigma_{\geq d - p}\Omega^\bullet)$. Pour $p = d$,
c'est l'assertion de la proposition.

\medskip\noindent
Le cas initial est $p = 0$. On obtient alors simplement l'accouplement
entre $H^i(X, \mathcal{O}_X)$ et $H^{d - i}(X, \Omega^d)$ du
lemme \ref{derham-lemma-duality-hodge}, d'où le résultat.

\medskip\noindent
Étape de récurrence. Supposons le résultat vrai pour $p$.
Considérons alors le triangle distingué
$$
\Omega^{p + 1}[-p - 1] \to
\sigma_{\leq p + 1}\Omega^\bullet \to
\sigma_{\leq p}\Omega^\bullet \to
\Omega^{p + 1}[-p]
$$
et le triangle distingué
$$
\sigma_{\geq d - p}\Omega^\bullet \to
\sigma_{\geq d - p - 1}\Omega^\bullet \to
\Omega^{d - p - 1}[-d + p + 1] \to
(\sigma_{\geq d - p}\Omega^\bullet)[1]
$$
Observons que tous deux sont des triangles distingués dans la catégorie homotopique
des complexes de faisceaux de $\underline{k}$-modules ; en particulier, les
morphismes $\sigma_{\leq p}\Omega^\bullet \to \Omega^{p + 1}[-p]$ et
$\Omega^{d - p - 1}[-d + p + 1] \to (\sigma_{\geq d - p}\Omega^\bullet)[1]$
sont donnés par de véritables morphismes de complexes, à l'aide de la différentielle
$\Omega^p \to \Omega^{p + 1}$ et de la différentielle
$\Omega^{d - p - 1} \to \Omega^{d - p}$.
Considérons les suites exactes longues de cohomologie associées à ces
triangles distingués
$$
\xymatrix{
H^{i - 1}(X, \sigma_{\leq p}\Omega^\bullet) \ar[d]_a \\
H^i(X, \Omega^{p + 1}[-p - 1]) \ar[d]_b \\
H^i(X, \sigma_{\leq p + 1}\Omega^\bullet) \ar[d]_c \\
H^i(X, \sigma_{\leq p}\Omega^\bullet) \ar[d]_d \\
H^{i + 1}(X, \Omega^{p + 1}[-p - 1])
}
\quad\quad
\xymatrix{
H^{2d - i  + 1}(X, \sigma_{\geq d - p}\Omega^\bullet) \\
H^{2d - i}(X, \Omega^{d - p - 1}[-d + p + 1]) \ar[u]_{a'} \\
H^{2d - i}(X, \sigma_{\geq d - p - 1}\Omega^\bullet) \ar[u]_{b'} \\
H^{2d - i}(X, \sigma_{\geq d - p}\Omega^\bullet) \ar[u]_{c'} \\
H^{2d - i - 1}(X, \Omega^{d - p - 1}[-d + p + 1]) \ar[u]_{d'}
}
$$
Par récurrence et d'après le lemme \ref{derham-lemma-duality-hodge},
on sait que les accouplements construits ci-dessus entre les
$k$-espaces vectoriels des première, deuxième, quatrième et cinquième
lignes sont parfaits. D'après le lemme des $5$, pour montrer que
l'accouplement entre les groupes de cohomologie de la ligne du milieu
est parfait, il suffit de montrer que les couples
$(a, a')$, $(b, b')$, $(c, c')$ et $(d, d')$
sont compatibles avec les accouplements donnés (voir ci-dessous).

\medskip\noindent
Démontrons-le pour le couple $(c, c')$. Il suffit d'observer
que l'on dispose d'un diagramme commutatif
$$
\xymatrix{
\text{Tot}(\sigma_{\leq p} \Omega^\bullet \otimes_{\underline{k}}
\sigma_{\geq d - p} \Omega^\bullet) \ar[d]_{\gamma_p} &
\text{Tot}(\sigma_{\leq p + 1} \Omega^\bullet \otimes_{\underline{k}}
\sigma_{\geq d - p} \Omega^\bullet) \ar[l] \ar[d] \\
\Omega^\bullet &
\text{Tot}(\sigma_{\leq p + 1} \Omega^\bullet \otimes_{\underline{k}}
\sigma_{\geq d - p - 1} \Omega^\bullet) \ar[l]_-{\gamma_{p + 1}}
}
$$
Ainsi, si l'on a $\alpha \in H^i(X, \sigma_{\leq p + 1}\Omega^\bullet)$
et $\beta \in H^{2d - i}(X, \sigma_{\geq d - p}\Omega^\bullet)$,
on obtient
$\gamma_p(\alpha \cup c'(\beta)) = \gamma_{p + 1}(c(\alpha) \cup \beta)$
par fonctorialité du cup-produit.

\medskip\noindent
De même, pour le couple $(b, b')$, on utilise le diagramme commutatif
$$
\xymatrix{
\text{Tot}(\sigma_{\leq p + 1} \Omega^\bullet \otimes_{\underline{k}}
\sigma_{\geq d - p - 1} \Omega^\bullet) \ar[d]_{\gamma_{p + 1}} &
\text{Tot}(\Omega^{p + 1}[-p - 1] \otimes_{\underline{k}}
\sigma_{\geq d - p - 1} \Omega^\bullet) \ar[l] \ar[d] \\
\Omega^\bullet &
\Omega^{p + 1}[-p - 1]
\otimes_{\underline{k}}
\Omega^{d - p - 1}[-d + p + 1] \ar[l]_-\wedge
}
$$
et l'on raisonne de la même manière.

\medskip\noindent
Pour le couple $(d, d')$, on utilise le diagramme commutatif
$$
\xymatrix{
\Omega^{p + 1}[-p] \otimes_{\underline{k}}
\Omega^{d - p - 1}[-d + p] \ar[d] &
\text{Tot}(\sigma_{\leq p}\Omega^\bullet \otimes_{\underline{k}}
\Omega^{d - p - 1}[-d + p]) \ar[l] \ar[d] \\
\Omega^\bullet &
\text{Tot}(\sigma_{\leq p}\Omega^\bullet \otimes_{\underline{k}}
\sigma_{\geq d - p}\Omega^\bullet) \ar[l]
}
$$
et l'on considère les classes de cohomologie dans
$H^i(X, \sigma_{\leq p}\Omega^\bullet)$ et
$H^{2d - i}(X, \Omega^{d - p - 1}[-d + p])$.
En remplaçant $i$ par $i - 1$, on obtient le résultat pour le couple $(a, a')$,
ce qui achève la démonstration du caractère parfait de nos accouplements.

\medskip\noindent
Nous omettons l'argument établissant l'unicité de $t$ à
précomposition près par la multiplication par une unité de $H^0(X, \mathcal{O}_X)$.
\end{proof}







\section{Classes de Chern}
\label{derham-section-chern-classes}

\noindent
Les résultats établis jusqu'ici suffisent pour appliquer la discussion de
Théories cohomologiques de Weil, section \ref{weil-section-chern},
afin de construire des classes de Chern en cohomologie de de Rham.

\begin{lemma}
\label{derham-lemma-chern-classes}
Il existe une unique règle qui associe à tout schéma quasi-compact et
quasi-séparé $X$ une classe de Chern totale
$$
c^{dR} :
K_0(\textit{Vect}(X))
\longrightarrow
\prod\nolimits_{i \geq 0} H^{2i}_{dR}(X/\mathbf{Z})
$$
possédant les propriétés suivantes :
\begin{enumerate}
\item on a $c^{dR}(\alpha + \beta) = c^{dR}(\alpha) c^{dR}(\beta)$
pour $\alpha, \beta \in K_0(\textit{Vect}(X))$,
\item si $f : X \to X'$ est un morphisme de schémas quasi-compacts et
quasi-séparés, alors $c^{dR}(f^*\alpha) = f^*c^{dR}(\alpha)$,
\item pour $\mathcal{L} \in \Pic(X)$, on a
$c^{dR}([\mathcal{L}]) = 1 + c_1^{dR}(\mathcal{L})$
\end{enumerate}
\end{lemma}

\noindent
La construction s'étend facilement à tous les schémas, mais il faut pour cela
généraliser légèrement la discussion de Théories cohomologiques de Weil,
section \ref{weil-section-chern}.

\begin{proof}
Nous allons appliquer Théories cohomologiques de Weil, proposition
\ref{weil-proposition-chern-class}, pour obtenir ce résultat.

\medskip\noindent
Soit $\mathcal{C}$ la catégorie de tous les schémas quasi-compacts et quasi-séparés.
Elle satisfait assurément aux conditions
(1), (2) et (3) (a), (b) et (c) de Théories cohomologiques de Weil,
section \ref{weil-section-chern}.

\medskip\noindent
Notre foncteur contravariant $A$ de $\mathcal{C}$ vers la
catégorie des algèbres graduées enverra $X$ sur
$A(X) = \bigoplus_{i \geq 0} H_{dR}^{2i}(X/\mathbf{Z})$,
muni de son cup-produit.
La fonctorialité est expliquée à la section \ref{derham-section-de-rham-cohomology}
et le cup-produit à la section \ref{derham-section-cup-product}.
Pour les applications additives $c_1^A$, on prend $c_1^{dR}$, construit
à la section \ref{derham-section-first-chern-class}.

\medskip\noindent
On obtient en fait des algèbres commutatives, d'après
le lemme \ref{derham-lemma-cup-product-graded-commutative},
ce qui établit l'axiome (1) pour $A$.

\medskip\noindent
Pour vérifier l'axiome (2) pour $A$, il suffit de vérifier que
$H^*_{dR}(X \coprod Y/\mathbf{Z}) =  H^*_{dR}(X/\mathbf{Z}) \times
H^*_{dR}(Y/\mathbf{Z})$.
Cela résulte du fait que la cohomologie de de Rham
se construit en prenant la cohomologie d'un faisceau d'algèbres différentielles
graduées (pour la topologie de Zariski).

\medskip\noindent
L'axiome (3) pour $A$ affirme simplement que la formation des premières classes
de Chern des modules inversibles est compatible avec les images inverses.
Cela résulte du lemme plus général \ref{derham-lemma-pullback-c1}.

\medskip\noindent
L'axiome (4) pour $A$ est la formule du fibré projectif,
démontrée dans la proposition \ref{derham-proposition-projective-space-bundle-formula}.

\medskip\noindent
Axiome (5). Soit $X$ un schéma quasi-compact et quasi-séparé, et
soit $\mathcal{E} \to \mathcal{F}$ une surjection de
$\mathcal{O}_X$-modules localement libres de rangs $r + 1$ et $r$. Notons
$i : P' = \mathbf{P}(\mathcal{F}) \to \mathbf{P}(\mathcal{E}) = P$
le morphisme d'inclusion correspondant. C'est un morphisme de schémas projectifs lisses
sur $X$ qui présente $P'$ comme un diviseur de Cartier effectif sur $P$.
Ainsi, d'après le lemme \ref{derham-lemma-check-log-smooth}, le complexe à pôles logarithmiques
pour $P' \subset P$ sur $\mathbf{Z}$ est défini.
Par conséquent, pour $a \in A(P)$ tel que $i^*a = 0$, on a
$a \cup c_1^A(\mathcal{O}_P(P')) = 0$ d'après
le lemme \ref{derham-lemma-log-complex-consequence}.
Cela achève la démonstration.
\end{proof}

\begin{remark}
\label{derham-remark-splitting-principle}
Les analogues de Théories cohomologiques de Weil, lemmes
\ref{weil-lemma-splitting-principle} (principe de décomposition) et
\ref{weil-lemma-chern-classes-E-tensor-L} (classes de Chern des produits tensoriels),
sont valables pour les classes de Chern de de Rham sur les schémas quasi-compacts et quasi-séparés.
C'est clair, puisque nous avons montré dans la démonstration du
lemme \ref{derham-lemma-chern-classes}
que tous les axiomes de Théories cohomologiques de Weil, section
\ref{weil-section-chern}, sont satisfaits.
\end{remark}

\noindent
Pour les schémas sur $\mathbf{Q}$, on peut construire un caractère de Chern.

\begin{lemma}
\label{derham-lemma-chern-character}
Il existe une unique règle qui associe à tout schéma quasi-compact et quasi-séparé
$X$ sur $\mathbf{Q}$ un « caractère de Chern »
$$
ch^{dR} : K_0(\textit{Vect}(X)) \longrightarrow
\prod\nolimits_{i \geq 0} H_{dR}^{2i}(X/\mathbf{Q})
$$
possédant les propriétés suivantes :
\begin{enumerate}
\item $ch^{dR}$ est un morphisme d'anneaux pour tout $X$,
\item si $f : X' \to X$ est un morphisme de schémas quasi-compacts et quasi-séparés
sur $\mathbf{Q}$, alors $f^* \circ ch^{dR} =  ch^{dR} \circ f^*$, et
\item pour $\mathcal{L} \in \Pic(X)$,
on a $ch^{dR}([\mathcal{L}]) = \exp(c_1^{dR}(\mathcal{L}))$.
\end{enumerate}
\end{lemma}

\noindent
La construction s'étend facilement à tous les schémas sur $\mathbf{Q}$,
mais il faut pour cela généraliser légèrement la discussion de
Théories cohomologiques de Weil, section \ref{weil-section-chern}.

\begin{proof}
Exactement comme dans la démonstration du lemme \ref{derham-lemma-chern-classes},
on montre que la catégorie des schémas quasi-compacts et quasi-séparés
sur $\mathbf{Q}$, munie du foncteur
$A^*(X) = \bigoplus_{i \geq 0} H_{dR}^{2i}(X/\mathbf{Q})$,
satisfait aux axiomes de
Théories cohomologiques de Weil, section \ref{weil-section-chern}.
De plus, dans ce cas, $A(X)$ est une $\mathbf{Q}$-algèbre pour
tout $X$. Le lemme résulte donc de
Théories cohomologiques de Weil, proposition
\ref{weil-proposition-chern-character}.
\end{proof}







\section{Une théorie cohomologique de Weil}
\label{derham-section-weil}

\noindent
Soit $k$ un corps de caractéristique $0$. Dans cette section, nous montrons que
le foncteur
$$
X \longmapsto H^*_{dR}(X/k)
$$
définit une théorie cohomologique de Weil sur $k$ à coefficients dans $k$, au sens
de Théories cohomologiques de Weil, définition
\ref{weil-definition-weil-cohomology-theory}.
Nous vérifierons que les constructions précédentes de ce
chapitre fournissent des données (D0), (D1) et (D2') satisfaisant
aux axiomes (A1) -- (A9) de
Théories cohomologiques de Weil, section \ref{weil-section-c1}.

\medskip\noindent
Dans toute la suite de cette section, fixons le corps $k$ de caractéristique
$0$ et posons $F = k$. Prenons ensuite les données suivantes :
\begin{enumerate}
\item[(D0)] Pour l'espace vectoriel de dimension $1$ sur $F$, $F(1)$, prenons
$F(1) = F = k$.
\item[(D1)] Pour le foncteur $H^*$, prenons celui qui envoie
un schéma projectif lisse $X$ sur $k$ sur $H^*_{dR}(X/k)$.
La fonctorialité est expliquée à la section \ref{derham-section-de-rham-cohomology}
et le cup-produit à la section \ref{derham-section-cup-product}.
On obtient des $F$-algèbres graduées commutatives d'après
le lemme \ref{derham-lemma-cup-product-graded-commutative}.
\item[(D2')] Pour les applications $c_1^H : \Pic(X) \to H^2(X)(1)$, on
utilise la première classe de Chern de de Rham introduite à la
section \ref{derham-section-first-chern-class}.
\end{enumerate}
Nous allons montrer que les axiomes (A1) -- (A9) sont satisfaits.

\medskip\noindent
Dans ce paragraphe, nous allons ramener la vérification des
axiomes au cas où $k$ est algébriquement clos, en
utilisant Théories cohomologiques de Weil, lemme \ref{weil-lemma-check-over-extension}.
Notons $k'$ une clôture algébrique de $k$.
Posons $F' = k'$. On obtient des données (D0), (D1), (D2') sur $k'$ à
corps de coefficients $F'$, exactement comme ci-dessus.
D'après le lemme \ref{derham-lemma-proper-smooth-de-Rham}, il existe des
isomorphismes fonctoriels
$$
H_{dR}^{2d}(X/k) \otimes_k k'
\longrightarrow
H_{dR}^{2d}(X_{k'}/k')
$$
pour $X$ lisse et projectif sur $k$. De plus, les diagrammes
$$
\xymatrix{
\Pic(X) \ar[r]_{c^{dR}_1} \ar[d] & H_{dR}^2(X/k) \ar[d] \\
\Pic(X_{k'}) \ar[r]^{c^{dR}_1} & H_{dR}^2(X_{k'}/k')
}
$$
sont commutatifs d'après le lemme \ref{derham-lemma-pullback-c1}.
Cela achève la réduction.

\medskip\noindent
Supposons $k$ algébriquement clos de caractéristique nulle.
Nous allons établir les axiomes (A1) -- (A9) pour les données (D0), (D1) et (D2')
ci-dessus.

\medskip\noindent
Axiome (A1). Il faut vérifier que
$H^*_{dR}(X \coprod Y/k) =  H^*_{dR}(X/k) \times H^*_{dR}(Y/k)$.
Cela résulte du fait que la cohomologie de de Rham
se construit en prenant la cohomologie d'un faisceau d'algèbres différentielles
graduées (pour la topologie de Zariski).

\medskip\noindent
Axiome (A2). Il affirme simplement que la formation des premières classes
de Chern des modules inversibles est compatible avec les images inverses.
Cela résulte du lemme plus général \ref{derham-lemma-pullback-c1}.

\medskip\noindent
Axiome (A3). Il résulte de la proposition
plus générale \ref{derham-proposition-projective-space-bundle-formula}.

\medskip\noindent
Axiome (A4). Il résulte du lemme
plus général \ref{derham-lemma-log-complex-consequence}.

\medskip\noindent
Dès à présent, en utilisant
Théories cohomologiques de Weil, lemmes \ref{weil-lemma-chern-classes} et
\ref{weil-lemma-cycle-classes}, on obtient un caractère de Chern et des
applications de classe de cycle
$$
\gamma :
\CH^*(X)
\longrightarrow
\bigoplus\nolimits_{i \geq 0} H^{2i}_{dR}(X/k)
$$
pour $X$ projectif lisse sur $k$, qui sont des homomorphismes d'anneaux gradués
compatibles avec les images inverses par les morphismes $f : X \to Y$
de schémas projectifs lisses sur $k$.

\medskip\noindent
Axiome (A5). On a $H_{dR}^*(\Spec(k)/k) = k = F$ en degré $0$.
On dispose de la formule de K\"unneth pour le produit de deux
$k$-schémas projectifs lisses d'après le lemme \ref{derham-lemma-kunneth-de-rham} (observons que les
produits tensoriels dérivés de l'énoncé ne posent pas de difficulté, puisque la
tensorisation s'effectue sur le corps $k$).

\medskip\noindent
Axiome (A7). Il résulte de la proposition \ref{derham-proposition-blowup-split}.

\medskip\noindent
Axiome (A8). Soit $X$ un schéma projectif lisse sur $k$.
Le texte explicatif accompagnant cet axiome dans
Théories cohomologiques de Weil, section \ref{weil-section-c1},
montre que $k' = H^0(X, \mathcal{O}_X)$ est une
$k$-algèbre finie séparable. Il en résulte que $H_{dR}^*(\Spec(k')/k) = k'$,
concentré en degré $0$, puisque $\Omega_{k'/k} = 0$. D'après
le lemme \ref{derham-lemma-bottom-part-degenerates},
on a aussi $H_{dR}^0(X, \mathcal{O}_X) = k'$, ce qui donne
l'axiome.

\medskip\noindent
Axiome (A6). Soit $X$ un schéma projectif lisse non vide sur $k$,
équidimensionnel de dimension $d$. Notons
$\Delta : X \to X \times_{\Spec(k)} X$
le morphisme diagonal de $X$ sur $k$. Il faut montrer qu'il
existe une application $k$-linéaire
$$
\lambda : H_{dR}^{2d}(X/k) \longrightarrow k
$$
telle que $(1 \otimes \lambda)\gamma([\Delta]) = 1$ dans $H^0_{dR}(X/k)$.
Écrivons
$$
\gamma = \gamma([\Delta]) = \gamma_0 + \ldots + \gamma_{2d}
$$,
où $\gamma_i \in H_{dR}^i(X/k) \otimes_k H_{dR}^{2d - i}(X/k)$
sont les composantes de K\"unneth. Il s'agit de montrer qu'il existe une
application linéaire $\lambda : H_{dR}^{2d}(X/k) \to k$ telle que
$(1 \otimes \lambda)\gamma_0 = 1$ dans $H^0_{dR}(X/k)$.

\medskip\noindent
Soit $X = \coprod X_i$ la décomposition de $X$ en composantes connexes,
donc irréductibles. On a alors la décomposition correspondante
$\Delta = \coprod \Delta_i$ avec $\Delta_i \subset X_i \times X_i$.
Il en résulte que
$$
\gamma([\Delta]) = \sum \gamma([\Delta_i])
$$,
et, de plus, $\gamma([\Delta_i])$ correspond à la classe de
$\Delta_i \subset X_i \times X_i$ par la décomposition
$$
H^*_{dR}(X \times X) = \prod\nolimits_{i, j} H^*_{dR}(X_i \times X_j)
$$
Nous omettons les détails ; une façon de le montrer consiste à utiliser, dans
$\CH^0(X \times X)$, les idempotents $e_{i, j}$ correspondant aux
sous-schémas ouverts et fermés $X_i \times X_j$, ainsi que le fait que
$\gamma$ est un morphisme d'anneaux envoyant $e_{i, j}$ sur
l'idempotent correspondant dans la décomposition en produit de la cohomologie affichée ci-dessus.
Si l'on peut trouver $\lambda_i : H_{dR}^{2d}(X_i/k) \to k$ tel que
$(1 \otimes \lambda_i)\gamma([\Delta_i]) = 1$ dans $H^0_{dR}(X_i/k)$,
alors $\lambda = \sum \lambda_i$ résoudra le problème pour $X$.
On peut donc supposer, et l'on suppose désormais, $X$ irréductible.

\medskip\noindent
Démonstration de l'axiome (A6) pour $X$ irréductible. Puisque $k$ est algébriquement
clos, on a $H^0_{dR}(X/k) = k$, car $H^0(X, \mathcal{O}_X) = k$,
$X$ étant une variété projective sur un corps algébriquement clos
(voir par exemple Variétés, lemme
\ref{varieties-lemma-proper-geometrically-reduced-global-sections}).
Soit $x \in X$ un point fermé quelconque.
Considérons le diagramme cartésien
$$
\xymatrix{
x \ar[d] \ar[r] & X \ar[d]^\Delta \\
X \ar[r]^-{x \times \text{id}} & X \times_{\Spec(k)} X
}
$$
La compatibilité de $\gamma$ avec les images inverses implique que
$\gamma([\Delta])$ s'envoie sur $\gamma([x])$ dans $H_{dR}^{2d}(X/k)$ ;
autrement dit, $\gamma_0 = 1 \otimes \gamma([x])$.
On en déduit deux faits : (a) la classe
$\gamma([x])$ est indépendante de $x$, (b) il suffit
de montrer que $\gamma([x])$ est non nulle, et donc (c)
il suffit de trouver un cycle de dimension zéro $\alpha$ sur $X$ tel que
$\gamma(\alpha) \not = 0$. Pour cela, choisissons un morphisme
fini
$$
f : X \longrightarrow \mathbf{P}^d_k
$$
Pour l'existence d'un tel morphisme, voir
Théorie de l'intersection, section \ref{intersection-section-projection},
en particulier le lemme \ref{intersection-lemma-projection-generically-finite}.
Observons que $f$ est syntomique fini (localement d'intersection complète
d'après Compléments sur les morphismes, lemme \ref{more-morphisms-lemma-lci-permanence},
et plat d'après Algèbre, lemme \ref{algebra-lemma-CM-over-regular-flat}).
D'après la proposition \ref{derham-proposition-Garel}, on dispose d'un morphisme trace
$$
\Theta_f :
f_*\Omega^\bullet_{X/k}
\longrightarrow
\Omega^\bullet_{\mathbf{P}^d_k/k}
$$
dont la composée avec le morphisme canonique
$$
\Omega^\bullet_{\mathbf{P}^d_k/k}
\longrightarrow
f_*\Omega^\bullet_{X/k}
$$
est la multiplication par le degré de $f$. On obtient donc une application
$$
\Theta : H_{dR}^{2d}(X/k) \to H_{dR}^{2d}(\mathbf{P}^d_k/k)
$$
telle que $\Theta \circ f^*$ soit la multiplication par un entier strictement positif.
Ainsi, si l'on peut trouver sur $\mathbf{P}^d_k$ un cycle de dimension zéro dont la classe
est non nulle, on conclut par la compatibilité de $\gamma$
avec les images inverses. Cela résulte du
lemme \ref{derham-lemma-de-rham-cohomology-projective-space} et achève
la démonstration de l'axiome (A6).

\medskip\noindent
Nous utiliserons les faits suivants sans plus les mentionner.
Premièrement, d'après Théories cohomologiques de Weil, remarque \ref{weil-remark-trace},
l'application $\lambda_X : H^{2d}_{dR}(X/k) \to k$ est unique.
Deuxièmement, dans la démonstration de l'axiome (A6), nous avons
vu que $\lambda_X(\gamma([x])) = 1$ lorsque $X$ est irréductible ; autrement dit,
la composée de l'application de classe de cycle
$\gamma : \CH^d(X) \to H_{dR}^{2d}(X/k)$ avec $\lambda_X$
est l'application degré.

\medskip\noindent
Axiome (A9). Soit $Y \subset X$ un diviseur lisse non vide sur un
schéma projectif lisse équidimensionnel non vide $X$ sur $k$,
de dimension $d$. Il faut montrer que le diagramme
$$
\xymatrix{
H_{dR}^{2d - 2}(X/k)
\ar[rrr]_{c^{dR}_1(\mathcal{O}_X(Y)) \cap -} \ar[d]_{restriction} & & &
H_{dR}^{2d}(X) \ar[d]^{\lambda_X} \\
H_{dR}^{2d - 2}(Y/k) \ar[rrr]^-{\lambda_Y} & & & k
}
$$
est commutatif, où $\lambda_X$ et $\lambda_Y$ sont comme dans l'axiome (A6).
Nous avons vu ci-dessus que, si l'on décompose $X = \coprod X_i$ en composantes connexes
(ou, de manière équivalente, irréductibles),
on a la décomposition correspondante $\lambda_X = \sum \lambda_{X_i}$.
De même, si l'on décompose $Y = \coprod Y_j$ en composantes connexes (ou, de manière équivalente,
irréductibles), on a $\lambda_Y = \sum \lambda_{Y_j}$.
De plus, dans ce cas,
$\mathcal{O}_X(Y) = \otimes_j \mathcal{O}_X(Y_j)$, d'où
$$
c_1^{dR}(\mathcal{O}_X(Y)) = \sum\nolimits_j
c^{dR}_1(\mathcal{O}_X(Y_j))
$$
dans $H_{dR}^2(X/k)$. Une chasse au diagramme directe montre qu'il suffit
de prouver la commutativité lorsque $X$ et $Y$ sont tous deux
irréductibles. Alors $H_{dR}^{2d - 2}(Y/k)$ est de dimension $1$,
car on dispose de la dualité de Poincaré pour $Y$ d'après
Théories cohomologiques de Weil, lemme \ref{weil-lemma-poincare-duality}.
D'après l'axiome (A4), le noyau de la restriction (flèche verticale de gauche)
est contenu dans le noyau du cup-produit avec $c^{dR}_1(\mathcal{O}_X(Y))$.
Il suffit donc de trouver une classe de cohomologie
$a \in H_{dR}^{2d - 2}(X)$ dont la restriction à $Y$ soit non nulle
et pour laquelle le diagramme soit commutatif, à savoir pour $a$.
Prenons un module inversible ample quelconque $\mathcal{L}$ et posons
$$
a = c^{dR}_1(\mathcal{L})^{d - 1}
$$
On sait alors que $a|_Y = c^{dR}_1(\mathcal{L}|_Y)^{d - 1}$,
d'où
$$
\lambda_Y(a|_Y) = \deg(c_1(\mathcal{L}|_Y)^{d - 1} \cap [Y])
$$
d'après la description de $\lambda_Y$ ci-dessus. C'est un entier strictement positif,
d'après Homologie de Chow, lemme
\ref{chow-lemma-degrees-and-numerical-intersections}, combiné avec
Variétés, lemme \ref{varieties-lemma-ample-positive}.
De même, on obtient
$$
\lambda_X(c^{dR}_1(\mathcal{O}_X(Y)) \cap a) =
\deg(c_1(\mathcal{O}_X(Y)) \cap c_1(\mathcal{L})^{d - 1} \cap [X])
$$
Puisque l'on sait que $c_1(\mathcal{O}_X(Y)) \cap [X] = [Y]$, presque
par définition, on dispose d'une égalité de cycles de dimension zéro
$$
(Y \to X)_*\left(c_1(\mathcal{L}|_Y)^{d - 1} \cap [Y]\right) =
c_1(\mathcal{O}_X(Y)) \cap c_1(\mathcal{L})^{d - 1} \cap [X]
$$
sur $X$. Ces cycles ont donc le même degré, ce qui achève la démonstration.

\begin{proposition}
\label{derham-proposition-de-rham-is-weil}
Soit $k$ un corps de caractéristique nulle. Le foncteur qui
envoie un schéma projectif lisse $X$ sur $k$ sur $H_{dR}^*(X/k)$
est une théorie cohomologique de Weil au sens de
Théories cohomologiques de Weil, définition
\ref{weil-definition-weil-cohomology-theory}.
\end{proposition}

\begin{proof}
Dans la discussion ci-dessus, nous avons montré que nos données (D0), (D1), (D2')
satisfont aux axiomes (A1) -- (A9) de Théories cohomologiques de Weil, section
\ref{weil-section-c1}. On conclut donc par
Théories cohomologiques de Weil, proposition \ref{weil-proposition-get-weil}.

\medskip\noindent
Ne lisez pas ce qui suit. Dans la démonstration des assertions, nous avons aussi utilisé les
lemmes \ref{derham-lemma-proper-smooth-de-Rham},
\ref{derham-lemma-pullback-c1},
\ref{derham-lemma-log-complex-consequence},
\ref{derham-lemma-kunneth-de-rham},
\ref{derham-lemma-bottom-part-degenerates} et
\ref{derham-lemma-de-rham-cohomology-projective-space},
les propositions
\ref{derham-proposition-projective-space-bundle-formula},
\ref{derham-proposition-blowup-split} et
\ref{derham-proposition-Garel},
Théories cohomologiques de Weil, lemmes
\ref{weil-lemma-check-over-extension},
\ref{weil-lemma-chern-classes},
\ref{weil-lemma-cycle-classes} et
\ref{weil-lemma-poincare-duality},
Théories cohomologiques de Weil, remarque \ref{weil-remark-trace},
Variétés, lemmes
\ref{varieties-lemma-proper-geometrically-reduced-global-sections} et
\ref{varieties-lemma-ample-positive},
Théorie de l'intersection, section \ref{intersection-section-projection} et
lemme \ref{intersection-lemma-projection-generically-finite},
Compléments sur les morphismes, lemme \ref{more-morphisms-lemma-lci-permanence},
Algèbre, lemme \ref{algebra-lemma-CM-over-regular-flat}, et
Homologie de Chow, lemme
\ref{chow-lemma-degrees-and-numerical-intersections}.
\end{proof}

\begin{remark}
\label{derham-remark-hodge-cohomology-is-weil}
Exactement de la même manière que ci-dessus, on peut montrer que
la cohomologie de Hodge $X \mapsto H_{Hodge}^*(X/k)$, munie
de $c_1^{Hodge}$, définit une théorie
cohomologique de Weil. Si nous en avons besoin, nous en donnerons ici
un énoncé précis et une démonstration. Il en résulte la conséquence
amusante suivante : si les nombres de Betti d'une théorie cohomologique
de Weil sont indépendants de la théorie cohomologique de Weil choisie
(sur notre corps $k$ de caractéristique $0$), alors
la suite spectrale de Hodge vers de Rham
dégénère en $E_1$ ! Bien entendu, la dégénérescence de
la suite spectrale de Hodge vers de Rham est connue
(voir par exemple \cite{Deligne-Illusie} pour une remarquable démonstration algébrique),
mais ce résultat n'est nullement facile ! Cela suggère que l'indépendance
des nombres de Betti est également difficile à démontrer
et, à notre connaissance, reste un problème ouvert. Voir
Théories cohomologiques de Weil, remarque
\ref{weil-remark-betti-numbers-in-some-sense} pour une question apparentée.
\end{remark}







\section{Morphismes de Gysin pour les immersions fermées}
\label{derham-section-gysin}

\noindent
Dans cette section, nous définissons le morphisme de Gysin pour les immersions fermées.

\begin{remark}
\label{derham-remark-gysin-equations}
Soit $X \to S$ un morphisme de schémas. Soient
$f_1, \ldots, f_c \in \Gamma(X, \mathcal{O}_X)$ et $Z \subset X$
le sous-schéma fermé défini par $f_1, \ldots, f_c$. Nous étudierons ci-dessous le {\it morphisme de Gysin}
\begin{equation}
\label{derham-equation-gysin}
\gamma^p_{f_1, \ldots, f_c} :
\Omega^p_{Z/S}
\longrightarrow
\mathcal{H}_Z^c(\Omega^{p + c}_{X/S})
\end{equation}
défini comme suit. Étant donnée une section locale $\omega$ de $\Omega^p_{Z/S}$
qui est la restriction d'une section $\tilde \omega$ de $\Omega^p_{X/S}$,
nous posons
$$
\gamma^p_{f_1, \ldots, f_c}(\omega) =
c_{f_1, \ldots, f_c}(\tilde \omega|_Z) \wedge
\text{d}f_1 \wedge \ldots \wedge \text{d}f_c
$$
où $c_{f_1, \ldots, f_c} : \Omega^p_{X/S} \otimes \mathcal{O}_Z \to
\mathcal{H}_Z^c(\Omega^p_{X/S})$ est le morphisme construit dans
Catégories dérivées des schémas, remarque
\ref{perfect-remark-supported-map-c-equations}.
Cette définition est indépendante des choix : pour $\omega$ donné, nous pouvons modifier le choix de
$\tilde \omega$ par des éléments de la forme
$\sum f_i \omega'_i + \sum \text{d}(f_i) \wedge \omega''_i$,
que la construction envoie sur zéro.
\end{remark}

\begin{lemma}
\label{derham-lemma-gysin-differential}
Le morphisme de Gysin (\ref{derham-equation-gysin}) est compatible avec les différentielles de de Rham
sur $\Omega^\bullet_{X/S}$ et $\Omega^\bullet_{Z/S}$.
\end{lemma}

\begin{proof}
Ceci résulte d'un calcul presque immédiat, une fois l'énoncé correctement interprété.
Rappelons d'abord que le foncteur
$\mathcal{H}^c_Z$ calculé dans la catégorie des $\mathcal{O}_X$-modules
coïncide avec le foncteur défini de manière analogue dans la catégorie des faisceaux abéliens
sur $X$ ; voir
Cohomologie, lemme \ref{cohomology-lemma-sections-support-abelian-unbounded}.
La différentielle $\text{d} : \Omega^p_{X/S} \to \Omega^{p + 1}_{X/S}$
induit donc un morphisme
$\mathcal{H}^c_Z(\Omega^p_{X/S}) \to \mathcal{H}^c_Z(\Omega^{p + 1}_{X/S})$.
De plus, la construction du complexe de {\v C}ech alterné augmenté dans
Catégories dérivées des schémas, remarque \ref{perfect-remark-support-c-equations},
s'applique à la catégorie des faisceaux abéliens. Le morphisme
$$
\Coker\left(\bigoplus \mathcal{F}_{1 \ldots \hat i \ldots c} \to
\mathcal{F}_{1 \ldots c}\right)
\longrightarrow
i_*\mathcal{H}^c_Z(\mathcal{F})
$$
utilisé pour construire $c_{f_1, \ldots, f_c}$ dans
Catégories dérivées des schémas, remarque
\ref{perfect-remark-supported-map-c-equations},
est bien défini et
fonctoriel dans la catégorie de tous les faisceaux abéliens sur $X$.
Le lemme résulte donc de l'égalité
$$
\text{d}\left(
\frac{\tilde \omega \wedge \text{d}f_1 \wedge \ldots \wedge
\text{d}f_c}{f_1 \ldots f_c}\right) =
\frac{\text{d}(\tilde \omega) \wedge
\text{d}f_1 \wedge \ldots \wedge \text{d}f_c}{f_1 \ldots f_c}
$$,
qui est claire.
\end{proof}

\begin{lemma}
\label{derham-lemma-gysin-global}
Soit $X \to S$ un morphisme de schémas. Soit $Z \to X$ une immersion fermée
de présentation finie dont le faisceau conormal $\mathcal{C}_{Z/X}$ est
localement libre de rang $c$. Il existe alors un morphisme canonique
$$
\gamma^p : \Omega^p_{Z/S} \to \mathcal{H}^c_Z(\Omega^{p + c}_{X/S})
$$
donné localement par les morphismes $\gamma^p_{f_1, \ldots, f_c}$
de la remarque \ref{derham-remark-gysin-equations}.
\end{lemma}

\begin{proof}
Les hypothèses impliquent que, pour tout $x \in Z \subset X$, il existe un
voisinage ouvert $U$ de $x$ sur lequel $Z$ est défini par $c$
éléments $f_1, \ldots, f_c \in \mathcal{O}_X(U)$. Il suffit donc de montrer que, si $f_1, \ldots, f_c$ et
$g_1, \ldots, g_c$ sont des éléments de $\mathcal{O}_X(U)$ définissant $Z \cap U$,
les morphismes $\gamma^p_{f_1, \ldots, f_c}$
et $\gamma^p_{g_1, \ldots, g_c}$ coïncident. Pour cela, quitte à restreindre
$U$, nous pouvons supposer que $g_j = \sum a_{ji} f_i$ pour certains
$a_{ji} \in \mathcal{O}_X(U)$. Nous avons alors
$c_{f_1, \ldots, f_c} = \det(a_{ji}) c_{g_1, \ldots, g_c}$ d'après
Catégories dérivées des schémas, lemme
\ref{perfect-lemma-supported-map-determinant}.
D'autre part, nous avons
$$
\text{d}(g_1) \wedge \ldots \wedge \text{d}(g_c) \equiv
\det(a_{ji}) \text{d}(f_1) \wedge \ldots \wedge \text{d}(f_c)
\bmod (f_1, \ldots, f_c)\Omega^c_{X/S}
$$
En combinant ces relations, un calcul direct donne
l'égalité voulue.
\end{proof}

\begin{lemma}
\label{derham-lemma-gysin-differential-global}
Soient $X \to S$ et $i : Z \to X$ comme dans le lemme \ref{derham-lemma-gysin-global}.
Le morphisme de Gysin $\gamma^p$ est compatible avec les différentielles de de Rham
sur $\Omega^\bullet_{X/S}$ et $\Omega^\bullet_{Z/S}$.
\end{lemma}

\begin{proof}
La vérification est locale, et l'assertion résulte alors du
lemme \ref{derham-lemma-gysin-differential}.
\end{proof}

\begin{lemma}
\label{derham-lemma-gysin-projection}
Soient $X \to S$ et $i : Z \to X$ comme dans le lemme \ref{derham-lemma-gysin-global}.
Pour $\alpha \in H^q(X, \Omega^p_{X/S})$, nous avons
$\gamma^p(\alpha|_Z) = i^{-1}\alpha \wedge \gamma^0(1)$ dans
$H^q(Z, \mathcal{H}^c_Z(\Omega^{p + c}_{X/S}))$.
Les notations sont expliquées dans la démonstration.
\end{lemma}

\begin{proof}
La restriction $\alpha|_Z$ est l'élément de $H^q(Z, \Omega^p_{Z/S})$
donné par la fonctorialité de la cohomologie de Hodge. En appliquant la fonctorialité
de la cohomologie au morphisme
$\gamma^p : \Omega^p_{Z/S} \to \mathcal{H}^c_Z(\Omega^{p + c}_{X/S})$,
nous obtenons $\gamma^p(\alpha|_Z)$ dans
$H^q(Z, \mathcal{H}^c_Z(\Omega^{p + c}_{X/S}))$.
Cela explique le membre de gauche de la formule.

\medskip\noindent
Pour expliquer le membre de droite, nous prenons d'abord l'image inverse par le morphisme
d'espaces annelés $i : (Z, i^{-1}\mathcal{O}_X) \to (X, \mathcal{O}_X)$,
ce qui donne l'élément $i^{-1}\alpha \in H^q(Z, i^{-1}\Omega^p_{X/S})$.
Soit $\gamma^0(1) \in H^0(Z, \mathcal{H}_Z^c(\Omega^c_{X/S}))$
l'image de $1 \in H^0(Z, \mathcal{O}_Z) = H^0(Z, \Omega^0_{Z/S})$
par $\gamma^0$. Le cup-produit fournit un élément
$$
i^{-1}\alpha \cup \gamma^0(1)
\in
H^{q + c}(Z, 
i^{-1}\Omega^p_{X/S} \otimes_{i^{-1}\mathcal{O}_X}
\mathcal{H}^c_Z(\Omega^c_{X/S}))
$$
En utilisant Cohomologie, remarque \ref{cohomology-remark-support-cup-product},
et le produit extérieur, nous obtenons des morphismes canoniques
$$
i^{-1}\Omega^p_{X/S} \otimes_{i^{-1}\mathcal{O}_X}^\mathbf{L}
R\mathcal{H}_Z(\Omega^c_{X/S}) \to
R\mathcal{H}_Z(\Omega^p_{X/S} \otimes_{\mathcal{O}_X}^\mathbf{L}
\Omega^c_{X/S}) \to
R\mathcal{H}_Z(\Omega^{p + c}_{X/S})
$$
D'après Catégories dérivées des schémas, lemme
\ref{perfect-lemma-supported-trivial-vanishing},
les faisceaux de cohomologie des objets $R\mathcal{H}_Z(\Omega^j_{X/S})$ sont nuls
en degrés $> c$. Ainsi, en prenant les faisceaux de cohomologie
en degré $c$, nous obtenons un morphisme
$$
i^{-1}\Omega^p_{X/S} \otimes_{i^{-1}\mathcal{O}_X}
\mathcal{H}^c_Z(\Omega^c_{X/S}) \longrightarrow
\mathcal{H}^c_Z(\Omega^{p + c}_{X/S})
$$
L'expression $i^{-1}\alpha \wedge \gamma^0(1)$ est l'image
du cup-produit $i^{-1}\alpha \cup \gamma^0(1)$ par la
fonctorialité de la cohomologie.

\medskip\noindent
Le sens de la formule étant ainsi précisé, les
propriétés générales des cup-produits
(Cohomologie, section \ref{cohomology-section-cup-product})
montrent qu'il suffit de prouver que le diagramme
$$
\xymatrix{
i^{-1}\Omega^p_X \otimes \Omega^0_Z \ar[rr]_{\text{id} \otimes \gamma^0}
\ar[d] & &
i^{-1}\Omega^p_X \otimes \mathcal{H}^c_Z(\Omega^c_X) \ar[d]^\wedge \\
\Omega^p_Z \otimes \Omega^0_Z \ar[r]^\wedge &
\Omega^p_Z \ar[r]^{\gamma^p} &
\mathcal{H}^c_Z(\Omega^{p + c}_X)
}
$$
est commutatif dans la catégorie des faisceaux sur $Z$ (avec un abus de
notation évident). Cela se ramène à un calcul simple pour les morphismes
$\gamma^j_{f_1, \ldots, f_c}$, que nous omettons ; ces morphismes sont en fait
choisis précisément pour assurer cette commutativité et pour que $1$ ait pour image
$\frac{\text{d}f_1 \wedge \ldots \wedge \text{d}f_c}{f_1 \ldots f_c}$.
\end{proof}

\begin{lemma}
\label{derham-lemma-gysin-transverse}
Soit $c \geq 0$ un entier. Soit
$$
\xymatrix{
Z' \ar[d]_h \ar[r] & X' \ar[d]_g \ar[r] & S' \ar[d] \\
Z \ar[r] & X \ar[r] & S
}
$$
un diagramme commutatif de schémas.
Supposons que
\begin{enumerate}
\item $Z \to X$ et $Z' \to X'$
vérifient les hypothèses du lemme \ref{derham-lemma-gysin-global} ;
\item le carré de gauche du diagramme soit cartésien ;
\item $h^*\mathcal{C}_{Z/X} \to \mathcal{C}_{Z'/X'}$
(Morphismes, lemme \ref{morphisms-lemma-conormal-functorial})
soit un isomorphisme.
\end{enumerate}
Alors le diagramme
$$
\xymatrix{
h^*\Omega^p_{Z/S} \ar[rr]_-{h^{-1}\gamma^p} \ar[d] & &
\mathcal{O}_{X'}|_{Z'} \otimes_{h^{-1}\mathcal{O}_X|_Z}
h^{-1}\mathcal{H}^c_Z(\Omega^{p + c}_{X/S}) \ar[d] \\
\Omega^p_{Z'/S'} \ar[rr]^{\gamma^p} & &
\mathcal{H}^c_{Z'}(\Omega^{p + c}_{X'/S'})
}
$$
est commutatif. La flèche verticale de gauche est donnée par la fonctorialité des modules de
différentielles et celle de droite fait intervenir
Cohomologie, remarque \ref{cohomology-remark-support-functorial}.
\end{lemma}

\begin{proof}
Plus précisément, considérons le composé
\begin{align*}
\mathcal{O}_{X'}|_{Z'} \otimes_{h^{-1}\mathcal{O}_X|_Z}^\mathbf{L}
h^{-1}R\mathcal{H}_Z(\Omega^{p + c}_{X/S})
& \to
R\mathcal{H}_{Z'}(Lg^*\Omega^{p + c}_{X/S}) \\
& \to
R\mathcal{H}_{Z'}(g^*\Omega^{p + c}_{X/S}) \\
& \to
R\mathcal{H}_{Z'}(\Omega^{p + c}_{X'/S'})
\end{align*}
dont la première flèche est donnée par
Cohomologie, remarque \ref{cohomology-remark-support-functorial},
et la dernière par la fonctorialité des différentielles.
L'annulation des faisceaux de cohomologie en degrés $> c$,
d'après Catégories dérivées des schémas, lemme
\ref{perfect-lemma-supported-trivial-vanishing},
fournit la flèche verticale de droite.
Nous pouvons vérifier la commutativité localement.
Nous pouvons donc supposer que $Z$ est défini par
$f_1, \ldots, f_c \in \Gamma(X, \mathcal{O}_X)$.
Alors $Z'$ est défini par $f'_i = g^\sharp(f_i)$.
Les morphismes $c_{f_1, \ldots, f_c}$ et $c_{f'_1, \ldots, f'_c}$
s'insèrent dans le diagramme commutatif
$$
\xymatrix{
h^*i^*\Omega^p_{X/S} \ar[rr]_-{h^{-1}c_{f_1, \ldots, f_c}} \ar[d] & &
\mathcal{O}_{X'}|_{Z'} \otimes_{h^{-1}\mathcal{O}_X|_Z}
h^{-1}\mathcal{H}^c_Z(\Omega^p_{X/S}) \ar[d] \\
(i')^*\Omega^p_{X'/S'} \ar[rr]^{c_{f'_1, \ldots, f'_c}} & &
\mathcal{H}^c_{Z'}(\Omega^p_{X'/S'})
}
$$
Voir Catégories dérivées des schémas, remarque
\ref{perfect-remark-supported-functorial}.
Rappelons que, pour une $p$-forme $\omega$ sur $Z$, nous définissons
$\gamma^p(\omega)$ en choisissant (localement sur $X$ et $Z$)
une $p$-forme $\tilde \omega$ sur $X$ relevant $\omega$, puis en posant
$\gamma^p(\omega) =
c_{f_1, \ldots, f_c}(\tilde \omega) \wedge
\text{d}f_1 \wedge \ldots \wedge \text{d}f_c$.
Puisque l'image inverse de la forme $\text{d}f_1 \wedge \ldots \wedge \text{d}f_c$
est
$\text{d}f'_1 \wedge \ldots \wedge \text{d}f'_c$, nous concluons.
\end{proof}

\begin{remark}
\label{derham-remark-how-to-use}
Soient $X \to S$, $i : Z \to X$ et $c \geq 0$ comme dans le
lemme \ref{derham-lemma-gysin-global}.
Soit $p \geq 0$ ; supposons que $\mathcal{H}^i_Z(\Omega^{p + c}_{X/S}) = 0$
pour $i = 0, \ldots, c - 1$. Cette annulation a lieu si $X \to S$ est lisse
et si $Z \to X$ est une immersion régulière au sens de Koszul ; voir
Catégories dérivées des schémas, lemme \ref{perfect-lemma-supported-vanishing}.
Nous obtenons alors un morphisme
$$
\gamma^{p, q} :
H^q(Z, \Omega^p_{Z/S})
\longrightarrow
H^{q + c}(X, \Omega^{p + c}_{X/S})
$$
en utilisant d'abord
$\gamma^p : \Omega^p_{Z/S} \to \mathcal{H}^c_Z(\Omega^{p + c}_{X/S})$
pour atteindre
$$
H^q(Z, \mathcal{H}^c_Z(\Omega^{p + c}_{X/S})) =
H^q(Z, R\mathcal{H}_Z(\Omega^{p + c}_{X/S})[c]) =
H^q(X, i_*R\mathcal{H}_Z(\Omega^{p + c}_{X/S})[c])
$$,
puis le morphisme d'adjonction
$i_*R\mathcal{H}_Z(\Omega^{p + c}_{X/S}) \to \Omega^{p + c}_{X/S}$
pour passer au module de cohomologie de Hodge voulu.
\end{remark}

\begin{lemma}
\label{derham-lemma-gysin-differential-hodge}
Soient $X \to S$ et $i : Z \to X$ comme dans le lemme \ref{derham-lemma-gysin-global}.
Supposons que $X \to S$ soit lisse et que $Z \to X$ soit régulière au sens de Koszul.
Les morphismes de Gysin $\gamma^{p, q}$ sont compatibles avec les différentielles de de Rham
sur $\Omega^\bullet_{X/S}$ et $\Omega^\bullet_{Z/S}$.
\end{lemma}

\begin{proof}
Cela résulte immédiatement du
lemme \ref{derham-lemma-gysin-differential-global}.
\end{proof}

\begin{lemma}
\label{derham-lemma-gysin-projection-global}
Soient $X \to S$, $i : Z \to X$ et $c \geq 0$ comme dans le
lemme \ref{derham-lemma-gysin-global}. Supposons que $X \to S$ soit lisse et
que $Z \to X$ soit régulière au sens de Koszul. Pour $\alpha \in H^q(X, \Omega^p_{X/S})$, nous avons
$\gamma^{p, q}(\alpha|_Z) = \alpha \cup \gamma^{0, 0}(1)$ dans
$H^{q + c}(X, \Omega^{p + c}_{X/S})$, où $\gamma^{a, b}$ est défini dans la
remarque \ref{derham-remark-how-to-use}.
\end{lemma}

\begin{proof}
Ce lemme résulte du lemme \ref{derham-lemma-gysin-projection}
et de Cohomologie, lemme \ref{cohomology-lemma-support-cup-product}.
Le lecteur peut omettre la discussion plus détaillée qui suit.

\medskip\noindent
Nous utiliserons sans le rappeler l'égalité
$R\mathcal{H}_Z(\Omega^j_{X/S}) = \mathcal{H}^c_Z(\Omega^j_{X/S})[-c]$
pour tout $j$, signalée dans la remarque \ref{derham-remark-how-to-use}.
Nous utiliserons également sans les mentionner les identifications
$H^{q + c}_Z(X, \Omega^j_{X/S}) = H^{q + c}(Z, R\mathcal{H}_Z(\Omega^j_{X/S}) =
H^q(Z, \mathcal{H}^c_Z(\Omega^j_{X/S}))$ ; voir
Cohomologie, lemme \ref{cohomology-lemma-local-to-global-sections-with-support},
pour la première. Avec ces identifications :
\begin{enumerate}
\item $\gamma^0(1) \in H^c_Z(X, \Omega^c_{X/S})$ a pour image $\gamma^{0, 0}(1)$
dans $H^c(X, \Omega^c_{X/S})$ ;
\item le membre de droite $i^{-1}\alpha \wedge \gamma^0(1)$
de l'égalité du lemme \ref{derham-lemma-gysin-projection}
est l'image par le produit extérieur du cup-produit de
Cohomologie, remarque \ref{cohomology-remark-support-cup-product-global},
des éléments $\alpha$ et $\gamma^0(1)$ ; autrement dit, les constructions
de la démonstration du lemme \ref{derham-lemma-gysin-projection} et de
Cohomologie, remarque \ref{cohomology-remark-support-cup-product-global}, coïncident ;
\item d'après Cohomologie, lemme \ref{cohomology-lemma-support-cup-product},
cet élément a pour image $\alpha \cup \gamma^{0, 0}(1)$ dans
$H^{q + c}(X, \Omega^p_{X/S} \otimes \Omega^c_{X/S})$ ;
\item le membre de gauche $\gamma^p(\alpha|_Z)$ de l'égalité du
lemme \ref{derham-lemma-gysin-projection} a pour image
$\gamma^{p, q}(\alpha|_Z)$.
\end{enumerate}
Cela achève la démonstration.
\end{proof}

\begin{lemma}
\label{derham-lemma-gysin-transverse-global}
Soient $c \geq 0$ et
$$
\xymatrix{
Z' \ar[d]_h \ar[r] & X' \ar[d]_g \ar[r] & S' \ar[d] \\
Z \ar[r] & X \ar[r] & S
}
$$
vérifiant les hypothèses du lemme \ref{derham-lemma-gysin-transverse}. Supposons de plus
que $X \to S$ et $X' \to S'$ soient lisses, et que
$Z \to X$ et $Z' \to X'$ soient des immersions régulières au sens de Koszul.
Alors le diagramme
$$
\xymatrix{
H^q(Z, \Omega^p_{Z/S}) \ar[rr]_-{\gamma^{p, q}} \ar[d] & &
H^{q + c}(X, \Omega^{p + c}_{X/S}) \ar[d] \\
H^q(Z', \Omega^p_{Z'/S'}) \ar[rr]^{\gamma^{p, q}} & &
H^{q + c}(X', \Omega^{p + c}_{X'/S'})
}
$$
est commutatif, où $\gamma^{p, q}$ est défini dans la remarque \ref{derham-remark-how-to-use}.
\end{lemma}

\begin{proof}
Cela résulte de la combinaison du lemme \ref{derham-lemma-gysin-transverse}
et de Cohomologie, lemme \ref{cohomology-lemma-support-functorial}.
\end{proof}

\begin{lemma}
\label{derham-lemma-class-of-a-point}
Soit $k$ un corps. Soit $X$ un schéma irréductible, lisse et propre sur $k$,
de dimension $d$. Soit $Z \subset X$ le sous-schéma fermé réduit constitué
d'un unique point $k$-rationnel $x$. Alors l'image de
$1 \in k = H^0(Z, \mathcal{O}_Z) = H^0(Z, \Omega^0_{Z/k})$
par le morphisme $H^0(Z, \Omega^0_{Z/k}) \to H^d(X, \Omega^d_{X/k})$
de la remarque \ref{derham-remark-how-to-use} est non nulle.
\end{lemma}

\begin{proof}
Le morphisme $\gamma^0 : \mathcal{O}_Z \to
\mathcal{H}^d_Z(\Omega^d_{X/k}) = R\mathcal{H}_Z(\Omega^d_{X/k})[d]$
a pour adjoint un morphisme
$$
g^0 : i_*\mathcal{O}_Z \longrightarrow \Omega^d_{X/k}[d]
$$
dans $D(\mathcal{O}_X)$. Rappelons que $\Omega^d_{X/k} = \omega_X$ est un
faisceau dualisant pour $X/k$ ; voir
Dualité pour les schémas, lemme \ref{duality-lemma-duality-proper-over-field}.
Le dual $k$-linéaire du morphisme de l'énoncé
est donc le morphisme
$$
H^0(X, \mathcal{O}_X) \to \Ext^d_X(i_*\mathcal{O}_Z, \omega_X)
$$
qui envoie $1$ sur $g^0$. Il suffit ainsi de montrer que $g^0$ est non nul.
Nous pouvons le vérifier dans un voisinage quelconque $U$ du point $x$. Choisissons $U$
de sorte qu'il existe $f_1, \ldots, f_d \in \mathcal{O}_X(U)$
s'annulant seulement en $x$ et engendrant l'idéal maximal
$\mathfrak m_x \subset \mathcal{O}_{X, x}$. Nous pouvons supposer
que $U = \Spec(R)$ est affine. En reprenant la
construction de $\gamma^0$, nous voyons que notre extension est donnée par
$$
k \to
(R \to \bigoplus\nolimits_{i_0} R_{f_{i_0}} \to
\bigoplus\nolimits_{i_0 < i_1} R_{f_{i_0}f_{i_1}} \to
\ldots \to R_{f_1\ldots f_r})[d] \to R[d]
$$,
où $1$ a pour image $1/f_1 \ldots f_c$ par le premier morphisme.
Cette extension est non nulle, car $1/f_1 \ldots f_c$ est, dans ce cas, un élément non nul
du groupe de cohomologie locale $H^d_{(f_1, \ldots, f_d)}(R)$.
\end{proof}






\section{Dualité de Poincaré relative}
\label{derham-section-relative-poincare-duality}

\noindent
Dans cette section, nous démontrons la dualité de Poincaré pour la cohomologie de de Rham relative
d'un schéma propre et lisse sur une base. Nous recommandons vivement au lecteur
de consulter d'abord la section \ref{derham-section-poincare-duality}.

\begin{situation}
\label{derham-situation-relative-duality}
Ici, $S$ est un schéma quasi-compact et quasi-séparé, et
$f : X \to S$ est un morphisme propre et lisse de schémas dont toutes les
fibres sont non vides et équidimensionnelles de dimension $n$.
\end{situation}

\begin{lemma}
\label{derham-lemma-relative-bottom-part-degenerates}
Dans la situation \ref{derham-situation-relative-duality}, l'image directe
$f_*\mathcal{O}_X$ est une $\mathcal{O}_S$-algèbre finie étale
et, localement sur $S$, nous avons $Rf_*\mathcal{O}_X = f_*\mathcal{O}_X \oplus P$
dans $D(\mathcal{O}_S)$, où $P$ est parfait d'amplitude de Tor dans $[1, \infty)$.
Le morphisme $\text{d} : f_*\mathcal{O}_X \to f_*\Omega_{X/S}$ est nul.
\end{lemma}

\begin{proof}
La première partie de l'énoncé résulte de
Catégories dérivées des schémas, lemme \ref{perfect-lemma-proper-flat-geom-red}.
En posant $S' = \underline{\Spec}_S(f_*\mathcal{O}_X)$, nous obtenons une factorisation
$X \to S' \to S$ (c'est la factorisation de Stein ; voir
Compléments sur les morphismes, section \ref{more-morphisms-section-stein-factorization},
bien que nous n'en ayons pas besoin ici),
et nous voyons que $\Omega_{X/S} = \Omega_{X/S'}$, par exemple d'après
Morphismes, lemmes \ref{morphisms-lemma-triangle-differentials} et
\ref{morphisms-lemma-etale-at-point}. Cela implique évidemment que
$\text{d} : f_*\mathcal{O}_X \to f_*\Omega_{X/S}$ est nul.
\end{proof}

\begin{lemma}
\label{derham-lemma-relative-duality-hodge}
Dans la situation \ref{derham-situation-relative-duality}, il existe un morphisme de
$\mathcal{O}_S$-modules
$$
t : Rf_*\Omega^n_{X/S}[n] \longrightarrow \mathcal{O}_S
$$,
unique à précomposition près par la multiplication par un élément inversible de
$H^0(X, \mathcal{O}_X)$, ayant la propriété suivante : pour tout $p$, l'accouplement
$$
Rf_*\Omega^p_{X/S}
\otimes_{\mathcal{O}_S}^\mathbf{L}
Rf_*\Omega^{n - p}_{X/S}[n]
\longrightarrow
\mathcal{O}_S
$$
donné par le cup-produit relatif composé avec $t$
est un accouplement parfait de complexes parfaits sur $S$.
\end{lemma}

\begin{proof}
Supposons d'abord que $S$ soit un schéma noethérien.
Soit $\omega^\bullet_{X/S}$ le complexe dualisant relatif de $X$ sur $S$, comme
dans Dualité pour les schémas, remarque \ref{duality-remark-relative-dualizing-complex},
et soit $Rf_*\omega_{X/S}^\bullet \to \mathcal{O}_S$ son morphisme trace. D'après
Dualité pour les schémas, lemme \ref{duality-lemma-smooth-proper}
(c'est ici que nous utilisons que $S$ est noethérien),
il existe un isomorphisme $\omega^\bullet_{X/S} \cong \Omega^n_{X/S}[n]$,
qui nous permet d'obtenir $t$. Les complexes $Rf_*\Omega^p_{X/S}$
sont parfaits d'après le lemme \ref{derham-lemma-proper-smooth-de-Rham}.
Puisque $\Omega^p_{X/S}$ est localement libre et que
$\Omega^p_{X/S} \otimes_{\mathcal{O}_X} \Omega^{n - p}_{X/S} \to
\Omega^n_{X/S}$ définit un isomorphisme $\Omega^p_{X/S} \cong
\SheafHom_{\mathcal{O}_X}(\Omega^{n - p}_{X/S}, \Omega^n_{X/S})$,
l'accouplement induit par le cup-produit relatif est parfait d'après
Dualité pour les schémas, remarque
\ref{duality-remark-relative-dualizing-complex-relative-cup-product}.

\medskip\noindent
Démonstration de l'existence dans le cas général.
Par approximation noethérienne absolue, nous pouvons trouver un diagramme cartésien
$$
\xymatrix{
X \ar[r] \ar[d]^f & X_0 \ar[d]^{f_0} \\
S \ar[r]^g & S_0
}
$$
avec $S_0$ noethérien et $f_0$ propre et lisse,
dont toutes les fibres sont non vides et équidimensionnelles de dimension $n$.
Voir Limites, lemmes \ref{limits-lemma-descend-finite-presentation},
\ref{limits-lemma-descend-smooth}, \ref{limits-lemma-descend-dimension-d},
\ref{limits-lemma-descend-surjective} et
\ref{limits-lemma-eventually-proper}. Choisissons un morphisme
$$
t_0 : Rf_{0, *}\Omega^n_{X_0/S_0}[n] \longrightarrow \mathcal{O}_{S_0}
$$
comme dans l'énoncé. Par cohomologie et changement de base, les
complexes $Rf_*(\Omega^p_{X/S})$ s'identifient à
$Lg^*(Rf_{0, *}\Omega^p_{X_0/S_0})$ pour tout $p \geq 0$ ; voir
le lemme \ref{derham-lemma-proper-smooth-de-Rham}. En particulier, l'image inverse
de $t_0$ par $g$ nous fournit un morphisme $t$ comme dans l'énoncé.
En effet, d'après Cohomologie, lemme \ref{cohomology-lemma-base-change-cup-product},
le morphisme de cup-produit
$$
Rf_*\Omega^p_{X/S}
\otimes_{\mathcal{O}_S}^\mathbf{L}
Rf_*\Omega^{n - p}_{X/S}[n]
\longrightarrow
Rf_*\Omega^n_{X/S}
$$
est l'image inverse par $g$ du morphisme de cup-produit correspondant pour $f_0$.
La propriété selon laquelle $t_0$ donne des accouplements parfaits se transporte donc par image inverse
en la même propriété pour $t$ sur $S$.

\medskip\noindent
Unicité de $t$. Choisissons un triangle distingué
$f_*\mathcal{O}_X \to Rf_*\mathcal{O}_X \to P \to f_*\mathcal{O}_X[1]$.
D'après le lemme \ref{derham-lemma-relative-bottom-part-degenerates},
l'objet $P$ est parfait d'amplitude de Tor dans $[1, \infty)$,
et le triangle est scindé localement sur $S$.
Ainsi, $R\SheafHom_{\mathcal{O}_X}(P, \mathcal{O}_X)$ est parfait
d'amplitude de Tor dans $(-\infty, -1]$. La dualité ci-dessus montre donc que,
localement sur $S$, nous avons
$$
Rf_*\Omega^n_{X/S}[n] \cong
R\SheafHom_{\mathcal{O}_S}(f_*\mathcal{O}_X, \mathcal{O}_S)
\oplus R\SheafHom_{\mathcal{O}_X}(P, \mathcal{O}_X)
$$
Il en résulte que $R^nf_*\Omega^n_{X/S}$ est localement libre de rang fini et
que nous obtenons un accouplement $\mathcal{O}_S$-bilinéaire parfait
$$
f_*\mathcal{O}_X \times R^nf_*\Omega^n_{X/S} \longrightarrow \mathcal{O}_S
$$
au moyen de $t$.
Cela implique que tout morphisme $\mathcal{O}_S$-linéaire
$t' : R^nf_*\Omega^n_{X/S} \to \mathcal{O}_S$ est de la forme
$t' = t \circ g$ pour un certain
$g \in \Gamma(S, f_*\mathcal{O}_X) = \Gamma(X, \mathcal{O}_X)$.
Pour que $t'$ définisse encore un accouplement parfait, $g$ doit
être inversible. Cela achève la démonstration.
\end{proof}

\begin{lemma}
\label{derham-lemma-relative-top-part-degenerates}
Dans la situation \ref{derham-situation-relative-duality}, le morphisme
$\text{d} : R^nf_*\Omega^{n - 1}_{X/S} \to R^nf_*\Omega^n_{X/S}$
est nul.
\end{lemma}

\noindent
Comme nous l'avons signalé dans la démonstration du lemme \ref{derham-lemma-top-part-degenerates},
ce lemme n'est pas une conséquence immédiate des lemmes
\ref{derham-lemma-relative-duality-hodge} et
\ref{derham-lemma-relative-bottom-part-degenerates}.

\begin{proof}[Démonstration lorsque $S$ est réduit]
Supposons $S$ réduit. Observons que
$\text{d} : R^nf_*\Omega^{n - 1}_{X/S} \to R^nf_*\Omega^n_{X/S}$
est un morphisme $\mathcal{O}_S$-linéaire de $\mathcal{O}_S$-modules quasi-cohérents.
Le $\mathcal{O}_S$-module $R^nf_*\Omega^n_{X/S}$ est localement libre de rang fini
(car il est dual du $\mathcal{O}_S$-module localement libre de rang fini
$f_*\mathcal{O}_X$, d'après les lemmes
\ref{derham-lemma-relative-duality-hodge} et
\ref{derham-lemma-relative-bottom-part-degenerates}).
Puisque $S$ est réduit, il suffit de montrer que
le germe de $\text{d}$ en chaque point générique $\eta \in S$
est nul ; on le voit en considérant les sections sur les ouverts affines,
en utilisant que le but de $\text{d}$ est localement libre et en appliquant
Algèbre, lemme \ref{algebra-lemma-reduced-ring-sub-product-fields} (2).
Puisque $S$ est réduit, nous avons $\mathcal{O}_{S, \eta} = \kappa(\eta)$ ; voir
Algèbre, lemme \ref{algebra-lemma-minimal-prime-reduced-ring}.
Ainsi, $\text{d}_\eta$ s'identifie au morphisme
$$
\text{d} :
H^n(X_\eta, \Omega^{n - 1}_{X_\eta/\kappa(\eta)})
\longrightarrow
H^n(X_\eta, \Omega^n_{X_\eta/\kappa(\eta)})
$$,
qui est nul d'après le lemme \ref{derham-lemma-top-part-degenerates}.
\end{proof}

\begin{proof}[Démonstration dans le cas général]
La question est locale sur $S$ pour la topologie plate : si $S' \to S$ est un morphisme
plat surjectif de schémas et si le morphisme considéré est nul après image inverse sur $S'$,
alors il est nul. De plus, sa formation commute aux changements de base
plats par le théorème de changement de base plat (Cohomologie des schémas, lemme
\ref{coherent-lemma-flat-base-change-cohomology}).

\medskip\noindent
Considérons la factorisation de Stein $X \to S' \to S$ de
Compléments sur les morphismes, théorème
\ref{more-morphisms-theorem-stein-factorization-general}.
D'après le lemme \ref{derham-lemma-relative-bottom-part-degenerates}, le morphisme
$\pi : S' \to S$ est fini étale.
Le morphisme $f : X \to S'$ est propre (d'après le théorème),
lisse (d'après Compléments sur les morphismes, lemme
\ref{more-morphisms-lemma-smooth-etale-permanence}) et à fibres géométriquement
connexes d'après le théorème de factorisation de Stein.
Dans la démonstration du lemme \ref{derham-lemma-relative-bottom-part-degenerates},
nous avons vu que $\Omega_{X/S} = \Omega_{X/S'}$ puisque $S' \to S$ est étale.
Ainsi, $\Omega^\bullet_{X/S} = \Omega^\bullet_{X/S'}$.
Nous avons
$$
R^qf_*\Omega^p_{X/S} = \pi_*R^qf'_*\Omega^p_{X/S'}
$$
pour tous $p, q$, d'après la suite spectrale de Leray
(Cohomologie, lemme \ref{cohomology-lemma-relative-Leray}),
le fait que $\pi$ soit fini, donc affine, et
Cohomologie des schémas, lemme \ref{coherent-lemma-relative-affine-vanishing}
(nous utilisons aussi, bien entendu, que $R^qf'_*\Omega^p_{X'/S}$ est
quasi-cohérent).
Le morphisme du lemme s'obtient donc en appliquant $\pi_*$ à
$\text{d} : R^nf'_*\Omega^{n - 1}_{X/S'} \to R^nf'_*\Omega^n_{X/S'}$.
Autrement dit, pour démontrer le lemme, nous pouvons remplacer
$f : X \to S$ par $f' : X \to S'$ et nous ramener au cas étudié
dans le paragraphe suivant.

\medskip\noindent
Supposons que $f$ ait des fibres géométriquement connexes et que
$f_*\mathcal{O}_X = \mathcal{O}_S$.
Pour chaque $s \in S$, nous pouvons choisir un voisinage étale
$(S', s') \to (S, s)$ tel que le changement de base $X' \to S'$ de $S$
admette une section. Voir Compléments sur les morphismes, lemme
\ref{more-morphisms-lemma-etale-nbhd-dominates-smooth}.
D'après les remarques initiales de la démonstration, cela nous ramène au cas
étudié dans le paragraphe suivant.

\medskip\noindent
Supposons que $f$ ait des fibres géométriquement connexes,
que $f_*\mathcal{O}_X = \mathcal{O}_S$ et que nous disposions
d'une section $s : S \to X$ de $f$. Nous pouvons supposer, et nous supposons, que $S = \Spec(A)$
est affine. Le morphisme
$s^* : R\Gamma(X, \mathcal{O}_X) \to R\Gamma(S, \mathcal{O}_S) = A$
est une rétraction du morphisme $A \to R\Gamma(X, \mathcal{O}_X)$. Nous pouvons donc écrire
$$
R\Gamma(X, \mathcal{O}_X) = A \oplus P
$$,
où $P$ est le « noyau » de $s^*$. D'après
le lemme \ref{derham-lemma-relative-bottom-part-degenerates}, l'objet $P$
de $D(A)$ est parfait d'amplitude de Tor dans $[1, n]$. Comme dans la démonstration
du lemme \ref{derham-lemma-relative-duality-hodge}, nous voyons que
$H^n(X, \Omega^n_{X/S})$ est un $A$-module localement libre de rang $1$
(et même le dual de $A$, donc libre de rang $1$ ; nous allons bientôt choisir
un générateur, mais nous ne souhaitons pas vérifier qu'il s'agit du même,
et cela ne sera d'ailleurs pas nécessaire).

\medskip\noindent
Notons $Z \subset X$ l'image de $s$, qui est un sous-schéma fermé de $X$ d'après
Schémas, lemme \ref{schemes-lemma-section-immersion}.
Observons que $Z \to X$ est une immersion fermée régulière (et, a fortiori, régulière au sens de Koszul d'après
Diviseurs, lemme \ref{divisors-lemma-regular-quasi-regular-immersion}),
en vertu de
Diviseurs, lemme \ref{divisors-lemma-section-smooth-regular-immersion}.
Bien entendu, $Z \to X$ est de codimension $n$. D'après
la remarque \ref{derham-remark-how-to-use},
nous pouvons donc considérer le morphisme
$$
\gamma^{0, 0} : H^0(Z, \Omega^0_{Z/S}) \longrightarrow H^n(X, \Omega^n_{X/S})
$$,
et nous posons $\xi = \gamma^{0, 0}(1) \in H^n(X, \Omega^n_{X/S})$.

\medskip\noindent
Nous affirmons que $\xi$ est un élément d'une base. En effet, puisque le changement de base est valable
en degré maximal (voir, par exemple, Limites, lemme
\ref{limits-lemma-higher-direct-images-zero-above-dimension-fibre}),
nous avons
$H^n(X, \Omega^n_{X/S}) \otimes_A k = H^n(X_k, \Omega^n_{X_k/k})$
pour tout morphisme d'anneaux $A \to k$. La compatibilité de
la construction de $\xi$ au changement de base,
voir le lemme \ref{derham-lemma-gysin-transverse-global},
montre que l'image de $\xi$ dans $H^n(X_k, \Omega^n_{X_k/k})$
est non nulle d'après le lemme \ref{derham-lemma-class-of-a-point} lorsque $k$ est un corps.
Ainsi, $\xi$ est une section partout non nulle d'un module inversible,
donc un générateur.

\medskip\noindent
Soit $\theta \in H^n(X, \Omega^{n - 1}_{X/S})$. Nous devons montrer que
$\text{d}(\theta)$ est nul dans $H^n(X, \Omega^n_{X/S})$.
Nous pouvons écrire $\text{d}(\theta) = a \xi$ pour un certain $a \in A$,
puisque $\xi$ est un élément d'une base. Il s'agit alors de montrer que $a = 0$.

\medskip\noindent
Considérons l'immersion fermée
$$
\Delta : X \to X \times_S X
$$
C'est aussi une section d'un morphisme lisse (à savoir l'une ou l'autre des projections),
donc une immersion à la fois régulière et régulière au sens de Koszul, de codimension $n$.
Nous pouvons ainsi considérer les morphismes
$$
\gamma^{p, q} :
H^q(X, \Omega^p_{X/S})
\longrightarrow
H^{q + n}(X \times_S X, \Omega^{p + n}_{X \times_S X/S})
$$
de la remarque \ref{derham-remark-how-to-use}. Considérons l'image
$$
\gamma^{n - 1, n}(\theta) \in
H^{2n}(X \times_S X, \Omega^{2n - 1}_{X \times_S X})
$$
D'après le lemme \ref{derham-lemma-de-rham-complex-product}, nous avons
$$
\Omega^{2n - 1}_{X \times_S X} =
\Omega^{n - 1}_{X/S} \boxtimes \Omega^n_{X/S} \oplus
\Omega^n_{X/S} \boxtimes \Omega^{n - 1}_{X/S}
$$
La formule de Künneth (Catégories dérivées des schémas, lemme \ref{perfect-lemma-kunneth}, ou
Catégories dérivées des schémas, lemme \ref{perfect-lemma-kunneth-single-sheaf})
donne
$$
H^{2n}(X \times_S X, \Omega^{n - 1}_{X/S} \boxtimes \Omega^n_{X/S}) =
H^n(X, \Omega^{n - 1}_{X/S}) \otimes_A H^n(X, \Omega^n_{X/S})
$$
et
$$
H^{2n}(X \times_S X, \Omega^n_{X/S} \boxtimes \Omega^{n - 1}_{X/S}) =
H^n(X, \Omega^n_{X/S}) \otimes_A H^n(X, \Omega^{n - 1}_{X/S})
$$
En effet, puisque nous nous plaçons en degré maximal, aucun groupe Tor supérieur
n'intervient. Comme $\xi$ est un générateur, cela signifie
que nous pouvons écrire
$$
\gamma^{n - 1, n}(\theta) = \theta_1 \otimes \xi + \xi \otimes \theta_2
$$
avec $\theta_1, \theta_2 \in H^n(X, \Omega^{n - 1}_{X/S})$.
En raisonnant exactement de la même manière, nous pouvons écrire
$$
\gamma^{n, n}(\xi) = b \xi \otimes \xi
$$
dans
$H^{2n}(X \times_S X, \Omega^{2n}_{X \times_S X/S}) =
H^n(X, \Omega^n_{X/S}) \otimes_A H^n(X, \Omega^n_{X/S})$
pour un certain $b \in H^0(S, \mathcal{O}_S)$.

\medskip\noindent
{\bf Affirmation :} $\theta_1 = \theta$, $\theta_2 = \theta$ et $b = 1$.
Montrons que cette affirmation implique le résultat voulu, $a = 0$.
En effet, d'après le lemme \ref{derham-lemma-gysin-differential-hodge},
nous avons
$$
\gamma^{n, n}(\text{d}(\theta)) = \text{d}(\gamma^{n - 1, n}(\theta))
$$
Nos choix précédents donnent alors
$$
a \xi \otimes \xi =
\gamma^{n, n}(a\xi) =
\text{d}(\theta \otimes \xi + \xi \otimes \theta) =
a \xi \otimes \xi + (-1)^n a \xi \otimes \xi
$$
La dernière égalité provient du fait que le morphisme
$\text{d} : \Omega^{2n - 1}_{X \otimes_S X/S} \to \Omega^{2n}_{X \times_S X/S}$
est, d'après le lemme \ref{derham-lemma-de-rham-complex-product},
la somme de la différentielle
$\text{d} \boxtimes 1 : \Omega^{n - 1}_{X/S} \boxtimes \Omega^n_{X/S}
\to \Omega^n_{X/S} \boxtimes \Omega^n_{X/S}$
et de la différentielle
$(-1)^n 1 \boxtimes \text{d} : \Omega^n_{X/S} \boxtimes \Omega^{n - 1}_{X/S}
\to \Omega^n_{X/S} \boxtimes \Omega^n_{X/S}$. Pour plus de détails, voir la discussion dans
la section \ref{derham-section-kunneth} et dans
Catégories dérivées des schémas, section
\ref{perfect-section-kunneth-complexes}.
Puisque $\xi \otimes \xi$ est une base du module libre de rang $1$ sur $A$
$H^n(X, \Omega^n_{X/S}) \otimes_A H^n(X, \Omega^n_{X/S})$,
nous concluons que
$$
a = a + (-1)^n a \Rightarrow a = 0
$$,
comme voulu.

\medskip\noindent
Dans la suite de la démonstration, nous établissons l'affirmation ci-dessus. Notons
$\eta = \gamma^{0, 0}(1) \in H^n(X \times_S X, \Omega^n_{X \times_S X/S})$.
Puisque $\Omega^n_{X \times_S X/S} =
\bigoplus_{p + p' = n} \Omega^p_{X/S} \boxtimes \Omega^{p'}_{X/S}$,
nous pouvons écrire
$$
\eta = \eta_0 + \eta_1 + \ldots + \eta_n
$$,
où $\eta_p$ appartient à
$H^n(X \times_S X, \Omega^p_{X/S} \boxtimes \Omega^{n - p}_{X/S})$.
Pour $p = 0$, nous pouvons écrire
\begin{align*}
H^n(X \times_S X, \mathcal{O}_X \boxtimes \Omega^n_{X/S})
& =
H^n(R\Gamma(X, \mathcal{O}_X) \otimes_A^\mathbf{L}
R\Gamma(X, \Omega^n_{X/S})) \\
& =
A \otimes_A H^n(X, \Omega^n_{X/S}) \oplus
H^n(P \otimes_A^\mathbf{L} R\Gamma(X, \Omega^n_{X/S}))
\end{align*}
grâce à la décomposition $R\Gamma(X, \mathcal{O}_X) = A \oplus P$ obtenue plus haut.
Considérons le morphisme $(s, \text{id}) : X \to X \times_S X$.
Alors $(s, \text{id})^{-1}(\Delta) = Z$ au sens schématique.
Ainsi, $(s, \text{id})^*\eta = \xi$ d'après
le lemme \ref{derham-lemma-gysin-transverse-global}. Cela signifie que
$$
\xi = (s, \text{id})^*\eta = (s^* \otimes \text{id})(\eta_0)
$$
Autrement dit, la première composante de $\eta_0$
dans la décomposition en somme directe ci-dessus est $\xi$. Nous pouvons donc écrire
$$
\eta_0 = 1 \otimes \xi + \eta'_0
$$
avec $\eta'_0 \in H^n(P \otimes_A^\mathbf{L} R\Gamma(X, \Omega^n_{X/S}))$.
Exactement de la même manière, pour $p = n$, nous pouvons écrire
\begin{align*}
H^n(X \times_S X, \Omega^n_{X/S} \boxtimes \mathcal{O}_X)
& =
H^n(R\Gamma(X, \Omega^n_{X/S}) \otimes_A^\mathbf{L}
R\Gamma(X, \mathcal{O}_X)) \\
& =
H^n(X, \Omega^n_{X/S}) \otimes_A A \oplus
H^n(R\Gamma(X, \Omega^n_{X/S}) \otimes_A^\mathbf{L} P)
\end{align*},
puis
$$
\eta_n = \xi \otimes 1 + \eta'_n
$$
avec $\eta'_n \in H^n(R\Gamma(X, \Omega^n_{X/S}) \otimes_A^\mathbf{L} P)$.

\medskip\noindent
Observons que $\text{pr}_1^*\theta = \theta \otimes 1$
et $\text{pr}_2^*\theta = 1 \otimes \theta$ sont des
classes de cohomologie de Hodge sur
$X \times_S X$ dont l'image inverse par la diagonale est $\theta$, cette diagonale étant notée $\Delta$.
D'après le lemme \ref{derham-lemma-gysin-projection-global}, nous avons donc
$$
\theta_1 \otimes \xi + \xi \otimes \theta_2 =
\gamma^{n - 1, n}(\theta) =
(\theta \otimes 1) \cup \eta =
(1 \otimes \theta) \cup \eta
$$
dans l'anneau de cohomologie de Hodge de $X \times_S X$ sur $S$.
En termes de la décomposition en somme directe des modules
de différentielles de $X \times_S X/S$, nous obtenons
$$
\theta_1 \otimes \xi =
(\theta \otimes 1) \cup \eta_0
\quad\text{et}\quad
\xi \otimes \theta_2 =
(1 \otimes \theta) \cup \eta_n
$$
La formule $\eta_0 = 1 \otimes \xi + \eta'_0$ obtenue plus haut
montre que, pour établir $\theta_1 = \theta$, il suffit de prouver
$$
(\theta \otimes 1) \cup \eta'_0 = 0
$$
Pour cela, observons que le cup-produit avec $\theta \otimes 1$ est donné
par l'action en cohomologie du morphisme
$$
(P \otimes_A^\mathbf{L} R\Gamma(X, \Omega^n_{X/S}))[-n]
\xrightarrow{\theta \otimes 1}
R\Gamma(X, \Omega^{n - 1}_{X/S}) \otimes_A^\mathbf{L}
R\Gamma(X, \Omega^n_{X/S})
$$
dans la catégorie dérivée ; voir Cohomologie, remarque
\ref{cohomology-remark-cup-with-element-map-total-cohomology}.
Ce morphisme est le produit tensoriel dérivé des deux morphismes
$$
\theta : P[-n] \to R\Gamma(X, \Omega^{n - 1}_{X/S})
\quad\text{et}\quad
1 : R\Gamma(X, \Omega^n_{X/S}) \to R\Gamma(X, \Omega^n_{X/S})
$$
d'après Catégories dérivées des schémas, remarque
\ref{perfect-remark-annoying-compatibility}.
Or le premier est nul dans $D(A)$, car il va
d'un complexe parfait d'amplitude de Tor dans $[n + 1, 2n]$ vers un complexe
dont la cohomologie n'est éventuellement non nulle qu'en degrés $0, 1, \ldots, n$ ; voir
Compléments d'algèbre, lemme \ref{more-algebra-lemma-splitting-unique}.
Un raisonnement analogue établit l'annulation de
$(1 \otimes \theta) \cup \eta'_n$. Enfin,
exactement de la même manière, nous obtenons
$$
b \xi \otimes \xi = \gamma^{n, n}(\xi) = (\xi \otimes 1) \cup \eta_0
$$,
et nous concluons comme précédemment en montrant que
$(\xi \otimes 1) \cup \eta'_0 = 0$ par le même raisonnement que ci-dessus.
Cela achève la démonstration.
\end{proof}

\begin{proposition}
\label{derham-proposition-relative-poincare-duality}
Soit $S$ un schéma quasi-compact et quasi-séparé. Soit $f : X \to S$
un morphisme propre et lisse de schémas dont toutes les fibres sont non vides
et équidimensionnelles de dimension $n$. Il existe un morphisme de
$\mathcal{O}_S$-modules
$$
t : R^{2n}f_*\Omega^\bullet_{X/S} \longrightarrow \mathcal{O}_S
$$,
unique à précomposition près par la multiplication par un élément inversible de
$H^0(X, \mathcal{O}_X)$, ayant la propriété suivante : l'accouplement
$$
Rf_*\Omega^\bullet_{X/S}
\otimes_{\mathcal{O}_S}^\mathbf{L}
Rf_*\Omega^\bullet_{X/S}[2n]
\longrightarrow
\mathcal{O}_S, \quad
(\xi, \xi') \longmapsto t(\xi \cup \xi')
$$
est un accouplement parfait de complexes parfaits sur $S$.
\end{proposition}

\begin{proof}
La démonstration est exactement la même que celle de la
proposition \ref{derham-proposition-poincare-duality}.

\medskip\noindent
La suite spectrale relative de Hodge vers de Rham
$$
E_1^{p, q} = R^qf_*\Omega^p_{X/S} \Rightarrow R^{p + q}f_*\Omega^\bullet_{X/S}
$$
(section \ref{derham-section-hodge-to-de-rham}), l'annulation
de $\Omega^i_{X/S}$ pour $i > n$, l'annulation donnée, par exemple, dans Limites, lemme
\ref{limits-lemma-higher-direct-images-zero-above-dimension-fibre},
et les résultats des lemmes \ref{derham-lemma-relative-bottom-part-degenerates} et
\ref{derham-lemma-relative-top-part-degenerates}
montrent que $R^0f_*\Omega_{X/S} = R^0f_*\mathcal{O}_X$
et $R^nf_*\Omega^n_{X/S} = R^{2n}f_*\Omega^\bullet_{X/S}$.
Plus précisément, ces identifications proviennent des morphismes
de complexes
$$
\Omega^\bullet_{X/S} \to \mathcal{O}_X[0]
\quad\text{et}\quad
\Omega^n_{X/S}[-n] \to \Omega^\bullet_{X/S}
$$
Choisissons $t : R^{2n}f_*\Omega_{X/S} \to \mathcal{O}_S$
qui corresponde, via cette identification, à un morphisme $t$ comme dans le
lemme \ref{derham-lemma-relative-duality-hodge}.

\medskip\noindent
Abrégeons les notations en posant $\Omega^\bullet = \Omega^\bullet_{X/S}$.
Considérons le morphisme (\ref{derham-equation-wedge}), qui s'écrit ici
$$
\wedge :
\text{Tot}(\Omega^\bullet \otimes_{f^{-1}\mathcal{O}_S} \Omega^\bullet)
\longrightarrow
\Omega^\bullet
$$
Pour chaque entier $p = 0, 1, \ldots, n$, ce morphisme s'annule sur le sous-complexe
$\text{Tot}(\sigma_{> p} \Omega^\bullet \otimes_{f^{-1}\mathcal{O}_S}
\sigma_{\geq n - p} \Omega^\bullet)$ pour des raisons de degré.
La restriction de $\wedge$ au sous-complexe
$\text{Tot}(\Omega^\bullet \otimes_{f^{-1}\mathcal{O}_S}
\geq_{n - p}\Omega^\bullet)$ se factorise donc par un morphisme de complexes
$$
\gamma_p :
\text{Tot}(\sigma_{\leq p} \Omega^\bullet \otimes_{f^{-1}\mathcal{O}_S}
\sigma_{\geq n - p} \Omega^\bullet)
\longrightarrow
\Omega^\bullet
$$
En procédant comme dans la section \ref{derham-section-cup-product}, nous obtenons
des cup-produits relatifs
$$
Rf_*\sigma_{\leq p} \Omega^\bullet
\otimes_{\mathcal{O}_S}^\mathbf{L}
Rf_*\sigma_{\geq n - p}\Omega^\bullet
\longrightarrow
Rf_*\Omega^\bullet
$$
Nous montrerons par récurrence sur $p$ que ces cup-produits, composés avec $t$,
induisent des accouplements parfaits entre $Rf_*\sigma_{\leq p} \Omega^\bullet$
et $Rf_*\sigma_{\geq n - p}\Omega^\bullet[2n]$. Pour $p = n$,
il s'agit de l'assertion de la proposition.

\medskip\noindent
Le cas initial est $p = 0$. Nous avons alors
$$
Rf_*\sigma_{\leq p}\Omega^\bullet = Rf_*\mathcal{O}_X
\quad\text{et}\quad
Rf_*\sigma_{\geq n - p}\Omega^\bullet[2n] = Rf_*(\Omega^n[-n])[2n] =
Rf_*\Omega^n[n]
$$
Nous obtenons simplement l'accouplement
entre $Rf_*\mathcal{O}_X$ et $Rf_*\Omega^n[n]$ du
lemme \ref{derham-lemma-relative-duality-hodge}, et le résultat est donc vrai.

\medskip\noindent
Étape de récurrence. Supposons le résultat vrai pour $p$.
Considérons alors le triangle distingué
$$
\Omega^{p + 1}[-p - 1] \to
\sigma_{\leq p + 1}\Omega^\bullet \to
\sigma_{\leq p}\Omega^\bullet \to
\Omega^{p + 1}[-p]
$$
et le triangle distingué
$$
\sigma_{\geq n - p}\Omega^\bullet \to
\sigma_{\geq n - p - 1}\Omega^\bullet \to
\Omega^{n - p - 1}[-n + p + 1] \to
(\sigma_{\geq n - p}\Omega^\bullet)[1]
$$
Observons que ce sont deux triangles distingués dans la catégorie homotopique
des complexes de faisceaux de $f^{-1}\mathcal{O}_S$-modules ; en particulier,
les morphismes $\sigma_{\leq p}\Omega^\bullet \to \Omega^{p + 1}[-p]$ et
$\Omega^{n - p - 1}[-d + p + 1] \to (\sigma_{\geq n - p}\Omega^\bullet)[1]$
sont donnés par de véritables morphismes de complexes, à savoir à l'aide de la différentielle
$\Omega^p \to \Omega^{p + 1}$ et de la différentielle
$\Omega^{n - p - 1} \to \Omega^{n - p}$.
Considérons les triangles distingués obtenus à partir de ceux-ci
en appliquant $Rf_*$ :
$$
\xymatrix{
Rf_*\sigma_{\leq p}\Omega^\bullet \ar[d]_a \\
Rf_*\Omega^{p + 1}[-p - 1] \ar[d]_b \\
Rf_*\sigma_{\leq p + 1}\Omega^\bullet \ar[d]_c \\
Rf_*\sigma_{\leq p}\Omega^\bullet \ar[d]_d \\
Rf_*\Omega^{p + 1}[-p - 1]
}
\quad\quad
\xymatrix{
Rf_*\sigma_{\geq n - p}\Omega^\bullet \\
Rf_*\Omega^{n - p - 1}[-n + p + 1] \ar[u]_{a'} \\
Rf_*\sigma_{\geq n - p - 1}\Omega^\bullet \ar[u]_{b'} \\
Rf_*\sigma_{\geq n - p}\Omega^\bullet \ar[u]_{c'} \\
Rf_*\Omega^{n - p - 1}[-n + p + 1] \ar[u]_{d'}
}
$$
Nous montrerons ci-dessous que les couples $(a, a')$, $(b, b')$, $(c, c')$ et
$(d, d')$ sont compatibles avec les accouplements donnés. Nous obtenons donc un
morphisme du triangle distingué de gauche vers le triangle distingué
obtenu en appliquant $R\SheafHom(-, \mathcal{O}_S)$ au triangle distingué
de droite. Par récurrence et d'après le lemme \ref{derham-lemma-duality-hodge},
les accouplements construits ci-dessus entre les
complexes des première, deuxième, quatrième et cinquième
lignes sont parfaits, c'est-à-dire qu'ils définissent des isomorphismes après dualisation.
D'après Catégories dérivées, lemme \ref{derived-lemma-third-isomorphism-triangle},
nous en concluons que l'accouplement entre les complexes de la ligne médiane
est parfait, comme voulu.

\medskip\noindent
Soient $e : K \to K'$ et $e' : M' \to M$ des morphismes d'objets
de $D(\mathcal{O}_S)$, et soient
$K \otimes_{\mathcal{O}_S}^\mathbf{L} M \to \mathcal{O}_S$ et
$K' \otimes_{\mathcal{O}_S}^\mathbf{L} M' \to \mathcal{O}_S$
des accouplements. Nous disons qu'ils sont compatibles si le
diagramme
$$
\xymatrix{
K' \otimes_{\mathcal{O}_S}^\mathbf{L} M' \ar[d] &
K \otimes_{\mathcal{O}_S}^\mathbf{L} M'
\ar[l]^{e \otimes 1} \ar[d]^{1 \otimes e'} \\
\mathcal{O}_S &
K \otimes_{\mathcal{O}_S}^\mathbf{L} M \ar[l]
}
$$
est commutatif. Cela signifie en effet que le diagramme
$$
\xymatrix{
K \ar[r] \ar[d]_e & R\SheafHom(M, \mathcal{O}_S) \ar[d]^{R\SheafHom(e', -)} \\
K' \ar[r] & R\SheafHom(M', \mathcal{O}_S)
}
$$
est commutatif, ce qui suffit pour nos besoins.

\medskip\noindent
Démontrons cette compatibilité pour le couple $(c, c')$. Il suffit d'observer
que nous avons un diagramme commutatif
$$
\xymatrix{
\text{Tot}(\sigma_{\leq p} \Omega^\bullet \otimes_{f^{-1}\mathcal{O}_S}
\sigma_{\geq n - p} \Omega^\bullet) \ar[d]_{\gamma_p} &
\text{Tot}(\sigma_{\leq p + 1} \Omega^\bullet \otimes_{f^{-1}\mathcal{O}_S}
\sigma_{\geq n - p} \Omega^\bullet) \ar[l] \ar[d] \\
\Omega^\bullet &
\text{Tot}(\sigma_{\leq p + 1} \Omega^\bullet \otimes_{f^{-1}\mathcal{O}_S}
\sigma_{\geq n - p - 1} \Omega^\bullet) \ar[l]_-{\gamma_{p + 1}}
}
$$
La fonctorialité du cup-produit donne alors la commutativité du
diagramme voulu.

\medskip\noindent
De même, pour le couple $(b, b')$, nous utilisons le diagramme commutatif
$$
\xymatrix{
\text{Tot}(\sigma_{\leq p + 1} \Omega^\bullet \otimes_{f^{-1}\mathcal{O}_S}
\sigma_{\geq n - p - 1} \Omega^\bullet) \ar[d]_{\gamma_{p + 1}} &
\text{Tot}(\Omega^{p + 1}[-p - 1] \otimes_{f^{-1}\mathcal{O}_S}
\sigma_{\geq n - p - 1} \Omega^\bullet) \ar[l] \ar[d] \\
\Omega^\bullet &
\Omega^{p + 1}[-p - 1]
\otimes_{f^{-1}\mathcal{O}_S}
\Omega^{n - p - 1}[-n + p + 1] \ar[l]_-\wedge
}
$$

\medskip\noindent
Pour les couples $(d, d')$ et $(a, a')$, nous utilisons le diagramme commutatif
$$
\xymatrix{
\Omega^{p + 1}[-p] \otimes_{f^{-1}\mathcal{O}_S}
\Omega^{n - p - 1}[-n + p] \ar[d] &
\text{Tot}(\sigma_{\leq p}\Omega^\bullet \otimes_{f^{-1}\mathcal{O}_S}
\Omega^{n - p - 1}[-n + p]) \ar[l] \ar[d] \\
\Omega^\bullet &
\text{Tot}(\sigma_{\leq p}\Omega^\bullet \otimes_{f^{-1}\mathcal{O}_S}
\sigma_{\geq n - p}\Omega^\bullet) \ar[l]
}
$$

\medskip\noindent
Nous omettons la démonstration de l'unicité de $t$ à
précomposition près par la multiplication par un élément inversible de $H^0(X, \mathcal{O}_X)$.
\end{proof}















% OK, start here.
%

\chapter{Cohomologie locale}



\phantomsection
\label{local-cohomology-section-phantom}


\section{Introduction}
\label{local-cohomology-section-introduction}

\noindent
Ce chapitre poursuit l'étude de la cohomologie locale.
Une référence est \cite{SGA2}.
La définition de la cohomologie locale se trouve dans
Complexes dualisants, section \ref{dualizing-section-local-cohomology}.
Pour les anneaux noethériens, prendre la cohomologie locale revient
à dériver un foncteur de torsion convenable, comme le montre
Complexes dualisants, section
\ref{dualizing-section-local-cohomology-noetherian}.
La relation avec la profondeur se trouve dans
Complexes dualisants, section
\ref{dualizing-section-depth}.

\medskip\noindent
Nous étudions des propriétés de finitude de la cohomologie locale qui conduisent
à la démonstration d'une version assez générale du
théorème de finitude de Grothendieck; voir Théorème \ref{local-cohomology-theorem-finiteness}
et Lemme \ref{local-cohomology-lemma-finiteness-Rjstar} (images directes supérieures
de modules cohérents par des immersions ouvertes).
Nos méthodes incorporent quelques arguments particulièrement élégants que le lecteur
trouvera dans des articles de Faltings; voir
\cite{Faltings-annulators} et \cite{Faltings-finiteness}.

\medskip\noindent
Comme applications, nous étudions le théorème d'annulation de
Hartshorne--Lichtenbaum. Nous étudions également
l'action de Frobenius et celle des opérateurs différentiels
sur la cohomologie locale.




\section{Généralités}
\label{local-cohomology-section-generalities}

\noindent
Le lemme suivant montre que le foncteur $R\Gamma_Z$
est lié à la cohomologie à supports.

\begin{lemma}
\label{local-cohomology-lemma-local-cohomology-is-local-cohomology}
Soit $A$ un anneau et soit $I$ un idéal de type fini.
Posons $Z = V(I) \subset X = \Spec(A)$. Pour $K \in D(A)$ correspondant
à $\widetilde{K} \in D_\QCoh(\mathcal{O}_X)$ par
Catégories dérivées de schémas, Lemme \ref{perfect-lemma-affine-compare-bounded},
il existe un isomorphisme fonctoriel
$$
R\Gamma_Z(K) = R\Gamma_Z(X, \widetilde{K})
$$
où le membre de gauche est celui de
Complexes dualisants, Équation (\ref{dualizing-equation-local-cohomology}),
et celui de droite est le foncteur de
Cohomologie, section \ref{cohomology-section-cohomology-support-bis}.
\end{lemma}

\begin{proof}
Par Cohomologie, Lemme \ref{cohomology-lemma-triangle-sections-with-support},
il existe un triangle distingué
$$
R\Gamma_Z(X, \widetilde{K})
\to R\Gamma(X, \widetilde{K})
\to R\Gamma(U, \widetilde{K})
\to R\Gamma_Z(X, \widetilde{K})[1]
$$
où $U = X \setminus Z$. Nous savons que $R\Gamma(X, \widetilde{K}) = K$
par Catégories dérivées de schémas, Lemme
\ref{perfect-lemma-affine-compare-bounded}.
Écrivons $I = (f_1, \ldots, f_r)$. Nous obtenons alors un recouvrement ouvert
affine fini $\mathcal{U} : U = D(f_1) \cup \ldots \cup D(f_r)$.
Par Catégories dérivées de schémas, Lemme
\ref{perfect-lemma-alternating-cech-complex-complex-computes-cohomology}
le complexe alterné de {\v C}ech
$\text{Tot}(\check{\mathcal{C}}_{alt}^\bullet(\mathcal{U},
\widetilde{K^\bullet}))$
calcule $R\Gamma(U, \widetilde{K})$, où $K^\bullet$ est un
complexe quelconque de $A$-modules représentant $K$. En déroulant les
définitions, nous obtenons
$$
R\Gamma(U, \widetilde{K}) =
\text{Tot}\left(
K^\bullet \otimes_A
(\prod\nolimits_{i_0} A_{f_{i_0}} \to
\prod\nolimits_{i_0 < i_1} A_{f_{i_0}f_{i_1}} \to
\ldots \to A_{f_1\ldots f_r})\right)
$$
Il est clair que
$K^\bullet = R\Gamma(X, \widetilde{K^\bullet}) \to
R\Gamma(U, \widetilde{K}^\bullet)$
est induit par l'application diagonale de $A$ dans $\prod A_{f_i}$.
Nous en concluons que
$$
R\Gamma_Z(X, \widetilde{K}) =
\text{Tot}\left(
K^\bullet \otimes_A
(A \to \prod\nolimits_{i_0} A_{f_{i_0}} \to
\prod\nolimits_{i_0 < i_1} A_{f_{i_0}f_{i_1}} \to
\ldots \to A_{f_1\ldots f_r})\right)
$$
Par Complexes dualisants, Lemme \ref{dualizing-lemma-local-cohomology-adjoint},
le complexe du membre de droite calcule $R\Gamma_Z(K)$,
ce qui démontre le lemme pour $K$ fixé.

\medskip\noindent
Terminons en montrant que l'isomorphisme peut être rendu fonctoriel en $K$.
Considérons l'application de complexes de $A$-modules
$$
M^\bullet = \text{Tot}\left(
K^\bullet \otimes_A
(A \to \prod\nolimits_{i_0} A_{f_{i_0}} \to
\prod\nolimits_{i_0 < i_1} A_{f_{i_0}f_{i_1}} \to
\ldots \to A_{f_1\ldots f_r})\right)
\longrightarrow
K^\bullet
$$
Les arguments précédents montrent que cette flèche devient un isomorphisme après
application du foncteur $R\Gamma_Z(X, \widetilde{\ })$ et que
$R\Gamma(U, \widetilde{M}^\bullet) = 0$ (détails omis). Nous avons alors
$$
R\Gamma_Z(X, \widetilde{K}^\bullet)
\leftarrow
R\Gamma_Z(X, \widetilde{M}^\bullet)
\rightarrow
R\Gamma(X, \widetilde{M}^\bullet) = M^\bullet = R\Gamma_Z(K^\bullet)
$$
et toutes les applications sont des isomorphismes dans $D(A)$, fonctoriels en $K^\bullet$,
comme souhaité.
\end{proof}

\begin{lemma}
\label{local-cohomology-lemma-local-cohomology}
Soit $A$ un anneau et soit $I \subset A$ un idéal de type fini.
Posons $X = \Spec(A)$, $Z = V(I)$, $U = X \setminus Z$, et soit $j : U \to X$
le morphisme d'inclusion. Soit $\mathcal{F}$ un
$\mathcal{O}_U$-module quasi-cohérent. Alors
\begin{enumerate}
\item il existe un $A$-module $M$ tel que $\mathcal{F}$ soit la
restriction de $\widetilde{M}$ à $U$,
\item pour un tel $M$, il existe une suite exacte
$$
0 \to H^0_Z(M) \to M \to H^0(U, \mathcal{F}) \to H^1_Z(M) \to 0
$$
et des isomorphismes $H^p(U, \mathcal{F}) = H^{p + 1}_Z(M)$ pour $p \geq 1$,
\item on peut prendre $M = H^0(U, \mathcal{F})$, auquel cas
$H^0_Z(M) = H^1_Z(M) = 0$.
\end{enumerate}
\end{lemma}

\begin{proof}
L'existence de $M$ résulte de
Propriétés, Lemme \ref{properties-lemma-extend-trivial},
et du fait que les faisceaux quasi-cohérents sur $X$ correspondent
aux $A$-modules (Schémas, Lemme \ref{schemes-lemma-equivalence-quasi-coherent}).
Considérons ensuite le triangle distingué
$$
R\Gamma_Z(X, \widetilde{M}) \to R\Gamma(X, \widetilde{M}) \to
R\Gamma(U, \widetilde{M}|_U) \to R\Gamma_Z(X, \widetilde{M})[1]
$$
de Cohomologie, Lemme \ref{cohomology-lemma-triangle-sections-with-support}.
Comme $X$ est affine, nous avons $R\Gamma(X, \widetilde{M}) = M$
par Cohomologie des schémas, Lemme
\ref{coherent-lemma-quasi-coherent-affine-cohomology-zero}.
Par notre choix de $M$, nous avons $\mathcal{F} = \widetilde{M}|_U$,
d'où une suite exacte
$$
0 \to H^0_Z(X, \widetilde{M}) \to M \to H^0(U, \mathcal{F}) \to
H^1_Z(X, \widetilde{M}) \to 0
$$
et des isomorphismes $H^p(U, \mathcal{F}) = H^{p + 1}_Z(X, \widetilde{M})$
pour $p \geq 1$. Par Lemme \ref{local-cohomology-lemma-local-cohomology-is-local-cohomology},
nous avons $H^i_Z(M) = H^i_Z(X, \widetilde{M})$ pour tout $i$.
Les assertions (1) et (2) sont donc démontrées.
Enfin, en posant $M' = H^0(U, \mathcal{F})$, nous voyons que
le noyau et le conoyau de $M \to M'$ sont de torsion par rapport aux puissances de $I$.
Par conséquent, $\widetilde{M}|_U \to \widetilde{M'}|_U$ est un isomorphisme
et nous pouvons bien utiliser $M'$ comme annoncé en (3). Il va sans dire
que $H^0_Z(M')$ et $H^0_Z(M')$ sont tous deux nuls.
\end{proof}

\begin{lemma}
\label{local-cohomology-lemma-already-torsion}
Soient $I, J \subset A$ des idéaux de type fini d'un anneau $A$.
Si $M$ est un module de torsion par rapport aux puissances de $I$, alors l'application
canonique
$$
H^i_{V(I) \cap V(J)}(M) \to H^i_{V(J)}(M)
$$
est un isomorphisme pour tout $i$.
\end{lemma}

\begin{proof}
La suite spectrale de
Complexes dualisants, Lemme \ref{dualizing-lemma-local-cohomology-ss},
ramène l'assertion à $R\Gamma_I(M) = M$, qui résulte immédiatement
de la construction de la cohomologie locale dans
Complexes dualisants, section \ref{dualizing-section-local-cohomology}.
\end{proof}

\begin{lemma}
\label{local-cohomology-lemma-multiplicative}
Soit $S \subset A$ une partie multiplicative d'un anneau $A$.
Soit $M$ un $A$-module tel que $S^{-1}M = 0$. Alors
$\colim_{f \in S} H^0_{V(f)}(M) = M$ et
$\colim_{f \in S} H^1_{V(f)}(M) = 0$.
\end{lemma}

\begin{proof}
L'assertion sur $H^0$ résulte directement des définitions.
Pour celle sur $H^1$, remarquons que $R\Gamma_{V(f)}$
et $H^1_{V(f)}$ commutent aux limites inductives. Nous pouvons donc supposer
que $M$ est annulé par un certain $f \in S$. Alors
$H^1_{V(ff')}(M) = 0$ pour tout $f' \in S$ (par exemple par
Lemme \ref{local-cohomology-lemma-already-torsion}).
\end{proof}

\begin{lemma}
\label{local-cohomology-lemma-elements-come-from-bigger}
Soit $I \subset A$ un idéal de type fini d'un anneau $A$.
Soit $\mathfrak p$ un idéal premier. Soit $M$ un $A$-module.
Soit $i \geq 0$ un entier et considérons l'application
$$
\Psi :
\colim_{f \in A, f \not \in \mathfrak p} H^i_{V((I, f))}(M)
\longrightarrow
H^i_{V(I)}(M)
$$
Alors
\begin{enumerate}
\item $\Im(\Psi)$ est l'ensemble des éléments dont l'image est nulle dans
$H^i_{V(I)}(M)_\mathfrak p$,
\item si $H^{i - 1}_{V(I)}(M)_\mathfrak p = 0$, alors $\Psi$ est injective,
\item si $H^{i - 1}_{V(I)}(M)_\mathfrak p = H^i_{V(I)}(M)_\mathfrak p = 0$,
alors $\Psi$ est un isomorphisme.
\end{enumerate}
\end{lemma}

\begin{proof}
Pour $f \in A$, $f \not \in \mathfrak p$, la suite spectrale de
Complexes dualisants, Lemme \ref{dualizing-lemma-local-cohomology-ss},
dégénère et fournit des suites exactes courtes
$$
0 \to H^1_{V(f)}(H^{i - 1}_{V(I)}(M)) \to
H^i_{V((I, f))}(M) \to H^0_{V(f)}(H^i_{V(I)}(M)) \to 0
$$
Ceci démontre (1), et (2) en résulte avec
Lemme \ref{local-cohomology-lemma-multiplicative}.
L'assertion (3) en est une conséquence formelle.
\end{proof}

\begin{lemma}
\label{local-cohomology-lemma-isomorphism}
Soient $I \subset I' \subset A$ des idéaux de type fini d'un
anneau noethérien $A$. Soit $M$ un $A$-module. Soit $i \geq 0$ un entier.
Considérons l'application
$$
\Psi : H^i_{V(I')}(M) \to H^i_{V(I)}(M)
$$
Les assertions suivantes sont vraies :
\begin{enumerate}
\item si $H^i_{\mathfrak pA_\mathfrak p}(M_\mathfrak p) = 0$
pour tout $\mathfrak p \in V(I) \setminus V(I')$, alors
$\Psi$ est surjective,
\item si $H^{i - 1}_{\mathfrak pA_\mathfrak p}(M_\mathfrak p) = 0$
pour tout $\mathfrak p \in V(I) \setminus V(I')$, alors
$\Psi$ est injective,
\item si $H^i_{\mathfrak pA_\mathfrak p}(M_\mathfrak p) =
H^{i - 1}_{\mathfrak pA_\mathfrak p}(M_\mathfrak p) = 0$
pour tout $\mathfrak p \in V(I) \setminus V(I')$, alors
$\Psi$ est un isomorphisme.
\end{enumerate}
\end{lemma}

\begin{proof}
Démonstration de (1).
Soit $\xi \in H^i_{V(I)}(M)$. Comme $A$ est noethérien, il existe un
plus grand idéal $I \subset I'' \subset I'$ tel que $\xi$ soit l'image
d'un certain $\xi'' \in H^i_{V(I'')}(M)$. Si $V(I'') = V(I')$, nous avons
terminé. Sinon, choisissons un point générique $\mathfrak p \in V(I'')$ hors de $V(I')$.
Nous avons alors $H^i_{V(I'')}(M)_\mathfrak p =
H^i_{\mathfrak pA_\mathfrak p}(M_\mathfrak p) = 0$ par hypothèse.
Par Lemme \ref{local-cohomology-lemma-elements-come-from-bigger}, nous pouvons agrandir $I''$,
ce qui contredit sa maximalité.

\medskip\noindent
Démonstration de (2). Soit $\xi' \in H^i_{V(I')}(M)$ dans le noyau de $\Psi$.
Comme $A$ est noethérien, il existe un
plus grand idéal $I \subset I'' \subset I'$ tel que l'image de $\xi'$
soit nulle dans $H^i_{V(I'')}(M)$. Si $V(I'') = V(I')$, nous avons
terminé. Sinon, choisissons un point générique $\mathfrak p  \in V(I'')$
hors de $V(I')$. Nous avons alors $H^{i - 1}_{V(I'')}(M)_\mathfrak p =
H^{i - 1}_{\mathfrak pA_\mathfrak p}(M_\mathfrak p) = 0$ par hypothèse.
Par Lemme \ref{local-cohomology-lemma-elements-come-from-bigger}, nous pouvons agrandir $I''$,
ce qui contredit sa maximalité.

\medskip\noindent
L'assertion (3) résulte formellement de (1) et (2).
\end{proof}





\section{Lemme de connexité de Hartshorne}
\label{local-cohomology-section-hartshorne-connectedness}

\noindent
Le titre de cette section renvoie au résultat suivant.

\begin{lemma}
\label{local-cohomology-lemma-depth-2-connected-punctured-spectrum}
\begin{reference}
\cite[Proposition 2.1]{Hartshorne-connectedness}
\end{reference}
\begin{slogan}
Connexité de Hartshorne
\end{slogan}
Soit $A$ un anneau local noethérien de profondeur $\geq 2$.
Alors les spectres épointés de $A$, $A^h$ et $A^{sh}$ sont connexes.
\end{lemma}

\begin{proof}
Soit $U$ le spectre épointé de $A$.
Si $U$ n'est pas connexe, alors
$\Gamma(U, \mathcal{O}_U)$ possède un idempotent non trivial.
Mais $A$, étant local, ne possède pas d'idempotent non trivial.
Ainsi $A \to \Gamma(U, \mathcal{O}_U)$ n'est pas un isomorphisme.
Par Lemme \ref{local-cohomology-lemma-local-cohomology},
nous en concluons que $H^0_\mathfrak m(A)$ ou $H^1_\mathfrak m(A)$
est non nul. Par conséquent, $\text{depth}(A) \leq 1$ par
Complexes dualisants, Lemme \ref{dualizing-lemma-depth}.
Pour obtenir le résultat pour $A^h$ et $A^{sh}$, utiliser
Compléments d'algèbre, Lemme \ref{more-algebra-lemma-henselization-depth}.
\end{proof}

\begin{lemma}
\label{local-cohomology-lemma-catenary-S2-equidimensional}
\begin{reference}
\cite[IV Corollary 5.10.9]{EGA}
\end{reference}
Soit $A$ un anneau local noethérien, caténaire et $(S_2)$.
Alors $\Spec(A)$ est équidimensionnel.
\end{lemma}

\begin{proof}
Posons $X = \Spec(A)$ et $d = \dim(A) = \dim(X)$. Dans $X$, considérons la
réunion $X_1$ des composantes irréductibles de dimension $d$ et la réunion
$X_2$ des composantes irréductibles de dimension $< d$. Bien entendu,
$X = X_1 \cup X_2$. Si $X_2 = \emptyset$,
le lemme est démontré. Sinon, $Z = X_1 \cap X_2$ est un fermé non vide
de $X$, car il contient au moins le point fermé de $X$.
Nous pouvons donc choisir un point générique $z \in Z$ d'une composante irréductible
de $Z$. Rappelons que le spectre de $\mathcal{O}_{Z, z}$ est l'ensemble des points
de $X$ qui se spécialisent en $z$. Puisque $z$ appartient à la fois à une
composante irréductible de dimension $d$ et à une composante irréductible
de dimension $< d$, nous obtenons des spécialisations non triviales $x_1 \leadsto z$ et
$x_2 \leadsto z$ telles que les adhérences de $x_1$ et $x_2$ aient des
dimensions différentes. Comme $X$ est caténaire, cela n'est possible que si au moins
l'une des spécialisations $x_1 \leadsto z$ et $x_2 \leadsto z$ n'est pas
immédiate. Ainsi $\dim(\mathcal{O}_{Z, z}) \geq 2$. Par conséquent,
$\text{depth}(\mathcal{O}_{Z, z}) \geq 2$, puisque $A$ est $(S_2)$.
Cependant, le spectre épointé $U$ de $\mathcal{O}_{Z, z}$ n'est pas connexe,
car les fermés $U \cap X_1$ et $U \cap X_2$ sont disjoints
(par notre choix de $z$) et recouvrent $U$. Cela contredit
Lemme \ref{local-cohomology-lemma-depth-2-connected-punctured-spectrum},
et la démonstration est achevée.
\end{proof}



\section{Dimension cohomologique}
\label{local-cohomology-section-cd}

\noindent
Une brève section sur la dimension cohomologique.

\begin{lemma}
\label{local-cohomology-lemma-cd}
Soit $I \subset A$ un idéal de type fini d'un anneau $A$.
Posons $Y = V(I) \subset X = \Spec(A)$. Soit $d \geq -1$ un entier.
Les conditions suivantes sont équivalentes :
\begin{enumerate}
\item $H^i_Y(A) = 0$ pour $i > d$,
\item $H^i_Y(M) = 0$ pour $i > d$ et tout $A$-module $M$, et
\item si $d = -1$, alors $Y = \emptyset$; si $d = 0$, alors
$Y$ est ouvert et fermé dans $X$; et si $d > 0$, alors
$H^i(X \setminus Y, \mathcal{F}) = 0$ pour $i \geq d$
pour tout $\mathcal{O}_{X \setminus Y}$-module quasi-cohérent $\mathcal{F}$.
\end{enumerate}
\end{lemma}

\begin{proof}
Remarquons que $R\Gamma_Y(-)$ est de dimension cohomologique finie, par exemple par
Complexes dualisants, Lemme \ref{dualizing-lemma-local-cohomology-adjoint}.
Il existe donc un entier $i_0$ tel que
$H^i_Y(M) = 0$ pour tout $A$-module $M$ et tout $i \geq i_0$.

\medskip\noindent
Montrons que (1) et (2) sont équivalentes. Il est immédiat que
(2) implique (1). Supposons (1). Par récurrence descendante sur $i > d$,
nous allons montrer que $H^i_Y(M) = 0$ pour tout $A$-module $M$.
Nous l'avons vu ci-dessus pour $i \geq i_0$. Pour l'étape de récurrence,
soit $i_0 > i > d$. Choisissons un $A$-module $M$ quelconque et insérons-le dans
une suite exacte courte $0 \to N \to F \to M \to 0$, où $F$ est un
$A$-module libre. Comme $R\Gamma_Y$ est un adjoint à droite,
$H^i_Y(-)$ commute aux sommes directes. Ainsi $H^i_Y(F) = 0$
puisque $i > d$, par l'hypothèse (1). Il s'ensuit que
$H^i_Y(M) = H^{i + 1}_Y(N) = 0$, comme voulu.

\medskip\noindent
Supposons $d = -1$ et (2) satisfaite. Alors $0 = H^0_Y(A/I) = A/I \Rightarrow A = I
\Rightarrow Y = \emptyset$. Ainsi (3) est satisfaite. Nous omettons la réciproque.

\medskip\noindent
Supposons $d = 0$ et (2) satisfaite. Posons
$J = H^0_I(A) = \{x \in A \mid I^nx = 0 \text{ pour un certain }n > 0\}$.
Alors
$$
H^1_Y(A) = \Coker(A \to \Gamma(X \setminus Y, \mathcal{O}_{X \setminus Y}))
\quad\text{et}\quad
H^1_Y(I) = \Coker(I \to \Gamma(X \setminus Y, \mathcal{O}_{X \setminus Y}))
$$
et le noyau de la première application est égal à $J$. Voir
Lemme \ref{local-cohomology-lemma-local-cohomology}.
Nous déduisons de (2) que $I(A/J) = A/J$.
Nous pouvons donc choisir $f \in I$
d'image $1$ dans $A/J$. Alors $1 - f \in J$, donc $I^n(1 - f) = 0$ pour un certain
$n > 0$. Ainsi $f^n = f^{n + 1}$. Donc $e = f^n \in I$ est un idempotent.
Considérons l'idempotent complémentaire $e' = 1 - f^n \in J$.
Pour tout élément $g \in I$, nous avons $g^m e' = 0$ pour un certain $m > 0$.
Ainsi $I$ est contenu dans le radical de l'idéal $(e) \subset I$.
Cela signifie que $Y = V(I) = V(e)$ est ouvert et fermé dans $X$, comme en (3).
Réciproquement, si $Y = V(I)$ est ouvert et fermé, le foncteur
$H^0_Y(-)$ est exact et ses foncteurs dérivés supérieurs sont nuls.

\medskip\noindent
Si $d > 0$, Lemme \ref{local-cohomology-lemma-local-cohomology} montre immédiatement
que (2) est équivalente à (3).
\end{proof}

\begin{definition}
\label{local-cohomology-definition-cd}
Soit $I \subset A$ un idéal de type fini d'un anneau $A$.
Le plus petit entier $d \geq -1$ satisfaisant les conditions équivalentes
de Lemme \ref{local-cohomology-lemma-cd} est appelé la
{\it dimension cohomologique de $I$ dans $A$} et est
noté $\text{cd}(A, I)$.
\end{definition}

\noindent
Ainsi $\text{cd}(A, I) = -1$ si
$I = A$, et $\text{cd}(A, I) = 0$ si $I$ est localement nilpotent
ou engendré par un idempotent.
Le lemme suivant montre que $\text{cd}(A, I)$ existe.

\begin{lemma}
\label{local-cohomology-lemma-bound-cd}
Soit $I \subset A$ un idéal de type fini d'un anneau $A$.
Alors
\begin{enumerate}
\item $\text{cd}(A, I)$ est au plus égal au nombre de
générateurs de $I$,
\item $\text{cd}(A, I) \leq r$ s'il existe $f_1, \ldots, f_r \in A$
tels que $V(f_1, \ldots, f_r) = V(I)$,
\item $\text{cd}(A, I) \leq c$ si $\Spec(A) \setminus V(I)$
peut être recouvert par $c$ ouverts affines.
\end{enumerate}
\end{lemma}

\begin{proof}
La description explicite de $R\Gamma_Y(-)$ donnée dans
Complexes dualisants, Lemme \ref{dualizing-lemma-local-cohomology-adjoint},
montre (1). Nous pouvons déduire (2) de (1) en utilisant le
fait que $R\Gamma_Z$ ne dépend que du fermé
$Z$, et non du choix de l'idéal de type fini
$I \subset A$ tel que $V(I) = Z$. Cela résulte soit de la
construction de la cohomologie locale dans
Complexes dualisants, section \ref{dualizing-section-local-cohomology},
combinée avec
Compléments d'algèbre, Lemme \ref{more-algebra-lemma-local-cohomology-closed},
soit de Lemme \ref{local-cohomology-lemma-local-cohomology-is-local-cohomology}.
Pour (3), nous utilisons Lemme \ref{local-cohomology-lemma-cd}
et le résultat d'annulation de Cohomologie des schémas, Lemme
\ref{coherent-lemma-vanishing-nr-affines}.
\end{proof}

\begin{lemma}
\label{local-cohomology-lemma-cd-sum}
Soient $I, J \subset A$ des idéaux de type fini d'un anneau $A$.
Alors $\text{cd}(A, I + J) \leq \text{cd}(A, I) + \text{cd}(A, J)$.
\end{lemma}

\begin{proof}
Utiliser la définition et Complexes dualisants, Lemme
\ref{dualizing-lemma-local-cohomology-ss}.
\end{proof}

\begin{lemma}
\label{local-cohomology-lemma-cd-change-rings}
Soit $A \to B$ un homomorphisme d'anneaux. Soit $I \subset A$ un idéal de type fini.
Alors $\text{cd}(B, IB) \leq \text{cd}(A, I)$. Si $A \to B$ est fidèlement
plat, l'égalité est satisfaite.
\end{lemma}

\begin{proof}
Utiliser la définition et
Complexes dualisants, Lemme \ref{dualizing-lemma-torsion-change-rings}.
\end{proof}

\begin{lemma}
\label{local-cohomology-lemma-cd-local}
Soit $I \subset A$ un idéal de type fini d'un anneau $A$.
Alors $\text{cd}(A, I) = \max \text{cd}(A_\mathfrak p, I_\mathfrak p)$.
\end{lemma}

\begin{proof}
Posons $Y = V(I)$ et $Y' = V(I_\mathfrak p) \subset \Spec(A_\mathfrak p)$.
Rappelons que
$R\Gamma_Y(A) \otimes_A A_\mathfrak p = R\Gamma_{Y'}(A_\mathfrak p)$
par Complexes dualisants, Lemme \ref{dualizing-lemma-torsion-change-rings}.
Nous concluons donc par Algèbre, Lemme \ref{algebra-lemma-characterize-zero-local}.
\end{proof}

\begin{lemma}
\label{local-cohomology-lemma-cd-dimension}
Soit $I \subset A$ un idéal de type fini d'un anneau $A$.
Si $M$ est un $A$-module de type fini, alors
$H^i_{V(I)}(M) = 0$ pour $i > \dim(\text{Supp}(M))$.
En particulier, $\text{cd}(A, I) \leq \dim(A)$.
\end{lemma}

\begin{proof}
Nous démontrons d'abord la seconde assertion.
Rappelons que $\dim(A)$ désigne la dimension de Krull. Par
Lemme \ref{local-cohomology-lemma-cd-local}, nous pouvons supposer $A$ local.
Si $V(I) = \emptyset$, le résultat est vrai.
Si $V(I) \not = \emptyset$, alors
$\dim(\Spec(A) \setminus V(I)) < \dim(A)$, car
le point fermé manque. Remarquons que
$U = \Spec(A) \setminus V(I)$ est un ouvert quasi-compact
de l'espace spectral $\Spec(A)$, donc lui-même un espace spectral.
Voir Algèbre, Lemme \ref{algebra-lemma-spec-spectral}, et
Topologie, Lemme \ref{topology-lemma-spectral-sub}.
Ainsi Cohomologie, Proposition
\ref{cohomology-proposition-cohomological-dimension-spectral}
implique $H^i(U, \mathcal{F}) = 0$ pour $i \geq \dim(A)$,
ce qui donne le résultat par Lemme \ref{local-cohomology-lemma-cd}.
Dans le cas noethérien, le lecteur peut utiliser
Cohomologie de Grothendieck, Proposition
\ref{cohomology-proposition-vanishing-Noetherian}.

\medskip\noindent
Nous allons déduire la première assertion de la seconde.
Soit $\mathfrak a$ l'annulateur du $A$-module de type fini $M$.
Posons $B = A/\mathfrak a$. Rappelons que $\Spec(B) = \text{Supp}(M)$; voir
Algèbre, Lemme \ref{algebra-lemma-support-closed}.
Posons $J = IB$. Alors $M$ est un $B$-module
et $H^i_{V(I)}(M) = H^i_{V(J)}(M)$; voir
Complexes dualisants, Lemme
\ref{dualizing-lemma-local-cohomology-and-restriction}.
Comme $\text{cd}(B, J) \leq \dim(B) = \dim(\text{Supp}(M))$
par la première partie, nous concluons.
\end{proof}

\begin{lemma}
\label{local-cohomology-lemma-cd-is-one}
Soit $I \subset A$ un idéal de type fini d'un anneau $A$. Si
$\text{cd}(A, I) = 1$, alors $\Spec(A) \setminus V(I)$ est affine non vide.
\end{lemma}

\begin{proof}
Cela résulte de Lemme \ref{local-cohomology-lemma-cd} et de
Cohomologie des schémas, Lemme
\ref{coherent-lemma-quasi-compact-h1-zero-covering}.
\end{proof}

\begin{lemma}
\label{local-cohomology-lemma-cd-maximal}
Soit $(A, \mathfrak m)$ un anneau local noethérien de dimension $d$.
Alors $H^d_\mathfrak m(A)$ est non nul et $\text{cd}(A, \mathfrak m) = d$.
\end{lemma}

\begin{proof}
Par l'une des caractérisations de la dimension, il existe
un idéal de définition de $A$ engendré par $d$ éléments; voir
Algèbre, Proposition \ref{algebra-proposition-dimension}.
Ainsi $\text{cd}(A, \mathfrak m) \leq d$ par
Lemme \ref{local-cohomology-lemma-bound-cd}. Donc $H^d_\mathfrak m(A)$ est
non nul si et seulement si $\text{cd}(A, \mathfrak m) = d$, si et seulement si
$\text{cd}(A, \mathfrak m) \geq d$.

\medskip\noindent
Soit $A \to A^\wedge$ l'application de $A$ dans son complété.
Remarquons que $A^\wedge$ est un anneau local noethérien de
même dimension que $A$, d'idéal maximal $\mathfrak m A^\wedge$.
Voir Algèbre, Lemmes
\ref{algebra-lemma-completion-Noetherian-Noetherian},
\ref{algebra-lemma-completion-complete},
\ref{algebra-lemma-completion-faithfully-flat} et
Compléments d'algèbre, Lemme \ref{more-algebra-lemma-completion-dimension}.
Par Lemme \ref{local-cohomology-lemma-cd-change-rings},
il suffit de démontrer le lemme pour $A^\wedge$.

\medskip\noindent
Par le paragraphe précédent, nous pouvons supposer que $A$ est
un anneau local complet. Alors $A$ possède un complexe dualisant normalisé
$\omega_A^\bullet$ (Complexes dualisants, Lemme
\ref{dualizing-lemma-ubiquity-dualizing}).
Le théorème de dualité locale (sous la forme de
Complexes dualisants, Lemme \ref{dualizing-lemma-special-case-local-duality})
nous dit que $H^d_\mathfrak m(A)$ est le dual de Matlis de
$\text{Ext}^{-d}(A, \omega_A^\bullet) = H^{-d}(\omega_A^\bullet)$
qui est non nul, par exemple par
Complexes dualisants, Lemme
\ref{dualizing-lemma-nonvanishing-generically-local}.
\end{proof}

\begin{lemma}
\label{local-cohomology-lemma-cd-bound-dim-local}
Soit $(A, \mathfrak m)$ un anneau local noethérien.
Soit $I \subset A$ un idéal propre.
Soit $\mathfrak p \subset A$ un idéal premier
tel que $V(\mathfrak p) \cap V(I) = \{\mathfrak m\}$.
Alors $\dim(A/\mathfrak p) \leq \text{cd}(A, I)$.
\end{lemma}

\begin{proof}
Par Lemme \ref{local-cohomology-lemma-cd-change-rings}, nous avons
$\text{cd}(A, I) \geq \text{cd}(A/\mathfrak p, I(A/\mathfrak p))$.
Comme $V(I) \cap V(\mathfrak p) = \{\mathfrak m\}$, nous avons
$\text{cd}(A/\mathfrak p, I(A/\mathfrak p)) =
\text{cd}(A/\mathfrak p, \mathfrak m/\mathfrak p)$.
Par Lemme \ref{local-cohomology-lemma-cd-maximal}, ceci est égal à $\dim(A/\mathfrak p)$.
\end{proof}

\begin{lemma}
\label{local-cohomology-lemma-cd-blowup}
Soit $A$ un anneau noethérien. Soit $I \subset A$ un idéal.
Soit $b : X' \to X = \Spec(A)$ l'éclatement de $I$.
Si les fibres de $b$ sont de dimension $\leq d - 1$, alors
$\text{cd}(A, I) \leq d$.
\end{lemma}

\begin{proof}
Posons $U = X \setminus V(I)$. Notons $j : U \to X'$ l'immersion ouverte
canonique; voir Diviseurs, section \ref{divisors-section-blowing-up}.
Comme le diviseur exceptionnel est un diviseur de Cartier effectif
(Diviseurs, Lemme
\ref{divisors-lemma-blowing-up-gives-effective-Cartier-divisor})
nous voyons que $j$ est affine; voir
Diviseurs, Lemme
\ref{divisors-lemma-complement-locally-principal-closed-subscheme}.
Soit $\mathcal{F}$ un $\mathcal{O}_U$-module quasi-cohérent.
Alors $R^pj_*\mathcal{F} = 0$ pour $p > 0$; voir
Cohomologie des schémas, Lemme
\ref{coherent-lemma-relative-affine-vanishing}.
Par ailleurs, $R^qb_*(j_*\mathcal{F}) = 0$ pour
$q \geq d$ par Limites, Lemme
\ref{limits-lemma-higher-direct-images-zero-above-dimension-fibre}.
Ainsi, par la suite spectrale de Leray
(Cohomologie, Lemme \ref{cohomology-lemma-relative-Leray}),
nous concluons que $R^n(b \circ j)_*\mathcal{F} = 0$ pour
$n \geq d$. Donc $H^n(U, \mathcal{F}) = 0$ pour $n \geq d$
(par Cohomologie, Lemme \ref{cohomology-lemma-apply-Leray}).
Par définition, cela signifie que $\text{cd}(A, I) \leq d$.
\end{proof}







\section{Supports plus généraux}
\label{local-cohomology-section-supports}

\noindent
Soit $A$ un anneau noethérien. Soit $M$ un $A$-module.
Soit $T \subset \Spec(A)$ une partie stable par spécialisation
(Topologie, Définition \ref{topology-definition-specialization}).
Définissons
$$
H^0_T(M) = \colim_{Z \subset T} H^0_Z(M)
$$
où la limite inductive est prise sur l'ensemble ordonné filtrant des
fermés $Z$ de $\Spec(A)$ contenus dans
$T$\footnote{Comme $T$ est stable par spécialisation,
nous avons $T = \bigcup_{Z \subset T} Z$; voir
Topologie, Lemme \ref{topology-lemma-stable-specialization}.}.
Autrement dit, un élément $m$ de $M$ appartient à $H^0_T(M) \subset M$
si et seulement si le support $V(\text{Ann}_R(m))$ de $m$
est contenu dans $T$.

\begin{lemma}
\label{local-cohomology-lemma-support}
Soit $A$ un anneau noethérien. Soit $T \subset \Spec(A)$ une partie stable
par spécialisation. Pour un $A$-module $M$, les conditions suivantes sont équivalentes :
\begin{enumerate}
\item $H^0_T(M) = M$, et
\item $\text{Supp}(M) \subset T$.
\end{enumerate}
La catégorie de ces $A$-modules est une sous-catégorie de Serre
de la catégorie des $A$-modules, stable par sommes directes.
\end{lemma}

\begin{proof}
L'équivalence vient de ce que le support d'un élément de $M$
est contenu dans le support de $M$ et, réciproquement, le support de
$M$ est la réunion des supports de ses éléments.
La catégorie de ces modules est une sous-catégorie de Serre
(Homologie, Définition \ref{homology-definition-serre-subcategory})
de $\text{Mod}_A$ par
Algèbre, Lemme \ref{algebra-lemma-support-quotient}.
Nous omettons la démonstration de l'assertion sur les sommes directes.
\end{proof}

\noindent
Soit $A$ un anneau noethérien. Soit $T \subset \Spec(A)$ une partie stable
par spécialisation. Notons $\text{Mod}_{A, T} \subset \text{Mod}_A$
la sous-catégorie de Serre décrite dans Lemme \ref{local-cohomology-lemma-support}.
Notons $D_T(A) \subset D(A)$ la
sous-catégorie triangulée strictement pleine et saturée de $D(A)$
(Catégories dérivées, Lemme \ref{derived-lemma-cohomology-in-serre-subcategory})
formée des complexes de $A$-modules dont les modules de cohomologie
appartiennent à $\text{Mod}_{A, T}$. Nous obtenons des foncteurs
$$
D(\text{Mod}_{A, T}) \to D_T(A) \to D(A)
$$
Voir la discussion dans
Catégories dérivées, section \ref{derived-section-triangulated-sub}.
Notons $RH^0_T : D(A) \to D(\text{Mod}_{A, T})$ le foncteur
dérivé droit de $H^0_T$. Nous noterons
$$
R\Gamma_T : D^+(A) \to D^+_T(A),
$$
la composée de $RH^0_T : D^+(A) \to D^+(\text{Mod}_{A, T})$ et de
$D^+(\text{Mod}_{A, T}) \to D^+_T(A)$. Si la dimension de $A$ est
finie\footnote{Si $\dim(A) = \infty$, la construction
peut avoir des propriétés inattendues sur les complexes non bornés.},
nous noterons
$$
R\Gamma_T : D(A) \to D_T(A)
$$
la composée de $RH^0_T$ et de
$D(\text{Mod}_{A, T}) \to D_T(A)$.

\begin{lemma}
\label{local-cohomology-lemma-adjoint}
Soit $A$ un anneau noethérien. Soit $T \subset \Spec(A)$
une partie stable par spécialisation. Le foncteur
$RH^0_T$ est l'adjoint à droite du foncteur
$D(\text{Mod}_{A, T}) \to D(A)$.
\end{lemma}

\begin{proof}
Cela résulte du fait que le foncteur $H^0_T(-)$ est
l'adjoint à droite du foncteur d'inclusion
$\text{Mod}_{A, T} \to \text{Mod}_A$; voir
Catégories dérivées, Lemme \ref{derived-lemma-derived-adjoint-functors}.
\end{proof}

\begin{lemma}
\label{local-cohomology-lemma-adjoint-ext}
Soit $A$ un anneau noethérien. Soit $T \subset \Spec(A)$
une partie stable par spécialisation.
Pour tout objet $K$ de $D(A)$, nous avons
$$
H^i(RH^0_T(K)) = \colim_{Z \subset T\text{ fermé}} H^i_Z(K)
$$
\end{lemma}

\begin{proof}
Soit $J^\bullet$ un complexe K-injectif représentant $K$.
Par définition, $RH^0_T$ est représenté par le complexe
$$
H^0_T(J^\bullet) = \colim H^0_Z(J^\bullet)
$$
où l'égalité résulte de notre définition de $H^0_T$.
Comme les limites inductives filtrantes sont exactes, la cohomologie de ce
complexe en degré $i$ est
$\colim H^i(H^0_Z(J^\bullet)) = \colim H^i_Z(K)$
comme voulu.
\end{proof}

\begin{lemma}
\label{local-cohomology-lemma-equal-plus}
Soit $A$ un anneau noethérien. Soit $T \subset \Spec(A)$ une partie stable
par spécialisation. Le foncteur $D^+(\text{Mod}_{A, T}) \to D^+_T(A)$
est une équivalence.
\end{lemma}

\begin{proof}
Soit $M$ un objet de $\text{Mod}_{A, T}$. Choisissons un plongement
$M \to J$ dans un $A$-module injectif. Par
Complexes dualisants, Proposition
\ref{dualizing-proposition-structure-injectives-noetherian}
le module $J$ est une somme directe d'enveloppes injectives de corps résiduels.
Soit $E$ une enveloppe injective du corps résiduel de $\mathfrak p$.
Comme $E$ est de torsion par rapport aux puissances de $\mathfrak p$, nous voyons que
$H^0_T(E) = 0$ si $\mathfrak p \not \in T$ et
$H^0_T(E) = E$ si $\mathfrak p \in T$.
Ainsi $H^0_T(J)$ est injectif, comme somme directe d'enveloppes injectives
(par la proposition), et nous avons un plongement $M \to H^0_T(J)$.
Ainsi tout objet $M$ de $\text{Mod}_{A, T}$ possède une résolution injective
$M \to J^\bullet$ avec $J^n$ également dans $\text{Mod}_{A, T}$. Il s'ensuit
que $RH^0_T(M) = M$.

\medskip\noindent
Supposons ensuite que $K \in D_T^+(A)$. Alors la suite spectrale
$$
R^qH^0_T(H^p(K)) \Rightarrow R^{p + q}H^0_T(K)
$$
(Catégories dérivées, Lemme \ref{derived-lemma-two-ss-complex-functor})
converge, et nous avons vu ci-dessus que seuls les termes avec $q = 0$
sont non nuls. Ainsi $RH^0_T(K) \to K$ est un isomorphisme.
Le foncteur $D^+(\text{Mod}_{A, T}) \to D^+_T(A)$
est donc une équivalence, de quasi-inverse $RH^0_T$.
\end{proof}

\begin{lemma}
\label{local-cohomology-lemma-equal-full}
Soit $A$ un anneau noethérien. Soit $T \subset \Spec(A)$ une partie stable
par spécialisation. Si $\dim(A) < \infty$, alors le foncteur
$D(\text{Mod}_{A, T}) \to D_T(A)$ est une équivalence.
\end{lemma}

\begin{proof}
Écrivons $\dim(A) = d$. Alors $H^i_Z(M) = 0$ pour $i > d$
pour tout fermé $Z$ de $\Spec(A)$; voir
Lemme \ref{local-cohomology-lemma-cd-dimension}.
Par Lemme \ref{local-cohomology-lemma-adjoint-ext}, nous obtenons que $H^0_T$ est de
dimension cohomologique bornée.

\medskip\noindent
Soit $K \in D_T(A)$. Nous affirmons que $RH^0_T(K) \to K$ est un
isomorphisme. Nous savons que c'est vrai lorsque $K$ est borné inférieurement; voir
Lemme \ref{local-cohomology-lemma-equal-plus}. Cependant, comme $H^0_T$ est de dimension
cohomologique bornée, la $i$-ième cohomologie de
$RH_T^0(K)$ ne dépend que de $\tau_{\geq -d + i}K$, et nous concluons.
Ainsi $D(\text{Mod}_{A, T}) \to D_T(A)$ est une équivalence, de
quasi-inverse $RH^0_T$.
\end{proof}

\begin{remark}
\label{local-cohomology-remark-upshot}
Soit $A$ un anneau noethérien. Soit $T \subset \Spec(A)$ une
partie stable par spécialisation.
La conclusion de la discussion précédente est que
$R\Gamma_T : D^+(A) \to D_T^+(A)$ est l'adjoint à droite
du foncteur d'inclusion $D_T^+(A) \to D^+(A)$.
Si $\dim(A) < \infty$, alors
$R\Gamma_T : D(A) \to D_T(A)$ est l'adjoint à droite
du foncteur d'inclusion $D_T(A) \to D(A)$.
Dans les deux cas, nous avons
$$
H^i_T(K) = H^i(R\Gamma_T(K)) = R^iH^0_T(K) =
\colim_{Z \subset T\text{ fermé}} H^i_Z(K)
$$
Cela résulte de la combinaison des
Lemmes \ref{local-cohomology-lemma-adjoint}, \ref{local-cohomology-lemma-adjoint-ext},
\ref{local-cohomology-lemma-equal-plus} et \ref{local-cohomology-lemma-equal-full}.
\end{remark}

\begin{lemma}
\label{local-cohomology-lemma-torsion-change-rings}
Soit $A \to B$ un homomorphisme plat d'anneaux noethériens.
Soit $T \subset \Spec(A)$ une partie stable par spécialisation.
Soit $T' \subset \Spec(B)$ l'image inverse de $T$.
Alors l'application canonique
$$
R\Gamma_T(K) \otimes_A^\mathbf{L} B
\longrightarrow
R\Gamma_{T'}(K \otimes_A^\mathbf{L} B)
$$
est un isomorphisme pour $K \in D^+(A)$. Si $A$ et $B$ sont de
dimension finie, c'est vrai pour $K \in D(A)$.
\end{lemma}

\begin{proof}
De l'application $R\Gamma_T(K) \to K$, nous tirons une application
$R\Gamma_T(K) \otimes_A^\mathbf{L} B \to K \otimes_A^\mathbf{L} B$.
Les modules de cohomologie de $R\Gamma_T(K) \otimes_A^\mathbf{L} B$
sont supportés sur $T'$, d'où la flèche du lemme.
Cette flèche est un isomorphisme si $T$ est un fermé de $\Spec(A)$, par
Complexes dualisants, Lemme \ref{dualizing-lemma-torsion-change-rings}.
Rappelons que $H^i_T(K)$ est la limite inductive des $H^i_Z(K)$, où $Z$ parcourt
l'ensemble filtrant des fermés de $T$; voir
Lemme \ref{local-cohomology-lemma-adjoint-ext}.
De même,
$H^i_{T'}(K \otimes_A^\mathbf{L} B) =
\colim H^i_{Z'}(K \otimes_A^\mathbf{L} B)$, où $Z'$ est l'image
inverse de $Z$. Le résultat s'ensuit, car $\otimes_A B$ commute
aux limites inductives filtrantes et il n'y a pas de Tor supérieurs.
\end{proof}

\begin{lemma}
\label{local-cohomology-lemma-local-cohomology-ss}
Soit $A$ un anneau et soient $T, T' \subset \Spec(A)$ des parties
stables par spécialisation. Pour $K \in D^+(A)$,
il existe une suite spectrale
$$
E_2^{p, q} = H^p_T(H^p_{T'}(K)) \Rightarrow H^{p + q}_{T \cap T'}(K)
$$
comme dans Catégories dérivées, Lemme
\ref{derived-lemma-grothendieck-spectral-sequence}.
\end{lemma}

\begin{proof}
Soit $E$ un objet de $D_{T \cap T'}(A)$. Alors nous avons
$$
\Hom(E, R\Gamma_T(R\Gamma_{T'}(K))) =
\Hom(E, R\Gamma_{T'}(K)) =
\Hom(E, K)
$$
La première égalité vient de la propriété d'adjonction de $R\Gamma_T$,
et la seconde de celle de $R\Gamma_{T'}$.
Par ailleurs, si $J^\bullet$ est un complexe borné inférieurement
d'injectifs représentant $K$, alors $H^0_{T'}(J^\bullet)$
est un complexe de $A$-modules injectifs représentant $R\Gamma_{T'}(K)$,
et donc $H^0_T(H^0_{T'}(J^\bullet))$ est un complexe représentant
$R\Gamma_T(R\Gamma_{T'}(K))$. Ainsi $R\Gamma_T(R\Gamma_{T'}(K))$
est un objet de $D^+_{T \cap T'}(A)$. En combinant ces deux
faits, nous obtenons $R\Gamma_{T \cap T'} = R\Gamma_T \circ R\Gamma_{T'}$.
Le lemme cité dans l'énoncé fournit alors la suite spectrale.
\end{proof}

\begin{lemma}
\label{local-cohomology-lemma-torsion-tensor-product}
Soit $A$ un anneau noethérien. Soit $T \subset \Spec(A)$ une partie
stable par spécialisation. Supposons $A$ de dimension finie. Alors
$$
R\Gamma_T(K) = R\Gamma_T(A) \otimes_A^\mathbf{L} K
$$
pour $K \in D(A)$. Pour $K, L \in D(A)$, nous avons
$$
R\Gamma_T(K \otimes_A^\mathbf{L} L) =
K \otimes_A^\mathbf{L} R\Gamma_T(L) =
R\Gamma_T(K) \otimes_A^\mathbf{L} L =
R\Gamma_T(K) \otimes_A^\mathbf{L} R\Gamma_T(L)
$$
Si $K$ ou $L$ appartient à $D_T(A)$, alors $K \otimes_A^\mathbf{L} L$ aussi.
\end{lemma}

\begin{proof}
Par construction, nous pouvons représenter $R\Gamma_T(A)$ par un complexe $J^\bullet$ de
$\text{Mod}_{A, T}$. Si nous représentons $K$ par un complexe K-plat $K^\bullet$,
alors $R\Gamma_T(A) \otimes_A^\mathbf{L} K$ est représenté
par le complexe $\text{Tot}(J^\bullet \otimes_A K^\bullet)$ dans
$\text{Mod}_{A, T}$. L'application $R\Gamma_T(A) \to A$ fournit
une application $R\Gamma_T(A) \otimes_A^\mathbf{L} K\to K$. Par la propriété
d'adjonction de $R\Gamma_T$, nous obtenons donc une application canonique
$$
R\Gamma_T(A) \otimes_A^\mathbf{L} K \longrightarrow R\Gamma_T(K)
$$
qui factorise l'application que nous venons de construire. Remarquons que $R\Gamma_T$ commute
aux sommes directes dans $D(A)$, par exemple par Lemme \ref{local-cohomology-lemma-adjoint-ext},
le fait que les limites inductives filtrantes commutent aux sommes directes et le
fait que la cohomologie locale usuelle commute aux sommes directes
(par exemple par Complexes dualisants, Lemme
\ref{dualizing-lemma-local-cohomology-adjoint}).
Ainsi, par Compléments d'algèbre, Remarque \ref{more-algebra-remark-P-resolution},
il suffit de vérifier que l'application est un isomorphisme pour
$K = A[k]$, où $k \in \mathbf{Z}$. C'est clair.

\medskip\noindent
Les dernières assertions résultent de ce que nous venons de démontrer
par transitivité des produits tensoriels dérivés.
\end{proof}





\section{Filtrations de la cohomologie locale}
\label{local-cohomology-section-filter-local-cohomology}

\noindent
Quelques procédés liés à la suite spectrale de
Lemme \ref{local-cohomology-lemma-local-cohomology-ss}.

\begin{lemma}
\label{local-cohomology-lemma-filter-local-cohomology}
Soit $A$ un anneau noethérien. Soit $T \subset \Spec(A)$
une partie stable par spécialisation. Soit $T' \subset T$
l'ensemble des idéaux premiers non minimaux de $T$. Alors $T'$
est une partie de $\Spec(A)$ stable par spécialisation,
et, pour tout $A$-module $M$, il existe une suite exacte
$$
0 \to
\colim_{Z, f} H^1_f(H^{i - 1}_Z(M)) \to
H^i_{T'}(M) \to H^i_T(M) \to
\bigoplus\nolimits_{\mathfrak p \in T \setminus T'}
H^i_{\mathfrak p A_\mathfrak p}(M_\mathfrak p)
$$
où la limite inductive porte sur les fermés $Z \subset T$
et les $f \in A$ tels que $V(f) \cap Z \subset T'$.
\end{lemma}

\begin{proof}
Pour tous $Z$ et $f$, la suite spectrale de
Complexes dualisants, Lemme \ref{dualizing-lemma-local-cohomology-ss},
dégénère et donne des suites exactes courtes
$$
0 \to H^1_f(H^{i - 1}_Z(M)) \to
H^i_{Z \cap V(f)}(M) \to H^0_f(H^i_Z(M)) \to 0
$$
Nous l'utiliserons ci-dessous sans autre mention.

\medskip\noindent
Soit $\xi \in H^i_T(M)$ dont l'image dans la somme directe est nulle.
Écrivons d'abord $\xi$ comme l'image d'un certain $\xi' \in H^i_Z(M)$
pour un fermé $Z \subset T$; voir Lemme \ref{local-cohomology-lemma-adjoint-ext}.
Alors l'image de $\xi'$ est nulle dans $H^i_{\mathfrak p A_\mathfrak p}(M_\mathfrak p)$
pour tout $\mathfrak p \in Z$, $\mathfrak p \not \in T'$.
Comme ces idéaux premiers sont en nombre fini,
nous pouvons choisir $f \in A$ n'appartenant à aucun d'eux
et tel que $f$ annule $\xi'$. Alors $\xi'$
est l'image d'un certain $\xi'' \in H^i_{Z'}(M)$,
où $Z' = Z \cap V(f)$. Par notre choix de $f$, nous avons
$Z' \subset T'$, ce qui donne l'exactitude à l'avant-dernière place.

\medskip\noindent
Soit $\xi \in H^i_{T'}(M)$ dont l'image dans $H^i_T(M)$ est nulle.
Choisissons des fermés $Z' \subset Z$, avec $Z' \subset T'$
et $Z \subset T$, tels que $\xi$ provienne de $\xi' \in H^i_{Z'}(M)$
et que son image dans $H^i_Z(M)$ soit nulle. Nous pouvons alors trouver $f \in A$
tel que $V(f) \cap Z = Z'$, et nous concluons.
\end{proof}

\begin{lemma}
\label{local-cohomology-lemma-zero}
Soit $A$ un anneau noethérien de dimension finie.
Soit $T \subset \Spec(A)$ une partie stable par spécialisation.
Soit $\{M_n\}_{n \geq 0}$ un système projectif de $A$-modules.
Soit $i \geq 0$ un entier. Supposons que, pour tout $m$, il
existe un entier $m'(m) \geq m$ tel que, pour tout
$\mathfrak p \in T$, l'application induite
$$
H^i_{\mathfrak p A_\mathfrak p}(M_{k, \mathfrak p})
\longrightarrow
H^i_{\mathfrak p A_\mathfrak p}(M_{m, \mathfrak p})
$$
soit nulle pour $k \geq m'(m)$. Soit $m'' : \mathbf{N} \to \mathbf{N}$
la composée itérée $2^{\dim(T)}$ fois de $m'$ avec elle-même. Alors l'application
$H^i_T(M_k) \to H^i_T(M_m)$ est nulle pour tout $k \geq m''(m)$.
\end{lemma}

\begin{proof}
Commençons par une remarque générale : supposons donnée une suite
exacte
$$
(A_n) \to (B_n) \to (C_n)
$$
de systèmes projectifs de groupes abéliens. Supposons que, pour tout
$m$, il existe un entier $m'(m) \geq m$ tel que
$$
A_k \to A_m
\quad\text{et}\quad
C_k \to C_m
$$
soient nulles pour $k \geq m'(m)$. Alors, pour $k \geq m'(m'(m))$,
l'application $B_k \to B_m$ est nulle.

\medskip\noindent
Nous démontrons le lemme par récurrence sur $\dim(T)$, qui est
finie puisque $\dim(A)$ l'est. Soit $T' \subset T$
l'ensemble des idéaux premiers non minimaux de $T$. Alors $T'$
est une partie de $\Spec(A)$ stable par spécialisation,
et les hypothèses du lemme s'appliquent à $T'$.
Comme $\dim(T') < \dim(T)$, le lemme vaut pour $T'$.
Pour tout $A$-module $M$, il existe une suite exacte
$$
H^i_{T'}(M) \to H^i_T(M) \to
\bigoplus\nolimits_{\mathfrak p \in T \setminus T'}
H^i_{\mathfrak p A_\mathfrak p}(M_\mathfrak p)
$$
par Lemme \ref{local-cohomology-lemma-filter-local-cohomology}.
Nous concluons donc par la remarque initiale de la démonstration.
\end{proof}

\begin{lemma}
\label{local-cohomology-lemma-essential-image}
Soit $A$ un anneau noethérien. Soit $T \subset \Spec(A)$ une partie
stable par spécialisation. Soit $\{M_n\}_{n \geq 0}$ un système projectif
de $A$-modules. Soit $i \geq 0$ un entier. Supposons la dimension de $A$
finie et que, pour tout $m$, il existe un entier $m'(m) \geq m$
tel que, pour tout $\mathfrak p \in T$, nous ayons :
\begin{enumerate}
\item $H^{i - 1}_{\mathfrak p A_\mathfrak p}(M_{k, \mathfrak p})
\to H^{i - 1}_{\mathfrak p A_\mathfrak p}(M_{m, \mathfrak p})$
est nulle pour $k \geq m'(m)$, et
\item $ H^i_{\mathfrak p A_\mathfrak p}(M_{k, \mathfrak p}) \to
H^i_{\mathfrak p A_\mathfrak p}(M_{m, \mathfrak p})$
est d'image $G(\mathfrak p, m)$ indépendante de $k \geq m'(m)$ et, de plus,
$G(\mathfrak p, m)$ s'injecte dans
$H^i_{\mathfrak p A_\mathfrak p}(M_{0, \mathfrak p})$.
\end{enumerate}
Alors il existe un entier $m_0$ tel que, pour tout $m \geq m_0$,
il existe un entier $m''(m) \geq m$ tel que,
pour $k \geq m''(m)$, l'image de $H^i_T(M_k) \to H^i_T(M_m)$
s'injecte dans $H^i_T(M_{m_0})$.
\end{lemma}

\begin{proof}
Commençons par une remarque générale : supposons donnée une suite
exacte
$$
(A_n) \to (B_n) \to (C_n) \to (D_n)
$$
de systèmes projectifs de groupes abéliens. Supposons qu'il existe
un entier $m_0$ tel que, pour tout $m \geq m_0$,
il existe un entier $m'(m) \geq m$ tel que les applications
$$
\Im(B_k \to B_m) \longrightarrow B_{m_0}
\quad\text{et}\quad
\Im(D_k \to D_m) \longrightarrow D_{m_0}
$$
soient injectives pour $k \geq m'(m)$ et que $A_k \to A_m$ soit nulle
pour $k \geq m'(m)$. Alors, pour $m \geq m'(m_0)$ et $k \geq m'(m'(m))$,
l'application
$$
\Im(C_k \to C_m) \to C_{m'(m_0)}
$$
est injective. En effet, soit $c_0 \in C_m$ l'image de $c_3 \in C_k$,
et supposons que l'image de $c_0$ dans $C_{m'(m_0)}$ soit nulle. Représentons la situation par
$$
C_k \to C_{m'(m'(m))} \to C_{m'(m)} \to C_m \to C_{m'(m_0)},\quad
c_3 \mapsto c_2 \mapsto c_1 \mapsto c_0 \mapsto 0
$$
Il faut montrer que $c_0 = 0$.
L'image $d_3$ de $c_3$ s'envoie sur zéro dans $C_{m_0}$, et donc
l'image $d_1 \in D_{m'(m)}$ est nulle.
Nous pouvons ainsi choisir $b_1 \in B_{m'(m)}$ dont l'image est
$c_1$. Comme $c_3$ s'envoie sur zéro dans
$C_{m'(m_0)}$, nous trouvons un élément $a_{-1} \in A_{m'(m_0)}$
dont l'image est l'image $b_{-1} \in B_{m'(m_0)}$ de $b_1$.
Comme $a_{-1}$ s'envoie sur zéro dans $A_{m_0}$, nous en concluons que
$b_1$ s'envoie sur zéro dans $B_{m_0}$. Ainsi l'image $b_0 \in B_m$
est nulle, ce qui implique bien entendu $c_0 = 0$, comme voulu.

\medskip\noindent
Nous démontrons le lemme par récurrence sur $\dim(T)$, qui est
finie puisque $\dim(A)$ l'est. Soit $T' \subset T$
l'ensemble des idéaux premiers non minimaux de $T$. Alors $T'$ est une partie
de $\Spec(A)$ stable par spécialisation et les hypothèses
du lemme s'appliquent à $T'$. Comme $\dim(T') < \dim(T)$, nous savons
que le lemme vaut pour $T'$. Pour tout $A$-module $M$, il existe une
suite exacte
$$
0 \to \colim_{Z, f} H^1_f(H^{i - 1}_Z(M)) \to
H^i_{T'}(M) \to H^i_T(M) \to
\bigoplus\nolimits_{\mathfrak p \in T \setminus T'}
H^i_{\mathfrak p A_\mathfrak p}(M_\mathfrak p)
$$
par Lemme \ref{local-cohomology-lemma-filter-local-cohomology}.
Nous concluons donc par la remarque initiale de la démonstration
et par le fait que le système de groupes
$$
\left\{\colim_{Z, f} H^1_f(H^{i - 1}_Z(M_n))\right\}_{n \geq 0}
$$
est pro-nul par Lemme \ref{local-cohomology-lemma-zero}; on utilise ici que la fonction
$m''(m)$ de ce lemme, pour $H^{i - 1}_Z(M)$, est indépendante de $Z$.
\end{proof}










\section{Finitude de la cohomologie locale, I}
\label{local-cohomology-section-finiteness}

\noindent
Nous suivrons l'approche de Faltings concernant la finitude des modules de
cohomologie locale; voir \cite{Faltings-annulators} et \cite{Faltings-finiteness}.
Voici un lemme qui montre qu'il suffit d'établir que les modules de
cohomologie locale ont un annulateur pour démontrer qu'ils
sont des modules de type fini.

\begin{lemma}
\label{local-cohomology-lemma-check-finiteness-local-cohomology-by-annihilator}
\begin{reference}
\cite[Lemma 3]{Faltings-annulators}
\end{reference}
Soit $A$ un anneau noethérien. Soit $T \subset \Spec(A)$ une partie stable
par spécialisation. Soit $M$ un $A$-module de type fini. Soit $n \geq 0$.
Les conditions suivantes sont équivalentes :
\begin{enumerate}
\item $H^i_T(M)$ est de type fini pour $i \leq n$,
\item il existe un idéal $J \subset A$ tel que $V(J) \subset T$
et que $J$ annule $H^i_T(M)$ pour $i \leq n$.
\end{enumerate}
Si $T = V(I) = Z$ pour un idéal $I \subset A$, ces conditions sont aussi
équivalentes à :
\begin{enumerate}
\item[(3)] il existe $e \geq 0$ tel que $I^e$ annule
$H^i_Z(M)$ pour $i \leq n$.
\end{enumerate}
\end{lemma}

\begin{proof}
Nous démontrons l'équivalence de (1) et (2) par récurrence sur $n$.
Pour $n = 0$, le module $H^0_T(M) \subset M$ est de type fini. Ainsi (1) est vraie.
Comme $H^0_T(M) = \colim H^0_{V(J)}(M)$, avec $J$ comme en (2), nous voyons
que (2) est vraie. Supposons désormais $n > 0$.

\medskip\noindent
Supposons (1) vraie. Rappelons que $H^i_J(M) = H^i_{V(J)}(M)$; voir
Complexes dualisants, Lemme \ref{dualizing-lemma-local-cohomology-noetherian}.
Ainsi $H^i_T(M) = \colim H^i_J(M)$, où la limite inductive porte sur les idéaux
$J \subset A$ tels que $V(J) \subset T$; voir
Lemme \ref{local-cohomology-lemma-adjoint-ext}. Comme $H^i_T(M)$ est de type fini
pour $i \leq n$, nous pouvons trouver $J \subset A$ comme en (2) tel que
$H^i_J(M) \to H^i_T(M)$ soit surjectif pour $i \leq n$.
La liste finie de générateurs est donc formée d'éléments annulés par des puissances de $J$,
et (2) est satisfaite après remplacement de $J$ par une puissance.

\medskip\noindent
Supposons donné $J$ comme en (2). Posons $N = H^0_T(M)$ et $M' = M/N$.
Par construction de $R\Gamma_T$, nous obtenons
$H^i_T(N) = 0$ pour $i > 0$ et $H^0_T(N) = N$; voir
Remarque \ref{local-cohomology-remark-upshot}. Il en résulte que
$H^0_T(M') = 0$ et $H^i_T(M') = H^i_T(M)$ pour $i > 0$.
Nous pouvons donc remplacer $M$ par $M'$.
Nous pouvons ainsi supposer que $H^0_T(M) = 0$.
Cela signifie que l'ensemble fini des idéaux premiers associés à $M$
ne rencontre pas $T$. Par évitement des idéaux premiers (Algèbre, Lemme \ref{algebra-lemma-silly}),
nous pouvons trouver $f \in J$ qui n'appartient à aucun des idéaux premiers associés à $M$.
La suite exacte longue de cohomologie locale associée à la suite
exacte courte
$$
0 \to M \to M \to M/fM \to 0
$$
se transforme alors en suites exactes courtes
$$
0 \to H^i_T(M) \to H^i_T(M/fM) \to H^{i + 1}_T(M) \to 0
$$
pour $i < n$. Nous en concluons que $J^2$ annule $H^i_T(M/fM)$
pour $i < n$. Par l'hypothèse de récurrence, $H^i_T(M/fM)$
est de type fini pour $i < n$. En utilisant encore une fois la suite exacte courte,
nous voyons que $H^{i + 1}_T(M)$ est de type fini pour $i < n$, comme voulu.

\medskip\noindent
Nous omettons la démonstration de l'équivalence de (2) et (3)
dans le cas $T = V(I)$.
\end{proof}

\noindent
Le résultat suivant de Faltings permet d'établir la finitude
de la cohomologie locale au niveau des anneaux locaux.

\begin{lemma}
\label{local-cohomology-lemma-check-finiteness-local-cohomology-locally}
\begin{reference}
C'est un cas particulier de \cite[Satz 1]{Faltings-finiteness}.
\end{reference}
Soient $A$ un anneau noethérien, $I \subset A$ un idéal, $M$ un
$A$-module de type fini et $n \geq 0$ un entier. Posons $Z = V(I)$.
Les conditions suivantes sont équivalentes :
\begin{enumerate}
\item les modules $H^i_Z(M)$ sont de type fini pour $i \leq n$, et
\item pour tout $\mathfrak p \in \Spec(A)$, les modules
$H^i_Z(M)_\mathfrak p$, $i \leq n$, sont des $A_\mathfrak p$-modules de type fini.
\end{enumerate}
\end{lemma}

\begin{proof}
L'implication (1) $\Rightarrow$ (2) est immédiate. Démontrons la réciproque
par récurrence sur $n$. Le cas $n = 0$ est clair, car (1) et
(2) sont toujours vraies dans ce cas.

\medskip\noindent
Supposons $n > 0$ et (2) vraie. Posons $N = H^0_Z(M)$ et $M' = M/N$.
Par Complexes dualisants, Lemme \ref{dualizing-lemma-divide-by-torsion},
nous pouvons remplacer $M$ par $M'$.
Nous pouvons donc supposer que $H^0_Z(M) = 0$.
Cela signifie que $\text{depth}_I(M) > 0$
(Complexes dualisants, Lemme \ref{dualizing-lemma-depth}).
Choisissons $f \in I$ qui soit un non-diviseur de zéro dans $M$ et considérons la suite
exacte courte
$$
0 \to M \to M \to M/fM \to 0
$$
qui produit une suite exacte longue
$$
0 \to H^0_Z(M/fM) \to H^1_Z(M) \to H^1_Z(M) \to H^1_Z(M/fM) \to
H^2_Z(M) \to \ldots
$$
et de même après localisation. Ainsi l'hypothèse (2) implique que
les modules $H^i_Z(M/fM)_\mathfrak p$ sont de type fini pour $i < n$. Par conséquent,
par l'hypothèse de récurrence, les $H^i_Z(M/fM)$ sont de type fini pour $i < n$.

\medskip\noindent
Soit $\mathfrak p$ un idéal premier de $A$ associé à
$H^i_Z(M)$ pour un certain $i \leq n$. Supposons que $\mathfrak p$ soit l'annulateur
de l'élément $x \in H^i_Z(M)$. Alors $\mathfrak p \in Z$, donc
$f \in \mathfrak p$. Ainsi $fx = 0$, et $x$ provient donc d'un
élément de $H^{i - 1}_Z(M/fM)$ par l'application de bord $\delta$ de la suite
exacte longue précédente. Il s'ensuit que $\mathfrak p$ est un idéal premier associé
au module de type fini $\Im(\delta)$. Nous en concluons que
$\text{Ass}(H^i_Z(M))$ est fini pour $i \leq n$; voir
Algèbre, Lemme \ref{algebra-lemma-finite-ass}.

\medskip\noindent
Rappelons que
$$
H^i_Z(M) \subset
\prod\nolimits_{\mathfrak p \in \text{Ass}(H^i_Z(M))}
H^i_Z(M)_\mathfrak p
$$
par Algèbre, Lemme \ref{algebra-lemma-zero-at-ass-zero}. Comme, par
hypothèse, les modules du membre de droite sont de type fini et de torsion par rapport aux puissances de $I$,
nous pouvons trouver des entiers $e_{\mathfrak p, i} \geq 0$, $i \leq n$,
$\mathfrak p \in \text{Ass}(H^i_Z(M))$, tels que
$I^{e_{\mathfrak p, i}}$ annule $H^i_Z(M)_\mathfrak p$. Nous en concluons
que $I^e$, avec $e = \max\{e_{\mathfrak p, i}\}$, annule $H^i_Z(M)$
pour $i \leq n$. Par
Lemme \ref{local-cohomology-lemma-check-finiteness-local-cohomology-by-annihilator},
nous voyons que $H^i_Z(M)$ est de type fini pour $i \leq n$.
\end{proof}

\begin{lemma}
\label{local-cohomology-lemma-annihilate-local-cohomology}
Soit $A$ un anneau et soient $J \subset I \subset A$ des idéaux de type fini.
Soit $i \geq 0$ un entier. Posons $Z = V(I)$. Si
$H^i_Z(A)$ est annulé par $J^n$ pour un certain $n$, alors
$H^i_Z(M)$ est annulé par $J^m$ pour un certain $m = m(M)$,
pour tout $A$-module $M$ de présentation finie tel que
$M_f$ soit un $A_f$-module libre de rang fini pour tout $f \in I$.
\end{lemma}

\begin{proof}
Considérons l'annulateur $\mathfrak a$ de $H^i_Z(M)$.
Soit $\mathfrak p \subset A$ tel que $\mathfrak p \not \in Z$.
Par hypothèse, il existe $f \in I$, $f \not \in \mathfrak p$,
et un isomorphisme $\varphi : A_f^{\oplus r} \to M_f$
de $A_f$-modules. En chassant les dénominateurs (et en utilisant que
$M$ est de présentation finie), nous trouvons des applications
$$
a : A^{\oplus r} \longrightarrow M
\quad\text{et}\quad
b : M \longrightarrow A^{\oplus r}
$$
telles que $a_f = f^N \varphi$ et $b_f = f^N \varphi^{-1}$ pour un certain $N$.
De plus, nous pouvons supposer que $a \circ b$ et $b \circ a$ sont égales à
la multiplication par $f^{2N}$. Ainsi $H^i_Z(M)$ est annulé par
$f^{2N}J^n$, c'est-à-dire $f^{2N}J^n \subset \mathfrak a$.

\medskip\noindent
Comme $U = \Spec(A) \setminus Z$ est quasi-compact, nous pouvons trouver un nombre fini de
$f_1, \ldots, f_t$ et $N_1, \ldots, N_t$ tels que $U = \bigcup D(f_j)$ et
$f_j^{2N_j}J^n \subset \mathfrak a$. Alors $V(I) = V(f_1, \ldots, f_t)$
et, comme $I$ est de type fini, nous en concluons que
$I^M \subset (f_1, \ldots, f_t)$ pour un certain $M$.
En définitive, $J^m \subset \mathfrak a$ pour
$m \gg 0$; par exemple, $m = M (2N_1 + \ldots + 2N_t) n$ convient.
\end{proof}

\begin{lemma}
\label{local-cohomology-lemma-local-finiteness-for-finite-locally-free}
Soit $A$ un anneau noethérien. Soit $I \subset A$ un idéal. Posons $Z = V(I)$.
Soit $n \geq 0$ un entier. Si $H^i_Z(A)$ est de type fini pour $0 \leq i \leq n$,
il en est de même de $H^i_Z(M)$, $0 \leq i \leq n$, pour
tout $A$-module de type fini $M$ tel que $M_f$ soit un
$A_f$-module libre de rang fini pour tout $f \in I$.
\end{lemma}

\begin{proof}
L'hypothèse que $H^i_Z(A)$ est de type fini pour $0 \leq i \leq n$
implique qu'il existe $e \geq 0$ tel que $I^e$ annule
$H^i_Z(A)$ pour $0 \leq i \leq n$; voir
Lemme \ref{local-cohomology-lemma-check-finiteness-local-cohomology-by-annihilator}.
Puis Lemme \ref{local-cohomology-lemma-annihilate-local-cohomology}
implique que $H^i_Z(M)$, $0 \leq i \leq n$, est annulé
par $I^m$ pour un certain $m = m(M, i)$. Nous pouvons prendre le même $m$
pour tout $0 \leq i \leq n$. Alors
Lemme \ref{local-cohomology-lemma-check-finiteness-local-cohomology-by-annihilator}
implique que $H^i_Z(M)$ est de type fini pour $0 \leq i \leq n$,
comme voulu.
\end{proof}





\section{Finitude des images directes, I}
\label{local-cohomology-section-finiteness-pushforward}

\noindent
Dans cette section, nous étudions le cas non trivial le plus simple du théorème de finitude, à savoir la finitude de la cohomologie locale de degré un, ou, de manière équivalente, celle de $j_*\mathcal{F}$, où $j : U \to X$ est une immersion ouverte, $X$ est localement noethérien et $\mathcal{F}$ est un faisceau cohérent sur $U$. En suivant une méthode de Koll\'ar (\cite{Kollar-variants} et \cite{Kollar-local-global-hulls}), nous obtenons une condition nécessaire et suffisante ; voir la proposition \ref{local-cohomology-proposition-kollar}. Le lecteur qui s'intéresse aux images directes supérieures ou aux groupes de cohomologie locale de degré supérieur peut passer directement à la section \ref{local-cohomology-section-finiteness-pushforward-II} ou à la section \ref{local-cohomology-section-finiteness-II} (développées indépendamment du reste de la présente section).

\begin{lemma}
\label{local-cohomology-lemma-check-finiteness-pushforward-on-associated-points}
Soit $X$ un schéma localement noethérien. Soit $j : U \to X$ l'inclusion d'un sous-schéma ouvert de complémentaire $Z$. Pour $x \in U$, soit $i_x : W_x \to U$ le sous-schéma fermé intègre de point générique $x$. Soit $\mathcal{F}$ un $\mathcal{O}_U$-module cohérent. Les conditions suivantes sont équivalentes :
\begin{enumerate}
\item pour tout $x \in \text{Ass}(\mathcal{F})$, le $\mathcal{O}_X$-module $j_*i_{x, *}\mathcal{O}_{W_x}$ est cohérent,
\item $j_*\mathcal{F}$ est cohérent.
\end{enumerate}
\end{lemma}

\begin{proof}
Montrons d'abord que (1) implique (2). Supposons la condition (1) vérifiée. L'énoncé est local sur $X$ ; on peut donc supposer $X$ affine. Alors $U$ est quasi-compact, donc $\text{Ass}(\mathcal{F})$ est fini (Diviseurs, lemme \ref{divisors-lemma-finite-ass}). On peut ainsi raisonner par récurrence sur le nombre de points associés. Soit $x \in U$ un point générique d'une composante irréductible du support de $\mathcal{F}$. D'après Diviseurs, lemme \ref{divisors-lemma-finite-ass}, on a $x \in \text{Ass}(\mathcal{F})$. Le choix de $x$ assure que $\dim(\mathcal{F}_x) = 0$ en tant que $\mathcal{O}_{X, x}$-module. Ainsi $\mathcal{F}_x$ est de longueur finie comme $\mathcal{O}_{X, x}$-module (Algèbre, lemme \ref{algebra-lemma-support-point}). On peut donc raisonner par récurrence sur cette longueur.

\medskip\noindent
Posons $\mathcal{G} = j_*i_{x, *}\mathcal{O}_{W_x}$. C'est un $\mathcal{O}_X$-module cohérent par hypothèse. On a $\mathcal{G}_x = \kappa(x)$. Choisissons une application non nulle $\varphi_x : \mathcal{F}_x \to \kappa(x) = \mathcal{G}_x$. D'après Cohomologie des schémas, lemme \ref{coherent-lemma-map-stalks-local-map}, il existe un ouvert $x \in V \subset U$ et une application $\varphi_V : \mathcal{F}|_V \to \mathcal{G}|_V$ dont le germe en $x$ est $\varphi_x$. Choisissons $f \in \Gamma(X, \mathcal{O}_X)$ ne s'annulant pas en $x$ et tel que $D(f) \subset V$. D'après Cohomologie des schémas, lemme \ref{coherent-lemma-homs-over-open} (par exemple), on voit que $\varphi_V$ se prolonge en $f^n\mathcal{F} \to \mathcal{G}|_U$ pour un certain $n$. En précomposant avec la multiplication par $f^n$, on obtient une application $\mathcal{F} \to \mathcal{G}|_U$ dont le germe en $x$ est non nul. Soit $\mathcal{F}' \subset \mathcal{F}$ son noyau. Notons que $\text{Ass}(\mathcal{F}') \subset \text{Ass}(\mathcal{F})$ ; voir Diviseurs, lemme \ref{divisors-lemma-ses-ass}. Puisque $\text{length}_{\mathcal{O}_{X, x}}(\mathcal{F}'_x) =
\text{length}_{\mathcal{O}_{X, x}}(\mathcal{F}_x) - 1$, l'hypothèse de récurrence permet de conclure que $j_*\mathcal{F}'$ est cohérent. Puisque $\mathcal{G} = j_*(\mathcal{G}|_U) = j_*i_{x, *}\mathcal{O}_{W_x}$ est cohérent, on peut considérer la suite exacte
$$
0 \to j_*\mathcal{F}' \to j_*\mathcal{F} \to \mathcal{G}
$$
D'après Schémas, lemme \ref{schemes-lemma-push-forward-quasi-coherent}, le faisceau $j_*\mathcal{F}$ est quasi-cohérent. L'image de $j_*\mathcal{F}$ dans $j_*(\mathcal{G}|_U)$ est donc cohérente, d'après Cohomologie des schémas, lemme \ref{coherent-lemma-coherent-Noetherian-quasi-coherent-sub-quotient}. Enfin, $j_*\mathcal{F}$ est cohérent d'après Cohomologie des schémas, lemme \ref{coherent-lemma-coherent-abelian-Noetherian}.

\medskip\noindent
Supposons la condition (2) vérifiée. Exactement comme ci-dessus, on se ramène au cas où $X$ est affine. Choisissons $x \in \text{Ass}(\mathcal{F})$ et posons $\mathcal{G} = j_*i_{x, *}\mathcal{O}_{W_x}$. Choisissons ensuite une application non nulle $\varphi_x : \mathcal{G}_x = \kappa(x) \to \mathcal{F}_x$, qui existe précisément parce que $x$ est un point associé de $\mathcal{F}$. En raisonnant exactement comme ci-dessus, on peut supposer que $\varphi_x$ se prolonge en un homomorphisme de $\mathcal{O}_U$-modules $\varphi : \mathcal{G}|_U \to \mathcal{F}$. Alors $\varphi$ est injective (par exemple d'après Diviseurs, lemme \ref{divisors-lemma-check-injective-on-ass}), et l'on obtient une application injective $\mathcal{G} = j_*(\mathcal{G}|_V) \to j_*\mathcal{F}$. Ainsi la condition (1) est vérifiée.
\end{proof}

\begin{lemma}
\label{local-cohomology-lemma-finiteness-pushforwards-and-H1-local}
Soient $A$ un anneau noethérien et $I \subset A$ un idéal. Posons $X = \Spec(A)$, $Z = V(I)$, $U = X \setminus Z$, et notons $j : U \to X$ le morphisme d'inclusion. Soit $\mathcal{F}$ un $\mathcal{O}_U$-module cohérent. Alors
\begin{enumerate}
\item il existe un $A$-module de type fini $M$ tel que $\mathcal{F}$ soit la restriction de $\widetilde{M}$ à $U$,
\item pour un tel $M$, on a une suite exacte
$$
0 \to H^0_Z(M) \to M \to H^0(U, \mathcal{F}) \to H^1_Z(M) \to 0
$$
et des isomorphismes $H^p(U, \mathcal{F}) = H^{p + 1}_Z(M)$ pour $p \geq 1$,
\item pour $M$ et $p \geq 0$ donnés, les conditions suivantes sont équivalentes :
\begin{enumerate}
\item $R^pj_*\mathcal{F}$ est cohérent,
\item $H^p(U, \mathcal{F})$ est un $A$-module de type fini,
\item $H^{p + 1}_Z(M)$ est un $A$-module de type fini,
\end{enumerate}
\item si les conditions équivalentes de (3) sont vérifiées pour $p = 0$, on peut prendre $M = \Gamma(U, \mathcal{F})$, auquel cas on a $H^0_Z(M) = H^1_Z(M) = 0$.
\end{enumerate}
\end{lemma}

\begin{proof}
D'après Propriétés, lemme \ref{properties-lemma-lift-finite-presentation}, il existe un $\mathcal{O}_X$-module cohérent $\mathcal{F}'$ dont la restriction à $U$ est isomorphe à $\mathcal{F}$. Supposons que $\mathcal{F}'$ corresponde au $A$-module de type fini $M$ comme en (1). Notons que $R^pj_*\mathcal{F}$ est quasi-cohérent (Cohomologie des schémas, lemme \ref{coherent-lemma-quasi-coherence-higher-direct-images}) et correspond au $A$-module $H^p(U, \mathcal{F})$. D'après le lemme \ref{local-cohomology-lemma-local-cohomology-is-local-cohomology} et la discussion de Cohomologie, sections \ref{cohomology-section-cohomology-support} et \ref{cohomology-section-cohomology-support-bis}, on obtient une suite exacte
$$
0 \to H^0_Z(M) \to M \to H^0(U, \mathcal{F}) \to H^1_Z(M) \to 0
$$
et des isomorphismes $H^p(U, \mathcal{F}) = H^{p + 1}_Z(M)$ pour $p \geq 1$. On utilise ici le fait que $H^j(X, \mathcal{F}') = 0$ pour $j > 0$, puisque $X$ est affine et $\mathcal{F}'$ est quasi-cohérent (Cohomologie des schémas, lemme \ref{coherent-lemma-quasi-coherent-affine-cohomology-zero}). Cela prouve (2). Les points (3) et (4) résultent directement de (2) ; voir aussi le lemme \ref{local-cohomology-lemma-local-cohomology}.
\end{proof}

\begin{lemma}
\label{local-cohomology-lemma-finiteness-pushforward}
Soit $X$ un schéma localement noethérien. Soit $j : U \to X$ l'inclusion d'un sous-schéma ouvert de complémentaire $Z$. Soit $\mathcal{F}$ un $\mathcal{O}_U$-module cohérent. Supposons que
\begin{enumerate}
\item $X$ soit de Nagata,
\item $X$ soit universellement caténaire,
\item pour $x \in \text{Ass}(\mathcal{F})$ et $z \in Z \cap \overline{\{x\}}$, on ait $\dim(\mathcal{O}_{\overline{\{x\}}, z}) \geq 2$.
\end{enumerate}
Alors $j_*\mathcal{F}$ est cohérent.
\end{lemma}

\begin{proof}
D'après le lemme \ref{local-cohomology-lemma-check-finiteness-pushforward-on-associated-points}, il suffit de montrer que $j_*i_{x, *}\mathcal{O}_{W_x}$ est cohérent pour $x \in \text{Ass}(\mathcal{F})$. Soit $\pi : Y \to X$ la normalisation de $X$ dans $\Spec(\kappa(x))$ ; voir Morphismes, section \ref{morphisms-section-normalization}. D'après Morphismes, lemme \ref{morphisms-lemma-nagata-normalization-finite-general}, le morphisme $\pi$ est fini. Puisque $\pi$ est fini, $\mathcal{G} = \pi_*\mathcal{O}_Y$ est un $\mathcal{O}_X$-module cohérent d'après Cohomologie des schémas, lemme \ref{coherent-lemma-finite-pushforward-coherent}. Observons que $W_x = U \cap \pi(Y)$. Ainsi $\pi|_{\pi^{-1}(U)} : \pi^{-1}(U) \to U$ se factorise par $i_x : W_x \to U$, et l'on obtient une application canonique
$$
i_{x, *}\mathcal{O}_{W_x}
\longrightarrow
(\pi|_{\pi^{-1}(U)})_*(\mathcal{O}_{\pi^{-1}(U)}) =
(\pi_*\mathcal{O}_Y)|_U = \mathcal{G}|_U
$$
Cette application est injective (par exemple d'après Diviseurs, lemme \ref{divisors-lemma-check-injective-on-ass}). On a donc $j_*i_{x, *}\mathcal{O}_{W_x} \subset j_*\mathcal{G}|_U$, et il suffit de montrer que $j_*\mathcal{G}|_U$ est cohérent.

\medskip\noindent
Il reste à montrer que $j_*(\mathcal{G}|_U)$ est cohérent. Montrons que Diviseurs, lemme \ref{divisors-lemma-depth-2-hartog}, s'applique à
$$
\mathcal{G} \longrightarrow j_*(\mathcal{G}|_U)
$$
ce qui achève la démonstration. Il suffit de montrer que $\text{depth}(\mathcal{G}_z) \geq 2$ pour $z \in Z$. Soient $y_1, \ldots, y_n \in Y$ les points s'envoyant sur $z$. D'après Algèbre, lemme \ref{algebra-lemma-depth-goes-down-finite}, il suffit de montrer que $\text{depth}(\mathcal{O}_{Y, y_i}) \geq 2$ pour $i = 1, \ldots, n$. Sinon, Propriétés, lemme \ref{properties-lemma-criterion-normal}, montre que $\dim(\mathcal{O}_{Y, y_i}) = 1$ pour un certain $i$. C'est impossible d'après la formule des dimensions (Morphismes, lemme \ref{morphisms-lemma-dimension-formula}) pour $\pi : Y \to \overline{\{x\}}$ et l'hypothèse (3).
\end{proof}

\begin{lemma}
\label{local-cohomology-lemma-sharp-finiteness-pushforward}
Soit $X$ un schéma intègre localement noethérien. Soit $j : U \to X$ l'inclusion d'un sous-schéma ouvert non vide de complémentaire $Z$. Supposons que, pour tout $z \in Z$ et tout idéal premier associé $\mathfrak p$ du complété $\mathcal{O}_{X, z}^\wedge$, on ait $\dim(\mathcal{O}_{X, z}^\wedge/\mathfrak p) \geq 2$. Alors $j_*\mathcal{O}_U$ est cohérent.
\end{lemma}

\begin{proof}
On peut supposer $X$ affine. Les lemmes \ref{local-cohomology-lemma-check-finiteness-local-cohomology-locally} et \ref{local-cohomology-lemma-finiteness-pushforwards-and-H1-local} permettent de se ramener à $X = \Spec(A)$, où $(A, \mathfrak m)$ est un anneau local noethérien intègre et $\mathfrak m \in Z$. On peut alors raisonner par récurrence sur $d = \dim(A)$. (Les cas initiaux $d = 0, 1$ ne se présentent pas en vertu de l'hypothèse sur les anneaux locaux.) Posons $V = \Spec(A) \setminus \{\mathfrak m\}$. Observons que les anneaux locaux de $V$ ont une dimension strictement inférieure à $d$. En reprenant les arguments pour $j' : U \to V$ et en utilisant la récurrence, on conclut que $j'_*\mathcal{O}_U$ est un $\mathcal{O}_V$-module cohérent. Choisissons un élément non nul $f \in A$ qui s'annule sur $Z$. Puisque $D(f) \cap V \subset U$, on trouve $n$ tel que la multiplication par $f^n$ sur $U$ se prolonge en une application $f^n : j'_*\mathcal{O}_U \to \mathcal{O}_V$ sur $V$ (par exemple d'après Cohomologie des schémas, lemme \ref{coherent-lemma-homs-over-open}). Cette application est injective ; on a donc une application injective
$$
j_*\mathcal{O}_U = j''_* j'_* \mathcal{O}_U \to j''_*\mathcal{O}_V
$$
sur $X$, où $j'' : V \to X$ est le morphisme d'inclusion. Il suffit donc de montrer que $j''_*\mathcal{O}_V$ est cohérent. Autrement dit, on peut supposer que $X$ est le spectre d'un anneau local noethérien intègre et que $Z$ est constitué du point fermé.

\medskip\noindent
Supposons $X = \Spec(A)$, avec $(A, \mathfrak m)$ local et $Z = \{\mathfrak m\}$. Soit $A^\wedge$ le complété de $A$. Posons $X^\wedge = \Spec(A^\wedge)$, $Z^\wedge = \{\mathfrak m^\wedge\}$, $U^\wedge = X^\wedge \setminus Z^\wedge$ et $\mathcal{F}^\wedge = \mathcal{O}_{U^\wedge}$. L'anneau $A^\wedge$ est universellement caténaire et de Nagata (Algèbre, remarque \ref{algebra-remark-Noetherian-complete-local-ring-universally-catenary} et lemme \ref{algebra-lemma-Noetherian-complete-local-Nagata}). De plus, la condition (3) du lemme \ref{local-cohomology-lemma-finiteness-pushforward} pour $X^\wedge, Z^\wedge, U^\wedge, \mathcal{F}^\wedge$ est vérifiée par hypothèse ! On voit donc que $(U^\wedge \to X^\wedge)_*\mathcal{O}_{U^\wedge}$ est cohérent. Puisque le morphisme $c : X^\wedge \to X$ est plat, l'image réciproque de $j_*\mathcal{O}_U$ est $(U^\wedge \to X^\wedge)_*\mathcal{O}_{U^\wedge}$ (Cohomologie des schémas, lemme \ref{coherent-lemma-flat-base-change-cohomology}). Enfin, puisque $c$ est fidèlement plat, on conclut que $j_*\mathcal{O}_U$ est cohérent d'après Descente, lemme \ref{descent-lemma-finite-type-descends}.
\end{proof}

\begin{remark}
\label{local-cohomology-remark-closure}
Soit $j : U \to X$ une immersion ouverte de schémas localement noethériens. Soit $x \in U$. Soit $i_x : W_x \to U$ le sous-schéma fermé intègre de point générique $x$, et soit $\overline{\{x\}}$ l'adhérence dans $X$. On a alors un diagramme commutatif
$$
\xymatrix{
W_x \ar[d]_{i_x} \ar[r]_{j'} & \overline{\{x\}} \ar[d]^i \\
U \ar[r]^j & X
}
$$
On a $j_*i_{x, *}\mathcal{O}_{W_x} = i_*j'_*\mathcal{O}_{W_x}$. Puisque la flèche verticale de gauche est une immersion fermée, $j_*i_{x, *}\mathcal{O}_{W_x}$ est cohérent si et seulement si $j'_*\mathcal{O}_{W_x}$ est cohérent.
\end{remark}

\begin{remark}
\label{local-cohomology-remark-no-finiteness-pushforward}
Soit $X$ un schéma localement noethérien. Soit $j : U \to X$ l'inclusion d'un sous-schéma ouvert de complémentaire $Z$. Soit $\mathcal{F}$ un $\mathcal{O}_U$-module cohérent. S'il existe $x \in \text{Ass}(\mathcal{F})$ et $z \in Z \cap \overline{\{x\}}$ tels que $\dim(\mathcal{O}_{\overline{\{x\}}, z}) \leq 1$, alors $j_*\mathcal{F}$ n'est pas cohérent. Pour le montrer, on peut effectuer un changement de base plat vers le spectre de $\mathcal{O}_{X, z}$. Posons $X' = \overline{\{x\}}$. L'hypothèse implique $\mathcal{O}_{X' \cap U} \subset \mathcal{F}$. Il suffit donc de voir que $j_*\mathcal{O}_{X' \cap U}$ n'est pas cohérent. Cela est clair, car $X' = \{x, z\}$ ; ainsi $j_*\mathcal{O}_{X' \cap U}$ correspond à $\kappa(x)$ en tant que $\mathcal{O}_{X, z}$-module, qui ne peut être de type fini puisque $x$ n'est pas un point fermé.

\medskip\noindent
En fait, la réciproque du lemme \ref{local-cohomology-lemma-sharp-finiteness-pushforward} est vraie : étant donnée une immersion ouverte $j : U \to X$ de schémas noethériens intègres, s'il existe $z \in X \setminus U$ et un idéal premier associé $\mathfrak p$ du complété $\mathcal{O}_{X, z}^\wedge$ tel que $\dim(\mathcal{O}_{X, z}^\wedge/\mathfrak p) = 1$, alors $j_*\mathcal{O}_U$ n'est pas cohérent. En effet, on peut passer à l'anneau local, agrandir $U$ au spectre épointé, passer au complété, puis appliquer l'argument ci-dessus pour obtenir la non-finitude.
\end{remark}

\begin{proposition}[Koll\'ar]
\label{local-cohomology-proposition-kollar}
\begin{reference}
Voir \cite{k-coherent}, et voir \cite[IV, Proposition 7.2.2]{EGA} pour un cas particulier.
\end{reference}
\begin{slogan}
Analogue faible du théorème de Hartogs : sur un schéma noethérien, la restriction d'un faisceau cohérent à un ouvert dont le complémentaire est de codimension 2 dans le support du faisceau est cohérente.
\end{slogan}
Soit $j : U \to X$ une immersion ouverte de schémas localement noethériens de complémentaire $Z$. Soit $\mathcal{F}$ un $\mathcal{O}_U$-module cohérent. Les conditions suivantes sont équivalentes :
\begin{enumerate}
\item $j_*\mathcal{F}$ est cohérent,
\item pour $x \in \text{Ass}(\mathcal{F})$, $z \in Z \cap \overline{\{x\}}$ et tout idéal premier associé $\mathfrak p$ du complété $\mathcal{O}_{\overline{\{x\}}, z}^\wedge$, on a $\dim(\mathcal{O}_{\overline{\{x\}}, z}^\wedge/\mathfrak p) \geq 2$.
\end{enumerate}
\end{proposition}

\begin{proof}
Si la condition (2) est vérifiée, on obtient (1) en combinant le lemme \ref{local-cohomology-lemma-check-finiteness-pushforward-on-associated-points}, la remarque \ref{local-cohomology-remark-closure} et le lemme \ref{local-cohomology-lemma-sharp-finiteness-pushforward}. Si la condition (2) n'est pas vérifiée, alors $j_*i_{x, *}\mathcal{O}_{W_x}$ n'est pas de type fini pour un certain $x \in \text{Ass}(\mathcal{F})$, d'après la discussion de la remarque \ref{local-cohomology-remark-no-finiteness-pushforward} (et la remarque \ref{local-cohomology-remark-closure}). Ainsi $j_*\mathcal{F}$ n'est pas cohérent d'après le lemme \ref{local-cohomology-lemma-check-finiteness-pushforward-on-associated-points}.
\end{proof}

\begin{lemma}
\label{local-cohomology-lemma-kollar-finiteness-H1-local}
Soient $A$ un anneau noethérien et $I \subset A$ un idéal. Posons $Z = V(I)$. Soit $M$ un $A$-module de type fini. Les conditions suivantes sont équivalentes :
\begin{enumerate}
\item $H^1_Z(M)$ est un $A$-module de type fini,
\item pour tous $\mathfrak p \in \text{Ass}(M)$, $\mathfrak p \not \in Z$ et tout $\mathfrak q \in V(\mathfrak p + I)$, le complété de $(A/\mathfrak p)_\mathfrak q$ n'a pas d'idéal premier associé de dimension $1$.
\end{enumerate}
\end{lemma}

\begin{proof}
Cela résulte immédiatement de la proposition \ref{local-cohomology-proposition-kollar}, par l'intermédiaire du lemme \ref{local-cohomology-lemma-finiteness-pushforwards-and-H1-local}.
\end{proof}

\noindent
La formulation du lemme suivant a l'avantage que les conditions (1) et (2) se transmettent aux schémas de type fini sur $X$. De plus, c'est cette forme de finitude que nous généraliserons aux images directes supérieures dans la section \ref{local-cohomology-section-finiteness-pushforward-II}.

\begin{lemma}
\label{local-cohomology-lemma-finiteness-pushforward-general}
Soit $X$ un schéma localement noethérien. Soit $j : U \to X$ l'inclusion d'un sous-schéma ouvert de complémentaire $Z$. Soit $\mathcal{F}$ un $\mathcal{O}_U$-module cohérent. Supposons que
\begin{enumerate}
\item $X$ soit universellement caténaire,
\item pour tout $z \in Z$, les fibres formelles de $\mathcal{O}_{X, z}$ vérifient $(S_1)$.
\end{enumerate}
Dans cette situation, les conditions suivantes sont équivalentes :
\begin{enumerate}
\item[(a)] pour $x \in \text{Ass}(\mathcal{F})$ et $z \in Z \cap \overline{\{x\}}$, on a $\dim(\mathcal{O}_{\overline{\{x\}}, z}) \geq 2$,
\item[(b)] $j_*\mathcal{F}$ est cohérent.
\end{enumerate}
\end{lemma}

\begin{proof}
Soit $x \in \text{Ass}(\mathcal{F})$. D'après la proposition \ref{local-cohomology-proposition-kollar}, il suffit de vérifier que $A = \mathcal{O}_{\overline{\{x\}}, z}$ satisfait à la condition de la proposition concernant les idéaux premiers associés de son complété si et seulement si $\dim(A) \geq 2$. Observons que $A$ est universellement caténaire (ce qui est clair) et que ses fibres formelles vérifient $(S_1)$, d'après Compléments d'algèbre, lemme \ref{more-algebra-lemma-formal-fibres-normal} et proposition \ref{more-algebra-proposition-finite-type-over-P-ring}. Soit $\mathfrak p' \subset A^\wedge$ un idéal premier associé. Puisque $A \to A^\wedge$ est plat, Algèbre, lemme \ref{algebra-lemma-bourbaki}, montre que $\mathfrak p'$ est au-dessus de $(0) \subset A$. La fibre formelle $A^\wedge \otimes_A F$ vérifie $(S_1)$, où $F$ est le corps des fractions de $A$. On en déduit que $\mathfrak p'$ est un idéal premier minimal ; voir Algèbre, lemme \ref{algebra-lemma-criterion-no-embedded-primes}. Puisque $A$ est universellement caténaire, il est formellement caténaire d'après Compléments d'algèbre, proposition \ref{more-algebra-proposition-ratliff}. Ainsi $\dim(A^\wedge/\mathfrak p') = \dim(A)$, ce qui prouve l'équivalence.
\end{proof}






\section{Profondeur et dimension}
\label{local-cohomology-section-dept-dimension}

\noindent
Quelques lemmes auxiliaires.

\begin{lemma}
\label{local-cohomology-lemma-ideal-depth-function}
Soit $A$ un anneau noethérien. Soit $I \subset A$ un idéal. Soit $M$ un $A$-module de type fini. Soit $\mathfrak p \in V(I)$ un idéal premier. Supposons $e = \text{depth}_{IA_\mathfrak p}(M_\mathfrak p) < \infty$. Alors il existe un ouvert non vide $U \subset V(\mathfrak p)$ tel que $\text{depth}_{IA_\mathfrak q}(M_\mathfrak q) \geq e$ pour tout $\mathfrak q \in U$.
\end{lemma}

\begin{proof}
Par définition de la profondeur, on a $IM_\mathfrak p \not = M_\mathfrak p$ et il existe une suite $M_\mathfrak p$-régulière $f_1, \ldots, f_e \in IA_\mathfrak p$. Après avoir remplacé $A$ par un localisé principal, on peut supposer que $f_1, \ldots, f_e \in I$ forment une suite $M$-régulière ; voir Algèbre, lemme \ref{algebra-lemma-regular-sequence-in-neighbourhood}. Considérons le module $M' = M/IM$. Puisque $\mathfrak p \in \text{Supp}(M')$ et que le support d'un module de type fini est fermé, on trouve $V(\mathfrak p) \subset \text{Supp}(M')$. Ainsi, pour $\mathfrak q \in V(\mathfrak p)$, on a $IM_\mathfrak q \not = M_\mathfrak q$. Par conséquent, l'exactitude de la localisation et la définition de la profondeur donnent $\text{depth}_{IA_\mathfrak q}(M_\mathfrak q) \geq e$ pour tout $\mathfrak q \in V(I)$.
\end{proof}

\begin{lemma}
\label{local-cohomology-lemma-depth-function}
Soit $A$ un anneau noethérien. Soit $M$ un $A$-module de type fini. Soit $\mathfrak p$ un idéal premier. Supposons $e = \text{depth}_{A_\mathfrak p}(M_\mathfrak p) < \infty$. Alors il existe un ouvert non vide $U \subset V(\mathfrak p)$ tel que $\text{depth}_{A_\mathfrak q}(M_\mathfrak q) \geq e$ pour tout $\mathfrak q \in U$ et tel que, pour tous les $\mathfrak q \in U$ sauf un nombre fini, on ait $\text{depth}_{A_\mathfrak q}(M_\mathfrak q) > e$.
\end{lemma}

\begin{proof}
Par définition de la profondeur, on a $\mathfrak p M_\mathfrak p \not = M_\mathfrak p$ et il existe une suite $M_\mathfrak p$-régulière $f_1, \ldots, f_e \in \mathfrak p A_\mathfrak p$. Après avoir remplacé $A$ par un localisé principal, on peut supposer que $f_1, \ldots, f_e \in \mathfrak p$ forment une suite $M$-régulière ; voir Algèbre, lemme \ref{algebra-lemma-regular-sequence-in-neighbourhood}. Considérons le module $M' = M/(f_1, \ldots, f_e)M$. Puisque $\mathfrak p \in \text{Supp}(M')$ et que le support d'un module de type fini est fermé, on trouve $V(\mathfrak p) \subset \text{Supp}(M')$. Ainsi, pour $\mathfrak q \in V(\mathfrak p)$, on a $\mathfrak q M_\mathfrak q \not = M_\mathfrak q$. Par conséquent, l'exactitude de la localisation et la définition de la profondeur donnent $\text{depth}_{A_\mathfrak q}(M_\mathfrak q) \geq e$ pour tout $\mathfrak q \in V(I)$. De plus, dès que $\mathfrak q$ n'est pas un idéal premier associé au module $M'$, la profondeur augmente. La dernière assertion résulte donc d'Algèbre, lemme \ref{algebra-lemma-finite-ass}.
\end{proof}

\begin{lemma}
\label{local-cohomology-lemma-finite-nr-points-next-S}
Soit $X$ un schéma noethérien muni d'un complexe dualisant $\omega_X^\bullet$. Soit $\mathcal{F}$ un $\mathcal{O}_X$-module cohérent. Soit $k \geq 0$ un entier. Supposons que $\mathcal{F}$ vérifie $(S_k)$. Alors les points $x \in X$ tels que
$$
\text{depth}(\mathcal{F}_x) = k
\quad\text{et}\quad
\dim(\text{Supp}(\mathcal{F}_x)) > k
$$
sont en nombre fini.
\end{lemma}

\begin{proof}
Démontrons ce lemme par récurrence sur $k$. Le cas initial $k = 0$ affirme que $\mathcal{F}$ possède un nombre fini de points associés immergés, ce qui résulte de Diviseurs, lemme \ref{divisors-lemma-finite-ass}.

\medskip\noindent
Supposons $k > 0$ et le résultat établi pour toute valeur plus petite de $k$. On peut recouvrir $X$ par un nombre fini d'ouverts affines ; on peut donc supposer $X = \Spec(A)$ affine. Alors $\mathcal{F}$ est le $\mathcal{O}_X$-module cohérent associé à un $A$-module de type fini $M$ qui vérifie $(S_k)$. Nous utiliserons Algèbre, lemmes \ref{algebra-lemma-one-equation-module} et \ref{algebra-lemma-depth-drops-by-one}, sans le mentionner à nouveau.

\medskip\noindent
Soit $f \in A$ un non-diviseur de zéro dans $M$. Alors $M/fM$ vérifie $(S_{k - 1})$. Par récurrence, les idéaux premiers $\mathfrak p \in V(f)$ tels que $\text{depth}((M/fM)_\mathfrak p) = k - 1$ et $\dim(\text{Supp}((M/fM)_\mathfrak p)) > k - 1$ sont en nombre fini. Ce sont exactement les idéaux premiers $\mathfrak p \in V(f)$ tels que $\text{depth}(M_\mathfrak p) = k$ et $\dim(\text{Supp}(M_\mathfrak p)) > k$. On peut donc remplacer $A$ par $A_f$ et $M$ par $M_f$ pour démontrer l'assertion de finitude.

\medskip\noindent
Puisque $M$ vérifie $(S_k)$ et que $k > 0$, le module $M$ n'a pas d'idéal premier associé immergé (Algèbre, lemme \ref{algebra-lemma-criterion-no-embedded-primes}). Ainsi $\text{Ass}(M)$ est l'ensemble des points génériques du support de $M$. Complexes dualisants, lemme \ref{dualizing-lemma-CM-open}, montre donc que l'ensemble $U = \{\mathfrak q \mid M_\mathfrak q\text{ est de Cohen-Macaulay}\}$ est un ouvert contenant $\text{Ass}(M)$. Par évitement des idéaux premiers (Algèbre, lemme \ref{algebra-lemma-silly}), on peut choisir $f \in A$ tel que $f \not \in \mathfrak p$ pour $\mathfrak p \in \text{Ass}(M)$ et tel que $D(f) \subset U$. Alors $f$ est un non-diviseur de zéro dans $M$ (Algèbre, lemme \ref{algebra-lemma-ass-zero-divisors}). Après avoir remplacé $A$ par $A_f$ et $M$ par $M_f$ (voir ci-dessus), on obtient que $M$ est de Cohen-Macaulay. Ainsi, pour tout $\mathfrak q \subset A$, on a $\dim(M_\mathfrak q) = \text{depth}(M_\mathfrak q)$ ; l'ensemble décrit dans le lemme est donc vide, et a fortiori fini.
\end{proof}

\begin{lemma}
\label{local-cohomology-lemma-sitting-in-degrees}
Soit $(A, \mathfrak m)$ un anneau local noethérien muni d'un complexe dualisant normalisé $\omega_A^\bullet$. Soit $M$ un $A$-module de type fini. Posons $E^i = \text{Ext}_A^{-i}(M, \omega_A^\bullet)$. Alors
\begin{enumerate}
\item $E^i$ est un $A$-module de type fini, non nul seulement pour $0 \leq i \leq \dim(\text{Supp}(M))$,
\item $\dim(\text{Supp}(E^i)) \leq i$,
\item $\text{depth}(M)$ est le plus petit entier $\delta \geq 0$ tel que $E^\delta \not = 0$,
\item $\mathfrak p \in \text{Supp}(E^0 \oplus \ldots \oplus E^i)
\Leftrightarrow
\text{depth}_{A_\mathfrak p}(M_\mathfrak p) + \dim(A/\mathfrak p) \leq i$,
\item l'annulateur de $E^i$ est égal à celui de $H^i_\mathfrak m(M)$.
\end{enumerate}
\end{lemma}

\begin{proof}
Les points (1), (2) et (3) reproduisent les énoncés de Complexes dualisants, lemme \ref{dualizing-lemma-sitting-in-degrees}. Pour un idéal premier $\mathfrak p$ de $A$, le complexe $(\omega_A^\bullet)_\mathfrak p[-\dim(A/\mathfrak p)]$ est un complexe dualisant normalisé pour $A_\mathfrak p$. Voir Complexes dualisants, lemme \ref{dualizing-lemma-dimension-function}. Ainsi

$$
E^i_\mathfrak p =
\text{Ext}^{-i}_A(M, \omega_A^\bullet)_\mathfrak p =
\text{Ext}^{-i + \dim(A/\mathfrak p)}_{A_\mathfrak p}
(M_\mathfrak p, (\omega_A^\bullet)_\mathfrak p[-\dim(A/\mathfrak p)])
$$
est nul pour $i - \dim(A/\mathfrak p) < \text{depth}_{A_\mathfrak p}(M_\mathfrak p)$ et non nul pour $i = \dim(A/\mathfrak p) + \text{depth}_{A_\mathfrak p}(M_\mathfrak p)$, d'après (3) appliqué sur $A_\mathfrak p$. Cela prouve (4). Si $E$ est une enveloppe injective du corps résiduel de $A$, alors on a
$$
\Hom_A(H^i_\mathfrak m(M), E) =
\text{Ext}^{-i}_A(M, \omega_A^\bullet)^\wedge =
(E^i)^\wedge =
E^i \otimes_A A^\wedge
$$
d'après le théorème de dualité locale (sous la forme de Complexes dualisants, lemme \ref{dualizing-lemma-special-case-local-duality}). Puisque $A \to A^\wedge$ est fidèlement plat, la dualité de Matlis (Complexes dualisants, proposition \ref{dualizing-proposition-matlis}) prouve (5).
\end{proof}






\section{Annulateurs de la cohomologie locale, I}
\label{local-cohomology-section-annihilators}

\noindent
Cette section présente un résultat dû à Faltings ; voir \cite{Faltings-annulators}.

\begin{proposition}
\label{local-cohomology-proposition-annihilator}
\begin{reference}
\cite{Faltings-annulators}.
\end{reference}
Soit $A$ un anneau noethérien possédant un complexe dualisant. Soient $T \subset T' \subset \Spec(A)$ des parties stables par spécialisation. Soit $s \geq 0$ un entier. Soit $M$ un $A$-module de type fini. Les conditions suivantes sont équivalentes :
\begin{enumerate}
\item il existe un idéal $J \subset A$ tel que $V(J) \subset T'$ et tel que $J$ annule $H^i_T(M)$ pour $i \leq s$,
\item pour tous $\mathfrak p \not \in T'$, $\mathfrak q \in T$ tels que $\mathfrak p \subset \mathfrak q$, on a
$$
\text{depth}_{A_\mathfrak p}(M_\mathfrak p) +
\dim((A/\mathfrak p)_\mathfrak q) > s
$$
\end{enumerate}
\end{proposition}

\begin{proof}
Soit $\omega_A^\bullet$ un complexe dualisant. Soit $\delta$ sa fonction de dimension ; voir Complexes dualisants, section \ref{dualizing-section-dimension-function}. Les $A$-modules de type fini
$$
E^i = \Ext_A^i(M, \omega_A^\bullet)
$$
joueront un rôle important. Pour $\mathfrak p \subset A$, nous écrirons $H^i_\mathfrak p$ pour désigner la cohomologie locale d'un $A_\mathfrak p$-module relativement à $\mathfrak pA_\mathfrak p$. On voit alors que le complété $\mathfrak pA_\mathfrak p$-adique de
$$
(E^i)_\mathfrak p =
\Ext^{\delta(\mathfrak p) + i}_{A_\mathfrak p}(M_\mathfrak p,
(\omega_A^\bullet)_\mathfrak p[-\delta(\mathfrak p)])
$$
est le dual de Matlis de
$$
H^{-\delta(\mathfrak p) - i}_{\mathfrak p}(M_\mathfrak p)
$$
d'après Complexes dualisants, lemme \ref{dualizing-lemma-special-case-local-duality}. On en déduit en particulier le fait suivant : un idéal $J \subset A$ annule $(E^i)_\mathfrak p$ si et seulement si $J$ annule $H^{-\delta(\mathfrak p) - i}_{\mathfrak p}(M_\mathfrak p)$.

\medskip\noindent
Posons $T_n = \{\mathfrak p \in T \mid \delta(\mathfrak p) \leq n\}$. Puisque $\delta$ est une fonction bornée, on voit que $T_a = \emptyset$ pour $a \ll 0$ et $T_b = T$ pour $b \gg 0$.

\medskip\noindent
Supposons (2). Démontrons l'existence de $J$ comme dans (1). Nous utiliserons une double récurrence. Pour $i \leq s$, considérons l'hypothèse de récurrence $IH_i$ : $H^a_T(M)$ est annulé par un idéal $J \subset A$ tel que $V(J) \subset T'$ pour $0 \leq a \leq i$. Le cas $IH_0$ est trivial, car $H^0_T(M)$ est un sous-module de $M$, donc de type fini, et est par conséquent annulé par un idéal $J$ tel que $V(J) \subset T$.

\medskip\noindent
Étape de récurrence. Supposons $IH_{i - 1}$ vérifiée pour un certain $0 < i \leq s$. Choisissons $J'$ tel que $V(J') \subset T'$, annulant $H^a_T(M)$ pour $0 \leq a \leq i - 1$ (ce que permet l'hypothèse de récurrence). Montrons par récurrence descendante sur $n$ qu'il existe un idéal $J$ tel que $V(J) \subset T'$ et tel que les idéaux premiers associés à $J H^i_T(M)$ appartiennent à $T_n$. Pour $n \ll 0$, cela implique $JH^i_T(M) = 0$ (Algèbre, lemme \ref{algebra-lemma-ass-zero}), et donc $IH_i$ sera vérifiée. Le cas initial $n \gg 0$ est trivial, car $T = T_n$ dans ce cas, et tous les idéaux premiers associés à $H^i_T(M)$ appartiennent à $T$.

\medskip\noindent
Supposons donc donné $J$ ayant la propriété pour $n$. Soit $\mathfrak q \in T_n$. Soit $T_\mathfrak q \subset \Spec(A_\mathfrak q)$ l'image réciproque de $T$. On a $H^j_T(M)_\mathfrak q = H^j_{T_\mathfrak q}(M_\mathfrak q)$ d'après le lemme \ref{local-cohomology-lemma-torsion-change-rings}. Considérons la suite spectrale
$$
H_\mathfrak q^p(H^q_{T_\mathfrak q}(M_\mathfrak q))
\Rightarrow
H^{p + q}_\mathfrak q(M_\mathfrak q)
$$
du lemme \ref{local-cohomology-lemma-local-cohomology-ss}. Nous trouverons ci-dessous un idéal $J'' \subset A$ tel que $V(J'') \subset T'$ et tel que $H^i_\mathfrak q(M_\mathfrak q)$ soit annulé par $J''$ pour tout $\mathfrak q \in T_n \setminus T_{n - 1}$. Affirmation : $J (J')^i J''$ convient pour $n - 1$. En effet, soit $\mathfrak q \in T_n \setminus T_{n - 1}$. La suite spectrale ci-dessus définit une filtration
$$
E_\infty^{0, i} = E_{i + 2}^{0, i} \subset \ldots \subset E_3^{0, i} \subset
E_2^{0, i} = H^0_\mathfrak q(H^i_{T_\mathfrak q}(M_\mathfrak q))
$$
Le module $E_\infty^{0, i}$ est annulé par $J''$. Les sous-quotients $E_j^{0, i}/E_{j + 1}^{0, i}$ pour $i + 1 \geq j \geq 2$ sont annulés par $J'$, car le but de $d_j^{0, i}$ est un sous-quotient de
$$
H^j_\mathfrak q(H^{i - j + 1}_{T_\mathfrak q}(M_\mathfrak q)) =
H^j_\mathfrak q(H^{i - j + 1}_T(M)_\mathfrak q)
$$
et $H^{i - j + 1}_T(M)_\mathfrak q$ est annulé par $J'$, par le choix de $J'$. Enfin, par le choix de $J$, on a $J H^i_T(M)_\mathfrak q \subset H^0_\mathfrak q(H^i_T(M)_\mathfrak q)$, puisque les points non fermés de $\Spec(A_\mathfrak q)$ ont des valeurs de $\delta$ plus grandes. Ainsi $\mathfrak q$ ne peut être un idéal premier associé à $J(J')^iJ'' H^i_T(M)$, comme voulu.

\medskip\noindent
D'après les remarques initiales, $J''$ doit annuler
$$
(E^{-\delta(\mathfrak q) - i})_\mathfrak q =
(E^{-n - i})_\mathfrak q
$$
pour tout $\mathfrak q \in T_n \setminus T_{n - 1}$. Mais si $J''$ convient pour un $\mathfrak q$, il convient pour tous les $\mathfrak q$ d'un voisinage ouvert de $\mathfrak q$, puisque les modules $E^{-n - i}$ sont de type fini. Toute partie de $\Spec(A)$ étant noethérienne pour la topologie induite (Topologie, lemme \ref{topology-lemma-Noetherian}), il suffit donc de prouver l'existence de $J''$ pour un seul $\mathfrak q$.

\medskip\noindent
Les modules Ext étant de type fini, l'existence de $J''$ équivaut à
$$
\text{Supp}(E^{-n - i}) \cap \Spec(A_\mathfrak q) \subset T'.
$$
Cela revient à montrer que le localisé de $E^{-n - i}$ en tout $\mathfrak p \subset \mathfrak q$, $\mathfrak p \not \in T'$ est nul. Par dualité locale sur $A_\mathfrak p$, il faut donc montrer que
$$
H^{i + n - \delta(\mathfrak p)}_\mathfrak p(M_\mathfrak p) =
H^{i - \dim((A/\mathfrak p)_\mathfrak q)}_\mathfrak p(M_\mathfrak p)
$$
est nul (on utilise ici le fait que $\delta$ est une fonction de dimension). Cette annulation résulte de l'hypothèse du lemme, de $i \leq s$ et de Complexes dualisants, lemme \ref{dualizing-lemma-depth}.

\medskip\noindent
Pour démontrer l'implication réciproque, supposons que la condition (2) ne soit pas vérifiée et reprenons les arguments précédents en sens inverse. Choisissons d'abord $\mathfrak q \in T$, $\mathfrak p \subset \mathfrak q$, avec $\mathfrak p \not \in T'$, tels que
$$
i = \text{depth}_{A_\mathfrak p}(M_\mathfrak p) +
\dim((A/\mathfrak p)_\mathfrak q) \leq s
$$
soit minimal. Alors $H^{i - \dim((A/\mathfrak p)_\mathfrak q)}_\mathfrak p(M_\mathfrak p)$ est non nul d'après la non-annulation de Complexes dualisants, lemme \ref{dualizing-lemma-depth}. Posons $n = \delta(\mathfrak q)$. Il n'existe alors aucun idéal $J \subset A$ tel que $V(J) \subset T'$ et $J(E^{-n - i})_\mathfrak q = 0$. Ainsi $H^i_\mathfrak q(M_\mathfrak q)$ n'est annulé par aucun idéal $J \subset A$ tel que $V(J) \subset T'$. Par minimalité de $i$, la suite spectrale ci-dessus montre que le module $H^i_T(M)_\mathfrak q$ n'est annulé par aucun idéal $J \subset A$ tel que $V(J) \subset T'$. Ainsi $H^i_T(M)$ n'est annulé par aucun idéal $J \subset A$ tel que $V(J) \subset T'$. Cela achève la démonstration de la proposition.
\end{proof}

\begin{lemma}
\label{local-cohomology-lemma-kill-local-cohomology-at-prime}
Soit $I$ un idéal d'un anneau noethérien $A$. Soient $M$ un $A$-module de type fini, $\mathfrak p \subset A$ un idéal premier et $s \geq 0$ un entier. Supposons que
\begin{enumerate}
\item $A$ possède un complexe dualisant,
\item $\mathfrak p \not \in V(I)$,
\item pour tous les idéaux premiers $\mathfrak p' \subset \mathfrak p$ et $\mathfrak q \in V(I)$ tels que $\mathfrak p' \subset \mathfrak q$, on ait
$$
\text{depth}_{A_{\mathfrak p'}}(M_{\mathfrak p'}) +
\dim((A/\mathfrak p')_\mathfrak q) > s
$$
\end{enumerate}
Alors il existe $f \in A$, $f \not \in \mathfrak p$, qui annule $H^i_{V(I)}(M)$ pour $i \leq s$.
\end{lemma}

\begin{proof}
Considérons les ensembles
$$
T = V(I)
\quad\text{et}\quad
T' = \bigcup\nolimits_{f \in A, f \not \in \mathfrak p} V(f)
$$
Ce sont des parties de $\Spec(A)$ stables par spécialisation. Observons que $T \subset T'$ et $\mathfrak p \not \in T'$. L'hypothèse (3) signifie que l'hypothèse (2) de la proposition \ref{local-cohomology-proposition-annihilator} est vérifiée. On peut donc trouver $J \subset A$ tel que $V(J) \subset T'$ et tel que $J H^i_{V(I)}(M) = 0$ pour $i \leq s$. Choisissons $f \in A$, $f \not \in \mathfrak p$, tel que $V(J) \subset V(f)$. Une puissance de $f$ annule $H^i_{V(I)}(M)$ pour $i \leq s$.
\end{proof}





\section{Finitude de la cohomologie locale, II}
\label{local-cohomology-section-finiteness-II}

\noindent
Nous poursuivons l'étude de la finitude de la cohomologie locale commencée à la section \ref{local-cohomology-section-finiteness}. Le théorème des annulateurs de Faltings permet de démontrer aisément le résultat fondamental suivant.

\begin{proposition}
\label{local-cohomology-proposition-finiteness}
\begin{reference}
\cite{Faltings-annulators}.
\end{reference}
Soit $A$ un anneau noethérien possédant un complexe dualisant. Soit $T \subset \Spec(A)$ une partie stable par spécialisation. Soit $s \geq 0$ un entier. Soit $M$ un $A$-module de type fini. Les conditions suivantes sont équivalentes :
\begin{enumerate}
\item $H^i_T(M)$ est un $A$-module de type fini pour $i \leq s$,
\item pour tous $\mathfrak p \not \in T$, $\mathfrak q \in T$ tels que $\mathfrak p \subset \mathfrak q$, on a
$$
\text{depth}_{A_\mathfrak p}(M_\mathfrak p) +
\dim((A/\mathfrak p)_\mathfrak q) > s
$$
\end{enumerate}
\end{proposition}

\begin{proof}
Conséquence formelle de la proposition \ref{local-cohomology-proposition-annihilator} et du lemme \ref{local-cohomology-lemma-check-finiteness-local-cohomology-by-annihilator}.
\end{proof}

\noindent
Outre quelques lemmes destinés à un usage ultérieur, le reste de cette section examine dans quelle mesure on peut affaiblir la condition de la proposition \ref{local-cohomology-proposition-finiteness} selon laquelle $A$ possède un complexe dualisant. La réponse est, en gros, qu'il faut supposer que les fibres formelles de $A$ vérifient $(S_n)$ pour $n$ assez grand.

\medskip\noindent
Soient $A$ un anneau noethérien et $I \subset A$ un idéal. Posons $X = \Spec(A)$ et $Z = V(I) \subset X$. Soit $M$ un $A$-module de type fini. Définissons
\begin{equation}
\label{local-cohomology-equation-cutoff}
s_{A, I}(M) =
\min \{
\text{depth}_{A_\mathfrak p}(M_\mathfrak p) + \dim((A/\mathfrak p)_\mathfrak q)
\mid
\mathfrak p \in X \setminus Z, \mathfrak q \in Z,
\mathfrak p \subset \mathfrak q
\}
\end{equation}
Notre convention sur la profondeur est que la profondeur de $0$ vaut $\infty$ ; il suffit donc de considérer les idéaux premiers $\mathfrak p$ du support de $M$. Nous verrons que $s_{A, I}(M)$ est un invariant important de la situation.

\begin{lemma}
\label{local-cohomology-lemma-cutoff}
Soit $A \to B$ un homomorphisme fini d'anneaux noethériens. Soit $I \subset A$ un idéal et posons $J = IB$. Soit $M$ un $B$-module de type fini. Si $A$ est universellement caténaire, alors $s_{B, J}(M) = s_{A, I}(M)$.
\end{lemma}

\begin{proof}
Soient $\mathfrak p \subset \mathfrak q \subset A$ des idéaux premiers tels que $I \subset \mathfrak q$ et $I \not \subset \mathfrak p$. Puisque $A \to B$ est fini, les idéaux premiers $\mathfrak p_i$ au-dessus de $\mathfrak p$ sont en nombre fini. D'après Algèbre, lemme \ref{algebra-lemma-depth-goes-down-finite}, on a
$$
\text{depth}(M_\mathfrak p) = \min \text{depth}(M_{\mathfrak p_i})
$$
Soient $\mathfrak p_i \subset \mathfrak q_{ij}$ des idéaux premiers au-dessus de $\mathfrak q$. Par le théorème de montée pour $A \to B$ (Algèbre, lemme \ref{algebra-lemma-integral-going-up}), il existe au moins un $\mathfrak q_{ij}$ pour chaque $i$. Alors on a
$$
\dim((B/\mathfrak p_i)_{\mathfrak q_{ij}}) =
\dim((A/\mathfrak p)_\mathfrak q)
$$
d'après la formule des dimensions ; voir Algèbre, lemme \ref{algebra-lemma-dimension-formula}. Cela implique que le minimum des quantités utilisées pour définir $s_{B, J}(M)$ pour les couples $(\mathfrak p_i, \mathfrak q_{ij})$ est égal à la quantité correspondant au couple $(\mathfrak p, \mathfrak q)$. Cela prouve le lemme.
\end{proof}

\begin{lemma}
\label{local-cohomology-lemma-change-completion}
Soit $A$ un anneau noethérien possédant un complexe dualisant. Soit $I \subset A$ un idéal. Soit $M$ un $A$-module de type fini. Soient $A', M'$ les complétés $I$-adiques de $A, M$. Soient $\mathfrak p' \subset \mathfrak q'$ des idéaux premiers de $A'$, avec $\mathfrak q' \in V(IA')$, au-dessus de $\mathfrak p \subset \mathfrak q$ dans $A$. Alors
$$
\text{depth}_{A_{\mathfrak p'}}(M'_{\mathfrak p'})
\geq
\text{depth}_{A_\mathfrak p}(M_\mathfrak p)
$$
et
$$
\text{depth}_{A_{\mathfrak p'}}(M'_{\mathfrak p'}) +
\dim((A'/\mathfrak p')_{\mathfrak q'}) =
\text{depth}_{A_\mathfrak p}(M_\mathfrak p) +
\dim((A/\mathfrak p)_\mathfrak q)
$$
\end{lemma}

\begin{proof}
On a
$$
\text{depth}(M'_{\mathfrak p'}) =
\text{depth}(M_\mathfrak p) +
\text{depth}(A'_{\mathfrak p'}/\mathfrak p A'_{\mathfrak p'})
\geq \text{depth}(M_\mathfrak p)
$$
par platitude de $A \to A'$ ; voir Algèbre, lemme \ref{algebra-lemma-apply-grothendieck-module}. Puisque les fibres de $A \to A'$ sont de Cohen-Macaulay (Complexes dualisants, lemme \ref{dualizing-lemma-dualizing-gorenstein-formal-fibres}, et Compléments d'algèbre, section \ref{more-algebra-section-properties-formal-fibres}), on voit que $\text{depth}(A'_{\mathfrak p'}/\mathfrak p A'_{\mathfrak p'}) =
\dim(A'_{\mathfrak p'}/\mathfrak p A'_{\mathfrak p'})$. On obtient donc
\begin{align*}
\text{depth}(M'_{\mathfrak p'}) +
\dim((A'/\mathfrak p')_{\mathfrak q'})
& =
\text{depth}(M_\mathfrak p) +
\dim(A'_{\mathfrak p'}/\mathfrak p A'_{\mathfrak p'}) +
\dim((A'/\mathfrak p')_{\mathfrak q'}) \\
& =
\text{depth}(M_\mathfrak p) +
\dim((A'/\mathfrak p A')_{\mathfrak q'}) \\
& =
\text{depth}(M_\mathfrak p) +
\dim((A/\mathfrak p)_\mathfrak q)
\end{align*}
La deuxième égalité vient du fait que $A'$ est caténaire et la troisième de Compléments d'algèbre, lemme \ref{more-algebra-lemma-completion-dimension}, puisque $(A/\mathfrak p)_\mathfrak q$ et $(A'/\mathfrak p A')_{\mathfrak q'}$ ont les mêmes complétés $I$-adiques.
\end{proof}





\begin{lemma}
\label{local-cohomology-lemma-cutoff-completion}
Soit $A$ un anneau local noethérien universellement caténaire. Soit $I \subset A$ un idéal. Soit $M$ un $A$-module de type fini. Alors
$$
s_{A, I}(M) \geq s_{A^\wedge, I^\wedge}(M^\wedge)
$$
Si les fibres formelles de $A$ vérifient $(S_n)$, alors $\min(n + 1, s_{A, I}(M)) \leq s_{A^\wedge, I^\wedge}(M^\wedge)$.
\end{lemma}

\begin{proof}
Écrivons $X = \Spec(A)$, $X^\wedge = \Spec(A^\wedge)$, $Z = V(I) \subset X$ et $Z^\wedge = V(I^\wedge)$. Soient $\mathfrak p' \subset \mathfrak q' \subset A^\wedge$ des idéaux premiers tels que $\mathfrak p' \not \in Z^\wedge$ et $\mathfrak q' \in Z^\wedge$. Soient $\mathfrak p \subset \mathfrak q$ les idéaux premiers correspondants de $A$. Alors $\mathfrak p \not \in Z$ et $\mathfrak q \in Z$. On a le diagramme
$$
\xymatrix{
\mathfrak p' \ar[r] & \mathfrak q' \ar[r] & A^\wedge \\
\mathfrak p \ar[r] \ar@{-}[u] &
\mathfrak q \ar[r] \ar@{-}[u] & A \ar[u]
}
$$
Écrivons
\begin{align*}
a & = \dim(A/\mathfrak p) = \dim(A^\wedge/\mathfrak pA^\wedge),\\
b & = \dim(A/\mathfrak q) = \dim(A^\wedge/\mathfrak qA^\wedge),\\
a' & = \dim(A^\wedge/\mathfrak p'),\\
b' & = \dim(A^\wedge/\mathfrak q')
\end{align*}
Les égalités résultent de Compléments d'algèbre, lemme \ref{more-algebra-lemma-completion-dimension}. Écrivons aussi
\begin{align*}
p & = \dim(A^\wedge_{\mathfrak p'}/\mathfrak p A^\wedge_{\mathfrak p'}) =
\dim((A^\wedge/\mathfrak p A^\wedge)_{\mathfrak p'}) \\
q & = \dim(A^\wedge_{\mathfrak q'}/\mathfrak p A^\wedge_{\mathfrak q'}) =
\dim((A^\wedge/\mathfrak q A^\wedge)_{\mathfrak q'})
\end{align*}
Puisque $A$ est universellement caténaire, $A^\wedge/\mathfrak pA^\wedge = (A/\mathfrak p)^\wedge$ est équidimensionnel de dimension $a$ (Compléments d'algèbre, proposition \ref{more-algebra-proposition-ratliff}). Ainsi $a = a' + p$. De même, $b = b' + q$. D'après Algèbre, lemme \ref{algebra-lemma-apply-grothendieck-module}, appliqué à l'homomorphisme local plat $A_\mathfrak p \to A^\wedge_{\mathfrak p'}$, on a
$$
\text{depth}(M^\wedge_{\mathfrak p'})
=
\text{depth}(M_\mathfrak p) +
\text{depth}(A^\wedge_{\mathfrak p'} / \mathfrak p A^\wedge_{\mathfrak p'})
$$
La quantité dont on prend le minimum pour $s_{A, I}(M)$ est
$$
s(\mathfrak p, \mathfrak q) =
\text{depth}(M_\mathfrak p) + \dim((A/\mathfrak p)_\mathfrak q) =
\text{depth}(M_\mathfrak p) + a - b
$$
(la dernière égalité vient du fait que $A$ est caténaire). La quantité dont on prend le minimum pour $s_{A^\wedge, I^\wedge}(M^\wedge)$ est
$$
s(\mathfrak p', \mathfrak q') =
\text{depth}(M^\wedge_{\mathfrak p'}) +
\dim((A^\wedge/\mathfrak p')_{\mathfrak q'}) =
\text{depth}(M^\wedge_{\mathfrak p'}) + a' - b'
$$
(la dernière égalité vient du fait que $A^\wedge$ est caténaire). Nous disposons maintenant des notations nécessaires pour commencer la démonstration.

\medskip\noindent
Soient $\mathfrak p \subset \mathfrak q \subset A$ des idéaux premiers tels que $\mathfrak p \not \in Z$, $\mathfrak q \in Z$ et $s_{A, I}(M) = s(\mathfrak p, \mathfrak q)$. On peut alors choisir $\mathfrak q'$ minimal au-dessus de $\mathfrak q A^\wedge$ et $\mathfrak p' \subset \mathfrak q'$ minimal au-dessus de $\mathfrak p A^\wedge$ (par le théorème de descente pour $A \to A^\wedge$). On dispose alors de quatre idéaux premiers comme ci-dessus, avec $p = 0$ et $q = 0$. De plus, on a $\text{depth}(A^\wedge_{\mathfrak p'} / \mathfrak p A^\wedge_{\mathfrak p'})=0$, là encore parce que $p = 0$. Cela signifie que $s(\mathfrak p', \mathfrak q') = s(\mathfrak p, \mathfrak q)$. On obtient ainsi la première inégalité.

\medskip\noindent
Supposons que les fibres formelles de $A$ vérifient $(S_n)$. Alors $\text{depth}(A^\wedge_{\mathfrak p'} / \mathfrak p A^\wedge_{\mathfrak p'})
\geq \min(n, p)$. D'où
$$
s(\mathfrak p', \mathfrak q') \geq
s(\mathfrak p, \mathfrak q) + q + \min(n, p) - p \geq
s_{A, I}(M) + q + \min(n, p) - p
$$
Le seul cas susceptible de poser problème est donc $p > n$. Dans ce cas,
\begin{align*}
s(\mathfrak p', \mathfrak q')
& =
\text{depth}(M^\wedge_{\mathfrak p'}) +
\dim((A^\wedge/\mathfrak p')_{\mathfrak q'}) \\
& =
\text{depth}(M_\mathfrak p) +
\text{depth}(A^\wedge_{\mathfrak p'} / \mathfrak p A^\wedge_{\mathfrak p'}) +
\dim((A^\wedge/\mathfrak p')_{\mathfrak q'}) \\
& \geq
0 + n + 1
\end{align*}
car $(A^\wedge/\mathfrak p')_{\mathfrak q'}$ possède au moins deux idéaux premiers. Cela prouve la seconde inégalité.
\end{proof}

\noindent
La méthode de démonstration du lemme suivant s'applique plus généralement, mais les résultats plus forts ainsi obtenus seront englobés par le théorème \ref{local-cohomology-theorem-finiteness} ci-dessous.

\begin{lemma}
\label{local-cohomology-lemma-local-annihilator}
\begin{reference}
C'est un cas particulier de \cite[Satz 1]{Faltings-annulators}.
\end{reference}
Soit $A$ un anneau local noethérien de Gorenstein. Soit $I \subset A$ un idéal et posons $Z = V(I) \subset \Spec(A)$. Soit $M$ un $A$-module de type fini. Soit $s = s_{A, I}(M)$ comme dans (\ref{local-cohomology-equation-cutoff}). Alors $H^i_Z(M)$ est de type fini pour $i < s$, mais $H^s_Z(M)$ n'est pas de type fini.
\end{lemma}

\begin{proof}
Puisqu'un anneau local de Gorenstein possède un complexe dualisant, c'est un cas particulier de la proposition \ref{local-cohomology-proposition-finiteness}. Il serait utile de disposer d'une démonstration courte de ce cas particulier, qui interviendra ci-dessous dans la démonstration d'un théorème général de finitude.
\end{proof}

\noindent
Observons que les hypothèses du théorème suivant sont satisfaites par les anneaux noethériens excellents (par définition), par les anneaux noethériens possédant un complexe dualisant (Complexes dualisants, lemme \ref{dualizing-lemma-universally-catenary}, et Complexes dualisants, lemme \ref{dualizing-lemma-dualizing-gorenstein-formal-fibres}), ainsi que par les quotients d'anneaux noethériens réguliers.

\begin{theorem}
\label{local-cohomology-theorem-finiteness}
\begin{reference}
C'est un cas particulier de \cite[Satz 2]{Faltings-finiteness}.
\end{reference}
Soient $A$ un anneau noethérien et $I \subset A$ un idéal. Posons $Z = V(I) \subset \Spec(A)$. Soit $M$ un $A$-module de type fini. Posons $s = s_{A, I}(M)$ comme dans (\ref{local-cohomology-equation-cutoff}). Supposons que
\begin{enumerate}
\item $A$ soit universellement caténaire,
\item les fibres formelles des anneaux locaux de $A$ soient de Cohen-Macaulay.
\end{enumerate}
Alors $H^i_Z(M)$ est de type fini pour $0 \leq i < s$ et $H^s_Z(M)$ n'est pas de type fini.
\end{theorem}

\begin{proof}
D'après le lemme \ref{local-cohomology-lemma-check-finiteness-local-cohomology-locally}, on peut supposer que $A$ est un anneau local.

\medskip\noindent
Si $A$ est un anneau local noethérien complet, on peut écrire $A$ comme quotient d'un anneau local régulier complet $B$ par le théorème de structure de Cohen (Algèbre, théorème \ref{algebra-theorem-cohen-structure-theorem}). À l'aide du lemme \ref{local-cohomology-lemma-cutoff} et de Complexes dualisants, lemme \ref{dualizing-lemma-local-cohomology-and-restriction}, on se ramène au cas d'un anneau local régulier, qui résulte du lemme \ref{local-cohomology-lemma-local-annihilator} puisqu'un anneau local régulier est de Gorenstein (Complexes dualisants, lemme \ref{dualizing-lemma-regular-gorenstein}).

\medskip\noindent
Soit $A$ un anneau local noethérien. Soit $\mathfrak m$ son idéal maximal. On peut supposer $I \subset \mathfrak m$, sinon le lemme est trivial. Soit $A^\wedge$ le complété de $A$, soit $Z^\wedge = V(IA^\wedge)$, et soit $M^\wedge = M \otimes_A A^\wedge$ le complété de $M$ (Algèbre, lemme \ref{algebra-lemma-completion-tensor}). Alors $H^i_Z(M) \otimes_A A^\wedge = H^i_{Z^\wedge}(M^\wedge)$ d'après Complexes dualisants, lemme \ref{dualizing-lemma-torsion-change-rings}, et la platitude de $A \to A^\wedge$ (Algèbre, lemme \ref{algebra-lemma-completion-flat}). Il suffit donc de montrer que $H^i_{Z^\wedge}(M^\wedge)$ est de type fini pour $i < s$ et n'est pas de type fini pour $i = s$ ; voir Algèbre, lemme \ref{algebra-lemma-descend-properties-modules}. Puisque le résultat est connu pour $A^\wedge$, il suffit de montrer que $s_{A, I}(M) = s_{A^\wedge, I^\wedge}(M^\wedge)$. Cela résulte du lemme \ref{local-cohomology-lemma-cutoff-completion}.
\end{proof}

\begin{remark}
\label{local-cohomology-remark-astute-reader}
Le lecteur attentif aura remarqué qu'une condition légèrement plus faible sur les fibres formelles des anneaux locaux de $A$ suffit. Plus précisément, dans la situation du théorème \ref{local-cohomology-theorem-finiteness}, supposons $A$ universellement caténaire, sans faire d'hypothèse sur les fibres formelles. Supposons donné $n$ et cherchons à montrer que les $H^i_Z(M)$ sont de type fini pour $i \leq n$. La même démonstration montre qu'il suffit d'avoir $s_{A, I}(M) > n$ et que les fibres formelles des anneaux locaux de $A$ vérifient $(S_n)$. D'autre part, si l'on veut montrer que $H^s_Z(M)$ n'est pas de type fini, où $s = s_{A, I}(M)$, nos arguments le prouvent lorsque les fibres formelles vérifient $(S_{s - 1})$.
\end{remark}







\section{Finitude des images directes, II}
\label{local-cohomology-section-finiteness-pushforward-II}

\noindent
Cette section poursuit la section \ref{local-cohomology-section-finiteness-pushforward}. Nous y recueillons les fruits du travail effectué à la section \ref{local-cohomology-section-finiteness-II}.

\begin{lemma}
\label{local-cohomology-lemma-finiteness-Rjstar}
Soit $X$ un schéma localement noethérien. Soit $j : U \to X$ l'inclusion d'un sous-schéma ouvert de complémentaire $Z$. Soit $\mathcal{F}$ un $\mathcal{O}_U$-module cohérent. Soit $n \geq 0$ un entier. Supposons que
\begin{enumerate}
\item $X$ soit universellement caténaire,
\item pour tout $z \in Z$, les fibres formelles de $\mathcal{O}_{X, z}$ vérifient $(S_n)$.
\end{enumerate}
Dans cette situation, les conditions suivantes sont équivalentes :
\begin{enumerate}
\item[(a)] pour $x \in \text{Supp}(\mathcal{F})$ et $z \in Z \cap \overline{\{x\}}$, on a $\text{depth}_{\mathcal{O}_{X, x}}(\mathcal{F}_x) +
\dim(\mathcal{O}_{\overline{\{x\}}, z}) > n$,
\item[(b)] $R^pj_*\mathcal{F}$ est cohérent pour $0 \leq p < n$.
\end{enumerate}
\end{lemma}

\begin{proof}
L'énoncé est local sur $X$ ; on peut donc supposer $X$ affine. Écrivons $X = \Spec(A)$ et $Z = V(I)$. Soit $M$ un $A$-module de type fini dont le $\mathcal{O}_X$-module cohérent associé se restreint à $\mathcal{F}$ sur $U$ ; voir le lemme \ref{local-cohomology-lemma-finiteness-pushforwards-and-H1-local}. Ce lemme dit aussi que $R^pj_*\mathcal{F}$ est cohérent si et seulement si $H^{p + 1}_Z(M)$ est un $A$-module de type fini. Observons que le minimum des expressions $\text{depth}_{\mathcal{O}_{X, x}}(\mathcal{F}_x) +
\dim(\mathcal{O}_{\overline{\{x\}}, z})$ est le nombre $s_{A, I}(M)$ de (\ref{local-cohomology-equation-cutoff}). Cela étant dit, le lemme résulte du théorème \ref{local-cohomology-theorem-finiteness}, avec les précisions apportées par la remarque \ref{local-cohomology-remark-astute-reader}.
\end{proof}

\begin{lemma}
\label{local-cohomology-lemma-finiteness-for-finite-locally-free}
Soit $X$ un schéma localement noethérien. Soit $j : U \to X$ l'inclusion d'un sous-schéma ouvert de complémentaire $Z$. Soit $n \geq 0$ un entier. Si $R^pj_*\mathcal{O}_U$ est cohérent pour $0 \leq p < n$, il en est de même de $R^pj_*\mathcal{F}$, $0 \leq p < n$, pour tout $\mathcal{O}_U$-module localement libre de type fini $\mathcal{F}$.
\end{lemma}

\begin{proof}
La question est locale sur $X$ ; on peut donc supposer $X$ affine. Écrivons $X = \Spec(A)$ et $Z = V(I)$. Par l'intermédiaire du lemme \ref{local-cohomology-lemma-finiteness-pushforwards-and-H1-local}, notre lemme résulte du lemme \ref{local-cohomology-lemma-local-finiteness-for-finite-locally-free}.
\end{proof}

\begin{lemma}
\label{local-cohomology-lemma-annihilate-Hp}
\begin{reference}
\cite[Lemma 1.9]{Bhatt-local}
\end{reference}
Soient $A$ un anneau et $J \subset I \subset A$ des idéaux de type fini. Soit $p \geq 0$ un entier. Posons $U = \Spec(A) \setminus V(I)$. Si $H^p(U, \mathcal{O}_U)$ est annulé par $J^n$ pour un certain $n$, alors $H^p(U, \mathcal{F})$ est annulé par $J^m$ pour un certain $m = m(\mathcal{F})$, pour tout $\mathcal{O}_U$-module localement libre de type fini $\mathcal{F}$.
\end{lemma}

\begin{proof}
Considérons l'annulateur $\mathfrak a$ de $H^p(U, \mathcal{F})$. Soit $u \in U$. Il existe un voisinage ouvert $u \in U' \subset U$ et un isomorphisme $\varphi : \mathcal{O}_{U'}^{\oplus r} \to \mathcal{F}|_{U'}$. Choisissons $f \in A$ tel que $u \in D(f) \subset U'$. Il existe des applications
$$
a : \mathcal{O}_U^{\oplus r} \longrightarrow \mathcal{F}
\quad\text{et}\quad
b : \mathcal{F} \longrightarrow \mathcal{O}_U^{\oplus r}
$$
dont les restrictions à $D(f)$ sont égales à $f^N \varphi$ et $f^N \varphi^{-1}$ pour un certain $N$. De plus, on peut supposer que $a \circ b$ et $b \circ a$ sont égales à la multiplication par $f^{2N}$. Cela résulte de Propriétés, lemme \ref{properties-lemma-section-maps-backwards}, puisque $U$ est quasi-compact ($I$ est de type fini), séparé, et que $\mathcal{F}$ et $\mathcal{O}_U^{\oplus r}$ sont de présentation finie. On voit donc que $H^p(U, \mathcal{F})$ est annulé par $f^{2N}J^n$, c'est-à-dire que $f^{2N}J^n \subset \mathfrak a$.

\medskip\noindent
Puisque $U$ est quasi-compact, on peut trouver un nombre fini de $f_1, \ldots, f_t$ et de $N_1, \ldots, N_t$ tels que $U = \bigcup D(f_i)$ et $f_i^{2N_i}J^n \subset \mathfrak a$. Alors $V(I) = V(f_1, \ldots, f_t)$ et, puisque $I$ est de type fini, on en déduit $I^M \subset (f_1, \ldots, f_t)$ pour un certain $M$. Finalement, on voit que $J^m \subset \mathfrak a$ pour $m \gg 0$ ; par exemple, $m = M (2N_1 + \ldots + 2N_t) n$ convient.
\end{proof}











\section{Annulateurs de la cohomologie locale, II}
\label{local-cohomology-section-annihilators-II}

\noindent
Nous étendons la discussion des annulateurs de la cohomologie locale de la section \ref{local-cohomology-section-annihilators} aux complexes bornés inférieurement dont les modules de cohomologie sont de type fini.

\begin{definition}
\label{local-cohomology-definition-depth-complex}
Soit $I$ un idéal d'un anneau noethérien $A$. Soit $K \in D^+_{\textit{Coh}}(A)$. On définit la {\it $I$-profondeur} de $K$, notée $\text{depth}_I(K)$, comme le plus grand $m \in \mathbf{Z} \cup \{\infty\}$ tel que $H^i_I(K) = 0$ pour tout $i < m$. Si $A$ est local d'idéal maximal $\mathfrak m$, on appelle simplement $\text{depth}_\mathfrak m(K)$ la {\it profondeur} de $K$.
\end{definition}

\noindent
Cette définition est compatible avec Algèbre, définition \ref{algebra-definition-depth}, d'après Complexes dualisants, lemme \ref{dualizing-lemma-depth}.

\begin{proposition}
\label{local-cohomology-proposition-annihilator-complex}
Soit $A$ un anneau noethérien possédant un complexe dualisant. Soient $T \subset T' \subset \Spec(A)$ des parties stables par spécialisation. Soit $s \in \mathbf{Z}$. Soit $K$ un objet de $D_{\textit{Coh}}^+(A)$. Les conditions suivantes sont équivalentes :
\begin{enumerate}
\item il existe un idéal $J \subset A$ tel que $V(J) \subset T'$ et tel que $J$ annule $H^i_T(K)$ pour $i \leq s$,
\item pour tous $\mathfrak p \not \in T'$, $\mathfrak q \in T$ tels que $\mathfrak p \subset \mathfrak q$, on a
$$
\text{depth}_{A_\mathfrak p}(K_\mathfrak p) +
\dim((A/\mathfrak p)_\mathfrak q) > s
$$
\end{enumerate}
\end{proposition}

\begin{proof}
Cette proposition est la généralisation naturelle de la proposition \ref{local-cohomology-proposition-annihilator}, dont le lecteur est invité à lire d'abord la démonstration. Soit $\omega_A^\bullet$ un complexe dualisant. Soit $\delta$ sa fonction de dimension ; voir Complexes dualisants, section \ref{dualizing-section-dimension-function}. Les $A$-modules de type fini
$$
E^i = \Ext_A^i(K, \omega_A^\bullet)
$$
joueront un rôle important. Pour $\mathfrak p \subset A$, nous écrirons $H^i_\mathfrak p$ pour désigner la cohomologie locale d'un objet de $D(A_\mathfrak p)$ relativement à $\mathfrak pA_\mathfrak p$. On voit alors que le complété $\mathfrak pA_\mathfrak p$-adique de
$$
(E^i)_\mathfrak p =
\Ext^{\delta(\mathfrak p) + i}_{A_\mathfrak p}(K_\mathfrak p,
(\omega_A^\bullet)_\mathfrak p[-\delta(\mathfrak p)])
$$
est le dual de Matlis de
$$
H^{-\delta(\mathfrak p) - i}_{\mathfrak p}(K_\mathfrak p)
$$
d'après Complexes dualisants, lemme \ref{dualizing-lemma-special-case-local-duality}. On en déduit en particulier le fait suivant : un idéal $J \subset A$ annule $(E^i)_\mathfrak p$ si et seulement si $J$ annule $H^{-\delta(\mathfrak p) - i}_{\mathfrak p}(K_\mathfrak p)$.

\medskip\noindent
Posons $T_n = \{\mathfrak p \in T \mid \delta(\mathfrak p) \leq n\}$. Puisque $\delta$ est une fonction bornée, on voit que $T_a = \emptyset$ pour $a \ll 0$ et $T_b = T$ pour $b \gg 0$.

\medskip\noindent
Supposons (2). Démontrons l'existence de $J$ comme dans (1). Nous utiliserons une double récurrence. Pour $i \leq s$, considérons l'hypothèse de récurrence $IH_i$ : $H^a_T(K)$ est annulé par un idéal $J \subset A$ tel que $V(J) \subset T'$ pour $a \leq i$. Le cas $IH_i$ est trivial pour $i$ assez petit, puisque $K$ est borné inférieurement.

\medskip\noindent
Étape de récurrence. Supposons $IH_{i - 1}$ vérifiée pour un certain $i \leq s$. Choisissons $J'$ tel que $V(J') \subset T'$, annulant $H^a_T(K)$ pour $a \leq i - 1$ (ce que permet l'hypothèse de récurrence). Montrons par récurrence descendante sur $n$ qu'il existe un idéal $J$ tel que $V(J) \subset T'$ et tel que les idéaux premiers associés à $J H^i_T(K)$ appartiennent à $T_n$. Pour $n \ll 0$, cela implique $JH^i_T(K) = 0$ (Algèbre, lemme \ref{algebra-lemma-ass-zero}), et donc $IH_i$ sera vérifiée. Le cas initial $n \gg 0$ est trivial, car $T = T_n$ dans ce cas, et tous les idéaux premiers associés à $H^i_T(K)$ appartiennent à $T$.

\medskip\noindent
Supposons donc donné $J$ ayant la propriété pour $n$. Soit $\mathfrak q \in T_n$. Soit $T_\mathfrak q \subset \Spec(A_\mathfrak q)$ l'image réciproque de $T$. On a $H^j_T(K)_\mathfrak q = H^j_{T_\mathfrak q}(K_\mathfrak q)$ d'après le lemme \ref{local-cohomology-lemma-torsion-change-rings}. Considérons la suite spectrale
$$
H_\mathfrak q^p(H^q_{T_\mathfrak q}(K_\mathfrak q))
\Rightarrow
H^{p + q}_\mathfrak q(K_\mathfrak q)
$$
du lemme \ref{local-cohomology-lemma-local-cohomology-ss}. Nous trouverons ci-dessous un idéal $J'' \subset A$ tel que $V(J'') \subset T'$ et tel que $H^i_\mathfrak q(K_\mathfrak q)$ soit annulé par $J''$ pour tout $\mathfrak q \in T_n \setminus T_{n - 1}$. Affirmation : $J (J')^i J''$ convient pour $n - 1$. En effet, soit $\mathfrak q \in T_n \setminus T_{n - 1}$. La suite spectrale ci-dessus définit une filtration
$$
E_\infty^{0, i} = E_{i + 2}^{0, i} \subset \ldots \subset E_3^{0, i} \subset
E_2^{0, i} = H^0_\mathfrak q(H^i_{T_\mathfrak q}(K_\mathfrak q))
$$
Le module $E_\infty^{0, i}$ est annulé par $J''$. Les sous-quotients $E_j^{0, i}/E_{j + 1}^{0, i}$ pour $i + 1 \geq j \geq 2$ sont annulés par $J'$, car le but de $d_j^{0, i}$ est un sous-quotient de
$$
H^j_\mathfrak q(H^{i - j + 1}_{T_\mathfrak q}(K_\mathfrak q)) =
H^j_\mathfrak q(H^{i - j + 1}_T(K)_\mathfrak q)
$$
et $H^{i - j + 1}_T(K)_\mathfrak q$ est annulé par $J'$, par le choix de $J'$. Enfin, par le choix de $J$, on a $J H^i_T(K)_\mathfrak q \subset H^0_\mathfrak q(H^i_T(K)_\mathfrak q)$, puisque les points non fermés de $\Spec(A_\mathfrak q)$ ont des valeurs de $\delta$ plus grandes. Ainsi $\mathfrak q$ ne peut être un idéal premier associé à $J(J')^iJ'' H^i_T(K)$, comme voulu.

\medskip\noindent
D'après les remarques initiales, $J''$ doit annuler
$$
(E^{-\delta(\mathfrak q) - i})_\mathfrak q =
(E^{-n - i})_\mathfrak q
$$
pour tout $\mathfrak q \in T_n \setminus T_{n - 1}$. Mais si $J''$ convient pour un $\mathfrak q$, il convient pour tous les $\mathfrak q$ d'un voisinage ouvert de $\mathfrak q$, puisque les modules $E^{-n - i}$ sont de type fini. Toute partie de $\Spec(A)$ étant noethérienne pour la topologie induite (Topologie, lemme \ref{topology-lemma-Noetherian}), il suffit donc de prouver l'existence de $J''$ pour un seul $\mathfrak q$.

\medskip\noindent
Les modules Ext étant de type fini, l'existence de $J''$ équivaut à
$$
\text{Supp}(E^{-n - i}) \cap \Spec(A_\mathfrak q) \subset T'.
$$
Cela revient à montrer que le localisé de $E^{-n - i}$ en tout $\mathfrak p \subset \mathfrak q$, $\mathfrak p \not \in T'$ est nul. Par dualité locale sur $A_\mathfrak p$, il faut donc montrer que
$$
H^{i + n - \delta(\mathfrak p)}_\mathfrak p(K_\mathfrak p) =
H^{i - \dim((A/\mathfrak p)_\mathfrak q)}_\mathfrak p(K_\mathfrak p)
$$
est nul (on utilise ici le fait que $\delta$ est une fonction de dimension). Cette annulation résulte de l'hypothèse du lemme, de $i \leq s$ et de notre définition de la profondeur, définition \ref{local-cohomology-definition-depth-complex}.

\medskip\noindent
Pour démontrer l'implication réciproque, supposons que la condition (2) ne soit pas vérifiée et reprenons les arguments précédents en sens inverse. Choisissons d'abord $\mathfrak q \in T$, $\mathfrak p \subset \mathfrak q$, avec $\mathfrak p \not \in T'$, tels que
$$
i = \text{depth}_{A_\mathfrak p}(K_\mathfrak p) +
\dim((A/\mathfrak p)_\mathfrak q) \leq s
$$
soit minimal. Alors $H^{i - \dim((A/\mathfrak p)_\mathfrak q)}_\mathfrak p(K_\mathfrak p)$ est non nul d'après notre définition de la profondeur, définition \ref{local-cohomology-definition-depth-complex}. Posons $n = \delta(\mathfrak q)$. Il n'existe alors aucun idéal $J \subset A$ tel que $V(J) \subset T'$ et $J(E^{-n - i})_\mathfrak q = 0$. Ainsi $H^i_\mathfrak q(K_\mathfrak q)$ n'est annulé par aucun idéal $J \subset A$ tel que $V(J) \subset T'$. Par minimalité de $i$, la suite spectrale ci-dessus montre que le module $H^i_T(K)_\mathfrak q$ n'est annulé par aucun idéal $J \subset A$ tel que $V(J) \subset T'$. Ainsi $H^i_T(K)$ n'est annulé par aucun idéal $J \subset A$ tel que $V(J) \subset T'$. Cela achève la démonstration de la proposition.
\end{proof}







\section{Finitude de la cohomologie locale, III}
\label{local-cohomology-section-finiteness-III}

\noindent
Nous étendons la discussion de la finitude de la cohomologie locale des sections \ref{local-cohomology-section-finiteness} et \ref{local-cohomology-section-finiteness-II} aux complexes bornés inférieurement dont les modules de cohomologie sont de type fini.

\begin{lemma}
\label{local-cohomology-lemma-check-finiteness-local-cohomology-by-annihilator-complex}
Soit $A$ un anneau noethérien. Soit $T \subset \Spec(A)$ une partie stable par spécialisation. Soit $K$ un objet de $D_{\textit{Coh}}^+(A)$. Soit $n \in \mathbf{Z}$. Les conditions suivantes sont équivalentes :
\begin{enumerate}
\item $H^i_T(K)$ est de type fini pour $i \leq n$,
\item il existe un idéal $J \subset A$ tel que $V(J) \subset T$ et tel que $J$ annule $H^i_T(K)$ pour $i \leq n$.
\end{enumerate}
Si $T = V(I) = Z$ pour un idéal $I \subset A$, ces conditions sont également équivalentes à
\begin{enumerate}
\item[(3)] il existe $e \geq 0$ tel que $I^e$ annule $H^i_Z(K)$ pour $i \leq n$.
\end{enumerate}
\end{lemma}

\begin{proof}
Ce lemme est la généralisation naturelle du lemme \ref{local-cohomology-lemma-check-finiteness-local-cohomology-by-annihilator}, dont le lecteur est invité à lire d'abord la démonstration. Supposons la condition (1) vérifiée. Rappelons que $H^i_J(K) = H^i_{V(J)}(K)$ ; voir Complexes dualisants, lemme \ref{dualizing-lemma-local-cohomology-noetherian}. Ainsi $H^i_T(K) = \colim H^i_J(K)$, où la limite inductive porte sur les idéaux $J \subset A$ tels que $V(J) \subset T$ ; voir le lemme \ref{local-cohomology-lemma-adjoint-ext}. Puisque $H^i_T(K)$ est de type fini pour $i \leq n$, on peut trouver $J \subset A$ comme dans (2) tel que $H^i_J(K) \to H^i_T(K)$ soit surjective pour $i \leq n$. Les éléments de la liste finie de générateurs sont donc de torsion par rapport aux puissances de $J$, et la condition (2) est vérifiée après remplacement de $J$ par une puissance convenable.

\medskip\noindent
Soit $a \in \mathbf{Z}$ un entier tel que $H^i(K) = 0$ pour $i < a$. Démontrons (2) $\Rightarrow$ (1) par récurrence descendante sur $a$. Si $a > n$, on a $H^i_T(K) = 0$ pour $i \leq n$ ; ainsi (1) et (2) sont tous deux vrais, et il n'y a rien à démontrer.

\medskip\noindent
Supposons donné $J$ comme dans (2). Observons que $N = H^a_T(K) = H^0_T(H^a(K))$ est de type fini, comme sous-module du $A$-module de type fini $H^a(K)$. Si $n = a$, la démonstration est terminée ; supposons donc désormais $a < n$. Par construction de $R\Gamma_T$, on a $H^i_T(N) = 0$ pour $i > 0$ et $H^0_T(N) = N$ ; voir la remarque \ref{local-cohomology-remark-upshot}. Choisissons un triangle distingué
$$
N[-a] \to K \to K' \to N[-a + 1]
$$
On voit alors que $H^a_T(K') = 0$ et $H^i_T(K) = H^i_T(K')$ pour $i > a$. On peut donc remplacer $K$ par $K'$. Ainsi on peut supposer $H^a_T(K) = 0$. Cela signifie qu'aucun des idéaux premiers associés à $H^a(K)$, qui sont en nombre fini, n'appartient à $T$. Par évitement des idéaux premiers (Algèbre, lemme \ref{algebra-lemma-silly}), on peut trouver $f \in J$ n'appartenant à aucun des idéaux premiers associés à $H^a(K)$. Choisissons un triangle distingué
$$
L \to K \xrightarrow{f} K \to L[1]
$$
Par construction, on a $H^i(L) = 0$ pour $i \leq a$. D'autre part, on a une suite exacte longue de cohomologie
$$
0 \to H^{a + 1}_T(L) \to H^{a + 1}_T(K) \xrightarrow{f}
H^{a + 1}_T(K) \to H^{a + 2}_T(L) \to H^{a + 2}_T(K) \xrightarrow{f} \ldots
$$
qui se décompose en l'identification $H^{a + 1}_T(L) = H^{a + 1}_T(K)$ et en suites exactes courtes
$$
0 \to H^{i - 1}_T(K) \to H^i_T(L) \to H^i_T(K) \to 0
$$
pour $i \leq n$, puisque $f \in J$. On en déduit que $J^2$ annule $H^i_T(L)$ pour $i \leq n$. Par l'hypothèse de récurrence appliquée à $L$, on voit que $H^i_T(L)$ est de type fini pour $i \leq n$. En utilisant une nouvelle fois la suite exacte courte, on voit que $H^i_T(K)$ est de type fini pour $i \leq n$, comme voulu.

\medskip\noindent
Nous omettons la démonstration de l'équivalence de (2) et (3) dans le cas $T = V(I)$.
\end{proof}

\begin{proposition}
\label{local-cohomology-proposition-finiteness-complex}
Soit $A$ un anneau noethérien possédant un complexe dualisant. Soit $T \subset \Spec(A)$ une partie stable par spécialisation. Soit $s \in \mathbf{Z}$. Soit $K \in D_{\textit{Coh}}^+(A)$. Les conditions suivantes sont équivalentes :
\begin{enumerate}
\item $H^i_T(K)$ est un $A$-module de type fini pour $i \leq s$,
\item pour tous $\mathfrak p \not \in T$, $\mathfrak q \in T$ tels que $\mathfrak p \subset \mathfrak q$, on a
$$
\text{depth}_{A_\mathfrak p}(K_\mathfrak p) +
\dim((A/\mathfrak p)_\mathfrak q) > s
$$
\end{enumerate}
\end{proposition}

\begin{proof}
Conséquence formelle de la proposition \ref{local-cohomology-proposition-annihilator-complex} et du lemme \ref{local-cohomology-lemma-check-finiteness-local-cohomology-by-annihilator-complex}.
\end{proof}






\section{Amélioration des modules cohérents}
\label{local-cohomology-section-improve}

\noindent
Des constructions analogues figurent dans \cite{EGA} et, plus récemment, dans \cite{Kollar-local-global-hulls} et \cite{Kollar-variants}.

\begin{lemma}
\label{local-cohomology-lemma-get-depth-1-along-Z}
Soit $X$ un schéma noethérien. Soit $T \subset X$ une partie stable par spécialisation. Soit $\mathcal{F}$ un $\mathcal{O}_X$-module cohérent. Alors il existe une unique application $\mathcal{F} \to \mathcal{F}'$ de $\mathcal{O}_X$-modules cohérents telle que
\begin{enumerate}
\item $\mathcal{F} \to \mathcal{F}'$ soit surjective,
\item $\mathcal{F}_x \to \mathcal{F}'_x$ soit un isomorphisme pour $x \not \in T$,
\item $\text{depth}_{\mathcal{O}_{X, x}}(\mathcal{F}'_x) \geq 1$ pour $x \in T$.
\end{enumerate}
Si $f : Y \to X$ est un morphisme plat, avec $Y$ noethérien, alors $f^*\mathcal{F} \to f^*\mathcal{F}'$ est le quotient correspondant pour $f^{-1}(T) \subset Y$ et $f^*\mathcal{F}$.
\end{lemma}

\begin{proof}
La condition (3) signifie simplement que $\text{Ass}(\mathcal{F}') \cap T = \emptyset$. Ainsi $\mathcal{F} \to \mathcal{F}'$ est le quotient de $\mathcal{F}$ par le sous-faisceau des sections dont le support est contenu dans $T$. Cela prouve l'unicité. L'assertion concernant les images réciproques résulte de Diviseurs, lemme \ref{divisors-lemma-bourbaki}, et de l'unicité.

\medskip\noindent
Existence de $\mathcal{F} \to \mathcal{F}'$. Par unicité, il suffit de prouver l'existence et l'unicité localement sur $X$ ; nous omettons ce petit détail. On peut donc supposer $X = \Spec(A)$ affine et $\mathcal{F}$ égal au module cohérent associé au $A$-module de type fini $M$. Posons $M' = M / H^0_T(M)$, avec $H^0_T(M)$ comme dans la section \ref{local-cohomology-section-supports}. Alors $M_\mathfrak p = M'_\mathfrak p$ pour $\mathfrak p \not \in T$, ce qui prouve (1). D'autre part, on a $H^0_T(M) = \colim H^0_Z(M)$, où $Z$ parcourt les parties fermées de $X$ contenues dans $T$. Ainsi, d'après Complexes dualisants, lemme \ref{dualizing-lemma-divide-by-torsion}, on a $H^0_T(M') = 0$, c'est-à-dire qu'aucun idéal premier associé à $M'$ n'appartient à $T$. Par conséquent, $\text{depth}(M'_\mathfrak p) \geq 1$ pour $\mathfrak p \in T$.
\end{proof}

\begin{lemma}
\label{local-cohomology-lemma-get-depth-2-along-Z}
Soit $j : U \to X$ une immersion ouverte de schémas noethériens. Soit $\mathcal{F}$ un $\mathcal{O}_X$-module cohérent. Supposons $\mathcal{F}' = j_*(\mathcal{F}|_U)$ cohérent. Alors $\mathcal{F} \to \mathcal{F}'$ est l'unique application de $\mathcal{O}_X$-modules cohérents telle que
\begin{enumerate}
\item $\mathcal{F}|_U \to \mathcal{F}'|_U$ soit un isomorphisme,
\item $\text{depth}_{\mathcal{O}_{X, x}}(\mathcal{F}'_x) \geq 2$ pour $x \in X$, $x \not \in U$.
\end{enumerate}
Si $f : Y \to X$ est un morphisme plat, avec $Y$ noethérien, alors $f^*\mathcal{F} \to f^*\mathcal{F}'$ est l'application correspondante pour $f^{-1}(U) \subset Y$.
\end{lemma}

\begin{proof}
On a $\text{depth}_{\mathcal{O}_{X, x}}(\mathcal{F}'_x) \geq 2$ d'après Diviseurs, lemme \ref{divisors-lemma-depth-pushforward}, point (3). L'unicité de $\mathcal{F} \to \mathcal{F}'$ résulte de Diviseurs, lemme \ref{divisors-lemma-depth-2-hartog}. La compatibilité aux images réciproques plates résulte du changement de base plat ; voir Cohomologie des schémas, lemme \ref{coherent-lemma-flat-base-change-cohomology}.
\end{proof}

\begin{lemma}
\label{local-cohomology-lemma-make-S2-along-Z}
Soit $X$ un schéma noethérien. Soit $Z \subset X$ un sous-schéma fermé. Soit $\mathcal{F}$ un $\mathcal{O}_X$-module cohérent. Supposons que $X$ soit universellement caténaire et que les fibres formelles de ses anneaux locaux vérifient $(S_1)$. Alors il existe une unique application $\mathcal{F} \to \mathcal{F}''$ de $\mathcal{O}_X$-modules cohérents telle que
\begin{enumerate}
\item $\mathcal{F}_x \to \mathcal{F}''_x$ soit un isomorphisme pour $x \in X \setminus Z$,
\item $\mathcal{F}_x \to \mathcal{F}''_x$ soit surjective et $\text{depth}_{\mathcal{O}_{X, x}}(\mathcal{F}''_x) = 1$ pour tout $x \in Z$ tel qu'il existe une spécialisation immédiate $x' \leadsto x$ avec $x' \not \in Z$ et $x' \in \text{Ass}(\mathcal{F})$,
\item $\text{depth}_{\mathcal{O}_{X, x}}(\mathcal{F}''_x) \geq 2$ pour les autres $x \in Z$.
\end{enumerate}
Si $f : Y \to X$ est un morphisme de Cohen-Macaulay, avec $Y$ noethérien, alors $f^*\mathcal{F} \to f^*\mathcal{F}''$ vérifie les mêmes propriétés relativement à $f^{-1}(Z) \subset Y$.
\end{lemma}

\begin{proof}
Soit $\mathcal{F} \to \mathcal{F}'$ l'application construite au lemme \ref{local-cohomology-lemma-get-depth-1-along-Z} pour la partie $Z$ de $X$. Rappelons que $\mathcal{F}'$ est le quotient de $\mathcal{F}$ par le sous-faisceau des sections à support dans $Z$.

\medskip\noindent
Démontrons d'abord l'unicité. Soit $\mathcal{F} \to \mathcal{F}''$ comme dans le lemme. On obtient une factorisation $\mathcal{F} \to \mathcal{F}' \to \mathcal{F}''$, puisque $\text{Ass}(\mathcal{F}'') \cap Z = \emptyset$ d'après (2) et (3). Soit $U \subset X$ le plus grand sous-schéma ouvert tel que $\mathcal{F}'|_U \to \mathcal{F}''|_U$ soit un isomorphisme. On voit que $U$ contient tous les points de (2). Alors Diviseurs, lemme \ref{divisors-lemma-depth-2-hartog}, permet de conclure que $\mathcal{F}'' = j_*(\mathcal{F}'|_U)$. On obtient ainsi l'unicité (petit détail : si l'on dispose de deux tels $\mathcal{F}''$, on prend l'intersection des ouverts $U$ obtenus pour chacun).

\medskip\noindent
Démonstration de l'existence. Rappelons que $\text{Ass}(\mathcal{F}') = \{x_1, \ldots, x_n\}$ est fini et que $x_i \not \in Z$. Soit $Y_i$ l'adhérence de $\{x_i\}$. Soient $Z_{i, j}$ les composantes irréductibles de $Z \cap Y_i$. Observons que $\text{Supp}(\mathcal{F}') \cap Z = \bigcup Z_{i, j}$. Soit $z_{i, j} \in Z_{i, j}$ le point générique. Posons
$$
d_{i, j} = \dim(\mathcal{O}_{\overline{\{x_i\}}, z_{i, j}})
$$
Si $d_{i, j} = 1$, alors $z_{i, j}$ est l'un des points de (2). Il n'est donc pas nécessaire de modifier $\mathcal{F}'$ en ces points. De plus, en supposant toujours $d_{i, j} = 1$, le lemme \ref{local-cohomology-lemma-depth-function} permet de trouver un voisinage ouvert $z_{i, j} \in V_{i, j} \subset X$ tel que $\text{depth}_{\mathcal{O}_{X, z}}(\mathcal{F}'_z) \geq 2$ pour $z \in Z_{i, j} \cap V_{i, j}$, $z \not = z_{i, j}$. Posons
$$
Z' = X \setminus
\left(
X \setminus Z \cup \bigcup\nolimits_{d_{i, j} = 1} V_{i, j})
\right)
$$
Notons $j' : X \setminus Z' \to X$. Le choix de $Z'$ assure que les hypothèses du lemme \ref{local-cohomology-lemma-finiteness-pushforward-general} sont satisfaites. On conclut en posant $\mathcal{F}'' = j'_*(\mathcal{F}'|_{X \setminus Z'})$ et en appliquant le lemme \ref{local-cohomology-lemma-get-depth-2-along-Z}.

\medskip\noindent
La dernière assertion résulte de la formule de variation de la profondeur par un homomorphisme local plat ; voir Algèbre, lemme \ref{algebra-lemma-apply-grothendieck-module}, et de l'hypothèse sur les fibres de $f$ contenue dans le fait que $f$ est de Cohen-Macaulay. Détails omis.
\end{proof}

\begin{lemma}
\label{local-cohomology-lemma-make-S2-along-T-simple}
Soit $X$ un schéma noethérien possédant localement un complexe dualisant. Soit $T' \subset X$ une partie stable par spécialisation. Soit $\mathcal{F}$ un $\mathcal{O}_X$-module cohérent. Supposons que, si $x \leadsto x'$ est une spécialisation immédiate de points de $X$ avec $x' \in T'$ et $x \not \in T'$, alors $\text{depth}(\mathcal{F}_x) \geq 1$. Il existe alors une unique application $\mathcal{F} \to \mathcal{F}''$ de $\mathcal{O}_X$-modules cohérents telle que
\begin{enumerate}
\item $\mathcal{F}_x \to \mathcal{F}''_x$ soit un isomorphisme pour $x \not \in T'$,
\item $\text{depth}_{\mathcal{O}_{X, x}}(\mathcal{F}''_x) \geq 2$ pour $x \in T'$.
\end{enumerate}
Si $f : Y \to X$ est un morphisme de Cohen-Macaulay, avec $Y$ noethérien, alors $f^*\mathcal{F} \to f^*\mathcal{F}''$ vérifie les mêmes propriétés relativement à $f^{-1}(T') \subset Y$.
\end{lemma}

\begin{proof}
Démonstration de l'unicité. Soit $\mathcal{F} \to \mathcal{F}''$ comme dans le lemme. Soit $U \subset X$ le plus grand sous-schéma ouvert tel que $\mathcal{F}|_U \to \mathcal{F}''|_U$ soit un isomorphisme. D'après (1), $U$ contient l'ensemble $X \setminus T'$. D'après (2) et Diviseurs, lemme \ref{divisors-lemma-depth-2-hartog}, on conclut que $\mathcal{F}'' = j_*(\mathcal{F}|_U)$, où $j : U \to X$ est le morphisme d'inclusion. On obtient ainsi l'unicité (petit détail : si l'on dispose de deux tels $\mathcal{F}''$, on prend l'intersection des ouverts $U$ obtenus pour chacun).

\medskip\noindent
Démonstration de l'existence. Soit $\mathcal{F} \to \mathcal{F}'$ le quotient de $\mathcal{F}$ construit au lemme \ref{local-cohomology-lemma-get-depth-1-along-Z} à l'aide de $T'$. Rappelons que $\mathcal{F}'$ est le quotient de $\mathcal{F}$ par le sous-faisceau des sections à support dans $T'$. Définissons
$$
\mathcal{F}'' = \colim j_*(\mathcal{F}'|_V)
$$
où $j : V \to X$ parcourt les sous-schémas ouverts tels que $X \setminus V \subset T'$. Observons que le système inductif est filtrant, puisque $T'$ est stable par spécialisation. Chacune des applications $\mathcal{F}' \to j_*(\mathcal{F}'|_V)$ est injective, puisque $\text{Ass}(\mathcal{F}')$ est disjoint de $T'$. Ainsi $\mathcal{F}' \to \mathcal{F}''$ est injective.

\medskip\noindent
Supposons $X = \Spec(A)$ affine et $\mathcal{F}$ correspondant au $A$-module de type fini $M$. Alors $\mathcal{F}'$ correspond à $M' = M / H^0_{T'}(M)$ ; voir la démonstration du lemme \ref{local-cohomology-lemma-get-depth-1-along-Z}. En appliquant les lemmes \ref{local-cohomology-lemma-local-cohomology} et \ref{local-cohomology-lemma-adjoint-ext}, on voit que $\mathcal{F}''$ correspond à un $A$-module $M''$ qui s'insère dans la suite exacte courte
$$
0 \to M' \to M'' \to H^1_{T'}(M') \to 0
$$
D'après la proposition \ref{local-cohomology-proposition-finiteness} et la condition de l'énoncé sur les spécialisations immédiates, $M''$ est un $A$-module de type fini. On voit ainsi que $\mathcal{F}''$ est cohérent.

\medskip\noindent
La cohérence de $\mathcal{F}''$ et l'injectivité des applications de transition dans la limite inductive filtrante ci-dessus impliquent que $\mathcal{F}'' = j_*(\mathcal{F}'|_V)$ pour tout $V$ assez petit comme ci-dessus. Les conditions (1) et (2) résultent alors du lemme \ref{local-cohomology-lemma-get-depth-2-along-Z}.

\medskip\noindent
La dernière assertion résulte de la formule de variation de la profondeur par un homomorphisme local plat ; voir Algèbre, lemme \ref{algebra-lemma-apply-grothendieck-module}, et de l'hypothèse sur les fibres de $f$ contenue dans le fait que $f$ est de Cohen-Macaulay. Détails omis.
\end{proof}

\begin{lemma}
\label{local-cohomology-lemma-make-S2-along-T}
Soit $X$ un schéma noethérien possédant localement un complexe dualisant. Soient $T' \subset T \subset X$ des parties stables par spécialisation telles que, si $x \leadsto x'$ est une spécialisation immédiate de points de $X$ et $x' \in T'$, alors $x \in T$. Soit $\mathcal{F}$ un $\mathcal{O}_X$-module cohérent. Il existe alors une unique application $\mathcal{F} \to \mathcal{F}''$ de $\mathcal{O}_X$-modules cohérents telle que
\begin{enumerate}
\item $\mathcal{F}_x \to \mathcal{F}''_x$ soit un isomorphisme pour $x \not \in T$,
\item $\mathcal{F}_x \to \mathcal{F}''_x$ soit surjective et $\text{depth}_{\mathcal{O}_{X, x}}(\mathcal{F}''_x) \geq 1$ pour $x \in T$, $x \not \in T'$,
\item $\text{depth}_{\mathcal{O}_{X, x}}(\mathcal{F}''_x) \geq 2$ pour $x \in T'$.
\end{enumerate}
Si $f : Y \to X$ est un morphisme de Cohen-Macaulay, avec $Y$ noethérien, alors $f^*\mathcal{F} \to f^*\mathcal{F}''$ vérifie les mêmes propriétés relativement à $f^{-1}(T') \subset f^{-1}(T) \subset Y$.
\end{lemma}

\begin{proof}
Soit d'abord $\mathcal{F} \to \mathcal{F}'$ le quotient de $\mathcal{F}$ construit au lemme \ref{local-cohomology-lemma-get-depth-1-along-Z} à l'aide de $T$. Soit ensuite $\mathcal{F}' \to \mathcal{F}''$ l'unique application de modules cohérents construite au lemme \ref{local-cohomology-lemma-make-S2-along-T-simple} à l'aide de $T'$. Alors $\mathcal{F} \to \mathcal{F}''$ possède les propriétés voulues.
\end{proof}










\section{Annulation de Hartshorne-Lichtenbaum}
\label{local-cohomology-section-Hartshorne-Lichtenbaum-vanishing}

\noindent
Ce résultat d'annulation est l'analogue local du théorème de Lichtenbaum présenté dans Dualité pour les schémas, section \ref{duality-section-lichtenbaum}. On le trouvera, ainsi que bien d'autres résultats, dans \cite{CD}.

\begin{lemma}
\label{local-cohomology-lemma-cd-top-vanishing}
Soit $A$ un anneau noethérien de dimension $d$. Soient $I \subset I' \subset A$ des idéaux. Si $I'$ est contenu dans le radical de Jacobson de $A$ et si $\text{cd}(A, I') < d$, alors $\text{cd}(A, I) < d$.
\end{lemma}

\begin{proof}
D'après le lemme \ref{local-cohomology-lemma-cd-dimension}, on sait que $\text{cd}(A, I) \leq d$. Nous utiliserons le lemme \ref{local-cohomology-lemma-isomorphism} pour montrer que
$$
H^d_{V(I')}(A) \to H^d_{V(I)}(A)
$$
est surjective, ce qui achèvera la démonstration. Choisissons $\mathfrak p \in V(I) \setminus V(I')$. L'hypothèse sur $I'$ montre que $\mathfrak p$ n'est pas un idéal maximal de $A$. Ainsi $\dim(A_\mathfrak p) < d$. Alors $H^d_{\mathfrak pA_\mathfrak p}(A_\mathfrak p) = 0$ d'après le lemme \ref{local-cohomology-lemma-cd-dimension}.
\end{proof}

\begin{lemma}
\label{local-cohomology-lemma-cd-top-vanishing-some-module}
Soit $A$ un anneau noethérien de dimension $d$. Soit $I \subset A$ un idéal. Si $H^d_{V(I)}(M) = 0$ pour un $A$-module de type fini dont le support contient toutes les composantes irréductibles de dimension $d$, alors $\text{cd}(A, I) < d$.
\end{lemma}

\begin{proof}
D'après le lemme \ref{local-cohomology-lemma-cd-dimension}, on sait que $\text{cd}(A, I) \leq d$. Ainsi, pour tout $A$-module de type fini $N$, on a $H^i_{V(I)}(N) = 0$ pour $i > d$. Disons que la propriété $\mathcal{P}$ est vérifiée pour le $A$-module de type fini $N$ si $H^d_{V(I)}(N) = 0$. L'une des hypothèses est que $\mathcal{P}(M)$ est vérifiée. Observons que $\mathcal{P}(N_1 \oplus N_2)
\Leftrightarrow (\mathcal{P}(N_1) \wedge \mathcal{P}(N_2))$. Si $N \to N'$ est surjective, alors $\mathcal{P}(N) \Rightarrow \mathcal{P}(N')$, puisque $H^{d + 1}_{V(I)}$ est nul (voir ci-dessus). Soient $\mathfrak p_1, \ldots, \mathfrak p_n$ les idéaux premiers minimaux de $A$ tels que $\dim(A/\mathfrak p_i) = d$. Observons que $\mathcal{P}(N)$ est vérifiée si le support de $N$ est disjoint de $\{\mathfrak p_1, \ldots, \mathfrak p_n\}$, pour des raisons de dimension ; voir le lemme \ref{local-cohomology-lemma-cd-dimension}. Pour chaque $i$, posons $M_i = M/\mathfrak p_i M$. C'est un $A$-module de type fini annulé par $\mathfrak p_i$, dont le support est égal à $V(\mathfrak p_i)$ (on utilise ici l'hypothèse sur le support de $M$). Enfin, si $J \subset A$ est un idéal, alors $\mathcal{P}(JM_i)$ est vérifiée, car $JM_i$ est un quotient d'une somme directe de copies de $M$. Il résulte donc de Cohomologie des schémas, lemme \ref{coherent-lemma-property-higher-rank-cohomological}, que $\mathcal{P}$ est vérifiée pour tout $A$-module de type fini.
\end{proof}

\begin{lemma}
\label{local-cohomology-lemma-top-coh-divisible}
Soit $A$ un anneau local noethérien de dimension $d$. Soit $f \in A$ un élément n'appartenant à aucun idéal premier minimal de dimension $d$. Alors $f : H^d_{V(I)}(M) \to H^d_{V(I)}(M)$ est surjective pour tout $A$-module de type fini $M$ et tout idéal $I \subset A$.
\end{lemma}

\begin{proof}
Le support de $M/fM$ est de dimension $< d$, par l'hypothèse sur $f$. Ainsi $H^d_{V(I)}(M/fM) = 0$ d'après le lemme \ref{local-cohomology-lemma-cd-dimension}. Par conséquent, $H^d_{V(I)}(fM) \to H^d_{V(I)}(M)$ est surjective. Puisque le lemme \ref{local-cohomology-lemma-cd-dimension} donne $\text{cd}(A, I) \leq d$, on voit aussi que la surjection $M \to fM$, $x \mapsto fx$ induit une surjection $H^d_{V(I)}(M) \to H^d_{V(I)}(fM)$.
\end{proof}

\begin{lemma}
\label{local-cohomology-lemma-cd-bound-dualizing}
Soit $A$ un anneau local noethérien muni d'un complexe dualisant normalisé $\omega_A^\bullet$. Soit $I \subset A$ un idéal. Si $H^0_{V(I)}(\omega_A^\bullet) = 0$, alors $\text{cd}(A, I) < \dim(A)$.
\end{lemma}

\begin{proof}
Posons $d = \dim(A)$. Soient $\mathfrak p_1, \ldots, \mathfrak p_n \subset A$ les idéaux premiers minimaux de dimension $d$. Rappelons que le $A$-module de type fini $H^{-i}(\omega_A^\bullet)$ est non nul seulement pour $i \in \{0, \ldots, d\}$ et que le support de $H^{-i}(\omega_A^\bullet)$ est de dimension $\leq i$ ; voir le lemme \ref{local-cohomology-lemma-sitting-in-degrees}. Posons $\omega_A = H^{-d}(\omega_A^\bullet)$. Par évitement des idéaux premiers (Algèbre, lemme \ref{algebra-lemma-silly}), on peut trouver $f \in A$, $f \not \in \mathfrak p_i$, qui annule $H^{-i}(\omega_A^\bullet)$ pour $i < d$. Considérons le triangle distingué
$$
\omega_A[d] \to \omega_A^\bullet \to
\tau_{\geq -d + 1}\omega_A^\bullet \to \omega_A[d + 1]
$$
Voir Catégories dérivées, remarque \ref{derived-remark-truncation-distinguished-triangle}. D'après Catégories dérivées, lemme \ref{derived-lemma-trick-vanishing-composition}, $f^d$ induit l'endomorphisme nul de $\tau_{\geq -d + 1}\omega_A^\bullet$. Les axiomes des catégories triangulées fournissent une application
$$
\omega_A^\bullet \to \omega_A[d]
$$
dont le composé avec $\omega_A[d] \to \omega_A^\bullet$ est la multiplication par $f^d$ sur $\omega_A[d]$. On en déduit que $f^d$ annule $H^d_{V(I)}(\omega_A)$. Le lemme \ref{local-cohomology-lemma-top-coh-divisible} donne alors $H^d_{V(I)}(\omega_A) = 0$. On conclut par le lemme \ref{local-cohomology-lemma-cd-top-vanishing-some-module} et le fait que $(\omega_A)_{\mathfrak p_i}$ est non nul (voir par exemple Complexes dualisants, lemme \ref{dualizing-lemma-nonvanishing-generically-local}).
\end{proof}

\begin{lemma}
\label{local-cohomology-lemma-inverse-system-symbolic-powers}
Soit $(A, \mathfrak m)$ un anneau local noethérien complet intègre. Soit $\mathfrak p \subset A$ un idéal premier de dimension $1$. Pour tout $n \geq 1$, il existe $m \geq n$ tel que $\mathfrak p^{(m)} \subset \mathfrak p^n$.
\end{lemma}

\begin{proof}
Rappelons que la puissance symbolique $\mathfrak p^{(m)}$ est définie comme le noyau de $A \to A_\mathfrak p/\mathfrak p^mA_\mathfrak p$. Puisque la localisation est exacte, on en déduit que, dans la suite exacte courte
$$
0 \to \mathfrak a_n \to A/\mathfrak p^n \to A/\mathfrak p^{(n)} \to 0
$$
le support de $\mathfrak a_n$ est contenu dans $\{\mathfrak m\}$. En particulier, le système projectif $(\mathfrak a_n)$ vérifie la condition de Mittag-Leffler, puisque chaque $\mathfrak a_n$ est un $A$-module artinien. Le lemme équivaut donc à la condition $\lim \mathfrak a_n = 0$. Soit $f \in \lim \mathfrak a_n$. Alors $f$ est un élément de $A = \lim A/\mathfrak p^n$ (on utilise ici la complétude de $A$) dont l'image dans le complété $A_\mathfrak p^\wedge$ de $A_\mathfrak p$ est nulle. Puisque $A_\mathfrak p \to A_\mathfrak p^\wedge$ est fidèlement plat, on voit que $f$ a une image nulle dans $A_\mathfrak p$. Puisque $A$ est intègre, $f$ est nul, comme voulu.
\end{proof}

\begin{proposition}
\label{local-cohomology-proposition-Hartshorne-Lichtenbaum-vanishing}
\begin{reference}
\cite[Theorem 3.1]{CD}
\end{reference}
Soit $A$ un anneau local noethérien de complété $A^\wedge$. Soit $I \subset A$ un idéal tel que
$$
\dim V(IA^\wedge + \mathfrak p) \geq 1
$$
pour tout idéal premier minimal $\mathfrak p \subset A^\wedge$ de dimension $\dim(A)$. Alors $\text{cd}(A, I) < \dim(A)$.
\end{proposition}

\begin{proof}
Puisque $A \to A^\wedge$ est fidèlement plat, on a $H^d_{V(I)}(A) \otimes_A A^\wedge = H^d_{V(IA^\wedge)}(A^\wedge)$ d'après Complexes dualisants, lemme \ref{dualizing-lemma-torsion-change-rings}. On peut donc supposer $A$ complet.

\medskip\noindent
Supposons $A$ complet. Soient $\mathfrak p_1, \ldots, \mathfrak p_n \subset A$ les idéaux premiers minimaux de dimension $d$. Considérons l'anneau local complet $A_i = A/\mathfrak p_i$. On a $H^d_{V(I)}(A_i) = H^d_{V(IA_i)}(A_i)$ d'après Complexes dualisants, lemme \ref{dualizing-lemma-local-cohomology-and-restriction}. D'après le lemme \ref{local-cohomology-lemma-cd-top-vanishing-some-module}, il suffit de prouver le lemme pour $(A_i, IA_i)$. On peut donc supposer que $A$ est un anneau local complet intègre.

\medskip\noindent
Supposons que $A$ soit un anneau local complet intègre. On peut choisir un idéal premier $\mathfrak p \supset I$ tel que $\dim(A/\mathfrak p) = 1$. D'après le lemme \ref{local-cohomology-lemma-cd-top-vanishing}, il suffit de prouver le lemme pour $\mathfrak p$.

\medskip\noindent
D'après le lemme \ref{local-cohomology-lemma-cd-bound-dualizing}, il suffit de montrer que $H^0_{V(\mathfrak p)}(\omega_A^\bullet) = 0$. Rappelons que
$$
H^0_{V(\mathfrak p)}(\omega_A^\bullet) =
\colim \text{Ext}^0_A(A/\mathfrak p^n, \omega_A^\bullet)
$$
D'après le lemme \ref{local-cohomology-lemma-inverse-system-symbolic-powers}, cette limite inductive est égale à
$$
\colim \text{Ext}^0_A(A/\mathfrak p^{(n)}, \omega_A^\bullet)
$$
Puisque $\text{depth}(A/\mathfrak p^{(n)}) = 1$, ces groupes Ext sont nuls d'après le lemme \ref{local-cohomology-lemma-sitting-in-degrees}, comme voulu.
\end{proof}

\begin{lemma}
\label{local-cohomology-lemma-affine-complement}
Soit $(A, \mathfrak m)$ un anneau local noethérien. Soit $I \subset A$ un idéal. Supposons $A$ excellent, normal, et $\dim V(I) \geq 1$. Alors $\text{cd}(A, I) < \dim(A)$. En particulier, si $\dim(A) = 2$, alors $\Spec(A) \setminus V(I)$ est affine.
\end{lemma}

\begin{proof}
D'après Compléments d'algèbre, lemme \ref{more-algebra-lemma-completion-normal-local-ring}, le complété $A^\wedge$ est normal, donc intègre. L'hypothèse de la proposition \ref{local-cohomology-proposition-Hartshorne-Lichtenbaum-vanishing} est donc vérifiée, et l'on conclut. L'assertion d'affinité résulte du lemme \ref{local-cohomology-lemma-cd-is-one}.
\end{proof}
















\section{Action de Frobenius}
\label{local-cohomology-section-frobenius}

\noindent
Soit $p$ un nombre premier. Soit $A$ un anneau tel que $p = 0$ dans $A$. L'{\it endomorphisme de Frobenius} de $A$ est l'application
$$
F : A \longrightarrow A,
\quad
a \longmapsto a^p
$$
Dans cette section, nous démontrons des lemmes sur les modules munis d'une action de Frobenius.

\begin{lemma}
\label{local-cohomology-lemma-annihilator-frobenius-module}
Soit $p$ un nombre premier. Soit $(A, \mathfrak m, \kappa)$ un anneau local noethérien tel que $p = 0$ dans $A$. Soit $M$ un $A$-module de type fini tel que $M \otimes_{A, F} A \cong M$. Alors $M$ est libre de type fini.
\end{lemma}

\begin{proof}
Choisissons une présentation $A^{\oplus m} \to A^{\oplus n} \to M$ qui induit un isomorphisme $\kappa^{\oplus n} \to M/\mathfrak m M$. Soit $T = (a_{ij})$ la matrice de l'application $A^{\oplus m} \to A^{\oplus n}$. Observons que $a_{ij} \in \mathfrak m$. En effectuant le changement de base par $F$ et en utilisant son exactitude à droite, on obtient une présentation $A^{\oplus m} \to A^{\oplus n} \to M$ dont la matrice est $T = (a_{ij}^p)$. On a donc une présentation avec $a_{ij} \in \mathfrak m^p$. En répétant cette construction, on trouve, pour chaque $e \geq 1$, une présentation avec $a_{ij} \in \mathfrak m^e$. Cela implique que les idéaux de Fitting (Compléments d'algèbre, définition \ref{more-algebra-definition-fitting-ideal}) $\text{Fit}_k(M)$ pour $k < n$ sont contenus dans $\bigcap_{e \geq 1} \mathfrak m^e$. Ce dernier étant nul par le théorème d'intersection de Krull (Algèbre, lemme \ref{algebra-lemma-intersect-powers-ideal-module-zero}), on conclut que $M$ est libre de rang $n$, d'après Compléments d'algèbre, lemme \ref{more-algebra-lemma-fitting-ideal-finite-locally-free}.
\end{proof}

\noindent
Dans cette section, nous dirons que des éléments $f_1, \ldots, f_r$ d'un anneau $A$ sont {\it indépendants} si $\sum a_if_i = 0$ implique $a_i \in (f_1, \ldots, f_r)$. Autrement dit, en posant $I = (f_1, \ldots, f_r)$, le module $I/I^2$ est libre sur $A/I$, de base $f_1, \ldots, f_r$.

\begin{lemma}
\label{local-cohomology-lemma-1}
\begin{reference}
Voir \cite{Lech-inequalities} et \cite[Lemma 1 page 299]{MatCA}.
\end{reference}
Soit $A$ un anneau. Si $f_1, \ldots, f_{r - 1}, f_rg_r$ sont indépendants, alors $f_1, \ldots, f_r$ sont indépendants.
\end{lemma}

\begin{proof}
Supposons $\sum a_if_i = 0$. Alors $\sum a_ig_rf_i = 0$. Ainsi $a_r \in (f_1, \ldots, f_{r - 1}, f_rg_r)$. Écrivons $a_r = \sum_{i < r} b_i f_i + b f_rg_r$. Alors $0 = \sum_{i < r} (a_i + b_if_r)f_i + bf_r^2g_r$. Par conséquent, $a_i + b_i f_r \in (f_1, \ldots, f_{r - 1}, f_rg_r)$, ce qui implique $a_i \in (f_1, \ldots, f_r)$, comme voulu.

\end{proof}

\begin{lemma}
\label{local-cohomology-lemma-2}
\begin{reference}
Voir \cite{Lech-inequalities} et \cite[Lemma 2 page 300]{MatCA}.
\end{reference}
Soit $A$ un anneau. Si $f_1, \ldots, f_{r - 1}, f_rg_r$ sont indépendants et si le $A$-module $A/(f_1, \ldots, f_{r - 1}, f_rg_r)$ est de longueur finie, alors
\begin{align*}
& \text{length}_A(A/(f_1, \ldots, f_{r - 1}, f_rg_r)) \\
& =
\text{length}_A(A/(f_1, \ldots, f_{r - 1}, f_r)) +
\text{length}_A(A/(f_1, \ldots, f_{r - 1}, g_r))
\end{align*}
\end{lemma}

\begin{proof}
Montrons qu'il existe une suite exacte
$$
0 \to
A/(f_1, \ldots, f_{r - 1}, g_r) \xrightarrow{f_r}
A/(f_1, \ldots, f_{r - 1}, f_rg_r) \to
A/(f_1, \ldots, f_{r - 1}, f_r) \to 0
$$
En effet, si $a f_r \in (f_1, \ldots, f_{r - 1}, f_rg_r)$, alors $\sum_{i < r} a_i f_i + (a + bg_r)f_r = 0$ pour certains $b, a_i \in A$. Ainsi $\sum_{i < r} a_i g_r f_i + (a + bg_r)g_rf_r = 0$, ce qui implique $a + bg_r \in (f_1, \ldots, f_{r - 1}, f_rg_r)$ et signifie que $a$ a une image nulle dans $A/(f_1, \ldots, f_{r - 1}, g_r)$. Cela prouve l'affirmation. Pour conclure, utiliser l'additivité des longueurs (Algèbre, lemme \ref{algebra-lemma-length-additive}).
\end{proof}

\begin{lemma}
\label{local-cohomology-lemma-3}
\begin{reference}
Voir \cite{Lech-inequalities} et \cite[Lemma 3 page 300]{MatCA}.
\end{reference}
Soit $(A, \mathfrak m)$ un anneau local. Si $\mathfrak m = (x_1, \ldots, x_r)$ et si $x_1^{e_1}, \ldots, x_r^{e_r}$ sont indépendants pour certains $e_i > 0$, alors $\text{length}_A(A/(x_1^{e_1}, \ldots, x_r^{e_r})) = e_1\ldots e_r$.
\end{lemma}

\begin{proof}
Utiliser les lemmes \ref{local-cohomology-lemma-1} et \ref{local-cohomology-lemma-2} et une récurrence.
\end{proof}

\begin{lemma}
\label{local-cohomology-lemma-flat-extension-independent}
Soit $\varphi : A \to B$ un homomorphisme plat d'anneaux. Si $f_1, \ldots, f_r \in A$ sont indépendants, alors $\varphi(f_1), \ldots, \varphi(f_r) \in B$ sont indépendants.
\end{lemma}

\begin{proof}
Posons $I = (f_1, \ldots, f_r)$ et $J = \varphi(I)B$. Par platitude, on a $I/I^2 \otimes_A B = J/J^2$. La liberté de $I/I^2$ sur $A/I$ implique donc celle de $J/J^2$ sur $B/J$.
\end{proof}

\begin{lemma}[Kunz]
\label{local-cohomology-lemma-frobenius-flat-regular}
\begin{reference}
\cite{Kunz-flat}
\end{reference}
Soit $p$ un nombre premier. Soit $A$ un anneau noethérien tel que $p = 0$. Les conditions suivantes sont équivalentes :
\begin{enumerate}
\item $A$ est régulier,
\item $F : A \to A$, $a \mapsto a^p$ est plat.
\end{enumerate}
\end{lemma}

\begin{proof}
Observons que $\Spec(F) : \Spec(A) \to \Spec(A)$ est l'application identique. La régularité se définit à l'aide des anneaux locaux et la platitude se vérifie sur les anneaux locaux ; voir Algèbre, lemme \ref{algebra-lemma-flat-localization}. On peut donc supposer, et l'on suppose, que $A$ est un anneau local noethérien d'idéal maximal $\mathfrak m$.

\medskip\noindent
Supposons $A$ régulier. Soit $x_1, \ldots, x_d$ un système de paramètres de $A$. En appliquant $F$, on obtient $F(x_1), \ldots, F(x_d) = x_1^p, \ldots, x_d^p$, qui est un système de paramètres de $A$. Ainsi $F$ est plat ; voir Algèbre, lemmes \ref{algebra-lemma-CM-over-regular-flat} et \ref{algebra-lemma-regular-ring-CM}.

\medskip\noindent
Réciproquement, supposons $F$ plat. Écrivons $\mathfrak m = (x_1, \ldots, x_r)$ avec $r$ minimal. Alors $x_1, \ldots, x_r$ sont indépendants au sens défini ci-dessus. Puisque $F$ est plat, $x_1^p, \ldots, x_r^p$ sont indépendants ; voir le lemme \ref{local-cohomology-lemma-flat-extension-independent}. Ainsi $\text{length}_A(A/(x_1^p, \ldots, x_r^p)) = p^r$ d'après le lemme \ref{local-cohomology-lemma-3}. Posons $\chi(n) = \text{length}_A(A/\mathfrak m^n)$ et rappelons qu'il s'agit d'un polynôme numérique de degré $\dim(A)$ ; voir Algèbre, proposition \ref{algebra-proposition-dimension}. Choisissons $n \gg 0$. Observons que
$$
\mathfrak m^{pn + pr} \subset F(\mathfrak m^n)A \subset \mathfrak m^{pn}
$$
comme on le voit en considérant les monômes en $x_1, \ldots, x_r$. On a
$$
A/F(\mathfrak m^n)A = A/\mathfrak m^n \otimes_{A, F} A
$$
Par platitude de $F$, ce module est de longueur $\chi(n) \text{length}_A(A/F(\mathfrak m)A)$ (Algèbre, lemme \ref{algebra-lemma-pullback-module}), qui vaut $p^r\chi(n)$ d'après ce qui précède. On en déduit
$$
\chi(pn + pr) \geq p^r\chi(n) \geq \chi(pn)
$$
La comparaison des termes dominants donne $r = \dim(A)$, c'est-à-dire que $A$ est régulier.
\end{proof}




\section{Structure de certains modules}
\label{local-cohomology-section-structure}

\noindent
Quelques résultats sur la structure de certains types de modules sur les anneaux locaux réguliers. On trouvera de tels résultats et bien davantage dans \cite{Huneke-Sharp}, \cite{Lyubeznik}, \cite{Lyubeznik2}.

\begin{lemma}
\label{local-cohomology-lemma-structure-torsion-D-module-regular}
\begin{reference}
Cas particulier de \cite[Theorem 2.4]{Lyubeznik}
\end{reference}
Soit $k$ un corps de caractéristique $0$. Soit $d \geq 1$. Soit $A = k[[x_1, \ldots, x_d]]$, d'idéal maximal $\mathfrak m$. Soit $M$ un module de torsion par rapport aux puissances de $\mathfrak m$ sur l'anneau $A$, muni d'opérateurs additifs $D_1, \ldots, D_d$ vérifiant la règle de Leibniz

$$
D_i(fz) = \partial_i(f) z + f D_i(z)
$$
pour $f \in A$ et $z \in M$. Ici $\partial_i$ désigne la dérivation par rapport à $x_i$. Alors $M$ est isomorphe à une somme directe de copies de l'enveloppe injective $E$ de $k$.
\end{lemma}

\begin{proof}
Choisissons un ensemble $J$ et un isomorphisme $M[\mathfrak m] \to \bigoplus_{j \in J} k$. Puisque $\bigoplus_{j \in J} E$ est injectif (Complexes dualisants, lemme \ref{dualizing-lemma-sum-injective-modules}), on peut prolonger cet isomorphisme en un homomorphisme de $A$-modules $\varphi : M \to \bigoplus_{j \in J} E$. Montrons que $\varphi$ est un isomorphisme, c'est-à-dire une bijection.

\medskip\noindent
Injectivité. Soit $z \in M$ non nul. Puisque $M$ est de torsion par rapport aux puissances de $\mathfrak m$, on peut choisir un élément $f \in A$ tel que $fz \in M[\mathfrak m]$ et $fz \not = 0$. Alors $\varphi(fz) = f\varphi(z)$ est non nul, donc $\varphi(z)$ est non nul.

\medskip\noindent
Surjectivité. Soit $z \in M$. Alors $x_1^n z = 0$ pour un certain $n \geq 0$. Montrons que $z \in x_1M$ par récurrence sur $n$. Si $n = 0$, alors $z = 0$, et le résultat est vrai. Si $n > 0$, l'application de $D_1$ donne $0 = n x_1^{n - 1} z + x_1^nD_1(z)$. Ainsi $x_1^{n - 1}(nz + x_1D_1(z)) = 0$. Par récurrence, on obtient $nz + x_1D_1(z) \in x_1M$. Puisque $n$ est inversible, on en déduit $z \in x_1M$. Ainsi $M$ est $x_1$-divisible. Si $\varphi$ n'est pas surjective, on peut choisir $e \in \bigoplus_{j \in J} E$ n'appartenant pas à $M$. En raisonnant comme ci-dessus, on peut supposer $\mathfrak m e \subset M$, en particulier $x_1 e \in M$. Il existe un élément $z_1 \in M$ tel que $x_1 z_1 = x_1 e$. D'où $x_1(z_1 - e) = 0$. En remplaçant $e$ par $e - z_1$, on peut supposer $e$ annulé par $x_1$. Il suffit donc de montrer que
$$
\varphi[x_1] :
M[x_1]
\longrightarrow
\left(\bigoplus\nolimits_{j \in J} E\right)[x_1] =
\bigoplus\nolimits_{j \in J} E[x_1]
$$
est surjective. Si $d = 1$, c'est vrai par construction de $\varphi$. Si $d > 1$, on observe que $E[x_1]$ est l'enveloppe injective du corps résiduel de $k[[x_2, \ldots, x_d]]$ ; voir Complexes dualisants, lemme \ref{dualizing-lemma-quotient}. Observons que $M[x_1]$, en tant que module sur $k[[x_2, \ldots, x_d]]$, est de torsion par rapport aux puissances de $\mathfrak m/(x_1)$ et est muni d'opérateurs $D_2, \ldots, D_d$ vérifiant la règle de Leibniz ci-dessus. Par récurrence sur $d$, on en déduit que $\varphi[x_1]$ est surjective, comme voulu.
\end{proof}

\begin{lemma}
\label{local-cohomology-lemma-structure-torsion-Frobenius-regular}
\begin{reference}
Résulte de \cite[Corollary 3.6]{Huneke-Sharp} moyennant quelques arguments supplémentaires. Résulte également directement de \cite[Theorem 1.4]{Lyubeznik2}.
\end{reference}
Soit $p$ un nombre premier. Soit $(A, \mathfrak m, k)$ un anneau local régulier tel que $p = 0$. Notons $F : A \to A$, $a \mapsto a^p$ l'endomorphisme de Frobenius. Soit $M$ un module de torsion par rapport aux puissances de $\mathfrak m$ tel que $M \otimes_{A, F} A \cong M$. Alors $M$ est isomorphe à une somme directe de copies de l'enveloppe injective $E$ de $k$.
\end{lemma}

\begin{proof}
Choisissons un ensemble $J$ et un homomorphisme de $A$-modules $\varphi : M \to \bigoplus_{j \in J} E$ qui envoie $M[\mathfrak m]$ isomorphiquement sur $(\bigoplus_{j \in J} E)[\mathfrak m] = \bigoplus_{j \in J} k$. Montrons que $\varphi$ est un isomorphisme, c'est-à-dire une bijection.

\medskip\noindent
Injectivité. Soit $z \in M$ non nul. Puisque $M$ est de torsion par rapport aux puissances de $\mathfrak m$, on peut choisir un élément $f \in A$ tel que $fz \in M[\mathfrak m]$ et $fz \not = 0$. Alors $\varphi(fz) = f\varphi(z)$ est non nul, donc $\varphi(z)$ est non nul.

\medskip\noindent
Surjectivité. Rappelons que $F$ est plat ; voir le lemme \ref{local-cohomology-lemma-frobenius-flat-regular}. Soit $x_1, \ldots, x_d$ un système minimal de générateurs de $\mathfrak m$. Notons
$$
M_n = M[x_1^{p^n}, \ldots, x_d^{p^n}]
$$
le sous-module de $M$ constitué des éléments annulés par $x_1^{p^n}, \ldots, x_d^{p^n}$. Ainsi $M_0 = M[\mathfrak m]$ est un espace vectoriel sur $k$. De plus, $M = \bigcup M_n$, puisque $M$ est de torsion par rapport aux puissances de $\mathfrak m$. Puisque $F^n$ est plat et que $F^n(x_i) = x_i^{p^n}$, on a
$$
M_n \cong (M \otimes_{A, F^n} A)[x_1^{p^n}, \ldots, x_d^{p^n}] =
M[x_1, \ldots, x_d] \otimes_{A, F^n} A =
M_0 \otimes_k A/(x_1^{p^n}, \ldots, x_d^{p^n})
$$
Ainsi $M_n$ est libre sur $A/(x_1^{p^n}, \ldots, x_d^{p^n})$. Un calcul montre que tout élément de $A/(x_1^{p^n}, \ldots, x_d^{p^n})$ annulé par $x_1^{p^n - 1}$ est divisible par $x_1$ ; par exemple, on peut utiliser le fait que $A/(x_1^{p^n}, \ldots, x_d^{p^n}) \cong
k[x_1, \ldots, x_d]/(x_1^{p^n}, \ldots, x_d^{p^n})$ d'après Algèbre, lemme \ref{algebra-lemma-regular-complete-containing-coefficient-field}. Il en est donc de même pour tout élément de $M_n$. Puisque tout élément de $M$ appartient à $M_n$ pour tout $n \gg 0$ et que tout élément de $M$ est annulé par une puissance de $x_1$, on conclut que $M$ est $x_1$-divisible.

\medskip\noindent
Posons $x = x_1$. Nous avons vu ci-dessus que $M$ est $x$-divisible. Si $\varphi$ n'est pas surjective, on peut choisir $e \in \bigoplus_{j \in J} E$ n'appartenant pas à $M$. En raisonnant comme ci-dessus, on peut supposer $\mathfrak m e \subset M$, en particulier $x e \in M$. Il existe un élément $z_1 \in M$ tel que $x z_1 = x e$. D'où $x(z_1 - e) = 0$. En remplaçant $e$ par $e - z_1$, on peut supposer $e$ annulé par $x$. Il suffit donc de montrer que
$$
\varphi[x] :
M[x]
\longrightarrow
\left(\bigoplus\nolimits_{j \in J} E\right)[x] =
\bigoplus\nolimits_{j \in J} E[x]
$$
est surjective. Si $d = 1$, c'est vrai par construction de $\varphi$. Si $d > 1$, on observe que $E[x]$ est l'enveloppe injective du corps résiduel de l'anneau régulier $A/xA$ ; voir Complexes dualisants, lemme \ref{dualizing-lemma-quotient}. Observons que $M[x]$, en tant que module sur $A/xA$, est de torsion par rapport aux puissances de $\mathfrak m/(x)$ et que l'on a
\begin{align*}
M[x] \otimes_{A/xA, F} A/xA
& = M[x] \otimes_{A, F} A \otimes_A A/xA \\
& = (M \otimes_{A, F} A)[x^p] \otimes_A A/xA \\
& \cong M[x^p] \otimes_A A/xA
\end{align*}
Raisonner à l'aide de la platitude de $F$, comme précédemment. Montrons que $M[x^p] \otimes_A A/xA \to M[x]$, $z \otimes 1 \mapsto x^{p - 1}z$ est un isomorphisme. Il suffit de le prouver pour chacun des modules $M_n$, $n > 0$ définis ci-dessus ; cela résulte alors du même résultat pour $A/(x_1^{p^n}, \ldots, x_d^{p^n})$ et $x = x_1$. Par récurrence sur $\dim(A)$, on en déduit que $\varphi[x]$ est surjective, comme voulu.
\end{proof}





\section{Structures supplémentaires sur la cohomologie locale}
\label{local-cohomology-section-additional}

\noindent
Voici un exemple de résultat.

\begin{lemma}
\label{local-cohomology-lemma-derivation}
Soit $A$ un anneau. Soit $I \subset A$ un idéal de type fini. Posons $Z = V(I)$. Pour toute dérivation $\theta : A \to A$, il existe un opérateur additif canonique $D$ sur les modules de cohomologie locale $H^i_Z(A)$, vérifiant la règle de Leibniz relativement à $\theta$.
\end{lemma}

\begin{proof}
Soient $f_1, \ldots, f_r$ des éléments engendrant $I$. Rappelons que $R\Gamma_Z(A)$ est calculé par le complexe

$$
A \to \prod\nolimits_{i_0} A_{f_{i_0}} \to
\prod\nolimits_{i_0 < i_1} A_{f_{i_0}f_{i_1}}
\to \ldots \to A_{f_1\ldots f_r}
$$
Voir Complexes dualisants, lemme \ref{dualizing-lemma-local-cohomology-adjoint}. Puisque $\theta$ se prolonge de manière unique en un opérateur additif sur tout localisé de $A$, vérifiant la règle de Leibniz relativement à $\theta$, le lemme est clair.
\end{proof}

\begin{lemma}
\label{local-cohomology-lemma-frobenius}
Soit $p$ un nombre premier. Soit $A$ un anneau tel que $p = 0$. Notons $F : A \to A$, $a \mapsto a^p$ l'endomorphisme de Frobenius. Soit $I \subset A$ un idéal de type fini. Posons $Z = V(I)$. Il existe un isomorphisme $R\Gamma_Z(A) \otimes_{A, F}^\mathbf{L} A \cong R\Gamma_Z(A)$.
\end{lemma}

\begin{proof}
Cela résulte de Complexes dualisants, lemme \ref{dualizing-lemma-torsion-change-rings}, et du fait que $Z = V(f_1^p, \ldots, f_r^p)$ si $I = (f_1, \ldots, f_r)$.
\end{proof}

\begin{lemma}
\label{local-cohomology-lemma-etale-derivation}
Soit $A$ un anneau. Soit $V \to \Spec(A)$ quasi-compact, quasi-séparé et \'etale. Pour toute dérivation $\theta : A \to A$, il existe un opérateur additif canonique $D$ sur $H^i(V, \mathcal{O}_V)$ vérifiant la règle de Leibniz relativement à $\theta$.
\end{lemma}

\begin{proof}
Si $V$ est séparé, on peut raisonner à l'aide d'un recouvrement ouvert affine $V = \bigcup_{j = 1, \ldots m} V_j$. En effet, puisque $V$ est séparé, on peut écrire $V_{j_0 \ldots j_p} = \Spec(B_{j_0 \ldots j_p})$. Voir Schémas, lemme \ref{schemes-lemma-characterize-separated}. On trouve alors que le $A$-module $H^i(V, \mathcal{O}_V)$ est le groupe de cohomologie en degré $i$ du complexe de {\v C}ech
$$
\prod B_{j_0} \to
\prod B_{j_0j_1} \to
\prod B_{j_0j_1j_2} \to \ldots
$$
Voir Cohomologie des schémas, lemme \ref{coherent-lemma-cech-cohomology-quasi-coherent}. Chaque $B = B_{j_0 \ldots j_p}$ est une algèbre \'etale sur $A$. Ainsi $\Omega_B = \Omega_A \otimes_A B$, et l'on en déduit que $\theta$ se prolonge de manière unique en une dérivation $\theta_B : B \to B$. Ces applications définissent un endomorphisme du complexe de {\v C}ech et les opérateurs voulus sur les groupes de cohomologie.

\medskip\noindent
Dans le cas général, on utilise un hyperrecouvrement de $V$ par des ouverts affines, exactement comme dans la première partie de la démonstration de Cohomologie des schémas, lemme \ref{coherent-lemma-hypercoverings}. Nous omettons les détails.
\end{proof}

\begin{remark}
\label{local-cohomology-remark-higher-order-operators}
On peut renforcer les lemmes \ref{local-cohomology-lemma-derivation} et \ref{local-cohomology-lemma-etale-derivation} pour y inclure les opérateurs différentiels d'ordre supérieur, à l'aide d'Algèbre, lemme \ref{algebra-lemma-formally-etale-principal-parts} et remarque \ref{algebra-remark-formally-etale-differential-operators}. Si le besoin s'en fait sentir, nous énoncerons et démontrerons ici un lemme précis.
\end{remark}

\begin{lemma}
\label{local-cohomology-lemma-etale-frobenius}
Soit $p$ un nombre premier. Soit $A$ un anneau tel que $p = 0$. Notons $F : A \to A$, $a \mapsto a^p$ l'endomorphisme de Frobenius. Si $V \to \Spec(A)$ est quasi-compact, quasi-séparé et \'etale, alors il existe un isomorphisme $R\Gamma(V, \mathcal{O}_V) \otimes_{A, F}^\mathbf{L} A \cong
R\Gamma(V, \mathcal{O}_V)$.
\end{lemma}

\begin{proof}
Observons que le morphisme de Frobenius relatif
$$
V \longrightarrow V \times_{\Spec(A), \Spec(F)} \Spec(A)
$$
de $V$ sur $A$ est un isomorphisme ; voir Morphismes \'etales, lemme \ref{etale-lemma-relative-frobenius-etale}. Le lemme résulte donc du changement de base en cohomologie ; voir Catégories dérivées des schémas, lemme \ref{perfect-lemma-compare-base-change}. Observons que, puisque $V$ est \'etale sur $A$, il est plat sur $A$.
\end{proof}





\section{Un peu d'uniformité, I}
\label{local-cohomology-section-uniform-vanishing}

\noindent
Le principal objectif de cette section est d'énoncer et de démontrer le lemme \ref{local-cohomology-lemma-characterize-vanishing-tor-ext-above-e}.

\begin{lemma}
\label{local-cohomology-lemma-map-tor-1-zero}
Soit $R$ un anneau. Soit $M \to M'$ un homomorphisme de $R$-modules, avec $M$ de présentation finie, tel que $\text{Tor}_1^R(M, N) \to \text{Tor}_1^R(M', N)$ soit nul pour tous les $R$-modules $N$. Alors $M \to M'$ se factorise par un $R$-module libre.
\end{lemma}

\begin{proof}
On peut choisir un morphisme de suites exactes courtes
$$
\xymatrix{
0 \ar[r] &
K \ar[r] \ar[d] &
R^{\oplus r} \ar[r] \ar[d] &
M \ar[r] \ar[d] &
0 \\
0 \ar[r] &
K' \ar[r] &
\bigoplus_{i \in I} R \ar[r] &
M' \ar[r] &
0
}
$$
dont la flèche verticale de droite est l'application donnée. Ce morphisme se factorise par la suite exacte courte
\begin{equation}
\label{local-cohomology-equation-pushout}
0 \to K' \to E \to M \to 0
\end{equation}
qui est la somme amalgamée de la première suite exacte courte par $K \to K'$. Une chasse au diagramme montre que l'hypothèse du lemme implique que l'homomorphisme de bord $\text{Tor}_1^R(M, N) \to K' \otimes_R N$ induit par (\ref{local-cohomology-equation-pushout}) est nul, c'est-à-dire que la suite (\ref{local-cohomology-equation-pushout}) est universellement exacte. Algèbre, lemme \ref{algebra-lemma-universally-exact-split}, implique alors que (\ref{local-cohomology-equation-pushout}) est scindée (c'est ici que l'on utilise le fait que $M$ est de présentation finie). Ainsi l'application $M \to M'$ se factorise par $\bigoplus_{i \in I} R$, ce qui conclut.
\end{proof}

\begin{lemma}
\label{local-cohomology-lemma-characterize-vanishing-tor-ext-above-e}
Soit $R$ un anneau. Soit $\alpha : M \to M'$ un homomorphisme de $R$-modules. Soient $P_\bullet \to M$ et $P'_\bullet \to M'$ des résolutions par des $R$-modules projectifs. Soit $e \geq 0$ un entier. Considérons les conditions suivantes :
\begin{enumerate}
\item Il existe un morphisme de complexes $a_\bullet : P_\bullet \to P'_\bullet$ induisant $\alpha$ en cohomologie, avec $a_i = 0$ pour $i > e$.
\item Il existe un morphisme de complexes $a_\bullet : P_\bullet \to P'_\bullet$ induisant $\alpha$ en cohomologie, avec $a_{e + 1} = 0$.
\item L'application $\Ext^i_R(M', N) \to \Ext^i_R(M, N)$ est nulle pour tous les $R$-modules $N$ et tout $i > e$.
\item L'application $\Ext^{e + 1}_R(M', N) \to \Ext^{e + 1}_R(M, N)$ est nulle pour tous les $R$-modules $N$.
\item Soit $N = \Im(P'_{e + 1} \to P'_e)$, et notons $\xi \in \Ext^{e + 1}_R(M', N)$ l'élément canonique (voir la démonstration). Alors $\xi$ a une image nulle dans $\Ext^{e + 1}_R(M, N)$.
\item L'application $\text{Tor}_i^R(M, N) \to \text{Tor}_i^R(M', N)$ est nulle pour tous les $R$-modules $N$ et tout $i > e$.
\item L'application $\text{Tor}_{e + 1}^R(M, N) \to \text{Tor}_{e + 1}^R(M', N)$ est nulle pour tous les $R$-modules $N$.
\end{enumerate}
On a toujours les implications
$$
(1) \Leftrightarrow (2) \Leftrightarrow (3) \Leftrightarrow (4)
\Leftrightarrow (5) \Rightarrow (6) \Leftrightarrow (7)
$$
Si $M$ est $(-e - 1)$-pseudo-cohérent (par exemple si $R$ est noethérien et $M$ est un $R$-module de type fini), alors toutes les conditions sont équivalentes.
\end{lemma}

\begin{proof}
Il est clair que (2) implique (1). Si $a_\bullet$ est comme dans (1), on peut considérer le morphisme de complexes $a'_\bullet : P_\bullet \to P'_\bullet$ défini par $a'_i = a_i$ pour $i \leq e + 1$ et $a'_i = 0$ pour $i \geq e + 1$, afin d'obtenir un morphisme de complexes comme dans (2). Ainsi (1) équivaut à (2).

\medskip\noindent
La construction des foncteurs $\Ext$ et $\text{Tor}$ à l'aide de résolutions (Algèbre, sections \ref{algebra-section-ext} et \ref{algebra-section-tor}) montre que (1) et (2) impliquent toutes les autres conditions.

\medskip\noindent
Il est clair que (3) implique (4), qui implique (5). Soit $N$ comme dans (5). L'application canonique $\tilde \xi : P'_{e + 1} \to N$ précomposée avec $P'_{e + 2} \to P'_{e + 1}$ est nulle. On peut donc considérer la classe $\xi$ de $\tilde \xi$ dans
$$
\Ext^{e + 1}_R(M', N) =
\frac{\Ker(\Hom(P'_{e + 1}, N \to \Hom(P'_{e + 2}, N)}{
\Im(\Hom(P'_e, N \to \Hom(P'_{e + 1}, N)}
$$
Choisissons un morphisme de complexes $a_\bullet : P_\bullet \to P'_\bullet$ relevant $\alpha$ ; voir Catégories dérivées, lemme \ref{derived-lemma-morphisms-lift-projective}. Si $\xi$ a une image nulle dans $\Ext^{e + 1}_R(M', N)$, on trouve alors une application $\varphi : P_e \to N$ telle que $\tilde \xi  \circ a_{e + 1} = \varphi \circ d$. On obtient ainsi un morphisme de complexes
$$
\xymatrix{
\ldots \ar[r] &
P_{e + 1} \ar[r] \ar[d]^0 &
P_e \ar[r] \ar[d]^{a_e - \varphi} &
P_{e - 1} \ar[r] \ar[d]^{a_{e - 1}} &
\ldots \\
\ldots \ar[r] &
P'_{e + 1} \ar[r] &
P'_e \ar[r] &
P'_{e - 1} \ar[r] &
\ldots
}
$$
comme dans (2). Ainsi (1) -- (5) sont équivalentes.

\medskip\noindent
L'équivalence de (6) et (7) résulte du décalage de dimension ; nous omettons les détails.

\medskip\noindent
Supposons que $M$ soit $(-e - 1)$-pseudo-cohérent. (L'assertion entre parenthèses du lemme résulte de Compléments d'algèbre, lemme \ref{more-algebra-lemma-Noetherian-pseudo-coherent}.) Montrons que (7) implique (4), ce qui achèvera la démonstration. Nous raisonnons par récurrence sur $e$. Le cas initial est $e = 0$. Alors $M$ est de présentation finie d'après Compléments d'algèbre, lemme \ref{more-algebra-lemma-n-pseudo-module}, et le lemme \ref{local-cohomology-lemma-map-tor-1-zero} montre que $M \to M'$ se factorise par un module libre. Bien entendu, si $M \to M'$ se factorise par un module libre, alors $\Ext^i_R(M', N) \to \Ext^i_R(M, N)$ est nulle pour tout $i > 0$, comme voulu. Supposons $e > 0$. On peut choisir un morphisme de suites exactes courtes
$$
\xymatrix{
0 \ar[r] &
K \ar[r] \ar[d] &
R^{\oplus r} \ar[r] \ar[d] &
M \ar[r] \ar[d] &
0 \\
0 \ar[r] &
K' \ar[r] &
\bigoplus_{i \in I} R \ar[r] &
M' \ar[r] &
0
}
$$
dont la flèche verticale de droite est l'application donnée. On obtient $\text{Tor}_{i + 1}^R(M, N) = \text{Tor}^R_i(K, N)$ et $\Ext^{i + 1}_R(M, N) = \Ext^i_R(K, N)$ pour $i \geq 1$ et pour tous les $R$-modules $N$, et de même pour $M', K'$. Ainsi $\text{Tor}_e^R(K, N) \to \text{Tor}_e^R(K', N)$ est nulle pour tous les $R$-modules $N$. D'après Compléments d'algèbre, lemme \ref{more-algebra-lemma-cone-pseudo-coherent}, $K$ est $(-e)$-pseudo-cohérent. Par récurrence, on conclut que $\Ext^e(K', N) \to \Ext^e(K, N)$ est nulle pour tous les $R$-modules $N$, ce qui donne le résultat voulu.

\end{proof}

\begin{lemma}
\label{local-cohomology-lemma-cd-sequence-Koszul}
Soit $I$ un idéal d'un anneau noethérien $A$. Pour tout $n \geq 1$, il existe $m > n$ tel que l'application $A/I^m \to A/I^n$ vérifie les conditions équivalentes du lemme \ref{local-cohomology-lemma-characterize-vanishing-tor-ext-above-e} avec $e = \text{cd}(A, I)$.
\end{lemma}

\begin{proof}
Soit $\xi \in \Ext^{e + 1}_A(A/I^n, N)$ l'élément construit au lemme \ref{local-cohomology-lemma-characterize-vanishing-tor-ext-above-e}, point (5). Puisque $e = \text{cd}(A, I)$, on a $0 = H^{e + 1}_Z(N) = H^{e + 1}_I(N) = \colim \Ext^{e + 1}(A/I^m, N)$ d'après Complexes dualisants, lemmes \ref{dualizing-lemma-local-cohomology-noetherian} et \ref{dualizing-lemma-local-cohomology-ext}. On peut donc choisir $m \geq n$ tel que $\xi$ ait une image nulle dans $\Ext^{e + 1}_A(A/I^m, N)$, comme voulu.

\end{proof}






\section{Un peu d'uniformité, II}
\label{local-cohomology-section-uniform-vanishing-bis}

\noindent
Soit $I$ un idéal d'un anneau noethérien $A$. Soit $M$ un $A$-module de type fini. Soit $i > 0$. D'après Compléments d'algèbre, lemme \ref{more-algebra-lemma-tor-strictly-pro-zero}, il existe $c = c(A, I, M, i)$ tel que $\text{Tor}^A_i(M, A/I^n) \to \text{Tor}^A_i(M, A/I^{n - c})$ soit nulle pour tout $n \geq c$. Dans cette section, nous présentons des résultats montrant que l'on peut parfois choisir une constante $c$ convenant simultanément à tous les $A$-modules $M$ (et à une plage d'indices $i$). Ce sujet est lié au théorème d'Artin-Rees uniforme, étudié dans \cite{Huneke-uniform} et \cite{AHS}.

\medskip\noindent
Dans la remarque \ref{local-cohomology-remark-strict-pro-isomorphism}, nous l'appliquerons pour montrer que divers systèmes projectifs liés à la complétion dérivée sont (ou ne sont pas) strictement pro-isomorphes.

\medskip\noindent
Le lemme suivant peut être sensiblement renforcé.

\begin{lemma}
\label{local-cohomology-lemma-maps-zero-fixed-torsion}
Soit $I$ un idéal d'un anneau noethérien $A$. Pour tous $m \geq 0$ et $i > 0$, il existe $c = c(A, I, m, i) \geq 0$ tel que, pour tout $A$-module $M$ annulé par $I^m$, l'application
$$
\text{Tor}^A_i(M, A/I^n) \to \text{Tor}^A_i(M, A/I^{n - c})
$$
soit nulle pour tout $n \geq c$.
\end{lemma}

\begin{proof}
Raisonnons par récurrence sur $i$. Cas initial $i = 1$. La suite exacte courte $0 \to I^n \to A \to A/I^n \to 0$ définit une injection $\text{Tor}_1^A(M, A/I^n) \subset I^n \otimes_A M$ ; voir Algèbre, remarque \ref{algebra-remark-Tor-ring-mod-ideal}. Puisque $M$ est annulé par $I^m$, l'application $I^n \otimes_A M \to I^{n - m} \otimes_A M$ est nulle pour $n \geq m$. Le résultat est donc vrai avec $c = m$.

\medskip\noindent
Étape de récurrence. Soit $i > 1$ et supposons que $c$ convienne pour $i - 1$. D'après Compléments d'algèbre, lemme \ref{more-algebra-lemma-tor-strictly-pro-zero}, appliqué à $M = A/I^m$, on peut choisir $c' \geq 0$ tel que $\text{Tor}_i(A/I^m, A/I^n) \to \text{Tor}_i(A/I^m, A/I^{n - c'})$ soit nulle pour $n \geq c'$. Soit $M$ annulé par $I^m$. Choisissons une suite exacte courte
$$
0 \to S \to \bigoplus\nolimits_{i \in I} A/I^m \to M \to 0
$$
La suite exacte longue de Tor correspondante fournit une suite exacte
$$
\text{Tor}_i^A(\bigoplus\nolimits_{i \in I} A/I^m, A/I^n) \to
\text{Tor}_i^A(M, A/I^n) \to
\text{Tor}_{i - 1}^A(S, A/I^n)
$$
pour tous les entiers $n \geq 0$. Si $n \geq c + c'$, alors l'application $\text{Tor}_{i - 1}^A(S, A/I^n) \to \text{Tor}_{i - 1}^A(S, A/I^{n - c})$ est nulle et l'application $\text{Tor}_i^A(A/I^m, A/I^{n - c}) \to 
\text{Tor}_i^A(A/I^m, A/I^{n - c - c'})$ est nulle. En combinant cela avec les suites exactes courtes, on voit que le résultat est vrai pour $i$ avec la constante $c + c'$.

\end{proof}

\begin{lemma}
\label{local-cohomology-lemma-annihilates-affine}
Soit $I = (a_1, \ldots, a_t)$ un idéal d'un anneau noethérien $A$. Posons $a = a_1$ et notons $B = A[\frac{I}{a}]$ l'algèbre d'éclatement affine. Il existe $c > 0$ tel que $\text{Tor}_i^A(B, M)$ soit annulé par $I^c$ pour tous les $A$-modules $M$ et tout $i \geq t$.
\end{lemma}

\begin{proof}
Rappelons que $B$ est le quotient de $A[x_2, \ldots, x_t]/(a_1x_2 - a_2, \ldots, a_1x_t - a_t)$ par son idéal de torsion par rapport aux puissances de $a_1$ ; voir Algèbre, lemme \ref{algebra-lemma-affine-blowup-quotient-description}. Soit

$$
B_\bullet = \text{complexe de Koszul sur }a_1x_2 - a_2, \ldots, a_1x_t - a_t
\text{ sur }A[x_2, \ldots, x_t]
$$
considéré comme complexe de chaînes concentré en degrés $(t - 1), \ldots, 0$. Le complexe $B_\bullet[1/a_1]$ est isomorphe au complexe de Koszul sur $x_2 - a_2/a_1, \ldots, x_t - a_t/a_1$, qui est une suite régulière dans $A[1/a_1][x_2, \ldots, x_t]$. Puisque les suites régulières sont Koszul-régulières, on en déduit que l'augmentation
$$
\epsilon : B_\bullet \longrightarrow B
$$
est un quasi-isomorphisme après inversion de $a_1$. Puisque les modules d'homologie du cône $C_\bullet$ de $\epsilon$ sont des $A[x_2, \ldots, x_n]$-modules de type fini et que $C_\bullet$ est borné, il existe $c \geq 0$ tel que $a_1^c$ les annule tous. D'après Catégories dérivées, lemme \ref{derived-lemma-trick-vanishing-composition}, cela implique, après remplacement éventuel de $c$ par un entier plus grand, que $a_1^c$ est nul sur $C_\bullet$ dans $D(A)$. La démonstration s'achève en considérant le triangle distingué
$$
B_\bullet \otimes_A^\mathbf{L} M \to
B \otimes_A^\mathbf{L} M \to
C_\bullet \otimes_A^\mathbf{L} M
$$
En effet, le premier terme est représenté par $B_\bullet \otimes_A M$, concentré en degrés homologiques $(t - 1), \ldots, 0$, puisque les termes du complexe de Koszul $B_\bullet$ sont des $A$-modules libres (donc plats). Ainsi $\text{Tor}_i^A(B, M) = H_i(C_\bullet \otimes_A^\mathbf{L} M)$ pour $i > t - 1$, et ce module est annulé par $a_1^c$. Puisque $a_1^cB = I^cB$ et que le module Tor est un module sur $B$, on conclut.
\end{proof}

\noindent
Pour le reste de cette section, fixons un anneau noethérien $A$ et un idéal $I \subset A$. Notons
$$
p : X \to \Spec(A)
$$
l'éclatement de $\Spec(A)$ en l'idéal $I$. Autrement dit, $X$ est le $\text{Proj}$ de l'algèbre de Rees $\bigoplus_{n \geq 0} I^n$. D'après Cohomologie des schémas, lemmes \ref{coherent-lemma-coherent-on-proj} et \ref{coherent-lemma-recover-tail-graded-module}, on peut choisir un entier $q(A, I) \geq 0$ tel que, pour tout $q \geq q(A, I)$, on ait $H^i(X, \mathcal{O}_X(q)) = 0$ pour $i > 0$ et $H^0(X, \mathcal{O}_X(q)) = I^q$.

\begin{lemma}
\label{local-cohomology-lemma-compute-tor-Iq}
Dans la situation ci-dessus, pour $q \geq q(A, I)$ et tout $A$-module $M$, on a
$$
R\Gamma(X, Lp^*\widetilde{M}(q)) \cong M \otimes_A^\mathbf{L} I^q
$$
dans $D(A)$.
\end{lemma}

\begin{proof}
Choisissons une résolution libre $F_\bullet \to M$. Alors $\widetilde{F}_\bullet$ est une résolution plate de $\widetilde{M}$. Ainsi $Lp^*\widetilde{M}$ est donné par le complexe $p^*\widetilde{F}_\bullet$. Donc $Lp^*\widetilde{M}(q)$ est donné par le complexe $p^*\widetilde{F}_\bullet(q)$. Puisque les $p^*\widetilde{F}_i(q)$ sont acycliques à droite pour $\Gamma(X, -)$, par le choix de $q \geq q(A, I)$, et que l'on a $\Gamma(X, p^*\widetilde{F}_i(q)) = I^qF_i$, par le choix de $q \geq q(A, I)$, on obtient que $R\Gamma(X, Lp^*\widetilde{M}(q))$ est donné par le complexe de termes $I^qF_i$, d'après Catégories dérivées des schémas, lemme \ref{perfect-lemma-acyclicity-lemma-global}. Le résultat s'ensuit, car le complexe $I^qF_\bullet$ calcule $M \otimes_A^\mathbf{L} I^q$ par définition.
\end{proof}

\begin{lemma}
\label{local-cohomology-lemma-annihilates}
Dans la situation ci-dessus, soit $t$ une borne supérieure du nombre de générateurs de $I$. Il existe un entier $c = c(A, I) \geq 0$ tel que, pour tout $A$-module $M$, les faisceaux de cohomologie $H^j(Lp^*\widetilde{M})$ soient annulés par $I^c$ pour $j \leq -t$.
\end{lemma}

\begin{proof}
Écrivons $I = (a_1, \ldots, a_t)$. La question est locale sur les ouverts affines de $X$. Pour $1 \leq i \leq t$, soit $B_i = A[\frac{I}{a_i}]$ l'algèbre d'éclatement affine. Alors $X$ possède un recouvrement ouvert affine par les spectres des anneaux $B_i$ ; voir Diviseurs, lemme \ref{divisors-lemma-blowing-up-affine}. D'après la description de l'image réciproque dérivée donnée dans Catégories dérivées des schémas, lemme \ref{perfect-lemma-quasi-coherence-pullback}, il suffit de montrer que, pour chaque $i$, il existe $c \geq 0$ tel que
$$
\text{Tor}_j^A(B_i, M)
$$
soit annulé par $I^c$ pour $j \geq t$. C'est le lemme \ref{local-cohomology-lemma-annihilates-affine}.
\end{proof}

\begin{lemma}
\label{local-cohomology-lemma-annihilates-tors}
Dans la situation ci-dessus, soit $t$ une borne supérieure du nombre de générateurs de $I$. Il existe un entier $c = c(A, I) \geq 0$ tel que, pour tout $A$-module $M$, les modules Tor $\text{Tor}_i^A(M, A/I^q)$ soient annulés par $I^c$ pour $i > t$ et tout $q \geq 0$.
\end{lemma}

\begin{proof}
Soit $q(A, I)$ comme ci-dessus. Pour $q \geq q(A, I)$, on a
$$
R\Gamma(X, Lp^*\widetilde{M}(q)) = M \otimes_A^\mathbf{L} I^q
$$
d'après le lemme \ref{local-cohomology-lemma-compute-tor-Iq}. On dispose d'une suite spectrale bornée et convergente
$$
H^a(X, H^b(Lp^*\widetilde{M}(q))) \Rightarrow
\text{Tor}_{-a - b}^A(M, I^q)
$$
d'après Catégories dérivées des schémas, lemme \ref{perfect-lemma-spectral-sequence}. Soit $d$ un entier comme dans Cohomologie des schémas, lemme \ref{coherent-lemma-vanishing-nr-affines-quasi-separated} (on peut en fait prendre $d = t$ ; voir Cohomologie des schémas, lemme \ref{coherent-lemma-vanishing-nr-affines}). On voit alors que $H^{-i}(X, Lp^*\widetilde{M}(q)) = \text{Tor}_i^A(M, I^q)$ possède une filtration finie d'au plus $d$ étapes, dont les gradués sont des sous-quotients des modules
$$
H^a(X, H^{- i - a}(Lp^*\widetilde{M})(q)),\quad
a = 0, 1, \ldots, d - 1
$$
Si $i \geq t$, tous ces modules sont annulés par $I^c$, où $c = c(A, I)$ est comme dans le lemme \ref{local-cohomology-lemma-annihilates}, car les faisceaux de cohomologie $H^{- i - a}(Lp^*\widetilde{M})$ sont tous annulés par $I^c$ d'après ce lemme. Ainsi $\text{Tor}_i^A(M, I^q)$ est annulé par $I^{dc}$ pour $q \geq q(A, I)$ et $i \geq t$. La suite exacte courte $0 \to I^q \to A \to A/I^q \to 0$ montre alors que $\text{Tor}_i(M, A/I^q)$ est annulé par $I^{dc}$ pour $q \geq q(A, I)$ et $i > t$. On conclut que $I^m$, avec $m = \max(dc, q(A, I) - 1)$, annule $\text{Tor}_i^A(M, A/I^q)$ pour tout $q \geq 0$ et tout $i > t$, comme voulu.
\end{proof}

\begin{lemma}
\label{local-cohomology-lemma-tor-maps-vanish}
Soit $I$ un idéal d'un anneau noethérien $A$. Soit $t \geq 0$ une borne supérieure du nombre de générateurs de $I$. Il existe $N, c \geq 0$ tels que les applications
$$
\text{Tor}_{t + 1}^A(M, A/I^n) \to \text{Tor}_{t + 1}^A(M, A/I^{n - c})
$$
soient nulles pour tout $A$-module $M$ et tout $n \geq N$.
\end{lemma}

\begin{proof}
Soit $c_1$ la constante obtenue au lemme \ref{local-cohomology-lemma-annihilates-tors}. Rappelons que cette constante $c_1$ convient simultanément pour $\text{Tor}_i$ pour tous les $i > t$.

\medskip\noindent
Écrivons $I = (a_1, \ldots, a_t)$. Pour un $A$-module $M$, posons
$$
\ell(M) =
\#\{i \mid 1 \leq i \leq t,\ a_i^{c_1}\text{ est nul sur }M\}
$$
C'est un élément de $\{0, 1, \ldots, t\}$. Démontrons par récurrence descendante sur $0 \leq s \leq t$ l'assertion suivante $H_s$ : il existe $N, c \geq 0$ tels que, pour tout module $M$ avec $\ell(M) \geq s$, les applications
$$
\text{Tor}_{t + 1 + i}^A(M, A/I^n) \to \text{Tor}_{t + 1 + i}^A(M, A/I^{n - c})
$$
soient nulles pour $i = 0, \ldots, s$ et tout $n \geq N$.

\medskip\noindent
Cas initial : $s = t$. Si $\ell(M) = t$, alors $M$ est annulé par $(a_1^{c_1}, \ldots, a_t^{c_1}\}$ et donc par $I^{t(c_1 - 1) + 1}$. Le lemme \ref{local-cohomology-lemma-maps-zero-fixed-torsion} montre que $H_t$ est vérifiée en prenant pour $c = N$ le maximum des entiers $c(A, I, t(c_1 - 1) + 1, t + 1), \ldots, c(A, I, t(c_1 - 1) + 1, 2t + 1)$ obtenus dans ce lemme.

\medskip\noindent
Étape de récurrence. Soit $0 \leq s < t$ et supposons donnés $N, c$ comme dans $H_{s + 1}$. Considérons un module $M$ tel que $\ell(M) = s$. On peut alors choisir $i$ tel que $a_i^{c_1}$ soit non nul sur $M$. Il s'ensuit que $\ell(M[a_i^c]) \geq s + 1$ et $\ell(M/a_i^{c_1}M) \geq s + 1$, et l'hypothèse de récurrence s'applique à ces modules. Considérons la suite exacte
$$
0 \to M[a_i^{c_1}] \to M \xrightarrow{a_i^{c_1}} M \to M/a_i^{c_1}M \to 0
$$
Notons $E \subset M$ l'image de la flèche centrale. Considérons le diagramme de modules Tor correspondant
$$
\xymatrix{
&
&
\text{Tor}_{i + 1}(M/a_i^{c_1}M, A/I^q) \ar[d] \\
\text{Tor}_i(M[a_i^{c_1}], A/I^q) \ar[r] &
\text{Tor}_i(M, A/I^q) \ar[r] \ar[rd]^0 &
\text{Tor}_i(E, A/I^q) \ar[d] \\
&
&
\text{Tor}_i(M, A/I^q)
}
$$
à lignes et colonnes exactes (pour tout $q$). La flèche dirigée vers le sud-est est nulle par le choix de $c_1$. On en déduit que le module $\text{Tor}_i(M, A/I^q)$ s'insère dans une suite exacte entre un quotient de $\text{Tor}_i(M[a_i^{c_1}], A/I^q)$ et un sous-module de $\text{Tor}_{i + 1}(M/a_i^{c_1}M, A/I^q)$ dans le diagramme. On conclut donc que $H_s$ est vérifiée en remplaçant $N$ par $N + c$ et $c$ par $2c$. Nous omettons certains détails.
\end{proof}

\begin{proposition}
\label{local-cohomology-proposition-uniform-artin-rees}
Soit $I$ un idéal d'un anneau noethérien $A$. Soit $t \geq 0$ une borne supérieure du nombre de générateurs de $I$. Il existe $N, c \geq 0$ tels que, pour $n \geq N$, les applications
$$
A/I^n \to A/I^{n - c}
$$
vérifient les conditions équivalentes du lemme \ref{local-cohomology-lemma-characterize-vanishing-tor-ext-above-e} avec $e = t$.
\end{proposition}

\begin{proof}
Conséquence immédiate des lemmes \ref{local-cohomology-lemma-tor-maps-vanish} et \ref{local-cohomology-lemma-characterize-vanishing-tor-ext-above-e}.
\end{proof}

\begin{remark}
\label{local-cohomology-remark-better-bound}
L'article \cite{AHS} montre, parmi bien d'autres résultats, que si $A$ est local, la proposition \ref{local-cohomology-proposition-uniform-artin-rees} reste vraie en remplaçant $e = t$ par $e = \dim(A)$. Le lemme \ref{local-cohomology-lemma-cd-sequence-Koszul} conduit naturellement à se demander si la proposition \ref{local-cohomology-proposition-uniform-artin-rees} reste vraie en remplaçant $e = t$ par $e = \text{cd}(A, I)$. Nous ne le savons pas.
\end{remark}

\begin{remark}
\label{local-cohomology-remark-strict-pro-isomorphism}
Soit $I$ un idéal d'un anneau noethérien $A$. Écrivons $I = (f_1, \ldots, f_r)$. Notons $K_n^\bullet$ le complexe de Koszul sur $f_1^n, \ldots, f_r^n$, comme dans Compléments d'algèbre, situation \ref{more-algebra-situation-koszul}, et notons $K_n \in D(A)$ l'objet correspondant. Soit $M^\bullet$ un complexe borné de $A$-modules de type fini et notons $M \in D(A)$ l'objet correspondant. Considérons les systèmes projectifs suivants dans $D(A)$ :

\begin{enumerate}
\item $M^\bullet/I^nM^\bullet$, c'est-à-dire le complexe de termes $M^i/I^nM^i$,
\item $M \otimes_A^\mathbf{L} A/I^n$,
\item $M \otimes_A^\mathbf{L} K_n$,
\item $M \otimes_P^\mathbf{L} P/J^n$ (voir ci-dessous).
\end{enumerate}
Tous ces systèmes projectifs sont isomorphes en tant que pro-objets : l'isomorphisme entre (2) et (3) résulte de Compléments d'algèbre, lemme \ref{more-algebra-lemma-sequence-Koszul-complexes}. L'isomorphisme entre (1) et (2) est donné dans Compléments d'algèbre, lemme \ref{more-algebra-lemma-derived-completion-plain-completion}. Pour le dernier, voir ci-dessous.

\medskip\noindent
On peut toutefois se demander si ces isomorphismes de systèmes projectifs sont « stricts » ; cette terminologie et cette question se rattachent à la discussion de \cite[pages 61, 62]{quillenhomology}. Plus précisément, étant donnée une catégorie $\mathcal{C}$, on peut définir une « pro-catégorie stricte » dont les objets sont les systèmes projectifs $(X_n)$ et dont les morphismes $(X_n) \to (Y_n)$ sont donnés par des familles $(c, \varphi_n)$ constituées de $c \geq 0$ et de morphismes $\varphi_n : X_n \to Y_{n - c}$ pour tout $n \geq c$, satisfaisant à une condition évidente de compatibilité, à une certaine équivalence près (donnée essentiellement par l'augmentation de $c$). On demande alors si les systèmes projectifs ci-dessus sont isomorphes dans cette pro-catégorie stricte.

\medskip\noindent
Cela ne peut manifestement pas être le cas pour (1) et (3), même lorsque $M = A[0]$. En effet, le système $H^0(K_n) = A/(f_1^n, \ldots, f_r^n)$ n'est en général pas strictement pro-isomorphe, dans la catégorie des modules, au système $A/I^n$. Par exemple, si l'on prend $A = \mathbf{Z}[x_1, \ldots, x_r]$ et $f_i = x_i$, alors $H^0(K_n)$ n'est pas annulé par $I^{r(n - 1)}$.\footnote{Bien entendu, on peut demander si ces systèmes projectifs sont isomorphes dans une catégorie dont les objets sont les systèmes projectifs et dont les applications sont données par des familles $(r, c, \varphi_n)$ constituées de $r \geq 1$, de $c \geq 0$ et d'applications $\varphi_n : X_{rn} \to Y_{n - c}$ pour $n \geq c$.}

\medskip\noindent
Il se trouve que les résultats précédents montrent que l'application naturelle de (2) vers (1), étudiée dans Compléments d'algèbre, lemme \ref{more-algebra-lemma-derived-completion-plain-completion}, est un pro-isomorphisme strict. Esquissons la démonstration. Des arguments usuels faisant intervenir les troncations bêtes permettent d'abord de se ramener au cas où $M^\bullet$ est donné par un seul $A$-module de type fini $M$ placé en degré $0$. Choisissons $N, c \geq 0$ comme dans la proposition \ref{local-cohomology-proposition-uniform-artin-rees}. La proposition implique que, pour $n \geq N$, on obtient des factorisations

$$
M \otimes_A^\mathbf{L} A/I^n
\to
\tau_{\geq -t}(M \otimes_A^\mathbf{L} A/I^n)
\to
M \otimes_A^\mathbf{L} A/I^{n - c}
$$
des applications de transition du système (2). D'autre part, d'après Compléments d'algèbre, lemme \ref{more-algebra-lemma-tor-strictly-pro-zero}, on peut trouver une autre constante $c' = c'(M) \geq 0$ telle que les applications $\text{Tor}_i^A(M, A/I^{n'}) \to \text{Tor}_i(M, A/I^{n' - c'})$ soient nulles pour $i = 1, 2, \ldots, t$ et $n' \geq c'$. Il résulte alors de Catégories dérivées, lemme \ref{derived-lemma-trick-vanishing-composition}, que l'application
$$
\tau_{\geq -t}(M \otimes_A^\mathbf{L} A/I^{n + tc'})
\to
\tau_{\geq -t}(M \otimes_A^\mathbf{L} A/I^n)
$$
se factorise par $M \otimes_A^\mathbf{L}A/I^{n + tc'} \to M/I^{n + tc'}M$. En combinant cela avec le résultat précédent, on obtient une factorisation
$$
M \otimes_A^\mathbf{L}A/I^{n + tc'} \to M/I^{n + tc'}M
\to M \otimes_A^\mathbf{L} A/I^{n - c}
$$
qui donne ce que l'on cherche. Si nous avons besoin de ce résultat, nous l'énoncerons soigneusement et en donnerons une démonstration détaillée.

\medskip\noindent
Pour (4), supposons donnés un anneau noethérien $P$, un homomorphisme d'anneaux $P \to A$ et un idéal $J \subset P$ tels que $I = JA$. D'après Compléments d'algèbre, section \ref{more-algebra-section-derived-base-change}, on obtient un foncteur $M \otimes_P^\mathbf{L} - : D(P) \to D(A)$ et un système projectif $M \otimes_P^\mathbf{L} P/J^n$ dans $D(A)$, comme dans (4). Si $P$ est noethérien, le système (4) est pro-isomorphe au système (1), par comparaison avec les complexes de Koszul. Si $P \to A$ est fini, le système (4) est strictement pro-isomorphe au système (2), car le système projectif $A \otimes_P^\mathbf{L} P/J^n$ est strictement pro-isomorphe au système projectif $A/I^n$ (d'après la discussion précédente) et l'on a
$$
M \otimes_P^\mathbf{L} P/J^n = M \otimes_A^\mathbf{L}
(A \otimes_P^\mathbf{L} P/J^n)
$$
d'après Compléments d'algèbre, lemme \ref{more-algebra-lemma-derived-base-change}.

\medskip\noindent
Un exemple usuel dans (4) consiste à prendre $P = \mathbf{Z}[x_1, \ldots, x_r]$, l'application $P \to A$ envoyant $x_i$ sur $f_i$, et $J = (x_1, \ldots, x_r)$. Dans ce cas, on montre que
$$
M \otimes_P^\mathbf{L} P/J^n =
M \otimes_{A[x_1, \ldots, x_r]}^\mathbf{L}
A[x_1, \ldots, x_r]/(x_1, \ldots, x_r)^n
$$
et l'on se ramène à l'un des cas étudiés ci-dessus (ce cas est toutefois strictement plus simple, car $A[x_1, \ldots, x_r]/(x_1, \ldots, x_r)^n$ est de dimension Tor au plus $r$ pour tout $n$ ; on peut donc éviter l'étape utilisant la proposition \ref{local-cohomology-proposition-uniform-artin-rees}). Ce cas est étudié dans la démonstration de \cite[Proposition 3.5.1]{BS}.
\end{remark}








\section{Un peu d'uniformité, III}
\label{local-cohomology-section-uniform}

\noindent
Dans cette section, fixons un anneau noethérien $A$ et un idéal $I \subset A$. Notre objectif est de démontrer le lemme \ref{local-cohomology-lemma-bound-two-term-complex}, que nous utiliserons dans un chapitre ultérieur pour résoudre un problème de relèvement ; voir Algébrisation des espaces formels, lemme \ref{restricted-lemma-get-morphism-general-better}.

\medskip\noindent
Dans toute cette section, nous notons
$$
p : X \to \Spec(A)
$$
l'éclatement de $\Spec(A)$ en l'idéal $I$. Autrement dit, $X$ est le $\text{Proj}$ de l'algèbre de Rees $\bigoplus_{n \geq 0} I^n$. Nous considérons aussi le produit fibré
$$
\xymatrix{
Y \ar[r] \ar[d] & X \ar[d]^p \\
\Spec(A/I) \ar[r] & \Spec(A)
}
$$
Alors $Y$ est le diviseur exceptionnel de l'éclatement, donc un diviseur de Cartier effectif sur $X$ tel que $\mathcal{O}_X(-1) = \mathcal{O}_X(Y)$. Puisque la formation de $\text{Proj}$ commute au changement de base, on a
$$
Y = \text{Proj}(\bigoplus\nolimits_{n \geq 0} I^n/I^{n + 1}) = \text{Proj}(S)
$$
où $S = \text{Gr}_I(A) = \bigoplus_{n \geq 0} I^n/I^{n + 1}$.

\medskip\noindent
Nous notons $d = d(S) = d(\text{Gr}_I(A)) = d(\bigoplus_{n \geq 0} I^n/I^{n + 1})$ le maximum des dimensions des fibres de $p$ (et nous lui attribuons la valeur $0$ si $X = \emptyset$). Cette définition a bien un sens. En effet, on a
\begin{enumerate}
\item $d \leq t - 1$ si $I = (a_1, \ldots, a_t)$, puisqu'alors $X \subset \mathbf{P}^{t - 1}_A$,
\item $d$ est aussi la dimension maximale des fibres de $\text{Proj}(S) \to \Spec(S_0)$ lorsque $Y$ est non vide, et $d = 0$ si $Y = \emptyset$ (ou, de manière équivalente, $S = 0$, ou encore $I = A$).
\end{enumerate}
Ainsi $d$ ne dépend que de la classe d'isomorphisme de $S = \text{Gr}_I(A)$. Observons que $H^i(X, \mathcal{F}) = 0$ pour tout $\mathcal{O}_X$-module cohérent $\mathcal{F}$ et tout $i > d$, d'après Cohomologie des schémas, lemmes \ref{coherent-lemma-higher-direct-images-zero-above-dimension-fibre} et \ref{coherent-lemma-quasi-coherence-higher-direct-images-application}. Bien entendu, il en est de même pour les modules cohérents sur $Y$.

\medskip\noindent
Nous notons $q = q(S) = q(\text{Gr}_I(A)) = q(\bigoplus_{n \geq 0} I^n/I^{n + 1})$ l'entier défini comme suit. Notons que l'algèbre $S = \bigoplus_{n \geq 0} I^n/I^{n + 1}$ est un anneau gradué noethérien engendré en degré $1$ sur le degré $0$. Ainsi, d'après Cohomologie des schémas, lemmes \ref{coherent-lemma-coherent-on-proj} et \ref{coherent-lemma-recover-tail-graded-module}, on peut définir $q(S)$ comme le plus petit entier $q(S) \geq 0$ tel que, pour tout $q \geq q(S)$, on ait $H^i(Y, \mathcal{O}_Y(q)) = 0$ pour $1 \leq i \leq d$ et $H^0(Y, \mathcal{O}_Y(q)) = I^q/I^{q + 1}$. (Si $S = 0$, alors $q(S) = 0$.)

\medskip\noindent
Pour $n \geq 1$, on peut considérer le diviseur de Cartier effectif $nY$, que nous noterons $Y_n$.

\begin{lemma}
\label{local-cohomology-lemma-bound-q-and-d}
Avec $q_0 = q(S)$ et $d = d(S)$ comme ci-dessus, on a
\begin{enumerate}
\item pour $n \geq 1$, $q \geq q_0$ et $i > 0$, on a $H^i(X, \mathcal{O}_{Y_n}(q)) = 0$,
\item pour $n \geq 1$ et $q \geq q_0$, on a $H^0(X, \mathcal{O}_{Y_n}(q)) = I^q/I^{q + n}$,
\item pour $q \geq q_0$ et $i > 0$, on a $H^i(X, \mathcal{O}_X(q)) = 0$,
\item pour $q \geq q_0$, on a $H^0(X, \mathcal{O}_X(q)) = I^q$.
\end{enumerate}
\end{lemma}

\begin{proof}
Si $I = A$, alors $X$ est affine et les assertions sont triviales. On peut donc supposer, et l'on suppose, $I \not = A$. Ainsi $Y$ et $X$ sont des schémas non vides.

\medskip\noindent
Démontrons (1) et (2) par récurrence sur $n$. Le cas initial $n = 1$ est la définition de $q_0$, puisque $Y_1 = Y$. Rappelons que $\mathcal{O}_X(1) = \mathcal{O}_X(-Y)$. On a donc une suite exacte courte
$$
0 \to \mathcal{O}_{Y_n}(1) \to \mathcal{O}_{Y_{n + 1}} \to \mathcal{O}_Y \to 0
$$
Ainsi, pour $i > 0$, on trouve
$$
H^i(X, \mathcal{O}_{Y_n}(q + 1)) \to
H^i(X, \mathcal{O}_{Y_{n + 1}}(q)) \to
H^i(X, \mathcal{O}_{Y}(q))
$$
et l'annulation donnée des termes extrêmes entraîne l'annulation voulue du terme central. Pour $i = 0$, on obtient un diagramme commutatif
$$
\xymatrix{
0 \ar[r] &
I^{q + 1}/I^{q + 1 + n} \ar[d] \ar[r] &
I^q/I^{q + 1 + n} \ar[d] \ar[r] &
I^q/I^{q + 1} \ar[d] \ar[r] &
0 \\
0 \ar[r] &
H^0(X, \mathcal{O}_{Y_n}(q + 1)) \ar[r] &
H^0(X, \mathcal{O}_{Y_{n + 1}}(q)) \ar[r] &
H^0(Y, \mathcal{O}_Y(q)) \ar[r] &
0
}
$$
à lignes exactes pour $q \geq q_0$ (pour la ligne du bas, observer que le terme suivant de la suite exacte longue de cohomologie est nul pour $q \geq q_0$). Puisque $q \geq q_0$, les flèches verticales de gauche et de droite sont des isomorphismes ; celle du milieu l'est donc aussi.

\medskip\noindent
Nous omettons les démonstrations de (3) et (4), qui sont analogues. On peut en fait déduire (3) et (4) de (1) et (2) à l'aide du théorème des fonctions formelles, mais ce serait un outil disproportionné.
\end{proof}

\noindent
Introduisons une notation : étant donnés $n \geq c \geq 0$, un {\it $(A, n, c)$-module} est un $A$-module de type fini $M$ annulé par $I^n$ et qui, en tant que $A/I^n$-module, est $I^c/I^n$-projectif ; voir Compléments d'algèbre, section \ref{more-algebra-section-near-projective}.

\medskip\noindent
Nous ferons l'abus de notation suivant : étant donné un $A$-module $M$, nous noterons $p^*M$ le module quasi-cohérent obtenu par image réciproque par $p$ du module quasi-cohérent $\widetilde{M}$ sur $\Spec(A)$ associé à $M$. On a par exemple $\mathcal{O}_{Y_n} = p^*(A/I^n)$. Pour une suite exacte courte $0 \to K \to L \to M \to 0$ de $A$-modules, on obtient une suite exacte
$$
p^*K \to p^*L \to p^*M \to 0
$$
puisque $\widetilde{\ }$ est un foncteur exact et $p^*$ un foncteur exact à droite.

\begin{lemma}
\label{local-cohomology-lemma-almost-exactness}
Soit $0 \to K \to L \to M \to 0$ une suite exacte courte de $A$-modules telle que $K$ et $L$ soient annulés par $I^n$ et que $M$ soit un $(A, n, c)$-module. Alors le noyau de $p^*K \to p^*L$ est à support schématique dans $Y_c$.
\end{lemma}

\begin{proof}
Soit $\Spec(B) \subset X$ un ouvert affine. La restriction de la suite exacte à $\Spec(B)$ correspond à la suite de $B$-modules
$$
K \otimes_A B \to L \otimes_A B \to M \otimes_A B \to 0
$$
qui est isomorphe à la suite
$$
K \otimes_{A/I^n} B/I^nB \to
L \otimes_{A/I^n} B/I^nB \to
M \otimes_{A/I^n} B/I^nB \to 0
$$
Le noyau de la première application est donc l'image du module $\text{Tor}_1^{A/I^n}(M, B/I^nB)$. Rappelons que le diviseur exceptionnel $Y$ est défini par $I\mathcal{O}_X$. Il suffit donc de montrer que $\text{Tor}_1^{A/I^n}(M, B/I^nB)$ est annulé par $I^c$. Puisque la multiplication par $a \in I^c$ sur $M$ se factorise par un $A/I^n$-module libre de type fini, cela est clair.

\end{proof}

\noindent
On dispose de l'application canonique $\mathcal{O}_X \to \mathcal{O}_X(1)$, qui s'annule exactement le long de $Y$. Ainsi, pour tout $\mathcal{O}_X$-module cohérent $\mathcal{F}$, on a toujours des applications canoniques $\mathcal{F}(q) \to \mathcal{F}(q + n)$ pour tous $q \in \mathbf{Z}$ et $n \geq 0$.

\begin{lemma}
\label{local-cohomology-lemma-annihilated}
Soit $\mathcal{F}$ un $\mathcal{O}_X$-module cohérent. Alors $\mathcal{F}$ est à support schématique dans $Y_c$ si et seulement si l'application canonique $\mathcal{F} \to \mathcal{F}(c)$ est nulle.
\end{lemma}

\begin{proof}
Cela vient du fait que $\mathcal{O}_X \to \mathcal{O}_X(1)$ s'annule exactement le long de $Y$.
\end{proof}

\begin{lemma}
\label{local-cohomology-lemma-vanishing-coh-almost-projective}
Avec $q_0 = q(S)$ et $d = d(S)$ comme ci-dessus, supposons donnés des entiers $n \geq c \geq 0$, un $(A, n, c)$-module $M$, un indice $i \in \{0, 1, \ldots, d\}$ et un entier $q$. Distinguons les cas suivants :
\begin{enumerate}
\item Dans le cas $i = d \geq 1$ et $q \geq q_0$, on a $H^d(X, p^*M(q)) = 0$.
\item Dans le cas $i = d - 1 \geq 1$ et $q \geq q_0$, on a $H^{d - 1}(X, p^*M(q)) = 0$.
\item Dans le cas $d - 1 > i > 0$ et $q \geq q_0 + (d - 1 - i)c$, l'application $H^i(X, p^*M(q)) \to H^i(X, p^*M(q - (d - 1 - i)c))$ est nulle.
\item Dans le cas $i = 0$, $d \in \{0, 1\}$ et $q \geq q_0$, il existe une surjection
$$
I^qM \longrightarrow H^0(X, p^*M(q))
$$
\item Dans le cas $i = 0$, $d > 1$ et $q \geq q_0 + (d - 1)c$, l'application
$$
H^0(X, p^*M(q)) \to H^0(X, p^*M(q - (d - 1)c))
$$
a son image contenue dans celle de l'application canonique $I^{q - (d - 1)c}M \to H^0(X, p^*M(q - (d - 1)c))$.
\end{enumerate}
\end{lemma}

\begin{proof}
Soit $M$ un $(A, n, c)$-module. Choisissons une suite exacte courte
$$
0 \to K \to (A/I^n)^{\oplus r} \to M \to 0
$$
Nous utiliserons ci-dessous le fait que $K$ est un $(A, n, c)$-module ; voir Compléments d'algèbre, lemme \ref{more-algebra-lemma-ses-near-projective}. Considérons la suite exacte correspondante
$$
p^*K \to (\mathcal{O}_{Y_n})^{\oplus r} \to p^*M \to 0
$$
Décomposons-la en suites exactes courtes
$$
0 \to \mathcal{F} \to p^*K \to \mathcal{G} \to 0
\quad\text{et}\quad
0 \to \mathcal{G} \to (\mathcal{O}_{Y_n})^{\oplus r} \to p^*M \to 0
$$
D'après le lemme \ref{local-cohomology-lemma-almost-exactness}, le module cohérent $\mathcal{F}$ est à support schématique dans $Y_c$.

\medskip\noindent
Démonstration de (1). Supposons $d > 0$. Il faut prouver que $H^d(X, p^*M(q)) = 0$ pour $q \geq q_0$. Par l'annulation de la cohomologie des tordus de $\mathcal{G}$ en degrés $> d$ et la suite exacte longue de cohomologie associée à la seconde suite exacte courte ci-dessus, il suffit de prouver que $H^d(X, \mathcal{O}_{Y_n}(q)) = 0$. Cela résulte du lemme \ref{local-cohomology-lemma-bound-q-and-d}.

\medskip\noindent
Démonstration de (2). Supposons $d > 1$. Il faut prouver que $H^{d - 1}(X, p^*M(q)) = 0$ pour $q \geq q_0$. En raisonnant comme au paragraphe précédent, il suffit de montrer que $H^d(X, \mathcal{G}(q)) = 0$. À l'aide de la première suite exacte courte et de l'annulation de la cohomologie des tordus de $\mathcal{F}$ en degrés $> d$, on voit qu'il suffit de montrer que $H^d(X, p^*K(q))$ est nul, ce qui résulte de (1) et du fait que $K$ est un $(A, n, c)$-module (voir ci-dessus).

\medskip\noindent
Démonstration de (3). Soit $0 < i < d - 1$, et supposons l'assertion établie pour $i + 1$, sauf dans le cas $i = d - 2$, où l'on dispose de (2). La suite exacte longue de cohomologie associée à la seconde suite exacte courte ci-dessus fournit une injection
$$
H^i(X, p^*M(q - (d - 1 - i)c)) \subset
H^{i + 1}(X, \mathcal{G}(q - (d - 1 - i)c))
$$
puisque $q - (d - 1 - i)c \geq q_0$ entraîne l'annulation de $H^i(X, \mathcal{O}_{Y_n}(q - (d - 1 - i)c))$ (voir ci-dessus). Il suffit donc de montrer que l'application $H^{i + 1}(X, \mathcal{G}(q)) \to H^{i + 1}(X, \mathcal{G}(q - (d - 1 - i)c))$ est nulle. Pour l'étudier, considérons les morphismes de suites exactes
$$
\xymatrix{
H^{i + 1}(X, p^*K(q)) \ar[r] \ar[d] &
H^{i + 1}(X, \mathcal{G}(q)) \ar[r] \ar[d] \ar@{..>}[ld] &
H^{i + 2}(X, \mathcal{F}(q)) \ar[d] \\
H^{i + 1}(X, p^*K(q - c)) \ar[r] \ar[d] &
H^{i + 1}(X, \mathcal{G}(q - c)) \ar[r] \ar[d] &
H^{i + 2}(X, \mathcal{F}(q - c)) \\
H^{i + 1}(X, p^*K(q - (d - 1 - i)c)) \ar[r] &
H^{i + 1}(X, \mathcal{G}(q - (d - 1 - i)c))
}
$$
Puisque $\mathcal{F}$ est à support schématique dans $Y_c$, l'application canonique $\mathcal{G}(q) \to \mathcal{G}(q - c)$ se factorise par $p^*K(q - c)$ d'après le lemme \ref{local-cohomology-lemma-annihilated}. Cela donne la flèche pointillée du diagramme. (En fait, pour la démonstration, il suffit de remarquer que la flèche verticale à l'extrême droite est nulle pour obtenir la flèche pointillée comme application d'ensembles.) Il suffit donc de montrer que $H^{i + 1}(X, p^*K(q - c)) \to
H^{i + 1}(X, p^*K(q - (d - 1 - i)c))$ est nulle. Si $i = d - 2$, la source de cette flèche est nulle d'après (2), puisque $q - c \geq q_0$ et que $K$ est un $(A, n, c)$-module. Si $i < d - 2$, alors, puisque $K$ est un $(A, n, c)$-module, l'hypothèse de récurrence montre que l'application est bien nulle, car $q - c - (q - (d - 1 - i)c) = (d - 2 - i)c = (d - 1 - (i + 1))c$ et $q - c \geq q_0 + (d - 1 - (i + 1))c$. Cela achève la démonstration de (3).

\medskip\noindent
Démonstration de (4). Supposons $d \in \{0, 1\}$ et $q \geq q_0$. La première suite exacte courte donne alors une surjection $H^1(X, p^*K(q)) \to H^1(X, \mathcal{G}(q))$, dont la source est nulle d'après le cas (1). Ainsi, pour tout $q \in \mathbf{Z}$, l'application
$$
H^0(X, (\mathcal{O}_{Y_n})^{\oplus r}(q))
\longrightarrow
H^0(X, p^*M(q))
$$
est surjective. Pour $q \geq q_0$, sa source est égale à $(I^q/I^{q + n})^{\oplus r}$ d'après le lemme \ref{local-cohomology-lemma-bound-q-and-d}, ce qui prouve aisément l'assertion.

\medskip\noindent
Démonstration de (5). Supposons $d > 1$. En raisonnant comme dans la démonstration de (4), il suffit de montrer que l'image de
$$
H^0(X, p^*M(q))
\longrightarrow
H^0(X, p^*M(q - (d - 1)c))
$$
est contenue dans l'image de
$$
H^0(X, (\mathcal{O}_{Y_n})^{\oplus r}(q - (d - 1)c))
\longrightarrow
H^0(X, p^*M(q - (d - 1)c))
$$
Pour montrer cette inclusion, il suffit de montrer que, pour $\sigma \in H^0(X, p^*M(q))$ de bord $\xi \in H^1(X, \mathcal{G}(q))$, l'image de $\xi$ dans $H^1(X, \mathcal{G}(q - (d - 1)c))$ est nulle. Cela résulte des mêmes arguments que dans la démonstration de (3).
\end{proof}

\begin{remark}
\label{local-cohomology-remark-duals}
Étant donné un couple $(M, n)$ constitué d'un entier $n \geq 0$ et d'un $A/I^n$-module de type fini $M$, posons $M^\vee = \Hom_{A/I^n}(M, A/I^n)$. Étant donné un couple $(\mathcal{F}, n)$ constitué d'un entier $n$ et d'un $\mathcal{O}_{Y_n}$-module cohérent $\mathcal{F}$, posons
$$
\mathcal{F}^\vee =
\SheafHom_{\mathcal{O}_{Y_n}}(\mathcal{F}, \mathcal{O}_{Y_n})
$$
Pour $(M, n)$ comme ci-dessus, il existe une application canonique
$$
can : p^*(M^\vee) \longrightarrow (p^*M)^\vee
$$
En effet, si l'on choisit une présentation $(A/I^n)^{\oplus s} \to (A/I^n)^{\oplus r} \to M \to 0$, on obtient une présentation $\mathcal{O}_{Y_n}^{\oplus s} \to \mathcal{O}_{Y_n}^{\oplus r} \to
p^*M \to 0$. En passant aux duaux, on obtient des suites exactes
$$
0 \to M^\vee \to (A/I^n)^{\oplus r} \to (A/I^n)^{\oplus s}
$$
et
$$
0 \to (p^*M)^\vee \to
\mathcal{O}_{Y_n}^{\oplus r} \to
\mathcal{O}_{Y_n}^{\oplus s}
$$
En prenant l'image réciproque de la première suite par $p$, on trouve l'application voulue $can$. La construction de cette application est fonctorielle par rapport au $A/I^n$-module de type fini $M$. Le noyau et le conoyau de $can$ sont à support schématique dans $Y_c$ si $M$ est un $(A, n, c)$-module. En effet, dans ce cas, pour $a \in I^c$, l'application $a : M \to M$ se factorise par un $A/I^n$-module libre de type fini, pour lequel $can$ est un isomorphisme. Ainsi $a$ annule le noyau et le conoyau de $can$.
\end{remark}

\begin{lemma}
\label{local-cohomology-lemma-factor-hom}
Avec $q_0 = q(S)$ et $d = d(S)$ comme ci-dessus, soient $M$ un $(A, n, c)$-module et $\varphi : M \to I^n/I^{2n}$ une application $A$-linéaire. Supposons $n \geq \max(q_0 + (1 + d)c, (2 + d)c)$ et, si $d = 0$, supposons $n \geq q_0 + 2c$. Alors le composé
$$
M \xrightarrow{\varphi} I^n/I^{2n} \to
I^{n - (1 + d)c}/I^{2n - (1 + d)c}
$$
est de la forme $\sum a_i \psi_i$, avec $a_i \in I^c$ et $\psi_i : M \to I^{n - (2 + d)c}/I^{2n - (2 + d)c}$.
\end{lemma}

\begin{proof}
Cas $d > 1$. Puisque l'on dispose d'un système compatible d'applications $p^*(I^q) \to \mathcal{O}_X(q)$ pour $q \geq 0$, il existe des applications canoniques $p^*(I^q/I^{q + \nu}) \to \mathcal{O}_{Y_\nu}(q)$ pour $\nu \geq 0$. En utilisant cela et en prenant l'image réciproque de $\varphi$, on obtient une application
$$
\chi : p^*M \longrightarrow \mathcal{O}_{Y_n}(n)
$$
telle que le composé $M \to H^0(X, p^*M) \to H^0(X, \mathcal{O}_{Y_n}(n))$ soit l'homomorphisme donné $\varphi$ suivi de l'application $I^n/I^{2n} \to H^0(X, \mathcal{O}_{Y_n}(n))$. Puisque $\mathcal{O}_{Y_n}(n)$ est inversible sur $Y_n$, l'application linéaire $\chi$ détermine une section
$$
\sigma \in \Gamma(X, (p^*M)^\vee(n))
$$
avec les notations de la remarque \ref{local-cohomology-remark-duals}. La discussion de la remarque \ref{local-cohomology-remark-duals} montre que le conoyau et le noyau de $can : p^*(M^\vee) \to (p^*M)^\vee$ sont à support schématique dans $Y_c$. D'après le lemme \ref{local-cohomology-lemma-annihilated}, l'application $(p^*M)^\vee(n) \to (p^*M)^\vee(n - 2c)$ se factorise par $p^*(M^\vee)(n - 2c)$ ; nous omettons un petit détail. L'image de $\sigma$ dans $\Gamma(X, (p^*M)^\vee(n - 2c))$ provient donc d'un élément
$$
\sigma' \in \Gamma(X, p^*(M^\vee)(n - 2c))
$$
D'après le lemme \ref{local-cohomology-lemma-vanishing-coh-almost-projective}, point (5), le fait que $M^\vee$ soit un $(A, n, c)$-module en vertu de Compléments d'algèbre, lemme \ref{more-algebra-lemma-dual-near-projective}, et le fait que $n \geq q_0 + (1 + d)c$, donc $n - 2c \geq q_0 + (d - 1)c$, on voit que l'image de $\sigma'$ dans $H^0(X, p^*M^\vee(n - (1 + d)c))$ est l'image d'un élément $\tau$ de $I^{n - (1 + d)c}M^\vee$. Écrivons $\tau = \sum a_i \tau_i$ avec $\tau_i \in I^{n - (2 + d)c}M^\vee$ ; cela a un sens puisque $n - (2 + d)c \geq 0$. Alors $\tau_i$ détermine un homomorphisme de modules $\psi_i : M \to I^{n - (2 + d)c}/I^{2n - (2 + d)c}$ à l'aide de l'application d'évaluation $M \otimes M^\vee \to A/I^n$.

\medskip\noindent
Montrons que cela convient\footnote{Nous espérons qu'un lecteur proposera une démonstration moins laborieuse de ce fait.}. Choisissons $z \in M$ et montrons que $\varphi(z)$ et $\sum a_i \psi_i(z)$ ont la même image dans $I^{n - (1 + d)c}/I^{2n - (1 + d)c}$. D'abord, l'élément $z$ détermine une application $p^*z : \mathcal{O}_{Y_n} \to p^*M$ dont le composé avec $\chi$ est l'application $\mathcal{O}_{Y_n} \to \mathcal{O}_{Y_n}(n)$ correspondant à $\varphi(z)$ par l'application $I^n/I^{2n} \to \Gamma(\mathcal{O}_{Y_n}(n))$. Ensuite, $z$ et $p^*z$ déterminent les applications d'évaluation $e_z : M^\vee \to A/I^n$ et $e_{p^*z} : (p^*M)^\vee \to \mathcal{O}_{Y_n}$. Puisque $\chi(p^*z)$ est la section correspondant à $\varphi(z)$, on voit que $e_{p^*z}(\sigma)$ est la section correspondant à $\varphi(z)$. Ici et dans la suite, nous faisons un abus de notation : pour une application $a : \mathcal{F} \to \mathcal{G}$ de modules sur $X$, nous notons aussi $a : \mathcal{F}(t) \to \mathcal{F}(t)$ l'application correspondante entre les modules tordus. Le diagramme
$$
\xymatrix{
p^*(M^\vee) \ar[d]_{can} \ar[r]_{p^*e_z} & \mathcal{O}_{Y_n} \ar@{=}[d] \\
(p^*M)^\vee \ar[r]^{e_{p^*z}} & \mathcal{O}_{Y_n}
}
$$
est commutatif par fonctorialité de la construction $can$. Ainsi $(p^*e_z)(\sigma')$, dans $\Gamma(Y_n, \mathcal{O}_{Y_n}(n - 2c))$, est la section correspondant à l'image de $\varphi(z)$ dans $I^{n - 2c}/I^{2n - 2c}$. L'étape suivante est que $\sigma'$ s'envoie sur l'image de $\sum a_i \tau_i$ dans $H^0(X, p^*M^\vee(n - (1 + d)c))$. Cela implique que $(p^*e_z)(\sum a_i \tau_i) = \sum a_i p^*e_z(\tau_i)$, dans $\Gamma(Y_n, \mathcal{O}_{Y_n}(n - (1 + d)c))$, est la section correspondant à l'image de $\varphi(z)$ dans $I^{n - (1 + d)c}/I^{2n - (1 + d)c}$. Rappelons que $\psi_i$ est défini à partir de $\tau_i$ par une application d'évaluation. Si l'on note

$$
\chi_i : p^*M \longrightarrow \mathcal{O}_{Y_n}(n - (2 + d)c)
$$
l'application obtenue à partir de $\psi_i$, le même raisonnement que ci-dessus montre que la section correspondant à $\psi_i(z)$ est $\chi_i(p^*z) = e_{p^*z}(\chi_i) = p^*e_z(\tau_i)$. On en déduit que l'image de $\varphi(z)$ dans $\Gamma(Y_n, \mathcal{O}_{Y_n}(n - (1 + d)c))$ est égale à celle de $\sum a_i\psi_i(z)$. Puisque $n - (1 + d)c \geq q_0$, on a $\Gamma(Y_n, \mathcal{O}_{Y_n}(n - (1 + d)c)) =
I^{n - (1 + d)c}/I^{2n - (1 + d)c}$ d'après le lemme \ref{local-cohomology-lemma-bound-q-and-d}, ce qui prouve la compatibilité voulue.

\medskip\noindent
Cas $d = 1$. On raisonne comme ci-dessus pour obtenir
$$
\chi : p^*M \longrightarrow \mathcal{O}_{Y_n}(n),\quad
\sigma \in \Gamma(X, (p^*M)^\vee(n)),\quad
\sigma' \in \Gamma(X, p^*(M^\vee)(n - 2c)),
$$
puis on utilise le lemme \ref{local-cohomology-lemma-vanishing-coh-almost-projective}, point (4), pour voir que $\sigma'$ est l'image d'un élément $\tau \in I^{n - 2c}M^\vee$. Le reste de l'argument est identique.

\medskip\noindent
Cas $d = 0$. L'argument est exactement le même que dans le cas $d = 1$.
\end{proof}

\begin{lemma}
\label{local-cohomology-lemma-bound-two-term-complex}
Avec $d = d(S)$ et $q_0 = q(S)$ comme ci-dessus,
\begin{enumerate}
\item pour des entiers $n \geq c \geq 0$ tels que $n \geq \max(q_0 + (1 + d)c, (2 + d)c)$,
\item pour $K$ dans $D(A/I^n)$ tel que $H^i(K) = 0$ pour $i \not = -1, 0$ et $H^i(K)$ de type fini pour $i = -1, 0$, et tel que $\Ext^1_{A/I^c}(K, N)$ soit annulé par $I^c$ pour tous les $A/I^n$-modules de type fini $N$,
\end{enumerate}
l'application
$$
\Ext^1_{A/I^n}(K, I^n/I^{2n})
\longrightarrow
\Ext^1_{A/I^n}(K, I^{n - (1 + d)c}/I^{2n - 2(1 + d)c})
$$
est nulle.
\end{lemma}

\begin{proof}
Cas $d > 0$. Soit $K^{-1} \to K^0$ un complexe représentant $K$ comme dans Compléments d'algèbre, lemme \ref{more-algebra-lemma-ext-1-annihilated-definite}, point (5), relativement à l'idéal $I^c/I^n$ de l'anneau $A/I^n$. En particulier, $K^{-1}$ est $I^c/I^n$-projectif, puisque la multiplication par les éléments de $I^c/I^n$ se factorise même par $K^0$. D'après Compléments d'algèbre, lemme \ref{more-algebra-lemma-map-out-of-almost-free}, point (1), on a
$$
\Ext^1_{A/I^n}(K, I^n/I^{2n}) =
\Coker(\Hom_{A/I^n}(K^0, I^n/I^{2n}) \to \Hom_{A/I^n}(K^{-1}, I^n/I^{2n}))
$$
et de même pour les autres groupes Ext. Toute classe $\xi$ dans $\Ext^1_{A/I^n}(K, I^n/I^{2n})$ provient donc d'un élément $\varphi \in \Hom_{A/I^n}(K^{-1}, I^n/I^{2n})$. Notons $\varphi'$ l'image de $\varphi$ dans $\Hom_{A/I^n}(K^{-1}, I^{n - (1 + d)c}/I^{2n - (1 + d)c})$. D'après le lemme \ref{local-cohomology-lemma-factor-hom}, on peut écrire $\varphi' = \sum a_i \psi_i$ avec $a_i \in I^c$ et $\psi_i \in \Hom_{A/I^n}(M, I^{n - (2 + d)c}/I^{2n - (2 + d)c})$. Choisissons $h_i : K^0 \to K^{-1}$ tel que $a_i \text{id}_{K^{-1}} = h_i \circ d_K^{-1}$. Posons $\psi = \sum \psi_i \circ h_i : K^0 \to I^{n - (2 + d)c}/I^{2n - (2 + d)c}$. Alors $\varphi' = \psi \circ \text{d}_K^{-1}$, et l'on en déduit que $\xi$ a déjà une image nulle dans $\Ext^1_{A/I^n}(K, I^{n - (1 + d)c}/I^{2n - (1 + d)c})$, et a fortiori dans $\Ext^1_{A/I^n}(K, I^{n - (1 + d)c}/I^{2n - 2(1 + d)c})$.


\medskip\noindent
Cas $d = 0$\footnote{L'argument donné pour $d > 0$ convient, mais fournit un résultat légèrement plus faible.}. Soient $\xi$ et $\varphi$ comme ci-dessus. Considérons le diagramme
$$
\xymatrix{
K^0 \\
K^{-1} \ar[u] \ar[r]^\varphi & I^n/I^{2n} \ar[r] & I^{n - c}/I^{2n - c}
}
$$
En prenant l'image réciproque sur $X$ et en utilisant l'application $p^*(I^n/I^{2n}) \to \mathcal{O}_{Y_n}(n)$, on obtient un diagramme à flèches pleines
$$
\xymatrix{
p^*K^0 \ar@{..>}[rrd] \\
p^*K^{-1} \ar[u] \ar[r] &
\mathcal{O}_{Y_n}(n) \ar[r] &
\mathcal{O}_{Y_n}(n - c)
}
$$
On peut recouvrir $X$ par des ouverts affines $U = \Spec(B)$ tels qu'il existe $a \in I$ vérifiant la propriété suivante : $IB = aB$ et $a$ est un non-diviseur de zéro dans $B$. En effet, on peut recouvrir $X$ par les spectres d'algèbres d'éclatement affines ; voir Diviseurs, lemme \ref{divisors-lemma-blowing-up-affine}. La restriction de $\mathcal{O}_{Y_n}(n) \to \mathcal{O}_{Y_n}(n - c)$ à $U$ est isomorphe à l'application de $\mathcal{O}_U$-modules quasi-cohérents correspondant à l'homomorphisme de $B$-modules $a^c : B/a^nB \to B/a^nB$. Puisque $a^c : K^{-1} \to K^{-1}$ se factorise par $K^0$, la flèche pointillée existe sur $U$. Autrement dit, on peut trouver la flèche pointillée localement sur $X$ ! Or le faisceau des flèches pointillées complétant le diagramme est un espace principal homogène sous
$$
\mathcal{F} = \SheafHom_{\mathcal{O}_X}(
\Coker(p^*K^{-1} \to p^*K^0), \mathcal{O}_{Y_n}(n - c))
$$
qui est un $\mathcal{O}_X$-module cohérent. L'obstruction à l'existence de la flèche pointillée est donc un élément de $H^1(X, \mathcal{F})$. Ce groupe de cohomologie est nul puisque $1 > d = 0$ ; voir la discussion qui suit la définition de $d = d(S)$. On peut donc trouver une flèche pointillée $\psi : p^*K^0 \to \mathcal{O}_{Y_n}(n - c)$ complétant le diagramme. Puisque $n - c \geq q_0$, on voit que $\psi$ induit une application $K^0 \to I^{n - c}/I^{2n - c}$. Une chasse au diagramme donne $\varphi' = \psi \circ \text{d}_K^{-1}$, et la démonstration s'achève comme précédemment.
\end{proof}










% OK, start here.
%

\chapter{Géométrie algébrique et formelle}



\phantomsection
\label{algebraization-section-phantom}


\section{Introduction}
\label{algebraization-section-introduction}

\noindent
Ce chapitre poursuit l'étude de la géométrie algébrique formelle,
et notamment de la question de savoir si un objet formel est
le complété d'un objet algébrique. Une référence fondamentale est \cite{SGA2}.
Voici une liste des résultats déjà abordés
dans le projet Champs :
\begin{enumerate}
\item Le théorème des fonctions formelles, voir
Cohomologie des schémas, section \ref{coherent-section-theorem-formal-functions}.
\item Les modules formels cohérents, voir
Cohomologie des schémas, section \ref{coherent-section-coherent-formal}.
\item Le théorème d'existence de Grothendieck, voir
Cohomologie des schémas, sections \ref{coherent-section-existence},
\ref{coherent-section-existence-proper} et
\ref{coherent-section-existence-proper-support}.
\item Le théorème d'algébrisation de Grothendieck, voir
Cohomologie des schémas, section \ref{coherent-section-algebraization}.
\item Le théorème d'existence de Grothendieck dans un cadre plus général, voir
Compléments sur la platitude, sections \ref{flat-section-existence} et
\ref{flat-section-existence-derived}.
\end{enumerate}
Donnons un aperçu du contenu de ce chapitre.

\medskip\noindent
Soient $X$ un schéma et $\mathcal{I} \subset \mathcal{O}_X$
un faisceau quasi-cohérent d'idéaux de type fini. De nombreuses questions
de ce chapitre portent sur des systèmes projectifs $(\mathcal{F}_n)$
de $\mathcal{O}_X$-modules quasi-cohérents tels que
$\mathcal{F}_n = \mathcal{F}_{n + 1}/\mathcal{I}^n\mathcal{F}_{n + 1}$.
Un cas particulier important est celui où $X$ est un schéma sur un anneau
noethérien $A$ et où $\mathcal{I} = I \mathcal{O}_X$ pour un idéal $I \subset A$.
Dans Cohomologie, sections \ref{cohomology-section-inverse-systems},
\ref{cohomology-section-inverse-systems-bis} et
\ref{cohomology-section-inverse-systems-tri}, nous avons établi des résultats
généraux. Dans ce chapitre, les sections \ref{algebraization-section-ML-degree-zero} et
\ref{algebraization-section-formal-functions-principal} contiennent des résultats propres
aux schémas et aux modules quasi-cohérents. Dans la
section \ref{algebraization-section-formal-sections-cd-one}, nous démontrons
que la topologie de limite sur $\lim H^p(X, \mathcal{F}_n)$ est $I$-adique
lorsque $\text{cd}(A, I) = 1$. L'un des thèmes
de ce chapitre sera de montrer que les résultats démontrés lorsque
l'idéal est principal, $I = (f)$,
restent valables sous la seule hypothèse $\text{cd}(A, I) = 1$.

\medskip\noindent
Dans la section \ref{algebraization-section-derived-completion}, nous étudions la complétion dérivée
des modules sur un site annelé $(\mathcal{C}, \mathcal{O})$
relativement à un faisceau d'idéaux de type fini $\mathcal{I}$.
Cette section prolonge naturellement la théorie de la complétion dérivée
en algèbre commutative, telle qu'elle est exposée dans
Compléments d'algèbre, section \ref{more-algebra-section-derived-completion}.
Le premier résultat principal affirme l'existence de la complétion dérivée.
Le second affirme que, pour un morphisme $f$ de sites annelés,
la complétion dérivée commute à l'image directe dérivée :
$$
(Rf_*K)^\wedge = Rf_*(K^\wedge)
$$
si le faisceau d'idéaux à la source est localement engendré par des sections provenant
de l'idéal sur la base, voir le
lemme \ref{algebraization-lemma-pushforward-commutes-with-derived-completion}.
Soulignons que ces deux résultats principaux sont très élémentaires lorsque
les idéaux en question sont globalement de type fini, ce qui sera
le cas dans toutes les applications de cette théorie faites dans ce chapitre.
L'égalité ci-dessus est la version « correcte » du
théorème des fonctions formelles ; voir la discussion de la
section \ref{algebraization-section-formal-functions}.

\medskip\noindent
Soient $A$ un anneau noethérien et $I, J$ deux idéaux de $A$.
Soit $M$ un $A$-module de type fini.
Le sujet suivant de ce chapitre est l'application
$$
R\Gamma_J(M) \longrightarrow R\Gamma_J(M)^\wedge
$$
de la cohomologie locale de $M$ dans son complété $I$-adique dérivé.
Il se trouve que, sous des conditions de profondeur convenables,
cette application induit un isomorphisme en cohomologie dans un intervalle de degrés.
Dans la section \ref{algebraization-section-algebraization-sections-general},
nous travaillons essentiellement dans le cadre général qui vient d'être indiqué.
Dans la section \ref{algebraization-section-algebraization-punctured},
nous supposons que $A$ est un anneau local et que $J = \mathfrak m$ est un idéal maximal.
Nous conseillons au lecteur de lire cette section avant les deux autres
de cette partie du chapitre.
Enfin, dans la section \ref{algebraization-section-bootstrap}, nous exploitons
le cas local pour obtenir des résultats plus forts dans le cas général.

\medskip\noindent
Dans la partie suivante de ce chapitre, nous utilisons les résultats sur
la complétion de la cohomologie locale pour obtenir une liste non exhaustive de résultats sur
la cohomologie du complété des modules cohérents.
Plus précisément, soient $A$ un anneau noethérien, $I \subset A$
un idéal et $U \subset \Spec(A)$ un sous-schéma ouvert.
Si $\mathcal{F}$ est un $\mathcal{O}_U$-module cohérent, alors
nous pouvons considérer les applications
$$
H^i(U, \mathcal{F}) \longrightarrow \lim H^i(U, \mathcal{F}/I^n\mathcal{F})
$$
et nous demander si elles sont des isomorphismes dans un certain intervalle de degrés.
Dans la section \ref{algebraization-section-algebraization-sections},
nous étudions quelques exemples où $U$ est le spectre épointé
d'un anneau local. Dans la section \ref{algebraization-section-algebraization-sections-coherent},
nous traitons le cas général.
Dans la section \ref{algebraization-section-connected}, nous appliquons certains des résultats
obtenus à des questions de connexité en géométrie algébrique.

\medskip\noindent
Les dernières sections de ce chapitre sont consacrées à l'étude
de l'algébrisation des modules formels cohérents. Autrement dit, étant donné
un système projectif de modules cohérents $(\mathcal{F}_n)$ sur $U$
comme ci-dessus, avec
$\mathcal{F}_n = \mathcal{F}_{n + 1}/I^n\mathcal{F}_{n + 1}$
nous demandons s'il existe un $\mathcal{O}_U$-module cohérent
$\mathcal{F}$ tel que
$\mathcal{F}_n = \mathcal{F}/I^n\mathcal{F}$
pour tout $n$. Nous conseillons au lecteur de consulter la
section \ref{algebraization-section-algebraization-modules}
pour un énoncé précis de la question, un résultat général utile
(lemme \ref{algebraization-lemma-when-done}) et une application non triviale
(lemme \ref{algebraization-lemma-algebraization-principal-variant}).
La démonstration d'un résultat dépassant substantiellement ce cas
exige de développer une théorie sensiblement plus étendue.
On trouvera dans la section \ref{algebraization-section-algebraization-modules-conclusion}
les résultats de ce type les plus forts obtenus dans ce chapitre.



\section{sections formelles, I}
\label{algebraization-section-ML-degree-zero}

\noindent
Nous conseillons de consulter d'abord
Cohomologie, section \ref{cohomology-section-inverse-systems}.

\begin{lemma}
\label{algebraization-lemma-properties-system}
Soient $X$ un schéma et $\mathcal{I} \subset \mathcal{O}_X$
un faisceau quasi-cohérent d'idéaux. Soit
$$
\ldots \to \mathcal{F}_3 \to \mathcal{F}_2 \to \mathcal{F}_1
$$
un système projectif de $\mathcal{O}_X$-modules quasi-cohérents
tel que
$\mathcal{F}_n = \mathcal{F}_{n + 1}/\mathcal{I}^n\mathcal{F}_{n + 1}$.
Posons $\mathcal{F} = \lim \mathcal{F}_n$. Alors
\begin{enumerate}
\item $\mathcal{F} = R\lim \mathcal{F}_n$,
\item pour tout ouvert affine $U \subset X$, on a
$H^p(U, \mathcal{F}) = 0$ pour $p > 0$, et
\item pour tout $p$, il existe une suite exacte courte
$0 \to R^1\lim H^{p - 1}(X, \mathcal{F}_n) \to
H^p(X, \mathcal{F}) \to \lim H^p(X, \mathcal{F}_n) \to 0$.
\end{enumerate}
Si, de plus, $\mathcal{I}$ est de type fini, alors
\begin{enumerate}
\item[(4)]
$\mathcal{F}_n = \mathcal{F}/\mathcal{I}^n\mathcal{F}$, et
\item[(5)]
$\mathcal{I}^n \mathcal{F} = \lim_{m \geq n} \mathcal{I}^n\mathcal{F}_m$.
\end{enumerate}
\end{lemma}

\begin{proof}
Les assertions (1), (2) et (3) sont des faits généraux sur les systèmes projectifs de
modules quasi-cohérents dont les applications de transition sont surjectives ; voir
Catégories dérivées des schémas, lemme \ref{perfect-lemma-Rlim-quasi-coherent}
et Cohomologie, lemme \ref{cohomology-lemma-RGamma-commutes-with-Rlim}.
Supposons maintenant que $\mathcal{I}$ est de type fini.
Soit $U \subset X$ un ouvert affine. Écrivons $U = \Spec(A)$ et supposons que $\mathcal{I}|_U$
correspond à $I \subset A$. Remarquons que $I$ est un idéal de type fini.
L'équivalence entre la catégorie des $\mathcal{O}_U$-modules quasi-cohérents
et celle des $A$-modules (Schémas, lemme \ref{schemes-lemma-equivalence-quasi-coherent})
montre que $M_n = \mathcal{F}_n(U)$ est un système projectif
de $A$-modules tel que $M_n = M_{n + 1}/I^nM_{n + 1}$. Ainsi
$$
M = \mathcal{F}(U) = \lim \mathcal{F}_n(U) = \lim M_n
$$
est un module $I$-adiquement complet, avec $M/I^nM = M_n$, d'après
Algèbre, lemme \ref{algebra-lemma-limit-complete}. Cela prouve (4).
L'assertion (5) se traduit par l'égalité
$\lim_{m \geq n} I^nM/I^mM = I^nM$.
Puisque $I^mM = I^{m - n} \cdot I^nM$, cela signifie simplement que
$I^mM$ est $I$-adiquement complet. Cela résulte
d'Algèbre, lemme \ref{algebra-lemma-hathat-finitely-generated},
et du fait que $M$ est complet.
\end{proof}





\section{sections formelles, II}
\label{algebraization-section-formal-functions-principal}

\noindent
Nous conseillons de consulter d'abord
Cohomologie, sections
\ref{cohomology-section-inverse-systems-bis} et
\ref{cohomology-section-inverse-systems-tri}.

\begin{lemma}
\label{algebraization-lemma-equivalent-f-good}
Soient $X$ un schéma et $f \in \Gamma(X, \mathcal{O}_X)$. Soit
$$
\ldots \to \mathcal{F}_3 \to \mathcal{F}_2 \to \mathcal{F}_1
$$
un système projectif de $\mathcal{O}_X$-modules quasi-cohérents.
Les conditions suivantes sont équivalentes :
\begin{enumerate}
\item pour tout $n \geq 1$, l'application
$f : \mathcal{F}_{n + 1} \to \mathcal{F}_{n + 1}$ se factorise
par $\mathcal{F}_{n + 1} \to \mathcal{F}_n$ et définit une
suite exacte courte
$0 \to \mathcal{F}_n \to \mathcal{F}_{n + 1} \to \mathcal{F}_1 \to 0$,
\item pour tout $n \geq 1$, l'application
$f^n : \mathcal{F}_{n + 1} \to \mathcal{F}_{n + 1}$
se factorise par $\mathcal{F}_{n + 1} \to \mathcal{F}_1$
et définit une suite exacte courte
$0 \to \mathcal{F}_1 \to \mathcal{F}_{n + 1} \to \mathcal{F}_n \to 0$
\item il existe un $\mathcal{O}_X$-module $\mathcal{G}$
qui est $f$-divisible et tel que $\mathcal{F}_n = \mathcal{G}[f^n]$.
\item il existe un $\mathcal{O}_X$-module $\mathcal{F}$ qui est
sans $f$-torsion et tel que $\mathcal{F}_n = \mathcal{F}/f^n\mathcal{F}$.
\end{enumerate}
\end{lemma}

\begin{proof}
L'équivalence de (1), (2), (3) et l'implication (4) $\Rightarrow$ (1)
sont démontrées dans Cohomologie, lemme \ref{cohomology-lemma-equivalent-f-good}.
Supposons (1) vérifiée. Posons
$\mathcal{F} = \lim \mathcal{F}_n$. D'après le lemme \ref{algebraization-lemma-properties-system},
assertion (4), on a $\mathcal{F}_n = \mathcal{F}/f^n\mathcal{F}$.
Soient $U \subset X$ un ouvert et
$s = (s_n) \in \mathcal{F}(U) = \lim \mathcal{F}_n(U)$.
Choisissons $n \geq 1$. Si $fs = 0$, alors $s_{n + 1}$ appartient au noyau de
$\mathcal{F}_{n + 1} \to \mathcal{F}_n$ par la condition (1).
Ainsi $s_n = 0$. Comme $n$ est arbitraire, on obtient $s = 0$.
Donc $\mathcal{F}$ est sans $f$-torsion.
\end{proof}

\begin{lemma}
\label{algebraization-lemma-formal-functions-principal}
\begin{reference}
Version légèrement améliorée de \cite[Lemma 1.6]{Bhatt-local}
\end{reference}
Soient $A$ un anneau et $f \in A$. Soit $X$ un schéma sur $A$.
Soit $\mathcal{F}$ un $\mathcal{O}_X$-module quasi-cohérent.
Supposons que la suite $\mathcal{F}[f^n] = \Ker(f^n : \mathcal{F} \to \mathcal{F})$
stationne. Alors
$$
R\Gamma(X, \lim \mathcal{F}/f^n\mathcal{F}) =
R\Gamma(X, \mathcal{F})^\wedge
$$
où le membre de droite désigne le complété dérivé
relativement à l'idéal $(f) \subset A$. Par conséquent, pour
$p \in \mathbf{Z}$, on obtient un diagramme commutatif
$$
\xymatrix{
& 0 & 0 \\
0 \ar[r] &
\widehat{H^p(X, \mathcal{F})} \ar[r] \ar[u] &
\lim H^p(X, \mathcal{F}/f^n\mathcal{F}) \ar[r] \ar[u] &
T_f(H^{p + 1}(X, \mathcal{F})) \ar[r] &
0 \\
0 \ar[r] &
H^0(H^p(X, \mathcal{F})^\wedge) \ar[r] \ar[u] &
H^p(X, \lim \mathcal{F}/f^n\mathcal{F}) \ar[r] \ar[u] &
T_f(H^{p + 1}(X, \mathcal{F})) \ar[r] \ar@{=}[u] &
0 \\
&
R^1\lim H^p(X, \mathcal{F})[f^n] \ar[u] \ar[r]^\cong &
R^1\lim H^{p - 1}(X, \mathcal{F}/f^n\mathcal{F}) \ar[u] \\
& 0 \ar[u] & 0 \ar[u]
}
$$
à lignes et colonnes exactes, où
$\widehat{H^p(X, \mathcal{F})} =
\lim H^p(X, \mathcal{F})/f^n H^p(X, \mathcal{F})$
est le complété $f$-adique usuel
et où $T_f(-)$ désigne le module de Tate $f$-adique comme dans
Compléments d'algèbre, exemple
\ref{more-algebra-example-spectral-sequence-principal}.
\end{lemma}

\begin{proof}
D'après le lemme \ref{algebraization-lemma-properties-system}, on a
$\lim \mathcal{F}/f^n\mathcal{F} = R\lim \mathcal{F}/f^n \mathcal{F}$.
Le reste découle de Cohomologie, exemple
\ref{cohomology-example-formal-functions-principal}.
\end{proof}









\section{sections formelles, III}
\label{algebraization-section-formal-sections-cd-one}

\noindent
Dans cette section, nous démontrons le lemme \ref{algebraization-lemma-topology-I-adic},
qui, dans le cadre des schémas noethériens et des modules cohérents,
est l'analogue de Cohomologie, lemme \ref{cohomology-lemma-topology-I-adic-f},
lorsque l'idéal $I$ n'est pas supposé principal, mais
vérifie $\text{cd}(A, I) = 1$.

\begin{lemma}
\label{algebraization-lemma-cd-one}
Soit $I = (f_1, \ldots, f_r)$ un idéal d'un anneau noethérien $A$.
Si $\text{cd}(A, I) = 1$, il existe $c \geq 1$ et des applications
$\varphi_j : I^c \to A$ telles que $\sum f_j \varphi_j : I^c \to I$
soit l'inclusion.
\end{lemma}

\begin{proof}
Puisque $\text{cd}(A, I) = 1$, le complémentaire $U = \Spec(A) \setminus V(I)$
est affine (Cohomologie locale, lemme \ref{local-cohomology-lemma-cd-is-one}).
Écrivons $U = \Spec(B)$. Alors $IB = B$
et l'on peut écrire $1 = \sum_{j = 1, \ldots, r} f_j b_j$
pour des $b_j \in B$. D'après
Cohomologie des schémas, lemme \ref{coherent-lemma-homs-over-open},
on peut représenter les $b_j$ par des applications $\varphi_j : I^c \to A$
pour un certain $c \geq 0$. Alors $\sum f_j \varphi_j : I^c \to I \subset A$
est l'inclusion canonique, quitte à remplacer $c$ par un entier
plus grand, d'après le même lemme.
\end{proof}

\begin{lemma}
\label{algebraization-lemma-cd-one-extend}
Soit $I = (f_1, \ldots, f_r)$ un idéal d'un anneau noethérien $A$
tel que $\text{cd}(A, I) = 1$. Soient $c \geq 1$ et $\varphi_j : I^c \to A$,
$j = 1, \ldots, r$, comme dans le lemme \ref{algebraization-lemma-cd-one}.
Il existe alors un unique homomorphisme de $A$-algèbres graduées
$$
\Phi : \bigoplus\nolimits_{n \geq 0} I^{nc} \to A[T_1, \ldots, T_r]
$$
tel que $\Phi(g) = \sum \varphi_j(g) T_j$ pour $g \in I^c$.
De plus, la composée de $\Phi$ avec l'application
$A[T_1, \ldots, T_r] \to \bigoplus_{n \geq 0} I^n$,
$T_j \mapsto f_j$, est l'inclusion
$\bigoplus_{n \geq 0} I^{nc} \to \bigoplus_{n \geq 0} I^n$.
\end{lemma}

\begin{proof}
Pour tout $j$ et tout $m \geq c$, la restriction de $\varphi_j$ à
$I^m$ est une application $\varphi_j : I^m \to I^{m - c}$.
Étant donnés $j_1, \ldots, j_n \in \{1, \ldots, r\}$, montrons que la
composée
$$
\varphi_{j_1} \ldots \varphi_{j_n} :
I^{nc} \to I^{(n - 1)c} \to \ldots \to I^c \to A
$$
est indépendante de l'ordre des indices $j_1, \ldots, j_n$.
En effet, si $g = g_1 \ldots g_n$ avec $g_i \in I^c$, on
voit que
$$
(\varphi_{j_1} \ldots \varphi_{j_n})(g) =
\varphi_{j_1}(g_1) \ldots \varphi_{j_n}(g_n)
$$
est indépendant de l'ordre, car la multiplication dans $A$ est commutative.
On peut donc définir $\Phi$ en envoyant $g \in I^{nc}$ sur
$$
\Phi(g) = \sum\nolimits_{e_1 + \ldots + e_r = n}
(\varphi_1^{e_1} \circ \ldots \circ \varphi_r^{e_r})(g)
T_1^{e_1} \ldots T_r^{e_r}
$$
On vérifie directement qu'il s'agit d'un homomorphisme de $A$-algèbres
graduées possédant la propriété voulue. L'unicité est immédiate,
ainsi que la dernière propriété. Le lemme est démontré.
\end{proof}

\begin{lemma}
\label{algebraization-lemma-cd-one-extend-to-module}
Soit $I = (f_1, \ldots, f_r)$ un idéal d'un anneau noethérien $A$
tel que $\text{cd}(A, I) = 1$. Soient $c \geq 1$ et $\varphi_j : I^c \to A$,
$j = 1, \ldots, r$, comme dans le lemme \ref{algebraization-lemma-cd-one}.
Soit $A \to B$ un homomorphisme d'anneaux, avec $B$ noethérien, et soit $N$
un $B$-module de type fini. Alors, quitte à augmenter $c$
et à modifier les $\varphi_j$ en conséquence, il existe un unique
homomorphisme de $B$-modules gradués
$$
\Phi_N : \bigoplus\nolimits_{n \geq 0} I^{nc}N \to N[T_1, \ldots, T_r]
$$
tel que $\Phi_N(g x) = \Phi(g) x$ pour $g \in I^{nc}$ et $x \in N$,
où $\Phi$ est comme dans le lemme \ref{algebraization-lemma-cd-one-extend}.
La composée de $\Phi_N$ avec l'application
$N[T_1, \ldots, T_r] \to \bigoplus_{n \geq 0} I^nN$,
$T_j \mapsto f_j$, est l'inclusion
$\bigoplus_{n \geq 0} I^{nc}N \to \bigoplus_{n \geq 0} I^nN$.
\end{lemma}

\begin{proof}
L'unicité résulte clairement de la formule et de l'unicité de $\Phi$ dans le
lemme \ref{algebraization-lemma-cd-one-extend}. Considérons la $A$-algèbre noethérienne
$B' = B \oplus N$, où $N$ est un idéal de carré nul. Pour établir
l'existence de $\Phi_N$, il suffit
(d'après le lemme \ref{algebraization-lemma-cd-one}) de montrer que $\varphi_j$ se prolonge en
une application $\varphi'_j : I^cB' \to B'$, quitte à augmenter $c$
en un certain $c'$ (et à remplacer $\varphi_j$ par la composée de l'inclusion
$I^{c'} \to I^c$ avec $\varphi_j$). Rappelons que $\varphi_j$ correspond à une
section
$$
h_j \in \Gamma(\Spec(A) \setminus V(I), \mathcal{O}_{\Spec(A)})
$$
voir Cohomologie des schémas, lemme \ref{coherent-lemma-homs-over-open}.
(C'est d'ailleurs ainsi que nous avons choisi les $\varphi_j$ dans la démonstration du
lemme \ref{algebraization-lemma-cd-one}.) Utilisons ce même lemme pour représenter l'image inverse
$$
h'_j \in \Gamma(\Spec(B') \setminus V(IB'), \mathcal{O}_{\Spec(B')})
$$
de $h_j$ par une application $B'$-linéaire
$\varphi'_j : I^{c'}B' \to B'$ pour un certain $c' \geq c$.
La compatibilité avec $\varphi_j$ sera vérifiée pour $c'$
assez grand, par une nouvelle application du lemme :
il suffit en effet de la vérifier sur une famille finie de générateurs de $I^{c'}$.
Nous omettons ce détail.
\end{proof}

\begin{lemma}
\label{algebraization-lemma-cd-is-one-for-system}
Soit $I = (f_1, \ldots, f_r)$ un idéal d'un anneau noethérien $A$ tel que
$\text{cd}(A, I) = 1$. Soient $c \geq 1$ et $\varphi_j : I^c \to A$,
$j = 1, \ldots, r$, comme dans le lemme \ref{algebraization-lemma-cd-one}.
Soit $X$ un schéma noethérien sur $\Spec(A)$. Soit
$$
\ldots \to \mathcal{F}_3 \to \mathcal{F}_2 \to \mathcal{F}_1
$$
un système projectif de $\mathcal{O}_X$-modules cohérents
tel que $\mathcal{F}_n = \mathcal{F}_{n + 1}/I^n\mathcal{F}_{n + 1}$.
Posons $\mathcal{F} = \lim \mathcal{F}_n$.
Alors, quitte à augmenter $c$ et à modifier les $\varphi_j$ en conséquence,
il existe un unique homomorphisme de $\mathcal{O}_X$-modules gradués
$$
\Phi_\mathcal{F} :
\bigoplus\nolimits_{n \geq 0} I^{nc}\mathcal{F}
\longrightarrow
\mathcal{F}[T_1, \ldots, T_r]
$$
tel que $\Phi_\mathcal{F}(g s) = \Phi(g) s$ pour $g \in I^{nc}$ et
$s$ une section locale de $\mathcal{F}$, où $\Phi$ est comme dans le
lemme \ref{algebraization-lemma-cd-one-extend}. La composée de $\Phi_\mathcal{F}$
avec l'application
$\mathcal{F}[T_1, \ldots, T_r] \to \bigoplus_{n \geq 0} I^n\mathcal{F}$,
$T_j \mapsto f_j$
est l'inclusion canonique
$\bigoplus_{n \geq 0} I^{nc}\mathcal{F} \to
\bigoplus_{n \geq 0} I^n\mathcal{F}$.
\end{lemma}

\begin{proof}
L'unicité résulte immédiatement de la $\mathcal{O}_X$-linéarité
et de la condition $\Phi_\mathcal{F}(g s) = \Phi(g) s$ pour
$g \in I^{nc}$ et $s$ une section locale de $\mathcal{F}$.
On peut donc supposer $X = \Spec(B)$ affine.
Remarquons que $(\mathcal{F}_n)$ est un objet de la catégorie
$\textit{Coh}(X, I\mathcal{O}_X)$ introduite
dans Cohomologie des schémas, section \ref{coherent-section-coherent-formal}.
Soit $B' = B^\wedge$ le complété $I$-adique de $B$.
D'après Cohomologie des schémas, lemme \ref{coherent-lemma-inverse-systems-affine},
l'objet $(\mathcal{F}_n)$ correspond à un $B'$-module de type fini $N$,
au sens où $\mathcal{F}_n$ est le module cohérent
associé au $B$-module de type fini $N/I^n N$.
En appliquant le lemme \ref{algebraization-lemma-cd-one-extend-to-module}
à $I \subset A \to B'$ et à $N$,
on voit que, quitte à augmenter $c$ et à modifier les
$\varphi_j$ en conséquence, on obtient des applications uniques
$$
\Phi_N : \bigoplus\nolimits_{n \geq 0} I^{nc}N \to N[T_1, \ldots, T_r]
$$
ayant les propriétés correspondantes. Notons qu'en degré $n$, on obtient
un système projectif d'applications $N/I^mN \to \bigoplus_{e_1 + \ldots + e_r = n}
N/I^{m - nc}N \cdot T_1^{e_1} \ldots T_r^{e_r}$ pour $m \geq nc$.
En revenant au langage des faisceaux
cohérents, on voit que $\Phi_N$ correspond à un système d'applications
$$
\Phi^n_m :
I^{nc}\mathcal{F}_m
\longrightarrow
\bigoplus\nolimits_{e_1 + \ldots + e_r = n}
\mathcal{F}_{m - nc} \cdot T_1^{e_1} \ldots T_r^{e_r}
$$
pour $m \geq nc$ et $n \geq 1$ variables. En prenant la limite projective
de ces applications suivant $m$, on obtient $\Phi_\mathcal{F} = \bigoplus_n \lim_m \Phi^n_m$.
Notons que $\lim_m I^t\mathcal{F}_m = I^t \mathcal{F}$, comme on peut le voir
en évaluant sur les ouverts affines, par exemple ; mais ce fait est inutile ici, car
il existe clairement une application $I^t\mathcal{F} \to \lim_m I^t\mathcal{F}_m$.
\end{proof}

\begin{lemma}
\label{algebraization-lemma-topology-I-adic}
Soient $I$ un idéal d'un anneau noethérien $A$ et $X$ un schéma noethérien
sur $\Spec(A)$. Soit
$$
\ldots \to \mathcal{F}_3 \to \mathcal{F}_2 \to \mathcal{F}_1
$$
un système projectif de $\mathcal{O}_X$-modules cohérents
tel que $\mathcal{F}_n = \mathcal{F}_{n + 1}/I^n\mathcal{F}_{n + 1}$.
Si $\text{cd}(A, I) = 1$, alors, pour tout $p \in \mathbf{Z}$, la topologie de limite sur
$\lim H^p(X, \mathcal{F}_n)$ est $I$-adique.
\end{lemma}

\begin{proof}
Il est d'abord clair que l'image de $I^t \lim H^p(X, \mathcal{F}_n)$
dans $H^p(X, \mathcal{F}_t)$ est nulle. La topologie $I$-adique
est donc plus fine que la topologie de limite. Pour la réciproque, posons
$\mathcal{F} = \lim \mathcal{F}_n$, choisissons des générateurs $f_1, \ldots, f_r$
de $I$, choisissons $c \geq 1$, puis choisissons
$\Phi_\mathcal{F}$ comme dans le lemme \ref{algebraization-lemma-cd-is-one-for-system}.
Nous utiliserons les résultats du lemme \ref{algebraization-lemma-properties-system}
sans autre mention. En particulier, on a une suite exacte
courte
$$
0 \to R^1\lim H^{p - 1}(X, \mathcal{F}_n) \to H^p(X, \mathcal{F})
\to \lim H^p(X, \mathcal{F}_n) \to 0
$$
On peut donc relever tout élément $\xi$ de $\lim H^p(X, \mathcal{F}_n)$
en un élément $\xi' \in H^p(X, \mathcal{F})$. Supposons que $\xi$ ait une image nulle
dans $H^p(X, \mathcal{F}_{nc})$ pour un certain $n$ ; autrement
dit, supposons que $\xi$ soit « petit » pour la topologie de limite. On a une
suite exacte courte
$$
0 \to I^{nc}\mathcal{F} \to \mathcal{F} \to \mathcal{F}_{nc} \to 0
$$
et l'hypothèse signifie donc que l'on peut relever $\xi'$ en un élément
$\xi'' \in H^p(X, I^{nc}\mathcal{F})$. En appliquant $\Phi_\mathcal{F}$,
on obtient
$$
\Phi_\mathcal{F}(\xi'') = \sum\nolimits_{e_1 + \ldots + e_r = n}
\xi'_{e_1, \ldots, e_r} \cdot T_1^{e_1} \ldots T_r^{e_r}
$$
pour certains $\xi'_{e_1, \ldots, e_r} \in H^p(X, \mathcal{F})$.
En notant $\xi_{e_1, \ldots, e_r} \in \lim H^p(X, \mathcal{F}_n)$
leurs images et en utilisant la dernière assertion du
lemme \ref{algebraization-lemma-cd-is-one-for-system},
on conclut que
$$
\xi = \sum f_1^{e_1} \ldots f_r^{e_r} \xi_{e_1, \ldots, e_r}
$$
appartient à $I^n \lim H^p(X, \mathcal{F}_n)$, comme voulu.
\end{proof}

\begin{example}
\label{algebraization-example-not-I-adic}
Soit $k$ un corps. Posons $A = k[x, y][[s, t]]/(xs - yt)$.
Posons $I = (s, t)$ et $\mathfrak a = (x, y, s, t)$.
Posons $X = \Spec(A) - V(\mathfrak a)$ et
$\mathcal{F}_n = \mathcal{O}_X/I^n\mathcal{O}_X$.
Remarquons que la fonction rationnelle
$$
g = \frac{t}{x} = \frac{s}{y}
$$
est régulière dans un voisinage ouvert $V \subset X$ de
$V(I\mathcal{O}_X)$. Toute puissance $g^e$ définit donc une section
$g^e \in M = \lim H^0(X, \mathcal{F}_n)$. Remarquons que
$g^e \to 0$ lorsque $e \to \infty$ pour la topologie de limite sur $M$,
puisque l'image de $g^e$ dans $\mathcal{F}_e$ est nulle.
En revanche, $g^e \not \in IM$ pour tout $e$,
comme le lecteur pourra le voir en calculant $H^0(U, \mathcal{F}_n)$ ;
nous omettons ce calcul. Remarquons que $\text{cd}(A, I) = 2$.
Le résultat du lemme \ref{algebraization-lemma-topology-I-adic} est donc optimal.
\end{example}








\section{Conditions de Mittag-Leffler}
\label{algebraization-section-ML}

\noindent
Lorsque l'on prend la cohomologie locale relativement à l'idéal maximal
d'un anneau local noethérien, la condition de Mittag-Leffler est souvent
automatiquement satisfaite. Il en va donc de même des groupes de cohomologie
supérieure d'un système projectif de faisceaux cohérents à applications de transition
surjectives sur le spectre épointé.

\begin{lemma}
\label{algebraization-lemma-descending-chain}
Soit $(A, \mathfrak m)$ un anneau local noethérien.
\begin{enumerate}
\item Soit $M$ un $A$-module de type fini. Alors le $A$-module
$H^i_\mathfrak m(M)$ satisfait à la condition des chaînes descendantes
pour tout $i$.
\item Soit $U = \Spec(A) \setminus \{\mathfrak m\}$ le
spectre épointé de $A$.
Soit $\mathcal{F}$ un $\mathcal{O}_U$-module cohérent.
Alors le $A$-module $H^i(U, \mathcal{F})$
satisfait à la condition des chaînes descendantes pour $i > 0$.
\end{enumerate}
\end{lemma}

\begin{proof}
Démontrons (1) par récurrence sur la dimension du support de $M$.
L'énoncé vaut si $M = 0$ ; on peut donc supposer, et l'on suppose, $M$ non nul.

\medskip\noindent
Initialisation de la récurrence.
Si $\dim(\text{Supp}(M)) = 0$, le support
de $M$ est $\{\mathfrak m\}$, et l'on voit que $H^0_\mathfrak m(M) = M$
et $H^i_\mathfrak m(M) = 0$ pour $i > 0$, comme il résulte de la
construction de la cohomologie locale ; voir
Complexes dualisants, section \ref{dualizing-section-local-cohomology}.
Puisque $M$ est de longueur finie (Algèbre, lemme \ref{algebra-lemma-length-finite}),
il satisfait à la condition des chaînes descendantes.

\medskip\noindent
Étape de récurrence. Supposons $\dim(\text{Supp}(M)) > 0$.
D'après le cas initial, le module de type fini $H^0_\mathfrak m(M) \subset M$
satisfait à la condition des chaînes descendantes.
D'après Complexes dualisants, lemme \ref{dualizing-lemma-divide-by-torsion},
on peut remplacer $M$ par $M/H^0_\mathfrak m(M)$.
Alors $H^0_\mathfrak m(M) = 0$, c'est-à-dire que $M$ est de profondeur $\geq 1$ ; voir
Complexes dualisants, lemme \ref{dualizing-lemma-depth}.
Choisissons $x \in \mathfrak m$ tel que $x : M \to M$ soit injective.
D'après Algèbre, lemme \ref{algebra-lemma-one-equation-module}, on a
$\dim(\text{Supp}(M/xM)) = \dim(\text{Supp}(M)) - 1$ et l'hypothèse
de récurrence s'applique. Fixons un indice $i$ et considérons la
suite exacte
$$
H^{i - 1}_\mathfrak m(M/xM) \to H^i_\mathfrak m(M) \xrightarrow{x}
H^i_\mathfrak m(M)
$$
issue de la suite exacte courte $0 \to M \xrightarrow{x} M \to M/xM \to 0$.
Il s'ensuit que le sous-module de $x$-torsion $H^i_\mathfrak m(M)[x]$
est quotient d'un module satisfaisant à la condition des chaînes descendantes,
et satisfait donc lui-même à cette condition. Ainsi, le sous-module
de $\mathfrak m$-torsion $H^i_\mathfrak m(M)[\mathfrak m]$ satisfait
à la condition des chaînes descendantes (et est donc de dimension finie
sur $A/\mathfrak m$). On en conclut que le module de torsion par les puissances de $\mathfrak m$
$H^i_\mathfrak m(M)$ satisfait à la condition des chaînes descendantes
d'après Complexes dualisants, lemme
\ref{dualizing-lemma-describe-categories}.

\medskip\noindent
L'assertion (2) découle de (1) par Cohomologie locale,
lemme \ref{local-cohomology-lemma-finiteness-pushforwards-and-H1-local}.
\end{proof}

\begin{lemma}
\label{algebraization-lemma-ML-local}
Soit $(A, \mathfrak m)$ un anneau local noethérien.
\begin{enumerate}
\item Soit $(M_n)$ un système projectif de $A$-modules de type fini. Alors le
système projectif $H^i_\mathfrak m(M_n)$ satisfait à la condition de Mittag-Leffler
pour tout $i$.
\item Soit $U = \Spec(A) \setminus \{\mathfrak m\}$ le
spectre épointé de $A$.
Soit $\mathcal{F}_n$ un système projectif de
$\mathcal{O}_U$-modules cohérents.
Alors le système projectif $H^i(U, \mathcal{F}_n)$
satisfait à la condition de Mittag-Leffler pour $i > 0$.
\end{enumerate}
\end{lemma}

\begin{proof}
Cela résulte immédiatement du lemme \ref{algebraization-lemma-descending-chain}.
\end{proof}

\begin{lemma}
\label{algebraization-lemma-terrific}
Soit $(A, \mathfrak m)$ un anneau local noethérien.
Soit $(M_n)$ un système projectif de $A$-modules de type fini.
Soit $M \to \lim M_n$ une application, où $M$ est un $A$-module de type fini,
telle que, pour un certain $i$, l'application
$H^i_\mathfrak m(M) \to \lim H^i_\mathfrak m(M_n)$
soit un isomorphisme.
Alors le système projectif $H^i_\mathfrak m(M_n)$
est essentiellement constant, de valeur $H^i_\mathfrak m(M)$.
\end{lemma}

\begin{proof}
D'après le lemme \ref{algebraization-lemma-ML-local}, le système projectif $H^i_\mathfrak m(M_n)$
satisfait à la condition de Mittag-Leffler. Soit $E_n \subset H^i_\mathfrak m(M_n)$
l'image de $H^i_\mathfrak m(M_{n'})$ pour $n' \gg n$.
Alors $(E_n)$ est un système projectif à applications de transition surjectives
et $H^i_\mathfrak m(M) = \lim E_n$. Puisque $H^i_\mathfrak m(M)$
satisfait à la condition des chaînes descendantes d'après le
lemme \ref{algebraization-lemma-descending-chain},
il ne peut y avoir qu'un nombre fini de noyaux non triviaux
parmi ceux des surjections $H^i_\mathfrak m(M) \to E_n$.
Ainsi $E_n \to E_{n - 1}$ est un isomorphisme pour tout $n \gg 0$,
comme voulu.
\end{proof}

\begin{lemma}
\label{algebraization-lemma-local-cohomology-derived-completion}
Soit $(A, \mathfrak m)$ un anneau local noethérien.
Soient $I \subset A$ un idéal et $M$ un $A$-module de type fini.
Alors
$$
H^i(R\Gamma_\mathfrak m(M)^\wedge) = \lim H^i_\mathfrak m(M/I^nM)
$$
pour tout $i$, où $R\Gamma_\mathfrak m(M)^\wedge$ désigne
le complété $I$-adique dérivé.
\end{lemma}

\begin{proof}
Appliquons Complexes dualisants, lemme \ref{dualizing-lemma-completion-local},
et le lemme \ref{algebraization-lemma-ML-local} pour obtenir l'annulation des termes $R^1\lim$.
\end{proof}








\section{Complétion dérivée sur un site annelé}
\label{algebraization-section-derived-completion}

\noindent
Nous recommandons au lecteur de sauter cette section à la première lecture.

\medskip\noindent
La version algébrique de ce qui suit figure dans
Compléments d'algèbre, section \ref{more-algebra-section-derived-completion}.
Soit $\mathcal{O}$ un faisceau d'anneaux sur un site $\mathcal{C}$.
Soit $f$ une section globale de $\mathcal{O}$. Notons
$\mathcal{O}_f$ le faisceau associé au préfaisceau de localisés
$U \mapsto \mathcal{O}(U)_f$.

\begin{lemma}
\label{algebraization-lemma-map-twice-localize}
Soit $(\mathcal{C}, \mathcal{O})$ un site annelé. Soit $f$ une section
globale de $\mathcal{O}$.
\begin{enumerate}
\item Pour $L, N \in D(\mathcal{O}_f)$, on a
$R\SheafHom_\mathcal{O}(L, N) = R\SheafHom_{\mathcal{O}_f}(L, N)$.
En particulier, les deux structures de $\mathcal{O}_f$-module sur
$R\SheafHom_\mathcal{O}(L, N)$ coïncident.
\item Pour $K \in D(\mathcal{O})$ et
$L \in D(\mathcal{O}_f)$, on a
$$
R\SheafHom_\mathcal{O}(L, K) =
R\SheafHom_{\mathcal{O}_f}(L, R\SheafHom_\mathcal{O}(\mathcal{O}_f, K))
$$
En particulier
$R\SheafHom_\mathcal{O}(\mathcal{O}_f,
R\SheafHom_\mathcal{O}(\mathcal{O}_f, K)) =
R\SheafHom_\mathcal{O}(\mathcal{O}_f, K)$.
\item Si $g$ est une seconde section
globale de $\mathcal{O}$, alors
$$
R\SheafHom_\mathcal{O}(\mathcal{O}_f, R\SheafHom_\mathcal{O}(\mathcal{O}_g, K))
= R\SheafHom_\mathcal{O}(\mathcal{O}_{gf}, K).
$$
\end{enumerate}
\end{lemma}

\begin{proof}
Démonstration de (1). Soit $\mathcal{J}^\bullet$ un complexe K-injectif de
$\mathcal{O}_f$-modules représentant $N$. D'après Cohomologie sur les sites, lemme
\ref{sites-cohomology-lemma-K-injective-flat}, il s'ensuit que
$\mathcal{J}^\bullet$ est aussi un complexe K-injectif de
$\mathcal{O}$-modules. Soit $\mathcal{F}^\bullet$ un complexe de
$\mathcal{O}_f$-modules représentant $L$. Alors
$$
R\SheafHom_\mathcal{O}(L, N) =
R\SheafHom_\mathcal{O}(\mathcal{F}^\bullet, \mathcal{J}^\bullet) =
R\SheafHom_{\mathcal{O}_f}(\mathcal{F}^\bullet, \mathcal{J}^\bullet)
$$
d'après
Modules sur les sites, lemme \ref{sites-modules-lemma-epimorphism-modules},
puisque $\mathcal{J}^\bullet$ est un complexe K-injectif de $\mathcal{O}$-modules
et de $\mathcal{O}_f$-modules.

\medskip\noindent
Démonstration de (2). Soit $\mathcal{I}^\bullet$ un complexe K-injectif de
$\mathcal{O}$-modules représentant $K$.
Alors $R\SheafHom_\mathcal{O}(\mathcal{O}_f, K)$ est représenté par
$\SheafHom_\mathcal{O}(\mathcal{O}_f, \mathcal{I}^\bullet)$, qui est
un complexe K-injectif de $\mathcal{O}_f$-modules et de
$\mathcal{O}$-modules d'après
Cohomologie sur les sites, lemmes \ref{sites-cohomology-lemma-hom-K-injective} et
\ref{sites-cohomology-lemma-K-injective-flat}.
Soit $\mathcal{F}^\bullet$ un complexe de $\mathcal{O}_f$-modules
représentant $L$. Alors
$$
R\SheafHom_\mathcal{O}(L, K) =
R\SheafHom_\mathcal{O}(\mathcal{F}^\bullet, \mathcal{I}^\bullet) =
R\SheafHom_{\mathcal{O}_f}(\mathcal{F}^\bullet,
\SheafHom_\mathcal{O}(\mathcal{O}_f, \mathcal{I}^\bullet))
$$
d'après Modules sur les sites, lemme \ref{sites-modules-lemma-adjoint-hom-restrict},
et puisque $\SheafHom_\mathcal{O}(\mathcal{O}_f, \mathcal{I}^\bullet)$ est un
complexe K-injectif de $\mathcal{O}_f$-modules.

\medskip\noindent
Démonstration de (3). Cela résulte du fait que
$R\SheafHom_\mathcal{O}(\mathcal{O}_g, \mathcal{I}^\bullet)$
est K-injectif comme complexe de $\mathcal{O}$-modules et du fait que
$\SheafHom_\mathcal{O}(\mathcal{O}_f,
\SheafHom_\mathcal{O}(\mathcal{O}_g, \mathcal{H})) = 
\SheafHom_\mathcal{O}(\mathcal{O}_{gf}, \mathcal{H})$
pour tout faisceau de $\mathcal{O}$-modules $\mathcal{H}$.
\end{proof}

\noindent
Soit $K \in D(\mathcal{O})$. Notons
$T(K, f)$ une limite dérivée (Catégories dérivées, définition
\ref{derived-definition-derived-limit}) du système projectif
$$
\ldots \to K \xrightarrow{f} K \xrightarrow{f} K
$$
dans $D(\mathcal{O})$.

\begin{lemma}
\label{algebraization-lemma-hom-from-Af}
Soit $(\mathcal{C}, \mathcal{O})$ un site annelé. Soit $f$ une section
globale de $\mathcal{O}$. Soit $K \in D(\mathcal{O})$.
Les conditions suivantes sont équivalentes :
\begin{enumerate}
\item $R\SheafHom_\mathcal{O}(\mathcal{O}_f, K) = 0$,
\item $R\SheafHom_\mathcal{O}(L, K) = 0$ pour tout $L$ dans $D(\mathcal{O}_f)$,
\item $T(K, f) = 0$.
\end{enumerate}
\end{lemma}

\begin{proof}
Il est clair que (2) implique (1). L'implication (1) $\Rightarrow$ (2)
résulte du lemme \ref{algebraization-lemma-map-twice-localize}.
Une résolution libre du $\mathcal{O}$-module $\mathcal{O}_f$ est donnée par
$$
0 \to \bigoplus\nolimits_{n \in \mathbf{N}} \mathcal{O} \to
\bigoplus\nolimits_{n \in \mathbf{N}} \mathcal{O}
\to \mathcal{O}_f \to 0
$$
où la première application envoie une section locale $(x_0, x_1, \ldots)$ sur
$(x_0, x_1 - fx_0, x_2 - fx_1, \ldots)$, et la seconde envoie
$(x_0, x_1, \ldots)$ sur $x_0 + x_1/f + x_2/f^2 + \ldots$.
En appliquant $\SheafHom_\mathcal{O}(-, \mathcal{I}^\bullet)$,
où $\mathcal{I}^\bullet$ est un complexe K-injectif de $\mathcal{O}$-modules
représentant $K$, on obtient une suite exacte courte de complexes
$$
0 \to \SheafHom_\mathcal{O}(\mathcal{O}_f, \mathcal{I}^\bullet) \to
\prod \mathcal{I}^\bullet \to \prod \mathcal{I}^\bullet \to 0
$$
puisque $\mathcal{I}^n$ est un $\mathcal{O}$-module injectif.
Les produits sont des produits dans $D(\mathcal{O})$ ; voir
Injectifs, lemme \ref{injectives-lemma-derived-products}.
Cela signifie que l'objet $T(K, f)$ est un représentant de
$R\SheafHom_\mathcal{O}(\mathcal{O}_f, K)$ dans $D(\mathcal{O})$.
D'où l'équivalence de (1) et (3).
\end{proof}

\begin{lemma}
\label{algebraization-lemma-ideal-of-elements-complete-wrt}
Soit $(\mathcal{C}, \mathcal{O})$ un site annelé. Soit $K \in D(\mathcal{O})$.
La règle qui associe à $U$ l'ensemble $\mathcal{I}(U)$
des sections $f \in \mathcal{O}(U)$ telles que $T(K|_U, f) = 0$
est un faisceau d'idéaux de $\mathcal{O}$.
\end{lemma}

\begin{proof}
Nous utiliserons les résultats du lemme \ref{algebraization-lemma-hom-from-Af} sans autre
mention. Si $f \in \mathcal{I}(U)$ et $g \in \mathcal{O}(U)$, alors
$\mathcal{O}_{U, gf}$ est un $\mathcal{O}_{U, f}$-module,
donc $R\SheafHom_\mathcal{O}(\mathcal{O}_{U, gf}, K|_U) = 0$, d'où
$gf \in \mathcal{I}(U)$. Supposons $f, g \in \mathcal{O}(U)$.
On a alors une suite exacte courte
$$
0 \to \mathcal{O}_{U, f + g} \to
\mathcal{O}_{U, f(f + g)} \oplus \mathcal{O}_{U, g(f + g)} \to
\mathcal{O}_{U, gf(f + g)} \to 0
$$
puisque $f, g$ engendrent l'idéal unité de $\mathcal{O}(U)_{f + g}$.
Cela résulte
d'Algèbre, lemme \ref{algebra-lemma-standard-covering},
et du fait immédiat que la dernière flèche est surjective.
Puisque $R\SheafHom_\mathcal{O}( - , K|_U)$ est un foncteur exact
entre catégories triangulées, l'annulation de
$R\SheafHom_{\mathcal{O}_U}(\mathcal{O}_{U, f(f + g)}, K|_U)$,
$R\SheafHom_{\mathcal{O}_U}(\mathcal{O}_{U, g(f + g)}, K|_U)$ et
$R\SheafHom_{\mathcal{O}_U}(\mathcal{O}_{U, gf(f + g)}, K|_U)$,
entraîne l'annulation de
$R\SheafHom_{\mathcal{O}_U}(\mathcal{O}_{U, f + g}, K|_U)$.
Nous omettons la vérification de la condition de faisceau.
\end{proof}

\noindent
On peut donner la définition suivante pour tout site annelé.

\begin{definition}
\label{algebraization-definition-derived-complete}
Soit $(\mathcal{C}, \mathcal{O})$ un site annelé.
Soit $\mathcal{I} \subset \mathcal{O}$ un faisceau d'idéaux.
Soit $K \in D(\mathcal{O})$. On dit que $K$ est
{\it complet au sens dérivé relativement à $\mathcal{I}$}
si, pour tout objet $U$ de $\mathcal{C}$ et tout $f \in \mathcal{I}(U)$,
l'objet $T(K|_U, f)$ de $D(\mathcal{O}_U)$ est nul.
\end{definition}

\noindent
Il est clair que la sous-catégorie pleine
$D_{comp}(\mathcal{O}) = D_{comp}(\mathcal{O}, \mathcal{I}) \subset
D(\mathcal{O})$ formée des objets complets au sens dérivé
est une sous-catégorie triangulée saturée ; voir
Catégories dérivées, définitions
\ref{derived-definition-triangulated-subcategory} et
\ref{derived-definition-saturated}. Cette sous-catégorie est stable
par produits et limites homotopiques dans $D(\mathcal{O})$.
En revanche, elle n'est en général pas stable par sommes directes dénombrables.

\begin{lemma}
\label{algebraization-lemma-derived-complete-internal-hom}
Soit $(\mathcal{C}, \mathcal{O})$ un site annelé.
Soit $\mathcal{I} \subset \mathcal{O}$ un faisceau d'idéaux.
Si $K \in D(\mathcal{O})$ et $L \in D_{comp}(\mathcal{O})$, alors
$R\SheafHom_\mathcal{O}(K, L) \in D_{comp}(\mathcal{O})$.
\end{lemma}

\begin{proof}
Soient $U$ un objet de $\mathcal{C}$ et $f \in \mathcal{I}(U)$.
Rappelons que
$$
\Hom_{D(\mathcal{O}_U)}(\mathcal{O}_{U, f}, R\SheafHom_\mathcal{O}(K, L)|_U)
=
\Hom_{D(\mathcal{O}_U)}(
K|_U \otimes_{\mathcal{O}_U}^\mathbf{L} \mathcal{O}_{U, f}, L|_U)
$$
d'après Cohomologie sur les sites, lemme \ref{sites-cohomology-lemma-internal-hom}.
Le membre de droite est nul par le lemme \ref{algebraization-lemma-hom-from-Af}
et la relation entre Hom interne et Hom ; voir
Cohomologie sur les sites, lemme \ref{sites-cohomology-lemma-section-RHom-over-U}.
La même annulation vaut pour tout $U'/U$. Ainsi, l'objet
$R\SheafHom_{\mathcal{O}_U}(\mathcal{O}_{U, f},
R\SheafHom_\mathcal{O}(K, L)|_U)$ de $D(\mathcal{O}_U)$ a un faisceau de cohomologie
de degré $0$ nul (loc. cit.). Il en va de même des autres
faisceaux de cohomologie, autrement dit $R\SheafHom_{\mathcal{O}_U}(\mathcal{O}_{U, f},
R\SheafHom_\mathcal{O}(K, L)|_U)$ est nul dans $D(\mathcal{O}_U)$.
On conclut par le lemme \ref{algebraization-lemma-hom-from-Af}.
\end{proof}

\begin{lemma}
\label{algebraization-lemma-restriction-derived-complete}
Soit $\mathcal{C}$ un site. Soit $\mathcal{O} \to \mathcal{O}'$
un homomorphisme de faisceaux d'anneaux. Soit $\mathcal{I} \subset \mathcal{O}$
un faisceau d'idéaux. L'image réciproque de $D_{comp}(\mathcal{O}, \mathcal{I})$
par le foncteur de restriction $D(\mathcal{O}') \to D(\mathcal{O})$ est
$D_{comp}(\mathcal{O}', \mathcal{I}\mathcal{O}')$.
\end{lemma}

\begin{proof}
En utilisant le lemme \ref{algebraization-lemma-ideal-of-elements-complete-wrt},
on voit que $K' \in D(\mathcal{O}')$ appartient à
$D_{comp}(\mathcal{O}', \mathcal{I}\mathcal{O}')$
si et seulement si $T(K'|_U, f)$ est nul pour toute section locale
$f \in \mathcal{I}(U)$. Remarquons que les faisceaux de cohomologie de
$T(K'|_U, f)$ se calculent dans la catégorie des faisceaux abéliens ;
il importe donc peu que l'on considère $f$ comme une section de
$\mathcal{O}$ ou que l'on prenne l'image de $f$ comme section de $\mathcal{O}'$.
Le lemme résulte immédiatement de cette remarque et de la
définition des objets complets au sens dérivé.
\end{proof}

\begin{lemma}
\label{algebraization-lemma-pushforward-derived-complete}
Soit $f : (\Sh(\mathcal{D}), \mathcal{O}') \to (\Sh(\mathcal{C}), \mathcal{O})$
un morphisme de topos annelés. Soient $\mathcal{I} \subset \mathcal{O}$
et $\mathcal{I}' \subset \mathcal{O}'$ des faisceaux d'idéaux tels
que $f^\sharp$ envoie $f^{-1}\mathcal{I}$ dans $\mathcal{I}'$.
Alors $Rf_*$ envoie $D_{comp}(\mathcal{O}', \mathcal{I}')$
dans $D_{comp}(\mathcal{O}, \mathcal{I})$.
\end{lemma}

\begin{proof}
On peut supposer que $f$ est donné par un morphisme de sites annelés correspondant
à un foncteur continu $\mathcal{C} \to \mathcal{D}$
(Modules sur les sites, lemme
\ref{sites-modules-lemma-morphism-ringed-topoi-comes-from-morphism-ringed-sites}
).
Soit $U$ un objet de $\mathcal{C}$ et soit $g$ une section de
$\mathcal{I}$ sur $U$. Il faut montrer que
$\Hom_{D(\mathcal{O}_U)}(\mathcal{O}_{U, g}, Rf_*K|_U) = 0$
lorsque $K$ est complet au sens dérivé relativement à $\mathcal{I}'$.
En effet, d'après Cohomologie sur les sites, lemme
\ref{sites-cohomology-lemma-section-RHom-over-U}
cette annulation, appliquée à tous les objets au-dessus de $U$ et à tous les décalages de $K$,
entraînera que $R\SheafHom_{\mathcal{O}_U}(\mathcal{O}_{U, g}, Rf_*K|_U)$
est nul, donc que $T(Rf_*K|_U, g)$ est nul
(lemme \ref{algebraization-lemma-hom-from-Af}), ce qui est l'assertion à démontrer
(définition \ref{algebraization-definition-derived-complete}).
Soit $V$ dans $\mathcal{D}$ l'image de $U$. Alors
$$
\Hom_{D(\mathcal{O}_U)}(\mathcal{O}_{U, g}, Rf_*K|_U) =
\Hom_{D(\mathcal{O}'_V)}(\mathcal{O}'_{V, g'}, K|_V) = 0
$$
où $g' = f^\sharp(g) \in \mathcal{I}'(V)$. La seconde égalité
résulte de ce que $K$ est complet au sens dérivé, et la première de ce que
l'image inverse dérivée de $\mathcal{O}_{U, g}$ est $\mathcal{O}'_{V, g'}$,
ainsi que de
Cohomologie sur les sites, lemme \ref{sites-cohomology-lemma-adjoint}.
\end{proof}

\noindent
Le lemme suivant donne le cas le plus simple où l'on dispose de la complétion dérivée.

\begin{lemma}
\label{algebraization-lemma-derived-completion}
Soit $(\mathcal{C}, \mathcal{O})$ un site annelé. Soient $f_1, \ldots, f_r$
des sections globales de $\mathcal{O}$. Soit $\mathcal{I} \subset \mathcal{O}$
le faisceau d'idéaux engendré par $f_1, \ldots, f_r$.
Alors le foncteur d'inclusion $D_{comp}(\mathcal{O}) \to D(\mathcal{O})$
admet un adjoint à gauche ; autrement dit, pour tout objet $K$ de $D(\mathcal{O})$,
il existe une application $K \to K^\wedge$, avec $K^\wedge$ dans $D_{comp}(\mathcal{O})$,
telle que l'application
$$
\Hom_{D(\mathcal{O})}(K^\wedge, E) \longrightarrow \Hom_{D(\mathcal{O})}(K, E)
$$
soit bijective dès que $E$ appartient à $D_{comp}(\mathcal{O})$. En fait,
on a
$$
K^\wedge =
R\SheafHom_\mathcal{O}
(\mathcal{O} \to \prod\nolimits_{i_0} \mathcal{O}_{f_{i_0}} \to
\prod\nolimits_{i_0 < i_1} \mathcal{O}_{f_{i_0}f_{i_1}} \to
\ldots \to \mathcal{O}_{f_1\ldots f_r}, K)
$$
fonctoriellement en $K$.
\end{lemma}

\begin{proof}
Définissons $K^\wedge$ par la dernière formule affichée du lemme.
On a un morphisme de complexes
$$
(\mathcal{O} \to \prod\nolimits_{i_0} \mathcal{O}_{f_{i_0}} \to
\prod\nolimits_{i_0 < i_1} \mathcal{O}_{f_{i_0}f_{i_1}} \to
\ldots \to \mathcal{O}_{f_1\ldots f_r}) \longrightarrow \mathcal{O}
$$
qui induit une application $K \to K^\wedge$. Il suffit de montrer que
$K^\wedge$ est complet au sens dérivé et que $K \to K^\wedge$ est un
isomorphisme si $K$ est complet au sens dérivé.

\medskip\noindent
Soit $f$ une section globale de $\mathcal{O}$.
D'après le lemme \ref{algebraization-lemma-map-twice-localize}, l'objet
$R\SheafHom_\mathcal{O}(\mathcal{O}_f, K^\wedge)$
est égal à
$$
R\SheafHom_\mathcal{O}(
(\mathcal{O}_f \to \prod\nolimits_{i_0} \mathcal{O}_{ff_{i_0}} \to
\prod\nolimits_{i_0 < i_1} \mathcal{O}_{ff_{i_0}f_{i_1}} \to
\ldots \to \mathcal{O}_{ff_1\ldots f_r}), K)
$$
Si $f = f_i$ pour un certain $i$, alors $f_1, \ldots, f_r$ engendrent
l'idéal unité de $\mathcal{O}_f$ ; le complexe de
{\v C}ech alterné augmenté
$$
\mathcal{O}_f \to \prod\nolimits_{i_0} \mathcal{O}_{ff_{i_0}} \to
\prod\nolimits_{i_0 < i_1} \mathcal{O}_{ff_{i_0}f_{i_1}} \to
\ldots \to \mathcal{O}_{ff_1\ldots f_r}
$$
est donc nul (et même homotope à zéro). On voit ainsi que $K^\wedge$
est complet au sens dérivé.

\medskip\noindent
Si $K$ est complet au sens dérivé, alors $R\SheafHom_\mathcal{O}(\mathcal{O}_f, K)$
est nul pour tout $f = f_{i_0} \ldots f_{i_p}$, $p \geq 0$. Ainsi
$K \to K^\wedge$ est un isomorphisme dans $D(\mathcal{O})$.
\end{proof}

\noindent
Expliquons maintenant en quoi la complétion dérivée est une complétion.

\begin{lemma}
\label{algebraization-lemma-derived-completion-koszul}
Soit $(\mathcal{C}, \mathcal{O})$ un site annelé. Soient $f_1, \ldots, f_r$
des sections globales de $\mathcal{O}$. Soit $\mathcal{I} \subset \mathcal{O}$
le faisceau d'idéaux engendré par $f_1, \ldots, f_r$. Soit $K \in D(\mathcal{O})$.
Le complété dérivé $K^\wedge$ du lemme \ref{algebraization-lemma-derived-completion}
est donné par la formule
$$
K^\wedge = R\lim K \otimes^\mathbf{L}_\mathcal{O} K_n
$$
où $K_n = K(\mathcal{O}, f_1^n, \ldots, f_r^n)$
est le complexe de Koszul de $f_1^n, \ldots, f_r^n$ sur $\mathcal{O}$.
\end{lemma}

\begin{proof}
Dans Compléments d'algèbre, lemme
\ref{more-algebra-lemma-extended-alternating-Cech-is-colimit-koszul}
nous avons vu que le complexe de {\v C}ech alterné augmenté
$$
\mathcal{O} \to \prod\nolimits_{i_0} \mathcal{O}_{f_{i_0}} \to
\prod\nolimits_{i_0 < i_1} \mathcal{O}_{f_{i_0}f_{i_1}} \to
\ldots \to \mathcal{O}_{f_1\ldots f_r}
$$
est une limite inductive des complexes de Koszul
$K^n = K(\mathcal{O}, f_1^n, \ldots, f_r^n)$ placés en
degrés $0, \ldots, r$. Notons que $K^n$ est un complexe de chaînes fini
de $\mathcal{O}$-modules libres de type fini, dont le dual est
$\SheafHom_\mathcal{O}(K^n, \mathcal{O}) = K_n$, où $K_n$ est le complexe de Koszul
de cochaînes placé en degrés $-r, \ldots, 0$ (comme d'habitude). D'après le
lemme \ref{algebraization-lemma-derived-completion},
le foncteur $E \mapsto E^\wedge$ s'obtient en prenant
$R\SheafHom$ du complexe de {\v C}ech alterné augmenté dans $E$ :
$$
E^\wedge = R\SheafHom(\colim K^n, E)
$$
Cette expression est égale à $R\lim (E \otimes_\mathcal{O}^\mathbf{L} K_n)$
d'après
Cohomologie sur les sites, lemme \ref{sites-cohomology-lemma-colim-and-lim-of-duals}.
\end{proof}

\begin{lemma}
\label{algebraization-lemma-all-rings}
Il existe une construction des données suivantes :
\begin{enumerate}
\item pour tout couple $(A, I)$ formé d'un anneau $A$ et d'un idéal
de type fini $I \subset A$, un complexe $K(A, I)$ de $A$-modules,
\item un morphisme $K(A, I) \to A$ de complexes de $A$-modules,
\item pour tout homomorphisme d'anneaux $A \to B$ et tout idéal de type fini $I \subset A$,
un morphisme de complexes $K(A, I) \to K(B, IB)$,
\end{enumerate}
telle que
\begin{enumerate}
\item[(a)] pour $A \to B$ et $I \subset A$ de type fini, le diagramme
$$
\xymatrix{
K(A, I) \ar[r] \ar[d] & A \ar[d] \\
K(B, IB) \ar[r] & B
}
$$
soit commutatif,
\item[(b)] pour $A \to B \to C$ et $I \subset A$ de type fini,
la composée des applications
$K(A, I) \to K(B, IB) \to K(C, IC)$ soit l'application $K(A, I) \to K(C, IC)$.
\item[(c)] pour $A \to B$ et un idéal de type fini $I \subset A$,
l'application induite $K(A, I) \otimes_A^\mathbf{L} B \to K(B, IB)$
soit un isomorphisme dans $D(B)$, et
\item[(d)] si $I = (f_1, \ldots, f_r) \subset A$, il existe un diagramme
commutatif
$$
\xymatrix{
(A \to \prod\nolimits_{i_0} A_{f_{i_0}} \to
\prod\nolimits_{i_0 < i_1} A_{f_{i_0}f_{i_1}} \to
\ldots \to A_{f_1\ldots f_r}) \ar[r] \ar[d] &  K(A, I) \ar[d] \\
A \ar[r]^1 & A
}
$$
dans $D(A)$ dont les flèches horizontales sont des isomorphismes.
\end{enumerate}
\end{lemma}

\begin{proof}
Soit $S$ l'ensemble des anneaux $A_0$ de la forme
$A_0 = \mathbf{Z}[x_1, \ldots, x_n]/J$.
Toute $\mathbf{Z}$-algèbre de type fini est isomorphe à
un élément de $S$. Soit $\mathcal{A}_0$ la catégorie dont les objets
sont les couples $(A_0, I_0)$, où $A_0 \in S$ et $I_0 \subset A_0$
est un idéal, et dont les morphismes $(A_0, I_0) \to (B_0, J_0)$ sont
les homomorphismes d'anneaux $\varphi : A_0 \to B_0$ tels que $J_0 = \varphi(I_0)B_0$.

\medskip\noindent
Supposons que l'on puisse construire $K(A_0, I_0) \to A_0$ fonctoriellement pour
les objets de $\mathcal{A}_0$, avec les propriétés (a), (b), (c) et (d).
On pose alors
$$
K(A, I) = \colim_{\varphi : (A_0, I_0) \to (A, I)} K(A_0, I_0)
$$
où la limite inductive porte sur les homomorphismes d'anneaux $\varphi : A_0 \to A$ tels
que $\varphi(I_0)A = I$, avec $(A_0, I_0)$ dans $\mathcal{A}_0$.
Un morphisme de $(A_0, I_0) \to (A, I)$ vers $(A_0', I_0') \to (A, I)$
est donné par un morphisme $(A_0, I_0) \to (A_0', I_0')$ dans $\mathcal{A}_0$
commutant aux applications vers $A$.
La catégorie de ces $(A_0, I_0) \to (A, I)$ est filtrante
(nous omettons les détails). De plus, $\colim_{\varphi : (A_0, I_0) \to (A, I)} A_0 = A$,
de sorte que $K(A, I)$ est un complexe de $A$-modules.
Enfin, étant donnés $\varphi : A \to B$ et $I \subset A$,
pour chaque $(A_0, I_0) \to (A, I)$ intervenant dans la limite inductive, la composée
$(A_0, I_0) \to (B, IB)$ intervient dans celle relative à $(B, IB)$.
On obtient ainsi une application entre les limites inductives. Les propriétés (a), (b), (c) et (d)
en découlent aussitôt, compte tenu des propriétés correspondantes
des complexes $K(A_0, I_0)$.

\medskip\noindent
Munissons $\mathcal{C}_0 = \mathcal{A}_0^{opp}$ de la topologie chaotique.
Munissons $\mathcal{C}_0$ du faisceau d'anneaux
$\mathcal{O} : (A, I) \mapsto A$. Les idéaux $I$ définissent ensemble un
faisceau d'idéaux $\mathcal{I} \subset \mathcal{O}$.
Choisissons une résolution injective $\mathcal{O} \to \mathcal{J}^\bullet$.
Considérons l'objet
$$
\mathcal{F}^\bullet = \bigcup\nolimits_n \mathcal{J}^\bullet[\mathcal{I}^n]
$$
Soit $U = (A, I) \in \Ob(\mathcal{C}_0)$.
Puisque la topologie de $\mathcal{C}_0$ est chaotique, la valeur
$\mathcal{J}^\bullet(U)$ est une résolution de $A$ par des
$A$-modules injectifs. La valeur $\mathcal{F}^\bullet(U)$ est donc un
objet de $D(A)$ représentant l'image de $R\Gamma_I(A)$ dans $D(A)$ ; voir
Complexes dualisants, section \ref{dualizing-section-local-cohomology}.
Choisissons un complexe de $\mathcal{O}$-modules $\mathcal{K}^\bullet$
et un diagramme commutatif
$$
\xymatrix{
\mathcal{O} \ar[r] & \mathcal{J}^\bullet \\
\mathcal{K}^\bullet \ar[r] \ar[u] & \mathcal{F}^\bullet \ar[u]
}
$$
dont les flèches horizontales sont des quasi-isomorphismes. Cela est possible
par la construction de la catégorie dérivée $D(\mathcal{O})$.
Posons $K(A, I) = \mathcal{K}^\bullet(U)$, où $U = (A, I)$.
Les propriétés (a) et (b) sont claires, et les propriétés (c) et (d)
résultent de Complexes dualisants, lemmes
\ref{dualizing-lemma-compute-local-cohomology-noetherian} et
\ref{dualizing-lemma-local-cohomology-change-rings}.
\end{proof}

\begin{lemma}
\label{algebraization-lemma-global-extended-cech-complex}
Soit $(\mathcal{C}, \mathcal{O})$ un site annelé. Soit
$\mathcal{I} \subset \mathcal{O}$ un faisceau d'idéaux de type fini.
Il existe une application $K \to \mathcal{O}$ dans $D(\mathcal{O})$
telle que, pour tout $U \in \Ob(\mathcal{C})$ tel que
$\mathcal{I}|_U$ soit engendré par $f_1, \ldots, f_r \in \mathcal{I}(U)$,
il existe un isomorphisme
$$
(\mathcal{O}_U \to \prod\nolimits_{i_0} \mathcal{O}_{U, f_{i_0}} \to
\prod\nolimits_{i_0 < i_1} \mathcal{O}_{U, f_{i_0}f_{i_1}} \to
\ldots \to \mathcal{O}_{U, f_1\ldots f_r}) \longrightarrow K|_U
$$
compatible avec les applications vers $\mathcal{O}_U$.
\end{lemma}

\begin{proof}
Soit $\mathcal{C}' \subset \mathcal{C}$ la sous-catégorie pleine
des objets $U$ tels que $\mathcal{I}|_U$ soit engendré par
un nombre fini de sections. Alors $\mathcal{C}' \to \mathcal{C}$
est un foncteur cocontinu spécial
(Sites, définition \ref{sites-definition-special-cocontinuous-functor}).
Il suffit donc de travailler avec $\mathcal{C}'$ ; voir
Sites, lemme \ref{sites-lemma-equivalence}.
Autrement dit, on peut supposer que, pour tout
objet $U$ de $\mathcal{C}$, il existe un idéal de type fini
$I \subset \mathcal{I}(U)$ tel que
$\mathcal{I}|_U = \Im(I \otimes \mathcal{O}_U \to \mathcal{O}_U)$.
Nous dirons que $I$ engendre $\mathcal{I}|_U$.
Attention : on ne sait pas si $\mathcal{I}(U)$ est un idéal de type fini
de $\mathcal{O}(U)$.

\medskip\noindent
Soient $U$ un objet et $I \subset \mathcal{O}(U)$ un idéal de type
fini qui engendre $\mathcal{I}|_U$.
Sur la catégorie $\mathcal{C}/U$, considérons le complexe de préfaisceaux
$$
K_{U, I}^\bullet : U'/U \longmapsto K(\mathcal{O}(U'), I\mathcal{O}(U'))
$$
où $K(-, -)$ est comme dans le lemme \ref{algebraization-lemma-all-rings}.
Montrons que son faisceau associé est indépendant
du choix de $I$. En effet, si $I' \subset \mathcal{O}(U)$
est un idéal de type fini engendrant lui aussi $\mathcal{I}|_U$,
il existe un recouvrement $\{U_j \to U\}$ tel que
$I\mathcal{O}(U_j) = I'\mathcal{O}(U_j)$. (Indication : cela tient au fait que
$I$ et $I'$ sont tous deux de type fini et engendrent $\mathcal{I}|_U$.)
Ainsi $K_{U, I}^\bullet$ et $K_{U, I'}^\bullet$ sont {\it identiques}
sur tout objet au-dessus de l'un des $U_j$. L'assertion
sur les faisceaux associés s'ensuit. Notons $K_U^\bullet$ leur valeur commune.

\medskip\noindent
L'indépendance par rapport au choix de $I$ montre également que
$K_U^\bullet|_{\mathcal{C}/U'} = K_{U'}^\bullet$
dès que l'on se donne un morphisme
$U' \to U$, et donc un morphisme de localisation
$\mathcal{C}/U' \to \mathcal{C}/U$. Ainsi, les complexes
$K_U^\bullet$ se recollent en un unique complexe bien défini $K^\bullet$
de $\mathcal{O}$-modules. L'existence de l'application $K^\bullet \to \mathcal{O}$
et du quasi-isomorphisme du lemme résulte immédiatement
des propriétés correspondantes des complexes $K(-, -)$ dans le
lemme \ref{algebraization-lemma-all-rings}.
\end{proof}

\begin{proposition}
\label{algebraization-proposition-derived-completion}
Soit $(\mathcal{C}, \mathcal{O})$ un site annelé.
Soit $\mathcal{I} \subset \mathcal{O}$ un faisceau d'idéaux
de type fini. Le foncteur d'inclusion
$D_{comp}(\mathcal{O}) \to D(\mathcal{O})$ admet un adjoint à gauche.
\end{proposition}

\begin{proof}
Soit $K \to \mathcal{O}$ dans $D(\mathcal{O})$ l'application construite dans le
lemme \ref{algebraization-lemma-global-extended-cech-complex}. Soit $E \in D(\mathcal{O})$.
Alors $E^\wedge = R\SheafHom(K, E)$, muni de l'application $E \to E^\wedge$,
convient. En effet, localement sur le site $\mathcal{C}$, on
retrouve l'adjoint du lemme \ref{algebraization-lemma-derived-completion}.
Cela montre que $E^\wedge$ est toujours complet au sens dérivé et que
$E \to E^\wedge$ est un isomorphisme si $E$ est complet au sens dérivé.
\end{proof}

\begin{remark}[Comparaison avec la complétion]
\label{algebraization-remark-compare-with-completion}
Soit $(\mathcal{C}, \mathcal{O})$ un site annelé.
Soit $\mathcal{I} \subset \mathcal{O}$ un faisceau d'idéaux
de type fini. Soit $K \mapsto K^\wedge$ le foncteur de complétion dérivée
de la proposition \ref{algebraization-proposition-derived-completion}.
Pour tout $n \geq 1$, l'objet
$K \otimes_\mathcal{O}^\mathbf{L} \mathcal{O}/\mathcal{I}^n$
est complet au sens dérivé, car il est annulé par des puissances des
sections locales de $\mathcal{I}$. On a donc une factorisation canonique
$$
K \to K^\wedge \to K \otimes_\mathcal{O}^\mathbf{L} \mathcal{O}/\mathcal{I}^n
$$
de l'application canonique
$K \to K \otimes_\mathcal{O}^\mathbf{L} \mathcal{O}/\mathcal{I}^n$.
Ces applications sont compatibles lorsque $n$ varie, et l'on obtient une application de comparaison
$$
K^\wedge
\longrightarrow
R\lim \left(K \otimes_\mathcal{O}^\mathbf{L} \mathcal{O}/\mathcal{I}^n\right)
$$
Le membre de droite s'interprète plus manifestement comme une sorte de complété.
En général, cette application de comparaison n'est pas un isomorphisme.
\end{remark}

\begin{remark}[Localisation et complétion dérivée]
\label{algebraization-remark-localization-and-completion}
Soit $(\mathcal{C}, \mathcal{O})$ un site annelé.
Soit $\mathcal{I} \subset \mathcal{O}$ un faisceau d'idéaux
de type fini. Soit $K \mapsto K^\wedge$ le foncteur de complétion dérivée
de la proposition \ref{algebraization-proposition-derived-completion}. La construction
donnée dans la démonstration de la proposition montre que $K^\wedge|_U$
est le complété dérivé de $K|_U$ pour tout $U \in \Ob(\mathcal{C})$.
On peut aussi le démontrer comme suit. La définition
des objets complets au sens dérivé montre que $K^\wedge|_U$ est complet au sens dérivé.
On obtient ainsi une application canonique $a : (K|_U)^\wedge \to K^\wedge|_U$.
D'autre part, si $E$ est un objet complet au sens dérivé de
$D(\mathcal{O}_U)$, alors $Rj_*E$ est un objet complet au sens dérivé de
$D(\mathcal{O})$ d'après le lemme \ref{algebraization-lemma-pushforward-derived-complete}.
Ici $j$ est le morphisme de localisation
(Modules sur les sites, section \ref{sites-modules-section-localize}).
On obtient donc aussi une application
canonique $b : K^\wedge \to Rj_*((K|_U)^\wedge)$. Nous omettons la vérification formelle
du fait que l'adjoint de $b$ est l'inverse de $a$.
\end{remark}

\begin{remark}[Produit tensoriel complété]
\label{algebraization-remark-completed-tensor-product}
Soit $(\mathcal{C}, \mathcal{O})$ un site annelé. Soit
$\mathcal{I} \subset \mathcal{O}$ un faisceau d'idéaux de type fini.
Notons $K \mapsto K^\wedge$ l'adjoint de la
proposition \ref{algebraization-proposition-derived-completion}.
Posons alors
$$
K \otimes^\wedge_\mathcal{O} L = (K \otimes_\mathcal{O}^\mathbf{L} L)^\wedge
$$
Ce {\it produit tensoriel complété} définit un foncteur
$D_{comp}(\mathcal{O}) \times D_{comp}(\mathcal{O}) \to D_{comp}(\mathcal{O})$
tel que l'on ait
$$
\Hom_{D_{comp}(\mathcal{O})}(K, R\SheafHom_\mathcal{O}(L, M))
=
\Hom_{D_{comp}(\mathcal{O})}(K \otimes_\mathcal{O}^\wedge L, M)
$$
pour $K, L, M \in D_{comp}(\mathcal{O})$. Notons que
$R\SheafHom_\mathcal{O}(L, M) \in D_{comp}(\mathcal{O})$ d'après le
lemme \ref{algebraization-lemma-derived-complete-internal-hom}.
\end{remark}

\begin{lemma}
\label{algebraization-lemma-map-identifies-koszul-and-cech-complexes}
Soit $\mathcal{C}$ un site.
Supposons que $\varphi : \mathcal{O} \to \mathcal{O}'$ soit un homomorphisme plat
de faisceaux d'anneaux. Soient $f_1, \ldots, f_r$ des sections globales
de $\mathcal{O}$ telles que $\mathcal{O}/(f_1, \ldots, f_r) \cong
\mathcal{O}'/(f_1, \ldots, f_r)\mathcal{O}'$.
Alors le morphisme de complexes de {\v C}ech alternés augmentés
$$
\xymatrix{
\mathcal{O} \to
\prod_{i_0} \mathcal{O}_{f_{i_0}} \to
\prod_{i_0 < i_1} \mathcal{O}_{f_{i_0}f_{i_1}} \to \ldots \to
\mathcal{O}_{f_1\ldots f_r} \ar[d] \\
\mathcal{O}' \to
\prod_{i_0} \mathcal{O}'_{f_{i_0}} \to
\prod_{i_0 < i_1} \mathcal{O}'_{f_{i_0}f_{i_1}} \to \ldots \to
\mathcal{O}'_{f_1\ldots f_r}
}
$$
est un quasi-isomorphisme.
\end{lemma}

\begin{proof}
Remarquons que le second complexe est le produit tensoriel du premier
avec $\mathcal{O}'$. On peut écrire le premier complexe
de {\v C}ech alterné augmenté comme limite inductive des complexes de Koszul
$K_n = K(\mathcal{O}, f_1^n, \ldots, f_r^n)$ ; voir
Compléments d'algèbre, lemme
\ref{more-algebra-lemma-extended-alternating-Cech-is-colimit-koszul}.
Il suffit donc de montrer que $K_n \to K_n \otimes_\mathcal{O} \mathcal{O}'$
est un quasi-isomorphisme. Puisque $\mathcal{O} \to \mathcal{O}'$ est plat,
il suffit de montrer que $H^i \to H^i \otimes_\mathcal{O} \mathcal{O}'$
est un isomorphisme, où $H^i$ est le faisceau de cohomologie de degré $i$,
$H^i = H^i(K_n)$. Ces faisceaux sont annulés par $f_1^n, \ldots, f_r^n$ ; voir
Compléments d'algèbre, lemme \ref{more-algebra-lemma-homotopy-koszul}.
Ils sont donc annulés par $(f_1, \ldots, f_r)^m$
pour un certain $m \gg 0$. Ainsi $H^i \to H^i \otimes_\mathcal{O} \mathcal{O}'$
est un isomorphisme d'après Modules sur les sites, lemme
\ref{sites-modules-lemma-neighbourhood-isomorphism}.
\end{proof}

\begin{lemma}
\label{algebraization-lemma-restriction-derived-complete-equivalence}
Soit $\mathcal{C}$ un site. Soit $\mathcal{O} \to \mathcal{O}'$ un
homomorphisme de faisceaux d'anneaux. Soit $\mathcal{I} \subset \mathcal{O}$
un faisceau d'idéaux de type fini.
Si $\mathcal{O} \to \mathcal{O}'$ est plat et
$\mathcal{O}/\mathcal{I} \cong \mathcal{O}'/\mathcal{I}\mathcal{O}'$,
alors le foncteur de restriction $D(\mathcal{O}') \to D(\mathcal{O})$
induit une équivalence
$D_{comp}(\mathcal{O}', \mathcal{I}\mathcal{O}') \to
D_{comp}(\mathcal{O}, \mathcal{I})$.
\end{lemma}

\begin{proof}
Le lemme \ref{algebraization-lemma-pushforward-derived-complete} entraîne que le foncteur de
restriction $r : D(\mathcal{O}') \to D(\mathcal{O})$
envoie $D_{comp}(\mathcal{O}', \mathcal{I}\mathcal{O}')$
dans $D_{comp}(\mathcal{O}, \mathcal{I})$. Nous allons construire un
quasi-inverse $E \mapsto E'$.

\medskip\noindent
Soit $K \to \mathcal{O}$ le morphisme de $D(\mathcal{O})$
construit dans le lemme \ref{algebraization-lemma-global-extended-cech-complex}.
Posons $K' = K \otimes_\mathcal{O}^\mathbf{L} \mathcal{O}'$ dans $D(\mathcal{O}')$.
Alors $K' \to \mathcal{O}'$ est une application dans $D(\mathcal{O}')$ qui
satisfait aux conclusions du lemme \ref{algebraization-lemma-global-extended-cech-complex}
relativement à $\mathcal{I}' = \mathcal{I}\mathcal{O}'$.
L'application $K \to r(K')$ est un quasi-isomorphisme d'après le
lemme \ref{algebraization-lemma-map-identifies-koszul-and-cech-complexes}.
Pour $E \in D_{comp}(\mathcal{O}, \mathcal{I})$, posons maintenant
$$
E' = R\SheafHom_\mathcal{O}(r(K'), E)
$$
considéré comme un objet de $D(\mathcal{O}')$ à l'aide de la structure de $\mathcal{O}'$-module
de $K'$. Puisque $E$ est complet au sens dérivé,
on a $E = R\SheafHom_\mathcal{O}(K, E)$ ; voir la
démonstration de la proposition \ref{algebraization-proposition-derived-completion}.
D'autre part, puisque $K \to r(K')$ est un isomorphisme,
on voit qu'il existe un isomorphisme
$E \to r(E')$ dans $D(\mathcal{O})$. Pour achever la démonstration, il
faut montrer que, si $E = r(M')$ pour un objet $M'$ de
$D_{comp}(\mathcal{O}', \mathcal{I}')$, alors
$E' \cong M'$. Pour obtenir une application, on utilise
$$
M' = R\SheafHom_{\mathcal{O}'}(\mathcal{O}', M') \to
R\SheafHom_\mathcal{O}(r(\mathcal{O}'), r(M')) \to
R\SheafHom_\mathcal{O}(r(K'), r(M')) = E'
$$
où la seconde flèche utilise l'application $K' \to \mathcal{O}'$.
Pour voir que c'est un isomorphisme, on montre que $r$ appliqué
à cette flèche coïncide avec l'isomorphisme $E \to r(E')$ ci-dessus.
Nous omettons les détails.
\end{proof}

\begin{lemma}
\label{algebraization-lemma-pushforward-derived-complete-adjoint}
Soit $f : (\Sh(\mathcal{D}), \mathcal{O}') \to (\Sh(\mathcal{C}), \mathcal{O})$
un morphisme de topos annelés. Soient $\mathcal{I} \subset \mathcal{O}$
et $\mathcal{I}' \subset \mathcal{O}'$
des faisceaux d'idéaux de type fini tels que $f^\sharp$ envoie
$f^{-1}\mathcal{I}$ dans $\mathcal{I}'$.
Alors $Rf_*$ envoie $D_{comp}(\mathcal{O}', \mathcal{I}')$
dans $D_{comp}(\mathcal{O}, \mathcal{I})$ et admet un adjoint à gauche
$Lf_{comp}^*$, qui est $Lf^*$ suivi de la complétion dérivée.
\end{lemma}

\begin{proof}
La première assertion a été établie dans le
lemme \ref{algebraization-lemma-pushforward-derived-complete}.
Notons que la seconde a un sens, car on dispose d'un foncteur de complétion
dérivée $D(\mathcal{O}') \to D_{comp}(\mathcal{O}', \mathcal{I}')$
d'après la proposition \ref{algebraization-proposition-derived-completion}.
Soient donc $K \in D_{comp}(\mathcal{O}, \mathcal{I})$
et $M \in D_{comp}(\mathcal{O}', \mathcal{I}')$. On a alors
$$
\Hom(K, Rf_*M) = \Hom(Lf^*K, M) = \Hom(Lf_{comp}^*K, M)
$$
par la propriété universelle de la complétion dérivée.
\end{proof}

\begin{lemma}
\label{algebraization-lemma-pushforward-commutes-with-derived-completion}
\begin{reference}
Généralisation de \cite[Lemma 6.5.9 (2)]{BS}. Comparer à
\cite[Theorem 6.5]{HL-P} dans le cadre des modules quasi-cohérents
et des morphismes de champs algébriques (dérivés).
\end{reference}
Soit $f : (\Sh(\mathcal{D}), \mathcal{O}') \to (\Sh(\mathcal{C}), \mathcal{O})$
un morphisme de topos annelés. Soit $\mathcal{I} \subset \mathcal{O}$
un faisceau d'idéaux de type fini. Soit $\mathcal{I}' \subset \mathcal{O}'$
l'idéal engendré par $f^\sharp(f^{-1}\mathcal{I})$.
Alors $Rf_*$ commute à la complétion dérivée, c'est-à-dire
$Rf_*(K^\wedge) = (Rf_*K)^\wedge$.
\end{lemma}

\begin{proof}
Les foncteurs de complétion dérivée existent d'après la proposition \ref{algebraization-proposition-derived-completion}.
D'après le lemme \ref{algebraization-lemma-pushforward-derived-complete}, l'objet
$Rf_*(K^\wedge)$ est complet au sens dérivé ; on obtient donc une application canonique
$(Rf_*K)^\wedge \to Rf_*(K^\wedge)$ par la propriété universelle de la complétion
dérivée. On peut vérifier que c'est un isomorphisme localement sur $\mathcal{C}$.
Ainsi, puisque la complétion dérivée commute à la localisation
(remarque \ref{algebraization-remark-localization-and-completion}), on peut supposer
que $\mathcal{I}$ est engendré par des sections globales $f_1, \ldots, f_r$.
Alors $\mathcal{I}'$ est engendré par $g_i = f^\sharp(f_i)$. D'après le
lemme \ref{algebraization-lemma-derived-completion-koszul},
il faut montrer que
$$
R\lim \left(
Rf_*K \otimes^\mathbf{L}_\mathcal{O} K(\mathcal{O}, f_1^n, \ldots, f_r^n)
\right)
=
Rf_*\left(
R\lim
K \otimes^\mathbf{L}_{\mathcal{O}'} K(\mathcal{O}', g_1^n, \ldots, g_r^n)
\right)
$$
Puisque $Rf_*$ commute à $R\lim$
(Cohomologie sur les sites, lemme
\ref{sites-cohomology-lemma-Rf-commutes-with-Rlim})
il suffit de montrer que
$$
Rf_*K \otimes^\mathbf{L}_\mathcal{O} K(\mathcal{O}, f_1^n, \ldots, f_r^n) =
Rf_*\left(
K \otimes^\mathbf{L}_{\mathcal{O}'} K(\mathcal{O}', g_1^n, \ldots, g_r^n)
\right)
$$
Cela résulte de la formule de projection (Cohomologie sur les sites, lemme
\ref{sites-cohomology-lemma-projection-formula}) et du fait que
$Lf^*K(\mathcal{O}, f_1^n, \ldots, f_r^n) =
K(\mathcal{O}', g_1^n, \ldots, g_r^n)$.
\end{proof}

\begin{lemma}
\label{algebraization-lemma-formal-functions-general}
Soient $A$ un anneau et $I \subset A$ un idéal de type fini.
Soit $\mathcal{C}$ un site et soit $\mathcal{O}$ un faisceau
de $A$-algèbres. Soit $\mathcal{F}$ un faisceau de $\mathcal{O}$-modules.
On a alors
$$
R\Gamma(\mathcal{C}, \mathcal{F})^\wedge =
R\Gamma(\mathcal{C}, \mathcal{F}^\wedge)
$$
dans $D(A)$, où $\mathcal{F}^\wedge$ est le complété
dérivé de $\mathcal{F}$ relativement à $I\mathcal{O}$, tandis que le membre de
gauche est le complété dérivé relativement à $I$.
On obtient ainsi deux suites spectrales
$$
E_2^{i, j} = H^i(H^j(\mathcal{C}, \mathcal{F})^\wedge)
\quad\text{et}\quad
E_2^{p, q} = H^p(\mathcal{C}, H^q(\mathcal{F}^\wedge))
$$
convergeant toutes deux vers
$H^*(R\Gamma(\mathcal{C}, \mathcal{F})^\wedge) =
H^*(\mathcal{C}, \mathcal{F}^\wedge)$
\end{lemma}

\begin{proof}
Appliquons le lemme \ref{algebraization-lemma-pushforward-commutes-with-derived-completion}
au morphisme de topos annelés $(\mathcal{C}, \mathcal{O}) \to (pt, A)$
et prenons la cohomologie pour obtenir la première assertion. La seconde suite spectrale
est la seconde suite spectrale de
Catégories dérivées, lemme \ref{derived-lemma-two-ss-complex-functor}.
La première suite spectrale est celle de
Compléments d'algèbre, exemple
\ref{more-algebra-example-derived-completion-spectral-sequence}
appliqué à $R\Gamma(\mathcal{C}, \mathcal{F})^\wedge$.
\end{proof}

\begin{remark}
\label{algebraization-remark-local-calculation-derived-completion}
Soit $(\mathcal{C}, \mathcal{O})$ un site annelé.
Soit $\mathcal{I} \subset \mathcal{O}$ un faisceau d'idéaux
de type fini. Soit $K \mapsto K^\wedge$ la complétion dérivée de la
proposition \ref{algebraization-proposition-derived-completion}.
Soit $U \in \Ob(\mathcal{C})$ un objet tel que $\mathcal{I}$
soit engendré comme faisceau d'idéaux par $f_1, \ldots, f_r \in \mathcal{I}(U)$.
Posons $A = \mathcal{O}(U)$ et $I = (f_1, \ldots, f_r) \subset A$.
Attention : il se peut que l'égalité $I = \mathcal{I}(U)$ ne soit pas vérifiée.
On a alors
$$
R\Gamma(U, K^\wedge) = R\Gamma(U, K)^\wedge
$$
où le membre de droite est le complété dérivé de
l'objet $R\Gamma(U, K)$ de $D(A)$ relativement à $I$.
Cela résulte de la commutation de la complétion dérivée à la localisation
(remarque \ref{algebraization-remark-localization-and-completion}) et du
lemme \ref{algebraization-lemma-formal-functions-general}.
\end{remark}








\section{Le théorème des fonctions formelles}
\label{algebraization-section-formal-functions}

\noindent
Interrompons le fil de l'exposé pour dire quelques mots de
la complétion dérivée dans le cadre des modules quasi-cohérents sur les schémas
et l'utiliser pour donner une démonstration quelque peu différente du théorème
des fonctions formelles. Quelques indications bibliographiques figurent dans la
remarque \ref{algebraization-remark-references}.

\medskip\noindent
Le lemme \ref{algebraization-lemma-pushforward-commutes-with-derived-completion} est une
version dérivée (très formelle) du théorème des fonctions formelles
(Cohomologie des schémas, théorème \ref{coherent-theorem-formal-functions}).
Pour l'expliciter, supposons que $f : X \to S$ soit un morphisme de schémas,
que $\mathcal{I} \subset \mathcal{O}_S$ soit un faisceau quasi-cohérent d'idéaux
de type fini
et que $\mathcal{F}$ soit un faisceau quasi-cohérent sur $X$. Le lemme affirme alors que
\begin{equation}
\label{algebraization-equation-formal-functions}
Rf_*(\mathcal{F}^\wedge) = (Rf_*\mathcal{F})^\wedge
\end{equation}
où $\mathcal{F}^\wedge$ est le complété dérivé de $\mathcal{F}$
relativement à $f^{-1}\mathcal{I} \cdot \mathcal{O}_X$, et le membre
de droite est le complété dérivé de $Rf_*\mathcal{F}$
relativement à $\mathcal{I}$. Pour retrouver ainsi le théorème
des fonctions formelles, il reste quelques vérifications à faire.

\begin{lemma}
\label{algebraization-lemma-sections-derived-completion-pseudo-coherent}
Soit $X$ un schéma localement noethérien. Soit $\mathcal{I} \subset \mathcal{O}_X$
un faisceau quasi-cohérent d'idéaux. Soit $K$ un
objet pseudo-cohérent de $D(\mathcal{O}_X)$, de complété dérivé
$K^\wedge$. Alors
$$
H^p(U, K^\wedge) = \lim H^p(U, K)/I^nH^p(U, K) =
H^p(U, K)^\wedge
$$
pour tout ouvert affine $U \subset X$,
où $I = \mathcal{I}(U)$ et où le membre de droite est le complété
dérivé relativement à $I$.
\end{lemma}

\begin{proof}
Écrivons $U = \Spec(A)$. L'anneau $A$ est noethérien,
donc $I \subset A$ est de type fini. On a alors
$$
R\Gamma(U, K^\wedge) = R\Gamma(U, K)^\wedge
$$
d'après la remarque \ref{algebraization-remark-local-calculation-derived-completion}.
Or $R\Gamma(U, K)$ est un complexe pseudo-cohérent de $A$-modules
(Catégories dérivées des schémas, lemme
\ref{perfect-lemma-pseudo-coherent-affine}).
D'après Compléments d'algèbre, lemme
\ref{more-algebra-lemma-derived-completion-pseudo-coherent}
on en conclut que le module de cohomologie de degré $p$ de $R\Gamma(U, K^\wedge)$
est égal au complété $I$-adique de $H^p(U, K)$.
Cela prouve la première égalité. La seconde, moins importante,
résulte immédiatement d'une nouvelle application du lemme que l'on vient d'utiliser.
\end{proof}

\begin{lemma}
\label{algebraization-lemma-derived-completion-pseudo-coherent}
Soit $X$ un schéma localement noethérien. Soit $\mathcal{I} \subset \mathcal{O}_X$
un faisceau quasi-cohérent d'idéaux.
Soit $K$ un objet de $D(\mathcal{O}_X)$. Alors
\begin{enumerate}
\item le complété dérivé $K^\wedge$ est égal à
$R\lim (K \otimes_{\mathcal{O}_X}^\mathbf{L} \mathcal{O}_X/\mathcal{I}^n)$.
\end{enumerate}
Soit $K$ un objet pseudo-cohérent de $D(\mathcal{O}_X)$. Alors
\begin{enumerate}
\item[(2)] le faisceau de cohomologie $H^q(K^\wedge)$ est égal à
$\lim H^q(K)/\mathcal{I}^nH^q(K)$.
\end{enumerate}
Soit $\mathcal{F}$ un $\mathcal{O}_X$-module cohérent\footnote{Par exemple
$H^q(K)$ pour $K$ pseudo-cohérent sur notre schéma localement noethérien $X$.}. Alors
\begin{enumerate}
\item[(3)] le complété dérivé $\mathcal{F}^\wedge$ est égal à
$\lim \mathcal{F}/\mathcal{I}^n\mathcal{F}$,
\item[(4)]
$\lim \mathcal{F}/\mathcal{I}^n \mathcal{F} =
R\lim \mathcal{F}/\mathcal{I}^n \mathcal{F}$,
\item[(5)] $H^p(U, \mathcal{F}^\wedge) = 0$ pour $p \not = 0$ et pour tout
ouvert affine $U \subset X$.
\end{enumerate}
\end{lemma}

\begin{proof}
Démonstration de (1). On a une application canonique
$$
K \longrightarrow
R\lim (K \otimes_{\mathcal{O}_X}^\mathbf{L} \mathcal{O}_X/\mathcal{I}^n),
$$
voir la remarque \ref{algebraization-remark-compare-with-completion}.
La complétion dérivée commute au passage à un sous-schéma ouvert
(remarque \ref{algebraization-remark-localization-and-completion}).
La formation de $R\lim$ commute au passage à un sous-schéma ouvert.
Pour vérifier que notre application est un isomorphisme, on peut donc travailler localement.
On peut ainsi supposer $X = U = \Spec(A)$.
Écrivons $I = (f_1, \ldots, f_r)$. Soit
$K_n = K(A, f_1^n, \ldots, f_r^n)$ le complexe de Koszul.
D'après Compléments d'algèbre, lemme \ref{more-algebra-lemma-sequence-Koszul-complexes},
nous avons vu que les pro-systèmes $\{K_n\}$ et
$\{A/I^n\}$ de $D(A)$ sont isomorphes.
En utilisant l'équivalence $D(A) = D_{\QCoh}(\mathcal{O}_X)$
de Catégories dérivées des schémas, lemme
\ref{perfect-lemma-affine-compare-bounded}
on voit que les pro-systèmes $\{K(\mathcal{O}_X, f_1^n, \ldots, f_r^n)\}$
et $\{\mathcal{O}_X/\mathcal{I}^n\}$ sont isomorphes dans
$D(\mathcal{O}_X)$. Cela prouve la seconde égalité dans
$$
K^\wedge = R\lim \left(
K \otimes_{\mathcal{O}_X}^\mathbf{L} K(\mathcal{O}_X, f_1^n, \ldots, f_r^n)
\right) =
R\lim (K \otimes_{\mathcal{O}_X}^\mathbf{L} \mathcal{O}_X/\mathcal{I}^n)
$$
La première égalité est donnée par le
lemme \ref{algebraization-lemma-derived-completion-koszul}.

\medskip\noindent
Supposons $K$ pseudo-cohérent. Pour un ouvert affine $U \subset X$,
on a $H^q(U, K^\wedge) = \lim H^q(U, K)/\mathcal{I}^n(U)H^q(U, K)$
d'après le lemme \ref{algebraization-lemma-sections-derived-completion-pseudo-coherent}.
Comme cela vaut pour tout $U$, on voit que
$H^q(K^\wedge) = \lim H^q(K)/\mathcal{I}^nH^q(K)$ comme faisceaux.
Cela prouve (2).

\medskip\noindent
L'assertion (3) est un cas particulier de (2).
Les assertions (4) et (5) résultent de
Catégories dérivées des schémas, lemme
\ref{perfect-lemma-Rlim-quasi-coherent}.
\end{proof}

\begin{lemma}
\label{algebraization-lemma-formal-functions}
Soient $A$ un anneau noethérien et $I \subset A$ un idéal. Soit $X$ un
schéma noethérien sur $A$. Soit $\mathcal{F}$ un
$\mathcal{O}_X$-module cohérent. Supposons que $H^p(X, \mathcal{F})$ soit
un $A$-module de type fini pour tout $p$. Il existe alors des suites exactes courtes
$$
0 \to R^1\lim H^{p - 1}(X, \mathcal{F}/I^n\mathcal{F}) \to
H^p(X, \mathcal{F})^\wedge \to \lim H^p(X, \mathcal{F}/I^n\mathcal{F}) \to 0
$$
de $A$-modules, où $H^p(X, \mathcal{F})^\wedge$ est le complété $I$-adique
usuel. Si $f$ est propre, le terme $R^1\lim$ est nul.
\end{lemma}

\begin{proof}
Considérons les deux suites spectrales du
lemme \ref{algebraization-lemma-formal-functions-general}.
La première dégénère d'après Compléments d'algèbre, lemme
\ref{more-algebra-lemma-derived-completion-pseudo-coherent}.
On obtient $H^p(X, \mathcal{F})^\wedge$ en degré $p$.
C'est ici qu'intervient l'hypothèse que $H^p(X, \mathcal{F})$ est
un $A$-module de type fini. La seconde dégénère puisque
$$
\mathcal{F}^\wedge = \lim \mathcal{F}/I^n\mathcal{F} =
R\lim \mathcal{F}/I^n\mathcal{F}
$$
est un faisceau d'après le lemme \ref{algebraization-lemma-derived-completion-pseudo-coherent}.
On obtient $H^p(X, \lim \mathcal{F}/I^n\mathcal{F})$ en degré $p$.
Puisque $R\Gamma(X, -)$ commute aux limites dérivées
(Injectifs, lemme \ref{injectives-lemma-RF-commutes-with-Rlim}),
on obtient aussi
$$
R\Gamma(X, \lim \mathcal{F}/I^n\mathcal{F}) =
R\Gamma(X, R\lim \mathcal{F}/I^n\mathcal{F}) =
R\lim R\Gamma(X, \mathcal{F}/I^n\mathcal{F})
$$
D'après Compléments d'algèbre, remarque
\ref{more-algebra-remark-how-unique}
on obtient des suites exactes
$$
0 \to
R^1\lim H^{p - 1}(X, \mathcal{F}/I^n\mathcal{F}) \to
H^p(X, \lim \mathcal{F}/I^n\mathcal{F}) \to
\lim H^p(X, \mathcal{F}/I^n\mathcal{F}) \to 0
$$
de $A$-modules. En réunissant ces résultats, on obtient la première assertion du lemme.
L'annulation du terme $R^1\lim$ résulte de
Cohomologie des schémas, lemme \ref{coherent-lemma-ML-cohomology-powers-ideal}.
\end{proof}

\begin{remark}
\label{algebraization-remark-references}
Voici quelques références bibliographiques sur des questions voisines.
Il semble qu'un « théorème dérivé des fonctions formelles » pour les morphismes propres
remonte à \cite[théorème 6.3.1]{lurie-thesis}.
On trouve aussi un exposé dans \cite{dag12}, en particulier
au chapitre 4, qui traite du cas affine ; voir
Compléments d'algèbre, section \ref{more-algebra-section-derived-completion}.
Dans \cite[section 2.9]{G-R}, le changement de base propre et la
complétion dérivée sont étudiés à l'aide de modules (ind-)cohérents.
Un analogue de (\ref{algebraization-equation-formal-functions})
pour les complexes de modules quasi-cohérents figure dans
\cite[théorème 6.5]{HL-P}
\end{remark}







\section{Algébrisation de la cohomologie locale, I}
\label{algebraization-section-algebraization-sections-general}

\noindent
Soit $A$ un anneau noethérien et soient $I$ et $J$ deux idéaux de $A$.
Soit $M$ un $A$-module de type fini. Dans cette section, nous étudions les
groupes de cohomologie de l'objet
$$
R\Gamma_J(M)^\wedge
\quad\text{de}\quad
D(A)
$$
où ${}^\wedge$ désigne la complétion $I$-adique dérivée. Rappelons que, dans
Complexes dualisants, lemme \ref{dualizing-lemma-completion-local-H0},
nous avons montré que, si $A$ est complet relativement à $I$,
il existe un isomorphisme
$$
\colim H^0_Z(M) \longrightarrow H^0(R\Gamma_J(M)^\wedge)
$$
où la limite inductive filtrante porte sur les parties fermées $Z = V(J')$
avec $J' \subset J$ et $V(J') \cap V(I) = V(J) \cap V(I)$.
La réunion de ces parties fermées est
\begin{equation}
\label{algebraization-equation-associated-subset}
T = \{\mathfrak p \in \Spec(A) :
V(\mathfrak p) \cap V(I) \subset V(J) \cap V(I)\}
\end{equation}
C'est une partie de $\Spec(A)$ stable par spécialisation.
Le résultat précédent signifie alors que
$$
H^0_T(M) \longrightarrow H^0(R\Gamma_J(M)^\wedge)
$$
est un isomorphisme si $A$ est complet relativement à $I$ ; voir
Cohomologie locale, lemme \ref{local-cohomology-lemma-adjoint-ext} et
remarque \ref{local-cohomology-remark-upshot}.
Notre méthode pour étendre cet isomorphisme aux groupes de cohomologie supérieure
repose sur le lemme suivant.

\begin{lemma}
\label{algebraization-lemma-kill-completion-general}
Soient $I, J$ des idéaux d'un anneau noethérien $A$.
Soit $M$ un $A$-module de type fini. Soit $\mathfrak p \subset A$ un idéal premier.
Soient $s$ et $d$ des entiers. Supposons que
\begin{enumerate}
\item $A$ possède un complexe dualisant,
\item $\mathfrak p \not \in V(J) \cap V(I)$,
\item $\text{cd}(A, I) \leq d$, et
\item pour tout idéal premier $\mathfrak p' \subset \mathfrak p$,
on ait
$\text{depth}_{A_{\mathfrak p'}}(M_{\mathfrak p'}) +
\dim((A/\mathfrak p')_\mathfrak q) > d + s$
pour tout $\mathfrak q \in V(\mathfrak p') \cap V(J) \cap V(I)$.
\end{enumerate}
Alors il existe un $f \in A$, $f \not \in \mathfrak p$, qui annule
$H^i(R\Gamma_J(M)^\wedge)$ pour $i \leq s$, où ${}^\wedge$
désigne la complétion $I$-adique.
\end{lemma}

\begin{proof}
Nous utiliserons l'égalité $R\Gamma_J = R\Gamma_{V(J)}$ et de même pour
$I + J$ ; voir
Complexes dualisants, lemme \ref{dualizing-lemma-local-cohomology-noetherian}.
Remarquons que
$R\Gamma_J(M)^\wedge = R\Gamma_I(R\Gamma_J(M))^\wedge =
R\Gamma_{I + J}(M)^\wedge$ ; voir
Complexes dualisants, lemmes
\ref{dualizing-lemma-complete-and-local} et
\ref{dualizing-lemma-local-cohomology-ss}.
On peut donc remplacer $J$ par $I + J$ et supposer $I \subset J$
et $\mathfrak p \not \in V(J)$.
Rappelons que
$$
R\Gamma_J(M)^\wedge = R\Hom_A(R\Gamma_I(A), R\Gamma_J(M))
$$
d'après la description de la complétion dérivée dans
Compléments d'algèbre, lemme \ref{more-algebra-lemma-derived-completion},
combinée à la description de la cohomologie locale dans
Complexes dualisants, lemme
\ref{dualizing-lemma-compute-local-cohomology-noetherian}.
L'hypothèse (3) signifie que la cohomologie de $R\Gamma_I(A)$ n'est non nulle
qu'en degrés $\leq d$. En utilisant les troncatures canoniques de
$R\Gamma_I(A)$, on voit qu'il suffit de montrer que
$$
\text{Ext}^i(N, R\Gamma_J(M))
$$
est annulé par un $f \in A$, $f \not \in \mathfrak p$, pour
$i \leq s + d$ et pour tout $A$-module $N$.
En utilisant à présent les troncatures canoniques de $R\Gamma_J(M)$,
on voit qu'il suffit de montrer que
$H^i_J(M)$ est annulé par un $f \in A$, $f \not \in \mathfrak p$,
pour $i \leq s + d$.
Cela résulte de Cohomologie locale, lemme
\ref{local-cohomology-lemma-kill-local-cohomology-at-prime}.
\end{proof}

\begin{lemma}
\label{algebraization-lemma-kill-colimit-weak-general}
Soient $I, J$ des idéaux d'un anneau noethérien. Soit $M$ un $A$-module de type fini.
Soient $s$ et $d$ des entiers. Avec $T$ comme dans
(\ref{algebraization-equation-associated-subset}), supposons que
\begin{enumerate}
\item $A$ possède un complexe dualisant,
\item si $\mathfrak p \in V(I)$, aucune condition n'est imposée,
\item si $\mathfrak p \not \in V(I)$, $\mathfrak p \in T$, alors
$\dim((A/\mathfrak p)_\mathfrak q) \leq d$ pour un certain
$\mathfrak q \in V(\mathfrak p) \cap V(J) \cap V(I)$,
\item si $\mathfrak p \not \in V(I)$, $\mathfrak p \not \in T$, alors
$$
\text{depth}_{A_\mathfrak p}(M_\mathfrak p) \geq s
\quad\text{ou}\quad
\text{depth}_{A_\mathfrak p}(M_\mathfrak p) +
\dim((A/\mathfrak p)_\mathfrak q) > d + s
$$
pour tout $\mathfrak q \in V(\mathfrak p) \cap V(J) \cap V(I)$.
\end{enumerate}
Alors il existe un idéal $J_0 \subset J$ tel que
$V(J_0) \cap V(I) = V(J) \cap V(I)$ et tel que, pour tout $J' \subset J_0$ avec
$V(J') \cap V(I) = V(J) \cap V(I)$, l'application
$$
R\Gamma_{J'}(M) \longrightarrow R\Gamma_{J_0}(M)
$$
induise un isomorphisme en cohomologie en degrés $\leq s$ ;
de plus, ces modules sont annulés par une puissance de $J_0I$.
\end{lemma}

\begin{proof}
Considérons l'ensemble
$$
B = \{\mathfrak p \not \in V(I),\ \mathfrak p \in T,\text{ et }
\text{depth}(M_\mathfrak p) \leq s\}
$$
Choisissons $J_0 \subset J$ tel que $V(J_0)$ soit l'adhérence de $B \cup V(J)$.

\medskip\noindent
Assertion I : $V(J_0) \cap V(I) = V(J) \cap V(I)$.

\medskip\noindent
Démonstration de l'assertion I. L'inclusion $\supset$ résulte de la construction.
Soit $\mathfrak p$ un idéal premier minimal de $V(J_0)$.
Si $\mathfrak p \in B \cup V(J)$, alors ou bien $\mathfrak p \in T$,
ou bien $\mathfrak p \in V(J)$, et, dans les deux cas,
$V(\mathfrak p) \cap V(I) \subset V(J) \cap V(I)$, comme voulu.
Si $\mathfrak p \not \in B \cup V(J)$, alors
$V(\mathfrak p) \cap B$ est dense, donc infini, et l'on en conclut que
$\text{depth}(M_\mathfrak p) < s$ d'après
Cohomologie locale, lemme \ref{local-cohomology-lemma-depth-function}.
En fait, posons
$V(\mathfrak p) \cap B = \{\mathfrak p_\lambda\}_{\lambda \in \Lambda}$.
Choisissons $\mathfrak q_\lambda \in V(\mathfrak p_\lambda) \cap V(J) \cap V(I)$
comme dans (3).
Soit $\delta : \Spec(A) \to \mathbf{Z}$ la fonction de dimension
associée à un complexe dualisant $\omega_A^\bullet$ de $A$.
Puisque $\Lambda$ est infini et que $\delta$ est bornée,
il existe une partie infinie $\Lambda' \subset \Lambda$ sur laquelle
$\delta(\mathfrak q_\lambda)$ est constante. Pour
$\lambda \in \Lambda'$, on a
$$
\text{depth}(M_{\mathfrak p_\lambda}) +
\delta(\mathfrak p_\lambda) - \delta(\mathfrak q_\lambda) =
\text{depth}(M_{\mathfrak p_\lambda}) +
\dim((A/\mathfrak p_\lambda)_{\mathfrak q_\lambda})
\leq d + s
$$
d'après (3) et la définition de $B$. Par la semi-continuité de
la fonction $\text{depth} + \delta$ établie dans
Dualité pour les schémas, lemme \ref{duality-lemma-sitting-in-degrees},
on en conclut que
$$
\text{depth}(M_\mathfrak p) +
\dim((A/\mathfrak p)_{\mathfrak q_\lambda}) =
\text{depth}(M_\mathfrak p) + \delta(\mathfrak p) - \delta(\mathfrak q_\lambda)
\leq d + s
$$
Puisque l'on a aussi $\mathfrak p \not \in V(I)$, on déduit de (4) que
$\mathfrak p \in T$, c'est-à-dire
$V(\mathfrak p) \cap V(I) \subset V(J) \cap V(I)$. Cela achève la
démonstration de l'assertion I.

\medskip\noindent
Assertion II : $H^i_{J_0}(M) \to H^i_J(M)$ est un isomorphisme pour $i \leq s$
et $J' \subset J_0$ tel que $V(J') \cap V(I) = V(J) \cap V(I)$.

\medskip\noindent
Démonstration de l'assertion II. Choisissons $\mathfrak p \in V(J')$ n'appartenant pas à $V(J_0)$.
Il suffit de montrer que $H^i_{\mathfrak pA_\mathfrak p}(M_\mathfrak p) = 0$
pour $i \leq s$ ; voir
Cohomologie locale, lemme \ref{local-cohomology-lemma-isomorphism}.
Remarquons que $\mathfrak p \in T$. Puisque $\mathfrak p$ n'appartient pas à $B$,
on a donc $\text{depth}(M_\mathfrak p) > s$, et les groupes s'annulent d'après
Complexes dualisants, lemme \ref{dualizing-lemma-depth}.

\medskip\noindent
Assertion III. La dernière assertion du lemme est vraie.

\medskip\noindent
D'après l'assertion II, pour $i \leq s$, on a
$$
H^i_T(M) = H^i_{J_0}(M) = H^i_{J'}(M)
$$
pour tout idéal $J' \subset J_0$ tel que $V(J')  \cap V(I) = V(J) \cap V(I)$.
Voir Cohomologie locale, lemme \ref{local-cohomology-lemma-adjoint-ext}.
Vérifions les hypothèses de Cohomologie locale,
proposition \ref{local-cohomology-proposition-annihilator},
pour les parties $T \subset T \cup V(I)$, le module $M$ et l'entier $s$.
Il faut montrer qu'étant donnés $\mathfrak p \subset \mathfrak q$
tels que $\mathfrak p \not \in T \cup V(I)$ et $\mathfrak q \in T$,
on a
$$
\text{depth}_{A_\mathfrak p}(M_\mathfrak p) +
\dim((A/\mathfrak p)_\mathfrak q) > s
$$
Si $\text{depth}(M_\mathfrak p) \geq s$, c'est vrai puisque
la dimension de $(A/\mathfrak p)_\mathfrak q$ est au moins $1$.
On peut donc supposer $\text{depth}(M_\mathfrak p) < s$.
Si $\mathfrak q \in V(I)$, alors $\mathfrak q \in V(J) \cap V(I)$
et l'inégalité résulte de (4). Si $\mathfrak q \not \in V(I)$,
on peut utiliser (3) pour choisir
$\mathfrak q' \in V(\mathfrak q) \cap V(J) \cap V(I)$ tel que
$\dim((A/\mathfrak q)_{\mathfrak q'}) \leq d$.
L'hypothèse (4) donne alors
$$
\text{depth}_{A_\mathfrak p}(M_\mathfrak p) +
\dim((A/\mathfrak p)_{\mathfrak q'}) > s + d
$$
Puisque $A$ est caténaire, on en déduit l'inégalité voulue.
En appliquant Cohomologie locale,
proposition \ref{local-cohomology-proposition-annihilator}, on
trouve $J'' \subset A$ tel que $V(J'') \subset T \cup V(I)$
et tel que $J''$ annule $H^i_T(M)$ pour $i \leq s$.
On peut alors écrire $V(J'') \cup V(J_0) \cup V(I) = V(J'I)$
pour un certain $J' \subset J_0$ avec $V(J') \cap V(I) = V(J) \cap V(I)$.
En remplaçant $J_0$ par $J'$, on achève la démonstration.
\end{proof}

\begin{lemma}
\label{algebraization-lemma-kill-colimit-general}
Dans le lemme \ref{algebraization-lemma-kill-colimit-weak-general}, si, au lieu de la condition
vide (2), on suppose que
\begin{enumerate}
\item[(2')] si $\mathfrak p \in V(I)$, $\mathfrak p \not \in V(J) \cap V(I)$,
alors
$\text{depth}_{A_\mathfrak p}(M_\mathfrak p) +
\dim((A/\mathfrak p)_\mathfrak q) > s$
pour tout $\mathfrak q \in V(\mathfrak p) \cap V(J) \cap V(I)$,
\end{enumerate}
ces conditions entraînent aussi que $H^i_{J_0}(M)$ est un
$A$-module de type fini pour $i \leq s$.
\end{lemma}

\begin{proof}
Rappelons que $H^i_{J_0}(M) = H^i_T(M)$ ; voir la démonstration du
lemme \ref{algebraization-lemma-kill-colimit-weak-general}. Il suffit donc de
vérifier que, pour $\mathfrak p \not \in T$ et $\mathfrak q \in T$
avec $\mathfrak p \subset \mathfrak q$, on a
$\text{depth}_{A_\mathfrak p}(M_\mathfrak p) +
\dim((A/\mathfrak p)_\mathfrak q) > s$ ; voir Cohomologie locale,
proposition \ref{local-cohomology-proposition-finiteness}.
La condition (2') l'assure pour $\mathfrak p \in V(I)$.
Puisque l'on sait que $H^i_T(M)$ est annulé par une puissance de $IJ_0$,
on sait que la condition vaut si $\mathfrak p \not \in V(IJ_0)$
d'après Cohomologie locale, proposition \ref{local-cohomology-proposition-annihilator}.
Tous les cas sont ainsi couverts, et la démonstration est achevée.
\end{proof}

\begin{lemma}
\label{algebraization-lemma-kill-colimit-support-general}
Si, dans le lemme \ref{algebraization-lemma-kill-colimit-weak-general}, on suppose en outre que
\begin{enumerate}
\item[(6)] si $\mathfrak p \not \in V(I)$, $\mathfrak p \in T$, alors
$\text{depth}_{A_\mathfrak p}(M_\mathfrak p) > s$,
\end{enumerate}
alors $H^i_{J_0}(M) = H^i_J(M) = H^i_{J + I}(M)$ pour $i \leq s$, et ces
modules sont annulés par une puissance de $I$.
\end{lemma}

\begin{proof}
Choisissons $\mathfrak p \in V(J)$ ou $\mathfrak p \in V(J_0)$, mais
$\mathfrak p \not \in V(J + I) = V(J_0 + I)$.
Il suffit de montrer que $H^i_{\mathfrak pA_\mathfrak p}(M_\mathfrak p) = 0$
pour $i \leq s$ ; voir
Cohomologie locale, lemme \ref{local-cohomology-lemma-isomorphism}.
Ces groupes s'annulent d'après la condition (6) et
Complexes dualisants, lemme \ref{dualizing-lemma-depth}.
La dernière assertion résulte de
Cohomologie locale, proposition \ref{local-cohomology-proposition-annihilator}.
\end{proof}

\begin{lemma}
\label{algebraization-lemma-algebraize-local-cohomology-general}
Soient $I, J$ des idéaux d'un anneau noethérien $A$.
Soit $M$ un $A$-module de type fini.
Soient $s$ et $d$ des entiers. Avec $T$ comme dans
(\ref{algebraization-equation-associated-subset}), supposons que
\begin{enumerate}
\item $A$ est complet pour la topologie $I$-adique et possède un complexe dualisant,
\item si $\mathfrak p \in V(I)$, aucune condition n'est imposée,
\item $\text{cd}(A, I) \leq d$,
\item si $\mathfrak p \not \in V(I)$, $\mathfrak p \not \in T$, alors
$$
\text{depth}_{A_\mathfrak p}(M_\mathfrak p) \geq s
\quad\text{ou}\quad
\text{depth}_{A_\mathfrak p}(M_\mathfrak p) +
\dim((A/\mathfrak p)_\mathfrak q) > d + s
$$
pour tout $\mathfrak q \in V(\mathfrak p) \cap V(J) \cap V(I)$,
\item si $\mathfrak p \not \in V(I)$, $\mathfrak p \not \in T$,
$V(\mathfrak p) \cap V(J) \cap V(I) \not = \emptyset$ et
$\text{depth}(M_\mathfrak p) < s$, alors l'une
des conditions suivantes est satisfaite\footnote{Notre méthode
impose cette condition supplémentaire. Nous y reviendrons
(insérer une référence ultérieure).} :
\begin{enumerate}
\item $\dim(\text{Supp}(M_\mathfrak p)) < s + 2$\footnote{Par exemple,
si $M$ satisfait la condition $(S_s)$ de Serre
sur le complémentaire de $V(I) \cup T$.}, ou
\item  $\delta(\mathfrak p) > d + \delta_{max} - 1$
où $\delta$ est une fonction de dimension et $\delta_{max}$
est le maximum de $\delta$ sur $V(J) \cap V(I)$, ou
\item $\text{depth}_{A_\mathfrak p}(M_\mathfrak p) +
\dim((A/\mathfrak p)_\mathfrak q) > d + s + \delta_{max} - \delta_{min} - 2$
pour tout $\mathfrak q \in V(\mathfrak p) \cap V(J) \cap V(I)$.
\end{enumerate}
\end{enumerate}
Alors il existe un idéal $J_0 \subset J$ tel que
$V(J_0) \cap V(I) = V(J) \cap V(I)$
et tel que, pour tout $J' \subset J_0$ avec
$V(J') \cap V(I) = V(J) \cap V(I)$, l'application
$$
R\Gamma_{J'}(M) \longrightarrow R\Gamma_J(M)^\wedge
$$
induise un isomorphisme en cohomologie en degrés $\leq s$.
Ici ${}^\wedge$ désigne la complétion $I$-adique dérivée.
\end{lemma}

\noindent
Nous invitons le lecteur à lire d'abord la démonstration dans le cas local
(lemme \ref{algebraization-lemma-algebraize-local-cohomology}), car elle explique la structure
de la démonstration sans avoir à traiter toutes les inégalités.

\begin{proof}
Pour un idéal $\mathfrak a \subset A$, on a
$R\Gamma_\mathfrak a = R\Gamma_{V(\mathfrak a)}$ ; voir
Complexes dualisants, lemme \ref{dualizing-lemma-local-cohomology-noetherian}.
Remarquons ensuite que
$$
R\Gamma_J(M)^\wedge =
R\Gamma_I(R\Gamma_J(M))^\wedge =
R\Gamma_{I + J}(M)^\wedge =
R\Gamma_{I + J'}(M)^\wedge =
R\Gamma_I(R\Gamma_{J'}(M))^\wedge =
R\Gamma_{J'}(M)^\wedge
$$
d'après Complexes dualisants, lemmes \ref{dualizing-lemma-local-cohomology-ss} et
\ref{dualizing-lemma-complete-and-local}.
Cela explique comment on définit la flèche de l'énoncé du lemme.

\medskip\noindent
Nous affirmons que les hypothèses du lemme \ref{algebraization-lemma-kill-colimit-weak-general}
résultent de nos hypothèses actuelles sur $M$.
La seule chose à vérifier est l'hypothèse (3).
Soit donc $\mathfrak p \not \in V(I)$, $\mathfrak p \in T$.
Alors $V(\mathfrak p) \cap V(I)$ est non vide, puisque $I$ est
contenu dans le radical de Jacobson de $A$
(Algèbre, lemme \ref{algebra-lemma-radical-completion}).
Puisque $\mathfrak p \in T$, on a
$V(\mathfrak p) \cap V(I) = V(\mathfrak p) \cap V(J) \cap V(I)$.
Soit $\mathfrak q \in V(\mathfrak p) \cap V(I)$ le
point générique d'une composante irréductible.
On a $\text{cd}(A_\mathfrak q, I_\mathfrak q) \leq d$
d'après Cohomologie locale, lemme \ref{local-cohomology-lemma-cd-local}.
On a $V(\mathfrak pA_\mathfrak q) \cap V(I_\mathfrak q) =
\{\mathfrak qA_\mathfrak q\}$ par notre choix de $\mathfrak q$,
et l'on en conclut que $\dim((A/\mathfrak p)_\mathfrak q) \leq d$
d'après Cohomologie locale, lemme \ref{local-cohomology-lemma-cd-bound-dim-local}.

\medskip\noindent
Remarquons que le lemme vaut pour $s < 0$. Ce cas n'est pas trivial, car
il n'est pas clair a priori que $H^i(R\Gamma_J(M)^\wedge)$
soit nul pour $i < 0$. Cette annulation a toutefois été établie dans
Complexes dualisants, lemme \ref{dualizing-lemma-completion-local}.
Nous démontrerons le lemme par récurrence pour $s \geq 0$.

\medskip\noindent
Le lemme pour $s = 0$ résulte immédiatement de
la conclusion du lemme \ref{algebraization-lemma-kill-colimit-weak-general}
et de Complexes dualisants, lemme \ref{dualizing-lemma-completion-local-H0}.

\medskip\noindent
Supposons $s > 0$ et supposons le lemme démontré pour les valeurs plus petites de $s$.
Soit $M' \subset M$ le sous-module maximal dont le support est contenu
dans $V(I) \cup T$. Alors $M'$ est un $A$-module de type fini dont le support
est contenu dans $V(J') \cup V(I)$ pour un certain idéal $J' \subset J$
tel que $V(J') \cap V(I) = V(J) \cap V(I)$.
Nous affirmons que
$$
R\Gamma_{J'}(M') \to R\Gamma_J(M')^\wedge
$$
est un isomorphisme quel que soit le choix de $J'$.
En effet, on peut choisir une suite exacte courte
$0 \to M_1 \oplus M_2 \to M' \to N \to 0$, où
$M_1$ est annulé par une puissance de $J'$, où $M_2$ est annulé
par une puissance de $I$, et où $N$ est annulé par une puissance de $I + J'$.
Il suffit donc de montrer que l'assertion vaut pour $M_1$, $M_2$ et $N$.
Dans le cas de $M_1$, on voit que $R\Gamma_{J'}(M_1) = M_1$ et,
puisque $M_1$ est un $A$-module de type fini et complet pour la topologie $I$-adique,
on a $M_1^\wedge = M_1$. Cela prouve l'assertion pour $M_1$
d'après les remarques initiales de la démonstration. Dans le cas de $M_2$, on voit que
$H^i_J(M_2) = H^i_{I + J}(M) = H^i_{I + J'}(M) = H^i_{J'}(M_2)$
est annulé par une puissance de $I$, et est donc complet au sens dérivé.
L'assertion vaut donc aussi dans ce cas. Pour $N$, on peut utiliser l'un ou l'autre
des arguments que l'on vient de donner. En considérant la suite exacte courte
$0 \to M' \to M \to M/M' \to 0$
on voit qu'il suffit de démontrer le lemme pour $M/M'$.
On peut donc supposer $\text{Ass}(M) \cap (V(I) \cup T) = \emptyset$.

\medskip\noindent
Soit $\mathfrak p \in \text{Ass}(M)$ tel que
$V(\mathfrak p) \cap V(J) \cap V(I) = \emptyset$.
Puisque $I$ est contenu dans le radical de Jacobson de $A$, cela entraîne
que $V(\mathfrak p) \cap V(J') = \emptyset$ pour tout
$J' \subset J$ tel que $V(J') \cap V(I) = V(J) \cap V(I)$.
En posant donc $N = H^0_\mathfrak p(M)$, on voit que
$R\Gamma_J(N) = R\Gamma_{J'}(N) = 0$ pour tout
$J' \subset J$ tel que $V(J') \cap V(I) = V(J) \cap V(I)$.
En particulier, $R\Gamma_J(N)^\wedge = 0$.
On peut donc remplacer $M$ par $M/N$, puisque cela ne modifie la
structure de $M$ qu'aux idéaux premiers qui ne jouent aucun
rôle dans les conditions (4) ou (5). En répétant l'opération, on peut supposer que
$V(\mathfrak p) \cap V(J) \cap V(I) \not = \emptyset$
pour tout $\mathfrak p \in \text{Ass}(M)$.

\medskip\noindent
Supposons $\text{Ass}(M) \cap (V(I) \cup T) = \emptyset$ et que
$V(\mathfrak p) \cap V(J) \cap V(I) \not = \emptyset$
pour tout $\mathfrak p \in \text{Ass}(M)$.
Soit $\mathfrak p \in \text{Ass}(M)$. Nous voulons montrer que l'on peut appliquer le
lemme \ref{algebraization-lemma-kill-completion-general}.
C'est dans cette vérification que nous utiliserons la condition supplémentaire
(5). Choisissons $\mathfrak p' \subset \mathfrak p$
et $\mathfrak q' \subset V(\mathfrak p) \cap V(J) \cap V(I)$.
\begin{enumerate}
\item Si $M_{\mathfrak p'} = 0$, alors
$\text{depth}(M_{\mathfrak p'}) = \infty$ et
$\text{depth}(M_{\mathfrak p'}) +
\dim((A/\mathfrak p')_{\mathfrak q'}) > d + s$.
\item Si $\text{depth}(M_{\mathfrak p'}) < s$, alors
$\text{depth}(M_{\mathfrak p'}) +
\dim((A/\mathfrak p')_{\mathfrak q'}) > d + s$ d'après (4).
\end{enumerate}
Dans les cas restants, on a $M_{\mathfrak p'} \not = 0$ et
$\text{depth}(M_{\mathfrak p'}) \geq s$. En particulier, on voit que
$\mathfrak p'$ appartient au support de $M$, et l'on peut choisir
$\mathfrak p'' \subset \mathfrak p'$ avec $\mathfrak p'' \in \text{Ass}(M)$.
\begin{enumerate}
\item[(a)] Remarquons que
$\dim((A/\mathfrak p'')_{\mathfrak p'}) \geq \text{depth}(M_{\mathfrak p'})$
d'après Algèbre, lemme \ref{algebra-lemma-depth-dim-associated-primes}.
Si l'égalité a lieu, alors on a
$$
\text{depth}(M_{\mathfrak p'}) + \dim((A/\mathfrak p')_{\mathfrak q'}) =
\text{depth}(M_{\mathfrak p''}) + \dim((A/\mathfrak p'')_{\mathfrak q'})
> s + d
$$
d'après (4) appliquée à $\mathfrak p''$, et l'on a terminé. Cela signifie que la
seule difficulté se présente lorsque
$\dim((A/\mathfrak p'')_{\mathfrak p'}) > \text{depth}(M_{\mathfrak p'})$.
Cela entraîne que $\dim(M_\mathfrak p) \geq s + 2$.
Ainsi, si (5)(a) est satisfaite, ce cas ne se présente pas.
\item[(b)] Si (5)(b) est satisfaite, on obtient
$$
\text{depth}(M_{\mathfrak p'}) + \dim((A/\mathfrak p')_{\mathfrak q'})
\geq s + \delta(\mathfrak p') - \delta(\mathfrak q')
\geq s + 1 + \delta(\mathfrak p) - \delta_{max}
> s + d
$$
comme voulu.
\item[(c)] Si (5)(c) est satisfaite, on obtient
\begin{align*}
\text{depth}(M_{\mathfrak p'}) + \dim((A/\mathfrak p')_{\mathfrak q'})
& \geq
s + \delta(\mathfrak p') - \delta(\mathfrak q') \\
& \geq
s + 1 + \delta(\mathfrak p) - \delta(\mathfrak q') \\
& =
s + 1 + \delta(\mathfrak p) - \delta(\mathfrak q) +
\delta(\mathfrak q) - \delta(\mathfrak q') \\
& >
s + 1 + (s + d + \delta_{max} - \delta_{min} - 2) +
\delta(\mathfrak q) - \delta(\mathfrak q') \\
& \geq 
2s + d - 1 \geq s + d
\end{align*}
comme voulu. Remarquons que cet argument fonctionne parce que
nous savons qu'un idéal premier $\mathfrak q \in V(\mathfrak p) \cap V(J) \cap V(I)$
existe.
\end{enumerate}
Nous sommes maintenant prêts à effectuer l'étape de récurrence.

\medskip\noindent
Choisissons un idéal $J_0$ comme dans le lemme \ref{algebraization-lemma-kill-colimit-weak-general}
et un entier $t > 0$ tel que $(J_0I)^t$ annule $H^s_J(M)$.
Les hypothèses du lemme \ref{algebraization-lemma-kill-completion-general}
sont satisfaites pour tout $\mathfrak p \in \text{Ass}(M)$
(voir le paragraphe précédent).
Ainsi, l'annulateur $\mathfrak a \subset A$ de
$H^s(R\Gamma_J(M)^\wedge)$
n'est pas contenu dans $\mathfrak p$ pour $\mathfrak p \in \text{Ass}(M)$.
On peut donc trouver un $f \in \mathfrak a(J_0I)^t$
qui n'appartient à aucun idéal premier associé de $M$ et qui annule
à la fois $H^s(R\Gamma_J(M)^\wedge)$ et $H^s_J(M)$.
Alors $f$ est un non-diviseur de zéro dans $M$, et l'on peut considérer la
suite exacte courte
$$
0 \to M \xrightarrow{f} M \to M/fM \to 0
$$
Notre choix de $f$ montre que l'on obtient
$$
\xymatrix{
H^{s - 1}_{J'}(M) \ar[d] \ar[r] &
H^{s - 1}_{J'}(M/fM) \ar[d] \ar[r] &
H^s_{J'}(M) \ar[d] \ar[r] & 0 \\
H^{s - 1}(R\Gamma_J(M)^\wedge) \ar[r] &
H^{s - 1}(R\Gamma_J(M/fM)^\wedge) \ar[r] &
H^s(R\Gamma_J(M)^\wedge) \ar[r] & 0
}
$$
pour tout $J' \subset J_0$ tel que $V(J') \cap V(I) = V(J) \cap V(I)$.
Ainsi, si l'on choisit $J'$ de sorte qu'il convienne pour
$M$, $M/fM$ et $s - 1$ (ce qui est possible par l'hypothèse de récurrence ---
voir le paragraphe suivant), on en conclut que le lemme est vrai.

\medskip\noindent
Pour achever la démonstration, il faut montrer que le module
$M/fM$ satisfait les hypothèses (4) et (5) pour $s - 1$.
Soit donc $\mathfrak p$ un idéal premier du support
de $M/fM$, avec $\text{depth}((M/fM)_\mathfrak p) < s - 1$
et avec $V(\mathfrak p) \cap V(J) \cap V(I)$ non vide.
Alors $\dim(M_\mathfrak p) = \dim((M/fM)_\mathfrak p) + 1$
et $\text{depth}(M_\mathfrak p) = \text{depth}((M/fM)_\mathfrak p) + 1$.
En particulier, on sait que (4) et (5) sont satisfaites pour $\mathfrak p$ et $M$
avec la valeur initiale de $s$.
Les inégalités voulues s'en déduisent alors immédiatement.
\end{proof}

\begin{example}
\label{algebraization-example-no-ML}
Dans le lemme \ref{algebraization-lemma-algebraize-local-cohomology-general},
nous ne savons pas si les systèmes projectifs $H^i_J(M/I^nM)$ satisfont la
condition de Mittag-Leffler.
Par exemple, supposons $A = \mathbf{Z}_p[[x, y]]$, $I = (p)$,
$J = (p, x)$ et $M = A/(xy - p)$. Alors l'image de
$H^0_J(M/p^nM) \to H^0_J(M/pM)$
est l'idéal engendré par $y^n$ dans $M/pM = A/(p, xy)$.
\end{example}





\section{Algébrisation de la cohomologie locale, II}
\label{algebraization-section-algebraization-punctured}

\noindent
Dans cette section, nous reprenons les arguments de la
section \ref{algebraization-section-algebraization-sections-general}
lorsque $(A, \mathfrak m)$ est un anneau local et que l'on considère la cohomologie locale
$R\Gamma_\mathfrak m$ relativement à $\mathfrak m$. Comme précédemment, notre
outil principal est le lemme suivant.

\begin{lemma}
\label{algebraization-lemma-kill-completion}
Soit $(A, \mathfrak m)$ un anneau local noethérien.
Soit $I \subset A$ un idéal. Soit $M$ un $A$-module de type fini et
soit $\mathfrak p \subset A$ un idéal premier. Soient $s$ et $d$ des entiers. Supposons que
\begin{enumerate}
\item $A$ possède un complexe dualisant,
\item $\text{cd}(A, I) \leq d$, et
\item
$\text{depth}_{A_\mathfrak p}(M_\mathfrak p) + \dim(A/\mathfrak p) > d + s$.
\end{enumerate}
Alors il existe un $f \in A \setminus \mathfrak p$ qui annule
$H^i(R\Gamma_\mathfrak m(M)^\wedge)$ pour $i \leq s$, où ${}^\wedge$
désigne la complétion $I$-adique.
\end{lemma}

\begin{proof}
D'après Cohomologie locale, lemme
\ref{local-cohomology-lemma-sitting-in-degrees}
la fonction
$$
\mathfrak p' \longmapsto
\text{depth}_{A_{\mathfrak p'}}(M_{\mathfrak p'}) + \dim(A/\mathfrak p')
$$
est semi-continue inférieurement sur $\Spec(A)$. La valeur
de cette fonction en tout $\mathfrak p' \subset \mathfrak p$
est donc $> s + d$. Notre lemme est ainsi un cas particulier du
lemme \ref{algebraization-lemma-kill-completion-general},
pourvu que $\mathfrak p \not = \mathfrak m$.
Si $\mathfrak p = \mathfrak m$,
alors $H^i_\mathfrak m(M) = 0$ pour $i \leq s + d$ d'après
la relation entre la profondeur et la cohomologie locale
(Complexes dualisants, lemme \ref{dualizing-lemma-depth}).
Ainsi, l'argument donné dans la démonstration du
lemme \ref{algebraization-lemma-kill-completion-general}
montre que $H^i(R\Gamma_\mathfrak m(M)^\wedge) = 0$
pour $i \leq s$ dans ce cas (dégénéré).
\end{proof}

\begin{lemma}
\label{algebraization-lemma-kill-colimit-weak}
Soit $(A, \mathfrak m)$ un anneau local noethérien.
Soit $I \subset A$ un idéal. Soit $M$ un $A$-module de type fini.
Soient $s$ et $d$ des entiers. Supposons que
\begin{enumerate}
\item $A$ possède un complexe dualisant,
\item si $\mathfrak p \in V(I)$, aucune condition n'est imposée,
\item si $\mathfrak p \not \in V(I)$ et
$V(\mathfrak p) \cap V(I) = \{\mathfrak m\}$, alors
$\dim(A/\mathfrak p) \leq d$,
\item si $\mathfrak p \not \in V(I)$ et
$V(\mathfrak p) \cap V(I) \not = \{\mathfrak m\}$, alors
$$
\text{depth}_{A_\mathfrak p}(M_\mathfrak p) \geq s
\quad\text{ou}\quad
\text{depth}_{A_\mathfrak p}(M_\mathfrak p) + \dim(A/\mathfrak p) > d + s
$$
\end{enumerate}
Alors il existe un idéal $J_0 \subset A$ tel que
$V(J_0) \cap V(I) = \{\mathfrak m\}$ et tel que, pour tout $J \subset J_0$ avec
$V(J) \cap V(I) = \{\mathfrak m\}$, l'application
$$
R\Gamma_J(M) \longrightarrow R\Gamma_{J_0}(M)
$$
induise un isomorphisme en cohomologie en degrés $\leq s$ ;
de plus, ces modules sont annulés par une puissance de $J_0I$.
\end{lemma}

\begin{proof}
C'est un cas particulier du lemme \ref{algebraization-lemma-kill-colimit-weak-general}.
\end{proof}

\begin{lemma}
\label{algebraization-lemma-kill-colimit}
Dans le lemme \ref{algebraization-lemma-kill-colimit-weak}, si, au lieu de la condition
vide (2), on suppose que
\begin{enumerate}
\item[(2')] si $\mathfrak p \in V(I)$ et $\mathfrak p \not = \mathfrak m$,
alors $\text{depth}_{A_\mathfrak p}(M_\mathfrak p) + \dim(A/\mathfrak p) > s$,
\end{enumerate}
alors ces conditions entraînent aussi que $H^i_{J_0}(M)$ est un
$A$-module de type fini pour $i \leq s$.
\end{lemma}

\begin{proof}
C'est un cas particulier du lemme \ref{algebraization-lemma-kill-colimit-general}.
\end{proof}

\begin{lemma}
\label{algebraization-lemma-kill-colimit-support}
Si, dans le lemme \ref{algebraization-lemma-kill-colimit-weak}, on suppose en outre que
\begin{enumerate}
\item[(6)] si $\mathfrak p \not \in V(I)$ et
$V(\mathfrak p) \cap V(I) = \{\mathfrak m\}$, alors
$\text{depth}_{A_\mathfrak p}(M_\mathfrak p) > s$,
\end{enumerate}
alors $H^i_{J_0}(M) = H^i_J(M) = H^i_\mathfrak m(M)$ pour $i \leq s$,
et ces modules sont annulés par une puissance de $I$.
\end{lemma}

\begin{proof}
C'est un cas particulier du lemme \ref{algebraization-lemma-kill-colimit-support-general}.
\end{proof}

\begin{lemma}
\label{algebraization-lemma-algebraize-local-cohomology}
Soit $(A, \mathfrak m)$ un anneau local noethérien.
Soit $I \subset A$ un idéal. Soit $M$ un $A$-module de type fini.
Soient $s$ et $d$ des entiers. Supposons que
\begin{enumerate}
\item $A$ est complet pour la topologie $I$-adique et possède un complexe dualisant,
\item si $\mathfrak p \in V(I)$, aucune condition n'est imposée,
\item $\text{cd}(A, I) \leq d$,
\item si $\mathfrak p \not \in V(I)$ et
$V(\mathfrak p) \cap V(I) \not = \{\mathfrak m\}$, alors
$$
\text{depth}_{A_\mathfrak p}(M_\mathfrak p) \geq s
\quad\text{ou}\quad
\text{depth}_{A_\mathfrak p}(M_\mathfrak p) + \dim(A/\mathfrak p) > d + s
$$
\end{enumerate}
Alors il existe un idéal $J_0 \subset A$ tel que
$V(J_0) \cap V(I) = \{\mathfrak m\}$ et tel que, pour tout $J \subset J_0$ avec
$V(J) \cap V(I) = \{\mathfrak m\}$, l'application
$$
R\Gamma_J(M) \longrightarrow
R\Gamma_J(M)^\wedge = R\Gamma_\mathfrak m(M)^\wedge
$$
induise un isomorphisme en cohomologie en degrés $\leq s$.
Ici ${}^\wedge$ désigne la complétion $I$-adique dérivée.
\end{lemma}

\begin{proof}
Ce lemme est un cas particulier du
lemme \ref{algebraization-lemma-algebraize-local-cohomology-general},
puisque la condition (5)(c) résulte de la condition (4),
car $\delta_{max} = \delta_{min} = \delta(\mathfrak m)$.
Nous allons donner la démonstration de ce cas particulier important,
car elle est un peu plus simple (il y a moins de choses à vérifier).

\medskip\noindent
Il n'y a pas de différence entre $R\Gamma_\mathfrak a$ et
$R\Gamma_{V(\mathfrak a)}$ dans la situation présente ; voir
Complexes dualisants, lemme \ref{dualizing-lemma-local-cohomology-noetherian}.
Remarquons ensuite que
$$
R\Gamma_\mathfrak m(M)^\wedge =
R\Gamma_I(R\Gamma_J(M))^\wedge =
R\Gamma_J(M)^\wedge
$$
d'après Complexes dualisants, lemmes \ref{dualizing-lemma-local-cohomology-ss} et
\ref{dualizing-lemma-complete-and-local},
ce qui explique le signe d'égalité dans l'énoncé du lemme.

\medskip\noindent
Remarquons que le lemme vaut pour $s < 0$. Ce cas n'est pas trivial, car
il n'est pas clair a priori que $H^s(R\Gamma_\mathfrak m(M)^\wedge)$
soit nul pour $s$ négatif. Cette annulation a toutefois été établie
dans le lemme \ref{algebraization-lemma-local-cohomology-derived-completion}.
Nous démontrerons le lemme par récurrence pour $s \geq 0$.

\medskip\noindent
Les hypothèses du lemme \ref{algebraization-lemma-kill-colimit-weak}
sont satisfaites d'après Cohomologie locale, lemme
\ref{local-cohomology-lemma-cd-bound-dim-local}.
Le lemme pour $s = 0$ résulte du lemme \ref{algebraization-lemma-kill-colimit-weak} et de
Complexes dualisants, lemme \ref{dualizing-lemma-completion-local-H0}.

\medskip\noindent
Supposons $s > 0$ et le lemme vrai pour les valeurs plus petites de $s$.
Soit $M' \subset M$ le sous-module des éléments dont le
support est contenu dans $V(I) \cup V(J)$ pour un certain
idéal $J$ tel que $V(J) \cap V(I) = \{\mathfrak m\}$.
Alors $M'$ est un $A$-module de type fini.
Nous affirmons que
$$
R\Gamma_J(M') \to R\Gamma_\mathfrak m(M')^\wedge
$$
est un isomorphisme quel que soit le choix de $J$.
En effet, pour tout module de ce type, il existe une suite exacte courte
$0 \to M_1 \oplus M_2 \to M' \to N \to 0$, où
$M_1$ est annulé par une puissance de $J$, où $M_2$ est annulé
par une puissance de $I$ et où $N$ est annulé par une puissance de $\mathfrak m$.
Dans le cas de $M_1$, on voit que $R\Gamma_J(M_1) = M_1$ et,
puisque $M_1$ est un $A$-module de type fini et complet pour la topologie $I$-adique,
on a $M_1^\wedge = M_1$. L'assertion vaut donc pour $M_1$.
Dans le cas de $M_2$, on voit que $H^i_J(M_2)$ est annulé
par une puissance de $I$, et est donc complet au sens dérivé. Ainsi, l'assertion vaut
pour $M_2$. Les mêmes arguments montrent que l'assertion vaut pour $N$,
et l'on en conclut qu'elle est vraie. En considérant la
suite exacte courte $0 \to M' \to M \to M/M' \to 0$,
on voit qu'il suffit de démontrer le lemme pour $M/M'$.
On peut donc supposer que $\mathfrak p \in \text{Ass}(M)$
entraîne $V(\mathfrak p) \cap V(I) \not = \{\mathfrak m\}$, c'est-à-dire que
$\mathfrak p$ est un idéal premier comme dans (4).

\medskip\noindent
Choisissons un idéal $J_0$ comme dans le lemme \ref{algebraization-lemma-kill-colimit-weak}
et un entier $t > 0$ tel que $(J_0I)^t$ annule $H^s_J(M)$.
Ici, $J$ désigne un idéal arbitraire $J \subset J_0$ tel que
$V(J) \cap V(I) = \{\mathfrak m\}$.
Les hypothèses du lemme \ref{algebraization-lemma-kill-completion}
sont satisfaites pour tout $\mathfrak p \in \text{Ass}(M)$
(voir le paragraphe précédent). Ainsi, l'annulateur $\mathfrak a \subset A$ de
$H^s(R\Gamma_\mathfrak m(M)^\wedge)$
n'est pas contenu dans $\mathfrak p$ pour $\mathfrak p \in \text{Ass}(M)$.
On peut donc trouver un $f \in \mathfrak a(J_0I)^t$
qui n'appartient à aucun idéal premier associé de $M$ et qui annule
à la fois $H^s(R\Gamma_\mathfrak m(M)^\wedge)$ et $H^s_J(M)$.
Alors $f$ est un non-diviseur de zéro dans $M$, et l'on peut considérer la
suite exacte courte
$$
0 \to M \xrightarrow{f} M \to M/fM \to 0
$$
Notre choix de $f$ montre que l'on obtient
$$
\xymatrix{
H^{s - 1}_J(M) \ar[d] \ar[r] &
H^{s - 1}_J(M/fM) \ar[d] \ar[r] &
H^s_J(M) \ar[d] \ar[r] & 0 \\
H^{s - 1}(R\Gamma_\mathfrak m(M)^\wedge) \ar[r] &
H^{s - 1}(R\Gamma_\mathfrak m(M/fM)^\wedge) \ar[r] &
H^s(R\Gamma_\mathfrak m(M)^\wedge) \ar[r] & 0
}
$$
pour tout $J \subset J_0$ tel que $V(J) \cap V(I) = \{\mathfrak m\}$.
Ainsi, si l'on choisit $J$ de sorte qu'il convienne pour
$M$, $M/fM$ et $s - 1$ (ce qui est possible par l'hypothèse de récurrence),
on en conclut que le lemme est vrai.
\end{proof}








\section{Algébrisation de la cohomologie locale, III}
\label{algebraization-section-bootstrap}

\noindent
Dans cette section, nous nous appuyons sur les résultats des
sections \ref{algebraization-section-algebraization-sections-general} et
\ref{algebraization-section-algebraization-punctured}
pour obtenir un résultat plus fort dans la situation suivante.

\begin{situation}
\label{algebraization-situation-bootstrap}
Ici, $A$ est un anneau noethérien. Nous avons un idéal $I \subset A$,
un $A$-module de type fini $M$ et un sous-ensemble $T \subset V(I)$ stable par
spécialisation. Nous avons des entiers $s$ et $d$. Nous supposons que
\begin{enumerate}
\item[(1)] $A$ possède un complexe dualisant,
\item[(3)] $\text{cd}(A, I) \leq d$,
\item[(4)] pour des idéaux premiers $\mathfrak p \subset \mathfrak r \subset \mathfrak q$
tels que $\mathfrak p \not \in V(I)$,
$\mathfrak r \in V(I) \setminus T$,
$\mathfrak q \in T$, on a
$$
\text{depth}_{A_\mathfrak p}(M_\mathfrak p) \geq s
\quad\text{ou}\quad
\text{depth}_{A_\mathfrak p}(M_\mathfrak p) +
\dim((A/\mathfrak p)_\mathfrak q) > d + s
$$
\item[(6)] pour $\mathfrak q \in T$, si
$A', \mathfrak m', I', M'$ désignent les complétions $I$-adiques usuelles
de $A_\mathfrak q, \mathfrak qA_\mathfrak q, I_\mathfrak q, M_\mathfrak q$,
on a
$$
\text{depth}(M'_{\mathfrak p'}) > s
$$
pour tout $\mathfrak p' \in \Spec(A') \setminus V(I')$ tel que
$V(\mathfrak p') \cap V(I') = \{\mathfrak m'\}$.
\end{enumerate}
\end{situation}

\noindent
Le lemme suivant explique pourquoi, dans la situation \ref{algebraization-situation-bootstrap},
il suffit de considérer les triplets
$\mathfrak p \subset \mathfrak r \subset \mathfrak q$ d'idéaux premiers de (4),
bien que l'hypothèse elle-même ne fasse intervenir que $\mathfrak p$ et $\mathfrak q$.

\begin{lemma}
\label{algebraization-lemma-helper-bootstrap}
Dans la situation \ref{algebraization-situation-bootstrap}, soient $\mathfrak p \subset \mathfrak q$
des idéaux premiers de $A$ tels que $\mathfrak p \not \in V(I)$ et
$\mathfrak q \in T$. S'il n'existe aucun
$\mathfrak r \in V(I) \setminus T$ tel que
$\mathfrak p \subset \mathfrak r \subset \mathfrak q$
alors $\text{depth}(M_\mathfrak p) > s$.
\end{lemma}

\begin{proof}
Choisissons $\mathfrak q' \in T$ tel que
$\mathfrak p \subset \mathfrak q' \subset \mathfrak q$
et qu'il n'existe aucun idéal premier de $T$ strictement
compris entre $\mathfrak p$ et $\mathfrak q'$. Pour démontrer le lemme,
on peut, et l'on va, remplacer $\mathfrak q$ par $\mathfrak q'$.
Ensuite, soit $\mathfrak p' \subset A_\mathfrak q$ l'idéal premier correspondant à
$\mathfrak p$. On obtient alors
$V(\mathfrak p') \cap V(IA_\mathfrak q) = \{\mathfrak q A_\mathfrak q\}$
du fait qu'il n'existe pas d'idéal premier $\mathfrak r$ comme dans le lemme.
Soient $A', I', \mathfrak m', M'$ les complétions $I$-adiques de
$A_\mathfrak q, I_\mathfrak q, \mathfrak qA_\mathfrak q, M_\mathfrak q$.
Comme $A_\mathfrak q \to A'$ est fidèlement plat
(Algèbre, lemme \ref{algebra-lemma-completion-faithfully-flat}),
on peut choisir $\mathfrak p'' \subset A'$ au-dessus de $\mathfrak p'$
tel que $\dim(A'_{\mathfrak p''}/\mathfrak p' A'_{\mathfrak p''}) = 0$.
On voit alors que
$$
\text{depth}(M'_{\mathfrak p''}) =
\text{depth}((M_\mathfrak q \otimes_{A_\mathfrak q} A')_{\mathfrak p''}) =
\text{depth}(M_\mathfrak p \otimes_{A_\mathfrak p} A'_{\mathfrak p''}) =
\text{depth}(M_\mathfrak p)
$$
par platitude de $A \to A'$ et par notre choix de $\mathfrak p''$, voir
Algèbre, lemme \ref{algebra-lemma-apply-grothendieck-module}.
Comme $\mathfrak p''$ est au-dessus de $\mathfrak p'$, on a
$V(\mathfrak p'') \cap V(I') = \{\mathfrak m'\}$. Ainsi, la
condition (6) de la situation \ref{algebraization-situation-bootstrap} entraîne que
$\text{depth}(M'_{\mathfrak p''}) > s$, ce qui achève la démonstration.
\end{proof}

\noindent
Le lemme fastidieux suivant explique les relations entre les diverses
familles de conditions que l'on pourrait imposer.

\begin{lemma}
\label{algebraization-lemma-bootstrap-inherited}
Dans la situation \ref{algebraization-situation-bootstrap}, on a
\begin{enumerate}
\item[(E)] si $T' \subset T$ est un sous-ensemble plus petit, stable par spécialisation, alors
$A, I, T', M$ satisfont les hypothèses de la situation \ref{algebraization-situation-bootstrap},
\item[(F)] si $S \subset A$ est une partie multiplicative, alors
$S^{-1}A, S^{-1}I, T', S^{-1}M$
satisfont les hypothèses de la situation \ref{algebraization-situation-bootstrap},
où $T' \subset V(S^{-1}I)$ est l'image réciproque de $T$,
\item[(G)] le quadruplet $A', I', T', M'$
satisfait les hypothèses de la situation \ref{algebraization-situation-bootstrap},
où $A', I', M'$ sont les complétions $I$-adiques usuelles de $A, I, M$
et $T' \subset V(I')$ est l'image réciproque de $T$.
\end{enumerate}
Soit $I \subset \mathfrak a \subset A$ un idéal tel que
$V(\mathfrak a) \subset T$. Alors
\begin{enumerate}
\item[(A)] si $I$ est contenu dans le radical de Jacobson de $A$,
alors toutes les hypothèses des
lemmes \ref{algebraization-lemma-kill-colimit-weak-general} et
\ref{algebraization-lemma-kill-colimit-support-general} sont satisfaites
pour $A, I, \mathfrak a, M$,
\item[(B)] si $A$ est complet pour la topologie $I$-adique, alors
toutes les hypothèses, sauf éventuellement (5), du
lemme \ref{algebraization-lemma-algebraize-local-cohomology-general}
sont satisfaites pour $A, I, \mathfrak a, M$,
\item[(C)] si $A$ est local d'idéal maximal $\mathfrak m = \mathfrak a$,
alors toutes les hypothèses des
lemmes \ref{algebraization-lemma-kill-colimit-weak} et \ref{algebraization-lemma-kill-colimit-support}
sont satisfaites pour $A, \mathfrak m, I, M$,
\item[(D)] si $A$ est local d'idéal maximal $\mathfrak m = \mathfrak a$
et complet pour la topologie $I$-adique, alors toutes les hypothèses du
lemme \ref{algebraization-lemma-algebraize-local-cohomology}
sont satisfaites pour $A, \mathfrak m, I, M$,
\end{enumerate}
\end{lemma}

\begin{proof}
Démonstration de (E). Il faut montrer que les hypothèses (1), (3), (4), (6)
de la situation \ref{algebraization-situation-bootstrap} sont satisfaites pour
$A, I, T, M$. Remplacer $T$ par $T'$
affaiblit l'hypothèse (6) et renforce l'hypothèse (4). Cependant, si l'on a
$\mathfrak p \subset \mathfrak r \subset \mathfrak q$ avec
$\mathfrak p \not \in V(I)$, $\mathfrak r \in V(I) \setminus T'$,
$\mathfrak q \in T'$ comme dans l'hypothèse (4) pour $A, I, T', M$, alors
soit on peut choisir $\mathfrak r \in V(I) \setminus T$ et
la condition (4) pour $A, I, T, M$ s'applique, soit on ne peut
trouver un tel $\mathfrak r$, auquel cas on obtient
$\text{depth}(M_\mathfrak p) > s$ par le lemme \ref{algebraization-lemma-helper-bootstrap}.
Cela montre, comme voulu, que (4) est satisfaite pour $A, I, T', M$.

\medskip\noindent
Démonstration de (F). C'est immédiat et nous omettons les détails.

\medskip\noindent
Démonstration de (G). Il faut montrer que les hypothèses (1), (3), (4), (6)
de la situation \ref{algebraization-situation-bootstrap} sont satisfaites pour les complétions $I$-adiques
$A', I', T', M'$. Gardons à l'esprit que
$\Spec(A') \to \Spec(A)$ induit un isomorphisme $V(I') \to V(I)$.

\medskip\noindent
Hypothèse (1) : L'anneau $A'$ possède un complexe dualisant, voir
Complexes dualisants, lemme \ref{dualizing-lemma-ubiquity-dualizing}.

\medskip\noindent
Hypothèse (3) : Comme $I' = IA'$, cela résulte de Cohomologie locale,
lemme \ref{local-cohomology-lemma-cd-change-rings}.

\medskip\noindent
Hypothèse (4) : Si l'on a des idéaux premiers
$\mathfrak p' \subset \mathfrak r' \subset \mathfrak q'$ de $A'$
tels que $\mathfrak p' \not \in V(I')$,
$\mathfrak r' \in V(I') \setminus T'$,
$\mathfrak q' \in T'$, alors leurs images
$\mathfrak p \subset \mathfrak r \subset \mathfrak q$ dans
le spectre de $A$
satisfont
$\mathfrak p \not \in V(I)$, $\mathfrak r \in V(I) \setminus T$,
$\mathfrak q \in T$.
On a alors
$$
\text{depth}_{A_\mathfrak p}(M_\mathfrak p) \geq s
\quad\text{ou}\quad
\text{depth}_{A_\mathfrak p}(M_\mathfrak p) +
\dim((A/\mathfrak p)_\mathfrak q) > d + s
$$
d'après l'hypothèse (4) pour $A, I, T, M$. On a
$\text{depth}(M'_{\mathfrak p'}) \geq \text{depth}(M_\mathfrak p)$ et
$\text{depth}(M'_{\mathfrak p'}) +
\dim((A'/\mathfrak p')_{\mathfrak q'}) =
\text{depth}(M_\mathfrak p) +
\dim((A/\mathfrak p)_\mathfrak q)$
d'après Cohomologie locale, lemme \ref{local-cohomology-lemma-change-completion}.
Ainsi, l'hypothèse (4) est satisfaite pour $A', I', T', M'$.

\medskip\noindent
Hypothèse (6) : Soit $\mathfrak q' \in T'$ au-dessus de l'idéal
premier $\mathfrak q \in T$. Alors $A'_{\mathfrak q'}$
et $A_\mathfrak q$ ont des complétions $I$-adiques isomorphes,
de même que $M_\mathfrak q$ et $M'_{\mathfrak q'}$.
Ainsi, l'hypothèse (6) pour $A', I', T', M'$ équivaut
à l'hypothèse (6) pour $A, I, T, M$.

\medskip\noindent
Démonstration de (A). Il faut vérifier les conditions (1), (2), (3), (4) et (6)
des lemmes \ref{algebraization-lemma-kill-colimit-weak-general} et
\ref{algebraization-lemma-kill-colimit-support-general} pour
$(A, I, \mathfrak a, M)$. Attention : l'ensemble $T$ de l'énoncé de
ces lemmes n'est pas le même que l'ensemble $T$ ci-dessus.

\medskip\noindent
Condition (1) : Cela vaut parce que nous avons supposé que $A$ possède un complexe dualisant dans la
situation \ref{algebraization-situation-bootstrap}.

\medskip\noindent
Condition (2) : Elle est vide.

\medskip\noindent
Condition (3) : Soit $\mathfrak p \subset A$ tel que
$V(\mathfrak p) \cap V(I) \subset V(\mathfrak a)$.
Comme $I$ est contenu dans le radical de Jacobson de $A$, on voit
que $V(\mathfrak p) \cap V(I) \not = \emptyset$.
Soit $\mathfrak q \in V(\mathfrak p) \cap V(I)$ un point générique.
Puisque $\text{cd}(A_\mathfrak q, I_\mathfrak q) \leq d$
(Cohomologie locale, lemme \ref{local-cohomology-lemma-cd-local}) et puisque
$V(\mathfrak p A_\mathfrak q) \cap V(I_\mathfrak q) =
\{\mathfrak q A_\mathfrak q\}$, on obtient
$\dim((A/\mathfrak p)_\mathfrak q) \leq d$ par Cohomologie locale,
lemme \ref{local-cohomology-lemma-cd-bound-dim-local}, ce qui prouve (3).

\medskip\noindent
Condition (4) : Supposons $\mathfrak p \not \in V(I)$ et
$\mathfrak q \in V(\mathfrak p) \cap V(\mathfrak a)$.
Il suffit de montrer
$$
\text{depth}_{A_\mathfrak p}(M_\mathfrak p) \geq s
\quad\text{ou}\quad
\text{depth}_{A_\mathfrak p}(M_\mathfrak p) +
\dim((A/\mathfrak p)_\mathfrak q) > d + s
$$
S'il existe un idéal premier $\mathfrak p \subset \mathfrak r \subset \mathfrak q$
tel que $\mathfrak r \in V(I) \setminus T$, cela résulte alors
immédiatement de l'hypothèse (4) de la situation \ref{algebraization-situation-bootstrap}.
Sinon, $\text{depth}(M_\mathfrak p) > s$ d'après le
lemme \ref{algebraization-lemma-helper-bootstrap}.

\medskip\noindent
Condition (6) : Soit $\mathfrak p \not \in V(I)$ tel que
$V(\mathfrak p) \cap V(I) \subset V(\mathfrak a)$.
Comme $I$ est contenu dans le radical de Jacobson de $A$, on voit
que $V(\mathfrak p) \cap V(I) \not = \emptyset$.
Choisissons $\mathfrak q \in V(\mathfrak p) \cap V(I) \subset V(\mathfrak a)$.
Il est clair qu'il n'existe pas d'idéal premier
$\mathfrak p \subset \mathfrak r \subset \mathfrak q$
tel que $\mathfrak r \in V(I) \setminus T$.
D'après le lemme \ref{algebraization-lemma-helper-bootstrap}, on a
$\text{depth}(M_\mathfrak p) > s$, ce qui prouve (6).

\medskip\noindent
Démonstration de (B). Il faut vérifier les conditions (1), (2), (3), (4) du
lemme \ref{algebraization-lemma-algebraize-local-cohomology-general}. Attention :
l'ensemble $T$ dans l'énoncé de
ce lemme n'est pas le même que l'ensemble $T$ ci-dessus.

\medskip\noindent
Condition (1) : Cela vaut parce que $A$ est complet et possède un complexe dualisant.

\medskip\noindent
Condition (2) : Elle est vide.

\medskip\noindent
Condition (3) : C'est l'hypothèse (3) de la
situation \ref{algebraization-situation-bootstrap}.

\medskip\noindent
Condition (4) : C'est l'hypothèse (4) du
lemme \ref{algebraization-lemma-kill-colimit-weak-general}, que nous avons démontrée dans (A).

\medskip\noindent
Démonstration de (C). Cela vaut parce que les hypothèses des
lemmes \ref{algebraization-lemma-kill-colimit-weak} et \ref{algebraization-lemma-kill-colimit-support}
sont les mêmes que celles des
lemmes \ref{algebraization-lemma-kill-colimit-weak-general} et
\ref{algebraization-lemma-kill-colimit-support-general} dans le cas local,
et nous avons montré dans (A) qu'elles sont satisfaites.

\medskip\noindent
Démonstration de (D). Cela vaut parce que les hypothèses du
lemme \ref{algebraization-lemma-algebraize-local-cohomology}
sont les hypothèses (1), (2), (3), (4) du
lemme \ref{algebraization-lemma-algebraize-local-cohomology-general},
et nous avons montré dans (B) qu'elles sont satisfaites.
\end{proof}

\begin{lemma}
\label{algebraization-lemma-algebraize-local-cohomology-bis}
Dans la situation \ref{algebraization-situation-bootstrap}, supposons que $A$ est local d'idéal
maximal $\mathfrak m$ et que $T = \{\mathfrak m\}$. Alors
$H^i_\mathfrak m(M) \to \lim H^i_\mathfrak m(M/I^nM)$
est un isomorphisme pour $i \leq s$, et ces modules sont
annulés par une puissance de $I$.
\end{lemma}

\begin{proof}
Soient $A', I', \mathfrak m', M'$ les complétions $I$-adiques usuelles
de $A, I, \mathfrak m, M$. Rappelons que l'on a
$H^i_\mathfrak m(M) \otimes_A A' = H^i_{\mathfrak m'}(M')$
par platitude de $A \to A'$ et par Complexes dualisants, lemme
\ref{dualizing-lemma-torsion-change-rings}.
Comme $H^i_\mathfrak m(M)$ est de torsion annulée par une puissance de $\mathfrak m$, on a
$H^i_\mathfrak m(M) = H^i_\mathfrak m(M) \otimes_A A'$ ; voir
Compléments d'algèbre, lemme \ref{more-algebra-lemma-neighbourhood-equivalence}.
On en conclut que $H^i_\mathfrak m(M) = H^i_{\mathfrak m'}(M')$.
Les mêmes arguments montrent exactement que
$H^i_\mathfrak m(M/I^nM) =  H^i_{\mathfrak m'}(M'/(I')^nM')$
pour tous $n$ et $i$.

\medskip\noindent
Les lemmes \ref{algebraization-lemma-algebraize-local-cohomology},
\ref{algebraization-lemma-kill-colimit-weak} et
\ref{algebraization-lemma-kill-colimit-support}
s'appliquent à $A', \mathfrak m', I', M'$ d'après les parties (C) et (D) du
lemme \ref{algebraization-lemma-bootstrap-inherited}.
On obtient donc un isomorphisme
$$
H^i_{\mathfrak m'}(M') \longrightarrow H^i(R\Gamma_{\mathfrak m'}(M')^\wedge)
$$
pour $i \leq s$, où ${}^\wedge$ désigne la complétion $I'$-adique dérivée, et ces
modules sont annulés par une puissance de $I'$.
D'après le lemme \ref{algebraization-lemma-local-cohomology-derived-completion},
on obtient des isomorphismes
$$
H^i_{\mathfrak m'}(M') \longrightarrow
\lim H^i_{\mathfrak m'}(M'/(I')^nM'))
$$
pour $i \leq s$. En les combinant à la comparaison déjà établie
avec la cohomologie locale sur $A$, on conclut que le lemme est vrai.
\end{proof}

\begin{lemma}
\label{algebraization-lemma-bootstrap-bis-bis}
Soient $I \subset \mathfrak a$ des idéaux d'un anneau noethérien $A$.
Soit $M$ un $A$-module de type fini. Soient $s$ et $d$ des entiers.
Supposons que
\begin{enumerate}
\item[(a)] $A$ possède un complexe dualisant,
\item[(b)] $\text{cd}(A, I) \leq d$,
\item[(c)] si $\mathfrak p \not \in V(I)$ et
$\mathfrak q \in V(\mathfrak p) \cap V(\mathfrak a)$, alors
$\text{depth}_{A_\mathfrak p}(M_\mathfrak p) > s$ ou
$\text{depth}_{A_\mathfrak p}(M_\mathfrak p) +
\dim((A/\mathfrak p)_\mathfrak q) > d + s$.
\end{enumerate}
Alors $A, I, V(\mathfrak a), M, s, d$ sont comme dans la
situation \ref{algebraization-situation-bootstrap}.
\end{lemma}

\begin{proof}
Il faut montrer que les hypothèses (1), (3), (4) et (6) de la
situation \ref{algebraization-situation-bootstrap} sont satisfaites.
Il est clair que (a) $\Rightarrow$ (1),
(b) $\Rightarrow$ (3) et (c) $\Rightarrow$ (4).
Pour achever la démonstration, nous montrons dans le paragraphe suivant que (6) est satisfaite.

\medskip\noindent
Soit $\mathfrak q \in V(\mathfrak a)$.
Notons $A', I', \mathfrak m', M'$
les complétions $I$-adiques de
$A_\mathfrak q, I_\mathfrak q, \mathfrak qA_\mathfrak q, M_\mathfrak q$.
Soit $\mathfrak p' \subset A'$ un idéal premier non maximal tel que
$V(\mathfrak p') \cap V(I') = \{\mathfrak m'\}$.
Remarquons que cela entraîne $\dim(A'/\mathfrak p') \leq d$
d'après Cohomologie locale, lemme \ref{local-cohomology-lemma-cd-bound-dim-local}.
Notons $\mathfrak p \subset A$ l'image de $\mathfrak p'$.
On a
$\text{depth}(M'_{\mathfrak p'}) \geq \text{depth}(M_\mathfrak p)$ et
$\text{depth}(M'_{\mathfrak p'}) +
\dim(A'/\mathfrak p') =
\text{depth}(M_\mathfrak p) +
\dim((A/\mathfrak p)_\mathfrak q)$
d'après Cohomologie locale, lemme \ref{local-cohomology-lemma-change-completion}.
D'après l'hypothèse (c), ou bien on a
$\text{depth}(M'_{\mathfrak p'}) \geq \text{depth}(M_\mathfrak p) > s$
et c'est terminé, ou bien on a
$\text{depth}(M'_{\mathfrak p'}) +
\dim(A'/\mathfrak p') > s + d$, ce qui entraîne
$\text{depth}(M'_{\mathfrak p'}) > s$ par l'inégalité déjà établie
$\dim(A'/\mathfrak p') \leq d$. Dans les deux cas, on
obtient le résultat voulu.
\end{proof}

\begin{lemma}
\label{algebraization-lemma-bootstrap}
Dans la situation \ref{algebraization-situation-bootstrap}, les systèmes projectifs
$\{H^i_T(I^nM)\}_{n \geq 0}$ sont essentiellement nuls pour $i \leq s$.
De plus, il existe un entier $m_0$ tel que, pour tout
$m \geq m_0$, il existe un entier $m'(m) \geq m$ tel que, pour
$k \geq m'(m)$, l'image de
$H^{s + 1}_T(I^kM) \to H^{s + 1}_T(I^mM)$
s'envoie injectivement dans $H^{s + 1}_T(I^{m_0}M)$.
\end{lemma}

\begin{proof}
Fixons $m$. Soit $\mathfrak q \in T$.
D'après les lemmes \ref{algebraization-lemma-bootstrap-inherited} et
\ref{algebraization-lemma-algebraize-local-cohomology-bis}
on voit que
$$
H^i_\mathfrak q(M_\mathfrak q)
\longrightarrow
\lim H^i_\mathfrak q(M_\mathfrak q/I^nM_\mathfrak q)
$$
est un isomorphisme pour $i \leq s$. Les systèmes projectifs
$\{H^i_\mathfrak q(I^nM_\mathfrak q)\}_{n \geq 0}$ et
$\{H^i_\mathfrak q(M/I^nM)\}_{n \geq 0}$
satisfont la condition de Mittag-Leffler pour tout $i$ ; voir le
lemme \ref{algebraization-lemma-ML-local}. En considérant donc le système projectif des
suites exactes longues
$$
0 \to H^0_\mathfrak q(I^nM_\mathfrak q) \to
H^0_\mathfrak q(M_\mathfrak q) \to
H^0_\mathfrak q(M_\mathfrak q/I^nM_\mathfrak q) \to
H^1_\mathfrak q(I^nM_\mathfrak q) \to
H^1_\mathfrak q(M_\mathfrak q) \to \ldots
$$
on conclut (certains détails sont omis) qu'il existe un entier
$m'(m, \mathfrak q) \geq m$ tel que, pour tout $k \geq m'(m, \mathfrak q)$,
l'application
$H^i_\mathfrak q(I^kM_\mathfrak q) \to H^i_\mathfrak q(I^mM_\mathfrak q)$
est nulle pour $i \leq s$, et l'image de
$H^{s + 1}_\mathfrak q(I^kM_\mathfrak q) \to
H^{s + 1}_\mathfrak q(I^mM_\mathfrak q)$
ne dépend pas de $k \geq m'(m, \mathfrak q)$ et
s'envoie injectivement dans $H^{s + 1}_\mathfrak q(M_\mathfrak q)$.

\medskip\noindent
Supposons que l'on puisse montrer que $m'(m, \mathfrak q)$ peut être choisi
indépendamment de $\mathfrak q \in T$.
Alors le lemme résulte immédiatement de
Cohomologie locale, lemmes \ref{local-cohomology-lemma-zero} et
\ref{local-cohomology-lemma-essential-image}.

\medskip\noindent
Soit $\omega_A^\bullet$ un complexe dualisant. Soit
$\delta : \Spec(A) \to \mathbf{Z}$ la fonction de dimension
correspondante. Rappelons que $\delta$ ne prend qu'un
nombre fini de valeurs ; voir
Complexes dualisants, lemme \ref{dualizing-lemma-universally-catenary}.
Assertion : pour tout $d \in \mathbf{Z}$, l'entier
$m'(m, \mathfrak q)$ peut être choisi indépendamment
de $\mathfrak q \in T$ tel que $\delta(\mathfrak q) = d$.
Cette assertion entraîne clairement le lemme d'après ce qui précède.

\medskip\noindent
Choisissons $\mathfrak q \in T$ tel que $\delta(\mathfrak q) = d$.
Considérons les modules d'Ext
$$
E(n, j) = \text{Ext}^j_A(I^nM, \omega_A^\bullet)
$$
La propriété essentielle que nous utiliserons est que ce sont des $A$-modules de type fini.
Rappelons que $(\omega_A^\bullet)_\mathfrak q[-d]$ est un complexe
dualisant normalisé pour $A_\mathfrak q$, par définition de la
fonction de dimension associée à un complexe dualisant ; voir
Complexes dualisants, section \ref{dualizing-section-dimension-function}.
Le théorème de dualité locale (Complexes dualisants, lemme
\ref{dualizing-lemma-special-case-local-duality}) nous dit que
le complété $\mathfrak qA_\mathfrak q$-adique de
$E(n, -d - i)_\mathfrak q$ est le dual de Matlis de
$H^i_\mathfrak q(I^nM_\mathfrak q)$. Le choix de
$m'(m, \mathfrak q)$ pour $i \leq s$ dans le premier paragraphe nous dit donc que,
pour $k \geq m'(m, \mathfrak q)$ et $j \geq -d - s$, l'application
$$
E(m, j)_\mathfrak q \to E(k, j)_\mathfrak q
$$
est nulle. Puisque ces modules sont de type fini et ne sont non nuls que
pour un nombre fini de valeurs possibles de $j$ (petit détail omis),
on peut trouver un voisinage ouvert $W \subset \Spec(A)$ de $\mathfrak q$
tel que
$$
E(m, j)_{\mathfrak q'} \to E(m'(m, \mathfrak q), j)_{\mathfrak q'}
$$
soit nulle pour $j \geq -d - s$ et tout $\mathfrak q' \in W$.
Alors, bien entendu, les applications $E(m, j)_{\mathfrak q'} \to E(k, j)_{\mathfrak q'}$
pour $k \geq m'(m, \mathfrak q)$ sont elles aussi nulles.

\medskip\noindent
Pour $i = s + 1$, correspondant à $j = - d - s - 1$, la dualité locale
et les résultats du premier paragraphe montrent que
$$
K_{k, \mathfrak q} =
\Ker(E(m, -d - s - 1)_\mathfrak q \to E(k, -d - s - 1)_\mathfrak q)
$$
est indépendant de $k \geq m'(m, \mathfrak q)$ et que
$$
E(0, -d - s - 1)_\mathfrak q \to
E(m, -d - s - 1)_\mathfrak q/K_{m'(m, \mathfrak q), \mathfrak q}
$$
est surjective. Pour $k \geq m'(m, \mathfrak q)$, posons
$$
K_k = \Ker(E(m, -d - s - 1) \to E(k, -d - s - 1))
$$
Puisque $K_k$ est une suite croissante de sous-modules du module de type fini
$E(m, -d - s - 1)$, on voit que, quitte à augmenter
un peu $m'(m, \mathfrak q)$, on peut supposer
$K_{m'(m, \mathfrak q)} = K_k$ pour $k \geq m'(m, \mathfrak q)$.
Après avoir encore rétréci $W$ si nécessaire, on peut aussi supposer que
$$
E(0, -d - s - 1)_{\mathfrak q'} \to
E(m, -d - s - 1)_{\mathfrak q'}/K_{m'(m, \mathfrak q), \mathfrak q'}
$$
est surjective pour tout $\mathfrak q' \in W$ (comme précédemment, on utilise que
ces modules sont de type fini
et que l'application est surjective après localisation en $\mathfrak q$).

\medskip\noindent
Toute partie, en particulier
$T_d = \{\mathfrak q \in T \text{ tel que }\delta(\mathfrak q) = d\}$,
de l'espace topologique noethérien $\Spec(A)$,
munie de la topologie induite, est noethérienne et donc quasi-compacte.
Nous avons vu ci-dessus que, pour tout $\mathfrak q \in T_d$,
il existe un voisinage ouvert $W$ sur lequel
$m'(m, \mathfrak q)$ convient à tout $\mathfrak q' \in T_d \cap W$.
On en conclut qu'il existe un entier $m'(m, d)$ tel que, pour tout
$\mathfrak q \in T_d$, on ait
$$
E(m, j)_\mathfrak q \to E(m'(m, d), j)_\mathfrak q
$$
qui soit nulle pour $j \geq -d - s$ et, en posant
$K_{m'(m, d)} = \Ker(E(m, -d - s - 1) \to E(m'(m, d), -d - s - 1))$,
on ait
$$
K_{m'(m, d), \mathfrak q} =
\Ker(E(m, -d - s - 1)_{\mathfrak q} \to E(k, -d - s - 1)_{\mathfrak q})
$$
pour tout $k \geq m'(m, d)$, et que l'application
$$
E(0, -d - s - 1)_\mathfrak q \to
E(m, -d - s - 1)_\mathfrak q/K_{m'(m, d), \mathfrak q}
$$
soit surjective. En utilisant de nouveau le théorème de dualité locale (dans le sens
opposé), on en conclut que l'assertion est vraie. Cela achève la démonstration.
\end{proof}

\begin{lemma}
\label{algebraization-lemma-final-bootstrap}
Dans la situation \ref{algebraization-situation-bootstrap}, il existe un entier $m_0 \geq 0$
tel que
\begin{enumerate}
\item $H^i_T(M) \to H^i_T(M/I^nM)$ soit injective pour tout $i \leq s$
et tout $n \geq m_0$,
\item $\{H^i_T(M/I^nM)\}_{n \geq 0}$
vérifie la condition de Mittag-Leffler pour $i < s$,
\item $\{H^i_T(I^{m_0}M/I^nM)\}_{n \geq m_0}$
vérifie la condition de Mittag-Leffler pour $i \leq s$,
\item $H^i_T(M) \to \lim H^i_T(M/I^nM)$
soit un isomorphisme pour $i < s$,
\item $H^s_T(I^{m_0}M) \to \lim H^s_T(I^{m_0}M/I^nM)$
soit un isomorphisme pour $i \leq s$,
\item $H^s_T(M) \to \lim H^s_T(M/I^nM)$ soit
injective, de conoyau annulé par $I^{m_0}$, et
\item $R^1\lim H^s_T(M/I^nM)$ soit annulé par $I^{m_0}$.
\end{enumerate}
\end{lemma}

\begin{proof}
Pour démontrer (1), soit $0 \leq i \leq s$ et choisissons des entiers $m > n > 0$ tels que
$H^i_T(I^mM) \to H^i_T(I^nM)$ soit nulle ; c'est possible d'après le
lemme \ref{algebraization-lemma-bootstrap}. Le diagramme commutatif à lignes exactes
$$
\xymatrix{
H^i_T(I^nM) \ar[r] &
H^i_T(M) \ar[r] &
H^i_T(M/I^nM) \\
H^i_T(I^mM) \ar[r] \ar[u]^0 &
H^i_T(M) \ar[r] \ar[u]^1 &
H^i_T(M/I^mM) \ar[u] \\
}
$$
montre que l'application $H^i_T(M) \to H^i_T(M/I^mM)$ est injective. Ainsi, (1) vaut
pour tout $m_0$ suffisamment grand.

\medskip\noindent
Considérons les suites exactes longues
$$
0 \to H^0_T(I^nM) \to H^0_T(M) \to
H^0_T(M/I^nM) \to H^1_T(I^nM) \to
H^1_T(M) \to \ldots
$$
Les assertions (2) et (4) en résultent avec le lemme \ref{algebraization-lemma-bootstrap}.

\medskip\noindent
Soient $m_0$ et $m'(-)$ comme dans le lemme \ref{algebraization-lemma-bootstrap}.
Pour $m \geq m_0$, considérons la suite exacte longue
$$
H^s_T(I^mM) \to H^s_T(I^{m_0}M) \to
H^s_T(I^{m_0}M/I^mM) \to H^{s + 1}_T(I^mM) \to
H^1_T(I^{m_0}M)
$$
Alors, pour $k \geq m'(m)$, l'image de
$H^{s + 1}_T(I^kM) \to H^{s + 1}_T(I^mM)$
se plonge dans $H^{s + 1}_T(I^{m_0}M)$.
Par conséquent, l'image de
$H^s_T(I^{m_0}M/I^kM) \to H^s_T(I^{m_0}M/I^mM)$
a une image nulle dans $H^{s + 1}_T(I^mM)$ pour tout $k \geq m'(m)$.
On en conclut que (3) et (5) sont vraies.

\medskip\noindent
Considérons les suites exactes courtes
$0 \to I^{m_0}M \to M \to M/I^{m_0} M \to 0$ et
$0 \to I^{m_0}M/I^nM \to M/I^nM \to M/I^{m_0} M \to 0$.
On obtient un diagramme
$$
\xymatrix{
H^{s - 1}_T(M/I^{m_0}M) \ar[r] &
\lim H^s_T(I^{m_0}M/I^nM) \ar[r] &
\lim H^s_T(M/I^nM) \ar[r] &
H^s_T(M/I^{m_0}M) \\
H^{s - 1}_T(M/I^{m_0}M) \ar[r] \ar@{=}[u] &
H^s_T(I^{m_0}M) \ar[r] \ar[u]_{\cong} &
H^s_T(M) \ar[r] \ar[u] &
H^s_T(M/I^{m_0}M) \ar@{=}[u]
}
$$
dont la ligne inférieure est exacte. La ligne supérieure est elle aussi exacte
(aux deux places médianes) d'après
Homologie, lemme \ref{homology-lemma-apply-Mittag-Leffler}.
L'assertion (6) en résulte.

\medskip\noindent
Posons $B_n = H^s_T(M/I^nM)$. Soit $A_n \subset B_n$
l'image de $H^s_T(I^{m_0}M/I^nM) \to H^s_T(M/I^nM)$.
Alors $(A_n)$ vérifie la condition de Mittag-Leffler d'après (3) et
Homologie, lemme \ref{homology-lemma-Mittag-Leffler}.
De plus, $C_n = B_n/A_n$ est annulé par $I^{m_0}$. Ainsi,
$R^1\lim B_n \cong R^1\lim C_n$ est annulé par $I^{m_0}$, ce qui donne (7).
\end{proof}

\begin{theorem}
\label{algebraization-theorem-final-bootstrap}
Dans la situation \ref{algebraization-situation-bootstrap}, pour $i \leq s$, le système projectif
$\{H^i_T(M/I^nM)\}_{n \geq 0}$ est essentiellement constant de valeur $H^i_T(M)$.
En particulier, $\{H^i_T(M/I^nM)\}_{n \geq 0}$ est de Mittag-Leffler,
$H^i_T(M)$ est annulé par une puissance de $I$, et
$H^i_T(M) = \lim H^i_T(M/I^nM)$.
\end{theorem}

\begin{proof}
Soit $0 \leq i \leq s$.
Si l'on montre que $\{H^i_T(M/I^nM)\}_{n \geq 0}$ est de Mittag-Leffler et que
$H^i_T(M) = \lim H^i_T(M/I^nM)$, alors l'injectivité de
$H^i_T(M) \to H^i_T(M/I^nM)$ pour $n \gg 0$, établie dans le
lemme \ref{algebraization-lemma-final-bootstrap}, entraînera que
$\{H^i_T(M/I^nM)\}_{n \geq 0}$ est essentiellement constant de valeur $H^i_T(M)$.
Nous omettons un petit détail.

\medskip\noindent
Nous savons déjà que $\{H^i_T(M/I^nM)\}_{n \geq 0}$ est de Mittag-Leffler et que
$H^i_T(M) = \lim H^i_T(M/I^nM)$ pour $i < s$ d'après le
lemme \ref{algebraization-lemma-final-bootstrap}. Nous savons aussi que ces assertions valent
pour $i = s$ et le module $I^{m_0}M$, pour un certain $m_0 > 0$.
Pour achever la démonstration, nous allons montrer qu'en fait ces
assertions pour $i = s$ valent pour $M$.

\medskip\noindent
Posons $M' = H^0_I(M)$ et $M'' = M/M'$, de sorte que l'on ait une suite
exacte courte
$$
0 \to M' \to M \to M'' \to 0
$$
et $M''$ vérifie $H^0_I(M') = 0$ d'après
Complexes dualisants, lemme \ref{dualizing-lemma-divide-by-torsion}.
D'après Artin-Rees (Algèbre, lemme \ref{algebra-lemma-Artin-Rees}),
on obtient des suites exactes courtes
$$
0 \to M' \to M/I^n M \to M''/I^n M'' \to 0
$$
pour $n$ assez grand. Considérons les suites exactes longues
$$
H^s_T(M') \to
H^s_T(M/I^nM) \to
H^s_T(M''/I^nM'') \to
H^{s + 1}_T(M')
$$
Il est alors aisé de voir que, si le système projectif
$\{H^s_T(M''/I^nM'')\}_{n \geq 0}$ est de Mittag-Leffler,
le système projectif
$\{H^s_T(M/I^nM)\}_{n \geq 0}$ l'est aussi.
(Remarquons que la condition ML pour un système projectif de groupes $G_n$
ne dépend que des valeurs du système projectif pour $n$ assez grand.)
De plus, la suite
$$
H^s_T(M') \to
\lim H^s_T(M/I^nM) \to
\lim H^s_T(M''/I^nM'') \to
H^{s + 1}_T(M')
$$
est exacte parce que la condition ML est satisfaite aux endroits voulus ; voir
Homologie, lemme \ref{homology-lemma-apply-Mittag-Leffler}.
Ainsi, si $H^s_T(M'') \to \lim H^s_T(M''/I^nM'')$
est un isomorphisme, alors
$H^s_T(M) \to \lim H^s_T(M/I^nM)$
est aussi un isomorphisme d'après le lemme des cinq
(Homologie, lemme \ref{homology-lemma-five-lemma}).
On est ainsi ramené au cas traité dans le paragraphe suivant.

\medskip\noindent
Supposons que $H^0_I(M) = 0$. Choisissons des générateurs
$f_1, \ldots, f_r$ de $I^{m_0}$, où $m_0$ est l'entier
obtenu pour $M$ dans le lemme \ref{algebraization-lemma-final-bootstrap}.
Considérons alors la suite exacte
$$
0 \to M \xrightarrow{f_1, \ldots, f_r}
(I^{m_0}M)^{\oplus r} \to Q \to 0
$$
qui définit $Q$. Faisons quelques remarques : la première application est injective
précisément parce que $H^0_I(M) = 0$. Le conoyau $Q$ de cette injection
est un $A$-module de type fini tel que, pour tout $1 \leq j \leq r$,
on ait $Q_{f_j} \cong (M_{f_j})^{\oplus r - 1}$.
En particulier, pour un idéal premier $\mathfrak p \subset A$
tel que $\mathfrak p \not \in V(I)$, on a
$Q_\mathfrak p \cong (M_\mathfrak p)^{\oplus r - 1}$.
De même, étant donné $\mathfrak q \in T$ et
$\mathfrak p' \subset A' = (A_\mathfrak q)^\wedge$
non contenu dans $V(IA')$, on a
$Q'_{\mathfrak p'} \cong (M'_{\mathfrak p'})^{\oplus r - 1}$
où $Q' = (Q_\mathfrak q)^\wedge$ et $M' = (M_\mathfrak q)^\wedge$.
Ainsi, les conditions de la situation \ref{algebraization-situation-bootstrap}
sont satisfaites pour $A, I, T, Q$. (Remarquons que $Q$ peut avoir
un $H^0_I(Q)$ non nul, mais cela n'aura pas d'importance.)

\medskip\noindent
Pour tout $n \geq 0$, posons $F^nM = M \cap I^n(I^{m_0}M)^{\oplus r}$,
de sorte que l'on obtienne des suites exactes courtes
$$
0 \to F^nM \to I^n(I^{m_0}M)^{\oplus r} \to I^nQ \to 0
$$
D'après Artin-Rees (Algèbre, lemme \ref{algebra-lemma-Artin-Rees}),
il existe un $c \geq 0$ tel que
$I^n M \subset F^nM \subset I^{n - c}M$ pour tout $n \geq c$.
Soient $m_0$ l'entier et $m'(m)$
la fonction définie pour $m \geq m_0$,
obtenus dans le lemme \ref{algebraization-lemma-bootstrap}
appliqué à $M$. Remarquons que l'entier $m_0$
est le même que notre entier $m_0$ choisi ci-dessus (il n'est pas nécessaire de
le vérifier : on peut simplement prendre le maximum des deux entiers si
l'on préfère). Enfin, d'après le lemme \ref{algebraization-lemma-bootstrap}
appliqué à $Q$, pour tout entier $m$, il existe un entier
$m''(m) \geq m$ tel que $H^s_T(I^kQ) \to H^s_T(I^mQ)$
soit nulle pour tout $k \geq m''(m)$.

\medskip\noindent
Fixons $m \geq m_0$. Choisissons $k \geq m'(m''(m + c))$.
Choisissons $\xi \in H^{s + 1}_T(I^kM)$
qui s'envoie sur zéro dans $H^{s + 1}_T(M)$.
Nous voulons montrer que $\xi$ s'envoie sur zéro dans $H^{s + 1}_T(I^mM)$.
En effet, cela montrera que $\{H^s_T(M/I^nM)\}_{n \geq 0}$
est de Mittag-Leffler, exactement comme dans la démonstration du lemme \ref{algebraization-lemma-final-bootstrap}.
Voici un diagramme qui permet de visualiser l'argument :
$$
\xymatrix{
&
H^{s + 1}_T(I^kM) \ar[r] \ar[d] &
H^{s + 1}_T(I^k(I^{m_0}M)^{\oplus r}) \ar[d] &
\\
H^s_T(I^{m''(m + c)}Q) \ar[r]_-\delta \ar[d] &
H^{s + 1}_T(F^{m''(m + c)}M) \ar[r] \ar[d] &
H^{s + 1}_T(I^{m''(m + c)}(I^{m_0}M)^{\oplus r}) \\
H^s_T(I^{m + c}Q) \ar[r] &
H^{s + 1}_T(F^{m + c}M) \ar[d] &
\\
&
H^{s + 1}_T(I^mM)
}
$$
L'image de $\xi$ dans $H^{s + 1}_T(I^k(I^{m_0}M)^{\oplus r})$
s'envoie sur zéro dans $H^{s + 1}_T((I^{m_0}M)^{\oplus r})$
et s'envoie donc sur zéro dans
$H^{s + 1}_T(I^{m''(m + c)}(I^{m_0}M)^{\oplus r})$
par le choix de $m'(-)$.
Ainsi, l'image $\xi' \in H^{s + 1}_T(F^{m''(m + c)}M)$
s'envoie sur zéro dans $H^{s + 1}_T(I^{m''(m + c)}(I^{m_0}M)^{\oplus r})$
et, par conséquent, $\xi' = \delta(\eta)$ pour un certain
$\eta \in H^s_T(I^{m''(m + c)}Q)$.
D'après notre choix de $m''(-)$, on voit que $\eta$ s'envoie sur
zéro dans $H^s_T(I^{m + c}Q)$.
Cela signifie à son tour que $\xi'$ s'envoie sur zéro dans
$H^{s + 1}_T(F^{m + c}M)$.
Comme $F^{m + c}M \subset I^mM$, on conclut.

\medskip\noindent
Enfin, démontrons l'assertion sur les limites. Considérons les suites
exactes courtes
$$
0 \to M/F^nM \to (I^{m_0}M)^{\oplus r}/I^n (I^{m_0}M)^{\oplus r}
\to Q/I^nQ \to 0
$$
On a $\lim H^s_T(M/I^nM) = \lim H^s_T(M/F^nM)$,
car ces systèmes projectifs sont pro-isomorphes. On obtient un diagramme commutatif
$$
\xymatrix{
H^{s - 1}_T(Q) \ar[r] \ar[d] &
\lim H^{s - 1}_T(Q/I^nQ) \ar[d] \\
H^s_T(M) \ar[r] \ar[d] &
\lim H^s_T(M/I^nM) \ar[d] \\
H^s_T((I^{m_0}M)^{\oplus r}) \ar[r] \ar[d] &
\lim H^s_T((I^{m_0}M)^{\oplus r}/I^n(I^{m_0}M)^{\oplus r}) \ar[d] \\
H^s_T(Q) \ar[r] &
\lim H^s_T(Q/I^nQ)
}
$$
La colonne de droite est exacte parce que la condition ML est satisfaite aux endroits voulus ; voir
Homologie, lemme \ref{homology-lemma-apply-Mittag-Leffler}.
La flèche horizontale la plus basse est injective (!) d'après la
partie (6) du lemme \ref{algebraization-lemma-final-bootstrap}.
La flèche horizontale située au-dessus est bijective d'après la
partie (5) du lemme \ref{algebraization-lemma-final-bootstrap}.
Les flèches en degrés cohomologiques $\leq s - 1$ sont des isomorphismes.
On conclut donc que $H^s_T(M) \to \lim H^s_T(M/I^nM)$
est un isomorphisme d'après le lemme des cinq
(Homologie, lemme \ref{homology-lemma-five-lemma}).
Cela achève la démonstration du théorème.
\end{proof}

\begin{lemma}
\label{algebraization-lemma-combine-two}
Soient $I \subset \mathfrak a \subset A$ des idéaux d'un anneau noethérien $A$,
et soit $M$ un $A$-module de type fini. Soient $s$ et $d$ des entiers.
Supposons que
\begin{enumerate}
\item $A, I, V(\mathfrak a), M$ satisfont les hypothèses de la
situation \ref{algebraization-situation-bootstrap} pour $s$ et $d$, et
\item $A, I, \mathfrak a, M$ satisfont les conditions du
lemme \ref{algebraization-lemma-algebraize-local-cohomology-general}
pour $s + 1$ et $d$, avec $J = \mathfrak a$.
\end{enumerate}
Alors il existe un idéal
$J_0 \subset \mathfrak a$ avec $V(J_0) \cap V(I) = V(\mathfrak a)$
tel que, pour tout $J \subset J_0$ avec $V(J) \cap V(I) = V(\mathfrak a)$,
l'application
$$
H^{s + 1}_J(M) \longrightarrow \lim H^{s + 1}_\mathfrak a(M/I^nM)
$$
soit un isomorphisme.
\end{lemma}

\begin{proof}
En effet, l'existence de $J_0$
et l'isomorphisme
$H^{s + 1}_J(M) = H^{s + 1}(R\Gamma_\mathfrak a(M)^\wedge)$
résultent du lemme \ref{algebraization-lemma-algebraize-local-cohomology-general} ;
on a une suite exacte courte
$$
0 \to R^1\lim H^s_\mathfrak a(M/I^nM) \to
H^{s + 1}(R\Gamma_\mathfrak a(M)^\wedge) \to
\lim H^{s + 1}_\mathfrak a(M/I^nM) \to 0
$$
d'après Complexes dualisants, lemme \ref{dualizing-lemma-completion-local},
et le module $R^1\lim H^s_\mathfrak a(M/I^nM)$ est nul parce que
$\{H^s_\mathfrak a(M/I^nM)\}_{n \geq 0}$ satisfait la condition de Mittag-Leffler
d'après le théorème \ref{algebraization-theorem-final-bootstrap}.
\end{proof}








\section{Algébrisation des sections formelles, I}
\label{algebraization-section-algebraization-sections}

\noindent
Dans cette section, nous étudions le problème de l'algébrisation des
sections formelles dans le cas local.
Soit $(A, \mathfrak m)$ un anneau local noethérien.
Soit $I \subset A$ un idéal. Posons
$$
X = \Spec(A) \supset U = \Spec(A) \setminus \{\mathfrak m\}
$$
et notons $Y = V(I)$ le sous-schéma fermé correspondant à $I$.
Soit $\mathcal{F}$ un $\mathcal{O}_U$-module cohérent.
Dans cette section, nous considérons les limites
$$
\lim_n H^i(U, \mathcal{F}/I^n\mathcal{F})
$$
Ceci est étroitement lié à la cohomologie de l'image inverse
de $\mathcal{F}$ sur la complétion formelle de $U$ le long de $Y$ ;
toutefois, comme nous n'avons pas encore introduit les schémas formels,
nous ne pouvons pas employer ici cette terminologie.

\begin{lemma}
\label{algebraization-lemma-compare-with-derived-completion}
Soit $U$ le spectre épointé d'un anneau local noethérien $A$.
Soit $\mathcal{F}$ un $\mathcal{O}_U$-module cohérent.
Soit $I \subset A$ un idéal. Alors
$$
H^i(R\Gamma(U, \mathcal{F})^\wedge) =
\lim H^i(U, \mathcal{F}/I^n\mathcal{F})
$$
pour tout $i$, où $R\Gamma(U, \mathcal{F})^\wedge$ désigne
la complétion $I$-adique dérivée.
\end{lemma}

\begin{proof}
D'après les lemmes \ref{algebraization-lemma-formal-functions-general} et
\ref{algebraization-lemma-derived-completion-pseudo-coherent}, on a
$$
R\Gamma(U, \mathcal{F})^\wedge =
R\Gamma(U, \mathcal{F}^\wedge) =
R\Gamma(U, R\lim \mathcal{F}/I^n\mathcal{F})
$$
On obtient donc des suites exactes courtes
$$
0 \to R^1\lim H^{i - 1}(U, \mathcal{F}/I^n\mathcal{F}) \to
H^i(R\Gamma(U, \mathcal{F})^\wedge) \to
\lim H^i(U, \mathcal{F}/I^n\mathcal{F}) \to 0
$$
d'après Cohomologie, lemme \ref{cohomology-lemma-RGamma-commutes-with-Rlim}.
Les termes $R^1\lim$ s'annulent parce que les systèmes projectifs de groupes
$H^i(U, \mathcal{F}/I^n\mathcal{F})$ satisfont la condition de Mittag-Leffler
d'après le lemme \ref{algebraization-lemma-ML-local}.
\end{proof}

\begin{theorem}
\label{algebraization-theorem-algebraization-formal-sections}
\begin{reference}
La méthode de démonstration suit approximativement celle de la
démonstration de \cite[Theorem 1]{Faltings-algebraisation}
et de \cite[Satz 2]{Faltings-uber}.
Le résultat est presque le même que
\cite[Theorem 1.1]{MRaynaud-paper} (cas du complémentaire affine) et
\cite[Theorem 3.9]{MRaynaud-book} (le complémentaire est réunion de peu d'affines).
\end{reference}
Soit $(A, \mathfrak m)$ un anneau local noethérien qui possède un
complexe dualisant et qui est complet pour la topologie définie par un idéal $I$.
Posons $X = \Spec(A)$, $Y = V(I)$ et $U = X \setminus \{\mathfrak m\}$.
Soit $\mathcal{F}$ un faisceau cohérent sur $U$.
Supposons que
\begin{enumerate}
\item $\text{cd}(A, I) \leq d$, c'est-à-dire que
$H^i(X \setminus Y, \mathcal{G}) = 0$ pour $i \geq d$ et tout
$\mathcal{G}$ quasi-cohérent sur $X$,
\item pour tout $x \in X \setminus Y$ dont l'adhérence $\overline{\{x\}}$
dans $X$ rencontre $U \cap Y$, on ait
$$
\text{depth}_{\mathcal{O}_{X, x}}(\mathcal{F}_x) \geq s
\quad\text{ou}\quad
\text{depth}_{\mathcal{O}_{X, x}}(\mathcal{F}_x)
+ \dim(\overline{\{x\}}) > d + s
$$
\end{enumerate}
Alors il existe un ouvert $V_0 \subset U$ contenant $U \cap Y$
tel que, pour tout ouvert $V \subset V_0$ contenant $U \cap Y$,
l'application
$$
H^i(V, \mathcal{F}) \to \lim H^i(U, \mathcal{F}/I^n\mathcal{F})
$$
soit un isomorphisme pour $i < s$. Si, de plus,
$
\text{depth}_{\mathcal{O}_{X, x}}(\mathcal{F}_x) +
\dim(\overline{\{x\}}) > s
$
pour tout $x \in U \cap Y$, alors ces groupes de cohomologie sont des $A$-modules de type fini.
\end{theorem}

\begin{proof}
Choisissons un $A$-module de type fini $M$ tel que $\mathcal{F}$ soit la
restriction à $U$ du
$\mathcal{O}_X$-module cohérent associé à $M$ ; voir Cohomologie locale,
lemme \ref{local-cohomology-lemma-finiteness-pushforwards-and-H1-local}.
Alors les hypothèses du
lemme \ref{algebraization-lemma-algebraize-local-cohomology}
sont satisfaites.
Choisissons $J_0$ comme dans ce lemme et posons $V_0 = X \setminus V(J_0)$.
Les ouverts $V \subset V_0$ contenant $U \cap Y$
correspondent alors $1$ à $1$ aux idéaux $J \subset J_0$ tels que
$V(J) \cap V(I) = \{\mathfrak m\}$.
De plus, pour un tel choix, on a un triangle distingué
$$
R\Gamma_J(M) \to M \to R\Gamma(V, \mathcal{F}) \to
R\Gamma_J(M)[1]
$$
On a de même un triangle distingué
$$
R\Gamma_\mathfrak m(M)^\wedge \to
M \to
R\Gamma(U, \mathcal{F})^\wedge \to
R\Gamma_\mathfrak m(M)^\wedge[1]
$$
faisant intervenir les complétions $I$-adiques dérivées.
Les groupes de cohomologie de $R\Gamma(U, \mathcal{F})^\wedge$ sont
égaux aux limites figurant dans l'énoncé du théorème d'après le
lemme \ref{algebraization-lemma-compare-with-derived-completion}.
L'application canonique entre ces triangles
et quelques arguments simples montrent que notre
théorème résulte du lemme principal \ref{algebraization-lemma-algebraize-local-cohomology}
(remarquons qu'ici on a $i < s$, tandis que dans le lemme on a
$i \leq s$ ; cela provient du décalage).
Le caractère de type fini des groupes de cohomologie
(sous l'hypothèse supplémentaire) résulte du
lemme \ref{algebraization-lemma-kill-colimit}.
\end{proof}

\begin{lemma}
\label{algebraization-lemma-application-theorem}
Soit $(A, \mathfrak m)$ un anneau local noethérien qui possède un
complexe dualisant et qui est complet pour la topologie définie par un idéal $I$.
Posons $X = \Spec(A)$, $Y = V(I)$ et $U = X \setminus \{\mathfrak m\}$.
Soit $\mathcal{F}$ un faisceau cohérent sur $U$.
Supposons que, pour tout point associé $x \in U$ de $\mathcal{F}$,
on ait $\dim(\overline{\{x\}}) > \text{cd}(A, I) + 1$,
où $\overline{\{x\}}$ est l'adhérence dans $X$.
Alors l'application
$$
\colim H^0(V, \mathcal{F})
\longrightarrow
\lim H^0(U, \mathcal{F}/I^n\mathcal{F})
$$
est un isomorphisme de $A$-modules de type fini,
où la limite inductive porte sur les ouverts $V \subset U$
contenant $U \cap Y$.
\end{lemma}

\begin{proof}
Appliquons le théorème \ref{algebraization-theorem-algebraization-formal-sections} avec $s = 1$
(on obtient également le caractère de type fini).
\end{proof}




\section{Algébrisation des sections formelles, II}
\label{algebraization-section-algebraization-sections-coherent}

\noindent
Il est assez difficile d'énoncer succinctement toutes les conséquences
possibles des résultats des
sections \ref{algebraization-section-algebraization-sections-general} et
\ref{algebraization-section-bootstrap}
pour la cohomologie des faisceaux cohérents sur les schémas quasi-affines
et leur complétion relativement à un idéal.
Cette section donne une liste non exhaustive
d'applications à $H^0$. La section suivante contient des
applications à la cohomologie supérieure.

\begin{lemma}
\label{algebraization-lemma-ses-H0}
Soient $I \subset \mathfrak a$ des idéaux d'un anneau noethérien $A$.
Soit $0 \to \mathcal{F}' \to \mathcal{F} \to \mathcal{F}'' \to 0$
une suite exacte courte de modules cohérents sur
$U = \Spec(A) \setminus V(\mathfrak a)$. Soit $\mathcal{V}$
l'ensemble des sous-schémas ouverts $V \subset U$ contenant
$U \cap V(I)$, ordonné par inclusion inverse.
Considérons le diagramme commutatif
$$
\xymatrix{
\colim_\mathcal{V} H^0(V, \mathcal{F}') \ar[d] \ar[r] &
\colim_\mathcal{V} H^0(V, \mathcal{F}) \ar[d] \ar[r] &
\colim_\mathcal{V} H^0(V, \mathcal{F}'') \ar[d] \\
\lim H^0(U, \mathcal{F}'/I^n\mathcal{F}') \ar[r] &
\lim H^0(U, \mathcal{F}'/I^n\mathcal{F}) \ar[r] &
\lim H^0(U, \mathcal{F}'/I^n\mathcal{F}'')
}
$$
Si les flèches verticales de gauche et de droite sont des isomorphismes, celle du milieu l'est aussi.
Si les flèches verticales du milieu et de gauche sont des isomorphismes, celle de gauche l'est aussi.
\end{lemma}

\begin{proof}
Les suites du diagramme sont exactes au milieu et la première
flèche est injective. La dernière assertion résulte donc d'une simple
chasse au diagramme. Pour le reste de la démonstration, supposons que les flèches verticales
de gauche et de droite soient des isomorphismes. Une chasse au diagramme montre
que la flèche verticale du milieu est injective. Il ne reste
qu'à montrer qu'elle est surjective.

\medskip\noindent
On peut choisir des $A$-modules de type fini $M$ et $M'$ tels que $\mathcal{F}$
et $\mathcal{F}'$ soient les restrictions
de $\widetilde{M}$ et $\widetilde{M}'$ à $U$ ; voir Cohomologie locale, lemme
\ref{local-cohomology-lemma-finiteness-pushforwards-and-H1-local}.
Après avoir remplacé $M'$ par $\mathfrak a^n M'$ pour un certain $n \geq 0$,
on peut supposer que $\mathcal{F}' \to \mathcal{F}$
correspond à un homomorphisme de modules $M' \to M$ ; voir
Cohomologie des schémas, lemme \ref{coherent-lemma-homs-over-open}.
Après avoir remplacé $M'$ par l'image de $M' \to M$
et posé $M'' = M/M'$, on voit que notre suite exacte courte
correspond à la restriction de la suite exacte courte de
modules cohérents associée à la suite exacte courte
$0 \to M' \to M \to M'' \to 0$ de $A$-modules.

\medskip\noindent
Soit $\hat s \in \lim H^0(U, \mathcal{F}/I^n\mathcal{F})$
d'image $\hat s'' \in \lim H^0(U, \mathcal{F}''/I^n\mathcal{F}'')$.
Par hypothèse, on trouve $V \in \mathcal{V}$
et une section $s'' \in \mathcal{F}''(V)$ qui s'envoie sur $\hat s''$.
Soit $J \subset A$ un idéal tel que $V(J) = \Spec(A) \setminus V$.
D'après Cohomologie des schémas, lemme \ref{coherent-lemma-homs-over-open},
après avoir remplacé $J$ par une puissance, on peut supposer
qu'il existe une application $A$-linéaire $\varphi : J \to M''$
correspondant à $s''$. Fixons ce choix de $J$ ; dans la suite
de la démonstration, nous remplacerons $V$ par un $V$ plus petit dans $\mathcal{V}$,
c'est-à-dire que nous aurons $V \cap V(J) = \emptyset$.

\medskip\noindent
Choisissons une présentation $A^{\oplus m} \to A^{\oplus n} \to J \to 0$.
Notons $g_1, \ldots, g_n \in J$ les images des vecteurs de base
de $A^{\oplus n}$, de sorte que $J = (g_1, \ldots, g_n)$. Soit
$A^{\oplus m} \to A^{\oplus n}$ donnée par la matrice $(a_{ji})$,
de sorte que $\sum a_{ji} g_i = 0$, $j = 1, \ldots, m$.
Comme $M \to M''$ est surjective, pour chaque $i$ on peut choisir $m_i \in M$
qui s'envoie sur $\varphi(g_i) \in M''$. Alors l'élément $g_i \hat s - m_i$
de $\lim H^0(U, \mathcal{F}/I^n\mathcal{F})$ appartient au
sous-module $\lim H^0(U, \mathcal{F}'/I^n\mathcal{F}')$.
Par hypothèse, après avoir rétréci $V$, on peut supposer qu'il existe des
$s'_i \in \mathcal{F}'(V)$, $i = 1, \ldots, n$,
tels que les $s'_i$ s'envoient sur $g_i \hat s - m_i$.
Posons $s_i = s'_i + m_i$ dans $\mathcal{F}(V)$.
Remarquons que $\sum a_{ji} s_i$ s'envoie sur $\sum a_{ji}g_i\hat s = 0$
par l'application
$$
\colim_\mathcal{V} \mathcal{F}(V')
\longrightarrow
\lim H^0(U, \mathcal{F}/I^n\mathcal{F})
$$
Comme cette application est injective (voir ci-dessus), après avoir rétréci $V$ on peut supposer
que $\sum a_{ji}s_i = 0$ dans $\mathcal{F}(V)$ pour tout $j = 1, \ldots, m$.
Il s'ensuit que l'on obtient une application de $A$-modules $J \to \mathcal{F}(V)$
qui envoie $g_i$ sur $s_i$. Par la propriété universelle de $\widetilde{J}$,
cette application de $A$-modules correspond à une application de $\mathcal{O}_V$-modules
$\widetilde{J}|_V \to \mathcal{F}$. Cependant, puisque $V(J) \cap V = \emptyset$,
on a $\widetilde{J}|_V = \mathcal{O}_V$. Nous avons donc construit
une section $s \in \mathcal{F}(V)$. Nous omettons le calcul qui
montre que $s$ s'envoie sur $\hat s$ par l'application affichée ci-dessus.
\end{proof}

\noindent
Le lemme suivant sera supplanté par la
proposition \ref{algebraization-proposition-application-H0}.

\begin{lemma}
\label{algebraization-lemma-application-H0-pre}
Soient $I \subset \mathfrak a$ des idéaux d'un anneau noethérien $A$.
Soit $\mathcal{F}$ un module cohérent sur
$U = \Spec(A) \setminus V(\mathfrak a)$.
Supposons que
\begin{enumerate}
\item $A$ soit complet pour la topologie $I$-adique et possède un complexe dualisant,
\item si $x \in \text{Ass}(\mathcal{F})$, $x \not \in V(I)$,
$\overline{\{x\}} \cap V(I) \not \subset V(\mathfrak a)$,
et $z \in \overline{\{x\}} \cap V(\mathfrak a)$, alors
$\dim(\mathcal{O}_{\overline{\{x\}}, z}) > \text{cd}(A, I) + 1$,
\item l'une des conditions suivantes soit satisfaite :
\begin{enumerate}
\item la restriction de $\mathcal{F}$ à $U \setminus V(I)$ est $(S_1)$,
\item la dimension de $V(\mathfrak a)$ est au plus $2$\footnote{Au
sens où la différence entre les valeurs maximale et minimale
sur $V(\mathfrak a)$ d'une fonction de dimension sur $\Spec(A)$ est au plus $2$.}.
\end{enumerate}
\end{enumerate}
Alors on obtient un isomorphisme
$$
\colim H^0(V, \mathcal{F})
\longrightarrow
\lim H^0(U, \mathcal{F}/I^n\mathcal{F})
$$
où la limite inductive porte sur les ouverts $V \subset U$ contenant $U \cap V(I)$.
\end{lemma}

\begin{proof}
Choisissons un $A$-module de type fini $M$ tel que $\mathcal{F}$ soit la restriction
à $U$ du module cohérent associé à $M$ ; voir Cohomologie locale,
lemme \ref{local-cohomology-lemma-finiteness-pushforwards-and-H1-local}.
Posons $d = \text{cd}(A, I)$.
Soit $\mathfrak p$ un idéal premier de $A$ n'appartenant pas à $V(I)$,
et soit $\mathfrak q \in V(\mathfrak p) \cap V(\mathfrak a)$.
Alors, ou bien $\mathfrak p$ n'est pas un idéal premier associé de $M$
et donc $\text{depth}(M_\mathfrak p) \geq 1$,
ou bien on a $\dim((A/\mathfrak p)_\mathfrak q) > d + 1$ d'après (2).
Ainsi, les hypothèses du
lemme \ref{algebraization-lemma-algebraize-local-cohomology-general}
sont satisfaites pour $s = 1$ et $d$ ; on utilise ici la condition (3).
On trouve donc un idéal
$J_0 \subset \mathfrak a$ avec $V(J_0) \cap V(I) = V(\mathfrak a)$
tel que, pour tout $J \subset J_0$ avec $V(J) \cap V(I) = V(\mathfrak a)$,
les applications
$$
H^i_J(M) \longrightarrow H^i(R\Gamma_\mathfrak a(M)^\wedge)
$$
soient des isomorphismes pour $i = 0, 1$. Considérons les morphismes de
triangles exacts
$$
\xymatrix{
R\Gamma_J(M)  \ar[d] \ar[r] &
M \ar[r] \ar[d] &
R\Gamma(V, \mathcal{F}) \ar[d] \ar[r] &
R\Gamma_J(M)[1]  \ar[d] \\
R\Gamma_J(M)^\wedge \ar[r] &
M \ar[r] &
R\Gamma(V, \mathcal{F})^\wedge \ar[r] &
R\Gamma_J(M)^\wedge[1] \\
R\Gamma_\mathfrak a(M)^\wedge \ar[r] \ar[u] &
M \ar[r] \ar[u] &
R\Gamma(U, \mathcal{F})^\wedge \ar[r] \ar[u] &
R\Gamma_\mathfrak a(M)^\wedge[1] \ar[u]
}
$$
où $V = \Spec(A) \setminus V(J)$. Rappelons que
$R\Gamma_\mathfrak a(M)^\wedge \to R\Gamma_J(M)^\wedge$
est un isomorphisme (parce que $\mathfrak a$, $\mathfrak a + I$ et $J + I$
définissent le même sous-schéma fermé ; voir par exemple la
démonstration du lemme \ref{algebraization-lemma-algebraize-local-cohomology-general}).
Par conséquent,
$R\Gamma(U, \mathcal{F})^\wedge = R\Gamma(V, \mathcal{F})^\wedge$.
On obtient ainsi un diagramme commutatif
$$
\xymatrix{
0 \ar[r] &
H^0_J(M) \ar[r] \ar[d] &
M \ar[r] \ar[d] \ar[r] &
\Gamma(V, \mathcal{F}) \ar[d] \ar[r] &
H^1_J(M) \ar[d] \ar[r] &
0 \\
0 \ar[r] &
H^0(R\Gamma_J(M)^\wedge) \ar[r] &
M \ar[r] &
H^0(R\Gamma(V, \mathcal{F})^\wedge) \ar[r] &
H^1(R\Gamma_J(M)^\wedge) \ar[r] &
0 \\
0 \ar[r] &
H^0(R\Gamma_\mathfrak a(M)^\wedge) \ar[r] \ar[u] &
M \ar[r] \ar[u] &
H^0(R\Gamma(U, \mathcal{F})^\wedge) \ar[r] \ar[u] &
H^1(R\Gamma_\mathfrak a(M)^\wedge) \ar[r] \ar[u] &
0
}
$$
dont les lignes sont exactes et dont les flèches verticales inférieures sont des isomorphismes. Ainsi,
on obtient un isomorphisme
$\Gamma(V, \mathcal{F}) \to H^0(R\Gamma(U, \mathcal{F})^\wedge)$.
D'après les lemmes \ref{algebraization-lemma-formal-functions-general}
et \ref{algebraization-lemma-derived-completion-pseudo-coherent}, on a
$$
R\Gamma(U, \mathcal{F})^\wedge =
R\Gamma(U, \mathcal{F}^\wedge) =
R\Gamma(U, R\lim \mathcal{F}/I^n\mathcal{F})
$$
et l'on trouve $H^0(R\Gamma(U, \mathcal{F})^\wedge) =
\lim H^0(U, \mathcal{F}/I^n\mathcal{F})$ d'après
Cohomologie, lemme \ref{cohomology-lemma-RGamma-commutes-with-Rlim}.
\end{proof}

\noindent
Nous appliquons maintenant un argument d'amorçage au lemme précédent afin d'éliminer la condition (3).

\begin{proposition}
\label{algebraization-proposition-application-H0}
Soient $I \subset \mathfrak a$ des idéaux d'un anneau noethérien $A$.
Soit $\mathcal{F}$ un module cohérent sur
$U = \Spec(A) \setminus V(\mathfrak a)$.
Supposons que
\begin{enumerate}
\item $A$ soit complet pour la topologie $I$-adique et possède un complexe dualisant,
\item si $x \in \text{Ass}(\mathcal{F})$, $x \not \in V(I)$,
$\overline{\{x\}} \cap V(I) \not \subset V(\mathfrak a)$,
et $z \in \overline{\{x\}} \cap V(\mathfrak a)$, alors
$\dim(\mathcal{O}_{\overline{\{x\}}, z}) > \text{cd}(A, I) + 1$.
\end{enumerate}
Alors on obtient un isomorphisme
$$
\colim H^0(V, \mathcal{F})
\longrightarrow
\lim H^0(U, \mathcal{F}/I^n\mathcal{F})
$$
où la limite inductive porte sur les ouverts $V \subset U$ contenant $U \cap V(I)$.
\end{proposition}

\begin{proof}
Soit $T \subset U$ l'ensemble des points $x$ tels que
$\overline{\{x\}} \cap V(I) \subset V(\mathfrak a)$.
Soit $\mathcal{F} \to \mathcal{F}'$ la surjection
de modules cohérents sur $U$ construite dans
Cohomologie locale, lemme \ref{local-cohomology-lemma-get-depth-1-along-Z}.
Comme $\mathcal{F} \to \mathcal{F}'$ est un isomorphisme
au-dessus d'un ouvert $V \subset U$ contenant $U \cap V(I)$,
il suffit de démontrer le lemme après avoir remplacé $\mathcal{F}$
par $\mathcal{F}'$. On peut donc supposer, et l'on suppose, que
pour $x \in U$ tel que $\overline{\{x\}} \cap V(I) \subset V(\mathfrak a)$,
on a $\text{depth}(\mathcal{F}_x) \geq 1$.

\medskip\noindent
Soit $\mathcal{V}$ l'ensemble des sous-schémas ouverts $V \subset U$
contenant $U \cap V(I)$, ordonné par l'inclusion inverse.
C'est un ensemble ordonné filtrant. Nous affirmons d'abord que
$$
\mathcal{F}(V)
\longrightarrow
\lim H^0(U, \mathcal{F}/I^n\mathcal{F})
$$
est injective pour tout $V \in \mathcal{F}$ (et, en particulier, l'application
du lemme est injective). En effet, tout point associé $x$ de $\mathcal{F}$
doit vérifier $\overline{\{x\}} \cap U \cap Y \not = \emptyset$
d'après le paragraphe précédent. Si $y \in \overline{\{x\}} \cap U \cap Y$, alors
$\mathcal{F}_x$ est un localisé de $\mathcal{F}_y$
et $\mathcal{F}_y \subset \lim \mathcal{F}_y/I^n \mathcal{F}_y$
d'après le théorème d'intersection de Krull
(Algèbre, lemme \ref{algebra-lemma-intersect-powers-ideal-module-zero}).
Cela prouve l'assertion, car une section $s \in \mathcal{F}(V)$
appartenant au noyau devrait avoir un support vide et serait donc nulle.

\medskip\noindent
Choisissons un $A$-module de type fini $M$ tel que $\mathcal{F}$ soit la restriction
de $\widetilde{M}$ à $U$ ; voir Cohomologie locale, lemme
\ref{local-cohomology-lemma-finiteness-pushforwards-and-H1-local}.
On peut supposer, et l'on suppose, que $H^0_\mathfrak a(M) = 0$.
Posons $\text{Ass}(M) \setminus V(I) = \{\mathfrak p_1, \ldots, \mathfrak p_n\}$.
Nous allons démontrer le lemme par récurrence sur $n$. Après réindexation, on
peut supposer que $\mathfrak p_n$ est un élément minimal de l'ensemble
$\{\mathfrak p_1, \ldots, \mathfrak p_n\}$ pour l'inclusion, c'est-à-dire que
$\mathfrak p_n$ est un point générique du support de $M$.
Posons
$$
M' = H^0_{\mathfrak p_1 \ldots \mathfrak p_{n - 1} I}(M)
$$
et $M'' = M/M'$. Soient $\mathcal{F}'$ et $\mathcal{F}''$ les
$\mathcal{O}_U$-modules cohérents correspondant à $M'$ et $M''$.
Complexes dualisants, lemme \ref{dualizing-lemma-divide-by-torsion},
entraîne que $M''$ ne possède qu'un seul idéal premier associé, à savoir $\mathfrak p_n$.
Ainsi, $\mathcal{F}''$ ne possède qu'un seul point associé, et
on voit que la condition (3)(a) du lemme \ref{algebraization-lemma-application-H0-pre} est satisfaite ;
l'application $\colim H^0(V, \mathcal{F}'')
\to \lim H^0(U, \mathcal{F}''/I^n\mathcal{F}'')$ est donc un isomorphisme.
D'autre part, puisque
$\mathfrak p_n \not \in V(\mathfrak p_1 \ldots \mathfrak p_{n - 1} I)$
on voit que $\mathfrak p_n$ n'est pas un idéal premier associé de $M'$.
L'hypothèse de récurrence s'applique donc à $M'$ ; remarquons
que, puisque $\mathcal{F}' \subset \mathcal{F}$,
la condition $\text{depth}(\mathcal{F}'_x) \geq 1$ est satisfaite aux points $x$ tels que
$\overline{\{x\}} \cap V(I) \subset V(\mathfrak a)$ ; voir
Algèbre, lemme \ref{algebra-lemma-depth-in-ses}.
Ainsi, l'application
$\colim H^0(V, \mathcal{F}')
\to \lim H^0(U, \mathcal{F}'/I^n\mathcal{F}')$ est également un isomorphisme.
On conclut par le lemme \ref{algebraization-lemma-ses-H0}.
\end{proof}

\begin{lemma}
\label{algebraization-lemma-application-H0}
Soient $I \subset \mathfrak a$ des idéaux d'un anneau noethérien $A$.
Soit $\mathcal{F}$ un module cohérent sur
$U = \Spec(A) \setminus V(\mathfrak a)$.
Supposons que
\begin{enumerate}
\item $A$ soit complet pour la topologie $I$-adique et possède un complexe dualisant,
\item si $x \in \text{Ass}(\mathcal{F})$, $x \not \in V(I)$,
$\overline{\{x\}} \cap V(I) \not \subset V(\mathfrak a)$, et
$z \in V(\mathfrak a) \cap \overline{\{x\}}$, alors
$\dim(\mathcal{O}_{\overline{\{x\}}, z}) > \text{cd}(A, I) + 1$,
\item pour $x \in U$ tel que $\overline{\{x\}} \cap V(I) \subset V(\mathfrak a)$,
on ait $\text{depth}(\mathcal{F}_x) \geq 2$,
\end{enumerate}
alors on obtient un isomorphisme
$$
H^0(U, \mathcal{F})
\longrightarrow
\lim H^0(U, \mathcal{F}/I^n\mathcal{F})
$$
\end{lemma}

\begin{proof}
Soit $\hat s \in \lim H^0(U, \mathcal{F}/I^n\mathcal{F})$.
D'après la proposition \ref{algebraization-proposition-application-H0},
on trouve que $\hat s$ est l'image d'un élément $s \in \mathcal{F}(V)$
pour un certain ouvert $V \subset U$ contenant $U \cap V(I)$.
Cependant, la condition (3) montre que $\text{depth}(\mathcal{F}_x) \geq 2$
pour tout $x \in U \setminus V$, et l'on trouve donc
$\mathcal{F}(V) = \mathcal{F}(U)$ d'après
Diviseurs, lemme \ref{divisors-lemma-depth-2-hartog},
ce qui achève la démonstration.
\end{proof}

\begin{lemma}
\label{algebraization-lemma-alternative-colim-H0}
Soit $A$ un anneau noethérien. Soit $f \in \mathfrak a \subset A$
un élément d'un idéal de $A$. Soit $M$ un $A$-module de type fini.
Supposons que
\begin{enumerate}
\item $A$ soit complet pour la topologie $f$-adique,
\item $f$ soit un non-diviseur de zéro dans $M$,
\item $H^1_\mathfrak a(M/fM)$ soit un $A$-module de type fini.
\end{enumerate}
Alors, pour $U = \Spec(A) \setminus V(\mathfrak a)$, l'application
$$
\colim_V \Gamma(V, \widetilde{M})
\longrightarrow
\lim \Gamma(U, \widetilde{M/f^nM})
$$
est un isomorphisme, où la limite inductive porte sur les ouverts $V \subset U$
contenant $U \cap V(f)$.
\end{lemma}

\begin{proof}
Posons $\mathcal{F} = \widetilde{M}|_U$.
Le fait que $H^1_\mathfrak a(M/fM)$ soit de type fini entraîne que
$H^0(U, \mathcal{F}/f\mathcal{F})$ est de type fini ; voir
Cohomologie locale, lemme
\ref{local-cohomology-lemma-finiteness-pushforwards-and-H1-local}.
D'après Cohomologie, lemme \ref{cohomology-lemma-limit-finite}
(qui s'applique puisque $f$ est un non-diviseur de zéro dans $\mathcal{F}$),
on voit que $N = \lim H^0(U, \mathcal{F}/f^n\mathcal{F})$
est un $A$-module de type fini, est sans $f$-torsion, et que
$N/fN \subset H^0(U, \mathcal{F}/f\mathcal{F})$.
D'autre part, on a une application $M \to N$ et une application compatible
$$
M/fM \longrightarrow H^0(U, \mathcal{F}/f\mathcal{F})
$$
Pour $g \in \mathfrak a$, on voit que $(M/fM)_g$ s'envoie isomorphiquement
sur $H^0(U \cap D(f), \mathcal{F}/f\mathcal{F})$, puisque
$\mathcal{F}/f\mathcal{F}$ est la restriction de $\widetilde{M/fM}$ à $U$.
On en conclut que $M_g \to N_g$ induit un isomorphisme
$$
M_g/fM_g = (M/fM)_g \to (N/fN)_g = N_g/fN_g
$$
Puisque $f$ est un non-diviseur de zéro à la fois dans $N$ et dans $M$, on en conclut
que $M_g \to N_g$ induit un isomorphisme sur les complétés $f$-adiques,
ce qui entraîne à son tour que $M_g \to N_g$ est un isomorphisme dans un voisinage
ouvert de $V(f) \cap D(g)$.
Comme $g \in \mathfrak a$ était arbitraire, on conclut que $M$ et $N$
définissent des modules cohérents isomorphes au-dessus d'un ouvert $V$
comme dans l'énoncé du lemme. Cela achève la démonstration.
\end{proof}

\begin{proposition}
\label{algebraization-proposition-alternative-colim-H0}
Soit $A$ un anneau noethérien. Soit $f \in \mathfrak a \subset A$
un élément d'un idéal de $A$. Soit $\mathcal{F}$ un module cohérent
sur $U = \Spec(A) \setminus V(\mathfrak a)$. Supposons que
\begin{enumerate}
\item $A$ soit complet pour la topologie $f$-adique et possède un complexe dualisant,
\item si $x \in \text{Ass}(\mathcal{F})$, $x \not \in V(f)$,
$\overline{\{x\}} \cap V(f) \not \subset V(\mathfrak a)$,
et $z \in \overline{\{x\}} \cap V(\mathfrak a)$, alors
$\dim(\mathcal{O}_{\overline{\{x\}}, z}) > 2$.
\end{enumerate}
Alors l'application
$$
\colim_V \Gamma(V, \mathcal{F})
\longrightarrow
\lim \Gamma(U, \mathcal{F}/f^n\mathcal{F})
$$
est un isomorphisme, où la limite inductive porte sur les ouverts $V \subset U$
contenant $U \cap V(f)$.
\end{proposition}

\begin{proof}[Première démonstration]
Rappelons que $A$ est universellement caténaire et à fibres formelles
de Gorenstein ; voir Complexes dualisants, lemmes
\ref{dualizing-lemma-dualizing-gorenstein-formal-fibres} et
\ref{dualizing-lemma-universally-catenary}. On peut donc considérer l'application
$\mathcal{F} \to \mathcal{F}'$ construite dans
Cohomologie locale, lemme \ref{local-cohomology-lemma-make-S2-along-Z}
pour la partie fermée $V(f) \cap U$ de $U$.
Remarquons que
\begin{enumerate}
\item Le noyau et le conoyau de $\mathcal{F} \to \mathcal{F}'$
ont leur support dans $V(f) \cap U$.
\item Le module $\mathcal{F}'$ est sans $f$-torsion, car ses germes sont de
profondeur $\geq 1$ en tout point de $V(f) \cap U$, c'est-à-dire que $\mathcal{F}'$
n'a aucun point associé dans $V(f) \cap U$.
\item Si $y \in V(f) \cap U$ est un point associé de
$\mathcal{F}'/f\mathcal{F}'$, alors
$\text{depth}(\mathcal{F}'_y) = 1$
et, par conséquent (d'après la construction de $\mathcal{F}'$),
il existe une spécialisation immédiate $x \leadsto y$,
où $x \not \in V(f)$ est un point associé de $\mathcal{F}$.
Il résulte de notre hypothèse (2) que $y$ ne peut admettre
dans $\Spec(A)$ une spécialisation immédiate
en un point $z \in V(\mathfrak a)$.
\item Il résulte de (3) que $H^0(U, \mathcal{F}'/f\mathcal{F}')$
est un $A$-module de type fini ; voir
Cohomologie locale, lemme \ref{local-cohomology-lemma-finiteness-Rjstar}.
\end{enumerate}
Ces remarques vont nous permettre d'achever la démonstration.

\medskip\noindent
Nous affirmons d'abord que le lemme vaut pour $\mathcal{F}'$.
Choisissons en effet un $A$-module de type fini $M'$ tel que $\mathcal{F}'$
soit la restriction
à $U$ du module cohérent associé à $M'$ ; voir Cohomologie locale,
lemme \ref{local-cohomology-lemma-finiteness-pushforwards-and-H1-local}.
Comme $\mathcal{F}'$ est sans $f$-torsion, on peut également supposer que $M'$
est sans $f$-torsion. La remarque (4) ci-dessus montre que
$H^1_\mathfrak a(M')$ est un $A$-module de type fini ; voir Cohomologie locale,
lemme \ref{local-cohomology-lemma-finiteness-pushforwards-and-H1-local}.
L'assertion résulte donc du lemme \ref{algebraization-lemma-alternative-colim-H0}.

\medskip\noindent
Deuxièmement, remarquons que le lemme vaut trivialement pour tout
$\mathcal{O}_U$-module cohérent à support dans $V(f) \cap U$.
Soient $\mathcal{K}$, resp.\ $\mathcal{G}$, resp.\ $\mathcal{Q}$
le noyau, resp.\ l'image, resp.\ le conoyau de l'application
$\mathcal{F} \to \mathcal{F}'$.
La suite exacte courte
$0 \to \mathcal{G} \to \mathcal{F}' \to \mathcal{Q} \to 0$
et le lemme \ref{algebraization-lemma-ses-H0} montrent que le résultat vaut pour
$\mathcal{G}$. On procède ensuite de même avec la suite exacte courte
$0 \to \mathcal{K} \to \mathcal{F} \to \mathcal{G} \to 0$
pour achever la démonstration.
\end{proof}

\begin{proof}[Deuxième démonstration]
La proposition est un cas particulier de la
proposition \ref{algebraization-proposition-application-H0}.
\end{proof}

\begin{lemma}
\label{algebraization-lemma-alternative-H0}
Soit $A$ un anneau noethérien. Soit $f \in \mathfrak a \subset A$
un élément d'un idéal de $A$. Soit $M$ un $A$-module de type fini.
Supposons que
\begin{enumerate}
\item $A$ soit complet pour la topologie $f$-adique,
\item $H^1_\mathfrak a(M)$ et $H^2_\mathfrak a(M)$ soient
annulés par une puissance de $f$.
\end{enumerate}
Alors, pour $U = \Spec(A) \setminus V(\mathfrak a)$, l'application
$$
\Gamma(U, \widetilde{M})
\longrightarrow
\lim \Gamma(U, \widetilde{M/f^nM})
$$
est un isomorphisme.
\end{lemma}

\begin{proof}
On peut appliquer le lemme \ref{algebraization-lemma-formal-functions-principal} à $U$ et à
$\mathcal{F} = \widetilde{M}|_U$, car $\mathcal{F}$ est un objet noethérien de
la catégorie des $\mathcal{O}_U$-modules cohérents.
Puisque $H^1(U, \mathcal{F}) = H^2_\mathfrak a(M)$
(Cohomologie locale, lemme
\ref{local-cohomology-lemma-finiteness-pushforwards-and-H1-local})
est annulé par une puissance de $f$, on voit que
son module de Tate $f$-adique est nul.
Le lemme montre donc que $\lim H^0(U, \mathcal{F}/f^n \mathcal{F})$
est égal au complété $f$-adique usuel de $H^0(U, \mathcal{F})$.
Considérons la suite exacte courte
$$
0 \to M/H^0_\mathfrak a(M) \to H^0(U, \mathcal{F}) \to H^1_\mathfrak a(M) \to 0
$$
de Cohomologie locale, lemme
\ref{local-cohomology-lemma-finiteness-pushforwards-and-H1-local}.
Comme $M/H^0_\mathfrak a(M)$ est un $A$-module de type fini, il est
complet ; voir Algèbre, lemme \ref{algebra-lemma-completion-tensor}.
Comme $H^1_\mathfrak a(M)$ est annulé par une puissance de $f$, on déduit
d'Algèbre, lemme \ref{algebra-lemma-completion-differ-by-torsion},
que $H^0(U, \mathcal{F})$ est également complet. Cela achève la démonstration.
\end{proof}







\section{Algébrisation des sections formelles, III}
\label{algebraization-section-algebraization-sections-coherent-III}

\noindent
La section suivante contient une liste non exhaustive
d'applications des résultats
sur la complétion de la cohomologie locale à la cohomologie supérieure
des modules cohérents sur les schémas quasi-affines et à leur
complétion relativement à un idéal.

\begin{proposition}
\label{algebraization-proposition-application-higher}
Soient $I \subset \mathfrak a$ des idéaux d'un anneau noethérien $A$.
Soit $\mathcal{F}$ un module cohérent sur
$U = \Spec(A) \setminus V(\mathfrak a)$.
Soit $s \geq 0$.
Supposons que
\begin{enumerate}
\item $A$ soit complet pour la topologie $I$-adique et possède un complexe dualisant,
\item si $x \in U \setminus V(I)$, alors
$\text{depth}(\mathcal{F}_x) > s$ ou
$$
\text{depth}(\mathcal{F}_x) +
\dim(\mathcal{O}_{\overline{\{x\}}, z}) > \text{cd}(A, I) + s + 1
$$
pour tout $z \in V(\mathfrak a) \cap \overline{\{x\}}$,
\item l'une des conditions suivantes soit satisfaite :
\begin{enumerate}
\item la restriction de $\mathcal{F}$ à $U \setminus V(I)$
soit $(S_{s + 1})$, ou
\item la dimension de $V(\mathfrak a)$ soit au plus $2$\footnote{Au
sens où la différence entre les valeurs maximale et minimale
sur $V(\mathfrak a)$ d'une fonction de dimension sur $\Spec(A)$ est au plus $2$.}.
\end{enumerate}
\end{enumerate}
Alors les applications
$$
H^i(U, \mathcal{F})
\longrightarrow
\lim H^i(U, \mathcal{F}/I^n\mathcal{F})
$$
sont des isomorphismes pour $i < s$. De plus, on a un isomorphisme
$$
\colim H^s(V, \mathcal{F})
\longrightarrow
\lim H^s(U, \mathcal{F}/I^n\mathcal{F})
$$
où la limite inductive porte sur les ouverts $V \subset U$ contenant $U \cap V(I)$.
\end{proposition}

\begin{proof}
On peut supposer $s > 0$, puisque le cas $s = 0$ a été traité dans la
proposition \ref{algebraization-proposition-application-H0}.

\medskip\noindent
Choisissons un $A$-module de type fini $M$ tel que $\mathcal{F}$ soit la restriction
à $U$ du module cohérent associé à $M$ ; voir Cohomologie locale,
lemme \ref{local-cohomology-lemma-finiteness-pushforwards-and-H1-local}.
Posons $d = \text{cd}(A, I)$.
Soit $\mathfrak p$ un idéal premier de $A$ n'appartenant pas à $V(I)$,
et soit $\mathfrak q \in V(\mathfrak p) \cap V(\mathfrak a)$.
Alors ou bien $\text{depth}(M_\mathfrak p) \geq s + 1 > s$,
ou bien $\dim((A/\mathfrak p)_\mathfrak q) > d + s + 1$ d'après (2).
D'après le lemme \ref{algebraization-lemma-bootstrap-bis-bis}, on en conclut que les
hypothèses de la situation \ref{algebraization-situation-bootstrap}
sont satisfaites pour $A, I, V(\mathfrak a), M, s, d$.
D'autre part, les hypothèses du
lemme \ref{algebraization-lemma-algebraize-local-cohomology-general}
sont satisfaites pour $s + 1$ et $d$ ; c'est ici que l'on utilise la condition (3).

\medskip\noindent
En appliquant le lemme \ref{algebraization-lemma-algebraize-local-cohomology-general},
on trouve qu'il existe un idéal
$J_0 \subset \mathfrak a$ tel que $V(J_0) \cap V(I) = V(\mathfrak a)$
et que, pour tout $J \subset J_0$ tel que $V(J) \cap V(I) = V(\mathfrak a)$,
les applications
$$
H^i_J(M) \longrightarrow H^i(R\Gamma_\mathfrak a(M)^\wedge)
$$
sont des isomorphismes pour $i \leq s + 1$.

\medskip\noindent
Pour $i \leq s$, l'application $H^i_\mathfrak a(M) \to H^i_J(M)$
est un isomorphisme d'après les lemmes \ref{algebraization-lemma-bootstrap-inherited} et
\ref{algebraization-lemma-kill-colimit-support-general}.
En utilisant la comparaison entre cohomologie et cohomologie locale
(Cohomologie locale, lemme \ref{local-cohomology-lemma-local-cohomology}),
on déduit que
$H^i(U, \mathcal{F}) \to H^i(V,\mathcal{F})$
est un isomorphisme pour $V = \Spec(A) \setminus V(J)$ et
$i < s$.

\medskip\noindent
D'après le théorème \ref{algebraization-theorem-final-bootstrap}, on a
$H^i_\mathfrak a(M) = \lim H^i_\mathfrak a(M/I^nM)$
pour $i \leq s$. D'après le lemme \ref{algebraization-lemma-combine-two}, on a
$H^{s + 1}_\mathfrak a(M) = \lim H^{s + 1}_\mathfrak a(M/I^nM)$.

\medskip\noindent
L'isomorphisme $H^0(U, \mathcal{F}) = H^0(V, \mathcal{F}) =
\lim H^0(U, \mathcal{F}/I^n\mathcal{F})$ résulte de ce qui précède et de la
proposition \ref{algebraization-proposition-application-H0}.
Pour $0 < i < s$, on obtient les isomorphismes voulus
$H^i(U, \mathcal{F}) = H^i(V, \mathcal{F}) =
\lim H^i(U, \mathcal{F}/I^n\mathcal{F})$ de la
même manière, en utilisant la relation entre cohomologie locale
et cohomologie ; c'est plus facile que le cas $i = 0$,
car, pour $i > 0$, on a
$$
H^i(U, \mathcal{F}) = H^{i + 1}_\mathfrak a(M),
\quad
H^i(V, \mathcal{F}) = H^{i + 1}_J(M),
\quad
H^i(R\Gamma(U, \mathcal{F})^\wedge) = 
H^{i + 1}(R\Gamma_\mathfrak a(M)^\wedge)
$$
On procède de même pour la dernière assertion.
\end{proof}

\begin{lemma}
\label{algebraization-lemma-alternative-higher}
Soit $A$ un anneau noethérien. Soit $f \in \mathfrak a \subset A$
un élément d'un idéal de $A$. Soit $M$ un $A$-module de type fini.
Soit $s \geq 0$. Supposons que
\begin{enumerate}
\item $A$ soit complet pour la topologie $f$-adique,
\item $H^i_\mathfrak a(M)$ soit annulé par une puissance de $f$
pour $i \leq s + 1$.
\end{enumerate}
Alors, pour $U = \Spec(A) \setminus V(\mathfrak a)$, l'application
$$
H^i(U, \widetilde{M})
\longrightarrow
\lim H^i(U, \widetilde{M/f^nM})
$$
est un isomorphisme pour $i < s$.
\end{lemma}

\begin{proof}
Procédons par récurrence sur $s$. Si $s = 0$, l'assertion est vide. Si $s = 1$, le
résultat est le lemme \ref{algebraization-lemma-alternative-H0}. Supposons $s > 1$. Par récurrence,
il suffit de démontrer le résultat pour $i = s - 1 \geq 1$.
On peut appliquer le lemme \ref{algebraization-lemma-formal-functions-principal}
à $U$ et à $\mathcal{F} = \widetilde{M}|_U$,
car $\mathcal{F}$ est un objet noethérien de
la catégorie des $\mathcal{O}_U$-modules cohérents.
Remarquons que $H^j(U, \mathcal{F}) = H^{j + 1}_\mathfrak a(M)$ pour tout $j$, d'après
Cohomologie locale, lemme
\ref{local-cohomology-lemma-finiteness-pushforwards-and-H1-local}.
Ainsi, pour $j = s = (s - 1) + 1$, ce module est
annulé par une puissance de $f$ d'après l'hypothèse.
Il résulte alors du
lemme \ref{algebraization-lemma-formal-functions-principal} que
$\lim H^{s - 1}(U, \mathcal{F}/f^n\mathcal{F})$
est le complété $f$-adique usuel de $H^{s - 1}(U, \mathcal{F})$.
En utilisant de nouveau le fait que ce module est annulé par une puissance de
$f$, on voit que le complété est simplement égal à $H^{s - 1}(U, \mathcal{F})$,
comme voulu.
\end{proof}





\section{Application à la connexité}
\label{algebraization-section-connected}

\noindent
Dans cette section, nous étudions le théorème de connexité de Grothendieck
et ses variantes ; la version originale se trouve dans
\cite[Exposee XIII, théorème 2.1]{SGA2}. Il existe dans la littérature une version
appelée théorème de connexité de Faltings ;
nous supposons qu'il s'agit de \cite[théorème 6]{Faltings-some}.
Énonçons et démontrons la version optimale pour les anneaux locaux
complets donnée dans \cite[théorème 1.6]{Varbaro}.

\begin{lemma}
\label{algebraization-lemma-punctured-still-connected}
\begin{reference}
\cite[Theorem 1.6]{Varbaro}
\end{reference}
Soit $(A, \mathfrak m)$ un anneau local noethérien complet.
Soit $I$ un idéal propre de $A$.
Posons $X = \Spec(A)$ et $Y = V(I)$.
Notons
\begin{enumerate}
\item $d$ la dimension minimale d'une composante irréductible de $X$, et
\item $c$ la dimension minimale d'une partie fermée $Z \subset X$
telle que $X \setminus Z$ soit non connexe.
\end{enumerate}
Alors, pour toute partie fermée $Z \subset Y$, l'espace $Y \setminus Z$ est connexe si
$\dim(Z) < \min(c, d - 1) - \text{cd}(A, I)$. En particulier, le spectre
épointé de $A/I$ est connexe si $\text{cd}(A, I) < \min(c, d - 1)$.
\end{lemma}

\begin{proof}
Démontrons d'abord la dernière assertion. Supposons premièrement que le spectre
épointé de $A/I$ soit vide. Alors
Cohomologie locale, lemme \ref{local-cohomology-lemma-cd-bound-dim-local},
montre que toute composante irréductible de $X$ est de dimension
$\leq \text{cd}(A, I)$, et l'on obtient $\min(c, d - 1) - \text{cd}(A, I) < 0$,
ce qui entraîne que le lemme vaut dans ce cas. On peut donc supposer que
$U \cap Y$ est non vide, où $U = X \setminus \{\mathfrak m\}$
est le spectre épointé de $A$. On peut remplacer $A$ par son réduit.
Remarquons que $A$ possède un complexe dualisant
(Complexes dualisants, lemme \ref{dualizing-lemma-ubiquity-dualizing})
et que $A$ est complet relativement à $I$
(Algèbre, lemme \ref{algebra-lemma-complete-by-sub}).
Si l'on suppose $d - 1 > \text{cd}(A, I)$, on peut appliquer le
lemme \ref{algebraization-lemma-application-theorem} pour voir que
$$
\colim H^0(V, \mathcal{O}_V)
\longrightarrow
\lim H^0(U, \mathcal{O}_U/I^n\mathcal{O}_U)
$$
est un isomorphisme, où la limite inductive porte sur les ouverts $V \subset U$
contenant $U \cap Y$. Si $U \cap Y$ n'est pas connexe, alors
son $n$-ième voisinage infinitésimal dans $U$ n'est pas connexe
pour tout $n$, et l'on trouve que le
membre de droite possède un idempotent non trivial (on utilise ici
que $U \cap Y$ est non vide).
On peut donc trouver un $V$ qui n'est pas connexe.
Posons $Z = X \setminus V$. D'après
Cohomologie locale, lemme \ref{local-cohomology-lemma-cd-bound-dim-local},
on voit que toute composante irréductible de $Z$ est de dimension
$\leq \text{cd}(A, I)$. Ainsi $c \leq \text{cd}(A, I)$, ce qui
démontre bien la dernière assertion.

\medskip\noindent
On peut déduire l'énoncé du lemme de ce que l'on vient de démontrer,
comme suit. Supposons que $Z \subset Y$ soit fermé, que $Y \setminus Z$ ne soit pas
connexe et que $\dim(Z) = e$. Rappelons que, par convention, un espace connexe est non vide.
On conclut donc que soit (a) $Y = Z$, soit (b)
$Y \setminus Z = W_1 \amalg W_2$, avec les $W_i$ non vides, ouverts et fermés
dans $Y \setminus Z$. Dans le cas (b), on peut choisir des points $w_i \in W_i$
qui sont fermés dans $U$ ; voir
Morphismes, lemme \ref{morphisms-lemma-ubiquity-Jacobson-schemes}.
On peut alors trouver $f_1, \ldots, f_e \in \mathfrak m$
tels que $V(f_1, \ldots, f_e) \cap Z = \{\mathfrak m\}$
et, dans le cas (b), on peut supposer $w_i \in V(f_1, \ldots, f_e)$.
En effet, par récurrence et en évitant les idéaux premiers, on peut
choisir $f_i$ tel que $\dim V(f_1, \ldots, f_i) \cap Z = e - i$
et tel que, dans le cas (b), on ait $w_1, w_2 \in V(f_i)$.
Il s'ensuit que le spectre épointé de $A/I + (f_1, \ldots, f_e)$
n'est pas connexe (on omet un petit détail). Puisque
$\text{cd}(A, I + (f_1, \ldots, f_e)) \leq \text{cd}(A, I) + e$ d'après
Cohomologie locale, lemmes \ref{local-cohomology-lemma-cd-sum} et
\ref{local-cohomology-lemma-bound-cd}, on conclut que
$$
\text{cd}(A, I) + e \geq \min(c, d - 1)
$$
d'après la première partie de la démonstration. Cela entraîne
$e \geq \min(c, d - 1) - \text{cd}(A, I)$, ce qu'il fallait démontrer.
\end{proof}

\begin{lemma}
\label{algebraization-lemma-connected}
Soient $I \subset \mathfrak a$ des idéaux d'un anneau noethérien $A$.
Supposons que
\begin{enumerate}
\item $A$ soit complet pour la topologie $I$-adique et possède un complexe dualisant,
\item si $\mathfrak p \subset A$ est un idéal premier minimal n'appartenant pas
à $V(I)$ et si $\mathfrak q \in V(\mathfrak p) \cap V(\mathfrak a)$, alors
$\dim((A/\mathfrak p)_\mathfrak q) > \text{cd}(A, I) + 1$,
\item tout ouvert non vide $V \subset \Spec(A)$ qui contient
$V(I) \setminus V(\mathfrak a)$ soit connexe\footnote{Par exemple,
si $A$ est un domaine.}.
\end{enumerate}
Alors $V(I) \setminus V(\mathfrak a)$ est vide ou connexe.
\end{lemma}

\begin{proof}
On peut remplacer $A$ par son réduit. On dispose alors de l'inégalité
de (2) pour tous les idéaux premiers associés de $A$. D'après la
proposition \ref{algebraization-proposition-application-H0}, on voit que
$$
\colim H^0(V, \mathcal{O}_V) = \lim H^0(T_n, \mathcal{O}_{T_n})
$$
où la limite inductive porte sur les ouverts $V$ comme dans (3), et où $T_n$ est le
$n$-ième voisinage infinitésimal de $T = V(I) \setminus V(\mathfrak a)$
dans $U = \Spec(A) \setminus V(\mathfrak a)$. Ainsi, $T$ est vide
ou connexe ; sinon, le membre de droite posséderait en effet un
idempotent non trivial, tandis que l'on a supposé que le membre de gauche n'en possède pas.
On omet quelques détails.
\end{proof}

\begin{lemma}
\label{algebraization-lemma-H0}
Soit $A$ un domaine noethérien qui possède un complexe dualisant
et qui est complet relativement à un élément non nul $f \in A$.
Soit $f \in \mathfrak a \subset A$ un idéal.
Supposons que toute composante irréductible de $Z = V(\mathfrak a)$
soit de codimension $> 2$ dans $X = \Spec(A)$, c'est-à-dire que toute
composante irréductible de $Z$ soit de codimension $> 1$ dans $Y = V(f)$.
Alors $Y \setminus Z$ est connexe.
\end{lemma}

\begin{proof}
C'est un cas particulier du lemme \ref{algebraization-lemma-connected} (dont la démonstration
repose sur la proposition \ref{algebraization-proposition-application-H0}). Nous le
démontrons ci-dessous à l'aide de la proposition \ref{algebraization-proposition-alternative-colim-H0}, plus facile.

\medskip\noindent
Posons $U = X \setminus Z$. D'après la proposition \ref{algebraization-proposition-alternative-colim-H0},
on a un isomorphisme
$$
\colim \Gamma(V, \mathcal{O}_V) \to
\lim_n \Gamma(U, \mathcal{O}_U/f^n \mathcal{O}_U)
$$
où la limite inductive porte sur les ouverts $V \subset U$ contenant $U \cap Y$.
Ainsi, si $U \cap Y$ n'est pas connexe, il existe pour un certain $V$
un idempotent non trivial dans $\Gamma(V, \mathcal{O}_V)$. C'est impossible,
car $V$ est un schéma intègre puisque $X$ est le spectre d'un domaine.
\end{proof}






\section{Le foncteur de complétion}
\label{algebraization-section-completion}

\noindent
Soit $X$ un schéma noethérien. Soit $Y \subset X$ un sous-schéma fermé
de faisceau quasi-cohérent d'idéaux $\mathcal{I} \subset \mathcal{O}_X$.
Dans cette section, nous considérons les systèmes projectifs de
$\mathcal{O}_X$-modules cohérents $(\mathcal{F}_n)$ tels que $\mathcal{F}_n$ soit
annulé par $I^n$ et que les applications de transition induisent des
isomorphismes $\mathcal{F}_{n + 1}/I^n\mathcal{F}_{n + 1} \to \mathcal{F}_n$.
La catégorie de ces systèmes projectifs a été notée
$$
\textit{Coh}(X, \mathcal{I})
$$
dans Cohomologie des schémas, section \ref{coherent-section-coherent-formal}.
Cette catégorie est équivalente à la catégorie des modules cohérents
sur le complété formel de $X$ le long de $Y$ ; cependant, comme nous n'avons
pas encore introduit les schémas formels ni les modules cohérents sur ceux-ci,
nous ne pouvons employer ici cette terminologie. Nous nous intéressons tout particulièrement
au foncteur de complétion
$$
\textit{Coh}(\mathcal{O}_X)
\longrightarrow
\textit{Coh}(X, \mathcal{I}),\quad
\mathcal{F} \longmapsto \mathcal{F}^\wedge
$$
Voir
Cohomologie des schémas, équation (\ref{coherent-equation-completion-functor}).

\begin{lemma}
\label{algebraization-lemma-completion-fully-faithful}
Soit $X$ un schéma noethérien et soit $Y \subset X$ un sous-schéma fermé.
Soit $Y_n \subset X$ le $n$-ième voisinage infinitésimal de $Y$ dans $X$.
Considérons les conditions suivantes :
\begin{enumerate}
\item $X$ est quasi-affine et
$\Gamma(X, \mathcal{O}_X) \to \lim \Gamma(Y_n, \mathcal{O}_{Y_n})$
est un isomorphisme,
\item $X$ possède un module inversible ample $\mathcal{L}$ et
$\Gamma(X, \mathcal{L}^{\otimes m}) \to
\lim \Gamma(Y_n, \mathcal{L}^{\otimes m}|_{Y_n})$
est un isomorphisme pour tout $m \gg 0$,
\item pour tout $\mathcal{O}_X$-module de type fini et localement libre
$\mathcal{E}$, l'application
$\Gamma(X, \mathcal{E}) \to \lim \Gamma(Y_n, \mathcal{E}|_{Y_n})$
est un isomorphisme, et
\item le foncteur de complétion
$\textit{Coh}(\mathcal{O}_X) \to \textit{Coh}(X, \mathcal{I})$
est pleinement fidèle sur la sous-catégorie pleine des
objets de type fini et localement libres.
\end{enumerate}
Alors (1) $\Rightarrow$ (2) $\Rightarrow$ (3) $\Rightarrow$ (4)
et (4) $\Rightarrow$ (3).
\end{lemma}

\begin{proof}
Preuve de (3) $\Rightarrow$ (4). Si $\mathcal{F}$ et $\mathcal{G}$
sont de type fini et localement libres sur $X$, alors, en considérant
$\mathcal{H} = \SheafHom_{\mathcal{O}_X}(\mathcal{G}, \mathcal{F})$
et en utilisant Cohomologie des schémas, lemme
\ref{coherent-lemma-completion-internal-hom}
on voit que (3) entraîne (4).

\medskip\noindent
Preuve de (2) $\rightarrow$ (3). Soit en effet $\mathcal{L}$ ample
sur $X$, et supposons que $\mathcal{E}$ soit un
$\mathcal{O}_X$-module de type fini et localement libre.
Nous affirmons que l'on peut trouver une suite universellement exacte
$$
0 \to \mathcal{E} \to
(\mathcal{L}^{\otimes p})^{\oplus r} \to
(\mathcal{L}^{\otimes q})^{\oplus s}
$$
pour certains $r, s \geq 0$ et $0 \ll p \ll q$. Si tel est le cas, alors,
en utilisant la suite exacte
$$
0 \to \lim \Gamma(\mathcal{E}|_{Y_n}) \to
\lim \Gamma((\mathcal{L}^{\otimes p})^{\oplus r}|_{Y_n}) \to
\lim \Gamma((\mathcal{L}^{\otimes q})^{\oplus s}|_{Y_n})
$$
et les isomorphismes de (2), on obtient l'isomorphisme de (3).
Pour démontrer l'assertion, considérons le module localement libre dual
$\SheafHom_{\mathcal{O}_X}(\mathcal{E}, \mathcal{O}_X)$
et appliquons
Propriétés, proposition \ref{properties-proposition-characterize-ample},
pour trouver une surjection
$$
(\mathcal{L}^{\otimes -p})^{\oplus r}
\longrightarrow
\SheafHom_{\mathcal{O}_X}(\mathcal{E}, \mathcal{O}_X)
$$
En prenant les duaux, on obtient la première application de la suite exacte
(elle est universellement injective, car le fait d'être une surjection est universel).
En répétant l'opération avec le conoyau, on obtient la seconde. On omet quelques détails.

\medskip\noindent
Preuve de (1) $\Rightarrow$ (2). C'est vrai parce que, si $X$ est quasi-affine,
alors $\mathcal{O}_X$ est un module inversible ample ; voir
Propriétés, lemme \ref{properties-lemma-quasi-affine-O-ample}.

\medskip\noindent
Nous omettons la preuve de (4) $\Rightarrow$ (3).
\end{proof}

\noindent
Étant donnés un schéma noethérien et un faisceau quasi-cohérent d'idéaux
$\mathcal{I} \subset \mathcal{O}_X$, nous dirons
qu'un objet $(\mathcal{F}_n)$ de $\textit{Coh}(X, \mathcal{I})$
est {\it de type fini et localement libre} si chaque $\mathcal{F}_n$ est un
$\mathcal{O}_X/\mathcal{I}^n$-module de type fini et localement libre.

\begin{lemma}
\label{algebraization-lemma-completion-fully-faithful-general}
Soit $X$ un schéma noethérien et soit $Y \subset X$ un sous-schéma fermé
de faisceau d'idéaux $\mathcal{I} \subset \mathcal{O}_X$.
Soit $Y_n \subset X$ le $n$-ième voisinage infinitésimal de $Y$ dans $X$.
Soit $\mathcal{V}$ l'ensemble des sous-schémas ouverts $V \subset X$ contenant $Y$,
ordonné par l'inclusion inverse.
\begin{enumerate}
\item $X$ est quasi-affine et
$$
\colim_\mathcal{V} \Gamma(V, \mathcal{O}_V)
\longrightarrow
\lim \Gamma(Y_n, \mathcal{O}_{Y_n})
$$
est un isomorphisme,
\item $X$ possède un module inversible ample $\mathcal{L}$ et
$$
\colim_\mathcal{V} \Gamma(V, \mathcal{L}^{\otimes m})
\longrightarrow
\lim \Gamma(Y_n, \mathcal{L}^{\otimes m}|_{Y_n})
$$
est un isomorphisme pour tout $m \gg 0$,
\item pour tout $V \in \mathcal{V}$ et tout
$\mathcal{O}_V$-module $\mathcal{E}$ de type fini et localement libre, l'application
$$
\colim_{V' \geq V} \Gamma(V', \mathcal{E}|_{V'})
\longrightarrow
\lim \Gamma(Y_n, \mathcal{E}|_{Y_n})
$$
est un isomorphisme, et
\item le foncteur de complétion
$$
\colim_\mathcal{V} \textit{Coh}(\mathcal{O}_V)
\longrightarrow
\textit{Coh}(X, \mathcal{I}),
\quad
\mathcal{F} \longmapsto \mathcal{F}^\wedge
$$
est pleinement fidèle sur la sous-catégorie pleine des
objets de type fini et localement libres (voir l'explication ci-dessus).
\end{enumerate}
Alors (1) $\Rightarrow$ (2) $\Rightarrow$ (3) $\Rightarrow$ (4)
et (4) $\Rightarrow$ (3).
\end{lemma}

\begin{proof}
Remarquons que $\mathcal{V}$ est un ensemble ordonné filtrant ; les limites inductives sont donc
comme dans Catégories, section \ref{categories-section-directed-colimits}.
La suite de l'argument est presque exactement la même que l'argument
de la démonstration du lemme \ref{algebraization-lemma-completion-fully-faithful} ; nous invitons
le lecteur à l'omettre.

\medskip\noindent
Preuve de (3) $\Rightarrow$ (4). Si $\mathcal{F}$ et $\mathcal{G}$
sont de type fini et localement libres sur $V \in \mathcal{V}$, alors, en considérant
$\mathcal{H} = \SheafHom_{\mathcal{O}_V}(\mathcal{G}, \mathcal{F})$
et en utilisant Cohomologie des schémas, lemme
\ref{coherent-lemma-completion-internal-hom}
on voit que (3) entraîne (4).

\medskip\noindent
Preuve de (2) $\Rightarrow$ (3). Soit $\mathcal{L}$ ample
sur $X$, et supposons que $\mathcal{E}$ soit un
$\mathcal{O}_V$-module de type fini et localement libre
pour un certain $V \in \mathcal{V}$.
Nous affirmons que l'on peut trouver une suite universellement exacte
$$
0 \to \mathcal{E} \to
(\mathcal{L}^{\otimes p})^{\oplus r}|_{V} \to
(\mathcal{L}^{\otimes q})^{\oplus s}|_{V}
$$
pour certains $r, s \geq 0$ et $0 \ll p \ll q$. Si tel est le cas,
l'isomorphisme de (2) entraînera l'isomorphisme de (3).
Pour démontrer l'assertion, rappelons que $\mathcal{L}|_V$ est ample ; voir
Propriétés, lemme \ref{properties-lemma-ample-on-locally-closed}.
Considérons le module localement libre dual
$\SheafHom_{\mathcal{O}_V}(\mathcal{E}, \mathcal{O}_V)$
et appliquons
Propriétés, proposition \ref{properties-proposition-characterize-ample},
pour trouver une surjection
$$
(\mathcal{L}^{\otimes -p})^{\oplus r}|_V \longrightarrow
\SheafHom_{\mathcal{O}_V}(\mathcal{E}, \mathcal{O}_V)
$$
(elle est universellement injective, car le fait d'être une surjection est universel).
En prenant les duaux, on obtient la première application de la suite exacte.
En répétant l'opération avec le conoyau, on obtient la seconde. On omet quelques détails.

\medskip\noindent
Preuve de (1) $\Rightarrow$ (2). C'est vrai parce que, si $X$ est quasi-affine,
alors $\mathcal{O}_X$ est un module inversible ample ; voir
Propriétés, lemme \ref{properties-lemma-quasi-affine-O-ample}.

\medskip\noindent
Nous omettons la preuve de (4) $\Rightarrow$ (3).
\end{proof}

\begin{lemma}
\label{algebraization-lemma-recognize-formal-coherent-modules}
Soit $X$ un schéma noethérien. Soit $\mathcal{I} \subset \mathcal{O}_X$
un faisceau quasi-cohérent d'idéaux. Le foncteur
$$
\textit{Coh}(X, \mathcal{I}) \longrightarrow \text{Pro-}\QCoh(\mathcal{O}_X)
$$
est pleinement fidèle ; voir Catégories, remarque \ref{categories-remark-pro-category}.
\end{lemma}

\begin{proof}
Soient $(\mathcal{F}_n)$ et $(\mathcal{G}_n)$ des objets de
$\textit{Coh}(X, \mathcal{I})$. Un morphisme de pro-objets
$\alpha$ de $(\mathcal{F}_n)$ vers $(\mathcal{G}_n)$ est donné
par un système d'applications
$\alpha_n : \mathcal{F}_{m(n)} \to \mathcal{G}_n$ compatibles avec
les applications de transition, où $\mathbf{N} \to \mathbf{N}$, $n \mapsto m(n)$,
est une fonction croissante (en particulier $m(n) \geq n$). Puisque
$\mathcal{F}_n = \mathcal{F}_{m(n)}/\mathcal{I}^n\mathcal{F}_{m(n)}$
et puisque $\mathcal{G}_n$ est annulé par $\mathcal{I}^n$,
on voit que $\alpha_n$ induit une application $\mathcal{F}_n \to \mathcal{G}_n$.
\end{proof}

\noindent
Nous donnons ensuite quelques exemples du type de résultat de pleine fidélité
que nous pourrons démontrer à l'aide du travail effectué plus haut dans ce chapitre.

\begin{lemma}
\label{algebraization-lemma-fully-faithful}
Soient $I \subset \mathfrak a$ des idéaux d'un anneau noethérien $A$.
Soit $U = \Spec(A) \setminus V(\mathfrak a)$. Supposons que
\begin{enumerate}
\item $A$ soit complet pour la topologie $I$-adique et possède un complexe dualisant,
\item pour tout idéal premier associé $\mathfrak p \subset A$ tel que
$\mathfrak p \not \in V(I)$ et
$V(\mathfrak p) \cap V(I) \not \subset V(\mathfrak a)$ et
$\mathfrak q \in V(\mathfrak p) \cap V(\mathfrak a)$, on ait
$\dim((A/\mathfrak p)_\mathfrak q) > \text{cd}(A, I) + 1$,
\item pour $\mathfrak p \subset A$ tel que $\mathfrak p \not \in V(I)$
et $V(\mathfrak p) \cap V(I) \subset V(\mathfrak a)$,
on ait $\text{depth}(A_\mathfrak p) \geq 2$.
\end{enumerate}
Alors le foncteur de complétion
$$
\textit{Coh}(\mathcal{O}_U)
\longrightarrow
\textit{Coh}(U, I\mathcal{O}_U),
\quad
\mathcal{F} \longmapsto \mathcal{F}^\wedge
$$
est pleinement fidèle sur la sous-catégorie pleine des
objets de type fini et localement libres.
\end{lemma}

\begin{proof}
D'après le lemme \ref{algebraization-lemma-completion-fully-faithful},
il suffit de montrer que
$$
\Gamma(U, \mathcal{O}_U) =
\lim \Gamma(U, \mathcal{O}_U/I^n\mathcal{O}_U)
$$
Cela résulte immédiatement du
lemme \ref{algebraization-lemma-application-H0}.
\end{proof}

\begin{lemma}
\label{algebraization-lemma-fully-faithful-alternative}
Soit $A$ un anneau noethérien. Soit $f \in \mathfrak a \subset A$
un élément d'un idéal de $A$. Soit $U = \Spec(A) \setminus V(\mathfrak a)$.
Supposons que
\begin{enumerate}
\item $A$ soit complet pour la topologie $f$-adique,
\item $H^1_\mathfrak a(A)$ et $H^2_\mathfrak a(A)$ soient
annulés par une puissance de $f$.
\end{enumerate}
Alors le foncteur de complétion
$$
\textit{Coh}(\mathcal{O}_U)
\longrightarrow
\textit{Coh}(U, I\mathcal{O}_U),
\quad
\mathcal{F} \longmapsto \mathcal{F}^\wedge
$$
est pleinement fidèle sur la sous-catégorie pleine des
objets de type fini et localement libres.
\end{lemma}

\begin{proof}
D'après le lemme \ref{algebraization-lemma-completion-fully-faithful},
il suffit de montrer que
$$
\Gamma(U, \mathcal{O}_U) =
\lim \Gamma(U, \mathcal{O}_U/I^n\mathcal{O}_U)
$$
Cela résulte immédiatement du
lemme \ref{algebraization-lemma-alternative-H0}.
\end{proof}

\begin{lemma}
\label{algebraization-lemma-fully-faithful-simple-two}
Soit $A$ un anneau noethérien. Soit $f \in \mathfrak a$ un élément d'un
idéal de $A$. Soit $U = \Spec(A) \setminus V(\mathfrak a)$. Supposons que
\begin{enumerate}
\item $A$ possède un complexe dualisant et soit complet relativement à $f$,
\item pour tout idéal premier $\mathfrak p \subset A$, $f \not \in \mathfrak p$
et $\mathfrak q \in V(\mathfrak p) \cap V(\mathfrak a)$, on ait
$\text{depth}(A_\mathfrak p) + \dim((A/\mathfrak p)_\mathfrak q) > 2$.
\end{enumerate}
Alors le foncteur de complétion
$$
\textit{Coh}(\mathcal{O}_U)
\longrightarrow
\textit{Coh}(U, I\mathcal{O}_U),
\quad
\mathcal{F} \longmapsto \mathcal{F}^\wedge
$$
est pleinement fidèle sur la sous-catégorie pleine des objets de type fini et localement libres.
\end{lemma}

\begin{proof}
Cela résulte du lemme \ref{algebraization-lemma-fully-faithful-alternative} et de
Cohomologie locale, proposition \ref{local-cohomology-proposition-annihilator}.
\end{proof}

\begin{lemma}
\label{algebraization-lemma-fully-faithful-general}
Soient $I \subset \mathfrak a \subset A$ des idéaux d'un anneau noethérien $A$.
Soit $U = \Spec(A) \setminus V(\mathfrak a)$. Soit $\mathcal{V}$
l'ensemble des sous-schémas ouverts de $U$ contenant $U \cap V(I)$,
ordonné par l'inclusion inverse. Supposons que
\begin{enumerate}
\item $A$ soit complet pour la topologie $I$-adique et possède un complexe dualisant,
\item pour tout idéal premier associé
$\mathfrak p \subset A$ tel que
$I \not \subset \mathfrak p$ et
$V(\mathfrak p) \cap V(I) \not \subset V(\mathfrak a)$
et $\mathfrak q \in V(\mathfrak p) \cap V(\mathfrak a)$, on ait
$\dim((A/\mathfrak p)_\mathfrak q) > \text{cd}(A, I) + 1$.
\end{enumerate}
Alors le foncteur de complétion
$$
\colim_\mathcal{V} \textit{Coh}(\mathcal{O}_V)
\longrightarrow
\textit{Coh}(U, I\mathcal{O}_U),
\quad
\mathcal{F} \longmapsto \mathcal{F}^\wedge
$$
est pleinement fidèle sur la sous-catégorie pleine des
objets de type fini et localement libres.
\end{lemma}

\begin{proof}
D'après le lemme \ref{algebraization-lemma-completion-fully-faithful-general},
il suffit de montrer que
$$
\colim_\mathcal{V} \Gamma(V, \mathcal{O}_V) =
\lim \Gamma(U, \mathcal{O}_U/I^n\mathcal{O}_U)
$$
Cela résulte immédiatement de la proposition \ref{algebraization-proposition-application-H0}.
\end{proof}

\begin{lemma}
\label{algebraization-lemma-fully-faithful-general-alternative}
Soit $A$ un anneau noethérien. Soit $f \in \mathfrak a \subset A$
un élément d'un idéal de $A$. Soit $U = \Spec(A) \setminus V(\mathfrak a)$.
Soit $\mathcal{V}$ l'ensemble des sous-schémas ouverts de $U$ contenant $U \cap V(f)$,
ordonné par l'inclusion inverse. Supposons que
\begin{enumerate}
\item $A$ soit complet pour la topologie $f$-adique,
\item $f$ soit un non-diviseur de zéro,
\item $H^1_\mathfrak a(A/fA)$ soit un $A$-module de type fini.
\end{enumerate}
Alors le foncteur de complétion
$$
\colim_\mathcal{V} \textit{Coh}(\mathcal{O}_V)
\longrightarrow
\textit{Coh}(U, f\mathcal{O}_U),
\quad
\mathcal{F} \longmapsto \mathcal{F}^\wedge
$$
est pleinement fidèle sur la sous-catégorie pleine des objets de type fini et localement libres.
\end{lemma}

\begin{proof}
D'après le lemme \ref{algebraization-lemma-completion-fully-faithful-general},
il suffit de montrer que
$$
\colim_\mathcal{V} \Gamma(V, \mathcal{O}_V) =
\lim \Gamma(U, \mathcal{O}_U/I^n\mathcal{O}_U)
$$
Cela résulte immédiatement du lemme \ref{algebraization-lemma-alternative-colim-H0}.
\end{proof}

\begin{lemma}
\label{algebraization-lemma-fully-faithful-very-general}
Soient $I \subset \mathfrak a \subset A$ des idéaux d'un anneau noethérien $A$.
Soit $U = \Spec(A) \setminus V(\mathfrak a)$. Soit $\mathcal{V}$ l'ensemble
des sous-schémas ouverts de $U$ contenant $U \cap V(I)$, ordonné par l'inclusion
inverse. Soient $\mathcal{F}$ et
$\mathcal{G}$ des $\mathcal{O}_V$-modules cohérents pour un certain
$V \in \mathcal{V}$. L'application
$$
\colim_{V' \geq V} \Hom_V(\mathcal{G}|_{V'}, \mathcal{F}|_{V'})
\longrightarrow
\Hom_{\textit{Coh}(U, I\mathcal{O}_U)}(\mathcal{G}^\wedge, \mathcal{F}^\wedge)
$$
est bijective si les hypothèses suivantes sont satisfaites :
\begin{enumerate}
\item $A$ soit complet pour la topologie $I$-adique et possède un complexe dualisant,
\item si $x \in \text{Ass}(\mathcal{F})$, $x \not \in V(I)$,
$\overline{\{x\}} \cap V(I) \not \subset V(\mathfrak a)$
et $z \in \overline{\{x\}} \cap V(\mathfrak a)$, alors
$\dim(\mathcal{O}_{\overline{\{x\}}, z}) > \text{cd}(A, I) + 1$.
\end{enumerate}
\end{lemma}

\begin{proof}
On peut choisir des $\mathcal{O}_U$-modules cohérents
$\mathcal{F}'$ et $\mathcal{G}'$ dont les restrictions à $V$
sont $\mathcal{F}$ et $\mathcal{G}$ ; voir
Propriétés, lemme \ref{properties-lemma-lift-finite-presentation}.
On peut modifier le choix de $\mathcal{F}'$ pour assurer que
$\text{Ass}(\mathcal{F}') \subset V$ ; voir par exemple
Cohomologie locale, lemme \ref{local-cohomology-lemma-get-depth-1-along-Z}.
On peut donc remplacer, et l'on remplace, $V$ par $U$, et $\mathcal{F}$ et $\mathcal{G}$
par $\mathcal{F}'$ et $\mathcal{G}'$.
Posons $\mathcal{H} = \SheafHom_{\mathcal{O}_U}(\mathcal{G}, \mathcal{F})$.
C'est un $\mathcal{O}_U$-module cohérent. On a
$$
\Hom_V(\mathcal{G}|_V, \mathcal{F}|_V) =
H^0(V, \mathcal{H})
\quad\text{et}\quad
\lim H^0(U, \mathcal{H}/\mathcal{I}^n\mathcal{H}) =
\Mor_{\textit{Coh}(U, I\mathcal{O}_U)}
(\mathcal{G}^\wedge, \mathcal{F}^\wedge)
$$
Voir Cohomologie des schémas, lemme \ref{coherent-lemma-completion-internal-hom}.
Ainsi, si l'on peut montrer que les hypothèses de la
proposition \ref{algebraization-proposition-application-H0}
sont satisfaites pour $\mathcal{H}$, la démonstration est achevée.
C'est le cas parce que
$\text{Ass}(\mathcal{H}) \subset \text{Ass}(\mathcal{F})$.
Voir Cohomologie des schémas, lemme
\ref{coherent-lemma-hom-into-depth}.
\end{proof}








\section{Algébrisation des modules formels cohérents, I}
\label{algebraization-section-algebraization-modules}

\noindent
La surjectivité essentielle du foncteur de complétion (voir ci-dessous)
a été étudiée systématiquement dans
\cite{SGA2}, \cite{MRaynaud-book} et \cite{MRaynaud-paper}.
Nous nous plaçons dans la situation affine suivante.

\begin{situation}
\label{algebraization-situation-algebraize}
Ici, $A$ est un anneau noethérien et $I \subset \mathfrak a \subset A$ sont des idéaux.
On pose $X = \Spec(A)$, $Y = V(I) = \Spec(A/I)$ et
$Z = V(\mathfrak a) = \Spec(A/\mathfrak a)$. De plus, $U = X \setminus Z$.
\end{situation}

\noindent
Dans cette section, nous cherchons des conditions assurant qu'un objet
de $\textit{Coh}(U, I\mathcal{O}_U)$ appartient à l'image du foncteur de complétion
$\textit{Coh}(\mathcal{O}_U) \to \textit{Coh}(U, I\mathcal{O}_U)$.
Voir Cohomologie des schémas, section \ref{coherent-section-coherent-formal} et
section \ref{algebraization-section-completion}.

\begin{lemma}
\label{algebraization-lemma-system-of-modules}
Dans la situation \ref{algebraization-situation-algebraize}.
Considérons un système projectif $(M_n)$ de $A$-modules tel
que
\begin{enumerate}
\item $M_n$ soit un $A$-module de type fini,
\item $M_n$ soit annulé par $I^n$,
\item le noyau et le conoyau de $M_{n + 1}/I^nM_{n + 1} \to M_n$
soient de torsion annulée par une puissance de $\mathfrak a$.
\end{enumerate}
Alors $(\widetilde{M}_n|_U)$ appartient à $\textit{Coh}(U, I\mathcal{O}_U)$.
Réciproquement, tout objet de $\textit{Coh}(U, I\mathcal{O}_U)$
s'obtient de cette manière.
\end{lemma}

\begin{proof}
Nous omettons de vérifier que $(\widetilde{M}_n|_U)$ appartient à
$\textit{Coh}(U, I\mathcal{O}_U)$. Soit $(\mathcal{F}_n)$
un objet de $\textit{Coh}(U, I\mathcal{O}_U)$.
D'après Cohomologie locale, lemme
\ref{local-cohomology-lemma-finiteness-pushforwards-and-H1-local}
on voit que $\mathcal{F}_n = \widetilde{M_n}$ pour un certain
$A/I^n$-module de type fini $M_n$. Après avoir quotienté $M_n$ par $H^0_\mathfrak a(M_n)$,
on peut supposer $M_n \subset H^0(U, \mathcal{F}_n)$ ; voir
Complexes dualisants, lemme \ref{dualizing-lemma-divide-by-torsion},
et le lemme déjà cité.
En remplaçant par récurrence $M_{n + 1}$ par l'image réciproque
de $M_n$ par l'application $M_{n + 1} \to H^0(U, \mathcal{F}_{n + 1})
\to H^0(U, \mathcal{F}_n)$, on peut supposer que $M_{n + 1}$ s'envoie dans
$M_n$. On obtient ainsi un système projectif $(M_n)$ satisfaisant (1) et (2)
tel que $\mathcal{F}_n = \widetilde{M_n}$. Pour voir que (3)
est satisfaite, on utilise le fait que $M_{n + 1}/I^nM_{n + 1} \to M_n$ est une application
de $A$-modules de type fini qui induit un isomorphisme après
application de $\widetilde{\ }$ et restriction à $U$
(on utilise ici une fois de plus le premier lemme cité).
\end{proof}

\noindent
Dans la situation \ref{algebraization-situation-algebraize}, on peut étudier le foncteur de complétion de
Cohomologie des schémas, équation (\ref{coherent-equation-completion-functor}),
\begin{equation}
\label{algebraization-equation-completion}
\textit{Coh}(\mathcal{O}_U)
\longrightarrow
\textit{Coh}(U, I\mathcal{O}_U),\quad
\mathcal{F} \longmapsto \mathcal{F}^\wedge
\end{equation}
Si $A$ est complet pour la topologie $I$-adique, ce foncteur est pleinement fidèle
sur des sous-catégories appropriées d'après nos travaux antérieurs sur l'algébrisation des
sections formelles ; voir la section \ref{algebraization-section-completion} et le
lemme \ref{algebraization-lemma-fully-faithful-inequalities} pour quelques exemples de résultats.
Soit ensuite $(\mathcal{F}_n)$ un objet de $\textit{Coh}(U, I\mathcal{O}_U)$.
En supposant toujours $A$ complet pour la topologie $I$-adique, on peut demander :
quand $(\mathcal{F}_n)$ appartient-il à l'image essentielle du foncteur de complétion
affiché ci-dessus ?

\begin{lemma}
\label{algebraization-lemma-essential-image-completion}
Dans la situation \ref{algebraization-situation-algebraize}, soit $(\mathcal{F}_n)$
un objet de $\textit{Coh}(U, I\mathcal{O}_U)$. Considérons les
conditions suivantes :
\begin{enumerate}
\item $(\mathcal{F}_n)$ appartient à l'image essentielle
du foncteur (\ref{algebraization-equation-completion}),
\item $(\mathcal{F}_n)$ est le complété d'un
$\mathcal{O}_U$-module cohérent,
\item $(\mathcal{F}_n)$ est le complété d'un
$\mathcal{O}_V$-module cohérent pour un ouvert $U \cap Y \subset V \subset U$,
\item $(\mathcal{F}_n)$ est le complété de
la restriction à $U$ d'un $\mathcal{O}_X$-module cohérent,
\item $(\mathcal{F}_n)$ est la restriction à $U$ du
complété d'un $\mathcal{O}_X$-module cohérent,
\item il existe un objet $(\mathcal{G}_n)$ de
$\textit{Coh}(X, I\mathcal{O}_X)$ dont la restriction
à $U$ est $(\mathcal{F}_n)$.
\end{enumerate}
Alors les conditions (1), (2), (3), (4) et (5) sont équivalentes et entraînent (6).
Si $A$ est complet pour la topologie $I$-adique, la condition (6) entraîne les autres.
\end{lemma}

\begin{proof}
Les parties (1) et (2) sont équivalentes, car le complété d'un
$\mathcal{O}_U$-module cohérent $\mathcal{F}$ est par définition l'image de
$\mathcal{F}$ par le foncteur (\ref{algebraization-equation-completion}).
Si $V \subset U$ est un sous-schéma ouvert contenant $U \cap Y$, alors on a
$$
\textit{Coh}(V, I\mathcal{O}_V) =
\textit{Coh}(U, I\mathcal{O}_U)
$$
car la catégorie des $\mathcal{O}_V$-modules cohérents à support dans
$V \cap Y$ est la même que la catégorie des $\mathcal{O}_U$-modules cohérents
à support dans $U \cap Y$. Ainsi, le complété d'un
$\mathcal{O}_V$-module cohérent est un objet de $\textit{Coh}(U, I\mathcal{O}_U)$.
Ceci étant dit, l'équivalence de (2), (3), (4) et (5)
vaut parce que les foncteurs
$\textit{Coh}(\mathcal{O}_X) \to \textit{Coh}(\mathcal{O}_U) \to
\textit{Coh}(\mathcal{O}_V)$ sont essentiellement surjectifs.
Voir Propriétés, lemme \ref{properties-lemma-lift-finite-presentation}.

\medskip\noindent
La condition (5) entraîne toujours (6). Supposons $A$ complet pour la topologie $I$-adique.
Alors tout objet de $\textit{Coh}(X, I\mathcal{O}_X)$ correspond à un
$A$-module de type fini d'après Cohomologie des schémas, lemme
\ref{coherent-lemma-inverse-systems-affine}.
On voit ainsi que (6) entraîne (5) dans ce cas.
\end{proof}

\begin{example}
\label{algebraization-example-not-algebraizable}
Soit $k$ un corps. Soit $A = k[x, y][[t]]$, avec $I = (t)$ et
$\mathfrak a = (x, y, t)$. Employons les notations de la
situation \ref{algebraization-situation-algebraize}. Remarquons que
$U \cap Y = (D(x) \cap Y) \cup (D(y) \cap Y)$ est un recouvrement
ouvert affine. Pour $n \geq 1$, considérons le module inversible
$\mathcal{L}_n$ sur $\mathcal{O}_U/t^n\mathcal{O}_U$
obtenu en recollant $A_x/t^nA_x$ et $A_y/t^nA_y$ au moyen de l'élément inversible
de $A_{xy}/t^nA_{xy}$ qui est l'image d'une série formelle quelconque
de la forme
$$
u = 1 + \frac{t}{xy} + \sum_{n \geq 2} a_n \frac{t^n}{(xy)^{\varphi(n)}}
$$
où $a_n \in k[x, y]$ et $\varphi(n) \in \mathbf{N}$.
Alors $(\mathcal{L}_n)$ est un objet inversible de
$\textit{Coh}(U, I\mathcal{O}_U)$ qui n'est pas le
complété d'un $\mathcal{O}_U$-module cohérent $\mathcal{L}$.
Nous ne faisons qu'esquisser l'argument et omettons la plupart des détails.
Soit $y \in U \cap Y$. Le complété du germe
$\mathcal{L}_y$ serait alors un module inversible, donc $\mathcal{L}_y$
est inversible. Il existerait ainsi un ouvert $V \subset U$
contenant $U \cap Y$ tel que $\mathcal{L}|_V$ soit inversible.
D'après Diviseurs, lemme \ref{divisors-lemma-extend-invertible-module},
on trouve un $A$-module inversible $M$ tel que
$\widetilde{M}|_V \cong \mathcal{L}|_V$. Cependant, l'anneau $A$
est factoriel ; on voit donc que $M \cong A$, ce qui entraînerait
$\mathcal{L}_n \cong \mathcal{O}_U/I^n\mathcal{O}_U$.
Comme $\mathcal{L}_2 \not \cong \mathcal{O}_U/I^2\mathcal{O}_U$
par construction, on obtient la contradiction voulue.

\medskip\noindent
Remarquons que, si l'on prend $a_n = 0$ pour $n \geq 2$, on voit
que $\lim H^0(U, \mathcal{L}_n)$ est non nul : dans ce cas,
la fonction $x$ sur $D(x)$ et la fonction $x + t/y$ sur $D(y)$
se recollent. D'autre part, si l'on prend $a_n = 1$ et $\varphi(n) = 2^n$,
ou même $\varphi(n) = n^2$, le lecteur peut montrer que
$\lim H^0(U, \mathcal{L}_n)$ est nul ; cela fournit une autre démonstration
du fait que $(\mathcal{L}_n)$ n'est pas algébrisable dans ce cas.
\end{example}

\noindent
Si, dans la situation \ref{algebraization-situation-algebraize}, l'anneau $A$ n'est pas
complet pour la topologie $I$-adique, le lemme \ref{algebraization-lemma-essential-image-completion}
suggère qu'il convient de demander si $(\mathcal{F}_n)$
appartient à l'image essentielle du foncteur de restriction
$$
\textit{Coh}(X, I\mathcal{O}_X)
\longrightarrow
\textit{Coh}(U, I\mathcal{O}_U)
$$
Cependant, on ne peut plus dire que cela signifie que $(\mathcal{F}_n)$
est algébrisable. Nous introduisons donc la terminologie suivante.

\begin{definition}
\label{algebraization-definition-algebraizable}
Dans la situation \ref{algebraization-situation-algebraize}, soit $(\mathcal{F}_n)$ un
objet de $\textit{Coh}(U, I\mathcal{O}_U)$. Nous dirons que
{\it $(\mathcal{F}_n)$ s'étend à $X$} s'il existe un objet
$(\mathcal{G}_n)$ de $\textit{Coh}(X, I\mathcal{O}_X)$ dont la restriction
à $U$ est isomorphe à $(\mathcal{F}_n)$.
\end{definition}

\noindent
Cette notion équivaut à être algébrisable sur la complétion.

\begin{lemma}
\label{algebraization-lemma-algebraizable}
Dans la situation \ref{algebraization-situation-algebraize}, soit $(\mathcal{F}_n)$ un
objet de $\textit{Coh}(U, I\mathcal{O}_U)$. Soient $A', I', \mathfrak a'$
les complétions $I$-adiques de $A, I, \mathfrak a$. Posons $X' = \Spec(A')$
et $U' = X' \setminus V(\mathfrak a')$. Les conditions suivantes sont équivalentes :
\begin{enumerate}
\item $(\mathcal{F}_n)$ s'étend à $X$, et
\item l'image réciproque de $(\mathcal{F}_n)$ sur $U'$ est le complété
d'un $\mathcal{O}_{U'}$-module cohérent.
\end{enumerate}
\end{lemma}

\begin{proof}
Rappelons que $A \to A'$ est un homomorphisme plat d'anneaux qui induit un isomorphisme
$A/I \to A'/I'$. Voir
Algèbre, lemmes \ref{algebra-lemma-completion-flat} et
\ref{algebra-lemma-completion-complete}.
Ainsi, $X' \to X$ est un
morphisme plat induisant un isomorphisme $Y' \to Y$. Ainsi, $U' \to U$
est un morphisme plat qui induit un isomorphisme $U' \cap Y' \to U \cap Y$.
Cela entraîne que, dans le diagramme commutatif
$$
\xymatrix{
\textit{Coh}(X', I\mathcal{O}_{X'}) \ar[r] &
\textit{Coh}(U', I\mathcal{O}_{U'}) \\
\textit{Coh}(X, I\mathcal{O}_X) \ar[u] \ar[r] &
\textit{Coh}(U, I\mathcal{O}_U) \ar[u]
}
$$
les foncteurs verticaux sont des équivalences. Voir
Cohomologie des schémas, lemme
\ref{coherent-lemma-inverse-systems-pullback-equivalence}.
Le lemme résulte formellement de ceci et des résultats du
lemme \ref{algebraization-lemma-essential-image-completion}.
\end{proof}

\noindent
Dans la situation \ref{algebraization-situation-algebraize}, soit $(\mathcal{F}_n)$ un
objet de $\textit{Coh}(U, I\mathcal{O}_U)$. Pour déterminer si
$(\mathcal{F}_n)$ s'étend à $X$, il est naturel de considérer le $A$-module
\begin{equation}
\label{algebraization-equation-guess}
M = \lim H^0(U, \mathcal{F}_n)
\end{equation}
Remarquons que $M$ possède une topologie de limite qui est (a priori) moins fine que
la topologie $I$-adique, puisque $M \to H^0(U, \mathcal{F}_n)$ annule $I^nM$.
Il existe des applications canoniques
$$
\widetilde{M}|_U \to \widetilde{M/I^nM}|_U \to
\widetilde{H^0(U, \mathcal{F}_n)}|_U \to
\mathcal{F}_n
$$
On pourrait espérer que $\widetilde{M}$ se restreigne en un module cohérent
sur $U$ et que $(\mathcal{F}_n)$ soit le complété de ce module.
C'est naïf, car cela n'a presque aucune chance d'être vrai
si $A$ n'est pas complet. Mais, même si $A$ est complet pour la topologie $I$-adique,
cette notion est très difficile à manier.
Une approche moins naïve consiste à considérer la condition suivante.

\begin{definition}
\label{algebraization-definition-canonically-algebraizable}
Dans la situation \ref{algebraization-situation-algebraize}, soit $(\mathcal{F}_n)$ un
objet de $\textit{Coh}(U, I\mathcal{O}_U)$. Nous dirons que
{\it $(\mathcal{F}_n)$ s'étend canoniquement à $X$} si le
système projectif
$$
\{\widetilde{H^0(U, \mathcal{F}_n)}\}_{n \geq 1}
$$
dans $\QCoh(\mathcal{O}_X)$ est pro-isomorphe à un objet
$(\mathcal{G}_n)$ de $\textit{Coh}(X, I\mathcal{O}_X)$.
\end{definition}

\noindent
Nous verrons dans le lemme \ref{algebraization-lemma-canonically-algebraizable}
que la condition de la définition \ref{algebraization-definition-canonically-algebraizable}
est plus forte que celle de la définition \ref{algebraization-definition-algebraizable}.

\begin{lemma}
\label{algebraization-lemma-canonically-algebraizable}
Dans la situation \ref{algebraization-situation-algebraize}, soit $(\mathcal{F}_n)$ un
objet de $\textit{Coh}(U, I\mathcal{O}_U)$. Si $(\mathcal{F}_n)$
s'étend canoniquement à $X$, alors
\begin{enumerate}
\item $(\widetilde{H^0(U, \mathcal{F}_n)})$ est pro-isomorphe à un
objet $(\mathcal{G}_n)$ de $\textit{Coh}(X, I \mathcal{O}_X)$,
unique à isomorphisme unique près,
\item la restriction de $(\mathcal{G}_n)$ à $U$ est isomorphe
à $(\mathcal{F}_n)$, c'est-à-dire que $(\mathcal{F}_n)$ s'étend à $X$,
\item le système projectif $\{H^0(U, \mathcal{F}_n)\}$
satisfait la condition de Mittag-Leffler, et
\item le module $M$ de (\ref{algebraization-equation-guess}) est de type fini sur le
complété $I$-adique de $A$, et la topologie de limite sur
$M$ est la topologie $I$-adique.
\end{enumerate}
\end{lemma}

\begin{proof}
L'existence de $(\mathcal{G}_n)$ dans (1) résulte de la
définition \ref{algebraization-definition-canonically-algebraizable}.
L'unicité de $(\mathcal{G}_n)$ dans (1) résulte du
lemme \ref{algebraization-lemma-recognize-formal-coherent-modules}.
Écrivons $\mathcal{G}_n = \widetilde{M_n}$.
Alors $\{M_n\}$ est un système projectif de $A$-modules de type fini
tel que $M_n = M_{n + 1}/I^n M_{n + 1}$.
D'après la définition \ref{algebraization-definition-canonically-algebraizable},
le système projectif $\{H^0(U, \mathcal{F}_n)\}$
est pro-isomorphe à $\{M_n\}$.
On voit donc que le système projectif $\{H^0(U, \mathcal{F}_n)\}$
satisfait la condition de Mittag-Leffler et que
$M = \lim M_n$ (comme modules topologiques).
Les propriétés de $M$ de (4) résultent ainsi d'
Algèbre, lemmes \ref{algebra-lemma-limit-complete},
\ref{algebra-lemma-finite-over-complete-ring} et
\ref{algebra-lemma-hathat-finitely-generated}.
Comme $U$ est quasi-affine, les applications canoniques
$$
\widetilde{H^0(U, \mathcal{F}_n)}|_U \to \mathcal{F}_n
$$
sont des isomorphismes (Propriétés, lemme
\ref{properties-lemma-quasi-coherent-quasi-affine}).
On en conclut que $(\mathcal{G}_n|_U)$ et $(\mathcal{F}_n)$ sont
pro-isomorphes, donc isomorphes d'après le
lemme \ref{algebraization-lemma-recognize-formal-coherent-modules}.
\end{proof}

\begin{lemma}
\label{algebraization-lemma-canonically-extend-base-change}
Dans la situation \ref{algebraization-situation-algebraize}, soit $(\mathcal{F}_n)$ un
objet de $\textit{Coh}(U, I\mathcal{O}_U)$. Soit $A \to A'$ un homomorphisme plat
d'anneaux. Posons $X' = \Spec(A')$, soit $U' \subset X'$ l'image réciproque de $U$,
et notons $g : U' \to U$ le morphisme induit. Posons
$(\mathcal{F}'_n) = (g^*\mathcal{F}_n)$ ; voir
Cohomologie des schémas, lemme \ref{coherent-lemma-inverse-systems-pullback}.
Si $(\mathcal{F}_n)$ s'étend canoniquement à $X$, alors
$(\mathcal{F}'_n)$ s'étend canoniquement à $X'$.
De plus, l'extension obtenue dans le lemme \ref{algebraization-lemma-canonically-algebraizable}
pour $(\mathcal{F}_n)$ donne par image réciproque l'extension de
$(\mathcal{F}'_n)$.
\end{lemma}

\begin{proof}
Soit $f : X' \to X$ le morphisme induit.
On a $H^0(U', \mathcal{F}'_n) = H^0(U, \mathcal{F}_n) \otimes_A A'$ par
changement de base plat ; voir Cohomologie des schémas, lemme
\ref{coherent-lemma-flat-base-change-cohomology}.
Ainsi, si $(\mathcal{G}_n)$ dans $\textit{Coh}(X, I\mathcal{O}_X)$
est pro-isomorphe à $(\widetilde{H^0(U, \mathcal{F}_n)})$, alors
$(f^*\mathcal{G}_n)$ est pro-isomorphe à
$$
(f^*\widetilde{H^0(U, \mathcal{F}_n)}) =
(\widetilde{H^0(U, \mathcal{F}_n) \otimes_A A'}) =
(\widetilde{H^0(U', \mathcal{F}'_n)})
$$
Cela achève la démonstration.
\end{proof}

\begin{lemma}
\label{algebraization-lemma-when-done}
Dans la situation \ref{algebraization-situation-algebraize}, soit $(\mathcal{F}_n)$ un objet
de $\textit{Coh}(U, I\mathcal{O}_U)$. Soit $M$ comme dans (\ref{algebraization-equation-guess}).
Supposons que
\begin{enumerate}
\item[(a)] le système projectif $H^0(U, \mathcal{F}_n)$ satisfait la condition de Mittag-Leffler,
\item[(b)] la topologie de limite sur $M$ coïncide avec la topologie $I$-adique, et
\item[(c)] l'image de $M \to H^0(U, \mathcal{F}_n)$ est un $A$-module de type fini
pour tout $n$.
\end{enumerate}
Alors $(\mathcal{F}_n)$ s'étend canoniquement à $X$.
En particulier, si $A$ est complet pour la topologie $I$-adique, alors
$(\mathcal{F}_n)$ est le complété d'un $\mathcal{O}_U$-module cohérent.
\end{lemma}

\begin{proof}
Puisque $H^0(U, \mathcal{F}_n)$ satisfait la condition de Mittag-Leffler
et puisque la topologie de limite sur $M$ est la topologie $I$-adique,
on voit que $\{M/I^nM\}$ et $\{H^0(U, \mathcal{F}_n)\}$
sont des systèmes projectifs pro-isomorphes de $A$-modules.
Ainsi, si l'on pose
$$
\mathcal{G}_n = \widetilde{M/I^n M}
$$
on voit que, pour vérifier la condition de la
définition \ref{algebraization-definition-canonically-algebraizable},
il suffit de montrer que $M$ est un module de type fini sur le
complété $I$-adique de $A$. Cela résulte
du fait que $M/I^n M$ est de type fini d'après la condition (c),
de ce qui précède et d'
Algèbre, lemme \ref{algebra-lemma-finite-over-complete-ring}.
\end{proof}

\noindent
Le résultat suivant est, en un certain sens, l'application la plus directe possible du
lemme \ref{algebraization-lemma-when-done} ci-dessus.

\begin{lemma}
\label{algebraization-lemma-algebraization-principal-variant}
Dans la situation \ref{algebraization-situation-algebraize}, soit
$(\mathcal{F}_n)$ un objet de $\textit{Coh}(U, I\mathcal{O}_U)$.
Supposons que
\begin{enumerate}
\item $I = (f)$ soit un idéal principal, pour un non-diviseur de zéro $f \in \mathfrak a$,
\item $\mathcal{F}_n$ soit un
$\mathcal{O}_U/f^n\mathcal{O}_U$-module de type fini et localement libre,
\item $H^1_\mathfrak a(A/fA)$ et $H^2_\mathfrak a(A/fA)$
soient des $A$-modules de type fini.
\end{enumerate}
Alors $(\mathcal{F}_n)$ s'étend canoniquement à $X$. En particulier, si $A$
est complet, alors $(\mathcal{F}_n)$ est le complété d'un
$\mathcal{O}_U$-module cohérent.
\end{lemma}

\begin{proof}
Nous allons le démontrer en vérifiant les hypothèses (a), (b) et (c) du
lemme \ref{algebraization-lemma-when-done}.

\medskip\noindent
Comme $\mathcal{F}_n$ est localement libre sur $\mathcal{O}_U/f^n\mathcal{O}_U$,
on voit que l'on a des suites exactes courtes
$0 \to \mathcal{F}_n \to \mathcal{F}_{n + 1} \to \mathcal{F}_1 \to 0$
pour tout $n$. La condition (b) est donc satisfaite d'après
Cohomologie, lemme \ref{cohomology-lemma-topology-I-adic-f}.

\medskip\noindent
Comme $f$ est un non-diviseur de zéro, on obtient des suites exactes courtes
$$
0 \to A/f^nA \xrightarrow{f} A/f^{n + 1}A \to A/fA \to 0
$$
et l'on a les suites exactes courtes correspondantes
$0 \to \mathcal{F}_n \to \mathcal{F}_{n + 1} \to \mathcal{F}_1 \to 0$.
Nous utiliserons
Cohomologie locale, lemme
\ref{local-cohomology-lemma-finiteness-pushforwards-and-H1-local}
sans autre mention. Nos hypothèses entraînent que
$H^0(U, \mathcal{O}_U/f\mathcal{O}_U)$ et
$H^1(U, \mathcal{O}_U/f\mathcal{O}_U)$
sont des $A$-modules de type fini. Il en va donc de même pour $\mathcal{F}_1$ ; voir
Cohomologie locale, lemme
\ref{local-cohomology-lemma-finiteness-for-finite-locally-free}.
Par récurrence et à l'aide des suites exactes courtes, on trouve que
$H^0(U, \mathcal{F}_n)$ sont des $A$-modules de type fini pour tout $n$.
On voit ainsi que l'hypothèse (c) est satisfaite.

\medskip\noindent
Enfin, comme $H^1(U, \mathcal{F}_1)$ est un $A$-module de type fini,
on peut appliquer Cohomologie, lemme \ref{cohomology-lemma-ML},
pour voir que l'hypothèse (a) est satisfaite.
\end{proof}

\begin{remark}
\label{algebraization-remark-interesting-case-variant}
Dans le lemme \ref{algebraization-lemma-algebraization-principal-variant},
si $A$ est universellement caténaire à fibres formelles
de Cohen--Macaulay (par exemple si $A$ possède un complexe dualisant), alors
la condition selon laquelle
$H^1_\mathfrak a(A/fA)$ et $H^2_\mathfrak a(A/fA)$
sont des $A$-modules de type fini équivaut à
$$
\text{depth}((A/f)_\mathfrak p) + \dim((A/\mathfrak p)_\mathfrak q) > 2
$$
pour tout $\mathfrak p \in V(f) \setminus V(\mathfrak a)$
et $\mathfrak q \in V(\mathfrak p) \cap V(\mathfrak a)$,
d'après Cohomologie locale, théorème \ref{local-cohomology-theorem-finiteness}.

\medskip\noindent
Par exemple, si $A/fA$ est $(S_2)$ et si toute composante irréductible
de $Z = V(\mathfrak a)$ est de codimension $\geq 3$
dans $Y = \Spec(A/fA)$, alors on obtient que
$H^1_\mathfrak a(A/fA)$ et $H^2_\mathfrak a(A/fA)$ sont de type fini.
Il faut comparer cela aux conditions légèrement plus faibles
obtenues dans le lemme \ref{algebraization-lemma-algebraization-principal}
(voir aussi la remarque \ref{algebraization-remark-interesting-case}).
\end{remark}






\section{Algébrisation des modules formels cohérents, II}
\label{algebraization-section-algebraization-modules-general}

\noindent
Nous poursuivons la discussion commencée dans la
section \ref{algebraization-section-algebraization-modules}.
Cette section peut être omise lors d'une première lecture.

\begin{lemma}
\label{algebraization-lemma-map-kernel-cokernel-on-closed}
Dans la situation \ref{algebraization-situation-algebraize}. Soit
$(\mathcal{F}_n) \to (\mathcal{F}'_n)$ un morphisme de
$\textit{Coh}(U, I\mathcal{O}_U)$
dont le noyau et le conoyau sont annulés par une puissance de $I$. Alors
\begin{enumerate}
\item $(\mathcal{F}_n)$ s'étend à $X$ si et seulement si
$(\mathcal{F}'_n)$ s'étend à $X$, et
\item $(\mathcal{F}_n)$ est le complété d'un $\mathcal{O}_U$-module cohérent
si et seulement si $(\mathcal{F}'_n)$ l'est.
\end{enumerate}
\end{lemma}

\begin{proof}
La partie (2) résulte immédiatement de
Cohomologie des schémas, lemme \ref{coherent-lemma-existence-easy}.
Pour voir la partie (1), on utilise d'abord le lemme \ref{algebraization-lemma-algebraizable}
pour se ramener au cas où $A$ est complet pour la topologie $I$-adique.
Cependant, dans ce cas, (1) se ramène à (2) d'après le
lemme \ref{algebraization-lemma-essential-image-completion}.
\end{proof}

\noindent
Les deux lemmes suivants étaient initialement utilisés dans la démonstration
du lemme \ref{algebraization-lemma-when-done}. Nous les conservons ici pour le lecteur qui
souhaite connaître les résultats intermédiaires que l'on peut obtenir.

\begin{lemma}
\label{algebraization-lemma-when-ML}
Dans la situation \ref{algebraization-situation-algebraize}, soit $(\mathcal{F}_n)$ un objet
de $\textit{Coh}(U, I\mathcal{O}_U)$. Si le système projectif
$H^0(U, \mathcal{F}_n)$ satisfait la condition de Mittag-Leffler, alors les applications canoniques
$$
\widetilde{M/I^nM}|_U \to \mathcal{F}_n
$$
sont surjectives pour tout $n$, où $M$ est comme dans (\ref{algebraization-equation-guess}).
\end{lemma}

\begin{proof}
La surjectivité peut se vérifier sur le germe en un point $y \in Y \setminus Z$.
Si $y$ correspond à l'idéal premier $\mathfrak q \subset A$, on peut
choisir $f \in \mathfrak a$, $f \not \in \mathfrak q$. Il suffit alors
de montrer que
$$
M_f \longrightarrow H^0(U, \mathcal{F}_n)_f = H^0(D(f), \mathcal{F}_n)
$$
est surjective, puisque $D(f)$ est affine (l'égalité vaut d'après Propriétés,
lemme \ref{properties-lemma-invert-f-sections}). Comme on dispose de la
propriété de Mittag-Leffler, on trouve que
$$
\Im(M \to H^0(U, \mathcal{F}_n)) =
\Im(H^0(U, \mathcal{F}_m) \to H^0(U, \mathcal{F}_n))
$$
pour un certain $m \geq n$. À l'aide de la suite exacte longue de cohomologie, on voit
que
$$
\Coker(H^0(U, \mathcal{F}_m) \to H^0(U, \mathcal{F}_n))
\subset
H^1(U, \Ker(\mathcal{F}_m \to \mathcal{F}_n))
$$
Comme $U = X \setminus V(\mathfrak a)$, ce $H^1$ est de torsion annulée par une puissance de
$\mathfrak a$. Ainsi, après inversion de $f$, le conoyau devient nul.
\end{proof}

\begin{lemma}
\label{algebraization-lemma-when-topology}
Dans la situation \ref{algebraization-situation-algebraize}, soit $(\mathcal{F}_n)$ un objet
de $\textit{Coh}(U, I\mathcal{O}_U)$. Soit $M$ comme dans (\ref{algebraization-equation-guess}).
Posons
$$
\mathcal{G}_n = \widetilde{M/I^nM}.
$$
Si la topologie de limite sur $M$ coïncide avec la topologie $I$-adique, alors
$\mathcal{G}_n|_U$ est un
$\mathcal{O}_U$-module cohérent et l'application de systèmes projectifs
$$
(\mathcal{G}_n|_U) \longrightarrow (\mathcal{F}_n)
$$
est injective dans la catégorie abélienne $\textit{Coh}(U, I\mathcal{O}_U)$.
\end{lemma}

\begin{proof}
Remarquons que $\mathcal{G}_n$ est un $\mathcal{O}_X$-module quasi-cohérent
annulé par $I^n$ et que
$\mathcal{G}_{n + 1}/I^n\mathcal{G}_{n + 1} = \mathcal{G}_n$.
Considérons
$$
M_n = \Im(M \longrightarrow H^0(U, \mathcal{F}_n))
$$
L'hypothèse dit que les systèmes projectifs $(M_n)$ et
$(M/I^nM)$ sont isomorphes comme pro-objets de $\text{Mod}_A$.
Choisissons $f \in \mathfrak a$ de sorte que $D(f) \subset U$ soit un ouvert affine. On a alors
$$
(M_n)_f \subset H^0(U, \mathcal{F}_n)_f = H^0(D(f), \mathcal{F}_n)
$$
L'égalité vaut d'après Propriétés, lemme \ref{properties-lemma-invert-f-sections}.
Ainsi, $\widetilde{M_n}|_U \to \mathcal{F}_n$ est injective.
Il s'ensuit que $\widetilde{M_n}|_U$ est un module cohérent
(Cohomologie des schémas, lemme
\ref{coherent-lemma-coherent-Noetherian-quasi-coherent-sub-quotient}).
Puisque $M \to M/I^nM$ est surjective et se factorise par
$M_{n'} \to M/I^nM$ pour un certain $n' \geq n$, on trouve que $\mathcal{G}_n|_U$
est cohérent comme quotient d'un module cohérent.
Avec les remarques initiales de la démonstration, on en conclut
que $(\mathcal{G}_n|_U)$ forme bien un objet
de $\textit{Coh}(U, I\mathcal{O}_U)$.
Enfin, pour montrer l'injectivité de l'application,
il suffit de montrer que
$$
\lim (M/I^nM)_f = \lim H^0(D(f), \mathcal{G}_n) \to
\lim H^0(D(f), \mathcal{F}_n)
$$
est injective ; voir Cohomologie des schémas, lemmes
\ref{coherent-lemma-inverse-systems-abelian} et
\ref{coherent-lemma-inverse-systems-affine}.
L'injectivité de $\lim (M_n)_f \to \lim H^0(D(f), \mathcal{F}_n)$
est claire (voir ci-dessus), et notre remarque sur les pro-systèmes donne
$\lim (M_n)_f = \lim (M/I^nM)_f$. Cela achève la démonstration.
\end{proof}





\section{Une fonction de distance}
\label{algebraization-section-distance}

\noindent
Soit $Y$ un schéma noethérien et soit $Z \subset Y$ une partie fermée.
Nous définissons une fonction
\begin{equation}
\label{algebraization-equation-delta-Z}
\delta^Y_Z = \delta_Z : Y \longrightarrow \mathbf{Z}_{\geq 0} \cup \{\infty\}
\end{equation}
qui mesure la « distance » d'un point de $Y$ à $Z$.
Pour une discussion informelle, voir la remarque \ref{algebraization-remark-discussion}.
Soit $y \in Y$. On pose $\delta_Z(y) = \infty$ si $y$ appartient
à une composante connexe de $Y$ qui ne rencontre pas $Z$.
Si $y$ appartient à une composante connexe de $Y$ qui rencontre $Z$,
on peut trouver $k \geq 0$ et un système
$$
V_0 \subset W_0 \supset V_1 \subset W_1 \supset \ldots \supset
V_k \subset W_k
$$
de sous-schémas fermés intègres de $Y$ tel que $V_0 \subset Z$
et que $y \in W_k$ soit le point générique. Posons
$c_i = \text{codim}(V_i, W_i)$ pour $i = 0, \ldots, k$
et $b_i = \text{codim}(V_{i + 1}, W_i)$ pour $i = 0, \ldots, k - 1$.
Pour un tel système, on pose
$$
\delta(V_0, W_0, V_1, \ldots, W_k) =
k +
\max_{i = 0, 1, \ldots, k}
(c_i + c_{i + 1} + \ldots + c_k - b_i - b_{i + 1} - \ldots - b_{k - 1})
$$
Ceci est $\geq k$, car on peut prendre $i = k$ et l'on a $c_k \geq 0$.
Enfin, on pose
$$
\delta_Z(y) = \min \delta(V_0, W_0, V_1, \ldots, W_k)
$$
où le minimum porte sur tous les systèmes de sous-schémas fermés intègres de $Y$
comme ci-dessus.

\begin{lemma}
\label{algebraization-lemma-discussion}
Soit $Y$ un schéma noethérien et soit $Z \subset Y$ une partie fermée.
\begin{enumerate}
\item Pour $y \in Y$, on a $\delta_Z(y) = 0 \Leftrightarrow y \in Z$.
\item Les parties $\{y \in Y \mid \delta_Z(y) \leq k\}$ sont
stables par spécialisation.
\item Pour $y \in Y$ et $z \in \overline{\{y\}} \cap Z$, on a
$\dim(\mathcal{O}_{\overline{\{y\}}, z}) \geq \delta_Z(y)$.
\item Si $\delta$ est une fonction de dimension sur $Y$, alors
$\delta_Z(y) \geq \delta(y) - \delta_{max}$, où $\delta_{max}$
est la valeur maximale de $\delta$ sur $Z$.
\item Si $Y = \Spec(A)$ est le spectre d'un anneau local noethérien caténaire
d'idéal maximal $\mathfrak m$ et si $Z = \{\mathfrak m\}$, alors
$\delta_Z(y) = \dim(\overline{\{y\}})$.
\item Si $Y' \subset Y$ est un sous-schéma ouvert, alors
$\delta^{Y'}_{Y' \cap Z}(y') \geq \delta^Y_Z(y')$ pour $y' \in Y'$.
\end{enumerate}
Supposons $Y$ caténaire. Alors
\begin{enumerate}
\setcounter{enumi}{6}
\item Soit $y' \leadsto y$ une spécialisation immédiate de points de $Y$.
Si $Y$ est caténaire, alors $\delta_Z(y') \leq \delta_Z(y) + 1$.
\item Étant donné un diagramme de spécialisations
$$
\xymatrix{
& y'_0 \ar@{~>}[ld] \ar@{~>}[rd] &
& y'_1 \ar@{~>}[ld] & \ldots
& y'_{k - 1} \ar@{~>}[rd] &
\\
y_0 & &
y_1 & &
\ldots & &
y_k = y
}
$$
entre points de $Y$, avec $y_0 \in Z$ et $y_i' \leadsto y_i$
une spécialisation immédiate, on a $\delta_Z(y_k) \leq k$.
\end{enumerate}
\end{lemma}

\begin{proof}
Preuve de (1). Si $y \in Z$, on peut prendre $k = 0$ et
$V_0 = W_0 = \overline{\{y\}}$, et l'on obtient $\delta(V_0, W_0) = 0$, donc
$\delta_Z(y) = 0$. Si $y \not \in Z$, alors, pour tout système
$V_0 \subset W_0 \supset V_1 \subset W_1 \supset \ldots \subset W_k$
pour $y$, ou bien $k = 0$ et $V_0 \not = W_0$, ou bien $k > 0$.
Dans les deux cas, $\delta(V_0, W_0, \ldots, W_k) > 0$. Ainsi, $\delta_Z(y) > 0$.

\medskip\noindent
Preuve de (2). Soit $y \leadsto y'$ une spécialisation non triviale et soit
$V_0 \subset W_0 \supset V_1 \subset W_1 \supset \ldots \subset W_k$
un système pour $y$. Il y a deux cas.
Cas I : $V_k = W_k$, c'est-à-dire $c_k = 0$. Dans ce cas,
on peut poser $V'_k = W'_k = \overline{\{y'\}}$.
Un calcul facile montre que
$\delta(V_0, W_0, \ldots, V'_k, W'_k) \leq
\delta(V_0, W_0, \ldots, V_k, W_k)$
car seul $b_{k - 1}$ est remplacé par un entier plus grand.
Cas II : $V_k \not = W_k$, c'est-à-dire $c_k > 0$.
Dans ce cas, en posant $V_{k + 1} = W_{k + 1} = \overline{\{y'\}}$,
on voit que $V_0 \subset W_0 \supset \ldots \subset W_k \supset V_{k + 1}
\subset W_{k + 1}$ est un système pour $y'$. Alors $c_{k + 1} = 0$ et
$b_k > 0$, de sorte que l'on obtient
\begin{align*}
& \delta(V_0, \ldots, W_{k + 1}) \\
& =
k + 1 +
\max_{i = 0, 1, \ldots, k + 1}
(c_i + c_{i + 1} + \ldots + c_k + c_{k + 1}
- b_i - b_{i + 1} - \ldots - b_{k - 1} - b_k) \\
& =
k + 1 +
\max_{i = 0, 1, \ldots, k + 1}
(c_i + c_{i + 1} + \ldots + c_k
- b_i - b_{i + 1} - \ldots - b_{k - 1} - b_k) \\
& \leq
k + \max_{i = 0, 1, \ldots, k}
(c_i + c_{i + 1} + \ldots + c_k - b_i - b_{i + 1} - \ldots - b_{k - 1}) \\
& = \delta(V_0, \ldots, W_k)
\end{align*}
L'inégalité vaut parce que $c_k > 0$ et $b_k > 0$, ce qui entraîne en particulier
que $\delta(V_0, \ldots, W_k) \geq k + c_k \geq k + 1$.

\medskip\noindent
Preuve de (3). Étant donnés $y \in Y$ et $z \in \overline{\{y\}} \cap Z$,
on obtient le système
$$
V_0 = \overline{\{z\}} \subset W_0 = \overline{\{y\}}
$$
et $c_0 = \text{codim}(V_0, W_0) = \dim(\mathcal{O}_{\overline{\{y\}}, z})$
d'après Propriétés, lemme \ref{properties-lemma-codimension-local-ring}.
On voit ainsi que $\delta(V_0, W_0) = 0 + c_0 = c_0$, ce qui prouve
ce que l'on veut.

\medskip\noindent
Preuve de (4). Soit $\delta$ une fonction de dimension sur $Y$.
Soit $V_0 \subset W_0 \supset V_1 \subset W_1 \supset \ldots \subset W_k$
un système pour $y$. Soient $y'_i \in W_i$ et $y_i \in V_i$ les
points génériques ; ainsi, $y_0 \in Z$ et $y_k = y$. On voit alors que
$$
\delta(y_i) - \delta(y_{i - 1}) =
\delta(y'_{i - 1}) - \delta(y_{i - 1}) - \delta(y'_{i - 1}) + \delta(y_i) =
c_{i - 1} - b_{i - 1}
$$
Enfin, on a $\delta(y'_k) - \delta(y_{k - 1}) = c_k$.
On voit donc que
$$
\delta(y) - \delta(y_0) =
c_0 + \ldots + c_k - b_0 - \ldots - b_{k - 1}
$$
On en conclut que
$\delta(V_0, W_0, \ldots, W_k) \geq k + \delta(y) - \delta(y_0)$
ce qui prouve ce que l'on veut.

\medskip\noindent
Preuve de (5). La fonction $\delta(y) = \dim(\overline{\{y\}})$
est une fonction de dimension. Ainsi, $\delta(y) \leq \delta_Z(y)$ d'après la
partie (4). D'après la partie (3), on a $\delta_Z(y) \leq \delta(y)$,
et l'on a terminé.

\medskip\noindent
Preuve de (6). C'est clair, puisqu'il y a moins de systèmes à considérer
dans le calcul de $\delta^{Y'}_{Y' \cap Z}$.

\medskip\noindent
Preuve de (7). Soit
$V_0 \subset W_0 \supset V_1 \subset W_1 \supset \ldots \subset W_k$
un système pour $y$. Posons $W'_k = \overline{\{y'\}}$.
Comme $Y$ est caténaire, on voit que
$\text{codim}(V_k, W'_k) = \text{codim}(V_k, W_k) + 1$.
Il s'ensuit facilement que $\delta(V_0, \ldots, V_k, W'_k) \leq
\delta(V_0, \ldots, V_k, W_k) + 1$, ce qui prouve ce que l'on veut.

\medskip\noindent
Preuve de (8) : combinons (7) et (2).
\end{proof}

\begin{lemma}
\label{algebraization-lemma-change-distance-function}
Soit $Y$ un schéma noethérien universellement caténaire. Soit $Z \subset Y$
un sous-schéma fermé. Soit $f : Y' \to Y$ un morphisme de type fini
dont toutes les fibres sont de dimension $\leq e$. Posons $Z' = f^{-1}(Z)$.
Alors
$$
\delta_Z(y) \leq \delta_{Z'}(y') + e - \text{trdeg}_{\kappa(y)}(\kappa(y'))
$$
pour $y' \in Y'$ d'image $y \in Y$.
\end{lemma}

\begin{proof}
Si $\delta_{Z'}(y') = \infty$, il n'y a rien à démontrer.
Si $\delta_{Z'}(y') < \infty$, choisissons un système
de sous-schémas fermés intègres
$$
V'_0 \subset W'_0 \supset V'_1 \subset W'_1 \supset \ldots \subset W'_k
$$
de $Y'$, avec $V'_0 \subset Z'$ et $y'$ pour point générique de $W'_k$,
tel que $\delta_{Z'}(y') = \delta(V'_0, W'_0, \ldots, W'_k)$.
Notons
$$
V_0 \subset W_0 \supset V_1 \subset W_1 \supset \ldots \subset W_k
$$
les images schématiques des schémas ci-dessus dans $Y$. Remarquons
que $y$ est le point générique de $W_k$ et que $V_0 \subset Z$.
Pour chaque $i$, considérons le diagramme
$$
\xymatrix{
V'_i \ar[r] \ar[d] & W'_i \ar[d] & V'_{i + 1} \ar[l] \ar[d] \\
V_i \ar[r] & W_i & V_{i + 1} \ar[l]
}
$$
Notons $n_i$ la dimension relative de $V'_i/V_i$ et
$m_i$ la dimension relative de $W'_i/W_i$ ; plus précisément,
ce sont les degrés de transcendance des extensions correspondantes
des corps de fonctions. Posons
$c_i = \text{codim}(V_i, W_i)$,
$c'_i = \text{codim}(V'_i, W'_i)$,
$b_i = \text{codim}(V_{i + 1}, W_i)$, et
$b'_i = \text{codim}(V'_{i + 1}, W'_i)$.
D'après la formule de dimension, on a
$$
c_i = c'_i + n_i - m_i
\quad\text{et}\quad
b_i = b'_i + n_{i + 1} - m_i
$$
Voir Morphismes, lemme \ref{morphisms-lemma-dimension-formula}.
Ainsi, $c_i - b_i = c'_i - b'_i + n_i - n_{i + 1}$. On voit donc que
\begin{align*}
& c_i + c_{i + 1} + \ldots + c_k - b_i - b_{i + 1} - \ldots - b_{k - 1} \\
& =
c'_i + c'_{i + 1} + \ldots + c'_k - b'_i - b'_{i + 1} - \ldots - b'_{k - 1}
+ n_i - n_k + c_k - c'_k \\
& =
c'_i + c'_{i + 1} + \ldots + c'_k - b'_i - b'_{i + 1} - \ldots - b'_{k - 1}
+ n_i - m_k
\end{align*}
On voit donc que
\begin{align*}
\max_{i = 0, \ldots, k}
& (c_i + c_{i + 1} + \ldots + c_k - b_i - b_{i + 1} - \ldots - b_{k - 1}) \\
& =
\max_{i = 0, \ldots, k}
(c'_i + c'_{i + 1} + \ldots + c'_k - b'_i - b'_{i + 1} - \ldots - b'_{k - 1}
+ n_i - m_k) \\
& =
\max_{i = 0, \ldots, k}
(c'_i + c'_{i + 1} + \ldots + c'_k - b'_i - b'_{i + 1} - \ldots - b'_{k - 1}
+ n_i) - m_k \\
& \leq
\max_{i = 0, \ldots, k}
(c'_i + c'_{i + 1} + \ldots + c'_k - b'_i - b'_{i + 1} - \ldots - b'_{k - 1})
+ e - m_k
\end{align*}
Comme $m_k = \text{trdeg}_{\kappa(y)}(\kappa(y'))$, on conclut que
$$
\delta(V_0, W_0, \ldots, W_k) \leq
\delta(V'_0, W'_0, \ldots, W'_k) + e - \text{trdeg}_{\kappa(y)}(\kappa(y'))
$$
comme voulu.
\end{proof}

\begin{remark}
\label{algebraization-remark-discussion}
Soit $Y$ un schéma noethérien caténaire et soit $Z \subset Y$
une partie fermée. D'après le lemme \ref{algebraization-lemma-discussion}, on a
$$
\delta_Z(y) \leq \min
\left\{ k \middle|
\begin{matrix}
\text{ il existe des spécialisations dans }Y \\
y_0 \leftarrow y'_0 \rightarrow y_1 \leftarrow y'_1 \rightarrow \ldots
\leftarrow y'_{k - 1} \rightarrow y_k = y \\
\text{ avec }y_0 \in Z\text{ et }y_i' \leadsto y_i
\text{ immédiates}
\end{matrix}
\right\}
$$
Nous affirmons que si $Y$ est de type fini sur un corps,
alors l'égalité a lieu. Si nous avons un jour besoin de ce résultat, nous
énoncerons ici un résultat précis et le démontrerons.
Cependant, en général, si nous définissons $\delta_Z$
par le membre de droite de cette inégalité, nous ne
savons pas si le lemme \ref{algebraization-lemma-change-distance-function} reste valable.
\end{remark}

\begin{example}
\label{algebraization-example-distance}
Soit $k$ un corps et $Y = \mathbf{A}^n_k$. Notons
$\delta : Y \to \mathbf{Z}_{\geq 0}$ la fonction de dimension usuelle.
\begin{enumerate}
\item Si $Z = \{z\}$ pour un point fermé $z$, alors
\begin{enumerate}
\item $\delta_Z(y) = \delta(y)$ si $y \leadsto z$, et
\item $\delta_Z(y) = \delta(y) + 1$ si $y \not \leadsto z$.
\end{enumerate}
\item Si $Z$ est une sous-variété fermée et $W = \overline{\{y\}}$, alors
\begin{enumerate}
\item $\delta_Z(y) = 0$ si $W \subset Z$,
\item $\delta_Z(y) = \dim(W) - \dim(Z)$ si $Z$ est contenue dans $W$,
\item $\delta_Z(y) = 1$ si $\dim(W) \leq \dim(Z)$ et $W \not \subset Z$,
\item $\delta_Z(y) = \dim(W) - \dim(Z) + 1$ si $\dim(W) > \dim(Z)$
et $Z \not \subset W$.
\end{enumerate}
\end{enumerate}
Une généralisation du cas (1) s'obtient lorsque $Y$ est de type fini sur un corps
et que $Z = \{z\}$ est un point fermé. Alors $\delta_Z(y) = \delta(y) + t$,
où $t$ est la longueur minimale d'une chaîne de courbes reliant
$z$ à un point fermé de $\overline{\{y\}}$.
\end{example}







\section{Algébrisation des modules formels cohérents, III}
\label{algebraization-section-uniqueness}

\noindent
Nous poursuivons la discussion commencée aux sections
\ref{algebraization-section-algebraization-modules} et
\ref{algebraization-section-algebraization-modules-general}.
Nous utiliserons la fonction de distance de la section \ref{algebraization-section-distance}
pour formuler certaines conditions naturelles sur les
modules formels cohérents dans la situation \ref{algebraization-situation-algebraize}.

\medskip\noindent
Dans la situation \ref{algebraization-situation-algebraize}, étant donné un point $y \in U \cap Y$,
nous pouvons considérer la complétion $I$-adique
$$
\mathcal{O}_{X, y}^\wedge = \lim \mathcal{O}_{X, y}/I^n\mathcal{O}_{X, y}
$$
C'est un anneau local noethérien complet pour la topologie
$I\mathcal{O}_{X, y}^\wedge$-adique, d'idéal maximal $\mathfrak m_y^\wedge$ ; voir
Algèbre, section \ref{algebra-section-completion-noetherian}.
Soit $(\mathcal{F}_n)$ un objet de $\textit{Coh}(U, I\mathcal{O}_U)$.
Définissons le « germe » de $(\mathcal{F}_n)$ en $y$ par la formule
$$
\mathcal{F}_y^\wedge = \lim \mathcal{F}_{n, y}
$$
C'est un module de type fini sur $\mathcal{O}_{X, y}^\wedge$. Voir
Algèbre, lemmes \ref{algebra-lemma-limit-complete} et
\ref{algebra-lemma-finite-over-complete-ring}.

\begin{definition}
\label{algebraization-definition-s-d-inequalities}
Dans la situation \ref{algebraization-situation-algebraize}, soit $(\mathcal{F}_n)$ un objet
de $\textit{Coh}(U, I\mathcal{O}_U)$. Soient $a, b$ des entiers.
Soit $\delta^Y_Z$ comme dans (\ref{algebraization-equation-delta-Z}).
Nous disons que
{\it $(\mathcal{F}_n)$ satisfait les inégalités $(a, b)$} si, pour
$y \in U \cap Y$ et un idéal premier $\mathfrak p \subset \mathcal{O}_{X, y}^\wedge$
tel que $\mathfrak p \not \in V(I\mathcal{O}_{X, y}^\wedge)$,
\begin{enumerate}
\item si $V(\mathfrak p) \cap V(I\mathcal{O}_{X, y}^\wedge) \not =
\{\mathfrak m_y^\wedge\}$, alors
$$
\text{depth}((\mathcal{F}^\wedge_y)_\mathfrak p) + \delta^Y_Z(y) \geq a
\quad\text{ou}\quad
\text{depth}((\mathcal{F}^\wedge_y)_\mathfrak p) +
\dim(\mathcal{O}_{X, y}^\wedge/\mathfrak p) + \delta^Y_Z(y) > b
$$
\item si $V(\mathfrak p) \cap V(I\mathcal{O}_{X, y}^\wedge) =
\{\mathfrak m_y^\wedge\}$, alors
$$
\text{depth}((\mathcal{F}^\wedge_y)_\mathfrak p) + \delta^Y_Z(y) > a
$$
\end{enumerate}
Nous disons que {\it $(\mathcal{F}_n)$ satisfait les inégalités strictes $(a, b)$}
si, pour $y \in U \cap Y$ et un idéal premier
$\mathfrak p \subset \mathcal{O}_{X, y}^\wedge$ tel que
$\mathfrak p \not \in V(I\mathcal{O}_{X, y}^\wedge)$,
on a
$$
\text{depth}((\mathcal{F}^\wedge_y)_\mathfrak p) + \delta^Y_Z(y) > a
\quad\text{ou}\quad
\text{depth}((\mathcal{F}^\wedge_y)_\mathfrak p) +
\dim(\mathcal{O}_{X, y}^\wedge/\mathfrak p) + \delta^Y_Z(y) > b
$$
\end{definition}

\noindent
Voici quelques observations élémentaires.

\begin{lemma}
\label{algebraization-lemma-elementary}
Dans la situation \ref{algebraization-situation-algebraize}, soit $(\mathcal{F}_n)$ un objet
de $\textit{Coh}(U, I\mathcal{O}_U)$. Soient $a, b$ des entiers.
\begin{enumerate}
\item Si $(\mathcal{F}_n)$ est annulé par une puissance de $I$, alors
$(\mathcal{F}_n)$ satisfait les inégalités $(a, b)$ pour tous $a, b$.
\item Si $(\mathcal{F}_n)$ satisfait les inégalités $(a + 1, b)$, alors
$(\mathcal{F}_n)$ satisfait les inégalités strictes $(a, b)$.
\end{enumerate}
Si $\text{cd}(A, I) \leq d$ et si $A$ possède un complexe dualisant, alors
\begin{enumerate}
\item[(3)] $(\mathcal{F}_n)$ satisfait les inégalités $(s, s + d)$
si et seulement si, pour tout $y \in U \cap Y$, le sextuplet
$\mathcal{O}_{X, y}^\wedge, I\mathcal{O}_{X, y}^\wedge,
\{\mathfrak m_y^\wedge\}, \mathcal{F}_y^\wedge, s - \delta^Y_Z(y), d$
satisfait les hypothèses de la situation \ref{algebraization-situation-bootstrap}.
\item[(4)]
Si $(\mathcal{F}_n)$ satisfait les inégalités strictes $(s, s + d)$, alors
$(\mathcal{F}_n)$ satisfait les inégalités $(s, s + d)$.
\end{enumerate}
\end{lemma}

\begin{proof}
C'est immédiat, sauf pour la partie (4), qui découle du
lemme \ref{algebraization-lemma-bootstrap-bis-bis} et de la reformulation de (3).
\end{proof}

\begin{lemma}
\label{algebraization-lemma-explain-2-3-cd-1}
Dans la situation \ref{algebraization-situation-algebraize}, soit $(\mathcal{F}_n)$ un objet
de $\textit{Coh}(U, I\mathcal{O}_U)$. Si $\text{cd}(A, I) = 1$, alors
$\mathcal{F}$ satisfait les inégalités $(2, 3)$ si et seulement si
$$
\text{depth}((\mathcal{F}^\wedge_y)_\mathfrak p) +
\dim(\mathcal{O}_{X, y}^\wedge/\mathfrak p) + \delta^Y_Z(y) > 3
$$
pour tout $y \in U \cap Y$ et tout $\mathfrak p \subset \mathcal{O}_{X, y}^\wedge$
tel que $\mathfrak p \not \in V(I\mathcal{O}_{X, y}^\wedge)$.
\end{lemma}

\begin{proof}
Observons que, pour un idéal premier $\mathfrak p \subset \mathcal{O}_{X, y}^\wedge$,
$\mathfrak p \not \in V(I\mathcal{O}_{X, y}^\wedge)$,
nous avons $V(\mathfrak p) \cap V(I\mathcal{O}_{X, y}^\wedge) =
\{\mathfrak m_y^\wedge\}
\Leftrightarrow \dim(\mathcal{O}_{X, y}^\wedge/\mathfrak p) = 1$
car $\text{cd}(A, I) = 1$.
Voir Cohomologie locale, lemmes
\ref{local-cohomology-lemma-cd-change-rings} et
\ref{local-cohomology-lemma-cd-bound-dim-local}.
Considérons les trois nombres
$\alpha = \text{depth}((\mathcal{F}^\wedge_y)_\mathfrak p) \geq 0$,
$\beta = \dim(\mathcal{O}_{X, y}^\wedge/\mathfrak p) \geq 1$, et
$\gamma = \delta^Y_Z(y) \geq 1$.
La définition \ref{algebraization-definition-s-d-inequalities} exige alors
\begin{enumerate}
\item si $\beta > 1$, alors
$\alpha + \gamma \geq 2$ ou $\alpha + \beta + \gamma > 3$, et
\item si $\beta = 1$, alors $\alpha + \gamma > 2$.
\end{enumerate}
Il est immédiat que cela équivaut à
$\alpha + \beta + \gamma > 3$.
\end{proof}

\noindent
Dans la suite de cette section, que nous suggérons au lecteur de sauter en première
lecture, nous montrerons que, lorsque $A$ est $I$-adiquement complet,
la catégorie des objets $(\mathcal{F}_n)$ de $\textit{Coh}(U, I\mathcal{O}_U)$
qui s'étendent à $X$ et satisfont aux
inégalités $(1, 1 + \text{cd}(A, I))$ strictes
est équivalente à une sous-catégorie pleine de la catégorie des
$\mathcal{O}_U$-modules cohérents.

\begin{lemma}
\label{algebraization-lemma-sanity}
Dans la situation \ref{algebraization-situation-algebraize}, soit $\mathcal{F}$ un
$\mathcal{O}_U$-module cohérent et soit $d \geq 1$. Supposons que
\begin{enumerate}
\item $A$ soit $I$-adiquement complet, possède un complexe dualisant, et que
$\text{cd}(A, I) \leq d$,
\item le complété $\mathcal{F}^\wedge$ de $\mathcal{F}$
satisfasse aux inégalités $(1, 1 + d)$ strictes.
\end{enumerate}
Soit $x \in X$ un point. Posons $W = \overline{\{x\}}$.
Si $W \cap Y$ a une composante irréductible contenue dans $Z$
et une qui ne l'est pas, alors $\text{depth}(\mathcal{F}_x) \geq 1$.
\end{lemma}

\begin{proof}
Soit $W \cap Y = W_1 \cup \ldots \cup W_n$ sa décomposition en
composantes irréductibles. Par hypothèse, après renumérotation, nous pouvons trouver
$0 < m < n$ tel que $W_1, \ldots, W_m  \subset Z$ et
$W_{m + 1}, \ldots, W_n \not \subset Z$. Nous en concluons que
$$
W \cap Y \setminus
\left((W_1 \cup \ldots \cup W_m) \cap (W_{m + 1} \cup \ldots \cup W_n)\right)
$$
n'est pas connexe. Par le lemme \ref{algebraization-lemma-connected}, nous pouvons trouver
$1 \leq i \leq m < j \leq n$ et
$z \in W_i \cap W_j$ tels que $\dim(\mathcal{O}_{W, z}) \leq d + 1$.
Choisissons une spécialisation immédiate $y \leadsto z$ avec
$y \in W_j$, $y \not \in Z$ ; l'existence de $y$ résulte de
Propriétés, lemme \ref{properties-lemma-complement-closed-point-Jacobson}.
Observons que $\delta^Y_Z(y) = 1$ et $\dim(\mathcal{O}_{W, y}) \leq d$.
Soit $\mathfrak p \subset \mathcal{O}_{X, y}$ l'idéal premier correspondant à $x$.
Soit $\mathfrak p' \subset \mathcal{O}_{X, y}^\wedge$ un idéal premier minimal
au-dessus de $\mathfrak p\mathcal{O}_{X, y}^\wedge$. Alors
$$
\text{depth}(\mathcal{F}_x) =
\text{depth}((\mathcal{F}^\wedge_y)_{\mathfrak p'})
\quad\text{et}\quad
\dim(\mathcal{O}_{W, y}) = \dim(\mathcal{O}_{X, y}^\wedge/\mathfrak p')
$$
Voir Algèbre, lemme \ref{algebra-lemma-apply-grothendieck-module}, et
Cohomologie locale, lemme \ref{local-cohomology-lemma-change-completion}.
La conclusion se lit alors directement sur les inégalités données.
\end{proof}

\begin{lemma}
\label{algebraization-lemma-recover}
Dans la situation \ref{algebraization-situation-algebraize}, soit $\mathcal{F}$ un
$\mathcal{O}_U$-module cohérent et soit $d \geq 1$. Supposons que
\begin{enumerate}
\item $A$ soit $I$-adiquement complet, possède un complexe dualisant, et que
$\text{cd}(A, I) \leq d$,
\item le complété $\mathcal{F}^\wedge$ de $\mathcal{F}$
satisfasse aux inégalités $(1, 1+ d)$ strictes, et
\item pour $x \in U$ tel que $\overline{\{x\}} \cap Y \subset Z$,
nous ayons $\text{depth}(\mathcal{F}_x) \geq 2$.
\end{enumerate}
Alors $H^0(U, \mathcal{F}) \to \lim H^0(U, \mathcal{F}/I^n\mathcal{F})$
est un isomorphisme.
\end{lemma}

\begin{proof}
Nous allons le démontrer en vérifiant que le lemme \ref{algebraization-lemma-application-H0} s'applique.
Soit donc $x \in \text{Ass}(\mathcal{F})$ avec $x \not \in Y$.
Posons $W = \overline{\{x\}}$.
La condition (3) montre que $W \cap Y \not \subset Z$.
Le lemme \ref{algebraization-lemma-sanity} montre qu'aucune composante
irréductible de $W \cap Y$ n'est contenue dans $Z$.
Ainsi, si $z \in W \cap Z$, il existe une
spécialisation immédiate $y \leadsto z$, avec $y \in W \cap Y$ et $y \not \in Z$.
Pour l'existence de $y$, utiliser
Propriétés, lemme \ref{properties-lemma-complement-closed-point-Jacobson}.
Alors $\delta^Y_Z(y) = 1$ et l'hypothèse
implique que $\dim(\mathcal{O}_{W, y}) > d$.
Par conséquent $\dim(\mathcal{O}_{W, z}) > 1 + d$, ce qui conclut.
\end{proof}

\begin{lemma}
\label{algebraization-lemma-fully-faithful-inequalities}
Dans la situation \ref{algebraization-situation-algebraize}, soit $\mathcal{F}$ un
$\mathcal{O}_U$-module cohérent et soit $d \geq 1$. Supposons que
\begin{enumerate}
\item $A$ soit $I$-adiquement complet, possède un complexe dualisant, et que
$\text{cd}(A, I) \leq d$,
\item le complété $\mathcal{F}^\wedge$ de $\mathcal{F}$
satisfasse aux inégalités $(1, 1 + d)$ strictes, et
\item pour $x \in U$ tel que $\overline{\{x\}} \cap Y \subset Z$,
nous ayons $\text{depth}(\mathcal{F}_x) \geq 2$.
\end{enumerate}
Alors l'application
$$
\Hom_U(\mathcal{G}, \mathcal{F})
\longrightarrow
\Hom_{\textit{Coh}(U, I\mathcal{O}_U)}(\mathcal{G}^\wedge, \mathcal{F}^\wedge)
$$
est bijective pour tout $\mathcal{O}_U$-module cohérent $\mathcal{G}$. 
\end{lemma}

\begin{proof}
Posons $\mathcal{H} = \SheafHom_{\mathcal{O}_U}(\mathcal{G}, \mathcal{F})$.
À l'aide de Cohomologie des schémas, lemme
\ref{coherent-lemma-hom-into-depth}, ou de
Compléments sur l'algèbre, lemme \ref{more-algebra-lemma-hom-into-depth},
nous voyons que le complété de $\mathcal{H}$
satisfait aux inégalités $(1, 1 + d)$ strictes et que, pour
$x \in U$ tel que $\overline{\{x\}} \cap Y \subset Z$,
nous avons $\text{depth}(\mathcal{H}_x) \geq 2$. Nous omettons les détails.
Par le lemme \ref{algebraization-lemma-recover}, nous avons donc
$$
\Hom_U(\mathcal{G}, \mathcal{F}) =
H^0(U, \mathcal{H}) =
\lim H^0(U, \mathcal{H}/\mathcal{I}^n\mathcal{H}) =
\Mor_{\textit{Coh}(U, I\mathcal{O}_U)}
(\mathcal{G}^\wedge, \mathcal{F}^\wedge)
$$
Voir Cohomologie des schémas, lemme \ref{coherent-lemma-completion-internal-hom},
pour la dernière égalité.
\end{proof}

\begin{lemma}
\label{algebraization-lemma-construct-unique}
Dans la situation \ref{algebraization-situation-algebraize}, soit $(\mathcal{F}_n)$ un
objet de $\textit{Coh}(U, I\mathcal{O}_U)$ et soit $d \geq 1$. Supposons que
\begin{enumerate}
\item $A$ soit $I$-adiquement complet, possède un complexe dualisant, et que
$\text{cd}(A, I) \leq d$,
\item $(\mathcal{F}_n)$ soit le complété d'un $\mathcal{O}_U$-module cohérent,
\item $(\mathcal{F}_n)$ satisfasse aux inégalités $(1, 1 + d)$ strictes.
\end{enumerate}
Alors il existe un unique $\mathcal{O}_U$-module cohérent $\mathcal{F}$
dont le complété est $(\mathcal{F}_n)$ et tel que, pour
$x \in U$ avec $\overline{\{x\}} \cap Y \subset Z$,
nous ayons $\text{depth}(\mathcal{F}_x) \geq 2$.
\end{lemma}

\begin{proof}
Choisissons un $\mathcal{O}_U$-module cohérent $\mathcal{F}$ dont le
complété est $(\mathcal{F}_n)$. Posons
$T = \{x \in U \mid \overline{\{x\}} \cap Y \subset Z\}$.
Nous construirons $\mathcal{F}$ en appliquant Cohomologie locale,
lemme \ref{local-cohomology-lemma-make-S2-along-T-simple},
à $\mathcal{F}$ et $T$.
L'unicité résultera alors de la propriété d'unicité des morphismes
du lemme \ref{algebraization-lemma-fully-faithful-inequalities}.

\medskip\noindent
Puisque $T$ est stable par spécialisation dans $U$, il suffit de
vérifier ce qui suit. Si $x' \leadsto x$ est une
spécialisation immédiate de points de $U$ avec $x \in T$
et $x' \not \in T$, alors $\text{depth}(\mathcal{F}_{x'}) \geq 1$.
Posons $W = \overline{\{x\}}$ et $W' = \overline{\{x'\}}$.
Comme $x' \not \in T$, l'ensemble $W' \cap Y$ n'est pas contenu dans $Z$.
Si $W' \cap Y$ contient une composante irréductible contenue dans $Z$,
le lemme \ref{algebraization-lemma-sanity} permet de conclure.
Sinon, choisissons une composante irréductible $W_1$ de $W \cap Y$ et
une composante irréductible $W'_1$ de $W' \cap Y$ telles que $W_1 \subset W'_1$.
Soit $z \in W_1$ le point générique. Soit $y \leadsto z$, avec $y \in W'_1$,
une spécialisation immédiate telle que $y \not \in Z$ ; l'existence de $y$
résulte de $W'_1 \not \subset Z$ (voir ci-dessus) et de
Propriétés, lemme \ref{properties-lemma-complement-closed-point-Jacobson}.
Nous avons alors $z \in Z$, $x \leadsto z$,
$x' \leadsto y \leadsto z$, $y \in Y \setminus Z$, et $\delta^Y_Z(y) = 1$.
Par Cohomologie locale, lemme \ref{local-cohomology-lemma-cd-bound-dim-local},
et puisque $z$
est un point générique de $W \cap Y$, nous avons
$\dim(\mathcal{O}_{W, z}) \leq d$.
Comme $x' \leadsto x$ est une spécialisation immédiate, nous avons
$\dim(\mathcal{O}_{W', z}) \leq d + 1$.
Puisque $y \not = z$, nous en concluons que
$\dim(\mathcal{O}_{W', y}) \leq d$.
Si $\text{depth}(\mathcal{F}_{x'}) = 0$, nous obtiendrions
une contradiction avec l'hypothèse (3) ; nous omettons les détails du passage
de $\mathcal{O}_{X, y}$ à son complété.
Cela achève la démonstration.
\end{proof}








\section{Algébrisation des modules formels cohérents, IV}
\label{algebraization-section-algebraization-modules-local}

\noindent
Dans cette section, nous démontrons deux versions plus fortes du
lemme \ref{algebraization-lemma-algebraization-principal-variant}
dans le cas local, à savoir les lemmes \ref{algebraization-lemma-algebraization-principal}
et \ref{algebraization-lemma-algebraization-principal-bis}.
Bien que ces lemmes soient rendus caducs par la proposition plus générale
\ref{algebraization-proposition-cd-1}, leurs démonstrations sont sensiblement
plus simples.

\begin{lemma}
\label{algebraization-lemma-algebraization-principal}
Dans la situation \ref{algebraization-situation-algebraize}, soit $(\mathcal{F}_n)$ un objet
de $\textit{Coh}(U, I\mathcal{O}_U)$. Supposons que
\begin{enumerate}
\item $A$ soit local et que $\mathfrak a = \mathfrak m$ soit l'idéal maximal,
\item $A$ possède un complexe dualisant,
\item $I = (f)$ soit un idéal principal engendré par un élément non diviseur de zéro $f \in \mathfrak m$,
\item $\mathcal{F}_n$ soit un
$\mathcal{O}_U/f^n\mathcal{O}_U$-module localement libre de type fini,
\item si $\mathfrak p \in V(f) \setminus \{\mathfrak m\}$, alors
$\text{depth}((A/f)_\mathfrak p) + \dim(A/\mathfrak p) > 1$, et
\item si $\mathfrak p \not \in V(f)$ et
$V(\mathfrak p) \cap V(f) \not = \{\mathfrak m\}$, alors
$\text{depth}(A_\mathfrak p) + \dim(A/\mathfrak p) > 3$.
\end{enumerate}
Alors $(\mathcal{F}_n)$ s'étend canoniquement à $X$. En particulier, si $A$
est complet, alors $(\mathcal{F}_n)$ est le complété d'un
$\mathcal{O}_U$-module cohérent.
\end{lemma}

\begin{proof}
Nous allons le démontrer en vérifiant les hypothèses (a), (b) et (c) du
lemme \ref{algebraization-lemma-when-done}.

\medskip\noindent
Puisque $\mathcal{F}_n$ est localement libre sur $\mathcal{O}_U/f^n\mathcal{O}_U$,
nous avons des suites exactes courtes
$0 \to \mathcal{F}_n \to \mathcal{F}_{n + 1} \to \mathcal{F}_1 \to 0$
pour tout $n$. La condition (b) résulte donc de
Cohomologie, lemme \ref{cohomology-lemma-topology-I-adic-f}.

\medskip\noindent
Par récurrence sur $n$ et à l'aide des suites exactes courtes
$0 \to A/f^n \to A/f^{n + 1} \to A/f \to 0$, nous voyons que
les idéaux premiers associés de $A/f^nA$ coïncident avec les idéaux premiers
associés de $A/fA$. Comme les points associés de $\mathcal{F}_n$
correspondent aux idéaux premiers associés de $A/f^nA$ distincts de $\mathfrak m$
par l'hypothèse (3), nous en concluons que
$M_n = H^0(U, \mathcal{F}_n)$ est un $A$-module de type fini d'après (5) et
Cohomologie locale, proposition \ref{local-cohomology-proposition-kollar}.
L'hypothèse (c) est donc satisfaite.

\medskip\noindent
Pour achever la démonstration, il suffit de montrer qu'il existe $n > 1$
tel que l'image de
$$
H^1(U, \mathcal{F}_n) \longrightarrow H^1(U, \mathcal{F}_1)
$$
soit de longueur finie comme $A$-module. Cela impliquera en effet l'hypothèse (a)
par Cohomologie, lemme \ref{cohomology-lemma-ML-better}. L'image est indépendante
de $n$ pour $n$ assez grand d'après le lemme \ref{algebraization-lemma-ML-local}.
Soit $\omega_A^\bullet$ un complexe dualisant normalisé pour $A$.
D'après le théorème de dualité locale et la dualité de Matlis
(Complexes dualisants, lemme \ref{dualizing-lemma-special-case-local-duality}
et proposition \ref{dualizing-proposition-matlis}),
notre assertion équivaut à dire que l'image de
$$
\text{Ext}^{-2}_A(M_1, \omega_A^\bullet) \to
\text{Ext}^{-2}_A(M_n, \omega_A^\bullet)
$$
est de longueur finie pour $n \gg 1$. Les modules considérés sont des
$A$-modules de type fini à support dans $V(f)$. Il suffit donc de montrer que cette
application est nulle après localisation en un idéal premier $\mathfrak q$
contenant $f$ et distinct de $\mathfrak m$.
Soit $\omega_{A_\mathfrak q}^\bullet$ un complexe
dualisant normalisé sur $A_\mathfrak q$ et rappelons que
$\omega_{A_\mathfrak q}^\bullet =
(\omega_A^\bullet)_\mathfrak q[\dim(A/\mathfrak q)]$ d'après
Complexes dualisants, lemme \ref{dualizing-lemma-dimension-function}.
En utilisant la structure locale de $\mathcal{F}_n$ donnée en (4),
nous voyons qu'il suffit de montrer l'annulation de
$$
\text{Ext}^{-2 + \dim(A/\mathfrak q)}_{A_\mathfrak q}(
A_\mathfrak q/f, \omega_{A_\mathfrak q}^\bullet)
\to
\text{Ext}^{-2 + \dim(A/\mathfrak q)}_{A_\mathfrak q}(
A_\mathfrak q/f^n, \omega_{A_\mathfrak q}^\bullet)
$$
pour $n$ assez grand. Si $\dim(A/\mathfrak q) > 3$, cela résulte immédiatement de
Cohomologie locale, lemme \ref{local-cohomology-lemma-sitting-in-degrees}.
Dans les autres cas, nous utiliserons la suite exacte longue
$$
\ldots
\xrightarrow{f^n}
H^{-1}(\omega_{A_\mathfrak q}^\bullet)
\to
\text{Ext}^{-1}_{A_\mathfrak q}(
A_\mathfrak q/f^n, \omega_{A_\mathfrak q}^\bullet) \to
H^0(\omega_{A_\mathfrak q}^\bullet)
\xrightarrow{f^n}
H^0(\omega_{A_\mathfrak q}^\bullet)
\to
\text{Ext}^0_{A_\mathfrak q}(
A_\mathfrak q/f^n, \omega_{A_\mathfrak q}^\bullet) \to 0
$$
Si $\dim(A/\mathfrak q) = 2$, alors
$H^0(\omega_{A_\mathfrak q}^\bullet) = 0$
car $\text{depth}(A_\mathfrak q) \geq 1$, puisque
$f$ est un élément non diviseur de zéro.
La suite exacte longue montre donc que la condition s'écrit
$$
f^{n - 1} :
H^{-1}(\omega_{A_\mathfrak q}^\bullet)/f \to
H^{-1}(\omega_{A_\mathfrak q}^\bullet)/f^n
$$
est nulle. Or $H^{-1}(\omega^\bullet_\mathfrak q)$ est un module de type fini dont le support
est contenu dans les idéaux premiers $\mathfrak p \subset A_\mathfrak q$ tels que
$\text{depth}(A_\mathfrak p) + \dim((A/\mathfrak p)_\mathfrak q) \leq 1$.
Puisque $\dim((A/\mathfrak p)_\mathfrak q) = \dim(A/\mathfrak p) - 2$,
la condition (6) nous dit que ces idéaux premiers sont contenus dans $V(f)$.
D'où l'annulation voulue pour $n$ assez grand.
Enfin, si $\dim(A/\mathfrak q) = 1$, la condition (5), combinée
au fait que $f$ est un élément non diviseur de zéro,
assure que $A_\mathfrak q$ est de profondeur au moins $2$. Par conséquent,
$H^0(\omega_{A_\mathfrak q}^\bullet) =
H^{-1}(\omega_{A_\mathfrak q}^\bullet) = 0$,
et la suite exacte longue montre que l'assertion
équivaut à l'annulation de
$$
f^{n - 1} :
H^{-2}(\omega_{A_\mathfrak q}^\bullet)/f \to
H^{-2}(\omega_{A_\mathfrak q}^\bullet)/f^n
$$
Or $H^{-2}(\omega^\bullet_\mathfrak q)$ est un module de type fini dont le support
est contenu dans les idéaux premiers $\mathfrak p \subset A_\mathfrak q$
tels que $\text{depth}(A_\mathfrak p) + \dim((A/\mathfrak p)_\mathfrak q)
\leq 2$. D'après la condition (6), tous ces idéaux premiers sont contenus dans $V(f)$.
D'où l'annulation voulue pour $n$ assez grand.
\end{proof}

\begin{remark}
\label{algebraization-remark-interesting-case}
Soit $(A, \mathfrak m)$ un anneau local noethérien, intègre, normal, complet,
de dimension $\geq 4$, et soit $f \in \mathfrak m$ non nul.
Les hypothèses (1), (2), (3), (5) et (6) du
lemme \ref{algebraization-lemma-algebraization-principal}
sont alors satisfaites. Les fibrés vectoriels
sur la complétion formelle de $U$ le long de $U \cap V(f)$
peuvent donc être algébrisés. Dans le lemme \ref{algebraization-lemma-algebraization-principal-bis},
nous généraliserons ceci à des modules formels cohérents plus généraux ;
comparer également avec la remarque \ref{algebraization-remark-interesting-case-bis}.
\end{remark}

\begin{lemma}
\label{algebraization-lemma-helper-algebraize}
Dans la situation \ref{algebraization-situation-algebraize}, soit $(M_n)$ un système projectif de
$A$-modules comme dans le lemme \ref{algebraization-lemma-system-of-modules}, et soit
$(\mathcal{F}_n)$ l'objet correspondant de
$\textit{Coh}(U, I\mathcal{O}_U)$. Soient $d \geq \text{cd}(A, I)$
et $s \geq 0$ des entiers.
Avec les notations ci-dessus, supposons que
\begin{enumerate}
\item $A$ soit local, d'idéal maximal $\mathfrak m = \mathfrak a$,
\item $A$ possède un complexe dualisant, et
\item $(\mathcal{F}_n)$ satisfasse aux inégalités $(s, s + d)$
(définition \ref{algebraization-definition-s-d-inequalities}).
\end{enumerate}
Soit $E$ une enveloppe injective du corps résiduel de $A$. Alors, pour $i \leq s$,
il existe un $A$-module de type fini $N$, annulé par une puissance
de $I$, et, pour $n \gg 0$, des applications compatibles
$$
H^i_\mathfrak m(M_n) \to \Hom_A(N, E)
$$
dont les conoyaux sont des $A$-modules de longueur finie et dont les noyaux $K_n$
forment un système projectif tel que $\Im(K_{n''} \to K_{n'})$ soit de longueur
finie pour $n'' \gg n' \gg 0$.
\end{lemma}

\begin{proof}
Soit $\omega_A^\bullet$ un complexe dualisant normalisé. Alors
$\delta^Y_Z = \delta$ est la fonction de dimension associée à
ce complexe dualisant.
Observons que $\Ext^{-i}_A(M_n, \omega_A^\bullet)$ est un $A$-module de type fini
annulé par $I^n$. Fixons $0 \leq i \leq s$.
Nous trouverons ci-dessous $n_1 > n_0 > 0$ tels qu'en posant
$$
N = \Im(\Ext^{-i}_A(M_{n_0}, \omega_A^\bullet) \to
\Ext^{-i}_A(M_{n_1}, \omega_A^\bullet))
$$
les noyaux des applications
$$
N \to \Ext^{-i}_A(M_n, \omega_A^\bullet),\quad n \geq n_1
$$
soient des $A$-modules de longueur finie et que les conoyaux $Q_n$ forment un
système tel que $\Im(Q_{n'} \to Q_{n''})$ soit de longueur finie
pour $n'' \gg n' \gg n_1$. Cela revient à dire que
le système $\{\Ext^{-i}_A(M_n, \omega_A^\bullet)\}_{n \geq 1}$
est essentiellement constant dans la catégorie quotient de la catégorie des
$A$-modules de type fini par la sous-catégorie de Serre des $A$-modules de longueur finie.
D'après le théorème de dualité locale
(Complexes dualisants, lemme \ref{dualizing-lemma-special-case-local-duality})
et la dualité de Matlis
(Complexes dualisants, proposition \ref{dualizing-proposition-matlis}),
nous en concluons qu'il existe des applications
$$
H^i_\mathfrak m(M_n) \to \Hom_A(N, E),\quad n \geq n_1
$$
comme dans l'énoncé du lemme.

\medskip\noindent
Choisissons $f \in \mathfrak m$. Soit $B = A_f^\wedge$ le complété $I$-adique
de la localisation $A_f$. Rappelons que
$\omega_{A_f}^\bullet = \omega_A^\bullet \otimes_A A_f$
et $\omega_B^\bullet = \omega_A^\bullet \otimes_A B$ sont des complexes
dualisants (Complexes dualisants, lemme \ref{dualizing-lemma-dualizing-localize}
et \ref{dualizing-lemma-completion-henselization-dualizing}).
Soit $M$ le $B$-module de type fini $\lim M_{n, f}$ (comparer avec la
discussion dans Cohomologie des schémas, lemme
\ref{coherent-lemma-inverse-systems-affine}). Alors
$$
\Ext^{-i}_A(M_n, \omega_A^\bullet)_f =
\Ext^{-i}_{A_f}(M_{n, f}, \omega_{A_f}^\bullet) =
\Ext^{-i}_B(M/I^n M, \omega_B^\bullet)
$$
Comme $\mathfrak m$ peut être engendré par un nombre fini d'éléments $f \in \mathfrak m$,
il suffit de montrer que, pour chacun de ces $f$, le système
$$
\{\Ext^{-i}_B(M/I^n M, \omega_B^\bullet)\}_{n \geq 1}
$$
est essentiellement constant. Nous omettons quelques détails.

\medskip\noindent
Soit $\mathfrak q \subset IB$ un idéal premier. Alors $\mathfrak q$ correspond
à un point $y \in U \cap Y$. Observons que
$\delta(\mathfrak q) = \dim(\overline{\{y\}})$
est aussi la valeur de la fonction de dimension associée à $\omega_B^\bullet$
(nous omettons les détails ; utiliser que $\omega_B^\bullet$ s'obtient de
$\omega_A^\bullet$ par extension des scalaires à $B$). L'hypothèse
(3) garantit, via le lemme \ref{algebraization-lemma-elementary},
que le lemme \ref{algebraization-lemma-algebraize-local-cohomology-bis}
s'applique à
$B_\mathfrak q, IB_\mathfrak q, \mathfrak qB_\mathfrak q, M_\mathfrak q$
en remplaçant $s$ par $s - \delta(y)$. Nous obtenons
$$
H^{i - \delta(\mathfrak q)}_{\mathfrak qB_\mathfrak q}(M_\mathfrak q) =
\lim H^{i - \delta(\mathfrak q)}_{\mathfrak qB_\mathfrak q}(
(M/I^nM)_\mathfrak q)
$$
et ce module est annulé par une puissance de $I$.
Le lemme \ref{algebraization-lemma-terrific} montre que les systèmes projectifs
$H^{i - \delta(\mathfrak q)}_{\mathfrak qB_\mathfrak q}((M/I^nM)_\mathfrak q)$
sont essentiellement constants, de valeur
$H^{i - \delta(\mathfrak q)}_{\mathfrak qB_\mathfrak q}(M_\mathfrak q)$.
Comme $(\omega_B^\bullet)_\mathfrak q[-\delta(\mathfrak q)]$ est un complexe
dualisant normalisé sur $B_\mathfrak q$, le théorème de dualité locale
montre que le système
$$
\Ext^{-i}_B(M/I^n M, \omega_B^\bullet)_\mathfrak q
$$
est essentiellement constant, de valeur
$\Ext^{-i}_B(M, \omega_B^\bullet)_\mathfrak q$.

\medskip\noindent
Pour achever la démonstration, passons au global comme dans la démonstration du
lemme \ref{algebraization-lemma-bootstrap} ; l'argument est ici plus simple,
car nous connaissons déjà la valeur de notre système. Considérons les applications
$$
\alpha_n :
\Ext^{-i}_B(M/I^n M, \omega_B^\bullet)
\longrightarrow
\Ext^{-i}_B(M, \omega_B^\bullet)
$$
lorsque $n$ varie. D'après ce qui précède, pour tout $\mathfrak q$, il existe
un $n$ tel que $\alpha_n$ soit surjective après localisation en $\mathfrak q$.
Comme $B$ est noethérien et $\Ext^{-i}_B(M, \omega_B^\bullet)$
un module de type fini, il existe un $n$ tel que $\alpha_n$ soit surjective.
Pour tout $n$ tel que $\alpha_n$ soit surjective, étant donné un idéal premier
$\mathfrak q \in V(IB)$, il existe un $n' > n$ tel que
$\Ker(\alpha_n)$ s'envoie sur zéro dans $\Ext^{-i}(M/I^{n'}M, \omega_B^\bullet)$,
au moins après localisation en $\mathfrak q$.
Comme $\Ker(\alpha_n)$ est un $A$-module de type fini et que les supports des
sections sont quasi-compacts, il existe un $n'$ tel que
$\Ker(\alpha_n)$ s'envoie sur zéro dans $\Ext^{-i}(M/I^{n'}M, \omega_B^\bullet)$.
On voit ainsi que $\Ext^{-i}(M/I^n M, \omega_B^\bullet)$
est essentiellement constant, de valeur $\Ext^{-i}(M, \omega_B^\bullet)$.
Ceci achève la démonstration.
\end{proof}

\noindent
Voici une version plus générale du lemme \ref{algebraization-lemma-algebraization-principal}.

\begin{lemma}
\label{algebraization-lemma-algebraization-principal-bis}
Dans la situation \ref{algebraization-situation-algebraize}, soit $(\mathcal{F}_n)$ un objet
de $\textit{Coh}(U, I\mathcal{O}_U)$. Supposons que
\begin{enumerate}
\item $A$ soit local et $\mathfrak a = \mathfrak m$ soit l'idéal maximal,
\item $A$ possède un complexe dualisant,
\item $I = (f)$ soit un idéal principal,
\item $(\mathcal{F}_n)$ vérifie les inégalités $(2, 3)$.
\end{enumerate}
Alors $(\mathcal{F}_n)$ s'étend à $X$. En particulier, si $A$ est
complet pour la topologie $I$-adique, alors $(\mathcal{F}_n)$ est le complété
d'un $\mathcal{O}_U$-module cohérent. 
\end{lemma}

\begin{proof}
Rappelons que $\textit{Coh}(U, I\mathcal{O}_U)$ est une catégorie abélienne, voir
Cohomologie des schémas, lemme \ref{coherent-lemma-inverse-systems-abelian}.
Sur les ouverts affines de $U$, l'objet $(\mathcal{F}_n)$
correspond à un module de type fini sur un anneau noethérien
(Cohomologie des schémas, lemme \ref{coherent-lemma-inverse-systems-affine}).
Ainsi, les noyaux des applications $f^N : (\mathcal{F}_n) \to (\mathcal{F}_n)$
se stabilisent pour $N$ assez grand. D'après les
lemmes \ref{algebraization-lemma-map-kernel-cokernel-on-closed} et
\ref{algebraization-lemma-essential-image-completion},
pour démontrer le lemme,
nous pouvons remplacer $(\mathcal{F}_n)$ par l'image d'une telle application.
Nous pouvons donc supposer que $f$ est injectif sur $(\mathcal{F}_n)$.
Après ce remplacement, les conditions équivalentes du
lemme \ref{algebraization-lemma-equivalent-f-good} sont vérifiées pour le système projectif
$(\mathcal{F}_n)$ sur $U$. Nous l'utiliserons sans autre mention
dans la suite de la démonstration.

\medskip\noindent
Nous allons vérifier les hypothèses (a), (b) et (c) du
lemme \ref{algebraization-lemma-when-done}. L'hypothèse (b) résulte de
Cohomologie, lemme \ref{cohomology-lemma-topology-I-adic-f}.

\medskip\noindent
Choisissons un système projectif de modules $\{M_n\}$ comme dans le
lemme \ref{algebraization-lemma-system-of-modules}.
Nous pouvons supposer que $H^0_\mathfrak m(M_n) = 0$ en remplaçant, si nécessaire, $M_n$ par
$M_n/H^0_\mathfrak m(M_n)$. Nous obtenons alors des suites
exactes courtes
$$
0 \to M_n \to H^0(U, \mathcal{F}_n) \to H^1_\mathfrak m(M_n) \to 0
$$
pour tout $n$. Soit $E$ une enveloppe injective du corps résiduel de $A$.
D'après le lemme \ref{algebraization-lemma-helper-algebraize} et notre hypothèse (4),
nous pouvons choisir un entier $m \geq 0$, des $A$-modules de type fini
$N_1$ et $N_2$ annulés par $f^c$ pour un certain $c \geq 0$, ainsi que des
systèmes compatibles d'applications
$$
H^i_\mathfrak m(M_n) \to \Hom_A(N_i, E), \quad i = 1, 2
$$
pour $n \geq m$,
ayant les propriétés énoncées dans le lemme.

\medskip\noindent
Nous savons que $M = \lim H^0(U, \mathcal{F}_n)$ est un $A$-module dont la
topologie de limite est la topologie $f$-adique. Ainsi, étant donné $n$, le module
$M/f^nM$ est un sous-quotient de $H^0(U, \mathcal{F}_N)$ pour un certain $N \gg n$.
Les informations obtenues ci-dessus montrent que
$f^cM/f^nM$ est un $A$-module de type fini. Comme $f$ est un élément non diviseur de zéro
dans $M$, nous en déduisons que $M/f^{n - c}M$ est un $A$-module de type fini.
On voit ainsi que l'hypothèse (c) du lemme \ref{algebraization-lemma-when-done} est vérifiée.

\medskip\noindent
Étudions ensuite le module
$$
Ob = \lim H^1(U, \mathcal{F}_n) = \lim H^2_\mathfrak m(M_n)
$$
Pour $n \geq m$, soit $K_n$ le noyau de l'application
$H^2_\mathfrak m(M_n) \to \Hom_A(N_2, E)$.
Posons $K = \lim K_n$. Nous obtenons une suite exacte
$$
0 \to K \to Ob \to \Hom_A(N_2, E)
$$
D'après ce qui précède, la topologie de limite sur $Ob = \lim H^2_\mathfrak m(M_n)$
est la topologie $f$-adique. Comme $N_2$ est annulé par $f^c$,
nous en déduisons qu'il en va de même pour la topologie de limite sur $K = \lim K_n$.
Ainsi, $K/fK$ est un sous-quotient de $K_n$ pour $n \gg 1$.
Cependant, puisque $\{K_n\}$ est pro-isomorphe à un système projectif de
$A$-modules de longueur finie (d'après la conclusion du
lemme \ref{algebraization-lemma-helper-algebraize}),
nous en déduisons que $K/fK$ est un sous-quotient d'un
$A$-module de longueur finie. Il s'ensuit que $K$ est un $A$-module de type fini, voir
Algèbre, lemme \ref{algebra-lemma-finite-over-complete-ring}.
(En fait, nous voyons même que $\dim(\text{Supp}(K)) = 1$, mais
nous n'en aurons pas besoin.)

\medskip\noindent
Étant donné $n \geq 1$, considérons l'application de bord
$$
\delta_n :
H^0(U, \mathcal{F}_n)
\longrightarrow
\lim_N H^1(U, f^n\mathcal{F}_N) \xrightarrow{f^{-n}} Ob
$$
(la seconde application est un isomorphisme)
provenant des suites exactes courtes
$$
0 \to f^n\mathcal{F}_N \to \mathcal{F}_N \to \mathcal{F}_n \to 0
$$
Pour tout $n$, posons
$$
P_n = \Im(H^0(U, \mathcal{F}_{n + m}) \to H^0(U, \mathcal{F}_n))
$$
où $m$ est comme ci-dessus. Observons que $\{P_n\}$ est un système
projectif et que l'application $f : \mathcal{F}_n \to \mathcal{F}_{n + 1}$
sur les sections globales envoie $P_n$ dans $P_{n + 1}$.
Si $p \in P_n$, alors $\delta_n(p) \in K \subset Ob$,
car $\delta_n(p)$ s'envoie sur zéro dans
$H^1(U, f^n\mathcal{F}_{n + m}) = H^2_\mathfrak m(M_m)$
et la composée de $\delta_n$ et de $Ob \to \Hom_A(N_2, E)$
se factorise par $H^2_\mathfrak m(M_m)$ en vertu de notre choix de $m$.
Par conséquent,
$$
\bigoplus\nolimits_{n \geq 0} \Im(P_n \to Ob)
$$
est un $A[T]$-module gradué de type fini, où $T$ agit par multiplication par $f$.
En effet, c'est un sous-module gradué de $K[T]$ et $K$ est de type fini sur $A$.
En raisonnant comme dans la démonstration de
Cohomologie, lemme \ref{cohomology-lemma-ML-general}\footnote{Choisir
des générateurs homogènes de la forme $\delta_{n_j}(p_j)$ du module
affiché. Si $k = \max(n_j)$, on voit alors que pour $n \geq k$
et tout $p \in P_n$, il existe des $a_j \in A$ tels que
$p - \sum a_j f^{n - n_j} p_j$ appartienne au noyau de $\delta_n$
et donc à l'image de $P_{n'}$ pour tout $n' \geq n$.
Ainsi, $\Im(P_n \to P_{n - k}) = \Im(P_{n'} \to P_{n - k})$
pour tout $n' \geq n$.},
nous voyons que le système projectif $\{P_n\}$ satisfait à la condition ML.
Comme $\{P_n\}$ est pro-isomorphe à $\{H^0(U, \mathcal{F}_n)\}$,
nous en déduisons que $\{H^0(U, \mathcal{F}_n)\}$ satisfait à la condition ML.
Ainsi, l'hypothèse (a) du lemme \ref{algebraization-lemma-when-done}
est vérifiée et la démonstration est achevée.
\end{proof}

\noindent
Nous pouvons expliciter la condition du
lemme \ref{algebraization-lemma-algebraization-principal-bis} comme suit.

\begin{lemma}
\label{algebraization-lemma-unwinding-conditions}
Dans la situation \ref{algebraization-situation-algebraize}, soit $(\mathcal{F}_n)$ un objet
de $\textit{Coh}(U, I\mathcal{O}_U)$. Supposons que
\begin{enumerate}
\item $A$ soit local d'idéal maximal $\mathfrak a = \mathfrak m$,
\item $\text{cd}(A, I) = 1$.
\end{enumerate}
Alors $(\mathcal{F}_n)$ vérifie les inégalités $(2, 3)$ si et seulement
si, pour tout $y \in U \cap Y$ tel que $\dim(\{y\}) = 1$ et tout idéal premier
$\mathfrak p \subset \mathcal{O}_{X, y}^\wedge$,
$\mathfrak p \not \in V(I\mathcal{O}_{X, y}^\wedge)$, on a
$$
\text{depth}((\mathcal{F}_y^\wedge)_\mathfrak p) +
\dim(\mathcal{O}_{X, y}^\wedge/\mathfrak p) > 2
$$
\end{lemma}

\begin{proof}
Nous utiliserons le lemme \ref{algebraization-lemma-explain-2-3-cd-1}
sans autre mention. En particulier, nous voyons que la condition est nécessaire.
Réciproquement, supposons que la condition soit vraie.
Notons que $\delta^Y_Z(y) = \dim(\overline{\{y\}})$ d'après le
lemme \ref{algebraization-lemma-discussion}. Notons $\delta$ cette fonction.
Soit $y \in U \cap Y$. Si $\delta(y) > 2$, alors l'inégalité
du lemme \ref{algebraization-lemma-explain-2-3-cd-1} est vérifiée.
Supposons enfin que $\delta(y) = 2$. Il faut montrer que
$$
\text{depth}((\mathcal{F}_y^\wedge)_\mathfrak p) +
\dim(\mathcal{O}_{X, y}^\wedge/\mathfrak p) > 1
$$
Choisissons une spécialisation $y \leadsto y'$ telle que $\delta(y') = 1$. Il existe alors
un homomorphisme d'anneaux $\mathcal{O}_{X, y'}^\wedge \to \mathcal{O}_{X, y}^\wedge$
qui identifie le but au complété de la localisation
de $\mathcal{O}_{X, y'}^\wedge$ en un idéal premier $\mathfrak q$
tel que $\dim(\mathcal{O}_{X, y'}^\wedge/\mathfrak q) = 1$.
De plus, nous obtenons alors
$$
\mathcal{F}_y^\wedge =
\mathcal{F}_{y'}^\wedge
\otimes_{\mathcal{O}_{X, y'}^\wedge}
\mathcal{O}_{X, y}^\wedge
$$
Soit $\mathfrak p' \subset \mathcal{O}_{X, y'}^\wedge$ l'image
de $\mathfrak p$.
D'après Cohomologie locale, lemme \ref{local-cohomology-lemma-change-completion},
nous avons
\begin{align*}
\text{depth}((\mathcal{F}_y^\wedge)_\mathfrak p) +
\dim(\mathcal{O}_{X, y}^\wedge/\mathfrak p)
& =
\text{depth}((\mathcal{F}_{y'}^\wedge)_{\mathfrak p'}) +
\dim((\mathcal{O}_{X, y}^\wedge/\mathfrak p)_{\mathfrak p'}) \\
& =
\text{depth}((\mathcal{F}_{y'}^\wedge)_{\mathfrak p'}) +
\dim(\mathcal{O}_{X, y}^\wedge/\mathfrak p') - 1
\end{align*}
la dernière égalité résultant du fait que la spécialisation est immédiate.
Le lemme résulte donc de l'inégalité supposée pour $y', \mathfrak p'$.
\end{proof}

\begin{lemma}
\label{algebraization-lemma-unwinding-conditions-bis}
Dans la situation \ref{algebraization-situation-algebraize}, soit $(\mathcal{F}_n)$ un objet
de $\textit{Coh}(U, I\mathcal{O}_U)$. Supposons que
\begin{enumerate}
\item $A$ soit local d'idéal maximal $\mathfrak a = \mathfrak m$,
\item $A$ possède un complexe dualisant,
\item $\text{cd}(A, I) = 1$,
\item pour $y \in U \cap Y$, le module $\mathcal{F}_y^\wedge$
soit localement libre de type fini hors de $V(I\mathcal{O}_{X, y}^\wedge)$,
par exemple si $\mathcal{F}_n$ est un
$\mathcal{O}_U/I^n\mathcal{O}_U$-module localement libre de type fini, et
\item l'une des conditions suivantes soit vérifiée :
\begin{enumerate}
\item $A_f$ soit $(S_2)$ et toute composante irréductible de $X$
non contenue dans $Y$ soit de dimension $\geq 4$, ou
\item si $\mathfrak p \not \in V(f)$ et
$V(\mathfrak p) \cap V(f) \not = \{\mathfrak m\}$, alors
$\text{depth}(A_\mathfrak p) + \dim(A/\mathfrak p) > 3$.
\end{enumerate}
\end{enumerate}
Alors $(\mathcal{F}_n)$ vérifie les inégalités $(2, 3)$.
\end{lemma}

\begin{proof}
Nous allons utiliser le critère du lemme \ref{algebraization-lemma-unwinding-conditions}.
Soit $y \in U \cap Y$ tel que $\dim(\overline{\{y\}} = 1$ et soit
$\mathfrak p$ un idéal premier $\mathfrak p \subset \mathcal{O}_{X, y}^\wedge$ tel que
$\mathfrak p \not \in V(I\mathcal{O}_{X, y}^\wedge)$.
La condition (4) montre que
$\text{depth}((\mathcal{F}_y^\wedge)_\mathfrak p) =
\text{depth}((\mathcal{O}_{X, y}^\wedge)_\mathfrak p)$.
Il suffit donc de démontrer que
$$
\text{depth}((\mathcal{O}_{X, y}^\wedge)_\mathfrak p) +
\dim(\mathcal{O}_{X, y}^\wedge/\mathfrak p) > 2
$$
Soit $\mathfrak p_0 \subset A$ l'image de $\mathfrak p$.
Soit $\mathfrak q \subset A$ l'idéal premier correspondant à $y$.
D'après Cohomologie locale, lemme
\ref{local-cohomology-lemma-change-completion},
nous avons
\begin{align*}
\text{depth}((\mathcal{O}_{X, y}^\wedge)_\mathfrak p) +
\dim(\mathcal{O}_{X, y}^\wedge/\mathfrak p)
& =
\text{depth}(A_{\mathfrak p_0}) + \dim((A/\mathfrak p_0)_\mathfrak q) \\
& =
\text{depth}(A_{\mathfrak p_0}) + \dim(A/\mathfrak p_0) - 1
\end{align*}
Si (5)(a) est vérifiée, cette quantité est
$$
\geq \min(2, \dim(A_{\mathfrak p_0})) + \dim(A/\mathfrak p_0) - 1
$$
Notons que, dans tous les cas, $\dim(A/\mathfrak p_0) \geq 2$. Par conséquent, si
le minimum vaut $2$, la conclusion est acquise. Sinon, nous obtenons
$$
\dim(A_{\mathfrak p_0}) + \dim(A/\mathfrak p_0) - 1 \geq 4 - 1
$$
car toute composante de $\Spec(A)$ passant par $\mathfrak p_0$
est de dimension $\geq 4$. Si (5)(b) est vérifiée, la conclusion est immédiate.
\end{proof}

\begin{remark}
\label{algebraization-remark-interesting-case-bis}
Soit $(A, \mathfrak m)$ un anneau local noethérien qui possède un
complexe dualisant et qui est complet pour la topologie définie par $f \in \mathfrak m$.
Soit $(\mathcal{F}_n)$ un objet de $\textit{Coh}(U, f\mathcal{O}_U)$,
où $U$ est le spectre épointé de $A$.
Posons $Y = V(f) \subset X = \Spec(A)$.
Supposons que, pour $y \in U \cap V(f)$ fermé dans $U$, c'est-à-dire tel que
$\dim(\overline{\{y\}}) = 1$, le
$\mathcal{O}_{X, y}^\wedge$-module $\mathcal{F}_y^\wedge$
vérifie les deux conditions suivantes :
\begin{enumerate}
\item $\mathcal{F}_y^\wedge[1/f]$ soit $(S_2)$ comme
$\mathcal{O}_{X, y}^\wedge[1/f]$-module, et
\item pour $\mathfrak p \in \text{Ass}(\mathcal{F}_y^\wedge[1/f])$,
on ait $\dim(\mathcal{O}_{X, y}^\wedge/\mathfrak p) \geq 3$.
\end{enumerate}
Alors $(\mathcal{F}_n)$ est le complété d'un module cohérent sur $U$.
Ceci résulte des lemmes \ref{algebraization-lemma-algebraization-principal-bis}
et \ref{algebraization-lemma-unwinding-conditions}.
\end{remark}
































\section{Amélioration des modules formels cohérents}
\label{algebraization-section-improving-formal-modules}

\noindent
Soit $X$ un schéma noethérien. Soit $Y \subset X$ un sous-schéma fermé,
dont le faisceau quasi-cohérent d'idéaux est $\mathcal{I} \subset \mathcal{O}_X$.
Soit $(\mathcal{F}_n)$ un objet de $\textit{Coh}(X, \mathcal{I})$.
Dans cette section, nous construisons des applications
$(\mathcal{F}_n) \to (\mathcal{F}'_n)$
analogues aux applications construites dans
Cohomologie locale, section \ref{local-cohomology-section-improve},
pour les modules cohérents. Pour un point $y \in Y$, posons
$$
\mathcal{O}_{X, y}^\wedge = \lim \mathcal{O}_{X, y}/\mathcal{I}^n_y, \quad
\mathcal{I}_y^\wedge = \lim \mathcal{I}_y/\mathcal{I}^n_y
\quad\text{et}\quad
\mathfrak m_y^\wedge = \lim \mathfrak m_y/\mathcal{I}_y^n
$$
Alors $\mathcal{O}_{X, y}^\wedge$ est un anneau local noethérien
d'idéal maximal $\mathfrak m_y^\wedge$, complet pour la topologie définie par
$\mathcal{I}_y^\wedge = \mathcal{I}_y\mathcal{O}_{X, y}^\wedge$.
Posons aussi
$$
\mathcal{F}_y^\wedge = \lim \mathcal{F}_{n, y}
$$
Alors $\mathcal{F}_y^\wedge$ est un module de type fini sur
$\mathcal{O}_{X, y}^\wedge$ tel que
$\mathcal{F}_y^\wedge/(\mathcal{I}_y^\wedge)^n\mathcal{F}_y^\wedge =
\mathcal{F}_{n, y}$ pour tout $n$, voir Algèbre, lemmes
\ref{algebra-lemma-limit-complete} et
\ref{algebra-lemma-finite-over-complete-ring}.

\begin{lemma}
\label{algebraization-lemma-divide-torsion-formal-coherent-module}
Dans la situation ci-dessus, supposons que $X$ possède localement un complexe dualisant.
Soit $T \subset Y$ une partie stable par spécialisation.
Supposons que, pour $y \in T$ et pour un idéal premier non maximal
$\mathfrak p \subset \mathcal{O}_{X, y}^\wedge$ tel que
$V(\mathfrak p) \cap V(\mathcal{I}^\wedge_y) = \{\mathfrak m_y^\wedge\}$,
on ait
$$
\text{depth}_{(\mathcal{O}_{X, y})_\mathfrak p}
((\mathcal{F}^\wedge_y)_\mathfrak p) > 0
$$
Alors il existe une application canonique
$(\mathcal{F}_n) \to (\mathcal{F}_n')$
de systèmes projectifs de $\mathcal{O}_X$-modules cohérents
ayant les propriétés suivantes :
\begin{enumerate}
\item pour $y \in T$, on a $\text{depth}(\mathcal{F}'_{n, y}) \geq 1$,
\item $(\mathcal{F}'_n)$ est isomorphe, en tant que pro-système, à un objet
$(\mathcal{G}_n)$ de $\textit{Coh}(X, \mathcal{I})$,
\item le morphisme induit
$(\mathcal{F}_n) \to (\mathcal{G}_n)$ dans
$\textit{Coh}(X, \mathcal{I})$ est surjectif, de noyau
annulé par une puissance de $\mathcal{I}$.
\end{enumerate}
\end{lemma}

\begin{proof}
Pour tout $n$, soit $\mathcal{F}_n \to \mathcal{F}'_n$ la surjection
construite dans
Cohomologie locale, lemme \ref{local-cohomology-lemma-get-depth-1-along-Z}.
Comme il s'agit du quotient de $\mathcal{F}_n$ par le sous-faisceau
des sections à support dans $T$, nous voyons qu'il existe des applications canoniques
$\mathcal{F}'_{n + 1} \to \mathcal{F}'_n$ telles que l'on obtienne
une application $(\mathcal{F}_n) \to (\mathcal{F}_n')$
de systèmes projectifs de $\mathcal{O}_X$-modules cohérents.
La propriété (1) est vérifiée par construction.

\medskip\noindent
Pour démontrer les propriétés (2) et (3), nous pouvons supposer que $X = \Spec(A_0)$ est
affine et que $A_0$ possède un complexe dualisant. Soit $I_0 \subset A_0$ l'idéal
correspondant à $Y$. Soient $A, I$ les complétés $I_0$-adiques de
$A_0, I_0$. Pour un usage ultérieur, observons que $A$ possède un complexe dualisant
(Complexes dualisants, lemme \ref{dualizing-lemma-ubiquity-dualizing}).
Soit $M$ le $A$-module de type fini correspondant à $(\mathcal{F}_n)$, voir
Cohomologie des schémas, lemme \ref{coherent-lemma-inverse-systems-affine}.
Alors $\mathcal{F}_n$ correspond à $M_n = M/I^nM$. Rappelons que
$\mathcal{F}'_n$ correspond au quotient $M'_n = M_n / H^0_T(M_n)$,
voir Cohomologie locale, lemme \ref{local-cohomology-lemma-get-depth-1-along-Z}
et sa démonstration.

\medskip\noindent
Posons $s = 0$ et $d = \text{cd}(A, I)$.
Nous affirmons que $A, I, T, M, s, d$ vérifient les hypothèses (1), (3), (4), (6)
de la situation \ref{algebraization-situation-bootstrap}.
En effet, (1) et (3) résultent immédiatement de ce qui précède, (4) est
une condition vide puisque $s = 0$, et (6) est l'hypothèse
faite dans l'énoncé du lemme.

\medskip\noindent
D'après le théorème \ref{algebraization-theorem-final-bootstrap}, nous voyons que $\{H^0_T(M_n)\}$
est essentiellement constant, de valeur $H^0_T(M)$. Ainsi, la limite des
suites exactes courtes $0 \to H^0_T(M_n) \to M_n \to M'_n \to 0$
est la suite exacte courte
$$
0 \to H^0_T(M) \to M \to \lim M'_n \to 0
$$
Posons $M' = \lim M'_n = M/H^0_T(M)$. C'est un $A$-module de type fini.
Soit $(\mathcal{G}_n)$ l'objet de $\textit{Coh}(X, \mathcal{I})$
correspondant à $M'$. Pour achever la démonstration, il faut montrer que l'application
canonique $\{M'/I^nM'\} \to \{M'_n\}$ est un pro-isomorphisme.
Ceci équivaut à dire que
$\{H^0_T(M) + I^nM\} \to \{\Ker(M \to M'_n)\}$ est un
pro-isomorphisme. En quotientant par $I^nM$, il suffit de montrer que
$\{H^0_T(M)/H^0_T(M) \cap I^nM\} \to \{H^0_T(M_n)\}$
est un pro-isomorphisme. Comme $H^0_T(M)$ est annulé par une puissance
de $I$, le lemme d'Artin-Rees montre que $H^0_T(M) \cap I^nM = 0$
pour tout $n \gg 0$. Nous obtenons donc le pro-isomorphisme voulu,
puisque nous avons vu ci-dessus que $\{H^0_T(M_n)\}$
est essentiellement constant, de valeur $H^0_T(M)$.
\end{proof}

\begin{lemma}
\label{algebraization-lemma-improvement-formal-coherent-module-better}
Dans la situation ci-dessus, supposons que $X$ possède localement un complexe dualisant.
Soient $T' \subset T \subset Y$ des parties stables par spécialisation.
Soit $d \geq 0$ un entier. Supposons que
\begin{enumerate}
\item[(a)] localement sur des ouverts affines, on ait $X = \Spec(A_0)$ et $Y = V(I_0)$
et $\text{cd}(A_0, I_0) \leq d$,
\item[(b)] pour $y \in T$ et un idéal premier non maximal
$\mathfrak p \subset \mathcal{O}_{X, y}^\wedge$ tel que
$V(\mathfrak p) \cap V(\mathcal{I}_y^\wedge) = \{\mathfrak m_y^\wedge\}$,
on ait
$$
\text{depth}_{(\mathcal{O}_{X, y})_\mathfrak p}
((\mathcal{F}^\wedge_y)_\mathfrak p) > 0
$$
\item[(c)] pour $y \in T'$ et pour un idéal premier
$\mathfrak p \subset \mathcal{O}_{X, y}^\wedge$ tel que
$\mathfrak p \not \in V(\mathcal{I}_y^\wedge)$
et $V(\mathfrak p) \cap V(\mathcal{I}_y^\wedge) \not =
\{\mathfrak m_y^\wedge\}$, on ait
$$
\text{depth}_{(\mathcal{O}_{X, y})_\mathfrak p}
((\mathcal{F}^\wedge_y)_\mathfrak p) \geq 1
\quad\text{ou}\quad
\text{depth}_{(\mathcal{O}_{X, y})_\mathfrak p}
((\mathcal{F}^\wedge_y)_\mathfrak p) +
\dim(\mathcal{O}_{X, y}^\wedge/\mathfrak p) > 1 + d
$$
\item[(d)] pour $y \in T'$ et un idéal premier non maximal
$\mathfrak p \subset \mathcal{O}_{X, y}^\wedge$ tel que
$V(\mathfrak p) \cap V(\mathcal{I}_y^\wedge) = \{\mathfrak m_y^\wedge\}$,
on ait
$$
\text{depth}_{(\mathcal{O}_{X, y})_\mathfrak p}
((\mathcal{F}^\wedge_y)_\mathfrak p) > 1
$$
\item[(e)] si $y \leadsto y'$ est une spécialisation immédiate et
$y' \in T'$, alors $y \in T$.
\end{enumerate}
Alors il existe une application canonique $(\mathcal{F}_n) \to (\mathcal{F}_n'')$
de systèmes projectifs de $\mathcal{O}_X$-modules cohérents
ayant les propriétés suivantes :
\begin{enumerate}
\item pour $y \in T$, on a $\text{depth}(\mathcal{F}''_{n, y}) \geq 1$,
\item pour $y' \in T'$, on a $\text{depth}(\mathcal{F}''_{n, y'}) \geq 2$,
\item $(\mathcal{F}''_n)$ est isomorphe, en tant que pro-système, à un objet
$(\mathcal{H}_n)$ de $\textit{Coh}(X, \mathcal{I})$,
\item le morphisme induit $(\mathcal{F}_n) \to (\mathcal{H}_n)$ dans
$\textit{Coh}(X, \mathcal{I})$ a un noyau et un conoyau
annulés par une puissance de $\mathcal{I}$.
\end{enumerate}
\end{lemma}

\begin{proof}
Comme dans le lemme \ref{algebraization-lemma-divide-torsion-formal-coherent-module} et sa démonstration,
pour tout $n$, soit $\mathcal{F}_n \to \mathcal{F}'_n$ la surjection
construite dans
Cohomologie locale, lemme \ref{local-cohomology-lemma-get-depth-1-along-Z},
en utilisant $T$.
Ensuite, soit $\mathcal{F}'_n \to \mathcal{F}''_n$ l'injection
construite dans
Cohomologie locale, lemme \ref{local-cohomology-lemma-make-S2-along-T},
ainsi que dans sa démonstration. Les constructions montrent qu'il existe des applications canoniques
$\mathcal{F}''_{n + 1} \to \mathcal{F}''_n$ telles que l'on obtienne des
applications
$$
(\mathcal{F}_n) \longrightarrow (\mathcal{F}_n') \longrightarrow
(\mathcal{F}''_n)
$$
de systèmes projectifs de $\mathcal{O}_X$-modules cohérents.
Les propriétés (1) et (2) sont vérifiées par construction.

\medskip\noindent
Pour démontrer les propriétés (3) et (4), nous pouvons supposer que $X = \Spec(A_0)$ est
affine et que $A_0$ possède un complexe dualisant. Soit $I_0 \subset A_0$ l'idéal
correspondant à $Y$. Soient $A, I$ les complétés $I_0$-adiques de
$A_0, I_0$. Pour un usage ultérieur, observons que $A$ possède un complexe dualisant
(Complexes dualisants, lemme \ref{dualizing-lemma-ubiquity-dualizing}).
Soit $M$ le $A$-module de type fini correspondant à $(\mathcal{F}_n)$, voir
Cohomologie des schémas, lemme \ref{coherent-lemma-inverse-systems-affine}.
Alors $\mathcal{F}_n$ correspond à $M_n = M/I^nM$. Rappelons que
$\mathcal{F}'_n$ correspond au quotient $M'_n = M_n / H^0_T(M_n)$.
Rappelons aussi que $M' = \lim M'_n$ est le quotient de $M$ par
$H^0_T(M)$, que $H^0_T(M)$ est annulé par une puissance de $I$ et que
$\{M'_n\}$ et $\{M'/I^nM'\}$ sont isomorphes en tant que pro-systèmes.
Enfin, nous voyons que $\mathcal{F}''_n$ correspond à une extension
$$
0 \to M'_n \to M''_n \to H^1_{T'}(M'_n) \to 0
$$
voir la démonstration de
Cohomologie locale, lemme \ref{local-cohomology-lemma-make-S2-along-T}.

\medskip\noindent
Posons $s = 1$. Nous affirmons que $A, I, T', M', s, d$ vérifient les hypothèses
(1), (3), (4), (6) de la situation \ref{algebraization-situation-bootstrap}. En effet, (1) et (3)
sont immédiates, (4) est impliquée par (c) et (6) résulte de (d).
Nous omettons les démonstrations de (c) $\Rightarrow$ (4) et (d) $\Rightarrow$ (6).

\medskip\noindent
D'après le théorème \ref{algebraization-theorem-final-bootstrap}, nous voyons que $\{H^1_{T'}(M'/I^nM')\}$
est essentiellement constant, de valeur $H^1_{T'}(M')$. Nous en déduisons que
$\{H^1_{T'}(M'_n)\}$ est essentiellement constant, de valeur $H^1_{T'}(M')$.
Ainsi, la limite des suites exactes courtes affichées ci-dessus
est la suite exacte courte
$$
0 \to M' \to \lim M''_n \to H^1_{T'}(M') \to 0
$$
Posons $M'' = \lim M''_n$. Il résulte de
Cohomologie locale, proposition \ref{local-cohomology-proposition-finiteness},
que $H^1_{T'}(M')$, et donc $M''$, sont des $A$-modules de type fini (pour vérifier les
conditions, utiliser que $M'$ est de profondeur au moins $1$ aux idéaux premiers appartenant à $T - T'$).
Soit $(\mathcal{H}_n)$ l'objet de $\textit{Coh}(X, \mathcal{I})$
correspondant au $A$-module de type fini $M''$. Pour achever la démonstration, il faut
montrer que l'application canonique $\{M''/I^nM''\} \to \{M''_n\}$ est un
pro-isomorphisme. Puisque nous savons déjà que $\{M'/I^nM'\}$ est pro-isomorphe
à $\{M'_n\}$, le lecteur vérifiera (détails omis) que ceci équivaut à demander que
$\{H^1_{T'}(M')/I^nH^1_{T'}(M')\} \to \{H^1_{T'}(M'_n)\}$
soit un pro-isomorphisme. Comme $H^1_{T'}(M')/I^nH^1_{T'}(M') = H^1_{T'}(M')$
pour $n$ assez grand, cela résulte du fait que $\{H^1_{T'}(M'_n)\}$
est essentiellement constant, de valeur $H^1_{T'}(M')$, comme vu ci-dessus.
\end{proof}

\begin{lemma}
\label{algebraization-lemma-improvement-application}
Dans la situation \ref{algebraization-situation-algebraize}, supposons que $A$ possède
un complexe dualisant. Soit $d \geq \text{cd}(A, I)$. Soit $(\mathcal{F}_n)$
un objet de $\textit{Coh}(U, I\mathcal{O}_U)$. Supposons que
$(\mathcal{F}_n)$ vérifie les inégalités $(2, 2 + d)$, voir la
définition \ref{algebraization-definition-s-d-inequalities}.
Alors il existe une application canonique $(\mathcal{F}_n) \to (\mathcal{F}_n'')$
de systèmes projectifs de $\mathcal{O}_U$-modules cohérents
ayant les propriétés suivantes :
\begin{enumerate}
\item $\text{depth}(\mathcal{F}''_{n, y}) + \delta^Y_Z(y) \geq 3$
pour tout $y \in U \cap Y$,
\item $(\mathcal{F}''_n)$ est isomorphe, en tant que pro-système, à un objet
$(\mathcal{H}_n)$ de $\textit{Coh}(U, I\mathcal{O}_U)$,
\item le morphisme induit $(\mathcal{F}_n) \to (\mathcal{H}_n)$ dans
$\textit{Coh}(U, I\mathcal{O}_U)$ a un noyau et un conoyau
annulés par une puissance de $I$,
\item les modules $H^0(U, \mathcal{F}''_n)$ et $H^1(U, \mathcal{F}''_n)$
sont des $A$-modules de type fini pour tout $n$.
\end{enumerate}
\end{lemma}

\begin{proof}
L'existence et les propriétés (1), (2), (3) résultent du
lemme \ref{algebraization-lemma-improvement-formal-coherent-module-better} appliqué
à $U$, $U \cap Y$, $T = \{y \in U \cap Y : \delta^Y_Z(y) \leq 2\}$,
$T' = \{y \in U \cap Y : \delta^Y_Z(y) \leq 1\}$ et $(\mathcal{F}_n)$.
Le fait que les modules $H^0(U, \mathcal{F}''_n)$ et
$H^1(U, \mathcal{F}''_n)$ soient de type fini résulte de
Cohomologie locale, lemme \ref{local-cohomology-lemma-finiteness-Rjstar},
et des propriétés élémentaires de la fonction $\delta^Y_Z(-)$
démontrées dans le lemme \ref{algebraization-lemma-discussion}.
\end{proof}











\section{Algébrisation des modules formels cohérents, V}
\label{algebraization-section-algebraization-modules-conclusion}

\noindent
Dans cette section, nous démontrons nos résultats les plus généraux sur l'algébrisation
des modules formels cohérents. Nous les démontrons d'abord dans le cas où
l'idéal est de dimension cohomologique $1$. Nous les appliquons ensuite
à un éclatement pour démontrer un résultat plus général.

\begin{lemma}
\label{algebraization-lemma-cd-1-canonical}
Dans la situation \ref{algebraization-situation-algebraize}, soit $(\mathcal{F}_n)$
un objet de $\textit{Coh}(U, I\mathcal{O}_U)$. Supposons que
\begin{enumerate}
\item $A$ possède un complexe dualisant et $\text{cd}(A, I) = 1$,
\item $(\mathcal{F}_n)$ soit pro-isomorphe à un système projectif
$(\mathcal{F}_n'')$ de $\mathcal{O}_U$-modules cohérents tel que
$\text{depth}(\mathcal{F}''_{n, y}) + \delta^Y_Z(y) \geq 3$
pour tout $y \in U \cap Y$.
\end{enumerate}
Alors $(\mathcal{F}_n)$ s'étend canoniquement à $X$, voir la
définition \ref{algebraization-definition-canonically-algebraizable}.
\end{lemma}

\begin{proof}
Nous allons vérifier les hypothèses (a), (b) et (c) du lemme \ref{algebraization-lemma-when-done}.
Avant de commencer, signalons que les modules
$H^0(U, \mathcal{F}''_n)$ et $H^1(U, \mathcal{F}''_n)$
sont des $A$-modules de type fini pour tout $n$ d'après
Cohomologie locale, lemme \ref{local-cohomology-lemma-finiteness-Rjstar}.

\medskip\noindent
Observons que, pour tout $p \geq 0$,
la topologie de limite sur $\lim H^p(U, \mathcal{F}_n)$
est la topologie $I$-adique d'après le lemme \ref{algebraization-lemma-topology-I-adic}.
En particulier, l'hypothèse (b) est vérifiée.

\medskip\noindent
Nous savons que $M = \lim H^0(U, \mathcal{F}_n)$ est un $A$-module dont la
topologie de limite est la topologie $I$-adique. Ainsi, étant donné $n$, le module
$M/I^nM$ est un sous-quotient de $H^0(U, \mathcal{F}_N)$ pour un certain $N \gg n$.
Comme le système projectif $\{H^0(U, \mathcal{F}_N)\}$ est pro-isomorphe à un
système projectif de $A$-modules de type fini, à savoir $\{H^0(U, \mathcal{F}''_N)\}$,
nous en déduisons que $M/I^nM$ est de type fini. Il s'ensuit que $M$ est de type fini, voir
Algèbre, lemme \ref{algebra-lemma-finite-over-complete-ring}.
En particulier, l'hypothèse (c) est vérifiée.

\medskip\noindent
Pour tout $n \geq 0$, notons $Ob_n = \lim_N H^1(U, I^n\mathcal{F}_N)$.
Un cas particulier est $Ob = Ob_0 = \lim_N H^1(U, \mathcal{F}_N)$.
En raisonnant exactement comme au paragraphe précédent, nous trouvons que $Ob$
est un $A$-module de type fini. (En fait, nous savons aussi que $Ob/I Ob$ est annulé
par une puissance de $\mathfrak a$, mais cela semble assez difficile à utiliser.)

\medskip\noindent
Posons $\mathcal{F} = \lim \mathcal{F}_n$, choisissons des générateurs
$f_1, \ldots, f_r$ de $I$, choisissons $c \geq 1$ et choisissons
$\Phi_\mathcal{F}$ comme dans le lemme \ref{algebraization-lemma-cd-is-one-for-system}.
Nous utiliserons les résultats du lemme \ref{algebraization-lemma-properties-system}
sans autre mention. En particulier, pour tout $n \geq 1$, il existe des applications
$$
\delta_n :
H^0(U, \mathcal{F}_n)
\longrightarrow
H^1(U, I^n\mathcal{F})
\longrightarrow
Ob_n
$$
La première provient de la suite exacte courte
$0 \to I^n\mathcal{F} \to \mathcal{F} \to \mathcal{F}_n \to 0$
et la seconde de $I^n\mathcal{F} = \lim I^n\mathcal{F}_N$.
Nous utiliserons plus loin que, si $\delta_n(s) = 0$ pour $s \in H^0(U, \mathcal{F}_n)$,
alors, pour tout $n' \geq n$, il existe $s' \in H^0(U, \mathcal{F}_{n'})$
qui s'envoie sur $s$.
Observons qu'il existe des diagrammes commutatifs
$$
\xymatrix{
H^0(U, \mathcal{F}_{nc}) \ar[r] \ar[dd] &
H^1(U, I^{nc}\mathcal{F}) \ar[dd] \ar[rd]^{\Phi_\mathcal{F}} \\
& &
\bigoplus_{e_1 + \ldots + e_r = n}
H^1(U, \mathcal{F}) \cdot T_1^{e_1} \ldots T_r^{e_r} \ar[ld] \\
H^0(U, \mathcal{F}_n) \ar[r] &
H^1(U, I^n\mathcal{F})
}
$$
Nous en déduisons que l'application d'obstruction
$H^0(U, \mathcal{F}_n) \to Ob_n$
envoie l'image de
$H^0(U, \mathcal{F}_{nc}) \to H^0(U, \mathcal{F}_n)$
dans le sous-module
$$
Ob'_n =
\Im\left(
\bigoplus\nolimits_{e_1 + \ldots + e_r = n}
Ob \cdot T_1^{e_1} \ldots T_r^{e_r} \to Ob_n
\right)
$$
où, sur le facteur $Ob \cdot T_1^{e_1} \ldots T_r^{e_r}$,
nous utilisons l'application en cohomologie provenant des réductions modulo des
puissances de $I$ de l'application de multiplication
$f_1^{e_1} \ldots f_r^{e_r} : \mathcal{F} \to I^n\mathcal{F}$.
Par construction,
$$
\bigoplus\nolimits_{n \geq 0} Ob'_n
$$
est un module gradué de type fini sur l'algèbre de Rees $\bigoplus_{n \geq 0} I^n$.
Pour tout $n$, posons
$$
M_n = \{s \in H^0(U, \mathcal{F}_n) \mid \delta_n(s) \in Ob'_n\}
$$
Observons que $\{M_n\}$ est un système projectif et que l'application
$f_j : \mathcal{F}_n \to \mathcal{F}_{n + 1}$ sur les
sections globales envoie $M_n$ dans $M_{n + 1}$.
Par exactement le même argument que dans la démonstration de
Cohomologie, lemme \ref{cohomology-lemma-ML-general},
nous trouvons que $\{M_n\}$ satisfait à la condition ML. En effet, puisque l'algèbre de Rees
est noethérienne, nous pouvons choisir un nombre fini de générateurs homogènes
de la forme $\delta_{n_j}(z_j)$ avec $z_j \in M_{n_j}$ du sous-module gradué
$\bigoplus_{n \geq 0} \Im(M_n \to Ob'_n)$.
Alors, si $k = \max(n_j)$, nous voyons que pour $n \geq k$
et tout $z \in M_n$, il existe des $a_j \in I^{n - n_j}$ tels que
$z - \sum a_j z_j$ appartienne au noyau de $\delta_n$
et donc à l'image de $M_{n'}$ pour tout $n' \geq n$
(car l'annulation de $\delta_n$ signifie que nous pouvons
relever $z - \sum a_j z_j$ en un élément $z' \in H^0(U, \mathcal{F}_{n'c})$
pour tout $n' \ge n$, puis que l'image de $z'$ dans $H^0(U, \mathcal{F}_{n'})$
appartient à $M_{n'}$ d'après ce que nous avons démontré ci-dessus).
Ainsi, $\Im(M_n \to M_{n - k}) = \Im(M_{n'} \to M_{n - k})$
pour tout $n' \geq n$.

\medskip\noindent
Choisissons $n$. Par la propriété de Mittag-Leffler de $\{M_n\}$ établie ci-dessus,
nous pouvons trouver $n' \geq n$ tel que l'image de $M_{n'} \to M_n$
coïncide avec l'image de $M' \to M_n$. D'après ce qui précède, nous voyons que
l'image de $M' \to M_n$ contient l'image de
$H^0(U, \mathcal{F}_{n'c}) \to H^0(U, \mathcal{F}_n)$.
Ainsi, nous voyons que $\{M_n\}$ et $\{H^0(U, \mathcal{F}_n)\}$
sont pro-isomorphes. Par conséquent, $\{H^0(U, \mathcal{F}_n)\}$
satisfait à la condition ML, et nous concluons enfin que l'hypothèse (a) est vérifiée.
Ceci achève la démonstration.
\end{proof}

\begin{proposition}[Algébrisation en dimension cohomologique 1]
\label{algebraization-proposition-cd-1}
\begin{reference}
Le cas local de ce résultat est \cite[IV Corollaire 2.9]{MRaynaud-book}.
\end{reference}
Dans la situation \ref{algebraization-situation-algebraize}, soit $(\mathcal{F}_n)$
un objet de $\textit{Coh}(U, I\mathcal{O}_U)$. Supposons que
\begin{enumerate}
\item $A$ possède un complexe dualisant et $\text{cd}(A, I) = 1$,
\item $(\mathcal{F}_n)$ satisfasse aux inégalités $(2, 3)$, voir la
définition \ref{algebraization-definition-s-d-inequalities}.
\end{enumerate}
Alors $(\mathcal{F}_n)$ s'étend à $X$. En particulier, si $A$ est
complet pour la topologie $I$-adique, alors $(\mathcal{F}_n)$ est le complété
d'un $\mathcal{O}_U$-module cohérent.
\end{proposition}

\begin{proof}
D'après le lemme \ref{algebraization-lemma-map-kernel-cokernel-on-closed},
nous pouvons remplacer $(\mathcal{F}_n)$ par l'objet $(\mathcal{H}_n)$
de $\textit{Coh}(U, I\mathcal{O}_U)$ obtenu dans le
lemme \ref{algebraization-lemma-improvement-application}.
Nous pouvons donc supposer que $(\mathcal{F}_n)$ est pro-isomorphe
à un système projectif $(\mathcal{F}_n'')$ possédant les propriétés
mentionnées dans le lemme \ref{algebraization-lemma-improvement-application}.
Dans le lemme \ref{algebraization-lemma-cd-1-canonical}, nous avons démontré que
$(\mathcal{F}_n)$ s'étend canoniquement à $X$.
La dernière assertion résulte du lemme \ref{algebraization-lemma-canonically-algebraizable}.
\end{proof}

\begin{lemma}
\label{algebraization-lemma-blowup}
Dans la situation \ref{algebraization-situation-algebraize}, soit $(\mathcal{F}_n)$
un objet de $\textit{Coh}(U, I\mathcal{O}_U)$. Supposons que
\begin{enumerate}
\item $A$ possède un complexe dualisant,
\item toutes les fibres de l'éclatement $b : X' \to X$ de $I$
soient de dimension $\leq d - 1$,
\item l'une des conditions suivantes soit vérifiée :
\begin{enumerate}
\item $(\mathcal{F}_n)$ satisfait aux inégalités $(d + 1, d + 2)$
(définition \ref{algebraization-definition-s-d-inequalities}), ou
\item pour $y \in U \cap Y$ et un idéal premier
$\mathfrak p \subset \mathcal{O}_{X, y}^\wedge$ tel que
$\mathfrak p \not \in V(I\mathcal{O}_{X, y}^\wedge)$,
nous ayons
$$
\text{depth}((\mathcal{F}^\wedge_y)_\mathfrak p) +
\dim(\mathcal{O}_{X, y}^\wedge/\mathfrak p) + \delta^Y_Z(y) > d + 2
$$
\end{enumerate}
\end{enumerate}
Alors $(\mathcal{F}_n)$ s'étend à $X$.
\end{lemma}

\begin{proof}
Soit $Y' \subset X'$ le diviseur exceptionnel.
Soit $Z' \subset Y'$ l'image réciproque de $Z \subset Y$.
Alors $U' = X' \setminus Z'$ est l'image réciproque de $U$.
Avec $\delta^{Y'}_{Z'}$ comme dans (\ref{algebraization-equation-delta-Z}), posons
$$
T' = \{y' \in Y' \mid \delta^{Y'}_{Z'}(y') = 1\}
\subset
T = \{y' \in Y' \mid \delta^{Y'}_{Z'}(y') = 1\text{ ou }2\}
$$
D'après les parties (1) et (2) du lemme \ref{algebraization-lemma-discussion}, ces sous-ensembles sont
stables par spécialisation dans $U' \cap Y' = Y' \setminus Z'$. Considérons
l'objet $(b|_{U'}^*\mathcal{F}_n)$ de $\textit{Coh}(U', I\mathcal{O}_{U'})$,
voir Cohomologie des schémas, lemme \ref{coherent-lemma-inverse-systems-pullback}.
Pour $y' \in U' \cap Y'$, notons
$$
\mathcal{F}_{y'}^\wedge = \lim (b|_{U'}^*\mathcal{F}_n)_{y'}
$$
le « germe » de cette image réciproque en $y'$. Nous affirmons que les conditions
(a), (b), (c), (d) et (e) du
lemme \ref{algebraization-lemma-improvement-formal-coherent-module-better}
sont vérifiées pour l'objet $(b|_{U'}^*\mathcal{F}_n)$ sur $U'$, avec $d$
remplacé par $1$, et les sous-ensembles $T' \subset T \subset U' \cap Y'$.
La condition (a) est vérifiée parce que $Y'$ est un diviseur de Cartier effectif
et est donc localement défini par $1$ équation. La condition (e) est vérifiée
d'après la partie (7) du lemme \ref{algebraization-lemma-discussion}.
Pour démontrer (b), (c) et (d), quelques préparatifs sont nécessaires.

\medskip\noindent
Soit $y' \in U' \cap Y'$ et soit
$\mathfrak p' \subset \mathcal{O}_{X', y'}^\wedge$
un idéal premier n'appartenant pas à $V(I\mathcal{O}_{X', y'}^\wedge)$.
Notons $y = b(y') \in U \cap Y$. Choisissons $f \in I$ tel que
$y'$ appartienne au spectre de l'algèbre d'éclatement affine
$A[\frac{I}{f}]$, voir Diviseurs, lemme \ref{divisors-lemma-blowing-up-affine}.
Pour toute $A$-algèbre $B$, notons $B' = B[\frac{IB}{f}]$ l'algèbre d'éclatement
affine correspondante. Désignons la complétion $I$-adique par ${\ }^\wedge$.
Par notre choix de $f$, nous obtenons un homomorphisme d'anneaux
$(\mathcal{O}_{X, y}^\wedge)' \to \mathcal{O}_{X', y'}^\wedge$.
Si $\mathfrak q' \subset (\mathcal{O}_{X, y}^\wedge)'$
désigne l'image réciproque de $\mathfrak m_{y'}^\wedge$, alors
nous voyons que
$((\mathcal{O}_{X, y}^\wedge)'_{\mathfrak q'})^\wedge =
\mathcal{O}_{X', y'}^\wedge$.
Soit $\mathfrak p \subset \mathcal{O}_{X, y}^\wedge$ l'idéal premier
correspondant. Nous avons à ce stade un diagramme commutatif
$$
\xymatrix{
\mathcal{O}_{X, y}^\wedge \ar[d] \ar[r] &
(\mathcal{O}_{X, y}^\wedge)' \ar[d]_\alpha \ar[r] &
(\mathcal{O}_{X, y}^\wedge)'_{\mathfrak q'} \ar[d] \ar[r]_\beta &
\mathcal{O}_{X', y'}^\wedge \ar[d] \\
\mathcal{O}_{X, y}^\wedge/\mathfrak p \ar[r] &
(\mathcal{O}_{X, y}^\wedge/\mathfrak p)' \ar[r] &
(\mathcal{O}_{X, y}^\wedge/\mathfrak p)'_{\mathfrak q'} \ar[r]^\gamma &
((\mathcal{O}_{X, y}^\wedge/\mathfrak p)'_{\mathfrak q'})^\wedge \ar[d] \\
 & & &
\mathcal{O}_{X', y'}^\wedge/\mathfrak p'
}
$$
dont les flèches verticales sont surjectives. D'après
Compléments d'algèbre, lemme \ref{more-algebra-lemma-completion-dimension},
et la formule de dimension
(Algèbre, lemme \ref{algebra-lemma-dimension-formula}),
nous avons
$$
\dim(((\mathcal{O}_{X, y}^\wedge/\mathfrak p)'_{\mathfrak q'})^\wedge) =
\dim((\mathcal{O}_{X, y}^\wedge/\mathfrak p)'_{\mathfrak q'}) =
\dim(\mathcal{O}_{X, y}^\wedge/\mathfrak p)
- \text{trdeg}(\kappa(y')/\kappa(y))
$$
En suivant les définitions des images réciproques, des germes, des localisations
et des complétions, nous trouvons
$$
(\mathcal{F}_y^\wedge)_{\mathfrak p}
\otimes_{(\mathcal{O}_{X, y}^\wedge)_\mathfrak p}
(\mathcal{O}_{X', y'}^\wedge)_{\mathfrak p'}
=
(\mathcal{F}_{y'}^\wedge)_{\mathfrak p'}
$$
Détails omis. Les homomorphismes d'anneaux $\beta$ et $\gamma$ du diagramme
sont plats à fibres de Gorenstein (donc de Cohen--Macaulay), car ce sont des
complétions d'anneaux possédant un complexe dualisant. Voir
Complexes dualisants, lemmes
\ref{dualizing-lemma-formal-fibres-gorenstein} et
\ref{dualizing-lemma-dualizing-gorenstein-formal-fibres},
ainsi que la discussion de Compléments d'algèbre, section
\ref{more-algebra-section-properties-formal-fibres}.
Observons que $(\mathcal{O}_{X, y}^\wedge)_\mathfrak p =
(\mathcal{O}_{X, y}^\wedge)'_{\tilde{\mathfrak p}}$
où $\tilde{\mathfrak p}$ est le noyau de $\alpha$
dans le diagramme. D'autre part,
$(\mathcal{O}_{X, y}^\wedge)'_{\tilde{\mathfrak p}}
\to (\mathcal{O}_{X', y'}^\wedge)_{\mathfrak p'}$
est plat à fibres CM d'après ce qui précède. Par conséquent,
$(\mathcal{O}_{X, y}^\wedge)_\mathfrak p  \to
(\mathcal{O}_{X', y'}^\wedge)_{\mathfrak p'}$ est plat à fibres CM.
En appliquant Algèbre, lemme \ref{algebra-lemma-apply-grothendieck-module},
nous voyons que
$$
\text{depth}((\mathcal{F}_{y'}^\wedge)_{\mathfrak p'}) =
\text{depth}((\mathcal{F}_y^\wedge)_{\mathfrak p}) +
\dim(F_\mathfrak r)
$$
où $F$ est la fibre formelle générique de
$(\mathcal{O}_{X, y}^\wedge/\mathfrak p)'_{\mathfrak q'}$
et $\mathfrak r$ est l'idéal premier correspondant à $\mathfrak p'$.
Puisque $(\mathcal{O}_{X, y}^\wedge/\mathfrak p)'_{\mathfrak q'}$
est un domaine local universellement caténaire, son complété $I$-adique
est équidimensionnel et (universellement) caténaire d'après le théorème de Ratliff
(Compléments d'algèbre, proposition \ref{more-algebra-proposition-ratliff}).
Il s'ensuit alors que
$$
\dim(((\mathcal{O}_{X, y}^\wedge/\mathfrak p)'_{\mathfrak q'})^\wedge) =
\dim(F_\mathfrak r) + \dim(\mathcal{O}_{X', y'}^\wedge/\mathfrak p')
$$
En combinant ceci avec le lemme \ref{algebraization-lemma-change-distance-function},
nous obtenons
\begin{equation}
\label{algebraization-equation-one}
\begin{aligned}
&
\text{depth}((\mathcal{F}_{y'}^\wedge)_{\mathfrak p'}) +
\delta^{Y'}_{Z'}(y') \\
& =
\text{depth}((\mathcal{F}_y^\wedge)_{\mathfrak p}) +
\dim(F_\mathfrak r) + \delta^{Y'}_{Z'}(y') \\
& \geq
\text{depth}((\mathcal{F}_y^\wedge)_{\mathfrak p}) + \delta^Y_Z(y) +
\dim(F_\mathfrak r) + \text{trdeg}(\kappa(y')/\kappa(y)) - (d - 1) \\
& =
\text{depth}((\mathcal{F}_y^\wedge)_{\mathfrak p}) + \delta^Y_Z(y) - (d - 1)
+ \dim(\mathcal{O}_{X, y}^\wedge/\mathfrak p) -
\dim(\mathcal{O}_{X', y'}^\wedge/\mathfrak p')
\end{aligned}
\end{equation}
Gardons à l'esprit que
$\dim(\mathcal{O}_{X, y}^\wedge/\mathfrak p) \geq
\dim(\mathcal{O}_{X', y'}^\wedge/\mathfrak p')$. En réécrivant ceci, nous obtenons
\begin{equation}
\label{algebraization-equation-two}
\begin{aligned}
&
\text{depth}((\mathcal{F}_{y'}^\wedge)_{\mathfrak p'}) +
\dim(\mathcal{O}_{X', y'}^\wedge/\mathfrak p') +
\delta^{Y'}_{Z'}(y') \\
& \geq
\text{depth}((\mathcal{F}_y^\wedge)_{\mathfrak p}) +
\dim(\mathcal{O}_{X, y}^\wedge/\mathfrak p) +
\delta^Y_Z(y) - (d - 1)
\end{aligned}
\end{equation}
Cette inégalité nous permettra de vérifier les conditions restantes.

\medskip\noindent
Conditions (b) et (d) du
lemme \ref{algebraization-lemma-improvement-formal-coherent-module-better}. Supposons que
$V(\mathfrak p') \cap V(I\mathcal{O}_{X', y'}^\wedge) =
\{\mathfrak m_{y'}^\wedge\}$.
Cela implique que $\dim(\mathcal{O}_{X', y'}^\wedge/\mathfrak p') = 1$
parce que $Z'$ est un diviseur de Cartier effectif.
La conjonction de (b) et (d) équivaut à
$$
\text{depth}((\mathcal{F}_{y'}^\wedge)_{\mathfrak p'}) + \delta^{Y'}_{Z'}(y')
> 2
$$
Si $(\mathcal{F}_n)$ satisfait aux inégalités de (3)(b),
nous concluons immédiatement que ceci est vrai en appliquant (\ref{algebraization-equation-two}).
Si $(\mathcal{F}_n)$ satisfait à (3)(a), c'est-à-dire aux
inégalités $(d + 1, d + 2)$, alors nous voyons que, dans tous les cas,
$$
\text{depth}((\mathcal{F}_y^\wedge)_{\mathfrak p}) + \delta^Y_Z(y)
\geq d + 1
\quad\text{ou}\quad
\text{depth}((\mathcal{F}_y^\wedge)_{\mathfrak p}) +
\dim(\mathcal{O}_{X, y}^\wedge/\mathfrak p) +
\delta^Y_Z(y) > d + 2
$$
En examinant (\ref{algebraization-equation-one}) et (\ref{algebraization-equation-two}) ci-dessus, ceci donne
ce que nous voulons, sauf éventuellement si
$\dim(\mathcal{O}_{X, y}^\wedge/\mathfrak p) = 1$.
Cependant, si $\dim(\mathcal{O}_{X, y}^\wedge/\mathfrak p) = 1$, alors nous avons
$V(\mathfrak p) \cap V(I\mathcal{O}_{X, y}^\wedge) = \{\mathfrak m_y^\wedge\}$
et nous voyons qu'en fait
$$
\text{depth}((\mathcal{F}_y^\wedge)_{\mathfrak p}) + \delta^Y_Z(y) > d + 1
$$
car $(\mathcal{F}_n)$ satisfait aux inégalités $(d + 1, d + 2)$, et nous
concluons de nouveau.

\medskip\noindent
Condition (c) du
lemme \ref{algebraization-lemma-improvement-formal-coherent-module-better}. Supposons que
$V(\mathfrak p') \cap V(I\mathcal{O}_{X', y'}^\wedge) \not =
\{\mathfrak m_{y'}^\wedge\}$. Alors la condition (c) équivaut à
$$
\text{depth}((\mathcal{F}_{y'}^\wedge)_{\mathfrak p'}) + \delta^{Y'}_{Z'}(y')
\geq 2
\quad\text{ou}\quad
\text{depth}((\mathcal{F}_{y'}^\wedge)_{\mathfrak p'}) +
\dim(\mathcal{O}_{X', y'}^\wedge/\mathfrak p') +
\delta^{Y'}_{Z'}(y') > 3
$$
Si $(\mathcal{F}_n)$ satisfait aux inégalités de (3)(b),
nous voyons que la seconde des deux inégalités affichées est vérifiée
en appliquant (\ref{algebraization-equation-two}). Si $(\mathcal{F}_n)$ satisfait à (3)(a), c'est-à-dire
aux inégalités $(d + 1, d + 2)$, cela résulte immédiatement de
(\ref{algebraization-equation-one}) et (\ref{algebraization-equation-two}).
Ceci achève la démonstration de notre assertion.

\medskip\noindent
Choisissons $(b|_{U'}^*\mathcal{F}_n) \to (\mathcal{F}_n'')$
et $(\mathcal{H}_n)$ dans $\textit{Coh}(U', I\mathcal{O}_{U'})$
comme dans le lemme \ref{algebraization-lemma-improvement-formal-coherent-module-better}.
Pour tout ouvert affine $W \subset X'$, observons que
$\delta^{W \cap Y'}_{W \cap Z'}(y') \geq \delta^{Y'}_{Z'}(y')$ d'après la
partie (6) du lemme \ref{algebraization-lemma-discussion}. Nous voyons donc que
$(\mathcal{H}_n|_W)$ satisfait aux hypothèses du
lemme \ref{algebraization-lemma-cd-1-canonical}.
Ainsi, $(\mathcal{H}_n|_W)$ s'étend canoniquement à $W$.
Soit $(\mathcal{G}_{W, n})$ dans $\textit{Coh}(W, I\mathcal{O}_W)$
l'extension canonique comme dans le
lemme \ref{algebraization-lemma-canonically-algebraizable}.
D'après le lemme \ref{algebraization-lemma-canonically-extend-base-change},
nous voyons que, pour $W' \subset W$, il existe un unique isomorphisme
$$
(\mathcal{G}_{W, n}|_{W'}) \longrightarrow
(\mathcal{G}_{W', n})
$$
compatible avec les isomorphismes donnés
$(\mathcal{G}_{W, n}|_{W \cap U}) \cong (\mathcal{H}_n|_{W \cap U})$.
Nous en concluons qu'il existe un objet
$(\mathcal{G}_n)$ de $\textit{Coh}(X', I\mathcal{O}_{X'})$
dont la restriction à $U$ est isomorphe à $(\mathcal{H}_n)$.

\medskip\noindent
Si $A$ est complet pour la topologie $I$-adique, nous pouvons achever la démonstration comme suit.
D'après le théorème d'existence de Grothendieck
(Cohomologie des schémas, lemme \ref{coherent-lemma-existence-projective}),
nous voyons que $(\mathcal{G}_n)$ est le complété d'un
$\mathcal{O}_{X'}$-module cohérent. Puis, d'après
Cohomologie des schémas, lemme \ref{coherent-lemma-existence-easy},
nous voyons que $(b|_{U'}^*\mathcal{F}_n)$
est le complété d'un $\mathcal{O}_{U'}$-module cohérent
$\mathcal{F}'$.
D'après Cohomologie des schémas, lemme \ref{coherent-lemma-inverse-systems-push-pull},
nous voyons qu'il existe une application
$$
(\mathcal{F}_n) \longrightarrow ((b|_{U'})_*\mathcal{F}')^\wedge
$$
dont le noyau et le conoyau sont annulés par une puissance de $I$.
Enfin, nous concluons en appliquant le
lemme \ref{algebraization-lemma-map-kernel-cokernel-on-closed}.

\medskip\noindent
Si $A$ n'est pas complet, alors, avant de commencer la démonstration, nous pouvons remplacer $A$
par son complété, voir le lemme \ref{algebraization-lemma-algebraizable}.
Après complétion, les hypothèses restent vérifiées : cela est immédiat
pour la condition (3), résulte du
lemme \ref{dualizing-lemma-ubiquity-dualizing} de Complexes dualisants
pour la condition (1), et du
lemme \ref{divisors-lemma-flat-base-change-blowing-up} de Diviseurs
pour la condition (2).
Ainsi, le cas complet implique le cas général.
\end{proof}

\begin{proposition}[Algébrisation pour les idéaux ayant peu de générateurs]
\label{algebraization-proposition-d-generators}
Dans la situation \ref{algebraization-situation-algebraize}, soit $(\mathcal{F}_n)$
un objet de $\textit{Coh}(U, I\mathcal{O}_U)$. Supposons que
\begin{enumerate}
\item $A$ possède un complexe dualisant,
\item $V(I) = V(f_1, \ldots, f_d)$ pour un certain $d \geq 1$ et
$f_1, \ldots, f_d \in A$,
\item l'une des conditions suivantes soit vérifiée :
\begin{enumerate}
\item $(\mathcal{F}_n)$ satisfait aux inégalités $(d + 1, d + 2)$
(définition \ref{algebraization-definition-s-d-inequalities}), ou
\item pour $y \in U \cap Y$ et un idéal premier
$\mathfrak p \subset \mathcal{O}_{X, y}^\wedge$ tel que
$\mathfrak p \not \in V(I\mathcal{O}_{X, y}^\wedge)$,
nous ayons
$$
\text{depth}((\mathcal{F}^\wedge_y)_\mathfrak p) +
\dim(\mathcal{O}_{X, y}^\wedge/\mathfrak p) + \delta^Y_Z(y) > d + 2
$$
\end{enumerate}
\end{enumerate}
Alors $(\mathcal{F}_n)$ s'étend à $X$. En particulier, si $A$ est
complet pour la topologie $I$-adique, alors $(\mathcal{F}_n)$ est le complété
d'un $\mathcal{O}_U$-module cohérent. 
\end{proposition}

\begin{proof}
Nous pouvons supposer que $I = (f_1, \ldots, f_d)$, voir
Cohomologie des schémas, lemme
\ref{coherent-lemma-inverse-systems-ideals-equivalence}.
Nous voyons alors que
toutes les fibres de l'éclatement de $X$ le long de $I$ sont de dimension au plus $d - 1$.
Nous obtenons donc l'extension par le lemme \ref{algebraization-lemma-blowup}.
La dernière assertion résulte du lemme \ref{algebraization-lemma-essential-image-completion}.
\end{proof}

\noindent
Comparons le lemme suivant avec les
remarques \ref{algebraization-remark-interesting-case-variant},
\ref{algebraization-remark-interesting-case},
\ref{algebraization-remark-interesting-case-bis} et
\ref{algebraization-remark-interesting-case-ter}.

\begin{lemma}
\label{algebraization-lemma-interesting-case-final}
Dans la situation \ref{algebraization-situation-algebraize}, soit $(\mathcal{F}_n)$
un objet de $\textit{Coh}(U, I\mathcal{O}_U)$. Supposons que
\begin{enumerate}
\item $A$ soit un anneau local possédant un complexe dualisant,
\item toutes les composantes irréductibles de $X$ aient la même dimension,
\item le schéma $X \setminus Y$ soit de Cohen--Macaulay,
\item $I$ soit engendré par $d$ éléments,
\item $\dim(X) - \dim(Z) > d + 2$, et
\item pour $y \in U \cap Y$, le module $\mathcal{F}_y^\wedge$
soit localement libre de type fini hors de $V(I\mathcal{O}_{X, y}^\wedge)$,
par exemple si $\mathcal{F}_n$ est un
$\mathcal{O}_U/I^n\mathcal{O}_U$-module localement libre de type fini.
\end{enumerate}
Alors $(\mathcal{F}_n)$ s'étend à $X$. En particulier, si $A$ est complet pour la topologie $I$-adique,
alors $(\mathcal{F}_n)$ est le complété d'un
$\mathcal{O}_U$-module cohérent.
\end{lemma}

\begin{proof}
Nous allons montrer que les hypothèses (1), (2) et (3)(b) de la
proposition \ref{algebraization-proposition-d-generators} sont vérifiées.
C'est clair pour (1) et (2).

\medskip\noindent
Soit $y \in U \cap Y$ et soit $\mathfrak p$ un idéal premier
$\mathfrak p \subset \mathcal{O}_{X, y}^\wedge$ tel que
$\mathfrak p \not \in V(I\mathcal{O}_{X, y}^\wedge)$.
La dernière condition montre que
$\text{depth}((\mathcal{F}_y^\wedge)_\mathfrak p) =
\text{depth}((\mathcal{O}_{X, y}^\wedge)_\mathfrak p)$.
Puisque $X \setminus Y$ est de Cohen--Macaulay, nous voyons que
$(\mathcal{O}_{X, y}^\wedge)_\mathfrak p$ est de Cohen--Macaulay.
Ainsi, nous voyons que
\begin{align*}
& \text{depth}((\mathcal{F}^\wedge_y)_\mathfrak p) +
\dim(\mathcal{O}_{X, y}^\wedge/\mathfrak p) + \delta^Y_Z(y) \\
& =
\dim((\mathcal{O}_{X, y}^\wedge)_\mathfrak p) +
\dim(\mathcal{O}_{X, y}^\wedge/\mathfrak p) + \delta^Y_Z(y) \\
& =
\dim(\mathcal{O}_{X, y}^\wedge) + \delta^Y_Z(y)
\end{align*}
La dernière égalité vaut parce que $\mathcal{O}_{X, y}$ est équidimensionnel
d'après la deuxième condition.
Soit $\delta(y) = \dim(\overline{\{y\}})$. C'est une fonction de dimension
puisque $A$ est un anneau local caténaire.
D'après le lemme \ref{algebraization-lemma-discussion},
nous avons $\delta^Y_Z(y) \geq \delta(y) - \dim(Z)$. Puisque $X$ est
équidimensionnel, nous obtenons
$$
\dim(\mathcal{O}_{X, y}^\wedge) + \delta^Y_Z(y)
\geq \dim(\mathcal{O}_{X, y}^\wedge) + \delta(y) - \dim(Z)
= \dim(X) - \dim(Z)
$$
Nous obtenons donc l'inégalité voulue, ce qui conclut.
\end{proof}

\begin{remark}
\label{algebraization-remark-question}
Nous ne savons ni démontrer ni réfuter l'analogue de la
proposition \ref{algebraization-proposition-d-generators}
où l'hypothèse selon laquelle $I$ possède $d$ générateurs
est remplacée par l'hypothèse $\text{cd}(A, I) \leq d$.
Si vous connaissez une démonstration ou un contre-exemple, veuillez écrire à
\href{mailto:stacks.project@gmail.com}{stacks.project@gmail.com}.
Une autre question évidente est de savoir dans quelle mesure les conditions de la
proposition \ref{algebraization-proposition-d-generators}
sont nécessaires.
\end{remark}




\section{Algébrisation des modules formels cohérents, VI}
\label{algebraization-section-algebraization-modules-yet-more}

\noindent
Dans cette section, nous ajoutons quelques cas plus faciles à démontrer.

\begin{proposition}
\label{algebraization-proposition-algebraization-regular-sequence}
Dans la situation \ref{algebraization-situation-algebraize}, soit
$(\mathcal{F}_n)$ un objet de $\textit{Coh}(U, I\mathcal{O}_U)$.
Supposons que
\begin{enumerate}
\item il existe $f_1, \ldots, f_d \in I$ tels que,
pour $y \in U \cap Y$, l'idéal $I\mathcal{O}_{X, y}$
soit engendré par $f_1, \ldots, f_d$ et que
$f_1, \ldots, f_d$ forment une suite $\mathcal{F}_y^\wedge$-régulière,
\item $H^0(U, \mathcal{F}_1)$ et $H^1(U, \mathcal{F}_1)$
soient des $A$-modules de type fini.
\end{enumerate}
Alors $(\mathcal{F}_n)$ s'étend canoniquement à $X$. En particulier, si $A$
est complet, alors $(\mathcal{F}_n)$ est le complété d'un
$\mathcal{O}_U$-module cohérent.
\end{proposition}

\begin{proof}
Nous allons le démontrer en vérifiant les hypothèses (a), (b) et (c) du
lemme \ref{algebraization-lemma-when-done}.
Pour tout $n$, nous avons une suite exacte courte
$$
0 \to I^n\mathcal{F}_{n + 1} \to \mathcal{F}_{n + 1} \to \mathcal{F}_n \to 0
$$
Puisque $f_1, \ldots, f_d$ forment une suite régulière (et donc
quasi-régulière, voir Algèbre, lemme \ref{algebra-lemma-regular-quasi-regular})
sur chacun des « germes » $\mathcal{F}_y^\wedge$, et puisque nous avons
$I\mathcal{F}_n = (f_1, \ldots, f_d)\mathcal{F}_n$ pour tout $n$,
nous trouvons que
$$
I^n\mathcal{F}_{n + 1} =
\bigoplus\nolimits_{e_1 + \ldots + e_d = n} \mathcal{F}_1 \cdot
f_1^{e_1} \ldots f_d^{e_d}
$$
en le vérifiant sur les germes. En utilisant l'hypothèse de finitude de
$H^0(U, \mathcal{F}_1)$ et une récurrence, nous concluons d'abord que
$M_n = H^0(U, \mathcal{F}_n)$ est un $A$-module de type fini pour tout $n$.
Nous voyons ainsi que la condition (c) du lemme \ref{algebraization-lemma-when-done} est vérifiée.
Nous voyons aussi que
$$
\bigoplus\nolimits_{n \geq 0} H^1(U, I^n\mathcal{F}_{n + 1})
$$
est un $R = \bigoplus I^n/I^{n +1}$-module gradué de type fini.
D'après Cohomologie, lemme \ref{cohomology-lemma-ML-general},
nous concluons que la condition (a) du
lemme \ref{algebraization-lemma-when-done} est vérifiée. Enfin, la condition (b) du
lemme \ref{algebraization-lemma-when-done} est vérifiée parce que
$\bigoplus H^0(U, I^n\mathcal{F}_{n + 1})$ est un $R$-module gradué de type fini
et que nous pouvons appliquer
Cohomologie, lemme \ref{cohomology-lemma-topology-I-adic-general}.
\end{proof}

\begin{remark}
\label{algebraization-remark-interesting-case-ter}
Dans la situation de la
proposition \ref{algebraization-proposition-algebraization-regular-sequence},
si nous supposons que $A$ possède un complexe dualisant, alors
la condition que $H^0(U, \mathcal{F}_1)$ et
$H^1(U, \mathcal{F}_1)$ soient de type fini équivaut à
$$
\text{depth}(\mathcal{F}_{1, y}) +
\dim(\mathcal{O}_{\overline{\{y\}}, z}) > 2
$$
pour tout $y \in U \cap Y$ et tout $z \in Z \cap \overline{\{y\}}$.
Voir Cohomologie locale, lemme \ref{local-cohomology-lemma-finiteness-Rjstar}.
Cela vaut par exemple si $\mathcal{F}_1$ est un
$\mathcal{O}_{U \cap Y}$-module localement libre de type fini, si $Y$ est $(S_2)$, et si
$\text{codim}(Z', Y') \geq 3$ pour toute paire de composantes irréductibles
$Y'$ de $Y$ et $Z'$ de $Z$ telles que $Z' \subset Y'$.
\end{remark}

\begin{proposition}
\label{algebraization-proposition-algebraization-flat}
Dans la situation \ref{algebraization-situation-algebraize}, soit
$(\mathcal{F}_n)$ un objet de $\textit{Coh}(U, I\mathcal{O}_U)$.
Supposons qu'il existe un anneau local noethérien $(R, \mathfrak m)$ et un homomorphisme
d'anneaux $R \to A$ tels que
\begin{enumerate}
\item $I = \mathfrak m A$,
\item pour $y \in U \cap Y$, le germe $\mathcal{F}_y^\wedge$ soit $R$-plat,
\item $H^0(U, \mathcal{F}_1)$ et $H^1(U, \mathcal{F}_1)$ soient des
$A$-modules de type fini.
\end{enumerate}
Alors $(\mathcal{F}_n)$ s'étend canoniquement à $X$. En particulier, si $A$
est complet, alors $(\mathcal{F}_n)$ est le complété d'un
$\mathcal{O}_U$-module cohérent.
\end{proposition}

\begin{proof}
La démonstration est exactement la même que celle de la
proposition \ref{algebraization-proposition-algebraization-regular-sequence}.
En effet, si $\kappa = R/\mathfrak m$, alors, pour $n \geq 0$,
il existe un isomorphisme
$$
I^n \mathcal{F}_{n + 1} \cong
\mathcal{F}_1 \otimes_\kappa \mathfrak m^n/\mathfrak m^{n + 1}
$$
et le membre de droite est une somme directe finie de copies
de $\mathcal{F}_1$. Cela se vérifie sur les germes.
Tout le reste est exactement identique.
\end{proof}

\begin{remark}
\label{algebraization-remark-interesting-case-quater}
La proposition \ref{algebraization-proposition-algebraization-flat} est une version locale
de \cite[théorème 2.10 (i)]{Baranovsky}. Il est immédiat de déduire
les résultats globaux du résultat local ; esquissons l'argument.
En effet, supposons que $(R, \mathfrak m)$
soit un anneau local noethérien complet et que $X \to \Spec(R)$ soit un morphisme propre.
Pour $n \geq 1$, posons $X_n = X \times_{\Spec(R)} \Spec(R/\mathfrak m^n)$.
Soit $Z \subset X_1$ un sous-ensemble fermé de la fibre spéciale.
Posons $U = X \setminus Z$ et notons $j : U \to X$ le morphisme d'inclusion.
Supposons donné un objet
$$
(\mathcal{F}_n) \text{ de } \textit{Coh}(U, \mathfrak m\mathcal{O}_U)
$$
qui soit plat sur $R$ au sens où $\mathcal{F}_n$ est plat sur
$R/\mathfrak m^n$ pour tout $n$.
Supposons que $j_*\mathcal{F}_1$ et $R^1j_*\mathcal{F}_1$ soient des modules
cohérents. Alors, localement sur les ouverts affines de $X$, nous obtenons une extension canonique
de $(\mathcal{F}_n)$ par la
proposition \ref{algebraization-proposition-algebraization-flat},
et la formation de cette extension commute à la localisation
(d'après le lemme \ref{algebraization-lemma-algebraization-principal-variant}).
Nous obtenons donc un objet global canonique $(\mathcal{G}_n)$ de
$\textit{Coh}(X, \mathfrak m\mathcal{O}_X)$
dont la restriction à $U$ est $(\mathcal{F}_n)$.
D'après le théorème d'existence de Grothendieck
(Cohomologie des schémas, proposition
\ref{coherent-proposition-existence-proper}),
nous voyons qu'il existe un $\mathcal{O}_X$-module cohérent
$\mathcal{G}$ dont le complété est $(\mathcal{G}_n)$.
Nous voyons ainsi que $(\mathcal{F}_n)$ est algébrisable, c'est-à-dire
qu'il est le complété d'un $\mathcal{O}_U$-module cohérent.

\medskip\noindent
Ajoutons que la cohérence de $j_*\mathcal{F}_1$ et de $R^1j_*\mathcal{F}_1$
est une condition sur la fibre spéciale. En effet, si nous notons
$j_1 : U_1 \to X_1$ la fibre spéciale de $j : U \to X$, alors nous pouvons
considérer $\mathcal{F}_1$ comme un faisceau cohérent sur $U_1$, et nous avons
$j_*\mathcal{F}_1 = j_{1, *}\mathcal{F}_1$ et
$R^1j_*\mathcal{F}_1 = R^1j_{1, *}\mathcal{F}_1$.
Ainsi, par exemple, si $X_1$ est $(S_2)$ et irréductible, si nous avons
$\dim(X_1) - \dim(Z) \geq 3$, et si $\mathcal{F}_1$ est un
$\mathcal{O}_{U_1}$-module localement libre, alors $j_{1, *}\mathcal{F}_1$ et
$R^1j_{1, *}\mathcal{F}_1$ sont des modules cohérents.
\end{remark}








\section{Application au foncteur de complétion}
\label{algebraization-section-completion-application}

\noindent
Dans cette section, nous combinons seulement quelques résultats déjà obtenus
afin de pouvoir les citer commodément. Nous pourrions énoncer ici de nombreux
résultats (plus forts).

\begin{lemma}
\label{algebraization-lemma-equivalence-better}
Dans la situation \ref{algebraization-situation-algebraize}, supposons que
\begin{enumerate}
\item $A$ possède un complexe dualisant et soit complet pour la topologie $I$-adique,
\item $I = (f)$ soit engendré par un seul élément,
\item $A$ soit local d'idéal maximal $\mathfrak a = \mathfrak m$,
\item l'une des conditions suivantes soit vérifiée :
\begin{enumerate}
\item $A_f$ est $(S_2)$ et, pour $\mathfrak p \subset A$
minimal tel que $f \not \in \mathfrak p$, nous avons $\dim(A/\mathfrak p) \geq 4$, ou
\item si $\mathfrak p \not \in V(f)$ et
$V(\mathfrak p) \cap V(f) \not = \{\mathfrak m\}$, alors
$\text{depth}(A_\mathfrak p) + \dim(A/\mathfrak p) > 3$.
\end{enumerate}
\end{enumerate}
Alors, avec $U_0 = U \cap V(f)$, le foncteur de complétion
$$
\colim_{U_0 \subset U' \subset U\text{ ouvert}}
\textit{Coh}(\mathcal{O}_{U'})
\longrightarrow
\textit{Coh}(U, f\mathcal{O}_U)
$$
est une équivalence sur les sous-catégories pleines d'objets localement libres de type fini.
\end{lemma}

\begin{proof}
Il résulte du lemme \ref{algebraization-lemma-fully-faithful-general}
que le foncteur est pleinement fidèle (détails omis).
Démontrons la surjectivité essentielle. Soit $(\mathcal{F}_n)$ un objet localement
libre de type fini de $\textit{Coh}(U, f\mathcal{O}_U)$. D'après le
lemme \ref{algebraization-lemma-algebraization-principal-bis} ou la
proposition \ref{algebraization-proposition-cd-1},
il existe un $\mathcal{O}_U$-module cohérent $\mathcal{F}$
tel que $(\mathcal{F}_n)$ soit le complété de $\mathcal{F}$.
En effet, pour appliquer l'un ou l'autre résultat, la seule chose à
vérifier est que $(\mathcal{F}_n)$ satisfait aux inégalités $(2, 3)$.
Cela est fait dans le lemme \ref{algebraization-lemma-unwinding-conditions-bis}. Si $y \in U_0$,
alors le complété $f$-adique du germe $\mathcal{F}_y$ est isomorphe à
un module libre de type fini sur le complété $f$-adique de $\mathcal{O}_{U, y}$.
Par conséquent, $\mathcal{F}$ est localement libre de type fini dans un voisinage ouvert
$U'$ de $U_0$. Ceci achève la démonstration.
\end{proof}

\begin{lemma}
\label{algebraization-lemma-equivalence}
Dans la situation \ref{algebraization-situation-algebraize}, supposons que
\begin{enumerate}
\item $I = (f)$ soit principal,
\item $A$ soit complet pour la topologie $f$-adique,
\item $f$ soit un non-diviseur de zéro,
\item $H^1_\mathfrak a(A/fA)$ et $H^2_\mathfrak a(A/fA)$
soient des $A$-modules de type fini.
\end{enumerate}
Alors, avec $U_0 = U \cap V(f)$, le foncteur de complétion
$$
\colim_{U_0 \subset U' \subset U\text{ ouvert}}
\textit{Coh}(\mathcal{O}_{U'})
\longrightarrow
\textit{Coh}(U, f\mathcal{O}_U)
$$
est une équivalence sur les sous-catégories pleines d'objets localement libres de type fini.
\end{lemma}

\begin{proof}
Le foncteur est pleinement fidèle d'après le
lemme \ref{algebraization-lemma-fully-faithful-general-alternative}.
La surjectivité essentielle résulte du
lemme \ref{algebraization-lemma-algebraization-principal-variant}.
\end{proof}








\section{Triplets cohérents}
\label{algebraization-section-coherent-triples}

\noindent
Soit $(A, \mathfrak m)$ un anneau local noethérien.
Soit $f \in \mathfrak m$ un non-diviseur de zéro. Posons
$X = \Spec(A)$, $X_0 = \Spec(A/fA)$, $U = X \setminus V(\mathfrak m)$ et
$U_0 = U \cap X_0$.
Nous disons que $(\mathcal{F}, \mathcal{F}_0, \alpha)$ est un {\it triplet cohérent}
si nous avons
\begin{enumerate}
\item $\mathcal{F}$ est un $\mathcal{O}_U$-module cohérent tel que
$f : \mathcal{F} \to \mathcal{F}$ soit injective,
\item $\mathcal{F}_0$ est un $\mathcal{O}_{X_0}$-module cohérent,
\item $\alpha : \mathcal{F}/f\mathcal{F} \to \mathcal{F}_0|_{U_0}$
est un isomorphisme.
\end{enumerate}
Il existe une notion évidente de {\it morphisme de triplets cohérents}
qui fait de la collection de tous les triplets cohérents une catégorie.

\medskip\noindent
La catégorie des triplets cohérents est additive mais non abélienne.
Cependant, la notion de suite exacte courte de triplets
cohérents est claire.

\medskip\noindent
Étant donnés deux triplets cohérents $(\mathcal{F}, \mathcal{F}_0, \alpha)$
et $(\mathcal{G}, \mathcal{G}_0, \beta)$, il se peut que
$(\mathcal{F} \otimes_{\mathcal{O}_U} \mathcal{G},
\mathcal{F}_0  \otimes_{\mathcal{O}_{X_0}} \mathcal{G}_0,
\alpha \otimes \beta)$ ne soit pas un triplet cohérent\footnote{En effet, il
n'est pas nécessairement vrai que $f$
soit injectif sur $\mathcal{F} \otimes_{\mathcal{O}_U} \mathcal{G}$.}.
Cependant, si les germes $\mathcal{G}_x$ sont
libres pour tout $x \in U_0$, alors c'est bien le cas.

\medskip\noindent
Nous dirons que le triplet cohérent $(\mathcal{G}, \mathcal{G}_0, \beta)$
est {\it localement libre}, resp.\ {\it inversible},
si $\mathcal{G}$ et $\mathcal{G}_0$
sont des modules localement libres, resp.\ inversibles. Dans ce cas, tensoriser
par $(\mathcal{G}, \mathcal{G}_0, \beta)$ a un sens (voir ci-dessus)
et transforme les suites exactes courtes de triplets cohérents en suites exactes courtes
de triplets cohérents.

\begin{lemma}
\label{algebraization-lemma-prepare-chi-triple}
Pour tout triplet cohérent $(\mathcal{F}, \mathcal{F}_0, \alpha)$,
il existe un $\mathcal{O}_X$-module cohérent $\mathcal{F}'$
tel que $f : \mathcal{F}' \to \mathcal{F}'$ soit injective,
un isomorphisme $\alpha' : \mathcal{F}'|_U \to \mathcal{F}$ et une application
$\alpha'_0 : \mathcal{F}'/f\mathcal{F}' \to \mathcal{F}_0$
tel que $\alpha \circ (\alpha' \bmod f) = \alpha'_0|_{U_0}$.
\end{lemma}

\begin{proof}
Choisissons un $A$-module de type fini $M$ tel que $\mathcal{F}$ soit la restriction
à $U$ du $\mathcal{O}_X$-module cohérent associé à $M$, voir
Cohomologie locale, lemme
\ref{local-cohomology-lemma-finiteness-pushforwards-and-H1-local}.
Puisque $\mathcal{F}$ est sans torsion par $f$, nous pouvons remplacer $M$ par
son quotient par la torsion annulée par une puissance de $f$.
D'autre part, posons $M_0 = \Gamma(X_0, \mathcal{F}_0)$,
de sorte que $\mathcal{F}_0$ soit le $\mathcal{O}_{X_0}$-module cohérent
associé au $A/fA$-module de type fini $M_0$.
D'après Cohomologie des schémas, lemme \ref{coherent-lemma-homs-over-open},
il existe un $n$ tel que
l'isomorphisme $\alpha_0$ corresponde à un homomorphisme de $A/fA$-modules
$\mathfrak m^n M/fM \to M_0$ (dont le noyau et le conoyau
sont annulés par une puissance de $\mathfrak m$, mais nous n'en avons pas besoin).
Ainsi, si nous prenons $M' = \mathfrak m^n M$ et si nous posons
$\mathcal{F}'$ égal au $\mathcal{O}_X$-module cohérent
associé à $M'$, alors le lemme est clair.
\end{proof}

\noindent
Soit $(\mathcal{F}, \mathcal{F}_0, \alpha)$ un triplet cohérent.
Choisissons $\mathcal{F}', \alpha', \alpha'_0$ comme dans le
lemme \ref{algebraization-lemma-prepare-chi-triple}. Posons
\begin{equation}
\label{algebraization-equation-chi-triple}
\chi(\mathcal{F}, \mathcal{F}_0, \alpha) =
\text{length}_A(\Coker(\alpha'_0)) - 
\text{length}_A(\Ker(\alpha'_0))
\end{equation}
L'expression du membre de droite a un sens, car $\alpha'_0$ est un isomorphisme
sur $U_0$ ; son noyau et son conoyau sont donc des modules cohérents supportés
par $\{\mathfrak m\}$, qui sont par conséquent de longueur finie
(Algèbre, lemme \ref{algebra-lemma-support-point}).

\begin{lemma}
\label{algebraization-lemma-well-defined-chi-triple}
La quantité $\chi(\mathcal{F}, \mathcal{F}_0, \alpha)$ de
(\ref{algebraization-equation-chi-triple}) ne dépend pas du choix de
$\mathcal{F}', \alpha', \alpha'_0$ comme dans le lemme \ref{algebraization-lemma-prepare-chi-triple}.
\end{lemma}

\begin{proof}
Soient $\mathcal{F}', \alpha', \alpha'_0$ et
$\mathcal{F}'', \alpha'', \alpha''_0$ deux tels choix.
Pour $n > 0$, posons $\mathcal{F}'_n = \mathfrak m^n \mathcal{F}'$.
D'après Cohomologie des schémas, lemme \ref{coherent-lemma-homs-over-open},
pour un certain $n$,
il existe un homomorphisme de $\mathcal{O}_X$-modules $\mathcal{F}'_n \to \mathcal{F}''$
qui coïncide avec l'identification
$\mathcal{F}''|_U = \mathcal{F}'|_U$ déterminée par $\alpha'$ et $\alpha''$.
Alors le diagramme
$$
\xymatrix{
\mathcal{F}'_n/f\mathcal{F}'_n \ar[r] \ar[d] &
\mathcal{F}'/f\mathcal{F}' \ar[d]^{\alpha_0'} \\
\mathcal{F}''/f\mathcal{F}'' \ar[r]^{\alpha_0''} &
\mathcal{F}_0
}
$$
est commutatif après restriction à $U_0$. Par conséquent, d'après
Cohomologie des schémas, lemme \ref{coherent-lemma-homs-over-open},
il est commutatif après restriction à
$\mathfrak m^l(\mathcal{F}'_n/f\mathcal{F}'_n)$ pour un certain $l > 0$. Puisque
$\mathcal{F}'_{n + l}/f\mathcal{F}'_{n + l} \to \mathcal{F}'_n/f\mathcal{F}'_n$
se factorise par $\mathfrak m^l(\mathcal{F}'_n/f\mathcal{F}'_n)$,
nous voyons qu'après avoir remplacé $n$ par $n + l$, le diagramme
est commutatif. En d'autres termes, nous avons trouvé un troisième choix
$\mathcal{F}''', \alpha''', \alpha'''_0$
tel qu'il existe des applications $\mathcal{F}''' \to \mathcal{F}''$
et $\mathcal{F}''' \to \mathcal{F}'$ sur $X$,
compatibles avec les applications sur $U$ et $X_0$. Nous sommes ainsi ramenés
au cas examiné dans le paragraphe suivant.

\medskip\noindent
Supposons disposer d'une application $\mathcal{F}'' \to \mathcal{F}'$ sur $X$, compatible avec
$\alpha', \alpha''$ sur $U$ et avec $\alpha'_0, \alpha''_0$ sur $X_0$.
Observons que $\mathcal{F}'' \to \mathcal{F}'$ est injective, car c'est un
isomorphisme sur $U$ et que $f : \mathcal{F}'' \to \mathcal{F}''$
est injective. Il est clair que $\mathcal{F}'/\mathcal{F}''$ est supporté par
$\{\mathfrak m\}$, donc est de longueur finie. Nous avons les applications
de $\mathcal{O}_{X_0}$-modules cohérents
$$
\mathcal{F}''/f\mathcal{F}'' \to
\mathcal{F}'/f\mathcal{F}' \xrightarrow{\alpha'_0}
\mathcal{F}_0
$$
dont la composée est $\alpha''_0$ et qui sont des isomorphismes sur $U_0$.
L'algèbre homologique élémentaire donne une suite exacte à $6$ termes
$$
\begin{matrix}
0 \to
\Ker(\mathcal{F}''/f\mathcal{F}'' \to \mathcal{F}'/f\mathcal{F}') \to
\Ker(\alpha''_0) \to
\Ker(\alpha'_0) \to \\
\Coker(\mathcal{F}''/f\mathcal{F}'' \to \mathcal{F}'/f\mathcal{F}') \to
\Coker(\alpha''_0) \to
\Coker(\alpha'_0) \to 0
\end{matrix}
$$
Par additivité des longueurs (Algèbre, lemme \ref{algebra-lemma-length-additive}),
nous trouvons qu'il suffit de montrer que
$$
\text{length}_A(
\Coker(\mathcal{F}''/f\mathcal{F}'' \to \mathcal{F}'/f\mathcal{F}')) -
\text{length}_A(
\Ker(\mathcal{F}''/f\mathcal{F}'' \to \mathcal{F}'/f\mathcal{F}')) = 0
$$
Cela résulte de l'application du lemme du serpent au
diagramme
$$
\xymatrix{
0 \ar[r] &
\mathcal{F}'' \ar[r]_f \ar[d] &
\mathcal{F}'' \ar[r] \ar[d] &
\mathcal{F}''/f\mathcal{F}'' \ar[r] \ar[d] &
0 \\
0 \ar[r] &
\mathcal{F}' \ar[r]^f &
\mathcal{F}' \ar[r] &
\mathcal{F}'/f\mathcal{F}' \ar[r] &
0
}
$$
et du fait que $\mathcal{F}'/\mathcal{F}''$ est de longueur finie.
\end{proof}

\begin{lemma}
\label{algebraization-lemma-ses-chi-triple}
Nous avons
$\chi(\mathcal{G}, \mathcal{G}_0, \beta) =
\chi(\mathcal{F}, \mathcal{F}_0, \alpha) +
\chi(\mathcal{H}, \mathcal{H}_0, \gamma)$ si
$$
0 \to
(\mathcal{F}, \mathcal{F}_0, \alpha) \to
(\mathcal{G}, \mathcal{G}_0, \beta) \to
(\mathcal{H}, \mathcal{H}_0, \gamma)
\to 0
$$
est une suite exacte courte de triplets cohérents.
\end{lemma}

\begin{proof}
Choisissons $\mathcal{G}', \beta', \beta'_0$ comme dans le
lemme \ref{algebraization-lemma-prepare-chi-triple}
pour le triplet $(\mathcal{G}, \mathcal{G}_0, \beta)$.
Notons $j : U \to X$ le morphisme d'inclusion.
Soit $\mathcal{F}' \subset \mathcal{G}'$
le noyau de la composée
$$
\mathcal{G}' \xrightarrow{\beta'} j_*\mathcal{G} \to j_*\mathcal{H}
$$
Observons que $\mathcal{H}' = \mathcal{G}'/\mathcal{F}'$
est un sous-faisceau cohérent de $j_*\mathcal{H}$ et que, par conséquent,
$f : \mathcal{H}' \to \mathcal{H}'$ est injective.
Le lemme du serpent donne donc une suite exacte courte
$$
0 \to \mathcal{F}'/f\mathcal{F}' \to
\mathcal{G}'/f\mathcal{G}' \to
\mathcal{H}'/f\mathcal{H}' \to 0
$$
Nous avons par construction les isomorphismes
$\alpha' : \mathcal{F}'|_U \to \mathcal{F}$,
$\beta' : \mathcal{G}'|_U \to \mathcal{G}$ et
$\gamma' : \mathcal{H}'|_U \to \mathcal{H}$.
Pour achever la démonstration, il nous faudra construire des applications
$\alpha'_0 : \mathcal{F}'/f\mathcal{F}' \to \mathcal{F}_0$ et
$\gamma'_0 : \mathcal{H}'/f\mathcal{H}' \to \mathcal{H}_0$ comme dans le
lemme \ref{algebraization-lemma-prepare-chi-triple} et s'insérant dans
un diagramme commutatif
$$
\xymatrix{
0 \ar[r] &
\mathcal{F}'/f\mathcal{F}' \ar[r] \ar@{..>}[d]^{\alpha'_0} &
\mathcal{G}'/f\mathcal{G}' \ar[r] \ar[d]^{\beta'_0} &
\mathcal{H}'/f\mathcal{H}' \ar[r] \ar@{..>}[d]^{\gamma'_0} &
0 \\
0 \ar[r] &
\mathcal{F}_0 \ar[r] &
\mathcal{G}_0 \ar[r] &
\mathcal{H}_0 \ar[r] &
0
}
$$
Cependant, cela peut être impossible avec notre choix initial de $\mathcal{G}'$.
Le diagramme affiché montre que l'obstruction est
exactement la composée
$$
\delta :
\mathcal{F}'/f\mathcal{F}' \to
\mathcal{G}'/f\mathcal{G}' \xrightarrow{\beta'_0}
\mathcal{G}_0 \to
\mathcal{H}_0
$$
Notons que la restriction de $\delta$ à $U_0$ est nulle par notre choix de
$\mathcal{F}'$ et $\mathcal{H}'$. Par conséquent, d'après
Cohomologie des schémas, lemme \ref{coherent-lemma-homs-over-open},
il existe un $k > 0$ tel que
$\delta$ s'annule sur $\mathfrak m^k \cdot (\mathcal{F}'/f\mathcal{F}')$.
Pour $n > k$, posons $\mathcal{G}'_n = \mathfrak m^n \mathcal{G}'$,
$\mathcal{F}'_n = \mathcal{G}'_n \cap \mathcal{F}'$ et
$\mathcal{H}'_n = \mathcal{G}'_n/\mathcal{F}'_n$.
Observons que $\beta'_0$ peut être composée avec
$\mathcal{G}'_n/f\mathcal{G}'_n \to \mathcal{G}'/f\mathcal{G}'$
pour donner une application
$\beta'_{n, 0} : \mathcal{G}'_n/f\mathcal{G}'_n \to \mathcal{G}_0$
comme dans le lemme \ref{algebraization-lemma-prepare-chi-triple}.
D'après Artin--Rees (Algèbre, lemme \ref{algebra-lemma-Artin-Rees}),
nous pouvons choisir $n$ tel que
$\mathcal{F}'_n \subset \mathfrak m^k \mathcal{F}'$.
Comme ci-dessus, les applications
$f : \mathcal{F}'_n \to \mathcal{F}'_n$,
$f : \mathcal{G}'_n \to \mathcal{G}'_n$ et
$f : \mathcal{H}'_n \to \mathcal{H}'_n$ sont injectives
et, comme ci-dessus, le lemme du serpent donne une suite exacte
courte
$$
0 \to \mathcal{F}'_n/f\mathcal{F}'_n \to
\mathcal{G}'_n/f\mathcal{G}'_n \to
\mathcal{H}'_n/f\mathcal{H}'_n \to 0
$$
Comme ci-dessus, nous avons les isomorphismes
$\alpha'_n : \mathcal{F}'_n|_U \to \mathcal{F}$,
$\beta'_n : \mathcal{G}'_n|_U \to \mathcal{G}$ et
$\gamma'_n : \mathcal{H}'_n|_U \to \mathcal{H}$.
Considérons l'obstruction
$$
\delta_n :
\mathcal{F}'_n/f\mathcal{F}'_n \to
\mathcal{G}'_n/f\mathcal{G}'_n
\xrightarrow{\beta'_{n, 0}}
\mathcal{G}_0 \to
\mathcal{H}_0
$$
comme précédemment. Cependant, le diagramme commutatif
$$
\xymatrix{
\mathcal{F}'_n/f\mathcal{F}'_n \ar[r] \ar[d] &
\mathcal{G}'_n/f\mathcal{G}'_n \ar[r]_{\beta'_{n, 0}} \ar[d] &
\mathcal{G}_0 \ar[r] \ar[d] &
\mathcal{H}_0 \ar[d] \\
\mathcal{F}'/f\mathcal{F}' \ar[r] &
\mathcal{G}'/f\mathcal{G}' \ar[r]^{\beta'_0} &
\mathcal{G}_0 \ar[r] &
\mathcal{H}_0
}
$$
notre choix de $n$ et notre observation concernant $\delta$
montrent que $\delta_n = 0$.
Ceci fournit les applications voulues
$\alpha'_{n, 0} : \mathcal{F}'_n/f\mathcal{F}'_n \to \mathcal{F}_0$ et
$\gamma'_{n, 0} : \mathcal{H}'_n/f\mathcal{H}'_n \to \mathcal{H}_0$.
Ainsi, nous pouvons utiliser
$\mathcal{F}'_n, \alpha'_n, \alpha'_{n, 0}$,
$\mathcal{G}'_n, \beta'_n, \beta'_{n, 0}$ et
$\mathcal{H}'_n, \gamma'_n, \gamma'_{n, 0}$
pour calculer
$\chi(\mathcal{F}, \mathcal{F}_0, \alpha)$,
$\chi(\mathcal{G}, \mathcal{G}_0, \beta)$ et
$\chi(\mathcal{H}, \mathcal{H}_0, \gamma)$.
Enfin, le lemme résulte d'une
application du lemme du serpent à
$$
\xymatrix{
0 \ar[r] &
\mathcal{F}'_n/f\mathcal{F}'_n \ar[r] \ar[d] &
\mathcal{G}'_n/f\mathcal{G}'_n \ar[r] \ar[d] &
\mathcal{H}'_n/f\mathcal{H}'_n \ar[r] \ar[d] &
0 \\
0 \ar[r] &
\mathcal{F}_0 \ar[r] &
\mathcal{G}_0 \ar[r] &
\mathcal{H}_0 \ar[r] &
0
}
$$
et de l'additivité des longueurs (Algèbre, lemme \ref{algebra-lemma-length-additive}).
\end{proof}

\begin{proposition}
\label{algebraization-proposition-hilbert-triple}
Soit $(\mathcal{F}, \mathcal{F}_0, \alpha)$ un triplet cohérent.
Soit $(\mathcal{L}, \mathcal{L}_0, \lambda)$ un triplet cohérent
inversible. Alors la fonction
$$
\mathbf{Z} \longrightarrow \mathbf{Z},\quad
n \longmapsto
\chi((\mathcal{F}, \mathcal{F}_0, \alpha) \otimes
(\mathcal{L}, \mathcal{L}_0, \lambda)^{\otimes n})
$$
est un polynôme de degré $\leq \dim(\text{Supp}(\mathcal{F}))$.
\end{proposition}

\noindent
Plus précisément, si $\mathcal{F} = 0$, alors la fonction est constante.
Si $\mathcal{F}$ est à support fini dans $U$, alors la fonction est constante.
Si le support de $\mathcal{F}$ dans $U$ est de dimension $1$, c'est-à-dire
si l'adhérence du support de $\mathcal{F}$ dans $X$ est de dimension $2$, alors
la fonction est linéaire, et ainsi de suite.

\begin{proof}
Nous allons le démontrer par récurrence sur la dimension du support de
$\mathcal{F}$.

\medskip\noindent
Le cas initial est celui où $\mathcal{F} = 0$. Alors soit
$\mathcal{F}_0$ est nul, soit son support est $\{\mathfrak m\}$.
Dans ce cas, nous avons
$$
(\mathcal{F}, \mathcal{F}_0, \alpha) \otimes
(\mathcal{L}, \mathcal{L}_0, \lambda)^{\otimes n} =
(0, \mathcal{F}_0 \otimes \mathcal{L}_0^{\otimes n}, 0) \cong
(0, \mathcal{F}_0, 0)
$$
Ainsi, la fonction de la proposition est constante, de valeur égale
à la longueur de $\mathcal{F}_0$.

\medskip\noindent
Étape de récurrence. Supposons le support de $\mathcal{F}$ non vide.
Soit $\mathcal{G}_0 \subset \mathcal{F}_0$ le sous-module
des sections supportées par $\{\mathfrak m\}$. Nous obtenons alors une suite
exacte courte
$$
0 \to (0, \mathcal{G}_0, 0) \to
(\mathcal{F}, \mathcal{F}_0, \alpha) \to
(\mathcal{F}, \mathcal{F}_0/\mathcal{G}_0, \alpha) \to 0
$$
Cette suite reste exacte si nous tensorisons par le triplet cohérent
inversible $(\mathcal{L}, \mathcal{L}_0, \lambda)$, voir la
discussion ci-dessus. Ainsi, par additivité de $\chi$
(lemme \ref{algebraization-lemma-ses-chi-triple})
et par le cas initial expliqué ci-dessus, il suffit de démontrer
l'étape de récurrence pour
$(\mathcal{F}, \mathcal{F}_0/\mathcal{G}_0, \alpha)$.
Nous voyons ainsi que nous pouvons supposer que $\mathfrak m$ n'est pas
un point associé de $\mathcal{F}_0$.

\medskip\noindent
Soit $T = \text{Ass}(\mathcal{F}) \cup \text{Ass}(\mathcal{F}/f\mathcal{F})$.
Puisque $U$ est quasi-affine, nous pouvons trouver $s \in \Gamma(U, \mathcal{L})$
qui ne s'annule en aucun $u \in T$, voir
Propriétés, lemme
\ref{properties-lemma-quasi-affine-invertible-nonvanishing-section}.
Après avoir multiplié $s$ par un élément convenable de $\mathfrak m$,
nous pouvons supposer que $\lambda(s \bmod f) = s_0|_{U_0}$ pour un certain
$s_0 \in \Gamma(X_0, \mathcal{L}_0)$ ; détails omis.
Nous obtenons un morphisme
$$
(s, s_0) :
(\mathcal{O}_U, \mathcal{O}_{X_0}, 1)
\longrightarrow
(\mathcal{L}, \mathcal{L}_0, \lambda)
$$
dans la catégorie des triplets cohérents. Posons
$\mathcal{G} = \Coker(s : \mathcal{F} \to \mathcal{F} \otimes \mathcal{L})$
et
$\mathcal{G}_0 = \Coker(s_0 : \mathcal{F}_0 \to
\mathcal{F}_0 \otimes \mathcal{L}_0)$. Observons que $s_0 : \mathcal{F}_0 \to
\mathcal{F}_0 \otimes \mathcal{L}_0$ est injective, car elle est injective
sur $U_0$ par notre choix de $s$ et que $\mathfrak m$ n'est pas un
point associé de $\mathcal{F}_0$. Par conséquent,
il existe un
isomorphisme $\beta : \mathcal{G}/f\mathcal{G} \to \mathcal{G}_0|_{U_0}$
tel que nous obtenions une suite exacte courte
$$
0 \to
(\mathcal{F}, \mathcal{F}_0, \alpha) \to
(\mathcal{F}, \mathcal{F}_0, \alpha) \otimes
(\mathcal{L}, \mathcal{L}_0, \lambda) \to
(\mathcal{G}, \mathcal{G}_0, \beta) \to 0
$$
Par récurrence sur la dimension du support, nous savons que la proposition
est vraie pour le triplet cohérent $(\mathcal{G}, \mathcal{G}_0, \beta)$.
En utilisant l'additivité du lemme \ref{algebraization-lemma-ses-chi-triple},
nous voyons que
$$
n \longmapsto 
\chi((\mathcal{F}, \mathcal{F}_0, \alpha) \otimes
(\mathcal{L}, \mathcal{L}_0, \lambda)^{\otimes n + 1})
-
\chi((\mathcal{F}, \mathcal{F}_0, \alpha) \otimes
(\mathcal{L}, \mathcal{L}_0, \lambda)^{\otimes n})
$$
est un polynôme. Nous concluons par une variante du
lemme \ref{algebra-lemma-numerical-polynomial} d'Algèbre
pour les fonctions définies sur tous les entiers (détails omis).
\end{proof}

\begin{lemma}
\label{algebraization-lemma-nonnegative-chi-triple}
Supposons que $\text{depth}(A) \geq 3$ ou, de manière équivalente, que
$\text{depth}(A/fA) \geq 2$. Soit $(\mathcal{L}, \mathcal{L}_0, \lambda)$
un triplet cohérent inversible. Alors
$$
\chi(\mathcal{L}, \mathcal{L}_0, \lambda) =
\text{length}_A \Coker(\Gamma(U, \mathcal{L}) \to \Gamma(U_0, \mathcal{L}_0))
$$
et, en particulier, ceci est $\geq 0$. De plus, 
$\chi(\mathcal{L}, \mathcal{L}_0, \lambda) = 0$ si et seulement si
$\mathcal{L} \cong \mathcal{O}_U$.
\end{lemma}

\begin{proof}
L'équivalence des conditions de profondeur résulte du
lemme \ref{algebra-lemma-depth-drops-by-one} d'Algèbre.
La condition de profondeur donne
$\Gamma(U, \mathcal{O}_U) = A$ et
$\Gamma(U_0, \mathcal{O}_{U_0}) = A/fA$, voir
Complexes dualisants, lemme \ref{dualizing-lemma-depth}, et
Cohomologie locale, lemme
\ref{local-cohomology-lemma-finiteness-pushforwards-and-H1-local}.
En utilisant Cohomologie locale, lemme
\ref{local-cohomology-lemma-finiteness-for-finite-locally-free},
nous trouvons que $M = \Gamma(U, \mathcal{L})$ est un $A$-module de type fini.
Ceci implique à son tour que $\text{depth}(M) \geq 2$, par exemple d'après la
partie (4) de Cohomologie locale, lemme
\ref{local-cohomology-lemma-finiteness-pushforwards-and-H1-local},
ou d'après Diviseurs, lemme \ref{divisors-lemma-depth-pushforward}.
De plus, nous avons $\mathcal{L}_0 \cong \mathcal{O}_{X_0}$,
car $X_0$ est un schéma local. Nous voyons donc aussi que
$M_0 = \Gamma(X_0, \mathcal{L}_0) = \Gamma(U_0, \mathcal{L}_0|_{U_0})$
et que ce module est isomorphe à $A/fA$.

\medskip\noindent
D'après ce qui précède, $\mathcal{F}' = \widetilde{M}$ est un
$\mathcal{O}_X$-module cohérent dont la restriction à $U$ est isomorphe à $\mathcal{L}$.
L'isomorphisme $\lambda : \mathcal{L}/f\mathcal{L} \to \mathcal{L}_0|_{U_0}$
détermine sur les sections globales une application $M/fM \to M_0$
qui est un isomorphisme sur $U_0$.
Puisque $\text{depth}(M) \geq 2$, nous voyons
que $H^0_\mathfrak m(M/fM) = 0$, et il s'ensuit que
$M/fM \to M_0$ est injective. Ainsi, par définition,
$$
\chi(\mathcal{L}, \mathcal{L}_0, \lambda) =
\text{length}_A \Coker(M/fM \to M_0)
$$
ce qui donne la première assertion du lemme.

\medskip\noindent
Enfin, si cette longueur vaut $0$, alors $M \to M_0$ est surjective.
Nous pouvons donc trouver $s \in M = \Gamma(U, \mathcal{L})$
qui s'envoie sur une section trivialisante de $\mathcal{L}_0$.
Considérons les $A$-modules de type fini $K$, $Q$ définis par la suite
exacte
$$
0 \to K \to A \xrightarrow{s} M \to Q \to 0
$$
Les supports de $K$ et de $Q$ ne rencontrent pas $U_0$, car $s$
est non nul aux points de $U_0$. En utilisant
Algèbre, lemme \ref{algebra-lemma-depth-in-ses},
nous voyons que $\text{depth}(K) \geq 2$ (observons que
$As \subset M$ est de $\text{depth} \geq 1$ comme sous-module de $M$).
Ainsi, s'il est non vide, le support de $K$ est de dimension $\geq 2$ d'après
Algèbre, lemme \ref{algebra-lemma-bound-depth}.
Ceci contredit $\text{Supp}(M) \cap V(f) \subset \{\mathfrak m\}$,
à moins que $K = 0$. Lorsque $K = 0$, nous trouvons que
$\text{depth}(Q) \geq 2$ et concluons
$Q = 0$ comme précédemment. Ainsi $A \cong M$ et
$\mathcal{L}$ est trivial.
\end{proof}







\section{Modules inversibles sur les spectres épointés, I}
\label{algebraization-section-local-lefschetz-for-pic}

\noindent
Dans cette section, nous démontrons quelques théorèmes de Lefschetz locaux pour le groupe de Picard.
Certaines idées sont tirées de
\cite{Kollar-pic}, \cite{Bhatt-local} et \cite{Kollar-map-pic}.

\begin{lemma}
\label{algebraization-lemma-injective-torsion-in-pic}
Soit $(A, \mathfrak m)$ un anneau local noethérien. Soit $f \in \mathfrak m$
un non-diviseur de zéro et supposons que $\text{depth}(A/fA) \geq 2$, ou, de manière équivalente,
$\text{depth}(A) \geq 3$. Soient $U$, resp.\ $U_0$, les spectres épointés
de $A$, resp.\ de $A/fA$. L'application
$$
\Pic(U) \to \Pic(U_0)
$$
est injective sur le sous-groupe de torsion.
\end{lemma}

\begin{proof}
Soit $\mathcal{L}$ un $\mathcal{O}_U$-module inversible.
Observons que $\mathcal{L}$ s'envoie sur $0$ dans $\Pic(U_0)$
si et seulement si nous pouvons prolonger $\mathcal{L}$ en un triplet cohérent
inversible $(\mathcal{L}, \mathcal{L}_0, \lambda)$
comme dans la section \ref{algebraization-section-coherent-triples}.
D'après la proposition \ref{algebraization-proposition-hilbert-triple},
la fonction
$$
n \longmapsto \chi((\mathcal{L}, \mathcal{L}_0, \lambda)^{\otimes n})
$$
est un polynôme. D'après le lemme \ref{algebraization-lemma-nonnegative-chi-triple},
la valeur de ce polynôme est nulle si et seulement si
$\mathcal{L}^{\otimes n}$ est trivial.
Ainsi, si $\mathcal{L}$ est de torsion, alors ce
polynôme a une infinité de zéros, donc est
identiquement nul, et par conséquent $\mathcal{L}$ est trivial.
\end{proof}

\begin{proposition}[Koll\'ar]
\label{algebraization-proposition-injective-pic}
\begin{reference}
\cite[Theorem 1.9]{Kollar-map-pic}
\end{reference}
Soit $(A, \mathfrak m)$ un anneau local noethérien. Soit $f \in \mathfrak m$.
Supposons que
\begin{enumerate}
\item $A$ possède un complexe dualisant,
\item $f$ soit un non-diviseur de zéro,
\item $\text{depth}(A/fA) \geq 2$, ou, de manière équivalente, $\text{depth}(A) \geq 3$,
\item si $f \in \mathfrak p \subset A$ est un idéal premier tel que
$\dim(A/\mathfrak p) = 2$, alors $\text{depth}(A_\mathfrak p) \geq 2$.
\end{enumerate}
Soient $U$, resp.\ $U_0$, les spectres épointés de $A$, resp.\ de $A/fA$. L'application
$$
\Pic(U) \to \Pic(U_0)
$$
est injective. Enfin, si (1), (2) et (3) sont vérifiées, si $A$ est $(S_2)$ et si
$\dim(A) \geq 4$, alors (4) est vérifiée.
\end{proposition}

\begin{proof}
Soit $\mathcal{L}$ un $\mathcal{O}_U$-module inversible.
Observons que $\mathcal{L}$ s'envoie sur $0$ dans $\Pic(U_0)$
si et seulement si nous pouvons prolonger $\mathcal{L}$ en un triplet cohérent
inversible $(\mathcal{L}, \mathcal{L}_0, \lambda)$
comme dans la section \ref{algebraization-section-coherent-triples}.
D'après la proposition \ref{algebraization-proposition-hilbert-triple},
la fonction
$$
n \longmapsto \chi((\mathcal{L}, \mathcal{L}_0, \lambda)^{\otimes n})
$$
est un polynôme $P$. D'après le lemme \ref{algebraization-lemma-nonnegative-chi-triple},
nous avons
$P(n) \geq 0$ pour tout $n \in \mathbf{Z}$, avec égalité si et seulement si
$\mathcal{L}^{\otimes n}$ est trivial. En particulier, $P(0) = 0$
et soit $P$ est identiquement nul, auquel cas nous concluons, soit $P$ est de degré pair $\geq 2$.

\medskip\noindent
Posons $M = \Gamma(U, \mathcal{L})$ et
$M_0 = \Gamma(X_0, \mathcal{L}_0) = \Gamma(U_0, \mathcal{L}_0)$.
Alors $M$ est un $A$-module de type fini de profondeur $\geq 2$
et $M_0 \cong A/fA$, voir la démonstration du lemme \ref{algebraization-lemma-nonnegative-chi-triple}.
Notons que $H^2_\mathfrak m(M)$ est un $A$-module de type fini d'après
Cohomologie locale, lemme
\ref{local-cohomology-lemma-local-finiteness-for-finite-locally-free},
et le fait que $H^i_\mathfrak m(A) = 0$ pour $i = 0, 1, 2$
puisque $\text{depth}(A) \geq 3$.
Considérons la suite exacte courte
$$
0 \to M/fM \to M_0 \to Q \to 0
$$
Le lemme \ref{algebraization-lemma-nonnegative-chi-triple} nous dit que $Q$ est de longueur finie
égale à $\chi(\mathcal{L}, \mathcal{L}_0, \lambda)$.
Nous obtenons $Q = H^1_\mathfrak m(M/fM)$ et
$H^i_\mathfrak m(M/fM) = H^i_\mathfrak m(M_0) \cong H^i_\mathfrak m(A/fA)$
pour $i > 1$ par la suite exacte longue de cohomologie locale
associée à la suite exacte courte affichée. Considérons la suite
exacte longue de cohomologie locale associée à la suite
$0 \to M \to M \to M/fM \to 0$. Elle commence par
$$
0 \to Q \to H^2_\mathfrak m(M) \to H^2_\mathfrak m(M) \to
H^2_\mathfrak m(A/fA)
$$
Par additivité des longueurs, nous voyons que
$\chi(\mathcal{L}, \mathcal{L}_0, \lambda)$
est égale à la longueur de l'image de
$H^2_\mathfrak m(M) \to H^2_\mathfrak m(A/fA)$.

\medskip\noindent
Démontrons la proposition dans un cas particulier afin d'éclairer la suite de la démonstration.
Supposons en effet un instant que $H^2_\mathfrak m(A/fA)$ soit
un module de longueur finie. Nous aurions alors
$P(1) \leq \text{length}_A H^2_\mathfrak m(A/fA)$.
Exactement le même argument, appliqué à $\mathcal{L}^{\otimes n}$, montre que
$P(n) \leq \text{length}_A H^2_\mathfrak m(A/fA)$ pour tout $n$.
Ainsi, $P$ ne peut pas être de degré positif, et nous concluons.
Dans la suite de la démonstration, nous allons modifier cet argument afin d'obtenir
une borne supérieure linéaire pour $P(n)$, ce qui suffit.

\medskip\noindent
Étudions l'application
$H^2_\mathfrak m(M) \to H^2_\mathfrak m(M_0) \cong H^2_\mathfrak m(A/fA)$.
Choisissons un complexe dualisant normalisé $\omega_A^\bullet$ pour $A$.
Par dualité locale
(Complexes dualisants, lemme \ref{dualizing-lemma-special-case-local-duality}),
cette application est, au sens de Matlis, duale de l'application
$$
\text{Ext}^{-2}_A(M, \omega_A^\bullet)
\longleftarrow
\text{Ext}^{-2}_A(M_0, \omega_A^\bullet)
$$
dont l'image a donc la même longueur (finie).
Le support (s'il est non vide) du $A$-module de type fini
$\text{Ext}^{-2}_A(M_0, \omega_A^\bullet)$ est constitué de
$\mathfrak m$ et d'un nombre fini d'idéaux premiers
$\mathfrak p_1, \ldots, \mathfrak p_r$ contenant $f$ et tels que
$\dim(A/\mathfrak p_i) = 1$. En effet, d'après
Cohomologie locale, lemme \ref{local-cohomology-lemma-sitting-in-degrees},
le support est contenu dans l'ensemble des idéaux premiers $\mathfrak p \subset A$ tels que
$\text{depth}_{A_\mathfrak p}(M_{0, \mathfrak p}) + \dim(A/\mathfrak p) \leq 2$.
Il suffit donc de montrer qu'il n'existe aucun idéal premier $\mathfrak p$ contenant $f$ tel que
$\dim(A/\mathfrak p) = 2$ et
$\text{depth}_{A_\mathfrak p}(M_{0, \mathfrak p}) = 0$.
Cependant, puisque $M_{0, \mathfrak p} \cong (A/fA)_\mathfrak p$,
cela donnerait $\text{depth}(A_\mathfrak p) = 1$, ce qui contredit
l'hypothèse (4).
Choisissons une section $t \in \Gamma(U, \mathcal{L}^{\otimes -1})$
qui ne s'annule en aucun des points $\mathfrak p_1, \ldots, \mathfrak p_r$, voir
Propriétés, lemme
\ref{properties-lemma-quasi-affine-invertible-nonvanishing-section}.
La multiplication par $t$ sur les sections globales détermine une application $t : M \to A$
qui définit un isomorphisme $M_{\mathfrak p_i} \to A_{\mathfrak p_i}$ pour
$i = 1, \ldots, r$. Notons $t_0 = t|_{U_0}$ la section correspondante
de $\Gamma(U_0, \mathcal{L}_0^{\otimes -1})$, qui détermine de même
une application $t_0 : M_0 \to A/fA$ compatible avec $t$.
Nous en concluons qu'il existe un diagramme commutatif
$$
\xymatrix{
\text{Ext}^{-2}_A(M, \omega_A^\bullet) &
\text{Ext}^{-2}_A(M_0, \omega_A^\bullet) \ar[l] \\
\text{Ext}^{-2}_A(A, \omega_A^\bullet) \ar[u]^t &
\text{Ext}^{-2}_A(A/fA, \omega_A^\bullet) \ar[l] \ar[u]_{t_0}
}
$$
Il s'ensuit que la longueur de l'image de l'application horizontale
supérieure est au plus la longueur de $\text{Ext}^{-2}_A(A/fA, \omega_A^\bullet)$
augmentée de la longueur du conoyau de $t_0$.

\medskip\noindent
Cependant, si nous remplaçons $\mathcal{L}$ par $\mathcal{L}^n$ pour $n > 1$,
alors nous pouvons utiliser
$$
t^n :
M_n = \Gamma(U, \mathcal{L}^{\otimes n})
\longrightarrow
\Gamma(U, \mathcal{O}_U) = A
$$
à la place de $t$. Ceci remplace $t_0$ par sa puissance $n$-ième.
Ainsi, la longueur de l'image de l'application
$\text{Ext}^{-2}_A(M_n, \omega_A^\bullet) \leftarrow
\text{Ext}^{-2}_A(M_{n, 0}, \omega_A^\bullet)$
est au plus la longueur de $\text{Ext}^{-2}_A(A/fA, \omega_A^\bullet)$
augmentée de la longueur du conoyau de
$$
t_0^n : 
\text{Ext}^{-2}_A(A/fA, \omega_A^\bullet)
\longrightarrow
\text{Ext}^{-2}_A(M_{n, 0}, \omega_A^\bullet)
$$
Via l'isomorphisme $M_0 \cong A/fA$, l'application $t_0$ devient
$g : A/fA \to A/fA$ pour un certain $g \in A/fA$ et, via les isomorphismes
correspondants $M_{n, 0} \cong A/fA$, l'application $t_0^n$ devient
$g^n : A/fA \to A/fA$. Ainsi, la longueur du conoyau ci-dessus
est la longueur du quotient de
$\text{Ext}^{-2}_A(A/fA, \omega_A^\bullet)$ par $g^n$.
Puisque $\text{Ext}^{-2}_A(A/fA, \omega_A^\bullet)$ est un $A$-module de type fini
de support $T$ de dimension $1$, et puisque $V(g) \cap T$
est constitué du point fermé par notre choix de $t$,
cette longueur croît linéairement en $n$ d'après
Algèbre, lemme \ref{algebra-lemma-support-dimension-d}.

\medskip\noindent
Pour achever la démonstration, prouvons la dernière assertion. Supposons que
$f \in \mathfrak m \subset A$ satisfasse
à (1), (2) et (3), que $A$ soit $(S_2)$ et que $\dim(A) \geq 4$.
La condition (1) implique que $A$ est caténaire, voir
Complexes dualisants, lemme \ref{dualizing-lemma-universally-catenary}.
Alors $\Spec(A)$ est équidimensionnel d'après Cohomologie locale, lemme
\ref{local-cohomology-lemma-catenary-S2-equidimensional}.
Ainsi, $\dim(A_\mathfrak p) + \dim(A/\mathfrak p) \geq 4$
pour tout idéal premier $\mathfrak p$ de $A$. Alors
$\text{depth}(A_\mathfrak p) \geq \min(2, \dim(A_\mathfrak p))
\geq \min(2, 4 - \dim(A/\mathfrak p))$, et donc (4) est vérifiée.
\end{proof}

\begin{remark}
\label{algebraization-remark-compare-SGA2}
Dans SGA 2, nous trouvons le résultat suivant. Soit $(A, \mathfrak m)$ un
anneau local noethérien. Soit $f \in \mathfrak m$. Supposons que $A$
soit un quotient d'un anneau régulier, que l'élément
$f$ soit un non-diviseur de zéro et que
\begin{enumerate}
\item[(a)] si $\mathfrak p \subset A$ est un idéal premier tel que
$\dim(A/\mathfrak p) = 1$, alors $\text{depth}(A_\mathfrak p) \geq 2$, et
\item[(b)] $\text{depth}(A/fA) \geq 3$, ou, de manière équivalente,
$\text{depth}(A) \geq 4$.
\end{enumerate}
Soient $U$, resp.\ $U_0$, les spectres épointés de $A$, resp.\ de $A/fA$. Alors
l'application
$$
\Pic(U) \to \Pic(U_0)
$$
est injective. C'est \cite[Exposee XI, lemme 3.16]{SGA2}\footnote{La condition
(a) résulte de la condition (b), voir
Algèbre, lemme \ref{algebra-lemma-depth-localization}.}. Ce résultat
de SGA 2 résulte de la proposition \ref{algebraization-proposition-injective-pic},
car
\begin{enumerate}
\item un quotient d'un anneau régulier possède un complexe dualisant (voir
Complexes dualisants, lemme \ref{dualizing-lemma-regular-gorenstein} et
proposition \ref{dualizing-proposition-dualizing-essentially-finite-type}), et
\item si $\text{depth}(A) \geq 4$, alors $\text{depth}(A_\mathfrak p) \geq 2$
pour tout idéal premier $\mathfrak p$ tel que $\dim(A/\mathfrak p) = 2$, voir
Algèbre, lemme \ref{algebra-lemma-depth-localization}.
\end{enumerate}
\end{remark}






\section{Modules inversibles sur les spectres épointés, II}
\label{algebraization-section-local-lefschetz-for-pic-surjective}

\noindent
Passons maintenant à la surjectivité dans le théorème de Lefschetz local pour le groupe de Picard.
Pour commencer, afin de prolonger un module inversible sur $U_0$ à un voisinage ouvert,
nous disposons du critère simple suivant.

\begin{lemma}
\label{algebraization-lemma-surjective-Pic-first}
Soient $(A, \mathfrak m)$ un anneau local noethérien et $f \in \mathfrak m$.
Supposons que
\begin{enumerate}
\item $A$ soit complet pour la topologie $f$-adique,
\item $f$ soit un non-diviseur de zéro,
\item $H^1_\mathfrak m(A/fA)$ et $H^2_\mathfrak m(A/fA)$
soient des $A$-modules de type fini, et
\item $H^3_\mathfrak m(A/fA) = 0$\footnote{Observons que (3) et (4) sont vérifiées
si $\text{depth}(A/fA) \geq 4$, ou, de manière équivalente, $\text{depth}(A) \geq 5$.}.
\end{enumerate}
Soient $U$, resp.\ $U_0$, les spectres épointés de $A$, resp.\ de $A/fA$.
Alors
$$
\colim_{U_0 \subset U' \subset U\text{ ouvert}} \Pic(U')
\longrightarrow
\Pic(U_0)
$$
est surjective.
\end{lemma}

\begin{proof}
Soit $U_0 \subset U_n \subset U$ le $n$-ième voisinage infinitésimal
de $U_0$. Observons que le faisceau d'idéaux de $U_n$ dans $U_{n + 1}$ est
isomorphe à $\mathcal{O}_{U_0}$, puisque $U_0 \subset U$ est le sous-schéma
fermé principal défini par le non-diviseur de zéro $f$. Nous avons donc
une suite exacte de groupes abéliens
$$
\Pic(U_{n + 1}) \to \Pic(U_n) \to
H^2(U_0, \mathcal{O}_{U_0}) = H^3_\mathfrak m(A/fA) = 0
$$
voir Plus sur les morphismes, lemme
\ref{more-morphisms-lemma-picard-group-first-order-thickening}.
Ainsi, tout $\mathcal{O}_{U_0}$-module inversible est la restriction
d'un module formel cohérent inversible, c'est-à-dire d'un objet inversible de
$\textit{Coh}(U, f\mathcal{O}_U)$. On conclut en appliquant le
lemme \ref{algebraization-lemma-equivalence}.
\end{proof}

\begin{remark}
\label{algebraization-remark-surjective-Pic-second}
Soient $(A, \mathfrak m)$ un anneau local noethérien et $f \in \mathfrak m$.
La conclusion du lemme \ref{algebraization-lemma-surjective-Pic-first} vaut si l'on suppose que
\begin{enumerate}
\item $A$ possède un complexe dualisant,
\item $A$ est complet pour la topologie $f$-adique,
\item $f$ est un non-diviseur de zéro,
\item l'une des conditions suivantes est vérifiée :
\begin{enumerate}
\item $A_f$ est $(S_2)$ et, pour tout $\mathfrak p \subset A$ minimal tel que
$f \not \in \mathfrak p$, on a $\dim(A/\mathfrak p) \geq 4$, ou
\item si $\mathfrak p \not \in V(f)$ et
$V(\mathfrak p) \cap V(f) \not = \{\mathfrak m\}$, alors
$\text{depth}(A_\mathfrak p) + \dim(A/\mathfrak p) > 3$.
\end{enumerate}
\item $H^3_{\mathfrak m}(A/fA) = 0$.
\end{enumerate}
La démonstration est exactement la même que celle du
lemme \ref{algebraization-lemma-surjective-Pic-first},
en utilisant le lemme \ref{algebraization-lemma-equivalence-better} à la place du
lemme \ref{algebraization-lemma-equivalence}.
Il convient de faire ici deux remarques : (a)
il semble difficile de trouver des exemples où l'on sait que
$H^3_{\mathfrak m}(A/fA) = 0$ sans supposer que
$\text{depth}(A/fA) \geq 4$, et
(b) la démonstration du lemme \ref{algebraization-lemma-equivalence-better} est
nettement plus difficile que celle du lemme \ref{algebraization-lemma-equivalence}.
\end{remark}

\begin{lemma}
\label{algebraization-lemma-surjective-Pic-first-better}
Soient $(A, \mathfrak m)$ un anneau local noethérien et $f \in \mathfrak m$.
Supposons que
\begin{enumerate}
\item les conditions du lemme \ref{algebraization-lemma-surjective-Pic-first} soient vérifiées, et
\item pour tout idéal maximal $\mathfrak p \subset A_f$,
le spectre épointé de $(A_f)_\mathfrak p$ ait un groupe de Picard trivial.
\end{enumerate}
Soient $U$, resp.\ $U_0$, les spectres épointés de $A$, resp.\ de $A/fA$.
Alors
$$
\Pic(U) \longrightarrow \Pic(U_0)
$$
est surjective.
\end{lemma}

\begin{proof}
Soit $\mathcal{L}_0 \in \Pic(U_0)$. D'après le
lemme \ref{algebraization-lemma-surjective-Pic-first},
il existe un ouvert $U_0 \subset U' \subset U$
et $\mathcal{L}' \in \Pic(U')$ dont la restriction
à $U_0$ est $\mathcal{L}_0$.
Puisque $U' \supset U_0$, on voit que $U \setminus U'$
est formé de points correspondant à des idéaux premiers
$\mathfrak p_1, \ldots, \mathfrak p_n$ comme en (2).
Par hypothèse, on peut trouver des modules inversibles
$\mathcal{L}'_i$ sur $\Spec(A_{\mathfrak p_i})$ qui coïncident avec
$\mathcal{L}'$ sur le spectre épointé
$U' \times_U \Spec(A_{\mathfrak p_i})$, puisque
les modules inversibles triviaux se prolongent toujours.
D'après Limites, lemme \ref{limits-lemma-glueing-near-closed-point-modules},
appliqué $n$ fois, on voit que $\mathcal{L}'$ se prolonge en un
module inversible sur $U$.
\end{proof}

\begin{lemma}
\label{algebraization-lemma-local-pic-to-completion}
Soit $(A, \mathfrak m)$ un anneau local noethérien de profondeur $\geq 2$.
Soit $A^\wedge$ son complété. Soient $U$, resp.\ $U^\wedge$,
les spectres épointés de $A$, resp.\ de $A^\wedge$. Alors
$\Pic(U) \to \Pic(U^\wedge)$ est injective.
\end{lemma}

\begin{proof}
Soit $\mathcal{L}$ un $\mathcal{O}_U$-module inversible,
et notons $\mathcal{L}^\wedge$ son tiré en arrière sur $U^\wedge$.
On a $H^0(U, \mathcal{O}_U) = A$ par notre hypothèse sur la profondeur et par
Complexes dualisants, lemme \ref{dualizing-lemma-depth}, ainsi que
Cohomologie locale, lemme
\ref{local-cohomology-lemma-finiteness-pushforwards-and-H1-local}.
Ainsi, $\mathcal{L}$ est trivial si et seulement si
$M = H^0(U, \mathcal{L})$ est isomorphe à $A$ comme $A$-module.
(Détails omis.) Puisque $A \to A^\wedge$ est plat,
on a $M \otimes_A A^\wedge = \Gamma(U^\wedge, \mathcal{L}^\wedge)$
par changement de base plat, voir
Cohomologie des schémas, lemme \ref{coherent-lemma-flat-base-change-cohomology}.
Enfin, il est facile de voir que $M \cong A$ si et seulement si
$M \otimes_A A^\wedge \cong A^\wedge$.
\end{proof}

\begin{lemma}
\label{algebraization-lemma-trivial-local-pic-regular}
Soit $(A, \mathfrak m)$ un anneau local régulier. Alors le groupe de Picard
du spectre épointé de $A$ est trivial.
\end{lemma}

\begin{proof}
Combinons Diviseurs, lemme \ref{divisors-lemma-extend-invertible-module},
avec Plus sur l'algèbre, lemme \ref{more-algebra-lemma-regular-local-UFD}.
\end{proof}

\noindent
Nous pouvons maintenant tirer parti des résultats précédents pour démontrer que
les groupes de Picard des spectres épointés
d'intersections complètes de dimension $\geq 4$ sont triviaux.
Rappelons qu'un anneau local noethérien est appelé une intersection
complète si son complété est le quotient d'un
anneau local régulier par l'idéal engendré par une suite régulière.
Voir la discussion dans Algèbre à puissances divisées, section \ref{dpa-section-lci}.

\begin{proposition}[Grothendieck]
\label{algebraization-proposition-trivial-local-pic-complete-intersection}
Soit $(A, \mathfrak m)$ un anneau local noethérien. Si $A$ est une
intersection complète de dimension $\geq 4$, alors le groupe de Picard
du spectre épointé de $A$ est trivial.
\end{proposition}

\begin{proof}
D'après le lemme \ref{algebraization-lemma-local-pic-to-completion}, on peut supposer que $A$ est
un anneau local complet. Par hypothèse, on peut écrire
$A = B/(f_1, \ldots, f_r)$, où $B$ est un anneau local régulier
complet et $f_1, \ldots, f_r$ une suite régulière.
Nous achevons la démonstration par récurrence sur $r$.
Le cas initial $r = 0$ résulte du
lemme \ref{algebraization-lemma-trivial-local-pic-regular}.

\medskip\noindent
Supposons que $A = B/(f_1, \ldots, f_r)$ et que la proposition
soit vraie pour $r - 1$. Posons $A' = B/(f_1, \ldots, f_{r - 1})$ et appliquons le
lemme \ref{algebraization-lemma-surjective-Pic-first-better} à $f_r \in A'$.
C'est possible, car :
\begin{enumerate}
\item la condition (1) du lemme \ref{algebraization-lemma-surjective-Pic-first} est vérifiée
parce que nos anneaux locaux sont complets,
\item la condition (2) du lemme \ref{algebraization-lemma-surjective-Pic-first} est vérifiée
puisque $f_1, \ldots, f_r$ est une suite régulière,
\item les conditions (3) et (4) du lemme \ref{algebraization-lemma-surjective-Pic-first} sont vérifiées
puisque $A = A'/f_r A'$ est de Cohen--Macaulay et de dimension $\dim(A) \geq 4$,
\item la condition (2) du lemme \ref{algebraization-lemma-surjective-Pic-first-better}
est vérifiée par l'hypothèse de récurrence, car
$\dim((A'_{f_r})_\mathfrak p) \geq 4$ pour tout idéal maximal
$\mathfrak p$ de $A'_{f_r}$ et car
$(A'_{f_r})_\mathfrak p = B_\mathfrak q/(f_1, \ldots, f_{r - 1})$
pour un certain idéal premier $\mathfrak q \subset B$, l'anneau $B_\mathfrak q$ étant régulier.
\end{enumerate}
Ceci achève la démonstration.
\end{proof}

\begin{example}
\label{algebraization-example-grothendieck-sharp}
La borne sur la dimension dans la
proposition \ref{algebraization-proposition-trivial-local-pic-complete-intersection}
est optimale. Par exemple, le groupe de Picard du spectre épointé de
$A = k[[x, y, z, w]]/(xy - zw)$ n'est pas trivial. En effet,
l'idéal $I = (x, z)$ définit un diviseur de Cartier effectif $D$ sur le spectre
épointé $U$ de $A$, car il est facile de voir que $I_x, I_y, I_z, I_w$
sont des idéaux inversibles de $A_x, A_y, A_z, A_w$. D'autre part,
$A/I$ est de profondeur $\geq 1$ (en fait $2$), donc $I$ est de profondeur $\geq 2$
(en fait $3$), d'où
$I = \Gamma(U, \mathcal{O}_U(-D))$. Ainsi, si $\mathcal{O}_U(-D)$
était trivial, on aurait $I \cong \Gamma(U, \mathcal{O}_U) = A$,
ce qui est faux puisque $I$ n'est pas engendré par $1$ élément.
\end{example}

\begin{example}
\label{algebraization-example-grothendieck-sharp-bis}
La proposition \ref{algebraization-proposition-trivial-local-pic-complete-intersection}
ne s'étend pas aux quotients
$$
A = B/(f_1, \ldots, f_r)
$$
où $B$ est régulier et $\dim(B) - r \geq 4$. Autrement dit, la condition selon laquelle
$f_1, \ldots, f_r$ est une suite régulière est, en général, nécessaire
à l'annulation du groupe de Picard du spectre épointé de $A$.
En effet, soit $k$ un corps et posons
$$
A = k[[a, b, x, y, z, u, v, w]]/(a^3, b^3, xa^2 + yab + zb^2, w^2)
$$
Observons que $A = A_0[w]/(w^2)$ avec
$A_0 = k[[a, b, x, y, z, u, v]]/(a^3, b^3, xa^2 + yab + zb^2)$.
Nous montrerons ci-dessous que $A_0$ est de profondeur $2$.
Notons $U$ le spectre épointé de $A$ et $U_0$ le spectre épointé
de $A_0$. Observons qu'il existe une suite exacte courte
$0 \to A_0 \to A \to A_0 \to 0$ dont la première flèche est
donnée par la multiplication par $w$.
D'après Plus sur les morphismes, lemme
\ref{more-morphisms-lemma-picard-group-first-order-thickening},
on obtient une suite exacte
$$
H^0(U, \mathcal{O}_U^*) \to
H^0(U_0, \mathcal{O}_{U_0}^*) \to
H^1(U_0, \mathcal{O}_{U_0}) \to
\Pic(U)
$$
Comme $A_0$, et donc $A$, est de profondeur $2$, on voit que
$H^0(U_0, \mathcal{O}_{U_0}) = A_0$ et $H^0(U, \mathcal{O}_U) = A$,
et que $H^1(U_0, \mathcal{O}_{U_0})$ est non nul, voir
Complexes dualisants, lemme \ref{dualizing-lemma-depth}, et
Cohomologie locale, lemme \ref{local-cohomology-lemma-local-cohomology}.
Ainsi, la dernière flèche affichée ci-dessus est non nulle, et l'on conclut
que $\Pic(U)$ n'est pas nul.

\medskip\noindent
Pour montrer que $A_0$ est de profondeur $2$, il suffit de montrer que
$A_1 = k[[a, b, x, y, z]]/(a^3, b^3, xa^2 + yab + zb^2)$
est de profondeur $0$. En effet, l'image de $a^2b^2$ est
un élément non nul de $A_1$ qui est annulé
par chacune des variables $a, b, x, y, z$. Par exemple,
$ya^2b^2 = (yab)(ab) = - (xa^2 + zb^2)(ab) = -xa^3b - yab^3 = 0$ dans $A_1$.
Les autres cas sont analogues.
\end{example}










\section{Application aux théorèmes de Lefschetz}
\label{algebraization-section-lefschetz}

\noindent
Dans cette section, nous étudions la relation entre les faisceaux cohérents sur un
schéma projectif $P$ et les modules cohérents sur la complétion formelle le long
d'un diviseur ample $Q$.

\medskip\noindent
Soit $k$ un corps. Soit $P$ un schéma propre sur $k$.
Soit $\mathcal{L}$ un $\mathcal{O}_P$-module inversible ample.
Soit $s \in \Gamma(P, \mathcal{L})$ une section\footnote{Nous n'exigeons pas
que $s$ soit une section régulière. Par conséquent, $Q$ est seulement
un sous-schéma fermé localement principal de $P$, et non nécessairement un diviseur
de Cartier effectif.}, et soit
$Q = Z(s)$ son schéma des zéros, voir
Diviseurs, définition \ref{divisors-definition-zero-scheme-s}.
Pour tout $n \geq 1$, nous notons $Q_n = Z(s^n)$
le $n$-ième voisinage infinitésimal de $Q$.
Si $\mathcal{F}$ est un $\mathcal{O}_P$-module cohérent, nous notons
$\mathcal{F}_n = \mathcal{F}|_{Q_n}$ sa restriction, c'est-à-dire
le tiré en arrière de $\mathcal{F}$ par l'immersion fermée
$Q_n \to P$.

\begin{proposition}
\label{algebraization-proposition-lefschetz}
Dans la situation ci-dessus, supposons qu'il existe un entier $\sigma$ tel que
pour tout point $p \in P \setminus Q$, on ait
$$
\text{depth}(\mathcal{F}_p) + \dim(\overline{\{p\}}) > \sigma
$$
Alors l'application
$$
H^i(P, \mathcal{F}) \longrightarrow \lim H^i(Q_n, \mathcal{F}_n)
$$
est un isomorphisme pour $0 \leq i < \sigma$.
\end{proposition}

\begin{proof}
Nous utiliserons Plus sur les morphismes, lemme \ref{more-morphisms-lemma-apply-proj-spec},
ainsi que les notations et les résultats de
Plus sur les morphismes, section \ref{more-morphisms-section-proj-spec},
sans plus de précisions ; cette démonstration ne peut être comprise sans
au moins comprendre l'énoncé du lemme.
Observons que, dans notre cas,
$A = \bigoplus_{m \geq 0} \Gamma(P, \mathcal{L}^{\otimes m})$
est une $k$-algèbre de type fini
dont toutes les composantes graduées sont des $k$-espaces vectoriels de dimension finie, voir
Cohomologie des schémas, lemme \ref{coherent-lemma-coherent-proper-ample}.

\medskip\noindent
Nous considérerons désormais $s$ comme un élément $f \in A_1 \subset A$, c'est-à-dire
comme un élément homogène de degré $1$ de $A$. Notons
$Y = V(f) \subset X$ le sous-schéma fermé défini par $f$.
Alors $U \cap Y = (\pi|_U)^{-1}(Q)$ au sens schématique.
Rappelons la notation
$\mathcal{F}_U = \pi^*\mathcal{F}|_U = (\pi|_U)^*\mathcal{F}$.
Il s'agit d'un $\mathcal{O}_U$-module cohérent.
Choisissons un $A$-module de type fini $M$ tel que
$\mathcal{F}_U = \widetilde{M}|_U$
(pour l'existence, voir Cohomologie locale, lemme
\ref{local-cohomology-lemma-finiteness-pushforwards-and-H1-local}).
Nous affirmons que $H^i_Z(M)$ est annulé par
une puissance de $f$ pour $i \leq \sigma + 1$.

\medskip\noindent
Pour démontrer cette assertion, nous appliquerons
Cohomologie locale, proposition \ref{local-cohomology-proposition-annihilator}.
En termes géométriques, on voit qu'il suffit de démontrer,
pour $u \in U$, $u \not \in Y$ et $z \in \overline{\{u\}} \cap Z$,
que
$$
\text{depth}(\mathcal{F}_{U, u}) +
\dim(\mathcal{O}_{\overline{\{u\}}, z}) > \sigma + 1
$$
Cela ne demande qu'un bref raisonnement.

\medskip\noindent
Observons que $Z = \Spec(A_0)$ est un ensemble fini de points fermés de $X$, car
$A_0$ est une $k$-algèbre de dimension finie.
(Le lecteur qui souhaiterait que $Z$ soit un singleton peut remplacer la
$k$-algèbre finie $A_0$ par $k$ ; cela ne change rien d'autre à la démonstration.)

\medskip\noindent
Le morphisme $\pi : L \to P$ et sa restriction $\pi|_U : U \to P$
sont lisses de dimension relative $1$.
Soient $u \in U$, $u \not \in Y$ et $z \in \overline{\{u\}} \cap Z$.
Soit $p = \pi(u) \in P \setminus Q$ son image.
Alors $u$ est soit un point générique
de la fibre de $\pi$ au-dessus de $p$, soit un point fermé de cette fibre.
Si $u$ est un point générique de la fibre, alors
$\text{depth}(\mathcal{F}_{U, u}) = \text{depth}(\mathcal{F}_p)$
et
$\dim(\overline{\{u\}}) = \dim(\overline{\{p\}}) + 1$.
Si $u$ est un point fermé de la fibre, alors
$\text{depth}(\mathcal{F}_{U, u}) = \text{depth}(\mathcal{F}_p) + 1$
et
$\dim(\overline{\{u\}}) = \dim(\overline{\{p\}})$.
Dans les deux cas, on a
$\dim(\overline{\{u\}}) = \dim(\mathcal{O}_{\overline{\{u\}}, z})$
parce que tout point de $Z$ est fermé. L'inégalité voulue
résulte donc de l'hypothèse de l'énoncé de la
proposition.

\medskip\noindent
Soit $A'$ le complété $f$-adique de $A$. Alors $A \to A'$ est plat d'après
Algèbre, lemme \ref{algebra-lemma-completion-flat}.
Notons $U' \subset X' = \Spec(A')$ l'image réciproque de
$U$, et de même pour $Y'$ et $Z'$. Soit $\mathcal{F}'$
sur $U'$ le tiré en arrière de $\mathcal{F}_U$, et posons
$M' = M \otimes_A A'$.
Par changement de base plat pour la cohomologie locale
(Cohomologie locale, lemme \ref{local-cohomology-lemma-torsion-change-rings}),
on a
$$
H^i_{Z'}(M') = H^i_Z(M) \otimes_A A'
$$
et l'on constate que, pour $i \leq \sigma + 1$,
ces groupes sont annulés par une puissance de $f$. Considérons le diagramme
$$
\xymatrix{
& H^i(U, \mathcal{F}_U) \ar[ld] \ar[d] \ar[r] &
\lim H^i(U, \mathcal{F}_U/f^n\mathcal{F}_U) \ar@{=}[d] \\
H^i(U, \mathcal{F}_U) \otimes_A A' \ar@{=}[r] &
H^i(U', \mathcal{F}') \ar[r] &
\lim H^i(U', \mathcal{F}'/f^n\mathcal{F}')
}
$$
La flèche horizontale inférieure est un isomorphisme pour $i < \sigma$ d'après le
lemme \ref{algebraization-lemma-alternative-higher} et la propriété de torsion
que nous venons de démontrer. Le signe d'égalité horizontal provient du changement de base plat
(Cohomologie des schémas, lemme \ref{coherent-lemma-flat-base-change-cohomology}),
et le signe d'égalité vertical vient de ce que $U \cap Y$ et $U' \cap Y'$
ainsi que leurs $n$-ièmes voisinages infinitésimaux
sont identifiés deux à deux par des isomorphismes (puisque nous
complétons par rapport à $f$).

\medskip\noindent
En appliquant Plus sur les morphismes, équation
(\ref{more-morphisms-equation-cohomology-torsor}),
on obtient des décompositions en somme directe compatibles
$$
\lim H^i(U, \mathcal{F}_U/f^n\mathcal{F}_U) =
\lim
\left(
\bigoplus\nolimits_{m \in \mathbf{Z}}
H^i(Q_n, \mathcal{F}_n \otimes \mathcal{L}^{\otimes m})
\right)
$$
et
$$
H^i(U, \mathcal{F}_U) =
\bigoplus\nolimits_{m \in \mathbf{Z}}
H^i(P, \mathcal{F} \otimes \mathcal{L}^{\otimes m})
$$
On conclut donc par Algèbre, lemme \ref{algebra-lemma-daniel-litt}.
\end{proof}

\begin{lemma}
\label{algebraization-lemma-lefschetz-addendum}
Soit $k$ un corps. Soit $X$ un schéma propre sur $k$.
Soit $\mathcal{L}$ un $\mathcal{O}_X$-module inversible ample.
Soit $s \in \Gamma(X, \mathcal{L})$. Soit $Y = Z(s)$ le
schéma des zéros de $s$, dont le $n$-ième voisinage infinitésimal est $Y_n = Z(s^n)$.
Soit $\mathcal{F}$ un $\mathcal{O}_X$-module cohérent.
Supposons que, pour tout $x \in X \setminus Y$, on ait
$$
\text{depth}(\mathcal{F}_x) + \dim(\overline{\{x\}}) > 1
$$
Alors $\Gamma(V, \mathcal{F}) \to \lim \Gamma(Y_n, \mathcal{F}|_{Y_n})$
est un isomorphisme pour tout sous-schéma ouvert $V \subset X$ contenant $Y$.
\end{lemma}

\begin{proof}
D'après la proposition \ref{algebraization-proposition-lefschetz}, c'est vrai pour $V = X$.
Il suffit donc de montrer que l'application
$\Gamma(V, \mathcal{F}) \to \lim \Gamma(Y_n, \mathcal{F}|_{Y_n})$
est injective. Si $\sigma \in \Gamma(V, \mathcal{F})$
s'envoie sur zéro, son support est disjoint de $Y$ (détails omis ; indication :
utiliser le théorème d'intersection de Krull). Alors l'adhérence $T \subset X$
de $\text{Supp}(\sigma)$ est disjointe de $Y$.
Ainsi, $T$ est propre sur $k$ (car fermé dans $X$)
et affine (car fermé dans le schéma affine $X \setminus Y$, voir
Morphismes, lemme \ref{morphisms-lemma-proper-ample-delete-affine}),
donc fini sur $k$
(Morphismes, lemme \ref{morphisms-lemma-finite-proper}).
Ainsi, $T$ est un ensemble fini de points fermés de $X$.
L'hypothèse donne $\text{depth}(\mathcal{F}_x) \geq 2$, donc a fortiori au moins $1$,
pour $x \in T$. On conclut que
$\Gamma(V, \mathcal{F}) \to \Gamma(V \setminus T, \mathcal{F})$
est injective et que $\sigma = 0$, ce qu'il fallait démontrer.
\end{proof}


\begin{example}
\label{algebraization-example-lefschetz}
Soit $k$ un corps et soit $X$ une variété propre sur $k$.
Soit $Y \subset X$ un diviseur de Cartier effectif tel que
$\mathcal{O}_X(Y)$ soit ample, et
notons $Y_n$ son $n$-ième voisinage infinitésimal.
Soit $\mathcal{E}$ un $\mathcal{O}_X$-module localement libre de type fini.
Voici quelques cas particuliers de la proposition \ref{algebraization-proposition-lefschetz}.
\begin{enumerate}
\item Si $X$ est une courbe, on n'obtient aucun renseignement.
\item Si $X$ est une surface de Cohen--Macaulay (par exemple normale), alors
$$
H^0(X, \mathcal{E}) \to \lim H^0(Y_n, \mathcal{E}|_{Y_n})
$$
est un isomorphisme.
\item Si $X$ est une variété tridimensionnelle de Cohen--Macaulay, alors
$$
H^0(X, \mathcal{E}) \to \lim H^0(Y_n, \mathcal{E}|_{Y_n})
\quad\text{et}\quad
H^1(X, \mathcal{E}) \to \lim H^1(Y_n, \mathcal{E}|_{Y_n})
$$
sont des isomorphismes.
\end{enumerate}
Le principe général est sans doute clair. Si $X$ est une variété tridimensionnelle normale, alors
on peut conclure pour $H^0$, mais pas pour $H^1$.
\end{example}

\noindent
Avant de démontrer le résultat principal suivant, nous avons besoin d'un lemme.

\begin{lemma}
\label{algebraization-lemma-Gm-equivariant-extend-canonically}
Dans la situation \ref{algebraization-situation-algebraize}, soit $(\mathcal{F}_n)$ un
objet de $\textit{Coh}(U, I\mathcal{O}_U)$. Supposons que
\begin{enumerate}
\item $A$ soit un anneau gradué, $\mathfrak a = A_+$, et
$I$ un idéal homogène,
\item $(\mathcal{F}_n) = (\widetilde{M_n}|_U)$, où $(M_n)$
est un système projectif de $A$-modules gradués, et
\item $(\mathcal{F}_n)$ s'étende canoniquement à $X$.
\end{enumerate}
Alors il existe un $A$-module gradué de type fini $N$ tel que
\begin{enumerate}
\item[(a)] les systèmes projectifs $(N/I^nN)$ et $(M_n)$ soient pro-isomorphes
dans la catégorie quotient des $A$-modules gradués par les modules annulés par une puissance de $A_+$,
et
\item[(b)] $(\mathcal{F}_n)$ soit le complété du module cohérent
associé à $N$.
\end{enumerate}
\end{lemma}

\begin{proof}
Soit $(\mathcal{G}_n)$ le prolongement canonique du
lemme \ref{algebraization-lemma-canonically-algebraizable}.
Les graduations de $A$ et de $M_n$ déterminent une action
$$
a : \mathbf{G}_m \times X \longrightarrow X
$$
du schéma en groupes $\mathbf{G}_m$ sur $X$ telle que
$(\widetilde{M_n})$ devienne un système projectif de
modules quasi-cohérents $\mathbf{G}_m$-équivariants sur $\mathcal{O}_X$, voir
Groupoïdes, exemple \ref{groupoids-example-Gm-on-affine}.
Puisque $\mathfrak a$ et $I$ sont des idéaux homogènes,
les sous-schémas $Z$ et $Y$ sont fermés et le sous-schéma $U$ est ouvert ;
tous trois sont des sous-schémas $\mathbf{G}_m$-invariants.
La restriction $(\mathcal{F}_n)$ de $(\widetilde{M_n})$
est un système projectif de modules cohérents $\mathbf{G}_m$-équivariants
qui sont des $\mathcal{O}_U$-modules. Autrement dit, $(\mathcal{F}_n)$
est un module formel cohérent $\mathbf{G}_m$-équivariant,
au sens où il existe un isomorphisme
$$
\alpha : (a^*\mathcal{F}_n) \longrightarrow (p^*\mathcal{F}_n)
$$
sur $\mathbf{G}_m \times U$ satisfaisant à une condition de cocycle convenable.
Comme $a$ et $p$ sont des morphismes plats de schémas affines,
le lemme \ref{algebraization-lemma-canonically-extend-base-change}
montre qu'il existe un unique isomorphisme
$$
\beta : (a^*\mathcal{G}_n) \longrightarrow (p^*\mathcal{G}_n)
$$
sur $\mathbf{G}_m \times X$ dont la restriction est $\alpha$ sur
$\mathbf{G}_m \times U$. L'unicité garantit que
$\beta$ satisfait à la condition de cocycle correspondante.
De cette manière, chaque $\mathcal{G}_n$ devient
un module cohérent $\mathbf{G}_m$-équivariant sur $\mathcal{O}_X$,
de façon compatible avec les applications de transition.

\medskip\noindent
D'après Groupoïdes, lemme \ref{groupoids-lemma-Gm-equivariant-module},
on voit que $\mathcal{G}_n$, muni de sa structure $\mathbf{G}_m$-équivariante,
correspond à un $A$-module gradué $N_n$. Les applications de transition
$N_{n + 1} \to N_n$ sont des homomorphismes de modules gradués. Notons que $N_n$ est un
$A$-module de type fini et que $N_n = N_{n + 1}/I^n N_{n + 1}$, puisque
$(\mathcal{G}_n)$ est un objet de $\textit{Coh}(X, I\mathcal{O}_X)$.
Soit $N$ le $A$-module gradué de type fini fourni par
Algèbre, lemme \ref{algebra-lemma-finiteness-graded}.
Alors $N_n = N/I^nN$, de sorte que $(\mathcal{G}_n)$
est le complété du module cohérent
associé à $N$, et l'on voit a fortiori que (b) est vraie.

\medskip\noindent
Pour voir (a), il faut expliciter un peu davantage la situation ci-dessus.
Observons d'abord que le noyau et le conoyau de $M_n \to H^0(U, \mathcal{F}_n)$
sont de torsion annulée par une puissance de $A_+$ (Cohomologie locale, lemme
\ref{local-cohomology-lemma-finiteness-pushforwards-and-H1-local}).
Observons que $H^0(U, \mathcal{F}_n)$ est muni d'une graduation naturelle
pour laquelle ces applications et les applications de transition du système sont
des homomorphismes gradués de $A$-modules ; on peut par exemple utiliser le fait que
$(U \to X)_*\mathcal{F}_n$ est un module $\mathbf{G}_m$-équivariant
sur $X$ et appliquer 
Groupoïdes, lemme \ref{groupoids-lemma-Gm-equivariant-module}.
Rappelons ensuite que $(N_n)$ et $(H^0(U, \mathcal{F}_n))$
sont pro-isomorphes d'après la définition \ref{algebraization-definition-canonically-algebraizable}
et le lemme \ref{algebraization-lemma-canonically-algebraizable}.
Nous omettons la vérification que les applications définissant ce
pro-isomorphisme sont des homomorphismes de modules gradués.
Ainsi, $(N_n)$ et $(M_n)$ sont pro-isomorphes dans la catégorie quotient des
$A$-modules gradués par les modules de torsion annulés par une puissance de $A_+$.
\end{proof}

\noindent
Soit $k$ un corps. Soit $P$ un schéma propre sur $k$.
Soit $\mathcal{L}$ un $\mathcal{O}_P$-module inversible ample.
Soit $s \in \Gamma(P, \mathcal{L})$ une section et soit
$Q = Z(s)$ le schéma des zéros, voir
Diviseurs, définition \ref{divisors-definition-zero-scheme-s}.
Soit $\mathcal{I} \subset \mathcal{O}_P$ le faisceau d'idéaux de $Q$.
Nous noterons $\textit{Coh}(P, \mathcal{I})$ la catégorie
des modules formels cohérents introduite dans
Cohomologie des schémas, section \ref{coherent-section-coherent-formal}.

\begin{proposition}
\label{algebraization-proposition-lefschetz-existence}
Dans la situation ci-dessus, soit $(\mathcal{F}_n)$ un objet de
$\textit{Coh}(P, \mathcal{I})$. Supposons que, pour tout $q \in Q$ et tout
idéal premier $\mathfrak p \in \mathcal{O}_{P, q}^\wedge$ tel que
$\mathfrak p \not \in V(\mathcal{I}_q^\wedge)$, on ait
$$
\text{depth}((\mathcal{F}_q^\wedge)_\mathfrak p) +
\dim(\mathcal{O}_{P, q}^\wedge/\mathfrak p) +
\dim(\overline{\{q\}}) > 2
$$
Alors $(\mathcal{F}_n)$ est le complété d'un
$\mathcal{O}_P$-module cohérent.
\end{proposition}

\begin{proof}
D'après Cohomologie des schémas, lemme \ref{coherent-lemma-existence-easy},
pour démontrer l'énoncé, on peut remplacer $(\mathcal{F}_n)$ par un objet
qui n'en diffère que par de la $\mathcal{I}$-torsion (voir ci-dessous pour plus de précision).
Posons $T' = \{q \in Q \mid \dim(\overline{\{q\}}) = 0\}$
et $T = \{q \in Q \mid \dim(\overline{\{q\}}) \leq 1\}$.
L'hypothèse de la proposition signifie exactement que
$Q \subset P$, $(\mathcal{F}_n)$ et $T' \subset T \subset Q$
satisfont aux conditions du
lemme \ref{algebraization-lemma-improvement-formal-coherent-module-better}
avec $d = 1$ ; outre des manipulations immédiates d'inégalités, on utilise que
$V(\mathfrak p) \cap V(\mathcal{I}^\wedge_y) = \{\mathfrak m^\wedge_y\}
\Leftrightarrow \dim(\mathcal{O}_{P, q}^\wedge/\mathfrak p) = 1$
puisque $\mathcal{I}_y^\wedge$ est engendré par $1$ élément.
En combinant ces deux remarques, on peut remplacer $(\mathcal{F}_n)$ par
l'objet $(\mathcal{H}_n)$ de $\textit{Coh}(P, \mathcal{I})$ fourni par le
lemme \ref{algebraization-lemma-improvement-formal-coherent-module-better}.
Ainsi, on peut supposer, et on suppose, que $(\mathcal{F}_n)$ est pro-isomorphe à
un système projectif $(\mathcal{F}_n'')$ de $\mathcal{O}_P$-modules cohérents
tel que $\text{depth}(\mathcal{F}''_{n, q}) + \dim(\overline{\{q\}}) \geq 2$
pour tout $q \in Q$.

\medskip\noindent
Nous utiliserons Plus sur les morphismes, lemme \ref{more-morphisms-lemma-apply-proj-spec},
ainsi que les notations et les résultats de
Plus sur les morphismes, section \ref{more-morphisms-section-proj-spec},
sans plus de précisions ; cette démonstration ne peut être comprise sans
au moins comprendre l'énoncé du lemme.
Observons que, dans notre cas,
$A = \bigoplus_{m \geq 0} \Gamma(P, \mathcal{L}^{\otimes m})$
est une $k$-algèbre de type fini
dont toutes les composantes graduées sont des $k$-espaces vectoriels de dimension finie, voir
Cohomologie des schémas, lemme \ref{coherent-lemma-coherent-proper-ample}.

\medskip\noindent
D'après Cohomologie des schémas, lemme \ref{coherent-lemma-inverse-systems-pullback},
le tiré en arrière par $\pi|_U : U \to P$ est un objet
$(\pi|_U^*\mathcal{F}_n)$ de $\textit{Coh}(U, f\mathcal{O}_U)$
qui est pro-isomorphe au système projectif
$(\pi|_U^*\mathcal{F}_n'')$ de $\mathcal{O}_U$-modules cohérents.
Nous affirmons que
$$
\text{depth}(\pi|_U^*\mathcal{F}''_{n, y}) + \delta_Z^Y(y) \geq 3
$$
pour tout $y \in U \cap Y$. Puisque tous les points de $Z$ sont fermés, on
voit que $\delta_Z^Y(y) \geq \dim(\overline{\{y\}})$ pour tout
$y \in U \cap Y$, voir le lemme \ref{algebraization-lemma-discussion}.
Soit $q \in Q$ l'image de $y$. Comme le morphisme $\pi : U \to P$ est
lisse de dimension relative $1$, on voit que $y$ est soit un point fermé
d'une fibre de $\pi$, soit un point générique.
On en déduit que
$$
\text{depth}(\pi^*\mathcal{F}''_{n, y}) + \delta_Z^Y(y)
\geq
\text{depth}(\pi^*\mathcal{F}''_{n, y}) + \dim(\overline{\{y\}}) =
\text{depth}(\mathcal{F}''_{n, q}) + \dim(\overline{\{q\}}) + 1
$$
car soit la profondeur, soit la dimension augmente de $1$.
Ceci démontre l'assertion.

\medskip\noindent
D'après le lemme \ref{algebraization-lemma-cd-1-canonical}, on conclut que
$(\pi|_U^*\mathcal{F}_n)$ se prolonge canoniquement sur $X$.
Observons que
$$
M_n = \Gamma(U, \pi|_U^*\mathcal{F}_n) =
\bigoplus\nolimits_{m \in \mathbf{Z}}
\Gamma(P, \mathcal{F}_n \otimes_{\mathcal{O}_P} \mathcal{L}^{\otimes m})
$$
est canoniquement un $A$-module gradué, voir
Plus sur les morphismes, équation (\ref{more-morphisms-equation-cohomology-torsor}).
D'après Propriétés, lemme \ref{properties-lemma-quasi-coherent-quasi-affine},
on a $\pi|_U^*\mathcal{F}_n = \widetilde{M_n}|_U$.
On peut donc appliquer le lemme \ref{algebraization-lemma-Gm-equivariant-extend-canonically}
pour trouver un $A$-module gradué de type fini $N$ tel que
$(M_n)$ et $(N/I^nN)$ soient pro-isomorphes dans la catégorie quotient
des $A$-modules gradués par les modules de $A_+$-torsion.
Soit $\mathcal{F}$ le $\mathcal{O}_P$-module cohérent
associé à $N$, voir
Cohomologie des schémas, proposition
\ref{coherent-proposition-coherent-modules-on-proj-general}.
La même proposition montre que $(\mathcal{F}/\mathcal{I}^n\mathcal{F})$
est pro-isomorphe à $(\mathcal{F}_n)$.
Puisque tous deux sont des objets de $\textit{Coh}(P, \mathcal{I})$,
on conclut par le lemme \ref{algebraization-lemma-recognize-formal-coherent-modules}.
\end{proof}

\begin{example}
\label{algebraization-example-lefschetz-existence}
Soit $k$ un corps et soit $X$ une variété propre sur $k$.
Soit $Y \subset X$ un diviseur de Cartier effectif tel que
$\mathcal{O}_X(Y)$ soit ample, et notons
$\mathcal{I} \subset \mathcal{O}_X$ le faisceau d'idéaux
correspondant. Soit $(\mathcal{E}_n)$ un objet
de $\textit{Coh}(X, \mathcal{I})$ tel que $\mathcal{E}_n$ soit
localement libre de type fini.
Voici quelques cas particuliers de la
proposition \ref{algebraization-proposition-lefschetz-existence}.
\begin{enumerate}
\item Si $X$ est une courbe ou une surface, on n'obtient aucun renseignement.
\item Si $X$ est une variété tridimensionnelle de Cohen--Macaulay, alors
$(\mathcal{E}_n)$ est le complété d'un
$\mathcal{O}_X$-module cohérent $\mathcal{E}$.
\item Plus généralement, si $\dim(X) \geq 3$ et si $X$ est $(S_3)$, alors
$(\mathcal{E}_n)$ est le complété d'un
$\mathcal{O}_X$-module cohérent $\mathcal{E}$.
\end{enumerate}
Bien sûr, si $\mathcal{E}$ existe, alors $\mathcal{E}$ est localement libre
de type fini dans un voisinage ouvert de $Y$.
\end{example}

\begin{proposition}
\label{algebraization-proposition-lefschetz-equivalence}
Soit $k$ un corps. Soit $X$ un schéma propre sur $k$.
Soit $\mathcal{L}$ un $\mathcal{O}_X$-module inversible ample,
et soit $s \in \Gamma(X, \mathcal{L})$. Soit $Y = Z(s)$
le schéma des zéros de $s$, et notons $\mathcal{I} \subset \mathcal{O}_X$
le faisceau d'idéaux correspondant.
Soit $\mathcal{V}$ l'ensemble des sous-schémas ouverts de $X$ contenant $Y$,
ordonné par l'inclusion inverse.
Supposons que, pour tout $x \in X \setminus Y$, on ait
$$
\text{depth}(\mathcal{O}_{X, x}) + \dim(\overline{\{x\}}) > 2
$$
Alors le foncteur de complétion
$$
\colim_\mathcal{V}
\textit{Coh}(\mathcal{O}_V)
\longrightarrow
\textit{Coh}(X, \mathcal{I})
$$
induit une équivalence entre les sous-catégories pleines d'objets localement libres de type fini.
\end{proposition}

\begin{proof}
Pour démontrer la pleine fidélité, il suffit de montrer que
$$
\colim_\mathcal{V} \Gamma(V, \mathcal{L}^{\otimes m})
\longrightarrow
\lim \Gamma(Y_n, \mathcal{L}^{\otimes m}|_{Y_n})
$$
est un isomorphisme pour tout $m$, voir le
lemme \ref{algebraization-lemma-completion-fully-faithful-general}.
Ceci résulte du lemme \ref{algebraization-lemma-lefschetz-addendum}.

\medskip\noindent
Surjectivité essentielle. Soit $(\mathcal{F}_n)$ un objet localement
libre de type fini de $\textit{Coh}(X, \mathcal{I})$. Alors, pour $y \in Y$,
$\mathcal{F}_y^\wedge = \lim \mathcal{F}_{n, y}$
est un $\mathcal{O}_{X, y}^\wedge$-module libre de type fini.
Soit $\mathfrak p \subset \mathcal{O}_{X, y}^\wedge$
un idéal premier tel que $\mathfrak p \not \in V(\mathcal{I}_y^\wedge)$.
Alors $\mathfrak p$ est au-dessus d'un idéal premier $\mathfrak p_0 \subset \mathcal{O}_{X, y}$
qui correspond à une spécialisation $x \leadsto y$ avec
$x \not \in Y$. D'après Cohomologie locale, lemme
\ref{local-cohomology-lemma-change-completion}
et des résultats de théorie de la dimension
(voir Variétés, section \ref{varieties-section-algebraic-schemes}),
on a
$$
\text{depth}((\mathcal{O}_{X, y}^\wedge)_\mathfrak p) +
\dim(\mathcal{O}_{X, y}^\wedge/\mathfrak p) =
\text{depth}(\mathcal{O}_{X, x}) +
\dim(\overline{\{x\}}) - \dim(\overline{\{y\}})
$$
Ainsi, nos hypothèses impliquent que les hypothèses de la
proposition \ref{algebraization-proposition-lefschetz-existence}
sont satisfaites, et l'on obtient que $(\mathcal{F}_n)$
est le complété d'un $\mathcal{O}_X$-module cohérent $\mathcal{F}$.
Il s'ensuit que $\mathcal{F}_y$ est libre de type fini pour tout $y \in Y$,
et donc que $\mathcal{F}$ est localement libre de type fini dans un
voisinage ouvert $V$ de $Y$. Ceci achève la démonstration.
\end{proof}













% OK, start here.
%

\chapter{Courbes algébriques}



\phantomsection
\label{curves-section-phantom}


\section{Introduction}
\label{curves-section-introduction}

\noindent
Dans ce chapitre, nous développons quelques éléments de la théorie des courbes algébriques. Pour les courbes algébriques sur le corps des nombres complexes, on pourra consulter l'ouvrage \cite{ACGH}.

\medskip\noindent
Rappel des résultats déjà établis. Outre les résultats généraux de géométrie algébrique, nous avons déjà démontré plusieurs résultats propres aux courbes algébriques. En voici la liste. \begin{enumerate} \item Nous avons étudié les ouverts affines et les faisceaux inversibles amples sur les schémas noethériens de dimension $1$ dans Variétés, section \ref{varieties-section-dimension-one}. \item Nous avons vu qu'une courbe est soit affine, soit projective dans Variétés, section \ref{varieties-section-curves}. \item Nous avons étudié les degrés des modules localement libres sur les courbes propres dans Variétés, section \ref{varieties-section-divisors-curves}. \item Nous avons étudié le schéma de Picard d'une courbe projective non singulière sur un corps algébriquement clos dans Schémas de Picard des courbes, section \ref{pic-section-introduction}.
\end{enumerate}





\section{Courbes et corps de fonctions}
\label{curves-section-curves-function-fields}

\noindent
Dans cette section, nous précisons les résultats de Variétés, section \ref{varieties-section-varieties-rational-maps} dans le cas des courbes.

\begin{lemma}
\label{curves-lemma-extend-over-dvr}
Soit $k$ un corps. Soient $X$ une courbe et $Y$ une variété propre. Soient $U \subset X$ un ouvert non vide et $f : U \to Y$ un morphisme. Si $x \in X$ est un point fermé tel que $\mathcal{O}_{X, x}$ soit un anneau de valuation discrète, il existe un ouvert $U \subset U' \subset X$ contenant $x$ et un morphisme de variétés $f' : U' \to Y$ prolongeant $f$.
\end{lemma}

\begin{proof}
C'est un cas particulier de Morphismes, lemme \ref{morphisms-lemma-extend-across}.
\end{proof}

\begin{lemma}
\label{curves-lemma-extend-over-normal-curve}
Soit $k$ un corps. Soient $X$ une courbe normale et $Y$ une variété propre. L'ensemble des applications rationnelles de $X$ dans $Y$ s'identifie à l'ensemble des morphismes $X \to Y$.
\end{lemma}

\begin{proof}
Toute application rationnelle de $X$ dans $Y$ se prolonge en un morphisme $X \to Y$ d'après le lemme \ref{curves-lemma-extend-over-dvr}, puisque tout anneau local est un anneau de valuation discrète (par exemple d'après Variétés, lemme \ref{varieties-lemma-regular-point-on-curve}). Réciproquement, si deux morphismes $f,g: X \to Y$ sont équivalents en tant qu'applications rationnelles, alors $f = g$ d'après Morphismes, lemme \ref{morphisms-lemma-equality-of-morphisms}.
\end{proof}

\begin{lemma}
\label{curves-lemma-flat}
Soit $k$ un corps. Soit $f : X \to Y$ un morphisme non constant de courbes sur $k$. Si $Y$ est normale, alors $f$ est plat.
\end{lemma}

\begin{proof}
Prenons $x \in X$ d'image $y \in Y$. Alors $\mathcal{O}_{Y, y}$ est soit un corps, soit un anneau de valuation discrète (Variétés, lemme \ref{varieties-lemma-regular-point-on-curve}). Puisque $f$ est non constant, il est dominant (car il envoie nécessairement le point générique de $X$ sur celui de $Y$). Il en résulte que $\mathcal{O}_{Y, y} \to \mathcal{O}_{X, x}$ est injectif (Morphismes, lemme \ref{morphisms-lemma-dominant-between-integral}). Ainsi $\mathcal{O}_{X, x}$ est sans torsion en tant que $\mathcal{O}_{Y, y}$-module; par conséquent, $\mathcal{O}_{X, x}$ est plat en tant que $\mathcal{O}_{Y, y}$-module d'après Compléments d'algèbre, lemme \ref{more-algebra-lemma-valuation-ring-torsion-free-flat}.
\end{proof}

\begin{lemma}
\label{curves-lemma-finite}
Soit $k$ un corps. Soit $f : X \to Y$ un morphisme de schémas sur $k$. Supposons que \begin{enumerate} \item $Y$ soit séparé sur $k$, \item $X$ soit propre de dimension $\leq 1$ sur $k$, \item $f(Z)$ contienne au moins deux points pour toute composante irréductible $Z \subset X$ de dimension $1$. \end{enumerate} Alors $f$ est fini.
\end{lemma}

\begin{proof}
Le morphisme $f$ est propre d'après Morphismes, lemme \ref{morphisms-lemma-image-proper-scheme-closed}. Ainsi $f(X)$ est fermé et les images des points fermés sont fermées. Soit $y \in Y$ l'image d'un point fermé de $X$. Alors $f^{-1}(\{y\})$ est une partie fermée de $X$ ne contenant aucun point générique d'une composante irréductible de dimension $1$, d'après la condition (3). Il s'ensuit que $f^{-1}(\{y\})$ est fini. Par conséquent, $f$ est fini au-dessus d'un voisinage ouvert de $y$ d'après Compléments sur les morphismes, lemme \ref{more-morphisms-lemma-proper-finite-fibre-finite-in-neighbourhood} (si $Y$ est noethérien, on peut utiliser le résultat plus élémentaire Cohomologie des schémas, lemme \ref{coherent-lemma-proper-finite-fibre-finite-in-neighbourhood}). Puisque nous avons vu qu'il y a suffisamment de tels points $y$, la démonstration est achevée.
\end{proof}

\begin{lemma}
\label{curves-lemma-extend-to-completion}
Soit $k$ un corps. Soit $X \to Y$ un morphisme de variétés, où $Y$ est propre et $X$ est une courbe. Il existe une factorisation $X \to \overline{X} \to Y$ dans laquelle $X \to \overline{X}$ est une immersion ouverte et $\overline{X}$ est une courbe projective.
\end{lemma}

\begin{proof}
Cela résulte immédiatement du lemme \ref{curves-lemma-extend-over-dvr} et de Variétés, lemme \ref{varieties-lemma-reduced-dim-1-projective-completion}.
\end{proof}

\noindent
Voici le théorème principal de cette section. Un morphisme $f : X \to Y$ de variétés sera dit {\it constant} si son image $f(X)$ se réduit à un point $y$ de $Y$. Dans ce cas, $y$ est un point fermé de $Y$ (car l'image d'un point fermé de $X$ est un point fermé de $Y$).

\begin{theorem}
\label{curves-theorem-curves-rational-maps}
Soit $k$ un corps. Les catégories suivantes sont canoniquement équivalentes : \begin{enumerate} \item La catégorie des extensions de corps de type fini $K/k$ de degré de transcendance $1$. \item La catégorie des courbes et des applications rationnelles dominantes. \item La catégorie des courbes projectives normales et des morphismes non constants. \item La catégorie des courbes projectives non singulières et des morphismes non constants. \item La catégorie des courbes projectives régulières et des morphismes non constants. \item La catégorie des courbes propres normales et des morphismes non constants.
\end{enumerate}
\end{theorem}

\begin{proof}
L'équivalence entre les catégories (1) et (2) est la restriction de l'équivalence de Variétés, théorème \ref{varieties-theorem-varieties-rational-maps}. En effet, une variété est une courbe si et seulement si son corps de fonctions est de degré de transcendance $1$; voir par exemple Variétés, lemme \ref{varieties-lemma-dimension-locally-algebraic}.

\medskip\noindent
Les catégories (3), (4), (5) et (6) coïncident. Tout d'abord, les termes « régulier » et « non singulier » sont synonymes; voir Propriétés, définition \ref{properties-definition-regular}. Pour les schémas noethériens de dimension $1$, normalité et régularité sont équivalentes (Propriétés, lemmes \ref{properties-lemma-regular-normal} et \ref{properties-lemma-normal-dimension-1-regular}). Voir Variétés, lemme \ref{varieties-lemma-regular-point-on-curve} pour le cas des courbes. Ainsi (3) coïncide avec (5). Enfin, (6) coïncide avec (3) d'après Variétés, lemme \ref{varieties-lemma-dim-1-proper-projective}.

\medskip\noindent
Si $f : X \to Y$ est un morphisme non constant de courbes projectives non singulières, alors $f$ envoie le point générique $\eta$ de $X$ sur le point générique $\xi$ de $Y$. On obtient donc un morphisme $k(Y) = \mathcal{O}_{Y, \xi} \to \mathcal{O}_{X, \eta} = k(X)$ dans la catégorie (1). Si deux morphismes $f,g: X \to Y$ induisent le même morphisme $k(Y) \to k(X)$, alors, par l'équivalence entre (1) et (2), $f$ et $g$ sont équivalents en tant qu'applications rationnelles; donc $f=g$ d'après le lemme \ref{curves-lemma-extend-over-normal-curve}. Réciproquement, supposons donné un morphisme $k(Y) \to k(X)$ dans la catégorie (1). On obtient alors un morphisme $U \to Y$ pour un ouvert non vide $U \subset X$. D'après le lemme \ref{curves-lemma-extend-over-dvr}, celui-ci se prolonge à $X$ tout entier et fournit un morphisme dans la catégorie (5). Nous obtenons ainsi un foncteur pleinement fidèle (5)$\to$(1).

\medskip\noindent
Pour achever la démonstration, il reste à montrer que tout objet $K/k$ de (1) est le corps de fonctions d'une courbe projective normale. Nous savons déjà que $K = k(X)$ pour une certaine courbe $X$. En remplaçant $X$ par sa normalisation (qui est une variété birationnelle à $X$), on peut supposer $X$ normale (Variétés, lemme \ref{varieties-lemma-normalization-locally-algebraic}). Choisissons alors $X \to \overline{X}$ avec $\overline{X} \setminus X = \{x_1, \ldots, x_n\}$ comme dans Variétés, lemme \ref{varieties-lemma-reduced-dim-1-projective-completion}. Puisque $X$ est normale et que chacun des anneaux locaux $\mathcal{O}_{\overline{X}, x_i}$ est normal, $\overline{X}$ est une courbe projective normale, comme voulu. (Remarque : on peut aussi commencer par compactifier à l'aide de Variétés, lemme \ref{varieties-lemma-dim-1-projective-completion}, puis normaliser à l'aide de Variétés, lemme \ref{varieties-lemma-normalization-locally-algebraic}. Cette méthode évite le recours au résultat quelque peu délicat Morphismes, lemme \ref{morphisms-lemma-relative-normalization-normal-codim-1}.)
\end{proof}

\begin{definition}
\label{curves-definition-normal-projective-model}
Soit $k$ un corps. Soit $X$ une courbe. Un {\it modèle projectif non singulier de $X$} est un couple $(Y, \varphi)$, où $Y$ est une courbe projective non singulière et $\varphi : k(X) \to k(Y)$ un isomorphisme de corps de fonctions.
\end{definition}

\noindent
Un modèle projectif non singulier est déterminé à un unique isomorphisme près par le théorème \ref{curves-theorem-curves-rational-maps}. On parle donc souvent du « modèle projectif non singulier ». On omet généralement $\varphi$ dans la notation. Attention : il se peut que $Y$ ne soit pas lisse sur $k$; le lemme \ref{curves-lemma-nonsingular-model-smooth} montre toutefois que cela ne peut se produire qu'en caractéristique positive.

\begin{lemma}
\label{curves-lemma-nonsingular-model-smooth}
Soit $k$ un corps. Soient $X$ une courbe et $Y$ le modèle projectif non singulier de $X$. Si $k$ est parfait, alors $Y$ est une courbe projective lisse.
\end{lemma}

\begin{proof}
Voir par exemple Variétés, lemme \ref{varieties-lemma-regular-point-on-curve}.
\end{proof}

\begin{lemma}
\label{curves-lemma-smooth-models}
Soit $k$ un corps. Soit $X$ une courbe géométriquement irréductible sur $k$. Pour une extension de corps $K/k$, notons $Y_K$ un modèle projectif non singulier de $(X_K)_{red}$. \begin{enumerate} \item Si $X$ est propre, alors $Y_K$ est la normalisation de $X_K$. \item Il existe une extension $K/k$ finie purement inséparable telle que $Y_K$ soit lisse. \item Dès que $Y_K$ est lisse\footnote{Ou même géométriquement réduit.}, on a $H^0(Y_K, \mathcal{O}_{Y_K}) = K$. \item Étant donné un diagramme commutatif $$
\xymatrix{
\Omega & K' \ar[l] \\
K \ar[u] & k \ar[l] \ar[u]
}
$$ de corps tel que $Y_K$ et $Y_{K'}$ soient lisses, on a $Y_\Omega = (Y_K)_\Omega = (Y_{K'})_\Omega$.
\end{enumerate}
\end{lemma}

\begin{proof}
Soit $X'$ un modèle projectif non singulier de $X$. Alors $X'$ et $X$ possèdent des sous-schémas ouverts non vides isomorphes. En particulier, $X'$ est géométriquement irréductible, puisque $X$ l'est (nous omettons quelques détails). On peut donc supposer $X$ projective.

\medskip\noindent
Supposons $X$ propre. Alors $X_K$ est propre; sa normalisation $(X_K)^\nu$ est donc propre, étant finie sur un schéma propre (Variétés, lemme \ref{varieties-lemma-normalization-locally-algebraic} et Morphismes, lemmes \ref{morphisms-lemma-finite-proper} et \ref{morphisms-lemma-composition-proper}). D'autre part, $X_K$ est irréductible puisque $X$ est géométriquement irréductible. Ainsi $X_K^\nu$ est propre, normale, irréductible et birationnelle à $(X_K)_{red}$. Cela prouve (1), car toute courbe propre est projective (Variétés, lemme \ref{varieties-lemma-dim-1-proper-projective}).

\medskip\noindent
Démontrons (2). Puisque $X$ est propre et que (1) est établi, on peut appliquer Variétés, lemme \ref{varieties-lemma-finite-extension-geometrically-normal} pour trouver une extension $K/k$ finie purement inséparable telle que $Y_K$ soit géométriquement normale. Alors $Y_K$ est géométriquement régulière, puisque normalité et régularité sont équivalentes pour les courbes (Propriétés, lemme \ref{properties-lemma-normal-dimension-1-regular}). Ainsi $Y_K$ est une variété lisse d'après Variétés, lemme \ref{varieties-lemma-geometrically-regular-smooth}.

\medskip\noindent
Si $Y_K$ est géométriquement réduite, alors $Y_K$ est géométriquement intègre (Variétés, lemme \ref{varieties-lemma-geometrically-integral}) et l'on a $H^0(Y_K, \mathcal{O}_{Y_K}) = K$ d'après Variétés, lemme \ref{varieties-lemma-regular-functions-proper-variety}. Cela prouve (3), car toute variété lisse est géométriquement réduite (et même géométriquement régulière; voir Variétés, lemme \ref{varieties-lemma-geometrically-regular-smooth}).

\medskip\noindent
Si $Y_K$ est lisse, alors, pour toute extension $\Omega/K$, le changement de base $(Y_K)_\Omega$ est lisse sur $\Omega$ (Morphismes, lemme \ref{morphisms-lemma-base-change-smooth}). Il est donc clair que $Y_\Omega = (Y_K)_\Omega$. Cela prouve (4).
\end{proof}






\section{Séries linéaires}
\label{curves-section-linear-series}

\noindent
Nous nous écartons de la présentation classique (voir la remarque \ref{curves-remark-classical-linear-series}) en définissant les séries linéaires de la manière suivante.

\begin{definition}
\label{curves-definition-linear-series}
Soit $k$ un corps. Soit $X$ un schéma propre de dimension $\leq 1$ sur $k$. Soient $d \geq 0$ et $r \geq 0$. Une {\it série linéaire de degré $d$ et de dimension $r$} est un couple $(\mathcal{L}, V)$, où $\mathcal{L}$ est un $\mathcal{O}_X$-module inversible de degré $d$ (Variétés, définition \ref{varieties-definition-degree-invertible-sheaf}) et $V \subset H^0(X, \mathcal{L})$ un sous-espace vectoriel sur $k$ de dimension $r + 1$. On dira, en abrégé, que $(\mathcal{L}, V)$ est une {\it $\mathfrak g^r_d$} sur $X$.
\end{definition}

\noindent
Nous utiliserons surtout cette notion lorsque $X$ est une courbe propre non singulière. En fait, la définition ci-dessus n'est qu'une façon de généraliser la définition classique d'une $\mathfrak g^r_d$. Par exemple, si $X$ est une courbe propre, on peut généraliser les séries linéaires en autorisant $\mathcal{L}$ à être un $\mathcal{O}_X$-module cohérent sans torsion de rang $1$. Sur une courbe non singulière, tout module cohérent sans torsion est localement libre; on retrouve donc notre notion pour les courbes propres non singulières.

\medskip\noindent
Le lemme suivant précise la signification géométrique des séries linéaires sur les courbes propres non singulières.

\begin{lemma}
\label{curves-lemma-linear-series}
Soit $k$ un corps. Soit $X$ une courbe propre non singulière sur $k$. Soit $(\mathcal{L}, V)$ une $\mathfrak g^r_d$ sur $X$. Il existe alors un morphisme $$
\varphi : X \longrightarrow \mathbf{P}^r_k = \text{Proj}(k[T_0, \ldots, T_r])
$$ de variétés sur $k$ et une application $\alpha : \varphi^*\mathcal{O}_{\mathbf{P}^r_k}(1) \to \mathcal{L}$ tels que $\varphi^*T_0, \ldots, \varphi^*T_r$ aient pour images une base de $V$ par $\alpha$.
\end{lemma}

\begin{proof}
Soit $s_0, \ldots, s_r \in V$ une base sur $k$. Puisque $X$ est non singulière, l'image $\mathcal{L}' \subset \mathcal{L}$ de l'application $s_0, \ldots, s_r : \mathcal{O}_X^{\oplus r + 1} \to \mathcal{L}$ est un $\mathcal{O}_X$-module inversible, par exemple d'après Diviseurs, lemme \ref{divisors-lemma-torsion-free-over-regular-dim-1}. On utilise ensuite Constructions, lemme \ref{constructions-lemma-projective-space} pour obtenir un morphisme $$
\varphi = \varphi_{(\mathcal{L}', (s_0, \ldots, s_r))} :
X \longrightarrow \mathbf{P}^r_k
$$ comme dans l'énoncé du lemme.
\end{proof}

\begin{lemma}
\label{curves-lemma-linear-series-trivial-existence}
Soit $k$ un corps. Soit $X$ un schéma propre de dimension $\leq 1$ sur $k$. Si $X$ possède une $\mathfrak g^r_d$, alors $X$ possède une $\mathfrak g^s_d$ pour tout $0 \leq s \leq r$.
\end{lemma}

\begin{proof}
En effet, un espace vectoriel $V$ de dimension $r + 1$ sur $k$ possède un sous-espace vectoriel de dimension $s + 1$ pour tout $0 \leq s \leq r$.
\end{proof}

\begin{lemma}
\label{curves-lemma-g1d}
Soit $k$ un corps. Soit $X$ une courbe propre non singulière sur $k$. Soit $(\mathcal{L}, V)$ une $\mathfrak g^1_d$ sur $X$. Alors le morphisme $\varphi : X \to \mathbf{P}^1_k$ du lemme \ref{curves-lemma-linear-series} vérifie l'une des deux propriétés suivantes : \begin{enumerate} \item il est non constant et de degré $\leq d$; \item il se factorise par un point fermé de $\mathbf{P}^1_k$ et, dans ce cas, $H^0(X, \mathcal{O}_X) \not = k$.
\end{enumerate}
\end{lemma}

\begin{proof}
D'après le lemme \ref{curves-lemma-linear-series}, on voit que $\mathcal{L}' = \varphi^*\mathcal{O}_{\mathbf{P}^1_k}(1)$ admet une application non nulle $\mathcal{L}' \to \mathcal{L}$. Ainsi, d'après Variétés, lemme \ref{varieties-lemma-check-invertible-sheaf-trivial}, on a $0 \leq \deg(\mathcal{L}') \leq d$. Si $\deg(\mathcal{L}') = 0$, le même lemme donne $\mathcal{L}' \cong \mathcal{O}_X$; comme on dispose de deux sections linéairement indépendantes, on est dans le cas (2). Si $\deg(\mathcal{L}') > 0$, alors $\varphi$ est non constant (car l'image réciproque d'un module inversible par un morphisme constant est triviale). On a donc $$
\deg(\mathcal{L}') =
\deg(X/\mathbf{P}^1_k) \deg(\mathcal{O}_{\mathbf{P}^1_k}(1))
$$ d'après Variétés, lemme \ref{varieties-lemma-degree-pullback-map-proper-curves}. Cela achève la démonstration, puisque le degré de $\mathcal{O}_{\mathbf{P}^1_k}(1)$ est $1$.
\end{proof}

\begin{lemma}
\label{curves-lemma-grd-inequalities}
Soit $k$ un corps. Soit $X$ une courbe propre sur $k$ telle que $H^0(X, \mathcal{O}_X) = k$. Si $X$ possède une $\mathfrak g^r_d$, alors $r \leq d$. En cas d'égalité, on a $H^1(X, \mathcal{O}_X) = 0$, c'est-à-dire que le genre de $X$ (définition \ref{curves-definition-genus}) est $0$.
\end{lemma}

\begin{proof}
Soit $(\mathcal{L}, V)$ une $\mathfrak g^r_d$. Puisque cela ne fait qu'augmenter $r$, on peut supposer $V = H^0(X, \mathcal{L})$. Choisissons un élément non nul $s \in V$. Le schéma des zéros de $s$ est alors un diviseur de Cartier effectif $D \subset X$, on a $\mathcal{L} = \mathcal{O}_X(D)$ et l'on dispose d'une suite exacte courte $$
0 \to \mathcal{O}_X \to \mathcal{L} \to \mathcal{L}|_D \to 0
$$; voir Diviseurs, lemme \ref{divisors-lemma-characterize-OD} et remarque \ref{divisors-remark-ses-regular-section}. D'après Variétés, lemme \ref{varieties-lemma-degree-effective-Cartier-divisor}, on a $\deg(D) = \deg(\mathcal{L}) = d$. Puisque $D$ est un schéma artinien, on a $\mathcal{L}|_D \cong \mathcal{O}_D$\footnote{Dans notre cas, cela résulte de Diviseurs, lemme \ref{divisors-lemma-finite-trivialize-invertible-upstairs}, car $D \to \Spec(k)$ est fini.}. Ainsi $$
\dim_k H^0(D, \mathcal{L}|_D) = \dim_k H^0(D, \mathcal{O}_D) = \deg(D) = d
$$ D'autre part, par hypothèse, $\dim_k H^0(X, \mathcal{O}_X) = 1$ et $\dim H^0(X, \mathcal{L}) = r + 1$. On en conclut que $r + 1 \leq 1 + d$, c'est-à-dire $r \leq d$, comme dans le lemme.

\medskip\noindent
Supposons qu'il y ait égalité. Alors $H^0(X, \mathcal{L}) \to H^0(X, \mathcal{L}|_D)$ est surjective. Si nous savions que $H^1(X, \mathcal{L})$ est nul, la suite exacte longue de cohomologie donnerait la nullité de $H^1(X, \mathcal{O}_X)$ et la démonstration serait achevée. Nous allons remplacer $\mathcal{L}$ par une puissance assez grande dont la cohomologie s'annule. Puisque $\mathcal{L}|_D$ est le module inversible trivial (voir ci-dessus), on peut trouver une section $t$ de $\mathcal{L}$ dont la restriction à $D$ engendre $\mathcal{L}|_D$. Considérons l'application de multiplication $$
\mu :
H^0(X, \mathcal{L}) \otimes_k H^0(X, \mathcal{L})
\longrightarrow
H^0(X, \mathcal{L}^{\otimes 2})
$$ et la suite exacte courte $$
0 \to \mathcal{L} \xrightarrow{s}
\mathcal{L}^{\otimes 2} \to \mathcal{L}^{\otimes 2}|_D \to 0
$$ Puisque $H^0(\mathcal{L}) \to H^0(\mathcal{L}|_D)$ est surjective et que $t$ a pour image une trivialisation de $\mathcal{L}|_D$, on voit que $\mu(H^0(X, \mathcal{L}) \otimes t)$ fournit un sous-espace de $H^0(X, \mathcal{L}^{\otimes 2})$ qui se projette surjectivement sur les sections globales de $\mathcal{L}^{\otimes 2}|_D$. On a donc $$
\dim H^0(X, \mathcal{L}^{\otimes 2}) = r + 1 + d = 2r + 1 =
\deg(\mathcal{L}^{\otimes 2}) + 1
$$ Ainsi $\mathcal{L}^{\otimes 2}$ possède la même propriété que $\mathcal{L}$ : la dimension de l'espace des sections globales est égale au degré augmenté de un. Puisque $\mathcal{L}$ est ample (Variétés, lemme \ref{varieties-lemma-ample-curve}), il existe un entier $n_0$ tel que $\mathcal{L}^{\otimes n}$ ait une cohomologie $H^1$ nulle pour tout $n \geq n_0$ (Cohomologie des schémas, lemme \ref{coherent-lemma-coherent-proper-ample}). En appliquant l'argument précédent à $\mathcal{L}^{\otimes n}$ avec $n = 2^m$ pour un entier $m$ assez grand, on obtient donc le lemme.
\end{proof}

\begin{remark}[Définition classique] \label{curves-remark-classical-linear-series} Soit $X$ une courbe projective lisse sur un corps algébriquement clos $k$. Deux diviseurs de Cartier effectifs $D, D' \subset X$ sont dits {\it linéairement équivalents} si et seulement si $\mathcal{O}_X(D) \cong \mathcal{O}_X(D')$ en tant que $\mathcal{O}_X$-modules. Puisque $\Pic(X) = \text{Cl}(X)$ (Diviseurs, lemme \ref{divisors-lemma-local-rings-UFD-c1-bijective}), on voit que $D$ et $D'$ sont linéairement équivalents si et seulement si les diviseurs de Weil associés à $D$ et $D'$ définissent le même élément de $\text{Cl}(X)$. Pour un diviseur de Cartier effectif $D \subset X$ de degré $d$, le {\it système linéaire complet}, ou la {\it série linéaire complète}, $|D|$ de $D$ est l'ensemble des diviseurs de Cartier effectifs $E \subset X$ linéairement équivalents à $D$. Autrement dit, $|D|$ est l'ensemble des points fermés de la fibre du morphisme $$
\gamma_d :
\underline{\Hilbfunctor}^d_{X/k}
\longrightarrow
\underline{\Picardfunctor}^d_{X/k}
$$ (Schémas de Picard des courbes, lemme \ref{pic-lemma-picard-pieces}) au-dessus du point fermé correspondant à $\mathcal{O}_X(D)$. Cela munit $|D|$ d'une structure naturelle de schéma, et l'on a $|D| \cong \mathbf{P}^m_k$ avec $m + 1 = h^0(\mathcal{O}_X(D))$. Plus canoniquement, on a en fait $$
|D| = \mathbf{P}(H^0(X, \mathcal{O}_X(D))^\vee)
$$, où $(-)^\vee$ désigne le dual sur $k$ et où $\mathbf{P}$ est défini dans Constructions, exemple \ref{constructions-example-projective-space}. Dans ce langage, un {\it système linéaire}, ou une {\it série linéaire}, sur $X$ est une sous-variété fermée $L \subset |D|$ définissable par des équations linéaires. Si $L$ est de dimension $r$, alors $L = \mathbf{P}(V^\vee)$, où $V \subset H^0(X, \mathcal{O}_X(D))$ est un sous-espace vectoriel de dimension $r + 1$. Ainsi la série linéaire classique $L \subset |D|$ correspond à la série linéaire $(\mathcal{O}_X(D), V)$ définie ci-dessus.
\end{remark}







\section{Dualité}
\label{curves-section-duality}

\noindent
Dans cette section, nous explicitons, pour les schémas propres de dimension $1$ sur un corps, les conséquences des résultats très généraux sur les complexes dualisants et la dualité. Pour une discussion analogue concernant les schémas propres de dimension supérieure sur un corps, voir Dualité pour les schémas, section \ref{duality-section-duality-proper-over-field}.

\begin{lemma}
\label{curves-lemma-duality-dim-1}
Soit $X$ un schéma propre de dimension $\leq 1$ sur un corps $k$. Il existe un complexe dualisant $\omega_X^\bullet$ possédant les propriétés suivantes : \begin{enumerate} \item $H^i(\omega_X^\bullet)$ n'est non nul que pour $i = -1, 0$; \item $\omega_X = H^{-1}(\omega_X^\bullet)$ est un module cohérent de Cohen-Macaulay dont le support est constitué des composantes irréductibles de dimension $1$; \item pour $x \in X$ fermé, le module $H^0(\omega_{X, x}^\bullet)$ est non nul si et seulement si l'une des conditions suivantes est vérifiée : \begin{enumerate} \item $\dim(\mathcal{O}_{X, x}) = 0$, ou \item $\dim(\mathcal{O}_{X, x}) = 1$ et $\mathcal{O}_{X, x}$ n'est pas de Cohen-Macaulay; \end{enumerate} \item pour $K \in D_\QCoh(\mathcal{O}_X)$, il existe des isomorphismes fonctoriels\footnote{Cette propriété caractérise $\omega_X^\bullet$ dans $D_\QCoh(\mathcal{O}_X)$ à un unique isomorphisme près, par le lemme de Yoneda. Puisque $\omega_X^\bullet$ appartient à $D^b_{\textit{Coh}}(\mathcal{O}_X)$, il suffit en fait de considérer $K \in D^b_{\textit{Coh}}(\mathcal{O}_X)$.} $$
\Ext^i_X(K, \omega_X^\bullet) = \Hom_k(H^{-i}(X, K), k)
$$ compatibles aux décalages et aux triangles distingués; \item il existe des isomorphismes fonctoriels $\Hom(\mathcal{F}, \omega_X) = \Hom_k(H^1(X, \mathcal{F}), k)$ pour $\mathcal{F}$ quasi-cohérent sur $X$; \item si $X \to \Spec(k)$ est lisse de dimension relative $1$, alors $\omega_X \cong \Omega_{X/k}$.
\end{enumerate}
\end{lemma}

\begin{proof}
Notons $f : X \to \Spec(k)$ le morphisme structural. Partons du complexe dualisant relatif $$
\omega_X^\bullet = \omega_{X/k}^\bullet = a(\mathcal{O}_{\Spec(k)})
$$ décrit dans Dualité pour les schémas, remarque \ref{duality-remark-relative-dualizing-complex}. La propriété (4) est alors vraie par construction, puisque $a$ est l'adjoint à droite de $f_* : D_\QCoh(\mathcal{O}_X) \to D(\mathcal{O}_{\Spec(k)})$. Puisque $f$ est propre, on a $f^!(\mathcal{O}_{\Spec(k)}) = a(\mathcal{O}_{\Spec(k)})$ par définition; voir Dualité pour les schémas, section \ref{duality-section-upper-shriek}. Ainsi $\omega_X^\bullet$ et $\omega_X$ sont ceux de Dualité pour les schémas, exemple \ref{duality-example-proper-over-local} et exemple \ref{duality-example-equidimensional-over-field}. Les assertions (1) et (2) résultent de Dualité pour les schémas, lemme \ref{duality-lemma-vanishing-good-dualizing}. Pour un point fermé $x \in X$, le complexe $\omega_{X, x}^\bullet$ est un complexe dualisant normalisé sur $\mathcal{O}_{X, x}$; voir Dualité pour les schémas, lemme \ref{duality-lemma-good-dualizing-normalized}. L'assertion (3) résulte alors de Complexes dualisants, lemme \ref{dualizing-lemma-apply-CM}. L'assertion (5) résulte de Dualité pour les schémas, lemme \ref{duality-lemma-dualizing-module-proper-over-A} lorsque $\mathcal{F}$ est cohérent; dans le cas général, on explicite (4) pour $K = \mathcal{F}[0]$ et $i = -1$. L'assertion (6) résulte de Dualité pour les schémas, lemme \ref{duality-lemma-smooth-proper}.
\end{proof}

\begin{lemma}
\label{curves-lemma-duality-dim-1-CM}
Soit $X$ un schéma propre sur un corps $k$, de Cohen-Macaulay et équidimensionnel de dimension $1$. Le module $\omega_X$ du lemme \ref{curves-lemma-duality-dim-1} possède les propriétés suivantes : \begin{enumerate} \item $\omega_X$ est un module dualisant sur $X$ (Dualité pour les schémas, section \ref{duality-section-dualizing-module}); \item $\omega_X$ est un module cohérent de Cohen-Macaulay dont le support est $X$; \item il existe des isomorphismes fonctoriels $\Ext^i_X(K, \omega_X[1]) = \Hom_k(H^{-i}(X, K), k)$ compatibles aux décalages pour $K \in D_\QCoh(X)$; \item il existe des isomorphismes fonctoriels $\Ext^{1 + i}(\mathcal{F}, \omega_X) = \Hom_k(H^{-i}(X, \mathcal{F}), k)$ pour $\mathcal{F}$ quasi-cohérent sur $X$.
\end{enumerate}
\end{lemma}

\begin{proof}
Rappelons que, d'après la démonstration du lemme \ref{curves-lemma-duality-dim-1}, le module $\omega_X$ est celui de Dualité pour les schémas, exemple \ref{duality-example-proper-over-local} et qu'il est donc dualisant. Les autres assertions résultent du lemme \ref{curves-lemma-duality-dim-1} et du fait que $\omega_X^\bullet = \omega_X[1]$, puisque $X$ est de Cohen-Macaulay (Dualité pour les schémas, lemme \ref{duality-lemma-dualizing-module-CM-scheme}).
\end{proof}

\begin{remark}
\label{curves-remark-rework-duality-locally-free}
Soit $X$ un schéma propre de dimension $\leq 1$ sur un corps $k$. Soient $\omega_X^\bullet$ et $\omega_X$ comme dans le lemme \ref{curves-lemma-duality-dim-1}. Si $\mathcal{E}$ est un $\mathcal{O}_X$-module fini localement libre de dual $\mathcal{E}^\vee$, on a des isomorphismes canoniques $$
\Hom_k(H^{-i}(X, \mathcal{E}), k) =
H^i(X, \mathcal{E}^\vee \otimes_{\mathcal{O}_X}^\mathbf{L} \omega_X^\bullet)
$$ Cela résulte du lemme et de Cohomologie, lemme \ref{cohomology-lemma-dual-perfect-complex}. Si $X$ est de Cohen-Macaulay et équidimensionnel de dimension $1$, on a des isomorphismes canoniques $$
\Hom_k(H^{-i}(X, \mathcal{E}), k) =
H^{1 + i}(X, \mathcal{E}^\vee \otimes_{\mathcal{O}_X} \omega_X)
$$ d'après le lemme \ref{curves-lemma-duality-dim-1-CM}. En particulier, si $\mathcal{L}$ est un $\mathcal{O}_X$-module inversible, on a $$
\dim_k H^0(X, \mathcal{L}) =
\dim_k H^1(X, \mathcal{L}^{\otimes -1} \otimes_{\mathcal{O}_X} \omega_X)
$$ et
$$
\dim_k H^1(X, \mathcal{L}) =
\dim_k H^0(X, \mathcal{L}^{\otimes -1} \otimes_{\mathcal{O}_X} \omega_X)
$$
\end{remark}

\noindent
Vérifions une propriété de cohérence du complexe dualisant.

\begin{lemma}
\label{curves-lemma-sanity-check-duality}
Soit $X$ un schéma propre de dimension $\leq 1$ sur un corps $k$. Soient $\omega_X^\bullet$ et $\omega_X$ comme dans le lemme \ref{curves-lemma-duality-dim-1}. \begin{enumerate} \item Si $X \to \Spec(k)$ se factorise sous la forme $X \to \Spec(k') \to \Spec(k)$ pour un corps $k'$, alors $\omega_X^\bullet$ et $\omega_X$ vérifient les propriétés (4), (5), (6) en remplaçant $k$ par $k'$. \item Si $K/k$ est une extension de corps, les images réciproques de $\omega_X^\bullet$ et $\omega_X$ par le changement de base $X_K$ sont celles du lemme \ref{curves-lemma-duality-dim-1} pour le morphisme $X_K \to \Spec(K)$.
\end{enumerate}
\end{lemma}

\begin{proof}
Notons $f : X \to \Spec(k)$ le morphisme structural. L'assertion (1) signifie précisément que $\omega_X^\bullet$ et $\omega_X$ sont ceux du lemme \ref{curves-lemma-duality-dim-1} pour le morphisme $f' : X \to \Spec(k')$. Dans la démonstration du lemme \ref{curves-lemma-duality-dim-1}, nous avons posé $\omega_X^\bullet = a(\mathcal{O}_{\Spec(k)})$, où $a$ est l'adjoint à droite de Dualité pour les schémas, lemme \ref{duality-lemma-twisted-inverse-image} pour $f$. Il faut donc montrer que $a(\mathcal{O}_{\Spec(k)}) \cong a'(\mathcal{O}_{\Spec(k')})$, où $a'$ est l'adjoint à droite de Dualité pour les schémas, lemme \ref{duality-lemma-twisted-inverse-image} pour $f'$. Puisque $k' \subset H^0(X, \mathcal{O}_X)$, l'extension $k'/k$ est finie (Cohomologie des schémas, lemme \ref{coherent-lemma-proper-over-affine-cohomology-finite}). Par unicité des adjoints, on a $a = a' \circ b$, où $b$ est l'adjoint à droite de Dualité pour les schémas, lemme \ref{duality-lemma-twisted-inverse-image} pour $g : \Spec(k') \to \Spec(k)$. Autrement dit, on a $f^! = (f')^! \circ g^!$. Il suffit donc de montrer que $\Hom_k(k', k) \cong k'$ en tant que $k'$-modules; voir Dualité pour les schémas, exemple \ref{duality-example-affine-twisted-inverse-image}. C'est le cas, car ce sont des espaces vectoriels sur $k'$ de même dimension, à savoir $1$.

\medskip\noindent
Démontrons (2). L'assertion résulte de la compatibilité de $a$ au changement de base, d'après Dualité pour les schémas, lemme \ref{duality-lemma-more-base-change}. Voir la discussion dans Dualité pour les schémas, remarque \ref{duality-remark-relative-dualizing-complex}.
\end{proof}

\begin{lemma}
\label{curves-lemma-closed-immersion-dim-1-CM}
Soit $X$ un schéma propre de dimension $\leq 1$ sur un corps $k$. Soit $i : Y \to X$ une immersion fermée. Soient $\omega_X^\bullet$, $\omega_X$, $\omega_Y^\bullet$, $\omega_Y$ comme dans le lemme \ref{curves-lemma-duality-dim-1}. Alors \begin{enumerate} \item $\omega_Y^\bullet = R\SheafHom(\mathcal{O}_Y, \omega_X^\bullet)$; \item $\omega_Y = \SheafHom(\mathcal{O}_Y, \omega_X)$ et $i_*\omega_Y = \SheafHom_{\mathcal{O}_X}(i_*\mathcal{O}_Y, \omega_X)$.
\end{enumerate}
\end{lemma}

\begin{proof}
Notons $g : Y \to \Spec(k)$ et $f : X \to \Spec(k)$ les morphismes structuraux. Alors $g = f \circ i$. Notons $a, b, c$ les adjoints à droite associés respectivement à $f, g, i$ dans Dualité pour les schémas, lemme \ref{duality-lemma-twisted-inverse-image}. Alors $b = c \circ a$ par unicité des adjoints à droite et parce que $Rg_* = Rf_* \circ Ri_*$. Dans la démonstration du lemme \ref{curves-lemma-duality-dim-1}, nous avons posé $\omega_X^\bullet = a(\mathcal{O}_{\Spec(k)})$ et $\omega_Y^\bullet = b(\mathcal{O}_{\Spec(k)})$. Ainsi $\omega_Y^\bullet = c(\omega_X^\bullet)$, ce qui entraîne (1) d'après Dualité pour les schémas, lemme \ref{duality-lemma-twisted-inverse-image-closed}. Puisque $\omega_X = H^{-1}(\omega_X^\bullet)$ et $\omega_Y = H^{-1}(\omega_Y^\bullet)$, on en conclut que $\omega_Y = \SheafHom(\mathcal{O}_Y, \omega_X)$. Cela implique $i_*\omega_Y = \SheafHom_{\mathcal{O}_X}(i_*\mathcal{O}_Y, \omega_X)$ d'après Dualité pour les schémas, lemme \ref{duality-lemma-sheaf-with-exact-support-ext}.
\end{proof}

\begin{lemma}
\label{curves-lemma-closed-subscheme-reduced-gorenstein}
Soit $X$ un schéma propre sur un corps $k$, de Gorenstein, réduit et équidimensionnel de dimension $1$. Soit $i : Y \to X$ une immersion fermée dont la source est réduite et équidimensionnelle de dimension $1$. Soit $j : Z \to X$ l'adhérence schématique de $X \setminus Y$. Alors \begin{enumerate} \item $Y$ et $Z$ sont de Cohen-Macaulay; \item si $\mathcal{I} \subset \mathcal{O}_X$ et $\mathcal{J} \subset \mathcal{O}_X$ désignent respectivement les faisceaux d'idéaux de $Y$ et de $Z$ dans $X$, alors $$
\mathcal{I} = j_*\mathcal{I}'
\quad\text{et}\quad
\mathcal{J} = i_*\mathcal{J}'
$$, où $\mathcal{I}' \subset \mathcal{O}_Z$ et $\mathcal{J}' \subset \mathcal{O}_Y$ désignent respectivement les faisceaux d'idéaux de $Y \cap Z$ dans $Z$ et dans $Y$; \item $\omega_Y = \mathcal{J}'(i^*\omega_X)$ et $i_*(\omega_Y) = \mathcal{J}\omega_X$; \item $\omega_Z = \mathcal{I}'(j^*\omega_X)$ et $j_*(\omega_Z) = \mathcal{I}\omega_X$; \item on a les suites exactes courtes suivantes : \begin{align*}
0 \to \omega_X \to i_*i^*\omega_X \oplus j_*j^*\omega_X \to
\mathcal{O}_{Y \cap Z} \to 0 \\
0 \to i_*\omega_Y \to \omega_X \to j_*j^*\omega_X \to 0 \\
0 \to j_*\omega_Z \to \omega_X \to i_*i^*\omega_X \to 0 \\
0 \to i_*\omega_Y \oplus j_*\omega_Z \to \omega_X \to
\mathcal{O}_{Y \cap Z} \to 0 \\
0 \to \omega_Y \to i^*\omega_X \to \mathcal{O}_{Y \cap Z} \to 0 \\
0 \to \omega_Z \to j^*\omega_X \to \mathcal{O}_{Y \cap Z} \to 0
\end{align*} \end{enumerate} Ici $\omega_X$, $\omega_Y$, $\omega_Z$ sont ceux du lemme \ref{curves-lemma-duality-dim-1}.
\end{lemma}

\begin{proof}
Un schéma noethérien réduit de dimension $1$ est de Cohen-Macaulay; (1) est donc vrai. Puisque $X$ est réduit, on a $X = Y \cup Z$ au sens schématique. Avec les notations de Morphismes, lemme \ref{morphisms-lemma-scheme-theoretic-union}, et d'après son énoncé, on a une suite exacte courte $$
0 \to \mathcal{O}_X \to
\mathcal{O}_Y \oplus \mathcal{O}_Z \to \mathcal{O}_{Y \cap Z} \to 0
$$ Puisque $\mathcal{J} = \Ker(\mathcal{O}_X \to \mathcal{O}_Z)$, $\mathcal{J}' = \Ker(\mathcal{O}_Y \to \mathcal{O}_{Y \cap Z})$, $\mathcal{I} = \Ker(\mathcal{O}_X \to \mathcal{O}_Y)$ et $\mathcal{I}' = \Ker(\mathcal{O}_Z \to \mathcal{O}_{Y \cap Z})$, une chasse au diagramme donne (2). Remarquons que $\mathcal{I} + \mathcal{J}$ est le faisceau d'idéaux de $Y \cap Z$ et que $\mathcal{I} \cap \mathcal{J} = 0$. On dispose donc des suites exactes suivantes : \begin{align*}
0 \to \mathcal{O}_X \to \mathcal{O}_Y \oplus \mathcal{O}_Z \to
\mathcal{O}_{Y \cap Z} \to 0 \\
0 \to \mathcal{J} \to \mathcal{O}_X \to \mathcal{O}_Z \to 0 \\
0 \to \mathcal{I} \to \mathcal{O}_X \to \mathcal{O}_Y \to 0 \\
0 \to \mathcal{J} \oplus \mathcal{I} \to \mathcal{O}_X \to
\mathcal{O}_{Y \cap Z} \to 0 \\
0 \to \mathcal{J}' \to \mathcal{O}_Y \to \mathcal{O}_{Y \cap Z} \to 0 \\
0 \to \mathcal{I}' \to \mathcal{O}_Z \to \mathcal{O}_{Y \cap Z} \to 0
\end{align*} Puisque $X$ est de Gorenstein, $\omega_X$ est un $\mathcal{O}_X$-module inversible (Dualité pour les schémas, lemme \ref{duality-lemma-gorenstein}). Puisque $Y \cap Z$ est de dimension $0$, on a $\omega_X|_{Y \cap Z} \cong \mathcal{O}_{Y \cap Z}$. Une fois (3) et (4) démontrés, on obtient donc les suites exactes courtes du lemme en tensorisant la suite exacte courte ci-dessus par le module inversible $\omega_X$. Par symétrie, il suffit de démontrer (3), et, d'après (2), il suffit de montrer que $i_*(\omega_Y) = \mathcal{J}\omega_X$.

\medskip\noindent
On a $i_*\omega_Y = \SheafHom_{\mathcal{O}_X}(i_*\mathcal{O}_Y, \omega_X)$ d'après le lemme \ref{curves-lemma-closed-immersion-dim-1-CM}. En utilisant de nouveau l'inversibilité de $\omega_X$, on se ramène finalement à montrer que $\SheafHom_{\mathcal{O}_X}(\mathcal{O}_X/\mathcal{I}, \mathcal{O}_X)$ s'envoie isomorphiquement sur $\mathcal{J}$ par évaluation en $1$. Autrement dit, $\mathcal{J}$ doit être l'annulateur de $\mathcal{I}$. Cela résulte de ce qui précède.
\end{proof}







\section{Riemann-Roch}
\label{curves-section-Riemann-Roch}

\noindent
Soit $k$ un corps. Soit $X$ un schéma propre de dimension $\leq 1$ sur $k$. Dans Variétés, section \ref{varieties-section-divisors-curves}, nous avons défini le degré d'un $\mathcal{O}_X$-module localement libre $\mathcal{E}$ de rang constant par la formule \begin{equation}
\label{curves-equation-degree}
\deg(\mathcal{E}) =
\chi(X, \mathcal{E}) - \text{rank}(\mathcal{E})\chi(X, \mathcal{O}_X)
\end{equation}; voir Variétés, définition \ref{varieties-definition-degree-invertible-sheaf}. Dans le chapitre Homologie de Chow, nous avons défini la première classe de Chern de $\mathcal{E}$ comme une opération sur les cycles (Homologie de Chow, section \ref{chow-section-intersecting-chern-classes}) et démontré que \begin{equation}
\label{curves-equation-degree-c1}
\deg(\mathcal{E}) = \deg(c_1(\mathcal{E}) \cap [X]_1)
\end{equation}; voir Homologie de Chow, lemme \ref{chow-lemma-degree-vector-bundle}. En combinant (\ref{curves-equation-degree}) et (\ref{curves-equation-degree-c1}), on obtient une première version de la formule de Riemann-Roch : \begin{equation}
\label{curves-equation-rr}
\chi(X, \mathcal{E}) =
\deg(c_1(\mathcal{E}) \cap [X]_1) +
\text{rank}(\mathcal{E})\chi(X, \mathcal{O}_X)
\end{equation} Si $\mathcal{L}$ est un $\mathcal{O}_X$-module inversible, on peut également considérer le nombre d'intersection $(\mathcal{L} \cdot X)$ défini dans Variétés, définition \ref{varieties-definition-intersection-number}. Cela n'apporte cependant rien de nouveau, puisque \begin{equation}
\label{curves-equation-numerical-degree}
(\mathcal{L} \cdot X) = \deg(\mathcal{L})
\end{equation} d'après Variétés, lemme \ref{varieties-lemma-intersection-numbers-and-degrees-on-curves}. Si $\mathcal{L}$ est ample, cet entier est positif et s'appelle le degré \begin{equation}
\label{curves-equation-degree-X}
\deg_\mathcal{L}(X) = (\mathcal{L} \cdot X) = \deg(\mathcal{L})
\end{equation} de $X$ relativement à $\mathcal{L}$; voir Variétés, définition \ref{varieties-definition-degree}.

\medskip\noindent
Pour obtenir un véritable théorème de Riemann-Roch, nous voudrions écrire $\chi(X, \mathcal{O}_X)$ comme le degré d'un zéro-cycle canonique sur $X$. Nous renvoyons à \cite{F} pour une version entièrement générale. Nous utiliserons la dualité pour obtenir une formule lorsque $X$ est de Gorenstein; c'est toutefois, en un sens, un artifice (notamment parce que cette méthode ne peut pas fonctionner en dimension supérieure).

\medskip\noindent
Utilisons d'abord les lemmes \ref{curves-lemma-duality-dim-1} et \ref{curves-lemma-duality-dim-1-CM} pour relier la caractéristique d'Euler de $\mathcal{O}_X$ à celle du complexe dualisant ou du module dualisant.

\begin{lemma}
\label{curves-lemma-euler}
Soit $X$ un schéma propre de dimension $\leq 1$ sur un corps $k$. Avec $\omega_X^\bullet$ et $\omega_X$ comme dans le lemme \ref{curves-lemma-duality-dim-1}, on a $$
\chi(X, \mathcal{O}_X) = \chi(X, \omega_X^\bullet)
$$ Si $X$ est de Cohen-Macaulay et équidimensionnel de dimension $1$, alors
$$
\chi(X, \mathcal{O}_X) = - \chi(X, \omega_X)
$$
\end{lemma}

\begin{proof}
Le membre de droite de la première formule est défini comme suit : $$
\chi(X, \omega_X^\bullet) =
\sum\nolimits_{i \in \mathbf{Z}} (-1)^i\dim_k H^i(X, \omega_X^\bullet)
$$ Cette expression est bien définie puisque $\omega_X^\bullet$ appartient à $D^b_{\textit{Coh}}(\mathcal{O}_X)$, mais aussi parce que $$
H^i(X, \omega_X^\bullet) =
\Ext^i(\mathcal{O}_X, \omega_X^\bullet) =
\Hom_k(H^{-i}(X, \mathcal{O}_X), k)
$$, qui est toujours de dimension finie et n'est non nul que si $i = 0, -1$. Cela démontre aussi, bien entendu, la première formule. La seconde en résulte, puisque $\omega_X^\bullet = \omega_X[1]$ dans le cas de Cohen-Macaulay; voir le lemme \ref{curves-lemma-duality-dim-1-CM}.
\end{proof}

\noindent
Nous utiliserons le lemme \ref{curves-lemma-euler} pour obtenir la formule voulue pour $\chi(X, \mathcal{O}_X)$ lorsque $\omega_X$ est inversible, c'est-à-dire lorsque $X$ est de Gorenstein. L'énoncé affirme que le produit de $-1/2$ de la première classe de Chern de $\omega_X$ avec le cycle $[X]_1$ associé à $X$ est un zéro-cycle naturel sur $X$ à coefficients demi-entiers, de degré $\chi(X, \mathcal{O}_X)$. L'apparition de fractions dans l'énoncé de Riemann-Roch est inévitable.

\begin{lemma}[Riemann-Roch] \label{curves-lemma-rr} Soit $X$ un schéma propre sur un corps $k$, de Gorenstein et équidimensionnel de dimension $1$. Soit $\omega_X$ comme dans le lemme \ref{curves-lemma-duality-dim-1}. Alors \begin{enumerate} \item $\omega_X$ est un $\mathcal{O}_X$-module inversible; \item $\deg(\omega_X) = -2\chi(X, \mathcal{O}_X)$; \item pour un $\mathcal{O}_X$-module localement libre $\mathcal{E}$ de rang constant, on a $$
\chi(X, \mathcal{E}) = \deg(\mathcal{E}) -
\textstyle{\frac{1}{2}} \text{rank}(\mathcal{E}) \deg(\omega_X)
$$ et $\dim_k(H^i(X, \mathcal{E})) =
\dim_k(H^{1 - i}(X, \mathcal{E}^\vee \otimes_{\mathcal{O}_X} \omega_X))$ pour tout $i \in \mathbf{Z}$.
\end{enumerate}
\end{lemma}

\noindent
Les courbes non singulières (normales) sont de Gorenstein; voir Dualité pour les schémas, lemme \ref{duality-lemma-regular-gorenstein}.

\begin{proof}
Rappelons que les schémas de Gorenstein sont de Cohen-Macaulay (Dualité pour les schémas, lemme \ref{duality-lemma-gorenstein-CM}); ainsi $\omega_X$ est un module dualisant sur $X$, d'après le lemme \ref{curves-lemma-duality-dim-1-CM}. L'inversibilité du faisceau dualisant résulte essentiellement de la définition de la propriété de Gorenstein; voir Dualité pour les schémas, section \ref{duality-section-gorenstein}. En appliquant (\ref{curves-equation-rr}) à $\omega_X$, on obtient
$$
\chi(X, \omega_X) = 
\deg(c_1(\omega_X) \cap [X]_1) + \chi(X, \mathcal{O}_X)
$$ Avec le lemme \ref{curves-lemma-euler}, cela donne $$
2\chi(X, \mathcal{O}_X) = - \deg(c_1(\omega_X) \cap [X]_1) = - \deg(\omega_X)
$$, la seconde égalité résultant de (\ref{curves-equation-degree-c1}). En reportant cette expression dans (\ref{curves-equation-rr}) pour $\mathcal{E}$, on obtient la formule affichée du lemme. La symétrie des dimensions est une conséquence de la dualité pour $X$; voir la remarque \ref{curves-remark-rework-duality-locally-free}.
\end{proof}





\section{Quelques résultats d'annulation}
\label{curves-section-vanishing}

\noindent
Cette section rassemble quelques résultats d'annulation très faibles. On trouvera dans la section \ref{curves-section-more-vanishing} d'autres résultats, plus intéressants.

\begin{lemma}
\label{curves-lemma-automatic}
Soit $k$ un corps. Soit $X$ un schéma propre sur $k$ de dimension $1$ tel que $H^0(X, \mathcal{O}_X) = k$. Alors $X$ est connexe, de Cohen-Macaulay et équidimensionnel de dimension $1$.
\end{lemma}

\begin{proof}
Puisque $\Gamma(X, \mathcal{O}_X) = k$ n'a pas d'idempotent non trivial, $X$ est connexe. Cela montre déjà que $X$ est équidimensionnel de dimension $1$ (toute composante irréductible de dimension $0$ serait une composante connexe). Soit $\mathcal{I} \subset \mathcal{O}_X$ le plus grand sous-module cohérent à support dans les points fermés. Alors $\mathcal{I}$ existe (Diviseurs, lemme \ref{divisors-lemma-remove-embedded-points}) et est engendré par ses sections globales (Variétés, lemme \ref{varieties-lemma-chi-tensor-finite}). Puisque $1 \in \Gamma(X, \mathcal{O}_X)$ n'est pas une section de $\mathcal{I}$, on en conclut que $\mathcal{I} = 0$. Ainsi $X$ n'a pas de point immergé (Diviseurs, lemme \ref{divisors-lemma-remove-embedded-points}). Par conséquent, $X$ vérifie $(S_1)$ d'après Diviseurs, lemme \ref{divisors-lemma-S1-no-embedded}. Il s'ensuit que $X$ est de Cohen-Macaulay.
\end{proof}

\noindent
Dans cette section, nous nous plaçons dans la situation suivante.

\begin{situation}
\label{curves-situation-Cohen-Macaulay-curve}
Ici $k$ est un corps et $X$ un schéma propre sur $k$ de dimension $1$ tel que $H^0(X, \mathcal{O}_X) = k$.
\end{situation}

\noindent
D'après le lemme \ref{curves-lemma-automatic}, le schéma $X$ est de Cohen-Macaulay et équidimensionnel de dimension $1$. Le module dualisant $\omega_X$ étudié dans les lemmes \ref{curves-lemma-duality-dim-1} et \ref{curves-lemma-duality-dim-1-CM} a une cohomologie $H^1$ non nulle, puisqu'on a en fait $\dim_k H^1(X, \omega_X) = \dim_k H^0(X, \mathcal{O}_X) = 1$. Il se trouve que tout objet légèrement plus « positif » que $\omega_X$ a une cohomologie $H^1$ nulle.

\begin{lemma}
\label{curves-lemma-vanishing}
Dans la situation \ref{curves-situation-Cohen-Macaulay-curve}, soit une suite exacte $$
\omega_X \to \mathcal{F} \to \mathcal{Q} \to 0
$$ de $\mathcal{O}_X$-modules cohérents telle que $H^1(X, \mathcal{Q}) = 0$ (par exemple si $\dim(\text{Supp}(\mathcal{Q})) = 0$). Alors on a soit $H^1(X, \mathcal{F}) = 0$, soit $\mathcal{F} = \omega_X \oplus \mathcal{Q}$.
\end{lemma}

\begin{proof}
(L'assertion entre parenthèses résulte de Cohomologie des schémas, lemme \ref{coherent-lemma-coherent-support-dimension-0}.) Puisque $H^0(X, \mathcal{O}_X) = k$ est le dual de $H^1(X, \omega_X)$ (voir la section \ref{curves-section-Riemann-Roch}), on a $\dim H^1(X, \omega_X) = 1$. Le faisceau $\omega_X$ représente le foncteur $\mathcal{F} \mapsto \Hom_k(H^1(X, \mathcal{F}), k)$ sur la catégorie des $\mathcal{O}_X$-modules cohérents (Dualité pour les schémas, lemme \ref{duality-lemma-dualizing-module-proper-over-A}). Considérons une suite exacte comme dans l'énoncé et supposons que $H^1(X, \mathcal{F}) \not = 0$. Puisque $H^1(X, \mathcal{Q}) = 0$, l'application $H^1(X, \omega_X) \to H^1(X, \mathcal{F})$ est un isomorphisme. Par la propriété universelle de $\omega_X$ rappelée ci-dessus, il existe une application $\mathcal{F} \to \omega_X$ dont l'action sur $H^1$ est l'inverse de cet isomorphisme. La composée $\omega_X \to \mathcal{F} \to \omega_X$ est l'identité (par la propriété universelle), ce qui démontre le lemme.
\end{proof}

\begin{lemma}
\label{curves-lemma-vanishing-twist}
Dans la situation \ref{curves-situation-Cohen-Macaulay-curve}, soit $\mathcal{L}$ un $\mathcal{O}_X$-module inversible engendré par ses sections globales et non isomorphe à $\mathcal{O}_X$. Alors $H^1(X, \omega_X \otimes \mathcal{L}) = 0$.
\end{lemma}

\begin{proof}
Par la dualité étudiée dans la section \ref{curves-section-Riemann-Roch}, il faut montrer que $H^0(X, \mathcal{L}^{\otimes - 1}) = 0$. Sinon, on peut choisir une section globale $t$ de $\mathcal{L}^{\otimes - 1}$ et une section globale $s$ de $\mathcal{L}$ telles que $st \not = 0$. Mais alors $st$ est un multiple constant de $1$, puisque nous avons supposé que $H^0(X, \mathcal{O}_X) = k$. Il s'ensuit que $\mathcal{L} \cong \mathcal{O}_X$, d'où une contradiction.
\end{proof}

\begin{lemma}
\label{curves-lemma-globally-generated}
Dans la situation \ref{curves-situation-Cohen-Macaulay-curve}, soit une suite exacte $$
\omega_X \to \mathcal{F} \to \mathcal{Q} \to 0
$$ de $\mathcal{O}_X$-modules cohérents telle que $\dim(\text{Supp}(\mathcal{Q})) = 0$ et $\dim_k H^0(X, \mathcal{Q}) \geq 2$, et telle qu'il n'existe aucun sous-module non nul $\mathcal{Q}' \subset \mathcal{F}$ pour lequel $\mathcal{Q}' \to \mathcal{Q}$ soit injective. Alors le sous-module de $\mathcal{F}$ engendré par les sections globales s'envoie surjectivement sur $\mathcal{Q}$.
\end{lemma}

\begin{proof}
Soit $\mathcal{F}' \subset \mathcal{F}$ le sous-module engendré par les sections globales et par l'image de $\omega_X \to \mathcal{F}$. Puisque $\dim_k H^0(X, \mathcal{Q}) \geq 2$ et $\dim_k H^1(X, \omega_X) = \dim_k H^0(X, \mathcal{O}_X) = 1$, on voit que $\mathcal{F}' \to \mathcal{Q}$ n'est pas nul et que $\omega_X \to \mathcal{F}'$ n'est pas un isomorphisme. Ainsi $H^1(X, \mathcal{F}') = 0$ d'après le lemme \ref{curves-lemma-vanishing} et l'hypothèse portant sur $\mathcal{F}$. Considérons la suite exacte courte $$
0 \to \mathcal{F}' \to \mathcal{F} \to
\mathcal{Q}/\Im(\mathcal{F}' \to \mathcal{Q}) \to 0
$$ Si le quotient de droite était non nul, on obtiendrait une contradiction, car la dimension de $H^0(X, \mathcal{F})$ serait alors strictement supérieure à celle de $H^0(X, \mathcal{F}')$.
\end{proof}

\noindent
Voici un exemple de résultat d'engendrement par les sections globales.

\begin{lemma}
\label{curves-lemma-globally-generated-curve}
Dans la situation \ref{curves-situation-Cohen-Macaulay-curve}, supposons $X$ intègre. Soit $0 \to \omega_X \to \mathcal{F} \to \mathcal{Q} \to 0$ une suite exacte courte de $\mathcal{O}_X$-modules cohérents telle que $\mathcal{F}$ soit sans torsion, $\dim(\text{Supp}(\mathcal{Q})) = 0$ et $\dim_k H^0(X, \mathcal{Q}) \geq 2$. Alors $\mathcal{F}$ est engendré par ses sections globales.
\end{lemma}

\begin{proof}
Considérons le sous-module $\mathcal{F}'$ engendré par les sections globales. D'après le lemme \ref{curves-lemma-globally-generated}, l'application $\mathcal{F}' \to \mathcal{Q}$ est surjective; en particulier, $\mathcal{F}' \not = 0$. Puisque $X$ est une courbe, $\mathcal{F}' \subset \mathcal{F}$ est une inclusion de faisceaux de rang $1$; ainsi $\mathcal{Q}' = \mathcal{F}/\mathcal{F}'$ est à support dans un nombre fini de points. Raisonnons par l'absurde en supposant $\mathcal{Q}'$ non nul. On a alors $H^1(X, \mathcal{F}') \not = 0$. La propriété universelle donne donc une application non nulle $\mathcal{F}' \to \omega_X$ (Dualité pour les schémas, lemme \ref{duality-lemma-dualizing-module-proper-over-A}). L'image de la composée $\mathcal{F}' \to \omega_X \to \mathcal{F}$ est engendrée par des sections globales, donc contenue dans $\mathcal{F}'$. On obtient ainsi un endomorphisme non nul $\mathcal{F}' \to \mathcal{F}'$. Puisque $\mathcal{F}'$ est sans torsion de rang $1$ sur une courbe propre, cet endomorphisme est nécessairement un automorphisme (nous omettons les détails). Mais cela entraîne que $\mathcal{F}'$ est contenu dans $\omega_X \subset \mathcal{F}$, contrairement à la surjectivité de $\mathcal{F}' \to \mathcal{Q}$.
\end{proof}

\begin{lemma}
\label{curves-lemma-tensor-omega-with-globally-generated-invertible}
Dans la situation \ref{curves-situation-Cohen-Macaulay-curve}, soit $\mathcal{L}$ un $\mathcal{O}_X$-module inversible très ample tel que $\deg(\mathcal{L}) \geq 2$. Alors $\omega_X \otimes_{\mathcal{O}_X} \mathcal{L}$ est engendré par ses sections globales.
\end{lemma}

\begin{proof}
Supposons $k$ algébriquement clos. Soit $x \in X$ un point fermé. Soient $C_i \subset X$ les composantes irréductibles et, pour chaque $i$, soit $x_i \in C_i$ le point générique. D'après Variétés, lemme \ref{varieties-lemma-very-ample-vanish-at-point}, on peut choisir une section $s \in H^0(X, \mathcal{L})$ telle que $s$ s'annule en $x$ mais pas en $x_i$, pour tout $i$. L'application de modules correspondante $s : \mathcal{O}_X \to \mathcal{L}$ est injective, de conoyau $\mathcal{Q}$ à support dans un nombre fini de points et tel que $\dim_k H^0(X, \mathcal{Q}) \geq 2$. Considérons la suite exacte correspondante $$
0 \to \omega_X \to \omega_X \otimes \mathcal{L} \to
\omega_X \otimes \mathcal{Q} \to 0
$$ D'après le lemme \ref{curves-lemma-globally-generated}, le module engendré par les sections globales s'envoie surjectivement sur $\omega_X \otimes \mathcal{Q}$. Puisque $x$ était arbitraire, cela démontre le lemme. Nous omettons quelques détails.

\medskip\noindent
Ramenons le cas où $k$ n'est pas algébriquement clos au cas algébriquement clos. Le lecteur peut omettre la suite de la démonstration. Choisissons une clôture algébrique $\overline{k}$ de $k$ et considérons le changement de base $X_{\overline{k}}$. Vérifions que $X_{\overline{k}} \to \Spec(\overline{k})$ relève de la situation \ref{curves-situation-Cohen-Macaulay-curve}. Par changement de base plat (Cohomologie des schémas, lemme \ref{coherent-lemma-flat-base-change-cohomology}), on a $H^0(X_{\overline{k}}, \mathcal{O}) = \overline{k}$. Le schéma $X_{\overline{k}}$ est propre sur $\overline{k}$ (Morphismes, lemme \ref{morphisms-lemma-base-change-proper}) et équidimensionnel de dimension $1$ (Morphismes, lemme \ref{morphisms-lemma-dimension-fibre-after-base-change}). L'image réciproque de $\omega_X$ sur $X_{\overline{k}}$ est le module dualisant de $X_{\overline{k}}$ d'après le lemme \ref{curves-lemma-sanity-check-duality}. L'image réciproque de $\mathcal{L}$ sur $X_{\overline{k}}$ est très ample (Morphismes, lemme \ref{morphisms-lemma-very-ample-base-change}). Le degré de l'image réciproque de $\mathcal{L}$ sur $X_{\overline{k}}$ est égal au degré de $\mathcal{L}$ sur $X$ (Variétés, lemme \ref{varieties-lemma-degree-base-change}). Enfin, $\omega_X \otimes \mathcal{L}$ est engendré par ses sections globales si et seulement si son image par changement de base l'est (Variétés, lemme \ref{varieties-lemma-globally-generated-base-change}). Le résultat découle donc du cas où le corps de base est algébriquement clos.
\end{proof}




\section{Faisceaux inversibles très amples}
\label{curves-section-very-ample}

\noindent
Un critère usuel pour qu'un module inversible $\mathcal{L}$ sur un schéma $X$ de type fini sur un corps algébriquement clos soit très ample est que les sections de $\mathcal{L}$ séparent les points et les vecteurs tangents (Variétés, section \ref{varieties-section-separating-points-tangent-vectors}). Voici un autre critère pour les courbes; on le comparera à Variétés, sous-section \ref{varieties-subsection-regularity}.

\begin{lemma}
\label{curves-lemma-criterion-very-ample}
Soit $k$ un corps. Soit $X$ un schéma propre sur $k$ de dimension $1$ tel que $H^0(X, \mathcal{O}_X) = k$. Soit $\mathcal{L}$ un $\mathcal{O}_X$-module inversible. Supposons que \begin{enumerate} \item $\mathcal{L}$ possède une section globale régulière, \item $H^1(X, \mathcal{L}) = 0$, et \item $\mathcal{L}$ soit ample. \end{enumerate} Alors $\mathcal{L}^{\otimes 6}$ est très ample sur $X$ relativement à $k$.
\end{lemma}

\begin{proof}
Soit $s$ une section globale régulière de $\mathcal{L}$. Soit $i : Z = Z(s) \to X$ le schéma des zéros de $s$; voir Diviseurs, section \ref{divisors-section-effective-Cartier-invertible}. La condition (3) entraîne $Z \not = \emptyset$ (nous omettons un petit détail). Considérons la suite exacte courte $$
0 \to \mathcal{O}_X \xrightarrow{s} \mathcal{L} \to
i_*(\mathcal{L}|_Z) \to 0
$$ En tensorisant par $\mathcal{L}$, on obtient $$
0 \to \mathcal{L} \to \mathcal{L}^{\otimes 2} \to
i_*(\mathcal{L}^{\otimes 2}|_Z) \to 0
$$ Remarquons que $Z$ est de dimension $0$ (Diviseurs, lemme \ref{divisors-lemma-effective-Cartier-makes-dimension-drop}), donc est le spectre d'un anneau artinien (Variétés, lemme \ref{varieties-lemma-algebraic-scheme-dim-0}); ainsi $\mathcal{L}|_Z \cong \mathcal{O}_Z$ (Algèbre, lemme \ref{algebra-lemma-locally-free-semi-local-free}). La suite exacte courte montre aussi que $H^1(X, \mathcal{L}^{\otimes 2}) = 0$ (on peut par exemple utiliser Variétés, lemme \ref{varieties-lemma-chi-tensor-finite} pour obtenir l'annulation du terme de droite). Par récurrence sur $n \geq 1$, à l'aide de la suite $$
0 \to \mathcal{L}^{\otimes n} \xrightarrow{s}
\mathcal{L}^{\otimes n + 1} \to
i_*(\mathcal{L}^{\otimes n + 1}|_Z) \to 0
$$, on voit que $H^1(X, \mathcal{L}^{\otimes n}) = 0$ pour $n > 0$ et qu'il existe une section globale $t_{n + 1}$ de $\mathcal{L}^{\otimes n + 1}$ qui fournit une trivialisation de $\mathcal{L}^{\otimes n + 1}|_Z \cong \mathcal{O}_Z$.

\medskip\noindent
Considérons l'application de multiplication $$
\mu_n :
H^0(X, \mathcal{L}) \otimes_k H^0(X, \mathcal{L}^{\otimes n})
\oplus
H^0(X, \mathcal{L}^{\otimes 2}) \otimes_k H^0(X, \mathcal{L}^{\otimes n - 1})
\longrightarrow
H^0(X, \mathcal{L}^{\otimes n + 1})
$$ Nous affirmons qu'elle est surjective pour $n \geq 3$. Pour le voir, considérons la suite exacte courte $$
0 \to \mathcal{L}^{\otimes n} \xrightarrow{s}
\mathcal{L}^{\otimes n + 1} \to i_*(\mathcal{L}^{\otimes n + 1}|_Z) \to 0
$$ Les sections de $\mathcal{L}^{\otimes n + 1}$ provenant du terme de gauche sont dans l'image de $\mu_n$. D'autre part, puisque $H^0(\mathcal{L}^{\otimes 2}) \to H^0(\mathcal{L}^{\otimes 2}|_Z)$ est surjective (voir ci-dessus) et que $t_{n - 1}$ a pour image une trivialisation de $\mathcal{L}^{\otimes n - 1}|_Z$, $\mu_n(H^0(X, \mathcal{L}^{\otimes 2}) \otimes t_{n - 1})$ fournit un sous-espace de $H^0(X, \mathcal{L}^{\otimes n + 1})$ qui se projette surjectivement sur les sections globales de $\mathcal{L}^{\otimes n + 1}|_Z$. Cela démontre l'assertion.

\medskip\noindent
L'assertion du paragraphe précédent montre que la $k$-algèbre graduée $$
S = \bigoplus\nolimits_{n \geq 0} H^0(X, \mathcal{L}^{\otimes n})
$$ est engendrée en degrés $0, 1, 2, 3$ sur $k$. Rappelons que $X = \text{Proj}(S)$; voir Morphismes, lemme \ref{morphisms-lemma-proper-ample-is-proj}. Ainsi $S^{(6)} = \bigoplus_{n} S_{6n}$ est engendrée en degré $1$. Cela signifie que $\mathcal{L}^{\otimes 6}$ est très ample, comme voulu.
\end{proof}

\begin{lemma}
\label{curves-lemma-criterion-very-ample-bis}
Soit $k$ un corps. Soit $X$ un schéma propre sur $k$ de dimension $1$ tel que $H^0(X, \mathcal{O}_X) = k$. Soit $\mathcal{L}$ un $\mathcal{O}_X$-module inversible. Supposons que \begin{enumerate} \item $\mathcal{L}$ soit engendré par ses sections globales, \item $H^1(X, \mathcal{L}) = 0$, et \item $\mathcal{L}$ soit ample. \end{enumerate} Alors $\mathcal{L}^{\otimes 2}$ est très ample sur $X$ relativement à $k$.
\end{lemma}

\begin{proof}
Choisissons une base $s_0, \ldots, s_n$ de $H^0(X, \mathcal{L}^{\otimes 2})$ sur $k$. La propriété (1) montre que $\mathcal{L}^{\otimes 2}$ est engendré par ses sections globales et donne un morphisme $$
\varphi_{\mathcal{L}^{\otimes 2}, (s_0, \ldots, s_n)} :
X \longrightarrow \mathbf{P}^n_k
$$ Voir Constructions, section \ref{constructions-section-projective-space}. Le lemme affirme que ce morphisme est une immersion fermée. Pour le vérifier, on peut remplacer $k$ par sa clôture algébrique; voir Descente, lemme \ref{descent-lemma-descending-property-closed-immersion}. On peut donc supposer $k$ algébriquement clos.

\medskip\noindent
Supposons $k$ algébriquement clos. Pour chaque point générique $\eta_i \in X$, soit $V_i \subset H^0(X, \mathcal{L})$ le sous-espace vectoriel sur $k$ des sections s'annulant en $\eta_i$. Puisque $\mathcal{L}$ est engendré par ses sections globales, on a $V_i \not = H^0(X, \mathcal{L})$. Puisque $X$ n'a qu'un nombre fini de composantes irréductibles et que $k$ est infini, on peut trouver $s \in H^0(X, \mathcal{L})$ ne s'annulant en $\eta_i$ pour aucun $i$. Alors $s$ est une section régulière de $\mathcal{L}$ (car $X$ est de Cohen-Macaulay d'après le lemme \ref{curves-lemma-automatic}, et $\mathcal{L}$ n'a donc pas de point associé immergé).

\medskip\noindent
En particulier, toutes les assertions de la démonstration du lemme \ref{curves-lemma-criterion-very-ample} sont valables pour cette section $s$. De plus, puisque $\mathcal{L}$ est engendré par ses sections globales, on peut trouver une section globale $t \in H^0(X, \mathcal{L})$ telle que $t|_Z$ ne s'annule pas (on raisonne comme ci-dessus en utilisant la finitude de l'ensemble des points de $Z$). Dans la démonstration du lemme \ref{curves-lemma-criterion-very-ample}, on peut alors utiliser $t$ pour voir que l'application de multiplication $$
\mu_2 :
H^0(X, \mathcal{L}) \otimes_k H^0(X, \mathcal{L}^{\otimes 2})
\longrightarrow
H^0(X, \mathcal{L}^{\otimes 3})
$$ est elle aussi surjective. Ainsi $$
S = \bigoplus\nolimits_{n \geq 0} H^0(X, \mathcal{L}^{\otimes n})
$$ est engendrée en degrés $0, 1, 2$ sur $k$. En raisonnant comme dans la démonstration du lemme \ref{curves-lemma-criterion-very-ample}, on voit que $S^{(2)} = \bigoplus_{n} S_{2n}$ est engendrée en degré $1$. Cela signifie que $\mathcal{L}^{\otimes 2}$ est très ample, comme voulu. Nous omettons quelques détails.
\end{proof}









\section{Genre d'une courbe}
\label{curves-section-genus}

\noindent
Si $X$ est une courbe projective lisse géométriquement irréductible sur un corps $k$, nous avons déjà défini le genre de $X$ comme la dimension de $H^1(X, \mathcal{O}_X)$; voir Schémas de Picard des courbes, définition \ref{pic-definition-genus}. Remarquons que $H^0(X, \mathcal{O}_X) = k$ dans ce cas; voir Variétés, lemme \ref{varieties-lemma-regular-functions-proper-variety}. Généralisons cette définition comme suit.

\begin{definition}
\label{curves-definition-genus}
Soit $k$ un corps. Soit $X$ un schéma propre sur $k$ de dimension $1$ tel que $H^0(X, \mathcal{O}_X) = k$. Le {\it genre} de $X$ est alors $g = \dim_k H^1(X, \mathcal{O}_X)$.
\end{definition}

\noindent
On l'appelle parfois le {\it genre arithmétique} de $X$. Dans la littérature, le genre arithmétique d'une courbe propre $X$ sur $k$ est parfois défini par $$
p_a(X) = 1 - \chi(X, \mathcal{O}_X) =
1 - \dim_k H^0(X, \mathcal{O}_X) + \dim_k H^1(X, \mathcal{O}_X)
$$ Cette définition coïncide avec la nôtre lorsque celle-ci s'applique, puisque nous supposons $H^0(X, \mathcal{O}_X) = k$. Il faut toutefois noter que \begin{enumerate} \item $p_a(X)$ peut être négatif, et \item $p_a(X)$ dépend du corps de base $k$ et devrait être noté $p_a(X/k)$. \end{enumerate} Par exemple, si $k = \mathbf{Q}$ et $X = \mathbf{P}^1_{\mathbf{Q}(i)}$, alors $p_a(X/\mathbf{Q}) = -1$ et $p_a(X/\mathbf{Q}(i)) = 0$.

\medskip\noindent
L'hypothèse $H^0(X, \mathcal{O}_X) = k$ dans notre définition a deux conséquences. D'une part, elle écarte toute ambiguïté sur le corps de base. D'autre part, elle entraîne que le schéma $X$ est de Cohen-Macaulay et équidimensionnel de dimension $1$ (lemme \ref{curves-lemma-automatic}). Si $\omega_X$ désigne le module dualisant comme dans les lemmes \ref{curves-lemma-duality-dim-1} et \ref{curves-lemma-duality-dim-1-CM}, on a \begin{equation}
\label{curves-equation-genus}
g = \dim_k H^1(X, \mathcal{O}_X) = \dim_k H^0(X, \omega_X)
\end{equation} par dualité; voir la remarque \ref{curves-remark-rework-duality-locally-free}.

\medskip\noindent
Si $X$ est propre sur $k$ de dimension $\leq 1$ et si $H^0(X, \mathcal{O}_X)$ n'est pas égal au corps de base $k$, au lieu du genre arithmétique $p_a(X)$ donné par la formule ci-dessus, nous utiliserons l'invariant $\chi(X, \mathcal{O}_X)$. En fait, il est préconisé dans \cite[page 276]{FAC} et \cite[Introduction]{Hirzebruch} d'appeler $\chi(X, \mathcal{O}_X)$ le genre arithmétique.

\begin{lemma}
\label{curves-lemma-genus-base-change}
Soit $k'/k$ une extension de corps. Soit $X$ un schéma propre sur $k$ de dimension $1$ tel que $H^0(X, \mathcal{O}_X) = k$. Alors $X_{k'}$ est un schéma propre sur $k'$ de dimension $1$ tel que $H^0(X_{k'}, \mathcal{O}_{X_{k'}}) = k'$. De plus, le genre de $X_{k'}$ est égal au genre de $X$.
\end{lemma}

\begin{proof}
Le schéma $X_{k'}$ est de dimension $1$, par exemple d'après Morphismes, lemme \ref{morphisms-lemma-dimension-fibre-after-base-change}. Le morphisme $X_{k'} \to \Spec(k')$ est propre d'après Morphismes, lemme \ref{morphisms-lemma-base-change-proper}. L'égalité $H^0(X_{k'}, \mathcal{O}_{X_{k'}}) = k'$ résulte de Cohomologie des schémas, lemme \ref{coherent-lemma-flat-base-change-cohomology}. L'égalité des genres résulte du même lemme.
\end{proof}

\begin{lemma}
\label{curves-lemma-genus-gorenstein}
Soit $k$ un corps. Soit $X$ un schéma propre sur $k$ de dimension $1$ tel que $H^0(X, \mathcal{O}_X) = k$. Si $X$ est de Gorenstein, alors $$
\deg(\omega_X) = 2g - 2
$$, où $g$ est le genre de $X$ et $\omega_X$ est celui du lemme \ref{curves-lemma-duality-dim-1}.
\end{lemma}

\begin{proof}
Cela résulte immédiatement du lemme \ref{curves-lemma-rr}.
\end{proof}

\begin{lemma}
\label{curves-lemma-genus-smooth}
Soit $X$ une courbe propre lisse sur un corps $k$ telle que $H^0(X, \mathcal{O}_X) = k$. Alors $$
\dim_k H^0(X, \Omega_{X/k}) = g
\quad\text{et}\quad
\deg(\Omega_{X/k}) = 2g - 2
$$, où $g$ est le genre de $X$.
\end{lemma}

\begin{proof}
D'après le lemme \ref{curves-lemma-duality-dim-1}, on a $\Omega_{X/k} = \omega_X$. Les formules résultent donc de (\ref{curves-equation-genus}) et du lemme \ref{curves-lemma-genus-gorenstein}.
\end{proof}






\section{Courbes planes}
\label{curves-section-plane-curves}

\noindent
Soit $k$ un corps. Une {\it courbe plane} sera une courbe $X$ isomorphe à un sous-schéma fermé de $\mathbf{P}^2_k$. Le plongement $X \to \mathbf{P}^2_k$ sera souvent considéré comme donné. D'après Diviseurs, exemple \ref{divisors-example-closed-subscheme-of-proj}, une courbe est déterminée par l'idéal homogène correspondant $$
I(X) =
\Ker\left(
k[T_0, T_1, T_2] \longrightarrow \bigoplus \Gamma(X, \mathcal{O}_X(n))
\right)
$$ Rappelons que, dans cette situation, on a $$
X = \text{Proj}(k[T_0, T_1, T_2]/I)
$$ en tant que sous-schémas fermés de $\mathbf{P}^2_k$. Pour des renseignements plus généraux sur ces constructions, nous renvoyons à Diviseurs, exemple \ref{divisors-example-closed-subscheme-of-proj} et aux références qui y figurent. On a en fait $I(X) = (F)$ pour un polynôme homogène $F \in k[T_0, T_1, T_2]$; voir le lemme \ref{curves-lemma-equation-plane-curve}. Puisque $X$ est irréductible, $F$ est irréductible; voir le lemme \ref{curves-lemma-plane-curve}. De plus, en considérant la suite exacte courte $$
0 \to \mathcal{O}_{\mathbf{P}^2_k}(-d) \xrightarrow{F}
\mathcal{O}_{\mathbf{P}^2_k} \to \mathcal{O}_X \to 0
$$, où $d = \deg(F)$, on trouve que $H^0(X, \mathcal{O}_X) = k$ et que $X$ est de genre $(d - 1)(d - 2)/2$; voir la démonstration du lemme \ref{curves-lemma-genus-plane-curve}.

\medskip\noindent
Pour trouver des courbes planes lisses, le plus simple est d'en écrire des équations explicites. Notons $p$ la caractéristique de $k$. Si $p$ ne divise pas $d$, on peut prendre $$
F = T_0^d + T_1^d + T_2^d
$$ La courbe correspondante $X = V_+(F)$ s'appelle la {\it courbe de Fermat} de degré $d$. Elle est lisse, car, sur chaque ouvert affine standard $D_+(T_i)$, on obtient une courbe isomorphe à la courbe affine $$
\Spec(k[x, y]/(x^d + y^d + 1))
$$ L'homomorphisme d'anneaux $k \to k[x, y]/(x^d + y^d + 1)$ est lisse d'après Algèbre, lemme \ref{algebra-lemma-relative-global-complete-intersection-smooth}, puisque $d x^{d - 1}$ et $d y^{d - 1}$ engendrent l'idéal unité de $k[x, y]/(x^d + y^d + 1)$. Si $p | d$ mais $p \not = 3$, on peut utiliser l'équation $$
F = T_0^{d - 1}T_1 + T_1^{d - 1}T_2 + T_2^{d - 1}T_0
$$ En effet, sur les ouverts affines, on obtient $x + x^{d - 1}y + y^{d - 1}$, dont les dérivées sont $1 - x^{d - 2}y$ et $x^{d - 1} - y^{d - 2}$; le lieu des zéros communs de ces trois expressions est vide\footnote{En effet, puisque $x^{d - 1} = y^{d - 2}$, on a $0 = x + x^{d - 1}y + y^{d - 1} =
x + 2 x^{d - 1} y$. Puisque $x \not = 0$ du fait que $1 = x^{d - 2}y$, on obtient $0 = 1 + 2x^{d - 2}y = 3$, ce qui est absurde sauf si $3 = 0$.}. Nous laissons au lecteur le soin de construire des exemples en caractéristique $3$.

\medskip\noindent
Plus généralement, pour tout corps $k$ et tous $n$ et $d$, il existe une hypersurface lisse de degré $d$ dans $\mathbf{P}^n_k$; voir par exemple \cite{Poonen}.

\medskip\noindent
Bien entendu, cette méthode ne fournit que des courbes lisses dont le genre est un nombre triangulaire. Pour obtenir des courbes lisses de genre arbitraire, on peut chercher des courbes lisses tracées sur $\mathbf{P}^1 \times \mathbf{P}^1$ (insérer ici une référence ultérieure).

\begin{lemma}
\label{curves-lemma-equation-plane-curve}
Soit $Z \subset \mathbf{P}^2_k$ un sous-schéma fermé équidimensionnel de dimension $1$ et sans point immergé (ce qui équivaut à dire que $Z$ est de Cohen-Macaulay). Alors l'idéal $I(Z) \subset k[T_0, T_1, T_2]$ correspondant à $Z$ est principal.
\end{lemma}

\begin{proof}
C'est un cas particulier de Diviseurs, lemme \ref{divisors-lemma-equation-codim-1-in-projective-space} (voir aussi Variétés, lemme \ref{varieties-lemma-equation-codim-1-in-projective-space}). L'assertion entre parenthèses résulte du fait qu'un schéma noethérien de dimension $1$ est de Cohen-Macaulay si et seulement s'il n'a pas de point immergé; voir Diviseurs, lemme \ref{divisors-lemma-noetherian-dim-1-CM-no-embedded-points}.
\end{proof}

\begin{lemma}
\label{curves-lemma-plane-curve}
Soit $Z \subset \mathbf{P}^2_k$ comme dans le lemme \ref{curves-lemma-equation-plane-curve} et soit $I(Z) = (F)$ pour un certain $F \in k[T_0, T_1, T_2]$. Alors $Z$ est une courbe si et seulement si $F$ est irréductible.
\end{lemma}

\begin{proof}
Si $F$ est réductible, disons $F = F' F''$, soit $Z'$ le sous-schéma fermé de $\mathbf{P}^2_k$ défini par $F'$. Il est clair que $Z' \subset Z$ et que $Z' \not = Z$. Puisque $Z'$ est lui aussi de dimension $1$, on en conclut que soit $Z$ n'est pas réduit, soit $Z$ n'est pas irréductible. Réciproquement, écrivons $Z = \sum a_i D_i$, où $D_i$ sont les composantes irréductibles de $Z$; voir Diviseurs, lemmes \ref{divisors-lemma-codim-1-part} et \ref{divisors-lemma-codimension-1-is-effective-Cartier}. Soit $F_i \in k[T_0, T_1, T_2]$ le polynôme homogène engendrant l'idéal de $D_i$. Il est alors clair que $F$ et $\prod F_i^{a_i}$ définissent le même sous-schéma fermé de $\mathbf{P}^2_k$. Ainsi $F = \lambda \prod F_i^{a_i}$ pour un certain $\lambda \in k^*$, puisque tous deux engendrent l'idéal de $Z$. On voit donc que, si $F$ est irréductible, $Z$ est un diviseur premier, c'est-à-dire une courbe.
\end{proof}

\begin{lemma}
\label{curves-lemma-genus-plane-curve}
Soit $Z \subset \mathbf{P}^2_k$ comme dans le lemme \ref{curves-lemma-equation-plane-curve} et soit $I(Z) = (F)$ pour un certain $F \in k[T_0, T_1, T_2]$. Alors $H^0(Z, \mathcal{O}_Z) = k$ et le genre de $Z$ est $(d - 1)(d - 2)/2$, où $d = \deg(F)$.
\end{lemma}

\begin{proof}
Posons $S = k[T_0, T_1, T_2]$. On a une suite exacte de modules gradués $$
0 \to S(-d) \xrightarrow{F} S \to S/(F) \to 0
$$ Notons $i : Z \to \mathbf{P}^2_k$ l'immersion fermée donnée. En appliquant le foncteur exact $\widetilde{\ }$ (Constructions, lemme \ref{constructions-lemma-proj-sheaves}), on obtient $$
0 \to \mathcal{O}_{\mathbf{P}^2_k}(-d) \to
\mathcal{O}_{\mathbf{P}^2_k} \to i_*\mathcal{O}_Z \to 0
$$, car $F$ engendre l'idéal de $Z$. Les groupes de cohomologie de $\mathcal{O}_{\mathbf{P}^2_k}(-d)$ et $\mathcal{O}_{\mathbf{P}^2_k}$ sont donnés dans Cohomologie des schémas, lemme \ref{coherent-lemma-cohomology-projective-space-over-ring}. D'autre part, on a $H^q(Z, \mathcal{O}_Z) = H^q(\mathbf{P}^2_k, i_*\mathcal{O}_Z)$ d'après Cohomologie des schémas, lemme \ref{coherent-lemma-relative-affine-cohomology}. La suite exacte longue de cohomologie montre d'abord que $$
k = H^0(\mathbf{P}^2_k, \mathcal{O}_{\mathbf{P}^2_k}) \longrightarrow
H^0(Z, \mathcal{O}_Z)
$$ est un isomorphisme, puis que l'homomorphisme de bord $$
H^1(Z, \mathcal{O}_Z) \longrightarrow
H^2(\mathbf{P}^2_k, \mathcal{O}_{\mathbf{P}^2_k}(-d)) \cong
k[T_0, T_1, T_2]_{d - 3}
$$ est un isomorphisme. Comme on voit facilement que sa dimension est $(d - 1)(d - 2)/2$, la démonstration est achevée.
\end{proof}

\begin{lemma}
\label{curves-lemma-smooth-plane-curve-point-over-separable}
Soit $Z \subset \mathbf{P}^2_k$ comme dans le lemme \ref{curves-lemma-equation-plane-curve} et soit $I(Z) = (F)$ pour un certain $F \in k[T_0, T_1, T_2]$. Si $Z \to \Spec(k)$ est lisse en au moins un point et si $k$ est infini, il existe un point fermé $z \in Z$ dans le lieu lisse tel que $\kappa(z)/k$ soit une extension finie séparable de degré au plus $d$.
\end{lemma}

\begin{proof}
Supposons que $z' \in Z$ soit un point où $Z \to \Spec(k)$ est lisse. Quitte à renuméroter les coordonnées, on peut supposer $z'$ contenu dans $D_+(T_0)$. Posons $f = F(1, x, y) \in k[x, y]$. Alors $Z \cap D_+(T_0)$ est isomorphe au spectre de $k[x, y]/(f)$. Soient $f_x, f_y$ les dérivées partielles de $f$ par rapport à $x, y$. Puisque $z'$ est un point lisse de $Z/k$, l'un des éléments $f_x$ ou $f_y$ est non nul dans $z'$ (voir la discussion dans Algèbre, section \ref{algebra-section-smooth}). Après renumérotation des coordonnées, on peut supposer $f_y$ non nul en $z'$. Il existe donc un sous-schéma ouvert non vide $V \subset Z \cap D_{+}(T_0)$ tel que la projection $$
p : V \longrightarrow \Spec(k[x])
$$ soit \'etale. Puisque le degré de $f$ comme polynôme en $y$ est au plus $d$, les degrés des fibres de la projection $p$ sont au plus $d$ (voir la discussion dans Morphismes, section \ref{morphisms-section-universally-bounded}). De plus, puisque $p$ est \'etale, l'image de $p$ est un ouvert $U \subset \Spec(k[x])$. Enfin, puisque $k$ est infini, l'ensemble des points $k$-rationnels $U(k)$ de $U$ est infini, donc non vide. Choisissons un point $t \in U(k)$ et soit $z \in V$ un point d'image $t$. Alors $z$ convient.
\end{proof}





\section{Courbes de genre zéro}
\label{curves-section-genus-zero}

\noindent
Nous aurons besoin plus loin de connaître la structure d'une courbe propre de genre zéro. Une courbe propre de Gorenstein de genre zéro est une courbe plane de degré $2$, c'est-à-dire une conique; voir le lemme \ref{curves-lemma-genus-zero}. Une courbe propre quelconque de genre zéro s'obtient à partir d'une courbe non singulière (sur un corps plus grand) par une construction de somme amalgamée; voir le lemme \ref{curves-lemma-genus-zero-singular}. Puisqu'une courbe non singulière est de Gorenstein, ces deux résultats couvrent tous les cas possibles.

\begin{lemma}
\label{curves-lemma-genus-zero-pic}
Soit $X$ une courbe propre sur un corps $k$ telle que $H^0(X, \mathcal{O}_X) = k$. Si $X$ est de genre $0$, tout $\mathcal{O}_X$-module inversible $\mathcal{L}$ de degré $0$ est trivial.
\end{lemma}

\begin{proof}
En effet, on a $\dim_k H^0(X, \mathcal{L}) \geq 0 + 1 - 0 = 1$ par Riemann-Roch (lemme \ref{curves-lemma-rr}); ainsi $\mathcal{L}$ possède une section non nulle et l'on a donc $\mathcal{L} \cong \mathcal{O}_X$ d'après Variétés, lemme \ref{varieties-lemma-check-invertible-sheaf-trivial}.
\end{proof}

\begin{lemma}
\label{curves-lemma-genus-zero-positive-degree}
Soit $X$ une courbe propre sur un corps $k$ telle que $H^0(X, \mathcal{O}_X) = k$. Supposons $X$ de genre $0$. Soit $\mathcal{L}$ un $\mathcal{O}_X$-module inversible de degré $d > 0$. Alors \begin{enumerate} \item $\dim_k H^0(X, \mathcal{L}) = d + 1$ et $\dim_k H^1(X, \mathcal{L}) = 0$; \item $\mathcal{L}$ est très ample et définit une immersion fermée dans $\mathbf{P}^d_k$.
\end{enumerate}
\end{lemma}

\begin{proof}
Par définition du degré et du genre, on a $$
\dim_k H^0(X, \mathcal{L}) - \dim_k H^1(X, \mathcal{L}) = d + 1
$$ Soit $s$ une section non nulle de $\mathcal{L}$. Le schéma des zéros de $s$ est alors un diviseur de Cartier effectif $D \subset X$, on a $\mathcal{L} = \mathcal{O}_X(D)$ et l'on dispose d'une suite exacte courte $$
0 \to \mathcal{O}_X \to \mathcal{L} \to \mathcal{L}|_D \to 0
$$; voir Diviseurs, lemme \ref{divisors-lemma-characterize-OD} et remarque \ref{divisors-remark-ses-regular-section}. Puisque $H^1(X, \mathcal{O}_X) = 0$ par hypothèse, $H^0(X, \mathcal{L}) \to H^0(X, \mathcal{L}|_D)$ est surjective. Comme $\mathcal{L}|_D$ est engendré par ses sections globales (car $\dim(D) = 0$; voir Variétés, lemme \ref{varieties-lemma-chi-tensor-finite}), le module inversible $\mathcal{L}$ est engendré par ses sections globales. En fait, puisque $D$ est un schéma artinien, on a $\mathcal{L}|_D \cong \mathcal{O}_D$\footnote{Dans notre cas, cela résulte de Diviseurs, lemme \ref{divisors-lemma-finite-trivialize-invertible-upstairs}, car $D \to \Spec(k)$ est fini.}; on peut donc trouver une section $t$ de $\mathcal{L}$ dont la restriction à $D$ engendre $\mathcal{L}|_D$. La suite exacte courte montre aussi que $H^1(X, \mathcal{L}) = 0$.

\medskip\noindent
Pour $n \geq 1$, considérons l'application de multiplication $$
\mu_n :
H^0(X, \mathcal{L}) \otimes_k H^0(X, \mathcal{L}^{\otimes n})
\longrightarrow
H^0(X, \mathcal{L}^{\otimes n + 1})
$$ Nous affirmons qu'elle est surjective. Pour le voir, considérons la suite exacte courte $$
0 \to \mathcal{L}^{\otimes n} \xrightarrow{s}
\mathcal{L}^{\otimes n + 1} \to \mathcal{L}^{\otimes n + 1}|_D \to 0
$$ Les sections de $\mathcal{L}^{\otimes n + 1}$ provenant du terme de gauche sont dans l'image de $\mu_n$. D'autre part, puisque $H^0(\mathcal{L}) \to H^0(\mathcal{L}|_D)$ est surjective et que $t^n$ a pour image une trivialisation de $\mathcal{L}^{\otimes n}|_D$, on voit que $\mu_n(H^0(X, \mathcal{L}) \otimes t^n)$ fournit un sous-espace de $H^0(X, \mathcal{L}^{\otimes n + 1})$ qui se projette surjectivement sur les sections globales de $\mathcal{L}^{\otimes n + 1}|_D$. Cela démontre l'assertion.

\medskip\noindent
Remarquons que $\mathcal{L}$ est ample d'après Variétés, lemme \ref{varieties-lemma-ample-curve}. Ainsi Morphismes, lemme \ref{morphisms-lemma-proper-ample-is-proj} fournit un isomorphisme $$
X \longrightarrow
\text{Proj}\left(
\bigoplus\nolimits_{n \geq 0} H^0(X, \mathcal{L}^{\otimes n})\right)
$$ Puisque les applications $\mu_n$ sont surjectives pour tout $n \geq 1$, l'algèbre graduée du membre de droite est un quotient de l'algèbre symétrique sur $H^0(X, \mathcal{L})$. En choisissant une base sur $k$, notée $s_0, \ldots, s_d$, de $H^0(X, \mathcal{L})$, on voit qu'il s'agit d'un quotient d'une algèbre de polynômes en $d + 1$ variables. Puisque les quotients d'anneaux gradués correspondent à des immersions fermées des $\text{Proj}$ (Constructions, lemme \ref{constructions-lemma-surjective-graded-rings-generated-degree-1-map-proj}), on obtient une immersion fermée $X \to \mathbf{P}^d_k$. Nous omettons de vérifier que ce morphisme est celui de Constructions, lemme \ref{constructions-lemma-projective-space} associé aux sections $s_0, \ldots, s_d$ de $\mathcal{L}$.
\end{proof}

\begin{lemma}
\label{curves-lemma-genus-zero}
Soit $X$ une courbe propre sur un corps $k$ telle que $H^0(X, \mathcal{O}_X) = k$. Si $X$ est de Gorenstein et de genre $0$, alors $X$ est isomorphe à une courbe plane de degré $2$.
\end{lemma}

\begin{proof}
Considérons le faisceau inversible $\mathcal{L} = \omega_X^{\otimes -1}$, où $\omega_X$ est celui du lemme \ref{curves-lemma-duality-dim-1}. On a $\deg(\omega_X) = -2$ d'après le lemme \ref{curves-lemma-genus-gorenstein}, et donc $\deg(\mathcal{L}) = 2$. D'après le lemme \ref{curves-lemma-genus-zero-positive-degree}, le choix d'une base $s_0, s_1, s_2$ de l'espace vectoriel sur $k$ des sections globales de $\mathcal{L}$ donne une immersion fermée $$
\varphi_{(\mathcal{L}, (s_0, s_1, s_2))} :
X \longrightarrow \mathbf{P}^2_k
$$ Ainsi $X$ est une courbe plane d'un certain degré $d$. Soit $F \in k[T_0, T_1, T_2]_d$ son équation (lemme \ref{curves-lemma-equation-plane-curve}). Puisque le genre de $X$ est $0$, on voit que $d$ vaut $1$ ou $2$ (lemme \ref{curves-lemma-genus-plane-curve}). Remarquons que $F$ se restreint à la section nulle sur $\varphi(X)$ et que, par conséquent, $F(s_0, s_1, s_2)$ est la section nulle de $\mathcal{L}^{\otimes d}$. Puisque $s_0, s_1, s_2$ sont linéairement indépendantes, $F$ ne peut pas être linéaire, c'est-à-dire que $d = \deg(F) \geq 2$. Ainsi $d = 2$, ce qui achève la démonstration.
\end{proof}

\begin{proposition}[Caractérisation de la droite projective] \label{curves-proposition-projective-line} Soit $k$ un corps. Soit $X$ une courbe propre sur $k$. Les conditions suivantes sont équivalentes : \begin{enumerate} \item $X \cong \mathbf{P}^1_k$; \item $X$ est lisse et géométriquement irréductible sur $k$, $X$ est de genre $0$ et $X$ possède un module inversible de degré impair; \item $X$ est géométriquement intègre sur $k$, $X$ est de genre $0$, $X$ est de Gorenstein et $X$ possède un faisceau inversible de degré impair; \item $H^0(X, \mathcal{O}_X) = k$, $X$ est de genre $0$, $X$ est de Gorenstein et $X$ possède un faisceau inversible de degré impair; \item $X$ est géométriquement intègre sur $k$, $X$ est de genre $0$ et $X$ possède un $\mathcal{O}_X$-module inversible de degré $1$; \item $H^0(X, \mathcal{O}_X) = k$, $X$ est de genre $0$ et $X$ possède un $\mathcal{O}_X$-module inversible de degré $1$; \item $H^1(X, \mathcal{O}_X) = 0$ et $X$ possède un $\mathcal{O}_X$-module inversible de degré $1$; \item $H^1(X, \mathcal{O}_X) = 0$ et $X$ possède des points fermés $x_1, \ldots, x_n$ tels que $\mathcal{O}_{X, x_i}$ soit normal et $\gcd([\kappa(x_i) : k]) = 1$; \item ajouter d'autres conditions ici.
\end{enumerate}
\end{proposition}

\begin{proof}
Nous montrerons que chacune des conditions (2) -- (8) entraîne (1), et nous omettons de vérifier que (1) entraîne (2) -- (8).

\medskip\noindent
Supposons (2). Un schéma lisse sur $k$ est géométriquement réduit (Variétés, lemme \ref{varieties-lemma-smooth-geometrically-normal}) et régulier (Variétés, lemme \ref{varieties-lemma-smooth-regular}). Ainsi $X$ est de Gorenstein (Dualité pour les schémas, lemme \ref{duality-lemma-regular-gorenstein}). On est donc ramené à (3).

\medskip\noindent
Supposons (3). Puisque $X$ est géométriquement intègre sur $k$, on a $H^0(X, \mathcal{O}_X) = k$ d'après Variétés, lemme \ref{varieties-lemma-regular-functions-proper-variety}, et l'on est ramené à (4).

\medskip\noindent
Supposons (4). Puisque $X$ est de Gorenstein, le module dualisant $\omega_X$ du lemme \ref{curves-lemma-duality-dim-1} est de degré $\deg(\omega_X) = -2$ d'après le lemme \ref{curves-lemma-genus-gorenstein}. Avec l'existence supposée d'un module inversible de degré impair, cela entraîne l'existence d'un module inversible de degré $1$. On est ainsi ramené à (6).

\medskip\noindent
Supposons (5). Puisque $X$ est géométriquement intègre sur $k$, on a $H^0(X, \mathcal{O}_X) = k$ d'après Variétés, lemme \ref{varieties-lemma-regular-functions-proper-variety}, et l'on est ramené à (6).

\medskip\noindent
Supposons (6). Alors $X \cong \mathbf{P}^1_k$ d'après le lemme \ref{curves-lemma-genus-zero-positive-degree}.

\medskip\noindent
Supposons (7). Remarquons que $\kappa = H^0(X, \mathcal{O}_X)$ est un corps fini sur $k$ d'après Variétés, lemme \ref{varieties-lemma-regular-functions-proper-variety}. Si $d = [\kappa : k] > 1$, le degré de tout faisceau inversible est divisible par $d$ et il ne peut donc exister de faisceau inversible de degré $1$. Ainsi $d = 1$ et l'on est ramené au cas (6).

\medskip\noindent
Supposons (8). Remarquons que $\kappa = H^0(X, \mathcal{O}_X)$ est un corps fini sur $k$ d'après Variétés, lemme \ref{varieties-lemma-regular-functions-proper-variety}. Puisque $\kappa \subset \kappa(x_i)$, l'hypothèse sur le pgcd des degrés entraîne $k = \kappa$. La même condition permet de trouver des entiers $a_i$ tels que $1 = \sum a_i[\kappa(x_i) : k]$. Puisque $x_i$ définit un diviseur de Cartier effectif sur $X$ d'après Variétés, lemme \ref{varieties-lemma-regular-point-on-curve}, on peut considérer le module inversible $\mathcal{L} = \mathcal{O}_X(\sum a_i x_i)$. Par le choix de $a_i$, le degré de $\mathcal{L}$ est $1$. Ainsi $X \cong \mathbf{P}^1_k$ d'après le lemme \ref{curves-lemma-genus-zero-positive-degree}.
\end{proof}

\begin{lemma}
\label{curves-lemma-genus-zero-singular}
Soit $X$ une courbe propre sur un corps $k$ telle que $H^0(X, \mathcal{O}_X) = k$. Supposons $X$ singulière et de genre $0$. Il existe alors un diagramme $$
\xymatrix{
x' \ar[d] \ar[r] & X' \ar[d]^\nu \ar[r] & \Spec(k') \ar[d] \\
x \ar[r] & X \ar[r] & \Spec(k)
}
$$ dans lequel \begin{enumerate} \item $k'/k$ est une extension finie non triviale; \item $X' \cong \mathbf{P}^1_{k'}$; \item $x'$ est un point $k'$-rationnel de $X'$; \item $x$ est un point $k$-rationnel de $X$; \item $X' \setminus \{x'\} \to X \setminus \{x\}$ est un isomorphisme; \item $0 \to \mathcal{O}_X \to \nu_*\mathcal{O}_{X'} \to k'/k \to 0$ est une suite exacte courte, où $k'/k = \kappa(x')/\kappa(x)$ désigne le faisceau gratte-ciel au point $x$.
\end{enumerate}
\end{lemma}

\begin{proof}
Soit $\nu : X' \to X$ la normalisation de $X$; voir Variétés, sections \ref{varieties-section-normalization} et \ref{varieties-section-normalization-one-dimensional}. Puisque $X$ est singulière, $\nu$ n'est pas un isomorphisme. Si l'on pose $k' = H^0(X', \mathcal{O}_{X'}),\ n = [k' : k]$, l'extension ainsi obtenue est finie sur $k$ (Variétés, lemme \ref{varieties-lemma-regular-functions-proper-variety}). La suite exacte courte $$
0 \to \mathcal{O}_X \to \nu_*\mathcal{O}_{X'} \to \mathcal{Q} \to 0
$$ et le fait que $\mathcal{Q}$ est à support dans un nombre fini de points fermés donnent \begin{enumerate} \item $H^1(X', \mathcal{O}_{X'}) = 0$, c'est-à-dire que $X'$ est de genre $0$ comme courbe sur $k'$; \item une suite exacte courte $0 \to k \to k' \to H^0(X, \mathcal{Q}) \to 0$. \end{enumerate} En particulier, $k'/k$ est une extension non triviale.

\medskip\noindent
Considérons ensuite ce qu'on appelle souvent l'{\it idéal conducteur} $$
\mathcal{I} = \SheafHom_{\mathcal{O}_X}(\nu_*\mathcal{O}_{X'}, \mathcal{O}_X)
$$ C'est un $\mathcal{O}_X$-module quasi-cohérent. On considère $\mathcal{I}$ comme un idéal de $\mathcal{O}_X$ par l'application $\varphi \mapsto \varphi(1)$. Ainsi $\mathcal{I}(U)$ est l'ensemble des éléments $f \in \mathcal{O}_X(U)$ tels que $f \left(\nu_*\mathcal{O}_{X'}(U)\right) \subset \mathcal{O}_X(U)$. Autrement dit, la condition est que $f$ annule $\mathcal{Q}$. On a donc une suite exacte qui le définit : $$
0 \to \mathcal{I} \to \mathcal{O}_X \to
\SheafHom_{\mathcal{O}_X}(\mathcal{Q}, \mathcal{Q})
$$ Soit $U \subset X$ un ouvert affine contenant le support de $\mathcal{Q}$. Alors $V = \mathcal{Q}(U) = H^0(X, \mathcal{Q})$ est un espace vectoriel sur $k$ de dimension $n - 1$. L'image de $\mathcal{O}_X(U) \to \Hom_k(V, V)$ est une sous-algèbre commutative, donc de dimension $\leq n - 1$ sur $k$ (c'est une propriété des sous-algèbres commutatives des algèbres de matrices; nous omettons les détails). On en déduit une suite exacte courte $$
0 \to \mathcal{I} \to \mathcal{O}_X \to \mathcal{A} \to 0
$$, où $\text{Supp}(\mathcal{A}) = \text{Supp}(\mathcal{Q})$ et $\dim_k H^0(X, \mathcal{A}) \leq n - 1$. D'autre part, la description $\mathcal{I} = \SheafHom_{\mathcal{O}_X}(\nu_*\mathcal{O}_{X'}, \mathcal{O}_X)$ munit $\mathcal{I}$ d'une structure de $\nu_*\mathcal{O}_{X'}$-module telle que l'inclusion $\mathcal{I} \to \nu_*\mathcal{O}_{X'}$ soit une application de $\nu_*\mathcal{O}_{X'}$-modules. On en conclut que $\mathcal{I} = \nu_*\mathcal{I}'$ pour un faisceau quasi-cohérent d'idéaux $\mathcal{I}' \subset \mathcal{O}_{X'}$; voir Morphismes, lemme \ref{morphisms-lemma-affine-equivalence-modules}. Définissons $\mathcal{A}'$ comme le conoyau : $$
0 \to \mathcal{I}' \to \mathcal{O}_{X'} \to \mathcal{A}' \to 0
$$ En combinant les suites exactes précédentes, on obtient une suite exacte courte $0 \to \mathcal{A} \to \nu_*\mathcal{A}' \to \mathcal{Q} \to 0$. L'estimation ci-dessus, jointe à $\dim_k H^0(X, \mathcal{Q}) = n - 1$, donne $$
\dim_k H^0(X', \mathcal{A}') =
\dim_k H^0(X, \mathcal{A}) + \dim_k H^0(X, \mathcal{Q}) \leq 2 n - 2
$$ Cependant, puisque $X'$ est une courbe sur $k'$, le membre de gauche est divisible par $n$ (Variétés, lemme \ref{varieties-lemma-divisible}). Comme $\mathcal{A}$ et $\mathcal{A}'$ ne peuvent être nuls, on en déduit que $\dim_k H^0(X', \mathcal{A}') = n$, ce qui signifie que $\mathcal{I}'$ est le faisceau d'idéaux d'un point $k'$-rationnel, noté $x'$. La proposition \ref{curves-proposition-projective-line} donne $X' \cong \mathbf{P}^1_{k'}$. En revenant aux égalités précédentes, on en conclut que $\dim_k H^0(X, \mathcal{A}) = 1$. Cela signifie que $\mathcal{I}$ est le faisceau d'idéaux d'un point $k$-rationnel, noté $x$. Alors $\mathcal{A} = \kappa(x) = k$ et $\mathcal{A}' = \kappa(x') = k'$ en tant que faisceaux gratte-ciel. En comparant les suites exactes données ci-dessus, on obtient immédiatement l'assertion du lemme relative aux faisceaux structuraux.
\end{proof}

\begin{example}
\label{curves-example-squish-on-P1}
En fait, la situation décrite dans le lemme \ref{curves-lemma-genus-zero-singular} se présente pour toute extension finie non triviale $k'/k$. En effet, on peut considérer $$
A = \{f \in k'[x] \mid f(0) \in k \}
$$ Le spectre de $A$ est une courbe affine que l'on peut recoller au spectre de $B = k'[y]$ au moyen de l'isomorphisme $A_x \cong B_y$ envoyant $x^{-1}$ sur $y$. On obtient une courbe propre $X$ telle que $H^0(X, \mathcal{O}_X) = k$ et dont le point singulier $x$ correspond à l'idéal maximal $A \cap (x)$. La normalisation de $X$ est $\mathbf{P}^1_{k'}$, exactement comme dans le lemme.
\end{example}








\section{Genre géométrique}
\label{curves-section-geometric-genus}

\noindent
Si $X$ est une courbe propre et {\bf lisse} sur $k$ telle que $H^0(X, \mathcal{O}_X) = k$, on appelle $$
p_g(X) = \dim_k H^0(X, \Omega_{X/k})
$$ le {\it genre géométrique} de $X$. D'après le lemme \ref{curves-lemma-genus-smooth}, le genre géométrique de $X$ coïncide avec son genre (arithmétique). En dimension supérieure, en revanche, le genre géométrique diffère du genre arithmétique; voir la remarque \ref{curves-remark-genus-higher-dimension}.

\medskip\noindent
Pour les courbes singulières, nous définirons le genre géométrique comme suit.

\begin{definition}
\label{curves-definition-geometric-genus}
Soit $k$ un corps. Soit $X$ une courbe géométriquement irréductible sur $k$. Le {\it genre géométrique} de $X$ est le genre d'un modèle projectif lisse de $X$, éventuellement défini sur une extension de $k$ comme dans le lemme \ref{curves-lemma-smooth-models}.
\end{definition}

\noindent
Si $k$ est parfait, le modèle projectif non singulier $Y$ de $X$ est lisse (lemme \ref{curves-lemma-nonsingular-model-smooth}) et le genre géométrique de $X$ est simplement le genre de $Y$. Cela peut être faux si $k$ n'est pas parfait. Dans ce cas, on choisit une extension $K/k$ telle que le modèle projectif non singulier $Y_K$ de $(X_K)_{red}$ soit une courbe projective lisse et l'on définit le genre géométrique de $X$ comme le genre de $Y_K$. Les lemmes \ref{curves-lemma-smooth-models} et \ref{curves-lemma-genus-base-change} assurent que cette définition a un sens.

\begin{remark}
\label{curves-remark-genus-higher-dimension}
Supposons que $X$ soit une variété propre lisse de dimension $d$ sur un corps algébriquement clos $k$. On définit alors souvent le {\it genre arithmétique} par $p_a(X) = (-1)^d(\chi(X, \mathcal{O}_X) - 1)$ et le {\it genre géométrique} par $p_g(X) = \dim_k H^0(X, \Omega^d_{X/k})$. Dans cette situation, le genre arithmétique et le genre géométrique ne coïncident plus, bien que l'on ait toujours $\omega_X \cong \Omega_{X/k}^d$. Par exemple, si $d = 2$, on a \begin{align*}
p_a(X) - p_g(X) & =
h^0(X, \mathcal{O}_X) - h^1(X, \mathcal{O}_X) + h^2(X, \mathcal{O}_X) - 1
- h^0(X, \Omega^2_{X/k}) \\
& =
- h^1(X, \mathcal{O}_X) + h^2(X, \mathcal{O}_X) - h^0(X, \omega_X) \\
& =
- h^1(X, \mathcal{O}_X)
\end{align*}, où $h^i(X, \mathcal{F}) = \dim_k H^i(X, \mathcal{F})$ et où la dernière égalité résulte de la dualité. Ainsi, pour une surface, la différence $p_g(X) - p_a(X)$ est toujours positive ou nulle; on l'appelle parfois l'irrégularité de la surface. Si $X = C_1 \times C_2$ est un produit de courbes projectives lisses de genres $g_1$ et $g_2$, son irrégularité est $g_1 + g_2$.
\end{remark}






\section{Riemann-Hurwitz}
\label{curves-section-riemann-hurewitz}

\noindent
Soit $k$ un corps. Soit $f : X \to Y$ un morphisme de courbes lisses sur $k$. On obtient une suite exacte canonique $$
f^*\Omega_{Y/k} \xrightarrow{\text{d}f} \Omega_{X/k}
\longrightarrow \Omega_{X/Y} \longrightarrow 0
$$ d'après Morphismes, lemme \ref{morphisms-lemma-triangle-differentials}. Puisque $X$ et $Y$ sont lisses, les faisceaux $\Omega_{X/k}$ et $\Omega_{Y/k}$ sont des modules inversibles; voir Morphismes, lemme \ref{morphisms-lemma-smooth-omega-finite-locally-free}. Supposons la première application non nulle, c'est-à-dire $f$ génériquement \'etale; voir le lemme \ref{curves-lemma-generically-etale}. Soit $R \subset X$ le sous-schéma fermé défini par la différente $\mathfrak{D}_f$ de $f$. D'après Discriminants, lemme \ref{discriminant-lemma-discriminant-quasi-finite-morphism-smooth}, c'est aussi le lieu d'annulation de $\text{d}f$; il s'agit d'un diviseur de Cartier effectif, et l'on obtient $$
f^*\Omega_{Y/k} \otimes_{\mathcal{O}_X} \mathcal{O}_X(R) = \Omega_{X/k}
$$ En particulier, si $X$, $Y$ sont projectives avec $k = H^0(Y, \mathcal{O}_Y) = H^0(X, \mathcal{O}_X)$ et si $X$, $Y$ sont de genres $g_X$, $g_Y$, on obtient la formule de Riemann-Hurwitz \begin{align*}
2g_X - 2 & =
\deg(\Omega_{X/k}) \\
& =
\deg(f^*\Omega_{Y/k} \otimes_{\mathcal{O}_X} \mathcal{O}_X(R)) \\
& =
\deg(f) \deg(\Omega_{Y/k}) + \deg(R) \\
& =
\deg(f) (2g_Y - 2) + \deg(R)
\end{align*} La première et la dernière égalité résultent du lemme \ref{curves-lemma-genus-smooth}; la deuxième, de l'isomorphisme de faisceaux inversibles ci-dessus; la troisième, de l'additivité des degrés (Variétés, lemme \ref{varieties-lemma-degree-tensor-product}), de la formule du degré d'une image réciproque (Variétés, lemme \ref{varieties-lemma-degree-pullback-map-proper-curves}) et, enfin, de la formule du degré de $\mathcal{O}_X(R)$ (Variétés, lemme \ref{varieties-lemma-degree-effective-Cartier-divisor}).

\medskip\noindent
Pour utiliser la formule de Riemann-Hurwitz, il faut calculer $\deg(R) = \dim_k \Gamma(R, \mathcal{O}_R)$. La structure des schémas de dimension zéro sur $k$ (voir par exemple Variétés, lemme \ref{varieties-lemma-algebraic-scheme-dim-0}) montre que $R$ est une réunion disjointe finie de spectres d'anneaux locaux artiniens $R = \coprod_{x \in R} \Spec(\mathcal{O}_{R, x})$, chacun des $\mathcal{O}_{R, x}$ étant de dimension finie sur $k$. Ainsi
$$
\deg(R) = \sum\nolimits_{x \in R} \dim_k \mathcal{O}_{R, x} =
\sum\nolimits_{x \in R} d_x [\kappa(x) : k]
$$, où $$
d_x = \text{length}_{\mathcal{O}_{R, x}} \mathcal{O}_{R, x} =
\text{length}_{\mathcal{O}_{X, x}} \mathcal{O}_{R, x}
$$ est la multiplicité de $x$ dans $R$ (voir Algèbre, lemme \ref{algebra-lemma-pushdown-module}). Soit $x \in X$ un point fermé d'image $y \in Y$. En passant aux germes, on obtient une suite exacte $$
\Omega_{Y/k, y} \to \Omega_{X/k, x} \to \Omega_{X/Y, x} \to 0
$$ En choisissant des générateurs locaux $\eta_x$ et $\eta_y$ des modules (libres de rang $1$) $\Omega_{X/k, x}$ et $\Omega_{Y/k, y}$, on voit que $
\eta_y \mapsto h \eta_x
$ pour un élément non nul $h \in \mathcal{O}_{X, x}$. La suite exacte montre que $\Omega_{X/Y, x} \cong \mathcal{O}_{X, x}/h\mathcal{O}_{X, x}$ en tant que $\mathcal{O}_{X, x}$-modules. Puisque le diviseur $R$ est défini par $h$ (voir ci-dessus), on a $\mathcal{O}_{R, x} = \mathcal{O}_{X, x}/h\mathcal{O}_{X, x}$. On obtient donc les égalités suivantes : \begin{align*}
d_x
& =
\text{length}_{\mathcal{O}_{X, x}}(\mathcal{O}_{R, x}) \\
& =
\text{length}_{\mathcal{O}_{X, x}}(\mathcal{O}_{X, x}/h\mathcal{O}_{X, x}) \\
& =
\text{length}_{\mathcal{O}_{X, x}}(\Omega_{X/Y, x}) \\
& =
\text{ord}_{\mathcal{O}_{X, x}}(h) \\
& =
\text{ord}_{\mathcal{O}_{X, x}}(``\eta_y/\eta_x")
\end{align*} La première résulte de la définition de $d_x$. Nous avons établi la deuxième et la troisième ci-dessus. La quatrième égalité est la définition de $\text{ord}$; voir Algèbre, définition \ref{algebra-definition-ord}. Puisque $\mathcal{O}_{X, x}$ est un anneau de valuation discrète, l'entier $\text{ord}_{\mathcal{O}_{X, x}}(h)$ est simplement la valuation de $h$. La cinquième égalité sert de moyen mnémotechnique.

\medskip\noindent
Voici un cas où l'on peut « calculer » la multiplicité $d_x$ en fonction d'autres invariants. Si $\kappa(x)$ est séparable sur $k$, on peut choisir $\eta_x = \text{d}s$ et $\eta_y = \text{d}t$, où $s$ et $t$ sont des uniformisantes de $\mathcal{O}_{X, x}$ et $\mathcal{O}_{Y, y}$ (lemme \ref{curves-lemma-uniformizer-works}). Alors $t \mapsto u s^{e_x}$ pour une unité $u \in \mathcal{O}_{X, x}$, où $e_x$ est l'indice de ramification de l'extension $\mathcal{O}_{Y, y} \subset \mathcal{O}_{X, x}$. On obtient donc $$
\eta_y = \text{d}t = \text{d}(u s^{e_x}) =
e_x s^{e_x - 1} u \text{d}s + s^{e_x} \text{d}u
$$ En écrivant $\text{d}u = w \text{d}s$ pour un certain $w \in \mathcal{O}_{X, x}$, on voit que $$
``\eta_y/\eta_x" = e_x s^{e_x - 1} u + s^{e_x} w = (e_x u + s w)s^{e_x - 1}
$$ Son ordre d'annulation est donc $e_x - 1$, sauf si la caractéristique de $\kappa(x)$ est $p > 0$ et si $p$ divise $e_x$, auquel cas l'ordre d'annulation est $> e_x - 1$.

\medskip\noindent
En combinant tout ce qui précède, on voit que, si $k$ est de caractéristique zéro, alors $$
2g_X - 2 = (2g_Y - 2)\deg(f) +
\sum\nolimits_{x \in X} (e_x - 1)[\kappa(x) : k]
$$, où $e_x$ est l'indice de ramification de $\mathcal{O}_{X, x}$ au-dessus de $\mathcal{O}_{Y, f(x)}$. Cette formule précise est valable si et seulement si toute la ramification est modérée, c'est-à-dire lorsque les extensions résiduelles $\kappa(x)/\kappa(y)$ sont séparables et que $e_x$ est premier à la caractéristique de $k$; les arguments précédents ne suffisent cependant pas à le démontrer. Nous renvoyons au lemme \ref{curves-lemma-rhe} et à sa démonstration.

\begin{lemma}
\label{curves-lemma-generically-etale}
\begin{slogan}
Un morphisme de courbes lisses est séparable si et seulement s'il est étale presque partout
\end{slogan}
Soit $k$ un corps. Soit $f : X \to Y$ un morphisme de courbes lisses sur $k$. Les conditions suivantes sont équivalentes : \begin{enumerate} \item $\text{d}f : f^*\Omega_{Y/k} \to \Omega_{X/k}$ est non nulle; \item $\Omega_{X/Y}$ est à support dans un fermé propre de $X$; \item il existe un ouvert non vide $U \subset X$ tel que $f|_U : U \to Y$ soit non ramifié; \item il existe un ouvert non vide $U \subset X$ tel que $f|_U : U \to Y$ soit \'etale; \item l'extension $k(X)/k(Y)$ de corps de fonctions est finie séparable.
\end{enumerate}
\end{lemma}

\begin{proof}
Puisque $X$ et $Y$ sont lisses, les faisceaux $\Omega_{X/k}$ et $\Omega_{Y/k}$ sont des modules inversibles; voir Morphismes, lemme \ref{morphisms-lemma-smooth-omega-finite-locally-free}. La suite exacte $$
f^*\Omega_{Y/k} \longrightarrow \Omega_{X/k}
\longrightarrow \Omega_{X/Y} \longrightarrow 0
$$ de Morphismes, lemme \ref{morphisms-lemma-triangle-differentials} montre que (1) et (2) sont équivalentes entre elles et que chacune équivaut à la non-nullité de $f^*\Omega_{Y/k} \to \Omega_{X/k}$ au point générique. L'équivalence de (2) et (3) résulte de Morphismes, lemme \ref{morphisms-lemma-unramified-omega-zero}. L'équivalence de (3) et (4) résulte de Morphismes, lemme \ref{morphisms-lemma-flat-unramified-etale} et de la platitude automatique (lemme \ref{curves-lemma-flat}). Pour l'équivalence de (5) et (4), on utilise Algèbre, lemme \ref{algebra-lemma-smooth-at-generic-point}. Nous omettons quelques détails.
\end{proof}

\begin{lemma}
\label{curves-lemma-rh}
Soit $f : X \to Y$ un morphisme de courbes propres lisses sur un corps $k$ satisfaisant aux conditions équivalentes du lemme \ref{curves-lemma-generically-etale}. Si $k = H^0(Y, \mathcal{O}_Y) = H^0(X, \mathcal{O}_X)$ et si $X$ et $Y$ sont de genres $g_X$ et $g_Y$, alors $$
2g_X - 2 = (2g_Y - 2) \deg(f) + \deg(R)
$$, où $R \subset X$ est le diviseur de Cartier effectif défini par la différente de $f$.
\end{lemma}

\begin{proof}
Voir la discussion ci-dessus; nous avons utilisé Discriminants, lemme \ref{discriminant-lemma-discriminant-quasi-finite-morphism-smooth}, le lemme \ref{curves-lemma-genus-smooth} et Variétés, lemmes \ref{varieties-lemma-degree-tensor-product} et \ref{varieties-lemma-degree-pullback-map-proper-curves}.
\end{proof}

\begin{lemma}
\label{curves-lemma-uniformizer-works}
Soit $X \to \Spec(k)$ lisse de dimension relative $1$ en un point fermé $x \in X$. Si $\kappa(x)$ est séparable sur $k$, alors, pour toute uniformisante $s$ de l'anneau de valuation discrète $\mathcal{O}_{X, x}$, l'élément $\text{d}s$ engendre librement $\Omega_{X/k, x}$ sur $\mathcal{O}_{X, x}$.
\end{lemma}

\begin{proof}
L'anneau $\mathcal{O}_{X, x}$ est un anneau de valuation discrète d'après Algèbre, lemme \ref{algebra-lemma-characterize-smooth-over-field}. Puisque $x$ est fermé, $\kappa(x)$ est fini sur $k$. Ainsi, si $\kappa(x)/k$ est séparable, toute uniformisante $s$ a une image non nulle dans $\Omega_{X/k, x} \otimes_{\mathcal{O}_{X, x}} \kappa(x)$ d'après Algèbre, lemme \ref{algebra-lemma-computation-differential}. Puisque $\Omega_{X/k, x}$ est libre de rang $1$ sur $\mathcal{O}_{X, x}$, le résultat s'ensuit.
\end{proof}

\begin{lemma}
\label{curves-lemma-rhe}
Reprenons les notations et les hypothèses du lemme \ref{curves-lemma-rh}. Pour un point fermé $x \in X$, soit $d_x$ la multiplicité de $x$ dans $R$. Alors $$
2g_X - 2 = (2g_Y - 2) \deg(f) + \sum\nolimits d_x [\kappa(x) : k]
$$ De plus, on a les résultats suivants : \begin{enumerate} \item $d_x = \text{length}_{\mathcal{O}_{X, x}}(\Omega_{X/Y, x})$; \item $d_x \geq e_x - 1$, où $e_x$ est l'indice de ramification de $\mathcal{O}_{X, x}$ au-dessus de $\mathcal{O}_{Y, y}$; \item $d_x = e_x - 1$ si et seulement si $\mathcal{O}_{X, x}$ est modérément ramifié au-dessus de $\mathcal{O}_{Y, y}$.
\end{enumerate}
\end{lemma}

\begin{proof}
D'après le lemme \ref{curves-lemma-rh} et la discussion ci-dessus (où l'on a utilisé Variétés, lemme \ref{varieties-lemma-algebraic-scheme-dim-0} et Algèbre, lemme \ref{algebra-lemma-pushdown-module}), il suffit de démontrer les assertions concernant la multiplicité $d_x$ de $x$ dans $R$. L'assertion (1) a été établie dans cette discussion. Nous n'y avons démontré (2) et (3) que lorsque $\kappa(x)$ est séparable sur $k$. Dans la suite, nous traiterons (2) et (3) de façon uniforme à l'aide des résultats sur les différentes des morphismes quasi-finis de Gorenstein.

\medskip\noindent
Remarquons d'abord que $f$ est un morphisme quasi-fini de Gorenstein. Cela résulte, par exemple, du fait que $f$ est un morphisme plat quasi-fini et que $X$ est de Gorenstein (voir Dualité pour les schémas, lemme \ref{duality-lemma-flat-morphism-from-gorenstein-scheme}); cela a aussi été démontré dans la preuve de Discriminants, lemme \ref{discriminant-lemma-discriminant-quasi-finite-morphism-smooth} (utilisé ci-dessus). Ainsi $\omega_{X/Y}$ est inversible d'après Discriminants, lemme \ref{discriminant-lemma-gorenstein-quasi-finite}, et il le reste après remplacement de $X$ par des ouverts et après changement de base par un morphisme $Y' \to Y$. Nous utiliserons ce fait sans autre mention.

\medskip\noindent
Choisissons des ouverts affines $U \subset X$ et $V \subset Y$ tels que $x \in U$, $y \in V$, $f(U) \subset V$, et tels que $x$ soit le seul point de $U$ au-dessus de $y$. Écrivons $U = \Spec(A)$ et $V = \Spec(B)$. Alors $R \cap U$ est la différente de $f|_U : U \to V$. D'après Discriminants, lemme \ref{discriminant-lemma-base-change-different}, sa formation commute dans notre cas à tout changement de base. Par le choix de $U$ et $V$, on a $$
A \otimes_B \kappa(y) =
\mathcal{O}_{X, x} \otimes_{\mathcal{O}_{Y, y}} \kappa(y) =
\mathcal{O}_{X, x}/(s^{e_x})
$$, où $e_x$ est l'indice de ramification de l'énoncé. Soit $C = \mathcal{O}_{X, x}/(s^{e_x})$, considéré comme une algèbre finie sur $\kappa(y)$. Soit $\mathfrak{D}_{C/\kappa(y)}$ la différente de $C$ sur $\kappa(y)$ au sens de Discriminants, définition \ref{discriminant-definition-different}. Il suffit de montrer que $\mathfrak{D}_{C/\kappa(y)}$ est non nulle si et seulement si l'extension $\mathcal{O}_{Y, y} \subset \mathcal{O}_{X, x}$ est modérément ramifiée et que, dans ce cas, $\mathfrak{D}_{C/\kappa(y)}$ est l'idéal engendré par $s^{e_x - 1}$ dans $C$. Rappelons que la ramification modérée signifie exactement que $\kappa(x)/\kappa(y)$ est séparable et que la caractéristique de $\kappa(y)$ ne divise pas $e_x$. D'autre part, la différente de $C/\kappa(y)$ est non nulle si et seulement si $\tau_{C/\kappa(y)} \in \omega_{C/\kappa(y)}$ est non nulle. En effet, puisque $\omega_{C/\kappa(y)}$ est un $C$-module inversible (obtenu par changement de base à partir de $\omega_{A/B}$), il est libre de rang $1$, avec un générateur que nous noterons $\lambda$. Écrivons $\tau_{C/\kappa(y)} = h\lambda$ pour un certain $h \in C$. Alors $\mathfrak{D}_{C/\kappa(y)} = (h) \subset C$, d'où l'assertion. D'après Discriminants, lemme \ref{discriminant-lemma-tau-nonzero}, on a $\tau_{C/\kappa(y)} \not = 0$ si et seulement si $\kappa(x)/\kappa(y)$ est séparable et $e_x$ premier à la caractéristique. Enfin, même si $\tau_{C/\kappa(y)}$ est non nulle, on a encore $s \tau_{C/\kappa(y)} = 0$, car $s\tau_{C/\kappa(y)} : C \to \kappa(y)$ envoie $c$ sur la trace de l'opérateur nilpotent $sc$, laquelle est nulle. Ainsi $sh = 0$, puis $h \in (s^{e_x - 1})$, ce qui prouve que $\mathfrak{D}_{C/\kappa(y)} \subset (s^{e_x - 1})$ dans tous les cas. Puisque $(s^{e_x - 1}) \subset C$ est le plus petit idéal non nul, la dernière assertion est démontrée.
\end{proof}






\section{Morphismes inséparables}
\label{curves-section-inseparable}

\noindent
Quelques remarques sur le comportement du genre par les applications inséparables.

\begin{lemma}
\label{curves-lemma-dominated-by-smooth}
Soit $k$ un corps. Soit $f : X \to Y$ un morphisme surjectif de courbes sur $k$. Si $X$ est lisse sur $k$ et si $Y$ est normale, alors $Y$ est lisse sur $k$.
\end{lemma}

\begin{proof}
Soit $y \in Y$. Choisissons $x \in X$ d'image $y$. D'après Variétés, lemme \ref{varieties-lemma-flat-under-smooth}, il suffit de montrer que $f$ est plat en $x$. Cela résulte du lemme \ref{curves-lemma-flat}.
\end{proof}

\begin{lemma}
\label{curves-lemma-purely-inseparable}
Soit $k$ un corps de caractéristique $p > 0$. Soit $f : X \to Y$ un morphisme non constant de courbes propres non singulières sur $k$. Si l'extension $k(X)/k(Y)$ de corps de fonctions est purement inséparable, il existe une factorisation $$
X = X_0 \to X_1 \to \ldots \to X_n = Y
$$ telle que chaque $X_i$ soit une courbe propre non singulière et que $X_i \to X_{i + 1}$ soit un morphisme de degré $p$ avec $k(X_{i + 1}) \subset k(X_i)$ inséparable.
\end{lemma}

\begin{proof}
Cela résulte du théorème \ref{curves-theorem-curves-rational-maps} et du fait qu'une extension finie purement inséparable de corps s'obtient toujours comme une succession d'extensions (inséparables) de degré $p$; voir Corps, lemme \ref{fields-lemma-finite-purely-inseparable}.
\end{proof}

\begin{lemma}
\label{curves-lemma-inseparable-deg-p-smooth}
Soit $k$ un corps de caractéristique $p > 0$. Soit $f : X \to Y$ un morphisme non constant de courbes propres non singulières sur $k$. Si $X$ est lisse et si $k(Y) \subset k(X)$ est inséparable de degré $p$, il existe un unique isomorphisme $Y = X^{(p)}$ tel que $f$ soit $F_{X/k}$.
\end{lemma}

\begin{proof}
Le morphisme de Frobenius relatif $F_{X/k} : X \to X^{(p)}$ est construit dans Variétés, section \ref{varieties-section-frobenius}. Remarquons que $X^{(p)}$ est une courbe propre lisse sur $k$, obtenue par changement de base à partir de $X$. Le morphisme $F_{X/k}$ est de degré $p$ d'après Variétés, lemme \ref{varieties-lemma-inseparable-deg-p-smooth}. Ainsi $k(X^{(p)})$ et $k(Y)$ sont deux sous-corps de $k(X)$ tels que $[k(X) : k(Y)] = [k(X) : k(X^{(p)})] = p$. Pour démontrer le lemme, il suffit de montrer que $k(Y) = k(X^{(p)})$ dans $k(X)$. Voir le théorème \ref{curves-theorem-curves-rational-maps}.

\medskip\noindent
Posons $K = k(X)$. Considérons l'application $\text{d} : K \to \Omega_{K/k}$. Il résulte du lemme \ref{curves-lemma-generically-etale} que $k(Y)$ est contenu dans le noyau de $\text{d}$. D'après Variétés, lemme \ref{varieties-lemma-relative-frobenius-omega}, on voit que $k(X^{(p)})$ est dans le noyau de $\text{d}$. Puisque $X$ est une courbe lisse, $\Omega_{K/k}$ est un espace vectoriel de dimension $1$ sur $K$. Compléments d'algèbre, lemme \ref{more-algebra-lemma-p-basis} entraîne alors que $\Ker(\text{d}) = kK^p$ et que $[K : kK^p] = p$. Ainsi $k(Y) = kK^p = k(X^{(p)})$ pour des raisons de degré.
\end{proof}

\begin{lemma}
\label{curves-lemma-purely-inseparable-smooth}
Soit $k$ un corps de caractéristique $p > 0$. Soit $f : X \to Y$ un morphisme non constant de courbes propres non singulières sur $k$. Si $X$ est lisse et si $k(Y) \subset k(X)$ est purement inséparable, il existe un unique $n \geq 0$ et un unique isomorphisme $Y = X^{(p^n)}$ tels que $f$ soit le Frobenius relatif itéré $n$ fois de $X/k$.
\end{lemma}

\begin{proof}
Le Frobenius relatif itéré $n$ fois de $X/k$ est défini dans Variétés, remarque \ref{varieties-remark-n-fold-relative-frobenius}. Le lemme s'obtient en combinant les lemmes \ref{curves-lemma-inseparable-deg-p-smooth} et \ref{curves-lemma-purely-inseparable}.
\end{proof}

\begin{lemma}
\label{curves-lemma-purely-inseparable-smooth-genus}
Soit $k$ un corps de caractéristique $p > 0$. Soit $f : X \to Y$ un morphisme non constant de courbes propres non singulières sur $k$. Supposons que \begin{enumerate} \item $X$ soit lisse, \item $H^0(X, \mathcal{O}_X) = k$, \item $k(X)/k(Y)$ soit purement inséparable. \end{enumerate} Alors $Y$ est lisse, $H^0(Y, \mathcal{O}_Y) = k$ et le genre de $Y$ est égal au genre de $X$.
\end{lemma}

\begin{proof}
D'après le lemme \ref{curves-lemma-purely-inseparable-smooth}, $Y = X^{(p^n)}$ s'obtient par changement de base à partir de $X$ par $F_{\Spec(k)}^n$. Ainsi $Y$ est lisse, et les assertions sur la cohomologie et le genre résultent du lemme \ref{curves-lemma-genus-base-change}.
\end{proof}

\begin{example}
\label{curves-example-inseparable}
Cet exemple montre que le genre peut changer par un morphisme purement inséparable de courbes projectives non singulières. Soit $k$ un corps de caractéristique $3$. Supposons qu'il existe un élément $a \in k$ qui ne soit pas un cube, c'est-à-dire une puissance $3$-ième. Par exemple, $k = \mathbf{F}_3(a)$ convient. Soit $X$ la courbe plane d'équation homogène $$
F = T_1^2T_0 - T_2^3 + aT_0^3
$$ comme dans la section \ref{curves-section-plane-curves}. Sur l'ouvert affine $D_+(T_0)$, avec les coordonnées $x = T_1/T_0$ et $y = T_2/T_0$, on obtient $x^2 - y^3 + a = 0$, qui définit une courbe affine non singulière. De plus, le point à l'infini $(0 : 1: 0)$ est lisse. Ainsi $X$ est une courbe projective non singulière de genre $1$ (lemme \ref{curves-lemma-genus-plane-curve}). Considérons d'autre part le morphisme $f : X \to \mathbf{P}^1_k$ qui, sur $D_+(T_0)$, envoie $(x, y)$ sur $x \in \mathbf{A}^1_k \subset \mathbf{P}^1_k$. Alors $f$ est un morphisme de courbes propres non singulières sur $k$ induisant une extension inséparable de corps de fonctions de degré $p = 3$, mais le genre de $X$ est $1$ et celui de $\mathbf{P}^1_k$ est $0$.
\end{example}

\begin{proposition}
\label{curves-proposition-unwind-morphism-smooth}
Soit $k$ un corps de caractéristique $p > 0$. Soit $f : X \to Y$ un morphisme non constant de courbes propres lisses sur $k$. On peut factoriser $f$ sous la forme $$
X \longrightarrow X^{(p^n)} \longrightarrow Y
$$, où $X^{(p^n)} \to Y$ est un morphisme non constant de courbes propres lisses induisant une extension séparable $k(X^{(p^n)})/k(Y)$, où $$
X^{(p^n)} = X \times_{\Spec(k), F_{\Spec(k)}^n} \Spec(k),
$$ et où $X \to X^{(p^n)}$ est le Frobenius relatif itéré $n$ fois de $X$.
\end{proposition}

\begin{proof}
D'après Corps, lemme \ref{fields-lemma-separable-first}, il existe une extension intermédiaire $k(X)/E/k(Y)$ telle que $k(X)/E$ soit purement inséparable et $E/k(Y)$ séparable. Par le théorème \ref{curves-theorem-curves-rational-maps}, elle correspond à une factorisation $X \to Z \to Y$ de $f$, où $Z$ est une courbe propre non singulière. On conclut en appliquant le lemme \ref{curves-lemma-purely-inseparable-smooth} au morphisme $X \to Z$.
\end{proof}

\begin{lemma}
\label{curves-lemma-inseparable-linear-system}
Soit $k$ un corps de caractéristique $p > 0$. Soit $X$ une courbe propre lisse sur $k$. Soit $(\mathcal{L}, V)$ une $\mathfrak g^r_d$ avec $r \geq 1$. Alors l'une des deux assertions suivantes est vraie : \begin{enumerate} \item il existe une $\mathfrak g^1_d$ dont le morphisme correspondant $X \to \mathbf{P}^1_k$ (lemme \ref{curves-lemma-linear-series}) est génériquement \'etale (c'est-à-dire satisfait au lemme \ref{curves-lemma-generically-etale}); \item il existe une $\mathfrak g^r_{d'}$ sur $X^{(p)}$, où $d' \leq d/p$.
\end{enumerate}
\end{lemma}

\begin{proof}
Choisissons deux éléments linéairement indépendants sur $k$, notés $s, t \in V$. Alors $f = s/t$ est la fonction rationnelle définissant le morphisme $X \to \mathbf{P}^1_k$ correspondant à la série linéaire $(\mathcal{L}, ks + kt)$. Si ce morphisme n'est pas génériquement \'etale, alors $f \in k(X^{(p)})$ d'après la proposition \ref{curves-proposition-unwind-morphism-smooth}. Choisissons maintenant une base $s_0, \ldots, s_r$ de $V$ et soit $\mathcal{L}' \subset \mathcal{L}$ le faisceau inversible engendré par $s_0, \ldots, s_r$. Posons $f_i = s_i/s_0$ dans $k(X)$. Si, pour chaque couple $(s_0, s_i)$, on a $f_i \in k(X^{(p)})$, le morphisme $$
\varphi = \varphi_{(\mathcal{L}', (s_0, \ldots, s_r))} :
X
\longrightarrow
\mathbf{P}^r_k = \text{Proj}(k[T_0, \ldots, T_r])
$$ se factorise par $X^{(p)}$, car c'est vrai au-dessus de l'ouvert affine $D_+(T_0)$, et le morphisme sur la partie affine se prolonge à la courbe lisse $X^{(p)}$ tout entière d'après le lemme \ref{curves-lemma-extend-over-normal-curve}. Notons cette factorisation $$
X \xrightarrow{F_{X/k}} X^{(p)} \xrightarrow{\psi} \mathbf{P}^r_k
$$ de $\varphi$. Alors $\mathcal{N} = \psi^*\mathcal{O}_{\mathbf{P}^r_k}(1)$ est un $\mathcal{O}_{X^{(p)}}$-module inversible tel que $\mathcal{L}' = F_{X/k}^*\mathcal{N}$ et tel que $\psi^*T_0, \ldots, \psi^*T_r$ soient linéairement indépendantes sur $k$ (leurs images réciproques étant $s_0, \ldots, s_r$ sur $X$). Enfin, on a $$
d = \deg(\mathcal{L}) \geq \deg(\mathcal{L}') =
\deg(F_{X/k}) \deg(\mathcal{N}) = p \deg(\mathcal{N})
$$, comme voulu. Nous avons utilisé ici Variétés, lemmes \ref{varieties-lemma-check-invertible-sheaf-trivial}, \ref{varieties-lemma-degree-pullback-map-proper-curves} et \ref{varieties-lemma-inseparable-deg-p-smooth}.
\end{proof}

\begin{lemma}
\label{curves-lemma-point-over-separable-extension}
Soit $k$ un corps. Soit $X$ une courbe propre lisse sur $k$ telle que $H^0(X, \mathcal{O}_X) = k$ et de genre $g \geq 2$. Il existe alors un point fermé $x \in X$ tel que $\kappa(x)/k$ soit séparable de degré $\leq 2g - 2$.
\end{lemma}

\begin{proof}
Posons $\omega = \Omega_{X/k}$. D'après le lemme \ref{curves-lemma-genus-smooth}, ce module est de degré $2g - 2$ et possède $g$ sections globales. Nous disposons donc d'une $\mathfrak g^{g - 1}_{2g - 2}$. Le lemme élémentaire \ref{curves-lemma-linear-series-trivial-existence} donne une $\mathfrak g^1_{2g - 2}$ et le lemme \ref{curves-lemma-g1d} fournit un morphisme $$
\varphi : X \longrightarrow \mathbf{P}^1_k
$$ d'un certain degré $d \leq 2g - 2$. Puisque $\varphi$ est plat (lemme \ref{curves-lemma-flat}) et fini (lemme \ref{curves-lemma-finite}), il est fini localement libre de degré $d$ (Morphismes, lemme \ref{morphisms-lemma-finite-flat}). Choisissons un point rationnel $t \in \mathbf{P}^1_k$ et un point $x \in X$ tel que $\varphi(x) = t$. Alors $$
d \geq [\kappa(x) : \kappa(t)] = [\kappa(x) : k]
$$, par exemple d'après Morphismes, lemmes \ref{morphisms-lemma-finite-locally-free-universally-bounded} et \ref{morphisms-lemma-characterize-universally-bounded}. Si $k$ est parfait (par exemple de caractéristique zéro ou fini), le lemme est donc démontré. On est ainsi ramené au cas du paragraphe suivant.

\medskip\noindent
Supposons que $k$ soit un corps infini de caractéristique $p > 0$. Comme ci-dessus, nous utiliserons le fait que $X$ possède une $\mathfrak g^{g - 1}_{2g - 2}$. La courbe propre lisse $X^{(p)}$ a le même genre que $X$. Son genre est donc $> 0$. On en conclut que $X^{(p)}$ ne possède aucune $\mathfrak g^{g - 1}_d$ pour $d \leq g - 1$ d'après le lemme \ref{curves-lemma-grd-inequalities}. En appliquant le lemme \ref{curves-lemma-inseparable-linear-system} à notre $\mathfrak g^{g - 1}_{2g - 2}$ (et en remarquant que $(2g - 2)/p \leq g - 1$), on conclut que l'éventualité (2) ne se produit pas. On obtient donc un morphisme $$
\varphi : X \longrightarrow \mathbf{P}^1_k
$$ génériquement \'etale (au sens du lemme) et de degré $\leq 2g - 2$. Soit $U \subset X$ le sous-schéma ouvert non vide où $\varphi$ est \'etale. Alors $\varphi(U) \subset \mathbf{P}^1_k$ est un ouvert de Zariski non vide et l'on peut choisir un point $k$-rationnel $t \in \varphi(U)$, puisque $k$ est infini. Soit $u \in U$ un point tel que $\varphi(u) = t$. Alors $\kappa(u)/\kappa(t)$ est séparable (Morphismes, lemme \ref{morphisms-lemma-etale-over-field}), $\kappa(t) = k$ et $[\kappa(u) : k] \leq 2g - 2$ comme précédemment.
\end{proof}

\noindent
Le lemme suivant n'appartient pas vraiment à cette section, mais nous ne lui connaissons pas de meilleure place.

\begin{lemma}
\label{curves-lemma-ramification-to-algebraic-closure}
Soit $X$ une courbe lisse sur un corps $k$. Soit $\overline{x} \in X_{\overline{k}}$ un point fermé d'image $x \in X$. L'indice de ramification de $\mathcal{O}_{X, x} \subset \mathcal{O}_{X_{\overline{k}}, \overline{x}}$ est le degré d'inséparabilité de $\kappa(x)/k$.
\end{lemma}

\begin{proof}
Quitte à restreindre $X$, on peut supposer qu'il existe un morphisme \'etale $\pi : X \to \mathbf{A}^1_k$; voir Morphismes, lemme \ref{morphisms-lemma-smooth-etale-over-affine-space}. On peut alors considérer le diagramme d'anneaux locaux $$
\xymatrix{
\mathcal{O}_{X_{\overline{k}}, \overline{x}} &
\mathcal{O}_{\mathbf{A}^1_{\overline{k}}, \pi(\overline{x})} \ar[l] \\
\mathcal{O}_{X, x} \ar[u] &
\mathcal{O}_{\mathbf{A}^1_k, \pi(x)} \ar[l] \ar[u]
}
$$ Les flèches horizontales sont d'indice de ramification $1$, car elles correspondent à des morphismes \'etales. De plus, l'extension $\kappa(x)/\kappa(\pi(x))$ est séparable; ainsi $\kappa(x)$ et $\kappa(\pi(x))$ ont le même degré d'inséparabilité sur $k$. Par multiplicativité des indices de ramification, il suffit de démontrer le résultat lorsque $x$ est un point de la droite affine.

\medskip\noindent
Supposons $X = \mathbf{A}^1_k$. Dans ce cas, l'anneau local de $X$ en $x$ est de la forme $$
\mathcal{O}_{X, x} = k[t]_{(P)}
$$, où $P$ est un polynôme unitaire irréductible sur $k$. Alors $P(t) = Q(t^q)$ pour un polynôme séparable $Q \in k[t]$; voir Corps, lemme \ref{fields-lemma-irreducible-polynomials}. Remarquons que $\kappa(x) = k[t]/(P)$ est de degré d'inséparabilité $q$ sur $k$. D'autre part, sur $\overline{k}$, on peut factoriser $Q(t) = \prod (t - \alpha_i)$ avec les $\alpha_i$ deux à deux distincts. Écrivons $\alpha_i = \beta_i^q$ pour un unique $\beta_i \in \overline{k}$. Notre point $\overline{x}$ correspond alors à l'un des $\beta_i$; on conclut, car l'indice de ramification de $$
k[t]_{(P)} \longrightarrow \overline{k}[t]_{(t - \beta_i)}
$$ est bien égal à $q$, puisque l'uniformisante $P$ s'envoie sur $(t - \beta_i)^q$ multiplié par une unité.
\end{proof}





\section{Sommes amalgamées}
\label{curves-section-pushouts}

\noindent
Soit $k$ un corps. Considérons un diagramme en traits pleins $$
\xymatrix{
Z' \ar[d] \ar[r]_{i'} & X' \ar@{..>}[d]^a \\
Z \ar@{..>}[r]^i & X
}
$$ de schémas sur $k$ satisfaisant aux conditions suivantes : \begin{enumerate} \item[(a)] $X'$ est séparé de type fini sur $k$, de dimension $\leq 1$; \item[(b)] $i' : Z' \to X'$ est une immersion fermée; \item[(c)] $Z'$ et $Z$ sont finis sur $\Spec(k)$; \item[(d)] $Z' \to Z$ est surjectif. \end{enumerate} Dans cette situation, toute partie finie de $X'$ est contenue dans un ouvert affine; voir Variétés, proposition \ref{varieties-proposition-finite-set-of-points-of-codim-1-in-affine}. Les hypothèses de Compléments sur les morphismes, proposition \ref{more-morphisms-proposition-pushout-along-closed-immersion-and-integral} sont donc satisfaites, et l'on obtient les propriétés suivantes : \begin{enumerate} \item la somme amalgamée $X = Z \amalg_{Z'} X'$ existe dans la catégorie des schémas; \item $i : Z \to X$ est une immersion fermée; \item $a : X' \to X$ est entier et surjectif; \item $X \to \Spec(k)$ est séparé d'après Compléments sur les morphismes, lemme \ref{more-morphisms-lemma-pushout-separated}; \item $X \to \Spec(k)$ est de type fini d'après Compléments sur les morphismes, lemmes \ref{more-morphisms-lemma-pushout-finite-type}; \item ainsi $a : X' \to X$ est fini d'après Morphismes, lemmes \ref{morphisms-lemma-finite-integral} et \ref{morphisms-lemma-permanence-finite-type}; \item si $X' \to \Spec(k)$ est propre, alors $X \to \Spec(k)$ est propre d'après Morphismes, lemme \ref{morphisms-lemma-image-proper-is-proper}.
\end{enumerate}

\noindent
Le lemme suivant admet des généralisations considérables.

\begin{lemma}
\label{curves-lemma-complete-local-ring-pushout}
Dans la situation ci-dessus, supposons $Z = \Spec(k')$, où $k'$ est un corps, et $Z' = \Spec(k'_1 \times \ldots \times k'_n)$, où les $k'_i/k'$ sont des extensions finies de corps. Soient $x \in X$ l'image de $Z \to X$ et $x'_i \in X'$ l'image de $\Spec(k'_i) \to X'$. On a alors un diagramme de produit fibré $$
\xymatrix{
\prod\nolimits_{i = 1, \ldots, n} k'_i &
\prod\nolimits_{i = 1, \ldots, n} \mathcal{O}_{X', x'_i}^\wedge \ar[l] \\
k' \ar[u] &
\mathcal{O}_{X, x}^\wedge \ar[u] \ar[l]
}
$$ dans lequel les flèches horizontales sont les applications vers les corps résiduels.
\end{lemma}

\begin{proof}
Choisissons un voisinage ouvert affine $\Spec(A)$ de $x$ dans $X$. Soit $\Spec(A') \subset X'$ son image réciproque. Par construction, on a un diagramme de produit fibré $$
\xymatrix{
\prod\nolimits_{i = 1, \ldots, n} k'_i &
A' \ar[l] \\
k' \ar[u] &
A \ar[u] \ar[l]
}
$$ Puisque tous les objets sont finis sur $A$, ce diagramme reste un diagramme de produit fibré après complétion par rapport à l'idéal maximal $\mathfrak m \subset A$ correspondant à $x$ (Algèbre, lemme \ref{algebra-lemma-completion-flat}). Enfin, on applique Algèbre, lemme \ref{algebra-lemma-completion-finite-extension} pour identifier le complété de $A'$.
\end{proof}




\section{Identification de points et écrasement de vecteurs tangents}
\sectionmark{Identification de points et vecteurs tangents}
\label{curves-section-glueing-squishing}

\noindent
Dans la suite, nous noterons $k[\epsilon]$ l'algèbre des nombres duaux sur $k$ définie dans Variétés, définition \ref{varieties-definition-dual-numbers}.

\begin{lemma}
\label{curves-lemma-no-in-between-over-k}
Soit $k$ un corps algébriquement clos. Soit $k \subset A$ une extension d'anneaux telle que $A$ possède exactement deux sous-algèbres sur $k$. Alors on a soit $A = k \times k$, soit $A = k[\epsilon]$.
\end{lemma}

\begin{proof}
L'hypothèse signifie que $k \not = A$ et que tout sous-anneau $k \subset C \subset A$ est égal à $k$ ou à $A$. Soit $t \in A$ avec $t \not \in k$. Alors $A$ est engendré par $t$ sur $k$. Ainsi $A = k[x]/I$ pour un idéal $I$. Si $I = (0)$, on dispose de la sous-algèbre $k[x^2]$, ce qui est exclu. Sinon, $I$ est engendré par un polynôme unitaire $P$. Écrivons $P = \prod_{i = 1}^d (x - a_i)$. Si $d > 2$, la sous-algèbre engendrée par $(t - a_1)(t - a_2)$ donne une contradiction. Ainsi $d = 2$. Si $a_1 \not = a_2$, alors $A = k \times k$; si $a_1 = a_2$, alors $A = k[\epsilon]$.
\end{proof}

\begin{example}[Identification de points] \label{curves-example-glue-points} Soit $k$ un corps algébriquement clos. Soit $f : X' \to X$ un morphisme de schémas algébriques sur $k$. On dit que $X$ est obtenu en identifiant $a$ et $b$ dans $X'$ si les conditions suivantes sont satisfaites : \begin{enumerate} \item $a, b \in X'(k)$ sont des points distincts ayant pour image un même point $x \in X(k)$; \item $f$ est fini et $f^{-1}(X \setminus \{x\}) \to X \setminus \{x\}$ est un isomorphisme; \item on a une suite exacte courte $$
0 \to \mathcal{O}_X \to f_*\mathcal{O}_{X'} \xrightarrow{a - b} x_*k \to 0
$$ dont la flèche de droite envoie une section locale $h$ de $f_*\mathcal{O}_{X'}$ sur la différence $h(a) - h(b) \in k$. \end{enumerate} Dans ce cas, on a aussi une suite exacte courte $$
0 \to \mathcal{O}_X^* \to f_*\mathcal{O}_{X'}^*
\xrightarrow{ab^{-1}} x_*k^* \to 0
$$ dont la flèche de droite envoie une section locale $h$ de $f_*\mathcal{O}_{X'}^*$ sur la différence multiplicative $h(a)h(b)^{-1} \in k^*$.
\end{example}

\begin{example}[Écrasement d'un vecteur tangent] \label{curves-example-squish-tangent-vector} Soit $k$ un corps algébriquement clos. Soit $f : X' \to X$ un morphisme de schémas algébriques sur $k$. On dit que $X$ est obtenu en écrasant le vecteur tangent $\vartheta$ dans $X'$ si les conditions suivantes sont satisfaites : \begin{enumerate} \item $\vartheta : \Spec(k[\epsilon]) \to X'$ est une immersion fermée sur $k$ telle que $f \circ \vartheta$ se factorise par un point $x \in X(k)$; \item $f$ est fini et $f^{-1}(X \setminus \{x\}) \to X \setminus \{x\}$ est un isomorphisme; \item on a une suite exacte courte $$
0 \to \mathcal{O}_X \to f_*\mathcal{O}_{X'} \xrightarrow{\vartheta} x_*k \to 0
$$ dont la flèche de droite envoie une section locale $h$ de $f_*\mathcal{O}_{X'}$ sur le coefficient de $\epsilon$ dans $\vartheta^\sharp(h) \in k[\epsilon]$. \end{enumerate} Dans ce cas, on a aussi une suite exacte courte $$
0 \to \mathcal{O}_X^* \to f_*\mathcal{O}_{X'}^*
\xrightarrow{\vartheta} x_*k \to 0
$$ dont la flèche de droite envoie une section locale $h$ de $f_*\mathcal{O}_{X'}^*$ sur $\text{d}\log(\vartheta^\sharp(h))$, où $\text{d}\log : k[\epsilon]^* \to k$ est l'homomorphisme de groupes abéliens envoyant $a + b\epsilon$ sur $b/a \in k$.
\end{example}

\begin{lemma}
\label{curves-lemma-factor-almost-isomorphism}
Soit $k$ un corps algébriquement clos. Soit $f : X' \to X$ un morphisme fini de schémas algébriques sur $k$ tel que $\mathcal{O}_X \subset f_*\mathcal{O}_{X'}$ et tel que $f$ soit un isomorphisme hors d'un ensemble fini de points. Il existe alors une factorisation $$
X' = X_n \to X_{n - 1} \to \ldots \to X_1 \to X_0 = X
$$ telle que chaque $X_i \to X_{i - 1}$ soit soit une identification de deux points, soit l'écrasement d'un vecteur tangent (voir les exemples \ref{curves-example-glue-points} et \ref{curves-example-squish-tangent-vector}).
\end{lemma}

\begin{proof}
Soit $U \subset X$ le plus grand ouvert au-dessus duquel $f$ est un isomorphisme. Alors $X \setminus U = \{x_1, \ldots, x_n\}$ avec $x_i \in X(k)$. Nous considérerons les factorisations $X' \to Y \to X$ de $f$ telles que les deux morphismes soient finis et que $$
\mathcal{O}_X \subset g_*\mathcal{O}_Y \subset f_*\mathcal{O}_{X'}
$$, où $g : Y \to X$ est le morphisme donné. Par hypothèse, $\mathcal{O}_{X, x} \to (f_*\mathcal{O}_{X'})_x$ est un isomorphisme sauf si $x = x_i$ pour un certain $i$. Le conoyau $$
f_*\mathcal{O}_{X'}/\mathcal{O}_X = \bigoplus \mathcal{Q}_i
$$ est donc une somme directe de faisceaux gratte-ciel $\mathcal{Q}_i$ de supports $x_1, \ldots, x_n$. Puisque le quotient affiché est un $\mathcal{O}_X$-module cohérent, $\mathcal{Q}_i$ est de longueur finie sur $\mathcal{O}_{X, x_i}$. On peut donc raisonner par récurrence sur la somme de ces longueurs, c'est-à-dire sur la longueur du conoyau tout entier.

\medskip\noindent
Si $n > 1$, on peut définir une $\mathcal{O}_X$-sous-algèbre, notée $\mathcal{A} \subset f_*\mathcal{O}_{X'}$, en prenant l'image réciproque de $\mathcal{Q}_1$. On obtient ainsi une factorisation non triviale et l'on conclut par récurrence.

\medskip\noindent
Supposons $n = 1$. Abrégeons la notation en posant $x = x_1$. Considérons l'extension finie de $k$-algèbres $$
A = \mathcal{O}_{X, x} \subset (f_*\mathcal{O}_{X'})_x = B
$$ Remarquons que $\mathcal{Q} = \mathcal{Q}_1$ est le faisceau gratte-ciel de valeur $B/A$. On dispose d'une sous-algèbre sur $k$, à savoir $A \subset A + \mathfrak m_A B \subset B$. Si les deux inclusions sont strictes, on obtient une factorisation non triviale et l'on conclut par récurrence comme ci-dessus. Si $A + \mathfrak m_A B = B$, alors $A = B$ par Nakayama; ainsi $f$ est un isomorphisme et il n'y a rien à démontrer. On peut donc supposer $A = A + \mathfrak m_A B$. Posons $C = B/\mathfrak m_A B$. Si $C$ possède plus de $2$ sous-algèbres sur $k$, on obtient une sous-algèbre strictement intermédiaire entre $A$ et $B$ en prenant l'image réciproque dans $B$. On peut donc supposer que $C$ possède exactement $2$ sous-algèbres sur $k$. Alors $C = k \times k$ ou $C = k[\epsilon]$ d'après le lemme \ref{curves-lemma-no-in-between-over-k}. Dans ces deux cas, $f$ est respectivement l'identification de deux points ou l'écrasement d'un vecteur tangent.
\end{proof}

\begin{lemma}
\label{curves-lemma-glue-points}
Soit $k$ un corps algébriquement clos. Si $f : X' \to X$ est l'identification de deux points $a, b$ comme dans l'exemple \ref{curves-example-glue-points}, on a une suite exacte $$
k^* \to \Pic(X) \to \Pic(X') \to 0
$$ La première application est nulle si $a$ et $b$ appartiennent à des composantes connexes distinctes de $X'$, et injective si $X'$ est propre et si $a$ et $b$ appartiennent à la même composante connexe de $X'$.
\end{lemma}

\begin{proof}
L'application $\Pic(X) \to \Pic(X')$ est surjective d'après Variétés, lemme \ref{varieties-lemma-surjective-pic-birational-finite}. La suite exacte courte $$
0 \to \mathcal{O}_X^* \to f_*\mathcal{O}_{X'}^*
\xrightarrow{ab^{-1}} x_*k^* \to 0
$$ donne $$
H^0(X', \mathcal{O}_{X'}^*) \xrightarrow{ab^{-1}} k^* \to
H^1(X, \mathcal{O}_X^*) \to H^1(X, f_*\mathcal{O}_{X'}^*)
$$ On a $H^1(X, f_*\mathcal{O}_{X'}^*) \subset H^1(X', \mathcal{O}_{X'}^*)$ (par exemple par la suite spectrale de Leray; voir Cohomologie, lemme \ref{cohomology-lemma-Leray}). Le noyau de $\Pic(X) \to \Pic(X')$ est donc le conoyau de $ab^{-1} : H^0(X', \mathcal{O}_{X'}^*) \to k^*$. Si $a$ et $b$ appartiennent à des composantes connexes distinctes de $X'$, alors $ab^{-1}$ est surjective. Puisque $k$ est algébriquement clos, toute fonction régulière sur un schéma réduit connexe propre sur $k$ provient d'un élément de $k$; voir Variétés, lemme \ref{varieties-lemma-proper-geometrically-reduced-global-sections}. Ainsi $ab^{-1}$ est nulle si $X'$ est propre et si $a$ et $b$ appartiennent à la même composante connexe.
\end{proof}

\begin{lemma}
\label{curves-lemma-squish-tangent-vector}
Soit $k$ un corps algébriquement clos. Si $f : X' \to X$ est l'écrasement d'un vecteur tangent $\vartheta$ comme dans l'exemple \ref{curves-example-squish-tangent-vector}, on a une suite exacte $$
(k, +) \to \Pic(X) \to \Pic(X') \to 0
$$ dont la première application est injective si $X'$ est propre et réduit.
\end{lemma}

\begin{proof}
L'application $\Pic(X) \to \Pic(X')$ est surjective d'après Variétés, lemme \ref{varieties-lemma-surjective-pic-birational-finite}. La suite exacte courte $$
0 \to \mathcal{O}_X^* \to f_*\mathcal{O}_{X'}^*
\xrightarrow{\vartheta} x_*k \to 0
$$ de l'exemple \ref{curves-example-squish-tangent-vector} donne $$
H^0(X', \mathcal{O}_{X'}^*) \xrightarrow{\vartheta} k \to
H^1(X, \mathcal{O}_X^*) \to H^1(X, f_*\mathcal{O}_{X'}^*)
$$ On a $H^1(X, f_*\mathcal{O}_{X'}^*) \subset H^1(X', \mathcal{O}_{X'}^*)$ (par exemple par la suite spectrale de Leray; voir Cohomologie, lemme \ref{cohomology-lemma-Leray}). Le noyau de $\Pic(X) \to \Pic(X')$ est donc le conoyau de l'application $\vartheta : H^0(X', \mathcal{O}_{X'}^*) \to k$. Puisque $k$ est algébriquement clos, toute fonction régulière sur un schéma réduit connexe propre sur $k$ provient d'un élément de $k$; voir Variétés, lemme \ref{varieties-lemma-proper-geometrically-reduced-global-sections}. D'où la dernière assertion du lemme.
\end{proof}


\section{Singularités en croisement multiple et singularités nodales}
\label{curves-section-multicross}

\noindent
Dans cette section, nous étudions les singularités de courbes les plus simples possibles.

\medskip\noindent
Soit $k$ un corps. Considérons la $k$-algèbre locale complète \begin{equation}
\label{curves-equation-multicross}
A = \{(f_1, \ldots, f_n) \in k[[t]] \times \ldots \times k[[t]] \mid
f_1(0) = \ldots = f_n(0)\}
\end{equation} Dans la terminologie introduite dans Variétés, définition \ref{varieties-definition-wedge}, $A$ est un bouquet de $n$ copies de l'anneau des séries formelles en $1$ variable sur $k$. Remarquons que $k[[t]] \times \ldots \times k[[t]]$ est la clôture intégrale de $A$ dans son anneau total des fractions. L'invariant $\delta$ de $A$ est donc $n - 1$. On a un isomorphisme $$
k[[x_1, \ldots, x_n]]/(\{x_ix_j\}_{i \not = j}) \longrightarrow A
$$ obtenu en envoyant $x_i$ sur $(0, \ldots, 0, t, 0, \ldots, 0)$ dans $A$. Il s'ensuit que $\dim(A) = 1$ et $\dim_k \mathfrak m/\mathfrak m^2 = n$. En particulier, $A$ est régulier si et seulement si $n = 1$.

\begin{lemma}
\label{curves-lemma-multicross-algebra}
Soit $k$ un corps séparablement clos. Soit $A$ une algèbre locale de Nagata réduite de dimension $1$ sur $k$, de corps résiduel $k$. Alors $$
\text{invariant }\delta\text{ de }A \geq \text{nombre de branches de }A - 1
$$ En cas d'égalité, $A^\wedge$ est de la forme (\ref{curves-equation-multicross}).
\end{lemma}

\begin{proof}
Puisque le corps résiduel de $A$ est séparablement clos, le nombre de branches de $A$ est égal au nombre de branches géométriques de $A$; voir Compléments d'algèbre, définition \ref{more-algebra-definition-number-of-branches}. L'inégalité résulte de Variétés, lemme \ref{varieties-lemma-delta-number-branches-inequality}. Supposons qu'il y ait égalité. On peut remplacer $A$ par le complété de $A$; cela ne modifie ni le nombre de branches ni l'invariant $\delta$, d'après Compléments d'algèbre, lemme \ref{more-algebra-lemma-one-dimensional-number-of-branches} et Variétés, lemme \ref{varieties-lemma-delta-same-after-completion}. Alors $A$ est strictement hensélien; voir Algèbre, lemme \ref{algebra-lemma-complete-henselian}. D'après Variétés, lemme \ref{varieties-lemma-delta-number-branches-inequality-sh}, $A$ est un bouquet d'anneaux de valuation discrète complets. Chacun d'eux est isomorphe à $k[[t]]$ d'après Algèbre, lemme \ref{algebra-lemma-regular-complete-containing-coefficient-field}. Ainsi $A$ est de la forme (\ref{curves-equation-multicross}).
\end{proof}

\begin{definition}
\label{curves-definition-multicross}
Soit $k$ un corps algébriquement clos. Soit $X$ un schéma algébrique de dimension $1$ sur $k$. Soit $x \in X$ un point fermé. On dit que $x$ définit une {\it singularité en croisement multiple} si le complété $\mathcal{O}_{X, x}^\wedge$ est isomorphe à (\ref{curves-equation-multicross}) pour un certain $n \geq 2$. On dit que $x$ est un {\it nœud}, ou un {\it point double ordinaire}, ou {\it définit une singularité nodale}, si $n = 2$.
\end{definition}

\noindent
Ces singularités sont, en un sens, les singularités les plus simples que puisse présenter une courbe sur un corps algébriquement clos.

\begin{lemma}
\label{curves-lemma-multicross}
Soit $k$ un corps algébriquement clos. Soit $X$ un schéma algébrique réduit de dimension $1$ sur $k$. Soit $x \in X$. Les conditions suivantes sont équivalentes : \begin{enumerate} \item $x$ définit une singularité en croisement multiple; \item l'invariant $\delta$ de $X$ en $x$ est égal au nombre de branches de $X$ en $x$ diminué de $1$; \item il existe une suite de morphismes $U_n \to U_{n - 1} \to \ldots \to U_0 = U \subset X$, où $U$ est un voisinage ouvert de $x$, où $U_n$ est non singulier et où chaque $U_i \to U_{i - 1}$ est une identification de deux points comme dans l'exemple \ref{curves-example-glue-points}.
\end{enumerate}
\end{lemma}

\begin{proof}
L'équivalence de (1) et (2) est le lemme \ref{curves-lemma-multicross-algebra}.

\medskip\noindent
Supposons (3). Montrons par récurrence descendante sur $i$ que toutes les singularités de $U_i$ sont des croisements multiples. C'est vrai pour $U_n$, puisque $U_n$ n'a pas de point singulier. Si $U_i$ s'obtient à partir de $U_{i + 1}$ en identifiant $a, b \in U_{i + 1}$ en un point $c \in U_i$, alors $$
\mathcal{O}_{U_i, c}^\wedge \subset
\mathcal{O}_{U_{i + 1}, a}^\wedge \times \mathcal{O}_{U_{i + 1}, b}^\wedge
$$ est l'ensemble des éléments ayant les mêmes classes résiduelles dans $k$. Le nombre de branches en $c$ est donc la somme des nombres de branches en $a$ et $b$, et l'invariant $\delta$ en $c$ est la somme des invariants $\delta$ en $a$ et $b$ augmentée de $1$ (car l'inclusion affichée est de codimension $1$). Cela établit (2), comme voulu.

\medskip\noindent
Supposons les conditions équivalentes (1) et (2). On peut choisir un ouvert $U \subset X$ tel que $x$ soit le seul point singulier de $U$. Appliquons alors le lemme \ref{curves-lemma-factor-almost-isomorphism} au morphisme de normalisation $$
U^\nu = U_n \to U_{n - 1} \to \ldots \to U_1 \to U_0 = U
$$ Il reste seulement à montrer qu'aucune étape n'écrase un vecteur tangent. Supposons que $U_{i + 1} \to U_i$ soit le premier morphisme, pour le plus petit $i$, qui écrase un vecteur tangent $\theta$ en $u' \in U_{i + 1}$ au-dessus de $u \in U_i$. En raisonnant comme ci-dessus, on voit que $u$ est une singularité en croisement multiple (car les morphismes $U_i \to \ldots \to U_0$ identifient des paires de points). Mais le nombre de branches en $u'$ et en $u$ est le même, tandis que l'invariant $\delta$ de $U_i$ en $u$ dépasse de $1$ l'invariant $\delta$ de $U_{i + 1}$ en $u'$. D'après le lemme \ref{curves-lemma-multicross-algebra}, $u$ ne peut donc pas être une singularité en croisement multiple, d'où une contradiction.
\end{proof}

\begin{lemma}
\label{curves-lemma-multicross-gorenstein-is-nodal}
Soit $k$ un corps algébriquement clos. Soit $X$ un schéma algébrique réduit de dimension $1$ sur $k$. Soit $x \in X$ un point où ce schéma présente une singularité en croisement multiple (définition \ref{curves-definition-multicross}). Si $X$ est de Gorenstein, alors $x$ est un nœud.
\end{lemma}

\begin{proof}
L'application $\mathcal{O}_{X, x} \to \mathcal{O}_{X, x}^\wedge$ est plate et non ramifiée au sens où $\kappa(x) = \mathcal{O}_{X, x}^\wedge/\mathfrak m_x \mathcal{O}_{X, x}^\wedge$. (Voir Compléments d'algèbre, section \ref{more-algebra-section-permanence-completion}.) Ainsi, si $X$ est de Gorenstein, $\mathcal{O}_{X, x}$ est de Gorenstein, puis $\mathcal{O}_{X, x}^\wedge$ est de Gorenstein d'après Complexes dualisants, lemme \ref{dualizing-lemma-flat-under-gorenstein}. Il suffit donc de montrer que l'anneau $A$ de (\ref{curves-equation-multicross}), avec $n \geq 2$, est de Gorenstein si et seulement si $n = 2$.

\medskip\noindent
Si $n = 2$, alors $A = k[[x, y]]/(xy)$ est une intersection complète, donc de Gorenstein. Cela résulte par exemple de Dualité pour les schémas, lemme \ref{duality-lemma-gorenstein-lci} appliqué à $k[[x, y]] \to A$, et du fait que l'anneau local régulier $k[[x, y]]$ est de Gorenstein d'après Complexes dualisants, lemme \ref{dualizing-lemma-regular-gorenstein}.

\medskip\noindent
Supposons $n > 2$. Si $A$ était de Gorenstein, $A$ serait un complexe dualisant sur $A$ (Dualité pour les schémas, définition \ref{duality-definition-gorenstein}). Alors $R\Hom(k, A)$ serait égal à $k[n]$ pour un certain $n \in \mathbf{Z}$; voir Complexes dualisants, lemme \ref{dualizing-lemma-find-function}. Il en résulterait que $\Ext^1_A(k, A) \cong k$ ou $\Ext^1_A(k, A) = 0$ (suivant la valeur de $n$; en fait, $n$ doit valoir $-1$, mais cela n'importe pas ici). La suite exacte $$
0 \to \mathfrak m_A \to A \to k \to 0
$$ donne $$
\Ext^1_A(k, A) = \Hom_A(\mathfrak m_A, A)/A
$$, où $A \to \Hom_A(\mathfrak m_A, A)$ est défini par $a \mapsto (a' \mapsto aa')$. Soit $e_i \in \Hom_A(\mathfrak m_A, A)$ l'élément qui envoie $(f_1, \ldots, f_n) \in \mathfrak m_A$ sur $(0, \ldots, 0, f_i, 0, \ldots, 0)$. Le lecteur vérifie aisément que $e_1, \ldots, e_{n - 1}$ sont linéairement indépendants sur $k$ dans $\Hom_A(\mathfrak m_A, A)/A$. Ainsi $\dim_k \Ext^1_A(k, A) \geq n - 1 \geq 2$, ce qui achève la démonstration. (Remarquons que $e_1 + \ldots + e_n$ est l'image de $1$ par l'application $A \to \Hom_A(\mathfrak m_A, A)$.)
\end{proof}




\section{Torsion du groupe de Picard}
\label{curves-section-torsion-in-pic}

\noindent
Dans cette section, nous majorons la torsion du groupe de Picard d'un schéma propre de dimension $1$ sur un corps. Nous utiliserons ce résultat dans l'étude de la réduction semi-stable des courbes.

\medskip\noindent
Il ne semble pas exister de démonstration élémentaire du résultat du lemme \ref{curves-lemma-torsion-picard-smooth-projective}. L'analyse de sa preuve fait apparaître deux ingrédients essentiels : (1) l'existence d'une variété abélienne classifiant les faisceaux inversibles de degré zéro sur une courbe projective lisse et (2) la détermination de la structure des points de torsion d'une variété abélienne.

\begin{lemma}
\label{curves-lemma-torsion-picard-smooth-projective}
Soit $k$ un corps algébriquement clos. Soit $X$ une courbe projective lisse de genre $g$ sur $k$. \begin{enumerate} \item Si $n \geq 1$ est inversible dans $k$, alors $\Pic(X)[n] \cong (\mathbf{Z}/n\mathbf{Z})^{\oplus 2g}$. \item Si la caractéristique de $k$ est $p > 0$, il existe un entier $0 \leq f \leq g$ tel que $\Pic(X)[p^m] \cong (\mathbf{Z}/p^m\mathbf{Z})^{\oplus f}$ pour tout $m \geq 1$.
\end{enumerate}
\end{lemma}

\begin{proof}
Notons $\Pic^0(X) \subset \Pic(X)$ le sous-groupe des faisceaux inversibles de degré $0$. Autrement dit, on a une suite exacte courte $$
0 \to \Pic^0(X) \to \Pic(X) \xrightarrow{\deg} \mathbf{Z} \to 0.
$$ Le groupe $\Pic^0(X)$ est le groupe des points à valeurs dans $k$ du schéma en groupes $\underline{\Picardfunctor}^0_{X/k}$; voir Schémas de Picard des courbes, lemme \ref{pic-lemma-picard-pieces}. Le même lemme affirme que $\underline{\Picardfunctor}^0_{X/k}$ est une variété abélienne de dimension $g$ sur $k$ au sens de Groupoïdes, définition \ref{groupoids-definition-abelian-variety}. On conclut donc par Groupoïdes, proposition \ref{groupoids-proposition-review-abelian-varieties}.
\end{proof}

\begin{lemma}
\label{curves-lemma-torsion-picard-becomes-visible}
Soit $k$ un corps. Soit $n$ premier à la caractéristique de $k$. Soit $X$ une courbe propre lisse sur $k$ telle que $H^0(X, \mathcal{O}_X) = k$ et de genre $g$. \begin{enumerate} \item Si $g = 1$, il existe une extension finie séparable $k'/k$ telle que $X_{k'}$ possède un point $k'$-rationnel et que $\Pic(X_{k'})[n] \cong (\mathbf{Z}/n\mathbf{Z})^{\oplus 2}$. \item Si $g \geq 2$, il existe une extension finie séparable $k'/k$ avec $[k' : k] \leq (2g - 2)(n^{2g})!$ telle que $X_{k'}$ possède un point $k'$-rationnel et que $\Pic(X_{k'})[n] \cong (\mathbf{Z}/n\mathbf{Z})^{\oplus 2g}$.
\end{enumerate}
\end{lemma}

\begin{proof}
Supposons $g \geq 2$. On peut d'abord choisir une extension finie séparable de degré au plus $2g - 2$ sur laquelle $X$ acquiert un point rationnel; voir le lemme \ref{curves-lemma-point-over-separable-extension}. On peut donc supposer que $X$ possède un point $k$-rationnel, noté $x \in X(k)$, mais il faut désormais démontrer le lemme avec $[k' : k] \leq (n^{2g})!$. Soit $k \subset k^{sep} \subset \overline{k}$ une clôture séparable contenue dans une clôture algébrique. D'après le lemme \ref{curves-lemma-torsion-picard-smooth-projective}, on a $$
\Pic(X_{\overline{k}})[n] \cong (\mathbf{Z}/n\mathbf{Z})^{\oplus 2g}
$$ D'après Schémas de Picard des courbes, lemme \ref{pic-lemma-torsion-descends}, on en déduit $$
\Pic(X_{k^{sep}})[n] \cong (\mathbf{Z}/n\mathbf{Z})^{\oplus 2g}
$$ D'après Schémas de Picard des courbes, lemme \ref{pic-lemma-torsion-descends}, il existe une action continue $$
\text{Gal}(k^{sep}/k)
\longrightarrow
\text{Aut}(\Pic(X_{k^{sep}})[n])
$$, et le lemme est vrai pour le corps fixe $k'$ du noyau de cette application. Ce noyau est ouvert par continuité de l'action, ce qui entraîne que $k'/k$ est finie. Par la théorie de Galois, $\text{Gal}(k'/k)$ est l'image de la flèche affichée. Puisque le groupe des permutations d'un ensemble de cardinal $n^{2g}$ a pour cardinal $(n^{2g})!$, la théorie de Galois donne $[k' : k] \leq (n^{2g})!$. (Bien entendu, cela démontre le lemme avec la borne $|\text{GL}_{2g}(\mathbf{Z}/n\mathbf{Z})|$, mais seule l'existence d'une borne nous importe ici.)

\medskip\noindent
Si le genre est $1$, il n'y a pas de majoration du degré d'une extension finie séparable sur laquelle $X$ acquiert un point rationnel (nous omettons les détails). Une telle extension existe néanmoins, par exemple d'après Variétés, lemme \ref{varieties-lemma-smooth-separable-closed-points-dense}. La suite de la démonstration est identique au cas $g \geq 2$.
\end{proof}

\begin{proposition}
\label{curves-proposition-torsion-picard-reduced-proper}
Soit $k$ un corps algébriquement clos. Soit $X$ un schéma propre sur $k$, réduit, connexe et de dimension $1$. Soit $g$ le genre de $X$ et soit $g_{geom}$ la somme des genres géométriques des composantes irréductibles de $X$. Pour tout nombre premier $\ell$ différent de la caractéristique de $k$, on a $$
\dim_{\mathbf{F}_\ell} \Pic(X)[\ell]
\leq g + g_{geom}
$$, avec égalité si et seulement si toutes les singularités de $X$ sont des croisements multiples.
\end{proposition}

\begin{proof}
Soit $\nu : X^\nu \to X$ la normalisation (Variétés, lemme \ref{varieties-lemma-prepare-delta-invariant}). Choisissons une factorisation $$
X^\nu = X_n \to X_{n - 1} \to \ldots \to X_1 \to X_0 = X
$$ comme dans le lemme \ref{curves-lemma-factor-almost-isomorphism}. Notons $h^0_i = \dim_k H^0(X_i, \mathcal{O}_{X_i})$ et $h^1_i = \dim_k H^1(X_i, \mathcal{O}_{X_i})$. D'après les lemmes \ref{curves-lemma-glue-points} et \ref{curves-lemma-squish-tangent-vector}, pour chaque $n > i \geq 0$, on est dans l'un des trois cas suivants : \begin{enumerate} \item $X_i$ s'obtient en identifiant les points $a, b \in X_{i + 1}$, situés sur des composantes connexes distinctes : dans ce cas, $\Pic(X_i) = \Pic(X_{i + 1})$, $h^0_{i + 1} = h^0_i + 1$, $h^1_{i + 1} = h^1_i$; \item $X_i$ s'obtient en identifiant les points $a, b \in X_{i + 1}$, situés sur une même composante connexe : dans ce cas, on a une suite exacte courte $$
0 \to k^* \to \Pic(X_i) \to \Pic(X_{i + 1}) \to 0,
$$ et $h^0_{i + 1} = h^0_i$, $h^1_{i + 1} = h^1_i - 1$; \item $X_i$ s'obtient en écrasant un vecteur tangent dans $X_{i + 1}$ : dans ce cas, on a une suite exacte courte $$
0 \to (k, +) \to \Pic(X_i) \to \Pic(X_{i + 1}) \to 0,
$$ et $h^0_{i + 1} = h^0_i$, $h^1_{i + 1} = h^1_i - 1$. \end{enumerate} Pour les assertions sur les dimensions des groupes de cohomologie du faisceau structural, on utilise les suites exactes des exemples \ref{curves-example-glue-points} et \ref{curves-example-squish-tangent-vector}. Puisque $k$ est algébriquement clos de caractéristique différente de $\ell$, les groupes $(k, +)$ et $k^*$ sont divisibles par $\ell$ et leur torsion d'ordre $\ell$ est respectivement $(k, +)[\ell] = 0$ et $k^*[\ell] \cong \mathbf{F}_\ell$. Ainsi $$
\dim_{\mathbf{F}_\ell} \Pic(X_{i + 1})[\ell] -
\dim_{\mathbf{F}_\ell}\Pic(X_i)[\ell]
$$ est nul, sauf dans le cas (2), où il vaut $-1$. À l'issue de ce processus, on obtient la normalisation $X^\nu = X_n$, qui est une réunion disjointe de courbes projectives lisses sur $k$. On a donc \begin{enumerate} \item $h^1_n = g_{geom}$ et \item $\dim_{\mathbf{F}_\ell} \Pic(X_n)[\ell] = 2g_{geom}$. \end{enumerate} La dernière égalité résulte du lemme \ref{curves-lemma-torsion-picard-smooth-projective}. Puisque $g = h^1_0$, le nombre d'étapes de types (2) et (3) est au plus $h^1_0 - h^1_n = g - g_{geom}$. Le calcul des différences de rangs donne $$
\dim_{\mathbf{F}_\ell} \Pic(X)[\ell] \leq
g - g_{geom} + 2g_{geom} = g + g_{geom}
$$, avec égalité si et seulement si aucune étape de type (3) n'intervient. Cela signifie bien que toutes les singularités de $X$ sont des croisements multiples, d'après le lemme \ref{curves-lemma-multicross}. Réciproquement, si toutes les singularités sont des croisements multiples, le lemme \ref{curves-lemma-multicross} assure l'existence d'une suite $X^\nu = X_n \to \ldots \to X_0 = X$ comme ci-dessus sans étape de type (3), et l'on obtient l'égalité dans le lemme (il suffit de recoller les suites locales en chaque point pour en obtenir une qui convienne à tous les points singuliers de $X$; nous omettons quelques détails).
\end{proof}





\section{Comparaison du genre et du genre géométrique}
\label{curves-section-genus-geometric-genus}

\noindent
Soit $k$ un corps, et notons $\overline{k}$ une clôture algébrique de ce corps. Soit $X$ un schéma propre de dimension $\leq 1$ sur $k$. Nous définissons $g_{geom}(X/k)$ comme la somme des genres géométriques des composantes irréductibles de $X_{\overline{k}}$ de dimension $1$.

\begin{lemma}
\label{curves-lemma-bound-geometric-genus}
Soit $k$ un corps. Soit $X$ un schéma propre de dimension $\leq 1$ sur $k$. Alors $$
g_{geom}(X/k) = \sum\nolimits_{C \subset X} g_{geom}(C/k)
$$, la somme portant sur les composantes irréductibles $C \subset X$ de dimension $1$.
\end{lemma}

\begin{proof}
Cela résulte immédiatement de la définition et du fait qu'une composante irréductible $\overline{Z}$ de $X_{\overline{k}}$ se projette surjectivement sur une composante irréductible $Z$ de $X$ (Variétés, lemme \ref{varieties-lemma-image-irreducible}) de même dimension (Morphismes, lemme \ref{morphisms-lemma-dimension-fibre-after-base-change} appliqué au point générique de $\overline{Z}$).
\end{proof}

\begin{lemma}
\label{curves-lemma-geometric-genus-normalization}
Soit $k$ un corps. Soit $X$ un schéma propre de dimension $\leq 1$ sur $k$. Alors \begin{enumerate} \item On a $g_{geom}(X/k) = g_{geom}(X_{red}/k)$. \item Si $X' \to X$ est un morphisme propre birationnel, alors $g_{geom}(X'/k) = g_{geom}(X/k)$. \item Si $X^\nu \to X$ est le morphisme de normalisation, alors $g_{geom}(X^\nu/k) = g_{geom}(X/k)$.
\end{enumerate}
\end{lemma}

\begin{proof}
L'assertion (1) résulte immédiatement du lemme \ref{curves-lemma-bound-geometric-genus}. Si $X' \to X$ est propre birationnel, il est fini et c'est un isomorphisme au-dessus d'un ouvert dense (voir Variétés, lemmes \ref{varieties-lemma-finite-in-codim-1} et \ref{varieties-lemma-modification-normal-iso-over-codimension-1}). Ainsi $X'_{\overline{k}} \to X_{\overline{k}}$ est un isomorphisme au-dessus d'un ouvert dense. Les composantes irréductibles de $X'_{\overline{k}}$ et $X_{\overline{k}}$ sont donc en bijection, et les composantes correspondantes ont des corps de fonctions isomorphes. Leurs modèles projectifs non singuliers sont donc isomorphes, et leurs genres géométriques égaux. Cela prouve (2). L'assertion (3) résulte de (1), de (2) et du fait que $X^\nu \to X_{red}$ est birationnel (Morphismes, lemme \ref{morphisms-lemma-normalization-birational}).
\end{proof}

\begin{lemma}
\label{curves-lemma-genus-goes-down}
Soit $k$ un corps. Soit $X$ un schéma propre de dimension $\leq 1$ sur $k$. Soit $f : Y \to X$ un morphisme fini tel qu'il existe un ouvert dense $U \subset X$ au-dessus duquel $f$ soit une immersion fermée. Alors
$$
\dim_k H^1(X, \mathcal{O}_X) \geq \dim_k H^1(Y, \mathcal{O}_Y)
$$
\end{lemma}

\begin{proof}
Considérons la suite exacte $$
0 \to \mathcal{G} \to \mathcal{O}_X \to f_*\mathcal{O}_Y \to \mathcal{F} \to 0
$$ de faisceaux cohérents sur $X$. Par hypothèse, $\mathcal{F}$ est à support dans un nombre fini de points fermés, et sa cohomologie supérieure est donc nulle (Variétés, lemme \ref{varieties-lemma-chi-tensor-finite}). D'autre part, on a $H^2(X, \mathcal{G}) = 0$ d'après Cohomologie, proposition \ref{cohomology-proposition-vanishing-Noetherian}. Il s'ensuit formellement que l'application induite $H^1(X, \mathcal{O}_X) \to H^1(X, f_*\mathcal{O}_Y)$ est surjective. Puisque $H^1(X, f_*\mathcal{O}_Y) = H^1(Y, \mathcal{O}_Y)$ (Cohomologie des schémas, lemme \ref{coherent-lemma-relative-affine-cohomology}), le lemme en résulte.
\end{proof}

\begin{lemma}
\label{curves-lemma-genus-normalization}
Soit $k$ un corps. Soit $X$ un schéma propre de dimension $\leq 1$ sur $k$. Si $X' \to X$ est un morphisme propre birationnel, alors $$
\dim_k H^1(X, \mathcal{O}_X) \geq \dim_k H^1(X', \mathcal{O}_{X'})
$$ Si $X$ est réduit, si $H^0(X, \mathcal{O}_X) \to H^0(X', \mathcal{O}_{X'})$ est surjective et s'il y a égalité, alors $X' = X$.
\end{lemma}

\begin{proof}
Si $f : X' \to X$ est propre birationnel, il est fini et c'est un isomorphisme au-dessus d'un ouvert dense (voir Variétés, lemmes \ref{varieties-lemma-finite-in-codim-1} et \ref{varieties-lemma-modification-normal-iso-over-codimension-1}). L'inégalité résulte donc du lemme \ref{curves-lemma-genus-goes-down}. Supposons $X$ réduit. Alors $\mathcal{O}_X \to f_*\mathcal{O}_{X'}$ est injective et l'on obtient une suite exacte courte $$
0 \to \mathcal{O}_X \to f_*\mathcal{O}_{X'} \to \mathcal{F} \to 0
$$ Sous les hypothèses de la seconde assertion, la suite exacte longue de cohomologie donne $H^0(X, \mathcal{F}) = 0$. Alors $\mathcal{F} = 0$, car $\mathcal{F}$ est engendré par ses sections globales (Variétés, lemme \ref{varieties-lemma-chi-tensor-finite}), et $\mathcal{O}_X = f_*\mathcal{O}_{X'}$. Puisque $f$ est affine, cela entraîne $X = X'$.
\end{proof}

\begin{lemma}
\label{curves-lemma-bound-geometric-genus-curve}
Soit $k$ un corps. Soit $C$ une courbe propre sur $k$. Posons $\kappa = H^0(C, \mathcal{O}_C)$. Alors
$$
[\kappa : k]_s \dim_\kappa H^1(C, \mathcal{O}_C) \geq g_{geom}(C/k)
$$
\end{lemma}

\begin{proof}
D'après Variétés, lemme \ref{varieties-lemma-regular-functions-proper-variety}, $\kappa$ est un corps et une extension finie de $k$. Par Corps, lemme \ref{fields-lemma-separable-degree}, on a $[\kappa : k]_s = |\Mor_k(\kappa, \overline{k})|$; ainsi $\Spec(\kappa \otimes_k \overline{k})$ possède $[\kappa : k]_s$ points ayant chacun pour corps résiduel $\overline{k}$. On a donc $$
C_{\overline{k}} =
\bigcup\nolimits_{t \in \Spec(\kappa \otimes_k \overline{k})} C_t
$$ (réunion ensembliste). Ici $C_t = C \times_{\Spec(\kappa), t} \Spec(\overline{k})$, où l'on considère $t$ comme un homomorphisme de $k$-algèbres $t : \kappa \to \overline{k}$. Il en résulte que $g_{geom}(C/k) = \sum_t g_{geom}(C_t/\overline{k})$, la somme étant indexée par un ensemble de cardinal $[\kappa : k]_s$. On a $$
H^0(C_t, \mathcal{O}_{C_t}) = \overline{k}
\quad\text{et}\quad
\dim_{\overline{k}} H^1(C_t, \mathcal{O}_{C_t}) =
\dim_\kappa H^1(C, \mathcal{O}_C)
$$ par cohomologie et changement de base (Cohomologie des schémas, lemme \ref{coherent-lemma-flat-base-change-cohomology}). Remarquons que la normalisation $C_t^\nu$ est la réunion disjointe des modèles projectifs non singuliers des composantes irréductibles de $C_t$ (Morphismes, lemme \ref{morphisms-lemma-normalization-in-terms-of-components}). Ainsi $\dim_{\overline{k}} H^1(C_t^\nu, \mathcal{O}_{C_t^\nu})$ est égal à $g_{geom}(C_t/\overline{k})$. D'après le lemme \ref{curves-lemma-genus-goes-down}, on a $$
\dim_{\overline{k}} H^1(C_t, \mathcal{O}_{C_t}) \geq
\dim_{\overline{k}} H^1(C_t^\nu, \mathcal{O}_{C_t^\nu})
$$, ce qui achève la démonstration.
\end{proof}

\begin{lemma}
\label{curves-lemma-bound-torsion-simple}
Soit $k$ un corps. Soit $X$ un schéma propre de dimension $\leq 1$ sur $k$. Soit $\ell$ un nombre premier inversible dans $k$. Alors $$
\dim_{\mathbf{F}_\ell} \Pic(X)[\ell] \leq
\dim_k H^1(X, \mathcal{O}_X) + g_{geom}(X/k)
$$, où $g_{geom}(X/k)$ est défini ci-dessus.
\end{lemma}

\begin{proof}
L'application $\Pic(X) \to \Pic(X_{\overline{k}})$ est injective d'après Variétés, lemme \ref{varieties-lemma-change-fields-pic}. D'après Cohomologie des schémas, lemme \ref{coherent-lemma-flat-base-change-cohomology}, $\dim_k H^1(X, \mathcal{O}_X)$ est égal à $\dim_{\overline{k}} H^1(X_{\overline{k}}, \mathcal{O}_{X_{\overline{k}}})$. On peut donc supposer $k$ algébriquement clos.

\medskip\noindent
Soit $X_{red}$ le réduit de $X$. La surjection $\mathcal{O}_X \to \mathcal{O}_{X_{red}}$ induit une surjection $H^1(X, \mathcal{O}_X) \to H^1(X, \mathcal{O}_{X_{red}})$, car la cohomologie des faisceaux quasi-cohérents s'annule en degrés $\geq 2$ d'après Cohomologie, proposition \ref{cohomology-proposition-vanishing-Noetherian}. Puisque $X_{red} \to X$ induit une bijection sur les composantes irréductibles au-dessus de $\overline{k}$ et un isomorphisme sur la torsion d'ordre $\ell$ des groupes de Picard (Schémas de Picard des courbes, lemme \ref{pic-lemma-torsion-descends}), on peut remplacer $X$ par $X_{red}$. On est ainsi ramené à la proposition \ref{curves-proposition-torsion-picard-reduced-proper}.
\end{proof}






\section{Courbes nodales}
\label{curves-section-nodal}

\noindent
Nous avons déjà défini les points doubles ordinaires sur un corps algébriquement clos; voir la définition \ref{curves-definition-multicross}. Si $x \in X$ est un point fermé d'un schéma algébrique de dimension $1$ sur un corps algébriquement clos $k$, alors $x$ est un point double ordinaire si et seulement si $$
\mathcal{O}_{X, x}^\wedge \cong k[[x, y]]/(xy)
$$ Voir la discussion qui suit (\ref{curves-equation-multicross}) dans la section \ref{curves-section-multicross}.

\begin{definition}
\label{curves-definition-nodal}
Soit $k$ un corps. Soit $X$ un schéma localement algébrique de dimension $1$ sur $k$. \begin{enumerate} \item Un point fermé $x \in X$ est appelé un {\it nœud}, ou un {\it point double ordinaire}, ou est dit {\it définir une singularité nodale}, s'il existe un point double ordinaire $\overline{x} \in X_{\overline{k}}$ d'image $x$. \item On dit que {\it les singularités de $X$ sont au plus nodales} si tout point fermé de $X$ appartient au lieu lisse du morphisme structural $X \to \Spec(k)$ ou est un point double ordinaire.
\end{enumerate}
\end{definition}

\noindent
On appelle souvent {\it courbe nodale} un schéma algébrique de dimension $1$, noté $X$, lorsque les singularités de $X$ sont au plus nodales. On exige parfois aussi qu'une courbe nodale soit propre. Puisqu'une courbe nodale ainsi définie n'est pas nécessairement irréductible, cette terminologie entre en conflit avec notre définition antérieure d'une courbe comme variété de dimension $1$.

\begin{lemma}
\label{curves-lemma-reduced-quotient-regular-ring-dim-2}
Soit $(A, \mathfrak m)$ un anneau local régulier de dimension $2$. Soit $I \subset \mathfrak m$ un idéal. \begin{enumerate} \item Si $A/I$ est réduit, alors $I = (0)$, $I = \mathfrak m$ ou $I = (f)$ pour un élément non nul $f \in \mathfrak m$. \item Si $A/I$ est de profondeur $1$, alors $I = (f)$ pour un élément non nul $f \in \mathfrak m$. 
\end{enumerate}
\end{lemma}

\begin{proof}
Supposons $I \not = 0$. Écrivons $I = (f_1, \ldots, f_r)$. Puisque $A$ est factoriel (Compléments d'algèbre, lemme \ref{more-algebra-lemma-regular-local-UFD}), on peut écrire $f_i = fg_i$, où $f$ est le pgcd des $f_1, \ldots, f_r$. Le pgcd des $g_1, \ldots, g_r$ est donc $1$, ce qui signifie qu'il n'y a pas d'idéal premier de hauteur $1$ contenant $g_1, \ldots, g_r$. Alors soit $(g_1, \ldots, g_r) = A$, ce qui entraîne $I = (f)$, soit, dans le cas contraire, $\dim(A) = 2$ entraîne $V(g_1, \ldots, g_r) = \{\mathfrak m\}$, c'est-à-dire $\mathfrak m = \sqrt{(g_1, \ldots, g_r)}$.

\medskip\noindent
Supposons $A/I$ réduit, c'est-à-dire $I$ radical. Si $f$ est une unité, puisque $I$ est radical, on a $I = \mathfrak m$. Si $f \in \mathfrak m$, on voit que $f^n$ s'envoie sur zéro dans $A/I$. Ainsi $f \in I$ par réduction, et l'on en conclut que $I = (f)$.

\medskip\noindent
Supposons $A/I$ de profondeur $1$. Alors $\mathfrak m$ n'est pas un idéal premier associé à $A/I$. Puisque la classe de $f$ modulo $I$ est annulée par $g_1, \ldots, g_r$, la classe de $f$ est nulle dans $A/I$. Ainsi $I = (f)$, comme voulu.
\end{proof}

\noindent
Soit $\kappa$ un corps et soit $V$ un espace vectoriel sur $\kappa$. On dira qu'un élément $q \in \text{Sym}^2_\kappa(V)$ définit une {\it forme quadratique non dégénérée} si l'application $\kappa$-linéaire induite $V^\vee \to V$ est un isomorphisme. Si $q = \sum_{i \leq j} a_{ij} x_i x_j$ pour une base sur $\kappa$, notée $x_1, \ldots, x_n$, de $V$, cela signifie que le déterminant de la matrice $$
\left(
\begin{matrix}
2a_{11} & a_{12} & \ldots \\
a_{12} & 2a_{22} & \ldots \\
\ldots & \ldots & \ldots
\end{matrix}
\right)
$$ est non nul. Cela équivaut à la condition que les dérivées partielles de $q$ par rapport aux $x_i$ définissent $0$ au sens schématique.

\begin{lemma}
\label{curves-lemma-nodal-algebraic}
Soit $k$ un corps. Soit $(A, \mathfrak m, \kappa)$ une $k$-algèbre locale noethérienne. Les conditions suivantes sont équivalentes : \begin{enumerate} \item $\kappa/k$ est séparable, $A$ est réduit, $\dim_\kappa(\mathfrak m/\mathfrak m^2) = 2$, et il existe $q \in \text{Sym}^2_\kappa(\mathfrak m/\mathfrak m^2)$ définissant une forme quadratique non dégénérée et s'envoyant sur zéro dans $\mathfrak m^2/\mathfrak m^3$; \item $\kappa/k$ est séparable, $\text{depth}(A) = 1$, $\dim_\kappa(\mathfrak m/\mathfrak m^2) = 2$, et il existe $q \in \text{Sym}^2_\kappa(\mathfrak m/\mathfrak m^2)$ définissant une forme quadratique non dégénérée et s'envoyant sur zéro dans $\mathfrak m^2/\mathfrak m^3$; \item $\kappa/k$ est séparable, $A^\wedge \cong \kappa[[x, y]]/(ax^2 + bxy + cy^2)$ en tant que $k$-algèbre, où $ax^2 + bxy + cy^2$ est une forme quadratique non dégénérée sur $\kappa$.
\end{enumerate}
\end{lemma}

\begin{proof}
Supposons (3). Alors $A^\wedge$ est réduit, car $ax^2 + bxy + cy^2$ est soit irréductible, soit produit de deux éléments premiers non associés. Ainsi $A \subset A^\wedge$ est réduit. Il s'ensuit que (1) est vraie.

\medskip\noindent
Supposons (1). Alors $A$ ne peut pas être artinien, car il ne serait pas réduit puisque $\mathfrak m \not = (0)$. Ainsi $\dim(A) \geq 1$, puis $\text{depth}(A) \geq 1$ d'après Algèbre, lemme \ref{algebra-lemma-criterion-reduced}. D'autre part, $\dim(A) = 2$ entraîne que $A$ est régulier, ce qui contredit l'existence de $q$ d'après Algèbre, lemme \ref{algebra-lemma-regular-graded}. Ainsi $\dim(A) \leq 1$, puis $\text{depth}(A) = 1$ d'après Algèbre, lemme \ref{algebra-lemma-bound-depth}. Il s'ensuit que (2) est vraie.

\medskip\noindent
Supposons (2). Puisque la profondeur de $A$ est égale à celle de $A^\wedge$ (Compléments d'algèbre, lemme \ref{more-algebra-lemma-completion-depth}) et que les autres conditions sont invariantes par complétion, on peut supposer $A$ complet. Choisissons $\kappa \to A$ comme dans Compléments d'algèbre, lemme \ref{more-algebra-lemma-lift-residue-field}. Puisque $\dim_\kappa(\mathfrak m/\mathfrak m^2) = 2$, on peut choisir $x_0, y_0 \in \mathfrak m$ dont les images forment une base. On obtient un homomorphisme continu de $\kappa$-algèbres $$
\kappa[[x, y]] \longrightarrow A
$$ par les règles $x \mapsto x_0$ et $y \mapsto y_0$. Soit $q$ la classe de $ax_0^2 + bx_0y_0 + cy_0^2$ dans $\text{Sym}^2_\kappa(\mathfrak m/\mathfrak m^2)$. Écrivons $Q(x, y) = ax^2 + bxy + cy^2$, considéré comme un polynôme en deux variables. On voit alors que $$
Q(x_0, y_0) = ax_0^2 + bx_0y_0 + cy_0^2 =
\sum\nolimits_{i + j = 3} a_{ij} x_0^iy_0^j
$$ pour certains $a_{ij}$ dans $A$. Nous voulons augmenter l'ordre d'annulation en modifiant le choix de $x_0$, $y_0$. Supposons que $x_1, y_1 \in \mathfrak m^2$. Alors $$
Q(x_0 + x_1, y_0 + y_1) = Q(x_0, y_0) +
(2ax_0 + by_0)x_1 + (bx_0 + 2cy_0)y_1 \bmod \mathfrak m^4
$$ La non-dégénérescence de $Q$ signifie précisément que $2ax_0 + by_0$ et $bx_0 + 2cy_0$ forment une base sur $\kappa$ de $\mathfrak m/\mathfrak m^2$; voir la discussion précédant le lemme. On peut donc choisir $x_1, y_1 \in \mathfrak m^2$ tels que $Q(x_0 + x_1, y_0 + y_1) \in \mathfrak m^4$. En poursuivant par récurrence, on trouve $x_i, y_i \in \mathfrak m^{i + 1}$ tels que $$
Q(x_0 + x_1 + \ldots + x_n, y_0 + y_1 + \ldots + y_n) \in \mathfrak m^{n + 3}
$$ Puisque $A$ est complet, on peut poser $x_\infty = \sum x_i$ et $y_\infty = \sum y_i$ et considérer l'application $\kappa[[x, y]] \longrightarrow A$ envoyant $x$ sur $x_\infty$ et $y$ sur $y_\infty$. Elle induit une surjection $\kappa[[x, y]]/(Q) \longrightarrow A$ d'après Algèbre, lemme \ref{algebra-lemma-completion-generalities}. D'après le lemme \ref{curves-lemma-reduced-quotient-regular-ring-dim-2}, le noyau de $\kappa[[x, y]] \to A$ est principal. Mais ce noyau ne peut contenir de diviseur propre de $Q$, car un tel diviseur serait de degré $1$ dans $x, y$, ce qui contredirait $\dim(\mathfrak m/\mathfrak m^2) = 2$. Ainsi $Q$ engendre le noyau, comme voulu.
\end{proof}

\begin{lemma}
\label{curves-lemma-2-branches-delta-1}
Soit $k$ un corps. Soit $(A, \mathfrak m, \kappa)$ une $k$-algèbre locale de Nagata. Les conditions suivantes sont équivalentes : \begin{enumerate} \item $k \to A$ est comme dans le lemme \ref{curves-lemma-nodal-algebraic}; \item $\kappa/k$ est séparable, $A$ est réduit de dimension $1$, l'invariant $\delta$ de $A$ vaut $1$, et $A$ possède $2$ branches géométriques. \end{enumerate} Si ces conditions sont satisfaites, la clôture intégrale $A'$ de $A$ dans son anneau total des fractions possède soit $1$, soit $2$ idéaux maximaux $\mathfrak m'$, et les extensions $\kappa(\mathfrak m')/k$ sont séparables.
\end{lemma}

\begin{proof}
Dans les deux cas, $A$ et $A^\wedge$ sont réduits. Dans le cas (2), cela tient au fait que le complété d'un anneau local de Nagata réduit est réduit (Compléments d'algèbre, lemme \ref{more-algebra-lemma-completion-reduced}). Dans les deux cas, $A$ et $A^\wedge$ sont de dimension $1$ (Compléments d'algèbre, lemme \ref{more-algebra-lemma-completion-dimension}). L'invariant $\delta$ et le nombre de branches géométriques de $A$ et $A^\wedge$ coïncident d'après Variétés, lemme \ref{varieties-lemma-delta-same-after-completion} et Compléments d'algèbre, lemme \ref{more-algebra-lemma-one-dimensional-number-of-branches}. Soit $A'$ la clôture intégrale de $A$ dans son anneau total des fractions, comme dans Variétés, lemme \ref{varieties-lemma-pre-delta-invariant}. D'après Variétés, lemme \ref{varieties-lemma-normalization-same-after-completion}, $A' \otimes_A A^\wedge$ joue le même rôle pour $A^\wedge$. On peut donc remplacer $A$ par $A^\wedge$ et supposer $A$ complet.

\medskip\noindent
Supposons (1). Il suffit de montrer que $A$ possède deux branches géométriques et que son invariant $\delta$ vaut $1$. On peut supposer $A = \kappa[[x, y]]/(ax^2 + bxy + cy^2)$, la forme $q = ax^2 + bxy + cy^2$ étant non dégénérée. Il y a deux cas.

\medskip\noindent
Cas I : $q$ se décompose sur $\kappa$. Après changement de coordonnées, on peut alors supposer $q = xy$. On voit que
$$
A' = \kappa[[x, y]]/(x) \times \kappa[[x, y]]/(y)
$$

\medskip\noindent
Cas II : $q$ ne se décompose pas. Dans ce cas, $c \not = 0$ et la non-dégénérescence signifie que $b^2 - 4ac \not = 0$. Ainsi $\kappa' = \kappa[t]/(a + bt + ct^2)$ est une extension séparable de degré $2$ de $\kappa$. Alors $t = y/x$ est entier sur $A$, et l'on en conclut que $$
A' = \kappa'[[x]]
$$, où $y$ s'envoie sur $tx$ dans le membre de droite.

\medskip\noindent
Dans les deux cas, on vérifie directement que l'invariant $\delta$ vaut $1$ et que le nombre de branches géométriques est $2$. Cela montre que (1) entraîne (2). La dernière assertion du lemme est également établie.

\medskip\noindent
Supposons (2). Compléments d'algèbre, lemme \ref{more-algebra-lemma-number-of-branches-1} montre que $A'$ possède soit deux idéaux maximaux, soit un seul idéal maximal pour $A'$ avec $[\kappa(\mathfrak m') : \kappa]_s = 2$.

\medskip\noindent
Cas I : $A'$ possède deux idéaux maximaux $\mathfrak m'_1$, $\mathfrak m'_2$, de corps résiduels $\kappa_1$, $\kappa_2$. Puisque l'invariant $\delta$ est la longueur de $A'/A$ et qu'il existe une surjection $A'/A \to (\kappa_1 \times \kappa_2)/\kappa$, on a $\kappa = \kappa_1 = \kappa_2$. Puisque $A$ est complet (et hensélien d'après Algèbre, lemme \ref{algebra-lemma-complete-henselian}) et que $A'$ est fini sur $A$, on a $A' = A_1 \times A_2$ (par Algèbre, lemme \ref{algebra-lemma-finite-over-henselian}). Puisque $A'$ est un anneau normal, $A_1$ et $A_2$ sont des anneaux de valuation discrète. Ainsi $A_1$ et $A_2$ sont isomorphes à $\kappa[[t]]$ (comme $k$-algèbres) d'après Compléments d'algèbre, lemme \ref{more-algebra-lemma-power-series-over-residue-field}. Puisque l'invariant $\delta$ vaut $1$, $A$ est le bouquet de $A_1$ et $A_2$ (Variétés, définition \ref{varieties-definition-wedge}). Il s'ensuit aisément que $A \cong \kappa[[x, y]]/(xy)$.

\medskip\noindent
Cas II : $A'$ possède un seul idéal maximal $\mathfrak m'$, de corps résiduel $\kappa'$, et $[\kappa' : \kappa]_s = 2$. En raisonnant exactement comme dans le cas I, on voit que $[\kappa' : \kappa] = 2$ et que $\kappa'$ est séparable sur $\kappa$. Puisque $A'$ est normal, $A'$ est isomorphe à $\kappa'[[t]]$ (voir la référence ci-dessus). Puisque $A'/A$ est de longueur $1$, on en conclut que $$
A = \{f \in \kappa'[[t]] \mid f(0) \in \kappa\}
$$ Un calcul simple montre alors que $A$ est comme dans le cas (1).
\end{proof}

\begin{lemma}
\label{curves-lemma-fitting-ideal-well-defined}
Soit $k$ un corps. Soit $A = k[[x_1, \ldots, x_n]]$. Soit $I = (f_1, \ldots, f_m) \subset A$ un idéal. Pour tout $r \geq 0$, l'idéal de $A/I$ engendré par les mineurs d'ordre $r \times r$ de la matrice $(\partial f_j/\partial x_i)$ est indépendant du choix des générateurs de $I$ et du système régulier de paramètres $x_1, \ldots, x_n$ de $A$.
\end{lemma}

\begin{proof}
La démonstration « naturelle » consiste à montrer que cet idéal est le $(n - r)$-ième idéal de Fitting d'un module de différentielles continues de $A/I$ sur $k$. Voici une démonstration directe. Si $g_1, \ldots g_l$ est un second système de générateurs de $I$, on peut écrire $g_s = \sum a_{sj}f_j$ et l'on a l'égalité de matrices $$
(\partial g_s/\partial x_i) = (a_{sj}) (\partial f_j/\partial x_i)
+ (\partial a_{sj}/\partial x_i f_j)
$$ Le dernier terme est nul dans $A/I$. Par la formule de Cauchy-Binet, l'idéal des mineurs associé aux $g_s$ est contenu dans celui associé aux $f_j$. Par symétrie, ces idéaux sont égaux. Si $y_1, \ldots, y_n \in \mathfrak m_A$ est un second système régulier de paramètres, la matrice $(\partial y_j/\partial x_i)$ est inversible et l'on peut utiliser la règle de dérivation des fonctions composées. Nous omettons quelques détails.
\end{proof}

\begin{lemma}
\label{curves-lemma-fitting-ideal}
Soit $k$ un corps. Soit $A = k[[x_1, \ldots, x_n]]$. Soit $I = (f_1, \ldots, f_m) \subset \mathfrak m_A$ un idéal. Les conditions suivantes sont équivalentes : \begin{enumerate} \item $k \to A/I$ est comme dans le lemme \ref{curves-lemma-nodal-algebraic}; \item $A/I$ est réduit et les mineurs d'ordre $(n - 1) \times (n - 1)$ de la matrice $(\partial f_j/\partial x_i)$ engendrent $I + \mathfrak m_A$; \item $\text{depth}(A/I) = 1$ et les mineurs d'ordre $(n - 1) \times (n - 1)$ de la matrice $(\partial f_j/\partial x_i)$ engendrent $I + \mathfrak m_A$.
\end{enumerate}
\end{lemma}

\begin{proof}
D'après le lemme \ref{curves-lemma-fitting-ideal-well-defined}, nous pouvons changer le système de coordonnées et le choix des générateurs au cours de la démonstration.

\medskip\noindent
Si (1) est satisfaite, nous pouvons changer de coordonnées de sorte que $x_1, \ldots, x_{n - 2}$ s'envoient sur zéro dans $A/I$ et que $A/I = k[[x_{n - 1}, x_n]]/(a x_{n - 1}^2 + b x_{n - 1}x_n + c x_n^2)$ pour une certaine forme quadratique non dégénérée $a x_{n - 1}^2 + b x_{n - 1}x_n + c x_n^2$. Un calcul explicite montre alors que (2) et (3) sont vraies.

\medskip\noindent
Supposons que les mineurs d'ordre $(n - 1) \times (n - 1)$ de la matrice $(\partial f_j/\partial x_i)$ engendrent $I + \mathfrak m_A$. Supposons que, pour certains $i$ et $j$, la dérivée partielle $\partial f_j/\partial x_i$ soit une unité dans $A$. Nous pouvons alors prendre le système de paramètres $f_j, x_1, \ldots, x_{i - 1}, \hat x_i, x_{i + 1}, \ldots, x_n$ et les générateurs $f_j, f_1, \ldots, f_{j - 1}, \hat f_j, f_{j + 1}, \ldots, f_m$ de $I$. Nous obtenons ainsi un système régulier de paramètres $x_1, \ldots, x_n$ et des générateurs $x_1, f_2, \ldots, f_m$ de $I$. Cherchons ensuite $i \geq 2$ et $j \geq 2$ tels que $\partial f_j/\partial x_i$ soit une unité dans $A$. Si une telle paire existe, nous pouvons effectuer un remplacement comme ci-dessus et supposer que nous disposons d'un système régulier de paramètres $x_1, \ldots, x_n$ et de générateurs $x_1, x_2, f_3, \ldots, f_m$ de $I$. En poursuivant, nous parvenons en un nombre fini d'étapes à un système régulier de paramètres $x_1, \ldots, x_n$ et à des générateurs $x_1, \ldots, x_t, f_{t + 1}, \ldots, f_m$ de $I$ tels que $\partial f_j/\partial x_i \in \mathfrak m_A$ pour tous $i, j \geq t + 1$.

\medskip\noindent
Dans ce cas, la matrice des dérivées partielles a la forme par blocs suivante : $$
\left(
\begin{matrix}
I_{t \times t} & * \\
0 & \mathfrak m_A
\end{matrix}
\right)
$$ Tout mineur d'ordre $(n - 1) \times (n - 1)$ appartient donc à $\mathfrak m_A^{n - 1 - t}$. Remarquons que $I \not = \mathfrak m_A$, sans quoi l'idéal des mineurs contiendrait $1$. Il s'ensuit que $n - 1 - t \leq 1$, car il existe un élément de $\mathfrak m_A \setminus \mathfrak m_A^2 + I$ (sinon $I = \mathfrak m_A$ par le lemme de Nakayama). Ainsi $t \geq n - 2$. Nous avons vu ci-dessus que $t \not = n$; de même, si $t = n - 1$, il existerait un mineur inversible d'ordre $(n - 1) \times (n - 1)$, ce qui est également exclu. Par conséquent, $t = n - 2$. Alors $A/I$ est un quotient de $k[[x_{n - 1}, x_n]]$, et le lemme \ref{curves-lemma-reduced-quotient-regular-ring-dim-2} implique, dans les deux cas (2) et (3), que $I$ est engendré par $x_1, \ldots, x_{n - 2}, f$ pour un certain $f = f(x_{n - 1}, x_n)$. Dans ce cas, la condition sur les mineurs signifie exactement que le terme quadratique de $f$ est non dégénéré, c'est-à-dire que $A/I$ est comme dans le lemme \ref{curves-lemma-nodal-algebraic}.
\end{proof}

\begin{lemma}
\label{curves-lemma-nodal}
Soit $k$ un corps. Soit $X$ un schéma algébrique de dimension $1$ sur $k$. Soit $x \in X$ un point fermé. Les conditions suivantes sont équivalentes : \begin{enumerate} \item $x$ est un nœud; \item $k \to \mathcal{O}_{X, x}$ est comme dans le lemme \ref{curves-lemma-nodal-algebraic}; \item tout $\overline{x} \in X_{\overline{k}}$ s'envoyant sur $x$ définit une singularité nodale; \item $\kappa(x)/k$ est séparable, $\mathcal{O}_{X, x}$ est réduit, et le premier idéal de Fitting de $\Omega_{X/k}$ engendre $\mathfrak m_x$ dans $\mathcal{O}_{X, x}$; \item $\kappa(x)/k$ est séparable, $\text{depth}(\mathcal{O}_{X, x}) = 1$, et le premier idéal de Fitting de $\Omega_{X/k}$ engendre $\mathfrak m_x$ dans $\mathcal{O}_{X, x}$; \item $\kappa(x)/k$ est séparable et $\mathcal{O}_{X, x}$ est réduit, a pour invariant $\delta$ la valeur $1$, et possède $2$ branches géométriques.
\end{enumerate}
\end{lemma}

\begin{proof}
Supposons d'abord que $k$ soit algébriquement clos. Dans ce cas, l'équivalence de (1) et (3) est immédiate. L'équivalence de (1) et (3) avec (2) tient au fait que, à changement de coordonnées près, la seule forme quadratique non dégénérée en deux variables est $xy$. L'équivalence de (1) et (6) est le lemme \ref{curves-lemma-multicross-algebra}. Après avoir remplacé $X$ par un voisinage affine de $x$, nous pouvons supposer qu'il existe une immersion fermée $X \to \mathbf{A}^n_k$ envoyant $x$ sur $0$. Soient $f_1, \ldots, f_m \in k[x_1, \ldots, x_n]$ des générateurs de l'idéal $I$ de $X$ dans $\mathbf{A}^n_k$. Alors $\Omega_{X/k}$ correspond au $R = k[x_1, \ldots, x_n]/I$-module $\Omega_{R/k}$ qui admet une présentation $$
R^{\oplus m} \xrightarrow{(\partial f_j/\partial x_i)}
R^{\oplus n} \to \Omega_{R/k} \to 0
$$ (voir Algèbre, sections \ref{algebra-section-differentials} et \ref{algebra-section-netherlander}). Le premier idéal de Fitting de $\Omega_{R/k}$ est donc l'idéal engendré par les mineurs d'ordre $(n - 1) \times (n - 1)$ de la matrice $(\partial f_j/\partial x_i)$. Le lemme \ref{curves-lemma-fitting-ideal}, appliqué au complété de $k[x_1, \ldots, x_n] \to R$ en l'idéal maximal $(x_1, \ldots, x_n)$, montre ainsi que (2), (4) et (5) sont équivalentes.

\medskip\noindent
Supposons maintenant que $k$ soit un corps quelconque. Dans les cas (2), (4), (5) et (6), le corps résiduel $\kappa(x)$ est séparable sur $k$. Montrons qu'il en va de même dans les cas (1) et (3). Soit donc $Z \subset X$ le sous-schéma fermé de $X$ défini par le premier idéal de Fitting de $\Omega_{X/k}$. La formation de $Z$ commute aux extensions de corps (Diviseurs, lemme \ref{divisors-lemma-base-change-and-fitting-ideal-omega}). Si (1) ou (3) est vraie, il existe un point $\overline{x}$ de $X_{\overline{k}}$ tel que $\overline{x}$ soit un point isolé de multiplicité $1$ de $Z_{\overline{k}}$ (puisque nous disposons de l'équivalence des conditions du lemme sur $\overline{k}$). En particulier, $Z_{\overline{k}}$ est géométriquement réduit en $\overline{x}$ (car $\overline{k}$ est algébriquement clos). Ainsi $Z$ est géométriquement réduit en $x$ (Variétés, lemme \ref{varieties-lemma-geometrically-reduced-upstairs}). En particulier, $Z$ est réduit en $x$; par conséquent, $Z = \Spec(\kappa(x))$ dans un voisinage de $x$ et $\kappa(x)$ est géométriquement réduit sur $k$. Cela signifie que $\kappa(x)/k$ est séparable (Algèbre, lemme \ref{algebra-lemma-characterize-separable-field-extensions}).

\medskip\noindent
L'argument du paragraphe précédent montre que, si (1) ou (3) est satisfaite, le premier idéal de Fitting de $\Omega_{X/k}$ engendre $\mathfrak m_x$. Puisque $\mathcal{O}_{X, x} \to \mathcal{O}_{X_{\overline{k}}, \overline{x}}$ est plat et que $\mathcal{O}_{X_{\overline{k}}, \overline{x}}$ est réduit et de profondeur $1$, on voit que (4) et (5) sont satisfaites (utiliser Algèbre, lemmes \ref{algebra-lemma-descent-reduced} et \ref{algebra-lemma-apply-grothendieck}). Réciproquement, (4) implique (5) d'après Algèbre, lemme \ref{algebra-lemma-criterion-reduced}. Si (5) est satisfaite, alors $Z$ est géométriquement réduit en $x$ (car $\kappa(x)/k$ est séparable et $Z$ est égal à $x$ dans un voisinage). Ainsi $Z_{\overline{k}}$ est réduit en tout point $\overline{x}$ de $X_{\overline{k}}$ situé au-dessus de $x$. Autrement dit, le premier idéal de Fitting de $\Omega_{X_{\overline{k}}/\overline{k}}$ engendre $\mathfrak m_{\overline{x}}$ dans $\mathcal{O}_{X_{\overline{k}, \overline{x}}}$. De plus, puisque $\mathcal{O}_{X, x} \to \mathcal{O}_{X_{\overline{k}}, \overline{x}}$ est plat, on voit que $\text{depth}(\mathcal{O}_{X_{\overline{k}}, \overline{x}}) = 1$ (voir la référence ci-dessus). La condition (5) est donc satisfaite pour $\overline{x} \in X_{\overline{k}}$, et nous concluons que (3) est satisfaite (par l'équivalence sur les corps algébriquement clos). Nous avons ainsi établi l'équivalence de (1), (3), (4) et (5).

\medskip\noindent
L'équivalence de (2) et (6) résulte du lemme \ref{curves-lemma-2-branches-delta-1}.

\medskip\noindent
Enfin, démontrons l'équivalence de (2) $=$ (6) avec (1) $=$ (3) $=$ (4) $=$ (5). Remarquons d'abord que le nombre géométrique de branches de $X$ en $x$ et celui de $X_{\overline{k}}$ en $\overline{x}$ sont égaux d'après Variétés, lemme \ref{varieties-lemma-geometric-branches-and-change-of-fields}. Les informations dont nous disposons à ce stade montrent que, dans tous les cas, ce nombre est égal à $2$. D'autre part, dans le cas (1), il est clair que $X$ est géométriquement réduit en $x$, et donc que $$
\text{invariant }\delta\text{ de }X\text{ en }x \leq
\text{invariant }\delta\text{ de }X_{\overline{k}}\text{ en }\overline{x}
$$ d'après Variétés, lemme \ref{varieties-lemma-delta-invariant-and-change-of-fields}. Puisque, dans le cas (1), le membre de droite vaut $1$, cela impose que l'invariant $\delta$ de $X$ en $x$ soit égal à $1$ (car, s'il était nul, $\mathcal{O}_{X, x}$ serait un anneau de valuation discrète d'après Variétés, lemme \ref{varieties-lemma-delta-invariant-is-zero}, donc serait unibranche, contradiction). Ainsi (6) est satisfaite. Réciproquement, si (2) $=$ (6) est vraie, alors les hypothèses (a), (b) et (c) de Variétés, lemme \ref{varieties-lemma-geometrically-normal-in-codim-1} sont satisfaites pour $x \in X$ d'après le lemme \ref{curves-lemma-2-branches-delta-1}. Variétés, lemme \ref{varieties-lemma-delta-invariant-and-change-of-fields-better} s'applique donc et donne l'égalité dans l'inégalité affichée ci-dessus. Nous en concluons que (6) est satisfaite pour $\overline{x} \in X_{\overline{k}}$, et l'équivalence des conditions sur un corps algébriquement clos nous ramène au cas (1).
\end{proof}

\begin{remark}[L'extension quadratique associée à un nœud] \label{curves-remark-quadratic-extension} Soit $k$ un corps. Soit $(A, \mathfrak m, \kappa)$ une $k$-algèbre locale noethérienne. Supposons soit que $(A, \mathfrak m, \kappa)$ soit comme dans le lemme \ref{curves-lemma-nodal-algebraic}, soit que $A$ soit de Nagata comme dans le lemme \ref{curves-lemma-2-branches-delta-1}, soit que $A$ soit complet et comme dans le lemme \ref{curves-lemma-fitting-ideal}. Alors $A$ définit canoniquement une algèbre séparable de degré $2$ sur $\kappa$, notée $\kappa'$, comme suit : \begin{enumerate} \item soit $q = ax^2 + bxy + cy^2$ une forme quadratique non dégénérée comme dans le lemme \ref{curves-lemma-nodal-algebraic}, les coordonnées $x, y$ étant choisies de sorte que $a \not = 0$, et posons $\kappa' = \kappa[x]/(ax^2 + bx + c)$; \item soit $A' \supset A$ la clôture intégrale de $A$ dans son anneau total des fractions, et posons $\kappa' = A'/\mathfrak m A'$; \item soit $\kappa'$ la $\kappa$-algèbre telle que $\text{Proj}(\bigoplus_{n \geq 0} \mathfrak m^n/\mathfrak m^{n + 1}) =
\Spec(\kappa')$. \end{enumerate} L'équivalence de (1) et (2) a été établie dans la démonstration du lemme \ref{curves-lemma-2-branches-delta-1}. Nous omettons la démonstration de leur équivalence avec (3). Si $X$ est un $k$-schéma localement noethérien et si $x \in X$ est un point tel que $\mathcal{O}_{X, x} = A$, alors (3) montre que $\Spec(\kappa') = X^\nu \times_X \Spec(\kappa)$, où $\nu : X^\nu \to X$ est le morphisme de normalisation.
\end{remark}

\begin{remark}[Extension quadratique triviale] \label{curves-remark-trivial-quadratic-extension} Soit $k$ un corps. Soit $(A, \mathfrak m, \kappa)$ comme dans la remarque \ref{curves-remark-quadratic-extension}, et soit $\kappa'/\kappa$ l'algèbre séparable de degré $2$ associée. Les conditions suivantes sont équivalentes : \begin{enumerate} \item $\kappa' \cong \kappa \times \kappa$ comme $\kappa$-algèbre; \item la forme $q$ du lemme \ref{curves-lemma-nodal-algebraic} peut être choisie égale à $xy$; \item $A$ possède deux branches; \item l'extension $A'/A$ du lemme \ref{curves-lemma-2-branches-delta-1} possède deux idéaux maximaux; \item $A^\wedge \cong \kappa[[x, y]]/(xy)$ comme $k$-algèbre. \end{enumerate} L'équivalence de ces conditions a été établie dans la démonstration du lemme \ref{curves-lemma-2-branches-delta-1}. Si $X$ est un $k$-schéma localement noethérien et si $x \in X$ est un point tel que $\mathcal{O}_{X, x} = A$, cela signifie exactement qu'il existe deux points $x_1, x_2$ de la normalisation $X^\nu$ au-dessus de $x$ et que $\kappa(x) = \kappa(x_1) = \kappa(x_2)$.
\end{remark}

\begin{definition}
\label{curves-definition-split-node}
Soit $k$ un corps. Soit $X$ un schéma algébrique de dimension $1$ sur $k$. Soit $x \in X$ un point fermé. On dit que $x$ est un {\it nœud déployé} si $x$ est un nœud, si $\kappa(x) = k$, et si les assertions équivalentes de la remarque \ref{curves-remark-trivial-quadratic-extension} sont satisfaites pour $A = \mathcal{O}_{X, x}$.
\end{definition}

\noindent
Formulons le lemme obligé qui rassemble ce que nous savons déjà de cette notion.

\begin{lemma}
\label{curves-lemma-split-node}
Soit $k$ un corps. Soit $X$ un schéma algébrique de dimension $1$ sur $k$. Soit $x \in X$ un point fermé. Les conditions suivantes sont équivalentes : \begin{enumerate} \item $x$ est un nœud déployé; \item $x$ est un nœud et il existe exactement deux points $x_1, x_2$ de la normalisation $X^\nu$ au-dessus de $x$, avec $k = \kappa(x_1) = \kappa(x_2)$; \item $\mathcal{O}_{X, x}^\wedge \cong k[[x, y]]/(xy)$ comme $k$-algèbre; \item ajouter davantage ici.
\end{enumerate}
\end{lemma}

\begin{proof}
Cela résulte de la discussion de la remarque \ref{curves-remark-trivial-quadratic-extension} et du lemme \ref{curves-lemma-nodal}.
\end{proof}

\begin{lemma}
\label{curves-lemma-node-field-extension}
Soit $K/k$ une extension de corps. Soit $X$ un $k$-schéma localement algébrique de dimension $1$. Soit $y \in X_K$ un point d'image $x \in X$. Les conditions suivantes sont équivalentes : \begin{enumerate} \item $x$ est un point fermé de $X$ et un nœud; \item $y$ est un point fermé de $X_K$ et un nœud.
\end{enumerate}
\end{lemma}

\begin{proof}
Si $x$ est un point fermé de $X$, alors $y$ l'est aussi (considérer les corps résiduels). La réciproque n'est toutefois pas nécessairement vraie : il peut arriver qu'un point fermé de $X_K$ s'envoie sur un point non fermé de $X$. Dans ce cas, cependant, $y$ ne peut pas être un nœud. En effet, $X$ serait alors géométriquement unibranche en $x$ (car $x$ serait un point générique de $X$, tandis que $\mathcal{O}_{X, x}$ serait artinien et que tout anneau local artinien est géométriquement unibranche); par conséquent, $X_K$ serait géométriquement unibranche en $y$ (Variétés, lemme \ref{varieties-lemma-geometrically-unibranch-and-change-of-fields}), ce qui signifie que $y$ ne peut pas être un nœud d'après le lemme \ref{curves-lemma-nodal}. Nous pouvons donc supposer, et nous supposerons, que $x$ et $y$ sont tous deux des points fermés.

\medskip\noindent
Choisissons des clôtures algébriques $\overline{k}$, $\overline{K}$ et une application $\overline{k} \to \overline{K}$ prolongeant l'application donnée $k \to K$. En utilisant l'équivalence de (1) et (3) dans le lemme \ref{curves-lemma-nodal}, nous nous ramenons au cas où $k$ et $K$ sont algébriquement clos. Dans ce cas, nous pouvons raisonner comme dans la démonstration du lemme \ref{curves-lemma-nodal} pour voir que le nombre géométrique de branches et les invariants $\delta$ de $X$ en $x$ et de $X_K$ en $y$ sont les mêmes. On peut aussi choisir un isomorphisme $k[[x_1, \ldots, x_n]]/(g_1, \ldots, g_m) \to \mathcal{O}_{X, x}^\wedge$ de $k$-algèbres comme dans Variétés, lemme \ref{varieties-lemma-complete-local-ring-structure-as-algebra}. D'après Variétés, lemme \ref{varieties-lemma-base-change-complete-local-ring}, il induit un isomorphisme $K[[x_1, \ldots, x_n]]/(g_1, \ldots, g_m) \to \mathcal{O}_{X_K, y}^\wedge$ de $K$-algèbres. Par définition, il faut montrer que $$
k[[x_1, \ldots, x_n]]/(g_1, \ldots, g_m) \cong k[[s, t]]/(st)
$$ si et seulement si $$
K[[x_1, \ldots, x_n]]/(g_1, \ldots, g_m) \cong K[[s, t]]/(st)
$$ Nous invitons le lecteur à le démontrer. Puisque $k$ et $K$ sont des corps algébriquement clos, cela revient à demander que ces anneaux soient comme dans le lemme \ref{curves-lemma-nodal-algebraic}. Par le lemme \ref{curves-lemma-fitting-ideal}, cette condition se traduit en une assertion sur les mineurs d'ordre $(n - 1) \times (n - 1)$ de la matrice $(\partial g_j/\partial x_i)$, manifestement indépendante du corps utilisé. Nous omettons les détails.
\end{proof}

\begin{lemma}
\label{curves-lemma-node-etale-local}
Soit $k$ un corps. Soit $X$ un $k$-schéma localement algébrique de dimension $1$. Soit $Y \to X$ un morphisme étale. Soit $y \in Y$ un point d'image $x \in X$. Les conditions suivantes sont équivalentes : \begin{enumerate} \item $x$ est un point fermé de $X$ et un nœud; \item $y$ est un point fermé de $Y$ et un nœud.
\end{enumerate}
\end{lemma}

\begin{proof}
D'après le lemme \ref{curves-lemma-node-field-extension}, nous pouvons effectuer un changement de base vers la clôture algébrique de $k$. Les corps résiduels de $x$ et $y$ sont alors $k$. L'application $\mathcal{O}_{X, x}^\wedge \to \mathcal{O}_{Y, y}^\wedge$ est donc un isomorphisme (par exemple d'après Morphismes étales, lemme \ref{etale-lemma-characterize-etale-completions} ou Compléments d'algèbre, lemme \ref{more-algebra-lemma-flat-unramified}). Le lemme est dès lors clair.
\end{proof}

\begin{lemma}
\label{curves-lemma-node-over-separable-extension}
Soit $k'/k$ une extension finie séparable de corps. Soit $X$ un $k'$-schéma localement algébrique de dimension $1$. Soit $x \in X$ un point fermé. Les conditions suivantes sont équivalentes : \begin{enumerate} \item $x$ est un nœud; \item $x$ est un nœud lorsque $X$ est considéré comme un $k$-schéma localement algébrique.
\end{enumerate}
\end{lemma}

\begin{proof}
Cela résulte immédiatement de la caractérisation des nœuds du lemme \ref{curves-lemma-nodal}.
\end{proof}

\begin{lemma}
\label{curves-lemma-nodal-lci}
Soit $k$ un corps. Soit $X$ un $k$-schéma localement algébrique équidimensionnel de dimension $1$. Les conditions suivantes sont équivalentes : \begin{enumerate} \item les singularités de $X$ sont au plus nodales; \item $X$ est une intersection complète locale sur $k$ et le sous-schéma fermé $Z \subset X$ découpé par le premier idéal de Fitting de $\Omega_{X/k}$ est non ramifié sur $k$.
\end{enumerate}
\end{lemma}

\begin{proof}
Nous invitons vivement le lecteur à trouver sa propre démonstration de ce lemme; celle qui suit ne fait que réunir des résultats antérieurs et risque de masquer ce qui se passe réellement.

\medskip\noindent
Supposons (2). Puisque $Z \to \Spec(k)$ est quasi-fini (Morphismes, lemme \ref{morphisms-lemma-unramified-quasi-finite}), les corps résiduels des points $x \in Z$ sont finis sur $k$ (et séparables), d'après Morphismes, lemme \ref{morphisms-lemma-residue-field-quasi-finite}. Chaque $x \in Z$ est donc un point fermé de $X$ d'après Morphismes, lemme \ref{morphisms-lemma-algebraic-residue-field-extension-closed-point-fibre}. L'anneau local $\mathcal{O}_{X, x}$ est de Cohen--Macaulay d'après Algèbre, lemme \ref{algebra-lemma-lci-CM}. Comme $\dim(\mathcal{O}_{X, x}) = 1$ par la théorie de la dimension (Variétés, section \ref{varieties-section-algebraic-schemes}), nous en concluons que $\text{depth}(\mathcal{O}_{X, x}) = 1$. Ainsi $x$ est un nœud d'après le lemme \ref{curves-lemma-nodal}. Si $x \in X$, $x \not \in Z$, alors $X \to \Spec(k)$ est lisse en $x$ d'après Diviseurs, lemme \ref{divisors-lemma-d-fitting-ideal-omega-smooth}.

\medskip\noindent
Supposons (1). Sous cette hypothèse, $X$ est géométriquement réduit en tout point fermé (voir Variétés, lemme \ref{varieties-lemma-geometrically-reduced-upstairs}). Ainsi $X \to \Spec(k)$ est lisse sur un ouvert dense d'après Variétés, lemme \ref{varieties-lemma-geometrically-reduced-dense-smooth-open}. Par conséquent, $Z$ est fermé et formé de points fermés. D'après Diviseurs, lemme \ref{divisors-lemma-d-fitting-ideal-omega-smooth}, le morphisme $X \setminus Z \to \Spec(k)$ est lisse. Ainsi $X \setminus Z$ est une intersection complète locale d'après Morphismes, lemme \ref{morphisms-lemma-smooth-syntomic} et la définition d'une intersection complète locale de Morphismes, définition \ref{morphisms-definition-syntomic}. D'après le lemme \ref{curves-lemma-nodal}, pour tout point $x \in Z$, l'anneau local $\mathcal{O}_{Z, x}$ est égal à $\kappa(x)$ et $\kappa(x)$ est séparable sur $k$. Ainsi $Z \to \Spec(k)$ est non ramifié (Morphismes, lemme \ref{morphisms-lemma-unramified-over-field}). Enfin, le lemme \ref{curves-lemma-nodal}, par la partie (3) du lemme \ref{curves-lemma-nodal-algebraic}, montre que $\mathcal{O}_{X, x}$ est une intersection complète au sens d'Algèbre à puissances divisées, définition \ref{dpa-definition-lci}. Or Algèbre à puissances divisées, lemme \ref{dpa-lemma-check-lci-agrees} et Morphismes, lemme \ref{morphisms-lemma-local-complete-intersection} montrent que cette notion coïncide avec celle utilisée pour définir un schéma intersection complète locale sur un corps, ce qui achève la démonstration.
\end{proof}

\begin{lemma}
\label{curves-lemma-facts-about-nodal-curves}
Soit $k$ un corps. Soit $X$ un $k$-schéma localement algébrique équidimensionnel de dimension $1$ dont les singularités sont au plus nodales. Alors $X$ est de Gorenstein et géométriquement réduit.
\end{lemma}

\begin{proof}
L'assertion de Gorenstein résulte du lemme \ref{curves-lemma-nodal-lci} et de Dualité pour les schémas, lemme \ref{duality-lemma-gorenstein-lci}. On peut aussi utiliser le fait qu'il suffit de vérifier l'assertion après passage à la clôture algébrique (Dualité pour les schémas, lemme \ref{duality-lemma-gorenstein-base-change}), puis qu'un anneau local noethérien est de Gorenstein si et seulement si son complété l'est (Complexes dualisants, lemme \ref{dualizing-lemma-flat-under-gorenstein}), et enfin démontrer directement que les anneaux locaux $k[[t]]$ et $k[[x, y]]/(xy)$ sont de Gorenstein.

\medskip\noindent
Pour voir que $X$ est géométriquement réduit, il suffit de montrer que $X_{\overline{k}}$ est réduit (Variétés, lemmes \ref{varieties-lemma-perfect-reduced} et \ref{varieties-lemma-geometrically-reduced}). Or $X_{\overline{k}}$ est une courbe nodale sur un corps algébriquement clos. Les anneaux locaux complets de $X_{\overline{k}}$ sont donc isomorphes soit à $\overline{k}[[t]]$, soit à $\overline{k}[[x, y]]/(xy)$, qui sont réduits, comme voulu.
\end{proof}

\begin{lemma}
\label{curves-lemma-closed-subscheme-nodal-curve}
Soit $k$ un corps. Soit $X$ un $k$-schéma localement algébrique équidimensionnel de dimension $1$ dont les singularités sont au plus nodales. Si $Y \subset X$ est un sous-schéma fermé réduit équidimensionnel de dimension $1$, alors \begin{enumerate} \item les singularités de $Y$ sont au plus nodales; \item si $Z \subset X$ est l'adhérence schématique de $X \setminus Y$, alors \begin{enumerate} \item l'intersection schématique $Y \cap Z$ est la réunion disjointe de spectres d'extensions finies séparables de $k$; \item tout point de $Y \cap Z$ est un nœud de $X$; \item $Y \to \Spec(k)$ est lisse en tout point de $Y \cap Z$.
\end{enumerate}
\end{enumerate}
\end{lemma}

\begin{proof}
Puisque $X$ et $Y$ sont réduits et équidimensionnels de dimension $1$, on voit que $Y$ est la réunion schématique d'un sous-ensemble des composantes irréductibles de $X$ (dans un anneau réduit, $(0)$ est l'intersection des idéaux premiers minimaux). Soit $y \in Y$ un point fermé. Si $y$ appartient au lieu lisse de $X \to \Spec(k)$, alors $y$ se trouve sur une unique composante irréductible de $X$, et $Y$ et $X$ coïncident dans un voisinage ouvert de $y$. Ainsi $Y \to \Spec(k)$ est lisse en $y$. Si $y$ est un nœud de $X$ mais appartient encore à une unique composante irréductible de $X$, le même raisonnement montre que $y$ est un nœud de $Y$. Supposons que $y$ appartienne à plus de $1$ composante irréductible de $X$. Puisque le nombre de branches géométriques de $X$ en $y$ est $2$ d'après le lemme \ref{curves-lemma-nodal}, au plus $2$ composantes irréductibles peuvent passer par $y$, d'après Propriétés, lemme \ref{properties-lemma-number-of-branches-irreducible-components}. Si $Y$ contient ces deux composantes, alors $Y = X$ dans un voisinage ouvert de $y$, et $y$ est un nœud de $Y$. Supposons enfin que $Y$ ne contienne qu'une des composantes irréductibles. Après avoir remplacé $X$ par un voisinage ouvert de $y$, nous pouvons supposer que $Y$ est l'une des deux composantes irréductibles et que $Z$ est l'autre. D'après Propriétés, lemme \ref{properties-lemma-number-of-branches-irreducible-components}, on voit encore que $X$ possède deux branches en $y$, c'est-à-dire que l'anneau local $\mathcal{O}_{X, y}$ possède deux branches et que celles-ci proviennent de $\mathcal{O}_{Y, y}$ et $\mathcal{O}_{Z, y}$. Écrivons $\mathcal{O}_{X, y}^\wedge \cong \kappa(y)[[u, v]]/(uv)$ comme dans la remarque \ref{curves-remark-trivial-quadratic-extension}. Le corps $\kappa(y)$ est fini séparable sur $k$, par exemple d'après le lemme \ref{curves-lemma-nodal}. Ainsi, après avoir éventuellement interverti les rôles de $u$ et $v$, le complété de l'application $\mathcal{O}_{X, y} \to \mathcal{O}_{Y, y}$ correspond à $\kappa(y)[[u, v]]/(uv) \to \kappa(y)[[u]]$ et celui de l'application $\mathcal{O}_{X, y} \to \mathcal{O}_{Z, y}$ correspond à $\kappa(y)[[u, v]]/(uv) \to \kappa(y)[[v]]$. L'intersection schématique de $Y \cap Z$ est découpée par la somme de leurs idéaux, qui devient $(u, v)$ dans le complété, c'est-à-dire l'idéal maximal. Les assertions (2)(a) et (2)(b) sont donc claires. Enfin, (2)(c) est satisfaite : le complété de $\mathcal{O}_{Y, y}$ est régulier, donc $\mathcal{O}_{Y, y}$ est régulier (Compléments d'algèbre, lemme \ref{more-algebra-lemma-completion-regular}), et $\kappa(y)/k$ est séparable; on obtient donc la lissité dans un voisinage ouvert d'après Algèbre, lemme \ref{algebra-lemma-separable-smooth}.
\end{proof}




\section{Familles de courbes nodales}
\label{curves-section-families-nodal}

\noindent
Dans le projet Champs, les courbes sont des variétés irréductibles de dimension $1$. Dans la littérature, toutefois, une « courbe semi-stable » ou une « courbe nodale » n'est généralement pas irréductible et est souvent supposée propre, en particulier dans des expressions telles que « famille de courbes semi-stables », « famille de courbes nodales » ou « famille nodale ». Il nous est donc quelque peu difficile de choisir une terminologie à la fois conforme à la littérature et cohérente en interne. De plus, nous souhaitons réellement commencer par étudier la notion introduite dans le lemme suivant, qui est locale sur la source.

\begin{lemma}
\label{curves-lemma-nodal-family}
Soit $f : X \to S$ un morphisme de schémas. Les conditions suivantes sont équivalentes : \begin{enumerate} \item $f$ est plat et localement de présentation finie, toute fibre non vide $X_s$ est équidimensionnelle de dimension $1$, et les singularités de $X_s$ sont au plus nodales; \item $f$ est syntomique de dimension relative $1$ et le sous-schéma fermé $\text{Sing}(f) \subset X$ défini par le premier idéal de Fitting de $\Omega_{X/S}$ est non ramifié sur $S$.
\end{enumerate}
\end{lemma}

\begin{proof}
Rappelons que la formation de $\text{Sing}(f)$ commute au changement de base; voir Diviseurs, lemme \ref{divisors-lemma-base-change-and-fitting-ideal-omega}. Le lemme résulte donc du lemme \ref{curves-lemma-nodal-lci} et de Morphismes, lemmes \ref{morphisms-lemma-syntomic-flat-fibres} et \ref{morphisms-lemma-unramified-etale-fibres}. (Nous utilisons aussi les triviaux Morphismes, lemmes \ref{morphisms-lemma-syntomic-locally-finite-presentation} et \ref{morphisms-lemma-syntomic-flat}.)
\end{proof}

\begin{definition}
\label{curves-definition-nodal-family}
Soit $f : X \to S$ un morphisme de schémas. On dit que $f$ est {\it à singularités au plus nodales de dimension relative $1$} si $f$ satisfait les conditions équivalentes du lemme \ref{curves-lemma-nodal-family}.
\end{definition}

\noindent
Voici quelques raisons qui expliquent cette terminologie pesante\footnote{Mais n'hésitez pas à écrire au responsable du projet Champs si vous avez une meilleure suggestion.}. Premièrement, nous voulons éviter toute confusion entre cette notion et les autres notions de la littérature (voir l'introduction de cette section). Deuxièmement, nous pouvons imaginer plusieurs généralisations de cette notion à des morphismes de dimension relative supérieure : on peut, par exemple, demander que les morphismes soient étale-localement des composés de morphismes à singularités au plus nodales, ou demander que leurs fibres soient de dimension supérieure mais n'aient au pire que des points doubles ordinaires.

\begin{lemma}
\label{curves-lemma-smooth-relative-dimension-1}
Un morphisme lisse de dimension relative $1$ est à singularités au plus nodales de dimension relative $1$.
\end{lemma}

\begin{proof}
Omis.
\end{proof}

\begin{lemma}
\label{curves-lemma-base-change-nodal-family}
Soit $f : X \to S$ à singularités au plus nodales de dimension relative $1$. Il en va de même après tout changement de base de $f$.
\end{lemma}

\begin{proof}
Cela tient au fait que le changement de base d'un morphisme syntomique est syntomique (Morphismes, lemme \ref{morphisms-lemma-base-change-syntomic}), que le changement de base d'un morphisme de dimension relative $1$ est de dimension relative $1$ (Morphismes, lemme \ref{morphisms-lemma-base-change-relative-dimension-d}), que la formation de $\text{Sing}(f)$ commute au changement de base (Diviseurs, lemme \ref{divisors-lemma-base-change-and-fitting-ideal-omega}), et que le changement de base d'un morphisme non ramifié est non ramifié (Morphismes, lemme \ref{morphisms-lemma-base-change-unramified}).
\end{proof}

\noindent
Le lemme suivant montre que l'on peut vérifier fibre par fibre qu'un morphisme est à singularités au plus nodales de dimension relative $1$.

\begin{lemma}
\label{curves-lemma-locus-where-nodal}
Soit $f : X \to S$ un morphisme de schémas, plat et localement de présentation finie. Il existe un plus grand sous-schéma ouvert $U \subset X$ tel que $f|_U : U \to S$ soit à singularités au plus nodales de dimension relative $1$. De plus, la formation de $U$ commute à tout changement de base.
\end{lemma}

\begin{proof}
D'après Morphismes, lemme \ref{morphisms-lemma-set-points-where-fibres-lci}, il existe un tel ouvert sur lequel $f$ est syntomique. Nous pouvons donc supposer que $f$ est un morphisme syntomique. En particulier, $f$ est un morphisme de Cohen--Macaulay (Dualité pour les schémas, lemmes \ref{duality-lemma-lci-gorenstein} et \ref{duality-lemma-gorenstein-CM-morphism}). Ainsi $X$ est la réunion disjointe de sous-schémas ouverts et fermés sur lesquels $f$ a une dimension relative fixée; voir Morphismes, lemme \ref{morphisms-lemma-flat-finite-presentation-CM-fibres-relative-dimension}. Cette décomposition est préservée par tout changement de base; voir Morphismes, lemme \ref{morphisms-lemma-base-change-relative-dimension-d}. En ne gardant qu'un morceau, nous pouvons supposer que $f$ est syntomique de dimension relative $1$. Soit $\text{Sing}(f) \subset X$ le sous-schéma fermé défini par le premier idéal de Fitting de $\Omega_{X/S}$. Il existe un plus grand sous-schéma ouvert $W \subset \text{Sing}(f)$ tel que $W \to S$ soit non ramifié, et sa formation commute au changement de base (Morphismes, lemme \ref{morphisms-lemma-set-points-where-fibres-unramified}). Puisque la formation de $\text{Sing}(f)$ commute elle aussi au changement de base (Diviseurs, lemme \ref{divisors-lemma-base-change-and-fitting-ideal-omega}), nous voyons que $$
U = (X \setminus \text{Sing}(f)) \cup W
$$ est le plus grand sous-schéma ouvert de $X$ tel que $f|_U : U \to S$ soit à singularités au plus nodales de dimension relative $1$, et que la formation de $U$ commute au changement de base.
\end{proof}

\begin{lemma}
\label{curves-lemma-nodal-family-precompose-etale}
Soit $f : X \to S$ à singularités au plus nodales de dimension relative $1$. Si $Y \to X$ est un morphisme étale, alors le composé $g : Y \to S$ est à singularités au plus nodales de dimension relative $1$.
\end{lemma}

\begin{proof}
Remarquons que $g$ est plat et localement de présentation finie, comme composé de morphismes plats et localement de présentation finie (utiliser Morphismes, lemmes \ref{morphisms-lemma-etale-locally-finite-presentation}, \ref{morphisms-lemma-etale-flat}, \ref{morphisms-lemma-composition-finite-presentation} et \ref{morphisms-lemma-composition-flat}). Il suffit donc de démontrer que ses fibres ont des singularités au plus nodales. Cela résulte du lemme \ref{curves-lemma-node-etale-local} (et du fait que le composé d'un morphisme étale et d'un morphisme lisse est lisse, d'après Morphismes, lemmes \ref{morphisms-lemma-etale-smooth-unramified} et \ref{morphisms-lemma-composition-smooth}).
\end{proof}

\begin{lemma}
\label{curves-lemma-nodal-family-postcompose-etale}
Soit $S' \to S$ un morphisme étale de schémas. Soit $f : X \to S'$ à singularités au plus nodales de dimension relative $1$. Alors le composé $g : X \to S$ est à singularités au plus nodales de dimension relative $1$.
\end{lemma}

\begin{proof}
Remarquons que $g$ est plat et localement de présentation finie, comme composé de morphismes plats et localement de présentation finie (utiliser Morphismes, lemmes \ref{morphisms-lemma-etale-locally-finite-presentation}, \ref{morphisms-lemma-etale-flat}, \ref{morphisms-lemma-composition-finite-presentation} et \ref{morphisms-lemma-composition-flat}). Il suffit donc de démontrer que les fibres de $g$ ont des singularités au plus nodales. Cela résulte du lemme \ref{curves-lemma-node-over-separable-extension} et du résultat analogue pour les points lisses.
\end{proof}

\begin{lemma}
\label{curves-lemma-nodal-family-etale-local-source}
Soit $f : X \to S$ un morphisme de schémas. Soit $\{U_i \to X\}$ un recouvrement étale. Les conditions suivantes sont équivalentes : \begin{enumerate} \item $f$ est à singularités au plus nodales de dimension relative $1$; \item chaque $U_i \to S$ est à singularités au plus nodales de dimension relative $1$. \end{enumerate} Autrement dit, la propriété d'être à singularités au plus nodales de dimension relative $1$ est locale pour la topologie étale sur la source.
\end{lemma}

\begin{proof}
Une implication a été établie dans le lemme \ref{curves-lemma-nodal-family-precompose-etale}. Pour l'autre, remarquons que les propriétés d'être localement de présentation finie, d'être plat et d'avoir une dimension relative $1$ sont locales pour la topologie étale sur la source (Descente, lemmes \ref{descent-lemma-locally-finite-presentation-fppf-local-source}, \ref{descent-lemma-flat-fpqc-local-source} et \ref{descent-lemma-dimension-at-point}). En prenant les fibres, nous nous ramenons au cas où $S$ est le spectre d'un corps. Le résultat découle alors du lemme \ref{curves-lemma-node-etale-local} (et du fait que la lissité est locale pour la topologie étale sur la source, d'après Descente, lemme \ref{descent-lemma-smooth-smooth-local-source}).
\end{proof}

\begin{lemma}
\label{curves-lemma-nodal-family-fpqc-local-target}
Soit $f : X \to S$ un morphisme de schémas. Soit $\{U_i \to S\}$ un recouvrement fpqc. Les conditions suivantes sont équivalentes : \begin{enumerate} \item $f$ est à singularités au plus nodales de dimension relative $1$; \item chaque $X \times_S U_i \to U_i$ est à singularités au plus nodales de dimension relative $1$. \end{enumerate} Autrement dit, la propriété d'être à singularités au plus nodales de dimension relative $1$ est locale pour la topologie fpqc sur la cible.
\end{lemma}

\begin{proof}
Une implication a été établie dans le lemme \ref{curves-lemma-base-change-nodal-family}. Pour l'autre, remarquons que les propriétés d'être localement de présentation finie, d'être plat et d'avoir une dimension relative $1$ sont locales pour la topologie fpqc sur la cible (Descente, lemmes \ref{descent-lemma-descending-property-locally-finite-presentation} et \ref{descent-lemma-descending-property-flat}, et Morphismes, lemme \ref{morphisms-lemma-dimension-fibre-after-base-change}). En prenant les fibres, nous nous ramenons au cas où $S$ est le spectre d'un corps. Le résultat découle alors du lemme \ref{curves-lemma-node-field-extension} (et du fait que la lissité est locale pour la topologie fpqc sur la cible, d'après Descente, lemme \ref{descent-lemma-descending-property-smooth}).
\end{proof}

\begin{lemma}
\label{curves-lemma-descend-nodal-family}
Soit $S = \lim S_i$ la limite d'un système inductif de schémas dont les morphismes de transition sont affines. Soit $0 \in I$, et soit $f_0 : X_0 \to Y_0$ un morphisme de schémas au-dessus de $S_0$. Supposons $S_0$, $X_0$ et $Y_0$ quasi-compacts et quasi-séparés. Soit $f_i : X_i \to Y_i$ le changement de base de $f_0$ à $S_i$, et soit $f : X \to Y$ le changement de base de $f_0$ à $S$. Si \begin{enumerate} \item $f$ est à singularités au plus nodales de dimension relative $1$; \item $f_0$ est localement de présentation finie, \end{enumerate} alors il existe $i \geq 0$ tel que $f_i$ soit à singularités au plus nodales de dimension relative $1$.
\end{lemma}

\begin{proof}
D'après Limites, lemme \ref{limits-lemma-descend-syntomic}, il existe $i$ tel que $f_i$ soit syntomique. Alors $X_i = \coprod_{d \geq 0} X_{i, d}$ est la réunion disjointe de sous-schémas ouverts et fermés tels que $X_{i, d} \to Y_i$ ait une dimension relative $d$; voir Morphismes, lemme \ref{morphisms-lemma-syntomic-relative-dimension}. Le comportement de la dimension des fibres par changement de base donné dans Morphismes, lemme \ref{morphisms-lemma-dimension-fibre-after-base-change} montre que $X \to X_i$ s'envoie dans $X_{i, 1}$. Il existe alors $i' \geq i$ tel que $X_{i'} \to X_i$ s'envoie dans $X_{i, 1}$; voir Limites, lemme \ref{limits-lemma-limit-contained-in-constructible}. Ainsi $f_{i'} : X_{i'} \to Y_{i'}$ est syntomique de dimension relative $1$ (de nouveau par Morphismes, lemme \ref{morphisms-lemma-dimension-fibre-after-base-change}). Considérons le morphisme $\text{Sing}(f_{i'}) \to Y_{i'}$. Nous savons que son changement de base à $Y$ est non ramifié. D'après Limites, lemme \ref{limits-lemma-descend-unramified}, après avoir augmenté $i'$, le morphisme $\text{Sing}(f_{i'}) \to Y_{i'}$ devient donc non ramifié. Cela achève la démonstration.
\end{proof}

\begin{lemma}
\label{curves-lemma-formal-local-structure-nodal-family}
Soit $f : T \to S$ un morphisme de schémas. Soit $t \in T$, d'image $s \in S$. Supposons que \begin{enumerate} \item $f$ soit plat en $t$; \item $\mathcal{O}_{S, s}$ soit noethérien; \item $f$ soit localement de type fini; \item $t$ soit un nœud déployé de la fibre $T_s$. \end{enumerate} Alors il existe $h \in \mathfrak m_s^\wedge$ et un isomorphisme $$
\mathcal{O}_{T, t}^\wedge \cong
\mathcal{O}_{S, s}^\wedge[[x, y]]/(xy - h)
$$ de $\mathcal{O}_{S, s}^\wedge$-algèbres.
\end{lemma}

\begin{proof}
Remplaçons $S$ par $\Spec(\mathcal{O}_{S, s})$ et $T$ par son changement de base à $\Spec(\mathcal{O}_{S, s})$. Alors $T$ est localement noethérien, et donc $\mathcal{O}_{T, t}$ est noethérien. Posons $A = \mathcal{O}_{S, s}^\wedge$, $\mathfrak m = \mathfrak m_A$ et $B = \mathcal{O}_{T, t}^\wedge$. D'après Compléments d'algèbre, lemme \ref{more-algebra-lemma-flat-completion}, $A \to B$ est plat. Puisque $\mathcal{O}_{T, t}/\mathfrak m_s \mathcal{O}_{T, t} = \mathcal{O}_{T_s, t}$, on voit que $B/\mathfrak m B = \mathcal{O}_{T_s, t}^\wedge$. L'hypothèse (4) et le lemme \ref{curves-lemma-split-node} montrent qu'il existe $\overline{u}, \overline{v} \in B/\mathfrak m B$ tels que l'application $$
(A/\mathfrak m)[[x, y]] \longrightarrow B/\mathfrak m B,\quad
x \longmapsto \overline{u},
y \longmapsto \overline{v}
$$ soit surjective, de noyau $(xy)$.

\medskip\noindent
Supposons donnés $n \geq 1$ et $u, v \in B$ s'envoyant sur $\overline{u}, \overline{v}$ tels que $$
u v = h + \delta
$$ pour certains $h \in A$ et $\delta \in \mathfrak m^nB$. Nous affirmons qu'il existe $u', v' \in B$ avec $u - u', v - v' \in \mathfrak m^n B$ tels que $$
u' v' = h' + \delta'
$$ pour certains $h' \in A$ et $\delta' \in \mathfrak m^{n + 1}B$. Pour le voir, écrivons $\delta = \sum f_i b_i$ avec $f_i \in \mathfrak m^n$ et $b_i \in B$. Écrivons ensuite $b_i = a_i + u b_{i, 1} + v b_{i, 2} + \delta_i$ avec $a_i \in A$, $b_{i, 1}, b_{i, 2} \in B$ et $\delta_i \in \mathfrak m B$. C'est possible parce que le corps résiduel de $B$ coïncide avec celui de $A$, et que les images de $u$ et $v$ dans $B/\mathfrak m B$ engendrent l'idéal maximal. Posons alors $$
u' = u - \sum b_{i, 2}f_i,\quad
v' = v - \sum b_{i, 1}f_i
$$; nous obtenons $$
u'v' = h + \delta - \sum (b_{i, 1}u + b_{i, 2}v)f_i + \sum
c_{ij}f_if_j =
h + \sum a_if_i + \sum f_i \delta_i + \sum c_{ij}f_if_j
$$ pour un certain $c_{i, j} \in B$. Nous avons donc une formule comme ci-dessus avec $h' = h + \sum a_if_i$ et $\delta' = \sum f_i \delta_i + \sum c_{ij}f_if_j$.

\medskip\noindent
En raisonnant par récurrence à partir de relèvements quelconques $u_1, v_1 \in B$ de $\overline{u}, \overline{v}$, le paragraphe précédent fournit des suites d'éléments $u_n, v_n \in B$ et $h_n \in A$ telles que $u_n - u_{n + 1} \in \mathfrak m^n B$, $v_n - v_{n + 1} \in \mathfrak m^n B$, $h_n - h_{n + 1} \in \mathfrak m^n$, et telles que $u_n v_n - h_n \in \mathfrak m^n B$. Puisque $A$ et $B$ sont complets, nous pouvons poser $u_\infty = \lim u_n$, $v_\infty = \lim v_n$ et $h_\infty = \lim h_n$, et nous obtenons alors $u_\infty v_\infty = h_\infty$ dans $B$. Nous disposons donc d'un homomorphisme de $A$-algèbres $$
A[[x, y]]/(xy - h_\infty) \longrightarrow B
$$ envoyant $x$ sur $u_\infty$ et $v$ sur $v_\infty$. C'est une application d'algèbres $A$-plates qui devient un isomorphisme après division par $\mathfrak m$. Elle est surjective modulo $\mathfrak m$, donc surjective par complétude et par Algèbre, lemme \ref{algebra-lemma-completion-generalities}. Nous pouvons alors appliquer Algèbre, lemme \ref{algebra-lemma-mod-injective} pour conclure que c'est un isomorphisme.
\end{proof}

\begin{lemma}
\label{curves-lemma-genus-in-nodal-family-of-curves}
Soit $f : X \to S$ un morphisme de schémas. Supposons que \begin{enumerate} \item $f$ soit propre; \item $f$ soit à singularités au plus nodales de dimension relative $1$; \item les fibres géométriques de $f$ soient connexes. \end{enumerate} Alors (a) $f_*\mathcal{O}_X = \mathcal{O}_S$, et cela reste vrai après tout changement de base; (b) $R^1f_*\mathcal{O}_X$ est un $\mathcal{O}_S$-module fini localement libre dont la formation commute à tout changement de base; et (c) $R^qf_*\mathcal{O}_X = 0$ pour $q \geq 2$.
\end{lemma}

\begin{proof}
La partie (a) résulte de Catégories dérivées des schémas, lemme \ref{perfect-lemma-proper-flat-geom-red-connected}. D'après Catégories dérivées des schémas, lemme \ref{perfect-lemma-proper-flat-h0}, localement sur $S$, nous pouvons écrire $Rf_*\mathcal{O}_X = \mathcal{O}_S \oplus P$, où $P$ est parfait et d'amplitude de Tor contenue dans $[1, \infty)$. Rappelons que la formation de $Rf_*\mathcal{O}_X$ commute à tout changement de base (Catégories dérivées des schémas, lemme \ref{perfect-lemma-flat-proper-perfect-direct-image-general}). Ainsi, pour $s \in S$, on a $$
H^i(P \otimes_{\mathcal{O}_S}^\mathbf{L} \kappa(s)) =
H^i(X_s, \mathcal{O}_{X_s})
\text{ pour }i \geq 1
$$ Ce groupe est nul sauf si $i = 1$, puisque $X_s$ est un schéma noethérien de dimension $1$; voir Cohomologie, proposition \ref{cohomology-proposition-vanishing-Noetherian}. Alors $P = H^1(P)[-1]$ et $H^1(P)$ est fini localement libre, par exemple d'après Compléments d'algèbre, lemme \ref{more-algebra-lemma-lift-perfect-from-residue-field}. Comme tout est compatible au changement de base, nous concluons.
\end{proof}





\section{Structure locale étale des familles nodales}
\label{curves-section-etale-local-nodal}


\noindent
Considérons le morphisme de schémas $$
\Spec(\mathbf{Z}[u, v, a]/(uv - a))
\longrightarrow
\Spec(\mathbf{Z}[a])
$$ Le lemme suivant montre que ce morphisme est un modèle de la structure locale étale d'une famille de courbes nodales.

\begin{lemma}
\label{curves-lemma-etale-local-structure-nodal-family}
Soit $f : X \to S$ un morphisme de schémas. Supposons que $f$ soit à singularités au plus nodales de dimension relative $1$. Soit $x \in X$ un point singulier de la fibre $X_s$. Il existe alors un diagramme commutatif de schémas $$
\xymatrix{
X \ar[d] &
U \ar[rr] \ar[l] \ar[rd] & &
W \ar[r] \ar[ld] &
\Spec(\mathbf{Z}[u, v, a]/(uv - a)) \ar[d] \\
S & &
V \ar[ll] \ar[rr] & & \Spec(\mathbf{Z}[a])
}
$$ dans lequel $X \leftarrow U$, $S \leftarrow V$ et $U \to W$ sont des morphismes étales, et où le carré de droite est cartésien, ainsi qu'un point $u \in U$ s'envoyant sur $x$ dans $X$.
\end{lemma}

\begin{proof}[Première démonstration par approximation d'Artin] Utilisons d'abord l'approximation noethérienne absolue pour nous ramener au cas de schémas de type fini sur $\mathbf{Z}$. La question est locale sur $X$ et $S$. Nous pouvons donc supposer que $X$ et $S$ sont affines. Écrivons alors $S = \Spec(R)$ et écrivons $R$ comme une colimite filtrante $R = \colim R_i$ de $\mathbf{Z}$-algèbres de type fini. D'après Limites, lemme \ref{limits-lemma-descend-finite-presentation}, nous pouvons trouver $i$ et un morphisme $f_i : X_i \to \Spec(R_i)$ dont le changement de base à $S$ est $f$. Après avoir augmenté $i$, nous pouvons supposer que $f_i$ est à singularités au plus nodales de dimension relative $1$; voir le lemme \ref{curves-lemma-descend-nodal-family}. L'image $x_i \in X_i$ de $x$ sera un point singulier de sa fibre, par exemple parce que la formation de $\text{Sing}(f)$ commute au changement de base (Diviseurs, lemme \ref{divisors-lemma-base-change-and-fitting-ideal-omega}). Si nous pouvons démontrer le lemme pour $f_i : X_i \to \Spec(R_i)$ et $x_i$, il en résulte pour $f : X \to S$ par changement de base. Nous sommes ainsi ramenés au cas étudié dans le paragraphe suivant.

\medskip\noindent
Supposons $S$ de type fini sur $\mathbf{Z}$. Soit $s \in S$ l'image de $x$. Rappelons que $\kappa(x)$ est une extension finie séparable de $\kappa(s)$, par exemple parce que $\text{Sing}(f) \to S$ est non ramifié, ou parce que $x$ est un nœud de la fibre $X_s$ et que nous pouvons appliquer le lemme \ref{curves-lemma-nodal}. Soit en outre $\kappa'/\kappa(x)$ l'algèbre séparable de degré $2$ associée à $\mathcal{O}_{X_s, x}$ dans la remarque \ref{curves-remark-quadratic-extension}. D'après Compléments sur les morphismes, lemme \ref{more-morphisms-lemma-realize-prescribed-residue-field-extension-etale}, nous pouvons choisir un voisinage étale $(V, v) \to (S, s)$ tel que l'extension $\kappa(v)/\kappa(s)$ réalise soit l'extension $\kappa(x)/\kappa(s)$ dans le cas $\kappa' \cong \kappa(x) \times \kappa(x)$, soit l'extension $\kappa'/\kappa(s)$ si $\kappa'$ est un corps. Après avoir remplacé $X$ par $X \times_S V$ et $S$ par $V$, nous nous ramenons à la situation décrite au paragraphe suivant.

\medskip\noindent
Supposons $S$ de type fini sur $\mathbf{Z}$ et $x \in X_s$ un nœud déployé; voir la définition \ref{curves-definition-split-node}. D'après le lemme \ref{curves-lemma-formal-local-structure-nodal-family}, il existe un isomorphisme de $\mathcal{O}_{S, s}$-algèbres $$
\mathcal{O}_{X, x}^\wedge \cong
\mathcal{O}_{S, s}^\wedge[[s, t]]/(st - h)
$$ pour un certain $h \in \mathfrak m_s^\wedge \subset \mathcal{O}_{S, s}^\wedge$. Autrement dit, si l'on considère l'homomorphisme $$
\sigma : \mathbf{Z}[a] \longrightarrow \mathcal{O}_{S, s}^\wedge
$$ envoyant $a$ sur $h$, il existe un isomorphisme de $\mathcal{O}_{S, s}$-algèbres $$
\mathcal{O}_{X, x}^\wedge
\longrightarrow
\mathcal{O}_{Y_\sigma, y_\sigma}^\wedge
$$, où $$
Y_\sigma = \Spec(\mathbf{Z}[u, v, a]/(uv - a))
\times_{\Spec(\mathbf{Z}[a]), \sigma} \Spec(\mathcal{O}_{S, s}^\wedge)
$$ et où $y_\sigma$ est le point de $Y_\sigma$ situé au-dessus du point fermé de $\Spec(\mathcal{O}_{S, s}^\wedge)$ dont les coordonnées $u, v$ sont nulles. Puisque $\mathcal{O}_{S, s}$ est un G-anneau d'après Compléments d'algèbre, proposition \ref{more-algebra-proposition-ubiquity-G-ring}, nous pouvons conclure par Compléments sur les morphismes, lemme \ref{more-morphisms-lemma-relative-map-approximation-pre}.
\end{proof}

\begin{proof}[Démonstration sans approximation d'Artin] Nous ne faisons qu'esquisser cette démonstration, due à Mohan Swaminathan et calquée sur un argument de \cite[Proposition X.2.1]{ACG}. Contrairement à la précédente, cette démonstration utilise les équations elles-mêmes.

\medskip\noindent
Exactement comme dans la première démonstration, nous nous ramenons au cas où $S$ est de type fini sur $\mathbf{Z}$ et où $x \in X_s$ est un nœud déployé; voir la définition \ref{curves-definition-split-node}. D'après le lemme \ref{curves-lemma-split-node}, il existe un isomorphisme de $\kappa(s)$-algèbres $$
\mathcal{O}_{X_s, x}^\wedge \cong \kappa(s)[[u, v]]/(uv)
$$ Remarquons que l'espace tangent de $X_s$ en $x$ est de dimension $2$. D'après Variétés, lemme \ref{varieties-lemma-immersion-into-affine}, après avoir rétréci au sens de Zariski, nous pouvons donc supposer qu'il existe une immersion fermée $X \to Y$ de schémas sur $S$, avec $Y \to S$ lisse de dimension relative $2$. Puisque $X$ est syntomique de dimension relative $1$ (par définition), $X$ est un diviseur de Cartier effectif de $Y$ (cela résulte de Diviseurs, lemme \ref{divisors-lemma-immersion-lci-into-smooth-regular-immersion}). Après avoir rétréci $Y$ et $X$, nous pouvons donc supposer que $Y = \Spec(B)$ est affine et qu'il existe un non-diviseur de zéro $g \in B$ tel que $X = \Spec(B/gB)$.

\medskip\noindent
Écrivons $S = \Spec(R)$ et notons $\mathfrak p \subset R$ l'idéal premier correspondant à $s$. Soit $\mathfrak q \subset B$ l'idéal premier correspondant à $x$, considéré comme un point de $Y$. Nous avons des applications $$
B_\mathfrak q / \mathfrak p B_\mathfrak q
\longrightarrow
B_\mathfrak q / g B_\mathfrak q + \mathfrak p B_\mathfrak q
\longrightarrow
\kappa(s)[[u, v]]/(uv)
$$ dont la seconde devient un isomorphisme après complétion. Après avoir remplacé $B$ par une localisation principale convenable, nous pouvons supposer qu'il existe $U, V \in B$ qui s'envoient sur $u$ et $v$ à $(u, v)^2$ près. Alors le complété de $B_\mathfrak q/\mathfrak p B_\mathfrak q$ est $\kappa(s)[[U, V]]$. De plus, après avoir rétréci $Y$ et multiplié $g$ par une unité, nous pouvons supposer que $g$ s'envoie sur $UV$ plus des termes d'ordre supérieur dans $\kappa(s)[[U, V]]$. Après avoir encore rétréci $Y$, nous pouvons supposer en outre que \begin{enumerate} \item $\mathfrak q = \mathfrak pB + (U, V)$ \item $\Omega_{B/R}$ est libre sur $\text{d}U$ et $\text{d}V$. \end{enumerate} Écrivons $\text{d}(g) = g_2 \text{d}U + g_1 \text{d}V$ avec $g_1, g_2 \in B$. Remarquons que l'idéal $J = (g_1, g_2) \subset B$ est indépendant du choix des « coordonnées » $U$ et $V$, pour $g$ fixé. En réduisant modulo $\mathfrak p B_\mathfrak q + (U, V)^3B_\mathfrak q$, on voit que $g_1$ et $g_2$ sont congrus à $U$ et $V$ modulo $\mathfrak p B_\mathfrak q + (U, V)^2B_\mathfrak q$. Nous pouvons donc reprendre le raisonnement ci-dessus en remplaçant $U$ par $g_1$ et $V$ par $g_2$. Nous obtenons alors $J = (U, V)$, ce qui implique $g_1, g_2 \in (U, V) = J$ (bien entendu, ces éléments diffèrent des $g_1$ et $g_2$ correspondant à notre choix précédent de coordonnées).

\medskip\noindent
Remarquons que $\Spec(B/J) \to S = \Spec(R)$ est étale en $x$. Après avoir remplacé $S$ par un voisinage étale et rétréci encore $Y$, nous pouvons donc supposer que $R \to B \to B/J$ est un isomorphisme. Notons $r \in R$ l'élément dont l'image dans $B/J$ est celle de $-g$. Nous pouvons alors écrire $$
g + r \equiv a U + b V \bmod (U, V)^2
$$ pour un certain $a, b \in R$. En dérivant et en utilisant $\text{d}g \in J\Omega_{B/R}$, nous obtenons $a = b = 0$. Nous pouvons donc écrire $$
g + r = \alpha U^2 + \beta UV + \gamma V^2
$$ pour un certain $\alpha, \beta, \gamma \in B$. Alors $\beta$ s'envoie sur $1$ dans $\kappa(x)$, tandis que $\alpha$ et $\gamma$ s'envoient sur $0$ dans $\kappa(x)$. Nous pouvons donc supposer que $\beta$ et $\alpha + \beta + \gamma$ sont inversibles dans $B$. Considérons le revêtement $Y'$ de $Y$ donné par $$
Y' = \Spec(B[t]/(t(1 - t) - \alpha \gamma / \beta^2))
$$, et soit $x'$ l'unique point au-dessus de $x$ tel que $t = 1$. Remarquons que $Y' \to Y$ est étale en $x'$. Considérons les éléments $$
T_1 = (\alpha + t \beta)U + ((1 - t)\beta + \gamma)V
\quad\text{et}\quad
T_2 =
\frac{(\alpha + (1 - t) \beta)U + (t\beta + \gamma)V}{\alpha + \beta + \gamma}
$$ de $\Gamma(Y', \mathcal{O}_{Y'})$. On voit alors que $$
g = T_1 T_2 - r
$$ Puisque $T_1$ et $T_2$ sont également des coordonnées au sens précédent, le morphisme $Y' \to \mathbf{A}^2_S$ donné par $T_1$ et $T_2$ est étale dans un voisinage de $x'$, et l'image réciproque $X' \subset Y'$ de $X$ est l'image réciproque du sous-schéma fermé de $\mathbf{A}^2_S = \Spec(R[x_1, x_2])$ défini par l'équation $x_1x_2 - r = 0$, quitte à rétrécir encore $Y'$. Nous laissons au lecteur le soin d'assembler ces éléments pour obtenir le diagramme voulu.
\end{proof}




\section{Autres résultats d'annulation}
\label{curves-section-more-vanishing}

\noindent
Suite de la section \ref{curves-section-vanishing}.

\begin{lemma}
\label{curves-lemma-h1-nonzero-degree-leq-2g-2}
Dans la situation \ref{curves-situation-Cohen-Macaulay-curve}, supposons $X$ intègre et de genre $g$. Soit $\mathcal{L}$ un $\mathcal{O}_X$-module inversible. Soit $Z \subset X$ un sous-schéma fermé de dimension $0$, de faisceau d'idéaux $\mathcal{I} \subset \mathcal{O}_X$. Si $H^1(X, \mathcal{I}\mathcal{L})$ est non nul, alors $$
\deg(\mathcal{L}) \leq 2g - 2 + \deg(Z)
$$, avec inégalité stricte sauf si $\mathcal{I}\mathcal{L} \cong \omega_X$.
\end{lemma}

\begin{proof}
Toute courbe, par exemple $X$, est de Cohen--Macaulay. Si $H^1(X, \mathcal{I}\mathcal{L})$ est non nul, il existe une application non nulle $\mathcal{I}\mathcal{L} \to \omega_X$; voir le lemme \ref{curves-lemma-duality-dim-1-CM}. Puisque $\mathcal{I}\mathcal{L}$ est sans torsion, cette application est injective. Comme un corps est de Gorenstein et que $X$ est réduit, le lieu de Gorenstein $U \subset X$ de $X$ est non vide; voir Dualité pour les schémas, lemme \ref{duality-lemma-gorenstein}. Ce lemme nous dit aussi que $\omega_X|_U$ est inversible. Nous obtenons ainsi une suite exacte courte $$
0 \to \mathcal{I}\mathcal{L} \to \omega_X \to \mathcal{Q} \to 0
$$ où le support de $\mathcal{Q}$ est de dimension zéro. Par conséquent, \begin{align*}
0 & \leq \dim \Gamma(X, \mathcal{Q})\\
& =
\chi(\mathcal{Q}) \\
& =
\chi(\omega_X) - \chi(\mathcal{I}\mathcal{L}) \\
& =
\chi(\omega_X) - \deg(\mathcal{L}) - \chi(\mathcal{I}) \\
& =
2g - 2 - \deg(\mathcal{L}) + \deg(Z)
\end{align*} d'après les lemmes \ref{curves-lemma-euler} et \ref{curves-lemma-rr}, d'après (\ref{curves-equation-genus}), et d'après Variétés, lemmes \ref{varieties-lemma-chi-tensor-finite} et \ref{varieties-lemma-degree-on-proper-curve}. Nous avons aussi utilisé $\deg(Z) = \dim_k \Gamma(Z, \mathcal{O}_Z) = \chi(\mathcal{O}_Z)$ et la suite exacte courte $0 \to \mathcal{I} \to \mathcal{O}_X \to \mathcal{O}_Z \to 0$. Le lemme en résulte.
\end{proof}

\begin{lemma}
\label{curves-lemma-degree-more-than-2g-2}
\begin{reference}
\cite[Lemma 2]{Jongmin}
\end{reference}
Dans la situation \ref{curves-situation-Cohen-Macaulay-curve}, supposons $X$ intègre et de genre $g$. Soit $\mathcal{L}$ un $\mathcal{O}_X$-module inversible. Soit $Z \subset X$ un sous-schéma fermé de dimension $0$, de faisceau d'idéaux $\mathcal{I} \subset \mathcal{O}_X$. Si $\deg(\mathcal{L}) > 2g - 2 + \deg(Z)$, alors $H^1(X, \mathcal{I}\mathcal{L}) = 0$ et l'une des possibilités suivantes se présente :
\begin{enumerate}
\item $H^0(X, \mathcal{I}\mathcal{L}) \not = 0$;
\item $g = 0$ et $\deg(\mathcal{L}) = \deg(Z) - 1$.
\end{enumerate}
Dans le cas (2), si $Z = \emptyset$, alors $X \cong \mathbf{P}^1_k$ et $\mathcal{L}$ correspond à $\mathcal{O}_{\mathbf{P}^1}(-1)$.
\end{lemma}

\begin{proof}
L'annulation de $H^1(X, \mathcal{I}\mathcal{L})$ résulte du lemme \ref{curves-lemma-h1-nonzero-degree-leq-2g-2}. Si $H^0(X, \mathcal{I}\mathcal{L}) = 0$, alors $\chi(\mathcal{I}\mathcal{L}) = 0$. La suite exacte courte $0 \to \mathcal{I}\mathcal{L} \to \mathcal{L} \to \mathcal{O}_Z \to 0$ donne $\deg(\mathcal{L}) = g - 1 + \deg(Z)$. Ainsi $g - 1 + \deg(Z) > 2g - 2 + \deg(Z)$, ce qui implique $g = 0$; le cas (2) est donc réalisé. Si $Z = \emptyset$ dans le cas (2), alors $\mathcal{L}^{-1}$ est un faisceau inversible de degré $1$. Il existe par conséquent un isomorphisme $X \to \mathbf{P}^1_k$, et $\mathcal{L}^{-1}$ est l'image réciproque de $\mathcal{O}_{\mathbf{P}^1}(1)$ d'après le lemme \ref{curves-lemma-genus-zero-positive-degree}.
\end{proof}

\begin{lemma}
\label{curves-lemma-degree-more-than-2g}
\begin{reference}
\cite[Lemma 3]{Jongmin}
\end{reference}
Dans la situation \ref{curves-situation-Cohen-Macaulay-curve}, supposons $X$ intègre et de genre $g$. Soit $\mathcal{L}$ un $\mathcal{O}_X$-module inversible. Si $\deg(\mathcal{L}) \geq 2g$, alors $\mathcal{L}$ est engendré par ses sections globales.
\end{lemma}

\begin{proof}
Soit $Z \subset X$ le sous-schéma fermé découpé par les sections globales de $\mathcal{L}$. D'après le lemme \ref{curves-lemma-degree-more-than-2g-2}, on voit que $Z \not = X$. Soit $\mathcal{I} \subset \mathcal{O}_X$ le faisceau d'idéaux définissant $Z$. Considérons la suite exacte courte $$
0 \to \mathcal{I}\mathcal{L}
\to \mathcal{L} \to \mathcal{O}_Z \to 0
$$ Si $Z \not = \emptyset$, alors $H^1(X, \mathcal{I}\mathcal{L})$ est non nul, comme le montre la suite exacte longue de cohomologie. D'après le lemme \ref{curves-lemma-duality-dim-1-CM}, cela donne une application non nulle, donc injective, $$
\mathcal{I}\mathcal{L}
\longrightarrow
\omega_X
$$ En particulier, nous obtenons une application injective $H^0(X, \mathcal{L}) = H^0(X, \mathcal{I}\mathcal{L})
\to H^0(X, \omega_X)$. C'est impossible, puisque $$
\dim_k H^0(X, \mathcal{L}) = \dim_k H^1(X, \mathcal{L}) +
\deg(\mathcal{L}) + 1 - g \geq g + 1
$$ et $\dim H^0(X, \omega_X) = g$ d'après (\ref{curves-equation-genus}).
\end{proof}

\begin{lemma}
\label{curves-lemma-degree-more-than-2g-1-and-Z}
Dans la situation \ref{curves-situation-Cohen-Macaulay-curve}, supposons $X$ intègre et de genre $g$. Soit $\mathcal{L}$ un $\mathcal{O}_X$-module inversible. Soit $Z \subset X$ un sous-schéma fermé non vide de dimension $0$. Si $\deg(\mathcal{L}) \geq 2g - 1 + \deg(Z)$, alors $\mathcal{L}$ est engendré par ses sections globales et $H^0(X, \mathcal{L}) \to H^0(X, \mathcal{L}|_Z)$ est surjective.
\end{lemma}

\begin{proof}
L'engendrement par les sections globales résulte du lemme \ref{curves-lemma-degree-more-than-2g}. Si $\mathcal{I} \subset \mathcal{O}_X$ est le faisceau d'idéaux de $Z$, alors $H^1(X, \mathcal{I}\mathcal{L}) = 0$ d'après le lemme \ref{curves-lemma-h1-nonzero-degree-leq-2g-2}. D'où la surjectivité.
\end{proof}

\begin{lemma}
\label{curves-lemma-degree-at-least-2g+1}
Dans la situation \ref{curves-situation-Cohen-Macaulay-curve}, supposons $X$ géométriquement intègre sur $k$ et de genre $g$. Soit $\mathcal{L}$ un $\mathcal{O}_X$-module inversible. Si $\deg(\mathcal{L}) \geq 2g + 1$, alors $\mathcal{L}$ est très ample.
\end{lemma}

\begin{proof}
D'après le lemme \ref{curves-lemma-degree-more-than-2g}, $\mathcal{L}$ est engendré par ses sections globales et détermine donc un morphisme $f : X \to \mathbf{P}^n_k$, où $n = h^0(X,\mathcal{L}) - 1$. Dire que $\mathcal{L}$ est très ample revient à montrer que $f$ est une immersion fermée. Il suffit de vérifier que le changement de base de $f$ à une clôture algébrique $\overline{k}$ de $k$ est une immersion fermée (Descente, lemme \ref{descent-lemma-descending-property-closed-immersion}). Nous pouvons donc supposer $k$ algébriquement clos; par hypothèse, $X$ reste intègre. Le lemme \ref{curves-lemma-degree-more-than-2g-1-and-Z} montre que, pour tout sous-schéma fermé de dimension $0$, noté $Z\subset X$, et de degré 2, l'application de restriction $H^0(X, \mathcal{L}) \to H^0(X, \mathcal{L}|_Z)$ est surjective. D'après Variétés, lemme \ref{varieties-lemma-variant-separate-points-tangent-vectors}, $\mathcal{L}$ est très ample.
\end{proof}


\begin{lemma}
\label{curves-lemma-vanishing-on-gorenstein}
\begin{reference}
Version faible de \cite[Lemma 4]{Jongmin}
\end{reference}
Soit $k$ un corps. Soit $X$ un schéma propre sur $k$, réduit, connexe et de dimension $1$. Soit $\mathcal{L}$ un $\mathcal{O}_X$-module inversible. Soit $Z \subset X$ un sous-schéma fermé de dimension $0$, de faisceau d'idéaux $\mathcal{I} \subset \mathcal{O}_X$. Si $H^1(X, \mathcal{I}\mathcal{L}) \not = 0$, il existe un sous-schéma fermé réduit connexe $Y \subset X$ de dimension $1$ tel que
$$
\deg(\mathcal{L}|_Y) \leq -2\chi(Y, \mathcal{O}_Y) + \deg(Z \cap Y)
$$, où $Z \cap Y$ est l'intersection schématique.
\end{lemma}

\begin{proof}
Si $H^1(X, \mathcal{I}\mathcal{L})$ est non nul, il existe une application non nulle $\varphi : \mathcal{I}\mathcal{L} \to \omega_X$; voir le lemme \ref{curves-lemma-duality-dim-1-CM}. Soit $Y \subset X$ la réunion des composantes irréductibles $C$ de $X$ telles que $\varphi$ soit non nul au point générique de $C$. Alors $Y$ est un sous-schéma fermé réduit. Soit $\mathcal{J} \subset \mathcal{O}_X$ le faisceau d'idéaux de $Y$. Puisque $\mathcal{J}\mathcal{I}\mathcal{L}$ n'a pas de point associé immergé, en tant que sous-module de $\mathcal{L}$, et que $\varphi$ est nul aux points génériques du support de $\mathcal{J}$, par le choix de $Y$ et puisque $X$ est réduit, on voit que $\varphi$ se factorise sous la forme $$
\mathcal{I}\mathcal{L} \to
\mathcal{I}\mathcal{L}/\mathcal{J}\mathcal{I}\mathcal{L} \to \omega_X
$$ Nous pouvons considérer $\mathcal{I}\mathcal{L}/\mathcal{J}\mathcal{I}\mathcal{L}$ comme l'image directe d'un faisceau cohérent sur $Y$, que nous noterons encore de la même manière par abus de notation. Puisque $\omega_Y = \SheafHom(\mathcal{O}_Y, \omega_X)$ d'après le lemme \ref{curves-lemma-closed-immersion-dim-1-CM}, nous obtenons une application $$
\mathcal{I}\mathcal{L}/
\mathcal{J}\mathcal{I}\mathcal{L} 
\to \omega_Y
$$ de $\mathcal{O}_Y$-modules, injective aux points génériques de $Y$. Soit $\mathcal{I}' \subset \mathcal{O}_Y$ le faisceau d'idéaux de $Z \cap Y$. Il existe une application $\mathcal{I}\mathcal{L}/\mathcal{J}\mathcal{I}\mathcal{L} \to
\mathcal{I}'\mathcal{L}|_Y$ dont le noyau est supporté en des points fermés. Comme $\omega_Y$ est un module de Cohen--Macaulay, l'application précédente se factorise par une application injective $\mathcal{I}'\mathcal{L}|_Y \to
\omega_Y$. Nous obtenons ainsi une suite exacte $$
0 \to \mathcal{I}'\mathcal{L}|_Y \to \omega_Y \to \mathcal{Q} \to 0
$$ de faisceaux cohérents sur $Y$, où $\mathcal{Q}$ est supporté en dimension $0$; on utilise ici le fait que $\omega_Y$ est un module inversible aux points génériques de $Y$. Nous concluons que $$
0 \leq \dim \Gamma(Y, \mathcal{Q}) =
\chi(\mathcal{Q}) = \chi(\omega_Y) - \chi(\mathcal{I}'\mathcal{L}) =
-2\chi(\mathcal{O}_Y) - \deg(\mathcal{L}|_Y) + \deg(Z \cap Y)
$$ d'après le lemme \ref{curves-lemma-euler} et Variétés, lemme \ref{varieties-lemma-chi-tensor-finite}. Si $Y$ est connexe, cela démontre le lemme. Sinon, nous répétons la dernière partie de l'argument pour l'une des composantes connexes de $Y$.
\end{proof}

\begin{lemma}
\label{curves-lemma-global-generation}
Soit $k$ un corps. Soit $X$ un schéma propre sur $k$, réduit, connexe et de dimension $1$. Soit $\mathcal{L}$ un $\mathcal{O}_X$-module inversible. Supposons que, pour tout sous-schéma fermé réduit connexe $Y \subset X$ de dimension $1$, on ait $$
\deg(\mathcal{L}|_Y) \geq 2\dim_k H^1(Y, \mathcal{O}_Y)
$$ Alors $\mathcal{L}$ est engendré par ses sections globales.
\end{lemma}

\begin{proof}
Raisonnons par récurrence sur le nombre de composantes irréductibles de $X$. Si $X$ est irréductible, le lemme résulte du lemme \ref{curves-lemma-degree-more-than-2g} appliqué à $X$ considéré comme un schéma sur le corps $k' = H^0(X, \mathcal{O}_X)$. Supposons $X$ non irréductible. Avant de poursuivre, si $k$ est fini, remplaçons $k$ par une extension purement transcendante $K$. Cela est permis par Variétés, lemmes \ref{varieties-lemma-globally-generated-base-change}, \ref{varieties-lemma-degree-base-change}, \ref{varieties-lemma-geometrically-reduced-any-base-change} et \ref{varieties-lemma-bijection-irreducible-components}, par Cohomologie des schémas, lemme \ref{coherent-lemma-flat-base-change-cohomology}, par le lemme \ref{curves-lemma-sanity-check-duality} et par le fait élémentaire que $K$ est géométriquement intègre sur $k$.

\medskip\noindent
Supposons, pour aboutir à une contradiction, que $\mathcal{L}$ ne soit pas engendré par ses sections globales. Nous pouvons alors choisir un faisceau cohérent d'idéaux $\mathcal{I} \subset \mathcal{O}_X$ tel que $H^0(X, \mathcal{I}\mathcal{L}) = H^0(X, \mathcal{L})$ et tel que $\mathcal{O}_X/\mathcal{I}$ soit non nul et de support de dimension $0$. On peut, par exemple, prendre pour $\mathcal{I}$ le faisceau d'idéaux d'un point fermé quelconque du lieu d'annulation commun des sections globales de $\mathcal{L}$. Considérons la suite exacte courte $$
0 \to \mathcal{I}\mathcal{L} \to \mathcal{L} \to 
\mathcal{L}/\mathcal{I}\mathcal{L} \to 0
$$ Puisque le support de $\mathcal{L}/\mathcal{I}\mathcal{L}$ est de dimension $0$, $\mathcal{L}/\mathcal{I}\mathcal{L}$ est engendré par ses sections globales (Variétés, lemme \ref{varieties-lemma-chi-tensor-finite}). La suite exacte courte et le fait que $H^0(X, \mathcal{I}\mathcal{L}) = H^0(X, \mathcal{L})$ donnent une injection $H^0(X, \mathcal{L}/\mathcal{I}\mathcal{L}) \to H^1(X, \mathcal{I}\mathcal{L})$.

\medskip\noindent
Rappelons que le $k$-espace vectoriel $H^1(X, \mathcal{I}\mathcal{L})$ est dual de $\Hom(\mathcal{I}\mathcal{L}, \omega_X)$. Choisissons $\varphi : \mathcal{I}\mathcal{L} \to \omega_X$. D'après le lemme \ref{curves-lemma-vanishing-on-gorenstein}, on a $H^1(X, \mathcal{L}) = 0$. Par conséquent, $$
\dim_k H^0(X, \mathcal{I}\mathcal{L}) = \dim_k H^0(X, \mathcal{L}) =
\deg(\mathcal{L}) + \chi(\mathcal{O}_X) > \dim_k H^1(X, \mathcal{O}_X) =
\dim_k H^0(X, \omega_X)
$$ Nous en concluons que $\varphi$ n'est pas injective sur les sections globales; en particulier, $\varphi$ n'est pas injective. Pour chaque point générique $\eta \in X$ d'une composante irréductible de $X$, notons $V_\eta \subset \Hom(\mathcal{I}\mathcal{L}, \omega_X)$ le sous-$k$-espace vectoriel formé des $\varphi$ qui sont nuls en $\eta$. Puisque tout point associé de $\mathcal{I}\mathcal{L}$ est un point générique de $X$, ce qui précède montre que $\Hom(\mathcal{I}\mathcal{L}, \omega_X) = \bigcup V_\eta$. Comme $X$ possède un nombre fini de points génériques et que $k$ est infini, nous concluons $\Hom(\mathcal{I}\mathcal{L}, \omega_X) = V_\eta$ pour un certain $\eta$. Soit $\eta \in C \subset X$ la composante irréductible correspondante. Soit $Y \subset X$ la réunion des autres composantes irréductibles de $X$. Alors $Y$ est un sous-schéma fermé réduit non vide, distinct de $X$. Soit $\mathcal{J} \subset \mathcal{O}_X$ le faisceau d'idéaux de $Y$. Gardons à l'esprit que le support de $\mathcal{J}$ est $C$.

\medskip\noindent
Soit $\varphi : \mathcal{I}\mathcal{L} \to \omega_X$ quelconque. Puisque $\mathcal{J}\mathcal{I}\mathcal{L}$ n'a pas de point associé immergé, en tant que sous-module de $\mathcal{L}$, et que $\varphi$ est nul au point générique $\eta$ du support de $\mathcal{J}$, on voit que $\varphi$ se factorise sous la forme $$
\mathcal{I}\mathcal{L} \to
\mathcal{I}\mathcal{L}/\mathcal{J}\mathcal{I}\mathcal{L} \to \omega_X
$$ Nous pouvons considérer $\mathcal{I}\mathcal{L}/\mathcal{J}\mathcal{I}\mathcal{L}$ comme l'image directe d'un faisceau cohérent sur $Y$, que nous noterons encore de la même manière par abus de notation. Puisque $\omega_Y = \SheafHom(\mathcal{O}_Y, \omega_X)$ d'après le lemme \ref{curves-lemma-closed-immersion-dim-1-CM}, nous obtenons une factorisation $$
\mathcal{I}\mathcal{L} \to
\mathcal{I}\mathcal{L}/
\mathcal{J}\mathcal{I}\mathcal{L} 
\xrightarrow{\varphi'} \omega_Y \to \omega_X
$$ de $\varphi$. Soit $\mathcal{I}' \subset \mathcal{O}_Y$ l'image de $\mathcal{I} \subset \mathcal{O}_X$. Il existe une application surjective $\mathcal{I}\mathcal{L}/\mathcal{J}\mathcal{I}\mathcal{L} \to
\mathcal{I}'\mathcal{L}|_Y$ dont le noyau est supporté en des points fermés. Comme $\omega_Y$ est un module de Cohen--Macaulay sur $Y$, l'application $\varphi'$ se factorise par une application $\varphi'' : \mathcal{I}'\mathcal{L}|_Y \to \omega_Y$. Nous avons donc des diagrammes commutatifs $$
\vcenter{
\xymatrix{
0 \ar[r] &
\mathcal{I}\mathcal{L} \ar[r] \ar[d] &
\mathcal{L} \ar[r] \ar[d] &
\mathcal{L}/\mathcal{I}\mathcal{L} \ar[r] \ar[d] &
0 \\
0 \ar[r] &
\mathcal{I}'\mathcal{L}|_Y \ar[r] &
\mathcal{L}|_Y \ar[r] &
\mathcal{L}|_Y/\mathcal{I}'\mathcal{L}|_Y \ar[r] &
0
}
}
\quad\text{et}\quad
\vcenter{
\xymatrix{
\mathcal{I}\mathcal{L} \ar[r]_\varphi \ar[d] & \omega_X \\
\mathcal{I}'\mathcal{L}|_Y \ar[r]^{\varphi''} & \omega_Y \ar[u]
}
}
$$ Nous pouvons maintenant achever la démonstration comme suit. Puisque, pour tout $\varphi$, nous avons $\varphi''$ et puisque $\omega_X \in \textit{Coh}(\mathcal{O}_X)$ représente le foncteur $\mathcal{F} \mapsto \Hom_k(H^1(X, \mathcal{F}), k)$, on voit que $H^1(X, \mathcal{I}\mathcal{L}) \to H^1(Y, \mathcal{I}'\mathcal{L}|_Y)$ est injective. Comme l'application de bord $H^0(X, \mathcal{L}/\mathcal{I}\mathcal{L}) \to H^1(X, \mathcal{I}\mathcal{L})$ est injective, le composé $$
H^0(X, \mathcal{L}/\mathcal{I}\mathcal{L}) \to
H^0(X, \mathcal{L}|_Y/\mathcal{I}'\mathcal{L}|_Y) \to
H^1(X, \mathcal{I}'\mathcal{L}|_Y)
$$ l'est aussi. Puisque $\mathcal{L}/\mathcal{I}\mathcal{L} \to
\mathcal{L}|_Y/\mathcal{I}'\mathcal{L}|_Y$ est une application surjective de modules cohérents dont les supports sont de dimension $0$, la première application $H^0(X, \mathcal{L}/\mathcal{I}\mathcal{L}) \to
H^0(X, \mathcal{L}|_Y/\mathcal{I}'\mathcal{L}|_Y)$ est surjective, donc bijective. Mais, par récurrence, $\mathcal{L}|_Y$ est engendré par ses sections globales, y compris si $Y$ est non connexe, et l'application de bord $$
H^0(X, \mathcal{L}|_Y/\mathcal{I}'\mathcal{L}|_Y) \to
H^1(X, \mathcal{I}'\mathcal{L}|_Y)
$$ ne peut donc pas être injective. Cette contradiction achève la démonstration.
\end{proof}





\section{Contraction des queues rationnelles}
\label{curves-section-contracting-rational-tails}

\noindent
Dans cette section, nous étudions le cas le plus simple de contraction d'un schéma visant à améliorer les propriétés de positivité de son faisceau canonique.

\begin{example}[Contraction d'une queue rationnelle] \label{curves-example-rational-tail} Soit $k$ un corps. Soit $X$ un schéma propre sur $k$, de dimension $1$ et tel que $H^0(X, \mathcal{O}_X) = k$. Supposons que les singularités de $X$ soient au plus nodales. Une {\it queue rationnelle} est une composante irréductible $C \subset X$, considérée comme un sous-schéma fermé intègre, qui possède les propriétés suivantes : \begin{enumerate} \item $X' \not = \emptyset$, où $X' \subset X$ est l'adhérence schématique de $X \setminus C$; \item l'intersection schématique $C \cap X'$ est un unique point réduit $x$; \item $H^0(C, \mathcal{O}_C)$ s'envoie isomorphiquement sur le corps résiduel de $x$; \item $C$ est de genre zéro. \end{enumerate} Comme au moins deux composantes irréductibles de $X$ passent par $x$, nous en concluons que $x$ est un nœud. Posons $k' = H^0(C, \mathcal{O}_C) = \kappa(x)$. Alors $k'/k$ est une extension finie séparable de corps (lemme \ref{curves-lemma-nodal}). Il existe un morphisme canonique $$
c : X \longrightarrow X'
$$ qui induit l'identité sur $X'$ et envoie $C$ sur $x \in X'$ par le morphisme canonique $C \to \Spec(k') = x$. Cela résulte de Morphismes, lemme \ref{morphisms-lemma-scheme-theoretic-union}, puisque $X$ est la réunion schématique de $C$ et $X'$, car $X$ est réduit. De plus, nous affirmons que $$
c_*\mathcal{O}_X = \mathcal{O}_{X'}
\quad\text{et}\quad
R^1c_*\mathcal{O}_X = 0
$$ Pour le voir, notons $i_C : C \to X$, $i_{X'} : X' \to X$ et $i_x : x \to X$ les immersions et utilisons la suite exacte $$
0 \to \mathcal{O}_X \to
i_{C, *}\mathcal{O}_C \oplus i_{X', *}\mathcal{O}_{X'} \to
i_{x, *}\kappa(x) \to 0
$$ de Morphismes, lemme \ref{morphisms-lemma-scheme-theoretic-union}. La suite exacte longue des images directes supérieures montre qu'il suffit d'établir $H^0(C, \mathcal{O}_C) = k'$ et $H^1(C, \mathcal{O}_C) = 0$, ce qui résulte des hypothèses. Remarquons que $X'$ est lui aussi un schéma propre sur $k$, de dimension $1$, dont les singularités sont au plus nodales (lemme \ref{curves-lemma-closed-subscheme-nodal-curve}), que $H^0(X', \mathcal{O}_{X'}) = k$, et que $X'$ a le même genre que $X$. On dit que $c : X \to X'$ est la {\it contraction d'une queue rationnelle}.
\end{example}

\begin{lemma}
\label{curves-lemma-rational-tail-negative}
Soit $k$ un corps. Soit $X$ un schéma propre sur $k$, de dimension $1$ et tel que $H^0(X, \mathcal{O}_X) = k$. Supposons que les singularités de $X$ soient au plus nodales. Soit $C \subset X$ une queue rationnelle (exemple \ref{curves-example-rational-tail}). Alors $\deg(\omega_X|_C) < 0$.
\end{lemma}

\begin{proof}
Soit $X' \subset X$ comme dans l'exemple. Nous avons une suite exacte courte $$
0 \to \omega_C \to \omega_X|_C \to \mathcal{O}_{C \cap X'} \to 0
$$ Voir les lemmes \ref{curves-lemma-closed-subscheme-reduced-gorenstein}, \ref{curves-lemma-facts-about-nodal-curves} et \ref{curves-lemma-closed-subscheme-nodal-curve}. Avec $k'$ comme dans l'exemple, on voit que $\deg(\omega_C) = -2[k' : k]$ puisque $C \cong \mathbf{P}^1_{k'}$, d'après la proposition \ref{curves-proposition-projective-line}, et que $\deg(C \cap X') = [k' : k]$. Ainsi $\deg(\omega_X|_C) = -[k' : k]$, qui est négatif.
\end{proof}

\begin{lemma}
\label{curves-lemma-rational-tail-field-extension}
Soit $k$ un corps. Soit $X$ un schéma propre sur $k$, de dimension $1$ et tel que $H^0(X, \mathcal{O}_X) = k$. Supposons que les singularités de $X$ soient au plus nodales. Soit $C \subset X$ une queue rationnelle (exemple \ref{curves-example-rational-tail}). Pour toute extension de corps $K/k$, le changement de base $C_K \subset X_K$ est une réunion disjointe finie de queues rationnelles.
\end{lemma}

\begin{proof}
Soient $x \in C$ et $k' = \kappa(x)$ comme dans l'exemple. Remarquons que $C \cong \mathbf{P}^1_{k'}$ d'après la proposition \ref{curves-proposition-projective-line}. Puisque $k'/k$ est fini séparable, $k' \otimes_k K = K'_1 \times \ldots \times K'_n$ est un produit fini d'extensions finies séparables $K'_i/K$. Posons $C_i = \mathbf{P}^1_{K'_i}$ et notons $x_i \in C_i$ l'image réciproque de $x$. Alors $C_K = \coprod C_i$ et $X'_K \cap C_i = x_i$, comme voulu.
\end{proof}

\begin{lemma}
\label{curves-lemma-no-rational-tail}
Soit $k$ un corps. Soit $X$ un schéma propre sur $k$, de dimension $1$ et tel que $H^0(X, \mathcal{O}_X) = k$. Supposons que les singularités de $X$ soient au plus nodales. Si $X$ ne possède pas de queue rationnelle (exemple \ref{curves-example-rational-tail}), alors, pour tout sous-schéma fermé réduit connexe $Y \subset X$, $Y \not = X$, de dimension $1$, on a $\deg(\omega_X|_Y) \geq \dim_k H^1(Y, \mathcal{O}_Y)$.
\end{lemma}

\begin{proof}
Soit $Y \subset X$ comme dans l'énoncé. Alors $k' = H^0(Y, \mathcal{O}_Y)$ est un corps et une extension finie de $k$, et $[k' : k]$ divise tous les invariants numériques ci-dessous associés à $Y$ et aux faisceaux cohérents sur $Y$; voir Variétés, lemme \ref{varieties-lemma-divisible}. Soit $Z \subset X$ comme dans le lemme \ref{curves-lemma-closed-subscheme-reduced-gorenstein}. Nous utiliserons sans autre mention les résultats de ce lemme et des lemmes \ref{curves-lemma-facts-about-nodal-curves} et \ref{curves-lemma-closed-subscheme-nodal-curve}. Nous obtenons une suite exacte courte $$
0 \to \omega_Y \to \omega_X|_Y \to \mathcal{O}_{Y \cap Z} \to 0
$$ Voir le lemme \ref{curves-lemma-closed-subscheme-reduced-gorenstein}. Nous en concluons que $$
\deg(\omega_X|_Y) = \deg(Y \cap Z) + \deg(\omega_Y) =
\deg(Y \cap Z) - 2\chi(Y, \mathcal{O}_Y)
$$ Par conséquent, si le lemme est faux, alors $$
2[k' : k] > \deg(Y \cap Z) + \dim_k H^1(Y, \mathcal{O}_Y)
$$ Puisque $Y \cap Z$ est non vide, et compte tenu de la divisibilité mentionnée ci-dessus, cela n'est possible que si $Y \cap Z$ est un unique point $k'$-rationnel du lieu lisse de $Y$ et si $H^1(Y, \mathcal{O}_Y) = 0$. Si $Y$ est irréductible, cela implique que $Y$ est une queue rationnelle. Si $Y$ est réductible, puisque $\deg(\omega_X|_Y) = -[k' : k]$, il existe une composante irréductible $C$ de $Y$ telle que $\deg(\omega_X|_C) < 0$; voir Variétés, lemme \ref{varieties-lemma-degree-in-terms-of-components}. L'analyse précédente appliquée à $C$ montre alors que $C$ est une queue rationnelle.
\end{proof}

\begin{lemma}
\label{curves-lemma-no-rational-tail-semiample-genus-geq-2}
Soit $k$ un corps. Soit $X$ un schéma propre sur $k$, de dimension $1$ et tel que $H^0(X, \mathcal{O}_X) = k$. Supposons que les singularités de $X$ soient au plus nodales et que $X$ ne possède pas de queue rationnelle (exemple \ref{curves-example-rational-tail}). Si \begin{enumerate} \item le genre de $X$ est $0$, alors $X$ est isomorphe à une conique plane irréductible et $\omega_X^{\otimes -1}$ est très ample; \item le genre de $X$ est $1$, alors $\omega_X \cong \mathcal{O}_X$; \item le genre de $X$ est $\geq 2$, alors $\omega_X^{\otimes m}$ est engendré par ses sections globales pour $m \geq 2$.
\end{enumerate}
\end{lemma}

\begin{proof}
D'après le lemme \ref{curves-lemma-facts-about-nodal-curves}, $X$ est de Gorenstein, c'est-à-dire que $\omega_X$ est un $\mathcal{O}_X$-module inversible.

\medskip\noindent
Si le genre de $X$ est nul, alors $\deg(\omega_X) < 0$. Par conséquent, si $X$ possède plus d'une composante irréductible, nous obtenons une contradiction avec le lemme \ref{curves-lemma-no-rational-tail}. Dans le cas irréductible, $X$ est isomorphe à une conique plane irréductible et $\omega_X^{\otimes -1}$ est très ample d'après le lemme \ref{curves-lemma-genus-zero}.

\medskip\noindent
Si le genre de $X$ est $1$, alors $\omega_X$ possède une section globale et $\deg(\omega_X|_C) = 0$ pour toute composante irréductible. En effet, $\deg(\omega_X|_C) \geq 0$ pour toute composante irréductible $C$ d'après le lemme \ref{curves-lemma-no-rational-tail}, la somme de ces nombres est $0$ d'après le lemme \ref{curves-lemma-genus-gorenstein}, et nous pouvons appliquer Variétés, lemme \ref{varieties-lemma-degree-in-terms-of-components}. Alors $\omega_X \cong \mathcal{O}_X$ d'après Variétés, lemme \ref{varieties-lemma-no-sections-dual-nef}.

\medskip\noindent
Supposons que le genre $g$ de $X$ soit supérieur ou égal à $2$. Si $X$ est irréductible, le résultat découle du lemme \ref{curves-lemma-degree-more-than-2g}. Supposons $X$ réductible. D'après le lemme \ref{curves-lemma-no-rational-tail}, les inégalités du lemme \ref{curves-lemma-global-generation} sont satisfaites pour tout $Y \subset X$ comme dans l'énoncé, sauf pour $Y = X$. En examinant la démonstration du lemme \ref{curves-lemma-global-generation}, on voit que, dans le cas réductible, les seules inégalités utilisées pour $Y = X$ sont $$
\deg(\omega_X^{\otimes m}) > -2 \chi(\mathcal{O}_X)
\quad\text{et}\quad
\deg(\omega_X^{\otimes m}) + \chi(\mathcal{O}_X) > \dim_k H^1(X, \mathcal{O}_X)
$$ Comme elles sont toutes deux satisfaites sous les hypothèses $g \geq 2$ et $m \geq 2$, nous concluons.
\end{proof}

\begin{lemma}
\label{curves-lemma-contracting-rational-tails}
Soit $k$ un corps. Soit $X$ un schéma propre sur $k$, de dimension $1$ et tel que $H^0(X, \mathcal{O}_X) = k$. Supposons que les singularités de $X$ soient au plus nodales. Considérons une suite $$
X = X_0 \to X_1 \to \ldots \to X_n = X'
$$ de contractions de queues rationnelles (exemple \ref{curves-example-rational-tail}), poursuivie jusqu'à ce qu'il n'en reste aucune. Alors \begin{enumerate} \item si le genre de $X$ est $0$, $X'$ est une conique plane irréductible; \item si le genre de $X$ est $1$, alors $\omega_{X'} \cong \mathcal{O}_{X'}$; \item si le genre de $X$ est $> 1$, alors $\omega_{X'}^{\otimes m}$ est engendré par ses sections globales pour $m \geq 2$. \end{enumerate} Si le genre de $X$ est $\geq 1$, le morphisme $X \to X'$ est indépendant des choix et sa formation commute aux extensions du corps de base.
\end{lemma}

\begin{proof}
Contractons successivement les queues rationnelles jusqu'à ce qu'il n'en reste aucune. Les assertions (1), (2) et (3) résultent alors du lemme \ref{curves-lemma-no-rational-tail-semiample-genus-geq-2}.

\medskip\noindent
Unicité. Pour voir que $f : X \to X'$ est indépendant des choix effectués, il suffit de montrer que toute queue rationnelle $C \subset X$ est envoyée sur un point par $X \to X'$; nous omettons quelques détails. Dans le cas contraire, il existe une section $s \in \Gamma(X', \omega_{X'}^{\otimes 2})$ qui ne s'annule pas au point générique de la composante irréductible $f(C)$. Puisque, pour chacune des contractions $X_i \to X_{i + 1}$, nous disposons d'une section $X_{i + 1} \to X_i$, il existe une section $X' \to X$ de $f$. Nous avons alors une suite exacte $$
0 \to \omega_{X'} \to \omega_X \to \omega_X|_{X''} \to 0
$$, où $X'' \subset X$ est la réunion des composantes irréductibles contractées par $f$. Voir le lemme \ref{curves-lemma-closed-subscheme-reduced-gorenstein}. Nous obtenons donc une application $\omega_{X'}^{\otimes 2} \to \omega_X^{\otimes 2}$; en prenant l'image de $s$, nous obtenons une section de $\omega_X^{\otimes 2}$ qui ne s'annule pas au point générique de $C$. Cela contredit le fait que la restriction de $\omega_X$ à une queue rationnelle est de degré négatif (lemme \ref{curves-lemma-rational-tail-negative}).

\medskip\noindent
L'assertion concernant les extensions du corps de base résulte du lemme \ref{curves-lemma-rational-tail-field-extension}. Nous omettons quelques détails.
\end{proof}





\section{Contraction des ponts rationnels}
\label{curves-section-contracting-rational-bridges}

\noindent
Dans cette section, nous étudions le deuxième cas le plus simple, après celui de la section \ref{curves-section-contracting-rational-tails}, de contraction d'un schéma visant à améliorer les propriétés de positivité de son faisceau canonique.

\begin{example}[Contraction d'un pont rationnel] \label{curves-example-rational-bridge} Soit $k$ un corps. Soit $X$ un schéma propre sur $k$, de dimension $1$ et tel que $H^0(X, \mathcal{O}_X) = k$. Supposons que les singularités de $X$ soient au plus nodales. Un {\it pont rationnel} est une composante irréductible $C \subset X$, considérée comme un sous-schéma fermé intègre, qui possède les propriétés suivantes : \begin{enumerate} \item $X' \not = \emptyset$, où $X' \subset X$ est l'adhérence schématique de $X \setminus C$; \item l'intersection schématique $C \cap X'$ est de degré $2$ sur $H^0(C, \mathcal{O}_C)$; \item $C$ est de genre zéro. \end{enumerate} Posons $k' = H^0(C, \mathcal{O}_C)$ et $k'' = H^0(C \cap X', \mathcal{O}_{C \cap X'})$. Alors $k'$ est un corps (Variétés, lemme \ref{varieties-lemma-proper-geometrically-reduced-global-sections}) et $\dim_{k'}(k'') = 2$. Comme au moins deux composantes irréductibles de $X$ passent par chaque point de $C \cap X'$, nous concluons que ces points sont des nœuds de $X$ et des points lisses à la fois de $C$ et de $X'$ (lemme \ref{curves-lemma-closed-subscheme-nodal-curve}). Ainsi $k'/k$ est une extension finie séparable de corps, et $k''/k'$ est soit une extension séparable de corps de degré $2$, soit $k'' = k' \times k'$ (lemme \ref{curves-lemma-nodal}). D'après la section \ref{curves-section-pushouts}, il existe un carré cocartésien $$
\xymatrix{
C \cap X' \ar[r] \ar[d] &
X' \ar[d]^a \\
\Spec(k') \ar[r] &
Y
}
$$ doté de nombreuses bonnes propriétés, que nous utiliserons toutes ci-dessous sans autre mention. Soit $y \in Y$ l'image de $\Spec(k') \to Y$. Alors $$
\mathcal{O}_{Y, y}^\wedge \cong k'[[s, t]]/(st)
\quad\text{ou}\quad
\mathcal{O}_{Y, y}^\wedge \cong
\{f \in k''[[s]] : f(0) \in k'\}
$$, selon que $C \cap X'$ possède $2$ ou $1$ points. Cela résulte du lemme \ref{curves-lemma-complete-local-ring-pushout} et du fait que $\mathcal{O}_{X', p} \cong \kappa(p)[[t]]$ pour $p \in C \cap X'$ d'après Compléments d'algèbre, lemme \ref{more-algebra-lemma-power-series-over-residue-field}. Nous voyons ainsi que $y \in Y$ est un nœud; voir les lemmes \ref{curves-lemma-nodal} et \ref{curves-lemma-2-branches-delta-1}, et en particulier la discussion du cas II dans la démonstration de (2) $\Rightarrow$ (1) du lemme \ref{curves-lemma-2-branches-delta-1}. Les singularités de $Y$ sont donc au plus nodales.

\medskip\noindent
Nous pouvons prolonger le diagramme commutatif ci-dessus en un diagramme $$
\xymatrix{
C \cap X' \ar[r] \ar[d] &
X' \ar[d]^a \ar[r] & X \ar[ld]^c &
C \ar[ld] \ar[l] \\
\Spec(k') \ar[r] &
Y &
\Spec(k') \ar[l]
}
$$ dans lequel les deux flèches horizontales inférieures sont identiques. En effet, $X$ est la réunion schématique de $X'$ et $C$, donc un carré cocartésien d'après Morphismes, lemme \ref{morphisms-lemma-scheme-theoretic-union}, et les morphismes $C \to Y$ et $X' \to Y$ coïncident sur $C \cap X'$. Enfin, nous affirmons que $$
c_*\mathcal{O}_X = \mathcal{O}_Y
\quad\text{et}\quad
R^1c_*\mathcal{O}_X = 0
$$ Pour le voir, utilisons la suite exacte $$
0 \to \mathcal{O}_X \to
\mathcal{O}_C \oplus \mathcal{O}_{X'} \to
\mathcal{O}_{C \cap X'} \to 0
$$ de Morphismes, lemme \ref{morphisms-lemma-scheme-theoretic-union}. La suite exacte longue des images directes supérieures est $$
0 \to c_*\mathcal{O}_X \to c_*\mathcal{O}_C \oplus c_*\mathcal{O}_{X'} \to
c_*\mathcal{O}_{C \cap X'} \to
R^1c_*\mathcal{O}_X \to
R^1c_*\mathcal{O}_C \oplus R^1c_*\mathcal{O}_{X'}
$$ Puisque $c|_{X'} = a$ est affine, on a $R^1c_*\mathcal{O}_{X'} = 0$. Comme $c|_C$ se factorise sous la forme $C \to \Spec(k') \to Y$ et que $C$ est de genre zéro, on obtient $R^1c_*\mathcal{O}_C = 0$. Puisque $\mathcal{O}_{X'} \to \mathcal{O}_{C \cap X'}$ est surjective et $c|_{X'}$ affine, $c_*\mathcal{O}_{X'} \to c_*\mathcal{O}_{C \cap X'}$ est surjective. Cela démontre $R^1c_*\mathcal{O}_X = 0$. Enfin, on a $\mathcal{O}_Y = c_*\mathcal{O}_X$ d'après la suite exacte et la description du faisceau structural du carré cocartésien dans Compléments sur les morphismes, proposition \ref{more-morphisms-proposition-pushout-along-closed-immersion-and-integral}.

\medskip\noindent
Tout cela signifie que $Y$ est lui aussi un schéma propre sur $k$, de dimension $1$ et tel que $H^0(Y, \mathcal{O}_Y) = k$, dont les singularités sont au plus nodales (lemme \ref{curves-lemma-closed-subscheme-nodal-curve}), et que $Y$ a le même genre que $X$. On dit que $c : X \to Y$ est la {\it contraction d'un pont rationnel}.
\end{example}

\begin{lemma}
\label{curves-lemma-rational-bridge-zero}
Soit $k$ un corps. Soit $X$ un schéma propre sur $k$, de dimension $1$ et tel que $H^0(X, \mathcal{O}_X) = k$. Supposons que les singularités de $X$ soient au plus nodales. Soit $C \subset X$ un pont rationnel (exemple \ref{curves-example-rational-bridge}). Alors $\deg(\omega_X|_C) = 0$.
\end{lemma}

\begin{proof}
Soit $X' \subset X$ comme dans l'exemple. Nous avons une suite exacte courte $$
0 \to \omega_C \to \omega_X|_C \to \mathcal{O}_{C \cap X'} \to 0
$$ Voir les lemmes \ref{curves-lemma-closed-subscheme-reduced-gorenstein}, \ref{curves-lemma-facts-about-nodal-curves} et \ref{curves-lemma-closed-subscheme-nodal-curve}. Avec $k''/k'/k$ comme dans l'exemple, on voit que $\deg(\omega_C) = -2[k' : k]$, puisque $C$ est de genre $0$ (lemme \ref{curves-lemma-rr}), et que $\deg(C \cap X') = [k'' : k] = 2[k' : k]$. Ainsi $\deg(\omega_X|_C) = 0$.
\end{proof}

\begin{lemma}
\label{curves-lemma-rational-bridge-field-extension}
Soit $k$ un corps. Soit $X$ un schéma propre sur $k$, de dimension $1$ et tel que $H^0(X, \mathcal{O}_X) = k$. Supposons que les singularités de $X$ soient au plus nodales. Soit $C \subset X$ un pont rationnel (exemple \ref{curves-example-rational-bridge}). Pour toute extension de corps $K/k$, le changement de base $C_K \subset X_K$ est une réunion disjointe finie de ponts rationnels.
\end{lemma}

\begin{proof}
Soit $k''/k'/k$ comme dans l'exemple. Puisque $k'/k$ est fini séparable, $k' \otimes_k K = K'_1 \times \ldots \times K'_n$ est un produit fini d'extensions finies séparables $K'_i/K$. La décomposition en produit correspondante $k'' \otimes_k K = \prod K''_i$ fournit des extensions d'algèbres séparables de degré $2$, notées $K''_i/K'_i$. Posons $C_i = C_{K'_i}$. Alors $C_K = \coprod C_i$, et chaque $C_i$ est donc de genre $0$, considéré comme une courbe sur $K'_i$, car $H^1(C_K, \mathcal{O}_{C_K}) = 0$ par changement de base plat. Enfin, $X'_K \cap C_i = \Spec(K''_i)$ est de degré $2$ sur $K'_i$, comme voulu.
\end{proof}

\begin{lemma}
\label{curves-lemma-rational-bridge-canonical}
Soit $c : X \to Y$ la contraction d'un pont rationnel (exemple \ref{curves-example-rational-bridge}). Alors $c^*\omega_Y \cong \omega_X$.
\end{lemma}

\begin{proof}
On peut le démontrer par un calcul direct, mais nous préférons utiliser la caractérisation de $\omega_X$ comme le $\mathcal{O}_X$-module cohérent qui représente le foncteur $\textit{Coh}(\mathcal{O}_X) \to \textit{Ensembles}$, $\mathcal{F} \mapsto \Hom_k(H^1(X, \mathcal{F}), k) =
H^1(X, \mathcal{F})^\vee$; voir le lemme \ref{curves-lemma-duality-dim-1-CM} ou Dualité pour les schémas, lemme \ref{duality-lemma-dualizing-module-proper-over-A}.

\medskip\noindent
Plus précisément, notons $\mathcal{C}_Y$ la catégorie dont les objets sont les $\mathcal{O}_Y$-modules inversibles et dont les flèches sont les homomorphismes de $\mathcal{O}_Y$-modules. Notons $\mathcal{C}_X$ la catégorie dont les objets sont les $\mathcal{O}_X$-modules inversibles $\mathcal{L}$ tels que $\mathcal{L}|_C \cong \mathcal{O}_C$ et dont les flèches sont les homomorphismes de $\mathcal{O}_X$-modules. Nous affirmons que le foncteur $$
c^* : \mathcal{C}_Y \to \mathcal{C}_X
$$ est une équivalence de catégories. En effet, d'après Compléments sur les morphismes, lemme \ref{more-morphisms-lemma-bijection-on-Pic}, il est essentiellement surjectif. La formule de projection (Cohomologie, lemme \ref{cohomology-lemma-projection-formula}) montre ensuite que $c_*c^*\mathcal{N} = \mathcal{N}$, et donc que $c^*$ est une équivalence dont un quasi-inverse est donné par $c_*$.

\medskip\noindent
Nous affirmons que $\omega_X$ est un objet de $\mathcal{C}_X$. En effet, nous avons une suite exacte courte $$
0 \to \omega_C \to \omega_X|_C \to \mathcal{O}_{C \cap X'} \to 0
$$ Voir le lemme \ref{curves-lemma-closed-subscheme-reduced-gorenstein}. En prenant les degrés, nous trouvons $\deg(\omega_X|_C) = 0$ (nous omettons un petit détail). Ainsi $\omega_X|_C$ est trivial d'après le lemme \ref{curves-lemma-genus-zero-pic}, et $\omega_X$ est un objet de $\mathcal{C}_X$.

\medskip\noindent
Puisque $R^1c_*\mathcal{O}_X = 0$, la formule de projection montre que $R^1c_*c^*\mathcal{N} = 0$ pour $\mathcal{N} \in \Ob(\mathcal{C}_Y)$. La suite spectrale de Leray (Cohomologie, lemme \ref{cohomology-lemma-apply-Leray}) montre donc que le diagramme $$
\xymatrix{
\mathcal{C}_Y \ar[rr]_{c^*} \ar[dr]_{H^1(Y, -)^\vee} & &
\mathcal{C}_X \ar[ld]^{H^1(X, -)^\vee} \\
& \textit{Ensembles}
}
$$ de catégories et de foncteurs est commutatif. Puisque $\omega_Y \in \Ob(\mathcal{C}_Y)$ représente la flèche sud-est et que $\omega_X \in \Ob(\mathcal{C}_X)$ représente la flèche sud-ouest, nous concluons par le lemme de Yoneda (Catégories, lemme \ref{categories-lemma-yoneda}).
\end{proof}

\begin{lemma}
\label{curves-lemma-no-rational-bridge-ample-genus-geq-2}
Soit $k$ un corps. Soit $X$ un schéma propre sur $k$, de dimension $1$ et tel que $H^0(X, \mathcal{O}_X) = k$. Supposons que \begin{enumerate} \item les singularités de $X$ soient au plus nodales; \item $X$ ne possède pas de queue rationnelle (exemple \ref{curves-example-rational-tail}); \item $X$ ne possède pas de pont rationnel (exemple \ref{curves-example-rational-bridge}); \item le genre $g$ de $X$ soit $\geq 2$. \end{enumerate} Alors $\omega_X$ est ample.
\end{lemma}

\begin{proof}
Il suffit de montrer que $\deg(\omega_X|_C) > 0$ pour toute composante irréductible $C$ de $X$; voir Variétés, lemme \ref{varieties-lemma-ampleness-in-terms-of-degrees-components}. Si $X = C$ est irréductible, cela résulte de $g \geq 2$ et du lemme \ref{curves-lemma-genus-gorenstein}. Sinon, posons $k' = H^0(C, \mathcal{O}_C)$. C'est un corps et une extension finie de $k$, et $[k' : k]$ divise tous les invariants numériques ci-dessous associés à $C$ et aux faisceaux cohérents sur $C$; voir Variétés, lemme \ref{varieties-lemma-divisible}. Soit $X' \subset X$ l'adhérence de $X \setminus C$ comme dans le lemme \ref{curves-lemma-closed-subscheme-reduced-gorenstein}. Nous utiliserons sans autre mention les résultats de ce lemme et des lemmes \ref{curves-lemma-facts-about-nodal-curves} et \ref{curves-lemma-closed-subscheme-nodal-curve}. Nous obtenons une suite exacte courte $$
0 \to \omega_C \to \omega_X|_C \to \mathcal{O}_{C \cap X'} \to 0
$$ Voir le lemme \ref{curves-lemma-closed-subscheme-reduced-gorenstein}. Nous en concluons que $$
\deg(\omega_X|_C) = \deg(C \cap X') + \deg(\omega_C) =
\deg(C \cap X') - 2\chi(C, \mathcal{O}_C)
$$ Par conséquent, si le lemme est faux, alors $$
2[k' : k] \geq \deg(C \cap X') + 2\dim_k H^1(C, \mathcal{O}_C)
$$ Puisque $C \cap X'$ est non vide, et compte tenu de la divisibilité mentionnée ci-dessus, cela n'est possible que si l'une des deux possibilités suivantes se présente : \begin{enumerate} \item[(a)] $C \cap X'$ est un unique point $k'$-rationnel de $C$ et $H^1(C, \mathcal{O}_C) = 0$; \item[(b)] $C \cap X'$ est de degré $2$ sur $k'$ et $H^1(C, \mathcal{O}_C) = 0$. \end{enumerate} La première possibilité signifie que $C$ est une queue rationnelle et la seconde que $C$ est un pont rationnel. Comme les deux sont exclues, la démonstration est achevée.
\end{proof}

\begin{lemma}
\label{curves-lemma-contracting-rational-bridges}
Soit $k$ un corps. Soit $X$ un schéma propre sur $k$, de dimension $1$, tel que $H^0(X, \mathcal{O}_X) = k$ et de genre $g \geq 2$. Supposons que les singularités de $X$ soient au plus nodales et que $X$ ne possède aucune queue rationnelle. Considérons une suite $$
X = X_0 \to X_1 \to \ldots \to X_n = X'
$$ de contractions de ponts rationnels (exemple \ref{curves-example-rational-bridge}), poursuivie jusqu'à ce qu'il n'en reste aucun. Alors $\omega_{X'}$ est ample. Le morphisme $X \to X'$ est indépendant des choix et sa formation commute aux extensions du corps de base.
\end{lemma}

\begin{proof}
Contractons successivement les ponts rationnels jusqu'à ce qu'il n'en reste aucun. Alors $\omega_{X'}$ est ample d'après le lemme \ref{curves-lemma-no-rational-bridge-ample-genus-geq-2}.

\medskip\noindent
Notons $f : X \to X'$ le composé. Le lemme \ref{curves-lemma-rational-bridge-canonical} et un raisonnement par récurrence donnent $f^*\omega_{X'} = \omega_X$. On a $f_*\mathcal{O}_X = \mathcal{O}_{X'}$, puisque cette égalité est vraie pour la contraction d'un pont rationnel. La formule de projection donne donc $f_*f^*\mathcal{L} = \mathcal{L}$ pour tout $\mathcal{O}_{X'}$-module inversible $\mathcal{L}$. Ainsi $$
\Gamma(X', \omega_{X'}^{\otimes m}) = \Gamma(X, \omega_X^{\otimes m})
$$ pour tout $m$. Puisque $X'$ est le Proj de la somme directe de ces modules, d'après Morphismes, lemme \ref{morphisms-lemma-proper-ample-is-proj}, nous concluons que le morphisme $X \to X'$ est entièrement canonique.

\medskip\noindent
Soit $K/k$ une extension de corps. Alors $\omega_{X_K}$ est l'image réciproque de $\omega_X$ (lemme \ref{curves-lemma-sanity-check-duality}), et $\Gamma(X, \omega_X^{\otimes m}) \otimes_k K$ est égal à $\Gamma(X_K, \omega_{X_K}^{\otimes m})$ d'après Cohomologie des schémas, lemme \ref{coherent-lemma-flat-base-change-cohomology}. Ainsi, la formation de $f : X \to X'$ commute au changement de base d'après $K/k$ et les arguments ci-dessus. Nous omettons quelques détails.
\end{proof}






\section{Contraction vers une courbe stable}
\label{curves-section-contracting-to-stable}

\noindent
Dans cette section, nous combinons les morphismes de contraction obtenus dans les sections \ref{curves-section-contracting-rational-tails} et \ref{curves-section-contracting-rational-bridges}. Supposons donc que $k$ soit un corps et soit $X$ un schéma propre sur $k$, de dimension $1$, tel que $H^0(X, \mathcal{O}_X) = k$ et de genre $g \geq 2$. Supposons que les singularités de $X$ soient au plus nodales. En composant le morphisme du lemme \ref{curves-lemma-contracting-rational-tails} avec celui du lemme \ref{curves-lemma-contracting-rational-bridges}, nous obtenons un morphisme $$
c : X \longrightarrow Y
$$ tel que $Y$ soit lui aussi un schéma propre sur $k$, de dimension $1$, dont les singularités sont au plus nodales, tel que $k = H^0(Y, \mathcal{O}_Y)$, de genre $g$, tel que $\mathcal{O}_Y = c_*\mathcal{O}_X$ et $R^1c_*\mathcal{O}_X = 0$, et tel que $\omega_Y$ soit ample sur $Y$. Le lemme \ref{curves-lemma-characterize-contraction-to-stable} montre que ces conditions caractérisent en fait ce morphisme.

\begin{lemma}
\label{curves-lemma-contract-gorenstein-canonical}
Soit $k$ un corps. Soit $c : X \to Y$ un morphisme de schémas propres sur $k$. Supposons que \begin{enumerate} \item $\mathcal{O}_Y = c_*\mathcal{O}_X$ et $R^1c_*\mathcal{O}_X = 0$; \item $X$ et $Y$ soient réduits, de Gorenstein et de dimension $1$; \item $\exists\ m \in \mathbf{Z}$, avec $H^1(X, \omega_X^{\otimes m}) = 0$ et $\omega_X^{\otimes m}$ engendré par ses sections globales. \end{enumerate} Alors $c^*\omega_Y \cong \omega_X$.
\end{lemma}

\begin{proof}
Les fibres de $c$ sont géométriquement connexes d'après Compléments sur les morphismes, théorème \ref{more-morphisms-theorem-stein-factorization-Noetherian}. En particulier, $c$ est surjectif. Il existe un nombre fini de points fermés $y = y_1, \ldots, y_r$ de $Y$ au-dessus desquels $X_y$ est de dimension $1$, et sur $Y \setminus \{y_1, \ldots, y_r\}$ le morphisme $c$ est un isomorphisme. Nous omettons quelques détails; indication : hors de $\{y_1, \ldots, y_r\}$, le morphisme $c$ est fini; voir Cohomologie des schémas, lemme \ref{coherent-lemma-characterize-finite}.

\medskip\noindent
Construisons soigneusement une application $b : c^*\omega_Y \to \omega_X$. Notons $f : X \to \Spec(k)$ et $g : Y \to \Spec(k)$ les morphismes structuraux. On a $f^!k = \omega_X[1]$ et $g^!k = \omega_Y[1]$; voir le lemme \ref{curves-lemma-duality-dim-1} et sa démonstration. Alors $f^! = c^! \circ g^!$, et donc $c^!\omega_Y = \omega_X$. Il existe ainsi un isomorphisme fonctoriel $$
\Hom_{D(\mathcal{O}_X)}(\mathcal{F}, \omega_X)
\longrightarrow
\Hom_{D(\mathcal{O}_Y)}(Rc_*\mathcal{F}, \omega_Y)
$$ pour tout $\mathcal{O}_X$-module cohérent $\mathcal{F}$, par définition de $c^!$\footnote{Comme restriction de l'adjoint à droite de Dualité pour les schémas, lemme \ref{duality-lemma-twisted-inverse-image}, à $D^+_\QCoh(\mathcal{O}_Y)$.}. Cet isomorphisme est induit par une application trace $t : Rc_*\omega_X \to \omega_Y$, la co-unité de l'adjonction. Par la formule de projection (Cohomologie, lemme \ref{cohomology-lemma-projection-formula}), l'application canonique $a : \omega_Y \to Rc_*c^*\omega_Y$ est un isomorphisme. En combinant ce qui précède, nous obtenons une application canonique $b : c^*\omega_Y \to \omega_X$ telle que $$
t \circ Rc_*(b) = a^{-1}
$$ En particulier, la restriction de $b$ à $c^{-1}(Y \setminus \{y_1, \ldots, y_r\})$ est un isomorphisme, car il s'agit d'une application entre modules inversibles dont le composé avec une autre application est l'isomorphisme $a^{-1}$.

\medskip\noindent
Choisissons $m \in \mathbf{Z}$ comme dans (3) et considérons l'application $$
b^{\otimes m} :
\Gamma(Y, \omega_Y^{\otimes m})
\longrightarrow
\Gamma(X, \omega_X^{\otimes m})
$$ Cette application est injective parce que $Y$ est réduit et par la dernière propriété de $b$ mentionnée lors de sa construction. Le théorème de Riemann--Roch (lemme \ref{curves-lemma-rr}) donne $\chi(X, \omega_X^{\otimes m}) =\chi(Y, \omega_Y^{\otimes m})$. Ainsi $$
\dim_k \Gamma(Y, \omega_Y^{\otimes m}) \geq
\dim_k \Gamma(X, \omega_X^{\otimes m}) = \chi(X, \omega_X^{\otimes m})
$$, et nous concluons que $b^{\otimes m}$ induit un isomorphisme sur les sections globales. Donc $b^{\otimes m} : c^*\omega_Y^{\otimes m} \to \omega_X^{\otimes m}$ est surjective, puisque des générateurs de $\omega_X^{\otimes m}$ appartiennent à son image. Par conséquent, $b^{\otimes m}$ est un isomorphisme, puis $b$ est un isomorphisme.
\end{proof}

\begin{lemma}
\label{curves-lemma-characterize-contraction-to-stable}
Soit $k$ un corps. Soit $X$ un schéma propre sur $k$, de dimension $1$, tel que $H^0(X, \mathcal{O}_X) = k$ et de genre $g \geq 2$. Supposons que les singularités de $X$ soient au plus nodales. Il existe un unique morphisme, à unique isomorphisme près, $$
c : X \longrightarrow Y
$$ de schémas sur $k$ possédant les propriétés suivantes : \begin{enumerate} \item $Y$ est propre sur $k$, $\dim(Y) = 1$, et les singularités de $Y$ sont au plus nodales; \item $\mathcal{O}_Y = c_*\mathcal{O}_X$ et $R^1c_*\mathcal{O}_X = 0$; \item $\omega_Y$ est ample sur $Y$.
\end{enumerate}
\end{lemma}

\begin{proof}
Existence. Un morphisme possédant toutes les propriétés énumérées existe en combinant les lemmes \ref{curves-lemma-contracting-rational-tails} et \ref{curves-lemma-contracting-rational-bridges}, comme expliqué dans l'introduction de cette section. De plus, il peut s'écrire comme un composé $$
X \to X_1 \to X_2 \ldots \to X_n \to X_{n + 1} \to \ldots \to X_{n + n'}
$$ dans lequel les $n$ premiers morphismes sont des contractions de queues rationnelles et les $n'$ derniers des contractions de ponts rationnels. Remarquons que la propriété (2) est satisfaite par toute contraction d'une queue rationnelle (exemple \ref{curves-example-rational-tail}) et par toute contraction d'un pont rationnel (exemple \ref{curves-example-rational-bridge}). On voit aisément que cette propriété est préservée par composition.

\medskip\noindent
Unicité. Soit $c : X \to Y$ un morphisme satisfaisant les conditions (1), (2) et (3). Nous allons montrer qu'il existe un unique isomorphisme $X_{n + n'} \to Y$ compatible avec les morphismes $X \to X_{n + n'}$ et $c$.

\medskip\noindent
Avant de commencer la démonstration, faisons quelques observations sur $c$. Tout d'abord, les fibres de $c$ sont géométriquement connexes d'après Compléments sur les morphismes, théorème \ref{more-morphisms-theorem-stein-factorization-Noetherian}. En particulier, $c$ est surjectif. Pour tout point fermé $y \in Y$, la fibre $X_y$ satisfait $$
H^1(X_y, \mathcal{O}_{X_y}) = 0
\quad\text{et}\quad
H^0(X_y, \mathcal{O}_{X_y}) = \kappa(y)
$$ La première égalité résulte de Compléments sur les morphismes, lemme \ref{more-morphisms-lemma-check-h1-fibre-zero}, et la seconde de Compléments sur les morphismes, lemme \ref{more-morphisms-lemma-h1-fibre-zero-check-h0-kappa}. Ainsi, soit $X_y = x$, où $x$ est l'unique point de $X$ s'envoyant sur $y$ et possède le même corps résiduel que $y$, soit $X_y$ est un schéma propre de dimension $1$ sur $\kappa(y)$. Dans le second cas, remarquons que $X_y$ est de Cohen--Macaulay (lemme \ref{curves-lemma-automatic}). Toutefois, puisque $X$ est réduit, $X_y$ doit être réduit en tous ses points génériques (nous omettons les détails), et donc $X_y$ est réduit d'après Propriétés, lemme \ref{properties-lemma-criterion-reduced}. Il s'ensuit que les singularités de $X_y$ sont au plus nodales (lemme \ref{curves-lemma-closed-subscheme-nodal-curve}). Remarquons que $X_y$ est de genre zéro, d'après ce qui précède. Enfin, il n'existe qu'un nombre fini de points $y$ au-dessus desquels la fibre $X_y$ est de dimension $1$; notons-les $\{y_1, \ldots, y_r\}$. Le sous-schéma $c^{-1}(Y \setminus \{y_1, \ldots, y_r\})$ s'envoie isomorphiquement sur $Y \setminus \{y_1, \ldots, y_r\}$ par $c$. Nous omettons quelques détails; indication : hors de $\{y_1, \ldots, y_r\}$, le morphisme $c$ est fini; voir Cohomologie des schémas, lemme \ref{coherent-lemma-characterize-finite}.

\medskip\noindent
Soit $C \subset X$ une queue rationnelle. Nous affirmons que $c$ envoie $C$ sur un point. Supposons le contraire pour aboutir à une contradiction. Alors l'image de $C$ est une composante irréductible $D \subset Y$. Rappelons que $H^0(C, \mathcal{O}_C) = k'$ est une extension finie séparable de $k$, et que $C$ possède un point $k'$-rationnel $x$ qui est aussi l'unique point d'intersection de $C$ avec le « reste » de $X$. La discussion générale ci-dessus montre que $C \setminus \{x\} \subset c^{-1}(Y \setminus \{y_1, \ldots, y_r\})$ s'envoie isomorphiquement sur un ouvert $V$ de $D$. Soit $y = c(x) \in D$. Remarquons que $y$ est l'unique point de $D$ qui rencontre le « reste » de $Y$. Si $y \not \in \{y_1, \ldots, y_r\}$, alors $C \cong D$, et il est clair que $D$ est une queue rationnelle de $Y$, ce qui contredit l'amplitude de $\omega_Y$ (lemme \ref{curves-lemma-rational-tail-negative}). Ainsi $y \in \{y_1, \ldots, y_r\}$ et $\dim(X_y) = 1$. Alors $x \in X_y \cap C$, et $x$ est un point lisse de $X_y$ et de $C$ (lemme \ref{curves-lemma-closed-subscheme-nodal-curve}). Si $y \in D$ est un point singulier de $D$, alors $y$ est un nœud et $Y = D$, car aucune autre composante de $Y$ ne peut passer par $y$ d'après le lemme \ref{curves-lemma-closed-subscheme-nodal-curve}. Alors $X = X_y \cup C$, ce qui signifie que $g = 0$, puisque ce nombre est égal au genre de $X_y$ d'après la discussion de l'exemple \ref{curves-example-rational-tail}; contradiction. Si $y \in D$ est un point lisse de $D$, alors $C \to D$ est un isomorphisme, car le modèle projectif non singulier est unique et $C$ et $D$ sont birationnels; voir la section \ref{curves-section-curves-function-fields}. Alors $D$ est une queue rationnelle de $Y$, ce qui contredit l'amplitude de $\omega_Y$.

\medskip\noindent
Supposons $n \geq 1$. Si $C \subset X$ est la queue rationnelle contractée par $X \to X_1$, le paragraphe précédent montre que $C$ est envoyé sur un point de $Y$. Ainsi $c : X \to Y$ se factorise par $X \to X_1$, puisque $X$ est le carré cocartésien de $C$ et $X_1$; voir la discussion de l'exemple \ref{curves-example-rational-tail}. Après avoir remplacé $X$ par $X_1$, nous avons diminué $n$. Par récurrence, nous pouvons supposer $n = 0$, c'est-à-dire que $X$ ne possède aucune queue rationnelle.

\medskip\noindent
Supposons $n = 0$, c'est-à-dire que $X$ ne possède aucune queue rationnelle. Alors $\omega_X^{\otimes 2}$ et $\omega_X^{\otimes 3}$ sont engendrés par leurs sections globales d'après le lemme \ref{curves-lemma-no-rational-tail-semiample-genus-geq-2}. Il s'ensuit que $H^1(X, \omega_X^{\otimes 3}) = 0$ d'après le lemme \ref{curves-lemma-vanishing-twist}. En appliquant le lemme \ref{curves-lemma-contract-gorenstein-canonical} avec $m = 3$, nous obtenons $c^*\omega_Y \cong \omega_X$. Nous avons aussi $\omega_X = (X \to X_{n'})^*\omega_{X_{n'}}$ d'après le lemme \ref{curves-lemma-rational-bridge-canonical} et par récurrence. En appliquant la formule de projection à $c$ et à $X \to X_{n'}$, nous concluons que $$
\Gamma(X_{n'}, \omega_{X_{n'}}^{\otimes m}) =
\Gamma(X, \omega_X^{\otimes m}) =
\Gamma(Y, \omega_Y^{\otimes m})
$$ pour tout $m$. Puisque $X_{n'}$ et $Y$ sont les Proj des sommes directes correspondantes, d'après Morphismes, lemme \ref{morphisms-lemma-proper-ample-is-proj}, nous obtenons un isomorphisme canonique $X_{n'} = Y$, comme voulu. Nous omettons la vérification que c'est l'unique isomorphisme rendant le diagramme commutatif.
\end{proof}

\begin{lemma}
\label{curves-lemma-tricanonical}
Soit $k$ un corps. Soit $X$ un schéma propre sur $k$, de dimension $1$, tel que $H^0(X, \mathcal{O}_X) = k$ et de genre $g \geq 2$. Supposons que les singularités de $X$ soient au plus nodales et que $\omega_X$ soit ample. Alors $\omega_X^{\otimes 3}$ est très ample et $H^1(X, \omega_X^{\otimes 3}) = 0$.
\end{lemma}

\begin{proof}
En combinant Variétés, lemme \ref{varieties-lemma-ampleness-in-terms-of-degrees-components}, et les lemmes \ref{curves-lemma-rational-tail-negative} et \ref{curves-lemma-rational-bridge-zero}, nous voyons que $X$ ne contient ni queue rationnelle ni pont rationnel. Alors $\omega_X^{\otimes 3}$ est engendré par ses sections globales d'après le lemme \ref{curves-lemma-contracting-rational-tails}. Choisissons une $k$-base $s_0, \ldots, s_n$ de $H^0(X, \omega_X^{\otimes 3})$. Nous obtenons un morphisme $$
\varphi_{\omega_X^{\otimes 3}, (s_0, \ldots, s_n)} :
X \longrightarrow \mathbf{P}^n_k
$$ Voir Constructions, section \ref{constructions-section-projective-space}. Le lemme affirme que ce morphisme est une immersion fermée. Pour le vérifier, nous pouvons remplacer $k$ par sa clôture algébrique; voir Descente, lemme \ref{descent-lemma-descending-property-closed-immersion}. Nous pouvons donc supposer $k$ algébriquement clos.

\medskip\noindent
Supposons $k$ algébriquement clos. Nous allons utiliser Variétés, lemme \ref{varieties-lemma-variant-separate-points-tangent-vectors}, pour démontrer le lemme. Soit $Z \subset X$ un sous-schéma fermé de degré $2$ sur $k$, de faisceau d'idéaux $\mathcal{I} \subset \mathcal{O}_X$. Il faut montrer que $$
H^0(X, \omega_X^{\otimes 3}) \to H^0(Z, \omega_X^{\otimes 3}|_Z)
$$ est surjective. Il suffit donc d'établir $H^1(X, \mathcal{I}\omega_X^{\otimes 3}) = 0$. Pour cela, nous utiliserons le lemme \ref{curves-lemma-vanishing-on-gorenstein}. Il suffit ainsi de montrer que $$
3\deg(\omega_X|_Y) > -2\chi(Y, \mathcal{O}_Y) + \deg(Z \cap Y)
$$ pour tout sous-schéma fermé réduit connexe $Y \subset X$. Puisque $k$ est algébriquement clos et que $Y$ est connexe et réduit, on a $H^0(Y, \mathcal{O}_Y) = k$ (Variétés, lemme \ref{varieties-lemma-proper-geometrically-reduced-global-sections}). Par conséquent, $\chi(Y, \mathcal{O}_Y) = 1 - \dim H^1(Y, \mathcal{O}_Y)$. Il reste donc à montrer $$
3\deg(\omega_X|_Y) > -2 + 2\dim H^1(Y, \mathcal{O}_Y) + \deg(Z \cap Y)
$$, ce qui résulte du lemme \ref{curves-lemma-no-rational-tail}, sauf éventuellement si $Y = X$ ou si $\deg(\omega_X|_Y) = 0$. Puisque $\omega_X$ est ample, la seconde possibilité ne se présente pas; voir le premier lemme cité dans cette démonstration. Enfin, si $Y = X$, le théorème de Riemann--Roch (lemme \ref{curves-lemma-rr}) et le fait que $g \geq 2$ montrent que l'inégalité est satisfaite. Le même argument avec $Z = \emptyset$ donne $H^1(X, \omega_X^{\otimes 3}) = 0$.
\end{proof}







\section{Champs de vecteurs}
\label{curves-section-vector-fields}

\noindent
Dans cette section, nous étudions l'espace des champs de vecteurs sur une courbe. Les champs de vecteurs correspondent aux automorphismes infinitésimaux; voir Compléments sur les morphismes, section \ref{more-morphisms-section-action-by-derivations}. Ils jouent donc un rôle important dans la théorie des espaces de modules.

\medskip\noindent
Soit $k$ un corps algébriquement clos. Soit $X$ un schéma de type fini sur $k$. Soit $x \in X$ un point fermé. Nous dirons qu'un élément $D \in \text{Der}_k(\mathcal{O}_X, \mathcal{O}_X)$ {\it fixe $x$} si $D(\mathcal{I}) \subset \mathcal{I}$, où $\mathcal{I} \subset \mathcal{O}_X$ est le faisceau d'idéaux de $x$.

\begin{lemma}
\label{curves-lemma-smooth-vector-fields}
Soit $k$ un corps algébriquement clos. Soit $X$ une courbe lisse, propre et connexe sur $k$. Soit $g$ le genre de $X$. \begin{enumerate} \item Si $g \geq 2$, alors $\text{Der}_k(\mathcal{O}_X, \mathcal{O}_X)$ est nul; \item si $g = 1$ et si $D \in \text{Der}_k(\mathcal{O}_X, \mathcal{O}_X)$ est non nul, alors $D$ ne fixe aucun point fermé de $X$; \item si $g = 0$ et si $D \in \text{Der}_k(\mathcal{O}_X, \mathcal{O}_X)$ est non nul, alors $D$ fixe au plus $2$ points fermés de $X$.
\end{enumerate}
\end{lemma}

\begin{proof}
Rappelons que nous avons une $k$-dérivation universelle $d : \mathcal{O}_X \to \Omega_{X/k}$, et donc $D = \theta \circ d$ pour une certaine application $\mathcal{O}_X$-linéaire $\theta : \Omega_{X/k} \to \mathcal{O}_X$. Rappelons que $\Omega_{X/k} \cong \omega_X$; voir le lemme \ref{curves-lemma-duality-dim-1}. Le théorème de Riemann--Roch donne $\deg(\omega_X) = 2g - 2$ (lemme \ref{curves-lemma-rr}). Ainsi $\theta$ est nécessairement nul si $g > 1$, d'après Variétés, lemme \ref{varieties-lemma-check-invertible-sheaf-trivial}. Cela démontre la partie (1). Si $g = 1$, un $\theta$ non nul ne s'annule nulle part; si $g = 0$, un $\theta$ non nul s'annule suivant un diviseur de degré $2$. Les parties (2) et (3) en résultent dès que l'on montre que l'annulation de $\theta$ en un point fermé $x \in X$ équivaut à dire que $D$ fixe $x$, au sens défini ci-dessus. Soit $z \in \mathcal{O}_{X, x}$ une uniformisante. Alors $dz$ est un élément de base de $\Omega_{X, x}$; voir le lemme \ref{curves-lemma-uniformizer-works}. Puisque $D(z) = \theta(dz)$, nous concluons.
\end{proof}

\begin{lemma}
\label{curves-lemma-nodal-vector-fields}
Soit $k$ un corps algébriquement clos. Soit $X$ un schéma propre connexe de dimension $1$ sur $k$, à singularités au plus nodales. Soit $\nu : X^\nu \to X$ la normalisation. Soit $S \subset X^\nu$ l'ensemble des points où $\nu$ n'est pas un isomorphisme. Alors
$$
\text{Der}_k(\mathcal{O}_X, \mathcal{O}_X) =
\{D' \in \text{Der}_k(\mathcal{O}_{X^\nu}, \mathcal{O}_{X^\nu}) \mid
D' \text{ fixe tout }x^\nu \in S\}
$$
\end{lemma}

\begin{proof}
Soit $x \in X$ un nœud. Soient $x', x'' \in X^\nu$ les images réciproques de $x$. Tout nœud est déployé puisque $k$ est algébriquement clos; voir la définition \ref{curves-definition-split-node} et le lemme \ref{curves-lemma-split-node}. Soient $u \in \mathcal{O}_{X^\nu, x'}$ et $v \in \mathcal{O}_{X^\nu, x''}$ des uniformisantes. Remarquons que nous avons une suite exacte $$
0 \to \mathcal{O}_{X, x} \to
\mathcal{O}_{X^\nu, x'} \times \mathcal{O}_{X^\nu, x''} \to k \to 0
$$ Cela résulte du lemme \ref{curves-lemma-multicross}. Nous pouvons donc considérer $u$ et $v$ comme des éléments de $\mathcal{O}_{X, x}$ tels que $uv = 0$.

\medskip\noindent
Soit $D \in \text{Der}_k(\mathcal{O}_X, \mathcal{O}_X)$. Alors $0 = D(uv) = vD(u) + uD(v)$. Puisque $(u)$ est l'annulateur de $v$ dans $\mathcal{O}_{X, x}$ et réciproquement, on voit que $D(u) \in (u)$ et $D(v) \in (v)$. Comme $\mathcal{O}_{X^\nu, x'} = k + (u)$, nous concluons que nous pouvons prolonger $D$ en $\mathcal{O}_{X^\nu, x'}$ et que, de plus, ce prolongement fixe $x'$. Cela produit un $D'$ dans le membre de droite de l'égalité. Réciproquement, étant donné $D'$ qui fixe $x'$ et $x''$, nous voyons que $D'$ préserve le sous-anneau $\mathcal{O}_{X, x} \subset
\mathcal{O}_{X^\nu, x'} \times \mathcal{O}_{X^\nu, x''}$; c'est ainsi que l'on passe du membre de droite au membre de gauche de l'égalité.
\end{proof}

\begin{lemma}
\label{curves-lemma-stable-vector-fields}
Soit $k$ un corps algébriquement clos. Soit $X$ un schéma propre connexe de dimension $1$ sur $k$, à singularités au plus nodales. Supposons que le genre de $X$ soit au moins $2$ et que $X$ ne possède ni queue rationnelle ni pont rationnel. Alors $\text{Der}_k(\mathcal{O}_X, \mathcal{O}_X) = 0$.
\end{lemma}

\begin{proof}
Soit $D \in \text{Der}_k(\mathcal{O}_X, \mathcal{O}_X)$. Soit $X^\nu$ la normalisation de $X$. Soit $D' \in \text{Der}_k(\mathcal{O}_{X^\nu}, \mathcal{O}_{X^\nu})$ l'élément correspondant à $D$ par le lemme \ref{curves-lemma-nodal-vector-fields}. Soit $C \subset X^\nu$ une composante irréductible. Si le genre de $C$ est $> 1$, alors $D'|_{\mathcal{O}_C} = 0$ d'après la partie (1) du lemme \ref{curves-lemma-smooth-vector-fields}. Si le genre de $C$ est $1$, il existe au moins un point fermé $c$ de $C$ qui s'envoie sur un nœud de $X$, car sinon $X \cong C$ serait de genre $1$. Par la correspondance, cela signifie que $D'|_{\mathcal{O}_C}$ fixe $c$, et il est donc nul d'après la partie (2) du lemme \ref{curves-lemma-smooth-vector-fields}. Enfin, si le genre de $C$ est nul, il existe au moins $3$ points fermés deux à deux distincts $c_1, c_2, c_3 \in C$ qui s'envoient sur des nœuds de $X$. Sinon, soit $X$ serait $C$ avec deux points recollés, deux points de $C$ s'envoyant sur le même nœud, soit $C$ serait un pont rationnel, deux points s'envoyant sur des nœuds distincts de $X$, soit $C$ serait une queue rationnelle, un point s'envoyant sur un nœud de $X$. Ces trois possibilités sont exclues puisque $X$ est de genre $\geq 2$ et ne possède ni pont rationnel ni queue rationnelle. Ainsi $D'|_{\mathcal{O}_C}$ fixe $c_1, c_2, c_3$ et est donc nul d'après la partie (3) du lemme \ref{curves-lemma-smooth-vector-fields}.
\end{proof}









% OK, start here.
%

\chapter{Résolution des surfaces}



\phantomsection
\label{resolve-section-phantom}


\section{Introduction} \label{resolve-section-introduction}

\noindent Ce chapitre traite de la résolution des singularités des surfaces en suivant Lipman \cite{Lipman} et, pour l'essentiel, l'exposé d'Artin dans \cite{Artin-Lipman}. Le résultat principal (théorème \ref{resolve-theorem-resolve}) affirme qu'un schéma noethérien de dimension $2$, noté $Y$, admet une résolution des singularités lorsqu'il possède une normalisation finie $Y^\nu \to Y$ ayant un nombre fini de points singuliers $y_i \in Y^\nu$ et que, pour tout $i$, le complété $\mathcal{O}_{Y^\nu, y_i}^\wedge$ est normal.

\medskip\noindent En particulier, si $Y$ est un schéma de dimension $2$ de type fini sur un anneau de base quasi-excellent $R$ (par exemple un corps ou un anneau de Dedekind dont le corps des fractions est de caractéristique $0$, tel que $\mathbf{Z}$), alors la normalisation de $Y$ est finie, possède un nombre fini de points singuliers et les complétés de ses anneaux locaux sont normaux. Voir la discussion dans Compléments d'algèbre, sections \ref{more-algebra-section-singular-locus}, \ref{more-algebra-section-G-ring} et \ref{more-algebra-section-excellent}, ainsi que Compléments d'algèbre, lemme \ref{more-algebra-lemma-normal-goes-up}. Un tel $Y$ admet donc une résolution des singularités.

\medskip\noindent Voici les grandes lignes de la démonstration. Soit $A$ un anneau local noethérien intègre de dimension $2$. Les étapes de la démonstration sont les suivantes : \begin{enumerate} \item[N] remplacer $A$ par son normalisé ; \item[V] démontrer le théorème de Grauert-Riemenschneider ; \item[B] montrer qu'il existe un maximum $g$ des longueurs de $H^1(X, \mathcal{O}_X)$ lorsque l'on parcourt toutes les modifications normales $X \to \Spec(A)$, et se ramener au cas $g = 0$ ; \item[R] on dit que $A$ définit une singularité rationnelle si $g = 0$ ; dans ce cas, après un nombre fini d'éclatements, on peut supposer que $A$ est de Gorenstein et que $g = 0$ ; \item[D] on dit que $A$ définit un point double rationnel si $g = 0$ et si $A$ est de Gorenstein ; dans ce cas, on résout explicitement les singularités. \end{enumerate} Chacune de ces étapes nécessite des hypothèses sur l'anneau $A$. Nous les examinerons successivement.

\medskip\noindent Étape N : il faut ici supposer que $A$ possède une normalisation finie (ce qui n'est pas automatique). Dans la majeure partie de ce chapitre, nous supposerons que notre schéma est de Nagata lorsqu'il nous faudra savoir qu'une normalisation est finie. Toutefois, être de Nagata est une condition légèrement plus forte que celle qui figure dans l'énoncé du théorème. Une solution à ce léger problème aurait consisté à utiliser le fait que notre anneau $A$ est formellement non ramifié (c'est-à-dire que son complété est réduit) et à appliquer le lemme \ref{resolve-lemma-formally-unramified}. Cependant, avec la démonstration adoptée ici, il s'avère plus simple d'utiliser le lemme \ref{resolve-lemma-normalization-completion} pour déduire de la finitude de la normalisation au-dessus du complété celle de la normalisation au-dessus de $A$.

\medskip\noindent Étape V : il s'agit de la proposition \ref{resolve-proposition-Grauert-Riemenschneider}, qui affirme, en substance, que pour une modification normale $f : X \to \Spec(A)$ on a $R^1f_*\omega_X = 0$, où $\omega_X$ est le module dualisant de $X/A$ (remarque \ref{resolve-remark-dualizing-setup}). En fait, par dualité, ce résultat équivaut à un énoncé (lemme \ref{resolve-lemma-R1-injective}) concernant l'objet $Rf_*\mathcal{O}_X$ de la {\it catégorie dérivée} $D(A)$. Cela dit, la démonstration utilise le fait classique que les composantes de la fibre spéciale ont des faisceaux conormaux positifs (lemme \ref{resolve-lemma-nontrivial-normal-bundle}).

\medskip\noindent Étape B : c'est, en un sens, la partie la plus délicate de la démonstration. Nous n'aurons finalement besoin du résultat de cette étape que lorsque $A$ est un anneau local noethérien complet, bien que l'exposé soit un peu plus général. La terminologie est fixée dans la définition \ref{resolve-definition-reduce-to-rational}. Si $g$ (défini ci-dessus) est borné, un argument simple montre que l'on peut trouver une modification normale $X \to \Spec(A)$ telle que tous les points singuliers de $X$ soient des singularités rationnelles ; voir le lemme \ref{resolve-lemma-reduce-to-rational}. Nous montrons que, pour une extension finie $A \subset B$, l'invariant $g$ est borné pour $B$ s'il l'est pour $A$ dans les deux cas suivants : (1) lorsque l'extension des corps des fractions est séparable, voir le lemme \ref{resolve-lemma-reduce-to-rational} ; (2) lorsque l'extension des corps des fractions est de degré $p$, que la caractéristique est $p$ et que $A$ est régulier et complet, voir le lemme \ref{resolve-lemma-go-up-degree-p}.

\medskip\noindent Étape R : nous ramenons ici le cas $g = 0$ au cas de Gorenstein. Le fait remarquable qui permet de mener à bien la démonstration est que l'éclatement d'une singularité rationnelle de surface est normal ; voir le lemme \ref{resolve-lemma-blow-up-normal-rational}.

\medskip\noindent Étape D : la résolution des points doubles rationnels se fait par des calculs assez explicites ; voir la section \ref{resolve-section-rational-double-points}. Un point double rationnel est une singularité d'hypersurface (c'est vrai, mais nous ne le démontrons pas, car nous n'en avons pas besoin). Son équation locale est de la forme $$
a_{11} x_1^2 + a_{12} x_1x_2 + a_{13}x_1x_3 + a_{22} x_2^2 +
a_{23} x_2x_3 + a_{33} x_3^2 =
\sum a_{ijk} x_ix_jx_k
$$ Le fait que la partie quadratique ne puisse être nulle, puisque la multiplicité vaut $2$ et reste égale à $2$ après tout éclatement, joint à la normalité de chaque éclatement, permet d'obtenir rapidement une résolution. Une particularité est que nous ne disposons pas d'un invariant décroissant à chaque éclatement ; nous nous appuyons sur les résultats relatifs aux arcs formels de la section \ref{resolve-section-arcs} pour démontrer que le procédé s'arrête.

\medskip\noindent Pour réunir tous ces éléments, un travail supplémentaire est nécessaire. La difficulté principale consiste à relever une résolution du complété $A^\wedge$ en une résolution de $\Spec(A)$. Pour cela, nous montrons d'abord que, s'il existe une résolution, il en existe une par éclatements normalisés (lemme \ref{resolve-lemma-existence-implies-existence-by-normalized-blowing-ups}). Une suite d'éclatements normalisés peut être relevée à partir du complété grâce au lemme \ref{resolve-lemma-normalized-blowup-completion}. Nous utilisons ensuite ce fait jusque dans la démonstration de la résolution des anneaux locaux complets $A$, car notre méthode procède par récurrence sur le degré d'une inclusion finie $A_0 \subset A$ avec $A_0$ régulier ; voir le lemme \ref{resolve-lemma-resolve-complete}. Un résultat plus fort à l'étape B (tel que celui démontré dans l'article de Lipman) permettrait d'éviter cette étape.




\section{Une application trace en caractéristique positive} \label{resolve-section-trace}

\noindent Certains résultats de cette section se déduisent de l'étude beaucoup plus générale des traces sur les formes différentielles dans Cohomologie de de Rham, section \ref{derham-section-trace}. Voir la remarque \ref{resolve-remark-compare-Garel} pour une discussion.

\medskip\noindent Fixons un nombre premier $p$. Soit $R$ une $\mathbf{F}_p$-algèbre. Étant donné $a \in R$, posons $S = R[x]/(x^p - a)$. Définissons une application $R$-linéaire $$
\text{Tr}_x : \Omega_{S/R} \longrightarrow \Omega_R
$$ par la règle $$
x^i\text{d}x \longmapsto
\left\{
\begin{matrix}
0 & \text{si} & 0 \leq i \leq p - 2, \\
\text{d}a & \text{si} & i = p - 1
\end{matrix}
\right.
$$ Cela a un sens, car $\Omega_{S/R}$ est un $R$-module libre de base $x^i\text{d}x$, $0 \leq i \leq p - 1$. Le lemme suivant implique que l'application trace est bien définie, c'est-à-dire indépendante du choix de la coordonnée $x$.

\begin{lemma} \label{resolve-lemma-trace-well-defined} Soit $\varphi : R[x]/(x^p - a) \to R[y]/(y^p - b)$ un homomorphisme de $R$-algèbres. Alors $\text{Tr}_x = \text{Tr}_y \circ \varphi$. \end{lemma}

\begin{proof} Écrivons $\varphi(x) = \lambda_0 + \lambda_1 y + \ldots + \lambda_{p - 1}y^{p - 1}$ avec $\lambda_i \in R$. Dire que l'envoi de $x$ sur $\lambda_0 + \lambda_1 y + \ldots + \lambda_{p - 1}y^{p - 1}$ induit un homomorphisme de $R$-algèbres $R[x]/(x^p - a) \to R[y]/(y^p - b)$ équivaut à la condition $$
a = \lambda_0^p + \lambda_1^p b + \ldots + \lambda_{p - 1}^pb^{p - 1}
$$ dans l'anneau $R$. Considérons l'anneau de polynômes $$
R_{univ} = \mathbf{F}_p[b, \lambda_0, \ldots, \lambda_{p - 1}]
$$ muni de l'élément $a = \lambda_0^p + \lambda_1^p b + \ldots + \lambda_{p - 1}^pb^{p - 1}$ Considérons l'homomorphisme universel d'algèbres $\varphi_{univ} : R_{univ}[x]/(x^p - a) \to R_{univ}[y]/(y^p - b)$ défini par l'envoi de $x$ sur $\lambda_0 + \lambda_1 y + \ldots + \lambda_{p - 1}y^{p - 1}$. On obtient une application canonique $$
R_{univ} \longrightarrow R
$$ envoyant $b, \lambda_i$ sur $b, \lambda_i$. Par construction, on a un diagramme commutatif $$
\xymatrix{
R_{univ}[x]/(x^p - a) \ar[r] \ar[d]_{\varphi_{univ}} &
R[x]/(x^p - a) \ar[d]^\varphi \\
R_{univ}[y]/(y^p - b) \ar[r] & R[y]/(y^p - b)
}
$$ dont les flèches horizontales sont compatibles aux applications trace. Il suffit donc de démontrer le lemme pour l'application $\varphi_{univ}$. On peut ainsi supposer que $R = \mathbf{F}_p[b, \lambda_0, \ldots, \lambda_{p - 1}]$ est un anneau de polynômes. Nous vérifierons le lemme dans ce cas en calculant $\text{Tr}_y(\varphi(x)^i\text{d}\varphi(x))$ pour $i = 0 , \ldots, p - 1$.

\medskip\noindent Cas $0 \leq i \leq p - 2$. Développons $$
(\lambda_0 + \lambda_1 y + \ldots + \lambda_{p - 1}y^{p - 1})^i
(\lambda_1 + 2 \lambda_2 y + \ldots + (p - 1)\lambda_{p - 1}y^{p - 2})
$$ dans l'anneau $R[y]/(y^p - b)$. Il faut montrer que le coefficient de $y^{p - 1}$ est nul. Pour cela, il suffit de montrer que, dans l'expression ci-dessus considérée comme polynôme en $y$, les coefficients des puissances $y^{pk - 1}$ sont nuls. Écrivons alors notre polynôme sous la forme $$
\frac{\text{d}}{(i + 1)\text{d}y}
(\lambda_0 + \lambda_1 y + \ldots + \lambda_{p - 1}y^{p - 1})^{i + 1}
$$ Les coefficients de $y^{kp - 1}$ sont effectivement tous nuls.

\medskip\noindent Cas $i = p - 1$. Développons $$
(\lambda_0 + \lambda_1 y + \ldots + \lambda_{p - 1}y^{p - 1})^{p - 1}
(\lambda_1 + 2 \lambda_2 y + \ldots + (p - 1)\lambda_{p - 1}y^{p - 2})
$$ dans l'anneau $R[y]/(y^p - b)$. Pour achever la démonstration, il faut montrer que le coefficient de $y^{p - 1}$ multiplié par $\text{d}b$ vaut $\text{d}a$. Nous utilisons ici le fait que $R$ est égal à $S/pS$, où $S = \mathbf{Z}[b, \lambda_0, \ldots, \lambda_{p - 1}]$. L'expression ci-dessus, considérée comme polynôme en $y$, vaut alors $$
\frac{\text{d}}{p\text{d}y}
(\lambda_0 + \lambda_1 y + \ldots + \lambda_{p - 1}y^{p - 1})^p
$$ Puisque $\frac{\text{d}}{\text{d}y}(y^{pk}) = pk y^{pk - 1}$, il suffit de déterminer les coefficients de $y^{pk}$ dans le polynôme $(\lambda_0 + \lambda_1 y + \ldots + \lambda_{p - 1}y^{p - 1})^p$ modulo $p$. La somme de ces termes donne $$
\lambda_0^p + \lambda_1^py^p + \ldots + \lambda_{p - 1}^py^{p(p - 1)}
\bmod p
$$ On voit donc qu'après application de l'opérateur $\frac{\text{d}}{p\text{d}y}$ et réduction modulo $y^p - b$, on obtient la valeur $$
\lambda_1^p + 2\lambda_2^pb + \ldots + (p - 1)\lambda_{p - 1}b^{p - 2}
$$ pour le coefficient de $y^{p - 1}$ que nous voulions calculer. Or, puisque $a = \lambda_0^p + \lambda_1^p b + \ldots + \lambda_{p - 1}^pb^{p - 1}$ dans $R$, on obtient $$
\text{d}a = (\lambda_1^p  + 2 \lambda_2^p b + \ldots +
(p - 1) \lambda_{p - 1}^p b^{p - 2}) \text{d}b
$$ dans $R$. Le coefficient de $y^{p - 1}$ est donc bien celui que l'on cherchait. \end{proof}

\begin{lemma} \label{resolve-lemma-trace-higher} Soient $\mathbf{F}_p \subset \Lambda \subset R \subset S$ des extensions d'anneaux ; supposons que $S$ soit isomorphe à $R[x]/(x^p - a)$ pour un certain $a \in R$. Il existe alors des applications canoniques $R$-linéaires $$
\text{Tr} :
\Omega^{t + 1}_{S/\Lambda}
\longrightarrow
\Omega_{R/\Lambda}^{t + 1}
$$ pour $t \geq 0$, telles que $$
\eta_1 \wedge \ldots \wedge \eta_t \wedge x^i\text{d}x
\longmapsto
\left\{
\begin{matrix}
0 & \text{si} & 0 \leq i \leq p - 2, \\
\eta_1 \wedge \ldots \wedge \eta_t \wedge \text{d}a & \text{si} & i = p - 1
\end{matrix}
\right.
$$ pour $\eta_i \in \Omega_{R/\Lambda}$ et telles que $\text{Tr}$ annule l'image de $S \otimes_R \Omega_{R/\Lambda}^{t + 1} \to \Omega_{S/\Lambda}^{t + 1}$. \end{lemma}

\begin{proof} Pour $t = 0$, on utilise le composé $$
\Omega_{S/\Lambda} \to \Omega_{S/R} \to \Omega_R \to \Omega_{R/\Lambda}
$$ où la deuxième application est celle du lemme \ref{resolve-lemma-trace-well-defined}. On a une suite exacte $$
H_1(L_{S/R}) \xrightarrow{\delta} \Omega_{R/\Lambda} \otimes_R S \to
\Omega_{S/\Lambda} \to \Omega_{S/R} \to 0
$$ (Algèbre, lemme \ref{algebra-lemma-exact-sequence-NL}). Le module $\Omega_{S/R}$ est libre sur $S$ de base $\text{d}x$, et le module $H_1(L_{S/R})$ est libre sur $S$ de base $x^p - a$, que $\delta$ envoie sur $-\text{d}a \otimes 1$ dans $\Omega_{R/\Lambda} \otimes_R S$. En particulier, si l'on pose $$
M = \Coker(R \to \Omega_{R/\Lambda}, 1 \mapsto -\text{d}a)
$$, on voit que $\Coker(\delta) = M \otimes_R S$. On obtient une application canonique $$
\Omega^{t + 1}_{S/\Lambda} \to
\wedge_S^t(\Coker(\delta)) \otimes_S \Omega_{S/R} =
\wedge^t_R(M) \otimes_R \Omega_{S/R}
$$ Or l'image de l'application $\text{Tr} : \Omega_{S/R} \to \Omega_{R/\Lambda}$ du lemme \ref{resolve-lemma-trace-well-defined} est contenue dans $R\text{d}a$ ; le produit extérieur par un élément de cette image annule donc $\text{d}a$. Il existe par conséquent une application canonique $$
\wedge^t_R(M) \otimes_R \Omega_{S/R} \to \Omega_{R/\Lambda}^{t + 1}
$$ envoyant $\overline{\eta}_1 \wedge \ldots \wedge \overline{\eta}_t \wedge \omega$ sur $\eta_1 \wedge \ldots \wedge \eta_t \wedge \text{Tr}(\omega)$. \end{proof}

\begin{remark} \label{resolve-remark-compare-Garel} Soient $\mathbf{F}_p \subset \Lambda \subset R \subset S$ et $\text{Tr}$ comme dans le lemme \ref{resolve-lemma-trace-higher}. D'après Cohomologie de de Rham, proposition \ref{derham-proposition-Garel}, il existe un morphisme canonique de complexes $$
\Theta_{S/R} :
\Omega_{S/\Lambda}^\bullet
\longrightarrow
\Omega_{R/\Lambda}^\bullet
$$ Le calcul de Cohomologie de de Rham, exemple \ref{derham-example-Garel}, montre que $\Theta_{S/R}(x^i \text{d}x) = \text{Tr}_x(x^i\text{d}x)$ pour tout $i$. Puisque $\text{Trace}_{S/R} = \Theta^0_{S/R}$ est identiquement nul et que $$
\Theta_{S/R}(a \wedge b) = a \wedge \Theta_{S/R}(b)
$$ pour $a \in \Omega^i_{R/\Lambda}$ et $b \in \Omega^j_{S/\Lambda}$, on en déduit $\text{Tr} = \Theta_{S/R}$. L'avantage de $\text{Tr}$ est que sa construction est nettement plus élémentaire. \end{remark}

\begin{lemma} \label{resolve-lemma-trace-extends} Soit $S$ un schéma sur $\mathbf{F}_p$. Soit $f : Y \to X$ un morphisme fini de schémas noethériens normaux intègres sur $S$. Supposons que \begin{enumerate} \item l'extension des corps de fonctions soit purement inséparable de degré $p$, et que \item $\Omega_{X/S}$ soit un $\mathcal{O}_X$-module cohérent (par exemple lorsque $X$ est de type fini sur $S$). \end{enumerate} Pour $i \geq 1$, il existe une application canonique $$
\text{Tr} : f_*\Omega^i_{Y/S} \longrightarrow (\Omega_{X/S}^i)^{**}
$$ dont la fibre au point générique de $X$ est l'application trace du lemme \ref{resolve-lemma-trace-higher}. \end{lemma}

\begin{proof} La suite exacte $f^*\Omega_{X/S} \to \Omega_{Y/S} \to \Omega_{Y/X} \to 0$ montre que $\Omega_{Y/S}$, et donc $f_*\Omega_{Y/S}$, sont également des modules cohérents. Il suffit donc de prouver que l'application trace au point générique se prolonge aux fibres en $x \in X$ tels que $\dim(\mathcal{O}_{X, x}) = 1$ ; voir Diviseurs, lemme \ref{divisors-lemma-describe-reflexive-hull}. On se ramène ainsi au cas traité dans le paragraphe suivant.

\medskip\noindent Supposons que $X = \Spec(A)$ et $Y = \Spec(B)$, où $A$ est un anneau de valuation discrète et $B$ est fini sur $A$. Puisque l'extension induite $L/K$ des corps des fractions est purement inséparable, $B$ est lui aussi local. Ainsi $B$ est également un anneau de valuation discrète. Alors, ou bien \begin{enumerate} \item $B/A$ a pour indice de ramification $p$, et donc $B = A[x]/(x^p - a)$, où $a \in A$ est une uniformisante ; ou bien \item $\mathfrak m_B = \mathfrak m_A B$ et le corps résiduel $B/\mathfrak m_A B$ est purement inséparable de degré $p$ sur $\kappa_A = A/\mathfrak m_A$. Choisissons un élément quelconque $x \in B$ dont la classe résiduelle n'appartienne pas à $\kappa_A$ ; on aura alors $B = A[x]/(x^p - a)$, où $a \in A$ est inversible. \end{enumerate} Soit $\Spec(\Lambda) \subset S$ un ouvert affine tel que $X$ ait son image dans $\Spec(\Lambda)$. On peut alors appliquer le lemme \ref{resolve-lemma-trace-higher} pour voir que l'application trace se prolonge en $\Omega^i_{B/\Lambda} \to \Omega^i_{A/\Lambda}$ pour tout $i \geq 1$. \end{proof}


















\section{Transformations quadratiques} \label{resolve-section-quadratic}

\noindent Dans cette section, nous étudions ce qui se produit lorsqu'on éclate un point non singulier d'une surface. Nous hésitons à définir formellement un tel morphisme comme une {\it transformation quadratique}, car, d'une part, on emploie souvent d'autres noms et, d'autre part, l'expression ``transformation quadratique'' est parfois utilisée dans un autre sens.

\begin{lemma} \label{resolve-lemma-blowup} Soit $(A, \mathfrak m, \kappa)$ un anneau local régulier de dimension $2$. Soit $f : X \to S = \Spec(A)$ l'éclatement de $A$ en $\mathfrak m$, de diviseur exceptionnel $E$. Il existe une immersion fermée $$
r : X \longrightarrow \mathbf{P}^1_S
$$ au-dessus de $S$ telle que \begin{enumerate} \item $r|_E : E \to \mathbf{P}^1_\kappa$ soit un isomorphisme, \item $\mathcal{O}_X(E) = \mathcal{O}_X(-1) =
r^*\mathcal{O}_{\mathbf{P}^1}(-1)$, et \item $\mathcal{C}_{E/X} = (r|_E)^*\mathcal{O}_{\mathbf{P}^1}(1)$ et $\mathcal{N}_{E/X} = (r|_E)^*\mathcal{O}_{\mathbf{P}^1}(-1)$. \end{enumerate} \end{lemma}

\begin{proof} Puisque $A$ est régulier de dimension $2$, on peut écrire $\mathfrak m = (x, y)$. Alors $x$ et $y$, placés en degré $1$, engendrent l'algèbre de Rees $\bigoplus_{n \geq 0} \mathfrak m^n$ sur $A$. Rappelons que $X = \text{Proj}(\bigoplus_{n \geq 0} \mathfrak m^n)$ ; voir Diviseurs, lemme \ref{divisors-lemma-blowing-up-affine}. La surjection $$
A[T_0, T_1] \longrightarrow \bigoplus\nolimits_{n \geq 0} \mathfrak m^n,
\quad
T_0 \mapsto x,\ T_1 \mapsto y
$$
\begingroup\emergencystretch=2em
de $A$-algèbres graduées induit donc une immersion fermée $r : X \to \mathbf{P}^1_S = \text{Proj}(A[T_0, T_1])$ telle que $\mathcal{O}_X(1) = r^*\mathcal{O}_{\mathbf{P}^1_S}(1)$ ; voir Constructions, lemme \ref{constructions-lemma-surjective-graded-rings-generated-degree-1-map-proj}. Cela prouve (2), car $\mathcal{O}_X(E) = \mathcal{O}_X(-1)$ d'après Diviseurs, lemme \ref{divisors-lemma-blowing-up-gives-effective-Cartier-divisor}.\par\endgroup

\medskip\noindent Pour prouver (1), remarquons que $$
\left(\bigoplus\nolimits_{n \geq 0} \mathfrak m^n\right) \otimes_A \kappa =
\bigoplus\nolimits_{n \geq 0} \mathfrak m^n/\mathfrak m^{n + 1} \cong
\kappa[\overline{x}, \overline{y}]
$$ est une algèbre de polynômes ; voir Algèbre, lemme \ref{algebra-lemma-regular-graded}. Cela prouve que la fibre de $X \to S$ au-dessus de $\Spec(\kappa)$ est égale à $\text{Proj}(\kappa[\overline{x}, \overline{y}]) = \mathbf{P}^1_\kappa$ ; voir Constructions, lemme \ref{constructions-lemma-base-change-map-proj}. Rappelons que $E$ est le sous-schéma fermé de $X$ défini par $\mathfrak m\mathcal{O}_X$, c'est-à-dire $E = X_\kappa$. Par le choix du morphisme $r$, on voit que $r|_E$ réalise en fait l'identification de $E = X_\kappa$ avec la fibre spéciale de $\mathbf{P}^1_S \to S$.

\medskip\noindent L'assertion (3) résulte de (1), de (2) et de Diviseurs, lemme \ref{divisors-lemma-conormal-effective-Cartier-divisor}. \end{proof}

\begin{lemma} \label{resolve-lemma-blowup-regular} Soit $(A, \mathfrak m, \kappa)$ un anneau local régulier de dimension $2$. Soit $f : X \to S = \Spec(A)$ l'éclatement de $A$ en $\mathfrak m$. Alors $X$ est un schéma régulier irréductible. \end{lemma}

\begin{proof} Remarquons que $X$ est intègre d'après Diviseurs, lemme \ref{divisors-lemma-blow-up-integral-scheme}, et Algèbre, lemme \ref{algebra-lemma-regular-domain}. Pour voir que $X$ est régulier, il suffit de vérifier que $\mathcal{O}_{X, x}$ est régulier aux points fermés $x \in X$ ; voir Propriétés, lemme \ref{properties-lemma-characterize-regular}. Soit $x \in X$ un point fermé. Puisque $f$ est propre, $x$ a pour image $\mathfrak m$ ; autrement dit, $x$ appartient au diviseur exceptionnel $E$. Alors $E$ est un diviseur de Cartier effectif et $E \cong \mathbf{P}^1_\kappa$. Ainsi, si $g \in \mathfrak m_x \subset \mathcal{O}_{X, x}$ est une équation locale de $E$, on a $\mathcal{O}_{X, x}/(g) \cong \mathcal{O}_{\mathbf{P}^1_\kappa, x}$. Puisque $\mathbf{P}^1_\kappa$ est recouvert par deux ouverts affines qui sont les spectres d'anneaux de polynômes sur $\kappa$, l'anneau $\mathcal{O}_{\mathbf{P}^1_\kappa, x}$ est régulier d'après Algèbre, lemme \ref{algebra-lemma-dim-affine-space}. On conclut par Algèbre, lemme \ref{algebra-lemma-regular-mod-x}. \end{proof}

\begin{lemma} \label{resolve-lemma-blowup-pic} Soit $(A, \mathfrak m, \kappa)$ un anneau local régulier de dimension $2$. Soit $f : X \to S = \Spec(A)$ l'éclatement de $A$ en $\mathfrak m$. Alors $\Pic(X) = \mathbf{Z}$, avec pour générateur $\mathcal{O}_X(E)$. \end{lemma}

\begin{proof} Rappelons que $E = \mathbf{P}^1_\kappa$ a pour groupe de Picard $\mathbf{Z}$, avec pour générateur $\mathcal{O}(1)$ ; voir Diviseurs, lemme \ref{divisors-lemma-Pic-projective-space-UFD}. D'après le lemme \ref{resolve-lemma-blowup}, le $\mathcal{O}_X$-module inversible $\mathcal{O}_X(E)$ se restreint en $\mathcal{O}(-1)$. Ainsi $\mathcal{O}_X(E)$ engendre un groupe cyclique infini dans $\Pic(X)$. Puisque $A$ est régulier, il est factoriel ; voir Compléments d'algèbre, lemme \ref{more-algebra-lemma-regular-local-UFD}. Le spectre épointé $U = S \setminus \{\mathfrak m\} = X \setminus E$ a donc un groupe de Picard trivial ; voir Diviseurs, lemme \ref{divisors-lemma-open-subscheme-UFD}. Par conséquent, pour tout $\mathcal{O}_X$-module inversible $\mathcal{L}$, il existe un isomorphisme $s : \mathcal{O}_U \to \mathcal{L}|_U$. Alors $s$ est une section méromorphe régulière de $\mathcal{L}$ et l'on voit que $\text{div}_\mathcal{L}(s) = nE$ pour un certain $n \in \mathbf{Z}$ (Diviseurs, définition \ref{divisors-definition-divisor-invertible-sheaf}). D'après Diviseurs, lemme \ref{divisors-lemma-normal-c1-injective} (et la normalité de $X$ donnée par le lemme \ref{resolve-lemma-blowup-regular}), on conclut que $\mathcal{L} = \mathcal{O}_X(nE)$. \end{proof}

\begin{lemma} \label{resolve-lemma-cohomology-of-blowup} Soit $(A, \mathfrak m, \kappa)$ un anneau local régulier de dimension $2$. Soit $f : X \to S = \Spec(A)$ l'éclatement de $A$ en $\mathfrak m$. Soit $\mathcal{F}$ un $\mathcal{O}_X$-module quasi-cohérent. \begin{enumerate} \item $H^p(X, \mathcal{F}) = 0$ pour $p \not \in \{0, 1\}$ ; \item $H^1(X, \mathcal{O}_X(n)) = 0$ pour $n \geq -1$ ; \item $H^1(X, \mathcal{F}) = 0$ si $\mathcal{F}$ ou $\mathcal{F}(1)$ est engendré par ses sections globales ; \item $H^0(X, \mathcal{O}_X(n)) = \mathfrak m^{\max(0, n)}$ ; \item $\text{length}_A H^1(X, \mathcal{O}_X(n)) = -n(-n - 1)/2$ si $n < 0$. \end{enumerate} \end{lemma}

\begin{proof} Si $\mathfrak m = (x, y)$, alors $X$ est recouvert par les spectres des algèbres d'éclatement affines $A[\frac{\mathfrak m}{x}]$ et $A[\frac{\mathfrak m}{y}]$, car $x$ et $y$, placés en degré $1$, engendrent l'algèbre de Rees $\bigoplus \mathfrak m^n$ sur $A$. Voir Diviseurs, lemme \ref{divisors-lemma-blowing-up-affine}, et Constructions, lemme \ref{constructions-lemma-proj-quasi-compact}. Puisque $X$ est séparé d'après Constructions, lemme \ref{constructions-lemma-proj-separated}, la cohomologie des faisceaux quasi-cohérents s'annule en degrés $\geq 2$ d'après Cohomologie des schémas, lemme \ref{coherent-lemma-vanishing-nr-affines}.

\medskip\noindent Soit $i : E \to X$ le diviseur exceptionnel ; voir Diviseurs, définition \ref{divisors-definition-blow-up}. Rappelons que $\mathcal{O}_X(-E) = \mathcal{O}_X(1)$ est relativement ample pour $f$ ; voir Diviseurs, lemme \ref{divisors-lemma-blowing-up-gives-effective-Cartier-divisor}. On sait donc que $H^1(X, \mathcal{O}_X(-nE)) = 0$ pour un certain $n > 0$ ; voir Cohomologie des schémas, lemme \ref{coherent-lemma-kill-by-twisting}. Considérons la filtration $$
\mathcal{O}_X(-nE) \subset \mathcal{O}_X(-(n - 1)E) \subset
\ldots \subset \mathcal{O}_X(-E) \subset \mathcal{O}_X \subset \mathcal{O}_X(E)
$$ Ses quotients successifs sont les faisceaux $$
\mathcal{O}_X(-t E)/\mathcal{O}_X(-(t + 1)E) =
\mathcal{O}_X(t)/\mathcal{I}(t) =
i_*\mathcal{O}_E(t)
$$ où $\mathcal{I} = \mathcal{O}_X(-E)$ est le faisceau d'idéaux de $E$. D'après le lemme \ref{resolve-lemma-blowup}, on a $E = \mathbf{P}^1_\kappa$, et $\mathcal{O}_E(1)$ correspond bien à la torsion de Serre usuelle du faisceau structural sur $\mathbf{P}^1$. La cohomologie de $\mathcal{O}_E(t)$ s'annule donc en degré $1$ pour $t \geq -1$ ; voir Cohomologie des schémas, lemme \ref{coherent-lemma-cohomology-projective-space-over-ring}. Comme elle est égale à $H^1(X, i_*\mathcal{O}_E(t))$ (d'après Cohomologie des schémas, lemme \ref{coherent-lemma-relative-affine-cohomology}), on en déduit que $H^1(X, \mathcal{O}_X(-(t + 1)E)) \to H^1(X, \mathcal{O}_X(-tE))$ est surjective pour $t \geq -1$. Ainsi $$
0 = H^1(X, \mathcal{O}_X(-nE))
\longrightarrow
H^1(X, \mathcal{O}_X(-tE)) = H^1(X, \mathcal{O}_X(t))
$$ est surjective pour $t \geq -1$, ce qui prouve (2).

\medskip\noindent Supposons $\mathcal{F}$ engendré par ses sections globales. Il existe donc une suite exacte courte $$
0 \to \mathcal{G} \to \bigoplus\nolimits_{i \in I} \mathcal{O}_X
\to \mathcal{F} \to 0
$$
\begingroup\emergencystretch=1em
Remarquons que $H^1(X, \bigoplus_{i \in I} \mathcal{O}_X) =
\bigoplus_{i \in I} H^1(X, \mathcal{O}_X)$ d'après Cohomologie, lemme \ref{cohomology-lemma-quasi-separated-cohomology-colimit}. Par (2), on a $H^1(X, \mathcal{O}_X) = 0$. Si $\mathcal{F}(1)$ est engendré par ses sections globales, on peut trouver une surjection $\bigoplus_{i \in I} \mathcal{O}_X(-1) \to \mathcal{F}$ et raisonner de même. Autrement dit, (3) résulte de (2).
\par\endgroup

\medskip\noindent Pour (4), remarquons que, pour tout $n$ assez grand, on a $\Gamma(X, \mathcal{O}_X(n)) = \mathfrak m^n$ ; voir Cohomologie des schémas, lemme \ref{coherent-lemma-recover-tail-graded-module}. Si $n \geq 0$, la suite exacte courte $$
0 \to \mathcal{O}_X(n) \to \mathcal{O}_X(n - 1) \to
i_*\mathcal{O}_E(n - 1) \to 0
$$ et l'annulation de $H^1$ pour le faisceau de gauche donnent un diagramme commutatif
\begingroup\small
$$
\xymatrix{
0 \ar[r] &
\mathfrak m^{\max(0, n)} \ar[r] \ar[d] &
\mathfrak m^{\max(0, n - 1)} \ar[r] \ar[d] &
\mathfrak m^{\max(0, n)}/\mathfrak m^{\max(0, n - 1)} \ar[r] \ar[d] & 0\\
0 \ar[r] &
\Gamma(X, \mathcal{O}_X(n)) \ar[r] &
\Gamma(X, \mathcal{O}_X(n - 1)) \ar[r] &
\Gamma(E, \mathcal{O}_E(n - 1)) \ar[r] & 0
}
$$
\endgroup
à lignes exactes. En fait, les lignes sont également exactes pour $n < 0$, car les groupes de droite sont alors nuls. Dans la démonstration du lemme \ref{resolve-lemma-blowup}, nous avons vu que la flèche verticale de droite est un isomorphisme (détails omis). Si celle de gauche est un isomorphisme, il en est donc de même de celle du milieu. On obtient ainsi (4) par récurrence descendante sur $n$.

\medskip\noindent Enfin, on prouve (5) par récurrence descendante sur $n$ à l'aide des suites $$
0 \to \mathcal{O}_X(n) \to \mathcal{O}_X(n - 1) \to
i_*\mathcal{O}_E(n - 1) \to 0
$$ En effet, pour $n \geq -1$, on sait déjà que $H^1(X, \mathcal{O}_X(n)) = 0$. Puisque $$
H^1(X, i_*\mathcal{O}_E(-2)) =
H^1(E, \mathcal{O}_E(-2)) =
H^1(\mathbf{P}^1_\kappa, \mathcal{O}(-2)) \cong \kappa
$$ d'après Cohomologie des schémas, lemme \ref{coherent-lemma-cohomology-projective-space-over-ring}, et que ce module est de longueur $1$ sur $A$, la suite exacte longue de cohomologie montre que (5) est vrai pour $n = -2$. On poursuit de même. \end{proof}

\begin{lemma} \label{resolve-lemma-blowup-improve} Soit $(A, \mathfrak m)$ un anneau local régulier de dimension $2$. Soit $f : X \to S = \Spec(A)$ l'éclatement de $A$ en $\mathfrak m$. Soit $\mathfrak m^n \subset I \subset \mathfrak m$ un idéal. Soit $d \geq 0$ le plus grand entier tel que $$
I \mathcal{O}_X \subset \mathcal{O}_X(-dE)
$$, où $E$ est le diviseur exceptionnel. Posons $\mathcal{I}' = I\mathcal{O}_X(dE) \subset \mathcal{O}_X$. Alors $d > 0$, le faisceau $\mathcal{O}_X/\mathcal{I}'$ a pour support un ensemble fini de points fermés $x_1, \ldots, x_r$ de $X$, et \begin{align*}
\text{length}_A(A/I)
& >
\text{length}_A \Gamma(X, \mathcal{O}_X/\mathcal{I}') \\
& \geq
\sum\nolimits_{i = 1, \ldots, r}
\text{length}_{\mathcal{O}_{X, x_i}}
(\mathcal{O}_{X, x_i}/\mathcal{I}'_{x_i})
\end{align*} \end{lemma}

\begin{proof} Puisque $I \subset \mathfrak m$, tout élément de $I$ s'annule sur $E$. On a donc $d \geq 1$. D'autre part, puisque $\mathfrak m^n \subset I$, on a $d \leq n$. Considérons la suite exacte courte $$
0 \to I\mathcal{O}_X \to \mathcal{O}_X \to \mathcal{O}_X/I\mathcal{O}_X \to 0
$$ Puisque $I\mathcal{O}_X$ est engendré par ses sections globales, on a $H^1(X, I\mathcal{O}_X) = 0$ d'après le lemme \ref{resolve-lemma-cohomology-of-blowup}. On obtient donc une surjection $A/I \to \Gamma(X, \mathcal{O}_X/I\mathcal{O}_X)$. Considérons la suite exacte courte $$
0 \to
\mathcal{O}_X(-dE)/I\mathcal{O}_X \to
\mathcal{O}_X/I\mathcal{O}_X \to
\mathcal{O}_X/\mathcal{O}_X(-dE) \to 0
$$ D'après Diviseurs, lemme \ref{divisors-lemma-codim-1-part}, le faisceau $\mathcal{O}_X(-dE)/I\mathcal{O}_X$ a pour support un ensemble fini de points fermés de $X$. En particulier, ses groupes de cohomologie supérieure sont nuls (détail omis). Dans le diagramme suivant
\begingroup\footnotesize
$$
\xymatrix{
& & A/I \ar[d] \\
0 \ar[r] &
\Gamma(X, \mathcal{O}_X(-dE)/I\mathcal{O}_X) \ar[r] &
\Gamma(X, \mathcal{O}_X/I\mathcal{O}_X) \ar[r] &
\Gamma(X, \mathcal{O}_X/\mathcal{O}_X(-dE)) \ar[r] & 0
}
$$\endgroup, la ligne inférieure est donc exacte et la flèche verticale surjective. On a $$
\text{length}_A \Gamma(X, \mathcal{O}_X(-dE)/I\mathcal{O}_X) <
\text{length}_A(A/I)
$$, car $\Gamma(X, \mathcal{O}_X/\mathcal{O}_X(-dE))$ est non nul. En effet, l'image de $1 \in \Gamma(X, \mathcal{O}_X)$ est non nulle puisque $d > 0$.

\medskip\noindent Pour achever la démonstration, traduisons les résultats ci-dessus en les assertions du lemme. Puisque $\mathcal{O}_X(dE)$ est inversible, on a $$
\mathcal{O}_X/\mathcal{I}' =
\mathcal{O}_X(-dE)/I\mathcal{O}_X \otimes_{\mathcal{O}_X} \mathcal{O}_X(dE).
$$ Ainsi $\mathcal{O}_X/\mathcal{I}'$ et $\mathcal{O}_X(-dE)/I\mathcal{O}_X$ ont pour support le même ensemble fini de points fermés, notés $x_1, \ldots, x_r \in E \subset X$. De plus, on obtient $$
\Gamma(X, \mathcal{O}_X(-dE)/I\mathcal{O}_X) =
\bigoplus \mathcal{O}_X(-dE)_{x_i}/I\mathcal{O}_{X, x_i}
\cong
\bigoplus \mathcal{O}_{X, x_i}/\mathcal{I}'_{x_i} =
\Gamma(X, \mathcal{O}_X/\mathcal{I}')
$$, car un module inversible sur un anneau local est trivial. On en déduit l'inégalité stricte. On obtient également la seconde inégalité, car $$
\text{length}_A(\mathcal{O}_{X, x_i}/\mathcal{I}'_{x_i}) \geq
\text{length}_{\mathcal{O}_{X, x_i}}(\mathcal{O}_{X, x_i}/\mathcal{I}'_{x_i})
$$, comme il résulte immédiatement de la définition de la longueur. \end{proof}

\begin{lemma} \label{resolve-lemma-differentials-of-blowup} Soit $(A, \mathfrak m, \kappa)$ un anneau local régulier de dimension $2$. Soit $f : X \to S = \Spec(A)$ l'éclatement de $A$ en $\mathfrak m$. Alors $\Omega_{X/S} = i_*\Omega_{E/\kappa}$, où $i : E \to X$ est l'immersion du diviseur exceptionnel. \end{lemma}

\begin{proof} Écrivons $\mathbf{P}^1 = \mathbf{P}^1_S$ et soit $r : X \to \mathbf{P}^1$ comme dans le lemme \ref{resolve-lemma-blowup}. On a alors une suite exacte $$
\mathcal{C}_{X/\mathbf{P}^1} \to r^*\Omega_{\mathbf{P}^1/S} \to
\Omega_{X/S} \to 0
$$ ; voir Morphismes, lemme \ref{morphisms-lemma-differentials-relative-immersion}. Puisque $\Omega_{\mathbf{P}^1/S}|_E = \Omega_{E/\kappa}$ d'après Morphismes, lemme \ref{morphisms-lemma-base-change-differentials}, il suffit de voir que la première flèche définit une surjection sur le noyau de l'application canonique $r^*\Omega_{\mathbf{P}^1/S} \to i_*\Omega_{E/\kappa}$. Cela se vérifie localement. Avec les notations de la démonstration du lemme \ref{resolve-lemma-blowup}, sur un ouvert affine de $X$, le morphisme $f$ correspond à l'homomorphisme d'anneaux $$
A \to A[t]/(xt - y)
$$, où $x, y \in \mathfrak m$ sont des générateurs. Ainsi $\text{d}(xt - y) = x\text{d}t$ et $y\text{d}t = t \cdot x \text{d}t$, ce qui prouve l'assertion. \end{proof}



\section{Domination par des transformations quadratiques} \label{resolve-section-dominating-by-quadratic}

\noindent Le résultat précédent permet de prouver que les éclatements de points dominent toute modification d'un schéma régulier de dimension $2$.

\medskip\noindent Soit $X$ un schéma. Soit $x \in X$ un point fermé. Comme d'habitude, on considère $i : x = \Spec(\kappa(x)) \to X$ comme un sous-schéma fermé. L'{\it éclatement $X' \to X$ de $X$ au point $x$} est l'éclatement de $X$ suivant le sous-schéma fermé $x \subset X$. Remarquons que, si $X$ est localement noethérien, alors $X' \to X$ est projectif (donc propre) d'après Diviseurs, lemme \ref{divisors-lemma-blowing-up-projective}.

\begin{lemma} \label{resolve-lemma-make-ideal-principal} Soit $X$ un schéma noethérien. Soit $T \subset X$ un ensemble fini de points fermés $x$ tels que $\mathcal{O}_{X, x}$ soit régulier de dimension $2$ pour $x \in T$. Soit $\mathcal{I} \subset \mathcal{O}_X$ un faisceau quasi-cohérent d'idéaux tel que $\mathcal{O}_X/\mathcal{I}$ ait son support dans $T$. Il existe alors une suite $$
X_n \to X_{n - 1} \to \ldots \to X_1 \to X_0 = X
$$, où $X_{i + 1} \to X_i$ est l'éclatement de $X_i$ en un point fermé situé au-dessus d'un point de $T$, telle que $\mathcal{I}\mathcal{O}_{X_n}$ soit un faisceau d'idéaux inversible. \end{lemma}

\begin{proof} Écrivons $T = \{x_1, \ldots, x_r\}$. Notons $I_i$ la fibre de $\mathcal{I}$ en $x_i$. Posons $$
n_i = \text{length}_{\mathcal{O}_{X, x_i}}(\mathcal{O}_{X, x_i}/I_i)
$$ Cet entier est fini, car $\mathcal{O}_X/\mathcal{I}$ a son support dans $T$ et le support de $\mathcal{O}_{X, x_i}/I_i$ est donc égal à $\{\mathfrak m_{x_i}\}$ (voir Algèbre, lemme \ref{algebra-lemma-support-point}). Raisonnons par récurrence sur $\sum n_i$. Si $n_i = 0$ pour tout $i$, alors $\mathcal{I} = \mathcal{O}_X$ et la démonstration est terminée.

\medskip\noindent Supposons que $n_i > 0$. Soit $X' \to X$ l'éclatement de $X$ en $x_i$ (voir la discussion précédant le lemme). Puisque $\Spec(\mathcal{O}_{X, x_i}) \to X$ est plat, $X' \times_X \Spec(\mathcal{O}_{X, x_i})$ est l'éclatement de l'anneau $\mathcal{O}_{X, x_i}$ suivant son idéal maximal ; voir Diviseurs, lemme \ref{divisors-lemma-flat-base-change-blowing-up}. Le carré du diagramme commutatif $$
\xymatrix{
\text{Proj}(\bigoplus\nolimits_{d \geq 0} \mathfrak m_{x_i}^d) \ar[r] \ar[d] &
X' \ar[d] \\
\Spec(\mathcal{O}_{X, x_i}) \ar[r] & X
}
$$ est donc cartésien. Soient $E \subset X'$ et $E' \subset \text{Proj}(\bigoplus\nolimits_{d \geq 0} \mathfrak m_{x_i}^d)$ les diviseurs exceptionnels. Soit $d \geq 1$ l'entier donné par le lemme \ref{resolve-lemma-blowup-improve} pour l'idéal $\mathcal{I}_i \subset \mathcal{O}_{X, x_i}$. Puisque les flèches horizontales du diagramme sont plates, que $E' \to E$ est surjectif et que $E'$ est l'image réciproque de $E$, on voit que $$
\mathcal{I}\mathcal{O}_{X'} \subset \mathcal{O}_{X'}(-dE)
$$ (certains détails sont omis). Posons $\mathcal{I}' = \mathcal{I}\mathcal{O}_{X'}(dE) \subset \mathcal{O}_{X'}$. Alors $\mathcal{O}_{X'}/\mathcal{I}'$ a pour support un ensemble fini de points fermés $T' \subset |X'|$, car c'est le cas au-dessus de $X \setminus \{x_i\}$ et après image réciproque sur $\text{Proj}(\bigoplus\nolimits_{d \geq 0} \mathfrak m_{x_i}^d)$. La dernière assertion du lemme \ref{resolve-lemma-blowup-improve} indique que la somme des longueurs des fibres $\mathcal{O}_{X', x'}/\mathcal{I}'\mathcal{O}_{X', x'}$, pour $x'$ situé au-dessus de $x_i$, est $< n_i$. La somme des longueurs a donc diminué.

\medskip\noindent Par l'hypothèse de récurrence, il existe une suite $$
X'_n \to \ldots \to X'_1 \to X'
$$ d'éclatements en des points fermés situés au-dessus de $T'$ telle que $\mathcal{I}'\mathcal{O}_{X'_n}$ soit inversible. Puisque $\mathcal{I}'\mathcal{O}_{X'}(-dE) = \mathcal{I}\mathcal{O}_{X'}$, on voit que $\mathcal{I}\mathcal{O}_{X'_n} =
\mathcal{I}'\mathcal{O}_{X'_n}(-d(f')^{-1}E)$, où $f' : X'_n \to X'$ est le composé. Remarquons que $(f')^{-1}E$ est un diviseur de Cartier effectif d'après Diviseurs, lemme \ref{divisors-lemma-blow-up-pullback-effective-Cartier}. On conclut donc par Diviseurs, lemme \ref{divisors-lemma-sum-effective-Cartier-divisors}. \end{proof}

\begin{lemma} \label{resolve-lemma-dominate-by-blowing-up-in-points} Soit $X$ un schéma noethérien. Soit $T \subset X$ un ensemble fini de points fermés $x$ tels que $\mathcal{O}_{X, x}$ soit un anneau local régulier de dimension $2$. Soit $f : Y \to X$ un morphisme propre de schémas qui soit un isomorphisme au-dessus de $U = X \setminus T$. Il existe alors une suite $$
X_n \to X_{n - 1} \to \ldots \to X_1 \to X_0 = X
$$, où $X_{i + 1} \to X_i$ est l'éclatement de $X_i$ en un point fermé $x_i$ situé au-dessus d'un point de $T$, ainsi qu'une factorisation $X_n \to Y \to X$ du composé. \end{lemma}

\begin{proof} D'après Compléments de platitude, lemme \ref{flat-lemma-dominate-modification-by-blowup}, il existe un éclatement $U$-admissible $X' \to X$ qui domine $Y \to X$. On peut donc supposer qu'il existe un faisceau d'idéaux $\mathcal{I} \subset \mathcal{O}_X$ tel que $\mathcal{O}_X/\mathcal{I}$ ait son support dans $T$ et que $Y$ soit l'éclatement de $X$ suivant $\mathcal{I}$. D'après le lemme \ref{resolve-lemma-make-ideal-principal}, il existe une suite $$
X_n \to X_{n - 1} \to \ldots \to X_1 \to X_0 = X
$$, où $X_{i + 1} \to X_i$ est l'éclatement de $X_i$ en un point fermé $x_i$ situé au-dessus d'un point de $T$, telle que $\mathcal{I}\mathcal{O}_{X_n}$ soit un faisceau d'idéaux inversible. La propriété universelle de l'éclatement (Diviseurs, lemme \ref{divisors-lemma-universal-property-blowing-up}) fournit la factorisation cherchée. \end{proof}

\begin{lemma} \label{resolve-lemma-extend-rational-map-blowing-up} Soit $S$ un schéma. Soit $X$ un schéma sur $S$, régulier et de dimension $2$. Soit $Y$ un schéma propre sur $S$. Étant donnée une application $S$-rationnelle $f : U \to Y$ de $X$ vers $Y$, il existe une suite $$
X_n \to X_{n - 1} \to \ldots \to X_1 \to X_0 = X
$$ et un $S$-morphisme $f_n : X_n \to Y$ tels que $X_{i + 1} \to X_i$ soit l'éclatement de $X_i$ en un point fermé non situé au-dessus de $U$ et que $f_n$ et $f$ coïncident. \end{lemma}

\begin{proof} On peut supposer que $U$ contient tous les points de codimension $1$ ; voir Morphismes, lemme \ref{morphisms-lemma-extend-across}. Le complémentaire $T \subset X$ de $U$ est donc un ensemble fini de points fermés dont les anneaux locaux sont réguliers de dimension $2$. En appliquant Diviseurs, lemme \ref{divisors-lemma-extend-rational-map-after-modification}, on trouve un morphisme propre $p : X' \to X$ qui est un isomorphisme au-dessus de $U$ et un morphisme $f' : X' \to Y$ coïncidant avec $f$ au-dessus de $U$. Appliquons le lemme \ref{resolve-lemma-dominate-by-blowing-up-in-points} au morphisme $p : X' \to X$. Le composé $X_n \to X' \to Y$ est le morphisme cherché. \end{proof}






\section{Domination par des éclatements normalisés} \label{resolve-section-normalized-blowups}

\noindent Dans cette section, nous prouvons qu'une modification d'une surface peut être dominée par une suite d'éclatements normalisés de points.

\begin{definition} \label{resolve-definition-normalized-blowup} Soit $X$ un schéma dont tout ouvert quasi-compact possède un nombre fini de composantes irréductibles. Soit $x \in X$ un point fermé. L'{\it éclatement normalisé de $X$ en $x$} est le composé $X'' \to X' \to X$, où $X' \to X$ est l'éclatement de $X$ en $x$ et $X'' \to X'$ est la normalisation de $X'$. \end{definition}

\noindent La normalisation $X'' \to X'$ est ici définie, car le schéma $X'$ possède un recouvrement par des ouverts ayant un nombre fini de composantes irréductibles d'après Diviseurs, lemme \ref{divisors-lemma-blow-up-and-irreducible-components}. Voir Morphismes, définition \ref{morphisms-definition-normalization}, pour la définition de la normalisation.

\medskip\noindent En général, l'éclatement normalisé n'est pas nécessairement propre, même si $X$ est noethérien. Rappelons qu'un schéma est de Nagata s'il possède un recouvrement par des ouverts affines qui sont des spectres d'anneaux de Nagata (Propriétés, définition \ref{properties-definition-nagata}).

\begin{lemma} \label{resolve-lemma-Nagata-normalized-blowup} Dans la définition \ref{resolve-definition-normalized-blowup}, si $X$ est de Nagata, alors l'éclatement normalisé de $X$ en $x$ est normal, de Nagata et propre sur $X$. \end{lemma}

\begin{proof} Le morphisme d'éclatement $X' \to X$ est propre (puisque $X$ est localement noethérien, on peut appliquer Diviseurs, lemme \ref{divisors-lemma-blowing-up-projective}). Ainsi $X'$ est de Nagata (Morphismes, lemme \ref{morphisms-lemma-finite-type-nagata}). La normalisation $X'' \to X'$ est donc finie (Morphismes, lemme \ref{morphisms-lemma-nagata-normalization}), d'où la propreté de $X'' \to X$ (Morphismes, lemmes \ref{morphisms-lemma-finite-proper} et \ref{morphisms-lemma-composition-proper}). Il en résulte que l'éclatement normalisé est un espace algébrique normal (Morphismes, lemme \ref{morphisms-lemma-normalization-normal}) et de Nagata. \end{proof}

\noindent Dans le lemme suivant, il faut supposer que $X$ est noethérien pour assurer qu'il possède un nombre fini de composantes irréductibles. La propreté de $f : Y \to X$ assure alors que $Y$ possède lui aussi un nombre fini de composantes irréductibles ; on peut donc demander que $f$ soit birationnel (Morphismes, définition \ref{morphisms-definition-birational}).

\begin{lemma} \label{resolve-lemma-dominate-by-normalized-blowing-up} Soit $X$ un schéma noethérien de Nagata de dimension $2$. Soit $f : Y \to X$ un morphisme propre birationnel. Il existe alors un diagramme commutatif $$
\xymatrix{
X_n \ar[r] \ar[d] &
X_{n - 1} \ar[r] &
\ldots \ar[r] &
X_1 \ar[r] &
X_0 \ar[d] \\
Y \ar[rrrr]  & & & & X
}
$$, où $X_0 \to X$ est la normalisation et où $X_{i + 1} \to X_i$ est l'éclatement normalisé de $X_i$ en un point fermé. \end{lemma}

\begin{proof} Nous utiliserons sans autre mention les résultats de Morphismes, sections \ref{morphisms-section-nagata}, \ref{morphisms-section-dimension-formula} et \ref{morphisms-section-normalization}. On peut remplacer $Y$ par sa normalisation. Soit $X_0 \to X$ la normalisation. Le morphisme $Y \to X$ se factorise par $X_0$. On peut donc supposer que $X$ et $Y$ sont tous deux normaux.

\medskip\noindent Supposons $X$ et $Y$ normaux. Le morphisme $f : Y \to X$ est un isomorphisme au-dessus d'un ouvert contenant tous les points de codimension $0$ et $1$ de $Y$, ainsi que tout point de $Y$ au-dessus duquel la fibre est finie ; voir Variétés, lemme \ref{varieties-lemma-modification-normal-iso-over-codimension-1}. Il existe donc un ensemble fini de points fermés $T \subset X$ tel que $f$ soit un isomorphisme au-dessus de $X \setminus T$. Pour tout $x \in T$, la fibre $Y_x$ est un schéma propre géométriquement connexe de dimension $1$ sur $\kappa(x)$ ; voir Compléments sur les morphismes, lemme \ref{more-morphisms-lemma-geometrically-connected-fibres-towards-normal}. Ainsi $$
BadCurves(f) = \{C \subset Y\text{ fermé} \mid
\dim(C) = 1, f(C) = \text{un point}\}
$$ est un ensemble fini. Nous prouverons le lemme par récurrence sur le nombre d'éléments de $BadCurves(f)$. Le cas initial est celui où $BadCurves(f)$ est vide ; alors $f$ est un isomorphisme.

\medskip\noindent Fixons $x \in T$. Soit $X' \to X$ l'éclatement normalisé de $X$ en $x$ et soit $Y'$ la normalisation de $Y \times_X X'$. On a le diagramme $$
\xymatrix{
Y' \ar[r]_{f'} \ar[d] & X' \ar[d] \\
Y \ar[r]^f & X
}
$$ Soit $x' \in X'$ un point fermé situé au-dessus de $x$ tel que la fibre $Y'_{x'}$ soit de dimension $\geq 1$. Soit $C' \subset Y'$ une composante irréductible de $Y'_{x'}$, c'est-à-dire $C' \in BadCurves(f')$. Puisque $Y' \to Y \times_X X'$ est fini, $C'$ a nécessairement pour image une composante irréductible $C \subset Y_x$. Il est clair que $C \in BadCurves(f)$. Puisque $Y' \to Y$ est birationnel, donc un isomorphisme au-dessus des points de codimension $1$ de $Y$, on obtient une application injective $$
BadCurves(f') \longrightarrow BadCurves(f)
$$ Il suffit donc de montrer qu'après un nombre fini de ces éclatements normalisés, on élimine au moins une des courbes exceptionnelles considérées, c'est-à-dire que l'application affichée n'est pas surjective.

\medskip\noindent Nous éliminerons une telle courbe au moyen d'un argument dû à Zariski. Choisissons $C \in BadCurves(f)$ situé au-dessus du point $x$ fixé. Notons $\mathcal{O}_{Y, C}$ l'anneau local de $Y$ au point générique de $C$. Choisissons un élément $u \in \mathcal{O}_{X, C}$ dont l'image dans le corps résiduel $R(C)$ soit transcendante sur $\kappa(x)$ (cela est possible, car $R(C)$ est de degré de transcendance $1$ sur $\kappa(x)$ d'après Variétés, lemme \ref{varieties-lemma-dimension-locally-algebraic}). On peut écrire $u = a/b$ avec $a, b \in \mathcal{O}_{X, x}$, car $\mathcal{O}_{Y, C}$ et $\mathcal{O}_{X, x}$ ont le même corps des fractions. Le choix de $u$ impose $a, b \in \mathfrak m_x$. Ainsi $$
N_{u, a, b} = \min
\{\text{ord}_{\mathcal{O}_{Y, C}}(a), \text{ord}_{\mathcal{O}_{Y, C}}(b)\} > 0
$$ On peut donc raisonner par récurrence descendante sur cet entier. Soit $X' \to X$ l'éclatement normalisé de $x$ et soit $Y'$ la normalisation de $X' \times_X Y$, comme ci-dessus. Nous montrerons que, si $C$ est l'image d'une courbe exceptionnelle $C' \subset Y'$ située au-dessus de $x' \in X'$, il existe un choix de $a', b' \mathcal{O}_{X', x'}$ tel que $N_{u, a', b'} < N_{u, a, b}$. Cela achèvera la démonstration. En effet, puisque $X' \to X$ se factorise par l'éclatement, il existe un élément non nul $d \in \mathfrak m_{x'}$ tel que $a = a' d$ et $b = b' d$ (prendre pour $d$ une équation locale du diviseur exceptionnel de l'éclatement). Puisque $Y' \to Y$ est un isomorphisme au-dessus d'un ouvert contenant le point générique de $C$ (comme on l'a vu), on a $\mathcal{O}_{Y', C'} = \mathcal{O}_{Y, C}$. Donc $$
\text{ord}_{\mathcal{O}_{Y, C}}(a) =
\text{ord}_{\mathcal{O}_{Y', C'}}(a' d) =
\text{ord}_{\mathcal{O}_{Y', C'}}(a') +
\text{ord}_{\mathcal{O}_{Y', C'}}(d) >
\text{ord}_{\mathcal{O}_{Y', C'}}(a')
$$ On raisonne de même pour $b$, et la démonstration est terminée. \end{proof}

\begin{lemma} \label{resolve-lemma-extend-rational-map-normalized-blowing-up} Soit $S$ un schéma. Soit $X$ un schéma sur $S$, noethérien, de Nagata et de dimension $2$. Soit $Y$ un schéma propre sur $S$. Étant donnée une application $S$-rationnelle $f : U \to Y$ de $X$ vers $Y$, il existe une suite $$
X_n \to X_{n - 1} \to \ldots \to X_1 \to X_0 \to X
$$ et un $S$-morphisme $f_n : X_n \to Y$ tels que $X_0 \to X$ soit la normalisation, que $X_{i + 1} \to X_i$ soit l'éclatement normalisé de $X_i$ en un point fermé et que $f_n$ et $f$ coïncident. \end{lemma}

\begin{proof} En appliquant Diviseurs, lemme \ref{divisors-lemma-extend-rational-map-after-modification}, on trouve un morphisme propre $p : X' \to X$ qui est un isomorphisme au-dessus de $U$ et un morphisme $f' : X' \to Y$ coïncidant avec $f$ au-dessus de $U$. Appliquons le lemme \ref{resolve-lemma-dominate-by-normalized-blowing-up} au morphisme $p : X' \to X$. Le composé $X_n \to X' \to Y$ est le morphisme cherché. \end{proof}





\section{Modifications au-dessus d'anneaux locaux} \label{resolve-section-modifications}

\noindent Soit $S$ un schéma. Soient $s_1, \ldots, s_n \in S$ des points fermés deux à deux distincts. Supposons que l'immersion ouverte $$
U = S \setminus \{s_1, \ldots, s_n\} \longrightarrow S
$$ soit quasi-compacte. Notons $FP_{S, \{s_1, \ldots, s_n\}}$ la catégorie des morphismes $f : X \to S$ de présentation finie qui induisent un isomorphisme $f^{-1}(U) \to U$. Les morphismes sont les morphismes de schémas au-dessus de $S$. Pour tout $i$, posons $S_i = \Spec(\mathcal{O}_{S, s_i})$ et soit $V_i = S_i \setminus \{s_i\}$. Notons $FP_{S_i, s_i}$ la catégorie des morphismes $g_i : Y_i \to S_i$ de présentation finie qui induisent un isomorphisme $g_i^{-1}(V_i) \to V_i$. Les morphismes sont les morphismes au-dessus de $S_i$. Le changement de base définit un foncteur \begin{equation}
\label{resolve-equation-equivalence}
F :
FP_{S, \{s_1, \ldots, s_n\}}
\longrightarrow
FP_{S_1, s_1} \times \ldots \times FP_{S_n, s_n}
\end{equation} Pour ramener au moins certains des problèmes de ce chapitre au cas des anneaux locaux, on dispose du lemme suivant.

\begin{lemma} \label{resolve-lemma-equivalence} Le foncteur $F$ (\ref{resolve-equation-equivalence}) est une équivalence. \end{lemma}

\begin{proof} Pour $n = 1$, c'est Limites, lemme \ref{limits-lemma-modifications}. Pour $n > 1$, on peut démontrer le lemme exactement de la même manière, ou le déduire de ce cas. Par exemple, supposons que $g_i : Y_i \to S_i$ soient des objets de $FP_{S_i, s_i}$. Le cas $n = 1$ fournit des morphismes $f'_i : X'_i \to S$ de présentation finie, isomorphes au-dessus de $S \setminus \{s_i\}$, dont le changement de base à $S_i$ est $g_i$. On peut alors poser $$
f : X = X'_1 \times_S \ldots \times_S X'_n \to S
$$ C'est un objet de $FP_{S, \{s_1, \ldots, s_n\}}$ dont le changement de base par $S_i \to S$ redonne $g_i$. Le foncteur est donc essentiellement surjectif. Nous omettons la preuve de sa pleine fidélité. \end{proof}

\begin{lemma} \label{resolve-lemma-equivalence-properties} Soit $S, s_i, S_i$ comme dans (\ref{resolve-equation-equivalence}). Si $f : X \to S$ correspond à $g_i : Y_i \to S_i$ par $F$, alors $f$ est séparé, propre, fini, si et seulement si $g_i$ l'est pour $i = 1, \ldots, n$. \end{lemma}

\begin{proof} Cela résulte de Limites, lemme \ref{limits-lemma-modifications-properties}. \end{proof}

\begin{lemma} \label{resolve-lemma-equivalence-fibre} Soit $S, s_i, S_i$ comme dans (\ref{resolve-equation-equivalence}). Si $f : X \to S$ correspond à $g_i : Y_i \to S_i$ par $F$, alors $X_{s_i} \cong (Y_i)_{s_i}$ en tant que schémas sur $\kappa(s_i)$. \end{lemma}

\begin{proof} C'est clair. \end{proof}

\begin{lemma} \label{resolve-lemma-equivalence-sequence-blowups} Soit $S, s_i, S_i$ comme dans (\ref{resolve-equation-equivalence}) ; supposons que $f : X \to S$ corresponde à $g_i : Y_i \to S_i$ par $F$. Il existe alors une factorisation $$
X = Z_m \to Z_{m - 1} \to \ldots \to Z_1 \to Z_0 = S
$$ de $f$, où $Z_{j + 1} \to Z_j$ est l'éclatement de $Z_j$ en un point fermé $z_j$ situé au-dessus de $\{s_1, \ldots, s_n\}$, si et seulement si, pour tout $i$, il existe une factorisation $$
Y_i = Z_{i, m_i} \to Z_{i, m_i - 1} \to \ldots \to Z_{i, 1} \to Z_{i, 0} = S_i
$$ de $g_i$, où $Z_{i, j + 1} \to Z_{i, j}$ est l'éclatement de $Z_{i, j}$ en un point fermé $z_{i, j}$ situé au-dessus de $s_i$. \end{lemma}

\begin{proof} Partons d'une suite d'éclatements $Z_m \to Z_{m - 1} \to \ldots \to Z_1 \to Z_0 = S$. Le premier morphisme $Z_1 \to S$ est obtenu en éclatant l'un des $s_i$, disons $s_1$. En appliquant $F$ à $Z_1 \to S$, on obtient l'éclatement $Z_{1, 1} \to S_1$ en $s_1$ ; il s'agit bien de l'éclatement en $s_1$, tandis que $Z_{i, 0} = S_i$ pour $i > 1$. À l'étape suivante, soit on éclate l'un des $s_i$, $i \geq 2$, sur $Z_1$, soit on choisit un point fermé $z_1$ de la fibre de $Z_1 \to S$ au-dessus de $s_1$. Dans le premier cas, la construction est claire ; dans le second, on utilise $(Z_1)_{s_1} \cong (Z_{1, 1})_{s_1}$ (lemme \ref{resolve-lemma-equivalence-fibre}) pour obtenir un point fermé $z_{1, 1} \in Z_{1, 1}$ correspondant à $z_1$. On prend alors pour $Z_{1, 2} \to Z_{1, 1}$ l'éclatement en $z_{1, 1}$. En poursuivant ainsi, on construit les factorisations de chaque $g_i$.

\medskip\noindent Réciproquement, étant données des suites d'éclatements $Z_{i, m_i} \to Z_{i, m_i - 1} \to \ldots \to Z_{i, 1} \to Z_{i, 0} = S_i$, on construit la suite d'éclatements de $S$ exactement de la même manière. \end{proof}

\noindent Voici l'analogue du lemme \ref{resolve-lemma-equivalence-sequence-blowups} pour les éclatements normalisés.

\begin{lemma} \label{resolve-lemma-equivalence-sequence-normalized-blowups} Soit $S, s_i, S_i$ comme dans (\ref{resolve-equation-equivalence}) ; supposons que $f : X \to S$ corresponde à $g_i : Y_i \to S_i$ par $F$. Supposons que tout ouvert quasi-compact de $S$ possède un nombre fini de composantes irréductibles. Il existe alors une factorisation $$
X = Z_m \to Z_{m - 1} \to \ldots \to Z_1 \to Z_0 = S
$$ de $f$, où $Z_{j + 1} \to Z_j$ est l'éclatement normalisé de $Z_j$ en un point fermé $z_j$ situé au-dessus de $\{x_1, \ldots, x_n\}$, si et seulement si, pour tout $i$, il existe une factorisation $$
Y_i = Z_{i, m_i} \to Z_{i, m_i - 1} \to \ldots \to Z_{i, 1} \to Z_{i, 0} = S_i
$$ de $g_i$, où $Z_{i, j + 1} \to Z_{i, j}$ est l'éclatement normalisé de $Z_{i, j}$ en un point fermé $z_{i, j}$ situé au-dessus de $s_i$. \end{lemma}

\begin{proof} L'hypothèse sur $S$ sert à assurer, de proche en proche, que les schémas que l'on normalise possèdent localement un nombre fini de composantes irréductibles, de sorte que l'énoncé a un sens. Cela dit, le lemme se déduit exactement du même argument que celui utilisé pour démontrer le lemme \ref{resolve-lemma-equivalence-sequence-blowups}. \end{proof}








\section{Annulation} \label{resolve-section-vanishing}

\noindent Dans cette section, nous travaillerons souvent dans le cadre suivant. Rappelons qu'une modification est un morphisme propre birationnel entre schémas intègres (Morphismes, définition \ref{morphisms-definition-modification}).

\begin{situation} \label{resolve-situation-vanishing} Ici, $(A, \mathfrak m, \kappa)$ est un anneau local noethérien normal intègre de dimension $2$. Soit $s$ le point fermé de $S = \Spec(A)$ et soit $U = S \setminus \{s\}$. Soit $f : X \to S$ une modification. Nous notons $C_1, \ldots, C_r$ les composantes irréductibles de la fibre spéciale $X_s$ de $f$. \end{situation}

\noindent D'après Variétés, lemme \ref{varieties-lemma-modification-normal-iso-over-codimension-1}, le morphisme $f$ définit un isomorphisme $f^{-1}(U) \to U$. La fibre spéciale $X_s$ est propre sur $\Spec(\kappa)$ et de dimension au plus $1$ d'après Variétés, lemme \ref{varieties-lemma-dimension-fibre-in-higher-codimension}. Par factorisation de Stein (Compléments sur les morphismes, lemme \ref{more-morphisms-lemma-geometrically-connected-fibres-towards-normal}), on a $f_*\mathcal{O}_X = \mathcal{O}_S$ et la fibre spéciale $X_s$ est géométriquement connexe sur $\kappa$. Si $X_s$ est de dimension $0$, alors $f$ est fini (Compléments sur les morphismes, lemme \ref{more-morphisms-lemma-proper-finite-fibre-finite-in-neighbourhood}), donc un isomorphisme (Morphismes, lemme \ref{morphisms-lemma-finite-birational-over-normal}). Nous écartons ce cas sans intérêt et concluons que $\dim(C_i) = 1$ pour $i = 1, \ldots, r$.

\begin{lemma} \label{resolve-lemma-dominate-by-scheme-modification} Dans la situation \ref{resolve-situation-vanishing}, il existe un éclatement $U$-admissible $X' \to S$ qui domine $X$. \end{lemma}

\begin{proof} C'est un cas particulier de Compléments sur la platitude, lemme \ref{flat-lemma-dominate-modification-by-blowup}. \end{proof}

\begin{lemma} \label{resolve-lemma-nice-meromorphic-function} Dans la situation \ref{resolve-situation-vanishing}, il existe un élément non nul $f \in \mathfrak m$ tel que, pour tout $i = 1, \ldots, r$, il existe \begin{enumerate} \item un point fermé $x_i \in C_i$ vérifiant $x_i \not \in C_j$ pour $j \not = i$, \item une factorisation $f = g_i f_i$ de $f$ dans $\mathcal{O}_{X, x_i}$ telle que $g_i \in \mathfrak m_{x_i}$ ait une image non nulle dans $\mathcal{O}_{C_i, x_i}$. \end{enumerate} \end{lemma}

\begin{proof} Nous utiliserons sans autre mention les observations faites après la situation \ref{resolve-situation-vanishing}. Choisissons un point fermé $x_i \in C_i$ qui n'appartienne à aucun $C_j$ pour $j \not = i$. Choisissons $g_i \in \mathfrak m_{x_i}$ dont l'image dans $\mathcal{O}_{C_i, x_i}$ soit non nulle. Puisque le corps des fractions de $A$ est celui de $\mathcal{O}_{X_i, x_i}$, on peut écrire $g_i = a_i/b_i$ pour certains $a_i, b_i \in A$. Prenons $f = \prod a_i$. \end{proof}

\begin{lemma} \label{resolve-lemma-nontrivial-normal-bundle} Dans la situation \ref{resolve-situation-vanishing}, supposons $X$ normal. Soit $Z \subset X$ un diviseur de Cartier effectif non vide tel que $Z \subset X_s$ ensemblistement. Alors le faisceau conormal de $Z$ n'est pas trivial. Plus précisément, il existe un $i$ tel que $C_i \subset Z$ et $\deg(\mathcal{C}_{Z/X}|_{C_i}) > 0$. \end{lemma}

\begin{proof} Nous utiliserons sans autre mention les observations faites après la situation \ref{resolve-situation-vanishing}. Soit $f$ une fonction comme dans le lemme \ref{resolve-lemma-nice-meromorphic-function}. Soit $\xi_i \in C_i$ le point générique. Soit $\mathcal{O}_i$ l'anneau local de $X$ en $\xi_i$. Alors $\mathcal{O}_i$ est un anneau de valuation discrète. Soit $e_i$ la valuation de $f$ dans $\mathcal{O}_i$, de sorte que $e_i > 0$. Soit $h_i \in \mathcal{O}_i$ une équation locale de $Z$ et soit $d_i$ sa valuation. Alors $d_i \geq 0$. Choisissons et fixons $i$ de sorte que $d_i/e_i$ soit maximal (alors $d_i > 0$ puisque $Z$ n'est pas vide). Remplaçons $f$ par $f^{d_i}$ et $Z$ par $e_iZ$. Cela est permis en vertu de la relation $\mathcal{O}_X(e_i Z) = \mathcal{O}_X(Z)^{\otimes e_i}$, de la relation entre le faisceau conormal et $\mathcal{O}_X(Z)$ (voir Diviseurs, lemmes \ref{divisors-lemma-invertible-sheaf-sum-effective-Cartier-divisors} et \ref{divisors-lemma-conormal-effective-Cartier-divisor}, et du fait que le degré est multiplié par $e_i$, voir Variétés, lemme \ref{varieties-lemma-degree-tensor-product}. Soit $\mathcal{I}$ le faisceau d'idéaux de $Z$, de sorte que $\mathcal{C}_{Z/X} = \mathcal{I}|_Z$. Considérons l'image $\overline{f}$ de $f$ dans $\Gamma(Z, \mathcal{O}_Z)$. D'après les choix effectués ci-dessus, $\overline{f}$ s'annule aux points génériques des composantes irréductibles de $Z$ (ce sont tous des points génériques de $C_j$ puisque $Z$ est contenu dans la fibre spéciale). D'autre part, $Z$ est $(S_1)$ d'après Diviseurs, lemme \ref{divisors-lemma-normal-effective-Cartier-divisor-S1}. Ainsi le schéma $Z$ n'a pas de point associé immergé et nous concluons que $\overline{f} = 0$ (Diviseurs, lemmes \ref{divisors-lemma-S1-no-embedded} et \ref{divisors-lemma-restriction-injective-open-contains-weakly-ass}). Par conséquent, $f$ est une section globale de $\mathcal{I}$ qui engendre $\mathcal{I}_{\xi_i}$ par construction. L'image $s_i$ de $f$ dans $\Gamma(C_i, \mathcal{I}|_{C_i})$ est donc non nulle. Toutefois, notre choix de $f$ garantit que $s_i$ possède un zéro en $x_i$. Le degré de $\mathcal{I}|_{C_i}$ est donc $> 0$ d'après Variétés, lemme \ref{varieties-lemma-check-invertible-sheaf-trivial}. \end{proof}

\begin{lemma} \label{resolve-lemma-H1-injective} Dans la situation \ref{resolve-situation-vanishing}, supposons $X$ normal et $A$ de Nagata. L'application $$
H^1(X, \mathcal{O}_X) \longrightarrow H^1(f^{-1}(U), \mathcal{O}_X)
$$ est injective. \end{lemma}

\begin{proof} Soit $0 \to \mathcal{O}_X \to \mathcal{E} \to \mathcal{O}_X \to 0$ l'extension correspondant à un élément non trivial $\xi$ de $H^1(X, \mathcal{O}_X)$ (Cohomologie, lemme \ref{cohomology-lemma-h1-extensions}). Soit $\pi : P = \mathbf{P}(\mathcal{E}) \to X$ le fibré projectif associé à $\mathcal{E}$. La surjection $\mathcal{E} \to \mathcal{O}_X$ définit une section $\sigma : X \to P$ dont le faisceau conormal est isomorphe à $\mathcal{O}_X$ (Diviseurs, lemme \ref{divisors-lemma-conormal-sheaf-section-projective-bundle}). Si la restriction de $\xi$ à $f^{-1}(U)$ est triviale, on obtient une application $\mathcal{E}|_{f^{-1}(U)} \to \mathcal{O}_{f^{-1}(U)}$ qui scinde l'injection $\mathcal{O}_X \to \mathcal{E}$. Cela définit une seconde section $\sigma' : f^{-1}(U) \to P$ disjointe de $\sigma$. Puisque $\xi$ est non trivial, nous concluons que $\sigma'$ ne peut pas se prolonger à tout $X$ tout en restant disjointe de $\sigma$. Soit $X' \subset P$ l'image schématique de $\sigma'$ (Morphismes, définition \ref{morphisms-definition-scheme-theoretic-image}). Considérons le diagramme $$
\xymatrix{
& X' \ar[rd]_g \ar[r] & P \ar[d]_\pi \\
f^{-1}(U) \ar[ru]_{\sigma'} \ar[rr] & & X \ar@/_/[u]_\sigma
}
$$ Le morphisme $P \setminus \sigma(X) \to X$ est affine. Si $X' \cap \sigma(X) = \emptyset$, alors $X' \to X$ est à la fois affine et propre, donc fini (Morphismes, lemme \ref{morphisms-lemma-finite-proper}), puis un isomorphisme (car $X$ est normal ; voir Morphismes, lemme \ref{morphisms-lemma-finite-birational-over-normal}). Cela est impossible, comme indiqué ci-dessus.

\medskip\noindent Soit $X^\nu$ la normalisation de $X'$. Puisque $A$ est de Nagata, $X^\nu \to X'$ est fini (Morphismes, lemmes \ref{morphisms-lemma-nagata-normalization} et \ref{morphisms-lemma-ubiquity-nagata}). Soit $Z \subset X^\nu$ l'image réciproque du diviseur de Cartier effectif $\sigma(X) \subset P$. D'après ce qui précède, $Z$ n'est pas vide et est contenu dans la fibre fermée de $X^\nu \to S$. Puisque $P \to X$ est lisse, $\sigma(X)$ est un diviseur de Cartier effectif (Diviseurs, lemme \ref{divisors-lemma-section-smooth-regular-immersion}). Par conséquent, $Z \subset X^\nu$ est lui aussi un diviseur de Cartier effectif. Puisque le faisceau conormal de $\sigma(X)$ dans $P$ est $\mathcal{O}_X$, le faisceau conormal de $Z$ dans $X^\nu$ (qui est a priori inversible) est $\mathcal{O}_Z$ d'après Morphismes, lemme \ref{morphisms-lemma-conormal-functorial-flat}. Cela contredit le lemme \ref{resolve-lemma-nontrivial-normal-bundle} et achève la démonstration. \end{proof}

\begin{lemma} \label{resolve-lemma-R1-injective} Dans la situation \ref{resolve-situation-vanishing}, supposons $X$ normal et $A$ de Nagata. Alors $$
\Hom_{D(A)}(\kappa[-1], Rf_*\mathcal{O}_X)
$$ est nul. On utilise ici $D(A) = D_\QCoh(\mathcal{O}_S)$ pour considérer $Rf_*\mathcal{O}_X$ comme un objet de $D(A)$. \end{lemma}

\begin{proof} Par adjonction entre $Rf_*$ et $Lf^*$, une telle application équivaut à une application $\alpha : Lf^*\kappa[-1] \to \mathcal{O}_X$. Remarquons que $$
H^i(Lf^*\kappa[-1]) =
\left\{
\begin{matrix}
0 & \text{si} & i > 1 \\
\mathcal{O}_{X_s} & \text{si} & i = 1 \\
\text{un }\mathcal{O}_{X_s}\text{-module} & \text{si} & i \leq 0
\end{matrix}
\right.
$$
\begingroup\emergencystretch=4em
Puisque $\Hom(H^0(Lf^*\kappa[-1]), \mathcal{O}_X) = 0$, car $\mathcal{O}_X$ est sans torsion, la suite spectrale pour $\Ext$ (Cohomologie sur les sites, exemple \ref{sites-cohomology-example-hom-complex-into-sheaf}) implique que $\Hom_{D(\mathcal{O}_X)}(Lf^*\kappa[-1], \mathcal{O}_X)$ est égal à $\Ext^1_{\mathcal{O}_X}(\mathcal{O}_{X_s}, \mathcal{O}_X)$. Nous concluons que $\alpha : Lf^*\kappa[-1] \to \mathcal{O}_X$ est donné par une extension\par\endgroup
$$
0 \to \mathcal{O}_X \to \mathcal{E} \to \mathcal{O}_{X_s} \to 0
$$ D'après le lemme \ref{resolve-lemma-H1-injective}, l'image réciproque de cette extension par la surjection $\mathcal{O}_X \to \mathcal{O}_{X_s}$ est nulle (car cette image réciproque est manifestement scindée sur $f^{-1}(U)$). Ainsi $1 \in \mathcal{O}_{X_s}$ se relève en une section globale $s$ de $\mathcal{E}$. En multipliant $s$ par le faisceau d'idéaux $\mathcal{I}$ de $X_s$, on obtient une application de $\mathcal{O}_X$-modules $c_s : \mathcal{I} \to \mathcal{O}_X$. En appliquant $f_*$, on obtient une application $A$-linéaire $f_*c_s : \mathfrak m \to A$. Comme $A$ est un anneau local noethérien normal intègre, cette application est donnée par la multiplication par un élément $a \in A$. En remplaçant $s$ par $s -  a$, on trouve que $s$ est annulé par $\mathcal{I}$ et que l'extension est triviale, comme souhaité. \end{proof}

\begin{remark} \label{resolve-remark-dualizing-setup} Soit $X$ un schéma intègre noethérien normal de dimension $2$. Dans ce cas, les propriétés suivantes sont équivalentes : \begin{enumerate} \item $X$ possède un complexe dualisant $\omega_X^\bullet$ ; \item il existe un $\mathcal{O}_X$-module cohérent $\omega_X$ tel que $\omega_X[n]$ soit un complexe dualisant, où $n$ peut être un entier quelconque. \end{enumerate} Cela résulte du fait que $X$ est de Cohen--Macaulay (Propriétés, lemme \ref{properties-lemma-normal-dimension-2-Cohen-Macaulay}) et de Dualité pour les schémas, lemme \ref{duality-lemma-dualizing-module-CM-scheme}. Dans cette situation, nous dirons que $\omega_X$ est un {\it module dualisant}, conformément à Dualité pour les schémas, section \ref{duality-section-dualizing-module}. En particulier, lorsque $A$ est un anneau local noethérien normal intègre de dimension $2$, nous dirons que {\it $A$ possède un module dualisant $\omega_A$} si la propriété ci-dessus est satisfaite. Dans ce cas, si $X \to \Spec(A)$ est une modification normale, alors $X$ possède aussi un module dualisant ; voir Dualité pour les schémas, exemple \ref{duality-example-proper-over-local}. Dans cette situation, nous notons toujours $\omega_X$ le module dualisant normalisé par rapport à $\omega_A$, c'est-à-dire tel que $\omega_X[2]$ soit le complexe dualisant normalisé relativement à $\omega_A[2]$. Voir Dualité pour les schémas, section \ref{duality-section-glue}. \end{remark}

\noindent Le théorème d'annulation de Grauert--Riemenschneider énoncé dans la proposition suivante est une conséquence formelle du lemme \ref{resolve-lemma-R1-injective} et de la théorie générale de la dualité.

\begin{proposition}[Grauert--Riemenschneider] \label{resolve-proposition-Grauert-Riemenschneider} Dans la situation \ref{resolve-situation-vanishing}, supposons que \begin{enumerate} \item $X$ soit un schéma normal, \item $A$ soit de Nagata et possède un complexe dualisant $\omega_A^\bullet$. \end{enumerate} Soit $\omega_X$ le module dualisant de $X$ (remarque \ref{resolve-remark-dualizing-setup}). Alors $R^1f_*\omega_X = 0$. \end{proposition}

\begin{proof} \begingroup\emergencystretch=3em
Dans cette démonstration, nous utiliserons l'identification $D(A) = D_\QCoh(\mathcal{O}_S)$ pour identifier les $\mathcal{O}_S$-modules quasi-cohérents aux $A$-modules. De plus, nous pouvons supposer $\omega_A^\bullet$ normalisé ; voir Complexes dualisants, section \ref{dualizing-section-dualizing-local}. Comme $X$ est un schéma noethérien normal de dimension $2$, il est de Cohen--Macaulay (Propriétés, lemme \ref{properties-lemma-normal-dimension-2-Cohen-Macaulay}). Ainsi $\omega_X^\bullet = \omega_X[2]$ (Dualité pour les schémas, lemme \ref{duality-lemma-dualizing-module-CM-scheme}, et la normalisation dans Dualité pour les schémas, exemple \ref{duality-example-proper-over-local}). Si la proposition est fausse, il existe une application non nulle $R^1f_*\omega_X \to \kappa$. Autrement dit, on obtient une application non nulle $\alpha : Rf_*\omega_X^\bullet \to \kappa[1]$. En appliquant $R\Hom_A(-, \omega_A^\bullet)$, on obtient une application non nulle\par\endgroup
$$
\beta : \kappa[-1] \longrightarrow Rf_*\mathcal{O}_X
$$, ce qui est impossible d'après le lemme \ref{resolve-lemma-R1-injective}. Pour voir que $R\Hom_A(-, \omega_A^\bullet)$ a bien l'effet annoncé, remarquons d'abord que $$
R\Hom_A(\kappa[1], \omega_A^\bullet) =
R\Hom_A(\kappa, \omega_A^\bullet)[-1] =
\kappa[-1]
$$ puisque $\omega_A^\bullet$ est normalisé et que l'on a $$
R\Hom_A(Rf_*\omega_X^\bullet, \omega_A^\bullet) =
Rf_*R\SheafHom_{\mathcal{O}_X}(\omega_X^\bullet, \omega_X^\bullet) =
Rf_*\mathcal{O}_X
$$ La première égalité résulte de Dualité pour les schémas, exemple \ref{duality-example-iso-on-RSheafHom-noetherian}, et du fait que $\omega_X^\bullet = f^!\omega_A^\bullet$ par construction ; la seconde, du fait que $\omega_X^\bullet$ est un complexe dualisant sur $X$ (ce qui remonte à Dualité pour les schémas, lemme \ref{duality-lemma-shriek-dualizing}). \end{proof}





\section{Bornitude} \label{resolve-section-bounded}

\noindent Dans cette section, nous entamons l'étude qui conduira à une réduction au cas des singularités rationnelles pour les schémas de dimension $2$.

\begin{lemma} \label{resolve-lemma-exact-sequence} Soit $(A, \mathfrak m, \kappa)$ un anneau local noethérien normal intègre de dimension $2$. Considérons un diagramme commutatif $$
\xymatrix{
X' \ar[rd]_{f'} \ar[rr]_g & & X \ar[ld]^f \\
& \Spec(A)
}
$$ où $f$ et $f'$ sont des modifications comme dans la situation \ref{resolve-situation-vanishing} et où $X$ est normal. On a alors une suite exacte courte $$
0 \to H^1(X, \mathcal{O}_X) \to H^1(X', \mathcal{O}_{X'}) \to
H^0(X, R^1g_*\mathcal{O}_{X'}) \to 0
$$ De plus, $\dim(\text{Supp}(R^1g_*\mathcal{O}_{X'})) = 0$ et $R^1g_*\mathcal{O}_{X'}$ est engendré par ses sections globales. \end{lemma}

\begin{proof} Nous utiliserons sans autre mention les observations faites après la situation \ref{resolve-situation-vanishing}. Comme $X$ est normal et $g$ est dominant et birationnel, on a $g_*\mathcal{O}_{X'} = \mathcal{O}_X$ ; voir par exemple Compléments sur les morphismes, lemme \ref{more-morphisms-lemma-geometrically-connected-fibres-towards-normal}. Puisque les fibres de $g$ sont de dimension $\leq 1$, on a $R^pg_*\mathcal{O}_{X'} = 0$ pour $p > 1$ ; voir par exemple Cohomologie des schémas, lemme \ref{coherent-lemma-higher-direct-images-zero-above-dimension-fibre}. Le support de $R^1g_*\mathcal{O}_{X'}$ est contenu dans l'ensemble des points de $X$ où les fibres de $g'$ sont de dimension $\geq 1$. Il est donc contenu dans l'ensemble des images des composantes irréductibles $C' \subset X'_s$ qui s'envoient sur des points de $X_s$ ; c'est un ensemble fini de points fermés (rappelons que $X'_s \to X_s$ est un morphisme entre schémas propres de dimension $1$ sur $\kappa$). Alors $R^1g_*\mathcal{O}_{X'}$ est engendré par ses sections globales d'après Cohomologie des schémas, lemme \ref{coherent-lemma-coherent-support-dimension-0}. En utilisant le morphisme $f : X \to S$ et les références ci-dessus, on obtient $H^p(X, \mathcal{F}) = 0$ pour $p > 1$ et pour tout $\mathcal{O}_X$-module cohérent $\mathcal{F}$. La suite exacte courte du lemme résulte donc de la suite spectrale de Leray pour $g$ et $\mathcal{O}_{X'}$ ; voir Cohomologie, lemme \ref{cohomology-lemma-Leray}. \end{proof}

\begin{lemma} \label{resolve-lemma-bound-a-torsion} Soit $(A, \mathfrak m, \kappa)$ un anneau local normal intègre de Nagata et de dimension $2$. Soit $a \in A$ non nul. Il existe un entier $N$ tel que, pour toute modification $f : X \to \Spec(A)$ avec $X$ normal, le $A$-module $$
M_{X, a} = \Coker(A \longrightarrow H^0(Z, \mathcal{O}_Z))
$$, où $Z \subset X$ est défini par $a$, soit de longueur bornée par $N$. \end{lemma}

\begin{proof} La suite exacte courte $
0 \to \mathcal{O}_X \xrightarrow{a} \mathcal{O}_X \to \mathcal{O}_Z \to 0
$ montre que \begin{equation}
\label{resolve-equation-a-torsion}
M_{X, a} = H^1(X, \mathcal{O}_X)[a]
\end{equation} Ici, $N[a] = \{n \in N \mid an = 0\}$ pour un $A$-module $N$. Ainsi, si $a$ divise $b$, alors $M_{X, a} \subset M_{X, b}$. Supposons que, pour un certain $c \in A$, les modules $M_{X, c}$ soient de longueurs bornées. Alors, pour tout $X$, on a une suite exacte $$
0 \to M_{X, c} \to M_{X, c^2} \to M_{X, c}
$$ où la deuxième flèche est donnée par la multiplication par $c$. Il s'ensuit que $M_{X, c^2}$ est également de longueur bornée. Il suffit donc de trouver un $c \in A$ pour lequel le lemme est vrai et tel que $a$ divise $c^n$ pour un certain $n > 0$. D'après Compléments d'algèbre, lemme \ref{more-algebra-lemma-divides-radical}, nous pouvons supposer que $A/(a)$ est un anneau réduit.

\medskip\noindent Supposons $A/(a)$ réduit. Soit $A/(a) \subset B$ la normalisation de $A/(a)$ dans son anneau total des fractions. Comme $A$ est de Nagata, $\Coker(A \to B)$ est fini. Nous affirmons que la longueur de ce module fini constitue une borne. Pour le voir, considérons $f : X \to \Spec(A)$ comme dans le lemme et soit $Z' \subset Z$ l'adhérence schématique de $Z \cap f^{-1}(U)$. Alors $Z' \to \Spec(A/(a))$ est fini, par exemple d'après Variétés, lemme \ref{varieties-lemma-finite-in-codim-1}. Par conséquent, $Z' = \Spec(B')$ avec $A/(a) \subset B' \subset B$. D'autre part, nous affirmons que l'application $$
H^0(Z, \mathcal{O}_Z) \to H^0(Z', \mathcal{O}_{Z'})
$$ est injective. En effet, si $s \in H^0(Z, \mathcal{O}_Z)$ appartient au noyau, alors la restriction de $s$ à $f^{-1}(U) \cap Z$ est nulle. L'image de $s$ dans $H^1(X, \mathcal{O}_X)$ s'annule donc dans $H^1(f^{-1}(U), \mathcal{O}_X)$. D'après le lemme \ref{resolve-lemma-H1-injective}, $s$ provient d'un élément $\tilde s$ de $A$. Mais, par hypothèse, $\tilde s$ s'envoie sur zéro dans $B'$, ce qui entraîne $s = 0$. En réunissant ces faits, on voit que $M_{X, a}$ est un sous-quotient de $B'/A$ : tout élément de $B'$ ne se prolonge pas nécessairement en une section globale de $\mathcal{O}_Z$, mais, dans tous les cas, la longueur de $M_{X, a}$ est bornée par celle de $B/A$. \end{proof}

\noindent Dans certains cas, la résolution des singularités se ramène au cas des singularités rationnelles.

\begin{definition} \label{resolve-definition-reduce-to-rational} Soit $(A, \mathfrak m, \kappa)$ un anneau local normal intègre de Nagata et de dimension $2$. \begin{enumerate} \item Nous disons que $A$ {\it définit une singularité rationnelle} si, pour toute modification normale $X \to \Spec(A)$, on a $H^1(X, \mathcal{O}_X) = 0$. \item Nous disons que {\it la réduction aux singularités rationnelles est possible pour $A$} si les longueurs des $A$-modules $$
H^1(X, \mathcal{O}_X)
$$ sont bornées pour toutes les modifications $X \to \Spec(A)$ avec $X$ normal. \end{enumerate} \end{definition}

\noindent Le sens de la terminologie employée en (2) est expliqué par le lemme \ref{resolve-lemma-reduce-to-rational}. Le lemme suivant affirme, en termes approximatifs, que les anneaux locaux des modifications de $\Spec(A)$, lorsque $A$ définit une singularité rationnelle, définissent eux aussi des singularités rationnelles.

\begin{lemma} \label{resolve-lemma-rational-propagates} Soit $(A, \mathfrak m, \kappa)$ un anneau local normal intègre de Nagata et de dimension $2$ qui définit une singularité rationnelle. Soit $A \subset B$ une extension locale d'anneaux intègres ayant le même corps des fractions, essentiellement de type fini, telle que $\dim(B) = 2$ et que $B$ soit normal. Alors $B$ définit une singularité rationnelle. \end{lemma}

\begin{proof} Choisissons une $A$-algèbre de type fini $C$ telle que $B = C_\mathfrak q$ pour un idéal premier $\mathfrak q \subset C$. Après avoir remplacé $C$ par l'image de $C$ dans $B$, nous pouvons supposer que $C$ est un anneau intègre dont le corps des fractions est celui de $A$. On peut alors choisir une immersion fermée $\Spec(C) \to \mathbf{A}^n_A$ et prendre l'adhérence dans $\mathbf{P}^n_A$ pour conclure que $B$ est isomorphe à $\mathcal{O}_{X, x}$ pour un certain point fermé $x \in X$ d'une modification projective $X \to \Spec(A)$. (Morphismes, lemme \ref{morphisms-lemma-dimension-formula}, montre que $\kappa(x)$ est fini sur $\kappa$, puis Morphismes, lemme \ref{morphisms-lemma-algebraic-residue-field-extension-closed-point-fibre}, montre que $x$ est un point fermé.) Soit $\nu : X^\nu \to X$ la normalisation. Comme $A$ est de Nagata, le morphisme $\nu$ est fini (Morphismes, lemme \ref{morphisms-lemma-nagata-normalization}). Ainsi $X^\nu$ est projectif sur $A$ d'après Compléments sur les morphismes, lemme \ref{more-morphisms-lemma-category-projective}. Puisque $B = \mathcal{O}_{X, x}$ est normal, on a $\mathcal{O}_{X, x} = (\nu_*\mathcal{O}_{X^\nu})_x$. Il existe donc un unique point $x^\nu \in X^\nu$ au-dessus de $x$, et $\mathcal{O}_{X^\nu, x^\nu} = \mathcal{O}_{X, x}$. Nous pouvons ainsi supposer $X$ normal et projectif sur $A$. Soit $Y \to \Spec(\mathcal{O}_{X, x}) = \Spec(B)$ une modification avec $Y$ normal. Il faut montrer que $H^1(Y, \mathcal{O}_Y) = 0$. D'après Limites, lemme \ref{limits-lemma-modifications}, il existe un morphisme de schémas $g : X' \to X$ qui est un isomorphisme au-dessus de $X \setminus \{x\}$ et tel que $X' \times_X \Spec(\mathcal{O}_{X, x})$ soit isomorphe à $Y$. Alors $g$ est une modification, car il est propre d'après Limites, lemme \ref{limits-lemma-modifications-properties}. L'anneau local de $X'$ en un point de $x'$ est soit isomorphe à l'anneau local de $X$ en $g(x')$ si $g(x') \not = x$, soit, si $g(x') = x$, l'anneau local de $X'$ en $x'$ est isomorphe à celui de $Y$ au point correspondant. Ainsi $X'$ est normal puisque $X$ et $Y$ le sont. Par conséquent, $H^1(X', \mathcal{O}_{X'}) = 0$ d'après notre hypothèse sur $A$. Le lemme \ref{resolve-lemma-exact-sequence} donne $R^1g_*\mathcal{O}_{X'} = 0$. Cela signifie manifestement que $H^1(Y, \mathcal{O}_Y) = 0$, comme souhaité. \end{proof}

\begin{lemma} \label{resolve-lemma-reduce-to-rational} Soit $(A, \mathfrak m, \kappa)$ un anneau local normal intègre de Nagata et de dimension $2$. Si la réduction aux singularités rationnelles est possible pour $A$, alors il existe une suite finie d'éclatements normalisés $$
X = X_n \to X_{n - 1} \to \ldots \to X_1 \to X_0 = \Spec(A)
$$ en des points fermés telle que, pour tout point fermé $x \in X$, l'anneau local $\mathcal{O}_{X, x}$ définisse une singularité rationnelle. En particulier, $X \to \Spec(A)$ est une modification et $X$ est un schéma normal projectif sur $A$. \end{lemma}

\begin{proof} Choisissons une modification $X \to \Spec(A)$ avec $X$ normal qui maximise la longueur de $H^1(X, \mathcal{O}_X)$. D'après le lemme \ref{resolve-lemma-exact-sequence}, pour toute modification ultérieure $g : X' \to X$ avec $X'$ normal, on a $R^1g_*\mathcal{O}_{X'} = 0$ et $H^1(X, \mathcal{O}_X) = H^1(X', \mathcal{O}_{X'})$.

\medskip\noindent Soit $x \in X$ un point fermé. Nous allons montrer que $\mathcal{O}_{X, x}$ définit une singularité rationnelle. Soit $Y \to \Spec(\mathcal{O}_{X, x})$ une modification avec $Y$ normal. Il faut montrer que $H^1(Y, \mathcal{O}_Y) = 0$. D'après Limites, lemme \ref{limits-lemma-modifications}, il existe un morphisme de schémas $g : X' \to X$ qui est un isomorphisme au-dessus de $X \setminus \{x\}$ et tel que $X' \times_X \Spec(\mathcal{O}_{X, x})$ soit isomorphe à $Y$. Alors $g$ est une modification, car il est propre d'après Limites, lemme \ref{limits-lemma-modifications-properties}. L'anneau local de $X'$ en un point de $x'$ est soit isomorphe à l'anneau local de $X$ en $g(x')$ si $g(x') \not = x$, soit, si $g(x') = x$, l'anneau local de $X'$ en $x'$ est isomorphe à celui de $Y$ au point correspondant. Ainsi $X'$ est normal puisque $X$ et $Y$ le sont. Par maximalité, on a $R^1g_*\mathcal{O}_{X'} = 0$ (voir le premier paragraphe). Cela signifie manifestement que $H^1(Y, \mathcal{O}_Y) = 0$, comme souhaité.

\medskip\noindent En conclusion, nous avons trouvé une modification normale $X$ de $\Spec(A)$ telle que tous les anneaux locaux de $X$ aux points fermés définissent des singularités rationnelles. Choisissons alors une suite d'éclatements normalisés $X_n \to \ldots \to X_1 \to \Spec(A)$ telle que $X_n$ domine $X$ ; voir le lemme \ref{resolve-lemma-dominate-by-normalized-blowing-up}. Pour un point fermé $x' \in X_n$ s'envoyant sur $x \in X$, on peut appliquer le lemme \ref{resolve-lemma-rational-propagates} à l'homomorphisme d'anneaux $\mathcal{O}_{X, x} \to \mathcal{O}_{X_n, x'}$ pour voir que $\mathcal{O}_{X_n, x'}$ définit une singularité rationnelle. \end{proof}

\begin{lemma} \label{resolve-lemma-go-up-separable} Soit $A \to B$ un homomorphisme local fini injectif entre anneaux locaux normaux intègres de Nagata et de dimension $2$. Supposons que l'extension induite des corps des fractions soit séparable. Si la réduction aux singularités rationnelles est possible pour $A$, elle l'est aussi pour $B$. \end{lemma}

\begin{proof} Soit $n$ le degré de l'extension de corps des fractions $L/K$. Soit $\text{Trace}_{L/K} : L \to K$ la trace. Puisque l'extension est finie séparable, la forme bilinéaire $(h, g) \mapsto \text{Trace}_{L/K}(fg)$ est non dégénérée sur $L$, considéré comme $K$-espace vectoriel. Voir Corps, lemme \ref{fields-lemma-separable-trace-pairing}. Choisissons $b_1, \ldots, b_n \in B$ qui forment une base de $L$ sur $K$. D'après ce qui précède, $d = \det(\text{Trace}_{L/K}(b_ib_j)) \in A$ est non nul.

\medskip\noindent Soit $Y \to \Spec(B)$ une modification avec $Y$ normal. Il existe un éclatement $U$-admissible $X'$ de $\Spec(A)$ tel que le transformé strict $Y'$ de $Y$ soit fini sur $X'$ ; voir Compléments sur la platitude, lemme \ref{flat-lemma-finite-after-blowing-up}. Considérons le diagramme $$
\xymatrix{
Y' \ar[d] \ar[r] & Y \ar[r] & \Spec(B) \ar[d] \\
X' \ar[rr] & & \Spec(A)
}
$$ Après avoir remplacé $X'$ et $Y'$ par leurs normalisations, nous pouvons supposer que $X'$ et $Y'$ sont des modifications normales de $\Spec(A)$ et $\Spec(B)$. On se ramène ainsi au cas où il existe un diagramme commutatif $$
\xymatrix{
Y \ar[d]_\pi \ar[r]_-g & \Spec(B) \ar[d] \\
X \ar[r]^-f & \Spec(A)
}
$$ dans lequel $X$ et $Y$ sont des modifications normales de $\Spec(A)$ et $\Spec(B)$ et où $\pi$ est fini.

\medskip\noindent L'application trace de $L$ sur $K$ se prolonge en une application de $\mathcal{O}_X$-modules $\text{Trace} : \pi_*\mathcal{O}_Y \to \mathcal{O}_X$. Considérons l'application $$
\Phi : \pi_*\mathcal{O}_Y \longrightarrow \mathcal{O}_X^{\oplus n},\quad
s \longmapsto (\text{Trace}(b_1s), \ldots, \text{Trace}(b_ns))
$$ Cette application est injective (car elle l'est au point générique) et il existe une application $$
\mathcal{O}_X^{\oplus n} \longrightarrow \pi_*\mathcal{O}_Y,\quad
(s_1, \ldots, s_n) \longmapsto \sum b_i s_i
$$ dont le composé avec $\Phi$ a pour matrice $\text{Trace}(b_ib_j)$. Le conoyau de $\Phi$ est donc annulé par $d$. On obtient ainsi une suite exacte $$
H^0(X, \Coker(\Phi)) \to H^1(Y, \mathcal{O}_Y) \to
H^1(X, \mathcal{O}_X)^{\oplus n}
$$ Comme le terme de droite est borné par hypothèse, il suffit de montrer que la $d$-torsion de $H^1(Y, \mathcal{O}_Y)$ est bornée. C'est le contenu du lemme \ref{resolve-lemma-bound-a-torsion} et de (\ref{resolve-equation-a-torsion}). \end{proof}

\begin{lemma} \label{resolve-lemma-regular-rational} Soit $A$ un anneau local régulier de Nagata et de dimension $2$. Alors $A$ définit une singularité rationnelle. \end{lemma}

\begin{proof} (L'hypothèse selon laquelle $A$ est de Nagata n'est pas nécessaire pour cette démonstration, mais nous n'avons défini la notion de singularité rationnelle que pour les anneaux locaux normaux intègres de Nagata et de dimension $2$.) Soit $X \to \Spec(A)$ une modification avec $X$ normal. D'après le lemme \ref{resolve-lemma-dominate-by-blowing-up-in-points}, on peut dominer $X$ par un schéma $X_n$ qui est le dernier terme d'une suite $$
X_n \to X_{n - 1} \to \ldots \to X_1 \to X_0 = \Spec(A)
$$ d'éclatements en des points fermés. D'après le lemme \ref{resolve-lemma-blowup-regular}, les schémas $X_i$ sont réguliers, donc en particulier normaux (Algèbre, lemme \ref{algebra-lemma-regular-normal}). Le lemme \ref{resolve-lemma-exact-sequence} donne $H^1(X, \mathcal{O}_X) \subset H^1(X_n, \mathcal{O}_{X_n})$. Il suffit donc de démontrer $H^1(X_n, \mathcal{O}_{X_n}) = 0$. En utilisant de nouveau le lemme \ref{resolve-lemma-exact-sequence}, il suffit de montrer $R^1(X_i \to X_{i - 1})_*\mathcal{O}_{X_i} = 0$ pour $i = 1, \ldots, n$. Cela résulte du lemme \ref{resolve-lemma-cohomology-of-blowup}. \end{proof}

\begin{lemma} \label{resolve-lemma-bound-dualizing-implies-bound} Soit $A$ un anneau local normal intègre de Nagata et de dimension $2$ qui possède un complexe dualisant $\omega_A^\bullet$. S'il existe un élément non nul $d \in A$ tel que, pour toute modification normale $X \to \Spec(A)$, le conoyau de l'application trace $$
\Gamma(X, \omega_X) \to \omega_A
$$ soit annulé par $d$, alors la réduction aux singularités rationnelles est possible pour $A$. \end{lemma}

\begin{proof} \begingroup\emergencystretch=2em
Pour $X \to \Spec(A)$ comme dans l'énoncé, il faut borner $H^1(X, \mathcal{O}_X)$. Soit $\omega_X$ le module dualisant de $X$ comme dans l'énoncé de Grauert--Riemenschneider (proposition \ref{resolve-proposition-Grauert-Riemenschneider}). L'application trace est l'application $Rf_*\omega_X \to \omega_A$ décrite dans Dualité pour les schémas, section \ref{duality-section-trace}. D'après Grauert--Riemenschneider, on a $Rf_*\omega_X = f_*\omega_X$ ; l'application trace induit donc bien une application $\Gamma(X, \omega_X) \to \omega_A$. Par dualité, on a $Rf_*\omega_X = R\Hom_A(Rf_*\mathcal{O}_X, \omega_A)$ (on utilise ici que $\omega_X[2]$ est le complexe dualisant sur $X$ normalisé relativement à $\omega_A[2]$ ; voir Dualité pour les schémas, lemme \ref{duality-lemma-duality-bootstrap}, ou plus directement la section \ref{duality-section-duality}, voire, encore plus directement, l'exemple \ref{duality-example-iso-on-RSheafHom-noetherian}). Le triangle distingué\par\endgroup
$$
A \to Rf_*\mathcal{O}_X \to R^1f_*\mathcal{O}_X[-1] \to A[1]
$$ est transformé par $R\Hom_A(-, \omega_A)$ en la suite exacte courte $$
0 \to f_*\omega_X \to \omega_A \to
\Ext_A^2(R^1f_*\mathcal{O}_X, \omega_A) \to 0
$$ (et $\Ext_A^i(R^1f_*\mathcal{O}_X, \omega_A) = 0$ pour $i \not = 2$ ; cela résulte aussi de la discussion ci-dessous). Puisque $R^1f_*\mathcal{O}_X$ est supporté dans $\{\mathfrak m\}$, le théorème de dualité locale affirme que $$
\Ext_A^2(R^1f_*\mathcal{O}_X, \omega_A) =
\Ext_A^0(R^1f_*\mathcal{O}_X, \omega_A[2]) =
\Hom_A(R^1f_*\mathcal{O}_X, E)
$$ est le dual de Matlis de $R^1f_*\mathcal{O}_X$ (et que les autres groupes d'extensions sont nuls) ; voir Complexes dualisants, lemme \ref{dualizing-lemma-special-case-local-duality}. Par l'équivalence de catégories inhérente à la dualité de Matlis (Complexes dualisants, proposition \ref{dualizing-proposition-matlis}), si $R^1f_*\mathcal{O}_X$ n'est pas annulé par $d$, alors le $\Ext^2$ ci-dessus ne l'est pas non plus. On voit donc que $H^1(X, \mathcal{O}_X)$ est annulé par $d$. La bornitude requise résulte alors du lemme \ref{resolve-lemma-bound-a-torsion} et de (\ref{resolve-equation-a-torsion}). \end{proof}

\begin{lemma} \label{resolve-lemma-compare-differentials-dualizing} Soit $p$ un nombre premier. Soit $A$ un anneau local régulier de dimension $2$ et de caractéristique $p$. Soit $A_0 \subset A$ un sous-anneau tel que $\Omega_{A/A_0}$ soit libre de rang $r < \infty$. Posons $\omega_A = \Omega^r_{A/A_0}$. Si $X \to \Spec(A)$ est obtenu par une suite d'éclatements en des points fermés, alors il existe une application $$
\varphi_X : (\Omega^r_{X/\Spec(A_0)})^{**} \longrightarrow \omega_X
$$ qui prolonge l'identification donnée au point générique. \end{lemma}

\begin{proof} Remarquons que $A$ est de Gorenstein (Complexes dualisants, lemme \ref{dualizing-lemma-regular-gorenstein}) ; le module inversible $\omega_A$ sert donc bien de module dualisant. En outre, tout $X$ comme dans le lemme possède un module dualisant inversible $\omega_X$, puisque $X$ est régulier (donc de Gorenstein) et propre sur $A$ ; voir la remarque \ref{resolve-remark-dualizing-setup} et le lemme \ref{resolve-lemma-blowup-regular}. Supposons construite l'application $\varphi_X : (\Omega^r_{X/A_0})^{**} \to \omega_X$ et supposons que $b : X' \to X$ soit l'éclatement en un point fermé. Posons $\Omega^r_X = (\Omega^r_{X/A_0})^{**}$ et $\Omega^r_{X'} = (\Omega^r_{X'/A_0})^{**}$. Puisque $\omega_{X'} = b^!(\omega_X)$, une application $\Omega^r_{X'} \to \omega_{X'}$ équivaut à une application $Rb_*(\Omega^r_{X'}) \to \omega_X$. Voir la discussion de la remarque \ref{resolve-remark-dualizing-setup} et Dualité pour les schémas, section \ref{duality-section-duality}. Il suffit donc à son tour de construire une application $$
Rb_*(\Omega^r_{X'}) \longrightarrow \Omega^r_X
$$ Les faisceaux $\Omega^r_{X'}$ et $\Omega^r_X$ sont inversibles ; voir Diviseurs, lemme \ref{divisors-lemma-reflexive-over-regular-dim-2}. Considérons la suite exacte $$
b^*\Omega_{X/A_0} \to \Omega_{X'/A_0} \to \Omega_{X'/X} \to 0
$$ Un calcul local montre que $\Omega_{X'/X}$ est isomorphe à un module inversible sur le diviseur exceptionnel $E$ ; voir le lemme \ref{resolve-lemma-differentials-of-blowup}. Il en résulte que l'on a l'une des deux possibilités $$
\Omega^r_{X'} \cong (b^*\Omega^r_X)(E)
\quad\text{ou}\quad
\Omega^r_{X'} \cong b^*\Omega^r_X
$$ ; voir Diviseurs, lemme \ref{divisors-lemma-wedge-product-ses}. (La seconde possibilité ne se produit jamais en caractéristique zéro, mais peut se produire en caractéristique $p$.) Dans les deux cas, on a $R^1b_*(\Omega^r_{X'}) = 0$ et $b_*(\Omega^r_{X'}) = \Omega^r_X$ d'après le lemme \ref{resolve-lemma-cohomology-of-blowup}. \end{proof}

\begin{lemma} \label{resolve-lemma-go-up-degree-p} Soit $p$ un nombre premier. Soit $A$ un anneau local régulier complet de dimension $2$ et de caractéristique $p$. Soit $L/K$ une extension inséparable de degré $p$ du corps des fractions $K$ de $A$. Soit $B \subset L$ la clôture intégrale de $A$. Alors la réduction aux singularités rationnelles est possible pour $B$. \end{lemma}

\begin{proof} On a $A = k[[x, y]]$. Écrivons $L = K[x]/(x^p - f)$ pour un certain $f \in A$ et notons $g \in B$ la classe de congruence de $x$, c'est-à-dire l'élément tel que $g^p = f$. D'après Algèbre, lemme \ref{algebra-lemma-derivative-zero-pth-power}, $\text{d}f$ est non nul dans $\Omega_{K/\mathbf{F}_p}$. D'après Compléments d'algèbre, lemme \ref{more-algebra-lemma-power-series-ring-subfields}, il existe un sous-corps $k^p \subset k' \subset k$ avec $p^e = [k : k'] < \infty$ tel que $\text{d}f$ soit non nul dans $\Omega_{K/K_0}$, où $K_0$ est le corps des fractions de $A_0 = k'[[x^p, y^p]] \subset A$. Alors $$
\Omega_{A/A_0} =
A \otimes_k \Omega_{k/k'} \oplus A \text{d}x \oplus A \text{d}y
$$ est libre fini de rang $e + 2$. Posons $\omega_A = \Omega^{e + 2}_{A/A_0}$. Considérons l'application canonique $$
\text{Tr} :
\Omega^{e + 2}_{B/A_0}
\longrightarrow
\Omega^{e + 2}_{A/A_0} = \omega_A
$$ du lemme \ref{resolve-lemma-trace-extends}. Par dualité, elle détermine une application $$
c : \Omega^{e + 2}_{B/A_0} \to \omega_B = \Hom_A(B, \omega_A)
$$ Affirmation : le conoyau de $c$ est annulé par un élément non nul de $B$.

\medskip\noindent Puisque $\text{d}f$ est non nul dans $\Omega_{A/A_0}$, on peut trouver $\eta_1, \ldots, \eta_{e + 1} \in \Omega_{A/A_0}$ tels que $\theta = \eta_1 \wedge \ldots \wedge \eta_{e + 1} \wedge \text{d}f$ soit non nul dans $\omega_A = \Omega^{e + 2}_{A/A_0}$. Pour démontrer l'affirmation, nous allons construire des éléments $\omega_i$ de $\Omega^{e + 2}_{B/A_0}$, $i = 0, \ldots, p - 1$, qui sont envoyés sur $\varphi_i \in \omega_B = \Hom_A(B, \omega_A)$ avec $\varphi_i(g^j) = \delta_{ij}\theta$ pour $j = 0, \ldots, p - 1$. Puisque $\{1, g, \ldots, g^{p - 1}\}$ est une base de $L/K$, cela démontre l'affirmation. Posons $\eta = \eta_1 \wedge \ldots \wedge \eta_{e + 1}$, de sorte que $\theta = \eta \wedge \text{d}f$. Posons $\omega_i = \eta \wedge g^{p - 1 - i}\text{d}g$. Par construction, on a alors $$
\varphi_i(g^j) = \text{Tr}(g^j \eta \wedge g^{p - 1 - i}\text{d}g) =
\text{Tr}(\eta \wedge g^{p - 1 - i + j}\text{d}g) = \delta_{ij} \theta
$$ d'après la description explicite de l'application trace dans le lemme \ref{resolve-lemma-trace-higher}.

\medskip\noindent Soit $Y \to \Spec(B)$ une modification normale. Exactement comme dans la démonstration du lemme \ref{resolve-lemma-go-up-separable}, on peut se ramener au cas où $Y$ est fini sur une modification $X$ de $\Spec(A)$. D'après le lemme \ref{resolve-lemma-dominate-by-blowing-up-in-points}, on peut même supposer que $X \to \Spec(A)$ est obtenu par une suite d'éclatements en des points fermés. Considérons le diagramme : $$
\xymatrix{
Y \ar[d]_\pi \ar[r]_-g & \Spec(B) \ar[d] \\
X \ar[r]^-f & \Spec(A)
}
$$ On peut appliquer le lemme \ref{resolve-lemma-trace-extends} à $\pi$, ce qui fournit la première flèche de $$
\pi_*(\Omega^{e + 2}_{Y/A_0})
\xrightarrow{\text{Tr}}
(\Omega^{e + 2}_{X/A_0})^{**}
\xrightarrow{\varphi_X}
\omega_X
$$ ; la seconde provient du lemme \ref{resolve-lemma-compare-differentials-dualizing} (puisque $f$ est une suite d'éclatements en des points fermés). Par dualité pour le morphisme fini $\pi$, cela correspond à une application $$
c_Y : \Omega^{e + 2}_{Y/A_0} \longrightarrow \omega_Y
$$ qui prolonge l'application $c$ ci-dessus. Ainsi, l'image de $\Gamma(Y, \omega_Y) \to \omega_B$ contient celle de $c$. Notre affirmation montre que le conoyau est annulé par un élément non nul fixé de $B$. On conclut par le lemme \ref{resolve-lemma-bound-dualizing-implies-bound}. \end{proof}






\section{Singularités rationnelles} \label{resolve-section-rational-singularities}

\noindent Dans cette section, nous ramenons le cas des points singuliers rationnels à celui des points singuliers rationnels de Gorenstein. Voir \cite{Lipman-rational} et \cite{Mattuck}.

\begin{situation} \label{resolve-situation-rational} Ici, $(A, \mathfrak m, \kappa)$ est un anneau local normal intègre de Nagata et de dimension $2$ qui définit une singularité rationnelle. Soit $s$ le point fermé de $S = \Spec(A)$ et soit $U = S \setminus \{s\}$. Soit $f : X \to S$ une modification avec $X$ normal. Nous notons $C_1, \ldots, C_r$ les composantes irréductibles de la fibre spéciale $X_s$ de $f$. \end{situation}

\begin{lemma} \label{resolve-lemma-globally-generated} Dans la situation \ref{resolve-situation-rational}, soit $\mathcal{F}$ un $\mathcal{O}_X$-module quasi-cohérent. Alors \begin{enumerate} \item $H^p(X, \mathcal{F}) = 0$ pour $p \not \in \{0, 1\}$ ; \item $H^1(X, \mathcal{F}) = 0$ si $\mathcal{F}$ est engendré par ses sections globales. \end{enumerate} \end{lemma}

\begin{proof} La partie (1) résulte de Cohomologie des schémas, lemme \ref{coherent-lemma-higher-direct-images-zero-above-dimension-fibre}. Si $\mathcal{F}$ est engendré par ses sections globales, il existe une surjection $\bigoplus_{i \in I} \mathcal{O}_X \to \mathcal{F}$. D'après la partie (1) et la suite exacte longue de cohomologie, celle-ci induit une surjection sur $H^1$. Puisque $H^1(X, \mathcal{O}_X) = 0$, car $S$ définit une singularité rationnelle, et puisque $H^1(X, -)$ commute aux sommes directes (Cohomologie, lemme \ref{cohomology-lemma-quasi-separated-cohomology-colimit}), on conclut. \end{proof}

\begin{lemma} \label{resolve-lemma-sections-powers-I-rational} Dans la situation \ref{resolve-situation-rational}, supposons que $E = X_s$ soit un diviseur de Cartier effectif. Soit $\mathcal{I}$ le faisceau d'idéaux de $E$. Alors $H^0(X, \mathcal{I}^n) = \mathfrak m^n$ et $H^1(X, \mathcal{I}^n) = 0$. \end{lemma}

\begin{proof} On a $H^0(X, \mathcal{O}_X) = A$ ; voir la discussion qui suit la situation \ref{resolve-situation-vanishing}. Alors $\mathfrak m \subset H^0(X, \mathcal{I}) \subset H^0(X, \mathcal{O}_X)$. La seconde inclusion n'est pas une égalité, car $X_s \not = \emptyset$. Ainsi $H^0(X, \mathcal{I}) = \mathfrak m$. Comme $\mathcal{I}^n = \mathfrak m^n\mathcal{O}_X$, notre lemme \ref{resolve-lemma-globally-generated} montre que $H^1(X, \mathcal{I}^n) = 0$.

\medskip\noindent Choisissons des générateurs $x_1, \ldots, x_{\mu + 1}$ de $\mathfrak m$. Ils définissent des sections globales de $\mathcal{I}$ qui l'engendrent. On obtient donc une suite exacte courte $$
0 \to \mathcal{F} \to \mathcal{O}_X^{\oplus \mu + 1} \to \mathcal{I} \to 0
$$ Alors $\mathcal{F}$ est un $\mathcal{O}_X$-module localement libre fini de rang $\mu$ et $\mathcal{F} \otimes \mathcal{I}$ est engendré par ses sections globales d'après Constructions, lemme \ref{constructions-lemma-globally-generated-omega-twist-1}. Par conséquent, $\mathcal{F} \otimes \mathcal{I}^n$ est engendré par ses sections globales pour tout $n \geq 1$. Ainsi, pour $n \geq 2$, on peut considérer la suite exacte $$
0 \to \mathcal{F} \otimes \mathcal{I}^{n - 1} \to
(\mathcal{I}^{n - 1})^{\oplus \mu + 1} \to
\mathcal{I}^n \to 0
$$ En appliquant la suite exacte longue de cohomologie et en utilisant $H^1(X, \mathcal{F} \otimes \mathcal{I}^{n - 1}) = 0$ d'après le lemme \ref{resolve-lemma-globally-generated}, on obtient que tout élément de $H^0(X, \mathcal{I}^n)$ est de la forme $\sum x_i a_i$ pour certains $a_i \in H^0(X, \mathcal{I}^{n - 1})$. Cela montre par récurrence que $H^0(X, \mathcal{I}^n) = \mathfrak m^n$. \end{proof}

\begin{lemma} \label{resolve-lemma-blow-up-normal-rational} Dans la situation \ref{resolve-situation-rational}, l'éclatement de $\Spec(A)$ en $\mathfrak m$ est normal. \end{lemma}

\begin{proof} Soit $X' \to \Spec(A)$ l'éclatement ; autrement dit, $$
X' = \text{Proj}(A \oplus \mathfrak m \oplus \mathfrak m^2 \oplus \ldots).
$$ est le Proj de l'algèbre de Rees. Cela montre en particulier que $X'$ est intègre et que $X' \to \Spec(A)$ est une modification projective. Soit $X$ la normalisation de $X'$. Puisque $A$ est de Nagata, $\nu : X \to X'$ est fini (Morphismes, lemme \ref{morphisms-lemma-nagata-normalization}). Soit $E' \subset X'$ le diviseur exceptionnel et soit $E \subset X$ son image réciproque. Soient $\mathcal{I}' \subset \mathcal{O}_{X'}$ et $\mathcal{I} \subset \mathcal{O}_X$ leurs faisceaux d'idéaux. Rappelons que $\mathcal{I}' = \mathcal{O}_{X'}(1)$ (Diviseurs, lemme \ref{divisors-lemma-blowing-up-projective}). Remarquons que $\mathcal{I} = \nu^*\mathcal{I}'$ et que $E$ est un diviseur de Cartier effectif (Diviseurs, lemme \ref{divisors-lemma-pullback-effective-Cartier-defined}). Nous cherchons à montrer que $\nu$ est un isomorphisme. Comme $\nu$ est fini, il suffit de montrer que $\mathcal{O}_{X'} \to \nu_*\mathcal{O}_X$ est un isomorphisme. Sinon, il existe un $n \geq 0$ tel que $$
H^0(X', (\mathcal{I}')^n) \not =
H^0(X', (\nu_*\mathcal{O}_X) \otimes (\mathcal{I}')^n)
$$, par exemple parce que l'on peut retrouver les $\mathcal{O}_{X'}$-modules quasi-cohérents à partir de leurs modules gradués associés ; voir Propriétés, lemme \ref{properties-lemma-ample-quasi-coherent}. La formule de projection donne $$
H^0(X', (\nu_*\mathcal{O}_X) \otimes (\mathcal{I}')^n) =
H^0(X, \nu^*(\mathcal{I}')^n) =
H^0(X, \mathcal{I}^n) = \mathfrak m^n
$$, la dernière égalité résultant du lemme \ref{resolve-lemma-sections-powers-I-rational}. D'autre part, il existe manifestement une injection $\mathfrak m^n \to H^0(X', (\mathcal{I}')^n)$. Puisque $H^0(X', (\mathcal{I}')^n)$ est sans torsion, nous concluons qu'il y a égalité pour tout $n$, donc que $X = X'$. \end{proof}

\begin{lemma} \label{resolve-lemma-cohomology-blow-up-rational} Dans la situation \ref{resolve-situation-rational}, soit $X$ l'éclatement de $\Spec(A)$ en $\mathfrak m$. Soit $E \subset X$ le diviseur exceptionnel. Avec $\mathcal{O}_X(1) = \mathcal{I}$ comme d'habitude et $\mathcal{O}_E(1) = \mathcal{O}_X(1)|_E$, on a : \begin{enumerate} \item $E$ est une courbe propre de Cohen--Macaulay sur $\kappa$ ; \item $\mathcal{O}_E(1)$ est très ample ; \item $\deg(\mathcal{O}_E(1)) \geq 1$, l'égalité n'ayant lieu que si $A$ est un anneau local régulier ; \item $H^1(E, \mathcal{O}_E(n)) = 0$ pour $n \geq 0$ ; \item $H^0(E, \mathcal{O}_E(n)) = \mathfrak m^n/\mathfrak m^{n + 1}$ pour $n \geq 0$. \end{enumerate} \end{lemma}

\begin{proof} Puisque $\mathcal{O}_X(1)$ est très ample par construction, sa restriction à la fibre spéciale $E$ l'est aussi. D'après le lemme \ref{resolve-lemma-blow-up-normal-rational}, le schéma $X$ est normal. Alors $E$ est de Cohen--Macaulay d'après Diviseurs, lemme \ref{divisors-lemma-normal-effective-Cartier-divisor-S1}. Le lemme \ref{resolve-lemma-sections-powers-I-rational} s'applique, et les suites exactes $$
0 \to \mathcal{I}^{n + 1} \to \mathcal{I}^n \to i_*\mathcal{O}_E(n) \to 0
$$ ainsi que la suite exacte longue de cohomologie donnent (4) et (5). En particulier, on a $$
\deg(\mathcal{O}_E(1)) = \chi(E, \mathcal{O}_E(1)) - \chi(E, \mathcal{O}_E) =
\dim(\mathfrak m/\mathfrak m^2) - 1
$$ d'après Variétés, définition \ref{varieties-definition-degree-invertible-sheaf}. Cela donne également (3). \end{proof}

\begin{lemma} \label{resolve-lemma-dualizing-rational} Dans la situation \ref{resolve-situation-rational}, supposons que $A$ possède un complexe dualisant $\omega_A^\bullet$. Si $\omega_X$ est le module dualisant de $X$, alors l'application trace $H^0(X, \omega_X) \to \omega_A$ est un isomorphisme et, par conséquent, il existe une application canonique $f^*\omega_A \to \omega_X$. \end{lemma}

\begin{proof} \begingroup\emergencystretch=2em
D'après Grauert--Riemenschneider (proposition \ref{resolve-proposition-Grauert-Riemenschneider}), on a $Rf_*\omega_X = f_*\omega_X$. Par dualité, on dispose d'une suite exacte courte\par\endgroup
$$
0 \to f_*\omega_X \to \omega_A \to
\Ext^2_A(R^1f_*\mathcal{O}_X, \omega_A) \to 0
$$ (voir par exemple la démonstration du lemme \ref{resolve-lemma-bound-dualizing-implies-bound}) et, puisque $A$ définit une singularité rationnelle, on obtient $f_*\omega_X = \omega_A$. \end{proof}

\begin{lemma} \label{resolve-lemma-dualizing-blow-up-rational} Dans la situation \ref{resolve-situation-rational}, supposons que $A$ possède un complexe dualisant $\omega_A^\bullet$ et ne soit pas régulier. Soit $X$ l'éclatement de $\Spec(A)$ en $\mathfrak m$, de diviseur exceptionnel $E \subset X$. Soit $\omega_X$ le module dualisant de $X$. Alors : \begin{enumerate} \item $\omega_E = \omega_X|_E \otimes \mathcal{O}_E(-1)$ ; \item $H^1(X, \omega_X(n)) = 0$ pour $n \geq 0$ ; \item l'application $f^*\omega_A \to \omega_X$ du lemme \ref{resolve-lemma-dualizing-rational} est surjective. \end{enumerate} \end{lemma}

\begin{proof} Nous utiliserons sans autre mention les résultats du lemme \ref{resolve-lemma-cohomology-blow-up-rational}. Remarquons que $\omega_E = \omega_X|_E \otimes \mathcal{O}_E(-1)$ d'après Dualité pour les schémas, lemmes \ref{duality-lemma-sheaf-with-exact-support-effective-Cartier} et \ref{duality-lemma-twisted-inverse-image-closed}. Ainsi $\omega_X|_E = \omega_E(1)$. Considérons les suites exactes courtes $$
0 \to \omega_X(n + 1) \to \omega_X(n) \to i_*\omega_E(n + 1) \to 0
$$ D'après Courbes algébriques, lemme \ref{curves-lemma-vanishing-twist}, on a $H^1(E, \omega_E(n + 1)) = 0$ pour $n \geq 0$. Les applications $$
\ldots \to H^1(X, \omega_X(2)) \to H^1(X, \omega_X(1)) \to H^1(X, \omega_X)
$$ sont donc surjectives. Puisque $H^1(X, \omega_X(n))$ est nul pour $n \gg 0$ (Cohomologie des schémas, lemme \ref{coherent-lemma-kill-by-twisting}), on en déduit (2).

\medskip\noindent D'après Courbes algébriques, lemme \ref{curves-lemma-tensor-omega-with-globally-generated-invertible}, $\omega_X|_E = \omega_E \otimes \mathcal{O}_E(1)$ est engendré par ses sections globales. Puisque nous avons vu ci-dessus que $H^1(X, \omega_X(1)) = 0$, l'application $H^0(X, \omega_X) \to H^0(E, \omega_X|_E)$ est surjective. Nous concluons que $\omega_X$ est engendré par ses sections globales ; cela donne (3), car $\Gamma(X, \omega_X) = \omega_A$ est utilisé dans le lemme \ref{resolve-lemma-dualizing-rational} pour définir l'application. \end{proof}

\begin{lemma} \label{resolve-lemma-rational-to-gorenstein} Soit $(A, \mathfrak m, \kappa)$ un anneau local normal intègre de Nagata et de dimension $2$ qui définit une singularité rationnelle. Supposons que $A$ possède un complexe dualisant. Il existe alors une suite finie d'éclatements en des points fermés singuliers $$
X = X_n \to X_{n - 1} \to \ldots \to X_1 \to X_0 = \Spec(A)
$$ telle que $X_i$ soit normal pour tout $i$ et que le faisceau dualisant $\omega_X$ de $X$ soit un $\mathcal{O}_X$-module inversible. \end{lemma}

\begin{proof} Le module dualisant $\omega_A$ est un $A$-module fini dont le germe au point générique est inversible. En effet, $\omega_A \otimes_A K$ est un module dualisant pour le corps des fractions $K$ de $A$, donc est de rang $1$. Il existe ainsi un éclatement $b : Y \to \Spec(A)$ tel que le transformé strict de $\omega_A$ relativement à $b$ soit un $\mathcal{O}_Y$-module inversible ; voir Diviseurs, lemme \ref{divisors-lemma-blowup-fitting-ideal}. D'après le lemme \ref{resolve-lemma-dominate-by-normalized-blowing-up}, on peut choisir une suite d'éclatements normalisés $$
X_n \to X_{n - 1} \to \ldots \to X_1 \to \Spec(A)
$$ telle que $X_n$ domine $Y$. Le lemme \ref{resolve-lemma-blow-up-normal-rational} et un raisonnement par récurrence montrent que chaque $X_i \to X_{i - 1}$ est simplement un éclatement.

\medskip\noindent Nous affirmons que $\omega_{X_n}$ est inversible. Puisque $\omega_{X_n}$ est un $\mathcal{O}_{X_n}$-module cohérent, il suffit de vérifier que ses germes sont des modules inversibles. Si $x \in X_n$ est un point régulier, cela résulte du fait que les schémas réguliers sont de Gorenstein (Complexes dualisants, lemme \ref{dualizing-lemma-regular-gorenstein}). Si $x$ est un point singulier de $X_n$, alors chacune des images $x_i \in X_i$ de $x$ est un point singulier (car l'éclatement d'un point régulier est régulier d'après le lemme \ref{resolve-lemma-blowup-regular}). Considérons l'application canonique $f_n^*\omega_A \to \omega_{X_n}$ du lemme \ref{resolve-lemma-dualizing-rational}. Pour tout $i$, le morphisme $X_{i + 1} \to X_i$ est soit l'éclatement de $x_i$, soit un isomorphisme en $x_i$. Puisque $x_i$ est toujours un point singulier, le lemme \ref{resolve-lemma-dualizing-blow-up-rational} et une récurrence montrent que les applications $f_i^*\omega_A \to \omega_{X_i}$ sont toujours surjectives sur les germes en $x_i$. Par conséquent, $$
(f_n^*\omega_A)_x \longrightarrow \omega_{X_n, x}
$$ est surjective. D'autre part, par notre choix de $b$, le quotient de $f_n^*\omega_A$ par son sous-module de torsion est un module inversible $\mathcal{L}$. De plus, le module dualisant est sans torsion (Dualité pour les schémas, lemme \ref{duality-lemma-dualizing-module}). Il s'ensuit que $\mathcal{L}_x \cong  \omega_{X_n, x}$, ce qui achève la démonstration. \end{proof}






\section{Arcs formels} \label{resolve-section-arcs}

\noindent Soit $X$ un schéma localement noethérien. Dans cette section, nous appelons {\it arc formel} dans $X$ un morphisme $a : T \to X$ où $T$ est le spectre d'un anneau de valuation discrète complet $R$ dont le corps résiduel $\kappa$ est identifié au corps résiduel de l'image $p$ du point fermé de $\Spec(R)$. Nous dirons alors que l'arc formel $a$ est {\it centré en $p$}. Nous disons que l'arc formel $T \to X$ est {\it non singulier} si l'application induite $\mathfrak m_p/\mathfrak m_p^2 \to \mathfrak m_R/\mathfrak m_R^2$ est surjective.

\medskip\noindent Soient $a : T \to X$ et $T = \Spec(R)$ un arc formel non singulier centré en un point fermé $p$ de $X$. Supposons $X$ localement noethérien. Soit $b : X_1 \to X$ l'éclatement de $X$ en $x$. Comme $a$ est non singulier, il existe un élément $f \in \mathfrak m_p$ qui s'envoie sur une uniformisante de $R$. En particulier, le point générique de $T$ s'envoie sur un point de $X$ distinct de $p$. Autrement dit, si $K$ désigne le corps des fractions de $R$, la restriction de $a$ définit un morphisme $\Spec(K) \to X \setminus \{p\}$. Puisque le morphisme $b$ est propre et est un isomorphisme au-dessus de $X \setminus \{x\}$, le critère valuatif de propreté fournit un unique morphisme $a_1$ qui rend commutatif le diagramme suivant $$
\xymatrix{
T \ar[r]_{a_1} \ar[rd]_a & X_1 \ar[d]^{b} \\
& X
}
$$ Soit $p_1 \in X_1$ l'image du point fermé de $T$. Remarquons que $p_1$ est un point fermé, car c'est un point $\kappa = \kappa(p)$-rationnel de la fibre de $X_1 \to X$ au-dessus de $x$. Comme on dispose d'une factorisation $$
\mathcal{O}_{X, x} \to \mathcal{O}_{X_1, p_1} \to R
$$, on voit que $a_1$ est lui aussi un arc formel non singulier.

\medskip\noindent On peut répéter le procédé et obtenir une suite d'éclatements $$
\xymatrix{
T \ar[d]_a \ar[rd]_{a_1} \ar[rrd]_{a_2} \ar[rrrd]^{a_3} \\
(X, p) & (X_1, p_1) \ar[l] & (X_2, p_2) \ar[l] &
(X_3, p_3) \ar[l] & \ldots \ar[l]
}
$$ Les suites d'éclatements de ce type se caractérisent comme suit.

\begin{lemma} \label{resolve-lemma-sequence-blowups} Soit $X$ un schéma localement noethérien. Soit $$
(X, p) = (X_0, p_0) \leftarrow (X_1, p_1) \leftarrow (X_2, p_2) \leftarrow
(X_3, p_3) \leftarrow \ldots
$$ une suite d'éclatements telle que : \begin{enumerate} \item $p_i$ soit fermé, s'envoie sur $p_{i - 1}$ et vérifie $\kappa(p_i) = \kappa(p_{i - 1})$ ; \item il existe un élément $x_1 \in \mathfrak m_p$ dont l'image dans $\mathfrak m_{p_i}$, pour $i > 0$, définit le diviseur exceptionnel $E_i \subset X_i$. \end{enumerate} Alors la suite est obtenue, comme ci-dessus, à partir d'un arc non singulier $a : T \to X$. \end{lemma}

\begin{proof} Écrivons $\mathcal{O}_n = \mathcal{O}_{X_n, p_n}$ et $\mathcal{O} = \mathcal{O}_{X, p}$. Notons $\mathfrak m \subset \mathcal{O}$ et $\mathfrak m_n \subset \mathcal{O}_n$ les idéaux maximaux.

\medskip\noindent Nous affirmons que $x_1^t \not \in \mathfrak m_n^{t + 1}$. En effet, si ce n'était pas le cas, alors, dans l'anneau local $\mathcal{O}_{n + 1}$, l'élément $x_1^t$ appartiendrait à l'idéal de $(t + 1)E_{n + 1}$. Cela contredit l'hypothèse selon laquelle $x_1$ définit $E_{n + 1}$.

\medskip\noindent Pour tout $n$, choisissons des générateurs $y_{n, 1}, \ldots, y_{n, t_n}$ de $\mathfrak m_n$. Comme $\mathfrak m_n \mathcal{O}_{n + 1} = x_1\mathcal{O}_{n + 1}$ d'après l'hypothèse (2), on peut écrire $y_{n, i} = a_{n, i} x_1$ pour certains $a_{n, i} \in \mathcal{O}_{n + 1}$. Puisque l'application $\mathcal{O}_n \to \mathcal{O}_{n + 1}$ induit, d'après (1), un isomorphisme sur les corps résiduels, on peut choisir $c_{n, i} \in \mathcal{O}_n$ ayant la même classe résiduelle que $a_{n, i}$. On voit alors que $$
\mathfrak m_n = (x_1, z_{n, 1}, \ldots, z_{n, t_n}),
\quad z_{n, i} = y_{n, i} - c_{n, i} x_1
$$ et que les éléments $z_{n, i}$ s'envoient sur des éléments de $\mathfrak m_{n + 1}^2$ dans $\mathcal{O}_{n + 1}$.

\medskip\noindent Considérons $$
J_n = \Ker(\mathcal{O} \to \mathcal{O}_n/\mathfrak m_n^{n + 1})
$$ Nous affirmons que $\mathcal{O}/J_n$ est de longueur $n + 1$ et que $\mathcal{O}/(x_1) + J_n$ est égal au corps résiduel. Pour $n = 0$, c'est immédiat. Supposons l'assertion vraie pour $n$. Soit $f \in J_n$. Alors, dans $\mathcal{O}_n$, on a $$
f = a x_1^{n + 1} + x_1^n A_1(z_{n, i}) +
x_1^{n - 1} A_2(z_{n, i}) + \ldots + A_{n + 1}(z_{n, i})
$$ pour un certain $a \in \mathcal{O}_n$ et certains $A_i$ homogènes de degré $i$ à coefficients dans $\mathcal{O}_n$. Puisque $\mathcal{O} \to \mathcal{O}_n$ identifie les corps résiduels, on peut choisir $a \in \mathcal{O}$ (raisonner comme dans la construction de $z_{n, i}$ ci-dessus). En prenant l'image dans $\mathcal{O}_{n + 1}$, on voit que $f$ et $a x_1^{n + 1}$ ont la même image modulo $\mathfrak m_{n + 1}^{n + 2}$. Puisque $x_n^{n + 1} \not \in \mathfrak m_{n + 1}^{n + 2}$, il s'ensuit que $J_n/J_{n + 1}$ est de longueur $1$, ce qui démontre l'affirmation.

\medskip\noindent Considérons $R = \lim \mathcal{O}/J_n$. C'est un quotient du complété $\mathfrak m$-adique de $\mathcal{O}$ ; c'est donc un anneau local noethérien complet. D'autre part, il n'est pas de longueur finie et $x_1$ engendre son idéal maximal. Ainsi $R$ est un anneau de valuation discrète complet. L'application $\mathcal{O} \to R$ se relève en un homomorphisme local $\mathcal{O}_n \to R$ pour tout $n$. Il y a deux façons de le montrer : (1) pour tout $n$, on peut employer un procédé analogue pour construire $\mathcal{O}_n \to R_n$, puis montrer que $\mathcal{O} \to \mathcal{O}_n \to R_n$ se factorise par un isomorphisme $R \to R_n$ ; ou (2) on peut utiliser Diviseurs, lemme \ref{divisors-lemma-characterize-affine-blowup}, pour montrer que $\mathcal{O}_n$ est un localisé d'une algèbre d'éclatement affine itérée et construire explicitement une application $\mathcal{O}_n \to R$. Cela étant dit, il est clair que notre suite d'éclatements provient de l'arc non singulier $a : T = \Spec(R) \to X$. \end{proof}

\noindent Le lemme suivant est une sorte de lemme de désingularisation de N\'eron.

\begin{lemma} \label{resolve-lemma-sequence-blowups-along-arc-becomes-nonsingular} Soit $(A, \mathfrak m, \kappa)$ un anneau local noethérien intègre de dimension $2$. Soit $A \to R$ une surjection sur un anneau de valuation discrète complet. Elle définit un arc non singulier $a : T = \Spec(R) \to \Spec(A)$. Soit $$
\Spec(A) = X_0 \leftarrow X_1 \leftarrow X_2 \leftarrow X_3 \leftarrow \ldots
$$ la suite d'éclatements construite à partir de $a$. Si $A_\mathfrak p$ est un anneau local régulier, où $\mathfrak p = \Ker(A \to R)$, alors, pour un certain $i$, le schéma $X_i$ est régulier en $x_i$. \end{lemma}

\begin{proof} Soit $x_1 \in \mathfrak m$ un élément qui s'envoie sur une uniformisante de $R$. Remarquons que $\kappa(\mathfrak p) = K$ est le corps des fractions de $R$. Écrivons $\mathfrak p = (x_2, \ldots, x_r)$ avec $r$ minimal. Si $r = 2$, alors $\mathfrak m = (x_1, x_2)$ et $A$ est régulier, ce qui démontre le lemme. Supposons $r > 2$. Après renumérotation si nécessaire, nous pouvons supposer que $x_2$ s'envoie sur une uniformisante de $A_\mathfrak p$. Alors $\mathfrak p/\mathfrak p^2 + (x_2)$ est annulé par une puissance de $x_1$. Pour $i > 2$, il existe $n_i \geq 0$ et $a_i \in A$ tels que $$
x_1^{n_i} x_i - a_i x_2 = \sum\nolimits_{2 \leq j \leq k} a_{jk} x_jx_k
$$ pour certains $a_{jk} \in A$. Si $n_i = 0$ pour un certain $i$, on peut supprimer $x_i$ de la liste des générateurs de $\mathfrak p$ et conclure par récurrence sur $r$. Si, pour un certain $i$, l'élément $a_i$ est une unité, on peut supprimer $x_2$ de la liste des générateurs de $\mathfrak p$ et conclure de la même manière. Ainsi, soit $a_i \in \mathfrak p$, soit $a_i = u_i x_1^{m_1} \bmod \mathfrak p$ pour certains $m_1 > 0$ et une unité $u_i \in A$. On se trouve donc dans l'un des deux cas $$
x_1^{n_i} x_i = \sum\nolimits_{2 \leq j \leq k} a_{jk} x_jx_k
\quad\text{ou}\quad
x_1^{n_i} x_i - u_i x_1^{m_i} x_2 =
\sum\nolimits_{2 \leq j \leq k} a_{jk} x_jx_k
$$ Nous allons montrer qu'après éclatement les entiers $n_i$ et $m_i$ diminuent, ce qui achèvera la démonstration.

\medskip\noindent Examinons ce que deviennent ces équations dans l'algèbre d'éclatement affine $A' = A[\mathfrak m/x_1]$. Comme $\mathfrak m = (x_1, \ldots, x_r)$, on voit que $A'$ est engendré sur $R$ par $y_i = x_i/x_1$ pour $i \geq 2$. Manifestement, $A \to R$ se prolonge en $A' \to R$, de noyau $(y_2, \ldots, y_r)$. On voit alors que l'on se trouve dans l'un des deux cas $$
x_1^{n_i - 1} y_i = \sum\nolimits_{2 \leq j \leq k} a_{jk} y_jy_k
\quad\text{ou}\quad
x_1^{n_i - 1} y_i - u_i x_1^{m_1 - 1} y_2 =
\sum\nolimits_{2 \leq j \leq k} a_{jk} y_jy_k
$$, ce qui achève la démonstration. \end{proof}




\section{Changement de base au complété} \label{resolve-section-aux}

\noindent Le lemme élémentaire suivant se révélera utile dans la suite.

\begin{lemma} \label{resolve-lemma-iso-completions} Soit $(A, \mathfrak m, \kappa)$ un anneau local dont l'idéal maximal $\mathfrak m$ est de type fini. Soit $X$ un schéma sur $A$. Soit $Y = X \times_{\Spec(A)} \Spec(A^\wedge)$, où $A^\wedge$ est le complété $\mathfrak m$-adique de $A$. Pour un point $q \in Y$, d'image $p \in X$ situé au-dessus du point fermé de $\Spec(A)$, l'homomorphisme d'anneaux locaux $\mathcal{O}_{X, p} \to \mathcal{O}_{Y, q}$ induit un isomorphisme sur les complétés. \end{lemma}

\begin{proof} Nous pouvons supposer $X$ affine. Écrivons alors $X = \Spec(B)$. Soit $\mathfrak q \subset B' = B \otimes_A A^\wedge$ l'idéal premier correspondant à $q$ et soit $\mathfrak p \subset B$ l'idéal premier correspondant à $p$. D'après Algèbre, lemme \ref{algebra-lemma-hathat-finitely-generated}, on a $$
B'/(\mathfrak m^\wedge)^n B' =
A^\wedge/(\mathfrak m^\wedge)^n \otimes_A B =
A/\mathfrak m^n \otimes_A B = B/\mathfrak m^n B
$$ pour tout $n$. Puisque $\mathfrak m B \subset \mathfrak p$ et $\mathfrak m^\wedge B' \subset \mathfrak q$, on voit que $B/\mathfrak p^n$ et $B'/\mathfrak q^n$ sont tous deux des quotients de l'anneau affiché ci-dessus par la $n$-ième puissance du même idéal premier. Le lemme en résulte. \end{proof}

\begin{lemma} \label{resolve-lemma-port-regularity-to-completion} Soit $(A, \mathfrak m, \kappa)$ un anneau local noethérien. Soit $X \to \Spec(A)$ un morphisme localement de type fini. Posons $Y = X \times_{\Spec(A)} \Spec(A^\wedge)$. Soit $y \in Y$, d'image $x \in X$. Alors : \begin{enumerate} \item si $\mathcal{O}_{Y, y}$ est régulier, alors $\mathcal{O}_{X, x}$ est régulier ; \item si $y$ appartient à la fibre fermée, alors $\mathcal{O}_{Y, y}$ est régulier $\Leftrightarrow \mathcal{O}_{X, x}$ est régulier ; \item si $X$ est propre sur $A$, alors $X$ est régulier si et seulement si $Y$ est régulier. \end{enumerate} \end{lemma}

\begin{proof} Puisque $A \to A^\wedge$ est fidèlement plat (Algèbre, lemme \ref{algebra-lemma-completion-faithfully-flat}), $Y \to X$ est plat. La partie (1) résulte donc d'Algèbre, lemme \ref{algebra-lemma-descent-regular}. Le lemme \ref{resolve-lemma-iso-completions} montre que le morphisme $Y \to X$ induit un isomorphisme sur les anneaux locaux complétés aux points des fibres spéciales. La partie (2) résulte alors de Compléments d'algèbre, lemme \ref{more-algebra-lemma-completion-regular}. Si $X$ est propre sur $A$, alors $Y$ est propre sur $A^\wedge$ (Morphismes, lemme \ref{morphisms-lemma-base-change-proper}), et tout point fermé de $X$ et de $Y$ appartient à la fibre fermée. Par conséquent, $Y$ est un schéma régulier si et seulement si $X$ l'est, d'après Propriétés, lemme \ref{properties-lemma-characterize-regular}. \end{proof}

\begin{lemma} \label{resolve-lemma-descend-admissible-blowup} Soit $(A, \mathfrak m)$ un anneau local noethérien de complété $A^\wedge$. Soient $U \subset \Spec(A)$ et $U^\wedge \subset \Spec(A^\wedge)$ les spectres épointés. Si $Y \to \Spec(A^\wedge)$ est un éclatement $U^\wedge$-admissible, alors il existe un éclatement $U$-admissible $X \to \Spec(A)$ tel que $Y = X \times_{\Spec(A)} \Spec(A^\wedge)$. \end{lemma}

\begin{proof} Par définition, il existe un idéal $J \subset A^\wedge$ tel que $V(J) = \{\mathfrak m A^\wedge\}$ et tel que $Y$ soit l'éclatement de $S^\wedge$ le long du sous-schéma fermé défini par $J$ ; voir Diviseurs, définition \ref{divisors-definition-admissible-blowup}. Comme $A^\wedge$ est noethérien, cela entraîne $\mathfrak m^n A^\wedge \subset J$ pour un certain $n$. Puisque $A^\wedge/\mathfrak m^n A^\wedge = A/\mathfrak m^n$, on trouve un idéal $\mathfrak m^n \subset I \subset A$ tel que $J = I A^\wedge$. Soit $X \to S$ l'éclatement en $I$. Comme $A \to A^\wedge$ est plat, on conclut que le changement de base de $X$ est $Y$ d'après Diviseurs, lemme \ref{divisors-lemma-flat-base-change-blowing-up}. \end{proof}

\begin{lemma} \label{resolve-lemma-blowup-still-good} Soit $(A, \mathfrak m, \kappa)$ un anneau local normal intègre de Nagata et de dimension $2$. Supposons que $A$ définisse une singularité rationnelle et que le complété $A^\wedge$ de $A$ soit normal. Alors : \begin{enumerate} \item $A^\wedge$ définit une singularité rationnelle ; \item si $X \to \Spec(A)$ est l'éclatement en $\mathfrak m$, alors, pour tout point fermé $x \in X$, le complété $\mathcal{O}_{X, x}$ est normal. \end{enumerate} \end{lemma}

\begin{proof} \begingroup\emergencystretch=1em
Soit $Y \to \Spec(A^\wedge)$ une modification avec $Y$ normal. Il faut montrer que $H^1(Y, \mathcal{O}_Y) = 0$. D'après Variétés, lemme \ref{varieties-lemma-modification-normal-iso-over-codimension-1}, $Y \to \Spec(A^\wedge)$ est un isomorphisme au-dessus du spectre épointé $U^\wedge = \Spec(A^\wedge) \setminus \{\mathfrak m^\wedge\}$. D'après le lemme \ref{resolve-lemma-dominate-by-scheme-modification}, il existe un éclatement $U^\wedge$-admissible $Y' \to \Spec(A^\wedge)$ qui domine $Y$. Le lemme \ref{resolve-lemma-descend-admissible-blowup} fournit un éclatement $U$-admissible $X \to \Spec(A)$ dont le changement de base à $A^\wedge$ domine $Y$. Puisque $A$ est de Nagata, on peut remplacer $X$ par sa normalisation ; après quoi $X \to \Spec(A)$ est une modification normale (mais peut ne plus être un éclatement $U$-admissible). Alors $H^1(X, \mathcal{O}_X) = 0$ puisque $A$ définit une singularité rationnelle. Il s'ensuit que $H^1(X \times_{\Spec(A)} \Spec(A^\wedge),
\mathcal{O}_{X \times_{\Spec(A)} \Spec(A^\wedge)}) = 0$ par changement de base plat (Cohomologie des schémas, lemme \ref{coherent-lemma-flat-base-change-cohomology}, et platitude de $A \to A^\wedge$ d'après Algèbre, lemme \ref{algebra-lemma-completion-flat}). On obtient $H^1(Y, \mathcal{O}_Y) = 0$ d'après le lemme \ref{resolve-lemma-exact-sequence}.\par\endgroup

\medskip\noindent Enfin, soit $X \to \Spec(A)$ l'éclatement de $\Spec(A)$ en $\mathfrak m$. Alors $Y = X \times_{\Spec(A)} \Spec(A^\wedge)$ est l'éclatement de $\Spec(A^\wedge)$ en $\mathfrak m^\wedge$. D'après le lemme \ref{resolve-lemma-blow-up-normal-rational}, $Y$ et $X$ sont tous deux normaux. D'autre part, $A^\wedge$ est excellent (Compléments d'algèbre, proposition \ref{more-algebra-proposition-ubiquity-excellent}) ; tout ouvert affine de $Y$ est donc le spectre d'un anneau normal excellent intègre (Compléments d'algèbre, lemme \ref{more-algebra-lemma-finite-type-over-excellent}). Ainsi, pour $y \in Y$, l'homomorphisme d'anneaux $\mathcal{O}_{Y, y} \to \mathcal{O}_{Y, y}^\wedge$ est régulier et Compléments d'algèbre, lemme \ref{more-algebra-lemma-normal-goes-up}, montre que $\mathcal{O}_{Y, y}^\wedge$ est normal. Si $x \in X$ est un point fermé de la fibre spéciale, il existe un unique point fermé $y \in Y$ au-dessus de $x$. Puisque $\mathcal{O}_{X, x} \to \mathcal{O}_{Y, y}$ induit un isomorphisme sur les complétés (lemme \ref{resolve-lemma-iso-completions}), on conclut. \end{proof}

\begin{lemma} \label{resolve-lemma-formally-unramified} Soit $(A, \mathfrak m)$ un anneau local noethérien. Soit $X$ un schéma sur $A$. Supposons que : \begin{enumerate} \item $A$ soit analytiquement non ramifié (Algèbre, définition \ref{algebra-definition-analytically-unramified}) ; \item $X$ soit localement de type fini sur $A$ ; \item $X \to \Spec(A)$ soit \'etale aux points génériques des composantes irréductibles de $X$. \end{enumerate} Alors la normalisation de $X$ est finie sur $X$. \end{lemma}

\begin{proof} Puisque $A$ est analytiquement non ramifié, il est réduit d'après Algèbre, lemme \ref{algebra-lemma-analytically-unramified-easy}. Comme la normalisation de $X$ ne dépend que du réduit de $X$, nous pouvons remplacer $X$ par son réduit $X_{red}$ ; remarquons que $X_{red} \to X$ est un isomorphisme sur l'ouvert $U$ où $X \to \Spec(A)$ est \'etale, car $U$ est réduit (Descente, lemme \ref{descent-lemma-reduced-local-smooth}), de sorte que la condition (3) reste vraie après ce remplacement. De plus, nous pouvons supposer, et nous le faisons, que $X = \Spec(B)$ est affine.

\medskip\noindent L'application $$
K = \prod\nolimits_{\mathfrak p \subset A\text{ minimal}} \kappa(\mathfrak p)
\longrightarrow
K^\wedge = \prod\nolimits_{\mathfrak p^\wedge \subset A^\wedge\text{ minimal}}
\kappa(\mathfrak p^\wedge)
$$ est injective car $A \to A^\wedge$ est fidèlement plat (Algèbre, lemme \ref{algebra-lemma-completion-faithfully-flat}) et induit donc une application surjective entre les ensembles d'idéaux premiers minimaux (par descente pour les homomorphismes plats, voir Algèbre, section \ref{algebra-section-going-up}). Les deux membres sont des produits finis de corps puisque nos anneaux sont noethériens. Posons $L = \prod_{\mathfrak q \subset B\text{ minimal}} \kappa(\mathfrak q)$. Notre hypothèse (3) entraîne que $L = B \otimes_A K$ et que $K \to L$ est un homomorphisme d'anneaux fini \'etale (cela tient au fait que $A \to B$ est génériquement fini ; on peut par exemple utiliser Algèbre, lemme \ref{algebra-lemma-generically-finite}, ou les résultats plus détaillés de Morphismes, section \ref{morphisms-section-generically-finite}). Puisque $B$ est réduit, on voit que $B \subset L$. Cela implique que $$
C = B \otimes_A A^\wedge \subset
L \otimes_A A^\wedge = L \otimes_K K^\wedge = M
$$. Alors $M$ est l'anneau total des fractions de $C$ et est un produit fini de corps, car c'est une algèbre séparable finie sur $K^\wedge$. Il s'ensuit que $C$ est réduit et que sa normalisation $C'$ est la clôture intégrale de $C$ dans $M$. La normalisation $B'$ de $B$ est la clôture intégrale de $B$ dans $L$. Par platitude de $A \to A^\wedge$, on obtient une application injective $B' \otimes_A A^\wedge \to M$ dont l'image est contenue dans $C'$. On a ainsi $$
B' \otimes_A A^\wedge \longrightarrow C'
$$. Comme $A^\wedge$ est de Nagata (d'après Algèbre, lemme \ref{algebra-lemma-Noetherian-complete-local-Nagata}), on voit que $C'$ est fini sur $C = B \otimes_A A^\wedge$ (voir Algèbre, lemmes \ref{algebra-lemma-Noetherian-complete-local-Nagata} et \ref{algebra-lemma-nagata-in-reduced-finite-type-finite-integral-closure}). Puisque $C$ est noethérien, on conclut que $B' \otimes_A A^\wedge$ est fini sur $C = B \otimes_A A^\wedge$. Par descente fidèlement plate (Algèbre, lemme \ref{algebra-lemma-descend-properties-modules}), on voit donc que $B'$ est fini sur $B$, ce qu'il fallait démontrer. \end{proof}

\begin{lemma} \label{resolve-lemma-normalization-completion} Soit $(A, \mathfrak m, \kappa)$ un anneau local noethérien. Soit $X \to \Spec(A)$ un morphisme localement de type fini. Posons $Y = X \times_{\Spec(A)} \Spec(A^\wedge)$. Si le complémentaire de la fibre spéciale dans $Y$ est normal, alors la normalisation $X^\nu \to X$ est finie et le changement de base de $X^\nu$ à $\Spec(A^\wedge)$ redonne la normalisation de $Y$. \end{lemma}

\begin{proof} On se ramène immédiatement au cas où $X = \Spec(B)$ est affine, avec $B$ une $A$-algèbre de type fini. Posons $C = B \otimes_A A^\wedge$, de sorte que $Y = \Spec(C)$. Puisque $A \to A^\wedge$ est fidèlement plat, pour tout idéal premier $\mathfrak q \subset B$ il existe un idéal premier $\mathfrak r \subset C$ au-dessus de $\mathfrak q$. Alors $B_\mathfrak q \to C_\mathfrak r$ est fidèlement plat. Par conséquent, si $\mathfrak q$ n'est pas situé au-dessus de $\mathfrak m$, alors $C_\mathfrak r$ est normal par hypothèse sur $Y$, et l'on conclut que $B_\mathfrak q$ est normal d'après Algèbre, lemme \ref{algebra-lemma-descent-normal}. On voit ainsi que $X$ est normal en dehors de la fibre spéciale.

\medskip\noindent Rappelons que l'anneau local noethérien complet $A^\wedge$ est de Nagata (Algèbre, lemme \ref{algebra-lemma-Noetherian-complete-local-Nagata}). La normalisation $Y^\nu \to Y$ est donc finie (Morphismes, lemme \ref{morphisms-lemma-nagata-normalization}) et est un isomorphisme en dehors de la fibre spéciale. Écrivons $Y^\nu = \Spec(C')$. Alors $C \to C'$ est fini et est un isomorphisme en dehors de $V(\mathfrak m C)$. Puisque $B \to C$ est plat et induit un isomorphisme $B/\mathfrak m B \to C/\mathfrak m C$, il existe un homomorphisme fini d'anneaux $B \to B'$ dont le changement de base à $C$ redonne $C \to C'$. Voir Compléments d'algèbre, lemme \ref{more-algebra-lemma-application-formal-glueing} et remarque \ref{more-algebra-remark-formal-glueing-algebras}. On obtient ainsi un morphisme fini $X' \to X$ qui est un isomorphisme en dehors de la fibre spéciale et dont le changement de base redonne $Y^\nu \to Y$. D'après la discussion du premier paragraphe, $X'$ est normal aux points qui ne sont pas sur la fibre spéciale. À un point $x \in X'$ de la fibre spéciale correspondent un point $y \in Y^\nu$ et une application plate $\mathcal{O}_{X', x} \to \mathcal{O}_{Y^\nu, y}$. Puisque $\mathcal{O}_{Y^\nu, y}$ est normal, $\mathcal{O}_{X', x}$ l'est aussi ; voir Algèbre, lemme \ref{algebra-lemma-descent-normal}. Ainsi $X'$ est normal, et il s'ensuit que c'est la normalisation de $X$. \end{proof}

\begin{lemma} \label{resolve-lemma-normalized-blowup-completion} Soit $(A, \mathfrak m, \kappa)$ un anneau local noethérien intègre dont le complété $A^\wedge$ est normal. Alors, pour toute suite $$
Y_n \to Y_{n - 1} \to \ldots \to Y_1 \to \Spec(A^\wedge)
$$ d'éclatements normalisés, il existe une suite d'éclatements normalisés (propres) $$
X_n \to X_{n - 1} \to \ldots \to X_1 \to \Spec(A)
$$ dont le changement de base à $A^\wedge$ redonne la suite donnée. \end{lemma}

\begin{proof} Étant donnée la suite $Y_n \to \ldots \to Y_1 \to Y_0 = \Spec(A^\wedge)$, construisons par récurrence $X_n \to \ldots \to X_1 \to X_0 = \Spec(A)$. Le cas initial est $i = 0$. Étant donné $X_i$ dont le changement de base est $Y_i$, soit $Y'_i \to Y_i$ l'éclatement au point fermé $y_i \in Y_i$ tel que $Y_{i + 1}$ soit la normalisation de $Y_i$. Puisque les fibres fermées de $Y_i$ et de $X_i$ sont isomorphes, le point $y_i$ correspond à un point fermé $x_i$ de la fibre spéciale de $X_i$. Soit $X'_i \to X_i$ l'éclatement de $X_i$ en $x_i$. Alors le changement de base de $X'_i$ à $\Spec(A^\wedge)$ est isomorphe à $Y'_i$. D'après le lemme \ref{resolve-lemma-normalization-completion}, la normalisation $X_{i + 1} \to X'_i$ est finie et son changement de base à $\Spec(A^\wedge)$ est isomorphe à $Y_{i + 1}$. \end{proof}














\section{Points doubles rationnels} \label{resolve-section-rational-double-points}

\noindent Dans la section \ref{resolve-section-rational-singularities}, nous avons montré que la résolution des singularités rationnelles de dimension $2$ se ramène au cas de Gorenstein. Une singularité rationnelle de surface qui est de Gorenstein est un point double rationnel. Nous allons les résoudre par des calculs explicites.

\medskip\noindent D'après la discussion d'Exemples, section \ref{examples-section-bad}, il existe un anneau local noethérien normal intègre $A$ dont le complété est isomorphe à $\mathbf{C}[[x, y, z]]/(z^2)$. Dans ce cas, on pourrait dire que $A$ définit un point double rationnel ; en revanche, $\Spec(A)$ n'admet pas de résolution des singularités. Ce type de comportement ne peut pas se produire si $A$ est un anneau de Nagata ; voir Algèbre, lemme \ref{algebra-lemma-local-nagata-domain-analytically-unramified}.

\medskip\noindent La situation est toutefois pire encore : il existe un anneau local normal intègre de Nagata $A$ dont le complété est $\mathbf{C}[[x, y, z]]/(yz)$, et un autre dont le complété est $\mathbf{C}[[x, y, z]]/(y^2 - z^3)$. Il s'agit de l'exemple 2.5 de \cite{Nishimura-few}. C'est pourquoi, dans cette section, nous devons supposer que le complété de notre anneau est normal.

\begin{situation} \label{resolve-situation-rational-double-point} Dans ce qui suit, $(A, \mathfrak m, \kappa)$ est un anneau local normal intègre de Nagata, de dimension $2$, qui définit une singularité rationnelle, dont le complété est normal et qui est de Gorenstein. Nous supposons que $A$ n'est pas régulier. \end{situation}

\noindent Les arguments de cette section montreront que, dans cette situation, l'éclatement répété des points singuliers résout $\Spec(A)$. Au cours de la démonstration, nous aurons besoin du lemme suivant.

\begin{lemma} \label{resolve-lemma-issquare} Soit $\kappa$ un corps. Soit $I \subset \kappa[x, y]$ un idéal. Supposons que $$
a + b x + c y + d x^2 + exy + f y^2 \in I^2
$$ pour des $a, b, c, d, e, f \in k$ non tous nuls. Si la colongueur de $I$ dans $\kappa[x, y]$ est $> 1$, alors $a + b x + c y + d x^2 + exy + f y^2 = j(g + hx + iy)^2$ pour certains $j, g, h, i \in \kappa$. \end{lemma}

\begin{proof} Considérons les dérivées partielles $b + 2dx + ey$ et $c + ex + 2fy$. D'après les règles de Leibniz, elles appartiennent à $I$. Si l'une d'elles est non nulle, alors, après un changement linéaire de coordonnées, c'est-à-dire de la forme $x \mapsto \alpha + \beta x + \gamma y$ et $y \mapsto \delta + \epsilon x + \zeta y$, nous pouvons supposer que $x \in I$. On voit alors que $I = (x)$ ou que $I = (x, F)$, où $F$ est un polynôme unitaire de degré $\geq 2$ en $y$. Dans le premier cas, l'assertion est claire. Dans le second, remarquons que tout élément de $I^2$ peut s'écrire sous la forme $$
A(x, y) x^2 + B(y) x F + C(y) F^2
$$ pour certains $A(x, y) \in \kappa[x, y]$ et $B, C \in \kappa[y]$. Ainsi $$
a + b x + c y + d x^2 + exy + f y^2 = A(x, y) x^2 + B(y) x F + C(y) F^2
$$ et, par des considérations de degré, on voit que $B = C = 0$ et que $A$ est une constante.

\medskip\noindent Pour achever la démonstration, il reste à traiter le cas où les deux dérivées partielles sont nulles. Cela ne peut se produire qu'en caractéristique $2$, et l'on obtient alors $$
a + d x^2 + f y^2 \in I^2
$$. Nous pouvons supposer $f$ non nul (sinon, intervertissons les rôles de $x$ et de $y$). Après division par $f$, nous nous ramenons au cas où la caractéristique de $\kappa$ est $2$ et où $$
a + d x^2 + y^2 \in I^2
$$. Si $a$ et $d$ sont des carrés dans $\kappa$, c'est terminé. Sinon, il existe une dérivation $\theta : \kappa \to \kappa$ telle que $\theta(a) \not = 0$ ou $\theta(d) \not = 0$ ; voir Algèbre, lemme \ref{algebra-lemma-derivative-zero-pth-power}. Nous pouvons la prolonger en une dérivation de $\kappa[x, y]$ en posant $\theta(x) = \theta(y) = 0$. On trouve alors que $$
\theta(a) + \theta(d) x^2 \in I
$$. Le cas $\theta(d) = 0$ est absurde. Nous pouvons donc supposer que $\alpha + x^2 \in I$ pour un certain $\alpha \in \kappa$. En combinant cela avec ce qui précède, on obtient $a + \alpha d + y^2 \in I$. Par conséquent, $$
J = (\alpha + x^2, a + \alpha d + y^2) \subset I
$$, de codimension au plus $2$. Remarquons que $J/J^2$ est libre sur $\kappa[x, y]/J$, de base $\alpha + x^2$ et $a + \alpha d + y^2$. Ainsi, $a + d x^2 + y^2 =
1 \cdot (a + \alpha d + y^2) + d \cdot (\alpha + x^2) \in I^2$ implique que l'inclusion $J \subset I$ est stricte. On trouve donc un élément non nul de la forme $g + hx + iy + jxy$ dans $I$. Si $j = 0$, alors $I$ contient une forme linéaire et l'on peut conclure comme dans le premier paragraphe. Ainsi $j \not = 0$ et $\dim_\kappa(I/J) = 1$ (sinon, on pourrait trouver dans $I$ un élément comme ci-dessus avec $j = 0$). On conclut que $I$ est de la forme $(\alpha + x^2, \beta + y^2, g + hx + iy + jxy)$, avec $j \not = 0$, et de colongueur $3$. Dans ce cas, $a + dx^2 + y^2 \in I^2$ est impossible. On peut le montrer par un calcul direct, mais nous préférons raisonner comme suit. Pour établir cette assertion, nous pouvons supposer que $\kappa$ est algébriquement clos. Nous pouvons alors effectuer le changement de coordonnées $x \mapsto \sqrt{\alpha} + x$ et $y \mapsto \sqrt{\beta} + y$ et supposer que $I = (x^2, y^2, g' + h'x + i'y + jxy)$, avec le même $j$. Alors $g' = h' = i' = 0$, faute de quoi la colongueur de $I$ ne serait pas $3$. On obtient ainsi $I = (x^2, y^2, xy)$, et le résultat est clair. \end{proof}

\noindent Soit $(A, \mathfrak m, \kappa)$ comme dans la situation \ref{resolve-situation-rational-double-point}. Soit $X \to \Spec(A)$ l'éclatement en $\mathfrak m$ de $\Spec(A)$. D'après le lemme \ref{resolve-lemma-blow-up-normal-rational}, $X$ est normal. Toutes les singularités de $X$ sont rationnelles d'après le lemme \ref{resolve-lemma-rational-propagates}. Puisque $\omega_A = A$, le lemme \ref{resolve-lemma-dualizing-blow-up-rational} montre que $\omega_X \cong \mathcal{O}_X$ (voir la discussion de la remarque \ref{resolve-remark-dualizing-setup} pour les conventions). Toutes les singularités de $X$ sont donc de Gorenstein. En outre, les anneaux locaux de $X$ aux points fermés ont des complétés normaux d'après le lemme \ref{resolve-lemma-blowup-still-good}. Autrement dit, en éclatant $\Spec(A)$, on obtient une surface normale $X$ dont les points singuliers sont comme dans la situation \ref{resolve-situation-rational-double-point}. Nous l'utiliserons ci-dessous sans plus de commentaires. (Remarque : nous verrons au cours de la discussion ci-dessous que ces points singuliers sont en nombre fini.)

\medskip\noindent Soit $E \subset X$ le diviseur exceptionnel. On a $\omega_E = \mathcal{O}_E(-1)$ d'après le lemme \ref{resolve-lemma-dualizing-blow-up-rational}. Le lemme \ref{resolve-lemma-cohomology-blow-up-rational} donne $\kappa = H^0(E, \mathcal{O}_E)$. Ainsi, $E$ est une courbe de Gorenstein et, par Riemann--Roch tel qu'il est présenté dans Courbes algébriques, section \ref{curves-section-Riemann-Roch}, on a $$
\chi(E, \mathcal{O}_E) = 1 - g = -(1/2) \deg(\omega_E) =
(1/2)\deg(\mathcal{O}_E(1))
$$, où $g = \dim_\kappa H^1(E, \mathcal{O}_E) \geq 0$. Puisque $\deg(\mathcal{O}_E(1))$ est positif d'après Variétés, lemme \ref{varieties-lemma-ampleness-in-terms-of-degrees-components}, on trouve que $g = 0$ et $\deg(\mathcal{O}_E(1)) = 2$. Il s'ensuit que $$
\dim_\kappa (\mathfrak m^n/\mathfrak m^{n + 1}) = 2n + 1
$$ d'après le lemme \ref{resolve-lemma-cohomology-blow-up-rational} et Riemann--Roch sur $E$.

\medskip\noindent \begingroup\emergencystretch=2em
Choisissons $x_1, x_2, x_3 \in \mathfrak m$ qui s'envoient sur une base de $\mathfrak m/\mathfrak m^2$. Comme $\dim_\kappa(\mathfrak m^2/\mathfrak m^3) = 5$, les images de $x_i x_j$, $i \geq j$ dans cet espace vectoriel sur $\kappa$ satisfont une relation. Autrement dit, on peut trouver $a_{ij} \in A$, $i \geq j$, non tous contenus dans $\mathfrak m$, tels que\par\endgroup
$$
a_{11} x_1^2 + a_{12} x_1x_2 + a_{13}x_1x_3 + a_{22} x_2^2 +
a_{23} x_2x_3 + a_{33} x_3^2 =
\sum a_{ijk} x_ix_jx_k
$$ pour un certain $a_{ijk} \in A$, où $i \leq j \leq k$. Notons $a \mapsto \overline{a}$ l'application $A \to \kappa$. La forme quadratique $q = \sum \overline{a}_{ij} t_i t_j \in \kappa[t_1, t_2, t_3]$ est bien définie, compte tenu de nos choix, à multiplication près par un élément de $\kappa^*$. Si, au cours de nos arguments, nous constatons que $\overline{a}_{ij} = 0$ dans $\kappa$, nous pouvons absorber le terme $a_{ij} x_i x_j$ dans le membre de droite et supposer que $a_{ij} = 0$ ; cette opération modifie les $a_{ijk}$, mais pas les autres $a_{i'j'}$.

\medskip\noindent L'éclatement est recouvert par $3$ cartes affines correspondant aux « variables » $x_1, x_2, x_3$. Par symétrie, il suffit d'étudier l'une de ces cartes. Pour cela, soit $$
A' = A[\mathfrak m/x_1]
$$ l'algèbre affine d'éclatement (comme dans Algèbre, section \ref{algebra-section-blow-up}). Puisque $x_1, x_2, x_3$ engendrent $\mathfrak m$, on voit que $A'$ est engendrée par $y_2 = x_2/x_1$ et $y_3 = x_3/x_1$ sur $A$. Nous utiliserons parfois $y_1 = 1$ pour simplifier les formules. De plus, en considérant la relation ci-dessus, on trouve que $$
a_{11} + a_{12} y_2 + a_{13} y_3 + a_{22} y_2^2 +
a_{23} y_2y_3 + a_{33} y_3^2 =
x_1 (\sum a_{ijk} y_iy_jy_k)
$$ dans $A'$. Rappelons que $x_1 \in A'$ définit le diviseur exceptionnel $E$ sur cet ouvert affine de $X$, lequel est donc donné schématiquement par $$
\kappa[y_2, y_3]/
(\overline{a}_{11} + \overline{a}_{12} y_2 + \overline{a}_{13} y_3 +
\overline{a}_{22} y_2^2 + \overline{a}_{23} y_2y_3 + \overline{a}_{33} y_3^2)
$$. Autrement dit, $E \subset \mathbf{P}^2_\kappa = \text{Proj}(\kappa[t_1, t_2, t_3])$ est le sous-schéma des zéros de la forme quadratique $q$ introduite ci-dessus.

\medskip\noindent La forme quadratique $q$ est un invariant important de la singularité définie par $A$. Disons que nous sommes dans le {\bf cas II} si $q$ est le carré d'une forme linéaire multiplié par un élément de $\kappa^*$, et dans le {\bf cas I} sinon. Remarquons que nous sommes dans le cas II exactement lorsque, après modification de notre choix de $x_1, x_2, x_3$, on a $$
x_3^2 = \sum a_{ijk}x_ix_jx_k
$$ dans l'anneau local $A$.

\medskip\noindent Soit $\mathfrak m' \subset A'$ un idéal maximal situé au-dessus de $\mathfrak m$, de corps résiduel $\kappa'$. Autrement dit, $\mathfrak m'$ correspond à un point fermé $p \in E$ du diviseur exceptionnel. Rappelons que la surjection $$
\kappa[y_2, y_3] \to \kappa'
$$ a un noyau engendré par deux éléments $f_2, f_3 \in \kappa[y_2, y_3]$ (voir par exemple Algèbre, exemple \ref{algebra-example-spec-kxy}, ou la démonstration d'Algèbre, lemme \ref{algebra-lemma-dim-affine-space}). Soient $z_2, z_3 \in A'$ des éléments qui s'envoient respectivement sur $f_2, f_3$ dans $\kappa[y_2, y_3]$. On voit alors que $\mathfrak m' = (x_1, z_2, z_3)$, car $x_2$ et $x_3$ deviennent divisibles par $x_1$ dans $A'$.

\medskip\noindent {\bf Assertion.} Si $X$ est singulier en $p$, alors $\kappa' = \kappa$ ou bien nous sommes dans le cas II. En effet, si $A'_{\mathfrak m'}$ est singulier, alors $\dim_{\kappa'} \mathfrak m'/(\mathfrak m')^2 = 3$, ce qui implique que $\dim_{\kappa'} \overline{\mathfrak m}'/(\overline{\mathfrak m}')^2 = 2$, où $\overline{m}'$ est l'idéal maximal de $\mathcal{O}_{E, p} = \mathcal{O}_{X, p}/x_1\mathcal{O}_{X, p}$. Cela entraîne que $$
q(1, y_2, y_3) =
\overline{a}_{11} + \overline{a}_{12} y_2 + \overline{a}_{13} y_3 +
\overline{a}_{22} y_2^2 + \overline{a}_{23} y_2y_3 + \overline{a}_{33} y_3^2
\in (f_2, f_3)^2
$$, sans quoi il existerait une relation entre les classes de $z_2$ et de $z_3$ dans $\overline{\mathfrak m}'/(\overline{\mathfrak m}')^2$. L'assertion résulte maintenant du lemme \ref{resolve-lemma-issquare}.

\medskip\noindent Résolution dans le cas I. D'après l'assertion, tout point singulier de $X$ est $\kappa$-rationnel. Choisissons un tel point singulier $p$. Nous pouvons choisir nos $x_1, x_2, x_3 \in \mathfrak m$ de sorte que $p$ appartienne à la carte décrite ci-dessus et ait pour coordonnées $y_2 = y_3 = 0$. Comme c'est un point singulier, un raisonnement analogue à celui de la démonstration de l'assertion donne $q(1, y_2, y_3) \in (y_2, y_3)^2$. Nous pouvons donc choisir $a_{11} = a_{12} = a_{13} = 0$ et $q(t_1, t_2, t_3) = q(t_2, t_3)$. Il s'ensuit que $$
E = V(q) \subset \mathbf{P}^1_\kappa
$$ est soit la réunion de deux droites distinctes qui se rencontrent en $p$, soit une courbe de degré $2$ possédant un unique point $\kappa$-rationnel (nous omettons un petit détail ; utiliser que $q$ n'est pas, à un scalaire près, le carré d'une forme linéaire). Dans les deux cas, on conclut que $X$ possède un unique point singulier $p$, qui est $\kappa$-rationnel. Nous avons besoin d'un peu plus d'information dans ce cas. Tout d'abord, en examinant les termes de degré supérieur dans l'expression ci-dessus, on trouve que $\overline{a}_{111} = 0$, car $p$ est singulier. On peut alors écrire $a_{111} = b_{111} x_1 \bmod (x_2, x_3)$ pour un certain $b_{111} \in A$. La forme quadratique en $p$ associée aux générateurs $x_1, y_2, y_3$ de $\mathfrak m'$ est alors $$
q' =
\overline{b}_{111} t_1^2 +
\overline{a}_{112} t_1 t_2 + \overline{a}_{113} t_1 t_3 +
\overline{a}_{22} t_2^2 +
\overline{a}_{23} t_2 t_3 +
\overline{a}_{33} t_3^2
$$. On voit que $E' = V(q')$ rencontre la droite $t_1 = 0$ soit en deux points, soit en un point de degré $2$. On conclut que $p$ relève du cas I.

\medskip\noindent Supposons que l'éclatement $X' \to X$ de $X$ en $p$ ait de nouveau un point singulier $p'$. On voit alors que $p'$ est un point $\kappa$-rationnel et l'on peut éclater pour obtenir $X'' \to X'$. Si ce procédé ne s'arrête pas, on obtient une suite d'éclatements $$
\Spec(A) \leftarrow X \leftarrow X' \leftarrow X'' \leftarrow \ldots
$$. Nous voulons montrer que le lemme \ref{resolve-lemma-sequence-blowups} s'applique à cette situation. Pour cela, nous devons préciser le choix de l'élément $x_1$ de $\mathfrak m$. Supposons que $A$ relève du cas I et que $X$ ait un point singulier. Nous dirons que $x_1 \in \mathfrak m$ est une {\it bonne coordonnée} si, pour tout choix de $x_2, x_3$ (ou, de manière équivalente, pour un certain choix), la forme quadratique $q(t_1, t_2, t_3)$ possède la propriété que $q(0, t_2, t_3)$ n'est pas un carré multiplié par un scalaire. Nous avons vu ci-dessus qu'il existe une bonne coordonnée. Si $x_1$ est une bonne coordonnée, alors le point singulier $p \in E$ de $X$ n'appartient pas à l'hypersurface $t_1 = 0$, car celle-ci soit ne possède pas de point rationnel, soit, si elle en possède un, n'est pas singulière en ce point sur $X$. Remarquons que cela équivaut à dire que l'image de $x_1$ dans $\mathcal{O}_{X, p}$ définit le diviseur exceptionnel $E$. Les calculs ci-dessus montrent maintenant que, si $x_1$ est une bonne coordonnée pour $A$, alors $x_1 \in \mathfrak m'\mathcal{O}_{X, p}$ est une bonne coordonnée pour $p$. Cela utilise bien sûr le fait que la notion de bonne coordonnée ne dépend pas du choix de $x_2$, $x_3$ employé pour effectuer le calcul. Ainsi, $x_1$ s'envoie sur une bonne coordonnée en $p'$, en $p''$, etc. Le lemme \ref{resolve-lemma-sequence-blowups} s'applique donc, et notre suite d'éclatements provient d'un arc non singulier $A \to R$. Alors l'application $A^\wedge \to R$ est surjective. Puisque le complété de $A$ est normal, le lemme \ref{resolve-lemma-sequence-blowups-along-arc-becomes-nonsingular} montre qu'après un nombre fini d'éclatements $$
\Spec(A^\wedge) \leftarrow X^\wedge \leftarrow (X')^\wedge \leftarrow \ldots
$$, le schéma obtenu $(X^{(n)})^\wedge$ est régulier. Comme $(X^{(n)})^\wedge \to X^{(n)}$ induit des isomorphismes sur les anneaux locaux complétés (lemme \ref{resolve-lemma-iso-completions}), on conclut qu'il en va de même pour $X^{(n)}$.

\medskip\noindent Résolution dans le cas II. On a ici $$
x_3^2 = \sum a_{ijk}x_ix_jx_k
$$ dans $A$ pour un certain choix de générateurs $x_1, x_2, x_3$ de $\mathfrak m$. Alors $q = t_3^2$ et $E = 2C$, où $C$ est une droite. Rappelons que, dans $A'$, on obtient $$
y_3^2 = x_1(\sum a_{ijk} y_iy_jy_k)
$$. Puisque nous savons que $X$ est normal, l'anneau $\mathcal{O}_{X, \xi}$ au point générique $\xi$ de $C$ est un anneau de valuation discrète. L'élément $y_3 \in A'$ s'envoie sur une uniformisante de $\mathcal{O}_{X, \xi}$. Comme $x_1$ définit schématiquement $E$, qui est $C$ avec multiplicité $2$, on voit que $x_1$ est égal à une unité multipliée par $y_3^2$ dans $\mathcal{O}_{X, \xi}$. En considérant l'égalité ci-dessus, on conclut que $$
h(y_2) = \overline{a}_{111} + \overline{a}_{112} y_2 +
\overline{a}_{122} y_2^2 +
\overline{a}_{222} y_2^3
$$ doit être non nul dans le corps résiduel de $\xi$. Supposons maintenant que $p \in C$ définisse un point singulier. Alors $y_3$ s'annule en $p$ et $p$ doit correspondre à un zéro de $h$, d'après le raisonnement employé dans la démonstration de l'assertion ci-dessus. Si $h$ n'a pas un zéro double en $p$, alors la forme quadratique $q'$ en $p$ n'est pas un carré, et l'on conclut que $p$ relève du cas I, traité ci-dessus\footnote{L'idéal maximal en $p$ dans $A'$ est engendré par $y_3, x_1$ et par un troisième élément $g$ dont l'image dans $\kappa[y_2]$ est le diviseur premier de $h$ correspondant à $p$. Si ce diviseur premier ne divise pas $h$ deux fois, on voit alors que la forme quadratique en $p$ est de la forme $$
y_3^2 - x_1((\text{un élément})x_1 + (\text{un élément})y_3 + (\text{une unité})g)
$$, ce qui ne peut jamais être un carré dans $\kappa[y_3, x_1, g]$.}. Comme le degré de $h$ est $3$, il y a au plus un point singulier $p \in C$ relevant du cas II ; de plus, ce point est $\kappa$-rationnel. Après modification de notre choix de $x_1, x_2, x_3$, nous pouvons supposer qu'il s'agit du point $y_2 = y_3 = 0$. Alors $h = \overline{a}_{122} y_2^2 + \overline{a}_{222} y_2^3$. En outre, il faut encore que $\overline{a}_{113} = 0$ pour que la forme quadratique $q'$ ait la forme requise. Ainsi, l'anneau local $\mathcal{O}_{X, p}$ définit une singularité comme au paragraphe suivant.

\medskip\noindent Le dernier cas à traiter est celui où l'on peut choisir nos générateurs $x_1, x_2, x_3$ de $\mathfrak m$ de sorte que $$
x_3^2 + x_1(a x_2^2 + b x_2x_3 + c x_3^2) \in \mathfrak m^4
$$ pour un certain $a, b, c \in A$. C'est une sous-classe du cas II. Si $\overline{a} = 0$, alors on peut écrire $a = a_1 x_1 + a_2 x_2 + a_3 x_3$ et, après éclatement, on obtient $$
y_3^2 + x_1(a_1 x_1 y_2^2 + a_2 x_1 y_2^3 + a_3 x_1 y_2^2 y_3 +
b y_2 y_3 + c y_3^2) = x_1^2 (\sum a_{ijkl}y_iy_jy_ky_l)
$$. Cela signifie que $X$ n'est pas normal\footnote{En effet, l'équation montre qu'il est singulier le long du lieu de dimension $1$ donné par $x_1 = y_3 = 0$, ce qui est impossible pour une surface normale.}, contradiction. D'après le résultat du paragraphe précédent, si l'éclatement $X$ possède un point singulier $p$ qui relève du cas II, alors ce point est unique et $\kappa$-rationnel. En calculant les algèbres affines d'éclatement $A[\frac{\mathfrak m}{x_2}]$ et $A[\frac{\mathfrak m}{x_3}]$, le lecteur voit aisément que $p$ ne peut pas appartenir aux ouverts correspondants de $X$. Ainsi, $p$ appartient au spectre de $A[\frac{\mathfrak m}{x_1}]$. En effectuant l'éclatement comme précédemment, on voit que $p$ doit être le point de coordonnées $y_2 = y_3 = 0$ et que la nouvelle équation est de la forme $$
y_3^2 + x_1(a y_2^2 + b y_2 y_3 + c y_3^2) \in (\mathfrak m')^4
$$, qui a la même forme qu'auparavant et possède la propriété que $x_1$ définit le diviseur exceptionnel. Ainsi, si le procédé ne s'arrête pas, on obtient une suite infinie d'éclatements, et sur chacun d'eux $x_1$ définit le diviseur exceptionnel dans l'anneau local du point singulier. On peut donc achever la démonstration à l'aide des lemmes \ref{resolve-lemma-sequence-blowups} et \ref{resolve-lemma-sequence-blowups-along-arc-becomes-nonsingular} et du même raisonnement que précédemment.

\begin{lemma} \label{resolve-lemma-resolve-rational-double-points} Soit $(A, \mathfrak m, \kappa)$ un anneau local normal intègre de Nagata, de dimension $2$, qui définit une singularité rationnelle, dont le complété est normal et qui est de Gorenstein. Il existe alors une suite finie d'éclatements en des points fermés singuliers $$
X_n \to X_{n - 1} \to \ldots \to X_1 \to X_0 = \Spec(A)
$$ telle que $X_n$ soit régulier et que chacun des schémas intermédiaires $X_i$ soit normal et ne possède qu'un nombre fini de points singuliers du même type. \end{lemma}

\begin{proof} C'est exactement ce qui a été démontré dans la discussion ci-dessus. \end{proof}




\section{Propriétés impliquées} \label{resolve-section-existence-gives}

\noindent Dans cette section, nous montrons que, pour un schéma intègre noethérien, l'existence d'une altération régulière a de nombreuses conséquences. Les lecteurs qui ne s'intéressent pas aux « mauvais » anneaux noethériens peuvent omettre cette section.

\begin{lemma} \label{resolve-lemma-regular-alteration-implies} Soit $Y$ un schéma intègre noethérien. Supposons qu'il existe une altération $f : X \to Y$ avec $X$ régulier. Alors la normalisation $Y^\nu \to Y$ est finie et $Y$ possède un ouvert dense régulier. \end{lemma}

\begin{proof} Il suffit de le démontrer lorsque $Y = \Spec(A)$, où $A$ est un anneau noethérien intègre. Soit $B$ la clôture intégrale de $A$ dans son corps des fractions. Posons $C = \Gamma(X, \mathcal{O}_X)$. D'après Cohomologie des schémas, lemme \ref{coherent-lemma-proper-over-affine-cohomology-finite}, on voit que $C$ est un $A$-module fini. Comme $X$ est normal (Propriétés, lemme \ref{properties-lemma-regular-normal}), on voit que $C$ est un anneau normal intègre (Propriétés, lemme \ref{properties-lemma-normal-integral-sections}). Ainsi $B \subset C$, et l'on conclut que $B$ est fini sur $A$, puisque $A$ est noethérien.

\medskip\noindent Il existe un ouvert non vide $V \subset Y$ tel que $f^{-1}V \to V$ soit fini ; voir Morphismes, définition \ref{morphisms-definition-alteration}. Après avoir rétréci $V$, nous pouvons supposer que $f^{-1}V \to V$ est plat (Morphismes, proposition \ref{morphisms-proposition-generic-flatness}). Ainsi $f^{-1}V \to V$ est fidèlement plat. Alors $V$ est régulier d'après Algèbre, lemme \ref{algebra-lemma-descent-regular}. \end{proof}

\begin{lemma} \label{resolve-lemma-algebra-helper} Soit $(A, \mathfrak m)$ un anneau local noethérien. Soient $B \subset C$ des $A$-algèbres finies. Supposons que (a) $B$ soit un anneau normal et que (b) le complété $\mathfrak m$-adique $C^\wedge$ soit un anneau normal. Alors $B^\wedge$ est un anneau normal. \end{lemma}

\begin{proof} Considérons le diagramme commutatif $$
\xymatrix{
B \ar[r] \ar[d] & C \ar[d] \\
B^\wedge \ar[r] & C^\wedge
}
$$. Rappelons que la complétion $\mathfrak m$-adique est exacte sur la catégorie des $A$-modules finis, car elle est donnée par le produit tensoriel avec la $A$-algèbre plate $A^\wedge$ (Algèbre, lemme \ref{algebra-lemma-completion-flat}). Nous allons utiliser le critère de Serre (Algèbre, lemme \ref{algebra-lemma-criterion-normal}) pour démontrer que l'anneau noethérien $B^\wedge$ est normal. Soit $\mathfrak q \subset B^\wedge$ un idéal premier situé au-dessus de $\mathfrak p \subset B$. Si $\dim(B_\mathfrak p) \geq 2$, alors $\text{depth}(B_\mathfrak p) \geq 2$ et, puisque $B_\mathfrak p \to B^\wedge_\mathfrak q$ est plat, on trouve que $\text{depth}(B^\wedge_\mathfrak q) \geq 2$ (Algèbre, lemme \ref{algebra-lemma-apply-grothendieck}). Si $\dim(B_\mathfrak p) \leq 1$, alors $B_\mathfrak p$ est soit un anneau de valuation discrète, soit un corps. Dans ce cas, $C_\mathfrak p$ est fidèlement plat sur $B_\mathfrak p$ (car il est fini et sans torsion). Par conséquent, $B^\wedge_\mathfrak p \to C^\wedge_\mathfrak p$ est fidèlement plat, et il en va de même après localisation en $\mathfrak q$. Comme $C^\wedge$, et donc toute localisation de cet anneau, est $(S_2)$, on conclut que $B^\wedge_\mathfrak p$ est $(S_2)$ d'après Algèbre, lemme \ref{algebra-lemma-descent-Sk}. Au total, on trouve que $(S_2)$ est vérifiée pour $B^\wedge$. Pour démontrer que $B^\wedge$ est $(R_1)$, il suffit de considérer les idéaux premiers $\mathfrak q \subset B^\wedge$ tels que $\dim(B^\wedge_\mathfrak q) \leq 1$. Puisque $\dim(B^\wedge_\mathfrak q) = \dim(B_\mathfrak p) +
\dim(B^\wedge_\mathfrak q/\mathfrak p B^\wedge_\mathfrak q)$ d'après Algèbre, lemme \ref{algebra-lemma-dimension-base-fibre-total}, on trouve que $\dim(B_\mathfrak p) \leq 1$ et l'on voit, comme précédemment, que $B^\wedge_\mathfrak q \to C^\wedge_\mathfrak q$ est fidèlement plat. On conclut à l'aide d'Algèbre, lemme \ref{algebra-lemma-descent-Rk}. \end{proof}

\begin{lemma} \label{resolve-lemma-regular-alteration-implies-local} Soit $(A, \mathfrak m, \kappa)$ un anneau local noethérien intègre. Supposons qu'il existe une altération $f : X \to \Spec(A)$ avec $X$ régulier. Alors : \begin{enumerate} \item il existe un élément non nul $f \in A$ tel que $A_f$ soit régulier ; \item la clôture intégrale $B$ de $A$ dans son corps des fractions est finie sur $A$ ; \item le complété $\mathfrak m$-adique de $B$ est un anneau normal, c'est-à-dire que les complétés de $B$ en ses idéaux maximaux sont des anneaux normaux intègres ; \item la fibre formelle générique de $A$ est régulière. \end{enumerate} \end{lemma}

\begin{proof} Les parties (1) et (2) résultent du lemme \ref{resolve-lemma-regular-alteration-implies}. Nous devons reprendre une partie de sa démonstration afin de fixer les notations nécessaires à celle de (3). Posons $C = \Gamma(X, \mathcal{O}_X)$. D'après Cohomologie des schémas, lemme \ref{coherent-lemma-proper-over-affine-cohomology-finite}, on voit que $C$ est un $A$-module fini. Comme $X$ est normal (Propriétés, lemme \ref{properties-lemma-regular-normal}), on voit que $C$ est un anneau normal intègre (Propriétés, lemme \ref{properties-lemma-normal-integral-sections}). Ainsi $B \subset C$, et l'on conclut que $B$ est fini sur $A$, puisque $A$ est noethérien. D'après le lemme \ref{resolve-lemma-algebra-helper}, pour prouver (3), il suffit de montrer que le complété $\mathfrak m$-adique $C^\wedge$ est normal.

\medskip\noindent D'après Algèbre, lemme \ref{algebra-lemma-completion-finite-extension}, le complété $C^\wedge$ est le produit des complétés de $C$ aux idéaux premiers de $C$ situés au-dessus de $\mathfrak m$. Ceux-ci sont en nombre fini et sont les idéaux maximaux $\mathfrak m_1, \ldots, \mathfrak m_r$ de $C$. (Le résultat correspondant pour $B$ explique la dernière assertion du lemme.) En remplaçant donc $A$ par $C_{\mathfrak m_i}$ et $X$ par $X_i = X \times_{\Spec(C)} \Spec(C_{\mathfrak m_i})$, on se ramène au cas traité au paragraphe suivant. (Remarquons que $\Gamma(X_i, \mathcal{O}) = C_{\mathfrak m_i}$ d'après Cohomologie des schémas, lemme \ref{coherent-lemma-flat-base-change-cohomology}.)

\medskip\noindent Dans ce paragraphe, $A$ est un anneau local noethérien normal intègre et $f : X \to \Spec(A)$ est une altération régulière avec $\Gamma(X, \mathcal{O}_X) = A$. Il faut montrer que le complété $A^\wedge$ de $A$ est un anneau normal intègre. D'après le lemme \ref{resolve-lemma-port-regularity-to-completion}, $Y = X \times_{\Spec(A)} \Spec(A^\wedge)$ est régulier. Puisque $\Gamma(Y, \mathcal{O}_Y) = A^\wedge$ d'après Cohomologie des schémas, lemme \ref{coherent-lemma-flat-base-change-cohomology}, on conclut comme précédemment que $A^\wedge$ est normal. En effet, $Y$ est normal d'après Propriétés, lemme \ref{properties-lemma-regular-normal}. Il est connexe puisque $\Gamma(Y, \mathcal{O}_Y) = A^\wedge$ est local. Ainsi $Y$ est normal et intègre (car un schéma noethérien connexe et normal est intègre). Par conséquent, $\Gamma(Y, \mathcal{O}_Y) = A^\wedge$ est un anneau normal intègre d'après Propriétés, lemme \ref{properties-lemma-normal-integral-sections}. Cela démontre (3).

\medskip\noindent Démonstration de (4). Notons $\eta \in \Spec(A)$ le point générique et indiquons par un indice $\eta$ le changement de base à $\eta$. Puisque $f$ est une altération, le schéma $X_\eta$ est fini et fidèlement plat sur $\eta$. Comme $Y = X \times_{\Spec(A)} \Spec(A^\wedge)$ est régulier d'après le lemme \ref{resolve-lemma-port-regularity-to-completion}, on voit que $Y_\eta$ est régulier (en tant que limite d'ouverts de $Y$). Alors $Y_\eta \to \Spec(A^\wedge \otimes_A \kappa(\eta))$ est fini et fidèlement plat sur la fibre formelle générique. On conclut par Algèbre, lemme \ref{algebra-lemma-descent-regular}. \end{proof}











\section{Résolution} \label{resolve-section-resolution}



\noindent Voici une définition.

\begin{definition} \label{resolve-definition-resolution} Soit $Y$ un schéma intègre noethérien. Une {\it résolution des singularités} de $Y$ est une modification $f : X \to Y$ telle que $X$ soit régulier. \end{definition}

\noindent Dans le cas des surfaces, nous souhaitons parfois disposer d'un peu plus d'information.

\begin{definition} \label{resolve-definition-resolution-surface} Soit $Y$ un schéma intègre noethérien de dimension $2$. Nous disons que $Y$ admet une {\it résolution des singularités par éclatements normalisés} s'il existe une suite $$
Y_n \to Y_{n - 1} \to \ldots \to Y_1 \to Y_0 \to Y
$$ telle que : \begin{enumerate} \item $Y_i$ soit propre sur $Y$ pour $i = 0, \ldots, n$ ; \item $Y_0 \to Y$ soit la normalisation ; \item $Y_i \to Y_{i - 1}$ soit un éclatement normalisé pour $i = 1, \ldots, n$ ; \item $Y_n$ soit régulier. \end{enumerate} \end{definition}

\noindent Remarquons que la condition (1) implique que la normalisation $Y_0$ de $Y$ est finie sur $Y$ et que les normalisations utilisées dans les éclatements normalisés sont elles aussi finies.

\begin{lemma} \label{resolve-lemma-existence-implies-existence-by-normalized-blowing-ups} Soit $(A, \mathfrak m, \kappa)$ un anneau local noethérien. Supposons que $A$ soit normal et de dimension $2$. Si $\Spec(A)$ admet une résolution des singularités, alors $\Spec(A)$ admet une résolution par éclatements normalisés. \end{lemma}

\begin{proof} D'après le lemme \ref{resolve-lemma-regular-alteration-implies-local}, le complété $A^\wedge$ de $A$ est normal. Le lemme \ref{resolve-lemma-port-regularity-to-completion} montre que $\Spec(A^\wedge)$ admet une résolution. D'après le lemme \ref{resolve-lemma-normalized-blowup-completion}, toute suite $Y_n \to Y_{n - 1} \to \ldots \to \Spec(A^\wedge)$ d'éclatements normalisés provient d'une suite d'éclatements normalisés $X_n \to \ldots \to \Spec(A)$. De plus, si $Y_n$ est régulier, alors $X_n$ est régulier d'après le lemme \ref{resolve-lemma-port-regularity-to-completion}. Il suffit donc de démontrer le lemme lorsque $A$ est complet.

\medskip\noindent Supposons en outre que $A$ soit complet. Nous utiliserons le fait que $A$ est de Nagata (Algèbre, proposition \ref{algebra-proposition-ubiquity-nagata}), excellent (Compléments d'algèbre, proposition \ref{more-algebra-proposition-ubiquity-excellent}) et possède un complexe dualisant (Complexes dualisants, lemme \ref{dualizing-lemma-ubiquity-dualizing}). Il en va de même pour tout anneau essentiellement de type fini sur $A$. Si $B$ est un anneau local normal intègre excellent, alors son complété $B^\wedge$ est normal (car $B \to B^\wedge$ est régulier et Compléments d'algèbre, lemme \ref{more-algebra-lemma-normal-goes-up}, s'applique). Nous l'utiliserons sans plus de commentaires dans la suite de la démonstration.

\medskip\noindent Soit $X \to \Spec(A)$ une résolution des singularités. Choisissons une suite d'éclatements normalisés $$
Y_n \to Y_{n - 1} \to \ldots \to Y_1 \to \Spec(A)
$$ qui domine $X$ (lemme \ref{resolve-lemma-dominate-by-normalized-blowing-up}). Le morphisme $Y_n \to X$ est un isomorphisme en dehors d'un nombre fini de points de $X$. Nous pouvons donc appliquer le lemme \ref{resolve-lemma-dominate-by-blowing-up-in-points} pour trouver une suite d'éclatements $$
X_m \to X_{m - 1} \to \ldots \to X
$$ en des points fermés telle que $X_m$ domine $Y_n$. On a le diagramme $$
\xymatrix{
& Y_n \ar[rd] \ar[rr] & & \Spec(A) \\
X_m \ar[rr] \ar[ru] & & X \ar[ru]
}
$$. Pour démontrer le lemme, il suffit de montrer qu'un nombre fini d'éclatements normalisés de $Y_n$ produit un schéma régulier. Le diagramme ci-dessus montre que $Y_n$ admet une résolution, à savoir $X_m$. Comme $Y_n$ est une surface normale, cela implique que $Y_n$ possède au plus un nombre fini de singularités $y_1, \ldots, y_t$ (car $X_m \to Y_n$ est un isomorphisme en dehors des fibres de dimension $1$ ; voir Variétés, lemme \ref{varieties-lemma-modification-normal-iso-over-codimension-1}).

\medskip\noindent Soit $x_a \in X$ l'image de $y_a$. Alors $\mathcal{O}_{X, x_a}$ est régulier et définit donc une singularité rationnelle (lemme \ref{resolve-lemma-regular-rational}). Appliquons le lemme \ref{resolve-lemma-rational-propagates} à $\mathcal{O}_{X, x_a} \to \mathcal{O}_{Y_n, y_a}$ : il en résulte que $\mathcal{O}_{Y_n, y_a}$ définit une singularité rationnelle. D'après le lemme \ref{resolve-lemma-rational-to-gorenstein}, il existe une suite finie d'éclatements en des points fermés singuliers $$
Y_{a, n_a} \to Y_{a, n_a - 1} \to \ldots \to \Spec(\mathcal{O}_{Y_n, y_a})
$$ telle que $Y_{a, n_a}$ soit de Gorenstein, c'est-à-dire possède un module dualisant inversible. D'après le lemme \ref{resolve-lemma-equivalence-sequence-blowups} (qui est essentiellement trivial), appliqué avec $n' = \sum n_a$, ces suites correspondent à une suite d'éclatements $$
Y_{n + n'} \to Y_{n + n' - 1} \to \ldots \to Y_n
$$ telle que $Y_{n + n'}$ soit normal et que les anneaux locaux de $Y_{n + n'}$ soient de Gorenstein. À l'aide des références données ci-dessus, nous pouvons dominer $Y_{n + n'}$ par une suite d'éclatements $X_{m + m'} \to \ldots \to X_m$ dominant $Y_{n + n'}$, comme dans le diagramme suivant $$
\xymatrix{
 & Y_{n + n'} \ar[rr] & & Y_n \ar[rd] \ar[rr] & & \Spec(A) \\
X_{m + m'} \ar[ru] \ar[rr] & & X_m \ar[rr] \ar[ru] & & X \ar[ru]
}
$$. Ainsi, $Y_{n + n'}$ possède de nouveau un nombre fini de points singuliers $y'_1, \ldots, y'_s$, mais cette fois les singularités sont des points doubles rationnels ; plus précisément, les anneaux locaux $\mathcal{O}_{Y_{n + n'}, y'_b}$ sont comme dans le lemme \ref{resolve-lemma-resolve-rational-double-points}. En raisonnant exactement comme ci-dessus, on conclut que le lemme est vrai. \end{proof}

\begin{lemma} \label{resolve-lemma-resolve-complete} Soit $(A, \mathfrak m, \kappa)$ un anneau local noethérien complet. Supposons que $A$ soit un anneau normal intègre de dimension $2$. Alors $\Spec(A)$ admet une résolution des singularités. \end{lemma}

\begin{proof} Un anneau local noethérien complet est J-2 (Compléments d'algèbre, proposition \ref{more-algebra-proposition-ubiquity-J-2}), de Nagata (Algèbre, proposition \ref{algebra-proposition-ubiquity-nagata}), excellent (Compléments d'algèbre, proposition \ref{more-algebra-proposition-ubiquity-excellent}) et possède un complexe dualisant (Complexes dualisants, lemme \ref{dualizing-lemma-ubiquity-dualizing}). Il en va de même pour tout anneau essentiellement de type fini sur $A$. Si $B$ est un anneau local normal intègre excellent, alors son complété $B^\wedge$ est normal (car $B \to B^\wedge$ est régulier et Compléments d'algèbre, lemme \ref{more-algebra-lemma-normal-goes-up}, s'applique). Autrement dit, les anneaux locaux rencontrés dans la suite de la démonstration possèdent les propriétés d'excellence dont nous avons besoin.

\medskip\noindent Choisissons $A_0 \subset A$ avec $A_0$ un anneau local régulier complet et $A_0 \to A$ fini ; voir Algèbre, lemme \ref{algebra-lemma-complete-local-Noetherian-domain-finite-over-regular}. Cela induit une extension finie de corps des fractions $K/K_0$. Nous allons raisonner par récurrence sur $[K : K_0]$. Le cas initial est celui où le degré vaut $1$ ; alors $A_0 = A$ et le résultat est vrai.

\medskip\noindent Supposons qu'il existe un corps intermédiaire $K_0 \subset L \subset K$, $K_0 \not = L \not = K$. Soit $B \subset A$ la clôture intégrale de $A_0$ dans $L$. Par récurrence, choisissons une résolution des singularités $Y \to \Spec(B)$. Soit $X$ la normalisation de $Y \times_{\Spec(B)} \Spec(A)$. On a le diagramme $$
\xymatrix{
X \ar[r] \ar[d] & \Spec(A) \ar[d] \\
Y \ar[r] & \Spec(B)
}
$$. Puisque $A$ est J-2, le lieu régulier de $X$ est ouvert. Comme $X$ est une surface normale, on conclut que $X$ possède au plus un nombre fini de points singuliers $x_1, \ldots, x_n$ ; ce sont des points fermés tels que $\dim(\mathcal{O}_{X, x_i}) = 2$. Pour tout $i$, soit $y_i \in Y$ son image. Comme $\mathcal{O}_{Y, y_i}^\wedge \to \mathcal{O}_{X, x_i}^\wedge$ est fini, de degré inférieur au précédent, l'hypothèse de récurrence montre que $\mathcal{O}_{X, x_i}^\wedge$ admet une résolution des singularités. D'après le lemme \ref{resolve-lemma-existence-implies-existence-by-normalized-blowing-ups}, il existe une suite $$
Z^\wedge_{i, n_i} \to \ldots \to Z^\wedge_{i, 1} \to
\Spec(\mathcal{O}_{X, x_i}^\wedge)
$$ d'éclatements normalisés telle que $Z^\wedge_{i, n_i}$ soit régulier. Le lemme \ref{resolve-lemma-normalized-blowup-completion} fournit une suite correspondante d'éclatements normalisés $$
Z_{i, n_i} \to \ldots \to Z_{i, 1} \to \Spec(\mathcal{O}_{X, x_i})
$$. Alors $Z_{i, n_i}$ est un schéma régulier d'après le lemme \ref{resolve-lemma-port-regularity-to-completion}. Le lemme \ref{resolve-lemma-equivalence-sequence-normalized-blowups} permet d'insérer ces éclatements normalisés dans une suite correspondante $$
Z_n \to Z_{n - 1} \to \ldots \to Z_1 \to X
$$, et bien sûr $Z_n$ est lui aussi régulier (il suffit de considérer les anneaux locaux). Cela démontre l'étape de récurrence.

\medskip\noindent Supposons qu'il n'existe aucun corps intermédiaire $K_0 \subset L \subset K$ avec $K_0 \not = L \not = K$. Alors, soit $K/K_0$ est séparable, soit la caractéristique de $K$ est $p$ et $[K : K_0] = p$. Le lemme \ref{resolve-lemma-go-up-separable} ou le lemme \ref{resolve-lemma-go-up-degree-p} implique alors que la réduction aux singularités rationnelles est possible. D'après le lemme \ref{resolve-lemma-reduce-to-rational}, il existe donc une modification normale $X \to \Spec(A)$ telle que, pour tout point singulier $x$ de $X$, l'anneau local $\mathcal{O}_{X, x}$ définisse une singularité rationnelle. Puisque $A$ est J-2, on trouve que $X$ ne possède qu'un nombre fini de points singuliers $x_1, \ldots, x_n$. D'après le lemme \ref{resolve-lemma-rational-to-gorenstein}, il existe une suite finie d'éclatements en des points fermés singuliers $$
X_{i, n_i} \to X_{i, n_i - 1} \to \ldots \to \Spec(\mathcal{O}_{X, x_i})
$$ telle que $X_{i, n_i}$ soit de Gorenstein, c'est-à-dire possède un module dualisant inversible. D'après le lemme \ref{resolve-lemma-equivalence-sequence-blowups} (qui est essentiellement trivial), appliqué avec $n = \sum n_a$, ces suites correspondent à une suite d'éclatements $$
X_n \to X_{n - 1} \to \ldots \to X
$$ telle que $X_n$ soit normal et que les anneaux locaux de $X_n$ soient de Gorenstein. De nouveau, $X_n$ possède un nombre fini de points singuliers $x'_1, \ldots, x'_s$, mais cette fois les singularités sont des points doubles rationnels ; plus précisément, les anneaux locaux $\mathcal{O}_{X_n, x'_i}$ sont comme dans le lemme \ref{resolve-lemma-resolve-rational-double-points}. En raisonnant exactement comme ci-dessus, on conclut que le lemme est vrai. \end{proof}

\noindent Nous arrivons enfin au théorème principal de ce chapitre.

\begin{theorem}[Lipman]
\label{resolve-theorem-resolve}
\begin{reference}
\cite[Théorème, page 151]{Lipman}
\end{reference}
Soit $Y$ un schéma intègre noethérien de dimension deux. Les assertions suivantes sont équivalentes :
\begin{enumerate}
\item il existe une altération $X \to Y$ avec $X$ régulier ;
\item il existe une résolution des singularités de $Y$ ;
\item $Y$ admet une résolution des singularités par éclatements normalisés ;
\item la normalisation $Y^\nu \to Y$ est finie, $Y^\nu$ ne possède qu'un nombre fini de points singuliers $y_1, \ldots, y_m$ et, pour tout $y_i$, le complété de $\mathcal{O}_{Y^\nu, y_i}$ est normal.
\end{enumerate}
\end{theorem}

\begin{proof} Les implications (3) $\Rightarrow$ (2) $\Rightarrow$ (1) sont immédiates.

\medskip\noindent Soit $X \to Y$ une altération avec $X$ régulier. Alors $Y^\nu \to Y$ est fini d'après le lemme \ref{resolve-lemma-regular-alteration-implies}. Considérons la factorisation $f : X \to Y^\nu$ de Morphismes, lemme \ref{morphisms-lemma-normalization-normal}. Le morphisme $f$ est fini au-dessus d'un ouvert $V \subset Y^\nu$ contenant tout point de codimension $\leq 1$ de $Y^\nu$, d'après Variétés, lemme \ref{varieties-lemma-finite-in-codim-1}. Alors $f$ est plat sur $V$ d'après Algèbre, lemme \ref{algebra-lemma-CM-over-regular-flat}, et parce qu'un anneau local normal de dimension $\leq 2$ est de Cohen--Macaulay par le critère de Serre (Algèbre, lemme \ref{algebra-lemma-criterion-normal}). Alors $V$ est régulier d'après Algèbre, lemme \ref{algebra-lemma-descent-regular}. Comme $Y^\nu$ est noethérien, on conclut que $Y^\nu \setminus V = \{y_1, \ldots, y_m\}$ est fini. D'après le lemme \ref{resolve-lemma-regular-alteration-implies-local}, le complété de $\mathcal{O}_{Y^\nu, y_i}$ est normal. On voit ainsi que (1) $\Rightarrow$ (4).

\medskip\noindent Supposons (4). Il faut démontrer (3). Nous pouvons immédiatement remplacer $Y$ par sa normalisation. Soient $y_1, \ldots, y_m \in Y$ les points singuliers. En appliquant les lemmes \ref{resolve-lemma-resolve-complete} et \ref{resolve-lemma-existence-implies-existence-by-normalized-blowing-ups}, on trouve qu'il existe une suite finie d'éclatements normalisés $$
Y_{i, n_i} \to Y_{i, n_i - 1} \to \ldots \to \Spec(\mathcal{O}^\wedge_{Y, y_i})
$$ telle que $Y_{i, n_i}$ soit régulier. D'après le lemme \ref{resolve-lemma-normalized-blowup-completion}, il existe une suite correspondante d'éclatements normalisés $$
X_{i, n_i} \to \ldots \to X_{i, 1} \to \Spec(\mathcal{O}_{Y, y_i})
$$. Alors $X_{i, n_i}$ est un schéma régulier d'après le lemme \ref{resolve-lemma-port-regularity-to-completion}. Le lemme \ref{resolve-lemma-equivalence-sequence-normalized-blowups} permet d'insérer ces éclatements normalisés dans une suite correspondante $$
X_n \to X_{n - 1} \to \ldots \to X_1 \to Y
$$, et bien sûr $X_n$ est lui aussi régulier (il suffit de considérer les anneaux locaux). Cela achève la démonstration. \end{proof}






\section{Résolution plongée} \label{resolve-section-embedded}

\noindent Étant donnée une courbe sur une surface, il existe un éclatement qui transforme la courbe en un diviseur à croisements normaux stricts. Dans cette section, nous utiliserons le fait qu'un schéma localement noethérien de dimension un est normal si et seulement s'il est régulier (Algèbre, lemme \ref{algebra-lemma-characterize-dvr}). Nous utiliserons également le fait que tout point d'un schéma localement noethérien se spécialise en un point fermé (Propriétés, lemme \ref{properties-lemma-locally-Noetherian-closed-point}).

\begin{lemma} \label{resolve-lemma-resolve-curve} Soit $Y$ un schéma intègre noethérien de dimension un. Les assertions suivantes sont équivalentes : \begin{enumerate} \item il existe une altération $X \to Y$ avec $X$ régulier ; \item il existe une résolution des singularités de $Y$ ; \item il existe une suite finie $Y_n \to Y_{n - 1} \to \ldots \to Y_1 \to Y$ d'éclatements en des points fermés telle que $Y_n$ soit régulier ; \item la normalisation $Y^\nu \to Y$ est finie. \end{enumerate} \end{lemma}

\begin{proof} Les implications (3) $\Rightarrow$ (2) $\Rightarrow$ (1) sont immédiates. L'implication (1) $\Rightarrow$ (4) résulte du lemme \ref{resolve-lemma-regular-alteration-implies}. Remarquons qu'un schéma normal de dimension un est régulier ; l'implication (4) $\Rightarrow$ (2) est donc elle aussi claire. Il reste ainsi à montrer que les conditions équivalentes (1), (2) et (4) impliquent (3).

\medskip\noindent Soit $f : X \to Y$ une résolution des singularités. Comme $Y$ est de dimension un, on voit que $f$ est fini d'après Variétés, lemme \ref{varieties-lemma-finite-in-codim-1}. Nous allons construire des factorisations $$
X \to \ldots \to Y_2 \to Y_1 \to Y
$$ dans lesquelles $Y_i \to Y_{i - 1}$ est l'éclatement d'un point fermé et n'est pas un isomorphisme tant que $Y_{i - 1}$ n'est pas régulier. Chacun de ces morphismes sera fini (pour la même raison que ci-dessus), et l'on obtiendra un système correspondant $$
f_*\mathcal{O}_X \supset \ldots \supset
f_{2, *}\mathcal{O}_{Y_2} \supset
f_{1, *}\mathcal{O}_{Y_1} \supset \mathcal{O}_Y
$$, où $f_i : Y_i \to Y$ est le morphisme structural. Puisque $Y$ est noethérien, cette suite croissante de sous-modules cohérents doit se stabiliser (Cohomologie des schémas, lemme \ref{coherent-lemma-acc-coherent}), ce qui montre que, pour un certain $n$, le schéma $Y_n$ est régulier, comme souhaité. Pour construire $Y_i$ à partir de $Y_{i - 1}$, choisissons un point fermé singulier $y_{i - 1} \in Y_{i - 1}$ et soit $Y_i \to Y_{i - 1}$ l'éclatement correspondant. Comme $X$ est régulier de dimension $1$ (ses anneaux locaux aux points fermés sont donc des anneaux de valuation discrète, en particulier des anneaux principaux), le faisceau d'idéaux $\mathfrak m_{y_{i - 1}} \cdot \mathcal{O}_X$ est inversible. La propriété universelle de l'éclatement (Diviseurs, lemme \ref{divisors-lemma-universal-property-blowing-up}) fournit alors une factorisation $X \to Y_i$. Enfin, $Y_i \to Y_{i - 1}$ n'est pas un isomorphisme, puisque $\mathfrak m_{y_{i - 1}}$ n'est pas un idéal inversible. \end{proof}

\begin{lemma} \label{resolve-lemma-blowup-curve} Soit $X$ un schéma noethérien. Soit $Y \subset X$ un sous-schéma fermé intègre de dimension $1$ qui satisfait les conditions équivalentes du lemme \ref{resolve-lemma-resolve-curve}. Il existe alors une suite finie $$
X_n \to X_{n - 1} \to \ldots \to X_1 \to X
$$ d'éclatements en des points fermés telle que le transformé strict de $Y$ dans $X_n$ soit une courbe régulière. \end{lemma}

\begin{proof} Soit $Y_n \to Y_{n - 1} \to \ldots \to Y_1 \to Y$ la suite d'éclatements fournie par le lemme \ref{resolve-lemma-resolve-curve}. Soit $X_n \to X_{n - 1} \to \ldots \to X_1 \to X$ la suite correspondante d'éclatements de $X$. Cela convient, car le transformé strict est l'éclatement d'après Diviseurs, lemme \ref{divisors-lemma-strict-transform}. \end{proof}

\noindent Soit $X$ un schéma localement noethérien. Soient $Y, Z \subset X$ des sous-schémas fermés. Soit $p \in Y \cap Z$ un point fermé. Supposons que $Y$ soit intègre de dimension $1$ et que le point générique de $Y$ ne soit pas contenu dans $Z$. Dans cette situation, on peut considérer l'invariant \begin{equation}
\label{resolve-equation-multiplicity}
m_p(Y \cap Z) =
\text{length}_{\mathcal{O}_{X, p}}(\mathcal{O}_{Y \cap Z, p})
\end{equation}. C'est un entier $\geq 1$. En effet, si $I, J \subset \mathcal{O}_{X, p}$ sont les idéaux correspondant à $Y, Z$, on voit que $\mathcal{O}_{Y \cap Z, p} = \mathcal{O}_{X, p}/I + J$ a pour support $\{\mathfrak m_p\}$, car nous avons supposé que $Y \cap Z$ ne contient pas l'unique point de $Y$ qui se spécialise en $p$. La longueur est donc finie d'après Algèbre, lemme \ref{algebra-lemma-support-point}.

\begin{lemma} \label{resolve-lemma-blowup-nonsingular-curves-meeting-at-point} Dans la situation ci-dessus, soit $X' \to X$ l'éclatement de $X$ en $p$. Soient $Y', Z' \subset X'$ les transformés stricts de $Y, Z$. Si $\mathcal{O}_{Y, p}$ est régulier, alors : \begin{enumerate} \item $Y' \to Y$ est un isomorphisme ; \item $Y'$ rencontre la fibre exceptionnelle $E \subset X'$ en un point $q$ et $m_q(Y \cap E) = 1$ ; \item si $q \in Z'$ l'est aussi, alors $m_q(Y \cap Z') < m_p(Y \cap Z)$. \end{enumerate} \end{lemma}

\begin{proof} Puisque $\mathcal{O}_{X, p} \to \mathcal{O}_{Y, p}$ est surjectif et que $\mathcal{O}_{Y, p}$ est un anneau de valuation discrète, nous pouvons choisir un élément $x_1 \in \mathfrak m_p$ qui s'envoie sur une uniformisante dans $\mathcal{O}_{Y, p}$. Choisissons un ouvert affine $U = \Spec(A)$ contenant $p$ tel que $x_1 \in A$. Soit $\mathfrak m \subset A$ l'idéal maximal correspondant à $p$. Soient $I, J \subset A$ les idéaux qui définissent $Y, Z$ dans $\Spec(A)$. Après avoir rétréci $U$, nous pouvons supposer que $\mathfrak m = I + (x_1)$, autrement dit que $V(x_1) \cap U \cap Y = \{p\}$ schématiquement. On conclut que $p$ est un diviseur de Cartier effectif sur $Y$ et, puisque $Y'$ est l'éclatement de $Y$ en $p$ (Diviseurs, lemme \ref{divisors-lemma-strict-transform}), on voit que $Y' \to Y$ est un isomorphisme d'après Diviseurs, lemme \ref{divisors-lemma-blow-up-effective-Cartier-divisor}. La relation $\mathfrak m = I + (x_1)$ implique que $\mathfrak m^n \subset I + (x_1^n)$ ; on peut donc définir une application $$
\psi : A[\textstyle{\frac{\mathfrak m}{x_1}}] \longrightarrow A/I
$$ en envoyant $y/x_1^n \in A[\frac{\mathfrak m}{x_1}]$ sur la classe de $a$ dans $A/I$, où $a$ est choisi de sorte que $y \equiv ax_1^n \bmod I$. Alors $\psi$ correspond au morphisme de $Y \cap U$ dans $X'$ au-dessus de $U$ donné par $Y' \cong Y$. Puisque l'image de $x_1$ dans $A[\frac{\mathfrak m}{x_1}]$ définit le diviseur exceptionnel, on conclut que $m_q(Y', E) = 1$. Enfin, comme $J \subset \mathfrak m$, l'idéal $J' \subset A[\frac{\mathfrak m}{x_1}]$ contient certainement les éléments $f/x_1$ pour $f \in J$. Ainsi, si l'on choisit $f \in J$ dont l'image $\overline{f}$ dans $A/I$ a pour valuation minimale $m_p(Y \cap Z)$, on voit que $\psi(f/x_1) = \overline{f}/x_1$ dans $A/I$ a une valuation inférieure d'une unité, ce qui démontre la dernière partie du lemme. \end{proof}

\begin{lemma} \label{resolve-lemma-blowup-curves} Soit $X$ un schéma noethérien. Soient $Y_i \subset X$, $i = 1, \ldots, n$ des sous-schémas fermés intègres, chacun de dimension $1$ et satisfaisant les conditions équivalentes du lemme \ref{resolve-lemma-resolve-curve}. Il existe alors une suite finie $$
X_n \to X_{n - 1} \to \ldots \to X_1 \to X
$$ d'éclatements en des points fermés telle que les transformés stricts $Y'_i \subset X_n$ de $Y_i$ dans $X_n$ soient des courbes régulières deux à deux disjointes. \end{lemma}

\begin{proof} Il résulte du lemme \ref{resolve-lemma-blowup-curve} que nous pouvons supposer que $Y_i$ est une courbe régulière pour $i = 1, \ldots, n$. Pour tous $i \not = j$ et $p \in Y_i \cap Y_j$, nous disposons de l'invariant $m_p(Y_i \cap Y_j)$ (\ref{resolve-equation-multiplicity}). Si le maximum de ces nombres est $> 1$, on peut le diminuer (lemme \ref{resolve-lemma-blowup-nonsingular-curves-meeting-at-point}) en éclatant tous les points $p$ où ce maximum est atteint. Si le maximum vaut $1$, on peut alors séparer les courbes au moyen du même lemme, en éclatant tous ces points $p$. \end{proof}

\noindent Lorsque notre courbe est contenue dans une surface régulière, nous souhaitons souvent la transformer en un diviseur à croisements normaux.

\begin{lemma} \label{resolve-lemma-turn-into-effective-Cartier} Soit $X$ un schéma régulier de dimension $2$. Soit $Z \subset X$ un sous-schéma fermé propre. Il existe une suite $$
X_n \to \ldots \to X_1 \to X
$$ d'éclatements en des points fermés telle que l'image inverse $Z_n$ de $Z$ dans $X_n$ soit un diviseur de Cartier effectif. \end{lemma}

\begin{proof} Soit $D \subset Z$ le plus grand diviseur de Cartier effectif contenu dans $Z$. Alors $\mathcal{I}_Z \subset \mathcal{I}_D$, et le quotient est supporté en des points fermés d'après Diviseurs, lemme \ref{divisors-lemma-codim-1-part}. Nous pouvons donc écrire $\mathcal{I}_Z = \mathcal{I}_{Z'} \mathcal{I}_D$, où $Z' \subset X$ est un sous-schéma fermé qui, ensemblistement, consiste en un nombre fini de points fermés. En appliquant le lemme \ref{resolve-lemma-make-ideal-principal}, on trouve une suite d'éclatements comme dans l'énoncé du lemme telle que $\mathcal{I}_{Z'}\mathcal{O}_{X_n}$ soit inversible. Cela démontre le lemme. \end{proof}

\begin{lemma} \label{resolve-lemma-embedded-resolution} Soit $X$ un schéma régulier de dimension $2$. Soit $Z \subset X$ un sous-schéma fermé propre tel que toute composante irréductible $Y \subset Z$ de dimension $1$ satisfasse les conditions équivalentes du lemme \ref{resolve-lemma-resolve-curve}. Il existe alors une suite $$
X_n \to \ldots \to X_1 \to X
$$ d'éclatements en des points fermés telle que l'image inverse $Z_n$ de $Z$ dans $X_n$ soit un diviseur de Cartier effectif supporté sur un diviseur à croisements normaux stricts. \end{lemma}

\begin{proof} Soit $X' \to X$ l'éclatement d'un point fermé $p$. Alors l'image inverse $Z' \subset X'$ de $Z$ est supportée sur le transformé strict de $Z$ et sur le diviseur exceptionnel. Le diviseur exceptionnel est une courbe régulière (lemme \ref{resolve-lemma-blowup}), et le transformé strict $Y'$ de toute composante irréductible $Y$ est soit égal à $Y$, soit l'éclatement de $Y$ en $p$. Ainsi, ce procédé ne produit pas de nouvelles composantes singulières de dimension $1$. Il résulte donc des lemmes \ref{resolve-lemma-turn-into-effective-Cartier} et \ref{resolve-lemma-blowup-curves} que nous pouvons supposer que $Z$ est un diviseur de Cartier effectif et que toutes les composantes irréductibles $Y$ de $Z$ sont régulières. (Bien sûr, nous ne pouvons pas supposer que les composantes irréductibles sont deux à deux disjointes, car chaque fois que nous éclatons un point de $Z$, nous ajoutons une nouvelle composante irréductible à $Z$, à savoir le diviseur exceptionnel.)

\medskip\noindent Supposons que $Z$ soit un diviseur de Cartier effectif dont les composantes irréductibles $Y_i$ sont régulières. Pour tous $i \not = j$ et $p \in Y_i \cap Y_j$, nous disposons de l'invariant $m_p(Y_i \cap Y_j)$ (\ref{resolve-equation-multiplicity}). Si le maximum de ces nombres est $> 1$, nous pouvons le diminuer (lemme \ref{resolve-lemma-blowup-nonsingular-curves-meeting-at-point}) en éclatant tous les points $p$ où ce maximum est atteint (remarquons que les « nouveaux » invariants $m_{q_i}(Y'_i \cap E)$ valent toujours $1$). Si le maximum vaut $1$ et si, par exemple, $p \in Y_1 \cap \ldots \cap Y_r$ pour un certain $r > 2$ mais pour aucun des autres, alors, après avoir éclaté $p$, on voit que $Y'_1, \ldots, Y'_r$ ne se rencontrent pas aux points situés au-dessus de $p$, et que $m_{q_i}(Y'_i, E) = 1$, où $Y'_i \cap E = \{q_i\}$. En continuant donc à éclater les points où plus de $3$ composantes de $Z$ se rencontrent, on atteint une situation dans laquelle, pour tout point fermé $p \in X$, on a l'un des trois cas suivants : (a) aucune courbe $Y_i$ ne passe par $p$ ; (b) exactement une courbe $Y_i$ passe par $p$ et $\mathcal{O}_{Y_i, p}$ est régulier ; (c) exactement deux courbes $Y_i$, $Y_j$ passent par $p$, les anneaux locaux $\mathcal{O}_{Y_i, p}$, $\mathcal{O}_{Y_j, p}$ sont réguliers et $m_p(Y_i \cap Y_j) = 1$. Cela signifie que $\sum Y_i$ est un diviseur à croisements normaux stricts sur la surface régulière $X$ ; voir Morphismes \'etales, lemme \ref{etale-lemma-strict-normal-crossings}. \end{proof}






\section{Contraction des courbes exceptionnelles} \label{resolve-section-minus-one}

\noindent Soit $X$ un schéma noethérien. Soit $E \subset X$ un sous-schéma fermé possédant les propriétés suivantes : \begin{enumerate} \item $E$ est un diviseur de Cartier effectif sur $X$ ; \item il existe un corps $k$ et un isomorphisme de schémas $\mathbf{P}^1_k \to E$ ; \item l'image inverse du faisceau normal $\mathcal{N}_{E/X}$ par cet isomorphisme est isomorphe à $\mathcal{O}_{\mathbf{P}^1}(-1)$. \end{enumerate} Un tel sous-schéma fermé est appelé une {\it courbe exceptionnelle de première espèce}.

\medskip\noindent Soit $X'$ un schéma noethérien et soit $x \in X'$ un point fermé tel que $\mathcal{O}_{X', x}$ soit régulier de dimension $2$. Soit $b : X \to X'$ l'éclatement de $X'$ en $x$. Dans ce cas, la fibre exceptionnelle $E \subset X$ est une courbe exceptionnelle de première espèce. Cela résulte du lemme \ref{resolve-lemma-blowup}.

\medskip\noindent Question : toute courbe exceptionnelle de première espèce est-elle obtenue comme fibre d'un éclatement comme ci-dessus ? Autrement dit, existe-t-il toujours un morphisme propre de schémas $X \to X'$ tel que $E$ soit envoyée sur un point fermé $x \in X'$, que $\mathcal{O}_{X', x}$ soit régulier de dimension $2$ et que $X$ soit l'éclatement de $X'$ en $x$ ? Si tel est le cas, nous disons qu'{\it il existe une contraction de $E$}.

\begin{lemma} \label{resolve-lemma-factor-through-contraction} Soit $X$ un schéma noethérien. Soit $E \subset X$ une courbe exceptionnelle de première espèce. S'il existe une contraction $X \to X'$ de $E$, celle-ci possède la propriété universelle suivante : pour tout morphisme $\varphi : X \to Y$ tel que $\varphi(E)$ soit un point, il existe une unique factorisation $X \to X' \to Y$ de $\varphi$. \end{lemma}

\begin{proof} Soit $b : X \to X'$ une contraction de $E$. En tant qu'espace topologique, $X'$ est le quotient de $X$ par la relation qui identifie tous les points de $E$ à un seul point. En effet, $b$ est propre (Diviseurs, lemme \ref{divisors-lemma-blowing-up-projective} et Morphismes, lemme \ref{morphisms-lemma-locally-projective-proper}) et surjectif ; il définit donc une application submersive d'espaces topologiques (Topologie, lemme \ref{topology-lemma-closed-morphism-quotient-topology}). D'autre part, l'application canonique $\mathcal{O}_{X'} \to b_*\mathcal{O}_X$ est un isomorphisme. Cela est clair sur le complémentaire du point $x \in X'$, image de $E$ ; au niveau des germes en $x$, l'application est un isomorphisme d'après la partie (4) du lemme \ref{resolve-lemma-cohomology-of-blowup}. Ainsi, le couple $(X', \mathcal{O}_{X'})$ s'obtient à partir de $X$ en prenant le quotient comme espace topologique et en le munissant de $b_*\mathcal{O}_X$ comme faisceau structural.

\medskip\noindent Étant donné $\varphi$, soit $\varphi' : X' \to Y$ l'unique application d'espaces topologiques telle que $\varphi = \varphi' \circ b$. L'application $$
\varphi^\sharp : \varphi^{-1}\mathcal{O}_Y =
b^{-1}((\varphi')^{-1}\mathcal{O}_Y) \to \mathcal{O}_X
$$ est alors adjointe à une application $$
(\varphi')^\sharp :
(\varphi')^{-1}\mathcal{O}_Y \to b_*\mathcal{O}_X = \mathcal{O}_{X'}
$$. Le couple $(\varphi', (\varphi')^\sharp)$ est un morphisme d'espaces annelés de $X'$ vers $Y$ qui fournit la factorisation voulue. Comme $\varphi$ est un morphisme d'espaces localement annelés, il en est de même de $\varphi'$. En effet, la seule chose à vérifier est que l'homomorphisme $\mathcal{O}_{Y, y} \to \mathcal{O}_{X', x}$ est local, où $y \in Y$ est l'image de $E$ par $\varphi$. Cela est vrai, car l'image inverse d'un élément $f \in \mathfrak m_y$ est une fonction sur $X$ qui s'annule en tout point de $E$ ; par conséquent, l'image inverse de $f$ sur $X'$ est une fonction définie sur un voisinage de $x$ dans $X'$ et possédant la même propriété. Il est alors clair que cette fonction doit s'annuler en $x$, comme souhaité. \end{proof}

\begin{lemma} \label{resolve-lemma-contraction-unique} Soit $X$ un schéma noethérien. Soit $E \subset X$ une courbe exceptionnelle de première espèce. S'il existe une contraction de $E$, alors elle est unique à isomorphisme unique près. \end{lemma}

\begin{proof} Cela résulte immédiatement de la propriété universelle du lemme \ref{resolve-lemma-factor-through-contraction}. \end{proof}

\begin{lemma} \label{resolve-lemma-exceptional-first-kind-local} Soit $X$ un schéma noethérien. Soit $E \subset X$ une courbe exceptionnelle de première espèce. Posons $E_n = nE$ et notons $\mathcal{O}_n$ son faisceau structural. Alors $$
A = \lim H^0(E_n, \mathcal{O}_n)
$$ est un anneau local noethérien complet et régulier de dimension $2$, et $\Ker(A \to H^0(E_n, \mathcal{O}_n))$ est la $n$-ième puissance de son idéal maximal. \end{lemma}

\begin{proof} Rappelons qu'il existe un isomorphisme $\mathbf{P}^1_k \to E$ tel que l'image inverse du faisceau normal de $E$ dans $X$ soit isomorphe à $\mathcal{O}(-1)$. Alors $H^0(E, \mathcal{O}_E) = k$. Nous noterons $\mathcal{O}_n(iE)$ la restriction du faisceau inversible $\mathcal{O}_X(iE)$ à $E_n$, pour tous $n \geq 1$ et $i \in \mathbf{Z}$. Rappelons que $\mathcal{O}_X(-nE)$ est le faisceau d'idéaux de $E_n$. Ainsi, pour $d \geq 0$, nous obtenons une suite exacte courte $$
0 \to \mathcal{O}_E(-(d + n)E) \to
\mathcal{O}_{n + 1}(-dE) \to
\mathcal{O}_n(-dE) \to 0
$$. Comme $\mathcal{O}_E(-(d + n)E) = \mathcal{O}_{\mathbf{P}^1_k}(d + n)$, le premier groupe de cohomologie s'annule pour tous $d \geq 0$ et $n \geq 1$. On en déduit que les applications de transition du système $H^0(E_n, \mathcal{O}_n(-dE))$ sont surjectives. Pour $d = 0$, nous obtenons un système projectif de surjections d'anneaux tel que le noyau de chaque application de transition soit un idéal nilpotent. Ainsi, $A = \lim H^0(E_n, \mathcal{O}_n)$ est un anneau local de corps résiduel $k$ et d'idéal maximal $$
\lim \Ker(H^0(E_n, \mathcal{O}_n) \to H^0(E, \mathcal{O}_E)) =
\lim H^0(E_n, \mathcal{O}_n(-E))
$$.
\begingroup\emergencystretch=3em
Choisissons $x, y$ dans ce noyau, dont les images forment une $k$-base de $H^0(E, \mathcal{O}_E(-E)) = H^0(\mathbf{P}^1_k, \mathcal{O}(1))$. Alors $x^d, x^{d - 1}y, \ldots, y^d$ sont des éléments de $\lim H^0(E_n, \mathcal{O}_n(-dE))$ dont les images forment une base de $H^0(E, \mathcal{O}_E(-dE)) = H^0(\mathbf{P}^1_k, \mathcal{O}(d))$. On voit ainsi que $A$ est séparé et complet pour la topologie linéaire définie par les noyaux\par\endgroup
$$
I_n = \Ker(A \longrightarrow H^0(E_n, \mathcal{O}_n))
$$. Nous avons $x, y \in I_1$, $I_d I_{d'} \subset I_{d + d'}$, et $I_d/I_{d + 1}$ est un $k$-module libre de base $x^d, x^{d - 1}y, \ldots, y^d$. Nous allons montrer que $I_d = (x, y)^d$. En effet, si $z_e \in I_e$ avec $e \geq d$, alors nous pouvons écrire $$
z_e = a_{e, 0} x^d + a_{e, 1} x^{d - 1}y + \ldots + a_{e, d}y^d + z_{e + 1}
$$, où $a_{e, j} \in (x, y)^{e - d}$ et $z_{e + 1} \in I_{e + 1}$ d'après notre description de $I_d/I_{d + 1}$. Ainsi, en partant d'un certain $z = z_d \in I_d$, nous pouvons procéder par récurrence et obtenir $$
z = \sum\nolimits_{e \geq d} \sum\nolimits_j a_{e, j} x^{d - j} y^j
$$ avec des $a_{e, j} \in (x, y)^{e - d}$. Alors $a_j = \sum_{e \geq d} a_{e, j}$ existe (par complétude et parce que $a_{e, j} \in I_{e - d}$), et nous avons $z = \sum a_{e, j} x^{d - j} y^j$. Par conséquent, $I_d = (x, y)^d$. Ainsi, $A$ est complet pour la topologie $(x, y)$-adique. Alors $A$ est noethérien d'après Algèbre, lemme \ref{algebra-lemma-completion-Noetherian}. Il est clair que sa dimension vaut $2$ d'après la description de $(x, y)^d/(x, y)^{d + 1}$ et Algèbre, proposition \ref{algebra-proposition-dimension}. Comme son idéal maximal est engendré par deux éléments, il est régulier. \end{proof}

\begin{lemma} \label{resolve-lemma-contraction} Soit $X$ un schéma noethérien. Soit $E \subset X$ une courbe exceptionnelle de première espèce. S'il existe un morphisme $f : X \to Y$ tel que \begin{enumerate} \item $Y$ est noethérien, \item $f$ est propre, \item $f$ envoie $E$ sur un point $y$ de $Y$, \item $f$ est quasi-fini en tout point hors de $E$, \end{enumerate} alors il existe une contraction de $E$, et celle-ci est la factorisation de Stein de $f$. \end{lemma}

\begin{proof} Nous appliquons Compléments sur les morphismes, théorème \ref{more-morphisms-theorem-stein-factorization-Noetherian}, pour obtenir une factorisation de Stein $X \to X' \to Y$. Alors $X \to X'$ satisfait toutes les hypothèses du lemme (certains détails sont omis). Après avoir remplacé $Y$ par $X'$, nous pouvons donc supposer en outre que $f_*\mathcal{O}_X = \mathcal{O}_Y$ et que les fibres de $f$ sont géométriquement connexes.

\medskip\noindent Supposons que $f_*\mathcal{O}_X = \mathcal{O}_Y$ et que les fibres de $f$ soient géométriquement connexes. Remarquons que $y \in Y$ est un point fermé, car $f$ est fermé et $E$ est fermé. La restriction $f^{-1}(Y \setminus \{y\}) \to Y \setminus \{y\}$ de $f$ est un morphisme fini (Compléments sur les morphismes, lemme \ref{more-morphisms-lemma-characterize-finite}). Cette restriction est donc un isomorphisme puisque $f_*\mathcal{O}_X = \mathcal{O}_Y$, les morphismes finis étant affines. Pour montrer que $\mathcal{O}_{Y, y}$ est régulier de dimension $2$, considérons l'isomorphisme $$
\mathcal{O}_{Y, y}^\wedge \longrightarrow
\lim H^0(X \times_Y \Spec(\mathcal{O}_{Y, y}/\mathfrak m_y^n), \mathcal{O})
$$ fourni par Cohomologie des schémas, lemme \ref{coherent-lemma-formal-functions-stalk}. Soit $E_n = nE$ comme dans le lemme \ref{resolve-lemma-exceptional-first-kind-local}. Observons que $$
E_n \subset X \times_Y \Spec(\mathcal{O}_{Y, y}/\mathfrak m_y^n)
$$ car $E \subset X_y = X \times_Y \Spec(\kappa(y))$. D'autre part, puisque $E = f^{-1}(\{y\})$ ensemblistement (car les fibres de $f$ sont géométriquement connexes), la fibre schématique $X_y$ est schématiquement contenue dans $E_n$ pour un certain $n > 0$. Plus précisément, appliquons le lemme \ref{coherent-lemma-power-ideal-kills-sheaf} de Cohomologie des schémas au $\mathcal{O}_X$-module cohérent $\mathcal{F} = \mathcal{O}_{X_y}$ et au faisceau d'idéaux $\mathcal{I}$ de $E$, puis utilisons le fait que $\mathcal{I}^n$ est le faisceau d'idéaux de $E_n$. Cela montre que $$
X \times_Y \Spec(\mathcal{O}_{Y, y}/\mathfrak m_y^m) \subset E_{nm}
$$ La limite projective affichée ci-dessus est donc égale à $\lim H^0(E_n, \mathcal{O}_n)$, qui est un anneau local régulier de dimension deux d'après le lemme \ref{resolve-lemma-exceptional-first-kind-local}. Ainsi, $\mathcal{O}_{Y, y}$ est un anneau local régulier de dimension deux, car son complété l'est (Compléments d'algèbre, lemme \ref{more-algebra-lemma-completion-regular} et \ref{more-algebra-lemma-completion-dimension}).

\medskip\noindent Il reste à démontrer que $f : X \to Y$ est l'éclatement $b : Y' \to Y$ de $Y$ en $y$. Nous invitons le lecteur à trouver sa propre démonstration. Tout d'abord, le lemme \ref{resolve-lemma-exceptional-first-kind-local} entraîne également l'égalité schématique $X_y = E$. Comme le faisceau d'idéaux de $E$ est inversible, cela montre que $f^{-1}\mathfrak m_y \cdot \mathcal{O}_X$ est inversible. La propriété universelle de l'éclatement fournit donc une factorisation $$
X \to Y' \to Y
$$ du morphisme $f$ ; voir Diviseurs, lemme \ref{divisors-lemma-universal-property-blowing-up}. Rappelons que la fibre exceptionnelle $E' \subset Y'$ est une courbe exceptionnelle de première espèce d'après le lemme \ref{resolve-lemma-blowup}. Soit $g : E \to E'$ le morphisme induit. Comme, tant pour $E'$ que pour $E$, le faisceau conormal est engendré par les images inverses de $a$ et $b$, l'application canonique $g^*\mathcal{C}_{E'/Y'} \to \mathcal{C}_{E/X}$ (Morphismes, lemme \ref{morphisms-lemma-conormal-functorial}) est surjective. Les deux faisceaux étant inversibles, cette application est un isomorphisme. Comme $\mathcal{C}_{E/X}$ est de degré positif, $g$ ne peut être un morphisme constant. Ainsi, $g$ est à fibres finies, et $g$ est donc un morphisme fini (même référence que ci-dessus). Or $Y'$ est régulier, donc normal, en tout point de $E'$, tandis que $X \to Y'$ est birationnel et est un isomorphisme hors de $E'$ ; on conclut alors que $X \to Y'$ est un isomorphisme d'après Variétés, lemme \ref{varieties-lemma-modification-normal-iso-over-codimension-1}. \end{proof}

\begin{lemma} \label{resolve-lemma-pic-blowup} Soit $b : X \to X'$ la contraction d'une courbe exceptionnelle de première espèce $E \subset X$. Il existe alors une suite exacte courte $$
0 \to \Pic(X') \to \Pic(X) \to \mathbf{Z} \to 0
$$ où la première application est l'image inverse par $b$ et la seconde associe à $\mathcal{L}$ le degré de $\mathcal{L}$ sur la courbe exceptionnelle $E$. Cette suite est scindée par l'application $n \mapsto \mathcal{O}_X(-nE)$. \end{lemma}

\begin{proof} Puisque $E = \mathbf{P}^1_k$, le groupe de Picard de $E$ est $\mathbf{Z}$ ; voir Diviseurs, lemme \ref{divisors-lemma-Pic-projective-space-UFD}. Nous pouvons donc considérer la dernière application comme $\mathcal{L} \mapsto \mathcal{L}|_E$. Par définition des courbes exceptionnelles de première espèce, le degré de la restriction de $\mathcal{O}_X(E)$ à $E$ est $-1$. En combinant ces remarques, il suffit de montrer que $\Pic(X') \to \Pic(X)$ est injective et que son image est formée des faisceaux inversibles dont la restriction est isomorphe au faisceau trivial $\mathcal{O}_E$ sur $E$.

\medskip\noindent Étant donné un $\mathcal{O}_{X'}$-module inversible $\mathcal{L}'$, nous affirmons que l'application $\mathcal{L}' \to b_*b^*\mathcal{L}'$ est un isomorphisme. Cela est clair partout, sauf éventuellement au point image $x \in X'$ de $E$. Pour le vérifier sur les germes en $x$, nous pouvons remplacer $X'$ par un voisinage ouvert de $x$ et supposer que $\mathcal{L}'$ est $\mathcal{O}_{X'}$. Il faut alors montrer que l'application $\mathcal{O}_{X'} \to b_*\mathcal{O}_X$ est un isomorphisme. Cela résulte de la partie (4) du lemme \ref{resolve-lemma-cohomology-of-blowup}.

\medskip\noindent Soit $\mathcal{L}$ un $\mathcal{O}_X$-module inversible tel que $\mathcal{L}|_E = \mathcal{O}_E$. Nous affirmons alors que (1) $b_*\mathcal{L}$ est inversible et que (2) $b^*b_*\mathcal{L} \to \mathcal{L}$ est un isomorphisme. Les assertions (1) et (2) sont claires au-dessus de $X' \setminus \{x\}$. Il suffit donc de les démontrer après le changement de base à $\Spec(\mathcal{O}_{X', x})$. La formation de $b_*$ commute aux changements de base plats (Cohomologie des schémas, lemme \ref{coherent-lemma-flat-base-change-cohomology}), et il en va de même pour $b^*$ et pour la formation de l'application d'adjonction. Or, si $X'$ est le spectre d'un anneau local régulier, alors $\mathcal{L}$ est trivial d'après la description du groupe de Picard donnée au lemme \ref{resolve-lemma-blowup-pic}. L'affirmation est donc démontrée.

\medskip\noindent En combinant les affirmations démontrées dans les deux paragraphes précédents, on voit que l'application $\mathcal{L} \mapsto b_*\mathcal{L}$ est inverse de l'application $$
\Pic(X') \longrightarrow \Ker(\Pic(X) \to \Pic(E))
$$, ce qui démontre le lemme. \end{proof}

\begin{remark} \label{resolve-remark-pic-blowup} Soit $b : X \to X'$ la contraction d'une courbe exceptionnelle de première espèce $E \subset X$. Le lemme \ref{resolve-lemma-pic-blowup} fournit une identification $$
\Pic(X) = \Pic(X') \oplus \mathbf{Z}
$$ dans laquelle $\mathcal{L}$ correspond au couple $(\mathcal{L}', n)$ si et seulement si $\mathcal{L} = (b^*\mathcal{L}')(-nE)$, c'est-à-dire $\mathcal{L}(nE) = b^*\mathcal{L}'$. En fait, la preuve du lemme \ref{resolve-lemma-pic-blowup} montre que $\mathcal{L}' = b_*\mathcal{L}(nE)$. Bien entendu, l'application $\mathcal{L} \mapsto \mathcal{L}'$ est un homomorphisme de groupes. \end{remark}

\begin{lemma} \label{resolve-lemma-lift-sections-and-h1} Soit $X$ un schéma noethérien. Soit $E \subset X$ une courbe exceptionnelle de première espèce. Soit $\mathcal{L}$ un $\mathcal{O}_X$-module inversible. Soit $n$ l'entier tel que $\mathcal{L}|_E$ soit de degré $n$ lorsqu'on le considère comme module inversible sur $\mathbf{P}^1$. Alors \begin{enumerate} \item Si $H^1(X, \mathcal{L}) = 0$ et $n \geq 0$, alors $H^1(X, \mathcal{L}(iE)) = 0$ pour $0 \leq i \leq n + 1$. \item Si $n \leq 0$, alors $H^1(X, \mathcal{L}) \subset H^1(X, \mathcal{L}(E))$. \end{enumerate} \end{lemma}

\begin{proof} Observons que $\mathcal{L}|_E = \mathcal{O}(n)$ d'après Diviseurs, lemme \ref{divisors-lemma-Pic-projective-space-UFD}. Procédons par récurrence en utilisant la suite exacte longue de cohomologie associée à la suite exacte courte $$
0 \to \mathcal{L} \to \mathcal{L}(E) \to \mathcal{L}(E)|_E \to 0,
$$, ainsi que le fait que $H^1(\mathbf{P}^1, \mathcal{O}(d)) = 0$ pour $d \geq -1$ et que $H^0(\mathbf{P}^1, \mathcal{O}(d)) = 0$ pour $d \leq -1$. Nous omettons quelques détails. \end{proof}

\begin{lemma} \label{resolve-lemma-contract-ample} Soit $S = \Spec(R)$ un schéma noethérien affine. Soit $X \to S$ un morphisme propre. Soit $\mathcal{L}$ un faisceau inversible ample sur $X$. Soit $E \subset X$ une courbe exceptionnelle de première espèce. Alors \begin{enumerate} \item il existe une contraction $b : X \to X'$ de $E$, \item $X'$ est propre sur $S$, et \item le $\mathcal{O}_{X'}$-module inversible $\mathcal{L}'$ est ample, où $\mathcal{L}'$ est défini comme dans la remarque \ref{resolve-remark-pic-blowup}. \end{enumerate} \end{lemma}

\begin{proof} Soit $n$ le degré de $\mathcal{L}|_E$ comme dans le lemme \ref{resolve-lemma-lift-sections-and-h1}. Remarquons que $n > 0$, car $\mathcal{L}$ est ample sur $E$ (Variétés, lemme \ref{varieties-lemma-ample-curve}, et Propriétés, lemme \ref{properties-lemma-ample-on-closed}). Après avoir remplacé $\mathcal{L}$ par une puissance, nous pouvons supposer que $H^i(X, \mathcal{L}^{\otimes e}) = 0$ pour tous $i > 0$ et $e > 0$, voir Cohomologie des schémas, lemme \ref{coherent-lemma-vanshing-gives-ample}. Enfin, après avoir remplacé $\mathcal{L}$ par une autre puissance, nous pouvons supposer qu'il existe des sections globales $t_0, \ldots, t_n$ de $\mathcal{L}$ qui définissent une immersion fermée $\psi : X \to \mathbf{P}^n_S$, voir Morphismes, lemme \ref{morphisms-lemma-finite-type-over-affine-ample-very-ample}.

\medskip\noindent Posons $\mathcal{M} = \mathcal{L}(nE)$. Alors $\mathcal{M}|_E \cong \mathcal{O}_E$. Comme nous avons la suite exacte courte $$
0 \to \mathcal{M}(-E) \to \mathcal{M} \to \mathcal{O}_E \to 0
$$ et que $H^1(X, \mathcal{M}(-E))$ est nul (d'après le lemme \ref{resolve-lemma-lift-sections-and-h1} et le fait que $n > 0$), nous pouvons choisir une section $s_{n + 1}$ de $\mathcal{M}$ qui engendre $\mathcal{M}|_E$. Enfin, notons $s_0, \ldots, s_n$ les sections de $\mathcal{M}$ obtenues à partir des sections $t_0, \ldots, t_n$ de $\mathcal{L}$ choisies ci-dessus au moyen de $\mathcal{L} \subset \mathcal{L}(nE) = \mathcal{M}$. Prises ensemble, les sections $s_0, \ldots, s_n, s_{n + 1}$ engendrent $\mathcal{M}$ en tout point de $X$ et définissent donc un morphisme $$
\varphi : X \longrightarrow \mathbf{P}^{n + 1}_S
$$ sur $S$, voir Constructions, lemme \ref{constructions-lemma-projective-space}.

\medskip\noindent Nous allons vérifier ci-dessous les conditions du lemme \ref{resolve-lemma-contraction}. Une fois cela fait, nous voyons que la factorisation de Stein $X \to X' \to \mathbf{P}^{n + 1}_S$ de $\varphi$ est la contraction cherchée, ce qui démontre (1). De plus, le morphisme $X' \to \mathbf{P}^{n + 1}_S$ est fini ; par conséquent, $X'$ est propre sur $S$ (Morphismes, lemmes \ref{morphisms-lemma-finite-proper} et \ref{morphisms-lemma-composition-proper}). Cela démontre (2). Remarquons que $X'$ possède un faisceau inversible ample. En effet, l'image inverse $\mathcal{M}'$ de $\mathcal{O}_{\mathbf{P}^{n + 1}_S}(1)$ est ample d'après Morphismes, lemme \ref{morphisms-lemma-pullback-ample-tensor-relatively-ample}. Remarquons que $\mathcal{M}'$ a pour image inverse $\mathcal{M}$ sur $X$ (d'après Constructions, lemme \ref{constructions-lemma-projective-space}). Enfin, $\mathcal{M} = \mathcal{L}(nE)$. Puisque, dans les arguments ci-dessus, nous avons remplacé le faisceau $\mathcal{L}$ initial par une puissance positive, nous concluons que le $\mathcal{O}_{X'}$-module inversible $\mathcal{L}'$ mentionné en (3) du lemme est ample sur $X'$ d'après Propriétés, lemme \ref{properties-lemma-ample-power-ample}.

\medskip\noindent Observations immédiates : $\mathbf{P}^{n + 1}_S$ est noethérien et $\varphi$ est propre. Nous omettons les détails.

\medskip\noindent Remarquons ensuite que tout point de $U = X \setminus E$ est envoyé dans le sous-schéma ouvert $W$ de $\mathbf{P}^{n + 1}_S$ où l'une des $n + 1$ premières coordonnées homogènes est non nulle. En revanche, tout point de $E$ est envoyé sur un point dont les $n + 1$ premières coordonnées homogènes sont toutes nulles, donc en particulier dans le complémentaire de $W$. De plus, il est clair qu'il existe une factorisation $$
U = \varphi^{-1}(W) \xrightarrow{\varphi|_U}
W \xrightarrow{\text{pr}} \mathbf{P}^n_S
$$ de $\psi|_U$, où $\text{pr}$ est la projection définie par les $n + 1$ premières coordonnées et $\psi : X \to \mathbf{P}^n_S$ est l'immersion choisie ci-dessus. Il s'ensuit que $\varphi|_U : U \to W$ est quasi-fini.

\medskip\noindent Enfin, considérons l'application $\varphi|_E : E \to \mathbf{P}^{n + 1}_S$. Pour tout point $x \in E$, l'image $\varphi(x)$ a ses $n + 1$ premières coordonnées nulles ; autrement dit, le morphisme $\varphi|_E$ se factorise par le sous-schéma fermé $\mathbf{P}^0_S \cong S$. Le morphisme $E \to S = \Spec(R)$ se factorise sous la forme $E \to \Spec(H^0(E, \mathcal{O}_E)) \to \Spec(R)$ d'après Schémas, lemme \ref{schemes-lemma-morphism-into-affine}. Comme, par hypothèse, $H^0(E, \mathcal{O}_E)$ est un corps, nous concluons que $E$ s'envoie sur un point de $S \subset \mathbf{P}^{n + 1}_S$, ce qui achève la démonstration. \end{proof}

\begin{lemma} \label{resolve-lemma-contract-when-quasi-projective} Soit $S$ un schéma noethérien. Soit $f : X \to S$ un morphisme de type fini. Soit $E \subset X$ une courbe exceptionnelle de première espèce contenue dans une fibre de $f$. \begin{enumerate} \item Si $X$ est projectif sur $S$, alors il existe une contraction $X \to X'$ de $E$ et $X'$ est projectif sur $S$. \item Si $X$ est quasi-projectif sur $S$, alors il existe une contraction $X \to X'$ de $E$ et $X'$ est quasi-projectif sur $S$. \end{enumerate} \end{lemma}

\begin{proof} Les deux cas résultent du lemme \ref{resolve-lemma-contract-ample} et des résultats classiques sur les modules inversibles amples et les morphismes (quasi-)projectifs.

\medskip\noindent Démonstration de (1). La projectivité de $f$ signifie que $f$ est propre et qu'il existe un module inversible $f$-ample $\mathcal{L}$, voir Morphismes, lemme \ref{morphisms-lemma-projective-is-quasi-projective-proper} et définition \ref{morphisms-definition-quasi-projective}. Soit $U \subset S$ un ouvert affine contenant l'image de $E$. D'après le lemme \ref{resolve-lemma-contract-ample}, il existe une contraction $c : f^{-1}(U) \to V'$ de $E$ et un module inversible ample $\mathcal{N}'$ sur $V'$ dont l'image inverse sur $f^{-1}(U)$ est égale à $\mathcal{L}(nE)|_{f^{-1}(U)}$. Soit $v \in V'$ le point fermé tel que $c$ soit l'éclatement de $v$. Nous pouvons alors recoller $V'$ et $X \setminus E$ le long de $f^{-1}(U) \setminus E = V' \setminus \{v\}$ pour obtenir un schéma $X'$ sur $S$. Les morphismes $c$ et $\text{id}_{X \setminus E}$ se recollent en un morphisme $b : X \to X'$ qui est la contraction de $E$. L'image inverse de $U$ dans $X'$ est propre sur $U$. D'autre part, la restriction de $X' \to S$ au complémentaire de l'image de $v$ dans $S$ est isomorphe à la restriction de $X \to S$ à cet ouvert. Ainsi, $X' \to S$ est propre, puisque la propreté est locale sur la base d'après Morphismes, lemme \ref{morphisms-lemma-proper-local-on-the-base}. Enfin, les restrictions de $\mathcal{N}'$ et de $\mathcal{L}|_{X \setminus E}$ à $f^{-1}(U) \setminus E = V' \setminus \{v\}$ sont isomorphes ; elles se recollent donc en un module inversible $\mathcal{L}'$ sur $X'$. La restriction de $\mathcal{L}'$ à l'image inverse de $U$ dans $X'$ est ample, car il en est ainsi de $\mathcal{N}'$. Sur les ouverts affines de $S$ qui évitent l'image de $v$, il en va de même puisque c'est vrai pour $\mathcal{L}$. Ainsi, $\mathcal{L}'$ est relativement ample pour $(X' \to S)$ d'après Morphismes, lemme \ref{morphisms-lemma-characterize-relatively-ample}, et (1) est démontré.

\medskip\noindent Démonstration de (2). Nous pouvons écrire $X$ comme un sous-schéma ouvert d'un schéma $\overline{X}$ projectif sur $S$ d'après Morphismes, lemme \ref{morphisms-lemma-quasi-projective-open-projective}. D'après (1), il existe une contraction $b : \overline{X} \to \overline{X}'$, et $\overline{X}'$ est projectif sur $S$. Nous notons alors $X' \subset \overline{X}$ l'image de $X \to \overline{X}'$ ; c'est un ouvert puisque $b$ est un isomorphisme hors de $E$. Alors $X \to X'$ est la contraction cherchée. Remarquons que $X'$ est quasi-projectif sur $S$, car il possède un module inversible ample relativement à $S$, par la construction effectuée dans la démonstration de la partie (1). \end{proof}

\begin{lemma} \label{resolve-lemma-regular-dim-2-quasi-projective} Soit $S$ un schéma noethérien. Soit $f : X \to S$ un morphisme séparé de type fini, avec $X$ régulier de dimension $2$. Alors $X$ est quasi-projectif sur $S$. \end{lemma}

\begin{proof} Par le lemme de Chow (Cohomologie des schémas, lemme \ref{coherent-lemma-chow-Noetherian}), il existe un morphisme propre $\pi : X' \to X$ qui est un isomorphisme au-dessus d'un ouvert dense $U \subset X$ tel que $X' \to S$ soit H-quasi-projectif. D'après le lemme \ref{resolve-lemma-extend-rational-map-blowing-up}, il existe une suite d'éclatements en des points fermés $$
X_n \to \ldots \to X_1 \to X_0 = X
$$ et un $S$-morphisme $X_n \to X'$ prolongeant l'application rationnelle $U \to X'$. Remarquons que $X_n \to X$ est projectif d'après Diviseurs, lemme \ref{divisors-lemma-blowing-up-projective}, et Morphismes, lemme \ref{morphisms-lemma-composition-projective}. Il s'ensuit que $X_n \to X'$ est projectif d'après Morphismes, lemme \ref{morphisms-lemma-projective-permanence}. Par conséquent, $X_n \to S$ est quasi-projectif d'après Morphismes, lemme \ref{morphisms-lemma-composition-quasi-projective} (et le fait qu'un morphisme projectif est quasi-projectif ; voir Morphismes, lemme \ref{morphisms-lemma-projective-quasi-projective}). D'après le lemme \ref{resolve-lemma-contract-when-quasi-projective} (et l'unicité des contractions, lemme \ref{resolve-lemma-contraction-unique}), nous concluons que $X_{n - 1}, \ldots, X_0 = X$ sont quasi-projectifs sur $S$, comme souhaité. \end{proof}

\begin{lemma} \label{resolve-lemma-regular-dim-2-projective} Soit $S$ un schéma noethérien. Soit $f : X \to S$ un morphisme propre tel que $X$ soit régulier de dimension $2$. Alors $X$ est projectif sur $S$. \end{lemma}

\begin{proof} Cela résulte du lemme \ref{resolve-lemma-regular-dim-2-quasi-projective} et de Morphismes, lemme \ref{morphisms-lemma-projective-is-quasi-projective-proper}. \end{proof}






\section{Factorisation des morphismes birationnels} \label{resolve-section-factorizing}

\noindent Les morphismes propres birationnels entre surfaces non singulières sont des composés de transformations quadratiques.

\begin{lemma} \label{resolve-lemma-proper-birational-regular-surfaces} Soit $f : X \to Y$ un morphisme propre birationnel entre schémas noethériens intègres réguliers de dimension $2$. Alors $f$ est un composé d'éclatements en des points fermés. \end{lemma}

\begin{proof} Soit $V \subset Y$ le plus grand ouvert au-dessus duquel $f$ est un isomorphisme. Alors $V$ contient tous les points de codimension $1$ de $V$ (Variétés, lemme \ref{varieties-lemma-modification-normal-iso-over-codimension-1}). Soit $y \in Y$ un point fermé qui n'appartient pas à $V$. Nous voulons montrer que $f$ se factorise par l'éclatement $b : Y' \to Y$ de $Y$ en $y$. En effet, si tel est le cas, au moins une composante (et, en fait, exactement une) de la fibre $f^{-1}(y)$ s'enverra isomorphiquement sur la courbe exceptionnelle de $Y'$, et le nombre de courbes contenues dans les fibres de $X \to Y'$ sera strictement inférieur au nombre de courbes contenues dans les fibres de $X \to Y$ ; nous concluons donc par récurrence. Certains détails sont omis.

\medskip\noindent D'après le lemme \ref{resolve-lemma-extend-rational-map-blowing-up}, il existe une suite d'éclatements $$
X' = X_n \to X_{n - 1} \to \ldots \to X_1 \to X_0 = X
$$ en des points fermés situés au-dessus de la fibre $f^{-1}(y)$ et un morphisme $X' \to Y'$ tels que $$
\xymatrix{
X' \ar[d]_{f'} \ar[r] & X \ar[d]^f \\
Y' \ar[r] & Y
}
$$ soit commutatif. Nous voulons montrer que le morphisme $X' \to Y'$ se factorise par $X$ ; nous pourrons alors raisonner par récurrence sur $n$ pour nous ramener au cas où $X' \to X$ est l'éclatement de $X$ en un point fermé $x \in X$ qui s'envoie sur $y$.

\medskip\noindent Soit $E \subset X'$ la fibre exceptionnelle de l'éclatement $X' \to X$. Si $E$ s'envoie sur un point de $Y'$, alors le lemme \ref{resolve-lemma-factor-through-contraction} fournit la factorisation souhaitée. Nous allons montrer que, dans le cas contraire, nous obtenons une contradiction. En effet, si $f'(E)$ n'est pas un point, alors $E' = f'(E)$ est nécessairement la courbe exceptionnelle de $Y'$. Considérons le diagramme $$
\xymatrix{
E \ar[r] \ar[d]_g & X' \ar[d]_{f'} \ar[r] & X \ar[d]^f \\
E' \ar[r] & Y' \ar[r] & Y
}
$$ En raisonnant comme précédemment, $f'$ est un isomorphisme dans un voisinage ouvert du point générique de $E'$. Ainsi, $g : E \to E'$ est un morphisme fini birationnel. L'inverse de $g$, qui est une application rationnelle, est alors partout défini d'après Morphismes, lemme \ref{morphisms-lemma-extend-across}, et $g$ est un isomorphisme. Considérons l'application $$
g^*\mathcal{C}_{E'/Y'} \longrightarrow \mathcal{C}_{E/X'}
$$ de Morphismes, lemme \ref{morphisms-lemma-conormal-functorial}. Comme la source et le but sont des modules inversibles de degré $1$ sur $E = E' = \mathbf{P}^1_\kappa$ et que l'application est non nulle (puisque $f'$ est un isomorphisme au point générique de $E$), nous concluons que c'est un isomorphisme. D'après Morphismes, lemme \ref{morphisms-lemma-two-immersions}, nous concluons que $\Omega_{X'/Y'}|_E = 0$. Cela signifie que $f'$ est non ramifié en tout point de $E$ (Morphismes, lemme \ref{morphisms-lemma-unramified-at-point}). Ainsi, $f'$ est quasi-fini en tout point de $E$ (Morphismes, lemme \ref{morphisms-lemma-unramified-quasi-finite}). Par conséquent, le plus grand ouvert $V' \subset Y'$ au-dessus duquel $f'$ est un isomorphisme contient $E'$ d'après Variétés, lemme \ref{varieties-lemma-modification-normal-iso-over-codimension-1}. Il s'ensuit que l'image inverse de $y$ dans $X'$ est $E'$. L'image inverse de $y$ dans $X$ est donc $x$. Par conséquent, $x \in X$ appartient au plus grand ouvert au-dessus duquel $f$ est un isomorphisme, d'après Variétés, lemme \ref{varieties-lemma-modification-normal-iso-over-codimension-1}. Cela contredit l'hypothèse selon laquelle $y$ n'appartient pas à cet ouvert. \end{proof}

\begin{lemma} \label{resolve-lemma-birational-regular-surfaces} Soit $S$ un schéma noethérien. Soient $X$ et $Y$ des schémas propres intègres sur $S$, réguliers de dimension $2$. Alors $X$ et $Y$ sont $S$-birationnels si et seulement s'il existe un diagramme de $S$-morphismes $$
X = X_0 \leftarrow X_1 \leftarrow \ldots \leftarrow X_n = Y_m
\to \ldots \to Y_1 \to Y_0 = Y
$$ dans lequel chaque morphisme est l'éclatement d'un point fermé. \end{lemma}

\begin{proof} Soit $U \subset X$ un ouvert et soit $f : U \to Y$ l'application $S$-rationnelle donnée (qui est inversible comme application $S$-rationnelle). D'après le lemme \ref{resolve-lemma-extend-rational-map-blowing-up}, nous pouvons factoriser $f$ sous la forme $X_n \to \ldots \to X_1 \to X_0 = X$ et $f_n : X_n \to Y$. Comme $X_n$ est propre sur $S$ et $Y$ séparé sur $S$, le morphisme $f_n$ est propre. Il est clair que $f_n$ est birationnel. Le morphisme $f_n$ est donc un composé de contractions d'après le lemme \ref{resolve-lemma-proper-birational-regular-surfaces}. Nous omettons la démonstration de la réciproque. \end{proof}










% OK, start here.
%

\chapter{Réduction semi-stable}



\phantomsection
\label{models-section-phantom}


\section{Introduction}
\label{models-section-introduction}

\noindent
Dans ce chapitre, nous démontrons le théorème de réduction semi-stable pour les courbes.
Nous utiliserons la méthode d'Artin et Winters exposée dans leur article
\cite{Artin-Winters}.

\medskip\noindent
On peut en fait démontrer le théorème de réduction semi-stable
pour les courbes sans aucun résultat de désingularisation. En effet, il
existe une méthode pour établir l'existence et la projectivité des espaces de modules
de courbes semi-stables au moyen de la théorie géométrique des invariants (GIT)
développée par Mumford, voir \cite{GIT}. Cette méthode a été
défendue par Gieseker, qui a démontré le résultat complet dans ses
notes de cours \cite{Gieseker}\footnote{Les notes de cours de Gieseker
sont rédigées sur un corps algébriquement clos, mais la même méthode
s'applique sur $\mathbf{Z}$.}. C'est là une prouesse
assez remarquable : il semble quelque peu contraire à l'intuition que l'on puisse
démontrer un tel résultat sans jamais étudier véritablement les familles de courbes sur
une base de dimension strictement positive.

\medskip\noindent
Historiquement, la première démonstration du théorème de réduction semi-stable
pour les courbes figure dans l'article \cite{DM} de Deligne et Mumford.
Elle démontre le théorème en ramenant le problème au cas des
variétés abéliennes, déjà connu à l'époque grâce
à Grothendieck et à d'autres, voir \cite{SGA7-I} et \cite{SGA7-II}).
Le théorème de réduction semi-stable pour les variétés abéliennes utilise la théorie
des modèles de Néron, qui repose à son tour sur une étude des lois de groupe
birationnelles sur une base.

\medskip\noindent
La méthode de l'article d'Artin et Winters repose sur la désingularisation
des surfaces pour obtenir des modèles réguliers. Une fois acquise l'existence
de modèles réguliers, la démonstration consiste à analyser les
possibilités pour la fibre spéciale, puis à conclure au moyen d'une inégalité
sur la torsion du groupe de Picard d'un schéma de dimension $1$ sur un corps.
Un argument analogue figure dans un article \cite{Saito} de Saito, qui utilise
directement la cohomologie étale et obtient un résultat plus fort, puisqu'il
peut caractériser la réduction semi-stable en termes de l'action de
l'inertie sur la cohomologie étale $\ell$-adique.

\medskip\noindent
Une autre approche pour démontrer le théorème consiste à utiliser
les techniques de géométrie analytique rigide. Nous renvoyons ici le lecteur à
\cite{vanderPut} et \cite{Arzdorf-Wewers}.

\medskip\noindent
L'article \cite{Temkin} de Temkin utilise des techniques de théorie des valuations
(et démontre bien d'autres résultats) ; son appendice A donne également
un bon aperçu des différentes démonstrations et de leur rapport avec
les désingularisations des schémas de dimension $2$.

\medskip\noindent
Le lecteur pourra également consulter l'article de synthèse
\cite{Abbes-ssr} d'Ahmed Abbes.






\section{Algèbre linéaire}
\label{models-section-linear-algebra}

\noindent
Voici quelques lemmes que nous utiliserons par la suite.

\begin{lemma}
\label{models-lemma-recurring}
\begin{reference}
\cite[Theorem I]{Taussky}
\end{reference}
Soit $A = (a_{ij})$ une matrice complexe $n \times n$.
\begin{enumerate}
\item Si $|a_{ii}| > \sum_{j \not = i} |a_{ij}|$ pour tout $i$, alors
$\det(A)$ est non nul.
\item S'il existe un vecteur réel $m = (m_1, \ldots, m_n)$
avec $m_i > 0$ tel que $|a_{ii} m_i| > \sum_{j \not = i} |a_{ij}m_j|$
pour tout $i$, alors $\det(A)$ est non nul.
\end{enumerate}
\end{lemma}

\begin{proof}
Si $A$ est comme en (1) et $\det(A) = 0$, il existe un vecteur non nul
$z$ tel que $Az = 0$. Choisissons $r$ tel que $|z_r|$ soit maximal. Alors
$$
|a_{rr} z_r| = |\sum\nolimits_{k \not = r} a_{rk}z_k| \leq
\sum\nolimits_{k \not = r} |a_{rk}||z_k| \leq
|z_r| \sum\nolimits_{k \not = r} |a_{rk}| < |a_{rr}||z_r|
$$
ce qui est contradictoire. Pour démontrer (2), appliquons (1) à la matrice
$(a_{ij}m_j)$, dont le déterminant est $m_1 \ldots m_n \det(A)$.
\end{proof}

\begin{lemma}
\label{models-lemma-recurring-real}
Soit $A = (a_{ij})$ une matrice réelle $n \times n$ telle que
$a_{ij} \geq 0$ pour $i \not = j$. Soit $m = (m_1, \ldots, m_n)$ un vecteur réel
avec $m_i > 0$. Pour $I \subset \{1, \ldots, n\}$, soit
$x_I \in \mathbf{R}^n$
le vecteur dont la coordonnée d'indice $i$ est $m_i$ si $i \in I$
et $0$ sinon. Si
\begin{equation}
\label{models-equation-ineq}
-a_{ii}m_i \geq \sum\nolimits_{j \not = i} a_{ij}m_j
\end{equation}
pour tout $i$, alors $\Ker(A)$ est l'espace vectoriel
engendré par les vecteurs $x_I$ tels que
\begin{enumerate}
\item $a_{ij} = 0$ pour $i \in I$, $j \not \in I$, et
\item on a l'égalité dans (\ref{models-equation-ineq}) pour $i \in I$.
\end{enumerate}
\end{lemma}

\begin{proof}
En remplaçant $a_{ij}$ par $a_{ij}m_j$, on peut supposer $m_i = 1$ pour tout $i$.
Si $I \subset \{1, \ldots, n\}$ est tel que (1) et (2) soient vérifiées,
un calcul immédiat montre que $x_I$ appartient au noyau de $A$.
Réciproquement, soit $x = (x_1, \ldots, x_n) \in \mathbf{R}^n$ un
vecteur non nul du noyau de $A$. Nous allons montrer, par récurrence
sur le nombre de coordonnées non nulles de $x$, que $x$ appartient au
sous-espace engendré par les vecteurs $x_I$ vérifiant (1) et (2). Soit
$I \subset \{1, \ldots, n\}$ l'ensemble des indices $r$ pour lesquels $|x_r|$ est maximal.
Pour $r \in I$, on a
$$
|a_{rr} x_r| = |\sum\nolimits_{k \not = r} a_{rk}x_k| \leq
\sum\nolimits_{k \not = r} a_{rk}|x_k| \leq
|x_r| \sum\nolimits_{k \not = r} a_{rk} \leq |a_{rr}||x_r|
$$
On a donc l'égalité partout. En particulier, on voit que
$a_{rk} = 0$ si $r \in I$, $k \not \in I$, et que l'égalité est vérifiée
dans (\ref{models-equation-ineq}) pour $r \in I$. On voit alors que l'on
peut retrancher à $x$ un multiple convenable de $x_I$ pour diminuer
le nombre de coordonnées non nulles.
\end{proof}

\begin{lemma}
\label{models-lemma-recurring-symmetric-real}
Soit $A = (a_{ij})$ une matrice réelle symétrique $n \times n$ telle que
$a_{ij} \geq 0$ pour $i \not = j$.
Soit $m = (m_1, \ldots, m_n)$ un vecteur réel avec $m_i > 0$.
Supposons que
\begin{enumerate}
\item $Am = 0$,
\item il n'existe aucune partie propre non vide $I \subset \{1, \ldots, n\}$
telle que $a_{ij} = 0$ pour $i \in I$ et $j \not \in I$.
\end{enumerate}
Alors $x^t A x \leq 0$, avec égalité si et seulement si $x = qm$
pour un certain $q \in \mathbf{R}$.
\end{lemma}

\begin{proof}[Première démonstration]
En remplaçant $a_{ij}$ par $a_{ij}m_im_j$, on peut supposer $m_i = 1$
pour tout $i$. La condition (1) signifie que $-a_{ii} = \sum_{j \not = i} a_{ij}$
pour tout $i$. Rappelons que $x^tAx = \sum_{i, j} x_ia_{ij}x_j$.
Alors
\begin{align*}
\sum\nolimits_{i \not = j} -a_{ij}(x_j - x_i)^2 & =
\sum\nolimits_{i \not = j} -a_{ij}x_j^2 + 2a_{ij}x_ix_i - a_{ij}x_i^2 \\
& =
\sum\nolimits_j a_{jj} x_j^2 +
\sum\nolimits_{i \not = j} 2a_{ij}x_ix_i +
\sum\nolimits_j a_{jj} x_i^2 \\
& = 2x^tAx
\end{align*}
Cette expression est manifestement $\leq 0$. S'il y a égalité, soit $I$ l'ensemble
des indices $i$ tels que $x_i \not = x_1$. Alors $a_{ij} = 0$ pour $i \in I$
et $j \not \in I$. Ainsi $I = \{1, \ldots, n\}$ par la condition (2), et
$x$ est un multiple de $m = (1, \ldots, 1)$.
\end{proof}

\begin{proof}[Seconde démonstration]
La matrice $A$ a des valeurs propres réelles, d'après le théorème spectral.
Nous affirmons que toutes ses valeurs propres sont $\leq 0$.
En effet, puisque la propriété (1) signifie que
$-a_{ii}m_i = \sum_{j \not = i} a_{ij}m_j$ pour tout $i$,
la matrice $A' = A - \lambda I$, pour $\lambda > 0$,
vérifie $|a'_{ii}m_i| > \sum a'_{ij}m_j = \sum |a'_{ij}m_j|$ pour tout $i$.
Ainsi $A'$ est inversible par le lemme \ref{models-lemma-recurring}.
Il en résulte que la forme bilinéaire symétrique $x^tAy$
est semi-définie négative, c'est-à-dire que $x^tAx \leq 0$ pour tout $x$.
Le noyau de $A$ est donc égal
à l'ensemble des vecteurs $x$ tels que $x^tAx = 0$.
La description du noyau dans le lemme \ref{models-lemma-recurring-real}
donne la dernière assertion du lemme.
\end{proof}

\begin{lemma}
\label{models-lemma-orthogonal-direct-sum}
Soit $L$ un $\mathbf{Z}$-module libre de type fini muni
d'une forme bilinéaire symétrique entière définie positive
$\langle\ ,\ \rangle : L \times L \to \mathbf{Z}$.
Soit $A \subset L$ un sous-module tel que $L/A$ soit sans torsion. Posons
$B = \{b \in L \mid \langle a, b\rangle = 0,\ \forall a \in A\}$.
On dispose alors d'applications injectives
$$
A^\#/A \leftarrow L/(A \oplus B) \rightarrow B^\#/B
$$
dont les conoyaux sont des quotients de $L^\#/L$. Ici
$A^\# = \{a' \in A \otimes \mathbf{Q} \mid
\langle a, a'\rangle \in \mathbf{Z},\ \forall a \in A\}$
et de même pour $B$ et $L$.
\end{lemma}

\begin{proof}
Observons que
$L \otimes \mathbf{Q} = A \otimes \mathbf{Q} \oplus B \otimes \mathbf{Q}$
car la forme est non dégénérée sur $A$ (par positivité).
On note $\pi_B : L \otimes \mathbf{Q} \to B \otimes \mathbf{Q}$
la projection. Observons que $\pi_B(x) \in B^\#$ pour $x \in L$,
car la forme est entière. On obtient ainsi une suite exacte
$$
0 \to A \to L \xrightarrow{\pi_B} B^\# \to Q \to 0
$$
où $Q$ est le conoyau de $L \to B^\#$. Observons que $Q$
est un quotient de $L^\#/L$, car l'application $L^\# \to B^\#$ est surjective,
étant la duale $\mathbf{Z}$-linéaire de $B \to L$, qui est scindée
comme application de $\mathbf{Z}$-modules.
En quotientant par $A \oplus B$, on obtient une suite exacte courte
$$
0 \to L/(A \oplus B) \to B^\#/B \to Q \to 0
$$
Ceci démontre le lemme.
\end{proof}

\begin{lemma}
\label{models-lemma-coker}
Soient $L_0$, $L_1$ des $\mathbf{Z}$-modules libres de type fini munis
de formes bilinéaires symétriques entières définies positives
$\langle\ ,\ \rangle : L_i \times L_i \to \mathbf{Z}$.
Soient $\text{d} : L_0 \to L_1$ et $\text{d}^* : L_1 \to L_0$
des applications adjointes. Si $\langle\ ,\ \rangle$ sur $L_0$ est unimodulaire,
il existe un isomorphisme
$$
\Phi :
\Coker(\text{d}^*\text{d})_{torsion}
\longrightarrow
\Im(\text{d})^\#/\Im(\text{d})
$$
avec les notations du lemme \ref{models-lemma-orthogonal-direct-sum}.
\end{lemma}

\begin{proof}
Soit $x \in L_0$ un élément représentant une classe de torsion
dans $\Coker(\text{d}^*\text{d})$.
Pour un certain $a > 0$, on peut donc écrire $ax = \text{d}^*\text{d}(y)$.
Pour tout $z \in \Im(\text{d})$, disons $z = \text{d}(y')$, on a
$$
\langle (1/a)\text{d}(y), z \rangle =
\langle (1/a)\text{d}(y), \text{d}(y') \rangle =
\langle x, y' \rangle \in \mathbf{Z}
$$
Ainsi $(1/a)\text{d}(y) \in \Im(\text{d})^\#$. On définit
$\Phi(x) = (1/a)\text{d}(y) \bmod \Im(\text{d})$.
Nous omettons la preuve que $\Phi$ est bien définie, additive et injective.

\medskip\noindent
Pour montrer que $\Phi$ est surjective, soit $z \in \Im(\text{d})^\#$.
Alors $z$ définit une application linéaire $L_0 \to \mathbf{Z}$
par la règle $x \mapsto \langle z, \text{d}(x)\rangle$.
L'accouplement sur $L_0$ étant unimodulaire par hypothèse,
on peut trouver un $x' \in L_0$ tel que
$\langle x', x \rangle = \langle z, \text{d}(x)\rangle$
pour tout $x \in L_0$. En particulier, on voit
que $x'$ est orthogonal à $\Ker(\text{d})$.
Puisque $\Im(\text{d}^*\text{d}) \otimes \mathbf{Q}$
est le supplémentaire orthogonal de $\Ker(\text{d}) \otimes \mathbf{Q}$,
cela signifie que $x'$ définit une classe de torsion dans
$\Coker(\text{d}^*\text{d})$. Nous affirmons que $\Phi(x') = z$.
En effet, écrivons $a x' = \text{d}^*\text{d}(y)$
pour certains $y \in L_0$ et $a > 0$.
Pour tout $x \in L_0$, on obtient
$$
\langle z, \text{d}(x)\rangle =
\langle x', x \rangle =
\langle (1/a)\text{d}^*\text{d}(y), x \rangle =
\langle (1/a)\text{d}(y),\text{d}(x) \rangle
$$
Ainsi $z = \Phi(x')$, ce qui achève la démonstration.
\end{proof}

\begin{lemma}
\label{models-lemma-recurring-symmetric-integer}
Soit $A = (a_{ij})$ une matrice entière symétrique $n \times n$ telle que
$a_{ij} \geq 0$ pour $i \not = j$. Soit $m = (m_1, \ldots, m_n)$ un
vecteur entier avec $m_i > 0$. Supposons que
\begin{enumerate}
\item $Am = 0$,
\item il n'existe aucune partie propre non vide $I \subset \{1, \ldots, n\}$
telle que $a_{ij} = 0$ pour $i \in I$ et $j \not \in I$.
\end{enumerate}
Soit $e$ le nombre de couples $(i, j)$ tels que $i < j$ et $a_{ij} > 0$.
Alors, pour un nombre premier $\ell$ premier à tous les $a_{ij}$ et $m_i$,
on a
$$
\dim_{\mathbf{F}_\ell}(\Coker(A)[\ell]) \leq 1 - n + e
$$
\end{lemma}

\begin{proof}
D'après le lemme \ref{models-lemma-recurring-symmetric-real}, le rang de $A$ est $n - 1$.
La composée
$$
\mathbf{Z}^{\oplus n} \xrightarrow{\text{diag}(m_1, \ldots, m_n)}
\mathbf{Z}^{\oplus n} \xrightarrow{(a_{ij})}
\mathbf{Z}^{\oplus n} \xrightarrow{\text{diag}(m_1, \ldots, m_n)}
\mathbf{Z}^{\oplus n}
$$
a pour matrice $a_{ij}m_im_j$. Puisque les conoyaux de la première et de la dernière
application sont de torsion d'ordre premier à $\ell$, par notre condition sur $\ell$,
il suffit de démontrer le lemme pour la matrice
de coefficients $a_{ij}m_im_j$. On peut donc supposer $m = (1, \ldots, 1)$.

\medskip\noindent
Supposons $m = (1, \ldots, 1)$. Posons $V = \{1, \ldots, n\}$ et
$E = \{(i, j) \mid i < j\text{ et }a_{ij} > 0\}$. Pour
$e = (i, j) \in E$, posons $a_e = a_{ij}$. Définissons des applications
$s, t : E \to V$ en posant $s(i, j) = i$ et $t(i, j) = j$.
Posons
$\mathbf{Z}(V) = \bigoplus_{i \in V} \mathbf{Z}i$ et
$\mathbf{Z}(E) = \bigoplus_{e \in E} \mathbf{Z}e$.
On définit des accouplements symétriques définis positifs à valeurs entières
sur $\mathbf{Z}(V)$ et $\mathbf{Z}(E)$ en posant
$$
\langle i, i \rangle = 1\text{ pour }i \in V, \quad
\langle e, e \rangle = a_e\text{ pour }e \in E
$$
et en prenant tous les autres accouplements nuls. Considérons les applications
$$
\text{d} : \mathbf{Z}(V) \to \mathbf{Z}(E), \quad
i \longmapsto
\sum\nolimits_{e \in E,\ s(e) = i} e - \sum\nolimits_{e \in E,\ t(e) = i} e
$$
et
$$
\text{d}^*(e) = a_e(s(e) - t(e))
$$
Un calcul montre que
$$
\langle d(x), y\rangle = \langle x, \text{d}^*(y) \rangle
$$
autrement dit, $\text{d}$ et $\text{d}^*$ sont adjointes. On calcule ensuite
\begin{align*}
\text{d}^*\text{d}(i)
& = 
\text{d}^*(
\sum\nolimits_{e \in E,\ s(e) = i} e - \sum\nolimits_{e \in E,\ t(e) = i} e) \\
& =
\sum\nolimits_{e \in E,\ s(e) = i} a_e(s(e) - t(e)) -
\sum\nolimits_{e \in E,\ t(e) = i} a_e(s(e) - t(e))
\end{align*}
Le coefficient de $i$ dans $\text{d}^*\text{d}(i)$ est
$$
\sum\nolimits_{e \in E,\ s(e) = i} a_e +
\sum\nolimits_{e \in E,\ t(e) = i} a_e = - a_{ii}
$$
car $\sum_j a_{ij} = 0$, et le coefficient de
$j \not = i$ dans $\text{d}^*\text{d}(i)$ est $-a_{ij}$.
Ainsi $\Coker(A) = \Coker(\text{d}^*\text{d})$.

\medskip\noindent
Considérons l'inclusion
$$
\Im(\text{d}) \oplus \Ker(\text{d}^*) \subset \mathbf{Z}(E)
$$
Le membre de gauche est une somme directe orthogonale. Il est clair que
$\mathbf{Z}(E)/\Ker(\text{d}^*)$ est sans torsion.
Nous affirmons que $\mathbf{Z}(E)/\Im(\text{d})$ est également sans torsion.
En effet, soient $x = \sum x_e e \in \mathbf{Z}(E)$ et $a > 1$ tels
que $ax = \text{d}y$ pour un certain $y = \sum y_i i \in \mathbf{Z}(V)$.
Alors $a x_e = y_{s(e)} - y_{t(e)}$. La propriété (2) entraîne
que tous les $y_i$ ont la même classe de congruence modulo $a$.
On peut donc écrire $y = a y' + (y_1, y_1, \ldots, y_1)$.
Puisque $\text{d}(y_1, y_1, \ldots, y_1) = 0$, on en déduit
que $x = \text{d}(y')$, ce qu'il fallait démontrer.

\medskip\noindent
On peut donc appliquer le lemme \ref{models-lemma-orthogonal-direct-sum}
pour obtenir des applications injectives
$$
\Im(\text{d})^\#/\Im(\text{d}) \leftarrow
\mathbf{Z}(E)/(\Im(\text{d}) \oplus \Ker(\text{d}^*)) \rightarrow
\Ker(\text{d}^*)^\#/\Ker(\text{d}^*)
$$
dont les conoyaux sont annulés par le produit des $a_e$
(qui est premier à $\ell$). Puisque $\Ker(\text{d}^*)$ est
un réseau de rang $1 - n + e$, il suffit, pour achever la démonstration,
de montrer qu'il existe un isomorphisme
$$
\Phi : M_{torsion} \longrightarrow \Im(\text{d})^\#/\Im(\text{d})
$$
Ceci est démontré dans le lemme \ref{models-lemma-coker}.
\end{proof}









\section{Types numériques}
\label{models-section-numerical-types}

\noindent
Une partie des arguments fera intervenir la combinatoire des données suivantes.

\begin{definition}
\label{models-definition-type}
Un {\it type numérique} $T$ est donné par
$$
n, m_i, a_{ij}, w_i, g_i
$$
où $n \geq 1$ est un entier et où $m_i$, $a_{ij}$, $w_i$, $g_i$
sont des entiers, pour $1 \leq i, j \leq n$, soumis aux conditions suivantes :
\begin{enumerate}
\item $m_i > 0$, $w_i > 0$, $g_i \geq 0$,
\item la matrice $A = (a_{ij})$ est symétrique et $a_{ij} \geq 0$
pour $i \not = j$,
\item il n'existe aucune partie propre non vide $I \subset \{1, \ldots, n\}$
telle que $a_{ij} = 0$ pour $i \in I$, $j \not \in I$,
\item pour tout $i$, on a $\sum_j a_{ij}m_j = 0$, et
\item $w_i | a_{ij}$.
\end{enumerate}
\end{definition}

\noindent
Cette sorte de structure est certes quelque peu malcommode à manipuler,
mais c'est exactement ce qui apparaît dans les fibres spéciales des modèles
propres réguliers des courbes lisses géométriquement connexes.
Naturellement, ces types ne nous intéressent qu'à permutation des indices près.

\begin{definition}
\label{models-definition-type-equivalent}
On dit que deux types numériques $n, m_i, a_{ij}, w_i, g_i$ et
$n', m'_i, a'_{ij}, w'_i, g'_i$ sont des {\it types équivalents} s'il
existe une permutation $\sigma$ de $\{1, \ldots, n\}$
telle que $m_i = m'_{\sigma(i)}$, $a_{ij} = a'_{\sigma(i)\sigma(j)}$,
$w_i = w'_{\sigma(i)}$ et $g_i = g'_{\sigma(i)}$.
\end{definition}

\noindent
Un type numérique possède un genre.

\begin{lemma}
\label{models-lemma-genus}
Soit $n, m_i, a_{ij}, w_i, g_i$ un type numérique. Alors l'expression
$$
g = 1 + \sum m_i(w_i(g_i - 1) - \frac{1}{2} a_{ii})
$$
est un entier.
\end{lemma}

\begin{proof}
Pour montrer que $g$ est un entier, il faut montrer que $\sum a_{ii}m_i$ est pair.
On le voit en calculant modulo $2$ comme suit :
\begin{align*}
\sum\nolimits_i a_{ii} m_i 
& \equiv
\sum\nolimits_{i,\ m_i\text{ impair}} a_{ii}m_i \\
& \equiv
\sum\nolimits_{i,\ m_i\text{ impair}} \sum\nolimits_{j \not = i} a_{ij}m_j \\
& \equiv
\sum\nolimits_{i,\ m_i\text{ impair}}
\sum\nolimits_{j \not = i,\ m_j\text{ impair}} a_{ij}m_j \\
& \equiv
\sum\nolimits_{i < j,\ m_i\text{ et }m_j\text{ impairs}} a_{ij}(m_i + m_j) \\
& \equiv
0
\end{align*}
où l'on a utilisé que $a_{ij} = a_{ji}$ et que $\sum_j a_{ij}m_j = 0$
pour tout $i$.
\end{proof}

\begin{definition}
\label{models-definition-genus}
On dit que $n, m_i, a_{ij}, w_i, g_i$ est un {\it type numérique de genre $g$}
si $g = 1 + \sum m_i(w_i(g_i - 1) - \frac{1}{2} a_{ii})$ est l'entier
du lemme \ref{models-lemma-genus}.
\end{definition}

\noindent
Nous démontrerons plus bas (lemme \ref{models-lemma-genus-nonnegative}) que le genre
est presque toujours $\geq 0$. Il peut toutefois exister des types numériques de
genre négatif.

\begin{lemma}
\label{models-lemma-irreducible}
Soit $n, m_i, a_{ij}, w_i, g_i$ un type numérique de genre $g$.
Si $n = 1$, alors $a_{11} = 0$ et $g = 1 + m_1w_1(g_1 - 1)$.
De plus, on peut classifier tous ces types numériques comme suit :
\begin{enumerate}
\item Si $g < 0$, alors $g_1 = 0$ et il n'y a qu'un nombre fini de
types numériques de genre $g$ avec $n = 1$, correspondant aux factorisations
$m_1w_1 = 1 - g$.
\item Si $g = 0$, alors $m_1 = 1$, $w_1 = 1$, $g_1 = 0$,
comme dans le lemme \ref{models-lemma-genus-zero}.
\item Si $g = 1$, on en déduit $g_1 = 1$, mais $m_1, w_1$ peuvent être des entiers
strictement positifs arbitraires ; c'est le cas
(\ref{models-item-one}) du lemme \ref{models-lemma-genus-one}.
\item Si $g > 1$, alors $g_1 > 1$ et il n'y a qu'un nombre fini de
types numériques de genre $g$ avec $n = 1$, correspondant aux
factorisations $m_1w_1(g_1 - 1) = g - 1$.
\end{enumerate}
\end{lemma}

\begin{proof}
Le lemme est immédiat.
\end{proof}

\begin{lemma}
\label{models-lemma-diagonal-negative}
Soit $n, m_i, a_{ij}, w_i, g_i$ un type numérique de genre $g$.
Si $n > 1$, alors $a_{ii} < 0$ pour tout $i$.
\end{lemma}

\begin{proof}
Le lemme \ref{models-lemma-recurring-symmetric-real} s'applique à la matrice $A$.
\end{proof}

\begin{lemma}
\label{models-lemma-minus-one}
Soit $n, m_i, a_{ij}, w_i, g_i$ un type numérique de genre $g$.
Supposons $n > 1$. Si $i$ est tel que la contribution
$m_i(w_i(g_i - 1) - \frac{1}{2} a_{ii})$
au genre $g$ soit $< 0$, alors $g_i = 0$ et $a_{ii} = -w_i$.
\end{lemma}

\begin{proof}
Cela résulte immédiatement du lemme \ref{models-lemma-diagonal-negative} et des relations
$w_i > 0$, $g_i \geq 0$ et $w_i | a_{ii}$.
\end{proof}

\begin{definition}
\label{models-definition-type-minus-one}
Soit $n, m_i, a_{ij}, w_i, g_i$ un type numérique.
On dit que $i$ est un {\it indice $(-1)$} si $g_i = 0$ et $a_{ii} = -w_i$.
\end{definition}

\noindent
On peut « contracter » les indices $(-1)$.

\begin{lemma}
\label{models-lemma-contract}
Soit $n, m_i, a_{ij}, w_i, g_i$ un type numérique $T$.
Supposons que $n$ soit un indice $(-1)$. Il existe alors un type
numérique $T'$ donné par $n', m'_i, a'_{ij}, w'_i, g'_i$, avec
\begin{enumerate}
\item $n' = n - 1$,
\item $m'_i = m_i$,
\item $a'_{ij} = a_{ij} - a_{in}a_{jn}/a_{nn}$,
\item $w'_i = w_i/2$ si $a_{in}/w_n$ est pair et $a_{in}/w_i$ impair,
et $w'_i = w_i$ sinon,
\item $g'_i =
\frac{w_i}{w'_i}(g_i - 1) + 1 + \frac{a_{in}^2 - w_na_{in}}{2w'_iw_n}$.
\end{enumerate}
De plus, on a $g = g'$.
\end{lemma}

\begin{proof}
Observons que $n > 1$, par exemple d'après le lemme \ref{models-lemma-irreducible},
et donc que $n' \geq 1$. Vérifions les conditions (1) -- (5) de la
définition \ref{models-definition-type} pour $n', m'_i, a'_{ij}, w'_i, g'_i$.

\medskip\noindent
La condition (1) est immédiate.

\medskip\noindent
Condition (2). La symétrie de $A' = (a'_{ij})$ est immédiate
et, puisque $a_{nn} < 0$ d'après le lemme \ref{models-lemma-diagonal-negative},
on voit que $a'_{ij} \geq a_{ij} \geq 0$ si $i \not = j$.

\medskip\noindent
Condition (3). Supposons que $I \subset \{1, \ldots, n - 1\}$ soit tel que
$a'_{ii'} = 0$ pour $i \in I$ et $i' \in \{1, \ldots, n - 1\} \setminus I$.
On voit alors que, pour tout $i \in I$ et tout $i' \in I'$, on a
$a_{in}a_{i'n} = 0$. Ainsi, ou bien $a_{in} = 0$ pour tout $i \in I$, et
$I \subset \{1, \ldots, n\}$ contredit la propriété (3) pour $T$,
ou bien $a_{i'n} = 0$ pour tout $i' \in \{1, \ldots, n - 1\} \setminus I$,
et $I \cup \{n\} \subset \{1, \ldots, n\}$ contredit
la propriété (3) de $T$. La condition (3) est donc vérifiée pour $T'$.

\medskip\noindent
Condition (4). On calcule
$$
\sum\nolimits_{j = 1}^{n - 1} a'_{ij}m_j =
\sum\nolimits_{j = 1}^{n - 1}
(a_{ij}m_j - \frac{a_{in}a_{jn}m_j}{a_{nn}}) =
- a_{in}m_n - \frac{a_{in}}{a_{nn}}(-a_{nn}m_n) = 0
$$
comme voulu.

\medskip\noindent
Condition (5). Il faut montrer que $w'_i$ divise $a_{in}a_{jn}/a_{nn}$.
C'est clair puisque $a_{nn} = -w_n$, $w_n | a_{jn}$ et $w_i | a_{in}$.

\medskip\noindent
Pour montrer que $g = g'$, écrivons d'abord
\begin{align*}
g
& =
1 + \sum\nolimits_{i = 1}^n m_i(w_i(g_i - 1) - \frac{1}{2}a_{ii}) \\
& =
1 + \sum\nolimits_{i = 1}^{n - 1} m_i(w_i(g_i - 1) - \frac{1}{2}a_{ii})
-\frac{1}{2}m_nw_n \\
& =
1 +  \sum\nolimits_{i = 1}^{n - 1}
m_i(w_i(g_i - 1) - \frac{1}{2}a_{ii} - \frac{1}{2}a_{in})
\end{align*}
En comparant avec l'expression de $g'$, on voit qu'il suffit que
$$
w'_i(g'_i - 1) - \frac{1}{2}a'_{ii} =
w_i(g_i - 1) - \frac{1}{2}a_{in} - \frac{1}{2}a_{ii}
$$
pour $i \leq n - 1$. Autrement dit, on a
$$
g'_i = \frac{2w_i(g_i - 1) - a_{in} - a_{ii} + a'_{ii} + 2w'_i}{2w'_i} =
\frac{w_i}{w'_i}(g_i - 1) + 1 + \frac{a_{in}^2 - w_na_{in}}{2w'_iw_n}
$$
On vérifie élémentairement qu'il s'agit d'un entier $\geq 0$
si l'on choisit $w'_i$ comme en (4).
\end{proof}

\begin{lemma}
\label{models-lemma-top-genus}
Soit $n, m_i, a_{ij}, w_i, g_i$ un type numérique.
Soit $e$ le nombre de couples $(i, j)$ tels que $i < j$ et $a_{ij} > 0$.
Alors l'expression $g_{top} = 1 - n + e$ est $\geq 0$.
\end{lemma}

\begin{proof}
Sinon, $e < n - 1$, ce qui signifie qu'il existe un $i$ tel que
$a_{ij} = 0$ pour tout $j \not = i$. Cela contredit l'hypothèse
(3) de la définition \ref{models-definition-type}.
\end{proof}

\begin{definition}
\label{models-definition-top-genus}
Soit $n, m_i, a_{ij}, w_i, g_i$ un type numérique $T$. Le
{\it genre topologique de $T$} est l'entier positif ou nul
$g_{top} = 1 - n + e$ du lemme \ref{models-lemma-top-genus}.
\end{definition}

\noindent
Nous voulons minorer le genre par le genre topologique. Toutefois, cela
ne sera pas toujours possible, par exemple pour les types numériques
avec $n = 1$ du lemme \ref{models-lemma-irreducible}. Ce sera cependant
le cas pour les types numériques minimaux, définis comme suit.

\begin{definition}
\label{models-definition-type-minimal}
On dit que le type numérique $n, m_i, a_{ij}, w_i, g_i$ de genre $g$
est {\it minimal} s'il n'existe aucun $i$
tel que $g_i = 0$ et $a_{ii} = -w_i$, autrement dit, s'il
n'existe aucun indice $(-1)$.
\end{definition}

\noindent
Nous démontrerons que le genre $g$ d'un type minimal avec $n > 1$
est supérieur ou égal à $\max(1, g_{top})$.

\begin{lemma}
\label{models-lemma-non-irreducible-minimal-type-genus-at-least-one}
Si $n, m_i, a_{ij}, w_i, g_i$ est un type numérique minimal
avec $n > 1$, alors $g \geq 1$.
\end{lemma}

\begin{proof}
Cela résulte de ce que $g = 1 + \sum \Phi_i$, avec
$\Phi_i = m_i(w_i(g_i - 1) - \frac{1}{2} a_{ii})$ positif ou nul,
d'après le lemme \ref{models-lemma-minus-one} et la définition des types minimaux.
\end{proof}

\begin{lemma}
\label{models-lemma-genus-nonnegative}
Si $n, m_i, a_{ij}, w_i, g_i$ est un type numérique minimal
avec $n > 1$, alors $g \geq g_{top}$.
\end{lemma}

\begin{proof}
Le lecteur qui ne s'intéresse qu'aux types numériques
associés aux modèles propres réguliers peut omettre cette démonstration, car nous
redémontrerons le résultat plus loin dans la situation géométrique.
On peut écrire
$$
g_{top} = 1 - n + \frac{1}{2}\sum\nolimits_{a_{ij} > 0} 1 =
1 + \sum\nolimits_i (-1 +
\frac{1}{2}\sum\nolimits_{j \not = i,\ a_{ij} > 0} 1) 
$$
D'autre part, on a
\begin{align*}
g & =
1 + \sum m_i(w_i(g_i - 1) - \frac{1}{2} a_{ii}) \\
& =
1 + \sum m_iw_ig_i - \sum m_iw_i +
\frac{1}{2} \sum\nolimits_{i \not = j} a_{ij}m_j \\
& =
1 + \sum\nolimits_i
m_iw_i(-1 + g_i + \frac{1}{2} \sum\nolimits_{j \not = i} \frac{a_{ij}}{w_i})
\end{align*}
La première égalité est la définition, la deuxième utilise que
$\sum a_{ij}m_j = 0$, et la dernière utilise
la relation $a_{ij} = a_{ji}$ ainsi que l'interversion de l'ordre
des sommations. En comparant avec la formule pour $g_{top}$, on conclut
que le lemme est vrai si
$$
\Psi_i =
m_iw_i(-1 + g_i + \frac{1}{2} \sum\nolimits_{j \not = i} \frac{a_{ij}}{w_i})
- (-1 + \frac{1}{2}\sum\nolimits_{j \not = i,\ a_{ij} > 0} 1)
$$
est $\geq 0$ pour tout $i$. Ce n'est cependant pas toujours le cas.
Examinons pour quels indices on peut avoir $\Psi_i < 0$.
Observons d'abord que
$$
(-1 + g_i + \frac{1}{2}\sum\nolimits_{j \not = i} \frac{a_{ij}}{w_i}) \geq
(-1 + \frac{1}{2}\sum\nolimits_{j \not = i,\ a_{ij} > 0} 1)
$$
car $a_{ij}/w_i$ est un entier positif ou nul. Puisque $m_iw_i$ est
un entier strictement positif, on en déduit que $\Psi_i \geq 0$ dès que
$m_iw_i = 1$ ou que le membre de gauche de l'inégalité est $\geq 0$,
ce qui se produit si $g_i > 0$, ou si $a_{ij} > 0$ pour au moins deux indices $j$, ou
s'il existe un $j$ tel que $a_{ij} > w_i$. Ainsi
$$
P = \{i : \Psi_i < 0\}
$$
est l'ensemble des indices $i$ tels que $m_iw_i > 1$, $g_i = 0$,
$a_{ij} > 0$ pour un unique $j$ et $a_{ij} = w_i$ pour ce $j$.
De plus,
$$
i \in P \Rightarrow \Psi_i = \frac{1}{2}(-m_iw_i + 1)
$$
La stratégie consiste à montrer que, pour $i \in P$, on peut prélever une part
de $\Psi_j$, où $j$ est le voisin de $i$, c'est-à-dire où $a_{ij} > 0$.
Toutefois, cela ne suffit pas tout à fait, car $j$ peut être un indice tel que $\Psi_j = 0$.

\medskip\noindent
Considérons l'ensemble
$$
Z = \{j : g_j = 0\text{ et }
j\text{ a exactement deux voisins }i, k\text{ avec }
a_{ij} = w_j = a_{jk}\}
$$
Pour $j \in Z$, on a $\Psi_j = 0$. Nous considérerons des suites
$M = (i, j_1, \ldots, j_s)$ où $s \geq 0$,
$i \in P$, $j_1, \ldots, j_s \in Z$, et
$a_{ij_1} > 0, a_{j_1j_2} > 0, \ldots, a_{j_{s - 1}j_s} > 0$.
Si notre type numérique est constitué de deux indices appartenant à $P$,
ou, plus généralement, si notre type numérique est constitué de
deux indices appartenant à $P$ et de tous les autres indices dans $Z$, alors
$g_{top} = 0$ et on conclut par le
lemme \ref{models-lemma-non-irreducible-minimal-type-genus-at-least-one}.
On peut donc écarter ces cas, ce que nous faisons.

\medskip\noindent
Soit $M = (i, j_1, \ldots, j_s)$ une suite maximale, et soit
$k$ le second voisin de $j_s$. (Si $s = 0$, alors $k$
est l'unique voisin de $i$.) Par maximalité, $k \not \in Z$
et, d'après ce qui précède, $k \not \in P$. Observons que
$w_i = a_{ij_1} = w_{j_1} = a_{j_1j_2} = \ldots = w_{j_s} =
a_{j_sk}$. En examinant la définition
d'un type numérique, on voit que
\begin{align*}
m_ia_{ii} + m_{j_1}w_i & = 0,\\
m_iw_i + m_{j_1}a_{j_1j_1} + m_{j_2}w_i & = 0,\\
\ldots & \ldots \\
m_{j_{s - 1}}w_i + m_{j_s}a_{j_sj_s} + m_kw_i & = 0
\end{align*}
La première égalité entraîne $m_{j_1} \geq 2m_i$, car le
type numérique est minimal. La deuxième égalité entraîne alors
$m_{j_2} \geq 3m_i$, et ainsi de suite. Dans tous les cas, on en déduit que
$m_k \geq 2m_i$ (y compris lorsque $s = 0$).

\medskip\noindent
Soit $k$ un indice pour lequel on dispose d'un $t > 0$ et de suites
maximales deux à deux distinctes $M_1, \ldots, M_t$ comme ci-dessus, avec
$M_b = (i_b, j_{b, 1}, \ldots, j_{b, s_b})$
telles que $k$ soit un voisin de $j_{b, s_b}$ pour $b = 1, \ldots, t$.
Nous montrerons que $\Phi_j + \sum_{b = 1, \ldots, t} \Phi_{i_b} \geq 0$.
Cela achèvera la démonstration du lemme, d'après ce qui précède.
Soit $M$ la réunion des indices figurant dans $M_b$, $b = 1, \ldots, t$.
Écrivons
$$
\Psi_k =
-\sum\nolimits_{b = 1, \ldots, t} \Psi_{i_b} + \Psi_k'
$$
où
\begin{align*}
\Psi_k' & =
m_kw_k\left(-1 + g_k +
\frac{1}{2} \sum\nolimits_{b = 1, \ldots t}
(\frac{a_{kj_{b, s_b}}}{w_k} - \frac{m_{i_b}w_{i_b}}{m_kw_k}) +
\frac{1}{2} \sum\nolimits_{l \not = k,\ l \not \in M}
\frac{a_{kl}}{w_k}
\right) \\
&
-\left(
-1 + \frac{1}{2}\sum\nolimits_{l \not = k,\ l \not \in M,\ a_{kl} > 0} 1
\right)
\end{align*}
Supposons $\Psi_k' < 0$ pour obtenir une contradiction.
Si l'ensemble $\{l : l \not = k,\ l \not \in M,\ a_{kl} > 0\}$ est vide,
alors $\{1, \ldots, n\} = M \cup \{k\}$ et $g_{top} = 0$,
car $e = n - 1$ dans ce cas, et le résultat découle du
lemme \ref{models-lemma-non-irreducible-minimal-type-genus-at-least-one}.
On peut donc supposer qu'il existe au moins un tel $l$, qui contribue
pour $(1/2)a_{kl}/w_k \geq 1/2$ à la somme entre les premières parenthèses.
Pour tout $b = 1, \ldots, t$, on a
$$
\frac{a_{kj_{b, s_b}}}{w_k} - \frac{m_{i_b}w_{i_b}}{m_kw_k} =
\frac{w_{i_b}}{w_k}(1 - \frac{m_{i_b}}{m_k})
$$
Cette expression est $\geq \frac{1}{2}$, car $m_k \geq 2m_{i_b}$
d'après le paragraphe précédent, et elle est $\geq 1$ si $w_k < w_{i_b}$.
Il s'ensuit que $\Psi_k' < 0$ entraîne $g_k = 0$.
Si $t \geq 2$, ou si $t = 1$ et $w_k < w_{i_1}$, alors $\Psi_k' \geq 0$
(on utilise ici l'existence d'un $l$ établie ci-dessus), ce qui
est encore contradictoire.
Ainsi $t = 1$ et $w_k = w_{i_1}$. S'il y a au moins deux termes non nuls
dans la somme sur $l$, ou s'il existe un tel $k$ et que $a_{kl} > w_k$, alors
on a également $\Psi_k' \geq 0$. La dernière possibilité est que $t = 1$ et
qu'il existe un $l$ tel que $a_{kl} = w_k$. Ceci est exclu, car cela
signifierait que $k \in Z$, contrairement à la maximalité de $M_1$.
\end{proof}

\begin{lemma}
\label{models-lemma-minus-two}
Soit $n, m_i, a_{ij}, w_i, g_i$ un type numérique de genre $g$.
Supposons $n > 1$. Si $i$ est tel que la contribution
$m_i(w_i(g_i - 1) - \frac{1}{2} a_{ii})$
au genre $g$ soit $0$, alors $g_i = 0$ et $a_{ii} = -2w_i$.
\end{lemma}

\begin{proof}
Cela résulte immédiatement du lemme \ref{models-lemma-diagonal-negative} et des relations
$w_i > 0$, $g_i \geq 0$ et $w_i | a_{ii}$.
\end{proof}

\noindent
Les indices vérifiant cette relation jouent en fait un
rôle important dans la structure des types numériques minimaux.
Nous leur donnons donc un nom.

\begin{definition}
\label{models-definition-type-minus-two}
Soit $n, m_i, a_{ij}, w_i, g_i$ un type numérique de genre $g$.
On dit que $i$ est un {\it indice $(-2)$} si $g_i = 0$ et $a_{ii} = -2w_i$.
\end{definition}

\noindent
Étant donné un type numérique minimal de genre $g$, les indices $(-2)$
sont exactement ceux dont la contribution au genre n'est pas strictement positive
dans la formule
$$
g = 1 + \sum m_i(w_i(g_i - 1) - \frac{1}{2} a_{ii})
$$
Il sera donc quelque peu délicat de borner les quantités associées
aux indices $(-2)$, comme nous le verrons plus loin.

\begin{remark}
\label{models-remark-genus-equality}
Soit $n, m_i, a_{ij}, w_i, g_i$ un type numérique minimal avec $n > 1$.
L'égalité $g = g_{top}$ peut avoir lieu dans le lemme \ref{models-lemma-genus-nonnegative}.
Par exemple, si $m_i = w_i = 1$ et $g_i = 0$ pour tout $i$, et si
$a_{ij} \in \{0, 1\}$ pour $i < j$.
\end{remark}


\section{Le groupe de Picard d'un type numérique}
\label{models-section-picard-group}

\noindent
Voici la définition.

\begin{definition}
\label{models-definition-picard-group}
Soit $n, m_i, a_{ij}, w_i, g_i$ un type numérique $T$. Le
{\it groupe de Picard de $T$} est le conoyau de la matrice
$(a_{ij}/w_i)$, plus précisément
$$
\Pic(T) =
\Coker\left(
\mathbf{Z}^{\oplus n} \to \mathbf{Z}^{\oplus n},\quad
e_i
\mapsto
\sum \frac{a_{ij}}{w_j}e_j
\right)
$$
où $e_i$ désigne le vecteur d'indice $i$ de la base canonique de $\mathbf{Z}^{\oplus n}$.
\end{definition}

\begin{lemma}
\label{models-lemma-picard-rank-1}
Soit $n, m_i, a_{ij}, w_i, g_i$ un type numérique $T$.
Le groupe de Picard de $T$ est un groupe abélien de type fini de rang $1$.
\end{lemma}

\begin{proof}
Si $n = 1$, alors $A = (a_{ij})$ est la matrice nulle et
le résultat est clair. Pour $n > 1$, la matrice $A$ est de rang
$n - 1$ d'après le lemme \ref{models-lemma-recurring-real} ou le
lemme \ref{models-lemma-recurring-symmetric-real}.
Bien entendu, le rang ne change pas lorsqu'on multiplie les lignes
par $1/w_i$. Ceci démontre le lemme.
\end{proof}

\begin{lemma}
\label{models-lemma-picard-T-and-A}
Soit $n, m_i, a_{ij}, w_i, g_i$ un type numérique $T$.
Alors $\Pic(T) \subset \Coker(A)$, où $A = (a_{ij})$.
\end{lemma}

\begin{proof}
Puisque $\Pic(T)$ est le conoyau de $(a_{ij}/w_i)$,
on dispose d'un diagramme commutatif
$$
\xymatrix{
0 \ar[r] &
\mathbf{Z}^{\oplus n} \ar[rr]_A & &
\mathbf{Z}^{\oplus n} \ar[rr] & &
\Coker(A) \ar[r] & 0 \\
0 \ar[r] &
\mathbf{Z}^{\oplus n} \ar[rr]^{(a_{ij}/w_i)} \ar[u]_{\text{id}} & &
\mathbf{Z}^{\oplus n} \ar[rr] \ar[u]_{\text{diag}(w_1, \ldots, w_n)} & &
\Pic(T) \ar[r] \ar[u] & 0
}
$$
dont les lignes sont exactes. Le lemme du serpent entraîne que
$\Pic(T) \subset \Coker(A)$.
\end{proof}

\begin{lemma}
\label{models-lemma-contract-picard-group}
Soit $n, m_i, a_{ij}, w_i, g_i$ un type numérique $T$.
Supposons que $n$ soit un indice $(-1)$. Soit $T'$ le type numérique
construit dans le lemme \ref{models-lemma-contract}. Il existe une
application injective
$$
\Pic(T) \to \Pic(T')
$$
dont le conoyau est un $2$-groupe abélien élémentaire.
\end{lemma}

\begin{proof}
Rappelons que $n' = n - 1$. Soit $e_i$, resp. $e'_i$, le vecteur d'indice $i$ de la
base de $\mathbf{Z}^{\oplus n}$, resp.\ $\mathbf{Z}^{\oplus n - 1}$.
Notons d'abord
$$
q : \mathbf{Z}^{\oplus n} \to \mathbf{Z}^{\oplus n - 1},
\quad e_n \mapsto 0\text{ et }e_i \mapsto e'_i\text{ pour }i \leq n - 1
$$
et posons
$$
p : \mathbf{Z}^{\oplus n} \to \mathbf{Z}^{\oplus n - 1},\quad
e_n \mapsto \sum\nolimits_{j = 1}^{n - 1} \frac{a_{nj}}{w'_j} e'_j
\text{ et }
e_i \mapsto \frac{w_i}{w'_i} e'_i\text{ pour }i \leq n - 1
$$
Un calcul (que nous omettons) montre que l'on a un diagramme commutatif
$$
\xymatrix{
\mathbf{Z}^{\oplus n} \ar[rr]_{(a_{ij}/w_i)} \ar[d]_q & &
\mathbf{Z}^{\oplus n} \ar[d]^p \\
\mathbf{Z}^{\oplus n'} \ar[rr]^{(a'_{ij}/w'_i)} & &
\mathbf{Z}^{\oplus n'}
}
$$
Le conoyau de la flèche supérieure étant
$\Pic(T)$ et celui de la flèche inférieure
étant $\Pic(T')$, on obtient l'homomorphisme voulu
entre les groupes de Picard. Puisque $\frac{w_i}{w'_i} \in \{1, 2\}$,
le conoyau de $\Pic(T) \to \Pic(T')$
est annulé par $2$ (car $2e'_i$ appartient à l'image de $p$
pour tout $i \leq n - 1$).
Montrons enfin que $\Pic(T) \to \Pic(T')$ est injective.
Soit $L = (l_1, \ldots, l_n)$ un représentant
d'un élément de $\Pic(T)$ d'image nulle dans $\Pic(T')$.
Puisque $q$ est surjective, une chasse au diagramme montre que l'on peut supposer
$L$ dans le noyau de $p$. Cela signifie que
$l_na_{ni}/w'_i + l_iw_i/w'_i = 0$, c'est-à-dire $l_i = - a_{ni}/w_i l_n$.
Ainsi $L$ est l'image de $-l_ne_n$ par l'application $(a_{ij}/w_j)$,
ce qui démontre le lemme.
\end{proof}

\begin{lemma}
\label{models-lemma-picard-group-genus-nonpositive}
Soit $n, m_i, a_{ij}, w_i, g_i$ un type numérique $T$.
Si le genre $g$ de $T$ est $\leq 0$, alors $\Pic(T) = \mathbf{Z}$.
\end{lemma}

\begin{proof}
On raisonne par récurrence sur $n$. Si $n = 1$, l'assertion est claire.
Si $n > 1$, alors $T$ n'est pas minimal d'après le
lemme \ref{models-lemma-non-irreducible-minimal-type-genus-at-least-one}.
En remplaçant $T$ par un type équivalent,
on peut supposer que $n$ est un indice $(-1)$.
D'après le lemme \ref{models-lemma-contract-picard-group},
on a $\Pic(T) \subset \Pic(T')$.
D'après le lemme \ref{models-lemma-contract}, le genre
de $T'$ est égal à celui de $T$, et l'on conclut par
récurrence.
\end{proof}








\section{Classification des sous-graphes propres}
\label{models-section-classify-proper-subgraphs}

\noindent
Dans cette section, on suppose donné un type numérique
$n, m_i, a_{ij}, w_i, g_i$ de genre $g$. Nous dresserons
une liste complète des « sous-graphes » possibles constitués entièrement
d'indices $(-2)$ (définition \ref{models-definition-type-minus-two}) et,
en même temps, nous classifierons tous les types numériques
minimaux possibles de genre $1$. Autrement dit, nous démontrons ici la
proposition \ref{models-proposition-classify-subgraphs} et le
lemme \ref{models-lemma-genus-one}

\medskip\noindent
Notre stratégie sera la suivante. Soit $n, m_i, a_{ij}, w_i, g_i$
un type numérique de genre $g$. Soit $I \subset \{1, \ldots, n\}$ une partie
constituée d'indices $(-2)$ telle qu'il n'existe aucune partie propre non vide
$J \subset I$ avec $a_{jj'} = 0$ pour $j \in J$, $j' \in I \setminus J$.
On raisonne par récurrence sur le cardinal $|I|$ de $I$. Si $I = \{i\}$
est constitué de $1$ indice, les seules contraintes sur $m_i$, $a_{ii}$
et $w_i$ sont $w_i | a_{ii}$, d'après la définition \ref{models-definition-type},
et $a_{ii} < 0$, d'après le lemme \ref{models-lemma-diagonal-negative},
ce qui fournira le cas initial. Dans l'étape de récurrence,
on applique d'abord l'hypothèse de récurrence aux parties $I' \subset I$
de cardinal $|I'| < |I|$. On obtient ainsi des contraintes sur les valeurs possibles
de $m_i, a_{ij}, w_i$, $i, j \in I$. En particulier, puisque $|I'| < |I| \leq n$,
il résultera de $\sum a_{ij}m_j = 0$ et du
lemme \ref{models-lemma-recurring-symmetric-real} que les
sous-matrices $(a_{ij})_{i, j \in I'}$ sont définies négatives
et que leur déterminant est de signe $(-1)^m$.
Pour chaque possibilité restante, on calcule le déterminant de
$(a_{ij})_{i, j \in I}$. Si ce déterminant est de signe $-(-1)^{|I|}$,
on peut écarter ce cas, car le théorème de Sylvester montre
que la matrice $(a_{ij})_{i, j \in I}$ n'est pas semi-définie négative.
Si le déterminant est de signe $(-1)^{|I|}$, alors $|I| < n$, et l'on
conclut (provisoirement) que ce cas peut se présenter comme sous-graphe propre
et on l'inscrit dans l'un des lemmes de cette section.
Si le déterminant est $0$, on doit avoir $|I| = n$
(toujours d'après le lemme \ref{models-lemma-recurring-symmetric-real}) et $g = 0$.
Dans ces cas, on détermine toutes les valeurs possibles de $m_i, a_{ij}, w_i$, $i, j \in I$,
et on les énumère dans le lemme \ref{models-lemma-genus-one}. Une fois l'argument
achevé, on obtient tous les types numériques minimaux possibles de genre
$1$ avec $n > 1$, car chacun d'eux est nécessairement constitué entièrement
d'indices $(-2)$ (et apparaîtra donc au cours de la récurrence),
d'après la formule du genre et les remarques de la section précédente.
En définitive, le lecteur pourra parcourir la liste des
possibilités du lemme \ref{models-lemma-genus-one} pour vérifier que toutes les
configurations de sous-graphes propres présentées ici comme possibles
se réalisent effectivement déjà pour des types numériques de genre $1$.

\medskip\noindent
Supposons que $i$ et $j$ soient des indices $(-2)$ tels que $a_{ij} > 0$.
La matrice $A = (a_{ij})$ étant semi-définie négative d'après le
lemme \ref{models-lemma-recurring-symmetric-real}, on voit que la matrice
$$
\left(
\begin{matrix}
-2w_i & a_{ij} \\
a_{ij} & -2w_j
\end{matrix}
\right)
$$
est définie négative, sauf si $n = 2$. Le cas $n = 2$ peut se produire :
le déterminant $4w_1w_2 - a_{12}^2$ est alors nul. En utilisant que
$\text{lcm}(w_1, w_2)$ divise $a_{12}$, le lecteur trouve aisément
que les seules possibilités sont
$$
(w_1, w_2, a_{12}) = (w, w, 2w), (w, 4w, 4w), \text{ ou }(4w, w, 4w)
$$
Observons que le cas $(4w, w, 4w)$ s'obtient à partir du cas
$(w, 4w, 4w)$ en échangeant les indices $i, j$.
Dans ces cas, $g = 1$. Cela conduit aux cas
(\ref{models-item-two-cycle}) et (\ref{models-item-up4}) du lemme \ref{models-lemma-genus-one}.
En supposant $n > 2$, on voit
que le déterminant $4w_iw_j - a_{ij}^2$ de la matrice ci-dessus
est $> 0$, et l'on en déduit que $a_{ij}^2/w_iw_j < 4$.
D'autre part, on sait que $\text{lcm}(w_i, w_j) | a_{ij}$,
donc $a_{ij}^2/w_iw_j$ est un entier.
Ainsi $a_{ij}^2/w_iw_j \in \{1, 2, 3\}$ et $w_i | w_j$ ou
inversement. Cela conduit aux possibilités suivantes :
$$
(w_1, w_2, a_{12}) = (w, w, w), (w, 2w, 2w), (w, 3w, 3w),
(2w, w, 2w), \text{ ou }(3w, w, 3w)
$$
Observons que le cas $(2w, w, 2w)$ s'obtient à partir du cas
$(w, 2w, 2w)$ en échangeant les indices $i, j$, et de même
pour les cas $(3w, w, 3w)$ et $(w, 3w, 3w)$. Les trois premières
solutions conduisent aux cas (\ref{models-item-A2}), (\ref{models-item-B2}) et
(\ref{models-item-G2}) du lemme \ref{models-lemma-two-by-two}. Dans ce lemme,
on a explicité les conséquences pour les entiers $m_i$ et $m_j$,
en utilisant que $\sum_l a_{kl}m_l = 0$ pour tout $k$ entraîne en particulier
$a_{ii}m_i + a_{ij}m_j \leq 0$ pour $k = i$ et
$a_{ij}m_i + a_{jj}m_j \leq 0$ pour $k = j$.

\begin{lemma}
\label{models-lemma-two-by-two}
Classification des sous-graphes propres de la forme
$$
\xymatrix{
\bullet \ar@{-}[r] & \bullet
}
$$
Si $n > 2$, pour une paire $i, j$ d'indices $(-2)$ telle que $a_{ij} > 0$,
on a, à permutation près, les possibilités suivantes pour les $m$, les $a$ et les $w$ :
\begin{enumerate}
\item
\label{models-item-A2}
ils sont donnés par
$$
\left(
\begin{matrix}
m_1 \\
m_2
\end{matrix}
\right),
\quad
\left(
\begin{matrix}
-2w & w \\
w & -2w
\end{matrix}
\right),
\quad
\left(
\begin{matrix}
w \\
w
\end{matrix}
\right)
$$
avec $w$ arbitraire, $2m_1 \geq m_2$ et $2m_2 \geq m_1$, ou
\item
\label{models-item-B2}
ils sont donnés par
$$
\left(
\begin{matrix}
m_1 \\
m_2
\end{matrix}
\right),
\quad
\left(
\begin{matrix}
-2w & 2w \\
2w & -4w
\end{matrix}
\right),
\quad
\left(
\begin{matrix}
w \\
2w
\end{matrix}
\right)
$$
avec $w$ arbitraire, $m_1 \geq m_2$ et $2m_2 \geq m_1$, ou
\item
\label{models-item-G2}
ils sont donnés par
$$
\left(
\begin{matrix}
m_1 \\
m_2
\end{matrix}
\right),
\quad
\left(
\begin{matrix}
-2w & 3w \\
3w & -6w
\end{matrix}
\right),
\quad
\left(
\begin{matrix}
w \\
3w
\end{matrix}
\right)
$$
avec $w$ arbitraire, $2m_1 \geq 3m_2$ et $2m_2 \geq m_1$.
\end{enumerate}
\end{lemma}

\begin{proof}
Voir la discussion ci-dessus.
\end{proof}

\noindent
Supposons que $i$, $j$ et $k$ soient trois indices $(-2)$ tels que $a_{ij} > 0$
et $a_{jk} > 0$. Autrement dit, l'indice $i$ « rencontre » $j$, et
$j$ « rencontre » $k$. Nous utiliserons sans le rappeler que chaque paire
$(i, j)$, $(i, k)$ et $(j, k)$ figure dans la liste du lemme \ref{models-lemma-two-by-two}.
Puisque la matrice $A = (a_{ij})$ est semi-définie
négative d'après le lemme \ref{models-lemma-recurring-symmetric-real}, la matrice
$$
\left(
\begin{matrix}
-2w_i & a_{ij} & a_{ik} \\
a_{ij} & -2w_j & a_{jk} \\
a_{ik} & a_{jk} & -2w_k
\end{matrix}
\right)
$$
est définie négative, sauf si $n = 3$. Le cas $n = 3$ peut se produire :
le déterminant\footnote{Il vaut
$-8w_iw_jw_k + 2a_{ij}^2w_k + 2a_{jk}^2w_i + 2a_{ik}^2w_j +
2a_{ij}a_{jk}a_{ik}$.} de la matrice est alors nul,
et l'on obtient l'égalité
$$
4 = \frac{a_{ij}^2}{w_iw_j} +
\frac{a_{jk}^2}{w_jw_k} +
\frac{a_{ik}^2}{w_iw_k} +
\frac{a_{ij}a_{ik}a_{jk}}{w_iw_jw_k}
$$
entre entiers. Le dernier terme du membre de droite de cette égalité est déterminé
par les autres, car
$$
\left(\frac{a_{ij}a_{ik}a_{jk}}{w_iw_jw_k}\right)^2 =
\frac{a_{ij}^2}{w_iw_j} \frac{a_{jk}^2}{w_jw_k} \frac{a_{ik}^2}{w_iw_k}
$$
Puisque nous avons vu ci-dessus que
$\frac{a_{ij}^2}{w_iw_j}, \frac{a_{jk}^2}{w_jw_k}$ appartiennent à
$\{1, 2, 3\}$ et que $\frac{a_{ik}^2}{w_iw_k}$ appartient à $\{0, 1, 2, 3\}$,
on en déduit que les seules possibilités sont
$$
(\frac{a_{ij}^2}{w_iw_j}, \frac{a_{jk}^2}{w_jw_k}, \frac{a_{ik}^2}{w_iw_k}) =
(1, 1, 1), (1, 3, 0), (2, 2, 0),\text{ ou } (3, 1, 0)
$$
Observons que le cas $(3, 1, 0)$ s'obtient à partir du cas $(1, 3, 0)$
en renversant l'ordre des indices $i, j, k$.
Dans chacun de ces cas, $g = 1$ ; le lecteur les retrouvera aux cas
(\ref{models-item-three-cycle}), (\ref{models-item-equal-up3}), (\ref{models-item-equal-down3}),
(\ref{models-item-up-up}), (\ref{models-item-up-down}), (\ref{models-item-down-up}) du
lemme \ref{models-lemma-genus-one},
avec un cas correspondant à $(1, 1, 1)$,
deux cas correspondant à $(1, 3, 0)$ et
trois cas correspondant à $(2, 2, 0)$.
En supposant $n > 3$,
on obtient l'inégalité
$$
4 > \frac{a_{ij}^2}{w_iw_j} + \frac{a_{ik}^2}{w_iw_k} +
\frac{a_{jk}^2}{w_jw_k} + \frac{a_{ij}a_{ik}a_{jk}}{w_iw_jw_k}
$$
entre entiers. En utilisant les restrictions sur les nombres données ci-dessus,
on voit que les seules possibilités sont
$$
(\frac{a_{ij}^2}{w_iw_j}, \frac{a_{jk}^2}{w_jw_k}, \frac{a_{ik}^2}{w_iw_k}) =
(1, 1, 0), (1, 2, 0),\text{ ou }(2, 1, 0)
$$
en particulier $a_{ik} = 0$ (rappelons que l'on suppose
$a_{ij} > 0$ et $a_{jk} > 0$). Observons que le cas
$(2, 1, 0)$ s'obtient à partir du cas $(1, 2, 0)$ en renversant
l'ordre des indices $i, j, k$. Les deux premières solutions conduisent
aux cas (\ref{models-item-A3}), (\ref{models-item-C3}) et (\ref{models-item-B3})
du lemme \ref{models-lemma-three-by-three}, où l'on a aussi explicité les conséquences
pour les entiers $m_i$, $m_j$ et $m_k$.

\begin{lemma}
\label{models-lemma-three-by-three}
Classification des sous-graphes propres de la forme
$$
\xymatrix{
\bullet \ar@{-}[r] &
\bullet \ar@{-}[r] &
\bullet
}
$$
Si $n > 3$, pour un triplet $i, j, k$ d'indices $(-2)$
dont au moins deux des $a_{ij}, a_{ik}, a_{jk}$ sont non nuls, on a,
à permutation près, les possibilités suivantes pour les $m$, les $a$ et les $w$ :
\begin{enumerate}
\item
\label{models-item-A3}
ils sont donnés par
$$
\left(
\begin{matrix}
m_1 \\
m_2 \\
m_3
\end{matrix}
\right),
\quad
\left(
\begin{matrix}
-2w & w & 0 \\
w & -2w & w \\
0 & w & -2w
\end{matrix}
\right),
\quad
\left(
\begin{matrix}
w \\
w \\
w
\end{matrix}
\right)
$$
avec $2m_1 \geq m_2$, $2m_2 \geq m_1 + m_3$, $2m_3 \geq m_2$, ou
\item
\label{models-item-C3}
ils sont donnés par
$$
\left(
\begin{matrix}
m_1 \\
m_2 \\
m_3
\end{matrix}
\right),
\quad
\left(
\begin{matrix}
-2w & w & 0 \\
w & -2w & 2w \\
0 & 2w & -4w
\end{matrix}
\right),
\quad
\left(
\begin{matrix}
w \\
w \\
2w
\end{matrix}
\right)
$$
avec $2m_1 \geq m_2$, $2m_2 \geq m_1 + 2m_3$, $2m_3 \geq m_2$, ou
\item
\label{models-item-B3}
ils sont donnés par
$$
\left(
\begin{matrix}
m_1 \\
m_2 \\
m_3
\end{matrix}
\right),
\quad
\left(
\begin{matrix}
-4w & 2w & 0 \\
2w & -4w & 2w \\
0 & 2w & -2w
\end{matrix}
\right),
\quad
\left(
\begin{matrix}
2w \\
2w \\
w
\end{matrix}
\right)
$$
avec $2m_1 \geq m_2$, $2m_2 \geq m_1 + m_3$, $m_3 \geq m_2$.
\end{enumerate}
\end{lemma}

\begin{proof}
Voir la discussion ci-dessus.
\end{proof}

\noindent
Supposons que $i$, $j$, $k$ et $l$ soient quatre indices $(-2)$ tels que
$a_{ij} > 0$, $a_{jk} > 0$ et $a_{kl} > 0$. Autrement dit,
l'indice $i$ « rencontre » $j$, $j$ « rencontre » $k$, et $k$ « rencontre » $l$.
Le lemme \ref{models-lemma-three-by-three} donne alors $a_{ik} = a_{jl} = 0$.
La matrice $A = (a_{ij})$ étant semi-définie négative, on voit que la
matrice
$$
\left(
\begin{matrix}
-2w_i & a_{ij} & 0 & a_{il} \\
a_{ij} & -2w_j & a_{jk} & 0 \\
0 & a_{jk} & -2w_k & a_{kl} \\
a_{il} & 0 & a_{kl} & -2w_l
\end{matrix}
\right)
$$
est définie négative, sauf si $n = 4$. Le cas $n = 4$ peut se produire :
le déterminant\footnote{Il vaut
$16w_iw_jw_kw_l - 4a_{ij}^2w_kw_l - 4a_{jk}^2w_iw_l - 4a_{kl}^2w_iw_j -
4a_{il}^2w_jw_k + a_{ij}^2a_{kl}^2 + a_{jk}^2a_{il}^2 -
2a_{ij}a_{il}a_{jk}a_{kl}$.} de la matrice est alors nul, et l'on obtient l'égalité
$$
16 +
\frac{a_{ij}^2}{w_iw_j}\frac{a_{kl}^2}{w_kw_l} +
\frac{a_{jk}^2}{w_jw_k}\frac{a_{il}^2}{w_iw_l} =
4\frac{a_{ij}^2}{w_iw_j} + 4\frac{a_{jk}^2}{w_jw_k} + 4\frac{a_{kl}^2}{w_kw_l}
+ 4\frac{a_{il}^2}{w_iw_l} + 2\frac{a_{ij}a_{il}a_{jk}a_{kl}}{w_iw_jw_kw_l}
$$
entre entiers positifs ou nuls. Le dernier terme du membre de droite de cette égalité est
déterminé par les autres, car
$$
\left(\frac{a_{ij}a_{il}a_{jk}a_{kl}}{w_iw_jw_kw_l}\right)^2 =
\frac{a_{ij}^2}{w_iw_j} \frac{a_{jk}^2}{w_jw_k}
\frac{a_{kl}^2}{w_kw_l} \frac{a_{il}^2}{w_iw_l}
$$
Puisque nous avons vu ci-dessus que
$\frac{a_{ij}^2}{w_iw_j}, \frac{a_{jk}^2}{w_jw_k}, \frac{a_{kl}^2}{w_kw_l}$
appartiennent à $\{1, 2\}$ et que $\frac{a_{il}^2}{w_iw_l}$ appartient à $\{0, 1, 2\}$,
on en déduit que les seules solutions possibles sont
$$
(\frac{a_{ij}^2}{w_iw_j}, \frac{a_{jk}^2}{w_jw_k}, \frac{a_{kl}^2}{w_kw_l},
\frac{a_{il}^2}{w_iw_l}) =
(1, 1, 1, 1) \text{ ou } (2, 1, 2, 0)
$$
et que dans ce cas $g = 1$ ; le lecteur les retrouvera aux cas
(\ref{models-item-four-cycle}), (\ref{models-item-up-equal-up}),
(\ref{models-item-up-equal-down}) et (\ref{models-item-down-equal-up})
du lemme \ref{models-lemma-genus-one}. En supposant $n > 4$,
on obtient l'inégalité
$$
16 +
\frac{a_{ij}^2}{w_iw_j}\frac{a_{kl}^2}{w_kw_l} +
\frac{a_{jk}^2}{w_jw_k}\frac{a_{il}^2}{w_iw_l} >
4\frac{a_{ij}^2}{w_iw_j} + 4\frac{a_{jk}^2}{w_jw_k} + 4\frac{a_{kl}^2}{w_kw_l}
+ 4\frac{a_{il}^2}{w_iw_l} + 2\frac{a_{ij}a_{il}a_{jk}a_{kl}}{w_iw_jw_kw_l}
$$
entre entiers positifs ou nuls. En utilisant les restrictions sur les nombres données ci-dessus,
on voit que les seules possibilités sont
$$
(\frac{a_{ij}^2}{w_iw_j}, \frac{a_{jk}^2}{w_jw_k}, \frac{a_{kl}^2}{w_kw_l},
\frac{a_{il}^2}{w_iw_l}) =
(1, 1, 1, 0), (1, 1, 2, 0), (1, 2, 1, 0), \text{ ou }(2, 1, 1, 0)
$$
en particulier $a_{il} = 0$ (rappelons que les trois autres étaient supposés
non nuls). Observons que le cas $(2, 1, 1, 0)$ s'obtient
à partir du cas $(1, 1, 2, 0)$ en renversant l'ordre
des indices $i, j, k, l$. Les trois premières solutions conduisent
aux cas (\ref{models-item-A4}), (\ref{models-item-C4}), (\ref{models-item-B4}) et
(\ref{models-item-F4}) du lemme \ref{models-lemma-four-by-four},
où l'on a aussi explicité les conséquences pour les entiers $m_i$, $m_j$, $m_k$
et $m_l$.

\begin{lemma}
\label{models-lemma-four-by-four}
Classification des sous-graphes propres de la forme
$$
\xymatrix{
\bullet \ar@{-}[r] &
\bullet \ar@{-}[r] &
\bullet \ar@{-}[r] &
\bullet
}
$$
Si $n > 4$, pour quatre indices $(-2)$ $i, j, k, l$
tels que $a_{ij}, a_{jk}, a_{kl}$ soient non nuls, on a,
à permutation près, les possibilités suivantes pour les $m$, les $a$ et les $w$ :
\begin{enumerate}
\item
\label{models-item-A4}
ils sont donnés par
$$
\left(
\begin{matrix}
m_1 \\
m_2 \\
m_3 \\
m_4
\end{matrix}
\right),
\quad
\left(
\begin{matrix}
-2w & w & 0 & 0 \\
w & -2w & w & 0 \\
0 & w & -2w & w \\
0 & 0 & w & -2w 
\end{matrix}
\right),
\quad
\left(
\begin{matrix}
w \\
w \\
w \\
w
\end{matrix}
\right)
$$
avec $2m_1 \geq m_2$, $2m_2 \geq m_1 + m_3$, $2m_3 \geq m_2 + m_4$,
et $2m_4 \geq m_3$, ou
\item
\label{models-item-C4}
ils sont donnés par
$$
\left(
\begin{matrix}
m_1 \\
m_2 \\
m_3 \\
m_4
\end{matrix}
\right),
\quad
\left(
\begin{matrix}
-2w & w & 0 & 0 \\
w & -2w & w & 0 \\
0 & w & -2w & 2w \\
0 & 0 & 2w & -4w 
\end{matrix}
\right),
\quad
\left(
\begin{matrix}
w \\
w \\
w \\
2w
\end{matrix}
\right)
$$
avec $2m_1 \geq m_2$, $2m_2 \geq m_1 + m_3$, $2m_3 \geq m_2 + 2m_4$,
et $2m_4 \geq m_3$, ou
\item
\label{models-item-B4}
ils sont donnés par
$$
\left(
\begin{matrix}
m_1 \\
m_2 \\
m_3 \\
m_4
\end{matrix}
\right),
\quad
\left(
\begin{matrix}
-4w & 2w & 0 & 0 \\
2w & -4w & 2w & 0 \\
0 & 2w & -4w & 2w \\
0 & 0 & 2w & -2w 
\end{matrix}
\right),
\quad
\left(
\begin{matrix}
2w \\
2w \\
2w \\
w
\end{matrix}
\right)
$$
avec $2m_1 \geq m_2$, $2m_2 \geq m_1 + m_3$, $2m_3 \geq m_2 + m_4$,
et $m_4 \geq m_3$, ou
\item
\label{models-item-F4}
ils sont donnés par
$$
\left(
\begin{matrix}
m_1 \\
m_2 \\
m_3 \\
m_4
\end{matrix}
\right),
\quad
\left(
\begin{matrix}
-2w & w & 0 & 0 \\
w & -2w & 2w & 0 \\
0 & 2w & -4w & 2w \\
0 & 0 & 2w & -4w 
\end{matrix}
\right),
\quad
\left(
\begin{matrix}
w \\
w \\
2w \\
2w
\end{matrix}
\right)
$$
avec $2m_1 \geq m_2$, $2m_2 \geq m_1 + 2m_3$, $2m_3 \geq m_2 + m_4$,
et $2m_4 \geq m_3$.
\end{enumerate}
\end{lemma}

\begin{proof}
Voir la discussion ci-dessus.
\end{proof}

\noindent
Supposons que $i$, $j$, $k$
et $l$ soient quatre indices $(-2)$ tels que $a_{ij} > 0$, $a_{ij} > 0$ et
$a_{il} > 0$. Autrement dit, l'indice $i$ « rencontre » les indices
$j$, $k$, $l$. Le lemme \ref{models-lemma-three-by-three} donne alors
$a_{jk} = a_{jl} = a_{kl} = 0$. Puisque la matrice $A = (a_{ij})$ est
semi-définie négative, on voit que la matrice
$$
\left(
\begin{matrix}
-2w_i & a_{ij} & a_{ik} & a_{il} \\
a_{ij} & -2w_j & 0 & 0 \\
a_{ik} & 0 & -2w_k & 0 \\
a_{il} & 0 & 0 & -2w_l
\end{matrix}
\right)
$$
est définie négative, sauf si $n = 4$. Le cas $n = 4$ peut se produire :
le déterminant\footnote{Il vaut
$16w_iw_jw_kw_l - 4a_{ij}^2w_kw_l - 4a_{ik}^2w_jw_l - 4a_{il}^2w_jw_k$.}
de la matrice est alors nul, et l'on obtient l'égalité
$$
4 = \frac{a_{ij}^2}{w_iw_j} + \frac{a_{ik}^2}{w_iw_k} + \frac{a_{il}^2}{w_jw_l}
$$
entre entiers positifs ou nuls. Puisque nous avons vu ci-dessus que
$\frac{a_{ij}^2}{w_iw_j}, \frac{a_{ik}^2}{w_iw_k}, \frac{a_{il}^2}{w_iw_l}$
appartiennent à $\{1, 2\}$, on en déduit que la seule possibilité
est, à permutation près : $4 = 1 + 1 + 2$. Dans chacun de ces
cas, $g = 1$ ; le lecteur les retrouvera aux cas
(\ref{models-item-triple-with-up}) et (\ref{models-item-triple-with-down}) du
lemme \ref{models-lemma-genus-one}. En supposant $n > 4$,
on obtient l'inégalité
$$
4 > \frac{a_{ij}^2}{w_iw_j} + \frac{a_{ik}^2}{w_iw_k} + \frac{a_{il}^2}{w_jw_l}
$$
entre entiers positifs ou nuls. Il en résulte que $\frac{a_{ij}^2}{w_iw_j} =
\frac{a_{ik}^2}{w_iw_k} = \frac{a_{il}^2}{w_jw_l} = 1$
et que $w_i = w_j = w_k = w_l$. Cela conduit au cas
(\ref{models-item-D4}) du lemme \ref{models-lemma-D4},
où l'on a aussi explicité les conséquences pour les entiers $m_i$, $m_j$, $m_k$
et $m_l$.

\begin{lemma}
\label{models-lemma-D4}
Classification des sous-graphes propres de la forme
$$
\xymatrix{
\bullet \ar@{-}[r] & \bullet \ar@{-}[r] \ar@{-}[d] & \bullet \\
& \bullet
}
$$
Si $n > 4$, pour quatre indices $(-2)$ $i, j, k, l$
tels que $a_{ij}, a_{ik}, a_{il}$ soient non nuls, on a,
à permutation près, les possibilités suivantes pour les $m$, les $a$ et les $w$ :
\begin{enumerate}
\item
\label{models-item-D4}
ils sont donnés par
$$
\left(
\begin{matrix}
m_1 \\
m_2 \\
m_3 \\
m_4
\end{matrix}
\right),
\quad
\left(
\begin{matrix}
-2w & w & w & w \\
w & -2w & 0 & 0 \\
w & 0 & -2w & 0 \\
w & 0 & 0 & -2w 
\end{matrix}
\right),
\quad
\left(
\begin{matrix}
w \\
w \\
w \\
w
\end{matrix}
\right)
$$
avec $2m_1 \geq m_2 + m_3 + m_4$, $2m_2 \geq m_1$, $2m_3 \geq m_1$,
$2m_4 \geq m_1$. Observons que cela entraîne $m_1 \geq \max(m_2, m_3, m_4)$.
\end{enumerate}
\end{lemma}

\begin{proof}
Voir la discussion ci-dessus.
\end{proof}

\noindent
Supposons que $h$, $i$, $j$, $k$ et $l$ soient cinq indices $(-2)$
tels que $a_{hi} > 0$, $a_{ij} > 0$, $a_{jk} > 0$ et $a_{kl} > 0$.
Autrement dit, l'indice $h$ « rencontre » $i$, $i$ « rencontre » $j$,
$j$ « rencontre » $k$ et $k$ « rencontre » $l$.
On peut alors appliquer les lemmes \ref{models-lemma-three-by-three} et
\ref{models-lemma-four-by-four} pour voir que
$a_{hj} = a_{hk} = a_{ik} = a_{il} = a_{jl} = 0$
et que les fractions
$\frac{a_{hi}^2}{w_hw_i}, \frac{a_{ij}^2}{w_iw_j}, \frac{a_{jk}^2}{w_jw_k},
\frac{a_{kl}^2}{w_kw_l}$ appartiennent à $\{1, 2\}$, tandis que la fraction
$\frac{a_{hl}^2}{w_hw_l} \in \{0, 1, 2\}$.
Puisque la matrice $A = (a_{ij})$ est semi-définie négative, on voit que la
matrice
$$
\left(
\begin{matrix}
-2w_h & a_{hi} & 0 & 0 & a_{hl} \\
a_{hi} & -2w_i & a_{ij} & 0 & 0 \\
0 & a_{ij} & -2w_j & a_{jk} & 0 \\
0 & 0 & a_{jk} & -2w_k & a_{kl} \\
a_{hl} & 0 & 0 & a_{kl} & -2w_l
\end{matrix}
\right)
$$
est définie négative, sauf si $n = 5$. Le cas $n = 5$ peut se produire :
le déterminant\footnote{Il vaut
$-32w_hw_iw_jw_kw_l +
8a_{hi}^2w_jw_kw_l +
8a_{ij}^2w_hw_kw_l +
8a_{jk}^2w_hw_iw_l +
8a_{kl}^2w_hw_iw_j +
8a_{hl}^2w_iw_jw_k -
2a_{hi}^2a_{jk}^2w_l -
2a_{hi}^2a_{kl}^2w_j -
2a_{ij}^2a_{kl}^2w_h -
2a_{hl}^2a_{ij}^2w_k -
2a_{hl}^2a_{jk}^2w_i +
2a_{hi}a_{ij}a_{jk}a_{kl}a_{hl}
$
.}
de la matrice est nul et l'on obtient l'égalité
\begin{align*}
16 + 
\frac{a_{hi}^2}{w_hw_i}\frac{a_{jk}^2}{w_jw_k} +
\frac{a_{hi}^2}{w_hw_i}\frac{a_{kl}^2}{w_kw_l} +
\frac{a_{ij}^2}{w_iw_j}\frac{a_{kl}^2}{w_kw_l} +
\frac{a_{hl}^2}{w_hw_l}\frac{a_{ij}^2}{w_iw_j} +
\frac{a_{hl}^2}{w_hw_l}\frac{a_{jk}^2}{w_jw_k} \\
= 4\frac{a_{hi}^2}{w_hw_i} +
4\frac{a_{ij}^2}{w_iw_j} +
4\frac{a_{jk}^2}{w_jw_k} +
4\frac{a_{kl}^2}{w_kw_l} +
4\frac{a_{hl}^2}{w_hw_l} +
\frac{a_{hi}a_{ij}a_{jk}a_{kl}a_{hl}}{w_hw_iw_jw_kw_l}
\end{align*}
entre entiers positifs ou nuls. Le dernier terme du membre de droite de cette égalité
est déterminé par les autres, car
$$
\left(\frac{a_{hi}a_{ij}a_{jk}a_{kl}a_{hl}}{w_hw_iw_jw_kw_l} \right)^2 =
\frac{a_{hi}^2}{w_hw_i} \frac{a_{ij}^2}{w_iw_j}
\frac{a_{jk}^2}{w_jw_k} \frac{a_{kl}^2}{w_kw_l} \frac{a_{hl}^2}{w_hw_l}
$$
On conclut que les seules solutions possibles sont
$$
(\frac{a_{hi}^2}{w_hw_i}, \frac{a_{ij}^2}{w_iw_j},
\frac{a_{jk}^2}{w_jw_k}, \frac{a_{kl}^2}{w_kw_l}, \frac{a_{hl}^2}{w_hw_l}) =
(1, 1, 1, 1, 1), (1, 1, 2, 1, 0), (1, 2, 1, 1, 0), \text{ ou }(2, 1, 1, 2, 0)
$$
Observons que le cas $(1, 2, 1, 1, 0)$ s'obtient à partir du cas
$(1, 1, 2, 1, 0)$ en renversant l'ordre des indices $h, i, j, k, l$.
Dans ces cas, $g = 1$ ; le lecteur les retrouvera aux cas
(\ref{models-item-five-cycle}),
(\ref{models-item-equal-equal-up-equal}), (\ref{models-item-equal-equal-down-equal}),
(\ref{models-item-up-equal-equal-up}), (\ref{models-item-up-equal-equal-down}) et
(\ref{models-item-down-equal-equal-up}) du lemme \ref{models-lemma-genus-one},
avec un cas correspondant à $(1, 1, 1, 1, 1)$,
deux cas correspondant à $(1, 1, 2, 1, 0)$ et
trois cas correspondant à $(2, 1, 1, 2, 0)$.
En supposant $n > 5$, on obtient l'inégalité
\begin{align*}
16 + 
\frac{a_{hi}^2}{w_hw_i}\frac{a_{jk}^2}{w_jw_k} +
\frac{a_{hi}^2}{w_hw_i}\frac{a_{kl}^2}{w_kw_l} +
\frac{a_{ij}^2}{w_iw_j}\frac{a_{kl}^2}{w_kw_l} +
\frac{a_{hl}^2}{w_hw_l}\frac{a_{ij}^2}{w_iw_j} +
\frac{a_{hl}^2}{w_hw_l}\frac{a_{jk}^2}{w_jw_k} \\
> 4\frac{a_{hi}^2}{w_hw_i} +
4\frac{a_{ij}^2}{w_iw_j} +
4\frac{a_{jk}^2}{w_jw_k} +
4\frac{a_{kl}^2}{w_kw_l} +
4\frac{a_{hl}^2}{w_hw_l} +
\frac{a_{hi}a_{ij}a_{jk}a_{kl}a_{hl}}{w_hw_iw_jw_kw_l}
\end{align*}
entre entiers positifs ou nuls. En utilisant les restrictions données ci-dessus
sur ces nombres, on voit que les seules possibilités sont
$$
(\frac{a_{hi}^2}{w_hw_i}, \frac{a_{ij}^2}{w_iw_j},
\frac{a_{jk}^2}{w_jw_k}, \frac{a_{kl}^2}{w_kw_l}, \frac{a_{hl}^2}{w_hw_l}) =
(1, 1, 1, 1, 0), (1, 1, 1, 2, 0), \text{ ou } (2, 1, 1, 1, 0)
$$
en particulier, $a_{hl} = 0$ (rappelons que nous avons supposé les quatre autres
non nuls). Observons que le cas $(1, 1, 1, 2, 0)$
s'obtient à partir du cas $(2, 1, 1, 1, 0)$ en renversant l'ordre
des indices $h, i, j, k, l$. Les deux premières solutions conduisent aux
cas (\ref{models-item-A5}), (\ref{models-item-C5}) et (\ref{models-item-B5}) du
lemme \ref{models-lemma-five-by-five}, où l'on a aussi explicité les
conséquences pour les entiers $m_h$, $m_i$, $m_j$, $m_k$ et $m_l$.

\begin{lemma}
\label{models-lemma-five-by-five}
Classification des sous-graphes propres de la forme
$$
\xymatrix{
\bullet \ar@{-}[r] &
\bullet \ar@{-}[r] &
\bullet \ar@{-}[r] &
\bullet \ar@{-}[r] &
\bullet
}
$$
Si $n > 5$, pour cinq indices $(-2)$ $h, i, j, k, l$
tels que $a_{hi}, a_{ij}, a_{jk}, a_{kl}$ soient non nuls, on a,
à permutation près, les possibilités suivantes pour les $m$, les $a$ et les $w$ :
\begin{enumerate}
\item
\label{models-item-A5}
ils sont donnés par
$$
\left(
\begin{matrix}
m_1 \\
m_2 \\
m_3 \\
m_4 \\
m_5
\end{matrix}
\right),
\quad
\left(
\begin{matrix}
-2w & w & 0 & 0 & 0 \\
w & -2w & w & 0 & 0 \\
0 & w & -2w & w & 0 \\
0 & 0 & w & -2w & w \\
0 & 0 & 0 & w & -2w
\end{matrix}
\right),
\quad
\left(
\begin{matrix}
w \\
w \\
w \\
w \\
w
\end{matrix}
\right)
$$
avec $2m_1 \geq m_2$, $2m_2 \geq m_1 + m_3$, $2m_3 \geq m_2 + m_4$,
$2m_4 \geq m_3 + m_5$ et $2m_5 \geq m_4$, ou
\item
\label{models-item-C5}
ils sont donnés par
$$
\left(
\begin{matrix}
m_1 \\
m_2 \\
m_3 \\
m_4 \\
m_5
\end{matrix}
\right),
\quad
\left(
\begin{matrix}
-2w & w & 0 & 0 & 0 \\
w & -2w & w & 0 & 0 \\
0 & w & -2w & w & 0 \\
0 & 0 & w & -2w & 2w \\
0 & 0 & 0 & 2w & -4w
\end{matrix}
\right),
\quad
\left(
\begin{matrix}
w \\
w \\
w \\
w \\
2w
\end{matrix}
\right)
$$
avec $2m_1 \geq m_2$, $2m_2 \geq m_1 + m_3$, $2m_3 \geq m_2 + 2m_4$,
$2m_4 \geq m_3 + m_5$ et $2m_5 \geq m_4$, ou
\item
\label{models-item-B5}
ils sont donnés par
$$
\left(
\begin{matrix}
m_1 \\
m_2 \\
m_3 \\
m_4 \\
m_5
\end{matrix}
\right),
\quad
\left(
\begin{matrix}
-4w & 2w & 0 & 0 & 0 \\
2w & -4w & 2w & 0 & 0 \\
0 & 2w & -4w & 2w & 0 \\
0 & 0 & 2w & -4w & 2w \\
0 & 0 & 0 & 2w & -2w
\end{matrix}
\right),
\quad
\left(
\begin{matrix}
2w \\
2w \\
2w \\
2w \\
w
\end{matrix}
\right)
$$
avec $2m_1 \geq m_2$, $2m_2 \geq m_1 + m_3$, $2m_3 \geq m_2 + m_4$,
$2m_4 \geq m_3 + m_5$ et $m_4 \geq m_3$.
\end{enumerate}
\end{lemma}

\begin{proof}
Voir la discussion ci-dessus.
\end{proof}

\noindent
Supposons que $h$, $i$, $j$, $k$ et $l$ soient cinq indices $(-2)$
tels que $a_{hi} > 0$, $a_{hj} > 0$, $a_{hk} > 0$ et $a_{hl} > 0$.
Autrement dit, l'indice $h$ « rencontre » les indices
$i$, $j$, $k$, $l$. Le lemme \ref{models-lemma-three-by-three} montre alors
que $a_{ij} = a_{ik} = a_{il} = a_{jk} = a_{jl} = a_{kl} = 0$ et
le lemme \ref{models-lemma-D4} que
$w_h = w_i = w_j = w_k = w_l = w$ pour un entier $w > 0$ et
$a_{hi} = a_{hj} = a_{hk} = a_{hl} = -2w$.
La matrice correspondante
$$
\left(
\begin{matrix}
-2w & w & w & w & w \\
w & -2w & 0 & 0 & 0 \\
w & 0 & -2w & 0 & 0 \\
w & 0 & 0 & -2w & 0 \\
w & 0 & 0 & 0 & -2w
\end{matrix}
\right)
$$
est singulière. Cela ne peut donc se produire que si $n = 5$ et $g = 1$.
Le lecteur retrouvera cette possibilité au cas
(\ref{models-item-quadruple}) du lemme \ref{models-lemma-genus-one}.

\begin{lemma}
\label{models-lemma-fourfold}
Non-existence de sous-graphes propres de la forme
$$
\xymatrix{
\bullet \ar@{-}[r] & \bullet \ar@{-}[ld] \ar@{-}[r] \ar@{-}[d] & \bullet \\
\bullet & \bullet
}
$$
Si $n > 5$, il n'existe {\bf pas} cinq indices $(-2)$
$h$, $i$, $j$, $k$ tels que $a_{hi} > 0$, $a_{hj} > 0$, $a_{hk} > 0$ et
$a_{hl} > 0$.
\end{lemma}

\begin{proof}
Voir la discussion ci-dessus.
\end{proof}

\noindent
Supposons que $h$, $i$, $j$, $k$ et $l$ soient cinq indices $(-2)$
tels que $a_{hi} > 0$, $a_{ij} > 0$, $a_{jk} > 0$ et $a_{jl} > 0$.
Autrement dit, l'indice $h$ « rencontre » $i$, et l'indice $j$ « rencontre »
les indices $i$, $k$, $l$. Le
lemme \ref{models-lemma-D4} montre alors que $a_{ik} = a_{il} = a_{kl} = 0$,
$w_i = w_j = w_k = w_l = w$ et $a_{ij} = a_{jk} = a_{jl} = w$
pour un entier $w$. En appliquant le lemme \ref{models-lemma-four-by-four} aux
deux quadruplets
$h, i, j, k$ et $h, i, j, l$, on voit que $a_{hj} = a_{hk} = a_{hl} = 0$,
que $w_h = \frac{1}{2}w$, $w$ ou $2w$, et que
respectivement $a_{hi} = w$, $w$ ou $2w$.
Puisque $A$ est semi-définie négative, on voit que la matrice
$$
\left(
\begin{matrix}
-2w_h & a_{hi} & 0 & 0 & 0 \\
a_{hi} & -2w & w & 0 & 0 \\
0 & w & -2w & w & w \\
0 & 0 & w & -2w & 0 \\
0 & 0 & w & 0 & -2w
\end{matrix}
\right)
$$
est définie négative, sauf si $n = 5$. Le lecteur vérifiera que le
déterminant de la matrice vaut $0$ lorsque $w_h = \frac{1}{2}w$ ou $2w$.
Cela conduit aux cas (\ref{models-item-triple-extended-up}) et
(\ref{models-item-triple-extended-down}) du lemme \ref{models-lemma-genus-one}.
Pour $w_h = w$, on obtient le cas (\ref{models-item-D5}) du
lemme \ref{models-lemma-D5}.

\begin{lemma}
\label{models-lemma-D5}
Classification des sous-graphes propres de la forme
$$
\xymatrix{
\bullet \ar@{-}[r] & \bullet \ar@{-}[r] &
\bullet \ar@{-}[r] \ar@{-}[d] & \bullet \\
& & \bullet
}
$$
Si $n > 5$, pour cinq indices $(-2)$ $h, i, j, k, l$
tels que $a_{hi}, a_{ij}, a_{jk}, a_{jl}$ soient non nuls, on a,
à permutation près, les possibilités suivantes pour les $m$, les $a$ et les $w$ :
\begin{enumerate}
\item
\label{models-item-D5}
ils sont donnés par
$$
\left(
\begin{matrix}
m_1 \\
m_2 \\
m_3 \\
m_4 \\
m_5
\end{matrix}
\right),
\quad
\left(
\begin{matrix}
-2w & w & 0 & 0 & 0 \\
w & -2w & w & 0 & 0 \\
0 & w & -2w & w & w \\
0 & 0 & w & -2w & 0 \\
0 & 0 & w & 0 & -2w
\end{matrix}
\right),
\quad
\left(
\begin{matrix}
w \\
w \\
w \\
w \\
w
\end{matrix}
\right)
$$
avec $2m_1 \geq m_2$, $2m_2 \geq m_1 + m_3$, $2m_3 \geq m_2 + m_4 + m_5$,
$2m_4 \geq m_3$ et $2m_5 \geq m_3$.
\end{enumerate}
\end{lemma}

\begin{proof}
Voir la discussion ci-dessus.
\end{proof}

\noindent
Supposons que $t > 5$ et que $i_1, \ldots, i_t$ soient $t$ indices $(-2)$ distincts
tels que $a_{i_ji_{j + 1}}$ soit non nul pour $j = 1, \ldots, t - 1$. Nous allons
montrer par récurrence sur $t$ que, si $n = t$, cela conduit aux possibilités
(\ref{models-item-n-cycle}), (\ref{models-item-up-chain-equal-up}),
(\ref{models-item-up-chain-equal-down}), (\ref{models-item-down-chain-equal-up})
du lemme \ref{models-lemma-genus-one}, et que, si $n > t$, cela conduit aux cas
(\ref{models-item-An}), (\ref{models-item-Cn}) et (\ref{models-item-Bn}) du
lemme \ref{models-lemma-long}.
Premièrement, si $a_{i_1i_t}$ est non nul, il résulte clairement
du lemme \ref{models-lemma-five-by-five} que
$w_{i_1} = \ldots = w_{i_t} = w$, que
$a_{i_ji_{j + 1}} = w$ pour $j = 1, \ldots, t - 1$ et que
$a_{i_1i_t} = w$. Le vecteur $(1, \ldots, 1)$ appartient alors au noyau
de la matrice $t \times t$ correspondante. On doit donc avoir $n = t$ ;
on voit que le genre vaut $1$ et que l'on se trouve
dans le cas (\ref{models-item-n-cycle}) du lemme \ref{models-lemma-genus-one}.
On peut donc supposer $a_{i_1i_t} = 0$.
L'hypothèse de récurrence (ou le lemme \ref{models-lemma-five-by-five} si $t = 6$)
montre que $a_{i_ji_k} = 0$ si $k > j + 1$.
De plus, on a $w_{i_1} = \ldots = w_{i_{t - 1}} = w$
pour un entier $w$, et $w_{i_1}, w_{i_t} \in \{\frac{1}{2}w, w, 2w\}$.
En outre, si $w_{i_1}$, respectivement $w_{i_t}$, vaut
$\frac{1}{2}w$, $w$ ou $2w$, alors
$a_{i_1i_2}$, respectivement $a_{i_{t - 1}i_t}$,
vaut $w$, $w$ ou $2w$. Cela donne $9$ possibilités.
Dans chaque cas, il est facile de déterminer ce qui se passe :
\begin{enumerate}
\item si $(w_{i_1}, w_{i_t}) = (\frac{1}{2}w, \frac{1}{2}w)$, alors
on se trouve dans le cas (\ref{models-item-up-chain-equal-down}) du
lemme \ref{models-lemma-genus-one} ;
\item si $(w_{i_1}, w_{i_t}) = (\frac{1}{2}w, w)$ ou $(w, \frac{1}{2}w)$,
alors on se trouve dans le cas (\ref{models-item-Bn}) du lemme \ref{models-lemma-long} ;
\item si $(w_{i_1}, w_{i_t}) = (\frac{1}{2}w, 2w)$ ou $(2w, \frac{1}{2}w)$,
alors on se trouve dans le cas (\ref{models-item-up-chain-equal-up}) du
lemme \ref{models-lemma-genus-one} ;
\item si $(w_{i_1}, w_{i_t}) = (w, w)$, alors on se trouve dans le cas
(\ref{models-item-An}) du lemme \ref{models-lemma-long} ;
\item si $(w_{i_1}, w_{i_t}) = (w, 2w)$ ou $(2w, w)$, alors on se trouve
dans le cas (\ref{models-item-Cn}) du lemme \ref{models-lemma-long} ;
\item enfin, si $(w_{i_1}, w_{i_t}) = (2w, 2w)$, alors on se trouve dans le cas
(\ref{models-item-down-chain-equal-up}) du lemme \ref{models-lemma-genus-one}.
\end{enumerate}

\begin{lemma}
\label{models-lemma-long}
Classification des sous-graphes propres de la forme
$$
\xymatrix{
\bullet \ar@{-}[r] &
\bullet \ar@{-}[r] &
\bullet \ar@{..}[r] &
\bullet \ar@{-}[r] &
\bullet \ar@{-}[r] &
\bullet
}
$$
Soient $t > 5$ et $n > t$. Alors, pour $t$ indices $(-2)$ distincts
$i_1, \ldots, i_t$ tels que $a_{i_ji_{j + 1}}$ soit non nul pour
$j = 1, \ldots, t - 1$, on a, à renversement près de l'ordre de ces indices,
les possibilités suivantes pour les $a$ et les $w$ :
\begin{enumerate}
\item
\label{models-item-An}
ils sont donnés par $w_{i_1} = w_{i_2} = \ldots = w_{i_t} = w$,
$a_{i_ji_{j + 1}} = w$ et $a_{i_ji_k} = 0$ si $k > j + 1$, ou
\item
\label{models-item-Cn}
ils sont donnés par $w_{i_1} = w_{i_2} = \ldots = w_{i_{t - 1}} = w$,
$w_{j_t} = 2w$, $a_{i_ji_{j + 1}} = w$ pour $j < t - 1$,
$a_{i_{t - 1}i_t} = 2w$ et $a_{i_ji_k} = 0$ si $k > j + 1$, ou
\item
\label{models-item-Bn}
ils sont donnés par $w_{i_1} = w_{i_2} = \ldots = w_{i_{t - 1}} = 2w$,
$w_{j_t} = w$, $a_{i_ji_{j + 1}} = 2w$ et
$a_{i_{t - 1}i_t} = 2w$, et $a_{i_ji_k} = 0$ si $k > j + 1$.
\end{enumerate}
\end{lemma}

\begin{proof}
Voir la discussion ci-dessus.
\end{proof}

\noindent
Supposons que $t > 4$ et que $i_1, \ldots, i_{t + 1}$ soient $t + 1$
indices $(-2)$ distincts tels que $a_{i_ji_{j + 1}} > 0$
pour $j = 1, \ldots, t - 1$ et tels que $a_{j_{t - 1}j_{t + 1}} > 0$.
Voir la figure du lemme \ref{models-lemma-Dn}. Nous allons montrer par récurrence sur $t$
que, si $n = t + 1$, cela conduit aux possibilités
(\ref{models-item-Dn-extended-up}) et (\ref{models-item-Dn-extended-down})
du lemme \ref{models-lemma-genus-one}, et que, si $n > t + 1$, cela conduit au
cas (\ref{models-item-Dn}) du lemme \ref{models-lemma-Dn}.
L'hypothèse de récurrence (ou le lemme \ref{models-lemma-D5} lorsque $t = 5$)
montre que $a_{i_ji_k}$ est nul en dehors des termes requis
non nuls pour $j, k \geq 2$.
De plus, on voit que $w_2 = \ldots = w_{t + 1} = w$ pour un entier
$w$ et que les $a_{i_ji_k}$ non nuls pour $j, k \geq 2$ sont
égaux à $w$. En appliquant le lemme \ref{models-lemma-long}
(ou le lemme \ref{models-lemma-five-by-five} si $t = 5$) à la suite
$i_1, \ldots, i_t$ et à la suite
$i_1, \ldots, i_{t - 1}, i_{t + 1}$, on conclut que
$a_{i_1 i_j} = 0$ pour $j \geq 3$, que
$w_1$ vaut $\frac{1}{2}w$, $w$ ou $2w$
et que respectivement $a_{i_1i_2}$ vaut $w, w, 2w$.
Cela donne $3$ possibilités. Dans chaque cas, il est facile de
déterminer ce qui se passe :
\begin{enumerate}
\item Si $w_1 = \frac{1}{2}w$, alors on se trouve dans le cas
(\ref{models-item-Dn-extended-down}) du lemme \ref{models-lemma-genus-one}.
\item Si $w_1 = w$, alors on se trouve dans le cas
(\ref{models-item-Dn}) du lemme \ref{models-lemma-Dn}.
\item Si $w_1 = 2w$, alors on se trouve dans le cas
(\ref{models-item-Dn-extended-up}) du lemme \ref{models-lemma-genus-one}.
\end{enumerate}

\begin{lemma}
\label{models-lemma-Dn}
Classification des sous-graphes propres de la forme
$$
\xymatrix{
\bullet \ar@{-}[r] & \bullet \ar@{..}[r] & \bullet \ar@{-}[r] &
\bullet \ar@{-}[r] \ar@{-}[d] & \bullet \\
& & & \bullet
}
$$
Soient $t > 4$ et $n > t + 1$. Alors, pour $t + 1$ indices
$(-2)$ distincts $i_1, \ldots, i_{t + 1}$ tels que $a_{i_ji_{j + 1}}$
soit non nul pour $j = 1, \ldots, t - 1$ et que $a_{i_{t - 1}i_{t + 1}}$
soit non nul, on a les possibilités suivantes pour les $a$ et les $w$ :
\begin{enumerate}
\item
\label{models-item-Dn}
ils sont donnés par $w_{i_1} = w_{i_2} = \ldots = w_{i_{t + 1}} = w$,
$a_{i_ji_{j + 1}} = w$ pour $j = 1, \ldots, t - 1$,
$a_{i_{t - 1}i_{t + 1}} = w$ et $a_{i_ji_k} = 0$ pour les autres
paires $(j, k)$ telles que $j > k$.
\end{enumerate}
\end{lemma}

\begin{proof}
Voir la discussion ci-dessus.
\end{proof}

\noindent
Supposons donnés $6$ indices $(-2)$ distincts $g, h, i, j, k, l$
tels que $a_{gh}, a_{hi}, a_{ij}, a_{jk}, a_{il}$ soient non nuls.
Voir la figure du lemme \ref{models-lemma-E6}. On peut alors appliquer le
lemme \ref{models-lemma-D5} pour voir que l'on se trouve nécessairement dans la situation
du lemme \ref{models-lemma-E6}. Comme le déterminant vaut $3w^6 > 0$,
on conclut que, dans ce cas, on n'a jamais $n = 6$ !

\begin{lemma}
\label{models-lemma-E6}
Classification des sous-graphes propres de la forme
$$
\xymatrix{
\bullet \ar@{-}[r] & \bullet \ar@{-}[r] & \bullet \ar@{-}[r] \ar@{-}[d] &
\bullet \ar@{-}[r] & \bullet \\
& & \bullet
}
$$
Soit $n > 6$. Alors, pour $6$ indices $(-2)$ distincts
$i_1, \ldots, i_6$ tels que
$a_{12}, a_{23}, a_{34},
a_{45}, a_{36}$ soient non nuls, on a
les possibilités suivantes pour les $m$, les $a$ et les $w$ :
\begin{enumerate}
\item
\label{models-item-E6}
ils sont donnés par
$$
\left(
\begin{matrix}
m_1 \\
m_2 \\
m_3 \\
m_4 \\
m_5 \\
m_6
\end{matrix}
\right),
\quad
\left(
\begin{matrix}
-2w & w & 0 & 0 & 0 & 0 \\
w & -2w & w & 0 & 0 & 0 \\
0 & w & -2w & w & 0 & w \\
0 & 0 & w & -2w & w & 0 \\
0 & 0 & 0 & w & -2w & 0 \\
0 & 0 & w & 0 & 0 & -2w
\end{matrix}
\right),
\quad
\left(
\begin{matrix}
w \\
w \\
w \\
w \\
w \\
w
\end{matrix}
\right)
$$
avec $2m_1 \geq m_2$, $2m_2 \geq m_1 + m_3$, $2m_3 \geq m_2 + m_4 + m_6$,
$2m_4 \geq m_3 + m_5$, $2m_5 \geq m_3$ et $2m_6 \geq m_3$.
\end{enumerate}
\end{lemma}

\begin{proof}
Voir la discussion ci-dessus.
\end{proof}

\noindent
Supposons que $t \geq 4$ et que $i_0, \ldots, i_{t + 1}$ soient
$t + 2$ indices $(-2)$ distincts tels que
$a_{i_ji_{j + 1}} > 0$ pour $j = 1, \ldots, t - 1$,
$a_{i_0i_2} > 0$ et $a_{i_{t - 1}i_{t + 1}} > 0$.
Voir la figure du lemme \ref{models-lemma-double-triple}.
On peut alors appliquer les lemmes \ref{models-lemma-D5} et \ref{models-lemma-Dn}
pour voir que tous les autres $a_{i_ji_k}$ pour $j < k$ sont nuls,
que $w_{i_0} = \ldots = w_{i_{t + 1}} = w$ pour un certain
entier $w$ et que les coefficients hors diagonale requis
non nuls de $A$ sont égaux à $w$.
Un calcul montre que le déterminant de la
matrice correspondante est nul. Ainsi, $n = t + 2$
et l'on se trouve dans le cas (\ref{models-item-double-triple}) du
lemme \ref{models-lemma-genus-one}.

\begin{lemma}
\label{models-lemma-double-triple}
Non-existence de sous-graphes propres de la forme
$$
\xymatrix{
\bullet \ar@{-}[r] & \bullet \ar@{..}[r] \ar@{-}[d] &
\bullet \ar@{-}[d] \ar@{-}[r] & \bullet \\
& \bullet & \bullet
}
$$
Supposons $t \geq 4$ et $n > t + 2$.
Il n'existe {\bf pas} $t + 2$ indices
$(-2)$ distincts $i_0, \ldots, i_{t + 1}$ tels que
$a_{i_ji_{j + 1}} > 0$ pour $j = 1, \ldots, t - 1$,
$a_{i_0i_2} > 0$ et $a_{i_{t - 1}i_{t + 1}} > 0$.
\end{lemma}

\begin{proof}
Voir la discussion ci-dessus.
\end{proof}

\noindent
Supposons donnés $7$ indices $(-2)$ distincts
$f, g, h, i,
j, k, l$
tels que les nombres
$a_{fg}, a_{gh}, a_{ij},
a_{jh}, a_{kl}, a_{lh}$ soient non nuls.
Voir la figure du lemme \ref{models-lemma-E6-completed}. On peut alors appliquer le
lemme \ref{models-lemma-D5} pour voir que la matrice correspondante est
$$
\left(
\begin{matrix}
-2w & w & 0 & 0 & 0 & 0 & 0 \\
w & -2w & w & 0 & 0 & 0 & 0 \\
0 & w & -2w & 0 & w & 0 & w \\
0 & 0 & 0 & -2w & w & 0 & 0 \\
0 & 0 & w & w & -2w & 0 & 0 \\
0 & 0 & 0 & 0 & 0 & -2w & w \\
0 & 0 & w & 0 & 0 & w & -2w
\end{matrix}
\right)
$$
Comme le déterminant vaut $0$, on conclut que nécessairement $n = 7$
et $g = 1$, et l'on obtient le cas (\ref{models-item-E6-completed})
du lemme \ref{models-lemma-genus-one}.

\begin{lemma}
\label{models-lemma-E6-completed}
Non-existence de sous-graphes propres de la forme
$$
\xymatrix{
\bullet \ar@{-}[r] & \bullet \ar@{-}[r] & \bullet \ar@{-}[r] \ar@{-}[d] &
\bullet \ar@{-}[r] & \bullet \\
& & \bullet \ar@{-}[d] \\
& & \bullet
}
$$
Supposons $n > 7$. Il n'existe {\bf pas} $7$ indices
$(-2)$ distincts $f, g, h, i, j, k, l$
tels que $a_{fg}, a_{gh}, a_{ij}, a_{jh}, a_{kl}, a_{lh}$ soient non nuls.
\end{lemma}

\begin{proof}
Voir la discussion ci-dessus.
\end{proof}

\noindent
Supposons donnés $7$ indices $(-2)$ distincts
$f, g, h, i,
j, k, l$
tels que les nombres
$a_{fg}, a_{gh}, a_{hi},
a_{ij}, a_{jk}, a_{il}$ soient non nuls.
Voir la figure du lemme \ref{models-lemma-E7}. On peut alors appliquer les
lemmes \ref{models-lemma-D5} et \ref{models-lemma-Dn}
pour voir que l'on se trouve nécessairement dans la situation
du lemme \ref{models-lemma-E7}. Comme le déterminant vaut $-8w^7 > 0$,
on conclut que, dans ce cas, on n'a jamais $n = 7$ !

\begin{lemma}
\label{models-lemma-E7}
Classification des sous-graphes propres de la forme
$$
\xymatrix{
\bullet \ar@{-}[r] & \bullet \ar@{-}[r] & \bullet \ar@{-}[r] &
\bullet \ar@{-}[r] \ar@{-}[d] & \bullet \ar@{-}[r] & \bullet \\
& & & \bullet
}
$$
Soit $n > 7$. Alors, pour $7$ indices $(-2)$ distincts
$i_1, \ldots, i_7$ tels que
$a_{12}, a_{23}, a_{34},
a_{45}, a_{56}, a_{47}$ soient non nuls,
on a les possibilités suivantes pour les $m$, les $a$ et les $w$ :
\begin{enumerate}
\item
\label{models-item-E7}
ils sont donnés par
$$
\left(
\begin{matrix}
m_1 \\
m_2 \\
m_3 \\
m_4 \\
m_5 \\
m_6 \\
m_7
\end{matrix}
\right),
\quad
\left(
\begin{matrix}
-2w & w & 0 & 0 & 0 & 0 & 0 \\
w & -2w & w & 0 & 0 & 0 & 0 \\
0 & w & -2w & w & 0 & 0 & 0 \\
0 & 0 & w & -2w & w & 0 & w \\
0 & 0 & 0 & w & -2w & w & 0 \\
0 & 0 & 0 & 0 & w & -2w & 0 \\
0 & 0 & 0 & w & 0 & 0 & -2w
\end{matrix}
\right),
\quad
\left(
\begin{matrix}
w \\
w \\
w \\
w \\
w \\
w \\
w
\end{matrix}
\right)
$$
avec $2m_1 \geq m_2$, $2m_2 \geq m_1 + m_3$, $2m_3 \geq m_2 + m_4$,
$2m_4 \geq m_3 + m_5 + m_7$, $2m_5 \geq m_4 + m_6$, $2m_6 \geq m_5$
et $2m_7 \geq m_4$.
\end{enumerate}
\end{lemma}

\begin{proof}
Voir la discussion ci-dessus.
\end{proof}

\noindent
Supposons donnés $8$ indices $(-2)$ distincts dont la configuration
des coefficients non nuls $a_{ij}$ de la matrice $A$ est de la forme
$$
\xymatrix{
\bullet \ar@{-}[r] & \bullet \ar@{-}[r] & \bullet \ar@{-}[r] &
\bullet \ar@{-}[r] & \bullet \ar@{-}[r] \ar@{-}[d] &
\bullet \ar@{-}[r] & \bullet \\
& & & & \bullet
}
$$
ou de la forme
$$
\xymatrix{
\bullet \ar@{-}[r] & \bullet \ar@{-}[r] & \bullet \ar@{-}[r] &
\bullet \ar@{-}[r] \ar@{-}[d] & \bullet \ar@{-}[r] &
\bullet \ar@{-}[r] & \bullet \\
& & & \bullet
}
$$
En raisonnant exactement comme dans la démonstration du lemme \ref{models-lemma-E7},
on voit que la première configuration conduit au cas
(\ref{models-item-E8}) du lemme \ref{models-lemma-E8}
et ne conduit pas à un nouveau cas dans le lemme \ref{models-lemma-genus-one}.
En raisonnant exactement comme dans la démonstration du lemme \ref{models-lemma-E6-completed},
on voit que la seconde configuration ne se présente pas si
$n > 8$, mais qu'elle conduit au cas (\ref{models-item-E7-completed})
du lemme \ref{models-lemma-genus-one} lorsque $n = 8$.

\begin{lemma}
\label{models-lemma-E8}
Classification des sous-graphes propres de la forme
$$
\xymatrix{
\bullet \ar@{-}[r] & \bullet \ar@{-}[r] & \bullet \ar@{-}[r] &
\bullet \ar@{-}[r] & \bullet \ar@{-}[r] \ar@{-}[d] &
\bullet \ar@{-}[r] & \bullet \\
& & & & \bullet
}
$$
Soit $n > 8$. Alors, pour $8$ indices $(-2)$ distincts
$i_1, \ldots, i_8$ tels que
$a_{12}, a_{23}, a_{34},
a_{45}, a_{56}, a_{65}, a_{57}$
soient non nuls, on a les possibilités suivantes pour les $m$, les $a$ et les $w$ :
\begin{enumerate}
\item
\label{models-item-E8}
ils sont donnés par
$$
\left(
\begin{matrix}
m_1 \\
m_2 \\
m_3 \\
m_4 \\
m_5 \\
m_6 \\
m_7 \\
m_8
\end{matrix}
\right),
\quad
\left(
\begin{matrix}
-2w & w & 0 & 0 & 0 & 0 & 0 & 0 \\
w & -2w & w & 0 & 0 & 0 & 0 & 0 \\
0 & w & -2w & w & 0 & 0 & 0 & 0 \\
0 & 0 & w & -2w & w & 0 & 0 & 0 \\
0 & 0 & 0 & w & -2w & w & 0 & w \\
0 & 0 & 0 & 0 & w & -2w & w & 0 \\
0 & 0 & 0 & 0 & 0 & w & -2w & 0 \\
0 & 0 & 0 & 0 & w & 0 & 0 & -2w
\end{matrix}
\right),
\quad
\left(
\begin{matrix}
w \\
w \\
w \\
w \\
w \\
w \\
w \\
w
\end{matrix}
\right)
$$
avec $2m_1 \geq m_2$, $2m_2 \geq m_1 + m_3$, $2m_3 \geq m_2 + m_4$,
$2m_4 \geq m_3 + m_5$, $2m_5 \geq m_4 + m_6 + m_8$, $2m_6 \geq m_5 + m_7$,
$2m_7 \geq m_6$ et $2m_8 \geq m_5$.
\end{enumerate}
\end{lemma}

\begin{proof}
Voir la discussion ci-dessus.
\end{proof}

\begin{lemma}
\label{models-lemma-E7-completed}
Non-existence de sous-graphes propres de la forme
$$
\xymatrix{
\bullet \ar@{-}[r] &
\bullet \ar@{-}[r] &
\bullet \ar@{-}[r] & \bullet \ar@{-}[r] \ar@{-}[d] &
\bullet \ar@{-}[r] & \bullet \ar@{-}[r] & \bullet \\
& & & \bullet
}
$$
Supposons $n > 8$. Il n'existe {\bf pas} $8$ indices
$(-2)$ distincts $e, f, g, h, i, j, k, l$
tels que $a_{ef}, a_{fg}, a_{gh}, a_{hi}, a_{ij}, a_{jk}, a_{lh}$
soient non nuls.
\end{lemma}

\begin{proof}
Voir la discussion ci-dessus.
\end{proof}

\noindent
Supposons donnés $9$ indices $(-2)$ distincts dont la configuration
des coefficients non nuls $a_{ij}$ de la matrice $A$ est de la forme
$$
\xymatrix{
\bullet \ar@{-}[r] & \bullet \ar@{-}[r] &
\bullet \ar@{-}[r] & \bullet \ar@{-}[r] &
\bullet \ar@{-}[r] & \bullet \ar@{-}[r] \ar@{-}[d] &
\bullet \ar@{-}[r] & \bullet \\
& & & & & \bullet
}
$$
En raisonnant exactement comme dans la démonstration du lemme \ref{models-lemma-E6-completed},
on voit que cette configuration ne se présente pas si
$n > 9$, mais qu'elle conduit au cas (\ref{models-item-E8-completed})
du lemme \ref{models-lemma-genus-one} lorsque $n = 9$.

\begin{lemma}
\label{models-lemma-E8-completed}
Non-existence de sous-graphes propres de la forme
$$
\xymatrix{
\bullet \ar@{-}[r] & \bullet \ar@{-}[r] &
\bullet \ar@{-}[r] & \bullet \ar@{-}[r] &
\bullet \ar@{-}[r] & \bullet \ar@{-}[r] \ar@{-}[d] &
\bullet \ar@{-}[r] & \bullet \\
& & & & & \bullet
}
$$
Supposons $n > 9$. Il n'existe {\bf pas} $9$ indices
$(-2)$ distincts $d, e, f, g, h, i, j, k, l$
tels que $a_{de}, a_{ef}, a_{fg}, a_{gh}, a_{hi}, a_{ij}, a_{jk}, a_{lh}$
soient non nuls.
\end{lemma}

\begin{proof}
Voir la discussion ci-dessus.
\end{proof}

\noindent
En réunissant toutes les informations, on obtient le résultat suivant.

\begin{proposition}
\label{models-proposition-classify-subgraphs}
Soient $n, m_i, a_{ij}, w_i, g_i$ un type numérique de genre $g$.
Soit $I \subset \{1, \ldots, n\}$ un sous-ensemble strict de cardinal $\geq 2$
constitué d'indices $(-2)$ et tel qu'il
n'existe aucun sous-ensemble strict non vide $I' \subset I$
avec $a_{i'i} = 0$ pour $i' \in I$, $i \in I \setminus I'$.
Alors, à réindexation près, les $m_i$, les $a_{ij}$ et les $w_i$
pour $i, j \in I$ sont ceux qui figurent dans les
lemmes \ref{models-lemma-two-by-two},
\ref{models-lemma-three-by-three},
\ref{models-lemma-four-by-four},
\ref{models-lemma-D4},
\ref{models-lemma-five-by-five},
\ref{models-lemma-D5},
\ref{models-lemma-long},
\ref{models-lemma-Dn},
\ref{models-lemma-E6},
\ref{models-lemma-E7} ou
\ref{models-lemma-E8}.
\end{proposition}

\begin{proof}
Cela résulte de la discussion ci-dessus ; voir la discussion au
début de la section \ref{models-section-classify-proper-subgraphs}.
\end{proof}













\section{Classification des types minimaux de genres zéro et un}
\label{models-section-classification-genus-one}

\noindent
Le titre de la section dit tout.

\begin{lemma}[Genre zéro]
\label{models-lemma-genus-zero}
Le seul type numérique minimal de genre zéro est
$n = 1$, $m_1 = 1$, $a_{11} = 0$, $w_1 = 1$, $g_1 = 0$.
\end{lemma}

\begin{proof}
Cela résulte des lemmes \ref{models-lemma-non-irreducible-minimal-type-genus-at-least-one}
et \ref{models-lemma-irreducible}.
\end{proof}

\begin{lemma}[Genre un]
\label{models-lemma-genus-one}
À équivalence près, les types numériques minimaux de genre un sont les suivants :
\begin{enumerate}
\item
\label{models-item-one}
$n = 1$, $a_{11} = 0$, $g_1 = 1$, et $m_1, w_1 \geq 1$ sont arbitraires ;
\item
\label{models-item-two-cycle}
$n = 2$, et $m_i, a_{ij}, w_i, g_i$ sont donnés par
$$
\left(
\begin{matrix}
m \\
m
\end{matrix}
\right),
\quad
\left(
\begin{matrix}
-2w & 2w \\
2w & -2w
\end{matrix}
\right),
\quad
\left(
\begin{matrix}
w \\
w
\end{matrix}
\right),
\quad
\left(
\begin{matrix}
0 \\
0
\end{matrix}
\right)
$$
où $w$ et $m$ sont arbitraires ;
\item
\label{models-item-up4}
$n = 2$, et $m_i, a_{ij}, w_i, g_i$ sont donnés par
$$
\left(
\begin{matrix}
2m \\
m
\end{matrix}
\right),
\quad
\left(
\begin{matrix}
-2w & 4w \\
4w & -8w
\end{matrix}
\right),
\quad
\left(
\begin{matrix}
w \\
4w
\end{matrix}
\right),
\quad
\left(
\begin{matrix}
0 \\
0
\end{matrix}
\right)
$$
où $w$ et $m$ sont arbitraires ;
\item
\label{models-item-three-cycle}
$n = 3$, et $m_i, a_{ij}, w_i, g_i$ sont donnés par
$$
\left(
\begin{matrix}
m \\
m \\
m
\end{matrix}
\right),
\quad
\left(
\begin{matrix}
-2w & w & w \\
w & -2w & w \\
w & w & -2w
\end{matrix}
\right),
\quad
\left(
\begin{matrix}
w \\
w \\
w
\end{matrix}
\right),
\quad
\left(
\begin{matrix}
0 \\
0 \\
0
\end{matrix}
\right)
$$
où $w$ et $m$ sont arbitraires ;
\item
\label{models-item-equal-up3}
$n = 3$, et $m_i, a_{ij}, w_i, g_i$ sont donnés par
$$
\left(
\begin{matrix}
m \\
2m \\
m
\end{matrix}
\right),
\quad
\left(
\begin{matrix}
-2w & w & 0 \\
w & -2w & 3w \\
0 & 3w & -6w
\end{matrix}
\right),
\quad
\left(
\begin{matrix}
w \\
w \\
3w
\end{matrix}
\right),
\quad
\left(
\begin{matrix}
0 \\
0 \\
0
\end{matrix}
\right)
$$
où $w$ et $m$ sont arbitraires ;
\item
\label{models-item-equal-down3}
$n = 3$, et $m_i, a_{ij}, w_i, g_i$ sont donnés par
$$
\left(
\begin{matrix}
m \\
2m \\
3m
\end{matrix}
\right),
\quad
\left(
\begin{matrix}
-6w & 3w & 0 \\
3w & -6w & 3w \\
0 & 3w & -2w
\end{matrix}
\right),
\quad
\left(
\begin{matrix}
3w \\
3w \\
w
\end{matrix}
\right),
\quad
\left(
\begin{matrix}
0 \\
0 \\
0
\end{matrix}
\right)
$$
où $w$ et $m$ sont arbitraires ;
\item
\label{models-item-up-up}
$n = 3$, et $m_i, a_{ij}, w_i, g_i$ sont donnés par
$$
\left(
\begin{matrix}
2m \\
2m \\
m
\end{matrix}
\right),
\quad
\left(
\begin{matrix}
-2w & 2w & 0 \\
2w & -4w & 4w \\
0 & 4w & -8w
\end{matrix}
\right),
\quad
\left(
\begin{matrix}
w \\
2w \\
4w
\end{matrix}
\right),
\quad
\left(
\begin{matrix}
0 \\
0 \\
0
\end{matrix}
\right)
$$
où $w$ et $m$ sont arbitraires ;
\item
\label{models-item-up-down}
$n = 3$, et $m_i, a_{ij}, w_i, g_i$ sont donnés par
$$
\left(
\begin{matrix}
m \\
m \\
m
\end{matrix}
\right),
\quad
\left(
\begin{matrix}
-2w & 2w & 0 \\
2w & -4w & 2w \\
0 & 2w & -2w
\end{matrix}
\right),
\quad
\left(
\begin{matrix}
w \\
2w \\
w
\end{matrix}
\right),
\quad
\left(
\begin{matrix}
0 \\
0 \\
0
\end{matrix}
\right)
$$
où $w$ et $m$ sont arbitraires ;
\item
\label{models-item-down-up}
$n = 3$, et $m_i, a_{ij}, w_i, g_i$ sont donnés par
$$
\left(
\begin{matrix}
m \\
2m \\
m
\end{matrix}
\right),
\quad
\left(
\begin{matrix}
-4w & 2w & 0 \\
2w & -2w & 2w \\
0 & 2w & -4w
\end{matrix}
\right),
\quad
\left(
\begin{matrix}
2w \\
w \\
2w
\end{matrix}
\right),
\quad
\left(
\begin{matrix}
0 \\
0 \\
0
\end{matrix}
\right)
$$
où $w$ et $m$ sont arbitraires ;
\item
\label{models-item-four-cycle}
$n = 4$, et $m_i, a_{ij}, w_i, g_i$ sont donnés par
$$
\left(
\begin{matrix}
m \\
m \\
m \\
m
\end{matrix}
\right),
\quad
\left(
\begin{matrix}
-2w & w & 0 & w \\
w & -2w & w & 0 \\
0 & w & -2w & w \\
w & 0 & w & -2w
\end{matrix}
\right),
\quad
\left(
\begin{matrix}
w \\
w \\
w \\
w
\end{matrix}
\right),
\quad
\left(
\begin{matrix}
0 \\
0 \\
0 \\
0
\end{matrix}
\right)
$$
où $w$ et $m$ sont arbitraires ;
\item
\label{models-item-up-equal-up}
$n = 4$, et $m_i, a_{ij}, w_i, g_i$ sont donnés par
$$
\left(
\begin{matrix}
2m \\
2m \\
2m \\
m
\end{matrix}
\right),
\quad
\left(
\begin{matrix}
-2w & 2w & 0 & 0 \\
2w & -4w & 2w & 0 \\
0 & 2w & -4w & 4w \\
0 & 0 & 4w & -8w
\end{matrix}
\right),
\quad
\left(
\begin{matrix}
w \\
2w \\
2w \\
4w
\end{matrix}
\right),
\quad
\left(
\begin{matrix}
0 \\
0 \\
0 \\
0
\end{matrix}
\right)
$$
où $w$ et $m$ sont arbitraires ;
\item
\label{models-item-up-equal-down}
$n = 4$, et $m_i, a_{ij}, w_i, g_i$ sont donnés par
$$
\left(
\begin{matrix}
m \\
m \\
m \\
m
\end{matrix}
\right),
\quad
\left(
\begin{matrix}
-2w & 2w & 0 & 0 \\
2w & -4w & 2w & 0 \\
0 & 2w & -4w & 2w \\
0 & 0 & 2w & -2w
\end{matrix}
\right),
\quad
\left(
\begin{matrix}
w \\
2w \\
2w \\
w
\end{matrix}
\right),
\quad
\left(
\begin{matrix}
0 \\
0 \\
0 \\
0
\end{matrix}
\right)
$$
où $w$ et $m$ sont arbitraires ;
\item
\label{models-item-down-equal-up}
$n = 4$, et $m_i, a_{ij}, w_i, g_i$ sont donnés par
$$
\left(
\begin{matrix}
m \\
2m \\
2m \\
m
\end{matrix}
\right),
\quad
\left(
\begin{matrix}
-4w & 2w & 0 & 0 \\
2w & -2w & w & 0 \\
0 & w & -2w & 2w \\
0 & 0 & 2w & -4w
\end{matrix}
\right),
\quad
\left(
\begin{matrix}
2w \\
w \\
w \\
2w
\end{matrix}
\right),
\quad
\left(
\begin{matrix}
0 \\
0 \\
0 \\
0
\end{matrix}
\right)
$$
où $w$ et $m$ sont arbitraires ;
\item
\label{models-item-triple-with-up}
$n = 4$, et $m_i, a_{ij}, w_i, g_i$ sont donnés par
$$
\left(
\begin{matrix}
2m \\
m \\
m \\
m
\end{matrix}
\right),
\quad
\left(
\begin{matrix}
-2w & w & w & 2w \\
w & -2w & 0 & 0 \\
w & 0 & -2w & 0 \\
2w & 0 & 0 & -4w
\end{matrix}
\right),
\quad
\left(
\begin{matrix}
w \\
w \\
w \\
2w
\end{matrix}
\right),
\quad
\left(
\begin{matrix}
0 \\
0 \\
0 \\
0
\end{matrix}
\right)
$$
où $w$ et $m$ sont arbitraires ;
\item
\label{models-item-triple-with-down}
$n = 4$, et $m_i, a_{ij}, w_i, g_i$ sont donnés par
$$
\left(
\begin{matrix}
2m \\
m \\
m \\
2m
\end{matrix}
\right),
\quad
\left(
\begin{matrix}
-4w & 2w & 2w & 2w \\
2w & -4w & 0 & 0 \\
2w & 0 & -4w & 0 \\
2w & 0 & 0 & -2w
\end{matrix}
\right),
\quad
\left(
\begin{matrix}
2w \\
2w \\
2w \\
w
\end{matrix}
\right),
\quad
\left(
\begin{matrix}
0 \\
0 \\
0 \\
0
\end{matrix}
\right)
$$
où $w$ et $m$ sont arbitraires ;
\item
\label{models-item-five-cycle}
$n = 5$, et $m_i, a_{ij}, w_i, g_i$ sont donnés par
$$
\left(
\begin{matrix}
m \\
m \\
m \\
m \\
m
\end{matrix}
\right),
\quad
\left(
\begin{matrix}
-2w & w & 0 & 0 & w \\
w & -2w & w & 0 & 0 \\
0 & w & -2w & w & 0 \\
0 & 0 & w & -2w & w \\
w & 0 & 0 & w & -2w \\
\end{matrix}
\right),
\quad
\left(
\begin{matrix}
w \\
w \\
w \\
w \\
w
\end{matrix}
\right),
\quad
\left(
\begin{matrix}
0 \\
0 \\
0 \\
0 \\
0
\end{matrix}
\right)
$$
où $w$ et $m$ sont arbitraires ;
\item
\label{models-item-equal-equal-up-equal}
$n = 5$, et $m_i, a_{ij}, w_i, g_i$ sont donnés par
$$
\left(
\begin{matrix}
m \\
2m \\
3m \\
2m \\
m
\end{matrix}
\right),
\quad
\left(
\begin{matrix}
-2w & w & 0 & 0 & 0 \\
w & -2w & w & 0 & 0 \\
0 & w & -2w & 2w & 0 \\
0 & 0 & 2w & -4w & 2w \\
0 & 0 & 0 & 2w & -4w \\
\end{matrix}
\right),
\quad
\left(
\begin{matrix}
w \\
w \\
w \\
2w \\
2w
\end{matrix}
\right),
\quad
\left(
\begin{matrix}
0 \\
0 \\
0 \\
0 \\
0
\end{matrix}
\right)
$$
où $w$ et $m$ sont arbitraires ;
\item
\label{models-item-equal-equal-down-equal}
$n = 5$, et $m_i, a_{ij}, w_i, g_i$ sont donnés par
$$
\left(
\begin{matrix}
m \\
2m \\
3m \\
4m \\
2m
\end{matrix}
\right),
\quad
\left(
\begin{matrix}
-4w & 2w & 0 & 0 & 0 \\
2w & -4w & 2w & 0 & 0 \\
0 & 2w & -4w & 2w & 0 \\
0 & 0 & 2w & -2w & w \\
0 & 0 & 0 & w & -2w \\
\end{matrix}
\right),
\quad
\left(
\begin{matrix}
2w \\
2w \\
2w \\
w \\
w
\end{matrix}
\right),
\quad
\left(
\begin{matrix}
0 \\
0 \\
0 \\
0 \\
0
\end{matrix}
\right)
$$
où $w$ et $m$ sont arbitraires ;
\item
\label{models-item-up-equal-equal-up}
$n = 5$, et $m_i, a_{ij}, w_i, g_i$ sont donnés par
$$
\left(
\begin{matrix}
2m \\
2m \\
2m \\
2m \\
m
\end{matrix}
\right),
\quad
\left(
\begin{matrix}
-2w & 2w & 0 & 0 & 0 \\
2w & -4w & 2w & 0 & 0 \\
0 & 2w & -4w & 2w & 0 \\
0 & 0 & 2w & -4w & 4w \\
0 & 0 & 0 & 4w & -8w \\
\end{matrix}
\right),
\quad
\left(
\begin{matrix}
w \\
2w \\
2w \\
2w \\
4w
\end{matrix}
\right),
\quad
\left(
\begin{matrix}
0 \\
0 \\
0 \\
0 \\
0
\end{matrix}
\right)
$$
où $w$ et $m$ sont arbitraires ;
\item
\label{models-item-up-equal-equal-down}
$n = 5$, et $m_i, a_{ij}, w_i, g_i$ sont donnés par
$$
\left(
\begin{matrix}
m \\
m \\
m \\
m \\
m
\end{matrix}
\right),
\quad
\left(
\begin{matrix}
-2w & 2w & 0 & 0 & 0 \\
2w & -4w & 2w & 0 & 0 \\
0 & 2w & -4w & 2w & 0 \\
0 & 0 & 2w & -4w & 2w \\
0 & 0 & 0 & 2w & -2w \\
\end{matrix}
\right),
\quad
\left(
\begin{matrix}
w \\
2w \\
2w \\
2w \\
w
\end{matrix}
\right),
\quad
\left(
\begin{matrix}
0 \\
0 \\
0 \\
0 \\
0
\end{matrix}
\right)
$$
où $w$ et $m$ sont arbitraires ;
\item
\label{models-item-down-equal-equal-up}
$n = 5$, et $m_i, a_{ij}, w_i, g_i$ sont donnés par
$$
\left(
\begin{matrix}
m \\
2m \\
2m \\
2m \\
m
\end{matrix}
\right),
\quad
\left(
\begin{matrix}
-4w & 2w & 0 & 0 & 0 \\
2w & -2w & w & 0 & 0 \\
0 & w & -2w & w & 0 \\
0 & 0 & w & -2w & 2w \\
0 & 0 & 0 & 2w & -4w \\
\end{matrix}
\right),
\quad
\left(
\begin{matrix}
2w \\
w \\
w \\
w \\
2w
\end{matrix}
\right),
\quad
\left(
\begin{matrix}
0 \\
0 \\
0 \\
0 \\
0
\end{matrix}
\right)
$$
où $w$ et $m$ sont arbitraires ;
\item
\label{models-item-quadruple}
$n = 5$, et $m_i, a_{ij}, w_i, g_i$ sont donnés par
$$
\left(
\begin{matrix}
2m \\
m \\
m \\
m \\
m
\end{matrix}
\right),
\quad
\left(
\begin{matrix}
-2w & w & w & w & w \\
w & -2w & 0 & 0 & 0 \\
w & 0 & -2w & 0 & 0 \\
w & 0 & 0 & -2w & 0 \\
w & 0 & 0 & 0 & -2w \\
\end{matrix}
\right),
\quad
\left(
\begin{matrix}
w \\
w \\
w \\
w \\
w
\end{matrix}
\right),
\quad
\left(
\begin{matrix}
0 \\
0 \\
0 \\
0 \\
0
\end{matrix}
\right)
$$
où $w$ et $m$ sont arbitraires ;
\item
\label{models-item-triple-extended-up}
$n = 5$, et $m_i, a_{ij}, w_i, g_i$ sont donnés par
$$
\left(
\begin{matrix}
m \\
2m \\
2m \\
m \\
m
\end{matrix}
\right),
\quad
\left(
\begin{matrix}
-4w & 2w & 0 & 0 & 0 \\
2w & -2w & w & 0 & 0 \\
0 & w & -2w & w & w \\
0 & 0 & w & -2w & 0 \\
0 & 0 & w & 0 & -2w \\
\end{matrix}
\right),
\quad
\left(
\begin{matrix}
2w \\
w \\
w \\
w \\
w
\end{matrix}
\right),
\quad
\left(
\begin{matrix}
0 \\
0 \\
0 \\
0 \\
0
\end{matrix}
\right)
$$
où $w$ et $m$ sont arbitraires ;
\item
\label{models-item-triple-extended-down}
$n = 5$, et $m_i, a_{ij}, w_i, g_i$ sont donnés par
$$
\left(
\begin{matrix}
2m \\
2m \\
2m \\
m \\
m
\end{matrix}
\right),
\quad
\left(
\begin{matrix}
-2w & 2w & 0 & 0 & 0 \\
2w & -4w & 2w & 0 & 0 \\
0 & 2w & -4w & 2w & 2w \\
0 & 0 & 2w & -4w & 0 \\
0 & 0 & 2w & 0 & -4w \\
\end{matrix}
\right),
\quad
\left(
\begin{matrix}
w \\
2w \\
2w \\
2w \\
2w
\end{matrix}
\right),
\quad
\left(
\begin{matrix}
0 \\
0 \\
0 \\
0 \\
0
\end{matrix}
\right)
$$
où $w$ et $m$ sont arbitraires ;
\item
\label{models-item-n-cycle}
$n \geq 6$ et l'on a un $n$-cycle qui généralise (\ref{models-item-five-cycle}) :
\begin{enumerate}
\item $m_1 = \ldots = m_n = m$,
\item $a_{12} = \ldots = a_{(n - 1) n} = w$, $a_{1n} = w$,
et, pour les autres $i < j$, on a $a_{ij} = 0$,
\item $w_1 = \ldots = w_n = w$
\end{enumerate}
où $w$ et $m$ sont arbitraires ;
\item
\label{models-item-up-chain-equal-up}
$n \geq 6$ et l'on a une chaîne qui généralise (\ref{models-item-up-equal-equal-up}) :
\begin{enumerate}
\item $m_1 = \ldots = m_{n - 1} = 2m$, $m_n = m$,
\item $a_{12} = \ldots = a_{(n - 2) (n - 1)} = 2w$, $a_{(n - 1) n} = 4w$,
et, pour les autres $i < j$, on a $a_{ij} = 0$,
\item $w_1 = w$, $w_2 = \ldots = w_{n - 1} = 2w$, $w_n = 4w$
\end{enumerate}
où $w$ et $m$ sont arbitraires ;
\item
\label{models-item-up-chain-equal-down}
$n \geq 6$ et l'on a une chaîne qui généralise (\ref{models-item-up-equal-equal-down}) :
\begin{enumerate}
\item $m_1 = \ldots = m_n = m$,
\item $a_{12} = \ldots = a_{(n - 1) n} = w$,
et, pour les autres $i < j$, on a $a_{ij} = 0$,
\item $w_1 = w$, $w_2 = \ldots = w_{n - 1} = 2w$, $w_n = w$
\end{enumerate}
où $w$ et $m$ sont arbitraires ;
\item
\label{models-item-down-chain-equal-up}
$n \geq 6$ et l'on a une chaîne qui généralise (\ref{models-item-down-equal-equal-up}) :
\begin{enumerate}
\item $m_1 = w$, $w_2 = \ldots = m_{n - 1} = 2m$, $m_n = m$,
\item $a_{12} = 2w$, $a_{23} = \ldots = a_{(n - 2) (n - 1)} = w$,
$a_{(n - 1) n} = 2w$ et, pour les autres $i < j$, on a $a_{ij} = 0$,
\item $w_1 = 2w$, $w_2 = \ldots = w_{n - 1} = w$, $w_n = 2w$
\end{enumerate}
où $w$ et $m$ sont arbitraires ;
\item
\label{models-item-Dn-extended-up}
$n \geq 6$ et l'on a un type qui généralise (\ref{models-item-triple-extended-up}) :
\begin{enumerate}
\item $m_1 = m$, $m_2 = \ldots = m_{n - 3} = 2m$, $m_{n - 1} = m_n = m$,
\item $a_{12} = 2w$, $a_{23} = \ldots = a_{(n - 2) (n - 1)} = w$,
$a_{(n - 2) n} = w$ et, pour les autres $i < j$, on a $a_{ij} = 0$,
\item $w_1 = 2w$, $w_2 = \ldots = w_n = w$
\end{enumerate}
où $w$ et $m$ sont arbitraires ;
\item
\label{models-item-Dn-extended-down}
$n \geq 6$ et l'on a un type qui généralise (\ref{models-item-triple-extended-down}) :
\begin{enumerate}
\item $m_1 = \ldots = m_{n - 3} = 2m$, $m_{n - 1} = m_n = m$,
\item $a_{12} = \ldots = a_{(n - 2) (n - 1)} = 2w$,
$a_{(n - 2) n} = 2w$ et, pour les autres $i < j$, on a $a_{ij} = 0$,
\item $w_1 = w$, $w_2 = \ldots = w_n = 2w$
\end{enumerate}
où $w$ et $m$ sont arbitraires ;
\item
\label{models-item-double-triple}
$n \geq 6$ et l'on a un type qui généralise (\ref{models-item-quadruple}) :
\begin{enumerate}
\item $m_1 = m_2 = m$, $m_3 = \ldots = m_{n - 2} = 2m$, $m_{n - 1} = m_n = m$,
\item $a_{13} = w$, $a_{23} = \ldots = a_{(n - 2) (n - 1)} = w$,
$a_{(n - 2) n} = w$ et, pour les autres $i < j$, on a $a_{ij} = 0$,
\item $w_1 = \ldots = w_n = w$,
\end{enumerate}
où $w$ et $m$ sont arbitraires ;
\item
\label{models-item-E6-completed}
$n = 7$, et $m_i, a_{ij}, w_i, g_i$ sont donnés par
$$
\left(
\begin{matrix}
m \\
2m \\
3m \\
m \\
2m \\
m \\
2m
\end{matrix}
\right),
\quad
\left(
\begin{matrix}
-2w & w & 0 & 0 & 0 & 0 & 0 \\
w & -2w & w & 0 & 0 & 0 & 0 \\
0 & w & -2w & 0 & w & 0 & w \\
0 & 0 & 0 & -2w & w & 0 & 0 \\
0 & 0 & w & w & -2w & 0 & 0 \\
0 & 0 & 0 & 0 & 0 & -2w & w \\
0 & 0 & w & 0 & 0 & w & -2w
\end{matrix}
\right),
\quad
\left(
\begin{matrix}
w \\
w \\
w \\
w \\
w \\
w \\
w
\end{matrix}
\right),
\quad
\left(
\begin{matrix}
0 \\
0 \\
0 \\
0 \\
0 \\
0 \\
0
\end{matrix}
\right)
$$
où $w$ et $m$ sont arbitraires ;
\item
\label{models-item-E7-completed}
$n = 8$, et $m_i, a_{ij}, w_i, g_i$ sont donnés par
$$
\left(
\begin{matrix}
m \\
2m \\
3m \\
4m \\
3m \\
2m \\
m \\
2m
\end{matrix}
\right),
\quad
\left(
\begin{matrix}
-2w & w & 0 & 0 & 0 & 0 & 0 & 0 \\
w & -2w & w & 0 & 0 & 0 & 0 & 0 \\
0 & w & -2w & w & 0 & 0 & 0 & 0 \\
0 & 0 & w & -2w & w & 0 & 0 & w \\
0 & 0 & 0 & w & -2w & w & 0 & 0 \\
0 & 0 & 0 & 0 & w & -2w & w & 0 \\
0 & 0 & 0 & 0 & 0 & w & -2w & 0 \\
0 & 0 & 0 & w & 0 & 0 & 0 & -2w \\
\end{matrix}
\right),
\quad
\left(
\begin{matrix}
w \\
w \\
w \\
w \\
w \\
w \\
w \\
w
\end{matrix}
\right),
\quad
\left(
\begin{matrix}
0 \\
0 \\
0 \\
0 \\
0 \\
0 \\
0 \\
0
\end{matrix}
\right)
$$
où $w$ et $m$ sont arbitraires ;
\item
\label{models-item-E8-completed}
$n = 9$, et $m_i, a_{ij}, w_i, g_i$ sont donnés par
$$
\left(
\begin{matrix}
m \\
2m \\
3m \\
4m \\
5m \\
6m \\
4m \\
2m \\
3m
\end{matrix}
\right),
\quad
\left(
\begin{matrix}
-2w & w & 0 & 0 & 0 & 0 & 0 & 0 & 0 \\
w & -2w & w & 0 & 0 & 0 & 0 & 0 & 0 \\
0 & w & -2w & w & 0 & 0 & 0 & 0 & 0 \\
0 & 0 & w & -2w & w & 0 & 0 & 0 & 0 \\
0 & 0 & 0 & w & -2w & w & 0 & 0 & 0 \\
0 & 0 & 0 & 0 & w & -2w & w & 0 & w \\
0 & 0 & 0 & 0 & 0 & w & -2w & w & 0 \\
0 & 0 & 0 & 0 & 0 & 0 & w & -2w & 0 \\
0 & 0 & 0 & 0 & 0 & w & 0 & 0 & -2w \\
\end{matrix}
\right),
\quad
\left(
\begin{matrix}
w \\
w \\
w \\
w \\
w \\
w \\
w \\
w \\
w
\end{matrix}
\right),
\quad
\left(
\begin{matrix}
0 \\
0 \\
0 \\
0 \\
0 \\
0 \\
0 \\
0 \\
0
\end{matrix}
\right)
$$
où $w$ et $m$ sont arbitraires.
\end{enumerate}
\end{lemma}

\begin{proof}
Cela est démontré dans la section \ref{models-section-classify-proper-subgraphs}.
Voir la discussion au
début de la section \ref{models-section-classify-proper-subgraphs}.
\end{proof}






\section{Bornes pour les invariants des types numériques}
\label{models-section-towards-classification}

\noindent
Dans notre démonstration de la réduction semi-stable pour les courbes, nous utiliserons une borne
sur les groupes de Picard des types numériques de genre $g$, que nous établirons
dans cette section.

\begin{lemma}
\label{models-lemma-bound-neighbours}
Soient $n, m_i, a_{ij}, w_i, g_i$ un type numérique de genre $g$.
Pour $i, j$ tels que $a_{ij} > 0$, on a
$m_ia_{ij} \leq m_j|a_{jj}|$ et $m_iw_i \leq m_j|a_{jj}|$.
\end{lemma}

\begin{proof}
Pour tout indice $j$, on a $m_j a_{jj} + \sum_{i \not = j} m_ia_{ij} = 0$.
Ainsi, si l'on dispose d'une borne supérieure pour $|a_{jj}|$ et $m_j$, on obtient aussi une
borne supérieure pour les $a_{ij}$ non nuls (donc positifs), ainsi que pour
$m_i$. En rappelant que $w_i$ divise $a_{ij}$, le lecteur vérifiera aisément
le lemme.
\end{proof}

\begin{lemma}
\label{models-lemma-bound-heart}
Fixons $g \geq 2$. Pour tout type numérique minimal $n, m_i, a_{ij}, w_i, g_i$
de genre $g$ avec $n > 1$, on a :
\begin{enumerate}
\item l'ensemble $J \subset \{1, \ldots, n\}$ des indices non $(-2)$
possède au plus $2g - 2$ éléments ;
\item pour $j \in J$, on a $g_j < g$ ;
\item pour $j \in J$, on a $m_j|a_{jj}| \leq 6g - 6$ ; et
\item pour $j \in J$ et $i \in \{1, \ldots, n\}$,
on a $m_ia_{ij} \leq 6g - 6$.
\end{enumerate}
\end{lemma}

\begin{proof}
Rappelons que $g = 1 + \sum m_j(w_j(g_j - 1) - \frac{1}{2} a_{jj})$.
Pour $j \in J$, la contribution $m_j(w_j(g_j - 1) - \frac{1}{2} a_{jj})$
au genre $g$ est $> 0$, donc $\geq 1/2$. On utilise ici le
lemme \ref{models-lemma-minus-one},
la définition \ref{models-definition-type-minus-one},
la définition \ref{models-definition-type-minimal},
le lemme \ref{models-lemma-minus-two} et
la définition \ref{models-definition-type-minus-two} ; dans la suite, nous utiliserons ces
résultats sans le mentionner davantage.
Ainsi, $J$ possède au plus $2(g - 1)$ éléments.
Cela démontre (1).

\medskip\noindent
Rappelons que $-a_{ii} > 0$ pour tout $i$ d'après le lemme \ref{models-lemma-diagonal-negative}.
Par conséquent, pour $j \in J$, la contribution $m_j(w_j(g_j - 1) - \frac{1}{2} a_{jj})$
au genre $g$ est $> m_jw_j(g_j - 1)$. Ainsi,
$$
g - 1 > m_jw_j(g_j - 1) \Rightarrow g_j < (g - 1)/m_jw_j + 1
$$
Cela implique bien $g_j < g$, ce qui démontre (2).

\medskip\noindent
Pour $j \in J$, si $g_j > 0$, alors la contribution
$m_j(w_j(g_j - 1) - \frac{1}{2} a_{jj})$ au genre $g$
est $\geq -\frac{1}{2}m_ja_{jj}$, et l'on conclut immédiatement
que $m_j|a_{jj}| \leq 2(g - 1)$. Sinon, $a_{jj} = -kw_j$
pour un entier $k \geq 3$ (car $j \in J$), et l'on obtient
$$
m_jw_j(-1 + \frac{k}{2}) \leq g - 1
\Rightarrow
m_jw_j \leq \frac{2(g - 1)}{k - 2}
$$
En reportant ceci dans $a_{jj} = -km_jw_j$, on obtient
$$
m_j|a_{jj}| \leq 2(g - 1) \frac{k}{k - 2} \leq 6(g - 1)
$$
Cela démontre (3).

\medskip\noindent
L'assertion (4) résulte du lemme \ref{models-lemma-bound-neighbours} et de (3).
\end{proof}

\begin{lemma}
\label{models-lemma-bound-wm}
Fixons $g \geq 2$. Pour tout type numérique minimal $n, m_i, a_{ij}, w_i, g_i$
de genre $g$, on a $m_i|a_{ij}| \leq 768g$.
\end{lemma}

\begin{proof}
D'après le lemme \ref{models-lemma-bound-neighbours}, il suffit de montrer que
$m_i|a_{ii}| \leq 768g$ pour tout $i$.
Soit $J \subset \{1, \ldots, n\}$ l'ensemble des indices non $(-2)$,
comme dans le lemme \ref{models-lemma-bound-heart}. Observons que $J$ est
non vide puisque $g \geq 2$. De plus, $m_j|a_{jj}| \leq 6g$ pour $j \in J$
d'après ce lemme.

\medskip\noindent
Supposons que l'on ait $j \in J$ et une suite $i_1, \ldots, i_7$
d'indices $(-2)$ telle que $a_{ji_1}$ et $a_{i_1i_2}$,
$a_{i_2i_3}$, $a_{i_3i_4}$, $a_{i_4i_5}$, $a_{i_5i_6}$ et
$a_{i_6i_7}$ soient non nuls. Le
lemme \ref{models-lemma-bound-neighbours} montre alors
que $m_{i_1}w_{i_1} \leq 6g$ et $m_{i_1}a_{ji_1} \leq 6g$.
Comme $i_1$ est un indice $(-2)$, on a $a_{i_1i_1} = -2w_{i_1}$,
et l'on conclut que $m_{i_1}|a_{i_1i_1}| \leq 12g$.
En répétant l'argument, on conclut que
$m_{i_2}w_{i_2} \leq 12g$ et $m_{i_2}a_{i_1i_2} \leq 12g$.
Puis $m_{i_2}|a_{i_2i_2}| \leq 24g$, et ainsi de suite.
Finalement, on conclut que
$m_{i_k}|a_{i_ki_k}| \leq 2^k(6g) \leq 768g$ pour $k = 1, \ldots, 7$.

\medskip\noindent
Soit $I \subset \{1, \ldots, n\} \setminus J$ un sous-ensemble connexe maximal.
Autrement dit, il n'existe aucun sous-ensemble strict non vide
$I' \subset I$ tel que
$a_{i'i} = 0$ pour $i' \in I'$ et $i \in I \setminus I'$,
et $I$ est maximal pour cette propriété. En particulier, puisqu'un
type numérique est connexe par définition, on voit qu'il
existe $j \in J$ et $i \in I$ tels que $a_{ij} > 0$.
En considérant la classification
de tels $I$ dans la proposition \ref{models-proposition-classify-subgraphs}
et en utilisant le résultat du paragraphe précédent, on voit que
$w_i|a_{ii}| \leq 768g$ pour tout $i \in I$, sauf si $I$ est décrit dans le
lemme \ref{models-lemma-long} ou le
lemme \ref{models-lemma-Dn}.
On peut donc supposer que le lieu de non-annulation des $a_{ii'}$, $i, i' \in I$,
a l'une des deux formes suivantes :
$$
\xymatrix{
\bullet \ar@{-}[r] &
\bullet \ar@{-}[r] &
\bullet \ar@{..}[r] &
\bullet \ar@{-}[r] &
\bullet \ar@{-}[r] &
\bullet
}
$$
(laquelle possède 3 sous-cas détaillés dans le lemme \ref{models-lemma-long}),
ou bien
$$
\xymatrix{
\bullet \ar@{-}[r] &
\bullet \ar@{-}[r] &
\bullet \ar@{..}[r] &
\bullet \ar@{-}[r] &
\bullet \ar@{-}[r] \ar@{-}[d] &
\bullet \\
& & & & \bullet
}
$$
Nous allons démontrer la borne pour le premier sous-cas du
lemme \ref{models-lemma-long} et laisser les autres cas au lecteur (l'argument
y est presque exactement le même).

\medskip\noindent
Après réindexation, on peut supposer $I = \{1, \ldots, t\} \subset \{1, \ldots, n\}$
et qu'il existe un entier $w$ tel que
$$
w = w_1 = \ldots = w_t =
a_{12} = \ldots = a_{(t - 1) t} =
-\frac{1}{2} a_{i_1i_2} = \ldots = -\frac{1}{2} a_{(t - 1) t}
$$
Les égalités $a_{ii}m_i + \sum_{j \not = i} a_{ij}m_j = 0$ impliquent
que l'on a
$$
2m_2 \geq m_1 + m_3, \ldots, 2m_{t - 1} \geq m_{t - 2} + m_t
$$
On a égalité dans $2m_i \geq m_{i - 1} + m_{i + 1}$
si et seulement si $i$ ne « rencontre » aucun indice autre que
$i - 1$ et $i + 1$. Et si $i$ rencontre effectivement
un autre indice, alors cet indice appartient à $J$ (par maximalité de $I$).
En particulier, l'application
$\{1, \ldots, t\} \to \mathbf{Z}$, $i \mapsto m_i$, est concave.

\medskip\noindent
Soit $m = \max(m_i, i \in \{1, \ldots, t\})$. Alors
$m_i|a_{ii}| \leq 2mw$ pour $i \leq t$, et
notre but est de montrer que $2mw \leq 768g$.
Soient $s$, respectivement $s'$, dans $\{1, \ldots, t\}$, le
plus petit, respectivement le plus grand, indice tel que $m_s = m = m_{s'}$.
La concavité montre que $m_i = m$ pour $s \leq i \leq s'$.
Si $s > 1$, alors on n'a pas égalité dans
$2m_s \geq m_{s - 1} + m_{s + 1}$,
et l'on voit que $s$ rencontre un indice de $J$.
Dans ce cas, $2mw \leq 12g$ d'après le résultat du deuxième paragraphe
de la démonstration.
De même, si $s' < t$, alors $s'$ rencontre un indice de $J$
et l'on obtient également $2mw \leq 12g$.
Mais si $s = 1$ et $s' = t$, alors on conclut
que $a_{ij} = 0$ pour tout $j \in J$ et $i \in \{2, \ldots, t - 1\}$.
Or, puisque nous avons vu qu'il doit exister une paire $(i, j) \in I \times J$
telle que $a_{ij} > 0$, on conclut que cela se produit soit pour
$i = 1$, soit pour $i = t$, et l'on conclut $2mw \leq 12g$
de la même manière qu'auparavant (car $m_1 = m = m_t$ dans ce cas).
\end{proof}

\begin{proposition}
\label{models-proposition-bound-picard-group}
Soit $g \geq 2$. Pour tout type numérique $T$ de genre $g$
et tout nombre premier $\ell > 768g$, on a
$$
\dim_{\mathbf{F}_\ell} \Pic(T)[\ell] \leq g
$$
où $\Pic(T)$ est défini dans la définition \ref{models-definition-picard-group}.
Si $T$ est minimal, on a même
$$
\dim_{\mathbf{F}_\ell} \Pic(T)[\ell] \leq g_{top} \leq g
$$
où $g_{top}$ est défini dans la définition \ref{models-definition-top-genus}.
\end{proposition}

\begin{proof}
Supposons que $T$ soit donné par $n, m_i, a_{ij}, w_i, g_i$.
Si $T$ n'est pas minimal, alors il existe un indice $(-1)$.
Après avoir remplacé $T$ par un type équivalent, on peut supposer que
$n$ est un indice $(-1)$. En appliquant le lemme \ref{models-lemma-contract-picard-group},
on obtient $\Pic(T) \subset \Pic(T')$, où $T'$
est un type numérique de genre $g$ (lemme \ref{models-lemma-contract})
à $n - 1$ indices. On conclut donc par récurrence sur $n$,
pourvu que le lemme soit démontré pour les types numériques minimaux.

\medskip\noindent
Supposons que $T$ soit un type numérique minimal de genre $\geq 2$.
Observons que $g_{top} \leq g$ d'après le lemme \ref{models-lemma-genus-nonnegative}.
Si $A = (a_{ij})$, on a $\Pic(T) \subset \Coker(A)$ d'après le
lemme \ref{models-lemma-picard-T-and-A}. Il suffit donc de démontrer le lemme
pour $\Coker(A)$.
Le lemme \ref{models-lemma-bound-wm} montre que $m_i|a_{ij}| \leq 768g$ pour
tous $i, j$.
Le résultat découle donc du lemme \ref{models-lemma-recurring-symmetric-integer}.
\end{proof}







\section{Modèles}
\label{models-section-models}

\noindent
Dans ce chapitre, $R$ désignera un anneau de valuation discrète et $K$ son
corps des fractions. Si nécessaire, nous noterons $\pi \in R$ une
uniformisante et $k = R/(\pi)$ son corps résiduel.

\medskip\noindent
Soit $V$ un $K$-schéma algébrique
(Variétés, définition \ref{varieties-definition-algebraic-scheme}).
Par {\it modèle} de $V$, on entendra un morphisme plat de type
fini\footnote{Il est parfois utile d'autoriser les modèles qui sont
localement de type fini sur $R$, mais nous nous en préoccuperons
le moment venu.}
$X \to \Spec(R)$ muni d'un
isomorphisme $V \to X_K = X \times_{\Spec(R)} \Spec(K)$. Souvent, nous
identifierons $V$ et la fibre générique $X_K$ de $X$, et
écrirons simplement $V = X_K$.
La fibre spéciale est $X_k = X \times_{\Spec(R)} \Spec(k)$.
Un {\it morphisme de modèles $X \to X'$ de $V$}
est un morphisme $X \to X'$ de schémas sur $R$ qui induit
l'identité sur $V$.

\medskip\noindent
On dira que {\it $X$ est un modèle propre de $V$} si $X$ est un modèle
de $V$ et si le morphisme structural $X \to \Spec(R)$ est propre.
De même pour les modèles séparés, les modèles lisses et d'autres encore,
à ajouter ici.
On dira que {\it $X$ est un modèle régulier de $V$} si $X$ est un modèle
de $V$ et si $X$ est un schéma régulier.
De même pour les modèles normaux, les modèles réduits et d'autres encore,
à ajouter ici.

\medskip\noindent
Soit $R \subset R'$ une extension d'anneaux de valuation discrète
(Compléments d'algèbre, définition
\ref{more-algebra-definition-extension-discrete-valuation-rings}).
Celle-ci induit une extension $K'/K$ des corps des fractions.
Étant donné un schéma algébrique $V$ sur $K$, notons $V'$ le
changement de base $V \times_{\Spec(K)} \Spec(K')$. On a alors
un foncteur
$$
\text{modèles de }V\text{ sur }R
\longrightarrow
\text{modèles de }V'\text{ sur }R'
$$
qui envoie $X$ sur $X \times_{\Spec(R)} \Spec(R')$.

\begin{lemma}
\label{models-lemma-closure-is-model}
Soit $V_1 \to V_2$ une immersion fermée de schémas algébriques sur $K$.
Si $X_2$ est un modèle de $V_2$, alors l'image schématique
de $V_1 \to X_2$ est un modèle de $V_1$.
\end{lemma}

\begin{proof}
En utilisant
Morphismes, lemme \ref{morphisms-lemma-quasi-compact-scheme-theoretic-image} et
l'exemple \ref{morphisms-example-scheme-theoretic-image},
on se ramène à l'énoncé algébrique suivant.
Soit $A_1$ une $R$-algèbre de type fini, plate sur $R$.
Soit $A_1 \otimes_R K \to B_2$ une surjection. Alors
$A_2 = A_1 / \Ker(A_1 \to B_2)$ est une $R$-algèbre de type fini,
plate sur $R$, telle que $B_2 = A_2 \otimes_R K$.
Nous omettons la démonstration détaillée ; utiliser
Compléments d'algèbre, lemme \ref{more-algebra-lemma-dedekind-torsion-free-flat},
pour montrer que $A_2$ est plate.
\end{proof}

\begin{lemma}
\label{models-lemma-normalization}
Soit $X$ un modèle d'une variété géométriquement normale $V$ sur $K$.
Alors la normalisation $\nu : X^\nu \to X$ est finie et
le changement de base de $X^\nu$ à la complétion $R^\wedge$
est la normalisation du changement de base de $X$. De plus, pour
tout $x \in X^\nu$, la complétion de $\mathcal{O}_{X^\nu, x}$
est normale.
\end{lemma}

\begin{proof}
Observons que $R^\wedge$ est un anneau de valuation discrète
(Compléments d'algèbre, lemme \ref{more-algebra-lemma-completion-dvr}).
Posons $Y = X \times_{\Spec(R)} \Spec(R^\wedge)$.
Puisque $R^\wedge$ est un anneau de valuation discrète, on voit que
$$
Y \setminus Y_k =
Y \times_{\Spec(R^\wedge)} \Spec(K^\wedge) =
V \times_{\Spec(K)} \Spec(K^\wedge)
$$
où $K^\wedge$ est le corps des fractions de $R^\wedge$.
Comme $V$ est géométriquement normale, on trouve qu'il s'agit
d'un schéma normal. La première partie du lemme résulte donc de
Résolution des surfaces, lemme \ref{resolve-lemma-normalization-completion}.

\medskip\noindent
Pour démontrer la seconde partie, on peut supposer $X$ et $Y$ normaux
(par la première partie). Si $x$ appartient à la fibre générique, alors
$\mathcal{O}_{X, x} = \mathcal{O}_{V, x}$ est un anneau local normal
essentiellement de type fini sur un corps. Un tel anneau est
excellent (Compléments d'algèbre, proposition
\ref{more-algebra-proposition-ubiquity-excellent}).
Si $x$ est un point de la fibre spéciale d'image $y \in Y$, alors
$\mathcal{O}_{X, x}^\wedge = \mathcal{O}_{Y, y}^\wedge$
par Résolution des surfaces, lemme \ref{resolve-lemma-iso-completions}.
Dans ce cas, $\mathcal{O}_{Y, y}$ est un domaine local normal excellent,
par la même référence que précédemment, puisque $R^\wedge$ est excellent.
Si $B$ est un domaine local normal excellent, alors la complétion
$B^\wedge$ est normale (car $B \to B^\wedge$ est régulier et le
lemme \ref{more-algebra-lemma-normal-goes-up} de Compléments d'algèbre s'applique).
Ceci achève la démonstration.
\end{proof}

\begin{lemma}
\label{models-lemma-regular}
Soit $X$ un modèle d'une courbe lisse $C$ sur $K$. Alors
il existe une résolution des singularités de $X$,
et toute résolution est un modèle de $C$.
\end{lemma}

\begin{proof}
Vérifions que la condition (4) du théorème de Lipman
(Résolution des surfaces, théorème \ref{resolve-theorem-resolve}) est satisfaite.
Cela découle clairement du lemme \ref{models-lemma-normalization},
à l'exception de l'assertion selon laquelle $X^\nu$ ne possède qu'un nombre fini de
points singuliers. Pour le voir, on peut utiliser le fait que $R$ est J-2 d'après
Compléments d'algèbre, proposition \ref{more-algebra-proposition-ubiquity-J-2},
de sorte que le lieu non singulier est ouvert dans $X^\nu$.
Comme $X^\nu$ est normal de dimension $\leq 2$, les points singuliers
sont fermés ; le fait que le lieu singulier soit fermé signifie donc
qu'ils sont en nombre fini (puisque $X$ est quasi-compact).
Observons que toute résolution de $X$ est une modification de $X$
(Résolution des surfaces, définition \ref{resolve-definition-resolution}).
Celle-ci est un isomorphisme au-dessus du lieu normal de $X$ d'après Variétés, lemme
\ref{varieties-lemma-modification-normal-iso-over-codimension-1}.
Puisque l'ensemble des points normaux contient
$C = X_K$, on conclut que toute résolution est un modèle de $C$.
\end{proof}

\begin{definition}
\label{models-definition-minimal-model}
Soit $C$ une courbe projective lisse sur $K$ telle que
$H^0(C, \mathcal{O}_C) = K$. Un {\it modèle minimal}
est un modèle propre régulier $X$ de $C$ tel que
$X$ ne contienne aucune courbe exceptionnelle de première espèce
(Résolution des surfaces, section \ref{resolve-section-minus-one}).
\end{definition}

\noindent
En toute rigueur, un tel objet devrait être appelé modèle propre régulier minimal,
voire modèle projectif régulier relativement minimal. Mais tant que
nous nous limitons aux modèles sur des anneaux de valuation discrète (comme nous le ferons
dans ce chapitre), aucune confusion ne devrait en résulter.

\medskip\noindent
Les modèles minimaux existent toujours
(proposition \ref{models-proposition-exists-minimal-model}) et sont uniques
lorsque le genre est $> 0$ (lemme \ref{models-lemma-minimal-model-unique}).

\begin{lemma}
\label{models-lemma-pre-exists-minimal-model}
\begin{slogan}
Un modèle propre régulier d'une courbe s'obtient par éclatements successifs
à partir d'un modèle minimal
\end{slogan}
Soit $C$ une courbe projective lisse sur $K$ telle que
$H^0(C, \mathcal{O}_C) = K$. Si $X$ est un modèle propre régulier
de $C$, alors il existe une suite de morphismes
$$
X = X_m \to X_{m - 1} \to \ldots \to X_1 \to X_0
$$
de modèles propres réguliers de $C$, telle que chaque morphisme soit une
contraction d'une courbe exceptionnelle de première espèce et que
$X_0$ soit un modèle minimal.
\end{lemma}

\begin{proof}
D'après Résolution des surfaces, lemme \ref{resolve-lemma-regular-dim-2-projective},
$X$ est projectif sur $R$. Par conséquent, $X$ possède un faisceau inversible
ample d'après
Compléments sur les morphismes, lemme \ref{more-morphisms-lemma-projective}
(nous l'utiliserons ci-dessous).
Soit $E \subset X$ une courbe exceptionnelle de première espèce.
Voir Résolution des surfaces, section \ref{resolve-section-minus-one}.
D'après Résolution des surfaces, lemme \ref{resolve-lemma-contract-ample},
on peut contracter $E$ par un morphisme $X \to X'$ tel que $X'$ soit
régulier et projectif sur $R$. Il est clair que le nombre de
composantes irréductibles de $X'_k$ est inférieur d'une unité au
nombre de composantes irréductibles de $X_k$. On ne peut donc
effectuer qu'un nombre fini de ces contractions avant
d'obtenir un modèle minimal.
\end{proof}

\begin{proposition}
\label{models-proposition-exists-minimal-model}
Soit $C$ une courbe projective lisse sur $K$ telle que
$H^0(C, \mathcal{O}_C) = K$. Il existe un modèle minimal.
\end{proposition}

\begin{proof}
Choisissons une immersion fermée $C \to \mathbf{P}^n_K$. Soit
$X$ l'image schématique de $C \to \mathbf{P}^n_R$.
Alors $X \to \Spec(R)$ est un modèle projectif de $C$ d'après le
lemme \ref{models-lemma-closure-is-model}.
D'après le lemme \ref{models-lemma-regular}, il existe une résolution
des singularités $X' \to X$, et $X'$ est un modèle de $C$.
Alors $X' \to \Spec(R)$ est propre comme composé de morphismes propres.
On peut alors appliquer le lemme \ref{models-lemma-pre-exists-minimal-model}
pour obtenir un modèle minimal.
\end{proof}





\section{La géométrie d'un modèle régulier}
\label{models-section-special-fibre}

\noindent
Dans cette section, nous décrivons la géométrie d'un modèle propre régulier $X$ d'une
courbe projective lisse $C$ sur $K$ telle que $H^0(C, \mathcal{O}_C) = K$.

\begin{lemma}
\label{models-lemma-divisor-special-fiber}
Soit $X$ un modèle régulier d'une courbe lisse $C$ sur $K$.
\begin{enumerate}
\item la fibre spéciale $X_k$ est un diviseur de Cartier effectif sur $X$,
\item chaque composante irréductible $C_i$ de $X_k$ est un diviseur de
Cartier effectif sur $X$,
\item $X_k = \sum m_i C_i$ (somme de diviseurs de Cartier effectifs),
où $m_i$ est la multiplicité de $C_i$ dans $X_k$,
\item $\mathcal{O}_X(X_k) \cong \mathcal{O}_X$.
\end{enumerate}
\end{lemma}

\begin{proof}
Rappelons que $R$ est un anneau de valuation discrète d'uniformisante $\pi$
et de corps résiduel $k = R/(\pi)$. Puisque $X \to \Spec(R)$ est plat,
l'élément $\pi$ est un non-diviseur de zéro localement sur les ouverts affines de $X$
(voir Compléments d'algèbre, lemme
\ref{more-algebra-lemma-dedekind-torsion-free-flat}). Ainsi,
si $U = \Spec(A) \subset X$ est un ouvert affine, alors
$$
X_k \cap U = U_k = \Spec(A \otimes_R k) = \Spec(A/\pi A)
$$
et $\pi$ est un non-diviseur de zéro dans $A$.
Par conséquent, $X_k = V(\pi)$ est un diviseur de Cartier effectif d'après
Diviseurs, lemme \ref{divisors-lemma-characterize-effective-Cartier-divisor}.
Cela démontre (1).

\medskip\noindent
La discussion précédente montre que la paire $(\mathcal{O}_X(X_k), 1)$
est isomorphe à la paire $(\mathcal{O}_X, \pi)$, ce qui démontre (4).

\medskip\noindent
D'après Diviseurs, lemme \ref{divisors-lemma-effective-Cartier-divisor-is-a-sum},
il existe des diviseurs de Cartier effectifs intègres deux à deux distincts
$D_i \subset X$ et des entiers $a_i \geq 0$ tels que $X_k = \sum a_i D_i$.
On peut éliminer les diviseurs $D_i$ tels que $a_i = 0$. Il est alors
clair (d'après la définition de l'addition des diviseurs de Cartier
effectifs) que $X_k = \bigcup D_i$ ensemblistement. Ainsi, les $C_i = D_i$
sont les composantes irréductibles de $X_k$, ce qui démontre (2).
Soit $\xi_i$ le point générique de $C_i$.
Alors $\mathcal{O}_{X, \xi_i}$ est un anneau de valuation discrète
(Diviseurs, lemme \ref{divisors-lemma-integral-effective-Cartier-divisor-dvr}).
L'uniformisante $\pi_i \in \mathcal{O}_{X, \xi_i}$ est une équation locale
de $C_i$, et l'image de $\pi$ est une équation locale de $X_k$.
Comme $X_k = \sum a_i C_i$, on voit que $\pi$ et $\pi_i^{a_i}$
engendrent le même idéal dans $\mathcal{O}_{X, \xi_i}$.
D'autre part, la multiplicité de $C_i$ dans $X_k$ est
$$
m_i =
\text{longueur}_{\mathcal{O}_{C_i, \xi_i}} \mathcal{O}_{X_k, \xi_i} =
\text{longueur}_{\mathcal{O}_{C_i, \xi_i}} \mathcal{O}_{X, \xi_i}/(\pi) =
\text{longueur}_{\mathcal{O}_{C_i, \xi_i}} \mathcal{O}_{X, \xi_i}/(\pi_i^{a_i}) =
a_i
$$
Voir Homologie de Chow, définition
\ref{chow-definition-cycle-associated-to-closed-subscheme}.
Ainsi $a_i = m_i$, ce qui démontre (3).
\end{proof}

\begin{lemma}
\label{models-lemma-gorenstein}
Soit $X$ un modèle régulier d'une courbe lisse $C$ sur $K$. Alors
\begin{enumerate}
\item $X \to \Spec(R)$ est un morphisme de Gorenstein de dimension relative $1$,
\item chacune des composantes irréductibles $C_i$ de $X_k$ est de Gorenstein.
\end{enumerate}
\end{lemma}

\begin{proof}
Comme $X \to \Spec(R)$ est plat, pour démontrer (1),
il suffit de montrer que les fibres sont de Gorenstein
(Dualité pour les schémas, lemme \ref{duality-lemma-gorenstein-morphism}).
La fibre générique est une courbe lisse, qui est régulière et donc de Gorenstein
(Dualité pour les schémas, lemme \ref{duality-lemma-regular-gorenstein}).
Pour la fibre spéciale $X_k$, on utilise le fait qu'elle est un diviseur de
Cartier effectif sur un schéma régulier (donc de Gorenstein), et qu'elle est donc
de Gorenstein, par exemple d'après Complexes dualisants, lemme
\ref{dualizing-lemma-gorenstein-divide-by-nonzerodivisor}.
Les courbes $C_i$ sont de Gorenstein par le même argument.
\end{proof}

\begin{situation}
\label{models-situation-regular-model}
Soit $R$ un anneau de valuation discrète de corps des fractions $K$,
de corps résiduel $k$ et d'uniformisante $\pi$.
Soit $C$ une courbe projective lisse sur $K$ telle que $H^0(C, \mathcal{O}_C) = K$.
Soit $X$ un modèle propre régulier de $C$.
Soient $C_1, \ldots, C_n$ les composantes irréductibles de la fibre
spéciale $X_k$. Écrivons $X_k = \sum m_i C_i$ comme dans le
lemme \ref{models-lemma-divisor-special-fiber}.
\end{situation}

\begin{lemma}
\label{models-lemma-regular-model-connected}
Dans la situation \ref{models-situation-regular-model}, la fibre spéciale $X_k$ est connexe.
\end{lemma}

\begin{proof}
C'est une conséquence de Compléments sur les morphismes, lemme
\ref{more-morphisms-lemma-geometrically-connected-fibres-towards-normal}.
\end{proof}

\begin{lemma}
\label{models-lemma-regular-model-pic}
Dans la situation \ref{models-situation-regular-model}, il existe une suite exacte
$$
0 \to \mathbf{Z} \to \mathbf{Z}^{\oplus n} \to
\Pic(X) \to \Pic(C) \to 0
$$
où la première application envoie $1$ sur $(m_1, \ldots, m_n)$ et la seconde
envoie le $i$-ème vecteur de base sur $\mathcal{O}_X(C_i)$.
\end{lemma}

\begin{proof}
Observons que $C \subset X$ est un sous-schéma ouvert. L'application de restriction
$\Pic(X) \to \Pic(C)$ est surjective d'après
Diviseurs, lemme \ref{divisors-lemma-extend-invertible-module}.
Soit $\mathcal{L}$ un $\mathcal{O}_X$-module inversible
tel qu'il existe un isomorphisme $s : \mathcal{O}_C \to \mathcal{L}|_C$.
Alors $s$ est une section méromorphe régulière de $\mathcal{L}$,
et l'on voit que $\text{div}_\mathcal{L}(s) = \sum a_i C_i$
pour certains $a_i \in \mathbf{Z}$
(Diviseurs, définition \ref{divisors-definition-divisor-invertible-sheaf}).
D'après Diviseurs, lemme \ref{divisors-lemma-normal-c1-injective}
(et le fait que $X$ est normal),
on conclut que $\mathcal{L} = \mathcal{O}_X(\sum a_iC_i)$.
Enfin, supposons que $\mathcal{O}_X(\sum a_i C_i) \cong \mathcal{O}_X$.
Il existe alors un élément $g$ du corps des fonctions de $X$
tel que $\text{div}_X(g) = \sum a_i C_i$. En particulier, la fonction rationnelle
$g$ n'a ni zéro ni pôle sur la fibre générique $C$ de $X$.
Comme $C$ est un schéma normal, ceci implique $g \in H^0(C, \mathcal{O}_C) = K$.
Ainsi $g = \pi^a u$ pour un certain $a \in \mathbf{Z}$ et un certain $u \in R^*$.
On conclut que $\text{div}_X(g) = a \sum m_i C_i$, ce qui achève
la démonstration.
\end{proof}

\noindent
Dans la situation \ref{models-situation-regular-model}, pour tout
$\mathcal{O}_X$-module inversible $\mathcal{L}$ et tout $i$, on obtient un entier
$$
\deg(\mathcal{L}|_{C_i}) =
\chi(C_i, \mathcal{L}|_{C_i}) - \chi(C_i, \mathcal{O}_{C_i})
$$
en prenant le degré de la restriction de $\mathcal{L}$ à $C_i$
relativement au corps de base $k$\footnote{Observons qu'il peut arriver
que le corps $\kappa_i = H^0(C_i, \mathcal{O}_{C_i})$ soit strictement plus grand
que $k$. Dans ce cas, tout module inversible sur $C_i$ a un
degré (au sens défini ci-dessus) divisible par $[\kappa_i : k]$.},
comme dans Variétés, section \ref{varieties-section-divisors-curves}.

\begin{lemma}
\label{models-lemma-intersection-pairing}
Dans la situation \ref{models-situation-regular-model}, étant donnés un
$\mathcal{O}_X$-module inversible $\mathcal{L}$ et
$a = (a_1, \ldots, a_n) \in \mathbf{Z}^{\oplus n}$, on définit
$$
\langle a, \mathcal{L} \rangle = \sum a_i\deg(\mathcal{L}|_{C_i})
$$
Alors $\langle , \rangle$ est bilinéaire et, pour
$b = (b_1, \ldots, b_n) \in \mathbf{Z}^{\oplus n}$, on a
$$
\left\langle a, \mathcal{O}_X(\sum b_i C_i) \right\rangle =
\left\langle b, \mathcal{O}_X(\sum a_i C_i) \right\rangle
$$
\end{lemma}

\begin{proof}
La bilinéarité résulte immédiatement de la définition et de
Variétés, lemme \ref{varieties-lemma-degree-tensor-product}.
Pour démontrer la symétrie, il suffit de supposer que $a$ et $b$ sont des
vecteurs de la base canonique de $\mathbf{Z}^{\oplus n}$.
Il suffit donc de démontrer que
$$
\deg(\mathcal{O}_X(C_j)|_{C_i}) = \deg(\mathcal{O}_X(C_i)|_{C_j})
$$
pour tous $1 \leq i, j \leq n$. Si $i = j$, il n'y a rien à démontrer.
Si $i \not = j$, alors la section canonique $1$ de $\mathcal{O}_X(C_j)$
se restreint en une section non nulle (donc régulière) de $\mathcal{O}_X(C_j)|_{C_i}$,
dont le schéma des zéros est exactement $C_i \cap C_j$ (intersection schématique).
Autrement dit, $C_i \cap C_j$ est un diviseur de Cartier effectif sur $C_i$,
et
$$
\deg(\mathcal{O}_X(C_j)|_{C_i}) = \deg(C_i \cap C_j)
$$
d'après Variétés, lemme \ref{varieties-lemma-degree-effective-Cartier-divisor}.
Par symétrie, on obtient la même (!) formule pour l'autre membre,
ce qui achève la démonstration.
\end{proof}

\noindent
Dans la situation \ref{models-situation-regular-model}, il est souvent commode de considérer
$\mathbf{Z}^{\oplus n}$ comme le groupe abélien libre sur l'ensemble
$\{C_1, \ldots, C_n\}$. Nous noterons un élément de ce groupe
sous la forme $\sum a_i C_i$ ; nous le considérons ici comme une somme formelle, bien que,
de manière équivalente, on puisse (et nous le ferons parfois)
considérer une telle somme comme un diviseur de Weil sur $X$
à support dans la fibre spéciale $X_k$. Le
lemme \ref{models-lemma-intersection-pairing}
nous permet alors de définir une forme bilinéaire symétrique $(\ \cdot\ )$
sur ce groupe abélien libre par la règle
\begin{equation}
\label{models-equation-form}
\left(\sum a_i C_i \cdot \sum b_j C_j\right) =
\left\langle a, \mathcal{O}_X(\sum b_j C_j) \right\rangle =
\left\langle b, \mathcal{O}_X(\sum a_i C_i) \right\rangle
\end{equation}
Nous allons établir quelques propriétés de cette forme bilinéaire.

\begin{lemma}
\label{models-lemma-properties-form}
Dans la situation \ref{models-situation-regular-model}, la forme bilinéaire symétrique
(\ref{models-equation-form}) possède les propriétés suivantes :
\begin{enumerate}
\item $(C_i \cdot C_j) \geq 0$ si $i \not = j$, avec égalité si et seulement
si $C_i \cap C_j = \emptyset$,
\item $(\sum m_i C_i \cdot C_j) = 0$,
\item il n'existe aucun sous-ensemble propre non vide $I \subset \{1, \ldots, n\}$
tel que $(C_i \cdot C_j) = 0$ pour $i \in I$, $j \not \in I$,
\item $(\sum a_i C_i \cdot \sum a_i C_i) \leq 0$, avec égalité si et
seulement s'il existe $q \in \mathbf{Q}$ tel que $a_i = qm_i$
pour $i = 1, \ldots, n$.
\end{enumerate}
\end{lemma}

\begin{proof}
Dans la démonstration du lemme \ref{models-lemma-intersection-pairing}, nous avons vu que
$(C_i \cdot C_j) = \deg(C_i \cap C_j)$ si $i \not = j$. Ce nombre est
$\geq 0$ et il est $> 0 $ si et seulement si $C_i \cap C_j \not = \emptyset$.
Ceci démontre (1).

\medskip\noindent
Démonstration de (2). Cela résulte du fait que, d'après le
lemme \ref{models-lemma-divisor-special-fiber}, le faisceau inversible associé à
$\sum m_i C_i$ est trivial et que le faisceau trivial est de degré zéro.

\medskip\noindent
Démonstration de (3). Cette assertion exprime le fait que $X_k$ est connexe
(lemme \ref{models-lemma-regular-model-connected})
au moyen de la description des produits d'intersection donnée dans la démonstration de (1).

\medskip\noindent
La partie (4) résulte de (1), (2) et (3) par le
lemme \ref{models-lemma-recurring-symmetric-real}.
\end{proof}

\begin{lemma}
\label{models-lemma-multiple-fibre-normal-bundle}
Dans la situation \ref{models-situation-regular-model}, posons $d = \gcd(m_1, \ldots, m_n)$
et soit $D = \sum (m_i/d)C_i$ considéré comme un diviseur de Cartier effectif.
Alors $\mathcal{O}_X(D)$ est d'ordre divisant $d$ dans $\Pic(X)$,
et $\mathcal{C}_{D/X}$ est un $\mathcal{O}_D$-module inversible
d'ordre divisant $d$ dans $\Pic(D)$.
\end{lemma}

\begin{proof}
On a
$$
\mathcal{O}_X(D)^{\otimes d} = \mathcal{O}_X(dD) =
\mathcal{O}_X(X_k) = \mathcal{O}_X
$$
d'après le lemme \ref{models-lemma-divisor-special-fiber}.
On conclut puisque $\mathcal{C}_{D/X}$ est l'image inverse de
$\mathcal{O}_X(-D)$.
\end{proof}

\begin{lemma}
\label{models-lemma-regular-model-field}
\begin{reference}
\cite[Lemma 2.6]{Artin-Winters}
\end{reference}
Dans la situation \ref{models-situation-regular-model}, posons $d = \gcd(m_1, \ldots, m_n)$.
Soit $D = \sum (m_i/d) C_i$ considéré comme un diviseur de Cartier effectif. Alors il existe
une suite de diviseurs de Cartier effectifs
$$
(X_k)_{red} = Z_0 \subset Z_1 \subset \ldots \subset Z_m = D
$$
telle que $Z_j = Z_{j - 1} + C_{i_j}$ pour un certain $i_j \in \{1, \ldots, n\}$
pour $j = 1, \ldots, m$, et telle que $H^0(Z_j, \mathcal{O}_{Z_j})$
soit un corps fini sur $k$ pour $j = 0, \ldots m$.
\end{lemma}

\begin{proof}
Le réduit $D_{red} = (X_k)_{red} = \sum C_i$ est connexe
(lemme \ref{models-lemma-regular-model-connected}) et propre sur $k$. Par conséquent,
$H^0(D_{red}, \mathcal{O})$ est un corps et une extension finie de
$k$ d'après Variétés, lemme
\ref{varieties-lemma-proper-geometrically-reduced-global-sections}.
Le résultat est donc vrai pour $Z_0 = D_{red} = (X_k)_{red}$.
Supposons que nous ayons déjà construit
$$
(X_k)_{red} = Z_0 \subset Z_1 \subset \ldots \subset Z_t \subset D
$$
avec $Z_j = Z_{j - 1} + C_{i_j}$ pour un certain $i_j \in \{1, \ldots, n\}$
pour $j = 1, \ldots, t$, et de sorte que $H^0(Z_j, \mathcal{O}_{Z_j})$
soit un corps fini sur $k$ pour $j = 0, \ldots, t$.
Écrivons $Z_t = \sum a_i C_i$ avec $1 \leq a_i \leq m_i/d$.
Si $a_i = m_i/d$ pour tout $i$, alors $Z_t = D$ et le lemme est démontré.
Sinon, $a_i < m_i/d$ pour un certain $i$, et il s'ensuit que
$(Z_t \cdot Z_t) < 0$ d'après le lemme \ref{models-lemma-properties-form}. Cela signifie que
$(D - Z_t \cdot Z_t) > 0$, puisque $(D \cdot Z_t) = 0$ d'après ce lemme.
On peut donc trouver un $i$ tel que $a_i < m_i/d$ et
$(C_i \cdot Z_t) > 0$. Posons $Z_{t + 1} = Z_t + C_i$ et $i_{t + 1} = i$.
Considérons la suite exacte courte
$$
0 \to \mathcal{O}_X(-Z_t)|_{C_i} \to \mathcal{O}_{Z_{t + 1}} \to
\mathcal{O}_{Z_t} \to 0
$$
de Diviseurs, lemme \ref{divisors-lemma-ses-add-divisor}.
Par notre choix de $i$, on voit que
$\mathcal{O}_X(-Z_t)|_{C_i}$ est un faisceau inversible de degré négatif
sur la courbe propre $C_i$ ; il ne possède donc aucune section globale non nulle
(Variétés, lemme \ref{varieties-lemma-check-invertible-sheaf-trivial}).
On conclut que $H^0(\mathcal{O}_{Z_{t + 1}}) \subset H^0(\mathcal{O}_{Z_t})$
est un corps (cela est clair, mais résulte également de
Algèbre, lemme \ref{algebra-lemma-integral-under-field})
et une extension finie de $k$. Nous avons ainsi prolongé la suite.
Comme le procédé doit s'arrêter, par exemple parce que $t \leq \sum (m_i/d - 1)$,
ceci achève la démonstration.
\end{proof}

\begin{lemma}
\label{models-lemma-regular-model-genus}
\begin{reference}
\cite[Lemma 2.6]{Artin-Winters}
\end{reference}
Dans la situation \ref{models-situation-regular-model}, posons $d = \gcd(m_1, \ldots, m_n)$.
Soit $D = \sum (m_i/d) C_i$ considéré comme un diviseur de Cartier effectif sur $X$. Alors
$$
1 - g_C = d [\kappa : k] (1 - g_D)
$$
où $g_C$ est le genre de $C$, $g_D$ le genre de $D$, et
$\kappa = H^0(D, \mathcal{O}_D)$.
\end{lemma}

\begin{proof}
D'après le lemme \ref{models-lemma-regular-model-field}, $\kappa$ est un corps
et une extension finie de $k$. Comme on a aussi $H^0(C, \mathcal{O}_C) = K$,
on voit que les genres de $C$ et de $D$ sont définis (voir
Courbes algébriques, définition \ref{curves-definition-genus}) et que
$g_C = \dim_K H^1(C, \mathcal{O}_C)$ et
$g_D = \dim_\kappa H^1(D, \mathcal{O}_D)$.
D'après Catégories dérivées de schémas, lemme
\ref{perfect-lemma-chi-locally-constant-geometric},
on a
$$
1 - g_C = \chi(C, \mathcal{O}_C) =
\chi(X_k, \mathcal{O}_{X_k}) = \dim_k H^0(X_k, \mathcal{O}_{X_k})
- \dim_k H^1(X_k, \mathcal{O}_{X_k})
$$
Nous affirmons que
$$
\chi(X_k, \mathcal{O}_{X_k}) = d \chi(D, \mathcal{O}_D)
$$
Cela démontrera le lemme, car
$$
\chi(D, \mathcal{O}_D) =
\dim_k H^0(D, \mathcal{O}_D) - \dim_k H^1(D, \mathcal{O}_D) =
[\kappa : k](1 - g_D)
$$
Observons que $X_k = dD$ comme diviseur de Cartier effectif.
Pour démontrer l'affirmation, montrons par récurrence sur $1 \leq r \leq d$ que
$\chi(rD, \mathcal{O}_{rD}) = r \chi(D, \mathcal{O}_D)$.
Le cas initial $r = 1$ est trivial. Si $1 \leq r < d$, considérons
la suite exacte courte
$$
0 \to \mathcal{O}_X(rD)|_D \to \mathcal{O}_{(r + 1)D} \to
\mathcal{O}_{rD} \to 0
$$
de Diviseurs, lemme \ref{divisors-lemma-ses-add-divisor}. Par additivité
des caractéristiques d'Euler
(Variétés, lemme \ref{varieties-lemma-euler-characteristic-additive}),
il suffit de démontrer que
$\chi(D, \mathcal{O}_X(rD)|_D) = \chi(D, \mathcal{O}_D)$.
C'est vrai parce que $\mathcal{O}_X(rD)|_D$ est un élément de torsion
de $\Pic(D)$ (lemme \ref{models-lemma-multiple-fibre-normal-bundle})
et parce que le degré d'un fibré en droites est additif
(Variétés, lemme \ref{varieties-lemma-degree-tensor-product}),
donc nul pour les faisceaux inversibles de torsion.
\end{proof}

\begin{lemma}
\label{models-lemma-exceptional-curves-dont-meet}
Dans la situation \ref{models-situation-regular-model}, soient deux indices $i, j$
tels que $C_i$ et $C_j$ soient des courbes exceptionnelles de première espèce
et que $C_i \cap C_j \not = \emptyset$. Alors
$n = 2$, $m_1 = m_2 = 1$, $C_1 \cong \mathbf{P}^1_k$,
$C_2 \cong \mathbf{P}^1_k$, $C_1$ et $C_2$ se rencontrent en un point $k$-rationnel,
et $C$ est de genre $0$.
\end{lemma}

\begin{proof}
Choisissons des isomorphismes $C_i = \mathbf{P}^1_{\kappa_i}$ et
$C_j = \mathbf{P}^1_{\kappa_j}$. Le schéma $C_i \cap C_j$
est un diviseur de Cartier effectif non vide tant dans $C_i$ que dans $C_j$.
Ainsi,
$$
(C_i \cdot C_j) = \deg(C_i \cap C_j) \geq \max([\kappa_i: k], [\kappa_j : k])
$$
La première égalité a été démontrée dans la preuve du
lemme \ref{models-lemma-intersection-pairing}.
D'autre part, l'auto-intersection $(C_i \cdot C_i)$ est égale
au degré de $\mathcal{O}_X(C_i)$ sur $C_i$, qui vaut $-[\kappa_i : k]$
puisque $C_i$ est une courbe exceptionnelle de première espèce. De même pour
$C_j$. D'après le lemme \ref{models-lemma-properties-form},
$$
0 \geq (C_i + C_j)^2 = -[\kappa_i : k] + 2(C_i \cdot C_j) - [\kappa_j : k]
$$
Ceci implique que $[\kappa_i : k] = \deg(C_i \cap C_j) = [\kappa_j : k]$
et que $(C_i + C_j)^2 = 0$. En appliquant de nouveau le lemme,
on conclut que $n = 2$, $\{1, 2\} = \{i, j\}$ et $m_1 = m_2$.
De plus, l'intersection schématique $C_i \cap C_j$ consiste en
un unique point $p$ de corps résiduel $\kappa$, et
$\kappa_i \to \kappa \leftarrow \kappa_j$ sont des isomorphismes.
Soit $D = C_1 + C_2$ considéré comme un diviseur de Cartier effectif sur $X$.
Observons que $D$ est la réunion schématique de $C_1$ et $C_2$
(Diviseurs, lemme \ref{divisors-lemma-sum-effective-Cartier-divisors-union}) ;
on a donc une suite exacte courte
$$
0 \to \mathcal{O}_D \to \mathcal{O}_{C_1} \oplus \mathcal{O}_{C_2} \to
\mathcal{O}_p \to 0
$$
d'après Morphismes, lemme \ref{morphisms-lemma-scheme-theoretic-union}.
Puisque nous connaissons la cohomologie de $C_i \cong \mathbf{P}^1_\kappa$
(Cohomologie des schémas, lemme
\ref{coherent-lemma-cohomology-projective-space-over-ring}),
la suite exacte longue de cohomologie donne
$H^0(D, \mathcal{O}_D) = \kappa$ et
$H^1(D, \mathcal{O}_D) = 0$. D'après le lemme \ref{models-lemma-regular-model-genus},
on conclut que
$$
1 - g_C = d[\kappa : k](1 - 0)
$$
où $d = m_1 = m_2$. Il s'ensuit que $g_C = 0$, $d = m_1 = m_2 = 1$
et $\kappa = k$.
\end{proof}





\section{Unicité du modèle minimal}
\label{models-section-uniqueness}

\noindent
Si le genre de la fibre générique est positif, les modèles minimaux sont uniques
(lemme \ref{models-lemma-minimal-model-unique}) et possèdent par conséquent une propriété
universelle convenable (lemme \ref{models-lemma-minimal-model-mapping-property}).

\begin{lemma}
\label{models-lemma-minimal-model-unique}
Soit $C$ une courbe projective lisse sur $K$ telle que
$H^0(C, \mathcal{O}_C) = K$ et de genre $> 0$.
Il existe un unique modèle minimal de $C$.
\end{lemma}

\begin{proof}
Nous avons déjà démontré la partie difficile du lemme, à savoir l'existence
d'un modèle minimal (dont la preuve repose sur la
résolution des singularités de surfaces) ; voir la
proposition \ref{models-proposition-exists-minimal-model}.
Pour démontrer l'unicité, supposons que $X$ et $Y$ soient deux
modèles minimaux. D'après
Résolution des surfaces, lemme \ref{resolve-lemma-birational-regular-surfaces},
il existe un diagramme de $S$-morphismes
$$
X = X_0 \leftarrow X_1 \leftarrow \ldots \leftarrow X_n = Y_m
\to \ldots \to Y_1 \to Y_0 = Y
$$
où chaque morphisme est un éclatement en un point fermé. La
fibre exceptionnelle du morphisme $X_n \to X_{n - 1}$ est une
courbe exceptionnelle de première espèce $E$. Nous affirmons que $E$ est
contractée en un point par le morphisme $X_n = Y_m \to Y$.
Si tel est le cas, alors $X_n \to Y$ se factorise par $X_{n - 1}$ d'après
Résolution des surfaces, lemme \ref{resolve-lemma-factor-through-contraction}.
Dans ce cas, le morphisme $X_{n - 1} \to Y$ est encore une suite de
contractions de courbes exceptionnelles d'après
Résolution des surfaces, lemme
\ref{resolve-lemma-proper-birational-regular-surfaces}.
On conclut donc par récurrence sur $n$. (Le cas initial $n = 0$ signifie
qu'il existe une suite de contractions
$X = Y_m \to \ldots \to Y_1 \to Y_0 = Y$
aboutissant à $Y$. Toutefois, comme $X$ est un modèle minimal, il ne contient
aucune courbe exceptionnelle de première espèce ; ainsi $m = 0$ et $X = Y$.)

\medskip\noindent
Démonstration de l'affirmation. Nous allons montrer par récurrence sur $m$ que toute courbe
exceptionnelle de première espèce $E \subset Y_m$ est envoyée sur un point
par le morphisme $Y_m \to Y$. Si $m = 0$, c'est clair puisque
$Y$ est un modèle minimal. Si $m > 0$, alors soit
$Y_m \to Y_{m - 1}$ contracte $E$ (et nous avons terminé), soit
la fibre exceptionnelle $E' \subset Y_m$ de $Y_m \to Y_{m - 1}$
est une seconde courbe exceptionnelle de première espèce.
Puisque $E$ et $E'$ sont toutes deux des composantes irréductibles de la fibre
spéciale, et puisque $g_C > 0$ par hypothèse, on conclut que
$E \cap E' = \emptyset$ d'après le
lemme \ref{models-lemma-exceptional-curves-dont-meet}.
L'image de $E$ dans $Y_{m - 1}$ est alors une courbe
exceptionnelle de première espèce (cela est clair, car le morphisme
$Y_m \to Y_{m - 1}$ est un isomorphisme au voisinage de $E$).
Par récurrence, on voit que $Y_{m - 1} \to Y$ contracte cette courbe,
ce qui achève la démonstration.
\end{proof}

\begin{lemma}
\label{models-lemma-minimal-model-mapping-property}
Soit $C$ une courbe projective lisse sur $K$ telle que $H^0(C, \mathcal{O}_C) = K$
et de genre $> 0$. Soit $X$ le modèle minimal de $C$
(lemme \ref{models-lemma-minimal-model-unique}).
Soit $Y$ un modèle propre régulier de $C$. Il existe alors un unique
morphisme de modèles $Y \to X$ qui soit une suite de contractions de
courbes exceptionnelles de première espèce.
\end{lemma}

\begin{proof}
L'existence et les propriétés du morphisme $X \to Y$
résultent immédiatement du lemme \ref{models-lemma-pre-exists-minimal-model}
et de l'unicité du modèle minimal.
Le morphisme $Y \to X$ est unique parce que
$C \subset Y$ est schématiquement dense et que $X$ est séparé
(voir Morphismes, lemme \ref{morphisms-lemma-equality-of-morphisms}).
\end{proof}

\begin{example}
\label{models-example-nonunique-in-genus-zero}
Si le genre de $C$ est $0$, les modèles minimaux ne sont effectivement pas uniques.
Considérons en effet le sous-schéma fermé
$$
X \subset \mathbf{P}^2_R
$$
défini par $T_1T_2 - \pi T_0^2 = 0$. Plus précisément, $X$ est défini
comme $\text{Proj}(R[T_0, T_1, T_2]/(T_1T_2 - \pi T_0^2))$. La
fibre spéciale $X_k$ est alors la réunion de deux courbes exceptionnelles $C_1$, $C_2$, toutes deux
isomorphes à $\mathbf{P}^1_k$
(exactement comme dans le lemme \ref{models-lemma-exceptional-curves-dont-meet}).
La projection depuis $(0 : 1 : 0)$ définit un morphisme $X \to \mathbf{P}^1_R$
qui contracte $C_2$ et induit un isomorphisme de $C_1$ avec la fibre spéciale
de $\mathbf{P}^1_R$. La projection depuis $(0 : 0 : 1)$ définit un
morphisme $X \to \mathbf{P}^1_R$ qui contracte $C_1$ et induit un
isomorphisme de $C_2$ avec la fibre spéciale de $\mathbf{P}^1_R$.
Plus précisément, ces morphismes correspondent aux homomorphismes de $R$-algèbres
graduées
$$
R[T_0, T_1] \longrightarrow
R[T_0, T_1, T_2]/(T_1T_2 - \pi T_0^2) \longleftarrow
R[T_0, T_2]
$$
Nous étudierons ce phénomène dans le lemme \ref{models-lemma-nonuniqueness}.
\end{example}









\section{Une formule pour le genre}
\label{models-section-genus-formula}

\noindent
La structure combinatoire issue d'un modèle propre régulier
est soumise à une restriction supplémentaire.

\begin{lemma}
\label{models-lemma-add-component}
Dans la situation \ref{models-situation-regular-model}, supposons donnés des
diviseurs de Cartier effectifs $D, D' \subset X$ tels que
$D' = D + C_i$ pour un certain $i \in \{1, \ldots, n\}$ et $D' \subset X_k$.
Alors
$$
\chi(X_k, \mathcal{O}_{D'}) - \chi(X_k, \mathcal{O}_D) =
\chi(X_k, \mathcal{O}_X(-D)|_{C_i}) =
-(D \cdot C_i) + \chi(C_i, \mathcal{O}_{C_i})
$$
\end{lemma}

\begin{proof}
La seconde égalité résulte de la définition de la forme bilinéaire
$(\ \cdot\ )$ dans (\ref{models-equation-form}) et du
lemme \ref{models-lemma-intersection-pairing}. Pour voir la première
égalité, distinguons deux cas.
En effet, si $C_i \not \subset D$, alors $D'$ est la réunion
schématique de $D$ et $C_i$ (d'après
Diviseurs, lemme \ref{divisors-lemma-sum-effective-Cartier-divisors-union}),
et l'on obtient une suite exacte courte
$$
0 \to \mathcal{O}_{D'} \to
\mathcal{O}_D \times \mathcal{O}_{C_i} \to
\mathcal{O}_{D \cap C_i} \to 0
$$
d'après Morphismes, lemme \ref{morphisms-lemma-scheme-theoretic-union}.
Comme on a aussi une suite exacte
$$
0 \to \mathcal{O}_X(-D)|_{C_i} \to
\mathcal{O}_{C_i} \to \mathcal{O}_{D \cap C_i} \to 0
$$
(Diviseurs, remarque \ref{divisors-remark-ses-regular-section}),
on conclut que l'affirmation est vraie
par additivité des caractéristiques d'Euler
(Variétés, lemme \ref{varieties-lemma-euler-characteristic-additive}).
D'autre part, si $C_i \subset D$, alors on obtient une
suite exacte
$$
0 \to \mathcal{O}_X(-D)|_{C_i} \to \mathcal{O}_{D'} \to \mathcal{O}_D \to 0
$$
d'après Diviseurs, lemme \ref{divisors-lemma-ses-add-divisor},
et l'on voit immédiatement que le lemme est vrai.
\end{proof}

\begin{lemma}
\label{models-lemma-genus-formula}
Dans la situation \ref{models-situation-regular-model}, on a
$$
g_C = 1 + \sum\nolimits_{i = 1, \ldots, n}
m_i\left([\kappa_i : k] (g_i - 1) - \frac{1}{2}(C_i \cdot C_i)\right)
$$
où $\kappa_i = H^0(C_i, \mathcal{O}_{C_i})$,
$g_i$ est le genre de $C_i$, et $g_C$ le genre de $C$.
\end{lemma}

\begin{proof}
Notre outil fondamental sera Catégories dérivées de schémas, lemme
\ref{perfect-lemma-chi-locally-constant-geometric},
qui montre que
$$
1 - g_C = \chi(C, \mathcal{O}_C) =
\chi(X_k, \mathcal{O}_{X_k})
$$
Choisissons une suite de diviseurs de Cartier effectifs
$$
X_k = D_m \supset D_{m - 1} \supset \ldots \supset D_1 \supset D_0 = \emptyset
$$
telle que $D_{j + 1} = D_j + C_{i_j}$ pour tout $j$. (Il est clair que
l'on peut choisir une telle suite en diminuant d'une unité, à chaque étape, une multiplicité
non nulle de $D_{j + 1}$.) En appliquant le lemme \ref{models-lemma-add-component},
en partant de $\chi(\mathcal{O}_{D_0}) = 0$, on obtient
\begin{align*}
1 - g_C
& =
\chi(X_k, \mathcal{O}_{X_k}) \\
& =
\sum\nolimits_j
\left(-(D_j \cdot C_{i_j}) +  \chi(C_{i_j}, \mathcal{O}_{C_{i_j}})\right) \\
& =
- \sum\nolimits_j
(C_{i_1} + C_{i_2} + \ldots + C_{i_{j - 1}} \cdot C_{i_j}) +
\sum\nolimits_j \chi(C_{i_j}, \mathcal{O}_{C_{i_j}}) \\
& =
-\frac{1}{2}\sum\nolimits_{j \not = j'} (C_{i_{j'}} \cdot C_{i_j}) +
\sum m_i \chi(C_i, \mathcal{O}_{C_i}) \\
& =
\frac{1}{2} \sum m_i(C_i \cdot C_i) + \sum m_i \chi(C_i, \mathcal{O}_{C_i})
\end{align*}
La dernière égalité mérite peut-être une explication. En effet, puisque
$\sum_j C_{i_j} = \sum m_i C_i$, on a
$(\sum_j C_{i_j} \cdot \sum_j C_{i_j}) = 0$ d'après le
lemme \ref{models-lemma-properties-form}. On voit donc que
$$
0 = \sum\nolimits_{j \not = j'} (C_{i_{j'}} \cdot C_{i_j}) +
\sum m_i(C_i \cdot C_i)
$$
en décomposant ce produit en termes « non diagonaux » et « diagonaux ».
Notons que $\kappa_i$ est un corps fini sur $k$ d'après
Variétés, lemme \ref{varieties-lemma-regular-functions-proper-variety}.
Le genre de $C_i$ est donc défini et l'on a
$\chi(C_i, \mathcal{O}_{C_i}) = [\kappa_i : k](1 - g_i)$.
En réunissant le tout et en réarrangeant les termes, on obtient
$$
g_C = - \frac{1}{2}\sum m_i(C_i \cdot C_i) +
\sum m_i[\kappa_i : k](g_i - 1) + 1
$$
ce qui est aussi l'assertion du lemme.
\end{proof}

\begin{lemma}
\label{models-lemma-numerical-type-of-model}
Dans la situation \ref{models-situation-regular-model}, avec
$\kappa_i = H^0(C_i, \mathcal{O}_{C_i})$ et $g_i$ le genre de $C_i$,
les données
$$
n, m_i, (C_i \cdot C_j), [\kappa_i : k], g_i
$$
forment un type numérique de genre égal au genre de $C$.
\end{lemma}

\begin{proof}
(Dans la démonstration du lemme \ref{models-lemma-genus-formula},
nous avons vu que les quantités
utilisées dans l'énoncé du lemme sont bien définies.)
Nous devons vérifier les conditions (1) -- (5) de la
définition \ref{models-definition-type}.

\medskip\noindent
La condition (1) est immédiate.

\medskip\noindent
Condition (2). La symétrie de la matrice $(C_i \cdot C_j)$ résulte de
l'équation (\ref{models-equation-form}) et du
lemme \ref{models-lemma-intersection-pairing}.
La positivité de $(C_i \cdot C_j)$ pour $i \not = j$
est la partie (1) du lemme \ref{models-lemma-properties-form}.

\medskip\noindent
La condition (3) est la partie (3) du lemme \ref{models-lemma-properties-form}.

\medskip\noindent
La condition (4) est la partie (2) du lemme \ref{models-lemma-properties-form}.

\medskip\noindent
La condition (5) résulte du fait que $(C_i \cdot C_j)$ est
le degré d'un module inversible sur $C_i$, lequel est divisible
par $[\kappa_i : k]$ ; voir Variétés, lemme \ref{varieties-lemma-divisible}.

\medskip\noindent
La formule du genre démontrée dans le lemme \ref{models-lemma-genus-formula}
nous dit que le type numérique possède le genre annoncé ; voir la
définition \ref{models-definition-genus}.
\end{proof}

\begin{definition}
\label{models-definition-numerical-type-model}
Dans la situation \ref{models-situation-regular-model}, le
{\it type numérique associé à $X$} est le type numérique
décrit dans le lemme \ref{models-lemma-numerical-type-of-model}.
\end{definition}

\noindent
Nous allons maintenant faire correspondre la minimalité du modèle à celle du type.

\begin{lemma}
\label{models-lemma-numerical-type-minimal-model}
Dans la situation \ref{models-situation-regular-model}, les assertions suivantes
sont équivalentes :
\begin{enumerate}
\item $X$ est un modèle minimal, et
\item le type numérique associé à $X$ est minimal.
\end{enumerate}
\end{lemma}

\begin{proof}
Si le type numérique est minimal, il n'existe aucun $i$ tel que
$g_i = 0$ et $(C_i \cdot C_i) = -[\kappa_i: k]$ ; voir la
définition \ref{models-definition-type-minimal}.
Ceci implique assurément qu'aucune des courbes $C_i$
n'est une courbe exceptionnelle de première espèce.

\medskip\noindent
Réciproquement, supposons que le type numérique ne soit pas minimal.
Il existe alors un $i$ tel que $g_i = 0$ et
$(C_i \cdot C_i) = -[\kappa_i: k]$.
Nous affirmons que ceci implique que $C_i$ est une courbe exceptionnelle
de première espèce. En effet, le faisceau inversible
$\mathcal{O}_X(-C_i)|_{C_i}$ est de degré $-(C_i \cdot C_i) = [\kappa_i : k]$
lorsque $C_i$ est vue comme une courbe propre sur $k$, et est donc
de degré $1$ lorsque $C_i$ est vue comme une courbe propre sur $\kappa_i$.
En appliquant
Courbes algébriques, proposition \ref{curves-proposition-projective-line},
on conclut que $C_i \cong \mathbf{P}^1_{\kappa_i}$ comme schémas
sur $\kappa_i$. Comme le groupe de Picard de $\mathbf{P}^1$
sur un corps est $\mathbf{Z}$, on voit que le faisceau normal
de $C_i$ dans $X$ est isomorphe à $\mathcal{O}_{\mathbf{P}_{\kappa_i}}(-1)$,
ce qui achève la démonstration.
\end{proof}

\begin{remark}
\label{models-remark-numerical-type-not-from-model}
Tout type numérique ne provient pas d'un modèle, pour la raison triviale
qu'il existe des types numériques dont le genre est négatif.
Il existe des types numériques minimaux de genre positif qui
ne sont pas le type numérique associé à un modèle (sur un anneau de valuation discrète)
d'une courbe projective lisse géométriquement irréductible (sur le
corps des fractions de cet anneau). Un exemple simple est
$n = 1$, $m_1 = 1$, $a_{11} = 0$, $w_1 = 6$, $g_1 = 1$.
En effet, dans ce cas, la fibre spéciale $X_k$ ne serait pas
géométriquement connexe, car elle serait définie sur une extension
$\kappa$ de $k$ de degré $6$. Cela contredit le
fait que la fibre générique est géométriquement connexe (voir
Compléments sur les morphismes, lemme
\ref{more-morphisms-lemma-geometrically-connected-fibres-towards-normal}).
De même, $n = 2$, $m_1 = m_2 = 1$, $-a_{11} = -a_{22} = a_{12} = a_{21} = 6$,
$w_1 = w_2 = 6$, $g_1 = g_2 = 1$ fournirait un exemple
pour la même raison (détails omis). Mais si le pgcd des
$w_i$ vaut $1$, nous ne disposons d'aucun exemple.
\end{remark}

\begin{lemma}
\label{models-lemma-numerical-type-rational-point}
Dans la situation \ref{models-situation-regular-model}, supposons que $C$ ait un point $K$-rationnel.
Alors
\begin{enumerate}
\item $X_k$ possède un point $k$-rationnel $x$ qui est un point lisse de $X_k$
sur $k$,
\item si $x \in C_i$, alors $H^0(C_i, \mathcal{O}_{C_i}) = k$ et
$m_i = 1$, et
\item $H^0(X_k, \mathcal{O}_{X_k}) = k$, et $X_k$ est de genre égal au
genre de $C$.
\end{enumerate}
\end{lemma}

\begin{proof}
Puisque $X \to \Spec(R)$ est propre, le point $K$-rationnel se prolonge en
un morphisme $a : \Spec(R) \to X$ par le critère valuatif de propreté
(Morphismes, lemme \ref{morphisms-lemma-characterize-proper}).
Soit $x \in X$ l'image par $a$ du point fermé de $\Spec(R)$.
Alors $a$ correspond à un homomorphisme de $R$-algèbres
$\psi : \mathcal{O}_{X, x} \to R$
(voir Schémas, section \ref{schemes-section-points}).
Il s'ensuit que $\pi \not \in \mathfrak m_x^2$ (puisque l'image
de $\pi$ dans $R$ n'appartient pas à $\mathfrak m_R^2$).
Par conséquent, $\mathcal{O}_{X_k, x} = \mathcal{O}_{X, x}/\pi \mathcal{O}_{X, x}$
est régulier (Algèbre, lemme \ref{algebra-lemma-regular-ring-CM}).
Alors $X_k \to \Spec(k)$ est lisse en $x$ d'après
Algèbre, lemme \ref{algebra-lemma-separable-smooth}.
Il s'ensuit que $x$ appartient à une unique composante irréductible
$C_i$ de $X_k$, que $\mathcal{O}_{C_i, x} = \mathcal{O}_{X_k, x}$
et que $m_i = 1$. Le fait que $C_i$ possède un
point $k$-rationnel implique que le corps
$\kappa_i = H^0(C_i, \mathcal{O}_{C_i})$
(Variétés, lemme \ref{varieties-lemma-regular-functions-proper-variety})
est égal à $k$. Ceci démontre (1). On a
$H^0(X_k, \mathcal{O}_{X_k}) = k$
parce que $H^0(X_k, \mathcal{O}_{X_k})$ est une extension de corps
de $k$ (lemme \ref{models-lemma-regular-model-field})
qui s'envoie dans $H^0(C_i, \mathcal{O}_{C_i}) = k$.
L'égalité des genres résulte du lemme \ref{models-lemma-regular-model-genus}.
\end{proof}

\begin{lemma}
\label{models-lemma-genus-reduction-smaller}
Dans la situation \ref{models-situation-regular-model}, supposons que $X$ soit un modèle minimal,
que $\gcd(m_1, \ldots, m_n) = 1$ et que $H^0((X_k)_{red}, \mathcal{O}) = k$. Alors
l'application
$$
H^1(X_k, \mathcal{O}_{X_k}) \to H^1((X_k)_{red}, \mathcal{O}_{(X_k)_{red}})
$$
est surjective et possède un noyau non trivial dès que $(X_k)_{red} \not = X_k$.
\end{lemma}

\begin{proof}
Par annulation de la cohomologie en degrés $\geq 2$ sur $X_k$
(Cohomologie, proposition \ref{cohomology-proposition-vanishing-Noetherian}),
toute surjection de faisceaux abéliens sur $X_k$ induit une surjection sur $H^1$.
Considérons la suite
$$
(X_k)_{red} = Z_0 \subset Z_1 \subset \ldots \subset Z_m = X_k
$$
du lemme \ref{models-lemma-regular-model-field}. Puisque les homomorphismes de corps
$H^0(Z_j, \mathcal{O}_{Z_j}) \to
H^0((X_k)_{red}, \mathcal{O}_{(X_k)_{red}}) = k$
sont injectifs, on conclut que $H^0(Z_j, \mathcal{O}_{Z_j}) = k$ pour
$j = 0, \ldots, m$. Il s'ensuit que
$H^0(X_k, \mathcal{O}_{X_k}) \to H^0(Z_{m - 1}, \mathcal{O}_{Z_{m - 1}})$
est surjectif. Soit $C = C_{i_m}$. Alors $X_k = Z_{m - 1} + C$.
Soit $\mathcal{L} = \mathcal{O}_X(-Z_{m - 1})|_C$.
Alors $\mathcal{L}$ est un $\mathcal{O}_C$-module inversible.
Comme dans la démonstration du lemme \ref{models-lemma-regular-model-field},
il existe une suite exacte
$$
0 \to \mathcal{L} \to \mathcal{O}_{X_k} \to \mathcal{O}_{Z_{m - 1}} \to 0
$$
de faisceaux cohérents sur $X_k$. On en déduit
une suite exacte courte
$$
0 \to
H^1(C, \mathcal{L}) \to H^1(X_k, \mathcal{O}_{X_k}) \to
H^1(Z_{m - 1}, \mathcal{O}_{Z_{m - 1}}) \to 0
$$
Le degré de $\mathcal{L}$ sur $C$ relativement à $k$ est
$$
(C \cdot -Z_{m - 1}) = (C \cdot C - X_k) = (C \cdot C)
$$
Posons $\kappa = H^0(C, \mathcal{O}_C)$ et $w = [\kappa : k]$.
Par définition du degré d'un faisceau inversible, on voit que
$$
\chi(C, \mathcal{L}) =
\chi(C, \mathcal{O}_C) + (C \cdot C) =
w(1 - g_C) + (C \cdot C)
$$
où $g_C$ est le genre de $C$. Cette expression est $< 0$ puisque $X$ est
minimal et que, par conséquent, $C$ n'est pas une courbe exceptionnelle de
première espèce (voir la preuve du lemme
\ref{models-lemma-numerical-type-minimal-model}). Ainsi
$\dim_k H^1(C, \mathcal{L}) > 0$, ce qui achève la preuve.
\end{proof}

\begin{lemma}
\label{models-lemma-genus-reduction-bigger-than}
Dans la situation \ref{models-situation-regular-model}, supposons que $X_k$ possède
un point $k$-rationnel $x$ qui est un point lisse de
$X_k \to \Spec(k)$. Alors
$$
\dim_k H^1((X_k)_{red}, \mathcal{O}_{(X_k)_{red}}) \geq
g_{top} + g_{geom}(X_k/k)
$$
où $g_{geom}$ est défini dans Courbes algébriques, section
\ref{curves-section-genus-geometric-genus}, et où $g_{top}$ est le genre
topologique (définition \ref{models-definition-top-genus}) du type numérique
associé à $X_k$ (définition \ref{models-definition-numerical-type-model}).
\end{lemma}

\begin{proof}
Nous allons démontrer l'inégalité
$$
\dim_k H^1(D, \mathcal{O}_D) \geq g_{top}(D) + g_{geom}(D/k)
$$
pour tout diviseur de Cartier effectif réduit connexe
$D \subset (X_k)_{red}$ contenant $x$, par récurrence sur le nombre de
composantes irréductibles de $D$.
Ici $g_{top}(D) = 1 - m + e$, où $m$ est le nombre de composantes
irréductibles de $D$ et $e$ le nombre de paires non ordonnées de composantes
de $D$ qui se rencontrent.

\medskip\noindent
Cas initial : $D$ possède une seule composante irréductible. Alors
$D = C_i$ est l'unique composante irréductible contenant $x$.
Dans ce cas, $\dim_k H^1(D, \mathcal{O}_D) = g_i$ et
$g_{top}(D) = 0$. Puisque $C_i$ possède un point lisse $k$-rationnel,
elle est géométriquement intègre (Variétés, lemme
\ref{varieties-lemma-variety-with-smooth-rational-point}). Il s'ensuit que
$g_i$ est le genre de $C_{i, \overline{k}}$ (Courbes algébriques, lemme
\ref{curves-lemma-genus-base-change}). Il s'ensuit aussi que
$g_{geom}(D/k)$ est le genre de la normalisation
$C_{i, \overline{k}}^\nu$ de $C_{i, \overline{k}}$. En appliquant Courbes
algébriques, lemme \ref{curves-lemma-genus-normalization}, au morphisme de
normalisation $C_{i, \overline{k}}^\nu \to C_{i, \overline{k}}$, on obtient
\begin{equation}
\label{models-equation-genus-change-special-component}
\text{genre de }C_{i, \overline{k}} \geq
\text{genre de }C_{i, \overline{k}}^\nu
\end{equation}
En combinant ce qui précède, on conclut que
$\dim_k H^1(D, \mathcal{O}_D) \geq g_{top}(D) + g_{geom}(D/k)$
dans ce cas.

\medskip\noindent
Étape de récurrence. Supposons que $D$ possède plus de $1$ composante
irréductible. On peut alors écrire $D = C_i + D'$, avec $x \in D'$ et
$D'$ encore connexe. C'est un exercice de théorie des graphes que nous
laissons au lecteur (indication : prendre pour $C_i$ la composante de $D$
la plus éloignée de $x$). Calculons la variation des invariants.
Comme $x \in D'$, on a
$H^0(D, \mathcal{O}_D) = H^0(D', \mathcal{O}_{D'}) = k$.
En considérant la suite exacte courte de faisceaux
$$
0 \to \mathcal{O}_D \to \mathcal{O}_{C_i} \oplus \mathcal{O}_{D'}
\to \mathcal{O}_{C_i \cap D'} \to 0
$$
(Morphismes, lemme \ref{morphisms-lemma-scheme-theoretic-union}) et en
utilisant l'additivité des caractéristiques d'Euler, on trouve
\begin{align*}
\dim_k H^1(D, \mathcal{O}_D) - \dim_k H^1(D', \mathcal{O}_{D'})
& =
-\chi(\mathcal{O}_{C_i}) + \chi(\mathcal{O}_{C_i \cap D'}) \\
& =
w_i(g_i - 1) + \sum\nolimits_{C_j \subset D'} a_{ij}
\end{align*}
Comme dans le lemme \ref{models-lemma-numerical-type-of-model}, on pose ici
$w_i = [\kappa_i : k]$, $\kappa_i = H^0(C_i, \mathcal{O}_{C_i})$ ;
$g_i$ est le genre de $C_i$, et $a_{ij} = (C_i \cdot C_j)$.
On a
$$
g_{top}(D) - g_{top}(D') = -1 +
\sum\nolimits_{C_j \subset D'\text{ rencontrant }C_i} 1
$$
On a
$$
g_{geom}(D/k) - g_{geom}(D'/k) = g_{geom}(C_i/k)
$$
d'après Courbes algébriques, lemme
\ref{curves-lemma-bound-geometric-genus}. En combinant ces égalités avec
l'hypothèse de récurrence, on conclut qu'il suffit de montrer que
$$
w_i g_i - g_{geom}(C_i/k) +
\sum\nolimits_{C_j \subset D'\text{ rencontrant } C_i} (a_{ij} - 1) - (w_i - 1)
$$
est positif ou nul. En effet, on a
\begin{equation}
\label{models-equation-genus-change}
w_i g_i \geq [\kappa_i : k]_s g_i \geq g_{geom}(C_i/k)
\end{equation}
La seconde inégalité résulte de Courbes algébriques, lemme
\ref{curves-lemma-bound-geometric-genus-curve}. D'autre part, puisque
$w_i$ divise $a_{ij}$ (Variétés, lemme
\ref{varieties-lemma-divisible}), il est clair que
\begin{equation}
\label{models-equation-change-intersections}
\sum\nolimits_{C_j \subset D'\text{ rencontrant } C_i} (a_{ij} - 1) - (w_i - 1)
\geq 0
\end{equation}
car il existe au moins une composante $C_j \subset D'$ rencontrant $C_i$.
\end{proof}

\begin{lemma}
\label{models-lemma-equality-genus-reduction-bigger-than}
Si l'égalité a lieu dans le lemme
\ref{models-lemma-genus-reduction-bigger-than}, alors
\begin{enumerate}
\item l'unique composante irréductible de $X_k$ contenant $x$ est une
courbe projective lisse géométriquement irréductible sur $k$ ;
\item si $C \subset X_k$ est une autre composante irréductible, alors
$\kappa = H^0(C, \mathcal{O}_C)$ est une extension finie séparable de $k$,
$C$ possède un point $\kappa$-rationnel et $C$ est lisse sur $\kappa$.
\end{enumerate}
\end{lemma}

\begin{proof}
En reprenant la preuve du lemme
\ref{models-lemma-genus-reduction-bigger-than}, on voit que, pour obtenir
l'égalité, les inégalités
(\ref{models-equation-genus-change-special-component}),
(\ref{models-equation-genus-change}) et
(\ref{models-equation-change-intersections})
doivent être des égalités.

\medskip\noindent
Soit $C_i$ la composante irréductible contenant $x$.
L'égalité dans (\ref{models-equation-genus-change-special-component}) montre, par
Courbes algébriques, lemme \ref{curves-lemma-genus-normalization}, que
$C_{i, \overline{k}}^\nu \to C_{i, \overline{k}}$ est un isomorphisme.
Par conséquent, $C_{i, \overline{k}}$ est lisse et l'assertion (1) est vraie.

\medskip\noindent
Soit ensuite $C_i \subset X_k$ une autre composante irréductible.
On peut alors supposer que $D = D' + C_i$ comme dans l'étape de récurrence
de la preuve du lemme \ref{models-lemma-genus-reduction-bigger-than}.
L'égalité dans (\ref{models-equation-genus-change}) implique immédiatement que
$\kappa_i/k$ est finie séparable. L'égalité dans
(\ref{models-equation-change-intersections}) implique soit que $a_{ij} = 1$ pour
un certain $j$, soit qu'il existe une unique composante
$C_j \subset D'$ rencontrant $C_i$ et que $a_{ij} = w_i$.
Dans les deux cas, on trouve que $C_i$ possède un point
$\kappa_i$-rationnel $c$ et que $c = C_i \cap C_j$ au sens schématique.
Puisque $\mathcal{O}_{X,c}$ est un anneau local régulier, cela implique que
les équations locales de $C_i$ et de $C_j$ forment un système régulier de
paramètres de l'anneau local $\mathcal{O}_{X,c}$. Alors
$\mathcal{O}_{C_i,c}$ est régulier d'après Algèbre, lemme
\ref{algebra-lemma-regular-ring-CM}. On conclut que
$C_i \to \Spec(\kappa_i)$ est lisse en $c$ (Algèbre, lemme
\ref{algebra-lemma-separable-smooth}). Il s'ensuit que $C_i$ est
géométriquement intègre sur $\kappa_i$ (Variétés, lemme
\ref{varieties-lemma-variety-with-smooth-rational-point}). Pour conclure,
il reste à montrer que $C_i$ est lisse sur $\kappa_i$. Observons que
$$
C_{i, \overline{k}} = C_i \times_{\Spec(k)} \Spec(\overline{k})
= \coprod\nolimits_{\kappa_i \to \overline{k}}
C_i \times_{\Spec(\kappa_i)} \Spec(\overline{k})
$$
où il y a $[\kappa_i : k]$ composantes. Ainsi, si $C_i$ n'est pas lisse
sur $\kappa_i$, chacune de ces courbes n'est pas lisse ; elles ne sont donc
pas normales, et le morphisme de normalisation fait baisser le genre
(Courbes algébriques, lemme \ref{curves-lemma-genus-normalization}). Cela
est impossible, car il ferait baisser le genre géométrique de $C_i/k$, en
contradiction avec $[\kappa_i : k]g_i = g_{geom}(C_i/k)$.
\end{proof}








\section{Contraction des courbes exceptionnelles}
\label{models-section-contract}

\noindent
Le lemme suivant décrit la variation des nombres d'intersection lorsque
l'on contracte une courbe exceptionnelle de première espèce dans un modèle
propre régulier. Nous le plaçons ici surtout pour le comparer aux
contractions numériques introduites dans le lemme \ref{models-lemma-contract}.
Nous comparerons les contractions géométriques et numériques dans la
remarque \ref{models-remark-compare-contractions}.

\begin{lemma}
\label{models-lemma-blowdown-regular-model}
Dans la situation \ref{models-situation-regular-model}, supposons que $C_n$ soit
une courbe exceptionnelle de première espèce. Soit $f : X \to X'$ la
contraction de $C_n$. Posons $C'_i = f(C_i)$ et écrivons
$X'_k = \sum m'_iC'_i$. Alors $X'$, $C'_i$, $i = 1, \ldots, n' = n - 1$,
et $m'_i = m_i$ sont comme dans la situation
\ref{models-situation-regular-model}, et l'on a
\begin{enumerate}
\item pour $i,j<n$,
$(C'_i \cdot C'_j) =
(C_i \cdot C_j) - (C_i \cdot C_n)(C_j \cdot C_n)/(C_n \cdot C_n)$ ;
\item pour $i<n$, si $C_i \cap C_n \not = \emptyset$, il existe des
applications $\kappa_i \leftarrow \kappa'_i \rightarrow \kappa_n$.
\end{enumerate}
Ici $\kappa_i = H^0(C_i, \mathcal{O}_{C_i})$ et
$\kappa'_i = H^0(C'_i, \mathcal{O}_{C'_i})$.
\end{lemma}

\begin{proof}
D'après Résolution des surfaces, lemme
\ref{resolve-lemma-contract-ample}, on peut contracter $C_n$ par un
morphisme $f : X \to X'$ tel que $X'$ soit régulier et projectif sur $R$.
On voit donc que $X'$ est comme dans la situation
\ref{models-situation-regular-model}. Soit $x \in X'$ l'image de $C_n$.
Puisque $f$ induit un isomorphisme
$X \setminus C_n \to X' \setminus \{x\}$, il est clair que
$m'_i = m_i$ pour $i<n$.

\medskip\noindent
L'assertion (2) du lemme résulte immédiatement de l'existence des
morphismes $C_i \to C'_i$ et $C_n \to x \to C'_i$.

\medskip\noindent
D'après Diviseurs, lemme
\ref{divisors-lemma-blow-up-pullback-effective-Cartier}, l'image réciproque
$f^{-1}C'_i$ est définie. D'après Diviseurs, lemme
\ref{divisors-lemma-effective-Cartier-divisor-is-a-sum}, on voit que
$f^{-1}C'_i = C_i + e_iC_n$ pour un certain $e_i \geq 0$. Puisque
$\mathcal{O}_X(C_i + e_iC_n) = \mathcal{O}_X(f^{-1}C'_i) =
f^*\mathcal{O}_{X'}(C'_i)$
(Diviseurs, lemme \ref{divisors-lemma-pullback-effective-Cartier-divisors})
et que l'image réciproque d'un faisceau inversible se restreint au faisceau
inversible trivial sur $C_n$, on voit que
$$
0 = \deg_{C_n}(\mathcal{O}_X(C_i + e_iC_n)) =
(C_i + e_iC_n \cdot C_n) = (C_i \cdot C_n) + e_i(C_n \cdot C_n)
$$
Comme $f_j = f|_{C_j} : C_j \to C_j$ est un morphisme propre birationnel
de courbes propres sur $k$, on voit que
$\deg_{C'_j}(\mathcal{O}_{X'}(C'_i)|_{C'_j})$ est égal à
$\deg_{C_j}(f_j^*\mathcal{O}_{X'}(C'_i)|_{C'_j})$
(Variétés, lemme \ref{varieties-lemma-degree-birational-pullback}).
En considérant le diagramme commutatif
$$
\xymatrix{
C_j \ar[r] \ar[d]_{f_j} & X \ar[d]^f \\
C'_j \ar[r] & X'
}
$$
et en utilisant Diviseurs, lemme
\ref{divisors-lemma-pullback-effective-Cartier-divisors}, on voit que
$$
(C'_i \cdot C'_j) = \deg_{C'_j}(\mathcal{O}_{X'}(C'_i)|_{C'_j}) =
\deg_{C_j}(\mathcal{O}_X(C_i + e_iC_n)) = (C_i + e_iC_n \cdot C_j)
$$
En y substituant la formule de $e_i$ obtenue ci-dessus, on voit que (1)
est vraie.
\end{proof}

\begin{remark}
\label{models-remark-genus-change}
Dans la situation du lemme \ref{models-lemma-blowdown-regular-model}, on peut aussi
décrire exactement la relation entre le genre $g_i$ de $C_i$ et le genre
$g'_i$ de $C'_i$. La formule est
$$
g'_i = \frac{w_i}{w'_i}(g_i - 1) + 1 +
\frac{(C_i \cdot C_n)^2 - w_n(C_i \cdot C_n)}{2w'_iw_n}
$$
où $w_i = [\kappa_i : k]$, $w_n = [\kappa_n : k]$ et
$w'_i = [\kappa'_i : k]$.
Pour le démontrer, considérons la suite exacte courte
$$
0 \to \mathcal{O}_{X'}(-C'_i) \to \mathcal{O}_{X'} \to
\mathcal{O}_{C'_i} \to 0
$$
et son image réciproque sur $X$, qui s'écrit
$$
0 \to \mathcal{O}_X(-C'_i - e_iC_n) \to \mathcal{O}_X \to
\mathcal{O}_{C_i + e_iC_n} \to 0
$$
où $e_i$ est comme dans la preuve du lemme
\ref{models-lemma-blowdown-regular-model}. Comme
$Rf_*f^*\mathcal{L} = \mathcal{L}$ pour tout module inversible
$\mathcal{L}$ sur $X'$ (détails omis), on conclut que
$$
Rf_*\mathcal{O}_{C_i + e_iC_n} = \mathcal{O}_{C'_i}
$$
en tant que complexes de faisceaux cohérents sur $X'_k$.
Les deux membres ont donc la même caractéristique d'Euler, laquelle coïncide
avec la caractéristique d'Euler de $\mathcal{O}_{C_i + e_iC_n}$ sur $X_k$.
En utilisant la suite exacte
$$
0 \to \mathcal{O}_{C_i + e_iC_n} \to
\mathcal{O}_{C_i} \oplus \mathcal{O}_{e_iC_n} \to
\mathcal{O}_{C_i \cap e_iC_n} \to 0
$$
et en filtrant encore $\mathcal{O}_{e_iC_n}$ (détails omis), on trouve
$$
\chi(\mathcal{O}_{C'_i}) =
\chi(\mathcal{O}_{C_i}) - {e_i + 1 \choose 2}(C_n \cdot C_n)
- e_i(C_i \cdot C_n)
$$
Comme $e_i = -(C_i \cdot C_n)/(C_n \cdot C_n)$ et
$(C_n \cdot C_n) = -w_n$, on en déduit la formule énoncée au début de
cette remarque. Si nous en avons un jour besoin, nous la formulerons comme
un lemme et en donnerons une preuve détaillée.
\end{remark}

\begin{remark}
\label{models-remark-compare-contractions}
Soit $f : X \to X'$ comme dans le lemme
\ref{models-lemma-blowdown-regular-model}. Soient $n, m_i, a_{ij}, w_i, g_i$ le type
numérique associé à $X$, et $n', m'_i, a'_{ij}, w'_i, g'_i$ celui associé à
$X'$. Il résulte clairement du lemme \ref{models-lemma-blowdown-regular-model} et
de la remarque \ref{models-remark-genus-change} que cela coïncide avec la
contraction des types numériques du lemme \ref{models-lemma-contract}, à
l'exception de la valeur de $w'_i$. Dans la situation géométrique, $w'_i$
est un entier positif divisant à la fois $w_i$ et $w_n$. Dans le cas
numérique, nous avons choisi pour $w'_i$ le plus grand entier possible
divisant $w_i$ et tel que $g'_i$ (donné par la formule) soit entier.
Ce choix convient bien dans le cadre numérique, en ce qu'il aide à comparer
les groupes de Picard des types numériques ; voir le lemme
\ref{models-lemma-contract-picard-group} (bien que seule l'injectivité soit
utilisée dans la suite, et que celle-ci reste vraie pour des $w'_i$ plus
petits).
\end{remark}

\begin{lemma}
\label{models-lemma-nonuniqueness}
Soit $C$ une courbe projective lisse sur $K$, avec
$H^0(C, \mathcal{O}_C) = K$ et de genre $0$. S'il existe plus d'un modèle
minimal de $C$, alors la fibre spéciale de tout modèle minimal est
isomorphe à $\mathbf{P}^1_k$.
\end{lemma}

\noindent
On peut renforcer ce lemme en affirmant que la transformation birationnelle
entre deux modèles minimaux non isomorphes se factorise en une suite de
transformations élémentaires comme dans l'exemple
\ref{models-example-nonunique-in-genus-zero}. Si nous en avons un jour besoin,
nous en donnerons ici un énoncé précis et une preuve.

\begin{proof}
Soit $X$ un modèle minimal de $C$. Le type numérique associé à $X$ est de
genre $0$ et minimal (définition \ref{models-definition-numerical-type-model} et
lemme \ref{models-lemma-numerical-type-minimal-model}). Le lemme
\ref{models-lemma-genus-zero} montre donc que $X_k$ est réduit, irréductible,
vérifie $H^0(X_k, \mathcal{O}_{X_k}) = k$ et est de genre $0$.
Soit $Y$ un second modèle minimal de $C$, non isomorphe à $X$.
D'après Résolution des surfaces, lemme
\ref{resolve-lemma-birational-regular-surfaces}, il existe un diagramme de
$S$-morphismes
$$
X = X_0 \leftarrow X_1 \leftarrow \ldots \leftarrow X_n = Y_m
\to \ldots \to Y_1 \to Y_0 = Y
$$
où chaque morphisme est un éclatement en un point fermé. Nous démontrerons
le lemme par récurrence sur $m$. Le cas initial est $m=0$ ; l'assertion y
est vraie puisque nous avons supposé $Y$ minimal : cela impliquerait
$n=0$, mais $X$ n'est pas isomorphe à $Y$, de sorte que ce cas ne se
présente pas et qu'il n'y a rien à vérifier.

\medskip\noindent
Avant de poursuivre, remarquons que $n+1=m+1$ est égal au nombre de
composantes irréductibles de la fibre spéciale de $X_n=Y_m$, puisque
$X_k$ et $Y_k$ sont toutes deux irréductibles. Nous utiliserons également
l'observation suivante : si $X' \to X''$ est un morphisme de modèles
propres réguliers de $C$, alors $X' \to X''$ est un isomorphisme au-dessus
d'un ouvert de $X''$ dont le complément est un ensemble fini de points
fermés de la fibre spéciale de $X''$ ; voir Variétés, lemme
\ref{varieties-lemma-modification-normal-iso-over-codimension-1}.
En fait, tout morphisme $X' \to X''$ de ce type est une suite d'éclatements
en des points fermés (Résolution des surfaces, lemme
\ref{resolve-lemma-proper-birational-regular-surfaces}), et le nombre
d'éclatements est la différence entre les nombres de composantes
irréductibles des fibres spéciales de $X'$ et de $X''$.

\medskip\noindent
Soit $E_i \subset Y_i$, $m \geq i \geq 1$, la courbe contractée par le
morphisme $Y_i \to Y_{i-1}$. Soit $i$ le plus grand indice tel que $E_i$
soit de multiplicité $>1$ dans la fibre spéciale de $Y_i$. Alors les
éclatements suivants $Y_m \to \ldots \to Y_{i+1} \to Y_i$ sont des
isomorphismes au-dessus de $E_i$ ; sinon, un $E_j$ avec $j>i$ aurait une
multiplicité $>1$. Soit $E \subset Y_m$ l'image réciproque de $E_i$.
D'après ce qui précède, $E \subset Y_m$ est une courbe exceptionnelle de
première espèce. Soit $Y_m \to Y'$ la contraction de $E$ (elle existe par
Résolution des surfaces, lemme
\ref{resolve-lemma-contract-when-quasi-projective}). Le morphisme
$Y_m \to X$ doit contracter $E$, puisque $X_k$ est réduit. Il existe donc
des morphismes $Y' \to Y$ et $Y' \to X$ (Résolution des surfaces, lemme
\ref{resolve-lemma-factor-through-contraction}) qui sont des composés d'au
plus $n-1=m-1$ contractions de courbes exceptionnelles (voir la discussion
ci-dessus). La récurrence sur $m$ permet de conclure. En définitive, on peut
supposer que les fibres spéciales de tous les $X_i$ et $Y_i$ sont réduites.

\medskip\noindent
Puisque les fibres de $X_i$ et $Y_i$ sont réduites, les éclatements
$X_i \to X_{i-1}$ et $Y_i \to Y_{i-1}$ ont nécessairement lieu en des
points fermés qui sont des points réguliers des fibres spéciales.
En effet, si $X''$ est un modèle régulier de $C$, si $x \in X''$ est un
point fermé de la fibre spéciale et si $\pi \in \mathfrak m_x^2$, alors la
fibre exceptionnelle $E$ de l'éclatement $X' \to X''$ en $x$ a une
multiplicité au moins $2$ dans la fibre spéciale de $X'$ (calcul local
omis). Ainsi
$\mathcal{O}_{X''_k,x}=\mathcal{O}_{X'',x}/\pi$ est régulier (Algèbre,
lemme \ref{algebra-lemma-regular-ring-CM}), comme annoncé. En particulier,
$x$ est un diviseur de Cartier sur l'unique composante irréductible
$Z'$ de $X''_k$ qui le contient (Variétés, lemme
\ref{varieties-lemma-regular-point-on-curve}). Il s'ensuit que la
transformée stricte $Z \subset X'$ de $Z'$ s'envoie isomorphiquement sur
$Z'$ (utiliser Diviseurs, lemmes \ref{divisors-lemma-strict-transform} et
\ref{divisors-lemma-blow-up-effective-Cartier-divisor}). Autrement dit, si
une composante irréductible $Z$ de $X_i$ n'est pas contractée par
$X_i \to X_j$ ($i>j$), elle s'envoie isomorphiquement sur son image.

\medskip\noindent
Nous pouvons maintenant démontrer le lemme. Soit $E \subset Y_m$ la courbe
exceptionnelle de première espèce contractée par $Y_m \to Y_{m-1}$.
Si $E$ est contractée par le morphisme $Y_m=X_n \to X$, il existe une
factorisation $Y_{m-1} \to X$ (Résolution des surfaces, lemme
\ref{resolve-lemma-factor-through-contraction}) ; de plus,
$Y_{m-1} \to X$ est une suite d'éclatements en des points fermés
(Résolution des surfaces, lemme
\ref{resolve-lemma-proper-birational-regular-surfaces}). Dans ce cas, on
diminue $m$ et l'on conclut par récurrence. Supposons enfin que $E$ ne soit
pas contractée par $Y_m \to X$. Alors $E \to X_k$ est surjective puisque
$X_k$ est irréductible ; ce qui précède montre qu'il s'agit d'un
isomorphisme. Ainsi $X_k$ est isomorphe à une droite projective, comme
voulu.
\end{proof}





\section{Groupes de Picard des modèles}
\label{models-section-picard-groups-models}

\noindent
Supposons que $R,K,k,\pi,C,X,n,C_1,\ldots,C_n,m_1,\ldots,m_n$ soient comme
dans la situation \ref{models-situation-regular-model}. Dans le lemme
\ref{models-lemma-regular-model-pic}, nous avons obtenu une suite exacte
$$
0 \to \mathbf{Z} \to \mathbf{Z}^{\oplus n} \to
\Pic(X) \to \Pic(C) \to 0
$$
Nous voulons utiliser cette suite pour étudier la $\ell$-torsion des
groupes de Picard pour des nombres premiers $\ell$ convenables.

\begin{lemma}
\label{models-lemma-characterize-trivial}
Dans la situation \ref{models-situation-regular-model}, posons
$d=\gcd(m_1,\ldots,m_n)$. Si $\mathcal{L}$ est un
$\mathcal{O}_X$-module inversible qui
\begin{enumerate}
\item se restreint au module inversible trivial sur $C$, et
\item est de degré $0$ sur chaque $C_i$,
\end{enumerate}
alors $\mathcal{L}^{\otimes d} \cong \mathcal{O}_X$.
\end{lemma}

\begin{proof}
D'après le lemme \ref{models-lemma-regular-model-pic}, on a
$\mathcal{L} \cong \mathcal{O}_X(\sum a_iC_i)$ pour certains
$a_i \in \mathbf{Z}$. Le degré de $\mathcal{L}|_{C_j}$ est
$\sum_j a_j(C_i \cdot C_j)$. En particulier,
$(\sum a_iC_i \cdot \sum a_iC_i)=0$. Le lemme
\ref{models-lemma-properties-form} montre donc que
$(a_1,\ldots,a_n)=q(m_1,\ldots,m_n)$ pour un certain $q\in\mathbf{Q}$.
Ainsi $\mathcal{L}=\mathcal{O}_X(lD)$ pour un certain
$l\in\mathbf{Z}$, où $D=\sum(m_i/d)C_i$ est comme dans le lemme
\ref{models-lemma-multiple-fibre-normal-bundle}, et la conclusion en résulte.
\end{proof}

\begin{lemma}
\label{models-lemma-canonical-map-of-pic}
Dans la situation \ref{models-situation-regular-model}, soit $T$ le type numérique
associé à $X$. Il existe une application canonique
$$
\Pic(C) \to \Pic(T)
$$
dont le noyau est exactement constitué des modules inversibles sur $C$ qui
sont les restrictions de modules inversibles $\mathcal{L}$ sur $X$ tels
que $\deg_{C_i}(\mathcal{L}|_{C_i})=0$ pour $i=1,\ldots,n$.
\end{lemma}

\begin{proof}
Rappelons que $w_i=[\kappa_i:k]$, où
$\kappa_i=H^0(C_i,\mathcal{O}_{C_i)})$, et que le degré de tout module
inversible sur $C_i$ est divisible par $w_i$ (Variétés, lemme
\ref{varieties-lemma-divisible}). On peut donc considérer l'application
$$
\frac{\deg}{w} : \Pic(X) \to \mathbf{Z}^{\oplus n}, \quad
\mathcal{L} \mapsto
(\frac{\deg(\mathcal{L}|_{C_1})}{w_1}, \ldots,
\frac{\deg(\mathcal{L}|_{C_n})}{w_n})
$$
L'image de $\mathcal{O}_X(C_j)$ par cette application est
$$
((C_j \cdot C_1)/w_1, \ldots, (C_j \cdot C_n)/w_n) =
(a_{1j}/w_1, \ldots, a_{nj}/w_n)
$$
qui est exactement l'image du $j$-ième vecteur de la base canonique par
l'application $(a_{ij}/w_i):\mathbf{Z}^{\oplus n}\to\mathbf{Z}^{\oplus n}$
définissant le groupe de Picard de $T$ ; voir la définition
\ref{models-definition-picard-group}. Ainsi l'application canonique du lemme
provient du diagramme commutatif
$$
\xymatrix{
\mathbf{Z}^{\oplus n} \ar[r] \ar[d]_{\text{id}} &
\Pic(X) \ar[r] \ar[d]^{\frac{\deg}{w}} &
\Pic(C) \ar[r] \ar[d] & 0 \\
\mathbf{Z}^{\oplus n} \ar[r]^{(a_{ij}/w_i)} &
\mathbf{Z}^{\oplus n} \ar[r] &
\Pic(T) \ar[r] & 0
}
$$
à lignes exactes (la ligne supérieure l'est par le lemme
\ref{models-lemma-regular-model-pic}). La description du noyau est claire.
\end{proof}

\begin{lemma}
\label{models-lemma-sequence-torsion}
Dans la situation \ref{models-situation-regular-model}, posons
$d=\gcd(m_1,\ldots,m_n)$ et soit $T$ le type numérique associé à $X$.
Soit $h\geq1$ un entier premier à $d$. Il existe une suite exacte
$$
0 \to \Pic(X)[h] \to \Pic(C)[h] \to \Pic(T)[h]
$$
\end{lemma}

\begin{proof}
En prenant la $h$-torsion dans la suite exacte du lemme
\ref{models-lemma-regular-model-pic}, on obtient l'exactitude de
$0\to\Pic(X)[h]\to\Pic(C)[h]$, car $h$ est premier à $d$.
L'application du lemme \ref{models-lemma-canonical-map-of-pic} donne une
application $\Pic(C)[h]\to\Pic(T)[h]$ qui annule les éléments de
$\Pic(X)[h]$. Réciproquement, si $\xi\in\Pic(C)[h]$ s'envoie sur zéro
dans $\Pic(T)[h]$, on peut trouver un $\mathcal{O}_X$-module inversible
$\mathcal{L}$ tel que $\deg(\mathcal{L}|_{C_i})=0$ pour tout $i$ et dont
la restriction à $C$ est $\xi$. Alors $\mathcal{L}^{\otimes h}$ est de
$d$-torsion d'après le lemme \ref{models-lemma-characterize-trivial}.
Soit $d'$ un entier tel que $dd'\equiv1\bmod h$. Un tel entier existe car
$h$ et $d$ sont premiers entre eux. Alors
$\mathcal{L}^{\otimes dd'}$ est un faisceau inversible de $h$-torsion sur
$X$ dont la restriction à $C$ est $\xi$.
\end{proof}

\begin{lemma}
\label{models-lemma-torsion-embeds}
Dans la situation \ref{models-situation-regular-model}, soit $h$ un entier premier
à la caractéristique de $k$. Alors l'application
$$
\Pic(X)[h] \longrightarrow \Pic((X_k)_{red})[h]
$$
est injective.
\end{lemma}

\begin{proof}
Observons que $X\times_{\Spec(R)}\Spec(R/\pi^n)$ est un épaississement
d'ordre fini de $(X_k)_{red}$ (cela résulte par exemple de Cohomologie des
schémas, lemme \ref{coherent-lemma-power-ideal-kills-sheaf}). Ainsi,
l'application canonique
$\Pic(X\times_{\Spec(R)}\Spec(R/\pi^n))\to\Pic((X_k)_{red})$ identifie les
sous-groupes de $h$-torsion d'après Compléments sur les morphismes, lemme
\ref{more-morphisms-lemma-torsion-pic-thickening}, et notre hypothèse sur
$h$. Par conséquent, si $\mathcal{L}$ est un faisceau inversible de
$h$-torsion sur $X$ dont la restriction à $(X_k)_{red}$ est triviale, alors
la restriction de $\mathcal{L}$ à
$X\times_{\Spec(R)}\Spec(R/\pi^n)$ est triviale pour tout $n$.
On obtient
\begin{align*}
H^0(X, \mathcal{L})^\wedge
& =
\lim H^0(X \times_{\Spec(R)} \Spec(R/\pi^n),
\mathcal{L}|_{X \times_{\Spec(R)} \Spec(R/\pi^n)}) \\
& \cong
\lim H^0(X \times_{\Spec(R)} \Spec(R/\pi^n),
\mathcal{O}_{X \times_{\Spec(R)} \Spec(R/\pi^n)}) \\
& =
R^\wedge
\end{align*}
On utilise ici le théorème des fonctions formelles (Cohomologie des schémas,
théorème \ref{coherent-theorem-formal-functions}) pour la première et la
dernière égalité, et, par exemple, Compléments d'algèbre, lemme
\ref{more-algebra-lemma-isomorphic-completions}, pour l'isomorphisme du
milieu. Puisque $H^0(X,\mathcal{L})$ est un $R$-module fini et que $R$ est
un anneau de valuation discrète, cela signifie que
$H^0(X,\mathcal{L})$ est libre de rang $1$ comme $R$-module.
Soit $s\in H^0(X,\mathcal{L})$ un élément de base. En remontant les
isomorphismes ci-dessus, on voit alors que
$s|_{X\times_{\Spec(R)}\Spec(R/\pi^n)}$ est une trivialisation pour tout
$n$. Comme le lieu des zéros de $s$ est fermé dans $X$ et que
$X\to\Spec(R)$ est propre, on conclut que le lieu des zéros de $s$ est vide,
comme voulu.
\end{proof}






\section{Réduction semi-stable}
\label{models-section-semistable-reduction}

\noindent
Dans cette section, nous définissons avec soin ce que nous entendons par
réduction semi-stable.

\begin{example}
\label{models-example-blowup}
Soit $R$ un anneau de valuation discrète d'uniformisante $\pi$.
Pour $n\geq0$, considérons le morphisme d'anneaux
$$
R \longrightarrow A = R[x, y]/(xy - \pi^n)
$$
Posons $X=\Spec(A)$ et $S=\Spec(R)$. Si $n=0$, le morphisme
$X\to S$ est lisse. Pour tout $n$, le morphisme $X\to S$ est à
singularités au plus nodales de dimension relative $1$, au sens de Courbes
algébriques, section \ref{curves-section-families-nodal}. Si $n=1$, alors
$X$ est régulier ; mais si $n>1$, $X$ n'est pas régulier, car $(x,y)$
n'engendre plus l'idéal maximal $\mathfrak m=(\pi,x,y)$. Pour améliorer la
situation lorsque $n>1$, considérons l'éclatement
$b:X'\to X$ de $X$ en $\mathfrak m$. Voir Diviseurs, section
\ref{divisors-section-blowing-up}. Par construction, $X'$ est recouvert par
trois ouverts affines correspondant aux algèbres d'éclatement
$A[\frac{\mathfrak m}{\pi}]$, $A[\frac{\mathfrak m}{x}]$ et
$A[\frac{\mathfrak m}{y}]$.

\medskip\noindent
L'algèbre $A[\frac{\mathfrak m}{\pi}]$ possède les générateurs
$x'=x/\pi$ et $y'=y/\pi$, et $x'y'=\pi^{n-2}$. Cette partie de $X'$ est
donc le spectre de $R[x',y'](x'y'-\pi^{n-2})$.

\medskip\noindent
L'algèbre $A[\frac{\mathfrak m}{x}]$ possède les générateurs $x$ et
$u=\pi/x$, soumis à la relation $xu-\pi$. Notons que cet anneau contient
$y/x=\pi^n/x^2=u^2\pi^{n-2}$. Cette partie de $X'$ est donc régulière.

\medskip\noindent
Par symétrie, le cas de l'algèbre $A[\frac{\mathfrak m}{y}]$ est identique
à celui de $A[\frac{\mathfrak m}{x}]$.

\medskip\noindent
On voit ainsi que $X'\to S$ est à singularités au plus nodales de dimension
relative $1$ et que $X'$ est régulier, sauf en un point qui possède un
voisinage ouvert affine exactement comme ci-dessus, mais avec $n$ remplacé
par $n-2$. Par récurrence sur $n$, on conclut qu'il existe une suite
d'éclatements en des points fermés
$$
X_{\lfloor n/2 \rfloor} \to \ldots \to X_1 \to X_0 = X
$$
telle que $X_{\lfloor n/2 \rfloor}\to S$ soit à singularités au plus
nodales de dimension relative $1$ et que
$X_{\lfloor n/2 \rfloor}$ soit régulier.
\end{example}

\begin{lemma}
\label{models-lemma-etale-local-at-worst-nodal}
Soit $R$ un anneau de valuation discrète. Soit $X$ un schéma à singularités
au plus nodales de dimension relative $1$ sur $R$. Soit $x\in X$ un point
de la fibre spéciale de $X$ sur $R$. Il existe alors un diagramme commutatif
$$
\xymatrix{
X \ar[d] &
U \ar[r] \ar[d] \ar[l] &
\Spec(A) \ar[dl] \\
\Spec(R) &
\Spec(R') \ar[l]
}
$$
où $R\subset R'$ est une extension étale d'anneaux de valuation discrète,
le morphisme $U\to X$ est étale, le morphisme $U\to\Spec(A)$ est étale,
il existe un point $x'\in U$ d'image $x$, et
$$
A = R'[u, v]/(uv)
\quad\text{ou}\quad
A = R'[u, v]/(uv - \pi^n)
$$
où $n\geq0$ et $\pi\in R'$ est une uniformisante.
\end{lemma}

\begin{proof}
Nous avons déjà démontré ce lemme dans une généralité bien plus grande ;
voir Courbes algébriques, lemme
\ref{curves-lemma-etale-local-structure-nodal-family}. Il suffit ici de
traduire l'énoncé qui y est donné sous la forme de celui ci-dessus.

\medskip\noindent
Tout d'abord, si $X\to\Spec(R)$ est lisse en $x$, on peut trouver un
morphisme étale $U\to\mathbf{A}^1_R=\Spec(R[u])$ pour un voisinage ouvert
affine $U\subset X$ de $x$. C'est le lemme
\ref{morphisms-lemma-smooth-etale-over-affine-space} de Morphismes.
Après avoir, si nécessaire, remplacé la coordonnée $u$ par $u+1$, on peut
supposer que $x$ s'envoie sur un point de l'ouvert standard
$D(u)\subset\mathbf{A}^1_R$. Alors $D(u)=\Spec(A)$, avec
$A=R[u,v]/(uv-1)$, et le résultat est vrai dans ce cas.

\medskip\noindent
Supposons ensuite que $x$ soit un point singulier de la fibre. On peut
alors appliquer Courbes algébriques, lemme
\ref{curves-lemma-etale-local-structure-nodal-family}, pour obtenir un
diagramme
$$
\xymatrix{
X \ar[d] &
U \ar[rr] \ar[l] \ar[rd] & &
W \ar[r] \ar[ld] &
\Spec(\mathbf{Z}[u, v, a]/(uv - a)) \ar[d] \\
\Spec(R) & &
V \ar[ll] \ar[rr] & & \Spec(\mathbf{Z}[a])
}
$$
ayant toutes les propriétés mentionnées dans l'énoncé du lemme cité.
Soit $x'\in U$ le point d'image $x$ fourni par ce lemme. Rétrécissons
d'abord $V$ en un voisinage affine de l'image de $x'$. Écrivons
$V=\Spec(R')$. Alors $R\to R'$ est étale. Puisque $R$ est un anneau de
valuation discrète, $R'$ est un produit fini de domaines de Dedekind
quasi-locaux (utiliser Compléments d'algèbre, lemme
\ref{more-algebra-lemma-Dedekind-etale-extension}). On trouve donc, par
exemple à l'aide de l'évitement des idéaux premiers, un ouvert standard
$D(f)\subset V=\Spec(R')$ contenant l'image de $x'$ et tel que $R'_f$
soit un anneau de valuation discrète. En remplaçant $R'$ par $R'_f$, on se
ramène à la situation où $V=\Spec(R')$, avec $R\subset R'$ une extension
étale d'anneaux de valuation discrète (les extensions d'anneaux de valuation
discrète sont définies dans Compléments d'algèbre, définition
\ref{more-algebra-definition-extension-discrete-valuation-rings}).

\medskip\noindent
Le morphisme $V\to\Spec(\mathbf{Z}[a])$ est déterminé par l'image $h$ de
$a$ dans $R'$. Alors $W=\Spec(R'[u,v]/(uv-h))$. Ainsi le lemme est vrai
avec $A=R'[u,v]/(uv-h)$. Si $h=0$, on obtient clairement le premier cas
du lemme. Si $h\not = 0$, on peut écrire $h=\epsilon\pi^n$ pour un certain
$n\geq0$, où $\epsilon$ est une unité de $R'$. En effectuant le changement
de coordonnées $u_{new}=\epsilon u$ et $v_{new}=v$, on obtient le second
type d'isomorphisme de $A$ indiqué dans le lemme.
\end{proof}

\begin{lemma}
\label{models-lemma-blowup-at-worst-nodal}
Soit $R$ un anneau de valuation discrète. Soit $X$ un schéma quasi-compact
à singularités au plus nodales de dimension relative $1$ sur $R$, dont la
fibre générique est lisse. Il existe alors $m\geq0$ et une suite
$$
X_m \to \ldots \to X_1 \to X_0 = X
$$
tels que
\begin{enumerate}
\item $X_{i+1}\to X_i$ soit l'éclatement d'un point fermé $x_i$ où $X_i$
est singulier ;
\item $X_i\to\Spec(R)$ soit à singularités au plus nodales de dimension
relative $1$ ;
\item $X_m$ soit régulier.
\end{enumerate}
\end{lemma}

\noindent
Un énoncé légèrement plus fort, également vrai, serait que, quelle que soit
la manière dont on éclate les points singuliers, on aboutit finalement à
une résolution et que tous les éclatements intermédiaires sont à
singularités au plus nodales de dimension relative $1$ sur $R$.

\begin{proof}
Comme $X$ est quasi-compact, la fibre spéciale $X_k$ est quasi-compacte.
Puisque les singularités de $X_k$ sont au plus nodales, $X_k$ possède un nombre fini
de nœuds et est lisse sur $k$ ailleurs. Comme $X\to\Spec(R)$ est plat à
fibre générique lisse, il s'ensuit que $X$ est lisse sur $R$ hors du nombre
fini de nœuds de $X_k$ (utiliser Morphismes, lemme
\ref{morphisms-lemma-smooth-at-point}). Il s'ensuit que $X$ est régulier en
tout point, sauf éventuellement aux nœuds de sa fibre spéciale (voir
Algèbre, lemme \ref{algebra-lemma-regular-goes-up}). Soit $x\in X$ un tel
nœud. Choisissons un diagramme
$$
\xymatrix{
X \ar[d] &
U \ar[r] \ar[d] \ar[l] &
\Spec(A) \ar[dl] \\
\Spec(R) &
\Spec(R') \ar[l]
}
$$
comme dans le lemme \ref{models-lemma-etale-local-at-worst-nodal}.
Le cas $A=R'[u,v]/(uv)$ ne peut pas se produire, car il impliquerait que la
fibre générique de $X/R$ est singulière (petit détail omis). Ainsi
$A=R'[u,v]/(uv-\pi^n)$ pour un certain $n\geq0$. Puisque $x$ est un point
singulier, on a $n\geq2$ ; voir la discussion de l'exemple
\ref{models-example-blowup}.

\medskip\noindent
Après avoir rétréci $U$, on peut supposer qu'il existe un unique point
$u\in U$ d'image $x$. Soit $w\in\Spec(A)$ l'image de $u$. On peut aussi
supposer que $u$ est l'unique point de $U$ s'envoyant sur $w$. Puisque les
deux flèches horizontales sont étales, on voit que $u$, considéré comme
sous-schéma fermé de $U$, est l'image réciproque schématique de $x\in X$
et celle de $w\in\Spec(A)$. Comme l'éclatement commute au changement de
base plat (Diviseurs, lemme
\ref{divisors-lemma-flat-base-change-blowing-up}), on obtient un diagramme
commutatif
$$
\xymatrix{
X' \ar[d] &
U' \ar[l] \ar[d] \ar[r] &
W' \ar[d] \\
X & U \ar[l] \ar[r] & \Spec(A)
}
$$
à carrés cartésiens, où les flèches verticales sont les éclatements de
$x,u,w$ dans $X,U,\Spec(A)$. Le schéma $W'$ a été décrit dans l'exemple
\ref{models-example-blowup}. Nous y avons vu que $W'$ est à singularités au plus
nodales de dimension relative $1$ sur $R'$. Ainsi, $W'$ est à singularités
au plus nodales de dimension relative $1$ sur $R$ (Courbes algébriques, lemme
\ref{curves-lemma-nodal-family-postcompose-etale}). Par conséquent, $U'$ est
à singularités au plus nodales de dimension relative $1$ sur $R$ (voir
Courbes algébriques, lemme
\ref{curves-lemma-nodal-family-etale-local-source}). Puisque $X'\to X$ est
un isomorphisme au-dessus du complémentaire de $x$, on conclut que
$X'/R$ possède la même propriété (en appliquant encore Courbes algébriques,
lemme \ref{curves-lemma-nodal-family-etale-local-source}).

\medskip\noindent
Enfin, il faut montrer qu'après un nombre fini de ces éclatements on atteint
un modèle régulier $X_m$. C'est assez clair, car l'« invariant » $n$
diminue de $2$ sous l'éclatement décrit ci-dessus ; voir le calcul de
l'exemple \ref{models-example-blowup}. Toutefois, pour éviter de définir
précisément cet invariant et d'en établir les propriétés, raisonnons comme
suit. Si $n=2$, alors $W'$ est régulier, donc $X'$ est régulier en tous les
points situés au-dessus de $x$, et le nombre de points singuliers de $X$ a
diminué de $1$. Si $n>2$, l'unique point singulier $w'$ de $W'$ situé
au-dessus de $w$ vérifie $\kappa(w)=\kappa(w')$. Ainsi $U'$ possède un
unique point singulier $u'$ situé au-dessus de $u$, avec
$\kappa(u)=\kappa(u')$. Il est clair que cela implique que $X'$ possède un
unique point singulier $x'$ situé au-dessus de $x$, à savoir l'image de
$u'$. On peut donc raisonner exactement comme ci-dessus et obtenir un
diagramme commutatif
$$
\xymatrix{
X'' \ar[d] &
U'' \ar[l] \ar[d] \ar[r] &
W'' \ar[d] \\
X' & U' \ar[l] \ar[r] & W'
}
$$
à carrés cartésiens, où les flèches verticales sont les éclatements de
$x',u',w'$ dans $X',U',W'$. En poursuivant ainsi, on obtient une suite
compatible d'éclatements qui s'arrête après $\lfloor n/2\rfloor$ étapes.
À l'issue de ce processus, le schéma
$X^{(\lfloor n/2\rfloor)}$ possède un point singulier de moins que $X$.
Une récurrence sur le nombre de points singuliers achève la preuve.
\end{proof}

\begin{lemma}
\label{models-lemma-blowdown-at-worst-nodal}
Soit $R$ un anneau de valuation discrète, de corps des fractions $K$ et de
corps résiduel $k$. Supposons que $X\to\Spec(R)$ soit à singularités au
plus nodales de dimension relative $1$ sur $R$. Soit $X\to X'$ la
contraction d'une courbe exceptionnelle de première espèce $E\subset X$.
Alors $X'$ est à singularités au plus nodales de dimension relative $1$
sur $R$.
\end{lemma}

\begin{proof}
Soit $x'\in X'$ l'image de $E$. Il suffit de voir que
$X'\to\Spec(R)$ est à singularités au plus nodales de dimension relative
$1$ dans un voisinage de $x'$. La fibre fermée de $X\to\Spec(R)$ est
réduite ; ainsi $\pi\in R$ s'annule à l'ordre $1$ sur $E$. Cela implique
immédiatement que $\pi$, considéré comme élément de
$\mathfrak m_{x'}\subset\mathcal{O}_{X',x'}$, n'appartient pas à
$\mathfrak m_{x'}^2$. Comme $\mathcal{O}_{X',x'}$ est régulier de
dimension $2$ (par la définition des contractions dans Résolution des
surfaces, section \ref{resolve-section-minus-one}), il s'ensuit que
$\mathcal{O}_{X'_k,x'}$ est régulier de dimension $1$ (Algèbre, lemme
\ref{algebra-lemma-regular-ring-CM}). D'autre part, la courbe $E$ doit
rencontrer au moins une autre composante, disons $C$, de la fibre fermée
$X_k$. Soit $x\in E\cap C$. Alors $x$ est un nœud de la fibre spéciale
$X_k$, donc $\kappa(x)/k$ est finie séparable ; voir Courbes algébriques,
lemme \ref{curves-lemma-nodal}. Comme $x\mapsto x'$, on conclut que
$\kappa(x')/k$ est finie séparable. D'après Algèbre, lemme
\ref{algebra-lemma-separable-smooth}, le morphisme
$X'_k\to\Spec(k)$ est donc lisse dans un voisinage ouvert de $x'$.
Combiné à la platitude, cela montre que $X'\to\Spec(R)$ est lisse dans un
voisinage de $x'$ (Morphismes, lemme
\ref{morphisms-lemma-smooth-at-point}). Ceci achève la preuve, car un
morphisme lisse de dimension relative $1$ est à singularités au plus
nodales de dimension relative $1$ (Courbes algébriques, lemme
\ref{curves-lemma-smooth-relative-dimension-1}).
\end{proof}

\begin{lemma}
\label{models-lemma-semistable}
Soit $R$ un anneau de valuation discrète de corps des fractions $K$.
Soit $C$ une courbe projective lisse sur $K$ telle que
$H^0(C,\mathcal{O}_C)=K$. Les conditions suivantes sont équivalentes :
\begin{enumerate}
\item il existe un modèle propre de $C$ qui est à singularités au plus
nodales de dimension relative $1$ sur $R$ ;
\item il existe un modèle minimal de $C$ qui est à singularités au plus
nodales de dimension relative $1$ sur $R$ ;
\item tout modèle minimal de $C$ est à singularités au plus nodales de
dimension relative $1$ sur $R$.
\end{enumerate}
\end{lemma}

\begin{proof}
Pour donner un sens à cet énoncé, rappelons qu'un modèle minimal est, par
définition, un modèle propre régulier sans courbes exceptionnelles de
première espèce (définition \ref{models-definition-minimal-model}), que les modèles
minimaux existent (proposition \ref{models-proposition-exists-minimal-model}), et
qu'ils sont uniques si le genre de $C$ est $>0$ (lemme
\ref{models-lemma-minimal-model-unique}). Cela étant, les implications
(2) $\Rightarrow$ (1) et (3) $\Rightarrow$ (2) sont claires.

\medskip\noindent
Supposons (1). Soit $X$ un modèle propre de $C$ à singularités au plus
nodales de dimension relative $1$ sur $R$. En appliquant le lemme
\ref{models-lemma-blowup-at-worst-nodal}, on voit que l'on peut en outre supposer
$X$ régulier. Soit
$$
X = X_m \to X_{m - 1} \to \ldots \to X_1 \to X_0
$$
comme dans le lemme \ref{models-lemma-pre-exists-minimal-model}. Le lemme
\ref{models-lemma-blowdown-at-worst-nodal} et une récurrence montrent alors que
$X_0$ est à singularités au plus nodales de dimension relative $1$ sur
$R$.

\medskip\noindent
Pour achever la preuve, il reste à montrer que (2) implique (3). Cela est
clair si le genre de $C$ est $>0$, puisque le modèle minimal est alors
unique (voir la discussion ci-dessus). D'autre part, si le modèle minimal
n'est pas unique, alors $X\to\Spec(R)$ est lisse pour tout modèle minimal,
car sa fibre spéciale est isomorphe à $\mathbf{P}^1_k$ d'après le lemme
\ref{models-lemma-nonuniqueness}.
\end{proof}

\begin{definition}
\label{models-definition-semistable}
Soit $R$ un anneau de valuation discrète de corps des fractions $K$.
Soit $C$ une courbe projective lisse sur $K$ telle que
$H^0(C,\mathcal{O}_C)=K$. On dit que $C$ a une
{\it réduction semi-stable} si les conditions équivalentes du lemme
\ref{models-lemma-semistable} sont satisfaites.
\end{definition}

\begin{lemma}
\label{models-lemma-good}
Soit $R$ un anneau de valuation discrète de corps des fractions $K$.
Soit $C$ une courbe projective lisse sur $K$ telle que
$H^0(C,\mathcal{O}_C)=K$. Les conditions suivantes sont équivalentes :
\begin{enumerate}
\item il existe un modèle propre et lisse de $C$ ;
\item il existe un modèle minimal de $C$ qui est lisse sur $R$ ;
\item tout modèle minimal est lisse sur $R$.
\end{enumerate}
\end{lemma}

\begin{proof}
Si $X$ est un modèle propre et lisse, sa fibre spéciale est connexe
(lemme \ref{models-lemma-regular-model-connected}) et lisse, donc irréductible.
Cela implique immédiatement que ce modèle est minimal. Ainsi (1) implique (2).
Pour achever la preuve, il reste à montrer que (2) implique (3). Cela est
clair si le genre de $C$ est $>0$, puisque le modèle minimal est alors
unique (lemme \ref{models-lemma-minimal-model-unique}). D'autre part, si le modèle
minimal n'est pas unique, le morphisme $X\to\Spec(R)$ est lisse pour tout
modèle minimal, car sa fibre spéciale est isomorphe à $\mathbf{P}^1_k$
d'après le lemme \ref{models-lemma-nonuniqueness}.
\end{proof}

\begin{definition}
\label{models-definition-good}
Soit $R$ un anneau de valuation discrète de corps des fractions $K$.
Soit $C$ une courbe projective lisse sur $K$ telle que
$H^0(C,\mathcal{O}_C)=K$. On dit que $C$ a une {\it bonne réduction} si
les conditions équivalentes du lemme \ref{models-lemma-good} sont satisfaites.
\end{definition}








\section{Réduction semi-stable en genre zéro}
\label{models-section-semistable-reduction-genus-zero}

\noindent
Dans cette section, nous démontrons le théorème de réduction semi-stable
(théorème \ref{models-theorem-semistable-reduction}) pour les courbes de genre
zéro.

\medskip\noindent
Soit $R$ un anneau de valuation discrète de corps des fractions $K$.
Soit $C$ une courbe projective lisse sur $K$ telle que
$H^0(C,\mathcal{O}_C)=K$. Si le genre de $C$ est $0$, alors $C$ est
isomorphe à une conique ; voir Courbes algébriques, lemme
\ref{curves-lemma-genus-zero}. Il existe donc une extension finie séparable
$K'/K$ de degré au plus $2$ telle que $C(K') \not = \emptyset$ ; voir Courbes
algébriques, lemme
\ref{curves-lemma-smooth-plane-curve-point-over-separable}. Soit
$R'\subset K'$ la clôture intégrale de $R$ ; voir la discussion de
Compléments d'algèbre, remarque
\ref{more-algebra-remark-finite-separable-extension}. Nous montrerons que
$C_{K'}$ a une réduction semi-stable sur $R'_{\mathfrak m}$ pour tout idéal
maximal $\mathfrak m$ de $R'$ (dans le cas présent, il y en a bien sûr au
plus deux). En remplaçant $R$ par $R'_{\mathfrak m}$ et $C$ par $C_{K'}$,
on se ramène au cas du paragraphe suivant.

\medskip\noindent
Dans ce paragraphe, $R$ est un anneau de valuation discrète de corps des
fractions $K$, $C$ est une courbe projective lisse sur $K$ telle que
$H^0(C,\mathcal{O}_C)=K$, de genre $0$, et $C$ possède un point
$K$-rationnel. Dans ce cas, $C\cong\mathbf{P}^1_K$ d'après Courbes
algébriques, proposition \ref{curves-proposition-projective-line}. On peut
donc prendre $\mathbf{P}^1_R$ comme modèle, ce qui montre que $C$ a à la
fois bonne réduction et réduction semi-stable.

\begin{example}
\label{models-example-extension-necessary-genus-zero}
Soit $R=\mathbf{R}[[\pi]]$ et considérons le schéma
$$
X = V(T_1^2 + T_2^2 - \pi T_0^2) \subset \mathbf{P}^2_R
$$
Le changement de base de $X$ à $\mathbf{C}[[\pi]]$ est isomorphe au schéma
défini dans l'exemple \ref{models-example-nonunique-in-genus-zero}, car on a la
factorisation $T_1^2+T_2^2=(T_1+iT_2)(T_1-iT_2)$ sur $\mathbf{C}$.
Ainsi $X$ est régulier et sa fibre spéciale est irréductible mais
singulière ; par conséquent, $X$ est l'unique modèle minimal de sa fibre
générique (utiliser le lemme \ref{models-lemma-nonuniqueness}). Il s'ensuit qu'une
extension est nécessaire même en genre $0$.
\end{example}



\section{Réduction semi-stable en genre un}
\label{models-section-semistable-reduction-genus-one}

\noindent
Dans cette section, nous démontrons le théorème de réduction semi-stable
(théorème \ref{models-theorem-semistable-reduction}) pour les courbes de genre un.
Nous conseillons au lecteur de lire d'abord la preuve dans le cas du genre
$\geq2$ (section \ref{models-section-semistable-reduction-genus-at-least-two}).
Nous utiliserons autant que possible la classification des types numériques
minimaux de genre $1$ donnée dans le lemme \ref{models-lemma-genus-one}.

\medskip\noindent
Soit $R$ un anneau de valuation discrète de corps des fractions $K$.
Soit $C$ une courbe projective lisse sur $K$ telle que
$H^0(C,\mathcal{O}_C)=K$. Supposons que le genre de $C$ soit $1$.
Choisissons un nombre premier $\ell\geq7$ distinct de la caractéristique
de $k$. Choisissons une extension finie séparable $K'/K$ telle que
$C(K') \not = \emptyset$ et
$\Pic(C_{K'})[\ell]\cong(\mathbf{Z}/\ell\mathbf{Z})^{\oplus2}$.
Voir Courbes algébriques, lemme
\ref{curves-lemma-torsion-picard-becomes-visible}. Soit $R'\subset K'$ la
clôture intégrale de $R$ ; voir la discussion de Compléments d'algèbre,
remarque \ref{more-algebra-remark-finite-separable-extension}. On peut
remplacer $R$ par $R'_{\mathfrak m}$ pour un idéal maximal $\mathfrak m$
de $R'$, et $C$ par $C_{K'}$. On se ramène ainsi au cas du paragraphe
suivant.

\medskip\noindent
Dans le reste de cette section, $R$ est un anneau de valuation discrète de
corps des fractions $K$, $C$ est une courbe projective lisse sur $K$ telle
que $H^0(C,\mathcal{O}_C)=K$, de genre $1$, possédant un point
$K$-rationnel et telle que
$\Pic(C)[\ell]\cong(\mathbf{Z}/\ell\mathbf{Z})^{\oplus2}$ pour un nombre
premier $\ell\geq7$ distinct de la caractéristique de $k$. Nous allons
montrer que $C$ a une réduction semi-stable.

\medskip\noindent
Soit $X$ un modèle minimal de $C$ ; voir la proposition
\ref{models-proposition-exists-minimal-model}. Soit
$T=(n, m_i, (a_{ij}), w_i, g_i)$ le type numérique associé à $X$
(définition \ref{models-definition-numerical-type-model}). Alors $T$ est un type
numérique minimal (lemme \ref{models-lemma-numerical-type-minimal-model}).
Puisque $C$ possède un point rationnel, il existe un $i$ tel que
$m_i=w_i=1$, d'après le lemme
\ref{models-lemma-numerical-type-rational-point}. En examinant la classification
des types numériques minimaux de genre $1$ du lemme
\ref{models-lemma-genus-one}, on voit que $m=w=1$ et que les cas
(\ref{models-item-up4}),
(\ref{models-item-equal-down3}),
(\ref{models-item-up-up}),
(\ref{models-item-down-up}),
(\ref{models-item-up-equal-up}),
(\ref{models-item-down-equal-up}),
(\ref{models-item-triple-with-down}),
(\ref{models-item-equal-equal-down-equal}),
(\ref{models-item-up-equal-equal-up}),
(\ref{models-item-down-equal-equal-up}),
(\ref{models-item-triple-extended-down}),
(\ref{models-item-up-chain-equal-up}),
(\ref{models-item-down-chain-equal-up}) et
(\ref{models-item-Dn-extended-down}) sont exclus (car il n'existe aucun indice
pour lequel $w_i$ et $m_i$ sont tous deux égaux à $1$). Soit $e$ le nombre de paires $(i,j)$
avec $i<j$ et $a_{ij}>0$. Pour les cas restants, on a
\begin{enumerate}
\item[(A)] $e=n-1$ dans les cas
(\ref{models-item-one}),
(\ref{models-item-two-cycle}),
(\ref{models-item-equal-up3}),
(\ref{models-item-up-down}),
(\ref{models-item-up-equal-down}),
(\ref{models-item-triple-with-up}),
(\ref{models-item-equal-equal-up-equal}),
(\ref{models-item-up-equal-equal-down}),
(\ref{models-item-quadruple}),
(\ref{models-item-triple-extended-up}),
(\ref{models-item-up-chain-equal-down}),
(\ref{models-item-Dn-extended-up}),
(\ref{models-item-double-triple}),
(\ref{models-item-E6-completed}),
(\ref{models-item-E7-completed}) et
(\ref{models-item-E8-completed}) ;
\item[(B)] $e=n$ dans les cas
(\ref{models-item-three-cycle}),
(\ref{models-item-four-cycle}),
(\ref{models-item-five-cycle}) et
(\ref{models-item-n-cycle}).
\end{enumerate}
Nous traiterons ces deux cas séparément.

\medskip\noindent
Cas (A). Dans ce cas, $\Pic(T)[\ell]$ est trivial (le groupe de Picard d'un
type numérique est défini dans la section \ref{models-section-picard-group}).
Cette annulation résulte de ce que $\Pic(T)\subset\Coker(A)$ (lemme
\ref{models-lemma-picard-T-and-A}) et de ce que $\Coker(A)[\ell]=0$ d'après le
lemme \ref{models-lemma-recurring-symmetric-integer} et le fait que $\ell$ a été
choisi premier aux $a_{ij}$ et aux $m_i$. D'après les lemmes
\ref{models-lemma-sequence-torsion} et \ref{models-lemma-torsion-embeds}, on conclut qu'il
existe un plongement
$$
(\mathbf{Z}/\ell \mathbf{Z})^{\oplus 2} \subset \Pic((X_k)_{red})[\ell].
$$
D'après Courbes algébriques, lemme
\ref{curves-lemma-bound-torsion-simple}, on obtient
$$
2 \leq \dim_k H^1((X_k)_{red}, \mathcal{O}_{(X_k)_{red}}) +
g_{geom}((X_k)_{red}/k)
$$
D'après Courbes algébriques, lemmes
\ref{curves-lemma-bound-geometric-genus} et
\ref{curves-lemma-bound-geometric-genus-curve}, on a
$g_{geom}((X_k)_{red}/k)\leq\sum w_ig_i$. Les hypothèses du lemme
\ref{models-lemma-genus-reduction-smaller} sont satisfaites d'après le lemme
\ref{models-lemma-numerical-type-rational-point}, et l'on conclut que
$\dim_kH^1((X_k)_{red},\mathcal{O}_{(X_k)_{red}})\leq g=1$.
En combinant ces inégalités, on obtient
$$
2 \leq 1 + \sum w_i g_i
$$
L'examen de la liste montre que le type numérique est donné par
$n=1$, $w_1=m_1=g_1=1$. Comme toutes les inégalités sont des égalités,
on voit que $g_{geom}(C_1/k)=1$. D'autre part, on sait que $C_1$ possède
un point $k$-rationnel $x$ tel que $C_1\to\Spec(k)$ soit lisse en $x$.
Il s'ensuit que $C_1$ est géométriquement intègre (Variétés, lemme
\ref{varieties-lemma-variety-with-smooth-rational-point}). Ainsi
$g_{geom}(C_1/k)=1$ est à la fois égal au genre de la normalisation de
$C_{1,\overline{k}}$ et au genre de $C_{1,\overline{k}}$. Il s'ensuit que
le morphisme de normalisation
$C_{1,\overline{k}}^\nu\to C_{1,\overline{k}}$ est un isomorphisme
(Courbes algébriques, lemme \ref{curves-lemma-genus-normalization}). On
conclut que $C_1$ est lisse sur $k$, comme voulu.

\medskip\noindent
Cas (B). Ici, on conclut seulement qu'il existe un plongement
$$
\mathbf{Z}/\ell \mathbf{Z} \subset \Pic(X_k)[\ell]
$$
La classification des types montre que $m_i=w_i=1$ et $g_i=0$ pour tout
$i$. Chaque $C_i$ est donc une courbe de genre zéro sur $k$. De plus, pour
tout $i$, il existe un $j$ tel que $C_i\cap C_j$ soit un point
$k$-rationnel. Par conséquent, $C_i\cong\mathbf{P}^1_k$ d'après Courbes
algébriques, proposition \ref{curves-proposition-projective-line}.
En particulier, puisque $X_k$ est la réunion schématique des $C_i$, on voit
que $X_{\overline{k}}$ est la réunion schématique des
$C_{i,\overline{k}}$. Ainsi $X_{\overline{k}}$ est un schéma propre réduit
connexe de dimension $1$ sur $\overline{k}$ tel que
$\dim_{\overline{k}}H^1(X_{\overline{k}},
\mathcal{O}_{X_{\overline{k}}})=1$. De plus, d'après Variétés, lemme
\ref{varieties-lemma-change-fields-pic}, et ce qui précède, on a encore
$$
\dim_{\mathbf{F}_\ell}(\Pic(X_{\overline{k}}) \geq 1
$$
D'après Courbes algébriques, proposition
\ref{curves-proposition-torsion-picard-reduced-proper}, on voit que
$X_{\overline{k}}$ n'a que des singularités en croisement multiple.
Mais $X_k$ est de Gorenstein (lemme \ref{models-lemma-gorenstein}), donc
$X_{\overline{k}}$ l'est aussi (Dualité pour les schémas, lemme
\ref{duality-lemma-gorenstein-base-change}). On conclut que les singularités
de $X_{\overline{k}}$ sont au plus nodales d'après Courbes algébriques,
lemme \ref{curves-lemma-multicross-gorenstein-is-nodal}. Ceci achève la
preuve dans le cas (B).

\begin{example}
\label{models-example-extension-necessary-genus-one}
Soit $k$ un corps algébriquement clos. Soit $Z$ une courbe projective lisse
sur $k$ de genre positif $g$. Soit $n\geq1$ un entier premier à la
caractéristique de $k$. Soit $\mathcal{L}$ un $\mathcal{O}_Z$-module
inversible d'ordre $n$ ; voir Courbes algébriques, lemme
\ref{curves-lemma-torsion-picard-smooth-projective}. Choisissons un
isomorphisme $\varphi:\mathcal{L}^{\otimes n}\to\mathcal{O}_Z$.
Posons $R=k[[\pi]]$, de corps des fractions $K=k((\pi))$. Notons $Z_R$ le
changement de base de $Z$ à $R$, et $\mathcal{L}_R$ l'image réciproque de
$\mathcal{L}$ sur $Z_R$. Considérons le morphisme fini plat
$$
p : X \longrightarrow Z_R
$$
tel que
$$
p_*\mathcal{O}_X =
\text{Sym}^*_{\mathcal{O}_{Z_R}}(\mathcal{L}_R)/(\varphi - \pi) =
\mathcal{O}_{Z_R} \oplus \mathcal{L}_R \oplus
\mathcal{L}_R^{\otimes 2} \oplus \ldots \oplus \mathcal{L}_R^{\otimes n - 1}
$$
Plus précisément, si $U=\Spec(A)\subset Z$ est un ouvert affine tel que
$\mathcal{L}|_U$ soit trivialisé par une section $s$ vérifiant
$\varphi(s^{\otimes n})=f$ (où $f$ est une unité), alors
$$
p^{-1}(U_R) = \Spec\left(
(A \otimes_R R[[\pi]])[x]/(x^n - \pi f)
\right)
$$
Le lecteur vérifie que le morphisme $X_K\to Z_K$ entre les fibres
génériques est fini étale. La description du faisceau structural montre
que $H^0(X,\mathcal{O}_X)=R$ et $H^0(X_K,\mathcal{O}_{X_K})=K$.
D'après Riemann--Hurwitz (Courbes algébriques, lemme
\ref{curves-lemma-rhe}), le genre de $X_K$ est $n(g-1)+1$. En particulier,
$X_K$ est de genre $1$ si $Z$ est de genre $1$. D'autre part, le schéma
$X$ est régulier d'après l'équation locale ci-dessus, et la fibre spéciale
$X_k$ est, comme diviseur de Cartier effectif, égale à $n$ fois sa fibre
spéciale réduite. Il s'ensuit que toute extension finie $K'/K$ sur laquelle
$X_K$ acquiert une réduction semi-stable doit être ramifiée avec un indice
de ramification au moins $n$ (certains détails sont omis). Il n'existe donc
pas de borne universelle pour le degré d'une extension sur laquelle une
courbe de genre $1$ acquiert une réduction semi-stable.
\end{example}






\section{Réduction semi-stable en genre au moins deux}
\label{models-section-semistable-reduction-genus-at-least-two}

\noindent
Dans cette section, nous démontrons le théorème de réduction semi-stable
(théorème \ref{models-theorem-semistable-reduction}) pour les courbes de genre
$\geq2$. Fixons $g\geq2$.

\medskip\noindent
Soit $R$ un anneau de valuation discrète de corps des fractions $K$.
Soit $C$ une courbe projective lisse sur $K$ telle que
$H^0(C,\mathcal{O}_C)=K$. Supposons que le genre de $C$ soit $g$.
Choisissons un nombre premier $\ell>768g$ distinct de la caractéristique de
$k$. Choisissons une extension finie séparable $K'/K$ telle que
$C(K') \not = \emptyset$ et
$\Pic(C_{K'})[\ell]\cong(\mathbf{Z}/\ell\mathbf{Z})^{\oplus2g}$.
Voir Courbes algébriques, lemme
\ref{curves-lemma-torsion-picard-becomes-visible}. Soit $R'\subset K'$ la
clôture intégrale de $R$ ; voir la discussion de Compléments d'algèbre,
remarque \ref{more-algebra-remark-finite-separable-extension}. On peut
remplacer $R$ par $R'_{\mathfrak m}$ pour un idéal maximal $\mathfrak m$
de $R'$, et $C$ par $C_{K'}$. On se ramène ainsi au cas du paragraphe
suivant.

\medskip\noindent
Dans le reste de cette section, $R$ est un anneau de valuation discrète de
corps des fractions $K$, $C$ est une courbe projective lisse sur $K$ telle
que $H^0(C,\mathcal{O}_C)=K$, de genre $g$, possédant un point
$K$-rationnel et telle que
$\Pic(C)[\ell]\cong(\mathbf{Z}/\ell\mathbf{Z})^{\oplus2g}$ pour un nombre
premier $\ell\geq768g$ distinct de la caractéristique de $k$. Nous allons
montrer que $C$ a une réduction semi-stable.

\medskip\noindent
Dans le reste de cette section, nous utiliserons sans autre mention le fait
que les conclusions du lemme \ref{models-lemma-numerical-type-rational-point}
sont vraies.

\medskip\noindent
Soit $X$ un modèle minimal de $C$ ; voir la proposition
\ref{models-proposition-exists-minimal-model}. Soit
$T=(n, m_i, (a_{ij}), w_i, g_i)$ le type numérique associé à $X$
(définition \ref{models-definition-numerical-type-model}). Alors $T$ est un type
numérique minimal de genre $g$ (lemme
\ref{models-lemma-numerical-type-minimal-model}). D'après la proposition
\ref{models-proposition-bound-picard-group}, on a
$$
\dim_{\mathbf{F}_\ell} \Pic(T)[\ell] \leq g_{top}
$$
D'après les lemmes \ref{models-lemma-sequence-torsion} et
\ref{models-lemma-torsion-embeds}, on conclut qu'il existe un plongement
$$
(\mathbf{Z}/\ell \mathbf{Z})^{\oplus 2g - g_{top}}
\subset \Pic((X_k)_{red})[\ell].
$$
D'après Courbes algébriques, lemme
\ref{curves-lemma-bound-torsion-simple}, on obtient
$$
2g - g_{top} \leq
\dim_k H^1((X_k)_{red}, \mathcal{O}_{(X_k)_{red}}) + g_{geom}(X_k/k)
$$
D'après les lemmes \ref{models-lemma-genus-reduction-smaller} et
\ref{models-lemma-genus-reduction-bigger-than}, on a
$$
g \geq \dim_k H^1((X_k)_{red}, \mathcal{O}_{(X_k)_{red}}) \geq
g_{top} + g_{geom}(X_k/k)
$$
La théorie élémentaire des nombres montre que ces $3$ inégalités ne
peuvent être satisfaites que si elles sont toutes des égalités.
L'examen du lemme \ref{models-lemma-genus-reduction-smaller} montre que
$m_i=1$ pour tout $i$. Celui du lemme
\ref{models-lemma-equality-genus-reduction-bigger-than} montre que toute composante
irréductible de $X_k$ est lisse sur $k$.

\medskip\noindent
En particulier, puisque $X_k$ est la réunion schématique de ses composantes
irréductibles $C_i$, on voit que $X_{\overline{k}}$ est la réunion
schématique des $C_{i,\overline{k}}$. Ainsi $X_{\overline{k}}$ est un
schéma propre réduit connexe de dimension $1$ sur $\overline{k}$, tel que
$\dim_{\overline{k}}H^1(X_{\overline{k}},
\mathcal{O}_{X_{\overline{k}}})=g$. De plus, d'après Variétés, lemme
\ref{varieties-lemma-change-fields-pic}, et ce qui précède, on a encore
$$
\dim_{\mathbf{F}_\ell}(\Pic(X_{\overline{k}})[\ell])
\geq 2g - g_{top} =
\dim_{\overline{k}} H^1(X_{\overline{k}}, \mathcal{O}_{X_{\overline{k}}})
+ g_{geom}(X_{\overline{k}})
$$
D'après Courbes algébriques, proposition
\ref{curves-proposition-torsion-picard-reduced-proper}, on voit que
$X_{\overline{k}}$ n'a que des singularités en croisement multiple.
Mais $X_k$ est de Gorenstein (lemme \ref{models-lemma-gorenstein}), donc
$X_{\overline{k}}$ l'est aussi (Dualité pour les schémas, lemme
\ref{duality-lemma-gorenstein-base-change}). On conclut que les singularités
de $X_{\overline{k}}$ sont au plus nodales d'après Courbes algébriques,
lemme \ref{curves-lemma-multicross-gorenstein-is-nodal}. Ceci achève la
preuve.









\section{Réduction semi-stable des courbes}
\label{models-section-semistable-reduction-theorem}

\noindent
Dans cette section, nous achevons la preuve du théorème. Pour $g\geq2$,
soient $768g<\ell'<\ell$ les deux premiers nombres premiers $>768g$, et
posons
\begin{equation}
\label{models-equation-bound}
B_g = (2g - 2)(\ell^{2g})!
\end{equation}
La forme précise de $B_g$ n'a pas d'importance ; le point essentiel est
qu'il ne dépend que de $g$.

\begin{theorem}
\label{models-theorem-semistable-reduction}
\begin{reference}
\cite[Corollary 2.7]{DM}
\end{reference}
Soit $R$ un anneau de valuation discrète de corps des fractions $K$.
Soit $C$ une courbe projective lisse sur $K$ telle que
$H^0(C,\mathcal{O}_C)=K$. Il existe alors une extension d'anneaux de
valuation discrète $R\subset R'$ qui induit une extension finie séparable
$K'/K$ des corps des fractions et telle que $C_{K'}$ ait une réduction
semi-stable. Plus précisément :
\begin{enumerate}
\item Si le genre de $C$ est zéro, il existe une extension séparable de
degré $2$, notée $K'/K$, telle que $C_{K'}\cong\mathbf{P}^1_{K'}$ ; ainsi
$C_{K'}$ est isomorphe à la fibre générique du schéma projectif lisse
$\mathbf{P}^1_{R'}$ sur la clôture intégrale $R'$ de $R$ dans $K'$.
\item Si le genre de $C$ est un, il existe une extension finie séparable
$K'/K$ telle que $C_{K'}$ ait une réduction semi-stable sur
$R'_\mathfrak m$ pour tout idéal maximal $\mathfrak m$ de la clôture
intégrale $R'$ de $R$ dans $K'$. De plus, la fibre spéciale du modèle
minimal (unique) de $C_{K'}$ sur $R'_\mathfrak m$ est soit une courbe
lisse de genre un, soit un cycle de courbes rationnelles.
\item Si le genre $g$ de $C$ est supérieur à un, il existe une extension
finie séparable $K'/K$ de degré au plus $B_g$ (\ref{models-equation-bound}) telle
que $C_{K'}$ ait une réduction semi-stable sur $R'_\mathfrak m$ pour
tout idéal maximal $\mathfrak m$ de la clôture intégrale $R'$ de $R$ dans
$K'$.
\end{enumerate}
\end{theorem}

\begin{proof}
Pour le cas du genre zéro, voir la section
\ref{models-section-semistable-reduction-genus-zero}. Pour celui du genre un, voir
la section \ref{models-section-semistable-reduction-genus-one}. Pour celui du
genre supérieur à un, voir la section
\ref{models-section-semistable-reduction-genus-at-least-two}. Pour obtenir la borne
sur le degré $[K':K]$, on peut utiliser la borne sur le degré de l'extension
nécessaire pour rendre visible toute la $\ell$-torsion ou toute la
$\ell'$-torsion, démontrée dans Courbes algébriques, lemme
\ref{curves-lemma-torsion-picard-becomes-visible}. (La raison d'utiliser
$\ell$ et $\ell'$ est qu'il faut éviter la caractéristique du corps
résiduel $k$.)
\end{proof}

\begin{remark}[Amélioration de la borne]
\label{models-remark-improving-bound}
Les résultats de la littérature suggèrent que l'on peut améliorer la borne
donnée dans l'énoncé du théorème \ref{models-theorem-semistable-reduction}.
Par exemple, \cite{DM} montre que la réduction semi-stable de $C$ équivaut
à celle de sa jacobienne lorsque le corps résiduel est parfait, et cela est
vraisemblablement encore vrai pour un corps résiduel quelconque.
Une variété abélienne a une réduction semi-stable si l'action de Galois sur
sa $\ell$-torsion est triviale pour un $\ell\geq3$ quelconque distinct de
la caractéristique résiduelle. On peut donc vraisemblablement choisir
$\ell=5$ dans la formule (\ref{models-equation-bound}) définissant $B_g$ (mais la
preuve demanderait beaucoup plus de travail ; si nous en avons un jour
besoin, nous donnerons ici un énoncé précis et une preuve).
\end{remark}









% OK, start here.
%

\chapter{Foncteurs et morphismes}



\phantomsection
\label{functors-section-phantom}


\section{Introduction}
\label{functors-section-introduction}

\noindent
Soient $X$ et $Y$ des schémas. Ce chapitre porte sur les rapports
entre les foncteurs $\QCoh(\mathcal{O}_Y) \to \QCoh(\mathcal{O}_X)$ et les
morphismes de schémas $X \to Y$. Plus généralement, nous étudions les
rapports entre $\QCoh(\mathcal{O}_X)$ et $X$ ou, si $X$ est noethérien,
entre $\textit{Coh}(\mathcal{O}_X)$ et $X$.
Ces rapports ont été étudiés dans \cite{Gabriel}.






\section{Foncteurs sur les catégories de modules}
\label{functors-section-preliminary}

\noindent
Pour un anneau $A$, notons $\text{Mod}^{fp}_A$ la catégorie des
$A$-modules de présentation finie.

\begin{lemma}
\label{functors-lemma-functor-on-fp-modules}
Soit $A$ un anneau. Soit $\mathcal{B}$ une catégorie admettant des limites
inductives filtrantes. Soit $F : \text{Mod}^{fp}_A \to \mathcal{B}$ un foncteur. Alors $F$
se prolonge de manière unique en un foncteur $F' : \text{Mod}_A \to \mathcal{B}$
qui commute aux limites inductives filtrantes.
\end{lemma}

\begin{proof}
Cela résulte de
Catégories, lemme \ref{categories-lemma-extend-functor-by-colim}.
Pour voir que ce lemme s'applique, observons que les
$A$-modules de présentation finie sont des
objets compacts au sens catégorique de $\text{Mod}_A$, d'après
Algèbre, lemme \ref{algebra-lemma-characterize-finitely-presented-module-hom}.
De plus, tout $A$-module est limite inductive filtrante
de $A$-modules de présentation finie, d'après
Algèbre, lemme \ref{algebra-lemma-module-colimit-fp}.
\end{proof}

\noindent
Si une catégorie $\mathcal{B}$ est additive et admet des limites inductives filtrantes,
alors $\mathcal{B}$ admet des sommes directes quelconques : toute somme directe s'écrit
comme limite inductive filtrante de sommes directes finies.

\begin{lemma}
\label{functors-lemma-functor-on-fp-modules-additive}
Soient $A$, $\mathcal{B}$, $F$ comme dans le lemme \ref{functors-lemma-functor-on-fp-modules}.
Supposons que $\mathcal{B}$ et $F$ soient additifs. Alors
$F'$ est additif et commute aux sommes directes quelconques.
\end{lemma}

\begin{proof}
Pour montrer que $F'$ est additif, il suffit de montrer
que $F'(M) \oplus F'(M') \to F'(M \oplus M')$ est un isomorphisme pour
tous $A$-modules $M$, $M'$, voir
Homologie, lemme \ref{homology-lemma-additive-functor}.
Écrivons $M = \colim_i M_i$ et $M' = \colim_j M'_j$ comme limites inductives filtrantes
de $A$-modules de présentation finie $M_i$. Alors
$F'(M) = \colim_i F(M_i)$, $F'(M') = \colim_j F(M'_j)$, et
\begin{align*}
F'(M \oplus M')
& =
F'(\colim_{i, j} M_i \oplus M'_j) \\
& =
\colim_{i, j} F(M_i \oplus M'_j) \\
& =
\colim_{i, j} F(M_i) \oplus F(M'_j) \\
& =
F'(M) \oplus F'(M')
\end{align*}
comme voulu. Pour montrer que $F'$ commute aux sommes directes, supposons
que $M = \bigoplus_{i \in I} M_i$. Alors
$M = \colim_{I' \subset I\text{ fini}} \bigoplus_{i \in I'} M_i$
est une limite inductive filtrante. On obtient
\begin{align*}
F'(M)
& =
\colim_{I' \subset I\text{ fini}}
F'(\bigoplus\nolimits_{i \in I'} M_i) \\
& =
\colim_{I' \subset I\text{ fini}}
\bigoplus\nolimits_{i \in I'} F'(M_i) \\
& =
\bigoplus\nolimits_{i \in I} F'(M_i)
\end{align*}
La deuxième égalité résulte de l'additivité de $F'$ déjà démontrée.
\end{proof}

\noindent
Si une catégorie $\mathcal{B}$ est additive, admet des limites inductives filtrantes et
des conoyaux, alors $\mathcal{B}$ admet des limites inductives quelconques, voir
la discussion ci-dessus et Catégories, lemme
\ref{categories-lemma-colimits-coproducts-coequalizers}.

\begin{lemma}
\label{functors-lemma-functor-on-fp-modules-right-exact}
Soient $A$, $\mathcal{B}$, $F$ comme dans le lemme \ref{functors-lemma-functor-on-fp-modules}.
Supposons que $\mathcal{B}$ soit additive, admette des conoyaux, et que $F$ soit exact à droite. Alors
$F'$ est additif, exact à droite et commute aux sommes directes quelconques.
\end{lemma}

\begin{proof}
Puisque $F$ est exact à droite, $F$ commute aux sommes de deux objets, qui sont
représentées par les sommes directes. Ainsi $F$ est additif, d'après
Homologie, lemme \ref{homology-lemma-additive-functor}.
Donc $F'$ est additif et commute aux sommes directes, d'après le
lemme \ref{functors-lemma-functor-on-fp-modules-additive}.
Nous invitons le lecteur à démontrer lui-même que $F'$ est exact à droite
avant de lire la démonstration ci-dessous.

\medskip\noindent
Pour montrer que $F'$ est exact à droite, il suffit de montrer que $F'$ commute
aux conoyaux de doubles flèches, voir
Catégories, lemme \ref{categories-lemma-characterize-right-exact}.
Or, si $a, b : K \to L$ sont des homomorphismes de $A$-modules, le
conoyau de la double flèche formée de $a$ et $b$ est le conoyau de $a - b : K \to L$.
Soit donc $K \to L \to M \to 0$ une suite exacte
de $A$-modules. Nous devons montrer que, dans
$$
F'(K) \to F'(L) \to F'(M) \to 0
$$
la deuxième flèche est un conoyau de la première dans $\mathcal{B}$
(si $\mathcal{B}$ était abélienne, nous dirions que la suite ci-dessus
est exacte).
Écrivons $M = \colim_{i \in I} M_i$ comme limite inductive filtrante de
$A$-modules de présentation finie, voir
Algèbre, lemme \ref{algebra-lemma-module-colimit-fp}.
Posons $L_i = L \times_M M_i$. 
On obtient un système de suites exactes $K \to L_i \to M_i \to 0$ indexé par $I$.
Puisque les limites inductives commutent aux limites inductives, d'après
Catégories, lemme \ref{categories-lemma-colimits-commute},
et puisque les conoyaux sont des cas particuliers des conoyaux de doubles flèches,
il suffit de montrer que $F'(L_i) \to F(M_i)$ est un conoyau de
$F'(K) \to F'(L_i)$ dans $\mathcal{B}$ pour tout $i \in I$. Autrement dit, nous pouvons
supposer $M$ de présentation finie. Écrivons $L = \colim_{i \in I} L_i$
comme limite inductive filtrante de $A$-modules de présentation finie,
de telle sorte que chaque $L_i$ s'envoie surjectivement sur $M$.
Posons $K_i = K \times_L L_i$. On obtient un système de suites exactes courtes
$K_i \to L_i \to M \to 0$ indexé par $I$.
En répétant l'argument précédent, on se ramène à montrer que
$F(L_i) \to F(M_i)$ est un conoyau de
$F'(K) \to F(L_i)$ dans $\mathcal{B}$ pour tout $i \in I$.
Autrement dit, nous pouvons supposer que
$L$ et $M$ sont des $A$-modules de présentation finie.
Dans ce cas, le module $\Ker(L \to M)$ est de type fini
(Algèbre, lemme \ref{algebra-lemma-extension}).
On peut donc écrire $K = \colim_{i \in I} K_i$ comme limite inductive filtrante
de $A$-modules de présentation finie s'envoyant tous surjectivement sur $\Ker(L \to M)$.
On obtient un système de suites exactes courtes
$K_i \to L \to M \to 0$ indexé par $I$.
En répétant l'argument précédent, on se ramène à montrer que
$F(L) \to F(M)$ est un conoyau de
$F(K_i) \to F(L)$ dans $\mathcal{B}$ pour tout $i \in I$.
Autrement dit, nous pouvons supposer que $K$, $L$ et $M$
sont des $A$-modules de présentation finie. Ce dernier cas résulte
de l'hypothèse que $F$ est exact à droite.
\end{proof}

\noindent
Si une catégorie $\mathcal{B}$ est additive et admet des noyaux,
alors $\mathcal{B}$ admet des limites projectives finies. En effet, les produits finis
sont des sommes directes, qui existent, et le noyau de la double flèche $a, b : L \to M$
est le noyau de $a - b : K \to L$, qui existe. Toutes les limites projectives
finies existent donc, d'après Catégories, lemme \ref{categories-lemma-finite-limits-exist}.

\begin{lemma}
\label{functors-lemma-functor-on-fp-modules-left-exact}
Soient $A$, $\mathcal{B}$, $F$ comme dans le lemme \ref{functors-lemma-functor-on-fp-modules}.
Supposons que $A$ soit un anneau cohérent (Algèbre, définition
\ref{algebra-definition-coherent}), que $\mathcal{B}$ soit additive, admette des noyaux,
que les limites inductives filtrantes commutent à la formation des noyaux et que $F$ soit exact à gauche. Alors
$F'$ est additif, exact à gauche et commute aux sommes directes quelconques.
\end{lemma}

\begin{proof}
Puisque $A$ est cohérent, la catégorie $\text{Mod}^{fp}_A$ est abélienne,
avec les mêmes noyaux et conoyaux que dans $\text{Mod}_A$, voir
Algèbre, lemmes \ref{algebra-lemma-coherent-ring} et
\ref{algebra-lemma-coherent}. Toutes les limites projectives finies existent donc dans
$\text{Mod}^{fp}_A$, et
Catégories, définition \ref{categories-definition-exact}, s'applique.
Puisque $F$ est exact à gauche, $F$ commute aux produits de deux objets, qui sont
représentés par les sommes directes. Ainsi $F$ est additif, d'après
Homologie, lemme \ref{homology-lemma-additive-functor}.
Donc $F'$ est additif et commute aux sommes directes, d'après le
lemme \ref{functors-lemma-functor-on-fp-modules-additive}.
Nous invitons le lecteur à démontrer lui-même que $F'$ est exact à gauche
avant de lire la démonstration ci-dessous.

\medskip\noindent
Pour montrer que $F'$ est exact à gauche, il suffit de montrer que $F'$ commute
aux noyaux de doubles flèches, voir
Catégories, lemme \ref{categories-lemma-characterize-left-exact}.
Or, si $a, b : L \to M$ sont des homomorphismes de $A$-modules, le
noyau de la double flèche formée de $a$ et $b$ est le noyau de $a - b : L \to M$.
Soit donc $0 \to K \to L \to M$ une suite exacte
de $A$-modules. Nous devons montrer que, dans
$$
0 \to F'(K) \to F'(L) \to F'(M)
$$
la flèche $F'(K) \to F'(L)$ est un noyau de $F'(L) \to F'(M)$ dans $\mathcal{B}$
(si $\mathcal{B}$ était abélienne, nous dirions que la suite ci-dessus
est exacte).
Écrivons $M = \colim_{i \in I} M_i$ comme limite inductive filtrante de
$A$-modules de présentation finie, voir
Algèbre, lemme \ref{algebra-lemma-module-colimit-fp}.
Posons $L_i = L \times_M M_i$. 
On obtient un système de suites exactes $0 \to K \to L_i \to M_i$
indexé par $I$. Puisque, par hypothèse, les limites inductives filtrantes commutent à la formation des noyaux
dans $\mathcal{B}$, 
il suffit de montrer que $F'(K) \to F'(L_i)$ est un noyau de
$F'(L_i) \to F(M_i)$ dans $\mathcal{B}$ pour tout $i \in I$. Autrement dit, nous pouvons
supposer $M$ de présentation finie. Écrivons $L = \colim_{i \in I} L_i$
comme limite inductive filtrante de $A$-modules de présentation finie.
Posons $K_i = K \times_L L_i$. On obtient un système de suites exactes courtes
$0 \to K_i \to L_i \to M$ indexé par $I$.
En répétant l'argument précédent, on se ramène à montrer que
$F'(K_i) \to F(L_i)$ est un noyau de
$F(L_i) \to F(M)$ dans $\mathcal{B}$ pour tout $i \in I$.
Autrement dit, nous pouvons supposer que
$L$ et $M$ sont des $A$-modules de présentation finie.
Puisque $A$ est cohérent, le $A$-module $K = \Ker(L \to M)$ est
de présentation finie, car la catégorie des
$A$-modules de présentation finie est abélienne (voir les références ci-dessus).
Autrement dit, les trois modules $K$, $L$ et $M$
sont des $A$-modules de présentation finie. Ce dernier cas résulte
de l'hypothèse que $F$ est exact à gauche.
\end{proof}

\noindent
Si une catégorie $\mathcal{B}$ est additive et admet des conoyaux,
alors $\mathcal{B}$ admet des limites inductives finies. En effet, les sommes finies
sont des sommes directes, qui existent, et le conoyau de la double flèche $a, b : K \to L$
est le conoyau de $a - b : K \to L$, qui existe. Toutes les limites inductives
finies existent donc, d'après Catégories, lemme \ref{categories-lemma-colimits-exist}.

\begin{lemma}
\label{functors-lemma-functor-on-modules-fp}
Soit $A$ un anneau. Soit $\mathcal{B}$ une catégorie additive
admettant des conoyaux. Il existe une équivalence de catégories entre
\begin{enumerate}
\item la catégorie des foncteurs $F : \text{Mod}^{fp}_A \to \mathcal{B}$
exacts à droite, et
\item la catégorie des couples $(K, \kappa)$ où $K \in \Ob(\mathcal{B})$
et où $\kappa : A \to \text{End}_\mathcal{B}(K)$ est un homomorphisme d'anneaux,
\end{enumerate}
donnée par la règle qui à $F$ associe $F(A)$ muni de son action naturelle de $A$.
\end{lemma}

\begin{proof}
Soit $(K, \kappa)$ comme dans (2). Nous allons construire un foncteur
$F : \text{Mod}^{fp}_A \to \mathcal{B}$ tel que $F(A) = K$,
muni de l'action donnée de $A$ par $\kappa$. Pour un
entier $n \geq 0$, posons
$$
F(A^{\oplus n}) = K^{\oplus n}
$$
Étant donnée une application $A$-linéaire $\varphi : A^{\oplus m} \to A^{\oplus n}$
de matrice $(a_{ij}) \in \text{Mat}(n \times m, A)$, définissons
$$
F(\varphi) :
F(A^{\oplus m}) = K^{\oplus m}
\longrightarrow
K^{\oplus n} = F(A^{\oplus n})
$$
comme l'application de matrice $(\kappa(a_{ij}))$. Cela définit un foncteur
additif $F$ de la sous-catégorie pleine de
$\text{Mod}^{fp}_A$ d'objets $0$, $A$, $A^{\oplus 2}$, $\ldots$
dans $\mathcal{B}$ ; nous omettons la vérification.

\medskip\noindent
Pour chaque objet $M$ de $\text{Mod}^{fp}_A$, choisissons une présentation
$$
A^{\oplus m_M} \xrightarrow{\varphi_M} A^{\oplus n_M} \to M \to 0
$$
de $M$ comme $A$-module. Prenons la présentation triviale
$0 \to A^{\oplus n} \xrightarrow{1} A^{\oplus n} \to 0$ si $M = A^{\oplus n}$
(ce n'est pas nécessaire, mais simplifie l'exposé).
Pour chaque morphisme $f : M \to N$ de
$\text{Mod}^{fp}_A$, nous pouvons choisir un diagramme commutatif
\begin{equation}
\label{functors-equation-map}
\vcenter{
\xymatrix{
A^{\oplus m_M} \ar[r]_{\varphi_M} \ar[d]_{\psi_f} &
A^{\oplus n_M} \ar[r] \ar[d]_{\chi_f} &
M \ar[r] \ar[d]_f & 0 \\
A^{\oplus m_N} \ar[r]^{\varphi_N} &
A^{\oplus n_N} \ar[r] &
N \ar[r] & 0
}
}
\end{equation}
Ces choix étant faits, nous pouvons définir ce qui suit : pour un objet
$M$ de $\text{Mod}^{fp}_A$, posons
$$
F(M) = \Coker(F(\varphi_M) : F(A^{\oplus m_M}) \to F(A^{\oplus n_M}))
$$
et, pour un morphisme $f : M \to N$ de $\text{Mod}^{fp}_A$, posons
$$
F(f) = \text{le morphisme }F(M) \to F(N)\text{ induit par }
F(\psi_f)\text{ et }F(\chi_f)\text{ sur les conoyaux}
$$
Notons que cette règle prolonge le foncteur $F$ donné sur
la sous-catégorie pleine constituée des modules libres $A^{\oplus n}$.
Il reste à montrer que $F$ est un foncteur, que $F$ est additif
et que $F$ est exact à droite.

\medskip\noindent
Soit $f : M \to N$ un morphisme de $\text{Mod}^{fp}_A$. Montrons que le morphisme
$F(f)$ définie ci-dessus est indépendante du choix de $\psi_f$ et de $\chi_f$
dans (\ref{functors-equation-map}). Supposons en effet que
$$
\xymatrix{
A^{\oplus m_M} \ar[r]_{\varphi_M} \ar[d]_\psi &
A^{\oplus n_M} \ar[r] \ar[d]_\chi &
M \ar[r] \ar[d]_f & 0 \\
A^{\oplus m_N} \ar[r]^{\varphi_N} &
A^{\oplus n_N} \ar[r] &
N \ar[r] & 0
}
$$
soit également commutatif. Notons $F(f)' : F(M) \to F(N)$ le morphisme
induit par $F(\psi)$ et $F(\chi)$. À la vue des diagrammes
commutatifs, un argument élémentaire d'algèbre commutative fournit un homomorphisme
$\omega : A^{\oplus n_M} \to A^{\oplus m_N}$ telle que
$\chi = \chi_f + \varphi_N \circ \omega$. En appliquant $F$,
on trouve $F(\chi) = F(\chi_f) + F(\varphi_N) \circ F(\omega)$.
Puisque $F(N)$ est le conoyau de $F(\varphi_N)$, on voit
que le morphisme $F(A^{\oplus n_M}) \to F(M)$ égalise $F(f)$ et $F(f)'$.
Un conoyau étant un épimorphisme, on en conclut que $F(f) = F(f)'$.

\medskip\noindent
Montrons que $F$ est un foncteur. Observons d'abord que
$F(\text{id}_M) = \text{id}_{F(M)}$, car nous pouvons choisir
les identités pour $\psi_f$ et $\chi_f$ dans le diagramme ci-dessus
lorsque $f = \text{id}_M$. Supposons ensuite donnés
$f : M \to N$ et $g : L \to M$. Alors
$\psi = \psi_f \circ \psi_g$ et $\chi = \chi_f \circ \chi_g$
conviennent dans (\ref{functors-equation-map}) pour $f \circ g$.
Ils induisent donc le morphisme voulu, ce qui exprime précisément
que $F(f) \circ F(g) = F(f \circ g)$.

\medskip\noindent
Montrons que $F$ est additif. Supposons donnés
$f, g : M \to N$. Alors $\psi = \psi_f + \psi_g$ et
$\chi = \chi_f + \chi_g$ conviennent dans (\ref{functors-equation-map}) pour $f + g$.
Ils induisent donc le morphisme voulu, ce qui exprime précisément
que $F(f) + F(g) = F(f + g)$.

\medskip\noindent
Montrons enfin que $F$ est exact à droite. Il suffit de montrer que $F$
commute aux conoyaux de doubles flèches, voir
Catégories, lemme \ref{categories-lemma-characterize-right-exact}.
Pour cela, il suffit de montrer que $F$ commute aux conoyaux.
Soit $K \to L \to M \to 0$ une suite exacte de $A$-modules
avec $K$, $L$, $M$ de présentation finie. Puisque $F$ est un foncteur
additif, on obtient bien un complexe
$$
F(K) \to F(L) \to F(M) \to 0
$$
et il faut montrer que la deuxième flèche est le conoyau de la première
dans $\mathcal{B}$. En tout cas, on obtient un morphisme
$\Coker(F(K) \to F(L)) \to F(M)$.
Un argument élémentaire d'algèbre commutative fournit un diagramme commutatif
$$
\xymatrix{
A^{\oplus m_M} \ar[r]_{\varphi_M} \ar[d]_\psi &
A^{\oplus n_M} \ar[r] \ar[d]_\chi &
M \ar[r] \ar[d]_1 & 0 \\
K \ar[r] &
L \ar[r] &
M \ar[r] & 0
}
$$
En appliquant $F$ à ce diagramme et en utilisant la construction de $F(M)$ comme
conoyau de $F(\varphi_M)$, on trouve un morphisme
$F(M) \to \Coker(F(K) \to F(L))$ inverse à droite
du morphisme $\Coker(F(K) \to F(L)) \to F(M)$. Il en résulte d'abord
que $F(L) \to F(M)$ est toujours un épimorphisme. Ensuite, ce qui précède donne
une décomposition
$$
\Coker(F(K) \to F(L)) = F(M) \oplus E
$$
compatible à la fois avec
$F(M) \to \Coker(F(K) \to F(L))$ et $\Coker(F(K) \to F(L)) \to F(M)$.
Mais alors l'épimorphisme $p : F(L) \to E$ devient nul aussi bien
après composition avec $F(K) \to F(L)$ qu'après composition
avec $F(A^{n_M}) \to F(L)$. Or, puisque $K \oplus A^{n_M} \to L$
est surjectif (nous omettons l'argument algébrique), on en conclut que
$F(K \oplus A^{n_M}) \to F(L)$ est un épimorphisme (d'après ce qui précède),
d'où $E = 0$. Cela achève la démonstration.
\end{proof}

\begin{lemma}
\label{functors-lemma-functor-on-modules}
Soit $A$ un anneau. Soit $\mathcal{B}$ une catégorie additive
admettant des sommes directes quelconques et des conoyaux. Il existe une équivalence
de catégories entre
\begin{enumerate}
\item la catégorie des foncteurs $F : \text{Mod}_A \to \mathcal{B}$
exacts à droite et commutant aux sommes directes quelconques, et
\item la catégorie des couples $(K, \kappa)$ où $K \in \Ob(\mathcal{B})$
et où $\kappa : A \to \text{End}_\mathcal{B}(K)$ est un homomorphisme d'anneaux,
\end{enumerate}
donnée par la règle qui à $F$ associe $F(A)$ muni de son action naturelle de $A$.
\end{lemma}

\begin{proof}
Combinons les lemmes \ref{functors-lemma-functor-on-modules-fp} et
\ref{functors-lemma-functor-on-fp-modules-right-exact}.
\end{proof}






\section{Foncteurs entre catégories de modules}
\label{functors-section-functors}

\noindent
Le lemme suivant est l'archétype des résultats de ce chapitre.

\begin{lemma}
\label{functors-lemma-functor}
Soient $A$ et $B$ des anneaux. Soit $F : \text{Mod}_A \to \text{Mod}_B$
un foncteur. Les propriétés suivantes sont équivalentes :
\begin{enumerate}
\item $F$ est isomorphe au foncteur $M \mapsto M \otimes_A K$
pour un certain $A \otimes_\mathbf{Z} B$-module $K$ ;
\item $F$ est exact à droite et commute à toutes les sommes directes ;
\item $F$ commute à toutes les limites inductives ;
\item $F$ admet un adjoint à droite $G$.
\end{enumerate}
\end{lemma}

\begin{proof}
Si (1) est satisfaite, alors (4) l'est aussi, car un adjoint à droite de $M \mapsto M \otimes_A K$
est $N \mapsto \Hom_B(K, N)$, voir
Algèbre différentielle graduée, lemme \ref{dga-lemma-tensor-hom-adjunction}.
Si (4) est satisfaite, alors (3) l'est d'après Catégories, lemme \ref{categories-lemma-adjoint-exact}.
L'implication (3) $\Rightarrow$ (2) résulte immédiatement des définitions.

\medskip\noindent
Supposons (2). Nous allons démontrer (1). D'après la discussion de
Homologie, section \ref{homology-section-functors},
le foncteur $F$ est additif. Ainsi, $F$ induit
un homomorphisme d'anneaux $A \to \text{End}_B(F(M))$, $a \mapsto F(a \cdot \text{id}_M)$
pour tout $A$-module $M$. Nous en concluons que $F(M)$ est un
$A \otimes_\mathbf{Z} B$-module, fonctoriellement en $M$.
Posons $K = F(A)$. Définissons
$$
M \otimes_A K = M \otimes_A F(A) \longrightarrow F(M),
\quad m \otimes k \longmapsto F(\varphi_m)(k)
$$
Ici, $\varphi_m : A \to M$ envoie $a \to am$. La règle
$(m, k) \mapsto F(\varphi_m)(k)$ est $A$-bilinéaire (et $B$-linéaire
à droite), comme requis pour obtenir l'homomorphisme
$A \otimes_\mathbf{Z} B$-linéaire affichée.
Cette construction est fonctorielle en $M$ ; elle définit donc une transformation
de foncteurs $- \otimes_A K \to F(-)$ qui est un isomorphisme lorsqu'on
l'évalue en $A$. Pour tout $A$-module $M$, nous pouvons choisir une suite exacte
$$
\bigoplus\nolimits_{j \in J} A \to
\bigoplus\nolimits_{i \in I} A \to
M \to 0
$$
En utilisant les morphismes construits ci-dessus, nous obtenons un diagramme commutatif
$$
\xymatrix{
(\bigoplus\nolimits_{j \in J} A) \otimes_A K \ar[r] \ar[d] &
(\bigoplus\nolimits_{i \in I} A) \otimes_A K \ar[r] \ar[d] &
M \otimes_A K \ar[r] \ar[d] &
0 \\
F(\bigoplus\nolimits_{j \in J} A) \ar[r] &
F(\bigoplus\nolimits_{i \in I} A) \ar[r] &
F(M) \ar[r] &  0
}
$$
La ligne inférieure est exacte puisque $F$ est exact à droite.
La ligne supérieure est exacte puisque le produit tensoriel par $K$ est exact à droite.
Comme $F$ commute aux sommes directes, les deux flèches verticales de gauche
sont bijectives. Nous pouvons donc conclure.
\end{proof}

\begin{example}
\label{functors-example-functor-modules}
Soit $R$ un anneau. Soient $A$ et $B$ des $R$-algèbres. Soit $K$ un
$A \otimes_R B$-module. Nous pouvons alors considérer le foncteur
\begin{equation}
\label{functors-equation-FM-modules}
F : \text{Mod}_A \longrightarrow \text{Mod}_B,\quad
M \longmapsto M \otimes_A K
\end{equation}
Ce foncteur est $R$-linéaire, exact à droite,
commute aux sommes directes quelconques, commute
à toutes les limites inductives et admet un adjoint à droite (lemme \ref{functors-lemma-functor}).
\end{example}

\begin{lemma}
\label{functors-lemma-functor-modules}
Soit $R$ un anneau. Soient $A$ et $B$ des $R$-algèbres. Il existe une
équivalence de catégories entre
\begin{enumerate}
\item la catégorie des foncteurs $R$-linéaires
$F : \text{Mod}_A \to \text{Mod}_B$ qui
sont exacts à droite et commutent aux sommes directes quelconques, et
\item la catégorie $\text{Mod}_{A \otimes_R B}$,
\end{enumerate}
donnée en associant à $K$ le foncteur $F$ de (\ref{functors-equation-FM-modules}).
\end{lemma}

\begin{proof}
Soit $F$ un objet de la première catégorie. D'après le
lemme \ref{functors-lemma-functor}, nous pouvons supposer que $F(M) = M \otimes_A K$
fonctoriellement en $M$, pour un certain $A \otimes_\mathbf{Z} B$-module $K$.
La $R$-linéarité de $F$ implique immédiatement que la structure de
$A \otimes_\mathbf{Z} B$-module sur $K$ provient
d'une (unique) structure de $A \otimes_R B$-module sur $K$.
Nous voyons donc que l'association de $K$ à $F$ comme dans (\ref{functors-equation-FM-modules})
est essentiellement surjective.

\medskip\noindent
Pour démontrer que notre foncteur est pleinement fidèle, nous devons montrer que,
étant donnés des $A \otimes_R B$-modules $K$ et $K'$, toute transformation
$t : F \to F'$ entre les foncteurs correspondants provient
d'un unique $\varphi : K \to K'$. Puisque $K = F(A)$ et $K' = F'(A)$,
nous pouvons prendre pour $\varphi$ la valeur $t_A : F(A) \to F'(A)$
de $t$ en $A$. Cet homomorphisme est $A \otimes_R B$-linéaire d'après la
définition de la structure de $A \otimes B$-module sur $F(A)$
et $F'(A)$ donnée dans la démonstration du lemme \ref{functors-lemma-functor}.
\end{proof}

\begin{remark}
\label{functors-remark-composition}
Soit $R$ un anneau. Soient $A$, $B$, $C$ des $R$-algèbres.
Soient $F : \text{Mod}_A \to \text{Mod}_B$ et
$F' : \text{Mod}_B \to \text{Mod}_C$ des foncteurs
$R$-linéaires exacts à droite qui commutent aux sommes directes quelconques.
Si, par l'équivalence du lemme \ref{functors-lemma-functor-modules}, l'objet
$K$ de $\text{Mod}_{A \otimes_R B}$ correspond à $F$ et l'objet
$K'$ de $\text{Mod}_{B \otimes_R C}$ correspond à $F'$, alors
$K \otimes_B K'$, vu comme un objet de
$\text{Mod}_{A \otimes_R C}$, correspond à $F' \circ F$.
\end{remark}

\begin{remark}
\label{functors-remark-exact-flat}
Dans la situation du lemme \ref{functors-lemma-functor-modules},
supposons que $F$ corresponde à $K$. Alors
$F$ est exact $\Leftrightarrow$ $K$ est plat sur $A$.
\end{remark}

\begin{remark}
\label{functors-remark-finite}
Dans la situation du lemme \ref{functors-lemma-functor-modules},
supposons que $F$ corresponde à $K$. Alors
$F$ envoie les $A$-modules finis sur des $B$-modules finis
$\Leftrightarrow$ $K$ est fini comme $B$-module.
\end{remark}

\begin{remark}
\label{functors-remark-finite-presentation}
Dans la situation du lemme \ref{functors-lemma-functor-modules},
supposons que $F$ corresponde à $K$. Alors
$F$ envoie les $A$-modules de présentation finie sur des $B$-modules de présentation finie
$\Leftrightarrow$ $K$ est de présentation finie comme $B$-module.
\end{remark}

\begin{lemma}
\label{functors-lemma-functor-equivalence}
Soient $A$ et $B$ des anneaux. Si
$$
F : \text{Mod}_A \longrightarrow \text{Mod}_B
$$
est une équivalence de catégories, alors il existe un isomorphisme
d'anneaux $A \to B$ et un $B$-module inversible $L$ tels que
$F$ soit isomorphe au foncteur $M \mapsto (M \otimes_A B) \otimes_B L$.
\end{lemma}

\begin{proof}
Puisqu'une équivalence commute à toutes les limites inductives, nous voyons que le
lemme \ref{functors-lemma-functor} s'applique. Soit $K$ le
$A \otimes_\mathbf{Z} B$-module tel que $F$ soit
isomorphe au foncteur $M \mapsto M \otimes_A K$.
Soit $K'$ le $B \otimes_\mathbf{Z} A$-module tel qu'un
quasi-inverse de $F$ soit
isomorphe au foncteur $N \mapsto N \otimes_B K'$.
D'après la remarque \ref{functors-remark-composition} et le
lemme \ref{functors-lemma-functor-modules}, nous avons un isomorphisme
$$
\psi : K \otimes_B K' \longrightarrow A
$$
de $A \otimes_\mathbf{Z} A$-modules.
De même, nous avons un isomorphisme
$$
\psi' : K' \otimes_A K \longrightarrow B
$$
de $B \otimes_\mathbf{Z} B$-modules. Choisissons un élément
$\xi = \sum_{i = 1, \ldots, n} x_i \otimes y_i \in K \otimes_B K'$
tel que $\psi(\xi) = 1$. Considérons les isomorphismes
$$
K \xrightarrow{\psi^{-1} \otimes \text{id}_K}
K \otimes_B K' \otimes_A K \xrightarrow{\text{id}_K \otimes \psi'} K
$$
Leur composé est un isomorphisme donné par
$$
k \longmapsto \sum x_i \psi'(y_i \otimes k)
$$
Nous en concluons que cet automorphisme se factorise comme
$$
K \to B^{\oplus n} \to K
$$
en tant qu'homomorphisme de $B$-modules. Il s'ensuit que $K$ est un
$B$-module projectif fini.

\medskip\noindent
Nous affirmons que $K$ est inversible comme $B$-module. Cela revient
à demander que le rang de $K$ comme $B$-module ait la valeur constante $1$ ;
voir Plus sur l'algèbre, lemme \ref{more-algebra-lemma-invertible}, et
Algèbre, lemme \ref{algebra-lemma-finite-projective}.
Sinon, il existe un idéal maximal $\mathfrak m \subset B$
tel que soit (a) $K \otimes_B B/\mathfrak m = 0$, soit
(b) il existe une surjection $K \to (B/\mathfrak m)^{\oplus 2}$ de
$B$-modules. Le cas (a) est absurde puisque $K' \otimes_A K \otimes_B N = N$
pour tout $B$-module $N$. Le cas (b) impliquerait l'existence d'une surjection
$$
A = K \otimes_B K' \longrightarrow (B/\mathfrak m \otimes_B K')^{\oplus 2}
$$
de $A$-modules (à droite). C'est impossible, car la cible est un $A$-module
qui nécessite au moins deux générateurs : $B/\mathfrak m \otimes_B K'$
est non nul, car il est l'image du module non nul $B/\mathfrak m$ par
le quasi-inverse de $F$.

\medskip\noindent
Puisque $K$ est inversible comme $B$-module, nous avons $\Hom_B(K, K) = B$.
Comme $K = F(A)$, l'action de $A$ sur $K$ définit un isomorphisme d'anneaux
$A \to B$. Le lemme en résulte.
\end{proof}

\begin{lemma}
\label{functors-lemma-functor-equivalence-linear}
Soit $R$ un anneau. Soient $A$ et $B$ des $R$-algèbres. Si
$$
F : \text{Mod}_A \longrightarrow \text{Mod}_B
$$
est une équivalence $R$-linéaire de catégories, alors il existe
un isomorphisme $A \to B$ de $R$-algèbres et un $B$-module inversible $L$ tels que
$F$ soit isomorphe au foncteur $M \mapsto (M \otimes_A B) \otimes_B L$.
\end{lemma}

\begin{proof}
Le lemme \ref{functors-lemma-functor-equivalence} nous fournit $A \to B$ et $L$.
Pour achever la démonstration, il reste à montrer que la $R$-linéarité
de $F$ force $A \to B$ à être un homomorphisme de $R$-algèbres. Nous omettons les détails.
\end{proof}

\begin{remark}
\label{functors-remark-monoidal}
Soient $A$ et $B$ des anneaux. Munissons $\text{Mod}_A$ et $\text{Mod}_B$
de la structure monoïdale usuelle donnée par les produits tensoriels de modules.
Soit $F : \text{Mod}_A \to \text{Mod}_B$ un foncteur de
catégories monoïdales, voir
Catégories, définition \ref{categories-definition-functor-monoidal-categories}.
Voici quelques remarques :
\begin{enumerate}
\item Puisque $F(A)$ est une unité (d'après nos définitions), nous avons $F(A) = B$.
\item Nous obtenons une application multiplicative $\varphi : A \to B$
en envoyant $a \in A$ sur son action sur $F(A) = B$.
\item Prenons $A = B$ et $F(M) = M \otimes_A M$. Dans ce cas, $\varphi(a) = a^2$.
\item Si $F$ est additif, alors $\varphi$ est un homomorphisme d'anneaux.
\item Prenons $A = B = \mathbf{Z}$ et $F(M) = M/\text{torsion}$. Alors
$\varphi = \text{id}_\mathbf{Z}$, mais $F$ n'est pas le foncteur identité.
\item Si $F$ est exact à droite et commute aux sommes directes,
alors $F(M) = M \otimes_{A, \varphi} B$ d'après le lemme \ref{functors-lemma-functor}.
\end{enumerate}
Autrement dit, les homomorphismes d'anneaux $A \to B$ sont en bijection avec les classes d'isomorphisme
de foncteurs de catégories monoïdales $\text{Mod}_A \to \text{Mod}_B$
qui commutent à toutes les limites inductives.
\end{remark}




\section{Prolongement des foncteurs sur des catégories de modules}
\label{functors-section-functors-extend}

\noindent
Pour un anneau $A$, notons $\text{Mod}^{fp}_A$ la catégorie des
$A$-modules de présentation finie.

\begin{lemma}
\label{functors-lemma-functor-fp-modules}
Soient $A$ et $B$ des anneaux. Soit
$F : \text{Mod}^{fp}_A \to \text{Mod}^{fp}_B$ un foncteur.
Alors $F$ se prolonge de manière unique en un foncteur
$F' : \text{Mod}_A \to \text{Mod}_B$
qui commute aux limites inductives filtrantes.
\end{lemma}

\begin{proof}
C'est un cas particulier du lemme \ref{functors-lemma-functor-on-fp-modules}.
\end{proof}

\begin{remark}
\label{functors-remark-monoidal-extension}
Conservons $A$, $B$, $F$ et $F'$ comme dans le lemme \ref{functors-lemma-functor-fp-modules}.
Observons que le produit tensoriel de deux modules de présentation finie est
de présentation finie ; voir Algèbre, lemme \ref{algebra-lemma-tensor-finiteness}.
Nous pouvons donc munir $\text{Mod}^{fp}_A$, $\text{Mod}^{fp}_B$,
$\text{Mod}_A$ et $\text{Mod}_B$ de la structure monoïdale usuelle
donnée par les produits tensoriels de modules. Dans ce cas, si $F$ est
un foncteur de catégories monoïdales, alors $F'$ l'est aussi. Cela résulte immédiatement
du fait que les produits tensoriels de modules commutent aux limites inductives
filtrantes.
\end{remark}

\begin{lemma}
\label{functors-lemma-functor-fp-modules-exact}
Conservons $A$, $B$, $F$ et $F'$ comme dans le lemme \ref{functors-lemma-functor-fp-modules}.
\begin{enumerate}
\item Si $F$ est additif, alors $F'$ est additif et
commute aux sommes directes quelconques ;
\item si $F$ est exact à droite, alors $F'$ est exact à droite.
\end{enumerate}
\end{lemma}

\begin{proof}
Cela résulte des lemmes \ref{functors-lemma-functor-on-fp-modules-additive} et
\ref{functors-lemma-functor-on-fp-modules-right-exact}.
\end{proof}

\begin{remark}
\label{functors-remark-monoidal-extension-exact}
En combinant les remarques \ref{functors-remark-monoidal} et \ref{functors-remark-monoidal-extension}
et le lemme \ref{functors-lemma-functor-fp-modules-exact},
nous obtenons ce qui suit. Étant donnés des anneaux $A$ et $B$, l'ensemble des homomorphismes d'anneaux $A \to B$
est en bijection avec l'ensemble des classes d'isomorphisme
de foncteurs de catégories monoïdales $\text{Mod}^{fp}_A \to \text{Mod}^{fp}_B$
qui sont exacts à droite.
\end{remark}

\begin{lemma}
\label{functors-lemma-functor-fp-modules-left-exact}
Conservons $A$, $B$, $F$ et $F'$ comme dans le lemme \ref{functors-lemma-functor-fp-modules}.
Supposons que $A$ soit un anneau cohérent
(Algèbre, définition \ref{algebra-definition-coherent}).
Si $F$ est exact à gauche, alors $F'$ est exact à gauche.
\end{lemma}

\begin{proof}
C'est un cas particulier du lemme \ref{functors-lemma-functor-on-fp-modules-left-exact}.
\end{proof}

\noindent
Pour un anneau $A$, notons $\text{Mod}^{fg}_A$ la catégorie des
$A$-modules de type fini (aussi appelés $A$-modules finis).

\begin{lemma}
\label{functors-lemma-functor-finite-modules}
Soient $A$ et $B$ des anneaux noethériens. Soit
$F : \text{Mod}^{fg}_A \to \text{Mod}^{fg}_B$ un foncteur.
Alors $F$ se prolonge de manière unique en un foncteur $F' : \text{Mod}_A \to \text{Mod}_B$
qui commute aux limites inductives filtrantes. Si $F$ est additif, alors
$F'$ est additif et commute aux sommes directes quelconques.
Si $F$ est exact, exact à gauche ou exact à droite, alors $F'$ l'est aussi.
\end{lemma}

\begin{proof}
Voir les lemmes \ref{functors-lemma-functor-fp-modules-exact} et
\ref{functors-lemma-functor-fp-modules-left-exact}.
Utilisons aussi le fait que les $A$-modules finis sont des $A$-modules de présentation finie ;
voir Algèbre, lemme
\ref{algebra-lemma-Noetherian-finite-type-is-finite-presentation},
ainsi que le fait que les anneaux noethériens sont cohérents ; voir
Algèbre, lemme \ref{algebra-lemma-Noetherian-coherent}.
\end{proof}









\section{Foncteurs entre catégories de modules quasi-cohérents}
\label{functors-section-functor-quasi-coherent}

\noindent
Dans cette section, nous étudions brièvement les foncteurs entre catégories de
modules quasi-cohérents.

\begin{example}
\label{functors-example-functor-quasi-coherent}
Soit $R$ un anneau. Soient $X$ et $Y$ des
schémas sur $R$, avec $X$ quasi-compact et quasi-séparé.
Soit $\mathcal{K}$ un $\mathcal{O}_{X \times_R Y}$-module quasi-cohérent.
Nous pouvons alors considérer le foncteur
\begin{equation}
\label{functors-equation-FM-QCoh}
F : \QCoh(\mathcal{O}_X) \longrightarrow \QCoh(\mathcal{O}_Y),\quad
\mathcal{F} \longmapsto
\text{pr}_{2, *}(\text{pr}_1^*\mathcal{F}
\otimes_{\mathcal{O}_{X \times_R Y}} \mathcal{K})
\end{equation}
Le morphisme $\text{pr}_2$ est quasi-compact et quasi-séparé
(Schémas, lemmes \ref{schemes-lemma-quasi-compact-preserved-base-change}
et \ref{schemes-lemma-separated-permanence}). Par conséquent, l'image directe par
ce morphisme préserve les modules quasi-cohérents ; voir
Schémas, lemme \ref{schemes-lemma-push-forward-quasi-coherent}.
De plus, notre foncteur est $R$-linéaire et commute aux sommes directes quelconques ;
voir Cohomologie des schémas, lemme \ref{coherent-lemma-colimit-cohomology}.
\end{example}

\noindent
Le lemme suivant est une généralisation naturelle du
lemme \ref{functors-lemma-functor-modules}.

\begin{lemma}
\label{functors-lemma-functor-quasi-coherent-from-affine}
Soit $R$ un anneau. Soient $X$ et $Y$ des schémas sur $R$, avec $X$ affine.
Il existe une équivalence de catégories entre
\begin{enumerate}
\item la catégorie des foncteurs $R$-linéaires
$F : \QCoh(\mathcal{O}_X) \to \QCoh(\mathcal{O}_Y)$
qui sont exacts à droite et commutent aux sommes directes quelconques, et
\item la catégorie $\QCoh(\mathcal{O}_{X \times_R Y})$,
\end{enumerate}
donnée en associant à $\mathcal{K}$ le foncteur $F$ de (\ref{functors-equation-FM-QCoh}).
\end{lemma}

\begin{proof}
Soit $\mathcal{K}$ un objet de $\QCoh(\mathcal{O}_{X \times_R Y})$
et soit $F_\mathcal{K}$ le foncteur (\ref{functors-equation-FM-QCoh}). D'après la discussion de
l'exemple \ref{functors-example-functor-quasi-coherent}, nous savons déjà que
$F$ est $R$-linéaire et commute aux sommes directes quelconques.
Comme $\text{pr}_2 : X \times_R Y \to Y$ est affine
(Morphismes, lemme \ref{morphisms-lemma-base-change-affine}), le foncteur
$\text{pr}_{2, *}$ est exact ; voir Cohomologie des schémas, lemme
\ref{coherent-lemma-relative-affine-vanishing}.
Ainsi, $F$ est aussi exact à droite ; autrement dit, $F$ est comme en (1).

\medskip\noindent
Soit $F$ comme en (1). Écrivons $X = \Spec(A)$. Considérons le
$\mathcal{O}_Y$-module quasi-cohérent $\mathcal{G} = F(\mathcal{O}_X)$.
Le foncteur $F$ induit une application $R$-linéaire
$A \to \text{End}_{\mathcal{O}_Y}(\mathcal{G})$,
$a \mapsto F(a \cdot \text{id})$. Ainsi, $\mathcal{G}$ est un module sur le faisceau d'anneaux
$$
A \otimes_R \mathcal{O}_Y = \text{pr}_{2, *}\mathcal{O}_{X \times_R Y}
$$
D'après Morphismes, lemme \ref{morphisms-lemma-affine-equivalence-modules},
nous trouvons qu'il existe un unique module quasi-cohérent $\mathcal{K}$
sur $X \times_R Y$ tel que $F(\mathcal{O}_X) = \mathcal{G} =
\text{pr}_{2, *}\mathcal{K}$, de manière compatible avec l'action
de $A$ et de $\mathcal{O}_Y$. Notons $F_\mathcal{K}$ le foncteur
donné par (\ref{functors-equation-FM-QCoh}). Il existe une équivalence
$\text{Mod}_A \to \QCoh(\mathcal{O}_X)$ envoyant $A$ sur $\mathcal{O}_X$ ; voir
Schémas, lemme \ref{schemes-lemma-equivalence-quasi-coherent}.
Nous obtenons donc un isomorphisme $F \cong F_\mathcal{K}$ d'après le
lemme \ref{functors-lemma-functor-on-modules}, car nous avons un isomorphisme
$F(\mathcal{O}_X) \cong F_\mathcal{K}(\mathcal{O}_X)$ compatible avec
l'action de $A$ par construction.

\medskip\noindent
Cela montre que le foncteur qui associe à $\mathcal{K}$ le foncteur $F_\mathcal{K}$
est essentiellement surjectif. Nous omettons la vérification de la pleine fidélité.
\end{proof}

\begin{remark}
\label{functors-remark-affine-morphism}
Nous utiliserons ci-dessous que, pour un morphisme affine
$h : T \to S$, nous avons $h_*\mathcal{G} \otimes_{\mathcal{O}_S} \mathcal{H} =
h_*(\mathcal{G} \otimes_{\mathcal{O}_T} h^*\mathcal{H})$ pour
$\mathcal{G} \in \QCoh(\mathcal{O}_T)$ et
$\mathcal{H} \in \QCoh(\mathcal{O}_S)$. Cela résulte
immédiatement de la traduction en algèbre.
\end{remark}

\begin{lemma}
\label{functors-lemma-functor-quasi-coherent-from-affine-compose}
Dans le lemme \ref{functors-lemma-functor-quasi-coherent-from-affine}, supposons que $F$
corresponde à $\mathcal{K}$ dans $\QCoh(\mathcal{O}_{X \times_R Y})$.
Nous avons :
\begin{enumerate}
\item Si $f : X' \to X$ est un morphisme affine, alors $F \circ f_*$
correspond à $(f \times \text{id}_Y)^*\mathcal{K}$.
\item Si $g : Y' \to Y$ est un morphisme plat, alors $g^* \circ F$ correspond à
$(\text{id}_X \times g)^*\mathcal{K}$.
\item Si $j : V \to Y$ est une immersion ouverte, alors $j^* \circ F$
correspond à $\mathcal{K}|_{X \times_R V}$.
\end{enumerate}
\end{lemma}

\begin{proof}
Démonstration de (1). Considérons le diagramme commutatif
$$
\xymatrix{
X' \times_R Y \ar[rrd]^{\text{pr}'_2} \ar[rd]_{f \times \text{id}_Y}
\ar[dd]_{\text{pr}'_1} \\
& X \times_R Y \ar[r]_{\text{pr}_2} \ar[d]_{\text{pr}_1} & Y \\
X' \ar[r]^f & X
}
$$
Soit $\mathcal{F}'$ un module quasi-cohérent sur $X'$. Nous avons
\begin{align*}
\text{pr}_{2, *}(\text{pr}_1^*f_*\mathcal{F}'
\otimes_{\mathcal{O}_{X \times_R Y}} \mathcal{K})
& =
\text{pr}_{2, *}((f \times \text{id}_Y)_*
(\text{pr}'_1)^*\mathcal{F}'
\otimes_{\mathcal{O}_{X \times_R Y}} \mathcal{K}) \\
& =
\text{pr}_{2, *}(f \times \text{id}_Y)_*
\left((\text{pr}'_1)^*\mathcal{F}'
\otimes_{\mathcal{O}_{X' \times_R Y}}
(f \times \text{id}_Y)^*\mathcal{K})\right)  \\
& =
\text{pr}'_{2, *}((\text{pr}'_1)^*\mathcal{F}'
\otimes_{\mathcal{O}_{X' \times_R Y}} (f \times \text{id}_Y)^*\mathcal{K})
\end{align*}
Ici, la première égalité est le changement de base affine pour le carré
de gauche du diagramme ; voir
Cohomologie des schémas, lemme \ref{coherent-lemma-affine-base-change}.
La deuxième égalité résulte de la remarque \ref{functors-remark-affine-morphism}.
La troisième égalité est la fonctorialité des images directes de modules.
Cela démontre (1).

\medskip\noindent
Démonstration de (2). Considérons le diagramme commutatif
$$
\xymatrix{
X \times_R Y' \ar[rr]_-{\text{pr}'_2} \ar[rd]^{\text{id}_X \times g}
\ar[rdd]_{\text{pr}'_1} & & Y' \ar[d]^g \\
& X \times_R Y \ar[r]_-{\text{pr}_2} \ar[d]^{\text{pr}_1} & Y \\
& X
}
$$
Nous avons
\begin{align*}
g^*\text{pr}_{2, *}(\text{pr}_1^*\mathcal{F}
\otimes_{\mathcal{O}_{X \times_R Y}} \mathcal{K})
& =
\text{pr}'_{2, *}(
(\text{id}_X \times g)^*(
\text{pr}_1^*\mathcal{F} \otimes_{\mathcal{O}_{X \times_R Y}} \mathcal{K})) \\
& =
\text{pr}'_{2, *}((\text{pr}'_1)^*\mathcal{F}
\otimes_{\mathcal{O}_{X \times_R Y'}}
(\text{id}_X \times g)^*\mathcal{K})
\end{align*}
La première égalité résulte du changement de base plat pour le carré du diagramme ; voir
Cohomologie des schémas, lemme \ref{coherent-lemma-flat-base-change-cohomology}.
La deuxième égalité résulte de la fonctorialité de l'image inverse et du fait que
l'image inverse d'un produit tensoriel est le produit tensoriel des images inverses.

\medskip\noindent
La partie (3) est un cas particulier de (2).
\end{proof}

\begin{lemma}
\label{functors-lemma-functor-quasi-coherent-from-affine-diagonal-pre}
Soit $R$ un anneau. Soient $X$ et $Y$ des schémas sur $R$. Supposons que $X$
soit quasi-compact à diagonale affine. Soit
$F : \QCoh(\mathcal{O}_X) \to \QCoh(\mathcal{O}_Y)$
un foncteur $R$-linéaire exact à droite qui commute
aux sommes directes quelconques. Nous pouvons alors construire
\begin{enumerate}
\item un module quasi-cohérent $\mathcal{K}$ sur $X \times_R Y$, et
\item une transformation naturelle $t : F \to F_\mathcal{K}$,
où $F_\mathcal{K}$ désigne le foncteur (\ref{functors-equation-FM-QCoh}),
\end{enumerate}
telle que $t : F \circ f_* \to F_\mathcal{K} \circ f_*$ soit un isomorphisme
pour tout morphisme $f : X' \to X$ dont le schéma source est affine.
\end{lemma}

\begin{proof}
Considérons un morphisme $f' : X' \to X$ avec $X'$ affine. Comme la
diagonale de $X$ est affine, nous voyons que $f'$ est un morphisme affine
(Morphismes, lemme \ref{morphisms-lemma-affine-permanence}).
Ainsi, $f'_* : \QCoh(\mathcal{O}_{X'}) \to \QCoh(\mathcal{O}_X)$
est un foncteur exact $R$-linéaire
(Cohomologie des schémas, lemme \ref{coherent-lemma-relative-affine-vanishing})
qui commute aux sommes directes
(Cohomologie des schémas, lemme \ref{coherent-lemma-colimit-cohomology}).
Ainsi, $F \circ f'_*$ est un foncteur $R$-linéaire exact à droite qui commute
aux sommes directes quelconques. Par conséquent,
$F \circ f'_* = F_{\mathcal{K}'}$ pour un certain $\mathcal{K}'$
sur $X' \times_R Y$, d'après le lemme \ref{functors-lemma-functor-quasi-coherent-from-affine}.
De plus, étant donné un morphisme $f'' : X'' \to X'$ avec $X''$ affine,
nous obtenons une identification canonique
$(f'' \times \text{id}_Y)^*\mathcal{K}' = \mathcal{K}''$
à l'aide des références déjà données, combinées au
lemme \ref{functors-lemma-functor-quasi-coherent-from-affine-compose}.
Ces identifications satisfont une condition de cocycle pour
un autre morphisme $f''' : X''' \to X''$, condition dont nous laissons
au lecteur le soin d'écrire les détails.

\medskip\noindent
Choisissons un recouvrement ouvert affine $X = \bigcup_{i = 1, \ldots, n} U_i$.
Comme la diagonale de $X$ est affine, nous voyons que les intersections
$U_{i_0 \ldots i_p} = U_{i_0} \cap \ldots \cap U_{i_p}$ sont affines.
Comme ci-dessus, les morphismes d'inclusion
$j_{i_0 \ldots i_p} : U_{i_0 \ldots i_p} \to X$ sont affines.
Notons $\mathcal{K}_{i_0 \ldots i_p}$ le module quasi-cohérent
sur $U_{i_0 \ldots i_p} \times_R Y$ correspondant à
$F \circ j_{i_0 \ldots i_p *}$ comme ci-dessus.
D'après ce qui précède, nous obtenons des identifications
$$
\mathcal{K}_{i_0 \ldots i_p} =
\mathcal{K}_{i_0 \ldots \hat i_j \ldots i_p}|_{U_{i_0 \ldots i_p} \times_R Y}
$$
qui satisfont les compatibilités usuelles de recollement. Autrement dit, nous obtenons
un unique module quasi-cohérent $\mathcal{K}$ sur $X \times_R Y$
dont la restriction à $U_{i_0 \ldots i_p} \times_R Y$ est
$\mathcal{K}_{i_0 \ldots i_p}$, de manière compatible avec les identifications affichées.

\medskip\noindent
Construisons ensuite la transformation $t$. Étant donné un
$\mathcal{O}_X$-module quasi-cohérent $\mathcal{F}$, notons $\mathcal{F}_{i_0 \ldots i_p}$
la restriction de $\mathcal{F}$ à $U_{i_0 \ldots i_p}$ et notons
$(\text{pr}_1^*\mathcal{F} \otimes \mathcal{K})_{i_0 \ldots i_p}$
la restriction de $\text{pr}_1^*\mathcal{F} \otimes \mathcal{K}$ à
$U_{i_0 \ldots i_p} \times_R Y$.
Observons que
\begin{align*}
F(j_{i_0 \ldots i_p *}\mathcal{F}_{i_0 \ldots i_p})
& =
\text{pr}_{i_0 \ldots i_p, 2, *}(
\text{pr}_{i_0 \ldots i_p, 1}^*\mathcal{F}_{i_0 \ldots i_p}
\otimes \mathcal{K}_{i_0 \ldots i_p}) \\
& =
\text{pr}_{i_0 \ldots i_p, 2, *}
(\text{pr}_1^*\mathcal{F} \otimes \mathcal{K})_{i_0 \ldots i_p}
\end{align*}
où $\text{pr}_{i_0 \ldots i_p, 2} : U_{i_0 \ldots i_p} \times_R Y \to Y$
est la projection, et de même pour l'autre projection. De plus, ces
identifications sont compatibles avec les identifications affichées
au paragraphe précédent. Rappelons, d'après Cohomologie des schémas, lemme
\ref{coherent-lemma-separated-case-relative-cech},
que le complexe de {\v C}ech relatif
$$
\bigoplus 
\text{pr}_{i_0, 2, *}
(\text{pr}_1^*\mathcal{F} \otimes \mathcal{K})_{i_0}
\to
\bigoplus 
\text{pr}_{i_0i_1, 2, *}
(\text{pr}_1^*\mathcal{F} \otimes \mathcal{K})_{i_0i_1}
\to
\bigoplus 
\text{pr}_{i_0i_1i_2, 2, *}
(\text{pr}_1^*\mathcal{F} \otimes \mathcal{K})_{i_0i_1i_2}
\to \ldots
$$
calcule $R\text{pr}_{2, *}(\text{pr}_1^*\mathcal{F} \otimes \mathcal{K})$.
Ainsi, le faisceau de cohomologie en degré $0$ est $F_\mathcal{K}(\mathcal{F})$.
Nous obtenons donc le morphisme souhaité
$t : F(\mathcal{F}) \to F_\mathcal{K}(\mathcal{F})$
en considérant le diagramme commutatif suivant
$$
\xymatrix{
&
F(\mathcal{F}) \ar[r] \ar@{..>}[d] &
\bigoplus F(j_{i_0*}\mathcal{F}_{i_0}) \ar[r] \ar[d] &
\bigoplus F(j_{i_0i_1*}\mathcal{F}_{i_0i_1}) \ar[d] \\
0 \ar[r] &
F_\mathcal{K}(\mathcal{F}) \ar[r] &
\bigoplus 
\text{pr}_{i_0, 2, *}
(\text{pr}_1^*\mathcal{F} \otimes \mathcal{K})_{i_0}
\ar[r] &
\bigoplus 
\text{pr}_{i_0i_1, 2, *}
(\text{pr}_1^*\mathcal{F} \otimes \mathcal{K})_{i_0i_1}
}
$$
Nous obtenons la ligne supérieure en appliquant $F$ au complexe (exact)
$0 \to \mathcal{F} \to \bigoplus j_{i_0*}\mathcal{F}_{i_0} \to
\bigoplus j_{i_0i_1*}\mathcal{F}_{i_0i_1}$ (mais, comme
$F$ n'est pas exact, la ligne supérieure est seulement un complexe et n'est pas
nécessairement exacte).
Les flèches verticales pleines sont les identifications ci-dessus.
Cela définit bien la flèche pointillée comme souhaité.
La flèche est fonctorielle en $\mathcal{F}$ ; nous omettons les détails.

\medskip\noindent
Il nous reste à démontrer la dernière assertion. Soit $f : X' \to X$
comme dans l'énoncé du lemme, et soit $\mathcal{K}'$
le module quasi-cohérent sur $X' \times_R Y$ construit
au premier paragraphe de la démonstration. Si le morphisme
$f : X' \to X$ est à valeurs dans l'un des ouverts $U_i$, alors le
résultat découle du
lemme \ref{functors-lemma-functor-quasi-coherent-from-affine-compose},
car, dans ce cas, nous savons
que $\mathcal{K}_i = \mathcal{K}|_{U_i \times_R Y}$
a pour image inverse $\mathcal{K}$. En général, nous obtenons un
recouvrement ouvert affine $X' = \bigcup U'_i$ avec $U'_i = f^{-1}(U_i)$,
ainsi que des isomorphismes
$\mathcal{K}'|_{U'_i} = f_i^*\mathcal{K}_i$, où
$f_i : U'_i \to U_i$ est le morphisme induit.
Ces morphismes satisfont les conditions de compatibilité nécessaires
pour se recoller en un isomorphisme $\mathcal{K}' = f^*\mathcal{K}$,
et nous concluons. Nous omettons certains détails.
\end{proof}

\begin{lemma}
\label{functors-lemma-coh-noetherian-from-affine-flat}
Dans le lemme \ref{functors-lemma-functor-quasi-coherent-from-affine}
ou dans le lemme \ref{functors-lemma-functor-quasi-coherent-from-affine-diagonal-pre},
si $F$ est un foncteur exact, alors l'objet correspondant
$\mathcal{K}$ de $\QCoh(\mathcal{O}_{X \times_R Y})$ est plat sur $X$.
\end{lemma}

\begin{proof}
Nous pouvons supposer que $X$ est affine ; nous sommes donc dans le cas du
lemme \ref{functors-lemma-functor-quasi-coherent-from-affine}.
D'après le lemme \ref{functors-lemma-functor-quasi-coherent-from-affine-compose},
nous pouvons supposer que $Y$ est affine. Dans le cas affine, l'énoncé
se traduit par la remarque \ref{functors-remark-exact-flat}.
\end{proof}

\begin{lemma}
\label{functors-lemma-functor-quasi-coherent-from-affine-diagonal}
Soit $R$ un anneau. Soient $X$ et $Y$ des schémas sur $R$. Supposons que $X$ soit
quasi-compact à diagonale affine.
Il existe une équivalence de catégories entre
\begin{enumerate}
\item la catégorie des foncteurs exacts $R$-linéaires
$F : \QCoh(\mathcal{O}_X) \to \QCoh(\mathcal{O}_Y)$
qui commutent aux sommes directes quelconques, et
\item la sous-catégorie pleine de $\QCoh(\mathcal{O}_{X \times_R Y})$ constituée
des $\mathcal{K}$ tels que
\begin{enumerate}
\item $\mathcal{K}$ est plat sur $X$ ;
\item pour $\mathcal{F} \in \QCoh(\mathcal{O}_X)$, nous avons
$R^q\text{pr}_{2, *}(\text{pr}_1^*\mathcal{F}
\otimes_{\mathcal{O}_{X \times_R Y}} \mathcal{K}) = 0$ pour $q > 0$.
\end{enumerate}
\end{enumerate}
donnée en associant à $\mathcal{K}$ le foncteur $F$ de (\ref{functors-equation-FM-QCoh}).
\end{lemma}

\begin{proof}
Soit $\mathcal{K}$ comme en (2). Le foncteur $F$ de
(\ref{functors-equation-FM-QCoh}) commute aux sommes directes.
Puisque, d'après (1) (a), le module $\mathcal{K}$ est plat sur $X$,
nous voyons que, pour toute suite exacte courte
$0 \to \mathcal{F}_1 \to \mathcal{F}_2 \to \mathcal{F}_3 \to 0$,
nous obtenons une suite exacte courte
$$
0 \to
\text{pr}_1^*\mathcal{F}_1 \otimes_{\mathcal{O}_{X \times_R Y}} \mathcal{K} \to
\text{pr}_1^*\mathcal{F}_2 \otimes_{\mathcal{O}_{X \times_R Y}} \mathcal{K} \to
\text{pr}_1^*\mathcal{F}_3 \otimes_{\mathcal{O}_{X \times_R Y}} \mathcal{K} \to
0
$$
Puisque, d'après (2)(b), l'image directe supérieure $R^1\text{pr}_{2, *}$
du premier terme est nulle, nous concluons que
$0 \to F(\mathcal{F}_1) \to F(\mathcal{F}_2) \to F(\mathcal{F}_3) \to 0$
est exacte, et nous voyons que $F$ est comme en (1).

\medskip\noindent
Soit $F$ comme en (1). Soient $\mathcal{K}$ et $t : F \to F_\mathcal{K}$ comme
dans le lemme \ref{functors-lemma-functor-quasi-coherent-from-affine-diagonal-pre}.
D'après le lemme \ref{functors-lemma-coh-noetherian-from-affine-flat}, nous voyons
que $\mathcal{K}$ est plat sur $X$. Pour achever la démonstration, nous devons
montrer que $t$ est un isomorphisme et établir l'assertion sur les images
directes supérieures. Ces deux faits résultent de ce que le
complexe de {\v C}ech relatif
$$
\bigoplus 
\text{pr}_{i_0, 2, *}
(\text{pr}_1^*\mathcal{F} \otimes \mathcal{K})_{i_0}
\to
\bigoplus 
\text{pr}_{i_0i_1, 2, *}
(\text{pr}_1^*\mathcal{F} \otimes \mathcal{K})_{i_0i_1}
\to
\bigoplus 
\text{pr}_{i_0i_1i_2, 2, *}
(\text{pr}_1^*\mathcal{F} \otimes \mathcal{K})_{i_0i_1i_2}
\to \ldots
$$
calcule $R\text{pr}_{2, *}(\text{pr}_1^*\mathcal{F} \otimes \mathcal{K})$.
Voir la démonstration du
lemme \ref{functors-lemma-functor-quasi-coherent-from-affine-diagonal-pre}
pour les notations et la raison de cette assertion. Dans la démonstration du
lemme \ref{functors-lemma-functor-quasi-coherent-from-affine-diagonal-pre},
nous avons aussi constaté que ce complexe est égal à l'image par $F$ du complexe
$$
\bigoplus j_{i_0*}\mathcal{F}_{i_0} \to
\bigoplus j_{i_0i_1*}\mathcal{F}_{i_0i_1} \to
\bigoplus j_{i_0i_1i_2*}\mathcal{F}_{i_0i_1i_2} \to \ldots
$$
Ce complexe est exact sauf en degré zéro, où son faisceau de cohomologie
est égal à $\mathcal{F}$. Ainsi, puisque $F$ est un foncteur exact,
nous concluons que $F = F_\mathcal{K}$ et que (2)(b) est satisfaite.

\medskip\noindent
Nous omettons de démontrer que la construction qui associe à $F$
le module $\mathcal{K}$ est fonctorielle et est un quasi-inverse du
foncteur qui associe à $\mathcal{K}$ le foncteur $F_\mathcal{K}$
déterminé par (\ref{functors-equation-FM-QCoh}).
\end{proof}

\begin{remark}
\label{functors-remark-characterize-FM-QCoh}
Soit $R$ un anneau. Soient $X$ et $Y$ des schémas sur $R$. Supposons que $X$
soit quasi-compact à diagonale affine.
Le lemme \ref{functors-lemma-functor-quasi-coherent-from-affine-diagonal} peut
se généraliser comme suit : les foncteurs
(\ref{functors-equation-FM-QCoh}) associés aux modules quasi-cohérents sur
$X \times_R Y$ sont exactement ceux
$F : \QCoh(\mathcal{O}_X) \to \QCoh(\mathcal{O}_Y)$
qui possèdent les propriétés suivantes :
\begin{enumerate}
\item $F$ est $R$-linéaire et commute aux sommes directes quelconques ;
\item $F \circ j_*$ est exact à droite lorsque $j : U \to X$ est
l'inclusion d'un ouvert affine ;
\item $0 \to F(\mathcal{F}) \to F(\mathcal{G}) \to F(\mathcal{H})$
est exacte chaque fois que $0 \to \mathcal{F} \to \mathcal{G} \to \mathcal{H} \to 0$
est une suite exacte telle que, pour tout $x \in X$, la suite des germes
$0 \to \mathcal{F}_x \to \mathcal{G}_x \to \mathcal{H}_x \to 0$
soit une suite exacte courte scindée.
\end{enumerate}
En effet, ces hypothèses suffisent pour construire une transformation
$t : F \to F_\mathcal{K}$ comme dans le
lemme \ref{functors-lemma-functor-quasi-coherent-from-affine-diagonal-pre}
et montrer qu'elle est un isomorphisme. De plus, les propriétés (1), (2) et (3)
sont bien satisfaites par les foncteurs (\ref{functors-equation-FM-QCoh}).
Si nous en avons un jour besoin, nous l'énoncerons et le démontrerons ici avec soin.
\end{remark}

\begin{lemma}
\label{functors-lemma-compose-FM-QCoh}
Soit $R$ un anneau. Soient $X$, $Y$, $Z$ des schémas sur $R$. Supposons que
$X$ et $Y$ soient quasi-compacts et à diagonale affine. Soient
$$
F : \QCoh(\mathcal{O}_X) \to \QCoh(\mathcal{O}_Y)
\quad\text{et}\quad
G : \QCoh(\mathcal{O}_Y) \to \QCoh(\mathcal{O}_Z)
$$
des foncteurs exacts $R$-linéaires qui commutent aux sommes directes quelconques.
Soient $\mathcal{K}$ dans $\QCoh(\mathcal{O}_{X \times_R Y})$
et $\mathcal{L}$ dans $\QCoh(\mathcal{O}_{Y \times_R Z})$
les « noyaux » correspondants ; voir le
lemme \ref{functors-lemma-functor-quasi-coherent-from-affine-diagonal}.
Alors $G \circ F$ correspond à
$\text{pr}_{13, *}(\text{pr}_{12}^*\mathcal{K}
\otimes_{\mathcal{O}_{X \times_R Y \times_R Z}}
\text{pr}_{23}^*\mathcal{L})$ dans $\QCoh(\mathcal{O}_{X \times_R Z})$.
\end{lemma}

\begin{proof}
Puisque $G \circ F : \QCoh(\mathcal{O}_X) \to \QCoh(\mathcal{O}_Z)$
est $R$-linéaire, exact et commute aux sommes directes quelconques,
le lemme \ref{functors-lemma-functor-quasi-coherent-from-affine-diagonal} montre
qu'il existe un objet $\mathcal{M}$ de
$\QCoh(\mathcal{O}_{X \times_R Z})$ correspondant à $G \circ F$.
D'autre part, notons
$\mathcal{E} = \text{pr}_{13, *}(\text{pr}_{12}^*\mathcal{K}
\otimes \text{pr}_{23}^*\mathcal{L})$. Ici et dans la suite de
la démonstration, nous omettons l'indice des produits tensoriels.
Soient $U \subset X$ et $W \subset Z$ des sous-schémas ouverts affines.
Pour démontrer le lemme, nous allons construire un isomorphisme
$$
\Gamma(U \times_R W, \mathcal{E})
\cong
\Gamma(U \times_R W, \mathcal{M})
$$
compatible avec les applications de restriction lorsque $U$ et $W$ varient.

\medskip\noindent
Tout d'abord, nous observons que
$$
\Gamma(U \times_R W, \mathcal{E}) =
\Gamma(U \times_R Y \times_R W,
\text{pr}_{12}^*\mathcal{K} \otimes \text{pr}_{23}^*\mathcal{L})
$$
par construction. Il nous faut donc montrer qu'il en va de même
pour $\mathcal{M}$.

\medskip\noindent
Écrivons $U = \Spec(A)$ et notons $j : U \to X$ le morphisme d'inclusion.
Rappelons, d'après la construction de $\mathcal{M}$ dans la démonstration du
lemme \ref{functors-lemma-functor-quasi-coherent-from-affine}, que
$$
\Gamma(U \times_R W, \mathcal{M}) =
\Gamma(W, G(F(j_*\mathcal{O}_U)))
$$
où la structure de $A$-module du membre de droite est induite par
l'action de $A$ sur $\mathcal{O}_U$. La correspondance entre
$F$ et $\mathcal{K}$ nous dit que
$F(j_*\mathcal{O}_U) = b_*(a^*j_*\mathcal{O}_U \otimes \mathcal{K})$,
où $a : X \times_R Y \to X$ et $b : X \times_R Y \to Y$ sont
les morphismes de projection. Comme $j$ est un morphisme affine, nous avons
$a^*j_*\mathcal{O}_U = (j \times \text{id}_Y)_*\mathcal{O}_{U \times_R Y}$
d'après Cohomologie des schémas, lemme
\ref{coherent-lemma-affine-base-change}.
Ensuite, nous avons
$(j \times \text{id}_Y)_*\mathcal{O}_{U \times_R Y} \otimes \mathcal{K} =
(j \times \text{id}_Y)_*\mathcal{K}|_{U \times_R Y}$,
par exemple d'après la remarque \ref{functors-remark-affine-morphism}.
En rassemblant ces égalités, nous obtenons
$$
F(j_*\mathcal{O}_U) =
(U \times_R Y \to Y)_*\mathcal{K}|_{U \times_R Y}
$$
muni de l'action évidente de $A$. (Cette formule est implicite dans la
démonstration du lemme \ref{functors-lemma-functor-quasi-coherent-from-affine}.)
En appliquant le foncteur $G$, nous obtenons
$$
G(F(j_*\mathcal{O}_U)) =
t_*(s^*((U \times_R Y \to Y)_*\mathcal{K}|_{U \times_R Y})
\otimes \mathcal{L})
$$
où $s : Y \times_R Z \to Y$ et $t : Y \times_R Z \to Z$ sont
les morphismes de projection. En utilisant de nouveau le changement de base
pour un morphisme affine (Cohomologie des schémas, lemme
\ref{coherent-lemma-affine-base-change}), cette fois pour le carré
$$
\xymatrix{
U \times_R Y \times_R Z \ar[r] \ar[d] & U \times_R Y \ar[d] \\
Y \times_R Z \ar[r] & Y
}
$$
nous obtenons
$$
s^*((U \times_R Y \to Y)_*\mathcal{K}|_{U \times_R Y}) =
(U \times_R Y \times_R Z \to Y \times_R Z)_*
\text{pr}_{12}^*\mathcal{K}|_{U \times_R Y \times_R Z}
$$
En utilisant encore la remarque \ref{functors-remark-affine-morphism}, nous trouvons
\begin{align*}
(U \times_R Y \times_R Z \to Y \times_R Z)_*
\text{pr}_{12}^*\mathcal{K}|_{U \times_R Y \times_R Z}
\otimes \mathcal{L} \\
=
(U \times_R Y \times_R Z \to Y \times_R Z)_*
\left(\text{pr}_{12}^*\mathcal{K} \otimes
\text{pr}_{23}^*\mathcal{L}\right)|_{U \times_R Y \times_R Z}
\end{align*}
En appliquant à cette égalité le foncteur
$\Gamma(W, t_*(-)) = \Gamma(Y \times_R W, -)$, nous obtenons
\begin{align*}
\Gamma(U \times_R W, \mathcal{M})
& =
\Gamma(W, G(F(j_*\mathcal{O}_U))) \\
& =
\Gamma(Y \times_R W, (U \times_R Y \times_R Z \to Y \times_R Z)_*
(\text{pr}_{12}^*\mathcal{K} \otimes
\text{pr}_{23}^*\mathcal{L})|_{U \times_R Y \times_R Z}) \\
& =
\Gamma(U \times_R Y \times_R W,
\text{pr}_{12}^*\mathcal{K} \otimes \text{pr}_{23}^*\mathcal{L})
\end{align*}
comme voulu. Nous omettons la vérification que ces isomorphismes sont
compatibles avec les applications de restriction.
\end{proof}

\begin{lemma}
\label{functors-lemma-persistence-exactness}
Soient $R$, $X$, $Y$ et $\mathcal{K}$ comme au point (2) du
lemme \ref{functors-lemma-functor-quasi-coherent-from-affine-diagonal}.
Alors, pour tout schéma $T$ sur $R$, nous avons
$$
R^q\text{pr}_{13, *}(\text{pr}_{12}^*\mathcal{F}
\otimes_{\mathcal{O}_{T \times_R X \times_R Y}}
\text{pr}_{23}^*\mathcal{K}) = 0
$$
pour tout module quasi-cohérent $\mathcal{F}$ sur $T \times_R X$ et tout
$q > 0$.
\end{lemma}

\begin{proof}
La question est locale sur $T$ ; nous pouvons donc supposer $T$ affine.
Dans ce cas, nous pouvons considérer le diagramme
$$
\xymatrix{
T \times_R X \ar[d] &
T \times_R X \times_R Y \ar[d] \ar[l] \ar[r] &
T \times_R Y \ar[d] \\
X &
X \times_R Y \ar[l] \ar[r] &
Y
}
$$
dont les flèches verticales sont affines. En particulier, le foncteur image
directe par $T \times_R Y \to Y$ est fidèle et exact (Cohomologie des schémas,
lemme \ref{coherent-lemma-relative-affine-vanishing} et
Morphismes, lemme \ref{morphisms-lemma-affine-equivalence-modules}).
En parcourant le diagramme et en utilisant l'annulation des images directes
supérieures par les morphismes affines (voir la référence ci-dessus), nous
voyons qu'il suffit de démontrer que
$$
R^q\text{pr}_{2, *}(
\text{pr}_{23, *}(\text{pr}_{12}^*\mathcal{F}
\otimes_{\mathcal{O}_{T \times_R X \times_R Y}}
\text{pr}_{23}^*\mathcal{K})) =
R^q\text{pr}_{2, *}(
\text{pr}_{23, *}(\text{pr}_{12}^*\mathcal{F})
\otimes_{\mathcal{O}_{X \times_R Y}}
\mathcal{K}))
$$
est nul, ce qui résulte de l'hypothèse sur $\mathcal{K}$.
L'égalité découle de la remarque \ref{functors-remark-affine-morphism}.
\end{proof}

\begin{lemma}
\label{functors-lemma-functor-quasi-coherent-from-separated}
Dans le lemme \ref{functors-lemma-functor-quasi-coherent-from-affine-diagonal},
supposons que $F$ et $\mathcal{K}$ se correspondent. Si $X$ est séparé et
plat sur $R$, alors il existe une surjection
$\mathcal{O}_X \boxtimes F(\mathcal{O}_X) \to \mathcal{K}$.
\end{lemma}

\begin{proof}
Soit $\Delta : X \to X \times_R X$ le morphisme diagonal et
posons $\mathcal{O}_\Delta = \Delta_*\mathcal{O}_X$.
Comme $\Delta$ est une immersion fermée, nous avons une suite exacte courte
$$
0 \to \mathcal{I} \to 
\mathcal{O}_{X \times_R X} \to \mathcal{O}_\Delta \to 0
$$
Comme $\mathcal{K}$ est plat sur $X$, son image inverse
$\text{pr}_{23}^*\mathcal{K}$ sur $X \times_R X \times_R Y$
est plate sur $X \times_R X$. Nous obtenons une suite exacte courte
$$
0 \to 
\text{pr}_{12}^*\mathcal{I}
\otimes
\text{pr}_{23}^*\mathcal{K} \to
\text{pr}_{23}^*\mathcal{K} \to
\text{pr}_{12}^*\mathcal{O}_\Delta
\otimes
\text{pr}_{23}^*\mathcal{K} \to 0
$$
sur $X \times_R X \times_R Y$ ; voir
Modules, lemme \ref{modules-lemma-pullback-tensor-flat-module}.
Ainsi, d'après le lemme \ref{functors-lemma-persistence-exactness},
nous obtenons une surjection
$$
\text{pr}_{13, *}(\text{pr}_{23}^*\mathcal{K})
\to
\text{pr}_{13, *}(
\text{pr}_{12}^*\mathcal{O}_\Delta
\otimes
\text{pr}_{23}^*\mathcal{K})
$$
D'après le changement de base plat (Cohomologie des schémas, lemme
\ref{coherent-lemma-flat-base-change-cohomology}), la source de cette flèche
est égale à $\text{pr}_2^*\text{pr}_{2, *}\mathcal{K}
= \mathcal{O}_X \boxtimes F(\mathcal{O}_X)$. D'autre part, le but est
égal à
$$
\text{pr}_{13, *}(
\text{pr}_{12}^*\mathcal{O}_\Delta
\otimes
\text{pr}_{23}^*\mathcal{K}) =
\text{pr}_{13, *} (\Delta \times \text{id}_Y)_* \mathcal{K} =
\mathcal{K}
$$
ce qui achève la démonstration. La première égalité résulte par exemple de
Cohomologie, lemme \ref{cohomology-lemma-projection-formula-closed-immersion}
et du fait que $\text{pr}_{12}^*\mathcal{O}_\Delta =
(\Delta \times \text{id}_Y)_*\mathcal{O}_{X \times_R Y}$.
\end{proof}









\section{Reconstruction de Gabriel--Rosenberg}
\label{functors-section-gabriel}

\noindent
Le titre de cette section renvoie à des résultats tels que la
proposition \ref{functors-proposition-gabriel-rosenberg}.
Outre l'article original de Gabriel \cite{Gabriel}, on pourra consulter
\cite{Brandenburg}, qui démontre le résultat pour les schémas quasi-séparés
et examine la littérature. Dans cette section, nous démontrerons seulement
la reconstruction de Gabriel--Rosenberg pour les schémas quasi-compacts et
quasi-séparés.

\begin{lemma}
\label{functors-lemma-categorically-compact-QCoh}
Soit $X$ un schéma quasi-compact et quasi-séparé.
Soit $\mathcal{F}$ un $\mathcal{O}_X$-module quasi-cohérent.
Alors $\mathcal{F}$ est un objet compact au sens catégorique de
$\QCoh(\mathcal{O}_X)$ si et seulement si $\mathcal{F}$ est de
présentation finie.
\end{lemma}

\begin{proof}
Voir Catégories, définition \ref{categories-definition-compact-object},
pour notre notion d'objet compact au sens catégorique dans une catégorie.
Si $\mathcal{F}$ est de présentation finie, alors elle est compacte au sens
catégorique d'après Modules, lemme
\ref{modules-lemma-finite-presentation-quasi-compact-colimit}.
Réciproquement, tout module quasi-cohérent $\mathcal{F}$ peut s'écrire
comme une limite inductive filtrante $\mathcal{F} = \colim \mathcal{F}_i$
de $\mathcal{O}_X$-modules de présentation finie (donc quasi-cohérents) ;
voir Propriétés, lemme
\ref{properties-lemma-directed-colimit-finite-presentation}.
Si $\mathcal{F}$ est compact au sens catégorique, il existe un indice $i$
et un morphisme $\mathcal{F} \to \mathcal{F}_i$ qui est inverse à droite du
morphisme donné $\mathcal{F}_i \to \mathcal{F}$.
Nous concluons que $\mathcal{F}$ est facteur direct d'un module de
présentation finie, et est donc lui-même de présentation finie.
\end{proof}

\begin{lemma}
\label{functors-lemma-supported-on-support-pre}
Soit $X$ un schéma affine. Soit $\mathcal{F}$ un
$\mathcal{O}_X$-module de présentation finie. Soit $\mathcal{E}$ un
$\mathcal{O}_X$-module quasi-cohérent non nul. Si
$\text{Supp}(\mathcal{E}) \subset \text{Supp}(\mathcal{F})$,
alors il existe un morphisme non nul $\mathcal{F} \to \mathcal{E}$.
\end{lemma}

\begin{proof}
Traduisons l'énoncé en algèbre. Soit $A$ un anneau. Soit $M$ un
$A$-module de présentation finie. Soit $N$ un $A$-module non nul. Supposons
$\text{Supp}(N) \subset \text{Supp}(M)$. Il s'agit de montrer que
$\Hom_A(M, N)$ est non nul.
Nous pouvons supposer $N = A/I$ cyclique (en remplaçant $N$ par n'importe
lequel de ses sous-modules cycliques non nuls). Choisissons une présentation
$$
A^{\oplus m} \xrightarrow{T} A^{\oplus n} \to M \to 0
$$
Rappelons que $\text{Supp}(M)$ est défini par $\text{Fit}_0(M)$, qui est
l'idéal engendré par les mineurs $n \times n$ de la matrice $T$.
Voir Plus sur l'algèbre, lemme
\ref{more-algebra-lemma-fitting-ideal-basics}.
L'hypothèse $\text{Supp}(N) \subset \text{Supp}(M)$ signifie alors que
les éléments de $\text{Fit}_0(M)$ sont nilpotents dans $A/I$.
Considérons la suite exacte
$$
0 \to \Hom_A(M, A/I) \to (A/I)^{\oplus n} \xrightarrow{T^t} (A/I)^{\oplus m}
$$
Nous devons montrer que $T^t$ ne peut être injective ; nous invitons le
lecteur à en trouver une démonstration à partir de la nilpotence des éléments
de $\text{Fit}_0(M)$ dans $A/I$. Voici la nôtre.
Comme $\text{Fit}_0(M)$ est de type fini, cette nilpotence implique que
l'annulateur $J \subset A/I$ de $\text{Fit}_0(M)$ dans $A/I$ est non nul.
Pour montrer que $T^t$ n'est pas injective, nous pouvons localiser en un idéal
premier. En choisissant convenablement cet idéal premier, nous pouvons
supposer $A$ local et $J$ toujours non nul. Alors $T^t$ a un noyau non nul
d'après Plus sur l'algèbre, lemme
\ref{more-algebra-lemma-exact-length-1}.
\end{proof}

\begin{lemma}
\label{functors-lemma-supported-on-support}
Soit $X$ un schéma quasi-compact et quasi-séparé.
Soit $\mathcal{F}$ un $\mathcal{O}_X$-module de présentation finie.
Les deux sous-catégories suivantes de $\QCoh(\mathcal{O}_X)$ sont égales :
\begin{enumerate}
\item la sous-catégorie pleine $\mathcal{A} \subset \QCoh(\mathcal{O}_X)$
dont les objets sont les modules quasi-cohérents dont le support est contenu,
au sens ensembliste, dans $\text{Supp}(\mathcal{F})$ ;
\item la plus petite sous-catégorie de Serre
$\mathcal{B} \subset \QCoh(\mathcal{O}_X)$ qui contient
$\mathcal{F}$ et qui est stable par extensions et par sommes directes
quelconques.
\end{enumerate}
\end{lemma}

\begin{proof}
Observons que l'énoncé a un sens puisque les $\mathcal{O}_X$-modules de
présentation finie sont quasi-cohérents.
Comme $\mathcal{A}$ est une sous-catégorie de Serre stable par extensions et
par sommes directes et comme $\mathcal{F}$ est un objet de $\mathcal{A}$,
nous voyons que $\mathcal{B} \subset \mathcal{A}$. Il reste donc à montrer
que $\mathcal{A}$ est contenue dans $\mathcal{B}$.

\medskip\noindent
Soit $\mathcal{E}$ un objet de $\mathcal{A}$. Il existe un sous-module
maximal $\mathcal{E}' \subset \mathcal{E}$ appartenant à $\mathcal{B}$.
En effet, soit $\mathcal{E}_i \subset \mathcal{E}$, $i \in I$, la famille
des sous-objets qui sont des objets de $\mathcal{B}$. Alors
$\bigoplus \mathcal{E}_i$ appartient à $\mathcal{B}$, de même que
$$
\mathcal{E}' = \Im(\bigoplus \mathcal{E}_i \longrightarrow \mathcal{E})
$$
C'est manifestement le sous-module maximal recherché.

\medskip\noindent
Supposons maintenant qu'il existe un morphisme non nul
$\mathcal{G} \to \mathcal{E}/\mathcal{E}'$
avec $\mathcal{G}$ dans $\mathcal{B}$. Alors
$\mathcal{G}' = \mathcal{E} \times_{\mathcal{E}/\mathcal{E}'} \mathcal{G}$
appartient à $\mathcal{B}$ : $\mathcal{E}'$ en est le sous-objet et
$\mathcal{G}$ le quotient. L'image de $\mathcal{G}' \to \mathcal{E}$ serait alors
strictement plus grande que $\mathcal{E}'$, ce qui contredirait la maximalité
de $\mathcal{E}'$. Il suffit donc de démontrer l'assertion du paragraphe
suivant.

\medskip\noindent
Soit $\mathcal{E}$ un objet non nul de $\mathcal{A}$. Nous affirmons qu'il
existe un morphisme non nul $\mathcal{G} \to \mathcal{E}$ avec
$\mathcal{G}$ dans $\mathcal{B}$. Nous allons le démontrer par récurrence sur
le nombre minimal $n$ d'ouverts affines $U_i$ de $X$ tels que
$\text{Supp}(\mathcal{E}) \subset U_1 \cup \ldots \cup U_n$.
Posons $U = U_n$ et notons $j : U \to X$ le morphisme d'inclusion.
Notons $\mathcal{E}' = \Im(\mathcal{E} \to j_*\mathcal{E}|_U)$.
Le noyau $\mathcal{E}''$ de la surjection
$\mathcal{E} \to \mathcal{E}'$ a son support contenu dans
$U_1 \cup \ldots \cup U_{n - 1}$. Ainsi, si $\mathcal{E}''$ est non nul,
le résultat découle de l'hypothèse de récurrence. Autrement dit, nous pouvons
supposer $\mathcal{E} \subset j_*\mathcal{E}|_U$.
En particulier, $\mathcal{E}|_U$ est non nul.
D'après le lemme \ref{functors-lemma-supported-on-support-pre}, il existe une
morphisme non nul $\mathcal{F}|_U \to \mathcal{E}|_U$.
Il correspond à un morphisme
$$
\varphi : \mathcal{F} \longrightarrow j_*(\mathcal{E}|_U)
$$
dont la restriction à $U$ est non nulle.
En posant $\mathcal{G} = \varphi^{-1}(\mathcal{E})$, nous concluons.
\end{proof}

\begin{lemma}
\label{functors-lemma-quotient-supported-on-closed}
Soit $X$ un schéma quasi-compact et quasi-séparé.
Soit $Z \subset X$ une partie fermée telle que $U = X \setminus Z$
soit quasi-compact. Soit $\mathcal{A} \subset \QCoh(\mathcal{O}_X)$
la sous-catégorie pleine dont les objets sont les modules quasi-cohérents
à support dans $Z$. Alors le foncteur de restriction
$\QCoh(\mathcal{O}_X) \to \QCoh(\mathcal{O}_U)$ induit
une équivalence $\QCoh(\mathcal{O}_X)/\mathcal{A} \cong \QCoh(\mathcal{O}_U)$.
\end{lemma}

\begin{proof}
D'après la propriété universelle de la construction quotient
(Homologie, lemme \ref{homology-lemma-serre-subcategory-is-kernel}),
nous obtenons bien un foncteur induit
$\QCoh(\mathcal{O}_X)/\mathcal{A} \cong \QCoh(\mathcal{O}_U)$.
Notons $j : U \to X$ le morphisme d'inclusion. Comme $j$ est quasi-compact
et quasi-séparé, nous obtenons un foncteur
$j_* : \QCoh(\mathcal{O}_U) \to \QCoh(\mathcal{O}_X)$.
Le lecteur vérifiera qu'il définit un quasi-inverse ; nous omettons les détails.
\end{proof}

\begin{lemma}
\label{functors-lemma-characterize-affine}
Soit $X$ un schéma quasi-compact et quasi-séparé.
Si $\QCoh(\mathcal{O}_X)$ est équivalente à la catégorie des modules
sur un anneau, alors $X$ est affine.
\end{lemma}

\begin{proof}
Soit $F : \text{Mod}_R \to \QCoh(\mathcal{O}_X)$ une équivalence.
Alors $\mathcal{F} = F(R)$ possède les propriétés suivantes :
\begin{enumerate}
\item c'est un $\mathcal{O}_X$-module de présentation finie
(lemme \ref{functors-lemma-categorically-compact-QCoh}) ;
\item $\Hom_X(\mathcal{F}, -)$ est exact ;
\item $\Hom_X(\mathcal{F}, \mathcal{F})$ est un anneau commutatif ;
\item tout objet de $\QCoh(\mathcal{O}_X)$ est quotient d'une somme directe
de copies de $\mathcal{F}$.
\end{enumerate}
Soit $x \in X$ un point fermé. Considérons la surjection
$$
\mathcal{O}_X \to i_*\kappa(x)
$$
où le but est l'image directe de $\kappa(x)$ par le morphisme d'inclusion
$i : x \to X$. Nous avons
$$
\Hom_X(\mathcal{F}, i_*\kappa(x)) =
\Hom_{\mathcal{O}_{X, x}}(\mathcal{F}_x, \kappa(x))
$$
Tout d'abord, (4) implique que $\mathcal{F}_x$ est non nul.
D'après (2), tout homomorphisme $\mathcal{F}_x \to \kappa(x)$ se relève en
un homomorphisme $\mathcal{F}_x \to \mathcal{O}_{X, x}$ (puisqu'il se
relève même en un morphisme global $\mathcal{F} \to \mathcal{O}_X$).
Comme $\mathcal{F}_x$ est un $\mathcal{O}_{X, x}$-module fini, il s'ensuit
que $\mathcal{F}_x$ est un $\mathcal{O}_{X, x}$-module libre non nul de rang fini.
Comme $\mathcal{F}$ est de présentation finie, cela implique que
$\mathcal{F}$ est libre de rang fini strictement positif sur un voisinage ouvert de $x$
(Modules, lemme \ref{modules-lemma-finite-presentation-stalk-free}).
Puisque toute partie fermée de $X$ contient un point fermé
(Topologie, lemme \ref{topology-lemma-quasi-compact-closed-point}),
il s'ensuit que $\mathcal{F}$ est localement libre de rang fini strictement positif.
De même, l'homomorphisme
$$
\Hom_X(\mathcal{F}, \mathcal{F}) \to
\Hom_X(\mathcal{F}, i_*i^*\mathcal{F}) =
\Hom_{\kappa(x)}(\mathcal{F}_x/\mathfrak m_x \mathcal{F}_x,
\mathcal{F}_x/\mathfrak m_x \mathcal{F}_x)
$$
est surjective. D'après la propriété (3), nous concluons que le rang de
$\mathcal{F}_x$ est nécessairement $1$. Ainsi $\mathcal{F}$ est un
$\mathcal{O}_X$-module inversible. Nous en déduisons que le foncteur
$$
\mathcal{H}
\longmapsto
\Gamma(X, \mathcal{H}) = \Hom_X(\mathcal{O}_X, \mathcal{H}) =
\Hom_X(\mathcal{F}, \mathcal{H} \otimes_{\mathcal{O}_X} \mathcal{F})
$$
sur $\QCoh(\mathcal{O}_X)$ est lui aussi exact. Il s'ensuit que le premier
groupe $\Ext$
$$
\Ext^1_{\QCoh(\mathcal{O}_X)}(\mathcal{O}_X, \mathcal{H}) = 0
$$
calculé dans la catégorie abélienne $\QCoh(\mathcal{O}_X)$ s'annule pour
tout $\mathcal{H}$ dans $\QCoh(\mathcal{O}_X)$. Or, comme
$\QCoh(\mathcal{O}_X) \subset \textit{Mod}(\mathcal{O}_X)$
est stable par extensions (Schémas, section
\ref{schemes-section-quasi-coherent}), nous voyons que le groupe $\Ext^1$
entre modules quasi-cohérents calculé dans $\QCoh(\mathcal{O}_X)$ est le même
que celui calculé dans $\textit{Mod}(\mathcal{O}_X)$. Nous concluons donc que
$$
H^1(X, \mathcal{H}) =
\Ext^1_{\textit{Mod}(\mathcal{O}_X)}(\mathcal{O}_X, \mathcal{H}) = 0
$$
pour tout $\mathcal{H}$ dans $\QCoh(\mathcal{O}_X)$.
Il en résulte que $X$ est affine, par exemple d'après
Cohomologie des schémas, lemme
\ref{coherent-lemma-quasi-compact-h1-zero-covering}.
\end{proof}

\begin{proposition}
\label{functors-proposition-gabriel-rosenberg}
\begin{reference}
Cas particulier de \cite[Theorem 1.2]{Brandenburg}
\end{reference}
Soient $X$ et $Y$ des schémas quasi-compacts et quasi-séparés.
Si $F : \QCoh(\mathcal{O}_X) \to \QCoh(\mathcal{O}_Y)$
est une équivalence, alors il existe un isomorphisme de schémas
$f : Y \to X$ et un $\mathcal{O}_Y$-module inversible
$\mathcal{L}$ tels que $F(\mathcal{F}) = f^*\mathcal{F} \otimes \mathcal{L}$.
\end{proposition}

\begin{proof}
Bien entendu, $F$ est additif, exact, commute à toutes les limites,
à toutes les limites inductives, aux sommes directes, etc.
Soit $U \subset X$ un sous-schéma ouvert affine.
Soit $\mathcal{I} \subset \mathcal{O}_X$ un faisceau quasi-cohérent
d'idéaux de type fini tel que $Z = V(\mathcal{I})$ soit le complémentaire
de $U$ dans $X$ ; voir Propriétés, lemme
\ref{properties-lemma-quasi-coherent-finite-type-ideals}.
Alors $\mathcal{O}_X/\mathcal{I}$ est un $\mathcal{O}_X$-module de
présentation finie. Ainsi $\mathcal{G} = F(\mathcal{O}_X/\mathcal{I})$
est un $\mathcal{O}_Y$-module de présentation finie d'après le
lemme \ref{functors-lemma-categorically-compact-QCoh}.
Notons $T \subset Y$ le support de $\mathcal{G}$ et posons
$V = Y \setminus T$. Comme $\mathcal{G}$ est de présentation finie,
le schéma $V$ est un ouvert quasi-compact de $Y$.
D'après le lemme \ref{functors-lemma-supported-on-support}, nous voyons que $F$ induit
une équivalence entre
\begin{enumerate}
\item la sous-catégorie pleine de $\QCoh(\mathcal{O}_X)$ formée des modules
à support dans $Z$, et
\item la sous-catégorie pleine de $\QCoh(\mathcal{O}_Y)$ formée des modules
à support dans $T$.
\end{enumerate}
D'après le lemme \ref{functors-lemma-quotient-supported-on-closed}, nous obtenons un
diagramme commutatif
$$
\xymatrix{
\QCoh(\mathcal{O}_X) \ar[r]_F \ar[d] &
\QCoh(\mathcal{O}_Y) \ar[d] \\
\QCoh(\mathcal{O}_U) \ar[r]^{F_U} &
\QCoh(\mathcal{O}_V)
}
$$
où les flèches verticales sont les foncteurs de restriction et les flèches
horizontales des équivalences. D'après le lemme
\ref{functors-lemma-characterize-affine}, nous concluons que $V$ est affine.
Le cas affine est traité par le lemme \ref{functors-lemma-functor-equivalence}.
Nous trouvons donc un isomorphisme $f_U : V \to U$ et un
$\mathcal{O}_V$-module inversible $\mathcal{L}_U$ tels que
$F_U$ soit le foncteur
$\mathcal{F} \mapsto f_U^*\mathcal{F} \otimes \mathcal{L}_U$.

\medskip\noindent
Pour achever la démonstration, il suffit d'observer que les diagrammes
ci-dessus satisfont une compatibilité évidente à l'égard des inclusions de
sous-schémas ouverts affines de $X$. Les morphismes $f_U$ et les modules
inversibles $\mathcal{L}_U$ se recollent donc. Nous omettons les détails.
\end{proof}









\section{Foncteurs entre catégories de modules cohérents}
\label{functors-section-functor-coherent}

\noindent
Le lemme suivant garantit que nous pouvons utiliser les résultats sur les
foncteurs entre catégories de modules quasi-cohérents lorsque nous disposons
d'un foncteur entre catégories de modules cohérents.

\begin{lemma}
\label{functors-lemma-functor-coherent}
Soient $X$ et $Y$ des schémas noethériens. Soit
$F : \textit{Coh}(\mathcal{O}_X) \to \textit{Coh}(\mathcal{O}_Y)$
un foncteur. Alors $F$ se prolonge de manière unique en un foncteur
$\QCoh(\mathcal{O}_X) \to \QCoh(\mathcal{O}_Y)$
qui commute aux limites inductives filtrantes.
Si $F$ est additif, son prolongement commute aux sommes directes quelconques.
Si $F$ est exact, exact à gauche ou exact à droite, il en va de même de son
prolongement.
\end{lemma}

\begin{proof}
L'existence et l'unicité du prolongement sont un fait général ; voir
Catégories, lemme \ref{categories-lemma-extend-functor-by-colim}.
Pour vérifier que ce lemme s'applique, observons que les modules cohérents
sont de présentation finie (Modules, lemme
\ref{modules-lemma-coherent-finite-presentation}) et sont donc des objets
compacts au sens catégorique de $\textit{Mod}(\mathcal{O}_X)$ d'après
Modules, lemme \ref{modules-lemma-finite-presentation-quasi-compact-colimit}.
Enfin, tout module quasi-cohérent est limite inductive filtrante de modules
cohérents, par exemple d'après Propriétés, lemme
\ref{properties-lemma-quasi-coherent-colimit-finite-type}.

\medskip\noindent
Supposons $F$ additif. Si
$\mathcal{F} = \bigoplus_{j \in J} \mathcal{H}_j$
avec les $\mathcal{H}_j$ quasi-cohérents, alors
$\mathcal{F} = \colim_{J' \subset J\text{ fini}}
\bigoplus_{j \in J'} \mathcal{H}_j$.
En notant encore $F$ le prolongement de $F$, nous obtenons
\begin{align*}
F(\mathcal{F})
& =
\colim_{J' \subset J\text{ fini}}
F(\bigoplus\nolimits_{j \in J'} \mathcal{H}_j) \\
& =
\colim_{J' \subset J\text{ fini}}
\bigoplus\nolimits_{j \in J'} F(\mathcal{H}_j) \\
& =
\bigoplus\nolimits_{j \in J} F(\mathcal{H}_j)
\end{align*}
Ainsi $F$ commute aux sommes directes quelconques.

\medskip\noindent
Supposons que
$0 \to \mathcal{F} \to \mathcal{F}' \to \mathcal{F}'' \to 0$
soit une suite exacte courte de $\mathcal{O}_X$-modules quasi-cohérents.
Écrivons $\mathcal{F}' = \bigcup \mathcal{F}'_i$ comme réunion de ses
sous-modules cohérents ; voir Propriétés, lemme
\ref{properties-lemma-quasi-coherent-colimit-finite-type}.
Notons $\mathcal{F}''_i \subset \mathcal{F}''$ l'image de $\mathcal{F}'_i$
et posons $\mathcal{F}_i = \mathcal{F} \cap \mathcal{F}'_i =
\Ker(\mathcal{F}'_i \to \mathcal{F}''_i)$. Il est alors clair que
$\mathcal{F} = \bigcup \mathcal{F}_i$ et
$\mathcal{F}'' = \bigcup \mathcal{F}''_i$, et que nous avons des suites
exactes courtes
$$
0 \to \mathcal{F}_i \to \mathcal{F}_i' \to \mathcal{F}_i'' \to 0
$$
Comme le prolongement commute aux limites inductives filtrantes, nous avons
$F(\mathcal{F}) = \colim_{i \in I} F(\mathcal{F}_i)$,
$F(\mathcal{F}') = \colim_{i \in I} F(\mathcal{F}'_i)$ et
$F(\mathcal{F}'') = \colim_{i \in I} F(\mathcal{F}''_i)$.
Puisque les limites inductives filtrantes sont exactes
(Modules, lemme \ref{modules-lemma-limits-colimits}), nous concluons que les
propriétés d'exactitude de $F$ se transmettent à son prolongement.
\end{proof}

\begin{lemma}
\label{functors-lemma-equivalence-coherent}
Soient $X$ et $Y$ des schémas noethériens. Soit
$F : \textit{Coh}(\mathcal{O}_X) \to \textit{Coh}(\mathcal{O}_Y)$
une équivalence de catégories. Alors il existe un isomorphisme $f : Y \to X$
et un $\mathcal{O}_Y$-module inversible $\mathcal{L}$ tels que
$F(\mathcal{F}) = f^*\mathcal{F} \otimes \mathcal{L}$.
\end{lemma}

\begin{proof}
D'après le lemme \ref{functors-lemma-functor-coherent}, nous obtenons un unique foncteur
$F' : \QCoh(\mathcal{O}_X) \to \QCoh(\mathcal{O}_Y)$ qui prolonge $F$.
Il en va de même pour un quasi-inverse de $F$ et, par unicité, nous concluons
que $F'$ est une équivalence. D'après la
proposition \ref{functors-proposition-gabriel-rosenberg}, il existe un isomorphisme
$f : Y \to X$ et un $\mathcal{O}_Y$-module inversible $\mathcal{L}$
tels que $F'(\mathcal{F}) = f^*\mathcal{F} \otimes \mathcal{L}$.
Alors $f$ et $\mathcal{L}$ conviennent également pour $F$.
\end{proof}

\begin{remark}
\label{functors-remark-equivalence-coherent-linear}
Dans le lemme \ref{functors-lemma-equivalence-coherent}, si $X$ et $Y$ sont définis
sur un même anneau de base $R$ et si $F$ est $R$-linéaire, alors
l'isomorphisme $f$ est un morphisme de schémas sur $R$.
\end{remark}

\begin{lemma}
\label{functors-lemma-characterize-finite}
Soit $f : V \to X$ un morphisme quasi-fini séparé de schémas noethériens.
S'il existe un $\mathcal{O}_V$-module cohérent $\mathcal{K}$ dont le support
est $V$, tel que $f_*\mathcal{K}$ soit cohérent et que
$R^qf_*\mathcal{K} = 0$, alors $f$ est fini.
\end{lemma}

\begin{proof}
D'après le théorème principal de Zariski, il existe une immersion ouverte
$j : V \to Y$ au-dessus de $X$ telle que $\pi : Y \to X$ soit fini ; voir
Plus sur les morphismes, lemme
\ref{more-morphisms-lemma-quasi-finite-separated-pass-through-finite}.
Comme $\pi$ est affine, le foncteur $\pi_*$ est exact et fidèle sur la
catégorie des $\mathcal{O}_X$-modules cohérents.
Nous voyons donc que $j_*\mathcal{K}$ est cohérent et que
$R^qj_*\mathcal{K}$ est nul pour $q > 0$.
Autrement dit, nous sommes ramenés au cas traité dans le paragraphe suivant.

\medskip\noindent
Supposons que $f$ soit une immersion ouverte. Nous pouvons remplacer $X$ par
l'adhérence schématique de $V$. Supposons $X \setminus V$ non vide afin
d'aboutir à une contradiction. Choisissons un point générique
$\xi \in X \setminus V$ d'une composante irréductible de $X \setminus V$.
En examinant la situation après le changement de base par
$\Spec(\mathcal{O}_{X, \xi}) \to X$, puis en utilisant le changement de base
plat et Cohomologie locale, lemme
\ref{local-cohomology-lemma-finiteness-pushforwards-and-H1-local},
nous sommes ramenés au problème algébrique du paragraphe suivant.

\medskip\noindent
Soit $(A, \mathfrak m)$ un anneau local noethérien. Soit $M$ un $A$-module
fini dont le support est $\Spec(A)$. Alors
$H^i_\mathfrak m(M) \not = 0$ pour un certain $i$. Cela résulte de
Complexes dualisants, lemme \ref{dualizing-lemma-depth}, et du fait que
$M$ est non nul, donc de profondeur finie.
\end{proof}

\noindent
Le lemme suivant se généralise au cas où $k$ est un anneau noethérien et où
$X$ est plat sur $k$ (toutes les autres hypothèses restant inchangées).

\begin{lemma}
\label{functors-lemma-functor-coherent-over-field}
Soit $k$ un corps. Soient $X$ et $Y$ des schémas de type fini sur $k$,
avec $X$ séparé. Il existe une équivalence de catégories entre
\begin{enumerate}
\item la catégorie des foncteurs exacts $k$-linéaires
$F : \textit{Coh}(\mathcal{O}_X) \to \textit{Coh}(\mathcal{O}_Y)$, et
\item la catégorie des $\mathcal{O}_{X \times Y}$-modules cohérents
$\mathcal{K}$ qui sont plats sur $X$ et dont le support est fini sur $Y$,
\end{enumerate}
obtenue en associant à $\mathcal{K}$ la restriction du foncteur
(\ref{functors-equation-FM-QCoh}) à $\textit{Coh}(\mathcal{O}_X)$.
\end{lemma}

\begin{proof}
Soit $\mathcal{K}$ comme en (2). D'après le
lemme \ref{functors-lemma-functor-quasi-coherent-from-affine-diagonal}, le foncteur
$F$ donné par (\ref{functors-equation-FM-QCoh}) est exact et $k$-linéaire.
De plus, $F$ envoie $\textit{Coh}(\mathcal{O}_X)$ dans
$\textit{Coh}(\mathcal{O}_Y)$, par exemple d'après
Cohomologie des schémas, lemme
\ref{coherent-lemma-support-proper-over-base-pushforward}.

\medskip\noindent
Construisons le quasi-inverse de cette construction. Soit $F$ comme en (1).
D'après le lemme \ref{functors-lemma-functor-coherent}, nous pouvons prolonger $F$ en
un foncteur exact $k$-linéaire sur les catégories de modules quasi-cohérents,
qui commute aux sommes directes quelconques.
D'après le lemme \ref{functors-lemma-functor-quasi-coherent-from-affine-diagonal}, ce
prolongement correspond à un unique module quasi-cohérent $\mathcal{K}$, plat
sur $X$, tel que
$R^q\text{pr}_{2, *}(\text{pr}_1^*\mathcal{F}
\otimes_{\mathcal{O}_{X \times Y}} \mathcal{K}) = 0$ pour $q > 0$
et pour tout $\mathcal{O}_X$-module quasi-cohérent $\mathcal{F}$.
Comme $F(\mathcal{O}_X)$ est un $\mathcal{O}_Y$-module cohérent, il résulte
du lemme \ref{functors-lemma-functor-quasi-coherent-from-separated} que
$\mathcal{K}$ est cohérent.

\medskip\noindent
Pour tout point fermé $x \in X$, notons $\mathcal{O}_x$ le faisceau gratte-ciel
en $x$ de valeur le corps résiduel de $x$. Nous avons
$$
F(\mathcal{O}_x) =
\text{pr}_{2, *}(\text{pr}_1^*\mathcal{O}_x \otimes \mathcal{K}) =
(x \times Y \to Y)_*(\mathcal{K}|_{x \times Y})
$$
Comme $x \times Y \to Y$ est fini, l'image directe par ce morphisme est
fidèle. Ainsi, si $y \in Y$ appartient à l'image du support de
$\mathcal{K}|_{x \times Y}$, alors $y$ appartient au support de
$F(\mathcal{O}_x)$.

\medskip\noindent
Soit $Z \subset X \times Y$ le support schématique $Z$ de $\mathcal{K}$ ;
voir Morphismes, définition
\ref{morphisms-definition-scheme-theoretic-support}.
Nous montrons d'abord que $Z \to Y$ est quasi-fini en montrant que ses fibres
au-dessus des points fermés sont finies. En effet, si la fibre de $Z \to Y$
au-dessus d'un point fermé $y \in Y$ est de dimension $> 0$, nous pouvons
trouver une infinité de points fermés deux à deux distincts
$x_1, x_2, \ldots$ dans l'image de $Z_y \to X$. Comme nous avons une
surjection
$\mathcal{O}_X \to \bigoplus_{i = 1, \ldots, n} \mathcal{O}_{x_i}$,
nous obtenons une surjection
$$
F(\mathcal{O}_X) \to \bigoplus\nolimits_{i = 1, \ldots, n} F(\mathcal{O}_{x_i})
$$
D'après ce qui précède, le point $y$ appartient au support de chacun des
modules cohérents $F(\mathcal{O}_{x_i})$. Comme $F(\mathcal{O}_X)$ est un
module cohérent, cela conduit à une contradiction, car le germe de
$F(\mathcal{O}_X)$ en $y$ sera engendré par $< n$ éléments si
$n$ est assez grand.
Ainsi $Z \to Y$ est quasi-fini.
Comme $\text{pr}_{2, *}\mathcal{K}$ est cohérent et que
$R^q\text{pr}_{2, *}\mathcal{K} = 0$ pour $q > 0$, nous concluons que
$Z \to Y$ est fini d'après le lemme \ref{functors-lemma-characterize-finite}.
\end{proof}

\begin{lemma}
\label{functors-lemma-pushforward-invertible-pre}
Soit $f : X \to Y$ un morphisme séparé de type fini de schémas. Soit
$\mathcal{F}$ un module quasi-cohérent de type fini sur $X$, dont le support
est fini sur $Y$ et tel que $\mathcal{L} = f_*\mathcal{F}$ soit un
$\mathcal{O}_X$-module inversible. Alors il existe une section
$s : Y \to X$ telle que $\mathcal{F} \cong s_*\mathcal{L}$.
\end{lemma}

\begin{proof}
Localement dans le cas affine, l'énoncé se traduit par le problème algébrique
suivant. Soit $A \to B$ un homomorphisme d'anneaux et soit $N$ un $B$-module
qui est inversible comme $A$-module. Alors l'annulateur $J$ de $N$ dans $B$
est tel que $A \to B/J$ soit un isomorphisme. Nous omettons les détails.
\end{proof}

\begin{lemma}
\label{functors-lemma-pushforward-invertible}
Soit $f : X \to Y$ un morphisme séparé de type fini de schémas, muni d'une
section $s : Y \to X$. Soit $\mathcal{F}$ un module quasi-cohérent de type
fini sur $X$, à support ensembliste dans $s(Y)$, et tel que
$\mathcal{L} = f_*\mathcal{F}$ soit un $\mathcal{O}_X$-module inversible.
Si $Y$ est réduit, alors $\mathcal{F} \cong s_*\mathcal{L}$.
\end{lemma}

\begin{proof}
D'après le lemme \ref{functors-lemma-pushforward-invertible-pre}, il existe une section
$s' : Y  \to X$ telle que $\mathcal{F} = s'_*\mathcal{L}$. Comme $s'(Y)$
et $s(Y)$ ont la même partie fermée sous-jacente et sont tous deux des
sous-schémas fermés réduits de $X$, ils sont égaux.
Ainsi $s = s'$ et le lemme est démontré.
\end{proof}

\begin{lemma}
\label{functors-lemma-equivalence-coherent-over-field}
\begin{reference}
Version faible du résultat de \cite{Gabriel} selon lequel la catégorie des
modules quasi-cohérents détermine la classe d'isomorphisme d'un schéma.
\end{reference}
Soit $k$ un corps. Soient $X$ et $Y$ des schémas de type fini sur $k$,
avec $X$ séparé et $Y$ réduit. S'il existe une équivalence $k$-linéaire
$F : \textit{Coh}(\mathcal{O}_X) \to \textit{Coh}(\mathcal{O}_Y)$
de catégories, alors il existe un isomorphisme $f : Y \to X$ sur $k$ et un
$\mathcal{O}_Y$-module inversible $\mathcal{L}$ tels que
$F(\mathcal{F}) = f^*\mathcal{F} \otimes \mathcal{L}$.
\end{lemma}

\begin{proof}[Démonstration par reconstruction de Gabriel--Rosenberg]
Ce lemme est une forme faible des résultats exposés dans le
lemme \ref{functors-lemma-equivalence-coherent} et la
remarque \ref{functors-remark-equivalence-coherent-linear}.
\end{proof}

\begin{proof}[Démonstration indépendante de la reconstruction de Gabriel--Rosenberg]
D'après le lemme \ref{functors-lemma-functor-coherent-over-field}, nous obtenons un
$\mathcal{O}_{X \times Y}$-module cohérent $\mathcal{K}$, plat sur $X$ et
dont le support est fini sur $Y$, tel que $F$ soit donné par la restriction
du foncteur (\ref{functors-equation-FM-QCoh}) à $\textit{Coh}(\mathcal{O}_X)$.
Si nous pouvons montrer que $F(\mathcal{O}_X)$ est un
$\mathcal{O}_Y$-module inversible, alors le
lemme \ref{functors-lemma-pushforward-invertible-pre} montre que
$\mathcal{K} = s_*\mathcal{L}$ pour une section
$s : Y \to X \times Y$ de $\text{pr}_2$ et un
$\mathcal{O}_Y$-module inversible $\mathcal{L}$.
Cela montrera que $F$ est de la forme indiquée, avec
$f = \text{pr}_1 \circ s$. Nous omettons certains détails.

\medskip\noindent
Il reste à montrer que $F(\mathcal{O}_X)$ est inversible. Nous n'en donnons
qu'une esquisse et omettons certains détails.
Pour un point fermé $x \in X$, nous notons $\mathcal{O}_x$ dans
$\textit{Coh}(\mathcal{O}_X)$ le faisceau gratte-ciel en $x$ de valeur
$\kappa(x)$. Observons d'abord que les seuls objets simples de la catégorie
$\textit{Coh}(\mathcal{O}_X)$ sont ces faisceaux gratte-ciel
$\mathcal{O}_x$. Il en va de même pour $Y$. Ainsi, pour tout point fermé
$y \in Y$, il existe un point fermé $x \in X$ tel que
$\mathcal{O}_y \cong F(\mathcal{O}_x)$. De plus, en considérant les
endomorphismes, nous trouvons $\kappa(x) \cong \kappa(y)$ comme extensions
finies de $k$. Alors
$$
\Hom_Y(F(\mathcal{O}_X), \mathcal{O}_y) \cong
\Hom_Y(F(\mathcal{O}_X), F(\mathcal{O}_x)) \cong
\Hom_X(\mathcal{O}_X, \mathcal{O}_x) \cong \kappa(x) \cong \kappa(y)
$$
Il en résulte que le germe du $\mathcal{O}_Y$-module cohérent
$F(\mathcal{O}_X)$ en $y \in Y$ peut être engendré par $1$ générateur (et
pas moins), pour tout point fermé $y \in Y$. Il s'ensuit immédiatement que
$F(\mathcal{O}_X)$ est localement engendré par $1$ élément (et pas moins) et,
comme $Y$ est réduit, cela montre bien qu'il s'agit d'un module inversible.
\end{proof}













% OK, commencez ici.
%

\chapter{Catégories dérivées des variétés}



\phantomsection
\label{equiv-section-phantom}


\section{Introduction}
\label{equiv-section-introduction}

\noindent
Dans ce chapitre, nous poursuivons la discussion entamée dans
Catégories dérivées des schémas, section \ref{perfect-section-introduction}.
Nous étudierons les transformations de Fourier-Mukai,
introduites par Mukai dans \cite{Mukai}.
Nous démontrerons le théorème d'Orlov sur les équivalences dérivées (\cite{Orlov-K3}).
Nous étudierons également la dénombrabilité des classes d'équivalence dérivée,
établie par Anel et Toën dans \cite{AT}.

\medskip\noindent
Le livre de Daniel Huybrechts \cite{Huybrechts} constitue une bonne
introduction à ce sujet. Parmi les autres articles qui ont contribué
à populariser ce domaine, citons
\begin{enumerate}
\item l'article de Bondal et Kapranov, voir \cite{Bondal-Kapranov},
\item l'article de Bondal et Orlov, voir \cite{Bondal-Orlov},
\item l'article de Bondal et Van den Bergh, voir \cite{BvdB},
\item les articles de Beilinson, voir
\cite{Beilinson} et \cite{Beilinson-derived},
\item l'article d'Orlov, voir \cite{Orlov-AV},
\item l'article d'Orlov, voir \cite{Orlov-motives},
\item l'article de Rouquier, voir \cite{Rouquier-dimensions},
\item et bien d'autres encore pourraient être cités ici.
\end{enumerate}




\section{Conventions et notations}
\label{equiv-section-conventions}

\noindent
Soit $k$ un corps. Une catégorie triangulée $k$-linéaire $\mathcal{T}$
est une catégorie triangulée (Catégories dérivées, section
\ref{derived-section-triangulated-definitions})
munie d'une structure $k$-linéaire
(Algèbre différentielle graduée, section \ref{dga-section-linear})
telle que les foncteurs de translation $[n] : \mathcal{T} \to \mathcal{T}$
soient $k$-linéaires pour tout $n \in \mathbf{Z}$.

\medskip\noindent
Soit $k$ un corps. Nous notons $\text{Vect}_k$ la catégorie des
espaces vectoriels sur $k$. Pour un espace vectoriel sur $k$, noté $V$, nous notons
$V^\vee$ le dual $k$-linéaire de $V$, c'est-à-dire $V^\vee = \Hom_k(V, k)$.

\medskip\noindent
Soit $X$ un schéma. Nous notons $D_{perf}(\mathcal{O}_X)$ la
sous-catégorie pleine de $D(\mathcal{O}_X)$ formée des complexes parfaits
(Cohomologie, section \ref{cohomology-section-perfect}).
Si $X$ est noethérien, alors
$D_{perf}(\mathcal{O}_X) \subset D^b_{\textit{Coh}}(\mathcal{O}_X)$, voir
Catégories dérivées des schémas, lemme \ref{perfect-lemma-perfect-on-noetherian}.
Si $X$ est noethérien et régulier, alors
$D_{perf}(\mathcal{O}_X) = D^b_{\textit{Coh}}(\mathcal{O}_X)$, voir
Catégories dérivées des schémas, lemme \ref{perfect-lemma-perfect-on-regular}.

\medskip\noindent
Soit $k$ un corps. Soient $X$ et $Y$ des schémas sur $k$. Dans cette
situation, nous écrirons $X \times Y$ au lieu de $X \times_{\Spec(k)} Y$.


\medskip\noindent
Soit $S$ un schéma. Soient $X$ et $Y$ des schémas sur $S$.
Soit $\mathcal{F}$ un $\mathcal{O}_X$-module et soit
$\mathcal{G}$ un $\mathcal{O}_Y$-module. Nous posons
$$
\mathcal{F} \boxtimes \mathcal{G} =
\text{pr}_1^*\mathcal{F} \otimes_{\mathcal{O}_{X \times_S Y}}
\text{pr}_2^*\mathcal{G}
$$
comme $\mathcal{O}_{X \times_S Y}$-module.
Si $K \in D(\mathcal{O}_X)$ et $M \in D(\mathcal{O}_Y)$, nous posons
$$
K \boxtimes M =
L\text{pr}_1^*K \otimes_{\mathcal{O}_{X \times_S Y}}^\mathbf{L} L\text{pr}_2^*M
$$
comme objet de $D(\mathcal{O}_{X \times_S Y})$.
Notre notation est donc potentiellement ambiguë, mais le contexte devrait
indiquer clairement laquelle des deux significations est visée.





\section{Foncteurs de Serre}
\label{equiv-section-Serre-functors}

\noindent
Le contenu de cette section est repris de \cite{Bondal-Kapranov}.

\begin{lemma}
\label{equiv-lemma-Serre-functor-exists}
Soit $k$ un corps. Soit $\mathcal{T}$ une catégorie triangulée
$k$-linéaire telle que $\dim_k \Hom_\mathcal{T}(X, Y) < \infty$
pour tous $X, Y \in \Ob(\mathcal{T})$. Les assertions suivantes sont équivalentes :
\begin{enumerate}
\item il existe une équivalence $k$-linéaire
$S : \mathcal{T} \to \mathcal{T}$ et des isomorphismes $k$-linéaires
$c_{X, Y} : \Hom_\mathcal{T}(X, Y) \to \Hom_\mathcal{T}(Y, S(X))^\vee$
fonctoriels en $X, Y \in \Ob(\mathcal{T})$ ;
\item pour tout $X \in \Ob(\mathcal{T})$,
le foncteur $Y \mapsto \Hom_\mathcal{T}(X, Y)^\vee$
est représentable et le foncteur $Y \mapsto \Hom_\mathcal{T}(Y, X)^\vee$
est coreprésentable.
\end{enumerate}
\end{lemma}

\begin{proof}
La condition (1) entraîne (2) car, étant donnés $(S, c)$ et $X \in \Ob(\mathcal{T})$,
l'objet $S(X)$ représente le foncteur
$Y \mapsto \Hom_\mathcal{T}(X, Y)^\vee$ et l'objet $S^{-1}(X)$ coreprésente
le foncteur $Y \mapsto \Hom_\mathcal{T}(Y, X)^\vee$.

\medskip\noindent
Supposons (2). Nous utiliserons à plusieurs reprises le lemme de Yoneda, voir
Catégories, lemme \ref{categories-lemma-yoneda}.
Pour tout $X$, notons $S(X)$ l'objet qui représente le
foncteur $Y \mapsto \Hom_\mathcal{T}(X, Y)^\vee$. Étant donné
$\varphi : X \to X'$, on obtient un unique morphisme $S(\varphi) : S(X) \to S(X')$
déterminé par la transformation de foncteurs correspondante
$\Hom_\mathcal{T}(X, -)^\vee \to \Hom_\mathcal{T}(X', -)^\vee$.
Ainsi, $S$ est un foncteur et l'on obtient les isomorphismes $c_{X, Y}$
par construction. Il reste à montrer que $S$ est une équivalence.
Pour tout $X$, notons $S'(X)$ l'objet qui coreprésente le
foncteur $Y \mapsto \Hom_\mathcal{T}(Y, X)^\vee$. En raisonnant comme
ci-dessus, on voit que $S'$ est un foncteur. Nous affirmons que $S'$
est un quasi-inverse de $S$. En effet, observons que
$$
\Hom_\mathcal{T}(X, Y) = \Hom_\mathcal{T}(Y, S(X))^\vee =
\Hom_\mathcal{T}(S'(S(X)), Y)
$$
bifonctoriellement, c'est-à-dire que $S' \circ S \cong \text{id}_\mathcal{T}$.
De même, on a
$$
\Hom_\mathcal{T}(Y, X) = \Hom_\mathcal{T}(S'(X), Y)^\vee =
\Hom_\mathcal{T}(Y, S(S'(X)))
$$
et l'on obtient $S \circ S' \cong \text{id}_\mathcal{T}$.
\end{proof}

\begin{definition}
\label{equiv-definition-Serre-functor}
Soit $k$ un corps. Soit $\mathcal{T}$ une catégorie triangulée
$k$-linéaire telle que $\dim_k \Hom_\mathcal{T}(X, Y) < \infty$
pour tous $X, Y \in \Ob(\mathcal{T})$. On dit qu'{\it il existe un foncteur
de Serre} si les conditions équivalentes du lemme \ref{equiv-lemma-Serre-functor-exists}
sont satisfaites. Dans ce cas, un {\it foncteur de Serre} est une équivalence $k$-linéaire
$S : \mathcal{T} \to \mathcal{T}$ munie d'isomorphismes $k$-linéaires
$c_{X, Y} : \Hom_\mathcal{T}(X, Y) \to \Hom_\mathcal{T}(Y, S(X))^\vee$
fonctoriels en $X, Y \in \Ob(\mathcal{T})$.
\end{definition}

\begin{remark}
\label{equiv-remark-matress}
Soient $X^0 \to X^1 \to X^2 \to X^0[1]$ et $Y^0 \to Y^1 \to Y^2 \to Y^0[1]$
des triangles distingués dans une catégorie triangulée. Pour
$p \in \mathbf{Z}$, écrivons $p = 3n + i$ avec $i \in \{0, 1, 2\}$ et
posons $X^p = X^i[n]$. De même pour $Y^q$. Considérons le complexe double
dont les termes sont
$$
K^{p, q} = \Hom(X^{-p}, Y^q)
$$
La différentielle $d_1 : K^{p, q} \to K^{p + 1, q}$ est induite par le morphisme
$X^{-p - 1} \to X^{-p}$ (égal au morphisme correspondant
du premier triangle distingué, à un décalage près),
et la différentielle $d_2 : K^{p, q} \to K^{p, q + 1}$ est de même induite
par le morphisme $Y^q \to Y^{q + 1}$. D'après
Catégories dérivées, lemme \ref{derived-lemma-representable-homological},
les lignes et les colonnes de ce complexe double sont exactes.
De plus, $K^{p, q} = K^{p + 3, q - 3}$, et ces
égalités sont compatibles avec les différentielles. Enfin, l'axiome TR3
entraîne une propriété supplémentaire : étant donnés
$\alpha \in K^{p, q}$ et $\beta \in K^{p - 1, q + 1}$ tels
que $d_2 \alpha = d_1 \beta$, il existe $\gamma \in K^{p - 2, q + 2}$
tel que
$d_1 \gamma = d_2 \beta $ dans $K^{p - 1, q + 2}$ et
$d_2 \gamma = d_1 \alpha $ dans $K^{p - 2, q + 3} = K^{p + 1, q}$.
(Indication : pour $p = q = 0$, c'est exactement l'énoncé de TR3 ;
pour les autres indices, on le démontre par décalage.)
Un complexe double ayant ces propriétés est appelé {\it matelas}
(voir \cite{Bondal-Kapranov}).
\end{remark}

\begin{remark}
\label{equiv-remark-dual-matress}
Soit $k$ un corps et soit $K^{\bullet, \bullet}$ un complexe double
d'espaces vectoriels de dimension finie sur $k$, qui est aussi un {\it matelas} au sens de la
remarque \ref{equiv-remark-matress}. Nous affirmons que le complexe double dual
$L^{p, q} = \Hom_k(K^{-q, -p}, k)$ est lui aussi un {\it matelas}. L'exactitude des
lignes et des colonnes ainsi que la 3-périodicité sont immédiates. Pour vérifier
la condition supplémentaire, considérons l'application linéaire
$$
\partial :
K^{p, q} \oplus K^{p - 1, q + 1} \oplus K^{p - 2, q + 2}
\longrightarrow
K^{p, q + 1} \oplus K^{p - 1, q + 2} \oplus K^{p - 2, q + 3}
$$
qui envoie $(\alpha, \beta, \gamma)$ sur
$(d_2 \alpha - d_1 \beta, d_2 \gamma - d_1 \alpha, d_1 \gamma - d_2 \beta)$,
avec les identifications de la remarque \ref{equiv-remark-matress}.
La condition d'être un {\it matelas} s'écrit
$$
\Im(\partial) \cap
\left( 0 \oplus K^{p - 1, q + 2} \oplus K^{p - 2, q + 3} \right)
=
\Im(\partial|_{0 \oplus 0 \oplus K^{p - 2, q + 2}})
$$
En notant ${}^\wedge$ le dual d'un espace vectoriel ou d'une application, la dualisation
donne
$$
\Ker(\partial^\wedge) +
\left(L^{-p, -q - 1} \oplus 0 \oplus 0 \right)
=
\Ker(\text{pr}_{L^{-p + 2, -q - 2}} \circ \partial^\wedge)
$$
où le terme $\text{pr}$ désigne la projection. Cela implique
que
$\Im(\partial^\wedge) \cap (L^{-p, -q} \oplus L^{-p + 1, -q - 1} \oplus 0)$
est égal à
$\Im(\partial^\wedge|_{L^{-p, -q - 1} \oplus 0 \oplus 0})$, ce qui
se traduit par la condition de {\it matelas} pour $L^{\bullet, \bullet}$.
Nous omettons certains détails.
\end{remark}

\begin{lemma}
\label{equiv-lemma-Serre-functor}
Dans la situation de la définition \ref{equiv-definition-Serre-functor},
s'il existe un foncteur de Serre, celui-ci est unique à isomorphisme unique près
et c'est un foncteur exact de catégories triangulées.
\end{lemma}

\begin{proof}
Étant donné un foncteur de Serre $S$, l'objet $S(X)$ représente
le foncteur $Y \mapsto \Hom_\mathcal{T}(X, Y)^\vee$.
Ainsi, l'objet $S(X)$, muni de l'identification fonctorielle
$\Hom_\mathcal{T}(X, Y)^\vee = \Hom_\mathcal{T}(Y, S(X))$,
est déterminé à isomorphisme unique près par le lemme de Yoneda
(Catégories, lemme \ref{categories-lemma-yoneda}).
De plus, pour $\varphi : X \to X'$, le morphisme $S(\varphi) : S(X) \to S(X')$
est uniquement déterminé par la transformation de foncteurs correspondante
$\Hom_\mathcal{T}(X, -)^\vee \to \Hom_\mathcal{T}(X', -)^\vee$.

\medskip\noindent
Pour des objets $X, Y$ de $\mathcal{T}$, on a
\begin{align*}
\Hom(Y, S(X)[1])^\vee
& =
\Hom(Y[-1], S(X))^\vee \\
& =
\Hom(X, Y[-1]) \\
& =
\Hom(X[1], Y) \\
& =
\Hom(Y, S(X[1]))^\vee
\end{align*}
Le lemme de Yoneda fournit donc un unique isomorphisme
$S(X[1]) \to S(X)[1]$ qui induit l'isomorphisme obtenu de la première à la dernière expression.
Comme chacun des isomorphismes ci-dessus est fonctoriel à la fois en $X$ et en $Y$,
on obtient ainsi un isomorphisme de foncteurs
$S \circ [1] \to [1] \circ S$.

\medskip\noindent
Soit $(A, B, C, f, g, h)$ un triangle distingué de $\mathcal{T}$.
Il faut montrer que le triangle $(S(A), S(B), S(C), S(f), S(g), S(h))$
est distingué. On utilise ici l'isomorphisme canonique $S(A[1]) \to S(A)[1]$
construit ci-dessus pour identifier le but $S(A[1])$ de $S(h)$ à $S(A)[1]$.
Observons d'abord que, pour tout $X$ dans $\mathcal{T}$,
le triangle $(S(A), S(B), S(C), S(f), S(g), S(h))$ induit
une suite exacte longue
$$
\ldots \to
\Hom(X, S(A)) \to
\Hom(X, S(B)) \to
\Hom(X, S(C)) \to
\Hom(X, S(A)[1]) \to \ldots
$$
d'espaces vectoriels de dimension finie sur $k$. En effet, cette suite est
le dual $k$-linéaire de la suite
$$
\ldots \leftarrow
\Hom(A, X) \leftarrow
\Hom(B, X) \leftarrow
\Hom(C, X) \leftarrow
\Hom(A[1], X) \leftarrow
\ldots
$$
qui est exacte d'après Catégories dérivées, lemme
\ref{derived-lemma-representable-homological}.
Choisissons ensuite un triangle distingué $(S(A), E, S(C), i, p, S(h))$ ;
c'est possible d'après les axiomes TR1 et TR2. Nous voulons construire la flèche
en pointillés qui rende commutatif le diagramme suivant :
$$
\xymatrix{
S(C)[-1] \ar[r]_-{S(h[-1])} &
S(A) \ar[r]_{S(f)} &
S(B) \ar[r]_{S(g)} &
S(C) \ar[r]_{S(h)} &
S(A)[1] \\
S(C)[-1] \ar[r]^-{S(h[-1])} \ar@{=}[u] &
S(A) \ar[r]^i \ar@{=}[u] &
E \ar[r]^p \ar@{..>}[u]^\varphi &
S(C) \ar[r]^{S(h)} \ar@{=}[u] &
S(A)[1] \ar@{=}[u]
}
$$
En effet, si l'on dispose de $\varphi$, alors, pour tout $X$, l'application obtenue
$\Hom(X, E) \to \Hom(X, S(B))$ est un isomorphisme d'espaces vectoriels
sur $k$. On obtient en effet un diagramme commutatif
$$
\xymatrix{
\Hom(X, S(C)[-1]) \ar[r] &
\Hom(X, S(A)) \ar[r] &
\Hom(X, S(B)) \ar[r] &
\Hom(X, S(C)) \ar[r] &
\Hom(X, S(A)[1]) \\
\Hom(X, S(C)[-1]) \ar[r] \ar@{=}[u] &
\Hom(X, S(A)) \ar[r] \ar@{=}[u] &
\Hom(X, E) \ar[r] \ar[u]^\varphi &
\Hom(X, S(C)) \ar[r] \ar@{=}[u] &
\Hom(X, S(A)[1]) \ar@{=}[u]
}
$$
dont les lignes sont exactes (voir ci-dessus), et le lemme des cinq
(Homologie, lemme \ref{homology-lemma-five-lemma}) montre
que le morphisme médian est un isomorphisme. Le lemme de Yoneda
permet alors de conclure que $\varphi$ est un isomorphisme.

\medskip\noindent
Pour construire $\varphi$, formons le {\it matelas} de la remarque \ref{equiv-remark-matress}
à partir des triangles distingués
$A \to B \to C \to A[1]$ et $S(A) \to E \to S(C) \to S(A)[1]$.
D'après la remarque \ref{equiv-remark-dual-matress}, son dual $k$-linéaire
est également un {\it matelas}. Comme $S$ est un foncteur de Serre, on
voit, comme ci-dessus, que ce {\it matelas} dual est le complexe double formé par
les triangles (dont le second n'est pas encore connu pour être distingué)
$S(A) \to E \to S(C) \to S(A)[1]$ et $S(A) \to S(B) \to S(C) \to S(A)[1]$.
Or la condition de {\it matelas} dit exactement que les morphismes
$S(A) \to S(A)$ et $S(C) \to S(C)$ s'insèrent dans un morphisme de triangles,
c'est-à-dire qu'il existe un $\varphi$ qui rend commutatif le diagramme ci-dessus.
\end{proof}






\section{Exemples de foncteurs de Serre}
\label{equiv-section-examples-Serre-functors}

\noindent
Le lemme ci-dessous est l'exemple classique.

\begin{lemma}
\label{equiv-lemma-Serre-functor-Gorenstein-proper}
Soit $k$ un corps. Soit $X$ un schéma propre sur $k$ qui est de Gorenstein.
Considérons le complexe $\omega_X^\bullet$ du
lemme \ref{duality-lemma-duality-proper-over-field} de Dualité pour les schémas.
Alors le foncteur
$$
S : D_{perf}(\mathcal{O}_X) \longrightarrow D_{perf}(\mathcal{O}_X),\quad
K \longmapsto S(K) = \omega_X^\bullet \otimes_{\mathcal{O}_X}^\mathbf{L} K
$$
est un foncteur de Serre.
\end{lemma}

\begin{proof}
L'énoncé a un sens puisque $\dim \Hom_X(K, L) < \infty$
pour $K, L \in D_{perf}(\mathcal{O}_X)$ d'après
Catégories dérivées des schémas, lemme \ref{perfect-lemma-ext-finite}.
Comme $X$ est de Gorenstein, le complexe dualisant $\omega_X^\bullet$
est un objet inversible de $D(\mathcal{O}_X)$, voir
Dualité pour les schémas, lemme \ref{duality-lemma-gorenstein}.
En particulier, localement sur $X$, le complexe $\omega_X^\bullet$
possède un unique faisceau de cohomologie non nul, lequel est un module inversible, voir
Cohomologie, lemme \ref{cohomology-lemma-invertible-derived}.
Ainsi, $S(K)$ appartient à $D_{perf}(\mathcal{O}_X)$.
D'autre part, l'inversibilité de $\omega_X^\bullet$
implique clairement que $S$ est une auto-équivalence de $D_{perf}(\mathcal{O}_X)$.
Enfin, il faut construire un isomorphisme
$$
c_{K, L} : \Hom_X(K, L) \longrightarrow
\Hom_X(L, \omega_X^\bullet \otimes_{\mathcal{O}_X}^\mathbf{L} K)^\vee
$$
bifonctoriel en $K, L$. Pour cela, utilisons les isomorphismes canoniques
$$
\Hom_X(K, L) = H^0(X, L \otimes_{\mathcal{O}_X}^\mathbf{L} K^\vee)
$$
et
$$
\Hom_X(L, \omega_X^\bullet \otimes_{\mathcal{O}_X}^\mathbf{L} K) =
H^0(X, 
\omega_X^\bullet \otimes_{\mathcal{O}_X}^\mathbf{L} K
\otimes_{\mathcal{O}_X}^\mathbf{L} L^\vee)
$$
donnés dans Cohomologie, lemme \ref{cohomology-lemma-dual-perfect-complex}.
Puisque $(L \otimes_{\mathcal{O}_X}^\mathbf{L} K^\vee)^\vee =
(K^\vee)^\vee \otimes_{\mathcal{O}_X}^\mathbf{L} L^\vee$
et qu'il existe un isomorphisme canonique $K \to (K^\vee)^\vee$,
ces espaces vectoriels sur $k$ sont canoniquement duaux d'après
Dualité pour les schémas, lemme
\ref{duality-lemma-duality-proper-over-field-perfect}.
On obtient ainsi les isomorphismes $c_{K, L}$.
Nous omettons la démonstration de leur fonctorialité.
\end{proof}





\section{Caractérisation des modules cohérents}
\label{equiv-section-coherent}

\noindent
Cette section prolonge, en un certain sens, la discussion menée
dans Catégories dérivées des schémas, section \ref{perfect-section-pseudo-coherent}
et Compléments sur les morphismes, section
\ref{more-morphisms-section-characterize-pseudo-coherent}.

\medskip\noindent
Avant de pouvoir énoncer le résultat, nous avons besoin de quelques notations.
Soit $k$ un corps. Soit $n \geq 0$ un entier.
Soit $S = k[X_0, \ldots, X_n]$. Pour un entier $e$, notons
$S_e \subset S$ l'espace des polynômes homogènes de degré $e$.
Considérons la $k$-algèbre (non commutative)
$$
R =
\left(
\begin{matrix}
S_0 & S_1 & S_2 & \ldots & \ldots \\
0 & S_0 & S_1 & \ldots & \ldots\\
0 & 0 & S_0 & \ldots & \ldots \\
\ldots & \ldots & \ldots & \ldots & \ldots \\
0 & \ldots & \ldots & \ldots & S_0
\end{matrix}
\right)
$$
(à $n + 1$ lignes et colonnes), munie de la multiplication et de l'addition évidentes.

\begin{lemma}
\label{equiv-lemma-perfect-for-R}
Avec $k$, $n$ et $R$ comme ci-dessus, pour un objet $K$ de $D(R)$,
les assertions suivantes sont équivalentes :
\begin{enumerate}
\item $\sum_{i \in \mathbf{Z}} \dim_k H^i(K) < \infty$, et
\item $K$ est un objet compact.
\end{enumerate}
\end{lemma}

\begin{proof}
Si $K$ est un objet compact, alors $K$ peut être représenté par un complexe
$M^\bullet$ qui, comme $R$-module gradué, est projectif de type fini, voir
Algèbre différentielle graduée, lemme \ref{dga-lemma-compact}.
Puisque $\dim_k R < \infty$, nous concluons que $\sum \dim_k M^i < \infty$
et, a fortiori, que $\sum \dim_k H^i(M^\bullet) < \infty$.
(On peut aussi déduire facilement cette implication de la proposition plus élémentaire
Algèbre différentielle graduée, proposition \ref{dga-proposition-compact}.)

\medskip\noindent
Supposons que $K$ satisfasse (1). Considérons le triangle distingué
de troncation $\tau_{\leq m}K \to K \to \tau_{\geq m + 1}K$, voir
Catégories dérivées, remarque
\ref{derived-remark-truncation-distinguished-triangle}.
Il est clair que $\tau_{\leq m}K$ et $\tau_{\geq m + 1} K$
satisfont (1). Si l'on peut montrer qu'ils sont tous deux compacts, alors $K$ l'est aussi, voir
Catégories dérivées, lemme \ref{derived-lemma-compact-objects-subcategory}.
Ainsi, par récurrence sur le nombre de modules de cohomologie non nuls de $K$,
nous pouvons supposer que $H^i(K)$ n'est non nul que pour une seule valeur de $i$.
En décalant, nous pouvons supposer que $K$ est donné par le complexe
constitué d'un seul $R$-module $M$ de dimension finie, placé
en degré $0$.

\medskip\noindent
Puisque $\dim_k(M) < \infty$, nous voyons que $M$ est artinien comme $R$-module.
Il suffit donc de montrer que tout $R$-module simple représente un
objet compact de $D(R)$. Observons que
$$
I =
\left(
\begin{matrix}
0 & S_1 & S_2 & \ldots & \ldots \\
0 & 0 & S_1 & \ldots & \ldots\\
0 & 0 & 0 & \ldots & \ldots \\
\ldots & \ldots & \ldots & \ldots & \ldots \\
0 & \ldots & \ldots & \ldots & 0
\end{matrix}
\right)
$$
est un idéal bilatère nilpotent de $R$ et que $R/I$
est une $k$-algèbre commutative isomorphe à un produit de $n + 1$ copies de
$k$ (placées sur la diagonale de la matrice, c'est-à-dire que $R/I$ peut être relevée
en une sous-$k$-algèbre de $R$). Il s'ensuit que $R$ possède exactement $n + 1$
classes d'isomorphisme de modules simples $M_0, \ldots, M_n$ (placés le long
de la diagonale). Considérons le $R$-module à droite $P_i$ de vecteurs lignes
$$
P_i =
\left(
\begin{matrix}
0 &
\ldots &
0 &
S_0 &
\ldots &
S_{i - 1} &
S_i
\end{matrix}
\right)
$$
muni de la multiplication évidente $P_i \times R \to P_i$. Nous voyons alors que
$R \cong P_0 \oplus \ldots \oplus P_n$ comme $R$-module à droite. Comme il est clair que
$R$ est un objet compact de $D(R)$, nous concluons que chaque $P_i$ est un objet compact
de $D(R)$. (Nous concluons bien sûr aussi que chaque $P_i$ est projectif
comme $R$-module, mais ce n'est pas ce que nous avons à montrer dans cette démonstration.)
Il est clair que $P_0 = M_0$ est le premier de nos modules simples sur $R$.
Pour $P_1$, nous avons une suite exacte courte
$$
0 \to P_0^{\oplus n + 1} \to P_1 \to M_1 \to 0
$$
qui montre que $M_1$ s'insère dans un triangle distingué dont
les autres membres sont des objets compacts ; ainsi, $M_1$ est un objet compact
de $D(R)$. Plus généralement, il existe une suite exacte courte
$$
0 \to C_i \to P_i \to M_i \to 0
$$
où $C_i$ est un $R$-module de dimension finie dont les facteurs de composition
sont isomorphes à $M_j$ pour $j < i$. Par récurrence, nous concluons d'abord que
$C_i$ détermine un objet compact de $D(R)$, puis que $M_i$
en détermine un également, comme souhaité.
\end{proof}

\begin{lemma}
\label{equiv-lemma-coherent-on-projective-space}
Soit $k$ un corps. Soit $n \geq 0$. Soit
$K \in D_\QCoh(\mathcal{O}_{\mathbf{P}^n_k})$.
Les assertions suivantes sont équivalentes :
\begin{enumerate}
\item $K$ appartient à $D^b_{\textit{Coh}}(\mathcal{O}_{\mathbf{P}^n_k})$,
\item $\sum_{i \in \mathbf{Z}}
\dim_k H^i(\mathbf{P}^n_k, E \otimes^\mathbf{L} K) < \infty$
pour tout objet parfait $E$ de
$D(\mathcal{O}_{\mathbf{P}^n_k})$,
\item $\sum_{i \in \mathbf{Z}}
\dim_k \Ext^i_{\mathbf{P}^n_k}(E, K) < \infty$
pour tout objet parfait $E$ de $D(\mathcal{O}_{\mathbf{P}^n_k})$,
\item $\sum_{i \in \mathbf{Z}} \dim_k H^i(\mathbf{P}^n_k,
K \otimes^\mathbf{L} \mathcal{O}_{\mathbf{P}^n_k}(d)) < \infty$
pour $d = 0, 1, \ldots, n$.
\end{enumerate}
\end{lemma}

\begin{proof}
Les assertions (2) et (3) sont équivalentes d'après
Cohomologie, lemme \ref{cohomology-lemma-dual-perfect-complex}.
Si (1) est satisfaite, alors, pour $E$ parfait, le produit tensoriel dérivé
$E \otimes^\mathbf{L} K$ appartient à
$D^b_{\textit{Coh}}(\mathcal{O}_{\mathbf{P}^n_k})$
et l'on voit que (2) est satisfaite d'après
Catégories dérivées des schémas, lemme \ref{perfect-lemma-direct-image-coherent}.
Il est clair que (2) implique (4), puisque $\mathcal{O}_{\mathbf{P}^n_k}(d)$
peut être considéré
comme un objet parfait de la catégorie dérivée de $\mathbf{P}^n_k$.
Ainsi, il suffit de montrer que (4) implique (1).

\medskip\noindent
Supposons (4). Soit $R$ comme dans le lemme \ref{equiv-lemma-perfect-for-R}.
Soit $P = \bigoplus_{d = 0, \ldots, n} \mathcal{O}_{\mathbf{P}^n_k}(-d)$.
Rappelons que $R = \text{End}_{\mathbf{P}^n_k}(P)$, tandis que tous les autres groupes
d'extensions de $P$ par lui-même sont nuls, et que $P$ détermine une équivalence
$- \otimes^\mathbf{L} P : D(R) \to D_\QCoh(\mathcal{O}_{\mathbf{P}^n_k})$
d'après Catégories dérivées des schémas, lemme \ref{perfect-lemma-Pn-module-category}.
Disons que $K$ correspond à $L$ dans $D(R)$. Alors
\begin{align*}
H^i(L)
& =
\Ext^i_{D(R)}(R, L) \\
& =
\Ext^i_{\mathbf{P}^n_k}(P, K) \\
& =
H^i(\mathbf{P}^n_k, K \otimes P^\vee) \\
& =
\bigoplus\nolimits_{d = 0, \ldots, n}
H^i(\mathbf{P}^n_k, K \otimes \mathcal{O}(d))
\end{align*}
d'après Algèbre différentielle graduée, lemme
\ref{dga-lemma-upgrade-tensor-with-complex-derived}
(et le fait que $- \otimes^\mathbf{L} P$ est une équivalence)
et Cohomologie, lemme \ref{cohomology-lemma-dual-perfect-complex}.
Ainsi, notre hypothèse (4) implique que $L$ satisfait la condition (2) du
lemme \ref{equiv-lemma-perfect-for-R} et est donc un objet compact de $D(R)$.
Par conséquent, $K$ est un objet compact de
$D_\QCoh(\mathcal{O}_{\mathbf{P}^n_k})$.
Ainsi, $K$ est parfait d'après
Catégories dérivées des schémas, proposition
\ref{perfect-proposition-compact-is-perfect}.
Puisque $D_{perf}(\mathcal{O}_{\mathbf{P}^n_k}) =
D^b_{\textit{Coh}}(\mathcal{O}_{\mathbf{P}^n_k})$
d'après
Catégories dérivées des schémas, lemme \ref{perfect-lemma-perfect-on-regular},
nous concluons que (1) est satisfaite.
\end{proof}

\begin{lemma}
\label{equiv-lemma-finiteness}
Soit $X$ un schéma propre sur un corps $k$. Soit
$K \in D^b_{\textit{Coh}}(\mathcal{O}_X)$ et soit $E$ dans $D(\mathcal{O}_X)$
un objet parfait. Alors
$\sum_{i \in \mathbf{Z}} \dim_k \Ext^i_X(E, K) < \infty$.
\end{lemma}

\begin{proof}
Cela résulte, par exemple, de la combinaison des
lemmes \ref{perfect-lemma-ext-finite} et
\ref{perfect-lemma-ext-from-perfect-into-bounded-QCoh} de Catégories dérivées des schémas.
Autre démonstration : combiner
les lemmes de Catégories dérivées des schémas
\ref{perfect-lemma-perfect-on-noetherian} et
\ref{perfect-lemma-direct-image-coherent}.
\end{proof}

\begin{lemma}
\label{equiv-lemma-characterize-dbcoh-projective}
\begin{reference}
Dans le cas projectif, il s'agit de \cite[Lemme 7.46]{Rouquier-dimensions}
et ce résultat est implicite dans \cite[Théorème A.1]{BvdB}.
\end{reference}
Soit $X$ un schéma propre sur un corps $k$. Soit
$K \in \Ob(D_\QCoh(\mathcal{O}_X))$. Les assertions suivantes sont équivalentes :
\begin{enumerate}
\item $K \in D^b_{\textit{Coh}}(\mathcal{O}_X)$, et
\item $\sum_{i \in \mathbf{Z}} \dim_k \Ext^i_X(E, K) < \infty$
pour tout objet parfait $E$ de $D(\mathcal{O}_X)$.
\end{enumerate}
\end{lemma}

\begin{proof}
L'implication (1) $\Rightarrow$ (2) découle du
lemme \ref{equiv-lemma-finiteness}.
L'implication (2) $\Rightarrow$ (1) découle de
Compléments sur les morphismes, lemme
\ref{more-morphisms-lemma-characterize-relatively-perfect}
(voir Catégories dérivées des schémas, exemple
\ref{perfect-example-relatively-perfect-field}
pour la signification d'un objet relativement parfait sur un corps) ;
la démonstration plus simple du cas projectif figure au paragraphe suivant.

\medskip\noindent
Supposons (2) et que $X$ soit projectif sur $k$.
Choisissons une immersion fermée $i : X \to \mathbf{P}^n_k$. Il suffit de montrer
que $Ri_*K$ appartient à $D^b_{\textit{Coh}}(\mathbf{P}^n_k)$, car un module quasi-cohérent
$\mathcal{F}$ sur $X$ est cohérent, resp.\ nul, si et seulement si
$i_*\mathcal{F}$ est cohérent, resp.\ nul. Pour un objet parfait $E$
de $D(\mathcal{O}_{\mathbf{P}^n_k})$, $Li^*E$ est un objet parfait de
$D(\mathcal{O}_X)$ et
$$
\Ext^q_{\mathbf{P}^n_k}(E, Ri_*K) = \Ext^q_X(Li^*E, K)
$$
D'après notre hypothèse, nous voyons donc que
$\sum_{q \in \mathbf{Z}} \dim_k \Ext^q_{\mathbf{P}^n_k}(E, Ri_*K) < \infty$.
Nous concluons par le lemme \ref{equiv-lemma-coherent-on-projective-space}.
\end{proof}





\section{Un théorème de représentabilité}
\label{equiv-section-bondal-van-den-bergh}

\noindent
Le contenu de cette section est tiré de \cite{BvdB}.

\medskip\noindent
Soit $\mathcal{T}$ une catégorie triangulée $k$-linéaire.
Dans cette section, nous considérons les foncteurs cohomologiques $k$-linéaires
$H$ de $\mathcal{T}$ vers la catégorie des espaces vectoriels sur $k$.
Cela signifiera que $H$ est un foncteur
$$
H : \mathcal{T}^{opp} \longrightarrow \text{Vect}_k
$$
qui est $k$-linéaire et tel que, pour tout triangle distingué
$X \to Y \to Z$ de $\mathcal{T}$, la suite $H(Z) \to H(Y) \to H(X)$
est une suite exacte d'espaces vectoriels sur $k$. Voir
Catégories dérivées, définition \ref{derived-definition-homological}
et Algèbre différentielle graduée, section \ref{dga-section-linear}.

\begin{lemma}
\label{equiv-lemma-maps-from-compact-filtered}
Soit $\mathcal{D}$ une catégorie triangulée. Soit
$\mathcal{D}' \subset \mathcal{D}$ une sous-catégorie triangulée pleine. Soit
$X \in \Ob(\mathcal{D})$. La catégorie des flèches $E \to X$ avec
$E \in \Ob(\mathcal{D}')$ est filtrante.
\end{lemma}

\begin{proof}
Vérifions les conditions de
Catégories, définition \ref{categories-definition-directed}.
La catégorie est non vide, car elle contient $0 \to X$.
Si $E_i \to X$, $i = 1, 2$ sont des objets, alors $E_1 \oplus E_2 \to X$
est un objet et il existe des morphismes $(E_i \to X) \to (E_1 \oplus E_2 \to X)$.
Enfin, supposons que $a, b : (E \to X) \to (E' \to X)$ soient des morphismes.
Choisissons un triangle distingué $E \xrightarrow{a - b} E' \to E''$
dans $\mathcal{D}'$. Par l'axiome TR3, on obtient un morphisme de triangles
$$
\xymatrix{
E \ar[r]_{a - b} \ar[d] &
E' \ar[d] \ar[r] & E'' \ar[d] \\
0 \ar[r] &
X \ar[r] &
X
}
$$
et l'on constate que la flèche obtenue $(E' \to X) \to (E'' \to X)$
égalise $a$ et $b$.
\end{proof}

\begin{lemma}
\label{equiv-lemma-van-den-bergh}
\begin{reference}
\cite[Lemme 2.14]{CKN}
\end{reference}
Soit $k$ un corps. Soit $\mathcal{D}$ une catégorie triangulée $k$-linéaire
qui admet les sommes directes et est compactement engendrée.
Notons $\mathcal{D}_c$ la sous-catégorie pleine
des objets compacts. Soit $H : \mathcal{D}_c^{opp} \to \text{Vect}_k$
un foncteur cohomologique $k$-linéaire tel que
$\dim_k H(X) < \infty$ pour tout $X \in \Ob(\mathcal{D}_c)$.
Alors $H$ est isomorphe au foncteur $X \mapsto \Hom(X, Y)$
pour un certain $Y \in \Ob(\mathcal{D})$.
\end{lemma}

\begin{proof}
Nous utiliserons Catégories dérivées, lemme
\ref{derived-lemma-compact-objects-subcategory} sans autre mention.
Notons $G : \mathcal{D}_c \to \text{Vect}_k$ le foncteur homologique
$k$-linéaire qui envoie $X$ sur $H(X)^\vee$. Pour tout objet $Y$ de $\mathcal{D}$,
posons
$$
G'(Y) = \colim_{X \to Y, X \in \Ob(\mathcal{D}_c)} G(X)
$$
Il s'agit d'une colimite filtrante d'après le lemme \ref{equiv-lemma-maps-from-compact-filtered}.
Nous affirmons que $G'$ est un foncteur homologique $k$-linéaire,
que la restriction de $G'$ à $\mathcal{D}_c$ est $G$, et que $G'$
transforme les sommes directes en sommes directes.

\medskip\noindent
En effet, supposons que $Y_1 \to Y_2 \to Y_3$ soit un triangle distingué.
Soit $\xi \in G'(Y_2)$ dont l'image dans $G'(Y_3)$ est nulle. Comme la colimite est
filtrante, $\xi$ est représenté par un certain $X \to Y_2$ avec
$X \in \Ob(\mathcal{D}_c)$ et $g \in G(X)$.
Le fait que l'image de $\xi$ dans $G'(Y_3)$ soit nulle signifie que le composé
$X \to Y_2 \to Y_3$ se factorise sous la forme $X \to X' \to Y_3$ avec $X' \in \mathcal{D}_c$
et que l'image de $g$ dans $G(X')$ est nulle. Choisissons un triangle distingué
$X'' \to X \to X'$. Alors $X'' \in \Ob(\mathcal{D}_c)$.
Comme $G$ est homologique, on voit que $g$ est l'image d'un certain
$g'' \in G'(X'')$. Par l'axiome TR3, les morphismes $X \to Y_2$ et $X' \to Y_3$ s'insèrent dans
un morphisme de triangles distingués
$(X'' \to X \to X') \to (Y_1 \to Y_2 \to Y_3)$,
et l'on voit que $\xi$ est bien l'image de
l'élément de $G'(Y_1)$ représenté par $X'' \to Y_1$ et $g'' \in G(X'')$.

\medskip\noindent
Si $Y \in \Ob(\mathcal{D}_c)$, alors $\text{id} : Y \to Y$ est l'objet final
de la catégorie des flèches $X \to Y$ avec $X \in \Ob(\mathcal{D}_c)$.
On voit donc que $G'(Y) = G(Y)$ dans ce cas, ce qui prouve l'assertion
sur la restriction. Soit $Y = \bigoplus_{i \in I} Y_i$
une somme directe. Soient $a : X \to Y$ avec $X \in \Ob(\mathcal{D}_c)$
et $g \in G(X)$ représentant un élément $\xi$ de $G'(Y)$.
Le morphisme $a : X \to Y$ s'écrit de façon unique comme une somme de morphismes
$a_i : X \to Y_i$ presque tous nuls, puisque $X$ est un objet compact de $\mathcal{D}$.
Soit $I' = \{i \in I \mid a_i \not = 0\}$. On peut alors factoriser
$a$ comme le composé
$$
X \xrightarrow{(1, \ldots, 1)}
\bigoplus\nolimits_{i \in I'} X
\xrightarrow{\bigoplus_{i \in I'} a_i}
\bigoplus\nolimits_{i \in I} Y_i = Y
$$
On en déduit que $\xi = \sum_{i \in I'} \xi_i$
est la somme des images des éléments
$\xi_i \in G'(Y_i)$ correspondant à $a_i : X \to Y_i$
et $g \in G(X)$. Ainsi, $\bigoplus G'(Y_i) \to G'(Y)$
est surjective. Nous omettons la vérification (triviale) de son injectivité.

\medskip\noindent
Il s'ensuit que le foncteur $Y \mapsto G'(Y)^\vee$ est cohomologique
et transforme les sommes directes en produits directs. La représentabilité de Brown,
voir Catégories dérivées, proposition \ref{derived-proposition-brown},
montre donc qu'il existe $Y \in \Ob(\mathcal{D})$
et un isomorphisme
$G'(Z)^\vee = \Hom(Z, Y)$ fonctoriel en $Z$.
Pour $X \in \Ob(\mathcal{D}_c)$, on a
$G'(X)^\vee = G(X)^\vee = (H(X)^\vee)^\vee = H(X)$
car $\dim_k H(X) < \infty$, ce qui achève la démonstration.
\end{proof}

\begin{theorem}
\label{equiv-theorem-bondal-van-den-bergh}
\begin{reference}
Dans le cas projectif, il s'agit de \cite[Théorème A.1]{BvdB}
\end{reference}
Soit $X$ un schéma propre sur un corps $k$.
Soit $F : D_{perf}(\mathcal{O}_X)^{opp} \to \text{Vect}_k$
un foncteur cohomologique $k$-linéaire tel que
$$
\sum\nolimits_{n \in \mathbf{Z}} \dim_k F(E[n]) < \infty
$$
pour tout $E \in D_{perf}(\mathcal{O}_X)$. Alors $F$ est isomorphe à un foncteur
de la forme $E \mapsto \Hom_X(E, K)$ pour un certain
$K \in D^b_{\textit{Coh}}(\mathcal{O}_X)$.
\end{theorem}

\begin{proof}
La catégorie dérivée $D_\QCoh(\mathcal{O}_X)$ admet les sommes directes,
est compactement engendrée, et $D_{perf}(\mathcal{O}_X)$ est la sous-catégorie pleine
des objets compacts, voir
Catégories dérivées des schémas, lemme
\ref{perfect-lemma-quasi-coherence-direct-sums},
théorème \ref{perfect-theorem-bondal-van-den-Bergh}, et
proposition \ref{perfect-proposition-compact-is-perfect}.
Par le lemme \ref{equiv-lemma-van-den-bergh}, nous pouvons supposer que
$F(E) = \Hom_X(E, K)$ pour un certain $K \in \Ob(D_\QCoh(\mathcal{O}_X))$.
Il s'ensuit alors que $K$ appartient à $D^b_{\textit{Coh}}(\mathcal{O}_X)$
par le lemme \ref{equiv-lemma-characterize-dbcoh-projective}.
\end{proof}

\begin{lemma}
\label{equiv-lemma-homological-representable}
Soit $X$ un schéma propre et régulier sur un corps $k$.
Soit $G : D_{perf}(\mathcal{O}_X) \to \text{Vect}_k$
un foncteur homologique $k$-linéaire tel que
$$
\sum\nolimits_{n \in \mathbf{Z}} \dim_k G(E[n]) < \infty
$$
pour tout $E \in D_{perf}(\mathcal{O}_X)$. Alors $G$ est isomorphe à un foncteur
de la forme $E \mapsto \Hom_X(K, E)$ pour un certain $K \in D_{perf}(\mathcal{O}_X)$.
\end{lemma}

\begin{proof}
Considérons le foncteur contravariant $E \mapsto E^\vee$
sur $D_{perf}(\mathcal{O}_X)$, voir
Cohomologie, lemme \ref{cohomology-lemma-dual-perfect-complex}.
Ce foncteur est une anti-auto-équivalence exacte de $D_{perf}(\mathcal{O}_X)$.
Nous pouvons donc appliquer le théorème \ref{equiv-theorem-bondal-van-den-bergh}
au foncteur $F(E) = G(E^\vee)$ pour trouver
$K \in D_{perf}(\mathcal{O}_X)$ tel que $G(E^\vee) = \Hom_X(E, K)$.
Il s'ensuit que $G(E) = \Hom_X(E^\vee, K) = \Hom_X(K^\vee, E)$,
et l'on conclut que $K^\vee$ convient.
\end{proof}






\section{Existence d'adjoints}
\label{equiv-section-adjoints}

\noindent
Les résultats de l'article de Bondal et van den Bergh ont pour conséquence
l'énoncé suivant sur l'existence automatique d'adjoints.

\begin{lemma}
\label{equiv-lemma-always-right-adjoints}
Soit $k$ un corps. Soient $X$ et $Y$ des schémas propres sur $k$.
Si $X$ est régulier, tout foncteur exact $k$-linéaire
$F : D_{perf}(\mathcal{O}_X) \to D_{perf}(\mathcal{O}_Y)$
admet un adjoint à droite exact et un adjoint à gauche exact.
\end{lemma}

\begin{proof}
S'il existe, un adjoint est un foncteur exact d'après le lemme très général
\ref{derived-lemma-adjoint-is-exact} de Catégories dérivées.

\medskip\noindent
Démontrons l'existence d'un adjoint à droite.
Pour cela, il suffit de montrer que, pour
$M \in D_{perf}(\mathcal{O}_Y)$, le foncteur contravariant
$K \mapsto \Hom_Y(F(K), M)$ est représentable.
Ce foncteur est contravariant, $k$-linéaire et cohomologique.
Par le théorème \ref{equiv-theorem-bondal-van-den-bergh},
il suffit donc de montrer que
$$
\sum\nolimits_{i \in \mathbf{Z}} \dim_k \Ext^i_Y(F(K), M) < \infty
$$
Cela résulte du lemme \ref{equiv-lemma-finiteness}.

\medskip\noindent
Pour l'existence de l'adjoint à gauche, on raisonne de même
en utilisant le lemme \ref{equiv-lemma-homological-representable}
au lieu du théorème \ref{equiv-theorem-bondal-van-den-bergh}.
\end{proof}






\section{Foncteurs de Fourier-Mukai}
\label{equiv-section-fourier-mukai}

\noindent
Ces foncteurs ont été introduits pour la première fois dans \cite{Mukai}.

\begin{definition}
\label{equiv-definition-fourier-mukai-functor}
Soit $S$ un schéma. Soient $X$ et $Y$ des schémas sur $S$.
Soit $K \in D(\mathcal{O}_{X \times_S Y})$. Le foncteur exact
$$
\Phi_K : D(\mathcal{O}_X) \longrightarrow D(\mathcal{O}_Y),\quad
M \longmapsto R\text{pr}_{2, *}(
L\text{pr}_1^*M \otimes_{\mathcal{O}_{X \times_S Y}}^\mathbf{L} K)
$$
de catégories triangulées est appelé un {\it foncteur de Fourier-Mukai},
et $K$ est appelé un {\it noyau de Fourier-Mukai} pour ce foncteur.
De plus,
\begin{enumerate}
\item si $\Phi_K$ envoie $D_\QCoh(\mathcal{O}_X)$ dans $D_\QCoh(\mathcal{O}_Y)$,
alors le foncteur exact obtenu
$\Phi_K : D_\QCoh(\mathcal{O}_X) \to D_\QCoh(\mathcal{O}_Y)$
est appelé un foncteur de Fourier-Mukai,
\item si $\Phi_K$ envoie $D_{perf}(\mathcal{O}_X)$ dans
$D_{perf}(\mathcal{O}_Y)$, alors le foncteur exact obtenu
$\Phi_K : D_{perf}(\mathcal{O}_X) \to D_{perf}(\mathcal{O}_Y)$
est appelé un foncteur de Fourier-Mukai, et
\item si $X$ et $Y$ sont noethériens et si $\Phi_K$ envoie
$D^b_{\textit{Coh}}(\mathcal{O}_X)$ dans $D^b_{\textit{Coh}}(\mathcal{O}_Y)$,
alors le foncteur exact obtenu
$\Phi_K : D^b_{\textit{Coh}}(\mathcal{O}_X) \to
D^b_{\textit{Coh}}(\mathcal{O}_Y)$
est appelé un foncteur de Fourier-Mukai.
De même pour $D_{\textit{Coh}}$, $D^+_{\textit{Coh}}$, $D^-_{\textit{Coh}}$.
\end{enumerate}
\end{definition}

\begin{lemma}
\label{equiv-lemma-fourier-Mukai-QCoh}
Soit $S$ un schéma. Soient $X$ et $Y$ des schémas sur $S$.
Soit $K \in D(\mathcal{O}_{X \times_S Y})$.
Le foncteur de Fourier-Mukai correspondant $\Phi_K$ envoie
$D_\QCoh(\mathcal{O}_X)$ dans $D_\QCoh(\mathcal{O}_Y)$
si $K$ appartient à $D_\QCoh(\mathcal{O}_{X \times_S Y})$ et si $X \to S$ est
quasi-compact et quasi-séparé.
\end{lemma}

\begin{proof}
Cela résulte du fait que l'image inverse dérivée préserve
$D_\QCoh$
(Catégories dérivées des schémas, lemme
\ref{perfect-lemma-quasi-coherence-pullback}),
que les produits tensoriels dérivés préservent $D_\QCoh$
(Catégories dérivées des schémas, lemme
\ref{perfect-lemma-quasi-coherence-tensor-product}),
que la projection $\text{pr}_2 : X \times_S Y \to Y$ est
quasi-compacte et quasi-séparée
(Schémas, lemmes
\ref{schemes-lemma-quasi-compact-preserved-base-change} et
\ref{schemes-lemma-separated-permanence}), et
que l'image directe totale par un morphisme quasi-séparé et quasi-compact
préserve $D_\QCoh$
(Catégories dérivées des schémas, lemme
\ref{perfect-lemma-quasi-coherence-direct-image}).
\end{proof}

\begin{lemma}
\label{equiv-lemma-compose-fourier-mukai}
Soit $S$ un schéma. Soient $X, Y, Z$ des schémas sur $S$. Supposons que
$X \to S$, $Y \to S$ et $Z \to S$ soient quasi-compacts et quasi-séparés.
Soit $K \in D_\QCoh(\mathcal{O}_{X \times_S Y})$.
Soit $K' \in D_\QCoh(\mathcal{O}_{Y \times_S Z})$.
Considérons les foncteurs de Fourier-Mukai
$\Phi_K : D_\QCoh(\mathcal{O}_X) \to D_\QCoh(\mathcal{O}_Y)$
et $\Phi_{K'} : D_\QCoh(\mathcal{O}_Y) \to D_\QCoh(\mathcal{O}_Z)$.
Si $X$ et $Z$ sont tor-indépendants sur $S$ et si $Y \to S$ est plat,
alors
$$
\Phi_{K'} \circ \Phi_K = \Phi_{K''} :
D_\QCoh(\mathcal{O}_X)
\longrightarrow
D_\QCoh(\mathcal{O}_Z)
$$
où
$$
K'' = R\text{pr}_{13, *}(
L\text{pr}_{12}^*K
\otimes_{\mathcal{O}_{X \times_S Y \times_S Z}}^\mathbf{L}
L\text{pr}_{23}^*K')
$$
dans $D_\QCoh(\mathcal{O}_{X \times_S Z})$.
\end{lemma}

\begin{proof}
L'énoncé a un sens d'après le lemme \ref{equiv-lemma-fourier-Mukai-QCoh}.
Nous allons utiliser
Catégories dérivées des schémas, lemmes
\ref{perfect-lemma-quasi-coherence-pullback},
\ref{perfect-lemma-quasi-coherence-tensor-product}, et
\ref{perfect-lemma-quasi-coherence-direct-image},
ainsi que Schémas, lemmes
\ref{schemes-lemma-quasi-compact-preserved-base-change} et
\ref{schemes-lemma-separated-permanence},
sans autre mention.
D'après Catégories dérivées des schémas, lemme
\ref{perfect-lemma-flat-base-change-tor-independent},
$X \times_S Y$ et $Y \times_S Z$ sont tor-indépendants
sur $Y$. Autrement dit, le changement de base s'applique au diagramme cartésien
$$
\xymatrix{
X \times_S Y \times_S Z \ar[d] \ar[r] &
Y \times_S Z \ar[d]^{p^{YZ}_Y} \\
X \times_S Y \ar[r]^{p^{XY}_Y} & Y
}
$$
pour les complexes à faisceaux de cohomologie quasi-cohérents, voir
Catégories dérivées des schémas, lemme \ref{perfect-lemma-compare-base-change}.
En posant $p^* = Lp^*$, $p_* = Rp_*$ et $\otimes = \otimes^\mathbf{L}$,
on a, pour $M \in D_\QCoh(\mathcal{O}_X)$, la suite d'égalités
\begin{align*}
\Phi_{K'}(\Phi_K(M))
& =
p^{YZ}_{Z, *}(p^{YZ, *}_Y p^{XY}_{Y, *}(p^{XY, *}_X M \otimes K) \otimes K') \\
& =
p^{YZ}_{Z, *}(\text{pr}_{23, *} \text{pr}_{12}^*(p^{XY, *}_X M \otimes K)
\otimes K') \\
& =
p^{YZ}_{Z, *}(\text{pr}_{23, *}(\text{pr}_1^*M \otimes \text{pr}_{12}^*K)
\otimes K') \\
& =
p^{YZ}_{Z, *}(\text{pr}_{23, *}(\text{pr}_1^*M \otimes \text{pr}_{12}^*K
\otimes \text{pr}_{23}^*K')) \\
& =
\text{pr}_{3, *}(\text{pr}_1^*M \otimes \text{pr}_{12}^*K
\otimes \text{pr}_{23}^*K') \\
& =
p^{XZ}_{Z, *}\text{pr}_{13, *}(\text{pr}_1^*M \otimes \text{pr}_{12}^*K
\otimes \text{pr}_{23}^*K') \\
& =
p^{XZ}_{Z, *} (p^{XZ, *}_X M \otimes \text{pr}_{13, *}(\text{pr}_{12}^*K
\otimes \text{pr}_{23}^*K'))
\end{align*}
comme voulu. Nous avons utilisé ici la remarque sur le changement de base dans la
deuxième égalité et Catégories dérivées des schémas, lemme
\ref{perfect-lemma-cohomology-base-change}, dans la $4$e et la
dernière égalité.
\end{proof}

\begin{lemma}
\label{equiv-lemma-fourier-mukai}
Soit $S$ un schéma. Soient $X$ et $Y$ des schémas sur $S$.
Soit $K \in D(\mathcal{O}_{X \times_S Y})$.
Le foncteur de Fourier-Mukai correspondant $\Phi_K$ envoie
$D_{perf}(\mathcal{O}_X)$ dans $D_{perf}(\mathcal{O}_Y)$ si au moins
l'une des conditions suivantes est satisfaite :
\begin{enumerate}
\item $S$ est noethérien, $X \to S$ et $Y \to S$ sont de type fini,
$K \in D^b_{\textit{Coh}}(\mathcal{O}_{X \times_S Y})$, le support de $H^i(K)$
est propre sur $Y$ pour tout $i$, et $K$ est de tor-dimension finie
comme objet de $D(\text{pr}_2^{-1}\mathcal{O}_Y)$,
\item $X \to S$ est de présentation finie et $K$ peut être représenté
par un complexe borné $\mathcal{K}^\bullet$ de modules de présentation finie sur
$\mathcal{O}_{X \times_S Y}$, plat sur $Y$, et dont le support
est propre sur $Y$,
\item $X \to S$ est un morphisme propre, plat et de présentation finie,
et $K$ est parfait,
\item $S$ est noethérien, $X \to S$ est plat et propre, et $K$ est parfait,
\item $X \to S$ est un morphisme propre, plat et de présentation finie,
et $K$ est $Y$-parfait,
\item $S$ est noethérien, $X \to S$ est plat et propre, et $K$ est
$Y$-parfait.
\end{enumerate}
\end{lemma}

\begin{proof}
Si $M$ est parfait sur $X$, alors $L\text{pr}_1^*M$
est parfait sur $X \times_S Y$, voir
Cohomologie, lemme \ref{cohomology-lemma-perfect-pullback}.
Nous utiliserons ce fait sans autre mention ci-dessous.
Nous utiliserons également le fait que, si $X \to S$ est de type fini, propre,
plat ou de présentation finie, il en va de même pour
le changement de base $\text{pr}_2 : X \times_S Y \to Y$, voir
Morphismes, lemmes
\ref{morphisms-lemma-base-change-finite-type},
\ref{morphisms-lemma-base-change-proper},
\ref{morphisms-lemma-base-change-flat}, et
\ref{morphisms-lemma-base-change-finite-presentation}.

\medskip\noindent
L'assertion (1) résulte de
Catégories dérivées des schémas, lemme \ref{perfect-lemma-perfect-direct-image},
combiné avec
Catégories dérivées des schémas, lemme \ref{perfect-lemma-perfect-on-noetherian}.

\medskip\noindent
L'assertion (2) résulte de
Catégories dérivées des schémas, lemme
\ref{perfect-lemma-base-change-tensor-perfect}.

\medskip\noindent
L'assertion (3) résulte de
Catégories dérivées des schémas, lemme
\ref{perfect-lemma-flat-proper-perfect-direct-image-general}.

\medskip\noindent
L'assertion (4) résulte de l'assertion (3) et du fait qu'un morphisme de type fini
entre schémas noethériens est de présentation finie d'après Morphismes, lemme
\ref{morphisms-lemma-noetherian-finite-type-finite-presentation}.

\medskip\noindent
L'assertion (5) résulte de
Catégories dérivées des schémas, lemme
\ref{perfect-lemma-derived-pushforward-rel-perfect}, combiné avec
Catégories dérivées des schémas, lemme
\ref{perfect-lemma-perfect-relatively-perfect}.

\medskip\noindent
L'assertion (6) résulte de l'assertion (5) de la même manière que l'assertion (4) résulte de
l'assertion (3).
\end{proof}

\begin{lemma}
\label{equiv-lemma-fourier-mukai-Coh}
Soit $S$ un schéma noethérien. Soient $X$ et $Y$ des schémas de type fini
sur $S$. Soit $K \in D^b_{\textit{Coh}}(\mathcal{O}_{X \times_S Y})$.
Le foncteur de Fourier-Mukai correspondant $\Phi_K$ envoie
$D^b_{\textit{Coh}}(\mathcal{O}_X)$ dans $D^b_{\textit{Coh}}(\mathcal{O}_Y)$
si au moins l'une des conditions suivantes est satisfaite :
\begin{enumerate}
\item le support de $H^i(K)$ est propre sur $Y$ pour tout $i$, et $K$
est de tor-dimension finie comme objet de $D(\text{pr}_1^{-1}\mathcal{O}_X)$,
\item $K$ peut être représenté par un complexe borné $\mathcal{K}^\bullet$
de $\mathcal{O}_{X \times_S Y}$-modules cohérents, plat sur $X$, et dont le support
est propre sur $Y$,
\item le support de $H^i(K)$ est propre sur $Y$ pour tout $i$
et $X$ est un schéma régulier,
\item $K$ est parfait, le support de $H^i(K)$ est propre sur $Y$ pour tout $i$,
et $Y \to S$ est plat.
\end{enumerate}
En outre, dans chacun des cas, la condition sur le support est automatique
si $X \to S$ est propre.
\end{lemma}

\begin{proof}
Soit $M$ un objet de $D^b_{\textit{Coh}}(\mathcal{O}_X)$.
Dans chaque cas, nous utiliserons Catégories dérivées des schémas, lemme
\ref{perfect-lemma-direct-image-coherent}, pour montrer que
$$
\Phi_K(M) = R\text{pr}_{2, *}(
L\text{pr}_1^*M
\otimes_{\mathcal{O}_{X \times_S Y}}^\mathbf{L}
K)
$$
appartient à $D^b_{\textit{Coh}}(\mathcal{O}_Y)$. Le produit tensoriel dérivé
$L\text{pr}_1^*M \otimes_{\mathcal{O}_{X \times_S Y}}^\mathbf{L} K$
est un objet pseudo-cohérent de $D(\mathcal{O}_{X \times_S Y})$
(d'après
Cohomologie, lemme \ref{cohomology-lemma-pseudo-coherent-pullback},
Catégories dérivées des schémas, lemme
\ref{perfect-lemma-identify-pseudo-coherent-noetherian}, et
Cohomologie, lemme \ref{cohomology-lemma-tensor-pseudo-coherent}),
et possède donc des faisceaux de cohomologie cohérents (à nouveau d'après
Catégories dérivées des schémas, lemme
\ref{perfect-lemma-identify-pseudo-coherent-noetherian}).
Dans chaque cas, les supports des faisceaux de cohomologie
$H^i(L\text{pr}_1^*M \otimes_{\mathcal{O}_{X \times_S Y}}^\mathbf{L} K)$
sont propres sur $Y$, car ces supports sont contenus dans la
réunion des supports des $H^i(K)$. Il suffit donc, dans chaque cas,
de prouver que ce produit tensoriel est borné inférieurement.

\medskip\noindent
Cas (1). D'après Cohomologie, lemme \ref{cohomology-lemma-variant-derived-pullback},
on a
$$
L\text{pr}_1^*M
\otimes_{\mathcal{O}_{X \times_S Y}}^\mathbf{L}
K
\cong
\text{pr}_1^{-1}M
\otimes_{\text{pr}_1^{-1}\mathcal{O}_X}^\mathbf{L}
K
$$
avec les notations évidentes. L'hypothèse de tor-dimension finie
et le fait que $M$ ne possède qu'un nombre fini de
faisceaux de cohomologie non nuls impliquent donc la borne voulue.

\medskip\noindent
Le cas (2) résulte du fait qu'ici l'hypothèse implique que $K$ est de
tor-dimension finie comme objet de $D(\text{pr}_1^{-1}\mathcal{O}_X)$ ;
l'argument du paragraphe précédent s'applique donc.

\medskip\noindent
Dans le cas (3), $K$ est également de tor-dimension finie
comme objet de $D(\text{pr}_1^{-1}\mathcal{O}_X)$. En effet, choisissons
des ouverts affines $U = \Spec(A)$ et $V = \Spec(B)$, respectivement de $X$ et $Y$, dont les images sont contenues dans
l'ouvert affine $W = \Spec(R)$ de $S$. Alors
$K|_{U \times V}$ est donné par un complexe borné de
$A \otimes_R B$-modules de type fini $M^\bullet$. Comme $A$ est un anneau régulier
de dimension finie, chaque $M^i$ est de dimension projective finie
comme $A$-module (Algèbre, lemme
\ref{algebra-lemma-finite-gl-dim-finite-dim-regular})
et donc de tor-dimension finie comme $A$-module.
Ainsi, $M^\bullet$ est de tor-dimension finie comme complexe de $A$-modules
(Compléments sur l'algèbre, lemme
\ref{more-algebra-lemma-complex-finite-tor-dimension-modules}).
Puisque $X \times Y$ est quasi-compact, il existe $[a, b]$
tel que, pour tout point $z \in X \times Y$, le germe $K_z$
soit d'amplitude de tor contenue dans $[a, b]$ sur $\mathcal{O}_{X, \text{pr}_1(z)}$.
Cela implique que $K$ est de tor-dimension bornée comme objet de
$D(\text{pr}_1^{-1}\mathcal{O}_X)$, voir
Cohomologie, lemme \ref{cohomology-lemma-tor-amplitude-stalk}.
On conclut comme dans les deux paragraphes précédents.

\medskip\noindent
Cas (4). Avec les notations ci-dessus, l'homomorphisme d'anneaux $R \to B$ est plat.
L'homomorphisme d'anneaux $A \to A \otimes_R B$ est donc plat. Tout module projectif
sur $A \otimes_R B$ est donc plat sur $A$. Ainsi, tout complexe parfait de
$A \otimes_R B$-modules est de tor-dimension finie comme complexe
de $A$-modules, et l'on conclut comme précédemment.
\end{proof}

\begin{example}
\label{equiv-example-diagonal-fourier-mukai}
Soit $X \to S$ un morphisme séparé de schémas. Alors la diagonale
$\Delta : X \to X \times_S X$ est une immersion fermée et, par conséquent,
$\mathcal{O}_\Delta = \Delta_*\mathcal{O}_X = R\Delta_*\mathcal{O}_X$
est un $\mathcal{O}_{X \times_S X}$-module quasi-cohérent de type fini
qui est plat sur $X$ (pour l'une ou l'autre projection). Le foncteur de Fourier-Mukai
$\Phi_{\mathcal{O}_\Delta}$ est égal à l'identité dans ce cas.
En effet, pour tout $M \in D(\mathcal{O}_X)$, on a
\begin{align*}
L\text{pr}_1^*M \otimes_{\mathcal{O}_{X \times_S X}}^\mathbf{L}
\mathcal{O}_\Delta
& =
L\text{pr}_1^*M \otimes_{\mathcal{O}_{X \times_S X}}^\mathbf{L}
R\Delta_*\mathcal{O}_X \\
& =
R\Delta_*(
L\Delta^*L\text{pr}_1^*M \otimes_{\mathcal{O}_X}^\mathbf{L} \mathcal{O}_X) \\
& =
R\Delta_*(M)
\end{align*}
Nous avons examiné la première égalité ci-dessus.
La deuxième égalité est donnée par Cohomologie, lemme
\ref{cohomology-lemma-projection-formula-closed-immersion}.
La troisième vient de ce que $\text{pr}_1 \circ \Delta = \text{id}_X$ et du
lemme \ref{cohomology-lemma-derived-pullback-composition} de Cohomologie.
Pour obtenir l'image directe vers $X$, appliquons à ce qui précède le foncteur $R\text{pr}_{2, *}$ ;
on obtient $M$ d'après
Cohomologie, lemme \ref{cohomology-lemma-derived-pushforward-composition},
et le fait que $\text{pr}_2 \circ \Delta = \text{id}_X$.
\end{example}

\begin{lemma}
\label{equiv-lemma-fourier-mukai-right-adjoint}
\begin{reference}
Comparer avec la discussion de \cite{Rizzardo}.
\end{reference}
Soient $X \to S$ et $Y \to S$ des morphismes de schémas quasi-compacts et quasi-séparés.
Soit $\Phi : D_\QCoh(\mathcal{O}_X) \to D_\QCoh(\mathcal{O}_Y)$
un foncteur de Fourier-Mukai à noyau pseudo-cohérent
$K \in D_\QCoh(\mathcal{O}_{X \times_S Y})$.
Soit $a : D_\QCoh(\mathcal{O}_Y) \to  D_\QCoh(\mathcal{O}_{X \times_S Y})$
l'adjoint à droite de $R\text{pr}_{2, *}$, voir
Dualité pour les schémas, lemme \ref{duality-lemma-twisted-inverse-image}.
Notons
$$
K' = (Y \times_S X \to X \times_S Y)^*
R\SheafHom_{\mathcal{O}_{X \times_S Y}}(K, a(\mathcal{O}_Y)) \in
D_\QCoh(\mathcal{O}_{Y \times_S X})
$$
et notons $\Phi' : D_\QCoh(\mathcal{O}_Y) \to D_\QCoh(\mathcal{O}_X)$
la transformée de Fourier-Mukai correspondante. Il existe une application canonique
$$
\Hom_X(M, \Phi'(N)) \longrightarrow \Hom_Y(\Phi(M), N)
$$
fonctorielle en $M$ dans $D_\QCoh(\mathcal{O}_X)$ et en $N$ dans
$D_\QCoh(\mathcal{O}_Y)$, qui est un isomorphisme si
\begin{enumerate}
\item $N$ est parfait, ou
\item $K$ est parfait et $X \to S$ est propre, plat et de présentation finie.
\end{enumerate}
\end{lemma}

\begin{proof}
Le lemme \ref{equiv-lemma-fourier-Mukai-QCoh} donne un foncteur $\Phi$
comme dans l'énoncé. Observons que $a(\mathcal{O}_Y)$ appartient à
$D^+_\QCoh(\mathcal{O}_{X \times_S Y})$ d'après Dualité pour les schémas,
lemme \ref{duality-lemma-twisted-inverse-image-bounded-below}.
Ainsi, pour $K$ pseudo-cohérent, on a
$K' \in D_\QCoh(\mathcal{O}_{Y \times_S X})$
d'après Catégories dérivées des schémas, lemme
\ref{perfect-lemma-quasi-coherence-internal-hom},
et l'on obtient $\Phi'$ comme indiqué.

\medskip\noindent
Nous abrégeons
$\otimes^\mathbf{L} = \otimes_{\mathcal{O}_{X \times_S Y}}^\mathbf{L}$
et
$\SheafHom = R\SheafHom_{\mathcal{O}_{X \times_S Y}}$.
Soit $M$ dans $D_\QCoh(\mathcal{O}_X)$ et soit
$N$ dans $D_\QCoh(\mathcal{O}_Y)$. On a
\begin{align*}
\Hom_Y(\Phi(M), N)
& =
\Hom_Y(R\text{pr}_{2, *}(L\text{pr}_1^*M \otimes^\mathbf{L} K), N) \\
& =
\Hom_{X \times_S Y}(L\text{pr}_1^*M \otimes^\mathbf{L} K, a(N)) \\
& =
\Hom_{X \times_S Y}(L\text{pr}_1^*M,
R\SheafHom(K, a(N))) \\
& =
\Hom_X(M, R\text{pr}_{1, *}R\SheafHom(K, a(N)))
\end{align*}
où nous avons utilisé Cohomologie, lemmes \ref{cohomology-lemma-internal-hom}
et \ref{cohomology-lemma-adjoint}. Il existe des applications canoniques
$$
L\text{pr}_2^*N \otimes^\mathbf{L} R\SheafHom(K, a(\mathcal{O}_Y))
\xrightarrow{\alpha}
R\SheafHom(K, L\text{pr}_2^*N \otimes^\mathbf{L} a(\mathcal{O}_Y))
\xrightarrow{\beta}
R\SheafHom(K, a(N))
$$
Ici, $\alpha$ est l'application de
Cohomologie, lemme \ref{cohomology-lemma-internal-hom-diagonal-better},
et $\beta$ celle de Dualité pour les schémas, équation
(\ref{duality-equation-compare-with-pullback}).
En composant toutes ces flèches, on obtient la flèche fonctorielle affichée
dans l'énoncé du lemme.

\medskip\noindent
La flèche $\alpha$ est un isomorphisme d'après
Catégories dérivées des schémas, lemme
\ref{perfect-lemma-internal-hom-evaluate-tensor-isomorphism},
dès que $K$ ou $N$ est parfait.
La flèche $\beta$ est un isomorphisme si $N$ est parfait, d'après
Dualité pour les schémas, lemme \ref{duality-lemma-compare-with-pullback-perfect},
ou, en général, si $X \to S$ est
plat, propre et de présentation finie, d'après
Dualité pour les schémas, lemme
\ref{duality-lemma-compare-with-pullback-flat-proper}.
\end{proof}

\begin{lemma}
\label{equiv-lemma-fourier-mukai-left-adjoint}
\begin{reference}
Comparer avec la discussion de \cite{Rizzardo}.
\end{reference}
Soit $S$ un schéma noethérien. Soit $Y \to S$ un morphisme plat, propre
et de Gorenstein, et soit $X \to S$ un morphisme de type fini.
Notons $\omega^\bullet_{Y/S}$ le complexe dualisant relatif de
$Y$ sur $S$. Soit $\Phi : D_\QCoh(\mathcal{O}_X) \to D_\QCoh(\mathcal{O}_Y)$
un foncteur de Fourier-Mukai à noyau parfait
$K \in D_\QCoh(\mathcal{O}_{X \times_S Y})$. Notons
$$
K' = (Y \times_S X \to X \times_S Y)^*(K^\vee
\otimes_{\mathcal{O}_{X \times_S Y}}^\mathbf{L}
L\text{pr}_2^*\omega^\bullet_{Y/S})
\in
D_\QCoh(\mathcal{O}_{Y \times_S X})
$$
et notons $\Phi' : D_\QCoh(\mathcal{O}_Y) \to D_\QCoh(\mathcal{O}_X)$
la transformée de Fourier-Mukai correspondante. Il existe un isomorphisme
canonique
$$
\Hom_Y(N, \Phi(M)) \longrightarrow \Hom_X(\Phi'(N), M)
$$
fonctoriel en $M$ dans $D_\QCoh(\mathcal{O}_X)$ et en $N$ dans
$D_\QCoh(\mathcal{O}_Y)$.
\end{lemma}

\begin{proof}
Le lemme \ref{equiv-lemma-fourier-Mukai-QCoh} donne un foncteur $\Phi$
comme dans l'énoncé.

\medskip\noindent
Observons que la formation du complexe dualisant relatif commute
au changement de base dans notre situation, voir Dualité pour les schémas,
remarque \ref{duality-remark-relative-dualizing-complex}.
Ainsi, $L\text{pr}_2^*\omega^\bullet_{Y/S} = \omega^\bullet_{X \times_S Y/X}$.
De plus, $\omega^\bullet_{Y/S}$ est un objet
inversible de la catégorie dérivée, voir Dualité pour les schémas, lemme
\ref{duality-lemma-affine-flat-Noetherian-gorenstein}, et il est a fortiori
parfait.

\medskip\noindent
Pour démontrer effectivement le lemme, nous allons recourir à un artifice. Nous allons
montrer qu'en échangeant les rôles de $X$ et $Y$, ainsi que ceux de $K$ et $K'$,
on se trouve dans la situation du lemme \ref{equiv-lemma-fourier-mukai-right-adjoint},
ce qui donne le résultat. Il est clair que $K'$ est parfait comme
produit tensoriel d'objets parfaits, de sorte que la discussion du
lemme \ref{equiv-lemma-fourier-mukai-right-adjoint} s'y applique.
Pour montrer que la procédure du
lemme \ref{equiv-lemma-fourier-mukai-right-adjoint}, appliquée à $K'$
sur $Y \times_S X$, produit un complexe isomorphe à $K$, il suffit
(détails omis) de montrer que
$$
R\SheafHom(R\SheafHom(K, \omega^\bullet_{X \times_S Y/X}),
\omega^\bullet_{X \times_S Y/X}) = K
$$
C'est clair, car $K$ est parfait et $\omega^\bullet_{X \times_S Y/X}$
est inversible ; nous omettons les détails. Ainsi, le
lemme \ref{equiv-lemma-fourier-mukai-right-adjoint} fournit une application
$$
\Hom_Y(N, \Phi(M)) \longrightarrow \Hom_X(\Phi'(N), M)
$$
fonctorielle en $M$ dans $D_\QCoh(\mathcal{O}_X)$ et en $N$ dans
$D_\QCoh(\mathcal{O}_Y)$, qui est un isomorphisme puisque
$K'$ est parfait. Cela achève la démonstration.
\end{proof}

\begin{lemma}
\label{equiv-lemma-fourier-mukai-flat-proper-over-noetherian}
Soit $S$ un schéma noethérien.
\begin{enumerate}
\item Pour $X$ et $Y$ propres et plats sur $S$ et $K$ dans
$D_{perf}(\mathcal{O}_{X \times_S Y})$, on obtient un foncteur de Fourier-Mukai
$\Phi_K : D_{perf}(\mathcal{O}_X) \to D_{perf}(\mathcal{O}_Y)$.
\item Pour $X$, $Y$ et $Z$ propres et plats sur $S$, $K \in
D_{perf}(\mathcal{O}_{X \times_S Y})$, $K' \in
D_{perf}(\mathcal{O}_{Y \times_S Z})$, le composé
$\Phi_{K'} \circ \Phi_K : D_{perf}(\mathcal{O}_X) \to D_{perf}(\mathcal{O}_Z)$
est égal à $\Phi_{K''}$ avec $K'' \in D_{perf}(\mathcal{O}_{X \times_S Z})$,
calculé comme dans le lemme \ref{equiv-lemma-compose-fourier-mukai},
\item Pour $X$, $Y$, $K$ et $\Phi_K$ comme en (1), si $X \to S$ est de Gorenstein, alors
$\Phi_{K'} : D_{perf}(\mathcal{O}_Y) \to D_{perf}(\mathcal{O}_X)$ est un adjoint
à droite de $\Phi_K$, où $K' \in D_{perf}(\mathcal{O}_{Y \times_S X})$
est l'image inverse de $L\text{pr}_1^*\omega_{X/S}^\bullet
\otimes_{\mathcal{O}_{X \times_S Y}}^\mathbf{L} K^\vee$ par
$Y \times_S X \to X \times_S Y$.
\item Pour $X$, $Y$, $K$ et $\Phi_K$ comme en (1), si $Y \to S$ est de Gorenstein, alors
$\Phi_{K''} : D_{perf}(\mathcal{O}_Y) \to D_{perf}(\mathcal{O}_X)$ est un adjoint
à gauche de $\Phi_K$, où $K'' \in D_{perf}(\mathcal{O}_{Y \times_S X})$
est l'image inverse de $L\text{pr}_2^*\omega_{Y/S}^\bullet
\otimes_{\mathcal{O}_{X \times_S Y}}^\mathbf{L} K^\vee$ par
$Y \times_S X \to X \times_S Y$.
\end{enumerate}
\end{lemma}

\begin{proof}
L'assertion (1) est immédiate d'après le lemme \ref{equiv-lemma-fourier-mukai}, assertion (4).

\medskip\noindent
L'assertion (2) résulte du lemme \ref{equiv-lemma-compose-fourier-mukai} et du
fait que
$K'' = R\text{pr}_{13, *}(
L\text{pr}_{12}^*K
\otimes_{\mathcal{O}_{X \times_S Y \times_S Z}}^\mathbf{L}
L\text{pr}_{23}^*K')$ est parfait, par exemple d'après
Catégories dérivées des schémas, lemme
\ref{perfect-lemma-flat-proper-perfect-direct-image}.

\medskip\noindent
L'adjonction de l'assertion (3), sur tous les complexes à faisceaux de cohomologie
quasi-cohérents, résulte du lemme \ref{equiv-lemma-fourier-mukai-right-adjoint}, avec
$K'$ égal à l'image inverse de
$R\SheafHom_{\mathcal{O}_{X \times_S Y}}(K, a(\mathcal{O}_Y))$
par $Y \times_S X \to X \times_S Y$, où $a$ est l'adjoint à droite
de $R\text{pr}_{2, *} : D_\QCoh(\mathcal{O}_{X \times_S Y}) \to
D_\QCoh(\mathcal{O}_Y)$. Notons $f : X \to S$ le morphisme structural de $X$.
Puisque $f$ est propre, le foncteur
$f^! : D_\QCoh^+(\mathcal{O}_S) \to D_\QCoh^+(\mathcal{O}_X)$
est la restriction à $D_\QCoh^+(\mathcal{O}_S)$
de l'adjoint à droite de
$Rf_* : D_\QCoh(\mathcal{O}_X) \to D_\QCoh(\mathcal{O}_S)$, voir
Dualité pour les schémas, section \ref{duality-section-upper-shriek}.
Ainsi, le complexe dualisant relatif $\omega_{X/S}^\bullet$, tel qu'il est défini dans
Dualité pour les schémas, remarque
\ref{duality-remark-relative-dualizing-complex},
est égal à $\omega_{X/S}^\bullet = f^!\mathcal{O}_S$.
Comme la formation du complexe dualisant relatif
commute au changement de base (voir Dualité pour les schémas, remarque
\ref{duality-remark-relative-dualizing-complex}), on voit que
$a(\mathcal{O}_Y) = L\text{pr}_1^*\omega_{X/S}^\bullet$.
Ainsi
$$
R\SheafHom_{\mathcal{O}_{X \times_S Y}}(K, a(\mathcal{O}_Y))
\cong
L\text{pr}_1^*\omega_{X/S}^\bullet
\otimes_{\mathcal{O}_{X \times_S Y}}^\mathbf{L} K^\vee
$$
d'après Cohomologie, lemme \ref{cohomology-lemma-dual-perfect-complex}.
Enfin, comme $X \to S$ est supposé de Gorenstein, le complexe dualisant relatif
est inversible : cela résulte de Dualité pour les schémas, lemme
\ref{duality-lemma-affine-flat-Noetherian-gorenstein}.
On en déduit que $\omega_{X/S}^\bullet$ est parfait
(Cohomologie, lemme \ref{cohomology-lemma-invertible-derived})
et que $K'$ est donc parfait.
Par conséquent, $\Phi_{K'}$ envoie bien $D_{perf}(\mathcal{O}_Y)$
dans $D_{perf}(\mathcal{O}_X)$, ce qui achève la démonstration de (3).

\medskip\noindent
La démonstration de (4) est la même que celle de (3), si ce n'est qu'on utilise
le lemme \ref{equiv-lemma-fourier-mukai-left-adjoint} au lieu du
lemme \ref{equiv-lemma-fourier-mukai-right-adjoint}.
\end{proof}














\section{Résolutions et bornes}
\label{equiv-section-tricks-smooth}

\noindent
La diagonale d'un schéma lisse et propre possède une résolution commode.

\begin{lemma}
\label{equiv-lemma-on-product}
Soit $R$ un anneau noethérien. Soient $X$ et $Y$ des schémas de type fini sur $R$
qui possèdent la propriété de résolution. Pour tout
$\mathcal{O}_{X \times_R Y}$-module cohérent $\mathcal{F}$, il existe
une surjection $\mathcal{E} \boxtimes \mathcal{G} \to \mathcal{F}$,
où $\mathcal{E}$ est un $\mathcal{O}_X$-module localement libre de type fini
et $\mathcal{G}$ un $\mathcal{O}_Y$-module localement libre de type fini.
\end{lemma}

\begin{proof}
Soient $U \subset X$ et $V \subset Y$ des sous-schémas ouverts affines. Soit
$\mathcal{I} \subset \mathcal{O}_X$ le faisceau d'idéaux définissant la
structure réduite de sous-schéma fermé induite sur $X \setminus U$.
De même, soit $\mathcal{I}' \subset \mathcal{O}_Y$ le faisceau d'idéaux définissant la
structure réduite de sous-schéma fermé induite sur $Y \setminus V$.
Alors le faisceau d'idéaux
$$
\mathcal{J} = \Im(\text{pr}_1^*\mathcal{I} \otimes_{\mathcal{O}_{X \times_R Y}}
\text{pr}_2^*\mathcal{I}' \to \mathcal{O}_{X \times_R Y})
$$
vérifie $V(\mathcal{J}) = X \times_R Y \setminus U \times_R V$.
Pour toute section $s \in \mathcal{F}(U \times_R V)$, on peut trouver un entier
$n > 0$ et un morphisme $\mathcal{J}^n \to \mathcal{F}$ dont la restriction à
$U \times_R V$ donne $s$, voir
Cohomologie des schémas, lemme \ref{coherent-lemma-homs-over-open}.
Par hypothèse, on peut choisir des surjections
$\mathcal{E} \to \mathcal{I}$ et $\mathcal{G} \to \mathcal{I}'$.
Elles donnent les surjections correspondantes
$$
\mathcal{E} \boxtimes \mathcal{G} \to \mathcal{J}
\quad\text{et}\quad
\mathcal{E}^{\otimes n} \boxtimes \mathcal{G}^{\otimes n} \to \mathcal{J}^n
$$
et donc un morphisme
$\mathcal{E}^{\otimes n} \boxtimes \mathcal{G}^{\otimes n} \to \mathcal{F}$
dont l'image contient la section $s$ sur $U \times_R V$.
Comme on peut recouvrir $X \times_R Y$ par un nombre fini d'ouverts affines
de la forme $U \times_R V$ et que $\mathcal{F}|_{U \times_R V}$
est engendré par un nombre fini de sections (Propriétés, lemme
\ref{properties-lemma-finite-type-module}),
on en déduit qu'il existe une surjection
$$
\bigoplus\nolimits_{j = 1, \ldots, N}
\mathcal{E}_j^{\otimes n_j} \boxtimes \mathcal{G}_j^{\otimes n_j}
\to \mathcal{F}
$$
où $\mathcal{E}_j$ est localement libre de type fini sur $X$ et
$\mathcal{G}_j$ localement libre de type fini sur $Y$.
En posant $\mathcal{E} = \bigoplus \mathcal{E}_j^{\otimes n_j}$
et $\mathcal{G} = \bigoplus \mathcal{G}_j^{\otimes n_j}$,
on conclut que le lemme est vrai.
\end{proof}

\begin{lemma}
\label{equiv-lemma-on-product-general}
Soit $R$ un anneau. Soient $X$ et $Y$ des schémas quasi-compacts et quasi-séparés
sur $R$ qui possèdent la propriété de résolution. Pour tout
$\mathcal{O}_{X \times_R Y}$-module quasi-cohérent de type fini $\mathcal{F}$,
il existe une surjection $\mathcal{E} \boxtimes \mathcal{G} \to \mathcal{F}$,
où $\mathcal{E}$ est un $\mathcal{O}_X$-module localement libre de type fini
et $\mathcal{G}$ un $\mathcal{O}_Y$-module localement libre de type fini.
\end{lemma}

\begin{proof}
Cela résulte du lemme \ref{equiv-lemma-on-product} par un argument de passage à la limite.
Nous conseillons vivement au lecteur de ne pas lire la démonstration.
Comme $X \times_R Y$ est un sous-schéma fermé de $X \times_\mathbf{Z} Y$,
on peut sans inconvénient remplacer $R$ par $\mathbf{Z}$.
On peut écrire $\mathcal{F}$ comme quotient d'un
$\mathcal{O}_{X \times_R Y}$-module de présentation finie d'après
Propriétés, lemme
\ref{properties-lemma-finite-directed-colimit-surjective-maps}.
On peut donc supposer que $\mathcal{F}$ est de
présentation finie. Ensuite, on peut écrire $X = \lim X_i$,
où $X_i$ est de présentation finie sur $\mathbf{Z}$, et de même
$Y = \lim Y_j$, voir Limites, proposition \ref{limits-proposition-approximate}.
Alors $\mathcal{F}$ descend en $\mathcal{F}_{ij}$ sur un certain $X_i \times_R Y_j$
(Limites, lemme \ref{limits-lemma-descend-modules-finite-presentation}), tout comme
la propriété de résolution
(Catégories dérivées des schémas, lemme
\ref{perfect-lemma-resolution-property-descends}).
On applique alors le lemme \ref{equiv-lemma-on-product}
à $\mathcal{F}_{ij}$, puis on prend l'image inverse.
\end{proof}

\begin{lemma}
\label{equiv-lemma-diagonal-resolution}
Soit $R$ un anneau noethérien. Soit $X$ un schéma séparé de type fini
sur $R$ qui possède la propriété de résolution. Posons
$\mathcal{O}_\Delta = \Delta_*(\mathcal{O}_X)$, où
$\Delta : X \to X \times_R X$ est la diagonale de $X/k$.
Il existe une résolution
$$
\ldots \to
\mathcal{E}_2 \boxtimes \mathcal{G}_2 \to
\mathcal{E}_1 \boxtimes \mathcal{G}_1 \to
\mathcal{E}_0 \boxtimes \mathcal{G}_0 \to
\mathcal{O}_\Delta \to 0
$$
où chaque $\mathcal{E}_i$ et chaque $\mathcal{G}_i$ est un
$\mathcal{O}_X$-module localement libre de type fini.
\end{lemma}

\begin{proof}
Comme $X$ est séparé, le morphisme diagonal $\Delta$ est une immersion
fermée, et $\mathcal{O}_\Delta$ est donc un
$\mathcal{O}_{X \times_R X}$-module cohérent (Cohomologie des schémas, lemme
\ref{coherent-lemma-i-star-equivalence}).
Le lemme résulte donc immédiatement du lemme \ref{equiv-lemma-on-product}.
\end{proof}

\begin{lemma}
\label{equiv-lemma-Ext-0-regular}
Soit $X$ un schéma noethérien régulier de dimension $d < \infty$. Alors
\begin{enumerate}
\item pour $\mathcal{F}$, $\mathcal{G}$ qui sont des modules cohérents sur $\mathcal{O}_X$,
on a $\Ext^n_X(\mathcal{F}, \mathcal{G}) = 0$ pour $n > d$, et
\item pour $K, L \in D^b_{\textit{Coh}}(\mathcal{O}_X)$ et $a \in \mathbf{Z}$,
si $H^i(K) = 0$ pour $i < a + d$ et $H^i(L) = 0$ pour $i \geq a$, alors
$\Hom_X(K, L) = 0$.
\end{enumerate}
\end{lemma}

\begin{proof}
Pour démontrer (1), on utilise la suite spectrale
$$
H^p(X, \SheafExt^q(\mathcal{F}, \mathcal{G})) \Rightarrow
\Ext^{p + q}_X(\mathcal{F}, \mathcal{G})
$$
de Cohomologie, section \ref{cohomology-section-ext}. Soit $x \in X$.
On a
$$
\SheafExt^q(\mathcal{F}, \mathcal{G})_x =
\SheafExt^q_{\mathcal{O}_{X, x}}(\mathcal{F}_x, \mathcal{G}_x)
$$
voir Cohomologie, lemme \ref{cohomology-lemma-stalk-internal-hom}
(on utilise également ici le fait que $\mathcal{F}$ est pseudo-cohérent d'après
Catégories dérivées des schémas, lemme
\ref{perfect-lemma-identify-pseudo-coherent-noetherian}).
Posons $d_x = \dim(\mathcal{O}_{X, x})$.
Comme $\mathcal{O}_{X, x}$ est régulier, l'anneau
$\mathcal{O}_{X, x}$ est de dimension globale $d_x$, voir
Algèbre, proposition \ref{algebra-proposition-regular-finite-gl-dim}.
Ainsi, $\SheafExt^q_{\mathcal{O}_{X, x}}(\mathcal{F}_x, \mathcal{G}_x)$
est nul pour $q > d_x$. Il s'ensuit que les modules
$\SheafExt^q(\mathcal{F}, \mathcal{G})$ ont un support
de dimension au plus $d - q$. On a donc
$H^p(X, \SheafExt^q(\mathcal{F}, \mathcal{G})) = 0$ pour $p > d - q$
d'après Cohomologie, proposition \ref{cohomology-proposition-vanishing-Noetherian}.
Ceci prouve (1).

\medskip\noindent
Démonstration de (2).
On peut raisonner par récurrence sur le nombre de faisceaux de cohomologie non nuls
de $K$ et $L$. Le cas où ces nombres valent $0, 1$ résulte
de (1). Si le nombre de faisceaux de cohomologie non nuls de $K$
est $> 1$, soit $i \in \mathbf{Z}$ minimal tel que
$H^i(K)$ soit non nul. On obtient un triangle distingué
$$
H^i(K)[-i] \to K \to \tau_{\geq i + 1}K
$$
(Catégories dérivées, remarque
\ref{derived-remark-truncation-distinguished-triangle}),
et la nullité de $\Hom(K, L)$ résulte de celle
de $\Hom(H^i(K)[-i], L)$ et $\Hom(\tau_{\geq i + 1}K, L)$
d'après Catégories dérivées, lemme \ref{derived-lemma-representable-homological}.
De même si $L$ possède plus d'un faisceau de cohomologie non nul.
\end{proof}

\begin{lemma}
\label{equiv-lemma-split-complex-regular}
Soit $X$ un schéma noethérien régulier de dimension $d < \infty$.
Soient $K \in D^b_{\textit{Coh}}(\mathcal{O}_X)$ et $a \in \mathbf{Z}$.
Si $H^i(K) = 0$ pour $a < i < a + d$, alors
$K = \tau_{\leq a}K \oplus \tau_{\geq a + d}K$.
\end{lemma}

\begin{proof}
On a $\tau_{\leq a}K = \tau_{\leq a + d - 1}K$ par la nullité supposée
des faisceaux de cohomologie. D'après Catégories dérivées, remarque
\ref{derived-remark-truncation-distinguished-triangle},
on a un triangle distingué
$$
\tau_{\leq a}K \to K \to \tau_{\geq a + d}K \xrightarrow{\delta}
(\tau_{\leq a}K)[1]
$$
D'après Catégories dérivées, lemme \ref{derived-lemma-split}, il
suffit de montrer que le morphisme $\delta$ est nul.
Cela résulte du lemme \ref{equiv-lemma-Ext-0-regular}.
\end{proof}

\begin{lemma}
\label{equiv-lemma-diagonal-trick}
Soit $k$ un corps. Soit $X$ un schéma quasi-compact, séparé et lisse sur $k$.
Il existe des $\mathcal{O}_X$-modules localement libres de type fini
$\mathcal{E}$ et $\mathcal{G}$ tels que
$$
\mathcal{O}_\Delta \in \langle \mathcal{E} \boxtimes \mathcal{G} \rangle
$$
dans $D(\mathcal{O}_{X \times X})$, où la notation est celle de
Catégories dérivées, section \ref{derived-section-generators}.
\end{lemma}

\begin{proof}
Rappelons que $X$ est régulier d'après
Variétés, lemme \ref{varieties-lemma-smooth-regular}.
Ainsi, $X$ possède la propriété de résolution d'après
Catégories dérivées des schémas, lemme
\ref{perfect-lemma-regular-resolution-property}.
On peut donc choisir une résolution comme dans le lemme \ref{equiv-lemma-diagonal-resolution}.
Posons $\dim(X) = d$. Comme $X \times X$ est lisse sur $k$, il est régulier.
Ainsi, $X \times X$ est un schéma noethérien régulier et l'on a
$\dim(X \times X) = 2d$. L'objet
$$
K = (\mathcal{E}_{2d} \boxtimes \mathcal{G}_{2d} \to
\ldots \to
\mathcal{E}_0 \boxtimes \mathcal{G}_0)
$$
de $D_{perf}(\mathcal{O}_{X \times X})$
a pour faisceaux de cohomologie $\mathcal{O}_\Delta$
en degré $0$ et $\Ker(\mathcal{E}_{2d} \boxtimes \mathcal{G}_{2d} \to
\mathcal{E}_{2d-1} \boxtimes \mathcal{G}_{2d-1})$ en degré $-2d$, et zéro
à tous les autres degrés.
Le lemme \ref{equiv-lemma-split-complex-regular} montre donc que
$\mathcal{O}_\Delta$ est un facteur direct de $K$ dans
$D_{perf}(\mathcal{O}_{X \times X})$.
Manifestement, l'objet $K$ appartient à
$$
\left\langle
\bigoplus\nolimits_{i = 0, \ldots, 2d} \mathcal{E}_i \boxtimes \mathcal{G}_i
\right\rangle
\subset
\left\langle
\left(\bigoplus\nolimits_{i = 0, \ldots, 2d} \mathcal{E}_i\right)
\boxtimes
\left(\bigoplus\nolimits_{i = 0, \ldots, 2d} \mathcal{G}_i\right)
\right\rangle
$$
ce qui achève la démonstration. (Le lecteur peut consulter
Catégories dérivées, lemmes \ref{derived-lemma-generated-by-E-explicit} et
\ref{derived-lemma-in-cone-n}, pour voir que notre objet appartient à cette
catégorie.)
\end{proof}

\begin{lemma}
\label{equiv-lemma-smooth-proper-strong-generator}
Soit $k$ un corps. Soit $X$ un schéma propre et lisse sur $k$.
Alors $D_{perf}(\mathcal{O}_X)$
possède un générateur fort.
\end{lemma}

\begin{proof}
À l'aide du lemme \ref{equiv-lemma-diagonal-trick}, choisissons des
$\mathcal{O}_X$-modules localement libres de type fini $\mathcal{E}$ et $\mathcal{G}$ tels que
$\mathcal{O}_\Delta \in \langle \mathcal{E} \boxtimes \mathcal{G} \rangle$
dans $D(\mathcal{O}_{X \times X})$. Nous affirmons que $\mathcal{G}$
est un générateur fort de $D_{perf}(\mathcal{O}_X)$. Avec les notations de
Catégories dérivées, section \ref{derived-section-operate-on-full},
choisissons $m, n \geq 1$ tels que
$$
\mathcal{O}_\Delta \in
smd(add(\mathcal{E} \boxtimes \mathcal{G}[-m, m])^{\star n})
$$
C'est possible d'après Catégories dérivées, lemme
\ref{derived-lemma-find-smallest-containing-E}.
Soit $K$ un objet de $D_{perf}(\mathcal{O}_X)$. Comme
$L\text{pr}_1^*K \otimes_{\mathcal{O}_{X \times X}}^\mathbf{L} -$
est un foncteur exact et que
$$
L\text{pr}_1^*K \otimes_{\mathcal{O}_{X \times X}}^\mathbf{L}
(\mathcal{E} \boxtimes \mathcal{G}) =
(K \otimes_{\mathcal{O}_X}^\mathbf{L} \mathcal{E}) \boxtimes \mathcal{G}
$$
on déduit de
Catégories dérivées, remarque \ref{derived-remark-operations-functor}, que
$$
L\text{pr}_1^*K
\otimes_{\mathcal{O}_{X \times X}}^\mathbf{L}
\mathcal{O}_\Delta
\in
smd(add(
(K \otimes_{\mathcal{O}_X}^\mathbf{L} \mathcal{E})
\boxtimes \mathcal{G}[-m, m])^{\star n})
$$
En appliquant le foncteur exact $R\text{pr}_{2, *}$ et en observant que
$$
R\text{pr}_{2, *}
\left((K \otimes_{\mathcal{O}_X}^\mathbf{L} \mathcal{E}) \boxtimes
\mathcal{G}\right) =
R\Gamma(X, K \otimes_{\mathcal{O}_X}^\mathbf{L} \mathcal{E})
\otimes_k \mathcal{G}
$$
d'après Catégories dérivées des schémas, lemme
\ref{perfect-lemma-cohomology-base-change}, on en déduit que
$$
K = R\text{pr}_{2, *}(L\text{pr}_1^*K
\otimes_{\mathcal{O}_{X \times X}}^\mathbf{L} \mathcal{O}_\Delta)
\in
smd(add(R\Gamma(X, K \otimes_{\mathcal{O}_X}^\mathbf{L} \mathcal{E})
\otimes_k \mathcal{G}[-m, m])^{\star n})
$$
L'égalité résulte de la discussion de
l'exemple \ref{equiv-example-diagonal-fourier-mukai}.
Comme $K$ est parfait, il existe $a \leq b$ tels que
$H^i(X, K)$ ne soit non nul que pour $i \in [a, b]$. Comme $X$ est propre,
chaque $H^i(X, K)$ est de dimension finie. On en déduit que
le membre de droite est contenu dans
$smd(add(\mathcal{G}[-m + a, m + b])^{\star n})$, qui est
lui-même contenu dans $\langle \mathcal{G} \rangle_n$ d'après l'une des
références données ci-dessus. Cela achève la démonstration.
\end{proof}

\begin{lemma}
\label{equiv-lemma-diagonal-trick-proper}
Soit $k$ un corps. Soit $X$ un schéma propre et lisse sur $k$.
Il existe des entiers $m, n \geq 1$ et un
$\mathcal{O}_X$-module localement libre de type fini $\mathcal{G}$ tels que tout
$\mathcal{O}_X$-module cohérent appartienne à $smd(add(\mathcal{G}[-m, m])^{\star n})$,
avec les notations de Catégories dérivées, section
\ref{derived-section-operate-on-full}.
\end{lemma}

\begin{proof}
Dans la démonstration du lemme \ref{equiv-lemma-smooth-proper-strong-generator},
nous avons montré qu'il existe $m', n \geq 1$ tels que, pour tout
$\mathcal{O}_X$-module cohérent $\mathcal{F}$,
$$
\mathcal{F} \in smd(add(\mathcal{G}[-m' + a, m' + b])^{\star n})
$$
pour tous $a \leq b$ tels que $H^i(X, \mathcal{F})$ ne soit non nul que
pour $i \in [a, b]$. On peut donc prendre $a = 0$ et $b = \dim(X)$.
Prendre $m = \max(m', m' + b)$ achève la démonstration.
\end{proof}

\noindent
Le lemme suivant est le résultat de bornitude auquel renvoie
le titre de cette section.

\begin{lemma}
\label{equiv-lemma-boundedness}
Soit $k$ un corps. Soit $X$ un schéma lisse et propre sur $k$.
Soit $\mathcal{A}$ une catégorie abélienne. Soit
$H : D_{perf}(\mathcal{O}_X) \to \mathcal{A}$ un foncteur homologique
(Catégories dérivées, définition \ref{derived-definition-homological})
tel que, pour tout $K$ dans $D_{perf}(\mathcal{O}_X)$, l'objet
$H^i(K)$ ne soit non nul que pour un nombre fini de $i \in \mathbf{Z}$.
Alors il existe un entier $m \geq 1$ tel que
$H^i(\mathcal{F}) = 0$ pour tout $\mathcal{O}_X$-module cohérent
$\mathcal{F}$ et tout $i \not \in [-m, m]$.
De même pour les foncteurs cohomologiques.
\end{lemma}

\begin{proof}
Combiner le lemme \ref{equiv-lemma-diagonal-trick-proper} avec
Catégories dérivées, lemme \ref{derived-lemma-forward-cone-n}.
\end{proof}

\begin{lemma}
\label{equiv-lemma-bounded-fibres}
Soit $k$ un corps. Soient $X$ et $Y$ des schémas de type fini sur $k$.
Soit $K_0 \to K_1 \to K_2 \to \ldots$ un système d'objets
de $D_{perf}(\mathcal{O}_{X \times Y})$, et soit $m \geq 0$ un entier tel que
\begin{enumerate}
\item $H^q(K_i)$ ne soit non nul que pour $q \leq m$,
\item pour tout $\mathcal{O}_X$-module cohérent $\mathcal{F}$ tel que
$\dim(\text{Supp}(\mathcal{F})) = 0$, l'objet
$$
R\text{pr}_{2, *}(
\text{pr}_1^*\mathcal{F} \otimes_{\mathcal{O}_{X \times Y}}^\mathbf{L}
K_n)
$$
n'ait de faisceaux de cohomologie non nuls qu'aux degrés appartenant à
$[-m, m] \cup [-m - n, m - n]$ et, pour $n > 2m$, les morphismes de transition
induisent des isomorphismes sur les faisceaux de cohomologie aux degrés appartenant à $[-m, m]$.
\end{enumerate}
Alors $K_n$ n'a de faisceaux de cohomologie non nuls qu'aux degrés appartenant à
$[-m, m] \cup [-m - n, m - n]$ et, pour $n > 2m$, les
morphismes de transition induisent des isomorphismes sur les faisceaux de cohomologie aux degrés appartenant à
$[-m, m]$. De plus, si $X$ et $Y$ sont lisses sur $k$, alors, pour $n$
assez grand, on a $K_n = K \oplus C_n$ dans
$D_{perf}(\mathcal{O}_{X \times Y})$,
où $K$ n'a de cohomologie qu'aux degrés $[-m, m]$ et $C_n$ qu'aux
degrés $[-m - n, m - n]$, et les morphismes de transition
définissent des isomorphismes entre les diverses copies de $K$.
\end{lemma}

\begin{proof}
Soit $Z$ le support schématique d'un module $\mathcal{F}$ comme en (2).
Alors $Z \to \Spec(k)$ est fini, donc $Z \times Y \to Y$ est fini.
Il s'ensuit que, pour un objet $M$ de $D_\QCoh(\mathcal{O}_{X \times Y})$
dont les faisceaux de cohomologie sont supportés par $Z \times Y$, on a
$H^i(R\text{pr}_{2, *}(M)) = \text{pr}_{2, *}H^i(M)$, et le foncteur
$\text{pr}_{2, *}$ est fidèle sur les modules quasi-cohérents supportés
par $Z \times Y$ ; nous omettons les détails. On voit donc que les objets
$$
\text{pr}_1^*\mathcal{F} \otimes_{\mathcal{O}_{X \times Y}}^\mathbf{L} K_n
$$
de $D_{perf}(\mathcal{O}_{X \times Y})$ n'ont de faisceaux de cohomologie non nuls
qu'aux degrés appartenant à $[-m, m] \cup [-m - n, m - n]$ et que, pour $n > 2m$, les morphismes de transition
induisent des isomorphismes sur les faisceaux de cohomologie aux degrés appartenant à $[-m, m]$.
Soit $z \in X \times Y$ un point fermé dont l'image est le point fermé
$x \in X$. Alors on sait que
$$
K_{n, z} \otimes_{\mathcal{O}_{X \times Y, z}}^\mathbf{L}
\mathcal{O}_{X \times Y, z}/\mathfrak m_x^t\mathcal{O}_{X \times Y, z}
$$
n'a de cohomologie non nulle que dans les intervalles
$[-m, m] \cup [-m - n, m - n]$.
On déduit de Compléments sur l'algèbre, lemme
\ref{more-algebra-lemma-kollar-kovacs-pseudo-coherent},
que $K_{n, z}$ n'a de cohomologie non nulle
qu'aux degrés $[-m, m] \cup [-m - n, m - n]$. Comme cela vaut en tout
point fermé de $X \times Y$, on en déduit que $K_n$ n'a de faisceaux de cohomologie non nuls
qu'aux degrés $[-m, m] \cup [-m - n, m - n]$.
De même, on voit que les morphismes $K_n \to K_{n + 1}$
sont des isomorphismes sur les faisceaux de cohomologie aux degrés appartenant à $[-m, m]$
pour $n > 2m$.

\medskip\noindent
Si $X$ et $Y$ sont lisses sur $k$, alors $X \times Y$ est lisse
sur $k$, donc régulier d'après
Variétés, lemme \ref{varieties-lemma-smooth-regular}.
Nous obtiendrons donc la décomposition en somme directe de $K_n$
dès que $n > 2m + \dim(X \times Y)$, d'après
le lemme \ref{equiv-lemma-split-complex-regular}. La dernière assertion
en résulte clairement.
\end{proof}





\section{Foncteurs apparentés}
\label{equiv-section-sibling}

\noindent
Dans cette section, nous démontrons un résultat catégorique portant sur la notion suivante.

\begin{definition}
\label{equiv-definition-siblings}
Soit $\mathcal{A}$ une catégorie abélienne. Soit $\mathcal{D}$ une
catégorie triangulée. Nous disons que deux foncteurs exacts de catégories triangulées
$$
F, F' : D^b(\mathcal{A}) \longrightarrow \mathcal{D}
$$
sont dits {\it apparentés}, ou que $F'$ est {\it apparenté} à $F$,
si les deux conditions suivantes sont satisfaites
\begin{enumerate}
\item les foncteurs $F \circ i$ et $F' \circ i$ sont isomorphes,
où $i : \mathcal{A} \to D^b(\mathcal{A})$ est le foncteur d'inclusion, et
\item $F(K) \cong F'(K)$ pour tout $K$ de $D^b(\mathcal{A})$.
\end{enumerate}
\end{definition}

\noindent
La seconde condition est parfois une conséquence de la première.

\begin{lemma}
\label{equiv-lemma-sibling-fully-faithful}
Soit $\mathcal{A}$ une catégorie abélienne. Soit $\mathcal{D}$ une
catégorie triangulée. Soient
$F, F' : D^b(\mathcal{A}) \longrightarrow \mathcal{D}$
des foncteurs exacts de catégories triangulées. Supposons que
\begin{enumerate}
\item les foncteurs $F \circ i$ et $F' \circ i$ sont isomorphes,
où $i : \mathcal{A} \to D^b(\mathcal{A})$ est le foncteur d'inclusion, et
\item pour tous $X, Y \in \Ob(\mathcal{A})$, nous avons
$\Ext^q_\mathcal{D}(F(X), F(Y)) = 0$ pour $q < 0$ (par exemple
si $F$ est pleinement fidèle).
\end{enumerate}
Alors $F$ et $F'$ sont apparentés.
\end{lemma}

\begin{proof}
Soit $K \in D^b(\mathcal{A})$. Nous allons montrer que $F(K)$ est isomorphe à $F'(K)$.
On peut représenter $K$ par un complexe borné $A^\bullet$ d'objets de
$\mathcal{A}$. Après avoir remplacé $K$ par un décalé, on peut
supposer que $A^i = 0$ pour $i > 0$. Choisissons $n \geq 0$ tel que $A^{-i} = 0$
pour $i > n$. Les objets
$$
M_i = (A^{-i} \to \ldots \to A^0)[-i],\quad i = 0, \ldots, n
$$
forment un système de Postnikov dans $D^b(\mathcal{A})$ pour le complexe
$A^\bullet = A^{-n} \to \ldots \to A^0$ de $D^b(\mathcal{A})$.
Voir Catégories dérivées, exemple \ref{derived-example-key-postnikov}.
Comme $F$ et $F'$ sont tous deux des foncteurs exacts de catégories triangulées,
$$
F(M_i)
\quad\text{et}\quad
F'(M_i)
$$
forment tous deux un système de Postnikov dans $\mathcal{D}$ pour le complexe
$$
F(A^{-n}) \to \ldots \to F(A^0) =
F'(A^{-n}) \to \ldots \to F'(A^0)
$$
Puisque tous les $\Ext$ négatifs entre ces objets sont nuls par hypothèse,
l'unicité des systèmes de Postnikov permet de conclure
(Catégories dérivées, lemme \ref{derived-lemma-existence-postnikov-system})
que $F(K) = F(M_n[n]) \cong F'(M_n[n]) = F'(K)$.
\end{proof}

\begin{lemma}
\label{equiv-lemma-sibling-faithful}
Soient $F$ et $F'$ des foncteurs apparentés comme dans la définition \ref{equiv-definition-siblings}.
Alors
\begin{enumerate}
\item si $F$ est essentiellement surjectif, alors $F'$ est essentiellement
surjectif,
\item si $F$ est pleinement fidèle, alors $F'$ est pleinement fidèle.
\end{enumerate}
\end{lemma}

\begin{proof}
L'assertion (1) est immédiate d'après la propriété (2) des foncteurs apparentés.

\medskip\noindent
Supposons $F$ pleinement fidèle. Notons $\mathcal{D}' \subset \mathcal{D}$
l'image essentielle de $F$, de sorte que $F : D^b(\mathcal{A}) \to \mathcal{D}'$
soit une équivalence. Puisque le foncteur $F'$ se factorise par $\mathcal{D}'$
d'après la propriété (2) des foncteurs apparentés, on peut considérer le foncteur
$H = F^{-1} \circ F' : D^b(\mathcal{A}) \to D^b(\mathcal{A})$.
Observons que $H$ est apparenté au foncteur identique.
Comme il suffit de prouver que $H$ est pleinement fidèle,
on se ramène au problème examiné dans le paragraphe suivant.

\medskip\noindent
Posons $\mathcal{D} = D^b(\mathcal{A})$. Il faut montrer qu'un foncteur
$F : \mathcal{D} \to \mathcal{D}$ apparenté au foncteur identique est pleinement fidèle.
Notons $a_X : X \to F(X)$ l'isomorphisme fonctoriel pour
$X \in \Ob(\mathcal{A})$ fourni par la définition \ref{equiv-definition-siblings}.
Pour tout $K$ de $\mathcal{D}$ et tout triangle distingué
$K_1 \to K_2 \to K_3$ de $\mathcal{D}$,
si les applications
$$
F : \Hom(K, K_i[n]) \to \Hom(F(K), F(K_i[n]))
$$
sont des isomorphismes pour tout $n \in \mathbf{Z}$ et $i = 1, 3$, alors il en
est de même pour $i = 2$ et tout $n \in \mathbf{Z}$. On utilise ici
le lemme des cinq, Homologie, lemme \ref{homology-lemma-five-lemma}, et
Catégories dérivées, lemme \ref{derived-lemma-representable-homological} ;
nous omettons les détails. De même, si les applications
$$
F : \Hom(K_i[n], K) \to \Hom(F(K_i[n]), F(K))
$$
sont des isomorphismes pour tout $n \in \mathbf{Z}$ et $i = 1, 3$, alors il en
est de même pour $i = 2$ et tout $n \in \mathbf{Z}$. En utilisant les troncations
canoniques et une récurrence sur le nombre d'objets de cohomologie non nuls,
on voit qu'il suffit de montrer que
$$
F : \Ext^q(X, Y) \to \Ext^q(F(X), F(Y))
$$
est bijective pour tous $X, Y \in \Ob(\mathcal{A})$ et tout $q \in \mathbf{Z}$.
Comme $F$ est apparenté à $\text{id}$, on a $F(X) \cong X$ et
$F(Y) \cong Y$ ; le membre de droite est donc nul pour $q < 0$.
Le cas $q = 0$ résulte de l'hypothèse que $F$ est apparenté au
foncteur identique. Il reste à traiter les cas $q > 0$.

\medskip\noindent
Le cas $q = 1$ : injectivité. Un élément $\xi$ de $\Ext^1(X, Y)$
donne lieu à un triangle distingué
$$
Y \to E \to X \xrightarrow{\xi} Y[1]
$$
Observons que $E \in \Ob(\mathcal{A})$. Comme $F$ est apparenté au
foncteur identique, on obtient un diagramme commutatif
$$
\xymatrix{
E \ar[d] \ar[r] & X \ar[d] \\
F(E) \ar[r] & F(X)
}
$$
dont les flèches verticales sont les isomorphismes $a_E$ et $a_X$.
Par TR3, le triangle distingué associé à $\xi$ dont on est parti
est isomorphe au triangle distingué
$$
F(Y) \to F(E) \to F(X) \xrightarrow{F(\xi)} F(Y[1]) = F(Y)[1]
$$
Ainsi, $\xi = 0$ si et seulement si $F(\xi)$ est nul, c'est-à-dire que
$F : \Ext^1(X, Y) \to \Ext^1(F(X), F(Y))$ est injective.

\medskip\noindent
Le cas $q = 1$ : surjectivité. Soit $\theta$ un élément de
$\Ext^1(F(X), F(Y))$. Il définit une extension de $F(X)$ par $F(Y)$
dans $\mathcal{A}$, que l'on peut écrire sous la forme $F(E)$
puisque $F$ est apparenté au foncteur identique. On obtient ainsi un triangle
distingué
$$
F(Y) \xrightarrow{F(\alpha)} F(E)
\xrightarrow{F(\beta)} F(X)
\xrightarrow{\theta} F(Y[1]) = F(Y)[1]
$$
pour certains morphismes $\alpha : Y \to E$ et $\beta : E \to X$.
Comme $F$ est apparenté au foncteur identique, la suite
$0 \to Y \to E \to X \to 0$
est une suite exacte courte de $\mathcal{A}$ ! On obtient donc un
triangle distingué
$$
Y \xrightarrow{\alpha} E \xrightarrow{\beta} X \xrightarrow{\delta} Y[1]
$$
pour un certain morphisme $\delta : X \to Y[1]$. En appliquant le foncteur exact
$F$, on obtient le triangle distingué
$$
F(Y) \xrightarrow{F(\alpha)} F(E) \xrightarrow{F(\beta)} F(X)
\xrightarrow{F(\delta)} F(Y)[1]
$$
En raisonnant comme ci-dessus, on voit que ces triangles sont isomorphes.
Il existe donc un diagramme commutatif
$$
\xymatrix{
F(X) \ar[d]^\gamma \ar[r]_{F(\delta)} & F(Y[1]) \ar[d]_\epsilon \\
F(X) \ar[r]^\theta & F(Y[1])
}
$$
pour certains isomorphismes $\gamma$, $\epsilon$ (on pourrait être plus précis, mais cela
ne sera pas nécessaire). On peut écrire $\gamma = F(\gamma')$ et
$\epsilon = F(\epsilon')$. On a alors
$\theta = F(\epsilon' \circ \delta \circ (\gamma')^{-1})$,
ce qui prouve la surjectivité.

\medskip\noindent
Le cas $q > 1$ : surjectivité. En utilisant les extensions de Yoneda, voir
Catégories dérivées, section \ref{derived-section-ext}, on voit que, pour tout
élément $\xi$ de $\Ext^q(F(X), F(Y))$, il existe
$F(X) = B_0, B_1, \ldots, B_{q - 1}, B_q = F(Y) \in \Ob(\mathcal{A})$ et des
éléments
$$
\xi_i \in \Ext^1(B_{i - 1}, B_i)
$$
tels que $\xi$ soit le composé $\xi_q \circ \ldots \circ \xi_1$.
Écrivons $B_i = F(A_i)$ (bien entendu, on a $A_i = B_i$, mais il n'est pas
nécessaire de l'utiliser), de sorte que
$$
\xi_i = F(\eta_i) \in \Ext^1(F(A_{i - 1}), F(A_i))
\quad\text{avec}\quad
\eta_i \in \Ext^1(A_{i - 1}, A_i)
$$
par surjectivité pour $q = 1$. Alors $\eta = \eta_q \circ \ldots \circ \eta_1$
est un élément de $\Ext^q(X, Y)$ tel que $F(\eta) = \xi$.

\medskip\noindent
Le cas $q > 1$ : injectivité. Un élément $\xi$ de $\Ext^q(X, Y)$
donne lieu à un triangle distingué
$$
Y[q - 1] \to E \to X \xrightarrow{\xi} Y[q]
$$
En appliquant $F$, on obtient un triangle distingué
$$
F(Y)[q - 1] \to F(E) \to F(X) \xrightarrow{F(\xi)} F(Y)[q]
$$
Si $F(\xi) = 0$, alors $F(E) \cong F(Y)[q - 1] \oplus F(X)$
dans $\mathcal{D}$, voir
Catégories dérivées, lemme \ref{derived-lemma-split}.
Comme $F$ est apparenté au foncteur identique, on a
$E \cong F(E)$, et donc
$$
E \cong F(E) \cong F(Y)[q - 1] \oplus F(X) \cong Y[q - 1] \oplus X
$$
Autrement dit, $E$ est isomorphe à la
somme directe de ses objets de cohomologie. Cela implique que le
triangle distingué initial est scindé, c'est-à-dire que $\xi = 0$.
\end{proof}

\noindent
Donnons une définition non standard. Soit $\mathcal{A}$ une catégorie
abélienne. Disons que $\mathcal{A}$ {\it possède assez d'objets négatifs}
si, pour tout $X \in \Ob(\mathcal{A})$, il existe un objet $N$ tel que
\begin{enumerate}
\item il existe une surjection $N \to X$, et
\item $\Hom(X, N) = 0$.
\end{enumerate}
Démontrons deux lemmes sur cette notion afin de faciliter
la démonstration de la proposition \ref{equiv-proposition-siblings-isomorphic}.

\begin{lemma}
\label{equiv-lemma-good-map}
Soit $\mathcal{A}$ une catégorie abélienne possédant assez d'objets négatifs.
Soit $X \in D^b(\mathcal{A})$. Soit $b \in \mathbf{Z}$ tel que
$H^i(X) = 0$ pour $i > b$. Alors
il existe un morphisme $N[-b] \to X$ tel que le morphisme induit
$N \to H^b(X)$ soit surjectif et que $\Hom(H^b(X), N) = 0$.
\end{lemma}

\begin{proof}
À l'aide des foncteurs de troncation, on peut représenter $X$ par un complexe
$A^a \to A^{a + 1} \to \ldots \to A^b$ d'objets de $\mathcal{A}$.
Choisissons $N$ dans $\mathcal{A}$ de sorte qu'il existe une surjection
$t : N \to A^b$ et que $\Hom(A^b, N) = 0$. Alors la surjection $t$
définit un morphisme $N[-b] \to X$ comme voulu.
\end{proof}

\begin{lemma}
\label{equiv-lemma-good-map-zero}
Soit $\mathcal{A}$ une catégorie abélienne possédant assez d'objets négatifs.
Soit $f : X \to X'$ un morphisme de $D^b(\mathcal{A})$. Soit $b \in \mathbf{Z}$
tel que $H^i(X) = 0$ pour $i > b$ et $H^i(X') = 0$ pour $i \geq b$.
Alors il existe un morphisme $N[-b] \to X$ tel que le morphisme induit
$N \to H^b(X)$ soit surjectif, que $\Hom(H^b(X), N) = 0$, et
que le composé $N[-b] \to X \to X'$ soit nul.
\end{lemma}

\begin{proof}
On peut représenter $f$ par un morphisme $f^\bullet : A^\bullet \to B^\bullet$
de complexes bornés d'objets de $\mathcal{A}$, voir par exemple
Catégories dérivées, lemme \ref{derived-lemma-bounded-derived}.
Considérons l'objet
$$
C = \Ker(A^b \to A^{b + 1}) \times_{\Ker(B^b \to B^{b + 1})} B^{b - 1}
$$
de $\mathcal{A}$. Comme $H^b(B^\bullet) = 0$, on voit que
$C \to H^b(A^\bullet)$ est surjectif. D'autre part, le morphisme
$C \to A^b \to B^b$ est le même que le morphisme $C \to B^{b - 1} \to B^b$,
et le composé $C[-b] \to X \to X'$ est donc nul.
Puisque $\mathcal{A}$ possède assez d'objets négatifs, on peut trouver un objet $N$
muni d'une surjection $N \to C \oplus H^b(X)$ tel que
$\Hom(C \oplus H^b(X), N) = 0$. Alors $N$, avec le morphisme
$N[-b] \to X$, est une solution au problème posé par le lemme.
\end{proof}

\noindent
Nous encourageons le lecteur à lire l'original
\cite[Proposition 2.16]{Orlov-K3} pour les idées remarquables
qui interviennent dans la démonstration de la proposition suivante.

\begin{proposition}
\label{equiv-proposition-siblings-isomorphic}
\begin{reference}
\cite[Proposition 2.16]{Orlov-K3} ; le fait qu'il ne soit pas nécessaire
de supposer la nullité de $\Ext^q(N, X)$ pour $q > 0$ dans la définition
des objets négatifs ci-dessus est dû à \cite{Canonaco-Stellari}.
\end{reference}
Soient $F$ et $F'$ des foncteurs apparentés comme dans la définition \ref{equiv-definition-siblings}.
Supposons que $F$ soit pleinement fidèle et que $\mathcal{A}$ possède assez
d'objets négatifs (voir ci-dessus). Alors $F$ et $F'$ sont des foncteurs isomorphes.
\end{proposition}

\begin{proof}
D'après l'assertion (2) de la définition \ref{equiv-definition-siblings}, l'image du foncteur
$F'$ est contenue dans l'image essentielle du foncteur $F$. Ainsi,
le foncteur $H = F^{-1} \circ F'$ est apparenté au foncteur identique.
Cela nous ramène au cas décrit dans le paragraphe suivant.

\medskip\noindent
Soit $\mathcal{D} = D^b(\mathcal{A})$. Il faut montrer qu'un foncteur
$F : \mathcal{D} \to \mathcal{D}$ apparenté au foncteur identique est
isomorphe au foncteur identique. Pour un objet $X$ de $\mathcal{D}$,
disons que $X$ est de {\it largeur} $w = w(X)$ si $w \geq 0$ est minimal
parmi les entiers tels qu'il existe $a \in \mathbf{Z}$ avec $H^i(X) = 0$
pour $i \not \in [a, a + w - 1]$. Puisque $F$ est apparenté au foncteur identique
et que $F \circ [n] = [n] \circ F$, on dispose déjà d'isomorphismes
$$
c_X : X \to F(X)
$$
pour $w(X) \leq 1$, compatibles aux décalages. De plus, si $X = A[-a]$ et
$X' = A'[-a]$ pour certains $A, A' \in \Ob(\mathcal{A})$, alors, pour tout morphisme
$f : X \to X'$, le diagramme
\begin{equation}
\label{equiv-equation-to-show}
\vcenter{
\xymatrix{
X \ar[d]_{c_X} \ar[r]_f &
X' \ar[d]^{c_{X'}} \\
F(X) \ar[r]^{F(f)} &
F(X')
}
}
\end{equation}
est commutatif.

\medskip\noindent
Montrons ensuite que, pour tout morphisme $f : X \to X'$ tel que
$w(X), w(X') \leq 1$, le diagramme (\ref{equiv-equation-to-show}) est commutatif.
Si $X$ ou $X'$ est nul, c'est clair. Sinon, on peut écrire
$X = A[-a]$ et $X' = A'[-a']$ pour d'uniques $A, A'$ dans $\mathcal{A}$
et $a, a' \in \mathbf{Z}$. Le cas $a = a'$ a été examiné ci-dessus.
Si $a' > a$, alors $f = 0$ (Catégories dérivées, lemme
\ref{derived-lemma-negative-exts}), et le résultat est clair.
Si $a' < a$, alors $f$ correspond à un élément
$\xi \in \Ext^q(A, A')$ avec $q = a - a'$. En utilisant les extensions de Yoneda, voir
Catégories dérivées, section \ref{derived-section-ext}, on peut trouver
$A = A_0, A_1, \ldots, A_{q - 1}, A_q = A' \in \Ob(\mathcal{A})$ et des
éléments
$$
\xi_i \in \Ext^1(A_{i - 1}, A_i)
$$
tels que $\xi$ soit le composé $\xi_q \circ \ldots \circ \xi_1$.
Autrement dit, en posant $X_i = A_i[-a + i]$,
on obtient des morphismes
$$
X = X_0 \xrightarrow{f_1} X_1 \to \ldots \to X_{q - 1}
\xrightarrow{f_q} X_q = X'
$$
dont le composé est $f$. Comme la commutativité de (\ref{equiv-equation-to-show})
pour $f_1, \ldots, f_q$ implique celle pour $f$, on se ramène au cas $q = 1$.
Dans ce cas, après décalage, on peut supposer que l'on a un triangle distingué
$$
A' \to E \to A \xrightarrow{f} A'[1]
$$
Observons que $E$ est un objet de $\mathcal{A}$. Considérons le
diagramme suivant
$$
\xymatrix{
E \ar[d]_{c_E} \ar[r] &
A \ar[d]_{c_A} \ar[r]_f &
A'[1] \ar[d]^{c_{A'}[1]}
\ar@{..>}@<-1ex>[d]_\gamma \ar@{..>}[ld]^\epsilon \ar[r] &
E[1] \ar[d]^{c_E[1]} \\
F(E) \ar[r] &
F(A) \ar[r]^{F(f)} &
F(A')[1] \ar[r] &
F(E)[1]
}
$$
dont les lignes sont des triangles distingués.
Le carré de droite est déjà commutatif, mais on ne sait pas encore que
le carré central l'est. D'après les axiomes d'une catégorie triangulée,
on peut trouver un morphisme $\gamma$ qui rende le diagramme commutatif.
Alors le composé de $\gamma - c_{A'}[1]$ avec
$F(A')[1] \to F(E)[1]$ est nul ; on peut donc
trouver $\epsilon : A'[1] \to F(A)$ tel que
$\gamma - c_{A'}[1] = F(f) \circ \epsilon$. Cependant, toute flèche
$A'[1] \to F(A)$ est nulle, car c'est une classe d'extension négative
entre objets de $\mathcal{A}$. Ainsi, $\gamma = c_{A'}[1]$,
et l'on conclut que le carré central est lui aussi commutatif, ce que l'on
voulait montrer.

\medskip\noindent
Pour achever la démonstration, nous allons raisonner par récurrence sur $w$
et construire des isomorphismes $c_X : X \to F(X)$ pour tout
$X$ tel que $w(X) \leq w$, compatibles avec tous les morphismes entre
de tels objets. Le cas initial $w = 1$ a été démontré ci-dessus. Supposons
le résultat connu pour un certain $w \geq 1$.

\medskip\noindent
Soit $X$ un objet tel que $w(X) = w + 1$. Choisissons $a \in \mathbf{Z}$ avec
$H^i(X) = 0$ pour $i \not \in [a, a + w]$. Posons $b = a + w$, de sorte que
$H^b(X)$ soit non nul. Choisissons $N[-b] \to X$ comme dans le lemme \ref{equiv-lemma-good-map}.
Choisissons un triangle distingué
$$
N[-b] \to X \to Y \to N[-b + 1]
$$
Le calcul de la suite exacte longue de cohomologie donne
$w(Y) \leq w$. Par récurrence, on obtient donc les flèches pleines
du diagramme suivant
$$
\xymatrix{
N[-b] \ar[r] \ar[d]_{c_N[-b]} &
X \ar[r] \ar@{..>}[d]_{c_{N[-b] \to X}} &
Y \ar[r] \ar[d]^{c_Y} &
N[-b + 1] \ar[d]^{c_N[-b + 1]} \\
F(N)[-b] \ar[r] &
F(X) \ar[r] &
F(Y) \ar[r] &
F(N)[-b + 1]
}
$$
On obtient la flèche pointillée $c_{N[-b] \to X}$.
D'après Catégories dérivées, lemme \ref{derived-lemma-uniqueness-third-arrow},
la flèche pointillée est unique, car $\Hom(X, F(N)[-b]) \cong \Hom(X, N[-b]) = 0$
par le choix de $N$. En fait, $c_{N[-b] \to X}$ est l'unique flèche
pointillée qui rende commutatif le carré de sommets $X, Y, F(X), F(Y)$.

\medskip\noindent
Soit $N'[-b] \to X$ un autre morphisme comme dans le lemme \ref{equiv-lemma-good-map},
et prouvons que $c_{N[-b] \to X} = c_{N'[-b] \to X}$.
Observons que le morphisme $(N \oplus N')[-b] \to X$ satisfait lui aussi les
conditions du lemme \ref{equiv-lemma-good-map}.
On peut donc supposer que $N'[-b] \to X$ se factorise
en $N'[-b] \to N[-b] \to X$ pour un certain morphisme $N' \to N$.
Choisissons des triangles distingués $N[-b] \to X \to Y \to N[-b + 1]$ et
$N'[-b] \to X \to Y' \to N'[-b + 1]$. D'après l'axiome TR3, il existe
un morphisme $g : Y' \to Y$ qui, avec $\text{id}_X$ et $N' \to N$,
forme un morphisme de triangles. Comme on dispose de
(\ref{equiv-equation-to-show}) pour $g$, on en déduit que
$$
(F(X) \to F(Y)) \circ c_{N'[-b] \to X} = (F(X) \to F(Y)) \circ c_{N[-b] \to X}
$$
L'unicité de $c_{N[-b] \to X}$ relevée dans la construction
ci-dessus montre alors que $c_{N'[-b] \to X} = c_{N[-b] \to X}$.

\medskip\noindent
On peut donc maintenant définir, pour $X$ de largeur $w + 1$, l'isomorphisme
$c_X : X \to F(X)$ comme la valeur commune des morphismes $c_{N[-b] \to X}$,
où $N[-b] \to X$ est comme dans le lemme \ref{equiv-lemma-good-map}. Pour achever
la démonstration, il faut montrer que les diagrammes (\ref{equiv-equation-to-show})
sont commutatifs pour tout morphisme $f : X \to X'$ entre objets tels que $w(X) \leq w + 1$
et $w(X') \leq w + 1$. Choisissons $a \leq b \leq a + w$ tel que
$H^i(X) = 0$ pour $i \not \in [a, b]$ et
$a' \leq b' \leq a' + w$ tel que $H^i(X') = 0$ pour
$i \not \in [a', b']$. Nous allons raisonner par récurrence sur
$(b' - a') + (b - a)$ pour démontrer l'assertion. (Le cas initial
est celui où ce nombre est nul, ce qui convient puisque $w \geq 1$.)
Distinguons deux cas.

\medskip\noindent
Cas I : $b' < b$. Dans ce cas, d'après le lemme \ref{equiv-lemma-good-map-zero},
on peut choisir $N[-b] \to X$ comme dans le lemme \ref{equiv-lemma-good-map},
de sorte que le composé $N[-b] \to X \to X'$ soit nul.
Choisissons un triangle distingué $N[-b] \to X \to Y \to N[-b + 1]$. Comme
$N[-b] \to X'$ est nul, on voit que $f$ se factorise
en $X \to Y \to X'$. Comme $H^i(Y)$ n'est non nul que pour $i \in [a, b - 1]$,
l'hypothèse de récurrence montre que (\ref{equiv-equation-to-show}) est commutatif pour
$Y \to X'$. Le diagramme (\ref{equiv-equation-to-show}) est commutatif pour
$X \to Y$ par construction si $w(X) = w + 1$, et par notre première
hypothèse de récurrence si $w(X) \leq w$.
Ainsi, (\ref{equiv-equation-to-show}) est commutatif pour $f$.

\medskip\noindent
Cas II : $b' \geq b$. Dans ce cas, choisissons $N'[-b'] \to X'$
comme dans le lemme \ref{equiv-lemma-good-map}.
On peut aussi supposer que $\Hom(H^{b'}(X), N') = 0$ (cela n'est
pertinent que si $b' = b$), car on peut, par exemple,
remplacer $N'$ par un objet $N''$ muni d'une surjection sur $N' \oplus H^{b'}(X)$
et tel que $\Hom(N' \oplus H^{b'}(X), N'') = 0$.
Choisissons un triangle distingué
$N'[-b'] \to X' \to Y' \to N'[-b' + 1]$. Comme
$\Hom(X, X') \to \Hom(X, Y')$ est injectif par notre choix de $N'$
(nous omettons les détails), il en va de même pour
$\Hom(X, F(X')) \to \Hom(X, F(Y'))$.
Il suffit donc, dans ce cas, de vérifier que
(\ref{equiv-equation-to-show}) est commutatif pour le composé $X \to Y'$
des morphismes $X \to X' \to Y'$.
Comme $H^i(Y')$ n'est non nul que pour $i \in [a', b' - 1]$,
on conclut par l'hypothèse de récurrence.
\end{proof}










\section{Établissement de la pleine fidélité}
\label{equiv-section-get-fully-faithful}

\noindent
Il nous sera utile de savoir quand un foncteur est pleinement fidèle ;
nous proposons la variante suivante de \cite[Lemme 2.15]{Orlov-K3}.

\begin{lemma}
\label{equiv-lemma-get-fully-faithful}
\begin{reference}
Variante de \cite[Lemme 2.15]{Orlov-K3}
\end{reference}
Soit $F : \mathcal{D} \to \mathcal{D}'$ un foncteur exact de
catégories triangulées. Soit $S \subset \Ob(\mathcal{D})$ un
ensemble d'objets. Supposons que
\begin{enumerate}
\item $F$ possède des adjoints à droite et à gauche,
\item pour $K \in \mathcal{D}$, si $\Hom(E, K[i]) = 0$ pour tous
$E \in S$ et $i \in \mathbf{Z}$, alors $K = 0$,
\item pour $K \in \mathcal{D}$, si $\Hom(K, E[i]) = 0$ pour tous
$E \in S$ et $i \in \mathbf{Z}$, alors $K = 0$,
\item l'application $\Hom(E, E'[i]) \to \Hom(F(E), F(E')[i])$ induite par $F$
est bijective pour tous $E, E' \in S$ et $i \in \mathbf{Z}$.
\end{enumerate}
Alors $F$ est pleinement fidèle.
\end{lemma}

\begin{proof}
Notons $F_r$ et $F_l$ les adjoints à droite et à gauche de $F$. Pour
$E \in S$, choisissons un triangle distingué
$$
E \to F_r(F(E)) \to C \to E[1]
$$
où la première flèche est l'unité de l'adjonction. Pour $E' \in S$, on a
$$
\Hom(E', F_r(F(E))[i]) = \Hom(F(E'), F(E)[i]) = \Hom(E', E[i])
$$
La dernière égalité vaut par l'hypothèse (4).
En appliquant donc le foncteur homologique $\Hom(E', -)$
(Catégories dérivées, lemme \ref{derived-lemma-representable-homological})
au triangle distingué ci-dessus, on en déduit que $\Hom(E', C[i]) = 0$
pour tout $i \in \mathbf{Z}$ et $E' \in S$. L'hypothèse (2) donne alors
$C = 0$ et $E = F_r(F(E))$.

\medskip\noindent
Pour $K \in \Ob(\mathcal{D})$, choisissons un triangle distingué
$$
F_l(F(K)) \to K \to C \to F_l(F(K))[1]
$$
où la première flèche est la counité de l'adjonction. Pour $E \in S$,
on a
$$
\Hom(F_l(F(K)), E[i]) = \Hom(F(K), F(E)[i]) =
\Hom(K, F_r(F(E))[i]) = \Hom(K, E[i])
$$
où la dernière égalité résulte du premier paragraphe.
On conclut donc comme précédemment que $\Hom(C, E[i]) = 0$ pour tout $E \in S$
et $i \in \mathbf{Z}$. Ainsi, $C = 0$ d'après l'hypothèse (3).
Le foncteur $F$ est donc pleinement fidèle d'après Catégories, lemme
\ref{categories-lemma-adjoint-fully-faithful}.
\end{proof}

\begin{lemma}
\label{equiv-lemma-duality-at-point}
Soit $k$ un corps. Soit $X$ un schéma de type fini sur $k$ qui
est régulier. Soit $x \in X$ un point fermé. Pour un
$\mathcal{O}_X$-module cohérent $\mathcal{F}$ à support en $x$, choisissons
un $\mathcal{O}_X$-module cohérent $\mathcal{F}'$ à support en $x$
tel que $\mathcal{F}_x$ et $\mathcal{F}'_x$ soient duaux de Matlis.
Alors il existe un isomorphisme
$$
\Hom_X(\mathcal{F}, M) =
H^0(X, M \otimes_{\mathcal{O}_X}^\mathbf{L} \mathcal{F}'[-d_x])
$$
où $d_x = \dim(\mathcal{O}_{X, x})$,
fonctoriel en $M$ dans $D_{perf}(\mathcal{O}_X)$.
\end{lemma}

\begin{proof}
Puisque $\mathcal{F}$ est à support en $x$, on a
$$
\Hom_X(\mathcal{F}, M) =
\Hom_{\mathcal{O}_{X, x}}(\mathcal{F}_x, M_x)
$$
et, de même, on a
$$
H^0(X, M \otimes_{\mathcal{O}_X}^\mathbf{L} \mathcal{F}'[-d_x]) =
\text{Tor}^{\mathcal{O}_{X, x}}_{d_x}(M_x, \mathcal{F}'_x)
$$
Il suffit donc de montrer que, pour un anneau local noethérien régulier $A$
de dimension $d$ et un $A$-module $N$ de longueur finie, si
$N'$ est le dual de Matlis de $N$, alors il existe un isomorphisme fonctoriel
$$
\Hom_A(N, K) = \text{Tor}^A_d(K, N')
$$
pour $K$ dans $D_{perf}(A)$. On peut écrire le membre de gauche sous la forme
$H^0(R\Hom_A(N, A) \otimes_A^\mathbf{L} K)$ d'après
Compléments sur l'algèbre, lemme \ref{more-algebra-lemma-dual-perfect-complex},
et le fait que $N$ détermine un objet parfait de $D(A)$.
La formule vaut donc parce que
$$
R\Hom_A(N, A) = R\Hom_A(N, A[d])[-d] = N'[-d]
$$
d'après Complexes dualisants, lemme \ref{dualizing-lemma-dualizing-finite-length},
et parce que $A[d]$ est un complexe dualisant normalisé sur $A$
($A$ est un anneau de Gorenstein d'après
Complexes dualisants, lemme \ref{dualizing-lemma-regular-gorenstein}).
\end{proof}

\begin{lemma}
\label{equiv-lemma-orthogonal-point-sheaf}
Soit $k$ un corps. Soit $X$ un schéma de type fini sur $k$ qui
est régulier. Soit $x \in X$ un point fermé et notons $\mathcal{O}_x$
le faisceau gratte-ciel en $x$ de valeur $\kappa(x)$. Soit $K$ dans
$D_{perf}(\mathcal{O}_X)$.
\begin{enumerate}
\item Si $\Ext^i_X(\mathcal{O}_x, K) = 0$, alors il existe un voisinage
ouvert $U$ de $x$ tel que $H^{i - d_x}(K)|_U = 0$, où
$d_x = \dim(\mathcal{O}_{X, x})$.
\item Si $\Hom_X(\mathcal{O}_x, K[i]) = 0$ pour tout
$i \in \mathbf{Z}$, alors $K$ est nul dans un voisinage ouvert de $x$.
\item Si $\Ext^i_X(K, \mathcal{O}_x) = 0$, alors il existe un voisinage
ouvert $U$ de $x$ tel que $H^i(K^\vee)|_U = 0$.
\item Si $\Hom_X(K, \mathcal{O}_x[i]) = 0$ pour tout
$i \in \mathbf{Z}$, alors $K$ est nul dans un voisinage ouvert de $x$.
\item Si $H^i(X, K \otimes_{\mathcal{O}_X}^\mathbf{L} \mathcal{O}_x) = 0$,
alors il existe un voisinage ouvert $U$ de $x$ tel que
$H^i(K)|_U = 0$.
\item Si $H^i(X, K \otimes_{\mathcal{O}_X}^\mathbf{L} \mathcal{O}_x) = 0$
pour $i \in \mathbf{Z}$, alors $K$ est nul dans un
voisinage ouvert de $x$.
\end{enumerate}
\end{lemma}

\begin{proof}
Observons que $H^i(X, K \otimes_{\mathcal{O}_X}^\mathbf{L} \mathcal{O}_x)$
est égal à $K_x  \otimes_{\mathcal{O}_{X, x}}^\mathbf{L} \kappa(x)$.
L'assertion (5) résulte donc de Compléments sur l'algèbre, lemme
\ref{more-algebra-lemma-cut-complex-in-two}.
L'assertion (6) résulte de l'assertion (5).
L'assertion (1) résulte de l'assertion (5), du lemme \ref{equiv-lemma-duality-at-point} et du
fait que le dual de Matlis de $\kappa(x)$ est $\kappa(x)$.
L'assertion (2) résulte de l'assertion (1).
L'assertion (3) résulte de l'assertion (5) et du fait que
$\Ext^i(K, \mathcal{O}_x) =
H^i(X, K^\vee \otimes_{\mathcal{O}_X}^\mathbf{L} \mathcal{O}_x)$ d'après
Cohomologie, lemme \ref{cohomology-lemma-dual-perfect-complex}.
L'assertion (4) résulte de l'assertion (3) et du fait que $K \cong (K^\vee)^\vee$
d'après le lemme qui vient d'être cité.
\end{proof}

\begin{lemma}
\label{equiv-lemma-hom-into-point-sheaf}
Soit $X$ un schéma noethérien. Soit $x \in X$ un point fermé et
notons $\mathcal{O}_x$ le faisceau gratte-ciel en $x$ de valeur $\kappa(x)$.
Soit $K$ dans $D^b_{\textit{Coh}}(\mathcal{O}_X)$. Soit $b \in \mathbf{Z}$.
Les conditions suivantes sont équivalentes :
\begin{enumerate}
\item $H^i(K)_x = 0$ pour tout $i > b$, et
\item $\Hom_X(K, \mathcal{O}_x[-i]) = 0$ pour tout $i > b$.
\end{enumerate}
\end{lemma}

\begin{proof}
Considérons le complexe $K_x$ dans $D^b_{\textit{Coh}}(\mathcal{O}_{X, x})$.
Il existe un entier $b_x \in \mathbf{Z}$ tel que $K_x$
puisse être représenté par un complexe borné supérieurement
$$
\ldots \to
\mathcal{O}_{X, x}^{\oplus n_{b_x - 2}} \to
\mathcal{O}_{X, x}^{\oplus n_{b_x - 1}} \to
\mathcal{O}_{X, x}^{\oplus n_{b_x}} \to 0 \to \ldots
$$
où $\mathcal{O}_{X, x}^{\oplus n_i}$ est placé en degré $i$
et où tous les morphismes de transition sont donnés par des matrices dont les
coefficients appartiennent à $\mathfrak m_x$. Voir
Compléments sur l'algèbre, lemme
\ref{more-algebra-lemma-lift-pseudo-coherent-from-residue-field}.
Le résultat s'en déduit facilement (et les conditions équivalentes
sont satisfaites si et seulement si $b \geq b_x$).
\end{proof}

\begin{lemma}
\label{equiv-lemma-get-fully-faithful-geometric}
Soit $k$ un corps. Soient $X$ et $Y$ des schémas propres sur $k$.
Supposons que $X$ soit régulier. Alors un foncteur exact $k$-linéaire
$F : D_{perf}(\mathcal{O}_X) \to D_{perf}(\mathcal{O}_Y)$
est pleinement fidèle si et seulement si
pour tous les points fermés $x, x' \in X$, les applications
$$
F : \Ext^i_X(\mathcal{O}_x, \mathcal{O}_{x'})
\longrightarrow
\Ext^i_Y(F(\mathcal{O}_x), F(\mathcal{O}_{x'}))
$$
sont des isomorphismes pour tout $i \in \mathbf{Z}$.
Ici, $\mathcal{O}_x$ est le faisceau gratte-ciel en $x$ de valeur $\kappa(x)$.
\end{lemma}

\begin{proof}
D'après le lemme \ref{equiv-lemma-always-right-adjoints}, le foncteur $F$
possède un adjoint à gauche et un adjoint à droite. On peut donc appliquer le critère
du lemme \ref{equiv-lemma-get-fully-faithful},
car les hypothèses (2) et (3) de ce lemme
résultent du lemme \ref{equiv-lemma-orthogonal-point-sheaf}.
\end{proof}

\begin{lemma}
\label{equiv-lemma-noah-pre}
\begin{reference}
Courriel de Noah Olander du 9 juin 2020
\end{reference}
Soit $k$ un corps. Soit $X$ un schéma propre sur $k$ qui est régulier.
Soit $F : D_{perf}(\mathcal{O}_X) \to D_{perf}(\mathcal{O}_X)$
un foncteur exact $k$-linéaire. Supposons que, pour tout
$\mathcal{O}_X$-module cohérent $\mathcal{F}$ tel que $\dim(\text{Supp}(\mathcal{F})) = 0$,
il existe un isomorphisme $\mathcal{F} \cong F(\mathcal{F})$.
Alors $F$ est pleinement fidèle.
\end{lemma}

\begin{proof}
D'après le lemme \ref{equiv-lemma-get-fully-faithful-geometric}, il suffit de montrer
que les applications
$$
F : \Ext^i_X(\mathcal{O}_x, \mathcal{O}_{x'})
\longrightarrow
\Ext^i_X(F(\mathcal{O}_x), F(\mathcal{O}_{x'}))
$$
sont des isomorphismes pour tout $i \in \mathbf{Z}$ et tous les points fermés
$x, x' \in X$. Par hypothèse, la source et le but sont isomorphes.
Si $x \not = x'$, alors les deux membres sont nuls et le résultat est vrai.
Si $x = x'$, il suffit alors de prouver que l'application est injective
ou surjective. Pour $i < 0$, les deux membres sont nuls et le résultat est vrai.
Pour $i = 0$, tout morphisme non nul $\alpha : \mathcal{O}_x \to \mathcal{O}_x$ de
$\mathcal{O}_X$-modules est un isomorphisme. Ainsi, $F(\alpha)$ est lui aussi un
isomorphisme, donc $F(\alpha)$ est non nul. Cela prouve le résultat pour $i = 0$.
Pour $i = 1$, un élément non nul $\xi$ de $\Ext^1(\mathcal{O}_x, \mathcal{O}_x)$
correspond à une suite exacte courte non scindée
$$
0 \to \mathcal{O}_x \to \mathcal{F} \to \mathcal{O}_x \to 0
$$
Comme $F(\mathcal{F}) \cong \mathcal{F}$, on voit que $F(\mathcal{F})$
est également une extension non scindée de $\mathcal{O}_x$ par $\mathcal{O}_x$.
Puisque $\mathcal{O}_x \cong F(\mathcal{O}_x)$ est un
$\mathcal{O}_X$-module simple et que $\mathcal{F} \cong F(\mathcal{F})$ est de
longueur $2$, on voit que, dans le triangle distingué
$$
F(\mathcal{O}_x) \to F(\mathcal{F}) \to F(\mathcal{O}_x)
\xrightarrow{F(\xi)} F(\mathcal{O}_x)[1]
$$
les deux premières flèches doivent former une suite exacte courte qui doit être
isomorphe à la suite exacte courte ci-dessus, et qui est donc non scindée.
Il s'ensuit que $F(\xi)$ est non nul, ce qui permet de conclure pour $i = 1$.
Pour $i > 1$, la composition des classes d'extensions définit une surjection
$$
\Ext^1(F(\mathcal{O}_x), F(\mathcal{O}_x)) \otimes \ldots \otimes
\Ext^1(F(\mathcal{O}_x), F(\mathcal{O}_x))
\longrightarrow
\Ext^i(F(\mathcal{O}_x), F(\mathcal{O}_x))
$$
Voir Dualité pour les schémas, lemme \ref{duality-lemma-regular-ideal-ext}.
La surjectivité en degré $1$ implique donc la surjectivité pour $i > 0$.
Cela achève la démonstration.
\end{proof}










\section{Foncteurs particuliers}
\label{equiv-section-special-functors}

\noindent
Dans cette section, nous démontrons quelques résultats sur des foncteurs d'un type particulier
que nous utiliserons plus loin dans ce chapitre.

\begin{definition}
\label{equiv-definition-siblings-geometric}
Soit $k$ un corps. Soient $X$ et $Y$ des schémas de type fini sur $k$.
Rappelons que
$D^b_{\textit{Coh}}(\mathcal{O}_X) = D^b(\textit{Coh}(\mathcal{O}_X))$
d'après Catégories dérivées des schémas, proposition \ref{perfect-proposition-DCoh}.
Nous disons que deux foncteurs exacts $k$-linéaires
$$
F, F' :
D^b_{\textit{Coh}}(\mathcal{O}_X) = D^b(\textit{Coh}(\mathcal{O}_X))
\longrightarrow
D^b_{\textit{Coh}}(\mathcal{O}_Y)
$$
sont dits {\it apparentés}, ou que $F'$ est {\it apparenté} à $F$, si $F$ et $F'$
sont apparentés au sens de la définition \ref{equiv-definition-siblings},
la catégorie abélienne étant $\textit{Coh}(\mathcal{O}_X)$.
Si $X$ est régulier, alors
$D_{perf}(\mathcal{O}_X) = D^b_{\textit{Coh}}(\mathcal{O}_X)$ d'après
Catégories dérivées des schémas, lemme \ref{perfect-lemma-perfect-on-noetherian},
et nous employons la même terminologie pour les foncteurs exacts $k$-linéaires
$F, F' : D_{perf}(\mathcal{O}_X) \to D_{perf}(\mathcal{O}_Y)$.
\end{definition}

\begin{lemma}
\label{equiv-lemma-exact-functor-preserving-Coh}
Soit $k$ un corps. Soient $X$ et $Y$ des schémas de type fini sur $k$, avec
$X$ séparé. Soit
$F : D^b_{\textit{Coh}}(\mathcal{O}_X) \to D^b_{\textit{Coh}}(\mathcal{O}_Y)$
un foncteur exact $k$-linéaire qui envoie
$\textit{Coh}(\mathcal{O}_X) \subset D^b_{\textit{Coh}}(\mathcal{O}_X)$
dans
$\textit{Coh}(\mathcal{O}_Y) \subset D^b_{\textit{Coh}}(\mathcal{O}_Y)$.
Alors il existe un foncteur de Fourier-Mukai
$F' : D^b_{\textit{Coh}}(\mathcal{O}_X) \to D^b_{\textit{Coh}}(\mathcal{O}_Y)$
dont le noyau est un $\mathcal{O}_{X \times Y}$-module cohérent $\mathcal{K}$,
plat sur $X$ et à support fini sur $Y$ ; le foncteur ainsi obtenu est apparenté à $F$.
\end{lemma}

\begin{proof}
Notons $H : \textit{Coh}(\mathcal{O}_X) \to \textit{Coh}(\mathcal{O}_Y)$
la restriction de $F$. Comme $F$ est un foncteur exact de catégories
triangulées, $H$ est un foncteur exact de catégories abéliennes.
Bien entendu, $H$ est $k$-linéaire puisque $F$ l'est. D'après
Foncteurs et morphismes, lemme \ref{functors-lemma-functor-coherent-over-field},
on obtient un $\mathcal{O}_{X \times Y}$-module cohérent
$\mathcal{K}$ qui est plat sur $X$ et à support fini sur $Y$.
Soit $F'$ le foncteur de Fourier-Mukai défini à l'aide de $\mathcal{K}$,
de sorte que $F'$ se restreint à $H$ sur $ \textit{Coh}(\mathcal{O}_X)$.
Le foncteur $F'$ envoie $D^b_{\textit{Coh}}(\mathcal{O}_X)$
dans $D^b_{\textit{Coh}}(\mathcal{O}_Y)$ d'après le
lemme \ref{equiv-lemma-fourier-mukai-Coh}.
Observons que $F$ et $F'$ satisfont les première et deuxième
conditions du lemme \ref{equiv-lemma-sibling-fully-faithful} et sont donc apparentés.
\end{proof}

\begin{remark}
\label{equiv-remark-difficult}
Si $F, F' : D^b_{\textit{Coh}}(\mathcal{O}_X) \to \mathcal{D}$ sont apparentés, si $F$
est pleinement fidèle et si $X$ est réduit et projectif sur $k$, alors
$F \cong F'$ ; cela résulte de la
proposition \ref{equiv-proposition-siblings-isomorphic} par l'argument
donné dans la démonstration du théorème \ref{equiv-theorem-fully-faithful}.
Cependant, en général, nous ne savons pas si des foncteurs apparentés sont isomorphes.
Même dans la situation du lemme \ref{equiv-lemma-exact-functor-preserving-Coh},
il semble difficile de prouver que les foncteurs apparentés $F$ et $F'$
sont des foncteurs isomorphes. Si $X$ est lisse et propre sur $k$
et si $F$ est pleinement fidèle, alors $F \cong F'$, comme le montre
\cite{Noah}.
Si vous disposez d'une démonstration ou d'un contre-exemple dans des situations plus générales,
veuillez écrire à
\href{mailto:stacks.project@gmail.com}{stacks.project@gmail.com}.
\end{remark}

\begin{lemma}
\label{equiv-lemma-two-functors-pre}
Soit $k$ un corps. Soient $X$ et $Y$ des schémas propres sur $k$. Supposons
que $X$ soit régulier. Soient
$F, G : D_{perf}(\mathcal{O}_X) \to D_{perf}(\mathcal{O}_Y)$
des foncteurs exacts $k$-linéaires tels que
\begin{enumerate}
\item $F(\mathcal{F}) \cong G(\mathcal{F})$ pour tout
$\mathcal{O}_X$-module cohérent $\mathcal{F}$ tel que $\dim(\text{Supp}(\mathcal{F})) = 0$,
\item $F$ soit pleinement fidèle.
\end{enumerate}
Alors l'image essentielle de $G$ est contenue dans l'image essentielle
de $F$.
\end{lemma}

\begin{proof}
Rappelons que $F$ et $G$ possèdent chacun un adjoint à gauche et un adjoint à droite, voir le
lemme \ref{equiv-lemma-always-right-adjoints}. En particulier,
l'image essentielle $\mathcal{A} \subset D_{perf}(\mathcal{O}_Y)$ de $F$
satisfait les conditions équivalentes de
Catégories dérivées, lemme \ref{derived-lemma-right-adjoint}.
Nous affirmons que $G$ se factorise par $\mathcal{A}$.
Comme $\mathcal{A} = {}^\perp(\mathcal{A}^\perp)$ d'après
Catégories dérivées, lemme \ref{derived-lemma-right-adjoint},
il suffit de montrer que $\Hom_Y(G(M), N) = 0$ pour
tout $M$ dans $D_{perf}(\mathcal{O}_X)$ et $N \in \mathcal{A}^\perp$.
On a
$$
\Hom_Y(G(M), N) = \Hom_X(M, G_r(N))
$$
où $G_r$ est l'adjoint à droite de $G$. Il suffit donc de prouver
que $G_r(N) = 0$. Comme
$G(\mathcal{F}) \cong F(\mathcal{F})$ pour $\mathcal{F}$ comme en (1),
on voit que
$$
\Hom_X(\mathcal{F}, G_r(N)) =
\Hom_Y(G(\mathcal{F}), N) =
\Hom_Y(F(\mathcal{F}), N) = 0
$$
puisque $N$ appartient à l'orthogonal à droite de l'image essentielle $\mathcal{A}$ de $F$.
Bien entendu, le groupe $\Hom_X(\mathcal{F}, G_r(N)[i])$ s'annule encore
pour tout $i \in \mathbf{Z}$. Ainsi, $G_r(N) = 0$ d'après le
lemme \ref{equiv-lemma-orthogonal-point-sheaf}, ce qui conclut.
\end{proof}

\begin{lemma}
\label{equiv-lemma-noah}
\begin{reference}
Courriel de Noah Olander du 8 juin 2020
\end{reference}
Soit $k$ un corps. Soit $X$ un schéma propre sur $k$ qui est régulier.
Soit $F : D_{perf}(\mathcal{O}_X) \to D_{perf}(\mathcal{O}_X)$
un foncteur exact $k$-linéaire. Supposons que, pour tout
$\mathcal{O}_X$-module cohérent $\mathcal{F}$ tel que $\dim(\text{Supp}(\mathcal{F})) = 0$,
il existe un isomorphisme $\mathcal{F} \cong F(\mathcal{F})$.
Alors il existe un automorphisme $f : X \to X$ sur $k$
qui induit l'identité
sur l'espace topologique sous-jacent\footnote{Cela force souvent $f$
à être l'identité, voir Variétés, lemme \ref{varieties-lemma-automorphism}.},
et un $\mathcal{O}_X$-module inversible $\mathcal{L}$
tels que $F$ et $F'(M) = f^*M \otimes_{\mathcal{O}_X}^\mathbf{L} \mathcal{L}$
soient apparentés.
\end{lemma}

\begin{proof}
D'après le lemme \ref{equiv-lemma-noah-pre}, le foncteur $F$ est pleinement fidèle.
D'après le lemme \ref{equiv-lemma-two-functors-pre}, l'image essentielle du
foncteur identique est contenue dans l'image essentielle de $F$, c'est-à-dire
que $F$ est essentiellement surjectif. Ainsi, $F$ est une équivalence.
Observons que le quasi-inverse $F^{-1}$ satisfait les mêmes hypothèses
que $F$.

\medskip\noindent
Soit $M \in D_{perf}(\mathcal{O}_X)$, et supposons $H^i(M) = 0$ pour $i > b$.
Comme $F$ est pleinement fidèle, on voit que
$$
\Hom_X(M, \mathcal{O}_x[-i]) =
\Hom_X(F(M), F(\mathcal{O}_x)[-i]) \cong
\Hom_X(F(M), \mathcal{O}_x[-i])
$$
pour tout $i \in \mathbf{Z}$ et tout point fermé $x$ de $X$.
Le lemme \ref{equiv-lemma-hom-into-point-sheaf} montre donc que les faisceaux de cohomologie de $F(M)$
sont nuls aux degrés $> b$.

\medskip\noindent
Soit $\mathcal{F}$ un $\mathcal{O}_X$-module cohérent. D'après
ce qui précède, les faisceaux de cohomologie non nuls de $F(\mathcal{F})$
ne se trouvent qu'aux degrés $\leq 0$.
Posons $\mathcal{G} = H^0(F(\mathcal{F}))$. Choisissons un triangle
distingué
$$
K \to F(\mathcal{F}) \to \mathcal{G} \to K[1]
$$
Alors les faisceaux de cohomologie non nuls de $K$ ne se trouvent qu'aux
degrés $\leq -1$.
En appliquant $F^{-1}$, on obtient un triangle distingué
$$
F^{-1}(K) \to \mathcal{F} \to F^{-1}(\mathcal{G}) \to F^{-1}(K')[1]
$$
Comme les faisceaux de cohomologie non nuls de $F^{-1}(K)$ ne se trouvent qu'aux
degrés $\leq -1$ (par le paragraphe précédent appliqué à $F^{-1}$),
on voit que la flèche $F^{-1}(K) \to \mathcal{F}$ est nulle
(Catégories dérivées, lemme \ref{derived-lemma-negative-exts}).
Ainsi, $K \to F(\mathcal{F})$ est nul, ce qui implique
que $F(\mathcal{F}) = \mathcal{G}$ par le choix du
premier triangle distingué.

\medskip\noindent
Le paragraphe précédent montre que $F$ préserve
$\textit{Coh}(\mathcal{O}_X)$ et définit même une équivalence
$H : \textit{Coh}(\mathcal{O}_X) \to \textit{Coh}(\mathcal{O}_X)$.
D'après Foncteurs et morphismes, lemme
\ref{functors-lemma-equivalence-coherent-over-field},
on obtient un automorphisme $f : X \to X$ sur $k$
et un $\mathcal{O}_X$-module inversible $\mathcal{L}$
tels que $H(\mathcal{F}) = f^*\mathcal{F} \otimes \mathcal{L}$.
Posons $F'(M) = f^*M \otimes_{\mathcal{O}_X}^\mathbf{L} \mathcal{L}$.
Le lemme \ref{equiv-lemma-sibling-fully-faithful}
montre que $F$ et $F'$ sont apparentés.
Pour voir que $f$ induit l'identité sur l'espace topologique
sous-jacent de $X$, on utilise $F(\mathcal{O}_x) \cong \mathcal{O}_x$
et le fait que le support de $\mathcal{O}_x$ est $\{x\}$.
Cela achève la démonstration.
\end{proof}

\begin{lemma}
\label{equiv-lemma-two-functors}
Soit $k$ un corps. Soient $X$ et $Y$ des schémas propres sur $k$.
Supposons que $X$ soit régulier.
Soient $F, G : D_{perf}(\mathcal{O}_X) \to D_{perf}(\mathcal{O}_Y)$
des foncteurs exacts $k$-linéaires tels que
\begin{enumerate}
\item $F(\mathcal{F}) \cong G(\mathcal{F})$ pour tout
$\mathcal{O}_X$-module cohérent $\mathcal{F}$ tel que $\dim(\text{Supp}(\mathcal{F})) = 0$,
\item $F$ soit pleinement fidèle, et
\item $G$ soit un foncteur de Fourier-Mukai dont le noyau appartient à
$D_{perf}(\mathcal{O}_{X \times Y})$.
\end{enumerate}
Alors il existe un foncteur de Fourier-Mukai
$F' : D_{perf}(\mathcal{O}_X) \to D_{perf}(\mathcal{O}_Y)$
dont le noyau appartient à $D_{perf}(\mathcal{O}_{X \times Y})$
et pour lequel $F$ et $F'$ sont apparentés.
\end{lemma}

\begin{proof}
L'image essentielle de $G$ est contenue dans l'image essentielle
de $F$ d'après le lemme \ref{equiv-lemma-two-functors-pre}.
Considérons le foncteur $H = F^{-1} \circ G$,
qui a un sens puisque $F$ est pleinement fidèle.
Le lemme \ref{equiv-lemma-noah} fournit un automorphisme $f : X \to X$
et un $\mathcal{O}_X$-module inversible $\mathcal{L}$ tels que
le foncteur $H' : K \mapsto f^*K \otimes \mathcal{L}$
soit apparenté à $H$. En particulier,
$H$ est une auto-équivalence d'après le lemme \ref{equiv-lemma-sibling-faithful},
et $H$ induit une auto-équivalence de
$\textit{Coh}(\mathcal{O}_X)$ (puisque cela vaut pour le foncteur $H'$ qui lui est apparenté).
Ainsi, les quasi-inverses $H^{-1}$ et $(H')^{-1}$ existent et sont apparentés
(nous omettons un petit détail), et $(H')^{-1}$ envoie $M$ sur
$(f^{-1})^*(M \otimes_{\mathcal{O}_X}^\mathbf{L} \mathcal{L}^{\otimes -1})$,
qui est un foncteur de Fourier-Mukai (nous omettons les détails).
Alors, bien entendu, $F = G \circ H^{-1}$ est apparenté à
$G \circ (H')^{-1}$. Comme les composés de foncteurs de Fourier-Mukai
sont de Fourier-Mukai d'après le
lemme \ref{equiv-lemma-compose-fourier-mukai},
on conclut.
\end{proof}





\section{Foncteurs pleinement fidèles}
\label{equiv-section-fully-faithful}

\noindent
Notre objectif est de prouver que les foncteurs pleinement fidèles entre catégories dérivées
sont apparentés à des foncteurs de Fourier-Mukai, en suivant
\cite{Orlov-K3} et \cite{Ballard}.

\begin{situation}
\label{equiv-situation-fully-faithful}
Ici, $k$ est un corps. Nous avons des schémas propres et lisses $X$ et $Y$ sur $k$.
Nous avons un foncteur exact, $k$-linéaire et pleinement fidèle
$F : D_{perf}(\mathcal{O}_X) \to D_{perf}(\mathcal{O}_Y)$.
\end{situation}

\noindent
Avant de poursuivre, il est utile de lire au moins une partie de
Catégories dérivées, section \ref{derived-section-postnikov}.

\medskip\noindent
Rappelons que $X$ est régulier et possède donc la propriété de résolution
(Variétés, lemme \ref{varieties-lemma-smooth-regular} et
Catégories dérivées des schémas, lemme
\ref{perfect-lemma-regular-resolution-property}). Ainsi,
sur $X \times X$, nous pouvons choisir une résolution
$$
\ldots \to
\mathcal{E}_2 \boxtimes \mathcal{G}_2 \to
\mathcal{E}_1 \boxtimes \mathcal{G}_1 \to
\mathcal{E}_0 \boxtimes \mathcal{G}_0 \to
\mathcal{O}_\Delta \to 0
$$
où chaque $\mathcal{E}_i$ et chaque $\mathcal{G}_i$ est un
$\mathcal{O}_X$-module localement libre de type fini, voir le
lemme \ref{equiv-lemma-diagonal-resolution}.
En utilisant le complexe
\begin{equation}
\label{equiv-equation-original-complex}
\ldots \to
\mathcal{E}_2 \boxtimes \mathcal{G}_2 \to
\mathcal{E}_1 \boxtimes \mathcal{G}_1 \to
\mathcal{E}_0 \boxtimes \mathcal{G}_0
\end{equation}
dans $D_{perf}(\mathcal{O}_{X \times X})$, comme dans
Catégories dérivées, exemple \ref{derived-example-key-postnikov},
si, pour tout $n$, nous posons
$$
M_n = (\mathcal{E}_n \boxtimes \mathcal{G}_n \to \ldots \to
\mathcal{E}_0 \boxtimes \mathcal{G}_0)[-n]
$$
nous obtenons un système de Postnikov infini pour le complexe
(\ref{equiv-equation-original-complex}). Cela signifie que
les morphismes $M_0 \to M_1[1] \to M_2[2] \to \ldots$, ainsi que
$M_n \to \mathcal{E}_n \boxtimes \mathcal{G}_n$ et
$\mathcal{E}_n \boxtimes \mathcal{G}_n \to M_{n - 1}$,
satisfont à certaines conditions consignées dans
Catégories dérivées, définition \ref{derived-definition-postnikov-system}.
Posons
$$
\mathcal{F}_n = \Ker(\mathcal{E}_n \boxtimes \mathcal{G}_n \to
\mathcal{E}_{n - 1} \boxtimes \mathcal{G}_{n - 1})
$$
Observons que, puisque $\mathcal{O}_\Delta$ est plat sur $X$ via $\text{pr}_1$,
il en va de même de $\mathcal{F}_n$ pour tout $n$ (c'est une observation
commode, mais non essentielle). Nous avons
$$
H^q(M_n[n]) = \left\{
\begin{matrix}
\mathcal{O}_\Delta & \text{si} & q = 0 \\
\mathcal{F}_n & \text{si} & q = -n \\
0 & \text{si} & q \not = 0, -n
\end{matrix}
\right.
$$
Ainsi, pour $n \geq \dim(X \times X)$, nous avons
$$
M_n[n] \cong \mathcal{O}_\Delta \oplus \mathcal{F}_n[n]
$$
dans $D_{perf}(\mathcal{O}_{X \times X})$, d'après le

lemme \ref{equiv-lemma-split-complex-regular}.
\medskip\noindent
Nous nous intéressons au complexe
\begin{equation}
\label{equiv-equation-complex}
\ldots \to
\mathcal{E}_2 \boxtimes F(\mathcal{G}_2) \to
\mathcal{E}_1 \boxtimes F(\mathcal{G}_1) \to
\mathcal{E}_0 \boxtimes F(\mathcal{G}_0)
\end{equation}
dans $D_{perf}(\mathcal{O}_{X \times Y})$,
car la « totalisation » de ce complexe devrait
nous donner le noyau du foncteur de Fourier-Mukai que nous cherchons à construire.
Pour tous $i, j \geq 0$, nous avons
\begin{align*}
\Ext^q_{X \times Y}(\mathcal{E}_i \boxtimes F(\mathcal{G}_i),
\mathcal{E}_j \boxtimes F(\mathcal{G}_j))
& =
\bigoplus\nolimits_p
\Ext^{q + p}_X(\mathcal{E}_i, \mathcal{E}_j) \otimes_k
\Ext^{-p}_Y(F(\mathcal{G}_i), F(\mathcal{G}_j)) \\
& =
\bigoplus\nolimits_p
\Ext^{q + p}_X(\mathcal{E}_i, \mathcal{E}_j) \otimes_k
\Ext^{-p}_X(\mathcal{G}_i, \mathcal{G}_j)
\end{align*}
La seconde égalité vaut parce que $F$ est pleinement fidèle, et la première
résulte de Catégories dérivées des schémas, lemme
\ref{perfect-lemma-kunneth-Ext}.
On obtient que ces $\Ext^q$ sont nuls pour $q < 0$.
Ainsi, d'après Catégories dérivées, lemme
\ref{derived-lemma-existence-postnikov-system},
nous pouvons construire un système de Postnikov infini $K_0, K_1, K_2, \ldots$
dans $D_{perf}(\mathcal{O}_{X \times Y})$ pour le complexe
(\ref{equiv-equation-complex}).
Parallèlement à ce qui se passe pour $M_0, M_1, M_2, \ldots$, cela signifie
que nous obtenons des morphismes
$K_0 \to K_1[1] \to K_2[2] \to \ldots$, ainsi que
$K_n \to \mathcal{E}_n \boxtimes F(\mathcal{G}_n)$ et
$\mathcal{E}_n \boxtimes F(\mathcal{G}_n) \to K_{n - 1}$
dans $D_{perf}(\mathcal{O}_{X \times Y})$,
qui satisfont à certaines conditions consignées dans

Catégories dérivées, définition \ref{derived-definition-postnikov-system}.
\medskip\noindent
Soit $\mathcal{F}$ un $\mathcal{O}_X$-module cohérent dont le support
est un ensemble fini de points, c'est-à-dire tel que $\dim(\text{Supp}(\mathcal{F})) = 0$.
Considérons le foncteur exact entre catégories triangulées
$$
D_{perf}(\mathcal{O}_{X \times Y})
\longrightarrow
D_{perf}(\mathcal{O}_Y),\quad
N \longmapsto R\text{pr}_{2, *}(\text{pr}_1^*\mathcal{F}
\otimes^\mathbf{L}_{\mathcal{O}_{X \times Y}} N)
$$
Il s'ensuit que les objets $R\text{pr}_{2, *}(\text{pr}_1^*\mathcal{F}
\otimes^\mathbf{L}_{\mathcal{O}_{X \times Y}} K_i)$
forment un système de Postnikov pour le complexe de
$D_{perf}(\mathcal{O}_Y)$ dont les termes sont
$$
R\text{pr}_{2, *}(
(\mathcal{F} \otimes \mathcal{E}_i) \boxtimes F(\mathcal{G}_i)) =
\Gamma(X, \mathcal{F} \otimes \mathcal{E}_i) \otimes_k F(\mathcal{G}_i) =
F(\Gamma(X, \mathcal{F} \otimes \mathcal{E}_i) \otimes_k \mathcal{G}_i)
$$
Nous avons utilisé ici que $\mathcal{F} \otimes \mathcal{E}_i$ a une
cohomologie supérieure nulle, puisque son support est de dimension $0$.
D'autre part, en appliquant le foncteur exact
$$
D_{perf}(\mathcal{O}_{X \times X})
\longrightarrow
D_{perf}(\mathcal{O}_Y),\quad
N \longmapsto F(R\text{pr}_{2, *}(\text{pr}_1^*\mathcal{F}
\otimes^\mathbf{L}_{\mathcal{O}_{X \times X}} N))
$$
nous constatons que les objets
$F(R\text{pr}_{2, *}(\text{pr}_1^*\mathcal{F}
\otimes^\mathbf{L}_{\mathcal{O}_{X \times X}} M_n))$
forment un second système de Postnikov infini
pour le complexe de $D_{perf}(\mathcal{O}_Y)$ dont les termes sont
$$
F(R\text{pr}_{2, *}(
(\mathcal{F} \otimes \mathcal{E}_i) \boxtimes \mathcal{G}_i)) =
F(\Gamma(X, \mathcal{F} \otimes \mathcal{E}_i) \otimes_k \mathcal{G}_i)
$$
C'est le même complexe que précédemment ! Par unicité des systèmes de Postnikov
(Catégories dérivées, lemme \ref{derived-lemma-existence-postnikov-system}),
laquelle s'applique puisque
$$
\Ext^q_Y(
F(\Gamma(X, \mathcal{F} \otimes \mathcal{E}_i) \otimes_k \mathcal{G}_i),
F(\Gamma(X, \mathcal{F} \otimes \mathcal{E}_j) \otimes_k \mathcal{G}_j)) = 0,
\quad q < 0
$$
car $F$ est pleinement fidèle, nous obtenons un système d'isomorphismes
$$
F(R\text{pr}_{2, *}(\text{pr}_1^*\mathcal{F}
\otimes^\mathbf{L}_{\mathcal{O}_{X \times X}} M_n[n]))
\cong
R\text{pr}_{2, *}(\text{pr}_1^*\mathcal{F}
\otimes^\mathbf{L}_{\mathcal{O}_{X \times Y}} K_n[n])
$$
dans $D_{perf}(\mathcal{O}_Y)$, compatible avec les morphismes de
$D_{perf}(\mathcal{O}_Y)$ induits par les morphismes
$$
M_{n - 1}[n - 1] \to M_n[n]
\quad\text{et}\quad
K_{n - 1}[n - 1] \to K_n[n]
$$
$$
M_n \to \mathcal{E}_n \boxtimes \mathcal{G}_n
\quad\text{et}\quad
K_n \to \mathcal{E}_n \boxtimes F(\mathcal{G}_n)
$$
$$
\mathcal{E}_n \boxtimes \mathcal{G}_n \to M_{n - 1}
\quad\text{et}\quad
\mathcal{E}_n \boxtimes F(\mathcal{G}_n) \to K_{n - 1}
$$
qui font partie de la structure des systèmes de Postnikov.
Pour $n$ suffisamment grand, nous obtenons une décomposition en somme directe
$$
F(R\text{pr}_{2, *}(\text{pr}_1^*\mathcal{F}
\otimes^\mathbf{L}_{\mathcal{O}_{X \times X}} M_n[n]))
=
F(\mathcal{F}) \oplus
F(R\text{pr}_{2, *}(
\text{pr}_1^*\mathcal{F} \otimes_{\mathcal{O}_{X \times Y}} \mathcal{F}_n
))[n]
$$
qui correspond à la décomposition en somme directe de $M_n$ construite ci-dessus
(nous utilisons la platitude de $\mathcal{F}_n$ sur $X$ via $\text{pr}_1$
pour écrire un produit tensoriel ordinaire dans la formule ci-dessus, mais cela
n'est pas essentiel à l'argument).
D'après le lemme \ref{equiv-lemma-boundedness}, il existe un entier $m \geq 0$
tel que le premier facteur de cette décomposition en somme directe n'ait de
faisceaux de cohomologie non nuls que dans l'intervalle $[-m, m]$, tandis que le
second facteur n'a de faisceaux de cohomologie non nuls que dans l'intervalle
$[-m - n, m + \dim(X) - n]$.
On en conclut que le système $K_0 \to K_1[1] \to K_2[2] \to \ldots$
de $D_{perf}(\mathcal{O}_{X \times Y})$ satisfait aux hypothèses du
lemme \ref{equiv-lemma-bounded-fibres}, quitte à remplacer $m$ par un entier plus grand. On peut donc écrire
$$
K_n[n] = K \oplus C_n
$$
pour $n \gg 0$, de manière compatible avec les morphismes de transition, et
$C_n$ n'a de faisceaux de cohomologie non nuls que dans l'intervalle $[-m - n, m - n]$.
Notons $G$ le foncteur de Fourier-Mukai associé à $K$.
En regroupant ce qui précède, nous obtenons
$$
\begin{matrix}
G(\mathcal{F}) \oplus
R\text{pr}_{2, *}(
\text{pr}_1^*\mathcal{F} \otimes_{\mathcal{O}_{X \times Y}}^\mathbf{L} C_n)
\cong \\
R\text{pr}_{2, *}(\text{pr}_1^*\mathcal{F}
\otimes^\mathbf{L}_{\mathcal{O}_{X \times Y}} K_n[n]) \cong \\
F(R\text{pr}_{2, *}(\text{pr}_1^*\mathcal{F}
\otimes^\mathbf{L}_{\mathcal{O}_{X \times X}} M_n[n]))
\cong \\
F(\mathcal{F}) \oplus
F(R\text{pr}_{2, *}(
\text{pr}_1^*\mathcal{F} \otimes_{\mathcal{O}_{X \times Y}} \mathcal{F}_n
))[n]
\end{matrix}
$$
En examinant les degrés où vivent ces objets, on conclut que, pour $n \gg m$,
nous obtenons un isomorphisme
$$
F(\mathcal{F}) \cong G(\mathcal{F})
$$
De plus, cela vaut pour tout faisceau cohérent $\mathcal{F}$ sur $X$
dont le support est de dimension $0$.

\begin{lemma}
\label{equiv-lemma-fully-faithful}
Soit $k$ un corps. Soient $X$ et $Y$ des schémas propres et lisses sur $k$.
Étant donné un foncteur exact, $k$-linéaire et pleinement fidèle
$F : D_{perf}(\mathcal{O}_X) \to D_{perf}(\mathcal{O}_Y)$,
il existe un foncteur de Fourier-Mukai
$F' : D_{perf}(\mathcal{O}_X) \to D_{perf}(\mathcal{O}_Y)$ dont le noyau
appartient à $D_{perf}(\mathcal{O}_{X \times Y})$ et qui est apparenté à $F$.
\end{lemma}

\begin{proof}
Appliquons le lemme \ref{equiv-lemma-two-functors} à $F$ et au foncteur
$G$ construit ci-dessus.
\end{proof}

\noindent
Le théorème suivant reste vrai sans supposer que $X$ soit projectif,
voir \cite{Noah}.

\begin{theorem}[Orlov]
\label{equiv-theorem-fully-faithful}
\begin{reference}
\cite[Théorème 2.2]{Orlov-K3} ; cela est démontré dans \cite{Noah}
sans l'hypothèse que $X$ soit projectif
\end{reference}
Soit $k$ un corps. Soient $X$ et $Y$ des schémas propres et lisses sur $k$,
avec $X$ projectif sur $k$. Tout foncteur exact, $k$-linéaire et pleinement fidèle
$F : D_{perf}(\mathcal{O}_X) \to D_{perf}(\mathcal{O}_Y)$
est un foncteur de Fourier-Mukai associé à un noyau de
$D_{perf}(\mathcal{O}_{X \times Y})$.
\end{theorem}

\begin{proof}
Soit $F'$ le foncteur de Fourier-Mukai qui est apparenté à $F$
comme dans le lemme \ref{equiv-lemma-fully-faithful}.
D'après la proposition \ref{equiv-proposition-siblings-isomorphic}, nous avons $F \cong F'$
pourvu que nous montrions que $\textit{Coh}(\mathcal{O}_X)$ possède assez
d'objets négatifs. Toutefois, si $X = \Spec(k)$, par exemple,
ce n'est pas le cas. Nous décomposons donc d'abord $X = \coprod X_i$
en ses composantes connexes (et irréductibles), puis nous
montrons qu'il suffit de prouver le résultat pour chacun des
foncteurs composés (pleinement fidèles)
$$
F_i :
D_{perf}(\mathcal{O}_{X_i}) \to
D_{perf}(\mathcal{O}_X) \to
D_{perf}(\mathcal{O}_Y)
$$
Nous omettons les détails. Nous pouvons donc supposer $X$ irréductible.

\medskip\noindent
Le cas $\dim(X) = 0$. Ici, $X$ est le spectre d'une extension finie (séparable)
$k'/k$, et $D_{perf}(\mathcal{O}_X)$ est donc équivalente à la catégorie
des espaces vectoriels gradués sur $k'$ dans laquelle
$\mathcal{O}_X$ correspond à
l'espace vectoriel trivial de dimension $1$, placé en degré $0$.
Il est immédiat que deux foncteurs apparentés
$F, F' : D_{perf}(\mathcal{O}_X) \to D_{perf}(\mathcal{O}_Y)$
sont isomorphes. En effet, nous disposons d'un isomorphisme
$F(\mathcal{O}_X) \cong F'(\mathcal{O}_X)$
compatible avec l'action de la $k$-algèbre
$k' = \text{End}_{D_{perf}(\mathcal{O}_X)}(\mathcal{O}_X)$,
qui se prolonge canoniquement en un isomorphisme sur tout espace vectoriel gradué sur $k'$.

\medskip\noindent
Le cas $\dim(X) > 0$. Ici, $X$ est une variété projective lisse
de dimension $> 1$. Soit $\mathcal{F}$ un
$\mathcal{O}_X$-module cohérent.
Nous devons montrer qu'il existe un module cohérent $\mathcal{N}$ tel que
\begin{enumerate}
\item il existe une surjection $\mathcal{N} \to \mathcal{F}$, et
\item $\Hom(\mathcal{F}, \mathcal{N}) = 0$.
\end{enumerate}
Choisissons un $\mathcal{O}_X$-module inversible ample $\mathcal{L}$.
Nous affirmons que $\mathcal{N} = (\mathcal{L}^{\otimes n})^{\oplus r}$
convient pour $n \ll 0$ et $r$ assez grand.
La condition (1) résulte de
Propriétés, proposition \ref{properties-proposition-characterize-ample}.
Enfin, nous avons
$$
\Hom(\mathcal{F}, \mathcal{L}^{\otimes n}) =
H^0(X, \SheafHom(\mathcal{F}, \mathcal{L}^{\otimes n})) =
H^0(X, \SheafHom(\mathcal{F}, \mathcal{O}_X) \otimes \mathcal{L}^{\otimes n})
$$
Comme le dual $\SheafHom(\mathcal{F}, \mathcal{O}_X)$ est sans torsion, ce groupe
est nul pour $n \ll 0$ d'après Variétés, lemme
\ref{varieties-lemma-vanishin-h0-negative}. Cela achève la démonstration.
\end{proof}

\begin{proposition}
\label{equiv-proposition-equivalence}
Soit $k$ un corps. Soient $X$ et $Y$ des schémas propres et lisses sur $k$.
Si $F : D_{perf}(\mathcal{O}_X) \to D_{perf}(\mathcal{O}_Y)$
est une équivalence exacte $k$-linéaire de catégories triangulées, alors
il existe un foncteur de Fourier-Mukai
$F' : D_{perf}(\mathcal{O}_X) \to D_{perf}(\mathcal{O}_Y)$ dont le
noyau appartient à $D_{perf}(\mathcal{O}_{X \times Y})$,
qui est une équivalence et est apparenté à $F$.
\end{proposition}

\begin{proof}
Le foncteur $F'$ du lemme \ref{equiv-lemma-fully-faithful}
est une équivalence d'après le lemme \ref{equiv-lemma-sibling-faithful}.
\end{proof}

\begin{lemma}
\label{equiv-lemma-uniqueness}
Soit $k$ un corps. Soit $X$ un schéma propre et lisse sur $k$.
Soit $K \in D_{perf}(\mathcal{O}_{X \times X})$. Si le foncteur de Fourier-Mukai
$\Phi_K : D_{perf}(\mathcal{O}_X) \to D_{perf}(\mathcal{O}_X)$
est isomorphe au foncteur identique, alors
$K \cong \Delta_*\mathcal{O}_X$ dans $_{perf}(\mathcal{O}_{X \times X})$.
\end{lemma}

\begin{proof}
Soit $i$ le plus petit entier tel que le faisceau de cohomologie $H^i(K)$ soit
non nul. Soient $\mathcal{E}$ et $\mathcal{G}$ des $\mathcal{O}_X$-modules
localement libres de type fini. Alors
\begin{align*}
H^i(X \times X, K \otimes_{\mathcal{O}_{X \times X}}^\mathbf{L}
(\mathcal{E} \boxtimes \mathcal{G}))
& =
H^i(X, R\text{pr}_{2, *}(K \otimes_{\mathcal{O}_{X \times X}}^\mathbf{L}
(\mathcal{E} \boxtimes \mathcal{G}))) \\
& =
H^i(X, \Phi_K(\mathcal{E}) \otimes_{\mathcal{O}_X}^\mathbf{L} \mathcal{G}) \\
& \cong
H^i(X, \mathcal{E} \otimes \mathcal{G})
\end{align*}
ce qui est nul si $i < 0$. D'autre part, nous pouvons choisir
$\mathcal{E}$ et $\mathcal{G}$ de sorte qu'il existe une surjection
$\mathcal{E}^\vee \boxtimes \mathcal{G}^\vee \to H^i(K)$
d'après le lemme \ref{equiv-lemma-on-product}.
Dans ce cas, le membre de gauche des égalités n'est pas nul.
On conclut donc que $H^i(K) = 0$ pour $i < 0$.

\medskip\noindent
Soit $i$ le plus grand entier tel que $H^i(K)$ soit non nul.
Le même argument, avec $\mathcal{E}$ et $\mathcal{G}$ à support de dimension $0$,
montre que $i \leq 0$.
On conclut donc que $K$ est donné par un unique
$\mathcal{O}_{X \times X}$-module cohérent $\mathcal{K}$ placé en degré $0$.

\medskip\noindent
Puisque $R\text{pr}_{2, *}(\text{pr}_1^*\mathcal{F} \otimes \mathcal{K})$
est $\mathcal{F}$, en prenant $\mathcal{F}$ à support en des points fermés,
nous voyons que le support de $\mathcal{K}$ est fini sur $X$ via
$\text{pr}_2$. Comme $R\text{pr}_{2, *}(\mathcal{K}) \cong \mathcal{O}_X$,
on conclut, d'après Foncteurs et morphismes, lemme
\ref{functors-lemma-pushforward-invertible-pre},
que $\mathcal{K} = s_*\mathcal{O}_X$ pour une section $s : X \to X \times X$
de la seconde projection. Alors $\Phi_K(M) = f^*M$, où
$f = \text{pr}_1 \circ s$, et cela ne peut se produire que si $s$
est le morphisme diagonal, comme souhaité.
\end{proof}








\section{Une catégorie de noyaux de Fourier-Mukai}
\label{equiv-section-category-Fourier-Mukai-kernels}

\noindent
Soit $S$ un schéma. Nous affirmons qu'il existe une catégorie
dont
\begin{enumerate}
\item les objets sont les schémas propres et lisses sur $S$ ;
\item les morphismes de $X$ vers $Y$ sont les classes d'isomorphisme
des objets de $D_{perf}(\mathcal{O}_{X \times_S Y})$ ;
\item la composée de la classe d'isomorphisme de
$K \in D_{perf}(\mathcal{O}_{X \times_S Y})$
et de la classe d'isomorphisme de $K'$ dans $D_{perf}(\mathcal{O}_{Y \times_S Z})$
est la classe d'isomorphisme de
$$
R\text{pr}_{13, *}(
L\text{pr}_{12}^*K
\otimes_{\mathcal{O}_{X \times_S Y \times_S Z}}^\mathbf{L}
L\text{pr}_{23}^*K')
$$
qui appartient à $D_{perf}(\mathcal{O}_{X \times_S Z})$ d'après
Catégories dérivées des schémas, lemme
\ref{perfect-lemma-flat-proper-perfect-direct-image-general} ;
\item le morphisme identique de $X$ vers $X$ est la
classe d'isomorphisme de $\Delta_{X/S, *}\mathcal{O}_X$,
qui appartient à $D_{perf}(\mathcal{O}_{X \times_S X})$
d'après Compléments sur les morphismes, lemme
\ref{more-morphisms-lemma-perfect-closed-immersion-perfect-direct-image},
et le fait que $\Delta_{X/S}$ est un morphisme parfait d'après
Diviseurs, lemme
\ref{divisors-lemma-immersion-smooth-into-smooth-regular-immersion} et
Compléments sur les morphismes, lemme \ref{more-morphisms-lemma-regular-immersion-perfect}.
\end{enumerate}
Vérifions l'associativité de la composition
des morphismes ; nous omettons de vérifier que les morphismes identiques
sont bien des identités. Pour cela, supposons donnés
$X, Y, Z, W$ et
$c \in D_{perf}(\mathcal{O}_{X \times_S Y})$,
$c' \in D_{perf}(\mathcal{O}_{Y \times_S Z})$, et
$c'' \in D_{perf}(\mathcal{O}_{Z \times_S W})$. Alors
\begin{align*}
c'' \circ (c' \circ c)
& \cong
\text{pr}^{134}_{14, *}(
\text{pr}^{134, *}_{13}
\text{pr}^{123}_{13, *}(\text{pr}^{123, *}_{12}c \otimes
\text{pr}^{123, *}_{23}c')
\otimes \text{pr}^{134, *}_{34}c'') \\
& \cong
\text{pr}^{134}_{14, *}(
\text{pr}^{1234}_{134, *}
\text{pr}^{1234, *}_{123}(\text{pr}^{123, *}_{12}c \otimes
\text{pr}^{123, *}_{23}c')
\otimes \text{pr}^{134, *}_{34}c'') \\
& \cong
\text{pr}^{134}_{14, *}(
\text{pr}^{1234}_{134, *}
(\text{pr}^{1234, *}_{12}c \otimes
\text{pr}^{1234, *}_{23}c')
\otimes \text{pr}^{134, *}_{34}c'') \\
& \cong
\text{pr}^{134}_{14, *}
\text{pr}^{1234}_{134, *}
((\text{pr}^{1234, *}_{12}c \otimes
\text{pr}^{1234, *}_{23}c')
\otimes \text{pr}^{1234, *}_{34}c'') \\
& \cong
\text{pr}^{1234}_{14, *}(
(\text{pr}^{1234, *}_{12}c \otimes
\text{pr}^{1234, *}_{23}c') \otimes
\text{pr}^{1234, *}_{34}c'')
\end{align*}
Nous employons ici la notation
$$
p^{1234}_{134} : X \times_S Y \times_S Z \times_S W
\to X \times_S Z \times_S W
\quad\text{et}\quad
p^{134}_{14} : X \times_S Z \times_S W \to X \times_S W
$$
pour les projections, et de même pour les autres indices.
Nous écrivons aussi $\text{pr}_*$ au lieu de $R\text{pr}_*$ et
$\text{pr}^*$ au lieu de $L\text{pr}^*$, et nous omettons
tous les indices supérieurs et inférieurs sur $\otimes$.
La première égalité est la définition de la composition.
La deuxième égalité vaut parce que
$\text{pr}^{134, *}_{13} \text{pr}^{123}_{13, *} =
\text{pr}^{1234}_{134, *} \text{pr}^{1234, *}_{123}$
par changement de base (Catégories dérivées des schémas, lemme
\ref{perfect-lemma-compare-base-change}).
La troisième égalité vaut parce que les images inverses se composent
de la manière attendue et commutent aux produits tensoriels, voir
Cohomologie, lemmes \ref{cohomology-lemma-derived-pullback-composition} et
\ref{cohomology-lemma-pullback-tensor-product}.
La quatrième égalité résulte de la « formule de projection » pour
$p^{1234}_{134}$, voir Catégories dérivées des schémas, lemme
\ref{perfect-lemma-cohomology-base-change}.
La cinquième égalité exprime la compatibilité de l'image directe propre
avec la composition, voir
Cohomologie, lemme \ref{cohomology-lemma-derived-pushforward-composition}.
Comme le produit tensoriel est associatif,
cela achève la démonstration de l'associativité de la composition.

\begin{lemma}
\label{equiv-lemma-base-change-is-functor}
Soit $S' \to S$ un morphisme de schémas.
La règle qui associe
\begin{enumerate}
\item à tout schéma propre et lisse $X$ sur $S$ le schéma $X' =  S' \times_S X$, et
\item à la classe d'isomorphisme d'un objet $K$
de $D_{perf}(\mathcal{O}_{X \times_S Y})$ la classe d'isomorphisme de
$L(X' \times_{S'} Y' \to X \times_S Y)^*K$
dans $D_{perf}(\mathcal{O}_{X' \times_{S'} Y'})$
\end{enumerate}
est un foncteur de la catégorie définie pour $S$ vers la catégorie
définie pour $S'$.
\end{lemma}

\begin{proof}
Pour le voir, supposons donnés $X, Y, Z$ et
$K \in D_{perf}(\mathcal{O}_{X \times_S Y})$ et
$M \in D_{perf}(\mathcal{O}_{Y \times_S Z})$.
Notons
$K' \in D_{perf}(\mathcal{O}_{X' \times_{S'} Y'})$ et
$M' \in D_{perf}(\mathcal{O}_{Y' \times_{S'} Z'})$
leurs images inverses comme dans l'énoncé du lemme.
Le carré
$$
\xymatrix{
X' \times_{S'} Y' \times_{S'} Z' \ar[r] \ar[d]_{\text{pr}'_{13}} &
X \times_S Y \times_S Z \ar[d]^{\text{pr}_{13}} \\
X' \times_{S'} Z' \ar[r] &
X \times_S Z
}
$$
est cartésien et $\text{pr}_{13}$ est propre et lisse.
D'après Catégories dérivées des schémas, lemme
\ref{perfect-lemma-flat-proper-perfect-direct-image-general},
l'image inverse dérivée par la flèche horizontale inférieure
du composé
$$
R\text{pr}_{13, *}(
L\text{pr}_{12}^*K
\otimes_{\mathcal{O}_{X \times_S Y \times_S Z}}^\mathbf{L}
L\text{pr}_{23}^*M)
$$
est bien (canoniquement) isomorphe à
$$
R\text{pr}'_{13, *}(
L(\text{pr}'_{12})^*K'
\otimes_{\mathcal{O}_{X' \times_{S'} Y' \times_{S'} Z'}}^\mathbf{L}
L(\text{pr}'_{23})^*M')
$$
comme voulu. Nous omettons quelques détails.
\end{proof}





\section{Équivalences relatives}
\label{equiv-section-relative-equivalences}

\noindent
Dans cette section, nous démontrons quelques lemmes sur la notion suivante.

\begin{definition}
\label{equiv-definition-relative-equivalence-kernel}
Soit $S$ un schéma. Soient $X \to S$ et $Y \to S$ des morphismes lisses et propres.
Un objet $K \in D_{perf}(\mathcal{O}_{X \times_S Y})$
est appelé {\it noyau de Fourier-Mukai d'une équivalence relative
de $X$ vers $Y$ au-dessus de $S$}
s'il existe un objet $K' \in D_{perf}(\mathcal{O}_{X \times_S Y})$
tel que
$$
\Delta_{X/S, *}\mathcal{O}_X \cong
R\text{pr}_{13, *}(L\text{pr}_{12}^*K
\otimes_{\mathcal{O}_{X \times_S Y \times_S X}}^\mathbf{L}
L\text{pr}_{23}^*K')
$$
dans $D(\mathcal{O}_{X \times_S X})$ et
$$
\Delta_{Y/S, *}\mathcal{O}_Y \cong
R\text{pr}_{13, *}(L\text{pr}_{12}^*K'
\otimes_{\mathcal{O}_{Y \times_S X \times_S Y}}^\mathbf{L}
L\text{pr}_{23}^*K)
$$
dans $D(\mathcal{O}_{Y \times_S Y})$. Autrement dit, la classe d'isomorphisme
de $K$ définit une flèche inversible dans la catégorie définie à la
section \ref{equiv-section-category-Fourier-Mukai-kernels}.
\end{definition}

\noindent
La terminologie est volontairement lourde.

\begin{lemma}
\label{equiv-lemma-equivalences-rek}
Avec les notations de la définition \ref{equiv-definition-relative-equivalence-kernel},
soit $K$ le noyau de Fourier-Mukai d'une équivalence relative de $X$
vers $Y$ au-dessus de $S$. Alors les foncteurs de Fourier-Mukai correspondants
$\Phi_K : D_\QCoh(\mathcal{O}_X) \to D_\QCoh(\mathcal{O}_Y)$
(lemme \ref{equiv-lemma-fourier-Mukai-QCoh})
et $\Phi_K : D_{perf}(\mathcal{O}_X) \to D_{perf}(\mathcal{O}_Y)$
(lemme \ref{equiv-lemma-fourier-mukai})
sont des équivalences.
\end{lemma}

\begin{proof}
Cela résulte immédiatement du lemme \ref{equiv-lemma-compose-fourier-mukai} et de
l'exemple \ref{equiv-example-diagonal-fourier-mukai}.
\end{proof}

\begin{lemma}
\label{equiv-lemma-base-change-rek}
Avec les notations de la définition \ref{equiv-definition-relative-equivalence-kernel},
soit $K$ le noyau de Fourier-Mukai d'une équivalence relative de $X$
vers $Y$ au-dessus de $S$. Soit $S_1 \to S$ un morphisme de schémas. Posons
$X_1 = S_1 \times_S X$ et $Y_1 = S_1 \times_S Y$. Alors l'image inverse
$K_1 = L(X_1 \times_{S_1} Y_1 \to X \times_S Y)^*K$ est
le noyau de Fourier-Mukai d'une équivalence relative de $X_1$
vers $Y_1$ au-dessus de $S_1$.
\end{lemma}

\begin{proof}
Soit $K' \in D_{perf}(\mathcal{O}_{Y \times_S X})$ l'objet dont l'existence est
supposée dans la définition \ref{equiv-definition-relative-equivalence-kernel}.
Notons $K'_1$ l'image inverse de $K'$ par
$Y_1 \times_{S_1} X_1 \to Y \times_S X$.
Il suffit alors de démontrer que
$$
\Delta_{X_1/S_1, *}\mathcal{O}_X \cong
R\text{pr}_{13, *}(L\text{pr}_{12}^*K_1
\otimes_{\mathcal{O}_{X_1 \times_{S_1} Y_1 \times_{S_1} X_1}}^\mathbf{L}
L\text{pr}_{23}^*K_1')
$$
dans $D(\mathcal{O}_{X_1 \times_{S_1} X_1})$, et de même pour l'autre
condition. Comme
$$
\xymatrix{
X_1 \times_{S_1} Y_1 \times_{S_1} X_1 \ar[r] \ar[d]_{\text{pr}_{13}} &
X \times_S Y \times_S X \ar[d]^{\text{pr}_{13}} \\
X_1 \times_{S_1} X_1 \ar[r] &
X \times_S X
}
$$
est cartésien, il suffit, d'après Catégories dérivées des schémas, lemme
\ref{perfect-lemma-flat-proper-perfect-direct-image-general},
de montrer que
$$
\Delta_{X_1/S_1, *}\mathcal{O}_{X_1}
\cong
L(X_1 \times_{S_1} X_1 \to X \times_S X)^*\Delta_{X/S, *}\mathcal{O}_X 
$$
Cela résultera à son tour de la tor-indépendance de $X$ et de
$X_1 \times_{S_1} X_1$ au-dessus de $X \times_S X$, voir
Catégories dérivées des schémas, lemme \ref{perfect-lemma-compare-base-change}.
Cette tor-indépendance se vérifie directement, mais résulte aussi
du résultat plus général de Compléments sur les morphismes, lemme
\ref{more-morphisms-lemma-case-of-tor-independence}, appliqué au carré
de sommets $X, X, X, S$ et à son changement de base par $S_1 \to S$.
\end{proof}

\begin{lemma}
\label{equiv-lemma-descend-rek}
Soit $S = \lim_{i \in I} S_i$ la limite d'un système inductif de schémas
dont les morphismes de transition affines sont $g_{i'i} : S_{i'} \to S_i$.
Supposons que $S_i$ soit quasi-compact et quasi-séparé pour tout $i \in I$.
Soit $0 \in I$. Soient $X_0 \to S_0$ et $Y_0 \to S_0$ des morphismes lisses et propres.
Posons $X_i = S_i \times_{S_0} X_0$ pour $i \geq 0$,
et $X = S \times_{S_0} X_0$ ; de même pour $Y_0$. Si $K$ est le
noyau de Fourier-Mukai d'une équivalence relative de $X$ vers $Y$ au-dessus de $S$,
alors, pour un certain $i \geq 0$, il existe un
noyau de Fourier-Mukai d'une équivalence relative de $X_i$ vers $Y_i$ au-dessus de $S_i$.
\end{lemma}

\begin{proof}
Soit $K' \in D_{perf}(\mathcal{O}_{Y \times_S X})$ l'objet dont l'existence est
supposée dans la définition \ref{equiv-definition-relative-equivalence-kernel}.
Puisque $X \times_S Y = \lim X_i \times_{S_i} Y_i$, il existe un
$i$ et des objets $K_i$ et $K'_i$ dans
$D_{perf}(\mathcal{O}_{Y_i \times_{S_i} X_i})$
dont les images inverses sur $Y \times_S X$ sont $K$ et $K'$.
Voir Catégories dérivées des schémas, lemme \ref{perfect-lemma-descend-perfect}.
D'après Catégories dérivées des schémas, lemme
\ref{perfect-lemma-flat-proper-perfect-direct-image-general},
l'objet
$$
R\text{pr}_{13, *}(L\text{pr}_{12}^*K_i
\otimes_{\mathcal{O}_{X_i \times_{S_i} Y_i \times_{S_i} X_i}}^\mathbf{L}
L\text{pr}_{23}^*K_i')
$$
est parfait, et son image inverse sur $X \times_S X$ est égale à
$$
R\text{pr}_{13, *}(L\text{pr}_{12}^*K
\otimes_{\mathcal{O}_{X \times_S Y \times_S X}}^\mathbf{L}
L\text{pr}_{23}^*K') \cong \Delta_{X/S, *}\mathcal{O}_X
$$
Voir la démonstration du lemme \ref{equiv-lemma-base-change-rek}.
D'autre part, puisque $X_i \to S$ est lisse et séparé,
l'objet
$$
\Delta_{i, *}\mathcal{O}_{X_i}
$$
de $D(\mathcal{O}_{X_i \times_{S_i} X_i})$ est lui aussi parfait
(d'après Compléments sur les morphismes, lemmes
\ref{more-morphisms-lemma-smooth-diagonal-perfect} et
\ref{more-morphisms-lemma-perfect-proper-perfect-direct-image}), et
son image inverse sur $X \times_S X$ est égale à
$$
\Delta_{X/S, *}\mathcal{O}_X
$$
Voir la démonstration du lemme \ref{equiv-lemma-base-change-rek}. Ainsi, d'après
Catégories dérivées des schémas, lemme \ref{perfect-lemma-descend-perfect},
quitte à augmenter $i$, on peut supposer que
$$
\Delta_{i, *}\mathcal{O}_{X_i} \cong
R\text{pr}_{13, *}(L\text{pr}_{12}^*K_i
\otimes_{\mathcal{O}_{X_i \times_{S_i} Y_i \times_{S_i} X_i}}^\mathbf{L}
L\text{pr}_{23}^*K_i')
$$
comme voulu. Le même raisonnement s'applique après interversion des rôles de $K$ et $K'$.
\end{proof}








\section{Absence de déformations}
\label{equiv-section-no-deformations}

\noindent
Le titre de cette section renvoie au lemme \ref{equiv-lemma-no-deformations}.

\begin{lemma}
\label{equiv-lemma-deform-koszul}
Soit $(R, \mathfrak m, \kappa) \to (A, \mathfrak n, \lambda)$
un homomorphisme local plat d'anneaux locaux
essentiellement de présentation finie.
Soit $\overline{f}_1, \ldots, \overline{f}_r \in \mathfrak n/\mathfrak m A
\subset A/\mathfrak m A$ une suite régulière. Soit $K \in D(A)$. Supposons que
\begin{enumerate}
\item $K$ soit parfait,
\item $K \otimes_A^\mathbf{L} A/\mathfrak m A$ soit isomorphe dans
$D(A/\mathfrak m A)$ au
complexe de Koszul associé à $\overline{f}_1, \ldots, \overline{f}_r$.
\end{enumerate}
Alors $K$ est isomorphe dans $D(A)$ à un complexe de Koszul associé à une suite
régulière $f_1, \ldots, f_r \in A$ qui relève les éléments donnés
$\overline{f}_1, \ldots, \overline{f}_r$. De plus, $A/(f_1, \ldots, f_r)$
est plat sur $R$.
\end{lemma}

\begin{proof}
Utilisons des complexes de chaînes dans la démonstration de ce lemme.
Le complexe de Koszul $K_\bullet(\overline{f}_1, \ldots, \overline{f}_r)$
est défini dans Compléments d'algèbre, définition
\ref{more-algebra-definition-koszul-complex}.
D'après Compléments d'algèbre, lemme \ref{more-algebra-lemma-lift-complex-stably-frees},
on peut représenter $K$ par un complexe
$$
K_\bullet :
A \to A^{\oplus r} \to \ldots \to A^{\oplus r} \to A
$$
dont le produit tensoriel avec $A/\mathfrak mA$ est égal (!)
à $K_\bullet(\overline{f}_1, \ldots, \overline{f}_r)$.
Notons $f_1, \ldots, f_r \in A$ les composantes de la
flèche $A^{\oplus r} \to A$. Ces $f_i$ relèvent les
$\overline{f}_i$. D'après Algèbre, lemme
\ref{algebra-lemma-grothendieck-regular-sequence-general}
$f_1, \ldots, f_r$ forment une suite régulière dans $A$ et $A/(f_1, \ldots, f_r)$
est plat sur $R$. Posons $J = (f_1, \ldots, f_r) \subset A$.
Considérons le diagramme
$$
\xymatrix{
K_\bullet \ar[rd] \ar@{..>}[rr]_{\varphi_\bullet} & &
K_\bullet(f_1, \ldots, f_r) \ar[ld] \\
& A/J
}
$$
Comme $f_1, \ldots, f_r$ forment une suite régulière, la flèche sud-ouest
est un quasi-isomorphisme (voir
Compléments d'algèbre, lemme \ref{more-algebra-lemma-regular-koszul-regular}).
On peut donc trouver la flèche pointillée qui fait commuter le
diagramme, par exemple grâce à
Algèbre, lemme \ref{algebra-lemma-compare-resolutions}.
En réduisant modulo $\mathfrak m$, on obtient un diagramme commutatif
$$
\xymatrix{
K_\bullet(\overline{f}_1, \ldots, \overline{f}_r)
\ar[rd] \ar[rr]_{\overline{\varphi}_\bullet} & &
K_\bullet(\overline{f}_1, \ldots, \overline{f}_r) \ar[ld] \\
& (A/\mathfrak m A)/(\overline{f}_1, \ldots, \overline{f}_r)
}
$$
par notre choix de $K_\bullet$. Ainsi, $\overline{\varphi}$ est un isomorphisme
dans la catégorie dérivée $D(A/\mathfrak m A)$. Il s'ensuit que
$\overline{\varphi} \otimes_{A/\mathfrak m A}^\mathbf{L} \lambda$
est un isomorphisme. Puisque $\overline{f}_i \in \mathfrak n / \mathfrak m A$,
on voit que
$$
\text{Tor}_i^{A/\mathfrak m A}(
K_\bullet(\overline{f}_1, \ldots, \overline{f}_r), \lambda)
=
K_i(\overline{f}_1, \ldots, \overline{f}_r) \otimes_{A/\mathfrak m A} \lambda
$$
Par conséquent, $\varphi_i \bmod \mathfrak n$ est inversible.
Comme $A$ est local, cela signifie que $\varphi_i$ est un
isomorphisme, ce qui achève la démonstration.
\end{proof}

\begin{lemma}
\label{equiv-lemma-limit-arguments}
Soit $R \to S$ un morphisme plat de type fini d'anneaux noethériens.
Soit $\mathfrak q \subset S$ un idéal premier au-dessus de
$\mathfrak p \subset R$. Soit $K \in D(S)$ parfait.
Soit $f_1, \ldots, f_r \in \mathfrak q S_\mathfrak q$
une suite régulière telle que $S_\mathfrak q/(f_1, \ldots, f_r)$
soit plat sur $R$ et telle que
$K \otimes_S^\mathbf{L} S_\mathfrak q$ soit isomorphe au
complexe de Koszul associé à $f_1, \ldots, f_r$. Alors il existe
$g \in S$, $g \not \in \mathfrak q$, tel que
\begin{enumerate}
\item $f_1, \ldots, f_r$ soient les images de
$f'_1, \ldots, f'_r \in S_g$,
\item $f'_1, \ldots, f'_r$ forment une suite régulière dans $S_g$,
\item $S_g/(f'_1, \ldots, f'_r)$ soit plat sur $R$,
\item $K \otimes_S^\mathbf{L} S_g$ soit isomorphe au
complexe de Koszul associé à $f_1, \ldots, f_r$.
\end{enumerate}
\end{lemma}

\begin{proof}
La définition des localisés permet de trouver $g \in S$, $g \not \in \mathfrak q$,
qui vérifie (1). Après avoir remplacé $g$ par
$gg'$ pour un certain $g' \in S$, $g' \not \in \mathfrak q$,
on peut supposer que (2) est vérifiée ; voir
Algèbre, lemme \ref{algebra-lemma-regular-sequence-in-neighbourhood}.
D'après Algèbre, théorème \ref{algebra-theorem-openness-flatness},
$S_g/(f'_1, \ldots, f'_r)$ est plat sur $R$
dans un voisinage ouvert de $\mathfrak q$.
Après avoir de nouveau remplacé $g$ par $gg'$ pour un certain
$g' \in S$, $g' \not \in \mathfrak q$, on peut donc supposer que (3) est aussi vérifiée.
Enfin, on obtient (4) après un remplacement supplémentaire grâce à
Compléments d'algèbre, lemme \ref{more-algebra-lemma-colimit-perfect-complexes}.
\end{proof}

\noindent
Pour une généralisation du lemme suivant, voir
Compléments sur les morphismes d'espaces, lemme
\ref{spaces-more-morphisms-lemma-where-isomorphism}.

\begin{lemma}
\label{equiv-lemma-isomorphism-in-neighbourhood}
Soit $S$ un schéma noethérien. Soit $s \in S$.
Soit $p : X \to Y$ un morphisme de schémas sur $S$.
Supposons que
\begin{enumerate}
\item $Y \to S$ et $X \to S$ soient propres,
\item $X$ soit plat sur $S$,
\item $X_s \to Y_s$ soit un isomorphisme.
\end{enumerate}
Alors il existe un voisinage ouvert $U \subset S$ de $s$
tel que le changement de base $X_U \to Y_U$ soit un isomorphisme.
\end{lemma}

\begin{proof}
Le morphisme $p$ est propre d'après Morphismes, lemme
\ref{morphisms-lemma-closed-immersion-proper}.
D'après Cohomologie des schémas, lemme
\ref{coherent-lemma-proper-finite-fibre-finite-in-neighbourhood}
il existe un ouvert $Y_s \subset V \subset Y$ tel que
$p|_{p^{-1}(V)} : p^{-1}(V) \to V$ soit fini.
D'après Compléments sur les morphismes, théorème
\ref{more-morphisms-theorem-criterion-flatness-fibre-Noetherian}
il existe un ouvert $X_s \subset U \subset X$ tel que
$p|_U : U \to Y$ soit plat. Après avoir retranché les images de
$X \setminus U$ et de $Y \setminus V$ (qui sont des parties fermées
ne contenant pas $s$), on peut supposer que $p$ est plat et fini.
Alors $p$ est ouvert (Morphismes, lemme \ref{morphisms-lemma-fppf-open})
et $Y_s \subset p(X) \subset Y$ ; après avoir rétréci $S$,
on peut donc supposer que $p$ est surjectif.
Comme $p_s : X_s \to Y_s$ est un isomorphisme, l'application
$$
p^\sharp : \mathcal{O}_Y \longrightarrow p_*\mathcal{O}_X
$$
de $\mathcal{O}_Y$-modules cohérents ($p$ est fini)
devient un isomorphisme après image inverse par $i : Y_s \to Y$
(par exemple d'après Cohomologie des schémas, lemme
\ref{coherent-lemma-affine-base-change}).
Le lemme de Nakayama implique alors que
$\mathcal{O}_{Y, y} \to (p_*\mathcal{O}_X)_y$ est surjective
pour tout $y \in Y_s$. Il existe donc un ouvert $Y_s \subset V \subset Y$
tel que $p^\sharp|_V$ soit surjective
(Modules, lemme \ref{modules-lemma-finite-type-surjective-on-stalk}).
Après avoir de nouveau rétréci $S$, on peut donc supposer que
$p^\sharp$ est surjective, ce qui signifie que $p$ est une immersion
fermée (puisque $p$ est déjà fini).
Ainsi, $p$ est une immersion fermée plate et surjective
de schémas noethériens, donc un isomorphisme ; voir
Morphismes, section \ref{morphisms-section-flat-closed-immersions}.
\end{proof}

\begin{lemma}
\label{equiv-lemma-no-deformations}
Soit $k$ un corps. Soit $S$ un schéma de type fini sur $k$
muni d'un point $k$-rationnel $s$. Soit $Y \to S$ un morphisme lisse et propre.
Soit $X = Y_s \times S \to S$ la famille constante de fibre
$Y_s$. Soit $K$ le noyau de Fourier-Mukai d'une équivalence relative
de $X$ vers $Y$ au-dessus de $S$. Supposons que l'on ait
$$
L(Y_s \times_S Y_s \to X \times_S Y)^*K \cong 
\Delta_{Y_s/k, *} \mathcal{O}_{Y_s}
$$
dans $D(\mathcal{O}_{Y_s \times Y_s})$. Alors il existe un voisinage ouvert
$s \in U \subset S$ tel que $Y|_U$ soit isomorphe à $Y_s \times U$ sur $U$.
\end{lemma}

\begin{proof}
Notons $i : Y_s \times Y_s = X_s \times Y_s \to X \times_S Y$
l'immersion fermée naturelle. (Désormais, nous noterons $Y_s$, et non $X_s$,
la fibre de $X$ au-dessus de $s$.) Soit
$z \in Y_s \times Y_s = (X \times_S Y)_s \subset X \times_S Y$
un point fermé. Comme indiqué, nous considérons $z$ à la fois comme un point fermé
de $Y_s \times Y_s$ et comme un point fermé de $X \times_S Y$.

\medskip\noindent
Cas I : $z \not \in \Delta_{Y_s/k}(Y_s)$. Notons $\mathcal{O}_z$
le $\mathcal{O}_{Y_s \times Y_s}$-module cohérent supporté en $z$
dont la fibre est $\kappa(z)$. Alors $i_*\mathcal{O}_z$ est le
$\mathcal{O}_{X \times_S Y}$-module cohérent supporté en $z$
dont la fibre est $\kappa(z)$. Notre hypothèse signifie que
$$
K \otimes_{\mathcal{O}_{X \times_S Y}}^\mathbf{L} i_*\mathcal{O}_z =
Li^*K \otimes_{\mathcal{O}_{Y_s \times Y_s}}^\mathbf{L} \mathcal{O}_z = 0
$$
D'après le lemme \ref{equiv-lemma-orthogonal-point-sheaf}, il existe donc
un voisinage ouvert $U(z) \subset X \times_S Y$ de $z$
tel que $K|_{U(z)} = 0$. Dans ce cas, on pose $Z(z) = \emptyset$
comme sous-schéma fermé de $U(z)$.

\medskip\noindent
Cas II : $z \in \Delta_{Y_s/k}(Y_s)$. Comme $Y_s$ est lisse sur $k$,
on sait que $\Delta_{Y_s/k} : Y_s \to Y_s \times Y_s$ est une
immersion régulière ; voir Compléments sur les morphismes, lemme
\ref{more-morphisms-lemma-smooth-diagonal-perfect}.
Choisissons une suite régulière $\overline{f}_1, \ldots, \overline{f}_r \in
\mathcal{O}_{Y_s \times Y_s, z}$ engendrant l'idéal de
$\Delta_{Y_s/k}(Y_s)$. Comme une suite régulière est régulière au sens de Koszul
(Compléments d'algèbre, lemme \ref{more-algebra-lemma-regular-koszul-regular}),
notre hypothèse signifie que
$$
K_z \otimes_{\mathcal{O}_{X \times_S Y, z}}^\mathbf{L}
\mathcal{O}_{Y_s \times Y_s, z}
\in D(\mathcal{O}_{Y_s \times Y_s, z})
$$
est représenté par le complexe de Koszul associé à
$\overline{f}_1, \ldots, \overline{f}_r$ sur
$\mathcal{O}_{Y_s \times Y_s, z}$.
En appliquant le lemme \ref{equiv-lemma-deform-koszul} à
$\mathcal{O}_{S, s} \to \mathcal{O}_{X \times_S Y, z}$,
on conclut que $K_z \in D(\mathcal{O}_{X \times_S Y, z})$ est
représenté par le complexe de Koszul associé à une suite régulière
$f_1, \ldots, f_r \in \mathcal{O}_{X \times_S Y, z}$
qui relève la suite régulière
$\overline{f}_1, \ldots, \overline{f}_r$
et que, de plus, $\mathcal{O}_{X \times_S Y}/(f_1, \ldots, f_r)$
est plat sur $\mathcal{O}_{S, s}$.
Par des arguments de passage à la limite (lemme \ref{equiv-lemma-limit-arguments}),
on conclut qu'il existe un voisinage ouvert affine
$U(z) \subset X \times_S Y$ de $z$ et un sous-schéma fermé
$Z(z) \subset U(z)$ tels que
\begin{enumerate}
\item $Z(z) \to U(z)$ soit une immersion fermée régulière,
\item $K|_{U(z)}$ soit quasi-isomorphe à $\mathcal{O}_{Z(z)}$,
\item $Z(z) \to S$ soit plat,
\item $Z(z)_s = \Delta_{Y_s/k}(Y_s) \cap U(z)_s$
comme sous-schémas fermés de $U(z)_s$.
\end{enumerate}

\noindent
D'après la propriété (2), pour $z, z' \in Y_s \times Y_s$, on
a $Z(z) \cap U(z') = Z(z') \cap U(z)$ en tant que sous-schémas fermés.
On obtient donc un voisinage ouvert
$$
U = \bigcup\nolimits_{z \in Y_s \times Y_s\text{ fermé}} U(z)
$$
de $Y_s \times Y_s$ dans $X \times_S Y$ et un sous-schéma fermé $Z \subset U$
tels que (1) $Z \to U$ soit une immersion fermée régulière,
(2) $Z \to S$ soit plat et (3) $Z_s = \Delta_{Y_s/k}(Y_s)$.
Comme $X \times_S Y \to S$ est propre, après avoir remplacé $S$
par un voisinage ouvert de $s$, on peut supposer que $U = X \times_S Y$.
Comme les projections $Z_s \to Y_s$ et $Z_s \to X_s$
sont des isomorphismes, on conclut qu'après avoir rétréci $S$,
on peut supposer que $Z \to Y$ et $Z \to X$ sont des isomorphismes ; voir
le lemme \ref{equiv-lemma-isomorphism-in-neighbourhood}.
Ceci achève la démonstration.
\end{proof}

\begin{lemma}
\label{equiv-lemma-no-deformations-better}
Soit $k$ un corps algébriquement clos. Soit $X$
un schéma lisse et propre sur $k$.
Soit $f : Y \to S$ un morphisme lisse et propre, où $S$ est de type fini sur $k$.
Soit $K$ le noyau de Fourier-Mukai d'une équivalence relative
de $X \times S$ vers $Y$ au-dessus de $S$. Alors $S$ peut être recouvert par des
sous-schémas ouverts $U$ tels qu'il existe un $U$-isomorphisme
$f^{-1}(U) \cong Y_0 \times U$ pour un certain $Y_0$ propre et lisse sur $k$.
\end{lemma}

\begin{proof}
Choisissons un point fermé $s \in S$. Comme $k$ est algébriquement clos,
c'est un point $k$-rationnel. Posons $Y_0 = Y_s$. La restriction
$K_0$ de $K$ à $X \times Y_0$ est le noyau de Fourier-Mukai d'une
équivalence relative de $X$ vers $Y_0$ au-dessus de $\Spec(k)$ d'après le
lemme \ref{equiv-lemma-base-change-rek}. Soit $K'_0$,
dans $D_{perf}(\mathcal{O}_{Y_0 \times X})$,
l'objet dont l'existence est postulée dans la
définition \ref{equiv-definition-relative-equivalence-kernel}.
Alors $K'_0$ est le noyau de Fourier-Mukai d'une
équivalence relative de $Y_0$ vers $X$ au-dessus de $\Spec(k)$,
par la symétrie inhérente à la
définition \ref{equiv-definition-relative-equivalence-kernel}.
Ainsi, d'après le
lemme \ref{equiv-lemma-base-change-rek},
on voit que l'image inverse
$$
M = (Y_0 \times X \times S \to Y_0 \times X)^*K'_0
$$
sur $(Y_0 \times S) \times_S (X \times S) = Y_0 \times X \times S$
est le noyau de Fourier-Mukai d'une
équivalence relative de $Y_0 \times S$ vers $X \times S$ au-dessus de $S$.
Considérons maintenant le noyau
$$
K_{new} =
R\text{pr}_{13, *}(L\text{pr}_{12}^*M
\otimes_{\mathcal{O}_{(Y_0 \times S) \times_S (X \times S)
\times_S Y}}^\mathbf{L}
L\text{pr}_{23}^*K)
$$
sur $(Y_0 \times S) \times_S Y$. C'est le noyau de Fourier-Mukai d'une
équivalence relative de $Y_0 \times S$ vers $Y$ au-dessus de $S$, car il est
la composée de deux flèches inversibles dans
la catégorie construite à la
section \ref{equiv-section-category-Fourier-Mukai-kernels}.
De plus, cette composition commute au changement de base
(lemme \ref{equiv-lemma-base-change-is-functor}).
On voit donc que l'image inverse de $K_{new}$ sur
$((Y_0 \times S) \times_S Y)_s = Y_0 \times Y_0$
est égale à la composée de $K_0$ et de $K'_0$,
et donc à l'identité dans cette catégorie.
Autrement dit, on a
$$
L(Y_0 \times Y_0 \to (Y_0 \times S) \times_S Y)^*K_{new}
\cong
\Delta_{Y_0/k, *}\mathcal{O}_{Y_0}
$$
D'après le lemme \ref{equiv-lemma-no-deformations}, on conclut donc que $Y \to S$
est isomorphe à $Y_0 \times S$ dans un voisinage ouvert de $s$.
Ceci achève la démonstration.
\end{proof}






\section{Dénombrabilité}
\label{equiv-section-countability}

\noindent
Dans cette section, nous démontrons quelques lemmes élémentaires sur la dénombrabilité
de certains ensembles. Soit $\mathcal{C}$ une catégorie. Dans cette section,
nous dirons que $\mathcal{C}$ est {\it dénombrable} si
\begin{enumerate}
\item pour tous $X, Y \in \Ob(\mathcal{C})$, l'ensemble
$\Mor_\mathcal{C}(X, Y)$ est dénombrable, et
\item l'ensemble des classes d'isomorphisme d'objets de $\mathcal{C}$
est dénombrable.
\end{enumerate}

\begin{lemma}
\label{equiv-lemma-countable-finite-type}
Soit $R$ un anneau noethérien dénombrable. Alors la catégorie des schémas de type
fini sur $R$ est dénombrable.
\end{lemma}

\begin{proof}
Omis.
\end{proof}

\begin{lemma}
\label{equiv-lemma-countable-abelian}
Soit $\mathcal{A}$ une catégorie abélienne dénombrable.
Alors $D^b(\mathcal{A})$ est dénombrable.
\end{lemma}

\begin{proof}
Il suffit de démontrer l'assertion pour $D(\mathcal{A})$, puisque les autres
en sont des sous-catégories pleines. Comme tout objet de $D(\mathcal{A})$
est un complexe d'objets de $\mathcal{A}$, il est immédiat que l'ensemble des
classes d'isomorphisme d'objets de $D^b(\mathcal{A})$ est dénombrable.
De plus, pour des complexes bornés $A^\bullet$ et $B^\bullet$ de $\mathcal{A}$,
il est clair que $\Hom_{K^b(\mathcal{A})}(A^\bullet, B^\bullet)$ est dénombrable.
On a
$$
\Hom_{D^b(\mathcal{A})}(A^\bullet, B^\bullet) =
\colim_{s : (A')^\bullet \to A^\bullet
\text{ qis et }(A')^\bullet\text{ borné}}
\Hom_{K^b(\mathcal{A})}((A')^\bullet, B^\bullet)
$$
d'après Catégories dérivées, lemme \ref{derived-lemma-bounded-derived}.
Il s'agit donc d'un ensemble dénombrable, en tant que colimite dénombrable d'ensembles dénombrables.
\end{proof}

\begin{lemma}
\label{equiv-lemma-countable-perfect}
Soit $X$ un schéma de type fini sur un anneau noethérien dénombrable.
Alors les catégories $D_{perf}(\mathcal{O}_X)$ et
$D^b_{\textit{Coh}}(\mathcal{O}_X)$ sont dénombrables.
\end{lemma}

\begin{proof}
Observons que $X$ est noethérien d'après
Morphismes, lemme \ref{morphisms-lemma-finite-type-noetherian}.
Ainsi, $D_{perf}(\mathcal{O}_X)$ est une sous-catégorie pleine de
$D^b_{\textit{Coh}}(\mathcal{O}_X)$ d'après
Catégories dérivées des schémas, lemme \ref{perfect-lemma-perfect-on-noetherian}.
Il suffit donc de démontrer
le résultat pour $D^b_{\textit{Coh}}(\mathcal{O}_X)$.
Rappelons que
$D^b_{\textit{Coh}}(\mathcal{O}_X) = D^b(\textit{Coh}(\mathcal{O}_X))$
d'après
Catégories dérivées des schémas, proposition \ref{perfect-proposition-DCoh}.
Par le lemme \ref{equiv-lemma-countable-abelian},
il suffit de démontrer que $\textit{Coh}(\mathcal{O}_X)$ est
dénombrable. Nous omettons cette vérification.
\end{proof}

\begin{lemma}
\label{equiv-lemma-countable-isos}
Soit $K$ un corps algébriquement clos.
Soit $S$ un schéma de type fini sur $K$.
Soient $X \to S$ et $Y \to S$ des morphismes de type fini.
Il existe un ensemble dénombrable $I$ et, pour tout $i \in I$, un couple
$(S_i \to S, h_i)$ possédant les propriétés suivantes :
\begin{enumerate}
\item $S_i \to S$ est un morphisme de type fini, et l'on pose
$X_i = X \times_S S_i$ et $Y_i = Y \times_S S_i$ ;
\item $h_i : X_i \to Y_i$ est un isomorphisme sur $S_i$ ;
\item pour tout point fermé $s \in S(K)$, si $X_s \cong Y_s$
sur $K = \kappa(s)$, alors $s$ appartient à l'image de $S_i \to S$
pour un certain $i$.
\end{enumerate}
\end{lemma}

\begin{proof}
Le corps $K$ est la réunion filtrante de ses sous-corps dénombrables.
Dualement, $\Spec(K)$ est la limite cofiltrante des spectres
des sous-corps dénombrables de $K$.
Le lemme \ref{limits-lemma-descend-finite-presentation} de Limites
assure donc que l'on peut trouver un sous-corps dénombrable
$k$ et des morphismes $X_0 \to S_0$ et $Y_0 \to S_0$
de schémas de type fini sur $k$ tels que
$X \to S$ et $Y \to S$ soient leurs changements de base.

\medskip\noindent
Par le lemme \ref{equiv-lemma-countable-finite-type}, il existe un ensemble dénombrable $I$ et
des couples $(S_{0, i} \to S_0, h_{0, i})$ tels que
\begin{enumerate}
\item $S_{0, i} \to S_0$ est un morphisme de type fini, et l'on pose
$X_{0, i} = X_0 \times_{S_0} S_{0, i}$ et
$Y_{0, i} = Y_0 \times_{S_0} S_{0, i}$ ;
\item $h_{0, i} : X_{0, i} \to Y_{0, i}$ est un isomorphisme sur $S_{0, i}$.
\end{enumerate}
de façon que tout couple $(T \to S_0, h_T)$, avec $T \to S_0$ de type fini
et $h_T : X_0 \times_{S_0} T \to Y_0 \times_{S_0} T$ un isomorphisme,
soit isomorphe à l'un d'eux.
Notons $(S_i \to S, h_i)$ le changement de base de $(S_{0, i} \to S_0, h_{0, i})$
par $\Spec(K) \to \Spec(k)$.
Nous affirmons que cela convient.

\medskip\noindent
Soit $s \in S(K)$ et soit $h_s : X_s \to Y_s$ un isomorphisme sur
$K = \kappa(s)$. On peut écrire $K$ comme réunion filtrante de ses
$k$-sous-algèbres de type fini. Par conséquent, d'après
Limites, proposition
\ref{limits-proposition-characterize-locally-finite-presentation}, et
le lemme \ref{limits-lemma-descend-finite-presentation},
on peut trouver une telle $k$-sous-algèbre de type fini
$K \supset A \supset k$ telle que
\begin{enumerate}
\item il existe un diagramme commutatif
$$
\xymatrix{
\Spec(K) \ar[d]_s \ar[r] &
\Spec(A) \ar[d]^{s'} \\
S \ar[r] &
S_0}
$$
pour un certain morphisme $s' : \Spec(A) \to S_0$ sur $k$ ;
\item $h_s$ est le changement de base d'un isomorphisme
$h_{s'} : X_0 \times_{S_0, s'} \Spec(A) \to
X_0 \times_{S_0, s'} \Spec(A)$ sur $A$.
\end{enumerate}
Alors $(s' : \Spec(A) \to S_0, h_{s'})$ est bien sûr isomorphe
au couple $(S_{0, i} \to S_0, h_{0, i})$ pour un certain $i \in I$.
Ceci achève la démonstration, car le diagramme commutatif
de (1) montre que $s$ appartient à l'image
du changement de base de $s'$ à $\Spec(K)$.
\end{proof}

\begin{lemma}
\label{equiv-lemma-countable-equivs}
Soit $K$ un corps algébriquement clos. Il existe un ensemble dénombrable $I$
et, pour tout $i \in I$, un système $(S_i/K, X_i \to S_i, Y_i \to S_i, M_i)$
possédant les propriétés suivantes :
\begin{enumerate}
\item $S_i$ est un schéma de type fini sur $K$ ;
\item $X_i \to S_i$ et $Y_i \to S_i$ sont des morphismes
propres et lisses ;
\item $M_i \in D_{perf}(\mathcal{O}_{X_i \times_{S_i} Y_i})$
est le noyau de Fourier-Mukai d'une équivalence relative de
$X_i$ vers $Y_i$ sur $S_i$, et
\item pour tous schémas lisses et propres $X$ et $Y$ sur $K$
tels qu'il existe une équivalence exacte et $K$-linéaire
$D_{perf}(\mathcal{O}_X) \to D_{perf}(\mathcal{O}_Y)$
il existe $i \in I$ et $s \in S_i(K)$
tels que $X \cong (X_i)_s$ et $Y \cong (Y_i)_s$.
\end{enumerate}
\end{lemma}

\begin{proof}
Choisissons un sous-corps dénombrable $k \subset K$, par exemple le sous-corps premier.
D'après les lemmes \ref{equiv-lemma-countable-finite-type} et \ref{equiv-lemma-countable-perfect},
il existe un ensemble dénombrable de classes d'isomorphisme de systèmes
sur $k$ satisfaisant aux assertions (1), (2) et (3) du lemme.
Nous pouvons donc choisir un ensemble dénombrable
$I$ et, pour tout $i \in I$, un tel système
$$
(S_{0, i}/k, X_{0, i} \to S_{0, i}, Y_{0, i} \to S_{0, i}, M_{0, i})
$$
sur $k$ tel que chaque classe d'isomorphisme soit représentée au moins une fois.
Notons $(S_i/K, X_i \to S_i, Y_i \to S_i, M_i)$ le changement de base
à $K$ du système affiché. Ce système possède les propriétés (1), (2) et (3),
voir le lemme \ref{equiv-lemma-base-change-rek}. Démontrons la propriété (4).

\medskip\noindent
Considérons des schémas lisses et propres $X$ et $Y$ sur $K$
tels qu'il existe une équivalence exacte et $K$-linéaire
$F : D_{perf}(\mathcal{O}_X) \to D_{perf}(\mathcal{O}_Y)$.
Par la proposition \ref{equiv-proposition-equivalence},
nous pouvons supposer qu'il existe un objet
$M \in D_{perf}(\mathcal{O}_{X \times Y})$
tel que $F = \Phi_M$ soit le foncteur de Fourier-Mukai correspondant.
Par le lemme \ref{equiv-lemma-fourier-mukai-flat-proper-over-noetherian},
il existe un objet $M'$ dans $D_{perf}(\mathcal{O}_{Y \times X})$
tel que $\Phi_{M'}$ soit l'adjoint à droite de $\Phi_M$.
Comme $\Phi_M$ est une équivalence, cela signifie que
$\Phi_{M'}$ est le quasi-inverse de $\Phi_M$.
Par le lemme \ref{equiv-lemma-fourier-mukai-flat-proper-over-noetherian},
les foncteurs de Fourier-Mukai définis par les objets
$$
A = R\text{pr}_{13, *}(
L\text{pr}_{12}^*M
\otimes_{\mathcal{O}_{X \times Y \times X}}^\mathbf{L}
L\text{pr}_{23}^*M')
$$
dans $D_{perf}(\mathcal{O}_{X \times X})$ et
$$
B = R\text{pr}_{13, *}(
L\text{pr}_{12}^*M'
\otimes_{\mathcal{O}_{Y \times X \times Y}}^\mathbf{L}
L\text{pr}_{23}^*M)
$$
dans $D_{perf}(\mathcal{O}_{Y \times Y})$
sont isomorphes à
$\text{id} : D_{perf}(\mathcal{O}_X) \to D_{perf}(\mathcal{O}_X)$
et à
$\text{id} : D_{perf}(\mathcal{O}_Y) \to D_{perf}(\mathcal{O}_Y)$
respectivement. Par conséquent,
$A \cong \Delta_{X/K, *}\mathcal{O}_X$ et
$B \cong \Delta_{Y/K, *}\mathcal{O}_Y$
d'après le lemme \ref{equiv-lemma-uniqueness}. Ainsi, $M$ est le
noyau de Fourier-Mukai d'une équivalence relative de $X$ vers $Y$
sur $K$, par définition.

\medskip\noindent
On peut écrire $K$ comme colimite filtrante de ses
$k$-sous-algèbres de type fini $A \subset K$. D'après
Limites, lemme \ref{limits-lemma-descend-finite-presentation},
on peut trouver $X_0, Y_0$ de type fini sur $A$ dont les
changements de base à $K$ donnent $X$ et $Y$.
D'après Limites, lemmes
\ref{limits-lemma-eventually-proper} et \ref{limits-lemma-descend-smooth},
quitte à agrandir $A$, on peut supposer $X_0$ et $Y_0$
lisses et propres sur $A$.
Par le lemme \ref{equiv-lemma-descend-rek},
quitte à agrandir $A$, on peut supposer que $M$ est l'image inverse d'un
objet $M_0  \in D_{perf}(\mathcal{O}_{X_0 \times_{\Spec(A)} Y_0})$
qui est le noyau de Fourier-Mukai d'une équivalence relative
de $X_0$ vers $Y_0$ sur $\Spec(A)$.
Ainsi, le système $(S_0/k, X_0 \to S_0, Y_0 \to S_0, M_0)$
est isomorphe à
$(S_{0, i}/k, X_{0, i} \to S_{0, i}, Y_{0, i} \to S_{0, i}, M_{0, i})$
pour un certain $i \in I$.
Comme $S_i = S_{0, i} \times_{\Spec(k)} \Spec(K)$,
on en déduit que (4) est vraie pour $s : \Spec(K) \to S_i$
induit par le morphisme $\Spec(K) \to \Spec(A) \cong S_{0, i}$
provenant de $A \subset K$.
\end{proof}








\section{Dénombrabilité des variétés équivalentes au sens dérivé}
\label{equiv-section-countable-derived-equivalent}

\noindent
Dans cette section, nous démontrons un résultat d'Anel et To\"en, voir \cite{AT}.

\begin{definition}
\label{equiv-definition-derived-equivalent}
Soit $k$ un corps. Soient $X$ et $Y$ des schémas projectifs lisses sur $k$.
On dit que $X$ et $Y$ sont {\it équivalents au sens dérivé} s'il existe une équivalence $k$-linéaire
exacte
$D_{perf}(\mathcal{O}_X) \to D_{perf}(\mathcal{O}_Y)$.
\end{definition}

\noindent
Voici le résultat.

\begin{theorem}
\label{equiv-theorem-countable}
\begin{reference}
Légère amélioration de \cite{AT}
\end{reference}
Soit $K$ un corps algébriquement clos. Soit $\mathbf{X}$ un schéma lisse et propre
sur $K$. Il existe au plus un nombre dénombrable de classes d'isomorphisme
de schémas lisses et propres $\mathbf{Y}$ sur $K$ qui sont
équivalents à $\mathbf{X}$ au sens dérivé.
\end{theorem}

\begin{proof}
Choisissons un ensemble dénombrable $I$ et, pour $i \in I$, des systèmes
$(S_i/K, X_i \to S_i, Y_i \to S_i, M_i)$ satisfaisant aux propriétés
(1), (2), (3) et (4) du lemme \ref{equiv-lemma-countable-equivs}.
Fixons $i \in I$ et posons $S = S_i$, $X = X_i$, $Y = Y_i$ et
$M = M_i$. Il suffit manifestement de montrer que
l'ensemble des classes d'isomorphisme des fibres $Y_s$ pour $s \in S(K)$
telles que $X_s \cong \mathbf{X}$ est dénombrable.
C'est ce que nous démontrons dans le paragraphe suivant.

\medskip\noindent
Soit $S$ un schéma de type fini sur $K$, et soient $X \to S$ et $Y \to S$
des morphismes propres et lisses ; soit $M \in D_{perf}(\mathcal{O}_{X \times_S Y})$
le noyau de Fourier-Mukai d'une équivalence relative de $X$
vers $Y$ sur $S$. Nous allons montrer que
l'ensemble des classes d'isomorphisme des fibres $Y_s$ pour $s \in S(K)$
telles que $X_s \cong \mathbf{X}$ est dénombrable.
Par le lemme \ref{equiv-lemma-countable-isos}, appliqué
aux familles $\mathbf{X} \times S \to S$ et $X \to S$,
il existe un ensemble dénombrable $I$ et, pour $i \in I$, un couple
$(S_i \to S, h_i)$ possédant les propriétés suivantes :
\begin{enumerate}
\item $S_i \to S$ est un morphisme de type fini, et l'on pose
$X_i = X \times_S S_i$ ;
\item $h_i : \mathbf{X} \times S_i \to X_i$
est un isomorphisme sur $S_i$ ;
\item pour tout point fermé $s \in S(K)$, si $\mathbf{X} \cong X_s$
sur $K = \kappa(s)$, alors $s$ appartient à l'image de $S_i \to S$
pour un certain $i$.
\end{enumerate}
Posons $Y_i = Y \times_S S_i$. Notons
$M_i \in D_{perf}(\mathcal{O}_{X_i \times_{S_i} Y_i})$
l'image inverse de $M$. Par le lemme \ref{equiv-lemma-base-change-rek},
$M_i$ est le noyau de Fourier-Mukai d'une équivalence relative de
$X_i$ vers $Y_i$ sur $S_i$. Comme $I$ est dénombrable, la
propriété (3) ramène la démonstration à prouver que
l'ensemble des classes d'isomorphisme des fibres $Y_{i, s}$ pour $s \in S_i(K)$
est dénombrable.
En fait, cet ensemble est fini d'après
le lemme \ref{equiv-lemma-no-deformations-better},
ce qui achève la démonstration.
\end{proof}












% OK, start here.
%

\chapter{Groupes fondamentaux des schémas}



\phantomsection
\label{pione-section-phantom}


\section{Introduction}
\label{pione-section-introduction}

\noindent
Dans ce chapitre, nous étudions le groupe fondamental d'un schéma au sens de
Grothendieck ainsi que ses applications. Une référence fondamentale est
\cite{SGA1}. On trouvera une bonne introduction dans \cite{Lenstra}.
Voir aussi \cite{Murre-lectures} et \cite{Grothendieck-Murre}.










\section{Schémas étales sur un point}
\label{pione-section-schemes-etale-point}

\noindent
Dans cette section, nous décrivons les schémas étales sur le spectre d'un
corps. Avant d'énoncer le résultat, introduisons la catégorie des
$G$-ensembles pour un groupe topologique $G$.

\begin{definition}
\label{pione-definition-G-set-continuous}
Soit $G$ un groupe topologique.
Un {\it $G$-ensemble}, parfois appelé {\it $G$-ensemble discret},
est un ensemble $X$ muni d'une action à gauche $a : G \times X \to X$
telle que $a$ soit continue lorsque $X$ est muni de la topologie discrète et
$G \times X$ de la topologie produit.
Un {\it morphisme de $G$-ensembles} $f : X \to Y$ est simplement une
application $G$-équivariante de $X$ dans $Y$.
La catégorie des $G$-ensembles est notée {\it $G\textit{-Ensembles}$}.
\end{definition}

\noindent
La continuité de $a : G \times X \to X$ signifie simplement que le
stabilisateur de tout $x \in X$ est ouvert dans $G$.
Si $G$ est un groupe abstrait (c'est-à-dire un groupe non muni d'une
topologie), ceci concorde avec notre définition précédente (voir par exemple
Sites, exemple \ref{sites-example-site-on-group}), pourvu que l'on munisse
$G$ de la topologie discrète.

\medskip\noindent
Rappelons que, si $L/K$ est une extension galoisienne infinie, le groupe de
Galois $G = \text{Gal}(L/K)$ est muni d'une topologie canonique ; voir
Corps, section \ref{fields-section-infinite-galois}.

\begin{lemma}
\label{pione-lemma-sheaves-point}
Soit $K$ un corps. Soit $K^{sep}$ une clôture séparable de $K$.
Considérons le groupe profini $G = \text{Gal}(K^{sep}/K)$.
Le foncteur
$$
\begin{matrix}
\text{schémas étales sur }K &
\longrightarrow &
G\textit{-Ensembles} \\
X/K & \longmapsto &
\Mor_{\Spec(K)}(\Spec(K^{sep}), X)
\end{matrix}
$$
est une équivalence de catégories.
\end{lemma}

\begin{proof}
Un schéma $X$ sur $K$ est étale sur $K$ si et seulement si
$X \cong \coprod_{i\in I} \Spec(K_i)$, où chaque $K_i$ est une extension
séparable finie de $K$
(Morphismes, lemme \ref{morphisms-lemma-etale-over-field}).
Le foncteur du lemme associe à $X$ le $G$-ensemble
$$
\coprod\nolimits_i \Hom_K(K_i, K^{sep})
$$
muni de son action à gauche naturelle de $G$. Par définition de la topologie
sur $G$, chaque élément a un stabilisateur ouvert. Réciproquement, tout
$G$-ensemble $S$ est réunion disjointe de ses orbites. Écrivons
$S = \coprod S_i$. Choisissons $s_i \in S_i$ et notons $G_i \subset G$ son
stabilisateur ouvert. D'après la théorie de Galois
(Corps, théorème \ref{fields-theorem-inifinite-galois-theory})
les corps $(K^{sep})^{G_i}$ sont des extensions séparables finies de $K$ ;
par conséquent, le schéma
$$
\coprod\nolimits_i \Spec((K^{sep})^{G_i})
$$
est étale sur $K$. On obtient ainsi un quasi-inverse du foncteur du lemme.
Nous omettons certains détails.
\end{proof}

\begin{remark}
\label{pione-remark-covering-surjective}
Sous la correspondance du lemme \ref{pione-lemma-sheaves-point}, les recouvrements
du petit site étale $\Spec(K)_\etale$ de $K$ correspondent aux familles
surjectives de morphismes dans $G\textit{-Ensembles}$.
\end{remark}









\section{Catégories galoisiennes}
\label{pione-section-galois}

\noindent
Dans cette section, nous présentons une partie du contenu que le lecteur
trouvera dans \cite[Exposé V, sections 4, 5 et 6]{SGA1}.

\medskip\noindent
Soit $F : \mathcal{C} \to \textit{Ensembles}$ un foncteur.
Rappelons que, suivant nos conventions, les catégories ont un ensemble
d'objets et, pour toute paire d'objets, un ensemble de morphismes. Il existe
une application injective canonique
\begin{equation}
\label{pione-equation-embedding-product}
\text{Aut}(F)
\longrightarrow
\prod\nolimits_{X \in \Ob(\mathcal{C})} \text{Aut}(F(X))
\end{equation}
Pour un ensemble $E$, nous munissons $\text{Aut}(E)$ de la topologie
compacte-ouverte ; voir Topologie, exemple
\ref{topology-example-automorphisms-of-a-set}. Il s'agit bien entendu de la
topologie discrète lorsque $E$ est fini, ce qui est le cas qui nous intéresse
dans cette section\footnote{Lorsque nous étudierons le groupe fondamental
pro-étale, le cas général nous intéressera.}. Nous munissons
$\text{Aut}(F)$ de la topologie induite par la topologie produit du membre de
droite de (\ref{pione-equation-embedding-product}). En particulier, les applications
d'action
$$
\text{Aut}(F) \times F(X) \longrightarrow F(X)
$$
sont continues lorsque $F(X)$ est muni de la topologie discrète, puisque
c'est le cas des applications d'action $\text{Aut}(E) \times E \to E$ pour
tout ensemble $E$. La propriété universelle de notre topologie sur
$\text{Aut}(F)$ est la suivante : supposons que $G$ soit un groupe topologique
et que $G \to \text{Aut}(F)$ soit un homomorphisme de groupes tel que les
actions induites $G \times F(X) \to F(X)$ soient continues pour tout
$X \in \Ob(\mathcal{C})$, où $F(X)$ est muni de la topologie discrète.
Alors $G \to \text{Aut}(F)$ est continue.

\medskip\noindent
Le lemme suivant affirme que le groupe des automorphismes d'un foncteur à
valeurs dans la catégorie des ensembles finis est automatiquement profini.

\begin{lemma}
\label{pione-lemma-aut-inverse-limit}
Soit $\mathcal{C}$ une catégorie et soit
$F : \mathcal{C} \to \textit{Ensembles}$ un foncteur. L'application
(\ref{pione-equation-embedding-product}) identifie $\text{Aut}(F)$ à un sous-groupe
fermé de
$\prod_{X \in \Ob(\mathcal{C})} \text{Aut}(F(X))$.
En particulier, si $F(X)$ est fini pour tout $X$, alors
$\text{Aut}(F)$ est un groupe profini.
\end{lemma}

\begin{proof}
Soit $\xi = (\gamma_X) \in \prod \text{Aut}(F(X))$ un élément qui
n'appartient pas à $\text{Aut}(F)$. Il existe alors un morphisme
$f : X \to X'$ de $\mathcal{C}$ et un élément $x \in F(X)$ tels que
$F(f)(\gamma_X(x)) \not = \gamma_{X'}(F(f)(x))$.
Considérons le voisinage ouvert
$U = \{\gamma \in \text{Aut}(F(X)) \mid \gamma(x) = \gamma_X(x)\}$
de $\gamma_X$ et le voisinage ouvert
$U' = \{\gamma' \in \text{Aut}(F(X')) \mid \gamma'(F(f)(x)) =
\gamma_{X'}(F(f)(x))\}$.
Alors
$U \times U' \times \prod_{X'' \not = X, X'} \text{Aut}(F(X''))$
est un voisinage ouvert de $\xi$ qui ne rencontre pas $\text{Aut}(F)$.
La dernière assertion résulte du fait que
$\prod \text{Aut}(F(X))$ est un espace profini si chaque $F(X)$ est fini.
\end{proof}

\begin{example}
\label{pione-example-galois-category-G-sets}
Soit $G$ un groupe topologique. Un exemple important sera le foncteur d'oubli
\begin{equation}
\label{pione-equation-forgetful}
\textit{Finite-}G\textit{-Ensembles} \longrightarrow \textit{Ensembles}
\end{equation}
où $\textit{Finite-}G\textit{-Ensembles}$ est la sous-catégorie pleine de
$G\textit{-Ensembles}$ dont les objets sont les $G$-ensembles finis.
La catégorie $G\textit{-Ensembles}$ des $G$-ensembles est définie dans la
définition \ref{pione-definition-G-set-continuous}.
\end{example}

\noindent
Soit $G$ un groupe topologique. La {\it complétion profinie} de $G$ est le
groupe profini
$$
G^\wedge =
\lim_{U \subset G\text{ ouvert, distingué, d'indice fini}} G/U
$$
muni de sa topologie profinie. Remarquons que la limite est cofiltrante,
car toute intersection finie de sous-groupes ouverts distingués d'indice fini
est encore de cette forme. La propriété universelle de la complétion profinie
affirme que toute application continue $G \to H$ vers un groupe profini $H$
se factorise canoniquement sous la forme $G \to G^\wedge \to H$.

\begin{lemma}
\label{pione-lemma-single-out-profinite}
Soit $G$ un groupe topologique. Le groupe des automorphismes du foncteur
(\ref{pione-equation-forgetful}), muni de la topologie profinie du lemme
\ref{pione-lemma-aut-inverse-limit}, est la complétion profinie de $G$.
\end{lemma}

\begin{proof}
Notons $F_G$ le foncteur (\ref{pione-equation-forgetful}). Tout morphisme
$X \to Y$ dans $\textit{Finite-}G\textit{-Ensembles}$ commute à l'action de $G$.
Ainsi, tout $g \in G$ définit un automorphisme de $F_G$, et l'on obtient un
homomorphisme canonique de groupes $G \to \text{Aut}(F_G)$.
Remarquons que tout $G$-ensemble fini $X$ est réunion disjointe finie de
$G$-ensembles de la forme $G/H_i$, munis de l'action canonique de $G$, où
$H_i \subset G$ est un sous-groupe ouvert d'indice fini. Alors
$U_i = \bigcap gH_ig^{-1}$ est ouvert, distingué et d'indice fini.
De plus, $U_i$ agit trivialement sur $G/H_i$ ; ainsi,
$U = \bigcap U_i$ agit trivialement sur $F_G(X)$. L'action
$G \times F_G(X) \to F_G(X)$ est donc continue. Par la propriété universelle
de la topologie sur $\text{Aut}(F_G)$, l'application
$G \to \text{Aut}(F_G)$ est continue.
D'après le lemme \ref{pione-lemma-aut-inverse-limit} et la propriété universelle
de la complétion profinie, on en déduit un homomorphisme continu de groupes
$$
G^\wedge \longrightarrow \text{Aut}(F_G)
$$
De plus, puisque $G/U$ agit fidèlement sur $G/U$, cette application est
injective. Si son image est dense, l'application est surjective, donc c'est
un homéomorphisme d'après Topologie, lemme
\ref{topology-lemma-bijective-map}.

\medskip\noindent
Soient $\gamma \in \text{Aut}(F_G)$ et $X \in \Ob(\mathcal{C})$.
Montrons qu'il existe $g \in G$ tel que $\gamma$ et $g$ induisent la même
action sur $F_G(X)$, ce qui achèvera la démonstration. Comme ci-dessus,
$X$ est une réunion disjointe finie de $G/H_i$. Avec les notations $U_i$ et
$U$ précédentes, le $G$-ensemble fini $Y = G/U$ se surjecte sur $G/H_i$ pour
tout $i$ ; il suffit donc de trouver $g \in G$ tel que $\gamma$ et $g$
induisent la même action sur $F_G(G/U) = G/U$. Soit $e \in G$ l'élément
neutre, et écrivons $\gamma(eU) = g_0U$ pour un certain $g_0 \in G$.
Pour tout $g_1 \in G$, le morphisme
$$
R_{g_1} : G/U \longrightarrow G/U,\quad gU \longmapsto gg_1U
$$
de $\textit{Finite-}G\textit{-Ensembles}$ commute à l'action de $\gamma$. Ainsi
$$
\gamma(g_1U) = \gamma(R_{g_1}(eU)) = R_{g_1}(\gamma(eU)) =
R_{g_1}(g_0U) = g_0g_1U
$$
On voit donc que $g = g_0$ convient.
\end{proof}

\noindent
Rappelons qu'un foncteur exact est un foncteur qui commute à toutes les
limites finies et à toutes les colimites finies. En particulier, un tel
foncteur commute aux égalisateurs, aux coégalisateurs, aux produits fibrés,
aux sommes amalgamées, etc.

\begin{lemma}
\label{pione-lemma-second-fundamental-functor}
Soit $G$ un groupe topologique. Soit
$F : \textit{Finite-}G\textit{-Ensembles} \to \textit{Ensembles}$
un foncteur exact tel que $F(X)$ soit fini pour tout $X$.
Alors $F$ est isomorphe au foncteur (\ref{pione-equation-forgetful}).
\end{lemma}

\begin{proof}
Soit $X$ un objet non vide de $\textit{Finite-}G\textit{-Ensembles}$. Le diagramme
$$
\xymatrix{
X \ar[r] \ar[d] & \{*\} \ar[d] \\
\{*\} \ar[r] & \{*\}
}
$$
est cocartésien. On en déduit que $F(X)$ est non vide.
Soit $U \subset G$ un sous-groupe ouvert distingué d'indice fini. Remarquons que
$$
G/U \times G/U = \coprod\nolimits_{gU \in G/U} G/U
$$
où le facteur correspondant à $gU$ correspond, dans le membre de gauche, à
l'orbite de $(eU, gU)$. On a alors
$$
F(G/U) \times F(G/U) = F(G/U \times G/U) = \coprod\nolimits_{gU \in G/U} F(G/U)
$$
Puisque $F(G/U)$ est non vide, il s'ensuit que
$|F(G/U)| = |G/U|$. Par conséquent,
$$
\lim_{U \subset G\text{ ouvert, distingué, d'indice fini}} F(G/U)
$$
est non vide (Catégories, lemme \ref{categories-lemma-nonempty-limit}).
Choisissons dans cette limite un élément $\gamma = (\gamma_U)$.
Notons $F_G$ le foncteur (\ref{pione-equation-forgetful}). On peut identifier
$F_G$ au foncteur
$$
X \longmapsto \colim_U \Mor(G/U, X)
$$
où $f : G/U \to X$ correspond à $f(eU) \in X = F_G(X)$
(nous omettons les détails). L'élément $\gamma$ détermine donc une
application bien définie
$$
t : F_G \longrightarrow F
$$
En effet, étant donné $x \in X$, choisissons $U$ et
$f : G/U \to X$ envoyant $eU$ sur $x$, puis posons
$t_X(x) = F(f)(\gamma_U)$. Montrons que $t$ induit une application bijective
$t_{G/U} : F_G(G/U) \to F(G/U)$ pour tout $U$. Il en résulte immédiatement
que $t$ est un isomorphisme (nous omettons les détails).
Comme $|F_G(G/U)| = |F(G/U)|$, il suffit de montrer que $t_{G/U}$ est
surjective. Son image contient au moins un élément, à savoir
$t_{G/U}(eU) = F(\text{id}_{G/U})(\gamma_U) = \gamma_U$.
Pour $g \in G$, notons $R_g : G/U \to G/U$ la multiplication à droite.
Si $g \not \in U$, l'ensemble des points fixes de
$F(R_g) : F(G/U) \to F(G/U)$ est égal à
$F(\emptyset) = \emptyset$, car $F$ commute aux égalisateurs. Il s'ensuit que,
si $g_1, \ldots, g_{|G/U|}$ est un système de représentants de $G/U$, les
éléments $F(R_{g_i})(\gamma_U)$ sont deux à deux distincts et constituent
donc tout $F(G/U)$. Enfin,
$$
t_{G/U}(g_iU) = F(R_{g_i})(\gamma_U)
$$
ce qui achève la démonstration.
\end{proof}

\begin{example}
\label{pione-example-from-C-F-to-G-sets}
Soit $\mathcal{C}$ une catégorie et soit
$F : \mathcal{C} \to \textit{Ensembles}$ un foncteur tel que $F(X)$ soit fini
pour tout $X \in \Ob(\mathcal{C})$. Le lemme
\ref{pione-lemma-aut-inverse-limit} montre que $G = \text{Aut}(F)$ est muni de
manière canonique d'une structure de groupe topologique profini. On obtient
un foncteur
\begin{equation}
\label{pione-equation-remember}
\mathcal{C} \longrightarrow \textit{Finite-}G\textit{-Ensembles},\quad
X \longmapsto F(X)
\end{equation}
où $F(X)$ est muni de l'action induite de $G$. Cette action est continue par
construction de la topologie sur $\text{Aut}(F)$.
\end{example}

\noindent
Le but de la définition des catégories galoisiennes est de distinguer les
couples $(\mathcal{C}, F)$ pour lesquels le foncteur (\ref{pione-equation-remember})
est une équivalence. Voici notre définition d'une catégorie galoisienne.

\begin{definition}
\label{pione-definition-galois-category}
\begin{reference}
Cette définition diffère de celle de \cite[Expos\'e V, Définition 5.1]{SGA1}.
Comparer avec \cite[Définition 7.2.1]{BS}.
\end{reference}
Soit $\mathcal{C}$ une catégorie et soit $F : \mathcal{C} \to \textit{Ensembles}$
un foncteur. Le couple $(\mathcal{C}, F)$ est une {\it catégorie galoisienne} si
\begin{enumerate}
\item $\mathcal{C}$ admet les limites finies et les colimites finies,
\item
\label{pione-item-connected-components}
tout objet de $\mathcal{C}$ est un coproduit fini (éventuellement vide)
d'objets connexes,
\item $F(X)$ est fini pour tout $X \in \Ob(\mathcal{C})$, et
\item $F$ est conservatif\footnote{Plus précisément, étant donné un morphisme
$f$ de $\mathcal{C}$, si $F(f)$ est un isomorphisme, alors
$f$ est un isomorphisme.} et est exact\footnote{Cela signifie que
$F$ commute aux limites finies et aux colimites finies ; voir
Catégories, section \ref{categories-section-exact-functor}.}.
\end{enumerate}
On dit ici que $X \in \Ob(\mathcal{C})$ est connexe s'il
n'est pas initial et si, pour tout monomorphisme $Y \to X$,
soit $Y$ est initial, soit $Y \to X$ est un isomorphisme.
\end{definition}

\noindent
{\bf Avertissement :} Cette définition n'est pas la même que celle donnée dans la plupart
des références, même si nous verrons finalement que les deux sont équivalentes.
En effet, dans \cite[Expos\'e V, Définition 5.1]{SGA1}, une catégorie galoisienne est
définie comme une catégorie équivalente à $\textit{Finite-}G\textit{-Ensembles}$
pour un certain groupe profini $G$. Grothendieck caractérise alors les
catégories galoisiennes par une liste d'axiomes (G1) -- (G6) plus faibles
que les nôtres ci-dessus. Notre choix vise à souligner
l'existence des limites finies et des colimites finies, ainsi que l'exactitude du
foncteur $F$. En contrepartie, il nous faudra plus tard travailler
un peu davantage pour appliquer les résultats de cette section.

\begin{lemma}
\label{pione-lemma-epi-mono}
Soit $(\mathcal{C}, F)$ une catégorie galoisienne. Soit
$X \to Y \in \text{Arrows}(\mathcal{C})$. Alors
\begin{enumerate}
\item $F$ est fidèle,
\item $X \to Y$ est un monomorphisme
$\Leftrightarrow F(X) \to F(Y)$ est injective,
\item $X \to Y$ est un épimorphisme
$\Leftrightarrow F(X) \to F(Y)$ est surjective,
\item un objet $A$ de $\mathcal{C}$ est initial si et seulement si
$F(A) = \emptyset$,
\item un objet $Z$ de $\mathcal{C}$ est final si et seulement si
$F(Z)$ est un ensemble à un élément,
\item si $X$ et $Y$ sont connexes, alors $X \to Y$ est un épimorphisme,
\item
\label{pione-item-one-element}
si $X$ est connexe et si $a, b : X \to Y$ sont deux morphismes,
alors $a = b$ dès que $F(a)$ et $F(b)$ coïncident sur un élément de $F(X)$,
\item si $X = \coprod_{i = 1, \ldots, n} X_i$ et
$Y = \coprod_{j = 1, \ldots, m} Y_j$, où les $X_i$, $Y_j$ sont connexes,
alors il existe une application $\alpha : \{1, \ldots, n\} \to \{1, \ldots, m\}$
telle que $X \to Y$ provienne d'une famille de morphismes
$X_i \to Y_{\alpha(i)}$.
\end{enumerate}
\end{lemma}

\begin{proof}
Démonstration de (1). Supposons que $a, b : X \to Y$ vérifient $F(a) = F(b)$.
Soit $E$ l'égalisateur de $a$ et $b$. Alors $F(E) = F(X)$,
et l'on a $E = X$ puisque $F$ reflète les isomorphismes.

\medskip\noindent
Démonstration de (2). Cela vient de ce que $F$ transforme le morphisme $X \to X \times_Y X$
en l'application $F(X) \to F(X) \times_{F(Y)} F(X)$ et que $F$ reflète les isomorphismes.

\medskip\noindent
Démonstration de (3). Cela vient de ce que $F$ transforme le morphisme $Y \amalg_X Y \to Y$
en l'application $F(Y) \amalg_{F(X)} F(Y) \to F(Y)$ et que $F$ reflète les isomorphismes.

\medskip\noindent
Démonstration de (4). Il existe un objet initial $A$ et, bien entendu,
$F(A) = \emptyset$. Réciproquement, si $X$ est un objet tel que
$F(X) = \emptyset$, alors l'unique morphisme $A \to X$ induit une bijection
$F(A) \to F(X)$ ; par conséquent, $A \to X$ est un isomorphisme.

\medskip\noindent
Démonstration de (5). Il existe un objet final $Z$ et, bien entendu,
$F(Z)$ est un ensemble à un élément. Réciproquement, si $X$ est un objet tel que
$F(X)$ soit un ensemble à un élément, alors l'unique morphisme $X \to Z$ induit une bijection
$F(X) \to F(Z)$ ; par conséquent, $X \to Z$ est un isomorphisme.

\medskip\noindent
Démonstration de (6). L'égalisateur $E$ des deux morphismes $Y \to Y \amalg_X Y$ n'est pas
un objet initial de $\mathcal{C}$, car $X \to Y$ se factorise par $E$
et $F(X) \not = \emptyset$. Ainsi $E = Y$, d'où la conclusion.

\medskip\noindent
Démonstration de (\ref{pione-item-one-element}).
L'égalisateur $E$ de $a$ et $b$ est muni d'un monomorphisme
$E \to X$, et $F(E) \subset F(X)$ est l'ensemble des éléments sur lesquels
$F(a)$ et $F(b)$ coïncident. Pour conclure, on utilise le fait que $E$ est initial,
ou que $E = X$.

\medskip\noindent
Démonstration de (8). Pour tous $i, j$, on voit que $E_{ij} = X_i \times_Y Y_j$
est soit initial, soit égal à $X_i$. En choisissant $s \in F(X_i)$,
on voit que $E_{ij} = X_i$ si et seulement si $s$ s'envoie sur un élément
de $F(Y_j) \subset F(Y)$, ce qui arrive pour un unique $j = \alpha(i)$.
\end{proof}

\noindent
Le lemme précédent montre que, pour tout objet connexe $X$ d'une
catégorie galoisienne $(\mathcal{C}, F)$, le groupe d'automorphismes
$\text{Aut}(X)$ est d'ordre au plus $|F(X)|$. En effet, si $s \in F(X)$
et $g \in \text{Aut}(X)$, alors $F(g)(s) = s$ si et seulement
si $g = \text{id}_X$, d'après (\ref{pione-item-one-element}).
On dit que $X$ est {\it galoisien} si l'égalité a lieu.
De manière équivalente, $X$ est galoisien s'il est connexe et si
$\text{Aut}(X)$ agit transitivement sur $F(X)$.

\begin{lemma}
\label{pione-lemma-galois}
Soit $(\mathcal{C}, F)$ une catégorie galoisienne. Pour tout objet connexe $X$
de $\mathcal{C}$, il existe un objet galoisien $Y$ et un morphisme $Y \to X$.
\end{lemma}

\begin{proof}
Nous utiliserons sans autre mention les résultats du lemme \ref{pione-lemma-epi-mono}.
Posons $n = |F(X)|$. Considérons $X^n$ muni de son action naturelle de
$S_n$. Soit
$$
X^n = \coprod\nolimits_{t \in T} Z_t
$$
sa décomposition en objets connexes. Choisissons $t$ de sorte que
$F(Z_t)$ contienne $(s_1, \ldots, s_n)$, où les $s_i$ sont deux à deux distincts.
Si $(s'_1, \ldots, s'_n) \in F(Z_t)$ est un autre élément, nous
affirmons que les $s'_i$ sont eux aussi deux à deux distincts. Dans le cas contraire, disons
$s'_i = s'_j$ ; alors $Z_t$ est l'image d'une composante connexe de
$X^{n - 1}$ par le morphisme diagonal
$$
\Delta_{ij} : X^{n - 1} \longrightarrow X^n
$$
Puisque les morphismes entre objets connexes sont des épimorphismes et induisent
des surjections après application de $F$, on en déduirait que $s_i = s_j$,
ce qui n'est pas le cas.

\medskip\noindent
Soit $G \subset S_n$ le sous-groupe des éléments tels que $g(Z_t) = Z_t$.
En considérant l'action de $S_n$ sur
$$
F(X)^n = F(X^n) = \coprod\nolimits_{t' \in T} F(Z_{t'})
$$
on voit que $G = \{g \in S_n \mid g(s_1, \ldots, s_n) \in F(Z_t)\}$.
Choisissons maintenant un second élément $(s'_1, \ldots, s'_n) \in F(Z_t)$.
Nous avons vu plus haut que les $s'_i$ sont deux à deux distincts. On peut donc
trouver $g \in S_n$ tel que $g(s_1, \ldots, s_n) = (s'_1, \ldots, s'_n)$.
Autrement dit, l'action de $G$ sur $F(Z_t)$ est transitive, ce qui
achève la démonstration.
\end{proof}

\noindent
Voici un lemme essentiel.

\begin{lemma}
\label{pione-lemma-tame}
\begin{reference}
Comparer avec \cite[Définition 7.2.4]{BS}.
\end{reference}
Soit $(\mathcal{C}, F)$ une catégorie galoisienne. Soit $G = \text{Aut}(F)$
comme dans l'exemple \ref{pione-example-from-C-F-to-G-sets}. Pour tout objet connexe
$X$ de $\mathcal{C}$, l'action de $G$ sur $F(X)$ est transitive.
\end{lemma}

\begin{proof}
Nous utiliserons sans autre mention les résultats du lemme \ref{pione-lemma-epi-mono}.
Soit $I$ l'ensemble des classes d'isomorphie d'objets galoisiens de $\mathcal{C}$.
Pour chaque $i \in I$, soit $X_i$ un représentant de la classe d'isomorphie.
Choisissons $\gamma_i \in F(X_i)$ pour tout $i \in I$.
Nous munissons $I$ d'une relation d'ordre en posant $i \geq i'$ si
et seulement s'il existe un morphisme $f_{ii'} : X_i \to X_{i'}$.
Étant donné un tel morphisme, on peut le composer à gauche par un automorphisme
$X_{i'} \to X_{i'}$ de façon que $F(f_{ii'})(\gamma_i) = \gamma_{i'}$.
Avec cette normalisation, le morphisme $f_{ii'}$ est unique.
Remarquons que $I$ est un ensemble ordonné filtrant :
(Catégories, définition \ref{categories-definition-directed-set})
si $i_1, i_2 \in I$, il existe un objet galoisien $Y$ et un morphisme
$Y \to X_{i_1} \times X_{i_2}$ par le lemme \ref{pione-lemma-galois}, appliqué
à une composante connexe de $X_{i_1} \times X_{i_2}$.
Alors $Y \cong X_i$ pour un certain $i \in I$, et $i \geq i_1$, $i \geq i_2$.

\medskip\noindent
Nous affirmons que le foncteur $F$ est isomorphe au foncteur $F'$
qui à $X$ associe
$$
F'(X) = \colim_I \Mor_\mathcal{C}(X_i, X)
$$
au moyen de la transformation de foncteurs $t : F' \to F$ définie comme suit :
pour $f : X_i \to X$, on pose $t_X(f) = F(f)(\gamma_i)$.
En utilisant (\ref{pione-item-one-element}), on constate que $t_X$ est injective.
Pour montrer la surjectivité, soit $\gamma \in F(X)$. On peut aussitôt
se ramener au cas où $X$ est connexe, par la définition d'une
catégorie galoisienne. On peut ensuite supposer $X$ galoisien grâce au
lemme \ref{pione-lemma-galois}. Dans ce cas, $X$ est isomorphe à $X_i$
pour un certain $i$, et l'on peut choisir l'isomorphisme $X_i \to X$ de sorte
que $\gamma_i$ s'envoie sur $\gamma$ (par définition des objets galoisiens).
On conclut que $t$ est un isomorphisme.

\medskip\noindent
Posons $A_i = \text{Aut}(X_i)$.
Nous affirmons que, si $i \geq i'$, il existe une application canonique
$h_{ii'} : A_i \to A_{i'}$ telle que, pour tout $a \in A_i$,
le diagramme
$$
\xymatrix{
X_i \ar[d]_a \ar[r]_{f_{ii'}} & X_{i'} \ar[d]^{h_{ii'}(a)} \\
X_i \ar[r]^{f_{ii'}} & X_{i'}
}
$$
soit commutatif. Il suffit de définir $h_{ii'}(a) = a' : X_{i'} \to X_{i'}$
comme l'unique automorphisme tel que
$F(a')(\gamma_{i'}) = F(f_{ii'} \circ a)(\gamma_i)$.
Comme précédemment, cela rend le diagramme commutatif et, de plus, ce choix
est unique.
Il s'ensuit que
$h_{i'i''} \circ h_{ii'} = h_{ii''}$
si $i \geq i' \geq i''$.
Puisque $F(X_i) \to F(X_{i'})$ est surjective, on voit que
$A_i \to A_{i'}$ est surjective.
En passant à la limite projective, on obtient un groupe
$$
A = \lim_I A_i
$$
C'est un groupe profini, puisque les groupes d'automorphismes sont finis.
L'application $A \to A_i$ est surjective pour tout $i$, d'après
Catégories, lemme \ref{categories-lemma-nonempty-limit}.

\medskip\noindent
Puisque les éléments de $A$ agissent sur le système projectif des $X_i$, on obtient une action de
$A$ (à droite) sur $F'$ par précomposition. Autrement dit, on obtient
un homomorphisme $A^{opp} \to G$. Puisque $A \to A_i$ est surjective, on conclut
que $G$ agit transitivement sur $F(X_i)$ pour tout $i$. Comme tout objet connexe
est dominé par l'un des $X_i$, le lemme en découle.
\end{proof}

\begin{proposition}
\label{pione-proposition-galois}
\begin{reference}
C'est une version faible de \cite[Expos\'e V]{SGA1}.
La démonstration est empruntée à \cite[Théorème 7.2.5]{BS}.
\end{reference}
Soit $(\mathcal{C}, F)$ une catégorie galoisienne. Soit $G = \text{Aut}(F)$
comme dans l'exemple \ref{pione-example-from-C-F-to-G-sets}. Le foncteur
$F : \mathcal{C} \to \textit{Finite-}G\textit{-Ensembles}$
(\ref{pione-equation-remember}) est une équivalence.
\end{proposition}

\begin{proof}
Nous utiliserons sans plus les mentionner les résultats du lemme \ref{pione-lemma-epi-mono}.
En particulier, nous savons que le foncteur est fidèle.
Par le lemme \ref{pione-lemma-tame}, nous savons que pour tout $X$ connexe,
l'action de $G$ sur $F(X)$ est transitive. Ainsi, pour $X$ quelconque, le foncteur
$F$ fait correspondre les composantes connexes de $X$ (dont l'existence résulte
de l'un des axiomes d'une catégorie galoisienne) aux orbites de $G$ sur $F(X)$.
Soient $X$ et $Y$ des objets et soit
$s : F(X) \to F(Y)$ une application équivariante. Alors le graphe
$\Gamma_s \subset F(X) \times F(Y)$ de $s$
est une réunion de composantes connexes. Il existe donc une
réunion de composantes connexes $Z$ de $X \times Y$,
munie d'un monomorphisme $Z \to X \times Y$,
telle que $F(Z) = \Gamma_s$. Puisque $F(Z) \to F(X)$ est bijective,
nous voyons que $Z \to X$ est un isomorphisme et nous concluons
que $s = F(f)$, où $f : X \cong Z \to Y$ est le composé.
Ainsi, $F$ est pleinement fidèle.

\medskip\noindent
Pour achever la démonstration, montrons que $F$ est essentiellement surjectif.
Il suffit de montrer que $G/H$ appartient à l'image essentielle pour
tout sous-groupe ouvert $H \subset G$ d'indice fini.
Par définition de la topologie sur $G$, il existe une famille finie
d'objets $X_i$ telle que
$$
\Ker(G \longrightarrow \prod\nolimits_i \text{Aut}(F(X_i)))
$$
soit contenu dans $H$. Nous pouvons supposer que $X_i$ est connexe
pour tout $i$. Nous pouvons choisir un objet galoisien $Y$ muni d'un morphisme
vers une composante connexe de $\prod X_i$ à l'aide du
lemme \ref{pione-lemma-galois}. Choisissons un isomorphisme $F(Y) = G/U$
dans $G\textit{-sets}$ pour un certain sous-groupe ouvert $U \subset G$.
Comme $Y$ est galoisien, le groupe
$\text{Aut}(Y) = \text{Aut}_{G\textit{-Ensembles}}(G/U)$ opère transitivement
sur $F(Y) = G/U$. Cela implique que $U$ est distingué. Puisque
le morphisme induit de $F(Y)$ vers $F(X_i)$ est surjectif pour tout $i$, nous voyons que
$U \subset H$. Soit $M \subset \text{Aut}(Y)$ le sous-groupe fini
correspondant à
$$
(H/U)^{opp} \subset (G/U)^{opp} = \text{Aut}_{G\textit{-Ensembles}}(G/U)
= \text{Aut}(Y).
$$
Posons $X = Y/M$, c'est-à-dire que $X$ est le coégalisateur
des flèches $m : Y \to Y$, $m \in M$.
Comme $F$ est exact, nous voyons que $F(X) = G/H$, ce qui achève
la démonstration.
\end{proof}

\begin{lemma}
\label{pione-lemma-functoriality-galois}
Soient $(\mathcal{C}, F)$ et $(\mathcal{C}', F')$ des catégories galoisiennes.
Soit $H : \mathcal{C} \to \mathcal{C}'$ un foncteur exact.
Il existe un isomorphisme $t : F' \circ H \to F$.
Le choix de $t$ détermine un homomorphisme continu
$h : G' = \text{Aut}(F') \to \text{Aut}(F) = G$ et
un diagramme $2$-commutatif
$$
\xymatrix{
\mathcal{C} \ar[r]_H \ar[d] & \mathcal{C}' \ar[d] \\
\textit{Finite-}G\textit{-Ensembles} \ar[r]^h &
\textit{Finite-}G'\textit{-Ensembles}
}
$$
L'homomorphisme $h$ ne dépend du choix de $t$ qu'à
un automorphisme intérieur de $G$ près.
Réciproquement, étant donné un homomorphisme continu $h : G' \to G$, il
existe un foncteur exact $H : \mathcal{C} \to \mathcal{C}'$ et un
isomorphisme $t$ permettant de retrouver $h$ comme ci-dessus.
\end{lemma}

\begin{proof}
D'après la proposition \ref{pione-proposition-galois} et
le lemme \ref{pione-lemma-single-out-profinite}, nous pouvons supposer que
$\mathcal{C} = \textit{Finite-}G\textit{-Ensembles}$ et que $F$ est le
foncteur d'oubli, et de même pour $\mathcal{C}'$. L'existence de
$t$ résulte alors du lemme \ref{pione-lemma-second-fundamental-functor}. L'homomorphisme $h$
s'obtient par transport de structure au moyen de $t$. La commutativité du
diagramme est évidente. L'homomorphisme $h$ est unique à automorphisme intérieur
de $G$ près, car le choix de $t$ est
unique à un élément de $G$ près. La dernière assertion est immédiate.
\end{proof}





\section{Foncteurs et homomorphismes}
\label{pione-section-translation}

\noindent
Soient $(\mathcal{C}, F)$, $(\mathcal{C}', F')$, $(\mathcal{C}'', F'')$
des catégories galoisiennes. Posons $G = \text{Aut}(F)$, $G' = \text{Aut}(F')$ et
$G'' = \text{Aut}(F'')$. Soient $H : \mathcal{C} \to \mathcal{C}'$
et $H' : \mathcal{C}' \to \mathcal{C}''$ des foncteurs exacts.
Soient $h : G' \to G$ et $h' : G'' \to G'$ les homomorphismes
continus correspondants comme dans le lemme \ref{pione-lemma-functoriality-galois}.
Dans cette section, nous considérons le diagramme $2$-commutatif correspondant
\begin{equation}
\label{pione-equation-translation}
\vcenter{
\xymatrix{
\mathcal{C} \ar[r]_H \ar[d] &
\mathcal{C}' \ar[r]_{H'} \ar[d] &
\mathcal{C}'' \ar[d] \\
\textit{Finite-}G\textit{-Ensembles} \ar[r]^h &
\textit{Finite-}G'\textit{-Ensembles} \ar[r]^{h'} &
\textit{Finite-}G''\textit{-Ensembles}
}
}
\end{equation}
et nous relions les propriétés d'exactitude de la suite
$1 \to G'' \to G' \to G \to 1$ aux propriétés des foncteurs $H$ et $H'$.

\begin{lemma}
\label{pione-lemma-functoriality-galois-surjective}
Dans le diagramme (\ref{pione-equation-translation}), les assertions suivantes sont équivalentes :
\begin{enumerate}
\item $h : G' \to G$ est surjectif,
\item $H : \mathcal{C} \to \mathcal{C}'$ est pleinement fidèle,
\item si $X \in \Ob(\mathcal{C})$ est connexe, alors $H(X)$ est connexe,
\item si $X \in \Ob(\mathcal{C})$ est connexe et s'il existe
un morphisme $*' \to H(X)$ dans $\mathcal{C}'$, alors
il existe un morphisme $* \to X$, et
\item pour tout objet $X$ de $\mathcal{C}$, l'application
$\Mor_\mathcal{C}(*, X) \to \Mor_{\mathcal{C}'}(*', H(X))$
est bijective.
\end{enumerate}
Ici, $*$ et $*'$ sont des objets finaux de $\mathcal{C}$ et $\mathcal{C}'$.
\end{lemma}

\begin{proof}
Les implications (5) $\Rightarrow$ (4) et (2) $\Rightarrow$ (5) sont évidentes.

\medskip\noindent
Supposons (3). Soit $X$ un objet connexe de $\mathcal{C}$ et soit
$*' \to H(X)$ un morphisme. Comme $H(X)$ est connexe d'après (3),
nous voyons que $*' \to H(X)$ est un isomorphisme. Ainsi, le $G'$-ensemble
correspondant à $H(X)$ possède exactement un élément, ce qui signifie que le
$G$-ensemble correspondant à $X$ possède un seul élément, donc que $X$ est
isomorphe à l'objet final de $\mathcal{C}$ ; en particulier,
il existe un morphisme $* \to X$. Nous avons ainsi montré que (3) $\Rightarrow$ (4).

\medskip\noindent
Si (1) est vraie, alors le foncteur
$\textit{Finite-}G\textit{-Ensembles} \to \textit{Finite-}G'\textit{-Ensembles}$
est pleinement fidèle : en effet, une application entre ensembles munis d'une action de $G$ commute avec
l'action de $G$ si et seulement si elle commute avec l'action de $G'$.
Ainsi, (1) $\Rightarrow$ (2).

\medskip\noindent
Si (1) est vraie, alors, pour un $G$-ensemble $X$, les orbites de $G$ et de $G'$
coïncident. Ainsi, (1) $\Rightarrow$ (3).

\medskip\noindent
Pour achever la démonstration, il suffit de montrer que (4) implique (1).
Si (1) est fausse, c'est-à-dire si $h$ n'est pas surjectif, il existe alors
un sous-groupe ouvert $U \subset G$ contenant $h(G')$ et distinct
de $G$. Le $G$-ensemble fini $M = G/U$ est alors muni d'une action transitive,
mais $G'$ y possède un point fixe. L'objet $X$ de $\mathcal{C}$
correspondant à $M$ contredirait (4). Nous voyons ainsi que
(4) $\Rightarrow$ (1), ce qui achève la démonstration.
\end{proof}

\begin{lemma}
\label{pione-lemma-composition-trivial}
Dans le diagramme (\ref{pione-equation-translation}), les assertions suivantes sont équivalentes :
\begin{enumerate}
\item $h \circ h'$ est trivial, et
\item l'image de $H' \circ H$ est formée d'objets isomorphes à des coproduits
finis d'objets finaux.
\end{enumerate}
\end{lemma}

\begin{proof}
Nous pouvons remplacer $H$ et $H'$ par les foncteurs canoniques
$\textit{Finite-}G\textit{-Ensembles} \to \textit{Finite-}G'\textit{-Ensembles}
\to \textit{Finite-}G''\textit{-Ensembles}$ déterminés par $h$ et $h'$.
Il s'agit alors de dire que l'action de $G''$ sur tout $G$-ensemble est triviale
si et seulement si l'homomorphisme $G'' \to G$ est trivial. Cela est évident.
\end{proof}

\begin{lemma}
\label{pione-lemma-functoriality-galois-ses}
Dans le diagramme (\ref{pione-equation-translation}), les assertions suivantes sont équivalentes :
\begin{enumerate}
\item la suite $G'' \xrightarrow{h'} G' \xrightarrow{h} G \to 1$
est exacte au sens suivant : $h$ est surjectif, $h \circ h'$ est trivial,
et $\Ker(h)$ est le plus petit sous-groupe fermé distingué contenant $\Im(h')$,
\item $H$ est pleinement fidèle et un objet $X'$ de $\mathcal{C}'$ appartient à
l'image essentielle de $H$ si et seulement si $H'(X')$ est isomorphe à un
coproduit fini d'objets finaux, et
\item $H$ est pleinement fidèle, $H' \circ H$ envoie tout objet sur un coproduit
fini d'objets finaux, et, pour tout objet $X'$ de $\mathcal{C}'$
tel que $H'(X')$ soit un coproduit fini d'objets finaux, il existe
un objet $X$ de $\mathcal{C}$ et un épimorphisme $H(X) \to X'$.
\end{enumerate}
\end{lemma}

\begin{proof}
D'après les lemmes \ref{pione-lemma-functoriality-galois-surjective} et
\ref{pione-lemma-composition-trivial}, nous pouvons supposer que
$H$ est pleinement fidèle, que $h$ est surjectif, que $H' \circ H$ envoie les
objets sur des coproduits de copies de l'objet final et que $h \circ h'$
est trivial. Soit $N \subset G'$ le plus petit sous-groupe fermé distingué
contenant l'image de $h'$. Il est clair que
$N \subset \Ker(h)$.
Nous pouvons supposer que les foncteurs $H$ et $H'$ sont les foncteurs canoniques
$\textit{Finite-}G\textit{-Ensembles} \to \textit{Finite-}G'\textit{-Ensembles}
\to \textit{Finite-}G''\textit{-Ensembles}$ déterminés par $h$ et $h'$.

\medskip\noindent
Supposons que (2) soit vraie. Cela signifie que, pour tout $G'$-ensemble fini $X'$
sur lequel $G''$ agit trivialement, l'action de $G'$ se factorise par $G$.
Appliquons ceci à $X' = G'/U'N$, où $U'$ est un sous-groupe ouvert suffisamment petit de $G'$.
Nous voyons alors que $\Ker(h) \subset U'N$ pour tout $U'$. Comme $N$ est fermé,
il s'ensuit que $\Ker(h) \subset N$, c'est-à-dire que (1) est vraie.

\medskip\noindent
Supposons que (1) soit vraie. Cela signifie que $N = \Ker(h)$. Soit $X'$ un
$G'$-ensemble fini sur lequel $G''$ agit trivialement. Cela signifie que
$\Ker(G' \to \text{Aut}(X'))$ est un sous-groupe fermé distingué contenant
$\Im(h')$. Ainsi, $N = \Ker(h)$ y est contenu et l'action de $G'$
sur $X'$ se factorise par $G$, c'est-à-dire que (2) est vraie.

\medskip\noindent
Supposons que (3) soit vraie. Cela signifie que, pour tout $G'$-ensemble fini $X'$
sur lequel $G''$ agit trivialement, il existe une surjection de $G'$-ensembles
$X \to X'$, où $X$ est un $G$-ensemble. Cela signifie manifestement que l'action de
$G'$ sur $X'$ se factorise par $G$, c'est-à-dire que (2) est vraie.

\medskip\noindent
L'implication (2) $\Rightarrow$ (3) est immédiate. Ceci achève la démonstration.
\end{proof}

\begin{lemma}
\label{pione-lemma-functoriality-galois-injective}
Dans le diagramme (\ref{pione-equation-translation}), les assertions suivantes sont équivalentes :
\begin{enumerate}
\item $h'$ est injectif, et
\item pour tout objet connexe $X''$ de $\mathcal{C}''$,
il existe un objet $X'$ de $\mathcal{C}'$ et un diagramme
$$
X'' \leftarrow Y'' \rightarrow H'(X')
$$
dans $\mathcal{C}''$, où $Y'' \to X''$ est un épimorphisme et
$Y'' \to H'(X')$ est un monomorphisme.
\end{enumerate}
\end{lemma}

\begin{proof}
Nous pouvons remplacer $H'$ par le foncteur correspondant entre les catégories
de $G'$-ensembles finis et de $G''$-ensembles finis.

\medskip\noindent
Supposons que $h' : G'' \to G'$ soit injectif. Soit $H'' \subset G''$
un sous-groupe ouvert. Puisque la topologie sur $G''$ est la topologie induite
par celle de $G'$, il existe un sous-groupe ouvert $H' \subset G'$
tel que $(h')^{-1}(H') \subset H''$.
Alors le diagramme voulu est
$$
G''/H'' \leftarrow G''/(h')^{-1}(H') \rightarrow G'/H'
$$
Réciproquement, supposons que (2) soit satisfaite pour le foncteur
$\textit{Finite-}G'\textit{-Ensembles} \to \textit{Finite-}G''\textit{-Ensembles}$.
Soit $g'' \in \Ker(h')$. Choisissons un sous-groupe ouvert quelconque $H'' \subset G''$.
Par hypothèse, il existe un $G'$-ensemble fini $X'$ et un diagramme
$$
G''/H'' \leftarrow Y'' \rightarrow X'
$$
de $G''$-ensembles, dont la flèche de gauche est surjective et celle de droite injective.
Comme $g''$ appartient au noyau de $h'$, nous voyons que $g''$ agit trivialement sur $X'$.
Ainsi, $g''$ agit trivialement sur $Y''$, et donc aussi sur $G''/H''$.
Par conséquent, $g'' \in H''$. Comme ceci vaut pour tout sous-groupe ouvert, nous concluons
que $g''$ est l'élément neutre, comme voulu.
\end{proof}

\begin{lemma}
\label{pione-lemma-functoriality-galois-normal}
Dans le diagramme (\ref{pione-equation-translation}), les assertions suivantes sont équivalentes :
\begin{enumerate}
\item l'image de $h'$ est distinguée, et
\item pour tout objet connexe $X'$ de $\mathcal{C}'$ tel
qu'il existe un morphisme de l'objet final de $\mathcal{C}''$
vers $H'(X')$, l'objet $H'(X')$ est isomorphe à un coproduit fini
d'objets finaux.
\end{enumerate}
\end{lemma}

\begin{proof}
Cela se traduit par l'énoncé suivant pour l'homomorphisme continu
de groupes $h' : G'' \to G'$ : l'image de $h'$ est distinguée
si et seulement si tout sous-groupe ouvert $U' \subset G'$ qui
contient $h'(G'')$ contient aussi chacun des conjugués de $h'(G'')$.
Le résultat s'en déduit aisément ; nous omettons certains détails.
\end{proof}






\section{Morphismes finis \'etales}
\label{pione-section-finite-etale}

\noindent
Dans cette section, nous établissons suffisamment de résultats élémentaires sur les morphismes finis \'etales
pour pouvoir construire le groupe fondamental \'etale.

\medskip\noindent
Soit $X$ un schéma. Nous noterons $\textit{F\'Et}_X$
la catégorie des schémas finis et \'etales sur $X$.
Ainsi,
\begin{enumerate}
\item un objet de $\textit{F\'Et}_X$ est un morphisme fini \'etale
$Y \to X$ ayant pour but $X$, et
\item un morphisme dans $\textit{F\'Et}_X$
de $Y \to X$ vers $Y' \to X$ est un morphisme $Y \to Y'$ qui rend
commutatif le diagramme
$$
\xymatrix{
Y \ar[rr] \ar[rd] & &  Y' \ar[ld] \\
& X
}
$$
ci-dessus.
\end{enumerate}
Nous appellerons souvent un objet de $\textit{F\'Et}_X$ un
{\it revêtement \'etale fini} de $X$ (même si $Y$ est vide).
Il existe en fait un champ $p : \textit{F\'Et} \to \Sch$
sur la catégorie des schémas, dont la fibre au-dessus de $X$ est la catégorie
$\textit{F\'Et}_X$ que nous venons de définir. Voir Exemples de champs, section
\ref{examples-stacks-section-finite-etale}.

\begin{example}
\label{pione-example-finite-etale-geometric-point}
Soient $k$ un corps algébriquement clos et $X = \Spec(k)$. Dans ce cas,
$\textit{F\'Et}_X$ est équivalente à la catégorie des ensembles finis. Ceci vaut
plus généralement lorsque $k$ est séparablement clos. La raison en est que
un schéma \'etale sur $k$ est une réunion disjointe de spectres de
corps qui sont des extensions finies séparables de $k$ ; voir
Morphismes, lemme \ref{morphisms-lemma-etale-over-field}.
\end{example}

\begin{lemma}
\label{pione-lemma-finite-etale-covers-limits-colimits}
Soit $X$ un schéma. La catégorie $\textit{F\'Et}_X$ admet les limites finies et
les colimites finies, et, pour tout morphisme $X' \to X$, le foncteur de changement de base
$\textit{F\'Et}_X \to \textit{F\'Et}_{X'}$ est exact.
\end{lemma}

\begin{proof}
Limites finies et exactitude à gauche. D'après
Catégories, lemme \ref{categories-lemma-finite-limits-exist},
il suffit de montrer que $\textit{F\'Et}_X$ admet un objet final
et des produits fibrés. Cela est clair, car la catégorie de
tous les schémas sur $X$ admet un objet final (à savoir $X$) et des produits fibrés.
De plus, les produits fibrés de schémas finis \'etales sur $X$ sont
finis \'etales sur $X$. Il est en outre clair que le changement
de base commute à ces opérations et que, par conséquent, le changement de base
est exact à gauche (Catégories, lemme
\ref{categories-lemma-characterize-left-exact}).

\medskip\noindent
Colimites finies et exactitude à droite. D'après
Catégories, lemme \ref{categories-lemma-colimits-exist},
il suffit de montrer que $\textit{F\'Et}_X$ admet des sommes
finies et des coégalisateurs. Les sommes finies sont données
par les réunions disjointes (la somme vide est le schéma vide).
Soient $a, b : Z \to Y$ deux morphismes de $\textit{F\'Et}_X$.
Comme $Z \to X$ et $Y \to X$ sont finis \'etales, nous pouvons écrire
$Z = \underline{\Spec}(\mathcal{C})$ et $Y = \underline{\Spec}(\mathcal{B})$
pour certaines $\mathcal{O}_X$-algèbres localement libres de rang fini $\mathcal{C}$
et $\mathcal{B}$. Les morphismes $a, b$ induisent deux homomorphismes
$a^\sharp, b^\sharp : \mathcal{B} \to \mathcal{C}$.
Soit $\mathcal{A} = \text{Eq}(a^\sharp, b^\sharp)$ leur
égalisateur. Si
$$
\underline{\Spec}(\mathcal{A}) \longrightarrow X
$$
est fini \'etale, il est clair qu'il s'agit du coégalisateur
(en effet, tout objet de $\textit{F\'Et}_X$ peut s'écrire
comme le spectre relatif d'un faisceau de $\mathcal{O}_X$-algèbres).
Nous pouvons le vérifier après avoir remplacé $X$ par les éléments d'un recouvrement \'etale
(Descente, lemmes \ref{descent-lemma-descending-property-finite}
et \ref{descent-lemma-descending-property-etale}).
D'après Morphismes \'etales, lemme \ref{etale-lemma-finite-etale-etale-local},
nous pouvons supposer que
$Y = \coprod_{i = 1, \ldots, n} X$ et $Z = \coprod_{j = 1, \ldots, m} X$.
Alors
$$
\mathcal{C} = \prod\nolimits_{1 \leq j \leq m} \mathcal{O}_X
\quad\text{et}\quad
\mathcal{B} = \prod\nolimits_{1 \leq i \leq n} \mathcal{O}_X
$$
Après un nouveau remplacement par les éléments d'un recouvrement ouvert,
nous pouvons supposer que $a, b$ correspondent à des
applications $a_s, b_s : \{1, \ldots, m\} \to \{1, \ldots, n\}$, c'est-à-dire que
la composante $X$ de $Z$ d'indice $j$ est envoyée dans
la composante $X$ de $Y$ d'indice $a_s(j)$, resp.\ $b_s(j)$,
par le morphisme $a$, resp.\ $b$.
Soit $\{1, \ldots, n\} \to T$ le coégalisateur de $a_s, b_s$.
Alors nous voyons que
$$
\mathcal{A} = \prod\nolimits_{t \in T} \mathcal{O}_X
$$
dont le spectre est bien entendu fini \'etale sur $X$. Nous
omettons la vérification de la compatibilité avec le changement de base.
Ainsi, le changement de base est un foncteur exact à droite.
\end{proof}

\begin{remark}
\label{pione-remark-colimits-commute-forgetful}
Soit $X$ un schéma. Considérons les foncteurs naturels
$F_1 : \textit{F\'Et}_X \to \Sch$ et $F_2 : \textit{F\'Et}_X \to \Sch/X$.
Alors
\begin{enumerate}
\item Les foncteurs $F_1$ et $F_2$ commutent aux colimites finies.
\item Le foncteur $F_2$ commute aux limites finies,
\item Le foncteur $F_1$ commute aux limites finies connexes, c'est-à-dire
aux égalisateurs et aux produits fibrés.
\end{enumerate}
Les assertions sur les limites résultent immédiatement de la discussion dans
la démonstration du lemme \ref{pione-lemma-finite-etale-covers-limits-colimits}
et de Catégories, lemme \ref{categories-lemma-connected-limit-over-X}.
Il est clair que $F_1$ et $F_2$ commutent aux sommes finies.
D'après l'énoncé dual de Catégories, lemme
\ref{categories-lemma-characterize-left-exact},
il suffit de montrer que $F_1$ et $F_2$ commutent aux coégalisateurs.
Dans la démonstration du lemme \ref{pione-lemma-finite-etale-covers-limits-colimits},
nous avons vu que les coégalisateurs dans $\textit{F\'Et}_X$ prennent localement pour la topologie \'etale
la forme suivante :
$$
\xymatrix{
\coprod_{j \in J} U \ar@<1ex>[r]^a \ar@<-1ex>[r]_b &
\coprod_{i \in I} U \ar[r] &
\coprod_{t \in \text{Coeq}(a, b)} U
}
$$
ce qui est manifestement un coégalisateur dans la catégorie des schémas.
L'assertion découle donc du fait que la propriété d'être un coégalisateur
est locale pour la topologie fpqc, sous la forme précise donnée dans
Descente, lemme \ref{descent-lemma-coequalizer-fpqc-local}.
\end{remark}

\begin{lemma}
\label{pione-lemma-internal-hom-finite-etale}
Soit $X$ un schéma. Étant donnés $U, V$ finis \'etales sur $X$, il
existe un schéma $W$ fini \'etale sur $X$ tel que
$$
\Mor_X(X, W) = \Mor_X(U, V)
$$
et tel que cette égalité reste vraie après tout changement de base.
\end{lemma}

\begin{proof}
D'après Compléments sur les morphismes, lemme
\ref{more-morphisms-lemma-hom-from-finite-locally-free-separated-lqf},
il existe un schéma $W$ qui représente $\mathit{Mor}_X(U, V)$.
(On utilise le fait qu'un morphisme \'etale est localement quasi-fini d'après
Morphismes, lemme \ref{morphisms-lemma-etale-locally-quasi-finite},
et qu'un morphisme fini est séparé.)
Ce schéma satisfait manifestement la formule après tout changement de base.
Pour achever la démonstration, il reste à montrer que $W \to X$ est fini \'etale.
Nous pouvons le faire après avoir remplacé $X$ par les éléments d'un recouvrement \'etale
(Descente, lemmes \ref{descent-lemma-descending-property-finite}
et \ref{descent-lemma-descending-property-etale}).
D'après Morphismes \'etales, lemme \ref{etale-lemma-finite-etale-etale-local},
nous pouvons supposer que $U = \coprod_{i = 1, \ldots, n} X$
et $V = \coprod_{j = 1, \ldots, m} X$.
Dans ce cas,
$W = \coprod_{\alpha : \{1, \ldots, n\} \to \{1, \ldots, m\}} X$
immédiatement (nous omettons les détails), ce qui achève la démonstration.
\end{proof}

\noindent
Soit $X$ un schéma. Un {\it point géométrique} de $X$ est un morphisme
$\Spec(k) \to X$, où $k$ est algébriquement clos. Un tel point est
généralement noté $\overline{x}$, c'est-à-dire par une lettre minuscule surlignée.
Nous utilisons souvent $\overline{x}$ pour désigner aussi bien le schéma $\Spec(k)$ que
le morphisme, et nous notons $\kappa(\overline{x})$
le corps $k$. Nous disons que $\overline{x}$ {\it est au-dessus de} $x$
pour indiquer que $x \in X$ est l'image de $\overline{x}$.
Nous reviendrons plus en détail sur cette notion dans
Cohomologie \'etale, section \ref{etale-cohomology-section-stalks}.
Étant donnés $\overline{x}$ et un morphisme \'etale $U \to X$, nous pouvons
considérer
$$
|U_{\overline{x}}| : \text{l'ensemble sous-jacent des points du
schéma }U_{\overline{x}} = U \times_X \overline{x}
$$
Puisque le schéma $U_{\overline{x}}$, considéré comme schéma sur $\overline{x}$,
est une réunion disjointe de copies de $\overline{x}$
(Morphismes, lemme \ref{morphisms-lemma-etale-over-field}),
nous pouvons également décrire cet ensemble comme suit :
$$
|U_{\overline{x}}| =
\left\{
\begin{matrix}
\text{diagrammes} \\
\text{commutatifs}
\end{matrix}
\vcenter{
\xymatrix{
\overline{x} \ar[rd]_{\overline{x}} \ar[r]_{\overline{u}} & U \ar[d] \\
& X
}
}
\right\}
$$
L'association $U \mapsto |U_{\overline{x}}|$ est un foncteur
souvent noté $F_{\overline{x}}$.

\begin{lemma}
\label{pione-lemma-finite-etale-connected-galois-category}
Soit $X$ un schéma connexe. Soit $\overline{x}$ un point géométrique.
Le foncteur
$$
F_{\overline{x}} : \textit{F\'Et}_X \longrightarrow \textit{Ensembles},\quad
Y \longmapsto |Y_{\overline{x}}|
$$
définit une catégorie galoisienne (définition \ref{pione-definition-galois-category}).
\end{lemma}

\begin{proof}
Après avoir identifié $\textit{F\'Et}_{\overline{x}}$ à la catégorie des
ensembles finis (exemple \ref{pione-example-finite-etale-geometric-point}),
on voit que notre foncteur $F_{\overline{x}}$ n'est autre que
le foncteur de changement de base pour le morphisme $\overline{x} \to X$.
Ainsi, $\textit{F\'Et}_X$ admet des limites finies et des colimites finies,
et $F_{\overline{x}}$ est exact d'après le
lemme \ref{pione-lemma-finite-etale-covers-limits-colimits}.
Nous utiliserons aussi le fait que les limites finies dans $\textit{F\'Et}_X$
coïncident avec les limites finies correspondantes dans la catégorie
des schémas sur $X$ ; voir la remarque \ref{pione-remark-colimits-commute-forgetful}.

\medskip\noindent
Si $Y' \to Y$ est un monomorphisme dans $\textit{F\'Et}_X$,
alors $Y' \to Y' \times_Y Y'$ est un isomorphisme ; par
conséquent, $Y' \to Y$ est un monomorphisme de schémas. Il s'ensuit que
$Y' \to Y$ est une immersion ouverte
(Morphismes étales, théorème \ref{etale-theorem-etale-radicial-open}). Puisque
$Y'$ est fini sur $X$ et que $Y$ est séparé sur $X$,
le morphisme $Y' \to Y$ est fini
(Morphismes, lemme \ref{morphisms-lemma-finite-permanence}), donc fermé
(Morphismes, lemme \ref{morphisms-lemma-finite-proper}) ;
c'est donc l'inclusion d'un sous-schéma ouvert et fermé de $Y$.
Il s'ensuit que $Y$ est un objet connexe de la catégorie
$\textit{F\'Et}_X$ (au sens de la définition \ref{pione-definition-galois-category})
si et seulement si $Y$ est connexe comme schéma. Le
lemme \ref{topology-lemma-finite-fibre-connected-components} de Topologie
montre alors que $Y$ est le coproduit fini de ses composantes connexes,
tant comme schéma qu'au sens de la
définition \ref{pione-definition-galois-category}.

\medskip\noindent
Soit $Y \to Z$ un morphisme de $\textit{F\'Et}_X$ qui induit une
bijection $F_{\overline{x}}(Y) \to F_{\overline{x}}(Z)$. Il faut
montrer que $Y \to Z$ est un isomorphisme. D'après ce qui précède, on peut
supposer $Z$ connexe. Puisque $Y \to Z$ est fini étale, donc fini
localement libre, il suffit de montrer que $Y \to Z$ est fini localement
libre de degré $1$. Cela est vrai au voisinage de tout point de
$Z$ situé au-dessus de $\overline{x}$ ; comme $Z$ est connexe et que
le degré est localement constant, on conclut.
\end{proof}



\section{Groupes fondamentaux}
\label{pione-section-fundamental-groups}

\noindent
Dans cette section, nous définissons le groupe fondamental algébrique de Grothendieck.
La définition suivante a un sens grâce au
lemme \ref{pione-lemma-finite-etale-connected-galois-category}.

\begin{definition}
\label{pione-definition-fundamental-group}
Soit $X$ un schéma connexe. Soit $\overline{x}$ un point géométrique
de $X$. Le {\it groupe fondamental} de $X$, de
{\it point-base} $\overline{x}$, est le groupe
$$
\pi_1(X, \overline{x}) = \text{Aut}(F_{\overline{x}})
$$
des automorphismes du foncteur fibre
$F_{\overline{x}} : \textit{F\'Et}_X \to \textit{Ensembles}$,
muni de sa topologie profinie canonique définie dans le
lemme \ref{pione-lemma-aut-inverse-limit}.
\end{definition}

\noindent
En combinant ce qui précède avec les résultats de la section \ref{pione-section-galois},
on obtient le théorème suivant.

\begin{theorem}
\label{pione-theorem-fundamental-group}
Soit $X$ un schéma connexe. Soit $\overline{x}$ un point géométrique
de $X$.
\begin{enumerate}
\item Le foncteur fibre $F_{\overline{x}}$ définit une équivalence de
catégories
$$
\textit{F\'Et}_X \longrightarrow
\textit{Finite-}\pi_1(X, \overline{x})\textit{-Ensembles}
$$
\item Étant donné un second point géométrique $\overline{x}'$ de $X$, il
existe un isomorphisme $t : F_{\overline{x}} \to F_{\overline{x}'}$.
Celui-ci donne un isomorphisme $\pi_1(X, \overline{x}) \to \pi_1(X, \overline{x}')$
compatible avec les équivalences de (1). Cet isomorphisme est
indépendant de $t$, à automorphisme intérieur près.
\item Étant donné un morphisme $f : X \to Y$ de schémas connexes, posons
$\overline{y} = f \circ \overline{x}$. Il existe un homomorphisme
continu canonique
$$
f_* : \pi_1(X, \overline{x}) \to \pi_1(Y, \overline{y})
$$
tel que le diagramme
$$
\xymatrix{
\textit{F\'Et}_Y \ar[r]_{\text{changement de base}} \ar[d]_{F_{\overline{y}}} &
\textit{F\'Et}_X \ar[d]^{F_{\overline{x}}} \\
\textit{Finite-}\pi_1(Y, \overline{y})\textit{-Ensembles} \ar[r]^{f_*} &
\textit{Finite-}\pi_1(X, \overline{x})\textit{-Ensembles}
}
$$
soit commutatif.
\end{enumerate}
\end{theorem}

\begin{proof}
L'assertion (1) résulte du lemme \ref{pione-lemma-finite-etale-connected-galois-category}
et de la proposition \ref{pione-proposition-galois}.
L'assertion (2) est un cas particulier du lemme \ref{pione-lemma-functoriality-galois}.
Pour (3), observons que le diagramme
$$
\xymatrix{
\textit{F\'Et}_Y \ar[r] \ar[d]_{F_{\overline{y}}} &
\textit{F\'Et}_X \ar[d]^{F_{\overline{x}}} \\
\textit{Ensembles} \ar@{=}[r] & \textit{Ensembles}
}
$$
est commutatif (au sens strict, et pas seulement $2$-commutatif), car
$\overline{y} = f \circ \overline{x}$. On peut donc
appliquer le lemme \ref{pione-lemma-functoriality-galois}, avec la transformation
de foncteurs qui s'en déduit, pour obtenir (3).
\end{proof}

\begin{lemma}
\label{pione-lemma-fundamental-group-Galois-group}
Soit $K$ un corps et posons $X = \Spec(K)$. Soit $\overline{K}$ une
clôture algébrique et notons $\overline{x} : \Spec(\overline{K}) \to X$
le point géométrique correspondant. Soit $K^{sep} \subset \overline{K}$
la clôture séparable.
\begin{enumerate}
\item Le foncteur du lemme \ref{pione-lemma-sheaves-point} induit une équivalence
$$
\textit{F\'Et}_X \longrightarrow
\textit{Finite-}\text{Gal}(K^{sep}/K)\textit{-Ensembles}.
$$
compatible avec $F_{\overline{x}}$ et avec le foncteur
$\textit{Finite-}\text{Gal}(K^{sep}/K)\textit{-Ensembles} \to \textit{Ensembles}$.
\item Elle induit un isomorphisme canonique
$$
\text{Gal}(K^{sep}/K) \longrightarrow \pi_1(X, \overline{x})
$$
de groupes topologiques profinis.
\end{enumerate}
\end{lemma}

\begin{proof}
Le foncteur du lemme \ref{pione-lemma-sheaves-point} est le même que le foncteur
$F_{\overline{x}}$, car pour tout $Y$ étale sur $X$ on a
$$
\Mor_X(\Spec(\overline{K}), Y) = \Mor_X(\Spec(K^{sep}), Y)
$$
En effet, comme on l'a vu dans la preuve du lemme \ref{pione-lemma-sheaves-point},
$Y = \coprod_{i \in I} \Spec(L_i)$, où $L_i/K$ est une extension séparable finie de $K$.
Par conséquent, tout homomorphisme de $K$-algèbres $L_i \to \overline{K}$ se factorise
par $K^{sep}$. Notons aussi que $F_{\overline{x}}(Y)$ est fini
si et seulement si $I$ est fini, si et seulement si $Y \to X$ est fini étale.
Ceci démontre (1).

\medskip\noindent
L'assertion (2) est une conséquence formelle de (1), du
lemme \ref{pione-lemma-functoriality-galois} et du
lemme \ref{pione-lemma-single-out-profinite}.
(Voir aussi la remarque ci-dessous.)
\end{proof}

\begin{remark}
\label{pione-remark-variance}
Dans la situation du lemme \ref{pione-lemma-fundamental-group-Galois-group},
donnons une construction plus explicite de l'isomorphisme
$\text{Gal}(K^{sep}/K) \to
\pi_1(X, \overline{x}) = \text{Aut}(F_{\overline{x}})$.
Observons que
$\text{Gal}(K^{sep}/K) = \text{Aut}(\overline{K}/K)$,
car $\overline{K}$ est la clôture parfaite de $K^{sep}$.
Puisque $F_{\overline{x}}(Y) = \Mor_X(\Spec(\overline{K}), Y)$,
on peut considérer l'application
$$
\text{Aut}(\overline{K}/K) \times F_{\overline{x}}(Y) \to F_{\overline{x}}(Y),
\quad
(\sigma, \overline{y}) \mapsto
\sigma \cdot \overline{y} = \overline{y} \circ \Spec(\sigma)
$$
Il s'agit d'une action, car
$$
\sigma\tau \cdot \overline{y} =
\overline{y} \circ \Spec(\sigma\tau) =
\overline{y} \circ \Spec(\tau) \circ \Spec(\sigma) =
\sigma \cdot (\tau \cdot \overline{y})
$$
L'action est fonctorielle en $Y \in \textit{F\'Et}_X$, et l'on
obtient ainsi l'application voulue.
\end{remark}





\section{Revêtements galoisiens des schémas connexes}
\label{pione-section-finite-etale-under-galois}

\noindent
Soit $X$ un schéma connexe muni d'un point géométrique $\overline{x}$.
Puisque $F_{\overline{x}} : \textit{F\'Et}_X \to \textit{Ensembles}$ est le foncteur fibre d'une
catégorie galoisienne (lemme \ref{pione-lemma-finite-etale-connected-galois-category}),
les résultats de la section \ref{pione-section-galois} s'appliquent.
Dans cette section, nous transposons explicitement une partie de la terminologie
et des résultats au cadre des schémas et des morphismes finis étales.

\medskip\noindent
Nous dirons qu'un morphisme fini étale $Y \to X$ est un
{\it revêtement galoisien} si $Y$ définit un objet galoisien de
$\textit{F\'Et}_X$.
Pour un morphisme fini étale $Y \to X$, avec $G = \text{Aut}_X(Y)$,
les assertions suivantes sont équivalentes :
\begin{enumerate}
\item $Y$ est un revêtement galoisien de $X$,
\item $Y$ est connexe et $|G|$ est égal au degré de $Y \to X$,
\item $Y$ est connexe et $G$ agit transitivement sur $F_{\overline{x}}(Y)$, et
\item $Y$ est connexe et $G$ agit simplement transitivement sur
$F_{\overline{x}}(Y)$.
\end{enumerate}
Ceci résulte immédiatement de la discussion de la section \ref{pione-section-galois}.

\medskip\noindent
Pour tout morphisme fini étale $f : Y \to X$ avec $Y$ connexe,
il existe un revêtement galoisien fini étale $Y' \to X$ qui domine $Y$
(lemme \ref{pione-lemma-galois}).

\medskip\noindent
Les objets galoisiens de $\textit{F\'Et}_X$ correspondent, par l'équivalence
$$
F_{\overline{x}} : \textit{F\'Et}_X \to
\textit{Finite-}\pi_1(X, \overline{x})\textit{-Ensembles}
$$
du théorème \ref{pione-theorem-fundamental-group},
aux $\pi_1(X, \overline{x})\textit{-Ensembles}$ finis
de la forme $G = \pi_1(X, \overline{x})/H$, où $H$ est un
sous-groupe ouvert distingué. De manière équivalente, si $G$ est un groupe fini
et si $\pi_1(X, \overline{x}) \to G$ est une surjection continue,
alors $G$, considéré comme un $\pi_1(X, \overline{x})$-ensemble, correspond
à un revêtement galoisien.

\medskip\noindent
Si $Y_i \to X$, $i = 1, 2$, sont des revêtements galoisiens finis étales
de groupes de Galois $G_i$, alors il existe un revêtement galoisien
fini étale $Y \to X$ dont le groupe de Galois est un sous-groupe de
$G_1 \times G_2$. En effet, considérons les homomorphismes continus
correspondants $\pi_1(X, \overline{x}) \to G_i$ et
prenons pour $G$ l'image de l'homomorphisme continu induit
$\pi_1(X, \overline{x}) \to G_1 \times G_2$.









\section{Invariance topologique du groupe fondamental}
\label{pione-section-topological-invariance}

\noindent
Le résultat principal de cette section est qu'un homéomorphisme universel
de schémas connexes induit un isomorphisme sur les groupes fondamentaux.
Voir la proposition \ref{pione-proposition-universal-homeomorphism}.

\medskip\noindent
Au lieu de démontrer directement que deux schémas ont le même groupe
fondamental, nous démontrons souvent que leurs catégories de revêtements
finis étales sont les mêmes. Cela implique bien sûr que
leurs groupes fondamentaux sont égaux, pourvu qu'ils soient connexes.

\begin{lemma}
\label{pione-lemma-what-equivalence-gives}
Soit $f : X \to Y$ un morphisme de schémas quasi-compacts et quasi-séparés
tel que le foncteur de changement de base $\textit{F\'Et}_Y \to \textit{F\'Et}_X$
soit une équivalence de catégories. Alors
\begin{enumerate}
\item $f$ induit un homéomorphisme $\pi_0(X) \to \pi_0(Y)$,
\item si $X$, ou de manière équivalente $Y$, est connexe, alors
$\pi_1(X, \overline{x}) = \pi_1(Y, \overline{y})$.
\end{enumerate}
\end{lemma}

\begin{proof}
Soit $Y = Y_0 \amalg Y_1$ une décomposition en sous-schémas ouverts et fermés
non vides. Nous affirmons que $f(X)$ rencontre chacun des $Y_i$. En effet, sinon,
disons que $f(X) \subset Y_1$, nous pouvons considérer le morphisme fini étale
$V = Y_1 \to Y$. Celui-ci n'est pas un isomorphisme,
mais $V \times_Y X \to X$ en est un, ce qui
est une contradiction.

\medskip\noindent
Supposons que $X = X_0 \amalg X_1$ soit une décomposition en sous-schémas
ouverts et fermés. Considérons le morphisme fini étale $U = X_1 \to X$. Alors
$U = X \times_Y V$ pour un certain morphisme fini étale $V \to Y$. Le degré
du morphisme $V \to Y$ est localement constant ; on obtient donc une décomposition
$Y = \coprod_{d \geq 0} Y_d$ en sous-schémas ouverts et fermés
telle que $V \to Y$ soit de degré $d$ au-dessus de $Y_d$. Puisque
$f^{-1}(Y_d) = \emptyset$ pour $d > 1$, on conclut comme ci-dessus que $Y_d = \emptyset$
pour $d > 1$. Et l'on conclut que $f^{-1}(Y_i) = X_i$
pour $i = 0, 1$.

\medskip\noindent
Il s'ensuit que $f^{-1}$ induit une bijection entre l'ensemble des parties
ouvertes et fermées de $Y$ et celui des parties ouvertes et fermées de $X$.
Notons que $X$ et $Y$ sont des espaces spectraux ; voir Propriétés, lemme
\ref{properties-lemma-quasi-compact-quasi-separated-spectral}.
D'après Topologie, lemme \ref{topology-lemma-connected-component-intersection},
le treillis des parties ouvertes et fermées d'un espace spectral
détermine l'ensemble de ses composantes connexes.
Ainsi $\pi_0(X) \to \pi_0(Y)$ est bijective. Puisque $\pi_0(X)$ et
$\pi_0(Y)$ sont des espaces profinis
(Topologie, lemme \ref{topology-lemma-pi0-profinite}),
on conclut que $\pi_0(X) \to \pi_0(Y)$ est un homéomorphisme d'après
Topologie, lemme \ref{topology-lemma-bijective-map}. Ceci démontre (1).
L'assertion (2) est immédiate.
\end{proof}

\noindent
Le lemme suivant nous dit que le groupe fondamental d'un couple hensélien
est le groupe fondamental de la partie fermée.

\begin{lemma}
\label{pione-lemma-gabber}
Soit $(A, I)$ un couple hensélien. Posons $X = \Spec(A)$ et $Z = \Spec(A/I)$.
Le foncteur
$$
\textit{F\'Et}_X \longrightarrow \textit{F\'Et}_Z,\quad
U \longmapsto U \times_X Z
$$
est une équivalence de catégories.
\end{lemma}

\begin{proof}
C'est la traduction géométrique du
lemme \ref{more-algebra-lemma-finite-etale-equivalence} de Compléments d'algèbre.
\end{proof}

\noindent
Le lemme suivant nous dit que le groupe fondamental d'un épaississement
est le même que celui du schéma de départ. Nous l'utiliserons
dans la preuve de la proposition forte sur les homéomorphismes universels
ci-dessous.

\begin{lemma}
\label{pione-lemma-thickening}
Soit $X \subset X'$ un épaississement de schémas. Le foncteur
$$
\textit{F\'Et}_{X'} \longrightarrow \textit{F\'Et}_X,\quad
U' \longmapsto U' \times_{X'} X
$$
est une équivalence de catégories.
\end{lemma}

\begin{proof}
Pour une discussion des épaississements, voir
Compléments sur les morphismes, section \ref{more-morphisms-section-thickenings}.
Soit $U' \to X'$ un morphisme étale tel que $U = U' \times_{X'} X \to X$
soit fini étale. Alors $U' \to X'$ est lui aussi fini étale.
Cela résulte par exemple de Compléments sur les morphismes, lemme
\ref{more-morphisms-lemma-properties-that-extend-over-thickenings}.
Maintenant, si $X \subset X'$ est un épaississement d'ordre fini, cette remarque,
combinée avec Morphismes étales, théorème
\ref{etale-theorem-remarkable-equivalence},
démontre le lemme. Nous allons démontrer ci-dessous le lemme pour les épaississements
généraux, mais nous suggérons au lecteur de sauter la preuve.

\medskip\noindent
Soit $X' = \bigcup X_i'$ un recouvrement ouvert affine. Posons
$X_i = X \times_{X'} X_i'$, $X_{ij}' = X'_i \cap X'_j$,
$X_{ij} = X \times_{X'} X_{ij}'$, $X_{ijk}' = X'_i \cap X'_j \cap X'_k$,
$X_{ijk} = X \times_{X'} X_{ijk}'$.
Supposons que nous puissions démontrer
le résultat pour chacun des épaississements
$X_i \subset X'_i$, $X_{ij} \subset X_{ij}'$ et $X_{ijk} \subset X_{ijk}'$.
Alors le résultat pour $X \subset X'$ découle du recollement relatif des
schémas ; voir
Constructions, section \ref{constructions-section-relative-glueing}.
Observons que les schémas $X_i'$, $X_{ij}'$, $X_{ijk}'$ sont
tous séparés, comme sous-schémas ouverts de schémas affines. En répétant
l'argument une fois de plus, on se ramène au cas où les schémas
$X'_i$, $X_{ij}'$, $X_{ijk}'$ sont affines.

\medskip\noindent
Dans le cas affine, on a $X' = \Spec(A')$ et $X = \Spec(A'/I')$,
où $I'$ est un idéal localement nilpotent. Alors $(A', I')$ est un
couple hensélien (Compléments d'algèbre, lemme
\ref{more-algebra-lemma-locally-nilpotent-henselian}),
et le résultat découle du lemme \ref{pione-lemma-gabber} (qui est
beaucoup plus facile dans ce cas).
\end{proof}

\noindent
La ``bonne'' manière de démontrer la proposition suivante serait de
la déduire de l'invariance du site étale ; voir
Cohomologie étale, théorème
\ref{etale-cohomology-theorem-topological-invariance}.

\begin{proposition}
\label{pione-proposition-universal-homeomorphism}
Soit $f : X \to Y$ un homéomorphisme universel de schémas. Alors
$$
\textit{F\'Et}_Y \longrightarrow \textit{F\'Et}_X,\quad
V \longmapsto V \times_Y X
$$
est une équivalence. Ainsi, si $X$ et $Y$ sont connexes, alors
$f$ induit un isomorphisme $\pi_1(X, \overline{x}) \to \pi_1(Y, \overline{y})$
de groupes fondamentaux.
\end{proposition}

\begin{proof}
Rappelons qu'un homéomorphisme universel est la même chose qu'un morphisme
entier, universellement injectif et surjectif ; voir
Morphismes, lemme \ref{morphisms-lemma-universal-homeomorphism}.
En particulier, la diagonale $\Delta : X \to X \times_Y X$ est un épaississement
d'après Morphismes, lemme \ref{morphisms-lemma-universally-injective}.
Ainsi, d'après le lemme \ref{pione-lemma-thickening},
étant donné un morphisme fini étale $U \to X$,
il existe un unique isomorphisme
$$
\varphi : U \times_Y X \to X \times_Y U
$$
de schémas finis étales sur $X \times_Y X$ dont l'image réciproque par
$\Delta$ est $\text{id} : U \to U$ sur $X$.
Puisque $X \to X \times_Y X \times_Y X$
est également un épaississement (c'est une immersion fermée bijective),
on conclut que $(U, \varphi)$ est une donnée de descente relativement à $X/Y$.
D'après Morphismes étales, proposition \ref{etale-proposition-effective},
on conclut que $U = X \times_Y V$ pour un certain $V \to Y$
quasi-compact, séparé et étale.
Nous omettons la preuve que $V \to Y$ est fini (indications :
le morphisme $U \to V$ est surjectif et $U \to Y$ est entier).
On conclut que $\textit{F\'Et}_Y \to \textit{F\'Et}_X$
est essentiellement surjectif.

\medskip\noindent
En raisonnant de la même manière que ci-dessus, on voit qu'étant donnés
$V_1 \to Y$ et $V_2 \to Y$ dans $\textit{F\'Et}_Y$, tout
morphisme $a : X \times_Y V_1 \to X \times_Y V_2$ sur $X$
est compatible avec les données de descente canoniques. Ainsi $a$
provient d'un morphisme $V_1 \to V_2$ sur $Y$ d'après
Morphismes étales, lemme \ref{etale-lemma-fully-faithful-cases}.
\end{proof}










\section{Revêtements finis étales des schémas propres}
\label{pione-section-finite-etale-over-proper}

\noindent
Dans cette section, nous montrons que le groupe fondamental d'un schéma propre
connexe sur un anneau local hensélien est le même que le groupe fondamental
de sa fibre spéciale. Nous démontrons aussi une variante de ce résultat pour
un couple hensélien.

\medskip\noindent
Nous montrons également que le groupe
fondamental d'un schéma propre connexe sur un corps algébriquement clos $k$
ne change pas si l'on remplace $k$ par une extension algébriquement close.

\medskip\noindent
Au lieu d'énoncer et de démontrer les résultats dans le cas connexe,
nous les démontrons en général et laissons au lecteur le soin d'en déduire
le résultat pour les groupes fondamentaux à l'aide du
lemme \ref{pione-lemma-what-equivalence-gives}.

\begin{lemma}
\label{pione-lemma-finite-etale-on-proper-over-henselian}
Soit $A$ un anneau local hensélien. Soit $X$ un schéma propre sur $A$,
de fibre fermée $X_0$. Alors le foncteur
$$
\textit{F\'Et}_X \to \textit{F\'Et}_{X_0},\quad
U \longmapsto U_0 = U \times_X X_0
$$
est une équivalence de catégories.
\end{lemma}

\begin{proof}
La preuve donnée ici est un exemple d'application de l'algébrisation et de
l'approximation. Nous procédons en plusieurs étapes.

\medskip\noindent
Surjectivité essentielle lorsque $A$ est un anneau local noethérien complet.
Soit $X_n = X \times_{\Spec(A)} \Spec(A/\mathfrak m^{n + 1})$.
D'après Morphismes étales, théorème \ref{etale-theorem-remarkable-equivalence},
les inclusions
$$
X_0 \to X_1 \to X_2 \to \ldots
$$
induisent des équivalences de catégories entre la catégorie
des schémas étales sur $X_0$ et la catégorie des schémas
étales sur $X_n$.
De plus, si $U_n \to X_n$ correspond à un morphisme fini étale
$U_0 \to X_0$, alors $U_n \to X_n$ est également fini, par exemple
d'après Compléments sur les morphismes, lemme
\ref{more-morphisms-lemma-thicken-property-morphisms-cartesian}.
Dans ce cas, le morphisme $U_0 \to \Spec(A/\mathfrak m)$
est propre puisque $X_0$ est propre sur $A/\mathfrak m$. Nous pouvons donc appliquer
le théorème d'algébrisation de Grothendieck
(sous la forme du
lemme de Cohomologie des schémas
\ref{coherent-lemma-algebraize-formal-scheme-finite-over-proper})
pour voir qu'il existe un morphisme fini $U \to X$ dont la restriction
à $X_0$ redonne $U_0$. D'après Compléments sur les morphismes, lemme
\ref{more-morphisms-lemma-check-smoothness-on-infinitesimal-nbhds},
on voit que $U \to X$ est étale en tout point de $U_0$.
Cependant, puisque tout point de $U$ se spécialise en un point de $U_0$
(car $U$ est propre sur $A$), on conclut que $U \to X$ est étale.
On conclut ainsi que le foncteur est essentiellement surjectif.

\medskip\noindent
Pleine fidélité lorsque $A$ est un anneau local noethérien complet.
Soient $U \to X$ et $V \to X$ des morphismes finis étales, et
soit $\varphi_0 : U_0 \to V_0$ un morphisme sur $X_0$. Considérons
le morphisme
$$
\Gamma_{\varphi_0} : U_0 \longrightarrow U_0 \times_{X_0} V_0
$$
Ce morphisme est à la fois fini étale et une immersion fermée.
En appliquant la surjectivité essentielle au schéma $U \times_X V$, on trouve
un morphisme fini étale $W \to U \times_X V$ dont la fibre spéciale
est isomorphe à $\Gamma_{\varphi_0}$. Considérons la projection
$W \to U$. Elle est finie étale et est un isomorphisme au-dessus de $U_0$ par
construction. D'après Morphismes étales, lemme
\ref{etale-lemma-finite-etale-one-point},
$W \to U$ est un isomorphisme dans un voisinage ouvert de $U_0$.
C'est donc un isomorphisme, et la composée $\varphi : U \cong W \to V$
est le relèvement voulu de $\varphi_0$.

\medskip\noindent
Surjectivité essentielle lorsque $A$ est un G-anneau local noethérien hensélien.
Soit $U_0 \to X_0$ un morphisme fini étale.
Soit $A^\wedge$ le complété de $A$ pour la topologie définie par l'idéal maximal.
Soit $X^\wedge$ le changement de base de $X$ à $A^\wedge$.
D'après le résultat précédent, il existe un morphisme fini étale
$V \to X^\wedge$ dont la fibre spéciale est $U_0$.
Écrivons $A^\wedge = \colim A_i$, où $A \to A_i$ est de type fini.
D'après Limites, lemme \ref{limits-lemma-descend-finite-presentation},
il existe un indice $i$ et un morphisme de présentation finie $U_i \to X_{A_i}$
dont le changement de base à $X^\wedge$ est $V$. En augmentant $i$,
on peut supposer que $U_i \to X_{A_i}$ est fini et étale
(Limites, lemmes \ref{limits-lemma-descend-finite-finite-presentation} et
\ref{limits-lemma-descend-etale}). En écrivant
$$
A_i = A[x_1, \ldots, x_n]/(f_1, \ldots, f_m)
$$
l'homomorphisme d'anneaux $A_i \to A^\wedge$ peut se réinterpréter comme une solution
$(a_1, \ldots, a_n)$ dans $A^\wedge$ du système d'équations $f_j = 0$.
D'après Lissage des morphismes d'anneaux, théorème \ref{smoothing-theorem-approximation-property},
on peut approcher cette solution (à l'ordre $11$, par exemple) par une solution
$(b_1, \ldots, b_n)$ dans $A$. En retraduisant, on trouve un homomorphisme de $A$-algèbres
$A_i \to A$ qui donne le même point fermé que l'homomorphisme initial
$A_i \to A^\wedge$ (car $11 > 1$). Le changement de base $U \to X$ de $U_i \to X_{A_i}$
par cet homomorphisme d'anneaux sera donc un morphisme fini étale dont la
fibre spéciale est isomorphe à $U_0$.

\medskip\noindent
Pleine fidélité lorsque $A$ est un G-anneau local noethérien hensélien.
Elle se déduit de la surjectivité essentielle exactement de la même
manière que dans le cas où $A$ est noethérien complet.

\medskip\noindent
Cas général. Soit $(A, \mathfrak m)$ un anneau local hensélien.
Posons $S = \Spec(A)$ et notons $s \in S$ le point fermé. D'après Limites, lemme
\ref{limits-lemma-proper-limit-of-proper-finite-presentation-noetherian},
on peut écrire $X \to \Spec(A)$ comme une limite cofiltrante de
morphismes propres $X_i \to S_i$, où $S_i$ est de type fini sur $\mathbf{Z}$.
Pour chaque $i$, soit $s_i \in S_i$ l'image de $s$.
Puisque $S = \lim S_i$ et $A = \mathcal{O}_{S, s}$, on a
$A = \colim \mathcal{O}_{S_i, s_i}$. L'anneau $A_i = \mathcal{O}_{S_i, s_i}$
est un G-anneau local noethérien (Compléments d'algèbre, proposition
\ref{more-algebra-proposition-ubiquity-G-ring}).
D'après Compléments d'algèbre, lemme \ref{more-algebra-lemma-henselization-colimit},
on voit que $A = \colim A_i^h$. D'après
Compléments d'algèbre, lemme \ref{more-algebra-lemma-henselization-G-ring},
les anneaux $A_i^h$ sont des G-anneaux. Ainsi $A = \colim A_i^h$ et
$$
X = \lim (X_i \times_{S_i} \Spec(A_i^h))
$$
comme schémas. La catégorie des schémas finis étales sur $X$ est la limite
des catégories de schémas finis étales sur
$X_i \times_{S_i} \Spec(A_i^h)$ (d'après
Limites, lemmes
\ref{limits-lemma-descend-finite-presentation},
\ref{limits-lemma-descend-finite-finite-presentation} et
\ref{limits-lemma-descend-etale}).
Il en va de même pour les schémas finis étales sur
$X_0 = \lim (X_i \times_{S_i} s_i)$.
Ainsi, on déduit formellement le résultat pour $X / \Spec(A)$
du résultat pour les $(X_i \times_{S_i} \Spec(A_i^h)) / \Spec(A_i^h)$
traité ci-dessus.
\end{proof}

\begin{lemma}
\label{pione-lemma-finite-etale-on-proper-over-henselian-pair}
Soit $(A, I)$ un couple hensélien. Soit $X$ un schéma propre sur $A$.
Posons $X_0 = X \times_{\Spec(A)} \Spec(A/I)$. Alors le foncteur
$$
\textit{F\'Et}_X \to \textit{F\'Et}_{X_0},\quad
U \longmapsto U_0 = U \times_X X_0
$$
est une équivalence de catégories.
\end{lemma}

\begin{proof}
La preuve de ce lemme est exactement la même que celle du
lemme \ref{pione-lemma-finite-etale-on-proper-over-henselian}.

\medskip\noindent
Surjectivité essentielle lorsque $A$ est noethérien et complet pour la topologie $I$-adique.
Soit $X_n = X \times_{\Spec(A)} \Spec(A/I^{n + 1})$.
D'après Morphismes étales, théorème \ref{etale-theorem-remarkable-equivalence},
les inclusions
$$
X_0 \to X_1 \to X_2 \to \ldots
$$
induisent des équivalences de catégories entre la catégorie
des schémas étales sur $X_0$ et la catégorie des schémas
étales sur $X_n$.
De plus, si $U_n \to X_n$ correspond à un morphisme fini étale
$U_0 \to X_0$, alors $U_n \to X_n$ est également fini, par exemple
d'après Compléments sur les morphismes, lemme
\ref{more-morphisms-lemma-thicken-property-morphisms-cartesian}.
Dans ce cas, le morphisme $U_0 \to \Spec(A/I)$
est propre puisque $X_0$ est propre sur $A/I$. Nous pouvons donc appliquer
le théorème d'algébrisation de Grothendieck
(sous la forme du
lemme de Cohomologie des schémas
\ref{coherent-lemma-algebraize-formal-scheme-finite-over-proper})
pour voir qu'il existe un morphisme fini $U \to X$ dont la restriction
à $X_0$ redonne $U_0$. D'après Compléments sur les morphismes, lemme
\ref{more-morphisms-lemma-check-smoothness-on-infinitesimal-nbhds},
on voit que $U \to X$ est étale en tout point de $U_0$.
Cependant, puisque tout point de $U$ se spécialise en un point de $U_0$
(car $U$ est propre sur $A$), on conclut que $U \to X$ est étale.
On conclut ainsi que le foncteur est essentiellement surjectif.

\medskip\noindent
Pleine fidélité lorsque $A$ est noethérien et complet pour la topologie $I$-adique.
Soient $U \to X$ et $V \to X$ des morphismes finis étales, et
soit $\varphi_0 : U_0 \to V_0$ un morphisme sur $X_0$. Considérons
le morphisme
$$
\Gamma_{\varphi_0} : U_0 \longrightarrow U_0 \times_{X_0} V_0
$$
Ce morphisme est à la fois fini étale et une immersion fermée.
En appliquant la surjectivité essentielle au schéma $U \times_X V$, on trouve
un morphisme fini étale $W \to U \times_X V$ dont la fibre spéciale
est isomorphe à $\Gamma_{\varphi_0}$. Considérons la projection
$W \to U$. Elle est finie étale et est un isomorphisme au-dessus de $U_0$ par
construction. D'après Morphismes étales, lemme
\ref{etale-lemma-finite-etale-one-point},
$W \to U$ est un isomorphisme dans un voisinage ouvert de $U_0$.
C'est donc un isomorphisme, et la composée $\varphi : U \cong W \to V$
est le relèvement voulu de $\varphi_0$.

\medskip\noindent
Surjectivité essentielle lorsque $(A, I)$ est un couple hensélien
et que $A$ est un G-anneau noethérien.
Soit $U_0 \to X_0$ un morphisme fini étale.
Soit $A^\wedge$ le complété de $A$ pour la topologie $I$-adique.
Observons que $A^\wedge$ est un anneau noethérien qui est
complet pour la topologie $IA^\wedge$-adique ; voir Algèbre, lemmes
\ref{algebra-lemma-completion-complete} et
\ref{algebra-lemma-completion-Noetherian-Noetherian}.
Soit $X^\wedge$ le changement de base de $X$ à $A^\wedge$.
D'après le résultat précédent, il existe un morphisme fini étale
$V \to X^\wedge$ dont la fibre spéciale est $U_0$.
Écrivons $A^\wedge = \colim A_i$, où $A \to A_i$ est de type fini.
D'après Limites, lemme \ref{limits-lemma-descend-finite-presentation},
il existe un indice $i$ et un morphisme de présentation finie $U_i \to X_{A_i}$
dont le changement de base à $X^\wedge$ est $V$. En augmentant $i$,
on peut supposer que $U_i \to X_{A_i}$ est fini et étale
(Limites, lemmes \ref{limits-lemma-descend-finite-finite-presentation} et
\ref{limits-lemma-descend-etale}). En écrivant
$$
A_i = A[x_1, \ldots, x_n]/(f_1, \ldots, f_m)
$$
l'homomorphisme d'anneaux $A_i \to A^\wedge$ peut se réinterpréter comme une solution
$(a_1, \ldots, a_n)$ dans $A^\wedge$ du système d'équations $f_j = 0$.
D'après Lissage des morphismes d'anneaux, lemme \ref{smoothing-lemma-henselian-pair},
on peut approcher cette solution (à l'ordre $11$, par exemple) par une solution
$(b_1, \ldots, b_n)$ dans $A$. En retraduisant, on trouve un homomorphisme de $A$-algèbres
$A_i \to A$ qui donne le même point fermé que l'homomorphisme initial
$A_i \to A^\wedge$ (car $11 > 1$). Le changement de base $U \to X$ de $U_i \to X_{A_i}$
par cet homomorphisme d'anneaux sera donc un morphisme fini étale dont la
fibre spéciale est isomorphe à $U_0$.

\medskip\noindent
Pleine fidélité lorsque $(A, I)$ est un couple hensélien et que $A$ est un
G-anneau noethérien. Elle se déduit de la surjectivité essentielle exactement de la même
manière que dans le cas où $A$ est noethérien complet.

\medskip\noindent
Cas général. Soit $(A, I)$ un couple hensélien.
Posons $S = \Spec(A)$ et notons $S_0 = \Spec(A/I)$. D'après Limites, lemme
\ref{limits-lemma-proper-limit-of-proper-finite-presentation-noetherian},
on peut écrire $X \to \Spec(A)$ comme une limite cofiltrante de
morphismes propres $X_i \to S_i$, où $S_i$ est affine et de type fini
sur $\mathbf{Z}$. Écrivons $S_i = \Spec(A_i)$ et notons
$I_i \subset A_i$ l'image réciproque de $I$ par l'application $A_i \to A$.
Posons $S_{i, 0} = \Spec(A_i/I_i)$. Puisque $S = \lim S_i$, on a
$A = \colim A_i$. Ainsi, on a aussi $I = \colim I_i$ et $A/I = \colim A_i/I_i$.
L'anneau $A_i$ est un G-anneau noethérien (Compléments d'algèbre, proposition
\ref{more-algebra-proposition-ubiquity-G-ring}).
Notons $(A_i^h, I_i^h)$ l'hensélisation de la paire $(A_i, I_i)$.
D'après Compléments d'algèbre, lemme \ref{more-algebra-lemma-henselization-colimit},
on voit que $A = \colim A_i^h$. D'après
Compléments d'algèbre, lemme \ref{more-algebra-lemma-henselization-pair-G-ring},
les anneaux $A_i^h$ sont des G-anneaux. Ainsi $A = \colim A_i^h$ et
$$
X = \lim (X_i \times_{S_i} \Spec(A_i^h))
$$
comme schémas. La catégorie des schémas finis étales sur $X$ est la limite
des catégories de schémas finis étales sur
$X_i \times_{S_i} \Spec(A_i^h)$ (d'après
Limites, lemmes
\ref{limits-lemma-descend-finite-presentation},
\ref{limits-lemma-descend-finite-finite-presentation} et
\ref{limits-lemma-descend-etale}).
Il en va de même pour les schémas finis étales sur
$X_0 = \lim (X_i \times_{S_i} S_{i, 0})$.
Ainsi, on déduit formellement le résultat pour $X / \Spec(A)$
du résultat pour les $(X_i \times_{S_i} \Spec(A_i^h)) / \Spec(A_i^h)$
traité ci-dessus.
\end{proof}

\begin{lemma}
\label{pione-lemma-finite-etale-invariant-over-proper}
Soit $k'/k$ une extension de corps algébriquement clos.
Soit $X$ un schéma propre sur $k$. Alors le foncteur
$$
U \longmapsto U_{k'}
$$
est une équivalence de catégories entre les schémas finis étales sur
$X$ et les schémas finis étales sur $X_{k'}$.
\end{lemma}

\begin{proof}
Démontrons que le foncteur est essentiellement surjectif.
Soit $U' \to X_{k'}$ un morphisme fini étale.
Écrivons $k' = \colim A_i$ comme colimite filtrante de $k$-algèbres de type fini.
D'après Limites, lemme \ref{limits-lemma-descend-finite-presentation},
il existe un indice $i$ et un morphisme de présentation finie $U_i \to X_{A_i}$
dont le changement de base à $X_{k'}$ est $U'$. En augmentant $i$,
on peut supposer que $U_i \to X_{A_i}$ est fini et étale
(Limites, lemmes \ref{limits-lemma-descend-finite-finite-presentation} et
\ref{limits-lemma-descend-etale}).
Puisque $k$ est algébriquement clos, on peut trouver un
point à valeurs dans $k$, disons $t$, de $\Spec(A_i)$. Soit $U = (U_i)_t$ la
fibre de $U_i$ au-dessus de $t$. Soit $A_i^h$
l'hensélisé de $(A_i)_{\mathfrak m}$, où $\mathfrak m$ est
l'idéal maximal correspondant au point $t$. D'après
le lemme \ref{pione-lemma-finite-etale-on-proper-over-henselian},
on voit que $(U_i)_{A_i^h} = U \times \Spec(A_i^h)$ comme schémas
sur $X_{A_i^h}$. Maintenant, puisque
$A_i^h$ est algébrique sur $A_i$ (voir par exemple la discussion de
Lissage des morphismes d'anneaux, exemple \ref{smoothing-example-describe-henselian})
et puisque $k'$ est algébriquement clos,
on peut trouver un homomorphisme d'anneaux $A_i^h \to k'$ qui prolonge
l'inclusion donnée $A_i \subset k'$. On conclut donc que $U'$
est isomorphe au changement de base de $U$.
La preuve de la pleine fidélité est exactement la même.
\end{proof}








\section{Connexité locale}
\label{pione-section-unibranch}

\noindent
Dans cette section, nous demandons quand $\pi_1(U) \to \pi_1(X)$ est surjectif
pour $U$ un ouvert dense d'un schéma $X$. Nous verrons que c'est
le cas (en gros) lorsque $U \cap B$ est connexe pour toute petite
``boule'' $B$ autour d'un point $x \in X \setminus U$.

\begin{lemma}
\label{pione-lemma-dense-faithful}
Soit $f : X \to Y$ un morphisme de schémas. Si $f(X)$ est dense dans $Y$,
alors le foncteur de changement de base $\textit{F\'Et}_Y \to \textit{F\'Et}_X$
est fidèle.
\end{lemma}

\begin{proof}
Puisque la catégorie des revêtements finis étales possède un
Hom interne (lemme \ref{pione-lemma-internal-hom-finite-etale}),
il suffit de démontrer ce qui suit : étant donné $W$ fini étale sur $Y$
et un morphisme $s : X \to W$ sur $Y$, il existe au plus une section
$t : Y \to W$ telle que $s = t \circ f$. Considérons deux sections
$t_1, t_2 : Y \to W$ telles que $s = t_1 \circ f = t_2 \circ f$.
Puisque l'égalisateur de $t_1$ et $t_2$ est fermé dans $Y$
(Schémas, lemme \ref{schemes-lemma-where-are-they-equal})
et puisque $f(X)$ est dense dans $Y$, on voit que $t_1$ et $t_2$
coïncident sur $Y_{red}$. Il s'ensuit que $t_1$ et $t_2$ ont
la même image, qui est un sous-schéma ouvert et fermé de $W$ s'envoyant
isomorphiquement sur $Y$
(Morphismes étales, proposition \ref{etale-proposition-properties-sections}) ;
ils sont donc égaux.
\end{proof}

\noindent
La condition du lemme suivant, selon laquelle le spectre épointé
du hensélisé strict est connexe, résulte par exemple de
l'hypothèse que l'anneau local est géométriquement unibranche ; voir
Compléments d'algèbre, lemme \ref{more-algebra-lemma-geometrically-unibranch}.
Il existe une réciproque partielle dans
Propriétés, lemme \ref{properties-lemma-geometrically-unibranch}.

\begin{lemma}
\label{pione-lemma-same-etale-extensions}
Soit $(A, \mathfrak m)$ un anneau local. Posons $X = \Spec(A)$
et soit $U = X \setminus \{\mathfrak m\}$. Si le spectre épointé
du hensélisé strict de $A$ est connexe, alors
$$
\textit{F\'Et}_X \longrightarrow \textit{F\'Et}_U,\quad
Y \longmapsto Y \times_X U
$$
est un foncteur pleinement fidèle.
\end{lemma}

\begin{proof}
Supposons $A$ strictement hensélien. Dans ce cas, tout revêtement fini étale
$Y$ de $X$ est isomorphe à une réunion disjointe finie de
copies de $X$. Il suffit donc de démontrer que tout morphisme
$U \to U \amalg \ldots \amalg U$ sur $U$ se prolonge de manière unique en un morphisme
$X \to X \amalg \ldots \amalg X$ sur $X$.
Si $U$ est connexe (donc en particulier non vide), cela est vrai.

\medskip\noindent
Cas général. Puisque la catégorie des revêtements finis étales possède un
Hom interne (lemme \ref{pione-lemma-internal-hom-finite-etale}),
il suffit de démontrer ce qui suit : étant donné $Y$ fini étale sur $X$,
tout morphisme $s : U \to Y$ sur $X$ se prolonge en un morphisme $t : X \to Y$
sur $X$. Soit $A^{sh}$ l'hensélisé strict de $A$ et notons
$X^{sh} = \Spec(A^{sh})$, $U^{sh} = U \times_X X^{sh}$,
$Y^{sh} = Y \times_X X^{sh}$. D'après le premier paragraphe et notre hypothèse
sur $A$, on peut prolonger le changement de base $s^{sh} : U^{sh} \to Y^{sh}$ de $s$ en
$t^{sh} : X^{sh} \to Y^{sh}$. Posons $A' = A^{sh} \otimes_A A^{sh}$.
Alors les deux images réciproques $t'_1, t'_2$ de $t^{sh}$ sur $X' = \Spec(A')$
sont des prolongements de l'image réciproque $s'$ de $s$ sur $U' = U \times_X X'$.
Comme $A \to A'$ est plat, on voit que $U' \subset X'$ est dense (topologiquement)
par descente des idéaux premiers pour $A \to A'$
(Algèbre, lemme \ref{algebra-lemma-flat-going-down}). Ainsi,
$t'_1 = t'_2$ d'après le lemme \ref{pione-lemma-dense-faithful}.
Par conséquent, $t^{sh}$ provient d'un morphisme $t : X \to Y$,
par exemple d'après
Descente, lemme \ref{descent-lemma-fpqc-universal-effective-epimorphisms}.
\end{proof}

\noindent
Au vu du lemme \ref{pione-lemma-same-etale-extensions},
il est intéressant de savoir quand le
spectre épointé d'un anneau (et celui de son hensélisé strict)
est connexe. Il existe un célèbre lemme dû à Hartshorne
qui donne une condition suffisante ; voir
Cohomologie locale, lemme
\ref{local-cohomology-lemma-depth-2-connected-punctured-spectrum}.

\begin{lemma}
\label{pione-lemma-quasi-compact-dense-open-connected-at-infinity-Noetherian}
Soit $X$ un schéma. Soit $U \subset X$ un ouvert dense. Supposons que
\begin{enumerate}
\item l'espace topologique sous-jacent à $X$ soit noethérien, et
\item pour tout $x \in X \setminus U$, le spectre épointé du
hensélisé strict de $\mathcal{O}_{X, x}$ soit connexe.
\end{enumerate}
Alors $\textit{F\'Et}_X \to \textit{F\'Et}_U$ est pleinement fidèle.
\end{lemma}

\begin{proof}
Soient $Y_1, Y_2$ finis étales sur $X$ et soit
$\varphi : (Y_1)_U \to (Y_2)_U$ un morphisme sur $U$. Il faut montrer que
$\varphi$ se prolonge de manière unique en un morphisme $Y_1 \to Y_2$ sur $X$.
L'unicité découle du lemme \ref{pione-lemma-dense-faithful}.

\medskip\noindent
Soit $x \in X \setminus U$ un point générique d'une composante irréductible
de $X \setminus U$. Posons $V = U \times_X \Spec(\mathcal{O}_{X, x})$.
Par notre choix de $x$, c'est le spectre épointé de
$\Spec(\mathcal{O}_{X, x})$. D'après le
lemme \ref{pione-lemma-same-etale-extensions},
on peut prolonger de manière unique le morphisme $\varphi_V : (Y_1)_V \to (Y_2)_V$
en un morphisme
$(Y_1)_{\Spec(\mathcal{O}_{X, x})} \to (Y_2)_{\Spec(\mathcal{O}_{X, x})}$.
D'après Limites, lemme \ref{limits-lemma-glueing-near-point},
on trouve un ouvert $U \subset U'$ contenant $x$ et un prolongement
$\varphi' : (Y_1)_{U'} \to (Y_2)_{U'}$ de $\varphi$.
Puisque l'espace topologique sous-jacent à $X$ est noethérien,
ceci achève la preuve par récurrence noethérienne sur le complémentaire
de l'ouvert sur lequel $\varphi$ est défini.
\end{proof}

\begin{lemma}
\label{pione-lemma-retrocompact-dense-open-connected-at-infinity-closed}
Soit $X$ un schéma. Soit $U \subset X$ un ouvert dense. Supposons que
\begin{enumerate}
\item $U \to X$ soit quasi-compact,
\item tout point de $X \setminus U$ soit fermé, et
\item pour tout $x \in X \setminus U$, le spectre épointé du
hensélisé strict de $\mathcal{O}_{X, x}$ soit connexe.
\end{enumerate}
Alors $\textit{F\'Et}_X \to \textit{F\'Et}_U$ est pleinement fidèle.
\end{lemma}

\begin{proof}
Soient $Y_1, Y_2$ finis étales sur $X$ et soit
$\varphi : (Y_1)_U \to (Y_2)_U$ un morphisme sur $U$. Il faut montrer que
$\varphi$ se prolonge de manière unique en un morphisme $Y_1 \to Y_2$ sur $X$.
L'unicité découle du lemme \ref{pione-lemma-dense-faithful}.

\medskip\noindent
Soit $x \in X \setminus U$. Posons $V = U \times_X \Spec(\mathcal{O}_{X, x})$.
Puisque tout point de $X \setminus U$ est fermé, $V$ est le spectre épointé
de $\Spec(\mathcal{O}_{X, x})$. D'après le
lemme \ref{pione-lemma-same-etale-extensions},
on peut prolonger de manière unique le morphisme $\varphi_V : (Y_1)_V \to (Y_2)_V$
en un morphisme
$(Y_1)_{\Spec(\mathcal{O}_{X, x})} \to (Y_2)_{\Spec(\mathcal{O}_{X, x})}$.
D'après Limites, lemme \ref{limits-lemma-glueing-near-point},
(ceci utilise que $U$ est rétrocompact dans $X$),
on trouve un ouvert $U \subset U'_x$ contenant $x$ et un prolongement
$\varphi'_x : (Y_1)_{U'_x} \to (Y_2)_{U'_x}$ de $\varphi$.
Notons qu'étant donnés deux points $x, x' \in X \setminus U$, les
morphismes $\varphi'_x$ et $\varphi'_{x'}$ coïncident sur
$U'_x \cap U'_{x'}$, car $U$ est dense dans cet ouvert
(lemme \ref{pione-lemma-dense-faithful}). Ainsi, on peut prolonger $\varphi$
à $\bigcup U'_x = X$, comme souhaité.
\end{proof}

\begin{lemma}
\label{pione-lemma-quasi-compact-dense-open-connected-at-infinity}
Soit $X$ un schéma. Soit $U \subset X$ un ouvert dense. Supposons que
\begin{enumerate}
\item tout ouvert quasi-compact de $X$ possède un nombre fini de
composantes irréductibles,
\item pour tout $x \in X \setminus U$, le spectre épointé du
hensélisé strict de $\mathcal{O}_{X, x}$ soit connexe.
\end{enumerate}
Alors $\textit{F\'Et}_X \to \textit{F\'Et}_U$ est pleinement fidèle.
\end{lemma}

\begin{proof}
Soient $Y_1, Y_2$ finis étales sur $X$ et soit
$\varphi : (Y_1)_U \to (Y_2)_U$ un morphisme sur $U$. Il faut montrer que
$\varphi$ se prolonge de manière unique en un morphisme $Y_1 \to Y_2$ sur $X$.
L'unicité découle du lemme \ref{pione-lemma-dense-faithful}.
Nous allons démontrer l'existence en montrant que nous pouvons agrandir $U$
si $U \not = X$, puis conclure à l'aide du lemme de Zorn.

\medskip\noindent
Soit $x \in X \setminus U$ un point générique d'une composante irréductible
de $X \setminus U$. Posons $V = U \times_X \Spec(\mathcal{O}_{X, x})$.
Par notre choix de $x$, c'est le spectre épointé de
$\Spec(\mathcal{O}_{X, x})$. D'après le
lemme \ref{pione-lemma-same-etale-extensions},
on peut prolonger le morphisme $\varphi_V : (Y_1)_V \to (Y_2)_V$
(de manière unique) en un morphisme
$(Y_1)_{\Spec(\mathcal{O}_{X, x})} \to (Y_2)_{\Spec(\mathcal{O}_{X, x})}$.
Choisissons un voisinage affine $W \subset X$ de $x$.
Puisque $U \cap W$ est dense dans $W$, il contient les points génériques
$\eta_1, \ldots, \eta_n$ de $W$. Choisissons un ouvert affine
$W' \subset W \cap U$ contenant $\eta_1, \ldots, \eta_n$.
Posons $V' = W' \times_X \Spec(\mathcal{O}_{X, x})$.
D'après Limites, lemme \ref{limits-lemma-glueing-near-point},
appliqué à $x \in W \supset W'$,
on trouve un ouvert $W' \subset W'' \subset W$ avec $x \in W''$
et un morphisme $\varphi'' : (Y_1)_{W''} \to (Y_2)_{W''}$
coïncidant avec $\varphi$ sur $W'$. Puisque $W'$ est dense dans
$W'' \cap U$, le lemme \ref{pione-lemma-dense-faithful} montre
que $\varphi$ et $\varphi''$ coïncident sur $U \cap W''$.
Ainsi, $\varphi$ et $\varphi''$ se recollent en un morphisme
$\varphi'$ sur $U' = U \cup W''$ qui coïncide avec $\varphi$ sur $U$.
Observons que $x \in U'$, de sorte que nous avons prolongé $\varphi$
à un ouvert strictement plus grand.

\medskip\noindent
Considérons l'ensemble $\mathcal{S}$ des couples $(U', \varphi')$ où $U \subset U'$
et où $\varphi'$ est un prolongement de $\varphi$. Munissons $\mathcal{S}$
de l'ordre partiel évident. Si les couples $(U'_i, \varphi'_i)$
forment une chaîne, celle-ci admet une borne supérieure $(U', \varphi')$.
Il suffit de prendre $U' = \bigcup U'_i$ et de définir
$\varphi' : (Y_1)_{U'} \to (Y_2)_{U'}$ comme le morphisme
qui coïncide avec $\varphi'_i$ sur $U'_i$. Le lemme de Zorn s'applique donc,
et $\mathcal{S}$ possède un élément maximal. D'après l'argument précédent,
on voit que cet élément maximal est un prolongement de $\varphi$
sur tout $X$.
\end{proof}

\begin{lemma}
\label{pione-lemma-local-exact-sequence}
Soit $(A, \mathfrak m)$ un anneau local. Posons $X = \Spec(A)$ et
$U = X \setminus \{\mathfrak m\}$. Soit $U^{sh}$ le spectre épointé
de l'hensélisé strict $A^{sh}$ de $A$.
Supposons $U$ quasi-compact et $U^{sh}$ connexe. Alors la suite
$$
\pi_1(U^{sh}, \overline{u}) \to \pi_1(U, \overline{u}) \to
\pi_1(X, \overline{u}) \to 1
$$
est exacte au sens du lemme \ref{pione-lemma-functoriality-galois-ses}, partie (1).
\end{lemma}

\begin{proof}
L'application $\pi_1(U) \to \pi_1(X)$ est surjective d'après les
lemmes \ref{pione-lemma-same-etale-extensions} et
\ref{pione-lemma-functoriality-galois-surjective}.

\medskip\noindent
Écrivons $X^{sh} = \Spec(A^{sh})$. Soit $Y \to X$ un morphisme fini étale.
Alors $Y^{sh} = Y \times_X X^{sh} \to X^{sh}$ est un morphisme fini étale.
Puisque $A^{sh}$ est strictement hensélien, on voit que $Y^{sh}$ est isomorphe
à une réunion disjointe de copies de $X^{sh}$. Il en va donc de même pour
$Y \times_X U^{sh}$. Il s'ensuit que la composée
$\pi_1(U^{sh}) \to \pi_1(U) \to \pi_1(X)$ est triviale ; voir
le lemme \ref{pione-lemma-composition-trivial}.

\medskip\noindent
Pour achever la preuve, il suffit, d'après le
lemme \ref{pione-lemma-functoriality-galois-ses},
de montrer ce qui suit : étant donné un morphisme fini étale
$V \to U$ tel que $V \times_U U^{sh}$ soit une réunion
disjointe de copies de $U^{sh}$, on peut trouver un morphisme fini étale
$Y \to X$ avec $V \cong Y \times_X U$ sur $U$.
L'hypothèse implique qu'il existe un morphisme fini étale
$Y^{sh} \to X^{sh}$ et un isomorphisme
$V \times_U U^{sh} \cong Y^{sh} \times_{X^{sh}} U^{sh}$.
Considérons le diagramme suivant
$$
\xymatrix{
U \ar[d] & U^{sh} \ar[d] \ar[l] &
U^{sh} \times_U U^{sh} \ar[d] \ar@<1ex>[l] \ar@<-1ex>[l] &
U^{sh} \times_U U^{sh} \times_U U^{sh}
\ar[d] \ar@<1ex>[l] \ar[l] \ar@<-1ex>[l] \\
X & X^{sh} \ar[l] &
X^{sh} \times_X X^{sh} \ar@<1ex>[l] \ar@<-1ex>[l] &
X^{sh} \times_X X^{sh} \times_X X^{sh} \ar@<1ex>[l] \ar[l] \ar@<-1ex>[l]
}
$$
Puisque l'immersion ouverte $U \subset X$ est quasi-compacte par hypothèse, toutes les
flèches verticales sont des immersions ouvertes quasi-compactes.
Soit $\xi \in X^{sh} \times_X X^{sh}$ un point qui n'appartient pas
à $U^{sh} \times_U U^{sh}$. Alors $\xi$ est au-dessus du point
fermé $x^{sh}$ de $X^{sh}$.
Considérons l'homomorphisme d'anneaux locaux
$$
A^{sh} = \mathcal{O}_{X^{sh}, x^{sh}} \to
\mathcal{O}_{X^{sh} \times_X X^{sh}, \xi}
$$
déterminé par la première projection $X^{sh} \times_X X^{sh}$.
C'est une colimite filtrante d'homomorphismes locaux obtenus par
localisation d'homomorphismes d'anneaux étales.
Puisque $A^{sh}$ est strictement hensélien, on conclut que c'est un
isomorphisme. Comme ceci vaut pour tout $\xi$ du complémentaire,
il n'y a pas de spécialisations entre ces points ;
chacun de ces $\xi$ est donc un point fermé (on peut aussi le démontrer
directement). Puisque l'anneau local en $\xi$ est isomorphe
à $A^{sh}$, il est strictement hensélien et son spectre épointé est connexe.
Il en va de même pour les points $\xi$ de $X^{sh} \times_X X^{sh} \times_X X^{sh}$ qui ne sont pas
dans $U^{sh} \times_U U^{sh} \times_U U^{sh}$. Il résulte du
lemme \ref{pione-lemma-retrocompact-dense-open-connected-at-infinity-closed}
que l'image réciproque par les flèches verticales induit des foncteurs pleinement fidèles sur
les catégories de schémas finis étales. Ainsi, la
donnée de descente canonique sur $V \times_U U^{sh}$ relativement au
recouvrement fpqc $\{U^{sh} \to U\}$ se transporte en une
donnée de descente pour $Y^{sh}$ relativement au recouvrement fpqc $\{X^{sh} \to X\}$.
Puisque $Y^{sh} \to X^{sh}$ est fini, donc affine, cette donnée de descente est
effective (Descente, lemme \ref{descent-lemma-affine}).
On obtient ainsi un morphisme affine $Y \to X$ et un isomorphisme
$Y \times_X X^{sh} \to Y^{sh}$ compatible avec les données de descente.
Par pleine fidélité des données de descente
(comme dans Descente, lemme \ref{descent-lemma-refine-coverings-fully-faithful}),
on obtient un isomorphisme $V \to U \times_X Y$.
Enfin, $Y \to X$ est fini étale puisque $Y^{sh} \to X^{sh}$ l'est ; voir
Descente, lemmes \ref{descent-lemma-descending-property-etale} et
\ref{descent-lemma-descending-property-finite}.
\end{proof}

\noindent
Soit $X$ un schéma irréductible. Soit $\eta \in X$ le point générique.
Le morphisme canonique $\eta \to X$ induit une application canonique
\begin{equation}
\label{pione-equation-inclusion-generic-point}
\text{Gal}(\kappa(\eta)^{sep}/\kappa(\eta)) = \pi_1(\eta, \overline{\eta})
\longrightarrow \pi_1(X, \overline{\eta})
\end{equation}
L'identification du membre de gauche est donnée par le
lemme \ref{pione-lemma-fundamental-group-Galois-group}.

\begin{lemma}
\label{pione-lemma-irreducible-geometrically-unibranch}
Soit $X$ un schéma irréductible et géométriquement unibranche.
Pour tout ouvert non vide $U \subset X$, l'application canonique
$$
\pi_1(U, \overline{u}) \longrightarrow \pi_1(X, \overline{u})
$$
est surjective. L'application (\ref{pione-equation-inclusion-generic-point})
$\pi_1(\eta, \overline{\eta}) \to \pi_1(X, \overline{\eta})$
est également surjective.
\end{lemma}

\begin{proof}
D'après le lemme \ref{pione-lemma-thickening}, on peut remplacer $X$ par son réduit.
On peut donc supposer que $X$ est un schéma intègre. D'après le
lemme \ref{pione-lemma-functoriality-galois-surjective},
l'assertion du lemme revient à affirmer que
les foncteurs $\textit{F\'Et}_X \to \textit{F\'Et}_U$ et
$\textit{F\'Et}_X \to \textit{F\'Et}_\eta$ sont pleinement fidèles.

\medskip\noindent
Le résultat pour $\textit{F\'Et}_X \to \textit{F\'Et}_U$ découle
du lemme \ref{pione-lemma-quasi-compact-dense-open-connected-at-infinity}
et du fait que, pour un anneau local $A$ qui est
géométriquement unibranche, son hensélisé strict possède un
spectre irréductible. Voir
Compléments d'algèbre, lemme \ref{more-algebra-lemma-geometrically-unibranch}.

\medskip\noindent
Observons que le corps résiduel $\kappa(\eta) = \mathcal{O}_{X, \eta}$
est la colimite filtrante des $\mathcal{O}_X(U)$, où $U \subset X$
parcourt les ouverts affines non vides. Ainsi, $\textit{F\'Et}_\eta$ est la colimite des
catégories $\textit{F\'Et}_U$ sur de tels $U$ ; voir
Limites, lemmes \ref{limits-lemma-descend-finite-presentation},
\ref{limits-lemma-descend-finite-finite-presentation} et
\ref{limits-lemma-descend-etale}.
Un argument formel montre alors que la pleine fidélité de
$\textit{F\'Et}_X \to \textit{F\'Et}_\eta$ découle de la
pleine fidélité des foncteurs $\textit{F\'Et}_X \to \textit{F\'Et}_U$.
\end{proof}

\begin{lemma}
\label{pione-lemma-exact-sequence-finite-nr-closed-pts}
Soit $X$ un schéma. Soient $x_1, \ldots, x_n \in X$ des points fermés en nombre fini
tels que
\begin{enumerate}
\item $U = X \setminus \{x_1, \ldots, x_n\}$ soit connexe et soit
un ouvert rétrocompact de $X$, et
\item pour chaque $i$, le spectre épointé $U_i^{sh}$ du
hensélisé strict de $\mathcal{O}_{X, x_i}$ soit connexe.
\end{enumerate}
Alors l'application $\pi_1(U) \to \pi_1(X)$ est surjective et son noyau
est le plus petit sous-groupe distingué fermé de $\pi_1(U)$ qui contienne
l'image de $\pi_1(U_i^{sh}) \to \pi_1(U)$ pour $i = 1, \ldots, n$.
\end{lemma}

\begin{proof}
La surjectivité découle des
lemmes \ref{pione-lemma-retrocompact-dense-open-connected-at-infinity-closed} et
\ref{pione-lemma-functoriality-galois-surjective}.
Nous pouvons considérer la suite d'applications
$$
\pi_1(U)  \to \ldots \to
\pi_1(X \setminus \{x_1, x_2\}) \to \pi_1(X \setminus \{x_1\}) \to \pi_1(X)
$$
Un argument de théorie des groupes montre alors qu'il suffit de démontrer l'assertion sur le
noyau dans le cas $n = 1$ (détails omis). Écrivons
$x = x_1$, $U^{sh} = U_1^{sh}$,
posons $A = \mathcal{O}_{X, x}$ et soit $A^{sh}$ l'hensélisé strict.
Considérons le diagramme
$$
\xymatrix{
U \ar[d] &
\Spec(A) \setminus \{\mathfrak m\} \ar[l] \ar[d] &
U^{sh} \ar[d] \ar[l] \\
X & \Spec(A) \ar[l] & \Spec(A^{sh}) \ar[l]
}
$$
D'après le lemme \ref{pione-lemma-functoriality-galois-ses},
il faut montrer que les morphismes finis étales
$V \to U$ dont les images réciproques sont des revêtements triviaux de $U^{sh}$
se prolongent en des schémas finis étales sur $X$.
D'après le lemme \ref{pione-lemma-local-exact-sequence},
nous connaissons l'énoncé correspondant
pour les schémas finis étales sur le spectre épointé de $A$.
Cependant, d'après Limites, lemme \ref{limits-lemma-glueing-near-closed-point},
les schémas de présentation finie sur $X$ s'identifient aux données formées de
schémas de présentation finie sur $U$ et sur $A$, recollés sur
le spectre épointé de $A$. Ceci achève la preuve.
\end{proof}









\section{Groupes fondamentaux des schémas normaux}
\label{pione-section-normal}

\noindent
Soit $X$ un schéma intègre et géométriquement unibranche. Dans la section précédente,
nous avons vu que le groupe fondamental de $X$ est un quotient du
groupe de Galois du corps des fonctions rationnelles $K$ de $X$. Puisque l'application est continue,
son noyau est un sous-groupe distingué fermé du groupe de Galois. Ce noyau
correspond donc à une extension galoisienne $M/K$ par la théorie de Galois
(Corps, théorème \ref{fields-theorem-inifinite-galois-theory}).
Dans cette section, nous allons déterminer $M$ lorsque $X$ est un schéma intègre normal.

\medskip\noindent
Soit $X$ un schéma intègre normal de corps des fonctions rationnelles $K$.
Soit $L/K$ une extension finie. Considérons la normalisation
$Y \to X$ de $X$ dans le morphisme $\Spec(L) \to X$, telle qu'elle est définie dans
Morphismes, section \ref{morphisms-section-normalization-X-in-Y}.
Nous dirons (dans cette situation) que {\it $X$ est non ramifié dans $L$}
si $Y \to X$ est un morphisme non ramifié de schémas. Dans le
lemme \ref{pione-lemma-unramified}, nous expliciterons cette condition.
Observons que la fibre schématique de $Y \to X$ au-dessus de $\Spec(K)$
est $\Spec(L)$. Ainsi, l'extension de corps $L/K$ est séparable si $X$ est
non ramifié dans $L$ ; voir
Morphismes, lemme \ref{morphisms-lemma-unramified-over-field}.

\begin{lemma}
\label{pione-lemma-unramified-in-L}
Dans la situation ci-dessus, les conditions suivantes sont équivalentes :
\begin{enumerate}
\item $X$ est non ramifié dans $L$,
\item $Y \to X$ est étale, et
\item $Y \to X$ est fini étale.
\end{enumerate}
\end{lemma}

\begin{proof}
Observons que $Y \to X$ est un morphisme entier.
Dans chacun des cas, le morphisme $Y \to X$ est localement de type fini
par définition.
On trouve donc que, dans chacun des cas, $Y \to X$ est fini d'après
Morphismes, lemme \ref{morphisms-lemma-finite-integral}.
En particulier, on voit que (2) est équivalent à (3).
Un morphisme étale est non ramifié ; ainsi (2) implique (1).

\medskip\noindent
Réciproquement, supposons $Y \to X$ non ramifié. Puisqu'un schéma normal est
géométriquement unibranche
(Propriétés, lemme \ref{properties-lemma-normal-geometrically-unibranch}),
on voit que le morphisme $Y \to X$ est étale d'après Compléments sur les morphismes, lemme
\ref{more-morphisms-lemma-global-unramified-dominant-unibranch-is-etale}.
Nous donnons aussi une preuve directe dans le paragraphe suivant.

\medskip\noindent
Soit $x \in X$.
On peut choisir un voisinage étale $(U, u) \to (X, x)$ tel que
$$
Y \times_X U = \coprod V_j \longrightarrow U
$$
est une réunion disjointe d'immersions fermées, voir
Morphismes étales, lemme \ref{etale-lemma-finite-unramified-etale-local}.
En rétrécissant, on peut supposer $U$ quasi-compact.
Alors $U$ possède un nombre fini de composantes irréductibles
(Descente, lemme \ref{descent-lemma-locally-finite-nr-irred-local-fppf}).
Puisque $U$ est normal
(Descente, lemme \ref{descent-lemma-normal-local-smooth}), les
composantes irréductibles de $U$ sont ouvertes et fermées
(Propriétés, lemme \ref{properties-lemma-normal-locally-finite-nr-irreducibles})
et on peut supposer $U$ irréductible. Alors $U$ est un schéma intègre
dont le point générique $\xi$ s'envoie sur le point générique de $X$.
D'autre part, nous savons que $Y \times_X U$
est la normalisation de $U$ dans $\Spec(L) \times_X U$
d'après Compléments sur les morphismes, lemme
\ref{more-morphisms-lemma-normalization-smooth-localization}.
Tout point de $\Spec(L) \times_X U$ s'envoie sur $\xi$.
Ainsi, tout $V_j$ contient un point s'envoyant sur $\xi$ d'après
Morphismes, lemme \ref{morphisms-lemma-normalization-generic}.
Ainsi $V_j \to U$ est un isomorphisme puisque $U = \overline{\{\xi\}}$.
Par conséquent, $Y \times_X U \to U$ est étale. D'après
Descente, lemme \ref{descent-lemma-descending-property-etale},
on conclut que $Y \to X$ est étale au-dessus de
l'image de $U \to X$ (un voisinage ouvert de $x$).
\end{proof}

\begin{lemma}
\label{pione-lemma-finite-etale-covering-normal-unramified}
Soit $X$ un schéma intègre normal de corps des fonctions rationnelles $K$.
Soit $Y \to X$ un morphisme fini étale. Si $Y$ est connexe,
alors $Y$ est un schéma intègre normal et $Y$ est la normalisation
de $X$ dans le corps des fonctions rationnelles de $Y$.
\end{lemma}

\begin{proof}
Le schéma $Y$ est normal d'après
Descente, lemme \ref{descent-lemma-normal-local-smooth}.
Puisque $Y \to X$ est plat, tout point générique de $Y$ s'envoie
sur le point générique de $X$ d'après
Morphismes, lemme \ref{morphisms-lemma-generalizations-lift-flat}.
Puisque $Y \to X$ est fini, on voit que $Y$ possède un nombre fini
de composantes irréductibles. Ainsi $Y$ est la réunion disjointe d'un
nombre fini de schémas intègres normaux d'après
Propriétés, lemme \ref{properties-lemma-normal-locally-finite-nr-irreducibles}.
Ainsi, si $Y$ est connexe, alors $Y$ est un schéma intègre normal.

\medskip\noindent
Soit $L$ le corps des fonctions rationnelles de $Y$ et soit $Y' \to X$ la normalisation
de $X$ dans $L$. D'après
Morphismes, lemme \ref{morphisms-lemma-characterize-normalization},
on obtient une factorisation $Y' \to Y \to X$ et $Y' \to Y$ est
la normalisation de $Y$ dans $L$. Puisque $Y$ est normal, il est clair
que $Y' = Y$ (cela peut aussi se déduire de
Morphismes, lemme \ref{morphisms-lemma-finite-birational-over-normal}).
\end{proof}

\begin{proposition}
\label{pione-proposition-normal}
Soit $X$ un schéma intègre normal de corps des fonctions rationnelles $K$.
Alors l'application canonique (\ref{pione-equation-inclusion-generic-point})
$$
\text{Gal}(K^{sep}/K) = \pi_1(\eta, \overline{\eta})
\longrightarrow \pi_1(X, \overline{\eta})
$$
s'identifie à la projection canonique
$\text{Gal}(K^{sep}/K) \to \text{Gal}(M/K)$ où $M \subset K^{sep}$
est la réunion des sous-extensions finies $L$
telles que $X$ est non ramifié dans $L$.
\end{proposition}

\begin{proof}
Le schéma normal $X$ est géométriquement unibranche
(Propriétés, lemme \ref{properties-lemma-normal-geometrically-unibranch}).
Le lemme \ref{pione-lemma-irreducible-geometrically-unibranch} s'applique donc à $X$.
Ainsi $\pi_1(\eta, \overline{\eta}) \to \pi_1(X, \overline{\eta})$
est surjectif et la flèche horizontale supérieure du diagramme commutatif
$$
\xymatrix{
\textit{F\'Et}_X \ar[r] \ar[d] \ar[rd]_c & \textit{F\'Et}_\eta \ar[d] \\
\textit{Finite-}\pi_1(X, \overline{\eta})\textit{-sets} \ar[r] &
\textit{Finite-}\text{Gal}(K^{sep}/K)\textit{-sets}
}
$$
est pleinement fidèle. La flèche verticale de gauche est l'équivalence du
théorème \ref{pione-theorem-fundamental-group}
et la flèche verticale de droite est l'équivalence du
lemme \ref{pione-lemma-fundamental-group-Galois-group}. La flèche
horizontale inférieure est induite par l'application de la proposition.
D'après les lemmes \ref{pione-lemma-unramified-in-L} et
\ref{pione-lemma-finite-etale-covering-normal-unramified},
on voit que l'image essentielle de $c$
est constituée des $\text{Gal}(K^{sep}/K)\textit{-Ensembles}$ isomorphes
à des ensembles de la forme
$$
S = \Hom_K(\prod\nolimits_{i = 1, \ldots, n} L_i, K^{sep}) =
\coprod\nolimits_{i = 1, \ldots, n} \Hom_K(L_i, K^{sep})
$$
où $L_i/K$ est finie séparable et où $X$ est non ramifié dans $L_i$.
Ainsi, si $M \subset K^{sep}$ est comme dans l'énoncé du lemme,
alors $\text{Gal}(K^{sep}/M)$ est exactement le sous-groupe de
$\text{Gal}(K^{sep}/K)$ qui agit trivialement sur tout objet
de l'image essentielle de $c$. D'autre part, l'image essentielle de $c$
est exactement la catégorie des $S$ tels que l'action de $\text{Gal}(K^{sep}/K)$
se factorise par la surjection
$\text{Gal}(K^{sep}/K) \to \pi_1(X, \overline{\eta})$.
On conclut que $\text{Gal}(K^{sep}/M)$ est le noyau.
Ainsi $\text{Gal}(K^{sep}/M)$ est un sous-groupe distingué, $M/K$ est galoisienne,
et on a une suite exacte courte
$$
1 \to \text{Gal}(K^{sep}/M) \to
\text{Gal}(K^{sep}/K) \to
\text{Gal}(M/K) \to 1
$$
d'après la théorie de Galois (Corps, théorème
\ref{fields-theorem-inifinite-galois-theory} et
lemme \ref{fields-lemma-ses-infinite-galois}). La démonstration est terminée.
\end{proof}

\begin{lemma}
\label{pione-lemma-local-exact-sequence-normal}
Soit $(A, \mathfrak m)$ un anneau local normal.
Posons $X = \Spec(A)$. Soit $A^{sh}$ l'hensélisé strict de $A$.
Soient $K$ et $K^{sh}$ les corps des fractions de $A$ et de $A^{sh}$.
Alors la suite
$$
\pi_1(\Spec(K^{sh})) \to \pi_1(\Spec(K)) \to \pi_1(X) \to 1
$$
est exacte au sens de la partie (1) du lemme \ref{pione-lemma-functoriality-galois-ses}.
\end{lemma}

\begin{proof}
Observons que $A^{sh}$ est un domaine normal, voir
Compléments d'algèbre, lemme \ref{more-algebra-lemma-henselization-normal}.
L'application $\pi_1(\Spec(K)) \to \pi_1(X)$ est surjective d'après la
proposition \ref{pione-proposition-normal}.

\medskip\noindent
Écrivons $X^{sh} = \Spec(A^{sh})$. Soit $Y \to X$ un morphisme fini étale.
Alors $Y^{sh} = Y \times_X X^{sh} \to X^{sh}$ est un morphisme fini étale.
Puisque $A^{sh}$ est strictement hensélien, on voit que $Y^{sh}$ est isomorphe
à une réunion disjointe de copies de $X^{sh}$. Il en va donc de même pour
$Y \times_X \Spec(K^{sh})$. Il s'ensuit que la composée
$\pi_1(\Spec(K^{sh})) \to \pi_1(X)$ est triviale, voir
lemme \ref{pione-lemma-composition-trivial}.

\medskip\noindent
Pour achever la démonstration, il suffit, d'après le
lemme \ref{pione-lemma-functoriality-galois-ses},
de montrer ce qui suit : étant donné un morphisme fini étale
$V \to \Spec(K)$ tel que $V \times_{\Spec(K)} \Spec(K^{sh})$
soit une réunion disjointe de copies de $\Spec(K^{sh})$, on peut trouver un
morphisme fini étale
$Y \to X$ avec $V \cong Y \times_X \Spec(K)$ sur $\Spec(K)$.
Écrivons $V = \Spec(L)$, de sorte que $L$ est un produit fini
d'extensions finies séparables de $K$.
Soit $B \subset L$ la clôture intégrale de $A$ dans $L$.
Si $A \to B$ est étale, alors on peut prendre $Y = \Spec(B)$
et la démonstration est terminée. D'après
Algèbre, lemme \ref{algebra-lemma-integral-closure-commutes-smooth}
(et un passage à la limite que nous omettons),
on voit que $B \otimes_A A^{sh}$ est la clôture intégrale de
$A^{sh}$ dans $L^{sh} = L \otimes_K K^{sh}$.
Notre hypothèse est que $L^{sh}$ est un produit de copies de
$K^{sh}$ et donc $B^{sh}$ est un produit de copies de $A^{sh}$.
Ainsi $A^{sh} \to B^{sh}$ est étale. Comme $A \to A^{sh}$ est
fidèlement plat, il s'ensuit que $A \to B$ est étale
(Descente, lemme \ref{descent-lemma-descending-property-etale}),
comme souhaité.
\end{proof}







\section{Actions de groupes et clôture intégrale}
\label{pione-section-group-actions-integral}

\noindent
Dans cette section, on poursuit la discussion de
Compléments d'algèbre, section \ref{more-algebra-section-group-actions-integral}.
Rappelons qu'un anneau local normal est un domaine par définition.

\begin{lemma}
\label{pione-lemma-get-algebraic-closure}
Soit $A$ un domaine normal dont le corps des fractions $K$ est séparablement
clos. Soit $\mathfrak p \subset A$ un idéal premier non nul.
Alors le corps résiduel $\kappa(\mathfrak p)$ est algébriquement clos.
\end{lemma}

\begin{proof}
Supposons le lemme faux afin d'obtenir une contradiction. Il existe alors un
polynôme unitaire irréductible $P(T) \in \kappa(\mathfrak p)[T]$ de
degré $d > 1$. Après avoir remplacé $P$ par $a^d P(a^{-1}T)$ pour un $a \in A$
convenable (afin d'éliminer les dénominateurs), on peut supposer que $P$ est l'image d'un
polynôme unitaire $Q$ dans $A[T]$. Observons que $Q$ est irréductible dans
$K[T]$. En effet, une factorisation sur $K$ conduit à une factorisation
sur $A$ d'après Algèbre, lemme \ref{algebra-lemma-polynomials-divide},
que l'on pourrait réduire modulo $\mathfrak p$ pour obtenir une factorisation de $P$.
Comme $K$ est séparablement clos, $Q$ n'est pas un polynôme séparable
(Corps, définition \ref{fields-definition-separable}).
La caractéristique de $K$ est alors $p > 0$ et $Q$ a
un terme linéaire nul (Corps, définition \ref{fields-definition-separable}).
Cependant, on peut alors remplacer $Q$ par
$Q + a T$, où $a \in \mathfrak p$ est non nul, pour obtenir une contradiction.
\end{proof}

\begin{lemma}
\label{pione-lemma-normal-local-domain-separablly-closed-fraction-field}
Un anneau local normal dont le corps des fractions est séparablement clos est
strictement hensélien.
\end{lemma}

\begin{proof}
Soit $(A, \mathfrak m, \kappa)$ un anneau local normal dont le corps des fractions
$K$ est séparablement clos. Si $A = K$, alors c'est terminé. Sinon,
le corps résiduel $\kappa$ est algébriquement clos
d'après le lemme \ref{pione-lemma-get-algebraic-closure} et il suffit de
vérifier que $A$ est hensélien.
Soit $f \in A[T]$ unitaire et soit $a_0 \in \kappa$ une racine
de multiplicité $1$ de la réduction $\overline{f} \in \kappa[T]$.
Soit $f = \prod f_i$ la factorisation dans $K[T]$.
D'après Algèbre, lemme \ref{algebra-lemma-polynomials-divide}, on a
$f_i \in A[T]$. Ainsi $a_0$ est une racine de $f_i$ pour un certain $i$.
Après avoir remplacé $f$ par $f_i$, on peut supposer $f$ irréductible.
Alors, puisque la dérivée $f'$ ne peut être nulle dans $A[T]$
car $a_0$ est une racine simple, on conclut que $f$ est linéaire
du fait que $K$ est séparablement clos.
Ainsi $A$ est hensélien, voir
Algèbre, définition \ref{algebra-definition-henselian}.
\end{proof}

\begin{lemma}
\label{pione-lemma-inertia-base-change}
Soit $G$ un groupe fini agissant sur un anneau $R$. Soit $R^G \to A$ un morphisme
d'anneaux. Soit $\mathfrak q' \subset A \otimes_{R^G} R$ un idéal premier au-dessus
de l'idéal premier $\mathfrak q \subset R$. Alors
$$
I_\mathfrak q = \{\sigma \in G \mid
\sigma(\mathfrak q) = \mathfrak q\text{ et }
\sigma \bmod \mathfrak q = \text{id}_{\kappa(\mathfrak q)}\}
$$
est égal à
$$
I_{\mathfrak q'} = \{\sigma \in G \mid
\sigma(\mathfrak q') = \mathfrak q'\text{ et }
\sigma \bmod \mathfrak q' = \text{id}_{\kappa(\mathfrak q')}\}
$$
\end{lemma}

\begin{proof}
Puisque $\mathfrak q$ est l'image réciproque de $\mathfrak q'$
et puisque $\kappa(\mathfrak q) \subset \kappa(\mathfrak q')$,
on obtient $I_{\mathfrak q'} \subset I_\mathfrak q$.
Réciproquement, si $\sigma \in I_\mathfrak q$, alors $\sigma$
agit trivialement sur l'anneau de la fibre $A \otimes_{R^G} \kappa(\mathfrak q)$.
Ainsi $\sigma$ fixe tous les idéaux premiers au-dessus de $\mathfrak q$
et induit l'identité sur leurs corps résiduels.
\end{proof}

\begin{lemma}
\label{pione-lemma-inertia-invariants-etale}
Soit $G$ un groupe fini agissant sur un anneau $R$. Soit $\mathfrak q \subset R$
un idéal premier. Posons
$$
I = \{\sigma \in G \mid \sigma(\mathfrak q) = \mathfrak q
\text{ et } \sigma \bmod \mathfrak q = \text{id}_{\kappa(\mathfrak q)}\}
$$
Alors $R^G \to R^I$ est étale en $R^I \cap \mathfrak q$.
\end{lemma}

\begin{proof}
La stratégie de la démonstration consiste à utiliser une localisation étale pour
se ramener au cas où $R^G \to R^I$ est un isomorphisme local en
$R^I \cap \mathfrak q$.
Soit $R^G \to A$ un morphisme étale d'anneaux. On affirme que si le résultat
vaut pour l'action de $G$ sur $A \otimes_{R^G} R$ et un idéal premier
$\mathfrak q'$ de $A \otimes_{R^G} R$ au-dessus de $\mathfrak q$, alors
le résultat est vrai.

\medskip\noindent
Pour le vérifier, observons que, puisque $R^G \to A$ est plat, on a
$A = (A \otimes_{R^G} R)^G$, voir Compléments d'algèbre,
lemme \ref{more-algebra-lemma-base-change-invariants}.
D'après le lemme \ref{pione-lemma-inertia-base-change}, le groupe $I$ ne change pas.
Une seconde application de Compléments d'algèbre,
lemme \ref{more-algebra-lemma-base-change-invariants},
montre alors que $A \otimes_{R^G} R^I = (A \otimes_{R^G} R)^I$
(parce que $R^I \to A \otimes_{R^G} R^I$ est plat).
Ainsi
$$
\xymatrix{
\Spec((A \otimes_{R^G} R)^I) \ar[d] \ar[r] & \Spec(R^I) \ar[d] \\
\Spec(A) \ar[r] & \Spec(R^G)
}
$$
est cartésien et les flèches horizontales sont étales. Ainsi, si la
flèche verticale de gauche est étale dans un voisinage ouvert $W$ de
$(A \otimes_{R^G} R)^I \cap \mathfrak q'$, alors la flèche verticale de droite
est étale aux points de l'image (ouverte) de $W$ dans
$\Spec(R^I)$, voir
Descente, lemme \ref{descent-lemma-smooth-permanence}. En particulier,
le morphisme $\Spec(R^I) \to \Spec(R^G)$ est étale en $R^I \cap \mathfrak q$.

\medskip\noindent
Soit $\mathfrak p = R^G \cap \mathfrak q$.
D'après Compléments d'algèbre, lemme \ref{more-algebra-lemma-one-orbit},
la fibre de $\Spec(R) \to \Spec(R^G)$ au-dessus de $\mathfrak p$ est
finie. De plus, les extensions de corps résiduels en ces points
sont algébriques, normales, à groupes d'automorphismes finis d'après
Compléments d'algèbre, lemme \ref{more-algebra-lemma-one-orbit-geometric}.
On peut donc appliquer
Compléments sur les morphismes,
lemme \ref{more-morphisms-lemma-etale-makes-integral-split},
au morphisme d'anneaux entier $R^G \to R$ et à l'idéal premier $\mathfrak p$.
Combiné à l'affirmation ci-dessus, cela nous ramène au cas où
$R = A_1 \times \ldots \times A_n$, chaque $A_i$ possédant un unique
idéal premier $\mathfrak q_i$ au-dessus de $\mathfrak p$, tel que les
extensions de corps résiduels $\kappa(\mathfrak q_i)/\kappa(\mathfrak p)$
soient purement inséparables. Bien entendu, $\mathfrak q$ est l'un de
ces idéaux premiers, disons $\mathfrak q = \mathfrak q_1$.

\medskip\noindent
Il se peut que $G$ ne permute pas les facteurs $A_i$
(ce serait vrai si le spectre de $A_i$ était connexe,
par exemple si $R^G$ était local). On peut y remédier comme suit ;
nous invitons le lecteur à le vérifier lui-même, éventuellement
en utilisant des idempotents plutôt que la topologie.
Rappelons que la décomposition en produit donne une décomposition correspondante
de $\Spec(R)$ en réunion disjointe de sous-ensembles ouverts et fermés
$U_i$. Puisque $G$ est fini, on peut raffiner ce recouvrement
en une décomposition en réunion disjointe finie
$\Spec(R) = \coprod_{j \in J} W_j$ par des sous-ensembles ouverts
et fermés $W_j$, telle que, pour tout $j \in J$, il existe
un $j' \in J$ avec $\sigma(W_j) = W_{j'}$. La réunion des
$W_j$ qui ne rencontrent pas $\{\mathfrak q_1, \ldots, \mathfrak q_n\}$
est un sous-ensemble fermé qui ne rencontre pas la fibre au-dessus de $\mathfrak p$,
donc s'envoie sur un sous-ensemble fermé de $\Spec(R^G)$ qui ne contient pas
$\mathfrak p$, car $\Spec(R) \to \Spec(R^G)$ est fermé.
Ainsi, après avoir remplacé $R^G$ par une localisation principale
(ce qui est permis par l'affirmation), on peut supposer que chaque $W_j$ rencontre
l'un des points $\mathfrak q_i$. On pose alors $U_i = W_j$
si $\mathfrak q_i \in W_j$. La décomposition en produit correspondante
$R = A_1 \times \ldots \times A_n$ est telle
que $G$ permute les facteurs $A_i$.

\medskip\noindent
On peut donc supposer que l'on dispose d'une décomposition en produit
$R = A_1 \times \ldots \times A_n$ compatible avec l'action de $G$,
où chaque $A_i$ possède un unique idéal premier $\mathfrak q_i$ au-dessus
de $\mathfrak p$ et les extensions de corps
$\kappa(\mathfrak q_i)/\kappa(\mathfrak p)$ sont purement inséparables.
Écrivons $A' = A_2 \times \ldots \times A_n$, de sorte que
$$
R = A_1 \times A'
$$
Puisque $\mathfrak q = \mathfrak q_1$, on trouve que tout
$\sigma \in I$ préserve la décomposition en produit ci-dessus.
Par conséquent,
$$
R^I = (A_1)^I \times (A')^I
$$
Observons que $I = D = \{\sigma \in G \mid \sigma(\mathfrak q) = \mathfrak q\}$
parce que $\kappa(\mathfrak q)/\kappa(\mathfrak p)$ est purement inséparable.
Puisque l'action de $G$ sur les idéaux premiers au-dessus de $\mathfrak p$ est transitive
(Compléments d'algèbre, lemme \ref{more-algebra-lemma-one-orbit}),
on conclut que l'indice de $I$ dans $G$ est $n$ et que l'on peut écrire
$G = eI \amalg \sigma_2I \amalg \ldots \amalg \sigma_nI$, de sorte que
$A_i = \sigma_i(A_1)$ pour $i = 2, \ldots, n$. Il s'ensuit que
$$
R^G = (A_1)^I.
$$
Ainsi, l'application $R^G \to R^I$ est étale en $R^I \cap \mathfrak q$
et la démonstration est terminée.
\end{proof}

\noindent
Le lemme suivant généralise
Compléments d'algèbre, lemme \ref{more-algebra-lemma-inertial-invariants-unramified}.

\begin{lemma}
\label{pione-lemma-inertial-invariants-unramified}
Soit $A$ un domaine normal de corps des fractions $K$.
Soit $L/K$ une extension galoisienne (éventuellement infinie).
Soit $G = \text{Gal}(L/K)$ et soit
$B$ la clôture intégrale de $A$ dans $L$.
Soit $\mathfrak q \subset B$. Posons
$$
I = \{\sigma \in G \mid
\sigma(\mathfrak q) = \mathfrak q \text{ et }
\sigma \bmod \mathfrak q = \text{id}_{\kappa(\mathfrak q)}\}
$$
Alors $(B^I)_{B^I \cap \mathfrak q}$ est une colimite filtrante
d'$A$-algèbres étales.
\end{lemma}

\begin{proof}
On peut écrire $L$ comme colimite filtrante d'extensions galoisiennes finies
de $K$. Il suffit donc de démontrer ce lemme lorsque $L/K$ est
une extension galoisienne finie, voir
Algèbre, lemme \ref{algebra-lemma-colimit-colimit-etale}.
Puisque $A = B^G$, car $A$ est intégralement
clos dans $K = L^G$, le résultat découle du
lemme \ref{pione-lemma-inertia-invariants-etale}.
\end{proof}







\section{Théorie de la ramification}
\label{pione-section-ramification}

\noindent
Dans cette section, on poursuit la discussion de
Compléments d'algèbre, section \ref{more-algebra-section-ramification},
et on la relie à notre discussion des groupes fondamentaux des schémas.

\medskip\noindent
Soit $(A, \mathfrak m, \kappa)$ un anneau local normal de
corps des fractions $K$. Choisissons une clôture séparable $K^{sep}$. Soit
$A^{sep}$ la clôture intégrale de $A$ dans $K^{sep}$.
Choisissons un idéal maximal $\mathfrak m^{sep} \subset A^{sep}$.
Soient $A \subset A^h \subset A^{sh}$ l'hensélisé et
l'hensélisé strict. Observons que $A^h$ et $A^{sh}$ sont eux aussi des anneaux normaux
(Compléments d'algèbre, lemme \ref{more-algebra-lemma-henselization-normal}).
Notons $K^h$ et $K^{sh}$ leurs corps des fractions.
Puisque $(A^{sep})_{\mathfrak m^{sep}}$ est strictement hensélien d'après le
lemme \ref{pione-lemma-normal-local-domain-separablly-closed-fraction-field},
on peut choisir un morphisme de $A$-algèbres $A^{sh} \to (A^{sep})_{\mathfrak m^{sep}}$.
En effet, choisissons d'abord un $\kappa$-plongement\footnote{Cela est possible
parce que $\kappa(\mathfrak m^{sh})$ est une clôture séparable
de $\kappa$ et que $\kappa(\mathfrak m^{sep})$ est une clôture algébrique
de $\kappa$ d'après le lemme \ref{pione-lemma-get-algebraic-closure}.}
$\kappa(\mathfrak m^{sh}) \to \kappa(\mathfrak m^{sep})$ et
prolongeons-le ensuite (de manière unique) en un homomorphisme de $A$-algèbres d'après
Algèbre, lemme \ref{algebra-lemma-strictly-henselian-functorial}.
On obtient le diagramme suivant
$$
\xymatrix{
K^{sep} & K^{sh} \ar[l] & K^h \ar[l] & K \ar[l] \\
(A^{sep})_{\mathfrak m^{sep}} \ar[u] &
A^{sh} \ar[u] \ar[l] &
A^h \ar[u] \ar[l] &
A \ar[u] \ar[l]
}
$$
On peut prendre les groupes fondamentaux des spectres de ces anneaux.
Bien entendu, puisque $K^{sep}$, $(A^{sep})_{\mathfrak m^{sep}}$ et
$A^{sh}$ sont strictement henséliens, on obtient pour eux des groupes triviaux.
La partie intéressante est donc la suivante
\begin{equation}
\label{pione-equation-inertia-diagram-pione}
\vcenter{
\xymatrix{
\pi_1(U^{sh}) \ar[r] \ar[rd]_1 & \pi_1(U^h) \ar[d] \ar[r] & \pi_1(U) \ar[d] \\
& \pi_1(X^h) \ar[r] & \pi_1(X)
}
}
\end{equation}
Ici, $X^h$ et $X$ sont les spectres de $A^h$ et $A$, et
$U^{sh}$, $U^h$, $U$ sont les spectres de $K^{sh}$, $K^h$ et $K$.
L'indice $1$ signifie que l'application est triviale ; cela résulte
de ce qu'elle se factorise par le groupe trivial $\pi_1(X^{sh})$.
D'autre part, le groupe profini $G = \text{Gal}(K^{sep}/K)$
agit sur $A^{sep}$ et on peut poser les définitions suivantes
$$
D = \{\sigma \in G \mid \sigma(\mathfrak m^{sep}) = \mathfrak m^{sep}\}
\supset
I = \{\sigma \in D \mid \sigma \bmod \mathfrak m^{sep} =
\text{id}_{\kappa(\mathfrak m^{sep})}\}
$$
Ces groupes sont parfois appelés le
{\it groupe de décomposition} et le {\it groupe d'inertie},
notamment lorsque $A$ est un anneau de valuation discrète.

\begin{lemma}
\label{pione-lemma-identify-inertia}
Dans la situation décrite ci-dessus, via l'isomorphisme
$\pi_1(U) = \text{Gal}(K^{sep}/K)$, le diagramme
(\ref{pione-equation-inertia-diagram-pione})
se traduit par le diagramme
$$
\xymatrix{
I \ar[r] \ar[rd]_1 & D \ar[d] \ar[r] & \text{Gal}(K^{sep}/K) \ar[d] \\
& \text{Gal}(\kappa(\mathfrak m^{sh})/\kappa) \ar[r] & \text{Gal}(M/K)
}
$$
où $K^{sep}/M/K$ est la sous-extension maximale non ramifiée
relativement à $A$. De plus, les flèches verticales sont surjectives,
le noyau de la flèche verticale de gauche est $I$ et le noyau de la
flèche verticale de droite est
le plus petit sous-groupe distingué fermé de $\text{Gal}(K^{sep}/K)$
qui contient $I$.
\end{lemma}

\begin{proof}
Par construction, le groupe $D$ agit sur $(A^{sep})_{\mathfrak m^{sep}}$
au-dessus de $A$. Par l'unicité de $A^{sh} \to (A^{sep})_{\mathfrak m^{sep}}$
une fois donnée l'application sur les corps résiduels
(Algèbre, lemme \ref{algebra-lemma-strictly-henselian-functorial}),
on voit que l'image de $A^{sh} \to (A^{sep})_{\mathfrak m^{sep}}$
est contenue dans $((A^{sep})_{\mathfrak m^{sep}})^I$.
D'autre part, le
lemme \ref{pione-lemma-inertial-invariants-unramified}
montre que $((A^{sep})_{\mathfrak m^{sep}})^I$
est une colimite filtrante d'extensions étales de $A$.
Puisque $A^{sh}$ est l'extension maximale de ce type, on conclut
que $A^{sh} = ((A^{sep})_{\mathfrak m^{sep}})^I$.
Ainsi $K^{sh} = (K^{sep})^I$.

\medskip\noindent
Rappelons que $I$ est le noyau d'une application surjective
$D \to \text{Aut}(\kappa(\mathfrak m^{sep})/\kappa)$, voir
Compléments d'algèbre, lemme \ref{more-algebra-lemma-one-orbit-geometric-galois}.
On a $\text{Aut}(\kappa(\mathfrak m^{sep})/\kappa) =
\text{Gal}(\kappa(\mathfrak m^{sh})/\kappa)$,
car on a vu ci-dessus que ces corps sont respectivement la clôture algébrique
et la clôture séparable de $\kappa$.
D'autre part, tout automorphisme de $A^{sh}$ au-dessus de $A$
est un automorphisme de $A^{sh}$ au-dessus de $A^h$ par l'unicité
dans Algèbre, lemme \ref{algebra-lemma-henselian-functorial}.
En outre, $A^{sh}$ est la colimite des extensions finies étales
$A^h \subset A'$ qui correspondent $1$-à-$1$
aux extensions finies séparables $\kappa'/\kappa$, voir
Algèbre, remarque \ref{algebra-remark-construct-sh-from-h}.
Ainsi
$$
\text{Aut}(A^{sh}/A) = \text{Aut}(A^{sh}/A^h) =
\text{Gal}(\kappa(\mathfrak m^{sh})/\kappa)
$$
Soit $\kappa'/\kappa$ une extension galoisienne finie de
groupe de Galois $G$. Soit $A^h \subset A'$ l'extension finie étale
qui correspond à $\kappa \subset \kappa'$ d'après
Algèbre, lemme \ref{algebra-lemma-henselian-cat-finite-etale}.
Il s'ensuit alors que
$(A')^G = A^h$ en considérant les corps des fractions et les degrés
(on omet ce petit détail). En prenant la colimite, on conclut que
$(A^{sh})^{\text{Gal}(\kappa(\mathfrak m^{sh})/\kappa)} = A^h$.
En combinant tout ce qui précède, on trouve $A^h = ((A^{sep})_{\mathfrak m^{sep}})^D$.
Ainsi $K^h = (K^{sep})^D$.

\medskip\noindent
Puisque $U$, $U^h$, $U^{sh}$ sont les spectres des corps
$K$, $K^h$, $K^{sh}$, on voit que les lignes supérieures des diagrammes
se correspondent via le
lemme \ref{pione-lemma-fundamental-group-Galois-group}.
D'après le lemme \ref{pione-lemma-gabber}, on a
$\pi_1(X^h) = \text{Gal}(\kappa(\mathfrak m^{sh})/\kappa)$.
L'exactitude de la suite
$1 \to I \to D \to \text{Gal}(\kappa(\mathfrak m^{sh})/\kappa) \to 1$
a été signalée ci-dessus.
D'après la proposition \ref{pione-proposition-normal},
on voit que $\pi_1(X) = \text{Gal}(M/K)$.
Enfin, l'assertion sur le noyau de
$\text{Gal}(K^{sep}/K) \to \text{Gal}(M/K) = \pi_1(X)$
découle du lemme \ref{pione-lemma-local-exact-sequence-normal}.
Cela achève la démonstration.
\end{proof}

\noindent
Soit $X$ un schéma intègre normal de corps des fonctions rationnelles $K$.
Soit $K^{sep}$ une clôture séparable de $K$.
Soit $X^{sep} \to X$ la normalisation de $X$ dans $K^{sep}$.
Puisque $G = \text{Gal}(K^{sep}/K)$ agit sur $K^{sep}$,
on obtient une action à droite de $G$ sur $X^{sep}$.
Pour $y \in X^{sep}$, posons
$$
D_y = \{\sigma \in G \mid \sigma(y) = y\} \supset
I_y = \{\sigma \in D_y \mid \sigma \bmod \mathfrak m_y =
\text{id}_{\kappa(y)} \}
$$
comme ci-dessus. D'autre part, pour $x \in X$,
soit $\mathcal{O}_{X, x}^{sh}$ un hensélisé strict,
soit $K_x^{sh}$ le corps des fractions de $\mathcal{O}_{X, x}^{sh}$
et choisissons un $K$-plongement $K_x^{sh} \to K^{sep}$.

\begin{lemma}
\label{pione-lemma-normal-pione-quotient-inertia}
Soit $X$ un schéma intègre normal de corps des fonctions rationnelles $K$.
Avec les notations ci-dessus, les trois sous-groupes suivants de
$\text{Gal}(K^{sep}/K) = \pi_1(\Spec(K))$
sont égaux :
\begin{enumerate}
\item le noyau de la surjection
$\text{Gal}(K^{sep}/K) \longrightarrow \pi_1(X)$,
\item le plus petit sous-groupe distingué fermé qui contient $I_y$
pour tout $y \in X^{sep}$, et
\item le plus petit sous-groupe distingué fermé qui contient
$\text{Gal}(K^{sep}/K_x^{sh})$ pour tout $x \in  X$.
\end{enumerate}
\end{lemma}

\begin{proof}
L'équivalence de (2) et (3) découle du
lemme \ref{pione-lemma-identify-inertia},
qui nous dit que $I_y$ est conjugué à $\text{Gal}(K^{sep}/K_x^{sh})$
si $y$ est au-dessus de $x$. D'après le lemme \ref{pione-lemma-local-exact-sequence-normal},
on voit que $\text{Gal}(K^{sep}/K_x^{sh})$ s'envoie trivialement dans
$\pi_1(\Spec(\mathcal{O}_{X, x}))$ et que, par conséquent, le sous-groupe
$N \subset G = \text{Gal}(K^{sep}/K)$
de (2) et (3) est contenu dans le noyau de
$G \longrightarrow \pi_1(X)$.

\medskip\noindent
Pour démontrer l'autre inclusion, puisque $N$ est distingué, il suffit de montrer ceci :
étant donné $N \subset U \subset G$ avec $U$ ouvert normal,
la projection canonique $G \to G/U$ se factorise par $\pi_1(X)$.
Autrement dit, si $L/K$ est l'extension galoisienne correspondant
à $U$, alors il faut montrer que $X$ est non ramifié dans $L$
(section \ref{pione-section-normal}, en particulier
proposition \ref{pione-proposition-normal}).
Il suffit de le faire lorsque $X$ est affine (on procède ainsi
afin de pouvoir invoquer des résultats algébriques dans la suite de la démonstration).
Soit $Y \to X$ la normalisation de $X$ dans $L$.
L'inclusion $L \subset K^{sep}$ induit un morphisme
$\pi : X^{sep} \to Y$. Pour $y \in X^{sep}$,
le groupe d'inertie de $\pi(y)$ dans $\text{Gal}(L/K)$
est l'image de $I_y$ dans $\text{Gal}(L/K)$ ; cela découle
de Compléments d'algèbre, lemme
\ref{more-algebra-lemma-one-orbit-geometric-galois-compare}.
Puisque $N \subset U$, tous ces groupes d'inertie sont triviaux.
On conclut que $Y \to X$ est étale en appliquant le
lemme \ref{pione-lemma-inertia-invariants-etale}.
(Autre possibilité : on peut utiliser le lemme \ref{pione-lemma-local-exact-sequence-normal}
pour voir que le changement de base de $Y$ à $\Spec(\mathcal{O}_{X, x})$ est
étale pour tout $x \in X$, puis conclure à partir de là
avec un peu plus de travail.)
\end{proof}

\begin{example}
\label{pione-example-bigger-codim}
Soit $X$ un schéma noethérien intègre normal de corps des fonctions rationnelles $K$.
La pureté du lieu de branchement (voir plus bas) nous dit que si $X$ est régulier, alors
il suffit, dans le lemme \ref{pione-lemma-normal-pione-quotient-inertia},
de considérer les groupes d'inertie $I = \pi_1(\Spec(K_x^{sh}))$
pour les points $x$ de codimension $1$ dans $X$.
En général, cela ne suffit cependant pas. En effet, soit
$Y = \mathbf{A}_k^n = \Spec(k[t_1, \ldots, t_n])$,
où $k$ est un corps de caractéristique différente de $2$.
Soit $G = \{\pm 1\}$ le groupe d'ordre $2$ agissant sur $Y$
par multiplication sur les coordonnées. Posons
$$
X = \Spec(k[t_it_j, i, j \in \{1, \ldots, n\}])
$$
Le plongement $k[t_it_j] \subset k[t_1, \ldots, t_n]$
définit un morphisme de degré $2$, à savoir $Y \to X$, qui est non ramifié partout
sauf au-dessus de l'idéal maximal $\mathfrak m = (t_it_j)$,
qui est un point de codimension $n$ dans $X$.
\end{example}

\begin{lemma}
\label{pione-lemma-unramified}
Soit $X$ un schéma intègre normal de corps des fonctions rationnelles $K$.
Soit $L/K$ une extension finie. Soit $Y \to X$ la normalisation
de $X$ dans $L$. Les conditions suivantes sont équivalentes :
\begin{enumerate}
\item $X$ est non ramifié dans $L$ au sens de la section \ref{pione-section-normal},
\item $Y \to X$ est un morphisme non ramifié de schémas,
\item $Y \to X$ est un morphisme étale de schémas,
\item $Y \to X$ est un morphisme fini étale de schémas,
\item pour $x \in X$, la projection
$Y \times_X \Spec(\mathcal{O}_{X, x}) \to \Spec(\mathcal{O}_{X, x})$
est non ramifiée,
\item la même condition que (5), mais avec $\mathcal{O}_{X, x}^h$,
\item la même condition que (5), mais avec $\mathcal{O}_{X, x}^{sh}$,
\item pour $x \in X$, la fibre schématique $Y_x$
est étale sur $x$ de degré $\geq [L : K]$.
\end{enumerate}
Si $L/K$ est galoisienne de groupe de Galois $G$, ces conditions sont aussi
équivalentes à
\begin{enumerate}
\item[(9)] pour $y \in Y$, le groupe
$I_y = \{g \in G \mid g(y) = y\text{ et }
g \bmod \mathfrak m_y = \text{id}_{\kappa(y)}\}$ est trivial.
\end{enumerate}
\end{lemma}

\begin{proof}
L'équivalence de (1) et (2) est la définition de (1).
L'équivalence de (2), (3) et (4) est le lemme \ref{pione-lemma-unramified-in-L}.
Il est immédiat de démontrer que (4) $\Rightarrow$ (5),
(5) $\Rightarrow$ (6), (6) $\Rightarrow$ (7).

\medskip\noindent
Supposons (7). Observons que $\mathcal{O}_{X, x}^{sh}$ est un domaine local normal
(Compléments d'algèbre, lemme \ref{more-algebra-lemma-henselization-normal}).
Soit $L^{sh} = L \otimes_K K_x^{sh}$, où $K_x^{sh}$ est le corps des fractions
de $\mathcal{O}_{X, x}^{sh}$. Alors $L^{sh} = \prod_{i = 1, \ldots, n} L_i$,
avec $L_i/K_x^{sh}$ finie séparable. D'après
Algèbre, lemme \ref{algebra-lemma-integral-closure-commutes-smooth}
(et un passage à la limite que nous omettons),
on voit que $Y \times_X \Spec(\mathcal{O}_{X, x}^{sh})$
est la clôture intégrale de $\Spec(\mathcal{O}_{X, x}^{sh})$ dans $L^{sh}$.
Par conséquent, d'après le lemme \ref{pione-lemma-unramified-in-L} (appliqué aux facteurs
$L_i$ de $L^{sh}$), on voit que
$Y \times_X \Spec(\mathcal{O}_{X, x}^{sh}) \to \Spec(\mathcal{O}_{X, x}^{sh})$
est fini étale. En considérant le point générique, on voit que
le degré est égal à $[L : K]$ et on voit donc que (8) est vraie.

\medskip\noindent
Supposons (8). Supposons que $x \in X$ et que la fibre schématique $Y_x$
est étale sur $x$ de degré $\geq [L : K]$. Observons que cela signifie
que $Y$ possède $\geq [L : K]$ points géométriques au-dessus de $x$.
Nous allons montrer que $Y \to X$ est fini étale au-dessus d'un voisinage de $x$.
Cela démontrera que (1) est vraie.
Pour le démontrer, on peut supposer $X = \Spec(R)$, que le point $x$ correspond à
l'idéal premier $\mathfrak p \subset R$ et que $Y = \Spec(S)$. On applique
Compléments sur les morphismes,
lemme \ref{more-morphisms-lemma-etale-makes-integral-split}, et on trouve un
voisinage étale $(U, u) \to (X, x)$ tel que
$Y \times_X U = V_1 \amalg \ldots \amalg V_m$, où $V_i$
possède un unique point $v_i$ au-dessus de $u$ avec $\kappa(v_i)/\kappa(u)$
purement inséparable. En rétrécissant $U$ si nécessaire, on peut supposer que $U$ est
un schéma intègre normal de point générique $\xi$ (utiliser
Descente, lemmes \ref{descent-lemma-locally-finite-nr-irred-local-fppf} et
\ref{descent-lemma-normal-local-smooth}, et
Propriétés, lemme \ref{properties-lemma-normal-locally-finite-nr-irreducibles}).
D'après notre remarque sur les points géométriques, on voit que $m \geq [L : K]$.
D'autre part, d'après Compléments sur les morphismes, lemme
\ref{more-morphisms-lemma-normalization-smooth-localization},
on voit que $\coprod V_i \to U$ est la normalisation de $U$ dans
$\Spec(L) \times_X U$. Comme $K \subset \kappa(\xi)$ est finie séparable,
on peut écrire $\Spec(L) \times_X U = \Spec(\prod_{i = 1, \ldots, n} L_i)$,
avec $L_i/\kappa(\xi)$ finie et $[L : K] = \sum [L_i : \kappa(\xi)]$.
Puisque $V_j$ est non vide pour tout $j$ et que $m \geq [L : K]$,
on conclut que $m = n$ et que $[L_i : \kappa(\xi)] = 1$
pour tout $i$. Alors $V_j \to U$ est en particulier un isomorphisme
étale, et donc $Y \times_X U \to U$ est étale. D'après
Descente, lemme \ref{descent-lemma-descending-property-etale},
on conclut que $Y \to X$ est étale au-dessus de
l'image de $U \to X$ (un voisinage ouvert de $x$).

\medskip\noindent
Supposons $L/K$ galoisienne et (9) vérifiée. Alors $Y \to X$ est étale
d'après le lemme \ref{pione-lemma-inertial-invariants-unramified}.
On omet la démonstration de ce que (1) implique (9).
\end{proof}

\noindent
Dans le cas d'extensions galoisiennes infinies d'anneaux de valuation discrète,
on peut en dire un tout petit peu plus. Pour cela, introduisons la notation suivante.
Un sous-ensemble $S \subset \mathbf{N}$ d'entiers est {\it filtrant pour la divisibilité}
si $1 \in S$ et si, pour $n, m \in S$, il existe $k \in S$ tel que
$n | k$ et $m | k$. Définissons sur $S$ un ordre partiel par la règle
$n \geq_S m$ si et seulement si $m | n$. Étant donné un corps $\kappa$, on obtient
un système projectif de groupes finis $\{\mu_n(\kappa)\}_{n \in S}$
dont les applications de transition sont
$$
\mu_n(\kappa) \longrightarrow \mu_m(\kappa),\quad
\zeta \longmapsto \zeta^{n/m}
$$
pour $n \geq_S m$. On peut alors former le groupe profini
$$
\lim_{n \in S} \mu_n(\kappa)
$$
Observons que la catégorie d'indices de cette limite est cofiltrante (car $S$ est filtrant).
La construction est fonctorielle en $\kappa$. En particulier,
$\text{Aut}(\kappa)$ agit sur ce groupe profini.
Par exemple, si $S = \{1, n\}$, on obtient alors $\mu_n(\kappa)$.
Si $S = \{1, \ell, \ell^2, \ell^3, \ldots\}$ pour un nombre premier
$\ell$ différent de la caractéristique de $\kappa$, cela produit
$\lim_n \mu_{\ell^n}(\kappa)$,
que l'on appelle parfois le module de Tate $\ell$-adique du groupe
multiplicatif de $\kappa$ (comparer avec
Compléments d'algèbre, exemple
\ref{more-algebra-example-spectral-sequence-principal}).

\begin{lemma}
\label{pione-lemma-structure-decomposition}
Soit $A$ un anneau de valuation discrète de corps des fractions $K$.
Soit $L/K$ une extension galoisienne (éventuellement infinie).
Soit $B$ la clôture intégrale de $A$ dans $L$.
Soit $\mathfrak m$ un idéal maximal de $B$ et posons $\kappa_A = A/(\mathfrak m \cap A)$.
Soient $G = \text{Gal}(L/K)$,
$D = \{\sigma \in G \mid \sigma(\mathfrak m) = \mathfrak m\}$, et
$I = \{\sigma \in D \mid \sigma \bmod \mathfrak m =
\text{id}_{\kappa(\mathfrak m)}\}$.
Le groupe de décomposition $D$ s'insère dans une suite exacte canonique
$$
1 \to I \to D \to \text{Aut}(\kappa(\mathfrak m)/\kappa_A) \to 1
$$
Le groupe d'inertie $I$ s'insère dans une suite exacte canonique
$$
1 \to P \to I \to I_t \to 1
$$
telle que
\begin{enumerate}
\item $P$ est un sous-groupe distingué de $D$,
\item $P$ est un pro-$p$-groupe si la caractéristique de
$\kappa_A$ est $p > 1$, et $P = \{1\}$ si la caractéristique de $\kappa_A$
est nulle (avec la convention $p = 1$ dans ce dernier cas),
\item il existe un $S \subset \mathbf{N}$ filtrant pour la divisibilité
tel que $\kappa(\mathfrak m)$ contient une racine $n$-ième primitive de l'unité
pour tout $n \in S$ (les éléments de $S$ sont premiers à $p$),
\item il existe une application surjective canonique
$$
\theta_{can} : I \to \lim_{n \in S} \mu_n(\kappa(\mathfrak m))
$$
dont le noyau est $P$, qui vérifie
$\theta_{can}(\tau \sigma \tau^{-1}) = \tau(\theta_{can}(\sigma))$
pour $\tau \in D$, $\sigma \in I$, et qui induit un isomorphisme
$I_t \to \lim_{n \in S} \mu_n(\kappa(\mathfrak m))$.
\end{enumerate}
\end{lemma}

\begin{proof}
Il s'agit essentiellement d'une reformulation des résultats sur les extensions galoisiennes finies
démontrés dans Compléments d'algèbre, section \ref{more-algebra-section-ramification}.
La surjectivité de l'application $D \to \text{Aut}(\kappa(\mathfrak m)/\kappa_A)$ est le
lemme \ref{more-algebra-lemma-one-orbit-geometric-galois} de Compléments d'algèbre.
Cela donne la première suite exacte.

\medskip\noindent
Pour construire la seconde suite exacte courte, soit $\Lambda$ l'ensemble
des sous-extensions galoisiennes finies, c'est-à-dire que $\lambda \in \Lambda$ correspond
à $L/L_\lambda/K$. Posons $G_\lambda = \text{Gal}(L_\lambda/K)$.
Rappelons que les $G_\lambda$ forment un système projectif de groupes finis à applications
de transition surjectives et que $G = \lim_{\lambda \in \Lambda} G_\lambda$, voir
Corps, lemme \ref{fields-lemma-infinite-galois-limit}.
On note $B_\lambda$ la clôture intégrale de $A$ dans $L_\lambda$.
On pose ensuite $\mathfrak m_\lambda = \mathfrak m \cap B_\lambda$
et on note $P_\lambda, I_\lambda, D_\lambda$ les groupes
d'inertie sauvage, d'inertie et de décomposition de
$\mathfrak m_\lambda$, voir Compléments d'algèbre, lemme
\ref{more-algebra-lemma-galois-inertia}.
Pour $\lambda \geq \lambda'$, la restriction définit
un diagramme commutatif
$$
\xymatrix{
P_\lambda \ar[d] \ar[r] &
I_\lambda \ar[d] \ar[r] &
D_\lambda \ar[d] \ar[r] &
G_\lambda \ar[d] \\
P_{\lambda'} \ar[r] &
I_{\lambda'} \ar[r] &
D_{\lambda'} \ar[r] &
G_{\lambda'}
}
$$
dont les flèches verticales sont surjectives, voir
Compléments d'algèbre, lemme \ref{more-algebra-lemma-compare-inertia}.

\medskip\noindent
Il résulte immédiatement des définitions
que $I = \lim I_\lambda$ et $D = \lim D_\lambda$
sous l'isomorphisme $G = \lim G_\lambda$ ci-dessus.
Puisque $L = \colim L_\lambda$, on a $B = \colim B_\lambda$
et $\kappa(\mathfrak m) = \colim \kappa(\mathfrak m_\lambda)$.
Puisque les applications de transition du système $D_\lambda$
sont compatibles avec les applications
$D_\lambda \to \text{Aut}(\kappa(\mathfrak m_\lambda)/\kappa)$
(voir Compléments d'algèbre, lemme \ref{more-algebra-lemma-compare-inertia}),
on voit que l'application $D \to \text{Aut}(\kappa(\mathfrak m)/\kappa)$
est la limite des applications
$D_\lambda \to \text{Aut}(\kappa(\mathfrak m_\lambda)/\kappa)$.

\medskip\noindent
Il existe des applications canoniques
$$
\theta_{\lambda, can} :
I_\lambda
\longrightarrow
\mu_{n_\lambda}(\kappa(\mathfrak m_\lambda))
$$
où $n_\lambda = |I_\lambda|/|P_\lambda|$, où
$\mu_{n_\lambda}(\kappa(\mathfrak m_\lambda))$ est
d'ordre $n_\lambda$, telles que
$\theta_{\lambda, can}(\tau \sigma \tau^{-1}) =
\tau(\theta_{\lambda, can}(\sigma))$ pour
$\tau \in D_\lambda$ et $\sigma \in I_\lambda$, et telles que
l'on obtient des diagrammes commutatifs
$$
\xymatrix{
I_\lambda \ar[r]_-{\theta_{\lambda, can}} \ar[d] &
\mu_{n_\lambda}(\kappa(\mathfrak m_\lambda))
\ar[d]^{(-)^{n_\lambda/n_{\lambda'}}} \\
I_{\lambda'} \ar[r]^-{\theta_{\lambda', can}} &
\mu_{n_{\lambda'}}(\kappa(\mathfrak m_{\lambda'}))
}
$$
voir
Compléments d'algèbre, remarque \ref{more-algebra-remark-canonical-inertia-character}.

\medskip\noindent
Soit $S \subset \mathbf{N}$ l'ensemble des entiers $n_\lambda$.
Puisque $\Lambda$ est filtrant, on voit que $S$ est filtrant pour la divisibilité.
À l'aide des diagrammes commutatifs affichés ci-dessus, on peut prendre les limites des
applications $\theta_{\lambda, can}$ afin d'obtenir
$$
\theta_{can} : I \to \lim_{n \in S} \mu_n(\kappa(\mathfrak m)).
$$
Cette application est continue (on omet ce petit détail). Puisque les applications de transition
du système des $I_\lambda$ sont surjectives
et que $\Lambda$ est filtrant, les projections $I \to I_\lambda$
sont surjectives. Pour tout $\lambda$, le diagramme
$$
\xymatrix{
I \ar[d] \ar[r]_-{\theta_{can}} &
\lim_{n \in S} \mu_n(\kappa(\mathfrak m)) \ar[d] \\
I_{\lambda} \ar[r]^-{\theta_{\lambda, can}} &
\mu_{n_\lambda}(\kappa(\mathfrak m_\lambda))
}
$$
est commutatif. L'image de $\theta_{can}$ se projette donc sur le groupe fini
$\mu_{n_\lambda}(\kappa(\mathfrak m)) =
\mu_{n_\lambda}(\kappa(\mathfrak m_\lambda))$ d'ordre $n_\lambda$
(voir ci-dessus). Il s'ensuit que l'image de $\theta_{can}$ est dense.
D'autre part, $\theta_{can}$ est continue et la
source est un groupe profini. Ainsi $\theta_{can}$ est surjective
d'après un argument topologique.

\medskip\noindent
La propriété $\theta_{can}(\tau \sigma \tau^{-1}) = \tau(\theta_{can}(\sigma))$
pour $\tau \in D$, $\sigma \in I$ découle des propriétés correspondantes
des applications $\theta_{\lambda, can}$ et de la compatibilité de l'application
$D \to \text{Aut}(\kappa(\mathfrak m))$ avec les applications
$D_\lambda \to \text{Aut}(\kappa(\mathfrak m_\lambda))$.
En posant $P = \Ker(\theta_{can})$, cela implique
que $P$ est un sous-groupe distingué de $D$. En posant $I_t = I/P$,
on obtient l'isomorphisme $I_t \to \lim_{n \in S} \mu_n(\kappa(\mathfrak m))$
par la surjectivité de $\theta_{can}$.

\medskip\noindent
Pour achever la démonstration, montrons que $P = \lim P_\lambda$, ce qui prouve
que $P$ est un pro-$p$-groupe. Rappelons que le groupe d'inertie modérée
$I_{\lambda, t} = I_\lambda/P_\lambda$ est d'ordre $n_\lambda$.
Puisque les applications de transition $P_\lambda \to P_{\lambda'}$ sont surjectives
et que $\Lambda$ est filtrant, on obtient une suite exacte courte
$$
1 \to \lim P_\lambda \to I \to \lim I_{\lambda, t} \to 1
$$
(on omet les détails). Puisque, pour tout $\lambda$, l'application $\theta_{\lambda, can}$
induit un isomorphisme
$I_{\lambda, t} \cong \mu_{n_\lambda}(\kappa(\mathfrak m))$,
le résultat souhaité en découle.
\end{proof}

\begin{lemma}
\label{pione-lemma-structure-decomposition-separable-closure}
Soit $A$ un anneau de valuation discrète de corps des fractions $K$.
Soit $K^{sep}$ une clôture séparable de $K$.
Soit $A^{sep}$ la clôture intégrale de $A$ dans $K^{sep}$.
Soit $\mathfrak m^{sep}$ un idéal maximal de $A^{sep}$.
Soit $\mathfrak m = \mathfrak m^{sep} \cap A$, soient
$\kappa = A/\mathfrak m$, et
$\overline{\kappa} = A^{sep}/\mathfrak m^{sep}$.
Alors $\overline{\kappa}$ est une clôture algébrique de $\kappa$.
Soient $G = \text{Gal}(K^{sep}/K)$,
$D = \{\sigma \in G \mid \sigma(\mathfrak m^{sep}) = \mathfrak m^{sep}\}$, et
$I = \{\sigma \in D \mid \sigma \bmod \mathfrak m^{sep} =
\text{id}_{\kappa(\mathfrak m^{sep})}\}$.
Le groupe de décomposition $D$ s'insère dans une suite exacte canonique
$$
1 \to I \to D \to \text{Gal}(\kappa^{sep}/\kappa) \to 1
$$
où $\kappa^{sep} \subset \overline{\kappa}$ est la clôture
séparable de $\kappa$.
Le groupe d'inertie $I$ s'insère dans une suite exacte canonique
$$
1 \to P \to I \to I_t \to 1
$$
telle que
\begin{enumerate}
\item $P$ est un sous-groupe distingué de $D$,
\item $P$ est un pro-$p$-groupe si la caractéristique de
$\kappa$ est $p > 1$, et $P = \{1\}$ si la caractéristique de $\kappa$
est nulle (avec la convention $p = 1$ dans ce dernier cas),
\item il existe une application surjective canonique
$$
\theta_{can} : I \to \lim_{n\text{ premier à }p} \mu_n(\kappa^{sep})
$$
dont le noyau est $P$, qui vérifie
$\theta_{can}(\tau \sigma \tau^{-1}) = \tau(\theta_{can}(\sigma))$
pour $\tau \in D$, $\sigma \in I$, et qui induit un isomorphisme
$I_t \to \lim_{n\text{ premier à }p} \mu_n(\kappa^{sep})$.
\end{enumerate}
\end{lemma}

\begin{proof}
Le corps $\overline{\kappa}$ est la clôture algébrique de $\kappa$ d'après le
lemme \ref{pione-lemma-get-algebraic-closure}.
La plupart des assertions découlent immédiatement des parties correspondantes
du lemme \ref{pione-lemma-structure-decomposition}. Par exemple, puisque
$\text{Aut}(\overline{\kappa}/\kappa) = \text{Gal}(\kappa^{sep}/\kappa)$,
on obtient la première suite.
Alors la seule autre assertion qui requiert une démonstration est le fait que,
avec $S$ comme dans le lemme \ref{pione-lemma-structure-decomposition}, la
limite $\lim_{n \in S} \mu_n(\overline{\kappa})$ est égale à
$\lim_{n\text{ premier à }p} \mu_n(\kappa^{sep})$. Pour le voir,
il suffit de montrer que tout entier $n$ premier à $p$
divise un élément de $S$.
Soit $\pi \in A$ une uniformisante et considérons le corps de décomposition
$L$ du polynôme $X^n - \pi$. Puisque ce polynôme
est séparable, on voit que $L$ est une extension galoisienne finie de $K$.
Choisissons un plongement $L \to K^{sep}$.
Observons que si $B$ est la clôture intégrale de $A$ dans $L$,
alors l'indice de ramification de $A \to B_{\mathfrak m^{sep} \cap B}$
est divisible par $n$ (car $\pi$ admet une racine $n$-ième dans $B$ ; en fait,
l'indice de ramification est égal à $n$, mais nous n'en avons pas besoin).
Il résulte alors de la construction de l'ensemble $S$ dans la démonstration du
lemme \ref{pione-lemma-structure-decomposition}
que $n$ divise un élément de $S$.
\end{proof}








\section{Groupes fondamentaux géométriques et arithmétiques}
\label{pione-section-galois-action}

\noindent
Dans cette section, nous étudions ce qui se passe lorsque l'on compare le
groupe fondamental d'un schéma $X$ sur un corps $k$ au
groupe fondamental de $X_{\overline{k}}$, où $\overline{k}$
est la clôture algébrique de $k$.

\begin{lemma}
\label{pione-lemma-limit}
Soit $I$ un ensemble ordonné filtrant. Soit $X_i$ un
système projectif de schémas quasi-compacts et quasi-séparés
indexé par $I$, dont les morphismes de transition sont affines.
Soit $X = \lim X_i$ comme dans Limites, section \ref{limits-section-limits}.
Il existe alors une équivalence de catégories
$$
\colim \textit{F\'Et}_{X_i} = \textit{F\'Et}_X
$$
Si $X_i$ est connexe pour tout $i$ suffisamment grand et si $\overline{x}$
est un point géométrique de $X$, alors
$$
\pi_1(X, \overline{x}) = \lim \pi_1(X_i, \overline{x})
$$
\end{lemma}

\begin{proof}
L'équivalence de catégories résulte de Limites, lemmes
\ref{limits-lemma-descend-finite-presentation},
\ref{limits-lemma-descend-finite-finite-presentation}, et
\ref{limits-lemma-descend-etale}.
La seconde assertion découle formellement de l'assertion sur les
catégories.
\end{proof}

\begin{lemma}
\label{pione-lemma-perfection}
Soit $k$ un corps et notons $k^{perf}$ sa clôture parfaite. Soit $X$ un schéma connexe
sur $k$. Alors $X_{k^{perf}}$ est connexe et
$\pi_1(X_{k^{perf}}) \to \pi_1(X)$ est un isomorphisme.
\end{lemma}

\begin{proof}
C'est un cas particulier de l'invariance topologique du groupe fondamental.
Voir la proposition \ref{pione-proposition-universal-homeomorphism}.
Pour voir que $\Spec(k^{perf}) \to \Spec(k)$ est un
homéomorphisme universel, on peut utiliser
Algèbre, lemme \ref{algebra-lemma-radicial-integral-bijective}.
\end{proof}

\begin{lemma}
\label{pione-lemma-ses-field}
Soit $k$ un corps dont $\overline{k}$ est une clôture algébrique.
Soit $X$ un schéma quasi-compact et quasi-séparé sur $k$.
Si le changement de base $X_{\overline{k}}$ est connexe, alors
il existe une suite exacte courte
$$
1 \to \pi_1(X_{\overline{k}}) \to \pi_1(X) \to \pi_1(\Spec(k)) \to 1
$$
de groupes topologiques profinis.
\end{lemma}

\begin{proof}
Les objets connexes de $\textit{F\'Et}_{\Spec(k)}$ sont de la forme
$\Spec(k') \to \Spec(k)$, où $k'/k$ est une extension finie séparable.
Alors $X_{\Spec(k')}$ est connexe, puisque le morphisme
$X_{\overline{k}} \to X_{\Spec(k')}$ est surjectif et
$X_{\overline{k}}$ est connexe par hypothèse. Ainsi
$\pi_1(X) \to \pi_1(\Spec(k))$ est surjectif d'après le
lemme \ref{pione-lemma-functoriality-galois-surjective}.

\medskip\noindent
Avant de poursuivre, remarquons que l'on peut supposer que $k$ est parfait.
En effet, on a $\pi_1(X_{k^{perf}}) = \pi_1(X)$ et
$\pi_1(\Spec(k^{perf})) = \pi_1(\Spec(k))$ d'après le lemme \ref{pione-lemma-perfection}.

\medskip\noindent
Il est clair que le composé des foncteurs
$\textit{F\'Et}_{\Spec(k)} \to \textit{F\'Et}_X \to
\textit{F\'Et}_{X_{\overline{k}}}$ envoie les objets sur des réunions disjointes
de copies de $X_{\Spec(\overline{k})}$. Par conséquent, le composé
$\pi_1(X_{\overline{k}}) \to \pi_1(X) \to \pi_1(\Spec(k))$
est l'homomorphisme trivial d'après le lemme \ref{pione-lemma-composition-trivial}.

\medskip\noindent
Soit $U \to X$ un morphisme fini \'etale, avec $U$ connexe.
Observons que $U \times_X X_{\overline{k}} = U_{\overline{k}}$.
Supposons que $U_{\overline{k}} \to X_{\overline{k}}$
admette une section $s : X_{\overline{k}} \to U_{\overline{k}}$.
Alors $s(X_{\overline{k}})$ est une composante connexe ouverte de
$U_{\overline{k}}$. Pour $\sigma \in \text{Gal}(\overline{k}/k)$,
notons $s^\sigma$ le changement de base de $s$ par $\Spec(\sigma)$.
Puisque $U_{\overline{k}} \to X_{\overline{k}}$ est fini \'etale,
il n'a qu'un nombre fini de sections. Ainsi
$$
\overline{T} = \bigcup_{\sigma \in \text{Gal}(\overline{k}/k)} s^\sigma(X_{\overline{k}})
$$
est une réunion finie et l'on voit que $\overline{T}$ est un
sous-ensemble ouvert et fermé stable par $\text{Gal}(\overline{k}/k)$.
D'après Variétés, lemme \ref{varieties-lemma-closed-fixed-by-Galois},
on voit que $\overline{T}$ est l'image réciproque d'un sous-ensemble fermé
$T \subset U$. Puisque $U_{\overline{k}} \to U$ est ouvert
(Morphismes, lemme \ref{morphisms-lemma-scheme-over-field-universally-open}),
on conclut que $T$ est également ouvert. Comme $U$ est connexe, on
voit que $T = U$. Par conséquent, $U_{\overline{k}}$ est une réunion disjointe (finie)
de copies de $X_{\overline{k}}$. D'après le
lemme \ref{pione-lemma-functoriality-galois-normal}, on conclut que l'image de
$\pi_1(X_{\overline{k}}) \to \pi_1(X)$ est un sous-groupe distingué.

\medskip\noindent
Soit $V \to X_{\overline{k}}$ un revêtement fini \'etale. Rappelons que
$\overline{k}$ est la réunion des extensions finies séparables de $k$.
D'après le lemme \ref{pione-lemma-limit}, il existe une extension finie séparable $k'/k$
et un morphisme fini \'etale $U \to X_{k'}$ tels que
$V = X_{\overline{k}} \times_{X_{k'}} U =
U \times_{\Spec(k')} \Spec(\overline{k})$.
Alors le composé $U \to X_{k'} \to X$ est fini \'etale
et $U \times_{\Spec(k)} \Spec(\overline{k})$
contient $V = U \times_{\Spec(k')} \Spec(\overline{k})$
comme sous-schéma ouvert et fermé. (En effet, $\Spec(\overline{k})$
est un sous-schéma ouvert et fermé de
$\Spec(k') \times_{\Spec(k)} \Spec(\overline{k})$ via
l'application de multiplication $k' \otimes_k \overline{k} \to \overline{k}$.) D'après le
lemme \ref{pione-lemma-functoriality-galois-injective},
on conclut que $\pi_1(X_{\overline{k}}) \to \pi_1(X)$ est injectif.

\medskip\noindent
Enfin, nous devons montrer que, pour tout morphisme fini \'etale
$U \to X$ tel que $U_{\overline{k}}$ soit une réunion disjointe
de copies de $X_{\overline{k}}$, il existe un morphisme fini \'etale
$V \to \Spec(k)$ et une surjection $V \times_{\Spec(k)} X \to U$.
Voir le lemme \ref{pione-lemma-functoriality-galois-ses}.
En raisonnant comme ci-dessus à l'aide du lemme \ref{pione-lemma-limit},
il existe une extension finie séparable $k'/k$
pour laquelle il existe un isomorphisme
$U_{k'} \cong \coprod_{i = 1, \ldots, n} X_{k'}$.
Ainsi, en posant $V = \coprod_{i = 1, \ldots, n} \Spec(k')$,
on conclut.
\end{proof}






\section{Suite exacte d'homotopie}
\label{pione-section-homotopy-exact-sequence}

\noindent
Dans cette section, nous examinons le résultat suivant.
Soit $f : X \to S$ un morphisme propre et plat de
présentation finie dont les
fibres géométriques sont connexes et réduites.
Supposons $S$ connexe et soit $\overline{s}$
un point géométrique de $S$. Il existe alors une suite
exacte
$$
\pi_1(X_{\overline{s}}) \to \pi_1(X) \to \pi_1(S) \to 1
$$
de groupes fondamentaux. Voir la
proposition \ref{pione-proposition-first-homotopy-sequence}.

\begin{lemma}
\label{pione-lemma-stein-factorization-etale}
\begin{reference}
\cite[Exposé X, proposition 1.2, p. 262]{SGA1}.
\end{reference}
Soit $f : X \to S$ un morphisme propre de schémas.
Soit $X \to S' \to S$ la factorisation de Stein de $f$ ; voir
Compléments sur les morphismes, théorème
\ref{more-morphisms-theorem-stein-factorization-general}.
Si $f$ est de présentation finie et plat, à fibres géométriquement
réduites, alors $S' \to S$ est fini \'etale.
\end{lemma}

\begin{proof}
Cela résulte de Catégories dérivées des schémas,
lemme \ref{perfect-lemma-proper-flat-geom-red}
et des informations contenues dans
Compléments sur les morphismes, théorème
\ref{more-morphisms-theorem-stein-factorization-general}.
\end{proof}

\begin{proposition}
\label{pione-proposition-first-homotopy-sequence}
Soit $f : X \to S$ un morphisme propre et plat de présentation finie dont les
fibres géométriques sont connexes et réduites. Supposons $S$ connexe et
soit $\overline{s}$ un point géométrique de $S$. Alors il existe une suite
exacte
$$
\pi_1(X_{\overline{s}}) \to \pi_1(X) \to \pi_1(S) \to 1
$$
de groupes fondamentaux.
\end{proposition}

\begin{proof}
Soit $Y \to X$ un morphisme fini \'etale. Considérons la factorisation de Stein
$$
\xymatrix{
Y \ar[d] \ar[r] & X \ar[d] \\
T \ar[r] & S
}
$$
de $Y \to S$. D'après le lemme \ref{pione-lemma-stein-factorization-etale},
le morphisme $T \to S$ est fini \'etale. On obtient ainsi
un foncteur $\textit{F\'Et}_X \to \textit{F\'Et}_S$.
Pour tout morphisme fini \'etale $U \to S$, se donner un morphisme
$Y \to U \times_S X$ au-dessus de $X$ revient à se donner un morphisme
$Y \to U$ au-dessus de $S$ ; un tel morphisme se factorise de manière unique par
la factorisation de Stein, c'est-à-dire correspond à un unique
morphisme $T \to U$
(par la construction de la factorisation de Stein comme
normalisation relative dans Compléments sur les morphismes, lemme
\ref{more-morphisms-lemma-stein-universally-closed}
et par la factorisation donnée par
Morphismes, lemme \ref{morphisms-lemma-characterize-normalization}).
Ainsi, on voit que les foncteurs
$\textit{F\'Et}_X \to \textit{F\'Et}_S$ et
$\textit{F\'Et}_S \to \textit{F\'Et}_X$ sont adjoints l'un de l'autre.
Remarquons que la factorisation de Stein de $U \times_S X \to S$ est
$U$, car les fibres de $U \times_S X \to U$ sont géométriquement connexes.

\medskip\noindent
D'après la discussion ci-dessus et
Catégories, lemme \ref{categories-lemma-adjoint-fully-faithful},
on conclut que
$\textit{F\'Et}_S \to \textit{F\'Et}_X$
est pleinement fidèle, c'est-à-dire que $\pi_1(X) \to \pi_1(S)$ est surjectif
(lemme \ref{pione-lemma-functoriality-galois-surjective}).

\medskip\noindent
Il est immédiat que le composé
$\textit{F\'Et}_S \to \textit{F\'Et}_X \to \textit{F\'Et}_{X_{\overline{s}}}$
envoie tout $U$ sur une réunion disjointe de copies de $X_{\overline{s}}$.
Ainsi la composée $\pi_1(X_{\overline{s}}) \to \pi_1(X) \to \pi_1(S)$ est triviale
d'après le lemme \ref{pione-lemma-composition-trivial}.

\medskip\noindent
Soit $Y \to X$ un morphisme fini \'etale, où $Y$ est connexe, tel que
$Y \times_X X_{\overline{s}}$ contienne une composante connexe $Z$
isomorphe à $X_{\overline{s}}$. Considérons la factorisation de Stein $T$
ci-dessus. Soit $\overline{t} \in T_{\overline{s}}$ le point correspondant
à la fibre $Z$. Observons que $T$ est connexe (en tant qu'image d'un schéma
connexe) et que, par la surjectivité ci-dessus, $T \times_S X$ est connexe.
Considérons maintenant la factorisation
$$
\pi : Y \longrightarrow T \times_S X
$$
Soit $\overline{x} \in X_{\overline{s}}$ un point fermé quelconque. Notons que
$\kappa(\overline{t}) = \kappa(\overline{s}) = \kappa(\overline{x})$
est un corps algébriquement clos.
Alors la fibre de $\pi$ au-dessus de $(\overline{t}, \overline{x})$ est constituée
d'un unique point, à savoir l'unique point $\overline{z} \in Z$
correspondant à $\overline{x} \in X_{\overline{s}}$ par
l'isomorphisme $Z \to X_{\overline{s}}$. On conclut que le morphisme fini
\'etale $\pi$ est de degré $1$ dans un voisinage de
$(\overline{t}, \overline{x})$. Puisque $T \times_S X$ est connexe,
ce morphisme est partout de degré $1$, et l'on trouve que $Y \cong T \times_S X$.
Ainsi $Y \times_X X_{\overline{s}}$ se scinde complètement.
En combinant tout ce qui précède, on voit que les
lemmes \ref{pione-lemma-functoriality-galois-ses} et
\ref{pione-lemma-functoriality-galois-normal}
s'appliquent tous deux, ce qui achève la démonstration.
\end{proof}




\section{Homomorphismes de spécialisation}
\label{pione-section-specialization-map}

\noindent
Dans cette section, nous construisons des homomorphismes de spécialisation.
Soit $f : X \to S$ un morphisme propre de schémas
à fibres géométriquement connexes.
Soit $s' \leadsto s$ une spécialisation de points de $S$.
Soient $\overline{s}$ et $\overline{s}'$ des points géométriques
au-dessus de $s$ et $s'$. Il existe alors un homomorphisme de spécialisation
$$
sp : \pi_1(X_{\overline{s}'}) \longrightarrow \pi_1(X_{\overline{s}})
$$
La construction de cet homomorphisme est la suivante. Soit $A$ le
hensélisé strict de $\mathcal{O}_{S, s}$ relativement à
$\kappa(s) \subset \kappa(s)^{sep} \subset \kappa(\overline{s})$, voir
Algèbre, définition \ref{algebra-definition-henselization}.
Puisque $s' \leadsto s$, le point $s'$ correspond à un point de
$\Spec(\mathcal{O}_{S, s})$ et il existe donc au moins un point
(et éventuellement plusieurs points)
de $\Spec(A)$ au-dessus de $s'$ dont le corps résiduel est une extension
algébrique séparable de $\kappa(s')$.
Puisque $\kappa(\overline{s}')$ est algébriquement clos, on peut choisir
un morphisme $\varphi : \overline{s}' \to \Spec(A)$ donnant lieu à un diagramme
commutatif
$$
\xymatrix{
\overline{s}' \ar[r]_-\varphi \ar[rd] &
\Spec(A) \ar[d] &
\overline{s} \ar[l] \ar[ld] \\
& S
}
$$
L'homomorphisme de spécialisation est la composée
$$
\pi_1(X_{\overline{s}'}) \longrightarrow
\pi_1(X_A) =
\pi_1(X_{\kappa(s)^{sep}}) =
\pi_1(X_{\overline{s}})
$$
où la première égalité est donnée par le
lemme \ref{pione-lemma-finite-etale-on-proper-over-henselian}
et la seconde découle des
lemmes \ref{pione-lemma-perfection} et
\ref{pione-lemma-finite-etale-invariant-over-proper}.
Par construction, l'homomorphisme de spécialisation s'insère dans un diagramme
commutatif
$$
\xymatrix{
\pi_1(X_{\overline{s}'}) \ar[rr]_{sp} \ar[rd] & &
\pi_1(X_{\overline{s}}) \ar[ld] \\
& \pi_1(X)
}
$$
pourvu que $X$ soit connexe. L'homomorphisme de spécialisation dépend du
choix de $\varphi : \overline{s}' \to \Spec(A)$ ci-dessus et nous
écrirons $sp_\varphi$ si nous voulons préciser ce choix.

\begin{lemma}
\label{pione-lemma-specialization-map-base-change}
Considérons un diagramme commutatif
$$
\xymatrix{
Y \ar[d]_g \ar[r] & X \ar[d]^f \\
T \ar[r] & S
}
$$
de schémas où $f$ et $g$ sont des morphismes propres à fibres géométriquement connexes.
Soit $t' \leadsto t$ une spécialisation de points de $T$
et considérons un homomorphisme de spécialisation
$sp : \pi_1(Y_{\overline{t}'}) \to \pi_1(Y_{\overline{t}})$ comme ci-dessus.
Il existe alors un diagramme commutatif
$$
\xymatrix{
\pi_1(Y_{\overline{t}'}) \ar[r]_{sp} \ar[d] & \pi_1(Y_{\overline{t}}) \ar[d] \\
\pi_1(X_{\overline{s}'}) \ar[r]^{sp} & \pi_1(X_{\overline{s}})
}
$$
d'homomorphismes de spécialisation, où $\overline{s}$ et $\overline{s}'$
sont les images de $\overline{t}$ et $\overline{t}'$.
\end{lemma}

\begin{proof}
Soit $B$ le hensélisé strict de $\mathcal{O}_{T, t}$ relativement à
$\kappa(t) \subset \kappa(t)^{sep} \subset \kappa(\overline{t})$.
Choisissons $\psi : \overline{t}' \to \Spec(B)$ relevant $\overline{t}' \to T$
comme dans la construction de l'homomorphisme de spécialisation.
Notons $s$ et $s'$ les images de $t$ et $t'$ dans $S$.
Soit $A$ le hensélisé strict de $\mathcal{O}_{S, s}$
relativement à
$\kappa(s) \subset \kappa(s)^{sep} \subset \kappa(\overline{s})$.
Puisque $\kappa(\overline{s}) = \kappa(\overline{t})$,
la fonctorialité de l'hensélisation stricte
(Algèbre, lemme \ref{algebra-lemma-strictly-henselian-functorial})
nous donne un morphisme d'anneaux $A \to B$ s'insérant dans le diagramme commutatif
$$
\xymatrix{
\overline{t}' \ar[r]_-\psi \ar[d] & \Spec(B) \ar[d] \ar[r] & T \ar[d] \\
\overline{s}' \ar[r]^-\varphi & \Spec(A) \ar[r] & S
}
$$
Ici, le morphisme $\varphi : \overline{s}' \to \Spec(A)$ est simplement pris
égal à la composée $\overline{t}' \to \Spec(B) \to \Spec(A)$.
En effectuant le changement de base, on obtient un diagramme commutatif
$$
\xymatrix{
Y_{\overline{t}'} \ar[r] \ar[d] & Y_B \ar[d] \\
X_{\overline{s}'} \ar[r] & X_A
}
$$
et, d'après la construction de l'homomorphisme de spécialisation, la commutativité
de ce diagramme implique celle du diagramme du lemme.
\end{proof}

\begin{lemma}
\label{pione-lemma-specialization-map-composition}
Soit $f : X \to S$ un morphisme propre à fibres géométriquement connexes.
Soient $s'' \leadsto s' \leadsto s$ des spécialisations de points de $S$.
Une composée d'homomorphismes de spécialisation
$\pi_1(X_{\overline{s}''}) \to \pi_1(X_{\overline{s}'}) \to
\pi_1(X_{\overline{s}})$ est un homomorphisme de spécialisation
$\pi_1(X_{\overline{s}''}) \to \pi_1(X_{\overline{s}})$.
\end{lemma}

\begin{proof}
Soit $\mathcal{O}_{S, s} \to A$ l'hensélisé strict
construit à l'aide de $\kappa(s) \to \kappa(\overline{s})$.
Soit $A \to \kappa(\overline{s}')$ le morphisme utilisé pour construire
le premier homomorphisme de spécialisation. Soit $\mathcal{O}_{S, s'} \to A'$
l'hensélisé strict construit à l'aide de
$\kappa(s') \subset \kappa(\overline{s}')$.
Par fonctorialité de l'hensélisation stricte, il existe un morphisme
$A \to A'$ tel que sa composée avec $A' \to \kappa(\overline{s}')$
soit le morphisme donné
(Algèbre, lemme \ref{algebra-lemma-map-into-henselian-colimit}).
Ensuite, soit $A' \to \kappa(\overline{s}'')$ le morphisme utilisé pour
construire le second homomorphisme de spécialisation. Il est alors clair que
la composée du premier et du second homomorphisme de spécialisation
est l'homomorphisme de spécialisation
$\pi_1(X_{\overline{s}''}) \to \pi_1(X_{\overline{s}})$
construit à l'aide de $A \to A' \to \kappa(\overline{s}'')$.
\end{proof}

\noindent
Soit $X \to S$ un morphisme propre à fibres géométriquement connexes.
Soit $R$ un anneau de valuation strictement hensélien dont le corps des fractions
est algébriquement clos et soit $\Spec(R) \to S$
un morphisme. Soient $\eta, s \in \Spec(R)$ les points générique et fermé.
On peut alors considérer l'homomorphisme de spécialisation
$$
sp_R : \pi_1(X_\eta) \to \pi_1(X_s)
$$
pour le changement de base $X_R/\Spec(R)$. Notons que cela a un sens puisque
$\eta$ et $s$ ont tous deux un corps résiduel algébriquement clos.

\begin{lemma}
\label{pione-lemma-specialization-map-valuation-ring}
Soit $f : X \to S$ un morphisme propre à fibres géométriquement connexes.
Soit $s' \leadsto s$ une spécialisation de points de $S$ et soit
$sp : \pi_1(X_{\overline{s}'}) \to \pi_1(X_{\overline{s}})$
un homomorphisme de spécialisation. Il existe alors un anneau de valuation strictement hensélien
$R$ au-dessus de $S$ dont le corps des fractions est algébriquement clos,
tel que $sp$ soit isomorphe à $sp_R$ défini ci-dessus.
\end{lemma}

\begin{proof}
Soit $\mathcal{O}_{S, s} \to A$ l'hensélisé strict
construit à l'aide de $\kappa(s) \to \kappa(\overline{s})$.
Soit $A \to \kappa(\overline{s}')$ le morphisme utilisé pour construire $sp$.
Soit $R \subset \kappa(\overline{s}')$ un anneau de valuation de
corps des fractions $\kappa(\overline{s}')$ qui domine l'image de $A$.
Voir Algèbre, lemme \ref{algebra-lemma-dominate}.
Observons que $R$ est strictement hensélien, par exemple d'après le
lemme \ref{pione-lemma-normal-local-domain-separablly-closed-fraction-field}
et Algèbre, lemme \ref{algebra-lemma-valuation-ring-normal}.
Le lemme est alors clair.
\end{proof}

\noindent
Soit $X \to S$ un morphisme propre à fibres géométriquement connexes.
Soit $R$ un anneau de valuation discrète strictement hensélien et soit
$\Spec(R) \to S$ un morphisme. Soient $\eta, s \in \Spec(R)$ les points
générique et fermé. On peut alors considérer l'homomorphisme de spécialisation
$$
sp_R : \pi_1(X_{\overline{\eta}}) \to \pi_1(X_s)
$$
pour le changement de base $X_R/\Spec(R)$. Notons que cela a un sens puisque $s$
a un corps résiduel algébriquement clos.

\begin{lemma}
\label{pione-lemma-specialization-map-discrete-valuation-ring}
Soit $f : X \to S$ un morphisme propre à fibres géométriquement connexes.
Soit $s' \leadsto s$ une spécialisation de points de $S$ et soit
$sp : \pi_1(X_{\overline{s}'}) \to \pi_1(X_{\overline{s}})$
un homomorphisme de spécialisation. Si $S$ est noethérien, alors
il existe un anneau de valuation discrète strictement hensélien
$R$ au-dessus de $S$ tel que $sp$ soit isomorphe à $sp_R$
défini ci-dessus.
\end{lemma}

\begin{proof}
Soit $\mathcal{O}_{S, s} \to A$ l'hensélisé strict
construit à l'aide de $\kappa(s) \to \kappa(\overline{s})$.
Soit $A \to \kappa(\overline{s}')$ le morphisme utilisé pour construire $sp$.
Soit $R \subset \kappa(\overline{s}')$ un anneau de valuation discrète
qui domine l'image de $A$, voir Algèbre, lemme \ref{algebra-lemma-exists-dvr}.
Choisissons un diagramme de corps
$$
\xymatrix{
\kappa(\overline{s}) \ar[r] & k \\
A/\mathfrak m_A \ar[r] \ar[u] & R/\mathfrak m_R \ar[u]
}
$$
où $k$ est algébriquement clos. Soit $R^{sh}$ l'hensélisé
strict de $R$ construit à l'aide de $R \to k$. On obtient
un morphisme $A \to R^{sh}$ d'après Algèbre, lemme
\ref{algebra-lemma-strictly-henselian-functorial}.
L'anneau $R^{sh}$ est un anneau de valuation discrète d'après
Compléments d'algèbre, lemme \ref{more-algebra-lemma-henselization-dvr}.
Notons $\eta, o$ les points générique et fermé de $\Spec(R^{sh})$.
Comme le diagramme de schémas
$$
\xymatrix{
\overline{\eta} \ar[d] \ar[r] & \Spec(R^{sh}) \ar[d] &
\Spec(k) \ar[d] \ar[l] \\
\overline{s}' \ar[r] & \Spec(A) & \overline{s} \ar[l]
}
$$
est commutatif, on obtient un diagramme commutatif
$$
\xymatrix{
\pi_1(X_{\overline{\eta}}) \ar[d] \ar[r]_{sp_{R^{sh}}} & \pi_1(X_o) \ar[d] \\
\pi_1(X_{\overline{s}'}) \ar[r]^{sp} & \pi_1(X_{\overline{s}})
}
$$
d'homomorphismes de spécialisation, par la construction de ces homomorphismes.
Comme les flèches verticales sont des isomorphismes
(lemme \ref{pione-lemma-finite-etale-invariant-over-proper}), cela démontre le lemme.
\end{proof}








\section{Restriction à un sous-schéma fermé}
\label{pione-section-lefschetz}

\noindent
Dans cette section, nous démontrons quelques résultats sur le foncteur de restriction
$$
\textit{F\'Et}_X \longrightarrow \textit{F\'Et}_Y,\quad
U \longmapsto V = U \times_X Y
$$
où $X$ est un schéma et $Y$ un sous-schéma fermé. Grâce à
l'invariance topologique du groupe fondamental, on peut ramener
l'étude de ce foncteur à celle du foncteur de complétion sur les
modules finis localement libres.

\medskip\noindent
Dans les lemmes suivants, nous utilisons la notion de modules formels cohérents
définie dans
Cohomologie des schémas, section \ref{coherent-section-coherent-formal}.
Étant donnés un schéma noethérien et un faisceau quasi-cohérent d'idéaux
$\mathcal{I} \subset \mathcal{O}_X$, nous dirons
qu'un objet $(\mathcal{F}_n)$ de $\textit{Coh}(X, \mathcal{I})$
est {\it fini localement libre} si chaque $\mathcal{F}_n$ est un
$\mathcal{O}_X/\mathcal{I}^n$-module fini localement libre.

\begin{lemma}
\label{pione-lemma-restriction-fully-faithful}
Soit $X$ un schéma noethérien et soit $Y \subset X$ un sous-schéma fermé
de faisceau d'idéaux $\mathcal{I} \subset \mathcal{O}_X$.
Supposons que le foncteur de complétion
$$
\textit{Coh}(\mathcal{O}_X)
\longrightarrow 
\textit{Coh}(X, \mathcal{I}),\quad
\mathcal{F} \longmapsto \mathcal{F}^\wedge
$$
soit pleinement fidèle sur la sous-catégorie pleine des objets finis
localement libres (voir ci-dessus).
Alors le foncteur de restriction $\textit{F\'Et}_X \to \textit{F\'Et}_Y$
est pleinement fidèle.
\end{lemma}

\begin{proof}
Comme la catégorie des revêtements finis étales possède un
Hom interne (lemme \ref{pione-lemma-internal-hom-finite-etale}),
il suffit de démontrer ceci : étant donnés $U$ fini étale sur $X$
et un morphisme $t : Y \to U$ sur $X$, il existe une unique section
$s : X \to U$ telle que $t = s|_Y$. Voici le diagramme :
$$
\xymatrix{
& U \ar[d]^f \\
Y \ar[r] \ar[ru] & X \ar@{..>}@/^1em/[u]
}
$$
Trouver la flèche pointillée $s$ revient à trouver un morphisme de
$\mathcal{O}_X$-algèbres
$$
s^\sharp : f_*\mathcal{O}_U \longrightarrow \mathcal{O}_X
$$
qui, modulo le faisceau d'idéaux de $Y$, se réduise au morphisme d'algèbres donné
$t^\sharp : f_*\mathcal{O}_U \to \mathcal{O}_Y$.
D'après le lemme \ref{pione-lemma-thickening}, on peut relever $t$ de façon unique en un système
compatible de morphismes $t_n : Y_n \to U$, et donc en un morphisme
$$
\lim t_n^\sharp : f_*\mathcal{O}_U \longrightarrow \lim \mathcal{O}_{Y_n}
$$
de faisceaux d'algèbres sur $X$.
Comme $f_*\mathcal{O}_U$ est un $\mathcal{O}_X$-module fini localement libre,
on en déduit un unique morphisme de $\mathcal{O}_X$-modules
$\sigma : f_*\mathcal{O}_U \to \mathcal{O}_X$ dont la complétion
est $\lim t_n^\sharp$. Pour voir que $\sigma$ est un homomorphisme d'algèbres,
il faut vérifier que le diagramme
$$
\xymatrix{
f_*\mathcal{O}_U \otimes_{\mathcal{O}_X} f_*\mathcal{O}_U
\ar[r] \ar[d]_{\sigma \otimes \sigma} &
f_*\mathcal{O}_U \ar[d]^\sigma \\
\mathcal{O}_X \otimes_{\mathcal{O}_X} \mathcal{O}_X \ar[r] &
\mathcal{O}_X
}
$$
est commutatif. Pour tout $n$, on sait que ce diagramme est commutatif après restriction
à $Y_n$, c'est-à-dire après application du foncteur de complétion.
La fidélité du foncteur de complétion permet donc de conclure.
\end{proof}

\begin{lemma}
\label{pione-lemma-restriction-equivalence}
Soit $X$ un schéma noethérien et soit $Y \subset X$ un sous-schéma fermé
de faisceau d'idéaux $\mathcal{I} \subset \mathcal{O}_X$.
Supposons que le foncteur de complétion
$$
\textit{Coh}(\mathcal{O}_X)
\longrightarrow 
\textit{Coh}(X, \mathcal{I}),\quad
\mathcal{F} \longmapsto \mathcal{F}^\wedge
$$
soit une équivalence sur les sous-catégories pleines d'objets finis
localement libres (voir ci-dessus).
Alors le foncteur de restriction $\textit{F\'Et}_X \to \textit{F\'Et}_Y$
est une équivalence.
\end{lemma}

\begin{proof}
Le foncteur de restriction est pleinement fidèle d'après le
lemme \ref{pione-lemma-restriction-fully-faithful}.

\medskip\noindent
Soit $U_1 \to Y$ un morphisme fini étale. Pour achever la démonstration,
nous allons montrer que $U_1$ appartient à l'image essentielle du
foncteur de restriction.

\medskip\noindent
Pour $n \geq 1$, soit $Y_n$ le $n$-ième voisinage infinitésimal de $Y$.
D'après le lemme \ref{pione-lemma-thickening},
il existe un unique morphisme fini étale
$\pi_n : U_n \to Y_n$ dont le changement de base à $Y = Y_1$
redonne $U_1 \to Y_1$.
Considérons les faisceaux $\mathcal{F}_n = \pi_{n, *}\mathcal{O}_{U_n}$.
Nous pouvons et allons considérer $\mathcal{F}_n$ comme un $\mathcal{O}_X$-module sur $X$
qui est localement isomorphe à
$(\mathcal{O}_X/\mathcal{I}^n)^{\oplus r}$.
Ce $(\mathcal{F}_n)$ est un objet fini localement libre de
$\textit{Coh}(X, \mathcal{I})$.
Par hypothèse, il existe un $\mathcal{O}_X$-module fini localement libre
$\mathcal{F}$ et un système compatible d'isomorphismes
$$
\mathcal{F}/\mathcal{I}^n\mathcal{F} \to \mathcal{F}_n
$$
de $\mathcal{O}_X$-modules.

\medskip\noindent
Pour munir $\mathcal{F}$ d'une structure d'algèbre, considérons les
morphismes de multiplication
$\mathcal{F}_n \otimes_{\mathcal{O}_X} \mathcal{F}_n \to \mathcal{F}_n$
provenant du fait que les $\mathcal{F}_n = \pi_{n, *}\mathcal{O}_{U_n}$
sont des faisceaux d'algèbres. Ceux-ci définissent un morphisme
$$
(\mathcal{F}\otimes_{\mathcal{O}_X} \mathcal{F})^\wedge
\longrightarrow
\mathcal{F}^\wedge
$$
dans la catégorie $\textit{Coh}(X, \mathcal{I})$. Par hypothèse,
on peut donc supposer qu'il existe un morphisme
$\mu : \mathcal{F}\otimes_{\mathcal{O}_X} \mathcal{F} \to \mathcal{F}$
dont la restriction à $Y_n$ donne les morphismes de multiplication ci-dessus.
La fidélité du foncteur de l'énoncé du lemme montre alors
que $\mu$ définit une structure de $\mathcal{O}_X$-algèbre commutative
sur $\mathcal{F}$, compatible avec les structures d'algèbres données
sur les $\mathcal{F}_n$.
En posant
$$
U = \underline{\Spec}_X((\mathcal{F}, \mu))
$$
on obtient un schéma fini localement libre $\pi : U \to X$ dont la restriction
à $Y$ est isomorphe à $U_1$. Le discriminant de $\pi$ est le lieu des zéros
de la section
$$
\det(Q_\pi) :
\mathcal{O}_X
\longrightarrow
\wedge^{top}(\pi_*\mathcal{O}_U)^{\otimes -2}
$$
construite dans
Discriminants, section \ref{discriminant-section-discriminant}.
Comme sa restriction à $Y_n$ est un isomorphisme pour tout $n$
d'après Discriminants, lemme \ref{discriminant-lemma-discriminant},
on en conclut que c'est un isomorphisme. Ainsi $\pi$ est étale d'après
Discriminants, lemme \ref{discriminant-lemma-discriminant}.
\end{proof}

\begin{lemma}
\label{pione-lemma-restriction-fully-faithful-general}
Soit $X$ un schéma noethérien et soit $Y \subset X$ un sous-schéma fermé
de faisceau d'idéaux $\mathcal{I} \subset \mathcal{O}_X$.
Soit $\mathcal{V}$ l'ensemble des sous-schémas ouverts $V \subset X$ contenant $Y$,
ordonné par inclusion inverse. Supposons que le foncteur de complétion
$$
\colim_\mathcal{V} \textit{Coh}(\mathcal{O}_V)
\longrightarrow
\textit{Coh}(X, \mathcal{I}),
\quad
\mathcal{F} \longmapsto \mathcal{F}^\wedge
$$
soit pleinement fidèle sur la sous-catégorie pleine des
objets finis localement libres (voir ci-dessus).
Alors le foncteur de restriction
$\colim_\mathcal{V} \textit{F\'Et}_V \to \textit{F\'Et}_Y$
est pleinement fidèle.
\end{lemma}

\begin{proof}
Observons que $\mathcal{V}$ est un ensemble filtrant, de sorte que les colimites sont
celles de Catégories, section \ref{categories-section-directed-colimits}.
La suite de l'argument est presque exactement la même que dans la démonstration
du lemme \ref{pione-lemma-restriction-fully-faithful}; nous invitons le lecteur
à l'omettre.

\medskip\noindent
Comme la catégorie des revêtements finis étales possède un
Hom interne (lemme \ref{pione-lemma-internal-hom-finite-etale}),
il suffit de démontrer ceci : étant donnés $U$ fini étale sur
$V \in \mathcal{V}$
et un morphisme $t : Y \to U$ sur $V$, il existe $V' \geq V$
et un morphisme $s : V' \to U$ sur $V$ tel que $t = s|_Y$. Voici le diagramme :
$$
\xymatrix{
& & U \ar[d]^f \\
Y \ar[r] \ar[rru] & V' \ar@{..>}[ru] \ar[r] & V
}
$$
Trouver la flèche pointillée $s$ revient à trouver un morphisme de
$\mathcal{O}_{V'}$-algèbres
$$
s^\sharp : f_*\mathcal{O}_U|_{V'} \longrightarrow \mathcal{O}_{V'}
$$
qui, modulo le faisceau d'idéaux de $Y$, se réduise au morphisme d'algèbres donné
$t^\sharp : f_*\mathcal{O}_U \to \mathcal{O}_Y$.
D'après le lemme \ref{pione-lemma-thickening}, on peut relever $t$ de façon unique en un système
compatible de morphismes $t_n : Y_n \to U$, et donc un morphisme
$$
\lim t_n^\sharp : f_*\mathcal{O}_U \longrightarrow \lim \mathcal{O}_{Y_n}
$$
de faisceaux d'algèbres sur $V$.
Observons que $f_*\mathcal{O}_U$ est un
$\mathcal{O}_V$-module fini localement libre. Il existe donc $V' \geq V$ et un morphisme
$\sigma : f_*\mathcal{O}_U|_{V'} \to \mathcal{O}_{V'}$
dont la complétion est $\lim t_n^\sharp$.
Pour voir que $\sigma$ est un homomorphisme d'algèbres, il faut vérifier
que le diagramme
$$
\xymatrix{
(f_*\mathcal{O}_U \otimes_{\mathcal{O}_V} f_*\mathcal{O}_U)|_{V'}
\ar[r] \ar[d]_{\sigma \otimes \sigma} &
f_*\mathcal{O}_U|_{V'} \ar[d]^\sigma \\
\mathcal{O}_{V'} \otimes_{\mathcal{O}_{V'}} \mathcal{O}_{V'} \ar[r] &
\mathcal{O}_{V'}
}
$$
est commutatif. Pour tout $n$, on sait que ce diagramme est commutatif après restriction
à $Y_n$, c'est-à-dire après application du foncteur de complétion.
La fidélité du foncteur de complétion permet donc
d'en déduire qu'il existe $V'' \geq V'$ tel que
$\sigma|_{V''}$ soit un homomorphisme d'algèbres, comme souhaité.
\end{proof}

\begin{lemma}
\label{pione-lemma-restriction-equivalence-general}
Soit $X$ un schéma noethérien et soit $Y \subset X$ un sous-schéma fermé
de faisceau d'idéaux $\mathcal{I} \subset \mathcal{O}_X$.
Soit $\mathcal{V}$ l'ensemble des sous-schémas ouverts $V \subset X$ contenant $Y$,
ordonné par inclusion inverse. Supposons que le foncteur de complétion
$$
\colim_\mathcal{V} \textit{Coh}(\mathcal{O}_V)
\longrightarrow
\textit{Coh}(X, \mathcal{I}),
\quad
\mathcal{F} \longmapsto \mathcal{F}^\wedge
$$
induise une équivalence des sous-catégories pleines d'objets
finis localement libres (voir l'explication ci-dessus).
Alors le foncteur de restriction
$$
\colim_\mathcal{V} \textit{F\'Et}_V \to \textit{F\'Et}_Y
$$
est une équivalence.
\end{lemma}

\begin{proof}
Observons que $\mathcal{V}$ est un ensemble filtrant, de sorte que les colimites sont
celles de Catégories, section \ref{categories-section-directed-colimits}.
La suite de l'argument est presque exactement la même que dans la démonstration
du lemme \ref{pione-lemma-restriction-equivalence}; nous invitons le lecteur
à l'omettre.

\medskip\noindent
Le foncteur de restriction est pleinement fidèle d'après le
lemme \ref{pione-lemma-restriction-fully-faithful-general}.

\medskip\noindent
Soit $U_1 \to Y$ un morphisme fini étale. Pour achever la démonstration,
nous allons montrer que $U_1$ appartient à l'image essentielle du
foncteur de restriction.

\medskip\noindent
Pour $n \geq 1$, soit $Y_n$ le $n$-ième voisinage infinitésimal de $Y$.
D'après le lemme \ref{pione-lemma-thickening},
il existe un unique morphisme fini étale
$\pi_n : U_n \to Y_n$ dont le changement de base à $Y = Y_1$
redonne $U_1 \to Y_1$.
Considérons les faisceaux $\mathcal{F}_n = \pi_{n, *}\mathcal{O}_{U_n}$.
Nous pouvons et allons considérer $\mathcal{F}_n$ comme un $\mathcal{O}_X$-module sur $X$
qui est localement isomorphe à
$(\mathcal{O}_X/\mathcal{I}^n)^{\oplus r}$.
Ce $(\mathcal{F}_n)$ est un objet fini localement libre de
$\textit{Coh}(X, \mathcal{I})$.
Par hypothèse, il existe $V \in \mathcal{V}$,
un $\mathcal{O}_V$-module fini localement libre $\mathcal{F}$
et un système compatible d'isomorphismes
$$
\mathcal{F}/\mathcal{I}^n\mathcal{F} \to \mathcal{F}_n
$$
de $\mathcal{O}_V$-modules.

\medskip\noindent
Pour munir $\mathcal{F}$ d'une structure d'algèbre, considérons les
morphismes de multiplication
$\mathcal{F}_n \otimes_{\mathcal{O}_V} \mathcal{F}_n \to \mathcal{F}_n$
provenant du fait que les $\mathcal{F}_n = \pi_{n, *}\mathcal{O}_{U_n}$
sont des faisceaux d'algèbres. Ceux-ci définissent un morphisme
$$
(\mathcal{F}\otimes_{\mathcal{O}_V} \mathcal{F})^\wedge
\longrightarrow
\mathcal{F}^\wedge
$$
dans la catégorie $\textit{Coh}(X, \mathcal{I})$. Par hypothèse,
après avoir rétréci $V$, on peut supposer qu'il existe un morphisme
$\mu : \mathcal{F}\otimes_{\mathcal{O}_V} \mathcal{F} \to \mathcal{F}$
dont la restriction à $Y_n$ donne les morphismes de multiplication ci-dessus.
Après l'avoir éventuellement rétréci davantage, on peut supposer que $\mu$
définit une structure de $\mathcal{O}_V$-algèbre commutative
sur $\mathcal{F}$, compatible avec les structures d'algèbres données
sur les $\mathcal{F}_n$.
En posant
$$
U = \underline{\Spec}_V((\mathcal{F}, \mu))
$$
on obtient un schéma fini localement libre sur $V$ dont la restriction
à $Y$ est isomorphe à $U_1$. Il s'ensuit que $U \to V$
est étale en tous les points situés au-dessus de $Y$; voir
Compléments sur les morphismes, lemme
\ref{more-morphisms-lemma-check-smoothness-on-infinitesimal-nbhds}.
Ainsi, après avoir rétréci $V$ une fois encore, on peut supposer que $U \to V$ est
fini étale. Cela achève la démonstration.
\end{proof}

\begin{lemma}
\label{pione-lemma-restriction-faithful}
Soit $X$ un schéma et soit $Y \subset X$ un sous-schéma fermé.
Si toute composante connexe de $X$ rencontre $Y$, alors
le foncteur de restriction $\textit{F\'Et}_X \to \textit{F\'Et}_Y$
est fidèle.
\end{lemma}

\begin{proof}
Soient $a, b : U \to U'$ deux morphismes de schémas finis étales sur $X$
dont les restrictions à $Y$ sont égales. L'image d'une composante connexe
de $U$ est une composante connexe de $X$; cela résulte de
Topologie, lemme \ref{topology-lemma-finite-fibre-connected-components},
appliqué à la restriction de $U \to X$ à une composante connexe de $U$.
Par conséquent, l'image de toute composante connexe de $U$ rencontre $Y$
par hypothèse. On en conclut que $a = b$ après restriction à chaque
composante connexe de $U$, d'après Morphismes étales, proposition
\ref{etale-proposition-equality}. Comme l'égalisateur de $a$ et $b$
est un sous-schéma ouvert de $U$ (puisque la diagonale de $U'$ sur $X$ est ouverte),
on conclut.
\end{proof}

\begin{lemma}
\label{pione-lemma-restriction-fully-faithful-special}
Soit $X$ un schéma noethérien et soit $Y \subset X$ un sous-schéma fermé.
Soit $Y_n \subset X$ le $n$-ième voisinage infinitésimal de $Y$ dans $X$.
Supposons que l'une des conditions suivantes soit satisfaite :
\begin{enumerate}
\item $X$ est quasi-affine et
$\Gamma(X, \mathcal{O}_X) \to \lim \Gamma(Y_n, \mathcal{O}_{Y_n})$
est un isomorphisme, ou bien
\item $X$ possède un module inversible ample $\mathcal{L}$ et
$\Gamma(X, \mathcal{L}^{\otimes m}) \to
\lim \Gamma(Y_n, \mathcal{L}^{\otimes m}|_{Y_n})$
est un isomorphisme pour tout $m \gg 0$, ou bien
\item pour tout $\mathcal{O}_X$-module fini localement libre
$\mathcal{E}$, le morphisme
$\Gamma(X, \mathcal{E}) \to \lim \Gamma(Y_n, \mathcal{E}|_{Y_n})$
est un isomorphisme.
\end{enumerate}
Alors le foncteur de restriction $\textit{F\'Et}_X \to \textit{F\'Et}_Y$
est pleinement fidèle.
\end{lemma}

\begin{proof}
Ce lemme résulte formellement du
lemme \ref{pione-lemma-restriction-fully-faithful} et de
Géométrie algébrique et formelle, lemme
\ref{algebraization-lemma-completion-fully-faithful}.
\end{proof}

\begin{lemma}
\label{pione-lemma-restriction-fully-faithful-general-special}
Soit $X$ un schéma noethérien et soit $Y \subset X$ un sous-schéma fermé.
Soit $Y_n \subset X$ le $n$-ième voisinage infinitésimal de $Y$ dans $X$.
Soit $\mathcal{V}$ l'ensemble des sous-schémas ouverts $V \subset X$ contenant $Y$,
ordonné par inclusion inverse. Supposons que l'une des conditions suivantes soit satisfaite :
\begin{enumerate}
\item $X$ est quasi-affine et
$$
\colim_\mathcal{V} \Gamma(V, \mathcal{O}_V)
\longrightarrow
\lim \Gamma(Y_n, \mathcal{O}_{Y_n})
$$
est un isomorphisme, ou bien
\item $X$ possède un module inversible ample $\mathcal{L}$ et
$$
\colim_\mathcal{V} \Gamma(V, \mathcal{L}^{\otimes m})
\longrightarrow
\lim \Gamma(Y_n, \mathcal{L}^{\otimes m}|_{Y_n})
$$
est un isomorphisme pour tout $m \gg 0$, ou bien
\item pour tout $V \in \mathcal{V}$ et tout
$\mathcal{O}_V$-module fini localement libre $\mathcal{E}$, le morphisme
$$
\colim_{V' \geq V} \Gamma(V', \mathcal{E}|_{V'})
\longrightarrow
\lim \Gamma(Y_n, \mathcal{E}|_{Y_n})
$$
est un isomorphisme.
\end{enumerate}
Alors le foncteur
$$
\colim_\mathcal{V} \textit{F\'Et}_V \to \textit{F\'Et}_Y
$$
est pleinement fidèle.
\end{lemma}

\begin{proof}
Ce lemme résulte formellement du
lemme \ref{pione-lemma-restriction-fully-faithful-general} et de
Géométrie algébrique et formelle, lemme
\ref{algebraization-lemma-completion-fully-faithful-general}.
\end{proof}






\section{Sommes amalgamées et groupes fondamentaux}
\label{pione-section-pushouts}

\noindent
Voici le résultat principal.

\begin{lemma}
\label{pione-lemma-pushout-along-closed-immersion-and-integral}
Dans Compléments sur les morphismes, situation
\ref{more-morphisms-situation-pushout-along-closed-immersion-and-integral},
par exemple si $Z \to Y$ et $Z \to X$ sont des immersions fermées de schémas,
il existe une équivalence de catégories
$$
\textit{F\'Et}_{Y \amalg_Z X}
\longrightarrow
\textit{F\'Et}_Y
\times_{\textit{F\'Et}_Z}
\textit{F\'Et}_X
$$
\end{lemma}

\begin{proof}
La somme amalgamée existe d'après
Compléments sur les morphismes, proposition
\ref{more-morphisms-proposition-pushout-along-closed-immersion-and-integral}.
Le foncteur associe à un schéma $U$ fini étale sur la
somme amalgamée les changements de base $Y' = U \times_{Y \amalg_Z X} Y$
et $X' = U \times_{Y \amalg_Z X} X$, munis de l'isomorphisme naturel
$Y' \times_Y Z \to X' \times_X Z$ sur $Z$. Pour montrer que ce foncteur
est une équivalence, nous utilisons
Compléments sur les morphismes, lemme
\ref{more-morphisms-lemma-pushout-functor-equivalence-flat}
pour construire un foncteur quasi-inverse.
Il reste seulement à montrer que, si l'on se donne un morphisme
$U \to Y \amalg_Z X$ séparé, quasi-fini et étale
tel que $X' \to X$ et $Y' \to Y$ soient finis,
alors $U \to Y \amalg_Z X$ est fini.
Cela se déduit soit du fait algébrique correspondant
(Compléments d'algèbre, lemme
\ref{more-algebra-lemma-finite-module-over-fibre-product}),
soit du fait que
$$
X' \amalg Y' \to U
$$
est surjectif et que $X'$ et $Y'$ sont propres sur $Y \amalg_Z X$
(on utilise ici la description de la somme amalgamée dans Compléments sur les morphismes, proposition
\ref{more-morphisms-proposition-pushout-along-closed-immersion-and-integral}),
après quoi on peut appliquer
Morphismes, lemme \ref{morphisms-lemma-scheme-theoretic-image-is-proper},
pour conclure que $U$ est propre sur $Y \amalg_Z X$.
Comme un morphisme quasi-fini et propre est fini
(Compléments sur les morphismes, lemme \ref{more-morphisms-lemma-characterize-finite}),
la démonstration est achevée.
\end{proof}





\section{Revêtements finis étales des spectres épointés, I}
\label{pione-section-pi1-punctured-spec}

\noindent
Nous démontrons d'abord quelques résultats à la Lefschetz.

\begin{situation}
\label{pione-situation-local-lefschetz}
Soient $(A, \mathfrak m)$ un anneau local noethérien et $f \in \mathfrak m$.
Posons $X = \Spec(A)$ et $X_0 = \Spec(A/fA)$, et
notons $U = X \setminus \{\mathfrak m\}$ et
$U_0 = X_0 \setminus \{\mathfrak m\}$ les spectres épointés de
$A$ et $A/fA$.
\end{situation}

\noindent
Rappelons que, pour un schéma $X$, la catégorie des schémas finis
étales sur $X$ est notée $\textit{F\'Et}_X$ ; voir la
section \ref{pione-section-finite-etale}.
Dans la situation \ref{pione-situation-local-lefschetz},
nous étudierons les foncteurs de changement de base
$$
\xymatrix{
\textit{F\'Et}_X \ar[d] \ar[r] & \textit{F\'Et}_U \ar[d] \\
\textit{F\'Et}_{X_0} \ar[r] & \textit{F\'Et}_{U_0}
}
$$
Dans de nombreux cas, la flèche verticale de droite est fidèle.

\begin{lemma}
\label{pione-lemma-faithful}
Dans la situation \ref{pione-situation-local-lefschetz},
supposons que l'une des conditions suivantes soit satisfaite :
\begin{enumerate}
\item $\dim(A/\mathfrak p) \geq 2$ pour tout idéal premier minimal
$\mathfrak p \subset A$ tel que $f \not \in \mathfrak p$, ou bien
\item toute composante connexe de $U$ rencontre $U_0$.
\end{enumerate}
Alors
$$
\textit{F\'Et}_U \longrightarrow \textit{F\'Et}_{U_0},\quad
V \longmapsto V_0 = V \times_U U_0
$$
est un foncteur fidèle.
\end{lemma}

\begin{proof}
Le cas (2) résulte immédiatement du lemme \ref{pione-lemma-restriction-faithful}.
L'hypothèse (1) implique que toute composante irréductible de $U$ rencontre $U_0$ ; voir
Algèbre, lemme \ref{algebra-lemma-one-equation}.
Ainsi (1) résulte de (2).
\end{proof}

\noindent
Avant de démontrer un résultat plus intéressant, il nous faut quelques lemmes.

\begin{lemma}
\label{pione-lemma-fill-in-missing}
Dans la situation \ref{pione-situation-local-lefschetz}, soit $V \to U$ un morphisme
fini. Soit $A^\wedge$ le complété $\mathfrak m$-adique de $A$,
soit $X' = \Spec(A^\wedge)$, et soient $U'$ et $V'$ les changements de base de
$U$ et $V$ à $X'$. S'il existe un morphisme fini $Y' \to X'$ tel que
$V' = Y' \times_{X'} U'$, alors il existe un morphisme fini $Y \to X$
tel que $V = Y \times_X U$ et $Y' = Y \times_X X'$.
\end{lemma}

\begin{proof}
C'est une application directe de
Compléments d'algèbre, proposition \ref{more-algebra-proposition-equivalence}.
Choisissons en effet des générateurs $f_1, \ldots, f_t$ de $\mathfrak m$.
Pour tout $i$, écrivons $V \times_U D(f_i) = \Spec(B_i)$.
Pour $1 \leq i, j \leq t$, on obtient un isomorphisme
$\alpha_{ij} : (B_i)_{f_j} \to (B_j)_{f_i}$ de $A_{f_if_j}$-algèbres,
car les spectres de ces deux algèbres représentent $V \times_U D(f_if_j)$.
Écrivons $Y' = \Spec(B')$. Comme $V \times_U U' = Y' \times_{X'} U'$
on obtient des isomorphismes $\alpha_i : B'_{f_i} \to B_i \otimes_A A^\wedge$.
Un argument direct montre que $(B', B_i, \alpha_i, \alpha_{ij})$
est un objet de $\text{Glue}(A \to A^\wedge, f_1, \ldots, f_t)$ ; voir
Compléments d'algèbre, remarque \ref{more-algebra-remark-glueing-data}.
En appliquant la proposition citée ci-dessus (et en utilisant
Compléments d'algèbre, remarque \ref{more-algebra-remark-formal-glueing-algebras},
pour obtenir la structure d'algèbre), on trouve une $A$-algèbre $B$ telle que
$\text{Can}(B)$ soit isomorphe à $(B', B_i, \alpha_i, \alpha_{ij})$.
En posant $Y = \Spec(B)$, on voit que $Y \to X$ est un morphisme
muni d'isomorphismes compatibles
$V \cong Y \times_X U$ et $Y' = Y \times_X X'$, comme souhaité.
\end{proof}

\begin{lemma}
\label{pione-lemma-fully-faithful-henselian-completion}
Dans la situation \ref{pione-situation-local-lefschetz}, supposons $A$ hensélien
ou, plus généralement, que $(A, (f))$ soit un couple hensélien.
Soit $A^\wedge$ le complété $\mathfrak m$-adique de $A$,
soit $X' = \Spec(A^\wedge)$, et soient $U'$ et $U'_0$ les changements de base de
$U$ et $U_0$ à $X'$. Si $\textit{F\'Et}_{U'} \to \textit{F\'Et}_{U'_0}$
est pleinement fidèle, alors $\textit{F\'Et}_U \to \textit{F\'Et}_{U_0}$
est pleinement fidèle.
\end{lemma}

\begin{proof}
Supposons $\textit{F\'Et}_{U'} \longrightarrow \textit{F\'Et}_{U'_0}$
pleinement fidèle. Comme $X' \to X$ est fidèlement plat, il est
immédiat que le foncteur $V \to V_0 = V \times_U U_0$ est fidèle.
Comme la catégorie des revêtements finis étales possède un Hom interne
(lemme \ref{pione-lemma-internal-hom-finite-etale}),
il suffit de démontrer ceci : pour $V$ fini étale sur $U$,
on a
$$
\Mor_U(U, V) = \Mor_{U_0}(U_0, V_0)
$$
Supposons donc donné un morphisme $s_0 : U_0 \to V_0$ sur $U_0$ ; nous allons
construire un morphisme $s : U \to V$ sur $U$.

\medskip\noindent
Par hypothèse, il existe bien un morphisme $s' : U' \to V'$
dont la restriction à $U'_0$ est le changement de base $s'_0$ de $s_0$.
Comme $V' \to U'$ est fini étale, cela signifie que $V' = s'(U') \amalg W'$
pour un certain morphisme $W' \to U'$ fini et étale.
Choisissons un morphisme fini $Z' \to X'$ tel que $W' = Z' \times_{X'} U'$.
C'est possible grâce au théorème principal de Zariski, sous la forme énoncée dans
Compléments sur les morphismes, lemme
\ref{more-morphisms-lemma-quasi-finite-separated-pass-through-finite}
(nous omettons un petit détail).
Alors
$$
V' = s'(U') \amalg W' \longrightarrow X' \amalg Z' = Y'
$$
est une immersion ouverte telle que $V' = Y' \times_{X'} U'$.
D'après le lemme \ref{pione-lemma-fill-in-missing}, il existe $Y \to X$ fini
tel que $V = Y \times_X U$ et $Y' = Y \times_X X'$.
Écrivons $Y = \Spec(B)$, de sorte que $Y' = \Spec(B \otimes_A A^\wedge)$.
Alors $B \otimes_A A^\wedge$ possède un idempotent $e'$
correspondant au sous-schéma ouvert et fermé $X'$ de $Y' = X' \amalg Z'$.

\medskip\noindent
Supposons d'abord $A$ hensélien (ce cas est légèrement plus simple). L'image $\overline{e}$
de $e'$ dans $B \otimes_A \kappa(\mathfrak m) = B/\mathfrak mB$ se relève en un
idempotent $e$ de $B$, puisque $A$ est hensélien (car $B$ est un produit d'anneaux
locaux d'après Algèbre, lemme \ref{algebra-lemma-characterize-henselian}).
On voit alors que $e$ s'envoie sur $e'$ par unicité du relèvement des idempotents
(en utilisant que $B \otimes_A A^\wedge$ est un produit d'anneaux locaux).
Soit $Y_1 \subset Y$ le sous-schéma ouvert et fermé correspondant à $e$.
Alors $Y_1 \times_X X' = s'(X')$, ce qui implique que $Y_1 \to X$ est
un isomorphisme (par descente fidèlement plate) et fournit la section voulue.

\medskip\noindent
Supposons maintenant que $(A, (f))$ soit un couple hensélien. Nous utilisons ici que $s'$ est
un relèvement de $s'_0$. Soit en effet $Y_{0, 1} \subset Y_0 = Y \times_X X_0$
l'adhérence de $s_0(U_0) \subset V_0 = Y_0 \times_{X_0} U_0$.
Comme $X' \to X$ est plat, le changement de base $Y'_{0, 1} \subset Y'_0$
est l'adhérence de $s'_0(U'_0)$, laquelle est égale à $X'_0 \subset Y'_0$
(voir Morphismes, lemme
\ref{morphisms-lemma-flat-base-change-scheme-theoretic-image}).
Comme $Y'_0 \to Y_0$ est submersif
(Morphismes, lemme \ref{morphisms-lemma-fpqc-quotient-topology}),
on conclut que $Y_{0, 1}$ est ouvert et fermé dans $Y_0$.
Soit $e_0 \in B/fB$ l'idempotent correspondant.
D'après Compléments d'algèbre, lemme
\ref{more-algebra-lemma-characterize-henselian-pair},
on peut relever $e_0$ en un idempotent $e \in B$.
On conclut alors comme précédemment.
\end{proof}

\noindent
Dans la situation \ref{pione-situation-local-lefschetz},
la pleine fidélité du foncteur de restriction
$\textit{F\'Et}_U \longrightarrow \textit{F\'Et}_{U_0}$
vaut sous des hypothèses assez faibles.
En particulier, ces hypothèses n'impliquent souvent pas que
$U$ soit un schéma connexe, mais la conclusion garantit
que $U$ et $U_0$ ont le même nombre de composantes connexes.

\begin{lemma}
\label{pione-lemma-fully-faithful-simple}
Dans la situation \ref{pione-situation-local-lefschetz}, supposons que
\begin{enumerate}
\item[(a)] $A$ possède un complexe dualisant,
\item[(b)] le couple $(A, (f))$ est hensélien,
\item[(c)] l'une des conditions suivantes est satisfaite :
\begin{enumerate}
\item[(i)] $A_f$ vérifie $(S_2)$ et toute composante irréductible de $X$
non contenue dans $X_0$ est de dimension $\geq 3$, ou bien
\item[(ii)] pour tout idéal premier
$\mathfrak p \subset A$ tel que $f \not \in \mathfrak p$, on ait
$\text{depth}(A_\mathfrak p) + \dim(A/\mathfrak p) > 2$.
\end{enumerate}
\end{enumerate}
Alors le foncteur de restriction
$\textit{F\'Et}_U \longrightarrow \textit{F\'Et}_{U_0}$
est pleinement fidèle.
\end{lemma}

\begin{proof}
Soit $A'$ le complété $\mathfrak m$-adique de $A$. Nous allons montrer que
les hypothèses restent vraies pour $A'$. C'est clair pour les conditions
(a) et (b). La condition (c)(ii) est préservée par
Cohomologie locale, lemme \ref{local-cohomology-lemma-change-completion}.
Supposons ensuite que (c)(i) soit satisfaite. Comme $A$ est universellement caténaire
(Complexes dualisants, lemme \ref{dualizing-lemma-universally-catenary}),
on voit que toute composante irréductible de $\Spec(A')$ non contenue dans $V(f)$
est de dimension $\geq 3$ ; voir
Compléments d'algèbre, proposition \ref{more-algebra-proposition-ratliff}.
Comme $A \to A'$ est plat à fibres de Gorenstein,
le fait que $A_f$ vérifie $(S_2)$ implique que $A'_f$ vérifie $(S_2)$.
Références utilisées :
Complexes dualisants, section \ref{dualizing-section-formal-fibres},
Compléments d'algèbre, section \ref{more-algebra-section-properties-formal-fibres},
et Algèbre, lemme \ref{algebra-lemma-Sk-goes-up}.
Ainsi, d'après le lemme \ref{pione-lemma-fully-faithful-henselian-completion},
on peut supposer que $A$ est un anneau local noethérien complet.

\medskip\noindent
Supposons, en plus des autres hypothèses, que $A$ soit un anneau local complet.
D'après le lemme \ref{pione-lemma-restriction-fully-faithful}, le résultat découle de
Géométrie algébrique et formelle, lemme
\ref{algebraization-lemma-fully-faithful-simple-two}.
\end{proof}

\begin{lemma}
\label{pione-lemma-fully-faithful-minimal}
\begin{reference}
\cite[corollaire 1.11]{Bhatt-local}
\end{reference}
Dans la situation \ref{pione-situation-local-lefschetz}, supposons que
\begin{enumerate}
\item $H^1_\mathfrak m(A)$ et $H^2_\mathfrak m(A)$ sont
annulés par une puissance de $f$, et
\item $A$ est hensélien ou, plus généralement, $(A, (f))$ est un couple hensélien.
\end{enumerate}
Alors le foncteur de restriction
$\textit{F\'Et}_U \longrightarrow \textit{F\'Et}_{U_0}$
est pleinement fidèle.
\end{lemma}

\begin{proof}
D'après le lemme \ref{pione-lemma-fully-faithful-henselian-completion},
on peut supposer que $A$ est un anneau local noethérien complet.
(Les hypothèses se transportent ; utiliser
Complexes dualisants, lemme \ref{dualizing-lemma-torsion-change-rings}.)
D'après le lemme \ref{pione-lemma-restriction-fully-faithful},
le résultat découle du
lemme de Géométrie algébrique et formelle
\ref{algebraization-lemma-fully-faithful-alternative}.
\end{proof}

\begin{lemma}
\label{pione-lemma-fully-faithful}
Dans la situation \ref{pione-situation-local-lefschetz}, supposons que $A$ soit de profondeur $\geq 3$
et que $A$ soit hensélien ou, plus généralement, que $(A, (f))$ soit un couple hensélien. Alors
le foncteur de restriction
$\textit{F\'Et}_U \to \textit{F\'Et}_{U_0}$
est pleinement fidèle.
\end{lemma}

\begin{proof}
L'hypothèse sur la profondeur impose
$H^1_\mathfrak m(A) = H^2_\mathfrak m(A) = 0$ ; voir
Complexes dualisants, lemme \ref{dualizing-lemma-depth}.
Le lemme \ref{pione-lemma-fully-faithful-minimal} s'applique donc.
\end{proof}



























\section{Pureté dans le cas local, I}
\label{pione-section-local-purity}

\noindent
Soit $(A, \mathfrak m)$ un anneau local noethérien. Posons $X = \Spec(A)$
et soit $U = X \setminus \{\mathfrak m\}$ le spectre épointé.
Nous disons que {\it la pureté vaut pour $(A, \mathfrak m)$} si le foncteur de restriction
$$
\textit{F\'Et}_X \longrightarrow \textit{F\'Et}_U
$$
est essentiellement surjectif. Dans cette section, nous cherchons à comprendre comment la
question se transforme lorsque l'on passe de $X$ à une hypersurface $X_0$ de $X$ ;
autrement dit, nous étudions une sorte de propriété de Lefschetz locale pour les
groupes fondamentaux des spectres épointés.
Ces résultats seront utiles pour procéder par récurrence sur la dimension
dans les démonstrations de nos principaux résultats sur la pureté locale, à savoir le
lemme \ref{pione-lemma-local-purity},
la proposition \ref{pione-proposition-purity-complete-intersection} et la
proposition \ref{pione-proposition-purity-smooth-over-depth2}.

\begin{lemma}
\label{pione-lemma-sections-over-punctured-spec}
Soit $(A, \mathfrak m)$ un anneau local noethérien. Posons $X = \Spec(A)$
et soit $U = X \setminus \{\mathfrak m\}$.
Soit $\pi : Y \to X$ un morphisme fini tel que
$\text{depth}(\mathcal{O}_{Y, y}) \geq 2$ pour tout point fermé
$y \in Y$.
Alors $Y$ est le spectre de $B = \mathcal{O}_Y(\pi^{-1}(U))$.
\end{lemma}

\begin{proof}
Posons $V = \pi^{-1}(U)$ et notons $\pi' : V \to U$ la restriction de $\pi$.
Considérons le morphisme de $\mathcal{O}_X$-modules
$$
\pi_*\mathcal{O}_Y \longrightarrow j_*\pi'_*\mathcal{O}_V
$$
où $j : U \to X$ est le morphisme d'inclusion. Nous affirmons que le
lemme de Diviseurs \ref{divisors-lemma-depth-2-hartog}
s'applique à ce morphisme. Si tel est le cas, alors $B = \Gamma(Y, \mathcal{O}_Y)$
et le lemme en résulte. Soit $x \in X$ le point fermé.
Il suffit de montrer que
$\text{depth}((\pi_*\mathcal{O}_Y)_x) \geq 2$.
Soient $y_1, \ldots, y_n \in Y$ les points au-dessus de $x$.
D'après Algèbre, lemme \ref{algebra-lemma-depth-goes-down-finite},
il suffit de montrer que
$\text{depth}(\mathcal{O}_{Y, y_i}) \geq 2$ pour $i = 1, \ldots, n$.
Comme c'est précisément l'hypothèse du lemme, la démonstration est achevée.
\end{proof}

\begin{lemma}
\label{pione-lemma-reformulate-purity}
Soit $(A, \mathfrak m)$ un anneau local noethérien. Posons $X = \Spec(A)$
et soit $U = X \setminus \{\mathfrak m\}$.
Soit $V$ fini étale
sur $U$. Supposons $A$ de profondeur $\geq 2$. Les assertions suivantes sont équivalentes :
\begin{enumerate}
\item $V = Y \times_X U$ pour un certain morphisme fini étale $Y \to X$,
\item $B = \Gamma(V, \mathcal{O}_V)$ est fini étale sur $A$.
\end{enumerate}
\end{lemma}

\begin{proof}
Notons $\pi : V \to U$ le morphisme fini étale donné.
Supposons qu'il existe $Y$ comme dans (1). Soit $x \in X$ le point
correspondant à $\mathfrak m$.
Soit $y \in Y$ un point au-dessus de $x$. Nous affirmons que
$\text{depth}(\mathcal{O}_{Y, y}) \geq 2$.
Cela est vrai parce que $Y \to X$ est étale et que
$A = \mathcal{O}_{X, x}$ et $\mathcal{O}_{Y, y}$ ont donc
la même profondeur (Algèbre, lemme \ref{algebra-lemma-apply-grothendieck}).
Le lemme \ref{pione-lemma-sections-over-punctured-spec}
s'applique donc et $Y = \Spec(B)$.

\medskip\noindent
L'implication (2) $\Rightarrow$ (1) est plus facile et nous en
omettons les détails.
\end{proof}

\begin{lemma}
\label{pione-lemma-reformulate-purity-normal}
Soit $(A, \mathfrak m)$ un anneau local noethérien. Posons $X = \Spec(A)$
et soit $U = X \setminus \{\mathfrak m\}$. Supposons $A$ normal
de dimension $\geq 2$. Le foncteur
$$
\textit{F\'Et}_U \longrightarrow
\left\{
\begin{matrix}
A\text{-algèbres finies normales }B\text{ telles que} \\
\Spec(B) \to X\text{ soit étale au-dessus de }U
\end{matrix}
\right\},
\quad
V \longmapsto \Gamma(V, \mathcal{O}_V)
$$
est une équivalence. De plus, $V = Y \times_X U$ pour un certain morphisme $Y \to X$
fini étale si et seulement si $B = \Gamma(V, \mathcal{O}_V)$
est fini étale sur $A$.
\end{lemma}

\begin{proof}
Remarquons que $\text{depth}(A) \geq 2$ puisque $A$ est normal
(critère de normalité de Serre, Algèbre, lemme
\ref{algebra-lemma-criterion-normal}).
La dernière assertion découle donc du lemme \ref{pione-lemma-reformulate-purity}.
Étant donné $\pi : V \to U$ fini étale, posons $B = \Gamma(V, \mathcal{O}_V)$.
Si nous montrons que $B$ est normal et fini sur $A$, nous
obtenons le foncteur affiché. Comme il existe un foncteur quasi-inverse
évident, c'est aussi tout ce qu'il reste à montrer.

\medskip\noindent
Comme $A$ est normal, le schéma $V$ est normal
(Descente, lemme \ref{descent-lemma-normal-local-smooth}).
Par conséquent, $V$ est une réunion disjointe finie de schémas intègres
(Propriétés, lemme \ref{properties-lemma-normal-Noetherian}).
On peut donc supposer $V$ intègre.
Dans ce cas, le corps des fonctions $L$ de $V$
(Morphismes, section \ref{morphisms-section-rational-maps})
est une extension finie séparable du corps des fractions de $A$
(car on l'obtient en considérant la fibre générique
de $V \to U$ et en utilisant Morphismes, lemme
\ref{morphisms-lemma-etale-over-field}).
D'après Algèbre, lemme
\ref{algebra-lemma-Noetherian-normal-domain-finite-separable-extension},
la clôture intégrale $B' \subset L$ de $A$ dans $L$ est finie sur $A$.
D'après Compléments d'algèbre, lemme \ref{more-algebra-lemma-integral-closure-reflexive},
$B'$ est un $A$-module réflexif, ce qui implique à son tour
que $\text{depth}_A(B') \geq 2$ d'après
Compléments d'algèbre, lemme \ref{more-algebra-lemma-reflexive-over-normal}.

\medskip\noindent
Soit $f \in \mathfrak m$. Alors $B_f = \Gamma(V \times_U D(f), \mathcal{O}_V)$
(Propriétés, lemme \ref{properties-lemma-invert-f-sections}).
On a donc $B'_f = B_f$, puisque $B_f$ est normal (voir ci-dessus),
fini sur $A_f$ et de corps des fractions $L$.
Il s'ensuit que $V = \Spec(B') \times_X U$.
Nous concluons alors que $B = B'$ grâce au
lemme \ref{pione-lemma-sections-over-punctured-spec}
appliqué à $\Spec(B') \to X$.
Ce lemme s'applique parce que les localisés $B'_{\mathfrak m'}$
de $B'$ aux idéaux maximaux $\mathfrak m' \subset B'$ situés au-dessus de
$\mathfrak m$ sont de profondeur $\geq 2$ d'après
Algèbre, lemme \ref{algebra-lemma-depth-goes-down-finite},
et la remarque sur la profondeur du paragraphe précédent.
\end{proof}

\begin{lemma}
\label{pione-lemma-purity-and-completion}
Soit $(A, \mathfrak m)$ un anneau local noethérien. Posons $X = \Spec(A)$
et soit $U = X \setminus \{\mathfrak m\}$.
Soit $V$ fini étale sur $U$.
Soit $A^\wedge$ le complété $\mathfrak m$-adique de $A$,
soit $X' = \Spec(A^\wedge)$, et soient $U'$ et $V'$ les changements de base de
$U$ et $V$ à $X'$. Les assertions suivantes sont équivalentes :
\begin{enumerate}
\item $V = Y \times_X U$ pour un certain morphisme fini étale $Y \to X$, et
\item $V' = Y' \times_{X'} U'$ pour un certain morphisme fini étale $Y' \to X'$.
\end{enumerate}
\end{lemma}

\begin{proof}
L'implication (1) $\Rightarrow$ (2) s'obtient en effectuant le changement de base
d'une solution $Y \to X$. Soit $Y' \to X'$ comme dans (2).
D'après le lemme \ref{pione-lemma-fill-in-missing}, il existe un morphisme fini $Y \to X$
tel que $V = Y \times_X U$ et $Y' = Y \times_X X'$.
Par descente, le morphisme $Y \to X$ est fini étale
(Algèbre, lemmes \ref{algebra-lemma-descend-properties-modules} et
\ref{algebra-lemma-etale}). Cela achève la démonstration.
\end{proof}

\noindent
L'intérêt des deux lemmes suivants est que leurs hypothèses n'imposent pas à
$A$ d'être de profondeur $\geq 3$. Par exemple, si $A$ est un anneau local complet
normal et intègre de dimension $\geq 3$ et si $f \in \mathfrak m$ est non nul,
alors les hypothèses sont satisfaites.

\begin{lemma}
\label{pione-lemma-lift-simple}
Dans la situation \ref{pione-situation-local-lefschetz}, soit $V$ fini
étale sur $U$. Supposons que
\begin{enumerate}
\item[(a)] $A$ possède un complexe dualisant,
\item[(b)] le couple $(A, (f))$ soit hensélien,
\item[(c)] l'une des conditions suivantes soit satisfaite :
\begin{enumerate}
\item[(i)] $A_f$ vérifie $(S_2)$ et toute composante irréductible de $X$
non contenue dans $X_0$ est de dimension $\geq 3$, ou bien
\item[(ii)] pour tout idéal premier $\mathfrak p \subset A$, $f \not \in \mathfrak p$
on ait $\text{depth}(A_\mathfrak p) + \dim(A/\mathfrak p) > 2$.
\end{enumerate}
\item[(d)] $V_0 = V \times_U U_0$ soit égal à $Y_0 \times_{X_0} U_0$
pour un certain morphisme fini étale $Y_0 \to X_0$.
\end{enumerate}
Alors $V = Y \times_X U$ pour un certain morphisme fini étale $Y \to X$.
\end{lemma}

\begin{proof}
On se ramène au cas complet à l'aide du lemme \ref{pione-lemma-purity-and-completion}.
(Les hypothèses se transportent ; voir la démonstration du
lemme \ref{pione-lemma-fully-faithful-simple}.)

\medskip\noindent
Dans le cas complet, on peut relever $Y_0 \to X_0$ en un morphisme
fini étale $Y \to X$ d'après le
lemme de Compléments d'algèbre \ref{more-algebra-lemma-finite-etale-equivalence} ;
remarquons que $(A, fA)$ est un couple hensélien d'après le
lemme de Compléments d'algèbre \ref{more-algebra-lemma-complete-henselian}.
On peut alors utiliser le lemme \ref{pione-lemma-fully-faithful-simple}
pour voir que $V$ est isomorphe à $Y \times_X U$ et
la démonstration est achevée.
\end{proof}

\begin{lemma}
\label{pione-lemma-lift-purity-general}
Dans la situation \ref{pione-situation-local-lefschetz},
soit $V$ fini étale sur $U$. Supposons que
\begin{enumerate}
\item $H^1_\mathfrak m(A)$ et $H^2_\mathfrak m(A)$
soient annulés par une puissance de $f$,
\item $V_0 = V \times_U U_0$ soit égal à $Y_0 \times_{X_0} U_0$
pour un certain morphisme fini étale $Y_0 \to X_0$.
\end{enumerate}
Alors $V = Y \times_X U$ pour un certain morphisme fini étale $Y \to X$.
\end{lemma}

\begin{proof}
On se ramène au cas complet à l'aide du lemme \ref{pione-lemma-purity-and-completion}.
(Les hypothèses se transportent ; utiliser Complexes dualisants, lemme
\ref{dualizing-lemma-torsion-change-rings}.)

\medskip\noindent
Dans le cas complet, on peut relever $Y_0 \to X_0$ en un morphisme
fini étale $Y \to X$ d'après le
lemme de Compléments d'algèbre \ref{more-algebra-lemma-finite-etale-equivalence} ;
remarquons que $(A, fA)$ est un couple hensélien d'après le
lemme de Compléments d'algèbre \ref{more-algebra-lemma-complete-henselian}.
On peut alors utiliser le lemme \ref{pione-lemma-fully-faithful-minimal}
pour voir que $V$ est isomorphe à $Y \times_X U$ et
la démonstration est achevée.
\end{proof}

\begin{lemma}
\label{pione-lemma-lift-purity}
Dans la situation \ref{pione-situation-local-lefschetz},
soit $V$ fini étale sur $U$. Supposons que
\begin{enumerate}
\item $A$ soit de profondeur $\geq 3$,
\item $V_0 = V \times_U U_0$ soit égal à $Y_0 \times_{X_0} U_0$
pour un certain morphisme fini étale $Y_0 \to X_0$.
\end{enumerate}
Alors $V = Y \times_X U$ pour un certain morphisme fini étale $Y \to X$.
\end{lemma}

\begin{proof}
L'hypothèse sur la profondeur impose
$H^1_\mathfrak m(A) = H^2_\mathfrak m(A) = 0$ ; voir
Complexes dualisants, lemme \ref{dualizing-lemma-depth}.
Le lemme \ref{pione-lemma-lift-purity-general} s'applique donc.
\end{proof}










\section{Pureté du lieu de branchement}
\label{pione-section-purity}

\noindent
Nous utiliserons le discriminant d'un morphisme fini localement libre. Voir
Discriminants, section \ref{discriminant-section-discriminant}.

\begin{lemma}
\label{pione-lemma-find-point-codim-1}
Soit $(A, \mathfrak m)$ un anneau local noethérien tel que $\dim(A) \geq 1$.
Soit $f \in \mathfrak m$. Alors il existe $\mathfrak p \in V(f)$ tel que
$\dim(A_\mathfrak p) = 1$.
\end{lemma}

\begin{proof}
Raisonnons par récurrence sur $\dim(A)$. Si $\dim(A) = 1$, alors $\mathfrak p = \mathfrak m$
convient. Si $\dim(A) > 1$, soit $Z \subset \Spec(A)$ une composante irréductible
de dimension $> 1$. Alors $V(f) \cap Z$ est de dimension $> 0$
(Algèbre, lemme \ref{algebra-lemma-one-equation}). Choisissons un idéal premier
$\mathfrak q \in V(f) \cap Z$, $\mathfrak q \not = \mathfrak m$,
correspondant à un point fermé du spectre épointé de $A$ ;
c'est possible d'après
Propriétés, lemme \ref{properties-lemma-complement-closed-point-Jacobson}.
Alors $\mathfrak q$ n'est pas le point générique de $Z$. Par conséquent,
$0 < \dim(A_\mathfrak q) < \dim(A)$ et $f \in \mathfrak q A_\mathfrak q$.
Par récurrence sur la dimension, on peut trouver
$f \in \mathfrak p \subset A_\mathfrak q$ tel que
$\dim((A_\mathfrak q)_\mathfrak p) = 1$.
Alors $\mathfrak p \cap A$ convient.
\end{proof}

\begin{lemma}
\label{pione-lemma-ramification-quasi-finite-flat}
Soit $f : X \to Y$ un morphisme de schémas localement noethériens.
Soit $x \in X$. Supposons que
\begin{enumerate}
\item $f$ soit plat,
\item $f$ soit quasi-fini en $x$,
\item $x$ ne soit pas le point générique d'une composante irréductible de $X$,
\item pour les spécialisations $x' \leadsto x$ telles que
$\dim(\mathcal{O}_{X, x'}) = 1$, le morphisme $f$ soit non ramifié en $x'$.
\end{enumerate}
Alors $f$ est étale en $x$.
\end{lemma}

\begin{proof}
Remarquons que l'ensemble des points où $f$ est non ramifié coïncide avec
l'ensemble des points où $f$ est étale, et que cet ensemble est ouvert.
Voir Morphismes, définitions \ref{morphisms-definition-unramified}
et \ref{morphisms-definition-etale}, ainsi que le
lemme \ref{morphisms-lemma-flat-unramified-etale}.
Pour vérifier que $f$ est étale en $x$, on peut travailler localement pour la topologie étale
sur la base et sur la
source (Descente, lemmes \ref{descent-lemma-descending-property-etale} et
\ref{descent-lemma-etale-etale-local-source}).
On peut donc appliquer Compléments sur les morphismes, lemme
\ref{more-morphisms-lemma-etale-makes-quasi-finite-finite-at-point},
et supposer que $f : X \to Y$ est fini et que $x$ est l'unique
point de $X$ situé au-dessus de $y = f(x)$.
Il s'ensuit que $f$ est fini localement libre
(Morphismes, lemme \ref{morphisms-lemma-finite-flat}).

\medskip\noindent
Supposons $f$ fini localement libre et $x$ l'unique point de
$X$ situé au-dessus de $y = f(x)$. D'après
Discriminants, lemme \ref{discriminant-lemma-discriminant},
il existe un sous-schéma fermé localement principal $D_f \subset Y$
tel que $y' \in D_f$ si et seulement s'il existe $x' \in X$
avec $f(x') = y'$ et $f$ ramifié en $x'$. Il faut donc montrer
que $y \not \in D_f$. Supposons $y \in D_f$ et cherchons une contradiction.

\medskip\noindent
D'après la condition (3), on a $\dim(\mathcal{O}_{X, x}) \geq 1$.
On a $\dim(\mathcal{O}_{X, x}) = \dim(\mathcal{O}_{Y, y})$ d'après
Algèbre, lemme \ref{algebra-lemma-dimension-base-fibre-equals-total}.
D'après le lemme \ref{pione-lemma-find-point-codim-1},
on peut trouver $y' \in D_f$ se spécialisant en $y$
tel que $\dim(\mathcal{O}_{Y, y'}) = 1$.
Choisissons $x' \in X$ tel que $f(x') = y'$ et où $f$ est ramifié. Comme $f$
est fini, il est fermé, et par conséquent $x' \leadsto x$.
On a $\dim(\mathcal{O}_{X, x'}) = \dim(\mathcal{O}_{Y, y'}) = 1$
comme précédemment. Cela contredit la propriété (4).
\end{proof}

\begin{lemma}
\label{pione-lemma-local-purity}
Soit $(A, \mathfrak m)$ un anneau local régulier de dimension $d \geq 2$.
Posons $X = \Spec(A)$ et $U = X \setminus \{\mathfrak m\}$. Alors
\begin{enumerate}
\item le foncteur $\textit{F\'Et}_X \to \textit{F\'Et}_U$
est essentiellement surjectif, c'est-à-dire que la pureté vaut pour $A$,
\item tout morphisme fini $A \to B$ avec $B$ normal qui
induit un morphisme fini étale sur les spectres épointés est étale.
\end{enumerate}
\end{lemma}

\begin{proof}
Rappelons qu'un anneau local régulier est normal d'après
Algèbre, lemme \ref{algebra-lemma-regular-normal}.
Par conséquent, (1) et (2) sont équivalents d'après le
lemme \ref{pione-lemma-reformulate-purity-normal}.
Démontrons le lemme par récurrence sur $d$.

\medskip\noindent
Le cas $d = 2$. Dans ce cas, $A \to B$ est plat.
En effet, la descente des idéaux premiers vaut pour $A \to B$ d'après
Algèbre, proposition \ref{algebra-proposition-going-down-normal-integral}.
Alors $\dim(B_{\mathfrak m'}) = 2$ pour tout idéal maximal
$\mathfrak m' \subset B$ d'après
Algèbre, lemme \ref{algebra-lemma-dimension-base-fibre-equals-total}.
Puis $B_{\mathfrak m'}$ est un anneau de Cohen--Macaulay d'après
Algèbre, lemme \ref{algebra-lemma-criterion-normal}.
Par conséquent, et c'est l'étape importante, le
lemme d'Algèbre \ref{algebra-lemma-CM-over-regular-flat}
s'applique et montre que $A \to B_{\mathfrak m'}$ est plat.
Le lemme d'Algèbre \ref{algebra-lemma-flat-localization}
montre alors que $A \to B$ est plat. On peut donc appliquer le
lemme \ref{pione-lemma-ramification-quasi-finite-flat}
(ou raisonner directement à l'aide du plus simple
lemme de Discriminants \ref{discriminant-lemma-discriminant})
pour voir que $A \to B$ est étale.

\medskip\noindent
Le cas $d \geq 3$. Soit $V \to U$ fini étale.
Soit $f \in \mathfrak m_A$, $f \not \in \mathfrak m_A^2$.
Alors $A/fA$ est un anneau local régulier de dimension $d - 1 \geq 2$ ; voir
Algèbre, lemme \ref{algebra-lemma-regular-ring-CM}.
Soit $U_0$ le spectre épointé de $A/fA$ et posons
$V_0 = V \times_U U_0$.
D'après le lemme \ref{pione-lemma-lift-purity},
il suffit de montrer que $V_0$ appartient à l'image essentielle
de $\textit{F\'Et}_{\Spec(A/fA)} \to \textit{F\'Et}_{U_0}$.
Cela résulte de l'hypothèse de récurrence.
\end{proof}

\begin{lemma}[Pureté du lieu de branchement]
\label{pione-lemma-purity}
\begin{reference}
\cite{Nagata-Purity} et \cite[exposé X, théorème 3.1]{SGA1}
\end{reference}
\begin{history}
Ce résultat fut d'abord énoncé et démontré par Zariski sous
forme géométrique dans \cite{Zariski-Purity}.
La généralisation aux corps de base non parfaits, due à Nagata,
fut publiée comme article suivant dans le même volume des
Comptes rendus de l'Académie nationale des sciences des États-Unis d'Amérique,
dans \cite{Nagata-Remarks-Purity}. L'année suivante, Nagata
démontra le résultat pour les anneaux locaux noethériens dans \cite{Nagata-Purity}.
Sa démonstration utilise un résultat de Chow qui est un théorème de Bertini pour les
anneaux locaux complets et intègres ; voir \cite{Chow-Bertini} ;
l'histoire des théorèmes de Bertini est étudiée dans
l'article historique de Kleiman \cite{Kleiman-Bertini}.
Quelques années plus tard, Auslander trouva une démonstration entièrement différente ;
voir \cite{Auslander-Purity}.
\end{history}
Soit $f : X \to Y$ un morphisme de schémas localement noethériens.
Soit $x \in X$ et posons $y = f(x)$. Supposons que
\begin{enumerate}
\item $\mathcal{O}_{X, x}$ soit normal,
\item $\mathcal{O}_{Y, y}$ soit régulier,
\item $f$ soit quasi-fini en $x$,
\item $\dim(\mathcal{O}_{X, x}) = \dim(\mathcal{O}_{Y, y}) \geq 1$
\item pour les spécialisations $x' \leadsto x$ telles que
$\dim(\mathcal{O}_{X, x'}) = 1$, le morphisme $f$ soit non ramifié en $x'$.
\end{enumerate}
Alors $f$ est étale en $x$.
\end{lemma}

\begin{proof}
Nous allons démontrer le lemme par récurrence sur
$d = \dim(\mathcal{O}_{X, x}) = \dim(\mathcal{O}_{Y, y})$.

\medskip\noindent
Un cas sans intérêt est celui où $d = 1$.
Dans ce cas, nous supposons que $f$ est non ramifié en $x$
et que $\mathcal{O}_{Y, y}$ est un anneau de valuation discrète
(Algèbre, lemme \ref{algebra-lemma-characterize-dvr}).
Alors $\mathcal{O}_{X, x}$ est plat sur $\mathcal{O}_{Y, y}$
(sinon le morphisme ne serait pas quasi-fini en $x$),
et l'on voit que $f$ est plat en $x$. Comme plat $+$
non ramifié implique étale, on conclut (certains détails sont omis).

\medskip\noindent
Le cas $d \geq 2$. On va utiliser la récurrence sur $d$ pour se ramener
au cas étudié dans le lemme \ref{pione-lemma-local-purity}.
Pour vérifier que $f$ est étale en $x$, on peut travailler localement pour la topologie étale
sur la base et sur la source
(Descente, lemmes \ref{descent-lemma-descending-property-etale} et
\ref{descent-lemma-etale-etale-local-source}).
On peut donc appliquer Compléments sur les morphismes, lemme
\ref{more-morphisms-lemma-etale-makes-quasi-finite-finite-at-point},
et supposer que $f : X \to Y$ est fini et que $x$ est l'unique
point de $X$ situé au-dessus de $y$. On utilise ici le fait que les extensions étales entre
anneaux locaux ne changent ni la dimension, ni la normalité, ni la régularité ; voir
Compléments d'algèbre, section \ref{more-algebra-section-permanence-etale},
et
Morphismes étales, section \ref{etale-section-properties-permanence}.

\medskip\noindent
Ensuite, on peut effectuer le changement de base par $\Spec(\mathcal{O}_{Y, y})$
et supposer que $Y$ est le spectre d'un anneau local régulier.
Il s'ensuit que $X = \Spec(\mathcal{O}_{X, x})$, puisque
tout point de $X$ se spécialise nécessairement en $x$.

\medskip\noindent
Le morphisme d'anneaux $\mathcal{O}_{Y, y} \to \mathcal{O}_{X, x}$ est
fini et nécessairement injectif (par égalité des dimensions).
On conclut que la descente des idéaux premiers vaut pour
$\mathcal{O}_{Y, y} \to \mathcal{O}_{X, x}$ d'après
Algèbre, proposition \ref{algebra-proposition-going-down-normal-integral}
(et le fait qu'un anneau régulier est normal d'après
Algèbre, lemme \ref{algebra-lemma-regular-normal}).
Choisissons $x' \in X$, $x' \not = x$, d'image $y' = f(x')$.
Alors $\mathcal{O}_{X, x'}$ est normal comme localisé
d'un anneau normal intègre. De même, $\mathcal{O}_{Y, y'}$ est
régulier (voir Algèbre, lemme
\ref{algebra-lemma-localization-of-regular-local-is-regular}).
On a $\dim(\mathcal{O}_{X, x'}) = \dim(\mathcal{O}_{Y, y'})$ d'après
Algèbre, lemme \ref{algebra-lemma-dimension-base-fibre-equals-total}
(on a vérifié ci-dessus la descente des idéaux premiers).
Bien entendu, ces dimensions sont strictement inférieures à $d$, puisque $x' \not = x$,
et l'hypothèse de récurrence sur $d$ montre que $f$ est étale en $x'$.

\medskip\noindent
On arrive donc à la situation suivante : on a un morphisme
local fini $A \to B$ d'anneaux locaux noethériens
de dimension $d \geq 2$, avec $A$ régulier et $B$ normal, qui
induit un morphisme fini étale $V \to U$ sur les spectres épointés.
Notre but est de montrer que $A \to B$ est étale.
Cela résulte du lemme \ref{pione-lemma-local-purity},
et la démonstration est achevée.
\end{proof}

\noindent
Le lemme suivant est parfois utile pour déterminer le plus grand
ouvert sur lequel un morphisme fini étale se prolonge.

\begin{lemma}
\label{pione-lemma-extend-S2}
Soit $j : U \to X$ une immersion ouverte de schémas localement noethériens
telle que $\text{depth}(\mathcal{O}_{X, x}) \geq 2$ pour $x \not \in U$.
Soit $\pi : V \to U$ fini étale. Alors
\begin{enumerate}
\item $\mathcal{B} = j_*\pi_*\mathcal{O}_V$ est une algèbre cohérente
réflexive sur $\mathcal{O}_X$ ; posons $Y = \underline{\Spec}_X(\mathcal{B})$,
\item $Y \to X$ est l'unique morphisme fini tel que
$V = Y \times_X U$ et $\text{depth}(\mathcal{O}_{Y, y}) \geq 2$
pour $y \in Y \setminus V$,
\item $Y \to X$ est étale en $y$ si et seulement si $Y \to X$ est plat en $y$, et
\item $Y \to X$ est étale si et seulement si $\mathcal{B}$
est fini localement libre comme $\mathcal{O}_X$-module.
\end{enumerate}
De plus, (a) la construction de $\mathcal{B}$ et de $Y \to X$ commute
au changement de base par les morphismes plats $X' \to X$ de schémas
localement noethériens, et (b), si $V' \to U'$ est un morphisme fini étale avec
$U \subset U' \subset X$ ouvert qui se restreint à $V \to U$ au-dessus de $U$,
alors il existe un unique isomorphisme $Y \times_X U' = V'$ au-dessus de $U'$.
\end{lemma}

\begin{proof}
Remarquons que $\pi_*\mathcal{O}_V$ est un module fini localement libre
sur $\mathcal{O}_U$, donc en particulier réflexif.
D'après Diviseurs, lemme \ref{divisors-lemma-reflexive-S2-extend},
le module $j_*\pi_*\mathcal{O}_V$ est l'unique
module cohérent réflexif sur $X$ dont la restriction à $U$
est $\pi_*\mathcal{O}_V$. Cela démontre (1).

\medskip\noindent
Par construction, $Y \times_X U = V$.
Comme $\mathcal{B}$ est cohérente, on voit que $Y \to X$ est fini.
On a $\text{depth}(\mathcal{B}_x) \geq 2$ pour $x \in X \setminus U$
d'après Diviseurs, lemme \ref{divisors-lemma-reflexive-S2}.
Par conséquent, $\text{depth}(\mathcal{O}_{Y, y}) \geq 2$ pour $y \in Y \setminus V$
d'après Algèbre, lemme \ref{algebra-lemma-depth-goes-down-finite}.
Réciproquement, supposons que $\pi' : Y' \to X$ soit un morphisme fini tel que
$V = Y' \times_X U$ et $\text{depth}(\mathcal{O}_{Y', y'}) \geq 2$
pour $y' \in Y' \setminus V$. Alors $\pi'_*\mathcal{O}_{Y'}$
a pour restriction $\pi_*\mathcal{O}_V$ sur $U$ et vérifie
$\text{depth}((\pi'_*\mathcal{O}_{Y'})_x) \geq 2$ pour
$x \in X \setminus U$ d'après le
lemme d'Algèbre \ref{algebra-lemma-depth-goes-down-finite}.
Alors $\pi'_*\mathcal{O}_{Y'}$ est canoniquement isomorphe
à $j_*\pi_*\mathcal{O}_V$, par exemple d'après
Diviseurs, lemme \ref{divisors-lemma-depth-2-hartog}.
Cela démontre (2).

\medskip\noindent
Si $Y \to X$ est étale en $y$, alors $Y \to X$ est plat en $y$.
Réciproquement, supposons que $Y \to X$ soit plat en $y$.
Si $y \in V$, alors $Y \to X$ est étale en $y$.
Si $y \not \in V$, vérifions que (1), (2), (3) et (4) du
lemme \ref{pione-lemma-ramification-quasi-finite-flat} sont satisfaits
pour voir que $Y \to X$ est étale en $y$. Les parties (1) et (2)
sont claires, de même que (3), puisque $\text{depth}(\mathcal{O}_{Y, y}) \geq 2$.
Si $y' \leadsto y$ est une spécialisation et $\dim(\mathcal{O}_{Y, y'}) = 1$,
alors $y' \in V$, car sinon la profondeur de cet anneau local
serait au moins égale à $2$, ce qui contredirait le
lemme d'Algèbre \ref{algebra-lemma-bound-depth}.
Ainsi, $Y \to X$ est étale en $y'$, et l'on conclut que (4) du
lemme \ref{pione-lemma-ramification-quasi-finite-flat} est également satisfait.
Cela achève la démonstration de (3).

\medskip\noindent
La partie (4) résulte de (3) et du fait que $((Y \to X)_*\mathcal{O}_Y)_x$
est un $\mathcal{O}_{X, x}$-module plat si et seulement si $\mathcal{O}_{Y, y}$
est un $\mathcal{O}_{X, x}$-module plat pour tout $y \in Y$ au-dessus de $x$ ; voir
Algèbre, lemme \ref{algebra-lemma-flat-localization}. On utilise aussi ici
le fait qu'un module fini plat sur un anneau noethérien est fini localement
libre ; voir Algèbre, lemme \ref{algebra-lemma-finite-projective}
(et
Algèbre, lemme
\ref{algebra-lemma-Noetherian-finite-type-is-finite-presentation}).

\medskip\noindent
Quant aux dernières assertions du lemme, la partie (a) résulte du
changement de base plat ; voir
Cohomologie des schémas, lemme \ref{coherent-lemma-flat-base-change-cohomology},
et la partie (b) découle de l'unicité dans (2), appliquée à la restriction
$Y \times_X U'$.
\end{proof}

\begin{lemma}
\label{pione-lemma-extend-pure}
Soit $j : U \to X$ une immersion ouverte de schémas noethériens
telle que la pureté vaille pour $\mathcal{O}_{X, x}$ pour tout $x \not \in U$.
Alors
$$
\textit{F\'Et}_X \longrightarrow \textit{F\'Et}_U
$$
est essentiellement surjectif.
\end{lemma}

\begin{proof}
Soit $V \to U$ un morphisme fini étale. Par récurrence noethérienne,
il suffit de prolonger $V \to U$ en un morphisme fini étale
au-dessus d'un ouvert strictement plus grand de $X$.
Soit $x \in X \setminus U$ le point générique d'une
composante irréductible de $X \setminus U$.
Alors l'image inverse $U_x$ de $U$ dans $\Spec(\mathcal{O}_{X, x})$
est le spectre épointé de $\mathcal{O}_{X, x}$.
Par hypothèse, $V_x = V \times_U U_x$ est la restriction
d'un morphisme fini étale $Y_x \to \Spec(\mathcal{O}_{X, x})$
à $U_x$.
D'après Limites, lemme \ref{limits-lemma-glueing-near-point},
il existe un sous-schéma ouvert $U \subset U' \subset X$
contenant $x$ et un morphisme de présentation finie $V' \to U'$
dont la restriction à $U$ redonne $V \to U$ et
dont la restriction à $\Spec(\mathcal{O}_{X, x})$ redonne $Y_x$.
Enfin, le morphisme $V' \to U'$ est fini étale
après avoir éventuellement remplacé $U'$ par un ouvert plus petit, d'après
Limites, lemme \ref{limits-lemma-glueing-near-point-properties}.
\end{proof}
















\section{Revêtements finis étales des spectres épointés, II}
\label{pione-section-pi1-punctured-spec-II}

\noindent
Dans cette section, nous démontrons quelques variantes des résultats étudiés
dans la section \ref{pione-section-pi1-punctured-spec}. Supposons
donné un anneau local noethérien $(A, \mathfrak m)$ et un élément $f \in \mathfrak m$.
Posons $X = \Spec(A)$ et $X_0 = \Spec(A/fA)$, et
posons $U = X \setminus \{\mathfrak m\}$ et
$U_0 = X_0 \setminus \{\mathfrak m\}$, les spectres épointés de
$A$ et $A/fA$. Tout ceci est exactement comme dans la
situation \ref{pione-situation-local-lefschetz}.
La différence est que nous considérerons le foncteur de restriction
$$
\colim_{U_0 \subset U' \subset U\text{ ouvert}} \textit{F\'Et}_{U'}
\longrightarrow
\textit{F\'Et}_{U_0}
$$
Autrement dit, nous ne chercherons pas à relever les revêtements finis étales
de $U_0$ à tout $U$, mais seulement à un voisinage ouvert
$U'$ de $U_0$ dans $U$.

\begin{lemma}
\label{pione-lemma-faithful-general}
Dans la situation \ref{pione-situation-local-lefschetz}, soit $U' \subset U$
un ouvert contenant $U_0$. Supposons que, pour tout idéal premier minimal $\mathfrak p \subset A$
tel que $\mathfrak p \in U'$ et $\mathfrak p \not \in U_0$, on ait
$\dim(A/\mathfrak p) \geq 2$. Alors
$$
\textit{F\'Et}_{U'} \longrightarrow \textit{F\'Et}_{U_0},\quad
V' \longmapsto V_0 = V' \times_{U'} U_0
$$
est un foncteur fidèle. De plus, il existe un $U'$ satisfaisant
l'hypothèse, et tout ouvert plus petit $U'' \subset U'$ contenant
$U_0$ satisfait également cette hypothèse. En particulier, le
foncteur de restriction
$$
\colim_{U_0 \subset U' \subset U\text{ ouvert}} \textit{F\'Et}_{U'}
\longrightarrow
\textit{F\'Et}_{U_0}
$$
est fidèle.
\end{lemma}

\begin{proof}
D'après Algèbre, lemme \ref{algebra-lemma-one-equation},
on voit que $V(\mathfrak p)$ rencontre $U_0$ pour
tout idéal premier $\mathfrak p$ de $A$ tel que $\dim(A/\mathfrak p) \geq 2$.
Ainsi, le foncteur affiché est fidèle pour un $U'$ comme dans l'énoncé,
d'après le lemme \ref{pione-lemma-restriction-faithful}.
Pour voir qu'un tel $U'$ existe, remarquons que, pour
$\mathfrak p \subset A$ tel que $\mathfrak p \in U$,
$\mathfrak p \not \in U_0$ et $\dim(A/\mathfrak p) = 1$,
alors $\mathfrak p$ correspond à un point fermé de $U$
et donc $V(\mathfrak p) \cap U_0 = \emptyset$.
On peut ainsi prendre pour $U'$ le complémentaire des composantes
irréductibles de $X$ qui ne rencontrent pas $U_0$ et sont de dimension $1$.
\end{proof}

\begin{lemma}
\label{pione-lemma-fully-faithful-general-better}
Dans la situation \ref{pione-situation-local-lefschetz}, supposons que
\begin{enumerate}
\item $A$ possède un complexe dualisant et soit complet pour la topologie $f$-adique,
\item toute composante irréductible de $X$ non contenue dans $X_0$
soit de dimension $\geq 3$.
\end{enumerate}
Alors le foncteur de restriction
$$
\colim_{U_0 \subset U' \subset U\text{ ouvert}} \textit{F\'Et}_{U'}
\longrightarrow
\textit{F\'Et}_{U_0}
$$
est pleinement fidèle.
\end{lemma}

\begin{proof}
Pour le démontrer, on peut remplacer $A$ par son réduit, grâce à l'invariance topologique
du groupe fondamental ; voir le lemme \ref{pione-lemma-thickening}.
Le résultat découle alors du
lemme \ref{pione-lemma-restriction-fully-faithful-general}
et du lemme de Géométrie algébrique et formelle
\ref{algebraization-lemma-fully-faithful-general}.
\end{proof}

\begin{lemma}
\label{pione-lemma-fully-faithful-general}
Dans la situation \ref{pione-situation-local-lefschetz}, supposons que
\begin{enumerate}
\item $A$ soit complet pour la topologie $f$-adique,
\item $f$ soit non diviseur de zéro,
\item $H^1_\mathfrak m(A/fA)$ soit un $A$-module fini.
\end{enumerate}
Alors le foncteur de restriction
$$
\colim_{U_0 \subset U' \subset U\text{ ouvert}} \textit{F\'Et}_{U'}
\longrightarrow
\textit{F\'Et}_{U_0}
$$
est pleinement fidèle.
\end{lemma}

\begin{proof}
Cela résulte du
lemme \ref{pione-lemma-restriction-fully-faithful-general} et du
lemme de Géométrie algébrique et formelle
\ref{algebraization-lemma-fully-faithful-general-alternative}.
\end{proof}








\section{Revêtements finis étales des spectres épointés, III}
\label{pione-section-pi1-punctured-spec-III}

\noindent
Dans cette section, nous étudions, dans la situation \ref{pione-situation-local-lefschetz}, quand
le foncteur de restriction
$$
\colim_{U_0 \subset U' \subset U\text{ ouvert}} \textit{F\'Et}_{U'}
\longrightarrow
\textit{F\'Et}_{U_0}
$$
est une équivalence de catégories.

\begin{lemma}
\label{pione-lemma-essentially-surjective-general-better}
Dans la situation \ref{pione-situation-local-lefschetz}, supposons que
\begin{enumerate}
\item $A$ possède un complexe dualisant et soit complet pour la topologie $f$-adique,
\item l'une des conditions suivantes soit satisfaite :
\begin{enumerate}
\item $A_f$ vérifie $(S_2)$ et toute composante irréductible de $X$
non contenue dans $X_0$ soit de dimension $\geq 4$, ou bien
\item si $\mathfrak p \not \in V(f)$ et
$V(\mathfrak p) \cap V(f) \not = \{\mathfrak m\}$, alors
$\text{depth}(A_\mathfrak p) + \dim(A/\mathfrak p) > 3$.
\end{enumerate}
\end{enumerate}
Alors le foncteur de restriction
$$
\colim_{U_0 \subset U' \subset U\text{ ouvert}} \textit{F\'Et}_{U'}
\longrightarrow
\textit{F\'Et}_{U_0}
$$
est une équivalence.
\end{lemma}

\begin{proof}
Cela résulte du lemme \ref{pione-lemma-restriction-equivalence-general}
et du
lemme de Géométrie algébrique et formelle
\ref{algebraization-lemma-equivalence-better}.
\end{proof}

\begin{lemma}
\label{pione-lemma-essentially-surjective-general}
Dans la situation \ref{pione-situation-local-lefschetz}, supposons que
\begin{enumerate}
\item $A$ soit complet pour la topologie $f$-adique,
\item $f$ soit non diviseur de zéro,
\item $H^1_\mathfrak m(A/fA)$ et $H^2_\mathfrak m(A/fA)$
soient des $A$-modules finis.
\end{enumerate}
Alors le foncteur de restriction
$$
\colim_{U_0 \subset U' \subset U\text{ ouvert}} \textit{F\'Et}_{U'}
\longrightarrow
\textit{F\'Et}_{U_0}
$$
est une équivalence.
\end{lemma}

\begin{proof}
Cela résulte du lemme \ref{pione-lemma-restriction-equivalence-general}
et du
lemme de Géométrie algébrique et formelle
\ref{algebraization-lemma-equivalence}.
\end{proof}

\begin{remark}
\label{pione-remark-combine}
Soient $(A, \mathfrak m)$ un anneau local noethérien complet et
$f \in \mathfrak m$ non nul. Supposons que $A_f$ vérifie $(S_2)$ et que
toute composante irréductible de $\Spec(A)$ soit de dimension $\geq 4$.
Alors le lemme \ref{pione-lemma-essentially-surjective-general-better}
nous dit que la catégorie
$$
\colim\nolimits_{U' \subset U\text{ ouvert, }U_0 \subset U'}
\text{ catégorie des schémas finis étales sur }U'
$$
est équivalente à la catégorie des schémas finis étales sur $U_0$.
Cela vaut par exemple si $A$ est un anneau normal intègre de dimension $\geq 4$ !
\end{remark}




\section{Revêtements finis étales des spectres épointés, IV}
\label{pione-section-pi1-punctured-spec-IV}

\noindent
Soient $X, X_0, U, U_0$ comme dans la situation \ref{pione-situation-local-lefschetz}.
Dans cette section, nous demandons quand le foncteur de restriction
$$
\textit{F\'Et}_U
\longrightarrow
\textit{F\'Et}_{U_0}
$$
est essentiellement surjectif.
Nous procéderons en reprenant les résultats de la
section \ref{pione-section-pi1-punctured-spec-III},
puis en comblant les lacunes à l'aide de la pureté. Rappelons que
nous disons que {\it la pureté vaut} pour un anneau local noethérien
$(A, \mathfrak m)$ si le foncteur de restriction
$\textit{F\'Et}_X \to \textit{F\'Et}_U$ est essentiellement
surjectif, où $X = \Spec(A)$ et $U = X \setminus \{\mathfrak m\}$.

\begin{lemma}
\label{pione-lemma-equivalence-better}
Dans la situation \ref{pione-situation-local-lefschetz}, supposons que
\begin{enumerate}
\item $A$ possède un complexe dualisant et soit complet pour la topologie $f$-adique,
\item l'une des conditions suivantes soit satisfaite :
\begin{enumerate}
\item $A_f$ vérifie $(S_2)$ et toute composante irréductible de $X$
non contenue dans $X_0$ soit de dimension $\geq 4$, ou bien
\item si $\mathfrak p \not \in V(f)$ et
$V(\mathfrak p) \cap V(f) \not = \{\mathfrak m\}$, alors
$\text{depth}(A_\mathfrak p) + \dim(A/\mathfrak p) > 3$.
\end{enumerate}
\item pour tout idéal maximal $\mathfrak p \subset A_f$,
la pureté vaille pour $(A_f)_\mathfrak p$.
\end{enumerate}
Alors le foncteur de restriction $\textit{F\'Et}_U \to \textit{F\'Et}_{U_0}$
est essentiellement surjectif.
\end{lemma}

\begin{proof}
Soit $V_0 \to U_0$ un morphisme fini étale. D'après le
lemme \ref{pione-lemma-essentially-surjective-general-better},
il existe un ouvert $U' \subset U$ contenant $U_0$ et
un morphisme fini étale $V' \to U'$ dont le changement de base à $U_0$
est isomorphe à $V_0 \to U_0$. Comme $U' \supset U_0$,
on voit que $U \setminus U'$ est formé de points correspondant
aux idéaux premiers $\mathfrak p_1, \ldots, \mathfrak p_n$ comme dans (3).
Par hypothèse, on peut trouver des morphismes finis étales
$V'_i \to \Spec(A_{\mathfrak p_i})$ qui coïncident avec
$V' \to U'$ sur $U' \times_U \Spec(A_{\mathfrak p_i})$.
D'après Limites, lemme \ref{limits-lemma-glueing-near-closed-point},
appliqué $n$ fois, on voit que $V' \to U'$ se prolonge en un morphisme fini étale
$V \to U$.
\end{proof}

\begin{lemma}
\label{pione-lemma-equivalence}
Soit $(A, \mathfrak m)$ un anneau local noethérien.
Soit $f \in \mathfrak m$. Supposons que
\begin{enumerate}
\item $A$ soit complet pour la topologie $f$-adique,
\item $f$ soit non diviseur de zéro,
\item $H^1_\mathfrak m(A/fA)$ et $H^2_\mathfrak m(A/fA)$ soient des
$A$-modules finis,
\item pour tout idéal maximal $\mathfrak p \subset A_f$,
la pureté vaille pour $(A_f)_\mathfrak p$.
\end{enumerate}
Alors le foncteur de restriction $\textit{F\'Et}_U \to \textit{F\'Et}_{U_0}$
est essentiellement surjectif.
\end{lemma}

\begin{proof}
La démonstration est identique à celle du
lemme \ref{pione-lemma-equivalence-better},
en utilisant le
lemme \ref{pione-lemma-essentially-surjective-general}
au lieu du
lemme \ref{pione-lemma-essentially-surjective-general-better}.
\end{proof}



\section{Pureté dans le cas local, II}
\label{pione-section-local-purity-II}

\noindent
Cette section prolonge la section \ref{pione-section-local-purity}.
Rappelons que nous disons que {\it la pureté vaut} pour un anneau local noethérien
$(A, \mathfrak m)$ si le foncteur de restriction
$\textit{F\'Et}_X \to \textit{F\'Et}_U$ est essentiellement
surjectif, où $X = \Spec(A)$ et $U = X \setminus \{\mathfrak m\}$.

\begin{lemma}
\label{pione-lemma-purity-inherited-by-hypersurface-better}
Soit $(A, \mathfrak m)$ un anneau local noethérien.
Soit $f \in \mathfrak m$. Supposons que
\begin{enumerate}
\item $A$ possède un complexe dualisant et soit complet pour la topologie $f$-adique,
\item l'une des conditions suivantes soit satisfaite :
\begin{enumerate}
\item $A_f$ vérifie $(S_2)$ et toute composante irréductible de $X$
non contenue dans $X_0$ soit de dimension $\geq 4$, ou bien
\item si $\mathfrak p \not \in V(f)$ et
$V(\mathfrak p) \cap V(f) \not = \{\mathfrak m\}$, alors
$\text{depth}(A_\mathfrak p) + \dim(A/\mathfrak p) > 3$.
\end{enumerate}
\item pour tout idéal maximal $\mathfrak p \subset A_f$,
la pureté vaille pour $(A_f)_\mathfrak p$, et
\item la pureté vaille pour $A$.
\end{enumerate}
Alors la pureté vaut pour $A/fA$.
\end{lemma}

\begin{proof}
Notons $X = \Spec(A)$ et $U = X \setminus \{\mathfrak m\}$
le spectre épointé. De même, on a $X_0 = \Spec(A/fA)$
et $U_0 = X_0 \setminus \{\mathfrak m\}$.
Soit $V_0 \to U_0$ un morphisme fini étale. D'après le
lemme \ref{pione-lemma-equivalence-better},
il existe un morphisme fini étale $V \to U$
dont le changement de base à $U_0$
est isomorphe à $V_0 \to U_0$.
D'après l'hypothèse (4), le morphisme $V \to U$ se prolonge
en un morphisme fini étale $Y \to X$. La restriction de
$Y$ à $X_0$ est alors le prolongement cherché de $V_0 \to U_0$.
\end{proof}

\begin{lemma}
\label{pione-lemma-purity-inherited-by-hypersurface}
Soit $(A, \mathfrak m)$ un anneau local noethérien.
Soit $f \in \mathfrak m$. Supposons que
\begin{enumerate}
\item $A$ soit complet pour la topologie $f$-adique,
\item $f$ soit non diviseur de zéro,
\item $H^1_\mathfrak m(A/fA)$ et $H^2_\mathfrak m(A/fA)$ soient des
$A$-modules finis,
\item pour tout idéal maximal $\mathfrak p \subset A_f$,
la pureté vaille pour $(A_f)_\mathfrak p$,
\item la pureté vaille pour $A$.
\end{enumerate}
Alors la pureté vaut pour $A/fA$.
\end{lemma}

\begin{proof}
La démonstration est identique à celle du
lemme \ref{pione-lemma-purity-inherited-by-hypersurface-better},
en utilisant le
lemme \ref{pione-lemma-equivalence}
au lieu du
lemme \ref{pione-lemma-equivalence-better}.
\end{proof}

\noindent
On peut maintenant amplifier les résultats précédents pour démontrer que
la pureté vaut pour les intersections complètes de dimension $\geq 3$.
Rappelons qu'un anneau local noethérien est appelé une intersection
complète si son complété est le quotient d'un
anneau local régulier par l'idéal engendré par une suite régulière.
Voir la discussion dans Algèbre à puissances divisées, section \ref{dpa-section-lci}.

\begin{proposition}
\label{pione-proposition-purity-complete-intersection}
Soit $(A, \mathfrak m)$ un anneau local noethérien. Si $A$ est une
intersection complète de dimension $\geq 3$, alors la pureté
vaut pour $A$, au sens où tout revêtement fini étale du
spectre épointé se prolonge.
\end{proposition}

\begin{proof}
D'après le lemme \ref{pione-lemma-purity-and-completion}, on peut supposer $A$
local complet. Par hypothèse, on peut écrire
$A = B/(f_1, \ldots, f_r)$, où $B$ est un anneau local régulier complet
et $f_1, \ldots, f_r$ une suite régulière.
On achève la démonstration par récurrence sur $r$.
Le cas initial est $r = 0$ et résulte du
lemme \ref{pione-lemma-local-purity}, qui s'applique aux
anneaux réguliers de dimension $\geq 2$.

\medskip\noindent
Supposons que $A = B/(f_1, \ldots, f_r)$ et que la proposition
vaille pour $r - 1$. Posons $A' = B/(f_1, \ldots, f_{r - 1})$ et appliquons le
lemme \ref{pione-lemma-purity-inherited-by-hypersurface} à $f_r \in A'$.
C'est permis :
la condition (2) est satisfaite puisque $f_1, \ldots, f_r$ est une suite régulière,
la condition (1) est satisfaite puisque $B$, et donc $A'$, est complet,
la condition (3) est satisfaite puisque $A = A'/f_r A'$ est de Cohen--Macaulay de dimension
$\dim(A) \geq 3$ ; voir Complexes dualisants, lemme \ref{dualizing-lemma-depth},
la condition (4) est satisfaite par l'hypothèse de récurrence, car
$\dim((A'_{f_r})_\mathfrak p) \geq 3$ pour un idéal premier maximal
$\mathfrak p$ de $A'_{f_r}$, et car
$(A'_{f_r})_\mathfrak p = B_\mathfrak q/(f_1, \ldots, f_{r - 1})$
pour un certain $\mathfrak q \subset B$,
la condition (5) est satisfaite par l'hypothèse de récurrence.
\end{proof}




\section{Pureté dans le cas local, III}
\label{pione-section-local-purity-III}

\noindent
Cette section prolonge la discussion des
sections \ref{pione-section-local-purity} et \ref{pione-section-local-purity-II}.

\begin{lemma}
\label{pione-lemma-fully-faithful-power-series-over-depth2}
Soit $(A, \mathfrak m)$ un anneau local noethérien de profondeur $\geq 2$.
Soit $B = A[[x_1, \ldots, x_d]]$ avec $d \geq 1$.
Posons $Y = \Spec(B)$ et $Y_0 = V(x_1, \ldots, x_d)$.
Pour tout sous-schéma ouvert $V \subset Y$ tel que
$V_0 = V \cap Y_0$ soit égal à $Y_0 \setminus \{\mathfrak m_B\}$,
le foncteur de restriction
$$
\textit{F\'Et}_V \longrightarrow \textit{F\'Et}_{V_0}
$$
est pleinement fidèle.
\end{lemma}

\begin{proof}
Posons $I = (x_1, \ldots, x_d)$. Posons $X = \Spec(A)$.
Si l'on utilise le morphisme $Y \to X$ pour identifier $Y_0$ à $X$,
alors $V_0$ s'identifie au spectre épointé $U$ de $A$.
En prenant les images directes des modules par ce morphisme affine, on obtient
\begin{align*}
\lim_n \Gamma(V_0, \mathcal{O}_V/I^n\mathcal{O}_V)
& =
\lim_n \Gamma(V_0, \mathcal{O}_Y/I^n\mathcal{O}_Y) \\
& =
\lim_n \Gamma(U, \mathcal{O}_U[x_1, \ldots, x_d]/(x_1, \ldots, x_d)^n) \\
& =
\lim_n A[x_1, \ldots, x_d]/(x_1, \ldots, x_d)^n \\
& =
B
\end{align*}
En effet, puisque la profondeur de $A$ est $\geq 2$, on a $\Gamma(U, \mathcal{O}_U) = A$ ;
voir Cohomologie locale, lemme
\ref{local-cohomology-lemma-finiteness-pushforwards-and-H1-local}.
Ainsi, pour tout ouvert $V \subset Y$ comme dans le lemme, on obtient
$$
B = \Gamma(Y, \mathcal{O}_Y) \to \Gamma(V, \mathcal{O}_V) \to
\lim_n \Gamma(V_0, \mathcal{O}_Y/I^n\mathcal{O}_Y) = B
$$
ce qui implique que les deux flèches sont des isomorphismes (on omet un petit détail).
D'après Géométrie algébrique et formelle, lemme
\ref{algebraization-lemma-completion-fully-faithful},
on conclut que
$\textit{Coh}(\mathcal{O}_V) \to \textit{Coh}(V, I\mathcal{O}_V)$
est pleinement fidèle sur la sous-catégorie pleine des objets finis localement libres.
On conclut donc par le
lemme \ref{pione-lemma-restriction-fully-faithful}.
\end{proof}

\begin{lemma}
\label{pione-lemma-purity-power-series-over-depth2}
\begin{slogan}
Ramanujam--Samuel pour les revêtements finis étales
\end{slogan}
Soit $(A, \mathfrak m)$ un anneau local noethérien de profondeur $\geq 2$. Soit
$B = A[[x_1, \ldots, x_d]]$ avec $d \geq 1$. Pour tout ouvert
$V \subset Y = \Spec(B)$ qui contient
\begin{enumerate}
\item tout idéal premier $\mathfrak q \subset B$ tel que
$\mathfrak q \cap A \not = \mathfrak m$,
\item l'idéal premier $\mathfrak m B$,
\end{enumerate}
le foncteur
$
\textit{F\'Et}_Y
\to
\textit{F\'Et}_V
$
est une équivalence. En particulier, la pureté vaut pour $B$.
\end{lemma}

\begin{proof}
Un idéal premier $\mathfrak q \subset B$ qui n'appartient pas à $V$
est au-dessus de $\mathfrak m$. Dans ce cas, $A \to B_\mathfrak q$
est un homomorphisme local plat, et donc $\text{depth}(B_\mathfrak q) \geq 2$
(Algèbre, lemme \ref{algebra-lemma-apply-grothendieck}).
Le foncteur est donc pleinement fidèle d'après le
lemme \ref{pione-lemma-quasi-compact-dense-open-connected-at-infinity-Noetherian},
combiné au lemme de Cohomologie locale
\ref{local-cohomology-lemma-depth-2-connected-punctured-spectrum}.

\medskip\noindent
Soit $W \to V$ un morphisme fini étale. Soit $B \to C$ l'unique morphisme fini
d'anneaux tel que $\Spec(C) \to Y$ soit le morphisme fini prolongeant
$W \to V$ construit dans le lemme \ref{pione-lemma-extend-S2}.
Remarquons que $C = \Gamma(W, \mathcal{O}_W)$.

\medskip\noindent
Posons $Y_0 = V(x_1, \ldots, x_d)$ et $V_0 = V \cap Y_0$. Posons $X = \Spec(A)$.
Si l'on utilise le morphisme $Y \to X$ pour identifier $Y_0$ à $X$,
alors $V_0$ s'identifie au spectre épointé $U$ de $A$.
On peut ainsi considérer $W_0 = W \times_Y Y_0$ comme un schéma fini étale
sur $U$. Alors
$$
W_0 \times_U (U \times_X Y)
\quad\text{et}\quad
W \times_V (U \times_X Y)
$$
sont des schémas finis étales sur $U \times_X Y$ dont les restrictions sont
des schémas finis étales isomorphes sur $V_0$. D'après le
lemme \ref{pione-lemma-fully-faithful-power-series-over-depth2},
appliqué à l'ouvert $U \times_X Y$, on obtient un isomorphisme
$$
W_0 \times_U (U \times_X Y) \longrightarrow W \times_V (U \times_X Y)
$$
sur $U \times_X Y$.

\medskip\noindent
Remarquons que $C_0 = \Gamma(W_0, \mathcal{O}_{W_0})$ est une $A$-algèbre finie
d'après le lemme \ref{pione-lemma-extend-S2} appliqué à $W_0 \to U \subset X$ (exactement
comme ci-dessus pour $B \to C$). Comme la construction du
lemme \ref{pione-lemma-extend-S2} est compatible au changement de base plat
et au changement d'ouverts, l'isomorphisme ci-dessus induit
un isomorphisme
$$
\Psi : C \longrightarrow C_0 \otimes_A B
$$
de $B$-algèbres finies. Cependant, on sait que $\Spec(C) \to Y$
est étale en tout point situé au-dessus d'au moins un point de $Y$ qui est
au-dessus de $\mathfrak m \in X$. Comme $\Psi$ est un isomorphisme, on
conclut que $\Spec(C_0) \to X$
est étale au-dessus de $\mathfrak m$ (on omet un petit détail).
Cela signifie bien sûr que $A \to C_0$ est fini
étale, et donc que $B \to C$ est fini étale.
\end{proof}

\begin{lemma}
\label{pione-lemma-purity-smooth-over-depth2}
Soit $f : X \to S$ un morphisme de schémas. Soit $U \subset X$
un sous-schéma ouvert. Supposons que
\begin{enumerate}
\item $f$ soit lisse,
\item $S$ soit noethérien,
\item pour $s \in S$ tel que $\text{depth}(\mathcal{O}_{S, s}) \leq 1$,
on ait $X_s = U_s$,
\item $U_s \subset X_s$ soit dense pour tout $s \in S$.
\end{enumerate}
Alors $\textit{F\'Et}_X \to \textit{F\'Et}_U$ est une équivalence.
\end{lemma}

\begin{proof}
Le foncteur est pleinement fidèle d'après le
lemme \ref{pione-lemma-quasi-compact-dense-open-connected-at-infinity-Noetherian},
combiné au lemme de Cohomologie locale
\ref{local-cohomology-lemma-depth-2-connected-punctured-spectrum}
(ainsi qu'à une application du
lemme d'Algèbre \ref{algebra-lemma-apply-grothendieck}
pour vérifier la condition de profondeur).

\medskip\noindent
Soit $\pi : V \to U$ un morphisme fini étale. Soit $Y \to X$
le morphisme fini construit dans le lemme \ref{pione-lemma-extend-S2}.
Il faut montrer que $Y \to X$ est fini étale.
Pour le montrer en tout point $x \in X$ s'envoyant sur un
point donné $s \in S$, on peut effectuer le changement de base par un morphisme plat
$S' \to S$ de schémas noethériens tel que $s$ appartienne
à l'image. Cela résulte de la compatibilité de la
construction du lemme \ref{pione-lemma-extend-S2} au changement de base plat.

\medskip\noindent
Après avoir agrandi $U$, on peut supposer que $U \subset X$ est
le plus grand ouvert au-dessus duquel $Y \to X$ est fini étale.
Soit $Z \subset X$ le complémentaire de $U$.
Pour obtenir une contradiction, supposons $Z \not = \emptyset$.
Soit $s \in S$ un point de l'image de $Z \to S$
tel qu'aucune générisation stricte de $s$ n'appartienne à cette image.
Après changement de base à $\Spec(\mathcal{O}_{S, s})$,
on voit alors que $S = \Spec(A)$, où $(A, \mathfrak m, \kappa)$ est
un anneau local noethérien de profondeur $\geq 2$, et que $Z$ est
contenu dans la fibre fermée $X_s$
et n'est dense nulle part dans $X_s$. Choisissons un point fermé $z \in Z$.
Alors $\kappa(z)/\kappa$ est fini (par le théorème des zéros de Hilbert ; voir
Algèbre, théorème \ref{algebra-theorem-nullstellensatz}).
Choisissons un morphisme fini plat $(S', s') \to (S, s)$ de schémas locaux
réalisant l'extension résiduelle $\kappa(z)/\kappa$ ; voir
Algèbre, lemme \ref{algebra-lemma-finite-free-given-residue-field-extension}.
Après avoir effectué le changement de base par $S' \to S$, on se ramène au cas
où $\kappa(z) = \kappa$.

\medskip\noindent
D'après Compléments sur les morphismes, lemme \ref{more-morphisms-lemma-slice-smooth},
il existe un sous-schéma localement fermé $S' \subset X$ passant par $z$
tel que $S' \to S$ soit étale en $z$. Après changement de base
par $S' \to S$, on peut supposer qu'il existe une section $\sigma : S \to X$
telle que $\sigma(s) = z$. Choisissons un voisinage affine
$\Spec(B) \subset X$ de $z$. Alors $A \to B$ est un morphisme lisse
d'anneaux qui possède une section $\sigma : B \to A$. Notons $I = \Ker(\sigma)$
et notons $B^\wedge$ le complété $I$-adique de $B$.
Alors $B^\wedge \cong A[[x_1, \ldots, x_d]]$ pour un certain $d \geq 0$ ; voir
Algèbre, lemme \ref{algebra-lemma-section-smooth}.
Remarquons que $d > 0$, car sinon on verrait que $X \to S$
est étale en $z$, ce qui impliquerait que $z$ est un point générique de
$X_s$, et donc que $z \in U$ d'après l'hypothèse (4).
De même, si $d > 0$, alors $\mathfrak m B^\wedge$ s'envoie dans
$U$ par le morphisme $\Spec(B^\wedge) \to X$.
Il suffit de démontrer que $Y \to X$ est fini étale après changement de base
à $\Spec(B^\wedge)$. Comme $B \to B^\wedge$ est plat
(Algèbre, lemme \ref{algebra-lemma-completion-flat}),
cela résulte du lemme \ref{pione-lemma-purity-power-series-over-depth2}
et de l'unicité dans la construction de $Y \to X$.
\end{proof}

\begin{proposition}
\label{pione-proposition-purity-smooth-over-depth2}
Soit $A \to B$ un homomorphisme local d'anneaux locaux noethériens.
Supposons $A$ de profondeur $\geq 2$, le morphisme $A \to B$ formellement lisse pour la
topologie $\mathfrak m_B$-adique, et $\dim(B) > \dim(A)$. Pour tout ouvert
$V \subset Y = \Spec(B)$ qui contient
\begin{enumerate}
\item tout idéal premier $\mathfrak q \subset B$ tel que
$\mathfrak q \cap A \not = \mathfrak m_A$,
\item l'idéal premier $\mathfrak m_A B$,
\end{enumerate}
le foncteur $\textit{F\'Et}_Y \to \textit{F\'Et}_V$
est une équivalence. En particulier, la pureté vaut pour $B$.
\end{proposition}

\begin{proof}
Un idéal premier $\mathfrak q \subset B$ qui n'appartient pas à $V$
est au-dessus de $\mathfrak m_A$. Dans ce cas, $A \to B_\mathfrak q$
est un homomorphisme local plat, et donc $\text{depth}(B_\mathfrak q) \geq 2$
(Algèbre, lemme \ref{algebra-lemma-apply-grothendieck}).
Le foncteur est donc pleinement fidèle d'après le
lemme \ref{pione-lemma-quasi-compact-dense-open-connected-at-infinity-Noetherian},
combiné au lemme de Cohomologie locale
\ref{local-cohomology-lemma-depth-2-connected-punctured-spectrum}.

\medskip\noindent
Notons $A^\wedge$ et $B^\wedge$ les complétés de $A$ et $B$
pour leurs idéaux maximaux. Remarquons que les hypothèses
de la proposition sont satisfaites par $A^\wedge \to B^\wedge$ ; voir
Compléments d'algèbre, lemmes
\ref{more-algebra-lemma-completion-dimension},
\ref{more-algebra-lemma-completion-depth} et
\ref{more-algebra-lemma-formally-smooth-completion}.
Par l'unicité et la compatibilité au changement de base plat
de la construction du lemme \ref{pione-lemma-extend-S2},
il suffit de démontrer la surjectivité essentielle pour
$A^\wedge \to B^\wedge$ et l'image inverse de $V$
(détails omis ; comparer au lemme \ref{pione-lemma-purity-and-completion}
dans le cas où $V$ est le spectre épointé).
D'après Compléments d'algèbre, proposition \ref{more-algebra-proposition-fs-regular},
cela signifie que l'on peut supposer $A \to B$ régulier.

\medskip\noindent
Soit $W \to V$ un morphisme fini étale.
D'après le théorème de Popescu
(Lissage des morphismes d'anneaux, théorème \ref{smoothing-theorem-popescu}),
on peut écrire $B = \colim B_i$ comme colimite filtrante
de $A$-algèbres lisses. On peut choisir un $i$ et un
ouvert $V_i \subset \Spec(B_i)$ dont l'image inverse est $V$
(Limites, lemme \ref{limits-lemma-descend-opens}).
Après avoir augmenté $i$, on peut supposer qu'il existe un morphisme
fini étale $W_i \to V_i$ dont le changement de base à $V$
est $W \to V$ ; voir
Limites, lemmes \ref{limits-lemma-descend-finite-presentation},
\ref{limits-lemma-descend-finite-finite-presentation} et
\ref{limits-lemma-descend-etale}.
On peut supposer que le complémentaire de $V_i$ est contenu
dans la fibre fermée de $\Spec(B_i) \to \Spec(A)$, puisque cela
est vrai pour $V$ (soit on choisit ainsi $V_i$, soit on utilise
le lemme précédent pour montrer que cela vaut pour $i$ assez grand).
Soit $\eta$ le point générique de la fibre fermée
de $\Spec(B) \to \Spec(A)$. Comme $\eta \in V$, l'image de
$\eta$ appartient à $V_i$. Ainsi, après avoir remplacé $V_i$ par un
voisinage ouvert affine de l'image du point fermé
de $\Spec(B)$, on peut supposer que la fibre fermée
de $\Spec(B_i) \to \Spec(A)$ est irréductible et que
son point générique appartient à $V_i$ (détails omis ; en effet,
un schéma lisse sur un corps est réunion disjointe de schémas irréductibles).
On peut alors appliquer le lemme \ref{pione-lemma-purity-smooth-over-depth2}
pour voir que $W_i \to V_i$ se prolonge en un morphisme fini étale
$\Spec(C_i) \to \Spec(B_i)$ et, en le ramenant
à $\Spec(B)$, on conclut que $W$ appartient à
l'image essentielle du foncteur
$\textit{F\'Et}_Y \to \textit{F\'Et}_V$,
comme souhaité.
\end{proof}





\section{Lefschetz pour le groupe fondamental}
\label{pione-section-global-lefschetz}

\noindent
Nous avons bien sûr déjà démontré plusieurs résultats de ce type
dans le cas local. Dans cette section, nous étudions le cas projectif.

\begin{proposition}
\label{pione-proposition-lefschetz-fully-faithful}
Soit $k$ un corps. Soit $X$ un schéma propre sur $k$.
Soit $\mathcal{L}$ un $\mathcal{O}_X$-module inversible ample.
Soit $s \in \Gamma(X, \mathcal{L})$. Soit $Y = Z(s)$ le
schéma des zéros de $s$. Supposons que, pour tout $x \in X \setminus Y$,
on ait
$$
\text{depth}(\mathcal{O}_{X, x}) + \dim(\overline{\{x\}}) > 1
$$
Alors le foncteur de restriction $\textit{F\'Et}_X \to \textit{F\'Et}_Y$
est pleinement fidèle. En fait, pour tout sous-schéma ouvert $V \subset X$
contenant $Y$, le foncteur de restriction
$\textit{F\'Et}_V \to \textit{F\'Et}_Y$
est pleinement fidèle.
\end{proposition}

\begin{proof}
La première assertion est une conséquence formelle du
lemme \ref{pione-lemma-restriction-fully-faithful-special}
et de la
proposition de Géométrie algébrique et formelle
\ref{algebraization-proposition-lefschetz}.
La seconde assertion découle du
lemme \ref{pione-lemma-restriction-fully-faithful-special}
et du
lemme de Géométrie algébrique et formelle
\ref{algebraization-lemma-lefschetz-addendum}.
\end{proof}

\begin{proposition}
\label{pione-proposition-lefschetz-equivalence-general}
Soit $k$ un corps. Soit $X$ un schéma propre sur $k$.
Soit $\mathcal{L}$ un $\mathcal{O}_X$-module inversible ample.
Soit $s \in \Gamma(X, \mathcal{L})$. Soit $Y = Z(s)$ le
schéma des zéros de $s$. Soit $\mathcal{V}$ l'ensemble des
sous-schémas ouverts de $X$ contenant $Y$, ordonné par inclusion inverse.
Supposons que, pour tout $x \in X \setminus Y$, on ait
$$
\text{depth}(\mathcal{O}_{X, x}) + \dim(\overline{\{x\}}) > 2
$$
Alors le foncteur de restriction
$$
\colim_\mathcal{V} \textit{F\'Et}_V \to \textit{F\'Et}_Y
$$
est une équivalence.
\end{proposition}

\begin{proof}
C'est une conséquence formelle du
lemme \ref{pione-lemma-restriction-equivalence-general} et de la
proposition de Géométrie algébrique et formelle
\ref{algebraization-proposition-lefschetz-equivalence}.
\end{proof}

\begin{proposition}
\label{pione-proposition-lefschetz-equivalence}
Soit $k$ un corps. Soit $X$ un schéma propre sur $k$.
Soit $\mathcal{L}$ un $\mathcal{O}_X$-module inversible ample.
Soit $s \in \Gamma(X, \mathcal{L})$. Soit $Y = Z(s)$ le
schéma des zéros de $s$.
Supposons que, pour tout $x \in X \setminus Y$, on ait
$$
\text{depth}(\mathcal{O}_{X, x}) + \dim(\overline{\{x\}}) > 2
$$
et que, lorsque $x \in X \setminus Y$ est fermé, la pureté vaille pour
$\mathcal{O}_{X, x}$. Alors le foncteur de restriction
$\textit{F\'Et}_X \to \textit{F\'Et}_Y$
est une équivalence. Si $X$, ou de manière équivalente $Y$, est connexe, alors
$$
\pi_1(Y, \overline{y}) \to \pi_1(X, \overline{y})
$$
est un isomorphisme pour tout point géométrique $\overline{y}$ de $Y$.
\end{proposition}

\begin{proof}
La pleine fidélité résulte de la
proposition \ref{pione-proposition-lefschetz-fully-faithful}.
D'après la proposition \ref{pione-proposition-lefschetz-equivalence-general},
tout objet de $\textit{F\'Et}_Y$
est isomorphe au produit fibré $U \times_V Y$ pour un certain
morphisme fini étale $U \to V$, où $V \subset X$
est un sous-schéma ouvert contenant $Y$.
Le complémentaire $T = X \setminus V$
est\footnote{En effet, $T$ est propre sur $k$ (car fermé dans $X$)
et affine (car fermé dans le schéma affine $X \setminus Y$ ; voir
Morphismes, lemme \ref{morphisms-lemma-proper-ample-delete-affine}),
et donc fini sur $k$
(Morphismes, lemme \ref{morphisms-lemma-finite-proper}).
Ainsi, $T$ est un ensemble fini de points fermés.}
un ensemble fini de points fermés de $X \setminus Y$.
Écrivons $T = \{x_1, \ldots, x_n\}$.
Par hypothèse, on peut trouver des morphismes finis étales
$V'_i \to \Spec(\mathcal{O}_{X, x_i})$ qui coïncident avec
$U \to V$ sur $V \times_X \Spec(\mathcal{O}_{X, x_i})$.
D'après Limites, lemme \ref{limits-lemma-glueing-near-closed-point},
appliqué $n$ fois, on voit que $U \to V$ se prolonge en un morphisme fini étale
$U' \to X$, comme souhaité.
Voir le lemme \ref{pione-lemma-what-equivalence-gives} pour la dernière assertion.
\end{proof}










\section{Pureté du lieu de ramification}
\label{pione-section-purity-ramification}

\noindent
Dans cette section, nous examinons l'analogue de la pureté du lieu de branchement pour
les morphismes génériquement finis. Ce résultat est apparemment dû à Gabber.
Un cas particulier est le théorème de pureté de van der Waerden pour le lieu où
un morphisme birationnel d'une variété normale vers une variété lisse n'est pas
un isomorphisme.

\begin{lemma}
\label{pione-lemma-characterize-rational-singularity}
Soit $A$ un domaine local noethérien normal de dimension $2$.
Supposons que $A$ soit un anneau de Nagata, possède un module dualisant $\omega_A$ et admette une
résolution des singularités $f : X \to \Spec(A)$.
Soit $\omega_X$ comme dans Résolution des surfaces,
remarque \ref{resolve-remark-dualizing-setup}.
Si $\omega_X \cong \mathcal{O}_X(E)$ pour un certain diviseur de
Cartier effectif $E \subset X$ porté par la fibre exceptionnelle,
alors $A$ définit une singularité rationnelle.
Si $f$ est une résolution minimale, alors $E = 0$.
\end{lemma}

\begin{proof}
Il existe un morphisme trace $Rf_*\omega_X \to \omega_A$, voir
Dualité pour les schémas, section \ref{duality-section-trace}.
D'après Grauert--Riemenschneider
(Résolution des surfaces,
proposition \ref{resolve-proposition-Grauert-Riemenschneider}),
on a $R^1f_*\omega_X = 0$.
Ainsi, le morphisme trace est un morphisme $f_*\omega_X \to \omega_A$.
On peut alors considérer
$$
\mathcal{O}_{\Spec(A)} = f_*\mathcal{O}_X \to f_*\omega_X \to \omega_A
$$
où le premier morphisme provient du morphisme
$\mathcal{O}_X \to \mathcal{O}_X(E) = \omega_X$ dont l'existence
est supposée dans l'énoncé du lemme.
La composée est un isomorphisme d'après Diviseurs, lemme
\ref{divisors-lemma-check-isomorphism-via-depth-and-ass},
car c'est un isomorphisme sur le spectre épointé de $A$
(par l'hypothèse du lemme et le fait que $f$ est un isomorphisme
sur le spectre épointé) et que $A$ et $\omega_A$
sont des $A$-modules de profondeur $2$ (d'après
Algèbre, lemme \ref{algebra-lemma-criterion-normal} et
Complexes dualisants, lemme \ref{dualizing-lemma-depth-dualizing-module}).
Par conséquent, $f_*\omega_X \to \omega_A$ est surjectif, donc est un isomorphisme.
Ainsi, $Rf_*\omega_X = \omega_A$, ce qui, par dualité, implique
$Rf_*\mathcal{O}_X = \mathcal{O}_{\Spec(A)}$.
Il s'ensuit que $H^1(X, \mathcal{O}_X) = 0$, ce qui implique que $A$
définit une singularité rationnelle (voir la discussion dans
Résolution des surfaces, section
\ref{resolve-section-bounded}, en particulier les
lemmes \ref{resolve-lemma-regular-rational} et
\ref{resolve-lemma-exact-sequence}).
Si $f$ est minimal, alors $E = 0$, car le morphisme
$f^*\omega_A \to \omega_X$ est surjectif d'après
une application répétée de Résolution des surfaces, lemme
\ref{resolve-lemma-dualizing-blow-up-rational},
et $\omega_A \cong A$, comme on l'a vu ci-dessus.
\end{proof}

\begin{lemma}
\label{pione-lemma-key-purity-ramification}
Posons $Y = \Spec(A)$ et soit $f : X \to Y$ un morphisme de type fini.
Soit $x \in X$ un point. Supposons que
\begin{enumerate}
\item $A$ soit un anneau local régulier excellent,
\item $\mathcal{O}_{X, x}$ soit normal de dimension $2$,
\item $f$ soit \'etale hors de $\overline{\{x\}}$.
\end{enumerate}
Alors $f$ est \'etale en $x$.
\end{lemma}

\begin{proof}
On remplace d'abord $X$ par un voisinage ouvert affine de $x$.
Notons que $\mathcal{O}_{X, x}$ est un anneau local excellent
(Compléments d'algèbre, lemme \ref{more-algebra-lemma-finite-type-over-excellent}).
On peut donc choisir une résolution minimale des singularités
$W \to \Spec(\mathcal{O}_{X, x})$, voir
Résolution des surfaces, théorème \ref{resolve-theorem-resolve}.
Après avoir éventuellement remplacé $X$ par un voisinage ouvert affine de $x$,
on peut trouver un morphisme propre $b : X' \to X$ tel que
$X' \times_X \Spec(\mathcal{O}_{X, x}) = W$, voir
Limites, lemme \ref{limits-lemma-glueing-near-closed-point}.
Après avoir encore rétréci $X$, on peut supposer que $X'$ est régulier.
En effet, on sait que $W$ est régulier et que $X'$ est excellent,
et le lieu régulier du spectre d'un anneau excellent est ouvert.
Puisque $W \to \Spec(\mathcal{O}_{X, x})$ est projectif
(comme composé d'une suite d'éclatements normalisés), on peut supposer,
après avoir rétréci $X$, que $b$ est projectif (détails omis).
Soit $U = X \setminus \overline{\{x\}}$.
Puisque $W \to \Spec(\mathcal{O}_{X, x})$ est un isomorphisme sur le
spectre épointé, on peut supposer que $b : X' \to X$ est un isomorphisme sur $U$.
Ainsi, on peut et on va aussi considérer $U$ comme un sous-schéma ouvert de $X'$.
Posons $f' = f \circ b : X' \to \Spec(A)$.

\medskip\noindent
Puisque $A$ est régulier, on voit que $\mathcal{O}_Y$ est un complexe dualisant pour $Y$.
Par conséquent, $f^!\mathcal{O}_Y$ est un complexe dualisant sur $X$
(Dualité pour les schémas, lemme \ref{duality-lemma-shriek-dualizing}).
Le lieu de Cohen--Macaulay de $X$ est ouvert d'après
Dualité pour les schémas, lemme \ref{duality-lemma-dualizing-module-CM-scheme}
(cela peut aussi se démontrer à l'aide de l'excellence).
Puisque $\mathcal{O}_{X, x}$ est de Cohen--Macaulay, après avoir rétréci
$X$, on peut supposer que $X$ est de Cohen--Macaulay.
Notons qu'un morphisme \'etale est une intersection complète locale.
Ainsi,
Dualité pour les schémas, lemme \ref{duality-lemma-fundamental-class-almost-lci},
s'applique avec $r = 0$ et donne un morphisme
$$
\mathcal{O}_X \longrightarrow \omega_{X/Y} = H^0(f^!\mathcal{O}_Y)
$$
qui est un isomorphisme sur $X \setminus \overline{\{x\}}$.
Puisque $\omega_{X/Y}$ est $(S_2)$ d'après
Dualité pour les schémas, lemme \ref{duality-lemma-shriek-over-CM},
on constate que ce morphisme est un isomorphisme d'après
Diviseurs, lemme \ref{divisors-lemma-check-isomorphism-via-depth-and-ass}.
Cela montre déjà que $X$, et en particulier $\mathcal{O}_{X, x}$, est
de Gorenstein.

\medskip\noindent
Posons $\omega_{X'/Y} = H^0((f')^!\mathcal{O}_Y)$. En raisonnant exactement de la
même manière que ci-dessus, on voit que $(f')^!\mathcal{O}_Y = \omega_{X'/Y}[0]$
est un complexe dualisant pour $X'$. Puisque $X'$ est régulier,
le morphisme $X' \to Y$ est un morphisme d'intersection complète locale, voir
Compléments sur les morphismes, lemme
\ref{more-morphisms-lemma-morphism-regular-schemes-is-lci}.
D'après Dualité pour les schémas, lemme \ref{duality-lemma-fundamental-class-lci},
il existe un morphisme
$$
\mathcal{O}_{X'} \longrightarrow \omega_{X'/Y}
$$
qui est un isomorphisme sur $U$. On en déduit que
$\omega_{X'/Y} = \mathcal{O}_{X'}(E)$ pour un certain diviseur de
Cartier effectif $E \subset X'$ disjoint de $U$.

\medskip\noindent
Puisque $\omega_{X/Y} = \mathcal{O}_X$, on voit que
$\omega_{X'/Y} = b^! f^!\mathcal{O}_Y = b^!\mathcal{O}_X$.
En revenant à $W \to \Spec(\mathcal{O}_{X, x})$,
on voit que $\omega_W = \mathcal{O}_W(E|_W)$.
D'après le lemme \ref{pione-lemma-characterize-rational-singularity},
on obtient $E|_W = 0$.
Cela signifie que $f' : X' \to Y$ est \'etale d'après le lemme (déjà utilisé)
de Dualité pour les schémas, lemme \ref{duality-lemma-fundamental-class-lci}.
Cela achève immédiatement la démonstration, car le caractère \'etale
de $f'$ force $b$ à être un isomorphisme.
\end{proof}

\begin{lemma}[Pureté du lieu de ramification]
\label{pione-lemma-purity-ramification}
\begin{reference}
Ce résultat pour les espaces complexes se trouve à la page 170 de \cite{Fischer}.
Dans le cas général, il s'agit de \cite[théorème 2.4]{Zong}, attribué à Gabber.
\end{reference}
Soit $f : X \to Y$ un morphisme de schémas localement noethériens.
Soit $x \in X$ et posons $y = f(x)$. Supposons que
\begin{enumerate}
\item $\mathcal{O}_{X, x}$ soit normal de dimension $\geq 1$,
\item $\mathcal{O}_{Y, y}$ soit régulier,
\item $f$ soit localement de type fini, et
\item pour les spécialisations $x' \leadsto x$ telles que
$\dim(\mathcal{O}_{X, x'}) = 1$, le morphisme $f$ soit \'etale en $x'$.
\end{enumerate}
Alors $f$ est \'etale en $x$.
\end{lemma}

\begin{proof}
On va démontrer le lemme par récurrence sur $d = \dim(\mathcal{O}_{X, x})$.

\medskip\noindent
Le cas $d = 1$ est immédiat, puisque dans ce cas le morphisme
$f$ est \'etale en $x$ par hypothèse. Supposons $d \geq 2$.

\medskip\noindent
On peut effectuer le changement de base par $\Spec(\mathcal{O}_{Y, y}) \to Y$
sans modifier la conclusion du lemme, voir
Morphismes, lemme \ref{morphisms-lemma-set-points-where-fibres-etale}.
Ainsi, on peut supposer $Y = \Spec(A)$, où $A$ est un anneau local régulier,
et que $y$ correspond à l'idéal maximal $\mathfrak m$ de $A$.

\medskip\noindent
Soit $x' \leadsto x$ une spécialisation avec $x' \not = x$.
Alors $\mathcal{O}_{X, x'}$ est normal, comme localisation de
$\mathcal{O}_{X, x}$. Si $x'$ n'est pas un point générique de $X$,
alors $1 \leq \dim(\mathcal{O}_{X, x'}) < d$ et on conclut que
$f$ est \'etale en $x'$ par l'hypothèse de récurrence.
On peut donc supposer que $f$ est \'etale en tout point se spécialisant en
$x$. Puisque l'ensemble des points où $f$ est \'etale est ouvert dans $X$
(par définition), on peut, après avoir remplacé $X$ par un voisinage ouvert de $x$,
supposer que $f$ est \'etale hors de $\overline{\{x\}}$.
En particulier, on voit que $f$ est \'etale sauf en des points
situés au-dessus du point fermé $y \in Y = \Spec(A)$.

\medskip\noindent
Soit $X' = X \times_{\Spec(A)} \Spec(A^\wedge)$. Soit $x' \in X'$
l'unique point au-dessus de $x$. D'après ce qui précède, on voit que
$X'$ est \'etale sur $\Spec(A^\wedge)$ hors de la fibre fermée et, par conséquent,
$X'$ est normal hors de la fibre fermée. Puisque $X$ est normal,
on en déduit que $X'$ est normal d'après
Résolution des surfaces, lemme \ref{resolve-lemma-normalization-completion}.
Alors, si l'on peut montrer que $X' \to \Spec(A^\wedge)$ est \'etale en $x'$,
le morphisme $f$ est \'etale en $x$ (d'après le lemme déjà cité de
Morphismes, lemme \ref{morphisms-lemma-set-points-where-fibres-etale}).
Ainsi, on peut et on va supposer que $A$ est un anneau local régulier complet.

\medskip\noindent
Le cas $d = 2$ résulte alors du lemme \ref{pione-lemma-key-purity-ramification}.

\medskip\noindent
Supposons $d > 2$. Soit $t \in \mathfrak m$, $t \not \in \mathfrak m^2$.
Posons $Y_0 = \Spec(A/tA)$ et $X_0 = X \times_Y Y_0$.
Alors $X_0 \to Y_0$ est \'etale hors de la fibre au-dessus du point fermé.
Puisque $d > 2$, on a $\dim(\mathcal{O}_{X_0, x}) = d - 1$, qui est $\geq 2$.
La normalisation $X_0' \to X_0$ est surjective et finie
(car on travaille sur un anneau local complet et que de tels anneaux sont de Nagata).
Soit $x' \in X_0'$ un point dont l'image est $x$. Par hypothèse de récurrence, le
morphisme $X'_0 \to Y_0$ est \'etale en $x'$. Des inclusions
$\kappa(y) \subset \kappa(x) \subset \kappa(x')$,
on déduit que $\kappa(x)$ est fini sur $\kappa(y)$.
Par conséquent, $x$ est un point fermé de la fibre de $X \to Y$
au-dessus de $y$. Mais comme $x$ est également un point générique de
cette fibre, on conclut que $f$ est quasi-fini en $x$,
et on se ramène au cas de la pureté du lieu de branchement, voir le
lemme \ref{pione-lemma-purity}.
\end{proof}




\section{Affinité du complémentaire du lieu de ramification}
\label{pione-section-stronger-purity}

\noindent
Soit $f : X \to Y$ un morphisme de type fini de schémas noethériens,
avec $X$ normal et $Y$ régulier. Soit $V \subset X$ le plus grand
sous-schéma ouvert sur lequel $f$ est \'etale.
La discussion dans \cite[chapitre IV, section 21.12]{EGA}
laisse penser que $V \to X$ pourrait être un morphisme affine.
Notons que, si $V \to X$ est affine, on en déduit la pureté du
lieu de ramification (lemme \ref{pione-lemma-purity-ramification}) en
utilisant Diviseurs, lemme \ref{divisors-lemma-complement-affine-open-immersion}.
Ainsi, l'affinité de $V \to X$ est une forme ``forte'' de la pureté
du lieu de ramification.
Dans cette section, nous démontrons que $V \to X$ est affine lorsque
$X$ et $Y$ sont équicaractéristiques et excellents, voir le
théorème \ref{pione-theorem-global}. Il paraît raisonnable
de conjecturer que le résultat reste vrai lorsque $X$ et $Y$ sont
de caractéristique mixte (mais toujours excellents).

\begin{lemma}
\label{pione-lemma-structure-cohomology}
Soit $(A, \mathfrak m)$ un anneau local régulier contenant un corps.
Soit $f : V \to \Spec(A)$ un morphisme \'etale et quasi-compact.
Supposons que $\mathfrak m \not \in f(V)$ et que
$g : V \to \Spec(A) \setminus \{\mathfrak m\}$ soit affine.
Alors $H^i(V, \mathcal{O}_V)$, $i > 0$, est isomorphe à une somme
directe de copies de l'enveloppe injective du corps résiduel de $A$.
\end{lemma}

\begin{proof}
Notons $U = \Spec(A) \setminus \{\mathfrak m\}$ le spectre épointé.
Ainsi, $g : V \to U$ est affine.
On a $H^i(V, \mathcal{O}_V) = H^i(U, g_*\mathcal{O}_V)$ d'après
Cohomologie des schémas, lemme \ref{coherent-lemma-relative-affine-cohomology}.
Le $\mathcal{O}_U$-module $g_*\mathcal{O}_V$ est quasi-cohérent d'après
Schémas, lemme \ref{schemes-lemma-push-forward-quasi-coherent}.
Pour tout $\mathcal{O}_U$-module quasi-cohérent $\mathcal{F}$,
la cohomologie $H^i(U, \mathcal{F})$, $i > 0$,
est de $\mathfrak m$-torsion, voir par exemple
Cohomologie locale, lemme \ref{local-cohomology-lemma-local-cohomology}.
En particulier, les $A$-modules $H^i(V, \mathcal{O}_V)$, $i > 0$,
sont de $\mathfrak m$-torsion.
Pour tout homomorphisme plat d'anneaux $A \to A'$, on a
$H^i(V, \mathcal{O}_V) \otimes_A A' = H^i(V', \mathcal{O}_{V'})$
où $V' = V \times_{\Spec(A)} \Spec(A')$, par changement de base plat, d'après
Cohomologie des schémas, lemme \ref{coherent-lemma-flat-base-change-cohomology}.
Si l'on prend pour $A'$ le complété de $A$ (qui est plat d'après
Compléments d'algèbre, section \ref{more-algebra-section-permanence-completion}),
alors on voit que
$$
H^i(V, \mathcal{O}_V) = H^i(V, \mathcal{O}_V) \otimes_A A' =
H^i(V', \mathcal{O}_{V'}),\quad\text{pour } i > 0
$$
La première égalité résulte de la propriété de torsion que nous venons d'établir et du
lemme de Compléments d'algèbre \ref{more-algebra-lemma-neighbourhood-equivalence}.
De plus, l'enveloppe injective du corps résiduel $k$
est la même pour $A$ et $A'$, voir
Complexes dualisants, lemme \ref{dualizing-lemma-compare}.
On se ramène ainsi au cas $A = k[[x_1, \ldots, x_d]]$, voir
Algèbre, section \ref{algebra-section-cohen-structure-theorem}.

\medskip\noindent
Supposons que la caractéristique de $k$ soit $p > 0$. Puisque $F : A \to A$,
$a \mapsto a^p$, est plat (Cohomologie locale, lemme
\ref{local-cohomology-lemma-frobenius-flat-regular})
et puisque $V \times_{\Spec(A), \Spec(F)} \Spec(A) \cong V$ comme schémas
sur $\Spec(A)$ d'après
Morphismes \'etales, lemme \ref{etale-lemma-relative-frobenius-etale},
ce qui précède donne
$H^i(V, \mathcal{O}_V) \otimes_{A, F} A \cong H^i(V, \mathcal{O}_V)$.
Le résultat découle donc du
lemme de Cohomologie locale
\ref{local-cohomology-lemma-structure-torsion-Frobenius-regular}.

\medskip\noindent
Supposons que la caractéristique de $k$ soit $0$. D'après
Cohomologie locale, lemme \ref{local-cohomology-lemma-etale-derivation},
il existe des opérateurs additifs $D_j$, $j = 1, \ldots, d$, sur
$H^i(V, \mathcal{O}_V)$ qui satisfont la règle de Leibniz par
rapport à $\partial_j = \partial/\partial x_j$.
Le résultat découle donc du lemme de Cohomologie locale
\ref{local-cohomology-lemma-structure-torsion-D-module-regular}.
\end{proof}

\begin{lemma}
\label{pione-lemma-conclude}
Dans la situation du lemme \ref{pione-lemma-structure-cohomology},
supposons que $H^i(V, \mathcal{O}_V) = 0$ pour $i \geq \dim(A) - 1$.
Alors $V$ est affine.
\end{lemma}

\begin{proof}
Soit $k = A/\mathfrak m$. Puisque $V \times_{\Spec(A)} \Spec(k) = \emptyset$,
la cohomologie et le changement de base donnent
$$
R\Gamma(V, \mathcal{O}_V) \otimes_A^\mathbf{L} k = 0
$$
Voir
Catégories dérivées de schémas, lemme \ref{perfect-lemma-compare-base-change}.
Il existe donc une suite spectrale
(Compléments d'algèbre, exemple \ref{more-algebra-example-tor})
$$
E_2^{p, q} = \text{Tor}_{-p}(k, H^q(V, \mathcal{O}_V)),\quad
d_2^{p, q} : E_2^{p, q} \to E_2^{p + 2, q - 1}
$$
et $d_r^{p, q} : E_r^{p, q} \to E_r^{p + r, q - r + 1}$,
qui converge vers zéro. D'après le lemme \ref{pione-lemma-structure-cohomology},
le lemme de Complexes dualisants \ref{dualizing-lemma-tor-injective-hull},
et notre hypothèse $H^i(V, \mathcal{O}_V) = 0$ pour $i \geq \dim(A) - 1$,
on conclut qu'aucune différentielle non nulle
n'aboutit à la case $(p, q) = (0, 0)$ ni n'en part. Ainsi,
$H^0(V, \mathcal{O}_V) \otimes_A k = 0$. Cela signifie que,
si $\mathfrak m = (x_1, \ldots, x_d)$, on dispose
d'un recouvrement ouvert $V = \bigcup V \times_{\Spec(A)} \Spec(A_{x_i})$
par les sous-schémas ouverts affines $V \times_{\Spec(A)} \Spec(A_{x_i})$
(car $V$ est affine sur le spectre épointé de $A$)
tel que $x_1, \ldots, x_d$ engendrent
l'idéal unité de $\Gamma(V, \mathcal{O}_V)$.
Cela implique que $V$ est affine d'après
Propriétés, lemme \ref{properties-lemma-characterize-affine}.
\end{proof}

\begin{theorem}
\label{pione-theorem-global}
Soit $Y$ un schéma régulier excellent sur un corps. Soit $f : X \to Y$
un morphisme de schémas de type fini, avec $X$ normal. Soit $V \subset X$
le plus grand sous-schéma ouvert où $f$ est \'etale. Alors le morphisme
d'inclusion $V \to X$ est affine.
\end{theorem}

\begin{proof}
Soit $x \in X$ d'image $y \in Y$. Il suffit de démontrer que
$V \cap W$ est affine pour un voisinage ouvert affine $W$ de $x$.
Puisque $\Spec(\mathcal{O}_{X, x})$ est la limite des schémas $W$,
cela équivaut à dire que
$$
V_x = V \times_X \Spec(\mathcal{O}_{X, x})
$$
est affine (Limites, lemme \ref{limits-lemma-limit-affine}).
Ainsi, si le théorème vaut pour le morphisme
$X \times_Y \Spec(\mathcal{O}_{Y, y}) \to \Spec(\mathcal{O}_{Y, y})$,
alors il vaut dans la situation initiale. En particulier, on peut supposer $Y$
régulier de dimension finie, ce qui permet de raisonner par récurrence
sur la dimension $d = \dim(Y)$. En combinant ceci avec le même argument,
on peut supposer que $Y$ est local, de point fermé $y$, et que
$V \cap (X \setminus f^{-1}(\{y\})) \to X \setminus f^{-1}(\{y\})$
est affine.

\medskip\noindent
Soit $x \in X$ un point au-dessus de $y$. Si $x \in V$, il n'y a
rien à démontrer. Observons que $f^{-1}(\{y\}) \cap V$
est un ensemble fini de points fermés (les fibres d'un morphisme \'etale
sont discrètes). Ainsi, quitte à remplacer $X$
par un voisinage ouvert affine de $x$, on peut supposer que
$y \not \in f(V)$. Il faut démontrer que $V$ est affine.

\medskip\noindent
Soit $e(V)$ le plus grand $i$ tel que $H^i(V, \mathcal{O}_V) \not = 0$.
Comme $X$ est affine, l'entier $e(V)$ est le maximum des nombres $e(V_x)$
pour $x \in X \setminus V$, voir
Cohomologie locale, lemme \ref{local-cohomology-lemma-cd-local},
et la caractérisation de la dimension cohomologique dans
Cohomologie locale, lemme \ref{local-cohomology-lemma-cd}.
On a $e(V_x) \leq \dim(\mathcal{O}_{X, x}) - 1$ d'après
Cohomologie locale, lemme \ref{local-cohomology-lemma-cd-dimension}.
Si $\dim(\mathcal{O}_{X, x}) \geq 2$, la pureté du
lieu de ramification (lemme \ref{pione-lemma-purity-ramification})
montre que $V_x$ est strictement contenu dans le spectre épointé de
$\mathcal{O}_{X, x}$. Puisque $\mathcal{O}_{X, x}$ est
normal et excellent, cela implique
$e(V_x) \leq \dim(\mathcal{O}_{X, x}) - 2$ d'après le théorème
d'annulation de Hartshorne--Lichtenbaum
(Cohomologie locale, lemme \ref{local-cohomology-lemma-affine-complement}).
D'autre part, puisque $X \to Y$ est de type fini
et que $V \subset X$ est dense (quitte à remplacer $X$
par l'adhérence de $V$), on voit que $\dim(\mathcal{O}_{X, x}) \leq d$
d'après la formule des dimensions
(Morphismes, lemme \ref{morphisms-lemma-dimension-formula}).
Par conséquent, $e(V) \leq \max(0, d -  2)$.
Ainsi, $V$ est affine d'après le lemme \ref{pione-lemma-conclude}
si $d \geq 2$. Si $d = 1$ ou $d = 0$, le spectre épointé de
$\mathcal{O}_{Y, y}$ est affine, donc $V$ est affine.
\end{proof}




\section{Homomorphismes de spécialisation dans le cas propre et lisse}
\sectionmark{Spécialisation dans le cas propre et lisse}
\label{pione-section-specialization-smooth-proper}

\noindent
Dans cette section, nous examinons le résultat suivant.
Soit $f : X \to S$ un morphisme propre et lisse de schémas.
Soit $s \leadsto s'$ une spécialisation de points de $S$.
Alors l'homomorphisme de spécialisation
$$
sp : \pi_1(X_{\overline{s}}) \longrightarrow \pi_1(X_{\overline{s}'})
$$
de la section \ref{pione-section-specialization-map}
est surjectif et
\begin{enumerate}
\item si la caractéristique de $\kappa(s')$ est nulle, c'est
un isomorphisme, ou bien
\item si la caractéristique de $\kappa(s')$ est $p > 0$, il
induit un isomorphisme sur les quotients maximaux premiers à $p$.
\end{enumerate}

\begin{lemma}
\label{pione-lemma-specialization-map-surjective}
Soit $f : X \to S$ un morphisme plat et propre à fibres géométriquement
connexes. Soit $s' \leadsto s$ une spécialisation.
Si $X_s$ est géométriquement réduit, alors l'homomorphisme de
spécialisation $sp : \pi_1(X_{\overline{s}'}) \to \pi_1(X_{\overline{s}})$
est surjectif.
\end{lemma}

\begin{proof}
Puisque $X_s$ est géométriquement réduit, on peut supposer que toutes les
fibres sont géométriquement réduites, quitte à rétrécir $S$, voir
Compléments sur les morphismes, lemme \ref{more-morphisms-lemma-geometrically-reduced-open}.
Soit $\mathcal{O}_{S, s} \to A \to \kappa(\overline{s}')$ comme
dans la construction de l'homomorphisme de spécialisation, voir la
section \ref{pione-section-specialization-map}.
Il suffit donc de montrer que
$$
\pi_1(X_{\overline{s}'}) \to \pi_1(X_A)
$$
est surjective. Cela résulte de la
proposition \ref{pione-proposition-first-homotopy-sequence}
et de $\pi_1(\Spec(A)) = \{1\}$.
\end{proof}

\begin{proposition}
\label{pione-proposition-specialization-map-isomorphism}
Soit $f : X \to S$ un morphisme propre et lisse à fibres géométriquement
connexes. Soit $s' \leadsto s$ une spécialisation.
Si la caractéristique de $\kappa(s)$ est nulle, alors l'homomorphisme de
spécialisation
$$
sp : \pi_1(X_{\overline{s}'}) \to \pi_1(X_{\overline{s}})
$$
est un isomorphisme.
\end{proposition}

\begin{proof}
L'homomorphisme est surjectif d'après le
lemme \ref{pione-lemma-specialization-map-surjective}.
Il reste donc à montrer qu'il est injectif.

\medskip\noindent
On peut supposer que $S$ est affine. Alors $S$ est une limite cofiltrante de schémas
affines de type fini sur $\mathbf{Z}$.
On peut donc supposer que $X \to S$ s'obtient par
changement de base à partir de $X_0 \to S_0$, où $S_0$ est le spectre d'une
$\mathbf{Z}$-algèbre de type fini et $X_0 \to S_0$ est propre et lisse.
Voir Limites, lemme \ref{limits-lemma-descend-finite-presentation},
\ref{limits-lemma-descend-smooth}, et
\ref{limits-lemma-eventually-proper}. D'après le
lemme \ref{pione-lemma-specialization-map-base-change},
on se ramène au cas où la base est noethérienne.

\medskip\noindent
En appliquant le lemme \ref{pione-lemma-specialization-map-discrete-valuation-ring},
on se ramène au cas où la base $S$ est le spectre d'un
anneau de valuation discrète strictement hensélien $A$ et où l'on
considère l'homomorphisme de spécialisation sur $A$.
Soit $K$ le corps des fractions de $A$.
Choisissons une clôture algébrique $\overline{K}$ qui
correspond à un point géométrique générique $\overline{\eta}$ de $\Spec(A)$.
Pour toute sous-extension finie séparable $L/K$ de $\overline{K}/K$, soit $B \subset L$ la
clôture intégrale de $A$ dans $L$. C'est un anneau de
valuation discrète d'après
Compléments d'algèbre, remarque \ref{more-algebra-remark-finite-separable-extension}.

\medskip\noindent
Soit $X \to \Spec(A)$ comme dans le paragraphe précédent.
Pour montrer l'injectivité de l'homomorphisme de spécialisation,
il suffit de prouver que tout revêtement fini
\'etale $V$ de $X_{\overline{\eta}}$ provient par changement
de base d'un revêtement fini \'etale $Y \to X$.
En effet, alors $\pi_1(X_{\overline{\eta}}) \to \pi_1(X) = \pi_1(X_s)$
est injective d'après le lemme \ref{pione-lemma-functoriality-galois-injective}.

\medskip\noindent
Étant donné $V$, on peut d'abord descendre $V$ en $V' \to X_{K^{sep}}$ d'après le
lemme \ref{pione-lemma-perfection}, puis en
$V'' \to X_L$ d'après le lemme \ref{pione-lemma-limit}.
Soit $Z \to X_B$ la normalisation de $X_B$ dans $V''$.
Observons que $Z$ est normal et que $Z_L = V''$ comme schémas
sur $X_L$. Ainsi, $Z \to X_B$ est fini \'etale sur
la fibre générique. Le problème est que nous ne savons pas si
$Z \to X_B$ est partout \'etale. Puisque $X \to \Spec(A)$
a des fibres lisses géométriquement connexes, on voit que
la fibre spéciale $X_s$ est géométriquement irréductible.
Ainsi, la fibre spéciale de $X_B \to \Spec(B)$ est irréductible ;
soit $\xi_B$ son point générique. Soient
$\xi_1, \ldots, \xi_r$ les points de $Z$ qui s'envoient sur
$\xi_B$. Notre premier problème (et, comme on le verra, le seul)
est alors que les extensions
$$
\mathcal{O}_{X_B, \xi_B} \subset \mathcal{O}_{Z, \xi_i}
$$
d'anneaux de valuation discrète peuvent être ramifiées. Soit $e_i$
l'indice de ramification de cette extension. Notons que, puisque la
caractéristique de $\kappa(s)$ est nulle, la ramification est modérée !

\medskip\noindent
Pour éliminer la ramification, nous allons choisir une nouvelle extension finie
séparable $K^{sep}/L'/L/K$ telle que l'indice de ramification
$e$ des extensions induites $B'/B$ soit divisible par $e_i$.
Considérons le changement de base normalisé $Z'$ de $Z$ relativement à
$\Spec(B') \to \Spec(B)$, voir la discussion dans
Compléments sur les morphismes, section \ref{more-morphisms-section-reduced-fibre-theorem}.
Soient $\xi_{i, j}$ les points de $Z'$ qui s'envoient sur $\xi_{B'}$
et sur $\xi_i$ dans $Z$. Alors les anneaux locaux
$$
\mathcal{O}_{Z', \xi_{i, j}}
$$
sont des localisés de la clôture intégrale de $\mathcal{O}_{Z, \xi_i}$
dans $L' \otimes_L F_i$, où $F_i$ est le corps des fractions de
$\mathcal{O}_{Z, \xi_i}$ ; les détails sont omis. Le lemme d'Abhyankar
(Compléments d'algèbre, lemme \ref{more-algebra-lemma-abhyankar})
nous dit donc que
$$
\mathcal{O}_{X_{B'}, \xi_{B'}} \subset \mathcal{O}_{Z', \xi_{i, j}}
$$
est non ramifiée. On conclut que le morphisme $Z' \to X_{B'}$
est \'etale en codimension $1$. La pureté du
lieu de branchement (lemme \ref{pione-lemma-purity})
montre donc que $Z' \to X_{B'}$ est fini \'etale !

\medskip\noindent
Cependant, puisque l'extension de corps résiduels induite par $A \to B'$
est triviale (le corps résiduel de $A$ étant algébriquement clos,
car séparablement clos de caractéristique nulle),
on conclut que $Z'$ provient par changement de base d'un revêtement fini \'etale
$Y \to X$ en appliquant deux fois le
lemme \ref{pione-lemma-finite-etale-on-proper-over-henselian}
(d'abord pour obtenir $Y$ sur $A$, puis pour démontrer que
le changement de base à $B'$ est isomorphe à $Z'$).
Ceci achève la démonstration.
\end{proof}

\noindent
Soit $G$ un groupe profini. Soit $p$ un nombre premier.
Le {\it quotient maximal premier à $p$} est par définition
$$
G' = \lim_{U \subset G\text{ ouvert, distingué, d'indice premier à }p} G/U
$$
Si $X$ est un schéma connexe et si $p$ est fixé, le quotient maximal
premier à $p$ de $\pi_1(X)$ est noté $\pi'_1(X)$.

\begin{theorem}
\label{pione-theorem-specialization-map-isomorphism-prime-to-p}
Soit $f : X \to S$ un morphisme propre et lisse à fibres géométriquement
connexes. Soit $s' \leadsto s$ une spécialisation.
Si la caractéristique de $\kappa(s)$ est $p$, alors l'application de
spécialisation
$$
sp : \pi_1(X_{\overline{s}'}) \to \pi_1(X_{\overline{s}})
$$
est surjective et induit un isomorphisme
$$
\pi'_1(X_{\overline{s}'}) \cong \pi'_1(X_{\overline{s}})
$$
des quotients maximaux premiers à $p$.
\end{theorem}

\begin{proof}
Cela se démontre exactement de la même manière que la
proposition \ref{pione-proposition-specialization-map-isomorphism},
avec les différences suivantes :
\begin{enumerate}
\item Étant donné $X/A$, on ne montre plus que le foncteur
$\textit{F\'Et}_X \to \textit{F\'Et}_{X_{\overline{\eta}}}$
est essentiellement surjectif. On montre seulement que les objets galoisiens
dont le groupe de Galois est d'ordre premier à $p$ appartiennent à l'image
essentielle. Cela suffit pour conclure à l'injectivité de
$\pi'_1(X_{\overline{s}'}) \to \pi'_1(X_{\overline{s}})$ par
exactement le même argument.
\item Les extensions
$\mathcal{O}_{X_B, \xi_B} \subset \mathcal{O}_{Z, \xi_i}$
sont modérément ramifiées, car l'extension associée des corps de
fractions est galoisienne, de groupe d'ordre premier à $p$. Voir
Compléments d'algèbre, lemme \ref{more-algebra-lemma-galois-conclusion}.
\item L'extension $\kappa_B/\kappa_A$ n'est plus
nécessairement triviale, mais elle est purement inséparable.
Ainsi, le morphisme $X_{\kappa_B} \to X_{\kappa_A}$
est un homéomorphisme universel et induit un isomorphisme
de groupes fondamentaux d'après la proposition \ref{pione-proposition-universal-homeomorphism}.
\end{enumerate}
\end{proof}













\section{Ramification modérée}
\label{pione-section-tame}

\noindent
Soit $f : X \to Y$ un morphisme fini \'etale de schémas de type fini
sur $\mathbf{Z}$. Il existe de nombreuses façons de définir ce que signifie pour $f$
d'être modérément ramifié à $\infty$. L'article \cite{Kerz-Schmidt}
examine dans quelle mesure ces notions coïncident.

\medskip\noindent
Dans cette section, nous examinons une question différente et plus élémentaire qui
précède la notion de ramification modérée à l'infini.
On comparera avec la discussion (légèrement différente)
de \cite{Grothendieck-Murre}.
Supposons donnés
\begin{enumerate}
\item un schéma localement noethérien $X$,
\item un ouvert dense $U \subset X$,
\item un morphisme fini \'etale $f : Y \to U$
\end{enumerate}
tels que, pour tout diviseur premier $Z \subset X$
tel que $Z \cap U = \emptyset$, l'anneau local $\mathcal{O}_{X, \xi}$
de $X$ au point générique $\xi$ de $Z$ soit un anneau de valuation discrète.
En prenant pour $K_\xi$ le corps des fractions de $\mathcal{O}_{X, \xi}$,
on obtient un carré cartésien
$$
\xymatrix{
\Spec(K_\xi) \ar[r] \ar[d] & U \ar[d] \\
\Spec(\mathcal{O}_{X, \xi}) \ar[r] & X
}
$$
de schémas. En particulier, on voit que $Y \times_U \Spec(K_\xi)$
est le spectre d'une algèbre finie séparable $L_\xi/K_\xi$.
On dit alors que
{\it $Y$ est non ramifié sur $X$ en codimension $1$},
resp.\ {\it $Y$ est modérément ramifié sur $X$ en codimension $1$},
si $L_\xi/K_\xi$ est non ramifiée, resp.\ modérément ramifiée,
relativement à $\mathcal{O}_{X, \xi}$ pour tout $(Z, \xi)$
comme ci-dessus ; voir Compléments d'algèbre, définition
\ref{more-algebra-definition-types-of-extensions}.
Plus précisément, on décompose $L_\xi$ en un produit d'extensions
de corps finies séparables de $K_\xi$ et l'on exige que chacune soit
non ramifiée, resp.\ modérément ramifiée, relativement à
$\mathcal{O}_{X, \xi}$.

\begin{lemma}
\label{pione-lemma-pullback-tame-codim1}
Soit $g : X' \to X$ un morphisme de schémas localement noethériens.
Soit $U \subset X$ un ouvert dense. Supposons que
\begin{enumerate}
\item $U' = g^{-1}(U)$ soit un ouvert dense de $X'$,
\item pour tout diviseur premier $Z \subset X$ tel que $Z \cap U = \emptyset$,
l'anneau local $\mathcal{O}_{X, \xi}$ de $X$ au point générique $\xi$
de $Z$ soit un anneau de valuation discrète,
\item pour tout diviseur premier $Z' \subset X'$
tel que $Z' \cap U' = \emptyset$, l'anneau local $\mathcal{O}_{X', \xi'}$
de $X'$ au point générique $\xi'$ de $Z'$ soit un anneau de valuation discrète,
\item si $\xi' \in X'$ est comme en (3), alors $\xi = g(\xi')$ soit comme en (2).
\end{enumerate}
Alors, si $f : Y \to U$ est fini \'etale et si
$Y$ est non ramifié, resp.\ modérément ramifié, sur $X$
en codimension $1$, alors $Y' = Y \times_X X' \to U'$ est fini \'etale
et $Y'$ est non ramifié, resp.\ modérément ramifié, sur $X'$ en codimension $1$.
\end{lemma}

\begin{proof}
Le seul fait intéressant de ce lemme est le résultat d'algèbre commutative
donné dans Compléments d'algèbre, lemme \ref{more-algebra-lemma-tame-goes-up}.
\end{proof}

\noindent
Avec la terminologie introduite ci-dessus, on peut reformuler les
résultats de pureté obtenus plus haut sous la forme agréable suivante.

\begin{lemma}
\label{pione-lemma-purity-one-divisor-general}
Soit $X$ un schéma localement noethérien. Soit $U \subset X$ ouvert et dense.
Soit $Y \to U$ un morphisme fini \'etale. Supposons que
\begin{enumerate}
\item $Y$ soit non ramifié sur $X$ en codimension $1$, et
\item $\mathcal{O}_{X, x}$ soit régulier pour tout $x \in X \setminus U$.
\end{enumerate}
Alors il existe un morphisme fini \'etale $Y' \to X$
dont la restriction à $U$ est $Y$.
\end{lemma}

\begin{proof}
Soit $\xi \in X \setminus U$ le point générique d'une composante irréductible
de $X \setminus U$ de codimension $1$. Alors $\mathcal{O}_{X, \xi}$ est un
anneau de valuation discrète. Comme dans la discussion ci-dessus, écrivons
$Y \times_U \Spec(K_\xi) = \Spec(L_\xi)$.
Notons $B_\xi$ la clôture intégrale de $\mathcal{O}_{X, \xi}$ dans
$L_\xi$. Notre hypothèse que $Y$ est non ramifié sur $X$ en codimension $1$
signifie que $\mathcal{O}_{X, \xi} \to B_\xi$ est fini \'etale.
On obtient ainsi $Y_\xi \to \Spec(\mathcal{O}_{X, \xi})$ fini
\'etale et un isomorphisme
$$
Y \times_U \Spec(K_\xi) \cong
Y_\xi \times_{\Spec(\mathcal{O}_{X, \xi})} \Spec(K_\xi)
$$
sur $\Spec(K_\xi)$.
D'après Limites, lemme \ref{limits-lemma-glueing-near-point},
on trouve un sous-schéma ouvert $U \subset U' \subset X$
contenant $\xi$ et un morphisme $Y' \to U'$ de présentation finie
dont la restriction à $U$ redonne $Y$ et
dont la restriction à $\Spec(\mathcal{O}_{X, \xi})$ redonne $Y_\xi$.
Enfin, le morphisme $Y' \to U'$ est fini \'etale
après avoir éventuellement remplacé $U'$ par un ouvert plus petit, d'après
Limites, lemme \ref{limits-lemma-glueing-near-point-properties}.
En répétant l'argument avec les autres points génériques de
$X \setminus U$ de codimension $1$, on peut supposer que l'on dispose
d'un morphisme fini \'etale $Y' \to U'$ prolongeant $Y \to U$
à un sous-schéma ouvert $U' \subset X$ contenant $U$
et tous les points de codimension $1$ de $X \setminus U$.
On termine en appliquant le lemme \ref{pione-lemma-extend-pure}
à $Y' \to U'$. En effet, tous les anneaux locaux $\mathcal{O}_{X, x}$
pour $x \in X \setminus U'$ sont réguliers et
vérifient $\dim(\mathcal{O}_{X, x}) \geq 2$. La pureté vaut donc
pour $\mathcal{O}_{X, x}$ d'après le lemme \ref{pione-lemma-local-purity}.
\end{proof}

\begin{lemma}
\label{pione-lemma-purity-one-divisor}
Soit $X$ un schéma localement noethérien. Soit $D \subset X$
un diviseur de Cartier effectif tel que $D$ soit un schéma régulier.
Soit $Y \to X \setminus D$ un morphisme fini \'etale.
Si $Y$ est non ramifié sur $X$ en codimension $1$, alors
il existe un morphisme fini \'etale $Y' \to X$
dont la restriction à $X \setminus D$ est $Y$.
\end{lemma}

\begin{proof}
C'est un cas particulier du lemme \ref{pione-lemma-purity-one-divisor-general}.
Premièrement, $D$ n'est dense nulle part dans $X$ (voir la discussion dans
Diviseurs, section \ref{divisors-section-effective-Cartier-divisors}),
donc $X \setminus D$ est dense dans $X$. Deuxièmement, l'anneau
$\mathcal{O}_{X, x}$ est un anneau local régulier pour tout $x \in D$ d'après
Algèbre, lemme \ref{algebra-lemma-regular-mod-x},
et notre hypothèse selon laquelle $\mathcal{O}_{D, x}$ est régulier.
\end{proof}

\begin{example}[Morphisme standard modérément ramifié]
\label{pione-example-tamely-ramified}
Soit $A$ un anneau noethérien. Soit $f \in A$ un non diviseur de zéro
tel que $A/fA$ soit réduit. Cela implique que $A_\mathfrak p$
est un anneau de valuation discrète d'uniformisante $f$ pour tout
idéal premier minimal $\mathfrak p$ contenant $f$. Soit $e \geq 1$ un entier
inversible dans $A$. Posons
$$
C = A[x]/(x^e - f)
$$
Alors $\Spec(C) \to \Spec(A)$ est un morphisme fini localement libre
qui est \'etale sur le spectre de $A_f$. Le morphisme fini \'etale
$$
\Spec(C_f) \longrightarrow \Spec(A_f)
$$
est modérément ramifié sur $\Spec(A)$ en codimension $1$. La
modération résulte immédiatement de la caractérisation des extensions
modérément ramifiées dans
Compléments d'algèbre, lemme \ref{more-algebra-lemma-characterize-tame}.
\end{example}

\noindent
Voici une version du lemme d'Abhyankar pour les diviseurs réguliers.

\begin{lemma}[Lemme d'Abhyankar pour un diviseur régulier]
\label{pione-lemma-abhyankar-one-divisor}
Soit $X$ un schéma localement noethérien. Soit $D \subset X$
un diviseur de Cartier effectif tel que $D$ soit un schéma régulier.
Soit $Y \to X \setminus D$ un morphisme fini \'etale.
Si $Y$ est modérément ramifié sur $X$ en codimension $1$, alors,
localement pour la topologie \'etale sur $X$, le morphisme $Y \to X \setminus D$ est la restriction
à $X \setminus D$ d'une réunion disjointe finie de morphismes standard modérément ramifiés
tels que ceux décrits dans l'exemple \ref{pione-example-tamely-ramified}.
\end{lemma}

\begin{proof}
Pour tout $x \in X$, nous allons trouver un voisinage \'etale
$(U, u) \to (X, x)$, avec $U = \Spec(A)$, tel que le changement de base
$Y \times_{X \setminus D} (U \setminus D) \to U \setminus D$ soit une réunion disjointe finie de restrictions de morphismes standard
modérément ramifiés comme dans l'exemple \ref{pione-example-tamely-ramified}.
On supposera $x \in D$ ; le cas $x \not \in D$ résulte de
Morphismes \'etales, lemme \ref{etale-lemma-finite-etale-etale-local},
en prenant $e = 1$ et $f = 1$ dans l'exemple \ref{pione-example-tamely-ramified}.

\medskip\noindent
Dans ce paragraphe, on se ramène au cas où $X$ est le spectre d'un
anneau local strictement hensélien et où $x$ est le point fermé. En effet, quitte à rétrécir
$X$, on peut supposer que $X = \Spec(A)$ et que $D \subset X$ est défini par un
non diviseur de zéro $f \in A$. Alors $Y$ est affine, car c'est un revêtement fini
\'etale de $\Spec(A_f)$. Écrivons $Y = \Spec(B)$.
Soit $A^{sh}$ l'hensélisé strict de $\mathcal{O}_{X, x}$.
Puisque $A \to A^{sh}$ est plat et que $x \in D$, on voit que $f$ s'envoie
sur un non diviseur de zéro appartenant à l'idéal maximal de $A^{sh}$.
Observons que $A^{sh}/fA^{sh} = (A/fA)^{sh}$
(Algèbre, lemme \ref{algebra-lemma-quotient-strict-henselization})
est un anneau local régulier, car c'est l'hensélisé strict de $\mathcal{O}_{D, x}$ ;
voir Compléments d'algèbre, lemme \ref{more-algebra-lemma-henselization-regular}.
D'après le lemme \ref{pione-lemma-pullback-tame-codim1},
le changement de base de $Y$ à $\Spec(A^{sh})$ est modérément ramifié en
codimension $1$. Supposons l'assertion démontrée pour le changement de base
de $Y$ à $\Spec(A^{sh})$. Puisque $A^{sh}$ est strictement hensélien, tout
voisinage \'etale possède une section, et l'on conclut que l'on a un isomorphisme
$$
A^{sh} \otimes_A B \cong
\prod\nolimits_{i \in I} A^{sh}_f[x]/(x^{e_i} - f)
$$
où $I$ est un ensemble fini et chaque $e_i$ un entier inversible dans $A^{sh}$.
Écrivons $A^{sh} = \colim A_\lambda$ comme une limite inductive filtrante d'algèbres \'etales
sur $A$, notées $A_\lambda$. L'isomorphisme affiché
descend en un isomorphisme
$$
A_\lambda \otimes_A B \cong \prod\nolimits_{i \in I}
(A_\lambda)_f[x]/(x^{e_i} - f)
$$
pour $\lambda$ assez grand, voir par exemple
Algèbre, lemme \ref{algebra-lemma-colimit-category-fp-algebras}.
Après avoir encore un peu augmenté $\lambda$,
on peut supposer que $e_i$ est inversible dans $A_\lambda$. Alors
$\Spec(A_\lambda) \to \Spec(A)$ est le voisinage \'etale recherché
de $x$.

\medskip\noindent
Supposons que $X = \Spec(A)$, où $A$ est un anneau local strictement hensélien,
et que $x \in X$ corresponde à l'idéal maximal $\mathfrak m$ de $A$.
Soit $f \in \mathfrak m$ le non diviseur de zéro qui définit $D$.
Alors $A/fA$ est régulier et, d'après Algèbre, lemme \ref{algebra-lemma-regular-mod-x},
on voit que $A$ est lui aussi régulier. Nous utiliserons certaines propriétés des anneaux
locaux réguliers, par exemple le fait qu'ils sont des anneaux intègres normaux ; voir
Algèbre, lemmes \ref{algebra-lemma-regular-domain} et
\ref{algebra-lemma-regular-normal}. En particulier, $A$ est intègre.
Comme ci-dessus, on voit que $Y = \Spec(B)$ est affine et que $A_f \to B$
est fini \'etale. Ainsi, $B$ est lui aussi normal
(Algèbre, lemme \ref{algebra-lemma-normal-goes-up}),
et l'on conclut que $Y$ est une réunion disjointe finie de spectres d'anneaux intègres normaux d'après
Algèbre, lemme \ref{algebra-lemma-characterize-reduced-ring-normal}.
On peut donc supposer que $B$ est lui aussi intègre.
Puisque $A$ et $A/fA$ sont intègres, on conclut que $f$ est un élément
premier de $A$ qui engendre un idéal premier de hauteur $1$, $\mathfrak p = fA$.
Observons que $A_\mathfrak p$ est un anneau de valuation discrète d'uniformisante $f$.

\medskip\noindent
Soit $K$ le corps des fractions de $A$ et soit $L$ le corps des
fractions de $B$. L'hypothèse de ramification modérée signifie que
$L$ est modérément ramifié relativement à $A_\mathfrak p$.
Soit $e \geq 1$ le produit des indices de ramification de $L$
sur $A_\mathfrak p$, comme dans
Compléments d'algèbre, remarque \ref{more-algebra-remark-finite-separable-extension}.
Alors $e$ est inversible dans $\kappa(\mathfrak p)$, mais
à ce stade, nous ne savons pas si $e$ est inversible dans $A/fA$
ou dans $A$.

\medskip\noindent
Considérons la $A$-algèbre finie et libre
$$
A' = A[x]/(x^e - f)
$$
Il s'agit d'une extension locale finie d'un anneau local
strictement hensélien ; elle est donc strictement hensélienne.
Observons que $f' = x$ est un non diviseur de zéro dans $A'$ et que
$A'/f'A' \cong A/fA$ est un anneau régulier. Comme précédemment, $A'$
est donc régulier et, a fortiori, intègre. Posons
$B' = B \otimes_A A' = B \otimes_{A_f} A'_{f'}$.
D'après le lemme d'Abhyankar
(Compléments d'algèbre, lemme \ref{more-algebra-lemma-abhyankar}),
on voit que $\Spec(B')$ est non ramifié sur $\Spec(A')$
en codimension $1$. En effet, d'après le lemme \ref{pione-lemma-pullback-tame-codim1},
on voit que $\Spec(B')$ est encore au moins modérément ramifié
sur $\Spec(A')$ en codimension $1$. Mais le lemme d'Abhyankar
nous dit que les indices de ramification sont tous devenus égaux à $1$.
D'après le lemme \ref{pione-lemma-purity-one-divisor}, on conclut que
$\Spec(B') \to \Spec(A'_{f'})$ se prolonge en un morphisme fini \'etale
$\Spec(C) \to \Spec(A')$. Cependant, puisque $A'$ est strictement hensélien,
on conclut que $\Spec(C)$ est une réunion disjointe finie de copies
de $\Spec(A')$. Conclusion : il existe au moins un morphisme
$\Spec(A'_{f'}) \to \Spec(B')$. On conclut qu'il existe
une inclusion $B \subset A'_{f'}$ de $A_f$-algèbres.

\medskip\noindent
Il s'ensuit que l'on a des inclusions $K \subset L \subset K[f^{1/e}]$.
D'après Compléments d'algèbre, lemme \ref{more-algebra-lemma-pull-root-uniformizer},
on conclut que $L$ est égal à $K[f^{1/n}]$ pour un diviseur $n$ de $e$.
La normalité de $B$ implique alors que $B \cong A_f[y]/(y^n - f)$ ; les détails
sont omis.
Cependant, le lieu de ramification de l'homomorphisme $A_f \to A_f[y]/(y^n - f)$
est défini par $ny^{n - 1}$, ce qui implique nécessairement que $n$ est une unité de $A_f$.
Puisque $f$ est un élément premier, cela signifie que $n = u f^s$ pour
une unité $u$ de $A$ et $s \in \mathbf{Z}$. Puisque $n$ s'envoie sur une unité de
$\kappa(\mathfrak p)$, on obtient $s = 0$, ce qui achève la démonstration.
\end{proof}

\begin{lemma}
\label{pione-lemma-extend-tame-covering-normal}
Dans la situation du lemme \ref{pione-lemma-abhyankar-one-divisor},
la normalisation de $X$ dans $Y$ est un morphisme fini localement libre
$\pi : Y' \to X$ tel que
\begin{enumerate}
\item la restriction de $Y'$ à $X \setminus D$ soit isomorphe à $Y$,
\item $D' = \pi^{-1}(D)_{red}$ soit un diviseur de Cartier effectif sur $Y'$, et
\item $D'$ soit un schéma régulier.
\end{enumerate}
De plus, localement pour la topologie \'etale sur $X$, le morphisme $Y' \to X$ est une réunion
disjointe finie de morphismes
$$
\Spec(A[x]/(x^e - f)) \to \Spec(A)
$$
où $A$ est un anneau noethérien, $f \in A$ un non diviseur de zéro tel que
$A/fA$ soit régulier, et $e \geq 1$ est inversible dans $A$.
\end{lemma}

\begin{proof}
Ce n'est qu'un complément au lemme \ref{pione-lemma-abhyankar-one-divisor},
et, en fait, l'énoncé du présent lemme découle presque immédiatement de
la démonstration de ce lemme. Mais on peut également déduire
le présent lemme du résultat du lemme \ref{pione-lemma-abhyankar-one-divisor}.
En effet, la normalisation de $X$ dans $Y$ commute au
changement de base \'etale, voir Compléments sur les morphismes, lemme
\ref{more-morphisms-lemma-normalization-smooth-localization}.
On voit donc que l'on peut démontrer les assertions sur la structure locale
de $Y' \to X$ localement pour la topologie \'etale sur $X$. Ainsi, d'après le
lemme \ref{pione-lemma-abhyankar-one-divisor}, on peut supposer que
$X = \Spec(A)$, où $A$ est un anneau noethérien, que l'on a un
non diviseur de zéro $f\in A$ tel que $A/fA$ soit régulier, et que $Y$
est une réunion disjointe finie de spectres d'anneaux $A_f[x]/(x^e - f)$
où $e$ est inversible dans $A$. Nous omettons la vérification que
la clôture intégrale de $A$ dans $A_f[x]/(x^e - f)$ est
égale à $A' = A[x]/(x^e - f)$. (Pour le voir, montrer que
les localisés de $A'$ aux idéaux premiers au-dessus de $(f)$ sont réguliers.)
Nous omettons les détails.
\end{proof}

\begin{lemma}
\label{pione-lemma-tame-covering-split}
Dans la situation du lemme \ref{pione-lemma-abhyankar-one-divisor},
soit $Y' \to X$ comme dans le lemme \ref{pione-lemma-extend-tame-covering-normal}.
Soit $R$ un anneau de valuation discrète de corps des fractions $K$.
Soit
$$
t : \Spec(R) \to X
$$
un morphisme tel que l'image réciproque schématique
$t^{-1}D$ soit le point fermé réduit de $\Spec(R)$.
\begin{enumerate}
\item Si $t|_{\Spec(K)}$ se relève en un point de $Y$, alors
on obtient un relèvement $t' : \Spec(R) \to Y'$ tel que $Y' \to X$
soit \'etale le long de $t'(\Spec(R))$.
\item Si $\Spec(K) \times_X Y$ est isomorphe à une réunion disjointe
de copies de $\Spec(K)$, alors $Y' \to X$ est fini \'etale
sur un voisinage ouvert de $t(\Spec(R))$.
\end{enumerate}
\end{lemma}

\begin{proof}
En appliquant le critère valuatif de propreté au
morphisme fini $Y' \to X$, on voit que les points à valeurs dans $\Spec(K)$
de $Y$ qui coïncident avec $t|_{\Spec(K)}$ comme morphismes vers $X$
se relèvent de manière unique en morphismes $t' : \Spec(R) \to Y'$.
Ainsi, l'énoncé (1) a un sens.

\medskip\noindent
Choisissons un voisinage \'etale $(U, u) \to (X, t(\mathfrak m_R))$
tel que $U = \Spec(A)$ et que $Y' \times_X U \to U$
admette une description comme dans le lemme \ref{pione-lemma-extend-tame-covering-normal}
pour un certain $f \in A$. Alors $\Spec(R) \times_X U \to \Spec(R)$ est \'etale
et surjectif. Si $R'$ désigne l'anneau local de
$\Spec(R) \times_X U$ situé au-dessus du point fermé de $\Spec(R)$,
alors $R'$ est un anneau de valuation discrète et $R \subset R'$
est une extension non ramifiée d'anneaux de valuation discrète
(Compléments d'algèbre, lemme \ref{more-algebra-lemma-Dedekind-etale-extension}).
L'hypothèse sur $t$ signifie que l'homomorphisme $A \to R'$
correspondant à
$$
\Spec(R') \to \Spec(R) \times_X U \to U
$$
envoie $f$ sur une uniformisante $\pi \in R'$. Supposons maintenant que
$$
Y' \times_X U =
\coprod\nolimits_{i \in I} \Spec(A[x]/(x^{e_i} - f))
$$
pour certains $e_i \geq 1$. Alors on voit que
$$
\Spec(R') \times_U (Y' \times_X U) =
\coprod\nolimits_{i \in I} \Spec(R'[x]/(x^{e_i} - \pi))
$$
Les anneaux $R'[x]/(x^{e_i} - \pi)$ sont des anneaux de valuation discrète
(Compléments d'algèbre, lemme \ref{more-algebra-lemma-pull-root-uniformizer})
et n'admettent donc aucun morphisme vers le corps des fractions de $R'$, sauf si $e_i = 1$.

\medskip\noindent
Démonstration de (1). Dans ce cas, le changement de base de l'application $t' : \Spec(R) \to Y'$
détermine un morphisme correspondant $t'' : \Spec(R') \to Y' \times_X U$
dont l'image doit être contenue dans une composante correspondant à $i \in I$ avec $e_i = 1$
d'après la discussion ci-dessus. On voit donc clairement que $Y' \times_X U \to U$
est \'etale le long de l'image de $t''$. Comme la propriété d'être \'etale
se vérifie après un changement de base \'etale, ceci démontre (1).

\medskip\noindent
Démonstration de (2). Dans ce cas, l'hypothèse implique que $e_i = 1$
pour tout $i \in I$. Ainsi, $Y' \times_X U \to U$ est fini \'etale
et l'on conclut comme ci-dessus.
\end{proof}

\begin{lemma}
\label{pione-lemma-extend-covering}
Soit $S$ un schéma noethérien intègre normal de point générique $\eta$.
Soit $f : X \to S$ un morphisme lisse à fibres géométriquement connexes.
Soit $\sigma : S \to X$ une section de $f$. Soit $Z \to X_\eta$ un
revêtement galoisien fini \'etale (section \ref{pione-section-finite-etale-under-galois})
de groupe $G$, dont l'ordre est inversible sur $S$, tel que
$Z$ possède un point $\kappa(\eta)$-rationnel qui s'envoie sur $\sigma(\eta)$.
Alors il existe un revêtement galoisien fini \'etale $Y \to X$ de groupe $G$
dont la restriction à $X_\eta$ est $Z$.
\end{lemma}

\begin{proof}
Supposons d'abord que $S = \Spec(R)$ soit le spectre d'un anneau de valuation discrète $R$
de point fermé $s \in S$. Alors $X_s$ est un diviseur de Cartier effectif
dans $X$, et $X_s$ est régulier, car c'est un schéma lisse sur un corps.
De plus, la fibre générique $X_\eta$ est le sous-schéma ouvert $X \setminus X_s$.
Il résulte de
Compléments d'algèbre, lemme \ref{more-algebra-lemma-galois-conclusion},
et de l'hypothèse sur $G$ que $Z$ est modérément ramifié
sur $X$ en codimension $1$. Soit $Z' \to X$ comme dans le
lemme \ref{pione-lemma-extend-tame-covering-normal}. Observons que
l'action de $G$ sur $Z$ se prolonge en une action de $G$ sur $Z'$.
D'après le lemme \ref{pione-lemma-tame-covering-split},
on voit que $Z' \to X$ est fini \'etale sur un voisinage
ouvert de $\sigma(s)$.
Puisque $X_s$ est irréductible, cela implique que $Z \to X_\eta$
est non ramifié sur $X$ en codimension $1$.
On obtient alors un morphisme fini \'etale $Y \to X$
dont la restriction à $X_\eta$ est $Z$, d'après le
lemme \ref{pione-lemma-purity-one-divisor}.
Bien sûr, $Y \cong Z'$ (détails omis ; indication : calculer localement pour la topologie \'etale),
et $Y$ est donc un revêtement galoisien de groupe $G$.

\medskip\noindent
Cas général. Soit $U \subset S$ un sous-schéma ouvert maximal
tel qu'il existe un revêtement galoisien fini \'etale
$Y \to X \times_S U$ de groupe $G$
dont la restriction à $X_\eta$ soit isomorphe à $Z$.
Supposons que $U \not = S$ afin d'obtenir une contradiction.
Soit $s \in S \setminus U$ un point générique d'une composante irréductible
de $S \setminus U$. Alors l'image inverse $U_s$
de $U$ dans $\Spec(\mathcal{O}_{S, s})$
est le spectre épointé de $\mathcal{O}_{S, s}$.
Nous affirmons que $Y \times_S U_s \to X \times_S U_s$
est la restriction d'un revêtement galoisien fini \'etale
$Y'_s \to X \times_S \Spec(\mathcal{O}_{S, s})$
de groupe $G$.

\medskip\noindent
Montrons d'abord que l'affirmation donne la contradiction recherchée.
D'après Limites, lemme \ref{limits-lemma-glueing-near-point},
on trouve un sous-schéma ouvert $U \subset U' \subset S$
contenant $s$ et un morphisme $Y'' \to U'$ de présentation finie
dont la restriction à $U$ redonne $Y \to U$ et
dont la restriction à $\Spec(\mathcal{O}_{S, s})$ redonne $Y'_s$.
De plus, par l'équivalence de catégories donnée dans le
lemme, on peut supposer, après avoir rétréci $U'$,
qu'il existe un morphisme $Y'' \to U' \times_S X$
et une action de $G$ sur $Y''$ au-dessus de $U' \times_S X$,
compatibles avec les morphismes et actions donnés après
changement de base à $U$ et à $\Spec(\mathcal{O}_{S, s})$.
Après avoir encore rétréci $U'$ si nécessaire, on peut
supposer que $Y'' \to U' \times_S X$ est fini \'etale, voir
Limites, lemme \ref{limits-lemma-glueing-near-point-properties}.
Cela signifie que nous avons trouvé un ouvert strictement plus grand de $S$ sur lequel
$Y$ se prolonge en un revêtement galoisien fini \'etale de groupe $G$,
ce qui donne la contradiction recherchée.

\medskip\noindent
Démonstration de l'affirmation. On peut remplacer, et l'on remplace, $S$ par $\Spec(\mathcal{O}_{S, s})$.
Alors $S = \Spec(A)$, où $(A, \mathfrak m)$ est un anneau intègre local normal.
De plus, $U \subset S$ est le spectre épointé et l'on dispose d'un revêtement
galoisien fini \'etale $Y \to X \times_S U$ de groupe $G$.
Si $\dim(A) = 1$, on peut construire le prolongement de $Y$ en
un revêtement galoisien de $X$ par le premier paragraphe de la démonstration.
On peut donc supposer $\dim(A) \geq 2$, et donc $\text{depth}(A) \geq 2$,
puisque $S$ est normal, voir Algèbre, lemme \ref{algebra-lemma-criterion-normal}.
Puisque $X \to S$ est plat, on conclut que
$\text{depth}(\mathcal{O}_{X, x}) \geq 2$ pour tout point $x \in X$
qui s'envoie sur $s$, voir
Algèbre, lemme \ref{algebra-lemma-apply-grothendieck}.
Soit
$$
Y' \longrightarrow X
$$
le morphisme fini construit dans le lemme \ref{pione-lemma-extend-S2}
à partir de $Y \to X \times_S U$. Observons que l'on obtient une action canonique
de $G$ sur $Y'$. Il reste donc seulement à montrer que $Y'$
est \'etale sur $X$. En fait, d'après le
lemme \ref{pione-lemma-purity-smooth-over-depth2} (par exemple),
il suffit même de montrer que $Y' \to X$ est \'etale au-dessus du
point générique (unique) de la fibre $X_s$. Nous le faisons par
un calcul local dans un voisinage (formel) de $\sigma(s)$.

\medskip\noindent
Choisissons un ouvert affine $\Spec(B) \subset X$ contenant $\sigma(s)$.
Alors $A \to B$ est un homomorphisme d'anneaux lisse qui admet une section $\sigma : B \to A$.
Notons $I = \Ker(\sigma)$ et notons $B^\wedge$ le
complété $I$-adique de $B$. Alors $B^\wedge \cong A[[x_1, \ldots, x_d]]$
pour un certain $d \geq 0$, voir
Algèbre, lemme \ref{algebra-lemma-section-smooth}.
Bien sûr, $B \to B^\wedge$ est plat
(Algèbre, lemme \ref{algebra-lemma-completion-flat}),
et l'image de $\Spec(B^\wedge) \to X$ contient le point générique de $X_s$.
Soit $V \subset \Spec(B^\wedge)$ l'image inverse de $U$.
Considérons le morphisme fini \'etale
$$
W = Y \times_{(X \times_S U)} V \longrightarrow V
$$
Par la compatibilité de la construction de $Y'$
au changement de base plat dans le lemme \ref{pione-lemma-extend-S2},
on voit que le changement de base
$Y' \times_X \Spec(B^\wedge) \to \Spec(B^\wedge)$
est construit à partir de $W \to V$ sur $\Spec(B^\wedge)$
par le procédé du lemme \ref{pione-lemma-extend-S2}.
Posons $V_0 = V \cap V(x_1, \ldots, x_d) \subset V$ et $W_0 = W \times_V V_0$.
Le schéma $V_0$ est intègre normal, s'envoie dans $\sigma(S)$
par le morphisme $\Spec(B^\wedge) \to X$ et s'identifie en fait
à $\sigma(U)$. On sait donc que $W_0 \to V_0 = U$ se scinde complètement,
car c'est le cas de sa fibre générique d'après notre hypothèse
selon laquelle $Z \to X_\eta$ possède un point $\kappa(\eta)$-rationnel au-dessus de
$\sigma(\eta)$ (et, bien sûr, l'action de $G$ implique alors que toute la
fibre $Z_{\sigma(\eta)}$ est une réunion disjointe de copies du
schéma $\eta = \Spec(\kappa(\eta))$).
Enfin, d'après le
lemme \ref{pione-lemma-fully-faithful-power-series-over-depth2},
on a
$$
W_0 \times_U V \cong W
$$
Cela montre que $W$ est une réunion disjointe de copies de $V$
et donc que $Y' \times_X \Spec(B^\wedge)$ est une réunion disjointe
de copies de $\Spec(B^\wedge)$, ce qui achève la démonstration.
\end{proof}

\begin{lemma}
\label{pione-lemma-extend-covering-general}
Soit $S$ un schéma intègre normal quasi-compact et quasi-séparé
de point générique $\eta$. Soit $f : X \to S$ un morphisme lisse quasi-compact et
quasi-séparé à fibres géométriquement connexes.
Soit $\sigma : S \to X$ une section de $f$. Soit $Z \to X_\eta$ un
revêtement galoisien fini \'etale (section \ref{pione-section-finite-etale-under-galois})
de groupe $G$, dont l'ordre est inversible sur $S$, tel que
$Z$ possède un point $\kappa(\eta)$-rationnel qui s'envoie sur $\sigma(\eta)$.
Alors il existe un revêtement galoisien fini \'etale $Y \to X$ de groupe $G$
dont la restriction à $X_\eta$ est $Z$.
\end{lemma}

\begin{proof}
Si $S$ est noethérien, c'est le résultat du
lemme \ref{pione-lemma-extend-covering}. Le cas général
en découle par un argument standard de passage à la limite.
Nous recommandons vivement au lecteur de sauter la démonstration.

\medskip\noindent
On peut écrire $S = \lim S_i$ comme limite projective d'un système filtrant
de schémas dont les morphismes de transition sont affines et dont les $S_i$ sont
de type fini sur $\mathbf{Z}$, voir
Limites, proposition \ref{limits-proposition-approximate}.
Pour chaque $i$, soit $S \to S'_i \to S_i$ la normalisation
de $S_i$ dans $S$, voir
Morphismes, section \ref{morphisms-section-normalization-X-in-Y}.
En combinant Algèbre, proposition \ref{algebra-proposition-ubiquity-nagata},
avec Morphismes, lemmes \ref{morphisms-lemma-nagata-normalization-finite} et
\ref{morphisms-lemma-normal-normalization},
on conclut que $S'_i$ est de type fini sur $\mathbf{Z}$,
fini sur $S_i$, et que $S'_i$ est un schéma intègre normal
tel que $S \to S'_i$ soit dominant.
D'après Morphismes, lemme \ref{morphisms-lemma-functoriality-normalization},
on obtient des morphismes de transition $S'_{i'} \to S'_i$ compatibles
avec les morphismes de transition $S_{i'} \to S_i$ et avec
les morphismes de source $S$.
Nous affirmons que $S = \lim S'_i$. Démonstration de l'affirmation omise (indication :
considérer les ouverts affines situés au-dessus d'un ouvert affine choisi de $S_i$
pour un certain $i$, afin de traduire ceci en un simple problème
d'algèbre). On conclut que l'on peut écrire $S = \lim S_i$
comme limite projective d'un système filtrant de schémas intègres normaux $S_i$
dont les morphismes de transition sont affines et dont les $S_i$ sont
de type fini sur $\mathbf{Z}$.

\medskip\noindent
Pour un certain $i$, on peut trouver un morphisme lisse $X_i \to S_i$
de présentation finie dont le changement de base à $S$ est $X \to S$.
Voir Limites, lemmes
\ref{limits-lemma-descend-finite-presentation} et
\ref{limits-lemma-descend-smooth}.
En augmentant $i$, on peut supposer que la section
$\sigma$ se relève en une section $\sigma_i : S_i \to X_i$
(par l'équivalence de catégories de
Limites, lemme \ref{limits-lemma-descend-finite-presentation}).
On peut remplacer $X_i$ par son sous-schéma ouvert $X_i^0$
étudié dans Compléments sur les morphismes, section
\ref{more-morphisms-section-connected-components}
car l'image de $X \to X_i$ s'y envoie manifestement
(l'ouverture résulte de Compléments sur les morphismes, lemme
\ref{more-morphisms-lemma-connected-along-section-open}).
On peut donc supposer que les fibres de $X_i \to S_i$ sont
géométriquement connexes.
En augmentant $i$, on peut supposer que $|G|$ est inversible
sur $S_i$.
Soit $\eta_i \in S_i$ le point générique.
Puisque $X_\eta$ est la limite des schémas
$X_{i, \eta_i}$, on peut employer exactement les mêmes arguments
pour faire descendre $Z \to X_\eta$ en un revêtement galoisien fini \'etale
$Z_i \to X_{i, \eta_i}$, quitte à augmenter $i$.
Voir le lemme \ref{pione-lemma-limit}.
Quitte à augmenter encore une fois $i$, on peut
supposer que $Z_i$ possède un point $\kappa(\eta_i)$-rationnel
s'envoyant sur $\sigma_i(\eta_i)$.
On applique alors le lemme dans le cas noethérien
et on effectue le changement de base à $X$ pour conclure.
\end{proof}











\section{Astuces en caractéristique positive}
\label{pione-section-achinger}

\noindent
Dans l'article de Piotr Achinger \cite{Achinger}, il est montré qu'un
schéma affine en caractéristique positive est toujours un $K(\pi, 1)$.
Dans cette section, nous expliquons les parties les plus élémentaires de \cite{Achinger}.
Plus précisément, nous montrons que, pour un corps $k$ de caractéristique positive,
un schéma affine \'etale sur $\mathbf{A}^n_k$ est en fait fini \'etale
sur $\mathbf{A}^n_k$ (par un autre morphisme). Nous montrons aussi
qu'une immersion fermée entre schémas affines connexes
en caractéristique positive induit un homomorphisme injectif entre les
groupes fondamentaux \'etales.

\medskip\noindent
Soit $k$ un corps de caractéristique $p > 0$.
Considérons
$$
k[x_1, \ldots, x_n] \longrightarrow A
$$
une surjection de $k$-algèbres de type fini dont la source est
l'algèbre de polynômes en $x_1, \ldots, x_n$. Notons
$I \subset k[x_1, \ldots, x_n]$ le noyau, de sorte que l'on ait
$A = k[x_1, \ldots, x_n]/I$. Nous ne supposons pas que $A$ soit non nul
(autrement dit, nous autorisons le cas où $A$ est l'anneau nul
et $I = k[x_1, \ldots, x_n]$). Enfin, nous supposons donné un
homomorphisme d'anneaux fini \'etale $\pi : A \to B$.

\medskip\noindent
Supposons donnés $k, n, k[x_1, \ldots, x_n] \to A, I, \pi : A \to B$.
Soit $C$ une $k$-algèbre. Considérons les
diagrammes commutatifs
$$
\xymatrix{
& B \\
C \ar[r] & C/\varphi(I)C \ar[u]^\tau \\
k[x_1, \ldots, x_n] \ar[u]^\varphi \ar[r] &
A \ar[u] \ar@/_3em/[uu]_\pi
}
$$
où $\varphi$ est un homomorphisme \'etale de $k$-algèbres et $\tau$ un
homomorphisme surjectif de $k$-algèbres.
Fixons $C, \varphi, \tau$. Pour tout $r \geq 0$ et tous
$y_1, \ldots, y_r \in C$ qui engendrent $C$ comme algèbre sur
$\Im(\varphi)$, soit $s = s(r, y_1, \ldots, y_r) \in \{0, \ldots, r\}$
le plus grand entier tel que $y_i$ soit entier sur $\Im(\varphi)$
pour $1 \leq i \leq s$. On définit $NF(C, \varphi, \tau)$
comme la valeur minimale de
$r - s = r - s(r, y_1, \ldots, y_r)$ pour tous les choix de
$r$ et $y_1, \ldots, y_r$ comme ci-dessus.
Observons que $NF(C, \varphi, \tau)$ vaut $0$ si et seulement si $\varphi$ est fini.

\begin{lemma}
\label{pione-lemma-lower-invariant}
Dans la situation ci-dessus, si $NF(C, \varphi, \tau) > 0$, alors il existe
un homomorphisme \'etale de $k$-algèbres $\varphi'$ et un homomorphisme surjectif de $k$-algèbres
$\tau'$ s'insérant dans le diagramme commutatif
$$
\xymatrix{
& B \\
C \ar[r] & C/\varphi'(I)C \ar[u]_{\tau'} \\
k[x_1, \ldots, x_n] \ar[u]^{\varphi'} \ar[r] &
A \ar[u] \ar@/_3em/[uu]_\pi
}
$$
tel que $NF(C, \varphi', \tau') < NF(C, \varphi, \tau)$.
\end{lemma}

\begin{proof}
Choisissons $r \geq 0$ et $y_1, \ldots, y_r \in C$ qui engendrent
$C$ sur $\Im(\varphi)$, et soit $0 \leq s \leq r$ tel que
$y_1, \ldots, y_s$ soient entiers sur $\Im(\varphi)$
et que $r - s = NF(C, \varphi, \tau) > 0$.
Comme $B$ est fini sur $A$, l'image de $y_{s + 1}$ dans $B$
annule un polynôme unitaire à coefficients dans $A$. On peut donc trouver
un entier $d \geq 1$ et des éléments $f_1, \ldots, f_d \in k[x_1, \ldots, x_n]$
tels que
$$
z = y_{s + 1}^d + \varphi(f_1) y_{s + 1}^{d - 1} + \ldots + \varphi(f_d) \in
J = \Ker(C \to C/\varphi(I)C \xrightarrow{\tau} B)
$$
Comme $\varphi : k[x_1, \ldots, x_n] \to C$ est \'etale, il existe un
polynôme non nul et non constant $g \in k[T_1, \ldots, T_{n + 1}]$
tel que
$$
g(\varphi(x_1), \ldots, \varphi(x_n), z) = 0
\quad\text{dans}\quad
C
$$
Pour le voir, on peut utiliser par exemple le fait que
$C \otimes_{\varphi, k[x_1, \ldots, x_n]} k(x_1, \ldots, x_n)$
est un produit fini d'extensions de corps finies séparables
de $k(x_1, \ldots, x_n)$
(voir Algèbre, lemme \ref{algebra-lemma-etale-over-field})
et que, par conséquent, $z$ annule un polynôme unitaire
à coefficients dans $k(x_1, \ldots, x_n)$. En chassant les dénominateurs,
on obtient $g$.

\medskip\noindent
L'existence de $g$ et le
lemme \ref{algebra-lemma-helper-polynomial} d'Algèbre
fournissent des entiers $e_1, e_2, \ldots, e_n \geq 1$ tels que
$z$ soit entier sur le sous-anneau $C'$ de $C$ engendré par
$t_1 = \varphi(x_1) + z^{pe_1}, \ldots, t_n = \varphi(x_n) + z^{pe_n}$.
Bien entendu, les éléments $\varphi(x_1), \ldots, \varphi(x_n)$
sont eux aussi entiers sur $C'$, de même que les éléments
$y_1, \ldots, y_s$. Enfin, par notre choix de $z$, l'élément
$y_{s + 1}$ est lui aussi entier sur $C'$.

\medskip\noindent
Considérons l'homomorphisme d'anneaux
$$
\varphi' : k[x_1, \ldots, x_n] \longrightarrow C, \quad
x_i \longmapsto t_i
$$
dont l'image est $C'$. Comme
$\text{d}(\varphi(x_i)) = \text{d}(t_i) = \text{d}(\varphi'(x_i))$
dans $\Omega_{C/k}$ (et c'est ici que nous utilisons l'hypothèse que la caractéristique de $k$
est $p > 0$), on conclut que $\varphi'$ est \'etale puisque $\varphi$ est
\'etale, voir
Algèbre, lemme \ref{algebra-lemma-characterize-etale-over-polynomial-ring}.
Observons que $\varphi'(x_i) - \varphi(x_i) = t_i - \varphi(x_i) = z^{pe_i}$
appartient au noyau $J$ de l'application $C \to C/\varphi(I)C \to B$ par notre
choix de $z$ comme élément de $J$.
Par conséquent, pour $f \in I$, l'élément
$$
\varphi'(f) =
f(t_1, \ldots, t_n) =
f(\varphi(x_1) + z^{pe_1}, \ldots, \varphi(x_n) + z^{pe_n}) =
\varphi(f) + \text{élément de }(z)
$$
appartient également à $J$. Autrement dit, $\varphi'(I)C \subset J$ et
nous obtenons une surjection
$$
\tau' : C/\varphi'(I)C \longrightarrow C/J \cong B
$$
d'algèbres \'etales sur $A$. Enfin, l'algèbre
$C$ est engendrée par les éléments
$\varphi(x_1), \ldots, \varphi(x_n), y_1, \ldots, y_r$
sur $C' = \Im(\varphi')$, avec
$\varphi(x_1), \ldots, \varphi(x_n), y_1, \ldots, y_{s + 1}$
entiers sur $C' = \Im(\varphi')$. Ainsi
$NF(C, \varphi', \tau') < r - s = NF(C, \varphi, \tau)$.
Ceci achève la démonstration.
\end{proof}

\begin{lemma}
\label{pione-lemma-affine-etale-over-affine-space}
Soient $k$ un corps de caractéristique $p > 0$ et $X \to \mathbf{A}^n_k$ un
morphisme \'etale, où $X$ est affine. Alors il existe un morphisme fini \'etale
$X \to \mathbf{A}^n_k$.
\end{lemma}

\begin{proof}
Écrivons $X = \Spec(C)$. Posons $A = 0$ et notons $I = k[x_1, \ldots, x_n]$.
Par hypothèse, il existe un homomorphisme de $k$-algèbres \'etale
$\varphi : k[x_1, \ldots, x_n] \to C$. Notons $\tau : C/\varphi(I)C \to 0$
l'unique surjection. Nous pouvons choisir $\varphi$ et $\tau$ de telle sorte
que $NF(C, \varphi, \tau)$ soit minimal. Par le lemme \ref{pione-lemma-lower-invariant},
nous obtenons $NF(C, \varphi, \tau) = 0$. Ainsi $\varphi$ est fini \'etale.
\end{proof}

\begin{lemma}
\label{pione-lemma-dominate-affine-space}
Soient $k$ un corps de caractéristique $p > 0$ et $Z \subset \mathbf{A}^n_k$
un sous-schéma fermé. Soit $Y \to Z$ un morphisme fini \'etale. Il existe un
morphisme fini \'etale $f : U \to \mathbf{A}^n_k$ tel qu'
il existe une immersion ouverte et fermée $Y \to f^{-1}(Z)$ au-dessus de $Z$.
\end{lemma}

\begin{proof}
Traduisons le problème en algèbre. Écrivons
$\mathbf{A}^n_k = \Spec(k[x_1, \ldots, x_n])$.
Alors $Z = \Spec(A)$, où $A = k[x_1, \ldots, x_n]/I$
pour un certain idéal $I \subset k[x_1, \ldots, x_n]$.
Écrivons $Y = \Spec(B)$, de sorte que $Y \to Z$ corresponde à
l'homomorphisme fini \'etale de $k$-algèbres $A \to B$.

\medskip\noindent
Par Algèbre, lemme \ref{algebra-lemma-lift-etale},
il existe un homomorphisme d'anneaux \'etale
$$
\varphi : k[x_1, \ldots, x_n] \to C
$$
et un homomorphisme surjectif de $A$-algèbres $\tau : C/\varphi(I)C \to B$.
(Nous pouvons même choisir $C, \varphi, \tau$ de sorte que $\tau$ soit un isomorphisme,
mais nous ne nous en servirons pas.) Nous pouvons choisir $\varphi$ et $\tau$ de telle sorte
que $NF(C, \varphi, \tau)$ soit minimal. Par le lemme \ref{pione-lemma-lower-invariant},
nous obtenons $NF(C, \varphi, \tau) = 0$. Ainsi $\varphi$ est fini \'etale.

\medskip\noindent
Soit $f : U = \Spec(C) \to \mathbf{A}^n_k$ le morphisme fini \'etale
correspondant à $\varphi$. Le morphisme
$Y \to f^{-1}(Z) = \Spec(C/\varphi(I)C)$ induit par $\tau$
est une immersion fermée, car $\tau$ est surjective, et une immersion ouverte, car
il s'agit d'un morphisme \'etale d'après Morphismes, lemme
\ref{morphisms-lemma-etale-permanence}. Ceci achève la démonstration.
\end{proof}

\noindent
Voici le résultat principal.

\begin{proposition}
\label{pione-proposition-injective}
Soient $p$ un nombre premier et $i : Z \to X$ une immersion fermée
de schémas affines connexes sur $\mathbf{F}_p$. Pour tout point géométrique
$\overline{z}$ de $Z$, l'application
$$
\pi_1(Z, \overline{z}) \to \pi_1(X, \overline{z})
$$
est injective.
\end{proposition}

\begin{proof}
Soit $Y \to Z$ un morphisme fini \'etale. Il suffit de construire
un morphisme fini \'etale $f : U \to X$ tel que $Y$ soit isomorphe à
un sous-schéma ouvert et fermé de $f^{-1}(Z)$, voir
le lemme \ref{pione-lemma-functoriality-galois-injective}.
Écrivons $Z = \Spec(A)$ et $X = \Spec(R)$, de sorte que l'immersion fermée
$Z \to X$ soit donnée par une surjection $R \to A$. Nous pouvons écrire
$A = \colim A_i$ comme la limite inductive filtrante de ses sous-$\mathbf{F}_p$-algèbres
de type fini. Par le lemme \ref{pione-lemma-limit},
nous pouvons trouver un $i$ et un morphisme fini \'etale $Y_i \to Z_i = \Spec(A_i)$
tel que $Y = Z \times_{Z_i} Y_i$.

\medskip\noindent
Choisissons une surjection $\mathbf{F}_p[x_1, \ldots, x_n] \to A_i$.
Elle détermine une immersion fermée
$$
Z_i = \Spec(A_i)
\longrightarrow
X_i = \mathbf{A}^n_{\mathbf{F}_p} = \Spec(\mathbf{F}_p[x_1, \ldots, x_n])
$$
Par la propriété universelle des algèbres de polynômes et puisque
$R \to A$ est surjective, nous pouvons trouver un diagramme commutatif
$$
\xymatrix{
\mathbf{F}_p[x_1, \ldots, x_n] \ar[r] \ar[d] &
A_i \ar[d] \\
R \ar[r] &
A
}
$$
de $\mathbf{F}_p$-algèbres. Nous avons donc un diagramme commutatif
$$
\xymatrix{
Y_i \ar[r] &
Z_i \ar[r] &
X_i \\
Y \ar[u] \ar[r] &
Z \ar[u] \ar[r] &
X \ar[u]
}
$$
dont le carré de gauche est cartésien. Il est clair que si nous pouvons
trouver $f_i : U_i \to X_i$ fini \'etale tel que $Y_i$ soit
isomorphe à un sous-schéma ouvert et fermé de $f_i^{-1}(Z_i)$,
alors le changement de base $f : U \to X$ de $f_i$ par $X \to X_i$
fournit une solution à notre problème. Nous concluons donc en appliquant
le lemme \ref{pione-lemma-dominate-affine-space} à
$Y_i \to Z_i \to X_i = \mathbf{A}^n_{\mathbf{F}_p}$.
\end{proof}















% OK, start here.
%

\chapter{Cohomologie \'etale}



\phantomsection
\label{etale-cohomology-section-phantom}



\section{Introduction}
\label{etale-cohomology-section-introduction}

\noindent
Ce chapitre est le premier d'une s\'erie de chapitres sur la cohomologie
\'etale des sch\'emas. Nous y pr\'esentons les rudiments de la topologie
\'etale et de la cohomologie des faisceaux ab\'eliens pour cette topologie. Une
grande partie des sujets abord\'es peut sans inconv\'enient \^etre omise lors
d'une premi\`ere lecture; on consultera la section suivante pour d\'ecider des
passages \`a laisser de c\^ot\'e.

\medskip\noindent
La version initiale de ce chapitre a \`a l'origine \'et\'e constitu\'ee par les notes
de la premi\`ere partie d'un cours de cohomologie \'etale donn\'e par Johan de Jong
\`a l'Universit\'e Columbia \`a l'automne 2009. Les premiers auteurs des notes
\'etaient Thibaut Pugin, Zachary Maddock et Min Lee. La seconde partie du
cours se trouve dans le chapitre sur la formule des traces; voir
La formule des traces, section \ref{trace-section-introduction}.




\section{Quelles sections omettre lors d'une premi\`ere lecture?}
\label{etale-cohomology-section-skip}

\noindent
Nous voulons utiliser le contenu de ce chapitre pour d\'evelopper la th\'eorie
des espaces alg\'ebriques, des champs de Deligne--Mumford, des champs
alg\'ebriques, etc. Nous avons donc ajout\'e \`a l'expos\'e initial de la
cohomologie \'etale des sch\'emas certains d\'eveloppements assez techniques.
Le lecteur les reconna\^itra \`a la fr\'equence du mot ``topos'', aux discussions
relevant de la th\'eorie des ensembles ou encore aux d\'emonstrations portant
sur des propri\'et\'es tr\`es g\'en\'erales des morphismes de sch\'emas. Certaines de
ces discussions peuvent \^etre omises lors d'une premi\`ere lecture.

\medskip\noindent
En particulier, nous sugg\'erons au lecteur d'omettre les sections suivantes:
\begin{enumerate}
\item Comparaison des grands et petits topos,
section \ref{etale-cohomology-section-compare}.
\item Reconstruction des morphismes,
section \ref{etale-cohomology-section-morphisms}.
\item Image directe et image inverse,
section \ref{etale-cohomology-section-monomorphisms}.
\item Propri\'et\'e (A),
section \ref{etale-cohomology-section-A}.
\item Propri\'et\'e (B),
section \ref{etale-cohomology-section-B}.
\item Propri\'et\'e (C),
section \ref{etale-cohomology-section-C}.
\item Invariance topologique du petit site \'etale,
section \ref{etale-cohomology-section-topological-invariance}.
\item Morphismes entiers universellement injectifs,
section \ref{etale-cohomology-section-integral-universally-injective}.
\item Grands sites et image directe,
section \ref{etale-cohomology-section-big}.
\item Exactitude du grand foncteur image directe \`a support propre,
section \ref{etale-cohomology-section-exactness-lower-shriek}.
\end{enumerate}
Outre ces sections, quelques r\'esultats isol\'es peuvent \^etre omis; le lecteur
les reconna\^itra aux mots-cl\'es indiqu\'es ci-dessus.



%9.08.09
\section{Prologue}
\label{etale-cohomology-section-prologue}

\noindent
Ces le\c{c}ons portent sur une autre th\'eorie cohomologique. Il faut tout
d'abord observer que la topologie de Zariski n'est pas enti\`erement
satisfaisante. L'une des principales raisons pour lesquelles elle ne donne pas
les r\'esultats souhait\'es est que, si $X$ est une vari\'et\'e complexe et
$\mathcal{F}$ un faisceau constant, alors
$$
H^i(X, \mathcal{F}) = 0, \quad \text{ pour tout } i > 0.
$$
En voici la raison. Dans un sch\'ema irr\'eductible (en particulier dans une
vari\'et\'e), deux ouverts non vides quelconques se rencontrent, de sorte que les
applications de restriction d'un faisceau constant sont surjectives. On dit
que le faisceau est {\it flasque}. Dans ce cas, tous les groupes de cohomologie
de {\v C}ech sup\'erieurs sont nuls, de m\^eme que tous les groupes de
cohomologie de Zariski sup\'erieurs. Autrement dit, la topologie de Zariski ne
poss\`ede ``pas assez'' d'ouverts pour d\'etecter cette cohomologie sup\'erieure.

\medskip\noindent
En revanche, si $X$ est une vari\'et\'e complexe projective lisse, alors
$$
H_{Betti}^{2 \dim X}(X (\mathbf{C}), \Lambda) = \Lambda \quad \text{ pour }
\Lambda = \mathbf{Z}, \ \mathbf{Z}/n\mathbf{Z},
$$
o\`u $X(\mathbf{C})$ d\'esigne l'ensemble des points complexes de $X$. Il serait
souhaitable de retrouver cette propri\'et\'e en g\'eom\'etrie alg\'ebrique, tout
particuli\`erement en caract\'eristique positive.




\section{La topologie \'etale}
\label{etale-cohomology-section-etale-topology}

\noindent
Il est tr\`es difficile de simplement ``ajouter'' des ouverts suppl\'ementaires
pour raffiner la topologie de Zariski. Une mani\`ere efficace de d\'efinir une
topologie consiste \`a consid\'erer non seulement des ouverts, mais aussi certains
sch\'emas au-dessus de ceux-ci. Pour d\'efinir la topologie \'etale, on consid\`ere
tous les morphismes \'etales $\varphi : U \to X$. Si $X$ est une vari\'et\'e
projective lisse sur $\mathbf{C}$, cela signifie que
\begin{enumerate}
\item $U$ est une r\'eunion disjointe de vari\'et\'es lisses, et
\item $\varphi$ est localement un isomorphisme (analytiquement).
\end{enumerate}
Le mot ``analytiquement'' renvoie \`a la topologie usuelle (transcendante) sur
$\mathbf{C}$. La seconde condition signifie donc que la diff\'erentielle de
$\varphi$ est partout de rang maximal (et, en particulier, que toutes les
composantes de $U$ ont la m\^eme dimension que $X$).

\medskip\noindent
Un rev\^etement double --- d\'efini de mani\`ere informelle comme une application
finie de degr\'e $2$ entre vari\'et\'es --- par exemple
$$
\Spec(\mathbf{C}[t])
\longrightarrow
\Spec(\mathbf{C}[t]),
\quad t \longmapsto t^2
$$
n'est pas un morphisme \'etale s'il poss\`ede une fibre r\'eduite \`a un seul point.
Dans cet exemple, cela se produit pour $t = 0$. Pour qu'une application finie
entre vari\'et\'es sur $\mathbf{C}$ soit \'etale, toutes ses fibres doivent avoir le
m\^eme nombre de points. En retirant le point $t = 0$ de la source de
l'application dans l'exemple, on rend le morphisme \'etale. Mais on peut aussi
retirer d'autres points de la source et le morphisme demeure \'etale. Pour
consid\'erer la topologie \'etale, il faut examiner tous ces morphismes. \`A la
diff\'erence de la topologie de Zariski, ils ne sont pas n\'ecessairement de
simples ouverts de $X$, bien que leurs images le soient toujours.

\begin{definition}
\label{etale-cohomology-definition-etale-covering-initial}
Une famille de morphismes $\{ \varphi_i : U_i \to X\}_{i \in I}$ est appel\'ee
un {\it recouvrement \'etale} si chaque $\varphi_i$ est un morphisme \'etale et
si leurs images recouvrent $X$, c'est-\`a-dire si
$X = \bigcup_{i \in I} \varphi_i(U_i)$.
\end{definition}

\noindent
Ceci ``d\'efinit'' la topologie \'etale. Autrement dit, nous pouvons maintenant
pr\'eciser ce que sont les faisceaux. Un {\it faisceau \'etale} $\mathcal{F}$
d'ensembles (resp.\ de groupes ab\'eliens, d'espaces vectoriels, etc.) sur $X$
est la donn\'ee de:
\begin{enumerate}
\item pour chaque morphisme \'etale $\varphi : U \to X$, un ensemble
(resp.\ un groupe ab\'elien, un espace vectoriel, etc.) $\mathcal{F}(U)$;
\item pour chaque paire $U, \ U'$ de sch\'emas \'etales sur $X$, et chaque
morphisme $U \to U'$ au-dessus de $X$ (qui est automatiquement \'etale), une
application de restriction
$\rho^{U'}_U : \mathcal{F}(U') \to \mathcal{F}(U)$
\end{enumerate}
Ces donn\'ees doivent satisfaire les conditions $\rho^U_U = \text{id}$ pour le
morphisme identique $U \to U$ et
$\rho^{U'}_U \circ \rho^{U''}_{U'} = \rho^{U''}_U$ pour des morphismes
$U \to U' \to U''$ de sch\'emas \'etales sur $X$, ainsi que l'{\it axiome de
faisceau} suivant:
\begin{itemize}
\item[(*)] pour tout recouvrement \'etale $\{ \varphi_i : U_i \to U\}_{i \in
I}$, le diagramme
$$
\xymatrix{
\mathcal{F} (U) \ar[r] &
\Pi_{i \in I} \mathcal{F} (U_i) \ar@<1ex>[r] \ar@<-1ex>[r] &
\Pi_{i, j \in I} \mathcal{F} (U_i \times_U U_j)
}
$$
est un diagramme d'\'egalisateur dans la cat\'egorie des ensembles
(resp.\ des groupes ab\'eliens, des espaces vectoriels, etc.).
\end{itemize}

\begin{remark}
\label{etale-cohomology-remark-i-is-j}
Dans le dernier \'enonc\'e, il est essentiel de ne pas oublier le cas $i = j$,
qui constitue en g\'en\'eral une condition loin d'\^etre triviale (contrairement
\`a ce qui se passe pour la topologie de Zariski). De fait, des recouvrements
importants ne comportent souvent qu'un seul \`el\'ement.
\end{remark}

\noindent
Puisque l'identit\'e est un morphisme \'etale, nous pouvons calculer les sections
globales d'un faisceau \'etale, et la cohomologie sera simplement donn\'ee par les
foncteurs d\'eriv\'es \`a droite correspondants. Autrement dit, une fois la th\'eorie
davantage d\'evelopp\'ee et les \'enonc\'es rendus pr\'ecis, rien ne s'opposera \`a la
d\'efinition de la cohomologie.




\section{Prouesses de la topologie \'etale}
\label{etale-cohomology-section-feats}

\noindent
Pour tout entier naturel $n \in \mathbf{N} = \{1, 2, 3, 4, \ldots\}$, on a
$$
H_\etale^2 (\mathbf{P}^1_\mathbf{C}, \mathbf{Z}/n\mathbf{Z}) =
\mathbf{Z}/n\mathbf{Z}.
$$
Plus g\'en\'eralement, si $X$ est une vari\'et\'e complexe, ses nombres de Betti
\'etales \`a coefficients dans un corps fini co\"incident avec les nombres de
Betti usuels de $X(\mathbf{C})$, c'est-\`a-dire que
$$
\dim_{\mathbf{F}_q} H_\etale^{2i} (X, \mathbf{F}_q) =
\dim_{\mathbf{F}_q} H_{Betti}^{2i} (X(\mathbf{C}), \mathbf{F}_q).
$$
C'est tout \`a fait satisfaisant. Toutefois, ces \'egalit\'es ne valent que pour des
coefficients de torsion, et non en g\'en\'eral. \`A coefficients entiers, on a
$$
H_\etale^2 (\mathbf{P}^1_\mathbf{C}, \mathbf{Z}) = 0.
$$
En revanche,
$H_{Betti}^2(\mathbf{P}^1(\mathbf{C}), \mathbf{Z}) = \mathbf{Z}$, car l'espace
topologique $\mathbf{P}^1(\mathbf{C})$ est hom\'eomorphe \`a une sph\`ere de
dimension $2$. Il existe des moyens de retrouver des coefficients non de
torsion \`a partir de coefficients de torsion par un proc\'ed\'e de passage \`a la
limite, que nous aborderons bient\^ot.




\section{Un calcul}
\label{etale-cohomology-section-computation}

\noindent
Comment calculer la cohomologie de $\mathbf{P}^1_\mathbf{C}$ \`a coefficients
dans $\Lambda = \mathbf{Z}/n\mathbf{Z}$?
Nous utilisons la cohomologie de {\v C}ech. Un recouvrement de
$\mathbf{P}^1_\mathbf{C}$ est donn\'e par les deux ouverts standards $U_0, U_1$,
tous deux isomorphes \`a $\mathbf{A}^1_\mathbf{C}$, et dont l'intersection est
isomorphe \`a
$\mathbf{A}^1_\mathbf{C} \setminus \{0\} = \mathbf{G}_{m, \mathbf{C}}$.
Il se trouve que la suite de Mayer--Vietoris vaut en cohomologie \'etale. On en
d\'eduit une suite exacte
$$
H_\etale^{i-1}(U_0\cap U_1, \Lambda) \to
H_\etale^i(\mathbf{P}^1_\mathbf{C}, \Lambda) \to
H_\etale^i(U_0, \Lambda) \oplus
H_\etale^i(U_1, \Lambda) \to H_\etale^i(U_0\cap U_1,
\Lambda).
$$
Pour obtenir la r\'eponse attendue, il faudrait montrer que la somme directe
figurant dans le troisi\`eme terme est nulle. En fait, comme pour la topologie
usuelle, on a bien
$$
H_\etale^q (\mathbf{A}^1_\mathbf{C}, \Lambda) = 0
\quad \text{ pour } q \geq 1,
$$
et
$$
H_\etale^q (\mathbf{A}^1_\mathbf{C} \setminus \{0\}, \Lambda) = \left\{
\begin{matrix}
\Lambda & \text{ si }q = 1\text{, et} \\
0 & \text{ pour }q \geq 2.
\end{matrix}
\right.
$$
Ces r\'esultats sont d\'ej\`a assez difficiles (quelle en serait une preuve
\'el\'ementaire?). Expliquons comment effectuer ce calcul une fois que nous
disposerons du formalisme de la cohomologie \'etale.

\medskip\noindent
{\bf Cohomologie sup\'erieure.} Elle est trait\'ee par le fait g\'en\'eral suivant:
si $X$ est une courbe affine sur $\mathbf{C}$, alors
$$
H_\etale^q (X, \mathbf{Z}/n\mathbf{Z}) = 0 \quad \text{ pour } q \geq 2.
$$
On le d\'emontre en consid\'erant le point g\'en\'erique de la courbe et en faisant
un peu de cohomologie galoisienne. Il ne reste donc \`a examiner que la
cohomologie en degr\'e 1.

\medskip\noindent
{\bf Cohomologie en degr\'e 1.} Nous utilisons les identifications suivantes:
\begin{eqnarray*}
H_\etale^1 (X, \mathbf{Z}/n\mathbf{Z}) = \left\{
\begin{matrix}
\text{faisceaux d'ensembles }\mathcal{F}\text{ sur le site \'etale }X_\etale
\text{ munis d'une} \\
\text{action }\mathbf{Z}/n\mathbf{Z} \times \mathcal{F} \to \mathcal{F}
\text{ telle que }\mathcal{F}\text{ soit un }\mathbf{Z}/n\mathbf{Z}\text{-torseur.}
\end{matrix}
\right\}
\Big/ \cong
\\
 = \left\{
\begin{matrix}
\text{morphismes }Y \to X\text{ finis \'etales, munis} \\
\text{ d'une action libre de }\mathbf{Z}/n\mathbf{Z}\text{ tels que }
X = Y/(\mathbf{Z}/n\mathbf{Z}).
\end{matrix}
\right\}
\Big/ \cong.
\end{eqnarray*}
La premi\`ere identification est tr\`es g\'en\'erale (elle vaut pour toute th\'eorie
cohomologique sur un site) et n'a rien \`a voir avec la topologie \'etale. La
seconde identification est une cons\'equence de la th\'eorie de la descente. Le
dernier ensemble d\'ecrit une collection d'objets g\'eom\'etriques que nous pouvons
manipuler concr\`etement.

\medskip\noindent
La courbe $\mathbf{A}^1_\mathbf{C}$ n'a aucun rev\^etement \'etale fini non
trivial; par cons\'equent,
$H_\etale^1 (\mathbf{A}^1_\mathbf{C}, \mathbf{Z}/n\mathbf{Z}) = 0$.
On peut le voir soit topologiquement, soit au moyen de l'argument du paragraphe
suivant.

\medskip\noindent
D\'ecrivons les rev\^etements \'etales finis
$\varphi : Y \to \mathbf{A}^1_\mathbf{C} \setminus \{0\}$.
Il suffit de consid\'erer le cas o\`u $Y$ est connexe, ce que nous supposons. Nous
allons d\'eterminer les possibilit\'es pour $Y$ en appliquant la formule de
Riemann--Hurwitz (c'est bien s\^ur un peu superflu, et le lecteur peut passer la
section suivante s'il le souhaite). Supposons que ce morphisme soit de degr\'e
$n$ et consid\'erons une compactification projective
$$
\xymatrix{
{Y\ } \ar@{^{(}->}[r] \ar[d]^\varphi &
{\bar Y} \ar[d]^{\bar\varphi} \\
{\mathbf{A}^1_\mathbf{C} \setminus \{0\}} \ar@{^{(}->}[r] &
{\mathbf{P}^1_\mathbf{C}}
}
$$
Bien que $\varphi$ soit \'etale et non ramifi\'e, $\bar{\varphi}$ peut \^etre
ramifi\'e en 0 et en $\infty$. Supposons que les ant\'ec\'edents de 0 soient les
points $y_1, \ldots, y_r$, d'indices de ramification $e_1, \ldots e_r$, et que
les ant\'ec\'edents de $\infty$ soient les points $y_1', \ldots, y_s'$, d'indices
de ramification $d_1, \ldots d_s$. En particulier,
$\sum e_i = n = \sum d_j$. La formule de Riemann--Hurwitz donne
$$
2 g_Y - 2 = -2n + \sum (e_i - 1) + \sum (d_j - 1)
$$
et donc $g_Y = 0$, $r = s = 1$ et $e_1 = d_1 = n$.
Ainsi $Y \cong {\mathbf{A}^1_\mathbf{C} \setminus \{0\}}$, et l'on voit
facilement que $\varphi(z) = \lambda z^n$ pour un certain
$\lambda \in \mathbf{C}^*$. Apr\`es avoir reparam\'etr\'e $Y$, nous pouvons
supposer $\lambda = 1$. Notre rev\^etement s'obtient donc en extrayant une
racine $n$-i\`eme de la coordonn\'ee sur
$\mathbf{A}^1_{\mathbf{C}} \setminus \{0\}$.

\medskip\noindent
Rappelons qu'il nous faut classifier les rev\^etements de
${\mathbf{A}^1_\mathbf{C} \setminus \{0\}}$ munis d'une action libre de
$\mathbf{Z}/n\mathbf{Z}$. Dans notre cas, toute action de ce type correspond
\`a un automorphisme de $Y$ envoyant $z$ sur $\zeta_n z$, o\`u $\zeta_n$ est une
racine primitive $n$-i\`eme de l'unit\'e. Il existe $\phi(n)$ telles actions (ici,
$\phi(n)$ d\'esigne l'indicatrice d'Euler). Il existe donc exactement $\phi(n)$
rev\^etements \'etales finis connexes munis d'une action libre donn\'ee de
$\mathbf{Z}/n\mathbf{Z}$, chacun correspondant \`a une racine primitive
$n$-i\`eme de l'unit\'e. Nous laissons au lecteur le soin de v\'erifier que les
rev\^etements \'etales finis non connexes de degr\'e $n$ de
$\mathbf{A}^1_{\mathbf{C}} \setminus \{0\}$, munis d'une action libre donn\'ee
de $\mathbf{Z}/n\mathbf{Z}$, correspondent bijectivement aux racines
$n$-i\`emes non primitives de $1$. Autrement dit, ce calcul montre que
$$
H_\etale^1 (\mathbf{A}^1_\mathbf{C} \setminus \{0\},
\mathbf{Z}/n\mathbf{Z}) =
\Hom(\mu_n(\mathbf{C}), \mathbf{Z}/n\mathbf{Z}) \cong \mathbf{Z}/n\mathbf{Z}.
$$
La premi\`ere identification est canonique, la seconde ne l'est pas; voir la
remarque \ref{etale-cohomology-remark-normalize-H1-Gm}.
Puisque la preuve de la formule de Riemann--Hurwitz ne fait pas appel au calcul
de la cohomologie, ce qui pr\'ec\`ede constitue effectivement une preuve (\`a
condition de compl\'eter les d\'etails concernant l'annulation, etc.).




\section{Coefficients non de torsion}
\label{etale-cohomology-section-nontorsion}

\noindent
Pour \'etudier les coefficients non de torsion, on pose la d\'efinition suivante:
$$
H_\etale^i (X, \mathbf{Q}_\ell) :=
\left( \lim_n H_\etale^i(X, \mathbf{Z}/\ell^n\mathbf{Z}) \right)
\otimes_{\mathbf{Z}_\ell} \mathbf{Q}_\ell.
$$
Le symbole $\lim_n$ d\'esigne la {\it limite} du syst\`eme de groupes de
cohomologie $H_\etale^i(X, \mathbf{Z}/\ell^n\mathbf{Z})$ index\'e par $n$; voir
Cat\'egories, section \ref{categories-section-posets-limits}.
Nous devrons donc \'etudier des syst\`emes de faisceaux satisfaisant certaines
conditions de compatibilit\'e.




\section{Th\'eorie des faisceaux}
\label{etale-cohomology-section-sheaf-theory}
%9.10.09

\noindent
Nous abordons maintenant v\'eritablement les sites et les faisceaux.
Une quantit\'e impressionnante de r\'esultats abstraits utiles pourrait trouver place
dans les quelques sections qui suivent. Une partie de ces r\'esultats est d\'evelopp\'ee dans des
chapitres ant\'erieurs, tels que ceux sur les sites, les modules sur les sites et la cohomologie
sur les sites. Nous nous effor\c{c}ons de ne pas ajouter trop de mati\`ere ici, mais seulement
ce qu'il faut pour que la suite de ce chapitre soit compr\'ehensible.




\section{Pr\'efaisceaux}
\label{etale-cohomology-section-presheaves}

\noindent
Une r\'ef\'erence pour cette section est
Sites, section \ref{sites-section-presheaves}.

\begin{definition}
\label{etale-cohomology-definition-presheaf}
Soit $\mathcal{C}$ une cat\'egorie. Un {\it pr\'efaisceau d'ensembles} (respectivement, un
{\it pr\'efaisceau ab\'elien}) sur $\mathcal{C}$ est un foncteur $\mathcal{C}^{opp} \to
\textit{Ensembles}$ (resp.\ $\textit{Ab}$).
\end{definition}

\noindent
{\bf Terminologie.} Si $U \in \Ob(\mathcal{C})$, les \'el\'ements de
$\mathcal{F}(U)$ sont appel\'es les {\it sections} de $\mathcal{F}$ au-dessus de
$U$. Pour $\varphi : V \to U$ dans $\mathcal{C}$, l'application
$\mathcal{F}(\varphi) : \mathcal{F}(U) \to \mathcal{F}(V)$
est appel\'ee l'{\it application de restriction} et est souvent not\'ee $s \mapsto s|_V$
ou parfois $s \mapsto \varphi^*s$. La notation $s|_V$ est ambigu\"e
car l'application de restriction d\'epend de $\varphi$, mais il s'agit d'un abus
de notation usuel. Nous employons aussi la notation
$\Gamma(U, \mathcal{F}) = \mathcal{F}(U)$.

\medskip\noindent
Dire que $\mathcal{F}$ est un foncteur signifie que, si
$W \to V \to U$ sont des morphismes de $\mathcal{C}$ et si
$s \in \Gamma(U, \mathcal{F})$, alors
$(s|_V)|_W = s |_W$, avec l'abus de
notation que nous venons de voir. De plus, les applications de restriction correspondant aux
morphismes identit\'e $\text{id}_U : U \to U$ sont l'identit\'e.

\medskip\noindent
La cat\'egorie des pr\'efaisceaux d'ensembles (respectivement, des pr\'efaisceaux ab\'eliens) sur
$\mathcal{C}$ est not\'ee $\textit{PSh} (\mathcal{C})$ (resp. $\textit{PAb}
(\mathcal{C})$). C'est la cat\'egorie des foncteurs de $\mathcal{C}^{opp}$ vers
$\textit{Ensembles}$ (resp. $\textit{Ab}$), c'est-\`a-dire que les morphismes de
pr\'efaisceaux sont les transformations naturelles de foncteurs. Nous ne consid\'erons les
cat\'egories $\textit{PSh}(\mathcal{C})$ et $\textit{PAb}(\mathcal{C})$
que lorsque la cat\'egorie $\mathcal{C}$ est petite. (Notre convention est qu'une cat\'egorie
est petite sauf mention contraire et, si elle ne l'est pas, elle doit \^etre
r\'epertori\'ee dans Cat\'egories, remarque \ref{categories-remark-big-categories}.)

\begin{example}
\label{etale-cohomology-example-representable-presheaf}
\'Etant donn\'e un objet $X \in \Ob(\mathcal{C})$, consid\'erons la r\`egle
$$
\begin{matrix}
U & \longmapsto & h_X(U) = \Mor_\mathcal{C}(U, X) \\
V \xrightarrow{\varphi} U & \longmapsto &
(\psi \mapsto \psi \circ \varphi) : h_X(U) \to h_X(V).
\end{matrix}
$$
Elle d\'efinit un foncteur $h_X : \mathcal{C}^{opp} \to \textit{Ensembles}$
et donc un pr\'efaisceau. On l'appelle le
{\it pr\'efaisceau repr\'esentable associ\'e \`a $X$.}
Il n'est pas vrai que les pr\'efaisceaux repr\'esentables soient des faisceaux pour toute topologie sur
tout site.
\end{example}

\begin{lemma}[Yoneda]
\label{etale-cohomology-lemma-yoneda}
\begin{slogan}
Les morphismes entre objets sont en bijection avec les transformations naturelles
entre les foncteurs qu'ils repr\'esentent.
\end{slogan}
Soit $\mathcal{C}$ une cat\'egorie, et soient $X, Y \in
\Ob(\mathcal{C})$. Il existe une bijection naturelle
$$
\begin{matrix}
\Mor_\mathcal{C}(X, Y) &
\longrightarrow &
\Mor_{\textit{PSh}(\mathcal{C})} (h_X, h_Y) \\
\psi &
\longmapsto &
h_\psi = \psi \circ - : h_X \to h_Y.
\end{matrix}
$$
\end{lemma}

\begin{proof}
Voir
Cat\'egories, lemme \ref{categories-lemma-yoneda}.
\end{proof}




\section{Sites}
\label{etale-cohomology-section-sites}


\begin{definition}
\label{etale-cohomology-definition-family-morphisms-fixed-target}
Soit $\mathcal{C}$ une cat\'egorie. Une {\it famille de morphismes \`a cible fix\'ee}
$\mathcal{U} = \{\varphi_i : U_i \to U\}_{i\in I}$ est la donn\'ee de
\begin{enumerate}
\item un objet $U \in \mathcal{C}$,
\item un ensemble $I$ (\'eventuellement vide), et
\item pour tout $i\in I$, un morphisme $\varphi_i : U_i \to U$ de $\mathcal{C}$
de cible $U$.
\end{enumerate}
\end{definition}

\noindent
Il existe une notion de {\it morphisme de familles de morphismes \`a cible
fix\'ee}. Un cas particulier est la notion de {\it raffinement}.
On pourra consulter \`a ce sujet
Sites, section \ref{sites-section-refinements}.

\begin{definition}
\label{etale-cohomology-definition-site}
Un {\it site}\footnote{Ce que nous appelons un site est appel\'e une cat\'egorie munie d'une
pr\'etopologie dans \cite[Expos\'e II, D\'efinition 1.3]{SGA4}.
Dans \cite{ArtinTopologies}, on l'appelle une cat\'egorie munie d'une topologie de
Grothendieck.} est la donn\'ee d'une cat\'egorie $\mathcal{C}$ et d'un ensemble
$\text{Cov}(\mathcal{C})$ form\'e de familles de morphismes \`a cible fix\'ee
appel\'ees {\it recouvrements}, qui satisfont aux axiomes suivants :
\begin{enumerate}
\item (isomorphisme) si $\varphi : V \to U$ est un isomorphisme dans $\mathcal{C}$,
alors $\{\varphi : V \to U\}$ est un recouvrement,
\item (localit\'e) si $\{\varphi_i : U_i \to U\}_{i\in I}$ est un recouvrement et si
pour tout $i \in I$ on se donne un recouvrement
$\{\psi_{ij} : U_{ij} \to U_i \}_{j\in I_i}$, alors
$$
\{
\varphi_i \circ \psi_{ij} : U_{ij} \to U
\}_{(i, j)\in \prod_{i\in I} \{i\} \times I_i}
$$
est aussi un recouvrement, et
\item (changement de base) si $\{U_i \to U\}_{i\in I}$
est un recouvrement et si $V \to U$ est un morphisme dans $\mathcal{C}$, alors
\begin{enumerate}
\item pour tout $i \in I$, le produit fibr\'e
$U_i \times_U V$ existe dans $\mathcal{C}$, et
\item $\{U_i \times_U V \to V\}_{i\in I}$ est un recouvrement.
\end{enumerate}
\end{enumerate}
\end{definition}

\noindent
Pour nous, la cat\'egorie sous-jacente \`a un site est toujours ``petite'', c.-\`a-d. que sa
collection d'objets forme un ensemble et que la collection des recouvrements d'un
site est aussi un ensemble (comme dans la d\'efinition ci-dessus). Dans ce chapitre,
nous omettrons le plus souvent les arguments qui ram\`enent \`a un ensemble la collection
des objets et des recouvrements. Pour une discussion plus approfondie, voir
Sites, remarque \ref{sites-remark-no-big-sites}.

\begin{example}
\label{etale-cohomology-example-site-topological-space}
Si $X$ est un espace topologique, il poss\`ede un site associ\'e $X_{Zar}$
d\'efini comme suit : les objets de $X_{Zar}$ sont les ouverts de $X$,
les morphismes entre eux sont les applications d'inclusion, et les recouvrements sont
les recouvrements topologiques usuels (surjectifs). Observons que si
$U, V \subset W \subset X$ sont des ouverts, alors $U \times_W V = U \cap V$
existe : cette cat\'egorie admet les produits fibr\'es. Toutes les v\'erifications sont imm\'ediates et
tout fonctionne comme pr\'evu.
\end{example}




\section{Faisceaux}
\label{etale-cohomology-section-sheaves}

\begin{definition}
\label{etale-cohomology-definition-sheaf}
Un pr\'efaisceau $\mathcal{F}$ d'ensembles (resp. un pr\'efaisceau ab\'elien) sur un site
$\mathcal{C}$ est dit {\it s\'epar\'e} si, pour tout recouvrement
$\{\varphi_i : U_i \to U\}_{i\in I} \in \text{Cov} (\mathcal{C})$,
l'application
$$
\mathcal{F}(U) \longrightarrow \prod\nolimits_{i\in I} \mathcal{F}(U_i)
$$
est injective. Ici, l'application est $s \mapsto (s|_{U_i})_{i\in I}$.
Le pr\'efaisceau $\mathcal{F}$ est un {\it faisceau} si, pour tout recouvrement
$\{\varphi_i : U_i \to U\}_{i\in I} \in \text{Cov} (\mathcal{C})$, le
diagramme
\begin{equation}
\label{etale-cohomology-equation-sheaf-axiom}
\xymatrix{
\mathcal{F}(U) \ar[r] &
\prod_{i\in I} \mathcal{F}(U_i) \ar@<1ex>[r] \ar@<-1ex>[r] &
\prod_{i, j \in I} \mathcal{F}(U_i \times_U U_j),
}
\end{equation}
o\`u la premi\`ere application est $s \mapsto (s|_{U_i})_{i\in I}$ et les deux
applications de droite sont
$(s_i)_{i\in I} \mapsto (s_i |_{U_i \times_U U_j})$ et
$(s_i)_{i\in I} \mapsto (s_j |_{U_i \times_U U_j})$,
est un diagramme \'egalisateur dans la cat\'egorie des ensembles (resp.\ des groupes ab\'eliens).
\end{definition}

\begin{remark}
\label{etale-cohomology-remark-empty-covering}
Pour le recouvrement vide (o\`u $I = \emptyset$), ceci implique que
$\mathcal{F}(\emptyset)$ est un produit vide, qui est un objet final de la
cat\'egorie correspondante (un singleton, tant pour $\textit{Ensembles}$ que pour
$\textit{Ab}$).
\end{remark}

\begin{example}
\label{etale-cohomology-example-sheaf-site-space}
En explicitant cette d\'efinition pour le site $X_{Zar}$ associ\'e \`a un espace
topologique, voir l'exemple \ref{etale-cohomology-example-site-topological-space}, on retrouve la notion usuelle
de faisceau.
\end{example}

\begin{definition}
\label{etale-cohomology-definition-category-sheaves}
Nous notons $\Sh(\mathcal{C})$ (resp.\ $\textit{Ab}(\mathcal{C})$)
la sous-cat\'egorie pleine de $\textit{PSh}(\mathcal{C})$
(resp.\ $\textit{PAb}(\mathcal{C})$) dont les objets sont les faisceaux. C'est la
{\it cat\'egorie des faisceaux d'ensembles} (resp.\ celle des {\it faisceaux ab\'eliens}) sur
$\mathcal{C}$.
\end{definition}




\section{L'exemple des G-ensembles}
\label{etale-cohomology-section-G-sets}

\noindent
Soit $G$ un groupe et d\'efinissons un site $\mathcal{T}_G$ comme suit : la cat\'egorie
sous-jacente est la cat\'egorie des $G$-ensembles, c'est-\`a-dire que ses objets sont des ensembles munis
d'une action \`a gauche de $G$ et que les morphismes sont les applications \'equivariantes ; les
recouvrements de $\mathcal{T}_G$ sont les familles
$\{\varphi_i : U_i \to U\}_{i\in I}$ satisfaisant
$U = \bigcup_{i\in I} \varphi_i(U_i)$.

\medskip\noindent
Il existe un objet particulier du site $\mathcal{T}_G$, \`a savoir le $G$-ensemble $G$
muni de son action naturelle par translations \`a gauche. Nous le notons ${}_G G$.
Observons qu'il existe un isomorphisme naturel de groupes
$$
\begin{matrix}
\rho : & G^{opp} & \longrightarrow & \text{Aut}_{G\textit{-Ensembles}}({}_G G) \\
& g & \longmapsto & (h \mapsto hg).
\end{matrix}
$$
En particulier, pour tout pr\'efaisceau $\mathcal{F}$, l'ensemble $\mathcal{F}({}_G G)$
h\'erite d'une action de $G$ au moyen de $\rho$. (Notons que, par contravariance de
$\mathcal{F}$, l'ensemble $\mathcal{F}({}_G G)$ est encore un $G$-ensemble \`a gauche.) En fait,
le foncteur
$$
\begin{matrix}
\Sh(\mathcal{T}_G) & \longrightarrow & G\textit{-Ensembles} \\
\mathcal{F} & \longmapsto & \mathcal{F}({}_G G)
\end{matrix}
$$
est une \'equivalence de cat\'egories. Son quasi-inverse est le foncteur $X \mapsto
h_X$. Sans donner la d\'emonstration compl\`ete (que l'on trouvera dans
Sites, section \ref{sites-section-example-sheaf-G-sets}),
essayons d'expliquer pourquoi il en est ainsi.
\begin{enumerate}
\item
Si $S$ est un $G$-ensemble, nous pouvons le d\'ecomposer en orbites $S = \coprod_{i\in I}
O_i$. L'axiome des faisceaux appliqu\'e au recouvrement $\{O_i \to S\}_{i\in I}$ affirme que
$$
\xymatrix{
\mathcal{F}(S) \ar[r] &
\prod_{i\in I} \mathcal{F}(O_i) \ar@<1ex>[r] \ar@<-1ex>[r] &
\prod_{i, j \in I} \mathcal{F}(O_i \times_S O_j)
}
$$
est un diagramme \'egalisateur. Comme les produits fibr\'es dans $G\textit{-Ensembles}$
sont induits par les produits fibr\'es dans $\textit{Ensembles}$, et comme
$\mathcal{F}(\emptyset)$ est un $G$-ensemble r\'eduit \`a un point, nous obtenons
$$
\prod_{i, j \in I} \mathcal{F}(O_i \times_S O_j) = \prod_{i \in I}
\mathcal{F}(O_i)
$$
et les deux applications ci-dessus sont en fait identiques. Ainsi, l'axiome des faisceaux affirme simplement
que $\mathcal{F}(S) = \prod_{i\in I} \mathcal{F}(O_i)$.
\item
Si $S$ est le $G$-ensemble $S= G/H$ et si $\mathcal{F}$ est un faisceau sur $\mathcal{T}_G$,
nous affirmons que
$$
\mathcal{F}(G/H) = \mathcal{F}({}_G G)^H
$$
et en particulier que $\mathcal{F}(\{*\}) = \mathcal{F}({}_G G)^G$. Pour le voir,
appliquons l'axiome des faisceaux au recouvrement $\{ {}_G G \to G/H \}$ de $S$. Nous
avons
\begin{eqnarray*}
{}_G G \times_{G/H} {}_G G & \cong & G \times H \\
(g_1, g_2) & \longmapsto & (g_1, g_1^{-1} g_2)
\end{eqnarray*}
Le membre de gauche est une r\'eunion disjointe de copies de ${}_G G$ (comme $G$-ensemble). L'axiome des faisceaux
s'\'ecrit donc
$$
\xymatrix{
\mathcal{F} (G/H) \ar[r] &
\mathcal{F}({}_G G) \ar@<1ex>[r] \ar@<-1ex>[r] &
\prod_{h\in H} \mathcal{F}({}_G G)
}
$$
o\`u les deux applications de droite sont $s \mapsto (s)_{h \in H}$ et $s \mapsto
(hs)_{h \in H}$. Par cons\'equent, $\mathcal{F}(G/H) = \mathcal{F}({}_G G)^H$, comme
annonc\'e.
\end{enumerate}
Ceci ne d\'emontre pas tout \`a fait l'\'equivalence de cat\'egories annonc\'ee, mais montre au
moins qu'un faisceau $\mathcal{F}$ est enti\`erement d\'etermin\'e par ses sections au-dessus de
${}_G G$. On trouvera les d\'etails (ainsi que des remarques ensemblistes) dans
Sites, section \ref{sites-section-example-sheaf-G-sets}.




\section{Passage au faisceau associ\'e}
\label{etale-cohomology-section-sheafification}

\begin{definition}
\label{etale-cohomology-definition-0-cech}
Soit $\mathcal{F}$ un pr\'efaisceau sur le site $\mathcal{C}$ et soit
$\mathcal{U} = \{U_i \to U\} \in \text{Cov} (\mathcal{C})$.
Nous d\'efinissons le {\it groupe de cohomologie de {\v C}ech de degr\'e z\'ero} de
$\mathcal{F}$ relativement \`a $\mathcal{U}$ par
$$
\check H^0 (\mathcal{U}, \mathcal{F}) =
\left\{
(s_i)_{i\in I} \in \prod\nolimits_{i\in I }\mathcal{F}(U_i)
\text{ tels que }
s_i|_{U_i \times_U U_j} = s_j |_{U_i \times_U U_j}
\right\}.
$$
\end{definition}

\noindent
Il existe une application canonique
$\mathcal{F}(U) \to \check H^0 (\mathcal{U}, \mathcal{F})$,
$s \mapsto (s |_{U_i})_{i\in I}$.
Nous appelons {\it morphisme de recouvrements} d'un recouvrement
$\mathcal{V} = \{V_j \to V\}_{j \in J}$ vers $\mathcal{U}$ un triplet
$(\chi, \alpha, \chi_j)$, o\`u
$\chi : V \to U$ est un morphisme,
$\alpha : J \to I$ est une application d'ensembles, et, pour tout
$j \in J$, le morphisme $\chi_j$ s'ins\`ere dans un diagramme commutatif
$$
\xymatrix{
V_j \ar[rr]_{\chi_j} \ar[d] & & U_{\alpha(j)} \ar[d] \\
V \ar[rr]^\chi & & U.
}
$$
\`A partir des donn\'ees $\chi, \alpha, \{\chi_j\}_{j \in J}$, nous d\'efinissons
\begin{eqnarray*}
\check H^0(\mathcal{U}, \mathcal{F}) & \longrightarrow &
\check H^0(\mathcal{V}, \mathcal{F}) \\
(s_i)_{i\in I} & \longmapsto &
\left(\chi_j^*\left(s_{\alpha(j)}\right)\right)_{j\in J}.
\end{eqnarray*}
Nous affirmons alors que
\begin{enumerate}
\item l'application est bien d\'efinie, et
\item elle ne d\'epend que de $\chi$ et est ind\'ependante du choix de
$\alpha, \{\chi_j\}_{j \in J}$.
\end{enumerate}
Nous omettons la d\'emonstration du premier point.
Pour voir (2), consid\'erons un autre triplet $(\psi, \beta, \psi_j)$ tel que
$\chi = \psi$. On a alors le diagramme commutatif
$$
\xymatrix{
V_j \ar[rrr]_{(\chi_j, \psi_j)} \ar[dd] & & &
U_{\alpha(j)} \times_U U_{\beta(j)} \ar[dl] \ar[dr] \\
& & U_{\alpha(j)} \ar[dr] & &
U_{\beta(j)} \ar[dl] \\
V \ar[rrr]^{\chi = \psi} & & & U.
}
$$
\'Etant donn\'ee une section $s \in \mathcal{F}(\mathcal{U})$, son image dans
$\mathcal{F}(V_j)$ par l'application d\'efinie par
$(\chi, \alpha, \{\chi_j\}_{j \in J})$
est $\chi_j^*s_{\alpha(j)}$, et
son image par l'application d\'efinie par $(\psi, \beta, \{\psi_j\}_{j \in J})$
est $\psi_j^*s_{\beta(j)}$. Ces deux \'el\'ements
sont \'egaux car, par hypoth\`ese, $s \in \check H^0(\mathcal{U}, \mathcal{F})$ ;
ils sont donc tous deux \'egaux au tir\'e en arri\`ere de la valeur commune
$$
s_{\alpha(j)}|_{U_{\alpha(j)} \times_U U_{\beta(j)}} =
s_{\beta(j)}|_{U_{\alpha(j)} \times_U U_{\beta(j)}}
$$
par l'application $(\chi_j, \psi_j)$ du diagramme.

\begin{theorem}
\label{etale-cohomology-theorem-sheafification}
Soient $\mathcal{C}$ un site et $\mathcal{F}$ un pr\'efaisceau sur $\mathcal{C}$.
\begin{enumerate}
\item La r\`egle
$$
U \mapsto \mathcal{F}^+(U) :=
\colim_{\mathcal{U} \text{ recouvrement de }U}
\check H^0(\mathcal{U}, \mathcal{F})
$$
est un pr\'efaisceau. De plus, la limite inductive est filtrante.
\item Il existe un morphisme canonique de pr\'efaisceaux $\mathcal{F} \to \mathcal{F}^+$.
\item Si $\mathcal{F}$ est un pr\'efaisceau s\'epar\'e, alors $\mathcal{F}^+$ est un faisceau
et le morphisme de (2) est injectif.
\item $\mathcal{F}^+$ est un pr\'efaisceau s\'epar\'e.
\item $\mathcal{F}^\# = (\mathcal{F}^+)^+$ est un faisceau, et le morphisme canonique
induit un isomorphisme fonctoriel
$$
\Hom_{\textit{PSh}(\mathcal{C})}(\mathcal{F}, \mathcal{G}) =
\Hom_{\Sh(\mathcal{C})}(\mathcal{F}^\#, \mathcal{G})
$$
pour tout $\mathcal{G} \in \Sh(\mathcal{C})$.
\end{enumerate}
\end{theorem}

\begin{proof}
Voir Sites, th\'eor\`eme \ref{sites-theorem-plus}.
\end{proof}

\noindent
Autrement dit, le passage au faisceau associ\'e, donn\'e par le morphisme naturel
$\mathcal{F} \to \mathcal{F}^\#$, est adjoint \`a gauche au foncteur d'oubli
$\Sh(\mathcal{C}) \to \textit{PSh}(\mathcal{C})$.




\section{Cohomologie}
\label{etale-cohomology-section-cohomology}

\noindent
Le r\'esultat fondamental suivant permet de d\'efinir la cohomologie
des faisceaux ab\'eliens sur les sites.

\begin{theorem}
\label{etale-cohomology-theorem-enough-injectives}
La cat\'egorie des faisceaux ab\'eliens sur un site est une cat\'egorie ab\'elienne
qui poss\`ede assez d'injectifs.
\end{theorem}

\begin{proof}
Voir
Modules sur les sites, lemme \ref{sites-modules-lemma-abelian-abelian} et
Injectifs, th\'eor\`eme \ref{injectives-theorem-sheaves-injectives}.
\end{proof}

\noindent
On peut donc d\'efinir la cohomologie au moyen des foncteurs d\'eriv\'es \`a droite du foncteur
des sections : si $U \in \Ob(\mathcal{C})$ et
$\mathcal{F} \in \textit{Ab}(\mathcal{C})$,
$$
H^p(U, \mathcal{F}) :=
R^p\Gamma(U, \mathcal{F}) =
H^p(\Gamma(U, \mathcal{I}^\bullet))
$$
o\`u $\mathcal{F} \to \mathcal{I}^\bullet$ est une r\'esolution injective. Pour ce faire,
il faut v\'erifier que le foncteur $\Gamma(U, -)$ est exact \`a gauche. C'est
le cas, et c'est en partie pourquoi la cat\'egorie $\textit{Ab}(\mathcal{C})$ est ab\'elienne ;
voir
Modules sur les sites, lemme \ref{sites-modules-lemma-abelian-abelian}.
Pour une \'etude plus g\'en\'erale de la cohomologie sur les sites (y compris le
foncteur des sections globales et ses foncteurs d\'eriv\'es \`a droite), voir
Cohomologie sur les sites, section \ref{sites-cohomology-section-cohomology-sheaves}.



\section{La topologie fpqc}
\label{etale-cohomology-section-fpqc}
%9.15.09

\noindent
Avant d'aborder la cohomologie \'etale, nous \'etudions quelque peu la topologie fpqc, car
elle se comporte bien avec les faisceaux quasi-coh\'erents.

\begin{definition}
\label{etale-cohomology-definition-fpqc-covering}
Soit $T$ un sch\'ema. Un {\it recouvrement fpqc} de $T$ est une famille
$\{ \varphi_i : T_i \to T\}_{i \in I}$ telle que
\begin{enumerate}
\item chaque $\varphi_i$ est un morphisme plat et
$\bigcup_{i\in I} \varphi_i(T_i) = T$, et
\item pour tout ouvert affine $U \subset T$, il existe un ensemble fini
$K$, une application $\mathbf{i} : K \to I$ et des ouverts affines
$U_{\mathbf{i}(k)} \subset T_{\mathbf{i}(k)}$ tels que
$U = \bigcup_{k \in K} \varphi_{\mathbf{i}(k)}(U_{\mathbf{i}(k)})$.
\end{enumerate}
\end{definition}

\begin{remark}
\label{etale-cohomology-remark-fpqc}
La premi\`ere condition correspond \`a fp, sigle fran\c{c}ais de
{\it fid\`element plat}, qui signifie ``faithfully flat'' en anglais, et
la seconde \`a qc, {\it quasi-compact}.
La seconde partie de la premi\`ere condition est superflue lorsque la seconde est satisfaite.
\end{remark}

\begin{example}
\label{etale-cohomology-example-fpqc-coverings}
Voici des exemples de recouvrements fpqc.
\begin{enumerate}
\item Tout recouvrement ouvert de Zariski de $T$ est un recouvrement fpqc.
\item La famille $\{\Spec(B) \to \Spec(A)\}$ est un recouvrement fpqc
si et seulement si $A \to B$ est un morphisme d'anneaux fid\`element plat.
\item Si $f: X \to Y$ est plat, surjectif et quasi-compact, alors $\{ f: X\to
Y\}$ est un recouvrement fpqc.
\item Le morphisme
$\varphi :
\coprod_{x \in \mathbf{A}^1_k} \Spec(\mathcal{O}_{\mathbf{A}^1_k, x})
\to \mathbf{A}^1_k$,
o\`u $k$ est un corps, est plat et surjectif. Il n'est pas quasi-compact, et
en fait la famille $\{\varphi\}$ n'est pas un recouvrement fpqc.
\item \'Ecrivons
$\mathbf{A}^2_k = \Spec(k[x, y])$. Notons $i_x : D(x) \to \mathbf{A}^2_k$
et $i_y : D(y) \to \mathbf{A}^2_k$ les immersions ouvertes standard.
Alors les familles
$\{i_x, i_y, \Spec(k[[x, y]]) \to \mathbf{A}^2_k\}$
et
$\{i_x, i_y, \Spec(\mathcal{O}_{\mathbf{A}^2_k, 0}) \to \mathbf{A}^2_k\}$
sont des recouvrements fpqc.
\end{enumerate}
\end{example}

\begin{lemma}
\label{etale-cohomology-lemma-site-fpqc}
La collection des recouvrements fpqc de la cat\'egorie des sch\'emas
satisfait les axiomes d'un site.
\end{lemma}

\begin{proof}
Voir Topologies, lemme \ref{topologies-lemma-fpqc}.
\end{proof}

\noindent
Il semble que ce lemme permette de d\'efinir le site fpqc de la cat\'egorie
des sch\'emas. Cependant, un probl\`eme ensembliste surgit lorsqu'on
consid\`ere la topologie fpqc ; voir
Topologies, section \ref{topologies-section-fpqc}.
Il provient de ce que nous exigeons que les sites soient ``petits'', alors qu'aucune
petite cat\'egorie de sch\'emas ne peut contenir un syst\`eme cofinal de recouvrements fpqc d'un
sch\'ema non vide donn\'e. Bien que cela ne nous emp\^eche pas, \`a proprement parler,
de d\'efinir des sites fpqc ``partiels'',
il ne semble pas prudent de le faire. La solution consiste \`a autoriser
la notion de faisceau pour la topologie fpqc (voir ci-dessous), mais \`a interdire de
consid\'erer la cat\'egorie de tous les faisceaux fpqc.

\begin{definition}
\label{etale-cohomology-definition-sheaf-property-fpqc}
Soit $S$ un sch\'ema. La cat\'egorie des sch\'emas sur $S$ est not\'ee
$\Sch/S$. Consid\'erons un foncteur
$\mathcal{F} : (\Sch/S)^{opp} \to \textit{Ensembles}$, autrement dit
un pr\'efaisceau d'ensembles. Nous dirons que $\mathcal{F}$
{\it v\'erifie la propri\'et\'e de faisceau pour la topologie fpqc}
si, pour tout recouvrement fpqc $\{U_i \to U\}_{i \in I}$ de sch\'emas sur $S$,
le diagramme (\ref{etale-cohomology-equation-sheaf-axiom}) est un diagramme \'egalisateur.
\end{definition}

\noindent
Nous dirons de m\^eme que $\mathcal{F}$
{\it v\'erifie la propri\'et\'e de faisceau pour la topologie de Zariski} si, pour
tout recouvrement ouvert $U = \bigcup_{i \in I} U_i$, le diagramme
(\ref{etale-cohomology-equation-sheaf-axiom}) est un diagramme \'egalisateur. Voir
Sch\'emas, d\'efinition \ref{schemes-definition-representable-by-open-immersions}.
Il est clair que cela revient \`a dire que, pour tout sch\'ema $T$ sur $S$, la
restriction de $\mathcal{F}$ aux ouverts de $T$ est un faisceau (au sens usuel).

\begin{lemma}
\label{etale-cohomology-lemma-fpqc-sheaves}
Soit $\mathcal{F}$ un pr\'efaisceau sur $\Sch/S$. Alors
$\mathcal{F}$ v\'erifie la propri\'et\'e de faisceau pour la topologie fpqc
si et seulement si
\begin{enumerate}
\item $\mathcal{F}$ v\'erifie la propri\'et\'e de faisceau relativement \`a la
topologie de Zariski, et
\item pour tout morphisme fid\`element plat $\Spec(B) \to \Spec(A)$
de sch\'emas affines sur $S$, l'axiome de faisceau est v\'erifi\'e pour le recouvrement
$\{\Spec(B) \to \Spec(A)\}$. Plus pr\'ecis\'ement, cela signifie que
$$
\xymatrix{
\mathcal{F}(\Spec(A)) \ar[r] &
\mathcal{F}(\Spec(B)) \ar@<1ex>[r] \ar@<-1ex>[r] &
\mathcal{F}(\Spec(B \otimes_A B))
}
$$
est un diagramme \'egalisateur.
\end{enumerate}
\end{lemma}

\begin{proof}
Voir Topologies, lemme \ref{topologies-lemma-sheaf-property-fpqc}.
\end{proof}

\noindent
Une autre mani\`ere de consid\'erer un pr\'efaisceau $\mathcal{F}$ sur
$\Sch/S$ satisfaisant la condition de faisceau pour la
topologie fpqc consiste \`a se donner les donn\'ees suivantes :
\begin{enumerate}
\item pour tout $T/S$, un faisceau usuel (c'est-\`a-dire pour la topologie de Zariski) $\mathcal{F}_T$ sur
$T_{Zar}$,
\item pour tout morphisme $f : T' \to T$ sur $S$, un morphisme de restriction
$f^{-1}\mathcal{F}_T \to \mathcal{F}_{T'}$
\end{enumerate}
satisfaisant
\begin{enumerate}
\item[(a)] les morphismes de restriction sont fonctoriels,
\item[(b)] si $f : T' \to T$ est une immersion ouverte, alors le morphisme de restriction
$f^{-1}\mathcal{F}_T \to \mathcal{F}_{T'}$ est un isomorphisme, et
\item[(c)] pour tout morphisme fid\`element plat
$\Spec(B) \to \Spec(A)$ sur $S$, le diagramme
$$
\xymatrix{
\mathcal{F}_{\Spec(A)}(\Spec(A)) \ar[r] &
\mathcal{F}_{\Spec(B)}(\Spec(B)) \ar@<1ex>[r] \ar@<-1ex>[r] &
\mathcal{F}_{\Spec(B \otimes_A B)}(\Spec(B \otimes_A B))
}
$$
est un diagramme \'egalisateur.
\end{enumerate}
Les donn\'ees (1) et (2), avec les conditions (a), (b), d\'efinissent un pr\'efaisceau
sur $\Sch/S$ satisfaisant la condition de faisceau pour la topologie de Zariski.
D'apr\`es le lemme \ref{etale-cohomology-lemma-fpqc-sheaves}, la condition (c) suffit alors pour obtenir la
condition de faisceau pour la topologie fpqc.

\begin{example}
\label{etale-cohomology-example-quasi-coherent}
Consid\'erons le pr\'efaisceau
$$
\begin{matrix}
\mathcal{F} : & (\Sch/S)^{opp} & \longrightarrow & \textit{Ab} \\
& T/S & \longmapsto & \Gamma(T, \Omega_{T/S}).
\end{matrix}
$$
La compatibilit\'e des diff\'erentielles avec la localisation entra\^ine que
$\mathcal{F}$ est un faisceau sur le site de Zariski.
Il ne v\'erifie toutefois pas la condition de faisceau pour la topologie fpqc.
Consid\'erons en effet le cas
$S = \Spec(\mathbf{F}_p)$ et le morphisme
$$
\varphi :
V = \Spec(\mathbf{F}_p[v])
\to
U = \Spec(\mathbf{F}_p[u])
$$
d\'efini par l'envoi de $u$ sur $v^p$. La famille $\{\varphi\}$ est un recouvrement fpqc,
mais le morphisme de restriction
$\mathcal{F}(U) \to \mathcal{F}(V)$
envoie le g\'en\'erateur $\text{d}u$ sur $\text{d}(v^p) = 0$, donc
c'est le morphisme nul, et le diagramme
$$
\xymatrix{
\mathcal{F}(U) \ar[r]^{0} &
\mathcal{F}(V) \ar@<1ex>[r] \ar@<-1ex>[r] &
\mathcal{F}(V \times_U V)
}
$$
n'est pas un diagramme \'egalisateur. Nous verrons plus loin que $\mathcal{F}$ donne en fait
lieu \`a un faisceau sur les sites \'etale et lisse.
\end{example}

\begin{lemma}
\label{etale-cohomology-lemma-representable-sheaf-fpqc}
Tout pr\'efaisceau repr\'esentable sur $\Sch/S$ satisfait \`a la
condition de faisceau pour la topologie fpqc.
\end{lemma}

\begin{proof}
Voir
Descente, lemme \ref{descent-lemma-fpqc-universal-effective-epimorphisms}.
\end{proof}

\noindent
Nous y reviendrons plus loin, car la d\'emonstration de ce fait utilise
la descente des faisceaux quasi-coh\'erents, que nous \'etudierons dans la
section suivante. Une mani\`ere \'el\'egante d'exprimer le lemme consiste \`a dire que
{\it la topologie fpqc est moins fine que la topologie canonique}, ou
que la topologie fpqc est {\it sous-canonique}. Dans le cadre des sites,
ceci est \'etudi\'e dans
Sites, section \ref{sites-section-representable-sheaves}.

\begin{remark}
\label{etale-cohomology-remark-fpqc-finest}
La topologie fpqc est plus fine que les topologies de Zariski, \'etale,
lisse, syntomique et fppf. Par cons\'equent, tout pr\'efaisceau
satisfaisant \`a la condition de faisceau pour la topologie fpqc est un
faisceau sur les sites de Zariski, \'etale, lisse, syntomique et fppf.
En particulier, les pr\'efaisceaux repr\'esentables sont des faisceaux sur
le site \'etale d'un sch\'ema,
par exemple.
\end{remark}

\begin{example}
\label{etale-cohomology-example-additive-group-sheaf}
Soit $S$ un sch\'ema.
Consid\'erons le sch\'ema en groupes additif $\mathbf{G}_{a, S} = \mathbf{A}^1_S$
sur $S$, voir
Groupo\"ides, exemple \ref{groupoids-example-additive-group}.
Le pr\'efaisceau repr\'esentable associ\'e est donn\'e par
$$
h_{\mathbf{G}_{a, S}}(T) =
\Mor_S(T, \mathbf{G}_{a, S}) =
\Gamma(T, \mathcal{O}_T).
$$
D'apr\`es ce qui pr\'ec\`ede, nous savons qu'il s'agit d'un pr\'efaisceau
d'ensembles satisfaisant \`a la condition de faisceau pour la topologie fpqc. Il s'agit manifestement
aussi d'un pr\'efaisceau d'anneaux. Nous pouvons donc le consid\'erer comme le foncteur
$$
\begin{matrix}
\mathcal{O} : &
(\Sch/S)^{opp} &
\longrightarrow &
\textit{Rings} \\
&
T/S &
\longmapsto &
\Gamma(T, \mathcal{O}_T)
\end{matrix}
$$
qui satisfait \`a la condition de faisceau pour la topologie fpqc.
On dispose de m\^eme d'une notion de $\mathcal{O}$-module, et ainsi de
suite.
\end{example}




\section{Descente fid\`element plate}
\label{etale-cohomology-section-fpqc-descent}

\noindent
Dans cette section, nous \'etudions la descente fid\`element plate des modules
quasi-coh\'erents. Plus pr\'ecis\'ement, nous d\'emontrerons que les modules quasi-coh\'erents
satisfont \`a la descente effective relativement aux recouvrements fpqc.

\begin{definition}
\label{etale-cohomology-definition-descent-datum}
Soit $\mathcal{U} = \{ t_i : T_i \to T\}_{i \in I}$ une famille de
morphismes de sch\'emas \`a cible fix\'ee. Une {\it donn\'ee de descente} pour les
faisceaux quasi-coh\'erents relativement \`a $\mathcal{U}$ est une famille
$((\mathcal{F}_i)_{i \in I}, (\varphi_{ij})_{i, j \in I})$ telle que
\begin{enumerate}
\item $\mathcal{F}_i$ est un faisceau quasi-coh\'erent sur $T_i$, et
\item $\varphi_{ij} : \text{pr}_0^* \mathcal{F}_i \to
\text{pr}_1^* \mathcal{F}_j$ est un isomorphisme de modules
sur $T_i \times_T T_j$,
\end{enumerate}
et que la {\it condition de cocycle} soit satisfaite : les diagrammes
$$
\xymatrix{
\text{pr}_0^*\mathcal{F}_i \ar[dr]_{\text{pr}_{02}^*\varphi_{ik}}
\ar[rr]^{\text{pr}_{01}^*\varphi_{ij}} & &
\text{pr}_1^*\mathcal{F}_j \ar[dl]^{\text{pr}_{12}^*\varphi_{jk}} \\
& \text{pr}_2^*\mathcal{F}_k
}
$$
commutent sur $T_i \times_T T_j \times_T T_k$.
Cette donn\'ee de descente est dite {\it effective} s'il existe un faisceau
quasi-coh\'erent $\mathcal{F}$ sur $T$ et des isomorphismes de $\mathcal{O}_{T_i}$-modules
$\varphi_i : t_i^* \mathcal{F} \cong \mathcal{F}_i$ compatibles avec
les applications $\varphi_{ij}$, \`a savoir
$$
\varphi_{ij} = \text{pr}_1^* (\varphi_j) \circ \text{pr}_0^* (\varphi_i)^{-1}.
$$
\end{definition}

\noindent
Dans cette section et la suivante, nous \'etudions quelques ingr\'edients de la
d\'emonstration du th\'eor\`eme ci-dessous, ainsi que des r\'esultats connexes.

\begin{theorem}
\label{etale-cohomology-theorem-descent-quasi-coherent}
Si $\mathcal{V} = \{T_i \to T\}_{i\in I}$ est un recouvrement fpqc, alors toutes
les donn\'ees de descente pour les faisceaux quasi-coh\'erents relativement \`a $\mathcal{V}$
sont effectives.
\end{theorem}

\begin{proof}
Voir
Descente, proposition \ref{descent-proposition-fpqc-descent-quasi-coherent}.
\end{proof}

\noindent
Autrement dit, la cat\'egorie fibr\'ee des faisceaux quasi-coh\'erents est un champ
sur le site fpqc.
La d\'emonstration du th\'eor\`eme comporte deux \'etapes. La premi\`ere consiste \`a observer
que, pour les recouvrements de Zariski, le r\'esultat est facile (ou bien connu), par
le recollement usuel des faisceaux (voir
Faisceaux, section \ref{sheaves-section-glueing-sheaves})
et le caract\`ere local de la quasi-coh\'erence. La seconde \'etape traite le cas d'un
recouvrement fpqc de la forme $\{\Spec(B) \to \Spec(A)\}$,
o\`u $A \to B$ est un homomorphisme d'anneaux fid\`element plat.
Il s'agit d'un lemme d'alg\`ebre, que nous \'enon\c{c}ons maintenant.

\medskip\noindent
{\bf Descente des modules.}
Si $A \to B$ est un homomorphisme d'anneaux, consid\'erons le complexe
$$
(B/A)_\bullet : B \to B \otimes_A B \to B \otimes_A B \otimes_A B \to \ldots
$$
o\`u $B$ est en degr\'e 0, $B \otimes_A B$ en degr\'e 1, etc., les applications \'etant
donn\'ees par
\begin{eqnarray*}
b & \mapsto & 1 \otimes b - b \otimes 1, \\
b_0 \otimes b_1 & \mapsto & 1 \otimes b_0 \otimes b_1 - b_0 \otimes 1 \otimes
b_1 + b_0 \otimes b_1 \otimes 1, \\
& \text{etc.}
\end{eqnarray*}

\begin{lemma}
\label{etale-cohomology-lemma-algebra-descent}
Si $A \to B$ est fid\`element plat, alors le complexe $(B/A)_\bullet$ est exact en
degr\'es strictement positifs, et $H^0((B/A)_\bullet) = A$.
\end{lemma}

\begin{proof}
Voir Descente, lemme \ref{descent-lemma-ff-exact}.
\end{proof}

\noindent
Grothendieck d\'emontre ceci en trois \'etapes. Premi\`erement, il suppose que l'application $A
\to B$ admet une section, et construit une homotopie explicite avec le complexe dont
$A$ est l'unique terme non nul, en degr\'e 0. Deuxi\`emement, il observe que, pour d\'emontrer
le r\'esultat, il suffit de le faire apr\`es un changement de base fid\`element plat $A \to
A'$, en rempla\c{c}ant $B$ par $B' = B \otimes_A A'$. Troisi\`emement, il effectue le
changement de base fid\`element plat $A \to A' = B$ et remarque que l'application
$A' = B \to B' = B \otimes_A B$ poss\`ede une section naturelle.

\medskip\noindent
La m\^eme strat\'egie d\'emontre le lemme suivant.

\begin{lemma}
\label{etale-cohomology-lemma-descent-modules}
Si $A \to B$ est fid\`element plat et si $M$ est un $A$-module, alors le
complexe $(B/A)_\bullet \otimes_A M$ est exact en degr\'es strictement positifs, et
$H^0((B/A)_\bullet \otimes_A M) = M$.
\end{lemma}

\begin{proof}
Voir Descente, lemme \ref{descent-lemma-ff-exact}.
\end{proof}

\begin{definition}
\label{etale-cohomology-definition-descent-datum-modules}
Soient $A \to B$ un homomorphisme d'anneaux et $N$ un $B$-module. Une {\it donn\'ee de descente} pour
$N$ relativement \`a $A \to B$ est un isomorphisme
$\varphi : N \otimes_A B \cong B \otimes_A N$ de $B \otimes_A B$-modules tel
que le diagramme de $B \otimes_A B \otimes_A B$-modules
$$
\xymatrix{
{N \otimes_A B \otimes_A B} \ar[dr]_{\varphi_{02}} \ar[rr]^{\varphi_{01}} & &
{B \otimes_A N \otimes_A B} \ar[dl]^{\varphi_{12}} \\
& {B \otimes_A B \otimes_A N}
}
$$
commute, o\`u $\varphi_{01} = \varphi \otimes \text{id}_B$, et de m\^eme
pour $\varphi_{12}$ et $\varphi_{02}$.
\end{definition}

\noindent
Si $N' = B \otimes_A M$ pour un certain $A$-module M, alors il est muni d'une donn\'ee de descente
canonique donn\'ee par l'application
$$
\begin{matrix}
\varphi_\text{can}: & N' \otimes_A B & \to & B \otimes_A N' \\
& b_0 \otimes m \otimes b_1 & \mapsto & b_0 \otimes b_1 \otimes m.
\end{matrix}
$$

\begin{definition}
\label{etale-cohomology-definition-effective-modules}
Une donn\'ee de descente $(N, \varphi)$ est dite {\it effective} s'il existe un
$A$-module $M$ tel que $(N, \varphi) \cong (B \otimes_A M,
\varphi_\text{can})$, avec la notion \'evidente d'isomorphisme de donn\'ees de descente.
\end{definition}

\noindent
Le th\'eor\`eme \ref{etale-cohomology-theorem-descent-quasi-coherent} est une cons\'equence du
r\'esultat suivant.

\begin{theorem}
\label{etale-cohomology-theorem-descent-modules}
Si $A \to B$ est fid\`element plat, alors les donn\'ees de descente relativement \`a $A\to B$
sont effectives.
\end{theorem}

\begin{proof}
Voir
Descente, proposition \ref{descent-proposition-descent-module}.
Voir aussi
Descente, remarque \ref{descent-remark-homotopy-equivalent-cosimplicial-algebras},
pour un autre point de vue sur la d\'emonstration.
\end{proof}

\begin{remarks}
\label{etale-cohomology-remarks-theorem-modules-exactness}
Les r\'esultats sur la descente des modules ont plusieurs applications :
\begin{enumerate}
\item L'exactitude du complexe de {\v C}ech en degr\'es strictement positifs pour
le recouvrement $\{\Spec(B) \to \Spec(A)\}$, o\`u $A \to B$ est
fid\`element plat. Ceci donnera des r\'esultats d'annulation en cohomologie.
\item Si $(N, \varphi)$ est une donn\'ee de descente relativement \`a un homomorphisme
fid\`element plat $A \to B$, alors l'$A$-module correspondant est donn\'e par
$$
M = \Ker \left(
\begin{matrix}
N & \longrightarrow & B \otimes_A N \\
n & \longmapsto & 1 \otimes n - \varphi(n \otimes 1)
\end{matrix}
\right).
$$
Voir
Descente, proposition \ref{descent-proposition-descent-module}.
\end{enumerate}
\end{remarks}




%9.17.09
\section{Faisceaux quasi-coh\'erents}
\label{etale-cohomology-section-quasi-coherent}

\noindent
Nous pouvons appliquer la descente des modules \`a l'\'etude des faisceaux quasi-coh\'erents.

\begin{proposition}
\label{etale-cohomology-proposition-quasi-coherent-sheaf-fpqc}
Pour tout faisceau quasi-coh\'erent $\mathcal{F}$ sur $S$, le pr\'efaisceau
$$
\begin{matrix}
\mathcal{F}^a : & \Sch/S & \to & \textit{Ab}\\
& (f: T \to S) & \mapsto & \Gamma(T, f^*\mathcal{F})
\end{matrix}
$$
est un $\mathcal{O}$-module qui satisfait \`a la condition de faisceau pour la
topologie fpqc.
\end{proposition}

\begin{proof}
Ceci est d\'emontr\'e dans
Descente, lemme \ref{descent-lemma-sheaf-condition-holds}.
Indiquons ici la d\'emonstration. Comme \'etabli dans le
lemme \ref{etale-cohomology-lemma-fpqc-sheaves},
il suffit de v\'erifier la propri\'et\'e de faisceau
sur les recouvrements de Zariski et les morphismes fid\`element plats de sch\'emas affines. La
propri\'et\'e de faisceau pour les recouvrements de Zariski rel\`eve de la th\'eorie usuelle des sch\'emas, puisque
$\Gamma(U, i^\ast \mathcal{F}) = \mathcal{F}(U)$ lorsque
$i : U \hookrightarrow S$ est une immersion ouverte.

\medskip\noindent
Pour $\left\{\Spec(B)\to \Spec(A)\right\}$, o\`u $A\to B$ est fid\`element
plat et o\`u
$\mathcal{F}|_{\Spec(A)} = \widetilde{M}$,
cela correspond au fait que
$M = H^0\left((B/A)_\bullet \otimes_A M \right)$, c'est-\`a-dire que
\begin{align*}
0 \to M \to B \otimes_A M \to B \otimes_A B \otimes_A M
\end{align*}
est exacte d'apr\`es le
lemme \ref{etale-cohomology-lemma-descent-modules}.
\end{proof}

\noindent
Il existe une notion abstraite de faisceau quasi-coh\'erent sur un site annel\'e.
Nous l'introduisons bri\`evement ici. Pour plus d'informations, consulter
Modules sur les sites, section \ref{sites-modules-section-local}.
Soient $\mathcal{C}$ une cat\'egorie et $U$ un objet de $\mathcal{C}$.
Alors $\mathcal{C}/U$ d\'esigne la cat\'egorie des objets au-dessus de $U$, voir
Cat\'egories, exemple \ref{categories-example-category-over-X}.
Si $\mathcal{C}$ est un site, alors $\mathcal{C}/U$ est \'egalement un site :
les recouvrements de $V/U$ sont les familles $\{V_i/U \to V/U\}$ de morphismes
de $\mathcal{C}/U$ \`a cible fix\'ee telles que
$\{V_i \to V\}$ soit un recouvrement de $\mathcal{C}$. De plus, pour tout
faisceau $\mathcal{F}$ sur $\mathcal{C}$, la {\it restriction}
$\mathcal{F}|_{\mathcal{C}/U}$ (d\'efinie de mani\`ere \'evidente)
est encore un faisceau. Voir
Sites, section \ref{sites-section-localize},
pour plus de d\'etails.

\begin{definition}
\label{etale-cohomology-definition-ringed-site}
Soit $\mathcal{C}$ un {\it site annel\'e}, c'est-\`a-dire un site muni d'un
faisceau d'anneaux $\mathcal{O}$. Un faisceau de $\mathcal{O}$-modules $\mathcal{F}$ sur
$\mathcal{C}$ est dit {\it quasi-coh\'erent} si, pour tout
$U \in \Ob(\mathcal{C})$, il existe un recouvrement
$\{U_i \to U\}_{i\in I}$ de $\mathcal{C}$ tel que la restriction
$\mathcal{F}|_{\mathcal{C}/U_i}$ soit isomorphe au conoyau d'une
application $\mathcal{O}$-lin\'eaire entre $\mathcal{O}$-modules libres
$$
\bigoplus\nolimits_{k \in K} \mathcal{O}|_{\mathcal{C}/U_i}
\longrightarrow
\bigoplus\nolimits_{l \in L} \mathcal{O}|_{\mathcal{C}/U_i}.
$$
La somme directe index\'ee par $K$ est le faisceau associ\'e au pr\'efaisceau
$V \mapsto \bigoplus_{k \in K} \mathcal{O}(V)$, et de m\^eme pour l'autre.
\end{definition}

\noindent
Bien qu'il soit utile de pouvoir donner une d\'efinition g\'en\'erale comme ci-dessus,
cette notion ne se comporte pas bien en g\'en\'eral.

\begin{remark}
\label{etale-cohomology-remark-final-object}
Dans le cas o\`u $\mathcal{C}$ poss\`ede un objet final, par exemple $S$, il
suffit de v\'erifier la condition de la d\'efinition pour
$U = S$ dans l'\'enonc\'e ci-dessus. Voir
Modules sur les sites, lemme \ref{sites-modules-lemma-local-final-object}.
\end{remark}

\begin{theorem}[M\'eta-th\'eor\`eme sur les faisceaux quasi-coh\'erents]
\label{etale-cohomology-theorem-quasi-coherent}
Soit $S$ un sch\'ema.
Soit $\mathcal{C}$ un site. Supposons que
\begin{enumerate}
\item la cat\'egorie sous-jacente $\mathcal{C}$ soit une
sous-cat\'egorie pleine de $\Sch/S$,
\item tout recouvrement de Zariski de $T \in \Ob(\mathcal{C})$
puisse \^etre raffin\'e par un recouvrement de $\mathcal{C}$,
\item $S/S$ soit un objet de $\mathcal{C}$,
\item tout recouvrement de $\mathcal{C}$ soit un recouvrement fpqc de sch\'emas.
\end{enumerate}
Alors le pr\'efaisceau $\mathcal{O}$ est un faisceau sur $\mathcal{C}$ et
tout $\mathcal{O}$-module quasi-coh\'erent sur $(\mathcal{C}, \mathcal{O})$
est de la forme $\mathcal{F}^a$ pour un certain faisceau quasi-coh\'erent
$\mathcal{F}$ sur $S$.
\end{theorem}

\begin{proof}
Apr\`es quelques arguments formels, c'est exactement le th\'eor\`eme
\ref{etale-cohomology-theorem-descent-quasi-coherent}. Nous omettons les d\'etails. Dans
Descente, Proposition \ref{descent-proposition-equivalence-quasi-coherent},
nous d\'emontrons une version plus pr\'ecise du th\'eor\`eme pour les
grands sites de Zariski, fppf, \'etale, lisse et syntomique de $S$,
ainsi que pour les petits sites de Zariski et \'etale de $S$.
\end{proof}

\noindent
Autrement dit, il n'y a aucune diff\'erence entre les modules quasi-coh\'erents
sur le sch\'ema $S$ et les $\mathcal{O}$-modules quasi-coh\'erents
sur les sites $\mathcal{C}$ comme dans le th\'eor\`eme. On trouvera des \'enonc\'es plus pr\'ecis
pour les grands et petits sites $(\Sch/S)_{fppf}$, $S_\etale$, etc.,
dans Descente, sections
\ref{descent-section-quasi-coherent-sheaves},
\ref{descent-section-quasi-coherent-cohomology}, et
\ref{descent-section-quasi-coherent-sheaves-bis}.
Dans ce chapitre, nous parlerons parfois d'un
``site comme dans le th\'eor\`eme \ref{etale-cohomology-theorem-quasi-coherent}''
afin d'\'enoncer commod\'ement des r\'esultats valables dans chacune de ces
situations.






\section{Cohomologie de {\v C}ech}
\label{etale-cohomology-section-cech-cohomology}

\noindent
Notre prochain objectif est d'utiliser la th\'eorie de la descente pour montrer que
$H^i(\mathcal{C}, \mathcal{F}^a) = H_{Zar}^i(S, \mathcal{F})$
pour tout faisceau quasi-coh\'erent $\mathcal{F}$ sur $S$, et
tout site $\mathcal{C}$ comme dans le th\'eor\`eme \ref{etale-cohomology-theorem-quasi-coherent}.
\`A cette fin, nous introduisons la cohomologie de {\v C}ech sur les sites.
Voir \cite{ArtinTopologies} et
Cohomologie sur les sites, sections \ref{sites-cohomology-section-cech},
\ref{sites-cohomology-section-cech-functor}
et \ref{sites-cohomology-section-cech-cohomology-cohomology}
pour plus de d\'etails.

\begin{definition}
\label{etale-cohomology-definition-cech-complex}
Soit $\mathcal{C}$ une cat\'egorie. Soit
$\mathcal{U} = \{U_i \to U\}_{i \in I}$ une famille de morphismes de
$\mathcal{C}$ \`a cible fix\'ee. Supposons que tous les produits fibr\'es
$U_{i_0} \times_U \ldots \times_U U_{i_p}$ existent dans $\mathcal{C}$.
Soit $\mathcal{F} \in \textit{PAb}(\mathcal{C})$ un
pr\'efaisceau ab\'elien. Nous d\'efinissons le {\it complexe de {\v C}ech}
$\check{\mathcal{C}}^\bullet(\mathcal{U}, \mathcal{F})$ par
$$
\prod_{i_0 \in I} \mathcal{F}(U_{i_0}) \to
\prod_{i_0, i_1 \in I} \mathcal{F}(U_{i_0} \times_U U_{i_1}) \to
\prod_{i_0, i_1, i_2 \in I}
\mathcal{F}(U_{i_0} \times_U U_{i_1} \times_U U_{i_2}) \to \ldots
$$
o\`u le premier terme est en degr\'e 0 et les applications sont usuelles.
Les {\it groupes de cohomologie de {\v C}ech} sont d\'efinis par
$$
\check{H}^p(\mathcal{U}, \mathcal{F}) =
H^p(\check{\mathcal{C}}^\bullet(\mathcal{U}, \mathcal{F})).
$$
\end{definition}

\noindent
Dans la d\'efinition ci-dessus, il est essentiel d'autoriser les indices
$(i_0, \ldots, i_p)$ \`a se r\'ep\'eter.

\begin{lemma}
\label{etale-cohomology-lemma-cech-presheaves}
Notations et hypoth\`eses comme dans la d\'efinition \ref{etale-cohomology-definition-cech-complex}.
Le foncteur $\check{\mathcal{C}}^\bullet(\mathcal{U}, -)$
est exact sur la cat\'egorie $\textit{PAb}(\mathcal{C})$.
\end{lemma}

\noindent
Autrement dit, si $0\to \mathcal{F}_1\to \mathcal{F}_2\to \mathcal{F}_3\to 0$
est une suite exacte courte de pr\'efaisceaux de groupes ab\'eliens, alors
$$
0 \to \check{\mathcal{C}}^\bullet\left(\mathcal{U}, \mathcal{F}_1\right)
\to\check{\mathcal{C}}^\bullet(\mathcal{U}, \mathcal{F}_2) \to
\check{\mathcal{C}}^\bullet(\mathcal{U}, \mathcal{F}_3)\to 0
$$
est une suite exacte courte de complexes.

\begin{proof}
Cela r\'esulte imm\'ediatement de la d\'efinition d'une suite exacte courte de
pr\'efaisceaux. En effet, puisque la cat\'egorie des pr\'efaisceaux ab\'eliens est la cat\'egorie des
foncteurs d'une certaine cat\'egorie \`a valeurs dans $\textit{Ab}$, c'est automatiquement une
cat\'egorie ab\'elienne : une suite $\mathcal{F}_1\to \mathcal{F}_2\to \mathcal{F}_3$
est exacte dans $\textit{PAb}$ si et seulement si, pour tout
$U \in \Ob(\mathcal{C})$, la suite
$\mathcal{F}_1(U) \to \mathcal{F}_2(U) \to \mathcal{F}_3(U)$ est exacte dans
$\textit{Ab}$. Ainsi, le complexe ci-dessus est simplement un produit de suites exactes courtes
en chaque degr\'e. Voir aussi
Cohomologie sur les sites, lemme \ref{sites-cohomology-lemma-cech-exact-presheaves}.
\end{proof}

\noindent
Cela montre que $\check{H}^\bullet(\mathcal{U}, -)$ est un $\delta$-foncteur.
Nous allons maintenant montrer qu'il s'agit d'un $\delta$-foncteur universel. Il nous faut donc
montrer que c'est un foncteur {\it effa\c{c}able}. Commen\c{c}ons par rappeler le
lemme de Yoneda.

\begin{lemma}[Lemme de Yoneda]
\label{etale-cohomology-lemma-yoneda-presheaf}
Pour tout pr\'efaisceau $\mathcal{F}$ sur une cat\'egorie $\mathcal{C}$
et tout $U \in \Ob(\mathcal{C})$, il existe un isomorphisme fonctoriel
$$
\Hom_{\textit{PSh}(\mathcal{C})}(h_U, \mathcal{F}) =
\mathcal{F}(U).
$$
\end{lemma}

\begin{proof}
Voir Catégories, lemme \ref{categories-lemma-yoneda}.
\end{proof}

\noindent
Pour un ensemble $E$, nous notons (dans cette section)
$\mathbf{Z}[E]$ le groupe ab\'elien libre engendr\'e par $E$. Sous forme de formule,
$\mathbf{Z}[E] = \bigoplus_{e \in E} \mathbf{Z}$, c'est-\`a-dire que $\mathbf{Z}[E]$ est
un $\mathbf{Z}$-module libre dont une base est form\'ee des \'el\'ements de $E$.
Avec cette notation, nous introduisons le pr\'efaisceau ab\'elien libre engendr\'e par un
pr\'efaisceau d'ensembles.

\begin{definition}
\label{etale-cohomology-definition-free-abelian-presheaf}
Soit $\mathcal{C}$ une cat\'egorie.
\'Etant donn\'e un pr\'efaisceau d'ensembles $\mathcal{G}$, nous d\'efinissons le
{\it pr\'efaisceau ab\'elien libre engendr\'e par $\mathcal{G}$},
not\'e $\mathbf{Z}_\mathcal{G}$, par la r\`egle
$$
\mathbf{Z}_\mathcal{G}(U)
=
\mathbf{Z}[\mathcal{G}(U)]
$$
pour $U \in \Ob(\mathcal{C})$,
les applications de restriction \'etant induites par celles de $\mathcal{G}$.
Dans le cas particulier o\`u $\mathcal{G} = h_U$, nous \'ecrivons simplement
$\mathbf{Z}_U = \mathbf{Z}_{h_U}$.
\end{definition}

\noindent
Le foncteur $\mathcal{G} \mapsto \mathbf{Z}_\mathcal{G}$ est adjoint \`a gauche du
foncteur d'oubli $\textit{PAb}(\mathcal{C}) \to \textit{PSh}(\mathcal{C})$.
Ainsi, pour tout pr\'efaisceau $\mathcal{F}$, il existe un isomorphisme canonique
$$
\Hom_{\textit{PAb}(\mathcal{C})}(\mathbf{Z}_U, \mathcal{F})
=
\Hom_{\textit{PSh}(\mathcal{C})}(h_U, \mathcal{F})
=
\mathcal{F}(U)
$$
o\`u la derni\`ere \'egalit\'e r\'esulte du lemme de Yoneda. En particulier, nous obtenons le
r\'esultat suivant.

\begin{lemma}
\label{etale-cohomology-lemma-cech-complex-describe}
Notations et hypoth\`eses comme dans la d\'efinition \ref{etale-cohomology-definition-cech-complex}.
Le complexe de {\v C}ech $\check{\mathcal{C}}^\bullet(\mathcal{U}, \mathcal{F})$
peut se d\'ecrire explicitement comme suit :
\begin{eqnarray*}
\check{\mathcal{C}}^\bullet(\mathcal{U}, \mathcal{F})
& = &
\left(
\prod_{i_0 \in I}
\Hom_{\textit{PAb}(\mathcal{C})}(\mathbf{Z}_{U_{i_0}}, \mathcal{F}) \to
\prod_{i_0, i_1 \in I}
\Hom_{\textit{PAb}(\mathcal{C})}(
\mathbf{Z}_{U_{i_0} \times_U U_{i_1}}, \mathcal{F}) \to \ldots
\right) \\
& = &
\Hom_{\textit{PAb}(\mathcal{C})}\left(
\left(
\bigoplus_{i_0 \in I} \mathbf{Z}_{U_{i_0}} \leftarrow
\bigoplus_{i_0, i_1 \in I} \mathbf{Z}_{U_{i_0} \times_U U_{i_1}} \leftarrow
\ldots
\right), \mathcal{F}\right)
\end{eqnarray*}
\end{lemma}

\begin{proof}
Cela r\'esulte de la formule ci-dessus. Voir
Cohomologie sur les sites, lemme \ref{sites-cohomology-lemma-cech-map-into}.
\end{proof}

\noindent
Nous sommes ainsi ramen\'es \`a l'\'etude du seul complexe situ\'e dans le premier argument du
dernier $\Hom$.

\begin{lemma}
\label{etale-cohomology-lemma-exact}
Notations et hypoth\`eses comme dans la d\'efinition \ref{etale-cohomology-definition-cech-complex}.
Le complexe de pr\'efaisceaux ab\'eliens
\begin{align*}
\mathbf{Z}_\mathcal{U}^\bullet \quad : \quad
\bigoplus_{i_0 \in I} \mathbf{Z}_{U_{i_0}} \leftarrow
\bigoplus_{i_0, i_1 \in I} \mathbf{Z}_{U_{i_0} \times_U U_{i_1}} \leftarrow
\bigoplus_{i_0, i_1, i_2 \in I}
\mathbf{Z}_{U_{i_0} \times_U U_{i_1} \times_U U_{i_2}} \leftarrow
\ldots
\end{align*}
est exact en tout degr\'e sauf en degr\'e $0$ dans $\textit{PAb}(\mathcal{C})$.
\end{lemma}

\begin{proof}
Pour tout $V\in \Ob(\mathcal{C})$, le complexe de groupes ab\'eliens
$\mathbf{Z}_\mathcal{U}^\bullet(V)$ est
$$
\begin{matrix}
\mathbf{Z}\left[
\coprod_{i_0\in I} \Mor_\mathcal{C}(V, U_{i_0})\right]
\leftarrow
\mathbf{Z}\left[
\coprod_{i_0, i_1 \in I}
\Mor_\mathcal{C}(V, U_{i_0} \times_U U_{i_1})\right]
\leftarrow \ldots = \\
\bigoplus_{\varphi : V \to U}
\left(
\mathbf{Z}\left[\coprod_{i_0 \in I} \Mor_\varphi(V, U_{i_0})\right]
\leftarrow
\mathbf{Z}\left[\coprod_{i_0, i_1\in I} \Mor_\varphi(V, U_{i_0}) \times
\Mor_\varphi(V, U_{i_1})\right]
\leftarrow
\ldots
\right)
\end{matrix}
$$
o\`u
$$
\Mor_{\varphi}(V, U_i)
=
\{ V \to U_i \text{ tels que } V \to U_i \to U \text{ est \'egal \`a } \varphi \}.
$$
Posons $S_\varphi = \coprod_{i\in I} \Mor_\varphi(V, U_i)$, de sorte que
$$
\mathbf{Z}_\mathcal{U}^\bullet(V)
=
\bigoplus_{\varphi : V \to U}
\left(
\mathbf{Z}[S_\varphi] \leftarrow
\mathbf{Z}[S_\varphi \times S_\varphi] \leftarrow
\mathbf{Z}[S_\varphi \times S_\varphi \times S_\varphi] \leftarrow
\ldots \right).
$$
Il suffit donc de montrer que, pour chaque $S = S_\varphi$, le complexe
\begin{align*}
\mathbf{Z}[S] \leftarrow
\mathbf{Z}[S \times S] \leftarrow
\mathbf{Z}[S \times S \times S] \leftarrow \ldots
\end{align*}
est exact en degr\'es n\'egatifs. Pour le voir, nous pouvons donner une homotopie explicite.
Fixons $s\in S$ et d\'efinissons $K: n_{(s_0, \ldots, s_p)} \mapsto n_{(s, s_0,
\ldots, s_p)}.$ On v\'erifie ais\'ement que $K$ est une homotopie \`a z\'ero pour l'op\'erateur
$$
\delta :
\eta_{(s_0, \ldots, s_p)}
\mapsto
\sum\nolimits_{i = 0}^p (-1)^p \eta_{(s_0, \ldots, \hat s_i, \ldots, s_p)}.
$$
Voir
Cohomologie sur les sites, lemme \ref{sites-cohomology-lemma-homology-complex},
pour plus de d\'etails.
\end{proof}

\begin{lemma}
\label{etale-cohomology-lemma-hom-injective}
Notations et hypoth\`eses comme dans la d\'efinition \ref{etale-cohomology-definition-cech-complex}.
Si $\mathcal{I}$ est un objet injectif de $\textit{PAb}(\mathcal{C})$,
alors $\check H^p(\mathcal{U}, \mathcal{I}) = 0$ pour tout $p > 0$.
\end{lemma}

\begin{proof}
Le complexe de {\v C}ech s'obtient en appliquant le foncteur
$\Hom_{\textit{PAb}(\mathcal{C})}(-, \mathcal{I}) $ au complexe $
\mathbf{Z}^\bullet_\mathcal{U} $, c'est-\`a-dire
$$
\check H^p(\mathcal{U}, \mathcal{I}) = H^p
(\Hom_{\textit{PAb}(\mathcal{C})} (\mathbf{Z}^\bullet_\mathcal{U},
\mathcal{I})).
$$
Mais nous venons de voir que $\mathbf{Z}^\bullet_\mathcal{U}$ est exact en
degr\'es n\'egatifs, et le foncteur $\Hom_{\textit{PAb}(\mathcal{C})}(-,
\mathcal{I})$ est exact ; par cons\'equent, $\Hom_{\textit{PAb}(\mathcal{C})}
(\mathbf{Z}^\bullet_\mathcal{U}, \mathcal{I})$ est exact en degr\'es positifs.
\end{proof}

\begin{theorem}
\label{etale-cohomology-theorem-cech-derived}
Notations et hypoth\`eses comme dans la d\'efinition \ref{etale-cohomology-definition-cech-complex}.
Sur $\textit{PAb}(\mathcal{C})$, les foncteurs $\check{H}^p(\mathcal{U}, -)$ sont
les foncteurs d\'eriv\'es \`a droite de $\check{H}^0(\mathcal{U}, -)$.
\end{theorem}

\begin{proof}
D'apr\`es le lemme \ref{etale-cohomology-lemma-hom-injective}, les foncteurs
$\check H^p(\mathcal{U}, -)$ sont des
$\delta$-foncteurs universels, puisqu'ils sont effa\c{c}ables.
Il en va de m\^eme des foncteurs d\'eriv\'es \`a droite de $\check H^0(\mathcal{U}, -)$. Puisqu'ils
co\"incident en degr\'e $0$, ils co\"incident par la propri\'et\'e universelle des
$\delta$-foncteurs universels. Pour plus de d\'etails, voir
Cohomologie sur les sites,
lemme \ref{sites-cohomology-lemma-cech-cohomology-derived-presheaves}.
\end{proof}

\begin{remark}
\label{etale-cohomology-remark-presheaves-no-topology}
Notons que tous les \'enonc\'es pr\'ec\'edents portent sur les pr\'efaisceaux ; nous n'avons donc
pas encore utilis\'e la topologie.
\end{remark}




\section{La suite spectrale de {\v C}ech \`a la cohomologie}
\label{etale-cohomology-section-cech-ss}

\noindent
Cette suite spectrale est essentielle pour \`etablir les r\'esultats fondamentaux sur la
cohomologie des faisceaux.

\begin{lemma}
\label{etale-cohomology-lemma-forget-injectives}
Le foncteur d'oubli $\textit{Ab}(\mathcal{C})\to \textit{PAb}(\mathcal{C})$
transforme les objets injectifs en objets injectifs.
\end{lemma}

\begin{proof}
C'est formel, puisque le foncteur d'oubli admet un adjoint \`a gauche,
\`a savoir le passage au faisceau associ\'e, qui est un foncteur exact. Pour plus de d\'etails, voir
Cohomologie sur les sites,
lemme \ref{sites-cohomology-lemma-injective-abelian-sheaf-injective-presheaf}.
\end{proof}

\begin{theorem}
\label{etale-cohomology-theorem-cech-ss}
Soit $\mathcal{C}$ un site. Pour tout recouvrement
$\mathcal{U} = \{U_i \to U\}_{i \in I}$ de $U \in \Ob(\mathcal{C})$
et tout faisceau ab\'elien $\mathcal{F}$ sur $\mathcal{C}$,
il existe une suite spectrale
$$
E_2^{p, q}
=
\check H^p(\mathcal{U}, \underline{H}^q(\mathcal{F}))
\Rightarrow
H^{p+q}(U, \mathcal{F}),
$$
o\`u $\underline{H}^q(\mathcal{F})$ est le pr\'efaisceau ab\'elien
$V \mapsto H^q(V, \mathcal{F})$.
\end{theorem}

\begin{proof}
Choisissons une r\'esolution injective $\mathcal{F}\to \mathcal{I}^\bullet$ dans
$\textit{Ab}(\mathcal{C})$ et consid\'erons le bicomplexe
$\check{\mathcal{C}}^\bullet(\mathcal{U}, \mathcal{I}^\bullet)$
et les applications
$$
\xymatrix{
\Gamma(U, I^\bullet) \ar[r] &
\check{\mathcal{C}}^\bullet(\mathcal{U}, \mathcal{I}^\bullet) \\
& \check{\mathcal{C}}^\bullet(\mathcal{U}, \mathcal{F}) \ar[u]
}
$$
Ici, l'application horizontale est l'application naturelle
$\Gamma(U, I^\bullet) \to
\check{\mathcal{C}}^0(\mathcal{U}, \mathcal{I}^\bullet)$
vers la colonne de gauche, et l'application verticale est induite par
$\mathcal{F}\to \mathcal{I}^0$ et aboutit \`a la ligne inf\'erieure.
Par hypoth\`ese, $\mathcal{I}^\bullet$ est un complexe d'injectifs dans
$\textit{Ab}(\mathcal{C})$ ; par le
lemme \ref{etale-cohomology-lemma-forget-injectives}, c'est donc un complexe d'injectifs dans
$\textit{PAb}(\mathcal{C})$. Ainsi, les lignes du bicomplexe sont
exactes en degr\'es positifs (lemme \ref{etale-cohomology-lemma-hom-injective}), et
le noyau de $\check{\mathcal{C}}^0(\mathcal{U}, \mathcal{I}^\bullet)
\to \check{\mathcal{C}}^1(\mathcal{U}, \mathcal{I}^\bullet)$
est \'egal \`a
$\Gamma(U, \mathcal{I}^\bullet)$, puisque $\mathcal{I}^\bullet$
est un complexe de faisceaux. En particulier, la cohomologie du complexe total
est la cohomologie usuelle
du foncteur des sections globales $H^0(U, \mathcal{F})$.

\medskip\noindent
Dans la direction verticale, le $q$-i\`eme groupe de cohomologie de la $p$-i\`eme colonne est
$$
\prod_{i_0, \ldots, i_p}
H^q(U_{i_0} \times_U \ldots \times_U U_{i_p}, \mathcal{F})
=
\prod_{i_0, \ldots, i_p}
\underline{H}^q(\mathcal{F})(U_{i_0} \times_U \ldots \times_U U_{i_p})
$$
dans le terme $E_1^{p, q}$. Il s'agit donc d'une suite spectrale standard de bicomplexe,
et sa page $E_2$ est celle indiqu\'ee. Pour plus de d\'etails, voir
Cohomologie sur les sites,
lemme \ref{sites-cohomology-lemma-cech-spectral-sequence}.
\end{proof}

\begin{remark}
\label{etale-cohomology-remark-grothendieck-ss}
C'est une suite spectrale de Grothendieck pour la compos\'ee des foncteurs
$$
\textit{Ab}(\mathcal{C}) \longrightarrow
\textit{PAb}(\mathcal{C}) \xrightarrow{\check H^0} \textit{Ab}.
$$
\end{remark}








\section{Grands et petits sites des sch\'emas}
\label{etale-cohomology-section-big-small}

\noindent
Soit $S$ un sch\'ema.
Soit $\tau$ l'une des topologies dont nous allons parler.
Ainsi, $\tau \in \{fppf, syntomic, smooth, \etale, Zariski\}$.
Bien entendu, si seule la topologie \'etale vous int\'eresse, vous pouvez
simplement supposer $\tau = \etale$ partout. En outre, nous traiterons
plus en d\'etail des morphismes \'etales, des recouvrements \'etales et des sites \'etales
\`a partir de la section \ref{etale-cohomology-section-etale-site}.
Pour poursuivre l'\'etude de la cohomologie des
faisceaux quasi-coh\'erents, il est commode d'introduire le
grand $\tau$-site et, lorsque $\tau \in \{\etale, Zariski\}$, le
petit $\tau$-site de $S$. Pour ce faire, nous introduisons d'abord
la notion de $\tau$-recouvrement.

\begin{definition}
\label{etale-cohomology-definition-tau-covering}
(Voir
Topologies, définitions
\ref{topologies-definition-fppf-covering},
\ref{topologies-definition-syntomic-covering},
\ref{topologies-definition-smooth-covering},
\ref{topologies-definition-etale-covering}, et
\ref{topologies-definition-zariski-covering}.)
Soit $\tau \in \{fppf, syntomic, smooth, \etale, Zariski\}$.
Une famille de morphismes de sch\'emas $\{f_i : T_i \to T\}_{i \in I}$ \`a cible
fix\'ee est appel\'ee un {\it $\tau$-recouvrement} si et seulement si
chaque $f_i$ est plat et de pr\'esentation finie, syntomique, lisse, \'etale,
respectivement une immersion ouverte, et si $\bigcup f_i(T_i) = T$.
\end{definition}

\noindent
La classe de tous les $\tau$-recouvrements satisfait aux axiomes
(1), (2) et (3) de la
d\'efinition \ref{etale-cohomology-definition-site} (notre d\'efinition d'un site), voir
Topologies, lemmes
\ref{topologies-lemma-fppf},
\ref{topologies-lemma-syntomic},
\ref{topologies-lemma-smooth},
\ref{topologies-lemma-etale}, et
\ref{topologies-lemma-zariski}.

\medskip\noindent
Introduisons les sites avec lesquels nous travaillerons.
Contrairement \`a ce qui se passe dans
\cite{SGA4}, nous ne voulons pas choisir un univers. Nous choisissons plut\^ot un ``univers
partiel'' (qui est un ensemble suffisamment grand comme dans
Ensembles, section \ref{sets-section-sets-hierarchy}) et consid\'erons tous les sch\'emas
contenus dans cet ensemble. Bien entendu, nous veillons \`a ce que notre sch\'ema de base privil\'egi\'e
$S$ appartienne \`a l'univers partiel. Apr\`es avoir choisi la cat\'egorie sous-jacente,
nous choisissons un ensemble suffisamment grand de $\tau$-recouvrements qui en fait un site.
Les d\'etails se trouvent dans le chapitre sur les topologies des sch\'emas ; les choix
laissent une grande libert\'e, mais, en fin de compte, les choix effectivement faits n'auront
aucune incidence sur la cohomologie \'etale (ou autre) de $S$ (tout comme, dans \cite{SGA4},
le choix effectif de l'univers n'a finalement aucune importance). En outre, la mani\`ere dont le
texte est r\'edig\'e permet au lecteur qui accepte les cardinaux fortement inaccessibles
(c'est-\`a-dire les univers) de les utiliser en remplacement.

\begin{definition}
\label{etale-cohomology-definition-tau-site}
Soit $S$ un sch\'ema.
Soit $\tau \in \{fppf, syntomic, smooth, \etale, \linebreak[0] Zariski\}$.
\begin{enumerate}
\item Un {\it grand $\tau$-site de $S$} est l'un quelconque des sites
$(\Sch/S)_\tau$ construits comme expliqu\'e ci-dessus et, plus en d\'etail, dans
Topologies, définitions
\ref{topologies-definition-big-small-fppf},
\ref{topologies-definition-big-small-syntomic},
\ref{topologies-definition-big-small-smooth},
\ref{topologies-definition-big-small-etale}, et
\ref{topologies-definition-big-small-Zariski}.
\item Si $\tau \in \{\etale, Zariski\}$, alors le
{\it petit $\tau$-site de $S$}
est la sous-cat\'egorie pleine $S_\tau$ de $(\Sch/S)_\tau$ dont les objets
sont les sch\'emas $T$ sur $S$ dont le morphisme structural $T \to S$ est \'etale,
respectivement une immersion ouverte. Un recouvrement dans $S_\tau$ est un recouvrement
$\{U_i \to U\}$ dans $(\Sch/S)_\tau$
tel que $U$ soit un objet de $S_\tau$.
\end{enumerate}
\end{definition}

\noindent
La cat\'egorie sous-jacente au site $(\Sch/S)_\tau$ poss\`ede des propri\'et\'es
raisonnables de ``stabilit\'e'' : \'etant donn\'e un sch\'ema $T$ qui lui appartient, tout sous-sch\'ema
localement ferm\'e de $T$ est isomorphe \`a un objet de $(\Sch/S)_\tau$.
Parmi les autres propri\'et\'es de ce type : stabilit\'e par produits fibr\'es de sch\'emas,
par sommes disjointes d\'enombrables,
par passage aux sch\'emas de type fini sur un sch\'ema donn\'e ; \'etant donn\'e un sch\'ema affine
$\Spec(R)$, on peut compl\'eter ou localiser $R$, ou en prendre le quotient
par un id\'eal, tout en restant dans la cat\'egorie, etc.
En revanche, par exemple, des sommes disjointes arbitraires
de sch\'emas de $(\Sch/S)_\tau$ feront sortir de cette cat\'egorie.
Notons aussi que, pour un objet $T$ de $(\Sch/S)_\tau$, il existe des
$\tau$-recouvrements $\{T_i \to T\}_{i \in I}$ (comme dans la
d\'efinition \ref{etale-cohomology-definition-tau-covering})
qui ne sont pas des recouvrements dans $(\Sch/S)_\tau$, par exemple
parce que les sch\'emas $T_i$ ne sont pas des objets de la cat\'egorie
$(\Sch/S)_\tau$. Mais notre choix des sites $(\Sch/S)_\tau$
est tel qu'il existe toujours
un recouvrement $\{U_j \to T\}_{j \in J}$ de $(\Sch/S)_\tau$ qui raffine
le recouvrement $\{T_i \to T\}_{i \in I}$, voir
Topologies, lemmes
\ref{topologies-lemma-fppf-induced},
\ref{topologies-lemma-syntomic-induced},
\ref{topologies-lemma-smooth-induced},
\ref{topologies-lemma-etale-induced}, et
\ref{topologies-lemma-zariski-induced}.
Dans ce chapitre, nous ignorerons pour l'essentiel ces questions.

\medskip\noindent
Si $\mathcal{F}$ est un faisceau sur $(\Sch/S)_\tau$ ou $S_\tau$, alors
nous notons
$$
H^p_\tau(U, \mathcal{F}), \text{ en particulier }
H^p_\tau(S, \mathcal{F})
$$
les groupes de cohomologie de $\mathcal{F}$ sur l'objet $U$ du site, voir
la section \ref{etale-cohomology-section-cohomology}. Nous avons ainsi
$H^p_{fppf}(S, \mathcal{F})$,
$H^p_{syntomic}(S, \mathcal{F})$,
$H^p_{smooth}(S, \mathcal{F})$,
$H^p_\etale(S, \mathcal{F})$, et
$H^p_{Zar}(S, \mathcal{F})$. Les deux derniers sont potentiellement ambigus, puisqu'ils
peuvent renvoyer soit au grand, soit au petit site \'etale ou de Zariski. Cependant,
cette ambigu\"it\'e est sans cons\'equence gr\^ace au lemme suivant.

\begin{lemma}
\label{etale-cohomology-lemma-compare-cohomology-big-small}
Soit $\tau \in \{\etale, Zariski\}$.
Si $\mathcal{F}$ est un faisceau ab\'elien d\'efini sur
$(\Sch/S)_\tau$, alors
les groupes de cohomologie de $\mathcal{F}$ sur $S$ co\"incident avec les groupes de cohomologie
de $\mathcal{F}|_{S_\tau}$ sur $S$.
\end{lemma}

\begin{proof}
D'apr\`es
Topologies, lemmes \ref{topologies-lemma-at-the-bottom} et
\ref{topologies-lemma-at-the-bottom-etale},
les foncteurs $S_\tau \to (\Sch/S)_\tau$
satisfont aux hypoth\`eses de
Sites, lemme \ref{sites-lemma-bigger-site}.
Notre lemme r\'esulte donc de
Cohomologie sur les sites, lemme \ref{sites-cohomology-lemma-cohomology-bigger-site}.
\end{proof}

\noindent
La cat\'egorie des faisceaux sur le grand ou le petit site \'etale de $S$ ne d\'epend que
de la sous-cat\'egorie pleine de $(\Sch/S)_\etale$ ou de $S_\etale$ form\'ee des
sch\'emas affines, et il suffit de consid\'erer entre eux les recouvrements \'etales standards
(d\'efinis ci-dessous). Cela donne les sites
$(\textit{Aff}/S)_\etale$ et $S_{affine, \etale}$, voir
Topologies, définition \ref{topologies-definition-big-small-etale}.
Les r\'esultats de comparaison sont d\'emontr\'es dans Topologies,
lemmes \ref{topologies-lemma-affine-big-site-etale} et
\ref{topologies-lemma-alternative}. Voici notre d\'efinition des
recouvrements standards pour certaines des topologies consid\'er\'ees dans ce chapitre.

\begin{definition}
\label{etale-cohomology-definition-standard-tau}
(Voir
Topologies, définitions
\ref{topologies-definition-standard-fppf},
\ref{topologies-definition-standard-syntomic},
\ref{topologies-definition-standard-smooth},
\ref{topologies-definition-standard-etale}, et
\ref{topologies-definition-standard-Zariski}.)
Soit $\tau \in \{fppf, syntomic, smooth, \etale, Zariski\}$.
Soit $T$ un sch\'ema affine.
Un {\it $\tau$-recouvrement standard} de $T$ est une famille
$\{f_j : U_j \to T\}_{j = 1, \ldots, m}$ o\`u chaque $U_j$ est affine,
et chaque $f_j$ est plat et de pr\'esentation finie,
syntomique standard, lisse standard, \'etale, respectivement l'immersion d'un
ouvert principal standard dans $T$, et $T = \bigcup f_j(U_j)$.
\end{definition}

\begin{lemma}
\label{etale-cohomology-lemma-tau-affine}
Soit $\tau \in \{fppf, syntomic, smooth, \etale, Zariski\}$.
Tout $\tau$-recouvrement d'un sch\'ema affine peut \^etre raffin\'e par un
$\tau$-recouvrement standard.
\end{lemma}

\begin{proof}
Voir
Topologies, lemmes
\ref{topologies-lemma-fppf-affine},
\ref{topologies-lemma-syntomic-affine},
\ref{topologies-lemma-smooth-affine},
\ref{topologies-lemma-etale-affine}, et
\ref{topologies-lemma-zariski-affine}.
\end{proof}

\noindent
Par souci d'exhaustivit\'e, nous \'enon\c{c}ons et d\'emontrons l'invariance, par rapport au choix de l'univers
partiel, des groupes de cohomologie que nous consid\'erons. Nous d\'emontrerons l'invariance
du petit topos \'etale dans le
lemme \ref{etale-cohomology-lemma-etale-topos-independent-partial-universe} ci-dessous.
Pour les notations et la terminologie employ\'ees dans ce lemme, nous renvoyons \`a
Topologies, section \ref{topologies-section-change-alpha}.

\begin{lemma}
\label{etale-cohomology-lemma-cohomology-enlarge-partial-universe}
Soit $\tau \in \{fppf, syntomic, smooth, \etale, Zariski\}$.
Soit $S$ un sch\'ema.
Soient $(\Sch/S)_\tau$ et $(\Sch'/S)_\tau$ deux
grands $\tau$-sites de $S$, et supposons que le premier soit contenu dans le second.
Dans ce cas,
\begin{enumerate}
\item pour tout faisceau ab\'elien $\mathcal{F}'$ d\'efini sur $(\Sch'/S)_\tau$ et
tout objet $U$ de $(\Sch/S)_\tau$, nous avons
$$
H^p_\tau(U, \mathcal{F}'|_{(\Sch/S)_\tau}) =
H^p_\tau(U, \mathcal{F}')
$$
Autrement dit : la cohomologie de $\mathcal{F}'$ sur $U$, calcul\'ee dans le site le plus grand,
co\"incide avec celle de la restriction de $\mathcal{F}'$ au site le plus petit,
sur $U$.
\item pour tout faisceau ab\'elien $\mathcal{F}$ sur $(\Sch/S)_\tau$, il existe un
faisceau ab\'elien $\mathcal{F}'$ sur $(\Sch/S)_\tau'$ dont la restriction \`a
$(\Sch/S)_\tau$ est isomorphe \`a $\mathcal{F}$.
\end{enumerate}
\end{lemma}

\begin{proof}
D'apr\`es Topologies, lemme \ref{topologies-lemma-change-alpha}, le foncteur d'inclusion
$(\Sch/S)_\tau \to (\Sch'/S)_\tau$ satisfait aux hypoth\`eses de
Sites, lemme \ref{sites-lemma-bigger-site}. Cela implique (2), et (1)
r\'esulte de
Cohomologie sur les sites, lemme \ref{sites-cohomology-lemma-cohomology-bigger-site}.
\end{proof}




\section{Le topos \'etale}
\label{etale-cohomology-section-etale-topos}

\noindent
Un {\it topos} est la cat\'egorie des faisceaux d'ensembles sur un site, voir
Sites, définition \ref{sites-definition-topos}. Il est donc d'usage
d'employer l'expression ``topos \'etale d'un sch\'ema'' pour d\'esigner
la cat\'egorie des faisceaux sur le petit site \'etale d'un sch\'ema.
En voici la d\'efinition formelle.

\begin{definition}
\label{etale-cohomology-definition-etale-topos}
Soit $S$ un sch\'ema.
\begin{enumerate}
\item Le {\it topos \'etale}, ou le {\it petit topos \'etale}
de $S$, est la cat\'egorie $\Sh(S_\etale)$ des faisceaux d'ensembles sur
le petit site \'etale de $S$.
\item Le {\it topos de Zariski}, ou le {\it petit topos de Zariski}
de $S$, est la cat\'egorie $\Sh(S_{Zar})$ des faisceaux d'ensembles sur le
petit site de Zariski de $S$.
\item Pour $\tau \in \{fppf, syntomic, smooth, \etale, Zariski\}$, un
{\it grand $\tau$-topos} est la cat\'egorie des faisceaux d'ensembles sur un
grand $\tau$-topos de $S$.
\end{enumerate}
\end{definition}

\noindent
Notons que le petit topos de Zariski de $S$ est simplement la cat\'egorie des faisceaux
d'ensembles sur l'espace topologique sous-jacent \`a $S$, voir
Topologies, lemme \ref{topologies-lemma-Zariski-usual}.
Alors que le petit topos \'etale ne d\'epend pas des choix faits lors de la
construction du petit site \'etale, les grands topos d\'ependent en g\'en\'eral
de ces choix.

\medskip\noindent
Il se trouve que le grand ou le petit topos \'etale ne d\'epend que de la
sous-cat\'egorie pleine de $(\Sch/S)_\etale$ ou de $S_\etale$ form\'ee des
sch\'emas affines, voir
Topologies, lemmes \ref{topologies-lemma-affine-big-site-etale} et
\ref{topologies-lemma-alternative}.
Nous l'utiliserons par exemple dans la d\'emonstration du lemme suivant.

\begin{lemma}
\label{etale-cohomology-lemma-etale-topos-independent-partial-universe}
Soit $S$ un sch\'ema. Le topos \'etale de $S$ est ind\'ependant
(\`a \'equivalence canonique pr\`es) de la construction du petit
site \'etale de la d\'efinition \ref{etale-cohomology-definition-tau-site}.
\end{lemma}

\begin{proof}
Nous devons montrer qu'\'etant donn\'es deux grands sites \'etales
$\Sch_\etale$ et $\Sch_\etale'$ contenant
$S$, on a $\Sh(S_\etale) \cong \Sh(S_\etale')$
avec les notations \'evidentes. D'apr\`es Topologies, lemme \ref{topologies-lemma-contained-in},
nous pouvons supposer $\Sch_\etale \subset \Sch_\etale'$.
D'apr\`es Ensembles, lemme \ref{sets-lemma-what-is-in-it},
tout sch\'ema affine \'etale sur $S$ est isomorphe \`a un objet
\`a la fois de $\Sch_\etale$ et de $\Sch_\etale'$.
Ainsi, le foncteur induit
$S_{affine, \etale} \to S_{affine, \etale}'$
est une \'equivalence. En outre, il est clair que ce foncteur
et un quasi-inverse transforment les recouvrements \'etales standards en
recouvrements \'etales standards.
Le r\'esultat d\'ecoule donc de
Topologies, lemme \ref{topologies-lemma-alternative}.
\end{proof}





\section{Cohomologie des faisceaux quasi-coh\'erents}
\label{etale-cohomology-section-cohomology-quasi-coherent}
%9.22.09

\noindent
Commen\c{c}ons par un lemme simple (qui reste vrai dans une g\'en\'eralit\'e plus grande que
celle de l'\'enonc\'e). Il affirme que le complexe de {\v C}ech d'un recouvrement standard est
\'egal au complexe de {\v C}ech d'un recouvrement fpqc de la forme
$\{\Spec(B) \to \Spec(A)\}$ o\`u $A \to B$ est fid\`element plat.

\begin{lemma}
\label{etale-cohomology-lemma-cech-complex}
Soit $\tau \in \{fppf, syntomic, smooth, \etale, Zariski\}$.
Soit $S$ un sch\'ema.
Soit $\mathcal{F}$ un faisceau ab\'elien sur $(\Sch/S)_\tau$, ou sur
$S_\tau$ dans le cas o\`u $\tau = \etale$, et soit
$\mathcal{U} = \{U_i \to U\}_{i \in I}$
un $\tau$-recouvrement standard de ce site.
Soit $V = \coprod_{i \in I} U_i$. Alors
\begin{enumerate}
\item $V$ est un sch\'ema affine,
\item $\mathcal{V} = \{V \to U\}$ est un recouvrement fpqc
et aussi un $\tau$-recouvrement, sauf si $\tau = Zariski$,
\item les complexes de {\v C}ech
$\check{\mathcal{C}}^\bullet (\mathcal{U}, \mathcal{F})$ et
$\check{\mathcal{C}}^\bullet (\mathcal{V}, \mathcal{F})$ co\"incident.
\end{enumerate}
\end{lemma}

\begin{proof}
La d\'efinition d'un $\tau$-recouvrement standard est donn\'ee dans
Topologies, définition
\ref{topologies-definition-standard-Zariski},
\ref{topologies-definition-standard-etale},
\ref{topologies-definition-standard-smooth},
\ref{topologies-definition-standard-syntomic}, et
\ref{topologies-definition-standard-fppf}.
Par d\'efinition, chacun des sch\'emas
$U_i$ est affine et $I$ est un ensemble fini. Ainsi, $V$ est un sch\'ema affine.
Il est clair que $V \to U$ est plat et surjectif ; par cons\'equent,
$\mathcal{V}$ est un recouvrement fpqc, voir
l'exemple \ref{etale-cohomology-example-fpqc-coverings}.
Sauf dans le cas de Zariski, le recouvrement $\mathcal{V}$
est aussi un $\tau$-recouvrement, voir
Topologies, définition
\ref{topologies-definition-etale-covering},
\ref{topologies-definition-smooth-covering},
\ref{topologies-definition-syntomic-covering}, et
\ref{topologies-definition-fppf-covering}.

\medskip\noindent
Notons que $\mathcal{U}$ est un raffinement de $\mathcal{V}$ ;
il existe donc un morphisme de complexes de {\v C}ech
$\check{\mathcal{C}}^\bullet (\mathcal{V}, \mathcal{F}) \to
\check{\mathcal{C}}^\bullet (\mathcal{U}, \mathcal{F})$, voir
Cohomologie sur les sites,
Equation (\ref{sites-cohomology-equation-map-cech-complexes}).
Ensuite, observons que si $T = \coprod_{j \in J} T_j$ est une
somme disjointe de sch\'emas du site sur lequel $\mathcal{F}$ est d\'efini,
alors la famille de morphismes \`a cible fix\'ee
$\{T_j \to T\}_{j \in J}$ est un recouvrement de Zariski, et donc
\begin{equation}
\label{etale-cohomology-equation-sheaf-coprod}
\mathcal{F}(T) =
\mathcal{F}(\coprod\nolimits_{j \in J} T_j) =
\prod\nolimits_{j \in J} \mathcal{F}(T_j)
\end{equation}
d'apr\`es la condition de faisceau pour $\mathcal{F}$.
Cela implique que le morphisme de complexes de {\v C}ech ci-dessus est un isomorphisme
en chaque degr\'e, car
$$
V \times_U \ldots \times_U V
=
\coprod\nolimits_{i_0, \ldots i_p} U_{i_0} \times_U \ldots \times_U U_{i_p}
$$
en tant que sch\'emas.
\end{proof}

\noindent
Notons que l'\'egalit\'e (\ref{etale-cohomology-equation-sheaf-coprod})
est fausse pour un pr\'efaisceau g\'en\'eral. M\^eme pour les faisceaux, elle n'est pas vraie sur tout
site, car les coproduits peuvent ne pas donner des recouvrements et ne pas \^etre disjoints.
Mais elle est vraie pour tous les sites usuels (du moins pour tous ceux que nous \'etudierons).

\begin{remark}
\label{etale-cohomology-remark-refinement}
Dans l'\'enonc\'e du lemme \ref{etale-cohomology-lemma-cech-complex}, le recouvrement $\mathcal{U}$
est un raffinement de $\mathcal{V}$, mais non l'inverse. Les recouvrements
de la forme $\{V \to U\}$ ne forment pas une sous-cat\'egorie initiale de la
cat\'egorie de tous les recouvrements de $U$. Il reste pourtant vrai que
nous pouvons calculer la cohomologie de {\v C}ech $\check H^n(U, \mathcal{F})$ (qui
est d\'efinie comme la limite inductive index\'ee par la cat\'egorie oppos\'ee \`a celle des
recouvrements $\mathcal{U}$ de $U$ des groupes de cohomologie de {\v C}ech de
$\mathcal{F}$ relativement \`a $\mathcal{U}$) au moyen des recouvrements
$\{V \to U\}$. Nous formulerons un lemme pr\'ecis (il ne vaut que pour les faisceaux)
et l'ajouterons ici si nous en avons un jour besoin.
\end{remark}

\begin{lemma}[Localit\'e de la cohomologie]
\label{etale-cohomology-lemma-locality-cohomology}
Soient $\mathcal{C}$ un site, $\mathcal{F}$ un faisceau ab\'elien sur $\mathcal{C}$,
$U$ un objet de $\mathcal{C}$, $p > 0$ un entier et $\xi \in
H^p(U, \mathcal{F})$. Alors il existe un recouvrement
$\mathcal{U} = \{U_i \to U\}_{i \in I}$ de $U$ dans $\mathcal{C}$
tel que $\xi |_{U_i} = 0$ pour tout $i \in I$.
\end{lemma}

\begin{proof}
Choisissons une r\'esolution injective $\mathcal{F} \to \mathcal{I}^\bullet$. Alors
$\xi$ est repr\'esent\'ee par un cocycle $\tilde{\xi} \in \mathcal{I}^p(U)$
tel que $d^p(\tilde{\xi}) = 0$. Par hypoth\`ese, la suite
$\mathcal{I}^{p - 1} \to \mathcal{I}^p \to \mathcal{I}^{p + 1}$ est exacte dans
$\textit{Ab}(\mathcal{C})$, ce qui signifie qu'il existe un recouvrement
$\mathcal{U} = \{U_i \to U\}_{i \in I}$ tel que
$\tilde{\xi}|_{U_i} = d^{p - 1}(\xi_i)$ pour un certain
$\xi_i \in \mathcal{I}^{p-1}(U_i)$. Puisque
la classe de cohomologie $\xi|_{U_i}$ est repr\'esent\'ee par le cocycle
$\tilde{\xi}|_{U_i}$, qui est un cobord, elle s'annule.
Pour plus de d\'etails, voir
Cohomologie sur les sites,
lemme \ref{sites-cohomology-lemma-kill-cohomology-class-on-covering}.
\end{proof}

\begin{theorem}
\label{etale-cohomology-theorem-zariski-fpqc-quasi-coherent}
Soit $S$ un sch\'ema et soit $\mathcal{F}$ un $\mathcal{O}_S$-module quasi-coh\'erent.
Soit $\mathcal{C}$ l'un des sites $(\Sch/S)_\tau$ pour
$\tau \in \{fppf, syntomic, smooth, \etale, Zariski\}$, ou
$S_\etale$. Alors
$$
H^p(S, \mathcal{F}) = H^p_\tau(S, \mathcal{F}^a)
$$
pour tout $p \geq 0$, o\`u
\begin{enumerate}
\item le membre de gauche d\'esigne la cohomologie usuelle du faisceau
$\mathcal{F}$ sur l'espace topologique sous-jacent au sch\'ema $S$, et
\item le membre de droite d\'esigne la cohomologie
du faisceau ab\'elien $\mathcal{F}^a$ (voir
la proposition \ref{etale-cohomology-proposition-quasi-coherent-sheaf-fpqc})
sur le site $\mathcal{C}$.
\end{enumerate}
\end{theorem}

\begin{proof}
Nous allons montrer que
$H^p(U, f^*\mathcal{F}) = H^p_\tau(U, \mathcal{F}^a)$
pour tout objet $f : U \to S$ du site $\mathcal{C}$.
Le r\'esultat est vrai pour $p = 0$ par la propri\'et\'e de faisceau.

\medskip\noindent
Supposons que $U$ soit affine. Nous voulons alors prouver que
$H^p_\tau(U, \mathcal{F}^a) = 0$ pour tout $p > 0$. Nous raisonnons par r\'ecurrence sur $p$.
\begin{enumerate}
\item[p = 1]
Prenons $\xi \in H^1_\tau(U, \mathcal{F}^a)$.
D'apr\`es le lemme \ref{etale-cohomology-lemma-locality-cohomology},
il existe un recouvrement fpqc $\mathcal{U} = \{U_i \to U\}_{i \in I}$
tel que $\xi|_{U_i} = 0$ pour tout $i \in I$. Quitte \`a raffiner
$\mathcal{U}$, nous pouvons supposer que $\mathcal{U}$ est un
$\tau$-recouvrement standard. En appliquant la suite spectrale du
th\'eor\`eme \ref{etale-cohomology-theorem-cech-ss},
nous voyons que $\xi$ provient d'une classe de cohomologie
$\check \xi \in \check H^1(\mathcal{U}, \mathcal{F}^a)$.
Consid\'erons le recouvrement $\mathcal{V} = \{\coprod_{i\in I} U_i \to U\}$. D'apr\`es le
lemme \ref{etale-cohomology-lemma-cech-complex},
$\check H^\bullet(\mathcal{U}, \mathcal{F}^a) =
\check H^\bullet(\mathcal{V}, \mathcal{F}^a)$.
D'autre part, puisque $\mathcal{V}$ est un recouvrement de la forme
$\{\Spec(B) \to \Spec(A)\}$ et $f^*\mathcal{F} = \widetilde{M}$
pour un certain $A$-module $M$, nous voyons que le complexe de {\v C}ech
$\check{\mathcal{C}}^\bullet(\mathcal{V}, \mathcal{F})$
n'est autre que le complexe $(B/A)_\bullet \otimes_A M$.
D'apr\`es le lemme \ref{etale-cohomology-lemma-descent-modules},
$H^p((B/A)_\bullet \otimes_A M) = 0$ pour $p > 0$ ; ainsi, $\check \xi = 0$
et donc $\xi = 0$.
\item[p > 1]
Prenons $\xi \in H^p_\tau(U, \mathcal{F}^a)$. D'apr\`es le
lemme \ref{etale-cohomology-lemma-locality-cohomology},
il existe un recouvrement fpqc $\mathcal{U} = \{U_i \to U\}_{i \in I}$
tel que $\xi|_{U_i} = 0$ pour tout $i \in I$. Quitte \`a raffiner
$\mathcal{U}$, nous pouvons supposer que $\mathcal{U}$ est un
$\tau$-recouvrement standard. Nous appliquons la suite spectrale du
th\'eor\`eme \ref{etale-cohomology-theorem-cech-ss}.
Observons que les intersections $U_{i_0} \times_U \ldots \times_U U_{i_p}$
sont affines ; par l'hypoth\`ese de r\'ecurrence, les groupes de cohomologie
$$
E_2^{p, q} = \check H^p(\mathcal{U}, \underline{H}^q(\mathcal{F}^a))
$$
s'annulent donc pour tout $0 < q < p$. Nous voyons que $\xi$ doit provenir d'un
$\check \xi \in \check H^p(\mathcal{U}, \mathcal{F}^a)$. En rempla\c{c}ant
$\mathcal{U}$ par le recouvrement $\mathcal{V}$ qui ne contient qu'un seul morphisme
et en utilisant de nouveau le lemme \ref{etale-cohomology-lemma-descent-modules},
nous voyons que la classe de cohomologie de {\v C}ech $\check \xi$ doit \^etre nulle,
donc $\xi = 0$.
\end{enumerate}
Supposons ensuite que $U$ soit s\'epar\'e. Choisissons des ouverts affines
$U = \bigcup_{i \in I} U_i$ formant un recouvrement de $U$. La famille
$\mathcal{U} = \{U_i \to U\}_{i \in I}$ est alors un recouvrement fpqc,
et toutes les intersections
$U_{i_0} \times_U \ldots \times_U U_{i_p}$ sont affines,
puisque $U$ est s\'epar\'e. Toutes les lignes de la suite spectrale du
th\'eor\`eme \ref{etale-cohomology-theorem-cech-ss}
sont donc nulles, sauf la ligne de degr\'e z\'ero. Par cons\'equent,
$$
H^p_\tau(U, \mathcal{F}^a) =
\check H^p(\mathcal{U}, \mathcal{F}^a) =
\check H^p(\mathcal{U}, \mathcal{F}) = H^p(U, \mathcal{F})
$$
o\`u la derni\`ere \'egalit\'e r\'esulte de la th\'eorie usuelle des sch\'emas, voir
Cohomologie des schémas, lemme
\ref{coherent-lemma-cech-cohomology-quasi-coherent}.

\medskip\noindent
Le cas g\'en\'eral est technique et (pour prolonger la d\'emonstration donn\'ee ici)
n\'ecessite une discussion sur les morphismes de suites spectrales ; nous ne le traiterons donc pas.
Il r\'esulte de
Descente, Proposition \ref{descent-proposition-same-cohomology-quasi-coherent},
(dont la d\'emonstration adopte une approche l\'eg\`erement diff\'erente), combin\'ee avec
Cohomologie sur les sites, lemme \ref{sites-cohomology-lemma-cohomology-of-open}.
\end{proof}

\begin{remark}
\label{etale-cohomology-remark-right-derived-global-sections}
Commentaire sur le th\'eor\`eme \ref{etale-cohomology-theorem-zariski-fpqc-quasi-coherent}.
Puisque $S$ est un objet final de la cat\'egorie $\mathcal{C}$, les groupes de cohomologie
du membre de droite sont simplement les foncteurs d\'eriv\'es \`a droite du
foncteur des sections globales. En fait, la d\'emonstration montre que
$H^p(U, f^*\mathcal{F}) = H^p_\tau(U, \mathcal{F}^a)$
pour tout objet $f : U \to S$ du site $\mathcal{C}$.
\end{remark}





\section{Exemples de faisceaux}
\label{etale-cohomology-section-examples-sheaves}

\noindent
Soient $S$ et $\tau$ comme dans la section \ref{etale-cohomology-section-big-small}.
Nous avons d\'ej\`a vu que tout pr\'efaisceau repr\'esentable est un faisceau sur
$(\Sch/S)_\tau$ ou $S_\tau$, voir
le lemme \ref{etale-cohomology-lemma-representable-sheaf-fpqc}
et la
remarque \ref{etale-cohomology-remark-fpqc-finest}.
Voici quelques cas particuliers.

\begin{definition}
\label{etale-cohomology-definition-additive-sheaf}
Sur l'un quelconque des sites $(\Sch/S)_\tau$ ou $S_\tau$ de la
section \ref{etale-cohomology-section-big-small} :
\begin{enumerate}
\item Le faisceau $T \mapsto \Gamma(T, \mathcal{O}_T)$ est not\'e
$\mathcal{O}_S$, ou $\mathbf{G}_a$, ou $\mathbf{G}_{a, S}$ si nous
voulons indiquer le sch\'ema de base.
\item De m\^eme, le faisceau
$T \mapsto \Gamma(T, \mathcal{O}^*_T)$ est not\'e $\mathcal{O}_S^*$, ou
$\mathbf{G}_m$, ou $\mathbf{G}_{m, S}$ si nous voulons
indiquer le sch\'ema de base.
\item Le {\it faisceau constant} $\underline{\mathbf{Z}/n\mathbf{Z}}$ sur un
site quelconque est le faisceau associ\'e au pr\'efaisceau constant
$U \mapsto \mathbf{Z}/n\mathbf{Z}$.
\end{enumerate}
\end{definition}

\noindent
Le premier est un faisceau, par exemple d'apr\`es le
th\'eor\`eme \ref{etale-cohomology-theorem-quasi-coherent}.
Le deuxi\`eme est un sous-pr\'efaisceau du premier, dont on v\'erifie facilement
qu'il est lui-m\^eme un faisceau. Le troisi\`eme est un faisceau par d\'efinition.
Notons que chacun de ces faisceaux est repr\'esentable.
Le premier et le deuxi\`eme le sont par les sch\'emas $\mathbf{G}_{a, S}$ et
$\mathbf{G}_{m, S}$, voir
Groupoïdes, section \ref{groupoids-section-group-schemes}.
Le troisi\`eme l'est par le sch\'ema en groupes fini \'etale $\mathbf{Z}/n\mathbf{Z}_S$,
parfois not\'e $(\mathbf{Z}/n\mathbf{Z})_S$,
qui consiste simplement en $n$ copies de $S$ munies
de l'\'evidente structure de sch\'ema en groupes sur $S$, voir
Groupoïdes, exemple \ref{groupoids-example-constant-group}
et la remarque suivante.

\begin{remark}
\label{etale-cohomology-remark-constant-locally-constant-maps}
Soit $G$ un groupe abstrait.
Sur l'un quelconque des sites $(\Sch/S)_\tau$ ou $S_\tau$ de la
section \ref{etale-cohomology-section-big-small},
le faisceau associ\'e $\underline{G}$
au pr\'efaisceau constant associ\'e \`a $G$ pour la
{\it topologie de Zariski} du site donne d\'ej\`a
$$
\Gamma(U, \underline{G}) =
\{\text{applications localement constantes pour la topologie de Zariski }U \to G\}
$$
Ce faisceau de Zariski est repr\'esentable par le sch\'ema en groupes $G_S$ d'apr\`es
Groupoïdes, exemple \ref{groupoids-example-constant-group}.
D'apr\`es le
lemme \ref{etale-cohomology-lemma-representable-sheaf-fpqc},
tout pr\'efaisceau repr\'esentable satisfait aussi la condition de faisceau pour la
topologie $\tau$ ; nous concluons donc que le faisceau associ\'e de Zariski
$\underline{G}$ ci-dessus est aussi le faisceau associ\'e pour la topologie $\tau$.
\end{remark}

\begin{definition}
\label{etale-cohomology-definition-structure-sheaf}
Soit $S$ un sch\'ema. Le {\it faisceau structural} de $S$ est le faisceau d'anneaux
$\mathcal{O}_S$
sur l'un quelconque des sites $S_{Zar}$, $S_\etale$ ou $(\Sch/S)_\tau$
consid\'er\'es ci-dessus.
\end{definition}

\noindent
S'il peut y avoir une ambigu\"it\'e quant au site sur lequel nous travaillons,
nous l'indiquerons au moyen d'indices. Par exemple, nous pouvons utiliser
$\mathcal{O}_{S_\etale}$ pour souligner que nous travaillons sur le
petit site \'etale de $S$.

\begin{remark}
\label{etale-cohomology-remark-special-case-fpqc-cohomology-quasi-coherent}
Dans la terminologie introduite ci-dessus, un cas particulier du
th\'eor\`eme \ref{etale-cohomology-theorem-zariski-fpqc-quasi-coherent}
est
$$
H_{fppf}^p(X, \mathbf{G}_a) =
H_\etale^p(X, \mathbf{G}_a) =
H_{Zar}^p(X, \mathbf{G}_a) =
H^p(X, \mathcal{O}_X)
$$
pour tout $p \geq 0$. De plus, nous pourrions utiliser la notation
$H^p_{fppf}(X, \mathcal{O}_X)$ pour d\'esigner la cohomologie du
faisceau structural sur le grand site fppf de $X$.
\end{remark}




\section{Groupes de Picard}
\label{etale-cohomology-section-picard-groups}

\noindent
Le th\'eor\`eme suivant est parfois appel\'e ``th\'eor\`eme 90 de Hilbert''.

\begin{theorem}
\label{etale-cohomology-theorem-picard-group}
Pour tout sch\'ema $X$, nous avons des identifications canoniques
\begin{align*}
H_{fppf}^1(X, \mathbf{G}_m) & = H^1_{syntomic}(X, \mathbf{G}_m) \\
& = H^1_{smooth}(X, \mathbf{G}_m) \\
& = H_\etale^1(X, \mathbf{G}_m) \\
& = H^1_{Zar}(X, \mathbf{G}_m) \\
& = \Pic(X) \\
& = H^1(X, \mathcal{O}_X^*)
\end{align*}
\end{theorem}

\begin{proof}
Soit $\tau$ l'une des topologies consid\'er\'ees dans la
section \ref{etale-cohomology-section-big-small}.
D'apr\`es
Cohomologie sur les sites, lemme
\ref{sites-cohomology-lemma-h1-invertible},
nous voyons que
$H^1_\tau(X, \mathbf{G}_m) =
H^1_\tau(X, \mathcal{O}_\tau^*) =
\Pic(\mathcal{O}_\tau)$,
o\`u $\mathcal{O}_\tau$ est le faisceau structural du site
$(\Sch/X)_\tau$. Or un $\mathcal{O}_\tau$-module inversible
est un $\mathcal{O}_\tau$-module quasi-coh\'erent.
D'apr\`es le th\'eor\`eme \ref{etale-cohomology-theorem-quasi-coherent}, ou plus pr\'ecis\'ement
Descente, Proposition \ref{descent-proposition-equivalence-quasi-coherent},
nous voyons que $\Pic(\mathcal{O}_\tau) = \Pic(X)$.
La derni\`ere \'egalit\'e se d\'emontre de la m\^eme mani\`ere.
\end{proof}






\section{Le site \'etale}
\label{etale-cohomology-section-etale-site}

\noindent
Nous commen\c{c}ons maintenant \`a examiner plus en d\'etail le site \'etale d'un
sch\'ema. Dans un premier temps, nous examinons bri\`evement la notion de
morphisme \'etale.





\section{Morphismes \'etales}
\label{etale-cohomology-section-etale-morphism}

\noindent
Pour plus de d\'etails, voir
Morphismes, section \ref{morphisms-section-etale}
pour la d\'efinition formelle, et
Morphismes étales, sections
\ref{etale-section-etale-morphisms},
\ref{etale-section-structure-etale-map},
\ref{etale-section-etale-smooth},
\ref{etale-section-topological-etale},
\ref{etale-section-functorial-etale}, et
\ref{etale-section-properties-permanence}
pour un aper\c{c}u de propri\'et\'es int\'eressantes des morphismes \'etales.

\medskip\noindent
Rappelons qu'une alg\`ebre $A$ sur un corps alg\'ebriquement clos $k$ est
{\it lisse} si elle est de type fini et si le module des diff\'erentielles
$\Omega_{A/k}$ est localement libre de type fini, de rang \`egal \`a la dimension.
Un sch\'ema $X$ sur $k$ est {\it lisse} sur $k$ s'il est localement de type
fini et si chaque ouvert affine est le spectre d'une $k$-alg\`ebre lisse.
Si $k$ n'est pas alg\'ebriquement clos, alors une $k$-alg\`ebre $A$ est
une $k$-alg\`ebre lisse si $A \otimes_k \overline{k}$ est une
$\overline{k}$-alg\`ebre lisse. Un morphisme d'anneaux $A \to B$ est lisse s'il est
plat, de pr\'esentation finie et si, pour tout id\'eal premier $\mathfrak p \subset A$,
l'anneau de la fibre $\kappa(\mathfrak p) \otimes_A B$ est lisse sur le corps
r\'esiduel $\kappa(\mathfrak p)$. Plus g\'en\'eralement, un morphisme de sch\'emas est
{\it lisse} s'il est plat, localement de pr\'esentation finie et si ses
fibres g\'eom\'etriques sont lisses.

\medskip\noindent
Pour ces faits, voir
Morphismes, section \ref{morphisms-section-smooth}.
Cela nous permet de d\'efinir un morphisme \'etale comme suit.

\begin{definition}
\label{etale-cohomology-definition-etale-morphism}
Un morphisme de sch\'emas est {\it \'etale} s'il est lisse de dimension relative 0.
\end{definition}

\noindent
En particulier, un morphisme de sch\'emas $X \to S$ est \'etale s'il est lisse
et si $\Omega_{X/S} = 0$.

\begin{proposition}
\label{etale-cohomology-proposition-etale-morphisms}
Propri\'et\'es des morphismes \'etales.
\begin{enumerate}
\item Soit $k$ un corps. Un morphisme de sch\'emas $U \to \Spec(k)$ est
\'etale si et seulement si $U \cong \coprod_{i \in I} \Spec(k_i)$,
o\`u, pour chaque $i \in I$,
le corps $k_i$ est une extension finie s\'eparable de $k$.
\item Soit $\varphi : U \to S$ un morphisme de sch\'emas. Les
conditions suivantes sont \`equivalentes :
\begin{enumerate}
\item $\varphi$ est \'etale,
\item $\varphi$ est localement de pr\'esentation finie, plat, et toutes ses fibres sont
\'etales,
\item $\varphi$ est plat, non ramifi\'e et localement de pr\'esentation finie.
\end{enumerate}
\item Un morphisme d'anneaux $A \to B$ est \'etale si et seulement si
$B \cong A[x_1, \ldots, x_n]/(f_1, \ldots, f_n)$
tel que $\Delta = \det \left( \frac{\partial f_i}{\partial x_j} \right)$
soit inversible dans $B$.
\item Le changement de base d'un morphisme \'etale est \'etale.
\item Les compos\'ees de morphismes \'etales sont \'etales.
\item Les produits fibr\'es et les produits de morphismes \'etales sont \'etales.
\item Un morphisme \'etale a pour dimension relative 0.
\item Soit $Y \to X$ un morphisme \'etale.
Si $X$ est r\'eduit (respectivement r\'egulier), alors $Y$ l'est aussi.
\item Les morphismes \'etales sont ouverts.
\item Si $X \to S$ et $Y \to S$ sont \'etales, alors tout
$S$-morphisme $X \to Y$ est lui aussi \'etale.
\end{enumerate}
\end{proposition}

\begin{proof}
Nous avons d\'emontr\'e ces faits (et davantage) dans les chapitres pr\'ec\'edents.
Voici une liste de r\'ef\'erences :
(1) Morphismes, lemme \ref{morphisms-lemma-etale-over-field}.
(2) Morphismes, lemmes \ref{morphisms-lemma-etale-flat-etale-fibres}
et \ref{morphisms-lemma-flat-unramified-etale}.
(3) Algèbre, lemme \ref{algebra-lemma-etale-standard-smooth}.
(4) Morphismes, lemme \ref{morphisms-lemma-base-change-etale}.
(5) Morphismes, lemme \ref{morphisms-lemma-composition-etale}.
(6) D\'ecoule formellement de (4) et (5).
(7) Morphismes, lemmes \ref{morphisms-lemma-etale-locally-quasi-finite}
et \ref{morphisms-lemma-locally-quasi-finite-rel-dimension-0}.
(8) Voir Algèbre, lemmes \ref{algebra-lemma-reduced-goes-up} et
\ref{algebra-lemma-Rk-goes-up} ; voir aussi d'autres r\'esultats de ce type
dans Morphismes étales, section \ref{etale-section-properties-permanence}.
(9) Voir Morphismes, lemme \ref{morphisms-lemma-fppf-open} et
\ref{morphisms-lemma-etale-flat}.
(10) Voir Morphismes, lemme \ref{morphisms-lemma-etale-permanence}.
\end{proof}

\begin{definition}
\label{etale-cohomology-definition-standard-etale}
Un morphisme d'anneaux $A \to B$ est dit {\it \'etale standard} si
$B \cong \left(A[t]/(f)\right)_g$ avec $f, g \in A[t]$, o\`u $f$ est unitaire,
et si $\text{d}f/\text{d}t$ est inversible dans $B$.
\end{definition}

\noindent
Il est vrai qu'un morphisme d'anneaux \'etale standard est \'etale. En effet, supposons
que $B = \left(A[t]/(f)\right)_g$ avec $f, g \in A[t]$, o\`u $f$ est unitaire,
et que $\text{d}f/\text{d}t$ soit inversible dans $B$. Alors $A[t]/(f)$ est un
$A$-module libre de rang \'egal au degr\'e du polyn\^ome unitaire $f$.
Par cons\'equent, $B$, en tant que localisation de cette alg\`ebre qui est libre comme module, est de pr\'esentation finie
et plat sur $A$. Pour achever la preuve que $B$ est \'etale, il suffit
de montrer que les anneaux des fibres
$$
\kappa(\mathfrak p) \otimes_A B
\cong
\kappa(\mathfrak p) \otimes_A (A[t]/(f))_g
\cong
\kappa(\mathfrak p)[t, 1/\overline{g}]/(\overline{f})
$$
sont des produits finis d'extensions de corps finies s\'eparables.
Ici, $\overline{f}, \overline{g} \in \kappa(\mathfrak p)[t]$ sont
les images de $f$ et $g$. Soit
$$
\overline{f} = \overline{f}_1 \ldots \overline{f}_a
\overline{f}_{a + 1}^{e_1} \ldots \overline{f}_{a + b}^{e_b}
$$
la factorisation de $\overline{f}$ en puissances de facteurs unitaires irr\'eductibles
deux \`a deux distincts $\overline{f}_i$, avec $e_1, \ldots, e_b > 1$.
Par hypoth\`ese, $\text{d}\overline{f}/\text{d}t$ est inversible dans
$\kappa(\mathfrak p)[t, 1/\overline{g}]$. Nous voyons donc que
au moins les $\overline{f}_i$, $i > a$, sont inversibles. Nous concluons
que
$$
\kappa(\mathfrak p)[t, 1/\overline{g}]/(\overline{f})
\cong
\prod\nolimits_{i \in I} \kappa(\mathfrak p)[t]/(\overline{f}_i)
$$
o\`u $I \subset \{1, \ldots, a\}$ est le sous-ensemble des indices $i$ tels que
$\overline{f}_i$ ne divise pas $\overline{g}$. De plus, l'image de
$\text{d}\overline{f}/\text{d}t$ dans le facteur
$\kappa(\mathfrak p)[t]/(\overline{f}_i)$ est clairement \'egale au produit d'une
unit\'e par $\text{d}\overline{f}_i/\text{d}t$. Nous en concluons que
$\kappa_i = \kappa(\mathfrak p)[t]/(\overline{f}_i)$ est une extension de corps
finie de $\kappa(\mathfrak p)$, engendr\'ee par un \'el\'ement dont le
polyn\^ome minimal est s\'eparable, c'est-\`a-dire que l'extension de corps
$\kappa_i/\kappa(\mathfrak p)$ est finie s\'eparable, comme souhait\'e.

\medskip\noindent
Il se trouve que tout morphisme d'anneaux \'etale est localement \'etale standard.
Pour formuler cela, introduisons la notation suivante.
Un morphisme d'anneaux $A \to B$ est {\it \'etale en un id\'eal premier $\mathfrak q$} de $B$ s'il
existe $h \in B$, $h \not \in \mathfrak q$, tel que $A \to B_h$ soit \'etale.
Voici le r\'esultat.

\begin{theorem}
\label{etale-cohomology-theorem-standard-etale}
Un morphisme d'anneaux $A \to B$ est \'etale en un id\'eal premier $\mathfrak q$ si et seulement s'il
existe $g \in B$, $g \not \in \mathfrak q$, tel que $B_g$ soit \'etale
standard sur $A$.
\end{theorem}

\begin{proof}
Voir
Algèbre, Proposition \ref{algebra-proposition-etale-locally-standard}.
\end{proof}





\section{Recouvrements \'etales}
\label{etale-cohomology-section-etale-covering}

\noindent
Rappelons la d\'efinition.

\begin{definition}
\label{etale-cohomology-definition-etale-covering}
Un {\it recouvrement \'etale} d'un sch\'ema $U$ est une famille de morphismes
de sch\'emas
$\{\varphi_i : U_i \to U\}_{i \in I}$ telle que
\begin{enumerate}
\item chaque $\varphi_i$ est un morphisme \'etale,
\item les $U_i$ recouvrent $U$, c'est-\`a-dire $U = \bigcup_{i\in I}\varphi_i(U_i)$.
\end{enumerate}
\end{definition}

\begin{lemma}
\label{etale-cohomology-lemma-etale-fpqc}
Tout recouvrement \'etale est un recouvrement fpqc.
\end{lemma}

\begin{proof}
(Voir aussi
Topologies,
lemme \ref{topologies-lemma-zariski-etale-smooth-syntomic-fppf-fpqc}.)
Soit $\{\varphi_i : U_i \to U\}_{i \in I}$ un recouvrement \'etale.
Puisqu'un morphisme \'etale est plat et que les \'el\'ements du recouvrement doivent
recouvrir sa cible, la propri\'et\'e fp (fid\`element plat) est satisfaite.
Pour v\'erifier la propri\'et\'e qc (quasi-compacit\'e), soit $V \subset U$ un ouvert
affine, et \'ecrivons $\varphi_i^{-1}(V) = \bigcup_{j \in J_i} V_{ij}$
pour certains ouverts affines $V_{ij} \subset U_i$. Puisque $\varphi_i$ est ouvert
(car les morphismes \'etales sont ouverts), nous voyons que
$V = \bigcup_{i\in I} \bigcup_{j \in J_i} \varphi_i(V_{ij})$
est un recouvrement ouvert de $V$.
De plus, puisque $V$ est quasi-compact, ce recouvrement admet un
raffinement fini.
\end{proof}

\noindent
Ainsi, tout \'enonc\'e vrai pour les recouvrements fpqc
reste vrai {\it a fortiori} pour les recouvrements \'etales. Par
exemple, le site \'etale est sous-canonique.

\begin{definition}
\label{etale-cohomology-definition-big-etale-site}
(Pour plus de d\'etails, voir la section \ref{etale-cohomology-section-big-small}, ou
Topologies, section \ref{topologies-section-etale}.)
Soit $S$ un sch\'ema.
Le {\it grand site \'etale de $S$} est le site
$(\Sch/S)_\etale$, voir
la d\'efinition \ref{etale-cohomology-definition-tau-site}.
Le {\it petit site \'etale de $S$} est le site $S_\etale$, voir
la d\'efinition \ref{etale-cohomology-definition-tau-site}.
Nous d\'efinissons de m\^eme le {\it grand} et le {\it petit site de Zariski} de $S$,
not\'es $(\Sch/S)_{Zar}$ et $S_{Zar}$.
\end{definition}

\noindent
Grosso modo, le grand site \'etale de $S$ est constitu\'e des sch\'emas sur $S$
et ses recouvrements sont les recouvrements \'etales. Le petit site \'etale de $S$ est constitu\'e
des sch\'emas \'etales sur $S$, avec pour recouvrements les recouvrements \'etales.
En fait, tout morphisme entre objets de $S_\etale$ est \'etale, en
vertu de la
proposition \ref{etale-cohomology-proposition-etale-morphisms} ;
par cons\'equent, pour v\'erifier que $\{U_i \to U\}_{i \in I}$ dans $S_\etale$
est un recouvrement, il suffit de v\'erifier que $\coprod U_i \to U$ est surjectif.

\medskip\noindent
Le petit site \'etale a moins d'objets que le grand site \'etale ; il
ne contient que les ``ouverts'' de la topologie \'etale sur $S$. C'est une sous-cat\'egorie
pleine du grand site \'etale, et sa topologie est induite par la
topologie du grand site. Ainsi, le foncteur de restriction
du grand site \'etale au petit est exact et transforme les objets injectifs en
objets injectifs. Cela a la cons\'equence suivante.

\begin{proposition}
\label{etale-cohomology-proposition-cohomology-restrict-small-site}
Soit $S$ un sch\'ema et soit $\mathcal{F}$ un faisceau ab\'elien sur
$(\Sch/S)_\etale$.
Alors $\mathcal{F}|_{S_\etale}$ est un faisceau sur $S_\etale$ et
$$
H^p_\etale(S, \mathcal{F}|_{S_\etale}) =
H^p_\etale(S, \mathcal{F})
$$
pour tout $p \geq 0$.
\end{proposition}

\begin{proof}
C'est un cas particulier du lemme \ref{etale-cohomology-lemma-compare-cohomology-big-small}.
\end{proof}

\noindent
Conform\'ement \`a la notation g\'en\'erale introduite dans la
section \ref{etale-cohomology-section-big-small},
nous notons $H_\etale^p(S, \mathcal{F})$ le groupe de cohomologie ci-dessus.





%9.24.09
\section{Th\'eorie de Kummer}
\label{etale-cohomology-section-kummer}

\noindent
Soit $n \in \mathbf{N}$ et consid\'erons le foncteur $\mu_n$ d\'efini par
$$
\begin{matrix}
\Sch^{opp} & \longrightarrow & \textit{Ab} \\
S & \longmapsto &
\mu_n(S)
=
\{t \in \Gamma(S, \mathcal{O}_S^*) \mid t^n = 1 \}.
\end{matrix}
$$
D'apr\`es
Groupoïdes, exemple \ref{groupoids-example-roots-of-unity},
ce foncteur est repr\'esentable, et le sch\'ema qui le repr\'esente
est lui aussi not\'e $\mu_n$. D'apr\`es le
lemme \ref{etale-cohomology-lemma-representable-sheaf-fpqc},
ce foncteur satisfait \`a la condition de faisceau pour la topologie fpqc
(en particulier, il satisfait aussi \`a la condition de faisceau pour les
topologies \'etale, de Zariski, etc.).

\begin{lemma}
\label{etale-cohomology-lemma-kummer-sequence}
Si $n\in \mathcal{O}_S^*$, alors
$$
0 \to
\mu_{n, S} \to
\mathbf{G}_{m, S} \xrightarrow{(\cdot)^n}
\mathbf{G}_{m, S} \to 0
$$
est une suite exacte courte de faisceaux aussi bien sur le petit que sur le
grand site \'etale de $S$.
\end{lemma}

\begin{proof}
Par d\'efinition, le faisceau $\mu_{n, S}$ est le noyau de l'application
$(\cdot)^n$. Il suffit donc de montrer que la derni\`ere application est surjective.
Soit $U$ un sch\'ema sur $S$. Soit
$f \in \mathbf{G}_m(U) = \Gamma(U, \mathcal{O}_U^*)$.
Nous devons montrer que nous pouvons trouver un recouvrement \'etale de
$U$ tel que la restriction de $f$ \`a chacun de ses membres soit une puissance $n$-i\`eme.
Posons
$$
U' =
\underline{\Spec}_U(\mathcal{O}_U[T]/(T^n-f))
\xrightarrow{\pi}
U.
$$
(Voir
Constructions, section \ref{constructions-section-spec-via-glueing} ou
\ref{constructions-section-spec}
pour une pr\'esentation du spectre relatif.)
Soit $\Spec(A) \subset U$ un ouvert affine, et supposons que $f|_{\Spec(A)}$ corresponde
\`a l'unit\'e $a \in A^*$. Alors $\pi^{-1}(\Spec(A)) = \Spec(B)$ avec
$B = A[T]/(T^n - a)$. Le morphisme d'anneaux $A \to B$ est fini libre de rang $n$,
donc fid\`element plat, et nous en concluons que
$\Spec(B) \to \Spec(A)$ est surjectif. Comme cela vaut pour tout
ouvert affine de $U$, nous concluons que $\pi$ est surjectif.
De plus, $n$ et $T^{n - 1}$ sont inversibles dans $B$, donc
$nT^{n-1} \in B^*$ et le morphisme d'anneaux $A \to B$ est \'etale standard,
en particulier \'etale. Comme cela vaut pour tout ouvert affine de $U$,
nous concluons que $\pi$ est \'etale. Par cons\'equent,
$\mathcal{U} = \{\pi : U' \to U\}$ est un recouvrement \'etale.
De plus, $f|_{U'} = (f')^n$, o\`u $f'$ est la classe de $T$
dans $\Gamma(U', \mathcal{O}_{U'}^*)$ ; ainsi, $\mathcal{U}$ poss\`ede la propri\'et\'e voulue.
\end{proof}

\begin{remark}
\label{etale-cohomology-remark-no-kummer-sequence-zariski}
Le lemme \ref{etale-cohomology-lemma-kummer-sequence} est faux lorsque ``\'etale'' est remplac\'e
par ``Zariski''.
Comme la topologie \'etale est moins fine que la topologie lisse, voir
Topologies, lemme \ref{topologies-lemma-zariski-etale-smooth},
il s'ensuit que la suite est aussi exacte pour la topologie lisse.
\end{remark}

\noindent
D'apr\`es le
th\'eor\`eme \ref{etale-cohomology-theorem-picard-group},
le
lemme \ref{etale-cohomology-lemma-kummer-sequence}
et les propri\'et\'es g\'en\'erales de la cohomologie, nous obtenons
la suite exacte longue de cohomologie
$$
\xymatrix{
0 \ar[r] &
H_\etale^0(S, \mu_{n, S}) \ar[r] &
\Gamma(S, \mathcal{O}_S^*) \ar[r]^{(\cdot)^n} &
\Gamma(S, \mathcal{O}_S^*) \ar@(rd, ul)[rdllllr]
\\
& H_\etale^1(S, \mu_{n, S}) \ar[r] &
\Pic(S) \ar[r]^{(\cdot)^n} &
\Pic(S) \ar@(rd, ul)[rdllllr] \\
& H_\etale^2(S, \mu_{n, S}) \ar[r] &
\ldots
}
$$
du moins si $n$ est inversible sur $S$. Lorsque $n$ n'est pas inversible sur $S$,
nous pouvons appliquer le lemme suivant.

\begin{lemma}
\label{etale-cohomology-lemma-kummer-sequence-syntomic}
Pour tout $n \in \mathbf{N}$, la suite
$$
0 \to
\mu_{n, S} \to
\mathbf{G}_{m, S} \xrightarrow{(\cdot)^n}
\mathbf{G}_{m, S} \to 0
$$
est une suite exacte courte de faisceaux sur les sites
$(\Sch/S)_{fppf}$ et $(\Sch/S)_{syntomic}$.
\end{lemma}

\begin{proof}
Par d\'efinition, le faisceau $\mu_{n, S}$ est le noyau de l'application
$(\cdot)^n$. Il suffit donc de montrer que la derni\`ere application est surjective.
Comme la topologie syntomique est moins fine que la topologie fppf, voir
Topologies, lemme \ref{topologies-lemma-zariski-etale-smooth-syntomic-fppf},
il suffit de le d\'emontrer pour la topologie syntomique.
Soit $U$ un sch\'ema sur $S$. Soit
$f \in \mathbf{G}_m(U) = \Gamma(U, \mathcal{O}_U^*)$.
Nous devons montrer que nous pouvons trouver un recouvrement syntomique de
$U$ tel que la restriction de $f$ \`a chacun de ses membres soit une puissance $n$-i\`eme.
Posons
$$
U' =
\underline{\Spec}_U(\mathcal{O}_U[T]/(T^n-f))
\xrightarrow{\pi}
U.
$$
(Voir
Constructions, section \ref{constructions-section-spec-via-glueing} ou
\ref{constructions-section-spec}
pour une pr\'esentation du spectre relatif.)
Soit $\Spec(A) \subset U$ un ouvert affine, et supposons que $f|_{\Spec(A)}$ corresponde
\`a l'unit\'e $a \in A^*$. Alors $\pi^{-1}(\Spec(A)) = \Spec(B)$ avec
$B = A[T]/(T^n - a)$. Le morphisme d'anneaux $A \to B$ est fini libre de rang $n$,
donc fid\`element plat, et nous en concluons que
$\Spec(B) \to \Spec(A)$ est surjectif. Comme cela vaut pour tout
ouvert affine de $U$, nous concluons que $\pi$ est surjectif.
De plus, $B$ est une intersection compl\`ete globale relative sur $A$, donc
le morphisme d'anneaux $A \to B$ est syntomique standard,
en particulier syntomique. Comme cela vaut pour tout ouvert affine de $U$,
nous concluons que $\pi$ est syntomique. Par cons\'equent,
$\mathcal{U} = \{\pi : U' \to U\}$ est un recouvrement syntomique.
De plus, $f|_{U'} = (f')^n$, o\`u $f'$ est la classe de $T$
dans $\Gamma(U', \mathcal{O}_{U'}^*)$ ; ainsi, $\mathcal{U}$ poss\`ede la propri\'et\'e voulue.
\end{proof}

\begin{remark}
\label{etale-cohomology-remark-no-kummer-sequence-smooth-etale-zariski}
Le lemme \ref{etale-cohomology-lemma-kummer-sequence-syntomic}
est faux pour les topologies lisse, \'etale ou de Zariski.
\end{remark}

\noindent
D'apr\`es le
th\'eor\`eme \ref{etale-cohomology-theorem-picard-group},
le
lemme \ref{etale-cohomology-lemma-kummer-sequence-syntomic}
et les propri\'et\'es g\'en\'erales de la cohomologie, nous obtenons
la suite exacte longue de cohomologie
$$
\xymatrix{
0 \ar[r] &
H_{fppf}^0(S, \mu_{n, S}) \ar[r] &
\Gamma(S, \mathcal{O}_S^*) \ar[r]^{(\cdot)^n} &
\Gamma(S, \mathcal{O}_S^*) \ar@(rd, ul)[rdllllr]
\\
& H_{fppf}^1(S, \mu_{n, S}) \ar[r] &
\Pic(S) \ar[r]^{(\cdot)^n} &
\Pic(S) \ar@(rd, ul)[rdllllr] \\
& H_{fppf}^2(S, \mu_{n, S}) \ar[r] &
\ldots
}
$$
pour tout sch\'ema $S$ et tout entier $n$. Bien entendu, il existe une suite analogue
en cohomologie syntomique.

\medskip\noindent
Soit $n \in \mathbf{N}$ et soit $S$ un sch\'ema quelconque.
Il existe une autre fa\c{c}on, plus directe, de d\'ecrire le premier groupe de cohomologie \`a
valeurs dans $\mu_n$. Consid\'erons les paires
$(\mathcal{L}, \alpha)$ o\`u $\mathcal{L}$ est un faisceau inversible sur $S$
et $\alpha : \mathcal{L}^{\otimes n} \to \mathcal{O}_S$ est une trivialisation
de la $n$-i\`eme puissance tensorielle de $\mathcal{L}$.
Soit $(\mathcal{L}', \alpha')$ une seconde paire de ce type.
Un isomorphisme $\varphi : (\mathcal{L}, \alpha) \to (\mathcal{L}', \alpha')$
est un isomorphisme $\varphi : \mathcal{L} \to \mathcal{L}'$ de faisceaux
inversibles tel que le diagramme
$$
\xymatrix{
\mathcal{L}^{\otimes n} \ar[d]_{\varphi^{\otimes n}} \ar[r]_\alpha &
\mathcal{O}_S \ar[d]^1 \\
(\mathcal{L}')^{\otimes n} \ar[r]^{\alpha'} &
\mathcal{O}_S \\
}
$$
soit commutatif. Nous avons donc
\begin{equation}
\label{etale-cohomology-equation-isomorphisms-pairs}
\mathit{Isom}_S((\mathcal{L}, \alpha), (\mathcal{L}', \alpha'))
=
\left\{
\begin{matrix}
\emptyset & \text{si} & \text{elles ne sont pas isomorphes} \\
H^0(S, \mu_{n, S})\cdot \varphi & \text{si} &
\varphi \text{ est un isomorphisme de paires}
\end{matrix}
\right.
\end{equation}
De plus, \'etant donn\'ees deux paires $(\mathcal{L}, \alpha)$, $(\mathcal{L}', \alpha')$,
le produit tensoriel
$$
(\mathcal{L}, \alpha) \otimes (\mathcal{L}', \alpha')
=
(\mathcal{L} \otimes \mathcal{L}', \alpha \otimes \alpha')
$$
est encore une paire. La paire $(\mathcal{O}_S, 1)$ est un \'el\'ement neutre pour cette
op\'eration de produit tensoriel, et un inverse est donn\'e par
$$
(\mathcal{L}, \alpha)^{-1} = (\mathcal{L}^{\otimes -1}, \alpha^{\otimes -1}).
$$
Ainsi, l'ensemble des classes d'isomorphisme de paires forme un groupe ab\'elien.
Notons que
$$
(\mathcal{L}, \alpha)^{\otimes n}
=
(\mathcal{L}^{\otimes n}, \alpha^{\otimes n})
\xrightarrow{\alpha}
(\mathcal{O}_S, 1)
$$
est un isomorphisme ;
par cons\'equent, tout \'el\'ement de ce groupe est d'ordre divisant $n$. Nous avertissons le lecteur
que ce groupe n'est en g\'en\'eral {\bf pas} la $n$-torsion de $\Pic(S)$.

\begin{lemma}
\label{etale-cohomology-lemma-describe-h1-mun}
Soit $S$ un sch\'ema. Il existe une identification canonique
$$
H_\etale^1(S, \mu_n) =
\text{groupe des paires }(\mathcal{L}, \alpha)\text{ \`a isomorphisme pr\`es comme ci-dessus}
$$
si $n$ est inversible sur $S$. En g\'en\'eral, nous avons
$$
H_{fppf}^1(S, \mu_n) =
\text{groupe des paires }(\mathcal{L}, \alpha)\text{ \`a isomorphisme pr\`es comme ci-dessus}.
$$
Le m\^eme r\'esultat vaut en rempla\c{c}ant fppf par syntomique.
\end{lemma}

\begin{proof}
Nous d\'emontrons d'abord le second isomorphisme.
Soit $(\mathcal{L}, \alpha)$ une paire comme ci-dessus.
Choisissons un recouvrement de $S = \bigcup U_i$ par des ouverts affines tel que
$\mathcal{L}|_{U_i} \cong \mathcal{O}_{U_i}$. Choisissons un g\'en\'erateur $s_i \in \mathcal{L}(U_i)$.
Alors $\alpha(s_i^{\otimes n}) = f_i \in \mathcal{O}_S^*(U_i)$.
En \'ecrivant $U_i = \Spec(A_i)$, nous voyons qu'il existe une
intersection compl\`ete globale relative $A_i \to B_i = A_i[T]/(T^n - f_i)$
telle que $f_i$ devienne une puissance $n$-i\`eme dans $B_i$. Autrement dit, en posant
$V_i = \Spec(B_i)$, nous obtenons un recouvrement syntomique
$\mathcal{V} = \{V_i \to S\}_{i \in I}$ et des trivialisations
$\varphi_i : (\mathcal{L}, \alpha)|_{V_i} \to (\mathcal{O}_{V_i}, 1)$.

\medskip\noindent
Nous utiliserons ce r\'esultat (l'existence du recouvrement $\mathcal{V}$)
pour associer \`a cette paire une classe de cohomologie dans
$H^1_{syntomic}(S, \mu_{n, S})$. Nous donnons deux constructions (\'equivalentes).

\medskip\noindent
Premi\`ere construction : au moyen de la cohomologie de {\v C}ech.
Sur les intersections doubles $V_i \times_S V_j$, nous avons l'isomorphisme
$$
(\mathcal{O}_{V_i \times_S V_j}, 1)
\xrightarrow{\text{pr}_0^*\varphi_i^{-1}}
(\mathcal{L}|_{V_i \times_S V_j}, \alpha|_{V_i \times_S V_j})
\xrightarrow{\text{pr}_1^*\varphi_j}
(\mathcal{O}_{V_i \times_S V_j}, 1)
$$
de paires. D'apr\`es (\ref{etale-cohomology-equation-isomorphisms-pairs}), il est donn\'e par un
\'el\'ement $\zeta_{ij} \in \mu_n(V_i \times_S V_j)$. Nous omettons la v\'erification
que ces $\zeta_{ij}$ forment un $1$-cocycle, c'est-\`a-dire qu'ils donnent
un \'el\'ement $(\zeta_{i_0i_1}) \in \check C(\mathcal{V}, \mu_n)$
tel que $d(\zeta_{i_0i_1}) = 0$. Sa classe est donc un \'el\'ement de
$\check H^1(\mathcal{V}, \mu_n)$ et, d'apr\`es le
th\'eor\`eme \ref{etale-cohomology-theorem-cech-ss},
elle a pour image une classe de cohomologie dans $H^1_{syntomic}(S, \mu_{n, S})$.

\medskip\noindent
Seconde construction : au moyen des torseurs. Consid\'erons le pr\'efaisceau
$$
\mu_n(\mathcal{L}, \alpha) :
U
\longmapsto
\mathit{Isom}_U((\mathcal{O}_U, 1), (\mathcal{L}, \alpha)|_U)
$$
sur $(\Sch/S)_{syntomic}$.
Nous pouvons le voir comme un sous-pr\'efaisceau de
$\SheafHom_\mathcal{O}(\mathcal{O}, \mathcal{L})$ (faisceau Hom interne,
voir
Modules sur les sites, section \ref{sites-modules-section-internal-hom}).
Comme les conditions qui d\'efinissent ce sous-pr\'efaisceau sont locales, nous voyons que c'est
un faisceau.
D'apr\`es (\ref{etale-cohomology-equation-isomorphisms-pairs}), ce faisceau est muni d'une action libre du
faisceau $\mu_{n, S}$. Il reste donc seulement \`a v\'erifier qu'il
admet localement des sections. C'est vrai gr\^ace \`a l'existence du
recouvrement trivialisant $\mathcal{V}$. Ainsi, $\mu_n(\mathcal{L}, \alpha)$
est un $\mu_{n, S}$-torseur et, d'apr\`es
Cohomologie sur les sites, lemme \ref{sites-cohomology-lemma-torsors-h1},
nous obtenons un \'el\'ement correspondant de $H^1_{syntomic}(S, \mu_{n, S})$.

\medskip\noindent
Il nous reste maintenant \`a montrer les points suivants :
\begin{enumerate}
\item Les deux constructions donnent la m\^eme classe de cohomologie.
\item Des paires isomorphes donnent la m\^eme classe de cohomologie.
\item La classe de cohomologie de
$(\mathcal{L}, \alpha) \otimes (\mathcal{L}', \alpha')$
est la somme des classes de cohomologie de
$(\mathcal{L}, \alpha)$ et $(\mathcal{L}', \alpha')$.
\item Si la classe de cohomologie est triviale, alors la paire est triviale.
\item Tout \'el\'ement de $H^1_{syntomic}(S, \mu_{n, S})$ est la
classe de cohomologie d'une paire.
\end{enumerate}
Nous omettons la preuve de (1). Le point (2) ressort clairement de la seconde construction,
puisque des torseurs isomorphes donnent la m\^eme classe de cohomologie.
Le point (3) ressort clairement de la premi\`ere construction, puisque les classes de
{\v C}ech obtenues s'additionnent. Le point (4) ressort clairement de la seconde construction,
puisqu'un torseur est trivial si et seulement s'il poss\`ede une section globale, voir
Cohomologie sur les sites, lemme \ref{sites-cohomology-lemma-trivial-torsor}.

\medskip\noindent
Le point (5) se voit comme suit (bien qu'une preuve directe soit
pr\'ef\'erable). Supposons que $\xi \in H^1_{syntomic}(S, \mu_{n, S})$.
Alors $\xi$ s'envoie sur un \'el\'ement
$\overline{\xi} \in H^1_{syntomic}(S, \mathbf{G}_{m, S})$
tel que $n \overline{\xi} = 0$. D'apr\`es le
th\'eor\`eme \ref{etale-cohomology-theorem-picard-group},
nous voyons que $\overline{\xi}$ correspond \`a un faisceau inversible $\mathcal{L}$
dont la $n$-i\`eme puissance tensorielle est isomorphe \`a $\mathcal{O}_S$.
Il existe donc une paire $(\mathcal{L}, \alpha')$ dont la classe de cohomologie
$\xi'$ a la m\^eme image $\overline{\xi'}$ dans
$H^1_{syntomic}(S, \mathbf{G}_{m, S})$. Il suffit donc de montrer
que $\xi - \xi'$ est la classe d'une paire. Par construction et d'apr\`es la
suite exacte longue de cohomologie ci-dessus, nous voyons que
$\xi - \xi' = \partial(f)$ pour un certain $f \in H^0(S, \mathcal{O}_S^*)$.
Consid\'erons la paire $(\mathcal{O}_S, f)$. Nous omettons la v\'erification
que la classe de cohomologie de cette paire est $\partial(f)$, ce qui
ach\`eve la preuve de la premi\`ere identification (avec fppf remplac\'ee
par syntomique).

\medskip\noindent
Pour voir la premi\`ere identification, notons que si $n$ est inversible sur $S$, alors le
recouvrement $\mathcal{V}$ construit dans la premi\`ere partie de la preuve
est en fait un recouvrement \'etale (comparer avec la preuve du
lemme \ref{etale-cohomology-lemma-kummer-sequence}). Le reste de la preuve est ind\'ependant
de la topologie, \`a l'exception du tout dernier argument, qui utilise l'exactitude de
la suite de Kummer, c'est-\`a-dire le lemme \ref{etale-cohomology-lemma-kummer-sequence}.
\end{proof}






\section{Voisinages, fibres et points}
\label{etale-cohomology-section-stalks}

\noindent
Nous pouvons associer \`a tout point g\'eom\'etrique de $S$ un foncteur fibre qui est
exact. Un morphisme de faisceaux sur $S_\etale$ est un isomorphisme si et seulement
s'il
induit un isomorphisme sur toutes ces fibres. Un complexe de faisceaux ab\'eliens est
exact si et seulement si le complexe des fibres est exact en tout point g\'eom\'etrique.
Ces faits, pris ensemble, signifient que le petit site \'etale d'un sch\'ema $S$
poss\`ede assez de points. Il se trouve aussi que tout point du petit topos \'etale
de $S$ (notion abstraite) est donn\'e par un point g\'eom\'etrique.
Ainsi, en un certain sens, le petit topos \'etale de $S$ peut se comprendre en
termes de points g\'eom\'etriques et de voisinages.

\begin{definition}
\label{etale-cohomology-definition-geometric-point}
Soit $S$ un sch\'ema.
\begin{enumerate}
\item Un {\it point g\'eom\'etrique} de $S$ est un morphisme
$\Spec(k) \to S$ o\`u $k$ est alg\'ebriquement clos.
Un tel point se note habituellement $\overline{s}$, autrement dit au moyen d'une
lettre minuscule surmont\'ee d'une barre. Nous employons souvent $\overline{s}$ pour d\'esigner aussi bien le sch\'ema
$\Spec(k)$ que le morphisme, et nous employons $\kappa(\overline{s})$
pour d\'esigner $k$.
\item Nous disons que $\overline{s}$ {\it est au-dessus de} $s$
pour indiquer que $s \in S$ est l'image de $\overline{s}$.
\item Un {\it voisinage \'etale} d'un point g\'eom\'etrique $\overline{s}$
de $S$ est un diagramme commutatif
$$
\xymatrix{
& U \ar[d]^\varphi \\
{\overline{s}} \ar[r]^{\overline{s}} \ar[ur]^{\bar u} & S
}
$$
o\`u $\varphi$ est un morphisme \'etale de sch\'emas.
Nous \'ecrivons $(U, \overline{u}) \to (S, \overline{s})$.
\item Un {\it morphisme de voisinages \'etales}
$(U, \overline{u}) \to (U', \overline{u}')$
est un $S$-morphisme $h: U \to U'$
tel que $\overline{u}' = h \circ \overline{u}$.
\end{enumerate}
\end{definition}

\begin{remark}
\label{etale-cohomology-remark-etale-between-etale}
Puisque $U$ et $U'$ sont \'etales sur $S$, tout $S$-morphisme
entre eux est lui aussi \'etale, voir la
proposition \ref{etale-cohomology-proposition-etale-morphisms}.
En particulier, tous les morphismes de voisinages \'etales sont \'etales.
\end{remark}

\begin{remark}
\label{etale-cohomology-remark-etale-neighbourhoods}
Soient $S$ un sch\'ema et $s \in S$ un point. Dans
Compléments sur les morphismes,
définition \ref{more-morphisms-definition-etale-neighbourhood},
nous avons d\'efini la notion de voisinage \'etale $(U, u) \to (S, s)$
de $(S, s)$. Si $\overline{s}$ est un point g\'eom\'etrique de $S$ au-dessus de
$s$, alors tout voisinage \'etale $(U, \overline{u}) \to (S, \overline{s})$
donne lieu \`a un voisinage \'etale $(U, u)$ de $(S, s)$ en prenant pour
$u \in U$ l'unique point de $U$ tel que $\overline{u}$
soit au-dessus de $u$. R\'eciproquement, \'etant donn\'e un voisinage \'etale $(U, u)$
de $(S, s)$, l'extension de corps r\'esiduels $\kappa(u)/\kappa(s)$
est finie s\'eparable (voir la
proposition \ref{etale-cohomology-proposition-etale-morphisms}) ;
nous pouvons donc trouver un plongement $\kappa(u) \subset \kappa(\overline{s})$
qui induit l'identit\'e sur $\kappa(s)$. Autrement dit, nous pouvons trouver un point g\'eom\'etrique
$\overline{u}$ de $U$ au-dessus de $u$ tel que $(U, \overline{u})$
soit un voisinage \'etale de $(S, \overline{s})$.
Nous utiliserons ces observations pour passer d'un type de
voisinage \'etale \`a l'autre.
\end{remark}

\begin{lemma}
\label{etale-cohomology-lemma-cofinal-etale}
Soit $S$ un sch\'ema, et soit $\overline{s}$ un point g\'eom\'etrique de $S$.
La cat\'egorie des voisinages \'etales est cofiltrante. Plus pr\'ecis\'ement :
\begin{enumerate}
\item Soient $(U_i, \overline{u}_i)_{i = 1, 2}$ deux voisinages \'etales de
$\overline{s}$ dans $S$. Alors il existe un troisi\`eme voisinage \'etale
$(U, \overline{u})$ et des morphismes
$(U, \overline{u}) \to (U_i, \overline{u}_i)$, $i = 1, 2$.
\item Soient $h_1, h_2: (U, \overline{u}) \to (U', \overline{u}')$ deux
morphismes entre voisinages \'etales de $\overline{s}$. Alors il existe un
voisinage \'etale $(U'', \overline{u}'')$ et un morphisme
$h : (U'', \overline{u}'') \to (U, \overline{u})$
qui \'egalise $h_1$ et $h_2$, c'est-\`a-dire tel que
$h_1 \circ h = h_2 \circ h$.
\end{enumerate}
\end{lemma}

\begin{proof}
Pour le point (1), consid\'erons le produit fibr\'e $U = U_1 \times_S U_2$.
Il est \'etale sur $U_1$ et sur $U_2$, puisque les morphismes \'etales sont
pr\'eserv\'es par changement de base, voir la
proposition \ref{etale-cohomology-proposition-etale-morphisms}.
Le morphisme $\overline{s} \to U$ d\'efini par $(\overline{u}_1, \overline{u}_2)$
lui conf\`ere la structure d'un voisinage \'etale muni de morphismes vers
$U_1$ et $U_2$. Pour le point (2), d\'efinissons $U''$ comme le produit fibr\'e
$$
\xymatrix{
U'' \ar[r] \ar[d] & U \ar[d]^{(h_1, h_2)} \\
U' \ar[r]^-\Delta & U' \times_S U'.
}
$$
Puisque les compos\'ees de $\overline{u}$ et $\overline{u}'$ vers $S$ sont toutes deux \'egales \`a $\overline{s}$,
nous voyons que $\overline{u}'' = (\overline{u}, \overline{u}')$ est un point g\'eom\'etrique
de $U''$. En particulier, $U'' \not = \emptyset$.
De plus, puisque $U'$ est \'etale sur $S$, le produit fibr\'e
$U'\times_S U'$ l'est aussi (voir la
proposition \ref{etale-cohomology-proposition-etale-morphisms}).
Ainsi, la fl\`eche verticale $(h_1, h_2)$ est \'etale d'apr\`es la
remarque \ref{etale-cohomology-remark-etale-between-etale} ci-dessus.
Par cons\'equent, $U''$ est \'etale sur $U'$ par changement de base, et donc aussi
\'etale sur $S$ (car les compos\'ees de morphismes \'etales sont \'etales).
Ainsi, $(U'', \overline{u}'')$ est une solution au probl\`eme.
\end{proof}

\begin{lemma}
\label{etale-cohomology-lemma-geometric-lift-to-cover}
Soit $S$ un sch\'ema.
Soit $\overline{s}$ un point g\'eom\'etrique de $S$.
Soit $(U, \overline{u})$ un voisinage \'etale de $\overline{s}$.
Soit $\mathcal{U} = \{\varphi_i : U_i \to U \}_{i\in I}$ un recouvrement \'etale.
Alors il existe $i \in I$ et $\overline{u}_i : \overline{s} \to U_i$
tels que $\varphi_i : (U_i, \overline{u}_i) \to (U, \overline{u})$
soit un morphisme de voisinages \'etales.
\end{lemma}

\begin{proof}
Comme $U = \bigcup_{i\in I} \varphi_i(U_i)$, le produit fibr\'e
$\overline{s} \times_{\overline{u}, U, \varphi_i} U_i$ n'est pas vide
pour un certain $i$. Consid\'erons alors le diagramme cart\'esien
$$
\xymatrix{
\overline{s} \times_{\overline{u}, U, \varphi_i} U_i
\ar[d]^{\text{pr}_1} \ar[r]_-{\text{pr}_2} & U_i
\ar[d]^{\varphi_i} \\
\Spec(k) = \overline{s} \ar@/^1pc/[u]^\sigma
\ar[r]^-{\overline{u}} & U
}
$$
La projection $\text{pr}_1$ est le changement de base d'un morphisme \'etale ; elle
est donc \'etale, voir la
proposition \ref{etale-cohomology-proposition-etale-morphisms}.
Par cons\'equent, $\overline{s} \times_{\overline{u}, U, \varphi_i} U_i$
est une somme disjointe d'extensions finies s\'eparables de $k$, d'apr\`es la
proposition \ref{etale-cohomology-proposition-etale-morphisms}. Ici,
$\overline{s} = \Spec(k)$. Mais $k$ est alg\'ebriquement clos, donc toutes
ces extensions sont triviales, et il existe une section $\sigma$ de
$\text{pr}_1$. La compos\'ee
$\text{pr}_2 \circ \sigma$ donne un morphisme compatible avec $\overline{u}$.
\end{proof}

\begin{definition}
\label{etale-cohomology-definition-stalk}
Soit $S$ un sch\'ema.
Soit $\mathcal{F}$ un pr\'efaisceau sur $S_\etale$.
Soit $\overline{s}$ un point g\'eom\'etrique de $S$.
La {\it fibre} de $\mathcal{F}$ en $\overline{s}$ est
$$
\mathcal{F}_{\overline{s}}
=
\colim_{(U, \overline{u})} \mathcal{F}(U)
$$
o\`u $(U, \overline{u})$ parcourt tous les voisinages
\'etales de $\overline{s}$ dans $S$.
\end{definition}

\noindent
D'apr\`es le lemme \ref{etale-cohomology-lemma-cofinal-etale}, cette limite inductive est index\'ee par une cat\'egorie
filtrante, \`a savoir l'oppos\'ee de la cat\'egorie des voisinages \'etales.
Autrement dit, un \'el\'ement de $\mathcal{F}_{\overline{s}}$ peut \'etre
consid\'er\'e comme un triplet $(U, \overline{u}, \sigma)$ o\`u
$\sigma \in \mathcal{F}(U)$. Deux triplets
$(U, \overline{u}, \sigma)$, $(U', \overline{u}', \sigma')$
d\'efinissent le m\^eme \'el\'ement de la fibre s'il existe un troisi\`eme
voisinage \'etale $(U'', \overline{u}'')$ et des morphismes de voisinages
\'etales $h : (U'', \overline{u}'') \to (U, \overline{u})$,
$h' : (U'', \overline{u}'') \to (U', \overline{u}')$ tels que
$h^*\sigma = (h')^*\sigma'$ dans $\mathcal{F}(U'')$. Voir
Catégories, section \ref{categories-section-directed-colimits}.

\begin{lemma}
\label{etale-cohomology-lemma-stalk-gives-point}
Soient $S$ un sch\'ema et $\overline{s}$ un point g\'eom\'etrique de $S$.
Consid\'erons le foncteur
\begin{align*}
u : S_\etale & \longrightarrow \textit{Ensembles}, \\
U & \longmapsto
|U_{\overline{s}}|
=
\{\overline{u} \text{ tel que }(U, \overline{u})
\text{ soit un voisinage \'etale de }\overline{s}\}.
\end{align*}
Ici, $|U_{\overline{s}}|$ d\'esigne l'ensemble sous-jacent \`a la fibre g\'eom\'etrique.
Alors $u$ d\'efinit un point $p$ du site $S_\etale$
(Sites, définition \ref{sites-definition-point})
et le foncteur fibre associ\'e $\mathcal{F} \mapsto \mathcal{F}_p$
(Sites, Equation \ref{sites-equation-stalk})
est le foncteur $\mathcal{F} \mapsto \mathcal{F}_{\overline{s}}$
d\'efini ci-dessus.
\end{lemma}

\begin{proof}
Dans la preuve du
lemme \ref{etale-cohomology-lemma-geometric-lift-to-cover},
nous avons vu que le sch\'ema $U_{\overline{s}}$ est une somme disjointe de
sch\'emas isomorphes \`a $\overline{s}$. Nous pouvons donc aussi consid\'erer
$|U_{\overline{s}}|$ comme l'ensemble des points g\'eom\'etriques de $U$ situ\'es au-dessus de
$\overline{s}$, c'est-\`a-dire comme l'ensemble des morphismes
$\overline{u} : \overline{s} \to U$ qui s'inscrivent dans le diagramme de la
d\'efinition \ref{etale-cohomology-definition-geometric-point}.
Il s'ensuit que $u(S)$ est un singleton, et que
$u(U \times_V W) = u(U) \times_{u(V)} u(W)$
chaque fois que $U \to V$ et $W \to V$ sont des morphismes dans $S_\etale$.
Et, \'etant donn\'e un recouvrement $\{U_i \to U\}_{i \in I}$ dans $S_\etale$,
nous voyons que $\coprod u(U_i) \to u(U)$ est surjectif par le
lemme \ref{etale-cohomology-lemma-geometric-lift-to-cover}.
Par cons\'equent,
Sites, Proposition \ref{sites-proposition-point-limits}
s'applique, de sorte que $p$ est un point du site $S_\etale$.
Enfin, notre foncteur $\mathcal{F} \mapsto \mathcal{F}_{\overline{s}}$
est donn\'e par exactement la m\^eme limite inductive que le foncteur
$\mathcal{F} \mapsto \mathcal{F}_p$ associ\'e \`a $p$ dans
Sites, Equation \ref{sites-equation-stalk},
ce qui d\'emontre la derni\`ere assertion.
\end{proof}

\begin{remark}
\label{etale-cohomology-remark-map-stalks}
Soit $S$ un sch\'ema et soient $\overline{s} : \Spec(k) \to S$
et $\overline{s}' : \Spec(k') \to S$ deux points g\'eom\'etriques de
$S$. Un {\it morphisme $a : \overline{s} \to \overline{s}'$ de points g\'eom\'etriques}
est simplement un morphisme $a : \Spec(k) \to \Spec(k')$ tel que
$\overline{s}' \circ a = \overline{s}$. \'Etant donn\'e un tel morphisme, nous obtenons
un foncteur de la cat\'egorie des voisinages \'etales de $\overline{s}'$
vers la cat\'egorie des voisinages \'etales de $\overline{s}$ par la r\`egle
$(U, \overline{u}') \mapsto (U, \overline{u}' \circ a)$. Nous obtenons donc
une application canonique
$$
\mathcal{F}_{\overline{s}'}
=
\colim_{(U, \overline{u}')} \mathcal{F}(U)
\longrightarrow
\colim_{(U, \overline{u})} \mathcal{F}(U)
=
\mathcal{F}_{\overline{s}}
$$
donn\'ee par Catégories, lemme \ref{categories-lemma-functorial-colimit}. En utilisant la
description des \'el\'ements des fibres par des triplets, elle envoie l'\'el\'ement de
$\mathcal{F}_{\overline{s}'}$ repr\'esent\'e par le triplet
$(U, \overline{u}', \sigma)$ sur l'\'el\'ement de $\mathcal{F}_{\overline{s}}$
repr\'esent\'e par le triplet $(U, \overline{u}' \circ a, \sigma)$.
Comme le foncteur ci-dessus est manifestement une \'equivalence, nous concluons que cette
application canonique est un isomorphisme de foncteurs fibres.

\medskip\noindent
Assurons-nous que l'application de fibres correspondant \`a $a$ est orient\'ee
dans le bon sens. Notons que ce qui pr\'ec\`ede signifie, d'apr\`es
Sites, définition \ref{sites-definition-morphism-points},
que $a$ d\'efinit un morphisme $a : p \to p'$ entre les points $p, p'$ du
site $S_\etale$ associ\'es \`a $\overline{s}, \overline{s}'$ par le
lemme \ref{etale-cohomology-lemma-stalk-gives-point}. Il existe des morphismes plus g\'en\'eraux de
points (correspondant \`a des sp\'ecialisations de points de $S$), que nous
d\'ecrirons plus loin et qui ne seront pas des isomorphismes, voir la section
\ref{etale-cohomology-section-specialization}.
\end{remark}

\begin{lemma}
\label{etale-cohomology-lemma-stalk-exact}
Soit $S$ un sch\'ema. Soit $\overline{s}$ un point g\'eom\'etrique de $S$.
\begin{enumerate}
\item Le foncteur fibre
$\textit{PAb}(S_\etale) \to \textit{Ab}$,
$\mathcal{F}  \mapsto  \mathcal{F}_{\overline{s}}$
est exact.
\item On a $(\mathcal{F}^\#)_{\overline{s}} = \mathcal{F}_{\overline{s}}$
pour tout pr\'efaisceau d'ensembles $\mathcal{F}$ sur $S_\etale$.
\item Le foncteur
$\textit{Ab}(S_\etale) \to \textit{Ab}$,
$\mathcal{F} \mapsto \mathcal{F}_{\overline{s}}$ est exact.
\item De m\^eme, les foncteurs
$\textit{PSh}(S_\etale) \to \textit{Ensembles}$ et
$\Sh(S_\etale) \to \textit{Ensembles}$ donn\'es par le foncteur fibre
$\mathcal{F} \mapsto \mathcal{F}_{\overline{s}}$ sont exacts (voir
Catégories, définition \ref{categories-definition-exact})
et commutent aux limites inductives arbitraires.
\end{enumerate}
\end{lemma}

\begin{proof}
Avant d'indiquer comment le d\'emontrer par des arguments directs,
notons que le r\'esultat d\'ecoule de la th\'eorie g\'en\'erale de
Modules sur les sites, section \ref{sites-modules-section-stalks}.
Cela est vrai parce que $\mathcal{F} \mapsto \mathcal{F}_{\overline{s}}$
provient d'un point du petit site \'etale de $S$, voir le
lemme \ref{etale-cohomology-lemma-stalk-gives-point}.
Nous ne donnerons une preuve directe que de (1), (2) et (3), et omettrons
une preuve directe de (4).

\medskip\noindent
L'exactitude comme foncteur sur $\textit{PAb}(S_\etale)$ d\'ecoule formellement du
fait que les limites inductives filtrantes commutent \`a toutes les limites inductives et aux
limites finies. L'identification des fibres dans (2) se fait au moyen de l'application
$$
\kappa :
\mathcal{F}_{\overline{s}}
\longrightarrow
(\mathcal{F}^\#)_{\overline{s}}
$$
induite par le morphisme naturel $\mathcal{F}\to \mathcal{F}^\#$, voir le
th\'eor\`eme \ref{etale-cohomology-theorem-sheafification}.
Nous affirmons que cette application est un isomorphisme de groupes ab\'eliens. Nous montrerons
l'injectivit\'e et omettrons la preuve de la surjectivit\'e.

\medskip\noindent
Soit $\sigma\in \mathcal{F}_{\overline{s}}$.
Il existe un voisinage \'etale
$(U, \overline{u})\to (S, \overline{s})$ tel que $\sigma$ soit l'image d'une
section $s \in \mathcal{F}(U)$. Si $\kappa(\sigma) = 0$ dans
$(\mathcal{F}^\#)_{\overline{s}}$, alors il existe un morphisme de voisinages \'etales
$(U', \overline{u}')\to (U, \overline{u})$ tel que
$s|_{U'}$ soit nul dans $\mathcal{F}^\#(U')$. Il s'ensuit qu'il
existe un recouvrement \'etale
$\{U_i'\to U'\}_{i\in I}$ tel que $s|_{U_i'}=0$ dans
$\mathcal{F}(U_i')$ pour tout $i$. Par le lemme \ref{etale-cohomology-lemma-geometric-lift-to-cover},
il existe $i \in I$ et un morphisme
$\overline{u}_i': \overline{s} \to U_i'$ tels que
$(U_i', \overline{u}_i') \to (U', \overline{u}')\to (U, \overline{u})$
soient des morphismes de voisinages \'etales. Ainsi $\sigma = 0$
puisque $(U_i', \overline{u}_i') \to (U, \overline{u})$
est un morphisme de voisinages \'etales tel que
l'on ait $s|_{U'_i}=0$. Cela d\'emontre que $\kappa$ est injectif.

\medskip\noindent
Pour montrer que le foncteur $\textit{Ab}(S_\etale) \to \textit{Ab}$ est
exact, consid\'erons une suite exacte courte quelconque dans $\textit{Ab}(S_\etale)$ :
$
0\to \mathcal{F}\to \mathcal{G}\to \mathcal H \to 0.
$
Elle nous donne la suite exacte de pr\'efaisceaux
$$
0 \to \mathcal{F} \to \mathcal{G} \to \mathcal H \to
\mathcal H/^p\mathcal{G} \to 0,
$$
o\`u $/^p$ d\'esigne le quotient dans $\textit{PAb}(S_\etale)$.
En prenant les fibres en
$\overline{s}$, nous voyons que $(\mathcal H /^p\mathcal{G})_{\bar{s}} =
(\mathcal H /\mathcal{G})_{\bar{s}} = 0$, puisque le faisceau associ\'e \`a
$\mathcal H/^p\mathcal{G}$ est $0$.
Par cons\'equent,
$$
0\to \mathcal{F}_{\overline{s}} \to \mathcal{G}_{\overline{s}} \to
\mathcal{H}_{\overline{s}} \to 0 = (\mathcal H/^p\mathcal{G})_{\overline{s}}
$$
est exacte, puisque prendre les fibres est exact comme foncteur sur les pr\'efaisceaux.
\end{proof}

\begin{theorem}
\label{etale-cohomology-theorem-exactness-stalks}
Soit $S$ un sch\'ema.
Un morphisme $a : \mathcal{F} \to \mathcal{G}$ de faisceaux d'ensembles est injectif
(resp.\ surjectif) si et seulement si le morphisme sur les fibres
$a_{\overline{s}} : \mathcal{F}_{\overline{s}} \to \mathcal{G}_{\overline{s}}$
est injectif (resp.\ surjectif) pour tout point g\'eom\'etrique de $S$.
Une suite de faisceaux ab\'eliens sur $S_\etale$ est exacte
si et seulement si elle est exacte sur toutes les fibres aux points g\'eom\'etriques de $S$.
\end{theorem}

\begin{proof}
La n\'ecessit\'e de l'exactitude sur les fibres d\'ecoule du
lemme \ref{etale-cohomology-lemma-stalk-exact}.
R\'eciproquement, il suffit de montrer qu'un morphisme de faisceaux est surjectif
(respectivement injectif) si et seulement s'il est surjectif (respectivement
injectif) sur toutes les fibres. Nous le d\'emontrons dans le cas de la surjectivit\'e et omettons
la preuve dans le cas de l'injectivit\'e.

\medskip\noindent
Soit $\alpha : \mathcal{F} \to \mathcal{G}$ un morphisme de faisceaux tel
que $\mathcal{F}_{\overline{s}} \to \mathcal{G}_{\overline{s}}$
soit surjectif pour tout point g\'eom\'etrique. Fixons
$U \in \Ob(S_\etale)$
et $s \in \mathcal{G}(U)$. Pour tout $u \in U$, choisissons un
$\overline{u} \to U$ situ\'e au-dessus de $u$ et un voisinage \'etale
$(V_u , \overline{v}_u) \to (U, \overline{u})$ tel que
$s|_{V_u} = \alpha(s_{V_u})$ pour un certain
$s_{V_u} \in \mathcal{F}(V_u)$.
C'est possible puisque $\alpha$ est surjectif sur les
fibres. Alors $\{V_u \to U\}_{u \in U}$
est un recouvrement \'etale sur lequel les restrictions de $s$
appartiennent \`a l'image du morphisme $\alpha$.
Ainsi, $\alpha$ est surjectif, voir
Sites, section \ref{sites-section-sheaves-injective}.
\end{proof}

\begin{remarks}
\label{etale-cohomology-remarks-enough-points}
Quelques remarques sur les points des sites g\'eom\'etriques.
\begin{enumerate}
\item Le th\'eor\`eme \ref{etale-cohomology-theorem-exactness-stalks} dit que la famille de points
de $S_\etale$ donn\'ee par les points g\'eom\'etriques de $S$
(lemme \ref{etale-cohomology-lemma-stalk-gives-point}) est conservative, voir
Sites, définition \ref{sites-definition-enough-points}.
En particulier, $S_\etale$ a assez de points.
\item Supposons que $\mathcal{F}$ soit un faisceau sur le grand site \'etale
\label{etale-cohomology-item-stalks-big}
de $S$. Soit $T \to S$ un objet du grand site \'etale de $S$,
et soit $\overline{t}$ un point g\'eom\'etrique de $T$. Nous d\'efinissons alors
$\mathcal{F}_{\overline{t}}$ comme la fibre
de la restriction $\mathcal{F}|_{T_\etale}$ de $\mathcal{F}$
au petit site \'etale de $T$. Autrement dit, nous pouvons d\'efinir
la fibre de $\mathcal{F}$ en tout point g\'eom\'etrique de tout
sch\'ema $T/S \in \Ob((\Sch/S)_\etale)$.
\item Le grand site \'etale de $S$ a lui aussi assez de points, en
consid\'erant tous les points g\'eom\'etriques de tous les objets de ce site, voir
(\ref{etale-cohomology-item-stalks-big}).
\end{enumerate}
\end{remarks}

\noindent
Le lemme suivant devrait \^etre omis en premi\`ere lecture.

\begin{lemma}
\label{etale-cohomology-lemma-points-small-etale-site}
Soit $S$ un sch\'ema.
\begin{enumerate}
\item Soit $p$ un point du petit site \'etale
$S_\etale$ de $S$ donn\'e par un foncteur
$u : S_\etale \to \textit{Ensembles}$.
Alors il existe un point g\'eom\'etrique $\overline{s}$ de $S$ tel que
$p$ soit isomorphe au point de $S_\etale$ associ\'e \`a
$\overline{s}$ dans le
lemme \ref{etale-cohomology-lemma-stalk-gives-point}.
\item Soit $p : \Sh(pt) \to \Sh(S_\etale)$ un point
du petit topos \'etale de $S$. Alors $p$ provient d'un point g\'eom\'etrique
de $S$, c'est-\`a-dire que le foncteur fibre $\mathcal{F} \mapsto \mathcal{F}_p$
est isomorphe \`a un foncteur fibre tel que d\'efini dans la
d\'efinition \ref{etale-cohomology-definition-stalk}.
\end{enumerate}
\end{lemma}

\begin{proof}
D'apr\`es
Sites, lemme \ref{sites-lemma-point-site-topos},
il existe une correspondance bijective entre les points du site et les points
du topos associ\'e ; il suffit donc de d\'emontrer (1).
D'apr\`es
Sites, Proposition \ref{sites-proposition-point-limits},
le foncteur $u$ a les propri\'et\'es suivantes :
(a) $u(S) = \{*\}$, (b) $u(U \times_V W) = u(U) \times_{u(V)} u(W)$, et
(c) si $\{U_i \to U\}$ est un recouvrement \'etale, alors
$\coprod u(U_i) \to u(U)$ est surjectif.
En particulier, si $U' \subset U$ est un sous-sch\'ema ouvert,
alors $u(U') \subset u(U)$. De plus, d'apr\`es
Sites, lemme \ref{sites-lemma-point-site-topos},
nous pouvons \'ecrire $u(U) = p^{-1}(h_U^\#)$ ; autrement dit, $u(U)$ est la
fibre du faisceau repr\'esentable $h_U$. Si
$U = V \amalg W$, alors nous voyons que $h_U = (h_V \amalg h_W)^\#$ et obtenons
$u(U) = u(V) \amalg u(W)$ puisque $p^{-1}$ est exact.

\medskip\noindent
Consid\'erons la restriction de $u$ \`a $S_{Zar}$. D'apr\`es
Sites, Exemples \ref{sites-example-point-topological} et
\ref{sites-example-point-topology},
il existe un unique point $s \in S$ tel que, pour $S' \subset S$ ouvert, nous
ayons $u(S') = \{*\}$ si $s \in S'$ et $u(S') = \emptyset$ si $s \not \in S'$.
Notons que si $\varphi : U \to S$ est un objet de $S_\etale$, alors
$\varphi(U) \subset S$ est ouvert (voir la
proposition \ref{etale-cohomology-proposition-etale-morphisms})
et $\{U \to \varphi(U)\}$ est un recouvrement \'etale. Nous concluons donc que
$u(U) = \emptyset \Leftrightarrow s \not \in \varphi(U)$.

\medskip\noindent
Choisissons un point g\'eom\'etrique $\overline{s} : \overline{s} \to S$ au-dessus de $s$, voir la
d\'efinition \ref{etale-cohomology-definition-geometric-point}
pour l'abus de notation usuel. Supposons que $\varphi : U \to S$ soit un objet
de $S_\etale$ avec $U$ affine. Notons que $\varphi$ est s\'epar\'e et
que la fibre $U_s$ de $\varphi$ au-dessus de $s$ est un sch\'ema affine sur
$\Spec(\kappa(s))$ et qu'elle est le spectre d'un produit fini
d'extensions finies s\'eparables $k_i$ de $\kappa(s)$. Nous pouvons donc appliquer
Morphismes étales, lemme \ref{etale-lemma-etale-etale-local-technical}
pour obtenir un voisinage \'etale $(V, \overline{v})$ de $(S, \overline{s})$
tel que
$$
U \times_S V = U_1 \amalg \ldots \amalg U_n \amalg W
$$
o\`u $U_i \to V$ est un isomorphisme et $W$ n'a aucun point au-dessus de
$\overline{v}$. Nous concluons donc que
$$
u(U) \times u(V) =
u(U \times_S V) =
u(U_1) \amalg \ldots \amalg u(U_n) \amalg u(W)
$$
et, bien s\^ur, aussi $u(U_i) = u(V)$. Apr\`es avoir un peu r\'etr\'eci $V$, nous pouvons
supposer que $V$ a exactement un point au-dessus de $s$ et que, par cons\'equent, $W$ n'a aucun
point au-dessus de $s$. D'apr\`es ce qui pr\'ec\`ede, cela donne alors $u(W) = \emptyset$.
Nous obtenons donc
$$
u(U) \times u(V) =
u(U_1) \amalg \ldots \amalg u(U_n) =
\coprod\nolimits_{i = 1, \ldots, n} u(V)
$$
o\`u les signes d'\'egalit\'e sont compatibles avec les applications vers l'ensemble $u(V)$.
Notons que $u(V) \not = \emptyset$ puisque $s$ est dans l'image de $V \to S$.
En particulier, nous voyons que, dans cette situation, $u(U)$ est un
ensemble fini \`a $n$ \'el\'ements : en effet, le membre de gauche a pour fibre $u(U)$ au-dessus de
tout \'el\'ement de $u(V)$ et le membre de droite a une fibre de cardinal $n$.

\medskip\noindent
Consid\'erons la limite
$$
\lim_{(V, \overline{v})} u(V)
$$
sur la cat\'egorie des voisinages \'etales $(V, \overline{v})$ de
$\overline{s}$. Il est clair que nous obtenons la m\^eme valeur en prenant
la limite sur la sous-cat\'egorie des $(V, \overline{v})$ tels que $V$ soit affine.
D'apr\`es le paragraphe pr\'ec\'edent (appliqu\'e en \'echangeant les r\^oles de $V$ et $U$),
nous voyons que, dans ce cas, $u(V)$ est toujours un ensemble fini non vide.
De plus, la cat\'egorie d'indexation est cofiltrante, voir le
lemme \ref{etale-cohomology-lemma-cofinal-etale}.
Par cons\'equent, d'apr\`es
Catégories, section \ref{categories-section-codirected-limits},
la limite est non vide. Choisissons un \'el\'ement $x$ de cette limite.
Cela signifie que nous obtenons un $x_{V, \overline{v}} \in u(V)$ pour
chaque voisinage \'etale $(V, \overline{v})$ de $(S, \overline{s})$
tel que, pour tout morphisme de voisinages \'etales
$\varphi : (V', \overline{v}') \to (V, \overline{v})$, on ait
$u(\varphi)(x_{V', \overline{v}'}) = x_{V, \overline{v}}$.

\medskip\noindent
Nous utiliserons le choix de $x$ pour construire une application bijective fonctorielle
$$
c : |U_{\overline{s}}| \longrightarrow u(U)
$$
pour $U \in \Ob(S_\etale)$, ce qui conclura la preuve. Voir le
lemme \ref{etale-cohomology-lemma-stalk-gives-point}
et sa preuve pour une description de $|U_{\overline{s}}|$.
Nous affirmons d'abord qu'il suffit de construire l'application pour $U$ affine.
Nous omettons la preuve de cette affirmation.
Supposons que $U \to S$ soit un objet de $S_\etale$, avec $U$ affine, et soit
$\overline{u} : \overline{s} \to U$ un \'el\'ement de $|U_{\overline{s}}|$.
Choisissons $(V, \overline{v})$ tel que $U \times_S V$ se d\'ecompose
comme dans le troisi\`eme paragraphe de la preuve.
Alors le couple $(\overline{u}, \overline{v})$ donne un point g\'eom\'etrique de
$U \times_S V$ au-dessus de $\overline{v}$ et d\'etermine l'une des
composantes $U_i$ de $U \times_S V$. Plus pr\'ecis\'ement, il existe
une section $\sigma : V \to U \times_S V$ de la projection $\text{pr}_U$
telle que $(\overline{u}, \overline{v}) = \sigma \circ \overline{v}$. Posons
$c(\overline{u}) = u(\text{pr}_U)(u(\sigma)(x_{V, \overline{v}})) \in u(U)$.
Nous devons v\'erifier que cela est ind\'ependant du choix de $(V, \overline{v})$. D'apr\`es le
lemme \ref{etale-cohomology-lemma-cofinal-etale},
la cat\'egorie des voisinages \'etales est cofiltrante.
Il suffit donc de montrer
que, \'etant donn\'es un morphisme de voisinages \'etales
$\varphi : (V', \overline{v}') \to (V, \overline{v})$ et le choix d'une
section $\sigma' : V' \to U \times_S V'$ de la projection telle que
$(\overline{u}, \overline{v'}) = \sigma' \circ \overline{v}'$,
on a $u(\sigma')(x_{V', \overline{v}'}) = u(\sigma)(x_{V, \overline{v}})$.
Consid\'erons le diagramme
$$
\xymatrix{
V' \ar[d]^{\sigma'} \ar[r]_\varphi & V \ar[d]^\sigma \\
U \times_S V' \ar[r]^{1 \times \varphi} &
U \times_S V
}
$$
Il se peut que ce diagramme ne soit pas commutatif. La raison en est
que les sch\'emas $V'$ et $V$ ne sont peut-\^etre pas connexes, et donc que
les d\'ecompositions utilis\'ees ci-dessus pour construire $\sigma'$ et $\sigma$ ne sont
peut-\^etre pas uniques. Mais nous savons que
$\sigma \circ \varphi \circ \overline{v}' =
(1 \times \varphi) \circ \sigma' \circ \overline{v}'$
par construction. Par cons\'equent, puisque $U \times_S V$ est \'etale sur $S$,
il existe un voisinage ouvert
$V'' \subset V'$ de $\overline{v'}$ tel que le diagramme
soit commutatif apr\`es restriction \`a $V''$, voir
Morphismes, lemme \ref{morphisms-lemma-value-at-one-point}.
Cela signifie que nous pouvons compl\'eter le diagramme ci-dessus en
$$
\xymatrix{
V'' \ar[r] \ar[d]^{\sigma'|_{V''}} &
V' \ar[d]^{\sigma'} \ar[r]_\varphi &
V \ar[d]^\sigma \\
U \times_S V'' \ar[r] &
U \times_S V' \ar[r]^{1 \times \varphi} &
U \times_S V
}
$$
de sorte que le carr\'e de gauche et le rectangle ext\'erieur soient commutatifs.
Comme $u$ est un foncteur, cela implique que
$x_{V'', \overline{v}'}$ s'envoie sur le m\^eme \'el\'ement de
$u(U \times_S V)$ quel que soit le chemin suivi dans le
diagramme. D'autre part, il s'envoie sur les \'el\'ements
$x_{V', \overline{v}'}$ et $x_{V, \overline{v}}$ dans
$u(V')$ et $u(V)$. Cela implique l'\'egalit\'e voulue
$u(\sigma')(x_{V', \overline{v}'}) = u(\sigma)(x_{V, \overline{v}})$.

\medskip\noindent
De mani\`ere analogue, on d\'emontre que la construction
$c : |U_{\overline{s}}| \to u(U)$ est fonctorielle en $U$ ;
nous omettons les d\'etails. Enfin, d'apr\`es les r\'esultats du
troisi\`eme paragraphe, il est clair que l'application $c$ est bijective,
ce qui termine la preuve du lemme.
\end{proof}







\section{Points dans d'autres topologies}
\label{etale-cohomology-section-points-topologies}

\noindent
Dans cette section, nous discutons bri\`evement de l'existence de points pour certains
sites autres que le site \'etale d'un sch\'ema. Nous renvoyons \`a
Sites, section \ref{sites-section-sites-enough-points}
et \`a
Topologies, section \ref{topologies-section-procedure} et suivantes
pour la terminologie employ\'ee dans cette section.
Tous les sites g\'eom\'etriques ont assez de points.

\begin{lemma}
\label{etale-cohomology-lemma-points-fppf}
Soit $S$ un sch\'ema. Tous les sites suivants ont assez de points :
$S_{affine, Zar}$, $S_{Zar}$, $S_{affine, \etale}$, $S_\etale$,
$(\Sch/S)_{Zar}$, $(\textit{Aff}/S)_{Zar}$,
$(\Sch/S)_\etale$, $(\textit{Aff}/S)_\etale$,
$(\Sch/S)_{smooth}$, $(\textit{Aff}/S)_{smooth}$,
$(\Sch/S)_{syntomic}$, $(\textit{Aff}/S)_{syntomic}$,
$(\Sch/S)_{fppf}$ et $(\textit{Aff}/S)_{fppf}$.
\end{lemma}

\begin{proof}
Pour chacun des grands sites, le topos associ\'e est \'equivalent au
topos d\'efini par le site $(\textit{Aff}/S)_\tau$, voir
Topologies, lemmes \ref{topologies-lemma-affine-big-site-Zariski},
\ref{topologies-lemma-affine-big-site-etale},
\ref{topologies-lemma-affine-big-site-smooth},
\ref{topologies-lemma-affine-big-site-syntomic} et
\ref{topologies-lemma-affine-big-site-fppf}.
Le r\'esultat pour les sites $(\textit{Aff}/S)_\tau$ d\'ecoule imm\'ediatement
du r\'esultat de Deligne,
Sites, lemme \ref{sites-lemma-criterion-points}.

\medskip\noindent
Le r\'esultat pour $S_{Zar}$ est clair. Celui pour $S_{affine, Zar}$
d\'ecoule du r\'esultat de Deligne. Le r\'esultat pour $S_\etale$
d\'ecoule soit (de la preuve) du
th\'eor\`eme \ref{etale-cohomology-theorem-exactness-stalks},
soit de Topologies, lemme \ref{topologies-lemma-alternative}
et du r\'esultat de Deligne appliqu\'e \`a $S_{affine, \etale}$.
\end{proof}

\noindent
Le lemme ci-dessus garantit l'existence de points, mais ne nous
dit pas \`a quoi ils ressemblent. Nous pouvons construire explicitement
{\it certains} points comme suit.
Supposons que $\overline{s} : \Spec(k) \to S$ soit un point
g\'eom\'etrique, avec $k$ alg\'ebriquement clos. Consid\'erons le foncteur
$$
u : (\Sch/S)_{fppf} \longrightarrow \textit{Ensembles},
\quad
u(U) = U(k) = \Mor_S(\Spec(k), U).
$$
Notons que $U \mapsto U(k)$ commute aux limites finies puisque
$S(k) = \{\overline{s}\}$ et
$(U_1 \times_U U_2)(k) = U_1(k) \times_{U(k)} U_2(k)$.
De plus, si $\{U_i \to U\}$ est un recouvrement fppf, alors
$\coprod U_i(k) \to U(k)$ est surjectif.
D'apr\`es
Sites, Proposition \ref{sites-proposition-point-limits},
nous voyons que $u$ d\'efinit un point $p$ de $(\Sch/S)_{fppf}$ dont les
fibres sont
$$
\mathcal{F}_p = \colim_{(U, x)} \mathcal{F}(U)
$$
o\`u la limite inductive porte, comme d'habitude, sur les couples $U \to S$, $x \in U(k)$.
Mais... cette cat\'egorie a un objet initial, \`a savoir
$(\Spec(k), \text{id})$ ; nous voyons donc que
$$
\mathcal{F}_p = \mathcal{F}(\Spec(k))
$$
ce qui n'est gu\`ere int\'eressant ! En fait, en g\'en\'eral, ces points ne
forment pas une famille conservative de points. Un type de point plus int\'eressant
est d\'ecrit dans la remarque suivante.

\begin{remark}
\label{etale-cohomology-remark-points-fppf-site}
\begin{reference}
Ceci est discut\'e dans \cite{Schroeer}.
\end{reference}
Soit $S = \Spec(A)$ un sch\'ema affine. Soit $(p, u)$ un point du
site $(\textit{Aff}/S)_{fppf}$, voir
Sites, sections \ref{sites-section-points} et
\ref{sites-section-construct-points}. Soit $B = \mathcal{O}_p$ la fibre
du faisceau structural au point $p$. Rappelons que
$$
B = \colim_{(U, x)} \mathcal{O}(U) =
\colim_{(\Spec(C), x_C)} C
$$
o\`u $x_C \in u(\Spec(C))$. Il peut arriver que
$\Spec(B)$ soit un objet de $(\textit{Aff}/S)_{fppf}$
et qu'il existe un \'el\'ement $x_B \in u(\Spec(B))$ qui s'envoie sur
le syst\`eme compatible $x_C$. Dans ce cas, le syst\`eme de voisinages
a un objet initial et il s'ensuit que
$\mathcal{F}_p = \mathcal{F}(\Spec(B))$ pour tout faisceau $\mathcal{F}$
sur $(\textit{Aff}/S)_{fppf}$. Il est imm\'ediat
de voir que, si $\mathcal{F} \mapsto \mathcal{F}(\Spec(B))$ d\'efinit un point
de $\Sh((\textit{Aff}/S)_{fppf})$, alors
$B$ doit \^etre une $A$-alg\`ebre locale telle que, pour tout morphisme d'anneaux
fid\`element plat et de pr\'esentation finie $B \to B'$, il existe une section $B' \to B$.
R\'eciproquement, pour toute $A$-alg\`ebre $B$ de ce type, le foncteur
$\mathcal{F} \mapsto \mathcal{F}(\Spec(B))$ est le foncteur fibre
d'un point. Nous omettons les d\'etails. On ne sait pas clairement \`a quoi ressemble un point quelconque du
site $(\textit{Aff}/S)_{fppf}$.
\end{remark}











\section{Supports des faisceaux ab\'eliens}
\label{etale-cohomology-section-support}

\noindent
Nous commen\c{c}ons par les supports des sections locales.

\begin{lemma}
\label{etale-cohomology-lemma-support-subsheaf-final}
Soit $S$ un sch\'ema. Soit $\mathcal{F}$ un sous-faisceau de l'objet
final du topos \'etale de $S$ (voir
Sites, exemple \ref{sites-example-singleton-sheaf}).
Alors il existe un unique ouvert
$W \subset S$ tel que $\mathcal{F} = h_W$.
\end{lemma}

\begin{proof}
La condition signifie que $\mathcal{F}(U)$ est un singleton ou est
vide pour tout $\varphi : U \to S$ dans $\Ob(S_\etale)$.
En particulier, les sections locales se recollent toujours. Si
$\mathcal{F}(U) \not = \emptyset$, alors
$\mathcal{F}(\varphi(U)) \not = \emptyset$ parce que
$\{\varphi : U \to \varphi(U)\}$ est un recouvrement.
Nous pouvons donc prendre
$W = \bigcup_{\varphi : U \to S, \mathcal{F}(U) \not = \emptyset} \varphi(U)$.
\end{proof}

\begin{lemma}
\label{etale-cohomology-lemma-zero-over-image}
Soit $S$ un sch\'ema.
Soit $\mathcal{F}$ un faisceau ab\'elien sur $S_\etale$.
Soit $\sigma \in \mathcal{F}(U)$ une section locale.
Il existe un ouvert $W \subset U$ tel que
\begin{enumerate}
\item $W \subset U$ soit le plus grand ouvert de Zariski de $U$ tel
que $\sigma|_W = 0$,
\item pour tout $\varphi : V \to U$ dans $S_\etale$, on ait
$$
\sigma|_V = 0 \Leftrightarrow \varphi(V) \subset W,
$$
\item pour tout point g\'eom\'etrique $\overline{u}$ de $U$, on ait
$$
(U, \overline{u}, \sigma) = 0\text{ dans }\mathcal{F}_{\overline{s}}
\Leftrightarrow
\overline{u} \in W
$$
o\`u $\overline{s} = (U \to S) \circ \overline{u}$.
\end{enumerate}
\end{lemma}

\begin{proof}
Comme $\mathcal{F}$ est un faisceau pour la topologie \'etale, la restriction de
$\mathcal{F}$ \`a $U_{Zar}$ est un faisceau sur $U$ pour la topologie de Zariski.
Il existe donc un ouvert de Zariski $W$ ayant la propri\'et\'e (1), voir
Modules, lemme \ref{modules-lemma-support-section-closed}. Soit
$\varphi : V \to U$ un morphisme de $S_\etale$. Notons que
$\varphi(V) \subset U$ est un ouvert et que
$\{V \to \varphi(V)\}$ est un recouvrement \'etale. Ainsi, si
$\sigma|_V = 0$, alors la condition de faisceau pour $\mathcal{F}$ montre
que $\sigma|_{\varphi(V)} = 0$. Cela d\'emontre (2).
Pour d\'emontrer (3), nous devons montrer que si $(U, \overline{u}, \sigma)$
d\'efinit l'\'el\'ement nul de $\mathcal{F}_{\overline{s}}$, alors
$\overline{u} \in W$. Cela est vrai parce que l'hypoth\`ese signifie
qu'il existe un morphisme de voisinages \'etales
$(V, \overline{v}) \to (U, \overline{u})$ tel que
$\sigma|_V = 0$. D'apr\`es (2), nous voyons donc que $V \to U$ est \`a valeurs dans $W$, et
par cons\'equent que $\overline{u} \in W$.
\end{proof}

\noindent
Soit $S$ un sch\'ema. Soit $s \in S$.
Soit $\mathcal{F}$ un faisceau sur $S_\etale$. D'apr\`es la
remarque \ref{etale-cohomology-remark-map-stalks},
la classe d'isomorphisme de la fibre du faisceau $\mathcal{F}$
en un point g\'eom\'etrique au-dessus de $s$ est bien d\'efinie.

\begin{definition}
\label{etale-cohomology-definition-support}
Soit $S$ un sch\'ema.
Soit $\mathcal{F}$ un faisceau ab\'elien sur $S_\etale$.
\begin{enumerate}
\item Le {\it support de $\mathcal{F}$} est l'ensemble des
points $s \in S$ tels que $\mathcal{F}_{\overline{s}} \not = 0$
pour tout (ou un) point g\'eom\'etrique $\overline{s}$ au-dessus de $s$.
\item Soit $\sigma \in \mathcal{F}(U)$ une section.
Le {\it support de $\sigma$} est le ferm\'e $U \setminus W$, o\`u
$W \subset U$ est le plus grand ouvert de $U$ sur lequel $\sigma$
est nulle (voir le
lemme \ref{etale-cohomology-lemma-zero-over-image}).
\end{enumerate}
\end{definition}

\noindent
En g\'en\'eral, le support d'un faisceau ab\'elien n'est pas ferm\'e.
Par exemple, supposons que $S = \mathbf{A}^1_{\mathbf{C}}$.
Soit $i_t : \Spec(\mathbf{C}) \to S$ l'inclusion du
point $t \in \mathbf{C}$.
Nous verrons plus loin que $\mathcal{F}_t = i_{t, *}(\mathbf{Z}/2\mathbf{Z})$
est un faisceau ab\'elien dont le support est exactement $\{t\}$, voir la
section \ref{etale-cohomology-section-closed-immersions}.
Alors
$$
\bigoplus\nolimits_{n \in \mathbf{N}} \mathcal{F}_n
$$
est un faisceau ab\'elien de support $\{1, 2, 3, \ldots\} \subset S$.
Cela est vrai parce que prendre les fibres commute aux limites inductives, voir le
lemme \ref{etale-cohomology-lemma-stalk-exact}.
Nous avons ainsi un exemple de faisceau ab\'elien dont le support n'est pas ferm\'e.
Voici quelques faits fondamentaux sur les supports des faisceaux et des sections.

\begin{lemma}
\label{etale-cohomology-lemma-support-section-closed}
Soit $S$ un sch\'ema.
Soit $\mathcal{F}$ un faisceau ab\'elien sur $S_\etale$.
Soient $U \in \Ob(S_\etale)$ et $\sigma \in \mathcal{F}(U)$.
\begin{enumerate}
\item Le support de $\sigma$ est ferm\'e dans $U$.
\item Le support de $\sigma + \sigma'$ est contenu dans la r\'eunion des
supports de $\sigma, \sigma' \in \mathcal{F}(U)$.
\item Si $\varphi : \mathcal{F} \to \mathcal{G}$ est un morphisme de
faisceaux ab\'eliens sur $S_\etale$, alors le support de $\varphi(\sigma)$
est contenu dans le support de $\sigma \in \mathcal{F}(U)$.
\item Le support de $\mathcal{F}$ est la r\'eunion des images des
supports de toutes les sections locales de $\mathcal{F}$.
\item Si $\mathcal{F} \to \mathcal{G}$ est surjectif, alors le support
de $\mathcal{G}$ est un sous-ensemble du support de $\mathcal{F}$.
\item Si $\mathcal{F} \to \mathcal{G}$ est injectif, alors le support
de $\mathcal{F}$ est un sous-ensemble du support de $\mathcal{G}$.
\end{enumerate}
\end{lemma}

\begin{proof}
L'assertion (1) est vraie par d\'efinition.
Les assertions (2) et (3) sont vraies parce qu'elles le sont pour les restrictions de
$\mathcal{F}$ et $\mathcal{G}$ \`a $U_{Zar}$, voir
Modules, lemme \ref{modules-lemma-support-section-closed}.
L'assertion (4) est une cons\'equence directe de la partie (3) du
lemme \ref{etale-cohomology-lemma-zero-over-image}.
Les assertions (5) et (6) d\'ecoulent des autres.
\end{proof}

\begin{lemma}
\label{etale-cohomology-lemma-support-sheaf-rings-closed}
Le support d'un faisceau d'anneaux sur $S_\etale$ est ferm\'e.
\end{lemma}

\begin{proof}
Cela est vrai parce que (selon nos conventions)
un anneau est $0$ si et seulement si
$1 = 0$ ; par cons\'equent, le support d'un faisceau d'anneaux
est le support de la section unit\'e.
\end{proof}




\section{Anneaux hens\'eliens}
\label{etale-cohomology-section-henselian-ring}

\noindent
Nous commen\c{c}ons par \'enoncer un th\'eor\`eme qui a d\'ej\`a \'et\'e utilis\'e de nombreuses fois
dans le projet Champs. Il existe de nombreuses versions de ce r\'esultat ; ici, nous
n'\'enon\c{c}ons que la version alg\'ebrique.

\begin{theorem}
\label{etale-cohomology-theorem-quasi-finite-etale-locally}
Soient $A\to B$ un morphisme d'anneaux de type fini et $\mathfrak p \subset A$ un
id\'eal premier. Alors il existe un morphisme d'anneaux \'etale $A \to A'$ et un id\'eal premier
$\mathfrak p' \subset A'$ au-dessus de $\mathfrak p$ tels que
\begin{enumerate}
\item
$\kappa(\mathfrak p) = \kappa(\mathfrak p')$,
\item
$ B \otimes_A A' = B_1\times \ldots \times B_r \times C$,
\item
$ A'\to B_i$ est fini et il existe un unique id\'eal premier $q_i\subset B_i$ au-dessus
de $\mathfrak p'$, et
\item toutes les composantes irr\'eductibles de la fibre
$\Spec(C \otimes_{A'} \kappa(\mathfrak p'))$ de $C$ au-dessus de $\mathfrak p'$
sont de dimension au moins 1.
\end{enumerate}
\end{theorem}

\begin{proof}
Voir Algèbre, lemme \ref{algebra-lemma-etale-makes-quasi-finite-finite}, ou
voir \cite[Th\'eor\`eme 18.12.1]{EGA4}. Pour de nombreuses versions en termes de
morphismes de sch\'emas, voir
Compléments sur les morphismes, section \ref{more-morphisms-section-etale-localization}.
\end{proof}

\noindent
Rappelons le lemme de Hensel.
Il existe de nombreuses versions de ce lemme. En voici deux :
\begin{enumerate}
\item[(f)] si $f\in \mathbf{Z}_p[T]$ est unitaire et
$f \bmod p = g_0 h_0$ avec $gcd(g_0, h_0) = 1$, alors $f$ se factorise
sous la forme $f = gh$ avec $\bar g = g_0$ et $\bar h = h_0$,
\item[(r)] si $f \in \mathbf{Z}_p[T]$ est unitaire, $a_0 \in \mathbf{F}_p$,
$\bar f(a_0) =0$ mais $\bar f'(a_0) \neq 0$,
alors il existe $a \in \mathbf{Z}_p$ tel que
$f(a) = 0$ et $\bar a = a_0$.
\end{enumerate}
Les deux versions sont vraies (nous le verrons plus loin). La premi\`ere version
demande des rel\`evements de factorisations en facteurs premiers entre eux,
et la seconde demande des rel\`evements de racines simples
modulo l'id\'eal maximal. Il se trouve que, pour un anneau local g\'en\'eral,
ces deux conditions sont \'equivalentes, et qu'elles sont
\'equivalentes \`a beaucoup d'autres conditions. Nous utilisons la propri\'et\'e de rel\`evement
des racines comme d\'efinition d'un anneau local hens\'elien, car c'est
souvent la plus facile \`a v\'erifier.

%10.01.09
\begin{definition}
\label{etale-cohomology-definition-henselian}
(Voir Algèbre, définition \ref{algebra-definition-henselian}.)
Un anneau local $(R, \mathfrak m, \kappa)$ est dit
{\it hens\'elien} si, pour tout
$f \in R[T]$ unitaire et tout $a_0 \in \kappa$ tel que
$\bar f(a_0) = 0$ et $\bar f'(a_0) \neq 0$, il existe
un $a \in R$ tel que $f(a) = 0$ et $a \bmod \mathfrak m = a_0$.
\end{definition}

\noindent
Les anneaux locaux complets constituent un bon exemple
d'anneaux locaux hens\'eliens \`a garder \`a l'esprit.
Rappelons
(Algèbre, définition \ref{algebra-definition-complete-local-ring})
qu'un anneau local complet est un anneau local $(R, \mathfrak m)$ tel que
$R \cong \lim_n R/\mathfrak m^n$, c'est-\`a-dire complet et s\'epar\'e
pour la topologie $\mathfrak m$-adique.

\begin{theorem}
\label{etale-cohomology-theorem-hensel}
Les anneaux locaux complets sont hens\'eliens.
\end{theorem}

\begin{proof}
M\'ethode de Newton. Voir
Algèbre, lemme \ref{algebra-lemma-complete-henselian}.
\end{proof}

\begin{theorem}
\label{etale-cohomology-theorem-henselian}
Soit $(R, \mathfrak m, \kappa)$ un anneau local. Les propri\'et\'es suivantes sont \'equivalentes :
\begin{enumerate}
\item $R$ est hens\'elien,
\item pour tout $f\in R[T]$ et toute factorisation $\bar f = g_0 h_0$ dans
$\kappa[T]$ avec $\gcd(g_0, h_0)=1$, il existe une factorisation $f = gh$ dans
$R[T]$ avec $\bar g = g_0$ et $\bar h = h_0$,
\item toute $R$-alg\`ebre finie $S$ est isomorphe \`a un produit fini
d'anneaux locaux finis sur $R$,
\item toute $R$-alg\`ebre de type fini $A$ est isomorphe \`a un produit
$A \cong A' \times C$ o\`u $A' \cong A_1 \times \ldots \times A_r$
est un produit de $R$-alg\`ebres locales finies et toutes les composantes
irr\'eductibles de $C \otimes_R \kappa$ sont de dimension au moins 1,
\item si $A$ est une $R$-alg\`ebre \'etale et $\mathfrak n$ un id\'eal maximal de
$A$ au-dessus de $\mathfrak m$ tel que $\kappa \cong A/\mathfrak n$, alors il
existe un isomorphisme $\varphi : A \cong R \times A'$ tel que
$\varphi(\mathfrak n) = \mathfrak m \times A' \subset R \times A'$.
\end{enumerate}
\end{theorem}

\begin{proof}
Il s'agit simplement d'une partie des r\'esultats de
Algèbre, lemme \ref{algebra-lemma-characterize-henselian}.
Notons que la partie (5) ci-dessus correspond \`a la partie (8) de
Algèbre, lemme \ref{algebra-lemma-characterize-henselian},
mais qu'elle est formul\'ee l\'eg\`erement diff\'eremment.
\end{proof}

\begin{lemma}
\label{etale-cohomology-lemma-finite-over-henselian}
Si $R$ est hens\'elien et $A$ est une $R$-alg\`ebre finie, alors $A$ est un produit
fini d'anneaux locaux hens\'eliens.
\end{lemma}

\begin{proof}
Voir
Algèbre, lemme \ref{algebra-lemma-finite-over-henselian}.
\end{proof}

\begin{definition}
\label{etale-cohomology-definition-strictly-henselian}
Un anneau local $R$ est dit {\it strictement hens\'elien} s'il est hens\'elien et si son
corps r\'esiduel est s\'eparablement clos.
\end{definition}

\begin{example}
\label{etale-cohomology-example-powerseries}
Dans le cas $R = \mathbf{C}[[t]]$, les $R$-alg\`ebres \'etales sont des produits finis
de l'extension triviale $R \to R$ et des extensions
$R \to R[X, X^{-1}]/(X^n-t)$.
Ces derni\`eres se factorisent par l'ouvert $D(t) \subset \Spec(R)$, de sorte que tout
recouvrement \'etale peut \^etre raffin\'e par le recouvrement
$\{\text{id} : \Spec(R) \to \Spec(R)\}$. Nous verrons ci-dessous que
c'est un fait assez g\'en\'eral sur les recouvrements \'etales des spectres d'anneaux
hens\'eliens. Cela montrera que la cohomologie \'etale sup\'erieure du spectre d'un anneau
strictement hens\'elien est nulle.
\end{example}

\begin{theorem}
\label{etale-cohomology-theorem-henselization}
Soient $(R, \mathfrak m, \kappa)$ un anneau local et
$\kappa\subset\kappa^{sep}$ une cl\^oture alg\'ebrique s\'eparable.
Il existe des morphismes locaux plats canoniques d'anneaux $R \to R^h \to R^{sh}$ tels que
\begin{enumerate}
\item $R^h$, $R^{sh}$ soient des limites inductives filtrantes de $R$-alg\`ebres \'etales,
\item $R^h$ soit hens\'elien, $R^{sh}$ soit strictement hens\'elien,
\item $\mathfrak m R^h$ (resp.\ $\mathfrak m R^{sh}$) soit
l'id\'eal maximal de $R^h$ (resp.\ de $R^{sh}$), et
\item $\kappa = R^h/\mathfrak m R^h$, et
$\kappa^{sep} = R^{sh}/\mathfrak m R^{sh}$ comme extensions de $\kappa$.
\end{enumerate}
\end{theorem}

\begin{proof}
La structure de $R^h$ et de $R^{sh}$ est d\'ecrite dans
Algèbre, lemmes \ref{algebra-lemma-henselization} et
\ref{algebra-lemma-strict-henselization}.
\end{proof}

\noindent
Les anneaux construits dans le th\'eor\`eme \ref{etale-cohomology-theorem-henselization}
sont appel\'es respectivement {\it l'hens\'elisation} et
{\it l'hens\'elisation stricte} de l'anneau local $R$, voir
Algèbre, définition \ref{algebra-definition-henselization}.
De nombreuses propri\'et\'es de $R$ se refl\`etent dans son hens\'elisation (stricte),
voir Compléments d'algèbre,
section \ref{more-algebra-section-permanence-henselization}.




\section{Fibres du faisceau structural}
\label{etale-cohomology-section-stalks-structure-sheaf}

\noindent
Dans cette section, nous identifions la fibre du faisceau structural en un point
g\'eom\'etrique \`a l'hens\'elisation stricte de l'anneau local au point
``usuel'' correspondant.

\begin{lemma}
\label{etale-cohomology-lemma-describe-etale-local-ring}
\begin{slogan}
La fibre du faisceau structural d'un sch\'ema
pour la topologie \'etale est l'hens\'elisation stricte.
\end{slogan}
Soit $S$ un sch\'ema.
Soit $\overline{s}$ un point g\'eom\'etrique de $S$ au-dessus de $s \in S$.
Posons $\kappa = \kappa(s)$ et notons
$\kappa \subset \kappa^{sep} \subset \kappa(\overline{s})$
la cl\^oture alg\'ebrique s\'eparable de $\kappa$ dans $\kappa(\overline{s})$.
Alors il existe une identification canonique
$$
(\mathcal{O}_{S, s})^{sh}
\cong
(\mathcal{O}_S)_{\overline{s}}
$$
o\`u le membre de gauche est l'hens\'elisation stricte de l'anneau local
$\mathcal{O}_{S, s}$ telle qu'elle est d\'ecrite dans le
th\'eor\`eme \ref{etale-cohomology-theorem-henselization},
et le membre de droite est la fibre du faisceau structural
$\mathcal{O}_S$ sur $S_\etale$ au
point g\'eom\'etrique $\overline{s}$.
\end{lemma}

\begin{proof}
Soit $\Spec(A) \subset S$ un voisinage affine de $s$.
Soit $\mathfrak p \subset A$ l'id\'eal premier correspondant \`a $s$.
Avec ces choix, nous avons des isomorphismes canoniques
$\mathcal{O}_{S, s} = A_{\mathfrak p}$ et $\kappa(s) = \kappa(\mathfrak p)$.
Nous avons donc
$\kappa(\mathfrak p) \subset \kappa^{sep} \subset \kappa(\overline{s})$.
Rappelons que
$$
(\mathcal{O}_S)_{\overline{s}} =
\colim_{(U, \overline{u})} \mathcal{O}(U)
$$
o\`u la limite inductive porte sur les voisinages \'etales de $(S, \overline{s})$.
Un syst\`eme cofinal est donn\'e par les voisinages \'etales $(U, \overline{u})$
tels que $U$ soit affine et que $U \to S$ se factorise par $\Spec(A)$.
Autrement dit, nous voyons que
$$
(\mathcal{O}_S)_{\overline{s}} = \colim_{(B, \mathfrak q, \phi)} B
$$
o\`u la limite inductive porte sur les $A$-alg\`ebres \'etales $B$ munies d'un id\'eal premier
$\mathfrak q$ au-dessus de $\mathfrak p$ et d'un
morphisme de $\kappa(\mathfrak p)$-alg\`ebres
$\phi : \kappa(\mathfrak q) \to \kappa(\overline{s})$.
Notons que, puisque $\kappa(\mathfrak q)$ est une extension finie s\'eparable de
$\kappa(\mathfrak p)$, l'image de $\phi$ est contenue dans $\kappa^{sep}$.
Moyennant ces identifications, le r\'esultat du lemme est \'equivalent
au r\'esultat de
Algèbre, lemme \ref{algebra-lemma-strict-henselization-different}.
\end{proof}

\begin{definition}
\label{etale-cohomology-definition-etale-local-rings}
Soit $S$ un sch\'ema. Soit $\overline{s}$ un point g\'eom\'etrique de $S$
au-dessus du point $s \in S$.
\begin{enumerate}
\item L'{\it anneau local \'etale de $S$ en $\overline{s}$}
est la fibre du faisceau structural $\mathcal{O}_S$ sur $S_\etale$
en $\overline{s}$. Nous appelons parfois cet anneau
{\it l'hens\'elisation stricte de $\mathcal{O}_{S, s}$} relative
au point g\'eom\'etrique $\overline{s}$.
Notation utilis\'ee : $\mathcal{O}_{S, \overline{s}}^{sh}$.
\item L'{\it hens\'elisation de $\mathcal{O}_{S, s}$} d\'esigne
l'hens\'elisation de l'anneau local de $S$ en $s$. Voir
Algèbre, définition \ref{algebra-definition-henselization},
et le
th\'eor\`eme \ref{etale-cohomology-theorem-henselization}.
Notation : $\mathcal{O}_{S, s}^h$.
\item L'{\it hens\'elisation stricte de $S$ en $\overline{s}$}
est le sch\'ema $\Spec(\mathcal{O}_{S, \overline{s}}^{sh})$.
\item L'{\it hens\'elisation de $S$ en $s$} est le sch\'ema
$\Spec(\mathcal{O}_{S, s}^h)$.
\end{enumerate}
\end{definition}

\noindent
Soit $f : T \to S$ un morphisme de sch\'emas. Soit $\overline{t}$
un point g\'eom\'etrique de $T$ d'image $\overline{s}$ dans $S$.
Soient $t \in T$ et $s \in S$ leurs images.
Nous obtenons alors un diagramme commutatif canonique
$$
\xymatrix{
\Spec(\mathcal{O}^{sh}_{T, \overline{t}}) \ar[r] \ar[d] &
\Spec(\mathcal{O}^h_{T, t}) \ar[r] \ar[d] &
T \ar[d]^f \\
\Spec(\mathcal{O}^{sh}_{S, \overline{s}}) \ar[r] &
\Spec(\mathcal{O}^h_{S, s}) \ar[r] &
S
}
$$
des hens\'elisations et hens\'elisations strictes de $T$ et de $S$. On peut
le d\'emontrer en choisissant des voisinages affines de $t$ et de $s$ et en utilisant
la fonctorialit\'e des hens\'elisations (strictes) donn\'ee par
Algèbre, lemmes \ref{algebra-lemma-henselian-functorial-improve} et
\ref{algebra-lemma-strictly-henselian-functorial-improve}.

\begin{lemma}
\label{etale-cohomology-lemma-describe-henselization}
Soit $S$ un sch\'ema. Soit $s \in S$. Alors
$$
\mathcal{O}_{S, s}^h =
\colim_{(U, u)} \mathcal{O}(U)
$$
o\`u la limite inductive porte sur la cat\'egorie filtrante des
voisinages \'etales $(U, u)$ de $(S, s)$ tels que
$\kappa(s) = \kappa(u)$.
\end{lemma}

\begin{proof}
Ce lemme est une copie de
Compléments sur les morphismes, lemme \ref{more-morphisms-lemma-describe-henselization}.
\end{proof}

\begin{remark}
\label{etale-cohomology-remark-henselization-Noetherian}
Soit $S$ un sch\'ema. Soit $s \in S$.
Si $S$ est localement noeth\'erien, alors $\mathcal{O}_{S, s}^h$
est lui aussi noeth\'erien et il a le m\^eme compl\'et\'e :
$$
\widehat{\mathcal{O}_{S, s}} \cong \widehat{\mathcal{O}_{S, s}^h}.
$$
En particulier,
$\mathcal{O}_{S, s} \subset
\mathcal{O}_{S, s}^h \subset
\widehat{\mathcal{O}_{S, s}}$.
L'hens\'elisation de $\mathcal{O}_{S, s}$ est en g\'en\'eral beaucoup
plus petite que son compl\'et\'e et h\'erite de nombre de ses propri\'et\'es.
Par exemple, si $\mathcal{O}_{S, s}$ est r\'eduit, alors
$\mathcal{O}_{S, s}^h$ l'est aussi, mais ce n'est en g\'en\'eral pas vrai du compl\'et\'e.
Ins\'erer ici des r\'ef\'erences futures.
\end{remark}

\begin{lemma}
\label{etale-cohomology-lemma-etale-site-locally-ringed}
Soit $S$ un sch\'ema. Le petit site \'etale $S_\etale$, muni de
son faisceau structural $\mathcal{O}_S$, est un site localement annel\'e, voir
Modules sur les sites, définition \ref{sites-modules-definition-locally-ringed}.
\end{lemma}

\begin{proof}
Cela r\'esulte du fait que les fibres
$(\mathcal{O}_S)_{\overline{s}} = \mathcal{O}^{sh}_{S, \overline{s}}$ sont
des anneaux locaux, et du fait que $S_\etale$ a assez de points, voir
le lemme \ref{etale-cohomology-lemma-describe-etale-local-ring},
le th\'eor\`eme \ref{etale-cohomology-theorem-exactness-stalks},
et les
remarques \ref{etale-cohomology-remarks-enough-points}.
Voir
Modules sur les sites, lemmes \ref{sites-modules-lemma-locally-ringed-stalk} et
\ref{sites-modules-lemma-ringed-stalk-not-zero}
pour le fait que cela implique que le petit site \'etale est localement annel\'e.
\end{proof}






\section{Fonctorialit\'e du petit topos \'etale}
\label{etale-cohomology-section-functoriality}

\noindent
Jusqu'ici, nous n'avons pas abord\'e la fonctorialit\'e du site \'etale,
autrement dit ce qui se passe lorsque l'on se donne un morphisme de sch\'emas. Une
discussion formelle pr\'ecise se trouve dans
Topologies, section \ref{topologies-section-etale}.
Dans cette section et les suivantes, nous examinons bri\`evement ces questions,
plus pr\'ecis\'ement dans le cadre des petits sites \'etales.

\medskip\noindent
Soit $f : X \to Y$ un morphisme de sch\'emas. Nous obtenons un foncteur
\begin{equation}
\label{etale-cohomology-equation-functorial}
u : Y_\etale \longrightarrow X_\etale, \quad
V/Y \longmapsto X \times_Y V/X.
\end{equation}
Ce foncteur poss\`ede les propri\'et\'es importantes suivantes :
\begin{enumerate}
\item $u(\text{objet final}) = \text{objet final}$,
\item $u$ pr\'eserve les produits fibr\'es,
\item si $\{V_j \to V\}$ est un recouvrement dans $Y_\etale$, alors
$\{u(V_j) \to u(V)\}$ est un recouvrement dans $X_\etale$.
\end{enumerate}
Chacune d'elles se v\'erifie facilement (nous omettons la v\'erification). Par cons\'equent, nous obtenons ce
que l'on appelle un {\it morphisme de sites}
$$
f_{small} : X_\etale \longrightarrow Y_\etale,
$$
voir
Sites, définition \ref{sites-definition-morphism-sites}
et
Sites, Proposition \ref{sites-proposition-get-morphism}.
Il n'est pas n\'ecessaire de conna\^itre en d\'etail la notion abstraite
pour travailler avec les faisceaux \'etales et la cohomologie \'etale.
Il suffit habituellement de savoir qu'il existe des foncteurs
$f_{small, *}$ (image directe) et $f_{small}^{-1}$ (image inverse)
sur les faisceaux \'etales, et de conna\^itre quelques-unes de leurs propri\'et\'es simples.
Nous examinerons ces propri\'et\'es dans les sections suivantes, mais nous nous
r\'ef\'ererons parfois au cadre plus abstrait pour les d\'emonstrations, car
c'est souvent le cadre naturel pour les \`etablir.


\section{Images directes}
\label{etale-cohomology-section-direct-image}

\noindent
D\'efinissons l'image directe d'un pr\'efaisceau.

\begin{definition}
\label{etale-cohomology-definition-direct-image-presheaf}
Soit $f: X\to Y$ un morphisme de sch\'emas.
Soit $\mathcal{F} $ un pr\'efaisceau d'ensembles sur $X_\etale$.
L'{\it image directe}, ou {\it pouss\'ee en avant}, de $\mathcal{F}$
(par $f$) est
$$
f_*\mathcal{F} : Y_\etale^{opp} \longrightarrow \textit{Ensembles}, \quad
(V/Y) \longmapsto \mathcal{F}(X \times_Y V/X).
$$
Nous \`ecrivons parfois $f_* = f_{small, *}$ pour distinguer ce foncteur des autres
foncteurs d'image directe (comme l'image directe usuelle pour la topologie de Zariski ou $f_{big, *}$).
\end{definition}

\noindent
C'est bien un pr\'efaisceau \'etale, puisque le changement de base d'un morphisme
\'etale est encore \'etale. Une fa\c{c}on plus cat\'egorique de l'exprimer consiste \`a dire que
$f_*\mathcal{F}$ est le foncteur compos\'e $\mathcal{F} \circ u$
o\`u $u$ est d\'efini par l'\'equation (\ref{etale-cohomology-equation-functorial}). Cela montre
clairement que la construction d\'epend fonctoriellement du pr\'efaisceau
$\mathcal{F}$ et nous obtenons donc un foncteur
$$
f_* = f_{small, *} :
\textit{PSh}(X_\etale)
\longrightarrow
\textit{PSh}(Y_\etale)
$$
Notons que si $\mathcal{F}$ est un pr\'efaisceau de groupes ab\'eliens, alors
$f_*\mathcal{F}$ est aussi un pr\'efaisceau de groupes ab\'eliens et nous obtenons
$$
f_* = f_{small, *} :
\textit{PAb}(X_\etale)
\longrightarrow
\textit{PAb}(Y_\etale)
$$
comme pr\'ec\'edemment (c'est-\`a-dire d\'efini par exactement la m\^eme r\`egle).

\begin{remark}
\label{etale-cohomology-remark-direct-image-sheaf}
Nous affirmons que l'image directe d'un faisceau est un faisceau.
En effet, si $\{V_j \to V\}$ est un recouvrement \'etale dans $Y_\etale$,
alors $\{X \times_Y V_j \to X \times_Y V\}$ est un recouvrement \'etale dans
$X_\etale$. Par cons\'equent, la condition de faisceau pour $\mathcal{F}$ relativement
\`a $\{X \times_Y V_i \to X \times_Y V\}$
est \`equivalente \`a la condition de faisceau pour $f_*\mathcal{F}$ relativement \`a
$\{V_i \to V\}$. Ainsi, si $\mathcal{F}$ est un faisceau,
$f_*\mathcal{F}$ l'est aussi.
\end{remark}

\begin{definition}
\label{etale-cohomology-definition-direct-image-sheaf}
Soit $f: X\to Y$ un morphisme de sch\'emas.
Soit $\mathcal{F} $ un faisceau d'ensembles sur $X_\etale$.
L'{\it image directe}, ou {\it pouss\'ee en avant}, de $\mathcal{F}$
(par $f$) est
$$
f_*\mathcal{F} : Y_\etale^{opp} \longrightarrow \textit{Ensembles}, \quad
(V/Y) \longmapsto \mathcal{F}(X \times_Y V/X)
$$
qui est un faisceau d'apr\`es la
remarque \ref{etale-cohomology-remark-direct-image-sheaf}.
Nous \`ecrivons parfois $f_* = f_{small, *}$ pour distinguer ce foncteur des autres
foncteurs d'image directe (comme l'image directe usuelle pour la topologie de Zariski ou $f_{big, *}$).
\end{definition}

\noindent
La discussion pr\'ec\'edente s'applique mot pour mot et nous obtenons des foncteurs
$$
f_* = f_{small, *} :
\Sh(X_\etale)
\longrightarrow
\Sh(Y_\etale)
$$
et
$$
f_* = f_{small, *} :
\textit{Ab}(X_\etale)
\longrightarrow
\textit{Ab}(Y_\etale)
$$
\`a nouveau appel\'es {\it foncteurs d'image directe}.

\medskip\noindent
Le foncteur $f_*$ sur les faisceaux ab\'eliens est exact \`a gauche. (Voir
Homology, section \ref{homology-section-functors}
pour savoir ce que signifie qu'un foncteur entre cat\'egories ab\'eliennes est exact \`a gauche.)
En effet, si
$0 \to \mathcal{F}_1 \to \mathcal{F}_2 \to \mathcal{F}_3$
est exacte sur $X_\etale$, alors, pour tout
$U/X \in \Ob(X_\etale)$,
la suite de groupes ab\'eliens
$0 \to \mathcal{F}_1(U) \to \mathcal{F}_2(U) \to \mathcal{F}_3(U)$
est exacte. Par cons\'equent, pour tout $V/Y \in \Ob(Y_\etale)$,
la suite de groupes ab\'eliens
$0 \to f_*\mathcal{F}_1(V) \to f_*\mathcal{F}_2(V) \to f_*\mathcal{F}_3(V)$
est exacte, car c'est la suite pr\'ec\'edente avec $U = X \times_Y V$.

\begin{definition}
\label{etale-cohomology-definition-higher-direct-images}
Soit $f: X \to Y$ un morphisme de sch\'emas.
Les foncteurs d\'eriv\'es droits $\{R^pf_*\}_{p \geq 1}$ de
$f_* : \textit{Ab}(X_\etale) \to \textit{Ab}(Y_\etale)$
sont appel\'es {\it images directes sup\'erieures}.
\end{definition}

\noindent
Les images directes sup\'erieures et leurs variantes dans la cat\'egorie d\'eriv\'ee
sont examin\'ees plus en d\'etail dans (ins\'erer ici une r\'ef\'erence future).



\section{Image inverse}
\label{etale-cohomology-section-inverse-image}

\noindent
Dans cette section, nous examinons bri\`evement l'image inverse des faisceaux sur les petits
sites \'etales. Sa construction pr\'ecise se trouve dans
Topologies, section \ref{topologies-section-etale}.

\begin{definition}
\label{etale-cohomology-definition-inverse-image}
Soit $f: X\to Y$ un morphisme de sch\'emas. Les foncteurs d'{\it image inverse}, ou
de {\it tir\'e en arri\`ere}\footnote{Nous utilisons la notation $f^{-1}$ pour les images inverses de
faisceaux d'ensembles ou de groupes ab\'eliens, et nous r\'eservons $f^*$ aux
images inverses de faisceaux de modules par un morphisme de sites ou de topos annel\'es.}
sont les foncteurs
$$
f^{-1} = f_{small}^{-1} :
\Sh(Y_\etale)
\longrightarrow
\Sh(X_\etale)
$$
et
$$
f^{-1} = f_{small}^{-1} :
\textit{Ab}(Y_\etale)
\longrightarrow
\textit{Ab}(X_\etale)
$$
qui sont adjoints \`a gauche de $f_* = f_{small, *}$. Ainsi,
$f^{-1}$ est caract\'eris\'e par le fait que
$$
\Hom_{{\Sh(X_\etale)}} (f^{-1}\mathcal{G}, \mathcal{F})
=
\Hom_{\Sh(Y_\etale)} (\mathcal{G}, f_*\mathcal{F})
$$
fonctoriellement, pour tout $\mathcal{F} \in \Sh(X_\etale)$ et tout
$\mathcal{G} \in \Sh(Y_\etale)$. De m\^eme, nous avons
$$
\Hom_{{\textit{Ab}(X_\etale)}} (f^{-1}\mathcal{G}, \mathcal{F})
=
\Hom_{\textit{Ab}(Y_\etale)} (\mathcal{G}, f_*\mathcal{F})
$$
pour $\mathcal{F} \in \textit{Ab}(X_\etale)$ et
$\mathcal{G} \in \textit{Ab}(Y_\etale)$.
\end{definition}

\noindent
L'existence d'un tel adjoint n'est pas triviale.
En revanche, il existe dans un cadre assez g\'en\'eral, voir la
remarque \ref{etale-cohomology-remark-functoriality-general}
ci-dessous. Le formalisme g\'en\'eral montre que $f^{-1}\mathcal{G}$
est le faisceau associ\'e au pr\'efaisceau
\begin{equation}
\label{etale-cohomology-equation-pullback}
U/X
\longmapsto
\colim_{U \to X \times_Y V} \mathcal{G}(V/Y)
\end{equation}
o\`u la limite inductive est prise sur la cat\'egorie des couples
$(V/Y, \varphi : U/X \to X \times_Y V/X)$.
Pour le voir, appliquons
Sites, Proposition \ref{sites-proposition-get-morphism}
au foncteur $u$ de l'\'equation (\ref{etale-cohomology-equation-functorial})
et utilisons la description de $u_s = (u_p\ )^\#$ dans
Sites, sections \ref{sites-section-continuous-functors} et
\ref{sites-section-functoriality-PSh}.
Nous utiliserons occasionnellement cette formule pour l'image inverse
afin de d\'emontrer certaines de ses propri\'et\'es fondamentales.

\begin{lemma}
\label{etale-cohomology-lemma-stalk-pullback}
Soit $f : X \to Y$ un morphisme de sch\'emas.
\begin{enumerate}
\item Le foncteur
$f^{-1} : \textit{Ab}(Y_\etale) \to \textit{Ab}(X_\etale)$
est exact.
\item Le foncteur
$f^{-1} : \Sh(Y_\etale) \to \Sh(X_\etale)$
est exact, c'est-\`a-dire qu'il commute aux limites et colimites finies, voir
Catégories, définition \ref{categories-definition-exact}.
\item Soit $\overline{x} \to X$ un point g\'eom\'etrique.
Soit $\mathcal{G}$ un faisceau sur $Y_\etale$.
Il existe alors une identification canonique
$$
(f^{-1}\mathcal{G})_{\overline{x}} = \mathcal{G}_{\overline{y}}.
$$
o\`u $\overline{y} = f \circ \overline{x}$.
\item Pour tout $V \to Y$ \'etale, nous avons $f^{-1}h_V = h_{X \times_Y V}$.
\end{enumerate}
\end{lemma}

\begin{proof}
L'exactitude de $f^{-1}$ sur les faisceaux d'ensembles r\'esulte de
Sites, Proposition \ref{sites-proposition-get-morphism}
appliqu\'ee \`a notre foncteur $u$ de l'\'equation (\ref{etale-cohomology-equation-functorial}).
En fait, l'exactitude de l'image inverse fait partie de la d\'efinition d'un
morphisme de topos (ou de sites, si l'on pr\'ef\`ere). Ainsi, (2) est v\'erifi\'e.
Cela implique (1), car, pour tout faisceau ab\'elien $\mathcal{G}$ sur
$Y_\etale$,
le faisceau d'ensembles sous-jacent \`a $f^{-1}\mathcal{F}$ est le m\^eme
que l'image par $f^{-1}$ du faisceau d'ensembles sous-jacent \`a $\mathcal{F}$, voir
Sites, section \ref{sites-section-sheaves-algebraic-structures}.
Voir aussi
Modules sur les sites, lemme \ref{sites-modules-lemma-flat-pullback-exact}.
Dans la litt\'erature, (1) et (2) sont parfois d\'eduits de (3) au moyen du
th\'eor\`eme \ref{etale-cohomology-theorem-exactness-stalks}.

\medskip\noindent
L'assertion (3) est un fait g\'en\'eral sur les fibres des images inverses, voir
Sites, lemme \ref{sites-lemma-point-morphism-sites}.
Nous allons aussi d\'emontrer directement (3) comme suit. Notons, d'apr\`es le
lemme \ref{etale-cohomology-lemma-stalk-exact},
que la prise des fibres commute \`a la faisceautisation.
Rappelons maintenant que $f^{-1}\mathcal{G}$ est le faisceau
associ\'e au pr\'efaisceau
$$
U \longrightarrow \colim_{U \to X \times_Y V} \mathcal{G}(V),
$$
voir l'\'equation (\ref{etale-cohomology-equation-pullback}).
Ainsi, nous avons
\begin{align*}
(f^{-1}\mathcal{G})_{\overline{x}}
& = \colim_{(U, \overline{u})} f^{-1}\mathcal{G}(U) \\
& = \colim_{(U, \overline{u})}
\colim_{a : U \to X \times_Y V} \mathcal{G}(V) \\
& = \colim_{(V, \overline{v})} \mathcal{G}(V) \\
& = \mathcal{G}_{\overline{y}}
\end{align*}
Dans la troisi\`eme \`egalit\'e, le couple $(U, \overline{u})$ et le morphisme
$a : U \to X \times_Y V$ correspondent au couple $(V, a \circ \overline{u})$.

\medskip\noindent
L'assertion (4) se d\'emontre de mani\`ere analogue en identifiant les limites inductives
qui d\'efinissent $f^{-1}h_V$. On peut aussi utiliser
le lemme de Yoneda (Catégories, lemme \ref{categories-lemma-yoneda})
et les \`egalit\'es fonctorielles
$$
\Mor_{\Sh(X_\etale)}(f^{-1}h_V, \mathcal{F}) =
\Mor_{\Sh(Y_\etale)}(h_V, f_*\mathcal{F}) =
f_*\mathcal{F}(V) = \mathcal{F}(X \times_Y V)
$$
ainsi que le fait que les pr\'efaisceaux repr\'esentables sont des faisceaux. Voir aussi
Sites, lemme \ref{sites-lemma-pullback-representable-sheaf}
pour un r\'esultat enti\`erement g\'en\'eral.
\end{proof}

\noindent
Le couple de foncteurs $(f_*, f^{-1})$ d\'efinit un morphisme de petits topos \'etales
$$
f_{small} :
\Sh(X_\etale)
\longrightarrow
\Sh(Y_\etale)
$$
De nombreux r\'esultats g\'en\'eraux sur la cohomologie des faisceaux valent pour les topos et les
morphismes de topos. Nous nous efforcerons de signaler ceux qui sont
g\'en\'eraux et ceux qui sont propres au topos \'etale.

\begin{remark}
\label{etale-cohomology-remark-functoriality-general}
Plus g\'en\'eralement, soient $\mathcal{C}_1, \mathcal{C}_2$ des sites, et
supposons qu'ils poss\`edent des objets finals et des produits fibr\'es. Soit
$u: \mathcal{C}_2 \to \mathcal{C}_1$ un foncteur satisfaisant aux conditions suivantes :
\begin{enumerate}
\item si $\{V_i \to V\}$ est un recouvrement dans $\mathcal{C}_2$, alors
$\{u(V_i) \to u(V)\}$ est un recouvrement dans $\mathcal{C}_1$ (nous
disons que $u$ est {\it continu}), et
\item $u$ commute aux limites finies (c'est-\`a-dire que $u$ est exact \`a gauche, ou encore que
$u$ pr\'eserve les produits fibr\'es et les objets finals).
\end{enumerate}
On peut alors d\'efinir
$f_*: \Sh(\mathcal{C}_1) \to \Sh(\mathcal{C}_2)$
par $ f_* \mathcal{F}(V) = \mathcal{F}(u(V))$.
De plus, il existe un foncteur exact $f^{-1}$ qui
est adjoint \`a gauche de $f_*$, voir
Sites, définition \ref{sites-definition-morphism-sites} et
Proposition \ref{sites-proposition-get-morphism}.
Attention : demander simplement que $u$ soit continu
et commute aux produits fibr\'es ne suffit pas pour obtenir un morphisme de topos.
\end{remark}




\section{Fonctorialit\'e des grands topos}
\label{etale-cohomology-section-functoriality-big-topoi}

\noindent
Pour tout morphisme de sch\'emas $f : X \to Y$, on dispose de toute une famille de
morphismes de topos associ\'es \`a $f$, voir
Topologies, section \ref{topologies-section-change-topologies}
pour une liste. Les plus utilis\'es sont peut-\^etre les morphismes de topos
$$
f_{big} = f_{big, \tau} :
\Sh((\Sch/X)_\tau)
\longrightarrow
\Sh((\Sch/Y)_\tau)
$$
o\`u $\tau \in \{Zariski, \etale, smooth, syntomic, fppf\}$.
Chacun d'eux correspond \`a un foncteur continu
$$
(\Sch/Y)_\tau \longrightarrow (\Sch/X)_\tau, \quad
V/Y \longmapsto X \times_Y V/X
$$
qui pr\'eserve les objets finals, les produits fibr\'es et les recouvrements, et d\'efinit donc
un morphisme de sites
$$
f_{big} : (\Sch/X)_\tau \longrightarrow (\Sch/Y)_\tau.
$$
Voir
Topologies, sections \ref{topologies-section-zariski},
\ref{topologies-section-etale},
\ref{topologies-section-smooth},
\ref{topologies-section-syntomic}, et
\ref{topologies-section-fppf}.
En particulier, l'image directe par $f_{big}$ est donn\'ee par la r\`egle
$$
(f_{big, *}\mathcal{F})(V/Y) = \mathcal{F}(X \times_Y V/X)
$$
Il se trouve que ces morphismes de topos poss\`edent un foncteur image
inverse $f_{big}^{-1}$ tr\`es facile \`a d\'ecrire.
Plus pr\'ecis\'ement, nous avons
$$
(f_{big}^{-1}\mathcal{G})(U/X) = \mathcal{G}(U/Y)
$$
o\`u le morphisme structural de $U/Y$ est la compos\'ee du
morphisme structural $U \to X$ et de $f$, voir
Topologies, lemmes \ref{topologies-lemma-morphism-big},
\ref{topologies-lemma-morphism-big-etale},
\ref{topologies-lemma-morphism-big-smooth},
\ref{topologies-lemma-morphism-big-syntomic}, et
\ref{topologies-lemma-morphism-big-fppf}.







\section{Fonctorialit\'e et faisceaux de modules}
\label{etale-cohomology-section-morphisms-modules}

\noindent
Dans cette section, nous allons reformuler une partie des r\'esultats expliqu\'es dans
Descente, sections \ref{descent-section-quasi-coherent-sheaves},
\ref{descent-section-quasi-coherent-cohomology}, et
\ref{descent-section-quasi-coherent-sheaves-bis}
dans le cadre des topologies \'etales. Soit $f : X \to Y$ un morphisme de
sch\'emas. Nous avons vu plus haut, voir les
sections \ref{etale-cohomology-section-functoriality}, \ref{etale-cohomology-section-direct-image} et
\ref{etale-cohomology-section-inverse-image},
qu'il induit un morphisme $f_{small}$ de petits sites \'etales. Dans
Descente, remarque \ref{descent-remark-change-topologies-ringed},
nous avons vu que $f$ induit aussi un morphisme naturel
$$
f_{small}^\sharp :
\mathcal{O}_{Y_\etale}
\longrightarrow
f_{small, *}\mathcal{O}_{X_\etale}
$$
de faisceaux d'anneaux sur
$Y_\etale$ tel que $(f_{small}, f_{small}^\sharp)$
soit un morphisme de sites annel\'es. Voir
Modules sur les sites, définition \ref{sites-modules-definition-ringed-site}
pour la d\'efinition d'un morphisme de sites annel\'es.
Rappelons simplement ici que $f_{small}^\sharp$ est d\'efini par le
syst\`eme compatible de morphismes
$$
\text{pr}_V^\sharp : \mathcal{O}(V) \longrightarrow \mathcal{O}(X \times_Y V)
$$
lorsque $V$ parcourt les objets de $Y_\etale$.

\medskip\noindent
Il est clair que cette construction est compatible avec les compositions de
morphismes de sch\'emas. Plus pr\'ecis\'ement, si $f : X \to Y$ et $g : Y \to Z$
sont des morphismes de sch\'emas, alors nous avons
$$
(g_{small}, g_{small}^\sharp)
\circ
(f_{small}, f_{small}^\sharp)
=
((g \circ f)_{small}, (g \circ f)_{small}^\sharp)
$$
comme morphismes de topos annel\'es. De plus, d'apr\`es
Modules sur les sites, définition \ref{sites-modules-definition-pushforward},
nous voyons que tout morphisme de sch\'emas $f : X \to Y$
donne des foncteurs bien d\'efinis d'image inverse et d'image directe
\begin{align*}
f_{small}^* :
\textit{Mod}(\mathcal{O}_{Y_\etale})
\longrightarrow
\textit{Mod}(\mathcal{O}_{X_\etale}), \\
f_{small, *} :
\textit{Mod}(\mathcal{O}_{X_\etale})
\longrightarrow
\textit{Mod}(\mathcal{O}_{Y_\etale})
\end{align*}
qui sont adjoints de la mani\`ere habituelle. Si $g : Y \to Z$ est un autre morphisme
de sch\'emas, alors nous avons
$(g \circ f)_{small}^* = f_{small}^* \circ g_{small}^*$
et $(g \circ f)_{small, *} = g_{small, *} \circ f_{small, *}$
en vertu de ce que nous avons dit sur les compositions.

\medskip\noindent
Il y a une diff\'erence assez importante entre la cat\'egorie
de tous les $\mathcal{O}_X$-modules sur $X$ et la cat\'egorie de tous les
$\mathcal{O}_{X_\etale}$-modules sur $X_\etale$. Mais les
r\'esultats de
Descente, sections \ref{descent-section-quasi-coherent-sheaves},
\ref{descent-section-quasi-coherent-cohomology}, et
\ref{descent-section-quasi-coherent-sheaves-bis}
nous disent qu'il y a peu de diff\'erence entre consid\'erer les modules quasi-coh\'erents
sur $S$ et les modules quasi-coh\'erents sur $S_\etale$.
(Nous l'avons d\'ej\`a vu dans le
th\'eor\`eme \ref{etale-cohomology-theorem-quasi-coherent},
par exemple.)
En particulier, si $f : X \to Y$ est un morphisme quelconque de sch\'emas, alors
les foncteurs d'image inverse $f_{small}^*$ et $f^*$ s'identifient pour les
faisceaux quasi-coh\'erents, voir
Descente,
Proposition \ref{descent-proposition-equivalence-quasi-coherent-functorial}.
De plus, il en va de m\^eme pour l'image directe pourvu que $f$ soit
quasi-compact et quasi-s\'epar\'e, voir
Descente, lemme \ref{descent-lemma-higher-direct-images-small-etale}.

\medskip\noindent
Disons quelques mots de la fonctorialit\'e du faisceau structural sur les grands sites.
Soit $f : X \to Y$ un morphisme de sch\'emas. Choisissons l'une des
topologies $\tau \in \{Zariski,\linebreak[0]
\etale,\linebreak[0] smooth,\linebreak[0] syntomic,\linebreak[0]
fppf\}$. Alors le morphisme
$f_{big} : (\Sch/X)_\tau \to (\Sch/Y)_\tau$
devient un morphisme de sites annel\'es au moyen d'un morphisme
$$
f_{big}^\sharp : \mathcal{O}_Y \longrightarrow f_{big, *}\mathcal{O}_X
$$
voir Descente, remarque \ref{descent-remark-change-topologies-ringed}.
En fait, il est donn\'e par la m\^eme construction que celle expos\'ee pour les petits
sites plus haut.












\section{Comparaison de topologies}
\label{etale-cohomology-section-compare-topologies}

\noindent
Dans cette section, nous commen\c{c}ons \`a \`etudier ce qui se produit lorsque l'on compare
les faisceaux relativement \`a diff\'erentes topologies.

\begin{lemma}
\label{etale-cohomology-lemma-where-sections-are-equal}
Soit $S$ un sch\'ema. Soit $\mathcal{F}$ un faisceau d'ensembles sur $S_\etale$.
Soient $s, t \in \mathcal{F}(S)$. Il existe alors un ouvert $W \subset S$
caract\'eris\'e par la propri\'et\'e suivante : un morphisme $f : T \to S$
se factorise par $W$ si et seulement si $s|_T = t|_T$ (la restriction est
l'image inverse par $f_{small}$).
\end{lemma}

\begin{proof}
Consid\'erons le pr\'efaisceau qui associe \`a $U \in \Ob(S_\etale)$ l'ensemble vide
si $s|_U \not = t|_U$, et un singleton sinon. Il est clair qu'il s'agit
d'un sous-faisceau de l'objet final de $\Sh(S_\etale)$. D'apr\`es le
lemme \ref{etale-cohomology-lemma-support-subsheaf-final},
nous trouvons un ouvert $W \subset S$ qui repr\'esente ce pr\'efaisceau.
Pour un point g\'eom\'etrique $\overline{x}$ de $S$, nous voyons que $\overline{x} \in W$
si et seulement si les germes de $s$ et de $t$ en $\overline{x}$ sont identiques.
D'apr\`es la description des fibres des images inverses donn\'ee dans le
lemme \ref{etale-cohomology-lemma-stalk-pullback},
nous voyons que $W$ poss\`ede la propri\'et\'e voulue.
\end{proof}

\begin{lemma}
\label{etale-cohomology-lemma-describe-pullback}
Soit $S$ un sch\'ema. Soit $\tau \in \{Zariski, \etale\}$. Consid\'erons le morphisme
$$
\pi_S : (\Sch/S)_\tau \longrightarrow S_\tau
$$
de Topologies, lemme \ref{topologies-lemma-at-the-bottom} ou
\ref{topologies-lemma-at-the-bottom-etale}. Soit $\mathcal{F}$ un faisceau sur
$S_\tau$. Alors $\pi_S^{-1}\mathcal{F}$ est donn\'e par la r\`egle
$$
(\pi_S^{-1}\mathcal{F})(T) = \Gamma(T_\tau, f_{small}^{-1}\mathcal{F})
$$
o\`u $f : T \to S$. De plus, $\pi_S^{-1}\mathcal{F}$ satisfait \`a la
condition de faisceau relativement aux recouvrements fpqc.
\end{lemma}

\begin{proof}
Observons que nous avons un morphisme $i_f : \Sh(T_\tau) \to \Sh(\Sch/S)_\tau)$
tel que $\pi_S \circ i_f = f_{small}$ comme morphismes
$T_\tau \to S_\tau$, voir
Topologies, lemmes \ref{topologies-lemma-put-in-T},
\ref{topologies-lemma-morphism-big-small},
\ref{topologies-lemma-put-in-T-etale}, et
\ref{topologies-lemma-morphism-big-small-etale}.
Comme l'image inverse est transitive, nous voyons que
$i_f^{-1} \pi_S^{-1}\mathcal{F} = f_{small}^{-1}\mathcal{F}$ comme voulu.

\medskip\noindent
Soit $\{g_i : T_i \to T\}_{i \in I}$ un recouvrement fpqc.
La derni\`ere assertion signifie ce qui suit : \`etant donn\'e un faisceau $\mathcal{G}$
sur $T_\tau$ et des sections
$s_i \in \Gamma(T_i, g_{i, small}^{-1}\mathcal{G})$ dont les images inverses
sur $T_i \times_T T_j$ concordent, il existe une unique section $s$ de $\mathcal{G}$
sur $T$ dont l'image inverse sur $T_i$ concorde avec $s_i$.

\medskip\noindent
Soit $V \to T$ un objet de $T_\tau$ et soit $t \in \mathcal{G}(V)$.
Pour tout $i$, il existe un plus grand ouvert $W_i \subset T_i \times_T V$
tel que les images inverses de $s_i$ et de $t$ concordent comme sections de l'image
inverse de $\mathcal{G}$ sur $W_i \subset T_i \times_T V$, voir le
lemme \ref{etale-cohomology-lemma-where-sections-are-equal}.
Comme $s_i$ et $s_j$ concordent sur $T_i \times_T T_j$, nous constatons
que les images inverses de $W_i$ et de $W_j$ sont le m\^eme ouvert au-dessus de
$T_i \times_T T_j \times_T V$. D'apr\`es
Descente, lemme \ref{descent-lemma-open-fpqc-covering},
nous trouvons un ouvert $W \subset V$ dont l'image inverse sur $T_i \times_T V$
redonne $W_i$.

\medskip\noindent
Par construction de $g_{i, small}^{-1}\mathcal{G}$, il existe
un recouvrement pour $\tau$ $\{T_{ij} \to T_i\}_{j \in J_i}$, pour chaque $j$ une
immersion ouverte ou un morphisme \'etale $V_{ij} \to T$, une section
$t_{ij} \in \mathcal{G}(V_{ij})$, et des diagrammes commutatifs
$$
\xymatrix{
T_{ij} \ar[r] \ar[d] & V_{ij} \ar[d] \\
T_i \ar[r] & T
}
$$
tels que $s_i|_{T_{ij}}$ soit l'image inverse de $t_{ij}$. Autrement dit,
apr\`es avoir remplac\'e le recouvrement $\{T_i \to T\}$ par $\{T_{ij} \to T\}$,
nous pouvons supposer qu'il existe des factorisations $T_i \to V_i \to T$ avec
$V_i \in \Ob(T_\tau)$ et des sections $t_i \in \mathcal{G}(V_i)$
dont les images inverses sont $s_i$ sur $T_i$.
D'apr\`es le r\'esultat du paragraphe pr\'ec\'edent, nous trouvons des ouverts $W_i \subset V_i$
tels que $t_i|_{W_i}$ ``concorde avec'' chaque $s_j$ sur $T_j \times_T W_i$.
Notons que $T_i \to V_i$ se factorise par $W_i$.
Ainsi, $\{W_i \to T\}$ est un recouvrement pour $\tau$ et le lemme est d\'emontr\'e.
\end{proof}

\begin{lemma}
\label{etale-cohomology-lemma-sections-upstairs}
Soit $S$ un sch\'ema. Soit $f : T \to S$ un morphisme tel que
\begin{enumerate}
\item $f$ est plat et quasi-compact, et
\item les fibres g\'eom\'etriques de $f$ sont connexes.
\end{enumerate}
Soit $\mathcal{F}$ un faisceau sur $S_\etale$.
Alors $\Gamma(S, \mathcal{F}) = \Gamma(T, f^{-1}_{small}\mathcal{F})$.
\end{lemma}

\begin{proof}
Il existe un morphisme canonique
$\Gamma(S, \mathcal{F}) \to \Gamma(T, f_{small}^{-1}\mathcal{F})$.
Comme $f$ est surjectif (puisque ses fibres sont connexes), nous voyons que
ce morphisme est injectif.

\medskip\noindent
Pour montrer que le morphisme est surjectif, soit
$\alpha \in \Gamma(T, f_{small}^{-1}\mathcal{F})$.
Comme $\{T \to S\}$ est un recouvrement fpqc, nous pouvons utiliser le
lemme \ref{etale-cohomology-lemma-describe-pullback} pour voir qu'il suffit de montrer que
les deux images inverses de $\alpha$ sur $T \times_S T$ sont la m\^eme section
par les deux projections. Soit $\overline{s} \to S$ un point g\'eom\'etrique.
Il suffit de montrer l'\'egalit\'e sur $(T \times_S T)_{\overline{s}}$,
car tout point g\'eom\'etrique de $T \times_S T$ est contenu dans l'une de
ces fibres g\'eom\'etriques. Autrement dit, nous voulons montrer que
$\alpha|_{T_{\overline{s}}}$ donne la m\^eme section par les deux projections sur
$$
(T \times_S T)_{\overline{s}} =
T_{\overline{s}} \times_{\overline{s}} T_{\overline{s}}
$$
vers $T_{\overline{s}}$.
Cependant, comme $\mathcal{F}|_{T_{\overline{s}}}$ est l'image
inverse de $\mathcal{F}|_{\overline{s}}$, c'est un faisceau constant
de valeur $\mathcal{F}_{\overline{s}}$. Comme $T_{\overline{s}}$
est connexe par hypoth\`ese, toute section d'un faisceau constant est constante.
Ainsi, $\alpha|_{T_{\overline{s}}}$ correspond \`a un \`el\'ement
de $\mathcal{F}_{\overline{s}}$. Les deux images inverses sur
$(T \times_S T)_{\overline{s}}$ correspondent donc toutes deux
\`a ce m\^eme \`el\'ement, d'o\`u la conclusion.
\end{proof}

\noindent
Voici une variante du lemme \ref{etale-cohomology-lemma-sections-upstairs}
o\`u nous ne supposons pas que le morphisme est plat.

\begin{lemma}
\label{etale-cohomology-lemma-sections-upstairs-submersive}
Soit $S$ un sch\'ema. Soit $f : X \to S$ un morphisme tel que
\begin{enumerate}
\item $f$ est submersif, et
\item les fibres g\'eom\'etriques de $f$ sont connexes.
\end{enumerate}
Soit $\mathcal{F}$ un faisceau sur $S_\etale$.
Alors $\Gamma(S, \mathcal{F}) = \Gamma(X, f^{-1}_{small}\mathcal{F})$.
\end{lemma}

\begin{proof}
Il existe un morphisme canonique
$\Gamma(S, \mathcal{F}) \to \Gamma(X, f_{small}^{-1}\mathcal{F})$.
Comme $f$ est surjectif (puisque ses fibres sont connexes), nous voyons que
ce morphisme est injectif.

\medskip\noindent
Pour montrer que le morphisme est surjectif, soit
$\tau \in \Gamma(X, f_{small}^{-1}\mathcal{F})$.
Il suffit de trouver
un recouvrement \'etale $\{U_i \to S\}$ et des
sections $\sigma_i \in \mathcal{F}(U_i)$
telles que l'image inverse de $\sigma_i$ soit $\tau|_{X \times_S U_i}$.
En effet, l'injectivit\'e montr\'ee ci-dessus garantit
que $\sigma_i$ et $\sigma_j$ induisent la m\^eme
section de $\mathcal{F}$ sur $U_i \times_S U_j$.
Nous obtenons donc une unique section $\sigma \in \mathcal{F}(S)$
qui se restreint \`a $\sigma_i$ sur $U_i$.
L'image inverse de $\sigma$ sur $X$ est alors $\tau$,
car c'est vrai localement.

\medskip\noindent
Soit $\overline{x}$ un point g\'eom\'etrique de $X$ d'image $\overline{s}$
dans $S$. Consid\'erons l'image de $\tau$ dans la fibre
$$
(f_{small}^{-1}\mathcal{F})_{\overline{x}} = \mathcal{F}_{\overline{s}}
$$
Voir le lemme \ref{etale-cohomology-lemma-stalk-pullback}.
Nous pouvons trouver un voisinage \'etale $U \to S$ de $\overline{s}$
et une section $\sigma \in \mathcal{F}(U)$ qui s'envoie sur cette image
dans la fibre. Ainsi, apr\`es avoir remplac\'e $S$ par $U$ et $X$ par $X \times_S U$,
nous pouvons supposer qu'il existe une section $\sigma$ de $\mathcal{F}$ sur $S$
dont l'image dans $(f_{small}^{-1}\mathcal{F})_{\overline{x}}$ est la
m\^eme que celle de $\tau$.

\medskip\noindent
D'apr\`es le lemme \ref{etale-cohomology-lemma-where-sections-are-equal},
il existe un ouvert maximal $W \subset X$ sur lequel
$f_{small}^{-1}\sigma$ et $\tau$ concordent sur $W$,
et la formation de $W$ commute aux changements de base ult\'erieurs.
Observons que l'image inverse de $\mathcal{F}$ sur la
fibre g\'eom\'etrique $X_{\overline{s}}$ est l'image inverse
de $\mathcal{F}_{\overline{s}}$, consid\'er\'e comme faisceau sur
$\overline{s}$, par $X_{\overline{s}} \to \overline{s}$.
Nous voyons donc que $\tau$ et $\sigma$ donnent des sections
du faisceau constant de valeur $\mathcal{F}_{\overline{s}}$
sur $X_{\overline{s}}$ qui concordent en un point. Comme
$X_{\overline{s}}$ est connexe par hypoth\`ese, nous concluons
que $W$ contient $X_s$. Le m\^eme argument appliqu\'e aux autres
fibres g\'eom\'etriques montre que $W$ contient toute fibre qu'il rencontre.
Comme $f$ est submersif, nous concluons que $W$ est l'image
inverse d'un voisinage ouvert de $s$ dans $S$.
Cela ach\`eve la d\'emonstration.
\end{proof}

\begin{lemma}
\label{etale-cohomology-lemma-sections-base-field-extension}
Soit $K/k$ une extension de corps o\`u $k$ est
s\'eparablement clos. Soit $S$ un sch\'ema sur $k$. Notons
$p : S_K = S \times_{\Spec(k)} \Spec(K) \to S$ la projection.
Soit $\mathcal{F}$ un faisceau sur $S_\etale$.
Alors $\Gamma(S, \mathcal{F}) = \Gamma(S_K, p^{-1}_{small}\mathcal{F})$.
\end{lemma}

\begin{proof}
Cela r\'esulte du lemme \ref{etale-cohomology-lemma-sections-upstairs}. En effet, il est clair
que $p$ est plat et quasi-compact comme changement de base de
$\Spec(K) \to \Spec(k)$. D'autre part, si $\overline{s} : \Spec(L) \to S$
est un point g\'eom\'etrique, alors la fibre de $p$ en $\overline{s}$
est le spectre de $K \otimes_k L$, qui est irr\'eductible, donc connexe, d'apr\`es
Algèbre, lemme \ref{algebra-lemma-separably-closed-irreducible}.
\end{proof}









\section{Reconstruction des morphismes}
\label{etale-cohomology-section-morphisms}

\noindent
Dans cette section, nous d\'emontrons que la r\`egle qui associe \`a un sch\'ema
son petit topos \'etale localement annel\'e est pleinement fid\`ele en un sens
convenable, voir le
th\'eor\`eme \ref{etale-cohomology-theorem-fully-faithful}.

\begin{lemma}
\label{etale-cohomology-lemma-morphism-locally-ringed}
Soit $f : X \to Y$ un morphisme de sch\'emas.
Le morphisme de sites annel\'es $(f_{small}, f_{small}^\sharp)$
associ\'e \`a $f$ est un morphisme de sites localement annel\'es, voir
Modules sur les sites,
définition \ref{sites-modules-definition-morphism-locally-ringed-topoi}.
\end{lemma}

\begin{proof}
Notons que l'assertion a un sens, puisque nous avons vu que
$(X_\etale, \mathcal{O}_{X_\etale})$ et
$(Y_\etale, \mathcal{O}_{Y_\etale})$
sont des sites localement annel\'es, voir le
lemme \ref{etale-cohomology-lemma-etale-site-locally-ringed}.
De plus, nous savons que $X_\etale$ a assez de points, voir le
th\'eor\`eme \ref{etale-cohomology-theorem-exactness-stalks}
et les
remarques \ref{etale-cohomology-remarks-enough-points}.
Il suffit donc de d\'emontrer que $(f_{small}, f_{small}^\sharp)$
satisfait \`a la condition (3) de
Modules sur les sites,
lemme \ref{sites-modules-lemma-locally-ringed-morphism}.
Pour le voir, prenons un point $p$ de $X_\etale$. D'apr\`es le
lemme \ref{etale-cohomology-lemma-points-small-etale-site},
$p$ correspond \`a un point g\'eom\'etrique $\overline{x}$ de $X$.
D'apr\`es le
lemme \ref{etale-cohomology-lemma-stalk-pullback},
le point $q = f_{small} \circ p$ correspond au
point g\'eom\'etrique $\overline{y} = f \circ \overline{x}$ de $Y$.
L'assertion \`a d\'emontrer est donc que le morphisme induit
sur les fibres
$$
(\mathcal{O}_Y)_{\overline{y}} \longrightarrow (\mathcal{O}_X)_{\overline{x}}
$$
est un morphisme local d'anneaux. Supposons que $a \in (\mathcal{O}_Y)_{\overline{y}}$
soit un \`el\'ement du membre de gauche qui s'envoie sur un \`el\'ement de l'id\'eal
maximal du membre de droite. Supposons que $a$ soit la classe d'\'equivalence
d'un triplet $(V, \overline{v}, a)$ o\`u $V \to Y$ est \'etale,
$\overline{v} : \overline{x} \to V$ au-dessus de $Y$, et $a \in \mathcal{O}(V)$.
Il s'envoie sur la classe d'\'equivalence de
$(X \times_Y V, \overline{x} \times \overline{v}, \text{pr}_V^\sharp(a))$
dans l'anneau local $(\mathcal{O}_X)_{\overline{x}}$. Mais il est clair que
le fait d'appartenir \`a l'id\'eal maximal signifie que l'image de $\text{pr}_V^\sharp(a)$
dans $\kappa(\overline{x})$ est nulle. Par cons\'equent, l'image de
$a$ dans $\kappa(\overline{x})$ est elle aussi nulle. Cela signifie que $a$ appartient \`a
l'id\'eal maximal de $(\mathcal{O}_Y)_{\overline{y}}$.
\end{proof}

\begin{lemma}
\label{etale-cohomology-lemma-2-morphism}
Soient $X$, $Y$ des sch\'emas. Soit $f : X \to Y$ un morphisme de sch\'emas.
Soit $t$ un $2$-morphisme de $(f_{small}, f_{small}^\sharp)$ vers lui-m\^eme, voir
Modules sur les sites,
définition \ref{sites-modules-definition-2-morphism-ringed-topoi}.
Alors $t = \text{id}$.
\end{lemma}

\begin{proof}
Cela signifie que $t : f^{-1}_{small} \to f^{-1}_{small}$
est une transformation naturelle telle que le diagramme
$$
\xymatrix{
f_{small}^{-1}\mathcal{O}_Y
\ar[rd]_{f_{small}^\sharp}  & &
f_{small}^{-1}\mathcal{O}_Y \ar[ll]^t \ar[ld]^{f_{small}^\sharp} \\
& \mathcal{O}_X
}
$$
soit commutatif. Supposons que $V \to Y$ soit \'etale avec $V$ affine. D'apr\`es
Morphismes, lemme \ref{morphisms-lemma-quasi-affine-finite-type-over-S},
nous pouvons choisir une immersion $i : V \to \mathbf{A}^n_Y$ au-dessus de $Y$.
En termes de faisceaux, cela signifie que $i$ induit une injection
$h_i : h_V \to \prod_{j = 1, \ldots, n} \mathcal{O}_Y$ de faisceaux.
Le changement de base $i'$ de $i$ \`a $X$ est une immersion
(Schémas, lemme \ref{schemes-lemma-base-change-immersion}).
Ainsi, $i' : X \times_Y V \to \mathbf{A}^n_X$ est une immersion, ce qui
signifie \`a son tour que
$h_{i'} : h_{X \times_Y V} \to \prod_{j = 1, \ldots, n} \mathcal{O}_X$
est une injection de faisceaux.
Via l'identification $f_{small}^{-1}h_V = h_{X \times_Y V}$ du
lemme \ref{etale-cohomology-lemma-stalk-pullback},
le morphisme $h_{i'}$ est \`egal \`a
$$
\xymatrix{
f_{small}^{-1}h_V \ar[r]^-{f^{-1}h_i} &
\prod_{j = 1, \ldots, n} f_{small}^{-1}\mathcal{O}_Y
\ar[r]^{\prod f^\sharp} &
\prod_{j = 1, \ldots, n} \mathcal{O}_X
}
$$
(nous omettons la v\'erification). Cela signifie que le morphisme
$t : f_{small}^{-1}h_V \to f_{small}^{-1}h_V$
s'ins\`ere dans le diagramme commutatif
$$
\xymatrix{
f_{small}^{-1}h_V \ar[r]^-{f^{-1}h_i} \ar[d]^t &
\prod_{j = 1, \ldots, n} f_{small}^{-1}\mathcal{O}_Y
\ar[r]^-{\prod f^\sharp} \ar[d]^{\prod t} &
\prod_{j = 1, \ldots, n} \mathcal{O}_X \ar[d]^{\text{id}}\\
f_{small}^{-1}h_V \ar[r]^-{f^{-1}h_i} &
\prod_{j = 1, \ldots, n} f_{small}^{-1}\mathcal{O}_Y
\ar[r]^-{\prod f^\sharp} &
\prod_{j = 1, \ldots, n} \mathcal{O}_X
}
$$
Le carr\'e de droite est commutatif par notre hypoth\`ese sur $t$
expliqu\'ee ci-dessus.
Comme la compos\'ee des fl\`eches horizontales est injective
d'apr\`es la discussion pr\'ec\'edente, nous concluons que la fl\`eche verticale de gauche
est elle aussi l'identit\'e. Tout faisceau d'ensembles sur
$Y_\etale$ est quotient d'un coproduit (\'enorme) de faisceaux
de la forme $h_V$ avec $V$ affine (combiner
Topologies, lemme \ref{topologies-lemma-alternative}
avec
Sites, lemme \ref{sites-lemma-sheaf-coequalizer-representable}).
Nous concluons donc que $t : f_{small}^{-1} \to f_{small}^{-1}$
est la transformation identit\'e, comme voulu.
\end{proof}

\begin{lemma}
\label{etale-cohomology-lemma-faithful}
Soient $X$, $Y$ des sch\'emas.
Deux morphismes de sch\'emas $a, b : X \to Y$
pour lesquels il existe un $2$-isomorphisme
$(a_{small}, a_{small}^\sharp) \cong (b_{small}, b_{small}^\sharp)$
dans la $2$-cat\'egorie des topos annel\'es sont \`egaux.
\end{lemma}

\begin{proof}
Donnons l'argument avec soin, car il est quelque peu d\'eroutant.
Soit $t : a_{small}^{-1} \to b_{small}^{-1}$ le $2$-isomorphisme.
Consid\'erons un ouvert quelconque $V \subset Y$. Notons que $h_V$ est un sous-faisceau du
faisceau final $*$. Ainsi, $a_{small}^{-1}h_V = h_{a^{-1}(V)}$
et $b_{small}^{-1}h_V = h_{b^{-1}(V)}$ sont tous deux des sous-faisceaux du faisceau final.
L'isomorphisme
$$
t : a_{small}^{-1}h_V = h_{a^{-1}(V)} \to b_{small}^{-1}h_V = h_{b^{-1}(V)}
$$
est donc n\'ecessairement l'identit\'e, et $a^{-1}(V) = b^{-1}(V)$.
Il s'ensuit que $a$ et $b$ sont \`egaux sur les espaces topologiques sous-jacents.
Prenons ensuite une section $f \in \mathcal{O}_Y(V)$. Elle d\'etermine, et est
d\'etermin\'ee par, un morphisme de faisceaux d'ensembles
$f : h_V \to \mathcal{O}_Y$.
Prenons son image inverse et appliquons $t$ pour obtenir un diagramme commutatif
$$
\xymatrix{
h_{b^{-1}(V)} \ar@{=}[r] &
b_{small}^{-1}h_V \ar[d]^{b_{small}^{-1}(f)} & &
a_{small}^{-1}h_V \ar[d]^{a_{small}^{-1}(f)} \ar[ll]^t &
h_{a^{-1}(V)} \ar@{=}[l]
\\
& b_{small}^{-1}\mathcal{O}_Y
\ar[rd]_{b^\sharp}  & &
a_{small}^{-1}\mathcal{O}_Y \ar[ll]^t \ar[ld]^{a^\sharp} \\
& & \mathcal{O}_X
}
$$
o\`u le triangle est commutatif par d\'efinition d'un $2$-isomorphisme dans
Modules sur les sites, section \ref{sites-modules-section-2-category}.
Nous avons vu plus haut que la compos\'ee des fl\`eches horizontales
sup\'erieures provient de l'identit\'e $a^{-1}(V) = b^{-1}(V)$.
La commutativit\'e du diagramme nous dit donc que
$a_{small}^\sharp(f) = b_{small}^\sharp(f)$ dans
$\mathcal{O}_X(a^{-1}(V)) = \mathcal{O}_X(b^{-1}(V))$.
Comme cela vaut pour tout ouvert $V$ et tout $f \in \mathcal{O}_Y(V)$,
nous concluons que $a = b$ comme morphismes de sch\'emas.
\end{proof}

\begin{lemma}
\label{etale-cohomology-lemma-morphism-ringed-etale-topoi-affines}
Soient $X$, $Y$ des sch\'emas affines.
Soit
$$
(g, g^\#) :
(\Sh(X_\etale), \mathcal{O}_X)
\longrightarrow
(\Sh(Y_\etale), \mathcal{O}_Y)
$$
un morphisme de topos localement annel\'es. Il existe alors un
unique morphisme de sch\'emas $f : X \to Y$ tel que
$(g, g^\#)$ soit $2$-isomorphe \`a $(f_{small}, f_{small}^\sharp)$,
voir
Modules sur les sites,
définition \ref{sites-modules-definition-2-morphism-ringed-topoi}.
\end{lemma}

\begin{proof}
Dans cette d\'emonstration, nous notons $\mathcal{O}_X$ le faisceau structural
du petit site \'etale $X_\etale$, et de m\^eme pour
$\mathcal{O}_Y$. Posons $Y = \Spec(B)$ et $X = \Spec(A)$. Comme
$B = \Gamma(Y_\etale, \mathcal{O}_Y)$,
$A = \Gamma(X_\etale, \mathcal{O}_X)$,
nous voyons que $g^\sharp$ induit un morphisme d'anneaux $\varphi : B \to A$.
Soit $f = \Spec(\varphi) : X \to Y$ le morphisme correspondant
de sch\'emas affines. Nous allons montrer que ce $f$ convient.

\medskip\noindent
Soit $V \to Y$ un sch\'ema affine \'etale sur $Y$. Nous pouvons donc \`ecrire
$V = \Spec(C)$, o\`u $C$ est une $B$-alg\`ebre \'etale. Nous pouvons \`ecrire
$$
C = B[x_1, \ldots, x_n]/(P_1, \ldots, P_n)
$$
avec des polyn\^omes $P_i$ tels que $\Delta = \det(\partial P_i/ \partial x_j)$
soit inversible dans $C$, voir par exemple
Algèbre, lemme \ref{algebra-lemma-etale-standard-smooth}.
Si $T$ est un sch\'ema sur $Y$, alors un $T$-point de $V$ est donn\'e par
$n$ sections de $\Gamma(T, \mathcal{O}_T)$ qui satisfont aux \`equations polynomiales
$P_1 = 0, \ldots, P_n = 0$. Autrement dit, le faisceau $h_V$
sur $Y_\etale$ est l'\'egalisateur des deux morphismes
$$
\xymatrix{
\prod\nolimits_{i = 1, \ldots, n} \mathcal{O}_Y
\ar@<1ex>[r]^a \ar@<-1ex>[r]_b &
\prod\nolimits_{j = 1, \ldots, n} \mathcal{O}_Y
}
$$
o\`u $b(h_1, \ldots, h_n) = 0$ et
$a(h_1, \ldots, h_n) =
(P_1(h_1, \ldots, h_n), \ldots, P_n(h_1, \ldots, h_n))$.
Comme $g^{-1}$ est exact, nous concluons que la ligne sup\'erieure de la
partie du diagramme commutatif suivant form\'ee par les fl\`eches pleines est elle aussi un diagramme d'\'egalisateur :
$$
\xymatrix{
g^{-1}h_V \ar[r] \ar@{..>}[d] &
\prod\nolimits_{i = 1, \ldots, n} g^{-1}\mathcal{O}_Y
\ar@<1ex>[r]^{g^{-1}a} \ar@<-1ex>[r]_{g^{-1}b} \ar[d]^{\prod g^\sharp} &
\prod\nolimits_{j = 1, \ldots, n} g^{-1}\mathcal{O}_Y \ar[d]^{\prod g^\sharp}\\
h_{X \times_Y V} \ar[r] &
\prod\nolimits_{i = 1, \ldots, n} \mathcal{O}_X
\ar@<1ex>[r]^{a'} \ar@<-1ex>[r]_{b'} &
\prod\nolimits_{j = 1, \ldots, n} \mathcal{O}_X  \\
}
$$
Ici, $b'$ est l'application nulle et $a'$ est l'application d\'efinie par les
images $P'_i = \varphi(P_i) \in A[x_1, \ldots, x_n]$ suivant la m\^eme
r\`egle
$a'(h_1, \ldots, h_n) =
(P'_1(h_1, \ldots, h_n), \ldots, P'_n(h_1, \ldots, h_n))$.
que celle qui d\'efinit $a$. La commutativit\'e du diagramme r\'esulte de
l'\'egalit\'e $\varphi = g^\sharp$ sur les sections globales. La ligne
inf\'erieure est \'egalement un diagramme \'egalisateur, exactement par les m\^emes arguments que
pr\'ec\'edemment, puisque $X \times_Y V$ est le sch\'ema affine
$\Spec(A \otimes_B C)$ et
$A \otimes_B C = A[x_1, \ldots, x_n]/(P'_1, \ldots, P'_n)$.
Nous obtenons donc une unique fl\`eche pointill\'ee
$g^{-1}h_V \to h_{X \times_Y V}$ qui s'ins\`ere dans le diagramme.

\medskip\noindent
Nous affirmons que le morphisme de faisceaux $g^{-1}h_V \to h_{X \times_Y V}$
est un isomorphisme. Comme le petit site \'etale de $X$ poss\`ede suffisamment de points
(Th\'eor\`eme \ref{etale-cohomology-theorem-exactness-stalks}),
il suffit de le d\'emontrer sur les germes. Soit donc $\overline{x}$ un
point g\'eom\'etrique de $X$, et notons $p$ le point associ\'e du
petit topos \'etale de $X$. Posons $q = g \circ p$. C'est un point du
petit topos \'etale de $Y$. D'apr\`es le
Lemme \ref{etale-cohomology-lemma-points-small-etale-site},
nous voyons que $q$ correspond \`a un point g\'eom\'etrique $\overline{y}$ de
$Y$. Consid\'erons le morphisme de germes
$$
(g^\sharp)_p :
(\mathcal{O}_Y)_{\overline{y}} =
\mathcal{O}_{Y, q} =
(g^{-1}\mathcal{O}_Y)_p
\longrightarrow
\mathcal{O}_{X, p} =
(\mathcal{O}_X)_{\overline{x}}
$$
Comme $(g, g^\sharp)$ est un morphisme de topos {\it localement} annel\'es,
$(g^\sharp)_p$ est un homomorphisme local entre anneaux locaux strictement
hens\'eliens. En localisant le grand diagramme commutatif ci-dessus et en appliquant
Alg\`ebre, Lemme \ref{algebra-lemma-strictly-henselian-solutions},
nous concluons que $(g^{-1}h_V)_p \to (h_{X \times_Y V})_p$ est un isomorphisme,
comme voulu.

\medskip\noindent
Nous affirmons que les isomorphismes $g^{-1}h_V \to h_{X \times_Y V}$ sont
fonctoriels. Plus pr\'ecis\'ement, supposons que $V_1 \to V_2$ soit un morphisme de sch\'emas affines
\'etales sur $Y$. \'Ecrivons
$V_i = \Spec(C_i)$ avec
$$
C_i = B[x_{i, 1}, \ldots, x_{i, n_i}]/(P_{i, 1}, \ldots, P_{i, n_i})
$$
Le morphisme $V_1 \to V_2$ est donn\'e par un homomorphisme de $B$-alg\`ebres $C_2 \to C_1$
qui, \`a son tour, est donn\'e par des polyn\^omes
$Q_j \in B[x_{1, 1}, \ldots, x_{1, n_1}]$, pour $j = 1, \ldots, n_2$.
Il est alors ais\'e de montrer que le diagramme de faisceaux
$$
\xymatrix{
h_{V_1} \ar[d] \ar[r] & \prod_{i = 1, \ldots, n_1} \mathcal{O}_Y
\ar[d]^{Q_1, \ldots, Q_{n_2}}\\
h_{V_2} \ar[r] & \prod_{i = 1, \ldots, n_2} \mathcal{O}_Y
}
$$
est commutatif et que, par image inverse sur $X_\etale$, nous obtenons le
diagramme commutatif \`a fl\`eches pleines
$$
\xymatrix{
g^{-1}h_{V_1} \ar@{..>}[dd] \ar[rrd] \ar[r] &
\prod_{i = 1, \ldots, n_1} g^{-1}\mathcal{O}_Y
\ar[dd]^{g^\sharp}
\ar[rrd]^{Q_1, \ldots, Q_{n_2}} \\
& & g^{-1}h_{V_2} \ar@{..>}[dd] \ar[r] &
\prod_{i = 1, \ldots, n_2} g^{-1}\mathcal{O}_Y
\ar[dd]^{g^\sharp} \\
h_{X \times_Y V_1} \ar[r] \ar[rrd] &
\prod\nolimits_{i = 1, \ldots, n_1} \mathcal{O}_X
\ar[rrd]^{Q'_1, \ldots, Q'_{n_2}} \\
& & h_{X \times_Y V_2} \ar[r] &
\prod\nolimits_{i = 1, \ldots, n_2} \mathcal{O}_X
}
$$
o\`u $Q'_j \in A[x_{1, 1}, \ldots, x_{1, n_1}]$ est l'image de
$Q_j$ par $\varphi$. Puisque les fl\`eches pointill\'ees existent, rendent les
deux carr\'es commutatifs, et que les fl\`eches horizontales sont injectives,
nous voyons que le diagramme entier est commutatif. Cela prouve la fonctorialit\'e
(ainsi que l'ind\'ependance de la construction de $g^{-1}h_V \to h_{X \times_Y V}$
\`a l'\'egard du choix de la pr\'esentation, bien qu'\`a
strictement parler nous n'ayons pas besoin de le montrer).

\medskip\noindent
\`A ce stade, nous pouvons montrer que $f_{small, *} \cong g_*$.
En effet, soit $\mathcal{F}$ un faisceau sur $X_\etale$. Pour tout objet affine
$V \in \Ob(X_\etale)$, nous avons
\begin{align*}
(g_*\mathcal{F})(V)
& =
\Mor_{\Sh(Y_\etale)}(h_V, g_*\mathcal{F}) \\
& =
\Mor_{\Sh(X_\etale)}(g^{-1}h_V, \mathcal{F}) \\
& =
\Mor_{\Sh(X_\etale)}(h_{X \times_Y V}, \mathcal{F}) \\
& =
\mathcal{F}(X \times_Y V) \\
& =
f_{small, *}\mathcal{F}(V)
\end{align*}
o\`u, dans la troisi\`eme \'egalit\'e, nous utilisons l'isomorphisme
$g^{-1}h_V \cong h_{X \times_Y V}$ construit ci-dessus. Ces isomorphismes
sont manifestement fonctoriels en $\mathcal{F}$ et en $V$,
car les isomorphismes $g^{-1}h_V \cong h_{X \times_Y V}$ sont fonctoriels.
Or tout faisceau sur $Y_\etale$ est d\'etermin\'e par sa restriction
\`a la sous-cat\'egorie des sch\'emas affines
(Topologies, Lemme \ref{topologies-lemma-alternative}) ;
nous obtenons donc un isomorphisme de foncteurs $f_{small, *} \cong g_*$,
comme voulu.

\medskip\noindent
Enfin, nous devons v\'erifier que, par l'isomorphisme
$f_{small, *} \cong g_*$ ci-dessus, les morphismes $f_{small}^\sharp$ et
$g^\sharp$ co\"incident. Par construction, c'est d\'ej\`a le cas sur les
sections globales de $\mathcal{O}_Y$, c'est-\`a-dire sur les \'el\'ements de $B$.
Il suffit de v\'erifier le r\'esultat sur les
sections au-dessus d'un ouvert affine $V$ \'etale sur $Y$ (encore d'apr\`es
Topologies, Lemme \ref{topologies-lemma-alternative}).
En \'ecrivant
$V = \Spec(C)$, $C = B[x_i]/(P_j)$ comme pr\'ec\'edemment, il suffit
de v\'erifier que les fonctions coordonn\'ees $x_i$ sont envoy\'ees sur
les m\^emes sections de $\mathcal{O}_X$ au-dessus de $X \times_Y V$.
Et c'est exactement ce que signifie la commutativit\'e du diagramme
$$
\xymatrix{
g^{-1}h_V \ar[r] \ar@{..>}[d] &
\prod\nolimits_{i = 1, \ldots, n} g^{-1}\mathcal{O}_Y
\ar[d]^{\prod g^\sharp} \\
h_{X \times_Y V} \ar[r] &
\prod\nolimits_{i = 1, \ldots, n} \mathcal{O}_X
}
$$
Le lemme est donc d\'emontr\'e.
\end{proof}

\noindent
Voici une version pour des sch\'emas quelconques.

\begin{theorem}
\label{etale-cohomology-theorem-fully-faithful}
Soient $X$, $Y$ des sch\'emas. Soit
$$
(g, g^\#) :
(\Sh(X_\etale), \mathcal{O}_X)
\longrightarrow
(\Sh(Y_\etale), \mathcal{O}_Y)
$$
un morphisme de topos localement annel\'es. Alors il existe un unique
morphisme de sch\'emas $f : X \to Y$ tel que
$(g, g^\#)$ soit isomorphe \`a $(f_{small}, f_{small}^\sharp)$.
Autrement dit, la construction
$$
\Sch \longrightarrow \textit{Topos localement annel\'es},
\quad
X \longrightarrow (X_\etale, \mathcal{O}_X)
$$
est pleinement fid\`ele (les morphismes du membre de droite sont consid\'er\'es \`a $2$-isomorphisme pr\`es).
\end{theorem}

\begin{proof}
On peut d\'emontrer ce th\'eor\`eme en adaptant soigneusement les arguments de
la d\'emonstration du
Lemme \ref{etale-cohomology-lemma-morphism-ringed-etale-topoi-affines}
au cadre global. Nous voulons toutefois indiquer comment
recoller le r\'esultat de ce lemme pour obtenir un morphisme global,
gr\^ace \`a la rigidit\'e fournie par le r\'esultat du
Lemme \ref{etale-cohomology-lemma-2-morphism}.
Malheureusement, les d\'etails sont quelque peu laborieux.

\medskip\noindent
Montrons l'existence lorsque $Y$ est affine. Choisissons alors un
recouvrement ouvert affine $X = \bigcup U_i$. Pour chaque $i$, le morphisme
d'inclusion $j_i : U_i \to X$ induit un morphisme de topos localement annel\'es
$(j_{i, small}, j_{i, small}^\sharp) :
(\Sh(U_{i, \etale}), \mathcal{O}_{U_i})
\to
(\Sh(X_\etale), \mathcal{O}_X)$
d'apr\`es le
Lemme \ref{etale-cohomology-lemma-morphism-locally-ringed}.
Nous pouvons le composer avec $(g, g^\sharp)$ pour obtenir un morphisme
de topos localement annel\'es
$$
(g, g^\sharp) \circ (j_{i, small}, j_{i, small}^\sharp) :
(\Sh(U_{i, \etale}), \mathcal{O}_{U_i})
\to
(\Sh(Y_\etale), \mathcal{O}_Y)
$$
voir
Modules sur les sites,
Lemme \ref{sites-modules-lemma-composition-morphisms-locally-ringed-topoi}.
D'apr\`es le
Lemme \ref{etale-cohomology-lemma-morphism-ringed-etale-topoi-affines},
il existe un unique morphisme de sch\'emas $f_i : U_i \to Y$
et un $2$-isomorphisme
$$
t_i :
(f_{i, small}, f_{i, small}^\sharp)
\longrightarrow
(g, g^\sharp) \circ (j_{i, small}, j_{i, small}^\sharp).
$$
Posons $U_{i, i'} = U_i \cap U_{i'}$, et notons $j_{i, i'} : U_{i, i'} \to U_i$
le morphisme d'inclusion. Puisque
$j_i \circ j_{i, i'} = j_{i'} \circ j_{i', i}$
nous avons
\begin{align*}
(g, g^\sharp) \circ
(j_{i, small}, j_{i, small}^\sharp) \circ
(j_{i, i', small}, j_{i, i', small}^\sharp)
= \\
(g, g^\sharp) \circ
(j_{i', small}, j_{i', small}^\sharp) \circ
(j_{i', i, small}, j_{i', i, small}^\sharp)
\end{align*}
Par unicit\'e (voir le
Lemme \ref{etale-cohomology-lemma-faithful}),
nous en concluons que
$f_i \circ j_{i, i'} = f_{i'} \circ j_{i', i}$ ; autrement dit, les
morphismes de sch\'emas $f_i = f \circ j_i$ sont les restrictions d'un
morphisme de sch\'emas global $f : X \to Y$. Consid\'erons le diagramme
de $2$-isomorphismes (dans lequel nous omettons les composantes ${}^\sharp$ pour all\'eger la
notation)
$$
\xymatrix{
g \circ j_{i, small} \circ j_{i, i', small}
\ar[rr]^{t_i \star \text{id}_{j_{i, i', small}}}
\ar@{=}[d] & &
f_{small} \circ j_{i, small} \circ j_{i, i', small} \ar@{=}[d] \\
g \circ j_{i', small} \circ j_{i', i, small}
\ar[rr]^{t_{i'} \star \text{id}_{j_{i', i, small}}} & &
f_{small} \circ j_{i', small} \circ j_{i', i, small}
}
$$
La notation $\star$ d\'esigne la composition horizontale ; voir
Catégories, D\'efinition \ref{categories-definition-2-category}
dans le cas g\'en\'eral et
Sites, section \ref{sites-section-2-category}
dans notre cas particulier. D'apr\`es le r\'esultat du
Lemme \ref{etale-cohomology-lemma-2-morphism},
ce diagramme est commutatif. Ainsi, pour tout faisceau $\mathcal{G}$
sur $Y_\etale$, les isomorphismes
$t_i : f_{small}^{-1}\mathcal{G}|_{U_i} \to g^{-1}\mathcal{G}|_{U_i}$
co\"incident au-dessus de $U_{i, i'}$ et nous obtenons un isomorphisme global
$t : f_{small}^{-1}\mathcal{G} \to g^{-1}\mathcal{G}$.
Il est clair que cet isomorphisme est fonctoriel en $\mathcal{G}$
et compatible avec les morphismes $f_{small}^\sharp$ et $g^\sharp$
(puisqu'il est localement compatible avec ces morphismes).
Cela d\'emontre le th\'eor\`eme lorsque $Y$ est affine.

\medskip\noindent
Dans le cas g\'en\'eral, soit $V \subset Y$ un ouvert affine.
Alors $h_V$ est un sous-faisceau du faisceau final $*$ sur $Y_\etale$.
Comme $g$ est exact, nous voyons que $g^{-1}h_V$ est un sous-faisceau du faisceau
final sur $X_\etale$. Ainsi, d'apr\`es le
Lemme \ref{etale-cohomology-lemma-support-subsheaf-final},
il existe un sous-sch\'ema ouvert $W \subset X$ tel que $g^{-1}h_V = h_W$. D'apr\`es
Modules sur les sites,
Lemme \ref{sites-modules-lemma-localize-morphism-locally-ringed-topoi},
il existe un diagramme commutatif de morphismes de topos localement
annel\'es
$$
\xymatrix{
(\Sh(W_\etale), \mathcal{O}_W) \ar[r] \ar[d]_{g'} &
(\Sh(X_\etale), \mathcal{O}_X) \ar[d]^g \\
(\Sh(V_\etale), \mathcal{O}_V) \ar[r] &
(\Sh(Y_\etale), \mathcal{O}_Y)
}
$$
o\`u les fl\`eches horizontales sont les morphismes de localisation
(induits par les morphismes d'inclusion $V \to Y$ et $W \to X$)
et o\`u $g'$ est induit par $g$. Le paragraphe pr\'ec\'edent
nous fournit un morphisme de sch\'emas $f' : W \to V$ et
un $2$-isomorphisme
$t : (f'_{small}, (f'_{small})^\sharp) \to (g', (g')^\sharp)$.
Exactement comme pr\'ec\'edemment, ces morphismes $f'$ (pour les divers ouverts affines $V \subset Y$)
co\"incident sur les intersections par unicit\'e, et donnent donc un morphisme $f : X \to Y$.
De plus, les $2$-isomorphismes $t$ sont compatibles sur les intersections, encore d'apr\`es le
Lemme \ref{etale-cohomology-lemma-2-morphism},
et nous obtenons un $2$-isomorphisme global
$(f_{small}, (f_{small})^\sharp) \to (g, (g)^\sharp)$.
comme voulu. Nous omettons certains d\'etails.
\end{proof}










\section{Image directe et image inverse}
\label{etale-cohomology-section-monomorphisms}

\noindent
Soit $f : X \to Y$ un morphisme de sch\'emas.
Voici une liste de conditions que nous consid\'ererons dans la suite :
\begin{enumerate}
\item[(A)] Pour tout morphisme \'etale $U \to X$ et tout $u \in U$, il existe
un morphisme \'etale $V \to Y$, une d\'ecomposition en union disjointe
$X \times_Y V = W \amalg W'$ et un morphisme $h : W \to U$ au-dessus de $X$
tel que $u$ appartienne \`a l'image de $h$.
\item[(B)] Pour tout morphisme \'etale $V \to Y$ et tout recouvrement \'etale
$\{U_i \to X \times_Y V\}$, il existe un recouvrement \'etale
$\{V_j \to V\}$ tel que, pour chaque $j$, on ait
$X \times_Y V_j = \coprod W_{ij}$, o\`u $W_{ij} \to X \times_Y V$
se factorise par $U_i \to X \times_Y V$ pour un certain $i$.
\item[(C)] Pour tout morphisme \'etale $U \to X$, il existe un morphisme \'etale $V \to Y$
et un morphisme surjectif $X \times_Y V \to U$ au-dessus de $X$.
\end{enumerate}
Il s'av\`ere que chacune de ces propri\'et\'es s'interpr\`ete en termes du
comportement du foncteur $f_{small, *}$. Nous expliciterons cela
dans les quelques sections suivantes.



\section{Propri\'et\'e (A)}
\label{etale-cohomology-section-A}

\noindent
Voir la section \ref{etale-cohomology-section-monomorphisms} pour la d\'efinition de la propri\'et\'e
(A).

\begin{lemma}
\label{etale-cohomology-lemma-property-A-implies}
Soit $f : X \to Y$ un morphisme de sch\'emas.
Supposons (A).
\begin{enumerate}
\item
$f_{small, *} :
\textit{Ab}(X_\etale)
\to
\textit{Ab}(Y_\etale)$
refl\`ete les injections et les surjections,
\item $f_{small}^{-1}f_{small, *}\mathcal{F} \to \mathcal{F}$
est surjectif pour tout faisceau ab\'elien $\mathcal{F}$ sur $X_\etale$,
\item
$f_{small, *} :
\textit{Ab}(X_\etale)
\to
\textit{Ab}(Y_\etale)$
est fid\`ele.
\end{enumerate}
\end{lemma}

\begin{proof}
Soit $\mathcal{F}$ un faisceau ab\'elien sur $X_\etale$.
Soit $U$ un objet de $X_\etale$. Par hypoth\`ese, nous pouvons trouver un
recouvrement $\{W_i \to U\}$ dans $X_\etale$ tel que chaque $W_i$ soit
un sous-sch\'ema ouvert et ferm\'e de $X \times_Y V_i$ pour un certain objet
$V_i$ de $Y_\etale$. La condition de faisceau montre que
$$
\mathcal{F}(U) \subset \prod \mathcal{F}(W_i)
$$
et que $\mathcal{F}(W_i)$ est un facteur direct de
$\mathcal{F}(X \times_Y V_i) = f_{small, *}\mathcal{F}(V_i)$.
Il est donc clair que $f_{small, *}$ refl\`ete les injections.

\medskip\noindent
Supposons ensuite que $a : \mathcal{G} \to \mathcal{F}$ soit un morphisme de
faisceaux ab\'eliens tel que $f_{small, *}a$ soit surjectif. Soit
$s \in \mathcal{F}(U)$, avec $U$ comme ci-dessus. Avec les $W_i$, $V_i$ ci-dessus,
nous voyons qu'il suffit de montrer que $s|_{W_i}$ est, localement pour la topologie \'etale,
l'image d'une section de $\mathcal{G}$ par $a$. Comme $\mathcal{F}(W_i)$
est un facteur direct de $\mathcal{F}(X \times_Y V_i)$,
il suffit de montrer que, pour tout $V \in \Ob(Y_\etale)$,
tout \'el\'ement $s \in \mathcal{F}(X \times_Y V)$
est, localement pour la topologie \'etale sur $X \times_Y V$, l'image d'une section de
$\mathcal{G}$ par $a$. Puisque
$\mathcal{F}(X \times_Y V) = f_{small, *}\mathcal{F}(V)$
nous voyons par hypoth\`ese qu'il existe un recouvrement $\{V_j \to V\}$ tel que
$s$ soit l'image de
$s_j \in f_{small, *}\mathcal{G}(V_j) = \mathcal{G}(X \times_Y V_j)$.
Cela prouve que $f_{small, *}$ refl\`ete les surjections.

\medskip\noindent
Les assertions (2) et (3) d\'ecoulent formellement de l'assertion (1) ; voir
Modules sur les sites, Lemme \ref{sites-modules-lemma-reflect-surjections}.
\end{proof}

\begin{lemma}
\label{etale-cohomology-lemma-locally-quasi-finite-A}
Soit $f : X \to Y$ un morphisme de sch\'emas s\'epar\'e et localement quasi-fini.
Alors la propri\'et\'e (A) ci-dessus est satisfaite.
\end{lemma}

\begin{proof}
Soient $U \to X$ un morphisme \'etale et $u \in U$.
L'\'enonc\'e g\'eom\'etrique (A) se r\'eduit directement au cas o\`u $U$ et $Y$
sont des sch\'emas affines. Notons $x \in X$ et $y \in Y$ les
images de $u$. Comme $X \to Y$ est localement quasi-fini et que $U \to X$ est
localement quasi-fini (voir
Morphismes, Lemme \ref{morphisms-lemma-etale-locally-quasi-finite}),
nous voyons que $U \to Y$ est localement quasi-fini (voir
Morphismes, Lemme \ref{morphisms-lemma-composition-quasi-finite}).
De plus, $X \to Y$ et $U \to Y$ sont tous deux s\'epar\'es. Ainsi, le
Compléments sur les morphismes, Lemme
\ref{more-morphisms-lemma-etale-splits-off-quasi-finite-part-technical-variant}
s'applique aux deux morphismes. Nous pouvons donc choisir un voisinage \'etale
$(V, v) \to (Y, y)$ tel que
$$
X \times_Y V = W \amalg R, \quad
U \times_Y V = W' \amalg R'
$$
et des points $w \in W$, $w' \in W'$ tels que
\begin{enumerate}
\item $W$, $R$ sont ouverts et ferm\'es dans $X \times_Y V$,
\item $W'$, $R'$ sont ouverts et ferm\'es dans $U \times_Y V$,
\item $W \to V$ et $W' \to V$ sont finis,
\item $w$, $w'$ s'envoient sur $v$,
\item $\kappa(v) \subset \kappa(w)$ et $\kappa(v) \subset \kappa(w')$
sont purement ins\'eparables, et
\item aucun autre point de $W$ ou de $W'$ ne s'envoie sur $v$.
\end{enumerate}
Voici un diagramme commutatif
$$
\xymatrix{
U \ar[d] & U \times_Y V \ar[l] \ar[d] & W' \amalg R' \ar[d] \ar[l] \\
X \ar[d] & X \times_Y V \ar[l] \ar[d] & W \amalg R \ar[l] \\
Y & V \ar[l]
}
$$
Apr\`es avoir r\'etr\'eci $V$, nous pouvons supposer que $W'$ s'envoie dans $W$ :
il suffit d'enlever l'image de l'image inverse de $R$ dans $W'$ ; c'est
un ferm\'e (puisque $W' \to V$ est fini) qui ne contient pas $v$.
Alors $W' \to W$ est fini, puisque $W \to V$ et $W' \to V$ sont finis.
Ainsi $W' \to W$ est fini \'etale, et la
fibre au-dessus de $w$ contient exactement un point, avec $\kappa(w) = \kappa(w')$. Par cons\'equent, $W' \to W$ est un
isomorphisme sur un voisinage ouvert $W^\circ$ de $w$ ; voir
Morphismes étales, Lemme \ref{etale-lemma-finite-etale-one-point}.
Comme $W \to V$ est fini, l'image de $W \setminus W^\circ$ est un sous-ensemble
ferm\'e $T$ de $V$ qui ne contient pas $v$. Apr\`es avoir remplac\'e $V$ par
$V \setminus T$, nous pouvons donc supposer que $W' \to W$ est un isomorphisme.
La d\'ecomposition $X \times_Y V = W \amalg R$ et le morphisme
$W \to U$ sont alors ceux recherch\'es, ce qui ach\`eve la d\'emonstration.
\end{proof}

\begin{lemma}
\label{etale-cohomology-lemma-integral-A}
Soit $f : X \to Y$ un morphisme entier de sch\'emas.
Alors la propri\'et\'e (A) est satisfaite.
\end{lemma}

\begin{proof}
Soient $U \to X$ \'etale et $u \in U$ un point.
Nous devons trouver un morphisme \'etale $V \to Y$, une d\'ecomposition en union disjointe
$X \times_Y V = W \amalg W'$ et un $X$-morphisme $W \to U$
tel que $u$ appartienne \`a son image. Nous pouvons r\'etr\'ecir $U$ et $Y$ et supposer
que $U$ et $Y$ sont affines. Alors $X$ est \'egalement affine, puisqu'un
morphisme entier est affine par d\'efinition. \'Ecrivons $Y = \Spec(A)$,
$X = \Spec(B)$ et $U = \Spec(C)$. Alors $A \to B$ est un
homomorphisme d'anneaux entier, et $B \to C$ un homomorphisme d'anneaux \'etale. D'apr\`es
Algèbre, Lemme \ref{algebra-lemma-etale},
nous pouvons trouver une sous-$A$-alg\`ebre finie $B' \subset B$ et un homomorphisme d'anneaux \'etale
$B' \to C'$ tels que $C = B \otimes_{B'} C'$. La question
se ram\`ene donc au morphisme \'etale
$U' = \Spec(C') \to X' = \Spec(B')$
au-dessus du morphisme fini $X' \to Y$. Dans ce cas, le r\'esultat d\'ecoule du
Lemme \ref{etale-cohomology-lemma-locally-quasi-finite-A}.
\end{proof}

\begin{lemma}
\label{etale-cohomology-lemma-when-push-pull-surjective}
Soit $f : X \to Y$ un morphisme de sch\'emas. Notons
$f_{small} :
\Sh(X_\etale)
\to
\Sh(Y_\etale)$
le morphisme associ\'e de petits topos \'etales. Supposons qu'au moins l'une
des hypoth\`eses suivantes soit satisfaite :
\begin{enumerate}
\item $f$ est entier, ou
\item $f$ est s\'epar\'e et localement quasi-fini.
\end{enumerate}
Alors le foncteur
$f_{small, *} :
\textit{Ab}(X_\etale)
\to
\textit{Ab}(Y_\etale)$
poss\`ede les propri\'et\'es suivantes :
\begin{enumerate}
\item le morphisme
$f_{small}^{-1}f_{small, *}\mathcal{F} \to \mathcal{F}$
est toujours surjectif,
\item $f_{small, *}$ est fid\`ele, et
\item $f_{small, *}$ refl\`ete les injections et les surjections.
\end{enumerate}
\end{lemma}

\begin{proof}
Il suffit de combiner les
Lemmes \ref{etale-cohomology-lemma-locally-quasi-finite-A},
\ref{etale-cohomology-lemma-integral-A}, et
\ref{etale-cohomology-lemma-property-A-implies}.
\end{proof}



\section{Propri\'et\'e (B)}
\label{etale-cohomology-section-B}

\noindent
Voir la section \ref{etale-cohomology-section-monomorphisms} pour la d\'efinition de la propri\'et\'e
(B).

\begin{lemma}
\label{etale-cohomology-lemma-property-B-implies}
Soit $f : X \to Y$ un morphisme de sch\'emas. Supposons la propri\'et\'e (B) satisfaite.
Alors le foncteur
$f_{small, *} :
\Sh(X_\etale)
\to
\Sh(Y_\etale)$
envoie les morphismes surjectifs sur des morphismes surjectifs.
\end{lemma}

\begin{proof}
Cela d\'ecoule de
Sites, Lemme \ref{sites-lemma-weaker}.
\end{proof}

\begin{lemma}
\label{etale-cohomology-lemma-simplify-B}
Soit $f : X \to Y$ un morphisme de sch\'emas. Supposons que
\begin{enumerate}
\item $V \to Y$ soit un morphisme \'etale de sch\'emas,
\item $\{U_i \to X \times_Y V\}$ soit un recouvrement \'etale, et
\item $v \in V$ soit un point.
\end{enumerate}
Supposons que, pour toute donn\'ee de ce type, il existe un voisinage \'etale
$(V', v') \to (V, v)$, une d\'ecomposition en union disjointe
$X \times_Y V' = \coprod W'_i$ et des morphismes $W'_i \to U_i$
au-dessus de $X \times_Y V$. Alors la propri\'et\'e (B) est satisfaite.
\end{lemma}

\begin{proof}
D\'emonstration omise.
\end{proof}

\begin{lemma}
\label{etale-cohomology-lemma-finite-B}
Soit $f : X \to Y$ un morphisme fini de sch\'emas.
Alors la propri\'et\'e (B) est satisfaite.
\end{lemma}

\begin{proof}
Consid\'erons un morphisme \'etale $V \to Y$, un recouvrement \'etale $\{U_i \to X \times_Y V\}$ et
un point $v \in V$. Nous devons trouver $V' \to V$, une d\'ecomposition et des morphismes comme dans le
Lemme \ref{etale-cohomology-lemma-simplify-B}.
Nous pouvons r\'etr\'ecir $V$ et $Y$, donc supposer que $V$ et $Y$ sont affines.
Comme $X$ est fini sur $Y$, il en r\'esulte aussi que $X$ est affine.
Au cours de la d\'emonstration, nous pouvons remplacer (un nombre fini de fois) $(V, v)$ par un
voisinage \'etale $(V', v')$ et, de mani\`ere correspondante, le recouvrement
$\{U_i \to X \times_Y V\}$ par $\{V' \times_V U_i \to X \times_Y V'\}$.

\medskip\noindent
Comme $X \times_Y V \to V$ est fini, il existe un nombre fini de
points deux \`a deux distincts $x_1, \ldots, x_n \in X \times_Y V$
qui s'envoient sur $v$. Nous pouvons appliquer
Compléments sur les morphismes, Lemme
\ref{more-morphisms-lemma-etale-splits-off-quasi-finite-part-technical-variant}
\`a $X \times_Y V \to V$ et aux points $x_1, \ldots, x_n$ situ\'es au-dessus de
$v$, et trouver un voisinage \'etale $(V', v') \to (V, v)$
tel que
$$
X \times_Y V' = R \amalg \coprod T_a
$$
o\`u $T_a \to V'$ est fini, avec exactement un point $p_a$ au-dessus de $v'$,
o\`u de plus $\kappa(v') \subset \kappa(p_a)$ est purement ins\'eparable, et
tel que la fibre de $R \to V'$ au-dessus de $v'$ soit vide.
Puisque $X \to Y$ est fini, $R \to V'$ l'est aussi. Apr\`es avoir
r\'etr\'eci $V'$, nous pouvons donc supposer que $R = \emptyset$. Ainsi, nous pouvons
supposer que $X \times_Y V = X_1 \amalg \ldots \amalg X_n$, avec
exactement un point $x_l \in X_l$ au-dessus de $v$, et tel que, de plus,
$\kappa(v) \subset \kappa(x_l)$ soit purement ins\'eparable. Remarquons que cette
propri\'et\'e est pr\'eserv\'ee par tout raffinement du voisinage \'etale
$(V, v)$.

\medskip\noindent
Pour chaque $l$, choisissons un $i_l$ et un point $u_l \in U_{i_l}$ qui s'envoie sur $x_l$.
Appliquons maintenant la propri\'et\'e (A) au morphisme fini
$X \times_Y V \to V$, aux morphismes \'etales
$U_{i_l} \to X \times_Y V$ et aux points $u_l$.
Cela est permis par le
Lemme \ref{etale-cohomology-lemma-integral-A}.
Nous obtenons ainsi un voisinage \'etale $(V', v') \to (V, v)$
et des d\'ecompositions
$$
X \times_Y V' = W_l \amalg R_l
$$
et des $X$-morphismes $a_l : W_l \to U_{i_l}$ dont l'image contient $u_{i_l}$.
Voici une figure :
$$
\xymatrix{
& & & U_{i_l} \ar[d] & \\
W_l \ar[rrru] \ar[r] & W_l \amalg R_l \ar@{=}[r] &
X \times_Y V' \ar[r] \ar[d] &
X \times_Y V \ar[r] \ar[d] & X \ar[d] \\
& & V' \ar[r] & V \ar[r] & Y
}
$$
Apr\`es avoir remplac\'e $(V, v)$ par $(V', v')$, nous concluons que chaque
$x_l$ appartient \`a un voisinage ouvert et ferm\'e $W_l$ tel que
le morphisme d'inclusion $W_l \to X \times_Y V$ se factorise par
$U_i \to X \times_Y V$ pour un certain $i$. En rempla\c{c}ant $W_l$ par $W_l \cap X_l$,
nous voyons que ces ouverts et ferm\'es sont disjoints et, de plus,
que $\{x_1, \ldots, x_n\} \subset W_1 \cup \ldots \cup W_n$.
Comme $X \times_Y V \to V$ est fini, nous pouvons r\'etr\'ecir $V$ et supposer que
$X \times_Y V = W_1 \amalg \ldots \amalg W_n$, comme voulu.
\end{proof}

\begin{lemma}
\label{etale-cohomology-lemma-integral-B}
Soit $f : X \to Y$ un morphisme entier de sch\'emas.
Alors la propri\'et\'e (B) est satisfaite.
\end{lemma}

\begin{proof}
Consid\'erons un morphisme \'etale $V \to Y$, un recouvrement \'etale $\{U_i \to X \times_Y V\}$ et
un point $v \in V$. Nous devons trouver $V' \to V$, une d\'ecomposition et des morphismes comme dans le
Lemme \ref{etale-cohomology-lemma-simplify-B}.
Nous pouvons r\'etr\'ecir $V$ et $Y$, donc supposer que $V$ et $Y$ sont affines.
Comme $X$ est entier sur $Y$, il en r\'esulte aussi que $X$ et
$X \times_Y V$ sont affines. Nous pouvons raffiner le recouvrement
$\{U_i \to X \times_Y V\}$ et donc supposer que
$\{U_i \to X \times_Y V\}_{i = 1, \ldots, n}$ est un recouvrement \'etale standard.
\'Ecrivons $Y = \Spec(A)$, $X = \Spec(B)$,
$V = \Spec(C)$ et $U_i = \Spec(B_i)$.
Alors $A \to B$ est un homomorphisme d'anneaux entier, et les $B \otimes_A C \to B_i$ sont des
homomorphismes d'anneaux \'etales. D'apr\`es
Algèbre, Lemme \ref{algebra-lemma-etale},
nous pouvons trouver une sous-$A$-alg\`ebre finie $B' \subset B$ et un homomorphisme d'anneaux \'etale
$B' \otimes_A C \to B'_i$, pour $i = 1, \ldots, n$,
tels que $B_i = B \otimes_{B'} B'_i$. La question
se ram\`ene donc au recouvrement \'etale
$\{\Spec(B'_i) \to X' \times_Y V\}_{i = 1, \ldots, n}$
o\`u $X' = \Spec(B')$ est fini sur $Y$.
Dans ce cas, le r\'esultat d\'ecoule du
Lemme \ref{etale-cohomology-lemma-finite-B}.
\end{proof}

\begin{lemma}
\label{etale-cohomology-lemma-what-integral}
Soit $f : X \to Y$ un morphisme de sch\'emas.
Supposons $f$ entier (par exemple fini).
Alors
\begin{enumerate}
\item $f_{small, *}$ envoie les morphismes surjectifs sur des morphismes surjectifs (pour les faisceaux
d'ensembles et pour les faisceaux ab\'eliens),
\item $f_{small}^{-1}f_{small, *}\mathcal{F} \to \mathcal{F}$
est surjectif pour tout faisceau ab\'elien $\mathcal{F}$ sur $X_\etale$,
\item
$f_{small, *} :
\textit{Ab}(X_\etale)
\to
\textit{Ab}(Y_\etale)$
est fid\`ele et refl\`ete les injections et les surjections, et
\item
$f_{small, *} :
\textit{Ab}(X_\etale)
\to
\textit{Ab}(Y_\etale)$
est exact.
\end{enumerate}
\end{lemma}

\begin{proof}
Nous avons vu les assertions (2) et (3) dans le
Lemme \ref{etale-cohomology-lemma-when-push-pull-surjective}.
L'assertion (1) d\'ecoule des
Lemmes \ref{etale-cohomology-lemma-integral-B} et \ref{etale-cohomology-lemma-property-B-implies}.
L'assertion (4) est une cons\'equence de l'assertion (1) ; voir
Modules sur les sites, Lemme \ref{sites-modules-lemma-exactness}.
\end{proof}








\section{Propri\'et\'e (C)}
\label{etale-cohomology-section-C}

\noindent
Voir la section \ref{etale-cohomology-section-monomorphisms} pour la d\'efinition de la propri\'et\'e
(C).

\begin{lemma}
\label{etale-cohomology-lemma-property-C-implies}
Soit $f : X \to Y$ un morphisme de sch\'emas. Supposons (C) satisfaite. Alors le foncteur
$f_{small, *} :
\Sh(X_\etale)
\to
\Sh(Y_\etale)$
refl\`ete les injections et les surjections.
\end{lemma}

\begin{proof}
Cela d\'ecoule de
Sites, Lemme \ref{sites-lemma-cover-from-below}.
Nous omettons la v\'erification du fait que la propri\'et\'e (C) implique que le foncteur
$Y_\etale \to X_\etale$, $V \mapsto X \times_Y V$
satisfait l'hypoth\`ese de
Sites, Lemme \ref{sites-lemma-cover-from-below}.
\end{proof}

\begin{remark}
\label{etale-cohomology-remark-property-C-strong}
La propri\'et\'e (C) est satisfaite si $f : X \to Y$ est une immersion ouverte. En effet, si
$U \in \Ob(X_\etale)$, nous pouvons aussi consid\'erer $U$ comme un objet
de $Y_\etale$ et $U \times_Y X = U$. Ainsi, la propri\'et\'e (C)
n'implique pas que $f_{small, *}$ soit exact, car ce n'est pas
le cas pour les immersions ouvertes (en g\'en\'eral).
\end{remark}

\begin{lemma}
\label{etale-cohomology-lemma-property-C-closed-implies}
Soit $f : X \to Y$ un morphisme de sch\'emas. Supposons que,
pour tout morphisme \'etale $V \to Y$, on ait
\begin{enumerate}
\item $X \times_Y V \to V$ poss\`ede la propri\'et\'e (C), et
\item $X \times_Y V \to V$ est ferm\'e.
\end{enumerate}
Alors le foncteur
$Y_\etale \to X_\etale$, $V \mapsto X \times_Y V$
est presque cocontinu ; voir
Sites, D\'efinition \ref{sites-definition-almost-cocontinuous}.
\end{lemma}

\begin{proof}
Soit $V \to Y$ un objet de $Y_\etale$ et soit
$\{U_i \to X \times_Y V\}_{i \in I}$ un recouvrement de $X_\etale$.
Par l'hypoth\`ese (1), pour chaque $i$, nous pouvons trouver un morphisme \'etale
$h_i : V_i \to V$ et un morphisme surjectif $X \times_Y V_i \to U_i$
au-dessus de $X \times_Y V$. Remarquons que $\bigcup h_i(V_i) \subset V$ est un
ouvert qui contient le ferm\'e $Z = \Im(X \times_Y V \to V)$.
Soit $h_0 : V_0 = V \setminus Z \to V$ l'immersion ouverte.
Il est clair que $\{V_i \to V\}_{i \in I \cup \{0\}}$ est un
recouvrement \'etale tel que, pour tout $i \in I \cup \{0\}$, on ait
soit $V_i \times_Y X = \emptyset$ (si $i = 0$), soit
que $V_i \times_Y X \to V \times_Y X$ se factorise par $U_i \to X \times_Y V$
(si $i \not = 0$). Le foncteur $Y_\etale \to X_\etale$ est donc
presque cocontinu.
\end{proof}

\begin{lemma}
\label{etale-cohomology-lemma-integral-homeo-onto-image-C}
Soit $f : X \to Y$ un morphisme entier de sch\'emas qui d\'efinit
un hom\'eomorphisme de $X$ sur un sous-ensemble ferm\'e de $Y$.
Alors la propri\'et\'e (C) est satisfaite.
\end{lemma}

\begin{proof}
Soit $g : U \to X$ un morphisme \'etale. Nous devons trouver un objet
$V \to Y$ de $Y_\etale$ et un morphisme surjectif $X \times_Y V \to U$
au-dessus de $X$. Supposons que, pour tout $u \in U$, nous puissions trouver un objet
$V_u \to Y$ de $Y_\etale$ et un morphisme $h_u : X \times_Y V_u \to U$
au-dessus de $X$ avec $u \in \Im(h_u)$. Nous pouvons alors prendre $V = \coprod V_u$
et $h = \coprod h_u$, ce qui donne le r\'esultat. Ainsi, \'etant donn\'e un point
$u \in U$, construisons une paire $(V_u, h_u)$ comme ci-dessus. Pour ce faire, nous pouvons
r\'etr\'ecir $U$ et supposer que $U$ est affine. Dans ce cas,
$g : U \to X$ est localement quasi-fini. Posons
$g^{-1}(g(\{u\})) = \{u, u_2, \ldots, u_n\}$. Puisqu'il n'existe aucune
sp\'ecialisations $u_i \leadsto u$, nous pouvons remplacer $U$ par un voisinage affine
de sorte que $g^{-1}(g(\{u\})) = \{u\}$.

\medskip\noindent
L'image $g(U) \subset X$ est ouverte ;
par cons\'equent, $f(g(U))$ est localement ferm\'e dans $Y$. Choisissons un ouvert $V \subset Y$ tel
que $f(g(U)) = f(X) \cap V$. Il en r\'esulte que $g$ se factorise par
$X \times_Y V$ et que le morphisme $\{U \to X \times_Y V\}$ ainsi obtenu est un recouvrement
\'etale. Comme $f$ poss\`ede la propri\'et\'e (B), voir le
Lemme \ref{etale-cohomology-lemma-integral-B},
nous voyons qu'il existe un recouvrement \'etale $\{V_j \to V\}$ tel que
$X \times_Y V_j \to X \times_Y V$ se factorise par $U$.
Il en r\'esulte que $V' = \coprod V_j$ est \'etale sur $Y$ et qu'il existe un
morphisme $h : X \times_Y V' \to U$ dont l'image
se projette surjectivement sur $g(U)$. Comme $u$ est le seul point de sa fibre, il doit
appartenir \`a l'image de $h$, ce qui donne le r\'esultat.
\end{proof}

\noindent
Nous invitons le lecteur \`a consid\'erer le lemme suivant comme une
\'etape\footnote{Une \'etape est un lieu o\`u l'on s'arr\^ete pour manger
et se reposer au cours d'un long voyage.} sur le chemin qui m\`ene au
r\'esultat ultime concernant $f_{small, *}$ pour les morphismes entiers et universellement
injectifs.

\begin{lemma}
\label{etale-cohomology-lemma-integral-universally-injective}
Soit $f : X \to Y$ un morphisme de sch\'emas. Supposons que $f$ soit
universellement injectif et entier (par exemple une immersion ferm\'ee).
Alors
\begin{enumerate}
\item
$f_{small, *} :
\Sh(X_\etale)
\to
\Sh(Y_\etale)$
refl\`ete les injections et les surjections,
\item
$f_{small, *} :
\Sh(X_\etale)
\to
\Sh(Y_\etale)$
commute aux sommes amalgam\'ees et aux co\'egalisateurs (et, plus g\'en\'eralement,
aux colimites connexes finies),
\item $f_{small, *}$ envoie les morphismes surjectifs sur des morphismes surjectifs (pour les faisceaux
d'ensembles et pour les faisceaux ab\'eliens),
\item le morphisme
$f_{small}^{-1}f_{small, *}\mathcal{F} \to \mathcal{F}$
est surjectif pour tout faisceau (d'ensembles ou de groupes ab\'eliens)
$\mathcal{F}$ sur $X_\etale$,
\item le foncteur $f_{small, *}$ est fid\`ele (sur les faisceaux d'ensembles et
sur les faisceaux ab\'eliens),
\item
$f_{small, *} :
\textit{Ab}(X_\etale)
\to
\textit{Ab}(Y_\etale)$
est exact, et
\item le foncteur
$Y_\etale \to X_\etale$, $V \mapsto X \times_Y V$ est
presque cocontinu.
\end{enumerate}
\end{lemma}

\begin{proof}
D'apr\`es les
Lemmes \ref{etale-cohomology-lemma-integral-A},
\ref{etale-cohomology-lemma-integral-B} et
\ref{etale-cohomology-lemma-integral-homeo-onto-image-C},
nous savons que le morphisme $f$ poss\`ede les propri\'et\'es (A), (B) et (C).
De plus, d'apr\`es le
Lemme \ref{etale-cohomology-lemma-property-C-closed-implies},
nous savons que le foncteur $Y_\etale \to X_\etale$ est
presque cocontinu. Nous avons maintenant :
\begin{enumerate}
\item la propri\'et\'e (C) implique (1) d'apr\`es le
Lemme \ref{etale-cohomology-lemma-property-C-implies},
\item la presque cocontinuit\'e implique (2) d'apr\`es
Sites, Lemme \ref{sites-lemma-morphism-of-sites-almost-cocontinuous},
\item la propri\'et\'e (B) implique (3) d'apr\`es le
Lemme \ref{etale-cohomology-lemma-property-B-implies}.
\end{enumerate}
Les propri\'et\'es (4), (5) et (6) d\'ecoulent formellement des trois premi\`eres ; voir
Sites, Lemme \ref{sites-lemma-exactness-properties}
et
Modules sur les sites, Lemme \ref{sites-modules-lemma-exactness}.
Nous avons vu la propri\'et\'e (7) ci-dessus.
\end{proof}




\section{Invariance topologique du petit site \'etale}
\label{etale-cohomology-section-topological-invariance}

\noindent
Dans le th\'eor\`eme suivant, nous montrons que le petit site \'etale est un invariant
topologique au sens suivant : si $f : X \to Y$ est un morphisme de sch\'emas
qui est un hom\'eomorphisme universel, alors $X_\etale \cong Y_\etale$
en tant que sites. Cela am\'eliore le r\'esultat de
Morphismes étales, Th\'eor\`eme \ref{etale-theorem-remarkable-equivalence}.
Nous d\'emontrons d'abord le r\'esultat pour les morphismes, puis nous \'enon\c{c}ons celui
pour les cat\'egories.

\begin{theorem}
\label{etale-cohomology-theorem-etale-topological}
Soient $X$ et $Y$ deux sch\'emas sur un sch\'ema de base $S$. Soit
$S' \to S$ un hom\'eomorphisme universel.
Notons $X'$ (resp.\ $Y'$) le changement de base \`a $S'$.
Si $X$ est \'etale sur $S$, alors l'application
$$
\Mor_S(Y, X) \longrightarrow \Mor_{S'}(Y', X')
$$
est bijective.
\end{theorem}

\begin{proof}
Apr\`es changement de base par $Y \to S$, nous pouvons supposer que $Y = S$.
Ainsi, nous pouvons et allons supposer que $X$ et $Y$ sont tous deux \'etales sur $S$.
Autrement dit, le th\'eor\`eme affirme que le foncteur de changement de base
est un foncteur pleinement fid\`ele de la cat\'egorie des sch\'emas \'etales
sur $S$ vers la cat\'egorie des sch\'emas \'etales sur $S'$.

\medskip\noindent
Consid\'erons le foncteur d'oubli
\begin{equation}
\label{etale-cohomology-equation-descent-etale-forget}
\begin{matrix}
\text{donn\'ees de descente }(X', \varphi')\text{ relativement \`a }S'/S \\
\text{ avec }X'\text{ \'etale sur }S'
\end{matrix}
\longrightarrow
\text{sch\'emas }X'\text{ \'etales sur }S'
\end{equation}
Nous affirmons que ce foncteur est une \'equivalence. D'autre part, le
foncteur
\begin{equation}
\label{etale-cohomology-equation-descent-etale}
\text{sch\'emas }X\text{ \'etales sur }S \longrightarrow
\begin{matrix}
\text{donn\'ees de descente }(X', \varphi')\text{ relativement \`a }S'/S \\
\text{ avec }X'\text{ \'etale sur }S'
\end{matrix}
\end{equation}
est pleinement fid\`ele d'apr\`es Morphismes étales, Lemme
\ref{etale-lemma-fully-faithful-cases}.
L'affirmation entra\^ine donc le th\'eor\`eme.

\medskip\noindent
D\'emonstration de l'affirmation.
Rappelons qu'un hom\'eomorphisme universel est la m\^eme chose qu'un morphisme
entier, universellement injectif et surjectif ; voir
Morphismes, Lemme \ref{morphisms-lemma-universal-homeomorphism}.
En particulier, la diagonale $\Delta : S' \to S' \times_S S'$ est un \'epaississement
d'apr\`es Morphismes, Lemme \ref{morphisms-lemma-universally-injective}.
Ainsi, d'apr\`es Morphismes étales, Th\'eor\`eme
\ref{etale-theorem-etale-topological},
nous voyons que, pour tout $X' \to S'$ \'etale, il existe un unique isomorphisme
$$
\varphi' : X' \times_S S' \to S' \times_S X'
$$
de sch\'emas \'etales sur $S' \times_S S'$ dont l'image inverse par
$\Delta$ est $\text{id} : X' \to X'$ sur $S'$.
Comme $S' \to S' \times_S S' \times_S S'$
est aussi un \'epaississement (c'est une immersion ferm\'ee bijective),
nous concluons que $(X', \varphi')$ est une donn\'ee de descente relativement \`a $S'/S$.
Le caract\`ere canonique de la construction de $\varphi'$ montre
qu'elle est compatible aux morphismes entre sch\'emas \'etales sur $S'$.
Autrement dit, nous obtenons un quasi-inverse
$X' \mapsto (X', \varphi')$ du foncteur
(\ref{etale-cohomology-equation-descent-etale-forget}). Cela d\'emontre l'affirmation et
ach\`eve la d\'emonstration du th\'eor\`eme.
\end{proof}

\begin{theorem}
\label{etale-cohomology-theorem-topological-invariance}
\begin{reference}
\cite[IV Theorem 18.1.2]{EGA}
\end{reference}
Soit $f : X \to Y$ un morphisme de sch\'emas.
Supposons $f$ entier, universellement injectif et surjectif
(c'est-\`a-dire que $f$ est un hom\'eomorphisme universel ; voir
Morphismes, Lemme \ref{morphisms-lemma-universal-homeomorphism}).
Le foncteur
$$
V \longmapsto V_X = X \times_Y V
$$
d\'efinit une \'equivalence de cat\'egories
$$
\{
\text{sch\'emas }V\text{ \'etales sur }Y
\}
\leftrightarrow
\{
\text{sch\'emas }U\text{ \'etales sur }X
\}
$$
\end{theorem}

\noindent
Nous donnons deux d\'emonstrations. La premi\`ere utilise l'effectivit\'e de la descente
des morphismes quasi-compacts, s\'epar\'es et \'etales relativement
aux morphismes entiers surjectifs. La seconde utilise les r\'esultats
sur les propri\'et\'es (A), (B) et (C) discut\'es plus haut dans le chapitre.

\begin{proof}[Premi\`ere d\'emonstration]
D'apr\`es le Th\'eor\`eme \ref{etale-cohomology-theorem-etale-topological},
nous voyons que le foncteur est pleinement fid\`ele.
Il reste \`a montrer que le foncteur est essentiellement surjectif.
Soit $U \to X$ un morphisme \'etale de sch\'emas.

\medskip\noindent
Supposons le r\'esultat vrai lorsque $U$ et $Y$ sont affines.
Dans ce cas, choisissons un recouvrement ouvert affine
$U = \bigcup U_i$ tel que chaque $U_i$ s'envoie
dans un ouvert affine de $Y$. Par hypoth\`ese (cas affine), nous pouvons
trouver des morphismes \'etales $V_i \to Y$ tels que $X \times_Y V_i \cong U_i$
comme sch\'emas sur $X$. Soit $V_{i, i'} \subset V_i$
le sous-sch\'ema ouvert dont l'espace topologique sous-jacent
correspond \`a $U_i \cap U_{i'}$. Puisque nous avons des isomorphismes
$$
X \times_Y V_{i, i'} \cong U_i \cap U_{i'} \cong X \times_Y V_{i', i}
$$
comme sch\'emas sur $X$, la pleine fid\'elit\'e montre que nous
obtenons des isomorphismes
$\theta_{i, i'} : V_{i, i'} \to V_{i', i}$ de sch\'emas sur $Y$.
Nous omettons de v\'erifier que ces isomorphismes satisfont \`a la
condition de cocycle de Schémas, section \ref{schemes-section-glueing-schemes}.
En appliquant Schémas, Lemme \ref{schemes-lemma-glue-schemes}
nous obtenons un sch\'ema $V \to Y$ par
recollement des sch\'emas $V_i$ suivant les identifications $\theta_{i, i'}$.
Il est clair que $V \to Y$ est \'etale et que $X \times_Y V \cong U$
par construction.

\medskip\noindent
Ainsi, il suffit de d\'emontrer le lemme dans le cas o\`u $U$ et $Y$ sont affines.
Rappelons que, dans la d\'emonstration du Th\'eor\`eme \ref{etale-cohomology-theorem-etale-topological},
nous avons montr\'e que $U$ est muni d'une unique donn\'ee de descente
$(U, \varphi)$ relativement \`a $X/Y$. D'apr\`es
Morphismes étales, Proposition \ref{etale-proposition-effective}
(qui s'applique car $U \to X$ est quasi-compact et s\'epar\'e
ainsi qu'\'etale, d'apr\`es notre r\'eduction au cas affine),
il existe un morphisme \'etale $V \to Y$ tel que
$X \times_Y V \cong U$, ce qui ach\`eve la d\'emonstration.
\end{proof}

\begin{proof}[Seconde d\'emonstration]
D'apr\`es le Th\'eor\`eme \ref{etale-cohomology-theorem-etale-topological},
nous voyons que le foncteur est pleinement fid\`ele.
Il reste \`a montrer que le foncteur est essentiellement surjectif.
Soit $U \to X$ un morphisme \'etale de sch\'emas.

\medskip\noindent
Supposons le r\'esultat vrai lorsque $U$ et $Y$ sont affines.
Dans ce cas, choisissons un recouvrement ouvert affine
$U = \bigcup U_i$ tel que chaque $U_i$ s'envoie
dans un ouvert affine de $Y$. Par hypoth\`ese (cas affine), nous pouvons
trouver des morphismes \'etales $V_i \to Y$ tels que $X \times_Y V_i \cong U_i$
comme sch\'emas sur $X$. Soit $V_{i, i'} \subset V_i$
le sous-sch\'ema ouvert dont l'espace topologique sous-jacent
correspond \`a $U_i \cap U_{i'}$. Puisque nous avons des isomorphismes
$$
X \times_Y V_{i, i'} \cong U_i \cap U_{i'} \cong X \times_Y V_{i', i}
$$
comme sch\'emas sur $X$, la pleine fid\'elit\'e montre que nous
obtenons des isomorphismes
$\theta_{i, i'} : V_{i, i'} \to V_{i', i}$ de sch\'emas sur $Y$.
Nous omettons de v\'erifier que ces isomorphismes satisfont \`a la
condition de cocycle de Schémas, section \ref{schemes-section-glueing-schemes}.
En appliquant Schémas, Lemme \ref{schemes-lemma-glue-schemes}
nous obtenons un sch\'ema $V \to Y$ par
recollement des sch\'emas $V_i$ suivant les identifications $\theta_{i, i'}$.
Il est clair que $V \to Y$ est \'etale et que $X \times_Y V \cong U$
par construction.

\medskip\noindent
Ainsi, il suffit de d\'emontrer que le foncteur
\begin{equation}
\label{etale-cohomology-equation-affine-etale}
\{
\text{sch\'emas affines }V\text{ \'etales sur }Y
\}
\leftrightarrow
\{
\text{sch\'emas affines }U\text{ \'etales sur }X
\}
\end{equation}
est essentiellement surjectif lorsque $X$ et $Y$ sont affines.

\medskip\noindent
Soit $U \to X$ un sch\'ema affine \'etale sur $X$.
Nous devons trouver $V \to Y$ \'etale (et affine) tel que $X \times_Y V$
soit isomorphe \`a $U$ sur $X$. Remarquons qu'un morphisme \'etale de sch\'emas affines
a des fibres universellement born\'ees, voir
Morphismes,
Lemmes \ref{morphisms-lemma-etale-locally-quasi-finite} et
\ref{morphisms-lemma-locally-quasi-finite-qc-source-universally-bounded}.
Nous pouvons donc raisonner par r\'ecurrence sur l'entier $n$ qui borne le degr\'e des fibres
de $U \to X$. Voir
Morphismes, Lemme \ref{morphisms-lemma-etale-universally-bounded}
pour une description de cet entier dans le cas d'un morphisme \'etale.
Si $n = 1$, alors $U \to X$ est une immersion ouverte (voir
Morphismes étales, Th\'eor\`eme \ref{etale-theorem-etale-radicial-open}),
et le r\'esultat est clair. Supposons $n > 1$.

\medskip\noindent
D'apr\`es
le Lemme \ref{etale-cohomology-lemma-integral-homeo-onto-image-C},
il existe un morphisme \'etale de sch\'emas $W \to Y$ et un
morphisme surjectif $W_X \to U$ sur $X$.
Comme $U$ est quasi-compact, nous pouvons remplacer $W$ par une somme disjointe d'un
nombre fini d'ouverts affines de $W$ ; nous pouvons donc supposer que $W$
est lui aussi affine. Voici un diagramme
$$
\xymatrix{
U \ar[d] & U \times_Y W \ar[l] \ar[d] & W_X \amalg R \ar@{=}[l]\\
X \ar[d] & W_X \ar[l] \ar[d] \\
Y & W \ar[l]
}
$$
La d\'ecomposition en somme disjointe provient du fait que, par construction, le
morphisme \'etale de sch\'emas affines $U \times_Y W \to W_X$ admet une section.
Nous voyons alors que le morphisme $R \to X \times_Y W$ est un morphisme \'etale
de sch\'emas affines dont les fibres ont un degr\'e universellement born\'e
par $n - 1$. Par l'hypoth\`ese de r\'ecurrence, il existe donc un sch\'ema
$V' \to W$ \'etale tel que $R \cong W_X \times_W V'$.
En posant $V'' = W \amalg V'$, nous obtenons un sch\'ema $V''$ \'etale sur $W$ dont le
changement de base \`a $W_X$ est isomorphe \`a $U \times_Y W$
sur $X \times_Y W$.

\medskip\noindent
\`A ce stade, nous pouvons utiliser la descente pour trouver $V$ sur $Y$ dont le
changement de base \`a $X$ est isomorphe \`a $U$ sur $X$. En effet, par la pleine
fid\'elit\'e du foncteur (\ref{etale-cohomology-equation-affine-etale})
correspondant \`a l'hom\'eomorphisme universel
$X \times_Y (W \times_Y W) \to (W \times_Y W)$
il existe un unique isomorphisme $\varphi : V'' \times_Y W \to W \times_Y V''$
dont le changement de base \`a $X \times_Y (W \times_Y W)$ est la donn\'ee de
descente canonique de $U \times_Y W$ sur $X \times_Y W$. En particulier,
$\varphi$ satisfait \`a la condition de cocycle. Ainsi, d'apr\`es
Descente, Lemme \ref{descent-lemma-affine},
nous voyons que $\varphi$ est effective (rappelons que tous les sch\'emas ci-dessus sont affines).
Nous obtenons donc $V \to Y$ et un isomorphisme $V'' \cong W \times_Y V$
tel que la donn\'ee de descente canonique sur $W \times_Y V/W/Y$ co\"incide
avec $\varphi$. Remarquons que $V \to Y$ est \'etale, d'apr\`es
Descente, Lemme \ref{descent-lemma-descending-property-etale}.
De plus, il existe un isomorphisme $V_X \cong U$ obtenu par
descente de l'isomorphisme
$$
V_X \times_X W_X =
X \times_Y V \times_Y W =
(X \times_Y W) \times_W (W \times_Y V) \cong
W_X  \times_W V'' \cong U \times_Y W
$$
dont nous disposons par construction. Nous omettons certains d\'etails.
\end{proof}

\begin{remark}
\label{etale-cohomology-remark-affine-inside-equivalence}
Dans la situation du
Th\'eor\`eme \ref{etale-cohomology-theorem-topological-invariance},
il est \`egalement vrai que $V \mapsto V_X$ induit une \'equivalence
entre les morphismes \'etales $V \to Y$ pour lesquels $V$ est affine et
les morphismes \'etales $U \to X$ pour lesquels $U$ est affine.
Cela r\'esulte par exemple de
Limits, Proposition \ref{limits-proposition-affine}.
\end{remark}

\begin{proposition}[Invariance topologique de la cohomologie \'etale]
\label{etale-cohomology-proposition-topological-invariance}
Soit $X_0 \to X$ un hom\'eomorphisme universel de sch\'emas
(par exemple l'immersion ferm\'ee d\'efinie par un faisceau d'id\'eaux nilpotent).
Alors
\begin{enumerate}
\item les sites \'etales $X_\etale$ et $(X_0)_\etale$ sont isomorphes,
\item les topos \'etales $\Sh(X_\etale)$ et $\Sh((X_0)_\etale)$
sont \'equivalents, et
\item $H^q_\etale(X, \mathcal{F}) = H^q_\etale(X_0, \mathcal{F}|_{X_0})$
pour tout $q$ et
pour tout faisceau ab\'elien $\mathcal{F}$ sur $X_\etale$.
\end{enumerate}
\end{proposition}

\begin{proof}
L'\'equivalence de cat\'egories $X_\etale \to (X_0)_\etale$ est
donn\'ee par le Th\'eor\`eme \ref{etale-cohomology-theorem-topological-invariance}. Nous omettons
de d\'emontrer que, sous cette \'equivalence, les recouvrements \'etales se correspondent.
Ainsi, (1) est d\'emontr\'e. Les assertions (2) et (3) d\'ecoulent formellement de (1).
\end{proof}






\section{Immersions ferm\'ees et image directe}
\label{etale-cohomology-section-closed-immersions}

\noindent
Avant d'\'enoncer et de d\'emontrer
Proposition \ref{etale-cohomology-proposition-closed-immersion-pushforward}
dans toute sa g\'en\'eralit\'e, nous l'\'enon\c{c}ons et la d\'emontrons bri\`evement pour les
immersions ferm\'ees. En effet, certains des arguments pr\'ec\'edents
sont beaucoup plus faciles \`a suivre dans le cas d'une immersion ferm\'ee et
nous les r\'ep\'etons donc ici sous leur forme simplifi\'ee.

\medskip\noindent
Dans la suite de cette section, $i : Z \to X$ est une immersion ferm\'ee.
Le foncteur
$$
\Sch/X \longrightarrow \Sch/Z, \quad
U \longmapsto U_Z = Z \times_X U
$$
sera not\'e $U \mapsto U_Z$, comme indiqu\'e. Comme la propri\'et\'e d'\^etre une immersion ferm\'ee
se conserve par tout changement de base, le sch\'ema $U_Z$ est un sous-sch\'ema ferm\'e
de $U$.

\begin{lemma}
\label{etale-cohomology-lemma-closed-immersion-almost-full}
Soit $i : Z \to X$ une immersion ferm\'ee de sch\'emas.
Soient $U, U'$ des sch\'emas \'etales sur $X$. Soit $h : U_Z \to U'_Z$
un morphisme sur $Z$. Alors il existe un diagramme
$$
\xymatrix{
U & W \ar[l]_a \ar[r]^b & U'
}
$$
dans $X_\etale$ tel que $a_Z : W_Z \to U_Z$
soit un isomorphisme et $h = b_Z \circ (a_Z)^{-1}$.
\end{lemma}

\begin{proof}
Consid\'erons le sch\'ema $M = U \times_X U'$. Le graphe $\Gamma_h \subset M_Z$
de $h$ est ouvert. Cela r\'esulte, par exemple, de ce que $\Gamma_h$ est l'image d'une
section du morphisme \'etale $\text{pr}_{1, Z} : M_Z \to U_Z$, voir
Morphismes étales, Proposition \ref{etale-proposition-properties-sections}.
Il existe donc un sous-sch\'ema ouvert $W \subset M$ dont l'intersection avec
la partie ferm\'ee $M_Z$ est $\Gamma_h$. Posons $a = \text{pr}_1|_W$
et $b = \text{pr}_2|_W$.
\end{proof}

\begin{lemma}
\label{etale-cohomology-lemma-closed-immersion-almost-essentially-surjective}
Soit $i : Z \to X$ une immersion ferm\'ee de sch\'emas.
Soit $V \to Z$ un morphisme \'etale de sch\'emas.
Il existe des morphismes \'etales $U_i \to X$ et des morphismes
$U_{i, Z} \to V$ tels que $\{U_{i, Z} \to V\}$
soit un recouvrement de Zariski de $V$.
\end{lemma}

\begin{proof}
Puisqu'il nous suffit de trouver un recouvrement de Zariski de $V$ form\'e de sch\'emas
de la forme $U_Z$, avec $U$ \'etale sur $X$, nous pouvons raisonner localement pour la topologie de Zariski sur $X$
et sur $V$. Nous pouvons donc supposer $X$ et $V$ affines. Dans le cas affine, cela r\'esulte de
Algèbre, Lemme \ref{algebra-lemma-lift-etale}.
\end{proof}

\noindent
Si $\overline{x} : \Spec(k) \to X$ est un point g\'eom\'etrique de $X$, alors
ou bien $\overline{x}$ se factorise (de fa\c{c}on unique) par le sous-sch\'ema ferm\'e $Z$, ou bien
$Z_{\overline{x}} = \emptyset$. Si $\overline{x}$ se factorise par $Z$,
nous disons que $\overline{x}$ est un point g\'eom\'etrique de $Z$ (ce qu'il est) et
nous employons la notation ``$\overline{x} \in Z$'' pour l'indiquer.

\begin{lemma}
\label{etale-cohomology-lemma-stalk-pushforward-closed-immersion}
Soit $i : Z \to X$ une immersion ferm\'ee de sch\'emas.
Soit $\mathcal{G}$ un faisceau d'ensembles sur $Z_\etale$.
Soit $\overline{x}$ un point g\'eom\'etrique de $X$.
Alors
$$
(i_{small, *}\mathcal{G})_{\overline{x}} =
\left\{
\begin{matrix}
* & \text{si} & \overline{x} \not \in Z \\
\mathcal{G}_{\overline{x}} & \text{si} & \overline{x} \in Z
\end{matrix}
\right.
$$
o\`u $*$ d\'esigne un ensemble r\'eduit \`a un \'el\'ement.
\end{lemma}

\begin{proof}
Posons $U = X \setminus Z$.
Notons que $i_{small, *}\mathcal{G}|_{U_\etale} = *$ est l'objet final
de la cat\'egorie des faisceaux \'etales sur $U$, c'est-\`a-dire le faisceau
qui associe un ensemble r\'eduit \`a un \'el\'ement \`a chaque sch\'ema \'etale sur $U$.
Cela explique la valeur de $(i_{small, *}\mathcal{G})_{\overline{x}}$
lorsque $\overline{x} \not \in Z$.

\medskip\noindent
Supposons maintenant que $\overline{x} \in Z$. Notons que
$$
(i_{small, *}\mathcal{G})_{\overline{x}}
=
\colim_{(U, \overline{u})} \mathcal{G}(U_Z)
$$
et, d'autre part,
$$
\mathcal{G}_{\overline{x}}
=
\colim_{(V, \overline{v})} \mathcal{G}(V).
$$
Soit $\mathcal{C}_1 = \{(U, \overline{u})\}^{opp}$ la cat\'egorie oppos\'ee \`a la
cat\'egorie des voisinages \'etales de $\overline{x}$ dans $X$, et soit
$\mathcal{C}_2 = \{(V, \overline{v})\}^{opp}$ la cat\'egorie oppos\'ee \`a la
cat\'egorie des voisinages \'etales de $\overline{x}$ dans $Z$. L'application canonique
$$
\mathcal{G}_{\overline{x}}
\longrightarrow
(i_{small, *}\mathcal{G})_{\overline{x}}
$$
correspond au foncteur $F : \mathcal{C}_1 \to \mathcal{C}_2$,
$F(U, \overline{u}) = (U_Z, \overline{x})$. Or les
lemmes \ref{etale-cohomology-lemma-closed-immersion-almost-essentially-surjective} et
\ref{etale-cohomology-lemma-closed-immersion-almost-full}
impliquent que $\mathcal{C}_1$ est cofinale dans $\mathcal{C}_2$, voir
Catégories, définition \ref{categories-definition-cofinal}.
Il en r\'esulte que la fl\`eche ci-dessus est un isomorphisme, voir
Catégories, lemme \ref{categories-lemma-cofinal}.
\end{proof}

\begin{proposition}
\label{etale-cohomology-proposition-closed-immersion-pushforward}
Soit $i : Z \to X$ une immersion ferm\'ee de sch\'emas.
\begin{enumerate}
\item Le foncteur
$$
i_{small, *} :
\Sh(Z_\etale)
\longrightarrow
\Sh(X_\etale)
$$
est pleinement fid\`ele et son image essentielle est form\'ee des faisceaux d'ensembles
$\mathcal{F}$ sur $X_\etale$ dont la restriction \`a $X \setminus Z$ est
isomorphe \`a $*$, et
\item le foncteur
$$
i_{small, *} :
\textit{Ab}(Z_\etale)
\longrightarrow
\textit{Ab}(X_\etale)
$$
est pleinement fid\`ele et son image essentielle est form\'ee des faisceaux ab\'eliens sur
$X_\etale$ dont le support est contenu dans $Z$.
\end{enumerate}
Dans les deux cas, $i_{small}^{-1}$ est un inverse \`a gauche du foncteur
$i_{small, *}$.
\end{proposition}

\begin{proof}
Traitons le cas des faisceaux d'ensembles.
Pour tout faisceau $\mathcal{G}$ sur $Z$, le morphisme
$i_{small}^{-1}i_{small, *}\mathcal{G} \to \mathcal{G}$
est un isomorphisme d'apr\`es
le lemme \ref{etale-cohomology-lemma-stalk-pushforward-closed-immersion}
(et
le th\'eor\`eme \ref{etale-cohomology-theorem-exactness-stalks}).
Il en r\'esulte formellement que $i_{small, *}$ est pleinement fid\`ele, voir
Sites, lemme \ref{sites-lemma-exactness-properties}.
Il est clair que $i_{small, *}\mathcal{G}|_{U_\etale} \cong *$
o\`u $U = X \setminus Z$. R\'eciproquement, supposons que $\mathcal{F}$
soit un faisceau d'ensembles sur $X$ tel que $\mathcal{F}|_{U_\etale} \cong *$.
Consid\'erons le morphisme d'adjonction
$$
\mathcal{F} \longrightarrow i_{small, *}i_{small}^{-1}\mathcal{F}
$$
En combinant
les lemmes \ref{etale-cohomology-lemma-stalk-pushforward-closed-immersion} et
\ref{etale-cohomology-lemma-stalk-pullback},
nous voyons que c'est un isomorphisme. Cela ach\`eve la d\'emonstration de (1).
La d\'emonstration de (2) est identique.
\end{proof}





\section{Morphismes entiers universellement injectifs}
\label{etale-cohomology-section-integral-universally-injective}

\noindent
Voici la version g\'en\'erale de la
proposition \ref{etale-cohomology-proposition-closed-immersion-pushforward}.

\begin{proposition}
\label{etale-cohomology-proposition-integral-universally-injective-pushforward}
Soit $f : X \to Y$ un morphisme entier de sch\'emas
et universellement injectif.
\begin{enumerate}
\item Le foncteur
$$
f_{small, *} :
\Sh(X_\etale)
\longrightarrow
\Sh(Y_\etale)
$$
est pleinement fid\`ele et son image essentielle est form\'ee des faisceaux d'ensembles
$\mathcal{F}$ sur $Y_\etale$ dont la restriction \`a $Y \setminus f(X)$ est
isomorphe \`a $*$, et
\item le foncteur
$$
f_{small, *} :
\textit{Ab}(X_\etale)
\longrightarrow
\textit{Ab}(Y_\etale)
$$
est pleinement fid\`ele et son image essentielle est form\'ee des faisceaux ab\'eliens sur
$Y_\etale$ dont le support est contenu dans $f(X)$.
\end{enumerate}
Dans les deux cas, $f_{small}^{-1}$ est un inverse \`a gauche du foncteur
$f_{small, *}$.
\end{proposition}

\begin{proof}
Nous pouvons factoriser $f$ sous la forme
$$
\xymatrix{
X \ar[r]^h & Z \ar[r]^i & Y
}
$$
o\`u $h$ est entier, universellement injectif et surjectif
et $i : Z \to Y$ est une immersion ferm\'ee.
Appliquons la
proposition \ref{etale-cohomology-proposition-closed-immersion-pushforward}
\`a $i$ et le
th\'eor\`eme \ref{etale-cohomology-theorem-topological-invariance}
\`a $h$.
\end{proof}









\section{Grands sites et image directe}
\label{etale-cohomology-section-big}

\noindent
Dans cette section, nous d\'emontrons quelques r\'esultats techniques sur $f_{big, *}$ pour
certains types de morphismes de sch\'emas.

\begin{lemma}
\label{etale-cohomology-lemma-monomorphism-big-push-pull}
Soit $\tau \in \{Zariski, \etale, smooth, syntomic, fppf\}$.
Soit $f : X \to Y$ un monomorphisme de sch\'emas.
Alors le morphisme canonique
$f_{big}^{-1}f_{big, *}\mathcal{F} \to \mathcal{F}$
est un isomorphisme pour tout faisceau $\mathcal{F}$ sur
$(\Sch/X)_\tau$.
\end{lemma}

\begin{proof}
Dans ce cas, le foncteur $(\Sch/X)_\tau \to (\Sch/Y)_\tau$
est continu, cocontinu et pleinement fid\`ele. Le r\'esultat d\'ecoule donc de
Sites, lemme \ref{sites-lemma-back-and-forth}.
\end{proof}

\begin{remark}
\label{etale-cohomology-remark-push-pull-shriek}
Dans la situation du
lemme \ref{etale-cohomology-lemma-monomorphism-big-push-pull},
le morphisme canonique
$\mathcal{F} \to f_{big}^{-1}f_{big!}\mathcal{F}$
est un isomorphisme pour tout faisceau d'ensembles $\mathcal{F}$ sur
$(\Sch/X)_\tau$. La d\'emonstration est la m\^eme. Cela
vaut aussi pour les faisceaux ab\'eliens. Cependant, notons
que le foncteur $f_{big!}$ pour les faisceaux ab\'eliens est d\'efini dans
Modules sur les sites, section \ref{sites-modules-section-exactness-lower-shriek}
et est en g\'en\'eral diff\'erent de $f_{big!}$ sur les faisceaux d'ensembles.
Le r\'esultat pour les faisceaux ab\'eliens d\'ecoule de
Modules sur les sites, lemme \ref{sites-modules-lemma-back-and-forth}.
\end{remark}

\begin{lemma}
\label{etale-cohomology-lemma-closed-immersion-cover-from-below}
Soit $f : X \to Y$ une immersion ferm\'ee de sch\'emas.
Soit $U \to X$ un morphisme syntomique (resp.\ lisse, resp.\ \'etale).
Il existe alors des morphismes syntomiques (resp.\ lisses, resp.\ \'etales)
$V_i \to Y$ et des morphismes $V_i \times_Y X \to U$ tels que
$\{V_i \times_Y X \to U\}$ forme un recouvrement de Zariski de $U$.
\end{lemma}

\begin{proof}
D\'emontrons le lemme dans le cas o\`u $\tau = syntomic$.
La question est locale sur $U$. Nous pouvons donc supposer que $U$ est
un sch\'ema affine dont l'image est contenue dans un ouvert affine de $Y$.
Nous sommes donc ramen\'es \`a d\'emontrer le cas suivant :
$Y = \Spec(A)$, $X = \Spec(A/I)$, et
$U = \Spec(\overline{B})$, o\`u
$A/I \to \overline{B}$ est un morphisme d'anneaux syntomique.
D'apr\`es Algèbre, lemme \ref{algebra-lemma-lift-syntomic}
nous pouvons trouver des \'el\'ements $\overline{g}_i \in \overline{B}$
tels que
$\overline{B}_{\overline{g}_i} = A_i/IA_i$ pour certains morphismes d'anneaux syntomiques
$A \to A_i$.
Cela d\'emontre le lemme dans le cas syntomique.
La d\'emonstration du cas lisse est identique, \`a ceci pr\`es qu'elle utilise
Algèbre, lemme \ref{algebra-lemma-lift-smooth}.
Dans le cas \'etale, on utilise
Algèbre, lemme \ref{algebra-lemma-lift-etale}.
\end{proof}

\begin{lemma}
\label{etale-cohomology-lemma-prepare-closed-immersion-almost-cocontinuous}
Soit $f : X \to Y$ une immersion ferm\'ee de sch\'emas.
Soit $\{U_i \to X\}$ un recouvrement syntomique (resp.\ lisse, resp.\ \'etale).
Il existe un recouvrement syntomique (resp.\ lisse, resp.\ \'etale) $\{V_j \to Y\}$
tel que, pour chaque $j$, soit $V_j \times_Y X = \emptyset$, soit le
morphisme $V_j \times_Y X \to X$ se factorise par $U_i$ pour un certain $i$.
\end{lemma}

\begin{proof}
Pour chaque $i$, nous pouvons choisir des morphismes syntomiques (resp.\ lisses, resp.\ \'etales)
$g_{ij} : V_{ij} \to Y$ et des morphismes $V_{ij} \times_Y X \to U_i$ au-dessus de $X$,
de sorte que $\{V_{ij} \times_Y X \to U_i\}$ forme un recouvrement de Zariski, voir
le lemme \ref{etale-cohomology-lemma-closed-immersion-cover-from-below}.
Cela implique en particulier que
$\bigcup_{ij} g_{ij}(V_{ij})$ contient le sous-ensemble ferm\'e $f(X)$.
Par cons\'equent, la famille de morphismes syntomiques (resp.\ lisses, resp.\ \'etales) $g_{ij}$
ainsi que l'immersion ouverte $Y \setminus f(X) \to Y$ forment le recouvrement
syntomique (resp.\ lisse, resp.\ \'etale) voulu de $Y$.
\end{proof}

\begin{lemma}
\label{etale-cohomology-lemma-closed-immersion-almost-cocontinuous}
Soit $f : X \to Y$ une immersion ferm\'ee de sch\'emas.
Soit $\tau \in \{syntomic, smooth, \etale\}$.
Le foncteur $V \mapsto X \times_Y V$ d\'efinit un foncteur presque
cocontinu (voir
Sites, D\'efinition \ref{sites-definition-almost-cocontinuous})
$(\Sch/Y)_\tau \to (\Sch/X)_\tau$ entre les
grands $\tau$-sites.
\end{lemma}

\begin{proof}
Il faut montrer ce qui suit : \'etant donn\'e un morphisme $V \to Y$
et un recouvrement syntomique quelconque (resp.\ lisse, resp.\ \'etale)
$\{U_i \to X \times_Y V\}$, il existe un
recouvrement lisse (resp.\ lisse, resp.\ \'etale) $\{V_j \to V\}$
tel que, pour tout $j$, soit $X \times_Y V_j$ est vide, soit
$X \times_Y V_j \to Z \times_Y V$ se factorise par l'un des
$U_i$. Cela r\'esulte de l'application du
Lemme \ref{etale-cohomology-lemma-prepare-closed-immersion-almost-cocontinuous}
ci-dessus \`a l'immersion ferm\'ee $X \times_Y V \to V$.
\end{proof}

\begin{lemma}
\label{etale-cohomology-lemma-closed-immersion-pushforward-exact}
Soit $f : X \to Y$ une immersion ferm\'ee de sch\'emas.
Soit $\tau \in \{syntomic, smooth, \etale\}$.
\begin{enumerate}
\item Le foncteur d'image directe
$f_{big, *} :
\Sh((\Sch/X)_\tau)
\to
\Sh((\Sch/Y)_\tau)$
commute aux co\'egalisateurs et aux sommes amalgam\'ees.
\item Le foncteur d'image directe
$f_{big, *} :
\textit{Ab}((\Sch/X)_\tau)
\to
\textit{Ab}((\Sch/Y)_\tau)$
est exact.
\end{enumerate}
\end{lemma}

\begin{proof}
Cela r\'esulte de
Sites, Lemme \ref{sites-lemma-morphism-of-sites-almost-cocontinuous},
Modules sur les sites,
Lemme \ref{sites-modules-lemma-morphism-ringed-sites-almost-cocontinuous},
et du
Lemme \ref{etale-cohomology-lemma-closed-immersion-almost-cocontinuous}
ci-dessus.
\end{proof}

\begin{remark}
\label{etale-cohomology-remark-fppf-closed-immersion-not-closed}
Dans le Lemme \ref{etale-cohomology-lemma-closed-immersion-pushforward-exact}, le cas $\tau = fppf$
n'est pas trait\'e. La raison en est qu'\'etant donn\'es un anneau $A$, un id\'eal $I$ et un
homomorphisme d'anneaux $A/I \to \overline{B}$ fid\`element plat et de pr\'esentation finie, il
n'y a aucune raison de penser que l'on puisse trouver {\it un quelconque} homomorphisme d'anneaux
$A \to B$, plat et de pr\'esentation finie, avec $B/IB \not = 0$, tel que $A/I \to B/IB$ se factorise par
$\overline{B}$. Ainsi, la preuve du
Lemme \ref{etale-cohomology-lemma-closed-immersion-almost-cocontinuous}
ne fonctionne pas pour la topologie fppf.
En fait, il est probablement faux que
$f_{big, *} : \textit{Ab}((\Sch/X)_{fppf})
\to \textit{Ab}((\Sch/Y)_{fppf})$
soit exact lorsque $f$ est une immersion ferm\'ee.
Si vous connaissez un exemple, veuillez envoyer un courriel \`a
\href{mailto:stacks.project@gmail.com}{stacks.project@gmail.com}.
\end{remark}












\section{Exactitude du grand foncteur image directe \`a support propre}
\sectionmark{Image directe à support propre}
\label{etale-cohomology-section-exactness-lower-shriek}

\noindent
Il s'agit simplement du r\'esultat technique suivant. Notons que le foncteur $f_{big!}$
n'a absolument rien \`a voir avec la cohomologie \`a support compact en
g\'en\'eral.

\begin{lemma}
\label{etale-cohomology-lemma-exactness-lower-shriek}
Soit $\tau \in \{Zariski, \etale, smooth, syntomic, fppf\}$.
Soit $f : X \to Y$ un morphisme de sch\'emas. Soit
$$
f_{big} :
\Sh((\Sch/X)_\tau)
\longrightarrow
\Sh((\Sch/Y)_\tau)
$$
le morphisme de topos correspondant comme dans
Topologies, Lemme
\ref{topologies-lemma-morphism-big},
\ref{topologies-lemma-morphism-big-etale},
\ref{topologies-lemma-morphism-big-smooth},
\ref{topologies-lemma-morphism-big-syntomic}, ou
\ref{topologies-lemma-morphism-big-fppf}.
\begin{enumerate}
\item Le foncteur
$f_{big}^{-1} : \textit{Ab}((\Sch/Y)_\tau) \to \textit{Ab}((\Sch/X)_\tau)$
admet un adjoint \`a gauche
$$
f_{big!} : \textit{Ab}((\Sch/X)_\tau) \to \textit{Ab}((\Sch/Y)_\tau)
$$
qui est exact.
\item Le foncteur
$f_{big}^* :
\textit{Mod}((\Sch/Y)_\tau, \mathcal{O})
\to
\textit{Mod}((\Sch/X)_\tau, \mathcal{O})$
admet un adjoint \`a gauche
$$
f_{big!} :
\textit{Mod}((\Sch/X)_\tau, \mathcal{O})
\to
\textit{Mod}((\Sch/Y)_\tau, \mathcal{O})
$$
qui est exact.
\end{enumerate}
De plus, les deux foncteurs $f_{big!}$ co\"incident sur les faisceaux sous-jacents
de groupes ab\'eliens.
\end{lemma}

\begin{proof}
Rappelons que $f_{big}$ est le morphisme de topos associ\'e au
foncteur continu et cocontinu
$u : (\Sch/X)_\tau \to (\Sch/Y)_\tau$, $U/X \mapsto U/Y$.
De plus, on a $f_{big}^{-1}\mathcal{O} = \mathcal{O}$.
Ainsi, l'existence de $f_{big!}$ r\'esulte de
Modules sur les sites, Lemme \ref{sites-modules-lemma-g-shriek-adjoint},
et, respectivement, de
Modules sur les sites, Lemme \ref{sites-modules-lemma-lower-shriek-modules}.
Notons que, si $U$ est un objet de $(\Sch/X)_\tau$, alors le foncteur
$u$ induit une \'equivalence de cat\'egories
$$
u' :
(\Sch/X)_\tau/U
\longrightarrow
(\Sch/Y)_\tau/U
$$
car les deux membres de la fl\`eche sont \'egaux \`a $(\Sch/U)_\tau$.
Ainsi, la co\"incidence de $f_{big!}$ sur les faisceaux ab\'eliens sous-jacents
r\'esulte de la discussion dans
Modules sur les sites, Remarque \ref{sites-modules-remark-when-shriek-equal}.
L'exactitude de $f_{big!}$ r\'esulte de
Modules sur les sites, Lemme \ref{sites-modules-lemma-exactness-lower-shriek}
puisque le foncteur $u$ ci-dessus commute aux produits fibr\'es et aux \'egalisateurs.
\end{proof}

\noindent
D\'emontrons ensuite un lemme technique qui sera utile plus loin pour comparer
les faisceaux de modules sur les diff\'erents sites associ\'es aux champs alg\'ebriques.

\begin{lemma}
\label{etale-cohomology-lemma-compare-structure-sheaves}
Soit $X$ un sch\'ema. Soit
$\tau \in \{Zariski, \etale, smooth, syntomic, fppf\}$.
Soient $\mathcal{C}_1 \subset \mathcal{C}_2 \subset (\Sch/X)_\tau$ des sous-cat\'egories pleines
ayant les propri\'et\'es suivantes :
\begin{enumerate}
\item Pour tout objet $U/X$ de $\mathcal{C}_t$,
\begin{enumerate}
\item si $\{U_i \to U\}$ est un recouvrement de $(\Sch/X)_\tau$, alors
$U_i/X$ est un objet de $\mathcal{C}_t$,
\item $U \times \mathbf{A}^1/X$ est un objet de $\mathcal{C}_t$.
\end{enumerate}
\item $X/X$ est un objet de $\mathcal{C}_t$.
\end{enumerate}
Munissons $\mathcal{C}_t$ de la structure de site dont les recouvrements sont
exactement les recouvrements $\{U_i \to U\}$ de $(\Sch/X)_\tau$ tels que
$U \in \Ob(\mathcal{C}_t)$. Alors
\begin{enumerate}
\item[(a)] Le foncteur $\mathcal{C}_1 \to \mathcal{C}_2$
est pleinement fid\`ele, continu et cocontinu.
\end{enumerate}
Notons $g : \Sh(\mathcal{C}_1) \to \Sh(\mathcal{C}_2)$ le
morphisme de topos correspondant. Notons $\mathcal{O}_t$ la restriction de $\mathcal{O}$
\`a $\mathcal{C}_t$. Notons $g_!$ le foncteur de
Modules sur les sites, D\'efinition \ref{sites-modules-definition-g-shriek}.
\begin{enumerate}
\item[(b)] Le morphisme canonique $g_!\mathcal{O}_1 \to \mathcal{O}_2$
est un isomorphisme.
\end{enumerate}
\end{lemma}

\begin{proof}
L'assertion (a) est imm\'ediate d'apr\`es les d\'efinitions.
Dans cette preuve, tous les sch\'emas sont des sch\'emas sur $X$ et tous les morphismes de
sch\'emas sont des morphismes de sch\'emas sur $X$. Notons que $g^{-1}$ est
donn\'e par restriction, de sorte que, pour tout objet $U$ de $\mathcal{C}_1$
on a $\mathcal{O}_1(U) = \mathcal{O}_2(U) = \mathcal{O}(U)$.
Rappelons que $g_!\mathcal{O}_1$ est le faisceau associ\'e au pr\'efaisceau
$g_{p!}\mathcal{O}_1$ qui associe \`a $V$ dans $\mathcal{C}_2$ le groupe
$$
\colim_{V \to U} \mathcal{O}(U)
$$
o\`u $U$ parcourt les objets de $\mathcal{C}_1$ et o\`u la colimite est
prise dans la cat\'egorie des groupes ab\'eliens. Nous utiliserons souvent ci-dessous
le fait que, si
$$
V \to U \to U'
$$
sont des morphismes avec $U, U' \in \Ob(\mathcal{C}_1)$
et si $f' \in \mathcal{O}(U')$ se restreint en $f \in \mathcal{O}(U)$,
alors $(V \to U, f)$ et $(V \to U', f')$ d\'efinissent le m\^eme \'el\'ement de la
colimite. De plus, $g_!\mathcal{O}_1 \to \mathcal{O}_2$ envoie l'\'el\'ement
$(V \to U, f)$ simplement sur la restriction de $f$ \`a $V$.

\medskip\noindent
Surjectivit\'e. Soit $V$ un sch\'ema et soit $h \in \mathcal{O}(V)$.
On obtient alors un morphisme $V \to X \times \mathbf{A}^1$ induit par $h$
et le morphisme structural $V \to X$. En \'ecrivant
$\mathbf{A}^1 = \Spec(\mathbf{Z}[x])$, on voit que l'\'el\'ement
$x \in \mathcal{O}(X \times \mathbf{A}^1)$ a pour image inverse
$h$. Puisque $X \times \mathbf{A}^1$ est un objet de $\mathcal{C}_1$
d'apr\`es les hypoth\`eses (1)(b) et (2), on obtient la surjectivit\'e voulue.

\medskip\noindent
Injectivit\'e. Soit $V$ un sch\'ema. Soit
$s = \sum_{i = 1, \ldots, n} (V \to U_i, f_i)$ un \'el\'ement de la colimite
affich\'ee ci-dessus. Pour tout $i$, on peut utiliser le morphisme
$f_i : U_i \to X \times \mathbf{A}^1$
pour voir que $(V \to U_i, f_i)$ d\'efinit le m\^eme \'el\'ement de la colimite que
$(f_i : V \to X \times \mathbf{A}^1, x)$. Nous pouvons alors consid\'erer
$$
f_1 \times \ldots \times f_n : V \to X \times \mathbf{A}^n
$$
et l'on voit que $s$ est \'equivalent dans la colimite \`a
$$
\sum\nolimits_{i = 1, \ldots, n}
(f_1 \times \ldots \times f_n : V \to X \times \mathbf{A}^n, x_i) =
(f_1 \times \ldots \times f_n : V \to X \times \mathbf{A}^n,
x_1 + \ldots + x_n)
$$
Maintenant, si $x_1 + \ldots + x_n$ se restreint \`a z\'ero sur $V$, alors on voit
que $f_1 \times \ldots \times f_n$ se factorise par
$X \times \mathbf{A}^{n - 1} = V(x_1 + \ldots + x_n)$. Ainsi, nous voyons
que $s$ est \'equivalent \`a z\'ero dans la limite inductive.
\end{proof}











%9.29.09
\section{Cohomologie \'etale}
\label{etale-cohomology-section-etale-cohomology}

\noindent
Dans les sections suivantes, nous d\'emontrons quelques r\'esultats \`el\'ementaires sur la cohomologie \'etale.
Voici une propri\'et\'e connue de la cohomologie des espaces
topologiques et qui vaut \`egalement pour la cohomologie \'etale.

\begin{lemma}[Mayer-Vietoris en cohomologie \'etale]
\label{etale-cohomology-lemma-mayer-vietoris}
Soit $X$ un sch\'ema. Supposons que $X = U \cup V$ soit la
r\'eunion de deux ouverts. Pour tout faisceau ab\'elien $\mathcal{F}$ sur $X_\etale$
il existe une suite exacte longue de cohomologie
$$
\begin{matrix}
0 \to
H^0_\etale(X, \mathcal{F}) \to
H^0_\etale(U, \mathcal{F}) \oplus H^0_\etale(V, \mathcal{F}) \to
H^0_\etale(U \cap V, \mathcal{F}) \phantom{\to \ldots} \\
\phantom{0} \to H^1_\etale(X, \mathcal{F}) \to
H^1_\etale(U, \mathcal{F}) \oplus H^1_\etale(V, \mathcal{F}) \to
H^1_\etale(U \cap V, \mathcal{F}) \to \ldots
\end{matrix}
$$
Cette suite exacte longue est fonctorielle en $\mathcal{F}$.
\end{lemma}

\begin{proof}
Observons que si $\mathcal{I}$ est un faisceau ab\'elien injectif, alors
$$
0 \to \mathcal{I}(X) \to \mathcal{I}(U) \oplus \mathcal{I}(V) \to
\mathcal{I}(U \cap V) \to 0
$$
est exacte. L'exactitude aux premier et second termes vient de ce que $\mathcal{I}$
est un faisceau. L'exactitude \`a droite vient de ce que
$\mathcal{I}(U) \to \mathcal{I}(U \cap V)$ est surjective d'apr\`es
Cohomologie sur les sites, Lemme
\ref{sites-cohomology-lemma-restriction-along-monomorphism-surjective}.
Une autre mani\`ere de le d\'emontrer serait de montrer que le conoyau
de l'application $\mathcal{I}(U) \oplus \mathcal{I}(V) \to
\mathcal{I}(U \cap V)$ est le premier groupe de cohomologie de {\v C}ech
de $\mathcal{I}$ relativement au recouvrement
$X = U \cup V$, lequel s'annule d'apr\`es les
lemmes \ref{etale-cohomology-lemma-hom-injective} et \ref{etale-cohomology-lemma-forget-injectives}.
Ainsi, si $\mathcal{F} \to \mathcal{I}^\bullet$ est une
r\'esolution injective, alors
$$
0 \to \mathcal{I}^\bullet(X) \to
\mathcal{I}^\bullet(U) \oplus \mathcal{I}^\bullet(V) \to
\mathcal{I}^\bullet(U \cap V) \to 0
$$
est une suite exacte courte de complexes et la suite exacte longue
de cohomologie associ\'ee est celle de l'\'enonc\'e du lemme.
\end{proof}

\begin{lemma}[Mayer-Vietoris relatif]
\label{etale-cohomology-lemma-relative-mayer-vietoris}
Soit $f : X \to Y$ un morphisme de sch\'emas. Supposons que $X = U \cup V$
soit une r\'eunion de deux sous-sch\'emas ouverts. Notons
$a = f|_U : U \to Y$, $b = f|_V : V \to Y$, et
$c = f|_{U \cap V} : U \cap V \to Y$.
Pour tout faisceau ab\'elien $\mathcal{F}$ sur $X_\etale$
il existe une suite exacte longue
$$
0 \to
f_*\mathcal{F} \to
a_*(\mathcal{F}|_U) \oplus b_*(\mathcal{F}|_V) \to
c_*(\mathcal{F}|_{U \cap V}) \to
R^1f_*\mathcal{F} \to \ldots
$$
sur $Y_\etale$.
Cette suite exacte longue est fonctorielle en $\mathcal{F}$.
\end{lemma}

\begin{proof}
Soit $\mathcal{F} \to \mathcal{I}^\bullet$ une r\'esolution injective
de $\mathcal{F}$ sur $X_\etale$. Nous affirmons qu'il existe
une suite exacte courte de complexes
$$
0 \to
f_*\mathcal{I}^\bullet \to
a_*\mathcal{I}^\bullet|_U \oplus b_*\mathcal{I}^\bullet|_V \to
c_*\mathcal{I}^\bullet|_{U \cap V} \to
0.
$$
En effet, pour tout $W$ de $Y_\etale$ et tout $n \geq 0$, la
suite correspondante des groupes de sections sur $W$
$$
0 \to
\mathcal{I}^n(W \times_Y X) \to
\mathcal{I}^n(W \times_Y U)
\oplus \mathcal{I}^n(W \times_Y V) \to
\mathcal{I}^n(W \times_Y (U \cap V)) \to
0
$$
est une suite exacte courte, comme on l'a montr\'e dans la d\'emonstration du lemme \ref{etale-cohomology-lemma-mayer-vietoris}.
Le lemme r\'esulte alors du passage aux faisceaux de cohomologie et du fait que
$\mathcal{I}^\bullet|_U$ est une r\'esolution injective de $\mathcal{F}|_U$
et de m\^eme pour $\mathcal{I}^\bullet|_V$, $\mathcal{I}^\bullet|_{U \cap V}$.
\end{proof}







\section{Limites inductives}
\label{etale-cohomology-section-colimit}

\noindent
Rappelons que si $(\mathcal{F}_i, \varphi_{ii'})$ est un diagramme de faisceaux sur
un site $\mathcal{C}$, sa limite inductive (dans la cat\'egorie des faisceaux) est le
faisceau associ\'e au pr\'efaisceau $U \mapsto \colim_i \mathcal{F}_i(U)$. Voir
Sites, Lemme \ref{sites-lemma-colimit-sheaves}.
Si le syst\`eme est filtrant, si $U$ est un objet quasi-compact de
$\mathcal{C}$ qui poss\`ede un syst\`eme cofinal de recouvrements par des objets
quasi-compacts, alors $\mathcal{F}(U) = \colim \mathcal{F}_i(U)$, voir
Sites, Lemme \ref{sites-lemma-directed-colimits-sections}.
Voir Cohomologie sur les sites, Lemme
\ref{sites-cohomology-lemma-colim-works-over-collection}
pour un r\'esultat concernant les groupes de cohomologie sup\'erieure des limites inductives
de faisceaux ab\'eliens.

\medskip\noindent
Dans Cohomologie sur les sites, Lemme \ref{sites-cohomology-lemma-colimit},
nous g\'en\'eralisons ce r\'esultat \`a un syst\`eme de faisceaux
sur un syst\`eme projectif de sites. Voici la notion correspondante
dans le cas d'un syst\`eme de faisceaux \'etales d\'efinis
sur un syst\`eme projectif de sch\'emas.

\begin{definition}
\label{etale-cohomology-definition-inverse-system-sheaves}
Soit $I$ un ensemble pr\'eordonn\'e. Soit $(X_i, f_{i'i})$ un syst\`eme
projectif de sch\'emas index\'e par $I$.
Un {\it syst\`eme $(\mathcal{F}_i, \varphi_{i'i})$ de faisceaux
sur $(X_i, f_{i'i})$} est la donn\'ee de
\begin{enumerate}
\item un faisceau $\mathcal{F}_i$ sur $(X_i)_\etale$ pour tout $i \in I$,
\item pour $i' \geq i$, un morphisme
$\varphi_{i'i} : f_{i'i}^{-1}\mathcal{F}_i \to \mathcal{F}_{i'}$
de faisceaux sur $(X_{i'})_\etale$
\end{enumerate}
de telle sorte que $\varphi_{i''i} = \varphi_{i''i'} \circ f_{i'' i'}^{-1}\varphi_{i'i}$
lorsque $i'' \geq i' \geq i$.
\end{definition}

\noindent
Dans la situation de la d\'efinition \ref{etale-cohomology-definition-inverse-system-sheaves},
supposons que $I$ soit filtrant et que les morphismes de transition $f_{i'i}$ soient affines.
Soit $X = \lim X_i$ la limite dans la cat\'egorie des sch\'emas, voir
Limites, section \ref{limits-section-limits}.
Notons $f_i : X \to X_i$ les morphismes de projection et consid\'erons les morphismes
$$
f_i^{-1}\mathcal{F}_i = f_{i'}^{-1}f_{i'i}^{-1}\mathcal{F}_i
\xrightarrow{f_{i'}^{-1}\varphi_{i'i}}
f_{i'}^{-1}\mathcal{F}_{i'}
$$
Cela fait des $f_i^{-1}\mathcal{F}_i$ un syst\`eme de faisceaux sur $X_\etale$
index\'e par $I$ (c'est un bon exercice de le v\'erifier).
Nous voulons souvent savoir si l'on a un isomorphisme
$$
H^q_\etale(X, \colim f_i^{-1}\mathcal{F}_i) =
\colim H^q_\etale(X_i, \mathcal{F}_i)
$$
Il s'av\'erera que cela est vrai si $X_i$ est quasi-compact et quasi-s\'epar\'e
pour tout $i$, voir le th\'eor\`eme \ref{etale-cohomology-theorem-colimit}.

\begin{lemma}
\label{etale-cohomology-lemma-colimit-affine-sites}
Soit $I$ un ensemble pr\'eordonn\'e filtrant. Soit $(X_i, f_{i'i})$ un syst\`eme
projectif de sch\'emas index\'e par $I$ dont les morphismes de transition sont affines.
Soit $X = \lim_{i \in I} X_i$. Avec les
notations de Topologies, Lemme \ref{topologies-lemma-alternative}, on a
$$
X_{affine, \etale} = \colim (X_i)_{affine, \etale}
$$
comme sites au sens de
Sites, Lemme \ref{sites-lemma-colimit-sites}.
\end{lemma}

\begin{proof}
D\'emontrons d'abord ceci lorsque $X$ et les $X_i$ sont quasi-compacts et
quasi-s\'epar\'es pour tout $i$ (c'est le cas dans toutes les situations qui
nous int\'eressent). Dans ce cas, tout objet de
$X_{affine, \etale}$, resp.\ de $(X_i)_{affine, \etale}$, est
de pr\'esentation finie sur $X$. De plus, la
cat\'egorie des sch\'emas de pr\'esentation finie sur $X$ est la
limite inductive des cat\'egories de sch\'emas de pr\'esentation finie
sur les $X_i$, voir
Limites, Lemme \ref{limits-lemma-descend-finite-presentation}.
Il en va de m\^eme des sous-cat\'egories d'objets affines et \'etales
sur $X$, d'apr\`es les lemmes de Limites
\ref{limits-lemma-limit-affine} et \ref{limits-lemma-descend-etale}.
Enfin, si $\{U^j \to U\}$ est un recouvrement de $X_{affine, \etale}$
et si $U_i^j \to U_i$ est un morphisme de sch\'emas affines \'etales sur
$X_i$ dont le changement de base \`a $X$ est $U^j \to U$, alors nous voyons que
le changement de base de $\{U^j_i \to U_i\}$ \`a un certain $X_{i'}$ est
un recouvrement pour $i'$ assez grand, voir
Limites, Lemme \ref{limits-lemma-descend-surjective}.

\medskip\noindent
Dans le cas g\'en\'eral, soit $U$ un objet de $X_{affine, \etale}$.
Alors $U \to X$ est \'etale et s\'epar\'e (puisque $U$ est s\'epar\'e),
mais n'est en g\'en\'eral pas quasi-compact. Toutefois, $U \to X$ est localement
de pr\'esentation finie et donc, d'apr\`es le
lemme \ref{limits-lemma-descend-finite-presentation-variant} de Limites,
il existe un $i$, un sch\'ema $U_i$ quasi-compact et quasi-s\'epar\'e, et
un morphisme $U_i \to X_i$ qui est localement de pr\'esentation finie
dont le changement de base \`a $X$ est $U \to X$. Alors $U = \lim_{i' \geq i} U_{i'}$
o\`u $U_{i'} = U_i \times_{X_i} X_{i'}$.
Quitte \`a augmenter $i$, nous pouvons supposer $U_i$ affine, voir
Limites, Lemme \ref{limits-lemma-limit-affine}.
Pour v\'erifier que $U_i \to X_i$ est \'etale pour $i$ assez grand,
choisissons un recouvrement ouvert affine fini $U_i = U_{i, 1} \cup \ldots \cup U_{i, m}$
tel que $U_{i, j} \to U_i \to X_i$ se factorise par un ouvert affine
$W_{i, j} \subset X_i$. On peut alors appliquer
Limites, Lemme \ref{limits-lemma-descend-etale},
pour voir que $U_{i, j} \to W_{i, j}$ est \'etale
apr\`es avoir, si n\'ecessaire, augment\'e $i$.
Nous voyons ainsi que le foncteur
$\colim (X_i)_{affine, \etale} \to X_{affine, \etale}$
est essentiellement surjectif. Sa pleine fid\'elit\'e r\'esulte
directement du r\'esultat d\'ej\`a utilis\'e :
Limites, Lemme \ref{limits-lemma-descend-finite-presentation-variant}.
L'assertion concernant les recouvrements se d\'emontre exactement de la m\^eme
mani\`ere que dans le premier paragraphe de la d\'emonstration.
\end{proof}

\noindent
Ce qui pr\'ec\`ede donne le r\'esultat g\'en\'eral suivant sur les limites inductives et la cohomologie.

\begin{theorem}
\label{etale-cohomology-theorem-colimit}
Soit $X = \lim_{i \in I} X_i$ la limite d'un syst\`eme projectif de sch\'emas
dont les morphismes de transition $f_{i'i} : X_{i'} \to X_i$ sont affines. Nous supposons
que $X_i$ est quasi-compact et quasi-s\'epar\'e pour tout $i \in I$.
Soit $(\mathcal{F}_i, \varphi_{i'i})$ un syst\`eme de faisceaux ab\'eliens
sur $(X_i, f_{i'i})$. Notons $f_i : X \to X_i$ la projection et posons
$\mathcal{F} = \colim f_i^{-1}\mathcal{F}_i$. Alors
$$
\colim_{i\in I} H_\etale^p(X_i, \mathcal{F}_i) = H_\etale^p(X, \mathcal{F}).
$$
pour tout $p \geq 0$.
\end{theorem}

\begin{proof}
D'apr\`es Topologies, lemme \ref{topologies-lemma-alternative},
nous pouvons calculer la cohomologie
de $\mathcal{F}$ sur $X_{affine, \etale}$.
Le r\'esultat d\'ecoule donc de la combinaison du
lemme \ref{etale-cohomology-lemma-colimit-affine-sites}
et du
lemme \ref{sites-cohomology-lemma-colimit} de Cohomologie sur les sites.
\end{proof}

\noindent
Les deux r\'esultats suivants sont des cas particuliers du th\'eor\`eme ci-dessus.

\begin{lemma}
\label{etale-cohomology-lemma-colimit}
Soit $X$ un sch\'ema quasi-compact et quasi-s\'epar\'e. Soit $I$
un ensemble ordonn\'e filtrant. Soit $(\mathcal{F}_i, \varphi_{ij})$ un syst\`eme
de faisceaux ab\'eliens sur $X_\etale$ index\'e par $I$. Alors
$$
\colim_{i\in I} H_\etale^p(X, \mathcal{F}_i) = H_\etale^p(X,
\colim_{i\in I} \mathcal{F}_i).
$$
\end{lemma}

\begin{proof}
C'est un cas particulier du th\'eor\`eme \ref{etale-cohomology-theorem-colimit}.
Nous esquissons \'egalement une d\'emonstration directe.
Nous le d\'emontrons simultan\'ement pour tout $X$, par r\'ecurrence sur $p$.
\begin{enumerate}
\item
Pour tout sch\'ema quasi-compact et quasi-s\'epar\'e $X$ et tout recouvrement \'etale
$\mathcal{U}$ de $X$, montrer qu'il existe un raffinement
$\mathcal{V} = \{V_j \to X\}_{j\in J}$ avec $J$ fini et chaque $V_j$
quasi-compact et quasi-s\'epar\'e, de telle sorte que tous les
$V_{j_0} \times_X \ldots \times_X V_{j_p}$ soient eux aussi
quasi-compacts et quasi-s\'epar\'es.
\item
En utilisant l'\'etape pr\'ec\'edente et la d\'efinition des limites inductives dans la cat\'egorie des
faisceaux, montrer que le th\'eor\`eme vaut pour $p = 0$ et tout $X$.
\item
En utilisant la localit\'e de la cohomologie
(lemme \ref{etale-cohomology-lemma-locality-cohomology}),
la suite spectrale de {\v C}ech vers la cohomologie
(th\'eor\`eme \ref{etale-cohomology-theorem-cech-ss}) et le fait que l'hypoth\`ese de r\'ecurrence
s'applique \`a tous les
$V_{j_0} \times_X \ldots \times_X V_{j_p}$
dans la situation ci-dessus, d\'emontrer l'\'etape de r\'ecurrence $p \to p + 1$.
\end{enumerate}
\end{proof}

\begin{lemma}
\label{etale-cohomology-lemma-directed-colimit-cohomology}
Soient $A$ un anneau, $(I, \leq)$ un ensemble ordonn\'e filtrant et $(B_i, \varphi_{ij})$ un
syst\`eme de $A$-alg\`ebres. Posons $B = \colim_{i\in I} B_i$. Soit $X \to \Spec(A)$
un morphisme de sch\'emas quasi-compact et quasi-s\'epar\'e. Soit
$\mathcal{F}$ un faisceau ab\'elien sur $X_\etale$.
Notons $Y_i = X \times_{\Spec(A)} \Spec(B_i)$,
$Y = X \times_{\Spec(A)} \Spec(B)$,
$\mathcal{G}_i = (Y_i \to X)^{-1}\mathcal{F}$ et
$\mathcal{G} = (Y \to X)^{-1}\mathcal{F}$. Alors
$$
H_\etale^p(Y, \mathcal{G}) =
\colim_{i\in I} H_\etale^p (Y_i, \mathcal{G}_i).
$$
\end{lemma}

\begin{proof}
C'est un cas particulier du th\'eor\`eme \ref{etale-cohomology-theorem-colimit}.
Nous indiquons \'egalement une d\'emonstration directe comme suit.
\begin{enumerate}
\item \'Etant donn\'e $V \to Y$ \'etale, avec $V$ quasi-compact et
quasi-s\'epar\'e, il existe $i\in I$ et $V_i \to Y_i$ tels que
$V = V_i \times_{Y_i} Y$.
Si tous les sch\'emas consid\'er\'es \'etaient affines, cela correspondrait \`a
l'\'enonc\'e alg\'ebrique suivant : si $B = \colim B_i$ et si $B \to C$ est \'etale,
alors il existe $i \in I$ et $B_i \to C_i$ \'etale tels que
$C \cong B \otimes_{B_i} C_i$.
Ceci est d\'emontr\'e dans Alg\`ebre, lemme \ref{algebra-lemma-etale}.
\item Dans la situation de (1), montrer que
$\mathcal{G}(V) = \colim_{i' \geq i} \mathcal{G}_{i'}(V_{i'})$
o\`u $V_{i'}$ est le changement de base de $V_i$ sur $Y_{i'}$.
\item D'apr\`es (1), on voit que pour tout recouvrement \'etale
$\mathcal{V} = \{V_j \to Y\}_{j\in J}$ avec $J$ fini et les
$V_j$ quasi-compacts et quasi-s\'epar\'es, il existe $i \in I$ et
un recouvrement \'etale $\mathcal{V}_i = \{V_{ij} \to Y_i\}_{j \in J}$
tel que $\mathcal{V} \cong \mathcal{V}_i \times_{Y_i} Y$.
\item Montrer que (2) et (3) impliquent
$$
\check H^*(\mathcal{V}, \mathcal{G})=
\colim_{i\in I} \check H^*(\mathcal{V}_i, \mathcal{G}_i).
$$
\item Utiliser judicieusement la suite spectrale de {\v C}ech vers la cohomologie
(th\'eor\`eme \ref{etale-cohomology-theorem-cech-ss}).
\end{enumerate}
\end{proof}

\begin{lemma}
\label{etale-cohomology-lemma-higher-direct-images}
Soit $f: X\to Y$ un morphisme de sch\'emas et soit $\mathcal{F}\in
\textit{Ab}(X_\etale)$. Alors $R^pf_*\mathcal{F}$ est le faisceau
associ\'e au pr\'efaisceau
$$
(V \to Y) \longmapsto H_\etale^p(X \times_Y V, \mathcal{F}|_{X \times_Y V}).
$$
Plus g\'en\'eralement, pour $K \in D(X_\etale)$, $R^pf_*K$ est le
faisceau associ\'e au pr\'efaisceau
$$
(V \to Y) \longmapsto H_\etale^p(X \times_Y V, K|_{X \times_Y V}).
$$
\end{lemma}

\begin{proof}
Ce lemme est valable pour les espaces topologiques, et la d\'emonstration dans ce cas est la
m\^eme. Voir Cohomologie sur les sites, lemme
\ref{sites-cohomology-lemma-higher-direct-images}
pour le cas d'un faisceau, et voir
Cohomologie sur les sites, lemme
\ref{sites-cohomology-lemma-sheafification-cohomology}
pour le cas d'un complexe de faisceaux ab\'eliens.
\end{proof}

\begin{lemma}
\label{etale-cohomology-lemma-relative-colimit}
Soit $S$ un sch\'ema. Soit $X = \lim_{i \in I} X_i$ une limite d'un
syst\`eme projectif de sch\'emas sur $S$ avec morphismes de transition affines
$f_{i'i} : X_{i'} \to X_i$. Nous supposons que les morphismes structuraux
$g_i : X_i \to S$ et $g : X \to S$ sont quasi-compacts et quasi-s\'epar\'es.
Soit $(\mathcal{F}_i, \varphi_{i'i})$ un syst\`eme de faisceaux ab\'eliens
sur $(X_i, f_{i'i})$. Notons $f_i : X \to X_i$ la projection et posons
$\mathcal{F} = \colim f_i^{-1}\mathcal{F}_i$. Alors
$$
\colim_{i\in I} R^p g_{i, *} \mathcal{F}_i = R^p g_* \mathcal{F}
$$
pour tout $p \geq 0$.
\end{lemma}

\begin{proof}
Rappelons (lemme \ref{etale-cohomology-lemma-higher-direct-images})
que $R^p g_{i, *} \mathcal{F}_i$ est le faisceau associ\'e au
pr\'efaisceau $U \mapsto H^p_\etale(U \times_S X_i, \mathcal{F}_i)$
et de m\^eme pour $R^pg_*\mathcal{F}$. De plus, la limite inductive d'un
syst\`eme de faisceaux est le faisceau associ\'e \`a la limite inductive calcul\'ee dans la cat\'egorie
des pr\'efaisceaux. Remarquons que tout objet de $S_\etale$ admet un recouvrement
par des objets quasi-compacts et quasi-s\'epar\'es (par exemple, des sch\'emas affines).
De plus, si $U$ est un objet quasi-compact et quasi-s\'epar\'e,
alors on a
$$
\colim H^p_\etale(U \times_S X_i, \mathcal{F}_i) =
H^p_\etale(U \times_S X, \mathcal{F})
$$
d'apr\`es le th\'eor\`eme \ref{etale-cohomology-theorem-colimit}. Le lemme en r\'esulte.
\end{proof}

\begin{lemma}
\label{etale-cohomology-lemma-relative-colimit-general}
Soit $I$ un ensemble ordonn\'e filtrant. Soit $g_i : X_i \to S_i$ un syst\`eme projectif de
morphismes de sch\'emas index\'e par $I$. Supposons que $g_i$ soit quasi-compact et
quasi-s\'epar\'e et que, pour $i' \geq i$, les morphismes de transition
$f_{i'i} : X_{i'} \to X_i$ et $h_{i'i} : S_{i'} \to S_i$ soient affines.
Soit $g : X \to S$ la limite des morphismes $g_i$, voir
Limites, section \ref{limits-section-limits}.
Notons $f_i : X \to X_i$ et $h_i : S \to S_i$ les projections.
Soit $(\mathcal{F}_i, \varphi_{i'i})$ un syst\`eme de faisceaux
sur $(X_i, f_{i'i})$. Posons $\mathcal{F} = \colim f_i^{-1}\mathcal{F}_i$. Alors
$$
R^p g_* \mathcal{F} =
\colim_{i \in I} h_i^{-1}R^p g_{i, *} \mathcal{F}_i
$$
pour tout $p \geq 0$.
\end{lemma}

\begin{proof}
Comment construit-on l'application du lemme ?
Pour $i' \geq i$, nous avons un diagramme commutatif
$$
\xymatrix{
X \ar[r]_{f_{i'}} \ar[d]_g &
X_{i'} \ar[r]_{f_{i'i}} \ar[d]_{g_{i'}} &
X_i \ar[d]^{g_i} \\
S \ar[r]^{h_{i'}} &
S_{i'} \ar[r]^{h_{i'i}} &
S_i
}
$$
Si nous composons l'application de changement de base
$h_{i'i}^{-1}Rg_{i, *}\mathcal{F}_i \to Rg_{i', *}f_{i'i}^{-1}\mathcal{F}_i$
(Cohomologie sur les sites, lemme
\ref{sites-cohomology-lemma-base-change-map-flat-case} ou
Remarque \ref{sites-cohomology-remark-base-change})
avec l'application $Rg_{i', *}\varphi_{i'i}$, nous obtenons alors
$\psi_{i'i} : h_{i' i}^{-1} R^p g_{i, *} \mathcal{F}_i \to
R^pg_{i', *} \mathcal{F}_{i'}$. De m\^eme, en utilisant le carr\'e de gauche
du diagramme, nous obtenons les applications
$\psi_i : h_i^{-1}R^pg_{i, *}\mathcal{F}_i \to R^pg_*\mathcal{F}$.
Les applications $h_{i'}^{-1}\psi_{i'i}$ et $\psi_i$ sont celles utilis\'ees dans
l'\'enonc\'e du lemme. Pour que cela ait un sens, nous devons v\'erifier que
$\psi_{i''i} = \psi_{i''i'} \circ h_{i''i'}^{-1}\psi_{i'i}$ et
$\psi_{i'} \circ h_{i'}^{-1}\psi_{i'i} = \psi_i$ ; cela r\'esulte
de Cohomologie sur les sites, remarque
\ref{sites-cohomology-remark-compose-base-change-horizontal}.

\medskip\noindent
D\'emonstration de l'\'egalit\'e. Premi\`ere d\'emonstration par
d\'ecalage de dimension\footnote{On peut aussi utiliser cette m\'ethode
pour construire les applications du lemme.}. Pour tout
$U$ affine et \'etale sur $X$, d'apr\`es le Th\'eor\`eme \ref{etale-cohomology-theorem-colimit},
nous avons
$$
g_*\mathcal{F}(U) =
H^0(U \times_S X, \mathcal{F}) =
\colim H^0(U_i \times_{S_i} X_i, \mathcal{F}_i) =
\colim g_{i, *}\mathcal{F}_i(U_i)
$$
o\`u la limite inductive porte sur les $i$ assez grands tels qu'
il existe un $i$ et un $U_i$ affine et \'etale sur $S_i$
dont le changement de base est $U$ sur $S$ (voir
Lemme \ref{etale-cohomology-lemma-colimit-affine-sites}).
Le membre de droite est \'egal \`a
$(\colim h_i^{-1}g_{i, *}\mathcal{F}_i)(U)$ d'apr\`es
Sites, Lemme \ref{sites-lemma-colimit}.
Cela d\'emontre le lemme pour $p = 0$.
Si $(\mathcal{G}_i, \varphi_{i'i})$ est un syst\`eme tel que
$\mathcal{G} = \colim f_i^{-1}\mathcal{G}_i$
et tel que $\mathcal{G}_i$ soit un faisceau ab\'elien injectif sur $X_i$
pour tout $i$, alors, pour tout $U$ affine et \'etale sur $X$, d'apr\`es
le Th\'eor\`eme \ref{etale-cohomology-theorem-colimit}, nous avons
$$
H^p(U \times_S X, \mathcal{G}) =
\colim H^p(U_i \times_{S_i} X_i, \mathcal{G}_i) = 0
$$
pour $p > 0$ (la m\^eme limite inductive que pr\'ec\'edemment). Ainsi $R^pg_*\mathcal{G} = 0$
et nous obtenons le r\'esultat pour $p > 0$ pour un tel syst\`eme.
En g\'en\'eral, nous pouvons choisir une suite exacte courte de syst\`emes
$$
0 \to (\mathcal{F}_i, \varphi_{i'i}) \to
(\mathcal{G}_i, \varphi_{i'i}) \to
(\mathcal{Q}_i, \varphi_{i'i}) \to 0
$$
o\`u $(\mathcal{G}_i, \varphi_{i'i})$ est comme ci-dessus, voir
Cohomologie sur les sites, Lemme
\ref{sites-cohomology-lemma-colim-sites-injective}.
Par r\'ecurrence, le lemme vaut pour $p - 1$ et, d'apr\`es ce qui pr\'ec\`ede, on a
l'annulation en degr\'e $p$ pour $(\mathcal{G}_i, \varphi_{i'i})$.
On en d\'eduit le r\'esultat pour $p$
et $(\mathcal{F}_i, \varphi_{i'i})$ par la suite exacte longue
de cohomologie.

\medskip\noindent
Seconde d\'emonstration.
Rappelons que $S_{affine, \etale} = \colim (S_i)_{affine, \etale}$, voir
le lemme \ref{etale-cohomology-lemma-colimit-affine-sites}. Ainsi, si $U$ est un objet de
$S_{affine, \etale}$, alors on peut \'ecrire $U = U_i \times_{S_i} S$
pour un certain $i$ et un certain $U_i$ de $(S_i)_{affine, \etale}$, et
$$
(\colim_{i \in I} h_i^{-1}R^p g_{i, *} \mathcal{F}_i)(U) =
\colim_{i' \geq i} (R^p g_{i', *}\mathcal{F}_{i'})(U_i \times_{S_i} S_{i'})
$$
d'apr\`es Sites, Lemme \ref{sites-lemma-colimit}, et la construction des
morphismes de transition du syst\`eme d\'ecrit ci-dessus. Puisque
$R^pg_{i', *}\mathcal{F}_{i'}$ est le faisceau associ\'e au pr\'efaisceau
$U_{i'} \mapsto H^p(U_{i'} \times_{S_{i'}} X_{i'}, \mathcal{F}_{i'})$
et puisque $R^pg_*\mathcal{F}$ est le faisceau associ\'e au pr\'efaisceau
$U \mapsto H^p(U \times_S X, \mathcal{F})$
(Lemme \ref{etale-cohomology-lemma-higher-direct-images})
on obtient un diagramme commutatif canonique
$$
\xymatrix{
\colim_{i' \geq i}
H^p(U_i \times_{S_i} X_{i'}, \mathcal{F}_{i'}) \ar[r] \ar[d] &
\colim_{i' \geq i}
(R^p g_{i', *}\mathcal{F}_{i'})(U_i \times_{S_i} S_{i'}) \ar[d] \\
H^p(U \times_S X, \mathcal{F}) \ar[r] &
R^pg_*\mathcal{F}(U)
}
$$
Remarquons que la fl\`eche verticale de gauche est un isomorphisme
d'apr\`es le th\'eor\`eme \ref{etale-cohomology-theorem-colimit}. Nous cherchons \`a montrer que
la fl\`eche verticale de droite est un isomorphisme. Cependant, nous
savons d\'ej\`a que la source et le but de cette fl\`eche
sont des faisceaux sur $S_{affine, \etale}$. Il suffit donc de
montrer : (1) tout \'el\'ement du but provient localement d'un
\'el\'ement de la source, et (2) tout \'el\'ement de la source
qui s'envoie sur z\'ero dans le but s'annule localement.
L'assertion (1) d\'ecoule imm\'ediatement de ce qui pr\'ec\`ede et du fait que
la fl\`eche horizontale inf\'erieure provient d'un morphisme de pr\'efaisceaux
qui devient un isomorphisme apr\`es faisceautisation.
Pour (2), soit $\xi \in  \colim_{i' \geq i}
(R^p g_{i', *}\mathcal{F}_{i'})(U_i \times_{S_i} S_{i'})$
dans le noyau. Choisissons $i' \geq i$ et
$\xi_{i'} \in (R^p g_{i', *}\mathcal{F}_{i'})(U_i \times_{S_i} S_{i'})$
repr\'esentant $\xi$.
Choisissons un recouvrement \'etale standard
$\{U_{i', k} \to U_i \times_{S_i} S_{i'}\}_{k = 1, \ldots, m}$
tel que $\xi_{i'}|_{U_{i', k}}$ provienne de
$\xi_{i', k} \in H^p(U_{i', k} \times_{S_{i'}} X_{i'}, \mathcal{F}_{i'})$.
Puisqu'il suffit de prouver que $\xi$ s'annule localement, nous
pouvons remplacer $U$ par les membres du recouvrement \'etale
$\{U_{i', k} \times_{S_{i'}} S \to U = U_i \times_{S_i} S\}$.
Apr\`es ce remplacement, on voit que $\xi$ est l'image
d'un \'el\'ement $\xi'$ du groupe
$\colim_{i' \geq i} H^p(U_i \times_{S_i} X_{i'}, \mathcal{F}_{i'})$
dans le diagramme ci-dessus. Puisque $\xi'$ s'envoie sur z\'ero dans $R^pg_*\mathcal{F}(U)$
nous pouvons effectuer un autre remplacement et supposer que $\xi'$ s'envoie
sur z\'ero dans $H^p(U \times_S X, \mathcal{F})$.
Toutefois, puisque la fl\`eche verticale de gauche est un isomorphisme,
on conclut alors que $\xi' = 0$, d'o\`u $\xi = 0$, comme voulu.
\end{proof}

\begin{lemma}
\label{etale-cohomology-lemma-linus-hamann}
Soit $X = \lim_{i \in I} X_i$ la limite d'un syst\`eme projectif de sch\'emas
dont les morphismes de transition $f_{i'i}$ sont affines et dont les morphismes de projection
sont $f_i : X \to X_i$. Soit $\mathcal{F}$ un faisceau sur $X_\etale$. Alors
\begin{enumerate}
\item il existe des morphismes canoniques
$\varphi_{i'i} : f_{i'i}^{-1}f_{i, *}\mathcal{F} \to f_{i', *}\mathcal{F}$
tels que $(f_{i, *}\mathcal{F}, \varphi_{i'i})$ soit un syst\`eme de
faisceaux sur $(X_i, f_{i'i})$ au sens de la
d\'efinition \ref{etale-cohomology-definition-inverse-system-sheaves}, et
\item $\mathcal{F} = \colim f_i^{-1}f_{i, *}\mathcal{F}$.
\end{enumerate}
\end{lemma}

\begin{proof}
Par Topologies, Lemme \ref{topologies-lemma-alternative}, et le
lemme \ref{etale-cohomology-lemma-colimit-affine-sites}, ceci est un cas particulier de
Sites, Lemme \ref{sites-lemma-colimit-push-pull}.
\end{proof}

\begin{lemma}
\label{etale-cohomology-lemma-compute-strangely}
Soit $I$ un ensemble ordonn\'e filtrant. Soit $g_i : X_i \to S_i$ un syst\`eme projectif de
morphismes de sch\'emas index\'e par $I$. Supposons que $g_i$ soit quasi-compact et
quasi-s\'epar\'e et que, pour $i' \geq i$, les morphismes de transition
$X_{i'} \to X_i$ et $S_{i'} \to S_i$ soient affines.
Soit $g : X \to S$ la limite des morphismes $g_i$, voir
Limites, section \ref{limits-section-limits}.
Notons $f_i : X \to X_i$ et $h_i : S \to S_i$ les projections.
Soit $\mathcal{F}$ un faisceau ab\'elien sur $X$. Alors on a
$$
R^pg_*\mathcal{F} = \colim_{i \in I} h_i^{-1}R^pg_{i, *}(f_{i, *}\mathcal{F})
$$
\end{lemma}

\begin{proof}
Combinaison formelle des lemmes \ref{etale-cohomology-lemma-relative-colimit-general}
et \ref{etale-cohomology-lemma-linus-hamann}.
\end{proof}











\section{Limites inductives et complexes}
\label{etale-cohomology-section-colimit-complexes}

\noindent
Dans cette section, nous \'etudions la cohomologie des syst\`emes de complexes
dans divers contextes, en poursuivant l'\'etude des faisceaux commenc\'ee
\`a la section \ref{etale-cohomology-section-colimit}.
Nous conseillons vivement au lecteur de ne pas lire cette section, sauf si cela est absolument
n\'ecessaire.

\begin{lemma}
\label{etale-cohomology-lemma-colimit-variant-complexes}
Soit $X = \lim_{i \in I} X_i$ la limite d'un syst\`eme projectif de sch\'emas
dont les morphismes de transition $f_{i'i} : X_{i'} \to X_i$ sont affines. Supposons
que $X_i$ soit quasi-compact et quasi-s\'epar\'e pour tout $i \in I$.
Soit $\mathcal{F}_i^\bullet$ un complexe de faisceaux ab\'eliens sur
$X_{i, \etale}$. Soit $\varphi_{i'i} : f_{i'i}^{-1}\mathcal{F}_i^\bullet \to
\mathcal{F}_{i'}^\bullet$ un morphisme de complexes sur $X_{i, \etale}$
tel que $\varphi_{i''i} = \varphi_{i''i'} \circ f_{i'' i'}^{-1}\varphi_{i'i}$
lorsque $i'' \geq i' \geq i$. Supposons qu'il existe un entier $a$ tel que
$\mathcal{F}_i^n = 0$ pour $n < a$ et tout $i \in I$.
Alors on a
$$
H^p_\etale(X, \colim f_i^{-1}\mathcal{F}_i^\bullet) =
\colim H^p_\etale(X_i, \mathcal{F}^\bullet_i)
$$
o\`u $f_i : X \to X_i$ est la projection.
\end{lemma}

\begin{proof}
Cela r\'esulte du th\'eor\`eme \ref{etale-cohomology-theorem-colimit}. Posons
$\mathcal{F}^\bullet = \colim f_i^{-1}\mathcal{F}_i^\bullet$.
Le th\'eor\`eme affirme que
$$
\colim_{i \in I} H_\etale^p(X_i, \mathcal{F}_i^n) =
H_\etale^p(X, \mathcal{F}^n)
$$
pour tous $n, p \in \mathbf{Z}$. Utilisons les suites spectrales
$$
E_{1, i}^{s, t} = H_\etale^t(X_i, \mathcal{F}_i^s) \Rightarrow
H_\etale^{s + t}(X_i, \mathcal{F}_i^\bullet)
$$
et
$$
E_1^{s, t} = H_\etale^t(X, \mathcal{F}^s) \Rightarrow
H_\etale^{s + t}(X, \mathcal{F}^\bullet)
$$
de Cat\'egories d\'eriv\'ees, Lemme \ref{derived-lemma-two-ss-complex-functor}.
Puisque $\mathcal{F}_i^n = 0$ pour $n < a$ (avec $a$ ind\'ependant de $i$),
on voit que seul un nombre fini de termes $E_{1, i}^{s, t}$
(ind\'ependant de $i$) et $E_1^{s, t}$ contribuent \`a
$H^q_\etale(X_i, \mathcal{F}_i^\bullet)$ et
$H^q_\etale(X, \mathcal{F}^\bullet)$ et $E_1^{s, t} = \colim E_{i, i}^{s, t}$.
Cela d\'emontre l'assertion. Certains d\'etails sont omis.
(Il existe un autre argument utilisant les troncatures ``b\^etes''
des complexes, qui \'evite d'utiliser les suites spectrales.)
\end{proof}

\begin{lemma}
\label{etale-cohomology-lemma-direct-sum-bounded-below-cohomology}
Soit $X$ un sch\'ema quasi-compact et quasi-s\'epar\'e. Soit
$K_i \in D(X_\etale)$, $i \in I$, une famille d'objets.
Supposons donn\'e $a \in \mathbf{Z}$ tel que $H^n(K_i) = 0$ pour $n < a$
et $i \in I$. Alors $R\Gamma(X, \bigoplus_i K_i) =
\bigoplus_i R\Gamma(X, K_i)$.
\end{lemma}

\begin{proof}
Il faut montrer que $H^p(X, \bigoplus_i K_i) =
\bigoplus_i H^p(X, K_i)$ pour tout $p \in \mathbf{Z}$.
Choisissons des complexes $\mathcal{F}_i^\bullet$ repr\'esentant $K_i$
tels que $\mathcal{F}_i^n = 0$ pour $n < a$.
La somme directe des complexes $\mathcal{F}_i^\bullet$
repr\'esente l'objet $\bigoplus K_i$ d'apr\`es
Injectifs, Lemme \ref{injectives-lemma-derived-products}.
Puisque $\bigoplus \mathcal{F}^\bullet$ est la limite inductive filtrante
des sommes directes finies, le r\'esultat d\'ecoule
du lemme \ref{etale-cohomology-lemma-colimit-variant-complexes}.
\end{proof}

\begin{lemma}
\label{etale-cohomology-lemma-colimit-complexes}
Soit $S$ un sch\'ema. Soit $X = \lim_{i \in I} X_i$ la limite d'un
syst\`eme projectif de sch\'emas sur $S$ avec pour morphismes de transition affines
$f_{i'i} : X_{i'} \to X_i$. Supposons que $X_i$ soit quasi-compact et
quasi-s\'epar\'e pour tout $i \in I$. Soit $K \in D^+(S_\etale)$. Alors
$$
\colim_{i \in I} H_\etale^p(X_i, K|_{X_i}) = H_\etale^p(X, K|_X).
$$
pour tout $p \in \mathbf{Z}$, o\`u $K|_{X_i}$ et $K|_X$
sont les images inverses de $K$ sur $X_i$ et $X$.
\end{lemma}

\begin{proof}
On peut repr\'esenter $K$ par un complexe born\'e inf\'erieurement $\mathcal{G}^\bullet$
de faisceaux ab\'eliens sur $S_\etale$. Disons que $\mathcal{G}^n = 0$ pour $n < a$.
Notons $\mathcal{F}^\bullet_i$ et $\mathcal{F}^\bullet$ les images inverses
de ce complexe sur $X_i$ et $X$. Ces complexes repr\'esentent les
objets $K|_{X_i}$ et $K|_X$, et l'on a
$\mathcal{F}^\bullet = \colim f_i^{-1}\mathcal{F}_i^\bullet$
terme \`a terme. Le lemme d\'ecoule donc du
lemme \ref{etale-cohomology-lemma-colimit-variant-complexes}.
\end{proof}

\begin{lemma}
\label{etale-cohomology-lemma-relative-colimit-general-complexes}
Soient $I$, $g_i : X_i \to S_i$, $g : X \to S$, $f_i$, $g_i$, $h_i$ comme dans le
lemme \ref{etale-cohomology-lemma-relative-colimit-general}.
Soient $0 \in I$ et $K_0 \in D^+(X_{0, \etale})$.
Pour $i \geq 0$, notons $K_i$ l'image inverse de $K_0$ sur $X_i$.
Notons $K$ l'image inverse de $K_0$ sur $X$. Alors
$$
R^pg_*K = \colim_{i \geq 0} h_i^{-1}R^pg_{i, *}K_i
$$
pour tout $p \in \mathbf{Z}$.
\end{lemma}

\begin{proof}
Fixons un entier $p_0 \in \mathbf{Z}$.
Soit $a$ un entier tel que $H^j(K_0) = 0$ pour $j < a$.
Nous allons d\'emontrer que la formule vaut pour tout $p \leq p_0$
par r\'ecurrence descendante sur $a$. Si
$a > p_0$, alors on voit que les membres de gauche et de droite
de la formule sont nuls pour $p \leq p_0$ par annulation triviale, voir
Cat\'egories d\'eriv\'ees, Lemme \ref{derived-lemma-negative-vanishing}.
Supposons $a \leq p_0$. Consid\'erons le triangle distingu\'e
$$
H^a(K_0)[-a] \to K_0 \to \tau_{\geq a + 1}K_0
$$
Le tir\'e en arri\`ere de ce triangle distingu\'e sur $X_i$ et $X$
donne des triangles distingu\'es compatibles pour $K_i$ et $K$.
Pour $p \leq p_0$, consid\'erons le diagramme commutatif
$$
\xymatrix{
\colim_{i \geq 0} h_i^{-1}R^{p - 1}g_{i, *}(\tau_{\geq a + 1}K_i)
\ar[r]_-\alpha \ar[d] &
R^{p - 1}g_*(\tau_{\geq a + 1}K) \ar[d] \\
\colim_{i \geq 0} h_i^{-1}R^pg_{i, *}(H^a(K_i)[-a])
\ar[r]_-\beta \ar[d] &
R^pg_*(H^a(K)[-a]) \ar[d] \\
\colim_{i \geq 0} h_i^{-1}R^pg_{i, *}K_i
\ar[r]_-\gamma \ar[d] &
R^pg_*K \ar[d] \\
\colim_{i \geq 0} R^pg_{i, *}\tau_{\geq a + 1}K_i
\ar[r]_-\delta \ar[d] &
R^pg_*\tau_{\geq a + 1}K \ar[d] \\
\colim_{i \geq 0} R^{p + 1}g_{i, *}(H^a(K_i)[-a])
\ar[r]^-\epsilon &
R^{p + 1}g_*(H^a(K)[-a])
}
$$
\`a colonnes exactes. Les fl\`eches $\beta$ et $\epsilon$ sont des isomorphismes d'apr\`es le
lemme \ref{etale-cohomology-lemma-relative-colimit-general}.
Les fl\`eches $\alpha$ et $\delta$ sont des isomorphismes
par l'hypoth\`ese de r\'ecurrence.
Ainsi, $\gamma$ est un isomorphisme, comme voulu.
\end{proof}

\begin{lemma}
\label{etale-cohomology-lemma-relative-colimit-general-really-complexes}
Soient $I$, $g_i : X_i \to S_i$, $g : X \to S$, $f_{ii'}$, $f_i$, $g_i$, $h_i$
comme dans le lemme \ref{etale-cohomology-lemma-relative-colimit-general}.
Soit $\mathcal{F}_i^\bullet$ un complexe de faisceaux ab\'eliens sur
$X_{i, \etale}$. Soit $\varphi_{i'i} : f_{i'i}^{-1}\mathcal{F}_i^\bullet \to
\mathcal{F}_{i'}^\bullet$ un morphisme de complexes sur $X_{i, \etale}$
tel que $\varphi_{i''i} = \varphi_{i''i'} \circ f_{i'' i'}^{-1}\varphi_{i'i}$
lorsque $i'' \geq i' \geq i$. Supposons qu'il existe un entier $a$ tel que
$\mathcal{F}_i^n = 0$ pour $n < a$ et tout $i \in I$. Alors
$$
R^pg_*(\colim f_i^{-1}\mathcal{F}_i^\bullet) =
\colim_{i \geq 0} h_i^{-1}R^pg_{i, *}\mathcal{F}_i^\bullet
$$
pour tout $p \in \mathbf{Z}$.
\end{lemma}

\begin{proof}
Ceci r\'esulte du lemme \ref{etale-cohomology-lemma-relative-colimit-general}. Posons
$\mathcal{F}^\bullet = \colim f_i^{-1}\mathcal{F}_i^\bullet$.
Le lemme affirme que
$$
\colim_{i \in I} h_i^{-1}R^pg_{i, *}\mathcal{F}_i^n = R^pg_*\mathcal{F}^n
$$
pour tous $n, p \in \mathbf{Z}$. Utilisons les suites spectrales
$$
E_{1, i}^{s, t} = R^tg_{i, *}\mathcal{F}_i^s \Rightarrow
R^{s + t}g_{i, *}\mathcal{F}_i^\bullet
$$
et
$$
E_1^{s, t} = R^tg_*\mathcal{F}^s \Rightarrow
R^{s + t}g_*\mathcal{F}^\bullet
$$
de Cat\'egories d\'eriv\'ees, Lemme \ref{derived-lemma-two-ss-complex-functor}.
Puisque $\mathcal{F}_i^n = 0$ pour $n < a$ (avec $a$ ind\'ependant de $i$),
on voit que seul un nombre fini fix\'e de termes $E_{1, i}^{s, t}$
(ind\'ependant de $i$) et $E_1^{s, t}$ contribuent et que
$E_1^{s, t} = \colim E_{i, i}^{s, t}$.
Cela d\'emontre l'assertion. Certains d\'etails sont omis.
(Il existe un autre argument utilisant les troncatures ``b\^etes''
des complexes, qui \'evite d'utiliser les suites spectrales.)
\end{proof}

\begin{lemma}
\label{etale-cohomology-lemma-direct-sum-bounded-below-pushforward}
Soit $f : X \to Y$ un morphisme de sch\'emas quasi-compact et quasi-s\'epar\'e.
Soit $K_i \in D(X_\etale)$, $i \in I$, une famille d'objets.
Supposons donn\'e $a \in \mathbf{Z}$ tel que $H^n(K_i) = 0$ pour $n < a$
et $i \in I$. Alors $Rf_*(\bigoplus_i K_i) = \bigoplus_i Rf_*K_i$.
\end{lemma}

\begin{proof}
Il faut montrer que $R^pf_*(\bigoplus_i K_i) =
\bigoplus_i R^pf_*K_i$ pour tout $p \in \mathbf{Z}$.
Choisissons des complexes $\mathcal{F}_i^\bullet$ repr\'esentant $K_i$
tels que $\mathcal{F}_i^n = 0$ pour $n < a$.
La somme directe des complexes $\mathcal{F}_i^\bullet$
repr\'esente l'objet $\bigoplus K_i$ d'apr\`es
Injectifs, Lemme \ref{injectives-lemma-derived-products}.
Puisque $\bigoplus \mathcal{F}^\bullet$ est la limite inductive filtrante
des sommes directes finies, le r\'esultat d\'ecoule
du lemme \ref{etale-cohomology-lemma-relative-colimit-general-really-complexes}.
\end{proof}














\section{Fibres des images directes sup\'erieures}
\label{etale-cohomology-section-stalks-direct-image}

\noindent
Les fibres des images directes sup\'erieures se calculent souvent comme suit.

\begin{theorem}
\label{etale-cohomology-theorem-higher-direct-images}
Soit $f: X \to S$ un morphisme de sch\'emas quasi-compact et quasi-s\'epar\'e,
$\mathcal{F}$ un faisceau ab\'elien sur $X_\etale$, et $\overline{s}$ un
point g\'eom\'etrique de $S$ au-dessus de $s \in S$. Alors
$$
\left(R^nf_* \mathcal{F}\right)_{\overline{s}} =
H_\etale^n( X \times_S \Spec(\mathcal{O}_{S, \overline{s}}^{sh}),
p^{-1}\mathcal{F})
$$
o\`u $p : X \times_S \Spec(\mathcal{O}_{S, \overline{s}}^{sh}) \to X$
est la projection. Pour $K \in D^+(X_\etale)$ et $n \in \mathbf{Z}$,
on a
$$
\left(R^nf_*K\right)_{\overline{s}} =
H_\etale^n(X \times_S \Spec(\mathcal{O}_{S, \overline{s}}^{sh}), p^{-1}K)
$$
En fait, on a
$$
\left(Rf_*K\right)_{\overline{s}}
=
R\Gamma_\etale(X \times_S \Spec(\mathcal{O}_{S, \overline{s}}^{sh}), p^{-1}K)
$$
dans $D^+(\textit{Ab})$.
\end{theorem}

\begin{proof}
Soit $\mathcal{I}$ la cat\'egorie des voisinages \'etales de $\overline{s}$
sur $S$. D'apr\`es le lemme \ref{etale-cohomology-lemma-higher-direct-images},
on a
$$
(R^nf_*\mathcal{F})_{\overline{s}} =
\colim_{(V, \overline{v}) \in \mathcal{I}^{opp}}
H_\etale^n(X \times_S V, \mathcal{F}|_{X \times_S V}).
$$
Nous pouvons remplacer $\mathcal{I}$ par la sous-cat\'egorie initiale form\'ee
des voisinages \'etales affines de $\overline{s}$. Remarquons que
$$
\Spec(\mathcal{O}_{S, \overline{s}}^{sh}) =
\lim_{(V, \overline{v}) \in \mathcal{I}} V
$$
d'apr\`es le lemme \ref{etale-cohomology-lemma-describe-etale-local-ring} et
Limites, Lemme
\ref{limits-lemma-directed-inverse-system-affine-schemes-has-limit}.
Puisque les produits fibr\'es commutent aux limites, on obtient aussi
$$
X \times_S \Spec(\mathcal{O}_{S, \overline{s}}^{sh}) =
\lim_{(V, \overline{v}) \in \mathcal{I}} X \times_S V
$$
On conclut par le lemme \ref{etale-cohomology-lemma-directed-colimit-cohomology}.
Pour la seconde variante, on emploie le m\^eme argument avec
le lemme \ref{etale-cohomology-lemma-colimit-complexes} au lieu du
lemme \ref{etale-cohomology-lemma-directed-colimit-cohomology}.

\medskip\noindent
Pour voir que la derni\`ere assertion est vraie, il suffit de construire un morphisme
$\left(Rf_*K\right)_{\overline{s}} \to
R\Gamma_\etale(X \times_S \Spec(\mathcal{O}_{S, \overline{s}}^{sh}), p^{-1}K)$
dans $D^+(\textit{Ab})$ qui r\'ealise les isomorphismes ci-dessus sur les groupes de cohomologie
en degr\'e $n$ pour tout $n$. Pour cela, choisissons un
complexe born\'e inf\'erieurement $\mathcal{J}^\bullet$ de faisceaux ab\'eliens injectifs
sur $X_\etale$ repr\'esentant $K$. Le complexe
$f_*\mathcal{J}^\bullet$ repr\'esente $Rf_*K$. Ainsi, le complexe
$$
(f_*\mathcal{J}^\bullet)_{\overline{s}} =
\colim_{(V, \overline{v}) \in \mathcal{I}^{opp}}
(f_*\mathcal{J}^\bullet)(V)
$$
repr\'esente $(Rf_*K)_{\overline{s}}$. Pour chaque $V$, on a des morphismes
$$
(f_*\mathcal{J}^\bullet)(V) =
\Gamma(X \times_S V, \mathcal{J}^\bullet)
\longrightarrow
\Gamma(X \times_S \Spec(\mathcal{O}_{S, \overline{s}}^{sh}),
p^{-1}\mathcal{J}^\bullet)
$$
et le complexe cible repr\'esente
$R\Gamma_\etale(X \times_S \Spec(\mathcal{O}_{S, \overline{s}}^{sh}), p^{-1}K)$
dans $D^+(\textit{Ab})$.
En prenant la limite inductive de ces morphismes, on obtient le r\'esultat.
\end{proof}

\begin{remark}
\label{etale-cohomology-remark-stalk-fibre}
Soit $f : X \to S$ un morphisme de sch\'emas.
Soit $K \in D(X_\etale)$.
Soit $\overline{s}$ un point g\'eom\'etrique de $S$.
Il existe toujours des morphismes canoniques
$$
(Rf_*K)_{\overline{s}}
\longrightarrow
R\Gamma(X \times_S \Spec(\mathcal{O}_{S, \overline{s}}^{sh}), p^{-1}K)
\longrightarrow
R\Gamma(X_{\overline{s}}, K|_{X_{\overline{s}}})
$$
o\`u $p : X \times_S \Spec(\mathcal{O}_{S, \overline{s}}^{sh}) \to X$
est la projection. Consid\'erons en effet le diagramme commutatif
$$
\xymatrix{
X_{\overline{s}} \ar[r] \ar[d]^{f_{\overline{s}}} &
X \times_S \Spec(\mathcal{O}_{S, \overline{s}}^{sh}) \ar[r]_-p \ar[d]^{f'} &
X \ar[d]^f \\
\overline{s} \ar[r]^-i &
\Spec(\mathcal{O}_{S, \overline{s}}^{sh}) \ar[r]^-j &
S
}
$$
On dispose des morphismes de changement de base
$$
i^{-1}Rf'_*(p^{-1}K) \to Rf_{\overline{s}, *}(K|_{X_{\overline{s}}})
\quad\text{et}\quad
j^{-1}Rf_*K \to Rf'_*(p^{-1}K)
$$
(Cohomologie sur les sites, Remarque \ref{sites-cohomology-remark-base-change})
associ\'es aux deux carr\'es de ce diagramme.
En prenant les sections globales, on obtient les morphismes voulus.
D'apr\`es Cohomologie sur les sites, Remarque
\ref{sites-cohomology-remark-compose-base-change-horizontal}
la compos\'ee de ces deux morphismes est le
morphisme usuel de changement de base
$(Rf_*K)_{\overline{s}} \to R\Gamma(X_{\overline{s}}, K|_{X_{\overline{s}}})$.
\end{remark}






\section{La suite spectrale de Leray}
\label{etale-cohomology-section-leray}

\begin{lemma}
\label{etale-cohomology-lemma-prepare-leray}
Soient $f: X \to Y$ un morphisme et $\mathcal{I}$ un objet injectif de
$\textit{Ab}(X_\etale)$. Soit $V \in \Ob(Y_\etale)$. Alors
\begin{enumerate}
\item pour tout recouvrement $\mathcal{V} = \{V_j\to V\}_{j \in J}$, on a
$\check H^p(\mathcal{V}, f_*\mathcal{I}) = 0$ pour tout $p > 0$ ;
\item $f_*\mathcal{I}$ est acyclique pour le foncteur $\Gamma(V, -)$ ; et
\item si $g : Y \to Z$, alors $f_*\mathcal{I}$ est acyclique pour $g_*$. 
\end{enumerate}
\end{lemma}

\begin{proof}
Remarquons que $\check{\mathcal{C}}^\bullet(\mathcal{V}, f_*\mathcal{I}) =
\check{\mathcal{C}}^\bullet(\mathcal{V} \times_Y X, \mathcal{I})$
dont les groupes de cohomologie sup\'erieure sont nuls d'apr\`es le lemme \ref{etale-cohomology-lemma-hom-injective}.
Cela d\'emontre (1). La deuxi\`eme assertion d\'ecoule du fait qu'un faisceau dont
la cohomologie sup\'erieure de {\v C}ech s'annule pour tout recouvrement a des groupes de
cohomologie sup\'erieure nuls. C'est un excellent exercice d'utilisation de la
suite spectrale de {\v C}ech vers la cohomologie, mais voir
Cohomologie sur les sites, Lemme \ref{sites-cohomology-lemma-cech-vanish-collection}
pour les d\'etails et pour un \'enonc\'e plus pr\'ecis et plus g\'en\'eral.
L'assertion (3) r\'esulte de (2) et de la description de
$R^pg_*$ dans le lemme \ref{etale-cohomology-lemma-higher-direct-images}.
\end{proof}

\noindent
Le formalisme des suites spectrales de Grothendieck donne alors le
r\'esultat suivant.

\begin{proposition}[Suite spectrale de Leray]
\label{etale-cohomology-proposition-leray}
Soit $f: X \to Y$ un morphisme de sch\'emas et soit $\mathcal{F}$ un faisceau \'etale sur
$X$. Il existe alors une suite spectrale
$$
E_2^{p, q} = H_\etale^p(Y, R^qf_*\mathcal{F}) \Rightarrow
H_\etale^{p+q}(X, \mathcal{F}).
$$
\end{proposition}

\begin{proof}
Voir le lemme \ref{etale-cohomology-lemma-prepare-leray} et
Cat\'egories d\'eriv\'ees, section
\ref{derived-section-composition-right-derived-functors}.
\end{proof}









\section{Annulation des images directes sup\'erieures finies}
\label{etale-cohomology-section-vanishing-finite-morphism}

\noindent
Le but suivant est de d\'emontrer que les images directes sup\'erieures d'un morphisme fini de
sch\'emas s'annulent.

\begin{lemma}
\label{etale-cohomology-lemma-vanishing-etale-cohomology-strictly-henselian}
Soit $R$ un anneau local strictement hens\'elien. Posons $S = \Spec(R)$ et soit
$\overline{s}$ son point ferm\'e. Alors le foncteur des
sections globales
$\Gamma(S, -) : \textit{Ab}(S_\etale) \to \textit{Ab}$
est exact. En fait, $\Gamma(S, \mathcal{F}) = \mathcal{F}_{\overline{s}}$
pour tout faisceau d'ensembles $\mathcal{F}$. En particulier,
$$
\forall p\geq 1, \quad H_\etale^p(S, \mathcal{F})=0
$$
pour tout $\mathcal{F}\in \textit{Ab}(S_\etale)$.
\end{lemma}

\begin{proof}
Si l'on montre que $\Gamma(S, \mathcal{F}) = \mathcal{F}_{\overline{s}}$,
alors $\Gamma(S, -)$ est exact, puisque le foncteur fibre est exact.
Soit $(U, \overline{u})$ un voisinage \'etale de $\overline{s}$.
Choisissons un voisinage ouvert affine $\Spec(A)$ de $\overline{u}$ dans $U$.
Alors $R \to A$ est \'etale et $\kappa(\overline{s}) = \kappa(\overline{u})$.
D'apr\`es le th\'eor\`eme \ref{etale-cohomology-theorem-henselian}, on voit que $A \cong R \times A'$
comme $R$-alg\`ebre, de mani\`ere compatible avec les morphismes vers
$\kappa(\overline{s}) = \kappa(\overline{u})$.
On obtient donc une section
$$
\xymatrix{
\Spec(A) \ar[r] & U \ar[d]\\
& S \ar[ul]
}
$$
Il s'ensuit que, dans le syst\`eme des voisinages \'etales de $\overline{s}$,
le morphisme identit\'e $(S, \overline{s}) \to (S, \overline{s})$ est cofinal.
Ainsi $\Gamma(S, \mathcal{F}) = \mathcal{F}_{\overline{s}}$.
La derni\`ere assertion du lemme r\'esulte du fait que les foncteurs d\'eriv\'es sup\'erieurs
d'un foncteur exact sont nuls, voir
Cat\'egories d\'eriv\'ees, Lemme \ref{derived-lemma-right-derived-exact-functor}.
\end{proof}

\begin{proposition}
\label{etale-cohomology-proposition-finite-higher-direct-image-zero}
Soit $f : X \to Y$ un morphisme fini de sch\'emas.
\begin{enumerate}
\item Pour tout point g\'eom\'etrique $\overline{y} : \Spec(k) \to Y$, on a
$$
(f_*\mathcal{F})_{\overline{y}} =
\prod\nolimits_{\overline{x} : \Spec(k) \to X,\ f(\overline{x}) =
\overline{y}} \mathcal{F}_{\overline{x}}.
$$
pour $\mathcal{F}$ dans $\Sh(X_\etale)$ et
$$
(f_*\mathcal{F})_{\overline{y}} =
\bigoplus\nolimits_{\overline{x} : \Spec(k) \to X,\ f(\overline{x}) =
\overline{y}} \mathcal{F}_{\overline{x}}.
$$
pour $\mathcal{F}$ dans $\textit{Ab}(X_\etale)$.
\item Pour tout $q \geq 1$, on a $R^q f_*\mathcal{F} = 0$
pour $\mathcal{F}$ dans $\textit{Ab}(X_\etale)$.
\end{enumerate}
\end{proposition}

\begin{proof}
Notons $X_{\overline{y}}^{sh}$ le produit fibr\'e
$X \times_Y \Spec(\mathcal{O}_{Y, \overline{y}}^{sh})$.
D'apr\`es le th\'eor\`eme \ref{etale-cohomology-theorem-higher-direct-images},
la fibre de $R^qf_*\mathcal{F}$ en $\overline{y}$ se calcule par
$H_\etale^q(X_{\overline{y}}^{sh}, \mathcal{F})$.
Puisque $f$ est fini, $X_{\bar y}^{sh}$ est fini sur
$\Spec(\mathcal{O}_{Y, \overline{y}}^{sh})$ ; ainsi,
$X_{\bar y}^{sh} = \Spec(A)$ pour un anneau $A$
fini sur $\mathcal{O}_{Y, \bar y}^{sh}$.
Comme ce dernier est strictement hens\'elien,
le lemme \ref{etale-cohomology-lemma-finite-over-henselian}
implique que $A$ est un produit fini d'anneaux locaux hens\'eliens
$A = A_1 \times \ldots \times A_r$. Puisque le corps r\'esiduel de
$\mathcal{O}_{Y, \overline{y}}^{sh}$ est s\'eparablement clos, il en va de
m\^eme pour chaque $A_i$. Ainsi $A_i$ est strictement hens\'elien.
Il s'ensuit que $X_{\overline{y}}^{sh} = \coprod_{i = 1}^r \Spec(A_i)$.
L'annulation du
lemme \ref{etale-cohomology-lemma-vanishing-etale-cohomology-strictly-henselian}
implique que $(R^qf_*\mathcal{F})_{\overline{y}} = 0$ pour $q > 0$ ;
cela implique (2) d'apr\`es le th\'eor\`eme \ref{etale-cohomology-theorem-exactness-stalks}.
L'assertion (1) r\'esulte de l'assertion correspondante du
lemme \ref{etale-cohomology-lemma-vanishing-etale-cohomology-strictly-henselian}.
\end{proof}

\begin{lemma}
\label{etale-cohomology-lemma-finite-pushforward-commutes-with-base-change}
Consid\'erons un carr\'e cart\'esien
$$
\xymatrix{
X' \ar[r]_{g'} \ar[d]_{f'} & X \ar[d]^f \\
Y' \ar[r]^g & Y
}
$$
de sch\'emas, avec $f$ un morphisme fini.
Pour tout faisceau d'ensembles $\mathcal{F}$ sur $X_\etale$, on a
$f'_*(g')^{-1}\mathcal{F} = g^{-1}f_*\mathcal{F}$.
\end{lemma}

\begin{proof}
Dans une grande g\'en\'eralit\'e, il existe un morphisme de changement de base
$g^{-1}f_*\mathcal{F} \to f'_*(g')^{-1}\mathcal{F}$, voir
Sites, section \ref{sites-section-pullback}.
Il suffit de v\'erifier l'assertion sur les fibres (th\'eor\`eme \ref{etale-cohomology-theorem-exactness-stalks}).
Soit $\overline{y}' : \Spec(k) \to Y'$ un point g\'eom\'etrique.
On a
\begin{align*}
(f'_*(g')^{-1}\mathcal{F})_{\overline{y}'}
& =
\prod\nolimits_{\overline{x}' : \Spec(k) \to X',\ f' \circ \overline{x}' =
\overline{y}'}
((g')^{-1}\mathcal{F})_{\overline{x}'} \\
& =
\prod\nolimits_{\overline{x}' : \Spec(k) \to X',\ f' \circ \overline{x}' =
\overline{y}'} \mathcal{F}_{g' \circ \overline{x}'} \\
& =
\prod\nolimits_{\overline{x} : \Spec(k) \to X,\ f \circ \overline{x} =
g \circ \overline{y}'} \mathcal{F}_{\overline{x}} \\
& =
(f_*\mathcal{F})_{g \circ \overline{y}'} \\
& =
(g^{-1}f_*\mathcal{F})_{\overline{y}'}
\end{align*}
La premi\`ere \'egalit\'e r\'esulte de la
proposition \ref{etale-cohomology-proposition-finite-higher-direct-image-zero}.
La deuxi\`eme \'egalit\'e r\'esulte du
lemme \ref{etale-cohomology-lemma-stalk-pullback}.
La troisi\`eme \'egalit\'e vaut parce que le diagramme est un carr\'e cart\'esien
et que, par cons\'equent, l'application
$$
\{\overline{x}' : \Spec(k) \to X',\ f' \circ \overline{x}' =
\overline{y}'\}
\longrightarrow
\{\overline{x} : \Spec(k) \to X,\ f \circ \overline{x} =
g \circ \overline{y}'\}
$$
qui envoie $\overline{x}'$ sur $g' \circ \overline{x}'$ est bijective.
La quatri\`eme \'egalit\'e r\'esulte de la
proposition \ref{etale-cohomology-proposition-finite-higher-direct-image-zero}.
La cinqui\`eme \'egalit\'e r\'esulte du
lemme \ref{etale-cohomology-lemma-stalk-pullback}.
\end{proof}

\begin{lemma}
\label{etale-cohomology-lemma-integral-pushforward-commutes-with-base-change}
Consid\'erons un carr\'e cart\'esien
$$
\xymatrix{
X' \ar[r]_{g'} \ar[d]_{f'} & X \ar[d]^f \\
Y' \ar[r]^g & Y
}
$$
de sch\'emas, avec $f$ un morphisme entier.
Pour tout faisceau d'ensembles $\mathcal{F}$ sur $X_\etale$, on a
$f'_*(g')^{-1}\mathcal{F} = g^{-1}f_*\mathcal{F}$.
\end{lemma}

\begin{proof}
La question est locale sur $Y$ ; nous pouvons donc supposer $Y$ affine.
On peut alors \'ecrire $X = \lim X_i$, o\`u $f_i : X_i \to Y$ est fini
(c'est imm\'ediat dans le cas affine, mais voir
Limites, Lemme \ref{limits-lemma-integral-limit-finite-and-finite-presentation}
pour une r\'ef\'erence). Notons $p_{i'i} : X_{i'} \to X_i$
les morphismes de transition et $p_i : X \to X_i$ les projections.
En posant $\mathcal{F}_i = p_{i, *}\mathcal{F}$, on obtient
du lemme \ref{etale-cohomology-lemma-linus-hamann}
un syst\`eme $(\mathcal{F}_i, \varphi_{i'i})$
tel que $\mathcal{F} = \colim p_i^{-1}\mathcal{F}_i$.
On obtient $f_*\mathcal{F} = \colim f_{i, *}\mathcal{F}_i$
par le lemme \ref{etale-cohomology-lemma-relative-colimit}.
Posons $X'_i = Y' \times_Y X_i$, de projections $f'_i$ et $g'_i$.
Alors $X' = \lim X'_i$, puisque les limites commutent aux limites.
Notons $p'_i : X' \to X'_i$ les projections. On a
\begin{align*}
g^{-1}f_*\mathcal{F}
& =
g^{-1} \colim f_{i, *}\mathcal{F}_i \\
& =
\colim g^{-1}f_{i, *}\mathcal{F}_i \\
& =
\colim f'_{i, *}(g'_i)^{-1}\mathcal{F}_i \\
& =
f'_*(\colim (p'_i)^{-1}(g'_i)^{-1}\mathcal{F}_i) \\
& =
f'_*(\colim (g')^{-1}p_i^{-1}\mathcal{F}_i) \\
& =
f'_*(g')^{-1} \colim p_i^{-1}\mathcal{F}_i \\
& =
f'_*(g')^{-1}\mathcal{F}
\end{align*}
comme voulu. Pour la premi\`ere \'egalit\'e, voir ci-dessus.
Pour la deuxi\`eme, on utilise le fait que l'image inverse commute aux limites inductives.
Pour la troisi\`eme, on utilise le cas fini, voir
le lemme \ref{etale-cohomology-lemma-finite-pushforward-commutes-with-base-change}.
Pour la quatri\`eme, on utilise le lemme \ref{etale-cohomology-lemma-relative-colimit}.
Pour la cinqui\`eme, on utilise $g'_i \circ p'_i = p_i \circ g'$.
Pour la sixi\`eme, on utilise le fait que l'image inverse commute aux limites inductives.
Pour la septi\`eme, on utilise $\mathcal{F} = \colim p_i^{-1}\mathcal{F}_i$.
\end{proof}

\noindent
Le lemme suivant est un cas de descente cohomologique portant sur
les faisceaux \'etales et les morphismes finis surjectifs. Nous g\'en\'eraliserons
consid\'erablement ce r\'esultat une fois d\'emontr\'e le th\'eor\`eme de changement de base propre.

\begin{lemma}
\label{etale-cohomology-lemma-cohomological-descent-finite}
Soit $f : X \to Y$ un morphisme fini surjectif de sch\'emas.
Posons $f_n : X_n \to Y$ \'egal au produit fibr\'e it\'er\'e $(n + 1)$ fois
de $X$ au-dessus de $Y$. Pour $\mathcal{F} \in \textit{Ab}(Y_\etale)$, posons
$\mathcal{F}_n = f_{n, *}f_n^{-1}\mathcal{F}$. Il existe une suite
exacte
$$
0 \to \mathcal{F} \to \mathcal{F}_0 \to \mathcal{F}_1 \to
\mathcal{F}_2 \to \ldots
$$
sur $Y_\etale$. De plus, il existe une suite spectrale
$$
E_1^{p, q} = H^q_\etale(X_p, f_p^{-1}\mathcal{F})
$$
convergeant vers $H^{p + q}(Y_\etale, \mathcal{F})$.
Cette suite spectrale est fonctorielle en $\mathcal{F}$.
\end{lemma}

\begin{proof}
Si nous d\'emontrons la premi\`ere assertion du lemme, nous obtenons une suite
spectrale de terme $E_1^{p, q} = H^q_\etale(Y, \mathcal{F}_p)$ convergeant
vers $H^{p + q}(Y_\etale, \mathcal{F})$, voir
Cat\'egories d\'eriv\'ees, lemme \ref{derived-lemma-two-ss-complex-functor}.
D'autre part, puisque
$R^if_{p, *}f_p^{-1}\mathcal{F} = 0$ pour $i > 0$
(proposition \ref{etale-cohomology-proposition-finite-higher-direct-image-zero})
on obtient
$$
H^q_\etale(X_p, f_p^{-1}\mathcal{F}) =
H^q_\etale(Y, f_{p, *}f_p^{-1} \mathcal{F}) =
H^q_\etale(Y, \mathcal{F}_p)
$$
d'apr\`es la proposition \ref{etale-cohomology-proposition-leray},
ce qui donne la suite spectrale du lemme.

\medskip\noindent
Pour d\'emontrer la premi\`ere assertion du lemme, observons que
les $X_n$ forment un sch\'ema simplicial au-dessus de $Y$, voir
Simplicial, exemple \ref{simplicial-example-fibre-products-simplicial-object}.
Observons de plus que, pour chacune des projections
$d_j : X_{n + 1} \to X_n$, il existe un morphisme
$d_j^{-1} f_n^{-1}\mathcal{F} \to f_{n + 1}^{-1}\mathcal{F}$.
Ces morphismes induisent des morphismes
$$
\delta_j : \mathcal{F}_n \to \mathcal{F}_{n + 1}
$$
pour $j = 0, \ldots, n + 1$. Nous utilisons la somme altern\'ee de ces morphismes
pour d\'efinir les diff\'erentielles $\mathcal{F}_n \to \mathcal{F}_{n + 1}$.
De m\^eme, il existe une augmentation canonique $\mathcal{F} \to \mathcal{F}_0$,
qui est simplement le morphisme canonique $\mathcal{F} \to f_*f^{-1}\mathcal{F}$.
Pour v\'erifier que cette suite de faisceaux est un complexe exact, il suffit
de le v\'erifier sur les fibres aux points g\'eom\'etriques (th\'eor\`eme \ref{etale-cohomology-theorem-exactness-stalks}).
Soit donc $\overline{y} : \Spec(k) \to Y$ un point g\'eom\'etrique. Posons
$E = \{\overline{x} : \Spec(k) \to X \mid f(\overline{x}) = \overline{y}\}$.
Alors $E$ est un ensemble fini non vide et l'on voit que
$$
(\mathcal{F}_n)_{\overline{y}} =
\bigoplus\nolimits_{e \in E^{n + 1}} \mathcal{F}_{\overline{y}}
$$
d'apr\`es la proposition \ref{etale-cohomology-proposition-finite-higher-direct-image-zero}
et le lemme \ref{etale-cohomology-lemma-stalk-pullback}.
Il reste donc \`a voir que, pour tout groupe ab\'elien $M$, la suite
$$
0 \to M \to \bigoplus\nolimits_{e \in E} M \to
\bigoplus\nolimits_{e \in E^2} M \to \ldots
$$
est exacte. Ici, le premier morphisme est le morphisme diagonal et le morphisme
$\bigoplus_{e \in E^{n + 1}} M  \to \bigoplus_{e \in E^{n + 2}} M$
est la somme altern\'ee des morphismes induits par les $(n + 2)$
projections $E^{n + 2} \to E^{n + 1}$. Cela se d\'emontre directement
ou se d\'eduit en appliquant Simplicial, lemme
\ref{simplicial-lemma-fibre-products-simplicial-object-w-section}
au morphisme $E \to \{*\}$.
\end{proof}

\begin{remark}
\label{etale-cohomology-remark-cohomological-descent-finite}
Dans la situation du lemme \ref{etale-cohomology-lemma-cohomological-descent-finite},
si $\mathcal{G}$ est un faisceau d'ensembles sur $Y_\etale$, alors on a
$$
\Gamma(Y, \mathcal{G}) =
\text{\'Egaliseur}(
\xymatrix{
\Gamma(X_0, f_0^{-1}\mathcal{G})
\ar@<1ex>[r] \ar@<-1ex>[r] &
\Gamma(X_1, f_1^{-1}\mathcal{G})
}
)
$$
Cela se d\'emontre exactement de la m\^eme mani\`ere, en montrant que
le faisceau $\mathcal{G}$ est l'\'egalisateur des deux morphismes
$f_{0, *}f_0^{-1}\mathcal{G} \to f_{1, *}f_1^{-1}\mathcal{G}$.
\end{remark}









%10.06.09
\section{Action galoisienne sur les fibres}
\label{etale-cohomology-section-galois-action-stalks}

\noindent
Dans cette section, nous d\'efinissons une action du groupe de Galois absolu du corps
r\'esiduel d'un point $s$ de $S$ sur le foncteur fibre en tout point g\'eom\'etrique situ\'e
au-dessus de $s$.

\medskip\noindent
Action galoisienne sur les fibres.
Soit $S$ un sch\'ema.
Soit $\overline{s}$ un point g\'eom\'etrique de $S$.
Soit $\sigma \in \text{Aut}(\kappa(\overline{s})/\kappa(s))$.
D\'efinissons comme suit une action de $\sigma$ sur la fibre $\mathcal{F}_{\overline{s}}$
d'un faisceau $\mathcal{F}$
\begin{equation}
\label{etale-cohomology-equation-galois-action}
\begin{matrix}
\mathcal{F}_{\overline{s}} &
\longrightarrow &
\mathcal{F}_{\overline{s}} \\
(U, \overline{u}, t) &
\longmapsto &
(U, \overline{u} \circ \Spec(\sigma), t).
\end{matrix}
\end{equation}
o\`u nous utilisons la description des \'el\'ements de la fibre au moyen de triplets,
comme dans la discussion qui suit la
d\'efinition \ref{etale-cohomology-definition-stalk}.
Il s'agit d'une action \`a gauche, car si
$\sigma_i \in \text{Aut}(\kappa(\overline{s})/\kappa(s))$
alors
\begin{align*}
\sigma_1 \cdot (\sigma_2 \cdot (U, \overline{u}, t))
& =
\sigma_1 \cdot (U, \overline{u} \circ \Spec(\sigma_2), t) \\
& =
(U, \overline{u} \circ \Spec(\sigma_2) \circ \Spec(\sigma_1), t) \\
& =
(U, \overline{u} \circ \Spec(\sigma_1 \circ \sigma_2), t) \\
& =
(\sigma_1 \circ \sigma_2) \cdot (U, \overline{u}, t)
\end{align*}
Il est clair que cette action est fonctorielle par rapport au faisceau $\mathcal{F}$.
Notons que nous aurions pu d\'efinir cette action en renvoyant directement \`a la
remarque \ref{etale-cohomology-remark-map-stalks}.

\begin{definition}
\label{etale-cohomology-definition-algebraic-geometric-point}
Soit $S$ un sch\'ema.
Soit $\overline{s}$ un point g\'eom\'etrique situ\'e au-dessus du point $s$ de $S$.
Notons $\kappa(s) \subset \kappa(s)^{sep} \subset \kappa(\overline{s})$
la cl\^oture s\'eparable de $\kappa(s)$ dans le corps alg\'ebriquement
clos $\kappa(\overline{s})$.
\begin{enumerate}
\item Dans cette situation, le {\it groupe de Galois absolu} de $\kappa(s)$
est $\text{Gal}(\kappa(s)^{sep}/\kappa(s))$. On le note parfois
$\text{Gal}_{\kappa(s)}$.
\item Le point g\'eom\'etrique $\overline{s}$ est dit
{\it alg\'ebrique} si $\kappa(s) \subset \kappa(\overline{s})$ est
une cl\^oture alg\'ebrique de $\kappa(s)$.
\end{enumerate}
\end{definition}

\begin{example}
\label{etale-cohomology-example-stupid}
Le point g\'eom\'etrique
$\Spec(\mathbf{C}) \to \Spec(\mathbf{Q})$
n'est pas alg\'ebrique.
\end{example}

\noindent
Soit $\kappa(s) \subset \kappa(s)^{sep} \subset \kappa(\overline{s})$
comme dans la d\'efinition. Puisque $\kappa(\overline{s})$ est alg\'ebriquement
clos, le morphisme
$$
\text{Aut}(\kappa(\overline{s})/\kappa(s))
\longrightarrow
\text{Gal}(\kappa(s)^{sep}/\kappa(s)) = \text{Gal}_{\kappa(s)}
$$
est surjectif. Supposons que $(U, \overline{u})$ soit un
voisinage \'etale de $\overline{s}$ et que $\overline{u}$ soit situ\'e au-dessus
du point $u$ de $U$. Puisque $U \to S$ est \'etale, l'extension de corps r\'esiduels
$\kappa(u)/\kappa(s)$ est finie s\'eparable.
Il s'ensuit ce qui suit
\begin{enumerate}
\item Si $\sigma \in \text{Aut}(\kappa(\overline{s})/\kappa(s)^{sep})$,
alors $\sigma$ agit trivialement sur $\mathcal{F}_{\overline{s}}$.
\item Plus pr\'ecis\'ement, l'action de
$\text{Aut}(\kappa(\overline{s})/\kappa(s))$
d\'etermine une action du groupe de Galois absolu
$\text{Gal}_{\kappa(s)}$ sur $\mathcal{F}_{\overline{s}}$, et est d\'etermin\'ee par elle.
\item Si $(U, \overline{u}, t)$ repr\'esente un \'el\'ement $\xi$ de
$\mathcal{F}_{\overline{s}}$, tout \'el\'ement de
$\text{Gal}(\kappa(s)^{sep}/K)$ agit trivialement, o\`u
$\kappa(s) \subset K \subset \kappa(s)^{sep}$ est l'image de
$\overline{u}^\sharp : \kappa(u) \to \kappa(\overline{s})$.
\end{enumerate}
Au total, nous voyons que $\mathcal{F}_{\overline{s}}$ devient un
$\text{Gal}_{\kappa(s)}$-ensemble (voir
Groupes fondamentaux, d\'efinition \ref{pione-definition-G-set-continuous}).
Nous pouvons donc consid\'erer le foncteur fibre comme un foncteur
$$
\Sh(S_\etale) \longrightarrow
\text{Gal}_{\kappa(s)}\textit{-Ensembles},
\quad
\mathcal{F} \longmapsto \mathcal{F}_{\overline{s}}
$$
et, dor\'enavant, c'est g\'en\'eralement ainsi que nous consid\'ererons le foncteur fibre.

\begin{theorem}
\label{etale-cohomology-theorem-equivalence-sheaves-point}
Soit $S = \Spec(K)$, o\`u $K$ est un corps.
Soit $\overline{s}$ un point g\'eom\'etrique de $S$.
Notons $G = \text{Gal}_{\kappa(s)}$ le groupe de Galois absolu.
Le passage aux fibres induit une \'equivalence de cat\'egories
$$
\Sh(S_\etale) \longrightarrow G\textit{-Ensembles},
\quad
\mathcal{F} \longmapsto \mathcal{F}_{\overline{s}}.
$$
\end{theorem}

\begin{proof}
Construisons un quasi-inverse de ce foncteur. Dans
Groupes fondamentaux, lemme \ref{pione-lemma-sheaves-point},
nous avons vu que, pour tout $G$-ensemble $M$, il existe un morphisme \'etale
$X \to \Spec(K)$
tel que $\Mor_K(\Spec(K^{sep}), X)$ soit
isomorphe \`a $M$ comme $G$-ensemble. Consid\'erons le faisceau
$\mathcal{F}$ sur $\Spec(K)_\etale$ d\'efini par
la r\`egle $U \mapsto \Mor_K(U, X)$. C'est un faisceau, car la topologie
\'etale est sous-canonique. On voit alors que
$\mathcal{F}_{\overline{s}} = \Mor_K(\Spec(K^{sep}), X) = M$
comme $G$-ensembles (d\'etails omis). Cela fournit un quasi-inverse du foncteur et
ach\`eve la d\'emonstration.
\end{proof}

\begin{remark}
\label{etale-cohomology-remark-every-sheaf-representable}
Une autre mani\`ere d'\'enoncer la conclusion du
th\'eor\`eme \ref{etale-cohomology-theorem-equivalence-sheaves-point} et de
Groupes fondamentaux, lemme \ref{pione-lemma-sheaves-point},
consiste \`a dire que tout faisceau sur $\Spec(K)_\etale$ est repr\'esentable
par un sch\'ema $X$ \'etale sur $\Spec(K)$.
Cela ne signifie pas que tout faisceau soit repr\'esentable au sens de
Sites, d\'efinition \ref{sites-definition-representable-sheaf}.
La raison est que, dans notre construction de $\Spec(K)_\etale$,
nous avons choisi un ensemble suffisamment grand de sch\'emas \'etales sur $\Spec(K)$,
tandis que les faisceaux sur $\Spec(K)_\etale$ forment une classe propre.
\end{remark}

\begin{lemma}
\label{etale-cohomology-lemma-global-sections-point}
Hypoth\`eses et notations comme dans le
th\'eor\`eme \ref{etale-cohomology-theorem-equivalence-sheaves-point}.
Il existe une bijection fonctorielle
$$
\Gamma(S, \mathcal{F}) = (\mathcal{F}_{\overline{s}})^G
$$
\end{lemma}

\begin{proof}
Nous pouvons le d\'emontrer par des arguments formels et \`a l'aide du
th\'eor\`eme \ref{etale-cohomology-theorem-equivalence-sheaves-point},
comme suit. Pour un faisceau $\mathcal{F}$ correspondant au
$G$-ensemble $M = \mathcal{F}_{\overline{s}}$, on a
\begin{eqnarray*}
\Gamma(S, \mathcal{F}) & = &
\Mor_{\Sh(S_\etale)}(h_{\Spec(K)},  \mathcal{F})
\\
& = & \Mor_{G\textit{-Ensembles}}(\{*\}, M) \\
& = & M^G
\end{eqnarray*}
Ici, la premi\`ere identification est expliqu\'ee dans
Sites, sections \ref{sites-section-presheaves} et
\ref{sites-section-representable-sheaves},
la deuxi\`eme r\'esulte du
th\'eor\`eme \ref{etale-cohomology-theorem-equivalence-sheaves-point}
et la troisi\`eme est claire. Nous donnerons aussi une preuve directe\footnote{Pour
les saint Thomas parmi nous.}.

\medskip\noindent
Supposons que $t \in \Gamma(S, \mathcal{F})$ soit une section globale.
Alors le triplet $(S, \overline{s}, t)$ d\'efinit un \'el\'ement de
$\mathcal{F}_{\overline{s}}$ qui est clairement invariant sous
l'action de $G$. R\'eciproquement, supposons que $(U, \overline{u}, t)$
d\'efinisse un \'el\'ement invariant de $\mathcal{F}_{\overline{s}}$.
Nous pouvons alors r\'etr\'ecir $U$ et supposer que $U = \Spec(L)$ pour une
extension de corps finie s\'eparable $L$ de $K$, voir la
proposition \ref{etale-cohomology-proposition-etale-morphisms}.
Dans ce cas, le morphisme $\mathcal{F}(U) \to \mathcal{F}_{\overline{s}}$
est injectif, car pour tout morphisme de voisinages \'etales
$(U', \overline{u}') \to (U, \overline{u})$, le morphisme de restriction
$\mathcal{F}(U) \to \mathcal{F}(U')$ est injectif puisque $U' \to U$
est un recouvrement de $S_\etale$.
Apr\`es avoir un peu agrandi $L$, nous pouvons supposer que $K \subset L$ est une extension
galoisienne finie. Nous utilisons alors l'\'egalit\'e
$$
\Spec(L) \times_{\Spec(K)} \Spec(L)
=
\coprod\nolimits_{\sigma \in \text{Gal}(L/K)} \Spec(L)
$$
o\`u les morphismes $\Spec(L) \to \Spec(L \otimes_K L)$
proviennent des morphismes d'anneaux $a \otimes b \mapsto a\sigma(b)$. Nous
voyons donc que l'invariance de $(U, \overline{u}, t)$
sous tout $G$ implique que $t \in \mathcal{F}(\Spec(L))$
a la m\^eme image dans
$\mathcal{F}(\Spec(L) \times_{\Spec(K)} \Spec(L))$
par restriction suivant chacune des deux projections (on utilise ici l'injectivit\'e mentionn\'ee
ci-dessus; d\'etails omis). La condition de faisceau pour $\mathcal{F}$
et le recouvrement \'etale $\{\Spec(L) \to \Spec(K)\}$ s'applique
donc, et nous concluons que $t$ provient d'une unique section de $\mathcal{F}$
sur $\Spec(K)$.
\end{proof}

\begin{remark}
\label{etale-cohomology-remark-stalk-pullback}
Soit $S$ un sch\'ema et soit $\overline{s} : \Spec(k) \to S$
un point g\'eom\'etrique de $S$. Par d\'efinition, cela signifie que $k$
est alg\'ebriquement clos. En particulier, le groupe de Galois absolu de $k$
est trivial. Par cons\'equent, d'apr\`es le
th\'eor\`eme \ref{etale-cohomology-theorem-equivalence-sheaves-point},
la cat\'egorie des faisceaux sur $\Spec(k)_\etale$ est \'equivalente
\`a la cat\'egorie des ensembles. L'\'equivalence est donn\'ee par la prise des
sections sur $\Spec(k)$. Cela nous fournit enfin une
autre d\'efinition du foncteur fibre. Plus pr\'ecis\'ement, le foncteur
$$
\Sh(S_\etale) \longrightarrow \textit{Ensembles}, \quad
\mathcal{F} \longmapsto \mathcal{F}_{\overline{s}}
$$
est isomorphe au foncteur
$$
\Sh(S_\etale)
\longrightarrow
\Sh(\Spec(k)_\etale) = \textit{Ensembles},
\quad
\mathcal{F} \longmapsto \overline{s}^*\mathcal{F}
$$
Pour le d\'emontrer rigoureusement, on peut utiliser
la partie (3) du lemme \ref{etale-cohomology-lemma-stalk-pullback}
avec $f = \overline{s}$. De plus, ceci \'etant dit, le cas g\'en\'eral de la
partie (3) du lemme \ref{etale-cohomology-lemma-stalk-pullback}
r\'esulte de la fonctorialit\'e des images inverses.
\end{remark}






\section{Cohomologie des groupes}
\label{etale-cohomology-section-group-cohomology}

\noindent
Dans la suite, lorsque nous \'ecrivons $H^i(G, M)$, nous entendons que $G$
est un groupe topologique, que $M$ est un $G$-module discret muni d'une
action continue de $G$, et que $H^i(G, -)$ est le $i$-i\`eme foncteur d\'eriv\'e
droit sur la cat\'egorie $\text{Mod}_G$ de ces $G$-modules, voir les
d\'efinitions \ref{etale-cohomology-definition-G-module-continuous} et
\ref{etale-cohomology-definition-galois-cohomology}. Cela comprend
le cas d'un groupe abstrait $G$, ce qui signifie simplement que l'on consid\`ere $G$
comme un groupe topologique muni de la topologie discr\`ete.

\medskip\noindent
Lorsque le module est muni d'une topologie non discr\`ete, nous emploierons la notation
$H^i_{cont}(G, M)$ pour d\'esigner les groupes de cohomologie continue
introduits dans \cite{Tate}, voir la
section \ref{etale-cohomology-section-continuous-group-cohomology}.

\begin{definition}
\label{etale-cohomology-definition-G-module-continuous}
Soit $G$ un groupe topologique.
\begin{enumerate}
\item Un {\it $G$-module}, parfois appel\'e {\it $G$-module discret},
est un groupe ab\'elien $M$ muni d'une action \`a gauche $a : G \times M \to M$
par homomorphismes de groupes, telle que $a$ soit continue lorsque $M$ est muni de la
topologie discr\`ete.
\item Un {\it morphisme de $G$-modules} $f : M \to N$ est un
homomorphisme $G$-\'equivariant de $M$ dans $N$.
\item La cat\'egorie des $G$-modules est not\'ee $\text{Mod}_G$.
\end{enumerate}
Soit $R$ un anneau.
\begin{enumerate}
\item Un {\it $R\text{-}G$-module} est un $R$-module $M$ muni
d'une action \`a gauche $a : G \times M \to M$ par applications $R$-lin\'eaires, telle que $a$
soit continue lorsque $M$ est muni de la topologie discr\`ete.
\item Un {\it morphisme de $R\text{-}G$-modules} $f : M \to N$ est une
application $R$-lin\'eaire $G$-\'equivariante de $M$ dans $N$.
\item La cat\'egorie des $R\text{-}G$-modules est not\'ee $\text{Mod}_{R, G}$.
\end{enumerate}
\end{definition}

\noindent
La continuit\'e de $a : G \times M \to M$ \'equivaut
\`a demander que le stabilisateur de tout $x \in M$ soit ouvert dans $G$.
Si $G$ est un groupe abstrait, cela correspond \`a la notion d'un
groupe ab\'elien muni d'une action de $G$, pourvu que l'on munisse $G$ de la
topologie discr\`ete. Observons que $\text{Mod}_{\mathbf{Z}, G} = \text{Mod}_G$.

\medskip\noindent
La cat\'egorie $\text{Mod}_G$ a assez d'injectifs, voir
Injectifs, lemme \ref{injectives-lemma-G-modules}.
Consid\'erons le foncteur exact \`a gauche
$$
\text{Mod}_G \longrightarrow \textit{Ab},
\quad
M \longmapsto M^G =
\{x \in M \mid g \cdot x = x\ \forall g \in G\}
$$
Nous notons parfois $M^G = H^0(G, M)$ et parfois
$M^G = \Gamma_G(M)$. Ce foncteur admet un foncteur d\'eriv\'e droit total
$R\Gamma_G(M)$ et, pour tout $i \geq 0$, un $i$-i\`eme foncteur d\'eriv\'e droit
$R^i\Gamma_G(M) = H^i(G, M)$.

\medskip\noindent
La m\^eme construction s'applique \`a
$H^0(G, -) : \text{Mod}_{R, G} \to \text{Mod}_R$. Nous verrons au
lemme \ref{etale-cohomology-lemma-modules-abelian} qu'elle co\"\i ncide avec la cohomologie
du $G$-module sous-jacent.

\begin{definition}
\label{etale-cohomology-definition-galois-cohomology}
Soit $G$ un groupe topologique. Soit $M$ un $G$-module discret
muni d'une action continue de $G$. Autrement dit, $M$ est un objet
de la cat\'egorie $\text{Mod}_G$ introduite dans la
d\'efinition \ref{etale-cohomology-definition-G-module-continuous}.
\begin{enumerate}
\item Les foncteurs d\'eriv\'es droits $H^i(G, M)$ de $H^0(G, M)$ sur la
cat\'egorie $\text{Mod}_G$ sont appel\'es les
{\it groupes de cohomologie continue de $G$ \`a coefficients dans $M$}.
\item Si $G$ est un groupe abstrait muni de la topologie discr\`ete,
les $H^i(G, M)$ sont appel\'es les {\it groupes de cohomologie de $G$ \`a coefficients dans $M$}.
\item Si $G$ est un groupe de Galois, les groupes $H^i(G, M)$ sont appel\'es
les {\it groupes de cohomologie galoisienne \`a coefficients dans $M$}.
\item Si $G$ est le groupe de Galois absolu d'un corps $K$, les groupes
$H^i(G, M)$ sont parfois appel\'es les {\it groupes de cohomologie galoisienne de $K$
\`a coefficients dans $M$}. Dans ce cas, nous \'ecrivons parfois
$H^i(K, M)$ au lieu de $H^i(G, M)$.
\end{enumerate}
\end{definition}

\begin{lemma}
\label{etale-cohomology-lemma-modules-abelian}
Soit $G$ un groupe topologique. Soit $R$ un anneau.
Pour tout $i \geq 0$, le diagramme
$$
\xymatrix{
\text{Mod}_{R, G} \ar[rr]_{H^i(G, -)} \ar[d] & &
\text{Mod}_R \ar[d] \\
\text{Mod}_G \ar[rr]^{H^i(G, -)} & &
\textit{Ab}
}
$$
dont les fl\`eches verticales sont les foncteurs d'oubli est commutatif.
\end{lemma}

\begin{proof}
Notons $F : \text{Mod}_{R, G} \to \text{Mod}_G$ le foncteur d'oubli.
Alors $F$ admet un adjoint \`a gauche $H : \text{Mod}_G \to \text{Mod}_{R, G}$
donn\'e par $H(M) = M \otimes_\mathbf{Z} R$. Observons que tout objet de
$\text{Mod}_G$ est un quotient d'une somme directe de modules de la forme
$\mathbf{Z}[G/U]$, o\`u $U \subset G$ est un sous-groupe ouvert.
Ici, $\mathbf{Z}[G/U]$ d\'esigne le $G$-module form\'e des
combinaisons $\mathbf{Z}$-lin\'eaires finies
des classes \`a droite modulo $U$ dans $G$, muni de l'action \`a gauche de $G$.
Ainsi, tout complexe born\'e sup\'erieurement dans $\text{Mod}_G$ est quasi-isomorphe
\`a un complexe born\'e sup\'erieurement dans $\text{Mod}_G$ dont les
termes sous-jacents sont des $\mathbf{Z}$-modules plats
(Cat\'egories d\'eriv\'ees, lemme \ref{derived-lemma-subcategory-left-resolution}).
Il est donc clair que $LH$ existe sur $D^-(\text{Mod}_G)$ et se calcule en
appliquant $H$ \`a tout complexe dont les termes sont des $\mathbf{Z}$-modules plats;
cela r\'esulte de
Cat\'egories d\'eriv\'ees, lemme \ref{derived-lemma-subcategory-left-acyclics} et
proposition \ref{derived-proposition-enough-acyclics}.
Nous concluons de Cat\'egories d\'eriv\'ees, lemme
\ref{derived-lemma-pre-derived-adjoint-functors}
que
$$
\text{Ext}^i(\mathbf{Z}, F(M)) = \text{Ext}^i(R, M)
$$
pour $M$ dans $\text{Mod}_{R, G}$.
Observons que $H^0(G, -) = \Hom(\mathbf{Z}, -)$ sur
$\text{Mod}_G$, o\`u $\mathbf{Z}$ d\'esigne le $G$-module
\`a action triviale. Par cons\'equent,
$H^i(G, -) = \text{Ext}^i(\mathbf{Z}, -)$ sur $\text{Mod}_G$.
De m\^eme, on a $H^i(G, -) = \text{Ext}^i(R, -)$ sur
$\text{Mod}_{R, G}$. En combinant ces faits, on obtient le lemme.
\end{proof}

\begin{lemma}
\label{etale-cohomology-lemma-ext-modules-hom}
Soit $G$ un groupe topologique. Soit $R$ un anneau.
Soient $M$, $N$ des $R\text{-}G$-modules. Si $M$ est projectif de type fini
comme $R$-module, alors
$\text{Ext}^i(M, N) = H^i(G, M^\vee \otimes_R N)$ (pour les notations,
voir la preuve).
\end{lemma}

\begin{proof}
Munissons le module $M^\vee = \Hom_R(M, R)$ de l'action contragr\'ediente
de $G$. Ainsi, $(g \cdot \lambda)(m) = \lambda(g^{-1} \cdot m)$
pour $g \in G$, $\lambda \in M^\vee$, $m \in M$. L'action de $G$ sur
$M^\vee \otimes_R N$ est l'action diagonale, c'est-\`a-dire qu'elle est donn\'ee par
$g \cdot (\lambda \otimes n) = g \cdot \lambda \otimes g \cdot n$.
Notons que, pour un troisi\`eme $R\text{-}G$-module $E$, on a
$\Hom(E, M^\vee \otimes_R N) = \Hom(M \otimes_R E, N)$.
En effet, cela est vrai au niveau des $R$-modules d'apr\`es
Alg\`ebre, lemmes \ref{algebra-lemma-hom-from-tensor-product} et
\ref{algebra-lemma-evaluation-map-iso-finite-projective}
et les d\'efinitions des actions de $G$ sont choisies de sorte que cela
reste vrai pour les $R\text{-}G$-modules. Il en r\'esulte que
$M^\vee \otimes_R N$ est un $R\text{-}G$-module injectif
si $N$ est un $R\text{-}G$-module injectif. Ainsi, si
$N \to N^\bullet$ est une r\'esolution injective, alors
$M^\vee \otimes_R N \to M^\vee \otimes_R N^\bullet$
est une r\'esolution injective. On a alors
$$
\Hom(M, N^\bullet) = \Hom(R, M^\vee \otimes_R N^\bullet) =
(M^\vee \otimes_R N^\bullet)^G
$$
Puisque le membre de gauche calcule $\text{Ext}^i(M, N)$ et celui de droite
calcule $H^i(G, M^\vee \otimes_R N)$, la preuve est termin\'ee.
\end{proof}

\begin{lemma}
\label{etale-cohomology-lemma-finite-dim-group-cohomology}
Soit $G$ un groupe topologique. Soit $k$ un corps.
Soit $V$ un $k\text{-}G$-module.
Si $G$ est topologiquement engendr\'e par un nombre fini d'\'el\'ements et si
$\dim_k(V) < \infty$, alors $\dim_k H^1(G, V) < \infty$.
\end{lemma}

\begin{proof}
Soient $g_1, \ldots, g_r \in G$ des \'el\'ements qui engendrent topologiquement $G$,
c'est-\`a-dire que le sous-groupe engendr\'e par $g_1, \ldots, g_r$ est dense.
D'apr\`es le lemme \ref{etale-cohomology-lemma-ext-modules-hom},
on voit que $H^1(G, V)$ est l'espace vectoriel sur $k$ des extensions
$$
0 \to V \to E \to k \to 0
$$
de $k\text{-}G$-modules. Choisissons $e \in E$ d'image $1 \in k$.
\'Ecrivons
$$
g_i \cdot e = v_i + e
$$
pour un $v_i \in V$. Cela est possible puisque $g_i \cdot 1 = 1$.
Nous affirmons que la liste des \'el\'ements $v_1, \ldots, v_r \in V$
d\'etermine la classe d'isomorphisme de l'extension $E$.
Une fois ceci d\'emontr\'e, le lemme en r\'esulte, car notre espace vectoriel
Ext est alors isomorphe \`a un sous-quotient de l'espace vectoriel sur $k$
$V^{\oplus r}$; quelques d\'etails sont omis.
Puisque $E$ est un objet de la cat\'egorie d\'efinie dans la
d\'efinition \ref{etale-cohomology-definition-G-module-continuous},
il existe un sous-groupe ouvert $U$ tel que
$u \cdot e = e$ pour tout $u \in U$.
Prenons maintenant un $g \in G$. Alors $gU$ contient un mot $w$ en
les \'el\'ements $g_1, \ldots, g_r$.
\'Ecrivons $gu = w$. Puisque l'\'el\'ement $w \cdot e$ est d\'etermin\'e par
$v_1, \ldots, v_r$, on voit que $g \cdot e = (gu) \cdot e = w \cdot e$
l'est lui aussi.
\end{proof}

\begin{lemma}
\label{etale-cohomology-lemma-profinite-group-cohomology-torsion}
Soit $G$ un groupe topologique profini.
Alors
\begin{enumerate}
\item $H^i(G, M)$ est de torsion pour $i > 0$ et tout $G$-module $M$, et
\item $H^i(G, M) = 0$ si $M$ est un $\mathbf{Q}$-espace vectoriel.
\end{enumerate}
\end{lemma}

\begin{proof}
Preuve de (1). Par d\'ecalage de dimension, il suffit
de montrer que $H^1(G, M)$ est de torsion pour tout $G$-module $M$.
Choisissons une suite exacte $0 \to M \to I \to N \to 0$, o\`u $I$
est un objet injectif de la cat\'egorie des $G$-modules.
Tout \'el\'ement de $H^1(G, M)$ est alors l'image d'un \'el\'ement
$y \in N^G$. Choisissons $x \in I$ d'image $y$.
Le stabilisateur $U \subset G$ de $x$ est ouvert, donc
d'indice fini $r$. Soit $g_1, \ldots, g_r \in G$ un syst\`eme
de repr\'esentants de $G/U$. Alors $\sum g_i(x)$ est un \'el\'ement
invariant de $I$ dont l'image est $ry$. Ainsi, $r$ annule l'\'el\'ement
de $H^1(G, M)$ dont nous \'etions partis. La partie (2) en r\'esulte, puisque
$H^i(G, M)$ est alors \`a la fois un $\mathbf{Q}$-espace vectoriel et de torsion.
\end{proof}





\section{Cohomologie continue de Tate}
\label{etale-cohomology-section-continuous-group-cohomology}

\noindent
La cohomologie continue de Tate (\cite{Tate}) est d\'efinie par le complexe des
cocha\^ines inhomog\`enes continues. On peut la d\'efinir lorsque $M$ est un
groupe ab\'elien topologique quelconque muni d'une action continue de $G$.
Plus pr\'ecis\'ement, consid\'erons le complexe
$$
C^\bullet_{cont}(G, M) :
M \to \text{Maps}_{cont}(G, M) \to
\text{Maps}_{cont}(G \times G, M) \to \ldots
$$
o\`u l'application de bord est d\'efinie, pour $n \geq 1$, par la formule
\begin{align*}
\text{d}(f)(g_1, \ldots, g_{n + 1})
& = g_1(f(g_2, \ldots, g_{n + 1})) \\
&
+ \sum\nolimits_{j = 1, \ldots, n}
(-1)^jf(g_1, \ldots, g_jg_{j + 1}, \ldots, g_{n + 1}) \\
&
+ (-1)^{n + 1}f(g_1, \ldots, g_n)
\end{align*}
et, pour $n = 0$, envoie $m \in M$ sur l'application $g \mapsto g(m) - m$. Posons
$$
H^i_{cont}(G, M) = H^i(C^\bullet_{cont}(G, M))
$$
Comme les termes du complexe font intervenir les applications continues de $G$ et
de ses puissances cart\'esiennes dans le module topologique $M$, il n'est pas clair
qu'une suite exacte courte de modules topologiques soit ainsi transform\'ee en
une suite exacte longue de cohomologie. Une autre difficult\'e tient \`a ce que la cat\'egorie
des groupes ab\'eliens topologiques n'est pas une cat\'egorie ab\'elienne !

\medskip\noindent
Cependant, une suite exacte courte de $G$-modules discrets fournit bien
une suite exacte courte de complexes de cocha\^ines continues
et donc une suite exacte longue de groupes de cohomologie
continue $H^i_{cont}(G, -)$.
Par cons\'equent, dans la cat\'egorie $\text{Mod}_G$ de la
d\'efinition \ref{etale-cohomology-definition-G-module-continuous}, les foncteurs
$H^i_{cont}(G, M)$ forment un $\delta$-foncteur cohomologique au sens de
Homologie, section \ref{homology-section-cohomological-delta-functor}.
Puisque la cohomologie $H^i(G, M)$
de la d\'efinition \ref{etale-cohomology-definition-galois-cohomology}
est un $\delta$-foncteur universel
(Cat\'egories d\'eriv\'ees, lemme \ref{derived-lemma-right-derived-delta-functor}),
on obtient des morphismes canoniques
$$
H^i(G, M) \longrightarrow H^i_{cont}(G, M)
$$
pour $M \in \text{Mod}_G$. On sait que ces morphismes sont des
isomorphismes lorsque $G$ est un groupe abstrait (c'est-\`a-dire muni de
la topologie discr\`ete) ou lorsque $G$ est un groupe profini
(ins\'erer ici une r\'ef\'erence ult\'erieure).
Si vous connaissez un exemple o\`u ce morphisme n'est pas un isomorphisme
pour un groupe topologique $G$ et $M \in \Ob(\text{Mod}_G)$,
veuillez l'envoyer par courriel \`a
\href{mailto:stacks.project@gmail.com}{stacks.project@gmail.com}.




\section{Cohomologie d'un point}
\label{etale-cohomology-section-cohomology-point}

\noindent
Comme cons\'equence de la discussion des sections pr\'ec\'edentes,
nous obtenons l'\'equivalence entre la cohomologie \'etale du spectre d'un
corps et la cohomologie galoisienne.

\begin{lemma}
\label{etale-cohomology-lemma-equivalence-abelian-sheaves-point}
Soit $S = \Spec(K)$, o\`u $K$ est un corps.
Soit $\overline{s}$ un point g\'eom\'etrique de $S$.
Notons $G = \text{Gal}_{\kappa(s)}$ le groupe de Galois absolu.
Le foncteur fibre induit une \'equivalence de cat\'egories
$$
\textit{Ab}(S_\etale) \longrightarrow \text{Mod}_G,
\quad
\mathcal{F} \longmapsto \mathcal{F}_{\overline{s}}.
$$
\end{lemma}

\begin{proof}
Dans le
th\'eor\`eme \ref{etale-cohomology-theorem-equivalence-sheaves-point},
nous avons vu l'\'equivalence entre les faisceaux d'ensembles et les $G$-ensembles.
Le pr\'esent lemme s'en d\'eduit formellement, car un faisceau ab\'elien n'est autre
qu'un faisceau d'ensembles muni d'une loi de groupe commutative, et un $G$-module
n'est autre qu'un $G$-ensemble muni d'une loi de groupe commutative.
\end{proof}

\begin{lemma}
\label{etale-cohomology-lemma-compare-cohomology-point}
Les notations et hypoth\`eses sont celles du
lemme \ref{etale-cohomology-lemma-equivalence-abelian-sheaves-point}.
Soit $\mathcal{F}$ un faisceau ab\'elien sur $\Spec(K)_\etale$
qui correspond au $G$-module $M$.
Alors
\begin{enumerate}
\item dans $D(\textit{Ab})$, on a un isomorphisme canonique
$R\Gamma(S, \mathcal{F}) = R\Gamma_G(M)$,
\item $H_\etale^0(S, \mathcal{F}) = M^G$, et
\item $H_\etale^q(S, \mathcal{F}) = H^q(G, M)$.
\end{enumerate}
\end{lemma}

\begin{proof}
Combiner le
lemme \ref{etale-cohomology-lemma-equivalence-abelian-sheaves-point}
avec le
lemme \ref{etale-cohomology-lemma-global-sections-point}.
\end{proof}

\begin{example}
\label{etale-cohomology-example-sheaves-point}
Faisceaux sur $\Spec(K)_\etale$.
Soit $G = \text{Gal}(K^{sep}/K)$ le groupe de Galois absolu de $K$.
\begin{enumerate}
\item Le faisceau constant $\underline{\mathbf{Z}/n\mathbf{Z}}$ correspond au
module $\mathbf{Z}/n\mathbf{Z}$ muni de l'action triviale de $G$,
\item le faisceau $\mathbf{G}_m|_{\Spec(K)_\etale}$ correspond \`a
$(K^{sep})^*$ muni de son action de $G$,
\item le faisceau $\mathbf{G}_a|_{\Spec(K^{sep})}$ correspond \`a
$(K^{sep}, +)$ muni de son action de $G$, et
\item le faisceau $\mu_n|_{\Spec(K^{sep})}$ correspond \`a
$\mu_n(K^{sep})$ muni de son action de $G$.
\end{enumerate}
D'apr\`es la
remarque \ref{etale-cohomology-remark-special-case-fpqc-cohomology-quasi-coherent}
et le
th\'eor\`eme \ref{etale-cohomology-theorem-picard-group},
on a les identifications suivantes pour les groupes de cohomologie :
\begin{align*}
H_\etale^0(S_\etale, \mathbf{G}_m) & =
\Gamma(S, \mathcal{O}_S^*) \\
H_\etale^1(S_\etale, \mathbf{G}_m) & =
H_{Zar}^1(S, \mathcal{O}_S^*) = \Pic(S) \\
H_\etale^i(S_\etale, \mathbf{G}_a) & =
H_{Zar}^i(S, \mathcal{O}_S)
\end{align*}
De plus, pour tout faisceau quasi-coh\'erent $\mathcal{F}$ sur $S_\etale$, on a
$$
H^i(S_\etale, \mathcal{F}) = H_{Zar}^i(S, \mathcal{F}),
$$
voir le
th\'eor\`eme \ref{etale-cohomology-theorem-zariski-fpqc-quasi-coherent}.
En particulier, on obtient la suite d'\'egalit\'es suivante
$$
0 =
\Pic(\Spec(K)) =
H_\etale^1(\Spec(K)_\etale, \mathbf{G}_m) =
H^1(G, (K^{sep})^*)
$$
qui n'est autre que le th\'eor\`eme 90 de Hilbert. De m\^eme, pour $i \geq 1$,
$$
0 = H^i(\Spec(K), \mathcal{O})
= H_\etale^i(\Spec(K)_\etale, \mathbf{G}_a)
= H^i(G, K^{sep})
$$
o\`u $K^{sep}$ est ici consid\'er\'e comme module galoisien, l'addition \'etant
sa loi de groupe. De la sorte, nous pouvons consid\'erer le travail accompli jusqu'ici comme
une mani\`ere compliqu\'ee de calculer des groupes de cohomologie galoisienne.
\end{example}

\noindent
Le r\'esultat suivant est une curiosit\'e et devrait \^etre omis lors d'une
premi\`ere lecture.

\begin{lemma}
\label{etale-cohomology-lemma-all-modules-quasi-coherent}
Soit $R$ un anneau local de dimension $0$. Soit $S = \Spec(R)$.
Alors tout $\mathcal{O}_S$-module sur $S_\etale$ est quasi-coh\'erent.
\end{lemma}

\begin{proof}
Soit $\mathcal{F}$ un $\mathcal{O}_S$-module sur $S_\etale$.
Nous devons montrer que $\mathcal{F}$ est d\'etermin\'e par le
$R$-module $M = \Gamma(S, \mathcal{F})$.
Plus pr\'ecis\'ement, si $\pi : X \to S$ est \'etale, nous devons
montrer que $\Gamma(X, \mathcal{F}) = \Gamma(X, \pi^*\widetilde{M})$.

\medskip\noindent
Soit $\mathfrak m \subset R$ l'id\'eal maximal et soit
$\kappa$ le corps r\'esiduel. D'apr\`es
Alg\`ebre, lemme \ref{algebra-lemma-local-dimension-zero-henselian},
l'anneau local $R$ est hens\'elien. Si $X \to S$ est \'etale,
alors l'espace topologique sous-jacent \`a $X$ est discret
d'apr\`es Morphismes, lemme \ref{morphisms-lemma-etale-over-field},
et $X$ est donc une r\'eunion disjointe de sch\'emas affines
ayant chacun un seul point. De plus, si $X = \Spec(A)$ est affine et
n'a qu'un point, alors $A$ est une $R$-alg\`ebre finie \'etale d'apr\`es
Alg\`ebre, lemme \ref{algebra-lemma-mop-up}.
Nous devons montrer que $\Gamma(X, \mathcal{F}) = M \otimes_R A$
dans ce cas.

\medskip\noindent
Le foncteur $A \mapsto A/\mathfrak m A$ d\'efinit une \'equivalence entre
la cat\'egorie des $R$-alg\`ebres finies \'etales
et celle des $\kappa$-alg\`ebres finies s\'eparables, d'apr\`es
Alg\`ebre, lemme \ref{algebra-lemma-henselian-cat-finite-etale}.
Consid\'erons d'abord le cas o\`u $A/\mathfrak m A$
est une extension galoisienne de $\kappa$, de groupe de Galois $G$.
Pour chaque $\sigma \in G$, notons $\sigma : A \to A$
l'automorphisme correspondant de $A$ au-dessus de $R$.
Posons $N = \Gamma(X, \mathcal{F})$.
Alors $\Spec(\sigma) : X \to X$ est un automorphisme au-dessus de $S$
et l'image inverse par celui-ci d\'efinit donc une application $\sigma : N \to N$
qui est $\sigma$-lin\'eaire : $\sigma(an) = \sigma(a) \sigma(n)$
pour $a \in A$ et $n \in N$.
Nous allons appliquer la descente galoisienne au module quasi-coh\'erent
$\widetilde{N}$ sur $X$, muni des isomorphismes
provenant de l'action de $\sigma$ sur $N$. Voir Descente, lemme
\ref{descent-lemma-galois-descent-more-general}.
Ce lemme fournit un isomorphisme $N = N^G \otimes_R A$.
D'autre part, il est clair que $N^G = M$ par la propri\'et\'e de faisceau
de $\mathcal{F}$. L'isomorphisme voulu est donc \'etabli.

\medskip\noindent
Le cas g\'en\'eral (o\`u $A$ est une $R$-alg\`ebre locale finie \'etale)
se d\'eduit du cas galoisien comme suit. Choisissons $A \to B$
finie \'etale, avec $B$ local et son corps r\'esiduel
galoisien sur $\kappa$. Posons $G = \text{Aut}(B/R) = \text{Gal}(\kappa_B/\kappa)$.
Soit $H \subset G$ le groupe de Galois correspondant \`a
l'extension galoisienne $\kappa_B/\kappa_A$. Alors, comme ci-dessus, on
montre que $\Gamma(X, \mathcal{F}) = \Gamma(\Spec(B), \mathcal{F})^H$.
En appliquant deux fois le r\'esultat relatif aux extensions galoisiennes, on obtient
$$
\Gamma(X, \mathcal{F}) = (M \otimes_R B)^H = M \otimes_R A
$$
comme voulu.
\end{proof}








\section{Cohomologie des courbes}
\label{etale-cohomology-section-cohomology-curves}

\noindent
La t\^ache suivante consiste \`a calculer la cohomologie \'etale d'une courbe lisse
sur un corps alg\'ebriquement clos, \`a coefficients de torsion, et en
particulier \`a montrer qu'elle s'annule en degr\'e au moins $3$. Pour cela, nous
calculerons la cohomologie au point g\'en\'erique, ce qui
revient \`a calculer de la cohomologie galoisienne.





\section{Groupes de Brauer}
\label{etale-cohomology-section-brauer-groups}

\noindent
Les groupes de Brauer des corps sont d\'efinis au moyen d'alg\`ebres centrales simples de dimension finie.
Dans cette section, nous rappelons les faits pertinents sur les groupes de Brauer, dont la plupart
sont expos\'es dans le chapitre
Groupes de Brauer, section \ref{brauer-section-introduction}.
Pour d'autres r\'ef\'erences, voir \cite{SerreCorpsLocaux},
\cite{SerreGaloisCohomology} ou \cite{Weil}.

\begin{theorem}
\label{etale-cohomology-theorem-central-simple-algebra}
Soit $K$ un corps. Pour une $K$-alg\`ebre $A$ unitaire et associative (non n\'ecessairement commutative),
les conditions suivantes sont \'equivalentes
\begin{enumerate}
\item $A$ est une $K$-alg\`ebre centrale simple de dimension finie,
\item $A$ est un $K$-espace vectoriel de dimension finie, $K$ est le centre de $A$,
et $A$ n'a pas d'id\'eal bilat\`ere non trivial,
\item il existe $d \geq 1$ tel que
$A \otimes_K \bar K \cong \text{Mat}(d \times d, \bar K)$,
\item il existe $d \geq 1$ tel que
$A \otimes_K K^{sep} \cong \text{Mat}(d \times d, K^{sep})$,
\item il existe $d \geq 1$ et une extension galoisienne finie $K'/K$
tels que
$A \otimes_K K' \cong \text{Mat}(d \times d, K')$,
\item il existe $n \geq 1$ et un corps gauche central $D$ de dimension finie
sur $K$ tel que $A \cong \text{Mat}(n \times n, D)$.
\end{enumerate}
L'entier $d$ est appel\'e le {\it degr\'e} de $A$.
\end{theorem}

\begin{proof}
Il s'agit d'une copie de
Groupes de Brauer, lemme \ref{brauer-lemma-finite-central-simple-algebra}.
\end{proof}

\begin{lemma}
\label{etale-cohomology-lemma-brauer-inverse}
Soit $A$ une alg\`ebre centrale simple de dimension finie sur $K$. Alors
$$
\begin{matrix}
A \otimes_K A^{opp} & \longrightarrow & \text{End}_K(A) \\
\ a \otimes a' & \longmapsto & (x \mapsto a x a')
\end{matrix}
$$
est un isomorphisme de $K$-alg\`ebres.
\end{lemma}

\begin{proof}
Voir
Groupes de Brauer, lemme \ref{brauer-lemma-inverse}.
\end{proof}

\begin{definition}
\label{etale-cohomology-definition-brauer-equivalent}
Deux alg\`ebres centrales simples de dimension finie $A_1$ et $A_2$ sur $K$ sont dites
{\it semblables}, ou {\it \'equivalentes}, s'il existe $m, n \geq 1$
tels que $\text{Mat}(n \times n, A_1)
\cong \text{Mat}(m \times m, A_2)$. On \'ecrit $A_1 \sim A_2$.
\end{definition}

\noindent
D'apr\`es Groupes de Brauer, lemme \ref{brauer-lemma-similar}, c'est une
relation d'\'equivalence.

\begin{definition}
\label{etale-cohomology-definition-brauer-group}
Soit $K$ un corps. Le {\it groupe de Brauer} de $K$ est l'ensemble $\text{Br} (K)$
des classes de similitude d'alg\`ebres centrales simples de dimension finie sur $K$, muni de
la loi de groupe induite par le produit tensoriel (sur $K$). La classe de $A$ dans
$\text{Br}(K)$ est not\'ee $[A]$. L'\'el\'ement neutre est
$[K] = [\text{Mat}(d \times d, K)]$ pour tout $d \geq 1$.
\end{definition}

\noindent
Le lemme pr\'ec\'edent entra\^ine l'existence des inverses et l'\'egalit\'e $-[A] = [A^{opp}]$.
Le groupe de Brauer d'un corps est toujours de torsion.
En fait, nous verrons que $[A]$ est d'ordre divisant $\deg(A)$
pour toute alg\`ebre centrale simple de dimension finie $A$ (voir le
lemme \ref{etale-cohomology-lemma-annihilated-by-degree}).
En g\'en\'eral, le groupe de Brauer n'est pas de type fini; par exemple,
le groupe de Brauer d'un corps local non archim\'edien est $\mathbf{Q}/\mathbf{Z}$.
Le groupe de Brauer de $\mathbf{C}(x, y)$ n'est pas d\'enombrable.

\begin{lemma}
\label{etale-cohomology-lemma-central-simple-algebra-pgln}
\begin{slogan}
Les alg\`ebres centrales simples sont class\'ees par la cohomologie galoisienne de PGL.
\end{slogan}
Soient $K$ un corps et $K^{sep}$ une cl\^oture s\'eparable.
Alors l'ensemble des classes d'isomorphisme d'alg\`ebres centrales simples de degr\'e
$d$ sur $K$ est en bijection avec la cohomologie non ab\'elienne
$H^1(\text{Gal}(K^{sep}/K), \text{PGL}_d(K^{sep}))$.
\end{lemma}

\begin{proof}[Esquisse de preuve.]
Le th\'eor\`eme de Skolem-Noether (voir
Groupes de Brauer, th\'eor\`eme \ref{brauer-theorem-skolem-noether})
entra\^ine que, pour tout corps $L$, le groupe
$\text{Aut}_{L\text{-Algebras}}(\text{Mat}_d(L))$
est \`egal \`a $\text{PGL}_d(L)$. D'apr\`es le
th\'eor\`eme \ref{etale-cohomology-theorem-central-simple-algebra}, on voit que les
alg\`ebres centrales simples de degr\'e $d$ correspondent
aux formes de la $K$-alg\`ebre $\text{Mat}_d(K)$.
En combinant ces faits, on voit que les classes d'isomorphisme d'alg\`ebres
centrales simples de degr\'e $d$ correspondent aux \'el\'ements de
$H^1(\text{Gal}(K^{sep}/K), \text{PGL}_d(K^{sep}))$.
Pour plus de d\'etails sur les formes tordues, voir par exemple
\cite{SilvermanEllipticCurves}.
\end{proof}

\noindent
Si $A$ est une alg\`ebre centrale simple de dimension finie et de degr\'e $d$ sur un corps $K$,
nous notons $\xi_A$ la classe de cohomologie correspondante dans
$H^1(\text{Gal}(K^{sep}/K), \text{PGL}_d(K^{sep}))$.
Consid\'erons la suite exacte courte
$$
1 \to (K^{sep})^* \to \text{GL}_d(K^{sep}) \to \text{PGL}_d(K^{sep}) \to 1,
$$
qui donne naissance \`a une suite exacte longue de cohomologie (jusqu'au degr\'e $2$), avec
pour homomorphisme bord
$$
\delta_d :
H ^1(\text{Gal}(K^{sep}/K), \text{PGL}_d(K^{sep}))
\longrightarrow
H^2(\text{Gal}(K^{sep}/K), (K^{sep})^*).
$$
Plus pr\'ecis\'ement, il est d\'efini comme suit : si $\xi$ est une classe de cohomologie
repr\'esent\'ee par le $1$-cocycle $(g_\sigma)$, alors $\delta_d(\xi)$ est la
classe du $2$-cocycle
\begin{equation}
\label{etale-cohomology-equation-two-cocycle}
(\sigma, \tau)
\longmapsto
\tilde g_\sigma^{-1} \tilde g_{\sigma \tau} \sigma(\tilde g_\tau^{-1})
\in (K^{sep})^*
\end{equation}
o\`u $\tilde g_\sigma \in \text{GL}_d(K^{sep})$ est un rel\`evement de $g_\sigma$.
Cela permet d'expliciter l'application
$$
\delta : \text{Br}(K) \longrightarrow H^2(\text{Gal}(K^{sep}/K), (K^{sep})^*),
\quad
[A] \longmapsto \delta_{\deg A} (\xi_A)
$$
comme suit. Supposons que $A$ soit de degr\'e $d$ sur $K$. Choisissons un isomorphisme
$\varphi : \text{Mat}_d(K^{sep}) \to A \otimes_K K^{sep}$. Pour
$\sigma \in \text{Gal}(K^{sep}/K)$, choisissons un \'el\'ement
$\tilde g_\sigma \in \text{GL}_d(K^{sep})$ tel que
$\varphi^{-1} \circ \sigma(\varphi)$ soit \'egal \`a l'application
$x \mapsto \tilde g_\sigma x \tilde g_\sigma^{-1}$. La classe dans $H^2$
est d\'efinie par le $2$-cocycle (\ref{etale-cohomology-equation-two-cocycle}).

\begin{theorem}
\label{etale-cohomology-theorem-brauer-delta}
Soit $K$ un corps de cl\^oture s\'eparable $K^{sep}$. L'application
$\delta : \text{Br}(K) \to H^2(\text{Gal}(K^{sep}/K), (K^{sep})^*)$
d\'efinie ci-dessus est un isomorphisme de groupes.
\end{theorem}

\begin{proof}[Esquisse de preuve]
Pour montrer que $\delta$ d\'efinit un homomorphisme de groupes, c'est-\`a-dire que
$\delta(A \otimes_K B) = \delta(A) + \delta(B)$, on calcule
directement sur les cocycles.

\medskip\noindent
Injectivit\'e de $\delta$. Dans le cas ab\'elien ($d = 1$), on a
l'identification
$$
H^1(\text{Gal}(K^{sep}/K), \text{GL}_d(K^{sep})) =
H_\etale^1(\Spec(K), \text{GL}_d(\mathcal{O}))
$$
dont le second membre est trivial par descente fpqc. Si cela \'etait vrai dans le
cas non ab\'elien, l'injectivit\'e de $\delta$ en r\'esulterait aussit\^ot. (Voir
\cite{SGA4.5}.) En fait, pour la d\'emontrer, on peut r\'einterpr\'eter $\delta([A])$ comme
l'obstruction \`a l'existence d'un $K$-espace vectoriel $V$ muni d'une structure de $A$-module
\`a gauche et tel que $\dim_K V = \deg A$. Lorsque $V$ existe, on
a $A \cong \text{End}_K(V)$.

\medskip\noindent
Pour la surjectivit\'e, choisissons une
classe de cohomologie $\xi \in H^2(\text{Gal}(K^{sep}/K), (K^{sep})^*)$ ;
il existe alors une extension galoisienne finie $K^{sep}/K'/K$
telle que $\xi$ soit l'image d'une classe
$\xi' \in H^2(\text{Gal}(K'|K), (K')^*)$. On construit alors
explicitement une alg\`ebre centrale simple sur $K$ \`a partir des donn\'ees $K', \xi'$.
\end{proof}










\section{Le groupe de Brauer d'un sch\'ema}
\label{etale-cohomology-section-brauer-scheme}

\noindent
Soit $S$ un sch\'ema. Une $\mathcal{O}_S$-alg\`ebre
$\mathcal{A}$ est dite {\it d'Azumaya} si elle est, localement pour la topologie \'etale, une alg\`ebre
de matrices, c'est-\`a-dire s'il existe un recouvrement \'etale
$\mathcal{U} = \{ \varphi_i : U_i \to S\}_{i \in I}$ tel que
$\varphi_i^*\mathcal{A} \cong \text{Mat}_{d_i}(\mathcal{O}_{U_i})$
pour un $d_i \geq 1$. Deux telles alg\`ebres
$\mathcal{A}$ et $\mathcal{B}$ sont dites {\it \'equivalentes} s'il existe
des $\mathcal{O}_S$-modules localement libres de type fini $\mathcal{F}$ et $\mathcal{G}$
de rang strictement positif en tout $s \in S$, tels que
$$
\mathcal{A} \otimes_{\mathcal{O}_S}
\SheafHom_{\mathcal{O}_S}(\mathcal{F}, \mathcal{F})
\cong
\mathcal{B} \otimes_{\mathcal{O}_S}
\SheafHom_{\mathcal{O}_S}(\mathcal{G}, \mathcal{G})
$$
comme $\mathcal{O}_S$-alg\`ebres. Le {\it groupe de Brauer} de
$S$ est l'ensemble $\text{Br}(S)$ des classes d'\'equivalence d'alg\`ebres
d'Azumaya sur $\mathcal{O}_S$, muni de l'op\'eration induite par le produit tensoriel (sur
$\mathcal{O}_S$).

\begin{lemma}
\label{etale-cohomology-lemma-end-unique-up-to-invertible}
Soit $S$ un sch\'ema. Soient $\mathcal{F}$ et $\mathcal{G}$ des faisceaux de
$\mathcal{O}_S$-modules localement libres de type fini et de rang strictement positif. S'il
existe un isomorphisme
$\SheafHom_{\mathcal{O}_S}(\mathcal{F}, \mathcal{F}) \cong
\SheafHom_{\mathcal{O}_S}(\mathcal{G}, \mathcal{G})$
de $\mathcal{O}_S$-alg\`ebres, alors il existe un faisceau inversible
$\mathcal{L}$ sur $S$ tel que
$\mathcal{F} \otimes_{\mathcal{O}_S} \mathcal{L} \cong \mathcal{G}$
et tel que cet isomorphisme induise l'isomorphisme donn\'e des
alg\`ebres d'endomorphismes.
\end{lemma}

\begin{proof}
Fixons un isomorphisme
$\SheafHom_{\mathcal{O}_S}(\mathcal{F}, \mathcal{F}) \to
\SheafHom_{\mathcal{O}_S}(\mathcal{G}, \mathcal{G})$.
Consid\'erons le faisceau $\mathcal{L} \subset \SheafHom(\mathcal{F}, \mathcal{G})$
engendr\'e, comme $\mathcal{O}_S$-module, par les isomorphismes locaux
$\varphi : \mathcal{F} \to \mathcal{G}$ tels que la conjugaison par
$\varphi$ soit l'isomorphisme donn\'e des alg\`ebres d'endomorphismes.
Un calcul local (qui se ram\`ene au cas o\`u $\mathcal{F}$ et $\mathcal{G}$
sont libres de type fini et $S$ est affine) montre que $\mathcal{L}$ est inversible.
Un autre calcul local montre que le morphisme d'\'evaluation
$$
\mathcal{F} \otimes_{\mathcal{O}_S} \mathcal{L} \longrightarrow \mathcal{G}
$$
est un isomorphisme.
\end{proof}

\noindent
L'argument donn\'e dans la preuve du lemme suivant se trouve dans
\cite{Saltman-torsion}.

\begin{lemma}
\label{etale-cohomology-lemma-annihilated-by-degree}
\begin{reference}
Argument tir\'e de \cite{Saltman-torsion}.
\end{reference}
Soit $S$ un sch\'ema. Soit $\mathcal{A}$ une alg\`ebre d'Azumaya qui est
localement libre de rang $d^2$ sur $S$. Alors la classe
de $\mathcal{A}$ dans le groupe de Brauer de $S$ est annul\'ee par $d$.
\end{lemma}

\begin{proof}
Choisissons un recouvrement \'etale $\{U_i \to S\}$ et des isomorphismes
$\mathcal{A}|_{U_i} \to \SheafHom(\mathcal{F}_i, \mathcal{F}_i)$
pour des $\mathcal{O}_{U_i}$-modules localement libres $\mathcal{F}_i$
de rang $d$. (On peut supposer $\mathcal{F}_i$ libre.) Consid\'erons la
compos\'ee
$$
p_i : \mathcal{F}_i^{\otimes d} \to
\wedge^d(\mathcal{F}_i) \to \mathcal{F}_i^{\otimes d}
$$
La premi\`ere fl\`eche est la projection usuelle et la seconde est
l'isomorphisme de la puissance ext\'erieure maximale de $\mathcal{F}_i$ sur
le sous-module des sections de $\mathcal{F}_i^{\otimes d}$ qui se transforment
suivant le caract\`ere signature sous l'action du groupe sym\'etrique
sur $d$ lettres. Alors $p_i^2 = d! p_i$ et le rang de $p_i$ vaut $1$.
En utilisant l'isomorphisme donn\'e
$\mathcal{A}|_{U_i} \to \SheafHom(\mathcal{F}_i, \mathcal{F}_i)$
et l'isomorphisme canonique
$$
\SheafHom(\mathcal{F}_i, \mathcal{F}_i)^{\otimes d} =
\SheafHom(\mathcal{F}_i^{\otimes d}, \mathcal{F}_i^{\otimes d})
$$
nous pouvons consid\'erer $p_i$ comme une section de $\mathcal{A}^{\otimes d}$
sur $U_i$. Nous affirmons que $p_i|_{U_i \times_S U_j} = p_j|_{U_i \times_S U_j}$
comme sections de $\mathcal{A}^{\otimes d}$. En effet, en appliquant le
lemme \ref{etale-cohomology-lemma-end-unique-up-to-invertible},
nous obtenons un faisceau inversible $\mathcal{L}_{ij}$ et un isomorphisme canonique
$$
\mathcal{F}_i|_{U_i \times_S U_j} \otimes \mathcal{L}_{ij}
\longrightarrow
\mathcal{F}_j|_{U_i \times_S U_j}.
$$
Cet isomorphisme montre que $p_i$ s'envoie sur $p_j$.
Comme $\mathcal{A}^{\otimes d}$ est un faisceau sur $S_\etale$
(proposition \ref{etale-cohomology-proposition-quasi-coherent-sheaf-fpqc}), on obtient une
section globale canonique $p \in \Gamma(S, \mathcal{A}^{\otimes d})$. Un calcul local
montre que
$$
\mathcal{H} =
\Im(\mathcal{A}^{\otimes d} \to \mathcal{A}^{\otimes d}, f \mapsto fp)
$$
est un module localement libre de rang $d^d$ et que la multiplication (\`a gauche)
par $\mathcal{A}^{\otimes d}$ induit un isomorphisme
$\mathcal{A}^{\otimes d} \to \SheafHom(\mathcal{H}, \mathcal{H})$.
Autrement dit, $\mathcal{A}^{\otimes d}$ est l'\'el\'ement neutre
du groupe de Brauer de $S$, comme voulu.
\end{proof}

\noindent
Dans ce cadre, l'analogue de l'isomorphisme $\delta$ du
th\'eor\`eme \ref{etale-cohomology-theorem-brauer-delta}
est une application
$$
\delta_S: \text{Br}(S) \to H_\etale^2(S, \mathbf{G}_m).
$$
L'application $\delta_S$ est injective. Si $S$ est quasi-compact ou
connexe, alors $\text{Br}(S)$ est un groupe de torsion; dans ce cas, l'image
de $\delta_S$ est contenue dans le {\it groupe de Brauer cohomologique} de $S$
$$
\text{Br}'(S) := H_\etale^2(S, \mathbf{G}_m)_\text{torsion}.
$$
Ainsi, si $S$ est quasi-compact ou connexe, on a une inclusion $\text{Br}(S)
\subset \text{Br}'(S)$. Ce n'est pas toujours une \'egalit\'e : il existe une
surface singuli\`ere non s\'epar\'ee $S$ pour laquelle $\text{Br}(S) \subset
\text{Br}'(S)$ est une inclusion stricte. Si $S$ est quasi-projectif, alors
$\text{Br}(S) = \text{Br}'(S)$. Cependant, on ignore si cela reste vrai pour
une vari\'et\'e lisse propre sur $\mathbf{C}$, par exemple.




\section{La suite d'Artin-Schreier}
\label{etale-cohomology-section-artin-schreier}

\noindent
Soit $p$ un nombre premier. Soit $S$ un sch\'ema de caract\'eristique $p$.
La suite {\it d'Artin-Schreier} est la suite exacte courte
$$
0 \longrightarrow \underline{\mathbf{Z}/p\mathbf{Z}}_S \longrightarrow
\mathbf{G}_{a, S} \xrightarrow{F-1} \mathbf{G}_{a, S} \longrightarrow 0
$$
o\`u $F - 1$ est l'application $x \mapsto x^p - x$.

\begin{lemma}
\label{etale-cohomology-lemma-vanishing-affine-char-p-p}
Soit $p$ un nombre premier. Soit $S$ un sch\'ema de caract\'eristique $p$.
\begin{enumerate}
\item Si $S$ est affine, alors
$H_\etale^q(S, \underline{\mathbf{Z}/p\mathbf{Z}}) = 0$ pour tout
$q \geq 2$.
\item Si $S$ est un sch\'ema quasi-compact et quasi-s\'epar\'e de
dimension $d$, alors $H_\etale^q(S, \underline{\mathbf{Z}/p\mathbf{Z}}) = 0$
pour tout $q \geq 2 + d$.
\end{enumerate}
\end{lemma}

\begin{proof}
Rappelons que la cohomologie \'etale du faisceau structural est \'egale
\`a sa cohomologie sur l'espace topologique sous-jacent
(théorème \ref{etale-cohomology-theorem-zariski-fpqc-quasi-coherent}).
La premi\`ere assertion r\'esulte de la suite exacte d'Artin-Schreier
et de l'annulation de la cohomologie du faisceau structural sur un sch\'ema
affine (Cohomologie des schémas, lemme
\ref{coherent-lemma-quasi-coherent-affine-cohomology-zero}).
La seconde assertion r\'esulte par le m\^eme argument de l'annulation
donn\'ee par Cohomologie, Proposition
\ref{cohomology-proposition-cohomological-dimension-spectral}
et du fait que $S$ est un espace spectral
(Properties, lemme
\ref{properties-lemma-quasi-compact-quasi-separated-spectral}).
\end{proof}

\begin{lemma}
\label{etale-cohomology-lemma-F-1}
Soit $k$ un corps alg\'ebriquement clos de caract\'eristique $p > 0$.
Soit $V$ un $k$-espace vectoriel de dimension finie. Soit $F : V \to V$
une application semi-lin\'eaire relativement au Frobenius, c'est-\`a-dire une application additive telle que
$F(\lambda v) = \lambda^p F(v)$ pour tous $\lambda \in k$ et $v \in V$.
Alors $F - 1 : V \to V$ est surjective, et son noyau est un
$\mathbf{F}_p$-espace vectoriel de dimension finie, celle-ci \'etant $\leq \dim_k(V)$.
\end{lemma}

\begin{proof}
Si $F = 0$, l'\'enonc\'e est vrai. Si $V$ est muni d'une filtration par des
sous-espaces vectoriels stables par $F$ telle que l'\'enonc\'e soit vrai pour chaque
quotient gradu\'e, alors il est vrai pour $(V, F)$. En combinant ces deux remarques,
nous pouvons supposer que le noyau de $F$ est nul.

\medskip\noindent
Choisissons une base $v_1, \ldots, v_n$ de $V$ et \'ecrivons
$F(v_i) = \sum a_{ij} v_j$. Remarquons que $v = \sum \lambda_i v_i$
appartient au noyau si et seulement si $\sum \lambda_i^p a_{ij} v_j = 0$.
Comme $k$ est alg\'ebriquement clos, cela implique que la matrice $(a_{ij})$
est inversible. Notons $(b_{ij})$ son inverse. Pour voir que $F - 1$
est surjective, choisissons $w = \sum \mu_i v_i \in V$ et cherchons \`a r\'esoudre
$$
(F - 1)(\sum \lambda_iv_i) =
\sum \lambda_i^p a_{ij} v_j - \sum \lambda_j v_j = \sum \mu_j v_j
$$
Cela \'equivaut \`a
$$
\sum \lambda_j^p v_j - \sum b_{ij} \lambda_i v_j = \sum b_{ij} \mu_i v_j
$$
autrement dit
$$
\lambda_j^p - \sum b_{ij} \lambda_i = \sum b_{ij} \mu_i,
\quad j = 1, \ldots, \dim(V).
$$
L'alg\`ebre
$$
A = k[x_1, \ldots, x_n]/
(x_j^p - \sum b_{ij} x_i - \sum b_{ij} \mu_i)
$$
est lisse standard sur $k$
(Algèbre, définition \ref{algebra-definition-standard-smooth})
car la matrice $(b_{ij})$ est inversible et les d\'eriv\'ees partielles
des $x_j^p$ sont nulles. Une base de $A$ sur $k$ est l'ensemble des mon\^omes
$x_1^{e_1} \ldots x_n^{e_n}$ avec $e_i < p$, donc $\dim_k(A) = p^n$.
Comme $k$ est alg\'ebriquement clos, on voit que $\Spec(A)$ poss\`ede exactement
$p^n$ points. Il s'ensuit que $F - 1$ est surjective et que chaque fibre
poss\`ede $p^n$ points, c'est-\`a-dire que le noyau de $F - 1$ est un groupe d'ordre $p^n$.
\end{proof}

\begin{lemma}
\label{etale-cohomology-lemma-F-2}
Dans la situation du lemme \ref{etale-cohomology-lemma-F-1}, soit $k'/k$ une extension
de corps alg\'ebriquement clos. Posons $V' = V \otimes_k k'$ et soit
$F' : V' \to V'$ l'application semi-lin\'eaire relativement au Frobenius induite par $F$.
Alors $\Ker(F - 1)$ s'envoie isomorphiquement sur $\Ker(F' - 1)$.
\end{lemma}

\begin{proof}
Soit $v_1, \ldots, v_n$ une base de $V$ et \'ecrivons $F(v_i) = \sum a_{ij}v_j$.
Alors $v = \sum \lambda_i v_i$ appartient \`a $\Ker(F - 1)$ si et seulement si
$\sum \lambda_i^pa_{ij} = \lambda_j$ pour $j = 1, \ldots, n$.
D'apr\`es le lemme \ref{etale-cohomology-lemma-F-1}, ces \'equations
d\'efinissent un sous-sch\'ema ferm\'e $Z$ de
$\mathbf{A}^n_k = \Spec(k[\lambda_1, \ldots, \lambda_n])$
fini sur $\Spec(k)$. Alors $Z(k) = Z(k')$, ce qui permet de conclure.
\end{proof}

\begin{lemma}
\label{etale-cohomology-lemma-top-cohomology-coherent}
Soit $X$ un sch\'ema s\'epar\'e de type fini sur un corps $k$.
Soit $\mathcal{F}$ un faisceau coh\'erent de $\mathcal{O}_X$-modules.
Alors $\dim_k H^d(X, \mathcal{F}) < \infty$, o\`u $d = \dim(X)$.
\end{lemma}

\begin{proof}
Nous allons le d\'emontrer par r\'ecurrence sur $d$. Le cas $d = 0$ est vrai car
dans ce cas $X$ est le spectre d'une $k$-alg\`ebre $A$ de dimension finie
(Variétés, lemme \ref{varieties-lemma-algebraic-scheme-dim-0})
et tout faisceau coh\'erent $\mathcal{F}$ correspond \`a un $A$-module de type fini
$M = H^0(X, \mathcal{F})$ qui v\'erifie $\dim_k M < \infty$.

\medskip\noindent
Supposons $d > 0$ et le r\'esultat d\'emontr\'e pour les sch\'emas s\'epar\'es
de type fini de dimension $< d$. Le sch\'ema $X$ est noeth\'erien. Consid\'erons
la propri\'et\'e $\mathcal{P}$ des faisceaux coh\'erents sur $X$ d\'efinie par la r\`egle
$$
\mathcal{P}(\mathcal{F}) \Leftrightarrow
\dim_k H^d(X, \mathcal{F}) < \infty
$$
Nous allons utiliser le r\'esultat de
Cohomologie des schémas, lemme \ref{coherent-lemma-property-initial}
pour montrer que $\mathcal{P}$ est satisfaite par tout faisceau coh\'erent sur $X$.

\medskip\noindent
Soit
$$
0 \to \mathcal{F}_1 \to \mathcal{F} \to \mathcal{F}_2 \to 0
$$
une suite exacte courte de faisceaux coh\'erents sur $X$.
Consid\'erons la suite exacte longue de cohomologie
$$
H^d(X, \mathcal{F}_1) \to
H^d(X, \mathcal{F}) \to
H^d(X, \mathcal{F}_2)
$$
Ainsi, si $\mathcal{P}$ est satisfaite par $\mathcal{F}_1$ et $\mathcal{F}_2$,
elle l'est par $\mathcal{F}$.

\medskip\noindent
Soit $Z \subset X$ un sous-sch\'ema ferm\'e int\`egre. Soit $\mathcal{I}$
un faisceau coh\'erent d'id\'eaux sur $Z$. Pour achever la d\'emonstration, il faut montrer
que $H^d(X, i_*\mathcal{I}) = H^d(Z, \mathcal{I})$ est de dimension finie.
Si $\dim(Z) < d$, le r\'esultat est vrai car ce groupe de cohomologie
est nul (Cohomologie, Proposition
\ref{cohomology-proposition-vanishing-Noetherian}).
Nous sommes ainsi ramen\'es \`a la situation examin\'ee au paragraphe suivant.

\medskip\noindent
Supposons que $X$ soit une vari\'et\'e de dimension $d$ et que
$\mathcal{F} = \mathcal{I}$ soit un faisceau coh\'erent d'id\'eaux. Dans ce
cas, nous avons une suite exacte courte
$$
0 \to \mathcal{I} \to \mathcal{O}_X \to i_*\mathcal{O}_Z \to 0
$$
o\`u $i : Z \to X$ est l'immersion ferm\'ee d\'efinie par $\mathcal{I}$.
L'hypoth\`ese de r\'ecurrence montre que
$H^{d - 1}(Z, \mathcal{O}_Z) = H^{d - 1}(X, i_*\mathcal{O}_Z)$ est
de dimension finie. Il suffit donc de d\'emontrer le r\'esultat
pour le faisceau structural.

\medskip\noindent
Nous pouvons appliquer le lemme de Chow
(Cohomologie des schémas, lemme \ref{coherent-lemma-chow-Noetherian})
au morphisme $X \to \Spec(k)$. Nous obtenons ainsi un diagramme
$$
\xymatrix{
X \ar[rd]_g & X' \ar[d]^-{g'} \ar[l]^\pi \ar[r]_i & \mathbf{P}^n_k \ar[dl] \\
& \Spec(k) &
}
$$
comme dans l'\'enonc\'e du lemme de Chow. Soit en outre $U \subset X$
le sous-sch\'ema ouvert dense tel que $\pi^{-1}(U) \to U$ soit un isomorphisme.
Nous pouvons aussi supposer que $X'$ est une vari\'et\'e, voir
Cohomologie des schémas, remarque \ref{coherent-remark-chow-Noetherian}.
Le morphisme $i' = (i, \pi) : X' \to \mathbf{P}^n_X$ est
une immersion ferm\'ee (loc. cit.). Par cons\'equent
$$
\mathcal{L} = i^*\mathcal{O}_{\mathbf{P}^n_k}(1) \cong
(i')^*\mathcal{O}_{\mathbf{P}^n_X}(1)
$$
est ample relativement \`a $\pi$ (par exemple d'apr\`es
Morphismes, lemme \ref{morphisms-lemma-characterize-ample-on-finite-type}).
Ainsi, d'apr\`es Cohomologie des schémas, lemme \ref{coherent-lemma-kill-by-twisting},
il existe un $n \geq 0$ tel que
$R^p\pi_*\mathcal{L}^{\otimes n} = 0$ pour tout $p > 0$.
Posons $\mathcal{G} = \pi_*\mathcal{L}^{\otimes n}$.
Choisissons une section globale non nulle $s$ de $\mathcal{L}^{\otimes n}$.
Puisque $\mathcal{G} = \pi_*\mathcal{L}^{\otimes n}$, la section $s$
correspond \`a une section de $\mathcal{G}$, c'est-\`a-dire \`a un morphisme
$\mathcal{O}_X \to \mathcal{G}$.
Comme $s|_U \not = 0$, puisque $X'$ est une vari\'et\'e et $\mathcal{L}$ est
inversible, on voit que $\mathcal{O}_X|_U \to \mathcal{G}|_U$
est non nul. Comme $\mathcal{G}|_U = \mathcal{L}^{\otimes n}|_{\pi^{-1}(U)}$
est inversible, nous obtenons une suite exacte courte
$$
0 \to \mathcal{O}_X \to \mathcal{G} \to \mathcal{Q} \to 0
$$
o\`u $\mathcal{Q}$ est coh\'erent et \`a support dans un sous-sch\'ema
ferm\'e strict de $X$. En raisonnant comme ci-dessus avec notre hypoth\`ese de
r\'ecurrence, on voit qu'il
suffit de d\'emontrer $\dim H^d(X, \mathcal{G}) < \infty$.

\medskip\noindent
La suite spectrale de Leray
(Cohomologie, lemme \ref{cohomology-lemma-apply-Leray})
montre que $H^d(X, \mathcal{G}) = H^d(X', \mathcal{L}^{\otimes n})$.
Soit $\overline{X}' \subset \mathbf{P}^n_k$ l'adh\'erence
de $X'$. Alors $\overline{X}'$ est une vari\'et\'e projective de dimension $d$
sur $k$ et $X' \subset \overline{X}'$ est un ouvert dense.
Le faisceau inversible $\mathcal{L}^{\otimes n}$ est la restriction de
$\mathcal{O}_{\overline{X}'}(n)$ \`a $X'$. D'apr\`es
Cohomologie, Proposition
\ref{cohomology-proposition-cohomological-dimension-spectral}
l'application
$$
H^d(\overline{X}', \mathcal{O}_{\overline{X}'}(n))
\longrightarrow
H^d(X', \mathcal{L}^{\otimes n})
$$
est surjective. Comme le groupe de cohomologie de gauche est
de dimension finie d'apr\`es
Cohomologie des schémas, lemme \ref{coherent-lemma-coherent-projective}
la d\'emonstration est achev\'ee.
\end{proof}

\begin{lemma}
\label{etale-cohomology-lemma-vanishing-variety-char-p-p}
Soit $X$ un sch\'ema s\'epar\'e de type fini sur un corps alg\'ebriquement clos
$k$ de caract\'eristique $p > 0$. Alors
$H_\etale^q(X, \underline{\mathbf{Z}/p\mathbf{Z}}) = 0$ pour
$q \geq dim(X) + 1$.
\end{lemma}

\begin{proof}
Posons $d = \dim(X)$. D'apr\`es l'annulation \'etablie dans le
lemme \ref{etale-cohomology-lemma-vanishing-affine-char-p-p},
il suffit de montrer que
$H_\etale^{d + 1}(X, \underline{\mathbf{Z}/p\mathbf{Z}}) = 0$.
D'apr\`es le lemme \ref{etale-cohomology-lemma-top-cohomology-coherent}, on voit que
$H^d(X, \mathcal{O}_X)$ est un espace vectoriel de dimension finie sur $k$.
Par cons\'equent, la suite exacte longue de cohomologie associ\'ee \`a la
suite d'Artin-Schreier se termine par
$$
H^d(X, \mathcal{O}_X) \xrightarrow{F - 1}
H^d(X, \mathcal{O}_X) \to H^{d + 1}_\etale(X, \mathbf{Z}/p\mathbf{Z}) \to 0
$$
D'apr\`es le lemme \ref{etale-cohomology-lemma-F-1}, l'application $F - 1$ de cette suite est surjective.
Cela d\'emontre le lemme.
\end{proof}

\begin{lemma}
\label{etale-cohomology-lemma-finiteness-proper-variety-char-p-p}
Soit $X$ un sch\'ema propre sur un corps alg\'ebriquement clos
$k$ de caract\'eristique $p > 0$. Alors
\begin{enumerate}
\item $H_\etale^q(X, \underline{\mathbf{Z}/p\mathbf{Z}})$
est un $\mathbf{Z}/p\mathbf{Z}$-module de type fini pour tout $q$, et
\item $H^q_\etale(X, \underline{\mathbf{Z}/p\mathbf{Z}}) \to
H^q_\etale(X_{k'}, \underline{\mathbf{Z}/p\mathbf{Z}})$
est un isomorphisme si $k'/k$ est une extension de corps alg\'ebriquement
clos.
\end{enumerate}
\end{lemma}

\begin{proof}
D'apr\`es Cohomologie des schémas, lemme
\ref{coherent-lemma-proper-over-affine-cohomology-finite}
et la comparaison de la cohomologie donn\'ee par le
th\'eor\`eme \ref{etale-cohomology-theorem-zariski-fpqc-quasi-coherent}, les groupes de cohomologie
$H^q_\etale(X, \mathbf{G}_a) = H^q(X, \mathcal{O}_X)$ sont
des espaces vectoriels de dimension finie sur $k$. Par cons\'equent, d'apr\`es le
lemme \ref{etale-cohomology-lemma-F-1}, la suite exacte longue de cohomologie
associ\'ee \`a la suite d'Artin-Schreier se d\'ecompose en
suites exactes courtes
$$
0 \to H_\etale^q(X, \underline{\mathbf{Z}/p\mathbf{Z}}) \to
H^q(X, \mathcal{O}_X) \xrightarrow{F - 1} H^q(X, \mathcal{O}_X) \to 0
$$
et, de plus, la dimension sur $\mathbf{F}_p$ des groupes de cohomologie
$H_\etale^q(X, \underline{\mathbf{Z}/p\mathbf{Z}})$ est inf\'erieure ou \'egale \`a la
dimension sur $k$ de l'espace vectoriel $H^q(X, \mathcal{O}_X)$.
Cela d\'emontre la premi\`ere assertion. La seconde r\'esulte
du lemme \ref{etale-cohomology-lemma-F-2}
car $H^q(X, \mathcal{O}_X) \otimes_k k' \to H^q(X_{k'}, \mathcal{O}_{X_{k'}})$
est un isomorphisme par changement de base plat
(Cohomologie des schémas,
lemme \ref{coherent-lemma-flat-base-change-cohomology}).
\end{proof}








\section{Faisceaux localement constants}
\label{etale-cohomology-section-locally-constant}

\noindent
Cette section est l'analogue de
Modules sur les sites, section \ref{sites-modules-section-locally-constant}
pour le site \'etale.

\begin{definition}
\label{etale-cohomology-definition-finite-locally-constant}
Soit $X$ un sch\'ema.
Soit $\mathcal{F}$ un faisceau d'ensembles sur $X_\etale$.
\begin{enumerate}
\item Soit $E$ un ensemble. Nous disons que $\mathcal{F}$ est le
{\it faisceau constant de valeur $E$} si $\mathcal{F}$ est le
faisceau associ\'e au pr\'efaisceau $U \mapsto E$.
Notation : $\underline{E}_X$ ou $\underline{E}$.
\item Nous disons que $\mathcal{F}$ est un {\it faisceau constant} s'il est
isomorphe \`a un faisceau comme en (1).
\item Nous disons que $\mathcal{F}$ est {\it localement constant} s'il existe un
recouvrement $\{U_i \to X\}$ tel que $\mathcal{F}|_{U_i}$ soit un faisceau constant.
\item Nous disons que $\mathcal{F}$ est {\it localement constant fini} s'il
est localement constant et si ses valeurs locales sont des ensembles finis.
\end{enumerate}
Soit $\mathcal{F}$ un faisceau de groupes ab\'eliens sur $X_\etale$.
\begin{enumerate}
\item Soit $A$ un groupe ab\'elien.
Nous disons que $\mathcal{F}$ est le {\it faisceau constant de valeur $A$} si
$\mathcal{F}$ est le faisceau associ\'e au pr\'efaisceau $U \mapsto A$.
Notation : $\underline{A}_X$ ou $\underline{A}$.
\item Nous disons que $\mathcal{F}$ est un {\it faisceau constant} s'il est isomorphe
comme faisceau ab\'elien \`a un faisceau comme en (1).
\item Nous disons que $\mathcal{F}$ est {\it localement constant} s'il existe un
recouvrement $\{U_i \to X\}$ tel que $\mathcal{F}|_{U_i}$ soit un faisceau constant.
\item Nous disons que $\mathcal{F}$ est {\it localement constant fini} s'il
est localement constant et si ses valeurs locales sont des groupes ab\'eliens finis.
\end{enumerate}
Soit $\Lambda$ un anneau. Soit $\mathcal{F}$ un faisceau de $\Lambda$-modules
sur $X_\etale$.
\begin{enumerate}
\item Soit $M$ un $\Lambda$-module.
Nous disons que $\mathcal{F}$ est le {\it faisceau constant de valeur $M$} si
$\mathcal{F}$ est le faisceau associ\'e au pr\'efaisceau $U \mapsto M$.
Notation : $\underline{M}_X$ ou $\underline{M}$.
\item Nous disons que $\mathcal{F}$ est un {\it faisceau constant} s'il est isomorphe
comme faisceau de $\Lambda$-modules \`a un faisceau comme en (1).
\item Nous disons que $\mathcal{F}$ est {\it localement constant} s'il existe un
recouvrement $\{U_i \to X\}$ tel que $\mathcal{F}|_{U_i}$ soit un faisceau constant.
\end{enumerate}
\end{definition}

\begin{lemma}
\label{etale-cohomology-lemma-pullback-locally-constant}
Soit $f : X \to Y$ un morphisme de sch\'emas. Si $\mathcal{G}$ est un
faisceau localement constant d'ensembles, de groupes ab\'eliens ou de $\Lambda$-modules
sur $Y_\etale$, alors $f^{-1}\mathcal{G}$ l'est aussi
sur $X_\etale$.
\end{lemma}

\begin{proof}
Cela vaut pour tout morphisme de topos, voir
Modules sur les sites, lemme \ref{sites-modules-lemma-pullback-locally-constant}.
\end{proof}

\begin{lemma}
\label{etale-cohomology-lemma-pushforward-locally-constant}
Soit $f : X \to Y$ un morphisme fini \'etale de sch\'emas.
Si $\mathcal{F}$ est un faisceau d'ensembles localement constant (fini),
un faisceau de groupes ab\'eliens localement constant (fini), ou
un faisceau de $\Lambda$-modules localement constant (de type fini)
sur $X_\etale$, alors $f_*\mathcal{F}$ l'est aussi
sur $Y_\etale$.
\end{lemma}

\begin{proof}
La construction de $f_*$ commute \`a la localisation \'etale.
Un morphisme fini \'etale est localement isomorphe \`a une somme disjointe
d'isomorphismes, voir
Morphismes étales, lemme \ref{etale-lemma-finite-etale-etale-local}.
Ainsi le lemme affirme que, si $\mathcal{F}_i$, $i = 1, \ldots, n$
sont des faisceaux d'ensembles localement constants (finis), alors
$\prod_{i = 1, \ldots, n} \mathcal{F}_i$ l'est aussi.
C'est clair. Il en va de m\^eme pour les faisceaux de groupes ab\'eliens et de modules.
\end{proof}

\begin{lemma}
\label{etale-cohomology-lemma-characterize-finite-locally-constant}
Soient $X$ un sch\'ema et $\mathcal{F}$ un faisceau d'ensembles sur $X_\etale$.
Alors les assertions suivantes sont \`equivalentes
\begin{enumerate}
\item $\mathcal{F}$ est localement constant fini, et
\item $\mathcal{F} = h_U$ pour un morphisme fini \'etale $U \to X$.
\end{enumerate}
\end{lemma}

\begin{proof}
Un morphisme fini \'etale est localement isomorphe \`a une somme disjointe
d'isomorphismes, voir
Morphismes étales, lemme \ref{etale-lemma-finite-etale-etale-local}.
Ainsi (2) implique (1). R\'eciproquement, si $\mathcal{F}$ est localement
constant fini, il existe un recouvrement \'etale $\{X_i \to X\}$ tel que
$\mathcal{F}|_{X_i}$ soit repr\'esentable par un morphisme fini \'etale $U_i \to X_i$.
En raisonnant exactement comme dans la d\'emonstration du
Descente, lemme \ref{descent-lemma-descent-data-sheaves}
nous obtenons une donn\'ee de descente de sch\'emas $(U_i, \varphi_{ij})$ relativement \`a
$\{X_i \to X\}$ (d\'etails omis). Cette donn\'ee de descente est effective, par
exemple d'apr\`es Descente, lemme \ref{descent-lemma-affine},
et le morphisme de sch\'emas obtenu $U \to X$ est fini \'etale
d'apr\`es Descente, lemmes \ref{descent-lemma-descending-property-finite} et
\ref{descent-lemma-descending-property-etale}.
\end{proof}

\begin{lemma}
\label{etale-cohomology-lemma-morphism-locally-constant}
Soit $X$ un sch\'ema.
\begin{enumerate}
\item Soit $\varphi : \mathcal{F} \to \mathcal{G}$ un morphisme
de faisceaux d'ensembles localement constants sur $X_\etale$.
Si $\mathcal{F}$ est localement constant fini, il existe un
recouvrement \'etale $\{U_i \to X\}$ tel que
$\varphi|_{U_i}$ soit le morphisme de faisceaux constants associ\'e \`a
une application d'ensembles.
\item Soit $\varphi : \mathcal{F} \to \mathcal{G}$ un morphisme
de faisceaux de groupes ab\'eliens localement constants sur $X_\etale$.
Si $\mathcal{F}$ est localement constant fini, il existe un recouvrement \'etale
$\{U_i \to X\}$ tel que $\varphi|_{U_i}$ soit le morphisme de
faisceaux ab\'eliens constants associ\'e \`a un homomorphisme de groupes ab\'eliens.
\item Soit $\Lambda$ un anneau.
Soit $\varphi : \mathcal{F} \to \mathcal{G}$ un morphisme
de faisceaux de $\Lambda$-modules localement constants sur $X_\etale$.
Si $\mathcal{F}$ est de type fini, alors il existe un recouvrement \'etale
$\{U_i \to X\}$ tel que $\varphi|_{U_i}$ soit le morphisme de
faisceaux constants de $\Lambda$-modules associ\'e \`a un homomorphisme de $\Lambda$-modules.
\end{enumerate}
\end{lemma}

\begin{proof}
Cela vaut sur tout site, voir
Modules sur les sites, lemme \ref{sites-modules-lemma-morphism-locally-constant}.
\end{proof}

\begin{lemma}
\label{etale-cohomology-lemma-kernel-finite-locally-constant}
Soit $X$ un sch\'ema.
\begin{enumerate}
\item La cat\'egorie des faisceaux d'ensembles localement constants finis
est stable par limites projectives et inductives finies dans $\Sh(X_\etale)$.
\item La cat\'egorie des faisceaux ab\'eliens localement constants finis est une
sous-cat\'egorie de Serre faible de $\textit{Ab}(X_\etale)$.
\item Soit $\Lambda$ un anneau noeth\'erien. La cat\'egorie des
faisceaux de $\Lambda$-modules localement constants de type fini sur
$X_\etale$ est une sous-cat\'egorie de Serre faible de
$\textit{Mod}(X_\etale, \Lambda)$.
\end{enumerate}
\end{lemma}

\begin{proof}
Cela vaut sur tout site, voir
Modules sur les sites, lemme
\ref{sites-modules-lemma-kernel-finite-locally-constant}.
\end{proof}

\begin{lemma}
\label{etale-cohomology-lemma-tensor-product-locally-constant}
Soit $X$ un sch\'ema. Soit $\Lambda$ un anneau.
Le produit tensoriel de deux faisceaux de $\Lambda$-modules localement constants
sur $X_\etale$ est un faisceau de $\Lambda$-modules localement constant.
\end{lemma}

\begin{proof}
Cela vaut sur tout site, voir
Modules sur les sites, lemme
\ref{sites-modules-lemma-tensor-product-locally-constant}.
\end{proof}

\begin{lemma}
\label{etale-cohomology-lemma-connected-locally-constant}
Soit $X$ un sch\'ema connexe. Soit $\Lambda$ un anneau et soit
$\mathcal{F}$ un faisceau de $\Lambda$-modules localement constant.
Alors il existe un $\Lambda$-module $M$ et un recouvrement \'etale
$\{U_i \to X\}$ tels que $\mathcal{F}|_{U_i} \cong \underline{M}|_{U_i}$.
\end{lemma}

\begin{proof}
Choisissons un recouvrement \'etale
$\{U_i \to X\}$ tel que $\mathcal{F}|_{U_i}$ soit constant, disons
$\mathcal{F}|_{U_i} \cong \underline{M_i}_{U_i}$.
Remarquons que $U_i \times_X U_j$ est vide si $M_i$ n'est pas isomorphe
\`a $M_j$.
Pour chaque classe d'isomorphisme de $\Lambda$-modules repr\'esent\'ee par $M$, posons $I_M = \{i \in I \mid M_i \cong M\}$.
Comme les morphismes \'etales sont ouverts, nous voyons que
$U_M = \bigcup_{i \in I_M} \Im(U_i \to X)$
est un ouvert de $X$. Alors $X = \coprod U_M$ est un recouvrement
ouvert disjoint de $X$. Comme $X$ est connexe, un seul $U_M$ est non vide
et le lemme en r\'esulte.
\end{proof}






\section{Faisceaux localement constants et groupe fondamental}
\label{etale-cohomology-section-pione}

\noindent
Dans certains cas, nous pouvons relier les faisceaux localement constants
au groupe fondamental d'un sch\'ema.

\begin{lemma}
\label{etale-cohomology-lemma-locally-constant-on-connected}
Soit $X$ un sch\'ema connexe. Soit $\overline{x}$ un point g\'eom\'etrique de $X$.
\begin{enumerate}
\item Il existe une \`equivalence de cat\'egories
$$
\left\{
\begin{matrix}
\text{faisceaux d'ensembles localement constants finis}\\
\text{sur }X_\etale
\end{matrix}
\right\}
\longleftrightarrow
\left\{
\begin{matrix}
\text{ensembles finis munis d'une action de }\pi_1(X, \overline{x})
\end{matrix}
\right\}
$$
\item Il existe une \`equivalence de cat\'egories
$$
\left\{
\begin{matrix}
\text{faisceaux de groupes ab\'eliens localement constants finis}\\
\text{sur }X_\etale
\end{matrix}
\right\}
\longleftrightarrow
\left\{
\begin{matrix}
\pi_1(X, \overline{x})\text{-modules finis}
\end{matrix}
\right\}
$$
\item Soit $\Lambda$ un anneau fini. Il existe une \`equivalence de cat\'egories
$$
\left\{
\begin{matrix}
\text{faisceaux localement constants de type fini}\\
\text{de }\Lambda\text{-modules sur }X_\etale
\end{matrix}
\right\}
\longleftrightarrow
\left\{
\begin{matrix}
\pi_1(X, \overline{x})\text{-modules finis munis}\\
\text{d'une structure commutante de }\Lambda\text{-module}
\end{matrix}
\right\}
$$
\end{enumerate}
\end{lemma}

\begin{proof}
Remarquons que $\pi_1(X, \overline{x})$ est un groupe topologique
profini, voir Fundamental Groups, définition
\ref{pione-definition-fundamental-group}.
Les cat\'egories des membres de gauche sont d\'efinies dans la
section \ref{etale-cohomology-section-locally-constant}.
La notation employ\'ee dans les membres de droite vient de
Fundamental Groups, définition \ref{pione-definition-G-set-continuous}
pour les ensembles et de la
d\'efinition \ref{etale-cohomology-definition-G-module-continuous} pour les groupes ab\'eliens.
Cela explique la notation.

\medskip\noindent
L'assertion (1) r\'esulte du
lemme \ref{etale-cohomology-lemma-characterize-finite-locally-constant}
et de Fundamental Groups, théorème \ref{pione-theorem-fundamental-group}.
Les parties (2) et (3) s'en d\'eduisent imm\'ediatement en munissant les
(faisceaux d')ensembles sous-jacents d'une structure suppl\'ementaire. Par exemple, un
faisceau de groupes ab\'eliens localement constant fini sur $X_\etale$ revient
\`a un faisceau d'ensembles localement constant fini $\mathcal{F}$
muni d'un morphisme $+ : \mathcal{F} \times \mathcal{F} \to \mathcal{F}$
satisfaisant aux axiomes usuels. L'\'equivalence de (1) pr\'eserve les produits
et transforme donc $+$ en une addition sur l'ensemble fini correspondant
muni d'une action de $\pi_1(X, \overline{x})$. Puisque les $\pi_1(X, \overline{x})$-modules
reviennent aux ensembles munis d'une action de $\pi_1(X, \overline{x})$ et d'une
structure compatible de groupe ab\'elien, nous obtenons (2). La partie (3) se d\'emontre
exactement de la m\^eme mani\`ere.
\end{proof}

\begin{lemma}
\label{etale-cohomology-lemma-locally-constant-on-connected-geom-unibranch}
Soit $X$ un sch\'ema irr\'eductible et g\'eom\'etriquement unibranche.
Soit $\overline{x}$ un point g\'eom\'etrique de $X$.
Soit $\Lambda$ un anneau. Il existe une \`equivalence de cat\'egories
$$
\left\{
\begin{matrix}
\text{faisceaux localement constants de type fini}\\
\text{de }\Lambda\text{-modules sur }X_\etale
\end{matrix}
\right\}
\longleftrightarrow
\left\{
\begin{matrix}
\Lambda\text{-modules de type fini }M\text{ munis}\\
\text{d'une action continue de }\pi_1(X, \overline{x})
\end{matrix}
\right\}
$$
\end{lemma}

\begin{proof}
La d\'emonstration du lemme \ref{etale-cohomology-lemma-locally-constant-on-connected}
ne s'applique pas, car un $\Lambda$-module de type fini $M$ peut ne pas avoir
un ensemble sous-jacent fini.

\medskip\noindent
Soit $\nu : X^\nu \to X$ le morphisme de normalisation. D'apr\`es Morphismes, lemme
\ref{morphisms-lemma-normalization-geom-unibranch-univ-homeo}
il s'agit d'un hom\'eomorphisme universel. D'apr\`es
Fundamental Groups, Proposition \ref{pione-proposition-universal-homeomorphism}
il induit un isomorphisme
$\pi_1(X^\nu, \overline{x}) \to \pi_1(X, \overline{x})$
et, d'apr\`es le th\'eor\`eme \ref{etale-cohomology-theorem-topological-invariance}, nous obtenons
une \`equivalence de cat\'egories entre les faisceaux localement constants de type fini
de $\Lambda$-modules sur $X_\etale$ et sur $X^\nu_\etale$.
Nous sommes ainsi ramen\'es au cas o\`u $X$ est un sch\'ema int\`egre normal.

\medskip\noindent
Supposons que $X$ soit un sch\'ema int\`egre normal. Soit $\eta \in X$ le point g\'en\'erique.
Soit $\overline{\eta}$ un point g\'eom\'etrique au-dessus de $\eta$.
D'apr\`es Fundamental Groups, Proposition \ref{pione-proposition-normal},
nous avons une surjection continue
$$
\text{Gal}(\kappa(\eta)^{sep}/\kappa(\eta)) = \pi_1(\eta, \overline{\eta})
\longrightarrow
\pi_1(X, \overline{\eta})
$$
dont le noyau est d\'ecrit dans Fundamental Groups, lemme
\ref{pione-lemma-normal-pione-quotient-inertia}.
Soit $\mathcal{F}$ un faisceau localement constant de type fini de $\Lambda$-modules
sur $X_\etale$. Soit $M = \mathcal{F}_{\overline{\eta}}$
la fibre de $\mathcal{F}$ en $\overline{\eta}$.
La section \ref{etale-cohomology-section-galois-action-stalks} nous donne une action continue de
$\text{Gal}(\kappa(\eta)^{sep}/\kappa(\eta))$ sur $M$.
Notre but est de montrer que cette action se factorise par la
surjection affich\'ee.
Puisque $\mathcal{F}$ est de type fini, $M$ est un $\Lambda$-module de type fini.
Puisque $\mathcal{F}$ est localement constant, pour tout $x \in X$,
la restriction de $\mathcal{F}$ \`a $\Spec(\mathcal{O}_{X, x}^{sh})$
est constante. L'action de $\text{Gal}(K^{sep}/K_x^{sh})$
(avec les notations de Fundamental Groups, lemme
\ref{pione-lemma-normal-pione-quotient-inertia})
sur $M$ est donc triviale. Nous obtenons la factorisation voulue.

\medskip\noindent
R\'eciproquement, supposons donn\'e un $\Lambda$-module de type fini $M$
muni d'une action continue de $\pi_1(X, \overline{\eta})$.
Nous allons construire un $\mathcal{F}$ tel que
$M \cong \mathcal{F}_{\overline{\eta}}$ comme
$\Lambda[\pi_1(X, \overline{\eta})]$-modules.
Choisissons des g\'en\'erateurs $m_1, \ldots, m_r \in M$.
Comme l'action de $\pi_1(X, \overline{\eta})$ sur $M$ est
continue, il existe pour tout $i$ un sous-groupe ouvert
$H_i$ du groupe profini $\pi_1(X, \overline{\eta})$
tel que tout $\gamma \in H_i$ fixe $m_i$. Nous en d\'eduisons que
tout \`el\'ement du sous-groupe ouvert $H = \bigcap_{i = 1, \ldots, r} H_i$
fixe tout \`el\'ement de $M$. En rempla\c cant $H$ par un sous-groupe plus petit, nous pouvons supposer que $H$
est un sous-groupe ouvert normal de $\pi_1(X, \overline{\eta})$.
Posons $G = \pi_1(X, \overline{\eta})/H$. Soit $f : Y \to X$ le
rev\^etement galoisien fini \'etale de groupe $G$ correspondant.
Nous pouvons voir $f_*\underline{\mathbf{Z}}$ comme un faisceau de
$\mathbf{Z}[G]$-modules sur $X_\etale$. Il suffit alors de poser
$$
\mathcal{F} =
f_*\underline{\mathbf{Z}} \otimes_{\underline{\mathbf{Z}[G]}} \underline{M}
$$
Nous laissons au lecteur le soin de calculer
$\mathcal{F}_{\overline{\eta}}$.
Nous omettons aussi de v\'erifier que cette construction
est inverse de celle du paragraphe pr\'ec\'edent.
\end{proof}

\begin{remark}
\label{etale-cohomology-remark-functorial-locally-constant-on-connected}
Les \`equivalences des lemmes \ref{etale-cohomology-lemma-locally-constant-on-connected} et
\ref{etale-cohomology-lemma-locally-constant-on-connected-geom-unibranch}
sont compatibles aux images inverses. Par exemple, soit $f : Y \to X$
un morphisme de sch\'emas connexes. Soit $\overline{y}$ un point g\'eom\'etrique
de $Y$ et posons $\overline{x} = f(\overline{y})$.
Alors le diagramme
$$
\xymatrix{
\substack{\text{faisceaux d'ensembles finis} \\
\text{localement constants sur }Y_\etale}
\ar[r] &
\substack{\text{ensembles finis} \\
\text{munis d'une action de }\pi_1(Y, \overline{y})} \\
\substack{\text{faisceaux d'ensembles finis} \\
\text{localement constants sur }X_\etale}
\ar[r] \ar[u]_{f^{-1}} &
\substack{\text{ensembles finis} \\
\text{munis d'une action de }\pi_1(X, \overline{x})} \ar[u]
}
$$
est commutatif, o\`u la fl\`eche verticale de droite provient de
l'homomorphisme continu
$\pi_1(Y, \overline{y}) \to \pi_1(X, \overline{x})$
induit par $f$. Cela r\'esulte imm\'ediatement du
diagramme commutatif de
Fundamental Groups, théorème \ref{pione-theorem-fundamental-group}.
Un r\'esultat analogue vaut dans les autres cas.
\end{remark}










\section{M\'ethode de la trace}
\label{etale-cohomology-section-trace-method}

\noindent
Une r\'ef\'erence pour cette section est \cite[Expos\'e IX, \S 5]{SGA4}.
Les r\'esultats qui suivent seront utilis\'es dans la d\'emonstration du
lemme \ref{etale-cohomology-lemma-vanishing-easier} ci-dessous.

\medskip\noindent
Soit $f : Y \to X$ un morphisme \'etale de sch\'emas. Il existe
une suite
$$
f_!, f^{-1}, f_*
$$
de foncteurs adjoints entre
$\textit{Ab}(X_\etale)$ et $\textit{Ab}(Y_\etale)$.
Le foncteur $f_!$ est \'etudi\'e dans la section \ref{etale-cohomology-section-extension-by-zero}.
Le morphisme d'adjonction $\text{id} \to f_* f^{-1}$ est appel\'e {\it restriction}.
Le morphisme d'adjonction $f_! f^{-1} \to \text{id}$ est souvent
appel\'e le {\it morphisme trace}. Si $f$ est fini \'etale, alors $f_* = f_!$
(lemme \ref{etale-cohomology-lemma-shriek-equals-star-finite-etale}) et
nous pouvons le voir comme un morphisme $f_*f^{-1} \to \text{id}$.

\begin{definition}
\label{etale-cohomology-definition-trace-map}
Soit $f : Y \to X$ un morphisme fini \'etale de sch\'emas.
Le morphisme $f_* f^{-1} \to \text{id}$ d\'ecrit ci-dessus et explicit\'e ci-dessous
est appel\'e la {\it trace}.
\end{definition}

\noindent
Soit $f : Y \to X$ un morphisme fini \'etale de sch\'emas. Le morphisme trace est
caract\'eris\'e par les deux propri\'et\'es suivantes :
\begin{enumerate}
\item il commute \`a la localisation \'etale sur $X$, et
\item si $Y = \coprod_{i = 1}^d X$, alors le morphisme trace est
l'application somme $f_*f^{-1} \mathcal{F} = \mathcal{F}^{\oplus d} \to \mathcal{F}$.
\end{enumerate}
D'apr\`es Morphismes étales, lemme \ref{etale-lemma-finite-etale-etale-local},
tout morphisme fini \'etale $f : Y \to X$ est, localement pour la topologie \'etale sur $X$,
de la forme donn\'ee en (2) pour un entier $d \geq 0$. Nous pouvons donc
d\'efinir le morphisme trace par cette caract\'erisation; en particulier,
nous n'avons besoin ni de l'existence de $f_!$, ni de l'\'egalit\'e
de $f_!$ et $f_*$ pour construire le morphisme trace.
Cette description montre que, si $f$ est de degr\'e constant $d$, alors
la compos\'ee
$$
\mathcal{F} \xrightarrow{res}
f_* f^{-1} \mathcal{F} \xrightarrow{trace}
\mathcal{F}
$$
est la multiplication par $d$. La « m\'ethode de la trace »
est l'observation suivante : si $\mathcal{F}$
est un faisceau ab\'elien sur $X_\etale$ tel que la multiplication par $d$
sur $\mathcal{F}$ soit un isomorphisme, alors l'application
$$
H^n_\etale(X, \mathcal{F}) \longrightarrow H^n_\etale(Y, f^{-1}\mathcal{F})
$$
est injective. En effet, nous avons
$$
H^n_\etale(Y, f^{-1}\mathcal{F}) = H^n_\etale(X, f_*f^{-1}\mathcal{F})
$$
par l'annulation des images directes sup\'erieures
(proposition \ref{etale-cohomology-proposition-finite-higher-direct-image-zero})
et la suite spectrale de Leray
(proposition \ref{etale-cohomology-proposition-leray}).
Nous pouvons donc consid\'erer les applications
$$
H^n_\etale(X, \mathcal{F}) \to
H^n_\etale(Y, f^{-1}\mathcal{F})= H^n_\etale(X, f_*f^{-1}\mathcal{F})
\xrightarrow{trace}
H^n_\etale(X, \mathcal{F})
$$
et leur compos\'ee est un isomorphisme (sous notre hypoth\`ese sur $\mathcal{F}$
et $f$). En particulier, si
$H_\etale^q(Y, f^{-1}\mathcal{F}) = 0$, alors
$H_\etale^q(X, \mathcal{F}) = 0$ aussi.
En effet, la multiplication par $d$ induit un
isomorphisme de $H_\etale^q(X, \mathcal{F})$ qui se factorise par
$H_\etale^q(Y, f^{-1}\mathcal{F})= 0$.

\medskip\noindent
On combine souvent ceci avec le lemme suivant.

\begin{lemma}
\label{etale-cohomology-lemma-pullback-filtered}
Soit $S$ un sch\'ema connexe. Soit $\ell$ un nombre premier. Soit
$\mathcal{F}$ un faisceau localement constant de type fini de
$\mathbf{F}_\ell$-espaces vectoriels sur $S_\etale$.
Alors il existe un morphisme fini \'etale
$f : T \to S$ de degr\'e premier \`a $\ell$ tel que $f^{-1}\mathcal{F}$
admette une filtration finie dont les quotients successifs sont
$\underline{\mathbf{Z}/\ell\mathbf{Z}}_T$.
\end{lemma}

\begin{proof}
Choisissons un point g\'eom\'etrique $\overline{s}$ de $S$.
Par l'\'equivalence du lemme \ref{etale-cohomology-lemma-locally-constant-on-connected},
le faisceau $\mathcal{F}$ correspond \`a un espace vectoriel
$V$ de dimension finie sur $\mathbf{F}_\ell$, muni d'une action continue de
$\pi_1(S, \overline{s})$.
Soit $G \subset \text{Aut}(V)$ l'image de l'homomorphisme
$\rho : \pi_1(S, \overline{s}) \to \text{Aut}(V)$ qui donne l'action.
Remarquons que $G$ est fini.
L'homomorphisme continu surjectif
$\overline{\rho} : \pi_1(S, \overline{s}) \to G$
correspond \`a un objet galoisien $Y \to S$ de
$\textit{F\'Et}_S$, de groupe d'automorphismes $G = \text{Aut}(Y/S)$, voir
Fundamental Groups, section \ref{pione-section-finite-etale-under-galois}.
Soit $H \subset G$ un $\ell$-sous-groupe de Sylow.
Nous affirmons que $T = Y/H \to S$ convient. En effet, soit $\overline{t} \in T$
un point g\'eom\'etrique au-dessus de $\overline{s}$. L'image de
$\pi_1(T, \overline{t}) \to \pi_1(S, \overline{s})$
est $(\overline{\rho})^{-1}(H)$ par fonctorialit\'e
des groupes fondamentaux. Par cons\'equent, l'action de $\pi_1(T, \overline{t})$
sur $V$ correspondant \`a $f^{-1}\mathcal{F}$ se fait par
l'application $\pi_1(T, \overline{t}) \to H$, voir la
remarque \ref{etale-cohomology-remark-functorial-locally-constant-on-connected}. Comme
$H$ est un $\ell$-groupe fini, les constituants irr\'eductibles de la
repr\'esentation $\rho|_{\pi_1(T, \overline{t})}$
sont tous triviaux de rang $1$ (c'est un lemme \`el\'ementaire de
th\'eorie des repr\'esentations des groupes finis).
Par l'\'equivalence du
lemme \ref{etale-cohomology-lemma-locally-constant-on-connected},
cela signifie que $f^{-1}\mathcal{F}$ est une extension successive de
faisceaux constants $\underline{\mathbf{Z}/\ell\mathbf{Z}}_T$.
De plus, le degr\'e de $T = Y/H \to S$ est premier \`a $\ell$,
car il est \'egal \`a l'indice de $H$ dans $G$.
\end{proof}

\begin{lemma}
\label{etale-cohomology-lemma-action-l-group}
Soit $\Lambda$ un anneau noeth\'erien.
Soit $\ell$ un nombre premier et $n \geq 1$.
Soit $H$ un $\ell$-groupe fini.
Soit $M$ un $\Lambda[H]$-module de type fini annul\'e par $\ell^n$.
Alors il existe une filtration finie
$0 = M_0 \subset M_1 \subset \ldots \subset M_t = M$
par des sous-$\Lambda[H]$-modules
telle que $H$ agisse trivialement sur $M_{i + 1}/M_i$ pour tout
$i = 0, \ldots, t - 1$.
\end{lemma}

\begin{proof}
Omis. Indication : montrer que l'id\'eal d'augmentation $\mathfrak m$
de l'anneau non commutatif $\mathbf{Z}/\ell^n\mathbf{Z}[H]$
est nilpotent.
\end{proof}

\begin{lemma}
\label{etale-cohomology-lemma-pullback-filtered-modules}
Soit $S$ un sch\'ema irr\'eductible et g\'eom\'etriquement unibranche.
Soit $\ell$ un nombre premier et $n \geq 1$. Soit $\Lambda$
un anneau noeth\'erien. Soit $\mathcal{F}$ un faisceau localement
constant de type fini de $\Lambda$-modules sur $S_\etale$
qui est annul\'e par $\ell^n$. Alors il existe un
morphisme fini \'etale $f : T \to S$ de degr\'e premier \`a $\ell$
tel que $f^{-1}\mathcal{F}$ admette une filtration finie dont les
quotients successifs sont de la forme $\underline{M}_T$
pour des $\Lambda$-modules de type fini $M$.
\end{lemma}

\begin{proof}
Choisissons un point g\'eom\'etrique $\overline{s}$ de $S$.
Par l'\'equivalence du
lemme \ref{etale-cohomology-lemma-locally-constant-on-connected-geom-unibranch},
le faisceau $\mathcal{F}$ correspond \`a un $\Lambda$-module de type fini
$M$ muni d'une action continue de $\pi_1(S, \overline{s})$.
Soit $G \subset \text{Aut}(M)$ l'image de l'homomorphisme
$\rho : \pi_1(S, \overline{s}) \to \text{Aut}(M)$ qui donne l'action.
Remarquons que $G$ est fini par continuit\'e et parce que $M$ est engendr\'e
par un nombre fini d'\'el\'ements comme $\Lambda$-module (voir la preuve du lemme \ref{etale-cohomology-lemma-locally-constant-on-connected-geom-unibranch}).
L'homomorphisme continu surjectif
$\overline{\rho} : \pi_1(S, \overline{s}) \to G$
correspond \`a un objet galoisien $Y \to S$ de
$\textit{F\'Et}_S$, de groupe d'automorphismes $G = \text{Aut}(Y/S)$, voir
Fundamental Groups, section \ref{pione-section-finite-etale-under-galois}.
Soit $H \subset G$ un $\ell$-sous-groupe de Sylow.
Nous affirmons que $T = Y/H \to S$ convient. En effet, soit $\overline{t} \in T$
un point g\'eom\'etrique au-dessus de $\overline{s}$. L'image de
$\pi_1(T, \overline{t}) \to \pi_1(S, \overline{s})$
est $(\overline{\rho})^{-1}(H)$ par fonctorialit\'e
des groupes fondamentaux. Par cons\'equent, l'action de $\pi_1(T, \overline{t})$
sur $M$ correspondant \`a $f^{-1}\mathcal{F}$ se fait par
l'application $\pi_1(T, \overline{t}) \to H$, voir la
remarque \ref{etale-cohomology-remark-functorial-locally-constant-on-connected}.
Soit $0 = M_0 \subset M_1 \subset \ldots \subset M_t = M$
comme dans le lemme \ref{etale-cohomology-lemma-action-l-group}.
Cela induit une filtration
$0 = \mathcal{F}_0 \subset \mathcal{F}_1 \subset \ldots
\subset \mathcal{F}_t = f^{-1}\mathcal{F}$
telle que les quotients successifs soient constants de
valeur $M_{i + 1}/M_i$.
Enfin, le degr\'e de $T = Y/H \to S$ est premier \`a $\ell$,
car il est \'egal \`a l'indice de $H$ dans $G$.
\end{proof}






\section{Cohomologie galoisienne}
\label{etale-cohomology-section-galois-cohomology}

\noindent
Dans cette section, nous d\'emontrons un r\'esultat de cohomologie galoisienne
(proposition \ref{etale-cohomology-proposition-serre-galois})
au moyen de la cohomologie \'etale et de l'argument de la
section \ref{etale-cohomology-section-trace-method}.
Cela nous permettra de d\'emontrer l'annulation des groupes de cohomologie \'etale sup\'erieure
sur le spectre d'un corps.

\begin{lemma}
\label{etale-cohomology-lemma-nonvanishing-inherited}
Soit $\ell$ un nombre premier et $n$ un entier $> 0$.
Soit $S$ un sch\'ema quasi-compact et quasi-s\'epar\'e.
Soit $X = \lim_{i \in I} X_i$ la limite projective d'un
syst\`eme projectif filtrant de $S$-sch\'emas, chaque $X_i \to S$
\'etant fini \'etale de degr\'e constant premier \`a $\ell$.
Les assertions suivantes sont \`equivalentes :
\begin{enumerate}
\item il existe un faisceau de torsion $\ell$-primaire
$\mathcal{G}$ sur $S$ tel que $H_\etale^n(S, \mathcal{G}) \neq 0$, et
\item il existe un faisceau de torsion $\ell$-primaire $\mathcal{F}$ sur $X$
tel que $H_\etale^n(X, \mathcal{F}) \neq 0$.
\end{enumerate}
En fait, \`a partir de
$\mathcal{G}$ nous pouvons prendre $\mathcal{F} = g^{-1}\mathcal{G}$,
et \`a partir de
$\mathcal{F}$ nous pouvons prendre $\mathcal{G} = g_*\mathcal{F}$.
\end{lemma}

\begin{proof}
Notons $g : X \to S$ et $g_i : X_i \to S$ les morphismes structuraux.
Fixons un faisceau de torsion $\ell$-primaire $\mathcal{G}$ sur $S$
tel que $H^n_\etale(S, \mathcal{G}) \not = 0$.
Le syst\`eme form\'e par $\mathcal{G}_i = g_i^{-1}\mathcal{G}$
satisfait aux conditions du th\'eor\`eme \ref{etale-cohomology-theorem-colimit}
et son faisceau limite inductive est $g^{-1}\mathcal{G}$. Nous en
d\'eduisons que :
$$
\colim_{i\in I} H^n_\etale(X_i, g_i^{-1}\mathcal{G}) =
H^n_\etale(X, g^{-1}\mathcal{G})
$$
Comme les $g_i$ sont des morphismes finis \'etales de degr\'e premier
\`a $\ell$, nous pouvons appliquer la « m\'ethode de la trace » et obtenons que
les applications
$$
H^n_\etale(S, \mathcal{G}) \to H^n_\etale(X_i, g_i^{-1}\mathcal{G})
$$
sont toutes injectives (et compatibles aux applications de transition).
Voir la section \ref{etale-cohomology-section-trace-method}. La limite inductive est donc non nulle, c'est-\`a-dire
$H^n(X,g^{-1}\mathcal{G}) \neq 0$, ce qui donne le r\'esultat voulu avec
$\mathcal{F} = g^{-1}\mathcal{G}$.

\medskip\noindent
R\'eciproquement, soit $\mathcal{F}$ un faisceau de torsion $\ell$-primaire sur $X$
tel que $H^n_\etale(X, \mathcal{F}) \not = 0$. Comme les $g_i$
sont des morphismes finis, leurs images directes sup\'erieures sont nulles
(proposition \ref{etale-cohomology-proposition-finite-higher-direct-image-zero}).
En appliquant alors le lemme \ref{etale-cohomology-lemma-relative-colimit},
nous pouvons conclure qu'il en va de m\^eme pour $g$.
L'annulation des images directes sup\'erieures donne
$H^n_\etale(X, \mathcal{F}) = H^n(S, g_*\mathcal{F}) \neq 0$
par Leray (proposition \ref{etale-cohomology-proposition-leray}),
ce qui donne le r\'esultat voulu avec $\mathcal{G} = g_*\mathcal{F}$.
\end{proof}

\begin{lemma}
\label{etale-cohomology-lemma-reduce-to-l-group}
Soit $\ell$ un nombre premier et $n$ un entier $> 0$.
Soit $K$ un corps, posons $G = \text{Gal}(K^{sep}/K)$, et soit
$H \subset G$ un $\ell$-sous-groupe de Sylow, l'extension de corps
correspondante \'etant not\'ee $L/K$. Alors
$H^n_\etale(\Spec(K), \mathcal{F}) = 0$ pour tout faisceau
de torsion $\ell$-primaire $\mathcal{F}$ si et seulement si
$H^n_\etale(\Spec(L), \underline{\mathbf{Z}/\ell\mathbf{Z}}) = 0$.
\end{lemma}

\begin{proof}
\'Ecrivons $L = \bigcup L_i$ comme r\'eunion de ses sous-extensions finies sur $K$.
Le choix de $H$ implique que $[L_i : K]$ est premier \`a $\ell$.
Ainsi $\Spec(L) = \lim_{i \in I} \Spec(L_i)$ comme dans le
lemme \ref{etale-cohomology-lemma-nonvanishing-inherited}.
Nous pouvons donc remplacer $K$ par $L$ et supposer que
le groupe de Galois absolu $G$ de $K$ est un
pro-$\ell$-groupe profini.

\medskip\noindent
Supposons $H^n(\Spec(K), \underline{\mathbf{Z}/\ell\mathbf{Z}}) = 0$.
Soit $\mathcal{F}$ un faisceau de torsion $\ell$-primaire sur $\Spec(K)_\etale$.
Nous allons montrer que $H^n_\etale(\Spec(K), \mathcal{F}) = 0$.
Par la correspondance du
lemme \ref{etale-cohomology-lemma-equivalence-abelian-sheaves-point},
notre faisceau $\mathcal{F}$ correspond \`a un $G$-module de torsion
$\ell$-primaire $M$. Par continuit\'e, tout ensemble fini d'\'el\'ements $x_1, \ldots, x_m \in M$
est fix\'e par un sous-groupe ouvert $U$.
Soit $M'$ le sous-module engendr\'e par les orbites de $x_1, \ldots, x_m$.
C'est un $\ell$-groupe ab\'elien fini,
car chaque $x_i$ est annul\'e par une puissance de $\ell$
et les orbites sont finies. Comme $M$ est la limite inductive filtrante de
ces sous-modules $M'$, nous voyons que $\mathcal{F}$ est la limite
inductive filtrante des sous-faisceaux correspondants $\mathcal{F}' \subset \mathcal{F}$.
En appliquant le th\'eor\`eme \ref{etale-cohomology-theorem-colimit} \`a cette limite inductive, nous nous ramenons
au cas o\`u $\mathcal{F}$ est un faisceau localement constant fini.

\medskip\noindent
Soit $M$ un $\ell$-groupe ab\'elien fini muni d'une action continue
du pro-$\ell$-groupe profini $G$. Il existe alors une filtration
stable par $G$
$$
0 = M_0 \subset M_1 \subset \ldots \subset M_r = M
$$
telle que $M_{i + 1}/M_i \cong \mathbf{Z}/\ell \mathbf{Z}$, avec
action triviale de $G$ (c'est un lemme \'el\'ementaire de th\'eorie des repr\'esentations
des groupes finis).
Le faisceau correspondant $\mathcal{F}$ admet donc une filtration
$$
0 = \mathcal{F}_0 \subset \mathcal{F}_1 \subset \ldots \subset
\mathcal{F}_r = \mathcal{F}
$$
dont les quotients successifs sont isomorphes \`a
$\underline{\mathbf{Z}/\ell \mathbf{Z}}$.
Nous concluons alors par r\'ecurrence et par la suite exacte longue de
cohomologie.
\end{proof}

\begin{lemma}
\label{etale-cohomology-lemma-reduce-to-l-group-higher}
Soit $\ell$ un nombre premier et $n$ un entier $> 0$.
Soit $K$ un corps, posons $G = \text{Gal}(K^{sep}/K)$, et soit
$H \subset G$ un $\ell$-sous-groupe de Sylow,
l'extension de corps correspondante \'etant not\'ee $L/K$.
Alors $H^q_\etale(\Spec(K),\mathcal{F}) = 0$ pour $q \geq n$ et tout
faisceau de torsion $\ell$-primaire $\mathcal{F}$ si et seulement si
$H^n_\etale(\Spec(L), \underline{\mathbf{Z}/\ell\mathbf{Z}}) = 0$.
\end{lemma}

\begin{proof}
L'implication directe est imm\'ediate; il suffit donc de d\'emontrer la r\'eciproque.
Nous raisonnons par r\'ecurrence sur $q$. Le cas $q = n$ est le
lemme \ref{etale-cohomology-lemma-reduce-to-l-group}. Soit maintenant
$\mathcal{F}$ un faisceau de torsion $\ell$-primaire sur $\Spec(K)$.
Soit $f : \Spec(K^{sep}) \rightarrow \Spec(K)$
le morphisme d'inclusion d'un point g\'eom\'etrique.
Consid\'erons alors la suite exacte :
$$
0 \rightarrow \mathcal{F} \xrightarrow{res}
f_* f^{-1} \mathcal{F} \rightarrow f_* f^{-1} \mathcal{F}/\mathcal{F} 
\rightarrow 0
$$
Remarquons que $K^{sep}$ s'\'ecrit comme limite inductive filtrante d'extensions
s\'eparables finies. Ainsi $f$
est la limite projective d'un syst\`eme projectif de morphismes finis \'etales.
Comme nous l'avons vu dans la d\'emonstration du
lemme \ref{etale-cohomology-lemma-nonvanishing-inherited}, nous pouvons conclure que les
images directes sup\'erieures de $f$ sont nulles. La cohomologie sup\'erieure de
$f_* f^{-1} \mathcal{F}$ s'identifie donc \`a celle du point g\'eom\'etrique, qui
est manifestement nulle. Comme tout reste de torsion $\ell$-primaire, nous pouvons appliquer
l'hypoth\`ese de r\'ecurrence conjointement avec la suite exacte longue de cohomologie
pour conclure au rang $q + 1$.
\end{proof}

\begin{proposition}
\label{etale-cohomology-proposition-serre-galois}
\begin{reference}
\cite[Chapter II, Section 3, Proposition 5]{SerreGaloisCohomology}
\end{reference}
Soit $K$ un corps muni d'une clôture séparable $K^{sep}$.
Supposons que, pour toute extension finie $K'$ de $K$, on ait
$\text{Br}(K') = 0$. Alors
\begin{enumerate}
\item $H^q(\text{Gal}(K^{sep}/K), (K^{sep})^*) = 0$
pour tout $q \geq 1$, et
\item $H^q(\text{Gal}(K^{sep}/K), M) = 0$
pour tout $\text{Gal}(K^{sep}/K)$-module de torsion $M$ et tout $q \geq 2$,
\end{enumerate}
\end{proposition}

\begin{proof}
Posons $p = \text{char}(K)$.
D'après le lemme \ref{etale-cohomology-lemma-compare-cohomology-point},
le théorème \ref{etale-cohomology-theorem-brauer-delta}
et l'exemple \ref{etale-cohomology-example-sheaves-point},
la proposition revient à montrer que, si
$H^2(\Spec(K'),\mathbf{G}_m|_{\Spec(K')_\etale}) = 0$
pour toute extension finie $K'/K$, alors :
\begin{itemize}
\item $H^q(\Spec(K),\mathbf{G}_m|_{\Spec(K)_\etale}) = 0$
pour tout $q \geq 1$, et
\item $H^q(\Spec(K),\mathcal{F}) = 0$
pour tout faisceau de torsion $\mathcal{F}$ et tout $q \geq 2$.
\end{itemize}
Commençons par démontrer le second point.
Comme $\mathcal{F}$ est un faisceau de torsion, la décomposition en
composantes $\ell$-primaires et la compatibilité de la cohomologie aux
limites inductives (ici, aux sommes directes; voir le théorème \ref{etale-cohomology-theorem-colimit})
nous ramènent à établir $H^q(\Spec(K),\mathcal{F}) = 0$, $q \geq 2$,
pour tout faisceau de torsion $\ell$-primaire et tout nombre premier $\ell$. Nous
pouvons ainsi étudier séparément chaque nombre premier.

\medskip\noindent
Supposons que $\ell \neq p$. Pour toute extension $K'/K$,
considérons la suite exacte de Kummer (lemme \ref{etale-cohomology-lemma-kummer-sequence})
$$
0 \to
\mu_{\ell, \Spec{K'}} \to
\mathbf{G}_{m, \Spec{K'}} \xrightarrow{(\cdot)^{\ell}}
\mathbf{G}_{m, \Spec{K'}} \to 0
$$
L'hypothèse donne $H^q(\Spec{K'},\mathbf{G}_m|_{\Spec(K')_\etale}) = 0$
pour $q = 2$; pour $q = 1$, la même annulation résulte du
théorème \ref{etale-cohomology-theorem-picard-group} conjointement à
$\Pic(K') = (0)$. La suite exacte longue de cohomologie
permet donc de conclure que $H^2(\Spec{K'}, \mu_\ell) = 0$ pour
toute extension finie $K'/K$. Soit maintenant $H$ un $\ell$-sous-groupe de Sylow
du groupe de Galois absolu de $K$, et soit $L$ l'extension
correspondante. On peut écrire $L$ comme limite inductive d'extensions finies;
en appliquant le théorème \ref{etale-cohomology-theorem-colimit} à cette limite, on voit que
$H^2(\Spec(L), \mu_\ell) = 0$.
Le faisceau $\mu_\ell$ est alors nécessairement constant. S'il ne l'était pas, il
existerait une extension galoisienne de degré premier à
$\ell$ de $L$, contrairement à la définition de $L$ (à savoir l'extension
obtenue en adjoignant les racines $\ell$-ièmes de l'unité à $L$).
Le lemme \ref{etale-cohomology-lemma-reduce-to-l-group-higher}
donne donc le résultat pour $\ell \neq p$.

\medskip\noindent
Supposons maintenant que $\ell = p$. Considérons la
suite exacte d'Artin-Schreier (section \ref{etale-cohomology-section-artin-schreier})
$$
0 \longrightarrow \underline{\mathbf{Z}/p\mathbf{Z}}_{\Spec{K}} \longrightarrow
\mathbf{G}_{a, \Spec{K}} \xrightarrow{F-1} \mathbf{G}_{a, \Spec{K}} 
\longrightarrow 0
$$
où $F - 1$ est le morphisme $x \mapsto x^p - x$. La cohomologie supérieure
de $\mathbf{G}_{a, \Spec{K}}$ est nulle, d'après la remarque
\ref{etale-cohomology-remark-special-case-fpqc-cohomology-quasi-coherent} et l'annulation de la
cohomologie supérieure du faisceau structural d'un schéma affine
(Cohomologie des schémas, lemme
\ref{coherent-lemma-quasi-coherent-affine-cohomology-zero}).
Cette observation vaut sur tout corps de
caractéristique $p$. En particulier, elle s'applique à l'extension $L$
correspondant à un $p$-sous-groupe de Sylow $H$. On en déduit
$H^n(\Spec{L}, \underline{\mathbf{Z}/p\mathbf{Z}}_{\Spec{L}}) = 0$
pour $n \geq 2$, d'où le résultat pour $\ell = p$, par le
lemme \ref{etale-cohomology-lemma-reduce-to-l-group-higher}.

\medskip\noindent
Pour achever la démonstration, il reste à montrer que
$H^q(\text{Gal}(K^{sep}/K), (K^{sep})^*) = 0$ pour tout $q \geq 1$.
Posons $G = \text{Gal}(K^{sep}/K)$ et $M = (K^{sep})^*$,
considéré comme un $G$-module. Nous avons déjà montré ci-dessus que
$H^1(G, M) = 0$ et $H^2(G, M) = 0$. Considérons la suite exacte
$$
0 \to A \to M \to M \otimes \mathbf{Q} \to B \to 0
$$
de $G$-modules. Ce qui précède donne $H^i(G, A) = 0$
et $H^i(G, B) = 0$ pour $i > 1$, car $A$ et $B$ sont des
$G$-modules de torsion. D'après le
lemme \ref{etale-cohomology-lemma-profinite-group-cohomology-torsion},
on a $H^i(G, M \otimes \mathbf{Q}) = 0$ pour $i > 0$.
Il est facile de vérifier que cela entraîne également
$H^i(G, M) = 0$ pour $i \geq 3$.
\end{proof}

%10.08.09

\begin{definition}
\label{etale-cohomology-definition-Cr}
Un corps $K$ est dit {\it $C_r$}
si, pour tout $0 < d^r < n$ et tout $f \in K[T_1,
\ldots, T_n]$ homogène de degré $d$, il existe $\alpha = (\alpha_1,
\ldots, \alpha_n)$, les $\alpha_i \in K$ n'étant pas tous nuls, tel que $f(\alpha) = 0$.
Un tel $\alpha$ est appelé une {\it solution non triviale} de $f$.
\end{definition}

\begin{example}
\label{etale-cohomology-example-algebraically-closed-field-Cr}
Tout corps algébriquement clos est $C_r$.
\end{example}

\noindent
En fait, on a le lemme élémentaire suivant.

\begin{lemma}
\label{etale-cohomology-lemma-algebraically-closed-find-solutions}
Soit $k$ un corps algébriquement clos. Soient
$f_1, \ldots, f_s \in k[T_1, \ldots, T_n]$
des polynômes homogènes de degrés $d_1, \ldots, d_s$, avec $d_i
> 0$. Si $s < n$, les équations $f_1 = \ldots = f_s = 0$ ont une solution
commune non triviale.
\end{lemma}

\begin{proof}
Cela résulte de la théorie de la dimension, par exemple sous la forme du
lemme \ref{varieties-lemma-intersection-in-affine-space} de Variétés,
appliqué $s - 1$ fois.
\end{proof}

\noindent
Le résultat suivant calcule le groupe de Brauer des corps $C_1$.

\begin{theorem}
\label{etale-cohomology-theorem-C1-brauer-group-zero}
Soit $K$ un corps $C_1$. Alors $\text{Br}(K) = 0$.
\end{theorem}

\begin{proof}
Soit $D$ un corps gauche de dimension finie sur $K$, de centre $K$. Nous
avons vu que
$$
D \otimes_K K^{sep} \cong \text{Mat}_d(K^{sep})
$$
à automorphisme intérieur près. Le déterminant $\det :
\text{Mat}_d(K^{sep}) \to K^{sep}$ est donc invariant par Galois et se descend en une
application homogène de degré $d$
$$
\det = N_\text{red} : D \longrightarrow K
$$
appelée la {\it norme réduite}. Comme $K$ est $C_1$, si $d > 1$, il existe
un élément non nul $x \in D$ tel que $N_\text{red}(x) = 0$. Cela implique clairement
que $x$ n'est pas inversible, contradiction. Par conséquent, $\text{Br}(K) = 0$.
\end{proof}

\begin{definition}
\label{etale-cohomology-definition-variety}
Soit $k$ un corps. Une {\it variété} est un schéma intègre, séparé et de
type fini sur $k$. Une {\it courbe} est une variété de dimension $1$.
\end{definition}

\begin{theorem}[Théorème de Tsen]
\label{etale-cohomology-theorem-tsen}
Le corps des fonctions d'une variété de dimension $r$ sur un corps algébriquement clos
$k$ est $C_r$.
\end{theorem}

\begin{proof}
Pour l'espace projectif, on peut montrer directement que le corps
$k(x_1, \ldots, x_r)$ est $C_r$ (exercice).

\medskip\noindent
Cas général. Sans perte de généralité, on peut supposer $X$ projectif.
Soit $f \in k(X)[T_1, \ldots, T_n]_d$ tel que $0 < d^r < n$.
Supposons que les coefficients de $f$ appartiennent à $\Gamma(X, \mathcal{O}_X(H))$
pour un diviseur ample $H \subset X$. Posons
$\mathbf{\alpha} = (\alpha_1, \ldots, \alpha_n)$, où $\alpha_i \in \Gamma(X,
\mathcal{O}_X(eH))$. Alors $f(\mathbf{\alpha}) \in \Gamma(X,
\mathcal{O}_X((de + 1)H))$. Considérons le système d'équations $f(\mathbf{\alpha})
=0$. D'après Riemann--Roch asymptotique
(Variétés, proposition \ref{varieties-proposition-asymptotic-riemann-roch}),
il existe un $c > 0$ tel que
\begin{itemize}
\item le nombre de variables est
$n\dim_k \Gamma(X, \mathcal{O}_X(eH)) \sim n e^r c$, et
\item le nombre d'équations est
$\dim_k \Gamma(X, \mathcal{O}_X((de + 1)H)) \sim (de + 1)^r c$.
\end{itemize}
Comme $n > d^r$, il y a plus de variables que d'équations. Celles-ci sont
homogènes; il existe donc une solution d'après le
lemme \ref{etale-cohomology-lemma-algebraically-closed-find-solutions}.
\end{proof}

\begin{lemma}
\label{etale-cohomology-lemma-curve-brauer-zero}
Soit $C$ une courbe sur un corps algébriquement clos $k$. Alors
le groupe de Brauer du corps des fonctions de $C$ est nul :
$\text{Br}(k(C)) = 0$.
\end{lemma}

\begin{proof}
Cela résulte immédiatement du théorème de Tsen
\ref{etale-cohomology-theorem-tsen} et du
théorème \ref{etale-cohomology-theorem-C1-brauer-group-zero}.
\end{proof}

\begin{lemma}
\label{etale-cohomology-lemma-cohomology-Gm-function-field-curve}
Soient $k$ un corps algébriquement clos et $K/k$ une extension de corps
de degré de transcendance 1. Alors, pour tout $q \geq 1$,
$H_\etale^q(\Spec(K), \mathbf{G}_m) = 0$.
\end{lemma}

\begin{proof}
Rappelons que
$H_\etale^q(\Spec(K), \mathbf{G}_m) = H^q(\text{Gal}(K^{sep}/K), (K^{sep})^*)$
d'après le lemme \ref{etale-cohomology-lemma-compare-cohomology-point}.
La proposition \ref{etale-cohomology-proposition-serre-galois} montre donc
qu'il suffit d'établir que, si $K'/K$ est une extension finie de corps, alors
$\text{Br}(K') = 0$. Or $K' = \colim K''$, où $K''$ parcourt
les extensions de type fini de $k$ contenues comme sous-corps dans $K'$ et de
degré de transcendance $1$.
On a $\text{Br}(K') = \colim \text{Br}(K'')$, ce qui nous ramène
à une extension de corps de type fini $K''/k$ de degré de transcendance
$1$. Un tel corps est le corps des fonctions d'une courbe sur $k$,
dont le groupe de Brauer est donc nul d'après le
lemme \ref{etale-cohomology-lemma-curve-brauer-zero}.
\end{proof}






\section{Annulation en degrés supérieurs pour le groupe multiplicatif}
\sectionmark{Annulation pour le groupe multiplicatif}
\label{etale-cohomology-section-higher-Gm}

\noindent
Dans cette section, on fixe un corps algébriquement clos $k$ et une courbe lisse
$X$ sur $k$. On note $i_x : x \hookrightarrow X$ l'inclusion d'un point fermé
de $X$ et $j : \eta \hookrightarrow X$ l'inclusion du point générique.
On note aussi $X_0$ l'ensemble des points fermés de $X$.

\begin{theorem}[La suite exacte fondamentale]
\label{etale-cohomology-theorem-fundamental-exact-sequence}
Il existe une suite exacte courte de faisceaux étales sur $X$
$$
0 \longrightarrow
\mathbf{G}_{m, X} \longrightarrow
j_* \mathbf{G}_{m, \eta} \longrightarrow
\bigoplus\nolimits_{x \in X_0} {i_x}_* \underline{\mathbf{Z}}
\longrightarrow 0.
$$
\end{theorem}

\begin{proof}
Soit $\varphi : U \to X$ un morphisme étale. D'après les propriétés des
morphismes étales (proposition \ref{etale-cohomology-proposition-etale-morphisms}),
$U = \coprod_i U_i$, où chaque $U_i$ est une courbe lisse munie d'un morphisme vers $X$.
La suite ci-dessus pour $U$ est le produit des suites correspondantes
pour les $U_i$; il suffit donc de traiter le cas où $U$ est connexe,
donc irréductible. Dans ce cas, on a la suite exacte bien connue
$$
1 \longrightarrow
\Gamma(U, \mathcal{O}_U^*) \longrightarrow
k(U)^* \longrightarrow \bigoplus\nolimits_{y \in U_0} \mathbf{Z}_y.
$$
Cela revient à une suite
$$
0 \longrightarrow \Gamma(U, \mathcal{O}_U^*) \longrightarrow
\Gamma(\eta \times_X U, \mathcal{O}_{\eta \times_X U}^*) \longrightarrow
\bigoplus\nolimits_{x \in X_0}
\Gamma(x \times_X U, \underline{\mathbf{Z}})
$$
qui, en développant les définitions, n'est autre que la suite
$$
0 \longrightarrow \mathbf{G}_m(U) \longrightarrow
j_* \mathbf{G}_{m, \eta}(U) \longrightarrow
\left(\bigoplus\nolimits_{x \in X_0} {i_x}_* \underline{\mathbf{Z}}\right)(U).
$$
Cette construction définit les applications de la suite exacte fondamentale et montre
qu'elle est exacte, sauf peut-être au dernier terme. Pour voir la surjectivité, rappelons que si
$U$ est une courbe non singulière et $D$ un diviseur sur $U$,
il existe un recouvrement ouvert de Zariski $\{U_j \to U\}$ de $U$
tel que $D|_{U_j} = \text{div}(f_j)$ pour un $f_j \in k(U)^*$.
\end{proof}

\begin{lemma}
\label{etale-cohomology-lemma-higher-direct-jstar-Gm}
Pour tout $q \geq 1$, $R^q j_*\mathbf{G}_{m, \eta} = 0$.
\end{lemma}

\begin{proof}
Il faut montrer que $(R^q j_*\mathbf{G}_{m, \eta})_{\bar x} = 0$ pour tout
point géométrique $\bar x$ de $X$.

\medskip\noindent
Supposons que $\bar x$ soit au-dessus d'un point fermé $x$ de $X$. Soit $\Spec(A)$
un voisinage ouvert affine de $x$ dans $X$, et soit $K$ le corps des fractions
de $A$. Alors
$$
\Spec(\mathcal{O}^{sh}_{X, \bar x}) \times_X \eta =
\Spec(\mathcal{O}^{sh}_{X, \bar x} \otimes_A K).
$$
L'anneau $\mathcal{O}^{sh}_{X, \bar x} \otimes_A K$ est un localisé de
l'anneau de valuation discrète $\mathcal{O}^{sh}_{X, \bar x}$; c'est donc soit
encore $\mathcal{O}^{sh}_{X, \bar x}$, soit son corps des fractions
$K^{sh}_{\bar x}$. Mais comme une uniformisante locale y devient inversible, c'est
nécessairement le second. Par conséquent,
$$
(R^q j_*\mathbf{G}_{m, \eta})_{(X, \bar x)} = H_\etale^q(\Spec
K^{sh}_{\bar x}, \mathbf{G}_m).
$$
Rappelons que $\mathcal{O}^{sh}_{X, \bar x} =
\colim_{(U, \bar u) \to \bar x} \mathcal{O}(U) = \colim_{A \subset B} B$
où $A \to B$ est étale; ainsi $K^{sh}_{\bar x}$ est une extension algébrique
de $K = k(X)$, et on peut appliquer le
lemme \ref{etale-cohomology-lemma-cohomology-Gm-function-field-curve}
pour obtenir l'annulation.

\medskip\noindent
Supposons que $\bar x = \bar \eta$ soit au-dessus du point générique $\eta$ de $X$ (en
fait, ce cas est superflu). Alors
$\mathcal{O}^{sh}_{X, \bar \eta} = \kappa(\eta)^{sep}$ et donc
\begin{eqnarray*}
(R^q j_*\mathbf{G}_{m, \eta})_{\bar \eta}
& = &
H_\etale^q(\Spec(\kappa(\eta)^{sep}) \times_X \eta, \mathbf{G}_m) \\
& = & H_\etale^q (\Spec(\kappa(\eta)^{sep}), \mathbf{G}_m) \\
& = & 0 \ \ \text{ pour } q \geq 1
\end{eqnarray*}
puisque le groupe de Galois correspondant est trivial.
\end{proof}

\begin{lemma}
\label{etale-cohomology-lemma-cohomology-jstar-Gm}
Pour tout $p \geq 1$, $H_\etale^p(X, j_*\mathbf{G}_{m, \eta}) = 0$.
\end{lemma}

\begin{proof}
La suite spectrale de Leray s'écrit
$$
E_2^{p, q} = H_\etale^p(X, R^qj_*\mathbf{G}_{m, \eta}) \Rightarrow
H_\etale^{p+q}(\eta, \mathbf{G}_{m, \eta}),
$$
Le membre de droite est nul pour $p + q \geq 1$ d'après le
lemme \ref{etale-cohomology-lemma-cohomology-Gm-function-field-curve}.
Le membre de gauche est nul pour $q \geq 1$ d'après le
lemme \ref{etale-cohomology-lemma-higher-direct-jstar-Gm}. On obtient donc l'annulation voulue.
\end{proof}

\begin{lemma}
\label{etale-cohomology-lemma-cohomology-istar-Z}
Pour tout $q \geq 1$, $H_\etale^q(X, \bigoplus_{x \in X_0} {i_x}_*
\underline{\mathbf{Z}}) = 0$.
\end{lemma}

\begin{proof}
Comme $X$ est quasi-compact et quasi-séparé, la cohomologie commute aux limites inductives;
il suffit donc de montrer l'annulation de $H_\etale^q(X, {i_x}_*
\underline{\mathbf{Z}})$. Mais l'inclusion $i_x$ d'un point fermé est
finie, donc $R^p {i_x}_* \underline{\mathbf{Z}} = 0$ pour tout $p \geq 1$ d'après la
proposition \ref{etale-cohomology-proposition-finite-higher-direct-image-zero}.
La suite spectrale de Leray donne alors
$H_\etale^q(X, {i_x}_* \underline{\mathbf{Z}}) =
H_\etale^q(x, \underline{\mathbf{Z}})$.
Enfin, comme $x$ est le spectre d'un
corps algébriquement clos, toute cohomologie supérieure de $x$ est nulle.
\end{proof}

\noindent
Cette série de lemmes donne le résultat suivant.

\begin{theorem}
\label{etale-cohomology-theorem-vanishing-cohomology-Gm-curve}
Soit $X$ une courbe lisse sur un corps algébriquement clos. Alors
$$
H_\etale^q(X, \mathbf{G}_m) = 0 \ \ \text{ pour tout } q \geq 2.
$$
\end{theorem}

\begin{proof}
Voir la discussion ci-dessus.
\end{proof}

\noindent
On obtient également la suite exacte longue de cohomologie
$$
0 \to
H_\etale^0(X, \mathbf{G}_m) \to
H_\etale^0(X, j_*\mathbf{G}_{m, \eta}) \to
H_\etale^0(X, \bigoplus\nolimits_{x \in X_0} {i_x}_*\underline{\mathbf{Z}}) \to
H_\etale^1(X, \mathbf{G}_m) \to 0
$$
bien qu'il s'agisse de la suite familière
$$
0 \to H_{Zar}^0(X, \mathcal{O}_X^*) \to k(X)^* \to \text{Div}(X)
\to \Pic(X) \to 0.
$$





\section{Groupes de Picard des courbes}
\label{etale-cohomology-section-pic-curves}

\noindent
L'étape suivante consiste à utiliser la suite de Kummer pour déduire des renseignements
sur le groupe de cohomologie d'une courbe à coefficients finis. Afin
d'obtenir une annulation dans la suite exacte longue, rappelons quelques faits sur les
groupes de Picard.

\medskip\noindent
Soit $X$ une courbe projective lisse sur un corps algébriquement clos $k$.
Soit $g = \dim_k H^1(X, \mathcal{O}_X)$ le genre de $X$.
Il existe une suite exacte courte
$$
0 \to \Pic^0(X) \to \Pic(X) \xrightarrow{\deg} \mathbf{Z} \to 0.
$$
Le groupe abélien $\Pic^0(X)$ s'identifie à
$\Pic^0(X) = \underline{\Picardfunctor}^0_{X/k}(k)$, c'est-à-dire
au groupe des points rationnels sur $k$ d'une variété abélienne
$\underline{\Picardfunctor}^0_{X/k}$ sur $k$ de dimension $g$.
Par conséquent, si $n \in k^*$, alors
$\Pic^0(X)[n] \cong (\mathbf{Z}/n\mathbf{Z})^{2g}$
en tant que groupes abéliens. Voir
Schémas de Picard des courbes, section \ref{pic-section-picard-curve},
et
Groupoïdes, section \ref{groupoids-section-abelian-varieties}.
Ce fait essentiel, à savoir la description de la torsion dans le
groupe de Picard d'une courbe projective lisse sur un corps
algébriquement clos, ne semble pas admettre de démonstration élémentaire.

\begin{lemma}
\label{etale-cohomology-lemma-cohomology-smooth-projective-curve}
Soit $X$ une courbe projective lisse de genre $g$ sur un
corps algébriquement clos $k$, et soit $n \geq 1$ inversible dans $k$.
Alors on a les identifications canoniques
$$
H_\etale^q(X, \mu_n) =
\left\{
\begin{matrix}
\mu_n(k) & \text{ si }q = 0, \\
\Pic^0(X)[n] & \text{ si }q = 1, \\
\mathbf{Z}/n\mathbf{Z} & \text{ si }q = 2, \\
0 & \text{ si }q \geq 3.
\end{matrix}
\right.
$$
Comme $\mu_n \cong \underline{\mathbf{Z}/n\mathbf{Z}}$, on en déduit des
identifications (non canoniques)
$$
H_\etale^q(X, \underline{\mathbf{Z}/n\mathbf{Z}}) \cong
\left\{
\begin{matrix}
\mathbf{Z}/n\mathbf{Z} & \text{ si }q = 0, \\
(\mathbf{Z}/n\mathbf{Z})^{2g} & \text{ si }q = 1, \\
\mathbf{Z}/n\mathbf{Z} & \text{ si }q = 2, \\
0 & \text{ si }q \geq 3.
\end{matrix}
\right.
$$
\end{lemma}

\begin{proof}
Les théorèmes \ref{etale-cohomology-theorem-picard-group} et
\ref{etale-cohomology-theorem-vanishing-cohomology-Gm-curve}
déterminent la cohomologie étale de $\mathbf{G}_m$ sur $X$
en fonction du groupe de Picard de $X$.
La suite exacte de Kummer $0\to \mu_{n, X} \to \mathbf{G}_{m, X}
\to \mathbf{G}_{m, X}\to 0$ (lemme \ref{etale-cohomology-lemma-kummer-sequence})
donne alors la suite exacte longue de cohomologie
$$
\xymatrix{
0 \ar[r] & \mu_n(k) \ar[r] &
k^* \ar[r]^{(\cdot)^n} &
k^* \ar@(rd, ul)[rdllllr] \\
& H_\etale^1(X, \mu_n) \ar[r] &
\Pic(X) \ar[r]^{(\cdot)^n} &
\Pic(X) \ar@(rd, ul)[rdllllr] \\
& H_\etale^2(X, \mu_n) \ar[r] & 0 \ar[r] & 0 \ldots
}
$$
L'application d'élévation à la puissance $n$, $k^* \to k^*$, est surjective puisque $k$ est
algébriquement clos. Il reste donc à calculer le noyau et le conoyau de l'application
$n : \Pic(X) \to \Pic(X)$. Considérons le diagramme commutatif à
lignes exactes
$$
\xymatrix{
0 \ar[r] &
\Pic^0(X) \ar[r] \ar@{>>}[d]^{(\cdot)^n} &
\Pic(X) \ar[r]_-\deg \ar[d]^{(\cdot)^n} &
\mathbf{Z} \ar[r] \ar@{^{(}->}[d]^n & 0 \\
0 \ar[r] &
\Pic^0(X) \ar[r] &
\Pic(X) \ar[r]^-\deg &
\mathbf{Z} \ar[r] & 0
}
$$
Le groupe $\Pic^0(X)$ est le groupe des points rationnels sur $k$ du
schéma en groupes $\underline{\Picardfunctor}^0_{X/k}$; voir
Schémas de Picard des courbes, lemme \ref{pic-lemma-picard-pieces}.
Le même lemme affirme que $\underline{\Picardfunctor}^0_{X/k}$
est une variété abélienne de dimension $g$ sur $k$ au sens de
Groupoïdes, définition \ref{groupoids-definition-abelian-variety}.
L'application verticale de gauche est donc surjective d'après
Groupoïdes, proposition \ref{groupoids-proposition-review-abelian-varieties}.
Le lemme du serpent donne alors les identifications canoniques énoncées
dans le lemme.

\medskip\noindent
Pour obtenir les identifications non canoniques du lemme, il faut
montrer que le noyau de $n : \Pic^0(X) \to \Pic^0(X)$
est isomorphe à $(\mathbf{Z}/n\mathbf{Z})^{\oplus 2g}$.
Cela fait aussi partie de Groupoïdes, proposition
\ref{groupoids-proposition-review-abelian-varieties}.
\end{proof}

\begin{lemma}
\label{etale-cohomology-lemma-pullback-on-h2-curve}
Soit $\pi : X \to Y$ un morphisme non constant de courbes projectives lisses
sur un corps algébriquement clos $k$, et soit $n \geq 1$ inversible dans $k$.
L'application
$$
\pi^* : H^2_\etale(Y, \mu_n) \longrightarrow H^2_\etale(X, \mu_n)
$$
est la multiplication par le degré de $\pi$.
\end{lemma}

\begin{proof}
L'énoncé a un sens puisque nous avons identifié les deux
groupes de cohomologie $ H^2_\etale(Y, \mu_n)$ et $H^2_\etale(X, \mu_n)$
à $\mathbf{Z}/n\mathbf{Z}$ dans le
lemme \ref{etale-cohomology-lemma-cohomology-smooth-projective-curve}.
En effet, si $\mathcal{L}$ est un faisceau inversible de degré $1$
sur $Y$, de classe $[\mathcal{L}] \in H^1_\etale(Y, \mathbf{G}_m)$,
alors l'image de $[\mathcal{L}]$ par l'homomorphisme bord est un générateur de
$H^2_\etale(Y, \mu_n)$. Il s'agit ici de l'homomorphisme bord
de la suite exacte longue de cohomologie associée à la suite exacte
de Kummer. Le résultat du lemme découle donc du fait que
le degré du faisceau inversible $\pi^*\mathcal{L}$ sur $X$ vaut $\deg(\pi)$.
Nous omettons quelques détails.
\end{proof}

\begin{lemma}
\label{etale-cohomology-lemma-vanishing-cohomology-mu-smooth-curve}
Soit $X$ une courbe affine lisse sur un corps algébriquement clos $k$, et
soit $n \in k^*$. Soit $X \subset \overline{X}$ une
compactification projective lisse
(Variétés, remarque \ref{varieties-remark-smooth-projective-compactification}).
Soit $g$ le genre de $\overline{X}$ et soit $r$ le nombre de
points de $\overline{X} \setminus X$. Alors
\begin{enumerate}
\item $H_\etale^0(X, \mu_n) = \mu_n(k)$;
\item $H_\etale^1(X, \mu_n) \cong (\mathbf{Z}/n\mathbf{Z})^{2g+r-1}$; et
\item $H_\etale^q(X, \mu_n) = 0$ pour tout $q \geq 2$.
\end{enumerate}
\end{lemma}

\begin{proof}
\'Ecrivons $X = \overline{X} - \{ x_1, \ldots, x_r\}$. Alors
$\Pic(X) = \Pic(\overline{X})/ R$, où $R$ est le sous-groupe engendré par
$\mathcal{O}_{\overline{X}}(x_i)$, $1 \leq i \leq r$.
Comme $r \geq 1$, on voit que $\Pic^0(\overline{X}) \to \Pic(X)$
est surjective; le groupe $\Pic(X)$ est donc divisible (voir la discussion dans la preuve
du lemme \ref{etale-cohomology-lemma-cohomology-smooth-projective-curve}).
La suite exacte de Kummer donne (1) et (3). Pour (2), on utilise
\begin{align*}
H_\etale^1(X, \mu_n)
& =
\{(\mathcal L, \alpha)\}/\cong \\
& =
(\{(\bar{\mathcal L}, D, \bar \alpha)\}/\cong)/\tilde{R}
\end{align*}
Pour le premier signe d'égalité et les notations, voir le lemme \ref{etale-cohomology-lemma-describe-h1-mun}.
Les triplets $(\bar{\mathcal L}, D, \bar \alpha)$ de la seconde ligne
sont formés d'un module inversible $\bar{\mathcal L}$ sur $\overline{X}$
de degré $0$, d'un diviseur $D$ sur $\overline{X}$ de degré $0$ supporté par
$\left\{x_1, \ldots, x_r\right\}$ et d'un isomorphisme
$\bar{\alpha}: \bar{\mathcal L}^{\otimes n} \cong \mathcal{O}_{\bar{X}}(D)$.
Les isomorphismes de triplets sont définis comme dans la discussion qui précède le
lemme \ref{etale-cohomology-lemma-describe-h1-mun}.
De plus, $\tilde{R}$ est le sous-groupe engendré par les triplets de la forme
$(\mathcal{O}_{\overline{X}}(D'), n D', 1^{\otimes n})$
où $D'$ est supporté par $\left\{x_1, \ldots, x_r\right\}$ et de degré 0.
La discussion ci-dessus montre que le second signe d'égalité est donné par
l'application qui envoie le triplet $(\bar{\mathcal L},\ D,\ \bar \alpha)$
sur la paire $(\bar{\mathcal L}|_X, \bar \alpha|_X)$. On obtient donc
une suite exacte
$$
0 \longrightarrow
H_\etale^1(\overline{X}, \mu_n) \longrightarrow
H_\etale^1(X, \mu_n) \longrightarrow
\bigoplus_{i = 1}^r \mathbf{Z}/n\mathbf{Z}
\xrightarrow{\ \sum\ }
\mathbf{Z}/n\mathbf{Z} \longrightarrow 0
$$
où l'application du milieu envoie la classe d'un triplet
$(\bar{ \mathcal L}, D, \bar \alpha)$ avec $D = \sum_{i = 1}^r a_i (x_i)$
sur le $r$-uplet $(a_i)_{i = 1}^r$. Il suffit alors d'utiliser le
lemme \ref{etale-cohomology-lemma-cohomology-smooth-projective-curve}
pour compter les rangs.
\end{proof}

\begin{remark}
\label{etale-cohomology-remark-natural-proof}
La manière ``naturelle'' de démontrer le lemme précédent consiste à appliquer l'excision à l'ouvert $X$ dans $\bar
X$. Cela est possible, mais nous n'avons pas développé cette théorie.
\end{remark}

\begin{remark}
\label{etale-cohomology-remark-normalize-H1-Gm}
Soit $k$ un corps algébriquement clos. Soit $n$ un entier premier à
la caractéristique de $k$. Rappelons que
$$
\mathbf{G}_{m, k} = \mathbf{A}^1_k \setminus \{0\} =
\mathbf{P}^1_k \setminus \{0, \infty\}
$$
Nous affirmons qu'il existe un isomorphisme canonique
$$
H^1_\etale(\mathbf{G}_{m, k}, \mu_n) = \mathbf{Z}/n\mathbf{Z}
$$
Que signifie cette affirmation? Elle signifie qu'il existe un élément $1_k$ de
$H^1_\etale(\mathbf{G}_{m, k}, \mu_n)$ tel que, pour
tout morphisme $\Spec(k') \to \Spec(k)$, l'application d'image inverse en
cohomologie étale associée à $\mathbf{G}_{m, k'} \to \mathbf{G}_{m, k}$
envoie $1_k$ sur $1_{k'}$. (En particulier, cet élément est
fixé par tout automorphisme de $k$.) Pour le voir, considérons le
torseur sous $\mu_{n, \mathbf{Z}}$
$\mathbf{G}_{m, \mathbf{Z}} \to \mathbf{G}_{m, \mathbf{Z}}$,
$x \mapsto x^n$. Par l'identification des classes de torseurs avec la
cohomologie de degré un, son image inverse fournit nos éléments canoniques $1_k$.
En revenant au faisceau constant, on obtient des identifications canoniques
$$
H^1_\etale(\mathbf{G}_{m, k}, \mathbf{Z}/n\mathbf{Z}) =
\Hom(\mu_n(k), \mathbf{Z}/n\mathbf{Z}),
$$
c'est-à-dire que ces isomorphismes sont compatibles avec les morphismes de
corps algébriquement clos, en particulier aux
automorphismes de $k$.
\end{remark}










\section{Prolongement par zéro}
\label{etale-cohomology-section-extension-by-zero}

\noindent
Les résultats généraux de
Modules sur les sites, section \ref{sites-modules-section-localize},
permettent de poser la définition suivante.

\begin{definition}
\label{etale-cohomology-definition-extension-zero}
Soit $j : U \to X$ un morphisme étale de schémas.
\begin{enumerate}
\item Le foncteur de restriction
$j^{-1} : \Sh(X_\etale) \to \Sh(U_\etale)$
admet un adjoint à gauche
$j_!^{Sh} : \Sh(U_\etale) \to \Sh(X_\etale)$.
\item Le foncteur de restriction
$j^{-1} : \textit{Ab}(X_\etale) \to \textit{Ab}(U_\etale)$
admet un adjoint à gauche, noté
$j_! : \textit{Ab}(U_\etale) \to \textit{Ab}(X_\etale)$
et appelé {\it prolongement par zéro}.
\item Soit $\Lambda$ un anneau. Le foncteur de restriction
$j^{-1} : \textit{Mod}(X_\etale, \Lambda) \to
\textit{Mod}(U_\etale, \Lambda)$
admet un adjoint à gauche, noté
$j_! : \textit{Mod}(U_\etale, \Lambda) \to
\textit{Mod}(X_\etale, \Lambda)$
et appelé {\it prolongement par zéro}.
\end{enumerate}
\end{definition}

\noindent
Si $\mathcal{F}$ est un faisceau abélien sur $U_\etale$, alors
$j_!\mathcal{F} \not = j_!^{Sh}\mathcal{F}$ en général. En revanche,
le $j_!$ des faisceaux de $\Lambda$-modules coïncide avec le $j_!$ des
faisceaux abéliens sous-jacents
(Modules sur les sites, remarque \ref{sites-modules-remark-localize-shriek-equal}).
Le foncteur $j_!$ est caractérisé par l'isomorphisme fonctoriel
$$
\Hom_X(j_!\mathcal{F}, \mathcal{G}) = \Hom_U(\mathcal{F}, j^{-1}\mathcal{G})
$$
pour tout $\mathcal{F} \in \textit{Ab}(U_\etale)$ et
$\mathcal{G} \in \textit{Ab}(X_\etale)$. Il en va de même pour les
faisceaux de $\Lambda$-modules.

\medskip\noindent
Pour décrire plus explicitement les foncteurs de la définition \ref{etale-cohomology-definition-extension-zero},
rappelons que $j^{-1}$ n'est autre que la restriction
le long du foncteur $U_\etale \to X_\etale$. Autrement dit,
$j^{-1}\mathcal{G}(U') = \mathcal{G}(U')$ pour $U'$ étale sur $U$.
Pour $\mathcal{F} \in \textit{Ab}(U_\etale)$, considérons le
préfaisceau
\begin{equation}
\label{etale-cohomology-equation-j-p-shriek}
j_{p!}\mathcal{F} : X_\etale \longrightarrow \textit{Ab},
\quad
V \longmapsto \bigoplus\nolimits_{V \to U} \mathcal{F}(V \to U)
\end{equation}
Alors $j_!\mathcal{F}$ est le faisceau associé au préfaisceau $j_{p!}\mathcal{F}$.
C'est démontré dans Modules sur les sites, lemme
\ref{sites-modules-lemma-extension-by-zero};
plus généralement, voir la discussion dans
Modules sur les sites, sections \ref{sites-modules-section-localize} et
\ref{sites-modules-section-exactness-lower-shriek}.

\begin{exercise}
\label{etale-cohomology-exercise-jshriek-direct}
Montrer directement que le foncteur $j_!$, défini comme le faisceau associé au
préfaisceau $j_{p!}$ donné dans
(\ref{etale-cohomology-equation-j-p-shriek}), est adjoint à gauche de $j^{-1}$.
\end{exercise}

\begin{proposition}
\label{etale-cohomology-proposition-describe-jshriek}
Soit $j : U \to X$ un morphisme étale de schémas.
Soit $\mathcal{F}$ un objet de $\textit{Ab}(U_\etale)$.
Si $\overline{x} : \Spec(k) \to X$ est un point géométrique de $X$, alors
$$
(j_!\mathcal{F})_{\overline{x}} =
\bigoplus\nolimits_{\overline{u} : \Spec(k) \to U,\ j(\overline{u}) =
\overline{x}} \mathcal{F}_{\bar{u}}.
$$
En particulier, $j_!$ est un foncteur exact.
\end{proposition}

\begin{proof}
L'exactitude de $j_!$ est très générale; voir Modules sur les sites,
lemme \ref{sites-modules-lemma-extension-by-zero-exact}.
Elle résulte bien sûr aussi de la description des fibres.
La formule de la fibre résulte de
Modules sur les sites, lemme \ref{sites-modules-lemma-stalk-j-shriek},
et de la description des points du petit site étale
en termes de points géométriques; voir le lemme \ref{etale-cohomology-lemma-points-small-etale-site}.

\medskip\noindent
Pour un usage ultérieur, notons que l'isomorphisme
\begin{align*}
(j_!\mathcal{F})_{\overline{x}}
& =
(j_{p!}\mathcal{F})_{\overline{x}} \\
& =
\colim_{(V, \overline{v})} j_{p!}\mathcal{F}(V) \\
& =
\colim_{(V, \overline{v})}
\bigoplus\nolimits_{\varphi : V \to U}
\mathcal{F}(V \xrightarrow{\varphi} U) \\
& \to
\bigoplus\nolimits_{\overline{u} : \Spec(k) \to U,\ j(\overline{u}) =
\overline{x}} \mathcal{F}_{\bar{u}}.
\end{align*}
construit dans Modules sur les sites, lemme \ref{sites-modules-lemma-stalk-j-shriek},
envoie $(V, \overline{v}, \varphi, s)$ sur la classe de $s$ dans la fibre
de $\mathcal{F}$ en $\overline{u} = \varphi(\overline{v})$.
\end{proof}

\begin{lemma}
\label{etale-cohomology-lemma-jshriek-open}
Soit $j : U \to X$ une immersion ouverte de schémas. Pour tout
faisceau abélien $\mathcal{F}$ sur $U_\etale$, les morphismes d'adjonction
$j^{-1}j_*\mathcal{F} \to \mathcal{F}$ et
$\mathcal{F} \to j^{-1}j_!\mathcal{F}$ sont des isomorphismes.
En fait, $j_!\mathcal{F}$ est l'unique faisceau abélien sur $X_\etale$
dont la restriction à $U$ est $\mathcal{F}$ et dont les fibres aux
points géométriques de $X \setminus U$ sont nulles.
\end{lemma}

\begin{proof}
Nous invitons le lecteur à démontrer la première assertion directement à partir des
définitions; nous utilisons ici seulement le fait qu'elle est un cas particulier du très général
lemme \ref{sites-modules-lemma-restrict-back} de Modules sur les sites.
Pour la seconde assertion, observons que si $\mathcal{G}$ est un faisceau abélien
sur $X_\etale$ dont la restriction à $U$ est $\mathcal{F}$, l'adjonction fournit
un morphisme $j_!\mathcal{F} \to \mathcal{G}$. Ce morphisme
induit alors un isomorphisme sur les fibres aux points géométriques de $U$ d'après la
proposition \ref{etale-cohomology-proposition-describe-jshriek}.
Ainsi, si les fibres de $\mathcal{G}$ aux points géométriques
de $X \setminus U$ sont nulles, alors $j_!\mathcal{F} \to \mathcal{G}$ est un
isomorphisme d'après le théorème \ref{etale-cohomology-theorem-exactness-stalks}.
\end{proof}

\begin{lemma}[Le prolongement par zéro commute au changement de base]
\label{etale-cohomology-lemma-shriek-base-change}
Soit $f: Y \to X$ un morphisme de schémas. Soit $j: V \to X$ un morphisme
étale. Considérons le produit fibré
$$
\xymatrix{
V' = Y \times_X V \ar[d]_{f'} \ar[r]_-{j'} & Y \ar[d]^f \\
V \ar[r]^j & X
}
$$
Alors on a $j'_! f'^{-1} = f^{-1} j_!$ sur les faisceaux abéliens et les
faisceaux de modules.
\end{lemma}

\begin{proof}
Cela résulte du fait que $j'_! f'^{-1}$ est adjoint à gauche de
$f'_* (j')^{-1}$, tandis que $f^{-1} j_!$ est adjoint à gauche de $j^{-1}f_*$.
De plus, $f'_* (j')^{-1} = j^{-1}f_*$, car $f_*$ commute à la
localisation étale (par construction). En fait, le lemme vaut très généralement
dans le cadre d'un morphisme de sites; voir
Modules sur les sites, lemme
\ref{sites-modules-lemma-localize-morphism-ringed-sites}.
\end{proof}

\begin{lemma}
\label{etale-cohomology-lemma-shriek-into-star-separated-etale}
Soit $j : U \to X$ séparé et étale. Il existe alors un morphisme
injectif fonctoriel $j_!\mathcal{F} \to j_*\mathcal{F}$
sur les faisceaux abéliens et les faisceaux de $\Lambda$-modules.
\end{lemma}

\begin{proof}
Démontrons-le dans le cas des faisceaux abéliens. Construisons un morphisme canonique
$$
j_{p!}\mathcal{F} \to j_*\mathcal{F}
$$
de préfaisceaux abéliens sur $X_\etale$, pour tout faisceau abélien
$\mathcal{F}$ sur $U_\etale$, où $j_{p!}$ est celui de
(\ref{etale-cohomology-equation-j-p-shriek}). Le faisceau associé à ce morphisme sera
le morphisme cherché $j_!\mathcal{F} \to j_*\mathcal{F}$.
En évaluant les deux membres sur $V \to X$ étale, on obtient
$$
j_{p!}\mathcal{F}(V) =
\bigoplus\nolimits_{\varphi : V \to U} \mathcal{F}(V \xrightarrow{\varphi} U)
\quad\text{et}\quad
j_*\mathcal{F}(V) = \mathcal{F}(V \times_X U)
$$
Pour chaque $\varphi$, on a une immersion ouverte et fermée
$$
\Gamma_\varphi = (1, \varphi) : V \longrightarrow V \times_X U
$$
au-dessus de $U$. Elle est ouverte, car c'est un morphisme entre schémas étales
sur $U$, et elle est fermée, car c'est une section d'un schéma séparé
sur $V$ (Schémas, lemme \ref{schemes-lemma-section-immersion}).
Ainsi, pour une section $s_\varphi \in \mathcal{F}(V \xrightarrow{\varphi} U)$,
il existe une unique section $s'_\varphi$ de $\mathcal{F}(V \times_X U)$
dont l'image inverse est $s_\varphi$ par $\Gamma_\varphi$
et dont la restriction au complémentaire de l'image de $\Gamma_\varphi$ est nulle.

\medskip\noindent
Pour montrer que notre morphisme est injectif, supposons que
$\sum_{i = 1}^n s_{\varphi_i}$ soit un élément de
$j_{p!}\mathcal{F}(V)$ dont l'image par la formule ci-dessus est nulle
dans $j_*\mathcal{F}(V)$. Il s'agit de montrer que la restriction de
$\sum_{i = 1}^n s_{\varphi_i}$ est nulle
sur les objets d'un recouvrement étale de $V$.
En considérant tous les égalisateurs deux à deux
(qui sont ouverts et fermés dans $V$) des morphismes
$\varphi_i : V \to U$ et en raisonnant localement sur $V$, on
peut supposer que les images des morphismes
$\Gamma_{\varphi_1}, \ldots, \Gamma_{\varphi_n}$
sont deux à deux disjointes. Comme notre hypothèse est que
$\sum_{i = 1}^n s'_{\varphi_i} = 0$,
on en déduit immédiatement que $s'_{\varphi_i} = 0$
pour tout $i$ (par disjonction des supports de ces
sections), d'où $s_{\varphi_i} = 0$ pour tout $i$, comme voulu.
\end{proof}

\begin{lemma}
\label{etale-cohomology-lemma-shriek-equals-star-finite-etale}
Soit $j : U \to X$ fini et étale. Alors le morphisme
$j_! \to j_*$ du lemme \ref{etale-cohomology-lemma-shriek-into-star-separated-etale}
est un isomorphisme
sur les faisceaux abéliens et les faisceaux de $\Lambda$-modules.
\end{lemma}

\begin{proof}
Il suffit de vérifier que $j_!\mathcal{F} \to j_*\mathcal{F}$
est un isomorphisme étale-localement sur $X$.
On peut donc supposer que $U \to X$ est une union disjointe finie
d'isomorphismes; voir
Morphismes étales, lemme \ref{etale-lemma-finite-etale-etale-local}.
Nous omettons la démonstration dans ce cas.
\end{proof}

\begin{lemma}
\label{etale-cohomology-lemma-ses-associated-to-open}
Soit $X$ un schéma. Soit $Z \subset X$ un sous-schéma fermé, et soit
$U \subset X$ son complémentaire. Notons $i : Z \to X$ et $j : U \to X$
les morphismes d'inclusion. Pour tout faisceau abélien $\mathcal{F}$ sur $X_\etale$,
il existe une suite exacte courte canonique
$$
0 \to j_!j^{-1}\mathcal{F} \to \mathcal{F} \to i_*i^{-1}\mathcal{F} \to 0
$$
sur $X_\etale$.
\end{lemma}

\begin{proof}
Les propriétés d'adjonction des foncteurs en jeu fournissent les
morphismes. Pour un point géométrique $\overline{x}$ de $X$, on a soit
$\overline{x} \in U$, auquel cas le morphisme de gauche induit
un isomorphisme sur les fibres et la fibre de $i_*i^{-1}\mathcal{F}$
est nulle, soit $\overline{x} \in Z$, auquel cas le morphisme de droite induit
un isomorphisme sur les fibres et la fibre de $j_!j^{-1}\mathcal{F}$
est nulle. Nous avons utilisé ici la description des fibres du
lemme \ref{etale-cohomology-lemma-stalk-pushforward-closed-immersion} et de la
proposition \ref{etale-cohomology-proposition-describe-jshriek}.
\end{proof}

\begin{lemma}
\label{etale-cohomology-lemma-compatible-shriek-push-finite}
Considérons un diagramme cartésien de schémas
$$
\xymatrix{
U \ar[d]_g \ar[r]_{j'} & X \ar[d]^f \\
V \ar[r]^j & Y
}
$$
où $f$ est fini et $j$ est une immersion ouverte.
Alors $f_* \circ j'_! = j_! \circ g_*$ comme foncteurs
$\textit{Ab}(U_\etale) \to \textit{Ab}(Y_\etale)$.
\end{lemma}

\begin{proof}
Soit $\mathcal{F}$ un objet de $\textit{Ab}(U_\etale)$.
Soit $\overline{y}$ un point géométrique de $Y$ qui n'appartient pas
à l'ouvert $V$. Alors
$$
(f_*j'_!\mathcal{F})_{\overline{y}} =
\bigoplus\nolimits_{\overline{x},\ f(\overline{x}) = \overline{y}}
(j'_!\mathcal{F})_{\overline{x}} = 0
$$
d'après la proposition \ref{etale-cohomology-proposition-finite-higher-direct-image-zero}, et
parce que la fibre de $j'_!\mathcal{F}$ en $\overline{x} \not \in U$
est nulle d'après le lemme \ref{etale-cohomology-lemma-jshriek-open}. D'autre part, on a
$$
j^{-1}f_*j'_!\mathcal{F} =
g_*(j')^{-1}j'_!\mathcal{F} =
g_*\mathcal{F}
$$
d'après les lemmes \ref{etale-cohomology-lemma-finite-pushforward-commutes-with-base-change} et
\ref{etale-cohomology-lemma-jshriek-open}.
La caractérisation de
$j_!$ dans le lemme \ref{etale-cohomology-lemma-jshriek-open} montre donc que
$f_*j'_!\mathcal{F} = j_!g_*\mathcal{F}$.
Nous omettons la vérification que cette identification
est fonctorielle en $\mathcal{F}$.
\end{proof}







\section{Faisceaux constructibles}
\label{etale-cohomology-section-constructible}

\noindent
Soit $X$ un schéma. Un {\it sous-schéma constructible localement fermé} de $X$
est un sous-schéma localement fermé $T \subset X$ dont l'espace
topologique sous-jacent à $T$ est une partie constructible de $X$.
Si $T, T' \subset X$ sont des sous-schémas localement
fermés ayant le même espace topologique sous-jacent, alors
$T_\etale \cong T'_\etale$ par l'invariance topologique
du site étale (théorème \ref{etale-cohomology-theorem-topological-invariance}).
Ainsi, dans la définition suivante, nous pouvons supposer réduits nos
sous-schémas localement fermés.

\begin{definition}
\label{etale-cohomology-definition-constructible}
Soit $X$ un schéma.
\begin{enumerate}
\item Un faisceau d'ensembles sur $X_\etale$ est {\it constructible}
si, pour tout ouvert affine $U \subset X$, il existe une décomposition finie
de $U$ en sous-schémas constructibles localement fermés $U = \coprod_i U_i$
telle que $\mathcal{F}|_{U_i}$ soit localement constant, de valeur finie, pour tout $i$.
\item Un faisceau de groupes abéliens sur $X_\etale$ est {\it constructible}
si, pour tout ouvert affine $U \subset X$, il existe une décomposition finie
de $U$ en sous-schémas constructibles localement fermés $U = \coprod_i U_i$
telle que $\mathcal{F}|_{U_i}$ soit localement constant, de valeur finie, pour tout $i$.
\item Soit $\Lambda$ un anneau noethérien. Un faisceau de $\Lambda$-modules
sur $X_\etale$ est {\it constructible} si, pour tout ouvert affine
$U \subset X$, il existe une décomposition finie
de $U$ en sous-schémas constructibles localement fermés
$U = \coprod_i U_i$ telle que
$\mathcal{F}|_{U_i}$ soit de type fini et localement constant pour tout $i$.
\end{enumerate}
\end{definition}

\noindent
Il semble que ce soit la définition admise. Une autre définition, qui se prête
plus aisément à des généralisations au-delà du site étale d'un schéma,
aurait consisté à définir les faisceaux constructibles en partant de
$h_U$, $j_{U!}\mathbf{Z}/n\mathbf{Z}$ et $j_{U!}\underline{\Lambda}$
où $U$ parcourt tous les objets quasi-compacts et quasi-séparés
de $X_\etale$, puis à prendre la plus petite sous-catégorie pleine de
$\Sh(X_\etale)$, $\textit{Ab}(X_\etale)$ et
$\textit{Mod}(X_\etale, \underline{\Lambda})$ qui les contienne
et soit stable par limites et colimites finies. Il résulte du
lemme \ref{etale-cohomology-lemma-constructible-abelian}
et des
lemmes \ref{etale-cohomology-lemma-category-constructible-sets},
\ref{etale-cohomology-lemma-category-constructible-abelian} et
\ref{etale-cohomology-lemma-category-constructible-modules}
que cette construction donne la même catégorie si $X$ est quasi-compact et
quasi-séparé. En général, elle ne donne cependant pas la même
catégorie.

\medskip\noindent
Une décomposition en union disjointe $U = \coprod U_i$ d'un schéma en
sous-schémas localement fermés sera appelée une {\it partition} de $U$
(comparer à Topologie, section \ref{topology-section-stratifications}).

\begin{lemma}
\label{etale-cohomology-lemma-constructible-quasi-compact-quasi-separated}
Soit $X$ un schéma quasi-compact et quasi-séparé. Soit $\mathcal{F}$
un faisceau d'ensembles sur $X_\etale$. Les conditions suivantes sont équivalentes :
\begin{enumerate}
\item $\mathcal{F}$ est constructible,
\item il existe un recouvrement ouvert $X = \bigcup U_i$ tel que
$\mathcal{F}|_{U_i}$ soit constructible, et
\item il existe une partition $X = \bigcup X_i$ en sous-schémas
constructibles localement fermés telle que $\mathcal{F}|_{X_i}$ soit
localement constant, de valeur finie.
\end{enumerate}
Un énoncé analogue vaut pour les faisceaux abéliens et les faisceaux de
$\Lambda$-modules si $\Lambda$ est noethérien.
\end{lemma}

\begin{proof}
Il est clair que (1) implique (2).

\medskip\noindent
Supposons (2). Pour tout $x \in X$, il existe un $i$ et un voisinage ouvert
affine $V_x \subset U_i$ de $x$. On peut donc trouver un recouvrement
ouvert affine fini $X = \bigcup V_j$ tel que, pour tout $j$, il
existe une décomposition finie $V_j = \coprod V_{j, k}$ en parties localement fermées
constructibles telle que $\mathcal{F}|_{V_{j, k}}$ soit localement constant,
de valeur finie. D'après Topologie, lemme
\ref{topology-lemma-quasi-compact-open-immersion-constructible-image}
chaque $V_{j, k}$ est une partie constructible de $X$.
D'après Topologie, lemme
\ref{topology-lemma-constructible-partition-refined-by-stratification}
il existe une stratification finie $X = \coprod X_l$ en strates constructibles
localement fermées telle que chacun des
$V_{j, k}$ soit une réunion de $X_l$. Ainsi, (3) est satisfaite.

\medskip\noindent
Supposons (3) satisfaite. Soit $U \subset X$ un ouvert affine.
Alors $U \cap X_i$ est une partie constructible localement fermée de $U$
(par exemple d'après Propriétés, lemme \ref{properties-lemma-locally-constructible})
et $U = \coprod U \cap X_i$ est une partition de $U$ comme dans la
définition \ref{etale-cohomology-definition-constructible}. Ainsi, (1) est satisfaite.
\end{proof}

\begin{lemma}
\label{etale-cohomology-lemma-constructible-constructible}
Soit $X$ un schéma quasi-compact et quasi-séparé.
Soit $\mathcal{F}$ un faisceau d'ensembles, de groupes abéliens ou de
$\Lambda$-modules (avec $\Lambda$ noethérien) sur $X_\etale$.
S'il existe des sous-schémas constructibles localement fermés $T_i \subset X$
tels que (a) $X = \bigcup T_j$ et que (b) $\mathcal{F}|_{T_j}$ soit
constructible, alors $\mathcal{F}$ est constructible.
\end{lemma}

\begin{proof}
Tout d'abord, on peut supposer le recouvrement fini puisque $X$ est quasi-compact
pour la topologie spectrale
(Topologie, lemme \ref{topology-lemma-constructible-hausdorff-quasi-compact}
et
Propriétés, lemme
\ref{properties-lemma-quasi-compact-quasi-separated-spectral}).
Remarquons que chaque $T_i$ est lui-même un schéma quasi-compact et quasi-séparé
(puisqu'il est constructible dans $X$; détails omis).
Il existe donc une partition finie $T_i = \coprod T_{i, j}$ en
parties localement fermées constructibles de $T_i$ telle que
$\mathcal{F}|_{T_{i, j}}$ soit localement constant, de valeur finie
(lemme \ref{etale-cohomology-lemma-constructible-quasi-compact-quasi-separated}).
D'après Topologie, lemme \ref{topology-lemma-constructible-in-constructible},
$T_{i, j}$ est un sous-schéma constructible localement fermé de $X$.
On peut alors appliquer Topologie, lemme
\ref{topology-lemma-constructible-partition-refined-by-stratification}
à $X = \bigcup T_{i, j}$ pour obtenir la partition cherchée de $X$.
\end{proof}

\begin{lemma}
\label{etale-cohomology-lemma-constructible-local}
Soit $X$ un schéma. La constructibilité d'un faisceau
d'ensembles, de groupes abéliens ou de $\Lambda$-modules (avec $\Lambda$ noethérien)
se vérifie localement pour la topologie de Zariski sur $X$.
\end{lemma}

\begin{proof}
L'énoncé signifie que, si $X = \bigcup U_i$ est un recouvrement ouvert
tel que $\mathcal{F}|_{U_i}$ soit constructible, alors $\mathcal{F}$
est constructible. Si $U \subset X$ est un ouvert affine, alors
$U = \bigcup U \cap U_i$ et $\mathcal{F}|_{U \cap U_i}$ est constructible
(il est évident que la restriction d'un faisceau constructible à
un ouvert est constructible). Il résulte du
lemme \ref{etale-cohomology-lemma-constructible-quasi-compact-quasi-separated}
que $\mathcal{F}|_U$ est constructible, c'est-à-dire qu'il existe une partition convenable
de $U$.
\end{proof}

\begin{lemma}
\label{etale-cohomology-lemma-pullback-constructible}
Soit $f : X \to Y$ un morphisme de schémas. Si $\mathcal{F}$ est un
faisceau constructible d'ensembles, de groupes abéliens ou de $\Lambda$-modules
(avec $\Lambda$ noethérien) sur $Y_\etale$, il en est de même
de $f^{-1}\mathcal{F}$ sur $X_\etale$.
\end{lemma}

\begin{proof}
D'après le lemme \ref{etale-cohomology-lemma-constructible-local}, on se ramène au cas
où $X$ et $Y$ sont affines. D'après le
lemme \ref{etale-cohomology-lemma-constructible-quasi-compact-quasi-separated},
il suffit de trouver une partition finie de $X$ en sous-schémas constructibles
localement fermés sur chacun desquels $f^{-1}\mathcal{F}$ soit localement
constant, de valeur finie.
Pour la trouver, il suffit de prendre l'image inverse de la partition de $Y$ adaptée à
$\mathcal{F}$ et d'utiliser le
lemme \ref{etale-cohomology-lemma-pullback-locally-constant}.
\end{proof}

\begin{lemma}
\label{etale-cohomology-lemma-constructible-abelian}
Soit $X$ un schéma.
\begin{enumerate}
\item La catégorie des faisceaux constructibles d'ensembles
est stable par limites et colimites finies dans $\Sh(X_\etale)$.
\item La catégorie des faisceaux abéliens constructibles est une
sous-catégorie de Serre faible de $\textit{Ab}(X_\etale)$.
\item Soit $\Lambda$ un anneau noethérien. La catégorie des
faisceaux constructibles de $\Lambda$-modules sur
$X_\etale$ est une sous-catégorie de Serre faible de
$\textit{Mod}(X_\etale, \Lambda)$.
\end{enumerate}
\end{lemma}

\begin{proof}
Nous démontrons (3). Nous utiliserons le critère donné dans
Homologie, lemme \ref{homology-lemma-characterize-weak-serre-subcategory}.
Supposons que $\varphi : \mathcal{F} \to \mathcal{G}$
soit un morphisme de faisceaux constructibles de $\Lambda$-modules.
Il faut montrer que $\mathcal{K} = \Ker(\varphi)$ et
$\mathcal{Q} = \Coker(\varphi)$ sont constructibles.
De même, supposons que
$0 \to \mathcal{F} \to \mathcal{E} \to \mathcal{G} \to 0$
soit une suite exacte courte de faisceaux de $\Lambda$-modules
avec $\mathcal{F}$ et $\mathcal{G}$ constructibles. Il faut montrer
que $\mathcal{E}$ est constructible.
Dans les deux cas, on peut remplacer $X$ par les membres d'un
recouvrement ouvert affine. On peut donc supposer $X$ affine.
On peut ensuite remplacer $X$ par les membres d'une
partition finie de $X$ en sous-schémas constructibles localement fermés
sur lesquels $\mathcal{F}$ et $\mathcal{G}$ sont de type fini et
localement constants. On conclut alors par le
lemme \ref{etale-cohomology-lemma-kernel-finite-locally-constant}.

\medskip\noindent
Les démonstrations de (1) et (2), très semblables, sont omises.
\end{proof}

\begin{lemma}
\label{etale-cohomology-lemma-support-constructible}
Soit $X$ un schéma quasi-compact et quasi-séparé.
\begin{enumerate}
\item Soit $\mathcal{F} \to \mathcal{G}$ un morphisme de faisceaux constructibles
d'ensembles sur $X_\etale$. Alors l'ensemble des points $x \in X$
où $\mathcal{F}_{\overline{x}} \to \mathcal{G}_{\overline{x}}$
est surjectif, resp.\ injectif, resp.\ isomorphe à une application donnée
d'ensembles, est constructible dans $X$.
\item Soit $\mathcal{F}$ un faisceau abélien constructible sur $X_\etale$.
Le support de $\mathcal{F}$ est constructible.
\item Soit $\Lambda$ un anneau noethérien.
Soit $\mathcal{F}$ un faisceau constructible de $\Lambda$-modules sur $X_\etale$.
Le support de $\mathcal{F}$ est constructible.
\end{enumerate}
\end{lemma}

\begin{proof}
Démontrons (1).
Soit $X = \coprod X_i$ une partition de $X$ en sous-schémas localement fermés
constructibles telle que $\mathcal{F}$ et $\mathcal{G}$ soient tous deux
localement constants, de valeur finie, sur les parties (appliquer le
lemme \ref{etale-cohomology-lemma-constructible-quasi-compact-quasi-separated}
à $\mathcal{F}$ et à $\mathcal{G}$, puis choisir un
raffinement commun). On applique alors le lemme \ref{etale-cohomology-lemma-morphism-locally-constant}
à la restriction du morphisme à chaque partie.

\medskip\noindent
Les démonstrations de (2) et (3) sont omises.
\end{proof}

\noindent
Le lemme suivant se révélera très utile par la suite.
Il affirme en substance que la catégorie des faisceaux constructibles
possède une forme faible de propriété ``noethérienne''.

\begin{lemma}
\label{etale-cohomology-lemma-colimit-constructible}
Soit $X$ un schéma quasi-compact et quasi-séparé. Soit
$\mathcal{F} = \colim_{i \in I} \mathcal{F}_i$ une limite inductive filtrante
de faisceaux d'ensembles, de faisceaux abéliens ou de faisceaux de modules.
\begin{enumerate}
\item Si $\mathcal{F}$ et les $\mathcal{F}_i$ sont des faisceaux constructibles
d'ensembles, alors l'ind-objet $\mathcal{F}_i$ est essentiellement constant de
valeur $\mathcal{F}$.
\item Si $\mathcal{F}$ et les $\mathcal{F}_i$ sont des faisceaux constructibles de
groupes abéliens, alors l'ind-objet $\mathcal{F}_i$ est essentiellement constant
de valeur $\mathcal{F}$.
\item Soit $\Lambda$ un anneau noethérien.
Si $\mathcal{F}$ et les $\mathcal{F}_i$ sont des faisceaux constructibles de
$\Lambda$-modules, alors l'ind-objet $\mathcal{F}_i$ est essentiellement constant
de valeur $\mathcal{F}$.
\end{enumerate}
\end{lemma}

\begin{proof}
Démontrons (1). Nous utiliserons sans plus le mentionner que les limites
et colimites finies de faisceaux constructibles sont constructibles
(lemme \ref{etale-cohomology-lemma-constructible-abelian}).
Pour tout $i$, soit $T_i \subset X$ l'ensemble des points $x \in X$
où $\mathcal{F}_{i, \overline{x}} \to \mathcal{F}_{\overline{x}}$
n'est pas surjectif. Puisque $\mathcal{F}_i$ et $\mathcal{F}$ sont
constructibles, $T_i$ est une partie constructible de $X$
(lemme \ref{etale-cohomology-lemma-support-constructible}).
Comme les fibres de $\mathcal{F}$ sont finies
et que $\mathcal{F} = \colim_{i \in I} \mathcal{F}_i$, on voit
que, pour tout $x \in X$, on a $x \not \in T_i$ pour $i$ assez grand.
Comme $X$ est un espace spectral d'après Propriétés, lemme
\ref{properties-lemma-quasi-compact-quasi-separated-spectral}
la topologie constructible sur $X$ est quasi-compact d'après
Topologie, lemme \ref{topology-lemma-constructible-hausdorff-quasi-compact}.
Ainsi, $T_i = \emptyset$ pour $i$ assez grand. Par conséquent,
$\mathcal{F}_i \to \mathcal{F}$ est surjectif pour $i$ assez grand.
Supposons maintenant que $\mathcal{F}_i \to \mathcal{F}$ soit surjectif pour tout $i$.
Fixons $i \in I$. Pour $i' \geq i$, notons $S_{i'} \subset X$ l'ensemble des
points $x$ pour lesquels le nombre d'éléments de
$\Im(\mathcal{F}_{i, \overline{x}} \to \mathcal{F}_{\overline{x}})$
n'est pas égal au nombre d'éléments de
$\Im(\mathcal{F}_{i, \overline{x}} \to \mathcal{F}_{i', \overline{x}})$.
Puisque $\mathcal{F}_i$, $\mathcal{F}_{i'}$ et $\mathcal{F}$ sont
constructibles, $S_{i'}$ est une partie constructible de $X$
(détails omis ; indication : utiliser le lemme \ref{etale-cohomology-lemma-support-constructible}).
Comme les fibres de $\mathcal{F}_i$ et de $\mathcal{F}$
sont finies et que $\mathcal{F} = \colim_{i' \geq i} \mathcal{F}_{i'}$,
on voit que, pour tout $x \in X$, on a $x \not \in S_{i'}$ pour $i'$
assez grand. Par le même argument que ci-dessus, il existe un $i'$ assez grand
tel que $S_{i'} = \emptyset$. Ainsi, $\mathcal{F}_i \to \mathcal{F}_{i'}$
se factorise par $\mathcal{F}$, comme souhaité.

\medskip\noindent
Démontrons (2). Remarquons qu'un faisceau abélien constructible est un faisceau
d'ensembles constructible. Le cas (2) résulte donc de (1).

\medskip\noindent
Démontrons (3). Nous utiliserons sans plus le mentionner que la catégorie des
faisceaux constructibles de $\Lambda$-modules est abélienne
(lemme \ref{etale-cohomology-lemma-constructible-abelian}).
Pour tout $i$, soit $\mathcal{Q}_i$ le conoyau du morphisme
$\mathcal{F}_i \to \mathcal{F}$. Le support $T_i$ de $\mathcal{Q}_i$
est une partie constructible de $X$ puisque $\mathcal{Q}_i$ est constructible
(lemme \ref{etale-cohomology-lemma-support-constructible}).
Comme les fibres de $\mathcal{F}$ sont des $\Lambda$-modules de type fini
et que $\mathcal{F} = \colim_{i \in I} \mathcal{F}_i$, on voit
que, pour tout $x \in X$, on a $x \not \in T_i$ pour $i$ assez grand.
Comme $X$ est un espace spectral d'après Propriétés, lemme
\ref{properties-lemma-quasi-compact-quasi-separated-spectral}
la topologie constructible sur $X$ est quasi-compact d'après
Topologie, lemme \ref{topology-lemma-constructible-hausdorff-quasi-compact}.
Ainsi, $T_i = \emptyset$ pour $i$ assez grand. Cela démontre la première
assertion. Pour la seconde, supposons maintenant que
$\mathcal{F}_i \to \mathcal{F}$ soit surjectif pour tout $i$.
Fixons $i \in I$. Pour $i' \geq i$, désignons par $\mathcal{K}_{i'}$
l'image de $\Ker(\mathcal{F}_i \to \mathcal{F})$ dans $\mathcal{F}_{i'}$.
Le support $S_{i'}$ de $\mathcal{K}_{i'}$
est une partie constructible de $X$ puisque $\mathcal{K}_{i'}$ est constructible.
Comme les fibres de $\Ker(\mathcal{F}_i \to \mathcal{F})$
sont des $\Lambda$-modules de type fini et que
$\mathcal{F} = \colim_{i' \geq i} \mathcal{F}_{i'}$, on voit
que, pour tout $x \in X$, on a $x \not \in S_{i'}$ pour $i'$ assez grand.
Par le même argument que ci-dessus, il existe un $i'$ assez grand
tel que $S_{i'} = \emptyset$. Ainsi, $\mathcal{F}_i \to \mathcal{F}_{i'}$
se factorise par $\mathcal{F}$, comme souhaité.
\end{proof}

\begin{lemma}
\label{etale-cohomology-lemma-tensor-product-constructible}
Soit $X$ un schéma. Soit $\Lambda$ un anneau noethérien.
Le produit tensoriel de deux faisceaux constructibles de $\Lambda$-modules
sur $X_\etale$ est un faisceau constructible de $\Lambda$-modules.
\end{lemma}

\begin{proof}
L'assertion se ramène immédiatement au cas où $X$ est affine.
Comme deux partitions quelconques de $X$ en strates constructibles localement
fermées admettent un raffinement commun du même type et
que les images inverses commutent au produit tensoriel, on se ramène au
lemme \ref{etale-cohomology-lemma-tensor-product-locally-constant}.
\end{proof}

\begin{lemma}
\label{etale-cohomology-lemma-tensor-constructible}
Soit $\Lambda \to \Lambda'$ un homomorphisme d'anneaux noethériens.
Soit $X$ un schéma. Soit $\mathcal{F}$ un faisceau constructible
de $\Lambda$-modules sur $X_\etale$. Alors
$\mathcal{F} \otimes_{\underline{\Lambda}} \underline{\Lambda'}$
est un faisceau constructible de $\Lambda'$-modules.
\end{lemma}

\begin{proof}
Omis. Indication : localement sur un ouvert affine, on peut utiliser la même stratification.
\end{proof}










\section{Lemmes auxiliaires sur les morphismes}
\label{etale-cohomology-section-stratify-morphisms}

\noindent
Voici quelques lemmes utiles pour d\'emontrer des propri\'et\'es de fonctorialit\'e
des faisceaux constructibles.

\begin{lemma}
\label{etale-cohomology-lemma-etale-stratified-finite}
Soit $U \to X$ un morphisme \'etale de sch\'emas quasi-compacts et quasi-s\'epar\'es
(par exemple un morphisme \'etale de sch\'emas noeth\'eriens). Alors il
existe une partition $X = \coprod_i X_i$ en sous-sch\'emas constructibles
localement ferm\'es telle que $X_i \times_X U \to X_i$ soit fini \'etale pour tout $i$.
\end{lemma}

\begin{proof}
Si $U \to X$ est s\'epar\'e, cela r\'esulte de
Compléments sur les morphismes, lemme \ref{more-morphisms-lemma-stratify-flat-fp-qf}.
En g\'en\'eral, on peut supposer $X$ affine. Choisissons un recouvrement ouvert
affine fini $U = \bigcup U_j$. Appliquons le cas pr\'ec\'edent \`a tous les morphismes
$U_j \to X$ et $U_j \cap U_{j'} \to X$, puis choisissons un raffinement
commun $X = \coprod X_i$ des partitions obtenues.
En raffinant encore la partition, on peut aussi supposer $X_i$ affine.
Fixons $i$ et posons $V = U \times_X X_i$. Les morphismes
$V_j = U_j \times_X X_i \to X_i$ et
$V_{jj'} = (U_j \cap U_{j'}) \times_X X_i \to X_i$ sont finis \'etales.
Ainsi, $V_j$ et $V_{jj'}$ sont des sch\'emas affines et $V_{jj'} \subset V_j$
est \`a la fois ferm\'e et ouvert (puisque $V_{jj'} \to X_i$ est propre, si bien que
Morphismes, lemme \ref{morphisms-lemma-image-proper-scheme-closed}
s'applique). Alors $V = \bigcup V_j$ est s\'epar\'e car
$\mathcal{O}(V_j) \to \mathcal{O}(V_{jj'})$ est surjectif, voir
Schémas, lemme \ref{schemes-lemma-characterize-separated}.
Le cas pr\'ec\'edent s'applique donc \`a $V \to X_i$, et l'on peut encore
raffiner la partition si n\'ecessaire (en fait, cela ne l'est pas, mais nous n'en
avons pas besoin).
\end{proof}

\noindent
Dans le cas noeth\'erien, on peut d\'emontrer le lemme pr\'ec\'edent par
r\'ecurrence noeth\'erienne \`a l'aide du petit lemme amusant suivant.

\begin{lemma}
\label{etale-cohomology-lemma-generically-finite}
Soit $f: X \to Y$ un morphisme de sch\'emas quasi-compact,
quasi-s\'epar\'e et localement de type fini. Si $\eta$ est un point g\'en\'erique
d'une composante irr\'eductible de $Y$ tel que $f^{-1}(\eta)$ soit fini, alors
il existe un ouvert $V \subset Y$ contenant $\eta$ tel que
$f^{-1}(V) \to V$ soit fini.
\end{lemma}

\begin{proof}
C'est Morphismes, lemme \ref{morphisms-lemma-generically-finite}.
\end{proof}

\noindent
L'\'enonc\'e du lemme suivant peut \^etre l\'eg\`erement renforc\'e.

\begin{lemma}
\label{etale-cohomology-lemma-decompose-quasi-finite-morphism}
Soit $f : Y \to X$ un morphisme quasi-fini et de pr\'esentation finie
de sch\'emas affines.
\begin{enumerate}
\item Il existe un morphisme surjectif de sch\'emas affines $X' \to X$ et un
sous-sch\'ema ferm\'e $Z' \subset Y' = X' \times_X Y$ tels que
\begin{enumerate}
\item l'immersion $Z' \subset Y'$ est un \'epaississement, et
\item $Z' \to X'$ est un morphisme fini \'etale.
\end{enumerate}
\item Il existe une partition finie $X = \coprod X_i$ en
strates affines, constructibles et localement ferm\'ees, et des morphismes surjectifs
finis localement libres $X'_i \to X_i$ tels que le morphisme r\'eduit associ\'e \`a
$Y'_i = X'_i \times_X Y \to X'_i$ soit isomorphe \`a
$\coprod_{j = 1}^{n_i} (X'_i)_{red} \to (X'_i)_{red}$ pour un certain $n_i$.
\end{enumerate}
\end{lemma}

\begin{proof}
En posant $X' = \coprod X'_i$, on voit que (2) implique (1).
\'Ecrivons $X = \Spec(A)$ et $Y = \Spec(B)$. Exprimons $A$ comme limite inductive filtrante
de $\mathbf{Z}$-alg\`ebres de type fini $A_i$. Comme $B$ est une $A$-alg\`ebre de
pr\'esentation finie, on voit qu'il existe $0 \in I$ et un
homomorphisme d'anneaux de type fini $A_0 \to B_0$ tels que $B = \colim B_i$, o\`u
$B_i = A_i \otimes_{A_0} B_0$, voir
Algèbre, lemme \ref{algebra-lemma-colimit-category-fp-algebras}.
Pour $i$ suffisamment grand, on voit que $A_i \to B_i$ est
quasi-fini, voir Limits, lemme \ref{limits-lemma-descend-quasi-finite}.
Nous sommes donc ramen\'es au cas des alg\`ebres de type fini sur $\mathbf{Z}$,
et en particulier au cas noeth\'erien. (D\'etails omis.)

\medskip\noindent
Supposons $X$ et $Y$ noeth\'eriens. Dans ce cas, tout sous-ensemble localement ferm\'e
de $X$ est constructible. Le lemme \ref{etale-cohomology-lemma-generically-finite}
et la r\'ecurrence noeth\'erienne montrent qu'il existe
une partition finie $X = \coprod X_i$ de $X$
en strates localement ferm\'ees telle que $Y \times_X X_i \to X_i$ soit fini.
On peut raffiner cette partition de sorte que les strates soient affines.
Ainsi, apr\`es avoir remplac\'e $X$ par $X' = \coprod X_i$, on peut supposer que
$Y \to X$ est fini.

\medskip\noindent
Supposons $X$ et $Y$ noeth\'eriens et $Y \to X$ fini.
Supposons pouvoir d\'emontrer (2) apr\`es changement de base par un morphisme
surjectif, plat et quasi-fini $U \to X$. On dispose donc d'une partition
$U = \coprod U_i$ et de morphismes finis localement libres $U'_i \to U_i$
tels que $U'_i \times_X Y \to U'_i$ soit isomorphe \`a
$\coprod_{j = 1}^{n_i} (U'_i)_{red} \to (U'_i)_{red}$ pour un certain $n_i$.
Alors, par l'argument du paragraphe pr\'ec\'edent, on peut trouver une
partition $X = \coprod X_j$ en strates affines localement ferm\'ees telle que
$X_j \times_X U_i \to X_j$ soit fini pour tous $i, j$. D'apr\`es
Morphismes, lemme \ref{morphisms-lemma-finite-flat},
chaque $X_j \times_X U_i \to X_j$ est fini localement libre.
Par cons\'equent, $X_j \times_X U'_i \to X_j$ est fini localement libre
(Morphismes, lemme \ref{morphisms-lemma-composition-finite-locally-free}).
Il s'ensuit que $X = \coprod X_j$ et $X_j' = \coprod_i X_j \times_X U'_i$
fournissent une solution pour $Y \to X$. Il suffit donc de d\'emontrer
le r\'esultat (dans le cas noeth\'erien) apr\`es un changement de base surjectif,
plat et quasi-fini.

\medskip\noindent
En appliquant Morphismes, lemme \ref{morphisms-lemma-massage-finite},
on voit qu'on peut supposer que $Y$ est un sous-sch\'ema ferm\'e d'un
sch\'ema affine $Z$ qui est, ensemblistement, une r\'eunion finie
$Z = \bigcup_{i \in I} Z_i$ de sous-sch\'emas ferm\'es s'envoyant isomorphiquement
sur $X$. Dans ce cas, on trouvera une partition finie $X = \coprod X_j$
en strates affines localement ferm\'ees qui convient (autrement dit $X'_j = X_j$).
Posons $T_i = Y \cap Z_i$. C'est un sous-sch\'ema ferm\'e de $X$.
Comme $X$ est noeth\'erien, on peut trouver une partition finie $X = \coprod X_j$
en sous-sch\'emas affines localement ferm\'es, telle que chaque
$X_j \times_X T_i$ soit vide ou soit, ensemblistement, \`egal \`a $X_j \times_X Z_i$.
En rempla\c cant $X$ par $X_j$, on voit qu'on peut supposer $I = I_1 \amalg I_2$,
avec $Z_i \subset Y$ pour $i \in I_1$ et $Z_i \cap Y = \emptyset$ pour
$i \in I_2$. En rempla\c cant $Z$ par $\bigcup_{i \in I_1} Z_i$, on voit qu'on
peut supposer $Y = Z$.
Enfin, on peut remplacer de nouveau $X$ par les membres d'une partition
comme ci-dessus telle que, pour tous $i, i' \in I$, l'intersection
$Z_i \cap Z_{i'}$ soit vide ou soit, ensemblistement, \`egale
\`a $Z_i$ et \`a $Z_{i'}$. Cela signifie clairement que $Y$ est, ensemblistement,
\'egal \`a une r\'eunion disjointe des $Z_i$, ce qu'il fallait d\'emontrer.
\end{proof}








\section{Compl\'ements sur les faisceaux constructibles}
\label{etale-cohomology-section-more-constructible}

\noindent
Soit $\Lambda$ un anneau noeth\'erien. Soit $X$ un sch\'ema.
Nous consid\'erons souvent $X_\etale$ comme un site annel\'e dont le
faisceau d'anneaux est $\underline{\Lambda}$. Pour les faisceaux ab\'eliens,
nous prenons souvent $\Lambda = \mathbf{Z}/n\mathbf{Z}$ pour un
entier $n$ convenable.

\begin{lemma}
\label{etale-cohomology-lemma-jshriek-constructible}
Soit $j : U \to X$ un morphisme \'etale de sch\'emas quasi-compacts et
quasi-s\'epar\'es.
\begin{enumerate}
\item Le faisceau $h_U$ est un faisceau d'ensembles constructible.
\item Le faisceau $j_!\underline{M}$ est un faisceau ab\'elien constructible
pour tout groupe ab\'elien fini $M$.
\item Si $\Lambda$ est un anneau noeth\'erien et $M$ un $\Lambda$-module de type fini,
alors $j_!\underline{M}$ est un faisceau constructible de $\Lambda$-modules
sur $X_\etale$.
\end{enumerate}
\end{lemma}

\begin{proof}
D'apr\`es le lemme \ref{etale-cohomology-lemma-etale-stratified-finite}, il existe une partition
$\coprod_i X_i$ telle que $\pi_i : j^{-1}(X_i) \to X_i$ soit fini \'etale.
La restriction de $h_U$ \`a $X_i$ est $h_{j^{-1}(X_i)}$, qui est fini et
localement constant d'apr\`es le lemme \ref{etale-cohomology-lemma-characterize-finite-locally-constant}.
Pour les cas (2) et (3), observons que
$$
j_!(\underline{M})|_{X_i} =
\pi_{i!}(\underline{M}) =
\pi_{i*}(\underline{M})
$$
par les lemmes \ref{etale-cohomology-lemma-shriek-base-change} et
\ref{etale-cohomology-lemma-shriek-equals-star-finite-etale}.
Il suffit donc de d\'emontrer le lemme lorsque $\pi : Y \to X$ est fini \'etale.
C'est le lemme \ref{etale-cohomology-lemma-pushforward-locally-constant}.
\end{proof}

\begin{lemma}
\label{etale-cohomology-lemma-torsion-colimit-constructible}
Soit $X$ un sch\'ema quasi-compact et quasi-s\'epar\'e.
\begin{enumerate}
\item Soit $\mathcal{F}$ un faisceau d'ensembles sur $X_\etale$.
Alors $\mathcal{F}$ est une limite inductive filtrante de
faisceaux d'ensembles constructibles.
\item Soit $\mathcal{F}$ un faisceau ab\'elien de torsion sur $X_\etale$.
Alors $\mathcal{F}$ est une limite inductive filtrante de faisceaux ab\'eliens constructibles.
\item Soient $\Lambda$ un anneau noeth\'erien et $\mathcal{F}$ un faisceau
de $\Lambda$-modules sur $X_\etale$. Alors
$\mathcal{F}$ est une limite inductive filtrante de faisceaux constructibles de
$\Lambda$-modules.
\end{enumerate}
\end{lemma}

\begin{proof}
Soit $\mathcal{B}$ la collection des objets quasi-compacts et quasi-s\'epar\'es
de $X_\etale$. D'apr\`es Modules sur les sites,
lemme \ref{sites-modules-lemma-module-filtered-colimit-constructibles},
tout faisceau d'ensembles est une limite inductive filtrante de faisceaux de la forme
$$
\text{Conoyau de la double fl\`eche}\left(
\xymatrix{
\coprod\nolimits_{j = 1}^m h_{V_j}
\ar@<1ex>[r] \ar@<-1ex>[r] &
\coprod\nolimits_{i = 1}^n h_{U_i}
}
\right)
$$
o\`u $V_j$ et $U_i$ sont des objets quasi-compacts et quasi-s\'epar\'es
de $X_\etale$. D'apr\`es les
lemmes \ref{etale-cohomology-lemma-jshriek-constructible} et \ref{etale-cohomology-lemma-constructible-abelian},
ces conoyaux de doubles fl\`eches sont constructibles. Cela d\'emontre (1).

\medskip\noindent
Soit $\Lambda$ un anneau noeth\'erien.
D'apr\`es Modules sur les sites,
lemme \ref{sites-modules-lemma-module-filtered-colimit-constructibles},
le faisceau de $\Lambda$-modules $\mathcal{F}$ est une limite inductive filtrante
de modules de la forme
$$
\Coker\left(
\bigoplus\nolimits_{j = 1}^m j_{V_j!}\underline{\Lambda}_{V_j}
\longrightarrow
\bigoplus\nolimits_{i = 1}^n j_{U_i!}\underline{\Lambda}_{U_i}
\right)
$$
o\`u $V_j$ et $U_i$ sont des objets quasi-compacts et quasi-s\'epar\'es
de $X_\etale$. D'apr\`es les
lemmes \ref{etale-cohomology-lemma-jshriek-constructible} et \ref{etale-cohomology-lemma-constructible-abelian},
ces conoyaux sont constructibles. Cela d\'emontre (3).

\medskip\noindent
D\'emontrons (2). \'Ecrivons d'abord $\mathcal{F} = \bigcup \mathcal{F}[n]$, o\`u
$\mathcal{F}[n]$ est le sous-faisceau de $n$-torsion. On peut alors consid\'erer
$\mathcal{F}[n]$ comme un faisceau de $\mathbf{Z}/n\mathbf{Z}$-modules
et appliquer (3).
\end{proof}

\begin{lemma}
\label{etale-cohomology-lemma-check-constructible}
Soit $f : X \to Y$ un morphisme surjectif de sch\'emas quasi-compacts et
quasi-s\'epar\'es.
\begin{enumerate}
\item Soit $\mathcal{F}$ un faisceau d'ensembles sur $Y_\etale$. Alors $\mathcal{F}$
est constructible si et seulement si $f^{-1}\mathcal{F}$ est constructible.
\item Soit $\mathcal{F}$ un faisceau ab\'elien sur $Y_\etale$. Alors $\mathcal{F}$
est constructible si et seulement si $f^{-1}\mathcal{F}$ est constructible.
\item Soit $\Lambda$ un anneau noeth\'erien.
Soit $\mathcal{F}$ un faisceau de $\Lambda$-modules sur $Y_\etale$.
Alors $\mathcal{F}$ est constructible si et seulement si $f^{-1}\mathcal{F}$
est constructible.
\end{enumerate}
\end{lemma}

\begin{proof}
Une implication r\'esulte du lemme \ref{etale-cohomology-lemma-pullback-constructible}.
Pour la r\'eciproque, supposons $f^{-1}\mathcal{F}$ constructible. Dans le cas ab\'elien,
la surjectivit\'e montre en outre que le faisceau initial est de torsion. \'Ecrivons $\mathcal{F} = \colim \mathcal{F}_i$ comme
limite inductive filtrante de faisceaux constructibles (d'ensembles, de groupes ab\'eliens ou de modules)
\`a l'aide du lemme \ref{etale-cohomology-lemma-torsion-colimit-constructible}.
Comme $f^{-1}$ est un adjoint \`a gauche, il commute aux limites inductives
(Catégories, lemme \ref{categories-lemma-adjoint-exact}), et l'on obtient
$f^{-1}\mathcal{F} = \colim f^{-1}\mathcal{F}_i$.
D'apr\`es le lemme \ref{etale-cohomology-lemma-colimit-constructible}, on voit que
$f^{-1}\mathcal{F}_i \to f^{-1}\mathcal{F}$
est surjectif pour tout $i$ suffisamment grand.
Comme $f$ est surjectif, on conclut (en examinant les fibres \`a l'aide du
lemme \ref{etale-cohomology-lemma-stalk-pullback} et du
th\'eor\`eme \ref{etale-cohomology-theorem-exactness-stalks})
que $\mathcal{F}_i \to \mathcal{F}$ est surjectif pour tout $i$ suffisamment grand.
Ainsi, $\mathcal{F}$ est le quotient d'un faisceau constructible $\mathcal{G}$.
En appliquant encore une fois l'argument \`a
$\mathcal{G} \times_\mathcal{F} \mathcal{G}$ ou
au noyau de $\mathcal{G} \to \mathcal{F}$,
on conclut en utilisant l'exactitude de $f^{-1}$ et le fait que la cat\'egorie des
faisceaux constructibles (d'ensembles, de groupes ab\'eliens ou de modules) est
stable par limites et colimites finies, ou par noyaux et conoyaux dans
$\Sh(Y_\etale)$, $\Sh(X_\etale)$, $\textit{Ab}(Y_\etale)$,
$\textit{Ab}(X_\etale)$, $\textit{Mod}(Y_\etale, \Lambda)$ et
$\textit{Mod}(X_\etale, \Lambda)$, voir le
lemme \ref{etale-cohomology-lemma-constructible-abelian}.
\end{proof}

\begin{lemma}
\label{etale-cohomology-lemma-pushforward-constructible}
Soit $f : X \to Y$ un morphisme fini \'etale de sch\'emas. Soit $\Lambda$ un
anneau noeth\'erien. Si $\mathcal{F}$ est un faisceau d'ensembles constructible,
un faisceau de groupes ab\'eliens constructible ou un faisceau constructible de
$\Lambda$-modules sur $X_\etale$, il en va de m\^eme de
$f_*\mathcal{F}$ sur $Y_\etale$.
\end{lemma}

\begin{proof}
D'apr\`es le lemme \ref{etale-cohomology-lemma-constructible-local}, il suffit de le v\'erifier
localement pour la topologie de Zariski sur $Y$ et, d'apr\`es le lemme \ref{etale-cohomology-lemma-check-constructible},
on peut remplacer $Y$ par un recouvrement \'etale (la construction de $f_*$
commute \`a la localisation \'etale). Un morphisme fini \'etale est
\'etale-localement isomorphe \`a une union disjointe d'isomorphismes, voir
Morphismes étales, lemme \ref{etale-lemma-finite-etale-etale-local}.
Ainsi, dans le cas des faisceaux d'ensembles, le lemme affirme que si
$\mathcal{F}_i$, $i = 1, \ldots, n$ sont des faisceaux d'ensembles constructibles, alors
$\prod_{i = 1}^n \mathcal{F}_i$ l'est aussi.
C'est clair. Il en va de m\^eme pour les faisceaux de groupes ab\'eliens et de modules.
\end{proof}

\begin{lemma}
\label{etale-cohomology-lemma-category-constructible-sets}
Soit $X$ un sch\'ema quasi-compact et quasi-s\'epar\'e. La cat\'egorie des
faisceaux d'ensembles constructibles est la sous-cat\'egorie pleine de $\Sh(X_\etale)$
form\'ee des faisceaux $\mathcal{F}$ qui sont les conoyaux de doubles fl\`eches
$$
\xymatrix{
\mathcal{F}_1
\ar@<1ex>[r] \ar@<-1ex>[r]
&
\mathcal{F}_0 \ar[r]
&
\mathcal{F}}
$$
pour lesquelles $\mathcal{F}_i$, pour $i = 0, 1$, est un coproduit fini de faisceaux de
la forme $h_U$, o\`u $U$ est un objet quasi-compact et quasi-s\'epar\'e
de $X_\etale$.
\end{lemma}

\begin{proof}
Dans la d\'emonstration du lemme \ref{etale-cohomology-lemma-torsion-colimit-constructible},
nous avons vu que les faisceaux de cette forme sont constructibles.
R\'eciproquement, supposons que, pour tout faisceau d'ensembles constructible
$\mathcal{F}$, on puisse trouver une surjection $\mathcal{F}_0 \to \mathcal{F}$
avec $\mathcal{F}_0$ comme dans le lemme. On trouve alors la surjection voulue
$\mathcal{F}_1 \to \mathcal{F}_0 \times_\mathcal{F} \mathcal{F}_0$,
car le dernier faisceau est constructible d'apr\`es le lemme \ref{etale-cohomology-lemma-constructible-abelian}.

\medskip\noindent
D'apr\`es Topologie, lemme
\ref{topology-lemma-constructible-partition-refined-by-stratification}
on peut choisir une stratification finie
$X = \coprod_{i \in I} X_i$ telle que $\mathcal{F}$ soit finie et localement
constante sur chaque strate. D\'emontrons le r\'esultat par r\'ecurrence sur
le cardinal de $I$. Soit $i \in I$ un \'el\'ement minimal pour l'ordre
partiel de $I$. Alors $X_i \subset X$ est ferm\'e.
Par r\'ecurrence, il existe un nombre fini d'objets quasi-compacts et quasi-s\'epar\'es
$U_\alpha$ de $(X \setminus X_i)_\etale$ et un morphisme surjectif
$\coprod h_{U_\alpha} \to \mathcal{F}|_{X \setminus X_i}$.
Ils d\'eterminent un morphisme
$$
\coprod h_{U_\alpha} \to \mathcal{F}
$$
qui est surjectif apr\`es restriction \`a $X \setminus X_i$. D'apr\`es le
lemme \ref{etale-cohomology-lemma-characterize-finite-locally-constant},
on a $\mathcal{F}|_{X_i} = h_V$ pour un sch\'ema $V$ fini \'etale
sur $X_i$. Soit $\overline{v}$ un point g\'eom\'etrique de $V$ au-dessus
de $\overline{x} \in X_i$. On peut consid\'erer $\overline{v}$ comme un
\'el\'ement de la fibre $\mathcal{F}_{\overline{x}} = V_{\overline{x}}$.
On peut donc trouver un voisinage \'etale $(U, \overline{u})$
de $\overline{x}$ et une section $s \in \mathcal{F}(U)$ dont la fibre en
$\overline{x}$ donne $\overline{v}$. En consid\'erant $s$ comme un morphisme
$s : h_U \to \mathcal{F}$ et en le restreignant \`a $X_i$, on obtient un morphisme
$s|_{X_i} : U \times_X X_i \to V$ au-dessus de $X_i$ qui envoie $\overline{u}$
sur $\overline{v}$. Comme $V$ est quasi-compact (il est fini sur le sous-sch\'ema
ferm\'e $X_i$ du sch\'ema quasi-compact $X$), un nombre fini
$s^{(1)}, \ldots, s^{(m)}$ de ces sections de $\mathcal{F}$ sur
$U^{(1)}, \ldots, U^{(m)}$ d\'eterminent conjointement le morphisme
surjectif
$$
\coprod s^{(j)}|_{X_i} : \coprod U^{(j)} \times_X X_i \longrightarrow V
$$
On obtient alors la surjection
$$
\coprod h_{U_\alpha} \amalg \coprod h_{U^{(j)}} \to \mathcal{F}
$$
voulue.
\end{proof}

\begin{lemma}
\label{etale-cohomology-lemma-category-constructible-modules}
Soit $X$ un schéma quasi-compact et quasi-séparé. Soit $\Lambda$
un anneau noethérien. La catégorie des faisceaux constructibles de
$\Lambda$-modules est exactement la catégorie des modules de la forme
$$
\Coker\left(
\bigoplus\nolimits_{j = 1}^m j_{V_j!}\underline{\Lambda}_{V_j}
\longrightarrow
\bigoplus\nolimits_{i = 1}^n j_{U_i!}\underline{\Lambda}_{U_i}
\right)
$$
où $V_j$ et $U_i$ sont des objets quasi-compacts et quasi-séparés de
$X_\etale$. En fait, on peut même supposer $U_i$ et $V_j$ affines.
\end{lemma}

\begin{proof}
Dans la démonstration du lemme \ref{etale-cohomology-lemma-torsion-colimit-constructible},
nous avons vu que les modules de cette forme sont constructibles. Comme la
catégorie des modules constructibles est abélienne
(lemme \ref{etale-cohomology-lemma-constructible-abelian}),
il suffit de montrer que, pour tout module constructible $\mathcal{F}$,
il existe une surjection
$$
\bigoplus\nolimits_{i = 1}^n j_{U_i!}\underline{\Lambda}_{U_i}
\longrightarrow \mathcal{F}
$$
pour certains objets affines $U_i$ de $X_\etale$. D'après
Modules sur les sites, lemme
\ref{sites-modules-lemma-module-filtered-colimit-constructibles}
il existe une surjection
$$
\Psi :
\bigoplus\nolimits_{i \in I} j_{U_i!}\underline{\Lambda}_{U_i}
\longrightarrow
\mathcal{F}
$$
où les $U_i$ sont affines et où la somme directe est prise sur un
ensemble d'indices $I$ éventuellement infini. Pour toute partie finie $I' \subset I$, posons
$$
T_{I'} = \text{Supp}(\Coker(
\bigoplus\nolimits_{i \in I'} j_{U_i!}\underline{\Lambda}_{U_i}
\longrightarrow \mathcal{F}))
$$
Par la définition même des faisceaux constructibles, l'ensemble $T_{I'}$
est une partie constructible de $X$. Nous voulons montrer que $T_{I'} = \emptyset$
pour un certain $I'$. Comme toute fibre $\mathcal{F}_{\overline{x}}$ est
un $\Lambda$-module de type fini et que $\Psi$ est surjective, pour
tout $x \in X$ il existe un $I'$ tel que $x \not \in T_{I'}$.
Autrement dit, on a
$\emptyset = \bigcap_{I' \subset I\text{ fini}} T_{I'}$. Comme
$X$ est un espace spectral d'après Properties, lemme
\ref{properties-lemma-quasi-compact-quasi-separated-spectral}
la topologie constructible sur $X$ est quasi-compacte d'après
Topologie, lemme \ref{topology-lemma-constructible-hausdorff-quasi-compact}.
Ainsi, $T_{I'} = \emptyset$ pour une partie finie $I' \subset I$,
ce qu'il fallait démontrer.
\end{proof}

\begin{lemma}
\label{etale-cohomology-lemma-category-constructible-abelian}
Soit $X$ un schéma quasi-compact et quasi-séparé. La catégorie des
faisceaux abéliens constructibles est exactement la catégorie des faisceaux
abéliens de la forme
$$
\Coker\left(
\bigoplus\nolimits_{j = 1}^m
j_{V_j!}\underline{\mathbf{Z}/m_j\mathbf{Z}}_{V_j}
\longrightarrow
\bigoplus\nolimits_{i = 1}^n
j_{U_i!}\underline{\mathbf{Z}/n_i\mathbf{Z}}_{U_i}
\right)
$$
où $V_j$ et $U_i$ sont des objets quasi-compacts et quasi-séparés de
$X_\etale$ et où $m_j$, $n_i$ sont des entiers positifs.
En fait, on peut même supposer $U_i$ et $V_j$ affines.
\end{lemma}

\begin{proof}
Cela résulte du lemme \ref{etale-cohomology-lemma-category-constructible-modules}
appliqué avec $\Lambda = \mathbf{Z}/n\mathbf{Z}$
et du fait que, puisque $X$ est quasi-compact, tout faisceau
abélien constructible est annulé par un entier positif $n$ (détails omis).
\end{proof}

\begin{lemma}
\label{etale-cohomology-lemma-constructible-is-compact}
Soit $X$ un schéma quasi-compact et quasi-séparé. Soit $\Lambda$ un
anneau noethérien. Soit $\mathcal{F}$ un faisceau constructible d'ensembles, de groupes
abéliens ou de $\Lambda$-modules sur $X_\etale$. Soit
$\mathcal{G} = \colim \mathcal{G}_i$ une limite inductive filtrante de faisceaux
d'ensembles, de groupes abéliens ou de $\Lambda$-modules. Alors
$$
\Mor(\mathcal{F}, \mathcal{G}) = \colim \Mor(\mathcal{F}, \mathcal{G}_i)
$$
dans la catégorie des faisceaux d'ensembles, de groupes abéliens ou de $\Lambda$-modules sur
$X_\etale$.
\end{lemma}

\begin{proof}
Traitons le cas des faisceaux d'ensembles. D'après le lemme \ref{etale-cohomology-lemma-category-constructible-sets},
il suffit de démontrer le lemme pour $h_U$, où $U$ est un objet quasi-compact
et quasi-séparé de $X_\etale$. Rappelons que
$\Mor(h_U, \mathcal{G}) = \mathcal{G}(U)$. Le résultat
découle donc de Sites, lemme \ref{sites-lemma-directed-colimits-sections}.

\medskip\noindent
Dans le cas des faisceaux abéliens ou des faisceaux de modules, le résultat se démontre
de la même manière à l'aide des
lemmes \ref{etale-cohomology-lemma-category-constructible-abelian} et
\ref{etale-cohomology-lemma-category-constructible-modules}.
Pour les faisceaux abéliens, ajoutons que
$\Mor(j_{U!}\underline{\mathbf{Z}/n\mathbf{Z}}, \mathcal{G})$
est égal à l'ensemble des éléments de $n$-torsion de $\mathcal{G}(U)$.
\end{proof}

\begin{lemma}
\label{etale-cohomology-lemma-finite-pushforward-constructible}
Soit $f : X \to Y$ un morphisme fini et de présentation finie de schémas.
Soit $\Lambda$ un anneau noethérien. Si $\mathcal{F}$ est un faisceau constructible
d'ensembles, de groupes abéliens ou de $\Lambda$-modules sur $X_\etale$,
alors $f_*\mathcal{F}$ l'est aussi.
\end{lemma}

\begin{proof}
Il suffit de le démontrer lorsque $X$ et $Y$ sont affines, d'après le
lemme \ref{etale-cohomology-lemma-constructible-local}.
D'après les
lemmes \ref{etale-cohomology-lemma-finite-pushforward-commutes-with-base-change} et
\ref{etale-cohomology-lemma-check-constructible}, on peut effectuer un changement de base par n'importe quel
schéma affine surjectif sur $Y$. Le
lemme \ref{etale-cohomology-lemma-decompose-quasi-finite-morphism}
nous ramène alors au cas d'un morphisme fini \'etale
(car un épaississement induit une équivalence des topos \'etales
et même des petits sites \'etales, voir le
théorème \ref{etale-cohomology-theorem-topological-invariance}).
Le cas fini \'etale est le
lemme \ref{etale-cohomology-lemma-pushforward-constructible}.
\end{proof}

\begin{lemma}
\label{etale-cohomology-lemma-category-is-colimit}
Soit $X = \lim_{i \in I} X_i$ la limite projective d'un système
projectif filtrant de schémas à morphismes de transition affines.
Nous supposons que $X_i$ est quasi-compact et quasi-séparé
pour tout $i \in I$.
\begin{enumerate}
\item La catégorie des faisceaux d'ensembles constructibles sur $X_\etale$
est la limite inductive des catégories des faisceaux d'ensembles constructibles
sur $(X_i)_\etale$.
\item La catégorie des faisceaux abéliens constructibles sur $X_\etale$
est la limite inductive des catégories des faisceaux abéliens constructibles
sur $(X_i)_\etale$.
\item Soit $\Lambda$ un anneau noethérien. La catégorie des faisceaux constructibles
de $\Lambda$-modules sur $X_\etale$ est la limite inductive des
catégories des faisceaux constructibles de $\Lambda$-modules sur $(X_i)_\etale$.
\end{enumerate}
\end{lemma}

\begin{proof}
Démontrons (1). Notons $f_i : X \to X_i$ les morphismes de projection.
La démonstration comporte trois parties, correspondant à ``fidèle'',
``pleinement fidèle'' et ``essentiellement surjectif''.

\medskip\noindent
Fidélité. Choisissons $0 \in I$ et soient $\mathcal{F}_0$, $\mathcal{G}_0$ des
faisceaux constructibles sur $X_0$. Supposons que
$a, b : \mathcal{F}_0 \to \mathcal{G}_0$ soient des morphismes tels que
$f_0^{-1}a = f_0^{-1}b$. Soit $E \subset X_0$ l'ensemble
des points $x \in X_0$ tels que $a_{\overline{x}} = b_{\overline{x}}$.
D'après le lemme \ref{etale-cohomology-lemma-support-constructible}, la partie
$E \subset X_0$ est constructible. Par hypothèse, $X \to X_0$ est à valeurs dans $E$.
D'après Limits, lemme \ref{limits-lemma-limit-contained-in-constructible},
il existe $i \geq 0$ tel que $X_i \to X_0$ soit à valeurs dans $E$.
Par conséquent, $f_{i0}^{-1}a = f_{i0}^{-1}b$.

\medskip\noindent
Pleine fidélité. Choisissons $0 \in I$ et soient $\mathcal{F}_0$, $\mathcal{G}_0$ des
faisceaux constructibles sur $X_0$. Supposons donné un morphisme
$a : f_0^{-1}\mathcal{F}_0 \to f_0^{-1}\mathcal{G}_0$.
Nous affirmons qu'il existe un $i$ et un morphisme
$a_i : f_{i0}^{-1}\mathcal{F}_0 \to f_{i0}^{-1}\mathcal{G}_0$
dont l'image inverse est $a$ sur $X$.
D'après le lemme \ref{etale-cohomology-lemma-category-constructible-sets},
on peut remplacer $\mathcal{F}_0$
par un coproduit fini de faisceaux représentés par des objets quasi-compacts
et quasi-séparés de $(X_0)_\etale$.
Il suffit donc d'établir ceci : si $U_0 \to X_0$ est un tel objet
de $(X_0)_\etale$, alors
$$
(f_0^{-1}\mathcal{G}_0)(U) = \colim_{i \geq 0} (f_{i0}^{-1}\mathcal{G}_0)(U_i)
$$
où $U = X \times_{X_0} U_0$ et $U_i = X_i \times_{X_0} U_0$.
C'est un cas particulier du théorème \ref{etale-cohomology-theorem-colimit}.

\medskip\noindent
Surjectivité essentielle. Il faut montrer que tout faisceau constructible $\mathcal{F}$
sur $X$ est isomorphe à $f_i^{-1}\mathcal{F}_i$ pour un faisceau constructible
$\mathcal{F}_i$ sur $X_i$. En appliquant le
lemme \ref{etale-cohomology-lemma-category-constructible-sets}
et en utilisant les résultats des deux paragraphes précédents, on voit ceci :
il suffit de le démontrer pour $h_U$, pour un certain objet quasi-compact
et quasi-séparé $U$ de $X_\etale$.
Dans ce cas, il faut montrer que $U$ s'obtient par changement de base à partir d'un
schéma quasi-compact et quasi-séparé \'etale sur $X_i$ pour un certain $i$.
Cela résulte des
lemmes \ref{limits-lemma-descend-finite-presentation} et
\ref{limits-lemma-descend-etale}.

\medskip\noindent
Démontrons (3). L'argument est très semblable à celui du cas des
faisceaux d'ensembles, mais on utilise le
lemme \ref{etale-cohomology-lemma-category-constructible-modules}
au lieu du
lemme \ref{etale-cohomology-lemma-category-constructible-sets}. Les détails sont omis.
L'assertion (2) résulte de (3), car tout faisceau abélien constructible
sur un schéma quasi-compact est un faisceau constructible de
$\mathbf{Z}/n\mathbf{Z}$-modules pour un certain $n$.
\end{proof}

\begin{lemma}
\label{etale-cohomology-lemma-category-loc-cst-is-colimit}
Soit $X = \lim_{i \in I} X_i$ la limite projective d'un système
projectif filtrant de schémas à morphismes de transition affines.
Nous supposons que $X_i$ est quasi-compact et quasi-séparé
pour tout $i \in I$.
\begin{enumerate}
\item La catégorie des faisceaux localement constants finis sur $X_\etale$
est la limite inductive des catégories des faisceaux localement constants finis
sur $(X_i)_\etale$.
\item La catégorie des faisceaux abéliens localement constants finis sur $X_\etale$
est la limite inductive des catégories des faisceaux abéliens localement constants finis
sur $(X_i)_\etale$.
\item Soit $\Lambda$ un anneau noethérien. La catégorie des faisceaux
localement constants de type fini de $\Lambda$-modules sur $X_\etale$
est la limite inductive des catégories des faisceaux localement constants
de type fini de $\Lambda$-modules sur $(X_i)_\etale$.
\end{enumerate}
\end{lemma}

\begin{proof}
D'après le lemme \ref{etale-cohomology-lemma-category-is-colimit}, le foncteur est dans chaque cas
pleinement fidèle. D'après le même lemme, il suffit, pour achever
la démonstration dans le cas (1), d'établir ceci : étant donné un faisceau constructible
$\mathcal{F}_i$ sur $X_i$ dont l'image inverse $\mathcal{F}$ sur $X$ est
localement constant fini, il existe $i' \geq i$ tel que
l'image inverse $\mathcal{F}_{i'}$ de $\mathcal{F}_i$ sur $X_{i'}$
soit localement constant fini. Par hypothèse, il existe un recouvrement \'etale
$\mathcal{U} = \{U_j \to X\}_{j \in J}$
tel que $\mathcal{F}|_{U_j} \cong \underline{S_j}$ pour un ensemble fini $S_j$.
On peut supposer $U_j$ affine pour tout $j \in J$.
Comme $X$ est quasi-compact, on peut supposer $J$ fini.
D'après le lemme \ref{etale-cohomology-lemma-colimit-affine-sites}, il existe $i' \geq i$
et un recouvrement \'etale $\mathcal{U}_{i'} = \{U_{i', j} \to X_{i'}\}_{j \in J}$
dont le changement de base à $X$ est $\mathcal{U}$. Alors
$\mathcal{F}_{i'}|_{U_{i', j}}$ et $\underline{S_j}$ sont des
faisceaux constructibles sur $(U_{i', j})_\etale$ dont les images inverses sur
$U_j$ sont isomorphes. Ainsi, quitte à augmenter $i'$, on obtient
que $\mathcal{F}_{i'}|_{U_{i', j}}$ et $\underline{S_j}$
sont isomorphes. Par conséquent, $\mathcal{F}_{i'}$ est localement constant fini.
La démonstration des cas (2) et (3) est exactement la même.
\end{proof}

\begin{lemma}
\label{etale-cohomology-lemma-irreducible-subsheaf-constant-zero}
Soit $X$ un schéma irréductible de point générique $\eta$.
\begin{enumerate}
\item Soit $S' \subset S$ une inclusion d'ensembles. Si l'on a
$\underline{S'} \subset \mathcal{G} \subset \underline{S}$
dans $\Sh(X_\etale)$ et si $S' = \mathcal{G}_{\overline{\eta}}$, alors
$\mathcal{G} = \underline{S'}$.
\item Soit $A' \subset A$ une inclusion de groupes abéliens. Si l'on a
$\underline{A'} \subset \mathcal{G} \subset \underline{A}$
dans $\textit{Ab}(X_\etale)$ et si $A' = \mathcal{G}_{\overline{\eta}}$, alors
$\mathcal{G} = \underline{A'}$.
\item Soit $M' \subset M$ une inclusion de modules sur un anneau $\Lambda$.
Si l'on a $\underline{M'} \subset \mathcal{G} \subset \underline{M}$
dans $\textit{Mod}(X_\etale, \underline{\Lambda})$
et si $M' = \mathcal{G}_{\overline{\eta}}$, alors
$\mathcal{G} = \underline{M'}$.
\end{enumerate}
\end{lemma}

\begin{proof}
En effet, pour tout morphisme \'etale $U \to X$
tel que $U \not = \emptyset$, le point $\eta$ appartient à l'image.
\end{proof}

\begin{lemma}
\label{etale-cohomology-lemma-push-constant-sheaf-from-open}
Soit $X$ un schéma intègre normal de corps des fonctions $K$.
Soit $E$ un ensemble.
\begin{enumerate}
\item Soit $g : \Spec(K) \to X$ l'inclusion du point générique.
Alors $g_*\underline{E} = \underline{E}$.
\item Soit $j : U \to X$ l'inclusion d'un ouvert non vide. Alors
$j_*\underline{E} = \underline{E}$.
\end{enumerate}
\end{lemma}

\begin{proof}
Démontrons (1). Soit $x \in X$ un point. Soit
$\mathcal{O}^{sh}_{X, \overline{x}}$
un hensélisé strict de $\mathcal{O}_{X, x}$. 
D'après Compléments d'algèbre, lemme \ref{more-algebra-lemma-henselization-normal},
on voit que $\mathcal{O}^{sh}_{X, \overline{x}}$ est un anneau intègre normal.
Ainsi, $\Spec(K) \times_X \Spec(\mathcal{O}^{sh}_{X, \overline{x}})$
est irréductible. Il s'ensuit que
la fibre $(g_*\underline{E})_{\overline{x}}$ est égale à $E$,
voir le théorème \ref{etale-cohomology-theorem-higher-direct-images}.

\medskip\noindent
Démontrons (2). Comme $g$ se factorise par $j$, il existe un morphisme
$j_*\underline{E} \to g_*\underline{E}$. Ce morphisme est injectif, car
pour tout schéma $V$ \'etale sur $X$, l'ensemble $\Spec(K) \times_X V$
est dense dans $U \times_X V$. D'autre part, il existe un morphisme
$\underline{E} \to j_*\underline{E}$, et on conclut.
\end{proof}

\begin{lemma}
\label{etale-cohomology-lemma-zero-in-generic-point}
Soit $X$ un schéma quasi-compact et quasi-séparé. Soit
$\eta \in X$ un point générique d'une composante irréductible de $X$.
\begin{enumerate}
\item Soit $\mathcal{F}$ un faisceau abélien de torsion sur $X_\etale$
dont la fibre $\mathcal{F}_{\overline{\eta}}$ est nulle.
Alors $\mathcal{F} = \colim \mathcal{F}_i$ est une limite inductive filtrante de
faisceaux abéliens constructibles $\mathcal{F}_i$ tels que, pour tout $i$,
le support de $\mathcal{F}_i$ soit contenu
dans un sous-schéma fermé ne contenant pas $\eta$.
\item Soit $\Lambda$ un anneau noethérien et $\mathcal{F}$ un faisceau
de $\Lambda$-modules sur $X_\etale$ dont la fibre
$\mathcal{F}_{\overline{\eta}}$ est nulle. Alors
$\mathcal{F} = \colim \mathcal{F}_i$
est une limite inductive filtrante de faisceaux constructibles de
$\Lambda$-modules $\mathcal{F}_i$ tels que, pour tout $i$,
le support de $\mathcal{F}_i$ soit contenu dans un sous-schéma fermé
ne contenant pas $\eta$.
\end{enumerate}
\end{lemma}

\begin{proof}
Démontrons (1). On peut écrire $\mathcal{F} = \colim_{i \in I} \mathcal{F}_i$,
où $\mathcal{F}_i$ est un faisceau abélien constructible, d'après le
lemme \ref{etale-cohomology-lemma-torsion-colimit-constructible}.
Choisissons $i \in I$. Comme $\mathcal{F}|_\eta$ est nul par hypothèse, on
voit qu'il existe $i'(i) \geq i$ tel que
$\mathcal{F}_i|_\eta \to \mathcal{F}_{i'(i)}|_\eta$ soit nul, voir le
lemme \ref{etale-cohomology-lemma-colimit-constructible}.
Alors $\mathcal{G}_i = \Im(\mathcal{F}_i \to \mathcal{F}_{i'(i)})$
est un faisceau abélien constructible (lemme \ref{etale-cohomology-lemma-constructible-abelian})
dont la fibre en $\eta$ est nulle.
Par conséquent, le support $E_i$ de $\mathcal{G}_i$ est une partie constructible
de $X$ ne contenant pas $\eta$. Comme
$\eta$ est un point générique d'une composante irréductible de
$X$, on voit que $\eta \not \in Z_i = \overline{E_i}$ d'après
Topologie, lemme \ref{topology-lemma-generic-point-in-constructible}.
Définissons un nouvel ensemble filtrant $I'$ ayant pour ensemble sous-jacent $I$ et
pour relation d'ordre la règle suivante :
$i_1$ est supérieur ou égal à $i_2$ si et seulement si $i_1 = i_2\ \text{ou}\ i_1 \geq i'(i_2)$.
Alors les faisceaux $\mathcal{G}_i$ forment un système indexé par $I'$
de limite inductive $\mathcal{F}$, ce qui achève la démonstration.

\medskip\noindent
La démonstration du cas (2) est exactement la même et nous l'omettons.
\end{proof}








\section{Faisceaux constructibles sur les schémas noethériens}
\label{etale-cohomology-section-noetherian-constructible}

\noindent
Si $X$ est un schéma noethérien, toute partie localement fermée est une
partie localement fermée constructible
(Topologie, lemme \ref{topology-lemma-constructible-Noetherian-space}).
Par conséquent, un faisceau abélien $\mathcal{F}$ sur $X_\etale$ est constructible
si et seulement s'il existe une partition finie $X = \coprod X_i$
telle que $\mathcal{F}|_{X_i}$ soit localement constant fini.
(Par convention, les parties d'une partition d'un espace topologique sont
localement fermées, voir Topologie, section \ref{topology-section-stratifications}.)
Autrement dit, on peut omettre l'adjectif ``constructible'' dans la
définition \ref{etale-cohomology-definition-constructible}. En fait, la catégorie
des faisceaux constructibles sur les
schémas noethériens possède quelques propriétés supplémentaires que nous
répertorions dans cette section.

\begin{proposition}
\label{etale-cohomology-proposition-constructible-over-noetherian}
Soit $X$ un schéma noethérien. Soit $\Lambda$ un anneau noethérien.
\begin{enumerate}
\item Tout sous-faisceau ou faisceau quotient d'un faisceau constructible
d'ensembles est constructible.
\item La catégorie des faisceaux abéliens constructibles sur $X_\etale$ est une
sous-catégorie de Serre (forte) de $\textit{Ab}(X_\etale)$. En particulier,
tout sous-faisceau et tout faisceau quotient d'un faisceau abélien constructible
sur $X_\etale$ sont constructibles.
\item La catégorie des faisceaux constructibles de $\Lambda$-modules
sur $X_\etale$ est une sous-catégorie de Serre (forte) de
$\textit{Mod}(X_\etale, \Lambda)$. En particulier, tout sous-module
et tout module quotient d'un faisceau constructible de $\Lambda$-modules
sur $X_\etale$ sont constructibles.
\end{enumerate}
\end{proposition}

\begin{proof}
Démontrons (1). Soit $\mathcal{G} \subset \mathcal{F}$, où $\mathcal{F}$ est
un faisceau constructible d'ensembles sur $X_\etale$. Soit $\eta \in X$ un point
générique d'une composante irréductible de $X$. Par récurrence noethérienne,
il suffit de trouver un voisinage ouvert $U$ de $\eta$ tel que
$\mathcal{G}|_U$ soit localement constant. Pour cela, on peut remplacer $X$
par un voisinage \'etale de $\eta$.
On peut donc supposer que $\mathcal{F}$ est constant et que $X$ est irréductible.

\medskip\noindent
Écrivons $\mathcal{F} = \underline{S}$ pour un ensemble fini $S$.
Alors $S' = \mathcal{G}_{\overline{\eta}} \subset S$;
écrivons $S' = \{s_1, \ldots, s_t\}$.
Choisissons un voisinage \'etale $(U, \overline{u})$ de $\overline{\eta}$
et des sections $\sigma_1, \ldots, \sigma_t \in \mathcal{G}(U)$ qui ont pour images
les $s_i$ dans $\mathcal{G}_{\overline{\eta}} \subset S$.
Comme $\sigma_i$ a pour image l'élément
$s_i \in S' \subset S = \Gamma(X, \mathcal{F})$,
on voit que les deux images inverses de $\sigma_i$ sur $U \times_X U$
sont égales en tant que sections de $\mathcal{G}$. Par la condition de faisceau
pour $\mathcal{G}$, on en déduit que $\sigma_i$ provient d'une section
de $\mathcal{G}$ sur l'ouvert $\Im(U \to X)$ de $X$.
Quitte à rétrécir $X$, on peut supposer que
$\underline{S'} \subset \mathcal{G} \subset \underline{S}$.
Alors $\underline{S'} = \mathcal{G}$ d'après le
lemme \ref{etale-cohomology-lemma-irreducible-subsheaf-constant-zero}.

\medskip\noindent
Soit $\mathcal{F} \to \mathcal{Q}$ un morphisme surjectif, où $\mathcal{F}$ est
un faisceau constructible d'ensembles sur $X_\etale$. Posons
$\mathcal{G} = \mathcal{F} \times_\mathcal{Q} \mathcal{F}$.
La première partie de la démonstration montre que $\mathcal{G}$ est
constructible comme sous-faisceau de $\mathcal{F} \times \mathcal{F}$.
Il en résulte que $\mathcal{Q}$ est constructible, voir le
lemme \ref{etale-cohomology-lemma-constructible-abelian}.

\medskip\noindent
Démontrons (3). Nous savons déjà que les faisceaux constructibles de modules
forment une sous-catégorie de Serre faible, voir le lemme \ref{etale-cohomology-lemma-constructible-abelian}.
Il suffit donc d'établir l'assertion concernant les sous-modules.

\medskip\noindent
Soit $\mathcal{G} \subset \mathcal{F}$ un sous-module d'un
faisceau constructible de $\Lambda$-modules sur $X_\etale$. Soit $\eta \in X$
un point générique d'une composante irréductible de $X$. Par récurrence noethérienne,
il suffit de trouver un voisinage ouvert $U$ de $\eta$ tel que
$\mathcal{G}|_U$ soit localement constant. Pour cela, on peut remplacer $X$
par un voisinage \'etale de $\eta$. On peut donc supposer que $\mathcal{F}$
est constant et que $X$ est irréductible.

\medskip\noindent
Écrivons $\mathcal{F} = \underline{M}$ pour un $\Lambda$-module de type fini $M$.
Alors $M' = \mathcal{G}_{\overline{\eta}} \subset M$. Choisissons un nombre fini
d'éléments $s_1, \ldots, s_t$ qui engendrent $M'$ comme $\Lambda$-module.
(C'est possible puisque $\Lambda$ est noethérien et que $M$ est de type fini.)
Choisissons un voisinage \'etale $(U, \overline{u})$ de $\overline{\eta}$
et des sections $\sigma_1, \ldots, \sigma_t \in \mathcal{G}(U)$ qui ont pour images
les $s_i$ dans $\mathcal{G}_{\overline{\eta}} \subset M$.
Comme $\sigma_i$ a pour image l'élément
$s_i \in M' \subset M = \Gamma(X, \mathcal{F})$,
on voit que les deux images inverses de $\sigma_i$ sur $U \times_X U$
sont égales en tant que sections de $\mathcal{G}$. Par la condition de faisceau
pour $\mathcal{G}$, on en déduit que $\sigma_i$ provient d'une section
de $\mathcal{G}$ sur l'ouvert $\Im(U \to X)$ de $X$.
Quitte à rétrécir $X$, on peut supposer que
$\underline{M'} \subset \mathcal{G} \subset \underline{M}$.
Alors $\underline{M'} = \mathcal{G}$ d'après le
lemme \ref{etale-cohomology-lemma-irreducible-subsheaf-constant-zero}.

\medskip\noindent
Démontrons (2). Cela se déduit de (3) de la manière habituelle. Nous omettons
les détails.
\end{proof}

\noindent
Le lemme suivant affirme que tout objet de la catégorie abélienne des
faisceaux constructibles sur $X$ est ``noethérien'', c'est-à-dire satisfait
à la condition de chaîne ascendante pour les sous-objets.

\begin{lemma}
\label{etale-cohomology-lemma-constructible-over-noetherian-noetherian}
Soit $X$ un schéma noethérien. Soit $\Lambda$ un anneau noethérien.
Considérons des inclusions
$$
\mathcal{F}_1 \subset \mathcal{F}_2 \subset \mathcal{F}_3 \subset \ldots
\subset \mathcal{F}
$$
dans la catégorie des faisceaux d'ensembles, de groupes abéliens ou de $\Lambda$-modules.
Si $\mathcal{F}$ est constructible, alors il existe $n$ tel que
l'on ait $\mathcal{F}_n = \mathcal{F}_{n + 1} = \mathcal{F}_{n + 2} = \ldots$.
\end{lemma}

\begin{proof}
D'après la proposition \ref{etale-cohomology-proposition-constructible-over-noetherian},
les $\mathcal{F}_i$ et $\colim \mathcal{F}_i$ sont constructibles.
Le lemme résulte alors du
lemme \ref{etale-cohomology-lemma-colimit-constructible}.
\end{proof}

\begin{lemma}
\label{etale-cohomology-lemma-constructible-maps-into-constant}
Soit $X$ un schéma noethérien.
\begin{enumerate}
\item Soit $\mathcal{F}$ un faisceau constructible d'ensembles sur $X_\etale$.
Il existe un monomorphisme de faisceaux
$$
\mathcal{F} \longrightarrow
\prod\nolimits_{i = 1}^n f_{i, *}\underline{E_i}
$$
où $f_i : Y_i \to X$ est un morphisme fini et $E_i$ un ensemble fini.
\item Soit $\mathcal{F}$ un faisceau abélien constructible sur $X_\etale$.
Il existe un monomorphisme de faisceaux abéliens
$$
\mathcal{F} \longrightarrow
\bigoplus\nolimits_{i = 1}^n f_{i, *}\underline{M_i}
$$
où $f_i : Y_i \to X$ est un morphisme fini et où
$M_i$ est un groupe abélien fini.
\item Soit $\Lambda$ un anneau noethérien.
Soit $\mathcal{F}$ un faisceau constructible de $\Lambda$-modules sur $X_\etale$.
Il existe un monomorphisme de faisceaux de modules
$$
\mathcal{F} \longrightarrow
\bigoplus\nolimits_{i = 1}^n f_{i, *}\underline{M_i}
$$
où $f_i : Y_i \to X$ est un morphisme fini et où
$M_i$ est un $\Lambda$-module de type fini.
\end{enumerate}
De plus, on peut supposer que chaque $Y_i$ est irréductible et réduit, et s'envoie
surjectivement sur un sous-schéma fermé irréductible et réduit $Z_i \subset X$ tel que
$Y_i \to Z_i$ soit fini \'etale au-dessus d'un ouvert non vide de $Z_i$.
\end{lemma}

\begin{proof}
Démontrons (1). Comme on dispose de la condition de chaîne ascendante pour les
sous-faisceaux de $\mathcal{F}$
(lemme \ref{etale-cohomology-lemma-constructible-over-noetherian-noetherian}), il
suffit de montrer que, pour tout point $x \in X$, on
peut trouver un morphisme $\varphi : \mathcal{F} \to f_*\underline{E}$, où
$f : Y \to X$ est fini et $E$ est un ensemble fini, tel que
$\varphi_{\overline{x}} : \mathcal{F}_{\overline{x}} \to
(f_*\underline{E})_{\overline{x}}$ soit injectif.
(On peut éviter cet argument en choisissant une partition de $X$ comme dans le
lemme \ref{etale-cohomology-lemma-constructible-quasi-compact-quasi-separated}
et en construisant un $Y_i \to X$ pour chaque composante irréductible
de chaque partie.)
Soit $Z \subset X$ le sous-schéma réduit induit
(Schémas, définition \ref{schemes-definition-reduced-induced-scheme})
sur $\overline{\{x\}}$.
Comme $\mathcal{F}$ est constructible, il existe une extension finie séparable
$K/\kappa(x)$ telle que
$\mathcal{F}|_{\Spec(K)}$ soit le faisceau constant de valeur $E$
pour un ensemble fini $E$. Soit $h : Y \to Z$ la normalisation
de $Z$ dans $\Spec(K)$, et notons $f : Y \to X$ le composé de $h$ avec $Z \to X$.
D'après Morphismes, lemme \ref{morphisms-lemma-normal-normalization},
on voit que $Y$ est un schéma intègre normal.
Comme $K/\kappa(x)$ est une extension finie, il est clair que $K$ est le corps
des fonctions de $Y$. Notons $g : \Spec(K) \to Y$ l'inclusion.
Le morphisme $\mathcal{F}|_{\Spec(K)} \to \underline{E}$ correspond par adjonction
à un morphisme $\mathcal{F}|_Y \to g_*\underline{E} = \underline{E}$
(lemme \ref{etale-cohomology-lemma-push-constant-sheaf-from-open}).
Celui-ci correspond à son tour par adjonction à un morphisme
$\varphi : \mathcal{F} \to f_*\underline{E}$.
Observons que le morphisme induit par $\varphi$ sur la fibre en un point géométrique
$\overline{x}$ est injectif : on peut choisir un relèvement $\overline{y} \in Y$
de $\overline{x}$, et le diagramme commutatif
$$
\xymatrix{
\mathcal{F}_{\overline{x}} \ar@{=}[r] \ar[d] &
(\mathcal{F}|_Y)_{\overline{y}} \ar@{=}[d] \\
(f_*\underline{E})_{\overline{x}} \ar[r] &
\underline{E}_{\overline{y}}
}
$$
prouve l'injectivité. Cependant, la démonstration n'est pas encore achevée, car le
morphisme $h : Y \to Z$ est entier, mais n'est pas fini en
général\footnote{Si $X$ est un schéma de Nagata, par exemple de type fini
sur un corps, alors $Y \to Z$ est fini.}.

\medskip\noindent
Pour résoudre le problème signalé dans la dernière phrase du paragraphe précédent,
écrivons $Y = \lim_{i \in I} Y_i$, où $Y_i$ est irréductible, intègre et
fini sur $Z$; notons $f_i : Y_i \to X$ les morphismes composés. Appliquons Properties, lemme
\ref{properties-lemma-integral-algebra-directed-colimit-finite},
à $h_*\mathcal{O}_Y$, considéré comme un faisceau de $\mathcal{O}_Z$-algèbres,
puis appliquons le foncteur $\underline{\Spec}_Z$.
Alors $f_*\underline{E} = \colim f_{i, *}\underline{E}$
d'après le lemme \ref{etale-cohomology-lemma-relative-colimit}.
D'après le lemme \ref{etale-cohomology-lemma-constructible-is-compact}, le morphisme
$\mathcal{F} \to f_*\underline{E}$
se factorise par $f_{i, *}\underline{E}$ pour un certain $i$.
Comme $Y_i \to Z$ est un morphisme fini entre schémas intègres
et que l'extension de corps de fonctions
induite par ce morphisme est finie séparable, on voit que le
morphisme est fini \'etale au-dessus d'un ouvert non vide de $Z$ (utiliser
Algèbre, lemme \ref{algebra-lemma-smooth-at-generic-point}; détails omis).
Ceci achève la démonstration de (1).

\medskip\noindent
Les démonstrations de (2) et (3) sont identiques à celle de (1).
\end{proof}

\noindent
Dans le lemme suivant, nous employons un procédé standard pour ramener une
assertion très générale au cas noethérien.

\begin{lemma}
\label{etale-cohomology-lemma-constructible-maps-into-constant-general}
\begin{reference}
\cite[Expos\'e IX, Proposition 2.14]{SGA4}
\end{reference}
Soit $X$ un schéma quasi-compact et quasi-séparé.
\begin{enumerate}
\item Soit $\mathcal{F}$ un faisceau constructible d'ensembles sur $X_\etale$.
Il existe un monomorphisme de faisceaux
$$
\mathcal{F} \longrightarrow
\prod\nolimits_{i = 1}^n f_{i, *}\underline{E_i}
$$
où $f_i : Y_i \to X$ est un morphisme fini et de présentation finie et où
$E_i$ est un ensemble fini.
\item Soit $\mathcal{F}$ un faisceau abélien constructible sur $X_\etale$.
Il existe un monomorphisme de faisceaux abéliens
$$
\mathcal{F} \longrightarrow
\bigoplus\nolimits_{i = 1}^n f_{i, *}\underline{M_i}
$$
où $f_i : Y_i \to X$ est un morphisme fini et de présentation finie et où
$M_i$ est un groupe abélien fini.
\item Soit $\Lambda$ un anneau noethérien.
Soit $\mathcal{F}$ un faisceau constructible de $\Lambda$-modules sur $X_\etale$.
Il existe un monomorphisme de faisceaux de modules
$$
\mathcal{F} \longrightarrow
\bigoplus\nolimits_{i = 1}^n f_{i, *}\underline{M_i}
$$
où $f_i : Y_i \to X$ est un morphisme fini et de présentation finie et où
$M_i$ est un $\Lambda$-module de type fini.
\end{enumerate}
\end{lemma}

\begin{proof}
Nous allons ramener ce lemme au cas noethérien par approximation noethérienne
absolue. Plus précisément, d'après la
proposition \ref{limits-proposition-approximate} de Limits,
on peut écrire $X = \lim_{t \in T} X_t$, où chaque $X_t$ est de type fini sur
$\Spec(\mathbf{Z})$ et les morphismes de transition sont affines. D'après le
lemme \ref{etale-cohomology-lemma-category-is-colimit},
la catégorie des faisceaux constructibles (d'ensembles, de groupes abéliens ou de
$\Lambda$-modules) sur $X_\etale$ est la limite inductive des
catégories correspondantes sur $X_t$. Ainsi, notre faisceau constructible $\mathcal{F}$
est l'image inverse d'un faisceau constructible du même type $\mathcal{F}_t$
sur $X_t$ pour un certain $t$. Appliquons alors le cas noethérien
(lemme \ref{etale-cohomology-lemma-constructible-maps-into-constant})
pour trouver un monomorphisme
$$
\mathcal{F}_t \longrightarrow
\prod\nolimits_{i = 1}^n f_{i, *}\underline{E_i}
\quad\text{ou}\quad
\mathcal{F}_t \longrightarrow
\bigoplus\nolimits_{i = 1}^n f_{i, *}\underline{M_i}
$$
sur $X_t$, pour certains morphismes finis $f_i : Y_i \to X_t$.
Comme $X_t$ est noethérien, les morphismes $f_i$ sont de présentation finie.
Comme l'image inverse est exacte et que la formation de $f_{i, *}$ commute
au changement de base
(lemme \ref{etale-cohomology-lemma-finite-pushforward-commutes-with-base-change}), on conclut.
\end{proof}

\begin{lemma}
\label{etale-cohomology-lemma-support-in-subset}
Soit $X$ un schéma noethérien. Soit $E \subset X$ une partie stable
par spécialisation.
\begin{enumerate}
\item Soit $\mathcal{F}$ un faisceau abélien de torsion sur $X_\etale$
dont le support est contenu dans $E$. Alors $\mathcal{F} = \colim \mathcal{F}_i$
est une limite inductive filtrante de faisceaux abéliens constructibles $\mathcal{F}_i$
tels que, pour tout $i$, le support de $\mathcal{F}_i$ soit contenu dans
une partie fermée elle-même contenue dans $E$.
\item Soient $\Lambda$ un anneau noethérien et $\mathcal{F}$ un faisceau
de $\Lambda$-modules sur $X_\etale$ dont le support est contenu dans $E$.
Alors $\mathcal{F} = \colim \mathcal{F}_i$
est une limite inductive filtrante de faisceaux constructibles de
$\Lambda$-modules $\mathcal{F}_i$ tels que, pour tout $i$,
le support de $\mathcal{F}_i$ soit contenu dans une partie fermée
elle-même contenue dans $E$.
\end{enumerate}
\end{lemma}

\begin{proof}
Démontrons (1). On peut écrire $\mathcal{F} = \colim_{i \in I} \mathcal{F}_i$,
où $\mathcal{F}_i$ est un faisceau abélien constructible, d'après le
lemme \ref{etale-cohomology-lemma-torsion-colimit-constructible}.
D'après la proposition \ref{etale-cohomology-proposition-constructible-over-noetherian},
l'image $\mathcal{F}'_i \subset \mathcal{F}$
du morphisme $\mathcal{F}_i \to \mathcal{F}$ est constructible.
Ainsi, $\mathcal{F} = \colim \mathcal{F}'_i$, et le support
de $\mathcal{F}'_i$ est contenu dans $E$.
Comme le support de $\mathcal{F}'_i$ est constructible
(par notre définition des faisceaux constructibles), on voit
que son adhérence est aussi contenue dans $E$, voir par exemple
Topologie, lemme
\ref{topology-lemma-constructible-stable-specialization-closed}.

\medskip\noindent
La démonstration du cas (2) est exactement la même et nous l'omettons.
\end{proof}








\section{Spécialisations et faisceaux étales}
\label{etale-cohomology-section-specialization}

\noindent
Image topologique : soit $X$ un espace topologique et soit $x' \leadsto x$
une spécialisation de points de $X$. Alors tout voisinage ouvert de $x$
contient $x'$. Par conséquent, pour tout faisceau $\mathcal{F}$ sur $X$, il existe un
{\it homomorphisme de spécialisation}
$$
sp : \mathcal{F}_x \longrightarrow \mathcal{F}_{x'}
$$
entre les fibres, qui envoie la classe d'équivalence du couple $(U, s)$ dans
$\mathcal{F}_x$ sur la classe d'équivalence du couple $(U, s)$ dans
$\mathcal{F}_{x'}$; voir Faisceaux, section \ref{sheaves-section-stalks},
pour la description des fibres en termes de classes d'équivalence de couples.
Cette application est bien entendu fonctorielle en $\mathcal{F}$, autrement dit $sp$
est une transformation de foncteurs.

\medskip\noindent
Pour les faisceaux de la topologie étale, on peut imiter cette construction; voir
\cite[Expos\'e VIII, 7.7, page 397]{SGA4}. Supposons donnés
un schéma $S$, un point géométrique $\overline{s}$ de $S$ et un point
géométrique $\overline{t}$ de $\Spec(\mathcal{O}^{sh}_{S, \overline{s}})$.
Pour tout faisceau $\mathcal{F}$ sur $S_\etale$, nous allons construire
l'{\it homomorphisme de spécialisation}
$$
sp : \mathcal{F}_{\overline{s}} \longrightarrow \mathcal{F}_{\overline{t}}
$$
Nous avons ici commis un abus de langage : au lieu d'écrire
$\mathcal{F}_{\overline{t}}$, il faudrait écrire $\mathcal{F}_{p(\overline{t})}$,
où $p : \Spec(\mathcal{O}^{sh}_{S, \overline{s}}) \to S$ est le morphisme
canonique. Rappelons que
$$
\mathcal{F}_{\overline{s}} = \colim_{(U, \overline{u})} \mathcal{F}(U)
$$
où la limite inductive porte sur tous les voisinages étales $(U, \overline{u})$
de $(S, \overline{s})$; voir section \ref{etale-cohomology-section-stalks}. Comme
$\mathcal{O}^{sh}_{S, \overline{s}}$ est la fibre du
faisceau structural, tout voisinage étale
$(U, \overline{u})$ de $(S, \overline{s})$ fournit un homomorphisme canonique
$\mathcal{O}_{U, u} \to \mathcal{O}^{sh}_{S, \overline{s}}$.
On obtient donc une unique factorisation
$$
\Spec(\mathcal{O}^{sh}_{S, \overline{s}}) \to U \to S
$$
Si $\overline{v}$ désigne l'image de $\overline{t}$ dans $U$, alors
$(U, \overline{v})$ est un voisinage étale de $(S, \overline{t})$.
Cette construction définit un foncteur de la catégorie des voisinages étales
de $(S, \overline{s})$ vers la catégorie des voisinages étales de
$(S, \overline{t})$. On peut donc définir l'application
$sp : \mathcal{F}_{\overline{s}} \to \mathcal{F}_{\overline{t}}$
en envoyant la classe d'équivalence de
$(U, \overline{u}, \sigma)$, où $\sigma \in \mathcal{F}(U)$,
sur la classe d'équivalence de $(U, \overline{v}, \sigma)$.

\medskip\noindent
Soit $K \in D(S_\etale)$. Pour $\overline{s}$ et $\overline{t}$
comme ci-dessus, on dispose de l'{\it homomorphisme de spécialisation}
$$
sp : K_{\overline{s}} \longrightarrow K_{\overline{t}}
\quad\text{dans}\quad D(\textit{Ab})
$$
En effet, si $K$ est représenté par le complexe $\mathcal{F}^\bullet$
de faisceaux abéliens, on prend simplement l'application
$$
K_{\overline{s}} = \mathcal{F}^\bullet_{\overline{s}}
\longrightarrow
\mathcal{F}^\bullet_{\overline{t}} = K_{\overline{t}}
$$
qui est donnée terme à terme par les homomorphismes de spécialisation des faisceaux
construits ci-dessus. Elle ne dépend pas du choix du complexe
représentant $K$, par exactitude des foncteurs fibres (autrement dit,
prendre les fibres d'un complexe est bien défini sur la catégorie dérivée).

\medskip\noindent
La construction est manifestement fonctorielle en le faisceau $\mathcal{F}$ sur
$S_\etale$. Si l'on considère les foncteurs fibres comme des morphismes de topos
$\overline{s}, \overline{t} : \textit{Ensembles} \to \Sh(S_\etale)$, alors
on peut considérer $sp$ comme un $2$-morphisme
$$
\xymatrix{
\textit{Ensembles}
\rrtwocell^{\overline{t}}_{\overline{s}}{\ sp}
&
&
\Sh(S_\etale)
}
$$
de topos.

\begin{remark}[Autre description de sp]
\label{etale-cohomology-remark-alternative-sp}
Soient $S$, $\overline{s}$ et $\overline{t}$ comme ci-dessus.
Une autre manière de décrire l'homomorphisme de spécialisation consiste à utiliser
$$
\mathcal{F}_{\overline{s}} =
\Gamma(\Spec(\mathcal{O}^{sh}_{S, \overline{s}}), p^{-1}\mathcal{F})
\quad\text{et}\quad
\mathcal{F}_{\overline{t}} = \Gamma(\overline{t},
\overline{t}^{-1}p^{-1}\mathcal{F})
$$
La première égalité résulte du Théorème \ref{etale-cohomology-theorem-higher-direct-images}
appliqué à $\text{id}_S : S \to S$, et la seconde résulte
du Lemme \ref{etale-cohomology-lemma-stalk-pullback}. On peut alors considérer $sp$
comme l'application
$$
sp :
\mathcal{F}_{\overline{s}} =
\Gamma(\Spec(\mathcal{O}^{sh}_{S, \overline{s}}), p^{-1}\mathcal{F})
\xrightarrow{\text{image inverse par }\overline{t}}
\Gamma(\overline{t}, \overline{t}^{-1}p^{-1}\mathcal{F}) =
\mathcal{F}_{\overline{t}}
$$
\end{remark}

\begin{remark}[Encore une description de sp]
\label{etale-cohomology-remark-another-sp}
Soient $S$, $\overline{s}$ et $\overline{t}$ comme ci-dessus.
Une autre possibilité consiste à utiliser l'unique morphisme
$$
c :
\Spec(\mathcal{O}^{sh}_{S, \overline{t}})
\longrightarrow
\Spec(\mathcal{O}^{sh}_{S, \overline{s}})
$$
au-dessus de $S$ qui est compatible avec le morphisme donné
$\overline{t} \to \Spec(\mathcal{O}^{sh}_{S, \overline{s}})$
et le morphisme
$\overline{t} \to \Spec(\mathcal{O}^{sh}_{S, \overline{t}})$.
L'existence et l'unicité de la flèche affichée
résultent d'Algèbre, Lemme \ref{algebra-lemma-map-into-henselian-colimit},
appliqué à $\mathcal{O}_{S, s}$, $\mathcal{O}^{sh}_{S, \overline{t}}$ et
$\mathcal{O}^{sh}_{S, \overline{s}} \to \kappa(\overline{t})$.
On obtient
$$
sp :
\mathcal{F}_{\overline{s}} =
\Gamma(\Spec(\mathcal{O}^{sh}_{S, \overline{s}}), \mathcal{F})
\xrightarrow{\text{image inverse par }c}
\Gamma(\Spec(\mathcal{O}^{sh}_{S, \overline{t}}), \mathcal{F}) =
\mathcal{F}_{\overline{t}}
$$
(avec les conventions de notation évidentes).
En fait, ce procédé fonctionne également pour les objets $K$ de $D(S_\etale)$ :
l'homomorphisme de spécialisation de $K$ est égal à l'application
$$
sp :
K_{\overline{s}} =
R\Gamma(\Spec(\mathcal{O}^{sh}_{S, \overline{s}}), K)
\xrightarrow{\text{image inverse par }c}
R\Gamma(\Spec(\mathcal{O}^{sh}_{S, \overline{t}}), K) =
K_{\overline{t}}
$$
Les signes d'égalité sont valables, car la prise des sections globales sur
les schémas strictement henséliens
$\Spec(\mathcal{O}^{sh}_{S, \overline{s}})$ et
$\Spec(\mathcal{O}^{sh}_{S, \overline{t}})$
est exacte (et revient à prendre les fibres en $\overline{s}$ et $\overline{t}$);
aucune difficulté liée au fait que $K$ puisse être non borné
n'apparaît donc.
\end{remark}

\begin{remark}[Relèvement des spécialisations]
\label{etale-cohomology-remark-can-lift}
Soit $S$ un schéma et soit $t \leadsto s$ une spécialisation de
points de $S$. Choisissons des points géométriques $\overline{t}$ et $\overline{s}$
au-dessus de $t$ et de $s$. Comme $t$ correspond à un point de
$\Spec(\mathcal{O}_{S, s})$ d'après
Schémas, Lemme \ref{schemes-lemma-specialize-points},
et comme $\mathcal{O}_{S, s} \to \mathcal{O}^{sh}_{S, \overline{s}}$
est fidèlement plat, on peut trouver un point
$t' \in \Spec(\mathcal{O}^{sh}_{S, \overline{s}})$ qui s'envoie sur $t$.
Puisque $\Spec(\mathcal{O}^{sh}_{S, \overline{s}})$ est une limite de
schémas étales sur $S$, l'extension $\kappa(t')/\kappa(t)$
est algébrique séparable (mais généralement pas finie, bien entendu).
Comme $\kappa(\overline{t})$ est algébriquement clos, on peut
choisir un plongement $\kappa(t') \to \kappa(\overline{t})$
d'extensions de $\kappa(t)$. Ce choix fournit un diagramme
commutatif
$$
\xymatrix{
\overline{t} \ar[d] \ar[r] &
\Spec(\mathcal{O}^{sh}_{S, \overline{s}}) \ar[d] &
\overline{s} \ar[l] \ar[d] \\
t \ar[r] & S & s \ar[l]
}
$$
de points et de points géométriques. Ainsi, si $t \leadsto s$,
on peut toujours « relever » $\overline{t}$ en un point géométrique
du hensélisé strict de $S$ en $\overline{s}$
et obtenir les homomorphismes de spécialisation ci-dessus.
\end{remark}

\begin{lemma}
\label{etale-cohomology-lemma-specialization-map-pullback}
Soit $g : S' \to S$ un morphisme de schémas. Soit $\mathcal{F}$ un faisceau
sur $S_\etale$. Soit $\overline{s}'$ un point géométrique de $S'$, et soit
$\overline{t}'$ un point géométrique de
$\Spec(\mathcal{O}^{sh}_{S', \overline{s}'})$. Posons
$\overline{s} = g(\overline{s}')$ et $\overline{t} = h(\overline{t}')$,
où $h : \Spec(\mathcal{O}^{sh}_{S', \overline{s}'}) \to
\Spec(\mathcal{O}^{sh}_{S, \overline{s}})$ est le morphisme canonique.
Pour tout faisceau $\mathcal{F}$ sur $S_\etale$, l'homomorphisme de spécialisation
$$
sp :
(g^{-1}\mathcal{F})_{\overline{s}'}
\longrightarrow
(g^{-1}\mathcal{F})_{\overline{t}'}
$$
est égal à l'homomorphisme de spécialisation
$sp : \mathcal{F}_{\overline{s}} \to \mathcal{F}_{\overline{t}}$
au moyen des identifications
$(g^{-1}\mathcal{F})_{\overline{s}'} = \mathcal{F}_{\overline{s}}$ et
$(g^{-1}\mathcal{F})_{\overline{t}'} = \mathcal{F}_{\overline{t}}$
du Lemme \ref{etale-cohomology-lemma-stalk-pullback}.
\end{lemma}

\begin{proof}
Omis.
\end{proof}

\begin{lemma}
\label{etale-cohomology-lemma-characterize-locally-constant}
Soit $S$ un schéma tel que tout ouvert quasi-compact de $S$ possède
un nombre fini de composantes irréductibles (par exemple, si $S$ a un
espace topologique sous-jacent noethérien ou si $S$ est localement noethérien).
Soit $\mathcal{F}$ un faisceau d'ensembles sur $S_\etale$.
Les conditions suivantes sont équivalentes :
\begin{enumerate}
\item $\mathcal{F}$ est localement constant fini;
\item toutes les fibres de $\mathcal{F}$ sont des ensembles finis et tous les homomorphismes
de spécialisation $sp : \mathcal{F}_{\overline{s}} \to \mathcal{F}_{\overline{t}}$
sont bijectifs.
\end{enumerate}
\end{lemma}

\begin{proof}
Supposons (2). Soit $\overline{s}$ un point géométrique de $S$ au-dessus de
$s \in S$. Pour démontrer (1), il faut trouver un voisinage étale
$(U, \overline{u})$ de $(S, \overline{s})$ tel que $\mathcal{F}|_U$
soit constant. On peut et on va supposer $S$ affine.

\medskip\noindent
Comme $\mathcal{F}_{\overline{s}}$ est fini, on peut choisir
$(U, \overline{u})$, un entier $n \geq 0$ et des éléments deux à deux distincts
$\sigma_1, \ldots, \sigma_n \in \mathcal{F}(U)$
tels que $\{\sigma_1, \ldots, \sigma_n\} \subset \mathcal{F}(U)$ s'envoie
bijectivement sur $\mathcal{F}_{\overline{s}}$ par l'application
$\mathcal{F}(U) \to \mathcal{F}_{\overline{s}}$.
Considérons le morphisme
$$
\varphi : \underline{\{1, \ldots, n\}} \longrightarrow \mathcal{F}|_U
$$
sur $U_\etale$ défini par $\sigma_1, \ldots, \sigma_n$.
Ce morphisme induit une bijection sur les fibres en $\overline{u}$
par construction. Considérons la partie
$$
E = \{u' \in U \mid \varphi_{\overline{u}'}\text{ est bijectif}\} \subset U
$$
Ici $\overline{u}'$ est un point géométrique quelconque de $U$ au-dessus de $u'$
(la condition ne dépend pas du choix, par Remarque \ref{etale-cohomology-remark-map-stalks}).
L'image $u \in U$ de $\overline{u}$ appartient à $E$.
Notre hypothèse sur les homomorphismes de spécialisation de $\mathcal{F}$,
la Remarque \ref{etale-cohomology-remark-can-lift} et le
Lemme \ref{etale-cohomology-lemma-specialization-map-pullback}
montrent que $E$ est stable par spécialisation
et par générisation dans l'espace topologique $U$.

\medskip\noindent
Après avoir rétréci $U$, on peut supposer $U$ affine lui aussi. Par
Descente, Lemme \ref{descent-lemma-locally-finite-nr-irred-local-fppf},
$U$ possède un nombre fini de composantes irréductibles.
Après avoir retiré celles qui ne passent pas
par $u$, on peut supposer que toute composante irréductible de $U$
passe par $u$. Puisque $U$ est un espace topologique sobre,
il s'ensuit que $E = U$ et que $\varphi$ est un isomorphisme par le
Théorème \ref{etale-cohomology-theorem-exactness-stalks}. Ainsi, (1) est vérifiée.

\medskip\noindent
Nous omettons la démonstration du fait que (1) implique (2).
\end{proof}

\begin{lemma}
\label{etale-cohomology-lemma-characterize-locally-constant-module}
Soit $S$ un schéma tel que tout ouvert quasi-compact de $S$ possède
un nombre fini de composantes irréductibles (par exemple, si $S$ a un
espace topologique sous-jacent noethérien ou si $S$ est localement noethérien).
Soit $\Lambda$ un anneau noethérien.
Soit $\mathcal{F}$ un faisceau de $\Lambda$-modules sur $S_\etale$.
Les conditions suivantes sont équivalentes :
\begin{enumerate}
\item $\mathcal{F}$ est un faisceau localement constant de type fini
de $\Lambda$-modules;
\item toutes les fibres de $\mathcal{F}$ sont des $\Lambda$-modules de type fini et
tous les homomorphismes de spécialisation
$sp : \mathcal{F}_{\overline{s}} \to \mathcal{F}_{\overline{t}}$
sont bijectifs.
\end{enumerate}
\end{lemma}

\begin{proof}
La démonstration est la même que celle du
Lemme \ref{etale-cohomology-lemma-characterize-locally-constant}.
Supposons (2). Soit $\overline{s}$ un point géométrique de $S$ au-dessus de
$s \in S$. Pour démontrer (1), il faut trouver un voisinage étale
$(U, \overline{u})$ de $(S, \overline{s})$ tel que $\mathcal{F}|_U$
soit constant. On peut et on va supposer $S$ affine.

\medskip\noindent
Comme $M = \mathcal{F}_{\overline{s}}$ est un $\Lambda$-module de type fini
et que $\Lambda$ est noethérien, on peut choisir une présentation
$$
\Lambda^{\oplus m} \xrightarrow{A} \Lambda^{\oplus n} \to M \to 0
$$
pour une matrice $A = (a_{ji})$ à coefficients dans $\Lambda$.
On peut choisir $(U, \overline{u})$ et des éléments
$\sigma_1, \ldots, \sigma_n \in \mathcal{F}(U)$
tels que $\sum a_{ji}\sigma_i = 0$ dans $\mathcal{F}(U)$ et
que les images des $\sigma_i$ dans $\mathcal{F}_{\overline{s}} = M$
soient celles des éléments de la base canonique de
$\Lambda^n$ dans la présentation de $M$ donnée ci-dessus.
Considérons le morphisme
$$
\varphi : \underline{M} \longrightarrow \mathcal{F}|_U
$$
sur $U_\etale$ défini par $\sigma_1, \ldots, \sigma_n$.
Ce morphisme induit une bijection sur les fibres en $\overline{u}$
par construction. Considérons la partie
$$
E = \{u' \in U \mid \varphi_{\overline{u}'}\text{ est bijectif}\} \subset U
$$
Ici $\overline{u}'$ est un point géométrique quelconque de $U$ au-dessus de $u'$
(la condition ne dépend pas du choix, par Remarque \ref{etale-cohomology-remark-map-stalks}).
L'image $u \in U$ de $\overline{u}$ appartient à $E$.
Notre hypothèse sur les homomorphismes de spécialisation de $\mathcal{F}$,
la Remarque \ref{etale-cohomology-remark-can-lift} et le
Lemme \ref{etale-cohomology-lemma-specialization-map-pullback}
montrent que $E$ est stable par spécialisation
et par générisation dans l'espace topologique $U$.

\medskip\noindent
Après avoir rétréci $U$, on peut supposer $U$ affine lui aussi. Par
Descente, Lemme \ref{descent-lemma-locally-finite-nr-irred-local-fppf},
$U$ possède un nombre fini de composantes irréductibles.
Après avoir retiré celles qui ne passent pas
par $u$, on peut supposer que toute composante irréductible de $U$
passe par $u$. Puisque $U$ est un espace topologique sobre,
il s'ensuit que $E = U$ et que $\varphi$ est un isomorphisme par le
Théorème \ref{etale-cohomology-theorem-exactness-stalks}. Ainsi, (1) est vérifiée.

\medskip\noindent
Nous omettons la démonstration du fait que (1) implique (2).
\end{proof}

\begin{lemma}
\label{etale-cohomology-lemma-specialization-map-pushforward}
Soit $f : X \to S$ un morphisme de schémas quasi-compact et quasi-séparé.
Soit $K \in D^+(X_\etale)$. Soit $\overline{s}$ un point géométrique de $S$
et soit $\overline{t}$ un point géométrique de
$\Spec(\mathcal{O}^{sh}_{S, \overline{s}})$. On a un diagramme
commutatif
$$
\xymatrix{
(Rf_*K)_{\overline{s}} \ar[r]_{sp} \ar@{=}[d] &
(Rf_*K)_{\overline{t}} \ar@{=}[d] \\
R\Gamma(X \times_S \Spec(\mathcal{O}^{sh}_{S, \overline{s}}), K)
\ar[r] &
R\Gamma(X \times_S \Spec(\mathcal{O}^{sh}_{S, \overline{t}}), K)
}
$$
où la flèche horizontale inférieure est l'image inverse par le morphisme
$\text{id}_X \times c$, où
$c : \Spec(\mathcal{O}^{sh}_{S, \overline{t}}) \to
\Spec(\mathcal{O}^{sh}_{S, \overline{s}})$
est le morphisme introduit dans la Remarque \ref{etale-cohomology-remark-another-sp}.
Les flèches verticales sont données par le
Théorème \ref{etale-cohomology-theorem-higher-direct-images}.
\end{lemma}

\begin{proof}
Cela résulte immédiatement de la description
de $sp$ donnée dans la Remarque \ref{etale-cohomology-remark-another-sp}.
\end{proof}

\begin{remark}
\label{etale-cohomology-remark-specialization-map-and-fibres}
Soit $f : X \to S$ un morphisme de schémas. Soit $K \in D(X_\etale)$.
Soit $\overline{s}$ un point géométrique de $S$ et soit $\overline{t}$
un point géométrique de $\Spec(\mathcal{O}^{sh}_{S, \overline{s}})$.
Soit $c$ comme dans la Remarque \ref{etale-cohomology-remark-another-sp}.
On peut toujours former un diagramme commutatif
$$
\xymatrix{
(Rf_*K)_{\overline{s}} \ar[r] \ar[d]_{sp} &
R\Gamma(X \times_S \Spec(\mathcal{O}_{S, \overline{s}}^{sh}), K)
\ar[r] \ar[d]_{(\text{id}_X \times c)^{-1}} &
R\Gamma(X_{\overline{s}}, K) \\
(Rf_*K)_{\overline{t}} \ar[r] &
R\Gamma(X \times_S \Spec(\mathcal{O}_{S, \overline{t}}^{sh}), K) \ar[r] &
R\Gamma(X_{\overline{t}}, K)
}
$$
où les flèches horizontales sont celles de la Remarque \ref{etale-cohomology-remark-stalk-fibre}.
En général, il n'existe pas de flèche verticale à droite
entre les cohomologies de $K$ sur les fibres qui complète ce
diagramme, même dans le cas du Lemme \ref{etale-cohomology-lemma-specialization-map-pushforward}.
\end{remark}















\section{Complexes à cohomologie constructible}
\label{etale-cohomology-section-Dc}

\noindent
Soit $\Lambda$ un anneau. Notons $D(X_\etale, \Lambda)$ la catégorie dérivée
des faisceaux de $\Lambda$-modules sur $X_\etale$.
On note $D^b(X_\etale, \Lambda)$ (respectivement $D^+$, $D^-$)
la sous-catégorie pleine des complexes bornés (respectivement bornés inférieurement,
supérieurement) de $D(X_\etale, \Lambda)$.

\begin{definition}
\label{etale-cohomology-definition-c}
Soit $X$ un schéma. Soit $\Lambda$ un anneau noethérien.
On note {\it $D_c(X_\etale, \Lambda)$} la sous-catégorie pleine
de $D(X_\etale, \Lambda)$ formée des complexes dont les faisceaux de cohomologie
sont des faisceaux constructibles de $\Lambda$-modules.
\end{definition}

\noindent
Cette définition a un sens par le Lemme \ref{etale-cohomology-lemma-constructible-abelian} et
Catégories dérivées, section \ref{derived-section-triangulated-sub}.
Ainsi, $D_c(X_\etale, \Lambda)$ est une sous-catégorie triangulée strictement
pleine et saturée de $D(X_\etale, \Lambda)$.

\begin{lemma}
\label{etale-cohomology-lemma-restrict-and-shriek-from-etale-c}
Soit $\Lambda$ un anneau noethérien.
Si $j : U \to X$ est un morphisme étale de schémas, alors
\begin{enumerate}
\item $K|_U \in D_c(U_\etale, \Lambda)$ si $K \in D_c(X_\etale, \Lambda)$;
\item $j_!M \in D_c(X_\etale, \Lambda)$ si $M \in D_c(U_\etale, \Lambda)$ et si
le morphisme $j$ est quasi-compact et quasi-séparé.
\end{enumerate}
\end{lemma}

\begin{proof}
La première assertion est claire. La seconde résulte de l'exactitude
de $j_!$ et du Lemme \ref{etale-cohomology-lemma-jshriek-constructible}.
\end{proof}

\begin{lemma}
\label{etale-cohomology-lemma-pullback-c}
Soit $\Lambda$ un anneau noethérien.
Soit $f : X \to Y$ un morphisme de schémas. Si $K \in D_c(Y_\etale, \Lambda)$,
alors $Lf^*K \in D_c(X_\etale, \Lambda)$.
\end{lemma}

\begin{proof}
Cela résulte de l'exactitude de $f^{-1} = f^*$ et du
Lemme \ref{etale-cohomology-lemma-pullback-constructible}.
\end{proof}

\begin{lemma}
\label{etale-cohomology-lemma-one-constructible}
Soit $X$ un schéma quasi-compact et quasi-séparé.
Soit $\Lambda$ un anneau noethérien. Soient $K \in D(X_\etale, \Lambda)$
et $b \in \mathbf{Z}$ tels que $H^b(K)$ soit constructible.
Il existe alors un faisceau $\mathcal{F}$ qui est une somme directe finie
de faisceaux $j_{U!}\underline{\Lambda}$, où $U \in \Ob(X_\etale)$ est affine, et
un morphisme $\mathcal{F}[-b] \to K$ dans $D(X_\etale, \Lambda)$
qui induit une surjection $\mathcal{F} \to H^b(K)$.
\end{lemma}

\begin{proof}
Représentons $K$ par un complexe $\mathcal{K}^\bullet$ de faisceaux de
$\Lambda$-modules. Considérons la surjection
$$
\Ker(\mathcal{K}^b \to \mathcal{K}^{b + 1})
\longrightarrow
H^b(K)
$$
Par Modules sur les sites, Lemme
\ref{sites-modules-lemma-module-quotient-direct-sum}
on peut choisir une surjection
$\bigoplus_{i \in I} j_{U_i!} \underline{\Lambda} \to
\Ker(\mathcal{K}^b \to \mathcal{K}^{b + 1})$
avec $U_i$ affine. Pour toute partie finie $I' \subset I$, notons
$\mathcal{H}_{I'} \subset H^b(K)$ l'image de
$\bigoplus_{i \in I'} j_{U_i!} \underline{\Lambda}$. Par le
Lemme \ref{etale-cohomology-lemma-colimit-constructible},
on a $\mathcal{H}_{I'} = H^b(K)$ pour une partie finie $I' \subset I$.
Le lemme s'ensuit en prenant
$\mathcal{F} = \bigoplus_{i \in I'} j_{U_i!} \underline{\Lambda}$.
\end{proof}

\begin{lemma}
\label{etale-cohomology-lemma-bounded-above-c}
Soit $X$ un schéma quasi-compact et quasi-séparé.
Soit $\Lambda$ un anneau noethérien. Soit $K \in D^-(X_\etale, \Lambda)$. Alors
les conditions suivantes sont équivalentes :
\begin{enumerate}
\item $K$ appartient à $D_c(X_\etale, \Lambda)$;
\item $K$ peut être représenté par un complexe borné supérieurement
dont les termes sont des sommes directes finies de $j_{U!}\underline{\Lambda}$,
où $U \in \Ob(X_\etale)$ est affine;
\item $K$ peut être représenté par un complexe borné supérieurement
de faisceaux constructibles plats de $\Lambda$-modules.
\end{enumerate}
\end{lemma}

\begin{proof}
Il est clair que (2) implique (3) et que (3) implique (1).
Supposons que $K$ appartienne à $D_c^-(X_\etale, \Lambda)$.
Supposons que $H^i(K) = 0$ pour $i > b$. Par récurrence sur $a$,
nous allons construire un complexe $\mathcal{F}^a \to \ldots \to \mathcal{F}^b$
tel que chaque $\mathcal{F}^i$ soit une somme directe finie
de $j_{U!}\underline{\Lambda}$, avec $U \in \Ob(X_\etale)$ affine,
et un morphisme $\mathcal{F}^\bullet \to K$ qui induise un isomorphisme
$H^i(\mathcal{F}^\bullet) \to H^i(K)$ pour $i > a$ et une surjection
$H^a(\mathcal{F}^\bullet) \to H^a(K)$.
Pour $a = b$, cela résulte du Lemme \ref{etale-cohomology-lemma-one-constructible}.
Une telle donnée étant construite, choisissons un triangle distingué
$$
\mathcal{F}^\bullet \to K \to L \to \mathcal{F}^\bullet[1]
$$
On voit alors que $H^i(L) = 0$ pour $i \geq a$. Choisissons
$\mathcal{F}^{a - 1}[-a +1] \to L$ comme dans le
Lemme \ref{etale-cohomology-lemma-one-constructible}. Le composé
$\mathcal{F}^{a - 1}[-a +1] \to L \to \mathcal{F}^\bullet$
correspond à un morphisme $\mathcal{F}^{a - 1} \to \mathcal{F}^a$
dont le composé avec $\mathcal{F}^a \to \mathcal{F}^{a + 1}$
est nul. Par TR4, on obtient un morphisme
$$
(\mathcal{F}^{a - 1} \to \ldots \to \mathcal{F}^b) \to K
$$
dans $D(X_\etale, \Lambda)$. Cela achève l'étape de récurrence et la
démonstration du lemme.
\end{proof}

\begin{lemma}
\label{etale-cohomology-lemma-tensor-c}
Soit $X$ un schéma. Soit $\Lambda$ un anneau noethérien.
Soient $K, L \in D_c^-(X_\etale, \Lambda)$. Alors
$K \otimes_\Lambda^\mathbf{L} L$ appartient à $D_c^-(X_\etale, \Lambda)$.
\end{lemma}

\begin{proof}
Cela résulte des Lemmes \ref{etale-cohomology-lemma-bounded-above-c} et
\ref{etale-cohomology-lemma-tensor-product-constructible}.
\end{proof}




\section{Tor-dimension finie et cohomologie constructible}
\label{etale-cohomology-section-ctf}

\noindent
Soit $X$ un schéma et soit $\Lambda$ un anneau noethérien. Une sous-catégorie
fréquemment utilisée de la catégorie dérivée $D_c(X_\etale, \Lambda)$ définie
à la section \ref{etale-cohomology-section-Dc} est la sous-catégorie pleine formée des objets
de tor-dimension (localement) finie. En voici la définition formelle.

\begin{definition}
\label{etale-cohomology-definition-ctf}
Soit $X$ un schéma. Soit $\Lambda$ un anneau noethérien. On note
{\it $D_{ctf}(X_\etale, \Lambda)$} la sous-catégorie pleine de
$D_c(X_\etale, \Lambda)$
formée des objets de tor-dimension localement finie.
\end{definition}

\noindent
C'est une sous-catégorie triangulée strictement pleine et saturée de
$D_c(X_\etale, \Lambda)$ et de $D(X_\etale, \Lambda)$. D'après nos conventions, voir
Cohomologie sur les sites, Définition \ref{sites-cohomology-definition-tor-amplitude},
on a
$$
D_{ctf}(X_\etale, \Lambda) \subset D^b_c(X_\etale, \Lambda)
\subset D^b(X_\etale, \Lambda)
$$
si $X$ est quasi-compact. Le Lemme \ref{etale-cohomology-lemma-when-ctf} donne une bonne
description des objets de $D_{ctf}(X_\etale, \Lambda)$.

\begin{remark}
\label{etale-cohomology-remark-different}
Les objets de la catégorie dérivée $D_{ctf}(X_\etale, \Lambda)$ possèdent, en un
certain sens, de meilleures propriétés globales que les objets parfaits de $D(\mathcal{O}_X)$.
En effet, un complexe de $\mathcal{O}_X$-modules peut être
localement quasi-isomorphe à un complexe fini dont les termes sont
des $\mathcal{O}_X$-modules localement libres de type fini, sans être
globalement quasi-isomorphe à un complexe borné de
$\mathcal{O}_X$-modules localement libres. Le lemme suivant montre que cela ne se produit pas pour
$D_{ctf}$ sur un schéma noethérien.
\end{remark}

\begin{lemma}
\label{etale-cohomology-lemma-when-ctf}
Soit $\Lambda$ un anneau noethérien. Soit $X$ un schéma quasi-compact
et quasi-séparé. Soit $K \in D(X_\etale, \Lambda)$. Les conditions suivantes
sont équivalentes :
\begin{enumerate}
\item $K \in D_{ctf}(X_\etale, \Lambda)$;
\item $K$ peut être représenté par un complexe fini de faisceaux
constructibles plats de $\Lambda$-modules.
\end{enumerate}
En fait, si $K$ est d'amplitude de Tor contenue dans $[a, b]$, on peut représenter
$K$ par un complexe $\mathcal{F}^a \to \ldots \to \mathcal{F}^b$ où
$\mathcal{F}^p$ est un faisceau constructible plat de $\Lambda$-modules.
\end{lemma}

\begin{proof}
Il est clair qu'un complexe fini de faisceaux constructibles
plats de $\Lambda$-modules est de tor-dimension finie.
Il est également clair qu'il appartient à $D_c(X_\etale, \Lambda)$.
Ainsi, (2) implique (1).

\medskip\noindent
Supposons (1). Choisissons $a, b \in \mathbf{Z}$ tels que
$H^i(K \otimes_\Lambda^\mathbf{L} \mathcal{G}) = 0$ si
$i \not \in [a, b]$, pour tout faisceau de $\Lambda$-modules $\mathcal{G}$.
Nous démontrerons la dernière assertion par récurrence sur $b - a$. Si
$a = b$, alors $K = H^a(K)[-a]$ est un faisceau constructible plat
et le résultat est acquis. Supposons maintenant $b > a$. Représentons $K$
par un complexe $\mathcal{K}^\bullet$ de faisceaux de $\Lambda$-modules.
Considérons la surjection
$$
\Ker(\mathcal{K}^b \to \mathcal{K}^{b + 1})
\longrightarrow
H^b(K)
$$
Par le Lemme \ref{etale-cohomology-lemma-category-constructible-modules},
on peut trouver un nombre fini de schémas affines $U_i$ étales sur $X$ et une
surjection $\bigoplus j_{U_i!}\underline{\Lambda}_{U_i} \to H^b(K)$.
Après avoir remplacé les $U_i$ par des recouvrements étales standard $\{U_{ij} \to U_i\}$,
on peut supposer que cette surjection se relève en un morphisme
$\mathcal{F} = \bigoplus j_{U_i!}\underline{\Lambda}_{U_i} \to
\Ker(\mathcal{K}^b \to \mathcal{K}^{b + 1})$.
Ce morphisme détermine un triangle distingué
$$
\mathcal{F}[-b] \to K \to L \to \mathcal{F}[-b + 1]
$$
dans $D(X_\etale, \Lambda)$. Comme $D_{ctf}(X_\etale, \Lambda)$ est une
sous-catégorie triangulée, $L$ lui appartient également. En fait, $L$ est
d'amplitude de Tor contenue dans $[a, b - 1]$, puisque $\mathcal{F}$ se surjecte sur
$H^b(K)$ (nous omettons les détails). Par l'hypothèse de récurrence, il existe
un complexe fini $\mathcal{F}^a \to \ldots \to \mathcal{F}^{b - 1}$
de faisceaux constructibles plats de $\Lambda$-modules qui représente $L$.
Le morphisme $L \to \mathcal{F}[-b + 1]$ correspond à un morphisme
$\mathcal{F}^{b-1} \to \mathcal{F}$ qui annule l'image
de $\mathcal{F}^{b-2} \to \mathcal{F}^{b-1}$. Posons $\mathcal{F}^b=\mathcal{F}$; il résulte alors
de l'axiome TR3 que $K$ est représenté par le complexe
$$
\mathcal{F}^a \to \ldots \to \mathcal{F}^{b - 1} \to \mathcal{F}^b
$$
ce qui achève la démonstration.
\end{proof}

\begin{remark}
\label{etale-cohomology-remark-projective-each-degree}
Soit $\Lambda$ un anneau noethérien. Soit $X$ un schéma.
Pour un complexe borné $K^\bullet$ de faisceaux constructibles plats de $\Lambda$-modules
sur $X_\etale$,
chaque fibre $K^p_{\overline{x}}$ est un $\Lambda$-module projectif de type fini.
Les fibres du complexe sont donc des complexes parfaits de $\Lambda$-modules.
\end{remark}

\begin{lemma}
\label{etale-cohomology-lemma-restrict-and-shriek-from-etale-ctf}
Soit $\Lambda$ un anneau noethérien.
Si $j : U \to X$ est un morphisme étale de schémas, alors
\begin{enumerate}
\item $K|_U \in D_{ctf}(U_\etale, \Lambda)$ si
$K \in D_{ctf}(X_\etale, \Lambda)$;
\item $j_!M \in D_{ctf}(X_\etale, \Lambda)$ si
$M \in D_{ctf}(U_\etale, \Lambda)$ et si
le morphisme $j$ est quasi-compact et quasi-séparé.
\end{enumerate}
\end{lemma}

\begin{proof}
Le moyen le plus simple de démontrer ce lemme est peut-être de se ramener au
cas où $X$ est affine, puis d'appliquer le Lemme \ref{etale-cohomology-lemma-when-ctf}
pour le reformuler comme un énoncé sur les complexes finis
de faisceaux constructibles plats de $\Lambda$-modules;
le résultat découle alors du Lemme \ref{etale-cohomology-lemma-jshriek-constructible}.
\end{proof}

\begin{lemma}
\label{etale-cohomology-lemma-pullback-ctf}
Soit $\Lambda$ un anneau noethérien. Soit $f : X \to Y$ un morphisme de schémas.
Si $K \in D_{ctf}(Y_\etale, \Lambda)$, alors
$Lf^*K \in D_{ctf}(X_\etale, \Lambda)$.
\end{lemma}

\begin{proof}
Appliquons le Lemme \ref{etale-cohomology-lemma-when-ctf} pour nous ramener à une question
sur les complexes finis de faisceaux constructibles plats de $\Lambda$-modules.
L'énoncé résulte alors de l'exactitude de $f^{-1} = f^*$ et du
Lemme \ref{etale-cohomology-lemma-pullback-constructible}.
\end{proof}

\begin{lemma}
\label{etale-cohomology-lemma-connected-ctf-locally-constant}
Soit $X$ un schéma connexe. Soit $\Lambda$ un anneau noethérien.
Soit $K \in D_{ctf}(X_\etale, \Lambda)$ à faisceaux de cohomologie localement constants.
Il existe alors un complexe fini de $\Lambda$-modules projectifs de type fini
$M^\bullet$ et un recouvrement étale $\{U_i \to X\}$ tels que
$K|_{U_i} \cong \underline{M^\bullet}|_{U_i}$ dans $D(U_{i, \etale}, \Lambda)$.
\end{lemma}

\begin{proof}
Choisissons un recouvrement étale $\{U_i \to X\}$ tel que $K|_{U_i}$
soit constant, disons $K|_{U_i} \cong \underline{M_i^\bullet}_{U_i}$
pour un complexe fini de $\Lambda$-modules $M_i^\bullet$ de type fini.
Voir Cohomologie sur les sites, Lemme
\ref{sites-cohomology-lemma-locally-constant}.
Remarquons que $U_i \times_X U_j$ est vide si $M_i^\bullet$
n'est pas isomorphe à $M_j^\bullet$ dans $D(\Lambda)$.
Pour chaque complexe de $\Lambda$-modules $M^\bullet$, posons
$I_{M^\bullet} =
\{i \in I \mid M_i^\bullet \cong M^\bullet\text{ dans }D(\Lambda)\}$.
Comme les morphismes étales sont ouverts,
$U_{M^\bullet} = \bigcup_{i \in I_{M^\bullet}} \Im(U_i \to X)$
est une partie ouverte de $X$. Alors $X = \coprod U_{M^\bullet}$ est un
recouvrement ouvert disjoint de $X$. Comme $X$ est connexe, un seul $U_{M^\bullet}$
est non vide. Puisque $K$ appartient à $D_{ctf}(X_\etale, \Lambda)$, $M^\bullet$
est un complexe parfait de $\Lambda$-modules; voir
Compléments d'algèbre, Lemme \ref{more-algebra-lemma-perfect}.
On peut donc supposer que $M^\bullet$ est un complexe fini de
$\Lambda$-modules projectifs de type fini.
\end{proof}








\section{Faisceaux de torsion}
\label{etale-cohomology-section-torsion}

\noindent
Voici une brève section sur les faisceaux abéliens de torsion et leur cohomologie
étale. Soit $\mathcal{C}$ un site. Nous avons montré dans
Cohomologie sur les sites, Lemme \ref{sites-cohomology-lemma-torsion},
que tout objet de $D(\mathcal{C})$ dont les faisceaux de cohomologie sont
de torsion peut être représenté par un complexe dont tous les termes
sont de torsion.

\begin{lemma}
\label{etale-cohomology-lemma-torsion-cohomology}
Soit $X$ un schéma quasi-compact et quasi-séparé.
\begin{enumerate}
\item Si $\mathcal{F}$ est un faisceau abélien de torsion sur $X_\etale$, alors
$H^n_\etale(X, \mathcal{F})$ est un groupe abélien de torsion pour tout $n$.
\item Si $K$ dans $D^+(X_\etale)$ a des faisceaux de cohomologie de torsion, alors
$H^n_\etale(X, K)$ est un groupe abélien de torsion pour tout $n$.
\end{enumerate}
\end{lemma}

\begin{proof}
Pour démontrer (1), écrivons $\mathcal{F} = \bigcup_d \mathcal{F}[d]$, où
$\mathcal{F}[d]$ est le sous-faisceau de $d$-torsion. Par le
Lemme \ref{etale-cohomology-lemma-colimit}, on a
$H^n_\etale(X, \mathcal{F}) = \colim H^n_\etale(X, \mathcal{F}[d])$.
Cela démontre (1), car $H^n_\etale(X, \mathcal{F}[d])$ est annulé par $d$.

\medskip\noindent
Pour démontrer (2), on peut utiliser la suite spectrale
$E_2^{p, q} = H^p_\etale(X, H^q(K))$ qui aboutit à $H^n_\etale(X, K)$
(Catégories dérivées, Lemme \ref{derived-lemma-two-ss-complex-functor})
et le résultat pour les faisceaux.
\end{proof}

\begin{lemma}
\label{etale-cohomology-lemma-torsion-direct-image}
Soit $f : X \to Y$ un morphisme de schémas quasi-compact et
quasi-séparé.
\begin{enumerate}
\item Si $\mathcal{F}$ est un faisceau abélien de torsion sur $X_\etale$, alors
$R^nf_*\mathcal{F}$ est un faisceau abélien de torsion sur $Y_\etale$ pour tout $n$.
\item Si $K$ dans $D^+(X_\etale)$ a des faisceaux de cohomologie de torsion, alors
$Rf_*K$ est un objet de $D^+(Y_\etale)$ dont les faisceaux de cohomologie sont
des faisceaux abéliens de torsion.
\end{enumerate}
\end{lemma}

\begin{proof}
Démontrons (1). Rappelons que $R^nf_*\mathcal{F}$ est le faisceau associé
au préfaisceau $V \mapsto H^n_\etale(X \times_Y V, \mathcal{F})$
sur $Y_\etale$. Voir Cohomologie sur les sites,
Lemme \ref{sites-cohomology-lemma-higher-direct-images}.
Si l'on choisit $V$ affine, alors $X \times_Y V$ est
quasi-compact et quasi-séparé puisque $f$ l'est; on peut donc
appliquer le Lemme \ref{etale-cohomology-lemma-torsion-cohomology} pour voir que
$H^n_\etale(X \times_Y V, \mathcal{F})$ est de torsion.

\medskip\noindent
Démontrons (2). Rappelons que $R^nf_*K$ est le faisceau associé
au préfaisceau $V \mapsto H^n_\etale(X \times_Y V, K)$
sur $Y_\etale$. Voir Cohomologie sur les sites,
Lemme \ref{sites-cohomology-lemma-unbounded-describe-higher-direct-images}.
Si l'on choisit $V$ affine, alors $X \times_Y V$ est
quasi-compact et quasi-séparé puisque $f$ l'est; on peut donc
appliquer le Lemme \ref{etale-cohomology-lemma-torsion-cohomology} pour voir que
$H^n_\etale(X \times_Y V, K)$ est de torsion.
\end{proof}









\section{Cohomologie à supports dans un sous-schéma fermé}
\label{etale-cohomology-section-cohomology-support}

\noindent
Soit $X$ un schéma et soit $Z \subset X$ un sous-schéma fermé.
Soit $\mathcal{F}$ un faisceau abélien sur $X_\etale$. Posons
$$
\Gamma_Z(X, \mathcal{F}) =
\{s \in \mathcal{F}(X) \mid \text{Supp}(s) \subset Z\}
$$
pour le groupe des sections à support dans $Z$ (Définition \ref{etale-cohomology-definition-support}).
C'est un foncteur exact à gauche, qui n'est pas exact en général.
On obtient donc un foncteur dérivé
$$
R\Gamma_Z(X, -) : D(X_\etale) \longrightarrow D(\textit{Ab})
$$
et les groupes de cohomologie à support dans $Z$ définis par
$H^q_Z(X, \mathcal{F}) = R^q\Gamma_Z(X, \mathcal{F})$.

\medskip\noindent
Soit $\mathcal{I}$ un faisceau abélien injectif sur $X_\etale$.
Posons $U = X \setminus Z$.
Alors l'application de restriction $\mathcal{I}(X) \to \mathcal{I}(U)$ est surjective
(Cohomologie sur les sites, Lemme
\ref{sites-cohomology-lemma-restriction-along-monomorphism-surjective})
et son noyau est $\Gamma_Z(X, \mathcal{I})$. Il s'ensuit immédiatement que,
pour $K \in D(X_\etale)$, il existe un triangle distingué
$$
R\Gamma_Z(X, K) \to R\Gamma(X, K) \to R\Gamma(U, K) \to R\Gamma_Z(X, K)[1]
$$
dans $D(\textit{Ab})$. On en déduit une longue suite exacte de
cohomologie
$$
\ldots \to H^i_Z(X, K) \to H^i(X, K) \to H^i(U, K) \to
H^{i + 1}_Z(X, K) \to \ldots
$$
pour tout $K$ de $D(X_\etale)$.

\medskip\noindent
Pour un faisceau abélien $\mathcal{F}$ sur $X_\etale$, on peut considérer
le {\it sous-faisceau des sections à support dans $Z$}, noté
$\mathcal{H}_Z(\mathcal{F})$, défini par la règle
$$
\mathcal{H}_Z(\mathcal{F})(U) =
\{s \in \mathcal{F}(U) \mid \text{Supp}(s) \subset U \times_X Z\}
$$
Nous utilisons ici le support d'une section défini dans la
Définition \ref{etale-cohomology-definition-support}. Grâce à l'équivalence de la
Proposition \ref{etale-cohomology-proposition-closed-immersion-pushforward},
on peut considérer $\mathcal{H}_Z(\mathcal{F})$ comme un faisceau abélien sur
$Z_\etale$. On obtient ainsi un foncteur
$$
\textit{Ab}(X_\etale) \longrightarrow \textit{Ab}(Z_\etale),\quad
\mathcal{F} \longmapsto \mathcal{H}_Z(\mathcal{F})
$$
qui est exact à gauche, mais qui n'est pas exact en général.

\begin{lemma}
\label{etale-cohomology-lemma-sections-with-support-acyclic}
Soit $i : Z \to X$ une immersion fermée de schémas.
Soit $\mathcal{I}$ un faisceau abélien injectif sur $X_\etale$.
Alors $\mathcal{H}_Z(\mathcal{I})$ est un faisceau abélien injectif
sur $Z_\etale$.
\end{lemma}

\begin{proof}
Remarquons que, pour tout faisceau abélien $\mathcal{G}$ sur $Z_\etale$,
on a
$$
\Hom_Z(\mathcal{G}, \mathcal{H}_Z(\mathcal{F})) =
\Hom_X(i_*\mathcal{G}, \mathcal{F})
$$
car toute section de $i_*\mathcal{G}$ a son support dans $Z$.
Comme $i_*$ est exact (section \ref{etale-cohomology-section-closed-immersions}) et que
$\mathcal{I}$ est injectif sur $X_\etale$, on en conclut que
$\mathcal{H}_Z(\mathcal{I})$ est injectif sur $Z_\etale$.
\end{proof}

\noindent
Notons
$$
R\mathcal{H}_Z : D(X_\etale) \longrightarrow D(Z_\etale)
$$
le foncteur dérivé. Posons
$\mathcal{H}^q_Z(\mathcal{F}) = R^q\mathcal{H}_Z(\mathcal{F})$, de sorte que
$\mathcal{H}^0_Z(\mathcal{F}) = \mathcal{H}_Z(\mathcal{F})$.
Par le lemme précédent, on dispose d'une suite spectrale de Grothendieck
$$
E_2^{p, q} = H^p(Z, \mathcal{H}^q_Z(\mathcal{F}))
\Rightarrow H^{p + q}_Z(X, \mathcal{F})
$$

\begin{lemma}
\label{etale-cohomology-lemma-cohomology-with-support-sheaf-on-support}
Soit $i : Z \to X$ une immersion fermée de schémas.
Soit $\mathcal{G}$ un faisceau abélien injectif sur $Z_\etale$.
Alors $\mathcal{H}^p_Z(i_*\mathcal{G}) = 0$ pour $p > 0$.
\end{lemma}

\begin{proof}
Cela est vrai parce que le foncteur $i_*$ est exact et transforme
les faisceaux abéliens injectifs en faisceaux abéliens injectifs
(Cohomologie sur les sites, Lemme
\ref{sites-cohomology-lemma-pushforward-injective-flat}).
\end{proof}

\begin{lemma}
\label{etale-cohomology-lemma-cohomology-with-support-triangle}
Soit $i : Z \to X$ une immersion fermée de schémas.
Soit $j : U \to X$ l'inclusion du complémentaire de $Z$.
Soit $\mathcal{F}$ un faisceau abélien sur $X_\etale$.
Il existe un triangle distingué
$$
i_*R\mathcal{H}_Z(\mathcal{F}) \to \mathcal{F} \to Rj_*(\mathcal{F}|_U) \to
i_*R\mathcal{H}_Z(\mathcal{F})[1]
$$
dans $D(X_\etale)$. Il donne une suite exacte
$$
0 \to i_*\mathcal{H}_Z(\mathcal{F}) \to \mathcal{F} \to
j_*(\mathcal{F}|_U) \to i_*\mathcal{H}^1_Z(\mathcal{F}) \to 0
$$
et des isomorphismes
$R^pj_*(\mathcal{F}|_U) \cong i_*\mathcal{H}^{p + 1}_Z(\mathcal{F})$
pour $p \geq 1$.
\end{lemma}

\begin{proof}
Pour obtenir le triangle distingué, choisissons une résolution injective
$\mathcal{F} \to \mathcal{I}^\bullet$. On obtient alors une suite exacte courte
de complexes
$$
0 \to
i_*\mathcal{H}_Z(\mathcal{I}^\bullet) \to \mathcal{I}^\bullet
\to j_*(\mathcal{I}^\bullet|_U) \to 0
$$
d'après la discussion précédente. Cela fournit le triangle distingué, par
Catégories dérivées, section \ref{derived-section-canonical-delta-functor}.
\end{proof}

\noindent
Soit $X$ un schéma et soit $Z \subset X$ un sous-schéma fermé.
On note $D_Z(X_\etale)$ la sous-catégorie triangulée strictement pleine et
saturée de $D(X_\etale)$ formée des complexes dont les faisceaux de cohomologie
sont à support dans $Z$. Remarquons que $D_Z(X_\etale)$ ne dépend que
de la partie fermée sous-jacente de $X$.

\begin{lemma}
\label{etale-cohomology-lemma-complexes-with-support-on-closed}
Soit $i : Z \to X$ une immersion fermée de schémas.
Le foncteur $Ri_{small, *} = i_{small, *} : D(Z_\etale) \to D(X_\etale)$
induit une équivalence $D(Z_\etale) \to D_Z(X_\etale)$ dont un quasi-inverse est
$$
i_{small}^{-1}|_{D_Z(X_\etale)} = R\mathcal{H}_Z|_{D_Z(X_\etale)}
$$
\end{lemma}

\begin{proof}
Rappelons que $i_{small}^{-1}$ et $i_{small, *}$ forment une paire de foncteurs
adjoints exacts telle que $i_{small}^{-1}i_{small, *}$ soit isomorphe au
foncteur identité sur les faisceaux abéliens. Voir la
Proposition \ref{etale-cohomology-proposition-closed-immersion-pushforward} et le
Lemme \ref{etale-cohomology-lemma-stalk-pullback}. Ainsi,
$i_{small, *} : D(Z_\etale) \to D_Z(X_\etale)$ est pleinement fidèle et
$i_{small}^{-1}$ en fournit
un inverse à gauche. D'autre part, supposons que $K$ soit un objet de
$D_Z(X_\etale)$ et considérons le morphisme d'adjonction
$K \to i_{small, *}i_{small}^{-1}K$.
Par exactitude de $i_{small, *}$ et de $i_{small}^{-1}$,
il induit sur les faisceaux de cohomologie les morphismes d'adjonction
$H^n(K) \to i_{small, *}i_{small}^{-1}H^n(K)$.
Comme ces faisceaux de cohomologie
sont à support dans $Z$, ces morphismes d'adjonction sont des isomorphismes;
on en conclut que $D(Z_\etale) \to D_Z(X_\etale)$ est une équivalence.

\medskip\noindent
Pour achever la démonstration, il faut montrer que $R\mathcal{H}_Z(K) = i_{small}^{-1}K$
si $K$ est un objet de $D_Z(X_\etale)$. Pour cela, on peut utiliser l'égalité
$K = i_{small, *}i_{small}^{-1}K$
que l'on vient de démontrer. On peut alors
choisir un représentant K-injectif $\mathcal{I}^\bullet$ de
$i_{small}^{-1}K$.
Comme $i_{small, *}$ est l'adjoint à droite du foncteur exact
$i_{small}^{-1}$, le
complexe $i_{small, *}\mathcal{I}^\bullet$ est K-injectif
(Catégories dérivées, Lemme \ref{derived-lemma-adjoint-preserve-K-injectives}).
On voit que $R\mathcal{H}_Z(K)$ est calculé par
$\mathcal{H}_Z(i_{small, *}\mathcal{I}^\bullet) = \mathcal{I}^\bullet$
comme voulu.
\end{proof}

\begin{lemma}
\label{etale-cohomology-lemma-cohomology-with-support-quasi-coherent}
Soit $X$ un schéma. Soit $Z \subset X$ un sous-schéma fermé, et soit $i : Z \to X$ l'immersion canonique.
Soit $\mathcal{F}$ un $\mathcal{O}_X$-module quasi-cohérent
et notons $\mathcal{F}^a$ le faisceau quasi-cohérent associé
sur le petit site étale de $X$
(Proposition \ref{etale-cohomology-proposition-quasi-coherent-sheaf-fpqc}). Alors
\begin{enumerate}
\item $H^q_Z(X, \mathcal{F})$ s'identifie à $H^q_Z(X_\etale, \mathcal{F}^a)$;
\item si le complémentaire de $Z$ est rétrocompact dans $X$, alors
$i_*\mathcal{H}^q_Z(\mathcal{F}^a)$ est un faisceau quasi-cohérent de
$\mathcal{O}_X$-modules égal à $(i_*\mathcal{H}^q_Z(\mathcal{F}))^a$.
\end{enumerate}
\end{lemma}

\begin{proof}
Soit $j : U \to X$ l'inclusion du complémentaire de $Z$.
L'assertion (1) sur les groupes de cohomologie résulte des longues
suites exactes de cohomologie à supports et des identifications
$H^q(X_\etale, \mathcal{F}^a) = H^q(X, \mathcal{F})$ et
$H^q(U_\etale, \mathcal{F}^a) = H^q(U, \mathcal{F})$; voir le
Théorème \ref{etale-cohomology-theorem-zariski-fpqc-quasi-coherent}.
Si $j : U \to X$ est un morphisme quasi-compact, c'est-à-dire si $U \subset X$
est rétrocompact, alors $R^qj_*$ transforme les faisceaux quasi-cohérents
en faisceaux quasi-cohérents
(Cohomologie des schémas, Lemme
\ref{coherent-lemma-quasi-coherence-higher-direct-images})
et commute à la formation du
faisceau associé sur les sites étales
(Descente, Lemme \ref{descent-lemma-higher-direct-images-small-etale}).
On conclut en appliquant le
Lemme \ref{etale-cohomology-lemma-cohomology-with-support-triangle}.
\end{proof}







\section{Schémas dont les anneaux locaux sont strictement henséliens}
\sectionmark{Anneaux locaux strictement henséliens}
\label{etale-cohomology-section-strictly-henselian-local-rings}

\noindent
Dans cette section, nous rassemblons quelques résultats sur la cohomologie étale
des schémas dont les anneaux locaux sont strictement henséliens. Voici par exemple
une généralisation utile du
Lemme \ref{etale-cohomology-lemma-vanishing-etale-cohomology-strictly-henselian}.

\begin{lemma}
\label{etale-cohomology-lemma-local-rings-strictly-henselian}
Soit $S$ un schéma dont tous les anneaux locaux sont strictement henséliens.
Alors, pour tout faisceau abélien $\mathcal{F}$ sur $S_\etale$, on a
$H^i(S_\etale, \mathcal{F}) = H^i(S_{Zar}, \mathcal{F})$.
\end{lemma}

\begin{proof}
Soit $\epsilon : S_\etale \to S_{Zar}$ le morphisme de sites défini
par le foncteur d'inclusion. Le faisceau de Zariski $R^p\epsilon_*\mathcal{F}$
est le faisceau associé au préfaisceau $U \mapsto H^p_\etale(U, \mathcal{F})$.
Sa fibre en $x \in S$ est donc
$\colim H^p_\etale(U, \mathcal{F}) =
H^p_\etale(\Spec(\mathcal{O}_{S, x}), \mathcal{G}_x)$,
où $\mathcal{G}_x$ désigne l'image inverse de $\mathcal{F}$
sur $\Spec(\mathcal{O}_{S, x})$; voir le
Lemme \ref{etale-cohomology-lemma-directed-colimit-cohomology}.
Ainsi, les faisceaux d'images directes supérieures $R^p\epsilon_*\mathcal{F}$ sont
nuls d'après le
Lemme \ref{etale-cohomology-lemma-vanishing-etale-cohomology-strictly-henselian},
et l'on conclut par la suite spectrale de Leray.
\end{proof}

\begin{lemma}
\label{etale-cohomology-lemma-gabber-criterion}
Soit $R$ un anneau dont tous les anneaux locaux sont strictement henséliens.
Soit $\mathcal{F}$ un faisceau sur $\Spec(R)_\etale$.
Supposons que, pour tous $f,g\in R$, le noyau de
$$
H^1_\etale(D(f + g), \mathcal{F})
\longrightarrow
H^1_\etale(D(f(f + g)), \mathcal{F}) \oplus
H^1_\etale(D(g(f + g)), \mathcal{F})
$$
soit nul. Alors $H^q_\etale(\Spec(R), \mathcal{F}) = 0$ pour $q > 0$.
\end{lemma}

\begin{proof}
D'après le Lemme \ref{etale-cohomology-lemma-local-rings-strictly-henselian}, la
cohomologie étale de $\mathcal{F}$ coïncide avec sa cohomologie de Zariski
sur tout ouvert de $\Spec(R)$. Nous démontrerons par récurrence sur $i$
l'assertion suivante : pour $h \in R$, on a $H^q_\etale(D(h), \mathcal{F}) = 0$
si $1 \leq q \leq i$. Le cas initial $i = 0$ est trivial. Supposons $i \geq 1$.

\medskip\noindent
Soient $\xi \in H^q_\etale(D(h), \mathcal{F})$, avec $1 \leq q \leq i$,
et $h \in R$. Si $q < i$, la récurrence conclut; supposons donc $q = i$.
Après avoir remplacé $R$ par $R_h$, on peut supposer que
$\xi \in H^i_\etale(\Spec(R), \mathcal{F})$; nous omettons quelques détails.
Soit $I \subset R$ l'ensemble des éléments $f \in R$ tels que
$\xi|_{D(f)} = 0$. Comme $\xi$ est localement nulle pour la topologie de Zariski,
pour tout idéal premier $\mathfrak p$ de $R$, il existe $f \in I$
tel que $f \not \in \mathfrak p$. Si l'on montre que $I$ est un
idéal, alors $1 \in I$ et la preuve sera achevée. Il est clair que
$f \in I$ et $r \in R$ impliquent $rf \in I$. Supposons donc $f,g \in I$
et montrons que $f + g \in I$. Si $q = i = 1$, c'est exactement
l'hypothèse du lemme, d'où le résultat pour $i = 1$. Si
$q = i > 1$, remarquons que
$$
D(f + g) = D(f(f + g)) \cup D(g(f + g))
$$
Par Mayer--Vietoris (Cohomologie, Lemme \ref{cohomology-lemma-mayer-vietoris},
qui s'applique puisque, sur les sous-schémas ouverts de $\Spec(R)$, la cohomologie
étale coïncide avec la cohomologie de Zariski), on dispose d'une suite exacte
$$
\xymatrix{
H^{i - 1}_\etale(D(fg(f + g)), \mathcal{F}) \ar[d] \\
H^i_\etale(D(f + g), \mathcal{F}) \ar[d] \\
H^i_\etale(D(f(f + g)), \mathcal{F}) \oplus
H^i_\etale(D(g(f + g)), \mathcal{F})
}
$$
et le résultat s'ensuit, puisque le premier groupe est nul par récurrence.
\end{proof}

\begin{lemma}
\label{etale-cohomology-lemma-affine-only-closed-points}
Soit $S$ un schéma affine tel que
(1) tous ses points soient fermés et (2) tous ses corps résiduels soient
séparablement clos. Alors,
pour tout faisceau abélien $\mathcal{F}$ sur $S_\etale$, on a
$H^i(S_\etale, \mathcal{F}) = 0$ pour $i > 0$.
\end{lemma}

\begin{proof}
La condition (1) implique que l'espace topologique sous-jacent
à $S$ est profini; voir
Algèbre, Lemme \ref{algebra-lemma-ring-with-only-minimal-primes}.
Les groupes de cohomologie supérieure d'un faisceau abélien sur l'espace
topologique $S$ (c'est-à-dire la cohomologie de Zariski) sont donc nuls; voir
Cohomologie, Lemme \ref{cohomology-lemma-vanishing-for-profinite}.
Les anneaux locaux sont strictement henséliens d'après
Algèbre, Lemme \ref{algebra-lemma-local-dimension-zero-henselian}.
La cohomologie étale de $S$ se calcule donc par la cohomologie de Zariski
d'après le Lemme \ref{etale-cohomology-lemma-local-rings-strictly-henselian},
ce qui achève la démonstration.
\end{proof}

\noindent
Le spectre d'un anneau absolument intégralement clos
est un exemple de schéma dont tous les anneaux locaux sont
strictement henséliens; voir Compléments d'algèbre, Lemme
\ref{more-algebra-lemma-absolutely-integrally-closed-strictly-henselian}.
Les domaines normaux à corps des fractions séparablement clos
possèdent même une propriété plus forte, comme l'explique le
lemme suivant.

\begin{lemma}
\label{etale-cohomology-lemma-normal-scheme-with-alg-closed-function-field}
Soit $X$ un schéma intègre normal à corps des fonctions
séparablement clos.
\begin{enumerate}
\item Tout morphisme étale séparé $U \to X$ est une
réunion disjointe d'immersions ouvertes.
\item Tous les anneaux locaux de $X$ sont strictement henséliens.
\end{enumerate}
\end{lemma}

\begin{proof}
Soit $R$ un domaine normal dont le corps des fractions est séparablement
clos. Soit $R \to A$ un homomorphisme d'anneaux étale. Alors
$A \otimes_R K$ est, comme $K$-algèbre, un produit fini
$\prod_{i=1}^n K$ de copies de $K$. Soient $e_i$, $i = 1, \ldots, n$,
les idempotents correspondants de $A \otimes_R K$. Comme $A$ est normal
(Algèbre, Lemme \ref{algebra-lemma-normal-goes-up}),
les idempotents $e_i$ appartiennent à $A$
(Algèbre, Lemme \ref{algebra-lemma-normal-ring-integrally-closed}).
Ainsi $A = \prod Ae_i$, et l'on peut supposer $A \otimes_R K = K$.
Puisque $A \subset A \otimes_R K = K$ (par platitude de $R \to A$ et
parce que $R \subset K$), on conclut que $A$ est un domaine.
Le même argument donne
$A \otimes_R A \subset (A \otimes_R A) \otimes_R K = K$.
L'application $A \otimes_R A \to A$ est donc
injective aussi bien que surjective. Ainsi, $R \to A$ définit une
immersion ouverte d'après
Morphismes, Lemme \ref{morphisms-lemma-universally-injective},
et
Morphismes étales, Théorème \ref{etale-theorem-etale-radicial-open}.

\medskip\noindent
Soit $f : U \to X$ un morphisme étale séparé. Soit $\eta \in X$
le point générique, et posons $f^{-1}(\{\eta\}) = \{\xi_i\}_{i \in I}$.
Le résultat du paragraphe précédent montre ceci :
pour tout ouvert affine $U' \subset U$ dont l'image dans $X$ est contenue dans
un ouvert affine, on a $U' = \coprod_{i \in I} U'_i$, où $U'_i$
est l'ensemble des points de $U'$ qui sont des spécialisations de $\xi_i$.
De plus, le morphisme $U'_i \to X$ est une immersion ouverte.
Il s'ensuit que $U_i = \overline{\{\xi_i\}}$ est un sous-schéma ouvert et fermé
de $U$, et que $U_i \to X$ est localement, sur la source,
un isomorphisme. D'après Morphismes,
Lemme \ref{morphisms-lemma-distinct-local-rings},
le fait que $U_i \to X$ soit séparé implique que
$U_i \to X$ est injectif; on conclut que $U_i \to X$
est une immersion ouverte, ce qui prouve (1).

\medskip\noindent
La partie (2) résulte de la partie (1) et de la description de l'hensélisé
strict de $\mathcal{O}_{X, x}$ comme anneau local en $\overline{x}$
sur le site étale de $X$ (Lemme \ref{etale-cohomology-lemma-describe-etale-local-ring}).
On peut aussi la démontrer directement; voir
Groupes fondamentaux, Lemme
\ref{pione-lemma-normal-local-domain-separablly-closed-fraction-field}.
\end{proof}

\begin{lemma}
\label{etale-cohomology-lemma-Rf-star-zero-normal-with-alg-closed-function-field}
Soit $f : X \to Y$ un morphisme de schémas, où $X$ est un schéma intègre
normal à corps des fonctions séparablement clos. Alors
$R^qf_*\underline{M} = 0$ pour $q > 0$ et tout groupe abélien $M$.
\end{lemma}

\begin{proof}
Rappelons que $R^qf_*\underline{M}$ est le faisceau associé
au préfaisceau $V \mapsto H^q_\etale(V \times_Y X, M)$ sur $Y_\etale$; voir le
Lemme \ref{etale-cohomology-lemma-higher-direct-images}.
Si $V$ est affine, alors $V \times_Y X \to X$ est séparé et étale.
Par conséquent, $V \times_Y X = \coprod U_i$ est une réunion disjointe de
sous-schémas ouverts $U_i$ de $X$; voir le
Lemme \ref{etale-cohomology-lemma-normal-scheme-with-alg-closed-function-field}.
D'après le Lemme \ref{etale-cohomology-lemma-local-rings-strictly-henselian},
$H^q_\etale(U_i, M)$ est égal à
$H^q_{Zar}(U_i, M)$. Ce dernier est nul d'après
Cohomologie, Lemme \ref{cohomology-lemma-irreducible-constant-cohomology-zero}.
\end{proof}

\begin{lemma}
\label{etale-cohomology-lemma-closed-of-affine-normal-scheme-with-alg-closed-function-field}
Soit $X$ un schéma affine intègre normal à corps des fonctions
séparablement clos. Soit $Z \subset X$ un sous-schéma fermé. Soit
$V \to Z$ un morphisme étale, avec $V$ affine. Alors $V$ est une réunion
disjointe finie de sous-schémas ouverts de $Z$. Si $V \to Z$ est
surjectif et fini étale, alors $V \to Z$ admet une section.
\end{lemma}

\begin{proof}
D'après Algèbre, Lemme \ref{algebra-lemma-lift-etale},
on peut relever $V$ en un schéma affine $U$ étale sur $X$.
Appliquons le Lemme \ref{etale-cohomology-lemma-normal-scheme-with-alg-closed-function-field}
à $U \to X$ pour obtenir la première assertion.

\medskip\noindent
La dernière assertion est une conséquence de la première.
Soit $V = \coprod_{i=1}^n V_i$ une décomposition finie
en sous-schémas ouverts et
fermés telle que $V_i \to Z$ soit une immersion ouverte.
Comme $V \to Z$ est fini, $V_i \to Z$ est aussi fermé.
Soit $U_i \subset Z$ son image. On a alors une décomposition
en sous-schémas ouverts et fermés
$$
Z =
\coprod\nolimits_{(A, B)}
\bigcap\nolimits_{i \in A} U_i \cap
\bigcap\nolimits_{i \in B} U_i^c
$$
où la réunion disjointe porte sur les partitions $\{1, \ldots, n\} = A \amalg B$
pour lesquelles $A$ est non vide.
Chacune des strates est contenue dans l'un des $U_i$, ce qui
fournit la section cherchée.
\end{proof}

\begin{lemma}
\label{etale-cohomology-lemma-gabber-for-h1-absolutely-algebraically-closed}
Soit $X$ un schéma affine intègre normal à corps des fonctions
séparablement clos. Soit $Z \subset X$ un sous-schéma fermé.
Pour tout groupe abélien fini $M$, on a $H^1_\etale(Z, \underline{M}) = 0$.
\end{lemma}

\begin{proof}
D'après Cohomologie sur les sites, Lemme \ref{sites-cohomology-lemma-torsors-h1},
un élément de $H^1_\etale(Z, \underline{M})$ correspond à un
$\underline{M}$-torseur $\mathcal{F}$ sur $Z_\etale$.
Un tel torseur est manifestement un faisceau localement constant fini.
Ainsi, $\mathcal{F}$ est représentable par un schéma $V$ fini
étale sur $Z$; voir le Lemme \ref{etale-cohomology-lemma-characterize-finite-locally-constant}.
Le morphisme $V \to Z$ est surjectif, puisque tout torseur est localement trivial.
Comme $V \to Z$ admet une section d'après le
Lemme \ref{etale-cohomology-lemma-closed-of-affine-normal-scheme-with-alg-closed-function-field},
la preuve est achevée.
\end{proof}

\begin{lemma}
\label{etale-cohomology-lemma-gabber-for-absolutely-algebraically-closed}
Soit $X$ un schéma affine intègre normal à corps des fonctions
séparablement clos. Soit $Z \subset X$ un sous-schéma fermé.
Pour tout groupe abélien fini $M$, on a
$H^q_\etale(Z, \underline{M}) = 0$ pour $q \geq 1$.
\end{lemma}

\begin{proof}
Écrivons $X = \Spec(R)$ et $Z = \Spec(R')$, de sorte que l'on dispose d'une surjection
d'anneaux $R \to R'$. Tous les anneaux locaux de $R'$ sont strictement henséliens d'après le
Lemme \ref{etale-cohomology-lemma-normal-scheme-with-alg-closed-function-field} et
Algèbre, Lemme \ref{algebra-lemma-quotient-strict-henselization}.
De plus, pour tout $f' \in R'$, il existe une surjection
$R_f \to R'_{f'}$, où $f \in R$ est un relèvement de $f'$.
Comme $R_f$ est un domaine normal à corps des fractions séparablement clos,
on a $H^1_\etale(D(f'), \underline{M}) = 0$ d'après le
Lemme \ref{etale-cohomology-lemma-gabber-for-h1-absolutely-algebraically-closed}.
On peut donc appliquer le Lemme \ref{etale-cohomology-lemma-gabber-criterion} à $Z = \Spec(R')$
pour conclure.
\end{proof}

\begin{lemma}
\label{etale-cohomology-lemma-integral-cover-trivial-cohomology}
Soit $X$ un schéma affine.
\begin{enumerate}
\item Il existe un morphisme entier surjectif $X' \to X$ tel que, pour
tout sous-schéma fermé $Z' \subset X'$, tout groupe abélien fini $M$ et
tout $q \geq 1$, on ait $H^q_\etale(Z', \underline{M}) = 0$.
\item Pour tout sous-schéma fermé $Z \subset X$, tout groupe abélien fini $M$,
tout $q \geq 1$ et tout $\xi \in H^q_\etale(Z, \underline{M})$, il existe un
morphisme fini surjectif $X' \to X$ de présentation finie tel que
l'image inverse de $\xi$ soit nulle dans $H^q_\etale(X' \times_X Z, \underline{M})$.
\end{enumerate}
\end{lemma}

\begin{proof}
Écrivons $X = \Spec(A)$. Écrivons $A = \mathbf{Z}[x_i]/J$ pour un idéal $J$.
Soit $R$ la clôture intégrale de $\mathbf{Z}[x_i]$ dans une clôture
algébrique du corps des fractions de $\mathbf{Z}[x_i]$. Posons
$A' = R/JR$ et $X' = \Spec(A')$. Le Lemme
\ref{etale-cohomology-lemma-gabber-for-absolutely-algebraically-closed} montre que cet exemple satisfait (1).

\medskip\noindent
Démontrons (2). Soit $X' \to X$ le morphisme entier surjectif construit
ci-dessus. L'image de $\xi$ dans
$H^q_\etale(X' \times_X Z, \underline{M})$ est bien nulle. On peut écrire $X'$ comme une
limite projective $X' = \lim X'_i$ de schémas finis et de présentation finie
sur $X$; cela se fait aisément dans le présent cas affine, mais c'est
un cas particulier du lemme plus général Limites, Lemme
\ref{limits-lemma-integral-limit-finite-and-finite-presentation}.
D'après le Lemme \ref{etale-cohomology-lemma-directed-colimit-cohomology},
l'image de $\xi$ dans $H^q_\etale(X'_i \times_X Z, \underline{M})$ est nulle
pour un indice $i$ assez grand.
\end{proof}








\section{Annulation pour les anneaux absolument intégralement clos}
\sectionmark{Anneaux absolument intégralement clos}
\label{etale-cohomology-section-aic-vanishing}

\noindent
Rappelons qu'un anneau $R$ est dit absolument intégralement clos
si tout polynôme unitaire à coefficients dans $R$ a une racine dans $R$
(Compléments d'algèbre, définition
\ref{more-algebra-definition-absolutely-integrally-closed}).
Dans cette section, nous démontrons que la cohomologie étale de
$\Spec(R)$ à coefficients dans un groupe de torsion fini
est nulle en degrés strictement positifs (proposition \ref{etale-cohomology-proposition-aic-vanishing}),
ce qui améliore légèrement le
lemme antérieur \ref{etale-cohomology-lemma-gabber-for-absolutely-algebraically-closed}.
Nous suggérons au lecteur de sauter cette section.

\begin{lemma}
\label{etale-cohomology-lemma-find-extension}
Soit $A$ un anneau. Soient $a, b \in A$ tels que
$aA + bA = A$ et que $a \bmod bA$ soit une racine de l'unité.
Il existe alors une extension monogène $A \subset B$
et un élément $y \in B$ tels que $u = a - by$ soit inversible.
\end{lemma}

\begin{proof}
Écrivons $a^n \equiv 1 \bmod bA$. En particulier, $a^i$ est inversible modulo
$b^mA$ pour tous $i, m \geq 1$. Nous affirmons qu'il existe des éléments
$a_1, \ldots, a_n \in A$ tels que
$$
1 = a^n + a_1 a^{n - 1}b + a_2 a^{n - 2}b^2 + \ldots + a_n b^n
$$
En effet, puisque $1 - a^n \in bA$, nous pouvons trouver un élément $a_1 \in A$
tel que $1 - a^n - a_1 a^{n - 1} b \in b^2A$, grâce à l'inversibilité
de $a^{n - 1}$ modulo $bA$. Ensuite, nous pouvons trouver $a_2 \in A$
tel que $1 - a^n - a_1 a^{n - 1} b - a_2 a^{n - 2} b^2 \in b^3A$.
Et ainsi de suite. Nous finissons par trouver $a_1, \ldots, a_{n - 1} \in A$
tels que $1 - (a^n + a_1 a^{n - 1}b + a_2 a^{n - 2}b^2 +
\ldots + a_{n - 1} ab^{n - 1}) \in b^nA$. Nous pouvons alors trouver
$a_n \in A$ de sorte que l'égalité affichée soit satisfaite.

\medskip\noindent
Pour $a_1, \ldots, a_n$ comme ci-dessus, nous affirmons que le choix suivant convient :
$$
B = A[y]/(y^n + a_1 y^{n - 1} + a_2 y^{n - 2} + \ldots + a_n)
$$
En effet, supposons que $\mathfrak q \subset B$ soit un idéal premier
situé au-dessus de $\mathfrak p \subset A$. Pour obtenir une contradiction,
supposons que $u = a - by$ appartienne à $\mathfrak q$. Si $b \in \mathfrak p$,
alors $a \not \in \mathfrak p$ puisque $aA + bA = A$, et donc $u$
n'appartient pas à $\mathfrak q$. Nous pouvons ainsi supposer que $b \not \in \mathfrak p$, c'est-à-dire
que $b \not \in \mathfrak q$. Il s'ensuit que $y \bmod \mathfrak q$
est égal à $a/b \bmod \mathfrak q$. Mais nous obtenons alors
$$
0 = y^n + a_1 y^{n - 1} + a_2 y^{n - 2} + \ldots + a_n =
b^{-n}(a^n + a_1 a^{n - 1}b + a_2 a^{n - 2}b^2 + \ldots + a_nb^n) =
b^{-n}
$$
une contradiction. Cela achève la preuve.
\end{proof}

\noindent
Pour expliquer la preuve, nous devons introduire quelques schémas en groupes.
Fixons un nombre premier $\ell$. Posons
$$
A = \mathbf{Z}[\zeta] =
\mathbf{Z}[x]/(x^{\ell - 1} + x^{\ell - 2} + \ldots + 1)
$$
Autrement dit, $A$ est l'extension monogène de $\mathbf{Z}$ engendrée par
une racine primitive $\ell$-ième de l'unité $\zeta$. Posons
$$
\pi = \zeta - 1
$$
Un calcul (omis) montre que $\ell$ est divisible par $\pi^{\ell - 1}$
dans $A$. Notre premier schéma en groupes sur $A$ est
$$
G = \Spec(A[s, \frac{1}{\pi s + 1}])
$$
dont la loi de groupe est donnée par la comultiplication
$$
\mu :
A[s, \frac{1}{\pi s + 1}]
\longrightarrow
A[s, \frac{1}{\pi s + 1}] \otimes_A A[s, \frac{1}{\pi s + 1}],\quad
s \longmapsto \pi s \otimes s + s \otimes 1 + 1 \otimes s
$$
Avec ce choix, nous avons
$$
\mu(\pi s + 1) = (\pi s + 1) \otimes (\pi s + 1)
$$
et nous avons donc bien le morphisme de $A$-algèbres indiqué. Nous omettons
la vérification que cela définit bien une loi de groupe. Notre second schéma
en groupes sur $A$ est
$$
H = \Spec(A[t, \frac{1}{\pi^\ell t + 1}])
$$
dont la loi de groupe est donnée par la comultiplication
$$
\mu : A[t, \frac{1}{\pi^\ell t + 1}]
\longrightarrow
A[t, \frac{1}{\pi^\ell t + 1}] \otimes_A A[t, \frac{1}{\pi^\ell t + 1}],\quad
t \longmapsto \pi^\ell t \otimes t + t \otimes 1 + 1 \otimes t
$$
La même vérification que précédemment montre que cela définit une loi de groupe.
Observons ensuite que le polynôme
$$
\Phi(s) = \frac{(\pi s + 1)^\ell - 1}{\pi^\ell}
$$
appartient à $A[s]$, est de degré $\ell$ et est unitaire en $s$. En effet, le coefficient
de $s^i$, pour $0 < i < \ell$, est égal à ${\ell \choose i}\pi^{i - \ell}$;
comme $\pi^{\ell - 1}$ divise $\ell$ dans $A$, c'est un élément de $A$.
Nous obtenons un morphisme d'anneaux
$$
A[t, \frac{1}{\pi^\ell t + 1}]
\longrightarrow
A[s, \frac{1}{\pi s + 1}],\quad
t \longmapsto \Phi(s)
$$
dont le lecteur vérifiera facilement la compatibilité avec les comultiplications.
Nous obtenons ainsi un morphisme de schémas en groupes
$$
f : G \to H
$$
Le lemme suivant montre en particulier que ce morphisme est fidèlement
plat (nous verrons en fait qu'il est fini étale surjectif).

\begin{lemma}
\label{etale-cohomology-lemma-monogenic-one}
Nous avons
$$
A[s, \frac{1}{\pi s + 1}] =
\left(A[t, \frac{1}{\pi^\ell t + 1}]\right)[s]/(\Phi(s) - t)
$$
En particulier, l'algèbre de Hopf de $G$ est une extension monogène
de l'algèbre de Hopf de $H$.
\end{lemma}

\begin{proof}
Cela résulte de la discussion ci-dessus et de la forme de $\Phi(s)$.
Notons en particulier que, par la relation $\Phi(s) = t$, l'élément
$\frac{1}{\pi^\ell t + 1}$ devient l'élément
$\frac{1}{(\pi s + 1)^\ell}$.
\end{proof}

\noindent
Calculons maintenant le noyau de $f$. Comme l'origine de $H$ est
donnée par $t = 0$ dans $H$, nous voyons que le noyau de $f$ est défini
par $\Phi(s) = 0$. Observons que les points à valeurs dans $A$
$\sigma_0, \ldots, \sigma_{\ell - 1}$
de $G$ donnés par
$$
\sigma_i :
s = \frac{\zeta^i - 1}{\pi} = \frac{\zeta^i - 1}{\zeta - 1} =
\zeta^{i - 1} + \zeta^{i - 2} + \ldots + 1,\quad
i = 0, 1, \ldots, \ell - 1
$$
appartiennent certainement à $\Ker(f)$. De plus, ils sont deux à deux
distincts dans {\bf toutes} les fibres de $G \to \Spec(A)$. Le lecteur
calculera aussi que $\sigma_i +_G \sigma_j = \sigma_{i + j \bmod \ell}$.
Nous obtenons donc une immersion fermée de schémas en groupes
$$
\underline{\mathbf{Z}/\ell \mathbf{Z}}_A \longrightarrow \Ker(f)
$$
qui envoie $i$ sur $\sigma_i$. Or, par construction, $\Ker(f)$
est fini et plat sur $\Spec(A)$, de degré $\ell$. Nous en concluons
que ce morphisme est un isomorphisme. En définitive, nous obtenons
une suite exacte courte
\begin{equation}
\label{etale-cohomology-equation-ses}
0 \to
\underline{\mathbf{Z}/\ell \mathbf{Z}}_A
\to G
\to H
\to 0
\end{equation}
de schémas en groupes sur $A$.

\begin{lemma}
\label{etale-cohomology-lemma-lift-points-H-to-G}
Soit $R$ une $A$-algèbre absolument intégralement close. Alors
$G(R) \to H(R)$ est surjectif.
\end{lemma}

\begin{proof}
Soit $h \in H(R)$ correspondant au morphisme de $A$-algèbres
$A[t, \frac{1}{\pi^\ell t + 1}] \to R$ qui envoie $t$ sur $a \in R$.
Comme $\Phi(s)$ est unitaire, nous pouvons trouver $b \in R$
tel que $\Phi(b) = a$. Par le lemme \ref{etale-cohomology-lemma-monogenic-one},
en envoyant $s$ sur $b$, nous obtenons un unique morphisme de $A$-algèbres
$A[s, \frac{1}{\pi s + 1}] \to R$ compatible
avec le morphisme $A[t, \frac{1}{\pi^\ell t + 1}] \to R$ ci-dessus.
Celui-ci correspond à son tour à un élément $g \in G(R)$
qui s'envoie sur $h \in H(R)$.
\end{proof}

\begin{lemma}
\label{etale-cohomology-lemma-interpolate}
Soit $R$ une $A$-algèbre absolument intégralement close.
Soient $I, J \subset R$ des idéaux tels que $I + J = R$.
Il existe un élément $g \in G(R)$ tel
que $g \bmod I = \sigma_0$ et $g \bmod J = \sigma_1$.
\end{lemma}

\begin{proof}
Choisissons $x \in I$ tel que $x \equiv 1 \bmod J$.
Nous pouvons remplacer, et remplaçons, $I$ par $xR$ et $J$ par $(x - 1)R$.
Nous cherchons alors un élément $s \in R$ tel que
\begin{enumerate}
\item $1 + \pi s$ soit inversible,
\item $s \equiv 0 \bmod xR$, et
\item $s \equiv 1 \bmod (x - 1)R$.
\end{enumerate}
Les deux dernières conditions signifient que $s = x + x(x - 1)y$ pour un $y \in R$.
La première exige que $1 + \pi s = 1 + \pi x + \pi x (x - 1) y$
soit inversible dans $R$.
Or $1 + \pi x$ et $\pi x (x - 1)$ engendrent l'idéal
unité de $R$, et $1 + \pi x$ est une racine $\ell$-ième de $1$
modulo $\pi x (x - 1)$\footnote{En effet, $1 + \pi x$ est congru
à $1$ modulo $\pi$, à $1$ modulo $x$ et à
$1 + \pi = \zeta$ modulo $x - 1$, et nous avons
$(\pi) \cap (x) \cap (x - 1) = (\pi x (x - 1))$ dans $A[x]$.}.
Le lemme \ref{etale-cohomology-lemma-find-extension} et le fait que $R$ soit
absolument intégralement clos donnent donc le résultat.
\end{proof}

\begin{proposition}
\label{etale-cohomology-proposition-aic-vanishing}
Soit $R$ un anneau absolument intégralement clos.
Soit $M$ un groupe abélien fini.
Alors $H^i_\etale(\Spec(R), \underline{M}) = 0$ pour $i > 0$.
\end{proposition}

\begin{proof}
Comme tout groupe abélien fini admet une filtration finie dont les
sous-quotients sont cycliques d'ordre premier, nous pouvons supposer
que $M = \mathbf{Z}/\ell\mathbf{Z}$, où $\ell$ est un nombre premier.

\medskip\noindent
Observons que tous les anneaux locaux de $R$ sont strictement henséliens; voir
Compléments d'algèbre, lemme
\ref{more-algebra-lemma-absolutely-integrally-closed-strictly-henselian}.
De plus, tout localisé de $R$ est encore absolument intégralement clos
par Compléments d'algèbre, lemme
\ref{more-algebra-lemma-absolutely-integrally-closed-quotient-localization}.
Ainsi, le lemme \ref{etale-cohomology-lemma-gabber-criterion} nous dit qu'il suffit de
montrer que le noyau de
$$
H^1_\etale(D(f + g), \mathbf{Z}/\ell\mathbf{Z})
\longrightarrow
H^1_\etale(D(f(f + g)), \mathbf{Z}/\ell\mathbf{Z}) \oplus
H^1_\etale(D(g(f + g)), \mathbf{Z}/\ell\mathbf{Z})
$$
est nul pour tous $f, g \in R$. Après avoir remplacé $R$ par $R_{f + g}$, nous
nous ramenons à l'affirmation suivante : si
$\xi \in H^1_\etale(\Spec(R), \mathbf{Z}/\ell\mathbf{Z})$
et si $\Spec(R) = U \cup V$ est un recouvrement par deux ouverts affines
tel que $\xi|_U$ et $\xi|_V$ soient triviaux, alors $\xi = 0$.

\medskip\noindent
Soit $A = \mathbf{Z}[\zeta]$ comme ci-dessus.
Puisque $\mathbf{Z} \subset A$ est monogène, nous pouvons trouver un morphisme
d'anneaux $A \to R$. Dorénavant, nous considérons $R$ comme une $A$-algèbre
et $\Spec(R)$ comme un schéma sur $\Spec(A)$.
Après changement de base de la suite exacte courte (\ref{etale-cohomology-equation-ses})
à $\Spec(R)$ et passage à la cohomologie étale, nous obtenons
$$
G(R) \to H(R) \to
H^1_\etale(\Spec(R), \mathbf{Z}/\ell\mathbf{Z}) \to
H^1_\etale(\Spec(R), G)
$$
Nous utiliserons cela dans la suite de la preuve.

\medskip\noindent
Soit $\tau \in \Gamma(U \cap V, \mathbf{Z}/\ell\mathbf{Z})$
une section dont l'image par l'homomorphisme bord de la suite de Mayer--Vietoris
(lemme \ref{etale-cohomology-lemma-mayer-vietoris}) est $\xi$.
Pour $i = 0, 1, \ldots, \ell - 1$, soit $A_i \subset U \cap V$
la partie ouverte et fermée où $\tau$ prend la valeur $i \bmod \ell$.
Nous avons donc une décomposition en union disjointe finie
$$
U \cap V = A_0 \amalg \ldots \amalg A_{\ell - 1}
$$
telle que $\tau$ soit constante sur chaque $A_i$. Pour
$i = 0, 1, \ldots, \ell - 1$, notons
$\tau_i \in H^0(U \cap V, \mathbf{Z}/\ell\mathbf{Z})$
l'élément égal à $1$ sur $A_i$ et
à $0$ sur $A_j$ pour $j \not = i$.
Alors $\tau$ est une somme de multiples des $\tau_i$\footnote{À d'éventuelles
erreurs de calcul près, nous avons $\tau = \sum i \tau_i$.}.
Il suffit donc de montrer que la classe de cohomologie correspondant
à $\tau_i$ est triviale. Nous sommes ainsi ramenés au cas où
$\tau$ ne prend que deux valeurs distinctes, à savoir $1$ et $0$.

\medskip\noindent
Supposons que $\tau$ ne prenne que les valeurs $1$ et $0$. Écrivons
$$
U \cap V = A \amalg B
$$
où $A$ est le lieu où $\tau = 0$ et $B$ le lieu
où $\tau = 1$. Alors $A$ et $B$ sont des parties fermées disjointes.
Notons $\overline{A}$ et $\overline{B}$ les adhérences de $A$ et $B$ dans
$\Spec(R)$. Nous avons alors une ``banane'', c'est-à-dire
que nous avons
$$
\overline{A} \cap \overline{B} = Z_1 \amalg Z_2
$$
où $Z_1 \subset U$ et $Z_2 \subset V$ sont des parties fermées disjointes.
Posons $T_1 = \Spec(R) \setminus V$ et $T_2 = \Spec(R) \setminus U$.
Observons que $Z_1 \subset T_1 \subset U$, que $Z_2 \subset T_2 \subset V$ et que
$T_1 \cap T_2 = \emptyset$. Au niveau des espaces topologiques, nous pouvons écrire
$$
\Spec(R) = \overline{A} \cup \overline{B} \cup T_1 \cup T_2
$$
Nous suggérons de faire un dessin pour visualiser cette situation.
Pour prouver que $\xi$ est nul, nous pouvons remplacer, et remplaçons, $R$ par
son réduit (proposition \ref{etale-cohomology-proposition-topological-invariance}).
Dans la suite, nous considérons $A$, $\overline{A}$, $B$, $\overline{B}$, $T_1$, $T_2$
comme des sous-schémas fermés réduits de $\Spec(R)$. Munissons ensuite $Z_1$
et $Z_2$ des structures de schémas
$$
Z_1 = \overline{A} \cap (\overline{B} \cup T_1)
\quad\text{et}\quad
Z_2 = \overline{A} \cap (\overline{B} \cup T_2)
$$
(réunions et intersections schématiques comme dans Morphismes, définition
\ref{morphisms-definition-scheme-theoretic-intersection-union}).

\medskip\noindent
Notons $X$ le torseur sous $G$ sur $\Spec(R)$ correspondant à l'image
de $\xi$ dans $H^1(\Spec(R), G)$. Si $X$ est trivial, alors $\xi$
provient d'un élément $h \in H(R)$ (voir la suite exacte de cohomologie
ci-dessus). Mais le lemme \ref{etale-cohomology-lemma-lift-points-H-to-G} permet alors
de relever $h$ en un élément de $G(R)$, et nous concluons
que $\xi = 0$, comme souhaité. Notre but est donc de prouver que $X$ est trivial.

\medskip\noindent
Rappelons que le plongement $\mathbf{Z}/\ell \mathbf{Z} \to G(R)$
envoie $i \bmod \ell$ sur $\sigma_i \in G(R)$. Observons que $\overline{A}$
est le spectre d'un anneau absolument intégralement clos (à savoir
un quotient de $R$). Par le
lemme \ref{etale-cohomology-lemma-interpolate}, nous pouvons trouver $g \in G(\overline{A})$ tel que
$g|_{\overline{A} \cap Z_1} = \sigma_0$ et
$g|_{\overline{A} \cap Z_2} = \sigma_1$ (schématiquement).
Nous pouvons alors définir
\begin{enumerate}
\item $g_1 \in G(U)$ qui vaut $g$ sur $\overline{A} \cap U$,
$\sigma_0$ sur $\overline{B} \cap U$ et $\sigma_0$ sur $T_1$, et
\item $g_2 \in G(V)$ qui vaut $g$ sur $\overline{A} \cap V$,
$\sigma_1$ sur $\overline{B} \cap V$ et $\sigma_1$ sur $T_2$.
\end{enumerate}
En effet, pour trouver $g_1$ comme dans (1), nous recollons la section
$\sigma_0$ sur $\Omega = (\overline{B} \cup T_1) \cap U$
avec la restriction de la section $g$ à $\Omega' = \overline{A} \cap U$.
Notons que $U = \Omega \cup \Omega'$ (schématiquement),
car $U$ est réduit et $\Omega \cap \Omega' = Z_1$ (schématiquement)
par notre choix de $Z_1$. Ainsi, d'après
Morphismes, lemme \ref{morphisms-lemma-scheme-theoretic-union},
$U$ est la somme amalgamée de $\Omega$ et $\Omega'$
au-dessus de $Z_1$. Nous pouvons donc trouver $g_1$.
L'existence de $g_2$ dans (2) se démontre de même.
Nous avons alors
$$
\tau = g_2|_{A \cup B} - g_1|_{A \cup B} \quad(\text{addition selon la loi de groupe})
$$
et nous voyons que $X$ est trivial, ce qui achève la preuve.
\end{proof}




















\section{Analogue affine du changement de base propre}
\label{etale-cohomology-section-gabber-affine-proper}

\noindent
Dans cette section, nous étudions un résultat d'Ofer Gabber, voir
\cite{gabber-affine-proper}. Il a aussi été démontré par Roland Huber, voir
\cite{Huber-henselian}. Nous avons déjà effectué une partie du travail nécessaire
à la démonstration de Gabber dans la section \ref{etale-cohomology-section-strictly-henselian-local-rings}.

\begin{lemma}
\label{etale-cohomology-lemma-efface-cohomology-on-closed-by-finite-cover}
Soit $X$ un schéma affine. Soit $\mathcal{F}$ un faisceau abélien de torsion
sur $X_\etale$. Soit $Z \subset X$ un sous-schéma fermé. Soit
$\xi \in H^q_\etale(Z, \mathcal{F}|_Z)$ pour un $q > 0$.
Il existe alors un morphisme injectif $\mathcal{F} \to \mathcal{F}'$
de faisceaux abéliens de torsion sur $X_\etale$ tel que
l'image de $\xi$ dans $H^q_\etale(Z, \mathcal{F}'|_Z)$ soit nulle.
\end{lemma}

\begin{proof}
D'après les lemmes \ref{etale-cohomology-lemma-torsion-colimit-constructible} et \ref{etale-cohomology-lemma-colimit},
on peut trouver un morphisme $\mathcal{G} \to \mathcal{F}$ où $\mathcal{G}$
est un faisceau abélien constructible et où $\xi$ provient d'un élément $\zeta$ de
$H^q_\etale(Z, \mathcal{G}|_Z)$. Supposons qu'on puisse trouver un morphisme injectif
$\mathcal{G} \to \mathcal{G}'$ de faisceaux abéliens de torsion sur $X_\etale$
tel que l'image de $\zeta$ dans $H^q_\etale(Z, \mathcal{G}'|_Z)$ soit nulle.
On peut alors prendre pour $\mathcal{F}'$ la somme amalgamée
$$
\mathcal{F}' = \mathcal{G}' \amalg_{\mathcal{G}} \mathcal{F}
$$
et l'on obtient la conclusion du lemme. (Observons que la restriction
à $Z$ est exacte, donc commute aux limites et colimites finies, et qu'en outre
elle commute aux colimites quelconques comme adjoint à gauche de l'image directe.)
On peut donc supposer $\mathcal{F}$ constructible.

\medskip\noindent
Supposons $\mathcal{F}$ constructible. D'après le
lemme \ref{etale-cohomology-lemma-constructible-maps-into-constant-general},
il suffit de démontrer le résultat lorsque $\mathcal{F}$
est de la forme $f_*\underline{M}$, où $M$ est un groupe abélien fini
et $f : Y \to X$ un morphisme fini de présentation finie
(de tels faisceaux sont encore constructibles d'après le
lemme \ref{etale-cohomology-lemma-finite-pushforward-constructible},
mais nous n'en aurons pas besoin).
Comme la formation de $f_*$ commute à tout changement de base
(lemme \ref{etale-cohomology-lemma-finite-pushforward-commutes-with-base-change}),
la restriction de $f_*\underline{M}$ à $Z$ est
égale à l'image directe de $\underline{M}$ par
$Y \times_X Z \to Z$. Par la suite spectrale de Leray
(proposition \ref{etale-cohomology-proposition-leray})
et l'annulation des images directes supérieures
(proposition \ref{etale-cohomology-proposition-finite-higher-direct-image-zero}),
on obtient
$$
H^q_\etale(Z, f_*\underline{M}|_Z) = H^q_\etale(Y \times_X Z, \underline{M}).
$$
D'après le lemme \ref{etale-cohomology-lemma-integral-cover-trivial-cohomology},
on peut trouver un morphisme fini surjectif $Y' \to Y$ de présentation finie
tel que l'image de $\xi$ dans $H^q(Y' \times_X Z, \underline{M})$ soit nulle.
Notons $f' : Y' \to X$ le composé $Y' \to Y \to X$. Nous affirmons que
le morphisme
$$
f_*\underline{M} \longrightarrow f'_*\underline{M}
$$
est injectif, ce qui achève la démonstration d'après ce qui précède.
Pour établir cette injectivité, il suffit de regarder les fibres. En effet,
si $\overline{x} : \Spec(k) \to X$ est un point géométrique, alors
$$
(f_*\underline{M})_{\overline{x}} =
\bigoplus\nolimits_{f(\overline{y}) = \overline{x}} M
$$
d'après la proposition \ref{etale-cohomology-proposition-finite-higher-direct-image-zero},
et de même pour l'autre faisceau.
Comme $Y' \to Y$ est surjectif et fini, l'application induite
entre les points géométriques relevant $\overline{x}$ est
elle aussi surjective, d'où la conclusion.
\end{proof}

\noindent
Le lemme précédent traitera les groupes de cohomologie supérieure dans
le résultat de Gabber. Le lemme suivant servira à traiter les
sections globales.

\begin{lemma}
\label{etale-cohomology-lemma-gabber-h0}
Soit $X$ un schéma quasi-compact et quasi-séparé.
Soit $i : Z \to X$ une immersion fermée. Supposons que
\begin{enumerate}
\item pour tout faisceau $\mathcal{F}$ sur $X_{Zar}$, l'application
$\Gamma(X, \mathcal{F}) \to \Gamma(Z, i^{-1}\mathcal{F})$
soit bijective, et
\item pour tout morphisme fini $X' \to X$, l'hypothèse (1) soit satisfaite
pour $Z \times_X X' \to X'$.
\end{enumerate}
Alors, pour tout faisceau $\mathcal{F}$ sur $X_\etale$, on a
$\Gamma(X, \mathcal{F}) = \Gamma(Z, i^{-1}_{small}\mathcal{F})$.
\end{lemma}

\begin{proof}
Soit $\mathcal{F}$ un faisceau sur $X_\etale$. Il existe un morphisme canonique
(de changement de base)
$$
i^{-1}(\mathcal{F}|_{X_{Zar}})
\longrightarrow
(i_{small}^{-1}\mathcal{F})|_{Z_{Zar}}
$$
de faisceaux sur $Z_{Zar}$. Nous montrerons son injectivité en regardant
les fibres. La fibre du membre de gauche en $z \in Z$
est celle de $\mathcal{F}|_{X_{Zar}}$ en $z$. La fibre du membre de droite
est la limite inductive portant sur tous les voisinages étales élémentaires
$(U, u) \to (X, z)$ tels que $U \times_X Z \to Z$ admette une section sur
un voisinage de $z$. Comme les morphismes étales sont ouverts, l'image de
$U \to X$ est un voisinage ouvert $U_0$ de $z$ dans $X$. L'application
$\mathcal{F}(U_0) \to \mathcal{F}(U)$ est injective par la condition de faisceau
pour $\mathcal{F}$ relativement au recouvrement étale $U \to U_0$.
En passant à la limite inductive sur tous les $U$ et $U_0$, on obtient l'injectivité sur les fibres.

\medskip\noindent
Il en résulte, avec l'hypothèse (1), que l'application
$\Gamma(X, \mathcal{F}) \to \Gamma(Z, i^{-1}_{small}\mathcal{F})$
est injective. Par (2), il en va de même sur tout $X'$ fini sur $X$.

\medskip\noindent
Avant de démontrer la surjectivité, introduisons une notation. Pour
tout objet $U$ de $X_\etale$ et tout $s \in \mathcal{F}(U)$, nous
notons $i^{-1}s \in (i_{small}^{-1}\mathcal{F})(U \times_X Z)$
l'image inverse (plus précisément, c'est l'image de $s$ par le morphisme
de sites du lemme \ref{sites-lemma-recover}, composé avec le morphisme du
préfaisceau image inverse vers le faisceau image inverse). Cette construction est
compatible avec les images inverses par les morphismes de $X_\etale$ et par les morphismes (finis)
$X' \to X$, d'une manière dont nous laissons au lecteur le soin de préciser les détails.

\medskip\noindent
Soit $t \in \Gamma(Z, i^{-1}_{small}\mathcal{F})$. Par construction de
$i^{-1}_{small}\mathcal{F}$, il existe un recouvrement étale
$\{V_j \to Z\}_{j \in J}$, des morphismes étales $U_j \to X$, des sections
$s_j \in \mathcal{F}(U_j)$ et des morphismes $V_j \to U_j$ au-dessus de $X$
tels que $t|_{V_j}$ soit l'image inverse de $i^{-1}s_j$ par
$V_j \to U_j \times_X Z$. Comme $V_j \to U_j \times_X Z$ est étale,
son image est un sous-schéma ouvert. Comme $U_j \times_X Z \subset U_j$ est
fermé, on peut, après avoir remplacé $U_j$ par un sous-schéma ouvert, supposer
que $V_j \to U_j \times_X Z$ est surjectif. Alors $\{V_j \to U_j \times_X Z\}$
est un recouvrement étale, d'où $t|_{U_j \times_X Z} = i^{-1}s_j$.
Observons que tout sous-schéma fermé non vide $T \subset X$ rencontre $Z$
par l'hypothèse (1), appliquée par exemple au faisceau $(T \to X)_*\underline{\mathbf{Z}}$.
On voit donc que $\coprod_{j \in J} U_j \to X$ est surjectif.
D'après le lemme
\ref{more-morphisms-lemma-there-is-a-scheme-integral-over} de Compléments sur les morphismes,
on peut trouver un morphisme fini surjectif $X' \to X$
tel que le morphisme $X' \to X$ se factorise, localement pour la topologie de Zariski, par $\coprod_{j \in J} U_j \to X$.
Écrivons $X' = \bigcup_{\alpha \in A} W_\alpha$, où cette réunion est un recouvrement ouvert,
et soit $j : A \to J$ une application telle qu'il existe des morphismes
$g_\alpha : W_\alpha \to U_{j(\alpha)}$ au-dessus de $X$. Notons
$\mathcal{F}'$ l'image inverse de $\mathcal{F}$ sur $X'_\etale$.
Notons $s'_\alpha \in \mathcal{F}'(W_\alpha)$ l'image inverse
de $s_{j(\alpha)}$ par $g_\alpha$.
Posons $Z' = X' \times_X Z$. Notons $i' : Z' \to X'$
le changement de base de $i$. Alors $(i'_{small})^{-1}\mathcal{F}' =
(Z' \to Z)_{small}^{-1}i_{small}^{-1}\mathcal{F}$. Par cette identification,
on peut noter $t' \in
\Gamma(Z', (i')^{-1}_{small}\mathcal{F}')$
l'image inverse de $t$ sur $Z'$.
Par construction, $t'|_{W_\alpha \times_X Z}$ est égal à $(i')^{-1}s'_\alpha$.
On en déduit que $t'$ est une section du
sous-faisceau
$$
(i')^{-1}(\mathcal{F}'|_{X'_{Zar}})
\subset
(i'_{small})^{-1}\mathcal{F}')|_{Z'_{Zar}}
$$
Par l'hypothèse (2), la section $t'$ provient d'une section
$s' \in \Gamma(X', \mathcal{F}'|_{X'_{Zar}}) = \Gamma(X', \mathcal{F}')$.
Par l'injectivité démontrée au deuxième alinéa,
on en déduit que les deux images inverses de $s'$ sur
$X' \times_X X'$ sont égales (car cela est vrai pour
les images inverses de $t'$ sur $Z' \times_Z Z'$). On en conclut que
$s'$ provient d'une section de $\mathcal{F}$ sur $X$ grâce à la
remarque \ref{etale-cohomology-remark-cohomological-descent-finite}.
\end{proof}

\begin{lemma}
\label{etale-cohomology-lemma-connected-topological}
Soit $Z \subset X$ une partie fermée d'un espace topologique $X$.
Supposons que
\begin{enumerate}
\item $X$ soit un espace spectral
(Topologie, définition \ref{topology-definition-spectral-space}), et
\item pour $x \in X$, l'intersection $Z \cap \overline{\{x\}}$
soit connexe (en particulier non vide).
\end{enumerate}
Si $Z = Z_1 \amalg Z_2$, où les $Z_i$ sont fermés dans $Z$,
il existe alors une décomposition $X = X_1 \amalg X_2$ où les
$X_i$ sont fermés dans $X$ et $Z_i = Z \cap X_i$.
\end{lemma}

\begin{proof}
Observons que $Z_i$ est quasi-compact. L'ensemble $W_i$ des points
qui se spécialisent en un point de $Z_i$ est donc fermé dans la topologie constructible
d'après le lemme \ref{topology-lemma-make-spectral-space} de Topologie.
L'hypothèse (2) implique que $X = W_1 \amalg W_2$.
Soit $x \in \overline{W_1}$. D'après la partie (1) du lemme
\ref{topology-lemma-constructible-stable-specialization-closed} de Topologie,
il existe une spécialisation $x_1 \leadsto x$ avec $x_1 \in W_1$.
Ainsi $\overline{\{x\}} \subset \overline{\{x_1\}}$, et l'on voit
que $x \in W_1$. Autrement dit, le choix $X_i = W_i$ convient.
\end{proof}

\begin{lemma}
\label{etale-cohomology-lemma-h0-topological}
Soit $Z \subset X$ une partie fermée d'un espace topologique $X$.
Supposons que
\begin{enumerate}
\item $X$ soit un espace spectral
(Topologie, définition \ref{topology-definition-spectral-space}), et
\item pour $x \in X$, l'intersection $Z \cap \overline{\{x\}}$
soit connexe (en particulier non vide).
\end{enumerate}
Alors, pour tout faisceau $\mathcal{F}$ sur $X$, on a
$\Gamma(X, \mathcal{F}) = \Gamma(Z, \mathcal{F}|_Z)$.
\end{lemma}

\begin{proof}
Si $x \leadsto x'$ est une spécialisation de points, il existe un
morphisme canonique $\mathcal{F}_{x'} \to \mathcal{F}_x$, compatible avec les
sections sur les ouverts et fonctoriel en $\mathcal{F}$. Comme tout point
de $X$ se spécialise en un point de $Z$, il s'ensuit que
$\Gamma(X, \mathcal{F}) \to \Gamma(Z, \mathcal{F}|_Z)$ est injective.
La difficulté est d'en démontrer la surjectivité.

\medskip\noindent
Notons $\mathcal{B}$ l'ensemble de tous les ouverts quasi-compacts de $X$.
Écrivons $\mathcal{F}$ comme une limite inductive filtrante $\mathcal{F} = \colim \mathcal{F}_i$,
où chaque $\mathcal{F}_i$ est comme dans
Modules, équation (\ref{modules-equation-towards-constructible-sets}).
Voir Modules, lemme \ref{modules-lemma-filtered-colimit-constructibles}.
Alors $\mathcal{F}|_Z = \colim \mathcal{F}_i|_Z$, car la restriction à $Z$
est un adjoint à gauche (Catégories, lemme \ref{categories-lemma-adjoint-exact}, et
Faisceaux, lemme \ref{sheaves-lemma-f-map}).
D'après le lemme \ref{sheaves-lemma-directed-colimits-sections} de Faisceaux,
les foncteurs $\Gamma(X, -)$ et $\Gamma(Z, -)$ commutent aux limites inductives filtrantes.
On peut donc supposer que notre faisceau $\mathcal{F}$ est comme dans
Modules, équation (\ref{modules-equation-towards-constructible-sets}).

\medskip\noindent
Supposons disposer d'une injection $\mathcal{F} \subset \mathcal{G}$.
On a alors
$$
\Gamma(X, \mathcal{F}) =
\Gamma(Z, \mathcal{F}|_Z) \cap \Gamma(X, \mathcal{G})
$$
l'intersection étant prise dans $\Gamma(Z, \mathcal{G}|_Z)$.
Cela résulte de la première remarque de la démonstration, car on peut vérifier
si une section globale de $\mathcal{G}$ appartient à $\mathcal{F}$ en regardant
les fibres, et tout point de $X$ se spécialise en un point de $Z$.

\medskip\noindent
D'après le lemme \ref{modules-lemma-constructible-in-constant} de Modules,
il existe une injection $\mathcal{F} \to \prod (Z_i \to X)_*\underline{S_i}$,
où le produit est fini, $Z_i \subset X$ est fermé et $S_i$ est fini.
Il suffit donc de démontrer la surjectivité pour les faisceaux
$(Z_i \to X)_*\underline{S_i}$. Observons que
$$
\Gamma(X, (Z_i \to X)_*\underline{S_i}) = \Gamma(Z_i, \underline{S_i})
\quad\text{et}\quad
\Gamma(Z, (Z_i \to X)_*\underline{S_i}|_Z) =
\Gamma(Z \cap Z_i, \underline{S_i})
$$
De plus, les conditions (1) et (2) sont héritées par $Z_i$; c'est clair
pour (2), et cela résulte du
lemme \ref{topology-lemma-spectral-sub} de Topologie pour (1). Il
suffit donc de démontrer le lemme dans le cas d'un faisceau constant de valeur finie.
Ce cas est une reformulation du lemme \ref{etale-cohomology-lemma-connected-topological},
ce qui achève la démonstration.
\end{proof}

\begin{example}
\label{etale-cohomology-example-quinard}
Le lemme \ref{etale-cohomology-lemma-h0-topological} est faux si $X$ n'est pas spectral.
Voici un exemple : soit $Y$ un espace topologique $T_1$, et soit
$y \in Y$ un point non ouvert. Soit $X = Y \amalg \{ x \}$, muni de
la topologie dont les fermés sont $\emptyset$, $\{y\}$ et tous les
$F \amalg \{ x \}$, où $F$ est une partie fermée de $Y$. Alors
$Z = \{x, y\}$ est une partie fermée de $X$ qui satisfait à l'hypothèse (2)
du lemme \ref{etale-cohomology-lemma-h0-topological}. Mais $X$ est connexe, tandis que $Z$ ne l'est pas.
La conclusion du lemme échoue donc pour le faisceau constant
de valeur $\{0, 1\}$ sur $X$.
\end{example}

\begin{lemma}
\label{etale-cohomology-lemma-h0-henselian-pair}
Soit $(A, I)$ un couple hensélien. Posons $X = \Spec(A)$ et
$Z = \Spec(A/I)$. Pour tout faisceau $\mathcal{F}$ sur $X_\etale$,
on a $\Gamma(X, \mathcal{F}) = \Gamma(Z, \mathcal{F}|_Z)$.
\end{lemma}

\begin{proof}
Rappelons que le spectre de tout anneau est un espace spectral, voir
Algèbre, lemme \ref{algebra-lemma-spec-spectral}. D'après le
lemme
\ref{more-algebra-lemma-irreducible-henselian-pair-connected} de Compléments d'algèbre,
on voit que $\overline{\{x\}} \cap Z$ est connexe pour tout $x \in X$.
Le lemme \ref{etale-cohomology-lemma-h0-topological} montre que l'énoncé
est vrai pour les faisceaux sur $X_{Zar}$. Pour tout morphisme fini $X' \to X$,
on a $X' = \Spec(A')$ et $Z \times_X X' = \Spec(A'/IA')$,
où $(A', IA')$ est un couple hensélien, voir le lemme
\ref{more-algebra-lemma-integral-over-henselian-pair} de Compléments d'algèbre,
et l'on obtient le même énoncé pour les faisceaux sur $(X')_{Zar}$.
On peut donc conclure en appliquant le lemme \ref{etale-cohomology-lemma-gabber-h0}.
\end{proof}

\noindent
Enfin, nous pouvons énoncer et démontrer le théorème de Gabber.

\begin{theorem}[Gabber]
\label{etale-cohomology-theorem-gabber}
Soit $(A, I)$ un couple hensélien. Posons $X = \Spec(A)$ et
$Z = \Spec(A/I)$. Pour tout faisceau abélien de torsion $\mathcal{F}$ sur $X_\etale$,
on a $H^q_\etale(X, \mathcal{F}) = H^q_\etale(Z, \mathcal{F}|_Z)$.
\end{theorem}

\begin{proof}
Le résultat vaut pour $q = 0$ par le lemme \ref{etale-cohomology-lemma-h0-henselian-pair}.
Soit $q \geq 1$. Supposons le résultat démontré en tout degré $< q$.
Soit $\mathcal{F}$ un faisceau abélien de torsion. Soit
$\mathcal{F} \to \mathcal{F}'$
un morphisme injectif de faisceaux abéliens de torsion (qui sera choisi plus tard),
de conoyau $\mathcal{Q}$, de sorte que l'on ait la suite exacte courte
$$
0 \to \mathcal{F} \to \mathcal{F}' \to \mathcal{Q} \to 0
$$
de faisceaux abéliens de torsion sur $X_\etale$. Cela donne un morphisme de suites exactes longues
de cohomologie sur $X$ et $Z$, dont une partie est de la forme
$$
\xymatrix{
H^{q - 1}_\etale(X, \mathcal{F}') \ar[d] \ar[r] &
H^{q - 1}_\etale(X, \mathcal{Q}) \ar[d] \ar[r] &
H^q_\etale(X, \mathcal{F}) \ar[d] \ar[r] &
H^q_\etale(X, \mathcal{F}') \ar[d] \\
H^{q - 1}_\etale(Z, \mathcal{F}'|_Z) \ar[r] &
H^{q - 1}_\etale(Z, \mathcal{Q}|_Z) \ar[r] &
H^q_\etale(Z, \mathcal{F}|_Z) \ar[r] &
H^q_\etale(Z, \mathcal{F}'|_Z)
}
$$
Nous allons terminer la démonstration à l'aide de ce diagramme commutatif de groupes abéliens
à lignes exactes.

\medskip\noindent
Injectivité pour $\mathcal{F}$. Soit $\xi$ un élément non nul de
$H^q_\etale(X, \mathcal{F})$.
D'après le lemme \ref{etale-cohomology-lemma-efface-cohomology-on-closed-by-finite-cover}, appliqué avec
$Z = X$ (!), on peut trouver $\mathcal{F} \subset \mathcal{F}'$ de sorte que
l'image de $\xi$ dans le groupe de droite soit nulle. Alors $\xi$ est l'image d'un
élément de $H^{q - 1}_\etale(X, \mathcal{Q})$, et la bijectivité
en degré $q - 1$ implique que l'image de $\xi$ n'est pas nulle dans
$H^q_\etale(Z, \mathcal{F}|_Z)$.

\medskip\noindent
Surjectivité pour $\mathcal{F}$. Soit $\xi$ un élément de
$H^q_\etale(Z, \mathcal{F}|_Z)$.
D'après le lemme \ref{etale-cohomology-lemma-efface-cohomology-on-closed-by-finite-cover}, appliqué avec
$Z = Z$, on peut trouver $\mathcal{F} \subset \mathcal{F}'$ de sorte que
l'image de $\xi$ dans le groupe de droite soit nulle. Alors $\xi$ est l'image d'un
élément de $H^{q - 1}_\etale(Z, \mathcal{Q}|_Z)$, et la bijectivité
en degré $q - 1$ implique que $\xi$ appartient à l'image du morphisme vertical.
\end{proof}

\begin{lemma}
\label{etale-cohomology-lemma-vanishing-restriction-injective}
Soit $X$ un schéma à diagonale affine pouvant être recouvert par
$n + 1$ ouverts affines. Soit $Z \subset X$ un sous-schéma fermé.
Soit $\mathcal{A}$ un faisceau d'anneaux de torsion sur $X_\etale$
et soit $\mathcal{I}$ un faisceau injectif de $\mathcal{A}$-modules
sur $X_\etale$.
Alors $H^q_\etale(Z, \mathcal{I}|_Z) = 0$ pour $q > n$.
\end{lemma}

\begin{proof}
Nous procédons par récurrence sur $n$. Si $n = 0$, alors $X$ est affine.
Écrivons $X = \Spec(A)$ et $Z = \Spec(A/I)$. Soit $A^h$ la limite inductive filtrante
des $A$-algèbres étales $B$ telles que $A/I \to B/IB$ soit un isomorphisme.
Alors $(A^h, IA^h)$ est un couple hensélien et $A/I = A^h/IA^h$, voir
le lemme \ref{more-algebra-lemma-henselization} de Compléments d'algèbre
et sa démonstration. Posons $X^h = \Spec(A^h)$.
D'après le théorème \ref{etale-cohomology-theorem-gabber},
on a
$$
H^q_\etale(Z, \mathcal{I}|_Z) = H^q_\etale(X^h, \mathcal{I}|_{X^h})
$$
D'après le théorème \ref{etale-cohomology-theorem-colimit}, on a
$$
H^q_\etale(X^h, \mathcal{I}|_{X^h}) =
\colim_{A \to B} H^q_\etale(\Spec(B), \mathcal{I}|_{\Spec(B)})
$$
où la limite inductive porte sur les $A$-algèbres $B$ ci-dessus.
Comme les morphismes $\Spec(B) \to \Spec(A)$ sont étales,
la restriction $\mathcal{I}|_{\Spec(B)}$ est un faisceau injectif
de $\mathcal{A}|_{\Spec(B)}$-modules
(Cohomologie sur les sites, lemme \ref{sites-cohomology-lemma-cohomology-of-open}).
Les groupes de cohomologie du membre de droite sont donc nuls, ce qui donne le
résultat dans ce cas.

\medskip\noindent
Étape de récurrence. On peut utiliser Mayer--Vietoris.
Supposons en effet que $X = U \cup V$, où $U$ est une réunion de $n$ ouverts
affines et où $V$ est affine. Puisque la diagonale de $X$ est affine,
on voit que $U \cap V$ est la réunion de $n$ ouverts affines. Mayer--Vietoris
fournit une suite exacte
$$
H^{q - 1}_\etale(U \cap V \cap Z, \mathcal{I}|_Z) \to
H^q_\etale(Z, \mathcal{I}|_Z) \to
H^q_\etale(U \cap Z, \mathcal{I}|_Z) \oplus
H^q_\etale(V \cap Z, \mathcal{I}|_Z)
$$
et l'hypothèse de récurrence donne l'annulation pour $q > n$, comme voulu.
\end{proof}





\section{Cohomologie des faisceaux de torsion sur les courbes}
\label{etale-cohomology-section-vanishing-torsion}

\noindent
Le but de cette section est d'établir les résultats fondamentaux de finitude et d'annulation
pour la cohomologie des faisceaux de torsion sur les courbes, voir le
théorème \ref{etale-cohomology-theorem-vanishing-affine-curves}.
Dans la section \ref{etale-cohomology-section-vanishing-torsion-coefficients},
nous étudierons les faisceaux constructibles de modules de torsion
sur un anneau noethérien.

\begin{situation}
\label{etale-cohomology-situation-what-to-prove}
Ici, $k$ est un corps algébriquement clos, $X$ est un schéma séparé de type fini
de dimension $\leq 1$ sur $k$, et $\mathcal{F}$ est un faisceau abélien
de torsion sur $X_\etale$.
\end{situation}

\noindent
Dans la situation \ref{etale-cohomology-situation-what-to-prove},
nous voulons démontrer les assertions suivantes :
\begin{enumerate}
\item
\label{etale-cohomology-item-vanishing}
$H^q_\etale(X, \mathcal{F}) = 0$ pour $q > 2$,
\item
\label{etale-cohomology-item-vanishing-affine}
$H^q_\etale(X, \mathcal{F}) = 0$ pour $q > 1$ si $X$ est affine,
\item
\label{etale-cohomology-item-vanishing-p-p}
$H^q_\etale(X, \mathcal{F}) = 0$ pour $q > 1$ si $p = \text{char}(k) > 0$
et si $\mathcal{F}$ est de torsion $p$-primaire,
\item
\label{etale-cohomology-item-finite-prime-to-p}
$H^q_\etale(X, \mathcal{F})$ est fini si $\mathcal{F}$ est constructible
et de torsion première à $\text{char}(k)$,
\item
\label{etale-cohomology-item-finite-proper}
$H^q_\etale(X, \mathcal{F})$ est fini si $X$ est propre et si
$\mathcal{F}$ est constructible,
\item
\label{etale-cohomology-item-base-change-prime-to-p}
$H^q_\etale(X, \mathcal{F}) \to
H^q_\etale(X_{k'}, \mathcal{F}|_{X_{k'}})$ est un isomorphisme
pour toute extension $k'/k$ de corps algébriquement clos
si $\mathcal{F}$ est de torsion première à $\text{char}(k)$,
\item
\label{etale-cohomology-item-base-change-proper}
$H^q_\etale(X, \mathcal{F}) \to
H^q_\etale(X_{k'}, \mathcal{F}|_{X_{k'}})$ est un isomorphisme
pour toute extension $k'/k$ de corps algébriquement clos
si $X$ est propre,
\item
\label{etale-cohomology-item-surjective}
$H^2_\etale(X, \mathcal{F}) \to H^2_\etale(U, \mathcal{F})$
est surjectif pour tout ouvert $U \subset X$.
\end{enumerate}
Étant donnée une situation \ref{etale-cohomology-situation-what-to-prove},
nous dirons que
« les assertions (\ref{etale-cohomology-item-vanishing}) -- (\ref{etale-cohomology-item-surjective}) sont vérifiées »
si celles qui s'appliquent à la situation considérée sont vraies.
Nous commençons la démonstration par la conséquence suivante de notre calcul
de la cohomologie à coefficients constants.

\begin{lemma}
\label{etale-cohomology-lemma-constant-smooth-statements}
Dans la situation \ref{etale-cohomology-situation-what-to-prove},
supposons que $X$ soit lisse et que $\mathcal{F} = \underline{\mathbf{Z}/\ell\mathbf{Z}}$
pour un nombre premier $\ell$. Alors les assertions
(\ref{etale-cohomology-item-vanishing}) -- (\ref{etale-cohomology-item-surjective}) sont vérifiées
pour $\mathcal{F}$.
\end{lemma}

\begin{proof}
Puisque $X$ est lisse, nous voyons que $X$ est une réunion disjointe finie de
courbes lisses et de points. La cohomologie étale supérieure de tout faisceau
sur un point (le spectre de $k$) est nulle, par exemple d'après le
lemme \ref{etale-cohomology-lemma-compare-cohomology-point}
ou le lemme \ref{etale-cohomology-lemma-affine-only-closed-points}.
Nous pouvons donc supposer que $X$ est une courbe lisse.

\medskip\noindent
Cas I : $\ell$ est distinct de la caractéristique de $k$.
Ce cas résulte du
lemme \ref{etale-cohomology-lemma-cohomology-smooth-projective-curve}
(cas projectif) et du
lemme \ref{etale-cohomology-lemma-vanishing-cohomology-mu-smooth-curve}
(cas affine). L'assertion (\ref{etale-cohomology-item-base-change-prime-to-p})
sur la cohomologie et l'extension du corps de base algébriquement clos
résulte de ce que le genre $g$ et le nombre $r$
de « points manquants » ne changent pas quand on passe de $k$ à $k'$.
L'assertion (\ref{etale-cohomology-item-surjective}) résulte de ce que $H^2_\etale(U, \mathcal{F})$
est nul dès que $U \not = X$, car $U$ est alors affine
(Variétés, lemmes \ref{varieties-lemma-proper-minus-point} et
\ref{varieties-lemma-curve-affine-projective}).

\medskip\noindent
Cas II : $\ell$ est égal à la caractéristique de $k$.
L'annulation résulte du lemme \ref{etale-cohomology-lemma-vanishing-variety-char-p-p}.
Les assertions (\ref{etale-cohomology-item-finite-proper}) et (\ref{etale-cohomology-item-base-change-proper})
résultent du
lemme \ref{etale-cohomology-lemma-finiteness-proper-variety-char-p-p}.
\end{proof}

\begin{remark}[Invariance par extension du corps de base algébriquement clos]
\label{etale-cohomology-remark-invariance}
Soit $k$ un corps algébriquement clos de caractéristique $p > 0$.
Dans la section \ref{etale-cohomology-section-artin-schreier}, nous avons vu qu'il existe
une suite exacte
$$
k[x] \to k[x] \to
H^1_\etale(\mathbf{A}^1_k, \mathbf{Z}/p\mathbf{Z}) \to 0
$$
où la première flèche envoie $f(x)$ sur $f^p - f$. Un système de représentants
du conoyau est formé des polynômes
$$
\sum\nolimits_{p \not | n} \lambda_n x^n
$$
avec $\lambda_n \in k$. (Si $k$ n'est pas algébriquement clos,
il faut aussi y ajouter certaines constantes.) En particulier,
lorsque $k'/k$ est une extension de corps algébriquement clos,
l'application
$$
H^1_\etale(\mathbf{A}^1_k, \mathbf{Z}/p\mathbf{Z})
\to
H^1_\etale(\mathbf{A}^1_{k'}, \mathbf{Z}/p\mathbf{Z})
$$
n'est pas un isomorphisme en général. En particulier, l'application
$\pi_1(\mathbf{A}^1_{k'}) \to \pi_1(\mathbf{A}^1_k)$
entre groupes fondamentaux étales (référence à ajouter ultérieurement)
n'est pas non plus un isomorphisme. Ainsi, le type d'homotopie étale
de la droite affine dépend du corps de base algébriquement clos.
Le lemme \ref{etale-cohomology-lemma-constant-smooth-statements} ci-dessus montre
que ce phénomène ne se produit qu'en caractéristique $p$
avec des coefficients de torsion $p$-primaire.
\end{remark}

\begin{lemma}
\label{etale-cohomology-lemma-ses-statements}
Soit $k$ un corps algébriquement clos. Soit $X$ un schéma séparé de type
fini sur $k$ et de dimension $\leq 1$. Soit
$0 \to \mathcal{F}_1 \to \mathcal{F} \to \mathcal{F}_2 \to 0$
une suite exacte courte de faisceaux abéliens de torsion sur $X$.
Si les assertions (\ref{etale-cohomology-item-vanishing}) -- (\ref{etale-cohomology-item-surjective}) sont vérifiées
pour $\mathcal{F}_1$ et $\mathcal{F}_2$, elles le sont aussi
pour $\mathcal{F}$.
\end{lemma}

\begin{proof}
Cela résulte pour l'essentiel immédiatement des définitions et de la suite exacte longue
de cohomologie. Observons aussi que $\mathcal{F}$ est constructible
(resp.\ de torsion première à la caractéristique de $k$) si et seulement si
$\mathcal{F}_1$ et $\mathcal{F}_2$ sont tous deux constructibles
(resp.\ de torsion première à la caractéristique de $k$). Voir la
proposition \ref{etale-cohomology-proposition-constructible-over-noetherian}.
Nous omettons quelques détails.
\end{proof}

\begin{lemma}
\label{etale-cohomology-lemma-finite-pushforward-statements}
Soit $k$ un corps algébriquement clos. Soit $f : X \to Y$ un
morphisme fini de schémas séparés de type fini sur $k$ et de
dimension $\leq 1$. Soit $\mathcal{F}$ un faisceau abélien de torsion sur $X$.
Si les assertions (\ref{etale-cohomology-item-vanishing}) -- (\ref{etale-cohomology-item-surjective}) sont vérifiées
pour $\mathcal{F}$, elles le sont aussi pour $f_*\mathcal{F}$.
\end{lemma}

\begin{proof}
En effet, nous avons $H^q_\etale(X, \mathcal{F}) = H^q_\etale(Y, f_*\mathcal{F})$
par l'annulation de $R^qf_*$ pour $q > 0$
(proposition \ref{etale-cohomology-proposition-finite-higher-direct-image-zero}) et par
la suite spectrale de Leray
(Cohomologie sur les sites, lemme \ref{sites-cohomology-lemma-apply-Leray}).
Pour (\ref{etale-cohomology-item-surjective}), on utilise que la formation de $f_*$ commute à
tout changement de base
(lemme \ref{etale-cohomology-lemma-finite-pushforward-commutes-with-base-change}).
\end{proof}

\begin{lemma}
\label{etale-cohomology-lemma-restrict-to-open}
Dans la situation \ref{etale-cohomology-situation-what-to-prove}, supposons $\mathcal{F}$
constructible. Soit $j : X' \to X$ l'inclusion d'un sous-schéma ouvert dense.
Alors les assertions
(\ref{etale-cohomology-item-vanishing}) -- (\ref{etale-cohomology-item-surjective}) sont vérifiées pour $\mathcal{F}$
si et seulement si elles le sont pour $j_!j^{-1}\mathcal{F}$.
\end{lemma}

\begin{proof}
Puisque $X'$ est dense, nous voyons que $Z = X \setminus X'$ est de dimension $0$
et est donc un ensemble fini $Z = \{x_1, \ldots, x_n\}$ de points $k$-rationnels.
Considérons la suite exacte courte
$$
0 \to j_!j^{-1}\mathcal{F} \to \mathcal{F} \to i_*i^{-1}\mathcal{F} \to 0
$$
du lemme \ref{etale-cohomology-lemma-ses-associated-to-open}. Observons que
$H^q_\etale(X, i_*i^{-1}\mathcal{F}) = H^q_\etale(Z, i^{-1}\mathcal{F})$.
En effet, $i : Z \to X$ est une immersion fermée, donc finie,
et nous avons donc l'annulation de $R^qi_*$ pour $q > 0$ d'après la
proposition \ref{etale-cohomology-proposition-finite-higher-direct-image-zero};
l'égalité résulte alors de la suite spectrale de Leray
(Cohomologie sur les sites, lemme \ref{sites-cohomology-lemma-apply-Leray}).
Puisque $Z$ est une réunion disjointe de spectres de corps algébriquement
clos, nous en déduisons que $H^q_\etale(Z, i^*\mathcal{F}) = 0$
pour $q > 0$ et que
$$
H^0_\etale(Z, i^{-1}\mathcal{F}) =
\bigoplus\nolimits_{i=1}^n \mathcal{F}_{x_i}
$$
est fini, car $\mathcal{F}_{x_i}$ est fini en raison de
l'hypothèse de constructibilité de $\mathcal{F}$.
La suite exacte longue de cohomologie donne une suite exacte
$$
0 \to
H^0_\etale(X, j_!j^{-1}\mathcal{F}) \to
H^0_\etale(X, \mathcal{F}) \to
H^0_\etale(Z, i^{-1}\mathcal{F}) \to
H^1_\etale(X, j_!j^{-1}\mathcal{F}) \to
H^1_\etale(X, \mathcal{F}) \to 0
$$
et des isomorphismes $H^q_\etale(X, j_!j^{-1}\mathcal{F}) \to
H^q_\etale(X, \mathcal{F})$ pour $q > 1$.

\medskip\noindent
Il est maintenant facile de déduire que chacune des assertions 
(\ref{etale-cohomology-item-vanishing}) -- (\ref{etale-cohomology-item-surjective})
est vérifiée pour $\mathcal{F}$ si et seulement si elle l'est pour
$j_!j^{-1}\mathcal{F}$. Faisons quelques remarques pour aider le lecteur :
(a) si $\mathcal{F}$ est de torsion première à la caractéristique
de $k$, alors $j_!j^{-1}\mathcal{F}$ l'est aussi,
(b) le faisceau $j_!j^{-1}\mathcal{F}$ est constructible,
(c) nous avons $H^0_\etale(Z, i^{-1}\mathcal{F}) =
H^0_\etale(Z_{k'}, i^{-1}\mathcal{F}|_{Z_{k'}})$, et (d)
si $U \subset X$ est un ouvert, alors $U' = U \cap X'$ est dense dans $U$.
\end{proof}

\begin{lemma}
\label{etale-cohomology-lemma-even-easier}
Dans la situation \ref{etale-cohomology-situation-what-to-prove}, supposons $X$ lisse.
Soit $j : U \to X$ une immersion ouverte. Soit $\ell$ un nombre premier.
Posons $\mathcal{F} = j_!\underline{\mathbf{Z}/\ell\mathbf{Z}}$.
Alors les assertions (\ref{etale-cohomology-item-vanishing}) -- (\ref{etale-cohomology-item-surjective}) sont vérifiées
pour $\mathcal{F}$.
\end{lemma}

\begin{proof}
Puisque $X$ est lisse, il est une réunion disjointe de courbes lisses;
nous pouvons donc supposer que $X$ est une courbe (c'est-à-dire irréductible).
Alors soit $U = \emptyset$, auquel cas il n'y a rien à démontrer,
soit $U \subset X$ est dense. Dans ce cas, le lemme résulte des
lemmes \ref{etale-cohomology-lemma-constant-smooth-statements} et
\ref{etale-cohomology-lemma-restrict-to-open}.
\end{proof}

\begin{lemma}
\label{etale-cohomology-lemma-somewhat-easier}
Dans la situation \ref{etale-cohomology-situation-what-to-prove}, supposons $X$ réduit.
Soit $j : U \to X$ une immersion ouverte. Soit $\ell$ un nombre premier
et posons $\mathcal{F} = j_!\underline{\mathbf{Z}/\ell\mathbf{Z}}$.
Alors les assertions (\ref{etale-cohomology-item-vanishing}) -- (\ref{etale-cohomology-item-surjective}) sont vérifiées
pour $\mathcal{F}$.
\end{lemma}

\begin{proof}
La différence avec le lemme \ref{etale-cohomology-lemma-even-easier} est qu'ici nous ne
supposons pas $X$ lisse. Soit $\nu : X^\nu \to X$ le morphisme de
normalisation. Alors $\nu$ est fini
(Variétés, lemme \ref{varieties-lemma-normalization-locally-algebraic})
et $X^\nu$ est lisse
(Variétés, lemme \ref{varieties-lemma-regular-point-on-curve}).
Soit $j^\nu : U^\nu \to X^\nu$ l'image inverse de $U$.
D'après le lemme \ref{etale-cohomology-lemma-even-easier}, le résultat vaut pour
$j^\nu_!\underline{\mathbf{Z}/\ell\mathbf{Z}}$.
D'après le lemme \ref{etale-cohomology-lemma-finite-pushforward-statements},
le résultat vaut pour $\nu_*j^\nu_!\underline{\mathbf{Z}/\ell\mathbf{Z}}$.
En général, il n'est pas vrai que
$\nu_*j^\nu_!\underline{\mathbf{Z}/\ell\mathbf{Z}}$ soit égal à
$j_!\underline{\mathbf{Z}/\ell\mathbf{Z}}$
mais nous pouvons contourner cette difficulté comme suit.
Puisque $X$ est réduit, le morphisme $\nu : X^\nu \to X$
est un isomorphisme au-dessus d'un ouvert dense $j' : X' \to X$
(Variétés, lemme \ref{varieties-lemma-normalization-locally-algebraic}).
Au-dessus de cet ouvert, les deux faisceaux coïncident :
$$
(j')^{-1}(\nu_*j^\nu_!\underline{\mathbf{Z}/\ell\mathbf{Z}})
=
(j')^{-1}(j_!\underline{\mathbf{Z}/\ell\mathbf{Z}})
$$
En appliquant deux fois le lemme \ref{etale-cohomology-lemma-restrict-to-open}
à $j' : X' \to X$ et aux faisceaux ci-dessus,
nous concluons.
\end{proof}

\begin{lemma}
\label{etale-cohomology-lemma-vanishing-easier}
Dans la situation \ref{etale-cohomology-situation-what-to-prove}, supposons $X$ réduit.
Soit $j : U \to X$ une immersion ouverte avec $U$ connexe. Soit
$\ell$ un nombre premier. Soit $\mathcal{G}$ un faisceau localement constant
fini de $\mathbf{F}_\ell$-espaces vectoriels sur $U$. Posons
$\mathcal{F} = j_!\mathcal{G}$. Alors les assertions
(\ref{etale-cohomology-item-vanishing}) -- (\ref{etale-cohomology-item-surjective}) sont vérifiées pour $\mathcal{F}$.
\end{lemma}

\begin{proof}
Soit $f : V \to U$ un morphisme fini étale de degré premier à $\ell$
comme dans le lemme \ref{etale-cohomology-lemma-pullback-filtered}. La discussion de la
section \ref{etale-cohomology-section-trace-method} donne des applications
$$
\mathcal{G} \to f_*f^{-1}\mathcal{G} \to \mathcal{G}
$$
dont le composé est un isomorphisme. Il suffit donc de démontrer le
lemme avec $\mathcal{F} = j_!f_*f^{-1}\mathcal{G}$.
En vertu du théorème principal de Zariski
(Compléments sur les morphismes, lemme
\ref{more-morphisms-lemma-quasi-finite-separated-pass-through-finite})
nous pouvons choisir un diagramme
$$
\xymatrix{
V \ar[r]_{j'} \ar[d]_f & Y \ar[d]^{\overline{f}} \\
U \ar[r]^j & X
}
$$
où $\overline{f} : Y \to X$ est fini et $j'$ une immersion ouverte
d'image dense. Nous pouvons remplacer $Y$ par son réduit (cela ne
change pas $V$, puisque $V$ est réduit car étale sur $U$).
Comme $f$ est fini et $V$ dense dans $Y$, nous avons $V = U \times_X Y$. D'après le
lemme \ref{etale-cohomology-lemma-compatible-shriek-push-finite}, nous avons
$$
j_!f_*f^{-1}\mathcal{G} = \overline{f}_*j'_!f^{-1}\mathcal{G}
$$
D'après le lemme \ref{etale-cohomology-lemma-finite-pushforward-statements}, il suffit de
considérer $j'_!f^{-1}\mathcal{G}$.
L'existence de la filtration donnée par le
lemme \ref{etale-cohomology-lemma-pullback-filtered},
le fait que $j'_!$ est exact et le
lemme \ref{etale-cohomology-lemma-ses-statements}
nous ramènent au cas
$\mathcal{F} = j'_!\underline{\mathbf{Z}/\ell\mathbf{Z}}$
qui est le lemme \ref{etale-cohomology-lemma-somewhat-easier}.
\end{proof}


%10.20.09

\begin{theorem}
\label{etale-cohomology-theorem-vanishing-affine-curves}
Si $k$ est un corps algébriquement clos, $X$ un schéma séparé de type fini
et de dimension $\leq 1$ sur $k$, et $\mathcal{F}$ un faisceau abélien
de torsion sur $X_\etale$, alors
\begin{enumerate}
\item
$H^q_\etale(X, \mathcal{F}) = 0$ pour $q > 2$,
\item
$H^q_\etale(X, \mathcal{F}) = 0$ pour $q > 1$ si $X$ est affine,
\item
$H^q_\etale(X, \mathcal{F}) = 0$ pour $q > 1$ si $p = \text{char}(k) > 0$
et si $\mathcal{F}$ est de torsion $p$-primaire,
\item
$H^q_\etale(X, \mathcal{F})$ est fini si $\mathcal{F}$ est constructible
et de torsion première à $\text{char}(k)$,
\item
$H^q_\etale(X, \mathcal{F})$ est fini si $X$ est propre et si
$\mathcal{F}$ est constructible,
\item
$H^q_\etale(X, \mathcal{F}) \to
H^q_\etale(X_{k'}, \mathcal{F}|_{X_{k'}})$ est un isomorphisme
pour toute extension $k'/k$ de corps algébriquement clos
si $\mathcal{F}$ est de torsion première à $\text{char}(k)$,
\item
$H^q_\etale(X, \mathcal{F}) \to
H^q_\etale(X_{k'}, \mathcal{F}|_{X_{k'}})$ est un isomorphisme
pour toute extension $k'/k$ de corps algébriquement clos
si $X$ est propre,
\item
$H^2_\etale(X, \mathcal{F}) \to H^2_\etale(U, \mathcal{F})$
est surjectif pour tout ouvert $U \subset X$.
\end{enumerate}
\end{theorem}

\begin{proof}
Le théorème affirme que, dans la situation \ref{etale-cohomology-situation-what-to-prove},
les assertions (\ref{etale-cohomology-item-vanishing}) -- (\ref{etale-cohomology-item-surjective}) sont vérifiées.
Nous commençons par remplacer $X$ par son réduit, ce qui est permis
par la proposition \ref{etale-cohomology-proposition-topological-invariance}.
D'après le lemme \ref{etale-cohomology-lemma-torsion-colimit-constructible}, nous pouvons écrire
$\mathcal{F}$ comme limite inductive filtrante de faisceaux abéliens constructibles.
La cohomologie commute aux limites inductives, voir le lemme \ref{etale-cohomology-lemma-colimit}.
De plus, l'image inverse par $X_{k'} \to X$ commute aux limites inductives en tant
qu'adjoint à gauche. Il suffit donc de démontrer les assertions pour un faisceau constructible.

\medskip\noindent
Dans ce paragraphe, nous utilisons le lemme \ref{etale-cohomology-lemma-ses-statements} sans le
mentionner davantage. Écrivons
$\mathcal{F} = \mathcal{F}_1 \oplus \ldots \oplus \mathcal{F}_r$
où $\mathcal{F}_i$ est $\ell_i$-primaire pour un nombre premier $\ell_i$; nous pouvons
supposer que $\ell^n$ annule $\mathcal{F}$ pour un nombre premier $\ell$. Considérons la
suite exacte
$$
0 \to \mathcal{F}[\ell] \to \mathcal{F} \to \mathcal{F}/\mathcal{F}[\ell] \to 0.
$$
Nous voyons ainsi qu'il suffit de supposer $\mathcal{F}$ de $\ell$-torsion.
Cela signifie que $\mathcal{F}$ est un faisceau constructible de
$\mathbf{F}_\ell$-espaces vectoriels pour un nombre premier $\ell$.

\medskip\noindent
Par définition, cela signifie qu'il existe un ouvert dense $U \subset X$
tel que $\mathcal{F}|_U$ soit un faisceau localement constant fini de
$\mathbf{F}_\ell$-espaces vectoriels. Comme $\dim(X) \leq 1$, nous pouvons
supposer, après avoir rétréci $U$, que $U = U_1 \amalg \ldots \amalg U_n$
est une réunion disjointe de schémas irréductibles (il suffit de supprimer les points
fermés appartenant aux intersections d'au moins $2$ composantes de $U$).
D'après le lemme \ref{etale-cohomology-lemma-restrict-to-open}, nous nous ramenons au cas
$\mathcal{F} = j_!\mathcal{G}$ où $\mathcal{G}$ est un faisceau
localement constant fini de $\mathbf{F}_\ell$-espaces vectoriels sur $U$.

\medskip\noindent
Comme nous avons choisi $U = U_1 \amalg \ldots \amalg U_n$ avec les $U_i$ irréductibles,
nous avons
$$
j_!\mathcal{G} =
j_{1!}(\mathcal{G}|_{U_1}) \oplus \ldots \oplus
j_{n!}(\mathcal{G}|_{U_n})
$$
où $j_i : U_i \to X$ est le morphisme d'inclusion.
Le cas de $j_{i!}(\mathcal{G}|_{U_i})$ est traité dans le
lemme \ref{etale-cohomology-lemma-vanishing-easier}.
\end{proof}

\begin{theorem}
\label{etale-cohomology-theorem-vanishing-curves}
Soit $X$ un schéma de type fini et de dimension $1$ sur un
corps algébriquement clos $k$. Soit $\mathcal{F}$ un faisceau de torsion
sur $X_\etale$. Alors
$$
H_\etale^q(X, \mathcal{F}) = 0, \quad \forall q \geq 3.
$$
Si $X$ est affine, alors $H_\etale^2(X, \mathcal{F}) = 0$ également.
\end{theorem}

\begin{proof}
Si $X$ est séparé, cela résulte immédiatement du théorème plus précis
\ref{etale-cohomology-theorem-vanishing-affine-curves}.
Si $X$ n'est pas séparé, choisissons un recouvrement ouvert affine
$X = X_1 \cup \ldots \cup X_n$. Par récurrence sur $n$, nous pouvons supposer
l'annulation vérifiée sur $U = X_1 \cup \ldots \cup X_{n - 1}$.
Alors Mayer--Vietoris (lemme \ref{etale-cohomology-lemma-mayer-vietoris}) donne
$$
H^2_\etale(U, \mathcal{F}) \oplus H^2_\etale(X_n, \mathcal{F}) \to
H^2_\etale(U \cap X_n, \mathcal{F}) \to
H^3_\etale(X, \mathcal{F}) \to 0
$$
Cependant, puisque $U \cap X_n$ est un ouvert d'un schéma affine,
il est affine en vertu de notre hypothèse de dimension. Le groupe
$H^2_\etale(U \cap X_n, \mathcal{F})$ s'annule donc
d'après le théorème \ref{etale-cohomology-theorem-vanishing-affine-curves}.
\end{proof}

\begin{lemma}
\label{etale-cohomology-lemma-base-change-dim-1-separably-closed}
Soit $k'/k$ une extension de corps séparablement clos.
Soit $X$ un schéma propre sur $k$ de dimension $\leq 1$.
Soit $\mathcal{F}$ un faisceau abélien de torsion sur $X$.
Alors l'application $H^q_\etale(X, \mathcal{F}) \to
H^q_\etale(X_{k'}, \mathcal{F}|_{X_{k'}})$ est un isomorphisme
pour $q \geq 0$.
\end{lemma}

\begin{proof}
Nous l'avons vu pour les corps algébriquement clos dans le
théorème \ref{etale-cohomology-theorem-vanishing-affine-curves}.
Étant donnés $k \subset k'$ comme dans l'énoncé du lemme, nous pouvons
choisir un diagramme
$$
\xymatrix{
k' \ar[r] & \overline{k}' \\
k \ar[u] \ar[r] & \overline{k} \ar[u]
}
$$
où $k \subset \overline{k}$ et $k' \subset \overline{k}'$ sont
les clôtures algébriques. Puisque $k$ et $k'$ sont séparablement clos,
les extensions de corps
$\overline{k}/k$ et $\overline{k}'/k'$
sont algébriques et purement inséparables. Dans ce cas, les morphismes
$X_{\overline{k}} \to X$ et $X_{\overline{k}'} \to X_{k'}$
sont des homéomorphismes universels. Ainsi, la cohomologie de $\mathcal{F}$
peut se calculer sur $X_{\overline{k}}$ et la cohomologie
de $\mathcal{F}|_{X_{k'}}$ peut se calculer sur $X_{\overline{k}'}$,
voir la proposition \ref{etale-cohomology-proposition-topological-invariance}.
Nous déduisons donc le cas général du cas des corps algébriquement
clos.
\end{proof}






\section{Cohomologie des modules de torsion sur les courbes}
\label{etale-cohomology-section-vanishing-torsion-coefficients}

\noindent
Dans cette section, nous reprenons les arguments de la
section \ref{etale-cohomology-section-vanishing-torsion}
pour les faisceaux constructibles de modules de torsion
sur un anneau noethérien. Nous commençons par l'étape la plus intéressante.

\begin{lemma}
\label{etale-cohomology-lemma-constant-statements-coefficients}
Soit $\Lambda$ un anneau noethérien, soit $M$ un $\Lambda$-module de type fini annulé
par un entier $n > 0$, soit $k$ un corps algébriquement clos et soit $X$ un
schéma séparé de type fini et de dimension $\leq 1$ sur $k$.
Alors
\begin{enumerate}
\item $H^q_\etale(X, \underline{M})$ est un $\Lambda$-module de type fini
si $n$ est premier à $\text{char}(k)$,
\item $H^q_\etale(X, \underline{M})$ est un $\Lambda$-module de type fini
si $X$ est propre.
\end{enumerate}
\end{lemma}

\begin{proof}
Si $n = \ell n'$ pour un nombre premier $\ell$, nous obtenons une suite
exacte courte $0 \to M[\ell] \to M \to M' \to 0$ de $\Lambda$-modules
finis, et $M'$ est annulé par $n'$.
Cela fournit une suite exacte courte correspondante de faisceaux constants,
laquelle donne à son tour une suite
exacte de modules de cohomologie
$$
H^q_\etale(X, \underline{M[\ell]}) \to
H^q_\etale(X, \underline{M}) \to
H^q_\etale(X, \underline{M'})
$$
Ainsi, si nous pouvons démontrer le résultat lorsque $M$ est annulé par
un nombre premier, une récurrence sur $n$ permet de conclure.

\medskip\noindent
Soit $\ell$ un nombre premier qui annule $M$.
Nous pouvons alors remplacer $\Lambda$ par la $\mathbf{F}_\ell$-algèbre
$\Lambda/\ell \Lambda$. En effet,
la cohomologie de $\underline{M}$ considéré comme faisceau de
$\Lambda$-modules est la même que sa cohomologie comme faisceau de
$\Lambda/\ell \Lambda$-modules, par exemple d'après
Cohomologie sur les sites, lemme
\ref{sites-cohomology-lemma-cohomology-modules-abelian-agree}.

\medskip\noindent
Supposons que $\ell$ soit un nombre premier qui annule $M$
et $\Lambda$.
Ramenons-nous au cas où $M$ est un $\Lambda$-module libre de type fini.
Choisissons pour cela une suite exacte courte
$$
0 \to N \to \Lambda^{\oplus m} \to M \to 0
$$
Elle détermine une suite exacte
$$
H^q_\etale(X, \underline{\Lambda^{\oplus m}}) \to
H^q_\etale(X, \underline{M}) \to
H^{q + 1}_\etale(X, \underline{N})
$$
Une récurrence descendante sur $q$ donne le résultat pour $M$
si nous le connaissons pour $\Lambda^{\oplus m}$.
Nous utilisons ici le fait que nos groupes de cohomologie
s'annulent en degrés $> 2$
d'après le théorème \ref{etale-cohomology-theorem-vanishing-affine-curves}.

\medskip\noindent
Soit $\ell$ un nombre premier et supposons que $\ell$ annule $\Lambda$.
Il reste à montrer que les groupes de cohomologie
$H^q_\etale(X, \underline{\Lambda})$ sont des $\Lambda$-modules de type fini.
Nous allons employer une astuce; l'argument « correct » utilise
un théorème sur les coefficients que nous démontrerons plus loin.
Choisissons une base $\Lambda = \bigoplus_{i \in I} \mathbf{F}_\ell e_i$
telle que $e_0 = 1$ pour un certain $0 \in I$.
Le choix de cette base détermine un isomorphisme
$$
\underline{\Lambda} = \bigoplus_{i \in I} \underline{\mathbf{F}_\ell} e_i
$$
de faisceaux sur $X_\etale$. Nous voyons donc que
$$
H^q_\etale(X, \underline{\Lambda}) =
H^q_\etale(X, \bigoplus_{i \in I} \underline{\mathbf{F}_\ell} e_i) =
\bigoplus_{i \in I} H^q_\etale(X, \underline{\mathbf{F}_\ell})e_i
$$
car la cohomologie sur $X$ commute aux sommes directes d'après le
théorème \ref{etale-cohomology-theorem-colimit} (ou le lemme \ref{etale-cohomology-lemma-colimit}, ou le
lemme \ref{etale-cohomology-lemma-direct-sum-bounded-below-cohomology}).
Comme nous savons déjà que $H^q_\etale(X, \underline{\mathbf{F}_\ell})$
est un $\mathbf{F}_\ell$-espace vectoriel de dimension finie
(théorème \ref{etale-cohomology-theorem-vanishing-affine-curves}),
nous voyons que $H^q_\etale(X, \underline{\Lambda})$
est libre sur $\Lambda$ et de même rang.
En effet, si $\xi_1, \ldots, \xi_m$ est une base de
$H^q_\etale(X, \underline{\mathbf{F}_\ell})$, alors
$\xi_1 e_0, \ldots, \xi_m e_0$ forment une $\Lambda$-base de
$H^q_\etale(X, \underline{\Lambda})$.
\end{proof}

\begin{lemma}
\label{etale-cohomology-lemma-finite-pushforward-coefficients}
Soit $\Lambda$ un anneau noethérien, soit $k$ un corps algébriquement clos,
soit $f : X \to Y$ un morphisme fini de schémas séparés de type fini
sur $k$ et de dimension $\leq 1$, et soit $\mathcal{F}$ un faisceau
de $\Lambda$-modules sur $X_\etale$. Si
$H^q_\etale(X, \mathcal{F})$ est un $\Lambda$-module de type fini, alors
$H^q_\etale(Y, f_*\mathcal{F})$ l'est aussi.
\end{lemma}

\begin{proof}
En effet, nous avons $H^q_\etale(X, \mathcal{F}) = H^q_\etale(Y, f_*\mathcal{F})$
par l'annulation de $R^qf_*$ pour $q > 0$
(proposition \ref{etale-cohomology-proposition-finite-higher-direct-image-zero}) et par
la suite spectrale de Leray
(Cohomologie sur les sites, lemme \ref{sites-cohomology-lemma-apply-Leray}).
\end{proof}

\begin{lemma}
\label{etale-cohomology-lemma-restrict-to-open-coefficients}
Soit $\Lambda$ un anneau noethérien, soit $k$ un corps algébriquement clos,
soit $X$ un schéma séparé de type fini sur $k$ et de dimension $\leq 1$,
soit $\mathcal{F}$ un faisceau constructible de $\Lambda$-modules sur $X_\etale$,
et soit $j : X' \to X$ l'inclusion d'un sous-schéma ouvert dense. Alors
$H^q_\etale(X, \mathcal{F})$ est un $\Lambda$-module de type fini si et seulement si
$H^q_\etale(X, j_!j^{-1}\mathcal{F})$ est un $\Lambda$-module de type fini.
\end{lemma}

\begin{proof}
Puisque $X'$ est dense, nous voyons que $Z = X \setminus X'$ est de dimension $0$
et est donc un ensemble fini $Z = \{x_1, \ldots, x_n\}$ de points $k$-rationnels.
Considérons la suite exacte courte
$$
0 \to j_!j^{-1}\mathcal{F} \to \mathcal{F} \to i_*i^{-1}\mathcal{F} \to 0
$$
du lemme \ref{etale-cohomology-lemma-ses-associated-to-open}. Observons que
$H^q_\etale(X, i_*i^{-1}\mathcal{F}) = H^q_\etale(Z, i^{-1}\mathcal{F})$.
En effet, $i : Z \to X$ est une immersion fermée, donc finie,
et nous avons donc l'annulation de $R^qi_*$ pour $q > 0$ d'après la
proposition \ref{etale-cohomology-proposition-finite-higher-direct-image-zero};
l'égalité résulte alors de la suite spectrale de Leray
(Cohomologie sur les sites, lemme \ref{sites-cohomology-lemma-apply-Leray}).
Puisque $Z$ est une réunion disjointe de spectres de corps algébriquement
clos, nous en déduisons que $H^q_\etale(Z, i^{-1}\mathcal{F}) = 0$
pour $q > 0$ et que
$$
H^0_\etale(Z, i^{-1}\mathcal{F}) =
\bigoplus\nolimits_{i=1}^n \mathcal{F}_{x_i}
$$
qui est un $\Lambda$-module de type fini, car $\mathcal{F}_{x_i}$ est de type fini en vertu de
l'hypothèse que $\mathcal{F}$ est un faisceau constructible de $\Lambda$-modules.
La suite exacte longue de cohomologie donne une suite exacte
$$
0 \to
H^0_\etale(X, j_!j^{-1}\mathcal{F}) \to
H^0_\etale(X, \mathcal{F}) \to
H^0_\etale(Z, i^{-1}\mathcal{F}) \to
H^1_\etale(X, j_!j^{-1}\mathcal{F}) \to
H^1_\etale(X, \mathcal{F}) \to 0
$$
et des isomorphismes $H^q_\etale(X, j_!j^{-1}\mathcal{F}) \to
H^q_\etale(X, \mathcal{F})$ pour $q > 1$. Le lemme en résulte aisément.
\end{proof}

\begin{lemma}
\label{etale-cohomology-lemma-somewhat-easier-coefficients}
Soit $\Lambda$ un anneau noethérien, soit $M$ un $\Lambda$-module de type fini annulé
par un entier $n > 0$, soit $k$ un corps algébriquement clos, soit $X$ un
schéma séparé de type fini et de dimension $\leq 1$ sur $k$, et soit
$j : U \to X$ une immersion ouverte. Alors
\begin{enumerate}
\item $H^q_\etale(X, j_!\underline{M})$ est un $\Lambda$-module de type fini
si $n$ est premier à $\text{char}(k)$,
\item $H^q_\etale(X, j_!\underline{M})$ est un $\Lambda$-module de type fini
si $X$ est propre.
\end{enumerate}
\end{lemma}

\begin{proof}
Comme $\dim(X) \leq 1$, il existe un ouvert $V \subset X$ disjoint de
$U$ tel que $X' = U \cup V$ soit ouvert dense dans $X$ (détails omis).
Si $j' : X' \to X$ désigne le morphisme d'inclusion, nous voyons que
$j_!\underline{M}$ est un facteur direct de $j'_!\underline{M}$.
Il suffit donc de démontrer le lemme lorsque $U$ est ouvert et dense dans $X$.
Ce cas résulte des lemmes \ref{etale-cohomology-lemma-restrict-to-open-coefficients}
et \ref{etale-cohomology-lemma-constant-statements-coefficients}.
\end{proof}

\begin{lemma}
\label{etale-cohomology-lemma-ses-statements-coefficients}
Soit $\Lambda$ un anneau noethérien, soit $k$ un corps algébriquement clos,
soit $X$ un schéma séparé de type fini sur $k$ et de dimension $\leq 1$, et soit
$0 \to \mathcal{F}_1 \to \mathcal{F} \to \mathcal{F}_2 \to 0$
une suite exacte courte de faisceaux de $\Lambda$-modules sur $X_\etale$. Si
les $H^q_\etale(X, \mathcal{F}_i)$, $i = 1, 2$, sont des $\Lambda$-modules de type fini,
alors $H^q_\etale(X, \mathcal{F})$ est un $\Lambda$-module de type fini.
\end{lemma}

\begin{proof}
Cela résulte immédiatement de la suite exacte longue de cohomologie.
\end{proof}

\begin{lemma}
\label{etale-cohomology-lemma-vanishing-easier-coefficients}
Soit $\Lambda$ un anneau noethérien, soit $k$ un corps algébriquement clos,
soit $X$ un schéma séparé de type fini et de dimension $\leq 1$ sur $k$,
soit $j : U \to X$ une immersion ouverte avec $U$ connexe, soit
$\ell$ un nombre premier, soit $n > 0$, et soit $\mathcal{G}$ un faisceau
localement constant de type fini de $\Lambda$-modules sur $U_\etale$, annulé par
$\ell^n$. Alors
\begin{enumerate}
\item $H^q_\etale(X, j_!\mathcal{G})$ est un $\Lambda$-module de type fini
si $\ell$ est premier à $\text{char}(k)$,
\item $H^q_\etale(X, j_!\mathcal{G})$ est un $\Lambda$-module de type fini
si $X$ est propre.
\end{enumerate}
\end{lemma}

\begin{proof}
Soit $f : V \to U$ un morphisme fini étale de degré premier à $\ell$
comme dans le lemme \ref{etale-cohomology-lemma-pullback-filtered-modules}. La discussion de la
section \ref{etale-cohomology-section-trace-method} donne des applications
$$
\mathcal{G} \to f_*f^{-1}\mathcal{G} \to \mathcal{G}
$$
dont le composé est un isomorphisme. Il suffit donc de démontrer la
finitude de $H^q_\etale(X, j_!f_*f^{-1}\mathcal{G})$.
En vertu du théorème principal de Zariski
(Compléments sur les morphismes, lemme
\ref{more-morphisms-lemma-quasi-finite-separated-pass-through-finite})
nous pouvons choisir un diagramme
$$
\xymatrix{
V \ar[r]_{j'} \ar[d]_f & Y \ar[d]^{\overline{f}} \\
U \ar[r]^j & X
}
$$
où $\overline{f} : Y \to X$ est fini et $j'$ une immersion ouverte d'image
dense. Comme $f$ est fini et $V$ dense dans $Y$, nous avons
$V = U \times_X Y$. D'après le lemme \ref{etale-cohomology-lemma-compatible-shriek-push-finite}, nous avons
$$
j_!f_*f^{-1}\mathcal{G} = \overline{f}_*j'_!f^{-1}\mathcal{G}
$$
D'après le lemme \ref{etale-cohomology-lemma-finite-pushforward-coefficients}, il suffit de
considérer $j'_!f^{-1}\mathcal{G}$. L'existence de la filtration donnée par le
lemme \ref{etale-cohomology-lemma-pullback-filtered-modules},
le fait que $j'_!$ est exact et le
lemme \ref{etale-cohomology-lemma-ses-statements-coefficients}
nous ramènent au cas
$\mathcal{F} = j'_!\underline{M}$ pour un $\Lambda$-module de type fini $M$,
qui est le lemme \ref{etale-cohomology-lemma-somewhat-easier-coefficients}.
\end{proof}

\begin{theorem}
\label{etale-cohomology-theorem-vanishing-affine-curves-coefficients}
Soit $\Lambda$ un anneau noethérien, soit $k$ un corps algébriquement clos,
soit $X$ un schéma séparé de type fini et de dimension $\leq 1$ sur $k$,
et soit $\mathcal{F}$ un faisceau constructible de $\Lambda$-modules
de torsion sur $X_\etale$. Alors
\begin{enumerate}
\item
\label{etale-cohomology-item-finite-prime-to-p-coefficients}
$H^q_\etale(X, \mathcal{F})$ est un $\Lambda$-module de type fini
si $\mathcal{F}$ est de torsion première à $\text{char}(k)$,
\item
\label{etale-cohomology-item-finite-proper-coefficients}
$H^q_\etale(X, \mathcal{F})$ est un $\Lambda$-module de type fini si $X$ est propre.
\end{enumerate}
\end{theorem}

\begin{proof}
Nous appliquerons le dévissage par suites exactes sans autre mention.
Écrivons $\mathcal{F} = \mathcal{F}_1 \oplus \ldots \oplus \mathcal{F}_r$
où $\mathcal{F}_i$ est annulé par $\ell_i^{n_i}$
pour un nombre premier $\ell_i$ et un entier $n_i > 0$.
D'après le lemme \ref{etale-cohomology-lemma-ses-statements-coefficients}, il suffit
de démontrer le théorème pour $\mathcal{F}_i$. Nous pouvons donc supposer,
et supposons, que $\ell^n$ annule $\mathcal{F}$ pour un nombre premier $\ell$
et un entier $n > 0$.

\medskip\noindent
Puisque $\mathcal{F}$ est constructible comme faisceau de $\Lambda$-modules,
il existe un ouvert dense $U \subset X$ tel que $\mathcal{F}|_U$ soit un
faisceau localement constant de type fini de $\Lambda$-modules.
Comme $\dim(X) \leq 1$, nous pouvons supposer, après avoir rétréci $U$, que
$U = U_1 \amalg \ldots \amalg U_n$ est une réunion disjointe de schémas
irréductibles (il suffit de supprimer les points fermés qui appartiennent aux
intersections d'au moins $2$ composantes de $U$). D'après le
lemme \ref{etale-cohomology-lemma-restrict-to-open-coefficients},
nous nous ramenons au cas
$\mathcal{F} = j_!\mathcal{G}$ où $\mathcal{G}$ est un faisceau localement
constant de type fini de $\Lambda$-modules sur $U$ (et annulé
par $\ell^n$).

\medskip\noindent
Comme nous avons choisi $U = U_1 \amalg \ldots \amalg U_n$ avec les $U_i$
irréductibles, nous avons
$$
j_!\mathcal{G} =
j_{1!}(\mathcal{G}|_{U_1}) \oplus \ldots \oplus
j_{n!}(\mathcal{G}|_{U_n})
$$
où $j_i : U_i \to X$ est le morphisme d'inclusion.
Le cas de $j_{i!}(\mathcal{G}|_{U_i})$ est traité dans le
lemme \ref{etale-cohomology-lemma-vanishing-easier-coefficients}.
\end{proof}








\section{Première cohomologie des schémas propres}
\label{etale-cohomology-section-finite-etale-over-proper}

\noindent
Dans Fundamental Groups, section \ref{pione-section-finite-etale-over-proper},
nous avons vu, en un certain sens, que la formation de
$R^1f_*\underline{G}$ commute au changement de base si $f : X \to Y$
est un morphisme propre et si $G$ est un groupe fini (non nécessairement
commutatif). Dans cette section,
nous déduisons une conséquence utile de ces résultats.

\begin{lemma}
\label{etale-cohomology-lemma-proper-over-henselian-and-h1}
Soit $A$ un anneau local hensélien. Soit $X$ un schéma propre sur $A$
de fibre fermée $X_0$. Soit $M$ un groupe abélien fini.
Alors $H^1_\etale(X, \underline{M}) = H^1_\etale(X_0, \underline{M})$.
\end{lemma}

\begin{proof}
D'après Cohomologie sur les sites, lemme \ref{sites-cohomology-lemma-torsors-h1},
un élément de $H^1_\etale(X, \underline{M})$ correspond à un
$\underline{M}$-torseur $\mathcal{F}$ sur $X_\etale$.
Un tel torseur est manifestement un faisceau localement constant fini.
Ainsi, $\mathcal{F}$ est représentable par un schéma $V$ fini
étale sur $X$, lemme \ref{etale-cohomology-lemma-characterize-finite-locally-constant}.
Réciproquement, un schéma $V$ fini étale sur $X$, muni d'une action de $M$
qui en fait un $M$-torseur sur $X$, donne une classe de cohomologie.
La même correspondance entre les classes de cohomologie sur $X_0$ et les
torseurs finis étales sur $X_0$ vaut.
Le lemme résulte donc de
l'équivalence de catégories de Fundamental Groups, lemme
\ref{pione-lemma-finite-etale-on-proper-over-henselian}.
\end{proof}

\noindent
Le lemme technique suivant est un ingrédient essentiel de la démonstration du
théorème de changement de base propre. L'argument vaut mot pour mot
pour tout schéma propre sur $A$ dont la fibre spéciale est de dimension
$\leq 1$, mais la conclusion sera en fait une conséquence du
théorème de changement de base propre, et nous n'avons besoin que de cette
version particulière pour le démontrer.

\begin{lemma}
\label{etale-cohomology-lemma-efface-cohomology-on-fibre-by-finite-cover}
Soit $A$ un anneau local hensélien. Soit $X = \mathbf{P}^1_A$.
Soit $X_0 \subset X$ la fibre fermée. Soit $\ell$ un nombre
premier. Soit $\mathcal{I}$ un faisceau injectif de
$\mathbf{Z}/\ell\mathbf{Z}$-modules sur $X_\etale$. Alors
$H^q_\etale(X_0, \mathcal{I}|_{X_0}) = 0$ pour $q > 0$.
\end{lemma}

\begin{proof}
Observons que $X$ est un schéma séparé que l'on peut recouvrir par $2$
ouverts affines. Pour $q > 1$, le résultat découle donc de la variante affine
du théorème de changement de base propre due à Gabber, voir le
lemme \ref{etale-cohomology-lemma-vanishing-restriction-injective}.
Nous pouvons donc supposer $q = 1$. Soit
$\xi \in H^1_\etale(X_0, \mathcal{I}|_{X_0})$.
Il s'agit de montrer que $\xi$ est nul.
D'après les lemmes \ref{etale-cohomology-lemma-torsion-colimit-constructible} et
\ref{etale-cohomology-lemma-colimit}, nous pouvons trouver un morphisme $\mathcal{F} \to \mathcal{I}$
où $\mathcal{F}$ est un faisceau constructible de
$\mathbf{Z}/\ell\mathbf{Z}$-modules
et où $\xi$ provient d'un élément $\zeta$ de
$H^1_\etale(X_0, \mathcal{F}|_{X_0})$. Supposons donné un morphisme injectif
$\mathcal{F} \to \mathcal{F}'$ de faisceaux de
$\mathbf{Z}/\ell\mathbf{Z}$-modules sur $X_\etale$.
Comme $\mathcal{I}$ est injectif, nous pouvons prolonger
le morphisme donné $\mathcal{F} \to \mathcal{I}$ en un morphisme $\mathcal{F}' \to \mathcal{I}$.
Dans cette situation, nous pouvons remplacer $\mathcal{F}$ par $\mathcal{F}'$
et $\zeta$ par son image dans $H^1_\etale(X_0, \mathcal{F}'|_{X_0})$.
De même, si $\mathcal{F} = \mathcal{F}_1 \oplus \mathcal{F}_2$ est une somme directe,
nous pouvons remplacer $\mathcal{F}$ par $\mathcal{F}_i$
et $\zeta$ par son image dans $H^1_\etale(X_0, \mathcal{F}_i|_{X_0})$.

\medskip\noindent
D'après le lemme \ref{etale-cohomology-lemma-constructible-maps-into-constant-general}
et les remarques ci-dessus, nous pouvons supposer que $\mathcal{F}$
est de la forme $f_*\underline{M}$, où $M$ est un
$\mathbf{Z}/\ell\mathbf{Z}$-module de type fini
et $f : Y \to X$ est un morphisme fini de présentation finie
(de tels faisceaux restent constructibles d'après le
lemme \ref{etale-cohomology-lemma-finite-pushforward-constructible},
mais nous n'en aurons pas besoin).
Puisque la formation de $f_*$ commute à tout changement de base
(lemme \ref{etale-cohomology-lemma-finite-pushforward-commutes-with-base-change}),
nous voyons que la restriction de $f_*\underline{M}$ à $X_0$ est
égale à l'image directe de $\underline{M}$ par le morphisme induit
$Y_0 \to X_0$ entre les fibres spéciales. La suite spectrale de Leray
(proposition \ref{etale-cohomology-proposition-leray})
et l'annulation des images directes supérieures
(proposition \ref{etale-cohomology-proposition-finite-higher-direct-image-zero})
donnent
$$
H^1_\etale(X_0, f_*\underline{M}|_{X_0}) = H^1_\etale(Y_0, \underline{M}).
$$
Puisque $Y \to \Spec(A)$ est propre, le
lemme \ref{etale-cohomology-lemma-proper-over-henselian-and-h1} montre que
$H^1_\etale(Y_0, \underline{M})$ est égal à
$H^1_\etale(Y, \underline{M})$. Ainsi, notre classe de cohomologie
$\zeta$ se relève en une classe de cohomologie
$$
\tilde\zeta \in H^1_\etale(Y, \underline{M}) = H^1_\etale(X, f_*\underline{M})
$$
Cependant, $\tilde \zeta$ a une image nulle dans
$H^1_\etale(X, \mathcal{I})$ puisque $\mathcal{I}$ est injectif;
par commutativité de
$$
\xymatrix{
H^1_\etale(X, f_*\underline{M}) \ar[r] \ar[d] &
H^1_\etale(X, \mathcal{I}) \ar[d] \\
H^1_\etale(X_0, (f_*\underline{M})|_{X_0}) \ar[r] &
H^1_\etale(X_0, \mathcal{I}|_{X_0})
}
$$
nous en concluons que l'image $\xi$ de $\zeta$ est nulle elle aussi.
\end{proof}











\section{Préliminaires sur le changement de base}
\label{etale-cohomology-section-base-change-preliminaries}

\noindent
Si seul le théorème de changement de base lisse ou le théorème de changement
de base propre vous intéresse, vous pouvez passer directement aux
sections correspondantes.
Dans cette section et les quelques suivantes, nous considérons des diagrammes commutatifs
$$
\xymatrix{
X \ar[d]_f & Y \ar[l]^h \ar[d]^e \\
S & T \ar[l]_g
}
$$
de schémas; nous supposons généralement ce diagramme cartésien,
c'est-à-dire $Y = X \times_S T$. Un diagramme commutatif comme ci-dessus
donne lieu à un diagramme commutatif
$$
\xymatrix{
X_\etale \ar[d]_{f_{small}} & Y_\etale \ar[d]^{e_{small}} \ar[l]^{h_{small}} \\
S_\etale & T_\etale \ar[l]_{g_{small}}
}
$$
de petits sites étales. Employons les notations
$$
f^{-1} = f_{small}^{-1}, \quad
g_* = g_{small, *}, \quad
e^{-1} = e_{small}^{-1}, \text{ et}\quad
h_* = h_{small, *}.
$$
D'après Sites, section \ref{sites-section-pullback},
nous obtenons un morphisme de changement de base ou d'image inverse
$$
f^{-1}g_*\mathcal{F}
\longrightarrow
h_*e^{-1}\mathcal{F}
$$
pour un faisceau $\mathcal{F}$ sur $T_\etale$. Si $\mathcal{F}$ est un faisceau
abélien sur $T_\etale$, nous obtenons alors un morphisme de changement de base dérivé
$$
f^{-1}Rg_*\mathcal{F}
\longrightarrow
Rh_*e^{-1}\mathcal{F}
$$
voir Cohomologie sur les sites, lemme
\ref{sites-cohomology-lemma-base-change-map-flat-case}.
Enfin, si $K$ est un objet arbitraire de $D(T_\etale)$,
il existe un morphisme de changement de base
$$
f^{-1}Rg_*K
\longrightarrow
Rh_*e^{-1}K
$$
voir
Cohomologie sur les sites, remarque \ref{sites-cohomology-remark-base-change}.

\begin{lemma}
\label{etale-cohomology-lemma-base-change-local}
Considérons un diagramme cartésien de schémas
$$
\xymatrix{
X \ar[d]_f & Y \ar[l]^h \ar[d]^e \\
S & T \ar[l]_g
}
$$
Soit $\{U_i \to X\}$ un recouvrement étale tel que $U_i \to S$
se factorise en $U_i \to V_i \to S$, avec $V_i \to S$ étale,
et considérons les diagrammes cartésiens
$$
\xymatrix{
U_i \ar[d]_{f_i} & U_i \times_X Y \ar[l]^{h_i} \ar[d]^{e_i} \\
V_i & V_i \times_S T \ar[l]_{g_i}
}
$$
Soit $\mathcal{F}$ un faisceau sur $T_\etale$. Soit $K \in D(T_\etale)$.
Posons $K_i = K|_{V_i \times_S T}$ et
$\mathcal{F}_i = \mathcal{F}|_{V_i \times_S T}$.
\begin{enumerate}
\item Si $f_i^{-1}g_{i, *}\mathcal{F}_i = h_{i, *}e_i^{-1}\mathcal{F}_i$
pour tout $i$, alors $f^{-1}g_*\mathcal{F} = h_*e^{-1}\mathcal{F}$.
\item Si $f_i^{-1}Rg_{i, *}K_i = Rh_{i, *}e_i^{-1}K_i$
pour tout $i$, alors $f^{-1}Rg_*K = Rh_*e^{-1}K$.
\item Si $\mathcal{F}$ est un faisceau abélien et si
$f_i^{-1}R^qg_{i, *}\mathcal{F}_i = R^qh_{i, *}e_i^{-1}\mathcal{F}_i$
pour tout $i$, alors
$f^{-1}R^qg_*\mathcal{F} = R^qh_*e^{-1}\mathcal{F}$.
\end{enumerate}
\end{lemma}

\begin{proof}
Démonstration de (1). Observons d'abord que
$$
(f^{-1}g_*\mathcal{F})|_{U_i} =
f_i^{-1}(g_*\mathcal{F}|_{V_i}) =
f_i^{-1}g_{i, *}\mathcal{F}_i
$$
La première égalité vient de ce que $U_i \to X \to S$ est égal à
$U_i \to V_i \to S$, et la seconde de ce que
$g_*\mathcal{F}|_{V_i} = g_{i, *}\mathcal{F}_i$ par le
lemme \ref{sites-lemma-localize-morphism-strong} de Sites.
De même, nous avons
$$
(h_*e^{-1}\mathcal{F})|_{U_i} =
h_{i, *}(e^{-1}\mathcal{F}|_{U_i \times_X Y}) =
h_{i, *}e_i^{-1}\mathcal{F}_i
$$
Ainsi, si les morphismes de changement de base
$f_i^{-1}g_{i, *}\mathcal{F}_i \to h_{i, *}e_i^{-1}\mathcal{F}_i$
sont des isomorphismes pour tout $i$, alors le morphisme de changement de base
$f^{-1}g_*\mathcal{F} \to h_*e^{-1}\mathcal{F}$
se restreint en un isomorphisme sur $U_i$ pour tout $i$,
et nous concluons qu'il est un isomorphisme puisque $\{U_i \to X\}$
est un recouvrement étale.

\medskip\noindent
Pour les deux autres assertions, nous remplaçons l'appel au
lemme \ref{sites-lemma-localize-morphism-strong} de Sites
par un appel au
lemme de Cohomologie sur les sites
\ref{sites-cohomology-lemma-restrict-direct-image-open}.
\end{proof}

\begin{lemma}
\label{etale-cohomology-lemma-base-change-compose}
Considérons une tour de diagrammes cartésiens de schémas
$$
\xymatrix{
W \ar[d]_i & Z \ar[l]^j \ar[d]^k \\
X \ar[d]_f & Y \ar[l]^h \ar[d]^e \\
S & T \ar[l]_g
}
$$
Soit $K \in D(T_\etale)$. Si
$$
f^{-1}Rg_*K \to Rh_*e^{-1}K
\quad\text{et}\quad
i^{-1}Rh_*e^{-1}K \to Rj_*k^{-1}e^{-1}K
$$
sont des isomorphismes, alors
$(f \circ i)^{-1}Rg_*K \to Rj_*(e \circ k)^{-1}K$
est un isomorphisme.
De même, si $\mathcal{F}$ est un faisceau abélien sur $T_\etale$ et si
$$
f^{-1}R^qg_*\mathcal{F} \to R^qh_*e^{-1}\mathcal{F}
\quad\text{et}\quad
i^{-1}R^qh_*e^{-1}\mathcal{F} \to R^qj_*k^{-1}e^{-1}\mathcal{F}
$$
sont des isomorphismes, alors
$(f \circ i)^{-1}R^qg_*\mathcal{F} \to R^qj_*(e \circ k)^{-1}\mathcal{F}$
est un isomorphisme.
\end{lemma}

\begin{proof}
Cela est formel, pourvu que l'on vérifie que la composée de ces
morphismes de changement de base est le morphisme de changement de base
du rectangle extérieur; voir la remarque de Cohomologie sur les sites
\ref{sites-cohomology-remark-compose-base-change-horizontal}.
\end{proof}

\begin{lemma}
\label{etale-cohomology-lemma-base-change-Rf-star-colim}
Soit $I$ un ensemble ordonné filtrant. Considérons un système inverse de
diagrammes cartésiens de schémas
$$
\xymatrix{
X_i \ar[d]_{f_i} & Y_i \ar[l]^{h_i} \ar[d]^{e_i} \\
S_i & T_i \ar[l]_{g_i}
}
$$
dont les morphismes de transition sont affines et où $g_i$ est quasi-compact et
quasi-séparé. Posons $X = \lim X_i$,
$S = \lim S_i$, $T = \lim T_i$ et $Y = \lim Y_i$ afin
d'obtenir le diagramme cartésien
$$
\xymatrix{
X \ar[d]_f & Y \ar[l]^h \ar[d]^e \\
S & T \ar[l]_g
}
$$
Soit $(\mathcal{F}_i, \varphi_{i'i})$ un système de faisceaux sur
$(T_i)$ comme dans la définition \ref{etale-cohomology-definition-inverse-system-sheaves}. Posons
$\mathcal{F} = \colim p_i^{-1}\mathcal{F}_i$ sur $T$,
où $p_i : T \to T_i$ est la projection.
Alors les assertions suivantes sont vraies.
\begin{enumerate}
\item Si $f_i^{-1}g_{i, *}\mathcal{F}_i = h_{i, *}e_i^{-1}\mathcal{F}_i$
pour tout $i$, alors
$f^{-1}g_*\mathcal{F} = h_*e^{-1}\mathcal{F}$.
\item Si $\mathcal{F}_i$ est un faisceau abélien pour tout $i$ et si
$f_i^{-1}R^qg_{i, *}\mathcal{F}_i = R^qh_{i, *}e_i^{-1}\mathcal{F}_i$
pour tout $i$, alors
$f^{-1}R^qg_*\mathcal{F} = R^qh_*e^{-1}\mathcal{F}$.
\end{enumerate}
\end{lemma}

\begin{proof}
Nous démontrons (2) et omettons la démonstration de (1). Nous utiliserons sans
autre mention le fait que l'image inverse des faisceaux commute aux limites inductives,
car c'est un adjoint à gauche. Remarquons que $h_i$ est quasi-compact et quasi-séparé en tant que
changement de base de $g_i$.
En notant $q_i : Y \to Y_i$ les projections, remarquons que
$e^{-1}\mathcal{F} = \colim e^{-1}p_i^{-1}\mathcal{F}_i =
\colim q_i^{-1}e_i^{-1}\mathcal{F}_i$.
Par le lemme \ref{etale-cohomology-lemma-relative-colimit-general},
cela donne
$$
R^qh_*e^{-1}\mathcal{F} = \colim r_i^{-1}R^qh_{i, *}e_i^{-1}\mathcal{F}_i
$$
où $r_i : X \to X_i$ est la projection.
De même, nous avons
$$
f^{-1}R^qg_*\mathcal{F} =
f^{-1}\colim s_i^{-1}R^qg_{i, *}\mathcal{F}_i =
\colim r_i^{-1}f_i^{-1}R^qg_{i, *}\mathcal{F}_i
$$
où $s_i : S \to S_i$ est la projection. Le lemme en résulte.
\end{proof}

\begin{lemma}
\label{etale-cohomology-lemma-base-change-Rf-star-colim-complexes}
Soient $I$, $X_i$, $Y_i$, $S_i$, $T_i$, $f_i$, $h_i$, $e_i$, $g_i$,
$X$, $Y$, $S$, $T$, $f$, $h$, $e$, $g$ comme dans l'énoncé
du lemme \ref{etale-cohomology-lemma-base-change-Rf-star-colim}.
Soient $0 \in I$ et $K_0 \in D^+(T_{0, \etale})$.
Pour $i \in I$, $i \geq 0$, notons $K_i$ l'image inverse de
$K_0$ sur $T_i$. Notons $K$ l'image inverse de $K_0$ sur $T$.
Si $f_i^{-1}Rg_{i, *}K_i = Rh_{i, *}e_i^{-1}K_i$
pour tout $i \geq 0$, alors $f^{-1}Rg_*K = Rh_*e^{-1}K$.
\end{lemma}

\begin{proof}
Il suffit de montrer que le morphisme de changement de base
$f^{-1}Rg_*K \to Rh_*e^{-1}K$ induit un isomorphisme
sur les faisceaux de cohomologie. Autrement dit, nous devons montrer
que $f^{-1}R^pg_*K \to R^ph_*e^{-1}K$ est un isomorphisme
pour tout $p \in \mathbf{Z}$, sachant que
$f_i^{-1}R^pg_{i, *}K_i \to R^ph_{i, *}e_i^{-1}K_i$
est un isomorphisme pour tous $i \geq 0$ et $p \in \mathbf{Z}$.
À ce point, nous pouvons raisonner exactement comme dans la démonstration du
lemme \ref{etale-cohomology-lemma-base-change-Rf-star-colim},
en remplaçant le renvoi au
lemme \ref{etale-cohomology-lemma-relative-colimit-general}
par un renvoi au
lemme \ref{etale-cohomology-lemma-relative-colimit-general-complexes}.
\end{proof}

\begin{lemma}
\label{etale-cohomology-lemma-base-change-f-star-general-stalks}
Considérons un diagramme cartésien de schémas
$$
\xymatrix{
X \ar[d]_f & Y \ar[l]^h \ar[d]^e \\
S & T \ar[l]_g
}
$$
où $g : T \to S$ est quasi-compact et quasi-séparé.
Soit $\mathcal{F}$ un
faisceau abélien sur $T_\etale$. Soit $q \geq 0$. Les assertions suivantes sont équivalentes.
\begin{enumerate}
\item Pour tout point géométrique $\overline{x}$ de $X$, d'image
$\overline{s} = f(\overline{x})$, nous avons
$$
H^q(\Spec(\mathcal{O}^{sh}_{X, \overline{x}}) \times_S T, \mathcal{F})
=
H^q(\Spec(\mathcal{O}^{sh}_{S, \overline{s}}) \times_S T, \mathcal{F})
$$
\item $f^{-1}R^qg_*\mathcal{F} \to R^qh_*e^{-1}\mathcal{F}$
est un isomorphisme.
\end{enumerate}
\end{lemma}

\begin{proof}
Puisque $Y = X \times_S T$, nous avons
$\Spec(\mathcal{O}^{sh}_{X, \overline{x}}) \times_X Y =
\Spec(\mathcal{O}^{sh}_{X, \overline{x}}) \times_S T$. Ainsi,
le morphisme de (1) est le morphisme sur les fibres en $\overline{x}$ associé au morphisme
de (2), par le théorème \ref{etale-cohomology-theorem-higher-direct-images} (et le
lemme \ref{etale-cohomology-lemma-stalk-pullback}).
Le résultat découle donc du théorème \ref{etale-cohomology-theorem-exactness-stalks}.
\end{proof}

\begin{lemma}
\label{etale-cohomology-lemma-check-stalks-better}
Soit $f : X \to S$ un morphisme de schémas.
Soit $\overline{x}$ un point géométrique de $X$, d'image $\overline{s}$ dans $S$.
Soit $\Spec(K) \to \Spec(\mathcal{O}^{sh}_{S, \overline{s}})$
un morphisme, où $K$ est un corps séparablement clos. Soit $\mathcal{F}$ un
faisceau abélien sur $\Spec(K)_\etale$. Soit $q \geq 0$. Les assertions suivantes sont
équivalentes.
\begin{enumerate}
\item
$H^q(\Spec(\mathcal{O}^{sh}_{X, \overline{x}}) \times_S \Spec(K), \mathcal{F}) =
H^q(\Spec(\mathcal{O}^{sh}_{S, \overline{s}}) \times_S \Spec(K), \mathcal{F})$
\item
$H^q(\Spec(\mathcal{O}^{sh}_{X, \overline{x}})
\times_{\Spec(\mathcal{O}^{sh}_{S, \overline{s}})} \Spec(K), \mathcal{F}) =
H^q(\Spec(K), \mathcal{F})$
\end{enumerate}
\end{lemma}

\begin{proof}
Remarquons que $\Spec(K) \times_S \Spec(\mathcal{O}^{sh}_{S, \overline{s}})$
est le spectre d'une limite inductive filtrante d'algèbres étales sur $K$.
Puisque $K$ est séparablement clos, toute $K$-algèbre étale
est un produit fini de copies de $K$. Nous pouvons donc écrire
$$
\Spec(K) \times_S \Spec(\mathcal{O}^{sh}_{S, \overline{s}}) =
\lim_{i \in I} \coprod\nolimits_{a \in A_i} \Spec(K)
$$
comme une limite cofiltrante dont chaque terme est une réunion disjointe
de copies de $\Spec(K)$ indexées par un ensemble fini $A_i$.
Remarquons que $A_i$ est non vide, puisque nous disposons de
$\Spec(K) \to \Spec(\mathcal{O}^{sh}_{S, \overline{s}})$.
Il s'ensuit que
\begin{align*}
\Spec(\mathcal{O}^{sh}_{X, \overline{x}}) \times_S \Spec(K)
& =
\Spec(\mathcal{O}^{sh}_{X, \overline{x}})
\times_{\Spec(\mathcal{O}^{sh}_{S, \overline{s}})}
\left(
\Spec(\mathcal{O}^{sh}_{S, \overline{s}})
\times_S \Spec(K)\right) \\
& =
\lim_{i \in I} \coprod\nolimits_{a \in A_i}
\Spec(\mathcal{O}^{sh}_{X, \overline{x}})
\times_{\Spec(\mathcal{O}^{sh}_{S, \overline{s}})} \Spec(K)
\end{align*}
Comme, dans notre situation, la cohomologie commute aux limites
de schémas (théorème \ref{etale-cohomology-theorem-colimit}), nous concluons.
\end{proof}













\section{Changement de base pour l'image directe}
\label{etale-cohomology-section-base-change-f-star}

\noindent
Cette section est préliminaire et devrait être omise en première lecture.
Nous y étudions pour quels morphismes $f : X \to S$ on a
$f^{-1}g_* = h_*e^{-1}$ sur tous les faisceaux (d'ensembles), pour tout
diagramme cartésien
$$
\xymatrix{
X \ar[d]_f & Y \ar[l]^h \ar[d]^e \\
S & T \ar[l]_g
}
$$
où $g$ est quasi-compact et quasi-séparé.

\begin{lemma}
\label{etale-cohomology-lemma-base-change-f-star-general}
Considérons le diagramme cartésien de schémas
$$
\xymatrix{
X \ar[d]_f & Y \ar[l]^h \ar[d]^e \\
S & T \ar[l]_g
}
$$
Supposons que $f$ soit plat et que tout objet $U$ de $X_\etale$ possède
un recouvrement $\{U_i \to U\}$ tel que $U_i \to S$
se factorise en $U_i \to V_i \to S$, avec $V_i \to S$
étale et $U_i \to V_i$ quasi-compact à
fibres géométriquement connexes.
Alors, pour tout faisceau d'ensembles $\mathcal{F}$ sur $T_\etale$, nous avons
$f^{-1}g_*\mathcal{F} = h_*e^{-1}\mathcal{F}$.
\end{lemma}

\begin{proof}
Soit $U \to X$ un morphisme étale tel que $U \to S$ se factorise en
$U \to V \to S$, avec $V \to S$ étale et $U \to V$ quasi-compact
à fibres géométriquement connexes. Remarquons que $U \to V$ est plat
(lemme \ref{flat-lemma-etale-flat-up-down} de Compléments sur la platitude).
Nous affirmons que
\begin{align*}
f^{-1}g_*\mathcal{F}(U)
& = g_*\mathcal{F}(V) \\
& = \mathcal{F}(V \times_S T) \\
& = e^{-1}\mathcal{F}(U \times_X Y) \\
& = h_*e^{-1}\mathcal{F}(U)
\end{align*}
En effet, en considérant $U$ comme un objet de $X_\etale$ et
$V$ comme un objet de $S_\etale$, nous voyons que la première égalité
découle du lemme \ref{etale-cohomology-lemma-sections-upstairs}\footnote{Plus
précisément, nous utilisons aussi le fait que la restriction de $f^{-1}g_*\mathcal{F}$
à $U_\etale$ est l'image inverse par $U \to V$ de la restriction de
$g_*\mathcal{F}$ à $V_\etale$. Voir le
lemme \ref{sites-lemma-localize-morphism-strong} de Sites.}.
En considérant $V \times_S T$ comme un objet de $T_\etale$,
la seconde égalité découle de la définition de $g_*$.
Remarquons que $U \times_X Y = U \times_S T$ (car $Y = X \times_S T$)
et que, par conséquent, $U \times_X Y \to V \times_S T$
est à fibres géométriquement connexes comme changement de base de $U \to V$.
En considérant $U \times_X Y$ comme un objet de $Y_\etale$, nous voyons que
la troisième égalité découle du lemme \ref{etale-cohomology-lemma-sections-upstairs},
comme précédemment. Enfin, la quatrième égalité découle de la définition
de $h_*$.

\medskip\noindent
Puisque, par hypothèse, tout objet de $X_\etale$ possède un recouvrement
étale auquel s'applique l'argument du paragraphe précédent,
nous voyons que le lemme est vrai.
\end{proof}

\begin{lemma}
\label{etale-cohomology-lemma-fppf-reduced-fibres-base-change-f-star}
Considérons un diagramme cartésien de schémas
$$
\xymatrix{
X \ar[d]_f & Y \ar[l]^h \ar[d]^e \\
S & T \ar[l]_g
}
$$
où $f$ est plat et localement de présentation finie,
à fibres géométriquement réduites.
Alors $f^{-1}g_*\mathcal{F} = h_*e^{-1}\mathcal{F}$
pour tout faisceau $\mathcal{F}$ sur $T_\etale$.
\end{lemma}

\begin{proof}
Combiner le lemme \ref{etale-cohomology-lemma-base-change-f-star-general} avec le
lemme de Compléments sur les morphismes
\ref{more-morphisms-lemma-cover-by-geometrically-connected}.
\end{proof}

\begin{lemma}
\label{etale-cohomology-lemma-base-change-f-star-field}
Considérons le diagramme cartésien de schémas
$$
\xymatrix{
X \ar[d]_f & Y \ar[l]^h \ar[d]^e \\
S & T \ar[l]_g
}
$$
Supposons que $S$ soit le spectre d'un corps séparablement clos.
Alors $f^{-1}g_*\mathcal{F} = h_*e^{-1}\mathcal{F}$
pour tout faisceau $\mathcal{F}$ sur $T_\etale$.
\end{lemma}

\begin{proof}
Nous pouvons travailler localement sur $X$. Nous pouvons donc supposer $X$ affine.
Nous pouvons alors écrire $X$ comme une limite cofiltrante de schémas affines
de type fini sur $S$. Par le lemme \ref{etale-cohomology-lemma-base-change-Rf-star-colim},
nous pouvons supposer que $X$ est de type fini sur $S$.
Le lemme \ref{etale-cohomology-lemma-base-change-f-star-general}
s'applique alors, car tout schéma de type
fini sur un corps séparablement clos est une réunion disjointe finie
de schémas connexes et géométriquement connexes
(voir le lemme de Variétés
\ref{varieties-lemma-separably-closed-field-connected-components}).
\end{proof}

\begin{lemma}
\label{etale-cohomology-lemma-base-change-f-star-valuation}
Considérons un diagramme cartésien de schémas
$$
\xymatrix{
X \ar[d]_f & Y \ar[l]^h \ar[d]^e \\
S & T \ar[l]_g
}
$$
Supposons que
\begin{enumerate}
\item $f$ soit plat et ouvert,
\item les corps résiduels de $S$ soient séparablement algébriquement clos,
\item étant donné un morphisme étale $U \to X$ avec $U$ affine,
nous puissions écrire $U$ comme une réunion disjointe finie de sous-schémas ouverts
de $X$ (par exemple si $X$ est un schéma intègre normal
dont le corps des fonctions est séparablement clos),
\item tout ouvert non vide d'une fibre $X_s$ de $f$ soit connexe
(par exemple si $X_s$ est irréductible ou vide).
\end{enumerate}
Alors, pour tout faisceau d'ensembles $\mathcal{F}$ sur $T_\etale$, nous avons
$f^{-1}g_*\mathcal{F} = h_*e^{-1}\mathcal{F}$.
\end{lemma}

\begin{proof}
Omis. Indication : les hypothèses impliquent presque trivialement
la condition du lemme \ref{etale-cohomology-lemma-base-change-f-star-general}.
L'exemple de la partie (3) découle du
lemme \ref{etale-cohomology-lemma-normal-scheme-with-alg-closed-function-field}.
\end{proof}

\noindent
Le lemme suivant n'a pas vraiment sa place ici, mais il ne semble
pas y avoir ailleurs de meilleur endroit où le placer.

\begin{lemma}
\label{etale-cohomology-lemma-fppf-reduced-fibres-pullback-products}
Soit $f : X \to S$ un morphisme de schémas plat et
localement de présentation finie, à fibres géométriquement réduites.
Alors $f^{-1} : \Sh(S_\etale) \to \Sh(X_\etale)$ commute
aux produits.
\end{lemma}

\begin{proof}
Soit $I$ un ensemble et soit $\mathcal{G}_i$ un faisceau sur $S_\etale$
pour $i \in I$.
Soit $U \to X$ un morphisme étale tel que $U \to S$ se factorise en
$U \to V \to S$, avec $V \to S$ étale et $U \to V$ plat de
présentation finie à fibres géométriquement connexes. Nous avons alors
\begin{align*}
f^{-1}(\prod \mathcal{G}_i)(U)
& =
(\prod \mathcal{G}_i)(V) \\
& =
\prod \mathcal{G}_i(V) \\
& =
\prod f^{-1}\mathcal{G}_i(U) \\
& =
(\prod f^{-1}\mathcal{G}_i)(U)
\end{align*}
où nous avons utilisé le lemme \ref{etale-cohomology-lemma-sections-upstairs}
dans les première et troisième égalités
(nous utilisons aussi le fait que la restriction de $f^{-1}\mathcal{G}$
à $U_\etale$ est l'image inverse par $U \to V$ de la restriction de
$\mathcal{G}$ à $V_\etale$; voir le
lemme \ref{sites-lemma-localize-morphism-strong} de Sites).
Par le lemme de Compléments sur les morphismes
\ref{more-morphisms-lemma-cover-by-geometrically-connected}
tout objet $U$ de $X_\etale$ possède un recouvrement étale
$\{U_i \to U\}$ tel que la discussion du paragraphe
précédent s'applique à $U_i$. Le lemme en résulte.
\end{proof}

\begin{lemma}
\label{etale-cohomology-lemma-base-change-f-star}
Soit $f : X \to S$ un morphisme plat de schémas tel
que, pour tout point géométrique $\overline{x}$ de $X$, le morphisme
$$
\mathcal{O}_{S, f(\overline{x})}^{sh}
\longrightarrow
\mathcal{O}_{X, \overline{x}}^{sh}
$$
soit à fibres géométriquement connexes. Alors, pour tout
diagramme cartésien de schémas
$$
\xymatrix{
X \ar[d]_f & Y \ar[l]^h \ar[d]^e \\
S & T \ar[l]_g
}
$$
où $g$ est quasi-compact et quasi-séparé, nous avons
$f^{-1}g_*\mathcal{F} = h_*e^{-1}\mathcal{F}$
pour tout faisceau d'ensembles $\mathcal{F}$ sur $T_\etale$.
\end{lemma}

\begin{proof}
Il suffit de vérifier l'égalité sur les fibres; voir le
théorème \ref{etale-cohomology-theorem-exactness-stalks}.
Par le théorème \ref{etale-cohomology-theorem-higher-direct-images}, nous avons
$$
(h_*e^{-1}\mathcal{F})_{\overline{x}} =
\Gamma(\Spec(\mathcal{O}_{X, \overline{x}}^{sh}) \times_X Y, e^{-1}\mathcal{F})
$$
et, de même, nous avons
$$
(f^{-1}g_*\mathcal{F})_{\overline{x}} =
(g_*\mathcal{F})_{f(\overline{x})} =
\Gamma(\Spec(\mathcal{O}_{S, f(\overline{x})}^{sh}) \times_S T, \mathcal{F})
$$
Ces ensembles sont égaux par application du lemme \ref{etale-cohomology-lemma-sections-upstairs}
au morphisme
$$
\Spec(\mathcal{O}_{X, \overline{x}}^{sh}) \times_X Y
\longrightarrow
\Spec(\mathcal{O}_{S, f(\overline{x})}^{sh}) \times_S T
$$
qui est un changement de base de
$\Spec(\mathcal{O}_{X, \overline{x}}^{sh}) \to
\Spec(\mathcal{O}_{S, f(\overline{x})}^{sh})$
car $Y = X \times_S T$.
\end{proof}







\section{Changement de base pour les images directes supérieures}
\label{etale-cohomology-section-base-change}

\noindent
Cette section est l'analogue de la section \ref{etale-cohomology-section-base-change-f-star}
pour les images directes supérieures.
Elle est préliminaire et devrait être omise en première lecture.

\begin{remark}
\label{etale-cohomology-remark-base-change-holds}
Soit $f : X \to S$ un morphisme de schémas. Soit $n$ un entier.
Nous dirons que $BC(f, n, q_0)$ est vraie si, pour tout diagramme commutatif
$$
\xymatrix{
X \ar[d]_f & X' \ar[l] \ar[d]_{f'} & Y \ar[l]^h \ar[d]^e \\
S & S' \ar[l] & T \ar[l]_g
}
$$
où $X' = X \times_S S'$ et $Y = X' \times_{S'} T$, où
$g$ est quasi-compact et quasi-séparé, et pour tout faisceau abélien
$\mathcal{F}$ sur $T_\etale$ annulé par $n$, le morphisme de changement de base
$$
(f')^{-1}R^qg_*\mathcal{F}
\longrightarrow
R^qh_*e^{-1}\mathcal{F}
$$
est un isomorphisme pour $q \leq q_0$.
\end{remark}

\begin{lemma}
\label{etale-cohomology-lemma-base-change-q-injective}
Avec $f : X \to S$ et $n$ comme dans la remarque \ref{etale-cohomology-remark-base-change-holds},
supposons que, pour un certain $q \geq 1$, on ait $BC(f, n, q - 1)$. Alors,
pour tout diagramme commutatif
$$
\xymatrix{
X \ar[d]_f & X' \ar[l] \ar[d]_{f'} & Y \ar[l]^h \ar[d]^e \\
S & S' \ar[l] & T \ar[l]_g
}
$$
où $X' = X \times_S S'$ et $Y = X' \times_{S'} T$, où
$g$ est quasi-compact et quasi-séparé, et pour tout faisceau abélien
$\mathcal{F}$ sur $T_\etale$ annulé par $n$, on a :
\begin{enumerate}
\item le morphisme de changement de base
$(f')^{-1}R^qg_*\mathcal{F}\to R^qh_*e^{-1}\mathcal{F}$
est injectif;
\item si $\mathcal{F} \subset \mathcal{G}$, où $\mathcal{G}$
sur $T_\etale$ est annulé par $n$, alors
$$
\Coker\left(
(f')^{-1}R^qg_*\mathcal{F}\to R^qh_*e^{-1}\mathcal{F}
\right)
\subset
\Coker\left(
(f')^{-1}R^qg_*\mathcal{G}\to R^qh_*e^{-1}\mathcal{G}
\right)
$$
\item si, dans (2), le faisceau $\mathcal{G}$ est un faisceau injectif
de $\mathbf{Z}/n\mathbf{Z}$-modules, alors
$$
\Coker\left((f')^{-1}R^qg_*\mathcal{F}\to R^qh_*e^{-1}\mathcal{F} \right)
\subset R^qh_*e^{-1}\mathcal{G}
$$
\end{enumerate}
\end{lemma}

\begin{proof}
Choisissons une suite exacte courte
$0 \to \mathcal{F} \to \mathcal{I} \to \mathcal{Q} \to 0$
où $\mathcal{I}$ est un faisceau injectif de $\mathbf{Z}/n\mathbf{Z}$-modules.
Considérons le diagramme induit
$$
\xymatrix{
(f')^{-1}R^{q - 1}g_*\mathcal{I} \ar[d]_{\cong} \ar[r] &
(f')^{-1}R^{q - 1}g_*\mathcal{Q} \ar[d]_{\cong} \ar[r] &
(f')^{-1}R^qg_*\mathcal{F} \ar[d] \ar[r] &
0 \ar[d] \\
R^{q - 1}h_*e^{-1}\mathcal{I} \ar[r] &
R^{q - 1}h_*e^{-1}\mathcal{Q} \ar[r] &
R^qh_*e^{-1}\mathcal{F} \ar[r] &
R^qh_*e^{-1}\mathcal{I}
}
$$
à lignes exactes. Le zéro figure dans le coin supérieur droit
car $\mathcal{I}$ est injectif. Les deux flèches verticales de gauche sont
des isomorphismes par $BC(f, n, q - 1)$. Nous concluons que (1) est vraie.
Ce qui précède montre aussi que
$$
\Coker\left(
(f')^{-1}R^qg_*\mathcal{F}\to R^qh_*e^{-1}\mathcal{F}
\right)
\subset
R^qh_*e^{-1}\mathcal{I}
$$
donc (3) est vraie. Pour démontrer (2), choisissons
$\mathcal{F} \subset \mathcal{G} \subset \mathcal{I}$.
\end{proof}

\begin{lemma}
\label{etale-cohomology-lemma-base-change-q-integral-top}
Avec $f : X \to S$ et $n$ comme dans la remarque \ref{etale-cohomology-remark-base-change-holds},
supposons que, pour un certain $q \geq 1$, on ait $BC(f, n, q - 1)$. Considérons
les diagrammes commutatifs
$$
\vcenter{
\xymatrix{
X \ar[d]_f &
X' \ar[d]_{f'} \ar[l] &
Y \ar[l]^h \ar[d]^e &
Y' \ar[l]^{\pi'} \ar[d]^{e'} \\
S &
S' \ar[l] &
T \ar[l]_g &
T' \ar[l]_\pi
}
}
\quad\text{et}\quad
\vcenter{
\xymatrix{
X' \ar[d]_{f'} & & Y' \ar[ll]^{h' = h \circ \pi'} \ar[d]^{e'} \\
S' & & T' \ar[ll]_{g' = g \circ \pi}
}
}
$$
où tous les carrés sont cartésiens, $g$ est quasi-compact et quasi-séparé, et
$\pi$ est entier surjectif. Soit $\mathcal{F}$ un faisceau abélien
sur $T_\etale$ annulé par $n$, et posons $\mathcal{F}' = \pi^{-1}\mathcal{F}$.
Si le morphisme de changement de base
$$
(f')^{-1}R^qg'_*\mathcal{F}' \longrightarrow R^qh'_*(e')^{-1}\mathcal{F}'
$$
est un isomorphisme, alors le morphisme de changement de base
$(f')^{-1}R^qg_*\mathcal{F} \to R^qh_*e^{-1}\mathcal{F}$
est un isomorphisme.
\end{lemma}

\begin{proof}
Remarquons que $\mathcal{F} \to \pi_*\mathcal{F}'$ est injectif
car $\pi$ est surjectif (cela se vérifie sur les fibres). Ainsi, par le
lemme \ref{etale-cohomology-lemma-base-change-q-injective},
nous voyons qu'il suffit de montrer que le morphisme de changement de base
$$
(f')^{-1}R^qg_*\pi_*\mathcal{F}'
\longrightarrow
R^qh_*e^{-1}\pi_*\mathcal{F}'
$$
est un isomorphisme. Cela découle de l'hypothèse, car
nous avons $R^qg_*\pi_*\mathcal{F}' = R^qg'_*\mathcal{F}'$,
nous avons $e^{-1}\pi_*\mathcal{F}' =\pi'_*(e')^{-1}\mathcal{F}'$, et
nous avons $R^qh_*\pi'_*(e')^{-1}\mathcal{F}' = R^qh'_*(e')^{-1}\mathcal{F}'$.
Ces égalités découlent des lemmes
\ref{etale-cohomology-lemma-integral-pushforward-commutes-with-base-change} et
\ref{etale-cohomology-lemma-what-integral} et de la suite spectrale de Leray relative
(lemme \ref{sites-cohomology-lemma-relative-Leray} de Cohomologie sur les sites).
\end{proof}

\begin{lemma}
\label{etale-cohomology-lemma-base-change-q-integral-bottom}
Avec $f : X \to S$ et $n$ comme dans la remarque \ref{etale-cohomology-remark-base-change-holds},
supposons que, pour un certain $q \geq 1$, on ait $BC(f, n, q - 1)$. Considérons
les diagrammes commutatifs
$$
\vcenter{
\xymatrix{
X \ar[d]_f &
X' \ar[d]_{f'} \ar[l] &
X'' \ar[l]^{\pi'} \ar[d]_{f''} &
Y \ar[l]^{h'} \ar[d]^e \\
S &
S' \ar[l] &
S'' \ar[l]_\pi &
T \ar[l]_{g'}
}
}
\quad\text{et}\quad
\vcenter{
\xymatrix{
X' \ar[d]_{f'} & & Y \ar[ll]^{h = h' \circ \pi'} \ar[d]^e \\
S' & & T \ar[ll]_{g = g' \circ \pi}
}
}
$$
où tous les carrés sont cartésiens, $g'$ est quasi-compact et quasi-séparé, et
$\pi$ est entier. Soit $\mathcal{F}$ un faisceau abélien
sur $T_\etale$ annulé par $n$. Si le morphisme de changement de base
$$
(f')^{-1}R^qg_*\mathcal{F} \longrightarrow R^qh_*e^{-1}\mathcal{F}
$$
est un isomorphisme, alors le morphisme de changement de base
$(f'')^{-1}R^qg'_*\mathcal{F} \to R^qh'_*e^{-1}\mathcal{F}$
est un isomorphisme.
\end{lemma}

\begin{proof}
Puisque $\pi$ et $\pi'$ sont entiers, nous avons $R\pi_* = \pi_*$ et
$R\pi'_* = \pi'_*$; voir le lemme \ref{etale-cohomology-lemma-what-integral}.
Nous avons aussi $(f')^{-1}\pi_* = \pi'_*(f'')^{-1}$. Nous voyons donc que
$\pi'_*(f'')^{-1}R^qg'_*\mathcal{F} = (f')^{-1}R^qg_*\mathcal{F}$
et
$\pi'_*R^qh'_*e^{-1}\mathcal{F} = R^qh_*e^{-1}\mathcal{F}$.
Ainsi, l'hypothèse signifie que notre morphisme devient un
isomorphisme après application du foncteur $\pi'_*$.
Nous voyons donc que c'est un isomorphisme par le lemme \ref{etale-cohomology-lemma-what-integral}.
\end{proof}

\begin{lemma}
\label{etale-cohomology-lemma-formal-argument}
Soit $T$ un schéma quasi-compact et quasi-séparé.
Soit $P$ une propriété des schémas quasi-compacts et quasi-séparés
sur $T$. Supposons que :
\begin{enumerate}
\item si $T'' \to T'$ est un épaississement de schémas quasi-compacts et
quasi-séparés sur $T$, alors $P(T'')$ si et seulement si $P(T')$;
\item si $T' = \lim T_i$ est la limite d'un système inverse de
schémas quasi-compacts et quasi-séparés sur $T$, à morphismes de
transition affines, et si $P(T_i)$ est vraie pour tout $i$, alors
$P(T')$ est vraie;
\item si $Z \subset T'$ est un sous-schéma fermé de complément
quasi-compact $V \subset T'$ et si $P(T')$ est vraie,
alors $P(V)$ ou $P(Z)$ est vraie.
\end{enumerate}
Alors $P(T)$ implique $P(\Spec(K))$ pour un certain morphisme $\Spec(K) \to T$,
où $K$ est un corps.
\end{lemma}

\begin{proof}
Considérons l'ensemble $\mathfrak T$ des sous-schémas fermés $T' \subset T$
tels que $P(T')$ soit vraie. Par l'hypothèse (2), cet ensemble possède un élément minimal,
noté $T'$. Par l'hypothèse (1), nous voyons que $T'$ est réduit.
Soit $\eta \in T'$ le point générique d'une composante
irréductible de $T'$. Alors $\eta = \Spec(K)$ pour un certain corps
$K$, et $\eta = \lim V$, où la limite porte sur les sous-schémas
ouverts affines $V \subset T'$ contenant $\eta$.
Par l'hypothèse (3) et la minimalité de $T'$, nous voyons
que $P(V)$ est vraie pour tous ces $V$. Ainsi $P(\eta)$ est vraie
par (2), ce qui achève la démonstration.
\end{proof}

\begin{lemma}
\label{etale-cohomology-lemma-base-change-does-not-hold-pre}
Avec $f : X \to S$ et $n$ comme dans la remarque \ref{etale-cohomology-remark-base-change-holds},
supposons que, pour un certain $q \geq 1$, $BC(f, n, q - 1)$ soit vraie, mais que
$BC(f, n, q)$ ne le soit pas. Il existe alors un diagramme commutatif
$$
\xymatrix{
X \ar[d]_f & X' \ar[d]_{f'} \ar[l] & Y \ar[l]^h \ar[d]^e \\
S & S' \ar[l] & \Spec(K) \ar[l]_g
}
$$
où $X' = X \times_S S'$, $Y = X' \times_{S'} \Spec(K)$,
$K$ est un corps et $\mathcal{F}$ est un faisceau abélien
sur $\Spec(K)$ annulé par $n$, tel que
$(f')^{-1}R^qg_*\mathcal{F} \to R^qh_*e^{-1}\mathcal{F}$
n'est pas un isomorphisme.
\end{lemma}

\begin{proof}
Choisissons un diagramme commutatif
$$
\xymatrix{
X \ar[d]_f & X' \ar[l] \ar[d]_{f'} & Y \ar[l]^h \ar[d]^e \\
S & S' \ar[l] & T \ar[l]_g
}
$$
où $X' = X \times_S S'$ et $Y = X' \times_{S'} T$, où
$g$ est quasi-compact et quasi-séparé, ainsi qu'un faisceau abélien
$\mathcal{F}$ sur $T_\etale$ annulé par $n$, tels que
le morphisme de changement de base
$(f')^{-1}R^qg_*\mathcal{F} \to R^qh_*e^{-1}\mathcal{F}$
n'est pas un isomorphisme. Nous pouvons bien entendu remplacer $S'$
par un ouvert affine de $S'$, et nous le faisons; cela implique que $T$ est quasi-compact
et quasi-séparé. Par le lemme \ref{etale-cohomology-lemma-base-change-q-injective}, nous voyons que
$(f')^{-1}R^qg_*\mathcal{F} \to R^qh_*e^{-1}\mathcal{F}$
est injectif. Choisissons un point géométrique $\overline{x}$ de $X'$
et un élément $\xi$ de $(R^qh_*e^{-1}\mathcal{F})_{\overline{x}}$
qui n'appartient pas à l'image du morphisme
$((f')^{-1}R^qg_*\mathcal{F})_{\overline{x}} \to
(R^qh_*e^{-1}\mathcal{F})_{\overline{x}}$.

\medskip\noindent
Considérons un morphisme $\pi : T' \to T$, avec $T'$ quasi-compact
et quasi-séparé, et notons $\mathcal{F}' = \pi^{-1}\mathcal{F}$.
Notons $\pi' : Y' = Y \times_T T' \to Y$ le changement de base de $\pi$
et $e' : Y' \to T'$ le changement de base de $e$. Nous avons les diagrammes
$$
\vcenter{
\xymatrix{
X' \ar[d]_{f'} & Y \ar[l]^h \ar[d]^e & Y' \ar[l]^{\pi'} \ar[d]^{e'} \\
S' & T \ar[l]_g & T' \ar[l]_\pi
}
}
\quad\text{et}\quad
\vcenter{
\xymatrix{
X' \ar[d]_{f'} & & Y' \ar[ll]^{h' = h \circ \pi'} \ar[d]^{e'} \\
S' & & T' \ar[ll]_{g' = g \circ \pi}
}
}
$$
À l'aide des morphismes d'image inverse, nous obtenons un diagramme commutatif canonique
$$
\xymatrix{
(f')^{-1}R^qg_*\mathcal{F} \ar[r] \ar[d] &
(f')^{-1}R^qg'_*\mathcal{F}' \ar[d] \\
R^qh_*e^{-1}\mathcal{F} \ar[r] &
R^qh'_*(e')^{-1}\mathcal{F}'
}
$$
de faisceaux abéliens sur $X'$. Soit $P(T')$ la propriété suivante :
\begin{itemize}
\item L'image $\xi'$ de $\xi$ dans
$(R^qh'_*(e')^{-1}\mathcal{F}')_{\overline{x}}$ n'appartient
pas à l'image du morphisme
$((f')^{-1}R^qg'_*\mathcal{F}')_{\overline{x}} \to
(R^qh'_*(e')^{-1}\mathcal{F}')_{\overline{x}}$.
\end{itemize}
Nous affirmons que les hypothèses (1), (2) et (3) du
lemme \ref{etale-cohomology-lemma-formal-argument} sont satisfaites par $P$,
ce qui démontre notre lemme.

\medskip\noindent
La condition (1) du lemme \ref{etale-cohomology-lemma-formal-argument}
est satisfaite par $P$, car la topologie étale d'un schéma
et celle d'un épaississement de ce schéma sont les mêmes. Voir la
proposition \ref{etale-cohomology-proposition-topological-invariance}.

\medskip\noindent
Supposons que $I$ soit un ensemble ordonné filtrant et que $(T_i)$
soit un système inverse indexé par $I$ de schémas quasi-compacts et quasi-séparés
sur $T$, à morphismes de transition affines.
Posons $T' = \lim T_i$. Notons $\mathcal{F}'$ et $\mathcal{F}_i$
les images inverses de $\mathcal{F}$ sur $T'$ et sur $T_i$, respectivement. Considérons
les diagrammes
$$
\vcenter{
\xymatrix{
X' \ar[d]_{f'} & Y \ar[l]^h \ar[d]^e & Y_i \ar[l]^{\pi_i'} \ar[d]^{e_i} \\
S' & T \ar[l]_g & T_i \ar[l]_{\pi_i}
}
}
\quad\text{et}\quad
\vcenter{
\xymatrix{
X' \ar[d]_{f'} & & Y_i \ar[ll]^{h_i = h \circ \pi_i'} \ar[d]^{e_i} \\
S' & & T_i \ar[ll]_{g_i = g \circ \pi_i}
}
}
$$
comme au paragraphe précédent. Il est clair que $\mathcal{F}'$ sur
$T'$ est la limite inductive des images inverses des $\mathcal{F}_i$ sur $T'$
et que $(e')^{-1}\mathcal{F}'$ est la limite inductive des images inverses
des $e_i^{-1}\mathcal{F}_i$ sur $Y'$.
Par le lemme \ref{etale-cohomology-lemma-relative-colimit-general},
nous avons
$$
R^qh'_*(e')^{-1}\mathcal{F}' = \colim R^qh_{i, *}e_i^{-1}\mathcal{F}_i
\quad\text{et}\quad
(f')^{-1}R^qg'_*\mathcal{F}' = \colim (f')^{-1}R^qg_{i, *}\mathcal{F}_i
$$
Il s'ensuit que, si $P(T_i)$ est vraie pour tout $i$, alors
$P(T')$ est vraie. Ainsi, la condition (2) du lemme \ref{etale-cohomology-lemma-formal-argument}
est satisfaite par $P$.

\medskip\noindent
La plus intéressante est la condition (3) du lemme \ref{etale-cohomology-lemma-formal-argument}.
Supposons que $T'$ soit un schéma quasi-compact et quasi-séparé sur $T$
tel que $P(T')$ soit vraie.
Soit $Z \subset T'$ un sous-schéma fermé de complément quasi-compact
$V \subset T'$. Considérons le diagramme
$$
\xymatrix{
Y' \times_{T'} Z \ar[d]_{e_Z} \ar[r]_{i'} &
Y' \ar[d]_{e'} &
Y' \times_{T'} V \ar[l]^{j'} \ar[d]^{e_V} \\
Z \ar[r]^i &
T' &
V \ar[l]_j
}
$$
Choisissons un morphisme injectif $j^{-1}\mathcal{F}' \to \mathcal{J}$,
où $\mathcal{J}$ est un faisceau injectif de $\mathbf{Z}/n\mathbf{Z}$-modules
sur $V$. En regardant les fibres, nous voyons que le morphisme
$$
\mathcal{F}' \to \mathcal{G} = j_*\mathcal{J} \oplus i_*i^{-1}\mathcal{F}'
$$
est injectif. Ainsi, $\xi'$ a pour image un élément non nul de
\begin{align*}
& \Coker\left(
((f')^{-1}R^qg'_*\mathcal{G})_{\overline{x}}
\to
(R^qh'_*(e')^{-1}\mathcal{G})_{\overline{x}}
\right) = \\
&
\Coker\left(
((f')^{-1}R^qg'_*j_*\mathcal{J})_{\overline{x}}
\to
(R^qh'_*(e')^{-1}j_*\mathcal{J})_{\overline{x}}
\right) \oplus \\
& \Coker\left(
((f')^{-1}R^qg'_*i_*i^{-1}\mathcal{F}')_{\overline{x}}
\to
(R^qh'_*(e')^{-1}i_*i^{-1}\mathcal{F}')_{\overline{x}}
\right)
\end{align*}
par la partie (2) du lemme \ref{etale-cohomology-lemma-base-change-q-injective}.
Si $\xi'$ ne s'envoie pas sur zéro dans le second facteur,
nous utilisons
$$
(f')^{-1}R^qg'_*i_*i^{-1}\mathcal{F}' =
(f')^{-1}R^q(g' \circ i)_*i^{-1}\mathcal{F}'
$$
(car $Ri_* = i_*$ par la
proposition \ref{etale-cohomology-proposition-finite-higher-direct-image-zero}) et
$$
R^qh'_*(e')^{-1}i_*i^{-1}\mathcal{F}' =
R^qh'_*i'_*e_Z^{-1}i^{-1}\mathcal{F}' =
R^q(h' \circ i')_*e_Z^{-1}i^{-1}\mathcal{F}'
$$
(la première égalité par le
lemme \ref{etale-cohomology-lemma-finite-pushforward-commutes-with-base-change}
et la seconde parce que
$Ri'_* = i'_*$ par la
proposition \ref{etale-cohomology-proposition-finite-higher-direct-image-zero}),
pour voir que $P(Z)$ est vraie.
Enfin, supposons que $\xi'$ ne s'envoie pas sur zéro dans le premier facteur.
Nous avons
$$
(e')^{-1}j_*\mathcal{J} = j'_*e_V^{-1}\mathcal{J}
\quad\text{et}\quad
R^aj'_*e_V^{-1}\mathcal{J} = 0, \quad a = 1, \ldots, q - 1
$$
par $BC(f, n, q - 1)$ appliquée au diagramme
$$
\xymatrix{
X' \ar[d]_{f'} & Y' \ar[l]^{h'} \ar[d]_{e'} & Y' \times_{T'} V \ar[l]^{j'} \ar[d]^{e_V} \\
S' & T' \ar[l]_{g'} & V \ar[l]_j
}
$$
et au fait que $\mathcal{J}$ est injectif.
Par la suite spectrale de Leray relative pour $h' \circ j'$
(lemme \ref{sites-cohomology-lemma-relative-Leray} de Cohomologie sur les sites),
nous en déduisons que
$$
R^qh'_*(e')^{-1}j_*\mathcal{J} =
R^qh'_*j'_*e_V^{-1}\mathcal{J}
\longrightarrow
R^q(h' \circ j')_* e_V^{-1}\mathcal{J}
$$
est injectif. Ainsi, $\xi'$ a pour image un élément non nul de
$(R^q(h' \circ j')_* e_V^{-1}\mathcal{J})_{\overline{x}}$.
En appliquant la partie (3) du lemme \ref{etale-cohomology-lemma-base-change-q-injective}
à l'injection $j^{-1}\mathcal{F}' \to \mathcal{J}$,
nous concluons que $P(V)$ est vraie.
\end{proof}

\begin{lemma}
\label{etale-cohomology-lemma-base-change-does-not-hold}
Avec $f : X \to S$ et $n$ comme dans la remarque \ref{etale-cohomology-remark-base-change-holds},
supposons que, pour un certain $q \geq 1$,
$BC(f, n, q - 1)$ soit vraie, mais que $BC(f, n, q)$ ne le soit pas.
Il existe alors un diagramme commutatif
$$
\xymatrix{
X \ar[d]_f & X' \ar[d] \ar[l] & Y \ar[l]^h \ar[d] \\
S & S' \ar[l] & \Spec(K) \ar[l]
}
$$
dont les deux carrés sont cartésiens, où :
\begin{enumerate}
\item $S'$ est affine, intègre et normal, à corps des fonctions
algébriquement clos;
\item $K$ est algébriquement clos et $\Spec(K) \to S'$
est dominant (autrement dit, $K$ est une extension du
corps des fonctions de $S'$).
\end{enumerate}
De plus, il existe un entier $d | n$
tel que $R^qh_*(\mathbf{Z}/d\mathbf{Z})$ soit non nul.
\end{lemma}

\noindent
Réciproquement, la non-annulation de $R^qh_*(\mathbf{Z}/d\mathbf{Z})$
dans le lemme implique que $BC(f, n, q)$ n'est pas vraie, puisque le
lemme \ref{etale-cohomology-lemma-Rf-star-zero-normal-with-alg-closed-function-field}
montre que $R^q(\Spec(K) \to S')_*\mathbf{Z}/d\mathbf{Z} = 0$.

\begin{proof}
Choisissons d'abord un diagramme et $\mathcal{F}$ comme dans le
lemme \ref{etale-cohomology-lemma-base-change-does-not-hold-pre}.
Nous pouvons supposer $S'$ affine, et nous le faisons (c'est évident, mais
voir la démonstration du lemme en cas de doute).
Par le lemme \ref{etale-cohomology-lemma-base-change-q-integral-top},
nous pouvons supposer $K$ algébriquement clos.
Alors $\mathcal{F}$ correspond à un $\mathbf{Z}/n\mathbf{Z}$-module.
Un tel module est une somme directe de copies de $\mathbf{Z}/d\mathbf{Z}$
pour divers $d | n$; nous pouvons donc supposer $\mathcal{F}$
constant, de valeur $\mathbf{Z}/d\mathbf{Z}$.
Par le lemme \ref{etale-cohomology-lemma-base-change-q-integral-bottom},
nous pouvons remplacer $S'$ par sa normalisation
dans $\Spec(K)$, ce qui achève la démonstration.
\end{proof}










\section{Changement de base lisse}
\label{etale-cohomology-section-smooth-base-change}

\noindent
Dans cette section, nous démontrons le théorème de changement de base lisse.

\begin{lemma}
\label{etale-cohomology-lemma-smooth-base-change-fields}
Soit $K/k$ une extension de corps. Soit $X$ une courbe affine lisse
sur $k$, munie d'un point rationnel $x \in X(k)$. Soit $\mathcal{F}$ un faisceau
abélien sur $\Spec(K)$ annulé par un entier $n$ inversible dans $k$.
Soit $q > 0$ et soit
$$
\xi \in H^q(X_K, (X_K \to \Spec(K))^{-1}\mathcal{F})
$$
Il existe :
\begin{enumerate}
\item des extensions finies $K'/K$ et $k'/k$, avec $k' \subset K'$;
\item un revêtement galoisien fini étale $Z \to X_{k'}$, de groupe $G$.
\end{enumerate}
L'ordre de $G$ divise une puissance de $n$, le morphisme
$Z \to X_{k'}$ est scindé au-dessus de $x_{k'}$, et
$\xi$ s'annule dans $H^q(Z_{K'}, (Z_{K'} \to \Spec(K))^{-1}\mathcal{F})$.
\end{lemma}

\begin{proof}
Pour $q > 1$, nous savons que $\xi$ s'annule dans
$H^q(X_{\overline{K}}, (X_{\overline{K}} \to \Spec(K))^{-1}\mathcal{F})$
(théorème \ref{etale-cohomology-theorem-vanishing-affine-curves}).
Par le lemme \ref{etale-cohomology-lemma-directed-colimit-cohomology}, nous voyons que
cela signifie qu'il existe une extension finie $K'/K$ telle que
$\xi$ s'annule dans $H^q(X_{K'}, (X_{K'} \to \Spec(K))^{-1}\mathcal{F})$.
Dans ce cas, nous pouvons donc prendre $k' = k$ et $Z = X$.

\medskip\noindent
Supposons $q = 1$. Rappelons que $\mathcal{F}$ correspond à un
module discret $M$ muni d'une action continue de $\text{Gal}_K$; voir le
lemme \ref{etale-cohomology-lemma-equivalence-abelian-sheaves-point}.
Puisque $M$ est de $n$-torsion, il est la réunion de ses sous-groupes finis
stables par $\text{Gal}_K$. Nous nous ramenons donc
au cas où $M$ est un groupe abélien fini annulé par $n$;
voir le lemme \ref{etale-cohomology-lemma-colimit}. Après avoir remplacé $K$ par une extension finie,
nous pouvons supposer que l'action de $\text{Gal}_K$ sur $M$ est triviale.
Nous pouvons donc supposer que $\mathcal{F} = \underline{M}$ est le faisceau
constant de valeur le groupe abélien fini $M$, annulé par $n$.

\medskip\noindent
Nous pouvons écrire $M$ comme une somme directe de groupes cycliques.
Deux revêtements galoisiens finis étales quelconques, dont les groupes de Galois
sont d'ordre inversible dans $k$, sont dominés par un troisième dont le
groupe de Galois est d'ordre inversible dans $k$
(section \ref{pione-section-finite-etale-under-galois} de Fundamental Groups).
Il suffit donc de démontrer le lemme lorsque
$M = \mathbf{Z}/d\mathbf{Z}$, où $d | n$.

\medskip\noindent
Supposons $M = \mathbf{Z}/d\mathbf{Z}$, où $d | n$.
Dans ce cas, $\overline{\xi} = \xi|_{X_{\overline{K}}}$ est un élément de
$$
H^1(X_{\overline{k}}, \mathbf{Z}/d\mathbf{Z}) =
H^1(X_{\overline{K}}, \mathbf{Z}/d\mathbf{Z})
$$
Voir le théorème \ref{etale-cohomology-theorem-vanishing-affine-curves}.
Ce groupe classifie les $\mathbf{Z}/d\mathbf{Z}$-torseurs; voir le
lemme \ref{sites-cohomology-lemma-torsors-h1} de Cohomologie sur les sites.
Le torseur correspondant à $\overline{\xi}$ (considéré comme un faisceau sur
$X_{\overline{k}, \etale}$) donne à son tour un morphisme fini étale
$T \to X_{\overline{k}}$ muni d'une action de $\mathbf{Z}/d\mathbf{Z}$,
transitive sur la fibre de $T$ au-dessus de $x_{\overline{k}}$; voir le
lemme \ref{etale-cohomology-lemma-characterize-finite-locally-constant}.
Choisissons une composante connexe $T' \subset T$ (si $\overline{\xi}$ est
d'ordre $d$, alors $T$ est déjà connexe).
Alors $T' \to X_{\overline{k}}$ est un revêtement galoisien
fini étale dont le groupe de Galois est un sous-groupe $G \subset \mathbf{Z}/d\mathbf{Z}$
(un petit détail est omis). En outre, l'élément $\overline{\xi}$ s'envoie sur zéro
par le morphisme $H^1(X_{\overline{k}}, \mathbf{Z}/d\mathbf{Z}) \to
H^1(T', \mathbf{Z}/d\mathbf{Z})$, car c'est l'une
des propriétés qui définissent $T$.

\medskip\noindent
Nous utilisons ensuite un argument de limite pour choisir une extension finie $k'/k$
contenue dans $\overline{k}$ telle que $T' \to X_{\overline{k}}$
descende en un revêtement galoisien fini étale $Z \to X_{k'}$, de groupe $G$.
Voir les lemmes de Limits \ref{limits-lemma-descend-finite-presentation},
\ref{limits-lemma-descend-finite-finite-presentation} et
\ref{limits-lemma-descend-etale}.
Après avoir agrandi $k'$, nous pouvons supposer que $Z$ est scindé au-dessus de $x_{k'}$.
Par construction, l'image de $\xi$ dans
$H^1(Z_{\overline{K}}, \mathbf{Z}/d\mathbf{Z})$ est nulle.
Ainsi, par le lemme \ref{etale-cohomology-lemma-directed-colimit-cohomology},
nous pouvons trouver une sous-extension finie $\overline{K}/K'/K$
contenant $k'$ telle que $\xi$ s'annule dans $H^1(Z_{K'}, \mathbf{Z}/d\mathbf{Z})$,
ce qui achève la démonstration.
\end{proof}

\begin{theorem}[Changement de base lisse]
\label{etale-cohomology-theorem-smooth-base-change}
Considérons un diagramme cartésien de schémas
$$
\xymatrix{
X \ar[d]_f & Y \ar[l]^h \ar[d]^e \\
S & T \ar[l]_g
}
$$
où $f$ est lisse et $g$ quasi-compact et quasi-séparé. Alors
$$
f^{-1}R^qg_*\mathcal{F} = R^qh_*e^{-1}\mathcal{F}
$$
pour tout $q$ et tout faisceau abélien $\mathcal{F}$
sur $T_\etale$ dont toutes les fibres aux points géométriques sont de torsion
d'ordres inversibles sur $S$.
\end{theorem}

\begin{proof}[Première démonstration du changement de base lisse]
Cette démonstration est très longue, mais plus directe (elle utilise moins de théorie générale)
que la seconde démonstration donnée ci-dessous.

\medskip\noindent
Le théorème est local sur $X_\etale$. Plus précisément, supposons donné
$U \to X$ étale tel que $U \to S$ se factorise en $U \to V \to S$,
avec $V \to S$ étale. Nous pouvons alors considérer le carré cartésien
$$
\xymatrix{
U \ar[d]_{f'} & U \times_X Y \ar[l]^{h'} \ar[d]^{e'} \\
V & V \times_S T \ar[l]_{g'}
}
$$
et, en posant $\mathcal{F}' = \mathcal{F}|_{V \times_S T}$,
nous avons $f^{-1}R^qg_*\mathcal{F}|_U = (f')^{-1}R^qg'_*\mathcal{F}'$
et $R^qh_*e^{-1}\mathcal{F}|_U = R^qh'_*(e')^{-1}\mathcal{F}'$
(cela découle de la compatibilité de la localisation avec les morphismes de sites; voir le
lemme \ref{sites-lemma-localize-morphism-strong} de Sites,
ainsi que le
lemme de Cohomologie sur les sites
\ref{sites-cohomology-lemma-restrict-direct-image-open}).
Il suffit donc de produire un recouvrement étale de $X$ par des
$U \to X$ et des factorisations $U \to V \to S$
comme ci-dessus, tels que le théorème vaille pour le diagramme formé de
$f'$, $h'$, $g'$ et $e'$.

\medskip\noindent
Par la structure locale des morphismes lisses, voir le
lemme \ref{morphisms-lemma-smooth-etale-over-affine-space} de Morphismes,
nous pouvons supposer $X$ et $S$ affines et que $X \to S$
se factorise par un morphisme étale $X \to \mathbf{A}^d_S$.
Si nous avons une tour de diagrammes cartésiens
$$
\xymatrix{
W \ar[d]_i & Z \ar[l]^j \ar[d]^k \\
X \ar[d]_f & Y \ar[l]^h \ar[d]^e \\
S & T \ar[l]_g
}
$$
et si le théorème vaut pour les carrés inférieur et supérieur, alors
il vaut pour le rectangle extérieur; cela est formel.
En écrivant $X \to S$ comme la composée
$$
X \to \mathbf{A}^{d - 1}_S \to \mathbf{A}^{d - 2}_S \to \ldots \to
\mathbf{A}^1_S \to S
$$
nous concluons qu'il suffit de démontrer le théorème lorsque $X$ et $S$
sont affines et que $X \to S$ est de dimension relative $1$.

\medskip\noindent
Pour tout $n \geq 1$ inversible sur $S$, soit $\mathcal{F}[n]$
le sous-faisceau des sections de $\mathcal{F}$ annulées par $n$. Alors
$\mathcal{F} = \colim \mathcal{F}[n]$ par notre hypothèse sur
les fibres de $\mathcal{F}$. Les foncteurs $e^{-1}$ et $f^{-1}$
commutent aux limites inductives car ce sont des adjoints à gauche. Les foncteurs
$R^qh_*$ et $R^qg_*$ commutent aux limites inductives filtrantes par le
lemme \ref{etale-cohomology-lemma-relative-colimit}.
Il suffit donc de démontrer le théorème pour $\mathcal{F}[n]$.
Désormais, nous fixons un entier $n$, travaillons avec des
faisceaux de $\mathbf{Z}/n\mathbf{Z}$-modules et
supposons que $S$ est un schéma sur $\Spec(\mathbf{Z}[1/n])$.

\medskip\noindent
Nous nous ramenons ensuite au cas où $T$ est affine. Puisque $g$ est quasi-compact
et quasi-séparé et que $S$ est affine, le schéma $T$ est quasi-compact et
quasi-séparé. Nous pouvons donc utiliser le principe de récurrence du
lemme \ref{coherent-lemma-induction-principle} de Cohomologie des schémas.
Par conséquent, il suffit de montrer que, si $T = W \cup W'$
est un recouvrement ouvert et si le théorème vaut pour les carrés
$$
\xymatrix{
X \ar[d] & e^{-1}(W) \ar[l]^i \ar[d] \\
S & W \ar[l]_a
}
\quad
\xymatrix{
X \ar[d] & e^{-1}(W') \ar[l]^j \ar[d] \\
S & W' \ar[l]_b
}
\quad
\xymatrix{
X \ar[d] & e^{-1}(W \cap W') \ar[l]^-k \ar[d] \\
S & W \cap W' \ar[l]_c
}
$$
alors le théorème vaut pour le diagramme initial. Pour le voir, considérons le
diagramme
$$
\xymatrix{
f^{-1}R^{q - 1}c_*\mathcal{F}|_{W \cap W'} \ar[d]^{\cong} \ar[r] &
f^{-1}R^qg_*\mathcal{F} \ar[d] \ar[r] &
f^{-1}R^qa_*\mathcal{F}|_W \oplus
f^{-1}R^qb_*\mathcal{F}|_{W'} \ar[d]_{\cong} \\
R^{q - 1}k_*e^{-1}\mathcal{F}|_{e^{-1}(W \cap W')} \ar[r] &
R^qh_*e^{-1}\mathcal{F} \ar[r] &
R^qi_*e^{-1}\mathcal{F}|_{e^{-1}(W)} \oplus
R^qj_*e^{-1}\mathcal{F}|_{e^{-1}(W')}
}
$$
dont les lignes sont les suites exactes longues du
lemme \ref{etale-cohomology-lemma-relative-mayer-vietoris}.
Le lemme des cinq donne donc la conclusion voulue.

\medskip\noindent
En résumé, nous pouvons supposer $S$, $X$, $T$ et $Y$ affines,
$\mathcal{F}$ de $n$-torsion, $X \to S$ lisse de dimension relative $1$,
et $S$ schéma sur $\mathbf{Z}[1/n]$.
Nous allons démontrer le théorème par récurrence sur $q$. Le cas initial $q = 0$
est traité par le lemme \ref{etale-cohomology-lemma-fppf-reduced-fibres-base-change-f-star}.
Supposons $q > 0$ et le théorème vrai en tout degré inférieur.
Choisissons une suite exacte courte
$0 \to \mathcal{F} \to \mathcal{I} \to \mathcal{Q} \to 0$
où $\mathcal{I}$ est un faisceau injectif de $\mathbf{Z}/n\mathbf{Z}$-modules.
Considérons le diagramme induit
$$
\xymatrix{
f^{-1}R^{q - 1}g_*\mathcal{I} \ar[d]_{\cong} \ar[r] &
f^{-1}R^{q - 1}g_*\mathcal{Q} \ar[d]_{\cong} \ar[r] &
f^{-1}R^qg_*\mathcal{F} \ar[d] \ar[r] &
0 \ar[d] \\
R^{q - 1}h_*e^{-1}\mathcal{I} \ar[r] &
R^{q - 1}h_*e^{-1}\mathcal{Q} \ar[r] &
R^qh_*e^{-1}\mathcal{F} \ar[r] &
R^qh_*e^{-1}\mathcal{I}
}
$$
à lignes exactes. Le zéro figure dans le coin supérieur droit
car $\mathcal{I}$ est injectif. Les deux flèches verticales de gauche sont
des isomorphismes par l'hypothèse de récurrence. Il suffit donc de démontrer
que $R^qh_*e^{-1}\mathcal{I} = 0$.

\medskip\noindent
Écrivons $S = \Spec(A)$ et $T = \Spec(B)$, et supposons que le morphisme
$T \to S$ soit donné par le morphisme d'anneaux $A \to B$. Nous pouvons écrire
$A \to B = \colim_{i \in I} (A_i \to B_i)$ comme une limite inductive
filtrante de morphismes d'anneaux de type fini sur $\mathbf{Z}[1/n]$
(voir le lemme \ref{algebra-lemma-limit-no-condition} d'Algèbre).
Pour $i \in I$, posons $S_i = \Spec(A_i)$ et $T_i = \Spec(B_i)$.
Pour $i$ assez grand, nous pouvons trouver un morphisme lisse $X_i \to S_i$
de dimension relative $1$ tel que $X = X_i \times_{S_i} S$; voir les
lemmes de Limits \ref{limits-lemma-descend-finite-presentation},
\ref{limits-lemma-descend-smooth} et
\ref{limits-lemma-descend-dimension-d}. Posons $Y_i = X_i \times_{S_i} T_i$
pour obtenir les carrés
$$
\xymatrix{
X_i \ar[d]_{f_i} & Y_i \ar[l]^{h_i} \ar[d]^{e_i} \\
S_i & T_i \ar[l]_{g_i}
}
$$
Remarquons que $\mathcal{I}_i = (T \to T_i)_*\mathcal{I}$
est un faisceau injectif de $\mathbf{Z}/n\mathbf{Z}$-modules sur
$T_i$; voir le lemme de Cohomologie sur les sites
\ref{sites-cohomology-lemma-pushforward-injective-flat}.
Nous avons $\mathcal{I} = \colim (T \to T_i)^{-1}\mathcal{I}_i$
par le lemme \ref{etale-cohomology-lemma-linus-hamann}. En prenant l'image inverse par $e$, nous obtenons
$e^{-1}\mathcal{I} = \colim (Y \to Y_i)^{-1}e_i^{-1}\mathcal{I}_i$.
Par le lemme \ref{etale-cohomology-lemma-relative-colimit-general}, appliqué au système
de morphismes $Y_i \to X_i$ de limite $Y \to X$, nous avons
$$
R^qh_*e^{-1}\mathcal{I} =
\colim (X \to X_i)^{-1} R^qh_{i, *} e_i^{-1}\mathcal{I}_i
$$
Cela nous ramène au cas où $T$ et $S$ sont affines de type
fini sur $\mathbf{Z}[1/n]$.

\medskip\noindent
En résumé, nous avons un entier $q \geq 1$ tel que le théorème vaille
en degrés $< q$, les schémas $S$ et $T$ sont affines de type fini
sur $\mathbf{Z}[1/n]$, le morphisme
$X \to S$ est lisse de dimension relative $1$, avec $X$ affine, et
$\mathcal{I}$ est un faisceau injectif de $\mathbf{Z}/n\mathbf{Z}$-modules;
nous devons montrer que $R^qh_*e^{-1}\mathcal{I} = 0$.
Nous allons le faire par récurrence sur $\dim(T)$.

\medskip\noindent
Le cas initial est $T = \emptyset$, c'est-à-dire $\dim(T) < 0$.
Si l'on préfère, on peut prendre comme cas initial
le cas $\dim(T) = 0$. Dans ce cas, $T \to S$ est fini
(en fait, même $T \to \Spec(\mathbf{Z}[1/n])$ est fini puisque la
cible est jacobsonienne; les détails sont omis), donc
$h$ est lui aussi fini et ses
images directes supérieures sont nulles (voir les références ci-dessous).

\medskip\noindent
Supposons $\dim(T) = d \geq 0$ et le résultat connu dans toutes les situations où $T$
est de dimension inférieure. Choisissons $U$ affine et étale sur $X$, ainsi qu'une section $\xi$
de $R^qh_*e^{-1}\mathcal{I}$ sur $U$. Nous devons montrer
que $\xi$ est nulle. Nous pouvons bien entendu remplacer $X$ par $U$
(et, en conséquence, $Y$ par $U \times_X Y$)
et supposer $\xi \in H^0(X, R^qh_*e^{-1}\mathcal{I})$.
De plus, puisque $R^qh_*e^{-1}\mathcal{I}$ est un faisceau,
il suffit de démontrer que $\xi$ est nulle localement sur $X$.
Nous pouvons donc remplacer $X$ par les membres d'un recouvrement étale.
En particulier, à l'aide du lemme \ref{etale-cohomology-lemma-higher-direct-images},
nous pouvons supposer que $\xi$ est l'image d'un élément
$\tilde \xi \in H^q(Y, e^{-1}\mathcal{I})$.
En termes de $\tilde \xi$, notre tâche est de montrer que
$\tilde \xi$ s'annule dans $H^q(U_i \times_X Y, e^{-1}\mathcal{I})$
pour un certain recouvrement étale $\{U_i \to X\}$.

\medskip\noindent
Par le lemme \ref{more-morphisms-lemma-cover-smooth-by-special} de Compléments sur les morphismes,
nous pouvons supposer que $X \to S$ se factorise en $X \to V \to S$,
où $V \to S$ est étale et $X \to V$ est un morphisme lisse
de schémas affines, de dimension relative $1$, muni d'une section et
à fibres géométriquement connexes. Remarquons que
$\dim(V \times_S T) \leq \dim(T) = d$
par exemple par le lemme de Compléments d'algèbre
\ref{more-algebra-lemma-dimension-etale-extension}.
Nous pouvons donc remplacer $S$ par $V$ et $T$ par $V \times_S T$
(exactement comme dans la discussion du premier paragraphe de la démonstration).
Ainsi, nous pouvons supposer $X \to S$ lisse de dimension relative $1$,
à fibres géométriquement connexes et muni d'une section $\sigma : S \to X$.

\medskip\noindent
Soit $\pi : T' \to T$ un morphisme fini surjectif.
Nous utiliserons ci-dessous que $\dim(T') \leq \dim(T) = d$; voir le
lemme \ref{algebra-lemma-integral-dim-up} d'Algèbre.
Choisissons un morphisme injectif $\pi^{-1}\mathcal{I} \to \mathcal{I}'$
vers un faisceau injectif de $\mathbf{Z}/n\mathbf{Z}$-modules.
Alors $\mathcal{I} \to \pi_*\mathcal{I}'$ est injectif
et possède donc une rétraction (car $\mathcal{I}$ est un faisceau injectif
de $\mathbf{Z}/n\mathbf{Z}$-modules).
Notons $\pi' : Y' = Y \times_T T' \to Y$ le changement de base de $\pi$
et $e' : Y' \to T'$ le changement de base de $e$. Nous avons le diagramme
$$
\xymatrix{
X \ar[d]_f & Y \ar[l]^h \ar[d]^e & Y' \ar[l]^{\pi'} \ar[d]^{e'} \\
S & T \ar[l]_g & T' \ar[l]_\pi
}
$$
Par la proposition \ref{etale-cohomology-proposition-finite-higher-direct-image-zero} et le
lemme \ref{etale-cohomology-lemma-finite-pushforward-commutes-with-base-change}, nous avons
$R\pi'_*(e')^{-1}\mathcal{I}' = e^{-1}\pi_*\mathcal{I}'$.
Ainsi, par la suite spectrale de Leray
(lemme \ref{sites-cohomology-lemma-Leray} de Cohomologie sur les sites),
nous avons
$$
H^q(Y', (e')^{-1}\mathcal{I}') =
H^q(Y, e^{-1}\pi_*\mathcal{I}') \supset H^q(Y, e^{-1}\mathcal{I})
$$
et cela reste vrai après changement de base par tout $U \to X$ étale.
Nous pouvons donc remplacer $T$ par $T'$, $\mathcal{I}$ par $\mathcal{I}'$
et $\tilde \xi$ par son image dans $H^q(Y', (e')^{-1}\mathcal{I}')$.

\medskip\noindent
Supposons donnée une factorisation $T \to S' \to S$, où
$\pi : S' \to S$ est fini.
En posant $X' = S' \times_S X$, nous pouvons considérer le diagramme induit
$$
\xymatrix{
X \ar[d]_f & X' \ar[l]^{\pi'} \ar[d]^{f'} & Y \ar[l]^{h'} \ar[d]^e \\
S & S' \ar[l]_\pi & T \ar[l]_g
}
$$
Puisque les images directes supérieures de $\pi'$ sont nulles, nous voyons que
$R^qh_*e^{-1}\mathcal{I} = \pi'_*R^qh'_*e^{-1}\mathcal{I}$
par la suite spectrale de Leray. Par conséquent,
$H^0(X, R^qh_*e^{-1}\mathcal{I}) = H^0(X', R^qh'_*e^{-1}\mathcal{I})$.
Ainsi, $\xi$ est nulle si et seulement si la section correspondante de
$R^qh'_*e^{-1}\mathcal{I}$ est nulle\footnote{Cette
étape peut aussi se voir autrement. En effet, nous devons montrer
qu'il existe un recouvrement étale $\{U_i \to X\}$ tel que
$\tilde \xi$ s'annule dans $H^q(U_i \times_X Y, e^{-1}\mathcal{I})$.
Cependant, si nous montrons qu'il existe un recouvrement étale
$\{U'_j \to X'\}$ tel que $\tilde \xi$ s'annule dans
$H^q(U'_j \times_{X'} Y, e^{-1}\mathcal{I})$, alors, par la propriété (B)
pour $X' \to X$ (lemme \ref{etale-cohomology-lemma-finite-B}), il existe un recouvrement étale
$\{U_i \to X\}$ tel que $U_i \times_X X'$ soit une réunion disjointe
de schémas sur $X'$, chacun se factorisant par un $U'_j$.
Nous voyons donc que $\tilde \xi$ s'annule dans $H^q(U_i \times_X Y, e^{-1}\mathcal{I})$,
comme voulu.}.
Nous pouvons donc remplacer $S$ par $S'$ et $X$ par $X'$.
Remarquons que le changement de base de $\sigma : S \to X$ est $\sigma' : S' \to X'$;
ainsi, après ce remplacement, il reste vrai que $X \to S$
possède une section $\sigma$ et des fibres géométriquement connexes.

\medskip\noindent
Nous utiliserons le fait que $S$ et $T$ sont des schémas de Nagata; voir la
proposition \ref{algebra-proposition-ubiquity-nagata} d'Algèbre.
Cela garantira
que diverses normalisations sont finies; voir les
lemmes de Morphismes \ref{morphisms-lemma-nagata-normalization-finite} et
\ref{morphisms-lemma-nagata-normalization}.
En particulier, nous pouvons d'abord remplacer $T$ par sa normalisation,
puis remplacer $S$ par la normalisation de $S$ dans $T$.
Alors $T \to S$ est une réunion disjointe de morphismes
dominants de schémas intègres normaux; voir le
lemme \ref{morphisms-lemma-normal-normalization} de Morphismes.
Nous pouvons clairement raisonner composante connexe
par composante connexe; nous pouvons donc supposer que $T \to S$ est un
morphisme dominant de schémas intègres normaux.

\medskip\noindent
Soient $s \in S$ et $t \in T$ les points génériques. Par le
lemme \ref{etale-cohomology-lemma-smooth-base-change-fields}, il existe des extensions finies
de corps $K/\kappa(t)$ et $k/\kappa(s)$ telles que $k$ soit contenu dans $K$,
ainsi qu'un revêtement galoisien fini étale $Z \to X_k$ de groupe de Galois $G$,
d'ordre divisant une puissance de $n$, scindé au-dessus de $\sigma(\Spec(k))$,
tel que $\tilde \xi$ s'envoie sur zéro dans $H^q(Z_K, e^{-1}\mathcal{I}|_{Z_K})$.
Soit $T' \to T$ la normalisation de $T$ dans $\Spec(K)$,
et soit $S' \to S$ la normalisation de $S$ dans $\Spec(k)$.
Nous obtenons alors un diagramme commutatif
$$
\xymatrix{
S' \ar[d] & T' \ar[l] \ar[d] \\
S & T \ar[l]
}
$$
dont les flèches verticales sont finies. Par les arguments donnés ci-dessus, nous
pouvons remplacer $S$ et $T$ par $S'$ et $T'$, et nous le faisons (et, en conséquence,
$X$ par $X \times_S S'$ et $Y$ par $Y \times_T T'$). Après ce remplacement,
nous concluons que nous avons un revêtement galoisien fini étale
$Z \to X_s$ de la fibre générique de $X \to S$,
de groupe de Galois $G$ d'ordre divisant une puissance de $n$,
scindé au-dessus de $\sigma(s)$, tel que $\tilde \xi$ s'envoie sur zéro dans
$H^q(Z_t, (Z_t \to Y)^{-1}e^{-1}\mathcal{I})$.
Ici, $Z_t = Z \times_S t = Z \times_s t = Z \times_{X_s} Y_t$.
Comme $n$ est inversible sur $S$,
le lemme \ref{pione-lemma-extend-covering} de Fundamental Groups
permet de trouver un morphisme fini étale $U \to X$ dont la restriction à $X_s$
est $Z$.

\medskip\noindent
À ce point, nous remplaçons $X$ par $U$ et $Y$ par $U \times_X Y$.
Après ce remplacement, il peut arriver que
les fibres de $X \to S$ ne soient plus géométriquement connexes
(il subsiste une section, mais nous ne l'utiliserons pas);
en revanche, ce remplacement assure que $\tilde \xi$ s'envoie sur zéro
dans $H^q(Y_t, e^{-1}\mathcal{I})$, c'est-à-dire que $\tilde \xi$ se restreint à
zéro sur la fibre générique de $Y \to T$.

\medskip\noindent
Rappelons que $t$ est le spectre du corps des fonctions de $T$, c'est-à-dire
que, comme schéma, $t$ est la limite des sous-schémas ouverts affines non vides de $T$.
Par le lemme \ref{etale-cohomology-lemma-directed-colimit-cohomology}, nous concluons qu'il existe
un sous-schéma ouvert non vide $V \subset T$ tel que $\tilde \xi$ s'envoie sur zéro dans
$H^q(Y \times_T V, e^{-1}\mathcal{I}|_{Y \times_T V})$.

\medskip\noindent
Notons $Z = T \setminus V$. Considérons le diagramme
$$
\xymatrix{
Y \times_T Z \ar[d]_{e_Z} \ar[r]_{i'} &
Y \ar[d]_e &
Y \times_T V \ar[l]^{j'} \ar[d]^{e_V} \\
Z \ar[r]^i &
T &
V \ar[l]_j
}
$$
Choisissons une injection $i^{-1}\mathcal{I} \to \mathcal{I}'$
dans un faisceau injectif de $\mathbf{Z}/n\mathbf{Z}$-modules sur $Z$.
En regardant les fibres, nous voyons que le morphisme
$$
\mathcal{I} \to j_*\mathcal{I}|_V \oplus i_*\mathcal{I}'
$$
est injectif et possède donc une rétraction, puisque $\mathcal{I}$ est un faisceau injectif
de $\mathbf{Z}/n\mathbf{Z}$-modules. Il suffit donc de montrer que
$\tilde \xi$ s'envoie sur zéro dans
$$
H^q(Y, e^{-1}j_*\mathcal{I}|_V) \oplus
H^q(Y, e^{-1}i_*\mathcal{I}')
$$
au moins après avoir remplacé $X$ par les membres d'un recouvrement étale.
Remarquons que
$$
e^{-1}j_*\mathcal{I}|_V = j'_*e_V^{-1}\mathcal{I}|_V,\quad
e^{-1}i_*\mathcal{I}' = i'_*e_Z^{-1}\mathcal{I}'
$$
Par l'hypothèse de récurrence sur $q$, nous voyons que
$$
R^aj'_*e_V^{-1}\mathcal{I}|_V = 0, \quad a = 1, \ldots, q - 1
$$
Par la suite spectrale de Leray pour $j'$ et l'annulation
ci-dessus, il s'ensuit que
$$
H^q(Y, j'_*(e_V^{-1}\mathcal{I}|_V))
\longrightarrow
H^q(Y \times_T V, e_V^{-1}(\mathcal{I}|_V)) =
H^q(Y \times_T V, e^{-1}\mathcal{I}|_{Y \times_T V})
$$
est injectif. Ainsi, l'image de $\tilde \xi$ dans le premier facteur
ci-dessus s'annule, car nous savons que $\tilde \xi$ s'annule
dans $H^q(Y \times_T V, e^{-1}\mathcal{I}|_{Y \times_T V})$.
Puisque $\dim(Z) < \dim(T) = d$, l'hypothèse de récurrence montre que l'image de $\tilde \xi$
dans le second facteur
$$
H^q(Y, e^{-1}i_*\mathcal{I}') =
H^q(Y, i'_*e_Z^{-1}\mathcal{I}') =
H^q(Y \times_T Z, e_Z^{-1}\mathcal{I}')
$$
s'annule après avoir remplacé $X$ par les membres d'un recouvrement étale convenable.
Cela achève la démonstration du théorème de changement de base lisse.
\end{proof}

\begin{proof}[Seconde démonstration du changement de base lisse]
Cette démonstration est la même que la première, plus longue;
elle n'est plus courte que parce que nous avons isolé dans plusieurs
lemmes les arguments utilisés.

\medskip\noindent
Le cas $q = 0$ est le
lemme \ref{etale-cohomology-lemma-fppf-reduced-fibres-base-change-f-star}.
Nous pouvons donc supposer $q > 0$ et le résultat vrai
en tout degré inférieur.

\medskip\noindent
Pour tout $n \geq 1$ inversible sur $S$, soit $\mathcal{F}[n]$
le sous-faisceau des sections de $\mathcal{F}$ annulées par $n$. Alors
$\mathcal{F} = \colim \mathcal{F}[n]$ par notre hypothèse sur
les fibres de $\mathcal{F}$. Les foncteurs $e^{-1}$ et $f^{-1}$
commutent aux limites inductives car ce sont des adjoints à gauche. Les foncteurs
$R^qh_*$ et $R^qg_*$ commutent aux limites inductives filtrantes par le
lemme \ref{etale-cohomology-lemma-relative-colimit}.
Il suffit donc de démontrer le théorème pour $\mathcal{F}[n]$.
Désormais, nous fixons un entier $n$ inversible sur $S$ et travaillons avec des
faisceaux de $\mathbf{Z}/n\mathbf{Z}$-modules.

\medskip\noindent
Par le lemme \ref{etale-cohomology-lemma-base-change-local}, la question est locale pour la topologie étale
sur $X$ et $S$. Par la structure locale des morphismes lisses, voir le
lemme \ref{morphisms-lemma-smooth-etale-over-affine-space} de Morphismes,
nous pouvons supposer $X$ et $S$ affines et que $X \to S$
se factorise par un morphisme étale $X \to \mathbf{A}^d_S$.
En écrivant $X \to S$ comme la composée
$$
X \to \mathbf{A}^{d - 1}_S \to \mathbf{A}^{d - 2}_S \to \ldots \to
\mathbf{A}^1_S \to S
$$
nous concluons du lemme \ref{etale-cohomology-lemma-base-change-compose}
qu'il suffit de démontrer le théorème lorsque $X$ et $S$
sont affines et que $X \to S$ est de dimension relative $1$.

\medskip\noindent
Par le lemme \ref{etale-cohomology-lemma-base-change-does-not-hold}, il suffit de montrer que
$R^qh_*\mathbf{Z}/d\mathbf{Z} = 0$ pour tout $d | n$, chaque fois que nous avons un
diagramme cartésien
$$
\xymatrix{
X \ar[d] & Y \ar[d] \ar[l]^h \\
S & \Spec(K) \ar[l]
}
$$
où $X \to S$ est affine et lisse de dimension relative $1$,
$S$ est le spectre d'un anneau intègre
normal $A$, de corps des fractions algébriquement clos $L$,
et $K/L$ est une extension de corps algébriquement clos.

\medskip\noindent
Rappelons que $R^qh_*\mathbf{Z}/d\mathbf{Z}$ est le faisceau associé
au préfaisceau
$$
U \longmapsto
H^q(U \times_X Y, \mathbf{Z}/d\mathbf{Z}) =
H^q(U \times_S \Spec(K), \mathbf{Z}/d\mathbf{Z})
$$
sur $X_\etale$ (lemme \ref{etale-cohomology-lemma-higher-direct-images}).
Il suffit donc de montrer ceci : étant donnés $U$ et
$\xi \in H^q(U \times_S \Spec(K), \mathbf{Z}/d\mathbf{Z})$,
il existe un recouvrement étale $\{U_i \to U\}$
tel que $\xi$ s'annule dans $H^q(U_i \times_S \Spec(K), \mathbf{Z}/d\mathbf{Z})$.

\medskip\noindent
Nous pouvons bien entendu prendre $U$ affine. Alors $U \times_S \Spec(K)$
est une courbe affine (lisse) sur $K$, et nous avons donc l'annulation
pour $q > 1$ par le théorème \ref{etale-cohomology-theorem-vanishing-affine-curves}.

\medskip\noindent
Dernier cas : $q = 1$. Nous pouvons remplacer $U$ par les membres d'un
recouvrement étale comme dans le lemme de Compléments sur les morphismes
\ref{more-morphisms-lemma-cover-smooth-by-special}.
Alors $U \to S$ se factorise en $U \to V \to S$, où
$U \to V$ est à fibres géométriquement connexes,
$U$ et $V$ sont affines, $V \to S$ est étale et il existe
une section $\sigma : V \to U$. Par le
lemme \ref{etale-cohomology-lemma-normal-scheme-with-alg-closed-function-field},
nous voyons que $V$ est isomorphe à une réunion disjointe (finie) de
sous-schémas ouverts (affines) de $S$.
Nous pouvons clairement remplacer $S$ par l'un d'eux
et $X$ par la composante correspondante de $U$.
Nous pouvons donc supposer $X \to S$ à fibres géométriquement connexes,
muni d'une section $\sigma$, et
$\xi \in H^1(Y, \mathbf{Z}/d\mathbf{Z})$.
Puisque $K$ et $L$ sont algébriquement clos, nous avons
$$
H^1(X_L, \mathbf{Z}/d\mathbf{Z}) =
H^1(Y, \mathbf{Z}/d\mathbf{Z})
$$
Voir le lemme \ref{etale-cohomology-lemma-base-change-dim-1-separably-closed}.
Il existe donc un revêtement galoisien fini étale
$Z \to X_L$, de groupe de Galois $G \subset \mathbf{Z}/d\mathbf{Z}$,
qui annule $\xi$.
On peut le voir soit dans l'énoncé, soit dans la
démonstration du lemme \ref{etale-cohomology-lemma-smooth-base-change-fields},
soit en utilisant directement que $\xi$ correspond à un
$\mathbf{Z}/d\mathbf{Z}$-torseur sur $X_L$.
Enfin, par le
lemme \ref{pione-lemma-extend-covering-general} de Fundamental Groups,
nous trouvons un morphisme fini étale
$X' \to X$ (nécessairement surjectif)
dont la restriction à $X_L$ est $Z \to X_L$.
Puisque $\xi$ s'annule sur $X'_K$, cela achève la démonstration.
\end{proof}

\noindent
La conséquence immédiate suivante du théorème de changement de base
lisse est celle que l'on utilise souvent en pratique.

\begin{lemma}
\label{etale-cohomology-lemma-smooth-base-change-general}
Soit $S$ un schéma. Soit $S' = \lim S_i$ une limite inverse
filtrante de schémas $S_i$ lisses sur $S$, à morphismes de transition
affines. Soit $f : X \to S$ quasi-compact et quasi-séparé,
et formons le carré cartésien
$$
\xymatrix{
X' \ar[d]_{f'} \ar[r]_{g'} & X \ar[d]^f \\
S' \ar[r]^g & S
}
$$
Alors
$$
g^{-1}Rf_*E = R(f')_*(g')^{-1}E
$$
pour tout $E \in D^+(X_\etale)$ dont les faisceaux de cohomologie $H^q(E)$
ont des fibres de torsion, d'ordres inversibles sur $S$.
\end{lemma}

\begin{proof}
Considérons les suites spectrales
$$
E_2^{p, q} = R^pf_*H^q(E)
\quad\text{et}\quad
{E'}_2^{p, q} = R^pf'_*H^q((g')^{-1}E) = R^pf'_*(g')^{-1}H^q(E)
$$
convergeant vers $R^nf_*E$ et $R^nf'_*(g')^{-1}E$.
Ces suites spectrales sont construites dans le
lemme \ref{derived-lemma-two-ss-complex-functor} de Catégories dérivées.
En combinant le théorème de changement de base lisse
(théorème \ref{etale-cohomology-theorem-smooth-base-change})
avec le lemme \ref{etale-cohomology-lemma-base-change-Rf-star-colim}, nous voyons que
$$
g^{-1}R^pf_*H^q(E) = R^p(f')_*(g')^{-1}H^q(E)
$$
En combinant tout ce qui précède, nous obtenons le lemme.
\end{proof}












\section{Applications du changement de base lisse}
\label{etale-cohomology-section-applications-smooth-base-change}

\noindent
Dans cette section, nous étudions quelques conséquences plus ou moins immédiates
du théorème de changement de base lisse.

\begin{lemma}
\label{etale-cohomology-lemma-base-change-field-extension}
Soit $L/K$ une extension de corps. Soit $g : T \to S$ un morphisme quasi-compact
et quasi-séparé de schémas sur $K$. Notons
$g_L : T_L \to S_L$ le changement de base de $g$ à $\Spec(L)$.
Soit $E \in D^+(T_\etale)$ dont les faisceaux de cohomologie ont des fibres
de torsion, d'ordres inversibles dans $K$. Soit $E_L$ l'image
inverse de $E$ sur $(T_L)_\etale$. Alors
$Rg_{L, *}E_L$ est l'image inverse de $Rg_*E$ sur $S_L$.
\end{lemma}

\begin{proof}
Si $L/K$ est séparable, alors $L$ est une limite inductive filtrante de
$K$-algèbres lisses; voir le
lemme \ref{algebra-lemma-colimit-syntomic} d'Algèbre.
Dans ce cas, le lemme découle donc immédiatement du
lemme \ref{etale-cohomology-lemma-smooth-base-change-general}.
Dans le cas général, soient $K'$ et $L'$ les clôtures parfaites
(définition \ref{algebra-definition-perfection} d'Algèbre)
de $K$ et $L$. Alors $\Spec(K') \to \Spec(K)$ et
$\Spec(L') \to \Spec(L)$ sont des homéomorphismes universels,
car $K'/K$ et $L'/L$ sont purement inséparables
(voir le lemme \ref{algebra-lemma-p-ring-map} d'Algèbre).
Nous avons donc $(T_{K'})_\etale = T_\etale$,
$(S_{K'})_\etale = S_\etale$,
$(T_{L'})_\etale = (T_L)_\etale$ et
$(S_{L'})_\etale = (S_L)_\etale$, par
l'invariance topologique de la cohomologie étale; voir la
proposition \ref{etale-cohomology-proposition-topological-invariance}.
Cela ramène le lemme au cas de l'extension de corps
$L'/K'$, qui est séparable (par définition des
corps parfaits; voir la définition \ref{algebra-definition-perfect} d'Algèbre).
\end{proof}

\begin{lemma}
\label{etale-cohomology-lemma-smooth-base-change-separably-closed}
Soit $K/k$ une extension de corps séparablement clos. Soit $X$
un schéma quasi-compact et quasi-séparé sur $k$.
Soit $E \in D^+(X_\etale)$ dont les faisceaux de cohomologie ont des fibres
de torsion, d'ordres inversibles dans $k$. Alors
\begin{enumerate}
\item les morphismes $H^q_\etale(X, E) \to H^q_\etale(X_K, E|_{X_K})$
sont des isomorphismes;
\item $E \to R(X_K \to X)_*E|_{X_K}$ est un isomorphisme.
\end{enumerate}
\end{lemma}

\begin{proof}
Démonstration de (1). Soient d'abord $\overline{k}$ et $\overline{K}$
les clôtures algébriques
de $k$ et $K$. Les morphismes $\Spec(\overline{k}) \to \Spec(k)$ et
$\Spec(\overline{K}) \to \Spec(K)$ sont des homéomorphismes universels,
car $\overline{k}/k$ et $\overline{K}/K$ sont purement inséparables
(voir le lemme \ref{algebra-lemma-p-ring-map} d'Algèbre).
Ainsi $H^q_\etale(X, E) =
H^q_\etale(X_{\overline{k}}, E|_{X_{\overline{k}}})$ par
l'invariance topologique de la cohomologie étale; voir la
proposition \ref{etale-cohomology-proposition-topological-invariance}.
De même pour $X_K$ et $X_{\overline{K}}$.
Nous pouvons donc supposer $k$ et $K$ algébriquement clos.
Dans ce cas, $K$ est une limite de $k$-algèbres lisses; voir le
lemme \ref{algebra-lemma-colimit-syntomic} d'Algèbre.
Nous concluons que notre lemme est un cas particulier du
théorème \ref{etale-cohomology-theorem-smooth-base-change}, sous la forme reformulée dans le
lemme \ref{etale-cohomology-lemma-smooth-base-change-general}.

\medskip\noindent
Démonstration de (2). Pour tout $U$ quasi-compact et quasi-séparé dans $X_\etale$,
ce qui précède montre que la restriction du morphisme
$E \to R(X_K \to X)_*E|_{X_K}$ détermine un isomorphisme en cohomologie.
Puisque tout objet de $X_\etale$ possède un recouvrement étale par de tels $U$,
cela démontre l'assertion voulue.
\end{proof}

\begin{lemma}
\label{etale-cohomology-lemma-base-change-does-not-hold-post}
Avec $f : X \to S$ et $n$ comme dans la remarque \ref{etale-cohomology-remark-base-change-holds},
supposons $n$ inversible sur $S$ et, pour un certain $q \geq 1$,
$BC(f, n, q - 1)$ vraie, mais $BC(f, n, q)$ fausse.
Il existe alors un diagramme commutatif
$$
\xymatrix{
X \ar[d]_f & X' \ar[d] \ar[l] & Y \ar[l]^h \ar[d] \\
S & S' \ar[l] & \Spec(K) \ar[l]
}
$$
dont les deux carrés sont cartésiens, où $S'$ est affine, intègre et normal,
de corps des fonctions algébriquement clos $K$, ainsi qu'un entier
$d | n$ tel que $R^qh_*(\mathbf{Z}/d\mathbf{Z})$ soit non nul.
\end{lemma}

\begin{proof}
Choisissons d'abord un diagramme et $\mathcal{F}$ comme dans le
lemme \ref{etale-cohomology-lemma-base-change-does-not-hold}.
Nous pouvons supposer $S'$ affine, et nous le faisons (c'est évident, mais
voir la démonstration du lemme en cas de doute).
Soit $K'$ le corps des fonctions de $S'$,
et soit $Y' = X' \times_{S'} \Spec(K')$; nous obtenons le diagramme
$$
\xymatrix{
X \ar[d]_f &
X' \ar[d] \ar[l] &
Y' \ar[l]^{h'} \ar[d] &
Y \ar[l] \ar[d] \\
S &
S' \ar[l] &
\Spec(K') \ar[l] &
\Spec(K) \ar[l]
}
$$
Par le lemme \ref{etale-cohomology-lemma-smooth-base-change-separably-closed},
l'image directe totale $R(Y \to Y')_*\mathbf{Z}/d\mathbf{Z}$
est isomorphe à $\mathbf{Z}/d\mathbf{Z}$ dans $D(Y'_\etale)$; nous
utilisons ici que $n$ est inversible sur $S$.
Ainsi $Rh'_*\mathbf{Z}/d\mathbf{Z} = Rh_*\mathbf{Z}/d\mathbf{Z}$
par la suite spectrale de Leray relative. Cela achève la démonstration.
\end{proof}










\section{Le théorème de changement de base propre}
\label{etale-cohomology-section-proper-base-change}

\noindent
Le théorème de changement de base propre est énoncé et démontré dans cette section.
Notre approche suit approximativement la démonstration de \cite[XII, théorème 5.1]{SGA4},
en utilisant les idées de Gabber (issues du cas affine) pour simplifier
légèrement les arguments.

\begin{lemma}
\label{etale-cohomology-lemma-zariski-h0-proper-over-henselian-pair}
Soit $(A, I)$ un couple hensélien. Soit $f : X \to \Spec(A)$ un morphisme propre
de schémas. Soit $Z = X \times_{\Spec(A)} \Spec(A/I)$. Pour tout
faisceau $\mathcal{F}$ sur l'espace topologique associé à $X$, nous
avons $\Gamma(X, \mathcal{F}) = \Gamma(Z, \mathcal{F}|_Z)$.
\end{lemma}

\begin{proof}
Nous utiliserons le lemme \ref{etale-cohomology-lemma-h0-topological} pour le démontrer. Remarquons d'abord
que l'espace topologique sous-jacent à $X$ est spectral par le lemme de Properties
\ref{properties-lemma-quasi-compact-quasi-separated-spectral}.
Soit $Y \subset X$ un sous-schéma fermé irréductible. Pour achever la démonstration,
montrons que $Y \cap Z = Y \times_{\Spec(A)} \Spec(A/I)$ est connexe.
En remplaçant $X$ par $Y$,
nous pouvons supposer $X$ irréductible, et nous devons montrer que $Z$
est connexe. Soit $X \to \Spec(B) \to \Spec(A)$ la factorisation de Stein
de $f$ (théorème de Compléments sur les morphismes
\ref{more-morphisms-theorem-stein-factorization-general}).
Alors $A \to B$ est entier et $(B, IB)$ est un couple hensélien
(lemme \ref{more-algebra-lemma-integral-over-henselian-pair} de Compléments d'algèbre).
Nous pouvons donc supposer que les fibres de $X \to \Spec(A)$ sont géométriquement
connexes. D'autre part, l'image $T \subset \Spec(A)$ de $f$
est irréductible et fermée puisque $X$ est propre sur $A$. Ainsi $T \cap V(I)$
est connexe par le lemme de Compléments d'algèbre
\ref{more-algebra-lemma-irreducible-henselian-pair-connected}.
Or $Y \times_{\Spec(A)} \Spec(A/I) \to T \cap V(I)$
est une application fermée surjective à fibres connexes.
Le résultat découle alors du lemme de Topologie
\ref{topology-lemma-connected-fibres-quotient-topology-connected-components}.
\end{proof}

\begin{lemma}
\label{etale-cohomology-lemma-h0-proper-over-henselian-pair}
Soit $(A, I)$ un couple hensélien. Soit $f : X \to \Spec(A)$ un morphisme propre
de schémas. Soit $i : Z \to X$ l'immersion fermée de
$X \times_{\Spec(A)} \Spec(A/I)$ dans $X$. Pour tout
faisceau $\mathcal{F}$ sur $X_\etale$, nous
avons $\Gamma(X, \mathcal{F}) = \Gamma(Z, i_{small}^{-1}\mathcal{F})$.
\end{lemma}

\begin{proof}
Cela découle des lemmes \ref{etale-cohomology-lemma-gabber-h0} et
\ref{etale-cohomology-lemma-zariski-h0-proper-over-henselian-pair}
ainsi que du fait que tout schéma fini sur $X$ est propre sur $\Spec(A)$.
\end{proof}

\begin{lemma}
\label{etale-cohomology-lemma-h0-proper-over-henselian-local}
Soit $A$ un anneau local hensélien. Soit $f : X \to \Spec(A)$
un morphisme propre de schémas. Soit $X_0 \subset X$ la fibre de
$f$ au-dessus du point fermé. Pour tout faisceau $\mathcal{F}$ sur $X_\etale$, nous
avons $\Gamma(X, \mathcal{F}) = \Gamma(X_0, \mathcal{F}|_{X_0})$.
\end{lemma}

\begin{proof}
C'est un cas particulier du lemme \ref{etale-cohomology-lemma-h0-proper-over-henselian-pair}.
\end{proof}

\noindent
Soit $f : X \to S$ un morphisme de schémas. Soit
$\overline{s} : \Spec(k) \to S$ un point géométrique. La fibre
de $f$ en $\overline{s}$ est le schéma
$X_{\overline{s}} = \Spec(k) \times_{\overline{s}, S} X$, considéré
comme un schéma sur $\Spec(k)$. Si $\mathcal{F}$ est un faisceau sur
$X_\etale$, notons
$\mathcal{F}_{\overline{s}} = p_{small}^{-1}\mathcal{F}$
l'image inverse de $\mathcal{F}$ sur $(X_{\overline{s}})_\etale$.
Dans la suite, nous considérerons l'ensemble
$$
\Gamma(X_{\overline{s}}, \mathcal{F}_{\overline{s}})
$$
Soit $s \in S$ le point image de $\overline{s}$. Soit $\kappa(s)^{sep}$
la clôture séparable de $\kappa(s)$ dans $k$, comme dans la
définition \ref{etale-cohomology-definition-algebraic-geometric-point}.
Par le lemme \ref{etale-cohomology-lemma-sections-base-field-extension},
l'image inverse définit une bijection
$$
\Gamma(X_{\kappa(s)^{sep}},
p_{sep}^{-1} \mathcal{F})
\longrightarrow
\Gamma(X_{\overline{s}}, \mathcal{F}_{\overline{s}})
$$
où $p_{sep} : X_{\kappa(s)^{sep}} = \Spec(\kappa(s)^{sep}) \times_S X \to X$
est la projection.

\begin{lemma}
\label{etale-cohomology-lemma-proper-pushforward-stalk}
Soit $f : X \to S$ un morphisme propre de schémas. Soit
$\overline{s} \to S$ un point géométrique.
Pour tout faisceau $\mathcal{F}$ sur $X_\etale$,
le morphisme canonique
$$
(f_*\mathcal{F})_{\overline{s}} \longrightarrow
\Gamma(X_{\overline{s}}, \mathcal{F}_{\overline{s}})
$$
est bijectif.
\end{lemma}

\begin{proof}
Par le théorème \ref{etale-cohomology-theorem-higher-direct-images} (pour les faisceaux d'ensembles),
nous avons
$$
(f_*\mathcal{F})_{\overline{s}} =
\Gamma(X \times_S \Spec(\mathcal{O}_{S, \overline{s}}^{sh}),
p_{small}^{-1}\mathcal{F})
$$
où $p : X \times_S \Spec(\mathcal{O}_{S, \overline{s}}^{sh}) \to X$
est la projection. Puisque le corps résiduel de l'anneau local strictement
hensélien $\mathcal{O}_{S, \overline{s}}^{sh}$ est $\kappa(s)^{sep}$,
le lemme résulte de la discussion précédente et du
lemme \ref{etale-cohomology-lemma-h0-proper-over-henselian-local}.
\end{proof}

\begin{lemma}
\label{etale-cohomology-lemma-proper-base-change-f-star}
Soit $f : X \to Y$ un morphisme propre de schémas. Soit $g : Y' \to Y$
un morphisme de schémas. Posons $X' = Y' \times_Y X$, de projections
$f' : X' \to Y'$ et $g' : X' \to X$. Soit $\mathcal{F}$ un faisceau quelconque sur
$X_\etale$. Alors $g^{-1}f_*\mathcal{F} = f'_*(g')^{-1}\mathcal{F}$.
\end{lemma}

\begin{proof}
Il existe un morphisme canonique $g^{-1}f_*\mathcal{F} \to f'_*(g')^{-1}\mathcal{F}$.
En effet, il est adjoint au morphisme
$$
f_*\mathcal{F} \longrightarrow
g_*f'_*(g')^{-1}\mathcal{F} = f_*g'_*(g')^{-1}\mathcal{F}
$$
qui est obtenu en appliquant $f_*$ au morphisme canonique
$\mathcal{F} \to g'_*(g')^{-1}\mathcal{F}$. Pour vérifier que ce morphisme est un
isomorphisme, nous pouvons calculer ce qui se passe sur les fibres.
Soit $y' : \Spec(k) \to Y'$ un point géométrique d'image $y$ dans $Y$.
Par le lemme \ref{etale-cohomology-lemma-proper-pushforward-stalk}, les fibres sont respectivement
$\Gamma(X'_{y'}, \mathcal{F}_{y'})$ et $\Gamma(X_y, \mathcal{F}_y)$.
Ici, les faisceaux $\mathcal{F}_y$ et $\mathcal{F}_{y'}$
sont les images inverses de $\mathcal{F}$ par les projections $X_y \to X$
et $X'_{y'} \to X$. Puisque $X_y = X'_{y'}$ comme schémas,
nous voyons que les fibres coïncident. Nous omettons la
vérification de la compatibilité de cet isomorphisme avec notre morphisme.
\end{proof}

\noindent
Nous commençons maintenant l'étude du théorème de changement de base propre.
Pour cela, introduisons quelques notations. Considérons un diagramme commutatif
\begin{equation}
\label{etale-cohomology-equation-base-change-diagram}
\vcenter{
\xymatrix{
X' \ar[r]_{g'} \ar[d]_{f'} & X \ar[d]^f \\
Y' \ar[r]^g & Y
}
}
\end{equation}
de morphismes de schémas. Nous obtenons alors un diagramme commutatif de sites
$$
\xymatrix{
X'_\etale \ar[r]_{g'_{small}} \ar[d]_{f'_{small}} &
X_\etale \ar[d]^{f_{small}} \\
Y'_\etale \ar[r]^{g_{small}} &
Y_\etale
}
$$
Pour tout objet $E$ de $D(X_\etale)$, nous obtenons un morphisme canonique de changement de base
\begin{equation}
\label{etale-cohomology-equation-base-change}
g_{small}^{-1}Rf_{small, *}E \longrightarrow Rf'_{small, *}(g'_{small})^{-1}E
\end{equation}
dans $D(Y'_\etale)$. Voir la remarque de Cohomologie sur les sites
\ref{sites-cohomology-remark-base-change}, où nous employons le faisceau
constant $\mathbf{Z}$ comme faisceau d'anneaux.
Nous omettrons généralement les indices ${}_{small}$ dans cette formule.
Par exemple, si $E = \mathcal{F}[0]$, où $\mathcal{F}$ est un faisceau
abélien sur $X_\etale$, le morphisme de changement de base est un morphisme
\begin{equation}
\label{etale-cohomology-equation-base-change-sheaf}
g^{-1}Rf_*\mathcal{F} \longrightarrow Rf'_*(g')^{-1}\mathcal{F}
\end{equation}
dans $D(Y'_\etale)$.

\medskip\noindent
Le morphisme (\ref{etale-cohomology-equation-base-change}) n'a aucune chance d'être un isomorphisme
dans la généralité ci-dessus. Le but est de montrer
qu'il est un isomorphisme si le diagramme (\ref{etale-cohomology-equation-base-change-diagram})
est cartésien, si $f : X \to Y$ est propre, si les faisceaux de cohomologie de $E$
sont de torsion et si $E$ est borné inférieurement.
Pour étudier cette question, introduisons la terminologie suivante.
Nous dirons que {\it la cohomologie commute au changement de base
pour $f : X \to Y$} si (\ref{etale-cohomology-equation-base-change-sheaf})
est un isomorphisme pour tout diagramme (\ref{etale-cohomology-equation-base-change-diagram})
où $X' = Y' \times_Y X$, et pour tout faisceau abélien de torsion $\mathcal{F}$.

\begin{lemma}
\label{etale-cohomology-lemma-proper-base-change-in-terms-of-injectives}
Soit $f : X \to Y$ un morphisme propre de schémas.
Les assertions suivantes sont équivalentes :
\begin{enumerate}
\item la cohomologie commute au changement de base pour $f$ (voir ci-dessus);
\item pour tout nombre premier $\ell$, tout faisceau injectif
de $\mathbf{Z}/\ell\mathbf{Z}$-modules $\mathcal{I}$
sur $X_\etale$ et tout diagramme (\ref{etale-cohomology-equation-base-change-diagram})
où $X' = Y' \times_Y X$, les faisceaux
$R^qf'_*(g')^{-1}\mathcal{I}$ sont nuls pour $q > 0$.
\end{enumerate}
\end{lemma}

\begin{proof}
Il est clair que (1) implique (2). Réciproquement, supposons (2) et soit
$\mathcal{F}$ un faisceau abélien de torsion sur $X_\etale$. Soit $Y' \to Y$
un morphisme de schémas et soit $X' = Y' \times_Y X$,
de projections $g' : X' \to X$ et $f' : X' \to Y'$, comme dans le
diagramme (\ref{etale-cohomology-equation-base-change-diagram}).
Nous voulons montrer que les morphismes de faisceaux
$$
g^{-1}R^qf_*\mathcal{F} \longrightarrow R^qf'_*(g')^{-1}\mathcal{F}
$$
sont des isomorphismes pour tout $q \geq 0$.

\medskip\noindent
Pour tout $n \geq 1$, soit $\mathcal{F}[n]$ le sous-faisceau des sections
de $\mathcal{F}$ annulées par $n$. Alors
$\mathcal{F} = \colim \mathcal{F}[n]$.
Les foncteurs $g^{-1}$ et $(g')^{-1}$ commutent aux limites inductives quelconques
(en tant qu'adjoints à gauche). La prise des images directes supérieures par $f$ ou $f'$
commute aux limites inductives filtrantes par le lemme \ref{etale-cohomology-lemma-relative-colimit}.
Nous voyons donc que
$$
g^{-1}R^qf_*\mathcal{F} = \colim g^{-1}R^qf_*\mathcal{F}[n]
\quad\text{et}\quad
R^qf'_*(g')^{-1}\mathcal{F} =
\colim R^qf'_*(g')^{-1}\mathcal{F}[n]
$$
Il suffit donc de démontrer le résultat lorsque $\mathcal{F}$ est
annulé par un entier positif $n$.

\medskip\noindent
Si $n = \ell n'$ pour un certain nombre premier $\ell$, nous obtenons une suite
exacte courte
$$
0 \to \mathcal{F}[\ell] \to \mathcal{F} \to
\mathcal{F}/\mathcal{F}[\ell] \to 0
$$
Remarquons que $\mathcal{F}/\mathcal{F}[\ell]$ est annulé par $n'$.
De plus, si le résultat vaut à la fois pour $\mathcal{F}[\ell]$ et
$\mathcal{F}/\mathcal{F}[\ell]$, alors le résultat découle de
la suite exacte longue des images directes supérieures (et du lemme des $5$).
Nous nous ramenons ainsi au cas où $\mathcal{F}$ est annulé
par un nombre premier $\ell$.

\medskip\noindent
Supposons que $\mathcal{F}$ soit annulé par un nombre premier $\ell$.
Choisissons une résolution injective $\mathcal{F} \to \mathcal{I}^\bullet$
dans $D(X_\etale, \mathbf{Z}/\ell\mathbf{Z})$. En appliquant l'hypothèse
(2) et le lemme d'acyclicité de Leray
(lemme \ref{derived-lemma-leray-acyclicity} de Catégories dérivées),
nous voyons que
$$
f'_*(g')^{-1}\mathcal{I}^\bullet
$$
calcule $Rf'_*(g')^{-1}\mathcal{F}$. Nous concluons en appliquant le
lemme \ref{etale-cohomology-lemma-proper-base-change-f-star}.
\end{proof}

\begin{lemma}
\label{etale-cohomology-lemma-sandwich}
Soient $f : X \to Y$ et $g : Y \to Z$ des morphismes propres de schémas. Supposons que :
\begin{enumerate}
\item la cohomologie commute au changement de base pour $f$;
\item la cohomologie commute au changement de base pour $g \circ f$;
\item $f$ est surjectif.
\end{enumerate}
Alors la cohomologie commute au changement de base pour $g$.
\end{lemma}

\begin{proof}
Nous utiliserons sans autre mention l'équivalence du
lemme \ref{etale-cohomology-lemma-proper-base-change-in-terms-of-injectives}.
Soit $\ell$ un nombre premier.
Soit $\mathcal{I}$ un faisceau injectif de
$\mathbf{Z}/\ell\mathbf{Z}$-modules sur $Y_\etale$.
Choisissons un morphisme injectif de faisceaux $f^{-1}\mathcal{I} \to \mathcal{J}$,
où $\mathcal{J}$ est un faisceau injectif de
$\mathbf{Z}/\ell\mathbf{Z}$-modules sur $X_\etale$.
Puisque $f$ est surjectif, le morphisme $\mathcal{I} \to f_*\mathcal{J}$
est injectif (il suffit de le vérifier sur les fibres aux points géométriques).
Puisque $\mathcal{I}$ est injectif, nous voyons que $\mathcal{I}$
est un facteur direct de $f_*\mathcal{J}$. Il suffit donc
de démontrer l'annulation voulue pour $f_*\mathcal{J}$.

\medskip\noindent
Soit $Z' \to Z$ un morphisme de schémas et posons
$Y' = Z' \times_Z Y$ et $X' = Z' \times_Z X = Y' \times_ Y X$.
Notons $a : X' \to X$, $b : Y' \to Y$ et $c : Z' \to Z$ les
projections. De même pour $f' : X' \to Y'$ et $g' : Y' \to Z'$.
Par le lemme \ref{etale-cohomology-lemma-proper-base-change-f-star}, nous avons
$b^{-1}f_*\mathcal{J} = f'_*a^{-1}\mathcal{J}$.
D'autre part, nous savons que $R^qf'_*a^{-1}\mathcal{J}$ et
$R^q(g' \circ f')_*a^{-1}\mathcal{J}$ sont nuls pour $q > 0$.
En utilisant la suite spectrale
(lemme \ref{sites-cohomology-lemma-relative-Leray} de Cohomologie sur les sites)
$$
R^pg'_* R^qf'_* a^{-1}\mathcal{J} \Rightarrow
R^{p + q}(g' \circ f')_* a^{-1}\mathcal{J}
$$
nous concluons que
$ R^pg'_*(b^{-1}f_*\mathcal{J}) = R^pg'_*(f'_*a^{-1}\mathcal{J}) = 0$
pour $p > 0$, comme voulu.
\end{proof}

\begin{lemma}
\label{etale-cohomology-lemma-composition}
Soient $f : X \to Y$ et $g : Y \to Z$ des morphismes propres de schémas. Supposons que :
\begin{enumerate}
\item la cohomologie commute au changement de base pour $f$;
\item la cohomologie commute au changement de base pour $g$.
\end{enumerate}
Alors la cohomologie commute au changement de base pour $g \circ f$.
\end{lemma}

\begin{proof}
Nous utiliserons sans autre mention l'équivalence du
lemme \ref{etale-cohomology-lemma-proper-base-change-in-terms-of-injectives}.
Soit $\ell$ un nombre premier.
Soit $\mathcal{I}$ un faisceau injectif de
$\mathbf{Z}/\ell\mathbf{Z}$-modules sur $X_\etale$.
Alors $f_*\mathcal{I}$ est un faisceau injectif de
$\mathbf{Z}/\ell\mathbf{Z}$-modules sur $Y_\etale$
(lemme de Cohomologie sur les sites
\ref{sites-cohomology-lemma-pushforward-injective-flat}).
Le résultat en découle formellement, mais nous allons aussi
l'expliciter.

\medskip\noindent
Soit $Z' \to Z$ un morphisme de schémas et posons
$Y' = Z' \times_Z Y$ et $X' = Z' \times_Z X = Y' \times_ Y X$.
Notons $a : X' \to X$, $b : Y' \to Y$ et $c : Z' \to Z$ les
projections. De même pour $f' : X' \to Y'$ et $g' : Y' \to Z'$.
Par le lemme \ref{etale-cohomology-lemma-proper-base-change-f-star}, nous avons
$b^{-1}f_*\mathcal{I} = f'_*a^{-1}\mathcal{I}$.
D'autre part, nous savons que $R^qf'_*a^{-1}\mathcal{I}$ et
$R^q(g')_*b^{-1}f_*\mathcal{I}$ sont nuls pour $q > 0$.
En utilisant la suite spectrale
(lemme \ref{sites-cohomology-lemma-relative-Leray} de Cohomologie sur les sites)
$$
R^pg'_* R^qf'_* a^{-1}\mathcal{I} \Rightarrow
R^{p + q}(g' \circ f')_* a^{-1}\mathcal{I}
$$
nous concluons que $R^p(g' \circ f')_*a^{-1}\mathcal{I} = 0$ pour
$p > 0$, comme voulu.
\end{proof}

\begin{lemma}
\label{etale-cohomology-lemma-finite}
\begin{slogan}
Le changement de base propre en cohomologie étale est valable pour les morphismes finis.
\end{slogan}
Soit $f : X \to Y$ un morphisme fini de schémas.
Alors la cohomologie commute au changement de base pour $f$.
\end{lemma}

\begin{proof}
Remarquons qu'un morphisme fini est propre; voir le
lemme \ref{morphisms-lemma-finite-proper} de Morphismes.
De plus, le changement de base d'un morphisme fini est fini; voir le
lemme \ref{morphisms-lemma-base-change-finite} de Morphismes.
Le résultat découle donc du
lemme \ref{etale-cohomology-lemma-proper-base-change-in-terms-of-injectives},
combiné à la
proposition \ref{etale-cohomology-proposition-finite-higher-direct-image-zero}.
\end{proof}

\begin{lemma}
\label{etale-cohomology-lemma-reduce-to-P1}
Pour démontrer que la cohomologie commute au changement de base pour
tout morphisme propre de schémas, il suffit de démontrer qu'elle
commute pour le morphisme $\mathbf{P}^1_S \to S$, pour tout schéma $S$.
\end{lemma}

\begin{proof}
Soit $f : X \to Y$ un morphisme propre de schémas.
Soit $Y = \bigcup Y_i$ un recouvrement ouvert affine,
et posons $X_i = f^{-1}(Y_i)$. Si nous pouvons démontrer
que la cohomologie commute au changement de base pour $X_i \to Y_i$,
alors elle commute au changement de base pour $f$.
En effet, la formation des images directes supérieures
commute à la localisation de Zariski (et même \'etale)
sur la base; voir le
lemme \ref{etale-cohomology-lemma-higher-direct-images}.
Nous pouvons donc supposer $Y$ affine.

\medskip\noindent
Soit $Y$ un schéma affine et soit $X \to Y$ un morphisme propre.
Par le lemme de Chow, il existe un diagramme commutatif
$$
\xymatrix{
X \ar[rd] & X' \ar[d] \ar[l]^\pi \ar[r] & \mathbf{P}^n_Y \ar[dl] \\
& Y &
}
$$
où $X' \to \mathbf{P}^n_Y$ est une immersion et
$\pi : X' \to X$ est propre et surjectif; voir le
lemme \ref{limits-lemma-chow-finite-type} de Limits.
Puisque $X \to Y$ est propre, on voit que $X' \to Y$ est propre
(lemme \ref{morphisms-lemma-composition-proper} de Morphismes).
Par conséquent, $X' \to \mathbf{P}^n_Y$ est une immersion fermée
(lemme \ref{morphisms-lemma-image-proper-scheme-closed} de Morphismes).
Il s'ensuit que $X' \to X \times_Y \mathbf{P}^n_Y = \mathbf{P}^n_X$
est une immersion fermée (en tant qu'immersion d'image fermée).

\medskip\noindent
Par le lemme \ref{etale-cohomology-lemma-sandwich},
il suffit de démontrer que la cohomologie commute au changement de base pour
$\pi$ et $X' \to Y$. Ces deux morphismes se factorisent en une immersion
fermée suivie d'une projection $\mathbf{P}^n_S \to S$ (pour un certain $S$).
Par le lemme \ref{etale-cohomology-lemma-finite}, le résultat vaut pour les immersions
fermées (car les immersions fermées sont finies).
Par le lemme \ref{etale-cohomology-lemma-composition}, il suffit de démontrer le
résultat pour les projections $\mathbf{P}^n_S \to S$.

\medskip\noindent
Pour tout $n \geq 1$, il existe un morphisme fini surjectif
$$
\mathbf{P}^1_S \times_S \ldots \times_S \mathbf{P}^1_S
\longrightarrow
\mathbf{P}^n_S
$$
donné en coordonnées par
$$
((x_1 : y_1), (x_2 : y_2), \ldots, (x_n : y_n))
\longmapsto
(F_0 : \ldots : F_n)
$$
où $F_0, \ldots, F_n$ sont les polynômes en $x_1, \ldots, y_n$
à coefficients entiers tels que
$$
\prod (x_i t + y_i) = F_0 t^n + F_1 t^{n - 1} + \ldots + F_n
$$
En appliquant encore une fois
les lemmes \ref{etale-cohomology-lemma-sandwich}, \ref{etale-cohomology-lemma-finite} et \ref{etale-cohomology-lemma-composition},
nous concluons que le lemme est vrai.
\end{proof}

\begin{theorem}[Changement de base propre]
\label{etale-cohomology-theorem-proper-base-change}
Soit $f : X \to Y$ un morphisme propre de schémas. Soit $g : Y' \to Y$
un morphisme de schémas. Posons $X' = Y' \times_Y X$
et considérons le diagramme cartésien
$$
\xymatrix{
X' \ar[r]_{g'} \ar[d]_{f'} & X \ar[d]^f \\
Y' \ar[r]^g & Y
}
$$
Soit $\mathcal{F}$ un faisceau abélien de torsion sur $X_\etale$.
Alors le morphisme de changement de base
$$
g^{-1}Rf_*\mathcal{F} \longrightarrow Rf'_*(g')^{-1}\mathcal{F}
$$
est un isomorphisme.
\end{theorem}

\begin{proof}
Dans la terminologie introduite ci-dessus, cela signifie que la cohomologie commute
au changement de base pour tout morphisme propre de schémas. Par le
lemme \ref{etale-cohomology-lemma-reduce-to-P1},
il suffit de démontrer que la cohomologie commute au changement de base
pour le morphisme $\mathbf{P}^1_S \to S$, pour tout schéma $S$.

\medskip\noindent
Soit $S$ le spectre d'un anneau local strictement hensélien, de point
fermé $s$. Posons $X = \mathbf{P}^1_S$ et $X_0 = X_s = \mathbf{P}^1_s$.
Soit $\mathcal{F}$ un faisceau de $\mathbf{Z}/\ell\mathbf{Z}$-modules
sur $X_\etale$. Le point clef de notre démonstration est l'égalité
$$
H^q_\etale(X, \mathcal{F}) = H^q_\etale(X_0, \mathcal{F}|_{X_0}).
$$
En effet, choisissons une résolution $\mathcal{F} \to \mathcal{I}^\bullet$
par des faisceaux injectifs de $\mathbf{Z}/\ell\mathbf{Z}$-modules.
Alors $\mathcal{I}^\bullet|_{X_0}$ est une résolution de $\mathcal{F}|_{X_0}$
par des objets acycliques à droite pour $H^0_\etale(X_0, -)$; voir le
lemme \ref{etale-cohomology-lemma-efface-cohomology-on-fibre-by-finite-cover}.
D'après le lemme d'acyclicité de Leray, le membre de droite est calculé par
le complexe $H^0_\etale(X_0, \mathcal{I}^\bullet|_{X_0})$,
qui est égal à $H^0_\etale(X, \mathcal{I}^\bullet)$ par le
lemme \ref{etale-cohomology-lemma-h0-proper-over-henselian-local}. Ce complexe
calcule le membre de gauche.

\medskip\noindent
Supposons maintenant $S$ quelconque et $\mathcal{F}$ un faisceau de
$\mathbf{Z}/\ell\mathbf{Z}$-modules sur $X_\etale$.
Soit $\overline{s} : \Spec(k) \to S$ un point géométrique
de $S$ au-dessus de $s \in S$. Nous avons
$$
(R^qf_*\mathcal{F})_{\overline{s}} =
H^q_\etale(\mathbf{P}^1_{\mathcal{O}_{S, \overline{s}}^{sh}},
\mathcal{F}|_{\mathbf{P}^1_{\mathcal{O}_{S, \overline{s}}^{sh}}}) =
H^q_\etale(\mathbf{P}^1_{\kappa(s)^{sep}},
\mathcal{F}|_{\mathbf{P}^1_{\kappa(s)^{sep}}})
$$
où $\kappa(s)^{sep}$ est le corps résiduel de
$\mathcal{O}_{S, \overline{s}}^{sh}$, c'est-à-dire la clôture algébrique
séparable de $\kappa(s)$ dans $k$.
La première égalité découle du théorème \ref{etale-cohomology-theorem-higher-direct-images},
et la seconde de la formule affichée au
paragraphe précédent.

\medskip\noindent
Enfin, considérons un morphisme quelconque de schémas $g : T \to S$, où
$S$ et $\mathcal{F}$ sont comme ci-dessus.
Soit $f' : \mathbf{P}^1_T \to T$ la projection et soit
$g' : \mathbf{P}^1_T \to \mathbf{P}^1_S$ le morphisme induit
par $g$. Considérons le morphisme de changement de base
$$
g^{-1}R^qf_*\mathcal{F}
\longrightarrow
R^qf'_*(g')^{-1}\mathcal{F}
$$
Soit $\overline{t}$ un point géométrique de $T$, d'image
$\overline{s} = g(\overline{t})$. D'après la discussion
ci-dessus, le morphisme induit sur les fibres en $\overline{t}$ est le morphisme
$$
H^q_\etale(\mathbf{P}^1_{\kappa(s)^{sep}},
\mathcal{F}|_{\mathbf{P}^1_{\kappa(s)^{sep}}})
\longrightarrow
H^q_\etale(\mathbf{P}^1_{\kappa(t)^{sep}},
\mathcal{F}|_{\mathbf{P}^1_{\kappa(t)^{sep}}})
$$
Puisque $\kappa(s)^{sep} \subset \kappa(t)^{sep}$, ce morphisme est un
isomorphisme par le lemme \ref{etale-cohomology-lemma-base-change-dim-1-separably-closed}.

\medskip\noindent
Cela démontre que la cohomologie commute au changement de base pour
$\mathbf{P}^1_S \to S$, avec des faisceaux de $\mathbf{Z}/\ell\mathbf{Z}$-modules.
En particulier, pour un faisceau injectif de $\mathbf{Z}/\ell\mathbf{Z}$-modules,
les images directes supérieures de tout changement de base sont nulles.
Autrement dit, la condition (2) du
lemme \ref{etale-cohomology-lemma-proper-base-change-in-terms-of-injectives}
est satisfaite, ce qui achève la démonstration.
\end{proof}

\begin{lemma}
\label{etale-cohomology-lemma-proper-base-change}
Soit $f : X \to Y$ un morphisme propre de schémas. Soit $g : Y' \to Y$
un morphisme de schémas. Posons $X' = Y' \times_Y X$ et notons
$f' : X' \to Y'$ et $g' : X' \to X$ les projections.
Soit $E \in D^+(X_\etale)$ à faisceaux de cohomologie de torsion.
Alors le morphisme de changement de base (\ref{etale-cohomology-equation-base-change})
$g^{-1}Rf_*E \to Rf'_*(g')^{-1}E$
est un isomorphisme.
\end{lemma}

\begin{proof}
C'est une conséquence simple du théorème de changement de base propre
(théorème \ref{etale-cohomology-theorem-proper-base-change}), en utilisant les suites
spectrales
$$
E_2^{p, q} = R^pf_*H^q(E)
\quad\text{et}\quad
{E'}_2^{p, q} = R^pf'_*(g')^{-1}H^q(E)
$$
convergeant vers $R^nf_*E$ et $R^nf'_*(g')^{-1}E$.
Les suites spectrales sont construites dans le
lemme \ref{derived-lemma-two-ss-complex-functor} de Catégories dérivées.
Certains détails sont omis.
\end{proof}

\begin{lemma}
\label{etale-cohomology-lemma-proper-base-change-stalk}
Soit $f : X \to Y$ un morphisme propre de schémas. Soit $\overline{y} \to Y$
un point géométrique.
\begin{enumerate}
\item Pour un faisceau abélien de torsion $\mathcal{F}$ sur $X_\etale$, nous avons
$(R^nf_*\mathcal{F})_{\overline{y}} =
H^n_\etale(X_{\overline{y}}, \mathcal{F}_{\overline{y}})$.
\item Pour $E \in D^+(X_\etale)$ à faisceaux de cohomologie de torsion, nous avons
$(R^nf_*E)_{\overline{y}} =
H^n_\etale(X_{\overline{y}}, E|_{X_{\overline{y}}})$.
\end{enumerate}
\end{lemma}

\begin{proof}
Dans l'énoncé, $\mathcal{F}_{\overline{y}}$ désigne l'image inverse
de $\mathcal{F}$ sur la fibre schématique
$X_{\overline{y}} = \overline{y} \times_Y X$.
Puisque l'image inverse par $\overline{y} \to Y$ fournit la
fibre de $\mathcal{F}$, la première assertion du lemme
est un cas particulier du théorème \ref{etale-cohomology-theorem-proper-base-change}.
La seconde est un cas particulier du lemme \ref{etale-cohomology-lemma-proper-base-change}.
\end{proof}






\section{Applications du changement de base propre}
\label{etale-cohomology-section-applications-proper-base-change}

\noindent
Dans cette section, nous étudions quelques conséquences plus ou moins immédiates
du théorème de changement de base propre.

\begin{lemma}
\label{etale-cohomology-lemma-base-change-separably-closed}
Soit $K/k$ une extension de corps séparablement clos.
Soit $X$ un schéma propre sur $k$.
Soit $\mathcal{F}$ un faisceau abélien de torsion sur $X_\etale$.
Alors le morphisme $H^q_\etale(X, \mathcal{F}) \to
H^q_\etale(X_K, \mathcal{F}|_{X_K})$ est un isomorphisme
pour $q \geq 0$.
\end{lemma}

\begin{proof}
L'examen des fibres montre qu'il s'agit
d'un cas particulier du
théorème \ref{etale-cohomology-theorem-proper-base-change}.
\end{proof}

\begin{lemma}
\label{etale-cohomology-lemma-cohomological-dimension-proper}
Soit $f : X \to Y$ un morphisme propre de schémas
dont toutes les fibres sont de dimension $\leq n$.
Alors, pour tout faisceau abélien de torsion $\mathcal{F}$ sur $X_\etale$,
nous avons $R^qf_*\mathcal{F} = 0$ pour $q > 2n$.
\end{lemma}

\begin{proof}
Nous le démontrons par récurrence sur $n$, pour tous les morphismes propres.

\medskip\noindent
Si $n = 0$, alors $f$ est un morphisme fini
(lemme \ref{more-morphisms-lemma-characterize-finite} de Compléments sur les morphismes),
et le
résultat est vrai par la proposition \ref{etale-cohomology-proposition-finite-higher-direct-image-zero}.

\medskip\noindent
Si $n > 0$, le lemme \ref{etale-cohomology-lemma-proper-base-change-stalk}
montre qu'il suffit de démontrer $H^i_\etale(X, \mathcal{F}) = 0$
pour $i > 2n$, lorsque $X$ est un schéma propre, $\dim(X) \leq n$,
sur un corps algébriquement clos $k$,
et que $\mathcal{F}$ est un faisceau abélien de torsion sur $X$.

\medskip\noindent
Si $n = 1$, cela découle du théorème \ref{etale-cohomology-theorem-vanishing-curves}.
Supposons $n > 1$. Par la proposition \ref{etale-cohomology-proposition-topological-invariance},
nous pouvons remplacer $X$ par sa réduction.
Soit $\nu : X^\nu \to X$ la normalisation.
C'est un morphisme fini birationnel surjectif
(voir le lemme \ref{varieties-lemma-normalization-locally-algebraic} de Variétés),
donc un isomorphisme au-dessus d'un ouvert dense $U \subset X$
(lemme \ref{morphisms-lemma-birational-birational} de Morphismes).
Nous voyons alors que
$c : \mathcal{F} \to \nu_*\nu^{-1}\mathcal{F}$ est injectif
(car $\nu$ est surjectif) et est un isomorphisme au-dessus de $U$.
Notons $i : Z \to X$ l'inclusion du complémentaire de $U$.
Puisque $U$ est dense dans $X$, nous avons $\dim(Z) < \dim(X) = n$.
Par la proposition \ref{etale-cohomology-proposition-closed-immersion-pushforward}, nous avons
$\Coker(c) = i_*\mathcal{G}$ pour un certain
faisceau abélien de torsion $\mathcal{G}$ sur $Z_\etale$.
Alors $H^q_\etale(X, \Coker(c)) = H^q_\etale(Z, \mathcal{G})$
(par la proposition \ref{etale-cohomology-proposition-finite-higher-direct-image-zero}
et la suite spectrale de Leray), et l'hypothèse de récurrence montre que
le conoyau de $c$ a sa cohomologie en degrés $\leq 2(n - 1)$.
Il suffit donc de démontrer le résultat pour $\nu_*\nu^{-1}\mathcal{F}$.
Comme $\nu$ est fini, cela nous ramène à montrer
que $H^i_\etale(X^\nu, \nu^{-1}\mathcal{F})$ est
nul pour $i > 2n$. Ce cas est traité au paragraphe suivant.

\medskip\noindent
Supposons que $X$ soit un schéma propre, intègre et normal sur $k$, de dimension $n$.
Choisissons une fonction rationnelle non constante $f$ sur $X$. Le graphe
$X' \subset X \times \mathbf{P}^1_k$ de $f$ s'insère dans un diagramme
$$
X \xleftarrow{b} X' \xrightarrow{f} \mathbf{P}^1_k
$$
Remarquons que $b$ est un isomorphisme au-dessus d'un sous-schéma ouvert
$U \subset X$ dont le complémentaire est un sous-schéma fermé
$Z \subset X$ de codimension $\geq 2$. En effet, $U$ est le
domaine de définition de $f$, qui contient tous les points de codimension $1$
de $X$; voir les lemmes de Morphismes
\ref{morphisms-lemma-rational-map-from-reduced-to-separated}
et \ref{morphisms-lemma-extend-across}
(combinés au critère de normalité de Serre; voir le
lemme \ref{properties-lemma-criterion-normal} de Properties).
De plus, les fibres de $b$ sont de dimension $\leq 1$ (comme sous-schémas fermés
de $\mathbf{P}^1$). Ainsi, par récurrence, $R^ib_*b^{-1}\mathcal{F}$ n'est non nul que
si $i \in \{0, 1, 2\}$. Choisissons un triangle distingué
$$
\mathcal{F} \to Rb_*b^{-1}\mathcal{F} \to Q \to \mathcal{F}[1]
$$
Comme précédemment, l'injectivité de $\mathcal{F} \to b_*b^{-1}\mathcal{F}$
et ce que nous venons de montrer impliquent que $Q$ a des
faisceaux de cohomologie non nuls seulement en degrés $0, 1, 2$, portés par $Z$.
De plus, ces faisceaux de cohomologie sont de torsion par le
lemme \ref{etale-cohomology-lemma-torsion-direct-image}.
Par récurrence et par la suite spectrale du
lemme \ref{derived-lemma-two-ss-complex-functor} de Catégories dérivées,
nous voyons que $H^i(X, Q)$ est nul pour
$i > 2 + 2\dim(Z) \leq 2 + 2(n - 2) = 2n - 2$. Il suffit donc
de démontrer que $H^i(X', b^{-1}\mathcal{F}) = 0$ pour
$i > 2n$. Nous utilisons alors le morphisme
$$
f : X' \to \mathbf{P}^1_k
$$
dont les fibres sont de dimension $< n$. Ainsi, par récurrence, nous voyons que
$R^if_*b^{-1}\mathcal{F} = 0$ pour $i > 2(n - 1)$.
Nous concluons par la suite spectrale de Leray
$$
H^i(\mathbf{P}^1_k, R^jf_*b^{-1}\mathcal{F})
\Rightarrow
H^{i + j}(X', b^{-1}\mathcal{F})
$$
et le fait que $\dim(\mathbf{P}^1_k) = 1$.
\end{proof}

\noindent
Avec des coefficients modulo $n$, nous pouvons appliquer le changement
de base propre aux complexes non bornés.

\begin{lemma}
\label{etale-cohomology-lemma-proper-base-change-mod-n}
Soit $f : X \to Y$ un morphisme propre de schémas. Soit $g : Y' \to Y$
un morphisme de schémas. Posons $X' = Y' \times_Y X$ et notons
$f' : X' \to Y'$ et $g' : X' \to X$ les projections.
Soit $n \geq 1$ un entier.
Soit $E \in D(X_\etale, \mathbf{Z}/n\mathbf{Z})$.
Alors le morphisme de changement de base (\ref{etale-cohomology-equation-base-change})
$g^{-1}Rf_*E \to Rf'_*(g')^{-1}E$
est un isomorphisme.
\end{lemma}

\begin{proof}
Il suffit de le démontrer lorsque $Y$ et $Y'$ sont quasi-compacts.
Par le lemme de Morphismes
\ref{morphisms-lemma-morphism-finite-type-bounded-dimension}
nous voyons que les dimensions des fibres de
$f : X \to Y$ et de $f' : X' \to Y'$ sont bornées. Ainsi, le
lemme \ref{etale-cohomology-lemma-cohomological-dimension-proper} implique que
$$
f_* :
\textit{Mod}(X_\etale, \mathbf{Z}/n\mathbf{Z})
\longrightarrow
\textit{Mod}(Y_\etale, \mathbf{Z}/n\mathbf{Z})
$$
et
$$
f'_* :
\textit{Mod}(X'_\etale, \mathbf{Z}/n\mathbf{Z})
\longrightarrow
\textit{Mod}(Y'_\etale, \mathbf{Z}/n\mathbf{Z})
$$
ont une dimension cohomologique finie au sens du
lemme \ref{derived-lemma-unbounded-right-derived} de Catégories dérivées.
Choisissons un complexe K-injectif $\mathcal{I}^\bullet$
de $\mathbf{Z}/n\mathbf{Z}$-modules dont chacun des
termes $\mathcal{I}^n$ est un faisceau injectif de
$\mathbf{Z}/n\mathbf{Z}$-modules; ce complexe représente $E$.
Voir le théorème d'Injectifs
\ref{injectives-theorem-K-injective-embedding-grothendieck}.
Par le théorème usuel de changement de base propre, on obtient
que $R^qf'_*(g')^{-1}\mathcal{I}^n = 0$ pour $q > 0$; voir le
théorème \ref{etale-cohomology-theorem-proper-base-change}.
Nous concluons donc, par le
lemme \ref{derived-lemma-unbounded-right-derived} de Catégories dérivées,
que nous pouvons calculer $Rf'_*(g')^{-1}E$ par le complexe
$f'_*(g')^{-1}\mathcal{I}^\bullet$. Une autre application
du théorème usuel de changement de base propre montre que
ce complexe est égal à $g^{-1}f_*\mathcal{I}^\bullet$, comme souhaité.
\end{proof}

\begin{lemma}
\label{etale-cohomology-lemma-pull-out-constant}
Soit $X$ un schéma quasi-compact et quasi-séparé.
Soient $E \in D^+(X_\etale)$ et $K \in D^+(\mathbf{Z})$.
Alors
$$
R\Gamma(X, E \otimes_\mathbf{Z}^\mathbf{L} \underline{K}) =
R\Gamma(X, E) \otimes_\mathbf{Z}^\mathbf{L} K
$$
\end{lemma}

\begin{proof}
Supposons $H^i(E) = 0$ pour $i \geq a$ et $H^j(K) = 0$ pour $j \geq b$.
Nous pouvons représenter $K$ par un complexe borné inférieurement
$K^\bullet$ de $\mathbf{Z}$-modules sans torsion.
(Choisissons un complexe K-plat $L^\bullet$ représentant $K$, puis
prenons $K^\bullet = \tau_{\geq b - 1}L^\bullet$. Cela fonctionne parce que
$\mathbf{Z}$ est de dimension globale $1$. Voir le
lemme \ref{more-algebra-lemma-last-one-flat} de Compléments d'algèbre.)
Nous pouvons représenter $E$ par un complexe borné inférieurement
$\mathcal{E}^\bullet$.
Alors $E \otimes_\mathbf{Z}^\mathbf{L} \underline{K}$ est
représenté par
$$
\text{Tot}(\mathcal{E}^\bullet \otimes_\mathbf{Z} \underline{K}^\bullet)
$$
En utilisant les triangles distingués
$$
\sigma_{\geq -b + n + 1}K^\bullet
\to K^\bullet \to
\sigma_{\leq -b + n}K^\bullet
$$
et l'annulation évidente
$$
H^n(X,
\text{Tot}(\mathcal{E}^\bullet \otimes_\mathbf{Z}
\sigma_{\geq -a + n + 1}\underline{K}^\bullet) = 0
$$
et
$$
H^n(R\Gamma(X, E)
\otimes_\mathbf{Z}^\mathbf{L} \sigma_{\geq -a + n + 1}K^\bullet) = 0
$$
nous nous ramenons au cas où $K^\bullet$ est un complexe borné de
$\mathbf{Z}$-modules plats. En répétant l'argument, nous nous ramenons au
cas où $K^\bullet$ est réduit à un unique $\mathbf{Z}$-module plat
placé en un certain degré. Ensuite, en utilisant les troncations bêtes
de $\mathcal{E}^\bullet$,
nous nous ramenons exactement de la même manière au cas où
$\mathcal{E}^\bullet$ est réduit à un unique faisceau abélien placé
en un certain degré. Il suffit donc de montrer que
$$
H^n(X, \mathcal{E} \otimes_\mathbf{Z} \underline{M}) =
H^n(X, \mathcal{E}) \otimes_\mathbf{Z} M
$$
Lorsque $M$ est un $\mathbf{Z}$-module plat et que $\mathcal{E}$ est
un faisceau abélien sur $X$,
on écrit $M$ comme une limite inductive filtrante de $\mathbf{Z}$-modules libres finis
(théorème de Lazard; voir le théorème \ref{algebra-theorem-lazard} d'Algèbre).
Par le théorème \ref{etale-cohomology-theorem-colimit}, cela nous ramène au cas d'un module libre fini
sur $\mathbf{Z}$, noté $M$, auquel cas le résultat est trivialement vrai.
\end{proof}

\begin{lemma}
\label{etale-cohomology-lemma-projection-formula-proper}
Soit $f : X \to Y$ un morphisme propre de schémas.
Soit $E \in D^+(X_\etale)$ à faisceaux de cohomologie de torsion.
Soit $K \in D^+(Y_\etale)$. Alors
$$
Rf_*E \otimes_\mathbf{Z}^\mathbf{L} K =
Rf_*(E \otimes_\mathbf{Z}^\mathbf{L} f^{-1}K)
$$
dans $D^+(Y_\etale)$.
\end{lemma}

\begin{proof}
Un morphisme canonique du membre de gauche vers celui de droite est fourni par la
section \ref{sites-cohomology-section-projection-formula} de Cohomologie sur les sites.
Nous allons vérifier l'égalité sur les fibres.
Rappelons que le calcul des produits tensoriels dérivés commute aux images inverses.
Voir le lemme de Cohomologie sur les sites
\ref{sites-cohomology-lemma-pullback-tensor-product}.
Nous avons donc
$$
(E \otimes_\mathbf{Z}^\mathbf{L} f^{-1}K)_{\overline{x}}
=
E_{\overline{x}} \otimes_\mathbf{Z}^\mathbf{L} K_{\overline{y}}
$$
où $\overline{y}$ est l'image de $\overline{x}$ dans $Y$.
Puisque $\mathbf{Z}$ est de dimension globale $1$, nous voyons que
la cohomologie de ce complexe est nulle en degrés $< - 1 + a + b$
si $H^i(E) = 0$ pour $i \geq a$ et $H^j(K) = 0$ pour $j \geq b$.
De plus, puisque $H^i(E)$ est un faisceau abélien
de torsion pour tout $i$, il en va de même des faisceaux de cohomologie
du complexe $E \otimes_\mathbf{Z}^\mathbf{L} K$.
En effet, nous avons
$$
(E \otimes_\mathbf{Z}^\mathbf{L} f^{-1}K)
\otimes_{\mathbf{Z}}^\mathbf{L} \mathbf{Q} =
(E \otimes_\mathbf{Z}^\mathbf{L} \mathbf{Q})
\otimes_{\mathbf{Q}}^\mathbf{L}
(f^{-1}K \otimes_{\mathbf{Z}}^\mathbf{L} \mathbf{Q})
$$
qui est nul dans la catégorie dérivée.
Le lemme \ref{etale-cohomology-lemma-proper-base-change-stalk} s'applique donc
aux deux membres, et il suffit de montrer
$$
R\Gamma(X_{\overline{y}},
E|_{X_{\overline{y}}} \otimes_\mathbf{Z}^\mathbf{L}
(X_{\overline{y}} \to \overline{y})^{-1}K_{\overline{y}}) =
R\Gamma(X_{\overline{y}},
E|_{X_{\overline{y}}}) \otimes_\mathbf{Z}^\mathbf{L} K_{\overline{y}}
$$
Cela résulte du lemme \ref{etale-cohomology-lemma-pull-out-constant}.
\end{proof}








\section{Acyclicité locale}
\label{etale-cohomology-section-local-acyclicity}

\noindent
Dans cette section, nous déduisons l'acyclicité locale des morphismes lisses
du théorème de changement de base lisse. Dans SGA 4 ou SGA 4.5, les auteurs
démontrent d'abord une version de l'acyclicité locale pour les morphismes lisses,
puis en déduisent le théorème de changement de base lisse.

\medskip\noindent
Nous utiliserons la formulation de l'acyclicité locale donnée par
Deligne \cite[définition 2.12, page 242]{SGA4.5}.
Soit $f : X \to S$ un morphisme de schémas.
Soit $\overline{x}$ un point géométrique de $X$, d'image
$\overline{s} = f(\overline{x})$ dans $S$.
Soit $\overline{t}$ un point géométrique de
$\Spec(\mathcal{O}^{sh}_{S, \overline{s}})$.
Nous obtenons un diagramme commutatif
$$
\xymatrix{
F_{\overline{x}, \overline{t}} =
\overline{t} \times_{\Spec(\mathcal{O}^{sh}_{S, \overline{s}})}
\Spec(\mathcal{O}^{sh}_{X, \overline{x}}) \ar[r] \ar[d] &
\Spec(\mathcal{O}^{sh}_{X, \overline{x}}) \ar[r] \ar[d] &
X \ar[d] \\
\overline{t} \ar[r] &
\Spec(\mathcal{O}^{sh}_{S, \overline{s}}) \ar[r] &
S
}
$$
Le schéma $F_{\overline{x}, \overline{t}}$ est appelé
une {\it variété de cycles évanescents de $f$ en $\overline{x}$}.
Soit $K$ un objet de $D(X_\etale)$. Pour tout morphisme de schémas
$g : Y\to X$, nous écrivons $R\Gamma(Y, K)$ au lieu de
$R\Gamma(Y_\etale, g_{small}^{-1}K)$. Puisque
$\mathcal{O}^{sh}_{X, \overline{x}}$ est strictement hensélien,
nous avons
$K_{\overline{x}} = R\Gamma(\Spec(\mathcal{O}^{sh}_{X, \overline{x}}), K)$.
Nous obtenons donc un morphisme canonique
\begin{equation}
\label{etale-cohomology-equation-alpha-K}
\alpha_{K, \overline{x}, \overline{t}} :
K_{\overline{x}}
\longrightarrow
R\Gamma(F_{\overline{x}, \overline{t}}, K)
\end{equation}
en prenant l'image inverse en cohomologie le long de
$F_{\overline{x}, \overline{t}} \to \Spec(\mathcal{O}^{sh}_{X, \overline{x}})$.

\begin{definition}
\label{etale-cohomology-definition-locally-acyclic}
\begin{reference}
\cite[définition 2.12, page 242]{SGA4.5} et
\cite[définition (1.3), page 54]{SGA4.5}
\end{reference}
Soit $f : X \to S$ un morphisme de schémas.
Soit $K$ un objet de $D(X_\etale)$.
\begin{enumerate}
\item
Soit $\overline{x}$ un point géométrique de $X$, d'image
$\overline{s} = f(\overline{x})$.
Nous dirons que $f$ est {\it localement acyclique en $\overline{x}$ relativement à $K$}
si, pour tout point géométrique $\overline{t}$ de
$\Spec(\mathcal{O}^{sh}_{S, \overline{s}})$, le
morphisme (\ref{etale-cohomology-equation-alpha-K}) est un isomorphisme\footnote{Nous ne
supposons pas que $\overline{t}$ soit un point géométrique algébrique de
$\Spec(\mathcal{O}^{sh}_{S, \overline{s}})$. Le
lemme \ref{etale-cohomology-lemma-smooth-base-change-separably-closed} permet souvent
de se ramener à ce cas.}.
\item Nous dirons que $f$ est {\it localement acyclique relativement à $K$}
si $f$ est localement acyclique en $\overline{x}$ relativement à $K$
pour tout point géométrique $\overline{x}$ de $X$.
\item Nous dirons que $f$ est {\it universellement localement acyclique relativement à $K$}
si, pour tout morphisme de schémas $S' \to S$, le changement de base $f' : X' \to S'$
est localement acyclique relativement à l'image inverse de $K$ sur $X'$.
\item Nous dirons que $f$ est {\it localement acyclique} si, pour tout point
géométrique $\overline{x}$ de $X$ et tout entier $n$ premier à la caractéristique
de $\kappa(\overline{x})$, le morphisme $f$ est localement acyclique
en $\overline{x}$ relativement au faisceau constant de valeur
$\mathbf{Z}/n\mathbf{Z}$.
\item Nous dirons que $f$ est {\it universellement localement acyclique} si,
pour tout morphisme de schémas $S' \to S$, le changement de base $f' : X' \to S'$
est localement acyclique.
\end{enumerate}
\end{definition}

\noindent
Soit $M$ un groupe abélien. Alors l'acyclicité locale de $f : X \to S$
relativement au faisceau constant $\underline{M}$ se ramène
à la condition
$$
H^q(F_{\overline{x}, \overline{t}}, \underline{M}) =
\left\{
\begin{matrix}
M & \text{si} & q = 0 \\
0 & \text{si} & q \not = 0
\end{matrix}
\right.
$$
pour tout point géométrique $\overline{x}$ de $X$ et tout point
géométrique $\overline{t}$ de $\Spec(\mathcal{O}^{sh}_{S, f(\overline{x})})$.
Nous voyons ainsi qu'être localement acyclique correspond à
l'annulation des groupes de cohomologie supérieure des fibres géométriques
$F_{\overline{x}, \overline{t}}$ des morphismes entre les
hensélisés stricts en $\overline{x}$ et $\overline{s}$.

\begin{proposition}
\label{etale-cohomology-proposition-smooth-locally-acyclic}
Soit $f : X \to S$ un morphisme lisse de schémas.
Alors $f$ est universellement localement acyclique.
\end{proposition}

\begin{proof}
Puisque le changement de base d'un morphisme lisse est lisse, il suffit
de montrer que les morphismes lisses sont localement acycliques.
Soit $\overline{x}$ un point géométrique de $X$, d'image
$\overline{s} = f(\overline{x})$. Soit
$\overline{t}$ un point géométrique de
$\Spec(\mathcal{O}^{sh}_{S, f(\overline{x})})$.
Puisque nous cherchons à démontrer une propriété du morphisme d'anneaux
$\mathcal{O}^{sh}_{S, \overline{s}} \to \mathcal{O}^{sh}_{X, \overline{x}}$
(voir la discussion qui suit la définition \ref{etale-cohomology-definition-locally-acyclic}),
nous pouvons remplacer $f : X \to S$ par le changement de base
$X \times_S \Spec(\mathcal{O}^{sh}_{S, \overline{s}}) \to
\Spec(\mathcal{O}^{sh}_{S, \overline{s}})$.
Nous pouvons donc supposer que $S$ est le spectre d'un anneau local
strictement hensélien et que $\overline{s}$ est au-dessus
du point fermé de $S$.

\medskip\noindent
Nous allons appliquer le lemme \ref{etale-cohomology-lemma-base-change-f-star-general-stalks}
au diagramme
$$
\xymatrix{
X \ar[d]_f & X_{\overline{t}} \ar[l]^h \ar[d]^e \\
S & \overline{t} \ar[l]_g
}
$$
et au faisceau $\mathcal{F} = \underline{M}$, où $M = \mathbf{Z}/n\mathbf{Z}$
pour un certain entier $n$ premier à la caractéristique du corps résiduel de
$\overline{x}$. Nous savons que le morphisme
$f^{-1}R^qg_*\mathcal{F} \to R^qh_*e^{-1}\mathcal{F}$
est un isomorphisme par changement de base lisse; voir le
théorème \ref{etale-cohomology-theorem-smooth-base-change} (l'hypothèse
sur la torsion est satisfaite par notre choix de $n$).
Le lemme \ref{etale-cohomology-lemma-base-change-f-star-general-stalks}
nous donne donc l'égalité médiane dans
$$
H^q(F_{\overline{x}, \overline{t}}, \underline{M}) =
H^q(\Spec(\mathcal{O}^{sh}_{X, \overline{x}}) \times_S \overline{t},
\underline{M})
=
H^q(\Spec(\mathcal{O}^{sh}_{S, \overline{s}}) \times_S \overline{t},
\underline{M}) =
H^q(\overline{t}, \underline{M})
$$
Pour les deux égalités extérieures, nous utilisons le fait que
$S = \Spec(\mathcal{O}^{sh}_{S, \overline{s}})$.
Puisque $\overline{t}$ est le spectre d'un corps séparablement
clos, nous concluons que
$$
H^q(F_{\overline{x}, \overline{t}}, \underline{M}) =
\left\{
\begin{matrix}
M & \text{si} & q = 0 \\
0 & \text{si} & q \not = 0
\end{matrix}
\right.
$$
ce qu'il fallait montrer (voir la discussion qui suit la
définition \ref{etale-cohomology-definition-locally-acyclic}).
\end{proof}

\begin{lemma}
\label{etale-cohomology-lemma-locally-acyclic-locally-constant}
Soit $f : X \to S$ un morphisme de schémas. Soit $\mathcal{F}$ un
faisceau abélien localement constant sur $X_\etale$ tel que, pour tout point
géométrique $\overline{x}$ de $X$, le groupe abélien
$\mathcal{F}_{\overline{x}}$ soit de torsion et que chacun de ses éléments soit
d'ordre premier à la caractéristique du
corps résiduel de $\overline{x}$. Si $f$ est localement acyclique, alors $f$ est
localement acyclique relativement à $\mathcal{F}$.
\end{lemma}

\begin{proof}
En effet, soit $\overline{x}$ un point géométrique de $X$.
Puisque $\mathcal{F}$ est localement constant, nous voyons que
la restriction de $\mathcal{F}$ à $\Spec(\mathcal{O}^{sh}_{X, \overline{x}})$
est isomorphe au faisceau constant $\underline{M}$, où
$M = \mathcal{F}_{\overline{x}}$. Par hypothèse, nous pouvons écrire
$M = \colim M_i$ comme une limite inductive filtrante de groupes abéliens finis
$M_i$ d'ordre premier à la caractéristique du corps résiduel de
$\overline{x}$. Considérons un point géométrique $\overline{t}$ de
$\Spec(\mathcal{O}^{sh}_{S, f(\overline{x})})$.
Puisque $F_{\overline{x}, \overline{t}}$ est affine, nous avons
$$
H^q(F_{\overline{x}, \overline{t}}, \underline{M}) =
\colim H^q(F_{\overline{x}, \overline{t}}, \underline{M_i})
$$
par le lemme \ref{etale-cohomology-lemma-colimit}.
Pour tout $i$, nous pouvons écrire $M_i = \bigoplus \mathbf{Z}/n_{i, j}\mathbf{Z}$
comme une somme directe finie, pour certains entiers $n_{i, j}$ premiers à la
caractéristique du corps résiduel de $\overline{x}$.
Puisque $f$ est localement acyclique, nous voyons que
$$
H^q(F_{\overline{x}, \overline{t}},
\underline{\mathbf{Z}/n_{i, j}\mathbf{Z}}) =
\left\{
\begin{matrix}
\mathbf{Z}/n_{i, j}\mathbf{Z} & \text{si} & q = 0 \\
0 & \text{si} & q \not = 0
\end{matrix}
\right.
$$
Voir la discussion qui suit la définition \ref{etale-cohomology-definition-locally-acyclic}.
En prenant les sommes directes puis la limite inductive, nous concluons que
$$
H^q(F_{\overline{x}, \overline{t}}, \underline{M}) =
\left\{
\begin{matrix}
M & \text{si} & q = 0 \\
0 & \text{si} & q \not = 0
\end{matrix}
\right.
$$
ce qui achève la démonstration.
\end{proof}

\begin{lemma}
\label{etale-cohomology-lemma-locally-acyclic-quasi-finite-base-change}
Soit
$$
\xymatrix{
X' \ar[r]_{g'} \ar[d]_{f'} & X \ar[d]^f \\
S' \ar[r]^g & S
}
$$
un diagramme cartésien de schémas. Soit $K$ un objet de $D(X_\etale)$.
Soit $\overline{x}'$ un point géométrique de $X'$, d'image $\overline{x}$
dans $X$. Si :
\begin{enumerate}
\item $f$ est localement acyclique en $\overline{x}$ relativement à $K$;
\item $g$ est localement quasi-fini, ou bien $S' = \lim S_i$
est une limite projective filtrante de schémas localement quasi-finis sur $S$
à morphismes de transition affines, ou bien $g : S' \to S$ est entier,
\end{enumerate}
alors $f'$ est localement acyclique en $\overline{x}'$ relativement à $(g')^{-1}K$.
\end{lemma}

\begin{proof}
Notons $\overline{s}'$ et $\overline{s}$ les images de
$\overline{x}'$ et $\overline{x}$ dans $S'$ et $S$.
Soit $\overline{t}'$ un point géométrique de
$\Spec(\mathcal{O}^{sh}_{S', \overline{s}'})$, et notons
$\overline{t}$ son image dans $\Spec(\mathcal{O}^{sh}_{S, \overline{s}})$.
Par le lemme d'Algèbre
\ref{algebra-lemma-base-change-strict-henselization-quasi-finite}
et nos hypothèses sur $g$, nous avons un isomorphisme
$$
\mathcal{O}^{sh}_{X, \overline{x}}
\otimes_{\mathcal{O}^{sh}_{S, \overline{s}}}
\mathcal{O}^{sh}_{S', \overline{s}'}
\longrightarrow
\mathcal{O}^{sh}_{X', \overline{x}'}
$$
Puisque, d'après nos conventions,
$\kappa(\overline{t}) = \kappa(\overline{t}')$,
nous concluons que
$$
F_{\overline{x}', \overline{t}'} =
\Spec\left(
\mathcal{O}^{sh}_{X', \overline{x}'}
\otimes_{\mathcal{O}^{sh}_{S', \overline{s}'}}
\kappa(\overline{t}')\right) =
\Spec\left(
\mathcal{O}^{sh}_{X, \overline{x}}
\otimes_{\mathcal{O}^{sh}_{S, \overline{s}}}
\kappa(\overline{t})\right) =
F_{\overline{x}, \overline{t}}
$$
Autrement dit, les variétés des cycles évanescents
de $f'$ en $\overline{x}'$ sont des exemples de variétés des cycles
évanescents de $f$ en $\overline{x}$. Le lemme découle
immédiatement de cela et des définitions.
\end{proof}










\section{Le morphisme de cospécialisation}
\label{etale-cohomology-section-cospecialization}

\noindent
Soit $f : X \to S$ un morphisme de schémas. Soit $\overline{x}$ un
point géométrique de $X$, d'image $\overline{s} = f(\overline{x})$ dans $S$.
Soit $\overline{t}$ un point géométrique de
$\Spec(\mathcal{O}^{sh}_{S, \overline{s}})$.
Soit $K \in D(X_\etale)$. Pour tout morphisme de schémas $g : Y \to X$,
nous écrivons $K|_Y$ au lieu de $g_{small}^{-1}K$ et $R\Gamma(Y, K)$
au lieu de $R\Gamma(Y_\etale, g_{small}^{-1}K)$.
Nous affirmons que si :
\begin{enumerate}
\item $K$ est borné inférieurement, c'est-à-dire $K \in D^+(X_\etale)$;
\item $f$ est localement acyclique relativement à $K$,
\end{enumerate}
alors il existe un {\it morphisme de cospécialisation}
$$
cosp :
R\Gamma(X_{\overline{t}}, K)
\longrightarrow
R\Gamma(X_{\overline{s}}, K)
$$
qui sera étroitement lié au morphisme de spécialisation
considéré dans la section \ref{etale-cohomology-section-specialization}, et en particulier
dans la remarque \ref{etale-cohomology-remark-specialization-map-and-fibres}.

\medskip\noindent
Pour construire ce morphisme, considérons les morphismes
$$
X_{\overline{t}}
\xrightarrow{h}
X \times_S \Spec(\mathcal{O}^{sh}_{S, \overline{s}})
\xleftarrow{i}
X_{\overline{s}}
$$
L'unité de l'adjonction entre $h^{-1}$ et $Rh_*$ donne un morphisme
$$
\beta_{K, \overline{s}, \overline{t}} :
K|_{X \times_S \Spec(\mathcal{O}^{sh}_{S, \overline{s}})}
\longrightarrow
Rh_*(K|_{X_{\overline{t}}})
$$
dans $D((X \times_S \Spec(\mathcal{O}^{sh}_{S, \overline{s}}))_\etale)$.
Le lemme \ref{etale-cohomology-lemma-beta-is-isomorphism} ci-dessous montre que l'image inverse
$i^{-1}\beta_{K, \overline{s}, \overline{t}}$
est un isomorphisme sous les hypothèses ci-dessus.
Nous pouvons donc définir le morphisme de cospécialisation comme la composée
\begin{align*}
R\Gamma(X_{\overline{t}}, K)
& =
R\Gamma(X \times_S \Spec(\mathcal{O}^{sh}_{S, \overline{s}}),
Rh_*(K|_{X_{\overline{t}}})) \\
&
\xrightarrow{i^{-1}}
R\Gamma(X_{\overline{s}}, i^{-1}Rh_*(K|_{X_{\overline{t}}})) \\
&
\xrightarrow{(i^{-1}\beta_{K, \overline{s}, \overline{t}})^{-1}}
R\Gamma(X_{\overline{s}},
i^{-1}(K|_{X \times_S \Spec(\mathcal{O}^{sh}_{S, \overline{s}})})) \\
& =
R\Gamma(X_{\overline{s}}, K)
\end{align*}

\begin{lemma}
\label{etale-cohomology-lemma-beta-is-isomorphism}
Le morphisme $i^{-1}\beta_{K, \overline{s}, \overline{t}}$ est un isomorphisme.
\end{lemma}

\begin{proof}
La construction des morphismes $h$, $i$ et $\beta_{K, \overline{s}, \overline{t}}$
ne dépend que du
changement de base de $X$ et $K$ à $\Spec(\mathcal{O}^{sh}_{S, \overline{s}})$.
Nous pouvons donc supposer que $S$ est un schéma strictement hensélien
de point fermé $\overline{s}$. Remarquons que l'acyclicité locale
de $f$ relativement à $K$ est préservée par ce changement de base (par exemple
par le lemme \ref{etale-cohomology-lemma-locally-acyclic-quasi-finite-base-change}, ou simplement
en comparant directement les anneaux strictement henséliens dans ce cas très particulier).

\medskip\noindent
Soit $\overline{x}$ un point géométrique de $X_{\overline{s}}$.
De manière équivalente, soit $\overline{x}$ un point géométrique dont l'image
par $f$ est $\overline{s}$. Calculons la fibre de
$i^{-1}\beta_{K, \overline{s}, \overline{t}}$
en $\overline{x}$. Tout d'abord, nous avons
$$
(i^{-1}\beta_{K, \overline{s}, \overline{t}})_{\overline{x}} =
(\beta_{K, \overline{s}, \overline{t}})_{\overline{x}}
$$
car l'image inverse préserve les fibres; voir le lemme \ref{etale-cohomology-lemma-stalk-pullback}.
Puisque nous sommes dans la situation $S = \Spec(\mathcal{O}^{sh}_{S, \overline{s}})$,
nous voyons que $h : X_{\overline{t}} \to X$ possède la propriété
$X_{\overline{t}} \times_X \Spec(\mathcal{O}^{sh}_{X, \overline{x}})
= F_{\overline{x}, \overline{t}}$. Nous voyons donc que
$$
(\beta_{K, \overline{s}, \overline{t}})_{\overline{x}} :
K_{\overline{x}}
\longrightarrow
Rh_*(K|_{X_{\overline{t}}})_{\overline{x}} =
R\Gamma(F_{\overline{x}, \overline{t}}, K)
$$
où l'égalité est donnée par le théorème \ref{etale-cohomology-theorem-higher-direct-images}.
Il s'ensuit que le morphisme $(\beta_{K, \overline{s}, \overline{t}})_{\overline{x}}$
n'est autre que le morphisme $\alpha_{K, \overline{x}, \overline{t}}$ employé
dans la définition \ref{etale-cohomology-definition-locally-acyclic}. Le résultat en découle, car nous pouvons
vérifier sur les fibres qu'un morphisme est un isomorphisme, par le
théorème \ref{etale-cohomology-theorem-exactness-stalks}.
\end{proof}

\noindent
Le morphisme de cospécialisation, lorsqu'il existe, est destiné à être l'inverse
du morphisme de spécialisation.

\begin{lemma}
\label{etale-cohomology-lemma-specialization-cospecialization}
Dans la situation ci-dessus, si de plus
$f$ est quasi-compact et quasi-séparé, alors le diagramme
$$
\xymatrix{
(Rf_*K)_{\overline{s}} \ar[r] \ar[d]_{sp} &
R\Gamma(X_{\overline{s}}, K) \\
(Rf_*K)_{\overline{t}} \ar[r] &
R\Gamma(X_{\overline{t}}, K) \ar[u]_{cosp}
}
$$
est commutatif.
\end{lemma}

\begin{proof}
Comme dans la démonstration du lemme \ref{etale-cohomology-lemma-beta-is-isomorphism},
nous pouvons remplacer $S$ par $\Spec(\mathcal{O}^{sh}_{S, \overline{s}})$.
Nos morphismes se simplifient alors en $h : X_{\overline{t}} \to X$,
$i : X_{\overline{s}} \to X$ et
$\beta_{K, \overline{s}, \overline{t}} : K \to Rh_*(K|_{X_{\overline{t}}})$.
En utilisant l'égalité $(Rf_*K)_{\overline{s}} = R\Gamma(X, K)$ donnée par le
théorème \ref{etale-cohomology-theorem-higher-direct-images},
la composée de $sp$ avec le morphisme de changement de base
$(Rf_*K)_{\overline{t}} \to R\Gamma(X_{\overline{t}}, K)$
n'est autre que l'image inverse en cohomologie le long de $h$.
C'est le même morphisme que
$$
R\Gamma(X, K) \xrightarrow{\beta_{K, \overline{s}, \overline{t}}}
R\Gamma(X, Rh_*(K|_{X_{\overline{t}}})) =
R\Gamma(X_{\overline{t}}, K)
$$
Or le morphisme $cosp$ inverse d'abord le signe $=$ dans cette
formule, puis prend l'image inverse le long de $i$, et enfin applique
l'inverse de $i^{-1}\beta_{K, \overline{s}, \overline{t}}$.
Nous obtenons donc la commutativité voulue.
\end{proof}

\begin{lemma}
\label{etale-cohomology-lemma-sp-isom-proper-torsion-loc-ac}
Soit $f : X \to S$ un morphisme de schémas. Soit $K \in D(X_\etale)$.
Supposons que :
\begin{enumerate}
\item $K$ est borné inférieurement, c'est-à-dire $K \in D^+(X_\etale)$;
\item $f$ est localement acyclique relativement à $K$;
\item $f$ est propre;
\item $K$ a des faisceaux de cohomologie de torsion.
\end{enumerate}
Alors, pour tout point géométrique $\overline{s}$ de $S$ et tout point
géométrique $\overline{t}$ de $\Spec(\mathcal{O}^{sh}_{S, \overline{s}})$,
le morphisme de spécialisation
$sp : (Rf_*K)_{\overline{s}} \to (Rf_*K)_{\overline{t}}$
et le morphisme de cospécialisation
$cosp : R\Gamma(X_{\overline{t}}, K) \to R\Gamma(X_{\overline{s}}, K)$
sont des isomorphismes.
\end{lemma}

\begin{proof}
Par le théorème de changement de base propre (sous la forme du
lemme \ref{etale-cohomology-lemma-proper-base-change-stalk}), nous avons
$(Rf_*K)_{\overline{s}} = R\Gamma(X_{\overline{s}}, K)$
et de même pour $\overline{t}$. La démonstration « correcte » consisterait
à montrer que l'argument du
lemme \ref{etale-cohomology-lemma-specialization-cospecialization}
établit que $sp$ et $cosp$ sont des isomorphismes inverses dans ce cas.
Nous montrerons plutôt directement que $cosp$ est un isomorphisme.
D'après la discussion ci-dessus, nous voyons que $cosp$ est un isomorphisme
si et seulement si l'image inverse par $i$
$$
R\Gamma(X \times_S \Spec(\mathcal{O}^{sh}_{S, \overline{s}}),
Rh_*(K|_{X_{\overline{t}}}))
\longrightarrow
R\Gamma(X_{\overline{s}}, i^{-1}Rh_*(K|_{X_{\overline{t}}}))
$$
est un isomorphisme dans $D^+(\textit{Ab})$. Cela résulte du théorème de
changement de base propre appliqué au morphisme propre
$f' : X \times_S \Spec(\mathcal{O}^{sh}_{S, \overline{s}}) \to
\Spec(\mathcal{O}^{sh}_{S, \overline{s}})$
et au morphisme $\overline{s} \to \Spec(\mathcal{O}^{sh}_{S, \overline{s}})$,
avec le complexe $K' = Rh_*(K|_{X_{\overline{t}}})$. Le complexe
$K'$ est borné inférieurement et a des faisceaux de cohomologie de torsion
par le lemme \ref{etale-cohomology-lemma-torsion-direct-image}.
Puisque $\Spec(\mathcal{O}^{sh}_{S, \overline{s}})$ est strictement hensélien,
avec $\overline{s}$ au-dessus du point fermé,
nous voyons que la source de la flèche affichée est égale à
$(Rf'_*K')_{\overline{s}}$, que la cible est égale à
$R\Gamma(X_{\overline{s}}, K')$ et que le morphisme affiché est un isomorphisme
par le lemme déjà utilisé
\ref{etale-cohomology-lemma-proper-base-change-stalk}.
Nous voyons donc que trois des quatre flèches du diagramme
du lemme \ref{etale-cohomology-lemma-specialization-cospecialization} sont des isomorphismes,
et nous concluons.
\end{proof}

\begin{lemma}
\label{etale-cohomology-lemma-sp-isom-proper-loc-cst-torsion}
Soit $f : X \to S$ un morphisme de schémas. Soit $\mathcal{F}$
un faisceau abélien sur $X_\etale$. Supposons que :
\begin{enumerate}
\item $f$ est lisse et propre;
\item $\mathcal{F}$ est localement constant;
\item $\mathcal{F}_{\overline{x}}$ est un groupe de torsion dont
tous les éléments sont d'ordre premier à la caractéristique résiduelle de
$\overline{x}$, pour tout point géométrique $\overline{x}$ de $X$.
\end{enumerate}
Alors, pour tout point géométrique $\overline{s}$ de $S$ et tout point
géométrique $\overline{t}$ de $\Spec(\mathcal{O}^{sh}_{S, \overline{s}})$,
le morphisme de spécialisation
$sp : (Rf_*\mathcal{F})_{\overline{s}} \to (Rf_*\mathcal{F})_{\overline{t}}$
est un isomorphisme.
\end{lemma}

\begin{proof}
Cela découle des lemmes \ref{etale-cohomology-lemma-sp-isom-proper-torsion-loc-ac}
et \ref{etale-cohomology-lemma-locally-acyclic-locally-constant}, ainsi que de la
proposition \ref{etale-cohomology-proposition-smooth-locally-acyclic}.
\end{proof}

















\section{Dimension cohomologique}
\label{etale-cohomology-section-cd}

\noindent
Nous pouvons déduire certaines bornes sur la dimension cohomologique des
schémas et sur la dimension cohomologique des corps en utilisant
les résultats de la section \ref{etale-cohomology-section-vanishing-torsion}
et une application, apparemment anodine, du théorème de changement
de base propre (dans la démonstration de la proposition \ref{etale-cohomology-proposition-cd-affine}).

\begin{definition}
\label{etale-cohomology-definition-cd}
Soit $X$ un schéma quasi-compact et quasi-séparé.
La {\it dimension cohomologique de $X$} est le plus petit
élément
$$
\text{cd}(X) \in \{0, 1, 2, \ldots\} \cup \{\infty\}
$$
tel que, pour tout faisceau abélien de torsion $\mathcal{F}$
sur $X_\etale$, nous ayons $H^i_\etale(X, \mathcal{F}) = 0$
pour $i > \text{cd}(X)$. Si $X = \Spec(A)$, nous l'appelons parfois
la dimension cohomologique de $A$.
\end{definition}

\noindent
Si le schéma est de caractéristique $p$, nous pouvons souvent
obtenir des bornes plus fines pour l'annulation de la cohomologie des
faisceaux de $p$-torsion. Nous traiterons cela ailleurs
(référence à ajouter ultérieurement).

\begin{lemma}
\label{etale-cohomology-lemma-cd-limit}
Soit $X = \lim X_i$ la limite projective d'un système filtrant de
schémas quasi-compacts et quasi-séparés dont les morphismes de
transition sont affines. Alors $\text{cd}(X) \leq \max \text{cd}(X_i)$.
\end{lemma}

\begin{proof}
Notons $f_i : X \to X_i$ les projections.
Soit $\mathcal{F}$ un faisceau abélien de torsion sur $X_\etale$.
Nous avons alors $\mathcal{F} = \lim f_i^{-1}f_{i, *}\mathcal{F}$
par le lemme \ref{etale-cohomology-lemma-linus-hamann}.
Ainsi, $H^q_\etale(X, \mathcal{F}) =
\colim H^q_\etale(X_i, f_{i, *}\mathcal{F})$
par le théorème \ref{etale-cohomology-theorem-colimit}.
Le lemme en découle.
\end{proof}

\begin{lemma}
\label{etale-cohomology-lemma-cd-curve-over-field}
Soit $K$ un corps. Soit $X$ un schéma affine de dimension $1$
et de type fini sur $K$. Alors
$\text{cd}(X) \leq 1 + \text{cd}(K)$.
\end{lemma}

\begin{proof}
Soit $\mathcal{F}$ un faisceau abélien de torsion sur $X_\etale$.
Considérons la suite spectrale de Leray du morphisme
$f : X \to \Spec(K)$. Nous obtenons
$$
E_2^{p, q} = H^p(\Spec(K), R^qf_*\mathcal{F})
$$
qui converge vers $H^{p + q}_\etale(X, \mathcal{F})$.
La fibre de $R^qf_*\mathcal{F}$ en un point géométrique
$\Spec(\overline{K}) \to \Spec(K)$ est la cohomologie de l'image
inverse de $\mathcal{F}$ sur $X_{\overline{K}}$.
Elle s'annule donc en degrés $\geq 2$ par le
théorème \ref{etale-cohomology-theorem-vanishing-affine-curves}.
\end{proof}

\begin{lemma}
\label{etale-cohomology-lemma-cd-field-extension}
Soit $L/K$ une extension de corps. Nous avons alors
$\text{cd}(L) \leq \text{cd}(K) + \text{trdeg}_K(L)$.
\end{lemma}

\begin{proof}
Si $\text{trdeg}_K(L) = \infty$, c'est clair.
Sinon, nous pouvons trouver une suite d'extensions
$L= L_r/L_{r - 1}/ \ldots /L_1/L_0 = K$ telle que
$\text{trdeg}_{L_i}(L_{i + 1}) = 1$ et $r = \text{trdeg}_K(L)$.
Il suffit donc de démontrer le lemme dans le cas $r = 1$.
Dans ce cas, nous pouvons écrire $L = \colim A_i$
comme limite inductive filtrante de ses sous-$K$-algèbres de type fini.
Par le lemme \ref{etale-cohomology-lemma-cd-limit}, il suffit de démontrer que
$\text{cd}(A_i) \leq 1 + \text{cd}(K)$. Cela découle
du lemme \ref{etale-cohomology-lemma-cd-curve-over-field}.
\end{proof}

\begin{lemma}
\label{etale-cohomology-lemma-strictly-henselian}
Soit $K$ un corps. Soit $X$ un schéma de type fini sur $K$.
Soit $x \in X$. Posons $a = \text{trdeg}_K(\kappa(x))$
et $d = \dim_x(X)$. Il existe alors un morphisme
$$
K(t_1, \ldots, t_a)^{sep} \longrightarrow \mathcal{O}_{X, x}^{sh}
$$
tel que
\begin{enumerate}
\item le corps résiduel de $\mathcal{O}_{X, x}^{sh}$ soit une extension purement inséparable
de $K(t_1, \ldots, t_a)^{sep}$,
\item $\mathcal{O}_{X, x}^{sh}$ soit une limite inductive filtrante de
$K(t_1, \ldots, t_a)^{sep}$-algèbres de type fini et de dimension $\leq d - a$.
\end{enumerate}
\end{lemma}

\begin{proof}
Nous pouvons supposer $X$ affine. Par normalisation de Noether, quitte à
rétrécir encore $X$, nous pouvons choisir un morphisme fini
$\pi : X \to \mathbf{A}^d_K$; voir le
lemme \ref{algebra-lemma-Noether-normalization-at-point} d'Algèbre.
Puisque $\kappa(x)$ est une extension finie du corps résiduel de $\pi(x)$,
ce corps résiduel est lui aussi de degré de transcendance $a$ sur $K$.
Nous pouvons donc trouver un morphisme fini $\pi' : \mathbf{A}^d_K \to \mathbf{A}^d_K$
tel que $\pi'(\pi(x))$ corresponde au point générique du sous-espace
linéaire $\mathbf{A}^a_K \subset \mathbf{A}^d_K$ défini en annulant
les $d - a$ dernières coordonnées. Par conséquent, le composé
$$
X \xrightarrow{\pi' \circ \pi} \mathbf{A}^d_K \xrightarrow{p} \mathbf{A}^a_K
$$
de $\pi' \circ \pi$ et de la projection $p$ sur les $a$ premières coordonnées
envoie $x$ sur le point générique $\eta \in \mathbf{A}^a_K$. Le morphisme induit
$$
K(t_1, \ldots, t_a)^{sep} =
\mathcal{O}_{\mathbf{A}^a_K, \eta}^{sh}
\longrightarrow \mathcal{O}_{X, x}^{sh}
$$
sur les anneaux locaux étales satisfait (1), car le corps résiduel
de $\mathcal{O}_{X, x}^{sh}$ est manifestement une extension algébrique du corps
séparablement clos $K(t_1, \ldots, t_a)^{sep}$. D'autre part, si $X = \Spec(B)$,
alors $\mathcal{O}_{X, x}^{sh} = \colim B_j$ est une limite inductive filtrante
de $B$-algèbres étales $B_j$. Remarquons que $B_j$ est quasi-finie sur
$K[t_1, \ldots, t_d]$, puisque $B$ est finie sur $K[t_1, \ldots, t_d]$.
Nous pouvons de même écrire
$K(t_1, \ldots, t_a)^{sep} = \colim A_i$ comme limite inductive filtrante
de $K[t_1, \ldots, t_a]$-algèbres étales. Pour tout $i$, nous pouvons
trouver un $j$ tel que
$A_i \to K(t_1, \ldots, t_a)^{sep} \to \mathcal{O}_{X, x}^{sh}$
se factorise par un morphisme $\psi_{i, j} : A_i \to B_j$. Alors $B_j$ est quasi-finie
sur $A_i[t_{a + 1}, \ldots, t_d]$. Par suite,
$$
B_{i, j} = B_j \otimes_{\psi_{i, j}, A_i} K(t_1, \ldots, t_a)^{sep}
$$
est de dimension $\leq d - a$, puisqu'elle est quasi-finie sur
$K(t_1, \ldots, t_a)^{sep}[t_{a + 1}, \ldots, t_d]$.
La démonstration de (2) est maintenant achevée, car $\mathcal{O}_{X, x}^{sh}$ est une
limite inductive filtrante\footnote{Soit $R$ un anneau.
Soit $A = \colim_{i \in I} A_i$
une limite inductive filtrante de $R$-algèbres de présentation finie.
Soit $B = \colim_{j \in J} B_j$ une limite inductive filtrante de $R$-algèbres.
Soit $A \to B$ un morphisme de $R$-algèbres.
Supposons que, pour tout $i \in I$, il existe $j \in J$ et un morphisme de $R$-algèbres
$\psi_{i, j} : A_i \to B_j$.
Posons $(i', j', \psi_{i', j'}) \geq (i, j, \psi_{i, j})$ si
$i' \geq i$, $j' \geq j$, et si $\psi_{i, j}$ et $\psi_{i', j'}$
sont compatibles. La collection des triplets forme alors un
ensemble filtrant et $B = \colim B_j \otimes_{\psi_{i, j}, A_i} A$.}
des algèbres $B_{i, j}$. Certains détails sont omis.
\end{proof}

\begin{proposition}
\label{etale-cohomology-proposition-cd-affine}
Soit $K$ un corps. Soit $X$ un schéma affine de type fini sur $K$.
Nous avons alors $\text{cd}(X) \leq \dim(X) + \text{cd}(K)$.
\end{proposition}

\begin{proof}
Nous allons le démontrer par récurrence sur $\dim(X)$.
Soit $\mathcal{F}$ un faisceau abélien de torsion sur $X_\etale$.

\medskip\noindent
Le cas $\dim(X) = 0$. Dans ce cas, le morphisme structural
$f : X \to \Spec(K)$ est fini. Nous voyons donc que $R^if_*\mathcal{F} = 0$
pour $i > 0$; voir la proposition \ref{etale-cohomology-proposition-finite-higher-direct-image-zero}.
Ainsi, $H^i_\etale(X, \mathcal{F}) = H^i_\etale(\Spec(K), f_*\mathcal{F})$
par la suite spectrale de Leray de $f$
(lemme \ref{sites-cohomology-lemma-Leray} de Cohomologie sur les sites),
et le résultat est clair.

\medskip\noindent
Le cas $\dim(X) = 1$. C'est le lemme \ref{etale-cohomology-lemma-cd-curve-over-field}.

\medskip\noindent
Supposons $d = \dim(X) > 1$ et la proposition vraie pour les schémas
affines de type fini et de dimension $< d$ sur des corps.
Par normalisation de Noether, voir par exemple le
lemme \ref{varieties-lemma-noether-normalization-affine} de Variétés,
il existe un morphisme fini $f : X \to \mathbf{A}^d_K$.
Rappelons que $R^if_*\mathcal{F} = 0$ pour $i > 0$ par la
proposition \ref{etale-cohomology-proposition-finite-higher-direct-image-zero}.
Par la suite spectrale de Leray de $f$
(lemme \ref{sites-cohomology-lemma-Leray} de Cohomologie sur les sites),
nous concluons qu'il suffit de démontrer le résultat pour $f_*\mathcal{F}$
sur $\mathbf{A}^d_K$.

\medskip\noindent
Interlude I. Soit $j : X \to Y$ une immersion ouverte de variétés lisses
de dimension $d$ sur $K$ (non nécessairement affines),
dont le complément est le support d'un diviseur de Cartier effectif $D$.
Les faisceaux $R^qj_*\mathcal{F}$ pour $q > 0$ sont supportés par $D$.
Nous affirmons que $(R^qj_*\mathcal{F})_{\overline{y}} = 0$
pour $a = \text{trdeg}_K(\kappa(y)) > d - q$. En effet, par le
théorème \ref{etale-cohomology-theorem-higher-direct-images}, nous avons
$$
(R^qj_*\mathcal{F})_{\overline{y}} =
H^q(\Spec(\mathcal{O}_{Y, y}^{sh}) \times_Y X, \mathcal{F})
$$
Choisissons une équation locale $f \in \mathfrak m_y \subset \mathcal{O}_{Y, y}$
de $D$. Nous avons alors
$$
\Spec(\mathcal{O}_{Y, y}^{sh}) \times_Y X =
\Spec(\mathcal{O}_{Y, y}^{sh}[1/f])
$$
À l'aide du lemme \ref{etale-cohomology-lemma-strictly-henselian}, nous obtenons un plongement
$$
K(t_1, \ldots, t_a)^{sep}(x) =
K(t_1, \ldots, t_a)^{sep}[x]_{(x)}[1/x]
\longrightarrow
\mathcal{O}_{Y, y}^{sh}[1/f]
$$
Puisque le degré de transcendance sur $K$ du corps des fractions de
$\mathcal{O}_{Y, y}^{sh}$ est $d$, nous voyons que $\mathcal{O}_{Y, y}^{sh}[1/f]$
est une limite inductive filtrante d'algèbres de type fini de dimension $(d - a - 1)$ sur
le corps $K(t_1, \ldots, t_a)^{sep}(x)$, qui est lui-même de
dimension cohomologique $1$ par le lemme \ref{etale-cohomology-lemma-cd-field-extension}. L'hypothèse
de récurrence et le lemme \ref{etale-cohomology-lemma-cd-limit}
donnent donc l'annulation voulue.

\medskip\noindent
Interlude II. Soit $Z$ une variété lisse sur $K$ de dimension $d - 1$.
Soit $E_a \subset Z$ l'ensemble des points $z \in Z$ tels que
$\text{trdeg}_K(\kappa(z)) \leq a$. Remarquons que $E_a$ est stable
par spécialisation; voir le
lemme \ref{varieties-lemma-dimension-locally-algebraic} de Variétés.
Supposons que $\mathcal{G}$ soit un faisceau abélien de torsion sur $Z$
dont le support est contenu dans $E_a$. Nous affirmons alors que
$H^b_\etale(Z, \mathcal{G}) = 0$ pour $b > a + \text{cd}(K)$.
En effet, nous pouvons écrire $\mathcal{G} = \colim \mathcal{G}_i$
avec $\mathcal{G}_i$ un faisceau abélien de torsion
supporté par un sous-schéma fermé $Z_i$ contenu dans $E_a$; voir le
lemme \ref{etale-cohomology-lemma-support-in-subset}.
L'hypothèse de récurrence implique alors
l'annulation voulue pour $\mathcal{G}_i$\footnote{Nous employons ici d'abord la
proposition \ref{etale-cohomology-proposition-closed-immersion-pushforward}
pour écrire $\mathcal{G}_i$ comme l'image directe d'un faisceau sur $Z_i$;
l'hypothèse de récurrence donne l'annulation de ce faisceau sur $Z_i$, et
la suite spectrale de Leray de $Z_i \to Z$ donne l'annulation
de $\mathcal{G}_i$.}. Enfin, nous
concluons par le théorème \ref{etale-cohomology-theorem-colimit}.

\medskip\noindent
Considérons le diagramme commutatif
$$
\xymatrix{
\mathbf{A}^d_K \ar[rd]_f \ar[rr]_-j & &
\mathbf{P}^1_K \times_K \mathbf{A}^{d - 1}_K \ar[ld]^g \\
& \mathbf{A}^{d - 1}_K
}
$$
Remarquons que $j$ est une immersion ouverte de variétés lisses de dimension $d$
dont le complément est un diviseur de Cartier effectif $D$. Nous
pouvons donc employer les résultats obtenus dans l'interlude I.
Nous allons étudier la suite spectrale de Leray relative
$$
E_2^{p, q} = R^pg_*R^qj_*\mathcal{F} \Rightarrow R^{p + q}f_*\mathcal{F}
$$
Comme $R^qj_*\mathcal{F}$ pour $q > 0$ est supporté par $D$ et que
$g|_D : D \to \mathbf{A}^{d - 1}_K$ est un isomorphisme, nous trouvons
$R^pg_*R^qj_*\mathcal{F} = 0$ pour $p > 0$ et $q > 0$. De plus, nous avons
$R^qj_*\mathcal{F} = 0$ pour $q > d$. D'autre part, $g$ est un
morphisme propre de dimension relative $1$. Ainsi, par le
lemme \ref{etale-cohomology-lemma-cohomological-dimension-proper},
nous voyons que $R^pg_*j_*\mathcal{F} = 0$ pour $p > 2$. La page $E_2$
de la suite spectrale a donc la forme suivante
$$
\begin{matrix}
g_*R^dj_*\mathcal{F} & 0 & 0 \\
\ldots & \ldots & \ldots \\
g_*R^2j_*\mathcal{F} & 0 & 0 \\
g_*R^1j_*\mathcal{F} & 0 & 0 \\
g_*j_*\mathcal{F} & R^1g_*j_*\mathcal{F} & R^2g_*j_*\mathcal{F}
\end{matrix}
$$
Nous concluons que $R^qf_*\mathcal{F} = g_*R^qj_*\mathcal{F}$ pour $q > 2$.
Par l'interlude I, nous voyons que le support de $R^qf_*\mathcal{F}$ pour $q > 2$
est contenu dans l'ensemble des points de $\mathbf{A}^{d - 1}_K$
dont le corps résiduel est de degré de transcendance $\leq d - q$.
Par l'interlude II,
$$
H^p(\mathbf{A}^{d - 1}_K, R^qf_*\mathcal{F}) = 0
\text{ pour }p > d - q + \text{cd}(K)\text{ et }q > 2
$$
 D'autre part, par le théorème \ref{etale-cohomology-theorem-higher-direct-images},
nous avons $R^2f_*\mathcal{F}_{\overline{\eta}} =
H^2(\mathbf{A}^1_{\overline{\eta}}, \mathcal{F}) = 0$
(annulation par le cas de dimension $1$), où $\eta$ est le
point générique de $\mathbf{A}^{d - 1}_K$.
Ainsi, de nouveau par l'interlude II, nous voyons que
$$
H^p(\mathbf{A}^{d - 1}_K, R^2f_*\mathcal{F}) = 0
\text{ pour }p > d - 2 + \text{cd}(K)
$$
Enfin, nous avons
$$
H^p(\mathbf{A}^{d - 1}_K, R^qf_*\mathcal{F}) = 0
\text{ pour }p > d - 1 + \text{cd}(K)\text{ et }q = 0, 1
$$
par l'hypothèse de récurrence. En combinant tout ce qui précède
avec la suite spectrale de Leray
$H^p(\mathbf{A}^{d - 1}_K, R^qf_*\mathcal{F}) \Rightarrow
H^{p + q}(\mathbf{A}^d_K, \mathcal{F})$, nous concluons.
\end{proof}

\begin{lemma}
\label{etale-cohomology-lemma-interlude-II}
Soit $K$ un corps. Soit $X$ un schéma affine de type fini sur $K$.
Soit $E_a \subset X$ l'ensemble des points
$x \in X$ tels que $\text{trdeg}_K(\kappa(x)) \leq a$.
Soit $\mathcal{F}$ un faisceau abélien de torsion sur $X_\etale$
dont le support est contenu dans $E_a$. Alors
$H^b_\etale(X, \mathcal{F}) = 0$ pour $b > a + \text{cd}(K)$.
\end{lemma}

\begin{proof}
Nous pouvons écrire $\mathcal{F} = \colim \mathcal{F}_i$
avec $\mathcal{F}_i$ un faisceau abélien de torsion
supporté par un sous-schéma fermé $Z_i$ contenu dans $E_a$; voir le
lemme \ref{etale-cohomology-lemma-support-in-subset}.
La proposition \ref{etale-cohomology-proposition-cd-affine} donne alors
l'annulation voulue pour $\mathcal{F}_i$. Nous omettons les détails;
indications : utiliser d'abord la proposition \ref{etale-cohomology-proposition-closed-immersion-pushforward}
pour écrire $\mathcal{F}_i$ comme l'image directe d'un faisceau sur $Z_i$,
employer l'annulation de ce faisceau sur $Z_i$, puis la
suite spectrale de Leray de $Z_i \to X$ pour obtenir l'annulation
de $\mathcal{F}_i$. Enfin, nous
concluons par le théorème \ref{etale-cohomology-theorem-colimit}.
\end{proof}

\begin{lemma}
\label{etale-cohomology-lemma-interlude-I}
Soit $f : X \to Y$ un morphisme affine de schémas de type
fini sur un corps $K$. Soit $E_a(X)$ l'ensemble des points $x \in X$
tels que $\text{trdeg}_K(\kappa(x)) \leq a$.
Soit $\mathcal{F}$ un faisceau abélien de torsion sur $X_\etale$
dont le support est contenu dans $E_a$. Alors
le support de $R^qf_*\mathcal{F}$ est contenu dans
$E_{a - q}(Y)$.
\end{lemma}

\begin{proof}
La question est locale sur $Y$; nous pouvons donc supposer $Y$ affine.
Alors $X$ est lui aussi affine, et nous pouvons choisir un diagramme
$$
\xymatrix{
X \ar[d]_f \ar[r]_i & \mathbf{A}^{n + m}_K \ar[d]^{\text{pr}} \\
Y \ar[r]^j & \mathbf{A}^n_K
}
$$
où les flèches horizontales sont des immersions fermées et la flèche
verticale de droite est la projection (détails omis).
Alors $j_*R^qf_*\mathcal{F} = R^q\text{pr}_*i_*\mathcal{F}$
par l'annulation des images directes supérieures de $i$ et $j$; voir la
proposition \ref{etale-cohomology-proposition-finite-higher-direct-image-zero}.
De plus, la description des fibres de $j_*$ dans cette proposition
montre qu'il suffit d'établir l'annulation pour $j_*R^qf_*\mathcal{F}$.
Nous pouvons donc supposer que $f$ est le morphisme de
projection $\text{pr} : \mathbf{A}^{n + m}_K \to \mathbf{A}^n_K$
et que $\mathcal{F}$ est un faisceau abélien de torsion sur $\mathbf{A}^{n + m}_K$
satisfaisant l'hypothèse de l'énoncé du lemme.

\medskip\noindent
Soit $y$ un point de $\mathbf{A}^n_K$.
Par le théorème \ref{etale-cohomology-theorem-higher-direct-images}, nous avons
$$
(R^q\text{pr}_*\mathcal{F})_{\overline{y}} =
H^q(\mathbf{A}^{n + m}_K \times_{\mathbf{A}^n_K} \Spec(\mathcal{O}_{Y, y}^{sh}),
\mathcal{F}) =
H^q(\mathbf{A}^m_{\mathcal{O}_{Y, y}^{sh}}, \mathcal{F})
$$
Posons $b = \text{trdeg}_K(\kappa(y))$. Le lemme \ref{etale-cohomology-lemma-strictly-henselian}
nous donne un plongement
$$
L = K(t_1, \ldots, t_b)^{sep} \longrightarrow \mathcal{O}_{Y, y}^{sh}
$$
Écrivons $\mathcal{O}_{Y, y}^{sh} = \colim B_i$ comme la limite
inductive filtrante de sous-$L$-algèbres de type fini $B_i \subset \mathcal{O}_{Y, y}^{sh}$
contenant l'anneau $K[T_1, \ldots, T_n]$ des fonctions régulières
sur $\mathbf{A}^n_K$. Nous obtenons alors
$$
\mathbf{A}^m_{\mathcal{O}_{Y, y}^{sh}} =
\lim \mathbf{A}^m_{B_i}
$$
Si $z \in \mathbf{A}^m_{B_i}$ est un point du support de
$\mathcal{F}$, alors l'image $x$ de $z$ dans $\mathbf{A}^{m + n}_K$
satisfait $\text{trdeg}_K(\kappa(x)) \leq a$ par l'hypothèse du lemme sur
$\mathcal{F}$.
Puisque $\mathcal{O}_{Y, y}^{sh}$ est une limite inductive filtrante
d'algèbres étales sur $K[T_1, \ldots, T_n]$ et que
$B_i \subset \mathcal{O}_{Y, y}^{sh}$
nous voyons que $\kappa(z)/\kappa(x)$ est algébrique
(certains détails sont omis). Alors $\text{trdeg}_K(\kappa(z)) \leq a$
et donc $\text{trdeg}_L(\kappa(z)) \leq a - b$.
Par le lemme \ref{etale-cohomology-lemma-interlude-II}, nous voyons que
$$
H^q(\mathbf{A}^m_{B_i}, \mathcal{F}) = 0\text{ pour }q > a - b
$$
Ainsi, par le théorème \ref{etale-cohomology-theorem-colimit}, nous obtenons
$(R^qf_*\mathcal{F})_{\overline{y}} = 0$ pour $q > a - b$,
comme voulu.
\end{proof}








\section{Dimension cohomologique finie}
\label{etale-cohomology-section-finite-cd}

\noindent
Nous poursuivons la discussion commencée dans la section \ref{etale-cohomology-section-cd}.

\begin{definition}
\label{etale-cohomology-definition-cd-f}
Soit $f : X \to Y$ un morphisme de schémas quasi-compact et
quasi-séparé.
La {\it dimension cohomologique de $f$} est le plus petit
élément
$$
\text{cd}(f) \in \{0, 1, 2, \ldots\} \cup \{\infty\}
$$
tel que, pour tout faisceau abélien de torsion $\mathcal{F}$
sur $X_\etale$, nous ayons $R^if_*\mathcal{F} = 0$
pour $i > \text{cd}(f)$.
\end{definition}

\begin{lemma}
\label{etale-cohomology-lemma-finite-cd}
Soit $K$ un corps.
\begin{enumerate}
\item Si $f : X \to Y$ est un morphisme de schémas de type fini sur $K$,
alors $\text{cd}(f) < \infty$.
\item Si $\text{cd}(K) < \infty$, alors $\text{cd}(X) < \infty$
pour tout schéma $X$ de type fini sur $K$.
\end{enumerate}
\end{lemma}

\begin{proof}
Démonstration de (1). Nous pouvons supposer $Y$ affine. Nous allons employer le
principe de récurrence du
lemme \ref{coherent-lemma-induction-principle} de Cohomologie des schémas
pour le démontrer. Si $X$ est lui aussi affine, le résultat vaut par le
lemme \ref{etale-cohomology-lemma-interlude-I}. Il suffit donc de montrer que, si
$X = U \cup V$ et si le résultat est vrai pour $U \to Y$, $V \to Y$ et
$U \cap V \to Y$, il est vrai pour $f$. Cela découle de la
suite de Mayer--Vietoris relative; voir le
lemme \ref{etale-cohomology-lemma-relative-mayer-vietoris}.

\medskip\noindent
Démonstration de (2). Nous allons employer le principe de récurrence du
lemme \ref{coherent-lemma-induction-principle} de Cohomologie des schémas
pour le démontrer. Si $X$ est affine, le résultat vaut par la
proposition \ref{etale-cohomology-proposition-cd-affine}. Il suffit donc de montrer que, si
$X = U \cup V$ et si le résultat est vrai pour $U$, $V$ et
$U \cap V $, il est vrai pour $X$. Cela découle de la
suite de Mayer--Vietoris; voir le
lemme \ref{etale-cohomology-lemma-mayer-vietoris}.
\end{proof}

\begin{lemma}
\label{etale-cohomology-lemma-finite-cd-mod-n-direct-sums}
Cohomologie et sommes directes. Soit $n \geq 1$ un entier.
\begin{enumerate}
\item Soit $f : X \to Y$ un morphisme de schémas quasi-compact et
quasi-séparé tel que $\text{cd}(f) < \infty$. Alors le foncteur
$$
Rf_* :
D(X_\etale, \mathbf{Z}/n\mathbf{Z})
\longrightarrow
D(Y_\etale, \mathbf{Z}/n\mathbf{Z})
$$
commute aux sommes directes.
\item Soit $X$ un schéma quasi-compact et quasi-séparé tel que
$\text{cd}(X) < \infty$. Alors le foncteur
$$
R\Gamma(X, -) :
D(X_\etale, \mathbf{Z}/n\mathbf{Z})
\longrightarrow
D(\mathbf{Z}/n\mathbf{Z})
$$
commute aux sommes directes.
\end{enumerate}
\end{lemma}

\begin{proof}
Démonstration de (1). Puisque $\text{cd}(f) < \infty$, nous voyons que
$$
f_* :
\textit{Mod}(X_\etale, \mathbf{Z}/n\mathbf{Z})
\longrightarrow
\textit{Mod}(Y_\etale, \mathbf{Z}/n\mathbf{Z})
$$
est de dimension cohomologique finie au sens du
lemme \ref{derived-lemma-unbounded-right-derived} de Catégories dérivées.
Soit $I$ un ensemble et, pour $i \in I$, soit $E_i$
un objet de $D(X_\etale, \mathbf{Z}/n\mathbf{Z})$.
Choisissons un complexe K-injectif $\mathcal{I}_i^\bullet$
de $\mathbf{Z}/n\mathbf{Z}$-modules dont chacun des
termes $\mathcal{I}_i^n$ est un faisceau injectif de
$\mathbf{Z}/n\mathbf{Z}$-modules et qui représente $E_i$.
Voir le théorème d'Injectifs
\ref{injectives-theorem-K-injective-embedding-grothendieck}.
Alors $\bigoplus E_i$ est représenté par le complexe
$\bigoplus \mathcal{I}_i^\bullet$ (somme directe terme à terme); voir le
lemme \ref{injectives-lemma-derived-products} d'Injectifs.
Par le lemme \ref{etale-cohomology-lemma-relative-colimit}, nous avons
$$
R^qf_*(\bigoplus \mathcal{I}_i^n) =
\bigoplus R^qf_*(\mathcal{I}_i^n) = 0
$$
pour $q > 0$ et tout $n$. Nous concluons donc par le
lemme \ref{derived-lemma-unbounded-right-derived} de Catégories dérivées
que nous pouvons calculer $Rf_*(\bigoplus E_i)$ par le complexe
$$
f_*(\bigoplus \mathcal{I}_i^\bullet) =
\bigoplus f_*(\mathcal{I}_i^\bullet)
$$
(égalité donnée de nouveau par le lemme \ref{etale-cohomology-lemma-relative-colimit}), qui
représente $\bigoplus Rf_*E_i$ par le lemme déjà employé
\ref{injectives-lemma-derived-products} d'Injectifs.

\medskip\noindent
Démonstration de (2). Elle est identique à celle de (1),
et nous l'omettons.
\end{proof}

\begin{lemma}
\label{etale-cohomology-lemma-proper-mod-n-direct-sums}
Soit $f : X \to Y$ un morphisme propre de schémas. Soit $n \geq 1$
un entier. Alors le foncteur
$$
Rf_* :
D(X_\etale, \mathbf{Z}/n\mathbf{Z})
\longrightarrow
D(Y_\etale, \mathbf{Z}/n\mathbf{Z})
$$
commute aux sommes directes.
\end{lemma}

\begin{proof}
Il suffit de le démontrer lorsque $Y$ est quasi-compact. Par le
lemme \ref{morphisms-lemma-morphism-finite-type-bounded-dimension} de Morphismes,
nous voyons que la dimension des fibres de
$f : X \to Y$ est bornée.
Le lemme \ref{etale-cohomology-lemma-cohomological-dimension-proper}
implique donc que $\text{cd}(f) < \infty$. Le résultat
découle alors du lemme \ref{etale-cohomology-lemma-finite-cd-mod-n-direct-sums}.
\end{proof}

\begin{lemma}
\label{etale-cohomology-lemma-pull-out-constant-mod-n}
Soit $X$ un schéma quasi-compact et quasi-séparé
tel que $\text{cd}(X) < \infty$. Soit $\Lambda$ un anneau de torsion.
Soient $E \in D(X_\etale, \Lambda)$ et $K \in D(\Lambda)$. Alors
$$
R\Gamma(X, E \otimes_\Lambda^\mathbf{L} \underline{K}) =
R\Gamma(X, E) \otimes_\Lambda^\mathbf{L} K
$$
\end{lemma}

\begin{proof}
Il existe un morphisme canonique de droite à gauche par la
section \ref{sites-cohomology-section-projection-formula} de Cohomologie sur les sites.
Soit $T(K)$ la propriété selon laquelle l'énoncé du
lemme vaut pour $K \in D(\Lambda)$.
Nous allons vérifier que les conditions (1), (2) et (3) de la
remarque \ref{more-algebra-remark-P-resolution} de Compléments d'algèbre
valent pour $T$, ce qui permettra de conclure.
La propriété (1) vaut parce que les deux membres de l'égalité
commutent aux sommes directes; voir le
lemme \ref{etale-cohomology-lemma-finite-cd-mod-n-direct-sums}.
La propriété (2) vaut parce que nous comparons des foncteurs exacts
entre catégories triangulées et que nous pouvons employer le
lemme \ref{derived-lemma-third-isomorphism-triangle} de Catégories dérivées.
La propriété (3) affirme que le lemme vaut lorsque $K = \Lambda[k]$
pour tout décalage $k \in \mathbf{Z}$, ce qui est évident.
\end{proof}

\begin{lemma}
\label{etale-cohomology-lemma-projection-formula-proper-mod-n}
Soit $f : X \to Y$ un morphisme propre de schémas. Soit $\Lambda$
un anneau de torsion. Soient $E \in D(X_\etale, \Lambda)$ et
$K \in D(Y_\etale, \Lambda)$. Alors
$$
Rf_*E \otimes_\Lambda^\mathbf{L} K =
Rf_*(E \otimes_\Lambda^\mathbf{L} f^{-1}K)
$$
dans $D(Y_\etale, \Lambda)$.
\end{lemma}

\begin{proof}
Il existe un morphisme canonique de gauche à droite par la
section \ref{sites-cohomology-section-projection-formula} de Cohomologie sur les sites.
Nous allons vérifier l'égalité sur les fibres en $\overline{y}$.
Par le changement de base propre (sous la forme du
lemme \ref{etale-cohomology-lemma-proper-base-change-mod-n}, où $Y' = \overline{y}$),
on se ramène au cas où $Y$ est le spectre d'un
corps algébriquement clos.
C'est le lemme \ref{etale-cohomology-lemma-pull-out-constant-mod-n},
où nous employons $\text{cd}(X) < \infty$, donné par le
lemme \ref{etale-cohomology-lemma-cohomological-dimension-proper}.
\end{proof}






\section{Formule de K\"unneth en cohomologie étale}
\label{etale-cohomology-section-kunneth}

\noindent
Nous démontrons d'abord une formule de K\"unneth lorsque l'un des facteurs
est propre. Nous l'employons ensuite pour démontrer une propriété de changement de base
pour les immersions ouvertes. Celle-ci donne alors un
« changement de base par des morphismes vers des spectres de corps »
(analogue au changement de base lisse).
Enfin, nous l'utilisons pour obtenir une formule de K\"unneth plus générale.

\begin{remark}
\label{etale-cohomology-remark-define-kunneth-map}
Considérons un diagramme cartésien dans la catégorie des schémas :
$$
\xymatrix{
X \times_S Y \ar[d]_p \ar[r]_q \ar[rd]_c & Y \ar[d]^g \\
X \ar[r]^f & S
}
$$
Soit $\Lambda$ un anneau, et soient $E \in D(X_\etale, \Lambda)$
et $K \in D(Y_\etale, \Lambda)$. Il existe alors un morphisme canonique
$$
Rf_*E \otimes_\Lambda^\mathbf{L} Rg_*K
\longrightarrow
Rc_*(p^{-1}E \otimes_\Lambda^\mathbf{L} q^{-1}K)
$$
Nous pouvons par exemple le définir au moyen des morphismes canoniques
$Rf_*E \to Rc_*p^{-1}E$ et $Rg_*K \to Rc_*q^{-1}K$, ainsi que du
cup-produit relatif défini dans la remarque de Cohomologie sur les sites
\ref{sites-cohomology-remark-cup-product}.
On peut aussi employer l'adjoint du morphisme
$$
c^{-1}(Rf_*E \otimes_\Lambda^\mathbf{L} Rg_*K)
=
p^{-1}f^{-1}Rf_*E \otimes_\Lambda^\mathbf{L} q^{-1} g^{-1}Rg_*K
\to
p^{-1}E \otimes_\Lambda^\mathbf{L} q^{-1}K
$$
qui emploie les morphismes d'adjonction $f^{-1}Rf_*E \to E$ et
$g^{-1}Rg_*K \to K$.
\end{remark}


\begin{lemma}
\label{etale-cohomology-lemma-kunneth-one-proper}
Soit $k$ un corps séparablement clos. Soit $X$ un schéma propre sur $k$.
Soit $Y$ un schéma quasi-compact et quasi-séparé sur $k$.
\begin{enumerate}
\item Si $E \in D^+(X_\etale)$ a des faisceaux de cohomologie de torsion et
$K \in D^+(Y_\etale)$, alors
$$
R\Gamma(X \times_{\Spec(k)} Y,
\text{pr}_1^{-1}E
\otimes_\mathbf{Z}^\mathbf{L}
\text{pr}_2^{-1}K
)
=
R\Gamma(X, E)
\otimes_\mathbf{Z}^\mathbf{L}
R\Gamma(Y, K)
$$
\item Si $n \geq 1$ est un entier, si $Y$ est de type fini sur $k$,
si $E \in D(X_\etale, \mathbf{Z}/n\mathbf{Z})$ et
$K \in D(Y_\etale, \mathbf{Z}/n\mathbf{Z})$, alors
$$
R\Gamma(X \times_{\Spec(k)} Y,
\text{pr}_1^{-1}E
\otimes_{\mathbf{Z}/n\mathbf{Z}}^\mathbf{L}
\text{pr}_2^{-1}K
)
=
R\Gamma(X, E)
\otimes_{\mathbf{Z}/n\mathbf{Z}}^\mathbf{L}
R\Gamma(Y, K)
$$
\end{enumerate}
\end{lemma}

\begin{proof}
Démonstration de (1). Par le lemme \ref{etale-cohomology-lemma-projection-formula-proper}, nous avons
$$
R\text{pr}_{2, *}(
\text{pr}_1^{-1}E
\otimes_\mathbf{Z}^\mathbf{L}
\text{pr}_2^{-1}K) =
R\text{pr}_{2, *}(\text{pr}_1^{-1}E)
\otimes_\mathbf{Z}^\mathbf{L}
K
$$
Par changement de base propre (sous la forme du lemme \ref{etale-cohomology-lemma-proper-base-change}),
cet objet est égal à
$$
\underline{R\Gamma(X, E)}
\otimes_\mathbf{Z}^\mathbf{L}
K
$$
de $D(Y_\etale)$. L'application de $R\Gamma(Y, -)$ à cet objet redonne le
membre de gauche de l'égalité de (1) par la suite spectrale de Leray
de $\text{pr}_2$. Nous concluons donc par le
lemme \ref{etale-cohomology-lemma-pull-out-constant}.

\medskip\noindent
Démonstration de (2). Elle est exactement la même que celle de (1),
sauf que nous employons les lemmes \ref{etale-cohomology-lemma-projection-formula-proper-mod-n},
\ref{etale-cohomology-lemma-proper-base-change-mod-n} et
\ref{etale-cohomology-lemma-pull-out-constant-mod-n}, ainsi que
$\text{cd}(Y) < \infty$ par le lemme \ref{etale-cohomology-lemma-finite-cd}.
\end{proof}

\begin{lemma}
\label{etale-cohomology-lemma-supported-in-closed-points}
Soit $K$ un corps séparablement clos. Soit $X$ un schéma de type
fini sur $K$. Soit $\mathcal{F}$ un faisceau abélien sur $X_\etale$
dont le support est contenu dans l'ensemble des points fermés de $X$.
Alors $H^q(X, \mathcal{F}) = 0$ pour $q > 0$ et $\mathcal{F}$
est engendré par ses sections globales.
\end{lemma}

\begin{proof}
(Si $\mathcal{F}$ est de torsion, l'annulation découle immédiatement
du lemme \ref{etale-cohomology-lemma-interlude-II}.)
Par le lemme \ref{etale-cohomology-lemma-support-in-subset}, nous pouvons écrire $\mathcal{F}$
comme limite inductive filtrante de faisceaux constructibles $\mathcal{F}_i$
de $\mathbf{Z}$-modules dont les supports $Z_i \subset X$ sont des ensembles
finis de points fermés. Par la
proposition \ref{etale-cohomology-proposition-closed-immersion-pushforward},
un tel faisceau est de la forme $(Z_i \to X)_*\mathcal{G}_i$,
où $\mathcal{G}_i$ est un faisceau sur $Z_i$. Comme $K$ est séparablement clos,
le schéma $Z_i$ est une réunion disjointe finie de spectres de corps
séparablement clos. Rappelons que $H^q(Z_i, \mathcal{G}_i) = H^q(X, \mathcal{F}_i)$
par la suite spectrale de Leray de $Z_i \to X$ et l'annulation des images
directes supérieures de ce morphisme
(proposition \ref{etale-cohomology-proposition-finite-higher-direct-image-zero}).
Par les lemmes \ref{etale-cohomology-lemma-equivalence-abelian-sheaves-point} et
\ref{etale-cohomology-lemma-compare-cohomology-point}
nous voyons que $H^q(Z_i, \mathcal{G}_i)$ est nul pour $q > 0$
et que $H^0(Z_i, \mathcal{G}_i)$ engendre $\mathcal{G}_i$.
Nous en déduisons l'annulation de $H^q(X, \mathcal{F}_i)$ pour $q > 0$
et le fait que $\mathcal{F}_i$ est engendré par ses sections globales.
Par le théorème \ref{etale-cohomology-theorem-colimit}, nous voyons que
$H^q(X, \mathcal{F}) = 0$ pour $q > 0$.
La démonstration est achevée, car une
limite inductive filtrante de faisceaux de groupes abéliens engendrés par leurs sections globales
est engendrée par ses sections globales (détails omis).
\end{proof}

\begin{lemma}
\label{etale-cohomology-lemma-vanishing-closed-points}
Soit $K$ un corps séparablement clos. Soit $X$ un schéma de type
fini sur $K$. Soit $Q \in D(X_\etale)$. Supposons que $Q_{\overline{x}}$
ne soit non nul que si $x$ est un point fermé de $X$. Alors
$$
Q = 0 \Leftrightarrow H^i(X, Q) = 0 \text{ pour tout }i
$$
\end{lemma}

\begin{proof}
L'implication de gauche à droite est triviale. Nous devons donc
démontrer l'implication réciproque.

\medskip\noindent
Supposons $Q$ borné inférieurement; ce cas suffit pour presque toutes les
applications. Si $Q$ n'est pas nul, considérons le plus petit $i$
tel que le faisceau de cohomologie $H^i(Q)$ soit non nul.
Par le lemme \ref{etale-cohomology-lemma-supported-in-closed-points}, nous avons
$H^i(X, Q) = H^0(X, H^i(Q)) \not = 0$, et nous concluons.

\medskip\noindent
Cas général. Soit $\mathcal{B} \subset \Ob(X_\etale)$ la sous-catégorie des
objets quasi-compacts. Par le lemme \ref{etale-cohomology-lemma-supported-in-closed-points},
les hypothèses du lemme de Cohomologie sur les sites
\ref{sites-cohomology-lemma-cohomology-over-U-trivial}
sont satisfaites. Nous concluons que $H^q(U, Q) = H^0(U, H^q(Q))$
pour tout $U \in \mathcal{B}$. En particulier, cela vaut pour $U = X$.
La conclusion découle donc du lemme \ref{etale-cohomology-lemma-supported-in-closed-points},
puisque $Q$ est nul dans $D(X_\etale)$ si et seulement si $H^q(Q)$ est nul
pour tout $q$.
\end{proof}

\begin{lemma}
\label{etale-cohomology-lemma-kunneth-localize-on-X}
Soit $K$ un corps. Soit $j : U \to X$ une immersion ouverte de
schémas de type fini sur $K$. Soit $Y$ un schéma de type fini
sur $K$. Considérons le diagramme
$$
\xymatrix{
Y \times_{\Spec(K)} X \ar[d]_q  &
Y \times_{\Spec(K)} U \ar[l]^h \ar[d]^p \\
X & U \ar[l]_j
}
$$
Alors le morphisme de changement de base $q^{-1}Rj_*\mathcal{F} \to Rh_*p^{-1}\mathcal{F}$
est un isomorphisme pour tout faisceau abélien $\mathcal{F}$ sur $U_\etale$
dont les fibres sont de torsion d'ordres inversibles dans $K$.
\end{lemma}

\begin{proof}
Écrivons $\mathcal{F} = \colim \mathcal{F}[n]$, où la limite porte
sur le système multiplicatif des entiers inversibles dans $K$.
Comme la cohomologie commute aux limites inductives filtrantes dans notre situation
(pour un renvoi précis, voir le lemme \ref{etale-cohomology-lemma-base-change-Rf-star-colim}),
il suffit de démontrer le lemme pour $\mathcal{F}[n]$. Nous pouvons donc
supposer que $\mathcal{F}$ est un faisceau de $\mathbf{Z}/n\mathbf{Z}$-modules
pour un entier $n$ inversible dans $K$ (nous ne l'emploierons qu'à la toute fin de
la démonstration). Dans la démonstration, nous notons brièvement $X \times_K Y$ le produit
fibré au-dessus de $\Spec(K)$.
Nous allons démontrer le lemme par récurrence sur $\dim(X) + \dim(Y)$.
Le lemme est trivial si $\dim(X) \leq 0$, car, dans ce cas,
$U$ est un sous-schéma ouvert et fermé de $X$.
Choisissons un point $z \in X \times_K Y$. Nous allons montrer
que le morphisme sur la fibre en $\overline{z}$ est un isomorphisme.

\medskip\noindent
Supposons que $z \mapsto x \in X$ et que $\text{trdeg}_K(\kappa(x)) > 0$.
Posons $X' = \Spec(\mathcal{O}_{X, x}^{sh})$
et notons $U' \subset X'$ l'image inverse de $U$.
Considérons le changement de base
$$
\xymatrix{
Y \times_K X' \ar[d]_{q'}  &
Y \times_K U' \ar[l]^{h'} \ar[d]^{p'} \\
X' & U' \ar[l]_{j'}
}
$$
de notre diagramme par $X' \to X$. Remarquons que $X' \to X$ est une limite
inductive filtrante de morphismes étales. Par changement de base lisse, sous la forme du
lemme \ref{etale-cohomology-lemma-smooth-base-change-general}, l'image inverse de
$q^{-1}Rj_*\mathcal{F} \to Rh_*p^{-1}\mathcal{F}$ sur $X'$, puis sur
$Y \times_K X'$, est le morphisme
$(q')^{-1}Rj'_*\mathcal{F}' \to Rh'_*(p')^{-1}\mathcal{F}'$, où
$\mathcal{F}'$ est l'image inverse de $\mathcal{F}$ sur $U'$.
(Il suffirait ici d'employer le changement de base étale, qui est
un résultat essentiellement trivial.) Il suffit donc de montrer
que $(q')^{-1}Rj'_*\mathcal{F}' \to Rh'_*(p')^{-1}\mathcal{F}'$
est un isomorphisme pour prouver que notre morphisme initial
est un isomorphisme sur les fibres en $\overline{z}$.
Par le lemme \ref{etale-cohomology-lemma-strictly-henselian},
il existe un corps séparablement clos $L/K$ tel que
$X' = \lim X_i$, où $X_i$ est affine de type fini sur $L$
et $\dim(X_i) < \dim(X)$. Pour $i$ assez grand, il
existe un ouvert $U_i \subset X_i$ dont l'image inverse dans $X'$ est $U'$.
Nous pouvons appliquer l'hypothèse de récurrence au diagramme
$$
\vcenter{
\xymatrix{
Y \times_K X_i \ar[d]_{q_i}  &
Y \times_K U_i \ar[l]^{h_i} \ar[d]^{p_i} \\
X_i & U_i \ar[l]_{j_i}
}
}
\quad\text{égal à}\quad
\vcenter{
\xymatrix{
Y_L \times_L X_i \ar[d]_{q_i}  &
Y_L \times_L U_i \ar[l]^{h_i} \ar[d]^{p_i} \\
X_i & U_i \ar[l]_{j_i}
}
}
$$
sur le corps $L$ et à l'image inverse de $\mathcal{F}$ sur ces diagrammes.
Par le lemme \ref{etale-cohomology-lemma-base-change-Rf-star-colim}, nous concluons que le morphisme
$(q')^{-1}Rj'_*\mathcal{F}' \to Rh'_*(p')^{-1}\mathcal{F}'$ est un isomorphisme.

\medskip\noindent
Supposons que $z \mapsto y \in Y$ et que $\text{trdeg}_K(\kappa(y)) > 0$.
Soit $Y' = \Spec(\mathcal{O}_{Y, y}^{sh})$.
Par le lemme \ref{etale-cohomology-lemma-strictly-henselian}, il existe un corps séparablement clos
$L/K$ tel que $Y' = \lim Y_i$, où $Y_i$ est affine de type fini sur $L$
et $\dim(Y_i) < \dim(Y)$. En particulier, $Y'$ est un schéma sur $L$.
Notons par un indice $L$ le changement de base des schémas sur $K$
aux schémas sur $L$. Considérons les diagrammes commutatifs
$$
\vcenter{
\xymatrix{
Y' \times_K X \ar[d]_f  &
Y' \times_K U \ar[l]^{h'} \ar[d]^{f'} \\
Y \times_K X \ar[d]_q  &
Y \times_K U \ar[l]^h \ar[d]^p \\
X & U \ar[l]_j
}
}
\quad\text{et}\quad
\vcenter{
\xymatrix{
Y' \times_L X_L \ar[d]_{q'}  &
Y' \times_L U_L \ar[l]^{h'} \ar[d]^{p'} \\
X_L \ar[d] &
U_L \ar[l]^{j_L} \ar[d] \\
X & U \ar[l]_j
}
}
$$
et remarquons que les lignes supérieure et inférieure sont les mêmes à gauche et à droite.
Par changement de base lisse, nous voyons que
$f^{-1}Rh_*p^{-1}\mathcal{F} = Rh'_*(f')^{-1}p^{-1}\mathcal{F}$
(comme au paragraphe précédent).
Par changement de base lisse pour $\Spec(L) \to \Spec(K)$
(lemme \ref{etale-cohomology-lemma-base-change-field-extension}),
nous voyons que $Rj_{L, *}\mathcal{F}_L$ est l'image inverse de
$Rj_*\mathcal{F}$ sur $X_L$. En combinant ces deux observations,
nous concluons qu'il suffit de montrer que le morphisme de changement de base
du carré supérieur du diagramme de droite
est un isomorphisme pour prouver que notre morphisme initial
est un isomorphisme sur les fibres en $\overline{z}$\footnote{Nous employons ici
le fait qu'une « composée verticale » de morphismes de changement de base est un morphisme de
changement de base, comme l'explique la remarque de Cohomologie sur les sites
\ref{sites-cohomology-remark-compose-base-change}.}.
Puis, en employant $Y' = \lim Y_i$ et en raisonnant exactement
comme au paragraphe précédent, nous voyons que l'hypothèse de
récurrence force notre morphisme sur $Y' \times_K X$ à être
un isomorphisme.

\medskip\noindent
Ainsi, tout contre-exemple pour lequel $\dim(X) + \dim(Y)$ est minimal
ne pourrait avoir de défaut d'isomorphisme de
$q^{-1}Rj_*\mathcal{F} \to Rh_*p^{-1}\mathcal{F}$
que sur les fibres en des points fermés de $X \times_K Y$ (car un point
$z$ de $X \times_K Y$ est fermé si et seulement si
ses images dans $X$ et dans $Y$ sont toutes deux fermées).
Comme il suffit de démontrer l'isomorphisme localement,
nous pouvons supposer $X$ et $Y$ affines. Nous pouvons alors
choisir une immersion ouverte dense $Y \to Y'$ avec $Y'$
projectif. (Choisir une immersion fermée $Y \to \mathbf{A}^n_K$
et prendre pour $Y'$ l'adhérence schématique de $Y$ dans
$\mathbf{P}^n_K$.) Alors $\dim(Y') = \dim(Y)$, et nous
obtenons donc un contre-exemple « minimal » avec $Y$ projectif sur $K$.
Au paragraphe suivant, nous montrons que c'est impossible.

\medskip\noindent
Considérons un diagramme comme dans l'énoncé du lemme, tel que
$q^{-1}Rj_*\mathcal{F} \to Rh_*p^{-1}\mathcal{F}$
soit un isomorphisme en tout point non fermé de $X \times_K Y$
et que $Y$ soit projectif.
La restriction du morphisme à $(X \times_K Y)_{K^{sep}}$
est le morphisme correspondant pour le diagramme du lemme
après changement de base à $K^{sep}$. Nous pouvons donc supposer, et supposons, $K$
séparablement clos. Choisissons un triangle distingué
$$
q^{-1}Rj_*\mathcal{F} \to Rh_*p^{-1}\mathcal{F} \to Q \to
(q^{-1}Rj_*\mathcal{F})[1]
$$
dans $D((X \times_K Y)_\etale)$. Comme $Q$ est supporté par des
points fermés, il suffit de démontrer
$H^i(X \times_K Y, Q) = 0$ pour tout $i$; voir le
lemme \ref{etale-cohomology-lemma-vanishing-closed-points}.
Il suffit donc de démontrer que
$q^{-1}Rj_*\mathcal{F} \to Rh_*p^{-1}\mathcal{F}$
induit un isomorphisme en cohomologie. Rappelons que $\mathcal{F}$
est annulé par un entier $n$ inversible dans $K$.
Par la formule de K\"unneth du lemme \ref{etale-cohomology-lemma-kunneth-one-proper},
nous avons
\begin{align*}
R\Gamma(X \times_K Y, q^{-1}Rj_*\mathcal{F})
&=
R\Gamma(X, Rj_*\mathcal{F})
\otimes_{\mathbf{Z}/n\mathbf{Z}}^\mathbf{L}
R\Gamma(Y, \mathbf{Z}/n\mathbf{Z}) \\
& =
R\Gamma(U, \mathcal{F})
\otimes_{\mathbf{Z}/n\mathbf{Z}}^\mathbf{L}
R\Gamma(Y, \mathbf{Z}/n\mathbf{Z})
\end{align*}
et
$$
R\Gamma(X \times_K Y, Rh_*p^{-1}\mathcal{F}) =
R\Gamma(U \times_K Y, p^{-1}\mathcal{F}) =
R\Gamma(U, \mathcal{F})
\otimes_{\mathbf{Z}/n\mathbf{Z}}^\mathbf{L}
R\Gamma(Y, \mathbf{Z}/n\mathbf{Z})
$$
Cela achève la démonstration.
\end{proof}

\begin{lemma}
\label{etale-cohomology-lemma-punctual-base-change}
Soit $K$ un corps. Pour tout diagramme commutatif
$$
\xymatrix{
X \ar[d] & X' \ar[l] \ar[d]_{f'} & Y \ar[l]^h \ar[d]^e \\
\Spec(K) & S' \ar[l] & T \ar[l]_g
}
$$
de schémas sur $K$ tel que
$X' = X \times_{\Spec(K)} S'$ et $Y = X' \times_{S'} T$, avec
$g$ quasi-compact et quasi-séparé, et pour tout faisceau abélien
$\mathcal{F}$ sur $T_\etale$ dont les fibres sont de torsion d'ordres
inversibles dans $K$, le morphisme de changement de base
$$
(f')^{-1}Rg_*\mathcal{F}
\longrightarrow
Rh_*e^{-1}\mathcal{F}
$$
est un isomorphisme.
\end{lemma}

\begin{proof}
La question est locale sur $X$; nous pouvons donc supposer $X$ affine.
Par le lemme \ref{limits-lemma-relative-approximation} de Limits,
nous pouvons écrire $X = \lim X_i$ comme une limite projective filtrante, à
morphismes de transition affines, de schémas $X_i$ de type fini sur $K$.
Notons $X'_i =  X_i \times_{\Spec(K)} S'$ et $Y_i = X'_i \times_{S'} T$.
Par le lemme \ref{etale-cohomology-lemma-base-change-Rf-star-colim},
il suffit de démontrer l'énoncé pour les carrés
de sommets $X_i, Y_i, S', T$.
Nous pouvons donc supposer $X$ de type fini sur $K$.
De même, nous pouvons écrire
$\mathcal{F} = \colim \mathcal{F}[n]$, où la limite porte
sur le système multiplicatif des entiers inversibles dans $K$.
Le même lemme employé ci-dessus nous ramène au cas où
$\mathcal{F}$ est un faisceau de $\mathbf{Z}/n\mathbf{Z}$-modules
pour un entier $n$ inversible dans $K$.

\medskip\noindent
Nous pouvons remplacer $K$ par sa clôture algébrique $\overline{K}$.
En effet, la formation de l'image directe commute au changement de base
à $\overline{K}$ d'après le
lemme \ref{etale-cohomology-lemma-base-change-field-extension} (pour $g$ comme pour
$h$). Et il suffit de démontrer l'égalité après
restriction à $X'_{\overline{K}}$. Nous pouvons ensuite remplacer
$X$ par sa réduction, grâce à l'invariance topologique de la cohomologie
étale; voir la
proposition \ref{etale-cohomology-proposition-topological-invariance}.
Après ce remplacement, le morphisme $X \to \Spec(K)$
est plat, de présentation finie et à fibres géométriquement réduites;
il en va de même après tout changement de base, notamment pour $X' \to S'$.
Ainsi, $(f')^{-1}Rg_*\mathcal{F} \to Rh_*e^{-1}\mathcal{F}$
est un isomorphisme par le lemme \ref{etale-cohomology-lemma-fppf-reduced-fibres-base-change-f-star}.

\medskip\noindent
Nous pouvons alors appliquer le
lemme \ref{etale-cohomology-lemma-base-change-does-not-hold-post}
pour voir qu'il suffit de démontrer ceci : étant donné un diagramme commutatif
$$
\xymatrix{
X \ar[d]_f & X' \ar[d] \ar[l] & Y \ar[l]^h \ar[d] \\
\Spec(K) & S' \ar[l] & \Spec(L) \ar[l]
}
$$
dont les deux carrés sont cartésiens, où $S'$ est affine, intègre et normal,
de corps des fonctions algébriquement clos $K$, alors
$R^qh_*(\mathbf{Z}/d\mathbf{Z})$ est nul pour $q > 0$ et tout $d | n$.
Remarquons que cette annulation équivaut à l'affirmation que
$$
(f')^{-1}R^q(\Spec(L) \to S')_*\mathbf{Z}/d\mathbf{Z}
\longrightarrow
R^qh_*\mathbf{Z}/d\mathbf{Z}
$$
est un isomorphisme, car le membre de gauche est nul, par exemple par le
lemme \ref{etale-cohomology-lemma-Rf-star-zero-normal-with-alg-closed-function-field}.

\medskip\noindent
Écrivons $S' = \Spec(B)$, de sorte que $L$ soit le corps des fractions de $B$.
Écrivons $B = \bigcup_{i \in I} B_i$ comme la réunion de ses
sous-$K$-algèbres de type fini $B_i$. Soit $J$ l'ensemble des couples $(i, g)$ où
$i \in I$ et $g \in B_i$ est non nul, ordonné par
$(i', g') \geq (i, g)$ si et seulement si $i' \geq i$ et si
l'image de $g$ est inversible dans $(B_{i'})_{g'}$.
Alors $L = \colim_{(i, g) \in J} (B_i)_g$.
Pour $j = (i, g) \in J$, posons $S_j = \Spec(B_i)$
et $U_j = \Spec((B_i)_g)$.
Alors
$$
\vcenter{
\xymatrix{
X' \ar[d] & Y \ar[l]^h \ar[d] \\
S' & \Spec(L) \ar[l]
}
}
\quad\text{est la limite projective des}\quad
\vcenter{
\xymatrix{
X \times_K S_j \ar[d] & X \times_K U_j \ar[l]^{h_j} \ar[d] \\
S_j & U_j \ar[l]
}
}
$$
Nous pouvons donc appliquer le lemme \ref{etale-cohomology-lemma-base-change-Rf-star-colim}
pour voir qu'il suffit de démontrer
le changement de base dans les diagrammes de droite, ce que nous avons
démontré dans le lemme \ref{etale-cohomology-lemma-kunneth-localize-on-X}.
\end{proof}

\begin{lemma}
\label{etale-cohomology-lemma-punctual-base-change-upgrade}
Soit $K$ un corps. Soit $n \geq 1$ inversible dans $K$.
Considérons un diagramme commutatif
$$
\xymatrix{
X \ar[d] & X' \ar[l]^p \ar[d]_{f'} & Y \ar[l]^h \ar[d]^e \\
\Spec(K) & S' \ar[l] & T \ar[l]_g
}
$$
de schémas tel que
$X' = X \times_{\Spec(K)} S'$ et $Y = X' \times_{S'} T$, avec
$g$ quasi-compact et quasi-séparé. Le morphisme canonique
$$
p^{-1}E \otimes_{\mathbf{Z}/n\mathbf{Z}}^\mathbf{L} (f')^{-1}Rg_*F
\longrightarrow
Rh_*(h^{-1}p^{-1}E \otimes_{\mathbf{Z}/n\mathbf{Z}}^\mathbf{L} e^{-1}F)
$$
est un isomorphisme si $E \in D^+(X_\etale, \mathbf{Z}/n\mathbf{Z})$
est d'amplitude de Tor dans $[a, \infty]$ pour un $a \in \mathbf{Z}$ et si
$F \in D^+(T_\etale, \mathbf{Z}/n\mathbf{Z})$.
\end{lemma}

\begin{proof}
Ce lemme généralise le lemme \ref{etale-cohomology-lemma-punctual-base-change}
aux objets de la catégorie dérivée; son assertion est vraie parce que,
dans le lemme \ref{etale-cohomology-lemma-punctual-base-change}, le schéma $X$ sur $K$
est arbitraire. Nous recommandons vivement au lecteur d'omettre cette démonstration laborieuse
(ou de n'en lire que le dernier paragraphe).

\medskip\noindent
Nous pouvons représenter $E$ par un complexe K-plat borné inférieurement
$\mathcal{E}^\bullet$ formé de $\mathbf{Z}/n\mathbf{Z}$-modules plats.
Voir le lemme de Cohomologie sur les sites
\ref{sites-cohomology-lemma-bounded-below-tor-amplitude}.
Choisissons un entier $b$ tel que $H^i(F) = 0$ pour $i < b$.
Choisissons un entier $N$ assez grand et considérons la suite exacte courte
$$
0 \to \sigma_{\geq N + 1}\mathcal{E}^\bullet \to
\mathcal{E}^\bullet \to
\sigma_{\leq N}\mathcal{E}^\bullet \to 0
$$
de troncatures bêtes. Elle produit un triangle distingué
$E'' \to E \to E' \to E''[1]$ dans $D(X_\etale, \mathbf{Z}/n\mathbf{Z})$.
Pour $F$ fixé, les deux membres de la flèche
de l'énoncé du lemme sont des foncteurs exacts en $E$. Remarquons que
$$
p^{-1}E'' \otimes_{\mathbf{Z}/n\mathbf{Z}}^\mathbf{L} (f')^{-1}Rg_*F
\quad\text{et}\quad
Rh_*(h^{-1}p^{-1}E'' \otimes_{\mathbf{Z}/n\mathbf{Z}}^\mathbf{L} e^{-1}F)
$$
sont concentrés en degrés $\geq N + b$. Si nous pouvons donc démontrer le lemme
pour l'objet $E'$, nous voyons qu'il vaut en degrés
$\leq N + b$, ce qui permet de conclure. Certains détails sont omis.
Nous pouvons ainsi supposer que $E$ est représenté
par un complexe borné de $\mathbf{Z}/n\mathbf{Z}$-modules plats.
Un autre argument de même nature permet de supposer
que $E$ est donné par un unique $\mathbf{Z}/n\mathbf{Z}$-module plat
$\mathcal{E}$.

\medskip\noindent
Nous appliquons ensuite les mêmes arguments à la variable $F$
pour nous ramener au cas où $F$ est donné par un unique
faisceau de $\mathbf{Z}/n\mathbf{Z}$-modules $\mathcal{F}$.
Supposons $\mathcal{F}$ annulé par un entier $m | n$.
Si $\ell$ est un nombre premier divisant $m$ et $m > \ell$,
nous pouvons considérer la suite exacte courte
$0 \to \mathcal{F}[\ell] \to \mathcal{F} \to
\mathcal{F}/\mathcal{F}[\ell] \to 0$
et nous ramener à un entier $m$ plus petit. Nous sommes finalement ramenés
au cas où $\mathcal{F}$ est annulé par un nombre
premier $\ell$ divisant $n$.
Dans ce cas, remarquons que
$$
p^{-1}\mathcal{E}
\otimes_{\mathbf{Z}/n\mathbf{Z}}^\mathbf{L}
(f')^{-1}Rg_*\mathcal{F}
=
p^{-1}(\mathcal{E}/\ell \mathcal{E})
\otimes_{\mathbf{F}_\ell}^\mathbf{L}
(f')^{-1}Rg_*\mathcal{F}
$$
par platitude de $\mathcal{E}$. Il en va de même pour l'autre terme.
Cela nous ramène au cas où nous travaillons avec des faisceaux
d'espaces vectoriels sur $\mathbf{F}_\ell$, traité ci-dessous.

\medskip\noindent
Supposons que $\ell$ soit un nombre premier inversible dans $K$.
Supposons que $\mathcal{E}$ et $\mathcal{F}$ soient des faisceaux
d'espaces vectoriels sur $\mathbf{F}_\ell$ sur $X_\etale$ et $T_\etale$.
Nous voulons montrer que
$$
p^{-1}\mathcal{E} \otimes_{\mathbf{F}_\ell} (f')^{-1}R^qg_*\mathcal{F}
\longrightarrow
R^qh_*(h^{-1}p^{-1}\mathcal{E} \otimes_{\mathbf{F}_\ell} e^{-1}\mathcal{F})
$$
est un isomorphisme pour tout $q \geq 0$. Cette question est locale sur $X$;
nous pouvons donc supposer $X$ affine. Nous pouvons écrire $\mathcal{E}$
comme limite inductive filtrante de faisceaux constructibles
d'espaces vectoriels sur $\mathbf{F}_\ell$ sur $X_\etale$; voir le
lemme \ref{etale-cohomology-lemma-torsion-colimit-constructible}.
Comme les produits tensoriels commutent
aux limites inductives filtrantes, et qu'il en va de même des
images directes supérieures (lemme \ref{etale-cohomology-lemma-relative-colimit}),
nous pouvons supposer que $\mathcal{E}$ est un faisceau constructible
d'espaces vectoriels sur $\mathbf{F}_\ell$ sur $X_\etale$.
Nous pouvons alors choisir un entier $m$ et des morphismes finis de présentation finie
$\pi_i : X_i \to X$, $i = 1, \ldots, m$,
tels qu'il existe un morphisme injectif
$$
\mathcal{E} \to
\bigoplus\nolimits_{i = 1}^m
\pi_{i, *}\mathbf{F}_\ell
$$
Voir le lemme \ref{etale-cohomology-lemma-constructible-maps-into-constant-general}.
Remarquons que la somme directe est également un faisceau constructible
(lemme \ref{etale-cohomology-lemma-finite-pushforward-constructible}).
Son conoyau est donc lui aussi constructible
(lemme \ref{etale-cohomology-lemma-constructible-abelian}).
Par décalage de dimension, c'est-à-dire par récurrence sur $q$,
dans la catégorie des faisceaux constructibles
d'espaces vectoriels sur $\mathbf{F}_\ell$ sur $X_\etale$, il suffit de démontrer le
résultat pour les faisceaux $\pi_{i, *}\mathbf{F}_\ell$
(détails omis; indication : commencer par démontrer l'injectivité pour $q = 0$
pour tout $\mathcal{E}$ constructible).
Pour traiter ce cas, complétons le diagramme du lemme en
$$
\xymatrix{
X_i \ar[d]^{\pi_i} &
X'_i \ar[l]^{p_i} \ar[d]^{\pi'_i} &
Y_i \ar[l]^{h_i} \ar[d]^{\rho_i} \\
X \ar[d] & X' \ar[l]^p \ar[d]_{f'} & Y \ar[l]^h \ar[d]^e \\
\Spec(K) & S' \ar[l] & T \ar[l]_g
}
$$
dont tous les carrés sont cartésiens. Dans les égalités ci-dessous, nous
allons employer $R\pi_{i, *} = \pi_{i, *}$, et les égalités analogues
pour $\pi'_i$ et $\rho_i$; nous allons employer la suite spectrale de Leray,
ainsi que
le lemme \ref{etale-cohomology-lemma-finite-pushforward-commutes-with-base-change}, et
nous allons employer
le lemme \ref{etale-cohomology-lemma-projection-formula-proper-mod-n}
(bien que ce dernier soit presque trivial pour les morphismes finis), appliqué à
$\pi_i$, $\pi'_i$ et $\rho_i$.
Nous voyons ainsi que
\begin{align*}
p^{-1}\pi_{i, *}\mathbf{F}_\ell
\otimes_{\mathbf{F}_\ell} (f')^{-1}R^qg_*\mathcal{F}
& =
\pi'_{i, *}\mathbf{F}_\ell
\otimes_{\mathbf{F}_\ell} (f')^{-1}R^qg_*\mathcal{F} \\
& =
\pi'_{i, *}((\pi'_i)^{-1} (f')^{-1}R^qg_*\mathcal{F})
\end{align*}
De même, nous avons
\begin{align*}
R^qh_*(h^{-1}p^{-1} \pi_{i, *}\mathbf{F}_\ell
\otimes_{\mathbf{F}_\ell} e^{-1}\mathcal{F})
& =
R^qh_*(\rho_{i, *}\mathbf{F}_\ell
\otimes_{\mathbf{F}_\ell} e^{-1}\mathcal{F}) \\
& =
R^qh_*\rho_{i, *}(\rho_i^{-1}e^{-1}\mathcal{F}) \\
& =
\pi'_{i, *}R^qh_{i, *} \rho_i^{-1}e^{-1}\mathcal{F}
\end{align*}
Puisque
$R^qh_{i, *} \rho_i^{-1}e^{-1}\mathcal{F} = 
(\pi'_i)^{-1} (f')^{-1}R^qg_*\mathcal{F}$ par le
lemme \ref{etale-cohomology-lemma-punctual-base-change},
nous concluons.
\end{proof}

\begin{lemma}
\label{etale-cohomology-lemma-punctual-base-change-upgrade-unbounded}
Soit $K$ un corps. Soit $n \geq 1$ inversible dans $K$.
Considérons un diagramme commutatif
$$
\xymatrix{
X \ar[d] & X' \ar[l]^p \ar[d]_{f'} & Y \ar[l]^h \ar[d]^e \\
\Spec(K) & S' \ar[l] & T \ar[l]_g
}
$$
de schémas de type fini sur $K$ tel que
$X' = X \times_{\Spec(K)} S'$ et $Y = X' \times_{S'} T$.
Le morphisme canonique
$$
p^{-1}E \otimes_{\mathbf{Z}/n\mathbf{Z}}^\mathbf{L} (f')^{-1}Rg_*F
\longrightarrow
Rh_*(h^{-1}p^{-1}E \otimes_{\mathbf{Z}/n\mathbf{Z}}^\mathbf{L} e^{-1}F)
$$
est un isomorphisme pour $E \in D(X_\etale, \mathbf{Z}/n\mathbf{Z})$
et $F \in D(T_\etale, \mathbf{Z}/n\mathbf{Z})$.
\end{lemma}

\begin{proof}
Nous allons nous ramener au
lemme \ref{etale-cohomology-lemma-punctual-base-change-upgrade}
en employant le fait que nos foncteurs commutent aux sommes directes.
Nous suggérons au lecteur d'omettre cette démonstration.
Rappelons que le produit tensoriel dérivé commute aux
sommes directes. Rappelons aussi que l'image inverse (dérivée) commute aux sommes directes.
Enfin, $Rh_*$ et $Rg_*$ commutent aux sommes directes; voir les
lemmes \ref{etale-cohomology-lemma-finite-cd} et
\ref{etale-cohomology-lemma-finite-cd-mod-n-direct-sums}
(c'est ici que nous employons le fait que nos schémas sont de type fini
sur $K$).

\medskip\noindent
Pour achever la démonstration, nous pouvons raisonner comme suit.
Écrivons d'abord $E = \text{hocolim} \tau_{\leq N} E$.
Comme nos foncteurs commutent aux sommes directes, ils commutent
aux colimites homotopiques. Il suffit donc de démontrer
le lemme pour $E$ borné supérieurement. De même, nous pouvons
supposer $F$ borné supérieurement.
Nous pouvons alors représenter $E$ par un complexe borné supérieurement
$\mathcal{E}^\bullet$ de faisceaux de $\mathbf{Z}/n\mathbf{Z}$-modules.
Alors
$$
\mathcal{E}^\bullet = \colim \sigma_{\geq -N}\mathcal{E}^\bullet
$$
(troncatures bêtes).
Nous pouvons donc supposer que $\mathcal{E}^\bullet$ est un complexe
borné de faisceaux de $\mathbf{Z}/n\mathbf{Z}$-modules.
Pour $F$, nous choisissons un complexe borné supérieurement de
faisceaux plats (!) de $\mathbf{Z}/n\mathbf{Z}$-modules.
Nous nous ramenons alors au cas où $F$ est représenté
par un complexe borné de faisceaux plats de $\mathbf{Z}/n\mathbf{Z}$-modules.
Le lemme \ref{etale-cohomology-lemma-punctual-base-change-upgrade}
s'applique, et nous concluons.
\end{proof}

\begin{lemma}
\label{etale-cohomology-lemma-kunneth}
Soit $k$ un corps séparablement clos. Soient $X$ et $Y$ des
schémas de type fini sur $k$. Soit $n \geq 1$ un entier
inversible dans $k$. Alors, pour
$E \in D(X_\etale, \mathbf{Z}/n\mathbf{Z})$ et
$K \in D(Y_\etale, \mathbf{Z}/n\mathbf{Z})$
nous avons
$$
R\Gamma(X \times_{\Spec(k)} Y,
\text{pr}_1^{-1}E
\otimes_{\mathbf{Z}/n\mathbf{Z}}^\mathbf{L}
\text{pr}_2^{-1}K
)
=
R\Gamma(X, E)
\otimes_{\mathbf{Z}/n\mathbf{Z}}^\mathbf{L}
R\Gamma(Y, K)
$$
\end{lemma}

\begin{proof}
Par le lemme \ref{etale-cohomology-lemma-punctual-base-change-upgrade-unbounded}, nous avons
$$
R\text{pr}_{1, *}(
\text{pr}_1^{-1}E
\otimes_{\mathbf{Z}/n\mathbf{Z}}^\mathbf{L}
\text{pr}_2^{-1}K) =
E \otimes_{\mathbf{Z}/n\mathbf{Z}}^\mathbf{L}
\underline{R\Gamma(Y, K)}
$$
Nous concluons par le
lemme \ref{etale-cohomology-lemma-pull-out-constant-mod-n},
que nous pouvons employer parce que
$\text{cd}(X) < \infty$ par le lemme \ref{etale-cohomology-lemma-finite-cd}.
\end{proof}














\section{Comparaison des topologies chaotique et de Zariski}
\label{etale-cohomology-section-compare-chaotic-Zariski}

\noindent
Lorsqu'on construit le faisceau structural d'un schéma affine, on en construit d'abord
les valeurs sur les ouverts affines, puis on l'étend à tous les ouverts.
Une construction analogue est souvent utile pour construire des complexes
de groupes abéliens sur un schéma $X$. Rappelons que $X_{affine, Zar}$
désigne la catégorie des ouverts affines de $X$, munie de la topologie définie
par les recouvrements de Zariski standard; voir la
définition \ref{topologies-definition-big-small-Zariski} de Topologies.
Rappelons au lecteur que le topos de $X_{affine, Zar}$
est le petit topos de Zariski de $X$; voir le lemme de Topologies
\ref{topologies-lemma-alternative-zariski}.
Dans cette section, nous notons $X_{affine}$ la même catégorie
sous-jacente munie de la topologie chaotique, c'est-à-dire telle que les faisceaux
coïncident avec les préfaisceaux. Nous obtenons un morphisme de sites
$$
\epsilon : X_{affine, Zar} \longrightarrow X_{affine}
$$
comme dans la section \ref{sites-cohomology-section-compare} de Cohomologie sur les sites.

\begin{lemma}
\label{etale-cohomology-lemma-check-zar}
Dans la situation ci-dessus, soit $K$ un objet de $D^+(X_{affine})$.
Alors $K$ appartient à l'image essentielle du foncteur (pleinement fidèle)
$R\epsilon_* ; D(X_{affine, Zar}) \to D(X_{affine})$ si et seulement
si les deux conditions suivantes sont satisfaites :
\begin{enumerate}
\item $R\Gamma(\emptyset, K)$ est nul dans $D(\textit{Ab})$, et
\item si $U = V \cup W$, où $U, V, W \subset X$ sont des ouverts affines et
$V, W \subset U$ des ouverts standard
(définition \ref{algebra-definition-Zariski-topology} d'Algèbre), alors
le morphisme $c^K_{U, V, W, V \cap W}$ du
lemme \ref{sites-cohomology-lemma-c-square} de Cohomologie sur les sites
est un quasi-isomorphisme.
\end{enumerate}
\end{lemma}

\begin{proof}
(Le foncteur $R\epsilon_*$ est pleinement fidèle d'après la discussion de la
section \ref{sites-cohomology-section-compare} de Cohomologie sur les sites.)
À un problème près concernant l'ensemble vide,
cela découle du très général lemme de Cohomologie sur les sites
\ref{sites-cohomology-lemma-descent-squares}, dont les hypothèses valent par le
lemme \ref{schemes-lemma-sheaf-on-affines} de Schémas et le
lemme de Cohomologie sur les sites
\ref{sites-cohomology-lemma-descent-squares-helper}.

\medskip\noindent
Pour contourner ce problème, notons $X_{affine, almost-chaotic}$
le site où l'objet vide $\emptyset$ possède deux recouvrements,
à savoir $\{\emptyset \to \emptyset\}$ et le recouvrement vide
(voir l'exemple \ref{sites-example-site-topological} de Sites pour une
discussion). Nous avons alors des morphismes de sites
$$
X_{affine, Zar} \to X_{affine, almost-chaotic} \to X_{affine}
$$
L'argument ci-dessus vaut pour la première flèche. Nous laissons alors
au lecteur le soin de vérifier qu'un objet $K$ de $D^+(X_{affine})$
appartient à l'image essentielle du foncteur (pleinement fidèle)
$D(X_{affine, almost-chaotic}) \to D(X_{affine})$ si et seulement
si $R\Gamma(\emptyset, K)$ est nul dans $D(\textit{Ab})$.
\end{proof}









\section{Comparaison des grands et petits topos}
\label{etale-cohomology-section-compare}

\noindent
Soit $S$ un schéma. Dans le
lemme \ref{topologies-lemma-at-the-bottom-etale} de Topologies,
nous avons introduit les morphismes de comparaison
$\pi_S : (\Sch/S)_\etale \to S_\etale$ et
$i_S : \Sh(S_\etale) \to \Sh((\Sch/S)_\etale)$,
avec $\pi_S \circ i_S = \text{id}$ et $\pi_{S, *} = i_S^{-1}$.
Plus généralement, si $f : T \to S$ est un objet de $(\Sch/S)_\etale$,
il existe un morphisme $i_f : \Sh(T_\etale) \to \Sh((\Sch/S)_\etale)$
tel que $f_{small} = \pi_S \circ i_f$; voir les
lemmes \ref{topologies-lemma-put-in-T-etale} et
\ref{topologies-lemma-morphism-big-small-etale} de Topologies. Dans la
remarque \ref{descent-remark-change-topologies-ringed} de Descente,
nous les avons étendus en un morphisme de sites annelés
$$
\pi_S : ((\Sch/S)_\etale, \mathcal{O}) \to (S_\etale, \mathcal{O}_S)
$$
et des morphismes de topos annelés
$$
i_S : (\Sh(S_\etale), \mathcal{O}_S) \to (\Sh((\Sch/S)_\etale), \mathcal{O})
$$
et
$$
i_f : (\Sh(T_\etale), \mathcal{O}_T) \to (\Sh((\Sch/S)_\etale, \mathcal{O}))
$$
Remarquons que la restriction $i_S^{-1} = \pi_{S, *}$ (voir la
définition \ref{topologies-definition-restriction-small-etale} de Topologies)
transforme $\mathcal{O}$ en $\mathcal{O}_S$.
De même, $i_f^{-1}$ transforme $\mathcal{O}$ en $\mathcal{O}_T$.
Voir la remarque \ref{descent-remark-change-topologies-ringed} de Descente.
Ainsi, $i_S^*\mathcal{F} = i_S^{-1}\mathcal{F}$ et
$i_f^*\mathcal{F} = i_f^{-1}\mathcal{F}$ pour tout $\mathcal{O}$-module
$\mathcal{F}$ sur $(\Sch/S)_\etale$. En particulier, $i_S^*$ et $i_f^*$
sont des foncteurs exacts. Le foncteur $i_S^*$ est souvent noté
$\mathcal{F} \mapsto \mathcal{F}|_{S_\etale}$ (ce qui n'entre pas
en conflit avec la notation de la
définition \ref{topologies-definition-restriction-small-etale} de Topologies).

\begin{lemma}
\label{etale-cohomology-lemma-compare-injectives}
Soit $S$ un schéma. Soit $T$ un objet de $(\Sch/S)_\etale$.
\begin{enumerate}
\item Si $\mathcal{I}$ est injectif dans $\textit{Ab}((\Sch/S)_\etale)$, alors
\begin{enumerate}
\item $i_f^{-1}\mathcal{I}$ est injectif dans $\textit{Ab}(T_\etale)$,
\item $\mathcal{I}|_{S_\etale}$ est injectif dans $\textit{Ab}(S_\etale)$,
\end{enumerate}
\item Si $\mathcal{I}^\bullet$ est un complexe K-injectif
dans $\textit{Ab}((\Sch/S)_\etale)$, alors
\begin{enumerate}
\item $i_f^{-1}\mathcal{I}^\bullet$ est un complexe K-injectif dans
$\textit{Ab}(T_\etale)$,
\item $\mathcal{I}^\bullet|_{S_\etale}$ est un complexe K-injectif dans
$\textit{Ab}(S_\etale)$,
\end{enumerate}
\end{enumerate}
Les énoncés correspondants ne valent pas pour les modules.
\end{lemma}

\begin{proof}
Les parties (1)(b) et (2)(b)
découlent formellement du fait que le foncteur de restriction
$\pi_{S, *} = i_S^{-1}$ est adjoint à droite du foncteur exact
$\pi_S^{-1}$; voir le
lemme \ref{homology-lemma-adjoint-preserve-injectives} d'Homology et le
lemme \ref{derived-lemma-adjoint-preserve-K-injectives} de Catégories dérivées.

\medskip\noindent
Les parties (1)(a) et (2)(a) se voient de deux manières. Première démonstration : on peut utiliser
le fait que $i_f^{-1}$ est adjoint à droite du foncteur exact $i_{f, !}$.
Ce foncteur est construit dans le
lemme \ref{topologies-lemma-put-in-T-etale} de Topologies
pour les faisceaux d'ensembles, et dans le
lemme \ref{sites-modules-lemma-g-shriek-adjoint} de Modules sur les sites pour les faisceaux abéliens.
Le lemme de Modules sur les sites
\ref{sites-modules-lemma-exactness-lower-shriek} montre qu'il est exact.
Seconde démonstration. Nous pouvons employer $i_f = i_T \circ f_{big}$, comme le montre
le lemme \ref{topologies-lemma-morphism-big-small-etale} de Topologies.
Puisque $f_{big}$ est une localisation, son image inverse
conserve les injectifs et les K-injectifs; voir les
lemmes \ref{sites-cohomology-lemma-cohomology-of-open} et
\ref{sites-cohomology-lemma-restrict-K-injective-to-open}.
Nous appliquons alors les parties (1)(b) et (2)(b) déjà démontrées
au foncteur $i_T^{-1}$, et nous concluons.

\medskip\noindent
Soit $S = \Spec(\mathbf{Z})$ et considérons le morphisme
$2 : \mathcal{O}_S \to \mathcal{O}_S$. C'est un morphisme injectif
de $\mathcal{O}_S$-modules sur $S_\etale$. Cependant, l'image inverse
$\pi_S^*(2) : \mathcal{O} \to \mathcal{O}$ n'est pas injective,
comme on le voit en l'évaluant sur $\Spec(\mathbf{F}_2)$. Choisissons maintenant
une injection $\alpha : \mathcal{O} \to \mathcal{I}$ dans un
$\mathcal{O}$-module injectif $\mathcal{I}$ sur $(\Sch/S)_\etale$.
Considérons alors le diagramme
$$
\xymatrix{
\mathcal{O}_S \ar[d]_2 \ar[rr]_{\alpha|_{S_\etale}} & &
\mathcal{I}|_{S_\etale} \\
\mathcal{O}_S \ar@{..>}[rru]
}
$$
La flèche pointillée ne peut exister dans la catégorie des $\mathcal{O}_S$-modules,
car cela signifierait
(par adjonction) que le morphisme injectif $\alpha$ se factorise par
le morphisme non injectif $\pi_S^*(2)$, ce qui est impossible.
Ainsi, $\mathcal{I}|_{S_\etale}$ n'est pas un $\mathcal{O}_S$-module injectif.
\end{proof}

\noindent
Soit $f : T \to S$ un morphisme de schémas. Le diagramme commutatif du
lemme \ref{topologies-lemma-morphism-big-small-etale} (3) de Topologies
donne un diagramme commutatif de sites annelés
$$
\xymatrix{
(T_\etale, \mathcal{O}_T) \ar[d]_{f_{small}} &
((\Sch/T)_\etale, \mathcal{O}) \ar[d]^{f_{big}} \ar[l]^{\pi_T} \\
(S_\etale, \mathcal{O}_S) &
((\Sch/S)_\etale, \mathcal{O}) \ar[l]_{\pi_S}
}
$$
comme on le voit aisément en explicitant les définitions de
$f_{small}^\sharp$, $f_{big}^\sharp$, $\pi_S^\sharp$ et $\pi_T^\sharp$.
Cela signifie en particulier que
\begin{equation}
\label{etale-cohomology-equation-compare-big-small}
(f_{big, *}\mathcal{F})|_{S_\etale} =
f_{small, *}(\mathcal{F}|_{T_\etale})
\end{equation}
pour tout faisceau $\mathcal{F}$ sur $(\Sch/T)_\etale$; si $\mathcal{F}$
est un faisceau de $\mathcal{O}$-modules, alors (\ref{etale-cohomology-equation-compare-big-small})
est un isomorphisme de $\mathcal{O}_S$-modules sur $S_\etale$.

\begin{lemma}
\label{etale-cohomology-lemma-compare-higher-direct-image}
Soit $f : T \to S$ un morphisme de schémas.
\begin{enumerate}
\item Pour $K$ dans $D((\Sch/T)_\etale)$, nous avons
$
(Rf_{big, *}K)|_{S_\etale} = Rf_{small, *}(K|_{T_\etale})
$
dans $D(S_\etale)$.
\item Pour $K$ dans $D((\Sch/T)_\etale, \mathcal{O})$, nous avons
$
(Rf_{big, *}K)|_{S_\etale} = Rf_{small, *}(K|_{T_\etale})
$
dans $D(\textit{Mod}(S_\etale, \mathcal{O}_S))$.
\end{enumerate}
Plus généralement, soit $g : S' \to S$ un objet de $(\Sch/S)_\etale$.
Considérons le produit fibré
$$
\xymatrix{
T' \ar[r]_{g'} \ar[d]_{f'} & T \ar[d]^f \\
S' \ar[r]^g & S
}
$$
Alors
\begin{enumerate}
\item[(3)] Pour $K$ dans $D((\Sch/T)_\etale)$, nous avons
$i_g^{-1}(Rf_{big, *}K) = Rf'_{small, *}(i_{g'}^{-1}K)$
dans $D(S'_\etale)$.
\item[(4)] Pour $K$ dans $D((\Sch/T)_\etale, \mathcal{O})$, nous avons
$i_g^*(Rf_{big, *}K) = Rf'_{small, *}(i_{g'}^*K)$
dans $D(\textit{Mod}(S'_\etale, \mathcal{O}_{S'}))$.
\item[(5)] Pour $K$ dans $D((\Sch/T)_\etale)$, nous avons
$g_{big}^{-1}(Rf_{big, *}K) = Rf'_{big, *}((g'_{big})^{-1}K)$
dans $D((\Sch/S')_\etale)$.
\item[(6)] Pour $K$ dans $D((\Sch/T)_\etale, \mathcal{O})$, nous avons
$g_{big}^*(Rf_{big, *}K) = Rf'_{big, *}((g'_{big})^*K)$
dans $D(\textit{Mod}(S'_\etale, \mathcal{O}_{S'}))$.
\end{enumerate}
\end{lemma}

\begin{proof}
La partie (1) découle du
lemme \ref{etale-cohomology-lemma-compare-injectives}
et de (\ref{etale-cohomology-equation-compare-big-small}),
en choisissant un complexe K-injectif de faisceaux abéliens représentant $K$.

\medskip\noindent
La partie (3) découle du lemme \ref{etale-cohomology-lemma-compare-injectives}
et du lemme de Topologies
\ref{topologies-lemma-morphism-big-small-cartesian-diagram-etale}
en choisissant un complexe K-injectif de faisceaux abéliens représentant $K$.

\medskip\noindent
La partie (5) est le lemme de Cohomologie sur les sites
\ref{sites-cohomology-lemma-localize-cartesian-square}.

\medskip\noindent
La partie (6) est le lemme de Cohomologie sur les sites
\ref{sites-cohomology-lemma-localize-cartesian-square-modules}.

\medskip\noindent
La partie (2) peut se démontrer comme suit. Nous avons vu ci-dessus
que $\pi_S \circ f_{big} = f_{small} \circ \pi_T$ comme morphismes
de sites annelés. Nous obtenons donc
$R\pi_{S, *} \circ Rf_{big, *} = Rf_{small, *} \circ R\pi_{T, *}$
par le lemme de Cohomologie sur les sites
\ref{sites-cohomology-lemma-derived-pushforward-composition}.
Puisque les foncteurs de restriction $\pi_{S, *}$ et $\pi_{T, *}$
sont exacts, nous concluons.

\medskip\noindent
La partie (4) découle de la partie (6) et de la partie (2) appliquée à $f' : T' \to S'$.
\end{proof}

\noindent
Soit $S$ un schéma et soit $\mathcal{H}$ un faisceau abélien sur
$(\Sch/S)_\etale$. Rappelons que $H^n_\etale(U, \mathcal{H})$ désigne
la cohomologie de $\mathcal{H}$ sur un objet $U$ de $(\Sch/S)_\etale$.

\begin{lemma}
\label{etale-cohomology-lemma-compare-cohomology}
Soit $f : T \to S$ un morphisme de schémas. Alors
\begin{enumerate}
\item Pour $K$ dans $D(S_\etale)$, nous avons
$H^n_\etale(S, \pi_S^{-1}K) = H^n(S_\etale, K)$.
\item Pour $K$ dans $D(S_\etale, \mathcal{O}_S)$, nous avons
$H^n_\etale(S, L\pi_S^*K) = H^n(S_\etale, K)$.
\item Pour $K$ dans $D(S_\etale)$, nous avons
$H^n_\etale(T, \pi_S^{-1}K) = H^n(T_\etale, f_{small}^{-1}K)$.
\item Pour $K$ dans $D(S_\etale, \mathcal{O}_S)$, nous avons
$H^n_\etale(T, L\pi_S^*K) = H^n(T_\etale, Lf_{small}^*K)$.
\item Pour $M$ dans $D((\Sch/S)_\etale)$, nous avons
$H^n_\etale(T, M) = H^n(T_\etale, i_f^{-1}M)$.
\item Pour $M$ dans $D((\Sch/S)_\etale, \mathcal{O})$, nous avons
$H^n_\etale(T, M) = H^n(T_\etale, i_f^*M)$.
\end{enumerate}
\end{lemma}

\begin{proof}
Pour démontrer (5), représentons $M$ par un complexe K-injectif de faisceaux abéliens,
appliquons le lemme \ref{etale-cohomology-lemma-compare-injectives}
et explicitons les définitions. La partie (3) en découle,
car $i_f^{-1}\pi_S^{-1} = f_{small}^{-1}$. La partie (1) est un cas
particulier de (3).

\medskip\noindent
La partie (6) découle du très général lemme de Cohomologie sur les sites
\ref{sites-cohomology-lemma-pullback-same-cohomology}. La partie
(4) en découle, car $Lf_{small}^* = i_f^* \circ L\pi_S^*$.
La partie (2) est un cas particulier de (4).
\end{proof}

\begin{lemma}
\label{etale-cohomology-lemma-cohomological-descent-etale}
Soit $S$ un schéma. Pour $K \in D(S_\etale)$, le morphisme
$$
K \longrightarrow R\pi_{S, *}\pi_S^{-1}K
$$
est un isomorphisme.
\end{lemma}

\begin{proof}
C'est vrai parce que $\pi_S^{-1}$ et $\pi_{S, *} = i_S^{-1}$
sont tous deux des foncteurs exacts, et que le composé $\pi_{S, *} \circ \pi_S^{-1}$
est le foncteur identité.
\end{proof}

\begin{lemma}
\label{etale-cohomology-lemma-compare-higher-direct-image-proper}
Soit $f : T \to S$ un morphisme propre de schémas. Nous avons alors
\begin{enumerate}
\item $\pi_S^{-1} \circ f_{small, *} = f_{big, *} \circ \pi_T^{-1}$
comme foncteurs $\Sh(T_\etale) \to \Sh((\Sch/S)_\etale)$,
\item $\pi_S^{-1}Rf_{small, *}K = Rf_{big, *}\pi_T^{-1}K$
pour $K$ dans $D^+(T_\etale)$ dont les faisceaux de cohomologie sont de torsion,
\item $\pi_S^{-1}Rf_{small, *}K = Rf_{big, *}\pi_T^{-1}K$
pour $K$ dans $D(T_\etale, \mathbf{Z}/n\mathbf{Z})$, et
\item $\pi_S^{-1}Rf_{small, *}K = Rf_{big, *}\pi_T^{-1}K$
pour tout $K$ dans $D(T_\etale)$ si $f$ est fini.
\end{enumerate}
\end{lemma}

\begin{proof}
Démonstration de (1). Soit $\mathcal{F}$ un faisceau sur $T_\etale$.
Soit $g : S' \to S$ un objet de $(\Sch/S)_\etale$. Considérons le
produit fibré
$$
\xymatrix{
T' \ar[r]_{f'} \ar[d]_{g'} & S' \ar[d]^g \\
T \ar[r]^f & S
}
$$
Nous avons alors
$$
(f_{big, *}\pi_T^{-1}\mathcal{F})(S') =
(\pi_T^{-1}\mathcal{F})(T') =
((g'_{small})^{-1}\mathcal{F})(T')  =
(f'_{small, *}(g'_{small})^{-1}\mathcal{F})(S')
$$
où la deuxième égalité vient du lemme \ref{etale-cohomology-lemma-describe-pullback}.
D'autre part,
$$
(\pi_S^{-1}f_{small, *}\mathcal{F})(S') =
(g_{small}^{-1}f_{small, *}\mathcal{F})(S')
$$
de nouveau par le lemme \ref{etale-cohomology-lemma-describe-pullback}.
Le changement de base propre pour les faisceaux d'ensembles
(lemme \ref{etale-cohomology-lemma-proper-base-change-f-star})
montre donc que les deux ensembles sont canoniquement isomorphes.
L'isomorphisme est compatible avec les morphismes de restriction
et définit un isomorphisme
$\pi_S^{-1}f_{small, *}\mathcal{F} = f_{big, *}\pi_T^{-1}\mathcal{F}$.
Il définit ainsi un isomorphisme de foncteurs
$\pi_S^{-1} \circ f_{small, *} = f_{big, *} \circ \pi_T^{-1}$.

\medskip\noindent
Démonstration de (2). Il existe un morphisme canonique de changement de base
$\pi_S^{-1}Rf_{small, *}K \to Rf_{big, *}\pi_T^{-1}K$
pour tout $K$ dans $D(T_\etale)$; voir la
remarque \ref{sites-cohomology-remark-base-change} de Cohomologie sur les sites.
Pour montrer que c'est un isomorphisme, il suffit de montrer que l'image inverse du
morphisme de changement de base par $i_g : \Sh(S'_\etale) \to \Sh((\Sch/S)_\etale)$
est un isomorphisme pour tout objet $g : S' \to S$ de $(\Sch/S)_\etale$.
Soient $T', g', f'$ comme au paragraphe précédent.
L'image inverse du morphisme de changement de base est
\begin{align*}
g_{small}^{-1}Rf_{small, *}K
& =
i_g^{-1}\pi_S^{-1}Rf_{small, *}K \\
& \to
i_g^{-1}Rf_{big, *}\pi_T^{-1}K \\
& =
Rf'_{small, *}(i_{g'}^{-1}\pi_T^{-1}K) \\
& =
Rf'_{small, *}((g'_{small})^{-1}K)
\end{align*}
où nous avons employé $\pi_S \circ i_g = g_{small}$,
$\pi_T \circ i_{g'} = g'_{small}$ et le
lemme \ref{etale-cohomology-lemma-compare-higher-direct-image}.
Ce morphisme est un isomorphisme par le théorème de changement de base pour un morphisme propre
(lemme \ref{etale-cohomology-lemma-proper-base-change}), pourvu que $K$ soit borné
inférieurement et que les faisceaux de cohomologie de $K$ soient de torsion.

\medskip\noindent
La démonstration de la partie (3) est la même que celle de la partie (2), sauf que
nous employons le lemme \ref{etale-cohomology-lemma-proper-base-change-mod-n}
au lieu du lemme \ref{etale-cohomology-lemma-proper-base-change}.

\medskip\noindent
Démonstration de (4). Si $f$ est fini, les foncteurs
$f_{small, *}$ et $f_{big, *}$ sont exacts. Cela découle
de la proposition \ref{etale-cohomology-proposition-finite-higher-direct-image-zero}
pour $f_{small}$. Comme tout changement de base $f'$ de $f$ est également fini,
la partie (3) du lemme \ref{etale-cohomology-lemma-compare-higher-direct-image} montre
que $f_{big, *}$ est lui aussi exact (ses foncteurs dérivés supérieurs étant nuls).
Ce cas découle donc de la partie (1).
\end{proof}




\section{Comparaison des topologies fppf et étale}
\label{etale-cohomology-section-fppf-etale}

\noindent
Un modèle pour cette section est la section consacrée à la comparaison de la
topologie usuelle et de la topologie qc sur les espaces topologiques
localement compacts, étudiée dans la section de Cohomologie sur les sites
\ref{sites-cohomology-section-cohomology-LC}.
Nous rappelons d'abord certains résultats des
sections de Topologies
\ref{topologies-section-change-topologies} et
\ref{topologies-section-etale}.

\medskip\noindent
Soit $S$ un schéma et soit $(\Sch/S)_{fppf}$ un site fppf.
Sur la même catégorie sous-jacente, nous avons une seconde topologie,
à savoir la topologie étale, et donc un second site
$(\Sch/S)_\etale$. Le foncteur identité
$(\Sch/S)_\etale \to (\Sch/S)_{fppf}$ est continu et définit
un morphisme de sites
$$
\epsilon_S : (\Sch/S)_{fppf} \longrightarrow (\Sch/S)_\etale
$$
Voir la section \ref{sites-cohomology-section-compare} de Cohomologie sur les sites.
Remarquons que $\epsilon_{S, *}$ est le foncteur identité sur les
préfaisceaux sous-jacents et que $\epsilon_S^{-1}$ associe à un faisceau étale son
faisceau fppf associé. Soit $S_\etale$ le petit site étale.
Il existe un morphisme de sites
$$
\pi_S : (\Sch/S)_\etale \longrightarrow S_\etale
$$
induit par le foncteur continu
$S_\etale \to (\Sch/S)_\etale$, $U \mapsto U$.
En effet, $S_\etale$ possède des produits fibrés et un objet final, le
foncteur ci-dessus y commute, et la
proposition \ref{sites-proposition-get-morphism} de Sites s'applique.

\begin{lemma}
\label{etale-cohomology-lemma-describe-pullback-pi-fppf}
Avec les notations ci-dessus.
Soit $\mathcal{F}$ un faisceau sur $S_\etale$. La règle
$$
(\Sch/S)_{fppf} \longrightarrow \textit{Ensembles},\quad
(f : X \to S) \longmapsto \Gamma(X, f_{small}^{-1}\mathcal{F})
$$
est un faisceau fppf, et a fortiori un faisceau sur $(\Sch/S)_\etale$.
En fait, ce faisceau est égal à
$\pi_S^{-1}\mathcal{F}$ sur $(\Sch/S)_\etale$ et à
$\epsilon_S^{-1}\pi_S^{-1}\mathcal{F}$ sur $(\Sch/S)_{fppf}$.
\end{lemma}

\begin{proof}
L'énoncé relatif à la topologie étale est le contenu
du lemme \ref{etale-cohomology-lemma-describe-pullback}. Pour achever la démonstration, il
suffit de montrer que $\pi_S^{-1}\mathcal{F}$ est un faisceau pour la topologie
fppf. Cela est démontré dans le lemme \ref{etale-cohomology-lemma-describe-pullback}
également.
\end{proof}

\noindent
Dans la situation du lemme \ref{etale-cohomology-lemma-describe-pullback-pi-fppf},
la composée de $\epsilon_S$ et $\pi_S$, ainsi que l'égalité,
déterminent un morphisme de sites
$$
a_S : (\Sch/S)_{fppf} \longrightarrow S_\etale
$$

\begin{lemma}
\label{etale-cohomology-lemma-push-pull-fppf-etale}
Avec les notations ci-dessus.
Soit $f : X \to Y$ un morphisme de $(\Sch/S)_{fppf}$.
Il existe alors des diagrammes commutatifs de topos
$$
\xymatrix{
\Sh((\Sch/X)_{fppf}) \ar[rr]_{f_{big, fppf}} \ar[d]_{\epsilon_X} & &
\Sh((\Sch/Y)_{fppf}) \ar[d]^{\epsilon_Y} \\
\Sh((\Sch/X)_\etale) \ar[rr]^{f_{big, \etale}} & &
\Sh((\Sch/Y)_\etale)
}
$$
et
$$
\xymatrix{
\Sh((\Sch/X)_{fppf}) \ar[rr]_{f_{big, fppf}} \ar[d]_{a_X} & &
\Sh((\Sch/Y)_{fppf}) \ar[d]^{a_Y} \\
\Sh(X_\etale) \ar[rr]^{f_{small}} & &
\Sh(Y_\etale)
}
$$
avec $a_X = \pi_X \circ \epsilon_X$ et $a_Y = \pi_X \circ \epsilon_X$.
\end{lemma}

\begin{proof}
La commutativité des diagrammes découle de la discussion de la
section \ref{topologies-section-change-topologies} de Topologies.
\end{proof}

\begin{lemma}
\label{etale-cohomology-lemma-proper-push-pull-fppf-etale}
Dans le lemme \ref{etale-cohomology-lemma-push-pull-fppf-etale}, si $f$ est propre, alors nous avons
$a_Y^{-1} \circ f_{small, *} = f_{big, fppf, *} \circ a_X^{-1}$.
\end{lemma}

\begin{proof}
On peut le démontrer en répétant la démonstration de la
partie (1) du lemme \ref{etale-cohomology-lemma-compare-higher-direct-image-proper};
nous allons plutôt l'en déduire.
Comme $\epsilon_{Y, *}$ est le foncteur identité sur les préfaisceaux sous-jacents,
il reflète les isomorphismes. La description
du lemme \ref{etale-cohomology-lemma-describe-pullback-pi-fppf}
montre que $\epsilon_{Y, *} \circ a_Y^{-1} = \pi_Y^{-1}$,
et de même pour $X$. Pour montrer que le morphisme canonique
$a_Y^{-1}f_{small, *}\mathcal{F} \to f_{big, fppf, *}a_X^{-1}\mathcal{F}$
est un isomorphisme, il suffit de montrer que
\begin{align*}
\pi_Y^{-1}f_{small, *}\mathcal{F}
& =
\epsilon_{Y, *}a_Y^{-1}f_{small, *}\mathcal{F} \\
& \to 
\epsilon_{Y, *}f_{big, fppf, *}a_X^{-1}\mathcal{F} \\
& =
f_{big, \etale, *} \epsilon_{X, *}a_X^{-1}\mathcal{F} \\
& =
f_{big, \etale, *}\pi_X^{-1}\mathcal{F}
\end{align*}
est un isomorphisme. C'est la partie
(1) du lemme \ref{etale-cohomology-lemma-compare-higher-direct-image-proper}.
\end{proof}

\begin{lemma}
\label{etale-cohomology-lemma-descent-sheaf-fppf-etale}
Dans le lemme \ref{etale-cohomology-lemma-push-pull-fppf-etale}, supposons
$f$ plat, localement de présentation finie et surjectif.
Alors le foncteur
$$
\Sh(Y_\etale) \longrightarrow
\left\{
(\mathcal{G}, \mathcal{H}, \alpha)
\middle|
\begin{matrix}
\mathcal{G} \in \Sh(X_\etale),\ \mathcal{H} \in \Sh((\Sch/Y)_{fppf}), \\
\alpha : a_X^{-1}\mathcal{G} \to f_{big, fppf}^{-1}\mathcal{H}
\text{ un isomorphisme}
\end{matrix}
\right\}
$$
qui associe à $\mathcal{F}$
$(f_{small}^{-1}\mathcal{F}, a_Y^{-1}\mathcal{F}, can)$ est une équivalence.
\end{lemma}

\begin{proof}
Le foncteur $a_X^{-1}$ est pleinement fidèle (car $a_{X, *}a_X^{-1} = \text{id}$ par le
lemme \ref{etale-cohomology-lemma-describe-pullback-pi-fppf}). Le foncteur d'oubli
$(\mathcal{G}, \mathcal{H}, \alpha) \mapsto \mathcal{H}$ identifie donc la
catégorie des triplets à une sous-catégorie pleine de $\Sh((\Sch/Y)_{fppf})$.
De plus, le foncteur $a_Y^{-1}$ est pleinement fidèle; le foncteur
du lemme est donc lui aussi pleinement fidèle.

\medskip\noindent
Supposons donné un recouvrement étale $\{Y_i \to Y\}$.
Soit $f_i : X_i \to Y_i$ le changement de base de $f$.
Notons $f_{ij} = f_i \times f_j : X_i \times_X X_j  \to Y_i \times_Y Y_j$.
Assertion : si le lemme est vrai pour $f_i$ et $f_{ij}$ pour tous $i, j$, alors
il est vrai pour $f$. Pour le voir, remarquons que le recouvrement étale donné
détermine un recouvrement étale de l'objet final dans chacun des
quatre sites $Y_\etale, X_\etale, (\Sch/Y)_{fppf}, (\Sch/X)_{fppf}$.
La catégorie des faisceaux équivaut donc à la catégorie des données de
recollement pour ce recouvrement
(lemme \ref{sites-lemma-mapping-property-glue} de Sites)
dans chacun des quatre cas. Un grand diagramme commutatif de
catégories achève alors la démonstration de l'affirmation. Nous omettons les détails.
L'assertion montre que nous pouvons travailler localement sur $Y$ pour la topologie étale.

\medskip\noindent
Remarquons que $\{X \to Y\}$ est un recouvrement fppf. En travaillant localement sur $Y$,
nous pouvons supposer qu'il existe un morphisme $s : X' \to X$
tel que le composé $f' = f \circ s : X' \to Y$ soit fini localement libre et surjectif ; voir le
lemme de Compléments sur les morphismes
\ref{more-morphisms-lemma-dominate-fppf-etale-locally}.
Assertion : si le lemme est vrai pour $f'$, alors il est vrai pour $f$.
En effet, étant donné un triplet $(\mathcal{G}, \mathcal{H}, \alpha)$
pour $f$, son image inverse par $s$ donne un triplet
$(s_{small}^{-1}\mathcal{G}, \mathcal{H}, s_{big, fppf}^{-1}\alpha)$
pour $f'$. Une solution de ce triplet donne un faisceau
$\mathcal{F}$ sur $Y_\etale$ tel que $a_Y^{-1}\mathcal{F} = \mathcal{H}$.
D'après le premier paragraphe de la démonstration, cela signifie que le triplet
appartient à l'image essentielle. Nous sommes ainsi ramenés
au cas décrit au paragraphe suivant.

\medskip\noindent
Supposons $f$ fini localement libre et surjectif.
Soit $(\mathcal{G}, \mathcal{H}, \alpha)$ un triplet.
Considérons dans ce cas le triplet
$$
(\mathcal{G}_1, \mathcal{H}_1, \alpha_1) =
(f_{small}^{-1}f_{small, *}\mathcal{G},
f_{big, fppf, *}f_{big, fppf}^{-1}\mathcal{H}, \alpha_1)
$$
où $\alpha_1$ est défini à partir des identifications
\begin{align*}
a_X^{-1}f_{small}^{-1}f_{small, *}\mathcal{G}
& =
f_{big, fppf}^{-1}a_Y^{-1}f_{small, *}\mathcal{G} \\
& =
f_{big, fppf}^{-1}f_{big, fppf, *}a_X^{-1}\mathcal{G} \\
& \to
f_{big, fppf}^{-1}f_{big, fppf, *}f_{big, fppf}^{-1}\mathcal{H}
\end{align*}
où la troisième égalité vient du lemme \ref{etale-cohomology-lemma-proper-push-pull-fppf-etale}
et la flèche est donnée par $\alpha$.
Ce triplet appartient à l'image de notre foncteur, car
$\mathcal{F}_1 = f_{small, *}\mathcal{F}$ constitue une solution
(pour le voir, employer de nouveau le lemme \ref{etale-cohomology-lemma-proper-push-pull-fppf-etale};
détails omis). Il existe un morphisme canonique de triplets
$$
(\mathcal{G}, \mathcal{H}, \alpha)
\to
(\mathcal{G}_1, \mathcal{H}_1, \alpha_1)
$$
qui emploie l'unité $\text{id} \to f_{big, fppf, *}f_{big, fppf}^{-1}$
sur la seconde composante (il suffit de prescrire les morphismes sur la
seconde composante, d'après le premier paragraphe de la démonstration). Puisque
$\{f : X \to Y\}$ est un recouvrement fppf, le morphisme
$\mathcal{H} \to \mathcal{H}_1$ est injectif (détails omis).
Posons
$$
\mathcal{G}_2 = \mathcal{G}_1 \amalg_\mathcal{G} \mathcal{G}_1\quad
\mathcal{H}_2 = \mathcal{H}_1 \amalg_\mathcal{H} \mathcal{H}_1
$$
et soit $\alpha_2$ l'isomorphisme induit (les foncteurs d'image inverse
sont exacts, donc cette construction a un sens). Alors $\mathcal{H}$ est
l'égalisateur des deux morphismes $\mathcal{H}_1 \to \mathcal{H}_2$.
En répétant les arguments ci-dessus pour le triplet
$(\mathcal{G}_2, \mathcal{H}_2, \alpha_2)$
nous trouvons un morphisme injectif de triplets
$$
(\mathcal{G}_2, \mathcal{H}_2, \alpha_2)
\to
(\mathcal{G}_3, \mathcal{H}_3, \alpha_3)
$$
tel que ce dernier triplet appartienne à l'image de notre foncteur.
Disons qu'il correspond à $\mathcal{F}_3$ dans $\Sh(Y_\etale)$.
Par pleine fidélité, nous obtenons deux morphismes
$\mathcal{F}_1 \to \mathcal{F}_3$, et nous pouvons prendre pour
$\mathcal{F}$ leur égalisateur.
Par exactitude des foncteurs d'image inverse considérés, nous
trouvons $a_Y^{-1}\mathcal{F} = \mathcal{H}$, comme voulu.
\end{proof}

\begin{lemma}
\label{etale-cohomology-lemma-compare-fppf-etale}
Considérons le morphisme de comparaison
$\epsilon : (\Sch/S)_{fppf} \to (\Sch/S)_\etale$.
Soit $\mathcal{P}$ la classe des morphismes finis de schémas.
Pour $X$ dans $(\Sch/S)_\etale$, notons
$\mathcal{A}'_X \subset \textit{Ab}((\Sch/X)_\etale)$
la sous-catégorie pleine formée des faisceaux de la forme
$\pi_X^{-1}\mathcal{F}$ avec $\mathcal{F}$ dans $\textit{Ab}(X_\etale)$.
Alors les propriétés de Cohomologie sur les sites
(\ref{sites-cohomology-item-base-change-P}),
(\ref{sites-cohomology-item-restriction-A}),
(\ref{sites-cohomology-item-A-sheaf}),
(\ref{sites-cohomology-item-A-and-P}) et
(\ref{sites-cohomology-item-refine-tau-by-P})
énumérées dans la situation de Cohomologie sur les sites
\ref{sites-cohomology-situation-compare} sont satisfaites.
\end{lemma}

\begin{proof}
Montrons d'abord que $\mathcal{A}'_X \subset \textit{Ab}((\Sch/X)_\etale)$
est une sous-catégorie de Serre faible en vérifiant les conditions (1), (2), (3) et (4)
du lemme \ref{homology-lemma-characterize-weak-serre-subcategory} d'Homology.
Les parties (1), (2), (3) sont immédiates, car $\pi_X^{-1}$ est exact et
pleinement fidèle, par exemple par le lemme \ref{etale-cohomology-lemma-cohomological-descent-etale}. Si
$0 \to \pi_X^{-1}\mathcal{F} \to \mathcal{G} \to \pi_X^{-1}\mathcal{F}' \to 0$
est une suite exacte courte dans $\textit{Ab}((\Sch/X)_\etale)$,
alors $0 \to \mathcal{F} \to \pi_{X, *}\mathcal{G} \to \mathcal{F}' \to 0$
est exacte par le lemme \ref{etale-cohomology-lemma-cohomological-descent-etale}.
Ainsi, $\mathcal{G} = \pi_X^{-1}\pi_{X, *}\mathcal{G}$ appartient à
$\mathcal{A}'_X$, ce qui vérifie la dernière condition.

\medskip\noindent
La propriété (\ref{sites-cohomology-item-base-change-P}) de Cohomologie sur les sites vaut
par l'existence des produits fibrés de schémas
et le fait que le changement de base d'un morphisme fini de
schémas est un morphisme fini de schémas; voir le
lemme \ref{morphisms-lemma-base-change-finite} de Morphismes.

\medskip\noindent
La propriété (\ref{sites-cohomology-item-restriction-A}) de Cohomologie sur les sites
découle du diagramme commutatif (3) du
lemme \ref{topologies-lemma-morphism-big-small-etale} de Topologies.

\medskip\noindent
La propriété (\ref{sites-cohomology-item-A-sheaf}) de Cohomologie sur les sites est le
lemme \ref{etale-cohomology-lemma-describe-pullback-pi-fppf}.

\medskip\noindent
La propriété (\ref{sites-cohomology-item-A-and-P}) de Cohomologie sur les sites vaut par la
partie (4) du lemme \ref{etale-cohomology-lemma-compare-higher-direct-image-proper}.

\medskip\noindent
La propriété (\ref{sites-cohomology-item-refine-tau-by-P}) de Cohomologie sur les sites
est impliquée par le
lemme de Compléments sur les morphismes
\ref{more-morphisms-lemma-dominate-fppf-etale-locally}.
\end{proof}

\begin{lemma}
\label{etale-cohomology-lemma-V-C-all-n-etale-fppf}
Avec les notations ci-dessus.
\begin{enumerate}
\item Pour $X \in \Ob((\Sch/S)_{fppf})$ et un faisceau abélien $\mathcal{F}$
sur $X_\etale$, nous avons
$\epsilon_{X, *}a_X^{-1}\mathcal{F} = \pi_X^{-1}\mathcal{F}$
et $R^i\epsilon_{X, *}(a_X^{-1}\mathcal{F}) = 0$ pour $i > 0$.
\item Pour un morphisme fini $f : X \to Y$ dans $(\Sch/S)_{fppf}$
et un faisceau abélien $\mathcal{F}$ sur $X$, nous avons
$a_Y^{-1}(R^if_{small, *}\mathcal{F}) =
R^if_{big, fppf, *}(a_X^{-1}\mathcal{F})$
pour tout $i$.
\item Pour un schéma $X$ et $K$ dans $D^+(X_\etale)$, le morphisme
$\pi_X^{-1}K \to R\epsilon_{X, *}(a_X^{-1}K)$ est un isomorphisme.
\item Pour un morphisme fini $f : X \to Y$ de schémas
et $K$ dans $D^+(X_\etale)$, nous avons
$a_Y^{-1}(Rf_{small, *}K) = Rf_{big, fppf, *}(a_X^{-1}K)$.
\item Pour un morphisme propre $f : X \to Y$ de schémas
et $K$ dans $D^+(X_\etale)$ à faisceaux de cohomologie de torsion, nous avons
$a_Y^{-1}(Rf_{small, *}K) = Rf_{big, fppf, *}(a_X^{-1}K)$.
\end{enumerate}
\end{lemma}

\begin{proof}
Par le lemme \ref{etale-cohomology-lemma-compare-fppf-etale}, tous les lemmes de la
section \ref{sites-cohomology-section-compare-general} de Cohomologie sur les sites
s'appliquent à notre situation. Pour traduire les résultats,
remarquons que la catégorie $\mathcal{A}_X$ du
lemme \ref{sites-cohomology-lemma-A} de Cohomologie sur les sites
est l'image essentielle de
$a_X^{-1} : \textit{Ab}(X_\etale) \to \textit{Ab}((\Sch/X)_{fppf})$.

\medskip\noindent
La partie (1) équivaut à $(V_n)$ pour tout $n$, ce qui vaut par le
lemme \ref{sites-cohomology-lemma-V-C-all-n-general} de Cohomologie sur les sites.

\medskip\noindent
La partie (2) découle de l'application de $\epsilon_Y^{-1}$ à la conclusion du
lemme \ref{sites-cohomology-lemma-V-implies-C-general} de Cohomologie sur les sites.

\medskip\noindent
La partie (3) découle du lemme de Cohomologie sur les sites
\ref{sites-cohomology-lemma-V-C-all-n-general}, partie (1),
car $\pi_X^{-1}K$ appartient à $D^+_{\mathcal{A}'_X}((\Sch/X)_\etale)$
et $a_X^{-1} = \epsilon_X^{-1} \circ a_X^{-1}$.

\medskip\noindent
La partie (4) découle du lemme de Cohomologie sur les sites
\ref{sites-cohomology-lemma-V-C-all-n-general}, partie (2),
pour la même raison.

\medskip\noindent
Partie (5). Nous employons
\begin{align*}
R\epsilon_{Y, *}Rf_{big, fppf, *}a_X^{-1}K
& =
Rf_{big, \etale, *}R\epsilon_{X, *}a_X^{-1}K \\
& =
Rf_{big, \etale, *}\pi_X^{-1}K \\
& =
\pi_Y^{-1}Rf_{small, *}K \\
& =
R\epsilon_{Y, *} a_Y^{-1}Rf_{small, *}K
\end{align*}
La première égalité vient du diagramme commutatif du
lemme \ref{etale-cohomology-lemma-push-pull-fppf-etale}
et du lemme de Cohomologie sur les sites
\ref{sites-cohomology-lemma-derived-pushforward-composition}.
La deuxième égalité est (3). La troisième est la partie (2) du
lemme \ref{etale-cohomology-lemma-compare-higher-direct-image-proper}.
La quatrième est de nouveau (3). Le morphisme de changement de base
$a_Y^{-1}(Rf_{small, *}K) \to Rf_{big, fppf, *}(a_X^{-1}K)$
induit donc un isomorphisme
$$
R\epsilon_{Y, *}a_Y^{-1}Rf_{small, *}K \to
R\epsilon_{Y, *}Rf_{big, fppf, *}a_X^{-1}K
$$
La remarque suivante achève la démonstration : un morphisme
$\alpha : a_Y^{-1}L \to M$, avec $L$ dans $D^+(Y_\etale)$
et $M$ dans $D^+((\Sch/Y)_{fppf})$, tel que $R\epsilon_{Y, *}\alpha$
soit un isomorphisme, est un isomorphisme. En effet,
nous montrons par récurrence sur $i$ que $H^i(\alpha)$ est un isomorphisme.
C'est vrai pour tout $i$ assez petit.
Si cela vaut pour $i \leq i_0$, nous voyons que
$R^j\epsilon_{Y, *}H^i(M) = 0$ pour $j > 0$ et $i \leq i_0$
par (1), puisque $H^i(M) = a_Y^{-1}H^i(L)$ dans cet intervalle.
Ainsi, $\epsilon_{Y, *}H^{i_0 + 1}(M) = H^{i_0 + 1}(R\epsilon_{Y, *}M)$
par un argument de suite spectrale.
Ainsi, $\epsilon_{Y, *}H^{i_0 + 1}(M) = \pi_Y^{-1}H^{i_0 + 1}(L) =
\epsilon_{Y, *}a_Y^{-1}H^{i_0 + 1}(L)$.
Cela implique que $H^{i_0 + 1}(\alpha)$ est un isomorphisme
(car $\epsilon_{Y, *}$ reflète les isomorphismes, puisqu'il est le
foncteur identité sur les préfaisceaux sous-jacents), comme voulu.
\end{proof}

\begin{lemma}
\label{etale-cohomology-lemma-cohomological-descent-etale-fppf}
Soit $X$ un schéma. Pour $K \in D^+(X_\etale)$, le morphisme
$$
K \longrightarrow Ra_{X, *}a_X^{-1}K
$$
est un isomorphisme, où $a_X : \Sh((\Sch/X)_{fppf}) \to \Sh(X_\etale)$
est comme ci-dessus.
\end{lemma}

\begin{proof}
Nous nous ramenons d'abord au cas où
$K$ est donné par un unique faisceau abélien. Représentons en effet $K$
par un complexe borné inférieurement $\mathcal{F}^\bullet$. Le cas d'un
faisceau montre que
$\mathcal{F}^n = a_{X, *} a_X^{-1} \mathcal{F}^n$
et que les faisceaux $R^qa_{X, *}a_X^{-1}\mathcal{F}^n$
sont nuls pour $q > 0$. Le lemme d'acyclicité de Leray
(lemme \ref{derived-lemma-leray-acyclicity} de Catégories dérivées),
appliqué à $a_X^{-1}\mathcal{F}^\bullet$
et au foncteur $a_{X, *}$, permet de conclure. Supposons désormais $K = \mathcal{F}$.

\medskip\noindent
Par le lemme \ref{etale-cohomology-lemma-describe-pullback-pi-fppf}, nous avons
$a_{X, *}a_X^{-1}\mathcal{F} = \mathcal{F}$. Il suffit donc de montrer que
$R^qa_{X, *}a_X^{-1}\mathcal{F} = 0$ pour $q > 0$.
Pour cela, nous pouvons employer $a_X = \epsilon_X \circ \pi_X$ et
la suite spectrale de Leray
(lemme \ref{sites-cohomology-lemma-relative-Leray} de Cohomologie sur les sites).
Par le lemme \ref{etale-cohomology-lemma-V-C-all-n-etale-fppf},
nous avons $R^i\epsilon_{X, *}(a_X^{-1}\mathcal{F}) = 0$ pour $i > 0$
et
$\epsilon_{X, *}a_X^{-1}\mathcal{F} = \pi_X^{-1}\mathcal{F}$.
Par le lemme \ref{etale-cohomology-lemma-cohomological-descent-etale}, nous avons
$R^j\pi_{X, *}(\pi_X^{-1}\mathcal{F}) = 0$ pour $j > 0$.
Cela achève la démonstration.
\end{proof}

\begin{lemma}
\label{etale-cohomology-lemma-compare-cohomology-etale-fppf}
Pour un schéma $X$ et $a_X : \Sh((\Sch/X)_{fppf}) \to \Sh(X_\etale)$
comme ci-dessus :
\begin{enumerate}
\item $H^q(X_\etale, \mathcal{F}) = H^q_{fppf}(X, a_X^{-1}\mathcal{F})$
pour un faisceau abélien $\mathcal{F}$ sur $X_\etale$,
\item $H^q(X_\etale, K) = H^q_{fppf}(X, a_X^{-1}K)$ pour $K \in D^+(X_\etale)$.
\end{enumerate}
Exemple : si $A$ est un groupe abélien, alors
$H^q_\etale(X, \underline{A}) = H^q_{fppf}(X, \underline{A})$.
\end{lemma}

\begin{proof}
Cela découle du lemme \ref{etale-cohomology-lemma-cohomological-descent-etale-fppf}
par la remarque \ref{sites-cohomology-remark-before-Leray} de Cohomologie sur les sites.
\end{proof}







\section{Comparaison des topologies fppf et étale : modules}
\label{etale-cohomology-section-fppf-etale-modules}

\noindent
Nous poursuivons la discussion de la section \ref{etale-cohomology-section-fppf-etale}, mais nous
examinons ici brièvement ce qui se passe pour les faisceaux de modules.

\medskip\noindent
Soit $S$ un schéma. Les morphismes de sites $\epsilon_S$, $\pi_S$ et
leur composé $a_S$, introduits dans la section \ref{etale-cohomology-section-fppf-etale},
s'étendent naturellement en morphismes de sites annelés. Le premier
s'écrit
$$
\epsilon_S :
((\Sch/S)_{fppf}, \mathcal{O})
\longrightarrow
((\Sch/S)_\etale, \mathcal{O})
$$
Remarquons que nous pouvons employer le même symbole pour le faisceau structural, car
ces faisceaux ont le même préfaisceau sous-jacent. Le deuxième est
$$
\pi_S :
((\Sch/S)_\etale, \mathcal{O})
\longrightarrow
(S_\etale, \mathcal{O}_S)
$$
Le troisième est le morphisme
$$
a_S :
((\Sch/S)_{fppf}, \mathcal{O})
\longrightarrow
(S_\etale, \mathcal{O}_S)
$$
Nous savons déjà que la catégorie des modules quasi-cohérents sur le
schéma $S$ est la même que celle des modules quasi-cohérents
sur $(S_\etale, \mathcal{O}_S)$; voir la
proposition \ref{descent-proposition-equivalence-quasi-coherent} de Descente.
Comme nous voulons énoncer une comparaison entre les cohomologies
étale et fppf, nous considérerons dans le reste de cette
section les faisceaux quasi-cohérents du point de vue du
petit site étale.
Rappelons ce que nous savons déjà des modules quasi-cohérents
sur ces sites.

\begin{lemma}
\label{etale-cohomology-lemma-review-quasi-coherent}
Soit $S$ un schéma. Soit $\mathcal{F}$ un
$\mathcal{O}_S$-module quasi-cohérent sur $S_\etale$.
\begin{enumerate}
\item La règle
$$
\mathcal{F}^a : (\Sch/S)_\etale \longrightarrow \textit{Ab},\quad
(f : T \to S) \longmapsto \Gamma(T, f_{small}^*\mathcal{F})
$$
satisfait la condition de faisceau pour les recouvrements fppf, et a fortiori étales,
\item $\mathcal{F}^a = \pi_S^*\mathcal{F}$ sur $(\Sch/S)_\etale$,
\item $\mathcal{F}^a = a_S^*\mathcal{F}$ sur $(\Sch/S)_{fppf}$,
\item la règle $\mathcal{F} \mapsto \mathcal{F}^a$ définit
une équivalence entre les $\mathcal{O}_S$-modules quasi-cohérents
et les modules quasi-cohérents sur
$((\Sch/S)_\etale, \mathcal{O})$,
\item la règle $\mathcal{F} \mapsto \mathcal{F}^a$ définit
une équivalence entre les $\mathcal{O}_S$-modules quasi-cohérents
et les modules quasi-cohérents sur
$((\Sch/S)_{fppf}, \mathcal{O})$,
\item nous avons $\epsilon_{S, *}a_S^*\mathcal{F} = \pi_S^*\mathcal{F}$
et $a_{S, *}a_S^*\mathcal{F} = \mathcal{F}$,
\item nous avons $R^i\epsilon_{S, *}(a_S^*\mathcal{F}) = 0$
et $R^ia_{S, *}(a_S^*\mathcal{F}) = 0$ pour $i > 0$.
\end{enumerate}
\end{lemma}

\begin{proof}
Nous invitons vivement le lecteur à trouver sa propre démonstration de ces résultats
à partir des matériaux des
sections \ref{descent-section-quasi-coherent-sheaves},
\ref{descent-section-quasi-coherent-cohomology} et
\ref{descent-section-quasi-coherent-sheaves-bis}.

\medskip\noindent
Expliquons d'abord pourquoi la notation de ce lemme est compatible avec notre
emploi antérieur de la notation $\mathcal{F}^a$ dans les
sections \ref{etale-cohomology-section-quasi-coherent} et
\ref{etale-cohomology-section-cohomology-quasi-coherent},
ainsi que dans la
section \ref{descent-section-quasi-coherent-sheaves} de Descente.
En effet, la
proposition \ref{descent-proposition-equivalence-quasi-coherent} de Descente
montre qu'il existe un module quasi-cohérent
$\mathcal{F}_0$ sur le schéma $S$ (autrement dit sur le petit
site de Zariski) tel que $\mathcal{F}$ soit la restriction de la
règle
$$
\mathcal{F}_0^a : (\Sch/S)_\etale \longrightarrow \textit{Ab},\quad
(f : T \to S) \longmapsto \Gamma(T, f^*\mathcal{F})
$$
à la sous-catégorie $S_\etale \subset (\Sch/S)_\etale$,
où $f^*$ désigne ici l'image inverse usuelle des faisceaux de modules sur les schémas.
Puisque $\mathcal{F}_0^a$ est l'image inverse par le morphisme de
sites annelés
$$
((\Sch/S)_\etale, \mathcal{O}) \longrightarrow (S_{Zar}, \mathcal{O}_{S_{Zar}})
$$
par la remarque \ref{descent-remark-change-topologies-ringed-sites} de Descente,
il découle immédiatement (par composition des images inverses) que
$\mathcal{F}^a = \mathcal{F}_0^a$. Cela démontre la propriété de faisceau,
même pour les recouvrements fpqc, par le
lemme \ref{descent-lemma-sheaf-condition-holds} de Descente (voir aussi la
proposition \ref{etale-cohomology-proposition-quasi-coherent-sheaf-fpqc}).
Les assertions (2) et (3) découlent alors
de nouveau de la remarque \ref{descent-remark-change-topologies-ringed-sites} de Descente,
et (4) et (5) de la
proposition \ref{descent-proposition-equivalence-quasi-coherent} de Descente
(voir aussi le méta-résultat du
théorème \ref{etale-cohomology-theorem-quasi-coherent}).

\medskip\noindent
La partie (6) découle immédiatement de la description du faisceau
$\mathcal{F}^a = \pi_S^*\mathcal{F} = a_S^*\mathcal{F}$.

\medskip\noindent
Pour tout faisceau abélien $\mathcal{H}$ sur $(\Sch/S)_{fppf}$, l'
image directe supérieure $R^p\epsilon_{S, *}\mathcal{H}$ est le faisceau
associé au préfaisceau $U \mapsto H^p_{fppf}(U, \mathcal{H})$
sur $(\Sch/S)_\etale$. Voir le
lemme \ref{sites-cohomology-lemma-higher-direct-images} de Cohomologie sur les sites.
Ainsi, pour démontrer
$R^p\epsilon_{S, *}a_S^*\mathcal{F} = R^p\epsilon_{S, *}\mathcal{F}^a = 0$
pour $p > 0$, il suffit de montrer que tout schéma $U$ sur $S$
possède un recouvrement étale $\{U_i \to U\}_{i \in I}$ tel que
$H^p_{fppf}(U_i, \mathcal{F}^a) = 0$ pour $p > 0$.
En prenant un recouvrement ouvert par des affines, l'
annulation requise découle de la comparaison avec la cohomologie usuelle
(proposition \ref{descent-proposition-same-cohomology-quasi-coherent} de Descente
ou
théorème \ref{etale-cohomology-theorem-zariski-fpqc-quasi-coherent}),
ainsi que de l'annulation de la cohomologie des faisceaux quasi-cohérents
sur les schémas affines donnée par le lemme de Cohomologie des schémas
\ref{coherent-lemma-quasi-coherent-affine-cohomology-zero}.

\medskip\noindent
Pour montrer que $R^pa_{S, *}a_S^{-1}\mathcal{F} = R^pa_{S, *}\mathcal{F}^a = 0$
pour $p > 0$, nous raisonnons exactement de la même manière. Cela achève la démonstration.
\end{proof}

\begin{lemma}
\label{etale-cohomology-lemma-cohomological-descent-etale-fppf-modules}
Soit $S$ un schéma. Étant donné $\mathcal{F}$, faisceau quasi-cohérent
de $\mathcal{O}_S$-modules sur $S_\etale$, les morphismes
$$
\pi_S^*\mathcal{F} \longrightarrow R\epsilon_{S, *}(a_S^*\mathcal{F})
\quad\text{et}\quad
\mathcal{F} \longrightarrow Ra_{S, *}(a_S^*\mathcal{F})
$$
sont des isomorphismes, où
$a_S : \Sh((\Sch/S)_{fppf}) \to \Sh(S_\etale)$ est comme ci-dessus.
\end{lemma}

\begin{proof}
C'est une conséquence immédiate des
parties (6) et (7) du
lemme \ref{etale-cohomology-lemma-review-quasi-coherent}.
\end{proof}

\begin{lemma}
\label{etale-cohomology-lemma-cohomological-descent-complex-modules}
Soit $S = \Spec(A)$ un schéma affine. Soit $M^\bullet$ un complexe
de $A$-modules. Considérons le complexe $\mathcal{F}^\bullet$ de
préfaisceaux de $\mathcal{O}$-modules sur
$(\textit{Aff}/S)_{fppf}$ donné par la règle
$$
(U/S) = (\Spec(B)/\Spec(A)) \longmapsto M^\bullet \otimes_A B
$$
C'est alors un complexe de modules, et le morphisme canonique
$$
M^\bullet \longrightarrow
R\Gamma((\textit{Aff}/S)_{fppf}, \mathcal{F}^\bullet)
$$
est un quasi-isomorphisme.
\end{lemma}

\begin{proof}
Chaque $\mathcal{F}^n$ est un faisceau de modules, car il coïncide avec la
restriction du module $\mathcal{G}^n = (\widetilde{M}^n)^a$ du
lemme \ref{etale-cohomology-lemma-review-quasi-coherent} à
$(\textit{Aff}/S)_{fppf} \subset (\Sch/S)_{fppf}$.
Puisque cette inclusion définit une équivalence de topos annelés
(lemme \ref{topologies-lemma-affine-big-site-fppf} de Topologies),
nous avons
$$
R\Gamma((\textit{Aff}/S)_{fppf}, \mathcal{F}^\bullet) =
R\Gamma((\Sch/S)_{fppf}, \mathcal{G}^\bullet)
$$
Remarquons que $M^\bullet = R\Gamma(S, \widetilde{M}^\bullet)$,
par exemple par le lemme de Catégories dérivées des schémas
\ref{perfect-lemma-affine-compare-bounded}. Nous cherchons donc
à montrer que le morphisme de comparaison
$$
R\Gamma(S, \widetilde{M}^\bullet)
\longrightarrow
R\Gamma((\Sch/S)_{fppf}, (\widetilde{M}^\bullet)^a)
$$
est un isomorphisme. Si $M^\bullet$ est borné inférieurement, cela
vaut par la proposition de Descente
\ref{descent-proposition-same-cohomology-quasi-coherent}
et la première suite spectrale du lemme de Catégories dérivées
\ref{derived-lemma-two-ss-complex-functor}.
Dans le cas général, écrivons $M^\bullet = \lim M_n^\bullet$
avec $M_n^\bullet = \tau_{\geq -n}M^\bullet$. Le système
$M_n^p$ est alors stationnaire de valeur $M^p$. Nous affirmons que
$$
(\widetilde{M}^\bullet)^a = R\lim (\widetilde{M}_n^\bullet)^a
$$
En effet, il suffit de montrer que le morphisme naturel de gauche à droite
induit un isomorphisme en cohomologie sur tout objet affine
$U = \Spec(B)$ de $(\Sch/S)_{fppf}$. Pour $i \in \mathbf{Z}$ et
$n > |i|$, nous avons
$$
H^i(U, (\widetilde{M}_n^\bullet)^a) =
H^i(\tau_{\geq -n}M^\bullet \otimes_A B) =
H^i(M^\bullet \otimes_A B)
$$
La première égalité vaut par le cas borné inférieurement traité ci-dessus.
Le lemme de Cohomologie sur les sites
\ref{sites-cohomology-lemma-RGamma-commutes-with-Rlim}
montre donc l'affirmation. Enfin, nous obtenons
\begin{align*}
R\Gamma((\Sch/S)_{fppf}, (\widetilde{M}^\bullet)^a)
& =
R\Gamma((\Sch/S)_{fppf}, R\lim (\widetilde{M}_n^\bullet)^a) \\
& =
R\lim R\Gamma((\Sch/S)_{fppf}, (\widetilde{M}_n^\bullet)^a) \\
& =
R\lim M_n^\bullet \\
& =
M^\bullet
\end{align*}
comme voulu. La deuxième égalité vaut parce que $R\lim$ commute à
$R\Gamma$; voir le lemme de Cohomologie sur les sites
\ref{sites-cohomology-lemma-RGamma-commutes-with-Rlim}.
\end{proof}










\section{Comparaison des topologies ph et étale}
\label{etale-cohomology-section-ph-etale}

\noindent
Un modèle pour cette section est la section consacrée à la comparaison de la
topologie usuelle et de la topologie qc sur les espaces topologiques
localement compacts, étudiée dans la section de Cohomologie sur les sites
\ref{sites-cohomology-section-cohomology-LC}.
Nous rappelons d'abord certains résultats des
sections de Topologies
\ref{topologies-section-change-topologies} et
\ref{topologies-section-etale}.

\medskip\noindent
Soit $S$ un schéma et soit $(\Sch/S)_{ph}$ un site ph.
Sur la même catégorie sous-jacente, nous avons une seconde topologie,
à savoir la topologie étale, et donc un second site
$(\Sch/S)_\etale$. Le foncteur identité
$(\Sch/S)_\etale \to (\Sch/S)_{ph}$ est continu
(par le lemme \ref{more-morphisms-lemma-fppf-ph} de Compléments sur les morphismes
et le lemme de Topologies
\ref{topologies-lemma-zariski-etale-smooth-syntomic-fppf})
et définit un morphisme de sites
$$
\epsilon_S : (\Sch/S)_{ph} \longrightarrow (\Sch/S)_\etale
$$
Voir la section \ref{sites-cohomology-section-compare} de Cohomologie sur les sites.
Remarquons que $\epsilon_{S, *}$ est le foncteur identité sur les
préfaisceaux sous-jacents et que $\epsilon_S^{-1}$ associe à un faisceau étale son
faisceau ph associé.
Soit $S_\etale$ le petit site étale.
Il existe un morphisme de sites
$$
\pi_S : (\Sch/S)_\etale \longrightarrow S_\etale
$$
induit par le foncteur continu
$S_\etale \to (\Sch/S)_\etale$, $U \mapsto U$.
En effet, $S_\etale$ possède des produits fibrés et un objet final, le
foncteur ci-dessus y commute, et la
proposition \ref{sites-proposition-get-morphism} de Sites s'applique.

\begin{lemma}
\label{etale-cohomology-lemma-describe-pullback-pi-ph}
Avec les notations ci-dessus.
Soit $\mathcal{F}$ un faisceau sur $S_\etale$. La règle
$$
(\Sch/S)_{ph} \longrightarrow \textit{Ensembles},\quad
(f : X \to S) \longmapsto \Gamma(X, f_{small}^{-1}\mathcal{F})
$$
est un faisceau ph, et a fortiori un faisceau sur $(\Sch/S)_\etale$.
En fait, ce faisceau est égal à
$\pi_S^{-1}\mathcal{F}$ sur $(\Sch/S)_\etale$ et à
$\epsilon_S^{-1}\pi_S^{-1}\mathcal{F}$ sur $(\Sch/S)_{ph}$.
\end{lemma}

\begin{proof}
L'énoncé relatif à la topologie étale est le contenu
du lemme \ref{etale-cohomology-lemma-describe-pullback}. Pour achever la démonstration, il
suffit de montrer que $\pi_S^{-1}\mathcal{F}$ est un faisceau pour la topologie
ph. Par le lemme \ref{topologies-lemma-characterize-sheaf} de Topologies,
il suffit de montrer que, pour tout morphisme propre surjectif
$V \to U$ de schémas sur $S$, nous avons un diagramme égalisateur
$$
\xymatrix{
(\pi_S^{-1}\mathcal{F})(U) \ar[r] &
(\pi_S^{-1}\mathcal{F})(V) \ar@<1ex>[r] \ar@<-1ex>[r] &
(\pi_S^{-1}\mathcal{F})(V \times_U V)
}
$$
Posons $\mathcal{G} = \pi_S^{-1}\mathcal{F}|_{U_\etale}$.
Considérons le diagramme commutatif
$$
\xymatrix{
V \times_U V \ar[r] \ar[rd]_g \ar[d] & V \ar[d]^f \\
V \ar[r]^f & U
}
$$
Nous avons
$$
(\pi_S^{-1}\mathcal{F})(V) = \Gamma(V, f^{-1}\mathcal{G}) =
\Gamma(U, f_*f^{-1}\mathcal{G})
$$
où nous employons $f_*$ et $f^{-1}$ pour désigner les fonctorialités entre
petits sites étales. D'autre part, nous avons
$$
(\pi_S^{-1}\mathcal{F})(V \times_U V) =
\Gamma(V \times_U V, g^{-1}\mathcal{G}) =
\Gamma(U, g_*g^{-1}\mathcal{G})
$$
Les deux morphismes du diagramme égalisateur proviennent des deux morphismes
$$
f_*f^{-1}\mathcal{G} \longrightarrow g_*g^{-1}\mathcal{G}
$$
Il suffit donc de démontrer que $\mathcal{G}$ est
l'égalisateur de ces deux morphismes de faisceaux.
Soit $\overline{u}$ un point géométrique de $U$. Posons
$\Omega = \mathcal{G}_{\overline{u}}$.
En prenant les fibres en $\overline{u}$ par le
lemme \ref{etale-cohomology-lemma-proper-pushforward-stalk},
nous obtenons les deux morphismes
$$
H^0(V_{\overline{u}}, \underline{\Omega}) \longrightarrow
H^0((V \times_U V)_{\overline{u}}, \underline{\Omega}) =
H^0(V_{\overline{u}} \times_{\overline{u}} V_{\overline{u}},
\underline{\Omega})
$$
où $\underline{\Omega}$ désigne le faisceau constant de valeur
$\Omega$. Ces morphismes sont bien sûr les images inverses par les projections.
Il est alors clair que les sections provenant de l'image inverse
par la projection sur le premier facteur sont constantes sur les fibres de
la première projection, tandis que celles provenant de l'image inverse
par la projection sur le second facteur sont constantes sur les fibres de
la seconde projection. Les sections dans l'intersection des images
de ces morphismes d'image inverse sont constantes sur tout
$V_{\overline{u}} \times_{\overline{u}} V_{\overline{u}}$, c'est-à-dire
qu'elles proviennent d'éléments de $\Omega$, comme voulu.
\end{proof}

\noindent
Dans la situation du lemme \ref{etale-cohomology-lemma-describe-pullback-pi-ph},
la composée de $\epsilon_S$ et $\pi_S$, ainsi que l'égalité,
déterminent un morphisme de sites
$$
a_S : (\Sch/S)_{ph} \longrightarrow S_\etale
$$

\begin{lemma}
\label{etale-cohomology-lemma-push-pull-ph-etale}
Avec les notations ci-dessus.
Soit $f : X \to Y$ un morphisme de $(\Sch/S)_{ph}$.
Il existe alors des diagrammes commutatifs de topos
$$
\xymatrix{
\Sh((\Sch/X)_{ph}) \ar[rr]_{f_{big, ph}} \ar[d]_{\epsilon_X} & &
\Sh((\Sch/Y)_{ph}) \ar[d]^{\epsilon_Y} \\
\Sh((\Sch/X)_\etale) \ar[rr]^{f_{big, \etale}} & &
\Sh((\Sch/Y)_\etale)
}
$$
et
$$
\xymatrix{
\Sh((\Sch/X)_{ph}) \ar[rr]_{f_{big, ph}} \ar[d]_{a_X} & &
\Sh((\Sch/Y)_{ph}) \ar[d]^{a_Y} \\
\Sh(X_\etale) \ar[rr]^{f_{small}} & &
\Sh(Y_\etale)
}
$$
avec $a_X = \pi_X \circ \epsilon_X$ et $a_Y = \pi_X \circ \epsilon_X$.
\end{lemma}

\begin{proof}
La commutativité des diagrammes découle de la discussion de la
section \ref{topologies-section-change-topologies} de Topologies.
\end{proof}

\begin{lemma}
\label{etale-cohomology-lemma-proper-push-pull-ph-etale}
Dans le lemme \ref{etale-cohomology-lemma-push-pull-ph-etale}, si $f$ est propre, alors nous avons
$a_Y^{-1} \circ f_{small, *} = f_{big, ph, *} \circ a_X^{-1}$.
\end{lemma}

\begin{proof}
On peut le démontrer en répétant la démonstration de la
partie (1) du lemme \ref{etale-cohomology-lemma-compare-higher-direct-image-proper};
nous allons plutôt l'en déduire.
Comme $\epsilon_{Y, *}$ est le foncteur identité sur les préfaisceaux sous-jacents,
il reflète les isomorphismes. La description
du lemme \ref{etale-cohomology-lemma-describe-pullback-pi-ph}
montre que $\epsilon_{Y, *} \circ a_Y^{-1} = \pi_Y^{-1}$,
et de même pour $X$. Pour montrer que le morphisme canonique
$a_Y^{-1}f_{small, *}\mathcal{F} \to f_{big, ph, *}a_X^{-1}\mathcal{F}$
est un isomorphisme, il suffit de montrer que
\begin{align*}
\pi_Y^{-1}f_{small, *}\mathcal{F}
& =
\epsilon_{Y, *}a_Y^{-1}f_{small, *}\mathcal{F} \\
& \to 
\epsilon_{Y, *}f_{big, ph, *}a_X^{-1}\mathcal{F} \\
& =
f_{big, \etale, *} \epsilon_{X, *}a_X^{-1}\mathcal{F} \\
& =
f_{big, \etale, *}\pi_X^{-1}\mathcal{F}
\end{align*}
est un isomorphisme. C'est la partie
(1) du lemme \ref{etale-cohomology-lemma-compare-higher-direct-image-proper}.
\end{proof}

\begin{lemma}
\label{etale-cohomology-lemma-compare-ph-etale}
Considérons le morphisme de comparaison
$\epsilon : (\Sch/S)_{ph} \to (\Sch/S)_\etale$.
Soit $\mathcal{P}$ la classe des morphismes propres de schémas.
Pour $X \in (\Sch/S)_\etale$, notons
$\mathcal{A}'_X \subset \textit{Ab}((\Sch/X)_\etale)$
la sous-catégorie pleine formée des faisceaux de la forme
$\pi_X^{-1}\mathcal{F}$, où $\mathcal{F}$ est un
faisceau abélien de torsion sur $X_\etale$.
Alors les propriétés de Cohomologie sur les sites
(\ref{sites-cohomology-item-base-change-P}),
(\ref{sites-cohomology-item-restriction-A}),
(\ref{sites-cohomology-item-A-sheaf}),
(\ref{sites-cohomology-item-A-and-P}) et
(\ref{sites-cohomology-item-refine-tau-by-P})
de la situation \ref{sites-cohomology-situation-compare} de Cohomologie sur les sites
sont satisfaites.
\end{lemma}

\begin{proof}
Montrons d'abord que $\mathcal{A}'_X \subset \textit{Ab}((\Sch/X)_\etale)$
est une sous-catégorie de Serre faible en vérifiant les conditions (1), (2), (3) et (4)
du lemme \ref{homology-lemma-characterize-weak-serre-subcategory} d'Homology.
Les parties (1), (2), (3) sont immédiates, car $\pi_X^{-1}$ est exact et
pleinement fidèle, par exemple par le lemme \ref{etale-cohomology-lemma-cohomological-descent-etale}. Si
$0 \to \pi_X^{-1}\mathcal{F} \to \mathcal{G} \to \pi_X^{-1}\mathcal{F}' \to 0$
est une suite exacte courte dans $\textit{Ab}((\Sch/X)_\etale)$,
alors $0 \to \mathcal{F} \to \pi_{X, *}\mathcal{G} \to \mathcal{F}' \to 0$
est exacte par le lemme \ref{etale-cohomology-lemma-cohomological-descent-etale}.
Nous voyons en particulier que $\pi_{X, *}\mathcal{G}$ est un faisceau
abélien de torsion sur $X_\etale$.
Ainsi, $\mathcal{G} = \pi_X^{-1}\pi_{X, *}\mathcal{G}$ appartient à
$\mathcal{A}'_X$, ce qui vérifie la dernière condition.

\medskip\noindent
La propriété (\ref{sites-cohomology-item-base-change-P}) de Cohomologie sur les sites vaut
par l'existence des produits fibrés de schémas
et le fait que le changement de base d'un morphisme propre de
schémas est un morphisme propre de schémas; voir le
lemme \ref{morphisms-lemma-base-change-proper} de Morphismes.

\medskip\noindent
La propriété (\ref{sites-cohomology-item-restriction-A}) de Cohomologie sur les sites
découle du diagramme commutatif (3) du
lemme \ref{topologies-lemma-morphism-big-small-etale} de Topologies.

\medskip\noindent
La propriété (\ref{sites-cohomology-item-A-sheaf}) de Cohomologie sur les sites est le
lemme \ref{etale-cohomology-lemma-describe-pullback-pi-ph}.

\medskip\noindent
La propriété (\ref{sites-cohomology-item-A-and-P}) de Cohomologie sur les sites vaut par la
partie (2) du lemme \ref{etale-cohomology-lemma-compare-higher-direct-image-proper}
et le fait que $R^if_{small, *}\mathcal{F}$
est de torsion si $\mathcal{F}$ est un faisceau
abélien de torsion sur $X_\etale$; voir le lemme \ref{etale-cohomology-lemma-torsion-direct-image}.

\medskip\noindent
La propriété (\ref{sites-cohomology-item-refine-tau-by-P}) de Cohomologie sur les sites
découle du lemme de Compléments sur les morphismes
\ref{more-morphisms-lemma-dominate-fppf-etale-locally}
combiné au fait qu'un morphisme fini est propre
et qu'un morphisme propre surjectif définit un recouvrement ph; voir le
lemme \ref{topologies-lemma-surjective-proper-ph} de Topologies.
\end{proof}

\begin{lemma}
\label{etale-cohomology-lemma-V-C-all-n-etale-ph}
Avec les notations ci-dessus.
\begin{enumerate}
\item Pour $X \in \Ob((\Sch/S)_{ph})$ et un faisceau abélien de torsion $\mathcal{F}$
sur $X_\etale$, nous avons
$\epsilon_{X, *}a_X^{-1}\mathcal{F} = \pi_X^{-1}\mathcal{F}$
et $R^i\epsilon_{X, *}(a_X^{-1}\mathcal{F}) = 0$ pour $i > 0$.
\item Pour un morphisme propre $f : X \to Y$ dans $(\Sch/S)_{ph}$
et un faisceau abélien de torsion $\mathcal{F}$ sur $X$, nous avons
$a_Y^{-1}(R^if_{small, *}\mathcal{F}) =
R^if_{big, ph, *}(a_X^{-1}\mathcal{F})$
pour tout $i$.
\item Pour un schéma $X$ et $K \in D^+(X_\etale)$ à faisceaux de
cohomologie de torsion, le morphisme
$\pi_X^{-1}K \to R\epsilon_{X, *}(a_X^{-1}K)$ est un isomorphisme.
\item Pour un morphisme propre $f : X \to Y$ de schémas
et $K \in D^+(X_\etale)$ à faisceaux de cohomologie de torsion, nous avons
$a_Y^{-1}(Rf_{small, *}K) = Rf_{big, ph, *}(a_X^{-1}K)$.
\end{enumerate}
\end{lemma}

\begin{proof}
Par le lemme \ref{etale-cohomology-lemma-compare-ph-etale}, tous les lemmes de la
section \ref{sites-cohomology-section-compare-general} de Cohomologie sur les sites
s'appliquent à notre situation. Pour traduire les résultats,
remarquons que la catégorie $\mathcal{A}_X$ du
lemme \ref{sites-cohomology-lemma-A} de Cohomologie sur les sites
est la sous-catégorie pleine de $\textit{Ab}((\Sch/X)_{ph})$
formée des faisceaux de la forme $a_X^{-1}\mathcal{F}$,
où $\mathcal{F}$ est un faisceau abélien de torsion sur $X_\etale$.

\medskip\noindent
La partie (1) équivaut à $(V_n)$ pour tout $n$, ce qui vaut par le
lemme \ref{sites-cohomology-lemma-V-C-all-n-general} de Cohomologie sur les sites.

\medskip\noindent
La partie (2) découle de l'application de $\epsilon_Y^{-1}$ à la conclusion du
lemme \ref{sites-cohomology-lemma-V-implies-C-general} de Cohomologie sur les sites.

\medskip\noindent
La partie (3) découle de la partie (1) du lemme de Cohomologie sur les sites
\ref{sites-cohomology-lemma-V-C-all-n-general},
car $\pi_X^{-1}K$ appartient à $D^+_{\mathcal{A}'_X}((\Sch/X)_\etale)$
et $a_X^{-1} = \epsilon_X^{-1} \circ \pi_X^{-1}$.

\medskip\noindent
La partie (4) découle de la partie (2) du lemme de Cohomologie sur les sites
\ref{sites-cohomology-lemma-V-C-all-n-general},
pour la même raison.
\end{proof}

\begin{lemma}
\label{etale-cohomology-lemma-cohomological-descent-etale-ph}
Soit $X$ un schéma. Pour $K \in D^+(X_\etale)$ à faisceaux de cohomologie
de torsion, le morphisme
$$
K \longrightarrow Ra_{X, *}a_X^{-1}K
$$
est un isomorphisme, où $a_X : \Sh((\Sch/X)_{ph}) \to \Sh(X_\etale)$ est comme ci-dessus.
\end{lemma}

\begin{proof}
Nous nous ramenons d'abord au cas où
$K$ est donné par un unique faisceau abélien. Représentons en effet $K$
par un complexe borné inférieurement $\mathcal{F}^\bullet$ de faisceaux
abéliens de torsion. C'est possible par le lemme de Cohomologie sur les sites
\ref{sites-cohomology-lemma-torsion}. Le cas d'un
faisceau montre que
$\mathcal{F}^n = a_{X, *} a_X^{-1} \mathcal{F}^n$
et que les faisceaux $R^qa_{X, *}a_X^{-1}\mathcal{F}^n$
sont nuls pour $q > 0$. Le lemme d'acyclicité de Leray
(lemme \ref{derived-lemma-leray-acyclicity} de Catégories dérivées),
appliqué à $a_X^{-1}\mathcal{F}^\bullet$
et au foncteur $a_{X, *}$, permet de conclure. Supposons désormais $K = \mathcal{F}$,
où $\mathcal{F}$ est un faisceau abélien de torsion.

\medskip\noindent
Par le lemme \ref{etale-cohomology-lemma-describe-pullback-pi-ph}, nous avons
$a_{X, *}a_X^{-1}\mathcal{F} = \mathcal{F}$. Il suffit donc de montrer que
$R^qa_{X, *}a_X^{-1}\mathcal{F} = 0$ pour $q > 0$.
Pour cela, nous pouvons employer $a_X = \pi_X \circ \epsilon_X$ et
la suite spectrale de Leray
(lemme \ref{sites-cohomology-lemma-relative-Leray} de Cohomologie sur les sites).
Par le lemme \ref{etale-cohomology-lemma-V-C-all-n-etale-ph},
nous avons $R^i\epsilon_{X, *}(a_X^{-1}\mathcal{F}) = 0$ pour $i > 0$
et $\epsilon_{X, *}a_X^{-1}\mathcal{F} = \pi_X^{-1}\mathcal{F}$.
Par le lemme \ref{etale-cohomology-lemma-cohomological-descent-etale}, nous avons
$R^j\pi_{X, *}(\pi_X^{-1}\mathcal{F}) = 0$ pour $j > 0$.
Cela achève la démonstration.
\end{proof}

\begin{lemma}
\label{etale-cohomology-lemma-compare-cohomology-etale-ph}
Pour un schéma $X$ et $a_X : \Sh((\Sch/X)_{ph}) \to \Sh(X_\etale)$
comme ci-dessus :
\begin{enumerate}
\item $H^q(X_\etale, \mathcal{F}) = H^q_{ph}(X, a_X^{-1}\mathcal{F})$
pour un faisceau abélien de torsion $\mathcal{F}$ sur $X_\etale$,
\item $H^q(X_\etale, K) = H^q_{ph}(X, a_X^{-1}K)$
pour $K \in D^+(X_\etale)$ à faisceaux de cohomologie de torsion.
\end{enumerate}
Exemple : si $A$ est un groupe abélien de torsion, alors
$H^q_\etale(X, \underline{A}) = H^q_{ph}(X, \underline{A})$.
\end{lemma}

\begin{proof}
Cela découle du lemme \ref{etale-cohomology-lemma-cohomological-descent-etale-ph}
par la remarque \ref{sites-cohomology-remark-before-Leray} de Cohomologie sur les sites.
\end{proof}












\section{Comparaison des topologies h et étale}
\label{etale-cohomology-section-h-etale}

\noindent
Un modèle pour cette section est la section consacrée à la comparaison de la
topologie usuelle et de la topologie qc sur les espaces topologiques
localement compacts, étudiée dans la section de Cohomologie sur les sites
\ref{sites-cohomology-section-cohomology-LC}.
De plus, cette section est presque mot pour mot identique à la
section comparant les topologies ph et étale.
Nous rappelons d'abord certains résultats des
sections de Topologies
\ref{topologies-section-change-topologies} et
\ref{topologies-section-etale}, ainsi que de la
section \ref{flat-section-h} de Compléments sur la platitude.

\medskip\noindent
Soit $S$ un schéma et soit $(\Sch/S)_h$ un site h.
Sur la même catégorie sous-jacente, nous avons une seconde topologie,
à savoir la topologie étale, et donc un second site
$(\Sch/S)_\etale$. Le foncteur identité
$(\Sch/S)_\etale \to (\Sch/S)_h$ est continu
(par le lemme \ref{flat-lemma-zariski-h} de Compléments sur la platitude
et le lemme de Topologies
\ref{topologies-lemma-zariski-etale-smooth-syntomic-fppf})
et définit un morphisme de sites
$$
\epsilon_S : (\Sch/S)_h \longrightarrow (\Sch/S)_\etale
$$
Voir la section \ref{sites-cohomology-section-compare} de Cohomologie sur les sites.
Remarquons que $\epsilon_{S, *}$ est le foncteur identité sur les
préfaisceaux sous-jacents et que $\epsilon_S^{-1}$ associe à un faisceau étale son
faisceau h associé. Soit $S_\etale$ le petit site étale.
Il existe un morphisme de sites
$$
\pi_S : (\Sch/S)_\etale \longrightarrow S_\etale
$$
induit par le foncteur continu
$S_\etale \to (\Sch/S)_\etale$, $U \mapsto U$.
En effet, $S_\etale$ possède des produits fibrés et un objet final, le
foncteur ci-dessus y commute, et la
proposition \ref{sites-proposition-get-morphism} de Sites s'applique.

\begin{lemma}
\label{etale-cohomology-lemma-describe-pullback-pi-h}
Avec les notations ci-dessus.
Soit $\mathcal{F}$ un faisceau sur $S_\etale$. La règle
$$
(\Sch/S)_h \longrightarrow \textit{Ensembles},\quad
(f : X \to S) \longmapsto \Gamma(X, f_{small}^{-1}\mathcal{F})
$$
est un faisceau h, et a fortiori un faisceau sur $(\Sch/S)_\etale$.
En fait, ce faisceau est égal à
$\pi_S^{-1}\mathcal{F}$ sur $(\Sch/S)_\etale$ et à
$\epsilon_S^{-1}\pi_S^{-1}\mathcal{F}$ sur $(\Sch/S)_h$.
\end{lemma}

\begin{proof}
L'énoncé relatif à la topologie étale est le contenu
du lemme \ref{etale-cohomology-lemma-describe-pullback}. Pour achever la démonstration, il
suffit de montrer que $\pi_S^{-1}\mathcal{F}$ est un faisceau pour la topologie
h. Or, dans le lemme \ref{etale-cohomology-lemma-describe-pullback-pi-ph},
nous avons montré que
$\pi_S^{-1}\mathcal{F}$ est même un faisceau pour la topologie ph,
qui est plus fine.
\end{proof}

\noindent
Dans la situation du lemme \ref{etale-cohomology-lemma-describe-pullback-pi-h},
la composée de $\epsilon_S$ et $\pi_S$, ainsi que l'égalité,
déterminent un morphisme de sites
$$
a_S : (\Sch/S)_h \longrightarrow S_\etale
$$

\begin{lemma}
\label{etale-cohomology-lemma-push-pull-h-etale}
Avec les notations ci-dessus.
Soit $f : X \to Y$ un morphisme de $(\Sch/S)_h$.
Il existe alors des diagrammes commutatifs de topos
$$
\xymatrix{
\Sh((\Sch/X)_h) \ar[rr]_{f_{big, h}} \ar[d]_{\epsilon_X} & &
\Sh((\Sch/Y)_h) \ar[d]^{\epsilon_Y} \\
\Sh((\Sch/X)_\etale) \ar[rr]^{f_{big, \etale}} & &
\Sh((\Sch/Y)_\etale)
}
$$
et
$$
\xymatrix{
\Sh((\Sch/X)_h) \ar[rr]_{f_{big, h}} \ar[d]_{a_X} & &
\Sh((\Sch/Y)_h) \ar[d]^{a_Y} \\
\Sh(X_\etale) \ar[rr]^{f_{small}} & &
\Sh(Y_\etale)
}
$$
avec $a_X = \pi_X \circ \epsilon_X$ et $a_Y = \pi_Y \circ \epsilon_Y$.
\end{lemma}

\begin{proof}
La commutativité des diagrammes découle de la discussion de la
section \ref{topologies-section-change-topologies} de Topologies.
\end{proof}

\begin{lemma}
\label{etale-cohomology-lemma-proper-push-pull-h-etale}
Dans le lemme \ref{etale-cohomology-lemma-push-pull-h-etale}, si $f$ est propre, alors nous avons
$a_Y^{-1} \circ f_{small, *} = f_{big, h, *} \circ a_X^{-1}$.
\end{lemma}

\begin{proof}
On peut le démontrer en répétant la démonstration de la
partie (1) du lemme \ref{etale-cohomology-lemma-compare-higher-direct-image-proper};
nous allons plutôt l'en déduire.
Comme $\epsilon_{Y, *}$ est le foncteur identité sur les préfaisceaux sous-jacents,
il reflète les isomorphismes. La description
du lemme \ref{etale-cohomology-lemma-describe-pullback-pi-h}
montre que $\epsilon_{Y, *} \circ a_Y^{-1} = \pi_Y^{-1}$,
et de même pour $X$. Pour montrer que le morphisme canonique
$a_Y^{-1}f_{small, *}\mathcal{F} \to f_{big, h, *}a_X^{-1}\mathcal{F}$
est un isomorphisme, il suffit de montrer que
\begin{align*}
\pi_Y^{-1}f_{small, *}\mathcal{F}
& =
\epsilon_{Y, *}a_Y^{-1}f_{small, *}\mathcal{F} \\
& \to 
\epsilon_{Y, *}f_{big, h, *}a_X^{-1}\mathcal{F} \\
& =
f_{big, \etale, *} \epsilon_{X, *}a_X^{-1}\mathcal{F} \\
& =
f_{big, \etale, *}\pi_X^{-1}\mathcal{F}
\end{align*}
est un isomorphisme. C'est la partie
(1) du lemme \ref{etale-cohomology-lemma-compare-higher-direct-image-proper}.
\end{proof}

\begin{lemma}
\label{etale-cohomology-lemma-compare-h-etale}
Considérons le morphisme de comparaison $\epsilon : (\Sch/S)_h \to (\Sch/S)_\etale$.
Soit $\mathcal{P}$ la classe des morphismes propres.
Pour $X \in (\Sch/S)_\etale$, notons
$\mathcal{A}'_X \subset \textit{Ab}((\Sch/X)_\etale)$
la sous-catégorie pleine formée des faisceaux de la forme
$\pi_X^{-1}\mathcal{F}$, où $\mathcal{F}$ est un
faisceau abélien de torsion sur $X_\etale$.
Alors les propriétés de Cohomologie sur les sites
(\ref{sites-cohomology-item-base-change-P}),
(\ref{sites-cohomology-item-restriction-A}),
(\ref{sites-cohomology-item-A-sheaf}),
(\ref{sites-cohomology-item-A-and-P}) et
(\ref{sites-cohomology-item-refine-tau-by-P})
de la situation de Cohomologie sur les sites
\ref{sites-cohomology-situation-compare} sont satisfaites.
\end{lemma}

\begin{proof}
Montrons d'abord que $\mathcal{A}'_X \subset \textit{Ab}((\Sch/X)_\etale)$
est une sous-catégorie de Serre faible en vérifiant les conditions (1), (2), (3) et (4)
du lemme \ref{homology-lemma-characterize-weak-serre-subcategory} d'Homology.
Les parties (1), (2), (3) sont immédiates, car $\pi_X^{-1}$ est exact et
pleinement fidèle, par exemple par le lemme \ref{etale-cohomology-lemma-cohomological-descent-etale}. Si
$0 \to \pi_X^{-1}\mathcal{F} \to \mathcal{G} \to \pi_X^{-1}\mathcal{F}' \to 0$
est une suite exacte courte dans $\textit{Ab}((\Sch/X)_\etale)$,
alors $0 \to \mathcal{F} \to \pi_{X, *}\mathcal{G} \to \mathcal{F}' \to 0$
est exacte par le lemme \ref{etale-cohomology-lemma-cohomological-descent-etale}.
Nous voyons en particulier que $\pi_{X, *}\mathcal{G}$ est un faisceau
abélien de torsion sur $X_\etale$.
Ainsi, $\mathcal{G} = \pi_X^{-1}\pi_{X, *}\mathcal{G}$ appartient à
$\mathcal{A}'_X$, ce qui vérifie la dernière condition.

\medskip\noindent
La propriété (\ref{sites-cohomology-item-base-change-P}) de Cohomologie sur les sites vaut
par l'existence des produits fibrés de schémas,
le fait que le changement de base d'un morphisme propre de
schémas est un morphisme propre de schémas; voir le
lemme \ref{morphisms-lemma-base-change-proper} de Morphismes; et
le fait que le changement de base d'un morphisme de présentation finie
est un morphisme de présentation finie; voir le
lemme \ref{morphisms-lemma-base-change-finite-presentation} de Morphismes.

\medskip\noindent
La propriété (\ref{sites-cohomology-item-restriction-A}) de Cohomologie sur les sites
découle du diagramme commutatif (3) du
lemme \ref{topologies-lemma-morphism-big-small-etale} de Topologies.

\medskip\noindent
La propriété (\ref{sites-cohomology-item-A-sheaf}) de Cohomologie sur les sites est le
lemme \ref{etale-cohomology-lemma-describe-pullback-pi-h}.

\medskip\noindent
La propriété (\ref{sites-cohomology-item-A-and-P}) de Cohomologie sur les sites vaut par la
partie (2) du lemme \ref{etale-cohomology-lemma-compare-higher-direct-image-proper}
et le fait que $R^if_{small,*}\mathcal{F}$
est de torsion si $\mathcal{F}$ est un faisceau
abélien de torsion sur $X_\etale$; voir le lemme \ref{etale-cohomology-lemma-torsion-direct-image}.

\medskip\noindent
La propriété (\ref{sites-cohomology-item-refine-tau-by-P}) de Cohomologie sur les sites
découle du lemme de Compléments sur les morphismes
\ref{more-morphisms-lemma-dominate-fppf-etale-locally}
combiné au fait qu'un morphisme fini localement libre surjectif
est surjectif, propre et de présentation finie, et définit donc
un recouvrement h par le lemme de Compléments sur la platitude
\ref{flat-lemma-surjective-proper-finite-presentation-h}.
\end{proof}

\begin{lemma}
\label{etale-cohomology-lemma-V-C-all-n-etale-h}
Avec les notations ci-dessus.
\begin{enumerate}
\item Pour $X \in \Ob((\Sch/S)_{h})$ et un faisceau abélien de torsion $\mathcal{F}$
sur $X_\etale$, nous avons
$\epsilon_{X, *}a_X^{-1}\mathcal{F} = \pi_X^{-1}\mathcal{F}$
et $R^i\epsilon_{X, *}(a_X^{-1}\mathcal{F}) = 0$ pour $i > 0$.
\item Pour un morphisme propre $f : X \to Y$ dans $(\Sch/S)_h$
et un faisceau abélien de torsion $\mathcal{F}$ sur $X$, nous avons
$a_Y^{-1}(R^if_{small, *}\mathcal{F}) =
R^if_{big, h, *}(a_X^{-1}\mathcal{F})$
pour tout $i$.
\item Pour un schéma $X$ et $K \in D^+(X_\etale)$ à faisceaux de
cohomologie de torsion, le morphisme
$\pi_X^{-1}K \to R\epsilon_{X, *}(a_X^{-1}K)$ est un isomorphisme.
\item Pour un morphisme propre $f : X \to Y$ de schémas
et $K \in D^+(X_\etale)$ à faisceaux de cohomologie de torsion, nous avons
$a_Y^{-1}(Rf_{small, *}K) = Rf_{big, h, *}(a_X^{-1}K)$.
\end{enumerate}
\end{lemma}

\begin{proof}
Par le lemme \ref{etale-cohomology-lemma-compare-h-etale}, tous les lemmes de la
section \ref{sites-cohomology-section-compare-general} de Cohomologie sur les sites
s'appliquent à notre situation. Pour traduire les résultats,
remarquons que la catégorie $\mathcal{A}_X$ du
lemme \ref{sites-cohomology-lemma-A} de Cohomologie sur les sites
est la sous-catégorie pleine de $\textit{Ab}((\Sch/X)_h)$
formée des faisceaux de la forme $a_X^{-1}\mathcal{F}$,
où $\mathcal{F}$ est un faisceau abélien de torsion sur $X_\etale$.

\medskip\noindent
La partie (1) équivaut à $(V_n)$ pour tout $n$, ce qui vaut par le
lemme \ref{sites-cohomology-lemma-V-C-all-n-general} de Cohomologie sur les sites.

\medskip\noindent
La partie (2) découle de l'application de $\epsilon_Y^{-1}$ à la conclusion du
lemme \ref{sites-cohomology-lemma-V-implies-C-general} de Cohomologie sur les sites.

\medskip\noindent
La partie (3) découle de la partie (1) du lemme de Cohomologie sur les sites
\ref{sites-cohomology-lemma-V-C-all-n-general},
car $\pi_X^{-1}K$ appartient à $D^+_{\mathcal{A}'_X}((\Sch/X)_\etale)$
et $a_X^{-1} = \epsilon_X^{-1} \circ \pi_X^{-1}$.

\medskip\noindent
La partie (4) découle de la partie (2) du lemme de Cohomologie sur les sites
\ref{sites-cohomology-lemma-V-C-all-n-general},
pour la même raison.
\end{proof}

\begin{lemma}
\label{etale-cohomology-lemma-cohomological-descent-etale-h}
Soit $X$ un schéma. Pour $K \in D^+(X_\etale)$ à faisceaux de cohomologie
de torsion, le morphisme
$$
K \longrightarrow Ra_{X, *}a_X^{-1}K
$$
est un isomorphisme, où $a_X : \Sh((\Sch/X)_h) \to \Sh(X_\etale)$ est comme ci-dessus.
\end{lemma}

\begin{proof}
Nous nous ramenons d'abord au cas où
$K$ est donné par un unique faisceau abélien. Représentons en effet $K$
par un complexe borné inférieurement $\mathcal{F}^\bullet$ de faisceaux
abéliens de torsion. C'est possible par le lemme de Cohomologie sur les sites
\ref{sites-cohomology-lemma-torsion}. Le cas d'un
faisceau montre que
$\mathcal{F}^n = a_{X, *} a_X^{-1} \mathcal{F}^n$
et que les faisceaux $R^qa_{X, *}a_X^{-1}\mathcal{F}^n$
sont nuls pour $q > 0$. Le lemme d'acyclicité de Leray
(lemme \ref{derived-lemma-leray-acyclicity} de Catégories dérivées),
appliqué à $a_X^{-1}\mathcal{F}^\bullet$
et au foncteur $a_{X, *}$, permet de conclure. Supposons désormais $K = \mathcal{F}$,
où $\mathcal{F}$ est un faisceau abélien de torsion.

\medskip\noindent
Par le lemme \ref{etale-cohomology-lemma-describe-pullback-pi-h}, nous avons
$a_{X, *}a_X^{-1}\mathcal{F} = \mathcal{F}$. Il suffit donc de montrer que
$R^qa_{X, *}a_X^{-1}\mathcal{F} = 0$ pour $q > 0$.
Pour cela, nous pouvons employer $a_X = \pi_X \circ \epsilon_X$ et
la suite spectrale de Leray
(lemme \ref{sites-cohomology-lemma-relative-Leray} de Cohomologie sur les sites).
Par le lemme \ref{etale-cohomology-lemma-V-C-all-n-etale-h},
nous avons $R^i\epsilon_{X, *}(a_X^{-1}\mathcal{F}) = 0$ pour $i > 0$
et $\epsilon_{X, *}a_X^{-1}\mathcal{F} = \pi_X^{-1}\mathcal{F}$.
Par le lemme \ref{etale-cohomology-lemma-cohomological-descent-etale}, nous avons
$R^j\pi_{X, *}(\pi_X^{-1}\mathcal{F}) = 0$ pour $j > 0$.
Cela achève la démonstration.
\end{proof}

\begin{lemma}
\label{etale-cohomology-lemma-compare-cohomology-etale-h}
Pour un schéma $X$ et $a_X : \Sh((\Sch/X)_h) \to \Sh(X_\etale)$
comme ci-dessus :
\begin{enumerate}
\item $H^q(X_\etale, \mathcal{F}) = H^q_h(X, a_X^{-1}\mathcal{F})$
pour un faisceau abélien de torsion $\mathcal{F}$ sur $X_\etale$,
\item $H^q(X_\etale, K) = H^q_h(X, a_X^{-1}K)$
pour $K \in D^+(X_\etale)$ à faisceaux de cohomologie de torsion.
\end{enumerate}
Exemple : si $A$ est un groupe abélien de torsion, alors
$H^q_\etale(X, \underline{A}) = H^q_h(X, \underline{A})$.
\end{lemma}

\begin{proof}
Cela découle du lemme \ref{etale-cohomology-lemma-cohomological-descent-etale-h}
par la remarque \ref{sites-cohomology-remark-before-Leray} de Cohomologie sur les sites.
\end{proof}











\section{Descente des faisceaux étales}
\label{etale-cohomology-section-glueing-etale}

\noindent
Nous démontrons que les faisceaux étales se recollent pour les topologies
fppf et h, ainsi que des résultats voisins. Nous avons déjà établi les
résultats connexes suivants :
\begin{enumerate}
\item le lemme \ref{etale-cohomology-lemma-describe-pullback} affirme qu'un faisceau
sur le petit site étale d'un schéma $S$ détermine un faisceau sur le
grand site étale de $S$ satisfaisant la condition de faisceau pour les
recouvrements fpqc (et a fortiori pour les recouvrements de Zariski, étales,
lisses, syntomiques et fppf);
\item le lemme \ref{etale-cohomology-lemma-describe-pullback-pi-fppf} reformule
le point précédent pour la topologie fppf;
\item le lemme \ref{etale-cohomology-lemma-describe-pullback-pi-ph} établit le même résultat pour la
topologie ph;
\item le lemme \ref{etale-cohomology-lemma-describe-pullback-pi-h} établit le même résultat pour la
topologie h;
\item le lemme \ref{etale-cohomology-lemma-descent-sheaf-fppf-etale} est une version de la
descente fppf pour les faisceaux étales; et
\item la remarque \ref{etale-cohomology-remark-cohomological-descent-finite} affirme que
les faisceaux étales descendent par les morphismes finis surjectifs
(nous allons préciser et renforcer ce résultat ci-dessous).
\end{enumerate}
Dans le chapitre sur les espaces simpliciaux, nous démontrerons d'autres
résultats à ce sujet; voir par exemple les sections de Espaces simpliciaux
\ref{spaces-simplicial-section-fppf-hypercovering}
et \ref{spaces-simplicial-section-proper-hypercovering-spaces}.

\medskip\noindent
Pour énoncer commodément nos résultats, introduisons quelques notations.
Soit $\mathcal{U} = \{f_i : X_i \to X\}$
une famille de morphismes de schémas de but fixé.
Une {\it donnée de descente} pour les faisceaux étales relativement à
$\mathcal{U}$ est une famille
$((\mathcal{F}_i)_{i \in I}, (\varphi_{ij})_{i, j \in I})$, où
\begin{enumerate}
\item $\mathcal{F}_i$ appartient à $\Sh(X_{i, \etale})$; et
\item $\varphi_{ij} :
\text{pr}_{0, small}^{-1} \mathcal{F}_i
\longrightarrow
\text{pr}_{1, small}^{-1} \mathcal{F}_j$
est un isomorphisme dans $\Sh((X_i \times_X X_j)_\etale)$.
\end{enumerate}
On exige la {\it condition de cocycle} : les diagrammes
$$
\xymatrix{
\text{pr}_{0, small}^{-1}\mathcal{F}_i
\ar[dr]_{\text{pr}_{02, small}^{-1}\varphi_{ik}}
\ar[rr]^{\text{pr}_{01, small}^{-1}\varphi_{ij}} & &
\text{pr}_{1, small}^{-1}\mathcal{F}_j
\ar[dl]^{\text{pr}_{12, small}^{-1}\varphi_{jk}} \\
& \text{pr}_{2, small}^{-1}\mathcal{F}_k
}
$$
commutent dans $\Sh((X_i \times_X X_j \times_X X_k)_\etale)$.
On dispose d'une notion évidente de {\it morphisme de données de descente},
et l'on obtient une catégorie de données de descente.
Une donnée de descente
$((\mathcal{F}_i)_{i \in I}, (\varphi_{ij})_{i, j \in I})$
est dite {\it effective} s'il existe un objet
$\mathcal{F}$ de $\Sh(X_\etale)$ et des isomorphismes
$\varphi_i : f_{i, small}^{-1} \mathcal{F} \to \mathcal{F}_i$
dans $\Sh(X_{i, \etale})$, compatibles aux $\varphi_{ij}$, c'est-à-dire
tels que
$$
\varphi_{ij} =
\text{pr}_{1, small}^{-1} (\varphi_j) \circ
\text{pr}_{0, small}^{-1} (\varphi_i^{-1})
$$
On peut également formuler cela comme suit. Tout objet $\mathcal{F}$
de $\Sh(X_\etale)$ fournit la {\it donnée de descente canonique}
$(f_{i, small}^{-1}\mathcal{F}, c_{ij})$, où $c_{ij}$
est l'isomorphisme canonique
$$
c_{ij} : \text{pr}_{0, small}^{-1} f_{i, small}^{-1}\mathcal{F}
\longrightarrow
\text{pr}_{1, small}^{-1} f_{j, small}^{-1}\mathcal{F}
$$
La donnée de descente
$((\mathcal{F}_i)_{i \in I}, (\varphi_{ij})_{i, j \in I})$
est effective si et seulement si elle est isomorphe à la donnée de descente
canonique associée à un certain $\mathcal{F}$ de $\Sh(X_\etale)$.

\medskip\noindent
Si la famille est réduite à un seul morphisme $\{X \to Y\}$,
nous voyons une donnée de descente comme un couple $(\mathcal{F}, \varphi)$,
où $\mathcal{F}$ est un objet de $\Sh(X_\etale)$ et
$\varphi$ un isomorphisme
$$
\text{pr}_{0, small}^{-1} \mathcal{F}
\longrightarrow
\text{pr}_{1, small}^{-1} \mathcal{F}
$$
dans $\Sh((X \times_Y X)_\etale)$ satisfaisant la condition de cocycle :
$$
\xymatrix{
\text{pr}_{0, small}^{-1}\mathcal{F}
\ar[dr]_{\text{pr}_{02, small}^{-1}\varphi}
\ar[rr]^{\text{pr}_{01, small}^{-1}\varphi} & &
\text{pr}_{1, small}^{-1}\mathcal{F}
\ar[dl]^{\text{pr}_{12, small}^{-1}\varphi} \\
& \text{pr}_{2, small}^{-1}\mathcal{F}
}
$$
commute dans $\Sh((X \times_Y X \times_Y X)_\etale)$.
Les notions de morphisme de données de descente et d'effectivité
sont exactement celles données ci-dessus.

\medskip\noindent
Nous démontrons d'abord la descente effective pour les morphismes entiers surjectifs.

\begin{lemma}
\label{etale-cohomology-lemma-glue-etale-sheaf-section}
Soit $f : X \to Y$ un morphisme de schémas qui possède une section. Alors le
foncteur
$$
\Sh(Y_\etale)
\longrightarrow
\text{données de descente pour les faisceaux étales relativement à }\{X \to Y\}
$$
qui associe à $\mathcal{G}$ dans $\Sh(Y_\etale)$ sa donnée de descente canonique
est une équivalence de catégories.
\end{lemma}

\begin{proof}
C'est formel et ne dépend que de la fonctorialité des foncteurs
d'image inverse. Nous omettons les détails. Indication : si $s : Y \to X$ est une section,
un quasi-inverse est le foncteur qui associe à $(\mathcal{F}, \varphi)$
l'objet $s_{small}^{-1}\mathcal{F}$.
\end{proof}

\begin{lemma}
\label{etale-cohomology-lemma-glue-etale-sheaf-integral-surjective}
Soit $f : X \to Y$ un morphisme entier surjectif de schémas.
Le foncteur
$$
\Sh(Y_\etale)
\longrightarrow
\text{données de descente pour les faisceaux étales relativement à }\{X \to Y\}
$$
est une équivalence de catégories.
\end{lemma}

\begin{proof}
Dans cette démonstration, nous omettons l'indice ${}_{small}$ de nos foncteurs
d'image inverse et d'image directe.
Notons $X_1 = X \times_Y X$ et notons $f_1 : X_1 \to Y$ le
morphisme $f \circ \text{pr}_0 = f \circ \text{pr}_1$.
Soit $(\mathcal{F}, \varphi)$ une donnée de descente pour $\{X \to Y\}$.
Posons $\mathcal{F}_1 = \text{pr}_0^{-1}\mathcal{F}$.
Nous pouvons voir $\varphi$ comme définissant un isomorphisme
$\mathcal{F}_1 \to \text{pr}_1^{-1}\mathcal{F}$.
Nous affirmons que la règle qui associe à une donnée de descente
$(\mathcal{F}, \varphi)$
le faisceau
$$
\mathcal{G} =
\text{Égalisateur}\left(
\xymatrix{
f_*\mathcal{F}
\ar@<1ex>[r] \ar@<-1ex>[r] &
f_{1, *}\mathcal{F}_1
}
\right)
$$
est un quasi-inverse du foncteur de l'énoncé du lemme.
Le premier des deux morphismes provient du morphisme
$$
f_*\mathcal{F} \to
f_*\text{pr}_{0, *}\text{pr}_0^{-1}\mathcal{F} =
f_{1, *}\mathcal{F}_1
$$
et le second provient du morphisme
$$
f_*\mathcal{F} \to
f_* \text{pr}_{1, *}\text{pr}_1^{-1}\mathcal{F}
\xleftarrow{\varphi}
f_* \text{pr}_{0, *} \text{pr}_0^{-1}\mathcal{F} =
f_{1, *}\mathcal{F}_1
$$
où le morphisme dirigé vers la gauche est inversible.
Pour justifier cette construction, il faut montrer
que le morphisme canonique $f^{-1}\mathcal{G} \to \mathcal{F}$
est un isomorphisme; nous omettons les détails. Pour ce faire, il
suffit de le vérifier après image inverse par une collection de morphismes
$\Spec(k) \to Y$, où $k$ est un corps algébriquement clos.
En effet, les changements de base correspondants $X_k \to X$ sont conjointement
surjectifs, et l'on peut vérifier qu'un morphisme de faisceaux sur $X_\etale$
est un isomorphisme en examinant ses fibres aux points géométriques; voir le
théorème \ref{etale-cohomology-theorem-exactness-stalks}.
Par le lemme \ref{etale-cohomology-lemma-integral-pushforward-commutes-with-base-change},
la construction de $\mathcal{G}$ à partir de la donnée de descente
$(\mathcal{F}, \varphi)$ commute à tout changement de base.
Nous pouvons donc supposer que $Y$ est le spectre d'un corps
algébriquement clos (remarquons que le changement de base préserve les propriétés
du morphisme $f$; voir les lemmes de Morphismes
\ref{morphisms-lemma-base-change-surjective}
et \ref{morphisms-lemma-base-change-finite}).
Dans ce cas, le morphisme $X \to Y$ possède une section, et nous
savons donc, par le lemme \ref{etale-cohomology-lemma-glue-etale-sheaf-section},
que le foncteur est une équivalence.
Or le lecteur peut vérifier que le foncteur est une équivalence
si et seulement si la construction ci-dessus est un quasi-inverse;
nous omettons les détails. Cela achève la démonstration.
\end{proof}

\begin{lemma}
\label{etale-cohomology-lemma-glue-etale-sheaf-proper-surjective}
Soit $f : X \to Y$ un morphisme propre surjectif de schémas.
Le foncteur
$$
\Sh(Y_\etale)
\longrightarrow
\text{données de descente pour les faisceaux étales relativement à }\{X \to Y\}
$$
est une équivalence de catégories.
\end{lemma}

\begin{proof}
La démonstration du
lemme \ref{etale-cohomology-lemma-glue-etale-sheaf-integral-surjective}
s'applique mot pour mot, si ce n'est que l'emploi du
lemme \ref{etale-cohomology-lemma-integral-pushforward-commutes-with-base-change}
doit être remplacé par celui du
lemme \ref{etale-cohomology-lemma-proper-base-change-f-star}.
\end{proof}

\begin{lemma}
\label{etale-cohomology-lemma-glue-etale-sheaf-check-after-base-change}
Soit $f : X \to Y$ un morphisme de schémas. Soit $Z \to Y$ un morphisme
entier surjectif de schémas ou un morphisme propre surjectif de schémas.
Si les foncteurs
$$
\Sh(Z_\etale)
\longrightarrow
\text{données de descente pour les faisceaux étales relativement à }\{X \times_Y Z \to Z\}
$$
et
$$
\Sh((Z \times_Y Z)_\etale)
\longrightarrow
\text{données de descente pour les faisceaux étales relativement à }
\{X \times_Y (Z \times_Y Z) \to Z \times_Y Z\}
$$
sont des équivalences de catégories, alors
$$
\Sh(Y_\etale)
\longrightarrow
\text{données de descente pour les faisceaux étales relativement à }\{X \to Y\}
$$
est une équivalence.
\end{lemma}

\begin{proof}
C'est une conséquence formelle des définitions et des
lemmes \ref{etale-cohomology-lemma-glue-etale-sheaf-integral-surjective} et
\ref{etale-cohomology-lemma-glue-etale-sheaf-proper-surjective}. Nous omettons les détails.
\end{proof}

\begin{lemma}
\label{etale-cohomology-lemma-glue-etale-sheaf-fppf-cover}
Soit $f : X \to Y$ un morphisme de schémas
surjectif, plat et localement de présentation finie.
Le foncteur
$$
\Sh(Y_\etale)
\longrightarrow
\text{données de descente pour les faisceaux étales relativement à }\{X \to Y\}
$$
est une équivalence de catégories.
\end{lemma}

\begin{proof}
Exactement comme dans la démonstration du
lemme \ref{etale-cohomology-lemma-glue-etale-sheaf-integral-surjective},
nous affirmons qu'un quasi-inverse est le foncteur qui associe
à $(\mathcal{F}, \varphi)$ le faisceau
$$
\mathcal{G} =
\text{Égalisateur}\left(
\xymatrix{
f_*\mathcal{F}
\ar@<1ex>[r] \ar@<-1ex>[r] &
f_{1, *}\mathcal{F}_1
}
\right)
$$
et, pour le démontrer, il suffit de montrer que
$f^{-1}\mathcal{G} \to \mathcal{F}$ est un isomorphisme.
Cette propriété se vérifie localement; nous pouvons donc supposer, et supposons, $Y$
affine. Il existe alors un morphisme fini surjectif
$Z \to Y$ et un recouvrement ouvert
$Z = \bigcup W_i$ tels que $W_i \to Y$ se factorise par $X$.
Voir le
lemme \ref{more-morphisms-lemma-dominate-fppf-finite} de Compléments sur les morphismes.
En appliquant le lemme
\ref{etale-cohomology-lemma-glue-etale-sheaf-check-after-base-change}
nous voyons qu'il suffit de démontrer
le lemme après avoir remplacé $Y$ par $Z$ et $Z \times_Y Z$, et $f$
par son changement de base. Nous pouvons donc supposer que $f$ possède localement des sections pour Zariski.
Puisque le problème est local sur $Y$, nous nous ramenons alors
au cas où nous disposons d'une section, lequel est le
lemme \ref{etale-cohomology-lemma-glue-etale-sheaf-section}.
\end{proof}

\begin{lemma}
\label{etale-cohomology-lemma-glue-etale-sheaf-fppf}
Soit $\{f_i : X_i \to X\}$ un recouvrement fppf de schémas.
Le foncteur
$$
\Sh(X_\etale)
\longrightarrow
\text{données de descente pour les faisceaux étales relativement à }\{f_i : X_i \to X\}
$$
est une équivalence de catégories.
\end{lemma}

\begin{proof}
Appliquons le lemme \ref{etale-cohomology-lemma-glue-etale-sheaf-fppf-cover}
au morphisme $f : \coprod X_i \to X$.
Un argument formel montre alors que les données de descente pour $f$
sont les mêmes que les données de descente pour le recouvrement; comparer
avec le lemme \ref{descent-lemma-family-is-one} de Descente.
Nous omettons les détails.
\end{proof}

\begin{lemma}
\label{etale-cohomology-lemma-glue-etale-sheaf-modification}
Soit $f : X' \to X$ un morphisme propre de schémas. Soit $i : Z \to X$
une immersion fermée. Posons $E = Z \times_X X'$. Voici le diagramme :
$$
\xymatrix{
E \ar[d]_g \ar[r]_j & X' \ar[d]^f \\
Z \ar[r]^i & X
}
$$
Si $f$ est un isomorphisme au-dessus de $X \setminus Z$, alors le foncteur
$$
\Sh(X_\etale)
\longrightarrow
\Sh(X'_\etale) \times_{\Sh(E_\etale)} \Sh(Z_\etale)
$$
est une équivalence de catégories.
\end{lemma}

\begin{proof}
Nous travaillons avec le $2$-produit fibré de catégories construit dans
l'exemple \ref{categories-example-2-fibre-product-categories} de Catégories.
Le foncteur associe à $\mathcal{F}$ le triplet
$(f^{-1}\mathcal{F}, i^{-1}\mathcal{F}, c)$, où
$c : j^{-1}f^{-1}\mathcal{F} \to g^{-1}i^{-1}\mathcal{F}$
est l'isomorphisme canonique. Construisons un foncteur quasi-inverse. Soit
$(\mathcal{F}', \mathcal{G}, \alpha)$ un objet
du membre de droite de la flèche.
Nous obtenons un isomorphisme
$$
i^{-1}f_*\mathcal{F}' = g_*j^{-1}\mathcal{F}'
\xrightarrow{g_*\alpha}
g_*g^{-1}\mathcal{G}
$$
La première égalité est le lemme \ref{etale-cohomology-lemma-proper-base-change-f-star}.
Nous en déduisons des morphismes
$i_*\mathcal{G} \to i_*g_*g^{-1}\mathcal{G}$
et $f_*\mathcal{F}' \to i_*g_*g^{-1}\mathcal{G}$. Posons
$$
\mathcal{F} = f_*\mathcal{F}' \times_{i_*g_*g^{-1}\mathcal{G}} i_*\mathcal{G}
$$
Nous affirmons que $\mathcal{F}$ est un objet du membre de gauche
de la flèche, dont l'image dans le membre de droite est isomorphe au
triplet initial. Calculons la fibre de $\mathcal{F}$
en un point géométrique $\overline{x}$ de $X$.

\medskip\noindent
Si $\overline{x}$ n'appartient pas
à $Z$, alors, d'une part, $\overline{x}$ provient d'un unique
point géométrique $\overline{x}'$ de $X'$ et
$\mathcal{F}'_{\overline{x}'} = (f_*\mathcal{F}')_{\overline{x}}$
et, d'autre part, $(i_*\mathcal{G})_{\overline{x}}$
et $(i_*g_*g^{-1}\mathcal{G})_{\overline{x}}$ sont des singletons.
Nous voyons donc que $\mathcal{F}_{\overline{x}}$ est égal à
$\mathcal{F}'_{\overline{x}'}$.

\medskip\noindent
Si $\overline{x}$ appartient à $Z$, c'est-à-dire si $\overline{x}$ est l'image d'un
point géométrique $\overline{z}$ de $Z$, alors nous obtenons
$(i_*\mathcal{G})_{\overline{x}} = \mathcal{G}_{\overline{z}}$
et
$$
(i_*g_*g^{-1}\mathcal{G})_{\overline{x}} =
(g_*g^{-1}\mathcal{G})_{\overline{z}} =
\Gamma(E_{\overline{z}}, g^{-1}\mathcal{G}|_{E_{\overline{z}}})
$$
(par le changement de base propre pour l'image directe employé ci-dessus),
et de même
$$
(f_*\mathcal{F}')_{\overline{x}} =
\Gamma(X'_{\overline{x}}, \mathcal{F}'|_{X'_{\overline{x}}})
$$
Puisque nous avons l'identification
$E_{\overline{z}} = X'_{\overline{x}}$ et que $\alpha$
définit un isomorphisme entre les faisceaux
$\mathcal{F}'|_{X'_{\overline{x}}}$ et
$g^{-1}\mathcal{G}|_{E_{\overline{z}}}$
nous en déduisons
$$
\mathcal{F}_{\overline{x}} = \mathcal{G}_{\overline{z}}
$$
dans ce cas.

\medskip\noindent
Pour achever la démonstration, observons qu'il existe des morphismes canoniques
$i^{-1}\mathcal{F} \to \mathcal{G}$ et $f^{-1}\mathcal{F} \to \mathcal{F}'$
compatibles avec $\alpha$ et induisent sur les fibres les isomorphismes
obtenus ci-dessus. Nous omettons la construction détaillée de ces morphismes.
\end{proof}

\begin{lemma}
\label{etale-cohomology-lemma-h-descent-etale-sheaves}
Soit $S$ un schéma. Alors la catégorie fibrée
$$
p : \mathcal{S} \longrightarrow (\Sch/S)_h
$$
dont la catégorie fibre au-dessus de $U$ est la catégorie $\Sh(U_\etale)$
des faisceaux sur le petit site étale de $U$ est un champ.
\end{lemma}

\begin{proof}
Pour démontrer le lemme, nous vérifions les conditions (1), (2) et (3) du
lemme \ref{flat-lemma-refine-check-h-stack} de Compléments sur la platitude.

\medskip\noindent
La condition (1) vaut parce que les faisceaux se recollent (et que les
recouvrements de Zariski sont des recouvrements étales). Voir le
lemme \ref{sites-lemma-glue-sheaves} de Sites.

\medskip\noindent
Pour vérifier la condition (2), supposons que $f : X \to Y$ soit un morphisme
surjectif, plat, propre et de présentation finie sur $S$, avec $Y$ affine.
Nous avons alors la descente pour $\{X \to Y\}$, soit par le
lemme \ref{etale-cohomology-lemma-glue-etale-sheaf-fppf-cover}, soit par le
lemme \ref{etale-cohomology-lemma-glue-etale-sheaf-proper-surjective}.

\medskip\noindent
La condition (3) découle immédiatement du résultat plus général qu'est le
lemme \ref{etale-cohomology-lemma-glue-etale-sheaf-modification}.
\end{proof}



















\section{Carrés d'éclatement et cohomologie étale}
\label{etale-cohomology-section-blow-up-square}

\noindent
Les carrés d'éclatement sont introduits dans la
section \ref{flat-section-blow-up-ph} de Compléments sur la platitude.
Le théorème de changement de base propre permet d'obtenir
un résultat de type Mayer--Vietoris pour les carrés d'éclatement.

\begin{lemma}
\label{etale-cohomology-lemma-blow-up-square-cohomological-descent}
Soit $X$ un schéma et soit $Z \subset X$ un sous-schéma fermé
défini par un idéal quasi-cohérent de type fini. Considérons le
carré d'éclatement correspondant
$$
\xymatrix{
E \ar[d]_\pi \ar[r]_j & X' \ar[d]^b \\
Z \ar[r]^i & X
}
$$
Pour $K \in D^+(X_\etale)$ à faisceaux de cohomologie de torsion,
nous avons un triangle distingué
$$
K \to Ri_*(K|_Z) \oplus Rb_*(K|_{X'}) \to Rc_*(K|_E) \to K[1]
$$
dans $D(X_\etale)$, où $c = i \circ \pi = b \circ j$.
\end{lemma}

\begin{proof}
La notation $K|_{X'}$ désigne $b_{small}^{-1}K$.
Choisissons un complexe borné inférieurement $\mathcal{F}^\bullet$
de faisceaux abéliens représentant $K$. Remarquons que
$i_*(\mathcal{F}^\bullet|_Z)$ représente $Ri_*(K|_Z)$,
car $i_*$ est exact
(proposition \ref{etale-cohomology-proposition-finite-higher-direct-image-zero}).
Choisissons un quasi-isomorphisme
$b_{small}^{-1}\mathcal{F}^\bullet \to \mathcal{I}^\bullet$
où $\mathcal{I}^\bullet$ est un complexe borné inférieurement de faisceaux
abéliens injectifs sur $X'_\etale$. Ce morphisme est adjoint à un morphisme
$\mathcal{F}^\bullet \to b_*(\mathcal{I}^\bullet)$, et
$b_*(\mathcal{I}^\bullet)$ représente $Rb_*(K|_{X'})$.
Nous avons $\pi_*(\mathcal{I}^\bullet|_E) = (b_*\mathcal{I}^\bullet)|_Z$
par le lemme \ref{etale-cohomology-lemma-proper-base-change-f-star}, et, par le
lemme \ref{etale-cohomology-lemma-proper-base-change}, ce complexe représente
$R\pi_*(K|_E)$. Par conséquent, le morphisme
$$
Ri_*(K|_Z) \oplus Rb_*(K|_{X'}) \to Rc_*(K|_E)
$$
est représenté par le morphisme surjectif de complexes bornés inférieurement
$$
i_*(\mathcal{F}^\bullet|_Z) \oplus
b_*(\mathcal{I}^\bullet)
\to
i_*\left(b_*(\mathcal{I}^\bullet)|_Z\right)
$$
Pour obtenir le triangle distingué cherché, il suffit de montrer que
le morphisme canonique
$\mathcal{F}^\bullet \to i_*(\mathcal{F}^\bullet|_Z) \oplus
b_*(\mathcal{I}^\bullet)$
induit un quasi-isomorphisme sur le noyau du morphisme de complexes
affiché ci-dessus (en effet, une suite exacte courte
de complexes détermine un triangle distingué dans la catégorie
dérivée; voir la
section \ref{derived-section-canonical-delta-functor} de Catégories dérivées).
On peut le vérifier sur les fibres en un point géométrique $\overline{x}$ de $X$.
Si $\overline{x}$ n'appartient pas à $Z$, alors $X' \to X$ est un isomorphisme
au-dessus d'un voisinage ouvert de $\overline{x}$. Si $\overline{x}'$
désigne le point géométrique correspondant de $X'$ dans ce cas, il
faut donc montrer que
$$
\mathcal{F}^\bullet_{\overline{x}} \to \mathcal{I}^\bullet_{\overline{x}'}
$$
est un quasi-isomorphisme. Cela résulte du choix de $\mathcal{I}^\bullet$.
Si $\overline{x}$ appartient à $Z$, alors
$b_*(\mathcal{I}^\bullet)_{\overline{x}} \to 
i_*\left(b_*(\mathcal{I}^\bullet)|_Z\right)_{\overline{x}}$
est un isomorphisme de complexes de groupes abéliens. Le
noyau est donc égal à
$i_*(\mathcal{F}^\bullet|_Z)_{\overline{x}} =
\mathcal{F}^\bullet_{\overline{x}}$, comme voulu.
\end{proof}

\begin{lemma}
\label{etale-cohomology-lemma-blow-up-square-etale-cohomology}
Soit $X$ un schéma et soit $K \in D^+(X_\etale)$ à
faisceaux de cohomologie de torsion. Soit $Z \subset X$ un sous-schéma fermé
défini par un idéal quasi-cohérent de type fini. Considérons le
carré d'éclatement correspondant
$$
\xymatrix{
E \ar[d] \ar[r] & X' \ar[d]^b \\
Z \ar[r] & X
}
$$
Il existe alors une suite exacte longue canonique
$$
H^p_\etale(X, K) \to
H^p_\etale(X', K|_{X'}) \oplus
H^p_\etale(Z, K|_Z) \to
H^p_\etale(E, K|_E) \to
H^{p + 1}_\etale(X, K)
$$
\end{lemma}

\begin{proof}[Première démonstration]
Cela découle immédiatement du
lemme \ref{etale-cohomology-lemma-blow-up-square-cohomological-descent}
et de l'égalité
$$
R\Gamma(X, Rb_*(K|_{X'})) = R\Gamma(X', K|_{X'})
$$
(voir la section \ref{sites-cohomology-section-leray} de Cohomologie sur les sites),
ainsi que des égalités analogues pour les autres termes.
\end{proof}

\begin{proof}[Seconde démonstration]
Par le lemme \ref{etale-cohomology-lemma-compare-cohomology-etale-ph},
ces groupes de cohomologie sont ceux de
$X, X', E, Z$ à valeurs dans un certain complexe de faisceaux abéliens
sur le site $(\Sch/X)_{ph}$. (Il s'agit de l'objet
$a_X^{-1}K$ de la catégorie dérivée; voir le
lemme \ref{etale-cohomology-lemma-describe-pullback-pi-ph} ci-dessus
et rappelons que $K|_{X'} = b_{small}^{-1}K$.)
Par le lemme \ref{flat-lemma-blow-up-square-ph} de Compléments sur la platitude,
le faisceau ph associé au diagramme de préfaisceaux
représentables est cocartésien. Le lemme découle donc du résultat très général qu'est le
lemme \ref{sites-cohomology-lemma-square-triangle} de Cohomologie sur les sites,
appliqué au site $(\Sch/X)_{ph}$ et au diagramme commutatif
du lemme.
\end{proof}

\begin{lemma}
\label{etale-cohomology-lemma-blow-up-square-equivalence}
Soit $X$ un schéma et soit $Z \subset X$ un sous-schéma fermé
défini par un idéal quasi-cohérent de type fini. Considérons le
carré d'éclatement correspondant
$$
\xymatrix{
E \ar[d]_\pi \ar[r]_j & X' \ar[d]^b \\
Z \ar[r]^i & X
}
$$
Supposons donnés
\begin{enumerate}
\item un objet $K'$ de $D^+(X'_\etale)$ à faisceaux de cohomologie de torsion;
\item un objet $L$ de $D^+(Z_\etale)$ à faisceaux de cohomologie de torsion; et
\item un isomorphisme $\gamma : K'|_E \to L|_E$.
\end{enumerate}
Il existe alors un objet $K$ de $D^+(X_\etale)$
et des isomorphismes $f : K|_{X'} \to K'$, $g : K|_Z \to L$ tels
que $\gamma = g|_E \circ f^{-1}|_E$.
De plus, supposons donnés
\begin{enumerate}
\item un objet $M$ de $D^+(X_\etale)$ à faisceaux de cohomologie de torsion;
\item un morphisme $\alpha : K' \to M|_{X'}$ de $D(X'_\etale)$;
\item un morphisme $\beta : L \to M|_Z$ de $D(Z_\etale)$;
\end{enumerate}
tels que
$$
\alpha|_E  = \beta|_E \circ \gamma.
$$
Alors il existe un morphisme $K \to M$ dans $D(X_\etale)$
dont la restriction à $X'$ est $\alpha \circ f$
et dont la restriction à $Z$ est $\beta \circ g$.
\end{lemma}

\begin{proof}
Si $K$ existe, le lemme \ref{etale-cohomology-lemma-blow-up-square-cohomological-descent}
donne un triangle distingué dans lequel il s'insère. Choisissons donc simplement
un triangle distingué
$$
K \to Ri_*(L) \oplus Rb_*(K') \to Rc_*(L|_E) \to K[1]
$$
où $c = i \circ \pi = b \circ j$. Ici, le morphisme $Ri_*(L) \to Rc_*(L|_E)$
est l'image par $Ri_*$ du morphisme d'adjonction $L \to R\pi_*(L|_E)$.
Le morphisme $Rb_*(K') \to Rc_*(L|_E)$ est la composée du morphisme canonique
$Rb_*(K') \to Rc_*(K'|_E)$ et de $Rc_*(\gamma)$.
Les morphismes $g$ et $f$ de l'énoncé du lemme sont les adjoints
de ces morphismes. Si l'on restreint ce triangle distingué à $Z$,
alors le morphisme $Rb_*(K') \to Rc_*(L|_E)$ devient un isomorphisme
par le théorème de changement de base (lemme \ref{etale-cohomology-lemma-proper-base-change}), et donc
le morphisme $g : K|_Z \to L$ est un isomorphisme.
L'examen du triangle distingué montre que $f : K|_{X'} \to K'$
est un isomorphisme au-dessus de $X' \setminus E = X \setminus Z$.
De plus, nous avons $\gamma \circ f|_E = g|_E$ par construction.
Comme $\gamma$ et $g$ sont des isomorphismes, nous en déduisons
que $f$ induit également des isomorphismes sur les fibres aux points géométriques de $E$.
Ainsi, $f$ est un isomorphisme.

\medskip\noindent
Pour la dernière assertion, on peut remplacer $K'$ par $K|_{X'}$,
$L$ par $K|_Z$ et $\gamma$ par l'identification canonique.
Observons que $\alpha$ et $\beta$ induisent un carré commutatif
$$
\xymatrix{
K \ar[r] \ar@{..>}[d] &
Ri_*(K|_Z) \oplus Rb_*(K|_{X'}) \ar[r] \ar[d]_{\beta \oplus \alpha} &
Rc_*(K|_E) \ar[r] \ar[d]_{\alpha|_E} &
K[1] \ar@{..>}[d] \\
M \ar[r] &
Ri_*(M|_Z) \oplus Rb_*(M|_{X'}) \ar[r] &
Rc_*(M|_E) \ar[r] &
M[1]
}
$$
Les axiomes d'une catégorie dérivée fournissent donc une flèche
pointillée qui définit un morphisme de triangles distingués.
\end{proof}














\section{Carrés presque d'éclatement et topologie h}
\label{etale-cohomology-section-blow-up-h}

\noindent
Dans cette section, nous poursuivons la discussion de la
section \ref{flat-section-blow-up-h} de Compléments sur la platitude.
Pour la commodité du lecteur, rappelons qu'un
carré presque d'éclatement est un diagramme commutatif
\begin{equation}
\label{etale-cohomology-equation-almost-blow-up-square}
\vcenter{
\xymatrix{
E \ar[d] \ar[r] & X' \ar[d]^b \\
Z \ar[r] & X
}
}
\end{equation}
de schémas satisfaisant les conditions suivantes :
\begin{enumerate}
\item $Z \to X$ est une immersion fermée de présentation finie;
\item $E = b^{-1}(Z)$ est un sous-schéma fermé localement principal de $X'$;
\item $b$ est propre et de présentation finie;
\item le sous-schéma fermé $X'' \subset X'$ défini par l'idéal quasi-cohérent
des sections de $\mathcal{O}_{X'}$ à support dans $E$
(lemme \ref{properties-lemma-sections-supported-on-closed-subset} de Propriétés)
est l'éclatement de $X$ le long de $Z$.
\end{enumerate}
Il s'ensuit que le morphisme $b$ induit un isomorphisme
$X' \setminus E \to X \setminus Z$.

\medskip\noindent
Nous allons donner un critère caractérisant le « caractère de faisceau h » des
objets de la catégorie dérivée du grand site fppf
$(\Sch/S)_{fppf}$ d'un schéma $S$. Sur la même catégorie sous-jacente,
nous disposons d'une seconde topologie, à savoir la topologie h
(section \ref{flat-section-h} de Compléments sur la platitude).
Rappelons que les recouvrements fppf sont des recouvrements h
(lemme \ref{flat-lemma-zariski-h} de Compléments sur la platitude). Nous pouvons donc
considérer le morphisme
$$
\epsilon : (\Sch/S)_h \longrightarrow (\Sch/S)_{fppf}
$$
Voir la section \ref{sites-cohomology-section-compare} de Cohomologie sur les sites.
En particulier, nous avons un foncteur pleinement fidèle
$$
R\epsilon_* : D((\Sch/S)_h) \longrightarrow D((\Sch/S)_{fppf})
$$
et nous pouvons demander : quelle est l'image essentielle de ce foncteur?

\begin{lemma}
\label{etale-cohomology-lemma-blow-up-square-h}
Avec les notations ci-dessus, si $K$ appartient à l'image essentielle
de $R\epsilon_*$, alors les morphismes $c^K_{X, Z, X', E}$ du
lemme \ref{sites-cohomology-lemma-c-square} de Cohomologie sur les sites
sont des quasi-isomorphismes.
\end{lemma}

\begin{proof}
Notons ${}^\#$ le passage au faisceau associé pour la topologie h.
Nous avons vu dans le lemme \ref{flat-lemma-blow-up-square-h} de Compléments sur la platitude
que $h_X^\# = h_Z^\# \amalg_{h_E^\#} h_{X'}^\#$. D'autre part,
le morphisme $h_E^\# \to h_{X'}^\#$ est injectif, car $E \to X'$ est un
monomorphisme. Ce lemme est donc un cas particulier du
lemme \ref{sites-cohomology-lemma-descent-squares-helper} de Cohomologie sur les sites
(qui est lui-même une conséquence formelle du
lemme \ref{sites-cohomology-lemma-square-triangle} de Cohomologie sur les sites).
\end{proof}

\begin{proposition}
\label{etale-cohomology-proposition-check-h}
Soit $K$ un objet de $D^+((\Sch/S)_{fppf})$.
Alors $K$ appartient à l'image essentielle de
$R\epsilon_* : D((\Sch/S)_h) \to D((\Sch/S)_{fppf})$
si et seulement si $c^K_{X, X', Z, E}$ est un quasi-isomorphisme
pour tout carré presque d'éclatement (\ref{etale-cohomology-equation-almost-blow-up-square})
dans $(\Sch/S)_h$ tel que $X$ soit affine.
\end{proposition}

\begin{proof}
Cela résulte de l'application du
lemme \ref{sites-cohomology-lemma-descent-squares} de Cohomologie sur les sites,
dont les hypothèses sont satisfaites par le
lemme \ref{etale-cohomology-lemma-blow-up-square-h} et la
proposition \ref{flat-proposition-check-h} de Compléments sur la platitude.
\end{proof}

\begin{lemma}
\label{etale-cohomology-lemma-refine-check-h}
Soit $K$ un objet de $D^+((\Sch/S)_{fppf})$. Alors $K$ appartient à
l'image essentielle de $R\epsilon_* : D((\Sch/S)_h) \to D((\Sch/S)_{fppf})$
si et seulement si $c^K_{X, X', Z, E}$ est un quasi-isomorphisme
pour tout carré presque d'éclatement comme dans les
exemples \ref{flat-example-one-generator} et
\ref{flat-example-two-generators}.
\end{lemma}

\begin{proof}
Cela résulte de l'application du
lemme \ref{sites-cohomology-lemma-descent-squares} de Cohomologie sur les sites,
dont les hypothèses sont satisfaites par le
lemme \ref{etale-cohomology-lemma-blow-up-square-h} et le
lemme \ref{flat-lemma-refine-check-h} de Compléments sur la platitude.
\end{proof}






\section{Cohomologie du faisceau structural pour la topologie h}
\label{etale-cohomology-section-cohomology-O-h}

\noindent
Soit $p$ un nombre premier. Soit $(\mathcal{C}, \mathcal{O})$ un site annelé
tel que $p\mathcal{O} = 0$. Nous notons alors $\colim_F \mathcal{O}$
la limite inductive, dans la catégorie des faisceaux d'anneaux, du système
$$
\mathcal{O} \xrightarrow{F} \mathcal{O} \xrightarrow{F}
\mathcal{O} \xrightarrow{F} \ldots
$$
où $F : \mathcal{O} \to \mathcal{O}$, $f \mapsto f^p$,
est l'endomorphisme de Frobenius.

\begin{lemma}
\label{etale-cohomology-lemma-h-sheaf-colim-F}
Soit $p$ un nombre premier. Soit $S$ un schéma sur $\mathbf{F}_p$.
Considérons le faisceau $\mathcal{O}^{perf} = \colim_F \mathcal{O}$
sur $(\Sch/S)_{fppf}$. Alors $\mathcal{O}^{perf}$ appartient à l'image
essentielle de $R\epsilon_* : D((\Sch/S)_h) \to D((\Sch/S)_{fppf})$.
\end{lemma}

\begin{proof}
Nous employons le critère du lemme \ref{etale-cohomology-lemma-refine-check-h}.
Avant d'en vérifier les conditions, remarquons que, pour un objet
quasi-compact et quasi-séparé $X$ de
$(\Sch/S)_{fppf}$, nous avons
$$
H^i_{fppf}(X, \mathcal{O}^{perf}) = \colim_F H^i_{fppf}(X, \mathcal{O})
$$
Voir le lemme de Cohomologie sur les sites
\ref{sites-cohomology-lemma-colim-works-over-collection}.
Nous emploierons également $H^i_{fppf}(X, \mathcal{O}) = H^i(X, \mathcal{O})$; voir la
proposition \ref{descent-proposition-same-cohomology-quasi-coherent} de Descente.

\medskip\noindent
Soient $A, f, J$ comme dans
l'exemple \ref{flat-example-one-generator} de Compléments sur la platitude,
et considérons le carré presque d'éclatement associé.
Puisque $X$, $X'$, $Z$, $E$ sont affines, la cohomologie
supérieure de $\mathcal{O}$ est nulle. Il suffit donc
de vérifier que
$$
0 \to
\mathcal{O}^{perf}(X) \to
\mathcal{O}^{perf}(X') \oplus \mathcal{O}^{perf}(Z) \to
\mathcal{O}^{perf}(E) \to 0
$$
est une suite exacte courte. Cela a été démontré dans (la démonstration du)
lemme \ref{flat-lemma-h-sheaf-colim-F} de Compléments sur la platitude.

\medskip\noindent
Soient $X, X', Z, E$ comme dans
l'exemple \ref{flat-example-two-generators} de Compléments sur la platitude.
Comme $X$ et $Z$ sont affines, nous avons
$H^i(X, \mathcal{O}_X) = H^i(Z, \mathcal{O}_Z) = 0$ pour $i > 0$.
Par le lemme \ref{flat-lemma-funny-blow-up} de Compléments sur la platitude,
nous avons $H^i(X', \mathcal{O}_{X'}) = 0$ pour $i > 0$.
Comme $E = \mathbf{P}^1_Z$ et que $Z$ est affine, nous avons aussi
$H^i(E, \mathcal{O}_E) = 0$ pour $i > 0$. Comme au paragraphe
précédent, nous nous ramenons à vérifier que
$$
0 \to
\mathcal{O}^{perf}(X) \to
\mathcal{O}^{perf}(X') \oplus \mathcal{O}^{perf}(Z) \to
\mathcal{O}^{perf}(E) \to 0
$$
est une suite exacte courte. Cela a été démontré dans (la démonstration du)
lemme \ref{flat-lemma-h-sheaf-colim-F} de Compléments sur la platitude.
\end{proof}

\begin{proposition}
\label{etale-cohomology-proposition-h-cohomology-structure-sheaf}
Soit $p$ un nombre premier. Soit $S$ un schéma quasi-compact et quasi-séparé
sur $\mathbf{F}_p$. Alors
$$
H^i((\Sch/S)_h, \mathcal{O}^h) =
\colim_F H^i(S, \mathcal{O})
$$
Dans le membre de gauche, $\mathcal{O}^h$ désigne
le faisceau h associé au faisceau structural.
\end{proposition}

\begin{proof}
C'est une reformulation du lemme \ref{etale-cohomology-lemma-h-sheaf-colim-F}.
Rappelons que
$\mathcal{O}^h = \mathcal{O}^{perf} = \colim_F \mathcal{O}$; voir le
lemme \ref{flat-lemma-char-p} de Compléments sur la platitude.
Le lemme \ref{etale-cohomology-lemma-h-sheaf-colim-F} montre que
$\mathcal{O}^{perf}$, vu comme objet de $D((\Sch/S)_{fppf})$,
est de la forme $R\epsilon_*K$ pour un certain $K \in D((\Sch/S)_h)$.
Alors $K = \epsilon^{-1}\mathcal{O}^{perf}$, qui est en fait
égal à $\mathcal{O}^{perf}$ puisque $\mathcal{O}^{perf}$ est un faisceau h. Voir la
section \ref{sites-cohomology-section-compare} de Cohomologie sur les sites.
Ainsi, $R\epsilon_*\mathcal{O}^{perf} = \mathcal{O}^{perf}$
(nous prions le lecteur d'excuser cette notation ambiguë).
Le résultat découle alors de la suite spectrale de Leray
$$
R\Gamma_h(S, \mathcal{O}^{perf}) =
R\Gamma_{fppf}(S, R\epsilon_*\mathcal{O}^{perf}) =
R\Gamma_{fppf}(S, \mathcal{O}^{perf})
$$
et du fait que
$$
H^i_{fppf}(S, \mathcal{O}^{perf}) =
H^i_{fppf}(S, \colim_F \mathcal{O}) =
\colim_F H^i_{fppf}(S, \mathcal{O})
$$
puisque $S$ est quasi-compact et quasi-séparé
(voir la démonstration du lemme \ref{etale-cohomology-lemma-h-sheaf-colim-F}).
\end{proof}















\bibliography{my}
\bibliographystyle{amsalpha}
\end{document}
